\chapter{Divisibility of periods}\label{chapter-period-divisibility}

In Stepanov's method a polynomial is shown to vanish to high order at many points.
In Alpöge's argument the order of vanishing is replaced by $p$-adic divisibility of the
periods of the iterated theta lifts of Chapter~\ref{chapter-unitary-theta}.
The divisibility comes from the $p$-integrality of Taylor jets of theta functions on CM
elliptic curves, and at a prime $p$ inert in the imaginary quadratic field $\mathbf Q(\sqrt{-d})$
of the host field it amounts to
about $a/p+a/p^{3}+\cdots$ for each of the $m$ lifts.
This chapter states the integrality of periods, deduces the multiplicity bound along the
chain of lifts, proves the integrality of periods for the definite-side data of the construction,
and states the lower bound for the gains at inert primes that enters the
inert-prime sum of Chapter~\ref{chapter-inert-prime-balance}.

\section{Theta functions on CM elliptic curves}\label{section-pd-cm-theta}

\begin{definition}\label{pd-cm-lattice}
Let $d\geq3$ be a squarefree integer with $d\equiv3\pmod4$, let
$K=\mathbf Q(\sqrt{-d})$ and $\omega_K=(1+i\sqrt d)/2$, so that
$\mathcal O_K=\mathbf Z[\omega_K]$. A lattice $L=\mathbf Z\omega_1+\mathbf Z\omega_2\subset\mathbf C$,
with $\omega_1,\omega_2$ linearly independent over $\mathbf R$, has \emph{complex
multiplication by $\mathcal O_K$} if $\omega_KL\subset L$. For a lattice $L$ we write
$g_2(L)=60\sum_{\omega\neq0}\omega^{-4}$ and $g_3(L)=140\sum_{\omega\neq0}\omega^{-6}$, the sums
over $\omega\in L\smallsetminus\{0\}$, and
\[
 j(L)=\frac{1728\,g_2(L)^3}{g_2(L)^3-27g_3(L)^2}.
\]
\end{definition}
\leanref{Formalism/LocalToGlobal/TorsionIntegrality/CMGoodModels/Statement.lean}{Formalism.L2.Th1.CM.HasCMBy}
\leanref{Formalism/LocalToGlobal/TorsionIntegrality/CMGoodModels/Statement.lean}{Formalism.L2.Th1.CM.jInv}

\begin{theorem}\label{pd-j-invariant-integral}
Let $d$ be a squarefree integer with $d\equiv3\pmod4$ and let $L\subset\mathbf C$ be a lattice
with complex multiplication by $\mathcal O_K$, $K=\mathbf Q(\sqrt{-d})$
(Definition~\ref{pd-cm-lattice}). Then $j(L)$ is an algebraic integer
\cite[Chapter~II, \S6]{Sil94}.
\end{theorem}
\leanref{Formalism/LocalToGlobal/TorsionIntegrality/CMGoodModels/Statement.lean}{Formalism.L2.Th1.CM.JInvariantIntegrality}
\leanref{Formalism/LocalToGlobal/TorsionIntegrality/CMGoodModels/CmAll.lean}{Formalism.L2.Th1.CM.cm1}

\begin{definition}\label{pd-good-model}
Let $L\subset\mathbf C$ be a lattice, let $l$ be a prime, let $\mathbf C_l$ be the completion of
an algebraic closure of $\mathbf Q_l$, let $\iota:\mathbf C\to\mathbf C_l$ be an isomorphism of
fields, and let $\widehat{\mathbf Q}_l^{\mathrm{ur}}\subset\mathbf C_l$ be the closure of the
maximal unramified extension of $\mathbf Q_l$. A \emph{good model} of $L$ at $(l,\iota)$ is a pair
$(g_2',g_3')$ of elements of $\widehat{\mathbf Q}_l^{\mathrm{ur}}$ with $|g_2'|\leq1$,
$|g_3'|\leq1$ and $|g_2'^3-27g_3'^2|=1$, such that $g_2'=\lambda^4\iota(g_2(L))$ and
$g_3'=\lambda^6\iota(g_3(L))$ for some $\lambda\in\mathbf C_l^{\times}$.
\end{definition}
\leanref{Formalism/LocalToGlobal/TorsionIntegrality/CMGoodModels/Statement.lean}{Formalism.L2.Th1.CM.IsGoodModel}

\begin{proposition}\label{pd-inert-formal-group}
Let $d$ be a squarefree integer with $d\equiv3\pmod4$ and let $L$ be a lattice with complex
multiplication by $\mathcal O_K$, $K=\mathbf Q(\sqrt{-d})$. Let $l\geq5$ be a prime with
$l\nmid d$ such that $-d$ is not a square modulo $l$, let $\iota:\mathbf C\to\mathbf C_l$ be an
isomorphism of fields, and let $(g_2',g_3')$ be a good model of $L$ at $(l,\iota)$
(Definition~\ref{pd-good-model}). Define $c_k\in\widehat{\mathbf Q}_l^{\mathrm{ur}}$ for $k\geq2$
by $c_2=g_2'/20$, $c_3=g_3'/28$ and
\[
 c_k=\frac{3}{(2k+1)(k-3)}\sum_{i=2}^{k-2}c_ic_{k-i}\qquad(k\geq4),
\]
so that $\wp(z)=z^{-2}+\sum_{k\geq2}c_kz^{2k-2}$ is the Weierstrass function of the curve with
invariants $g_2',g_3'$, and let
\[
 u(z)=z\Big(1+\sum_{k\geq2}c_kz^{2k}\Big)\Big(1-\sum_{k\geq2}(k-1)c_kz^{2k}\Big)^{-1}
 =-\frac{2\wp(z)}{\wp'(z)}
\]
be the formal parameter $-2x/y$ written in the coordinate of the invariant differential.
Let $\ell(X)=\sum_{j\geq0}X^{l^{2j}}/l^{j}$. Then there is
$\varepsilon\in\widehat{\mathbf Q}_l^{\mathrm{ur}}$ with $|\varepsilon|=1$ such that every
coefficient of the power series $u(\varepsilon\,\ell(X))$ lies in
$\widehat{\mathbf Q}_l^{\mathrm{ur}}$ and has absolute value at most $1$.
\end{proposition}
\leanref{Formalism/LocalToGlobal/TorsionIntegrality/CMGoodModels/Statement.lean}{Formalism.L2.Th1.CM.CM3}
\leanref{Formalism/LocalToGlobal/TorsionIntegrality/CMGoodModels/CmAll.lean}{Formalism.L2.Th1.CM.cm3}

The reduction of the curve is supersingular because $l$ is inert in $K$ \cite{Deu41}. The formal
group of the model is then a formal $\mathcal O_{K}\otimes\mathbf Z_l$-module of height one, hence
a Lubin-Tate group for the uniformizer $l$ of the unramified quadratic extension of
$\mathbf Q_l$, and it is isomorphic over the integers of $\widehat{\mathbf Q}_l^{\mathrm{ur}}$
to the Lubin-Tate group with logarithm $\ell$ \cite{LT65}, \cite{Haz78}; the
unit $\varepsilon$ is the derivative of that isomorphism at the origin.

\begin{definition}\label{pd-theta-function}
Let $n$ be a finite set, let $h=(h_{ij})$ be an $n\times n$ complex matrix, and let
$H(z,w)=\sum_{i,j}z_ih_{ij}\overline{w_j}$ on $\mathbf C^n$, with $E(z,w)=\operatorname{Im}H(z,w)$.
\begin{enumerate}
\item A subgroup $\Lambda\subset\mathbf C^n$ is \emph{self-dual} for $H$ if
  $\Lambda=\{x\in\mathbf C^n:E(x,y)\in\mathbf Z\text{ for all }y\in\Lambda\}$.
\item For $x\in\mathbf C^n$ and a function $F:\mathbf C^n\to\mathbf C$, let
  $(U_xF)(z)=F(z+x)\exp\big(-\pi H(z,x)-\pi H(x,x)/2\big)$.
\item For $c\in\mathbf R$ let $\alpha_c(x)=\exp\big(\pi i\,c\operatorname{Re}H(x,x)\big)$. A
  \emph{theta function} for $(H,\Lambda)$ with normalization $c$ is an entire function
  $F:\mathbf C^n\to\mathbf C$ with $\alpha_c(x)\,(U_xF)(z)=F(z)$ for all $x\in\Lambda$ and
  $z\in\mathbf C^n$.
\end{enumerate}
For a lattice $L\subset\mathbf C$ with basis $\omega_1,\omega_2$ we use the same notation for
$n=\{1\}$ and $H_L(z,w)=z\overline w/A(L)$, where $A(L)=|\operatorname{Im}(\overline{\omega_1}\omega_2)|$
is the covolume of $L$; we write $U^L_x$ for the corresponding operators. Let $\sigma_L$ and
$\zeta_L$ be the Weierstrass functions of $L$, $\eta_L:L\to\mathbf C$ the quasi-periods
$\eta_L(\omega_k)=2\zeta_L(\omega_k/2)$, extended $\mathbf R$-linearly, and
\[
 s_2(L)=\frac{\eta_L(\omega_1)\overline{\omega_2}-\eta_L(\omega_2)\overline{\omega_1}}
 {\omega_1\overline{\omega_2}-\omega_2\overline{\omega_1}},\qquad
 \sigma_L^{*}(z)=\sigma_L(z)\exp\big(-s_2(L)z^2/2\big).
\]
For $t\in\mathbf C$ with $(U^L_t\sigma_L^*)(0)\neq0$, the \emph{normalized theta function of $L$
translated by $t$} is $\vartheta_{L,t}(z)=(U^L_t\sigma^*_L)(z)/(U^L_t\sigma^*_L)(0)$.
\end{definition}
\leanref{Formalism/LocalToGlobal/TorsionIntegrality/QuadraticSemicharacters/IsSemiChar.lean}{Formalism.L2.Th1.H.IsSelfDual}
\leanref{Formalism/LocalToGlobal/TorsionIntegrality/QuadraticSemicharacters/IsSemiChar.lean}{Formalism.L2.Th1.H.opU}
\leanref{Formalism/LocalToGlobal/TorsionIntegrality/QuadraticSemicharacters/IsSemiChar.lean}{Formalism.L2.Th1.H.Th}
\leanref{Formalism/LocalToGlobal/TorsionIntegrality/QuadraticSemicharacters/ProductThetaMem.lean}{Formalism.L2.Th1.H.alphaQ}
\leanref{Formalism/LocalToGlobal/TorsionIntegrality/EisensteinTranslationProducts/SemiChar.lean}{Formalism.L2.Th1.F.th0}

\begin{definition}\label{pd-cm-datum}
A \emph{CM datum} is $(a_0,b_0,c_0,d,\lambda)$ with $a_0\geq1$ odd, $b_0,c_0\in\mathbf Z$, $d$ a
prime with $d\equiv7\pmod8$, $4a_0c_0=b_0^2+d$, and $\lambda\in\mathbf C^{\times}$. Its
\emph{lattice} is $L=\lambda(\mathbf Za_0+\mathbf Z\theta_K)$, where
$\theta_K=(b_0+i\sqrt d)/2$; it has complex multiplication by $\mathcal O_K$,
$K=\mathbf Q(\sqrt{-d})$.

Let $(a_0,b_0,c_0,d,\lambda)$ be a CM datum with lattice $L$, let $n$, $h$, $H$ be as in
Definition~\ref{pd-theta-function}, let $\Lambda\subset\mathbf C^n$ be a subgroup and $\theta$ a
function on $\mathbf C^n$. Let $p$ be a prime and $\iota:\mathbf C\to\mathbf C_p$ an isomorphism of
fields. On $\mathbf C^{4\times n}=(\mathbf C^n)^4$ let $H^{\oplus4}$ be the orthogonal sum of four
copies of $H$, let $H_{\mathrm{std}}(z,w)=A(L)^{-1}\sum_{a,i}z_{a,i}\overline{w_{a,i}}$, and let
$L^{4\times n}$ be the lattice of vectors with all coordinates in $L$. Then $\theta$ is
\emph{locally normalized} at $(p,\iota)$ if there are a $\mathbf C$-linear automorphism $g$ of
$\mathbf C^{4\times n}$ and $\varepsilon\in\mathbf C$ such that:
\begin{enumerate}
\item $H_{\mathrm{std}}(gz,gw)=H^{\oplus4}(z,w)$ for all $z,w$;
\item with $\Lambda'=g(\Lambda^4)$, the index $[\Lambda':\Lambda'\cap L^{4\times n}]$ is finite and
  prime to $p$, and the index $[L^{4\times n}:\Lambda'\cap L^{4\times n}]$ is finite;
\item $|\iota(\varepsilon)|\leq1$;
\item for all $z=(z_1,z_2,z_3,z_4)\in(\mathbf C^n)^4$,
  \[
   \theta(z_1)\theta(z_2)\theta(z_3)\theta(z_4)=\varepsilon\,
   [\Lambda':\Lambda'\cap L^{4\times n}]^{-1/2}\sum_{c}\beta(c)\,(U_cF_0)(gz),
  \]
  where $c$ runs over a set of representatives of $\Lambda'/(\Lambda'\cap L^{4\times n})$, the
  operators $U_c$ are those of $H_{\mathrm{std}}$, $\beta=\alpha_{\sqrt d/2}$ is formed with
  $H_{\mathrm{std}}$, and $F_0(y)=\prod_{a,i}\vartheta_{L,\lambda a_0/2}(y_{a,i})$.
\end{enumerate}
\end{definition}
\leanref{Formalism/LocalToGlobal/TorsionIntegrality/QuadraticSemicharacters/QuadraticSemicharactersFactAltForm.lean}{Formalism.L2.Th1.H.CmOK}
\leanref{Formalism/LocalToGlobal/TorsionIntegrality/QuadraticSemicharacters/QuadraticSemicharactersFactAltForm.lean}{Formalism.L2.Th1.H.cmL}
\leanref{Formalism/LocalToGlobal/TorsionIntegrality/GaussianEigenFrames/StatementC.lean}{Formalism.L2.Th1.GD.LocNorm}
\leanref{Formalism/LocalToGlobal/TorsionIntegrality/QuadraticSemicharacters/ProductThetaMem.lean}{Formalism.L2.Th1.H.kernel4}

\begin{definition}\label{pd-integral-jets}
Let $K$ be a field complete for a non-archimedean absolute value and let $\nu$ be a finite set.
A power series $G=\sum_\kappa g_\kappa X^\kappa\in K[[X_x:x\in\nu]]$ has \emph{integral jets} if
$|\kappa!\,g_\kappa|\leq1$ for every multi-index $\kappa$, where
$\kappa!=\prod_x\kappa_x!$. For a polynomial $P=\sum_\kappa p_\kappa X^\kappa\in K[X_x:x\in\nu]$
the \emph{jet pairing} is $\langle P,G\rangle=\sum_\kappa p_\kappa\,\kappa!\,g_\kappa$.
\end{definition}
\leanref{Formalism/GaussJetBounds/StatementN.lean}{Formalism.DworkPowerSeriesWeights.JetsIntegral}
\leanref{Formalism/GaussJetBounds/StatementI.lean}{Formalism.DworkPowerSeriesWeights.jetPair}

When $G$ is the Taylor series of a function at $0$, the coefficient $\kappa!\,g_\kappa$ is the
partial derivative $\partial^\kappa$ at $0$, and $\langle P,G\rangle$ is the value at $0$ of the
differential operator $P(\partial)$ applied to the function.

\begin{lemma}\label{pd-jets-closure}
Let $K$ be a field of characteristic $0$, complete for a non-archimedean absolute value, with
valuation ring $\mathcal O$. Call a power series over $K$ \emph{integral} if its coefficients lie in
$\mathcal O$.
\begin{enumerate}
\item Sums and products of power series with integral jets have integral jets, and so does the
  product of one with a constant in $\mathcal O$.
\item If $F$ has integral jets and constant term $0$, then $\exp(F)$ has integral jets.
\item If $F$ has integral jets in variables $Y_w$, and $C=(C_{w,x})$ is a matrix with entries in
  $\mathcal O$, then $F(CX)$, obtained by substituting $Y_w=\sum_xC_{w,x}X_x$, has integral jets.
\item If $G$ is integral and $u_w$ are power series with integral jets and constant term $0$, then
  $G(u)$, obtained by substituting $u_w$ for the $w$-th variable, has integral jets.
\item Let $F$ be a power series in variables $X_x$, $x\in\nu$, and let $F^{\otimes4}$ be the series
  $F(X^{(1)})F(X^{(2)})F(X^{(3)})F(X^{(4)})$ in four disjoint copies of the variables. If
  $F^{\otimes4}$ has integral jets, so does $F$.
\end{enumerate}
\end{lemma}
\leanref{Formalism/LocalToGlobal/TorsionIntegrality/GaussianEigenFrames/GramFrameComposition/JetsLin.lean}{Formalism.L2.Th1.GD.JL.fourth_root_mem}
\leanref{Development/AbelianVarietyTheory/Definitions/GaloisGeometry/FourthPowerTheta.lean}{GlobalTheory.ThLoc.LocJets.jets_comb}

\begin{proof}
A series has integral jets exactly when it is $\sum_\kappa a_\kappa X^\kappa/\kappa!$ with all
$a_\kappa\in\mathcal O$. (1) The sum is clear. For products, $X^\kappa/\kappa!\cdot
X^\lambda/\lambda!=\binom{\kappa+\lambda}{\kappa}X^{\kappa+\lambda}/(\kappa+\lambda)!$, and the
binomial coefficient is an integer. (2) Write $F=\sum_{\kappa\neq0}a_\kappa X^\kappa/\kappa!$. By
(1) it suffices that $F^n/n!$ has integral jets for every $n$. Expanding $F^n$ by the multinomial
theorem, $F^n/n!$ is a sum of products $\prod_\kappa (a_\kappa X^\kappa/\kappa!)^{n_\kappa}/n_\kappa!$
with $\sum n_\kappa=n$, and $(X^\kappa/\kappa!)^{r}/r!=\frac{(r\kappa)!}{r!\,(\kappa!)^r}\,
X^{r\kappa}/(r\kappa)!$, where $\frac{(r\kappa)!}{r!(\kappa!)^r}$ is an integer for
$\kappa\neq0$ (it counts the partitions of a set of size $r|\kappa|$, coloured by $\kappa$, into $r$
blocks of type $\kappa$). (3) Each $\big(\sum_xC_{w,x}X_x\big)^{k}/k!$ has integral jets by the
multinomial theorem, so (3) follows from (1). (4) Each monomial $u^\kappa$ has integral jets by (1),
and the sum converges coefficientwise because each $u_w$ has constant term $0$. (5) Write
$F=\sum_\kappa a_\kappa X^\kappa/\kappa!$. The coefficient of
$X^{(1)\kappa}X^{(2)\kappa}X^{(3)\kappa}X^{(4)\kappa}/(\kappa!)^4$ in $F^{\otimes4}$ is $a_\kappa^4$,
because the four copies of the variables are disjoint. So $|a_\kappa|^4\leq1$, and $|a_\kappa|\leq1$.
\end{proof}

\begin{definition}\label{pd-normalizing-series}
Let $L\subset\mathbf C$ be a lattice, $p$ a prime and $\iota:\mathbf C\to\mathbf C_p$ an isomorphism
of fields such that the Weierstrass curve
\[
 \mathcal E_L:\ y^2=x^3-\tfrac14\iota(g_2(L))\,x-\tfrac14\iota(g_3(L))
\]
has coefficients of absolute value at most $1$ and discriminant of absolute value $1$. Let
$X=-x/y$ be the formal parameter at the origin of $\mathcal E_L$, let
$\omega=(1+\cdots)\,dX$ be the invariant differential, $\log_{\mathcal E}(X)=\int_0^X\omega$
the formal logarithm and $\exp_{\mathcal E}$ its inverse, and let $F_{\mathcal E}(X,Y)=
\exp_{\mathcal E}(\log_{\mathcal E}X+\log_{\mathcal E}Y)$ be the formal group law. For
$a\in\mathbf C_p$ put $[a](X)=\exp_{\mathcal E}(a\log_{\mathcal E}X)$. Let $\partial$ be the
invariant derivation $\omega(X)^{-1}\frac{d}{dX}$, the derivative with respect to
$z=\log_{\mathcal E}(X)$.
\begin{enumerate}
\item The \emph{normalizing series} of $L$ at $(p,\iota)$ is the series $\phi_L\in\mathbf C_p[[X]]$
  with constant term $0$ such that
  \[
   X\,e^{-\phi_L(X)}=\iota(\sigma^*_L)\big(\log_{\mathcal E}X\big),
  \]
  where $\iota(\sigma^*_L)$ is the Taylor series of the function $\sigma^*_L$ of
  Definition~\ref{pd-theta-function} at $0$ with $\iota$ applied to its coefficients. Explicitly,
  if $\iota(\wp_L)(z)=z^{-2}+\sum_{n\geq0}w_nz^n$, then
  \[
   \phi_L(X)=\sum_{n\geq2}\frac{w_{n-2}}{n(n-1)}\,z^n+\frac{\iota(s_2(L))}{2}\,z^2-
   \log\frac{z}{X},\qquad z=\log_{\mathcal E}(X).
  \]
  Its \emph{jets} are the numbers $(\partial^i\phi_L)(0)$, $i\geq1$.
\item For power series $f,g$ over $\mathbf C_p$ in a finite set of variables, $f\equiv g$ means
  that $f-g$ has constant term $0$ and $\exp(f-g)$ is integral. A \emph{point} is an integral
  series with constant term $0$.
\item For points $x,y$ put $\beta_L(x,y)=\phi_L(F_{\mathcal E}(x,y))-\phi_L(x)-\phi_L(y)$.
\end{enumerate}
\end{definition}
\leanref{Formalism/LocalToGlobal/TorsionIntegrality/UniformGaussModels/GoodModelAll.lean}{Formalism.L2.Th1.Top.phC}
\leanref{Formalism/LocalToGlobal/TorsionIntegrality/SecondCoefficientNorms/HasSecondLogDeriv.lean}{Formalism.L2.Th1.C.IsPh}
\leanref{Formalism/LocalToGlobal/TorsionIntegrality/SecondCoefficientNorms/JetExpShape.lean}{Formalism.L2.Th1.C.jet}
\leanref{Formalism/LocalToGlobal/TorsionIntegrality/SecondCoefficientNorms/HasSecondLogDeriv.lean}{Formalism.L2.Th1.C.Cong}

The two descriptions of $\phi_L$ agree: $(\log\sigma_L)''=-\wp_L$ and $\sigma_L(z)=z+O(z^5)$ give
$\log(\sigma_L(z)/z)=-\sum_{n\geq2}w_{n-2}z^n/(n(n-1))$, and
$\sigma^*_L(z)=\sigma_L(z)e^{-s_2(L)z^2/2}$. Since the reduction of $\mathcal E_L$ is
an elliptic curve, the formal group law $F_{\mathcal E}$ is integral, and so is $\omega(X)$, with constant
term $1$ \cite[Chapter~IV]{Sil09}. Hence $\omega^{-1}$ is integral, and $\partial$ maps integral
series to integral series.

\begin{lemma}\label{pd-cm-isogeny}
Let $d$ be a squarefree integer with $d\equiv3\pmod4$, let $L$ be a lattice with complex
multiplication by $\mathcal O_K$, $K=\mathbf Q(\sqrt{-d})$ (Definition~\ref{pd-cm-lattice}), and
let $\alpha=m+n\omega_K$ with $m,n\in\mathbf Z$ and $N=\alpha\bar\alpha\neq0$. Let $S\subset\mathbf C$
be a set of representatives of the nonzero classes of $\alpha^{-1}L/L$. Then:
\begin{enumerate}
\item $|S|+1=N$; $NP\in L$ and $P\notin L$ for $P\in S$; and there is an involution $P\mapsto P'$
  of $S$ with $P+P'\in L$, so that $\wp_L(P')=\wp_L(P)$ and $\wp_L'(P')=-\wp_L'(P)$;
\item as Laurent series at $z=0$,
  $\alpha^2\wp_L(\alpha z)=\wp_L(z)+\sum_{P\in S}\big(\wp_L(z+P)-\wp_L(P)\big)$;
\item $\alpha(\alpha-\bar\alpha)\,s_2(L)=\sum_{P\in S}\wp_L(P)$;
\item for $P\in S$, the Taylor series $f_P$ at $z=0$ of $z^2\big(\wp_L(z)-\wp_L(P)\big)$, which has
  constant term $1$, satisfies
  \[
   f_P''f_P-(f_P')^2=\big(2\wp^0_L(z)-\wp_L(z+P)-\wp_L(z+P')\big)f_P^2,
  \]
  where $\wp^0_L(z)=\wp_L(z)-z^{-2}$ and all terms are Taylor series at $z=0$.
\end{enumerate}
\end{lemma}
\leanref{Formalism/LocalToGlobal/TorsionIntegrality/GaussSumParity/KummerThetaLogarithms/SecondLogDerivativeUnique.lean}{Formalism.L2.Th1.D.CurveNormEndomorphisms.cdata}
\leanref{Formalism/LocalToGlobal/TorsionIntegrality/SecondCoefficientNorms/HasSecondLogDeriv.lean}{Formalism.L2.Th1.C.IsEndoZ}

\begin{proof}
(1) Since $\omega_KL\subset L$, multiplication by $\alpha$ and by $\bar\alpha=(m+n)-n\omega_K$ maps
$L$ into itself. Let $\gamma\in M_2(\mathbf Z)$ be the matrix of multiplication by $\alpha$ in the
basis $\omega_1,\omega_2$. Then $\alpha$ is a root of the characteristic polynomial of $\gamma$, which
has integer coefficients. If $\alpha\notin\mathbf R$ its roots are $\alpha$ and $\bar\alpha$, so
$\det\gamma=\alpha\bar\alpha$; if $\alpha\in\mathbf R$ then $\gamma=\alpha\cdot1$, because
$\omega_1,\omega_2$ are linearly independent over $\mathbf R$, and $\det\gamma=\alpha^2$. In both
cases $|L/\alpha L|=|\det\gamma|=N$. Multiplication by $\alpha$ induces a bijection from
$\alpha^{-1}L/L$ onto $L/\alpha L$, so $\alpha^{-1}L/L$ has $N$ elements and $|S|+1=N$. For $P\in S$,
$P\notin L$ because its class is nonzero, and $NP=\bar\alpha(\alpha P)\in L$. Also
$-P\in\alpha^{-1}L$ and $-P\notin L$, so $-P$ is congruent modulo $L$ to exactly one $P'\in S$. Then
$P+P'\in L$, and $P''=P$ because both are congruent to $-P'$. Since $\wp_L$ is even, $\wp_L'$ is odd
and both are periodic for $L$, $\wp_L(P')=\wp_L(P)$ and $\wp_L'(P')=-\wp_L'(P)$.

(2) Put $T=S\cup\{0\}$ and $g(z)=\alpha^2\wp_L(\alpha z)-\sum_{P\in T}\wp_L(z+P)$. Since
$\alpha L\subset L$, $g$ is periodic for $L$. If $\alpha z_0\notin L$, then $z_0+P\notin L$ for
$P\in T$, because $\alpha P\in L$, so $g$ is holomorphic at $z_0$. If $\alpha z_0\in L$, then, as in
(1), exactly one $P_0\in T$ has $z_0+P_0\in L$. Near $z_0$, both $\alpha^2\wp_L(\alpha z)$ and
$\wp_L(z+P_0)$ are $(z-z_0)^{-2}$ plus a function holomorphic at $z_0$, and the other terms are
holomorphic at $z_0$. So $g$ extends to an elliptic function without poles, and is constant. Near
$0$, with $P_0=0$, $g(z)=\alpha^2\wp^0_L(\alpha z)-\wp^0_L(z)-\sum_{P\in S}\wp_L(z+P)$, and
$\wp^0_L(0)=0$; so the constant is $-\sum_{P\in S}\wp_L(P)$. Hence
\[
 \alpha^2\wp_L(\alpha z)=\wp_L(z)+\sum_{P\in S}\big(\wp_L(z+P)-\wp_L(P)\big)\qquad(\alpha z\notin L),
 \tag{$*$}
\]
and in particular for small $z\neq0$, which gives (2).

(3) The set $U=\{z\in\mathbf C:\alpha z\notin L\}$ is the complement of the countable set
$\alpha^{-1}L$, so it is open and connected. On $U$ put
\[
 \psi(z)=\alpha\zeta_L(\alpha z)-\zeta_L(z)-\sum_{P\in S}\zeta_L(z+P)-z\sum_{P\in S}\wp_L(P).
\]
Since $\zeta_L'=-\wp_L$, $(*)$ gives $\psi'=0$ on $U$, so $\psi$ is constant on $U$. Let
$\omega\in L$ and $x=\omega_1/(2\alpha)$. Then $\alpha x=\omega_1/2\notin L$, and $x+\omega\in U$
because $\alpha\omega\in L$. Since $\zeta_L(z+\omega')=\zeta_L(z)+\eta_L(\omega')$ for
$\omega'\in L$ and $z\notin L$, the equality $\psi(x+\omega)=\psi(x)$ and $|S|+1=N$ give
\[
 \alpha\eta_L(\alpha\omega)-N\eta_L(\omega)=\omega\sum_{P\in S}\wp_L(P).
\]
Since $\eta_L$ is $\mathbf R$-linear, $\eta_L(\omega)=s_2(L)\omega+c\,\bar\omega$ for all
$\omega\in\mathbf C$, with $s_2(L)$ as in Definition~\ref{pd-theta-function} and a constant $c$ (the
Legendre relation gives $c=\pi/A(L)$). Take $\omega=\omega_1$. The terms in $c$ give
$c(\alpha\bar\alpha-N)\overline{\omega_1}=0$, and the others give
$(\alpha^2-N)s_2(L)\omega_1=\omega_1\sum_P\wp_L(P)$. Since $\omega_1\neq0$ and
$\alpha^2-N=\alpha(\alpha-\bar\alpha)$, this is (3).

(4) Let $P\in S$. For $z\neq0$ near $0$ the addition formula
$\sigma_L(z+P)\sigma_L(z-P)=-\sigma_L(z)^2\sigma_L(P)^2\big(\wp_L(z)-\wp_L(P)\big)$
\cite[\S20.53]{WW27} holds, with all factors nonzero. Its logarithmic derivative in $z$, with
$\zeta_L=\sigma_L'/\sigma_L$, is
$\wp_L'(z)=\big(\wp_L(z)-\wp_L(P)\big)\big(\zeta_L(z+P)+\zeta_L(z-P)-2\zeta_L(z)\big)$. So the
function $f_P(z)=z^2\big(\wp_L(z)-\wp_L(P)\big)$ has $f_P'/f_P=2/z+\zeta_L(z+P)+\zeta_L(z-P)-2\zeta_L(z)$,
and differentiating once more, with $\zeta_L'=-\wp_L$,
\[
 \frac{f_P''f_P-(f_P')^2}{f_P^2}=-\frac2{z^2}-\wp_L(z+P)-\wp_L(z-P)+2\wp_L(z)
 =2\wp^0_L(z)-\wp_L(z+P)-\wp_L(z+P'),
\]
since $\wp_L(z-P)=\wp_L(z+P')$ by (1). Both sides are holomorphic at $0$, with $f_P(0)=1$, so the
identity holds for the Taylor series.
\end{proof}

\begin{remark}\label{pd-cm-isogeny-classical}
Classically, (1) is the index $[L:\alpha L]=|\alpha|^2$, the factor by which multiplication by
$\alpha$ scales areas, and (3) comes from integrating (2) once:
$\alpha\zeta_L(\alpha z)=\zeta_L(z)+\sum_{P\in S}\big(\zeta_L(z+P)-\zeta_L(P)\big)+z\sum_{P\in S}\wp_L(P)$,
where the constant of integration is $0$ because both sides have constant term $0$ in their Laurent
expansions at $z=0$. Replacing $z$ by $z+\omega$ with $\omega\in L$ gives
$\alpha\eta_L(\alpha\omega)=N\eta_L(\omega)+\omega\sum_{P\in S}\wp_L(P)$ directly. The proof of
Lemma~\ref{pd-cm-isogeny} does not determine this constant; it uses only that $\psi$ is constant on
the connected set $U$.
\end{remark}

\begin{lemma}\label{pd-sigma-congruences}
Let $L$, $p\neq2$ and $\iota$ be as in Definition~\ref{pd-normalizing-series}, with normalizing
series $\phi_L$ and form $\beta_L$.
\begin{enumerate}
\item $\phi_L([-1]X)\equiv\phi_L(X)$.
\item For points $x,x',y$, $\beta_L(F_{\mathcal E}(x,x'),y)\equiv\beta_L(x,y)+\beta_L(x',y)$.
\item Suppose that $d$ is a squarefree integer with $d\equiv3\pmod4$, that $L$ has complex
  multiplication by $\mathcal O_K$, $K=\mathbf Q(\sqrt{-d})$, and that $\alpha=m+n\omega_K$
  ($m,n\in\mathbf Z$) has $N=\alpha\bar\alpha$ prime to $p$ and $|\iota(\alpha)|=1$. Write $\alpha$
  also for $\iota(\alpha)$ and let $S$ be as in Lemma~\ref{pd-cm-isogeny}. Then $[\alpha]$,
  $[\alpha^{-1}]$ and $[\bar\alpha]$ are integral and $\phi_L([\alpha]X)\equiv N\phi_L(X)$.
\end{enumerate}
\end{lemma}
\leanref{Formalism/LocalToGlobal/TorsionIntegrality/RamifiedCurveKernels/StatementR.lean}{Formalism.L2.Th1.G.WeierstrassBiadditiveModEquiv}
\leanref{Formalism/LocalToGlobal/TorsionIntegrality/SecondCoefficientNorms/FirstLevelFromUniversal.lean}{Formalism.L2.Th1.C.EndIntFromDatum}
\leanref{Formalism/LocalToGlobal/TorsionIntegrality/SecondCoefficientNorms/FirstLevelFromUniversal.lean}{Formalism.L2.Th1.C.PhNormFromDatum}
\leanref{Formalism/LocalToGlobal/TorsionIntegrality/SecondCoefficientNorms/HasSecondLogDeriv.lean}{Formalism.L2.Th1.C.IsogenyFormula}

\begin{proof}
The coordinates of $\mathcal E_L$ have integral expansions $x=X^{-2}(1+\cdots)$ and
$y=-X^{-3}(1+\cdots)$, and $X\,e^{-\phi_L(X)}=\iota(\sigma^*_L)(z)$ with $z=\log_{\mathcal E}X$.

(1) The function $\sigma^*_L$ is odd, and $[-1](X)$ is the parameter of $-z$. So
$[-1](X)\,e^{-\phi_L([-1]X)}=-X\,e^{-\phi_L(X)}$, and
$\exp\big(\phi_L([-1]X)-\phi_L(X)\big)=-[-1](X)/X$, which is integral because
$[-1](X)=-X+\cdots$ is integral.

(2) Write $x\oplus y=F_{\mathcal E}(x,y)$ and $x\ominus y=F_{\mathcal E}(x,[-1]y)$, and for points
$x,y$ put
\[
 \rho(x,y)=\phi_L(x\oplus y)+\phi_L(x\ominus y)-2\phi_L(x)-2\phi_L(y).
\]
Two rules are used throughout. If $G\equiv0$ in some variables and points are substituted for
them, the result is still $\equiv0$, since $\exp$ commutes with the substitution. If $f\equiv g$ and
$f'\equiv g'$, then $f\pm f'\equiv g\pm g'$, since the exponential of a sum is the product of the
exponentials and the inverse of an integral $1$-unit is integral.

\emph{Step 1: the parallelogram congruence $\rho(X_0,X_1)\equiv0$.} The function $\sigma_L$
satisfies $\sigma_L(u+v)\sigma_L(u-v)=\sigma_L(u)^2\sigma_L(v)^2\big(\wp_L(v)-\wp_L(u)\big)$
\cite[\S20.53]{WW27}, and so does $\sigma^*_L$, because $(u+v)^2+(u-v)^2=2u^2+2v^2$. Substitute
$u=\log_{\mathcal E}X_0$ and $v=\log_{\mathcal E}X_1$. Since $X_0\oplus X_1$ and $X_0\ominus X_1$
are the parameters of $u+v$ and $u-v$, and $Xe^{-\phi_L(X)}=\iota(\sigma^*_L)(\log_{\mathcal E}X)$,
this gives
\[
 (X_0\oplus X_1)(X_0\ominus X_1)\,e^{-\rho(X_0,X_1)}=X_0^2\,e(X_1)-X_1^2\,e(X_0),
\]
where $e(X)=X^2x(X)$ and $x(X)=\iota(\wp_L)(\log_{\mathcal E}X)$ is the $x$-coordinate of
$\mathcal E_L$ in the parameter $X$. The series $e$ is integral with constant term $1$, because the
coefficients of $\mathcal E_L$ are integral. For the Weierstrass formal group over any commutative
ring there is a power series $c$ over that ring, with constant term $-1$, such that
$X_1^2e(X_0)-X_0^2e(X_1)=c\,(X_0\oplus X_1)(X_0\ominus X_1)$: the difference $x(X_0)-x(X_1)$
vanishes on $X_0=X_1$ and on $X_0=[-1]X_1$. Over the integral coefficients of $\mathcal E_L$, $c$ is
integral. Hence $e^{-\rho(X_0,X_1)}=-c$, an integral series with constant term $1$, so
$e^{\rho(X_0,X_1)}=(-c)^{-1}$ is integral and $\rho(X_0,X_1)\equiv0$.

\emph{Step 2: halving.} The series $[2](X)=F_{\mathcal E}(X,X)=2X+\cdots$ is integral, and its
linear coefficient $2$ is a unit because $|2|=1$. So its compositional inverse
$[1/2](X)=\exp_{\mathcal E}\big(\tfrac12\log_{\mathcal E}X\big)$ is integral, and it carries points
to points.

\emph{Step 3: the cube.} In three variables put $B=[1/2](X_1\oplus X_2)$,
$C=[1/2](X_1\ominus X_2)$ and $A=X_0\oplus B$; by Step 2 these are points. In the coordinates
$z_j=\log_{\mathcal E}X_j$ they have logarithms $(z_1+z_2)/2$, $(z_1-z_2)/2$ and
$z_0+(z_1+z_2)/2$, so
\[
 A\oplus B=X_0\oplus X_1\oplus X_2,\quad A\ominus B=X_0,\quad A\oplus C=X_0\oplus X_1,\quad
 A\ominus C=X_0\oplus X_2,
\]
\[
 B\oplus C=X_1,\quad B\ominus C=X_2,\quad B\oplus B=X_1\oplus X_2,\quad B\ominus B=0.
\]
Substituting these into the definition of $\rho$, and using $\phi_L(0)=0$, the terms in
$\phi_L(A)$, $\phi_L(B)$ and $\phi_L(C)$ cancel, and
\[
 \rho(A,B)-\rho(A,C)+\rho(B,C)-\rho(B,B)=\Delta(X_0,X_1,X_2),
\]
where
\[
 \Delta=\phi_L(X_0\oplus X_1\oplus X_2)-\phi_L(X_0\oplus X_1)-\phi_L(X_0\oplus X_2)
 -\phi_L(X_1\oplus X_2)+\phi_L(X_0)+\phi_L(X_1)+\phi_L(X_2).
\]
By Step 1 and the substitution rule, each term on the left is $\equiv0$. By the additivity rule,
$\Delta\equiv0$.

\emph{Step 4.} Expanding the definition of $\beta_L$,
$\beta_L(X_0\oplus X_1,X_2)-\beta_L(X_0,X_2)-\beta_L(X_1,X_2)=\Delta(X_0,X_1,X_2)$. Substituting the
points $x,x',y$ for $X_0,X_1,X_2$ gives (2). Part (2) used only that the coefficients of
$\mathcal E_L$ are integral and that $|2|=1$.

(3) By Lemma~\ref{pd-cm-isogeny}(2) and its derivative, the coordinates of $[\alpha]X$ are
$\alpha^{-2}\big(x(X)+\sum_P(x(X\oplus P)-x(P))\big)$ and
$\alpha^{-3}\big(y(X)+\sum_Py(X\oplus P)\big)$, where $X\oplus P$ is the sum on $\mathcal E_L$. The
points $P\in S$ have order dividing $N$, prime to $p$, so their coordinates are integral and they
do not reduce to the origin; by the addition law the coordinates of $X\oplus P$ are integral series in
$X$. Since $|\alpha|=1$, the $y$-coordinate of $[\alpha]X$ is $X^{-3}$ times a unit of the ring of
integral series, so $[\alpha]X=-x/y$ is integral. The same holds for $[\alpha^{-1}]$, the inverse
of $[\alpha]$, and $[\bar\alpha]$ is the difference, in the formal group, of the integral
endomorphisms $[\alpha+\bar\alpha]$ and $[\alpha]$. Integrating Lemma~\ref{pd-cm-isogeny}(2)
twice, with Lemma~\ref{pd-cm-isogeny}(4) for the factors $z^2(\wp_L(z)-\wp_L(P))$, and using
Lemma~\ref{pd-cm-isogeny}(3) to remove the quadratic term gives
\[
 \sigma^*_L(\alpha z)^2z^{2|S|}=\alpha^2\sigma^*_L(z)^{2N}\prod_{P\in S}z^2\big(\wp_L(z)-\wp_L(P)\big).
\]
In the coordinate $X$, $z^2(\wp_L(z)-\wp_L(P))=(z/X)^2\cdot u_P(X)$ with $u_P=X^2(x(X)-x(P))$ an
integral $1$-unit, and $\sigma^*_L(\alpha z)=[\alpha](X)e^{-\phi_L([\alpha]X)}$. Hence
\[
 \exp\big(2\phi_L([\alpha]X)-2N\phi_L(X)\big)=\Big(\frac{[\alpha](X)}{\alpha X}\Big)^2\prod_{P\in S}u_P(X)^{-1}.
\]
The right side is the square of an integral $1$-unit, because $[\alpha](X)/(\alpha X)$ is one and
square roots of integral $1$-units are integral when $|2|=1$. So
$\exp\big(\phi_L([\alpha]X)-N\phi_L(X)\big)$ is integral.
\end{proof}

\begin{proposition}\label{pd-jets-abstract}
Let $K$ be a field of characteristic $0$, complete for a non-archimedean absolute value with
$|2|=1$ and valuation ring $\mathcal O$. Let $l\in K[[X]]$ with $l(0)=0$ and $l'(0)=1$, such that
$F(X,Y)=l^{-1}(l(X)+l(Y))$ is integral, and use the notation $[a]$, $\partial$, $\equiv$, points
and $\beta$ of Definition~\ref{pd-normalizing-series} for $l$ and a series $\phi\in K[[X]]$ with
$\phi(0)=0$. Suppose that:
\begin{enumerate}
\item $\phi([-1]X)\equiv\phi(X)$;
\item $\beta(F(x,x'),y)\equiv\beta(x,y)+\beta(x',y)$ for all points $x,x',y$;
\item there are $a,b\in K$ with $ab=N$ a positive integer, $[a]$, $[a^{-1}]$ and $[b]$ integral,
  $\phi([a]X)\equiv N\phi(X)$, $a\neq b$ and $|a-b|=1$.
\end{enumerate}
Then $(\partial^i\phi)(0)\in\mathcal O$ for every $i\geq1$.
\end{proposition}
\leanref{Formalism/LocalToGlobal/TorsionIntegrality/SecondCoefficientNorms/ThirdJetIntegral.lean}{Formalism.L2.Th1.C.JetsUnramified}
\leanref{Formalism/LocalToGlobal/TorsionIntegrality/SecondCoefficientNorms/BinomialSeriesCoefficients/IntegralThirdInvariantDerivative.lean}{Formalism.L2.Th1.C.jets_unramified}
\leanref{Formalism/LocalToGlobal/TorsionIntegrality/SecondCoefficientNorms/IsotropicVanish.lean}{Formalism.L2.Th1.C.AdjFromNorm}

\begin{proof}
Put $\phi^\circ(z)=\phi(l^{-1}(z))$, so that $(\partial^i\phi)(0)=(\phi^\circ)^{(i)}(0)$, and note that
$l'=(\partial_YF(X,0))^{-1}$ is integral with constant term $1$.

\emph{A remark on congruences.} If $G$ has constant term $0$ and $e^G$ is integral, then $e^G$ is a
$1$-unit with integral inverse $e^{-G}$, so every partial derivative
$\partial_{X_j}G=e^{-G}\partial_{X_j}e^G$ is integral. In particular the coefficients of the
monomials of degree one in $G$, and the coefficient of $X_jX_k$ for $j\neq k$, lie in
$\mathcal O$, the latter because it is the constant term of $\partial_{X_k}\partial_{X_j}G$.

\emph{Step 1: adjunction.} For points $x,y$, since $[a]$ is an endomorphism of $F$,
$\beta([a]x,[a]y)=\phi([a]F(x,y))-\phi([a]x)-\phi([a]y)\equiv N\beta(x,y)$ by (3). By (2) and the
symmetry of $\beta$, $N\beta(x,y)\equiv\beta(x,[N]y)=\beta(x,[b][a]y)$. Replacing $y$ by the point
$[a^{-1}]y$ gives $\beta([a]x,y)\equiv\beta(x,[b]y)$.

\emph{Step 2: the first jet.} By (1), $G(X)=\phi([-1]X)-\phi(X)$ has $e^G$ integral. Since
$[-1]X=-X+O(X^2)$ and $l'(0)=1$, the coefficient of $X$ in $G$ is $-2(\phi^\circ)'(0)$. It lies in
$\mathcal O$, and $|2|=1$, so $(\phi^\circ)'(0)\in\mathcal O$.

\emph{Step 3: the second jet.} By Step 1 with $x=X_0$ and $y=X_1$, the series
$G=\beta([a]X_0,X_1)-\beta(X_0,[b]X_1)$ has $e^G$ integral. In the coordinates $z_j=l(X_j)$,
$\beta(x,y)=\phi^\circ(z+w)-\phi^\circ(z)-\phi^\circ(w)$, whose coefficient of $zw$ is $(\phi^\circ)''(0)$; since $[a]X_0=aX_0+
O(X_0^2)$, the coefficient of $X_0X_1$ in $G$ is $(a-b)(\phi^\circ)''(0)$. It lies in $\mathcal O$, and
$|a-b|=1$, so $(\phi^\circ)''(0)\in\mathcal O$.

\emph{Step 4: the higher jets.} By (2) the series
$\Delta=\beta(F(X_0,X_1),X_2)-\beta(X_0,X_2)-\beta(X_1,X_2)$ has $e^\Delta$ integral, so
$\partial_{X_0}\partial_{X_1}\partial_{X_2}\Delta$ is integral. In the coordinates $z_j$, $\Delta$ is
$\phi^\circ(z_0+z_1+z_2)$ plus terms that each depend on at most two of the $z_j$, so
$\partial_{X_0}\partial_{X_1}\partial_{X_2}\Delta=l'(X_0)l'(X_1)l'(X_2)\,(\partial^3\phi)\big(F(F(X_0,X_1),X_2)\big)$.
Setting $X_1=X_2=0$ shows that $l'(X_0)\,(\partial^3\phi)(X_0)$ is integral, and $l'^{-1}$ is
integral, so $\partial^3\phi$ is integral. Since $\partial$ preserves integral series,
$\partial^i\phi$ is integral for every $i\geq3$, and its constant term $(\phi^\circ)^{(i)}(0)$ lies in
$\mathcal O$.
\end{proof}

\begin{proposition}\label{pd-cm-jets}
Let $d$ be a squarefree integer with $d\equiv3\pmod4$, $L$ a lattice with complex multiplication by
$\mathcal O_K$, $K=\mathbf Q(\sqrt{-d})$, $p\neq2$ a prime with $p\nmid d$, and
$\iota:\mathbf C\to\mathbf C_p$ an isomorphism of fields such that $\mathcal E_L$ has coefficients
of absolute value at most $1$ and discriminant of absolute value $1$. Then every jet
$(\partial^i\phi_L)(0)$, $i\geq1$, of the normalizing series of $L$ at $(p,\iota)$
(Definition~\ref{pd-normalizing-series}) has absolute value at most $1$.
\end{proposition}
\leanref{Formalism/LocalToGlobal/TorsionIntegrality/UniformGaussModels/GoodModelAll.lean}{Formalism.L2.Th1.Top.JetsCM}
\leanref{Formalism/LocalToGlobal/TorsionIntegrality/UniformGaussModels/OrthogonalGaussModels/NormLeOneOfQuad.lean}{Formalism.L2.Th1.Top.JetsCM_of}
\leanref{Formalism/LocalToGlobal/TorsionIntegrality/UniformGaussModels/OrthogonalGaussModels/NormLeOneOfQuad.lean}{Formalism.L2.Th1.Top.CMDatum_proof}
\leanref{Formalism/LocalToGlobal/TorsionIntegrality/RamifiedCurveKernels/StatementV.lean}{Formalism.L2.Th1.G.CurveJetsUnram}

\begin{proof}
Write $p=2m+1$ and $\alpha=m+\omega_K$. Then $\alpha+\bar\alpha=2m+1=p$ and
$N=\alpha\bar\alpha=m^2+m+(1+d)/4=(p^2+d)/4$. If $p$ divided $N$, it would divide
$4N-p^2=d$; so $p\nmid N$. Since $\alpha\bar\alpha=N$ and $\iota(\alpha),\iota(\bar\alpha)$ are
integral, $|\iota(\alpha)|=1$. Moreover $\alpha-\bar\alpha=2\omega_K-1=i\sqrt d$, and
$|\iota(i\sqrt d)|^2=|d|=1$. Lemma~\ref{pd-cm-isogeny} applies to $\alpha$, so
Lemma~\ref{pd-sigma-congruences} gives conditions (1) to (3) of
Proposition~\ref{pd-jets-abstract} for $\phi_L$, with $l=\log_{\mathcal E}$, $a=\alpha$ and
$b=\bar\alpha$. That proposition gives the result.
\end{proof}

\begin{definition}\label{pd-translation-system}
Let $K$ be a field complete for a non-archimedean absolute value, with valuation ring
$\mathcal O$, and let
$\mathcal E:\ y^2+a_1xy+a_3y=x^3+a_2x^2+a_4x+a_6$ be a Weierstrass curve with $a_i\in\mathcal O$
and discriminant of absolute value $1$. Let $X=-x/y$ be the formal parameter at the origin, and
for a point $P\in\mathcal E(K)$ let $x_2(P)\in K[[X]]$ be $X^2x(X)$ if $P=0$ and
$X^2\,x(X\oplus P)$ otherwise, where $x(X\oplus P)$ is the $x$-coordinate of the sum of $P$ and the
point with parameter $X$, expanded as a power series in $X$. Let $T$ be an abelian group and
$\pi\colon T\to\mathcal E(K)$ a homomorphism. A \emph{translation system} for $\pi$ is a map
$t\mapsto G_t$ from $T$ to $K[[X]]$ such that:
\begin{enumerate}
\item if $\pi(t)=0$, then $G_t=\varepsilon X$ for some $\varepsilon\in K$ with $|\varepsilon|=1$;
\item if $\pi(t)\neq0$, then $G_t(0)\neq0$;
\item for $s,t\in T$ with $\pi(t)\neq0$,
  \[
   X^2\,G_{s+t}\,G_{s-t}=G_t(0)^2\big(x(\pi(t))\,X^2-x_2(\pi(s))\big)\,G_s^2 .
  \]
\end{enumerate}
\end{definition}
\leanref{Formalism/LocalToGlobal/TorsionIntegrality/TorsionTranslationPowers/IsTransSys.lean}{Formalism.L2.Th1.E.IsTransSys}
\leanref{Formalism/LocalToGlobal/TorsionIntegrality/TorsionTranslationPowers/IsTransSys.lean}{Formalism.L2.Th1.E.x2}

\begin{lemma}\label{pd-translation-units}
Let $K$, $\mathcal O$ and $\mathcal E$ be as in Definition~\ref{pd-translation-system}, let $T$ be
an abelian group and $\pi\colon T\to\mathcal E(K)$ a homomorphism such that every nonzero
$\pi(t)$ has coordinates in $\mathcal O$, and let $(G_t)$ be a translation system for $\pi$. Let
$t\in T$ and $N\geq1$ with $|N|=1$ and $N\pi(t)=0$. Then:
\begin{enumerate}
\item $G_t$ has coefficients in $\mathcal O$, at least one of absolute value $1$;
\item if $\pi(t)\neq0$, then $|G_t(0)|=1$ and $G_t/G_t(0)$ is an integral series with constant
  term $1$.
\end{enumerate}
\end{lemma}
\leanref{Formalism/LocalToGlobal/TorsionIntegrality/TorsionTranslationPowers/StatementT.lean}{Formalism.L2.Th1.E.TransUnit}
\leanref{Formalism/LocalToGlobal/TorsionIntegrality/TorsionTranslationPowers/StatementT.lean}{Formalism.L2.Th1.E.TransGauss}
\leanref{Formalism/LocalToGlobal/TorsionIntegrality/TorsionTranslationPowers/StatementT.lean}{Formalism.L2.Th1.E.TransPower}
\leanref{Formalism/LocalToGlobal/TorsionIntegrality/TorsionTranslationPowers/StatementT.lean}{Formalism.L2.Th1.E.FacShape}

\begin{proof}
If $\pi(t)=0$, then $G_t=\varepsilon X$ with $|\varepsilon|=1$, and (1) holds. Suppose
$\pi(t)\neq0$ and put $c=G_t(0)$, which is nonzero by condition (2) of
Definition~\ref{pd-translation-system}. Let $n=k+1$ be the order of $\pi(t)$; it divides $N$, so
$|n|=1$, and $k\geq1$. Call a series a \emph{unit at order $e$} if it is $aX^eV$ with $|a|=1$ and $V$
an integral series with constant term $1$; a product of units at orders $e$ and $e'$ is a unit at
order $e+e'$.

\emph{Step 1: the factors.} Let $P,Q$ be nonzero points with coordinates in $\mathcal O$ such that
$P+Q$ and $P-Q$ are $0$ or have $x$-coordinate in $\mathcal O$. The series $x(X\oplus P)$ is
integral, with constant term $x(P)$, linear coefficient
$\psi_2(P)=2y(P)+a_1x(P)+a_3$, and, when $\psi_2(P)=0$, quadratic coefficient
$\kappa(P)=3x(P)^2+2a_2x(P)+a_4-a_1y(P)$; these are the first derivatives of $x$ along the invariant
differential at $P$. The reduction of $\mathcal E$ is an elliptic curve, and a point with integral
coordinates reduces to a point other than the origin. If $P\neq\pm Q$, then $P\pm Q$ reduce to points
other than the origin, so $P$ and $\pm Q$ have different reductions and $|x(Q)-x(P)|=1$. If
$P=\pm Q$ and $2Q\neq0$, then $x(Q)=x(P)$, and $2P$ does not reduce to the origin, so
$|\psi_2(P)|=1$. If $P=Q$ and $2Q=0$, then $\psi_2(P)=0$, and $|\kappa(P)|=1$ because the reduction
of $P$ is a nonsingular point of the reduced curve. In each case $x(Q)-x(X\oplus P)$ is integral and
its first coefficient that is not $0$ has absolute value $1$, so it is a unit at order $0$, $1$ or $2$.

\emph{Step 2: the power formula.} For $1\leq j\leq k$, $\pi(jt)\neq0$, so
$x_2(\pi(jt))=X^2x(X\oplus\pi(jt))$, and condition (3) with $s=jt$, divided by $X^2$, gives
$G_{(j+1)t}G_{(j-1)t}=c^2\big(x(\pi(t))-x(X\oplus\pi(jt))\big)G_{jt}^2$. Write $G_0=\varepsilon X$
by condition (1). Induction on $k$ gives
\[
 \varepsilon^kX^k\,G_{(k+1)t}=G_t^{k+1}\prod_{j=0}^{k-1}\Big(c^2\big(x(\pi(t))-
 x(X\oplus\pi((j+1)t))\big)\Big)^{k-j}.
\]

\emph{Step 3: the constant term.} Since $\pi((k+1)t)=0$, $G_{(k+1)t}=\varepsilon'X$ with
$|\varepsilon'|=1$. The points $P=\pi((j+1)t)$, $j<k$, and $Q=\pi(t)$ are nonzero, and $P\pm Q$ are
$0$ or nonzero multiples of $\pi(t)$, so Step 1 applies to them by the hypothesis on $\pi$. So the
product in Step 2 is $c^r aX^mV$ for some $r,m\geq0$, $|a|=1$ and an integral series $V$ with
constant term $1$, and
\[
 \varepsilon^k\varepsilon'X^{k+1}=G_t^{k+1}\,c^r a\,X^m\,V.
\]
The lowest term on the right is $c^{k+1+r}aX^m$, since $G_t(0)=c\neq0$. Hence $m=k+1$ and
$\varepsilon^k\varepsilon'=c^{k+1+r}a$, so $|c|=1$. Dividing,
$V\cdot(G_t/c)^{k+1}=1$, so $(G_t/c)^{k+1}=V^{-1}$ is integral with constant term $1$.

\emph{Step 4: the root.} If $g$ has constant term $1$ and $g^n$ is integral with $|n|=1$, then
$g=(g^n)^{1/n}=\sum_{i\geq0}\binom{1/n}{i}(g^n-1)^i$ is integral, since the rational numbers
$\binom{1/n}{i}$ have absolute value at most $1$ when $|n|=1$. Applied to $g=G_t/c$ and $n=k+1$,
this gives (2), and (1) follows because $|c|=1$.
\end{proof}

\begin{lemma}\label{pd-analytic-translation-system}
Let $L$, $p$ and $\iota$ be as in Definition~\ref{pd-normalizing-series}. For $t\in\mathbf C$ let
$P_t$ be the point $\big(\iota(\wp_L(t)),\iota(\wp_L'(t))/2\big)$ of $\mathcal E_L(\mathbf C_p)$ if
$t\notin L$, and the origin if $t\in L$, and let
\[
 G_t(X)=e^{\phi_L(X)}\cdot\iota\big(U^L_t\sigma^*_L\big)\big(\log_{\mathcal E}X\big)\in
 \mathbf C_p[[X]],
\]
with $\iota(\cdot)$ the Taylor series at $0$ with $\iota$ applied to its coefficients. Then:
\begin{enumerate}
\item $t\mapsto P_t$ is a homomorphism $\mathbf C\to\mathcal E_L(\mathbf C_p)$ with kernel $L$;
\item $(G_t)_{t\in\mathbf C}$ is a translation system for $t\mapsto P_t$
  (Definition~\ref{pd-translation-system}), with $G_t=\pm X$ for $t\in L$, and
  $G_t(0)=\iota\big((U^L_t\sigma^*_L)(0)\big)$;
\item if $t\notin L$ and $Nt\in L$ for some $N\geq1$ with $p\nmid N$, then the coordinates of $P_t$
  have absolute value at most $1$;
\item if $t,t_0\in\mathbf C$ and $N\geq1$ satisfy $Nt\in L$, $Nt_0\in L$ and $t_0\notin L$, then
  \[
   U^L_t\vartheta_{L,t_0}=\zeta\,\big((U^L_{t_0}\sigma^*_L)(0)\big)^{-1}\,U^L_{t+t_0}\sigma^*_L,
  \]
  where $\zeta=\exp\big(-\pi i\operatorname{Im}H_L(t,t_0)\big)$ satisfies $\zeta^{2N^2}=1$.
\end{enumerate}
\end{lemma}
\leanref{Formalism/LocalToGlobal/TorsionIntegrality/TorsionTranslationPowers/StatementA.lean}{Formalism.L2.Th1.E.PtOfAdd}
\leanref{Formalism/LocalToGlobal/TorsionIntegrality/TorsionTranslationPowers/StatementA.lean}{Formalism.L2.Th1.E.PtOfZero}
\leanref{Formalism/LocalToGlobal/TorsionIntegrality/TorsionTranslationPowers/StatementA.lean}{Formalism.L2.Th1.E.AnalyticSys}
\leanref{Formalism/LocalToGlobal/TorsionIntegrality/TorsionTranslationPowers/StatementA.lean}{Formalism.L2.Th1.E.PtOfInt}
\leanref{Formalism/LocalToGlobal/TorsionIntegrality/TorsionTranslationPowers/StatementA.lean}{Formalism.L2.Th1.E.OpUTh0Root}

\begin{proof}
(1) The map $t\mapsto(\wp_L(t),\wp_L'(t)/2)$ is the analytic parametrization of the curve
$y^2=x^3-\frac14g_2(L)x-\frac14g_3(L)$, a homomorphism from $\mathbf C/L$ onto its points
\cite[Proposition~VI.3.6]{Sil09}; applying $\iota$ to the coordinates gives (1).

(2) Since $Xe^{-\phi_L(X)}=\iota(\sigma^*_L)(\log_{\mathcal E}X)$, $G_t$ is the expansion of
$(U^L_t\sigma^*_L)(z)/\sigma^*_L(z)$ times $X$, in the coordinate $z=\log_{\mathcal E}X$. Its
constant term is $\iota((U^L_t\sigma^*_L)(0))=\iota\big(\sigma^*_L(t)e^{-\pi H_L(t,t)/2}\big)$,
which is $0$ if and only if $t\in L$, because the zeros of $\sigma_L$ are the points of $L$. This
is condition (2). For $t\in L$, the quasi-periodicity of $\sigma_L$ gives
$U^L_t\sigma^*_L=\psi(t)\sigma^*_L$ with $\psi(t)=\pm1$, so $G_t=\pm X$; this is condition (1).
For condition (3), the addition formula of Step 1 in the proof of
Lemma~\ref{pd-sigma-congruences}, applied to $u=z+s$ and $v=t$, together with the computation of the
exponential factors of the operators $U^L$, gives
\[
 (U^L_{s+t}\sigma^*_L)(z)\,(U^L_{s-t}\sigma^*_L)(z)=(U^L_t\sigma^*_L)(0)^2\big(\wp_L(t)-
 \wp_L(z+s)\big)\,(U^L_s\sigma^*_L)(z)^2
\]
near $z=0$; the exponents agree because
$H_L(z,s+t)+H_L(z,s-t)=2H_L(z,s)$ and $H_L(s+t,s+t)+H_L(s-t,s-t)=2H_L(s,s)+2H_L(t,t)$. If
$s\notin L$, the expansion of $\iota(\wp_L(z+s))$ in $X$ is $x(X\oplus P_s)$, by (1) and the
compatibility of the analytic and formal group laws; multiplying by $X^2e^{2\phi_L}$ gives
condition (3). If $s\in L$, then $\wp_L(z+s)=\wp_L(z)$, whose expansion is $x(X)$, and the same
computation gives condition (3) with $x_2(0)=X^2x(X)$.

(3) The point $P_t$ is nonzero and $NP_t=0$. Since $\mathcal E_L$ has integral coefficients and
$N$ is prime to $p$, its coordinates are integral \cite[Theorem~VII.3.4]{Sil09}.

(4) Expanding the definition of $U^L$, $U^L_tU^L_{t_0}=\zeta\,U^L_{t+t_0}$ with
$\zeta=\exp\big(-\pi(H_L(t,t_0)-H_L(t_0,t))/2\big)=\exp\big(-\pi i\operatorname{Im}H_L(t,t_0)\big)$.
Apply this to $\vartheta_{L,t_0}=((U^L_{t_0}\sigma^*_L)(0))^{-1}U^L_{t_0}\sigma^*_L$. Since
$Nt,Nt_0\in L$ and $\operatorname{Im}H_L$ takes integer values on $L\times L$,
$N^2\operatorname{Im}H_L(t,t_0)\in\mathbf Z$, so $\zeta^{2N^2}=1$.
\end{proof}

\begin{lemma}\label{pd-theta-translates}
Let $L$, $p\neq2$ and $\iota$ be as in Definition~\ref{pd-normalizing-series}, and let
$\omega_0\in L$ with $\omega_0/2\notin L$. Then:
\begin{enumerate}
\item the series $e^{\phi_L(X)}\cdot\iota\big(\vartheta_{L,\omega_0/2}\big)\big(\log_{\mathcal E}X\big)$
  is integral with constant term $1$;
\item for $s\in\mathbf C$ and $N\geq1$ with $p\nmid N$ and $Ns\in L$, the series
  \[
   e^{\phi_L(X)}\cdot\iota\big(U^L_s\vartheta_{L,\omega_0/2}\big)\big(\log_{\mathcal E}X\big)
  \]
  is integral, and at least one of its coefficients has absolute value $1$.
\end{enumerate}
Here $\iota(\cdot)$ is the Taylor series at $0$ with $\iota$ applied to its coefficients
(Definition~\ref{pd-theta-function}).
\end{lemma}
\leanref{Formalism/LocalToGlobal/TorsionIntegrality/QuadraticSemicharacters/QuadraticSemicharactersFactVersion.lean}{Formalism.L2.Th1.H.TransInt}
\leanref{Formalism/LocalToGlobal/TorsionIntegrality/GaussianEigenFrames/StatementC.lean}{Formalism.L2.Th1.GD.TransIntC}
\leanref{Formalism/LocalToGlobal/TorsionIntegrality/TorsionTranslationPowers/StatementA.lean}{Formalism.L2.Th1.E.ThetaZeroAnalytic}

\begin{proof}
Let $(G_t)$ and $(P_t)$ be as in Lemma~\ref{pd-analytic-translation-system}. For $t\in\mathbf C$
and $N\geq1$ with $p\nmid N$ and $Nt\in L$, let $T_t=\mathbf Zt$ and restrict $t'\mapsto P_{t'}$
and $G$ to $T_t$. Every element $t'$ of $T_t$ has $Nt'\in L$, so by
Lemma~\ref{pd-analytic-translation-system}(3) every nonzero $P_{t'}$ has integral coordinates, and
by Lemma~\ref{pd-analytic-translation-system}(2) the restriction of $G$ is a translation system.
So Lemma~\ref{pd-translation-units} applies to $t$, with $|N|=1$ and $NP_t=0$.

(1) Put $t_0=\omega_0/2$ and $c=(U^L_{t_0}\sigma^*_L)(0)$. Then $2t_0\in L$, $t_0\notin L$ and
$p\nmid2$, so by Lemma~\ref{pd-translation-units}(2), $|\iota(c)|=|G_{t_0}(0)|=1$ and
$G_{t_0}/\iota(c)$ is integral with constant term $1$. Since
$\vartheta_{L,t_0}=c^{-1}U^L_{t_0}\sigma^*_L$, this series is
$e^{\phi_L(X)}\iota(\vartheta_{L,t_0})(\log_{\mathcal E}X)$.

(2) By Lemma~\ref{pd-analytic-translation-system}(4) with $2N$ in place of $N$,
$U^L_s\vartheta_{L,t_0}=\zeta c^{-1}U^L_{s+t_0}\sigma^*_L$ with $\zeta$ a root of unity, so the
series of the statement is $\iota(\zeta c^{-1})\,G_{s+t_0}$. Here $|\iota(\zeta)|=1$, and
$|\iota(c)|=1$ by the proof of (1). The point $s+t_0$ satisfies $2N(s+t_0)\in L$ and $p\nmid2N$,
so by Lemma~\ref{pd-translation-units}(1), $G_{s+t_0}$ is integral with a coefficient of
absolute value $1$. Multiplying by a constant of absolute value $1$ keeps both properties.
\end{proof}

\begin{lemma}\label{pd-fourfold-product}
Let $(a_0,b_0,c_0,d,\lambda)$ be a CM datum with lattice $L$ (Definition~\ref{pd-cm-datum}). Let
$n$ be a finite set, $h$ an $n\times n$ complex matrix with form $H$, and
$\Lambda\subset\mathbf C^n$ a subgroup that is self-dual for $H$ and satisfies
$\theta_K\Lambda\subset\Lambda$, and let $\theta$ be a theta function for $(H,\Lambda)$ with
normalization $\sqrt d/2$. Let $p\neq2$ be a prime and $\iota:\mathbf C\to\mathbf C_p$ an
isomorphism of fields such that $\mathcal E_L$ has coefficients of absolute value at most $1$ and
discriminant of absolute value $1$. Let $g$ and $\varepsilon$ satisfy conditions (1) to (4) of
Definition~\ref{pd-cm-datum}, and let $\gamma$ be the $(4\times n)\times(4\times n)$ matrix with
$\gamma_{w,w'}=\iota\big((g^{-1}e_{w'})_w\big)$. Let $t\in\mathbf C^n$ and $N\geq1$ with $p\nmid N$
and $Nt\in\Lambda$, and let $T$ be the Taylor series of $U_t\theta$ at $0$. Then there is an
integral series $G$ in the variables $X_w$, $w\in4\times n$, such that
\[
 \iota(T)^{\otimes4}(\gamma Z)=e^{-\sum_w\phi_L(X_w)}\,G(X)\Big|_{X_w=\exp_{\mathcal E}(Z_w)},
\]
where $\iota(T)^{\otimes4}$ is the product of four copies of $\iota(T)$ in the four blocks of
variables (Lemma~\ref{pd-jets-closure}(5)).
\end{lemma}
\leanref{Formalism/LocalToGlobal/TorsionIntegrality/GaussianEigenFrames/StatementC.lean}{Formalism.L2.Th1.GD.TensStd}
\leanref{Formalism/LocalToGlobal/TorsionIntegrality/GaussianEigenFrames/GramFrameComposition/LocNormFreeProof.lean}{Formalism.L2.Th1.GD.TensStd_of}
\leanref{Formalism/LocalToGlobal/TorsionIntegrality/QuadraticSemicharacters/QuadraticSemicharactersFactVersion.lean}{Formalism.L2.Th1.H.quadraticSemicharacterIdentity}
\leanref{Formalism/LocalToGlobal/TorsionIntegrality/GaussianEigenFrames/StatementA.lean}{Formalism.L2.Th1.GD.ShapeToZ}

\begin{proof}
Let $\Lambda'=g(\Lambda^4)$, let $\Lambda_0=L^{4\times n}$, and let $c$ run over a set of
representatives of $\Lambda'/(\Lambda'\cap\Lambda_0)$, whose number $I$ is prime to $p$ by
condition~(2). Put $\omega_0=\lambda a_0$; then $\omega_0\in L$ and $\omega_0/2\notin L$.

\emph{Step 1: the transfer identity.} Translating condition~(4) of Definition~\ref{pd-cm-datum} by
$(t,t,t,t)$ and using $H_{\mathrm{std}}(gz,gw)=H^{\oplus4}(z,w)$ gives, for $y\in\mathbf C^{4\times n}$,
\[
 \prod_{a=1}^4(U_t\theta)\big((g^{-1}y)_a\big)=\sum_c\varepsilon I^{-1/2}\beta(c)\,
 e^{-\pi iE_{\mathrm{std}}(g(t,t,t,t),c)}\prod_{w}\big(U^L_{g(t,t,t,t)_w+c_w}\vartheta_{L,\omega_0/2}\big)(y_w).
\]

\emph{Step 2: the constants.} $|\iota(\varepsilon)|\leq1$ and $I$ is prime to $p$. The number
$\beta(c)$ is a root of unity: $\beta$ is a semi-character of $\Lambda'$ for $H_{\mathrm{std}}$
with values $\pm1$ on $\Lambda'\cap\Lambda_0$, and $Ic\in\Lambda'\cap\Lambda_0$. The number
$e^{-\pi iE_{\mathrm{std}}(g(t,t,t,t),c)}$ is a root of unity because $Ng(t,t,t,t)\in\Lambda'$ and
$\Lambda'$ is self-dual for $H_{\mathrm{std}}$, so that $E_{\mathrm{std}}(g(t,t,t,t),c)\in N^{-1}\mathbf Z$.
So each coefficient has absolute value at most $1$ under $\iota$.

\emph{Step 3: the translates.} Each point $s=g(t,t,t,t)_w+c_w$ satisfies $NIs\in L$, because
$NIg(t,t,t,t)$ and $Ic$ lie in $\Lambda'\cap\Lambda_0$; and $p\nmid NI$. By Lemma~\ref{pd-theta-translates}(2), the Taylor series of
$U^L_s\vartheta_{L,\omega_0/2}$, with $\iota$ applied and written in the variable
$X_w=\exp_{\mathcal E}(y_w)$, is $e^{-\phi_L(X_w)}$ times an integral series.

\emph{Step 4.} Taking Taylor series at $y=0$ in Step 1 and applying $\iota$, the left side becomes
$\iota(T)^{\otimes4}(\gamma Z)$ with $Z=y$, and the right side becomes
$e^{-\sum_w\phi_L(X_w)}$ times a combination with coefficients of absolute value at most $1$ of
products of integral series in the $X_w$. Call that combination $G$.
\end{proof}

\begin{lemma}\label{pd-multiplier-coordinates}
In the situation of Lemma~\ref{pd-fourfold-product}, let $J$ and $\nu$ be finite sets, let
$v_j\in\mathbf C^n$ ($j\in J$) satisfy $\omega v_j\in\Lambda$ for all $\omega\in L$, and let
$A_{j,x}\in\mathbf C_p$ ($j\in J$, $x\in\nu$) have absolute value at most $1$. Put
$M_{i,x}=\sum_j\iota(v_{j,i})A_{j,x}$. Then there is a $(4\times n)\times(4\times\nu)$ matrix $C$
with entries of absolute value at most $1$ such that $\gamma C$ is the block diagonal matrix with
four blocks $M$.
\end{lemma}
\leanref{Formalism/LocalToGlobal/TorsionIntegrality/GaussianEigenFrames/StatementB.lean}{Formalism.L2.Th1.GD.GramTransfer}
\leanref{Formalism/LocalToGlobal/TorsionIntegrality/GaussianEigenFrames/GramFrameComposition/MatrixFormDiagBlocksBlockEmbed.lean}{Formalism.L2.Th1.GD.GramTransfer_of}
\leanref{Formalism/LocalToGlobal/TorsionIntegrality/GaussianEigenFrames/StatementB.lean}{Formalism.L2.Th1.GD.MulRing}
\leanref{Formalism/LocalToGlobal/TorsionIntegrality/GaussianEigenFrames/StatementA.lean}{Formalism.L2.Th1.GD.FrameCompose}

\begin{proof}
For $b\in\{1,\dots,4\}$ and $j\in J$ let $e_b(v_j)\in\mathbf C^{4\times n}$ be $v_j$ in the
$b$-th block and $0$ elsewhere, and put $Y_{b,j}=g(e_b(v_j))$. For $\omega\in L$,
$\omega Y_{b,j}=g(e_b(\omega v_j))\in g(\Lambda^4)=\Lambda'$. With $I$ as in the proof of
Lemma~\ref{pd-fourfold-product}, $I\Lambda'\subset\Lambda'\cap L^{4\times n}$, so
$\omega\cdot IY_{b,j,w}\in L$ for every coordinate $w$ and every $\omega\in L$. The lattice $L$ has
multiplier ring $\mathcal O_K$, so $IY_{b,j,w}=m+n'\omega_K$ with $m,n'\in\mathbf Z$. Hence
$D_{w,(b,j)}=\iota(Y_{b,j,w})$ has absolute value at most $|I|^{-1}=1$. By construction
$\sum_w\gamma_{(a,i),w}D_{w,(b,j)}=e_b(v_j)_{(a,i)}$ under $\iota$, which is $\iota(v_{j,i})$ if
$a=b$ and $0$ otherwise. Put $C_{w,(b,x)}=\sum_jD_{w,(b,j)}A_{j,x}$; its entries have absolute value
at most $1$, and $(\gamma C)_{(a,i),(b,x)}=\sum_j\big(\sum_w\gamma_{(a,i),w}D_{w,(b,j)}\big)A_{j,x}$
is $M_{i,x}$ if $a=b$ and $0$ otherwise.
\end{proof}

\begin{lemma}\label{pd-jets-from-fourth-power}
Let $L$, $p\neq2$ and $\iota$ be as in Definition~\ref{pd-normalizing-series}, and suppose that
every jet of $\phi_L$ has absolute value at most $1$. Let $n,\nu$ be finite sets, $T$ a power
series over $\mathbf C_p$ in variables indexed by $n$, $\gamma$ a $(4\times n)\times(4\times n)$
matrix and $G$ an integral series such that
$T^{\otimes4}(\gamma Z)=e^{-\sum_w\phi_L(X_w)}G(X)\big|_{X_w=\exp_{\mathcal E}(Z_w)}$. Let $M$ be an
$n\times\nu$ matrix and $C$ a $(4\times n)\times(4\times\nu)$ matrix with entries of absolute value
at most $1$ such that $\gamma C$ is block diagonal with four blocks $M$. Then $T(MX)$ has integral
jets.
\end{lemma}
\leanref{Formalism/LocalToGlobal/TorsionIntegrality/GaussianEigenFrames/StatementA.lean}{Formalism.L2.Th1.GD.JetsLin}
\leanref{Formalism/LocalToGlobal/TorsionIntegrality/GaussianEigenFrames/GramFrameComposition/JetsLin.lean}{Formalism.L2.Th1.GD.JetsLin_proof}

\begin{proof}
\emph{Step 1.} The series $\phi_L(\exp_{\mathcal E}(Z))$ has the jets $(\partial^i\phi_L)(0)$ as its
coefficients times $i!$, so it has integral jets and constant term $0$; by
Lemma~\ref{pd-jets-closure}(2) so does $e^{-\sum_w\phi_L(\exp_{\mathcal E}Z_w)}$.
\emph{Step 2.} The series $\exp_{\mathcal E}(Z)$ has integral jets: its derivative is
$\omega(\exp_{\mathcal E}Z)^{-1}$, an integral series in $\exp_{\mathcal E}(Z)$, so by induction
every derivative of $\exp_{\mathcal E}$ is an integral polynomial in $\exp_{\mathcal E}$ and
$\omega^{-1}(\exp_{\mathcal E})$, and its value at $0$ is integral. By
Lemma~\ref{pd-jets-closure}(4), $G(\exp_{\mathcal E}(Z))$ has integral jets.
\emph{Step 3.} By Steps 1 and 2 and Lemma~\ref{pd-jets-closure}(1), $T^{\otimes4}(\gamma Z)$ has
integral jets, and by Lemma~\ref{pd-jets-closure}(3) so does $T^{\otimes4}(\gamma CX)$.
\emph{Step 4.} Since $\gamma C$ is block diagonal with blocks $M$,
$T^{\otimes4}(\gamma CX)=\big(T(MX)\big)^{\otimes4}$. Lemma~\ref{pd-jets-closure}(5) gives the
result.
\end{proof}

\begin{theorem}\label{pd-theta-jets-integral}
Let $(a_0,b_0,c_0,d,\lambda)$ be a CM datum with lattice $L$ and
$\theta_K=(b_0+i\sqrt d)/2$ (Definition~\ref{pd-cm-datum}). Let $n$ be a finite set, $h$ an
$n\times n$ complex matrix with form $H$, and $\Lambda\subset\mathbf C^n$ a subgroup that is
self-dual for $H$ and satisfies $\theta_K\Lambda\subset\Lambda$. Let $\theta$ be a theta function
for $(H,\Lambda)$ with normalization $\sqrt d/2$ (Definition~\ref{pd-theta-function}). Let
$p\neq2$ be a prime with $p\nmid d$ and let $\iota:\mathbf C\to\mathbf C_p$ be an isomorphism of
fields such that the Weierstrass curve
$y^2=x^3-\tfrac14\iota(g_2(L))x-\tfrac14\iota(g_3(L))$ has coefficients of absolute value at
most $1$ and discriminant of absolute value $1$, and such that $\theta$ is locally normalized at
$(p,\iota)$. Let:
\begin{itemize}
\item $r_k\in\mathbf C_p$ for $k$ in a finite set $\kappa$, with $|r_k|\leq1$;
\item $N\geq1$ with $p\nmid N$, and $t_k\in\mathbf C^n$ with $Nt_k\in\Lambda$ for $k\in\kappa$;
\item $v_j\in\mathbf C^n$ for $j$ in a finite set $J$, with $\omega v_j\in\Lambda$ for all
  $\omega\in L$;
\item $A_{j,x}\in\mathbf C_p$ for $j\in J$ and $x$ in a finite set $\nu$, with
  $|A_{j,x}|\leq1$.
\end{itemize}
Let $M_{i,x}=\sum_j\iota(v_{j,i})A_{j,x}$, and let
$\Theta\in\mathbf C_p[[X_x:x\in\nu]]$ be $\sum_kr_k\,\iota(T_k)\big(MX\big)$, where $T_k$ is the
Taylor series at $0$ of $U_{t_k}\theta$, $\iota(T_k)$ applies $\iota$ to its coefficients, and
$(MX)_i=\sum_xM_{i,x}X_x$. Then $\Theta$ has integral jets, and $|\langle P,\Theta\rangle|\leq1$
for every $P\in\mathbf C_p[X_x:x\in\nu]$ whose coefficients have absolute value at most $1$
(Definition~\ref{pd-integral-jets}).
\end{theorem}
\leanref{Formalism/LocalToGlobal/TorsionIntegrality/GaussianEigenFrames/StatementC.lean}{Formalism.L2.Th1.GD.IntOpLoc}
\leanref{Development/AbelianVarietyTheory/Definitions/GaloisGeometry/FourthPowerTheta.lean}{GlobalTheory.ThLoc.LocJets.jets_odd_of}
\leanref{Development/AbelianVarietyTheory/Definitions/GaloisGeometry/FourthPowerTheta.lean}{GlobalTheory.ThLoc.LocJets.jetsOdd}

\begin{proof}
\emph{Step 1: the normalizing series.} The conditions on the CM datum make $d$ squarefree with
$d\equiv3\pmod4$, and $L$ has complex multiplication by $\mathcal O_K$. By
Proposition~\ref{pd-cm-jets}, every jet of the normalizing series $\phi_L$ at $(p,\iota)$ has
absolute value at most $1$.

\emph{Step 2: the fourfold products.} Let $g$ and $\varepsilon$ be as in the definition of local
normalization and let $\gamma$ be the matrix of Lemma~\ref{pd-fourfold-product}. For each
$k\in\kappa$, that lemma with $t=t_k$ gives an integral series $G_k$ with
$\iota(T_k)^{\otimes4}(\gamma Z)=e^{-\sum_w\phi_L(X_w)}G_k(X)|_{X_w=\exp_{\mathcal E}(Z_w)}$.

\emph{Step 3: the coordinates of the vectors $v_j$.} Lemma~\ref{pd-multiplier-coordinates} gives a
matrix $C$ with entries of absolute value at most $1$ such that $\gamma C$ is block diagonal with
four blocks $M$.

\emph{Step 4: one translate.} By Steps 1 to 3 and Lemma~\ref{pd-jets-from-fourth-power}, each
$\iota(T_k)(MX)$ has integral jets.

\emph{Step 5: the combination.} By Lemma~\ref{pd-jets-closure}(1), $\Theta$ has integral jets. If
the coefficients $p_\kappa$ of $P$ have absolute value at most $1$, every term
$p_\kappa\,\kappa!\,\vartheta_\kappa$ of $\langle P,\Theta\rangle$, where $\vartheta_\kappa$ is the
coefficient of $X^\kappa$ in $\Theta$, has absolute value at most $1$, and so does the sum by the
ultrametric inequality.
\end{proof}

\section{Periods of lifted vectors and their integrality}\label{section-pd-periods}

The periods are those of the iterated lifts of a definite-side datum
(Definition~\ref{th-definite-datum}): for a word $u$, a tuple $\mu$ of Hecke characters,
$\varphi\in\mathcal M_u$ and a top form $F\in\mathcal A_\mu$, the period
$\operatorname{per}_{u,\mu}(\varphi,F)$ is $\mathcal P_u(\operatorname{ev}_u(\operatorname{col}\varphi),F)$,
with $\varphi$ identified with an element of $M_{u,\mathrm{ram}}$ and with the iterated lift and
the period map of the column chain $C_{u,\mu}$ twisted by $\mu$. Their integrality is measured
with the following further parts of the datum.

\begin{definition}\label{pd-period-data}
Let $D$ be a definite-side datum (Definition~\ref{th-definite-datum}) with coefficient field $L$
and embedding $\iota:L\to\mathbf C$, base word $u_0$, column chains $C_u$ and twisted column chains
$C_{u,\mu}$.
\begin{enumerate}
\item For a finite place $w$ of $L$, $v_w$ is its normalized additive valuation, $e_w$ its
  ramification index over $\mathbf Q$ and $p_w$ its residue characteristic.
\item $D_\Gamma(u,\mu)$ is the stabilizer denominator of $C_{u,\mu}$
  (Definition~\ref{th-column-chain}) and $D_\Gamma(u)=D_\Gamma(u,\mu_0)$. The stabilizer
  denominators of $D$ are \emph{defined} if $D_\Gamma(u,\mu)=D_\Gamma(u_0)$ for all $u$ and $\mu$,
  and $|\Gamma_c|$ divides $D_\Gamma(u_0)$ for every class $c$ at every level $i\leq m$ of every
  chain $C_{u,\mu}$.
\item Each space $\mathcal A_\mu$ of top forms carries a finite spanning set $S_\mu$ and a
  self-adjoint idempotent $e_\mu$. With $\mathcal O$ the ring of coefficients of $D$, whose image in
  $L$ consists of algebraic integers and contains the image of $\mathcal O_E$, the \emph{integral
  top forms} for $\mu$ are the elements $\sum_{s\in S_\mu}\iota(c_s)s$, $c_s\in\mathcal O$, that
  lie in the image of $e_\mu$.
\item For each word $u$ and level $j$, the lattice of forms is an $\mathcal O_L$-lattice
  $\Lambda_{u,j}$ in an $L$-vector space $V_{u,j}$, with an $\iota$-semilinear injection into the
  functions on $\mathrm{cls}_j\times Y_j$ (Definition~\ref{th-definite-datum}(7)). The \emph{point
  lattice} of $u$ is an $\mathcal O_L$-lattice $\Lambda^{\mathrm{pt}}_u$ in an $L$-vector space
  $V^{\mathrm{pt}}_u$, with an $\iota$-semilinear injection $j^{\mathrm{pt}}_u$ into the space
  $\mathcal W_u$ of the period map of $u$; it comes with a finite set of \emph{point classes}, each
  with a finite stabilizer, a map from the point classes to $\mathrm{cls}_m$ and injective
  homomorphisms from their stabilizers to those of the image classes. The \emph{evaluation} of $u$
  is defined over $L$ and
  integral if $\operatorname{ev}_u$ maps the image of $V_{u,m}$ into $j^{\mathrm{pt}}_u(V^{\mathrm{pt}}_u)$
  and the image of $\Lambda_{u,m}$ into $j^{\mathrm{pt}}_u(\Lambda^{\mathrm{pt}}_u)$.
\item \emph{Twisted forms} for $D$ give the same data for the chains $C_{u,\mu}$,
  $\mu\neq\mu_0$: lattices of forms $\Lambda_{u,\mu,j}$, an evaluation $\operatorname{ev}_{u,\mu}$,
  a period pairing $\mathcal P_{u,\mu}$ of its target with $\mathcal A_\mu$, and a twisted point
  lattice $\Lambda^{\mathrm{pt}}_{u,\mu}$ with injection $j^{\mathrm{pt}}_{u,\mu}$.
\item Each place $w$ of $L$ has a tuple $\mu_w$, equal to $\mu_0$ or to the twisting tuple of
  $w$; a tuple is \emph{used} if it is $\mu_0$ or $\mu_w$ for some $w$. Divisibility data for $D$
  attach to each tuple $\mu$ a finite set $Q(\mu)$ of primes, the \emph{exceptional primes} of
  $\mu$, with $Q(\mu_0)=\emptyset$ and $p_w\notin Q(\mu_w)$ for every $w$.
\item For each word $u$ the division lattice $N_u$ (Definition~\ref{th-definite-datum}(9)) lies in
  an $L$-vector space $V(N_u)$, with an $\iota$-semilinear map $j_N:V(N_u)\to\mathcal M_u$ that
  carries a basis to a basis, and an injective $L$-linear map $\nu_u:V(N_u)\to V_{u,0}$ onto the
  vectors whose image lies in $M_{u,\mathrm{ram}}$, compatible with $j_N$ and the isomorphism
  $\mathcal M_u\cong M_{u,\mathrm{ram}}$. We put $M^{\mathrm{ram}}_u=\nu_u^{-1}(\Lambda_{u,0})$.
\end{enumerate}
\end{definition}
\leanref{Formalism/LocalRingCardinality/DataErr.lean}{Formalism.BasicNotions.DGamDef}
\leanref{Formalism/LocalRingCardinality/DataRows.lean}{Formalism.BasicNotions.Rows.Data.V0}
\leanref{Formalism/LocalRingCardinality/DataForms.lean}{Formalism.BasicNotions.MemberForms}
\leanref{Formalism/LocalRingCardinality/DataKerP.lean}{Formalism.BasicNotions.TwForms}
\leanref{Formalism/LocalRingCardinality/DataDIV.lean}{Formalism.BasicNotions.DivX}
\leanref{Formalism/SharpRestriction/Divisibility/LocalQuadraticParts.lean}{Formalism.MainNotions.LocalQuadraticParts.UsedT}
\leanref{Formalism/LocalRingCardinality/DataDIVDefs.lean}{Formalism.BasicNotions.DV.MramL}

\begin{definition}\label{pd-period-integrality}
Let $D$ be a definite-side datum with the data of Definition~\ref{pd-period-data}.
\emph{Period integrality} holds for $D$ if, whenever the stabilizer denominators of $D$ are
defined, the following holds for every word $u$, every $x\in\Lambda^{\mathrm{pt}}_u$ and every
integral top form $F$ for $\mu_0$: there is $y\in L$ with
$\iota(y)=\mathcal P_u\big(j^{\mathrm{pt}}_u(x),F\big)$ such that, for every finite place $w$ of
$L$,
\[
 y=0\quad\text{or}\quad v_w(y)\geq-e_w\,v_{p_w}(D_\Gamma(u)).
\]
\end{definition}
\leanref{Formalism/SharpRestriction/Divisibility/WeightDegreeSizes.lean}{Formalism.MainNotions.WeightDegreeSizes.PeriodClassSumWeyl}

Since $v_w(D_\Gamma(u))=e_w\,v_{p_w}(D_\Gamma(u))$, the condition says that the period is $0$ or
$\iota(y)$ with $D_\Gamma(u)\,y\in\mathcal O_L$.

\begin{definition}\label{pd-twisted-period-integrality}
Let $D$ be a definite-side datum with divisibility data and twisted forms
(Definition~\ref{pd-period-data}). \emph{Twisted period integrality} holds if, whenever the
stabilizer denominators of $D$ are defined, the following holds for every word $u$, every used
tuple $\mu\neq\mu_0$, every $x\in\Lambda^{\mathrm{pt}}_{u,\mu}$ and every integral top form $F$
for $\mu$: there is $y\in L$ with
$\iota(y)=\mathcal P_{u,\mu}\big(j^{\mathrm{pt}}_{u,\mu}(x),F\big)$ such that, for every finite
place $w$ of $L$ with $p_w\notin Q(\mu)$,
\[
 y=0\quad\text{or}\quad v_w(y)\geq-e_w\,v_{p_w}(D_\Gamma(u)).
\]
\end{definition}
\leanref{Formalism/LocalToGlobal/GroupLattices/AdmissibleAtGeneratorPlaces.lean}{Formalism.L2.G.Adm2InU.TwistedPeriodAdmissible}

Both conditions hold for the definite-side data built from true columns whose period rows are
integral (Theorem~\ref{pd-period-integrality-from-rows}), and hence for the members of the families
of the construction (Theorem~\ref{pd-period-integrality-along-families} and
Remark~\ref{pd-integrality-inputs}).

\section{The multiplicity bound}\label{section-pd-multiplicity}

\subsection{Size data and local models}\label{subsection-pd-size-data-local-models}

\begin{definition}\label{pd-size-data}
Let $D$ be a definite-side datum with ranks $k,m,t,n$, $f=2n$ and $N(j)=k+j(t+1)$
(Definition~\ref{th-ranks}), coefficient field $L$, and host with auxiliary primes $q,d$, integer
$Q_S$, level primes $g_0,g_1$ and top classes $T=T_{\mathrm{top}}$ (Definition~\ref{hf-host}). A prime $p$ is
\emph{inert} for $D$ if $p\nmid Q_S$, $p\neq q,d$ and $-d$ is not a square modulo $p$. An inert
prime $p$ is a \emph{signal} prime if the first sign of $\operatorname{art}(p)$
(Definition~\ref{hf-frame}) is $1$, and is of \emph{inert class} otherwise. \emph{Size data} for
$D$ consist of the following, subject to further conditions that this chapter does not use.
\begin{enumerate}
\item A scale $(x,\varepsilon,J)$ (Definition~\ref{hf-windows}), with $b=\exp(x^{\varepsilon})$,
  $X'=\exp(x^{\varepsilon/3})$, $\delta_1=x^{-\varepsilon/8}$ and the further derived quantities
  $a=xb$, $b'=a/X'$, $w''=b'/(xn)$, $X''=a/w''=xnX'$ and $p_0=k^2x^{26}$.
\item Exponents $a_t,b_t\in\mathbf Z$ for $t\in T$ with $a\leq a_t\leq a+a/(kx^{13})$ and
  $b_t\geq b'$.
\item A finite set of places of $L$, the \emph{support}, and real numbers $\mathrm{gain}_j(w)$,
  $F_j(w)$, $E_{u,j}(w)$, $\tau_w$, $\iota_w$ for lifts $1\leq j\leq m$, words $u$ and places $w$,
  all zero off the support; $\mathrm{gain}_j(w)=0$ when $p_w\in\{g_0,g_1\}$; $F_j(w)=0$ when
  $\mu_w=\mu_0$; and $\mathrm{gain}_j(w)=\tau_w=\iota_w=0$ when $p_w$ is not inert for $D$.
\item A set of \emph{dropped} primes, containing no auxiliary prime $T_u$. We put
  $\mathrm{gain}^{\flat}_j(w)=0$ if $p_w$ is dropped and $\mathrm{gain}^\flat_j(w)=\mathrm{gain}_j(w)$
  otherwise, and
  \[
   s(u,w)=\sum_{j=1}^{m}\big(\mathrm{gain}^\flat_j(w)+F_j(w)-E_{u,j}(w)\big).
  \]
\item For each word $u$ an $\mathcal O_L$-submodule $N^{\circ}_u$ of $V(N_u)$.
\item For each place $w$ a subset $T_w\subset T$, equal to $T$ when $p_w$ is a signal prime or
  $p_w\geq X'$.
\item Weights $w_{j,t,p,i}\geq1$ for lifts $j$, $t\in T$, primes $p$ and integers $i\geq0$.
\item Rational numbers $\operatorname{len}_{u,w}$, and real numbers $\eta>0$ (the width of the
  windows), $S_c=4k\Delta_S$ (the cloud size, with $\Delta_S$ as in
  Definition~\ref{th-definite-datum}(8)) and $\Lambda_G$.
\item A constant $C_\Gamma>0$ such that $\log\Gamma_P=C_\Gamma\,m^5f\,x^{\varepsilon}$, where
  $\Gamma_P>1$ is the number of Definition~\ref{th-definite-datum}(3) at every word $u$.
\item The splitting bound $B$ of the host (Definition~\ref{hf-host}) satisfies
  $\lfloor x^{\varepsilon/(5J)}\rfloor\leq B$.
\end{enumerate}
\end{definition}
\leanref{Formalism/LocalRingCardinality/DataParams₂.lean}{Formalism.BasicNotions.SizeData₂}
\leanref{Formalism/LocalRingCardinality/DataParams₂.lean}{Formalism.BasicNotions.SizeData₂.sg}
\leanref{Formalism/LocalRingCardinality/DataHost.lean}{Formalism.BasicNotions.primeClass}

The quantity $\mathrm{gain}_j(w)$ is the divisibility won by the $j$-th lift at $w$, $F_j(w)$ is
its local term at a twisted place and $E_{u,j}(w)$ its local error. All are measured in units of
$v_{p_w}$, so a valuation bound at $w$ multiplies them by $e_w$.

\begin{definition}\label{pd-multiplicity-bound}
Let $D$ be a definite-side datum with size data (Definition~\ref{pd-size-data}). The
\emph{multiplicity bound} holds for $D$ if for every word $u$, every finite place $w$ of $L$, every
$\varphi\in N^{\circ}_u$ and every integral top form $F$ for $\mu_w$ there is $y\in L$ with
$\iota(y)=\operatorname{per}_{u,\mu_w}(j_N\varphi,F)$ and
\[
 y=0\quad\text{or}\quad v_w(y)\geq\big\lfloor e_w\,s(u,w)\big\rfloor .
\]
\end{definition}
\leanref{Formalism/SharpRestriction/Divisibility/LocalQuadraticParts.lean}{Formalism.MainNotions.LocalQuadraticParts.adm2}

The multiplicity bound plays the role of the order of vanishing in Stepanov's method: each lift
contributes $\mathrm{gain}_j(w)$ to the $w$-adic divisibility of every period of a lifted division
vector, up to the local terms $F_j(w)-E_{u,j}(w)$.

\begin{definition}\label{pd-lift-valuation}
Let $D$ be a definite-side datum with twisted forms, $u$ a word, $\mu$ a tuple and $w$ a finite
place of $L$. For $x\in V_{u,0}$ and $s\in\mathbf R$, the iterated lift of $x$ along $C_{u,\mu}$
\emph{has $w$-valuation at least $s$} if it is the image of $a z$ for some $a\in L$ and some $z$ in
the span of $\Lambda_{u,\mu,m}$ over the valuation ring of $w$, with $a=0$ or $v_w(a)\geq e_ws$.
\end{definition}
\leanref{Formalism/SharpRestriction/GrowthBounds/WeilJacquetShalikaLaw.lean}{Formalism.MainNotions.WeilJacquetShalikaLaw.ColGe}
\leanref{Formalism/SharpRestriction/GrowthBounds/WeilJacquetShalikaLaw.lean}{Formalism.MainNotions.WeilJacquetShalikaLaw.VGe}

\begin{lemma}\label{pd-period-of-divisible-lift}
Let $D$ be a definite-side datum with divisibility data and twisted forms whose stabilizer
denominators are defined, and suppose that period integrality and twisted period integrality hold
(Definitions~\ref{pd-period-integrality} and~\ref{pd-twisted-period-integrality}). Suppose that the
evaluation of every word is defined over $L$ and integral, and that the same holds for the twisted
chains at used tuples $\mu\neq\mu_0$. Let $u$ be a word, $w$ a finite place, $\varphi\in V(N_u)$,
$s\in\mathbf R$, and $\mu$ a used tuple with $p_w\notin Q(\mu)$ such that the iterated lift of
$\nu_u(\varphi)$ along $C_{u,\mu}$ has $w$-valuation at least $s$. Then for every integral top form
$F$ for $\mu$ there is $y\in L$ with $\iota(y)=\operatorname{per}_{u,\mu}(j_N\varphi,F)$ and $y=0$
or $v_w(y)\geq e_w\big(s-v_{p_w}(D_\Gamma(u))\big)$.
\end{lemma}
\leanref{Formalism/LocalToGlobal/GroupLattices/AdmissibleAtGeneratorPlaces.lean}{Formalism.L2.G.Adm2InU.period_geTU}
\leanref{Formalism/LocalToGlobal/GroupLattices/AdmissibleAtGeneratorPlaces.lean}{Formalism.L2.G.Adm2InU.pair0U}

\begin{proof}
Take $\mu=\mu_0$; the twisted case is the same with $\mathcal P_{u,\mu}$ and the twisted point
lattice. By hypothesis the iterated lift of $\nu_u(\varphi)$ is the image of $az$, with $a=0$ or
$v_w(a)\geq e_ws$, and $z$ in the $w$-local span of $\Lambda_{u,m}$. The period is linear, so
$\operatorname{per}_{u,\mu_0}(j_N\varphi,F)=\iota(a)\,\mathcal P_u(\operatorname{ev}_u z,F)$, where
$z$ also denotes its image.

For $z\in\Lambda_{u,m}$, integrality of the evaluation gives $\operatorname{ev}_uz=j^{\mathrm{pt}}_u(x)$
with $x\in\Lambda^{\mathrm{pt}}_u$, so period integrality gives $y_0\in L$ with
$\iota(y_0)=\mathcal P_u(\operatorname{ev}_uz,F)$ and $y_0=0$ or
$v_w(y_0)\geq-e_wv_{p_w}(D_\Gamma(u))$. The set of $z$ for which such a $y_0$ exists contains $0$,
is closed under addition by the ultrametric inequality, and is closed under multiplication by
elements of $L$ with $v_w\geq0$. It is therefore the whole $w$-local span. Put $y=ay_0$; then
$y=0$ or $v_w(y)=v_w(a)+v_w(y_0)\geq e_w\big(s-v_{p_w}(D_\Gamma(u))\big)$.
\end{proof}

\begin{definition}\label{pd-family}
Let $0<\varepsilon<1/52$. A \emph{family of definite-side data} with exponent $\varepsilon$
consists of:
\begin{enumerate}
\item a threshold $x_1>1$ and a real number $J\geq1$;
\item a set $W\subset\mathbf R$ of \emph{scales}, unbounded above, and for every real $x$ a set
  $I_x$ of \emph{members at $x$}, nonempty for $x\in W$ and for $x=0$;
\item for every $x$ and every member $i\in I_x$ a definite-side datum $D_{x,i}$
  (Definition~\ref{th-definite-datum}) with size data $L_{x,i}$
  (Definition~\ref{pd-size-data}), whose scale is $(x,\varepsilon,J)$ when $x>x_1$;
\item for every $x>x_1$ and every member, the further choices to which the links refer:
  parameter choices, sizes and ties, divisibility data, twisted forms, gain data and error data;
\item a finite list of numerical constants, the same for all members at all scales.
\end{enumerate}
All members at one scale $x$ have the same ranks, the same images of their top classes in
$\mathbf Z$, the same set $T_2$, the same primes $q,d$, integer $Q_S$, Artin map, level primes
$g_0,g_1,p_{de},p'_{de}$ and discriminant term, and the same constant $C_\Gamma$ of their size data;
their words are identified with one set; and $o\leq C_o$ and $Q_S\leq C_on$ at scales $x>x_1$ for one
of the constants $C_o$. At every scale $x>x_1$ the members \emph{cover the base points}: for every
member $i_0$ at $x$ and every $u_0\in A$ there is a member $i$ at $x$ whose host has base point $u_0$
(Definition~\ref{hf-host}(6)) and whose exponents $a_t,b_t$ agree with those of $i_0$ at the top
classes with the same number.

A \emph{class} of families is a property of families of every exponent. The constants of a family
$\mathcal F$ are arranged in lists $c_{\mathrm{pre}}(\mathcal F)$, $c_{\mathrm{out}}(\mathcal F)$,
$c_{\mathrm{mid}}(\mathcal F,C_H)$ and $c_{\mathrm{last}}(\mathcal F)$, where the real number
$C_H$, the \emph{height constant}, enters the middle list. A \emph{largeness threshold}
$\kappa=(\kappa_{\mathrm{pre}},\kappa_{\mathrm{mid}},\kappa_{\mathrm{last}})$ consists of a
real number and two functions from finite lists of real numbers to $\mathbf R$; the family is
\emph{large for $(\kappa,C_H)$} if every constant in $c_{\mathrm{pre}}$ is at least
$\kappa_{\mathrm{pre}}$, every constant in $c_{\mathrm{mid}}(\mathcal F,C_H)$ is at least
$\kappa_{\mathrm{mid}}(c_{\mathrm{out}})$, and every constant in $c_{\mathrm{last}}$ is at least
$\kappa_{\mathrm{last}}(c_{\mathrm{out}},c_{\mathrm{mid}}(\mathcal F,C_H))$. A property
$P(\mathcal F,C_H)$ of a family and a real number holds \emph{eventually along} a class
$\mathcal C$ if there is $\varepsilon_1>0$ such that for every $\varepsilon\in(0,\varepsilon_1]$
there is a largeness threshold $\kappa$ such that $P(\mathcal F,C_H)$ holds for every
$\mathcal F\in\mathcal C$ of exponent $\varepsilon$ and every $C_H$ for which $\mathcal F$ is
large for $(\kappa,C_H)$.

A property of the members of a family holds \emph{at all large scales} if there is $x_0$ such
that it holds for $D_{x,i}$ and $L_{x,i}$ for every $x\geq x_0$ with $x>x_1$ and every member $i$
at $x$.
\end{definition}
\leanref{Formalism/SharpRestriction/FamCore.lean}{Formalism.MainNotions.FamCore}
\leanref{Formalism/SharpRestriction/ParConsts.lean}{Formalism.MainNotions.FamP}
\leanref{Formalism/SharpRestriction/ParamsHeightRankConstant.lean}{Formalism.MainNotions.FamClass}
\leanref{Formalism/SharpRestriction/ParamsHeightRankConstant.lean}{Formalism.MainNotions.FamP.Large}
\leanref{Formalism/SharpRestriction/ParamsHeightRankConstant.lean}{Formalism.MainNotions.Closed}

For a class $\mathcal A$ of families, $\mathcal D_3=\operatorname{Des}(\mathcal A_N,3)$ is its class
of designed families with constant $3$ (Definition~\ref{asm-designed}).

\begin{definition}\label{pd-hecke-at-T}
Let $D$ be a definite-side datum with ranks $k,m,t,n$, $f=2n$ and $N(j)=k+j(t+1)$, coefficient field
$L$ and the lattices of forms $\Lambda_{u,j}\subset V_{u,j}$ of Definition~\ref{pd-period-data}, with
their injections $\operatorname{emb}_j$ into the functions on $\mathrm{cls}_j\times Y_j$. Let $u$ be a
word, let $\mathbf F_v$ ($v\in V_u$) be the $f$ finite fields with $T_u$ elements of
Definition~\ref{th-definite-datum}(3), and let $\mathcal T_{\mathrm{top}}\subset\mathcal T$ be the
\emph{top embeddings}, those $\tau$ whose element of $G_\Phi$ is $\tau_1(t)$ for a top class $t$.
The lattices of forms of $u$ come with the following \emph{Hecke data at} $T_u$: for every level
$i\leq m$, every $v\in V_u$ and every class $c\in\mathrm{cls}_i$, an $\mathbf F_v$-vector space
$\ell_{i,v}(c)$ of dimension $N(i)$ with a nonzero linear form $f_{i,v,c}$, and for every
$x\in\ell_{i,v}(c)$ a class $c_x\in\mathrm{cls}_i$ and a linear automorphism $\rho_x$ of the
functions on $Y_i$; for every finite place $w$ of $L$ a bijection $v\mapsto\tau_w(v)$ from $V_u$ to
$\mathcal T$; and a set of places $w$ over $T_u$, those that \emph{see} $\pi_u$.
\begin{enumerate}
\item The \emph{Hecke operator} at $v\in V_u$ on functions $\varphi$ on $\mathrm{cls}_i\times Y_i$ is
  \[
   (U_v\varphi)(c)=\sum_{x\in\ell_{i,v}(c),\ f_{i,v,c}(x)=1}\rho_x\big(\varphi(c_x)\big).
  \]
\item For a place $w$ over $T_u$ let $V_{\mathrm{top}}(w)$ be the set of $v\in V_u$ with
  $\tau_w(v)\in\mathcal T_{\mathrm{top}}$ if $w$ sees $\pi_u$, and the empty set otherwise. Put
  $U_w=\prod_{v\in V_{\mathrm{top}}(w)}U_v$, the composite in a fixed order, and
  \[
   e_T(u,w)=m\sum_{v\in V_{\mathrm{top}}(w)}b_{\mathrm{top}}(\tau_w(v)),\qquad
   b_{\mathrm{top}}(\tau)=-\lambda_{u,0}(\tau)_{N(0)},
  \]
  minus the last entry of the dominant weight $\lambda_{u,0}(\tau)$ of
  Definition~\ref{th-definite-datum}(7).
\item The \emph{local lattice at} $T_u$ is the $\mathcal O_{L,w}$-module $N_\lambda(u,w)$ of the
  $y\in V_{u,0}$ with $\operatorname{emb}_0(y)=T_u^{\,e_T(u,w)}\,U_w^m\big(\operatorname{emb}_0(x)\big)$
  for some $x$ in the $w$-local span of $\Lambda_{u,0}$ with $\operatorname{emb}_0(x)\in
  M_{u,\mathrm{ram}}$.
\item The \emph{exponent at} $T_u$ is $E_T(u,w)=0$ if $w$ does not see $\pi_u$, and otherwise
  \[
   E_T(u,w)=2\sum_{\tau\in\mathcal T}W_\tau(u)+f\,N(m)+\tfrac14|\mathcal T_{\mathrm{top}}|(t+1)(m+1),
  \]
  where $W_\tau(u)$ is the largest of the numbers $\big|\lambda_{u,i}(\tau)_a-i(t+1)/2\big|$ over
  $0\leq i\leq m$ and $1\leq a\leq N(i)$, leaving out $a=1$ when $\tau\in\mathcal T_{\mathrm{top}}$.
\end{enumerate}
\end{definition}
\leanref{Formalism/LocalRingCardinality/DataForms.lean}{Formalism.BasicNotions.MemberForms}
\leanref{Formalism/LocalRingCardinality/DataLatDefs.lean}{Formalism.BasicNotions.LT.Uop}
\leanref{Formalism/LocalRingCardinality/DataLatDefs.lean}{Formalism.BasicNotions.LT.VVT}
\leanref{Formalism/LocalRingCardinality/DataLatDefs.lean}{Formalism.BasicNotions.LT.Uprod}
\leanref{Formalism/LocalRingCardinality/DataLatDefs.lean}{Formalism.BasicNotions.LT.expT}
\leanref{Formalism/LocalRingCardinality/DataLatDefs.lean}{Formalism.BasicNotions.LT.Nlam}
\leanref{Formalism/LocalRingCardinality/DataErr.lean}{Formalism.BasicNotions.EatT}
\leanref{Formalism/LocalRingCardinality/DataErr.lean}{Formalism.BasicNotions.Wtau}

\begin{definition}\label{pd-orbit-places}
Let $D$ be a definite-side datum with ranks $k,m,t,n$, $N(j)=k+j(t+1)$, Galois datum with group
$A$, top classes $T=T_{\mathrm{top}}$, parameters $\lambda_u(\tau)$ with entries $u_i(\tau)$ and
$e_b(\tau)$ (Definition~\ref{th-definite-datum}(2)), and the elements $h_p\in A$ of
Definition~\ref{pd-gain-data}.
\begin{enumerate}
\item The \emph{places} of $D$ are the pairs $z=(\epsilon,t)\in\{0,1\}\times A$. The place $z$
  corresponds to the embedding $\tau_z\in\mathcal T$ whose element of $G_\Phi$ is $\tau_1(t)$ if
  $\epsilon=0$ and $\tau_2(t)$ if $\epsilon=1$ (Definition~\ref{hf-frame}); it is a \emph{top place}
  if $\epsilon=0$ and $t\in T$. The \emph{weight} of the $j$-th lift at $z$ is
  $w_j(z)=e_j(\tau_z)$.
\item For a word $u$, a lift $1\leq j\leq m$, $\tau\in\mathcal T$ and $0\leq r<N(j-1)$, the
  \emph{row exponent} is
  \[
   n_{u,j,r}(\tau)=\pi_{u,r}(\tau)+e_j(\tau)-\tfrac12\big(N(j-1)+N(j)-1-2r\big),
  \]
  where $\pi_{u,r}(\tau)=u_{r+1}(\tau)$ for $r<k$, and
  $\pi_{u,r}(\tau)=-e_b(\tau)+\frac12(t+2-2s)$ for $r=k+(b-1)(t+1)+s-1$ with $1\leq s\leq t+1$.
  The row exponents are nonnegative integers, non-increasing in $r$. For $0\leq i<m$ and a place
  $z$, $\nu_{u,i}(z)$ is the Young diagram whose $r$-th row has length $n_{u,i+1,r}(\tau_z)$,
  $0\leq r<N(i)$; we write $|\nu|$ for the number of its cells and $\nu_r$ for the length of its
  $r$-th row, and the \emph{cap} of $\nu_{u,i}(z)$ is $\nu_1$ if $z$ is a top place and $\nu_0$
  otherwise.
\item For a prime $p$, the position $e\geq0$ steps along the orbit of $z=(\epsilon,t)$ is
  $\operatorname{pos}_p(z,e)=(\epsilon',t+eh_p)$, where $\epsilon'=\epsilon+1$ modulo $2$ if $p$ is a
  signal prime and $e$ is odd, and $\epsilon'=\epsilon$ otherwise. The \emph{orbit length}
  $\ell_p$ is twice the order of $2h_p$ if $p$ is a signal prime and the order of $h_p$ otherwise;
  $\operatorname{pos}_p(z,\ell_p)=z$, so $\operatorname{pos}_p(z,e)$ is defined for
  $e\in\mathbf Z/\ell_p\mathbf Z$. A set $R$ of places is a set of \emph{orbit representatives} at
  $p$ if every place is $\operatorname{pos}_p(r,e)$ for exactly one $r\in R$ and some $e$.
\end{enumerate}
\end{definition}
\leanref{Formalism/LocalRingCardinality/DataForced.lean}{Formalism.BasicNotions.tauPl}
\leanref{Formalism/LocalRingCardinality/DataForced.lean}{Formalism.BasicNotions.wPl}
\leanref{Formalism/LocalRingCardinality/DataDIVDefs.lean}{Formalism.BasicNotions.DV.nuRow}
\leanref{Formalism/LocalRingCardinality/DataKerI.lean}{Formalism.BasicNotions.nuD}
\leanref{Formalism/LocalRingCardinality/DataForced.lean}{Formalism.BasicNotions.posPl}
\leanref{Formalism/SharpRestriction/GrowthBounds/IsResidueClassZero.lean}{Formalism.MainNotions.G.nP}
\leanref{Formalism/SharpRestriction/Divisibility/OrbitPlaces.lean}{Formalism.MainNotions.OrbitPlaces.IsReps}

\begin{definition}\label{pd-minor-valuation}
Let $l$ be a prime, $\Lambda$ a finite set, $N_0,N_1\geq1$, and $\nu=(\nu_s)_{s\in\Lambda}$ Young
diagrams with at most $N_0$ rows. Consider variables $X_{s,a,b}$ and $Y_{s,r,a}$ for $s\in\Lambda$,
$0\leq a<N_0$, $0\leq b<N_1$ and $r\geq0$.
\begin{enumerate}
\item A \emph{column choice} into $N$ assigns to each $s\in\Lambda$ and each column $c$ of
  $\nu_s$, of length $|c|$, a strictly increasing map $I_{s,c}$ from $\{0,\dots,|c|-1\}$ to
  $\{0,\dots,N-1\}$.
\item For a column choice $J$ into $N_0$, $\nabla_J=\prod_s\prod_c\det\big(Y_{s,r,J_{s,c}(q)}\big)_{0\leq
  r,q<|c|}$. For a column choice $I$ into $N_1$, $G_I$ is the polynomial with integer coefficients
  \[
   G_I=\prod_{s\in\Lambda}\prod_c\det\Big(\sum_{a<N_0}Y_{s,r,a}\,X_{s,a,I_{s,c}(q)}\Big)_{0\leq
   r,q<|c|}.
  \]
\item For a power series $\Theta=\sum_\kappa\theta_\kappa X^\kappa$ over $\mathbf C_l$ and a
  polynomial $P=\sum_{\mu,\kappa}P_{\mu,\kappa}Y^\mu X^\kappa$, put
  $\langle P,\Theta\rangle_X=\sum_{\mu,\kappa}P_{\mu,\kappa}\,\kappa!\,\theta_\kappa\,Y^\mu$, the jet
  pairing (Definition~\ref{pd-integral-jets}) in the variables $X$.
\item For $B\in\mathbf R$, $\Theta$ has \emph{$\nu$-valuation at least $B$} if for every column
  choice $I$ into $N_1$ every coefficient of $\langle G_I,\Theta\rangle_X$ has absolute value at
  most $l^{-B}$.
\end{enumerate}
\end{definition}
\leanref{Formalism/SharpRestriction/Divisibility/ProductZerothFirstParts.lean}{Formalism.MainNotions.product_zeroth_first_parts.ColChoice}
\leanref{Formalism/SharpRestriction/Divisibility/ProductZerothFirstParts.lean}{Formalism.MainNotions.product_zeroth_first_parts.Gxz}
\leanref{Formalism/SharpRestriction/Divisibility/ProductZerothFirstParts.lean}{Formalism.MainNotions.product_zeroth_first_parts.applyAt0}
\leanref{Formalism/SharpRestriction/Divisibility/ProductZerothFirstParts.lean}{Formalism.MainNotions.product_zeroth_first_parts.KvalGe}
\leanref{Formalism/SharpRestriction/Divisibility/ProductZerothFirstParts.lean}{Formalism.MainNotions.product_zeroth_first_parts.nablaX}

\begin{definition}\label{pd-kernel-entries}
Let $D$ be a definite-side datum with divisibility data, twisted forms and the notation of
Definitions~\ref{pd-period-data} and~\ref{pd-orbit-places}. A place $w$ of $L$ is \emph{free} for a
word $u$ and a tuple $\mu$ if $p_w\neq T_u$, $p_w\notin\{g_0,g_1\}$ and $p_w\notin Q(\mu)$.
\emph{Kernel entries} for $D$ consist of the following.
\begin{enumerate}
\item For every finite place $w$ of $L$, an embedding $\iota_w\colon L\to\mathbf C_{p_w}$ that
  factors through an embedding of $\overline{\mathbf Q}$ compatible with $\iota$, such that
  $x\in\mathcal O_{L,w}$ if and only if $|\iota_w(x)|\leq1$.
\item For every word $u$, tuple $\mu$, place $w$, level $i<m$ and classes $c\in\mathrm{cls}_i$,
  $c'\in\mathrm{cls}_{i+1}$ of $C_{u,\mu}$: a \emph{kernel series} $\Theta=\Theta_{u,\mu,w,i}(c,c')$
  over $\mathbf C_{p_w}$ in the variables $X_{z,a,b}$, $z$ a place, $a<N(i)$, $b<N(i+1)$, and a
  nonzero constant $\eta_{u,\mu,w,i}(c,c')\in\mathbf C_{p_w}$.
\item For every $u,\mu$, level $i$ and class $c\in\mathrm{cls}_i$ of $C_{u,\mu}$, an $L$-vector
  space $Q_{u,\mu,i}(c)$ of \emph{rational vectors} with an injective $\iota$-semilinear map
  $\operatorname{ev}$ into the functions on $Y_i$; and, for every place $w$, $\iota_w$-semilinear
  \emph{source coordinates} $v\mapsto S_w(v)$, from $Q_{u,\mu,i}(c)$ to the linear forms on the
  polynomials over $\mathbf C_{p_w}$ in the variables $Y_{z,r,a}$, $a<N(i)$, and \emph{target
  coordinates} $v'\mapsto(T_w(v')(I))_I$, from $Q_{u,\mu,i+1}(c')$ to the families indexed by the
  column choices $I$ into $N(i+1)$ for the diagrams $\nu_{u,i}$.
\end{enumerate}
Write $k_i(c,c')$ for the map
$\varphi\mapsto\rho_i(c)\rho'_i(c')\prod_\tau\delta_{i+1}^{-s'_i(\tau)}\,\vartheta_i(c,c')\big(\prod_\tau
\delta_i^{s_i(\tau)}\varphi\big)$ of Definition~\ref{th-iterated-lift}, so that
$(K_i\varphi)_{c'}=\sum_c|\Gamma_c|^{-1}k_i(c,c')(\varphi_c)$. The kernel entries are subject to the
following conditions, at every place $w$ free for $(u,\mu)$ and every level $i<m$, where the
lattices of forms are those of $C_u$ for $\mu=\mu_0$ and the twisted ones otherwise.
\begin{enumerate}
\item[(a)] A function $\varphi$ on $\mathrm{cls}_i\times Y_i$ is the image of a vector of the
  $w$-local span of the lattice of forms at level $i$ if and only if $\varphi\in F_{u,i}$ and every
  $\varphi_c$ is $\operatorname{ev}(v)$ for some $v\in Q_{u,\mu,i}(c)$ with $|S_w(v)(\nabla_J)|\leq1$
  for every column choice $J$ into $N(i)$ (Definition~\ref{pd-minor-valuation}).
\item[(b)] A function $\varphi'$ on $\mathrm{cls}_{i+1}\times Y_{i+1}$ is the image of a vector of
  the $w$-local span of the lattice of forms at level $i+1$ if and only if $\varphi'\in F_{u,i+1}$
  and every $\varphi'_{c'}$ is $\operatorname{ev}(v')$ for some $v'\in Q_{u,\mu,i+1}(c')$ with
  $|T_w(v')(I)|\leq1$ for every $I$.
\item[(c)] Every $|T_w(v')(I)|$ is $|\iota_w(b)|$ for some $b\in L$.
\item[(d)] If $\operatorname{ev}(v')=k_i(c,c')(\operatorname{ev}(v))$, then for every column
  choice $I$
  \[
   T_w(v')(I)=\eta_{u,\mu,w,i}(c,c')\,S_w(v)\big(\langle G_I,\Theta_{u,\mu,w,i}(c,c')\rangle_X\big).
  \]
\end{enumerate}
\end{definition}
\leanref{Formalism/LocalRingCardinality/DataKerI.lean}{Formalism.BasicNotions.KerEnt}
\leanref{Formalism/SharpRestriction/Divisibility/OrbitPlaces.lean}{Formalism.MainNotions.OrbitPlaces.LevelFree}

\begin{definition}\label{pd-orbit-expansions}
Let $l\geq5$ be a prime. Let $W\subset\mathbf C_l$ be the closure of the field generated over
$\mathbf Q_l$ by the roots of unity of order prime to $l$, and $W^\circ$ its elements of absolute
value at most $1$. Let $n\geq1$, let $\Lambda$ be a finite set with a map
$\operatorname{pos}\colon\Lambda\times\mathbf Z/n\mathbf Z\to\Lambda$ such that
$\operatorname{pos}(z,0)=z$ and $\operatorname{pos}(\operatorname{pos}(z,e),e')=\operatorname{pos}(z,e+e')$,
and let $\beta$ be a finite set of labels.
\begin{enumerate}
\item The \emph{orbit logarithm} is the family of series
  \[
   \lambda_{z,b}(X)=\sum_{e\geq0}l^{-e}\,X_{\operatorname{pos}(z,-2e),b}^{\,l^{2e}}\qquad(z\in\Lambda,\
   b\in\beta).
  \]
  For $\Lambda=\mathbf Z/n\mathbf Z$ with $\operatorname{pos}(s,e)=s+e$ we use the same name.
\item For even $n$, a series $\beta_0\in\mathbf C_l[[Y_s:s\in\mathbf Z/n\mathbf Z]]$ satisfies:
  \begin{enumerate}
  \item[(B1)] if its coefficients lie in $W$;
  \item[(B2)] if $\beta_0(\tau_u(Y))=\beta_0(Y)$ for every $u$ with $u^{l^n-1}=1$, where
    $\tau_u(Y)_s=u^{l^s}Y_s$ for even $s$ and $u^{-l^s}Y_s$ for odd $s$, $0\leq s<n$;
  \item[(B3)] if every monomial with nonzero coefficient has positive degree in the even-indexed
    variables and in the odd-indexed ones, and is not of degree $(1,1)$ in them;
  \item[(B4)] if the two series $\exp(\delta_{\mathrm{ev}}\beta_0)$ and $\exp(\delta_{\mathrm{od}}\beta_0)$
    in the variables $Y_s,Z_s$, with the orbit logarithm substituted for each family of variables,
    have coefficients in $W^\circ$ and constant term $1$; here
    $\delta_{\mathrm{ev}}\beta_0=\beta_0(Y+Z_{\mathrm{ev}})-\beta_0(Y)-\beta_0(Z_{\mathrm{ev}}+Y_{\mathrm{od}})$,
    where $Y+Z_{\mathrm{ev}}$ adds $Z_s$ to $Y_s$ for even $s$ only and $Z_{\mathrm{ev}}+Y_{\mathrm{od}}$
    has $Z_s$ in the even and $Y_s$ in the odd positions, and $\delta_{\mathrm{od}}$ is the same with
    the parities exchanged.
  \end{enumerate}
\item Let $\Theta$ be a power series over $\mathbf C_l$ in variables $X_{z,a,b}$, $z\in\Lambda$,
  $a<N_0$, $b<N_1$; let $R\subset\Lambda$ be a set of orbit representatives, $\varepsilon\in\mathbf
  C_l$ with $|\varepsilon|=1$, signs $\sigma_r=\pm1$ and series $\beta_r$ in variables indexed by
  $\mathbf Z/n\mathbf Z$ ($r\in R$). Then $\Theta$ has the \emph{orbit expansion} with these data if
  there is an integral series $U$ with
  \[
   \Theta\big(\varepsilon\lambda(X)\big)=U(X)\cdot\exp\Big(\sum_{r\in R}\sigma_r\sum_{a,b}
   \beta_r\big((Y_{\operatorname{pos}(r,s),a,b})_{s\in\mathbf Z/n\mathbf Z}\big)\Big)\Big|_{Y=\lambda(X)}.
  \]
\item For $n=1$, $\Theta$ has the \emph{expansion at a split orbit} if there are
  $\varepsilon_0$ and $\varepsilon_z$ ($z\in\Lambda$) of absolute value $1$, signs $\sigma_z$, series
  $\beta_z\in\mathbf C_l[[Y_0,Y_1]]$ satisfying (B1) to (B4) for $n=2$, and an integral series $U$
  with
  \[
   \Theta\big(\varepsilon_0\lambda(X)\big)=U(X)\cdot\exp\Big(\sum_{z\in\Lambda}\frac{\sigma_z}2\sum_{a,b}
   \beta_z\big(Y_{z,a,b},\varepsilon_zY_{z,a,b}\big)\Big)\Big|_{Y=\lambda(X)}.
  \]
\item For odd $n$, $\Theta$ has the \emph{theta product form} if there are a subfield $K_0$ of
  $\mathbf C_l$ on which every nonzero absolute value is a power of $|l|$, a primitive
  $(l^{2n}-1)$-th root of unity $\zeta$, a set $R$ of orbit representatives, $\varepsilon$ with
  $|\varepsilon|=1$, signs $\sigma_r$ and series $\beta_r$ in variables indexed by
  $\mathbf Z/2n\mathbf Z$ that satisfy (B1) to (B4) for $2n$, with (B1) replaced by coefficients in
  $K_0$ and (B2) by invariance under $\tau_\zeta$, and an integral series $U$ with
  \[
   \Theta\big(\varepsilon\lambda(X)\big)=U(X)\cdot\exp\Big(\sum_{r\in R}\frac{\sigma_r}2\sum_{a,b}
   \beta_r\big((Y_{\operatorname{pos}(r,s),a,b})_{s\in\mathbf Z/2n\mathbf Z}\big)\Big)\Big|_{Y=\lambda(X)},
  \]
  where each orbit is traversed twice.
\end{enumerate}
\end{definition}
\leanref{Formalism/LocalRingCardinality/DataJet.lean}{Formalism.BasicNotions.jetLogG}
\leanref{Formalism/SharpRestriction/Divisibility/OrbitPlaces.lean}{Formalism.MainNotions.OrbitPlaces.logPl}
\leanref{Formalism/SharpRestriction/Divisibility/ProductZerothFirstParts.lean}{Formalism.MainNotions.product_zeroth_first_parts.B1}
\leanref{Formalism/SharpRestriction/Divisibility/ProductZerothFirstParts.lean}{Formalism.MainNotions.product_zeroth_first_parts.B2}
\leanref{Formalism/SharpRestriction/Divisibility/ProductZerothFirstParts.lean}{Formalism.MainNotions.product_zeroth_first_parts.B3}
\leanref{Formalism/SharpRestriction/Divisibility/ProductZerothFirstParts.lean}{Formalism.MainNotions.product_zeroth_first_parts.B4}
\leanref{Formalism/SharpRestriction/Divisibility/OrbitPlaces.lean}{Formalism.MainNotions.OrbitPlaces.ThetaExpansionAllPlaces}
\leanref{Formalism/SharpRestriction/Divisibility/OrbitPlaces.lean}{Formalism.MainNotions.OrbitPlaces.ThetaExpansionRealPlaces}
\leanref{Development/AutomorphicArithmetic/OrbitLabellings/ColDefs.lean}{UpperTheory.G2.HasThetaProductForm}

\begin{definition}\label{pd-local-model}
Let $p\geq5$ be a prime.
\begin{enumerate}
\item A \emph{Hida setting} of degree $n'$ at $p$, with $n'\geq2$ even, consists of a complete
  discrete valuation ring $W$ with uniformizer $p$, residue field of $p^{n'}$ elements and fraction
  field $Q$; an automorphism $\sigma$ of $Q$ with $\sigma^{n'}=1$ and $\sigma(x)\equiv x^p\pmod{pW}$
  for $x\in W$; a complete discrete valuation ring $\mathcal R$ with fraction field $\mathcal K$ and an
  element $\varpi_0\in\mathcal R$ with $\varpi_0^2=p$; and an embedding $\iota\colon Q\to\mathcal K$
  with $\iota(W)\subset\mathcal R$ that makes $\mathcal R$ finite over $W$.
\item A \emph{local model} at $p$ consists of $\mathcal R\subset\mathcal K$ and $\varpi_0$ as in (1);
  a finite set $V$ of \emph{orbits} and, for each $v\in V$, a Hida setting
  $(W_v,Q_v,\sigma_v,\iota_v)$ of degree $n'_v$ with these $\mathcal R$, $\mathcal K$ and
  $\varpi_0$; integers $k\geq1$ and $r\geq0$; dominant weights
  $\lambda_0(v,s)\in\mathbf Z^{k}$, $\lambda_0(v,s)_1\geq\cdots\geq\lambda_0(v,s)_k$, for the
  \emph{positions} $(v,s)$, $s\in\mathbf Z/n'_v\mathbf Z$; and a finite set of \emph{components}
  $c$, each with a \emph{level}: a group $\Gamma_c$ with homomorphisms $\Gamma_c\to GL_k(Q_v)$ for
  $v\in V$ and a character $\chi_c\colon\Gamma_c\to\mathcal R^\times$.
\end{enumerate}
In a local model, write $s$ also for its representative in $\{0,\dots,n'_v-1\}$, and put
$\lambda(v,s)=\lambda_0(v,s)$ if $s$ is even and $\lambda(v,s)=(-\lambda_0(v,s)_k,\dots,
-\lambda_0(v,s)_1)$ if $s$ is odd. A \emph{lattice tuple} is a family $L=(L_v)_{v\in V}$ of
$W_v$-lattices $L_v\subset Q_v^{k}$. Let $\mathcal X=\mathcal K[x_{v,s,a,b}]$, with $1\leq a,b\leq k$,
be the polynomial ring on one $k\times k$ matrix $x_{v,s}$ of variables for each position. For a
lattice tuple $L$ let $H^{\Delta}(L)\subset\mathcal X$ be the $\mathcal R$-span of the products
$\prod_{v,s}f_{v,s}(x_{v,s})$ with $f_{v,s}\in\Delta_{\lambda(v,s)}(M^\vee_{v,s})$, where $M_{v,s}$ is
the $\mathcal R$-span of $\iota_v\sigma_v^{s}(L_v)$,
$M^\vee_{v,s}=\{\phi\in\mathcal K^k:\phi\cdot y\in\mathcal R\text{ for all }y\in M_{v,s}\}$ is its
dual, and $\Delta_\lambda(M)$ is formed over $\mathcal R$ as in
Definition~\ref{pd-weight-lattices}(5); $H^{\nabla}(L)$ is defined in the same way from the lattices
$\nabla_{\lambda(v,s)}(M^\vee_{v,s})$. Equivalently, $f\in\Delta_\lambda(M^\vee)$ if and only if
$\det(g)^{\lambda_k}f(gx)$ lies in $\Delta_{\mathcal R}$ of the natural weight of $\lambda$ for every
$g$ whose columns form an $\mathcal R$-basis of $M$, and likewise for $\nabla$; so a matrix $g$ moves
these lattices by $f\mapsto\det(g)^{-\lambda_k}f(g^{-1}x)$, as in (3) below, and not by
$\rho_\lambda(g)$ of Definition~\ref{pd-weight-lattices}(4).
\begin{enumerate}
\item[(3)] An element $\gamma$ of a level acts on lattice tuples through the homomorphisms to
  $GL_k(Q_v)$, and on $\mathcal X$ by multiplying by
  $\prod_{v,s}\det(\iota_v\sigma_v^{s}\gamma_v)^{-\lambda(v,s)_k}$ and substituting
  $\iota_v\sigma_v^s(\gamma_v^{-1})\,x_{v,s}$ for $x_{v,s}$. The \emph{integral module} of the level
  $c$ is the $\mathcal R$-module $M_c$ of the functions $F$ from all tuples of $W_v$-submodules to
  $\mathcal X$ such that $F(L)\in H^{\Delta}(L)$ for every lattice tuple $L$, $F(L)=0$ for every
  other tuple, and $F(\gamma L)=\chi_c(\gamma)\,\gamma\cdot F(L)$ for all $\gamma\in\Gamma_c$. The
  \emph{function space} $\mathrm{Fn}$ of the model consists of the functions from the pairs of a
  component and a tuple to $\mathcal X$, and its integral module is $M=\prod_cM_c$.
\item[(4)] An \emph{edge} at $v$ is a pair of $W_v$-lattices $L''\subset L_v$ with
  $pL_v\subset L''$ and $|L_v/L''|=p^{n'_vr}$. The \emph{Hecke operator} at $v$ is
  \[
   (y_vF)(c,L)=p^{A_v}\sum_{L''}F\big(c,L\text{ with }L_v\text{ replaced by }L''\big),\qquad
   A_v=\sum_{s}\sum_{i\leq r}\lambda(v,s)_i,
  \]
  the sum over the edges $L''\subset L_v$ at $v$.
\item[(5)] \emph{Fills} of the model consist of sets $T_v\subset\mathbf Z/n'_v\mathbf Z$ of \emph{top
  positions}, integers $a(v,s)$, $b(v,s)$ and $c\geq0$ such that $\lambda_0(v,s)_1=a(v,s)$,
  $\lambda_0(v,s)_k=-b(v,s)$ and $|\lambda_0(v,s)_i|\leq c$ for $1<i<k$ if $s\in T_v$, and
  $|\lambda_0(v,s)_i|\leq c$ for all $i$ if $s\notin T_v$; integers $t,m_*\geq0$; and integers
  $w(v,j,s)$ for $j\geq1$ with $\max(c,-\lambda_0(v,s)_k)\leq w(v,1,s)\leq w(v,2,s)\leq\cdots$. The
  \emph{weights at level} $j$ are $\lambda_j(v,s)=\big(\lambda_0(v,s),-w(v,1,s)^{(t+1)},\dots,
  -w(v,j,s)^{(t+1)}\big)$, of length $N(j)=k+j(t+1)$, where $y^{(t+1)}$ is $y$ repeated $t+1$ times;
  they are required to be dominant. The model at level $j$ is the local model with the same Hida
  settings and the weights $\lambda_j$.
\item[(6)] \emph{Levels and a chain} over fills consist of, for each $j\geq1$, a finite set of
  components with levels for the model at level $j$, the components at level $0$ being those of the
  model; and, writing $\mathrm{Fm}_j$ for the function space at level $j$ and $\mathrm{Int}_j$ for
  its integral module, $\mathcal K$-linear maps $K_j\colon\mathrm{Fm}_j\to\mathrm{Fm}_{j+1}$. Put
  $K^{(j)}=K_{j-1}\circ\cdots\circ K_0$.
\item[(7)] Let $e_{\mathcal R}$ be the valuation of $p$ in $\mathcal R$. For $F\in\mathrm{Fm}_j$ and
  $g\in\mathbf R$, $F$ \emph{has valuation at least} $g$ if $F\in\varpi^{\lceil ge_{\mathcal R}\rceil}
  \mathrm{Int}_j$ for a uniformizer $\varpi$ of $\mathcal R$.
\end{enumerate}
\end{definition}
\leanref{Formalism/SharpRestriction/GrowthBounds/UnramifiedQuadraticSetting.lean}{Formalism.MainNotions.G.HidaSetting}
\leanref{Formalism/SharpRestriction/GrowthBounds/DataG.lean}{Formalism.MainNotions.G.Places}
\leanref{Formalism/SharpRestriction/GrowthBounds/DataG.lean}{Formalism.MainNotions.G.Level}
\leanref{Formalism/SharpRestriction/GrowthBounds/DataG.lean}{Formalism.MainNotions.G.GAtP}
\leanref{Formalism/SharpRestriction/GrowthBounds/DataGChain.lean}{Formalism.MainNotions.G.Fills}
\leanref{Formalism/SharpRestriction/GrowthBounds/DataGChain.lean}{Formalism.MainNotions.G.Levels}
\leanref{Formalism/SharpRestriction/GrowthBounds/DataGChain.lean}{Formalism.MainNotions.G.Chain}
\leanref{Formalism/SharpRestriction/GrowthBounds/IsResidueClassZero.lean}{Formalism.MainNotions.G.vGeLat}

\begin{definition}\label{pd-lattice-instance}
Let $D$ be a definite-side datum with ranks $k,m,t,n$, $N(j)=k+j(t+1)$, Galois datum with group $A$,
top classes $T$, size data and gain data (Definitions~\ref{pd-size-data} and~\ref{pd-gain-data}),
with the elements $h_p$, the representatives at $p$ and the places $z=(\epsilon,t)$ of
Definition~\ref{pd-orbit-places}. Put $c_G=4S_c$, $\kappa_j=\frac12(k+(2j-1)(t+1))$ and, for the
offsets $o_j(z)$ of the gain data,
\[
 \mathrm{fill}_j(z)=b_t\,[z=(0,t)\text{ is a top place}]+o_j(z)-\kappa_j .
\]
A finite place $w$ of $L$ is \emph{admissible} for a word $u$ if $p=p_w$ is a prime of inert class,
$p\geq5$, the $A$-component $h_p$ of the Frobenius class of $p$ has even order, $p$ is not a level
prime, and $p\nmid T_u$.
\begin{enumerate}
\item The datum carries, for every word $u$ and every place $w$ admissible for $u$, a local model
  $G(u,w)$ at $p_w$ (Definition~\ref{pd-local-model}). A \emph{tie} of these local models
  identifies, for all such $(u,w)$, the positions of $G(u,w)$ with the set $\mathcal T$ of
  embeddings, so that the rank of $G(u,w)$ is $k$; $n'_v$ is the order of the $A$-component of an
  arithmetic Frobenius at a prime of $E$ over $p_w$; the element of $G_\Phi$ of the position
  $(v,s+1)$ is that of $(v,s)$ times this Frobenius and complex conjugation; and $\lambda_0(v,s)$ is
  the weight $\lambda_{u,0}(\tau)$ of the corresponding embedding $\tau$
  (Definition~\ref{th-definite-datum}(7)).
\item A \emph{lattice instance} attaches to every word $u$ and place $w$ admissible for $u$ a map
  $t\mapsto v(t)$ from $A$ to the orbits of $G(u,w)$, fills of $G(u,w)$, levels and a chain over
  them, and an $\mathcal R$-submodule $V_{\mathrm{ram}}$ of its function space.
\item A \emph{comparison} attaches to every such $(u,w)$ an embedding $\epsilon_w\colon L\to\mathcal K$
  with $v_w(y)\geq0$ if and only if $\epsilon_w(y)\in\mathcal R$, and an $\epsilon_w$-semilinear
  additive map $\operatorname{cmp}_{u,w}$ from $V_{u,0}$ to the function space of $G(u,w)$ that
  carries a basis of $\Lambda_{u,0}$ to vectors linearly independent over $\mathcal K$, such that
  the $\mathcal R$-span of $\operatorname{cmp}_{u,w}(\Lambda_{u,0})$ is the integral module $M$ of
  $G(u,w)$; here $\Lambda_{u,0}\subset V_{u,0}$ is the lattice of forms in degree $0$
  (Definition~\ref{pd-period-data}(4)).
\item A lattice instance is \emph{tied to the gain data} by a tie of the local models if, at every
  word $u$ and every place $w$ admissible for $u$, with $p=p_w$: $m_*=m$, the integer $t$ of the
  fills is the rank $t$, $c\leq c_G$, and $r=r(\varepsilon)$ (Definition~\ref{pd-gain-link}) if
  $p<p_0$ and $r=1$ otherwise; $t\mapsto v(t)$ is injective on the representatives at $p$; every
  position whose element of $G_\Phi$ is $\tau_2(t)$ for some $t\in A$ is not a top position and has
  $w(v,j,s)=\mathrm{fill}_j((1,t))$ for every lift $j$; and for every representative $t_0$ at $p$,
  $n'_{v(t_0)}$ is the order of $h_p$, and for every $s$ the position $(v(t_0),s)$ has element
  $\tau_1(t_0+sh_p)$, it is a top position if and only if $t_0+sh_p\in T$, in which case
  $a(v(t_0),s)=a_{t_0+sh_p}$ and $b(v(t_0),s)=b_{t_0+sh_p}$, and
  $w(v(t_0),j,s)=\mathrm{fill}_j((0,t_0+sh_p))$ for every lift $j$.
\end{enumerate}
\end{definition}
\leanref{Formalism/LocalRingCardinality/RhoDetPrimeEqMuTildeMulChiPow.lean}{Formalism.BasicNotions.GIdx}
\leanref{Formalism/LocalRingCardinality/DataGatTies.lean}{Formalism.BasicNotions.GatTies}
\leanref{Formalism/SharpRestriction/GrowthBounds/IsResidueClassZero.lean}{Formalism.MainNotions.G.LatInst}
\leanref{Formalism/LocalRingCardinality/DataLatDefs.lean}{Formalism.BasicNotions.LatCmp}
\leanref{Formalism/SharpRestriction/GrowthBounds/IsResidueClassZero.lean}{Formalism.MainNotions.G.LatTied}
\leanref{Formalism/SharpRestriction/GrowthBounds/IsResidueClassZero.lean}{Formalism.MainNotions.G.fillG}

\begin{definition}\label{pd-formal-group-system}
Let $p\geq5$ be a prime, $W(\overline{\mathbf F}_p)$ the ring of Witt vectors with Frobenius
$\sigma$ and fraction field $K_0$, and $\ell(X)=\sum_{j\geq0}X^{p^{2j}}/p^{j}$. For power series
$f,g$ over $K_0$ in a set of variables, write $f\simeq g$ if $f-g$ has constant term $0$ and
$\exp\big((f-g)(\ell(X))\big)$, with $\ell(X_i)$ substituted for each variable $X_i$, is $1$ plus a
series with coefficients in $W(\overline{\mathbf F}_p)$. A series $\beta\in K_0[[Y_0,Y_1]]$ is
\emph{bi-additive} if each of its monomials has positive degree in $Y_0$ and in $Y_1$,
$\beta(X_0+X_1,X_2)\simeq\beta(X_0,X_2)+\beta(X_1,X_2)$,
$\beta(X_0,X_1+X_2)\simeq\beta(X_0,X_1)+\beta(X_0,X_2)$, and
$\beta(\alpha X_0,X_1)\simeq\beta(X_0,\sigma(\alpha)X_1)$ for every $\alpha\in W(\overline{\mathbf F}_p)$
with $\sigma^2(\alpha)=\alpha$.

A \emph{formal group system} at $p$ consists of a local model at $p$ with fills, levels and a chain
over them (Definition~\ref{pd-local-model}), with $2r\leq k$, $r\geq1$ and $m_*\geq1$; a uniformizer
$\varpi$ of $\mathcal R$; integers $d_\Gamma(j)$ such that $v_p$ of the order of the stabilizer of
every lattice tuple in every level at every level $j\leq m_*$ is at most $d_\Gamma(j)$; for each orbit
$v$ elements $\omega_1,\dots,\omega_{n'_v/2}$ and $\omega^*_1,\dots,\omega^*_{n'_v/2}$ of
$W(\overline{\mathbf F}_p)$ fixed by $\sigma^{n'_v}$ with
$\sum_{i<n'_v/2}\sigma^{2i}(\omega_r\omega^*_{r'})=\delta_{rr'}$; and \emph{kernel data}, described
below. For $i\geq0$, an \emph{edge of index $i$} at an orbit $v$ is a pair $L''\subset L$ of
$W_v$-lattices with $pL\subset L''$ and $|L/L''|=p^{n'_vi}$; the edges of
Definition~\ref{pd-local-model}(4) have index $r$. For a tuple $L$ and a $W_v$-lattice $L_1$, $L[v\mapsto
L_1]$ is $L$ with its component at $v$ replaced by $L_1$, and $T_{v,i}$ is the operator on functions
of a component and a tuple with $(T_{v,i}F)(c,L)=\sum_{L''}F(c,L[v\mapsto L''])$, the sum over the
edges $L''\subset L_v$ of index $i$. The kernel data are:
\begin{enumerate}
\item a field $\mathcal K_c$ containing $K_0$ and $\mathcal K$, with a valuation ring $\mathcal R_c$ of
  rank one that contains $W(\overline{\mathbf F}_p)$ and meets $\mathcal K$ in $\mathcal R$; we write
  $v(y)\geq g$ for $y\in\mathcal K_c$ if $y^{b}p^{-a}\in\mathcal R_c$ whenever $b>0$ and $a\leq gb$;
\item a series $\beta\in K_0[[Y_0,Y_1]]$;
\item for each $j$ a finite set $\mathrm{Rep}_j$ of pairs $(c,L)$ of a component and a lattice
  tuple at level $j$, which for $j\leq m_*$ contains, for every component $c$, exactly one
  representative $(c,L_0)$ of each orbit of lattice tuples under the level of $c$; and the stabilizer
  of every pair of a component and a lattice tuple at a level $j\leq m_*$ is finite;
\item for each $j$, each pair $(c,L)$ of a component and a tuple at level $j$ and each pair $(c',L')$
  at level $j+1$, a set of \emph{frames} $\phi$, and for each frame a \emph{kernel series}
  $\Theta_\phi$ over $\mathcal K_c$ in variables $X_e,X'_e$ indexed by $e=(v,r',a,b)$ with
  $r'\leq n'_v/2$, $a\leq N(j)$ and $b\leq N(j+1)$, and a matrix $B_\phi$ over $\mathcal K_c$;
\item for all pairs $(c,L)$ and $(c',L')$ as in (4), the \emph{kernel in orbit coordinates}
  $\Theta^{Z}(c,L;c',L')$, a series in variables
  $Z_{v,s,a,b}$, equal for every frame $\phi$ to $\Theta^Z_\phi$ followed by the linear substitution
  $B_\phi$, where $\Theta^Z_\phi$ is $\Theta_\phi$ with
  $X_{(v,r',a,b)}=\sum_{s\ \mathrm{even}}\sigma^{s}(\omega^*_{r'})Z_{v,s,a,b}$ and
  $X'_{(v,r',a,b)}=\sum_{s\ \mathrm{odd}}\sigma^{s}(\omega_{r'})Z_{v,s,a,b}$;
\item for each $j$, a $\mathcal K$-bilinear \emph{kernel operator} $\mathrm{Del}_j$ from pairs of a
  polynomial at level $j$ and a series in the $Z$-variables to polynomials at level $j+1$ over
  $\mathcal K_c$, which vanishes on every series that has no monomial whose degree in the variables
  of each position $(v,s)$ is the number of cells of the natural weight of $\lambda_{j+1}(v,s)$;
  integer \emph{lengths} $\ell_v(L)$ of tuples, with $\ell_v(L[v\mapsto L''])=\ell_v(L)-i$ for an edge
  $L''\subset L_v$ of index $i$, and $\ell_v(L[v'\mapsto L''])=\ell_v(L)$ for $v'\neq v$; and the
  \emph{kernel}
  $\ker_j(c,L;c',L')(G)=\prod_vp^{\mathrm{vm}_v(j+1)(\ell_v(L)-\ell_v(L'))}\,
  \mathrm{Del}_j\big(G,\Theta^Z(c,L;c',L')\big)$, where $\mathrm{vm}_v(j)$ is the sum of the weights
  $w(v,j,s)$ over even $s$ minus their sum over odd $s$;
\item \emph{Weil index data}: for each orbit $v$ and level $j$, a set $\mathrm{Aw}_{v,j}$, a map
  $\mathrm{aw}_{v,j}$ from the pairs $((c,L),(c',L'))$ at levels $j$ and $j+1$ to $\mathrm{Aw}_{v,j}$
  that does not change when the component at $v$ of $L$ or of $L'$ is changed, and for each
  $\alpha\in\mathrm{Aw}_{v,j}$ an additive map $\Xi_{v,j,\alpha}$ from the $\mathbf Q$-valued functions
  on the $N(j+1)\times N(j)$ matrices over $Q_v$ to series over $\mathcal K_c$ in the $Z$-variables.
  For tuples $L,L'$ at levels $j,j+1$, $\operatorname{ind}_{v,j}(L,L')$ is the function that is $1$ at
  the matrices $y$ with $yL_v\subset L'_v$ and $0$ at the others.
\end{enumerate}
For $j<m_*$, an orbit $v$ and $g\in\mathbf R$, the system has \emph{lower bound $g$ at $(j,v)$} if for
every series $U$ in the variables $X_e,X'_e$ such that $U(\ell(X))$ has coefficients in
$\mathcal R_c$, every coefficient of $\langle G^{(v)}_I,\Theta^Z_U\rangle_Z$ has $v\geq g$, for every
column choice $I$ into $N(j+1)$ (Definition~\ref{pd-minor-valuation}) for the diagrams $\nu_{v,s}$,
$s\in\mathbf Z/n'_v\mathbf Z$, given by the natural weights of $\lambda_{j+1}(v,s)$; here
$G^{(v)}_I$ is the product of the minors of Definition~\ref{pd-minor-valuation}(2) at the positions
of $v$, and $\Theta^Z_U$ is the series $\exp\big(-\sum_{e\text{ at }v}\beta(X_e,X'_e)\big)\,U$ in the
orbit coordinates of (5).

The system is \emph{admitted} if:
\begin{enumerate}
\item[(a)] (Dwork's lemma) for every prime $q\neq2$ and every series $\Psi$ over $K_0$ with constant
  term $0$, $\exp(\Psi)$ has coefficients in $W(\overline{\mathbf F}_q)$ if and only if
  $\Psi^{\sigma}(X^q)-q\Psi(X)\in qW(\overline{\mathbf F}_q)[[X]]$, where $\Psi^\sigma$ applies the
  Frobenius to the coefficients and $X^q$ is substituted for each variable \cite{Haz78};
\item[(b)] $\beta$ is bi-additive, and every kernel series, and its restriction to the variables of
  each orbit $v$, is $\exp(-\Phi)\,U$ with $U(\ell(X))$ integral over $\mathcal R_c$, $\Phi(-X)\simeq\Phi(X)$ and
  $\Phi\simeq\sum_e\beta(X_e,X'_e)$, the sum over the indices $e$ of the variables (at $v$);
\item[(c)] every coefficient of $\langle G_I,\Theta^Z_\phi\rangle_Z$ has $v\geq0$, for every frame
  $\phi$ and every column choice $I$ into $N(j+1)$ for the diagrams $\nu_{v,s}$ at all positions;
\item[(d)] $K_j$ carries the $\mathcal K$-span of $\mathrm{Int}_j$ into that of $\mathrm{Int}_{j+1}$;
\item[(e)] (expansion) for each $j<m_*$ there is $c_K\in\mathcal R_c^\times$ such that, if
  $j+1<m_*$,
  \[
   K_jF(c',L')=\sum_{(c,L)\in\mathrm{Rep}_j}|\mathrm{Stab}(c,L)|^{-1}c_K\ker_j(c,L;c',L')(F(c,L))
  \]
  for every $F$ in the $\mathcal K$-span of $\mathrm{Int}_j$, every component $c'$ at level $j+1$ and
  every lattice tuple $L'$; and, if $j+1=m_*$, there are finite sets $C$ and $Y$, integers
  $\mathrm{wt}(x)>0$ with $v_p(\mathrm{wt}(x))\leq d_\Gamma(j)$ for $x\in C$, and pairs $(c_{xy},L_{xy})$
  of a component and a lattice tuple at level $j$ for $(x,y)\in C\times Y$, such that
  \[
   K_jF(c',L')=\sum_{x\in C}\mathrm{wt}(x)^{-1}c_K\sum_{y\in Y}\ker_j(c_{xy},L_{xy};c',L')(F(c_{xy},L_{xy}))
  \]
  for all $F$, $c'$, $L'$ as before, and
  \[
   \sum_x\mathrm{wt}(x)^{-1}\sum_y\sum_{L_1}\ker_j\big(c_{xy},L_{xy}[v\mapsto L_1];c',L'\big)
   (F(c_{xy},L_{xy}))=\sum_x\mathrm{wt}(x)^{-1}\sum_y\ker_j(c_{xy},L_{xy};c',L')\big((T_{v,i}F)(c_{xy},L_{xy})\big)
  \]
  for every orbit $v$, $i\geq0$, $F$ in the $\mathcal K$-span of $\mathrm{Int}_j$ and $(c',L')$ with
  $L'$ a lattice tuple, the sum over the $L_1$ such that $(L_{xy})_v\subset L_1$ is an edge of index
  $i$;
\item[(f)] (Weil index) for each orbit $v$ and $j<m_*$ there are $u,u'\in\mathcal K$ of valuation
  $0$ such that for all pairs $(c,L)$, $(c',L')$ of a component and a lattice tuple at levels $j$
  and $j+1$,
  \[
   \Theta^Z(c,L;c',L')=u^{\ell_v(L)}u'^{\ell_v(L')}p^{-\frac{n'_v}2(N(j)\ell_v(L')-N(j+1)\ell_v(L))}\,
   \Xi_{v,j,\mathrm{aw}_{v,j}((c,L),(c',L'))}\big(\operatorname{ind}_{v,j}(L,L')\big);
  \]
\item[(g)] (equivariance) $\ker_j(c,\gamma L;c',L')(\chi_c(\gamma)\gamma\cdot G)=\ker_j(c,L;c',L')(G)$
  for $\gamma$ in the level of $c$;
\item[(h)] (adjunction) for every level $j\leq m_*$, orbit $v$, $i\geq0$ and function $\varphi$ of a
  component and two tuples invariant under the levels, the sum over $(c,L)\in\mathrm{Rep}_j$ with
  weights $|\mathrm{Stab}(c,L)|^{-1}$ of the sum of $\varphi(c,L,L[v\mapsto L_1])$ over the edges
  $L_v\subset L_1$ of index $i$ equals the same sum of the sum of $\varphi(c,L[v\mapsto L_0],L)$ over
  the edges $L_0\subset L_v$ of index $i$;
\item[(i)] (Hecke edges) let $j<m_*$, let $(c,L)$ be a component and a lattice tuple at level $j$, $c'$
  a component at level $j+1$, $L'$ and $L''$ tuples at level $j+1$ with $L'$ a lattice tuple, and
  $i_0$ a function from the orbits to $\{0,\dots,r\}$ such that for every orbit $v$ either $i_0(v)=0$
  and $L''_v=L'_v$, or $i_0(v)\geq1$ and $L''_v\subset L'_v$ is an edge of index $i_0(v)$. For a
  column choice $I$ into $N(j+1)$ for the diagrams $\nu_{v,s}$ of the natural weights of
  $\lambda_{j+1}(v,s)$, let $G^{i_0}_I$ be the product, over the orbits $v$, the positions $s$ and the
  columns $c$ of $\nu_{v,s}$, of the minors of Definition~\ref{pd-minor-valuation}(2) for the column
  $c$, each multiplied by $p^{e}$, where $e$ is the number of $q<|c|$ with $I_{v,s,c}(q)<i_0(v)$ if
  $s$ is odd, and $\min(|c|,i_0(v))$ minus that number if $s$ is even. Then there is a frame $\phi$
  for $((c,L),(c',L''))$ such that, for every series $f$ in the $Z$-variables and every $g\in\mathbf R$
  for which all coefficients of all $\langle G^{i_0}_I,f\rangle_Z$ have $v\geq g$, and for every
  $G\in H^\Delta(L)$, the polynomial
  \[
   \prod_vp^{N'_v}\cdot\prod_vp^{\mathrm{vm}_v(j+1)(\ell_v(L)-\ell_v(L''))}\,
   \mathrm{Del}_j(G,f_\phi)
  \]
  lies in the additive span of the products $y\,z$ with $y\in\mathcal K_c$, $v(y)\geq g$ and
  $z\in H^\Delta(L')$; here $f_\phi$ is $f$ followed by the linear substitution $B_\phi$, and
  $N'_v=\sum_{s\ \mathrm{even}}\sum_{a=1}^{i_0(v)}\lambda_0(v,s)_a+i_0(v)\sum_{s\ \mathrm{odd}}w(v,j+1,s)$;
\item[(j)] $y_vy_{v'}F=y_{v'}y_vF$ for all orbits $v,v'$ and every function $F$ of a component and a
  tuple, and $y_v$ maps $M$ into $M$ for every orbit $v$;
\item[(k)] the system has lower bound $0$ at every $(j,v)$ with $j<m_*$.
\end{enumerate}
\end{definition}
\leanref{Formalism/SharpRestriction/GrowthBounds/FormalGroupSystem.lean}{Formalism.MainNotions.FormalGroupSystem.Simeq}
\leanref{Formalism/SharpRestriction/GrowthBounds/FormalGroupSystem.lean}{Formalism.MainNotions.FormalGroupSystem.BeProp}
\leanref{Formalism/SharpRestriction/GrowthBounds/FormalGroupSystem.lean}{Formalism.MainNotions.FormalGroupSystem.Sys}
\leanref{Formalism/SharpRestriction/GrowthBounds/FormalGroupSystem.lean}{Formalism.MainNotions.FormalGroupSystem.Pend}
\leanref{Formalism/SharpRestriction/GrowthBounds/FormalGroupSystem.lean}{Formalism.MainNotions.FormalGroupSystem.Pend.LowAt}
\leanref{Formalism/SharpRestriction/GrowthBounds/FormalGroupSystem.lean}{Formalism.MainNotions.FormalGroupSystem.admitted}
\leanref{Formalism/SharpRestriction/GrowthBounds/FormalGroupSystem.lean}{Formalism.MainNotions.FormalGroupSystem.ArchimedeanSignatureMatching}
\leanref{Formalism/SharpRestriction/GrowthBounds/FormalGroupSystem.lean}{Formalism.MainNotions.FormalGroupSystem.ExpansionDecomposition}
\leanref{Formalism/SharpRestriction/GrowthBounds/FormalGroupSystem.lean}{Formalism.MainNotions.FormalGroupSystem.adm5}
\leanref{Formalism/SharpRestriction/GrowthBounds/FormalGroupSystem.lean}{Formalism.MainNotions.FormalGroupSystem.adm7}
\leanref{Formalism/SharpRestriction/GrowthBounds/FormalGroupSystem.lean}{Formalism.MainNotions.FormalGroupSystem.adm8}
\leanref{Formalism/SharpRestriction/GrowthBounds/FormalGroupSystem.lean}{Formalism.MainNotions.FormalGroupSystem.Sys.ind}
\leanref{Formalism/SharpRestriction/GrowthBounds/FormalGroupSystem.lean}{Formalism.MainNotions.FormalGroupSystem.Sys.PI}
\leanref{Formalism/SharpRestriction/GrowthBounds/FormalGroupSystem.lean}{Formalism.MainNotions.FormalGroupSystem.Nprime}
\leanref{Formalism/SharpRestriction/GrowthBounds/DataGChain.lean}{Formalism.MainNotions.G.Places.Tminus}

\begin{definition}\label{pd-served-selection}
Let $\mathcal C$ be a class of families with a twist selection. At every member $D$ of a family of
$\mathcal C$ at every scale beyond a threshold of the family, the selection chooses kernel entries
for $D$ (Definition~\ref{pd-kernel-entries}), with kernel series
$\Theta_{u,\mu,w,i}(c,c')$. Use the notation of Definitions~\ref{pd-size-data},
\ref{pd-gain-data}, \ref{pd-orbit-places} and~\ref{pd-lattice-instance}, with $c_G$, $\kappa_j$ and
$\mathrm{fill}_j$ as there. For a prime $p$, a lift $j$ and a representative $t_0$ at $p$, the
\emph{local gain} is
\[
 g_0(p,j,t_0)=\sum_{s}\ \sum_{e\ \mathrm{odd},\ e\ \mathrm{not\ excluded}}
 \max\Big(0,\frac{a_s-2c_G}{p^e}-\mathrm{fill}_j\big(\operatorname{pos}_p((0,s),e)\big)-c_G\Big)
 -\delta(p,t_0),
\]
where $s$ runs over the top classes with representative $t_0$, $e$ is \emph{excluded} for $s$ if,
for some top class $s'\neq s$, $e$ is the least positive integer with
$\operatorname{pos}_p((0,s),e)=(0,s')$, and $\delta(p,t_0)=K_{\delta,1}\ell_pN(m)^2\log_p(2a)+K_{\delta,2}N(p,t_0)(\log_p2a)^2$ with
constants $K_{\delta,1},K_{\delta,2}$ of $K_G$. The selection is \emph{served} along $\mathcal C$ if,
eventually along $\mathcal C$, every member at all large scales satisfies the following.
\begin{enumerate}
\item \emph{Twisted comparison}: the member carries a lattice instance with a comparison
  $\operatorname{cmp}$ (Definition~\ref{pd-lattice-instance}), and there are, for every tuple
  $\mu$, every word $u$ and every place $w$ admissible for $u$, levels and a chain over the fills of
  the lattice instance at $(u,w)$, equal to those of the lattice instance when $\mu=\mu_w$; write
  $K^{(m)}_\mu$ for the composite of the first $m$ maps of the chain of $\mu$. These \emph{local
  chains of the twists} satisfy:
  \begin{enumerate}
  \item[(a)] for every word $u$, place $w$ admissible for $u$, used tuple $\mu$ with
    $p_w\notin Q(\mu)$, $x\in V_{u,0}$ and $s\in\mathbf R$: if $K^{(m)}_\mu(\operatorname{cmp}_{u,w}x)$
    has valuation at least $s$ (Definition~\ref{pd-local-model}(7)), then the iterated lift of $x$
    along $C_{u,\mu}$ has $w$-valuation at least $s$ (Definition~\ref{pd-lift-valuation});
  \item[(b)] if the error data of the member are consistent, that is, their constants are large,
    their level data and the places that see $\pi_u$ are those of the datum, the stabilizer
    denominators are defined (Definition~\ref{pd-period-data}(2)) and every local error
    $E_{u,j}(w)$ is the one the error data define, then the following holds for every used tuple
    $\mu$, every tie of the local models under which the lattice instance with the local chains of
    $\mu$ is tied to the gain data, every word $u$ and every place $w$ admissible for $u$ with
    $p_w\notin Q(\mu)$. There are an admitted formal group system at $p_w$
    (Definition~\ref{pd-formal-group-system}) on $G(u,w)$ with the fills of the lattice instance and
    the levels and chain of $\mu$, and real numbers $F_v(j)$ for the orbits $v$ that are not of the
    form $v(t_0)$ with $t_0$ a representative at $p_w$, such that the system has lower bound
    $F_v(j+1)$ at $(j,v)$ for every such $v$ and $j<m$, and
    \[
     \sum_{j=1}^{m}\big(F_j(w)-E_{u,j}(w)\big)+v_{p_w}(D_\Gamma(u))\leq
     \sum_{j=1}^{m}\sum_vF_v(j)-\sum_{j<m}d_\Gamma(j).
    \]
  \end{enumerate}
\item \emph{Units and rationality}: at every place $w$ free for $(u,\mu)$ the constants
  $\eta_{u,\mu,w,i}(c,c')$ have absolute value $1$; for used tuples $\mu$, free $w$ and $i<m$, the
  map $k_i(c,c')$ of $C_{u,\mu}$ carries $\operatorname{ev}(Q_{u,\mu,i}(c))$ into
  $\operatorname{ev}(Q_{u,\mu,i+1}(c'))$; the lifts $K_i$ carry $F_{u,i}$ into $F_{u,i+1}$, for $C_u$
  and for the used $C_{u,\mu}$; and the classes of level $0$ of these chains have trivial
  stabilizers.
\item \emph{Inert expansions}: there are $\varepsilon_w$ with $|\varepsilon_w|=1$, signs
  $\sigma_w(z)$ and series $\beta_w(z)$ in variables indexed by $\mathbf Z/\ell_{p_w}\mathbf Z$ such
  that the following holds at every place $w$ free for $(u,\mu)$, with $p=p_w$, $\Lambda$ the set
  of places and $\operatorname{pos}=\operatorname{pos}_p$. If $p$ is a signal prime, or of inert class
  with $h_p$ of even order, then for every set $R$ of orbit representatives, $\sigma_w(r)=\pm1$ and
  $\beta_w(r)$ satisfies (B1) to (B4) for $n=\ell_p$ ($r\in R$), and every
  $\Theta_{u,\mu,w,i}(c,c')$ has the orbit expansion with $(R,\varepsilon_w,\sigma_w,\beta_w)$
  (Definition~\ref{pd-orbit-expansions}). If $p$ is of inert class and $h_p=0$, every
  $\Theta_{u,\mu,w,i}(c,c')$ has the expansion at a split orbit. If $p$ is of inert class and $h_p$
  has odd order at least $3$, every $\Theta_{u,\mu,w,i}(c,c')$ has the theta product form for
  $n=\ell_p$.
\item \emph{Orbit bounds}: let $w$ be a place with $p=p_w\geq5$ of inert class, $\ell_p$ even, $h_p\neq0$
  and $p\leq X''$, let $u$ be a word, $i<m$ and $t_0$ a representative at $p$. Then every series
  $\Theta$ in variables $X_{s,a,b}$, $s\in\mathbf Z/\ell_p\mathbf Z$, $a<N(i)$, $b<N(i+1)$, that has
  the orbit expansion with the single representative $0$, the constant $\varepsilon_w$, the sign
  $\sigma_w((0,t_0))$ and the series $\beta_w((0,t_0))$ has $\nu$-valuation at least
  $\max(g_0(p,i+1,t_0),0)$ (Definition~\ref{pd-minor-valuation}), for
  $\nu_s=\nu_{u,i}(\operatorname{pos}_p((0,t_0),s))$.
\item \emph{Jet integrality}: at every place $w$ free for $(u,\mu)$ with $p_w\neq d$, and every
  $i<m$, every $\Theta_{u,\mu,w,i}(c,c')$ has $\nu_{u,i}$-valuation at least $0$.
\item \emph{Shape at $d$}: at every place $w$ free for $(u,\mu)$ with $p_w=d$, every
  $\Theta_{u,\mu,w,i}(c,c')$ is
  $\exp\big(-S\sum_{t\in A}u_t\sum_{a,b}X_{(0,t),a,b}X_{(1,t),a,b}\big)$ times a series with
  integral jets, for some $S$ with $|S|\leq p_w^{1/2}$ and $u_t$ with $|u_t|\leq1$.
\item \emph{Shape of the row exponents}: for every word $u$, $i<m$ and place $z$, with
  $w=w_{i+1}(z)$ and $\nu=\nu_{u,i}(z)$: $w>0$; $0\leq|\nu|-(k-1)w\leq4w$ if $z$ is not a top place;
  $0\leq|\nu|-\nu_0-(k-3)w\leq4w$ and $2w\leq\nu_0$ if $z$ is a top place; $|\nu|\leq a^2$; the cap
  of $\nu$ is at most $b+\eta+1+3S_c$ and at most $\mathrm{fill}_{i+1}(z)+c_G$; and
  $a_t-2c_G\leq\nu_0-\nu_1$ when $z=(0,t)$ with $t\in T$. Moreover $2\log_5(a^2)+2\leq k-1$,
  $X''\leq a$ and $a>1$.
\end{enumerate}
\end{definition}
\leanref{Development/AutomorphicArithmetic/ThetaCorrespondence/ThetaPairLifts/ThirdPlaceSeries/GServT.lean}{Formalism.L2.G.ServT.SelectionServed}
\leanref{Formalism/SharpRestriction/GrowthBounds/IsResidueClassZero.lean}{Formalism.MainNotions.G.g0}
\leanref{Formalism/SharpRestriction/Divisibility/OrbitPlaces.lean}{Formalism.MainNotions.OrbitPlaces.EvenConditionsInert}
\leanref{Formalism/SharpRestriction/Divisibility/OrbitPlaces.lean}{Formalism.MainNotions.OrbitPlaces.EvenConditionsAtC}
\leanref{Formalism/SharpRestriction/Divisibility/OrbitPlaces.lean}{Formalism.MainNotions.OrbitPlaces.OrbitBoundWithoutOperator}
\leanref{Formalism/SharpRestriction/Divisibility/OrbitPlaces.lean}{Formalism.MainNotions.OrbitPlaces.JetIntegralityBis}
\leanref{Formalism/SharpRestriction/Divisibility/OrbitPlaces.lean}{Formalism.MainNotions.OrbitPlaces.ThetaKernelShapeAtD}
\leanref{Formalism/SharpRestriction/Divisibility/OrbitPlaces.lean}{Formalism.MainNotions.OrbitPlaces.Shape}
\leanref{Formalism/SharpRestriction/GrowthBounds/WeilJacquetShalikaLaw.lean}{Formalism.MainNotions.WeilJacquetShalikaLaw.CmpV}
\leanref{Formalism/SharpRestriction/GrowthBounds/WeilJacquetShalikaLaw.lean}{Formalism.MainNotions.WeilJacquetShalikaLaw.TwLoc}
\leanref{Formalism/SharpRestriction/SmoothDesignExistence/SumOverPrimesHypotheses.lean}{Formalism.MainNotions.Pk.SysTw}
\leanref{Formalism/SharpRestriction/GrowthBounds/WeilJacquetShalikaLaw.lean}{Formalism.MainNotions.WeilJacquetShalikaLaw.lowS}
\leanref{Formalism/LocalRingCardinality/DataErr.lean}{Formalism.BasicNotions.ErrDefs}

\begin{lemma}\label{pd-kernel-form}
Let $D$ be a definite-side datum with divisibility data and twisted forms whose stabilizer
denominators are defined, with kernel entries (Definition~\ref{pd-kernel-entries}). Suppose that
conditions (2) of Definition~\ref{pd-served-selection} hold: at free places the constants $\eta$ have
absolute value $1$ and the maps $k_i(c,c')$ of the used chains carry rational vectors to rational
vectors, and the lifts $K_i$ carry $F_{u,i}$ into $F_{u,i+1}$. Let $u$ be a word, $\mu$ a used tuple,
$w$ a place free for $(u,\mu)$, $i<m$ and $B\in\mathbf R$, and suppose that every
$\Theta_{u,\mu,w,i}(c,c')$ has $\nu_{u,i}$-valuation at least $B$
(Definition~\ref{pd-minor-valuation}). Put $s=B$ if $i=0$ and the classes of level $0$ of $C_u$ and of
the used $C_{u,\mu}$ have trivial stabilizers, and $s=B-v_{p_w}(D_\Gamma(u))$ otherwise. Then for
every vector $x$ of the lattice of forms of $C_{u,\mu}$ at level $i$ there are $a\in L$ and $z$ in the
$w$-local span of the lattice of forms at level $i+1$ such that $K_i(x)$ is the image of $az$, and
$a=0$ or $v_w(a)\geq e_ws$.
\end{lemma}
\leanref{Formalism/SharpRestriction/Divisibility/OrbitPlaces.lean}{Formalism.MainNotions.OrbitPlaces.kform_of_ker}
\leanref{Formalism/SharpRestriction/Divisibility/OrbitPlaces.lean}{Formalism.MainNotions.OrbitPlaces.kgeTw_of}
\leanref{Formalism/SharpRestriction/Divisibility/ChainOperators/GenericFramesLetters.lean}{Formalism.MainNotions.DIV.Lemmas.theta_isometry_facts_hold}

\begin{proof}
Write $p=p_w$, $\Theta=\Theta_{u,\mu,w,i}(c,c')$ and $\nu=\nu_{u,i}$.

\emph{Step 1: expansion in products of minors.} Let $I$ be a column choice into $N(i+1)$. By the
Cauchy-Binet formula, each determinant in $G_I$ is
$\sum_{J'}\det(Y_{s,r,J'(q)})\det(X_{s,J'(q'),I_{s,c}(q)})$, the sum over the strictly increasing maps
$J'$ of the right length into $N(i)$, so $G_I=\sum_J\nabla_J\,M_{I,J}(X)$ with polynomials $M_{I,J}$
in the $X$-variables, and $P=\langle G_I,\Theta\rangle_X$ lies in the $\mathbf C_p$-span of the
$\nabla_J$. For every commutative ring $R$ the $R$-span of the products $\nabla_J$ is a direct summand
of the polynomial ring over $R$ in the $Y$-variables; this is the standard basis theorem for
bitableaux \cite{DEP80}. Take $R$ to be the valuation ring of $\mathbf C_p$ and let
$\pi\in\mathbf C_p$ have $|\pi|=p^{-B}$. The coefficients of $\pi^{-1}P$ have absolute value at most
$1$ by hypothesis, so $\pi^{-1}P$ lies in the $R$-span of the $\nabla_J$. Hence
$P=\sum_J\gamma_J\nabla_J$ with $|\gamma_J|\leq p^{-B}$.

\emph{Step 2: the target coordinates.} Let $x$ be a vector of the lattice of forms at level $i$ and
$\varphi$ its image. By condition (a) of Definition~\ref{pd-kernel-entries}, $\varphi\in F_{u,i}$
and $\varphi_c=\operatorname{ev}(v_c)$ with $|S_w(v_c)(\nabla_J)|\leq1$ for all $J$. By the
rationality of the maps $k_i(c,c')$, $k_i(c,c')(\varphi_c)=\operatorname{ev}(v'_{c,c'})$ for some
rational $v'_{c,c'}$. By condition (d), $|\eta|=1$ and Step 1,
\[
 |T_w(v'_{c,c'})(I)|=\Big|S_w(v_c)\Big(\sum_J\gamma_J\nabla_J\Big)\Big|\leq\max_J|\gamma_J|\leq p^{-B}
\]
for every $c,c'$ and $I$.

\emph{Step 3: scaling.} Take $a\in L$ with $v_w(a)=\lceil e_wB\rceil$. By condition (c) each
$T_w(v'_{c,c'})(I)$ has the absolute value of some $b\in L$, and $|\iota_w(b)|\leq p^{-B}$ gives
$v_w(b)\geq e_wB$, hence $v_w(b)\geq v_w(a)$. So the target coordinates of $a^{-1}v'_{c,c'}$ have
absolute value at most $1$.

\emph{Step 4: the stabilizers.} Let $n=1$ in the first case of the statement and
$n=D_\Gamma(u,\mu)$ in the second; then $|\Gamma_c|$ divides $n$ for every class $c$ of level $i$.
Write $n=p^\kappa n'$ with $p\nmid n'$. Then $|\Gamma_c|^{-1}=p^{-\kappa}r_c$ with
$r_c=(n/|\Gamma_c|)/n'$, and $v_w(r_c)\geq0$. By Definition~\ref{th-iterated-lift},
\[
 (K_i\varphi)_{c'}=\sum_c|\Gamma_c|^{-1}k_i(c,c')(\varphi_c)=p^{-\kappa}a\sum_cr_c\,a^{-1}
 \operatorname{ev}(v'_{c,c'}).
\]
The function $\psi=p^\kappa a^{-1}K_i\varphi$ lies in $F_{u,i+1}$, because $K_i\varphi$ does, and
$\psi_{c'}=\operatorname{ev}\big(\sum_cr_ca^{-1}v'_{c,c'}\big)$ has target coordinates of absolute
value at most $1$, by Step 3 and the ultrametric inequality. By condition (b), $\psi$ is the image of
a vector $z$ of the $w$-local span of the lattice of forms at level $i+1$. So $K_i(x)$ is the image of
$(p^{-\kappa}a)z$, and $v_w(p^{-\kappa}a)\geq e_wB-e_w\kappa$. In the first case $\kappa=0$. In the
second, $D_\Gamma(u,\mu)=D_\Gamma(u_0)=D_\Gamma(u)$ because the stabilizer denominators are defined,
so $\kappa=v_p(D_\Gamma(u))$. Finally, the twisted forms at $\mu_0$ are identified with the
lattices of forms of $C_u$ by isomorphisms compatible with the images and with the lifts, so the
bound for $C_u$ follows from the bound for the twisted chain at $\mu_0$.
\end{proof}

\subsection{Weighted Dwork bounds}\label{subsection-pd-weighted-dwork}

\begin{definition}\label{pd-odd-level-programme}
Let $l\geq3$ and $d\geq2$ be integers with $d$ even, and put $\gamma_l=l/(l^2-1)$. For
$D'\colon\mathbf Z/d\mathbf Z\to\mathbf R_{\geq0}$, a \emph{flow} for $D'$ is a family of real numbers
$x_{s,i}\geq0$, for $s\in\mathbf Z/d\mathbf Z$ and odd $i<d$, with
$\sum_il^ix_{s,i}+\sum_ix_{s-i,i}\leq D'(s)$ for every $s$. Let $M(D')$ be the largest value of
$\sum_{s,i}x_{s,i}$ over the flows, and $F(D')=\gamma_l\sum_sD'(s)-M(D')$. A \emph{dual vector} is a
function $\sigma\colon\mathbf Z/d\mathbf Z\to[0,1]$ with $l^i\sigma(s)+\sigma(s+i)\geq1$ for every $s$
and every odd $i<d$; it is \emph{optimal} for $D'$ if
$F(D')=\sum_sD'(s)\big(\gamma_l-\sigma(s)\big)$.
\end{definition}
\leanref{Formalism/IncomingBounds/IncomingBound.lean}{Formalism.IncomingBounds.IsFlow}
\leanref{Formalism/IncomingBounds/IncomingBound.lean}{Formalism.IncomingBounds.Fnum}
\leanref{Formalism/IncomingBounds/IncomingBound.lean}{Formalism.IncomingBounds.IsDual}
\leanref{Formalism/IncomingBounds/IncomingBound.lean}{Formalism.IncomingBounds.IsOptDual}
\leanref{Formalism/LocalRingCardinality/DataParams₂.lean}{Formalism.BasicNotions.lpF}

\begin{lemma}\label{pd-lp-duality}
Let $l\geq3$, $d\geq2$ even, and $D',E'\colon\mathbf Z/d\mathbf Z\to\mathbf R_{\geq0}$
(Definition~\ref{pd-odd-level-programme}).
\begin{enumerate}
\item $F(D')$ is the largest value of $\sum_sD'(s)(\gamma_l-\sigma(s))$ over the dual vectors
  $\sigma$, and it is attained.
\item If $\sigma$ is optimal for $D'$, then
  $F(D')-\sum_sE'(s)\max(\sigma(s)-\gamma_l,0)\leq F(D'+E')\leq F(D')+\gamma_l\sum_sE'(s)$.
\item If $\sigma$ is optimal for $D'$, then the sum of $D'(s)$ over the $s$ with
  $\sigma(s)\geq2/l$ is at most $\frac1l\sum_sD'(s)$.
\end{enumerate}
\end{lemma}
\leanref{Formalism/IncomingBounds/IncomingBound.lean}{Formalism.IncomingBounds.LemmaFa}
\leanref{Formalism/IncomingBounds/IncomingBound.lean}{Formalism.IncomingBounds.LemmaFc}
\leanref{Formalism/IncomingBounds/IncomingBound.lean}{Formalism.IncomingBounds.IncomingBound}
\leanref{Formalism/IncomingBounds/FlowDegreeBounds/LemmaFd.lean}{Formalism.IncomingBounds.incomingBound_of}

\begin{proof}
(1) The programme defining $M(D')$ is a linear programme with a bounded nonempty feasible set (the
zero flow is feasible, and $x_{s,i}\leq D'(s)$), and its dual is to minimize $\sum_sD'(s)\sigma(s)$
over the $\sigma\geq0$ with $l^i\sigma(s)+\sigma(s+i)\geq1$. By the duality theorem of linear
programming the two optima are equal and attained. Replacing $\sigma$ by $\min(\sigma,1)$ keeps the
constraints and does not increase the objective, so the minimum is attained at a dual vector, and
$F(D')=\gamma_l\sum D'-\min_\sigma\sum D'\sigma$ is the stated maximum.

(2) By (1) applied to $D'+E'$, $F(D'+E')\geq\sum_s(D'+E')(s)(\gamma_l-\sigma(s))=F(D')+
\sum_sE'(s)(\gamma_l-\sigma(s))$, and $\gamma_l-\sigma(s)\geq-\max(\sigma(s)-\gamma_l,0)$. For the upper
bound let $\sigma'$ be optimal for $D'+E'$. Then
$F(D'+E')=\sum_sD'(s)(\gamma_l-\sigma'(s))+\sum_sE'(s)(\gamma_l-\sigma'(s))$; the first sum is at most
$F(D')$ by (1), and the second at most $\gamma_l\sum_sE'(s)$ since $\sigma'\geq0$.

(3) An optimal $\sigma$ satisfies $\sum_sD'(s)\sigma(s)=M(D')$, so it is an optimal solution of the
dual programme. Let $x$ be a flow with $\sum_{s,i}x_{s,i}=M(D')$. By complementary slackness, if
$\sigma(s)>0$ the constraint of $x$ at $s$ is an equality, and if $x_{s,i}>0$ then
$l^i\sigma(s)+\sigma(s+i)=1$. Suppose $\sigma(s)\geq2/l$. For odd $i\geq1$, $l^i\sigma(s)\geq2>1$, so
$x_{s,i}=0$ for all $i$, and the equality at $s$ reads $D'(s)=\sum_ix_{s-i,i}$. Summing over these $s$,
their sum is at most $\sum_{t,i}x_{t,i}$. For every $t$,
$l\sum_ix_{t,i}\leq\sum_il^ix_{t,i}\leq D'(t)$, so $\sum_{t,i}x_{t,i}\leq\frac1l\sum_tD'(t)$.
\end{proof}

\begin{lemma}\label{pd-lp-comparison}
Let $l\geq3$, $d\geq2$ even, $T\subset\mathbf Z/d\mathbf Z$ and $k\geq3$. Let
$w,E,D\colon\mathbf Z/d\mathbf Z\to\mathbf R$ with $w(s)>0$, $0\leq E(s)\leq4w(s)$, $D(s)=(k-1)w(s)+E(s)$
for $s\notin T$ and $D(s)=(k-3)w(s)+E(s)$ for $s\in T$. Then, with
$F$ as in Definition~\ref{pd-odd-level-programme},
\[
 F\big((k-1)w\big)-\frac{11}l\sum_sw(s)\leq F\big((k-1)w\big)-2\gamma_l\sum_{s\in T}w(s)-
 \frac8l\sum_sw(s)\leq F(D).
\]
\end{lemma}
\leanref{Formalism/GaussJetBounds/OneLetterGaussJets/StatementN.lean}{Formalism.DworkPowerSeriesWeights.OddLevelGainBounds.PropSix7}
\leanref{Formalism/GaussJetBounds/OneLetterGaussJets/OneLetterGaussProofs/CompanionTheory.lean}{Formalism.DworkPowerSeriesWeights.OddLevelGainBounds.proposition_six}

\begin{proof}
Put $S=\sum_sw(s)$, $S_T=\sum_{s\in T}w(s)$, $W'=(k-1)w$ and $E'(s)=2w(s)$ for $s\in T$, $E'(s)=0$
otherwise. Then $D\geq0$, $E'\geq0$ and $D+E'=W'+E$.

\emph{Step 1.} Let $\sigma_D$ be optimal for $D$ (Lemma~\ref{pd-lp-duality}(1)). By
Lemma~\ref{pd-lp-duality}(2) for $D$ and $E'$, $F(W'+E)=F(D+E')\leq F(D)+2\gamma_lS_T$.

\emph{Step 2.} Let $\sigma$ be optimal for $W'$. By Lemma~\ref{pd-lp-duality}(2) for $W'$ and $E$,
$F(W')-\sum_sE(s)\max(\sigma(s)-\gamma_l,0)\leq F(W'+E)$.

\emph{Step 3.} By Lemma~\ref{pd-lp-duality}(3) for $W'$, the sum of $(k-1)w(s)$ over the $s$ with
$\sigma(s)\geq2/l$ is at most $(k-1)S/l$, so the sum of $w(s)$ over these $s$ is at most $S/l$.
Since $\gamma_l\geq1/l$ and $0\leq\sigma\leq1$: if $\sigma(s)\geq2/l$, then
$E(s)\max(\sigma(s)-\gamma_l,0)\leq E(s)\leq4w(s)$; otherwise $\max(\sigma(s)-\gamma_l,0)\leq1/l$ and
$E(s)\max(\sigma(s)-\gamma_l,0)\leq4w(s)/l$. Summing,
$\sum_sE(s)\max(\sigma(s)-\gamma_l,0)\leq4S/l+4S/l=8S/l$.

\emph{Step 4.} Steps 1 to 3 give $F(W')-8S/l\leq F(D)+2\gamma_lS_T$, which is the second inequality.
For the first, $2\gamma_l=2l/(l^2-1)\leq3/l$ because $l^2\geq3$, and $S_T\leq S$, so
$2\gamma_lS_T+8S/l\leq11S/l$.
\end{proof}

\begin{definition}\label{pd-weights}
Let $K$ be a field complete for a non-archimedean absolute value and $l$ a prime with
$0<|l|<1$. Put $\rho_l=1/(l^2-1)$, so that $\gamma_l=l/(l^2-1)=1/(l-1)-\rho_l$
(Definition~\ref{pd-odd-level-programme}). Let $\iota$ be a set, $r\colon\iota\to\mathbf R$ a
family of \emph{weights} and $w\in\mathbf R$.
\begin{enumerate}
\item A coefficient $c\in K$ at the monomial $X^n=\prod_iX_i^{n_i}$ \emph{has weight at least $w$}
  if $|c|\leq|l|^{\,w-\sum_in_ir_i}$. A power series $f=\sum_nc_nX^n$ over $K$ in the variables
  $X_i$, $i\in\iota$, \emph{has weight at least $w$} if every $c_n$ has, and \emph{has weight at
  least $w$ up to} an exponent $\kappa$ if $c_n$ has for every $n\leq\kappa$.
\item A power series is \emph{integral} if its coefficients have absolute value at most $1$.
\item For $\kappa\in\mathbf N^\iota$ we write $|\kappa|=\sum_i\kappa_i$ and
  $\kappa!=\prod_i\kappa_i!$.
\end{enumerate}
\end{definition}
\leanref{Formalism/GaussJetBounds/Statement.lean}{Formalism.DworkPowerSeriesWeights.MonoWtGe}
\leanref{Formalism/GaussJetBounds/Statement.lean}{Formalism.DworkPowerSeriesWeights.WtGe}
\leanref{Formalism/GaussJetBounds/GenLogInverse.lean}{Formalism.DworkPowerSeriesWeights.WtGeLe}
\leanref{Formalism/GaussJetBounds/Statement.lean}{Formalism.DworkPowerSeriesWeights.Integral}
\leanref{Formalism/IncomingBounds/IncomingBound.lean}{Formalism.IncomingBounds.gam}

With the weights $\rho_l$ on every variable, a series of weight at least $w$ converges on the
polydisc of radius $|l|^{\rho_l}$, and the weight measures its size there.

\begin{lemma}\label{pd-weight-calculus}
Let $K$ be a field of characteristic $0$, complete for a non-archimedean absolute value, and $l$ a
prime with $0<|l|<1$ (Definition~\ref{pd-weights}). Let $\sigma$ and $\iota$ be finite sets.
\begin{enumerate}
\item Let $r\colon\sigma\to\mathbf R$ and $r'\colon\iota\to\mathbf R$ be weights, $f$ a series in
  variables indexed by $\sigma$ of weight at least $w$ for $r$, and $a_j$ ($j\in\sigma$) series in
  variables indexed by $\iota$ with constant term $0$ such that $a_j$ has weight at least $r_j$
  for $r'$. Then the series $f(a)$, obtained by substituting $a_j$ for the $j$-th variable, has
  weight at least $w$ for $r'$.
\item If $r\leq r'$ pointwise, a series of weight at least $w$ (up to $\kappa$) for $r$ has
  weight at least $w$ (up to $\kappa$) for $r'$.
\item If $f$ and $g$ have weights at least $w_1$ and $w_2$ up to $\kappa$ for the same weights,
  then $fg$ has weight at least $w_1+w_2$ up to $\kappa$.
\item If $B$ has constant term $0$ and weight at least $1/(l-1)$ up to $\kappa$, then $\exp(B)$ has
  weight at least $0$ up to $\kappa$.
\item If $D$ has constant term $0$ and $\exp(D)-1$ has weight at least $1/(l-1)$, then $D$ has
  weight at least $1/(l-1)$.
\item For $\kappa\in\mathbf N^\iota$ and an integer $a\geq1$ with $\kappa_i\leq a$ for every $i$,
  \[
   |\kappa!|\leq|l|^{\,|\kappa|/(l-1)-|\iota|(\log_la+1)};
  \]
  the same holds with $a=|\kappa|$, where $\log_l0$ is read as $0$.
\end{enumerate}
\end{lemma}
\leanref{Formalism/GaussJetBounds/GaussJetProofs/SubstitutionWeightBound.lean}{Formalism.DworkPowerSeriesWeights.substitution_weight_bound}
\leanref{Formalism/GaussJetBounds/GaussJetProofs/SubstitutionWeightBound.lean}{Formalism.DworkPowerSeriesWeights.radical_weight_monotone}
\leanref{Formalism/GaussJetBounds/GaussJetProofs/NormNatCastEqPowPadicValNat.lean}{Formalism.DworkPowerSeriesWeights.multiplicative_weight_le}
\leanref{Formalism/GaussJetBounds/GaussJetProofs/NormNatCastEqPowPadicValNat.lean}{Formalism.DworkPowerSeriesWeights.exponential_weight_le}
\leanref{Formalism/GaussJetBounds/GaussJetProofs/LogStep.lean}{Formalism.DworkPowerSeriesWeights.logarithm_weight_induction}
\leanref{Formalism/GaussJetBounds/GaussJetProofs/LemmaRA.lean}{Formalism.DworkPowerSeriesWeights.factValA}
\leanref{Formalism/GaussJetBounds/GaussJetProofs/NormNatCastEqPowPadicValNat.lean}{Formalism.DworkPowerSeriesWeights.factorial_valuation_bound}

\begin{proof}
(1) to (3) follow from the ultrametric inequality, since a monomial $a^n$ has weight at least
$\sum_jn_jr_j$ for $r'$. For (4) and (5) use $v_l(k!)\leq(k-1)/(l-1)$: in (4), $B^k/k!$ has weight
at least $k/(l-1)-(k-1)/(l-1)>0$ up to $\kappa$; in (5), the coefficients of $D$ are bounded by
induction on the degree, from $D=(\exp(D)-1)-\sum_{k\geq2}D^k/k!$. For (6), Legendre's formula
gives $v_l(k!)\geq k/(l-1)-\log_lk-1$ for $k\geq1$, and the product has at most $|\iota|$ factors
other than $1$.
\end{proof}

\begin{definition}\label{pd-letter-series}
Let $K$ be a field of characteristic $0$, complete for a non-archimedean absolute value, $l$ an odd
prime with $0<|l|<1$, and $d\geq2$ an even integer.
\begin{enumerate}
\item For a finite set $\tau$ and a map $q\colon\tau\to\tau$, the \emph{logarithm along $q$} is the
  family of series $\ell_q(X)_x=\sum_{e\geq0}l^{-e}X_{q^e(x)}^{\,l^{2e}}$, $x\in\tau$. For
  $\tau=\mathbf Z/d\mathbf Z$ (or $\{0,1\}\times\mathbf Z/d\mathbf Z$, or
  $\mathbf Z/d\mathbf Z\times\{0,\dots,N_0-1\}\times\{0,\dots,N_1-1\}$) and $q$ the shift by $-2$
  in the $\mathbf Z/d\mathbf Z$ coordinate, this is the orbit logarithm of
  Definition~\ref{pd-orbit-expansions}(1); for $q$ the identity it is the \emph{one-place logarithm}
  $\sum_el^{-e}X_x^{\,l^{2e}}$.
\item A monomial $\prod_sY_s^{n_s}$ in variables indexed by $\mathbf Z/d\mathbf Z$ is a
  \emph{bipower} if it is $Y_s^{\,l^e}Y_{s'}^{\,l^f}$ with $s$ even and $s'$ odd (as
  representatives in $\{0,\dots,d-1\}$), $(e,f)\neq(0,0)$, and $s+e=s'+f$ in
  $\mathbf Z/d\mathbf Z$.
\item Let $K_0\subset K$ be a subfield on which every nonzero absolute value is an integer power
  of $|l|$, and $\zeta\in K$ a primitive $(l^d-1)$-th root of unity. A series
  $\beta\in K[[Y_s:s\in\mathbf Z/d\mathbf Z]]$ is a \emph{letter series} for $(K_0,\zeta)$ if its
  coefficients lie in $K_0$, $\beta(\tau_\zeta Y)=\beta(Y)$ with $\tau_\zeta$ as in condition (B2)
  of Definition~\ref{pd-orbit-expansions}, $\beta$ satisfies condition (B3) there, and the series
  $\exp(\delta_{\mathrm{ev}}\beta)$ and $\exp(\delta_{\mathrm{od}}\beta)$ of condition (B4), with
  the orbit logarithm substituted for each family of variables, are integral with constant term
  $1$.
\end{enumerate}
\end{definition}
\leanref{Formalism/GaussJetBounds/GenLogInverse.lean}{Formalism.DworkPowerSeriesWeights.genLog}
\leanref{Formalism/GaussJetBounds/Statement.lean}{Formalism.DworkPowerSeriesWeights.letterLog}
\leanref{Formalism/GaussJetBounds/Statement.lean}{Formalism.DworkPowerSeriesWeights.IsBiPower}
\leanref{Formalism/GaussJetBounds/Statement.lean}{Formalism.DworkPowerSeriesWeights.BHyp}
\leanref{Formalism/GaussJetBounds/Statement.lean}{Formalism.DworkPowerSeriesWeights.Unram}

\begin{lemma}\label{pd-log-inverse}
Let $K$ be a field of characteristic $0$, complete for a non-archimedean absolute value, $l$ a prime
with $0<|l|<1$, $\tau$ a finite set and $q\colon\tau\to\tau$ a map, with logarithm $\ell_q$
(Definition~\ref{pd-letter-series}). There is a family $g=(g_x)_{x\in\tau}$ of series with constant
term $0$ such that $\ell_q(g)=X$ and $g(\ell_q)=X$. If $g$ is any family of series with constant
term $0$ and $\ell_q(g)=X$, then each $g_x$ has weight at least $\rho_l$ for the weights $\rho_l$ on
every variable (Definition~\ref{pd-weights}).
\end{lemma}
\leanref{Formalism/GaussJetBounds/GenLogInverse.lean}{Formalism.DworkPowerSeriesWeights.GenLogInverse}
\leanref{Formalism/GaussJetBounds/GenLogInverse.lean}{Formalism.DworkPowerSeriesWeights.GenLogWeight}
\leanref{Formalism/GaussJetBounds/GaussJetProofs/LogarithmWeightComparison.lean}{Formalism.DworkPowerSeriesWeights.generalized_log_inverse}
\leanref{Formalism/GaussJetBounds/GaussJetProofs/LogarithmWeightComparison.lean}{Formalism.DworkPowerSeriesWeights.generalized_log_weight}

\begin{proof}
The linear part of $\ell_q$ is the identity, so $g$ exists and is a two-sided inverse. For the
weight, induct on the degree of the monomial. The identity $\ell_q(g)=X$ reads
$g_x=X_x-\sum_{e\geq1}l^{-e}g_{q^e(x)}^{\,l^{2e}}$. In degree $N$, the coefficient of
$g_{q^e(x)}^{\,l^{2e}}$ involves only coefficients of $g$ of degree below $N$, which have weight at
least $\rho_l$ by induction; so it has weight at least $l^{2e}\rho_l$, and after division by $l^e$
at least $l^{2e}\rho_l-e\geq\rho_l$, because $(l^{2e}-1)/(l^2-1)\geq e$.
\end{proof}

\begin{lemma}\label{pd-binomial-extraction}
Let $K$ be a field of characteristic $0$, complete for a non-archimedean absolute value, $l$ an odd
prime with $0<|l|<1$, $d\geq2$ even and $\zeta\in K$ a primitive $(l^d-1)$-th root of unity. Let
$\beta\in K[[Y_s:s\in\mathbf Z/d\mathbf Z]]$ satisfy $\beta(\tau_\zeta Y)=\beta(Y)$ and condition
(B3) of Definition~\ref{pd-orbit-expansions}, and suppose that $\delta_{\mathrm{ev}}\beta$ and
$\delta_{\mathrm{od}}\beta$ have weight at least $1/(l-1)$ for the weights $\rho_l$
(Definition~\ref{pd-weights}). Then $\beta$ has weight at least $1/(l-1)-1$ for the weights
$\rho_l$, and its coefficient at every monomial that is not a bipower
(Definition~\ref{pd-letter-series}) has weight at least $1/(l-1)$.
\end{lemma}
\leanref{Formalism/GaussJetBounds/GenLogInverse.lean}{Formalism.DworkPowerSeriesWeights.BinomExtract}
\leanref{Formalism/GaussJetBounds/GaussJetProofs/SwappedParitySeriesBinomialForm.lean}{Formalism.DworkPowerSeriesWeights.binomial_coefficient_extraction}

This is a self-contained computation. A coefficient of $\beta$ appears in the expansions of
$\delta_{\mathrm{ev}}\beta$ and $\delta_{\mathrm{od}}\beta$ multiplied by the binomial coefficients
that split the exponent of one variable between $Y_s$ and $Z_s$. If one of them is a unit, the
coefficient has weight at least $1/(l-1)$. Otherwise the monomial is $Y_s^{\,l^e}Y_{s'}^{\,l^f}$
with $s$ even and $s'$ odd, a binomial coefficient of valuation $1$ gives weight at least
$1/(l-1)-1$, and the torus invariance gives $s+e=s'+f$.

\begin{lemma}\label{pd-letter-weights}
Let $K$, $l$, $d$, $K_0$ and $\zeta$ be as in Definition~\ref{pd-letter-series}(3), and let $\beta$
be a letter series for $(K_0,\zeta)$. Then $\beta$ has weight at least $1/(l-1)-1$ for the weights
$\rho_l$, and its coefficient at every monomial that is not a bipower has weight at least
$1/(l-1)$.
\end{lemma}
\leanref{Formalism/GaussJetBounds/Statement.lean}{Formalism.DworkPowerSeriesWeights.LemmaS}
\leanref{Formalism/GaussJetBounds/GaussJetProofs/ConstantCoeffEvenDerivativeEqZero.lean}{Formalism.DworkPowerSeriesWeights.lemmaS_of}
\leanref{Formalism/GaussJetBounds/GaussJetProofs/ConstantCoeffEvenDerivativeEqZero.lean}{Formalism.DworkPowerSeriesWeights.asmS_diff}
\leanref{Formalism/GaussJetBounds/GenLogInverse.lean}{Formalism.DworkPowerSeriesWeights.ExpDeInv}
\leanref{Formalism/GaussJetBounds/GaussJetProofs/OneLetterRadiusIdentity.lean}{Formalism.DworkPowerSeriesWeights.weightQuantization}

\begin{proof}
Let $D$ be $\delta_{\mathrm{ev}}\beta$ or $\delta_{\mathrm{od}}\beta$, a series in the variables
$Y_s,Z_s$, $s\in\mathbf Z/d\mathbf Z$, with constant term $0$ by (B3). For a monomial $Y^mZ^{m'}$ put
$h(m,m')=\sum_s\pm(m_s+m'_s)l^{s}$, with $s$ read in $\{0,\dots,d-1\}$ and the sign $+$ for even
$s$ and $-$ for odd $s$.

\emph{Step 1: the torus.} The coefficients of $\exp(D)$ lie in $K_0$. Since $\beta(\tau_\zeta Y)=
\beta(Y)$, $D$ and $\exp(D)$ are invariant under the substitution that multiplies $Y_s$ and $Z_s$ by $\zeta^{\pm l^s}$, with the
same signs. So a nonzero coefficient of $\exp(D)$ at $Y^mZ^{m'}$ has $\zeta^{h(m,m')}=1$, that is,
$l^d-1$ divides $h(m,m')$.

\emph{Step 2: the logarithm.} Let $g$ be the inverse of the orbit logarithm $\ell$ on the variables
$Y_s,Z_s$ (Lemma~\ref{pd-log-inverse}). By the last condition of a letter series,
$\exp(D)(\ell)-1$ is integral with constant term $0$, so it has weight at least $\rho_l$ for the
weights $\rho_l$. Each $g_x$ has weight at least $\rho_l$, so by Lemma~\ref{pd-weight-calculus}(1)
the series $\exp(D)-1=(\exp(D)(\ell)-1)(g)$ has weight at least $\rho_l$.

\emph{Step 3: quantization.} Let $c\neq0$ be a coefficient of $\exp(D)-1$ at $Y^mZ^{m'}$, of total
degree $N$. Write $v(c)$ for the integer with $|c|=|l|^{v(c)}$, and $O$ for the degree in the
variables with odd index. Since $l^2\equiv1$ modulo $l^2-1$, which divides $l^d-1$, Step 1 gives
$N\equiv(l+1)O$ modulo $l^2-1$, so $v(c)+\rho_lN$ lies in $O/(l-1)+\mathbf Z$. By Step 2 it is at
least $\rho_l>0$, and the positive elements of $O/(l-1)+\mathbf Z$ are at least $1/(l-1)$. So
$\exp(D)-1$ has weight at least $1/(l-1)$.

\emph{Step 4.} By Lemma~\ref{pd-weight-calculus}(5), $D$ has weight at least $1/(l-1)$. This holds
for $\delta_{\mathrm{ev}}\beta$ and $\delta_{\mathrm{od}}\beta$, and Lemma~\ref{pd-binomial-extraction}
gives the result.
\end{proof}

\begin{lemma}\label{pd-letter-substitution}
Let $K$ be a field of characteristic $0$, complete for a non-archimedean absolute value, $l$ a prime
with $0<|l|<1$, $d\geq2$ even, $N_0,N_1\geq1$, and $\beta\in K[[Y_s:s\in\mathbf Z/d\mathbf Z]]$ with
constant term $0$ that has weight at least $1/(l-1)-1$ for the weights $\rho_l$ and whose coefficient
at every monomial that is not a bipower has weight at least $1/(l-1)$
(Definitions~\ref{pd-weights} and~\ref{pd-letter-series}). Let $\iota$ be a finite set and
$M=(M_{y,i})$, for $y=(s,a,b)\in\mathbf Z/d\mathbf Z\times\{0,\dots,N_0-1\}\times\{0,\dots,N_1-1\}$
and $i\in\iota$, a matrix with $|M_{y,i}|\leq1$. Put $L_y=\sum_iM_{y,i}X_i$ and
\[
 B_M=\sum_{a<N_0,\ b<N_1}\beta\big((L_{(s,a,b)})_{s\in\mathbf Z/d\mathbf Z}\big).
\]
\begin{enumerate}
\item $B_M$ has weight at least $1/(l-1)-1$ for the weights $\rho_l$, and its coefficient at every
  monomial that is not of the form $X_i^{\,l^e}X_j^{\,l^f}$ ($i,j\in\iota$, $e,f\geq0$)
  has weight at least $1/(l-1)$.
\item Suppose that $\iota$ has a map $\operatorname{pos}\colon\iota\to\mathbf Z/d\mathbf Z$ with
  $M_{y,i}=0$ unless $\operatorname{pos}(i)$ is the first coordinate of $y$. If the coefficient of
  $B_M$ at $X_i^{\,l^e}X_j^{\,l^f}$ does not have weight at least $1/(l-1)$, then
  $\operatorname{pos}(i)\neq\operatorname{pos}(j)$, $\operatorname{pos}(i)+e=\operatorname{pos}(j)+f$
  and $e+f$ is odd.
\item Let $x=(x_{s,r,a})$, $s\in\mathbf Z/d\mathbf Z$, $r,a<N_0$, have $|x_{s,r,a}|\leq1$, let
  $\iota=\mathbf Z/d\mathbf Z\times\{0,\dots,N_0-1\}\times\{0,\dots,N_1-1\}$, and let $M^x$ be the
  matrix with $L_{(s,a,b)}=\sum_rx_{s,r,a}X_{s,r,b}$. Let $s\neq s'$, $k<N_0$, $b,b'<N_1$ and
  $e,f\geq0$. If $b\neq b'$, the coefficient of $B_{M^x}$ at
  $X_{s,0,b}^{\,l^e}X_{s',k,b'}^{\,l^f}$ is $0$. If
  $|\sum_ax_{s,0,a}^{\,l^{e-f}}x_{s',k,a}|\leq|l|$ when $f\leq e$, or
  $|\sum_ax_{s,0,a}x_{s',k,a}^{\,l^{f-e}}|\leq|l|$ when $e<f$, then the coefficient
  at $X_{s,0,b}^{\,l^e}X_{s',k,b}^{\,l^f}$ has weight at least $1/(l-1)$.
\end{enumerate}
\end{lemma}
\leanref{Formalism/GaussJetBounds/GenLogInverse.lean}{Formalism.DworkPowerSeriesWeights.PureCross}
\leanref{Formalism/GaussJetBounds/GaussJetProofs/PureCrossTermWeight.lean}{Formalism.DworkPowerSeriesWeights.pure_cross_term_weight}
\leanref{Formalism/GaussJetBounds/GaussJetProofs/ThmAJet.lean}{Formalism.DworkPowerSeriesWeights.jet_collapse}
\leanref{Formalism/GaussJetBounds/GaussJetProofs/ThmAJet.lean}{Formalism.DworkPowerSeriesWeights.jet_bipower}
\leanref{Formalism/GaussJetBounds/StatementC.lean}{Formalism.DworkPowerSeriesWeights.FramePure}
\leanref{Formalism/GaussJetBounds/GaussJetProofs/FrameCoeffMulLin.lean}{Formalism.DworkPowerSeriesWeights.frame_pure_weight}
\leanref{Formalism/GaussJetBounds/GaussJetProofs/FrameA.lean}{Formalism.DworkPowerSeriesWeights.frameA_frob_norm}

\begin{proof}
Substituting linear forms preserves the total degree. (1) A monomial $\prod_sY_s^{n_s}$ of $\beta$
becomes $\prod_sL_{(s,a,b)}^{\,n_s}$, whose coefficients have absolute value at most $1$. If the
monomial is a bipower $Y_s^{\,l^e}Y_{s'}^{\,l^f}$, the multinomial coefficients of
$L^{\,l^e}_{(s,a,b)}L^{\,l^f}_{(s',a,b)}$ are divisible by $l$ at every monomial not of the
form $X_i^{\,l^e}X_j^{\,l^f}$, and this factor $l$ makes up the difference between the two
weights. (2) Under the hypothesis, the
coefficient of $B_M$ at $X_i^{\,l^e}X_j^{\,l^f}$ is a combination, with coefficients of
absolute value at most $1$, of the coefficient of $\beta$ at $Y_{\operatorname{pos}(i)}^{\,l^e}
Y_{\operatorname{pos}(j)}^{\,l^f}$. That monomial is then a bipower, which forces the two
positions to have different parities, hence to differ, and $e-f$ to be odd because $d$ is
even. (3) The variables with last index $b$ occur only in the summands of $B_{M^x}$ with that index,
which gives the first claim. For $b=b'$, the only monomial of $\beta$ that contributes is
$Y_s^{\,l^e}Y_{s'}^{\,l^f}$, with factor $\sum_ax_{s,0,a}^{\,l^e}x_{s',k,a}^{\,l^f}$.
For $f\leq e$ this sum is congruent modulo $l$ to
$(\sum_ax_{s,0,a}^{\,l^{e-f}}x_{s',k,a})^{l^f}$, so it has absolute value at most $|l|$;
the case $e<f$ is the same.
\end{proof}

\begin{proposition}\label{pd-weighted-dwork}
Let $K$ be a field of characteristic $0$, complete for a non-archimedean absolute value, $l$ an odd
prime with $0<|l|<1$, $d\geq2$ even, $N_0,N_1\geq1$, $K_0$ and $\zeta$ as in
Definition~\ref{pd-letter-series}(3), and $\beta$ a letter series for $(K_0,\zeta)$. Let $\iota$ be a
finite set with a map $\operatorname{pos}\colon\iota\to\mathbf Z/d\mathbf Z$, and $M$ a matrix as
in Lemma~\ref{pd-letter-substitution}(2), with $L_y$ and $B_M$ as there. Let $g$ be a family of
series with constant term $0$ in the variables $X_y$,
$y\in\mathbf Z/d\mathbf Z\times\{0,\dots,N_0-1\}\times\{0,\dots,N_1-1\}$, such that substituting $g$
into the orbit logarithm gives $X$, let $U$ be an integral series in these variables and
$\varepsilon=\pm1$. Let $\sigma\colon\iota\to\mathbf R_{\geq0}$, $a\geq1$ an integer and
$\kappa\in\mathbf N^\iota$ with $\kappa_i\leq a$ for all $i$. Suppose that for all $i,j\in\iota$ and
$e,f\geq0$ such that $X_i^{\,l^e}X_j^{\,l^f}$ divides $X^\kappa$ and the coefficient
of $B_M$ at $X_i^{\,l^e}X_j^{\,l^f}$ does not have weight at least $1/(l-1)$ for the
weights $\rho_l$,
\[
 l^e\sigma_i+l^f\sigma_j\geq1 .
\]
Let $\theta_\kappa$ be the coefficient of $X^\kappa$ in $U(g(L))\cdot\exp(\varepsilon B_M)$, where
$U(g(L))$ is $U$ with $g_y(L)$ substituted for $X_y$. Then
\[
 |\kappa!\,\theta_\kappa|\leq|l|^{\,\sum_i\kappa_i(\gamma_l-\sigma_i)-|\iota|(\log_la+1)}.
\]
\end{proposition}
\leanref{Formalism/GaussJetBounds/StatementG.lean}{Formalism.DworkPowerSeriesWeights.LemmaRA}
\leanref{Formalism/GaussJetBounds/GaussJetProofs/LemmaRA.lean}{Formalism.DworkPowerSeriesWeights.lemmaRA_of}
\leanref{Formalism/GaussJetBounds/GaussJetProofs/LemmaRA.lean}{Formalism.DworkPowerSeriesWeights.lemma_R_log_loss}

\begin{proof}
Use the weights $r_i=\rho_l+\sigma_i\geq\rho_l$ on $\iota$.

\emph{Step 1: the exponential factor.} By Lemma~\ref{pd-letter-weights}, $\beta$ satisfies the
hypotheses of Lemma~\ref{pd-letter-substitution}, so $B_M$ has weight at least $1/(l-1)-1$ for the
weights $\rho_l$, and at least $1/(l-1)$ at the monomials not of the form
$X_i^{\,l^e}X_j^{\,l^f}$. Take $n\leq\kappa$. If the coefficient of $\varepsilon B_M$ at
$X^n$ has weight at least $1/(l-1)$ for $\rho_l$, it has for $r$ by
Lemma~\ref{pd-weight-calculus}(2). Otherwise $X^n=X_i^{\,l^e}X_j^{\,l^f}$, and the
hypothesis gives $\sum_kn_k\sigma_k=l^e\sigma_i+l^f\sigma_j\geq1$, so the bound with weight
$1/(l-1)-1$ for $\rho_l$ is a bound with weight $1/(l-1)$ for $r$. So $\varepsilon B_M$ has weight
at least $1/(l-1)$ up to $\kappa$ for $r$, and $\exp(\varepsilon B_M)$ has weight at least $0$ up to
$\kappa$ by Lemma~\ref{pd-weight-calculus}(4).

\emph{Step 2: the integral factor.} The series $U$ is integral, so it has weight at least $0$ for
$\rho_l$. Each $g_y$ has weight at least $\rho_l$ by Lemma~\ref{pd-log-inverse}, and each $L_y$ has
weight at least $\rho_l$ for the weights $\rho_l$ on $\iota$, since $|M_{y,i}|\leq1$. Two
applications of Lemma~\ref{pd-weight-calculus}(1) show that $U(g(L))$ has weight at least $0$ for
$\rho_l$, hence for $r$ by Lemma~\ref{pd-weight-calculus}(2).

\emph{Step 3.} By Lemma~\ref{pd-weight-calculus}(3), the product has weight at least $0$ up to
$\kappa$ for $r$: $|\theta_\kappa|\leq|l|^{-\sum_i\kappa_i(\rho_l+\sigma_i)}$. With
Lemma~\ref{pd-weight-calculus}(6) and $1/(l-1)-\rho_l=\gamma_l$ this is the bound.
\end{proof}

\begin{proposition}\label{pd-one-place-weighted}
Let $K$ be a field of characteristic $0$, complete for a non-archimedean absolute value, $l$ an odd
prime with $0<|l|<1$, $K_0\subset K$ a subfield on which every nonzero absolute value is an integer
power of $|l|$, $\zeta\in K$ a primitive $(l^2-1)$-th root of unity, and $\beta\in K[[Y_0,Y_1]]$ a
letter series for $(K_0,\zeta)$ with $d=2$ (Definition~\ref{pd-letter-series}). Let
$\varepsilon_1\in K$ with $|\varepsilon_1|=1$, $N_0,N_1\geq1$, $\iota$ a finite set and $M=(M_{x,i})$,
$x\in\{0,\dots,N_0-1\}\times\{0,\dots,N_1-1\}$, $i\in\iota$, with $|M_{x,i}|\leq1$, and put
$L_x=\sum_iM_{x,i}X_i$. Let $g$ be a family of series with constant term $0$ in the variables $X_x$
such that substituting $g$ into the one-place logarithm gives $X$, $U$ an integral series and
$\varepsilon=\pm1$. Let $\theta_\kappa$ be the coefficient of $X^\kappa$, $\kappa\in\mathbf N^\iota$,
in
\[
 U(g(L))\cdot\exp\Big(\frac\varepsilon2\sum_x\beta(L_x,\varepsilon_1L_x)\Big).
\]
Then $|\kappa!\,\theta_\kappa|\leq|l|^{\,\rho_l|\kappa|-|\iota|(\log_l|\kappa|+1)}$, with $\log_l0$
read as $0$.
\end{proposition}
\leanref{Formalism/GaussJetBounds/Statement.lean}{Formalism.DworkPowerSeriesWeights.OneLetterR}
\leanref{Formalism/GaussJetBounds/GaussJetProofs/AsmOne.lean}{Formalism.DworkPowerSeriesWeights.oneLetterR_of}
\leanref{Formalism/GaussJetBounds/OneLetterRadius.lean}{Formalism.DworkPowerSeriesWeights.oneLetterR}
\leanref{Formalism/GaussJetBounds/GaussJetProofs/OneLetterRadiusIdentity.lean}{Formalism.DworkPowerSeriesWeights.oneLetterRadiusIdentity}

\begin{proof}
Use the weight $\gamma_l$ on every variable. For $e+f$ odd,
$l^e+l^f\geq l+1$, and $\gamma_l-1/(l+1)=\rho_l$.

\emph{Step 1: the letter series.} By Lemma~\ref{pd-letter-weights}, a coefficient of $\beta$ at a
monomial that is not a bipower has weight at least $1/(l-1)$ for $\rho_l$, hence for $\gamma_l\geq
\rho_l$. A bipower $Y_0^{\,l^e}Y_1^{\,l^f}$ has $e+f$ odd, because $0+e=1+f$
in $\mathbf Z/2\mathbf Z$; its coefficient has weight at least $1/(l-1)-1$ for $\rho_l$, and
$(l^e+l^f)(\gamma_l-\rho_l)\geq1$, so it has weight at least $1/(l-1)$ for $\gamma_l$. So
$\beta$ has weight at least $1/(l-1)$ for $\gamma_l$.

\emph{Step 2.} The series $L_x$ and $\varepsilon_1L_x$ have weight at least $\gamma_l$ for
$\gamma_l$, so by Lemma~\ref{pd-weight-calculus}(1) each $\beta(L_x,\varepsilon_1L_x)$ has weight
at least $1/(l-1)$, and so does the sum times $\varepsilon/2$, since $|2|=1$. Its exponential has
weight at least $0$ by Lemma~\ref{pd-weight-calculus}(4).

\emph{Step 3.} As in Step 2 of the proof of Proposition~\ref{pd-weighted-dwork}, $U(g(L))$ has
weight at least $0$ for $\rho_l$, hence for $\gamma_l$. The product has weight at least $0$, so
$|\theta_\kappa|\leq|l|^{-\gamma_l|\kappa|}$, and Lemma~\ref{pd-weight-calculus}(6) with
$1/(l-1)-\gamma_l=\rho_l$ gives the bound.
\end{proof}

\begin{lemma}\label{pd-dual-radii}
Let $l\geq3$, $d\geq2$ even and $T\subset\mathbf Z/d\mathbf Z$, and for $s\in\mathbf Z/d\mathbf Z$ let
$d_T(s)$ be the least odd $e\geq1$ with $s+e\in T$, or $\infty$. Let $\sigma$ be a dual vector
(Definition~\ref{pd-odd-level-programme}) and $\omega\colon\mathbf Z/d\mathbf Z\to\mathbf R_{\geq0}$
with $\omega(s)\geq l^{-d_T(s)}$ for $s\in T$ with $d_T(s)<\infty$.
\begin{enumerate}
\item $l^i\sigma(s)+\sigma(s+i)\geq1$ for every $s$ and every odd $i\geq1$.
\item If $e+f$ is odd, $s+e=s'+f$, $\tau\geq\sigma(s)$ and $\tau'\geq\sigma(s')$,
  then $l^e\tau+l^f\tau'\geq1$.
\item If $j\geq1$ and $\tau\geq1-l^j\omega(s)$, then $l^{f+j}\omega(s)+l^f\tau\geq1$ for
  every $f\geq0$.
\item If $j\geq1$ and $\tau\geq l^{-j}$, then $l^e\omega(s)+l^{e+j}\tau\geq1$ for every
  $e\geq0$.
\item If $s,s'\in T$, $e+f$ is odd and $s+e=s'+f$, then
  $l^e\omega(s)+l^f\omega(s')\geq1$.
\end{enumerate}
\end{lemma}
\leanref{Formalism/IncomingBounds/IncomingBound.lean}{Formalism.IncomingBounds.IsDualAll}
\leanref{Formalism/IncomingBounds/WeightIdentities.lean}{Formalism.IncomingBounds.lemmaFa}
\leanref{Formalism/GaussJetBounds/Statement.lean}{Formalism.DworkPowerSeriesWeights.ThmARadii}
\leanref{Formalism/GaussJetBounds/GaussJetProofs/OneLetterRadiusIdentity.lean}{Formalism.DworkPowerSeriesWeights.dualRadiiInequalities}

\begin{proof}
(1) Let $i_0<d$ be the residue of $i$ modulo $d$; it is odd because $d$ is even, $s+i=s+i_0$, and
$l^i\sigma(s)\geq l^{i_0}\sigma(s)$. (2) If $f\leq e$, then $i=e-f$ is odd and
$s'=s+i$, so $l^e\tau+l^f\tau'\geq l^f(l^i\sigma(s)+\sigma(s+i))\geq1$ by (1); the other
case is symmetric. (3) and (4) are $l^f(l^j\omega(s)+\tau)\geq l^f$ and
$l^{e+j}\tau\geq l^e$. (5) If $f\leq e$, then $s+(e-f)=s'\in T$ with
$e-f$ odd, so $d_T(s)\leq e-f$ and $l^e\omega(s)\geq l^{e-d_T(s)}\geq
l^f\geq1$; the other case is symmetric.
\end{proof}

\begin{lemma}\label{pd-row-radii}
Let $l\geq3$, $d\geq2$ even, $T\subset\mathbf Z/d\mathbf Z$ and $\nu_0,\nu_1,n\colon\mathbf Z/d\mathbf Z
\to\mathbf N$ with $\nu_1\leq\nu_0\leq n$. Put $\operatorname{cap}(s)=\nu_1(s)$ if $s\in T$ and
$\nu_0(s)$ otherwise; $D(s)=n(s)-\nu_0(s)$ if $s\in T$ and $n(s)$ otherwise; $d_T(s)$ the least odd
$e\geq1$ with $s+e\in T$, or $\infty$; $W(s)$ the largest $\operatorname{cap}(s-i)$ over the odd $i$
with $l^i\leq n(s-i)$, or $0$ if there is none; and
\[
 \operatorname{top}(s)=\sum_{\substack{e\ \mathrm{odd}\\ e<d_T(s)}}\max\Big(\frac{\nu_0(s)}{l^e}-
 \operatorname{cap}(s+e),0\Big)+\sum_{\substack{e\ \mathrm{odd}\\ e>d_T(s)}}\frac{\nu_0(s)}{l^e}.
\]
These are the quantities of Proposition~\ref{pd-orbit-letter-bounds} with $\nu_{s,0}=\nu_0(s)$,
$\nu_{s,1}=\nu_1(s)$ and $|\nu_s|=n(s)$. Let $\sigma$ be optimal for $D$
(Definition~\ref{pd-odd-level-programme}). Let $\iota$ be a finite set with maps
$\operatorname{pos}\colon\iota\to\mathbf Z/d\mathbf Z$, a subset $\iota_{\mathrm{top}}\subset\iota$ of
\emph{top} elements, and a partial map $\operatorname{lab}$ from $\iota\smallsetminus
\iota_{\mathrm{top}}$ to pairs $(s_0,j)$, and let $\kappa\in\mathbf N^\iota$, such that:
\begin{enumerate}
\item every top $i$ has $\operatorname{pos}(i)\in T$, and for $s\in T$ the sum of $\kappa_i$ over the
  top $i$ with $\operatorname{pos}(i)=s$ is $\nu_0(s)$;
\item for every $s$, the sum of $\kappa_i$ over the $i$ that are not top with
  $\operatorname{pos}(i)=s$ is $D(s)$;
\item if $\operatorname{lab}(i)=(s_0,j)$, then $s_0\in T$, $j$ is odd,
  $\operatorname{pos}(i)=s_0+j$, $l^j\leq\nu_0(s_0)$ if $j>0$ and $l^{-j}\leq n(\operatorname{pos}(i))$
  if $j<0$;
\item for every pair $(s_0,j)$, the sum of $\kappa_i$ over the $i$ with $\operatorname{lab}(i)=
  (s_0,j)$ is at most $\operatorname{cap}(s_0+j)$.
\end{enumerate}
Then there is $\omega\colon\mathbf Z/d\mathbf Z\to\mathbf R_{\geq0}$ with $\omega(s)\geq
l^{-d_T(s)}$ for $s\in T$ with $d_T(s)<\infty$, such that, with the \emph{row radii}
\[
 r(i)=\begin{cases}
 \omega(\operatorname{pos}(i)),&i\text{ top},\\
 \max\big(\sigma(\operatorname{pos}(i)),1-l^j\omega(s_0)\big),&\operatorname{lab}(i)=(s_0,j),\ j>0,\\
 \max\big(\sigma(\operatorname{pos}(i)),l^{j}\big),&\operatorname{lab}(i)=(s_0,j),\ j<0,\\
 \sigma(\operatorname{pos}(i)),&\text{otherwise},
 \end{cases}
\]
we have
\[
 \sum_{s\in T}\big(\operatorname{top}(s)-\gamma_lW(s)\big)+F(D)\leq\sum_i\kappa_i\big(\gamma_l-r(i)\big).
\]
\end{lemma}
\leanref{Formalism/GaussJetBounds/StatementA.lean}{Formalism.DworkPowerSeriesWeights.rowRad}
\leanref{Formalism/GaussJetBounds/StatementA.lean}{Formalism.DworkPowerSeriesWeights.ThmAAssembly}
\leanref{Formalism/GaussJetBounds/GaussJetProofs/CapacityLeReachCapacity.lean}{Formalism.DworkPowerSeriesWeights.rowRadiusAssembly}

This is a self-contained computation. Write $r(i)\leq\sigma(\operatorname{pos}(i))+c(i)$ for the
$i$ that are not top, with the \emph{cost} $c(i)=\max(1-l^j\omega(s_0),0)$ or $l^j$ of the label
$(s_0,j)$ of $i$, and $c(i)=0$ for unlabelled $i$. At each $s\in T$ the number $\omega(s)$ is chosen
so that $\operatorname{top}(s)+\nu_0(s)\omega(s)+\sum_{j>0}\operatorname{cap}(s+j)
\max(1-l^j\omega(s),0)\leq\gamma_l\nu_0(s)$, the sum over the odd $j$ that occur in labels. A
label $(s_0,j)$ with $j<0$ has $\operatorname{cap}(s_0+j)\leq W(s_0)$, so these labels cost at most
$W(s_0)\sum_{i\ \mathrm{odd}}l^{-i}=\gamma_lW(s_0)$ in all. The rows that are not top contribute
$\sum_sD(s)(\gamma_l-\sigma(s))=F(D)$, because $\sigma$ is optimal.

\begin{proposition}\label{pd-single-monomial}
Let $K$ be a field of characteristic $0$, complete for a non-archimedean absolute value, $l\geq3$ a
prime with $0<|l|<1$, $d\geq2$ even, $N_0,N_1\geq1$, $K_0$ and $\zeta$ as in
Definition~\ref{pd-letter-series}(3), and $\beta$ a letter series for $(K_0,\zeta)$. Let
$x=(x_{s,r,a})$, $s\in\mathbf Z/d\mathbf Z$, $r,a<N_0$, be a \emph{frame}: $|x_{s,r,a}|\leq1$; its
\emph{row} $(s,r)$ is the vector $(x_{s,r,a})_a$. Let $g$, $U$ and $\varepsilon$ be as in
Proposition~\ref{pd-weighted-dwork}. Let $T\subset\mathbf Z/d\mathbf Z$; the \emph{top rows} are the
rows $(s,0)$ with $s\in T$. Let $\nu_0,\nu_1,n$, $\operatorname{cap}$, $D$, $W$ and
$\operatorname{top}$ be as in Lemma~\ref{pd-row-radii}, $a\geq1$ an integer with $n\leq a$, and
$\operatorname{lab}$ a partial map from the rows that are not top to pairs $(s_0,j)$ such that a row
$(s,k)$ with label $(s_0,j)$ satisfies condition (3) of Lemma~\ref{pd-row-radii} with
$\operatorname{pos}=s$. Let $\kappa$ be an exponent in the variables $X_{s,r,b}$
($s\in\mathbf Z/d\mathbf Z$, $r<N_0$, $b<N_1$) such that:
\begin{enumerate}
\item for $s\in T$, $\sum_b\kappa_{s,0,b}=\nu_0(s)$, and for every $s$ the sum of $\kappa_{s,r,b}$ over
  the rows $(s,r)$ that are not top is $D(s)$;
\item for every pair $(s_0,j)$, the sum of $\kappa_{s,r,b}$ over the rows $(s,r)$ with label
  $(s_0,j)$ is at most $\operatorname{cap}(s_0+j)$;
\item for every $s_0\in T$ and every row $(s,k)$ that is not top, with $s=s_0+j$, $j$ odd,
  $l^j\leq\nu_0(s_0)$ if $j>0$ and $l^{-j}\leq n(s)$ if $j<0$, and label different from $(s_0,j)$:
  $|\sum_ax_{s_0,0,a}^{\,l^j}x_{s,k,a}|\leq|l|$ if $j>0$, and
  $|\sum_ax_{s_0,0,a}x_{s,k,a}^{\,l^{-j}}|\leq|l|$ if $j<0$.
\end{enumerate}
Let $\theta_\kappa$ be the coefficient of $X^\kappa$ in $U(g(x^{\mathsf T}X))\cdot
\exp\big(\varepsilon\sum_{a,b}\beta((\sum_rx_{s,r,a}X_{s,r,b})_s)\big)$, where $x^{\mathsf T}X$
substitutes $\sum_rx_{s,r,a}X_{s,r,b}$ for $X_{s,a,b}$. Then
\[
 |\kappa!\,\theta_\kappa|\leq|l|^{\,\sum_{s\in T}(\operatorname{top}(s)-\gamma_lW(s))+F(D)
 -dN_0N_1(\log_la+1)}.
\]
\end{proposition}
\leanref{Formalism/GaussJetBounds/StatementG.lean}{Formalism.DworkPowerSeriesWeights.ThmAJetA}
\leanref{Formalism/GaussJetBounds/MonoShapeTop.lean}{Formalism.DworkPowerSeriesWeights.thmAJetA}
\leanref{Formalism/GaussJetBounds/GaussJetProofs/TheoremAOneMonomialLogLoss.lean}{Formalism.DworkPowerSeriesWeights.ThmAJetA_of}

\begin{proof}
Put $\iota=\mathbf Z/d\mathbf Z\times\{0,\dots,N_0-1\}\times\{0,\dots,N_1-1\}$ with
$\operatorname{pos}(s,r,b)=s$; the top elements are the $(s,0,b)$ with $s\in T$, and
$\operatorname{lab}(s,r,b)$ is the label of the row $(s,r)$. The matrix $M^x$ of
Lemma~\ref{pd-letter-substitution}(3) has $L=x^{\mathsf T}X$, and it satisfies the hypothesis of
Lemma~\ref{pd-letter-substitution}(2).

\emph{Step 1: the radii.} Let $\sigma$ be an optimal dual vector for $D$
(Lemma~\ref{pd-lp-duality}(1)). Conditions (1) to (4) of Lemma~\ref{pd-row-radii} hold: (1), (2) and
(4) are conditions (1) and (2) here, and (3) holds by hypothesis. That lemma gives $\omega$ and
radii $r(i)\geq0$ with
$\sum_{s\in T}(\operatorname{top}(s)-\gamma_lW(s))+F(D)\leq\sum_i\kappa_i(\gamma_l-r(i))$.

\emph{Step 2: the pairs.} Let $X_i^{\,l^e}X_j^{\,l^f}$ divide $X^\kappa$ with a coefficient
in $B_{M^x}$ that does not have weight at least $1/(l-1)$. We show
$l^e r(i)+l^f r(j)\geq1$. By Lemma~\ref{pd-letter-weights}, $\beta$ satisfies the
hypotheses of Lemma~\ref{pd-letter-substitution}, whose part (2) gives
$\operatorname{pos}(i)\neq\operatorname{pos}(j)$, $\operatorname{pos}(i)+e=\operatorname{pos}(j)+f$
and $e+f$ odd. Every radius is at least the corresponding value of $\sigma$ or of
$\omega$. If neither $i$ nor $j$ is top, apply Lemma~\ref{pd-dual-radii}(2); if both are top,
Lemma~\ref{pd-dual-radii}(5). Suppose $i=(s_0,0,b)$ is top and $j=(s,k,b')$ is not (the other case is
symmetric). By Lemma~\ref{pd-letter-substitution}(3), $b=b'$. Put $j_0=e-f$, which is odd,
with $s=s_0+j_0$. Since $l^e\leq\kappa_i\leq\nu_0(s_0)$ and $l^f\leq\kappa_j\leq n(s)$, the
row $(s,k)$ satisfies the size conditions of hypothesis (3) for $(s_0,j_0)$. If its label were
different from $(s_0,j_0)$, hypothesis (3) and Lemma~\ref{pd-letter-substitution}(3) would give the
coefficient weight at least $1/(l-1)$. So its label is $(s_0,j_0)$, and
$l^e r(i)+l^f r(j)=l^e\omega(s_0)+l^f r(j)\geq1$ by Lemma~\ref{pd-dual-radii}(3)
if $j_0>0$, with $e=f+j_0$, and by Lemma~\ref{pd-dual-radii}(4) if $j_0<0$, with
$f=e-j_0$.

\emph{Step 3.} Every $\kappa_i$ is at most $a$, by condition (1) and $D(s)\leq n(s)\leq a$. Step 2 is
the hypothesis of Proposition~\ref{pd-weighted-dwork} with $\sigma_i=r(i)$, whose bound, with
$|\iota|=dN_0N_1$ and Step 1, is the bound of the statement.
\end{proof}

\begin{proposition}\label{pd-orbit-letter-bounds}
Let $l\geq5$ be a prime, $d\geq2$ even, $N_0<N_1$, $\Theta$ a power series over $\mathbf C_l$ in
variables $X_{s,a,b}$ ($s\in\mathbf Z/d\mathbf Z$, $a<N_0$, $b<N_1$), $\nu_s$ Young diagrams with at most
$N_0$ rows, and $T\subset\mathbf Z/d\mathbf Z$ a set of \emph{top positions}. Suppose that $\Theta$ has
the orbit expansion (Definition~\ref{pd-orbit-expansions}) on $\mathbf Z/d\mathbf Z$ with the single
representative $0$, some $\varepsilon$, a sign and a series $\beta$ that satisfies (B1) to (B4). Let
$a\geq l$ with $|\nu_s|\leq a$ for all $s$ and $2\log_la+2\leq N_0-1$. Put $\gamma_l=l/(l^2-1)$ and,
for $s\in\mathbf Z/d\mathbf Z$: $A(s)=\nu_{s,0}$; $\operatorname{cap}(s)=\nu_{s,1}$ if $s\in T$ and
$\nu_{s,0}$ otherwise; $D(s)=|\nu_s|-A(s)$ if $s\in T$ and $|\nu_s|$ otherwise; $d_T(s)$ the least odd
$e\geq1$ with $s+e\in T$, or $\infty$; $W(s)$ the largest $\operatorname{cap}(s-i)$ over odd $i$ with
$l^i\leq|\nu_{s-i}|$, or $0$ if there is none;
\[
 \operatorname{top}(s)=\sum_{\substack{e\ \mathrm{odd}\\ e<d_T(s)}}\max\Big(\frac{A(s)}{l^e}-
 \operatorname{cap}(s+e),0\Big)+\sum_{\substack{e\ \mathrm{odd}\\ e>d_T(s)}}\frac{A(s)}{l^e},
 \qquad
 \operatorname{base}(s)=\sum_{\substack{e\ \mathrm{odd}\\ e\neq d_T(s)}}\max\Big(\frac{\nu_{s,0}-\nu_{s,1}}{l^e}-
 \operatorname{cap}(s+e),0\Big);
\]
and let $F$ be the function of Definition~\ref{pd-odd-level-programme}. Let
$E_0=dN_0N_1(\log_la+1)$. Then:
\begin{enumerate}
\item $\Theta$ has $\nu$-valuation at least
  $\sum_{s\in T}\big(\operatorname{top}(s)-\gamma_lW(s)\big)+F(D)-E_0$;
\item if $k\geq3$, $w_s>0$ and $0\leq E_s\leq4w_s$ satisfy $D(s)=(k-1)w_s+E_s$ for $s\notin T$ and
  $D(s)=(k-3)w_s+E_s$ for $s\in T$, then $\Theta$ has $\nu$-valuation at least
  \[
   \sum_{s\in T}\operatorname{base}(s)+F\big((k-1)w\big)-\frac{11}l\sum_sw_s-\gamma_l\sum_{s\in T}W(s)-E_0;
  \]
\item if instead $\Theta$ is a series in variables $X_{a,b}$ with the expansion at a split orbit
  (Definition~\ref{pd-orbit-expansions}(4), one place) and $\nu$ is one Young diagram with
  $|\nu|\leq a$, then $\Theta$ has $\nu$-valuation at least
  $|\nu|/(l^2-1)-N_0N_1(\log_la+1)$.
\end{enumerate}
\end{proposition}
\leanref{Formalism/SharpRestriction/Divisibility/JetPolynomialLogs/LocalBoundElevenComparison.lean}{Formalism.MainNotions.DIV.Bridge.DworkPowerSeriesWeights.split_bound_eleven_one}
\leanref{Formalism/SharpRestriction/Divisibility/ProductZerothFirstParts.lean}{Formalism.MainNotions.product_zeroth_first_parts.Loc.ThmA}
\leanref{Formalism/SharpRestriction/Divisibility/ProductZerothFirstParts.lean}{Formalism.MainNotions.product_zeroth_first_parts.Loc.Cor}
\leanref{Formalism/SharpRestriction/Divisibility/ProductZerothFirstParts.lean}{Formalism.MainNotions.product_zeroth_first_parts.Loc1.ThmOne}
\leanref{Formalism/SharpRestriction/Divisibility/JetPolynomialLogs/DworkPowerSeriesWeights.lean}{Formalism.MainNotions.DIV.Bridge.DworkPowerSeriesWeights.thmA_of}
\leanref{Formalism/SharpRestriction/Divisibility/JetPolynomialLogs/LocalBoundElevenComparison.lean}{Formalism.MainNotions.DIV.Bridge.DworkPowerSeriesWeights.corollary_split_eleven_one}
\leanref{Formalism/SharpRestriction/Divisibility/JetPolynomialLogs/DworkPowerSeriesWeights.lean}{Formalism.MainNotions.DIV.Bridge.DworkPowerSeriesWeights.thmOne_of}
\leanref{Formalism/GaussJetBounds/GaussJetProofs/ThmAGaussOf.lean}{Formalism.DworkPowerSeriesWeights.thmAGauss_of}
\leanref{Formalism/GaussJetBounds/GaussJetProofs/Generic.lean}{Formalism.DworkPowerSeriesWeights.theorem_A_generic}
\leanref{Formalism/GaussJetBounds/GaussJetProofs/FrameA.lean}{Formalism.DworkPowerSeriesWeights.theorem_A_frame}
\leanref{Formalism/GaussJetBounds/GaussJetProofs/FrameA.lean}{Formalism.DworkPowerSeriesWeights.frameA_letter}
\leanref{Formalism/GaussJetBounds/GaussJetProofs/ConstantCoeffAevalLinP.lean}{Formalism.DworkPowerSeriesWeights.thmAGen}
\leanref{Formalism/GaussJetBounds/GaussJetProofs/ConstantCoeffAevalLinP.lean}{Formalism.DworkPowerSeriesWeights.jet_pair_linearity}
\leanref{Formalism/GaussJetBounds/GaussJetProofs/ThmAOp.lean}{Formalism.DworkPowerSeriesWeights.thmAOp}
\leanref{Formalism/GaussJetBounds/GaussJetProofs/ThmAOp.lean}{Formalism.DworkPowerSeriesWeights.op_prod_shape}
\leanref{Formalism/GaussJetBounds/GaussJetProofs/OneGauss.lean}{Formalism.DworkPowerSeriesWeights.one_letter_gauss_sum}
\leanref{Formalism/SharpRestriction/Divisibility/JetPolynomialLogs/UniformBoundAllData.lean}{Formalism.MainNotions.DIV.Bridge.OneLetterJetKernel.tieB_of}
\leanref{Formalism/SharpRestriction/Divisibility/JetPolynomialLogs/LetterLogIntegrality.lean}{Formalism.MainNotions.DIV.Bridge.LetterLogIntegrality.integrality_of_letter_log}

\begin{proof}
\emph{Structure.} Part (1) is proved in three layers. First, three reductions: the conditions (B1)
to (B4) are restated over an unramified field with a root of unity (Step 1); the orbit expansion is
rewritten with the inverse of the orbit logarithm (Step 2); and the $\nu$-valuation becomes a bound
on the coefficients of the polynomials $\langle G_I,\Theta\rangle_X$, the \emph{frame jets}
(Step 3). Second, these coefficients are bounded by the value of the frame jet at a generic point,
which is the jet pairing of the minors with $\Theta$ in a \emph{generic frame} (Steps 4 and 5).
Third, the jet pairing in a frame is bounded (Step 6): row operations reduce it to single
monomials, which Proposition~\ref{pd-single-monomial} bounds; this is the Dwork-type estimate, and
it produces the terms $\operatorname{top}(s)$ of the odd levels, the corrections $\gamma_lW(s)$ and
the linear programme $F(D)$. Step 7 assembles (1). Part (2) follows from (1) and the linear programme
comparison of Lemma~\ref{pd-lp-comparison} (Step 8). Part (3) has the same three layers for one
place (Step 9). Write $\Theta=\Theta(X)$, $\varepsilon$, $\sigma=\pm1$ and $\beta$ for the data of the
orbit expansion, and $Y_{s,r,a}$ ($r,a<N_0$) for the variables of the frame jets.

\emph{Step 1: the field of $\beta$.} Let $W\subset\mathbf C_l$ be the field of (B1) and $\zeta\in W$ a
primitive $(l^d-1)$-th root of unity. Every nonzero element of $W$ has absolute value a power of
$|l|$, since $W$ is unramified over $\mathbf Q_l$. By (B1) the coefficients of $\beta$ lie in $W$; by
(B2) with $u=\zeta$, $\beta(\tau_\zeta Y)=\beta(Y)$; (B3) is unchanged; and (B4) gives that
$\exp(\delta_{\mathrm{ev}}\beta)$ and $\exp(\delta_{\mathrm{od}}\beta)$, with the orbit logarithm
substituted, are integral with constant term $1$. Here the orbit logarithm on
$\mathbf Z/d\mathbf Z$ is $\sum_el^{-e}X_{s-2e}^{\,l^{2e}}$ in each of the two families of
variables, which is the logarithm along the shift by $-2$ of Definition~\ref{pd-letter-series}(1).
So $\beta$ is a letter series for $(W,\zeta)$ (Definition~\ref{pd-letter-series}).

\emph{Step 2: the inverse logarithm.} The orbit logarithm $\lambda$ of
Definition~\ref{pd-orbit-expansions}(1) on the variables $X_{s,a,b}$ has linear part the identity,
so it has a compositional inverse $g=(g_{s,a,b})$, a family of series with constant term $0$ with
$g(\lambda(X))=X$ and $\lambda(g(X))=X$ (Lemma~\ref{pd-log-inverse}). Substituting
$g(X)$ for $X$ in the orbit expansion gives
\[
 \Theta(\varepsilon X)=U\big(g(X)\big)\cdot\exp\Big(\sigma\sum_{a,b}\beta\big((X_{s,a,b})_{s\in\mathbf Z/d}
 \big)\Big),
\]
with $U$ integral.

\emph{Step 3: frame jets.} By Definition~\ref{pd-minor-valuation}(4), $\Theta$ has $\nu$-valuation at
least $B$ if and only if every coefficient of every frame jet $\langle G_I,\Theta\rangle_X$ has
absolute value at most $l^{-B}$. Each $G_I$ is homogeneous of degree $\sum_s|\nu_s|$ in the
$X$-variables, so $\langle G_I,\Theta(\varepsilon X)\rangle_X=\varepsilon^{\sum_s|\nu_s|}\langle
G_I,\Theta\rangle_X$, and $|\varepsilon|=1$. It therefore suffices to bound the coefficients of the
frame jets of the series $\Theta(\varepsilon X)$ of Step 2.

\emph{Step 4: a generic point.} Let $\mathcal K\supset\mathbf C_l$ be a complete field with the
absolute value extending that of $\mathbf C_l$ and elements $Y^\circ_y$, $y=(s,r,a)$, such that
$|P(Y^\circ)|$ is the largest absolute value of a coefficient of $P$ for every polynomial $P$ over
$\mathbf C_l$ (the completion of $\mathbf C_l(Y)$ for the Gauss norm). If $k_y\geq1$ are integers,
the monomials $\prod_y(Y^\circ_y)^{k_yn_y}$ are distinct for distinct exponents $n$, so for every
polynomial $P$ and every exponent $n$, $|\mathrm{coeff}_nP|\leq|P((Y^\circ)^k)|$. The hypotheses of
Steps 1 and 2 hold over $\mathcal K$ as well. Let $m\geq0$ be the integer of Step 5 and $x$ the
\emph{generic frame}: $x_{s,r,a}=(Y^\circ_{s,r,a})^{l^m}$ if $s\in T$ and $r=0$, and
$x_{s,r,a}=Y^\circ_{s,r,a}$ otherwise. Then for every column choice $I$ and exponent $n$,
\[
 \big|\mathrm{coeff}_n\langle G_I,\Theta(\varepsilon X)\rangle_X\big|\leq
 \big|\langle G_I,\Theta(\varepsilon X)\rangle_X(x)\big|=
 \Big|\Big\langle\prod_s\prod_c\det\big(X_{s,q,I_{s,c}(q')}\big)_{q,q'<|c|},\ \Theta(\varepsilon\,
 x^{\mathsf T}X)\Big\rangle_X\Big|,
\]
where $x^{\mathsf T}X$ substitutes $\sum_rx_{s,r,a}X_{s,r,b}$ for $X_{s,a,b}$. The equality is the
adjunction of the jet pairing with a linear substitution: $\langle P(xX),G\rangle=\langle
P,G(x^{\mathsf T}X)\rangle$, where $xX$ substitutes $\sum_ax_{s,r,a}X_{s,a,b}$ for $X_{s,r,b}$, as in
the proof of Lemma~\ref{is-fock-facts}(3). We call $(x_{s,r,a})_a$ the \emph{row} $(s,r)$ of the
frame, and the rows $(s,0)$ with $s\in T$ its \emph{top rows}.

\emph{Step 5: labels.} For each $s$, the \emph{labels} at $s$ are the pairs $(s_0,j)$ with $s_0\in T$,
$j$ odd, $s=s_0+j$ in $\mathbf Z/d$, $l^j\leq\nu_{s_0,0}$ if $j>0$ and $l^{-j}\leq|\nu_s|$ if $j<0$.
Since $j$ determines $s_0=s-j$ and $l^{|j|}\leq a$, there are at most $\log_la+1$ labels at $s$, and
the hypothesis $2\log_la+2\leq N_0-1$ gives (number of labels at $s$)$+[s\in T]\leq N_0$. Number the
labels at $s$ by $J_s$, and let $m$ be the largest $-j$ over all labels. To the $i$-th label
$(s_0,j)$ at $s$ attach the row $e_s(i)=i+[s\in T]$ of the frame at $s$ and the vector
$\lambda_{s,i}(a)=(Y^\circ_{s_0,0,a})^{l^{m+j}}$. Then $|Y^\circ_y|=1$ for all $y$; if $j>0$,
$\lambda_{s,i}=x_{s_0,0,\cdot}^{\,l^{j}}$, and if $j<0$, $\lambda_{s,i}^{\,l^{-j}}=x_{s_0,0,\cdot}$; and the
$J_s\times J_s$ matrix $\big(\sum_a\lambda_{s,i}(a)\,x_{s,e_s(i'),a}\big)_{i,i'}$ has determinant of
absolute value $1$. Indeed, expand the determinant over permutations $\pi$ and indices $a_i$: the term
of $(\pi,(a_i))$ is $\pm\prod_i(Y^\circ_{s_0(i),0,a_i})^{l^{m+j(i)}}Y^\circ_{s,e_s(\pi i),a_i}$, with
$(s_0(i),j(i))$ the $i$-th label. Here $s_0(i)\neq s$, since $j(i)$ is odd and $d$ even, and distinct
labels have distinct pairs $(s_0(i),m+j(i))$, so the exponent of each $Y^\circ_{s_0,0,a}$ is a sum of
distinct powers of $l$, from which the $a_i$ can be read off. Take $\pi$ the identity and the $a_i$
distinct, which is possible since $|J_s|\leq N_0$; then $\pi$ is also read off from the factors
$Y^\circ_{s,e,a_i}$, so this monomial has coefficient $\pm1$. The determinant is thus a polynomial in
the $Y^\circ$ with integer coefficients, one of them $\pm1$, and its absolute value is $1$ by the
choice of $Y^\circ$.

\emph{Step 6: the bound in a frame.} Let $x$ be a frame with entries of absolute value at most $1$,
labelled as in Step 5: the row $e_s(i)$ carries the $i$-th label $(s_0,j)$ at $s$, with a vector
$\lambda_{s,i}$ of absolute value at most $1$ such that $\lambda_{s,i}=x_{s_0,0}^{\,l^j}$ if $j>0$ and
$\lambda_{s,i}^{\,l^{-j}}=x_{s_0,0}$ if $j<0$ (coordinatewise), and the determinants of Step 5 have
absolute value $1$. We show
\[
 \Big|\Big\langle\prod_s\prod_c\det\big(X_{s,q,I_{s,c}(q')}\big),\ \Theta(\varepsilon\,x^{\mathsf T}X)
 \Big\rangle_X\Big|\leq|l|^{\,\sum_{s\in T}(\operatorname{top}(s)-\gamma_lW(s))+F(D)-dN_0N_1(\log_la+1)}.
\]
(a) \emph{Row operations.} Fix $s$. The matrix $A=(\sum_a\lambda_{s,i}(a)x_{s,e_s(i'),a})_{i,i'}$ has
entries of absolute value at most $1$ and determinant of absolute value $1$, so $A^{-1}$ has entries of
absolute value at most $1$. Put $z_{i'}=\sum_{i''}(A^{-1})_{i'',i'}\,x_{s,e_s(i'')}$, so that
$\sum_a\lambda_{s,i}(a)z_{i',a}=\delta_{ii'}$, and $c_{k,i}=\sum_a\lambda_{s,i}(a)x_{s,k,a}$. Let $x'$
be the frame with the same top rows, $x'_{s,e_s(i)}=z_i$, and $x'_{s,k}=x_{s,k}-\sum_ic_{k,i}z_i$ for
the other rows that are not top. Then $|x'|\leq1$; $\sum_a\lambda_{s,i}(a)x'_{s,k,a}=0$ for every row
$k\neq e_s(i)$ that is not top; and $x_{s,k}=\sum_{k'}T_s(k,k')\,x'_{s,k'}$ for the rows $k$ that are
not top, the sum over the rows $k'$ that are not top, where $T_s(k,k')$ is $c_{k,i}$ if $k'=e_s(i)$,
plus $1$ if $k'=k$ is not a row $e_s(i)$; so $|T_s|\leq1$. Consequently, if $s_0\in T$, $(s_0,j)$ is a
label at $s=s_0+j$, carried by the row $e_s(i)$, and $(s,k)$ is another row that is not top, then
$\sum_ax_{s_0,0,a}^{\,l^j}x'_{s,k,a}=\sum_a\lambda_{s,i}(a)x'_{s,k,a}=0$ if $j>0$; and if $j<0$,
$\sum_ax_{s_0,0,a}x'^{\,l^{-j}}_{s,k,a}=\sum_a(\lambda_{s,i}(a)x'_{s,k,a})^{l^{-j}}$ is congruent modulo
$l$ to $(\sum_a\lambda_{s,i}(a)x'_{s,k,a})^{l^{-j}}=0$. So the frame $x'$, whose top rows are those of
$x$, satisfies hypothesis (3) of Proposition~\ref{pd-single-monomial}, with the label of $e_s(i)$
the $i$-th label at $s$.

(b) \emph{Change of frame.} Let $\tilde T_s(r,k')$ be $T_s(r,k')$ for the rows $r,k'$ that are not
top, $1$ for $r=k'$ a top row, and $0$ otherwise. Then $x^{\mathsf T}X=x'^{\mathsf T}(\tilde T^{\mathsf T}X)$,
and the adjunction of Step 4 gives
\[
 \Big\langle\prod_{s,c}\det\big(X_{s,q,I_{s,c}(q')}\big),\Theta(\varepsilon x^{\mathsf T}X)\Big\rangle_X
 =\Big\langle\prod_{s,c}\det\big((\tilde TX)_{s,q,I_{s,c}(q')}\big),\Theta(\varepsilon x'^{\mathsf T}X)
 \Big\rangle_X,
\]
where $(\tilde TX)_{s,r,b}=\sum_{k'}\tilde T_s(r,k')X_{s,k',b}$: the top rows are unchanged, and the
other rows are combinations, with coefficients of absolute value at most $1$, of the rows that are not
top.

(c) \emph{Monomials.} The polynomial $P'=\prod_{s,c}\det((\tilde TX)_{s,q,I_{s,c}(q')})_{q,q'<|c|}$ has
coefficients of absolute value at most $1$. By the Cauchy-Binet formula each factor is a combination of
the minors $\det(X_{s,k_q,I_{s,c}(q')})$ with distinct rows $k_0,\dots,k_{|c|-1}$, and for $s\in T$
these contain the top row, since the top row of $\tilde T_s$ is a unit vector and the other rows of
$\tilde T_s$ vanish in the top column. So every monomial $X^\kappa$ of $P'$ has: degree $\nu_{s,0}$ in
the top row at $s\in T$; total degree $|\nu_s|$ at $s$, hence total degree $D(s)$ in the rows at $s$
that are not top; and degree at most $\operatorname{cap}(s)$ in each row at $s$ that is not top,
because such a row occurs at most once in each column, and for $s\in T$ only in the columns of length
at least $2$. Each label is carried by one row, so condition (2) of
Proposition~\ref{pd-single-monomial} holds. By the ultrametric inequality the jet pairing of (b) is
at most the largest $|\kappa!\,\theta_\kappa|$ over these $\kappa$, where $\theta_\kappa$ is the
coefficient of $X^\kappa$ in $\Theta(\varepsilon x'^{\mathsf T}X)$.

(d) \emph{One monomial.} By Steps 1 and 2, $\Theta(\varepsilon x'^{\mathsf T}X)=U(g(x'^{\mathsf T}X))\cdot
\exp(\sigma\sum_{a,b}\beta((\sum_rx'_{s,r,a}X_{s,r,b})_s))$ with $\beta$ a letter series. By (a) and
(c), Proposition~\ref{pd-single-monomial} applies, with the sign $\sigma$, to $x'$, the labels, $\nu_0(s)=\nu_{s,0}$,
$\nu_1(s)=\nu_{s,1}$, $n(s)=|\nu_s|$ and every $\kappa$ of (c), and gives the bound.

\emph{Step 7: part (1).} By Steps 3 to 6, every coefficient of every frame jet of $\Theta$ has absolute
value at most $l^{-B}$ with $B=\sum_{s\in T}(\operatorname{top}(s)-\gamma_lW(s))+F(D)-
dN_0N_1(\log_l\lfloor a\rfloor+1)$; Steps 4 to 6 are applied with $\lfloor a\rfloor$, which satisfies the
same hypotheses as $a$ since $|\nu_s|$ is an integer. Since $\log_l\lfloor a\rfloor\leq\log_la$, this
gives (1).

\emph{Step 8: part (2).} Termwise, $\operatorname{base}(s)\leq\operatorname{top}(s)$: for odd $e<d_T(s)$
because $\nu_{s,0}-\nu_{s,1}\leq A(s)$, and for odd $e>d_T(s)$ because $\operatorname{cap}\geq0$ gives
$\max((\nu_{s,0}-\nu_{s,1})/l^e-\operatorname{cap}(s+e),0)\leq A(s)/l^e$. By
Lemma~\ref{pd-lp-comparison}, $F((k-1)w)-\frac{11}l\sum_sw_s\leq F(D)$. So the bound of (2) is at most
that of (1), and (2) follows from (1).

\emph{Step 9: part (3).} Write $\Theta(\varepsilon_0\lambda(X))=U(X)\exp\big(\frac{\sigma}2\sum_{a,b}
\beta(Y_{a,b},\varepsilon_1Y_{a,b})\big)|_{Y=\lambda(X)}$ for the expansion at a split orbit, with
$\lambda$ the one-place logarithm. As in Step 1, with $d=2$, $\beta$ is a letter series for $(W,\zeta)$,
$\zeta$ a primitive $(l^2-1)$-th root of unity. As in Step 2, with $g$ the inverse of the one-place
logarithm (Lemma~\ref{pd-log-inverse}), $\Theta(\varepsilon_0X)=U(g(X))\exp(\frac\sigma2\sum_{a,b}
\beta(X_{a,b},\varepsilon_1X_{a,b}))$. As in Step 3, it suffices to bound the coefficients of
$\langle G_I,\Theta(\varepsilon_0X)\rangle_X$, and as in Step 4, with the frame $x=Y^\circ$, each is at
most $|\langle\prod_c\det(X_{q,I_c(q')}),\Theta(\varepsilon_0x^{\mathsf T}X)\rangle_X|$. The product of
minors has integer coefficients and is homogeneous of degree $|\nu|$, so by the ultrametric inequality
this is at most the largest $|\kappa!\,\theta_\kappa|$ over exponents with $|\kappa|=|\nu|$, where
$\theta_\kappa$ is the coefficient of $X^\kappa$ in $\Theta(\varepsilon_0x^{\mathsf T}X)$. By
Proposition~\ref{pd-one-place-weighted}, applied with $\iota=\{0,\dots,N_0-1\}\times\{0,\dots,N_1-1\}$
and $L_{a,b}=\sum_rx_{r,a}X_{r,b}$, this is at most
$|l|^{|\nu|/(l^2-1)-N_0N_1(\log_l|\nu|+1)}$. Since $|\nu|\leq\lfloor a\rfloor\leq a$, this gives (3).
\end{proof}

\subsection{Served lattices and Hecke chains}\label{subsection-pd-served-lattices-hecke-chains}

\begin{definition}\label{pd-served-lattice}
Let $D$ be a definite-side datum with ranks $k,m,t,n$ and $N(j)=k+j(t+1)$, Galois datum with group
$A$, top classes $T$ and coefficient field $L$. Let it carry size data
(Definition~\ref{pd-size-data}) with scale $(x,\varepsilon,J)$, exponents $a_t,b_t$, width $\eta$,
local terms $F_j(w)$, local errors $E_{u,j}(w)$ and the derived quantities $a,b',w'',X'',p_0$; gain
data (Definition~\ref{pd-gain-data}) with representatives, numbers $N(p,t_0)$, served
representatives $V(p)$ and a constant $K_{\mathrm{Pr}}$ of $K_G$; and divisibility data. Let $c_G$
and $\mathrm{fill}_j$ be as in Definition~\ref{pd-lattice-instance}, $r(\varepsilon)$ as in
Definition~\ref{pd-gain-link}, and $N=Q_Sqd$ the conductor of the host. Let $p$ be a prime, with
$h_p$, $\ell_p$ and $\operatorname{pos}_p$ as in Definition~\ref{pd-orbit-places}, and put
$\gamma_p=p/(p^2-1)$. For a representative $t_0$ at $p$, let $T(p,t_0)$ be the set of top classes
with representative $t_0$; an element $s\in T(p,t_0)$ is at an \emph{even position} if
$s=t_0+ih_p$ for an even $i$ with $0\leq i<\ell_p$, and at an \emph{odd position} otherwise. For a
lift $j$, let $g_0(p,j,t_0)$ be the local gain of Definition~\ref{pd-served-selection}, with its
term $\delta(p,t_0)$, and put $g_0^+(p,j,t_0)=\max(g_0(p,j,t_0),0)$.
\begin{enumerate}
\item The \emph{unit gain} is
  \[
   g_U(p,j,t_0)=\gamma_p\sum_{s\in T(p,t_0)}(a_s-2c_G)-\sum_{s}\sum_{e}
   \min\Big(\frac{a_s-2c_G}{p^e},\ \mathrm{fill}_j\big(\operatorname{pos}_p((0,s),e)\big)+c_G\Big)
   -\delta(p,t_0),
  \]
  where $s$ runs over the elements of $T(p,t_0)$ at odd positions and $e$ over the odd integers
  $e\geq1$.
\item The \emph{Hecke loss} is
  \[
   \mathrm{Pr}(p,j,t_0)=\sum_{s}\big(\mathrm{fill}_j((0,s))-b_s\big)+K_{\mathrm{Pr}}\,\ell_p
   \big(w''+\eta+c_G+N(m)\big)
  \]
  if $p\geq p_0$, with $s$ as in (1), and
  $\mathrm{Pr}(p,j,t_0)=r(\varepsilon)\big(N(p,t_0)(2b'+c_G)+\ell_p(w''+\eta+c_G+N(m))\big)$ if
  $p<p_0$.
\item The surplus $S(p,t_0)$ of Definition~\ref{pd-gain-data}(f) is the integer
  \[
   S(p,t_0)=\Big\lceil\sum_{j=1}^{m}\max\big(0,\ g_U(p,j,t_0)-\mathrm{Pr}(p,j,t_0)-g_0^+(p,j,t_0)
   \big)\Big\rceil\geq0 .
  \]
\item The \emph{gain} of the $j$-th lift at $p$ is $\mathrm{gain}_j(p)=0$ if $p>X''$ or if $p$ is
  of inert class with $h_p=0$. Otherwise
  \[
   \mathrm{gain}_j(p)=\sum_{t_0\in V(p)}\max\big(g_U(p,j,t_0)-\mathrm{Pr}(p,j,t_0),\
   g_0^+(p,j,t_0)\big)+\sum_{t_0\notin V(p)}g_0^+(p,j,t_0),
  \]
  the second sum over the representatives at $p$ outside $V(p)$. In both cases
  $\mathrm{gain}_j(p)\geq0$.
\item The \emph{served gain} is $\mathrm{gain}^{\mathrm s}_j(p)=\mathrm{gain}_j(p)$ if $p$ is a
  signal prime or of inert class, $p\nmid Ng_0g_1$, and $p$ is not of inert class with $h_p$ of
  odd order at least $3$; otherwise $\mathrm{gain}^{\mathrm s}_j(p)=0$. The \emph{total local
  gain} is $\bar g_0(p,j)=0$ if $p$ is of inert class with $h_p$ of odd order at least $3$, and
  otherwise $\bar g_0(p,j)=\sum_{t_0}g_0^+(p,j,t_0)$, the sum over all representatives at $p$.
\item Let $e_p$ be the Weyl exponent of the gain data at $p$
  (Definition~\ref{pd-gain-data}(5)). The gain data also attach to $p$ and $t_0$ a real number
  $t(p,t_0)$ and an integer $D_p\geq0$ determined by $K_G$, the size data, $\ell_p$ and the
  positions $P(p,t_0)$. The \emph{served exponent} is $\tau_p=0$ if $V(p)=\emptyset$, and otherwise
  $\tau_p=e_p+m\sum_{t_0\in V(p)}\big(\max(\lceil t(p,t_0)\rceil,0)+D_p+1\big)$.
\end{enumerate}
Let $u$ be a word and $w$ a place admissible for $u$ (Definition~\ref{pd-lattice-instance}), with
$p=p_w$, and suppose that $D$ carries a lattice instance, with map $t\mapsto v(t)$ and submodule
$V_{\mathrm{ram}}$ at $(u,w)$, and a comparison $\operatorname{cmp}$. Let $\mathcal R$ and $M$ be
the ring and the integral module of the local model $G(u,w)$, and $y_v$ its Hecke operators
(Definition~\ref{pd-local-model}). For a set $V'$ of representatives at $p$, $y_{V'}$ is the
composite of the operators $y_{v(t_0)}^m$, $t_0\in V'$, in a fixed order. The \emph{served local
lattice} $N_P(u,w)$ is the intersection of $p^{\max(\lceil\tau_p\rceil,0)}M$, $V_{\mathrm{ram}}$ and
the $\mathcal R$-span of the functions
\[
 y_{V'}\big(p^{S(V(p)\smallsetminus V')}F\big),\qquad V'\subset V(p),\ F\in M,
\]
where $S(B)=\sum_{t_0\in B}S(p,t_0)$.

Now let $u$ be a word and $w$ any finite place of $L$, with $p=p_w$. A vector $x\in V_{u,0}$
\emph{lies in the served local lattice at} $w$ if either $w$ is admissible for $u$ and
$\operatorname{cmp}_{u,w}(x)\in N_P(u,w)$, or $V(p)=\emptyset$, $x$ lies in the $w$-local span of
$\Lambda_{u,0}$ and its image lies in $M_{u,\mathrm{ram}}$. The \emph{base exponent} at $(u,w)$ is
\[
 \mathrm{low}(u,w)=\sum_{j=1}^{m}\big(F_j(w)-E_{u,j}(w)\big)+v_p\big(D_\Gamma(u)\big).
\]
\end{definition}
\leanref{Formalism/SharpRestriction/GrowthBounds/IsResidueClassZero.lean}{Formalism.MainNotions.G.gU}
\leanref{Formalism/SharpRestriction/GrowthBounds/IsResidueClassZero.lean}{Formalism.MainNotions.G.Pr}
\leanref{Formalism/SharpRestriction/GrowthBounds/IsResidueClassZero.lean}{Formalism.MainNotions.G.Sv}
\leanref{Formalism/SharpRestriction/GrowthBounds/IsResidueClassZero.lean}{Formalism.MainNotions.G.gainP}
\leanref{Formalism/SharpRestriction/GrowthBounds/WeilJacquetShalikaLaw.lean}{Formalism.MainNotions.WeilJacquetShalikaLaw.gainS}
\leanref{Formalism/SharpRestriction/GrowthBounds/WeilJacquetShalikaLaw.lean}{Formalism.MainNotions.WeilJacquetShalikaLaw.g0S}
\leanref{Formalism/SharpRestriction/GrowthBounds/IsResidueClassZero.lean}{Formalism.MainNotions.G.tauP}
\leanref{Formalism/SharpRestriction/GrowthBounds/IsResidueClassZero.lean}{Formalism.MainNotions.G.yProd}
\leanref{Formalism/SharpRestriction/GrowthBounds/IsResidueClassZero.lean}{Formalism.MainNotions.G.NP}
\leanref{Formalism/SharpRestriction/GrowthBounds/WeilJacquetShalikaLaw.lean}{Formalism.MainNotions.WeilJacquetShalikaLaw.InNP}
\leanref{Formalism/SharpRestriction/GrowthBounds/WeilJacquetShalikaLaw.lean}{Formalism.MainNotions.WeilJacquetShalikaLaw.lowS}

\begin{definition}\label{pd-formal-group-numbers}
Let $p\geq5$ be a prime and take a formal group system at $p$
(Definition~\ref{pd-formal-group-system}), with rank $k$, integers $r$, $t$ and $m_*$, uniformizer
$\varpi$ of $\mathcal R$, valuation $e_{\mathcal R}$ of $p$ in $\mathcal R$, integral module $M$ of
the model, Hecke operators $y_v$ (Definition~\ref{pd-local-model}(4)), and fills with top positions
$T_v$ and weights $\lambda_0(v,s)$, $\lambda_j(v,s)$ and $w(v,j,s)$. Put $N(j)=k+j(t+1)$ and
$\gamma_p=p/(p^2-1)$, and for an integer $y$ let $\lg_p(y)=0$ if $y\leq0$ and
$\lg_p(y)=\lfloor\log_py\rfloor+1$ if $y>0$. Let $v$ be an orbit and $n'=n'_v$; a position
$s\in\mathbf Z/n'\mathbf Z$ is \emph{even} or \emph{odd} as its representative in
$\{0,\dots,n'-1\}$ is. For a position $s$ and an integer $j\geq1$, let $\alpha(s)=\lambda_0(v,s)_1-
\lambda_0(v,s)_2$ if $s\in T_v$ and $\alpha(s)=0$ otherwise. Let $\nu_j(s)=\lambda_j(v,s)+w(v,j,s)$ be
the natural weight of $\lambda_j(v,s)$, let $|\nu_j(s)|$ be the sum of its entries, and let its
\emph{cap} $c_j(s)$ be its second entry if $s\in T_v$ and its first entry otherwise. The
\emph{digit error} is
\[
 d_j=\sum_{s}\Big(N(j)\lg_p\alpha(s)+N(j-1)N(j)\lg_p\big(|\nu_j(s)|-\alpha(s)\big)\Big)
 +\sum_{s\in T_v}\big(\lg_p\alpha(s)\big)^2,
\]
the first sum over all positions. For $s\in T_v$ let $\Delta(s)\subset\{0,\dots,n'-1\}$ be the set
of residues of $s'-s$ for the $s'\in T_v$ of the other parity.
\begin{enumerate}
\item The \emph{orbit gain} is
  \[
   g_B(v,j)=\gamma_p\sum_{s\in T_v}\alpha(s)-\sum_{s}\sum_{e\geq0}\min\Big(\frac{\alpha(s)}{p^{2e+1}},\
   c_j(s+2e+1)\Big)-d_j,
  \]
  the middle sum over the odd $s\in T_v$.
\item The \emph{orbit local gain} is
  \[
   g_A(v,j)=\sum_{s\in T_v}\ \sum_{e\geq0,\ 2e+1\notin\Delta(s)}\max\Big(\frac{\alpha(s)}{p^{2e+1}}
   -c_j(s+2e+1),\ 0\Big)-d_j .
  \]
\item The \emph{orbit loss} is
  \[
   \mathrm{Pr}_B(v,j)=\sum_{s\ \mathrm{odd}}\Big(r\,w(v,j,s)+\sum_{a=k-r+1}^{k}\lambda_0(v,s)_a\Big)
   +\frac{n'rj(t+1)}{2}.
  \]
\item For a function $S$ from the orbits to $\mathbf R$ and a set $V$ of orbits, the \emph{product
  module} $I(S,V)$ is the $\mathcal R$-span, in the function space of the model, of the functions
  $y_{v_1}^{m_*}\cdots y_{v_l}^{m_*}(\varpi^nF)$, where $v_1,\dots,v_l$ are distinct orbits forming
  a set $B\subset V$, $F\in M$ and $n=\sum_{v\in V\smallsetminus B}\lceil S(v)e_{\mathcal R}\rceil$.
\end{enumerate}
\end{definition}
\leanref{Formalism/SharpRestriction/GrowthBounds/FormalGroupSystem.lean}{Formalism.MainNotions.FormalGroupSystem.de}
\leanref{Formalism/SharpRestriction/GrowthBounds/FormalGroupSystem.lean}{Formalism.MainNotions.FormalGroupSystem.gB}
\leanref{Formalism/SharpRestriction/GrowthBounds/FormalGroupSystem.lean}{Formalism.MainNotions.FormalGroupSystem.gA}
\leanref{Formalism/SharpRestriction/GrowthBounds/FormalGroupSystem.lean}{Formalism.MainNotions.FormalGroupSystem.Pr}
\leanref{Formalism/SharpRestriction/GrowthBounds/FormalGroupSystem.lean}{Formalism.MainNotions.FormalGroupSystem.Sys.Iprod}

\begin{definition}\label{pd-hecke-chain-operators}
Let $p\geq5$ be a prime and take a formal group system at $p$ (Definition~\ref{pd-formal-group-system}),
with rank $k$, integers $r$, $t$ and $m_*$, uniformizer $\varpi$ of $\mathcal R$, valuation
$e_{\mathcal R}$ of $p$ in $\mathcal R$, function spaces $\mathrm{Fm}_J$, edge operators $T_{v,i}$ and
the numbers of Definition~\ref{pd-formal-group-numbers}. An element $x\in\mathcal K$ \emph{has valuation
at least} $g\in\mathbf R$ if $x\in\varpi^{\lceil ge_{\mathcal R}\rceil}\mathcal R$. Let $v$ be an orbit, put
$Q_v=p^{n'_v/2}$ and $q_x=Q_v^{\,t-2x}$ for $0\leq x\leq t$, and let $\operatorname{vm}_v(J)$ be as in
Definition~\ref{pd-formal-group-system}(6).
\begin{enumerate}
\item For $J\geq1$ and $0\leq i\leq r$ let
  \[
   N'_v(J,i)=\sum_{s\ \mathrm{even}}\sum_{a=1}^{i}\lambda_0(v,s)_a+i\sum_{s\ \mathrm{odd}}w(v,J,s).
  \]
  The \emph{weighted edge operator} $T^\natural_{v,i}$ on $\mathrm{Fm}_J$ is the identity if $i=0$ and
  $p^{N'_v(J,i)}T_{v,i}$ if $i\geq1$. For every $J\geq0$ and $i\geq0$, the \emph{normalized edge
  operator} on $\mathrm{Fm}_J$ is $E_{v,i}=Q_v^{-i(N(J)-i)}T_{v,i}$.
\item The \emph{even top sum} is $\alpha_{\mathrm{ev}}(v)=\sum_s\alpha(s)$, the sum over the even
  $s\in T_v$.
\item Let $J\geq1$, and let $u_1,\dots,u_J$ be nonzero and $m_1,\dots,m_J$ arbitrary elements of
  $\mathcal K$. For $0\leq l\leq r$ put
  \[
   c_l(J)=p^{N'_v(J,r)-N'_v(J,r-l)}\,Q_v^{\,l(N(J)-2r+l)}\,h_l\big(m_{l'}q_x:1\leq l'\leq J,\
   0\leq x\leq t\big),
  \]
  where $h_l$ is the complete homogeneous symmetric polynomial of degree $l$ in the $J(t+1)$ entries
  of the family. The \emph{orbit operator} on $\mathrm{Fm}_J$ is
  \[
   \mathcal H_{v,J}=\sum_{l=0}^{r}(-1)^l\,c_l(J)\,(u_1\cdots u_J)^{-(r-l)}\,T^\natural_{v,r-l}.
  \]
\end{enumerate}
\end{definition}
\leanref{Formalism/SharpRestriction/GrowthBounds/FormalGroupSystem.lean}{Formalism.MainNotions.FormalGroupSystem.Sys.vGeK}
\leanref{Formalism/SharpRestriction/GrowthBounds/FormalGroupSystem.lean}{Formalism.MainNotions.FormalGroupSystem.Sys.qh}
\leanref{Formalism/SharpRestriction/GrowthBounds/FormalGroupSystem.lean}{Formalism.MainNotions.FormalGroupSystem.Sys.Tt}
\leanref{Formalism/SharpRestriction/GrowthBounds/FormalGroupSystem.lean}{Formalism.MainNotions.FormalGroupSystem.Sys.ecal}
\leanref{Formalism/SharpRestriction/GrowthBounds/FormalGroupSystem.lean}{Formalism.MainNotions.FormalGroupSystem.Sys.alEven}
\leanref{Formalism/SharpRestriction/GrowthBounds/FormalGroupSystem.lean}{Formalism.MainNotions.FormalGroupSystem.Sys.FF}
\leanref{Formalism/SharpRestriction/GrowthBounds/FormalGroupSystem.lean}{Formalism.MainNotions.FormalGroupSystem.Sys.cl}
\leanref{Formalism/SharpRestriction/GrowthBounds/FormalGroupSystem.lean}{Formalism.MainNotions.FormalGroupSystem.Sys.Ttu}
\leanref{Formalism/SharpRestriction/GrowthBounds/FormalGroupSystem.lean}{Formalism.MainNotions.FormalGroupSystem.Sys.Zop}

\begin{definition}\label{pd-kernel-orbit-index}
Take a formal group system at $p$ (Definition~\ref{pd-formal-group-system}), let $j<m_*$, and let $I$ be
a column choice into $N(j+1)$ for the diagrams $\nu_{v,s}$ of the natural weights of $\lambda_{j+1}(v,s)$
at all positions.
\begin{enumerate}
\item For a function $i_0$ from the orbits to $\{0,\dots,r\}$ let $G^{i_0}_I$ be the polynomial of
  condition~(i) of Definition~\ref{pd-formal-group-system}. For an orbit $v$ and $0\leq i\leq r$,
  $G^{(v),i}_I$ is the product of the factors of $G^{i_0}_I$ at the positions of $v$, for any $i_0$ with
  $i_0(v)=i$. Thus $G^{i_0}_I=\prod_vG^{(v),i_0(v)}_I$, and $G^{(v),0}_I$ is the polynomial
  $G^{(v)}_I$ of the definition of a lower bound.
\item Let $v$ be an orbit, $0\leq i\leq r$ and $g\in\mathbf R$. The system \emph{has lower bound $g$ at
  $(j,v)$ with index $i$} if for every series $U$ in the variables $X_e,X'_e$ such that $U(\ell(X))$ has
  coefficients in $\mathcal R_c$, and every column choice $I$ as above, every coefficient of
  $\langle G^{(v),i}_I,\Theta^Z_U\rangle_Z$ has $v\geq g$, where $\Theta^Z_U$ is the series of the
  definition of a lower bound. With index $0$ this is a lower bound $g$ at $(j,v)$.
\end{enumerate}
\end{definition}
\leanref{Formalism/SharpRestriction/GrowthBounds/FormalGroupSystem.lean}{Formalism.MainNotions.FormalGroupSystem.Sys.PIv}
\leanref{Formalism/SharpRestriction/GrowthBounds/FormalGroupSystem.lean}{Formalism.MainNotions.FormalGroupSystem.Sys.colExp}
\leanref{Formalism/SharpRestriction/GrowthBounds/MultisetValuation/AuxUOfFirstTheoremPlaceBound.lean}{Formalism.MainNotions.FormalGroupSystem.Pf.PIvS}
\leanref{Formalism/SharpRestriction/GrowthBounds/MultisetValuation/AuxUOfFirstTheoremPlaceBound.lean}{Formalism.MainNotions.FormalGroupSystem.Pf.LowAtI}

\begin{definition}\label{pd-kernel-orbit-series}
Let $p$ be a prime and take a formal group system at $p$ (Definition~\ref{pd-formal-group-system}). Let
$v$ be an orbit and $j<m_*$, write $Z$ for the variables $Z_{v',s,a,b}$ of the orbit coordinates at level
$j$, and let $K'$ be a field of characteristic $0$.
\begin{enumerate}
\item For a family $u=(u_s)$ of elements of $K'$ indexed by the positions of $v$, put
  $S_{s,a,b}=\sum_{l\geq0}(-1)^lZ_{v,s-2l,a,b}^{\,p^{2l}}$ and
  \[
   H_u=\sum_su_s\sum_{a,b}S_{s,a,b}^{\,p}\,S_{s+1,a,b}.
  \]
\item Let $R'\subset K'$ be a subring and $H$ a power series over $K'$ with constant term $0$. A power
  series $F$ over $K'$ in the variables $Z$, or in some of them, \emph{lies in the class of $H$ over
  $R'$} if there are power series $U_0,U_1,\dots$ with coefficients in $R'$ such that for every exponent
  $\alpha$ the coefficient of $Z^\alpha$ in $F$ is the sum over $0\leq i\leq|\alpha|$ of $p^{-i}$ times the
  coefficient of $Z^\alpha$ in $H^iU_i$.
\item For $\pi\in K'$, $F(\pi Z)$ is the series $F$ with $\pi Z_{v',s,a,b}$ substituted for each
  variable $Z_{v',s,a,b}$.
\end{enumerate}
\end{definition}
\leanref{Formalism/SharpRestriction/GrowthBounds/FormalGroupSystem.lean}{Formalism.MainNotions.FormalGroupSystem.Sys.sSA}
\leanref{Formalism/SharpRestriction/GrowthBounds/FormalGroupSystem.lean}{Formalism.MainNotions.FormalGroupSystem.Sys.hhA}
\leanref{Formalism/SharpRestriction/GrowthBounds/FormalGroupSystem.lean}{Formalism.MainNotions.FormalGroupSystem.atPi}
\leanref{Formalism/SharpRestriction/GrowthBounds/MultisetValuation/EntryAOneParity.lean}{Formalism.MainNotions.FormalGroupSystem.Pf.ea_InClassA}

\begin{lemma}\label{pd-hecke-chain-edge-scaling}
Let $p\geq5$ be a prime and take a formal group system at $p$ (Definition~\ref{pd-formal-group-system}).
Let $J\leq m_*$, let $v$ be an orbit and $0\leq i\leq N(J)$, and put
$A_J(v,i)=\sum_s\sum_{a=1}^{i}\lambda(v,s)_a$, where $\lambda(v,s)$ is formed as in
Definition~\ref{pd-local-model} from the weights $\lambda_J(v,s)$ of the model at level $J$.
\begin{enumerate}
\item If $L$ is a lattice tuple at level $J$ and $L''\subset L_v$ is an edge of index $i$ at $v$, then
  $p^{A_J(v,i)}H^\Delta(L[v\mapsto L''])\subset H^\Delta(L)$.
\item Let $F\in\mathrm{Fm}_J$ vanish at every pair $(c,L)$ with $L$ not a lattice tuple, and satisfy
  $F(c,\gamma L)=\chi_c(\gamma)\,\gamma\cdot F(c,L)$ for every component $c$, every $\gamma$ in its level
  and every tuple $L$. Then $T_{v,i}F$ has the same two properties. The operator $p^{A_J(v,i)}T_{v,i}$
  maps $\mathrm{Int}_J$ into itself, and $T_{v,i}$ maps the $\mathcal K$-span of $\mathrm{Int}_J$ into
  itself.
\end{enumerate}
\end{lemma}
\leanref{Formalism/SharpRestriction/GrowthBounds/MultisetValuation/RestrictedHeckeIntegralOfEdges.lean}{Formalism.MainNotions.FormalGroupSystem.Pf.RestWeylEdge}
\leanref{Formalism/SharpRestriction/GrowthBounds/MultisetValuation/RestrictedHeckeIntegralOfEdges.lean}{Formalism.MainNotions.FormalGroupSystem.Pf.restWeylEdge}
\leanref{Formalism/SharpRestriction/GrowthBounds/MultisetValuation/GShiftRowGenMonomial.lean}{Formalism.MainNotions.FormalGroupSystem.Pf.dl_pos}
\leanref{Formalism/SharpRestriction/GrowthBounds/MultisetValuation/SmulUpdate.lean}{Formalism.MainNotions.FormalGroupSystem.Pf.restTmEquiv}
\leanref{Formalism/SharpRestriction/GrowthBounds/MultisetValuation/RestrictedHeckeIntegralOfEdges.lean}{Formalism.MainNotions.FormalGroupSystem.Pf.TmInt}
\leanref{Formalism/SharpRestriction/GrowthBounds/MultisetValuation/RestrictedHeckeIntegralOfEdges.lean}{Formalism.MainNotions.FormalGroupSystem.Pf.restTmFQ}

\begin{proof}
(1) The tuples $L[v\mapsto L'']$ and $L$ differ only at $v$, so it suffices to show, for each position
$s$ of $v$, that $p^{\sum_{a\leq i}\lambda_a}f\in\Delta_\lambda(M^\vee)$ for every
$f\in\Delta_\lambda(M''^\vee)$, where $\lambda=\lambda(v,s)$, $N=N(J)$, and $M$ and $M''$ are the
$\mathcal R$-spans of $\iota_v\sigma_v^s(L_v)$ and $\iota_v\sigma_v^s(L'')$. By the elementary divisor
theorem there is a matrix $g$ whose columns form a $W_v$-basis of $L_v$ such that the columns of $gD$
form a basis of $L''$, where $D$ is diagonal with entries $p$ ($i$ times) and $1$. Since $\sigma_v$ fixes
$p$, the same holds for $g'=\iota_v\sigma_v^s(g)$, $M$ and $M''$. Let $\nu=\lambda-\lambda_N$ be the
natural weight and $F(x)=\det(g'D)^{\lambda_N}f(g'Dx)$, which lies in $\Delta_{\mathcal R}(\nu)$ by the
description of $\Delta_\lambda(M''^\vee)$ in Definition~\ref{pd-local-model}. Then
\[
 \det(g')^{\lambda_N}\,p^{\sum_{a\leq i}\lambda_a}f(g'x)=p^{\sum_{a\leq i}\nu_a}F(D^{-1}x).
\]
The substitution $x\mapsto D^{-1}x$ divides a monomial of degree $d$ in the first $i$ rows of $x$ by
$p^d$. The module $\Delta_{\mathcal R}(\nu)$ is spanned by polynomials that are homogeneous in each row of
$x$, since $h_\nu$ is and the operators $E^{(k)}_{ab}$ of Definition~\ref{pd-weight-lattices} preserve
this. Every monomial of an element of $\Delta_{\mathcal R}(\nu)\subset\nabla_{\mathcal R}(\nu)$ has degree
at most $\nu_1+\cdots+\nu_i$ in the first $i$ rows: by the standard basis theorem for bitableaux
\cite{DEP80} it is a combination of products, over the columns $c$ of $\nu$, of minors of $x$ with $|c|$
distinct rows, and such a product has degree at most $\sum_c\min(|c|,i)=\nu_1+\cdots+\nu_i$ in the first
$i$ rows. Hence $p^{\sum_{a\leq i}\nu_a}F(D^{-1}x)\in\Delta_{\mathcal R}(\nu)$. The condition for the basis
$g'$ of $M$ gives it for every basis, since $\Delta_{\mathcal R}(\nu)$ is stable under $x\mapsto hx$ for
$h\in GL_N(\mathcal R)$; so $p^{\sum_{a\leq i}\lambda_a}f\in\Delta_\lambda(M^\vee)$. The product over the
positions of $v$ of these powers of $p$ is $p^{A_J(v,i)}$.

(2) If $L$ is not a lattice tuple, then either $L_v$ is not a lattice, and there is no edge at $L_v$, or
$L[v\mapsto L'']$ is not a lattice tuple for any edge; in both cases $T_{v,i}F(c,L)=0$. An element
$\gamma$ of a level carries the edges $L''\subset L_v$ of index $i$ bijectively to the edges of index $i$
at $(\gamma L)_v$, with $\gamma(L[v\mapsto L''])=(\gamma L)[v\mapsto\gamma L'']$, which gives the
equivariance. For $F\in\mathrm{Int}_J$ and a lattice tuple $L$, (1) gives
$p^{A_J(v,i)}T_{v,i}F(c,L)=\sum_{L''}p^{A_J(v,i)}F(c,L[v\mapsto L''])\in H^\Delta(L)$, a finite sum. So
$p^{A_J(v,i)}T_{v,i}F\in\mathrm{Int}_J$, and the last assertion follows.
\end{proof}

\begin{lemma}\label{pd-hecke-chain-edge-count}
Let $W$ be a discrete valuation ring with uniformizer $p$, residue field of $q$ elements and fraction
field $Q$. Let $N,g,i\geq0$, let $A\subset Q^N$ and $B\subset Q^{N+g}$ be $W$-lattices, and let $x$ be an
$(N+g)\times N$ matrix over $Q$. Then the number of $W$-lattices $B_1$ with $pB\subset B_1\subset B$,
$|B/B_1|=q^i$ and $xA\subset B_1$ is
\[
 \sum_{l=0}^{i}q^{(i-l)(g-l)}\binom{g}{l}_{\!q}\cdot\#\big\{A_1:\ A\subset A_1,\ pA_1\subset A,\
 |A_1/A|=q^{i-l},\ xA_1\subset B\big\},
\]
where $\binom gl_q$ is the Gaussian binomial coefficient, which is $0$ for $l>g$.
\end{lemma}
\leanref{Formalism/SharpRestriction/GrowthBounds/MultisetValuation/CountLoc.lean}{Formalism.MainNotions.FormalGroupSystem.Pf.CountLoc}
\leanref{Formalism/SharpRestriction/GrowthBounds/MultisetValuation/SmulPsiGUp.lean}{Formalism.MainNotions.FormalGroupSystem.Pf.countLoc}
\leanref{Formalism/SharpRestriction/GrowthBounds/MultisetValuation/SmulPsiGUp.lean}{Formalism.MainNotions.FormalGroupSystem.Pf.Cnt.core}

\begin{proof}
If $xA\not\subset B$, both sides are $0$: on the left because $B_1\subset B$, on the right because
$A\subset A_1$. Suppose $xA\subset B$. Choosing bases, $A=W^N$, $B=W^{N+g}$ and $x$ has entries in $W$; let
$f\colon\mathbf F^N\to\mathbf F^{N+g}$ be its reduction, $\mathbf F=W/pW$, and $\rho$ its rank. The map
$B_1\mapsto B_1/pB$ is a bijection from the lattices counted on the left to the subspaces of
$\mathbf F^{N+g}$ of codimension $i$ that contain the image of $f$. The map $A_1\mapsto A_1/A\subset
p^{-1}A/A\cong\mathbf F^N$ is a bijection from the lattices $A_1$ on the right to the subspaces of
dimension $i-l$, and $xA_1\subset B$ if and only if this subspace lies in the kernel of $f$, because
$x(p^{-1}y)\in B$ if and only if $xy\in pB$. So the left side is $\binom{N+g-\rho}{i}_q$ and the right side
is $\sum_lq^{(i-l)(g-l)}\binom gl_q\binom{N-\rho}{i-l}_q$. These are equal by the $q$-Vandermonde
identity $\binom{m+n}{h}_q=\sum_k\binom{m}{h-k}_q\binom{n}{k}_q\,q^{(m-h+k)k}$ \cite[Chapter~3]{And76},
with $m=g$, $n=N-\rho$, $h=i$ and $k=i-l$.
\end{proof}

\begin{lemma}\label{pd-hecke-chain-relation}
Let $p\geq5$ be a prime and take an admitted formal group system at $p$
(Definition~\ref{pd-formal-group-system}), with the notation of
Definition~\ref{pd-hecke-chain-operators}. Let $v$ be an orbit and $j<m_*$, and let $a,a'\in\mathcal K$ be
the elements of valuation $0$ that condition~(f) gives for $v$ and $j$ (there called $u$ and $u'$). Then
for every $i\geq0$ and every $F$ in the $\mathcal K$-span of $\mathrm{Int}_j$,
\[
 E_{v,i}(K_jF)=(aa')^{-i}\,K_j\Big(\sum_{l=0}^{i}e_l\big(a\,p^{\operatorname{vm}_v(j+1)}q_x:0\leq x\leq
 t\big)\,E_{v,i-l}F\Big),
\]
where $E_{v,i}$ acts on $\mathrm{Fm}_{j+1}$ on the left and $E_{v,i-l}$ on $\mathrm{Fm}_j$ on the right,
and $e_l$ is the elementary symmetric polynomial of degree $l$ in the $t+1$ entries.
\end{lemma}
\leanref{Formalism/SharpRestriction/GrowthBounds/MultisetValuation/HeckeValueEstimateOfRecordValue.lean}{Formalism.MainNotions.FormalGroupSystem.Pf.RestHecke}
\leanref{Formalism/SharpRestriction/GrowthBounds/MultisetValuation/HeckeConditionOfAdm.lean}{Formalism.MainNotions.FormalGroupSystem.Pf.heckeCondition_of_adm}
\leanref{Formalism/SharpRestriction/GrowthBounds/MultisetValuation/HeckeConditionOfAdm.lean}{Formalism.MainNotions.FormalGroupSystem.Pf.hk_point}
\leanref{Formalism/SharpRestriction/GrowthBounds/MultisetValuation/CoupledSumOfEntries.lean}{Formalism.MainNotions.FormalGroupSystem.Pf.ker_dn_sum}
\leanref{Formalism/SharpRestriction/GrowthBounds/MultisetValuation/CoupledSumOfEntries.lean}{Formalism.MainNotions.FormalGroupSystem.Pf.hk_adj}
\leanref{Formalism/SharpRestriction/GrowthBounds/MultisetValuation/LatAtZeroUpd.lean}{Formalism.MainNotions.FormalGroupSystem.Pf.countG}
\leanref{Formalism/SharpRestriction/GrowthBounds/MultisetValuation/HkEsymmSmul.lean}{Formalism.MainNotions.FormalGroupSystem.Pf.hk_coef}
\leanref{Formalism/SharpRestriction/GrowthBounds/MultisetValuation/HkEsymmSmul.lean}{Formalism.MainNotions.FormalGroupSystem.Pf.nC_cast}
\leanref{Formalism/ThetaLogCoordinates/UniformThetaSubstitution/FrobeniusLiftNorms/GaussElementarySymmetric.lean}{Formalism.InvariantFactorPaths.G7.gauss_elementary_symmetric}

\begin{proof}
Put $N=N(j)$, so that $N(j+1)=N+t+1$, and write $Q=Q_v$ and $\operatorname{vm}=\operatorname{vm}_v(j+1)$.

\emph{Step 1: lattice tuples.} By condition~(d) and Lemma~\ref{pd-hecke-chain-edge-scaling}(2), both
sides lie in the $\mathcal K$-span of $\mathrm{Int}_{j+1}$, so both vanish at the tuples that are not
lattice tuples. Fix a component $c'$ and a lattice tuple $L'$ at level $j+1$.

\emph{Step 2: the kernel along edges.} Let $(c,L)$ be a component and a lattice tuple at level $j$, and
$G$ a polynomial at level $j$. For a $\mathbf Q$-valued function $\varphi$ on the $N(j+1)\times N$
matrices over $Q_v$ put
\[
 \kappa_{L,L'}(\varphi)=\prod_wp^{\operatorname{vm}_w(j+1)(\ell_w(L)-\ell_w(L'))}a^{\ell_v(L)}a'^{\ell_v(L')}
 Q^{-(N\ell_v(L')-(N+t+1)\ell_v(L))}\,\mathrm{Del}_j\big(G,\Xi_{v,j,\alpha}(\varphi)\big),
\]
with $\alpha=\operatorname{aw}_{v,j}((c,L),(c',L'))$. It is additive in $\varphi$, and by condition~(f)
and the definition of the kernel, $\ker_j(c,L;c',L')(G)=\kappa_{L,L'}(\operatorname{ind}_{v,j}(L,L'))$.
If $L_1\subset L'_v$ is an edge of index $i$, replacing $L'$ by $L'[v\mapsto L_1]$ lowers $\ell_v(L')$ by
$i$ and changes neither $\alpha$ nor the other lengths, so
$\kappa_{L,L'[v\mapsto L_1]}=C'_i\,\kappa_{L,L'}$ with $C'_i=p^{\operatorname{vm}\,i}a'^{-i}Q^{Ni}$. If
$L_v\subset A_1$ is an edge of index $k$, replacing $L$ by $L[v\mapsto A_1]$ raises $\ell_v(L)$ by $k$, so
$\kappa_{L[v\mapsto A_1],L'}=C_k\,\kappa_{L,L'}$ with $C_k=p^{\operatorname{vm}\,k}a^kQ^{(N+t+1)k}$.

\emph{Step 3: counting.} The residue field of $W_v$ has $Q^2$ elements. By
Lemma~\ref{pd-hecke-chain-edge-count} with $A=L_v$, $B=L'_v$ and $g=t+1$, evaluated at each matrix,
\[
 \sum_{L_1}\operatorname{ind}_{v,j}(L,L'[v\mapsto L_1])=\sum_{l=0}^{i}n(i,l)\sum_{A_1}
 \operatorname{ind}_{v,j}(L[v\mapsto A_1],L'),\qquad n(i,l)=Q^{2(i-l)(t+1-l)}\binom{t+1}{l}_{\!Q^2},
\]
where $L_1$ runs over the edges $L_1\subset L'_v$ of index $i$ and $A_1$ over the edges $L_v\subset A_1$
of index $i-l$. Applying $\kappa_{L,L'}$ and Step~2,
\[
 \sum_{L_1}\ker_j(c,L;c',L'[v\mapsto L_1])(G)=\sum_{l=0}^{i}C'_i\,n(i,l)\,C_{i-l}^{-1}\sum_{A_1}
 \ker_j(c,L[v\mapsto A_1];c',L')(G).
\]

\emph{Step 4: the expansion.} Write the expansion of condition~(e) as
$K_jF(c',L'')=\sum_xw_x^{-1}c_K\sum_y\ker_j(c_{xy},L_{xy};c',L'')(F(c_{xy},L_{xy}))$ for lattice tuples
$L''$: if $j+1<m_*$, $x$ runs over $\mathrm{Rep}_j$, $w_x$ is the order of the stabilizer and $y$ takes one
value; if $j+1=m_*$, these are the data of~(e). By Step~3 for each $(c_{xy},L_{xy})$,
\[
 (T_{v,i}K_jF)(c',L')=\sum_{L_1}K_jF(c',L'[v\mapsto L_1])=\sum_{l=0}^{i}C'_in(i,l)C_{i-l}^{-1}c_K
 \sum_xw_x^{-1}\sum_y\sum_{A_1}\ker_j(c_{xy},L_{xy}[v\mapsto A_1];c',L')(F(c_{xy},L_{xy})).
\]

\emph{Step 5: adjunction.} For $k\geq0$,
\[
 \sum_xw_x^{-1}\sum_y\sum_{A_1}\ker_j(c_{xy},L_{xy}[v\mapsto A_1];c',L')(F(c_{xy},L_{xy}))
 =\sum_xw_x^{-1}\sum_y\ker_j(c_{xy},L_{xy};c',L')\big((T_{v,k}F)(c_{xy},L_{xy})\big),
\]
the inner sum on the left over the edges $(L_{xy})_v\subset A_1$ of index $k$. If $j+1=m_*$ this is the
last identity of~(e). If $j+1<m_*$, apply condition~(h), coefficient by coefficient, to the function
$\varphi(c,L,L_1)=\ker_j(c,L_1;c',L')(F(c,L))$, set to $0$ when $L_1$ is not a lattice tuple; it is
invariant under the levels by condition~(g) and the equivariance of $F$, and the right side of~(h) is the
right side above by the linearity of the kernel. By Lemma~\ref{pd-hecke-chain-edge-scaling}(2),
$T_{v,k}F$ lies in the $\mathcal K$-span of $\mathrm{Int}_j$, so by~(e) the right side is
$c_K^{-1}K_j(T_{v,k}F)(c',L')$. With Step~4,
\[
 T_{v,i}(K_jF)=\sum_{l=0}^{i}C'_i\,n(i,l)\,C_{i-l}^{-1}\,K_j(T_{v,i-l}F)\quad\text{at }(c',L').
\]

\emph{Step 6: coefficients.} By the $q$-binomial theorem \cite[Chapter~3]{And76},
$e_l(1,q,\dots,q^{t})=q^{l(l-1)/2}\binom{t+1}{l}_q$; with $q=Q^2$ this gives
$\binom{t+1}l_{Q^2}=Q^{l(t+1-l)}e_l(q_x:0\leq x\leq t)$. A computation with the exponents then gives
\[
 Q^{-i(N+t+1-i)}\,C'_i\,n(i,l)\,C_{i-l}^{-1}=(aa')^{-i}\,e_l\big(a\,p^{\operatorname{vm}}q_x:0\leq x\leq
 t\big)\,Q^{-(i-l)(N-(i-l))}.
\]
Multiplying the identity of Step~5 by $Q^{-i(N+t+1-i)}$ gives the lemma.
\end{proof}

\begin{lemma}\label{pd-hecke-chain-inversion}
Let $\mathcal K$ be a field, $M_0,M_1,\dots$ vector spaces over $\mathcal K$ with linear maps
$K_j\colon M_j\to M_{j+1}$, and $K^{(j)}=K_{j-1}\circ\cdots\circ K_0$. Let $Q,m_1,m_2,\dots\in\mathcal K$,
$t\geq0$ and $J\geq0$, and for each $j$ let $E_{j,0},E_{j,1},\dots$ be endomorphisms of $M_j$. Let
$\beta_l$ be the family $(m_lQ^{t-2x})_{0\leq x\leq t}$ and $\beta_{\leq j}$ the family of the entries of
$\beta_1,\dots,\beta_j$. Suppose that
\[
 E_{j+1,i}K_j=K_j\sum_{l=0}^{i}e_l(\beta_{j+1})\,E_{j,i-l}
\]
for every $j<J$ and $i\geq0$. Then for every $j\leq J$ and $r\geq0$,
\[
 \sum_{l=0}^{r}(-1)^lh_l(\beta_{\leq j})\,E_{j,r-l}K^{(j)}=K^{(j)}E_{0,r}.
\]
\end{lemma}
\leanref{Formalism/ThetaLogCoordinates/UniformThetaSubstitution/StatementC.lean}{Formalism.InvariantFactorPaths.G7.Intertwine}
\leanref{Formalism/ThetaLogCoordinates/UniformThetaSubstitution/StatementC.lean}{Formalism.InvariantFactorPaths.G7.Inversion}
\leanref{Formalism/ThetaLogCoordinates/UniformThetaSubstitution/FrobeniusLiftNorms/GaussElementarySymmetric.lean}{Formalism.InvariantFactorPaths.G7.inversion_formula}
\leanref{Formalism/ThetaLogCoordinates/UniformThetaSubstitution/FrobeniusLiftNorms/GaussElementarySymmetric.lean}{Formalism.InvariantFactorPaths.G7.iterate_statement}
\leanref{Formalism/ThetaLogCoordinates/UniformThetaSubstitution/FrobeniusLiftNorms/GaussElementarySymmetric.lean}{Formalism.InvariantFactorPaths.G7.hecke_convolution}

\begin{proof}
The elementary symmetric polynomials of a family that is the union of two families $\alpha$ and $\beta$
are $e_n=\sum_{l+l'=n}e_l(\alpha)e_{l'}(\beta)$, and $\sum_{l=0}^n(-1)^lh_l(\beta)e_{n-l}(\beta)=0$ for
$n\geq1$ \cite[Chapter~I, (2.6)]{Mac95}. First, $E_{j,i}K^{(j)}=K^{(j)}\sum_le_l(\beta_{\leq j})E_{0,i-l}$
for $j\leq J$ and all $i$, by induction on $j$: for $j=0$ both sides are $E_{0,i}$, and
\[
 E_{j+1,i}K_jK^{(j)}=K_j\sum_le_l(\beta_{j+1})E_{j,i-l}K^{(j)}=K^{(j+1)}\sum_{l,l'}e_l(\beta_{j+1})
 e_{l'}(\beta_{\leq j})E_{0,i-l-l'}=K^{(j+1)}\sum_ne_n(\beta_{\leq j+1})E_{0,i-n}.
\]
Then
\[
 \sum_{l=0}^r(-1)^lh_l(\beta_{\leq j})E_{j,r-l}K^{(j)}=K^{(j)}\sum_{n=0}^{r}\Big(\sum_{l=0}^n(-1)^l
 h_l(\beta_{\leq j})e_{n-l}(\beta_{\leq j})\Big)E_{0,r-n}=K^{(j)}E_{0,r}.\qedhere
\]
\end{proof}

\begin{lemma}\label{pd-hecke-chain-commutation}
Let $p\geq5$ be a prime and take an admitted formal group system at $p$
(Definition~\ref{pd-formal-group-system}), with the notation of Definitions~\ref{pd-formal-group-numbers}
and~\ref{pd-hecke-chain-operators} and the Hecke operators $y_v$ of
Definition~\ref{pd-local-model}(4). Let $v$ be an orbit. There are $u_l\in\mathcal R^\times$ and
$m_l\in p^{\operatorname{vm}_v(l)}\mathcal R^\times$, for $1\leq l\leq m_*$, such that the orbit operators
$\mathcal H_{v,J}$ and the elements $c_l(J)$ formed from them satisfy, for every $j<m_*$:
\begin{enumerate}
\item $\mathcal H_{v,j+1}\big(K^{(j+1)}F\big)=p^{\mathrm{Pr}_B(v,j+1)}K^{(j+1)}(y_vF)$ for every $F$ in
  the $\mathcal K$-span of $M$, where $\mathrm{Pr}_B(v,j+1)$ is an integer because $n'_v$ is even;
\item $c_l(j+1)$ has valuation at least $0$ for $0\leq l\leq r$, and $c_r(j+1)$ has valuation at least
  $\alpha_{\mathrm{ev}}(v)$.
\end{enumerate}
\end{lemma}
\leanref{Formalism/SharpRestriction/GrowthBounds/FormalGroupSystem.lean}{Formalism.MainNotions.FormalGroupSystem.KernelCommutationUpToUnits}
\leanref{Formalism/SharpRestriction/GrowthBounds/FormalGroupSystem.lean}{Formalism.MainNotions.FormalGroupSystem.Conclusion4}
\leanref{Formalism/SharpRestriction/GrowthBounds/MultisetValuation/HeckeValueEstimateOfRecordValue.lean}{Formalism.MainNotions.FormalGroupSystem.Pf.alternate_expansion_of_relations}
\leanref{Formalism/SharpRestriction/GrowthBounds/MultisetValuation/TheoremOfAlternateForm.lean}{Formalism.MainNotions.FormalGroupSystem.Pf.theorem_of_alternate_form}
\leanref{Formalism/SharpRestriction/GrowthBounds/MultisetValuation/HeckeValueEstimateOfRecordValue.lean}{Formalism.MainNotions.FormalGroupSystem.Pf.Ci}
\leanref{Formalism/SharpRestriction/GrowthBounds/MultisetValuation/TheoremOfAlternateForm.lean}{Formalism.MainNotions.FormalGroupSystem.Pf.Cii}
\leanref{Formalism/ThetaLogCoordinates/UniformThetaSubstitution/FrobeniusLiftNorms/GaussElementarySymmetric.lean}{Formalism.InvariantFactorPaths.G7.theorem_c_i}

\begin{proof}
Write $Q=Q_v$ and $N'=N'_v$.

\emph{Step 1: the relation.} For $j<m_*$ let $a_j,a'_j$ be the elements of condition~(f) for $v$ and $j$,
and put $u_{j+1}=(a_ja'_j)^{-1}$ and $n_{j+1}=a_jp^{\operatorname{vm}_v(j+1)}$, so that
$u_l\in\mathcal R^\times$ and $n_l\in p^{\operatorname{vm}_v(l)}\mathcal R^\times$. By
Lemma~\ref{pd-hecke-chain-relation},
$E_{v,i}K_jF=u_{j+1}^{\,i}K_j\sum_le_l(n_{j+1}q_x)E_{v,i-l}F$ for $F$ in the $\mathcal K$-span of
$\mathrm{Int}_j$.

\emph{Step 2: inversion.} Let $M_j$ be the $\mathcal K$-span of $\mathrm{Int}_j$. The map $K_j$ carries
$M_j$ into $M_{j+1}$ by condition~(d), and $T_{v,i}$ carries $M_j$ into itself by
Lemma~\ref{pd-hecke-chain-edge-scaling}(2). Put $E'_{j,i}=(u_1\cdots u_j)^{-i}E_{v,i}$ on $M_j$ and
$m_l=(u_1\cdots u_{l-1})^{-1}n_l$, so that $m_l\in p^{\operatorname{vm}_v(l)}\mathcal R^\times$. Since
$e_l$ is homogeneous of degree $l$, Step~1 becomes $E'_{j+1,i}K_j=K_j\sum_le_l(m_{j+1}q_x)E'_{j,i-l}$.
Lemma~\ref{pd-hecke-chain-inversion} with $J=m_*$ gives, for $j<m_*$ and on $M_0$,
\[
 \sum_{l=0}^{r}(-1)^lh_l\big(m_{l'}q_x:l'\leq j+1,\ x\leq t\big)(u_1\cdots u_{j+1})^{-(r-l)}E_{v,r-l}
 K^{(j+1)}=K^{(j+1)}E_{v,r}.
\]

\emph{Step 3: part (1).} Multiply this identity by $p^{N'(j+1,r)}Q^{r(N(j+1)-r)}$. Since
$r(N-r)-(r-l)(N-r+l)=l(N-2r+l)$, $N'(j+1,0)=0$ and $T_{v,0}$ is the identity, the $l$-th term on the left
becomes $(-1)^lc_l(j+1)(u_1\cdots u_{j+1})^{-(r-l)}T^\natural_{v,r-l}K^{(j+1)}$, so the left side is
$\mathcal H_{v,j+1}K^{(j+1)}$. On the right, $E_{v,r}=Q^{-r(N(0)-r)}T_{v,r}$ at level $0$ and
$y_v=p^{A_v}T_{v,r}$ with $A_v=\sum_{s\ \mathrm{even}}\sum_{a\leq r}\lambda_0(v,s)_a-\sum_{s\
\mathrm{odd}}\sum_{a=k-r+1}^k\lambda_0(v,s)_a$. Since $N(j+1)-N(0)=(j+1)(t+1)$, the right side becomes
$p^{e}K^{(j+1)}y_v$ with
\[
 e=N'(j+1,r)+\tfrac{n'_v}2r(j+1)(t+1)-A_v=\sum_{s\ \mathrm{odd}}\Big(r\,w(v,j+1,s)+\sum_{a=k-r+1}^k
 \lambda_0(v,s)_a\Big)+\tfrac{n'_v}2r(j+1)(t+1)=\mathrm{Pr}_B(v,j+1).
\]

\emph{Step 4: part (2).} The element $m_{l'}q_x$ has valuation
$\operatorname{vm}_v(l')+\frac{n'_v}2(t-2x)$. For $l'\leq j+1$ the fills give $w(v,1,s)\leq
w(v,l',s)\leq w(v,j+1,s)$, so this is at least
$\mu=\sum_{s\ \mathrm{even}}w(v,1,s)-\sum_{s\ \mathrm{odd}}w(v,j+1,s)-\frac{n'_v}2t$, and $h_l$ of the
family has valuation at least $l\mu$. Hence $c_l(j+1)$ has valuation at least
\[
 N'(j+1,r)-N'(j+1,r-l)+\tfrac{n'_v}2l\big(N(j+1)-2r+l\big)+l\mu=\sum_{s\ \mathrm{even}}\Big(\sum_{a=r-l+1}^r
 \lambda_0(v,s)_a+l\,w(v,1,s)\Big)+\tfrac{n'_v}2l\big(N(j+1)-2r+l-t\big).
\]
By dominance and the fills, $\lambda_0(v,s)_a\geq\lambda_0(v,s)_k\geq-w(v,1,s)$, so the sum is at least
$0$; and $N(j+1)-t=k+j(t+1)+1>2r$, so the last term is at least $0$. For $l=r$ and an even $s\in T_v$, the
summand is $\lambda_0(v,s)_1+w(v,1,s)+\sum_{a=2}^r(\lambda_0(v,s)_a+w(v,1,s))\geq\lambda_0(v,s)_1-
\lambda_0(v,s)_2=\alpha(s)$, since $w(v,1,s)\geq-\lambda_0(v,s)_k\geq-\lambda_0(v,s)_2$; at the other
even $s$ it is at least $0=\alpha(s)$. So $c_r(j+1)$ has valuation at least $\alpha_{\mathrm{ev}}(v)$.
\end{proof}

\begin{lemma}\label{pd-kernel-orbit-descent}
Take a formal group system at $p$ (Definition~\ref{pd-formal-group-system}), let $J\leq m_*$, let $L$ be
a lattice tuple at level $J$ and $g\in\mathbf R$. Let $y$ be a polynomial over $\mathcal K$ in the
variables of $\mathcal X$ at level $J$, and suppose that $y=\sum_ka_kz_k$, a finite sum in the polynomial
ring over $\mathcal K_c$, with $v(a_k)\geq g$ and $z_k\in H^\Delta(L)$. Then
$y\in\varpi^{\lceil ge_{\mathcal R}\rceil}H^\Delta(L)$.
\end{lemma}
\leanref{Formalism/SharpRestriction/GrowthBounds/MultisetValuation/InLatCSmulC.lean}{Formalism.MainNotions.FormalGroupSystem.Pf.descent}
\leanref{Formalism/SharpRestriction/GrowthBounds/MultisetValuation/DescentCore.lean}{Formalism.MainNotions.FormalGroupSystem.Pf.descent_core}

\begin{proof}
By induction on the number $n$ of terms; for $n=0$, $y=0$. If the $z_k$ are linearly dependent over
$\mathcal K$, take a relation $\sum_kb_kz_k=0$ with $b_k\in\mathcal K$ not all $0$, and $k_0$ with
$b_k/b_{k_0}\in\mathcal R$ for all $k$, which exists because $\mathcal R$ is a valuation ring. Then
$y=\sum_{k\neq k_0}(a_k-(b_k/b_{k_0})a_{k_0})z_k$, with coefficients of $v\geq g$ because
$\mathcal R\subset\mathcal R_c$, and induction applies. If the $z_k$ are linearly independent over
$\mathcal K$, they stay so over $\mathcal K_c$, and since $y$ and the $z_k$ have coefficients in
$\mathcal K$, Cramer's rule on a nonzero $n\times n$ minor of their coefficients gives
$a_k\in\mathcal K$. Since $\mathcal R_c\cap\mathcal K=\mathcal R$ and $\mathcal R$ is a discrete valuation
ring, an element $a\in\mathcal K$ with $v(a)\geq g$ has order at least $ge_{\mathcal R}$ at $\varpi$, so
$a\in\varpi^{\lceil ge_{\mathcal R}\rceil}\mathcal R$. Hence $y\in\varpi^{\lceil ge_{\mathcal
R}\rceil}H^\Delta(L)$.
\end{proof}

\begin{lemma}\label{pd-kernel-chain-bound}
Let $p\geq5$ be a prime and take an admitted formal group system at $p$
(Definition~\ref{pd-formal-group-system}), with the notation of Definitions~\ref{pd-hecke-chain-operators}
and~\ref{pd-kernel-orbit-index}. Let $j<m_*$, let $i_0$ be a function from the orbits to
$\{0,\dots,r\}$, and $g\in\mathbf R$. Suppose that for all components and lattice tuples $(c,L)$ at level
$j$ and $(c',L')$ at level $j+1$, every frame $\phi$ for them and every column choice $I$ into $N(j+1)$,
every coefficient of $\langle G^{i_0}_I,\Theta^Z_\phi\rangle_Z$ has $v\geq g$. Let $T^\natural_{i_0}$ be
the composite, in any order, of the weighted edge operators $T^\natural_{v,i_0(v)}$ on $\mathrm{Fm}_{j+1}$
over the orbits $v$. Then for every $G\in\mathrm{Int}_j$ the function $T^\natural_{i_0}K_jG$ has valuation
at least $g-d_\Gamma(j)$.
\end{lemma}
\leanref{Formalism/SharpRestriction/GrowthBounds/MultisetValuation/PlaceB.lean}{Formalism.MainNotions.FormalGroupSystem.Pf.FormB}
\leanref{Formalism/SharpRestriction/GrowthBounds/MultisetValuation/SmulUpdate.lean}{Formalism.MainNotions.FormalGroupSystem.Pf.formB}
\leanref{Formalism/SharpRestriction/GrowthBounds/MultisetValuation/KernelSumOfClassData.lean}{Formalism.MainNotions.FormalGroupSystem.Pf.formB_of}
\leanref{Formalism/SharpRestriction/GrowthBounds/MultisetValuation/KernelSumOfClassData.lean}{Formalism.MainNotions.FormalGroupSystem.Pf.entry_mem}
\leanref{Formalism/SharpRestriction/GrowthBounds/MultisetValuation/InLatCSmulC.lean}{Formalism.MainNotions.FormalGroupSystem.Pf.wt_inv}

\begin{proof}
\emph{Step 1: reduction.} By condition~(d) and Lemma~\ref{pd-hecke-chain-edge-scaling}(2),
$T^\natural_{i_0}K_jG$ lies in the $\mathcal K$-span of $\mathrm{Int}_{j+1}$; it vanishes at the tuples
that are not lattice tuples and has the equivariance of $\mathrm{Int}_{j+1}$. Such a function lies in
$\varpi^n\mathrm{Int}_{j+1}$ as soon as its value at every component $c'$ and lattice tuple $L'$ lies in
$\varpi^nH^\Delta(L')$. Fix $c'$ and $L'$, and put $n=\lceil(g-d_\Gamma(j))e_{\mathcal R}\rceil$.

\emph{Step 2: the edges.} By the definition of the $T_{v,i}$,
\[
 (T^\natural_{i_0}K_jG)(c',L')=\sum_{L''}\prod_vp^{N'_v(j+1,i_0(v))}\,(K_jG)(c',L''),
\]
the sum over the tuples $L''$ with $L''_v=L'_v$ if $i_0(v)=0$ and $L''_v\subset L'_v$ an edge of index
$i_0(v)$ otherwise. The product of powers of $p$ is the factor $\prod_vp^{N'_v}$ of condition~(i).

\emph{Step 3: the kernels.} Write the expansion of condition~(e) as
$K_jG(c',L'')=\sum_xw_x^{-1}c_K\sum_y\ker_j(c_{xy},L_{xy};c',L'')(G(c_{xy},L_{xy}))$, as in Step~4 of the
proof of Lemma~\ref{pd-hecke-chain-relation}. Let $\phi$ be the frame for
$((c_{xy},L_{xy}),(c',L''))$ that condition~(i) gives for $L'$, $L''$ and $i_0$, and apply~(i) with
$f=\Theta^Z_\phi$, whose pairings with the $G^{i_0}_I$ have coefficients of $v\geq g$ by hypothesis. By
(5), $f_\phi=\Theta^Z(c_{xy},L_{xy};c',L'')$, so the polynomial of~(i) is
$\prod_vp^{N'_v}\ker_j(c_{xy},L_{xy};c',L'')(G(c_{xy},L_{xy}))$, with $G(c_{xy},L_{xy})\in
H^\Delta(L_{xy})$ since $G\in\mathrm{Int}_j$. So it is a sum of products $yz$ with $v(y)\geq g$ and $z\in
H^\Delta(L')$.

\emph{Step 4: denominators.} The element $c_K$ is a unit of $\mathcal R_c$, and $v_p(w_x)\leq
d_\Gamma(j)$, by the choice of $d_\Gamma$ if $j+1<m_*$ and by~(e) if $j+1=m_*$. Integers prime to $p$ are
units of $W(\overline{\mathbf F}_p)\subset\mathcal R_c$, so $v(w_x^{-1}c_K)\geq-d_\Gamma(j)$. Since
$\mathcal R_c$ has rank one, the conditions $v(y)\geq g$ add in products. Hence
$(T^\natural_{i_0}K_jG)(c',L')$ is a sum of products $yz$ with $v(y)\geq g-d_\Gamma(j)$ and $z\in
H^\Delta(L')$, and Lemma~\ref{pd-kernel-orbit-descent} puts it in $\varpi^nH^\Delta(L')$.
\end{proof}

\begin{lemma}\label{pd-kernel-orbit-product}
Let $p\geq5$ be a prime and take an admitted formal group system at $p$
(Definition~\ref{pd-formal-group-system}), with the notation of Definition~\ref{pd-kernel-orbit-index}.
Let $j<m_*$, let $i_0$ be a function from the orbits to $\{0,\dots,r\}$, and for each orbit $v$ let
$g_v\in\mathbf Q$ be such that the system has lower bound $g_v$ at $(j,v)$ with index $i_0(v)$. Then for
all components and lattice tuples $(c,L)$ at level $j$ and $(c',L')$ at level $j+1$, every frame $\phi$
for them and every column choice $I$ into $N(j+1)$, every coefficient of
$\langle G^{i_0}_I,\Theta^Z_\phi\rangle_Z$ has $v\geq\sum_vg_v$.
\end{lemma}
\leanref{Formalism/SharpRestriction/GrowthBounds/MultisetValuation/RestEntOf.lean}{Formalism.MainNotions.FormalGroupSystem.Pf.kernel_product_form}
\leanref{Formalism/SharpRestriction/GrowthBounds/MultisetValuation/KernelProductForm.lean}{Formalism.MainNotions.FormalGroupSystem.Pf.restPlaces_of}
\leanref{Formalism/SharpRestriction/GrowthBounds/MultisetValuation/KernelProductForm.lean}{Formalism.MainNotions.FormalGroupSystem.Pf.en_places_add}

\begin{proof}
\emph{Step 1: the product form.} By condition~(b), $\Theta_\phi=\exp(-\Phi)U$ with $U(\ell(X))$ integral
over $\mathcal R_c$ and $\Phi\simeq\sum_e\beta(X_e,X'_e)$. Let $\Phi_v=\sum_e\beta(X_e,X'_e)$, the sum
over the $e$ at $v$, and $\Psi=\Phi-\sum_v\Phi_v$. Since $\Psi\simeq0$, $\exp(\Psi(\ell(X)))$ is $1$ plus
a series with coefficients in $W(\overline{\mathbf F}_p)$, and so is its inverse
$\exp(-\Psi(\ell(X)))$. Hence $\Theta_\phi=\prod_v\exp(-\Phi_v)\cdot U_1$ with $U_1=\exp(-\Psi)U$ and
$U_1(\ell(X))$ integral over $\mathcal R_c$.

\emph{Step 2: expansion of $U_1$.} Let $\mu$ be the series with $\ell(\mu(X))=X$, and write
$U_1(\ell(X))=\sum_\gamma c_\gamma X^\gamma$ with $c_\gamma\in\mathcal R_c$. Then
$U_1=\sum_\gamma c_\gamma\prod_v\mu(X)^{\gamma_v}$, where $\gamma_v$ is the part of $\gamma$ at the
variables of $v$; since $\mu(X)^\gamma$ has order $|\gamma|$, only finitely many $\gamma$ contribute to
a given coefficient.

\emph{Step 3: orbit by orbit.} The orbit coordinates substitute for the variables at $v$ combinations of
the $Z_{v,s,a,b}$, so $\Theta^Z_\phi=\sum_\gamma c_\gamma\prod_vF_{\gamma,v}$ with
$F_{\gamma,v}=(\exp(-\Phi_v)\mu(X)^{\gamma_v})^Z$, a series in the variables at $v$. For polynomials
$P_v$ and series $F_v$ in disjoint sets of variables,
$\langle\prod_vP_v,\prod_vF_v\rangle_Z=\prod_v\langle P_v,F_v\rangle_Z$, because the factorials in the
jet pairing are multiplicative. As $G^{i_0}_I=\prod_vG^{(v),i_0(v)}_I$,
\[
 \langle G^{i_0}_I,\Theta^Z_\phi\rangle_Z=\sum_\gamma c_\gamma\prod_v\big\langle
 G^{(v),i_0(v)}_I,F_{\gamma,v}\big\rangle_Z.
\]
The series $\mu(X)^{\gamma_v}$ becomes $X^{\gamma_v}$ after substituting $\ell(X)$, so $F_{\gamma,v}$ is
$\Theta^Z_U$ for an admissible $U$, and the coefficients of the $v$-th factor have $v\geq g_v$.

\emph{Step 4: the bounds add.} Let $x,y\in\mathcal K_c$ with $v(x)\geq q_1$ and $v(y)\geq q_2$, where
$q_1,q_2\in\mathbf Q$. Then $v(xy)\geq q_1+q_2$. Indeed, let $b>0$ and $a\leq(q_1+q_2)b$ be integers,
let $M>0$ be the product of the denominators of $q_1$ and $q_2$, and put $a_i=q_ibM\in\mathbf Z$. Then
$x^{bM}p^{-a_1}$ and $y^{bM}p^{-a_2}$ lie in $\mathcal R_c$, and $a_1+a_2-aM\geq0$, so
\[
 \big((xy)^bp^{-a}\big)^M=x^{bM}p^{-a_1}\cdot y^{bM}p^{-a_2}\cdot p^{a_1+a_2-aM}\in\mathcal R_c .
\]
If $z^M\in\mathcal R_c$ then $z\in\mathcal R_c$: otherwise $z^{-1}\in\mathcal R_c$, since $\mathcal R_c$
is a valuation ring, and $z=z^M(z^{-1})^{M-1}\in\mathcal R_c$. So $(xy)^bp^{-a}\in\mathcal R_c$. By
induction the bound $\sum_vq_v$ holds for a product over the orbits when the $v$-th factor has bound
$q_v\in\mathbf Q$. Multiplication by $c_\gamma\in\mathcal R_c$ keeps a bound, and so does a sum: if
$v(x),v(y)\geq g$ and, say, $x/y\in\mathcal R_c$ (one of $x/y$ and $y/x$ is, as $\mathcal R_c$ is a
valuation ring), then $(x+y)^bp^{-a}=y^bp^{-a}(1+x/y)^b\in\mathcal R_c$. So every coefficient of $\langle G^{i_0}_I,\Theta^Z_\phi\rangle_Z$ has $v\geq\sum_vg_v$.
\end{proof}

\subsection{Classes of bi-additive series}\label{subsection-pd-bi-additive-classes}

\begin{definition}\label{pd-orbit-class-law}
Let $p$ be an odd prime, $W(\overline{\mathbf F}_p)$ the ring of Witt vectors with Frobenius $\sigma$ and
fraction field $K_0$, and $\ell(X)=\sum_{j\geq0}X^{p^{2j}}/p^j$. A power series over $K_0$ is
\emph{integral} if its coefficients lie in $W(\overline{\mathbf F}_p)$, and $f^\sigma$ is the series $f$
with $\sigma$ applied to its coefficients.
\begin{enumerate}
\item The \emph{odd logarithm} is $\varrho(X)=\sum_{j\geq0}X^{p^{2j+1}}/p^{j+1}$. For $c,c'$ in a field
  containing $K_0$, the \emph{explicit series} is
  \[
   G_{c,c'}(X_0,X_1)=c\,A(X_0,X_1)-c'A(X_1,X_0),\qquad A(X_0,X_1)=\sum_{0\leq i\leq j}
   \frac{X_0^{p^{2i}}X_1^{p^{2j+1}}}{p^{i+j+1}}.
  \]
\item The \emph{formal group law} is $F(X_0,X_1)=\ell^{-1}\big(\ell(X_0)+\ell(X_1)\big)$, where $\ell^{-1}$
  is the compositional inverse of $\ell$.
\item For power series $f,g$ over $K_0$ in the same variables, $f\simeq_0g$ if $f-g$ has constant term
  $0$ and $\exp(f-g)$ is integral. Thus $f\simeq g$ (Definition~\ref{pd-formal-group-system}) if and
  only if $f(\ell(X))\simeq_0g(\ell(X))$.
\item For $f\in K_0[[X]]$ and $H\in K_0[[X_0,X_1]]$ the \emph{coboundaries} are
  \[
   \delta_Ff=f(F(X_0,X_1))-f(X_0)-f(X_1),\qquad
   \delta_FH=H(F(X_0,X_1),X_2)-H(X_0,X_2)-H(X_1,X_2),
  \]
  and $H^{\mathrm{sw}}(X_0,X_1)=H(X_1,X_0)$. The series $H$ \emph{vanishes on the axes} if
  $H(X_0,0)=H(0,X_1)=0$. It has \emph{balanced support} if every monomial $X_0^{n_0}X_1^{n_1}$ of $H$ has
  $n_0,n_1\geq1$ and $n_0\equiv pn_1\pmod{p^2-1}$.
\item For $H\in K_0[[X_0,X_1]]$ and $\gamma,\gamma'\in K_0$ put
  \[
   \mathrm D_\sigma H=p^{-1}\big(H^\sigma(X_0^p,X_1^p)-pH\big),\qquad
   \mathrm Q_{\gamma,\gamma'}H=\mathrm D_\sigma H-\gamma\,\varrho(X_0)\ell(X_1)-\gamma'\ell(X_0)\varrho(X_1).
  \]
  The series $H$ \emph{has derivative} $c\in K_0$ if $\partial H/\partial X_0-c\,\ell'(X_0)\varrho(X_1)$ is
  integral.
\end{enumerate}
\end{definition}
\leanref{Formalism/ThetaLogCoordinates/UniformThetaSubstitution/OddLogRelation.lean}{Formalism.InvariantFactorPaths.G7.g0F}
\leanref{Formalism/ThetaLogCoordinates/UniformThetaSubstitution/OddLogRelation.lean}{Formalism.InvariantFactorPaths.G7.Aser}
\leanref{Formalism/ThetaLogCoordinates/UniformThetaSubstitution/OddLogRelation.lean}{Formalism.InvariantFactorPaths.G7.Gser}
\leanref{Formalism/ThetaLogCoordinates/UniformThetaSubstitution/IsLawOfLog.lean}{Formalism.InvariantFactorPaths.G7.IsLawOfLog}
\leanref{Formalism/ThetaLogCoordinates/UniformThetaSubstitution/CyclicPowerClassOfExpEquiv.lean}{Formalism.InvariantFactorPaths.G7.Simeq}
\leanref{Formalism/ThetaLogCoordinates/UniformThetaSubstitution/IsLawOfLog.lean}{Formalism.InvariantFactorPaths.G7.cob}
\leanref{Formalism/ThetaLogCoordinates/UniformThetaSubstitution/IsLawOfLog.lean}{Formalism.InvariantFactorPaths.G7.d1}
\leanref{Formalism/ThetaLogCoordinates/UniformThetaSubstitution/IsLawOfLog.lean}{Formalism.InvariantFactorPaths.G7.swap2}
\leanref{Formalism/ThetaLogCoordinates/UniformThetaSubstitution/IsLawOfLog.lean}{Formalism.InvariantFactorPaths.G7.DivAB}
\leanref{Formalism/ThetaLogCoordinates/UniformThetaSubstitution/FirstDeriv.lean}{Formalism.InvariantFactorPaths.G7.SuppOK}
\leanref{Formalism/ThetaLogCoordinates/UniformThetaSubstitution/FirstDeriv.lean}{Formalism.InvariantFactorPaths.G7.DeS}
\leanref{Formalism/ThetaLogCoordinates/UniformThetaSubstitution/FirstDeriv.lean}{Formalism.InvariantFactorPaths.G7.Hser}
\leanref{Formalism/ThetaLogCoordinates/UniformThetaSubstitution/FirstDeriv.lean}{Formalism.InvariantFactorPaths.G7.DerivIs}

\begin{lemma}\label{pd-orbit-class-facts}
Let $p$ be an odd prime, with the notation of Definition~\ref{pd-orbit-class-law}.
\begin{enumerate}
\item The formal group law $F$ is integral, and it is the only series without constant term with
  $\ell(F)=\ell(X_0)+\ell(X_1)$. It satisfies $F\equiv X_0+X_1$ modulo terms of degree $2$, $F(X_0,0)=X_0$,
  $F^\sigma=F$, and $F(\zeta X_0,\zeta X_1)=\zeta F(X_0,X_1)$ for every $\zeta\in K_0$ with
  $\zeta^{p^2-1}=1$.
\item $\varrho(X)=\ell(X^p)/p$, $\varrho(X^p)=\ell(X)-X$, the derivative $\ell'$ is integral, and
  $A(X_0^p,X_1^p)=p\big(\varrho(X_0)\ell(X_1)-A(X_1,X_0)\big)$. Hence, for $c,c'\in K_0$,
  \[
   G_{c',c}(X_0^p,X_1^p)-pG_{c,c'}=p\big(c'\varrho(X_0)\ell(X_1)-c\,\ell(X_0)\varrho(X_1)\big).
  \]
\item (Dwork's lemma \cite{Haz78}) A power series $f$ over $K_0$ without constant term has $\exp(f)$
  integral if and only if $f^\sigma(X^p)-pf$ is $p$ times an integral series, where $X^p$ is substituted
  for each variable.
\item If $f$ is integral, then $f^\sigma(X^p)-f^p$ is $p$ times an integral series.
\item The relation $\simeq_0$ is an equivalence relation. If $f\simeq_0g$ and $f'\simeq_0g'$, then
  $f+f'\simeq_0g+g'$, and $f(b)\simeq_0g(b)$ for every family $b$ of integral series without constant
  term. If $f\simeq_00$ and $N$ is an integer prime to $p$, then $N^{-1}f\simeq_00$.
\item If $f\simeq_00$ and $b$ is a family of integral series without constant term with $f(b)=0$, then
  $(\partial f/\partial X_e)(b)$ is integral for every variable $X_e$.
\item Let $\Phi$ be a power series over $K_0$ with $\partial\Phi/\partial X_x$ integral, $\phi$ a family
  of integral series without constant term, one for each variable of $\Phi$, and $Y$ an integral series
  without constant term. Let $\phi'$ be $\phi$ with $\phi_x$ replaced by $\phi_x+pY$. Then
  $\Phi(\phi')-\Phi(\phi)$ is $p$ times an integral series.
\item $\delta_F\ell=0$, and $\delta_F\varrho$ is integral.
\end{enumerate}
\end{lemma}
\leanref{Formalism/ThetaLogCoordinates/UniformThetaSubstitution/FrobeniusLiftNorms/CoeffSubstLgFOfDegreeOne.lean}{Formalism.InvariantFactorPaths.G7.formal_law_exists}
\leanref{Formalism/ThetaLogCoordinates/UniformThetaSubstitution/FrobeniusLiftNorms/CoeffSubstLgFOfDegreeOne.lean}{Formalism.InvariantFactorPaths.G7.formal_law_symmetric}
\leanref{Formalism/ThetaLogCoordinates/UniformThetaSubstitution/FrobeniusLiftNorms/TopsChoicePartA.lean}{Formalism.InvariantFactorPaths.G7.g_zero_relation_holds}
\leanref{Formalism/ThetaLogCoordinates/UniformThetaSubstitution/FrobeniusLiftNorms/TopsChoicePartA.lean}{Formalism.InvariantFactorPaths.G7.identity_b}
\leanref{Formalism/ThetaLogCoordinates/UniformThetaSubstitution/FrobeniusLiftNorms/TopsChoicePartA.lean}{Formalism.InvariantFactorPaths.G7.frobenius_expression_g}
\leanref{Formalism/ThetaLogCoordinates/UniformThetaSubstitution/FrobeniusLiftNorms/LocalComponentGL.lean}{Formalism.InvariantFactorPaths.G7.dwork_multivariable}
\leanref{Formalism/ThetaLogCoordinates/UniformThetaSubstitution/FrobeniusLiftNorms/LocalComponentGL.lean}{Formalism.InvariantFactorPaths.G7.DworkIntegrality.frob_congr}
\leanref{Formalism/ThetaLogCoordinates/UniformThetaSubstitution/FrobeniusLiftNorms/LogarithmEquivalence.lean}{Formalism.InvariantFactorPaths.G7.LogarithmEquivalence.simeqRules}
\leanref{Formalism/ThetaLogCoordinates/UniformThetaSubstitution/FrobeniusLiftNorms/ScalingSwapLemmas.lean}{Formalism.InvariantFactorPaths.G7.LamModule_of}
\leanref{Formalism/ThetaLogCoordinates/UniformThetaSubstitution/FrobeniusLiftNorms/ConstantCoeffLawFirstTwoOfThree.lean}{Formalism.InvariantFactorPaths.G7.GramFrameCompositionLinFacts.integral_lin}
\leanref{Formalism/ThetaLogCoordinates/UniformThetaSubstitution/FrobeniusLiftNorms/IntegralPderivLgF.lean}{Formalism.InvariantFactorPaths.G7.divided_derivative}
\leanref{Formalism/ThetaLogCoordinates/UniformThetaSubstitution/FrobeniusLiftNorms/IntegralPderivLgF.lean}{Formalism.InvariantFactorPaths.G7.DeIntegral_of}

\begin{proof}
(1) The functional equation lemma \cite{Haz78}, applied to $\ell$, shows that $F$ is integral, and $F$ is
the only solution because $\ell$ has a compositional inverse. As $\ell(X)\equiv X$ modulo degree $2$,
$F\equiv X_0+X_1$. The series $F(X_0,0)$ satisfies $\ell(F(X_0,0))=\ell(X_0)$, so it is $X_0$. The series
$F^\sigma$ satisfies the defining equation of $F$ because $\ell^\sigma=\ell$, and so does
$\zeta^{-1}F(\zeta X_0,\zeta X_1)$ because $\zeta^{p^{2j}}=\zeta$ gives $\ell(\zeta X)=\zeta\ell(X)$; by
uniqueness both are $F$.

(2) The first two identities and the integrality of $\ell'(X)=\sum_jp^jX^{p^{2j}-1}$ are read off the
coefficients. For the third, $A(X_0^p,X_1^p)=\sum_{i\leq j}p^{-(i+j+1)}X_0^{p^{2i+1}}X_1^{p^{2j+2}}$ is the
part with $j'>i$ of $p\,\varrho(X_0)\ell(X_1)=\sum_{i,j'\geq0}p^{-(i+j')}X_0^{p^{2i+1}}X_1^{p^{2j'}}$, and
the part with $j'\leq i$ is $pA(X_1,X_0)$. Exchanging $X_0$ and $X_1$ gives $A(X_1^p,X_0^p)=
p(\ell(X_0)\varrho(X_1)-A(X_0,X_1))$, and the identity for $G$ follows by linearity.

(3) is Dwork's lemma, as cited.

(4) For a polynomial, $(\sum_ma_mX^m)^p\equiv\sum_ma_m^pX^{pm}$ modulo $p$, and $a^p\equiv\sigma(a)$
modulo $pW(\overline{\mathbf F}_p)$ for $a\in W(\overline{\mathbf F}_p)$. A given coefficient of $f^p$ and
of $f^\sigma(X^p)$ depends only on a truncation of $f$.

(5) If $\exp(f-g)$ is integral, so is its inverse $\exp(g-f)$, since its constant term is $1$; and $\exp$
turns sums into products and commutes with substitutions. For the last claim,
$(N^{-1}f)^\sigma(X^p)-pN^{-1}f=N^{-1}(f^\sigma(X^p)-pf)$ with $N^{-1}\in W(\overline{\mathbf F}_p)$, so
(3) applies in both directions.

(6) $\partial\exp(f)/\partial X_e=\exp(f)\,\partial f/\partial X_e$, and $\exp(f)(b)=\exp(f(b))=1$. So
$(\partial f/\partial X_e)(b)=(\partial\exp(f)/\partial X_e)(b)$, which is integral because $\exp(f)$ is.

(7) The contribution of a monomial $X^d$ of $\Phi$, with coefficient $\Phi_d$, is
$\Phi_d\big((\phi_x+pY)^{d_x}-\phi_x^{d_x}\big)\prod_{s\neq x}\phi_s^{d_s}$, and
$(\phi_x+pY)^m-\phi_x^m=\sum_{k=1}^m\binom mkp^kY^k\phi_x^{m-k}$. Here
$\Phi_d\binom mkp^k=m\Phi_d\binom{m-1}{k-1}p^k/k$, where $m\Phi_d$ with $m=d_x$ is a coefficient of
$\partial\Phi/\partial X_x$, hence in $W(\overline{\mathbf F}_p)$, and $p^k/k\in pW(\overline{\mathbf F}_p)$
because $v_p(k)<k$.

(8) $\delta_F\ell=0$ is the definition of $F$. By (2), $\delta_F\varrho=p^{-1}\big(\ell(F^p)-\ell(X_0^p)
-\ell(X_1^p)\big)=p^{-1}\big(\ell(F^p)-\ell(F(X_0^p,X_1^p))\big)$. By (1) and (4),
$F(X_0^p,X_1^p)=F^\sigma(X_0^p,X_1^p)=F^p+pr$ with $r$ integral and without constant term. Since $\ell'$ is
integral, (7) with $\Phi=\ell$ shows that $\ell(F^p+pr)-\ell(F^p)$ is $p$ times an integral series, so
$\delta_F\varrho$ is integral.
\end{proof}

\begin{lemma}\label{pd-orbit-class-one-variable}
Let $p$ be an odd prime, with the notation of Definition~\ref{pd-orbit-class-law}, and let $f\in K_0[[X]]$
have constant term $0$ and $\delta_Ff$ integral. Then $f=\alpha\varrho+\beta\ell+R$ with
$\alpha\in W(\overline{\mathbf F}_p)$, $\beta\in K_0$ and $R$ integral.
\end{lemma}
\leanref{Formalism/ThetaLogCoordinates/UniformThetaSubstitution/FrobeniusLiftNorms/NotDvdChoosePowMul.lean}{Formalism.InvariantFactorPaths.G7.LemmaOne_of}
\leanref{Formalism/ThetaLogCoordinates/UniformThetaSubstitution/FrobeniusLiftNorms/NotDvdChoosePowMul.lean}{Formalism.InvariantFactorPaths.G7.LemOne.exists_limit}
\leanref{Formalism/ThetaLogCoordinates/UniformThetaSubstitution/FrobeniusLiftNorms/NotDvdChoosePowMul.lean}{Formalism.InvariantFactorPaths.G7.LemOne.good_succ}
\leanref{Formalism/ThetaLogCoordinates/UniformThetaSubstitution/FrobeniusLiftNorms/NotDvdChoosePowMul.lean}{Formalism.InvariantFactorPaths.G7.LemOne.norm_coeff_mul_choose}

\begin{proof}
Write $h_n$ for the coefficient of $X^n$ in a series $h$, and $W=W(\overline{\mathbf F}_p)$.

\emph{Step 1: one coefficient.} Let $h\in K_0[[X]]$ have constant term $0$, $\delta_Fh$ integral and
$h_d\in W$ for $d<n$, and let $0<i<n$. The monomial $X_0^iX_1^{n-i}$ does not occur in $h(X_0)$ or
$h(X_1)$, so its coefficient in $\delta_Fh$ is $\sum_{d\leq n}h_d\,c_d$, with $c_d$ its coefficient in
$F^d$; here $F^d$ has order $d$ and is integral. As $F\equiv X_0+X_1$ modulo degree $2$
(Lemma~\ref{pd-orbit-class-facts}(1)), $c_n=\binom ni$. Hence $\binom nih_n\in W$. If $n\geq2$ is not a
power of $p$, write $n=p^km$ with $p\nmid m$ and $m\geq2$; then $\binom n{p^k}$ is prime to $p$, so
$h_n\in W$. If $n=p^{k+1}$, then $\binom n{p^k}$ has $p$-adic valuation $1$, so $ph_n\in W$. By induction
on the degree, this gives: (a) if $h_d\in W$ for $d\leq p^k$, then $h_d\in W$ for $d<p^{k+1}$; and (b) if
$h_d\in W$ for $d<p^{k+1}$, then $ph_{p^{k+1}}\in W$.

\emph{Step 2: approximation.} For $a\in W$ and $b\in K_0$ put $r_{a,b}=f-a\varrho-b\ell$. By
Lemma~\ref{pd-orbit-class-facts}(8), $\delta_Fr_{a,b}=\delta_Ff-a\,\delta_F\varrho$ is integral. Call
$(a,b)$ \emph{good at} $j$ if $(r_{a,b})_d\in W$ for $d<p^{2j+1}$. The pair $(0,f_1)$ is good at $0$:
$r_{0,f_1}$ has coefficients $0$ in degrees $0$ and $1$, and (a) with $k=0$ applies. Let $(a,b)$ be good
at $j$. By (b), $\mu=p\,(r_{a,b})_{p^{2j+1}}\in W$; put $a'=a+\mu p^j$. The coefficient of $\varrho$ at
$X^{p^{2j+1}}$ is $p^{-(j+1)}$, and $p^j\varrho_d\in W$ for $d<p^{2j+1}$. So $r_{a',b}=r_{a,b}-\mu
p^j\varrho$ has coefficients in $W$ in degrees $<p^{2j+1}$ and coefficient $0$ in degree $p^{2j+1}$, and
by (a) in degrees $<p^{2j+2}$. By (b), $\nu=p\,(r_{a',b})_{p^{2j+2}}\in W$; put $b'=b+\nu p^j$. The
coefficient of $\ell$ at $X^{p^{2j+2}}$ is $p^{-(j+1)}$, and $p^j\ell_d\in W$ for $d<p^{2j+2}$, so in the
same way $(a',b')$ is good at $j+1$. Moreover $a'-a$ and $b'-b$ lie in $p^jW$.

\emph{Step 3: the limit.} Starting from $(0,f_1)$, Step 2 gives pairs $(a_j,b_j)$ good at $j$ with
$a_{j+1}-a_j$ and $b_{j+1}-b_j$ in $p^jW$. Since $K_0$ is complete, they converge to some $\alpha\in W$
and $\beta\in K_0$. For each $n$, $(r_{a_j,b_j})_n\in W$ as soon as $p^{2j+1}>n$, and it tends to
$(r_{\alpha,\beta})_n$. So $R=r_{\alpha,\beta}$ is integral, and $f=\alpha\varrho+\beta\ell+R$.
\end{proof}

\begin{lemma}\label{pd-orbit-class-two-variables}
Let $p$ be an odd prime, with the notation of Definition~\ref{pd-orbit-class-law}, and let
$H\in K_0[[X_0,X_1]]$ vanish on the axes, with $\delta_FH$ and $\delta_F(H^{\mathrm{sw}})$ integral. Then
$H=\beta\,\ell(X_0)\ell(X_1)+R$ with $\beta\in K_0$ and $R$ integral.
\end{lemma}
\leanref{Formalism/ThetaLogCoordinates/UniformThetaSubstitution/FrobeniusLiftNorms/LemmaTwoOf.lean}{Formalism.InvariantFactorPaths.G7.LemmaTwo_of}
\leanref{Formalism/ThetaLogCoordinates/UniformThetaSubstitution/FrobeniusLiftNorms/LemmaTwoOf.lean}{Formalism.InvariantFactorPaths.G7.LemTwo.coeff_partial_coboundary}

\begin{proof}
Write $W=W(\overline{\mathbf F}_p)$, $\ell_n$ and $\varrho_n$ for the coefficients of $X^n$ in $\ell$ and
$\varrho$, so that $\ell_{p^{2k}}=p^{-k}$, $\varrho_{p^{2k+1}}=p^{-(k+1)}$ and all others vanish, and
$H=\sum_mh_m(X_0)X_1^m=\sum_ih'_i(X_1)X_0^i$. The variable $X_2$ enters $\delta_FH$ only through the second
argument of $H$, so the coefficient of $X_0^aX_1^bX_2^m$ in $\delta_FH$ is that of $X_0^aX_1^b$ in
$\delta_Fh_m$. Hence $\delta_Fh_m$ is integral, and likewise $\delta_Fh'_i$, from $\delta_F(H^{\mathrm{sw}})$.
As $H$ vanishes on the axes, the $h_m$ and $h'_i$ have constant term $0$. By
Lemma~\ref{pd-orbit-class-one-variable}, $h_m=\alpha_m\varrho+\beta_m\ell+R_m$ and
$h'_i=\alpha'_i\varrho+\beta'_i\ell+R'_i$ with $\alpha_m,\alpha'_i\in W$ and $R_m$, $R'_i$ integral.
Comparing the two expressions for the coefficient of $X_0^iX_1^m$ in $H$,
\[
 (\ast)\qquad\alpha_m\varrho_i+\beta_m\ell_i-\alpha'_i\varrho_m-\beta'_i\ell_m\in W\quad\text{for all }i,m.
\]
An element $x\in K_0$ with $p^{-k}x\in p^{-M}W$ for a fixed $M$ and infinitely many $k$ is $0$.

(a) For $i=p^{2k+1}$ and $m=p^{2N}$, $(\ast)$ reads $\alpha_mp^{-(k+1)}-\beta'_ip^{-N}\in W$. As
$\alpha_m\in W$, $p^{-N}\beta'_i\in p^{-(k+1)}W$ for every $N$, so $\beta'_{p^{2k+1}}=0$.

(b) For $i=p^{2N+1}$ and any $m$, $(\ast)$ and (a) give $\alpha_mp^{-(N+1)}-\alpha'_i\varrho_m\in W$, so
$p^{-(N+1)}\alpha_m\in W+\varrho_mW$ for every $N$, and $\alpha_m=0$.

(c) If $m$ is not a power of $p^2$, then for $i=p^{2N}$, $(\ast)$ and (b) give
$\beta_mp^{-N}-\alpha'_i\varrho_m\in W$ for every $N$, so $\beta_m=0$.

(d) For $i=p^{2N}$, $(\ast)$ with $m=p^{2l}$ and with $m=1$ gives
$\beta_{p^{2l}}p^{-N}-\beta'_ip^{-l}\in W$ and $\beta_1p^{-N}-\beta'_i\in W$. Hence
$(\beta_{p^{2l}}-p^{-l}\beta_1)p^{-N}\in p^{-l}W$ for every $N$, and $\beta_{p^{2l}}=p^{-l}\beta_1$.

By (b), (c) and (d), $h_m=\beta_1\ell_m\ell+R_m$ for every $m$. So $H-\beta_1\ell(X_0)\ell(X_1)=\sum_m
R_m(X_0)X_1^m$ is integral.
\end{proof}

\begin{lemma}\label{pd-orbit-class-average}
Let $p$ be an odd prime, with the notation of Definition~\ref{pd-orbit-class-law}, and let
$b\in K_0[[X_0,X_1]]$ vanish on the axes, with $\delta_Fb\simeq_00$, $\delta_F(b^{\mathrm{sw}})\simeq_00$
and $b(\zeta X_0,X_1)\simeq_0b(X_0,\zeta^pX_1)$ for every $\zeta\in K_0$ with $\zeta^{p^2-1}=1$. Then there
is $G\in K_0[[X_0,X_1]]$ with balanced support such that $G\simeq_0b$, $\delta_FG\simeq_00$ and
$\delta_F(G^{\mathrm{sw}})\simeq_00$.
\end{lemma}
\leanref{Formalism/ThetaLogCoordinates/UniformThetaSubstitution/FrobeniusLiftNorms/ScalingSwapLemmas.lean}{Formalism.InvariantFactorPaths.G7.StepZero_of}

\begin{proof}
Let $N=p^2-1$, which is prime to $p$, and let $\zeta\in W(\overline{\mathbf F}_p)$ be a primitive $N$-th
root of unity, a Teichm\"uller representative. Put
\[
 G=N^{-1}\sum_{k=0}^{N-1}b_k,\qquad b_k(X_0,X_1)=b(\zeta^kX_0,\zeta^{-pk}X_1).
\]
\emph{Support.} The coefficient of $X_0^{n_0}X_1^{n_1}$ in $G$ is $N^{-1}\sum_k(\zeta^{n_0-pn_1})^k$ times
that in $b$. The sum vanishes unless $\zeta^{n_0-pn_1}=1$, that is $n_0\equiv pn_1\pmod N$. As $b$
vanishes on the axes, $G$ has balanced support.

\emph{$G\simeq_0b$.} The hypothesis for $\zeta^k$ gives $b(\zeta^kX_0,X_1)\simeq_0b(X_0,\zeta^{pk}X_1)$.
Substituting $\zeta^{-pk}X_1$ for $X_1$ gives $b_k\simeq_0b$ (Lemma~\ref{pd-orbit-class-facts}(5)). So
$G-b=N^{-1}\sum_k(b_k-b)\simeq_00$, by the same lemma.

\emph{Coboundaries.} Since $F(\zeta^kX_0,\zeta^kX_1)=\zeta^kF(X_0,X_1)$ (Lemma~\ref{pd-orbit-class-facts}(1)),
$\delta_Fb_k$ is $\delta_Fb$ with $\zeta^kX_0$, $\zeta^kX_1$, $\zeta^{-pk}X_2$ substituted for $X_0$,
$X_1$, $X_2$, so $\delta_Fb_k\simeq_00$, and $\delta_FG=N^{-1}\sum_k\delta_Fb_k\simeq_00$. Likewise
$G^{\mathrm{sw}}=N^{-1}\sum_kb^{\mathrm{sw}}(\zeta^{-pk}X_0,\zeta^kX_1)$, and $\zeta^{-pk}$ is also a
$(p^2-1)$-th root of unity, so $\delta_F(G^{\mathrm{sw}})\simeq_00$.
\end{proof}

\begin{lemma}\label{pd-orbit-class-derivative}
Let $p$ be an odd prime, with the notation of Definition~\ref{pd-orbit-class-law}, and let
$G\in K_0[[X_0,X_1]]$ have balanced support, with $\delta_FG\simeq_00$ and
$\delta_F(G^{\mathrm{sw}})\simeq_00$. Then $G$ has derivative $c$ for some
$c\in W(\overline{\mathbf F}_p)$.
\end{lemma}
\leanref{Formalism/ThetaLogCoordinates/UniformThetaSubstitution/FrobeniusLiftNorms/ConstantCoeffLawFirstTwoOfThree.lean}{Formalism.InvariantFactorPaths.G7.StepI_of}
\leanref{Formalism/ThetaLogCoordinates/UniformThetaSubstitution/FrobeniusLiftNorms/ConstantCoeffLawFirstTwoOfThree.lean}{Formalism.InvariantFactorPaths.G7.FirstDeriv_of}
\leanref{Formalism/ThetaLogCoordinates/UniformThetaSubstitution/FrobeniusLiftNorms/ConstantCoeffLawFirstTwoOfThree.lean}{Formalism.InvariantFactorPaths.G7.CobLin_of}

\begin{proof}
Let $g_{n_0,n_1}$ be the coefficients of $G$, and $f(X)=\sum_jg_{1,j}X^j=(\partial G/\partial X_0)(0,X)$.
As $G$ has balanced support, it vanishes on the axes.

\emph{Step 1.} The substitution of $(X_0,0,X_1)$ for $(X_0,X_1,X_2)$ kills $\delta_FG$, since
$F(X_0,0)=X_0$ and $G(0,X_1)=0$. Differentiating $\delta_FG$ in $X_1$ and then substituting,
\[
 \frac{\partial G}{\partial X_0}(X_0,X_1)\,\phi(X_0)-f(X_1),\qquad
 \phi(X_0)=\frac{\partial F}{\partial X_1}(X_0,0),
\]
is integral by Lemma~\ref{pd-orbit-class-facts}(6). Differentiating $\ell(F(X_0,X_1))=\ell(X_0)+\ell(X_1)$
in $X_1$ at $X_1=0$ gives $\ell'(X_0)\phi(X_0)=1$. So $\partial G/\partial X_0-\ell'(X_0)f(X_1)$ is
$\ell'(X_0)$ times this series, and it is integral since $\ell'$ is.

\emph{Step 2.} The substitution of $(X_0,X_1,0)$ kills $\delta_F(G^{\mathrm{sw}})$, as $G(0,X_1)=0$. The
derivative of $\delta_F(G^{\mathrm{sw}})$ in $X_2$, at $X_2=0$, is $f(F(X_0,X_1))-f(X_0)-f(X_1)=\delta_Ff$.
By Lemma~\ref{pd-orbit-class-facts}(6), $\delta_Ff$ is integral, and $f$ has constant term $g_{1,0}=0$.

\emph{Step 3.} By Lemma~\ref{pd-orbit-class-one-variable}, $f=c\varrho+\beta\ell+R$ with
$c\in W(\overline{\mathbf F}_p)$ and $R$ integral. The coefficient of $X^j$ in $f-c\varrho$ is integral for
every $j$. If $j\equiv p\pmod{p^2-1}$, then $\ell$ has no monomial $X^j$, as $p^{2i}\equiv1\not\equiv p$,
and the coefficient is that of $R$. Otherwise $g_{1,j}=0$ by balanced support, since $1\equiv pj$ would
give $j\equiv p$, and $\varrho$ has no monomial $X^j$, as $p^{2i+1}\equiv p$. Hence
$\partial G/\partial X_0-c\,\ell'(X_0)\varrho(X_1)=\big(\partial G/\partial X_0-\ell'(X_0)f(X_1)\big)+
\ell'(X_0)(f-c\varrho)(X_1)$ is integral.
\end{proof}

\begin{lemma}\label{pd-orbit-class-frobenius-twist}
Let $p$ be an odd prime, with the notation of Definition~\ref{pd-orbit-class-law}, let
$c\in W(\overline{\mathbf F}_p)$, and let $G\in K_0[[X_0,X_1]]$ vanish on the axes, have derivative $c$,
and satisfy $\delta_FG\simeq_00$. Then
\[
 \delta_F(\mathrm D_\sigma G)-\sigma(c)\,(\delta_F\varrho)(X_0,X_1)\,\ell(X_2)
\]
is integral.
\end{lemma}
\leanref{Formalism/ThetaLogCoordinates/UniformThetaSubstitution/FrobeniusLiftNorms/PartialCoboundaryLinear.lean}{Formalism.InvariantFactorPaths.G7.StepII_of}

\begin{proof}
Write $F^p=F(X_0,X_1)^p$ and $\tilde F=F(X_0^p,X_1^p)$. By Lemma~\ref{pd-orbit-class-facts}(1) and (4),
$\tilde F=F^\sigma(X_0^p,X_1^p)=F^p+pr$ with $r$ integral without constant term. Put
$\Phi=G^\sigma-\sigma(c)\ell(X_0)\varrho(X_1)$.

(a) As $\ell$ and $\varrho$ have rational coefficients, $\partial\Phi/\partial X_0=\big(\partial G/\partial
X_0-c\,\ell'(X_0)\varrho(X_1)\big)^\sigma$, which is integral because $G$ has derivative $c$. By
Lemma~\ref{pd-orbit-class-facts}(7) with $\phi=(F^p,X_2^p)$ and $Y=r$, the series
$E_1=\Phi(\tilde F,X_2^p)-\Phi(F^p,X_2^p)$ is $p$ times an integral series.

(b) As $F^\sigma=F$, $(\delta_FG)^\sigma(X_0^p,X_1^p,X_2^p)=G^\sigma(\tilde F,X_2^p)-G^\sigma(X_0^p,X_2^p)
-G^\sigma(X_1^p,X_2^p)$. Dwork's lemma (Lemma~\ref{pd-orbit-class-facts}(3)) for $\delta_FG\simeq_00$
shows that $E_2=(\delta_FG)^\sigma(X_0^p,X_1^p,X_2^p)-p\,\delta_FG$ is $p$ times an integral series.

(c) By definition, $\delta_F(\mathrm D_\sigma G)=p^{-1}\big(G^\sigma(F^p,X_2^p)-G^\sigma(X_0^p,X_2^p)
-G^\sigma(X_1^p,X_2^p)\big)-\delta_FG$. Comparing with (b) and with the definition of $\Phi$,
\[
 \delta_F(\mathrm D_\sigma G)=p^{-1}E_2-p^{-1}E_1-\sigma(c)\,p^{-1}\big(\ell(\tilde F)-\ell(F^p)\big)
 \varrho(X_2^p).
\]
By Lemma~\ref{pd-orbit-class-facts}(2), $\varrho(X_2^p)=\ell(X_2)-X_2$, and
$p^{-1}\big(\ell(F^p)-\ell(\tilde F)\big)=p^{-1}\big(\ell(F^p)-\ell(X_0^p)-\ell(X_1^p)\big)=\delta_F\varrho$.
Hence
\[
 \delta_F(\mathrm D_\sigma G)-\sigma(c)\,\delta_F\varrho\,\ell(X_2)=p^{-1}(E_2-E_1)-\sigma(c)\,\delta_F\varrho
 \,X_2,
\]
which is integral by (a), (b) and Lemma~\ref{pd-orbit-class-facts}(8).
\end{proof}

\begin{lemma}\label{pd-orbit-class-integral}
Let $p$ be an odd prime, with the notation of Definition~\ref{pd-orbit-class-law}, let
$\gamma,\gamma'\in K_0$, and let $G\in K_0[[X_0,X_1]]$ have balanced support, with
$\delta_F(\mathrm D_\sigma G)-\gamma\,\delta_F\varrho\,\ell(X_2)$ and
$\delta_F(\mathrm D_\sigma G^{\mathrm{sw}})-\gamma'\,\delta_F\varrho\,\ell(X_2)$ integral. Then
$\mathrm Q_{\gamma,\gamma'}G$ is integral.
\end{lemma}
\leanref{Formalism/ThetaLogCoordinates/UniformThetaSubstitution/FrobeniusLiftNorms/ConstantCoeffFirstTwoOfThree.lean}{Formalism.InvariantFactorPaths.G7.StepH_of}
\leanref{Formalism/ThetaLogCoordinates/UniformThetaSubstitution/FrobeniusLiftNorms/ScalingSwapLemmas.lean}{Formalism.InvariantFactorPaths.G7.swap_rules}

\begin{proof}
Put $Q=\mathrm Q_{\gamma,\gamma'}G$. It has balanced support: the monomials of $G$ have it; those of
$G^\sigma(X_0^p,X_1^p)$ have exponents $(pn_0,pn_1)$ with $n_0\equiv pn_1$, so $pn_0\equiv p\cdot pn_1$;
those of $\varrho(X_0)\ell(X_1)$ and $\ell(X_0)\varrho(X_1)$ are $(p^{2i+1},p^{2j})$ and
$(p^{2i},p^{2j+1})$, which satisfy the congruence as $p^2\equiv1$. In particular $Q$ vanishes on the axes.
As $\delta_F$ is linear, $\delta_F(\varrho(X_0)\ell(X_1))=\delta_F\varrho\,\ell(X_2)$ and
$\delta_F(\ell(X_0)\varrho(X_1))=\delta_F\ell\,\varrho(X_2)=0$, so
$\delta_FQ=\delta_F(\mathrm D_\sigma G)-\gamma\,\delta_F\varrho\,\ell(X_2)$ is integral. As $\mathrm
D_\sigma$ commutes with the exchange of the variables, $Q^{\mathrm{sw}}=\mathrm Q_{\gamma',\gamma}
(G^{\mathrm{sw}})$, and $\delta_F(Q^{\mathrm{sw}})$ is integral in the same way. By
Lemma~\ref{pd-orbit-class-two-variables}, $Q=\beta\ell(X_0)\ell(X_1)+R$ with $R$ integral. The monomials
$X_0^{p^{2i}}X_1^{p^{2j}}$ of $\ell(X_0)\ell(X_1)$ have $n_0\equiv1\not\equiv p\equiv pn_1$, so they do
not occur in $Q$, and every coefficient of $Q$ is $0$ or the corresponding coefficient of $R$.
\end{proof}

\begin{lemma}\label{pd-orbit-class-constants}
Let $p$ be an odd prime, with the notation of Definition~\ref{pd-orbit-class-law}, let
$c,c'\in W(\overline{\mathbf F}_p)$, and let $G\in K_0[[X_0,X_1]]$ have derivative $c$, with
$\mathrm Q_{\sigma(c),\sigma(c')}G$ integral. Then $c+\sigma(c')=0$.
\end{lemma}
\leanref{Formalism/ThetaLogCoordinates/UniformThetaSubstitution/FrobeniusLiftNorms/ConstantCoeffFirstTwoOfThree.lean}{Formalism.InvariantFactorPaths.G7.c_add_sigma_c_eq_zero}

\begin{proof}
Let $i\geq0$, $m=p^{2i+1}$, and $g$ the coefficient of $X_0X_1^m$ in $G$. The coefficient of $X_1^m$ in
$\partial G/\partial X_0-c\,\ell'(X_0)\varrho(X_1)$ is $g-cp^{-(i+1)}$, as $\ell'(0)=1$ and $\varrho$ has
coefficient $p^{-(i+1)}$ at $X^m$. The coefficient of $X_0X_1^m$ in $\mathrm Q_{\sigma(c),\sigma(c')}G$ is
$-g-\sigma(c')p^{-(i+1)}$, since $G^\sigma(X_0^p,X_1^p)$ has no monomial of degree $1$ in $X_0$, $\varrho$
has no linear term and $\ell$ has linear term $X$. Both coefficients lie in $W(\overline{\mathbf F}_p)$,
so their sum $-(c+\sigma(c'))p^{-(i+1)}$ does, for every $i$. Hence $c+\sigma(c')=0$.
\end{proof}

\begin{lemma}\label{pd-orbit-class-comparison}
Let $p$ be an odd prime, with the notation of Definition~\ref{pd-orbit-class-law}, let $c\in K_0$ with
$\sigma^2(c)=c$, and let $G\in K_0[[X_0,X_1]]$ vanish on the axes, with $\mathrm Q_{\sigma(c),-c}G$
integral. Then $G\simeq_0G_{c,\sigma(c)}$.
\end{lemma}
\leanref{Formalism/ThetaLogCoordinates/UniformThetaSubstitution/FrobeniusLiftNorms/ConstantCoeffFirstTwoOfThree.lean}{Formalism.InvariantFactorPaths.G7.StepIV_of}

\begin{proof}
Put $f=G-G_{c,\sigma(c)}$; it has constant term $0$. As $\sigma^2(c)=c$, $G_{c,\sigma(c)}^\sigma=
G_{\sigma(c),c}$, so by Lemma~\ref{pd-orbit-class-facts}(2) with $c'=\sigma(c)$,
\[
 f^\sigma(X_0^p,X_1^p)-pf=p\,\mathrm D_\sigma G-p\big(\sigma(c)\varrho(X_0)\ell(X_1)-c\,\ell(X_0)
 \varrho(X_1)\big)=p\,\mathrm Q_{\sigma(c),-c}G,
\]
which is $p$ times an integral series. By Dwork's lemma (Lemma~\ref{pd-orbit-class-facts}(3)),
$\exp(f)$ is integral.
\end{proof}

\begin{theorem}\label{pd-orbit-class-classification}
Let $p$ be an odd prime, with the notation of Definition~\ref{pd-orbit-class-law}, and let
$b\in K_0[[X_0,X_1]]$ vanish on the axes, with $\delta_Fb\simeq_00$, $\delta_F(b^{\mathrm{sw}})\simeq_00$
and $b(\zeta X_0,X_1)\simeq_0b(X_0,\zeta^pX_1)$ for every $\zeta\in K_0$ with $\zeta^{p^2-1}=1$. Then there
is $c\in W(\overline{\mathbf F}_p)$ with $\sigma^2(c)=c$ and $b\simeq_0G_{c,\sigma(c)}$.
\end{theorem}
\leanref{Formalism/ThetaLogCoordinates/UniformThetaSubstitution/FinalS.lean}{Formalism.InvariantFactorPaths.G7.Top.betaIsG}
\leanref{Formalism/ThetaLogCoordinates/UniformThetaSubstitution/FrobeniusLiftNorms/NormLiftLeOne.lean}{Formalism.InvariantFactorPaths.G7.BetaIsG_of}

\begin{proof}
By Lemma~\ref{pd-orbit-class-average} there is $G$ with balanced support, $G\simeq_0b$,
$\delta_FG\simeq_00$ and $\delta_F(G^{\mathrm{sw}})\simeq_00$. The series $G^{\mathrm{sw}}$ has the same
properties, since $n_0\equiv pn_1$ implies $n_1\equiv p^2n_1\equiv pn_0$. By
Lemma~\ref{pd-orbit-class-derivative}, $G$ has derivative $c$ and $G^{\mathrm{sw}}$ has derivative $c'$, for
some $c,c'\in W(\overline{\mathbf F}_p)$. By Lemma~\ref{pd-orbit-class-frobenius-twist}, applied to $G$ and
to $G^{\mathrm{sw}}$, the hypotheses of Lemma~\ref{pd-orbit-class-integral} hold with $\gamma=\sigma(c)$
and $\gamma'=\sigma(c')$, so $\mathrm Q_{\sigma(c),\sigma(c')}G$ is integral. By
Lemma~\ref{pd-orbit-class-constants}, $c+\sigma(c')=0$. The series $\mathrm Q_{\sigma(c'),\sigma(c)}
(G^{\mathrm{sw}})=(\mathrm Q_{\sigma(c),\sigma(c')}G)^{\mathrm{sw}}$ is integral too, and
Lemma~\ref{pd-orbit-class-constants} for $G^{\mathrm{sw}}$ gives $c'+\sigma(c)=0$. Hence $\sigma(c')=-c$
and $\sigma(c')=-\sigma^2(c)$, so $\sigma^2(c)=c$, and $\mathrm Q_{\sigma(c),-c}G$ is integral. By
Lemma~\ref{pd-orbit-class-comparison}, $G\simeq_0G_{c,\sigma(c)}$, and then $b\simeq_0G_{c,\sigma(c)}$ by
Lemma~\ref{pd-orbit-class-facts}(5).
\end{proof}

\begin{lemma}\label{pd-orbit-class-beta}
Let $p\geq5$ be a prime, with the notation of Definition~\ref{pd-orbit-class-law}, and let
$\beta\in K_0[[Y_0,Y_1]]$ be bi-additive (Definition~\ref{pd-formal-group-system}). Then there is
$c\in W(\overline{\mathbf F}_p)$ with $\sigma^2(c)=c$ and
$\beta(\ell(X_0),\ell(X_1))\simeq_0G_{c,\sigma(c)}(X_0,X_1)$.
\end{lemma}
\leanref{Formalism/SharpRestriction/GrowthBounds/MultisetValuation/DworkExponentialRestrictions.lean}{Formalism.MainNotions.FormalGroupSystem.Pf.RestBeta}
\leanref{Formalism/SharpRestriction/GrowthBounds/MultisetValuation/BiadditiveSecondVariable.lean}{Formalism.MainNotions.FormalGroupSystem.Pf.restBeta}
\leanref{Formalism/SharpRestriction/GrowthBounds/MultisetValuation/BiadditiveSecondVariable.lean}{Formalism.MainNotions.FormalGroupSystem.Pf.sx_betaHyp}

\begin{proof}
Let $b(X_0,X_1)=\beta(\ell(X_0),\ell(X_1))$; it satisfies the hypotheses of
Theorem~\ref{pd-orbit-class-classification}. Every monomial of $\beta$ has positive degree in $Y_0$ and in
$Y_1$, and $X$ divides $\ell(X)$, so $b$ vanishes on the axes. Since $\ell(F(X_0,X_1))=\ell(X_0)+\ell(X_1)$,
\[
 \delta_Fb=\big(\beta(Y_0+Y_1,Y_2)-\beta(Y_0,Y_2)-\beta(Y_1,Y_2)\big)\big|_{Y_i=\ell(X_i)},
\]
and this is $\simeq_00$ because the series in the $Y_i$ is $\simeq0$. The additivity of $\beta$ in the
second variable gives $\delta_F(b^{\mathrm{sw}})\simeq_00$ in the same way. Let $\zeta\in K_0$ with
$\zeta^{p^2-1}=1$. Then $\zeta\in W(\overline{\mathbf F}_p)$ is a Teichm\"uller representative, so
$\sigma(\zeta)=\zeta^p$ and $\sigma^2(\zeta)=\zeta$. As $\ell(\zeta X)=\zeta\ell(X)$ and
$\ell(\zeta^pX)=\zeta^p\ell(X)$, the series $b(\zeta X_0,X_1)-b(X_0,\zeta^pX_1)$ is
$\beta(\zeta Y_0,Y_1)-\beta(Y_0,\sigma(\zeta)Y_1)$ at $Y_i=\ell(X_i)$, which is $\simeq_00$ by
bi-additivity. Theorem~\ref{pd-orbit-class-classification} gives $c$.
\end{proof}

\begin{definition}\label{pd-orbit-class-series}
Let $p$ be a prime, $W(\overline{\mathbf F}_p)$ the ring of Witt vectors with fraction field $K_0$,
$\ell(X)=\sum_{j\geq0}X^{p^{2j}}/p^j$, $K'$ a field containing $K_0$ and $\pi\in K'$ with
$\pi^{p^2-1}=p$.
\begin{enumerate}
\item The \emph{rescaled logarithm} is $\lambda_\pi(X)=\ell(\pi X)/\pi=\sum_{j\geq0}\pi^{p^{2j}-1}p^{-j}
  X^{p^{2j}}$. Its linear term is $X$, and $\mu_\pi$ is its compositional inverse, the series without
  constant term with $\lambda_\pi(\mu_\pi(X))=X$.
\item The \emph{alternating series} is $\mathrm{alt}(X)=\sum_{l\geq0}(-1)^lX^{p^{2l}}$.
\item $R_\pi\subset K'$ is the subring generated by $\pi$ and $W(\overline{\mathbf F}_p)$. For power series
  $f,g$ over $K'$ in the same variables, $f\equiv g\pmod{pR_\pi}$ means $f-g=pY$ for a power series $Y$
  with coefficients in $R_\pi$.
\end{enumerate}
\end{definition}
\leanref{Formalism/ThetaLogCoordinates/UniformThetaSubstitution/CyclicPowerClassOfExpEquiv.lean}{Formalism.InvariantFactorPaths.G7.lamF}
\leanref{Formalism/ThetaLogCoordinates/UniformThetaSubstitution/CyclicPowerClassOfExpEquiv.lean}{Formalism.InvariantFactorPaths.G7.altF}
\leanref{Formalism/SharpRestriction/GrowthBounds/FormalGroupSystem.lean}{Formalism.MainNotions.FormalGroupSystem.Rp}
\leanref{Formalism/SharpRestriction/GrowthBounds/MultisetValuation/ThPosSubPosAdd.lean}{Formalism.MainNotions.FormalGroupSystem.Pf.Cg}

\begin{lemma}\label{pd-orbit-class-inverse}
Let $p$ be a prime, and let $\pi$, $\lambda_\pi$, $\mu_\pi$ and $\mathrm{alt}$ be as in
Definition~\ref{pd-orbit-class-series}. Then $\lambda_\pi$ and $\mu_\pi$ have coefficients in $\mathbf Z$,
$\lambda_\pi\equiv X+X^{p^2}$ and $\mu_\pi\equiv\mathrm{alt}$, coefficient by coefficient modulo $p$.
\end{lemma}
\leanref{Formalism/SharpRestriction/GrowthBounds/MultisetValuation/SxMuAltInt.lean}{Formalism.MainNotions.FormalGroupSystem.Pf.sx_lam_int}
\leanref{Formalism/SharpRestriction/GrowthBounds/MultisetValuation/SxMuAltInt.lean}{Formalism.MainNotions.FormalGroupSystem.Pf.sx_mu_alt_int}
\leanref{Formalism/SharpRestriction/GrowthBounds/MultisetValuation/TwoSumsCondition.lean}{Formalism.MainNotions.FormalGroupSystem.Pf.th_mu_int}

\begin{proof}
The coefficient of $X^n$ in $\lambda_\pi$ is $0$ unless $n=p^{2j}$, when it is $\pi^{p^{2j}-1}p^{-j}$. Put
$g_j=\sum_{i<j}p^{2i}$, so that $p^{2j}-1=(p^2-1)g_j$ and $\pi^{p^{2j}-1}=p^{g_j}$. Then $g_0=0$, $g_1=1$
and $g_{j+1}=p^2g_j+1\geq j+2$ for $j\geq1$, so the coefficient is the integer $p^{g_j-j}$, equal to $1$
for $j\leq1$ and divisible by $p$ for $j\geq2$.

Let $z_d$ and $m_n$ be the coefficients of $X^d$ in $\lambda_\pi$ and of $X^n$ in $\mu_\pi$, so $z_1=1$ and
$m_0=0$. Comparing coefficients of $X^n$ in $\lambda_\pi(\mu_\pi(X))=X$ gives
\[
 m_n=[n=1]-\sum_{2\leq d\leq n}z_d\,(\mu_\pi^d)_n,
\]
where $(\mu_\pi^d)_n$, for $d\geq2$, is a polynomial with integer coefficients in $m_1,\dots,m_{n-1}$. By
induction on $n$, $m_n\in\mathbf Z$. Modulo $p$, $z_d\equiv[d=p^2]$ for $d\geq2$, and
$(\sum_km_kX^k)^{p^2}\equiv\sum_km_kX^{kp^2}$, so $(\mu_\pi^{p^2})_n\equiv m_{n/p^2}$ if $p^2\mid n$ and
$\equiv0$ otherwise. Hence $m_n+[p^2\mid n]\,m_{n/p^2}\equiv[n=1]\pmod p$ for $n\geq1$. The coefficients
$a_n$ of $\mathrm{alt}$ satisfy $a_n+[p^2\mid n]\,a_{n/p^2}=[n=1]$: for $n=p^{2l}$ with $l\geq1$ the two
terms are $(-1)^l$ and $(-1)^{l-1}$, and if $n\neq1$ is not a power of $p^2$, neither is $n/p^2$. By
induction on $n$, $m_n\equiv a_n\pmod p$.
\end{proof}

\begin{lemma}\label{pd-orbit-class-frobenius}
Let $p$ be a prime, and let $K'$, $\pi$, $R_\pi$ and $\mu_\pi$ be as in
Definition~\ref{pd-orbit-class-series}.
\begin{enumerate}
\item If $x_1,\dots,x_k$ are power series with coefficients in $R_\pi$ and
  $\theta_1,\dots,\theta_k\in W(\overline{\mathbf F}_p)$, then $\big(\sum_i\theta_ix_i\big)^p\equiv\sum_i
  \sigma(\theta_i)x_i^p\pmod{pR_\pi}$.
\item Let $n'\geq2$ be even and $\theta\in W(\overline{\mathbf F}_p)$ with $\sigma^{n'}(\theta)=\theta$, so
  that $\sigma^t(\theta)$ is defined for $t\in\mathbf Z/n'\mathbf Z$. Let $P\subset\mathbf Z/n'\mathbf Z$ be
  the set of even residues or the set of odd residues, $Z_t$ ($t\in\mathbf Z/n'\mathbf Z$) variables, and
  $S_t=\sum_{l\geq0}(-1)^lZ_{t-2l}^{\,p^{2l}}$. Then
  \[
   \mu_\pi\Big(\sum_{t\in P}\sigma^t(\theta)Z_t\Big)\equiv\sum_{t\in P}\sigma^t(\theta)S_t\pmod{pR_\pi}.
  \]
\end{enumerate}
\end{lemma}
\leanref{Formalism/SharpRestriction/GrowthBounds/MultisetValuation/ThLinKConst.lean}{Formalism.MainNotions.FormalGroupSystem.Pf.th_frob_lin}
\leanref{Formalism/SharpRestriction/GrowthBounds/MultisetValuation/ThLinKConst.lean}{Formalism.MainNotions.FormalGroupSystem.Pf.th_claimF}
\leanref{Formalism/SharpRestriction/GrowthBounds/MultisetValuation/ThLinKConst.lean}{Formalism.MainNotions.FormalGroupSystem.Pf.th_claimM}
\leanref{Formalism/SharpRestriction/GrowthBounds/MultisetValuation/FrobeniusCongruenceWitt.lean}{Formalism.MainNotions.FormalGroupSystem.Pf.sx_frob_sum}
\leanref{Formalism/SharpRestriction/GrowthBounds/MultisetValuation/FrobeniusCongruenceWitt.lean}{Formalism.MainNotions.FormalGroupSystem.Pf.frobenius_congruence_witt}

\begin{proof}
(1) In any commutative ring, $(\sum_iy_i)^p=\sum_iy_i^p+pr$ for some $r$, by the binomial theorem and
induction on the number of terms. In the ring of power series with coefficients in $R_\pi$ this gives
$(\sum_i\theta_ix_i)^p\equiv\sum_i\theta_i^px_i^p$. For $\theta\in W(\overline{\mathbf F}_p)$,
$\sigma(\theta)-\theta^p\in pW(\overline{\mathbf F}_p)$, since $\sigma$ induces $x\mapsto x^p$ on
$W(\overline{\mathbf F}_p)/p=\overline{\mathbf F}_p$. So $\theta_i^px_i^p\equiv\sigma(\theta_i)x_i^p$.

(2) Put $L=\sum_{t\in P}\sigma^t(\theta)Z_t$. By induction on $k$,
$L^{p^k}\equiv\sum_{t\in P}\sigma^{t+k}(\theta)Z_t^{p^k}$: if it holds for $k$, raising it to the $p$-th
power keeps the congruence, as both sides have coefficients in $R_\pi$ and $f^p-g^p=(f-g)\sum_if^ig^{p-1-i}$,
and (1) applies to the right side. By Lemma~\ref{pd-orbit-class-inverse}, $\mu_\pi=\mathrm{alt}+pm$ with
$m$ having coefficients in $\mathbf Z$, and $L$ has coefficients in $W(\overline{\mathbf F}_p)$ and no
constant term, so $\mu_\pi(L)\equiv\mathrm{alt}(L)=\sum_l(-1)^lL^{p^{2l}}$. A coefficient of degree $D$
receives contributions only from the $l$ with $p^{2l}\leq D$, so coefficient by coefficient
\[
 \mu_\pi(L)\equiv\sum_{l\geq0}(-1)^l\sum_{s\in P}\sigma^{s+2l}(\theta)Z_s^{p^{2l}}.
\]
Put $t=s+2l$. As $n'$ is even, $s\mapsto s+2l$ is a bijection of $P$, and $\sigma^{s+2l}(\theta)=
\sigma^t(\theta)$ because $\sigma^{n'}(\theta)=\theta$. The right side is therefore
$\sum_{t\in P}\sigma^t(\theta)\sum_l(-1)^lZ_{t-2l}^{p^{2l}}=\sum_{t\in P}\sigma^t(\theta)S_t$.
\end{proof}

\begin{lemma}\label{pd-orbit-class-dual}
Let $p$ be a prime, $n'\geq2$ even, and $\omega_1,\dots,\omega_{n'/2}$,
$\omega^*_1,\dots,\omega^*_{n'/2}\in W(\overline{\mathbf F}_p)$ with
$\sum_{i<n'/2}\sigma^{2i}(\omega_r\omega^*_{r'})=\delta_{rr'}$ for all $r,r'$. Then
$\sum_r\sigma^a(\omega^*_r)\sigma^b(\omega_r)=\delta_{ab}$ for all $0\leq a,b<n'$ with $a\equiv b\pmod2$.
\end{lemma}
\leanref{Formalism/SharpRestriction/GrowthBounds/MultisetValuation/FrobeniusCongruenceWitt.lean}{Formalism.MainNotions.FormalGroupSystem.Pf.sx_dual}
\leanref{Formalism/SharpRestriction/GrowthBounds/MultisetValuation/FrobeniusCongruenceWitt.lean}{Formalism.MainNotions.FormalGroupSystem.Pf.sx_dual_even}

\begin{proof}
Let $A$ and $B$ be the square matrices of size $n'/2$ with entries $A_{ri}=\sigma^{2i}(\omega_r)$ and
$B_{ir'}=\sigma^{2i}(\omega^*_{r'})$, for $0\leq i<n'/2$. As $\sigma^{2i}$ is a ring homomorphism, the
hypothesis says $AB=1$. Then $BA=1$, that is $\sum_r\sigma^{2i}(\omega^*_r)\sigma^{2i'}(\omega_r)=
\delta_{ii'}$, which is the claim for $a=2i$ and $b=2i'$. For $a=2i+1$ and $b=2i'+1$, apply $\sigma$ to
this identity.
\end{proof}

\begin{lemma}\label{pd-orbit-class-two-sums}
Let $p\geq5$ be a prime and take a formal group system at $p$ (Definition~\ref{pd-formal-group-system}).
Let $v$ be an orbit, $j<m_*$, and let $K'$, $\pi$, $R_\pi$ and $\mu_\pi$ be as in
Definition~\ref{pd-orbit-class-series}. For the indices $e$ of the variables at $v$, let $X_e(Z)$ and
$X'_e(Z)$ be the linear forms of the orbit coordinates (Definition~\ref{pd-formal-group-system}(5)), and
put $s_e=\mu_\pi(X_e(Z))$ and $s'_e=\mu_\pi(X'_e(Z))$. Let $c\in W(\overline{\mathbf F}_p)$ with
$\sigma^2(c)=c$, and let $u_s=c$ for the odd positions $s$ of $v$ and $u_s=-\sigma(c)$ for the even ones.
Then the $s_e$ and $s'_e$ have coefficients in $R_\pi$, and
\[
 \sum_e\big(c\,s_es_e'^{\,p}-\sigma(c)\,s'_es_e^{\,p}\big)\equiv H_u\pmod{pR_\pi},
\]
with $H_u$ as in Definition~\ref{pd-kernel-orbit-series}(1).
\end{lemma}
\leanref{Formalism/SharpRestriction/GrowthBounds/MultisetValuation/TwoSumsCondition.lean}{Formalism.MainNotions.FormalGroupSystem.Pf.TwoSumsCondition}
\leanref{Formalism/SharpRestriction/GrowthBounds/MultisetValuation/TripleSumComm.lean}{Formalism.MainNotions.FormalGroupSystem.Pf.restTwoSums_of}
\leanref{Formalism/SharpRestriction/GrowthBounds/MultisetValuation/TripleSumComm.lean}{Formalism.MainNotions.FormalGroupSystem.Pf.th_sumR}
\leanref{Formalism/SharpRestriction/GrowthBounds/MultisetValuation/CommonEigenvectorThird.lean}{Formalism.MainNotions.FormalGroupSystem.Pf.restTwoSums}

\begin{proof}
Write $n'=n'_v$, fix $r'$, $a$, $b$, and write $e=(v,r',a,b)$, $Z_s=Z_{v,s,a,b}$ and $S_s=S_{s,a,b}$. Then
$X_e(Z)=\sum_{s\ \mathrm{even}}\sigma^s(\omega^*_{r'})Z_s$ and $X'_e(Z)=\sum_{s\ \mathrm{odd}}
\sigma^s(\omega_{r'})Z_s$, and $\sigma^{n'}$ fixes $\omega_{r'}$ and $\omega^*_{r'}$. By
Lemma~\ref{pd-orbit-class-frobenius}(2),
\[
 s_e\equiv x_e=\sum_{t\ \mathrm{even}}\sigma^t(\omega^*_{r'})S_t,\qquad s'_e\equiv z_e=\sum_{t\
 \mathrm{odd}}\sigma^t(\omega_{r'})S_t.
\]
The series $x_e$ and $z_e$ have coefficients in $W(\overline{\mathbf F}_p)$, so $s_e$ and $s'_e$ have
coefficients in $R_\pi$. A congruence modulo $pR_\pi$ between series with coefficients in $R_\pi$ is kept
by products and powers, so by Lemma~\ref{pd-orbit-class-frobenius}(1)
\[
 c\,s_es_e'^{\,p}\equiv c\,x_ez_e^p\equiv c\sum_{t\ \mathrm{even}}\ \sum_{t'\ \mathrm{odd}}
 \sigma^t(\omega^*_{r'})\,\sigma^{t'+1}(\omega_{r'})\,S_tS_{t'}^p.
\]
Sum over $r'$. For even $t$ and odd $t'$, both $t$ and the residue of $t'+1$ are even, so by
Lemma~\ref{pd-orbit-class-dual} the sum over $r'$ of the coefficient is $1$ if $t=t'+1$ and $0$ otherwise.
Hence $\sum_{r'}c\,s_es_e'^{\,p}\equiv c\sum_{t'\ \mathrm{odd}}S_{t'}^pS_{t'+1}$. In the same way, with the
even and odd positions exchanged, $\sum_{r'}\sigma(c)\,s'_es_e^{\,p}\equiv\sigma(c)\sum_{t'\
\mathrm{even}}S_{t'}^pS_{t'+1}$. The difference is $\sum_su_sS_s^pS_{s+1}$, and summing over $a$ and $b$
gives $H_u$.
\end{proof}

\begin{lemma}\label{pd-orbit-class-leading}
Let $p$ be a prime, let $K'$, $\pi$ and $R_\pi$ be as in Definition~\ref{pd-orbit-class-series}, and let
$c,c'\in R_\pi$. Then, with $G_{c,c'}$ as in Definition~\ref{pd-orbit-class-law}(1),
\[
 G_{c,c'}(\pi X,\pi Y)=\frac{\pi^{p+1}}p\big(c\,XY^p-c'\,YX^p\big)+\pi^{p+1}R(X,Y)
\]
for a power series $R$ with coefficients in $R_\pi$ and constant term $0$.
\end{lemma}
\leanref{Formalism/SharpRestriction/GrowthBounds/MultisetValuation/ExistsAtPiG1EqCPowMulIntegralSeries.lean}{Formalism.MainNotions.FormalGroupSystem.Pf.exists_atPi_G1_eq_C_pow_mul_integral_series}
\leanref{Formalism/SharpRestriction/GrowthBounds/MultisetValuation/ExistsAtPiG1EqCPowMulIntegralSeries.lean}{Formalism.MainNotions.FormalGroupSystem.Pf.th_wt_AA}

\begin{proof}
The terms of $A(X,Y)$ and of $A(Y,X)$ with $i=j=0$ are $XY^p/p$ and $YX^p/p$, and they give the first term.
Every other term has $i\leq j$ and $j\geq1$, and at $(\pi X,\pi Y)$ its coefficient is
$\pi^dp^{-(i+j+1)}$ with $d=p^{2i}+p\cdot p^{2j}$. With $g_i=\sum_{k<i}p^{2k}$, $p^{2i}=(p^2-1)g_i+1$ and
$g_i\geq i$, so $d=(p^2-1)(g_i+pg_j)+p+1$ and $\pi^dp^{-(i+j+1)}=\pi^{p+1}p^{g_i+pg_j-i-j-1}$. Here
$g_i+pg_j\geq i+2j\geq i+j+1$, so the coefficient is $\pi^{p+1}$ times an integer, and multiplying by $c$ or
$c'$ keeps it in $\pi^{p+1}R_\pi$.
\end{proof}

\begin{lemma}\label{pd-orbit-class-explicit}
Let $p\geq5$ be a prime and take a formal group system at $p$ (Definition~\ref{pd-formal-group-system}).
Let $v$ be an orbit, $j<m_*$, let $K'$, $\pi$, $R_\pi$ and $\mu_\pi$ be as in
Definition~\ref{pd-orbit-class-series}, let $c\in W(\overline{\mathbf F}_p)$ with $\sigma^2(c)=c$, and
let $s_e$, $s'_e$ and $u$ be as in Lemma~\ref{pd-orbit-class-two-sums}. Then
\[
 \exp\Big(-\sum_eG_{c,\sigma(c)}(\pi s_e,\pi s'_e)\Big)=\exp\big(-(\pi^{p+1}/p)H_u\big)\,V
\]
for a power series $V$ in the variables $Z$ with coefficients in $R_\pi$.
\end{lemma}
\leanref{Formalism/SharpRestriction/GrowthBounds/MultisetValuation/DworkExponentialRestrictions.lean}{Formalism.MainNotions.FormalGroupSystem.Pf.RestGc}
\leanref{Formalism/SharpRestriction/GrowthBounds/MultisetValuation/TwoSumsCondition.lean}{Formalism.MainNotions.FormalGroupSystem.Pf.gcCondition_of_twoSums}
\leanref{Formalism/SharpRestriction/GrowthBounds/MultisetValuation/CommonEigenvectorThird.lean}{Formalism.MainNotions.FormalGroupSystem.Pf.restGc}
\leanref{Formalism/SharpRestriction/GrowthBounds/MultisetValuation/ThRestGcVOf.lean}{Formalism.MainNotions.FormalGroupSystem.Pf.th_restGcV_of}

\begin{proof}
By Lemma~\ref{pd-orbit-class-leading} with $c'=\sigma(c)$, both in $W(\overline{\mathbf F}_p)\subset
R_\pi$, and $(X,Y)=(s_e,s'_e)$, which have coefficients in $R_\pi$ and no constant term,
\[
 \sum_eG_{c,\sigma(c)}(\pi s_e,\pi s'_e)=\frac{\pi^{p+1}}p\sum_e\big(c\,s_es_e'^{\,p}-\sigma(c)\,s'_e
 s_e^{\,p}\big)+\pi^{p+1}W,
\]
where $W=\sum_eR(s_e,s'_e)$ has coefficients in $R_\pi$ and constant term $0$. By
Lemma~\ref{pd-orbit-class-two-sums} the sum is $H_u+pY$ with $Y$ having coefficients in $R_\pi$, and $Y$
has constant term $0$ because the other terms have. So the left side of the claim is
$\exp\big(-(\pi^{p+1}/p)H_u-\pi^{p+1}(Y+W)\big)=\exp\big(-(\pi^{p+1}/p)H_u\big)\,V$ with
$V=\exp\big(-\pi^{p+1}(Y+W)\big)=\sum_{k\geq0}\frac{(-\pi^{p+1})^k}{k!}(Y+W)^k$, as both exponents have
constant term $0$. Writing $k!=p^dm$ with $p\nmid m$, $\pi^{(p+1)k}/k!=\pi^{(p+1)k-(p^2-1)d}m^{-1}$, where
$m^{-1}\in W(\overline{\mathbf F}_p)$ and $(p^2-1)d\leq(p+1)k$ because $d\leq k/(p-1)$ by Legendre's
formula. So $V$ has coefficients in $R_\pi$.
\end{proof}

\begin{lemma}\label{pd-kernel-orbit-class}
Let $p\geq5$ be a prime and take an admitted formal group system at $p$
(Definition~\ref{pd-formal-group-system}), with the notation of Definition~\ref{pd-kernel-orbit-series}.
Let $v$ be an orbit and $j<m_*$, and let $\Phi_v=\sum_e\beta(X_e,X'_e)$, the sum over the indices $e$ of
the variables at $v$. Let $K'$ be a field containing $\mathcal K_c$, $\pi\in K'$ with
$\pi^{p^2-1}=p$, $R_\pi\subset K'$ the subring generated by $\pi$ and $W(\overline{\mathbf F}_p)$, and
$R'\subset K'$ a subring containing $R_\pi$ and $\mathcal R_c$.
\begin{enumerate}
\item There are a family $u=(u_s)$ in $R_\pi$ and a power series $V$ with coefficients in $R_\pi$ such
  that the series $\exp(-\Phi_v)^Z$ in the orbit coordinates satisfies
  \[
   \exp(-\Phi_v)^Z(\pi Z)=\exp\big(-(\pi^{p+1}/p)H_u\big)\,V.
  \]
\item With this $u$, for every series $U$ in the variables $X_e,X'_e$ such that $U(\ell(X))$ has
  coefficients in $\mathcal R_c$, the series $\Theta^Z_U(\pi Z)$ lies in the class of $H_u$ over $R'$,
  where $\Theta^Z_U=(\exp(-\Phi_v)U)^Z$ is the series of the definition of a lower bound.
\end{enumerate}
\end{lemma}
\leanref{Formalism/SharpRestriction/GrowthBounds/MultisetValuation/EaRestSExpV.lean}{Formalism.MainNotions.FormalGroupSystem.Pf.ea_RestSExpV}
\leanref{Formalism/SharpRestriction/GrowthBounds/MultisetValuation/EaRestSExpVOf.lean}{Formalism.MainNotions.FormalGroupSystem.Pf.ea_restSExpV_of}
\leanref{Formalism/SharpRestriction/GrowthBounds/MultisetValuation/EntryAOneParity.lean}{Formalism.MainNotions.FormalGroupSystem.Pf.ea_RestSU}
\leanref{Formalism/SharpRestriction/GrowthBounds/MultisetValuation/EaRestSExpV.lean}{Formalism.MainNotions.FormalGroupSystem.Pf.ea_restSU_of}
\leanref{Formalism/SharpRestriction/GrowthBounds/MultisetValuation/AtPiEqRescale.lean}{Formalism.MainNotions.FormalGroupSystem.Pf.restSIntA}
\leanref{Formalism/SharpRestriction/GrowthBounds/MultisetValuation/InClassOfExp.lean}{Formalism.MainNotions.FormalGroupSystem.Pf.inClass_of_exp}

\begin{proof}
(1) By condition~(b), $\beta$ is bi-additive. Let $c\in W(\overline{\mathbf F}_p)$ with $\sigma^2(c)=c$ be
as in Lemma~\ref{pd-orbit-class-beta}, let $\mu_\pi$ be as in Definition~\ref{pd-orbit-class-series}, and
write $L$ for the family of linear forms $X_e(Z)$, $X'_e(Z)$ of the orbit coordinates at $v$, so that
$\exp(-\Phi_v)^Z(\pi Z)=\exp(-\Phi_v)(\pi L)$.

\emph{Step 1: the logarithm.} As $\ell(\pi\mu_\pi(X))=\pi\lambda_\pi(\mu_\pi(X))=\pi X$, substituting
$\pi\mu_\pi(L)$ into $f(\ell(X))$ gives $f(\pi L)$, for every power series $f$ over $K_0$ in the variables
$X_e,X'_e$ at $v$. For $f=\exp(-\Phi_v)$, $f(\ell(X))=\exp\big(-\sum_eb(X_e,X'_e)\big)$ with
$b(X_0,X_1)=\beta(\ell(X_0),\ell(X_1))$.

\emph{Step 2: the explicit series.} By Lemma~\ref{pd-orbit-class-beta},
$\exp\big(b-G_{c,\sigma(c)}\big)=1+h$ with $h$ integral and without constant term, so
$\exp\big(-(b-G_{c,\sigma(c)})\big)=(1+h)^{-1}$ is integral. Hence
\[
 \exp\Big(-\sum_eb(X_e,X'_e)\Big)=\exp\Big(-\sum_eG_{c,\sigma(c)}(X_e,X'_e)\Big)\,w,\qquad
 w=\prod_e\big(1+h(X_e,X'_e)\big)^{-1},
\]
with $w$ integral.

\emph{Step 3: substitution.} The forms in $L$ have coefficients $\sigma^s(\omega_{r'})$ and
$\sigma^s(\omega^*_{r'})$ in $W(\overline{\mathbf F}_p)$, and $\mu_\pi$ has coefficients in $\mathbf Z$ by
Lemma~\ref{pd-orbit-class-inverse}. So $\pi\mu_\pi(L)$ is a family of series with coefficients in $R_\pi$
and no constant term, and $V'=w(\pi\mu_\pi(L))$ has coefficients in $R_\pi$. With $s_e=\mu_\pi(X_e(Z))$
and $s'_e=\mu_\pi(X'_e(Z))$, Steps~1 and~2 give
\[
 \exp(-\Phi_v)^Z(\pi Z)=\exp\Big(-\sum_eG_{c,\sigma(c)}(\pi s_e,\pi s'_e)\Big)\,V'.
\]
By Lemma~\ref{pd-orbit-class-explicit}, the first factor is $\exp(-(\pi^{p+1}/p)H_u)V''$, where $u_s$ is
$c$ or $-\sigma(c)$, both in $R_\pi$, and $V''$ has coefficients in $R_\pi$. Take $V=V''V'$.

(2) The substitution $Z\mapsto\pi Z$ is a ring homomorphism, so
$\Theta^Z_U(\pi Z)=\exp(-\Phi_v)^Z(\pi Z)\cdot U^Z(\pi Z)$.

\emph{The factor $U^Z(\pi Z)$.} The series $\ell(\pi X)/\pi=\sum_j\pi^{p^{2j}-1}p^{-j}X^{p^{2j}}$ has
coefficients in $R_\pi$, because $\pi^{p^{2j}-1}p^{-j}=\pi^{p^{2j}-1-j(p^2-1)}$ and
$p^{2j}-1\geq j(p^2-1)$. Its linear term is $X$, so its compositional inverse $\mu_\pi$ also has
coefficients in $R_\pi$, and $\ell(\pi\mu_\pi(X))=\pi X$. Hence $U(\pi L)=(U\circ\ell)(\pi\mu_\pi(L))$ for
any family $L$ of series without constant term. The orbit coordinates are linear forms with coefficients
$\sigma^s(\omega_{r'})$, $\sigma^s(\omega^*_{r'})$ in $W(\overline{\mathbf F}_p)$, and $U(\ell(X))$ is
integral over $\mathcal R_c$, so $U^Z(\pi Z)$ has coefficients in $R'$.

\emph{The class.} By (1), $\Theta^Z_U(\pi Z)=\exp(-(\pi^{p+1}/p)H_u)V'$ with $V'=V\cdot U^Z(\pi Z)$
integral over $R'$, and
\[
 \exp\big(-(\pi^{p+1}/p)H_u\big)V'=\sum_{i\geq0}p^{-i}H_u^{\,i}\,\frac{(-\pi^{p+1})^i}{i!}V'.
\]
The element $\pi^{(p+1)i}/i!$ lies in $R_\pi$: writing $i!=p^{d}n$ with $p\nmid n$, it is
$\pi^{(p+1)i-(p^2-1)d}n^{-1}$, where $n^{-1}\in W(\overline{\mathbf F}_p)$ and $(p+1)i-(p^2-1)d\geq0$
since $d=v_p(i!)\leq i/(p-1)$ by Legendre's formula. As $H_u$ has no constant term, $H_u^i$ has order at
least $i$, so the coefficient of $Z^\alpha$ is the sum over $i\leq|\alpha|$, and $\Theta^Z_U(\pi Z)$ lies
in the class of $H_u$ over $R'$ with $U_i=((-\pi^{p+1})^i/i!)V'$.
\end{proof}

\subsection{Formal group bounds and the served law}\label{subsection-pd-formal-group-bounds-served-law}

\begin{definition}\label{pd-orbit-generic-data}
Let $p$ be a prime, $n'\geq2$ an even integer and $N_0,N_1\geq0$ integers. A position
$s\in\mathbf Z/n'\mathbf Z$ is \emph{even} or \emph{odd} as its representative in $\{0,\dots,n'-1\}$ is.
Consider variables $Z_{s,a,b}$ and $Y_{s,r,a}$ for positions $s$, $0\leq a,r<N_0$ and $0\leq b<N_1$. Put
$\gamma_p=p/(p^2-1)$ and $\rho_p=1/(p^2-1)$, and let $\lg_p(y)=0$ for $y\leq0$ and
$\lg_p(y)=\lfloor\log_py\rfloor+1$ for $y>0$.
\begin{enumerate}
\item A \emph{column datum} $I$ consists of a finite set $C_s$ for each position $s$ and, for each
  $c\in C_s$, an integer $h_{s,c}$ with $0\leq h_{s,c}\leq N_0$, the \emph{height} of $c$, and a map
  $I_{s,c}\colon\{0,\dots,h_{s,c}-1\}\to\{0,\dots,N_1-1\}$. Its \emph{product of minors} is
  \[
   G_I=\prod_s\prod_{c\in C_s}\det\Big(\sum_{a<N_0}Y_{s,r,a}Z_{s,a,I_{s,c}(q)}\Big)_{0\leq r,q<h_{s,c}}.
  \]
  For a power series $F=\sum_\kappa f_\kappa Z^\kappa$ and a polynomial
  $P=\sum_{\mu,\kappa}P_{\mu,\kappa}Y^\mu Z^\kappa$ over a field,
  $\langle P,F\rangle_Z=\sum_{\mu,\kappa}P_{\mu,\kappa}\,\kappa!\,f_\kappa Y^\mu$. The same formula
  defines the \emph{jet pairing} $\langle P,F\rangle$ of a polynomial and a power series in any finite
  set of variables.
\item Let $T$ be a set of positions, the \emph{top positions}, and let $A$, $\mathrm{cap}$ and
  $\mathrm{low}$ be functions from the positions to the nonnegative integers with $A(s)=0$ for
  $s\notin T$. For a position $s$ let $\Delta(s)\subset\{0,\dots,n'-1\}$ be the set of representatives
  of the $s'-s$ with $s'\in T$ of the other parity. The \emph{digit error} is
  \[
   d=\sum_s\big(N_1\lg_pA(s)+N_0N_1\lg_p\mathrm{low}(s)\big)+\sum_{s\in T}\big(\lg_pA(s)\big)^2,
  \]
  the first sum over all positions. The \emph{excess gain} is
  \[
   g^{\mathrm{ex}}=\sum_{s\in T}\ \sum_{e\geq0,\ 2e+1\notin\Delta(s)}\max\Big(\frac{A(s)}{p^{2e+1}}-
   \mathrm{cap}(s+2e+1),\ 0\Big)-d,
  \]
  and the \emph{capped gain} is
  \[
   g^{\mathrm{cap}}=\gamma_p\sum_{s\in T}A(s)-\sum_s\sum_{e\geq0}\min\Big(\frac{A(s)}{p^{2e+1}},\
   \mathrm{cap}(s+2e+1)\Big)-d,
  \]
  the middle sum over the odd $s\in T$.
\item For a column datum $I$ and a set $T$ of positions, $A(s)$ is the number of $c\in C_s$ with
  $h_{s,c}=1$ if $s\in T$, and $A(s)=0$ if $s\notin T$; $\mathrm{cap}(s)=|C_s|-A(s)$; and
  $\mathrm{low}(s)=\sum_{c\in C_s}h_{s,c}-A(s)$. The digit error and the gains of $(I,T)$ are those of
  (2) for these functions. For $i\geq1$ the \emph{edge exponent} $e_I(i)$ is the sum over the positions
  $s$ and the $c\in C_s$ of $n_{s,c}(i)$ if $s$ is odd and of $\min(h_{s,c},i)-n_{s,c}(i)$ if $s$ is
  even, where $n_{s,c}(i)$ is the number of elements of the image of $I_{s,c}$ that are less than $i$;
  thus $n_{s,c}(i)\leq\min(h_{s,c},i)$.
\item Let $K$ be a field in which $p\neq0$ and $R\subset K$ a subring. For a power series $H$ over $K$
  with constant term $0$ in a finite set of variables, a power series $F$ over $K$ in these variables
  \emph{lies in the class of $H$ over $R$} if there are power series $U_0,U_1,\dots$ with coefficients in
  $R$ such that for every exponent $\alpha$ the coefficient of the monomial of exponent $\alpha$ in $F$
  is the sum over $0\leq i\leq|\alpha|$ of $p^{-i}$ times its coefficient in $H^iU_i$. For a family
  $u=(u_s)$ in $K$ indexed by the positions, put $S_{s,a,b}=\sum_{l\geq0}(-1)^lZ_{s-2l,a,b}^{\,p^{2l}}$
  and $H_u=\sum_su_s\sum_{a,b}S_{s,a,b}^{\,p}\,S_{s+1,a,b}$.
\end{enumerate}
\end{definition}
\leanref{Formalism/GaussJetBounds/StatementP.lean}{Formalism.DworkPowerSeriesWeights.frameJet}
\leanref{Formalism/GaussJetBounds/StatementH.lean}{Formalism.DworkPowerSeriesWeights.minorProd}
\leanref{Formalism/ThetaLogCoordinates/UniformThetaSubstitution/StatementG.lean}{Formalism.InvariantFactorPaths.G7.alG}
\leanref{Formalism/ThetaLogCoordinates/UniformThetaSubstitution/StatementG.lean}{Formalism.InvariantFactorPaths.G7.capG}
\leanref{Formalism/ThetaLogCoordinates/UniformThetaSubstitution/StatementG.lean}{Formalism.InvariantFactorPaths.G7.lsG}
\leanref{Formalism/ThetaLogCoordinates/UniformThetaSubstitution/StatementE.lean}{Formalism.InvariantFactorPaths.G7.dSet}
\leanref{Formalism/ThetaLogCoordinates/UniformThetaSubstitution/StatementE.lean}{Formalism.InvariantFactorPaths.G7.oddT}
\leanref{Formalism/ThetaLogCoordinates/UniformThetaSubstitution/StatementAB.lean}{Formalism.InvariantFactorPaths.G7.lgN}
\leanref{Formalism/ThetaLogCoordinates/UniformThetaSubstitution/StatementAB.lean}{Formalism.InvariantFactorPaths.G7.deLoss}
\leanref{Formalism/GrowthSums/StatementA.lean}{Formalism.GrowthSums.gamP}
\leanref{Formalism/GrowthSums/StatementA.lean}{Formalism.GrowthSums.rhoP}
\leanref{Formalism/GrowthSums/StatementA.lean}{Formalism.GrowthSums.excessGain}
\leanref{Formalism/GrowthSums/StatementA.lean}{Formalism.GrowthSums.gainG}
\leanref{Formalism/ThetaLogCoordinates/UniformThetaSubstitution/OddLogRelation.lean}{Formalism.InvariantFactorPaths.G7.InClass}
\leanref{Formalism/ThetaLogCoordinates/UniformThetaSubstitution/CyclicPowerClassOfExpEquiv.lean}{Formalism.InvariantFactorPaths.G7.sPos}
\leanref{Formalism/ThetaLogCoordinates/UniformThetaSubstitution/CyclicPowerClassOfExpEquiv.lean}{Formalism.InvariantFactorPaths.G7.hhS}

\begin{definition}\label{pd-orbit-generic-entry}
Let $p$ be a prime, $n'\geq2$ an even integer, $N_0,N_1\geq0$ integers and $K$ a field, with the positions
and the variables $Z_{s,a,b}$ of Definition~\ref{pd-orbit-generic-data}.
\begin{enumerate}
\item Let $X$ be a finite set of variables, $Q$ a polynomial over $K$ in them, $(b_k)_{k\in\kappa}$ a
  finite family of functions from $X$ to $K$ and $\beta\in\mathbf N^\kappa$. The \emph{Hasse derivative}
  $D^\beta_bQ$ is the coefficient of $\prod_kt_k^{\beta_k}$ in $Q\big((x+\sum_kt_kb_k(x))_{x\in X}\big)$,
  where the $t_k$ are new variables; it is a polynomial in $X$. Put $L_{b_k}=\sum_{x\in X}b_k(x)\,x$ and
  $|\beta|=\sum_k\beta_k$. For a position $s$, a vector $\zeta\in K^{N_0}$ and $0\leq b<N_1$, the
  \emph{column direction} $\delta(s,\zeta,b)$ is the function on the variables $Z$ with value $\zeta_a$ at
  $Z_{s,a,b}$ and $0$ at the others.
\item \emph{Entry data} over $K$ consist of a set $T$ of positions, elements $\xi_{s,a}\in K$, integers
  $e(s,b)\geq0$ with $e(s,b)=0$ for $s\notin T$, and a finite set $\mathcal C$ of \emph{lower columns},
  each $c\in\mathcal C$ with a position $s_c$, a length $l_c\geq0$, rows $R_c(r,a)\in K$ for $r<l_c$ and
  $a<N_0$, and a map $I_c\colon\{0,\dots,l_c-1\}\to\{0,\dots,N_1-1\}$. Their \emph{lower product} is
  $Q=\prod_{c\in\mathcal C}Q_c$ with $Q_c=\det\big(\sum_aR_c(r,a)Z_{s_c,a,I_c(q)}\big)_{0\leq r,q<l_c}$,
  and their \emph{entry polynomial} is
  \[
   P=\prod_{s,b}f_{s,b}^{\,e(s,b)}\cdot Q,\qquad f_{s,b}=\sum_{a<N_0}\xi_{s,a}Z_{s,a,b}.
  \]
  Their \emph{counts} are $A(s)=\sum_be(s,b)$, the number $\mathrm{cap}(s)$ of $c\in\mathcal C$ with
  $s_c=s$, and $\mathrm{low}(s)=\sum_{s_c=s}l_c$; their digit error and gains are those of
  Definition~\ref{pd-orbit-generic-data}(2) for $T$ and these functions.
\end{enumerate}
\end{definition}
\leanref{Formalism/ThetaLogCoordinates/UniformThetaSubstitution/StatementAB.lean}{Formalism.InvariantFactorPaths.G7.hasseDir}
\leanref{Formalism/ThetaLogCoordinates/UniformThetaSubstitution/StatementAB.lean}{Formalism.InvariantFactorPaths.G7.colDir}
\leanref{Formalism/ThetaLogCoordinates/UniformThetaSubstitution/StatementAB.lean}{Formalism.InvariantFactorPaths.G7.minorAt}
\leanref{Formalism/ThetaLogCoordinates/UniformThetaSubstitution/StatementE.lean}{Formalism.InvariantFactorPaths.G7.topForm}
\leanref{Formalism/ThetaLogCoordinates/UniformThetaSubstitution/StatementE.lean}{Formalism.InvariantFactorPaths.G7.lowOp}
\leanref{Formalism/ThetaLogCoordinates/UniformThetaSubstitution/StatementE.lean}{Formalism.InvariantFactorPaths.G7.entryOp}
\leanref{Formalism/ThetaLogCoordinates/UniformThetaSubstitution/StatementE.lean}{Formalism.InvariantFactorPaths.G7.alN}
\leanref{Formalism/ThetaLogCoordinates/UniformThetaSubstitution/StatementE.lean}{Formalism.InvariantFactorPaths.G7.capN}
\leanref{Formalism/ThetaLogCoordinates/UniformThetaSubstitution/StatementE.lean}{Formalism.InvariantFactorPaths.G7.lsN}

\begin{lemma}\label{pd-orbit-generic-hasse}
Let $K$ be a field, $X$ a finite set of variables, $(b_k)_{k\in\kappa}$ a finite family of functions from
$X$ to $K$, $\beta\in\mathbf N^\kappa$ and $Q$ a polynomial over $K$ in $X$, with the Hasse derivatives
and the linear forms $L_{b_k}$ of Definition~\ref{pd-orbit-generic-entry}(1) and the jet pairing of
Definition~\ref{pd-orbit-generic-data}(1).
\begin{enumerate}
\item For every power series $G$ over $K$ in $X$,
  $\big\langle Q,\prod_kL_{b_k}^{\beta_k}\cdot G\big\rangle=\prod_k\beta_k!\cdot\langle D^\beta_bQ,G\rangle$.
\item $D^\beta_b(Q_1Q_2)=\sum_{\beta'+\beta''=\beta}D^{\beta'}_bQ_1\cdot D^{\beta''}_bQ_2$ for polynomials
  $Q_1,Q_2$. Hence, if $Q_1,\dots,Q_n$ are polynomials, $S\subset\kappa$, and $m_1,\dots,m_n$ are integers
  such that $D^{\beta'}_bQ_c=0$ whenever $\sum_{k\in S}\beta'_k>m_c$, then $D^\beta_b(Q_1\cdots Q_n)=0$
  whenever $\sum_{k\in S}\beta_k>\sum_cm_c$.
\item If $Q$ is homogeneous of degree $D$, then $D^\beta_bQ$ is homogeneous of degree $D-|\beta|$, and it
  is $0$ if $|\beta|>D$. For every monomial $x^n$ of $D^\beta_bQ$ there is a monomial $x^{n_0}$ of $Q$
  with $n\leq n_0$ and $|n|+|\beta|=|n_0|$. If $K$ has an ultrametric absolute value and the coefficients
  of $Q$ and the values of the $b_k$ have absolute value at most $1$, then so do the coefficients of
  $D^\beta_bQ$.
\item Let $W$ be a second finite set of variables, $(A_{w,x})$ a matrix over $K$ indexed by $W\times X$,
  $P$ a polynomial over $K$ in $W$ and $F$ a power series over $K$ in $X$. Let $F^A$ be $F$ with
  $\sum_wA_{w,x}w$ substituted for each $x$, and $P^A$ be $P$ with $\sum_xA_{w,x}x$ substituted for each
  $w$. Then $\langle P,F^A\rangle=\langle P^A,F\rangle$, the jet pairings in $W$ and in $X$.
\end{enumerate}
\end{lemma}
\leanref{Formalism/ThetaLogCoordinates/UniformThetaSubstitution/FrobeniusLiftNorms/FinsuppStepInduction.lean}{Formalism.InvariantFactorPaths.G7.e_one_identity_of}
\leanref{Formalism/ThetaLogCoordinates/UniformThetaSubstitution/FrobeniusLiftNorms/FinsuppStepInduction.lean}{Formalism.InvariantFactorPaths.G7.hasseDir_mul}
\leanref{Formalism/ThetaLogCoordinates/UniformThetaSubstitution/FrobeniusLiftNorms/FinsuppStepInduction.lean}{Formalism.InvariantFactorPaths.G7.e_one_vanish_of}
\leanref{Formalism/ThetaLogCoordinates/UniformThetaSubstitution/FrobeniusLiftNorms/FinsuppStepInduction.lean}{Formalism.InvariantFactorPaths.G7.e_one_integral_of}
\leanref{Formalism/ThetaLogCoordinates/UniformThetaSubstitution/FrobeniusLiftNorms/MinorProdBlockDegree.lean}{Formalism.InvariantFactorPaths.G7.hasseDerivSupport}
\leanref{Formalism/GaussJetBounds/GaussJetProofs/ConstantCoeffAevalLinP.lean}{Formalism.DworkPowerSeriesWeights.jet_pair_linearity}

\begin{proof}
For a variable $x$, a polynomial $P$ and a power series $G$, $\langle P,xG\rangle=\langle\partial
P/\partial x,G\rangle$, and $\langle Px,G\rangle=\langle P,\partial G/\partial x\rangle$ when $G$ is a
polynomial, because the monomial of exponent $\mu+\epsilon_x$ has factorial $(\mu_x+1)\mu!$; here
$\epsilon_x$ is the unit exponent at $x$.

(1) By the first identity, $\langle Q,L_{b_k}G\rangle=\langle\partial_kQ,G\rangle$ with
$\partial_k=\sum_xb_k(x)\partial/\partial x$. By the chain rule, the substitution
$x\mapsto x+\sum_jt_jb_j(x)$ carries $\partial_kQ$ to the derivative in $t_k$ of the image of $Q$, so
$D^\beta_b(\partial_kQ)=(\beta_k+1)D^{\beta+\epsilon_k}_bQ$. Induction on $|\beta|$, from $D^0_bQ=Q$,
gives (1).

(2) The substitution $x\mapsto x+\sum_kt_kb_k(x)$ is a ring homomorphism; compare the coefficients of
$\prod_kt_k^{\beta_k}$. For the consequence argue by induction on $n$: in the expansion of
$D^\beta_b(Q_1\cdot Q_2\cdots Q_n)$ by the first part, each term has $\sum_{k\in S}\beta'_k>m_1$ or
$\sum_{k\in S}\beta''_k>m_2+\dots+m_n$, so it vanishes.

(3) The image of a monomial $x^{n_0}$ is $\prod_x(x+\sum_kt_kb_k(x))^{n_0(x)}$. Its coefficient of
$\prod_kt_k^{\beta_k}$ is a combination of monomials $x^n$ with $n\leq n_0$ and $|n|=|n_0|-|\beta|$,
whose coefficients are sums of integers times products of values of the $b_k$. This gives the first two
assertions, and the ultrametric inequality gives the last.

(4) Both sides are linear in $F$ and involve finitely many of its coefficients, so let $F$ be a
polynomial. For $P$ constant both sides are $P\cdot F(0)$. If the identity holds for $P$ and every
polynomial $F$, then for $w\in W$, by the second identity and the chain rule,
\[
 \langle Pw,F^A\rangle=\Big\langle P,\frac{\partial F^A}{\partial w}\Big\rangle=\sum_xA_{w,x}\Big\langle
 P,\Big(\frac{\partial F}{\partial x}\Big)^A\Big\rangle=\sum_xA_{w,x}\Big\langle P^A,\frac{\partial F}
 {\partial x}\Big\rangle=\Big\langle P^A\sum_xA_{w,x}x,F\Big\rangle,
\]
which is $\langle(Pw)^A,F\rangle$. Induction on $P$ gives (4).
\end{proof}

\begin{lemma}\label{pd-orbit-generic-minors}
Let $p$ be a prime, $n'\geq2$ an even integer, $N_0,N_1\geq0$ integers and $K$ a field, with the
notation of Definitions~\ref{pd-orbit-generic-data} and~\ref{pd-orbit-generic-entry}. Let
$(b_k)_{k\in\kappa}$ be a finite family of column directions $b_k=\delta(s_k,\zeta_k,j_k)$ and
$\beta\in\mathbf N^\kappa$.
\begin{enumerate}
\item Let $s$ be a position, $l\geq0$, $R(r,a)\in K$ for $r<l$ and $a<N_0$,
  $I\colon\{0,\dots,l-1\}\to\{0,\dots,N_1-1\}$ a map, and $Q_0=\det\big(\sum_aR(r,a)Z_{s,a,I(q)}\big)_{r,q<l}$.
  Then $D^\beta_bQ_0=0$ in each of the following cases: $\beta_k\geq2$ for some $k$; there are $k\neq k'$
  with $\beta_k,\beta_{k'}\geq1$, $s_k=s_{k'}$, and $\zeta_k=\zeta_{k'}$ or $j_k=j_{k'}$; or
  $\beta_k\geq1$ for some $k$ with $s_k\neq s$ or $j_k$ not in the image of $I$.
\item Let $Q$ be the lower product of entry data over $K$ with lower columns $\mathcal C$ and counts
  $\mathrm{cap}$, let $S\subset\kappa$ and let $s_0$ be a position with $s_k=s_0$ for every $k\in S$.
  If the $\zeta_k$, $k\in S$, are all equal and $\sum_{k\in S}\beta_k>\mathrm{cap}(s_0)$, then
  $D^\beta_bQ=0$. If $J\subset\{0,\dots,N_1-1\}$, $j_k\in J$ for every $k\in S$, and $\sum_{k\in
  S}\beta_k$ is larger than the sum over the $c\in\mathcal C$ with $s_c=s_0$ of the number of elements of
  $J$ in the image of $I_c$, then $D^\beta_bQ=0$.
\end{enumerate}
\end{lemma}
\leanref{Formalism/ThetaLogCoordinates/UniformThetaSubstitution/FrobeniusLiftNorms/FinsuppStepInduction.lean}{Formalism.InvariantFactorPaths.G7.minor_hits}
\leanref{Formalism/ThetaLogCoordinates/UniformThetaSubstitution/FrobeniusLiftNorms/MinorProdBlockDegree.lean}{Formalism.InvariantFactorPaths.G7.HitsCapPos_of}
\leanref{Formalism/ThetaLogCoordinates/UniformThetaSubstitution/FrobeniusLiftNorms/MinorProdBlockDegree.lean}{Formalism.InvariantFactorPaths.G7.HitsColumnsPos_of}

\begin{proof}
(1) If $I$ is not injective, $Q_0$ has two equal columns and vanishes. Let $I$ be injective. After the
substitution $x\mapsto x+\sum_kt_kb_k(x)$, the entry $(r,q)$ of the matrix of $Q_0$ is
$\sum_aR(r,a)Z_{s,a,I(q)}$ plus the sum of $t_k\sum_aR(r,a)\zeta_{k,a}$ over the $k$ with $s_k=s$ and
$j_k=I(q)$.

\emph{Columns.} Let $S\subset\kappa$ and $J\subset\{0,\dots,N_1-1\}$ with $j_k\in J$ for the $k\in S$ with
$s_k=s$. The column $q$ involves the $t_k$, $k\in S$, only if $I(q)\in J$, and then linearly. The
determinant is linear in each column, so its degree in the $t_k$, $k\in S$, is at most the number of
elements of $J$ in the image of $I$, and $D^\beta_bQ_0=0$ if $\sum_{k\in S}\beta_k$ is larger.

\emph{Rows.} Let $S\subset\kappa$ with $\zeta_k=\zeta$ for every $k\in S$. The terms in the $t_k$, $k\in
S$, form the matrix $vw^{\mathsf T}$ with $v_r=\sum_aR(r,a)\zeta_a$ and $w_q=\sum t_k$, the sum over the
$k\in S$ with $s_k=s$ and $j_k=I(q)$. Expanding the determinant by columns, a term that takes two columns
from $vw^{\mathsf T}$ vanishes, since these columns are proportional. So the degree in the $t_k$, $k\in
S$, is at most $1$, and $D^\beta_bQ_0=0$ if $\sum_{k\in S}\beta_k\geq2$.

The cases of (1) follow: $\beta_k\geq2$ by rows with $S=\{k\}$; $\zeta_k=\zeta_{k'}$ by rows with
$S=\{k,k'\}$; $j_k=j_{k'}$ by columns with $S=\{k,k'\}$ and $J=\{j_k\}$; $s_k\neq s$ by columns with
$S=\{k\}$ and $J=\emptyset$; and $j_k$ outside the image of $I$ by columns with $S=\{k\}$ and
$J=\{j_k\}$.

(2) Apply Lemma~\ref{pd-orbit-generic-hasse}(2) to the factors $Q_c$ of $Q$. For the first assertion take
$m_c=1$ if $s_c=s_0$ and $m_c=0$ otherwise. If $s_c=s_0$ and $\sum_{k\in S}\beta'_k\geq2$, then some
$\beta'_k\geq2$, or two distinct $k,k'\in S$ have $\beta'_k,\beta'_{k'}\geq1$, $s_k=s_{k'}$ and
$\zeta_k=\zeta_{k'}$; so $D^{\beta'}_bQ_c=0$ by (1). If $s_c\neq s_0$ and $\sum_{k\in S}\beta'_k\geq1$,
some $k\in S$ with $\beta'_k\geq1$ has $s_k=s_0\neq s_c$, and again $D^{\beta'}_bQ_c=0$. For the second
assertion take $m_c$ the number of elements of $J$ in the image of $I_c$ if $s_c=s_0$, and $m_c=0$
otherwise. If $s_c=s_0$ and $D^{\beta'}_bQ_c\neq0$, then by (1) $\beta'_k\leq1$ for $k\in S$, and the
$k\in S$ with $\beta'_k=1$ have distinct columns $j_k$, which lie in $J$ and in the image of $I_c$; so
$\sum_{k\in S}\beta'_k\leq m_c$. The case $s_c\neq s_0$ is as before.
\end{proof}

\begin{lemma}\label{pd-orbit-generic-decomposition}
Let $p>2$ be a prime and $K$ a field with an ultrametric absolute value such that $|p|=p^{-1}$, with unit
ball $R$, and take $n'$, $N_0$, $N_1$, the variables $Z_{s,a,b}$ and the series $H_u$ of
Definition~\ref{pd-orbit-generic-data}. Let $u=(u_s)$ be a family in $R$, $T$ a set of positions and
$\xi_{s,a}\in R$ for positions $s$ and $a<N_0$. Let $W_{s,b}$, for positions $s$ and $b<N_1$, be new
variables, and for a power series $G$ in the $Z$ let $G'$ be the series in the $W$ and the $Z$ obtained by
substituting $Z_{s,a,b}+\xi_{s,a}W_{s,b}$ for $Z_{s,a,b}$ if $s\in T$, and keeping $Z_{s,a,b}$ if
$s\notin T$.
\begin{enumerate}
\item If $H$ is a power series in the $Z$ with constant term $0$ and $G$ lies in the class of $H$ over
  $R$, then $G'$ lies in the class of $H'$ over $R$.
\item Let $E\geq0$ and $D<p^{2E}$ be integers. There are power series $Z^\sharp$ and $Y^\sharp$ over $R$
  in the $W$ and the $Z$, every monomial of $Z^\sharp$ of degree at least $p$ in the variables $Z$, such
  that $H'_u$ and
  \[
   \sum_\phi A_\phi W^{d_\phi}+\sum_kB_kW^{d_k}L_k+Z^\sharp+pY^\sharp
  \]
  have the same coefficients in degrees at most $D$. Here $\phi$ runs over the tuples $(s,l,l',a,b)$ and
  $k$ over the triples $(s,l,b)$ with $s$ a position, $0\leq l,l'<E$, $a<N_0$ and $b<N_1$;
  $W^{d_\phi}=W_{s-2l,b}^{\,p^{1+2l}}W_{s+1-2l',b}^{\,p^{2l'}}$, and
  $A_\phi=(-1)^{l+l'}u_s\,\xi_{s-2l,a}^{\,p^{1+2l}}\,\xi_{s+1-2l',a}^{\,p^{2l'}}$ if $s-2l$ and $s+1-2l'$
  lie in $T$ and $A_\phi=0$ otherwise; $W^{d_k}=W_{s,b}^{\,p^{2l+1}}$,
  $L_k=\sum_a\xi_{s,a}^{\,p^{2l+1}}Z_{s+2l+1,a,b}$, and $B_k=(-1)^lu_{s+2l}$ if $s\in T$ and $B_k=0$
  otherwise.
\end{enumerate}
\end{lemma}
\leanref{Formalism/ThetaLogCoordinates/UniformThetaSubstitution/FrobeniusLiftNorms/TopsChoicePartB.lean}{Formalism.InvariantFactorPaths.G7.class_substitution}
\leanref{Formalism/ThetaLogCoordinates/UniformThetaSubstitution/FrobeniusLiftNorms/TopsChoicePartB.lean}{Formalism.InvariantFactorPaths.G7.inClass_iff}
\leanref{Formalism/ThetaLogCoordinates/UniformThetaSubstitution/FrobeniusLiftNorms/EDecomp.lean}{Formalism.InvariantFactorPaths.G7.decomposition_of_hh}
\leanref{Formalism/ThetaLogCoordinates/UniformThetaSubstitution/StatementE.lean}{Formalism.InvariantFactorPaths.G7.zSub}
\leanref{Formalism/ThetaLogCoordinates/UniformThetaSubstitution/StatementE.lean}{Formalism.InvariantFactorPaths.G7.aAof}
\leanref{Formalism/ThetaLogCoordinates/UniformThetaSubstitution/StatementE.lean}{Formalism.InvariantFactorPaths.G7.dAof}
\leanref{Formalism/ThetaLogCoordinates/UniformThetaSubstitution/StatementE.lean}{Formalism.InvariantFactorPaths.G7.aBof}
\leanref{Formalism/ThetaLogCoordinates/UniformThetaSubstitution/StatementE.lean}{Formalism.InvariantFactorPaths.G7.dBof}
\leanref{Formalism/ThetaLogCoordinates/UniformThetaSubstitution/StatementE.lean}{Formalism.InvariantFactorPaths.G7.bBof}

\begin{proof}
(1) Since $H^i$ has no monomial of degree less than $i$, $G$ lies in the class of $H$ over $R$ if and only
if there are series $U_i$ over $R$ such that, for every $M$, $G-\sum_{i<M}p^{-i}H^iU_i$ has no monomial of
degree less than $M$. The map $G\mapsto G'$ is a ring homomorphism, carries series over $R$ to series over
$R$ since $\xi_{s,a}\in R$, and does not lower degrees, since it substitutes linear forms for the
variables. Apply it to this condition, with the series $U'_i$.

(2) \emph{Truncation.} Let $S^E_{s,a,b}=\sum_{l<E}(-1)^lZ_{s-2l,a,b}^{\,p^{2l}}$ and let $H^E_u$ be
$H_u$ formed with the $S^E$. The differences $S_{s,a,b}-S^E_{s,a,b}$ have only monomials of degree at
least $p^{2E}>D$, hence so do $H_u-H^E_u$ and $H'_u-(H^E_u)'$. So it suffices to prove (2) for
$(H^E_u)'$.

\emph{Congruences.} Write $z_{s,a,b}=Z_{s,a,b}$, and $t_{s,a,b}=\xi_{s,a}W_{s,b}$ if $s\in T$ and
$t_{s,a,b}=0$ otherwise. For series $f,g$ over $R$, $(f+g)^p-f^p-g^p$ is $p$ times a series over $R$; by
induction the same holds for $(f+g)^{p^m}-f^{p^m}-g^{p^m}$ and for $(\sum_lf_l)^p-\sum_lf_l^p$. Since $p$
is odd, $((-1)^l)^p=(-1)^l$. Hence, up to $p$ times series over $R$,
\[
 \big((S^E_{s,a,b})'\big)^p\equiv\sum_{l<E}(-1)^l\big(z_{s-2l,a,b}^{\,p^{1+2l}}+t_{s-2l,a,b}^{\,p^{1+2l}}
 \big)=Z_1+T_1,\qquad (S^E_{s+1,a,b})'\equiv\sum_{l<E}(-1)^l\big(z_{s+1-2l,a,b}^{\,p^{2l}}+
 t_{s+1-2l,a,b}^{\,p^{2l}}\big)=Z_2+T_2,
\]
where $Z_1,Z_2$ collect the terms in $z$ and $T_1,T_2$ those in $t$. As $u_s\in R$,
$(H^E_u)'=\sum_su_s\sum_{a,b}(Z_1+T_1)(Z_2+T_2)+pY^\sharp$ with $Y^\sharp$ over $R$.

\emph{Sorting.} Write $(Z_1+T_1)(Z_2+T_2)=T_1T_2+T_1Z_{s+1,a,b}+\big(Z_1(Z_2+T_2)+T_1(Z_2-Z_{s+1,a,b})
\big)$. The series $Z_1$ and $Z_2-Z_{s+1,a,b}$ have only monomials of degree at least $p$ in the $Z$, so
the last terms, summed with the factors $u_s$, give $Z^\sharp$. Expanding $T_1T_2$ gives the terms
$A_\phi W^{d_\phi}$; a pair $(l,l')$ with $s-2l\notin T$ or $s+1-2l'\notin T$ contributes $0$. The terms
of $u_sT_1Z_{s+1,a,b}$ are $(-1)^lu_s\,\xi_{s-2l,a}^{\,p^{1+2l}}W_{s-2l,b}^{\,p^{1+2l}}Z_{s+1,a,b}$ with
$s-2l\in T$; replacing $s$ by $s+2l$ and summing over $a$ gives $\sum_kB_kW^{d_k}L_k$.
\end{proof}

\begin{lemma}\label{pd-orbit-generic-term}
Let $p$ be a prime and $K$ a field with an ultrametric absolute value such that $|p|=p^{-1}$, with unit
ball $R$. Let $W$ and $Z$ be finite sets of variables, and $h$, $g$ power series over $K$ in $W\cup Z$,
$h$ with constant term $0$ and $g$ in the class of $h$ over $R$
(Definition~\ref{pd-orbit-generic-data}(4)). Let $(A_\phi)_{\phi\in\Phi}$ and $(B_k)_{k\in\kappa}$ be
finite families in $R$ with exponents $d_\phi,d_k\in\mathbf N^W$, let $(b_k)_{k\in\kappa}$ be functions
from $Z$ to $R$ with linear forms $L_{b_k}=\sum_zb_k(z)z$, let $Z^\sharp$ and $Y^\sharp$ be series over
$R$, every monomial of $Z^\sharp$ of degree at least $p$ in the variables $Z$, and let $D\geq0$ be such
that $h$ and
\[
 \sum_\phi A_\phi W^{d_\phi}+\sum_kB_kW^{d_k}L_{b_k}+Z^\sharp+pY^\sharp
\]
have the same coefficients in degrees at most $D$. Let $Q$ be a polynomial in $Z$ with coefficients in
$R$, homogeneous of degree $D_Q$, and $e\in\mathbf N^W$ with $|e|+D_Q=D$. Then
$\langle W^eQ,g\rangle=0$, or there are $\alpha\in\mathbf N^\Phi$, $\beta\in\mathbf N^\kappa$, an integer
$i_2\geq0$ and a monomial $Z^n$ of $D^\beta_bQ$ (Definition~\ref{pd-orbit-generic-entry}(1)) such that
$A_\phi\neq0$ if $\alpha_\phi\neq0$, $B_k\neq0$ if $\beta_k\neq0$,
$\sum_\phi\alpha_\phi d_\phi+\sum_k\beta_kd_k\leq e$ in each coordinate, $p\,i_2\leq|n|$, and
\[
 |\langle W^eQ,g\rangle|\leq p^{\,\sum_\phi\alpha_\phi+|\beta|+i_2}\,|e!|\,\big|\,|\beta|!\,\big|\,|n!|.
\]
\end{lemma}
\leanref{Formalism/ThetaLogCoordinates/UniformThetaSubstitution/FrobeniusLiftNorms/NormNatCastLeOfDvd.lean}{Formalism.InvariantFactorPaths.G7.TermBound_of}
\leanref{Formalism/ThetaLogCoordinates/UniformThetaSubstitution/FrobeniusLiftNorms/NormNatCastLeOfDvd.lean}{Formalism.InvariantFactorPaths.G7.ETerm.term_bound}

\begin{proof}
\emph{Expansion.} Let $\Sigma$ be the displayed sum and $U_i$ the series of the class. The polynomial
$W^eQ$ is homogeneous of degree $D$, so $\langle W^eQ,g\rangle$ involves only coefficients of $g$ of
degree $D$. Each of these is the sum over $i\leq D$ of $p^{-i}$ times the coefficient of $h^iU_i$, and in
degrees at most $D$ the series $h^iU_i$ and $\Sigma^iU_i$ agree. Expanding $\Sigma^i$ by the multinomial
theorem over its terms, with exponents $\alpha_\phi$, $\beta_k$, $i_2$ for $Z^\sharp$ and $i_3$ for
$pY^\sharp$, of total $i$, writes $\langle W^eQ,g\rangle$ as a finite sum of terms
\[
 p^{-i}\binom{i}{\alpha,\beta,i_2,i_3}\Big\langle W^eQ,\prod_kL_{b_k}^{\beta_k}\cdot c\,W^{d}V\Big\rangle,
\]
with $c=\prod_\phi A_\phi^{\alpha_\phi}\prod_kB_k^{\beta_k}p^{i_3}$,
$d=\sum_\phi\alpha_\phi d_\phi+\sum_k\beta_kd_k$ and $V=(Z^\sharp)^{i_2}(Y^\sharp)^{i_3}U_i$. By the
ultrametric inequality, if $\langle W^eQ,g\rangle\neq0$, its absolute value is at most that of a nonzero
term, and it suffices to bound a nonzero term.

\emph{One term.} Extend the $b_k$ by $0$ at the variables $W$. Then $D^\beta_b(W^eQ)=W^eD^\beta_bQ$, and
by Lemma~\ref{pd-orbit-generic-hasse}(1), $\langle W^eQ,\prod_kL_{b_k}^{\beta_k}\cdot cW^dV\rangle=
\prod_k\beta_k!\cdot\langle W^eD^\beta_bQ,cW^dV\rangle$. Write $D^\beta_bQ=\sum_nq_nZ^n$. Then
$\langle W^eD^\beta_bQ,cW^dV\rangle=\sum_nq_n\,e!\,n!\,c\,v_n$, where $v_n$ is the coefficient of
$W^{e-d}Z^n$ in $V$ if $d\leq e$, and $v_n=0$ otherwise. If the term is nonzero, the ultrametric
inequality gives an $n$ whose summand is nonzero and has absolute value at least that of the sum. Then
$c\neq0$, so $A_\phi\neq0$ when $\alpha_\phi\neq0$ and $B_k\neq0$ when $\beta_k\neq0$; $d\leq e$; and
$v_n\neq0$, so $p\,i_2\leq|n|$, since every monomial of $(Z^\sharp)^{i_2}$ has degree at least $p\,i_2$
in the $Z$. The absolute value of the term is at most that of this summand times
$|p^{-i}\binom{i}{\alpha,\beta,i_2,i_3}\prod_k\beta_k!|$. Now $|q_n|\leq1$ by
Lemma~\ref{pd-orbit-generic-hasse}(3), $|cp^{-i_3}|\leq1$ and $|v_n|\leq1$; $|p^{-i}p^{i_3}|=
p^{\sum\alpha_\phi+|\beta|+i_2}$; and $\binom{i}{\alpha,\beta,i_2,i_3}\prod_k\beta_k!$ is $|\beta|!$
times the multinomial coefficient $\binom{i}{\alpha,|\beta|,i_2,i_3}$, an integer, so its absolute value
is at most $|\,|\beta|!\,|$. This gives the bound.
\end{proof}

\begin{lemma}\label{pd-orbit-generic-factorials}
Let $p$ be a prime, $S$ a finite set, $T\subset S$, $N_0,N_1\geq0$ integers and $E$ a finite set of
nonnegative integers, with $\gamma_p$, $\rho_p$ and $\lg_p$ as in Definition~\ref{pd-orbit-generic-data}.
For $s\in S$ let $A(s)$, $\mathrm{low}(s)$, $e(s,b)$ for $b<N_1$, $n(s,x)$ for
$x\in\{0,\dots,N_0-1\}\times\{0,\dots,N_1-1\}$, and $i(s,l)$ for $l\in E$ be nonnegative integers, and let
$i_0,i_2\geq0$ be integers, such that: $A(s)=0$ and $i(s,l)=0$ for $s\notin T$; $\sum_be(s,b)=A(s)$;
$\sum_xn(s,x)\leq\mathrm{low}(s)$; $p^{2l+1}i(s,l)\leq A(s)$; $\sum_{s,x}n(s,x)+\sum_{s,l}i(s,l)=
\sum_s\mathrm{low}(s)$; and $p\,i_2\leq\sum_s\mathrm{low}(s)$. Let
$d=\sum_{s\in S}(N_1\lg_pA(s)+N_0N_1\lg_p\mathrm{low}(s))+\sum_{s\in T}(\lg_pA(s))^2$. Then
\begin{multline*}
 \gamma_p\sum_{s\in T}A(s)-i_0-\sum_{s,l}i(s,l)-d\leq\sum_{s,b}v_p(e(s,b)!)+\sum_{s,l}v_p(i(s,l)!)+
 \sum_{s,x}v_p(n(s,x)!)\\-\rho_p\sum_s\big(A(s)+\mathrm{low}(s)\big)-\Big(i_0+\sum_{s,l}i(s,l)+i_2\Big).
\end{multline*}
\end{lemma}
\leanref{Formalism/ThetaLogCoordinates/UniformThetaSubstitution/FrobeniusLiftNorms/ABArith.lean}{Formalism.InvariantFactorPaths.G7.gainFactorialInequality}
\leanref{Formalism/ThetaLogCoordinates/UniformThetaSubstitution/FrobeniusLiftNorms/ABArith.lean}{Formalism.InvariantFactorPaths.G7.factorialValuationLowerBound}
\leanref{Formalism/GrowthSums/FactorialBaseGrowth/QuadraticExponentIdentity.lean}{Formalism.GrowthSums.oddPowerInvSumIdentity}

\begin{proof}
By Legendre's formula $v_p(m!)=(m-\sigma_p(m))/(p-1)$, where the digit sum $\sigma_p(m)$ is at most
$(p-1)\lg_pm$; so $v_p(m!)\geq m/(p-1)-\lg_pm$, and $\lg_p$ is non-decreasing. As $e(s,b)\leq A(s)$,
$\sum_bv_p(e(s,b)!)\geq A(s)/(p-1)-N_1\lg_pA(s)$; as $n(s,x)\leq\mathrm{low}(s)$,
$\sum_xv_p(n(s,x)!)\geq\sum_xn(s,x)/(p-1)-N_0N_1\lg_p\mathrm{low}(s)$. If $i(s,l)\neq0$, then
$p^{2l+1}\leq A(s)$, so $l<\lg_pA(s)$, and $\lg_pi(s,l)\leq\lg_pA(s)$; hence
$\sum_lv_p(i(s,l)!)\geq\sum_li(s,l)/(p-1)-(\lg_pA(s))^2$, and the terms with $s\notin T$ vanish. Adding,
and using $\sum_{s,x}n(s,x)+\sum_{s,l}i(s,l)=\sum_s\mathrm{low}(s)$ and $1/(p-1)=\gamma_p+\rho_p$, the
right side is at least $\gamma_p\sum_sA(s)+\gamma_p\sum_s\mathrm{low}(s)-d-i_0-\sum_{s,l}i(s,l)-i_2$.
Finally $i_2\leq\sum_s\mathrm{low}(s)/p\leq\gamma_p\sum_s\mathrm{low}(s)$.
\end{proof}

\begin{lemma}\label{pd-orbit-generic-star}
Let $p>2$ be a prime and $K$ a field with an ultrametric absolute value such that $|p|=p^{-1}$, with unit
ball $R$. Take entry data over $K$ (Definition~\ref{pd-orbit-generic-entry}) with $\xi_{s,a}\in R$ and
$R_c(r,a)\in R$, with entry polynomial $P$, lower product $Q$, counts $A$, $\mathrm{low}$ and digit error
$d$. Let $u=(u_s)$ be a family in $R$ and $F$ a power series over $K$ in the $Z$ that lies in the class
of $H_u$ over $R$ (Definition~\ref{pd-orbit-generic-data}(4)). For $E\geq1$ let $\phi$, $k$, $d_\phi$ and
$d_k$ be as in Lemma~\ref{pd-orbit-generic-decomposition}(2), and let
$b_k=\delta(s+2l+1,\xi_s^{\,p^{2l+1}},b)$ for $k=(s,l,b)$, where
$\xi_s^{\,p^{2l+1}}=(\xi_{s,a}^{\,p^{2l+1}})_a$. Then $\langle P,F\rangle_Z=0$, or there are $E\geq1$
and families $\alpha=(\alpha_\phi)$ and $\beta=(\beta_k)$ of nonnegative integers such that:
\begin{enumerate}
\item[(a)] if $\alpha_\phi\neq0$ for $\phi=(s,l,l',a,b)$, then $s-2l$ and $s+1-2l'$ lie in $T$; if
  $\beta_k\neq0$ for $k=(s,l,b)$, then $s\in T$;
\item[(b)] $\sum_\phi\alpha_\phi d_\phi+\sum_k\beta_kd_k\leq e$ in each coordinate, where $e$ is the
  exponent with $e(s,b)$ at $W_{s,b}$;
\item[(c)] $D^\beta_bQ\neq0$;
\item[(d)] $|\langle P,F\rangle_Z|\leq p^{-v-\rho_p\sum_s(A(s)+\mathrm{low}(s))}$ with
  $v=\gamma_p\sum_{s\in T}A(s)-\sum_\phi\alpha_\phi-|\beta|-d$.
\end{enumerate}
\end{lemma}
\leanref{Formalism/ThetaLogCoordinates/UniformThetaSubstitution/FrobeniusLiftNorms/NormLeOfDvd.lean}{Formalism.InvariantFactorPaths.G7.EntryStar_of}
\leanref{Formalism/ThetaLogCoordinates/UniformThetaSubstitution/FrobeniusLiftNorms/NormLeOfDvd.lean}{Formalism.InvariantFactorPaths.G7.StarEntryOperators.star_norm}
\leanref{Formalism/ThetaLogCoordinates/UniformThetaSubstitution/FrobeniusLiftNorms/MinorProdBlockDegree.lean}{Formalism.InvariantFactorPaths.G7.minorProdBlockDegree}

\begin{proof}
Let $D=\sum_{s,b}e(s,b)+\sum_s\mathrm{low}(s)$ and $E=D+1$, so that $D<p^{2E}$; use the variables $W$ and
the substitution $G\mapsto G'$ of Lemma~\ref{pd-orbit-generic-decomposition}.

\emph{Step 1: the substitution.} Apply Lemma~\ref{pd-orbit-generic-hasse}(4) to the matrix of
$G\mapsto G'$, with rows indexed by the $W$ and the $Z$: substituting $f_{s,b}$ for $W_{s,b}$ with
$s\in T$, $0$ for $W_{s,b}$ with $s\notin T$, and $Z$ for $Z$ carries $W^eQ$ to $P$, since $e(s,b)=0$ for
$s\notin T$. Hence $\langle W^eQ,F'\rangle=\langle P,F\rangle_Z$.

\emph{Step 2: the term bound.} By Lemma~\ref{pd-orbit-generic-decomposition}, $F'$ lies in the class of
$H'_u$ over $R$, and $H'_u$ agrees in degrees at most $D$ with the decomposition there, whose $A_\phi$
and $B_k$ lie in $R$ and whose $L_k$ is $L_{b_k}$. Each $Q_c$ is homogeneous of degree $l_c$, in the
variables at $s_c$, with coefficients in $R$, so $Q$ is homogeneous of degree $\sum_s\mathrm{low}(s)$
with coefficients in $R$, and every monomial of $Q$ has degree $\mathrm{low}(s)$ in the variables at each
position $s$. Lemma~\ref{pd-orbit-generic-term} applies with $h=H'_u$ and $g=F'$. If
$\langle P,F\rangle_Z\neq0$, it gives $\alpha$, $\beta$, $i_2$ and a monomial $Z^n$ of $D^\beta_bQ$;
then (a) holds because $A_\phi=0$ and $B_k=0$ in the other cases, (b) holds, and (c) holds.

\emph{Step 3: the factorials.} By Lemma~\ref{pd-orbit-generic-hasse}(3) there is a monomial $Z^{n_0}$ of
$Q$ with $n\leq n_0$ and $|n|+|\beta|=|n_0|=\sum_s\mathrm{low}(s)$. Put $n(s,x)=n_{s,a,b}$ for
$x=(a,b)$, $i(s,l)=\sum_b\beta_{(s,l,b)}$ and $i_0=\sum_\phi\alpha_\phi$. Then
$\sum_xn(s,x)\leq\mathrm{low}(s)$; $\sum_{s,x}n(s,x)+\sum_{s,l}i(s,l)=\sum_s\mathrm{low}(s)$;
$p^{2l+1}i(s,l)\leq\sum_be(s,b)=A(s)$ by (b), since $d_k$ is $p^{2l+1}$ at $W_{s,b}$ for $k=(s,l,b)$;
$i(s,l)=0$ for $s\notin T$ by (a); and $p\,i_2\leq|n|\leq\sum_s\mathrm{low}(s)$. So
Lemma~\ref{pd-orbit-generic-factorials} applies. For a nonzero integer $m$, $|m|=p^{-v_p(m)}$, since
$|p|=p^{-1}$ and integers prime to $p$ have absolute value $1$. So $|e!|=p^{-\sum v_p(e(s,b)!)}$,
$|n!|=p^{-\sum v_p(n(s,x)!)}$, and $|\,|\beta|!\,|\leq p^{-\sum v_p(i(s,l)!)}$, because
$\prod_{s,l}i(s,l)!$ divides $(\sum_{s,l}i(s,l))!=|\beta|!$. Inserting these into the bound of
Lemma~\ref{pd-orbit-generic-term} and using Lemma~\ref{pd-orbit-generic-factorials} gives (d).
\end{proof}

\begin{lemma}\label{pd-orbit-generic-gain-sums}
Let $p>1$ be a real number, $n'\geq2$ an even integer and $T$ a set of positions in
$\mathbf Z/n'\mathbf Z$. Let $A$ and $\mathrm{cap}$ be functions from the positions to the nonnegative real
numbers and $d$ a real number, and define $\Delta$, $\gamma_p$, $g^{\mathrm{ex}}$ and $g^{\mathrm{cap}}$
by the formulas of Definition~\ref{pd-orbit-generic-data}. Let $E$ be a finite set of nonnegative integers
and $i(s,e)$, $\mathrm{ch}(s,e)$ nonnegative real numbers that vanish for $e\notin E$.
\begin{enumerate}
\item Suppose that $\mathrm{ch}(s,e)=0$ when $2e+1\notin\Delta(s)$, and that for $s\in T$ and $e\geq0$,
  $\mathrm{ch}(s,e)+i(s,e)\leq A(s)/p^{2e+1}$ if $2e+1\in\Delta(s)$, and
  $i(s,e)\leq\min(A(s)/p^{2e+1},\mathrm{cap}(s+2e+1))$ if $2e+1\notin\Delta(s)$. Then
  \[
   g^{\mathrm{ex}}\leq\gamma_p\sum_{s\in T}A(s)-\sum_{s\in T}\sum_e\mathrm{ch}(s,e)-\sum_{s\in
   T}\sum_ei(s,e)-d.
  \]
\item Suppose that $i(s,e)\leq\min(A(s)/p^{2e+1},\mathrm{cap}(s+2e+1))$ for the odd $s\in T$ and all
  $e$, and let $x$ be a real number at least the sum of the $i(s,e)$ over the even $s\in T$ and all $e$.
  Then $g^{\mathrm{cap}}\leq\gamma_p\sum_{s\in T}A(s)-\sum_{s\in T}\sum_ei(s,e)-d+x$.
\end{enumerate}
\end{lemma}
\leanref{Formalism/ThetaLogCoordinates/UniformThetaSubstitution/FrobeniusLiftNorms/ABArith.lean}{Formalism.InvariantFactorPaths.G7.excessGainBound}
\leanref{Formalism/ThetaLogCoordinates/UniformThetaSubstitution/FrobeniusLiftNorms/ABArith.lean}{Formalism.InvariantFactorPaths.G7.cappedGainBound}
\leanref{Formalism/GrowthSums/FactorialBaseGrowth/PathBasisOfCount.lean}{Formalism.GrowthSums.cappedGainSubLE}

\begin{proof}
(1) Since $\sum_{e\geq0}p^{-2e-1}=\gamma_p$, $\gamma_pA(s)=\sum_eA(s)/p^{2e+1}$. Fix $s\in T$ and $e$. If
$2e+1\in\Delta(s)$, the $e$-th term of $g^{\mathrm{ex}}$ is $0\leq A(s)/p^{2e+1}-\mathrm{ch}(s,e)-i(s,e)$.
Otherwise $\mathrm{ch}(s,e)=0$, and $\max(A(s)/p^{2e+1}-\mathrm{cap}(s+2e+1),0)\leq A(s)/p^{2e+1}-i(s,e)$,
since $i(s,e)$ is at most both $\mathrm{cap}(s+2e+1)$ and $A(s)/p^{2e+1}$. Summing over $e$ and $s\in T$
gives (1).

(2) For odd $s\in T$, $\sum_ei(s,e)\leq\sum_{e\geq0}\min(A(s)/p^{2e+1},\mathrm{cap}(s+2e+1))$, the series
converging since its terms are at most $A(s)/p^{2e+1}$. With the hypothesis on $x$,
$\sum_{s\in T}\sum_ei(s,e)-x$ is at most the middle sum in $g^{\mathrm{cap}}$, which gives (2).
\end{proof}

\begin{lemma}\label{pd-orbit-generic-entry-bounds}
Let $p>2$ be a prime and $K$ a field with an ultrametric absolute value such that $|p|=p^{-1}$, with unit
ball $R$. Take entry data over $K$ (Definition~\ref{pd-orbit-generic-entry}) with $\xi_{s,a}\in R$ and
$R_c(r,a)\in R$, entry polynomial $P$, counts $A$, $\mathrm{low}$ and gains $g^{\mathrm{ex}}$ and
$g^{\mathrm{cap}}$. Let $u=(u_s)$ be a family in $R$ and $F$ a power series over $K$ in the $Z$ that lies
in the class of $H_u$ over $R$, and put $\delta=\sum_s(A(s)+\mathrm{low}(s))$. Then:
\begin{enumerate}
\item $|\langle P,F\rangle_Z|\leq p^{-g^{\mathrm{ex}}-\rho_p\delta}$;
\item $|\langle P,F\rangle_Z|\leq p^{-g^{\mathrm{cap}}+\sum_sA(s)-\rho_p\delta}$, the sum over the even
  $s\in T$;
\item for every $J\subset\{0,\dots,N_1-1\}$, $|p^{N_J}\langle P,F\rangle_Z|\leq
  p^{-g^{\mathrm{cap}}-\rho_p\delta}$, where $N_J$ is the sum of the $e(s,b)$ over the pairs $(s,b)$ with
  $s$ even and $b\notin J$ or $s$ odd and $b\in J$, plus the sum over the lower columns $c$ with $s_c$ odd
  of the number of elements of $J$ in the image of $I_c$.
\end{enumerate}
\end{lemma}
\leanref{Formalism/ThetaLogCoordinates/UniformThetaSubstitution/FrobeniusLiftNorms/EntryAOneParity.lean}{Formalism.InvariantFactorPaths.G7.entryExcessBound_of}
\leanref{Formalism/ThetaLogCoordinates/UniformThetaSubstitution/FrobeniusLiftNorms/EntryATwoAndB.lean}{Formalism.InvariantFactorPaths.G7.entryCappedBound_of}
\leanref{Formalism/ThetaLogCoordinates/UniformThetaSubstitution/FrobeniusLiftNorms/EntryATwoAndB.lean}{Formalism.InvariantFactorPaths.G7.EntryB_of}
\leanref{Formalism/ThetaLogCoordinates/UniformThetaSubstitution/FrobeniusLiftNorms/ABArith.lean}{Formalism.InvariantFactorPaths.G7.cyclicShiftLeast}
\leanref{Formalism/ThetaLogCoordinates/UniformThetaSubstitution/KOneConsequence.lean}{Formalism.InvariantFactorPaths.G7.Top.entry_a_one}
\leanref{Formalism/ThetaLogCoordinates/UniformThetaSubstitution/KOneConsequence.lean}{Formalism.InvariantFactorPaths.G7.Top.entry_a_two}
\leanref{Formalism/ThetaLogCoordinates/UniformThetaSubstitution/KOneConsequence.lean}{Formalism.InvariantFactorPaths.G7.Top.entryB}

\begin{proof}
If $\langle P,F\rangle_Z=0$ there is nothing to prove. Otherwise take $E$, $\alpha$, $\beta$ and $v$ from
Lemma~\ref{pd-orbit-generic-star}, so that $|\langle P,F\rangle_Z|\leq p^{-v-\rho_p\delta}$. For a
position $s$ and $l\geq0$ let $i(s,l)=\sum_b\beta_{(s,l,b)}$, with $i(s,l)=0$ for $l\geq E$.

\emph{Step 1: the counts of $\beta$.} By (b), $p^{2l+1}i(s,l)\leq A(s)$, since $d_k$ is $p^{2l+1}$ at
$W_{s,b}$ for $k=(s,l,b)$. The $b_k$ with $k=(s,l,b)$, $b<N_1$, have the same position $s+2l+1$ and the
same vector $\xi_s^{\,p^{2l+1}}$; if $i(s,l)>\mathrm{cap}(s+2l+1)$, then $D^\beta_bQ=0$ by
Lemma~\ref{pd-orbit-generic-minors}(2), against (c). So
$i(s,l)\leq\min(A(s)/p^{2l+1},\mathrm{cap}(s+2l+1))$ for all $s$ and $l$. By (a),
$\sum_{s\in T}\sum_li(s,l)=|\beta|$.

\emph{Step 2: part (1).} Let $\phi=(s,l,l',a,b)$ with $\alpha_\phi\neq0$, and put $s_1=s-2l$ and
$s_2=s+1-2l'$, so that $d_\phi$ is $p^{1+2l}$ at $W_{s_1,b}$ and $p^{2l'}$ at $W_{s_2,b}$, and
$s_1+(1+2l)=s_2+2l'$. If $l'\leq l$, let $\sigma_\phi=s_1$ and $\sigma'_\phi=s_2$; the least integer
$m\geq1$ with $s_1+m=s_2$ is at most $(1+2l)-2l'\leq1+2l$. If $l'>l$, let $\sigma_\phi=s_2$ and
$\sigma'_\phi=s_1$; then $s_2+2l'=s_1+(1+2l)$ with $1+2l<2l'$, and the least $m\geq1$ with $s_2+m=s_1$ is
at most $2l'$. In both cases $\sigma'_\phi-\sigma_\phi$ is odd in $\mathbf Z/n'\mathbf Z$, so that least
$m$ is its odd representative $2e_\phi+1$, the positions $\sigma_\phi$ and $\sigma'_\phi$ have different
parities and lie in $T$ by (a), and $p^{2e_\phi+1}$ is at most the exponent of $d_\phi$ at
$W_{\sigma_\phi,b}$. Let $\mathrm{ch}(s,e)$ be the sum of the $\alpha_\phi$ with $\sigma_\phi=s$ and
$e_\phi=e$. Summing (b) over $b$ at the position $s$ gives $p^{2e+1}(\mathrm{ch}(s,e)+i(s,e))\leq A(s)$.
If $\mathrm{ch}(s,e)\neq0$, some $\phi$ with $\alpha_\phi\neq0$ has $\sigma_\phi=s$ and $e_\phi=e$, and
the representative $2e+1$ of $\sigma'_\phi-s$ comes from $\sigma'_\phi\in T$ of the other parity; so
$2e+1\in\Delta(s)$. As $\sigma_\phi\in T$ when $\alpha_\phi\neq0$,
$\sum_{s\in T}\sum_e\mathrm{ch}(s,e)=\sum_\phi\alpha_\phi$. With Step 1,
Lemma~\ref{pd-orbit-generic-gain-sums}(1) gives $g^{\mathrm{ex}}\leq v$, hence (1).

\emph{Step 3: part (2).} If $\alpha_\phi\neq0$, the positions $s_1$ and $s_2$ lie in $T$ and have
different parities, so one of them is even, and $d_\phi$ is at least $1$ at the corresponding
$W_{s_i,b}$. Summing (b) over the $W_{s,b}$ with $s\in T$ even gives
$\sum_\phi\alpha_\phi+\sum_{s,e}i(s,e)\leq\sum_sA(s)$, the sums over the even $s\in T$.
Lemma~\ref{pd-orbit-generic-gain-sums}(2), with $x=\sum_sA(s)-\sum_\phi\alpha_\phi$ (the sum over the even
$s\in T$) and Step 1, gives $g^{\mathrm{cap}}\leq v+\sum_sA(s)$, hence (2).

\emph{Step 4: part (3).} Let $\mathcal P$ be the set of pairs $(s,b)$ with $s$ even and $b\notin J$ or $s$
odd and $b\in J$. If $\alpha_\phi\neq0$, then $s_1$ and $s_2$ have different parities, so one of
$(s_1,b)$ and $(s_2,b)$ lies in $\mathcal P$. Summing (b) over the $W_{s,b}$ with $(s,b)\in\mathcal P$
gives $\sum_\phi\alpha_\phi+\sum\beta_k\leq\sum_{(s,b)\in\mathcal P}e(s,b)$, the middle sum over the
$k=(s,l,b)$ with $(s,b)\in\mathcal P$. For $k=(s,l,b)$ with $s$ even and $b\in J$, the position
$s+2l+1$ of $b_k$ is odd. For an odd position $y$, the $k$ with $s$ even, $b\in J$ and $s+2l+1=y$ have
$b_k$ at $y$ with columns in $J$; by Lemma~\ref{pd-orbit-generic-minors}(2) and (c), the sum of their
$\beta_k$ is at most the sum over the lower columns $c$ with $s_c=y$ of the number of elements of $J$ in
the image of $I_c$. Summing over $y$, and since every $k=(s,l,b)$ with $s\in T$ even has
$(s,b)\in\mathcal P$ or $b\in J$, $\sum_\phi\alpha_\phi+\sum_{s,e}i(s,e)\leq N_J$, the middle sum over the
even $s\in T$. Lemma~\ref{pd-orbit-generic-gain-sums}(2) with $x=N_J-\sum_\phi\alpha_\phi$ gives
$g^{\mathrm{cap}}\leq v+N_J$. So $|p^{N_J}\langle P,F\rangle_Z|=p^{-N_J}|\langle P,F\rangle_Z|\leq
p^{-N_J-v-\rho_p\delta}\leq p^{-g^{\mathrm{cap}}-\rho_p\delta}$.
\end{proof}

\begin{lemma}\label{pd-orbit-generic-reduction}
Let $p$ be a prime and $K'$ a field of characteristic $0$ with an ultrametric absolute value such that
$|p|=p^{-1}$, with unit ball $R'$. Take $n'$, $N_0$, $N_1$, a column datum $I$ and a set $T$ of positions
(Definition~\ref{pd-orbit-generic-data}), a family $u=(u_s)$ in $R'$, and a power series $F$ over $K'$ in
the $Z$ that lies in the class of $H_u$ over $R'$. Let $K''$ be the field of fractions of the polynomial
ring $K'[Y]$ in the variables $Y_{s,r,a}$, with the absolute value whose value at a polynomial is the
largest absolute value of its coefficients, and $R''$ its unit ball. Form the entry data over $K''$ with
the set $T$, $\xi_{s,a}=Y_{s,0,a}$, $e(s,b)$ the number of $c\in C_s$ with $h_{s,c}=1$ and
$I_{s,c}(0)=b$ if $s\in T$ and $e(s,b)=0$ if $s\notin T$, and as lower columns the pairs $(s,c)$ with
$c\in C_s$ and not both $s\in T$ and $h_{s,c}=1$, each with position $s$, length $h_{s,c}$, rows
$Y_{s,r,a}$ and map $I_{s,c}$. Then:
\begin{enumerate}
\item this absolute value on $K''$ is ultrametric and extends that of $K'$, the $Y_{s,r,a}$ lie in $R''$,
  and $F$ lies in the class of $H_u$ over $R''$;
\item every coefficient $y$ of $\langle G_I,F\rangle_Z$ satisfies $|y|\leq|\langle P,F\rangle_Z|$, where
  $P$ is the entry polynomial;
\item the counts $A$, $\mathrm{cap}$ and $\mathrm{low}$ of the entry data are those of $(I,T)$;
\item for $i\geq1$ and $J=\{0,\dots,i-1\}$, the integer $N_J$ of
  Lemma~\ref{pd-orbit-generic-entry-bounds}(3) is at most the edge exponent $e_I(i)$.
\end{enumerate}
\end{lemma}
\leanref{Formalism/ThetaLogCoordinates/UniformThetaSubstitution/FrobeniusLiftNorms/EGaussK.lean}{Formalism.InvariantFactorPaths.G7.EGaussK.reduce}
\leanref{Formalism/ThetaLogCoordinates/UniformThetaSubstitution/FrobeniusLiftNorms/EGaussK.lean}{Formalism.InvariantFactorPaths.G7.generic_elements_exist}
\leanref{Formalism/ThetaLogCoordinates/UniformThetaSubstitution/FrobeniusLiftNorms/EGauss.lean}{Formalism.InvariantFactorPaths.G7.EGauss.struct}
\leanref{Formalism/ThetaLogCoordinates/UniformThetaSubstitution/FrobeniusLiftNorms/EGauss.lean}{Formalism.InvariantFactorPaths.G7.EGauss.inClass_map}
\leanref{Formalism/GaussJetBounds/GaussJetProofs/GField.lean}{Formalism.DworkPowerSeriesWeights.gAbs}
\leanref{Formalism/GaussJetBounds/GaussJetProofs/GField.lean}{Formalism.DworkPowerSeriesWeights.gN_mul}

\begin{proof}
(1) The Gauss norm on $K'[Y]$ is multiplicative and ultrametric \cite[\S5.1.2]{BGR84}, so it extends to
an ultrametric absolute value on $K''$; on constants it is the absolute value of $K'$, and
$|Y_{s,r,a}|=1$. The class of $H_u$ is defined by identities between coefficients, and the series $U_i$
over $R'$ are series over $R''$.

(2) The frame jet $\langle G_I,F\rangle_Z$ is a polynomial in the $Y$ over $K'$, so its absolute value as
an element of $K''$ is the largest absolute value of its coefficients, and it is at least $|y|$. As an
element of $K''$ it is the jet pairing $\langle G_I,F\rangle_Z$ with $G_I$ read as a polynomial in the $Z$
over $K''$, and this polynomial is $P$. Indeed, a column $c\in C_s$ with $s\in T$ and $h_{s,c}=1$ gives
the $1\times1$ minor $\sum_aY_{s,0,a}Z_{s,a,I_{s,c}(0)}=f_{s,I_{s,c}(0)}$, and grouping equal factors
gives $\prod_{s,b}f_{s,b}^{e(s,b)}$; the other columns give the factors $Q_c$ of the lower product.

(3) $\sum_be(s,b)$ is the number of columns of height $1$ at $s$ if $s\in T$, and $0$ otherwise; the
lower columns at $s$ are the other columns of $C_s$, and their lengths add up to $\sum_ch_{s,c}-A(s)$.

(4) Compare the contributions of each column $c\in C_s$. If $s\in T$ and $h_{s,c}=1$, with
$I_{s,c}(0)=b$, the column contributes $1$ to $N_J$ if $s$ is even and $b\geq i$ or $s$ is odd and
$b<i$, and $0$ otherwise. To $e_I(i)$ it contributes $\min(1,i)-n_{s,c}(i)$ if $s$ is even, which is $1$
if $b\geq i$ and $0$ otherwise, and $n_{s,c}(i)$ if $s$ is odd, which is $1$ if $b<i$ and $0$ otherwise.
These agree. A lower column at an odd position contributes $n_{s,c}(i)$ to both, and a lower column at an
even position contributes $0$ to $N_J$ and $\min(h_{s,c},i)-n_{s,c}(i)\geq0$ to $e_I(i)$.
\end{proof}

\begin{proposition}\label{pd-orbit-generic-bounds}
Let $p>2$ be a prime and $K'$ a field of characteristic $0$ with an ultrametric absolute value such that
$|p|=p^{-1}$, with unit ball $R'$, and let $\pi\in K'$ with $\pi^{p^2-1}=p$. Take $n'$, $N_0$, $N_1$, a
column datum $I$ and a set $T$ of positions, with the counts $A$, the gains $g^{\mathrm{ex}}$ and
$g^{\mathrm{cap}}$ and the edge exponents $e_I(i)$ of $(I,T)$ (Definition~\ref{pd-orbit-generic-data}).
Let $u=(u_s)$ be a family in $R'$ and $F$ a power series over $K'$ in the $Z$ such that $F(\pi Z)$, the
series $F$ with $\pi Z_{s,a,b}$ substituted for each $Z_{s,a,b}$, lies in the class of $H_u$ over $R'$.
Then every coefficient $y$ of $\langle G_I,F\rangle_Z$ satisfies:
\begin{enumerate}
\item $|y|\leq p^{-g^{\mathrm{ex}}}$;
\item $|y|\leq p^{-g^{\mathrm{cap}}+\sum_sA(s)}$, the sum over the even $s\in T$;
\item $|p^{e_I(i)}y|\leq p^{-g^{\mathrm{cap}}}$ for every $i\geq1$.
\end{enumerate}
\end{proposition}
\leanref{Formalism/ThetaLogCoordinates/UniformThetaSubstitution/FrobeniusLiftNorms/EGaussK.lean}{Formalism.InvariantFactorPaths.G7.kernel_a_one_generic}
\leanref{Formalism/ThetaLogCoordinates/UniformThetaSubstitution/FrobeniusLiftNorms/EGaussK.lean}{Formalism.InvariantFactorPaths.G7.kernel_a_two_generic}
\leanref{Formalism/ThetaLogCoordinates/UniformThetaSubstitution/FrobeniusLiftNorms/EGaussK.lean}{Formalism.InvariantFactorPaths.G7.kernel_b_generic}
\leanref{Formalism/ThetaLogCoordinates/UniformThetaSubstitution/FrobeniusLiftNorms/EGaussK.lean}{Formalism.InvariantFactorPaths.G7.gauss_a_one_generic}
\leanref{Formalism/ThetaLogCoordinates/UniformThetaSubstitution/FrobeniusLiftNorms/EGaussK.lean}{Formalism.InvariantFactorPaths.G7.gauss_a_two_generic}
\leanref{Formalism/ThetaLogCoordinates/UniformThetaSubstitution/FrobeniusLiftNorms/EGaussK.lean}{Formalism.InvariantFactorPaths.G7.gauss_b_generic}
\leanref{Formalism/ThetaLogCoordinates/UniformThetaSubstitution/FrobeniusLiftNorms/EGaussK.lean}{Formalism.InvariantFactorPaths.G7.EGaussK.descend}
\leanref{Formalism/ThetaLogCoordinates/UniformThetaSubstitution/FrobeniusLiftNorms/EKernel.lean}{Formalism.InvariantFactorPaths.G7.DeterminantJetScaling.frameJet_scale}

\begin{proof}
\emph{Step 1: series in the class.} Let $F_1$ be a power series over $K'$ in the $Z$ in the class of
$H_u$ over $R'$, and $\delta=\sum_s\sum_{c\in C_s}h_{s,c}=\sum_s(A(s)+\mathrm{low}(s))$. Every
coefficient $y_1$ of $\langle G_I,F_1\rangle_Z$ satisfies (1) to (3) with the right sides multiplied by
$p^{-\rho_p\delta}$. To see this, form $K''$ and the entry data of Lemma~\ref{pd-orbit-generic-reduction}.
By its parts (1) and (3), Lemma~\ref{pd-orbit-generic-entry-bounds} applies over $K''$ to $F_1$, with the
counts and gains of $(I,T)$, and by its part (2), $|y_1|\leq|\langle P,F_1\rangle_Z|$. Parts (1) and (2)
of Lemma~\ref{pd-orbit-generic-entry-bounds} give (1) and (2). For (3) let $J=\{0,\dots,i-1\}$; since
$N_J\leq e_I(i)$ by Lemma~\ref{pd-orbit-generic-reduction}(4) and $|p|<1$,
\[
 |p^{e_I(i)}y_1|\leq p^{-N_J}|\langle P,F_1\rangle_Z|=|p^{N_J}\langle P,F_1\rangle_Z|\leq
 p^{-g^{\mathrm{cap}}-\rho_p\delta}
\]
by Lemma~\ref{pd-orbit-generic-entry-bounds}(3).

\emph{Step 2: scaling.} The polynomial $G_I$ is homogeneous of degree $\delta$ in the $Z$, so only the
coefficients of $F$ of degree $\delta$ enter $\langle G_I,F\rangle_Z$, and those of $F(\pi Z)$ are
$\pi^\delta$ times them; hence $\langle G_I,F(\pi Z)\rangle_Z=\pi^\delta\langle G_I,F\rangle_Z$. As
$|\pi|^{p^2-1}=p^{-1}$, $|\pi^\delta|=p^{-\rho_p\delta}$. Step 1 with $F_1=F(\pi Z)$ bounds $|\pi^\delta y|$
by the bounds (1) to (3) times $p^{-\rho_p\delta}$, and dividing by $|\pi^\delta|$ gives (1) to (3).
\end{proof}

\begin{proposition}\label{pd-kernel-orbit-generic}
Let $p\geq5$ be a prime and take a formal group system at $p$ (Definition~\ref{pd-formal-group-system}),
with the notation of Definitions~\ref{pd-formal-group-numbers}, \ref{pd-hecke-chain-operators},
\ref{pd-kernel-orbit-index} and~\ref{pd-kernel-orbit-series}. Let $v$ be an orbit and $j<m_*$. Let $K'$
be a field of characteristic $0$ with an ultrametric absolute value such that $|p|=p^{-1}$, $R'$ the ring
of $x\in K'$ with $|x|\leq1$, $\pi\in K'$ with $\pi^{p^2-1}=p$, and $u=(u_s)$ a family in $R'$ indexed
by the positions of $v$. Let $F$ be a power series over $K'$ in the variables $Z_{v,s,a,b}$ at $v$ such
that $F(\pi Z)$ lies in the class of $H_u$ over $R'$. Let $I$ be a column choice into $N(j+1)$ for the
natural weights of $\lambda_{j+1}(v,s)$, $0\leq i\leq r$, and $y$ a coefficient of
$\langle G^{(v),i}_I,F\rangle_Z$.
\begin{enumerate}
\item If $i=0$, then $|y|\leq p^{-g_A(v,j+1)}$ and $|y|\leq p^{-g_B(v,j+1)+\alpha_{\mathrm{ev}}(v)}$.
\item If $1\leq i\leq r$, then $|y|\leq p^{-g_B(v,j+1)}$.
\end{enumerate}
\end{proposition}
\leanref{Formalism/ThetaLogCoordinates/UniformThetaSubstitution/FrobeniusLiftNorms/EGaussK.lean}{Formalism.InvariantFactorPaths.G7.kernel_a_one_generic}
\leanref{Formalism/ThetaLogCoordinates/UniformThetaSubstitution/FrobeniusLiftNorms/EGaussK.lean}{Formalism.InvariantFactorPaths.G7.kernel_a_two_generic}
\leanref{Formalism/ThetaLogCoordinates/UniformThetaSubstitution/FrobeniusLiftNorms/EGaussK.lean}{Formalism.InvariantFactorPaths.G7.kernel_b_generic}
\leanref{Formalism/SharpRestriction/GrowthBounds/MultisetValuation/GainOfCapacityLoss.lean}{Formalism.MainNotions.FormalGroupSystem.Pf.ea_brNum}
\leanref{Formalism/SharpRestriction/GrowthBounds/MultisetValuation/EbBrJet.lean}{Formalism.MainNotions.FormalGroupSystem.Pf.eb_brJet}

\begin{proof}
Apply Proposition~\ref{pd-orbit-generic-bounds} with $n'=n'_v$, with $N_0$ and $N_1$ the numbers of
values of the indices $a$ and $b$ of the variables $Z_{v,s,a,b}$, with the top positions $T_v$, and with
the column datum whose columns at a position $s$ are the columns of the natural weight of
$\lambda_{j+1}(v,s)$, their heights the column lengths and their maps those of $I$. The pairing
$\langle G^{(v),0}_I,F\rangle_Z$ is the pairing $\langle G_I,F\rangle_Z$ of that proposition. For these
diagrams the gains $g^{\mathrm{ex}}$ and $g^{\mathrm{cap}}$ are $g_A(v,j+1)$ and $g_B(v,j+1)$, and the
sum of the counts $A(s)$ over the even $s\in T_v$ is $\alpha_{\mathrm{ev}}(v)$; this is a computation with
the weights of the fills. So parts (1) and (2) of the proposition give (1). In (2),
$G^{(v),i}_I=p^{e}G^{(v),0}_I$, where $e$ is the sum over the positions $s$ of $v$ and the columns $c$ of
the exponents of condition~(i), and the bound for $G^{(v),i}_I$ is a bound for $p^{e}$ times the
coefficients of $\langle G^{(v),0}_I,F\rangle_Z$. The maps $I_{v,s,c}$ are strictly increasing, so the
number of $q<|c|$ with $I_{v,s,c}(q)<i$ is the number of elements of the image of $I_{v,s,c}$ less than
$i$, and $e$ is the edge exponent $e_I(i)$. Part (3) of the proposition gives (2).
% Lean: kernel_a_one_generic, kernel_a_two_generic, kernel_b_generic (EGaussK.lean), applied in
% Pf.ea_A1U, Pf.ea_A2U and Pf.ea_BU (EaGeCOfNorm.lean). The reduction to frame jets of the restricted
% series is eb_brJet, ea_brClass and ea_brScale; the gains are identified by ea_brNum.
\end{proof}

\begin{lemma}\label{pd-kernel-orbit-bounds}
Let $p\geq5$ be a prime and take an admitted formal group system at $p$
(Definition~\ref{pd-formal-group-system}), with the notation of Definitions~\ref{pd-formal-group-numbers},
\ref{pd-hecke-chain-operators} and~\ref{pd-kernel-orbit-index}. Let $v$ be an orbit and $j<m_*$. Then
the system has lower bound $g_A(v,j+1)$ at $(j,v)$, lower bound $g_B(v,j+1)-\alpha_{\mathrm{ev}}(v)$ at
$(j,v)$, and lower bound $g_B(v,j+1)$ at $(j,v)$ with index $i$ for every $1\leq i\leq r$.
\end{lemma}
\leanref{Formalism/SharpRestriction/GrowthBounds/MultisetValuation/FirstTheoremPlaceBound.lean}{Formalism.MainNotions.FormalGroupSystem.Pf.first_theorem_place_bound}
\leanref{Formalism/SharpRestriction/GrowthBounds/MultisetValuation/FirstTheoremPlaceBound.lean}{Formalism.MainNotions.FormalGroupSystem.Pf.second_lemma_place_bound}
\leanref{Formalism/SharpRestriction/GrowthBounds/MultisetValuation/FirstTheoremPlaceBound.lean}{Formalism.MainNotions.FormalGroupSystem.Pf.hecke_place_bound}
\leanref{Formalism/SharpRestriction/GrowthBounds/MultisetValuation/EaGeCOfNorm.lean}{Formalism.MainNotions.FormalGroupSystem.Pf.ea_lowN}
\leanref{Formalism/SharpRestriction/GrowthBounds/MultisetValuation/NmRestNormOfPlace.lean}{Formalism.MainNotions.FormalGroupSystem.Pf.nm_restNorm_of_place}

\begin{proof}
\emph{Step 1: an absolute value.} The ring $\mathcal R_c$ has rank one, and $p$ is not a unit in it,
because $p$ is not a unit in $\mathcal R$, which is finite over $W_v$, and
$\mathcal R_c\cap\mathcal K=\mathcal R$. So there is an ultrametric absolute value on $\mathcal K_c$ with
$|p|=p^{-1}$ and $|x|\leq1$ if and only if $x\in\mathcal R_c$; for $y\in\mathcal K_c$, $|y|\leq p^{-g}$
implies $v(y)\geq g$. It extends to an algebraic closure $K'$ of $\mathcal K_c$
\cite[Chapter~II, \S8]{Neu99}. Let $R'$ be the ring of $x\in K'$ with $|x|\leq1$, and $\pi\in K'$ with
$\pi^{p^2-1}=p$; then $R'$ contains $\mathcal R_c$, $W(\overline{\mathbf F}_p)$ and $\pi$.

\emph{Step 2.} Let $U$ be a series in the variables $X_e,X'_e$ with $U(\ell(X))$ integral over
$\mathcal R_c$. By Lemma~\ref{pd-kernel-orbit-class}(2), $\Theta^Z_U(\pi Z)$ lies in the class of $H_u$
over $R'$ for a family $u$ in $R'$. Let $F$ be $\Theta^Z_U$ with the variables at the other orbits set to
$0$. As $H_u$ involves only the variables at $v$, $F(\pi Z)$ also lies in the class of $H_u$ over $R'$,
and $\langle G^{(v),i}_I,F\rangle_Z=\langle G^{(v),i}_I,\Theta^Z_U\rangle_Z$ because $G^{(v),i}_I$
involves only the variables at $v$. By Proposition~\ref{pd-kernel-orbit-generic} every coefficient $y$ of
this pairing has $|y|\leq p^{-g}$, with $g=g_A(v,j+1)$ or $g=g_B(v,j+1)-\alpha_{\mathrm{ev}}(v)$ for
$i=0$ and $g=g_B(v,j+1)$ for $1\leq i\leq r$. Hence $v(y)\geq g$.
\end{proof}

\begin{lemma}\label{pd-hecke-chain-step}
Let $p\geq5$ be a prime and take an admitted formal group system at $p$
(Definition~\ref{pd-formal-group-system}) with $e_{\mathcal R}\geq1$, with the notation of
Definitions~\ref{pd-formal-group-numbers} and~\ref{pd-hecke-chain-operators}; for each orbit $v$ fix
$u_l$ and $m_l$ as in Lemma~\ref{pd-hecke-chain-commutation} and form the orbit operators $\mathcal
H_{v,J}$. Let $j<m_*$, let $B$ be a set of orbits, and for each orbit $v$ let $g_v\in\mathbf R$, with
$g_v\leq g_B(v,j+1)$ for $v\in B$, and such that the system has lower bound $g_v$ at $(j,v)$ for
$v\notin B$. Let $\mathcal H_B$ be the composite of the $\mathcal H_{v,j+1}$, $v\in B$, in a fixed order.
Then for every $G\in\mathrm{Int}_j$ the function $\mathcal H_B(K_jG)$ has valuation at least
$\sum_vg_v-d_\Gamma(j)$.
\end{lemma}
\leanref{Formalism/SharpRestriction/GrowthBounds/MultisetValuation/AofUpdate.lean}{Formalism.MainNotions.FormalGroupSystem.Pf.ZL_bound}
\leanref{Formalism/SharpRestriction/GrowthBounds/MultisetValuation/AofUpdate.lean}{Formalism.MainNotions.FormalGroupSystem.Pf.Zop_gen}
\leanref{Formalism/SharpRestriction/GrowthBounds/MultisetValuation/RamifiedLatticeOfChainEstimates.lean}{Formalism.MainNotions.FormalGroupSystem.Pf.form_of_place}
\leanref{Formalism/SharpRestriction/GrowthBounds/MultisetValuation/RestEntOf.lean}{Formalism.MainNotions.FormalGroupSystem.Pf.restEnt_of}

\begin{proof}
Each $\mathcal H_{v,j+1}$ is the sum over $0\leq i\leq r$ of
$(-1)^{r-i}c_{r-i}(j+1)(u_1\cdots u_{j+1})^{-i}T^\natural_{v,i}$, with the $u_l$ and $c_l$ of $v$.
Expanding the composite,
\[
 \mathcal H_B(K_jG)=\sum_{i_0}a_{i_0}\,T^\natural_{i_0}K_jG,\qquad a_{i_0}=\prod_{v\in
 B}(-1)^{r-i_0(v)}c_{r-i_0(v)}(j+1)(u_1\cdots u_{j+1})^{-i_0(v)},
\]
the sum over the functions $i_0$ from the orbits to $\{0,\dots,r\}$ that vanish outside $B$, where
$T^\natural_{i_0}$ is the composite of the $T^\natural_{v,i_0(v)}$ in the order of $\mathcal H_B$. By
Lemma~\ref{pd-hecke-chain-commutation}(2), $a_{i_0}$ has valuation at least the sum of
$\alpha_{\mathrm{ev}}(v)$ over the $v\in B$ with $i_0(v)=0$.

Fix $i_0$, and put $g'_v=g_v-\alpha_{\mathrm{ev}}(v)$ if $v\in B$ and $i_0(v)=0$, and $g'_v=g_v$
otherwise. The system has lower bound $g'_v$ at $(j,v)$ with index $i_0(v)$ for every orbit $v$: for
$v\notin B$ by hypothesis; for $v\in B$ with $i_0(v)=0$ by Lemma~\ref{pd-kernel-orbit-bounds}, since
$g'_v\leq g_B(v,j+1)-\alpha_{\mathrm{ev}}(v)$; and for $v\in B$ with $i_0(v)\geq1$ by
Lemma~\ref{pd-kernel-orbit-bounds}, since $g'_v\leq g_B(v,j+1)$. Here a lower bound $g$ implies every
lower bound $g''\leq g$. Let $g''<\sum_vg'_v$, let $N$ be the number of orbits and
$\epsilon=(\sum_vg'_v-g'')/(N+1)$, and choose rationals $q_v$ with $g'_v-\epsilon<q_v<g'_v$. Then
$\sum_vq_v\geq g''$, and the system has lower bound $q_v$ at $(j,v)$ with index $i_0(v)$. By
Lemma~\ref{pd-kernel-orbit-product} the hypothesis of Lemma~\ref{pd-kernel-chain-bound} holds with
$\sum_vq_v$, hence with $g''$. So $T^\natural_{i_0}K_jG$ has valuation at least $g''-d_\Gamma(j)$ for
every $g''<\sum_vg'_v$. Put $y=(\sum_vg'_v-d_\Gamma(j))e_{\mathcal R}$. For $g''<\sum_vg'_v$ close
enough to $\sum_vg'_v$, $(g''-d_\Gamma(j))e_{\mathcal R}>\lceil y\rceil-1$, so
$\lceil(g''-d_\Gamma(j))e_{\mathcal R}\rceil=\lceil y\rceil$, and $T^\natural_{i_0}K_jG$ has valuation at
least $\sum_vg'_v-d_\Gamma(j)$. An element of $\mathcal K$ of valuation at least $x$ times a function of
valuation at least $y$ has valuation at least $x+y$, since $\lceil xe_{\mathcal R}\rceil+\lceil
ye_{\mathcal R}\rceil\geq\lceil(x+y)e_{\mathcal R}\rceil$. So $a_{i_0}T^\natural_{i_0}K_jG$ has valuation
at least $\sum_vg'_v+\sum_{v\in B,\,i_0(v)=0}\alpha_{\mathrm{ev}}(v)-d_\Gamma(j)=\sum_vg_v-d_\Gamma(j)$,
and so has the sum over $i_0$.
\end{proof}

\begin{theorem}\label{pd-formal-group-bounds}
Let $p\geq5$ be a prime and take an admitted formal group system at $p$
(Definition~\ref{pd-formal-group-system}), with the integers $m_*$ and $d_\Gamma(j)$, the composite
$K^{(m_*)}$ of the first $m_*$ maps of its chain, and the numbers and modules of
Definition~\ref{pd-formal-group-numbers}.
\begin{enumerate}
\item For every orbit $v$ and every $j<m_*$, the system has lower bound $g_A(v,j+1)$ at $(j,v)$.
\item Let $G$, $G_0$ and $P$ be real functions of an orbit $v$ and an integer $j$, let $S$ be a real
  function of the orbits, and let $V$ be a set of orbits. Suppose that $G(v,j)\leq g_B(v,j)$ and
  $\mathrm{Pr}_B(v,j)\leq P(v,j)$ for $v\in V$ and $1\leq j\leq m_*$; that the system has lower
  bound $G_0(v,j+1)$ at $(j,v)$ for every orbit $v$ and every $j<m_*$; and that
  $\sum_{j=1}^{m_*}\max\big(G(v,j)-P(v,j)-G_0(v,j),0\big)\leq S(v)$ for $v\in V$. Then for every
  $F\in I(S,V)$ the function $K^{(m_*)}F$ has valuation at least
  \[
   \sum_{j=1}^{m_*}\Big(\sum_{v\in V}\max\big(G(v,j)-P(v,j),\ G_0(v,j)\big)+\sum_{v\notin V}G_0(v,j)
   \Big)-\sum_{j<m_*}d_\Gamma(j)
  \]
  (Definition~\ref{pd-local-model}(7)).
\end{enumerate}
\end{theorem}
\leanref{Formalism/SharpRestriction/GrowthBounds/FormalGroupSystem.lean}{Formalism.MainNotions.FormalGroupSystem.VolumeRatio₂}
\leanref{Formalism/SharpRestriction/GrowthBounds/FormalGroupSystem.lean}{Formalism.MainNotions.FormalGroupSystem.LowTermBound}
\leanref{Formalism/SharpRestriction/GrowthBounds/MultisetValuation/PrvAndKernelCommutationUpToUnitsOfRests.lean}{Formalism.MainNotions.FormalGroupSystem.Pf.packet_statements_hold}
\leanref{Formalism/SharpRestriction/GrowthBounds/MultisetValuation/PrvAndKernelCommutationUpToUnitsOfRests.lean}{Formalism.MainNotions.FormalGroupSystem.Pf.stmt_proved}

\begin{proof}
\emph{The case $e_{\mathcal R}=0$.} Then there are no orbits, since for an orbit $v$ the ring
$\mathcal R$ is finite over $W_v$, in which $p$ is not a unit. Part~(1) is empty. In part~(2),
$I(S,V)=M$, and a function has valuation at least $g$ if and only if it lies in the integral module. By
Lemma~\ref{pd-kernel-chain-bound} with $i_0=0$ and $g=0$, whose hypothesis is condition~(c), $K_j$ maps
$\mathrm{Int}_j$ into $\mathrm{Int}_{j+1}$ for $j<m_*$, so $K^{(m_*)}$ maps $M$ into
$\mathrm{Int}_{m_*}$. Suppose from now on that $e_{\mathcal R}\geq1$, and use the notation of
Definition~\ref{pd-hecke-chain-operators}.

\emph{Part (1).} This is the first assertion of Lemma~\ref{pd-kernel-orbit-bounds}.

\emph{Part (2): reduction.} For each orbit fix $u_l$ and $m_l$ as in
Lemma~\ref{pd-hecke-chain-commutation}. The functions of valuation at least a given number form an
$\mathcal R$-submodule, so it suffices to treat a generator $F=y_{v_1}^{m_*}\cdots y_{v_l}^{m_*}(\varpi^nF_0)$
of $I(S,V)$, with $B=\{v_1,\dots,v_l\}\subset V$, $F_0\in M$ and $n=\sum_{v\in V\smallsetminus
B}\lceil S(v)e_{\mathcal R}\rceil$. For a set $C$ of orbits let $Y_C$ be the composite of the $y_v$,
$v\in C$; these commute and map $M$ into $M$ by condition~(j), so $F=\varpi^nY_B^{m_*}F_0$. For $j<m_*$
let $B_j$ be the set of $v\in B$ with $G_0(v,j+1)\leq G(v,j+1)-P(v,j+1)$, and put
\[
 b_j=\sum_{v\in B_j}\big(G(v,j+1)-\mathrm{Pr}_B(v,j+1)\big)+\sum_{v\notin B_j}G_0(v,j+1)-d_\Gamma(j).
\]

\emph{Part (2): induction.} We show that for $0\leq J\leq m_*$ and $F_1\in M$, the function
$K^{(J)}(Y_B^JF_1)$ has valuation at least $\sigma_J=\sum_{j<J}b_j$. For $J=0$ this is $F_1\in M$. Let
$J<m_*$, put $C=B\smallsetminus B_J$ and $W=Y_B^J(Y_CF_1)$, so that $W\in M$ and $Y_B^{J+1}F_1=Y_{B_J}W$.
By the induction hypothesis applied to $Y_CF_1\in M$, $K^{(J)}W=\varpi^{\lceil\sigma_Je_{\mathcal
R}\rceil}G'$ with $G'\in\mathrm{Int}_J$. Applying Lemma~\ref{pd-hecke-chain-commutation}(1) for the
orbits of $B_J$ in turn, to functions in the $\mathcal K$-span of $M$,
\[
 \mathcal H_{B_J}\big(K^{(J+1)}W\big)=p^{\sum_{v\in B_J}\mathrm{Pr}_B(v,J+1)}K^{(J+1)}(Y_{B_J}W),
\]
where $\mathcal H_{B_J}$ is the composite of the $\mathcal H_{v,J+1}$, $v\in B_J$. Hence
\[
 K^{(J+1)}(Y_B^{J+1}F_1)=p^{-\sum_{v\in B_J}\mathrm{Pr}_B(v,J+1)}\varpi^{\lceil\sigma_Je_{\mathcal
 R}\rceil}\mathcal H_{B_J}(K_JG').
\]
Apply Lemma~\ref{pd-hecke-chain-step} with the set $B_J$, $g_v=G(v,J+1)$ for $v\in B_J$, which is at most
$g_B(v,J+1)$ by hypothesis since $B_J\subset V$, and $g_v=G_0(v,J+1)$ for the other orbits, which is a
lower bound at $(J,v)$ by hypothesis. So $\mathcal H_{B_J}(K_JG')$ has valuation at least
$\sum_vg_v-d_\Gamma(J)$, and $K^{(J+1)}(Y_B^{J+1}F_1)$ has valuation at least
$\sigma_J+\sum_vg_v-d_\Gamma(J)-\sum_{v\in B_J}\mathrm{Pr}_B(v,J+1)=\sigma_{J+1}$.

\emph{Part (2): the bound.} With $J=m_*$ and $F_1=F_0$, $K^{(m_*)}F=\varpi^nK^{(m_*)}(Y_B^{m_*}F_0)$ lies in
$\varpi^{n+\lceil\sigma_{m_*}e_{\mathcal R}\rceil}\mathrm{Int}_{m_*}$. Let $T$ be the bound of~(2). We show
$T\leq\sigma_{m_*}+\sum_{v\in V\smallsetminus B}S(v)$, comparing the terms at $j+1$ orbit by orbit. For
$v\in B_j$, $\max(G-P,G_0)=G-P\leq G-\mathrm{Pr}_B$, since $P\geq\mathrm{Pr}_B$ on $V$. For $v\in
B\smallsetminus B_j$, $\max(G-P,G_0)=G_0$. For $v\in V\smallsetminus B$,
$\max(G-P,G_0)\leq\max(G-P-G_0,0)+G_0$, and the sum over $j$ of the first term is at most $S(v)$. For
$v\notin V$ both terms are $G_0$. Hence $\lceil Te_{\mathcal R}\rceil\leq\lceil\sigma_{m_*}e_{\mathcal
R}\rceil+n$, and $K^{(m_*)}F$ has valuation at least $T$.
\end{proof}

\begin{lemma}\label{pd-served-numerics}
Let $D$ be a definite-side datum with ranks $k,m,t,n$, size data with scale $(x,\varepsilon,J)$
and derived quantity $a$, and gain data with the constants $K_{\delta,1}$, $K_{\delta,2}$
(Definition~\ref{pd-served-selection}) and $K_{\mathrm{Pr}}$ (Definition~\ref{pd-served-lattice})
of $K_G$. Use the notation of Definitions~\ref{pd-served-lattice}
and~\ref{pd-formal-group-numbers}.
\begin{enumerate}
\item Suppose that $K_{\delta,1}\geq5$, $K_{\delta,2}\geq4$, $K_{\mathrm{Pr}}\geq1$,
  $r(\varepsilon)<k$ and $x\geq2$, that the size data satisfy the further conditions of
  Definition~\ref{pd-size-data}, and that $D$ carries a lattice instance tied to the gain data by
  a tie of the local models (Definition~\ref{pd-lattice-instance}). Let $u$ be a word, $w$ a place
  admissible for $u$ with $p=p_w\leq2a$, $t_0$ a representative at $p$ and $1\leq j\leq m$, and
  give the orbit $v=v(t_0)$ of $G(u,w)$ the fills of the lattice instance. Then
  \[
   g_U(p,j,t_0)\leq g_B(v,j),\qquad g_0(p,j,t_0)\leq g_A(v,j),\qquad
   \mathrm{Pr}_B(v,j)\leq\mathrm{Pr}(p,j,t_0).
  \]
\item Suppose that the constants of $K_G$ are nonnegative, that $K_{\mathrm{Pr}}\geq1$, and that
  $D$ has parameter choices for its size data (Definition~\ref{pd-family}(4)). Let $w$ be a place
  with $p=p_w>2a$, $t_0$ a representative at $p$ and $1\leq j\leq m$. Then
  $g_U(p,j,t_0)\leq\mathrm{Pr}(p,j,t_0)$ and $g_0(p,j,t_0)\leq0$.
\end{enumerate}
\end{lemma}
\leanref{Formalism/SharpRestriction/GrowthBounds/FormalGroupSystem.lean}{Formalism.MainNotions.FormalGroupSystem.SmallResidueCharComparison}
\leanref{Formalism/SharpRestriction/GrowthBounds/ProofLW.lean}{Formalism.MainNotions.FormalGroupSystem.Pf.numeric_comparison_of_constants}
\leanref{Formalism/SharpRestriction/GrowthBounds/TopChoiceBounds.lean}{Formalism.MainNotions.G.TopChoiceBounds.hLarge_proved}

\begin{proof}
(1) Under the tie, the positions of $v$ are the residues $s$ modulo $n'=\ell_p$, the position $s$
having element $\tau_1(t_0+sh_p)$; it is a top position if and only if $t_0+sh_p\in T$, and then
$\lambda_0(v,s)_1=a_{t_0+sh_p}$; and $w(v,j,s)=\mathrm{fill}_j((0,t_0+sh_p))$. Since $p$ is of inert
class, $\operatorname{pos}_p((0,t_0+sh_p),e)=(0,t_0+(s+e)h_p)$. Put $\alpha'(s)=a_{t_0+sh_p}-2c_G$ and
$z_j(s)=\mathrm{fill}_j((0,t_0+sh_p))+c_G$. Then $g_U(p,j,t_0)$ and $g_0(p,j,t_0)$ are the
expressions of Definition~\ref{pd-formal-group-numbers}(1) and~(2) with $\alpha'$ in place of
$\alpha$, $z_j$ in place of the caps $c_j$, the excluded odd steps in place of the sets $\Delta(s)$,
and $\delta(p,t_0)$ in place of $d_j$. The tie and the bounds of the fills give $\alpha'(s)\leq
\alpha(s)$ at the top positions, $\alpha(s)\geq0$ and $0\leq c_j(s)\leq z_j(s)$ at all positions, and
that every odd step in $\Delta(s)$ is excluded; with $K_{\delta,1}\geq5$, $K_{\delta,2}\geq4$,
$p\leq2a$ and the size conditions they give $d_j\leq\delta(p,t_0)$. Since
$\gamma_p=\sum_{e\geq0}p^{-2e-1}$,
\[
 \gamma_p\alpha-\sum_{e\geq0}\min\Big(\frac{\alpha}{p^{2e+1}},z_e\Big)
 =\sum_{e\geq0}\max\Big(0,\frac{\alpha}{p^{2e+1}}-z_e\Big)
\]
for real numbers $\alpha$ and $z_e$, and this is nondecreasing in $\alpha$ and nonincreasing in each
$z_e$, as is $\max(\alpha/p^{2e+1}-z,0)$. Comparing term by term gives the first two inequalities.
The third is a computation with the weights of the fills, using $K_{\mathrm{Pr}}\geq1$,
$r(\varepsilon)<k$, $x\geq2$ and the size conditions.

(2) This is a computation from the bounds that the parameter choices put on the exponents $a_s$ and
on $b'$, $w''$, $\eta$ and $c_G$, using $p>2a$.
\end{proof}

\begin{lemma}\label{pd-served-unit-bound}
Let $D$ be a definite-side datum with ranks $k,m,t,n$, size data, gain data, divisibility data,
twisted forms and error data, carrying a lattice instance and local chains of the twists that
satisfy condition~(1)(b) of Definition~\ref{pd-served-selection}, and use the notation of
Definition~\ref{pd-served-lattice}. Suppose that the error data are consistent; that the
conclusion of Lemma~\ref{pd-served-numerics}(1) holds at every word, every place $w$ admissible for
it with $p_w\leq2a$, every representative and every lift, under every tie under which the lattice
instance is tied to the gain data; and that the conclusion of Lemma~\ref{pd-served-numerics}(2)
holds at every place $w$ with $p_w>2a$. Let $\mu$ be a used tuple, fix a tie under which the lattice
instance is tied to the gain data, let $u$ be a word and $w$ a place admissible for $u$ with
$p=p_w\notin Q(\mu)$, and let $K^{(m)}_\mu$ be the composite of the first $m$ maps of the local
chain of $\mu$ at $(u,w)$. Let $V'$ be a set of representatives at $p$ and $F\in M$. Then
$K^{(m)}_\mu(y_{V'}F)$ has valuation at least
\[
 \sum_{j=1}^{m}\Big(\sum_{t_0\in V'}\max\big(g_U-\mathrm{Pr},\ g_0^+\big)+\sum_{t_0\notin V'}g_0^+
 \Big)+\mathrm{low}(u,w),
\]
where $g_U$, $\mathrm{Pr}$ and $g_0^+$ are taken at $(p,j,t_0)$, the second inner sum runs over the
representatives at $p$ outside $V'$, and valuations are those of the levels of $\mu$ at level $m$
(Definition~\ref{pd-local-model}(7)).
\end{lemma}
\leanref{Formalism/SharpRestriction/GrowthBounds/JointsTw.lean}{Formalism.MainNotions.G.Joints.UAt_of}
\leanref{Formalism/SharpRestriction/GrowthBounds/JointsTw.lean}{Formalism.MainNotions.G.Joints.unit_bound_of_sum_over_primes}
\leanref{Formalism/SharpRestriction/GrowthBounds/WeilJacquetShalikaLaw.lean}{Formalism.MainNotions.WeilJacquetShalikaLaw.UAt}
\leanref{Formalism/SharpRestriction/GrowthBounds/WeilJacquetShalikaLaw.lean}{Formalism.MainNotions.WeilJacquetShalikaLaw.unitBoundAtTwistedPlaces}

\begin{proof}
Let $R$ be the set of representatives at $p$ and $t\mapsto v(t)$ the map of the lattice instance.
By condition~(1)(b) of Definition~\ref{pd-served-selection} there are an admitted formal group
system at $p$ on $G(u,w)$, with the fills of the lattice instance and the levels and chain of $\mu$,
and numbers $F_v(j)$ for the orbits $v\notin v(R)$, such that the system has lower bound
$F_v(j+1)$ at $(j,v)$ for these $v$ and all $j<m$, and
\[
 \mathrm{low}(u,w)\leq\sum_{j=1}^{m}\sum_{v\notin v(R)}F_v(j)-\sum_{j<m}d_\Gamma(j).
\]
The tie gives $m_*=m$ and that $t\mapsto v(t)$ is injective on $R$. Let $A$ be the statement
$p\leq2a$. For an orbit $v$ and a lift $j$ define $G(v,j)$, $P(v,j)$ and $G_0(v,j)$ as follows. If
$v=v(t_0)$ with $t_0\in V'$, put $G(v,j)=g_U(p,j,t_0)$ and $P(v,j)=\mathrm{Pr}(p,j,t_0)$ if $A$
holds, and $G(v,j)=\min(g_B(v,j),\mathrm{Pr}_B(v,j))$ and $P(v,j)=\mathrm{Pr}_B(v,j)$ otherwise;
for the other orbits put $G(v,j)=P(v,j)=0$. If $v=v(t_0)$ with $t_0\in R$, put
$G_0(v,j)=g_0^+(p,j,t_0)$ if $A$ holds and $G_0(v,j)=0$ otherwise; if $v\notin v(R)$, put
$G_0(v,j)=F_v(j)$. Let $V=v(V')$ and $S(v)=\sum_{j=1}^m\max(G(v,j)-P(v,j)-G_0(v,j),0)$.

The hypotheses of Theorem~\ref{pd-formal-group-bounds}(2) hold. For $v\in V$, $G\leq g_B$ and
$\mathrm{Pr}_B\leq P$, by Lemma~\ref{pd-served-numerics}(1) if $A$ holds and by definition
otherwise. For $v=v(t_0)$ with $t_0\in R$ and $j<m$, $G_0(v,j+1)\leq\max(g_A(v,j+1),0)$, by
Lemma~\ref{pd-served-numerics}(1) if $A$ holds and trivially otherwise. The system has lower bound
$g_A(v,j+1)$ at $(j,v)$
by Theorem~\ref{pd-formal-group-bounds}(1), and lower bound $0$ by condition~(k) of
Definition~\ref{pd-formal-group-system}; a lower bound $g$ implies every lower bound $g'\leq g$, so
the system has lower bound $G_0(v,j+1)$ at $(j,v)$. For $v\notin v(R)$ this is the choice of the
$F_v$. The condition on $S$ holds by its definition. Finally, $y_{V'}F$ lies in $I(S,V)$: take
$B=V$, listed in the order of the composite $y_{V'}$, so that $n=0$.

So $K^{(m)}_\mu(y_{V'}F)$ has valuation at least
\[
 \sum_{j=1}^{m}\Big(\sum_{v\in V}\max\big(G(v,j)-P(v,j),\ G_0(v,j)\big)+\sum_{v\notin V}G_0(v,j)
 \Big)-\sum_{j<m}d_\Gamma(j).
\]
Compare this with the claimed bound, one lift at a time. For $t_0\in V'$ and $v=v(t_0)$,
$\max(g_U-\mathrm{Pr},g_0^+)\leq\max(G-P,G_0)$: with equality if $A$ holds; otherwise
Lemma~\ref{pd-served-numerics}(2) gives $g_U-\mathrm{Pr}\leq0$ and $g_0\leq0$, so the left side is
$0$, while the right side is at least $G_0(v,j)=0$. For $t_0\in R\smallsetminus V'$,
$g_0^+(p,j,t_0)\leq G_0(v(t_0),j)$ in the same way. The orbits outside $v(R)$ contribute
$\sum_{v\notin v(R)}F_v(j)$, and the inequality from condition~(1)(b) bounds their total, less
$\sum_{j<m}d_\Gamma(j)$, below by $\mathrm{low}(u,w)$. A valuation bound implies every smaller one,
which gives the claim.
\end{proof}

\begin{lemma}\label{pd-served-lattice-bound}
Let $D$ be a definite-side datum with ranks $k,m,t,n$, size data, gain data and a lattice
instance, and use the notation of Definition~\ref{pd-served-lattice}. Let $u$ be a word and $w$ a
place admissible for $u$, with $p=p_w$, and let $K$ be the composite of the first $m$ maps of some
levels and chain over the fills of the lattice instance at $(u,w)$
(Definition~\ref{pd-local-model}(6)). Suppose that for every $V'\subset V(p)$ and every $F\in M$
the function $K(y_{V'}F)$ has valuation at least
\[
 \sum_{j=1}^{m}\Big(\sum_{t_0\in V'}\max\big(g_U-\mathrm{Pr},\ g_0^+\big)+\sum_{t_0\notin V'}g_0^+
 \Big)+\mathrm{low}(u,w),
\]
with $g_U$, $\mathrm{Pr}$ and $g_0^+$ taken at $(p,j,t_0)$ and the second inner sum over the
representatives at $p$ outside $V'$. Then for every $F\in N_P(u,w)$ the function $KF$ has valuation
at least $\sum_{j=1}^m\mathrm{gain}_j(p)+\mathrm{low}(u,w)$.
\end{lemma}
\leanref{Formalism/SharpRestriction/GrowthBounds/AdmissibleIntegers/StatementsOfRemainingBound.lean}{Formalism.MainNotions.WeilJacquetShalikaLaw.Proof.first_in_local_models}
\leanref{Formalism/SharpRestriction/GrowthBounds/TopChoiceBounds.lean}{Formalism.MainNotions.G.TopChoiceBounds.bg_gain_le}

\begin{proof}
Write $B_j(V')$ for the $j$-th term of the outer sum in the hypothesis, and put
$g=\sum_j\mathrm{gain}_j(p)+\mathrm{low}(u,w)$. The functions $F$ for which $KF$ has valuation at
least $g$ form an $\mathcal R$-module, since $K$ is linear and $\varpi^{\lceil ge_{\mathcal R}\rceil}
\mathrm{Int}_m$ is an $\mathcal R$-module. As $N_P(u,w)$ lies in the $\mathcal R$-span of the
functions $y_{V'}(p^nF)$ with $V'\subset V(p)$, $F\in M$ and $n=S(V(p)\smallsetminus V')$, it
suffices to treat these. The Hecke operators are $\mathcal R$-linear, so
$K(y_{V'}(p^nF))=p^nK(y_{V'}F)$. By hypothesis $K(y_{V'}F)$ has valuation at least
$\sum_jB_j(V')+\mathrm{low}(u,w)$, and multiplying by $p^n$ raises a valuation bound by $n$, because
$p$ is a unit times $\varpi^{e_{\mathcal R}}$.

It remains to show $\sum_j\mathrm{gain}_j(p)\leq\sum_jB_j(V')+n$. If $\mathrm{gain}_j(p)=0$ for all
$j$, this holds because every term of $B_j(V')$ is nonnegative and $n\geq0$. Otherwise
$\mathrm{gain}_j(p)=B_j(V(p))$. For $t_0\in V(p)\smallsetminus V'$,
\[
 \sum_j\max\big(g_U-\mathrm{Pr},g_0^+\big)\leq\sum_jg_0^++\sum_j\max\big(0,g_U-\mathrm{Pr}-g_0^+\big)
 \leq\sum_jg_0^++S(p,t_0),
\]
and summing over these $t_0$ gives $\sum_jB_j(V(p))\leq\sum_jB_j(V')+n$.
\end{proof}

\begin{theorem}\label{pd-served-law}
Let $D$ be a definite-side datum with ranks $k,m,t,n$, coefficient field $L$, size data, gain
data, divisibility data, twisted forms and error data, and a lattice instance with a comparison
$\operatorname{cmp}$ and local chains of the twists that satisfy condition~(1)(a) of
Definition~\ref{pd-served-selection}. Use the notation of Definition~\ref{pd-served-lattice}, and
let $\mathcal S_{u,w}$ be the set of vectors of the $w$-local span of $\Lambda_{u,0}$ whose image
lies in $M_{u,\mathrm{ram}}$. Suppose that the error data are consistent, that the local terms
$F_j(w)$ are the ones the size data define, and that the following hold for every word $u$, every
used tuple $\mu$ and every place $w$ with $p=p_w\neq T_u$ and $p\notin Q(\mu)$.
\begin{enumerate}
\item[(U)] Under every tie under which the lattice instance is tied to the gain data, if $w$ is
  admissible for $u$, then for every $V'\subset V(p)$ and every $F\in M$ the function
  $K^{(m)}_\mu(y_{V'}F)$ has valuation at least
  $\sum_{j=1}^m\big(\sum_{t_0\in V'}\max(g_U-\mathrm{Pr},g_0^+)+\sum_{t_0\notin V'}g_0^+\big)
  +\mathrm{low}(u,w)$, with $g_U$, $\mathrm{Pr}$ and $g_0^+$ taken at $(p,j,t_0)$ and the second
  inner sum over the representatives at $p$ outside $V'$.
\item[(K)] If $p\notin\{d,g_0,g_1\}$, then for every level $i<m$ and every $x$ in the lattice of
  forms at level $i$ of $C_u$, and of $C_{u,\mu}$ if $\mu\neq\mu_0$, $K_i(x)$ is the image of $az$
  for some $a\in L$ and some $z$ in the $w$-local span of the lattice of forms at level $i+1$, with
  $a=0$ or $v_w(a)\geq e_ws_i$, where $s_0=0$ and $s_i=-v_p(D_\Gamma(u))$ for $i\geq1$.
\item[(O)] If $p\leq X''$, $p\nmid Ng_0g_1$, and $p$ is a signal prime or of inert class with
  $h_p\neq0$, then for every $x\in\mathcal S_{u,w}$ the iterated lift of $x$ along $C_{u,\mu}$ has
  $w$-valuation at least $\sum_{j=1}^m\bar g_0(p,j)+\mathrm{low}(u,w)$.
\item[(C)] If $p\leq X''$, $p\nmid Ng_0g_1$, and $p$ is of inert class with $h_p=0$, then for every
  $x\in\mathcal S_{u,w}$ the iterated lift of $x$ along $C_{u,\mu}$ has $w$-valuation at least
  $\mathrm{low}(u,w)$.
\item[(D)] If $p=d$ and $p\notin\{g_0,g_1\}$, then for every $x\in\mathcal S_{u,w}$ the iterated
  lift of $x$ along $C_{u,\mu}$ has $w$-valuation at least $\mathrm{low}(u,w)$.
\end{enumerate}
Then for every word $u$, every used tuple $\mu$ and every place $w$ free for $(u,\mu)$
(Definition~\ref{pd-kernel-entries}):
\begin{enumerate}
\item[(i)] for every $x\in\mathcal S_{u,w}$, the iterated lift of $x$ along $C_{u,\mu}$ has
  $w$-valuation at least $\mathrm{low}(u,w)$ (Definition~\ref{pd-lift-valuation});
\item[(ii)] if the lattice instance is tied to the gain data by a tie of the local models, then for
  every $x$ that lies in the served local lattice at $w$, the iterated lift of $x$ along
  $C_{u,\mu}$ has $w$-valuation at least $\sum_{j=1}^m\mathrm{gain}^{\mathrm s}_j(p_w)+
  \mathrm{low}(u,w)$.
\end{enumerate}
\end{theorem}
\leanref{Formalism/SharpRestriction/GrowthBounds/AdmissibleIntegers.lean}{Formalism.MainNotions.WeilJacquetShalikaLaw.Proof.stmt_proved}
\leanref{Formalism/SharpRestriction/GrowthBounds/AdmissibleIntegers/StatementsOfRemainingBound.lean}{Formalism.MainNotions.WeilJacquetShalikaLaw.Proof.statements_of_remaining_bound}
\leanref{Formalism/SharpRestriction/GrowthBounds/WeilJacquetShalikaLaw.lean}{Formalism.MainNotions.WeilJacquetShalikaLaw.serv_of}
\leanref{Formalism/SharpRestriction/GrowthBounds/WeilJacquetShalikaLaw.lean}{Formalism.MainNotions.WeilJacquetShalikaLaw.Conclusion1}
\leanref{Formalism/SharpRestriction/GrowthBounds/WeilJacquetShalikaLaw.lean}{Formalism.MainNotions.WeilJacquetShalikaLaw.ColumnValuationAwayFromTypePrimes}
\leanref{Formalism/SharpRestriction/GrowthBounds/WeilJacquetShalikaLaw.lean}{Formalism.MainNotions.WeilJacquetShalikaLaw.IntegralKernelValuation}
\leanref{Formalism/SharpRestriction/GrowthBounds/WeilJacquetShalikaLaw.lean}{Formalism.MainNotions.WeilJacquetShalikaLaw.OrbitsWithoutOperator}
\leanref{Formalism/SharpRestriction/GrowthBounds/WeilJacquetShalikaLaw.lean}{Formalism.MainNotions.WeilJacquetShalikaLaw.MediumInertPrimeLengthBound}
\leanref{Formalism/SharpRestriction/GrowthBounds/WeilJacquetShalikaLaw.lean}{Formalism.MainNotions.WeilJacquetShalikaLaw.LossAtDiscriminant}

\begin{proof}
Write $p=p_w$. For a level $i$ and $s\in\mathbf R$, say that $y\in V_{u,i}$ has $w$-valuation at
least $s$ if $y=az$ with $z$ in the $w$-local span of $\Lambda_{u,i}$ and $a=0$ or $v_w(a)\geq
e_ws$, and in the same way for the twisted lattices. A sum of two vectors of $w$-valuation at least
$s$ has $w$-valuation at least $s$: if $y_1=a_1z_1$ and $y_2=a_2z_2$ with $a_1\neq0$ and
$v_w(a_1)\leq v_w(a_2)$, then $y_1+y_2=a_1(z_1+(a_2/a_1)z_2)$. If $y$ has $w$-valuation at least $s$
and $b=0$ or $v_w(b)\geq e_ws'$, then $by$ has $w$-valuation at least $s+s'$.

\emph{Step 1: places away from $d$ and the level primes.} Let $u$ be a word, $\mu$ a used tuple and
$w$ a place with $p\notin\{T_u,d,g_0,g_1\}$ and $p\notin Q(\mu)$, and let $x\in\mathcal S_{u,w}$.
Take $\mu=\mu_0$ first. By (K), $K_i(z)$ has $w$-valuation at least $s_i$ for $z\in\Lambda_{u,i}$.
By the two properties above and the linearity of $K_i$, the same holds for $z$ in the $w$-local
span of $\Lambda_{u,i}$, and if $y$ has $w$-valuation at least $s$, then $K_i(y)$ has $w$-valuation
at least $s+s_i$. By induction on $i$, $K_{i-1}\cdots K_0(x)$ has $w$-valuation at least
$s_0+\cdots+s_{i-1}$; for $i=m$ this says that the iterated lift of $x$ has $w$-valuation at least
$-(m-1)v_p(D_\Gamma(u))$. For $\mu\neq\mu_0$, the twisted forms carry the lattice of forms of $C_u$
at level $0$ onto that of $C_{u,\mu}$ by an $L$-linear isomorphism compatible with the images, so
$x$ gives a vector of the $w$-local span of the latter, and the same induction along $C_{u,\mu}$
gives the same bound.

\emph{Step 2: (i).} Let $w$ be free for $(u,\mu)$ and $x\in\mathcal S_{u,w}$. If $p=d$, (D) gives
(i). Suppose that $p\neq d$, that $p\leq X''$, and that $p$ is a signal prime or of inert class.
Such a prime is inert, so it does not divide $N=Q_Sqd$, and $p\notin\{g_0,g_1\}$ because $w$ is
free; so $p\nmid Ng_0g_1$. If $p$ is of inert class with $h_p=0$, (C) gives (i). Otherwise (O) gives
(i), since $\bar g_0(p,j)\geq0$. In the remaining case $p\notin\{T_u,g_0,g_1,d\}$, and either
$p>X''$ or $p$ is neither a signal prime nor of inert class. At such a place the size data define
$F_j(w)=0$, and the consistent error data define $E_{u,j}(w)=v_p(D_\Gamma(u))$, for every lift $j$.
So $\mathrm{low}(u,w)=-(m-1)v_p(D_\Gamma(u))$, and Step 1 gives (i).

\emph{Step 3: (ii).} Fix a tie under which the lattice instance is tied to the gain data, let $w$ be
free for $(u,\mu)$, and let $x$ lie in the served local lattice at $w$. Since $\mathrm{gain}_j(p)
\geq0$ and $\mathrm{gain}^{\mathrm s}_j(p)$ is $0$ or $\mathrm{gain}_j(p)$, we have
$\mathrm{gain}^{\mathrm s}_j(p)\leq\mathrm{gain}_j(p)$.

Suppose first that $w$ is admissible for $u$ and $\operatorname{cmp}_{u,w}(x)\in N_P(u,w)$. By (U)
and Lemma~\ref{pd-served-lattice-bound}, applied with $K=K^{(m)}_\mu$, the function
$K^{(m)}_\mu(\operatorname{cmp}_{u,w}x)$ has valuation at least $\sum_j\mathrm{gain}_j(p)+
\mathrm{low}(u,w)$. By condition~(1)(a) of Definition~\ref{pd-served-selection}, the iterated lift
of $x$ along $C_{u,\mu}$ has $w$-valuation at least this number, which is at least
$\sum_j\mathrm{gain}^{\mathrm s}_j(p)+\mathrm{low}(u,w)$.

Suppose next that $V(p)=\emptyset$ and $x\in\mathcal S_{u,w}$. By (i), the iterated lift of $x$ has
$w$-valuation at least $\mathrm{low}(u,w)$. This is (ii) when every $\mathrm{gain}^{\mathrm s}_j(p)$
vanishes, which is the case unless $p\leq X''$, $p\nmid Ng_0g_1$, and $p$ is a signal prime or of
inert class with $h_p\neq0$. In that case, as $V(p)=\emptyset$,
$\mathrm{gain}_j(p)=\sum_{t_0}g_0^+(p,j,t_0)$ over all representatives, so
$\mathrm{gain}^{\mathrm s}_j(p)=\bar g_0(p,j)$, and (O) gives (ii).
\end{proof}

\begin{theorem}\label{pd-multiplicity-from-integrality}
Let $\mathcal A$ be a class of families with a twist selection, and let
$\mathcal D_3=\operatorname{Des}(\mathcal A_N,3)$. Suppose that each of the following holds
eventually along $\mathcal D_3$.
\begin{enumerate}
\item \emph{Point integrality} at all large scales: for every word $u$, the point classes map
  injectively to $\mathrm{cls}_m$ with bijective maps of stabilizers; $\operatorname{ev}_u$ maps
  the span of the image of $\Lambda_{u,m}$ into the span of $j^{\mathrm{pt}}_u(\Lambda^{\mathrm{pt}}_u)$;
  the evaluation of $u$ is defined over $L$ and integral (Definition~\ref{pd-period-data}); and the
  evaluation multiplies the metrics at every infinite place, and the Petersson norm, by at most
  $\prod_\tau T_u^{-\deg_\tau}\beta_\tau^2$, the product over $\tau\in\mathcal T$, with
  $\deg_\tau$ and $\beta_\tau\geq0$ the degree and the bound of the forms of $u$ at $\tau$.
\item \emph{Twisted point integrality} at all large scales: for every word $u$ and every used
  tuple $\mu\neq\mu_0$, the twisted evaluation is defined over $L$ and integral.
\item Period integrality at all large scales (Definition~\ref{pd-period-integrality}).
\item Twisted period integrality at all large scales, for the divisibility data and twisted forms
  of the family (Definition~\ref{pd-twisted-period-integrality}).
\item \emph{Sharp column bounds at the auxiliary prime} at all large scales: for every word $u$
  and every place $w$ over $T_u$, the iterated lift along $C_u$, and along $C_{u,\mu}$ for every
  used tuple $\mu\neq\mu_0$ with $p_w\notin Q(\mu)$, of every vector of the local lattice
  $N_\lambda(u,w)$ has $w$-valuation at least $-m\,E_T(u,w)$
  (Definition~\ref{pd-lift-valuation}), where $N_\lambda(u,w)$ is the local lattice and $E_T(u,w)$
  the exponent at $T_u$ (Definition~\ref{pd-hecke-at-T}).
\item \emph{Bounds at the level primes}: there is $C_1$ such that, at all large scales, for every
  place $w$ over $g_i$ ($i=0,1$), every word $u$ with $p_w\neq T_u$, every used tuple $\mu$, every
  $\varphi\in M^{\mathrm{ram}}_u$ and every $F\in\mathcal A_\mu$ of the form
  $\sum_{s\in S_\mu}\iota(c_s)s$ with $c_s\in\mathcal O$, the period
  $\operatorname{per}_{u,\mu}(j_N\varphi,F)$ is $\iota(y)$ with $y=0$ or
  $v_w(y)/e_w\geq-\sum_j\nu_{u,j}(w)-C_1c_{\mathrm{err}}$, and
  $\sum_j(F_j(w)-E_{u,j}(w))\leq-\sum_j\nu_{u,j}(w)-C_1c_{\mathrm{err}}$. Here
  $\nu_{u,j}(w)=e_w^{-1}\sum_\tau n_{u,j,\tau}\,v_w\big(\tau(\overline{\pi_i})\big)$, where
  $\pi_i\in\mathcal O_E$ with $\pi_i\overline{\pi_i}=g_i$ is fixed by the error data of the
  family and $n_{u,j,\tau}$ is the sum of the row exponents of the $j$-th lift of $u$ at $\tau$;
  and $c_{\mathrm{err}}=f\big(mN(m)(N(m)+t)+m^3\eta\big)$.
\item The twist selection is served along $\mathcal D_3$ (Definition~\ref{pd-served-selection}).
\end{enumerate}
Then, eventually along $\mathcal D_3$, the multiplicity bound
(Definition~\ref{pd-multiplicity-bound}) holds at all large scales.
\end{theorem}
\leanref{Development/AutomorphicArithmetic/ThetaCorrespondence/ThetaPairLifts/ThirdPlaceSeries/ClosedFamilyAdmissibleTwistLink.lean}{GlobalTheory.Top.closed_family_admissible_twist_link_alt}
\leanref{Formalism/LocalToGlobal/GroupLattices/AdmissibleAtGeneratorPlaces.lean}{Formalism.L2.G.Adm2InU.admissible_family_of_inclusion}
\leanref{Formalism/LocalToGlobal/GroupLattices/AdmissibleAtGeneratorPlaces.lean}{Formalism.L2.G.Adm2InU.admissible_level_two}

\begin{proof}
Two consequences of the design come first.

\emph{Step 1: $T_u\nmid D_\Gamma(u)$.} Let $\mathcal F\in\mathcal D_3$, and take $x$ large enough
for conditions (2) and (4) of Definition~\ref{asm-designed}. Let $D$ be a member at $x$ and $u$ a
word, and suppose that the prime $T_u$ divides $D_\Gamma(u)$. By
Definition~\ref{asm-designed}(4), $T_u\leq2fN(m)+1$. The inequalities of a rank datum
(Definition~\ref{th-ranks}) give $f\leq k\leq t$, $m\leq t$ and $t\geq7$, so
\[
 2fN(m)+1\leq2t\big(t+t(t+1)\big)+1<t^5.
\]
The size conditions of the family (Definition~\ref{pd-family}(4)) include
$(Q_Sqd)^{36}t^{150}\leq b^{1/2}$, so $t^5\leq t^{300}\leq b$. Since $b\geq1$ and $m\geq2$,
$b\leq b^{m/2}\leq T_u$ by Definition~\ref{asm-designed}(2). Together these give $T_u<T_u$, a
contradiction. So $T_u\nmid D_\Gamma(u)$ at all large scales.

\emph{Step 2: the served column bounds.} At all large scales every member satisfies the following
\emph{served column bounds}. Let $u$ be a word, $\mu$ a used tuple and $w$ a place free for
$(u,\mu)$ (Definition~\ref{pd-kernel-entries}), and write $p=p_w$.
\begin{enumerate}
\item[(i)] For every vector $x$ of the $w$-local span of the lattice of forms at level $0$ whose
  image lies in $M_{u,\mathrm{ram}}$, the iterated lift of $x$ along $C_{u,\mu}$ has
  $w$-valuation at least $\sum_j\big(F_j(w)-E_{u,j}(w)\big)+v_p(D_\Gamma(u))$
  (Definition~\ref{pd-lift-valuation}).
\item[(ii)] For every vector $x$ that lies in the served local lattice at $w$
  (Definition~\ref{pd-served-lattice}), the iterated lift of $x$ along
  $C_{u,\mu}$ has $w$-valuation at least
  $\sum_j\big(\mathrm{gain}^{\mathrm s}_j(p)+F_j(w)-E_{u,j}(w)\big)+v_p(D_\Gamma(u))$, where
  $\mathrm{gain}^{\mathrm s}_j(p)$ is the served gain (Definition~\ref{pd-served-lattice}). It is
  $0$ unless $p$ is a signal prime or of inert class, $p\leq X''$, $p$ does not divide $g_0g_1$ or the conductor of the
  host, and $h_p$ is nonzero of even order if $p$ is of inert class; in that case it is a
  sum over the representatives $t_0$ at $p$ of terms at least $\max(g_0(p,j,t_0),0)$
  (Definition~\ref{pd-served-selection}).
\end{enumerate}
They are proved as follows, for a family of $\mathcal D_3$ large enough, a large scale and a member
$D$. The choices of the family give consistent error data, the local terms $F_j(w)$ that the size
data define, and a lattice instance tied to the gain data (Definition~\ref{pd-family}(4)).
Condition~(1) of Definition~\ref{pd-served-selection} gives the comparison of the twisted lattices
of forms with the local lattices of the gain data, and admitted formal group systems at the places
of the gain data. In a family the constants of $K_G$ are nonnegative, with $K_{\delta,1}=5$,
$K_{\delta,2}=4$ and $K_{\mathrm{Pr}}=1$, and at large scales $r(\varepsilon)<k$, $x\geq2$ and the
size conditions hold; so Lemma~\ref{pd-served-numerics} applies, its part~(2) through the parameter
choices of the family. With Theorem~\ref{pd-formal-group-bounds}, Lemma~\ref{pd-served-unit-bound}
then gives the unit bound~(U) of Theorem~\ref{pd-served-law}. Condition~(2), with
Lemma~\ref{pd-kernel-form}, gives the kernel bound of that lemma at every level, with no loss at
level $0$ because the classes of level $0$ have trivial stabilizers. Since the host has split bound
at least $3$, every signal prime and every prime of inert class with $h_p$ of even order is at least
$5$ and has even orbit length $\ell_p\geq2$, and every prime of inert class is at least $5$; and the
level primes differ from $d$. With conditions~(3) to~(7) and
Proposition~\ref{pd-orbit-letter-bounds}, this bounds below the $\nu_{u,i}$-valuation of every theta
kernel at every free place: by the sum over the orbit representatives of the local gains and the
local terms at signal primes and at primes of inert class with even orbits, by
$|\nu|/(p^2-1)$ less an error at split orbits, by the shape of condition~(6) at $d$, and by $0$
elsewhere. Lemma~\ref{pd-kernel-form} at each level, summed over the $m$ levels, turns these bounds
into lower bounds for the $w$-valuation of the iterated lift. The kernel bound at each level is
hypothesis~(K) of Theorem~\ref{pd-served-law}, and these lower bounds are its hypotheses (O), (C)
and~(D); with the unit bound and the comparison of condition~(1)(a), the theorem gives (i) and
(ii). The size condition
on the stabilizer denominators (Definition~\ref{asm-designed}(3)) gives the size conditions at the
level primes that accompany (i) and (ii).

The parameter choices of the family give $\mathrm{gain}^\flat_j(w)\leq\mathrm{gain}^{\mathrm s}_j(p_w)$,
and place $N^\circ_u$ locally at each place as used below.
% Lean: Step 2 is Formalism.MainNotions.WeilJacquetShalikaLaw.ServF (Serv: TypePrimeSizeConds,
% Conclusion1 for (ii), ColumnValuationAwayFromTypePrimes for (i)), proved by
% closed_servF_of_servT_of_dGamSize = Pk.sum_over_primes_service. The unit bound is
% G.Joints.unit_bound_of_sum_over_primes (Lemma pd-served-unit-bound) with
% FormalGroupSystem.Pf.stmt_proved (Theorem pd-formal-group-bounds), num_re of G.num_closed and
% TopChoiceBounds.hLarge_proved (Lemma pd-served-numerics); (K), (O), (C) and (D) are
% J.four_orbit_conditions_of_proof after OrbitPlaces.lift_bounds_of_essential,
% DIV.Lemmas.OrbitEntryBounds.orbit_entry_bounds_hold and DIV.Inh.orbit_admissible_of_split_bound;
% the law is WeilJacquetShalikaLaw.Proof.stmt_proved with serv_of (Theorem pd-served-law); and
% TypePrimeSizeConds is admissible_size_family. The local orbit bounds are Proposition
% pd-orbit-letter-bounds and the kernel bound is Lemma pd-kernel-form.

Fix a family of $\mathcal D_3$ large enough for all the conditions, a large scale $x$ and a member
$i$, and write $D$ and $L$ for its datum and size data. Since $v_w(y)$ is an integer, it suffices to
show $v_w(y)\geq e_ws(u,w)$. Fix $u$, $w$, $\varphi\in N^\circ_u$ and an integral top form $F$ for
$\mu_w$, and recall that $p_w\notin Q(\mu_w)$. There are three cases.

\emph{Places over $T_u$.} By the parameter choices of the family, $\nu_u(N^\circ_u)$ lies locally
at $w$ in $N_\lambda(u,w)$, and $s(u,w)=-mE_T(u,w)$: at $T_u$ the gains vanish and the local terms
are $F_j(w)=0$ and $E_{u,j}(w)=E_T(u,w)$. By condition~(5), the iterated lift of $\nu_u(\varphi)$
along $C_{u,\mu_w}$ has $w$-valuation at least $-mE_T(u,w)=s(u,w)$. Since $T_u\nmid D_\Gamma(u)$,
Lemma~\ref{pd-period-of-divisible-lift}, whose hypotheses on evaluations are conditions~(1)
and~(2), gives $v_w(y)\geq e_ws(u,w)$.

\emph{Places over $g_0$ or $g_1$, not over $T_u$.} The gains vanish at the level primes $g_0,g_1$,
so $s(u,w)=\sum_j(F_j(w)-E_{u,j}(w))$. Locally at $w$, $N^\circ_u$ lies in $M^{\mathrm{ram}}_u$:
there is $\sigma\in\mathcal O_L$ with $v_w(\sigma)=0$ and $\sigma\varphi\in M^{\mathrm{ram}}_u$. By
condition~(6), the period of $j_N(\sigma\varphi)=\iota(\sigma)j_N\varphi$ is $\iota(y')$ with
$y'=0$ or $v_w(y')\geq e_w\big(-\sum_j\nu_{u,j}(w)-C_1c_{\mathrm{err}}\big)\geq e_ws(u,w)$. The
period is linear in its first argument, so $y=\sigma^{-1}y'$, and $v_w(y)=v_w(y')$.

\emph{Other places.} Locally at $w$, $\nu_u(\varphi)$ lies in the served local lattice when $p_w$
has a served orbit, and in the $w$-local span of $\Lambda_{u,0}\cap\nu_u(V(N_u))$ otherwise. In
both cases the served column bound gives the iterated lift of $\nu_u(\varphi)$ $w$-valuation at
least
$\sum_j(\mathrm{gain}^{\mathrm s}_j(p_w)+F_j(w)-E_{u,j}(w))+v_{p_w}(D_\Gamma(u))\geq
s(u,w)+v_{p_w}(D_\Gamma(u))$. Lemma~\ref{pd-period-of-divisible-lift} gives
$v_w(y)\geq e_ws(u,w)$.
\end{proof}

\section{Integrality of periods along the families}\label{section-pd-period-integrality}

\begin{definition}\label{pd-weight-lattices}
Let $K$ be a number field, $\mathcal T$ a finite set and $n\geq1$, and let $X$ and $X_\tau$
($\tau\in\mathcal T$) be $n\times n$ matrices of variables. For integers
$\nu_1\geq\cdots\geq\nu_n\geq0$ let $h_\nu=\prod_{s=1}^n\delta_s^{\nu_s-\nu_{s+1}}$, where
$\delta_s$ is the leading principal $s\times s$ minor of $X$ and $\nu_{n+1}=0$.
\begin{enumerate}
\item For a commutative ring $R$, the \emph{dual Weyl module} $\nabla_R(\nu)\subset R[X]$ consists of
  the polynomials $f$ of degree $\nu_s$ in the entries of the $s$-th column of $X$, for each $s$,
  such that $f(X(1+yE_{lj}))=f(X)$ in $R[X,y]$ for all $l<j$, where $E_{lj}$ is the matrix unit.
  The \emph{Weyl module} $\Delta_R(\nu)\subset R[X]$ is the $R$-span of the polynomials obtained from
  $h_\nu$ by repeatedly applying the operators $E^{(k)}_{ab}$ ($1\leq a,b\leq n$, $k\geq0$), where
  $E^{(k)}_{ab}f$ is the coefficient of $y^k$ in $f$ with the $b$-th row of $X$ replaced by its sum
  with $y$ times the $a$-th row.
\item The \emph{Fock pairing} of two polynomials in the entries of the $X_\tau$ is
  $[f,g]=\sum_\kappa\kappa!\,f_\kappa g_\kappa$, where $f_\kappa$ and $g_\kappa$ are the
  coefficients of the monomial with exponent vector $\kappa$, and $\kappa!$ is the product of the
  factorials of the entries of $\kappa$.
\item Let $\lambda=(\lambda_\tau)_{\tau\in\mathcal T}$ be dominant weights,
  $\lambda_{\tau,1}\geq\cdots\geq\lambda_{\tau,n}$ in $\mathbf Z$, and put
  $\nu_{\tau,i}=\lambda_{\tau,i}-\lambda_{\tau,n}$. The space
  $V_\lambda(K)\subset K[X_\tau:\tau\in\mathcal T]$ is the $K$-span of the products
  $\prod_\tau f_\tau(X_\tau)$ with $f_\tau\in\nabla_K(\nu_\tau)$. The \emph{standard lattices}
  $\nabla_\lambda$ and $\Delta_\lambda$ are the $\mathcal O_K$-spans of such products with
  $f_\tau\in\nabla_{\mathcal O_K}(\nu_\tau)$, respectively $f_\tau\in\Delta_{\mathcal O_K}(\nu_\tau)$.
  The \emph{pairing} on $V_\lambda(K)$ is
  $\langle f,g\rangle_\lambda=[f,g]\prod_\tau[h_{\nu_\tau},h_{\nu_\tau}]^{-1}$.
\item The group $GL_n(K)^{\mathcal T}$ acts on $V_\lambda(K)$ by
  $(\rho_\lambda(g)f)\big((X_\tau)_\tau\big)=\prod_\tau\det(g_\tau)^{\lambda_{\tau,n}}
  f\big(({}^tg_\tau X_\tau)_\tau\big)$.
\item Let $M=(M_\tau)_{\tau\in\mathcal T}$ be $\mathcal O_K$-lattices in $K^n$. Then
  $\nabla_\lambda(M)$ is the set of $x\in V_\lambda(K)$ such that, for every finite place $w$ of $K$
  and all $B_\tau\in GL_n(K)$ whose columns form a basis of the localization of $M_\tau$ at $w$,
  $\rho_\lambda(B)^{-1}x$ lies in the localization of $\nabla_\lambda$ at $w$. The lattice
  $\Delta_\lambda(M)$ is defined in the same way from $\Delta_\lambda$.
\item Let $E$ be a number field with embeddings $\tau\colon E\to K$ for $\tau\in\mathcal T$, and let
  $L'\subset E^n$ be an $\mathcal O_E$-lattice. Let $M_\tau$ be the $\mathcal O_K$-span of
  $\tau(L')$ and $M^\vee_\tau=\{\phi\in K^n:\phi\cdot v\in\mathcal O_K\text{ for all }v\in M_\tau\}$.
  Then $\nabla_\lambda(L')=\nabla_\lambda(M)$, $\Delta_\lambda(L')=\Delta_\lambda(M)$,
  $\nabla_\lambda(L'^\vee)=\nabla_\lambda(M^\vee)$ and $\Delta_\lambda(L'^\vee)=\Delta_\lambda(M^\vee)$.
\end{enumerate}
\end{definition}
\leanref{Formalism/GradedWeylDivisibility/StatementA.lean}{Formalism.GradedWeylDivisibility.dualWeyl}
\leanref{Formalism/GradedWeylDivisibility/StatementA.lean}{Formalism.GradedWeylDivisibility.weyl}
\leanref{Formalism/LocalToGlobal/Divisibility/CoefMod.lean}{Formalism.L2.DIV.VV}
\leanref{Formalism/LocalToGlobal/Columns/CoefModB.lean}{Formalism.L2.DIV.pairB}
\leanref{Formalism/LocalToGlobal/Divisibility/CoefMod.lean}{Formalism.L2.DIV.nabP}
\leanref{Formalism/LocalToGlobal/Divisibility/CoefMod.lean}{Formalism.L2.DIV.nabAt}
\leanref{Formalism/LocalToGlobal/Divisibility/CoefMod.lean}{Formalism.L2.DIV.delAtr}

\begin{proposition}\label{pd-weight-duality}
Let $E$ and $K$ be number fields, $\mathcal T$ a finite set with embeddings $\tau\colon E\to K$ for
$\tau\in\mathcal T$, $n\geq1$, $\lambda=(\lambda_\tau)$ dominant weights, and $L'\subset E^n$ an
$\mathcal O_E$-lattice that spans $E^n$ (Definition~\ref{pd-weight-lattices}). Let
$x\in V_\lambda(K)$.
\begin{enumerate}
\item $x\in\Delta_\lambda(L'^\vee)$ if and only if $\langle x,y\rangle_\lambda\in\mathcal O_K$ for
  every $y\in\nabla_\lambda(L')$.
\item $x\in\nabla_\lambda(L'^\vee)$ if and only if $\langle x,y\rangle_\lambda\in\mathcal O_K$ for
  every $y\in\Delta_\lambda(L')$.
\end{enumerate}
\end{proposition}
\leanref{Formalism/LocalToGlobal/Columns/CoefModB.lean}{Formalism.L2.DIV.PairBDual}
\leanref{Development/AbelianVarietyTheory/Proofs/GlobalInterpolation/IsClassSumPairingFormsLat.lean}{GlobalTheory.PeriodClassSum.pairBDual}

\begin{proof}
Let $M_\tau$ be the $\mathcal O_K$-span of $\tau(L')$; it is a lattice in $K^n$ because $L'$ spans
$E^n$. Membership in the four lattices and in $\mathcal O_K$ is decided at the finite places of $K$
one at a time. For the Fock pairing, $[f({}^tgX),h(X)]=[f(X),h(gX)]$ for $g\in GL_n(K)$, and the
determinant factors cancel, so
\[
 \langle\rho_\lambda(g)x,\rho_\lambda({}^tg^{-1})y\rangle_\lambda=\langle x,y\rangle_\lambda
 \qquad(g\in GL_n(K)^{\mathcal T}).
\]
Fix a finite place $w$ and matrices $B_\tau$ whose columns form a basis of $M_\tau$ at $w$. The
columns of ${}^tB_\tau^{-1}$ form a basis of $M^\vee_\tau$ at $w$. Hence $y\mapsto\rho_\lambda(B)^{-1}y$
and $x\mapsto\rho_\lambda({}^tB)x$ carry the lattices of the statement at $w$ to the standard
lattices at $w$, and they preserve the pairing. This reduces both statements to the standard
lattices. For a single $\nu$ there are polynomials $b_i$ and $d_i$ with integer coefficients,
indexed by the semistandard tableaux $i$ of shape $\nu$, with $[b_i,d_j]=\delta_{ij}[h_\nu,h_\nu]$,
such that the $b_i$ span $\nabla_{\mathcal O_K}(\nu)$ and the $d_i$ span $\Delta_{\mathcal O_K}(\nu)$
over $\mathcal O_K$. The products of such families over $\tau\in\mathcal T$ are $\mathcal O_K$-bases
of $\nabla_\lambda$ and $\Delta_\lambda$ that are dual for $\langle\ ,\ \rangle_\lambda$, and both
statements follow.
\end{proof}

\begin{definition}\label{pd-point-data}
Let $E$ and $L$ be number fields, $\mathcal T$ a finite set with embeddings $\tau\colon E\to L$ for
$\tau\in\mathcal T$, $\iota\colon L\to\mathbf C$ an embedding, $n\geq1$, and
$\sigma=(\sigma_\tau)_{\tau\in\mathcal T}$ dominant weights (Definition~\ref{pd-weight-lattices}).
Write $V_\sigma(\mathbf C)$ for the complex span of $V_\sigma(L)$, on which
$\langle\ ,\ \rangle_\sigma$ extends $\mathbf C$-bilinearly, and $\iota$ also for its application to
coefficients. \emph{Point data} for $\sigma$ consist of:
\begin{enumerate}
\item a finite set of \emph{point classes} $c$, each with an $\mathcal O_E$-lattice $L_c\subset E^n$
  that spans $E^n$ and a finite group $\Gamma_c$, its \emph{stabilizer}, acting $L$-linearly on
  $V_\sigma(L)$;
\item a complex vector space of \emph{top forms} and, for each point class $c$, a linear map
  $F\mapsto R_c(F)$ from the top forms to $V_\sigma(\mathbf C)$, the \emph{row} of $F$ at $c$.
\end{enumerate}
The \emph{point space} consists of the families $z=(z_c)$ with $z_c\in V_\sigma(\mathbf C)$, one for
each point class. The \emph{point lattice} $\Lambda^{\mathrm{pt}}$ consists of the families
$x=(x_c)$ with $x_c\in V_\sigma(L)$ fixed by $\Gamma_c$ and $x_c\in\Delta_\sigma(L_c^\vee)$; we write
$\iota(x)=(\iota(x_c))$. The \emph{period} of $z$ against a top form $F$ is
\[
 P(z,F)=\sum_c\frac1{|\Gamma_c|}\,\langle z_c,R_c(F)\rangle_\sigma .
\]
The row of $F$ at $c$ is \emph{integral} if $R_c(F)=\iota(y)$ for some $y\in\nabla_\sigma(L_c)$.
\end{definition}
\leanref{Development/AbelianVarietyTheory/Definitions/AutomorphicMembership/PredMemGamRanks.lean}{GlobalTheory.Mem.Col2.GStep}
\leanref{Development/AbelianVarietyTheory/Definitions/AutomorphicMembership/PredMemGamRanks.lean}{GlobalTheory.Mem.Col2.GStep.hlat}
\leanref{Development/AbelianVarietyTheory/Definitions/AutomorphicMembership/PredMemGamRanks.lean}{GlobalTheory.Mem.Col2.PerIn}
\leanref{Development/AbelianVarietyTheory/Definitions/AutomorphicMembership/PredMemGamRanks.lean}{GlobalTheory.Mem.Col2.PerIn.per}

\begin{lemma}\label{pd-class-term-integral}
Let point data for dominant weights $\sigma$ be given (Definition~\ref{pd-point-data}), let $F$ be a
top form and $c$ a point class at which the row of $F$ is integral. For $x\in\Lambda^{\mathrm{pt}}$
let $\iota(x)_{[c]}$ be the vector of the point space equal to $\iota(x_c)$ at $c$ and to $0$ at the
other point classes. Then $|\Gamma_c|\,P(\iota(x)_{[c]},F)=\iota(a)$ for some $a\in\mathcal O_L$.
\end{lemma}
\leanref{Development/AbelianVarietyTheory/Proofs/GlobalInterpolation/ClassPairingIntegralFormRowIsIntegralStatement.lean}{GlobalTheory.T.classPairing_integral_form_row_isIntegral}

\begin{proof}
Only the class $c$ contributes to $P(\iota(x)_{[c]},F)$, so
$|\Gamma_c|\,P(\iota(x)_{[c]},F)=\langle\iota(x_c),R_c(F)\rangle_\sigma$. Write $R_c(F)=\iota(y)$
with $y\in\nabla_\sigma(L_c)$. The Fock pairing has rational coefficients and $[h_\nu,h_\nu]$ is a
positive integer, so $\langle\iota(x_c),\iota(y)\rangle_\sigma=\iota(\langle x_c,y\rangle_\sigma)$.
Since $x_c\in\Delta_\sigma(L_c^\vee)$, Proposition~\ref{pd-weight-duality}(1) gives
$a=\langle x_c,y\rangle_\sigma\in\mathcal O_L$.
\end{proof}

\begin{definition}\label{pd-true-columns}
Let a base for definite-side data be given: a host with CM field $E$, ranks $k,m,t,n$ with
$N(j)=k+j(t+1)$, words $u$ with parameters $\lambda_u$, elements $\pi_u\in\mathcal O_E$ with
$\pi_u\bar\pi_u=T_u$ as in Definition~\ref{th-definite-datum}(4), a coefficient field $L$ with
embeddings $\tau\colon E\to L$ for $\tau\in\mathcal T$ and $\iota\colon L\to\mathbf C$, and
\emph{row weights} $\sigma=(\sigma_\tau)$, dominant weights for $GL_{N(m)-1}$.
\begin{enumerate}
\item \emph{True columns} over the base consist of the following data, subject to the conditions
  stated with them. For a Hecke character $\chi$ of $E$ and $n\colon\mathcal T\to\mathbf Z$, $\chi$
  has \emph{infinity type} $n$ if its archimedean component is
  $k\mapsto\prod_\tau(\tau(k)/|\tau(k)|)^{n(\tau)}$ on $E^\times$.
  \begin{enumerate}
  \item[(a)] \emph{Characters of the host.} Primes $\bar\pi_0,\bar\pi_1$ of the subfield $K$ of the
    host over $g_0,g_1$ and injective characters $\omega_k$ of $(\mathbf Z/g_k)^\times$; unitary Hecke
    characters $\mu_{0,j}$ of $E$ ($1\leq j\leq m$), trivial on the ideles of $F$, tame, unramified
    outside the places over $g_0g_1$, of infinity type $-2e_j(\tau)$, with
    $\mu_{0,j,w}(a)=\omega_k(a)^{-2j}$ at the places $w$ of $E$ over $\bar\pi_k$ for every integer $a$
    prime to $g_k$; a unitary everywhere unramified Hecke character $\eta$ of infinity type $1$ whose
    restriction to $F$ is the quadratic character of $E/F$; a unitary everywhere unramified Hecke
    character $\chi_\Psi$, trivial on $F$, of infinity type $2$; and Hecke characters $\delta,\delta'$
    of finite order, trivial on $F$, tame, unramified outside $p_{de}$ (respectively $p'_{de}$), whose
    restriction to the units at each place over $p_{de}$ (respectively $p'_{de}$) has order $15$.
  \item[(b)] \emph{Characters at the level primes.} Let $S_G$ be the set of places of $E$ over $g_0$
    or $g_1$; they are the places over $\bar\pi_0,\bar\pi_1$. For each $w\in S_G$, a nontrivial
    additive character $\psi_{0,w}$ of $\mathcal O_E/w$ and characters $\chi_{j,w}$ of
    $(\mathcal O_E/w)^\times$, $1\leq j\leq m$, with values in $L^\times$; and for every finite place
    $v$ of $E$, elements $\mu_j(v)$ and $\eta(v)$ of $L^\times$. They satisfy: $j\mapsto\chi_{j,w}$ is
    injective; for $w,w'\in S_G$ over the same $g_k$, $\psi_{0,w}(a)=\psi_{0,w'}(a)$ for $a\in\mathbf Z$
    and $\chi_{j,w}(a)=\chi_{j,w'}(a)$ for $a\in\mathbf Z$ prime to $g_k$; for each $k$ there is an
    injective character $\omega'_k$ of $(\mathbf Z/g_k)^\times$ with $\chi_{j,w}(a)=\omega'_k(a)^{2j}$
    for $w$ over $g_k$ and $a\in\mathbf Z$ prime to $g_k$; $\mu_{0,j,w}(y)=\iota(\chi_{j,w}(y\bmod w))^{-1}$
    for $w\in S_G$ and $y\in\mathcal O_E$ prime to $w$; $\iota(\mu_j(v))=\mu_{0,j}(\varpi_v)$ for the
    places $v$ prime to $g_0g_1$, with $\varpi_v$ a uniformizer at $v$; $\iota(\eta(v))=\eta(\varpi_v)
    \mathrm N(v)^{-1/2}$; and, with $\xi_0\in E$ the host element (Definition~\ref{th-definite-datum}),
    $\iota(\psi_{0,w}(c))=\exp(2\pi ib/\mathrm N(w))$ whenever $c\in\mathcal O_E/w$ and $b\in\mathbf Z$
    satisfy $\xi_0c\equiv b$ modulo $w$.
  \item[(c)] \emph{Tuples.} A selection of a place of $E$ over each place of $F$; a set $S$ of primes
    containing the $T_u$, $g_0$, $g_1$ and the primes dividing the conductor of the host; a finite set
    of tuples $\mu$, each a character of the finite ideles of $E$ with values in a field $K_\mu$ with
    an embedding $\iota$ into $\mathbf C$, unramified at the places of $F$ that are not \emph{good for}
    $S$ (a place is good for $S$ if it splits in $E$, has degree one and has residue characteristic
    outside $S$), injectively indexed, with $\mu_0=1$ and two distinguished
    tuples $\mu_*,\mu'_*$; for each tuple a Hecke character $\xi_\mu$ of $E$ read off from $\mu$, that
    is $\xi_\mu(a)=\iota(\mu(a_f))$ for every idele $a$ with finite part $a_f$, with
    $\xi_{\mu_*}=\delta$ and $\xi_{\mu'_*}=\delta'$; and for each tuple sharp local constants and index
    data that satisfy their conditions at level $m-2$. The tuple $\mu$ is \emph{ramified} at a place
    $v$ of $F$ if it is nontrivial on the units of $E_w$ for some place $w$ of $E$ over $v$. The
    \emph{twisting places} $Q_E(\mu)$ are the selected places over the places of $F$ where $\mu$ is
    ramified; these places are good for $S$.
  \item[(d)] \emph{Ambient data}: a complex number $\delta_A$ that is a twelfth root of the
    discriminant of the period lattice of the host, a vacuum vector for it, and periods
    $c_D(\tau)$, $\tau\in\mathcal T$, that generate in the required sense; the spaces of the base
    are diagonal; and for every word $u$, $d^*_u$ (Definition~\ref{th-definite-datum}(3)) is the mass
    of the standard Hermitian lattice of rank $k$ over $\mathcal O_E$, times the $f$-th powers of the
    local index factors at $g_0$ and $g_1$, times the product over $\tau\in\mathcal T$ of the
    dimensions of the representations with highest weight $(u_1(\tau),\dots,u_k(\tau))$.
  \item[(e)] \emph{Classes.} Let $S_T$ be the set of places of $E$ over $T_u$, $S_G$ those over the
    level primes $g_0,g_1$, and $Q=Q_E(\mu)$. A \emph{decorated lattice} of rank $n$ is a unimodular
    $\mathcal O_E$-lattice $N\subset E^n$ for the standard Hermitian form, with a nonzero linear form
    $N_P\to\mathcal O_E/P$ for each $P\in S_T$, an isomorphism $N/wN\cong(\mathcal O_E/w)^n$ for each
    $w\in S_G$, and for each $P\in Q$ a \emph{frame}, a linear automorphism of $E^n$ carrying the
    standard local lattice at $P$ to $N_P$. For every word $u$, tuple $\mu$ and level $i\leq m$ the
    true columns give finitely many \emph{indices} $b$, and for each a group $G_b$, the product of the
    isometry group of the standard Hermitian form on $E^{N(i)}$, a level subgroup of $\prod_{w\in S_G}GL_{N(i)}(\mathcal O_E/w)$ and a
    level subgroup of $\prod_{P\in Q}GL_{N(i)}(\mathcal O_{E,P})$, with a character $G_b\to K_\mu^\times$;
    $G_b$ acts on the decorated lattices of rank $N(i)$, at level $m-1$ together with a point of the
    coupling set of $N$ (Definition~\ref{is-theta-kernel}(5)). The \emph{classes} $\mathrm{cls}_i$ of
    $C_{u,\mu}$ are the pairs of an index and an orbit, each with a chosen representative; its
    stabilizer $\Gamma_c$ is finite and acts on $E^{N(i)}$ through unitary matrices.
  \item[(f)] \emph{Theta kernels.} For classes $c\in\mathrm{cls}_i$ and $c'\in\mathrm{cls}_{i+1}$ with
    representatives of underlying lattices $N\subset E^{N(i)}$ and $N'\subset E^{N(i+1)}$, the theta
    kernel $\vartheta_i(c,c')$ of $C_{u,\mu}$ (Definition~\ref{th-iterated-lift}) is the jet map
    $J_{j_{c_D}(G),0}$ of Definition~\ref{is-jet-operator}(3), where $j_{c_D}(G)$ are the jets
    (Definition~\ref{is-jet-operator}(2)) with the periods $c_D$ of (d) of the entire function
    $G=\sum_x\Xi(x)\,\Theta_{N,N',c,x}$ on $V_{N(i),N(i+1)}$: the combination, over representatives
    $x$ of the cosets, of the coset theta functions of Definition~\ref{is-theta-kernel}(2), with
    the coefficient function $\Xi$ of the true columns at $(u,\mu,i)$, which is that of an entry, of
    the coupling or of the last step according as $i+2<m$, $i+2=m$ or $i+1=m$. The twists are
    $\rho_i(c)=\tilde\mu_{i+1}(d(c))^{-1}\tilde\xi_\mu(d(c))^{-1}\tilde\chi_\Psi(d(c))^{-N(i+1)/2}$ and
    the corresponding $\rho'_i(c')$, with $d(c)$ the determinant class of the representative and
    $\tilde\mu_{i+1}$, $\tilde\xi_\mu$, $\tilde\chi_\Psi$ the characters of $\mu_{0,i+1}$, $\xi_\mu$,
    $\chi_\Psi$ on determinant classes.
  \item[(g)] \emph{Point data.} Let $\mu$ be a tuple and $Q=Q_E(\mu)$. The true columns give a group
    $G$, the product of the isometry group of the standard Hermitian form on $E^{N(m)-1}$, a level
    subgroup of $\prod_{w\in S_G}GL_{N(m)-1}(\mathcal O_E/w)$ and a level subgroup of
    $\prod_{P\in Q}GL_{N(m)-1}(\mathcal O_{E,P})$, with a character $\chi\colon G\to L^\times$; it acts on
    the set $\mathcal D$ of decorated lattices of rank $N(m)-1$ for the empty set in place of $S_T$.
    Write $\gamma(g)\in GL_{N(m)-1}(E)$ for the first component of $g\in G$, $L_d$ for the underlying
    lattice of $d\in\mathcal D$, and $\rho_\sigma(h)$ for $\rho_\sigma((\tau(h))_\tau)$ when
    $h\in GL_{N(m)-1}(E)$ (Definition~\ref{pd-weight-lattices}(4)), acting on $V_\sigma(L)$ and,
    through $\iota$, on $V_\sigma(\mathbf C)$. The \emph{point data of $\mu$}, for $\sigma$ with
    $n=N(m)-1$ (Definition~\ref{pd-point-data}), are as follows. The point classes are the $G$-orbits
    $c$ on $\mathcal D$, each with a chosen representative $d_c$; $L_c=L_{d_c}$, and $\Gamma_c$ is the
    stabilizer of $d_c$ in $G$, acting on $V_\sigma(L)$ by $x\mapsto\chi(g)\rho_\sigma({}^t\gamma(g)^{-1})x$.
    The top forms are the functions $F$ on $\mathcal D$ with values in $V_\sigma(\mathbf C)$ such that
    $F(gd)=\iota(\chi(g))^{-1}\rho_\sigma(\gamma(g))F(d)$ for $g\in G$ and $d\in\mathcal D$, and
    $R_c(F)=F(d_c)$. For every word $u$, level $m$ of $C_{u,\mu}$ has an index $b_m$, a map
    $d\mapsto\tilde d$ from $\mathcal D$ to the decorated lattices of rank $N(m)$ on which $G_{b_m}$
    acts, where $\tilde d$ has underlying lattice $L_d\oplus\bar\pi_u^{-1}\mathcal O_E$, and a
    homomorphism $g\mapsto\tilde g$ from $G$ to $G_{b_m}$ such that $\gamma(\tilde g)=\gamma(g)\oplus1$
    and the character of $G_{b_m}$ takes the value $\chi(g)$ at $\tilde g$. The point class $c$ is
    sent to the class $\tilde c$ of $\tilde d_c$, and the homomorphism $\Gamma_c\to\Gamma_{\tilde c}$
    is $g\mapsto h^{-1}\tilde gh$, where $h\in G_{b_m}$ carries the representative of $\tilde c$ to
    $\tilde d_c$.
  \item[(h)] \emph{Laws and pins.} Let $u$ be a word and $\mu$ a tuple. Each group $G_b$ of each
    level of $C_{u,\mu}$, and the group $G$ of (g), acting on its decorated lattices $d$ with first
    components $\gamma(g)$ and character $\chi$, has finitely many orbits and finite stabilizers;
    ${}^t\overline{\gamma(g)}\gamma(g)=1$ and $\chi(g)$ is a root of unity for every $g$; $\gamma(g)$
    carries the underlying lattice of $d$ onto that of $gd$, and these underlying modules are
    lattices; and $g=1$ whenever $gd=d$ and $\gamma(g)=1$. Moreover $|\rho_i(c)|=|\rho'_i(c')|=1$ for
    $i<m$, $|\Gamma_c|$ divides the stabilizer denominator of $C_{u,\mu}$ for every class $c$ of level
    $i\leq m$, $\widetilde{gd}=\tilde g\tilde d$ for $g\in G$ and $d\in\mathcal D$, the homomorphism
    $g\mapsto\tilde g$ of (g) is injective, and the representative of the class of $\tilde d$ is
    $\tilde d'$ for some $d'\in\mathcal D$. The data of characters and levels at each $(u,\mu)$
    satisfy their conditions, the periods lie over the levels, $\xi_{\mu_0}=1$, the index data are
    nonempty at level $0$, the classes of level $0$ for different tuples correspond bijectively, and
    the underlying modules of the representatives are lattices; the inclusion of $X'$ into $X_m$
    carries the standard basis of $X'$ to the first vectors of the standard basis of $X_m$, and
    $c_u\pi_u=\pi_{u_0}$ for every word $u$.
  \end{enumerate}
\item A definite-side datum $D$ (Definition~\ref{th-definite-datum}) is \emph{built from true
  columns} if it is assembled from true columns and from later stages of the construction (row,
  division, size, norm and jet data). Its column chains are then the $C_{u,\mu}$. For every word
  $u$ the point classes and the point lattice $\Lambda^{\mathrm{pt}}_u$ of
  Definition~\ref{pd-period-data}(4) are those of the point data of $\mu_0$, and
  $j^{\mathrm{pt}}_u(x)$ is $\iota(x)$. For every word $u$ and tuple $\mu$ the space $\mathcal A_\mu$
  is identified with the top forms of the point data of $\mu$.
\item The \emph{period rows} of $D$ are \emph{integral} if for every word $u$, every tuple $\mu$,
  every element $F=\sum_{s\in S_\mu}\iota(c_s)s$ of $\mathcal A_\mu$ with $c_s\in\mathcal O$
  (Definition~\ref{pd-period-data}(3)) and every point class $c$ of $\mu$, the row of $F$ at $c$ is
  integral.
\item The true columns satisfy the \emph{degree-zero condition} if for every word $u$ and tuple $\mu$
  there is an $L$-linear isomorphism from the space of the lattice of forms in degree $0$ of $C_u$
  to that of $C_{u,\mu}$ that maps the one lattice onto the other and is compatible with the
  identification of functions in degree $0$ of Definition~\ref{th-definite-datum}(5). The lattices
  of forms and the point data of the chains $C_{u,\mu}$ are then twisted forms for $D$
  (Definition~\ref{pd-period-data}(5)), the \emph{twisted forms of the true columns}.
\item A \emph{tie} of the tuples of $D$ attaches to each tuple $\mu$ a Hecke character $\xi'_\mu$ of
  $E$ of finite order, trivial on $F$ and tame, injective in $\mu$, with $\xi'_{\mu_0}=1$ and with
  prescribed values at the two distinguished tuples and at the twisting tuples of the places of
  $L$. The tie is \emph{tied to the true columns} if $\xi'_\mu=\xi_\mu$ for every tuple $\mu$.
\item A slot (Definition~\ref{is-slot-rows}) is \emph{over true columns} if at each of its cores the
  definite-side datum is built from true columns, its period rows are integral, its twisted forms
  are those of the true columns, its tuples have a tie tied to the true columns, and the jet
  conditions of the true columns hold.
\end{enumerate}
\end{definition}
\leanref{Development/AbelianVarietyTheory/Definitions/AutomorphicMembership/TrueColumnChoicesHeckeCharReferenceEqOne.lean}{GlobalTheory.Mem.Col2.ColIn}
\leanref{Development/AbelianVarietyTheory/Definitions/AutomorphicMembership/LatAtLattice.lean}{GlobalTheory.Mem.Col2.TC.rec7}
\leanref{Development/AbelianVarietyTheory/Definitions/AutomorphicMembership/IsRationalUnitaryDeterminantMapRationalUnitaryTo.lean}{GlobalTheory.Mem.Chr.HFrame}
\leanref{Formalism/LocalToGlobal/UnitaryNorms/CharHost.lean}{Formalism.L2.Host.IsMu0}
\leanref{Formalism/LocalToGlobal/UnitaryNorms/CharHost.lean}{Formalism.L2.Host.IsEtaT}
\leanref{Formalism/LocalToGlobal/UnitaryNorms/CharHost.lean}{Formalism.L2.Host.IsChiPsi}
\leanref{Formalism/LocalToGlobal/UnitaryNorms/CharHost.lean}{Formalism.L2.Host.IsDelta}
\leanref{Formalism/LocalToGlobal/Columns/AbelianHostFieldConditions.lean}{Formalism.L2.Col.Chars}
\leanref{Formalism/LocalToGlobal/Columns/AbelianHostFieldConditions.lean}{Formalism.L2.Col.Chars.Ok}
\leanref{Formalism/LocalToGlobal/Columns/AbelianHostFieldConditions.lean}{Formalism.L2.Col.Chars.Tied}
\leanref{Formalism/LocalToGlobal/Columns/FiniteSumNonzeroLargeDepth.lean}{Formalism.L2.Col.Chars.WeilTied}
\leanref{Formalism/LocalToGlobal/Columns/AbelianHostFieldConditions.lean}{Formalism.L2.Col.IsMu0}
\leanref{Formalism/LocalToGlobal/Columns/RowForms.lean}{Formalism.L2.Col.IsSel}
\leanref{Formalism/LocalToGlobal/Columns/StdLatUnimodular.lean}{Formalism.L2.Col.Deco}
\leanref{Formalism/LocalToGlobal/Columns/FiniteSumNonzeroLargeDepth.lean}{Formalism.L2.Col.DecoF}
\leanref{Development/AbelianVarietyTheory/Definitions/AutomorphicMembership/PredMemGamRanks.lean}{GlobalTheory.Mem.Col2.GStep}
\leanref{Development/AbelianVarietyTheory/Definitions/AutomorphicMembership/PredMemGamRanks.lean}{GlobalTheory.Mem.Col2.stepAt}
\leanref{Development/AbelianVarietyTheory/Definitions/AutomorphicMembership/XiAtPenultimateEqXiCpl1.lean}{GlobalTheory.Mem.Col2.thAt}
\leanref{Development/AbelianVarietyTheory/Definitions/AutomorphicMembership/LatAtLattice.lean}{GlobalTheory.Mem.Col2.TC.chIn}
\leanref{Development/AbelianVarietyTheory/Definitions/AutomorphicMembership/ReferenceCharactersFactorTame.lean}{GlobalTheory.Mem.Col2.rhoAt}
\leanref{Development/AbelianVarietyTheory/Definitions/GlobalTwisting/LowerTieTwAt.lean}{GlobalTheory.PerQ.RowVal}
\leanref{Development/AbelianVarietyTheory/Definitions/GlobalTwisting/LowerTieTwAt.lean}{GlobalTheory.PerQ.RowInt}
\leanref{Development/AbelianVarietyTheory/Definitions/AutomorphicMembership/LatAtLattice.lean}{GlobalTheory.Mem.Col2.TC.Law0}
\leanref{Development/AbelianVarietyTheory/Definitions/AutomorphicMembership/LatAtLattice.lean}{GlobalTheory.Mem.Col2.TC.tw}
\leanref{Formalism/LocalToGlobal/UnitaryNorms/CharTied.lean}{Formalism.L2.Host.TupTied}
\leanref{Development/AbelianVarietyTheory/Definitions/GlobalTwisting/AndLeft.lean}{GlobalTheory.Fc.Cv.OverCols}

\begin{lemma}\label{pd-standard-comparison}
Let $D$ be a definite-side datum built from true columns (Definition~\ref{pd-true-columns}), let
$u$ be a word and $\mu$ a tuple, and if $\mu\neq\mu_0$ suppose that the true columns satisfy the
degree-zero condition. Let the point data of $\mu$ (Definition~\ref{pd-true-columns}(1)(g)) be given
by the group $G$ with character $\chi$ and first components $\gamma(g)$, acting on the decorated
lattices $d\in\mathcal D$, and have point classes $c$ with representatives $d_c$, stabilizers
$\Gamma_c$, rows $R_c$, point lattice
$\Lambda^{\mathrm{pt}}$ and period $P$, and let $\mathcal P_{u,\mu}$ be the period pairing of
$(u,\mu)$: $\mathcal P_u$ if $\mu=\mu_0$, and that of the twisted forms of the true columns
otherwise. For $y\in V_\sigma(\mathbf C)$ let $\bar y$ be the polynomial with the complex conjugate
coefficients, and for a top form $F$ let $j(F)$ be the vector of the point space with entries
$j(F)_c=\overline{R_c(F)}$. Let $V^\Gamma$ be the space of vectors $z$ of the point space with
$\iota(\chi(g))\rho_\sigma({}^t\gamma(g)^{-1})z_c=z_c$ for every point class $c$ and every
$g\in\Gamma_c$, and for vectors $z,z'$ of the point space put
\[
 \langle\!\langle z',z\rangle\!\rangle=\sum_c\frac1{|\Gamma_c|}\,\langle z_c,\overline{z'_c}\rangle_\sigma .
\]
\begin{enumerate}
\item The pairing $\mathcal P_{u,\mu}$ is the period $P$: it pairs the point space of the point data
  of $\mu$ with their top forms, which form $\mathcal A_\mu$, and $\mathcal P_{u,\mu}(z,F)=P(z,F)$.
  The point lattice of $(u,\mu)$ is $\Lambda^{\mathrm{pt}}$, with the injection $x\mapsto\iota(x)$.
\item The map $j$ is conjugate-linear, it maps the space of top forms onto $V^\Gamma$, and
  $P(z,F)=\langle\!\langle j(F),z\rangle\!\rangle$ for every vector $z$ of the point space and every
  top form $F$.
\item Suppose that the tuples of $D$ have a tie tied to the true columns, with characters $\xi'_\mu$
  (Definition~\ref{pd-true-columns}(5)). Let $v$ be a finite place of the totally real subfield of
  $E$. Then $\xi'_\mu$ is ramified at some place $w$ of $E$ over $v$, that is, nontrivial on the
  units of $E_w$, if and only if $v$ lies under a place of $Q_E(\mu)$.
\end{enumerate}
\end{lemma}
\leanref{Development/AbelianVarietyTheory/Definitions/AutomorphicMembership/LatAtLattice.lean}{GlobalTheory.Mem.Col2.TC.cols}
\leanref{Development/AbelianVarietyTheory/Proofs/IwasawaDecomposition/CevInjective.lean}{GlobalTheory.Mem.Col2.PerConj.period_conj}
\leanref{Development/AbelianVarietyTheory/Proofs/IwasawaDecomposition/TrueColumnChoiceExistsStandardChainData.lean}{GlobalTheory.Mem.Col2.TC.underQ_iff}
\leanref{Development/AbelianVarietyTheory/Proofs/IwasawaDecomposition/ExistsStandardColumnComparisonsTrueColumnsStatement.lean}{GlobalTheory.T.exists_standardColumnComparisons_trueColumns}
\leanref{Development/AbelianVarietyTheory/Proofs/IwasawaDecomposition/TrueColumnChoiceExistsStandardChainDataStatement.lean}{GlobalTheory.T.trueColumnChoice_exists_standardChainData}

\begin{proof}
(1) The datum is assembled from the true columns so that the period map of the member $u$ is that
of the point data of $\mu_0$: its space $\mathcal W_u$ is their point space, its space
$\mathcal M^{\mathrm t}_u$ is the space of their top forms, identified with $\mathcal A_{\mu_0}$,
and $\mathcal P_u=P$. In the same way the period map of the twisted chain $C_{u,\mu}$ is that of the
point data of $\mu$. The point lattice of $u$ is $\Lambda^{\mathrm{pt}}$ with
$j^{\mathrm{pt}}_u=\iota$ by Definition~\ref{pd-true-columns}(2). For $\mu\neq\mu_0$ the twisted
forms of the true columns are the lattices of forms and the point data of the chains $C_{u,\mu}$
(Definition~\ref{pd-true-columns}(4)), so the twisted point lattice is $\Lambda^{\mathrm{pt}}$ with
$j^{\mathrm{pt}}_{u,\mu}=\iota$.

(2) The conditions defining $\nabla(\nu)$ have integer coefficients, so $V_\sigma(\mathbf C)$ is
spanned by polynomials with rational coefficients and $y\mapsto\bar y$ is a conjugate-linear
involution of it. Hence $j$ is conjugate-linear, and
\[
 \langle\!\langle j(F),z\rangle\!\rangle=\sum_c\frac1{|\Gamma_c|}\,\langle z_c,R_c(F)\rangle_\sigma=P(z,F).
\]
Two rules move the bar past the actions. First,
$\overline{\rho_\sigma(h)y}=\rho_\sigma(\bar h)\bar y$ for $h\in GL_{N(m)-1}(\mathbf C)^{\mathcal T}$,
with $\bar h$ the entrywise conjugate, since
$(\rho_\sigma(h)y)(X)=\prod_\tau\det(h_\tau)^{\sigma_{\tau,N(m)-1}}y\big(({}^th_\tau X_\tau)_\tau\big)$.
Second, $E$ is a CM field, so $\iota(\tau(\bar a))=\overline{\iota(\tau(a))}$ for $a\in E$ and
$\tau\in\mathcal T$; for $g\in G$ the condition ${}^t\overline{\gamma(g)}\gamma(g)=1$ of
Definition~\ref{pd-true-columns}(1)(h) then gives
$\overline{\iota\tau(\gamma(g))}=\iota\tau({}^t\gamma(g)^{-1})$. Since $\chi(g)$ is a root of unity,
$\overline{\iota(\chi(g))}^{-1}=\iota(\chi(g))$. Together,
\[
 \overline{\iota(\chi(g))^{-1}\rho_\sigma(\gamma(g))y}=\iota(\chi(g))\,\rho_\sigma({}^t\gamma(g)^{-1})\,\bar y
 \qquad(g\in G,\ y\in V_\sigma(\mathbf C)).\tag{$*$}
\]
Let $F$ be a top form, $c$ a point class and $g\in\Gamma_c$. Since $gd_c=d_c$,
$R_c(F)=F(gd_c)=\iota(\chi(g))^{-1}\rho_\sigma(\gamma(g))R_c(F)$, and by $(*)$ the vector
$\overline{R_c(F)}$ is fixed by $\iota(\chi(g))\rho_\sigma({}^t\gamma(g)^{-1})$. So $j(F)\in V^\Gamma$.

Conversely, let $z\in V^\Gamma$. For $d\in\mathcal D$ in the class $c$ choose $h\in G$ with
$hd_c=d$, and put $F(d)=\iota(\chi(h))^{-1}\rho_\sigma(\gamma(h))\overline{z_c}$. If also $h'd_c=d$,
then $g=h^{-1}h'\in\Gamma_c$, and conjugating $z_c=\iota(\chi(g))\rho_\sigma({}^t\gamma(g)^{-1})z_c$
with $(*)$ gives $\overline{z_c}=\iota(\chi(g))^{-1}\rho_\sigma(\gamma(g))\overline{z_c}$; applying
$\iota(\chi(h))^{-1}\rho_\sigma(\gamma(h))$ shows that $F(d)$ does not depend on $h$. Since $\chi$ and
$g\mapsto\rho_\sigma(\gamma(g))$ are homomorphisms,
$F(g'd)=\iota(\chi(g'))^{-1}\rho_\sigma(\gamma(g'))F(d)$ for $g'\in G$, so $F$ is a top form. Taking
$h=1$ gives $F(d_c)=\overline{z_c}$, that is, $j(F)=z$.

(3) The tie is tied to the true columns, so $\xi'_\mu=\xi_\mu$. Let $w$ be a place of $E$. Since
$\xi_\mu(a)=\iota(\mu(a_f))$ for every idele $a$ (Definition~\ref{pd-true-columns}(1)(c)) and
$\iota$ is injective, $\xi_\mu$ is nontrivial on the units of $E_w$ if and only if $\mu$ is. So
$\xi'_\mu$ is ramified at some place over $v$ if and only if $\mu$ is ramified at $v$. The places of
$Q_E(\mu)$ are the selected places over the places where $\mu$ is ramified, and the selected place
over a place lies over it, so $v$ lies under a place of $Q_E(\mu)$ if and only if $\mu$ is ramified
at $v$.
\end{proof}

\begin{theorem}\label{pd-period-integrality-from-rows}
Let $D$ be a definite-side datum built from true columns whose period rows are integral
(Definition~\ref{pd-true-columns}).
\begin{enumerate}
\item Period integrality holds for $D$ (Definition~\ref{pd-period-integrality}).
\item Suppose moreover that the true columns satisfy the degree-zero condition, and that the tuples of
  $D$ have a tie tied to the true columns. Then for every choice of divisibility data for $D$,
  twisted period integrality holds for $D$ with the twisted forms of the true columns
  (Definition~\ref{pd-twisted-period-integrality}).
\end{enumerate}
\end{theorem}
\leanref{Development/AbelianVarietyTheory/Proofs/GlobalInterpolation/PeriodAdmissibleOfIntegralPeriodRowsStatement.lean}{GlobalTheory.T.periodClassSumWeyl_of_rowInt}
\leanref{Development/AbelianVarietyTheory/Proofs/GlobalInterpolation/AdmissiblePeriodTwistOfIntegralRowsStatement.lean}{GlobalTheory.T.admissible_period_twist_of_integral_rows}

\begin{proof}
% Lean: the comparisons are exists_standardColumnComparisons_trueColumns (= TC.colsTiedSym_std) and,
% for the twisted tuples, trueColumnChoice_exists_standardChainData (= TC.colsTwPTiedS_std, which uses
% the tie through underQ_iff); the class terms are perDec0At_std and perDecTwT_of_std, summed by
% petersson_eq_sum_restrict in PerRed.admissiblePeriodZero_of_classPeriodDecomposition and
% PerRedTw.perDecTw_of_at.
(1) The construction of $y$ below does not use that the stabilizer denominators are defined. Let
$u$ be a word, $x\in\Lambda^{\mathrm{pt}}_u$ and $F$ an integral top form for $\mu_0$. Apply
Lemma~\ref{pd-standard-comparison} with $\mu=\mu_0$. By its part (1),
$\mathcal P_u(j^{\mathrm{pt}}_u(x),F)=P(\iota(x),F)$ for the point data of $\mu_0$, and by its
part (2) this is the Hermitian product $\langle\!\langle j(F),\iota(x)\rangle\!\rangle$, with
$j(F)\in V^\Gamma$. The vector $\iota(x)$ also lies in $V^\Gamma$: each $x_c$ is fixed by $\Gamma_c$
(Definition~\ref{pd-point-data}), which acts by $x\mapsto\chi(g)\rho_\sigma({}^t\gamma(g)^{-1})x$
(Definition~\ref{pd-true-columns}(1)(g)). For a vector $z$ of the point space let $z_{[c]}$ be the
vector equal to $z_c$ at $c$ and to $0$ at the other point classes. The Hermitian product is the sum
over the point classes $c$ of the products of the restrictions to $c$, so
\[
 \mathcal P_u(j^{\mathrm{pt}}_u(x),F)=\sum_c\langle\!\langle j(F)_{[c]},\iota(x)_{[c]}\rangle\!\rangle .
\]
Fix a point class $c$. Only $c$ contributes to the products with $\iota(x)_{[c]}$, so by part (2)
of the lemma
$\langle\!\langle j(F)_{[c]},\iota(x)_{[c]}\rangle\!\rangle=\langle\!\langle j(F),\iota(x)_{[c]}\rangle\!\rangle
=P(\iota(x)_{[c]},F)$. The rows of $F$ are integral, so Lemma~\ref{pd-class-term-integral} gives
$a_c\in\mathcal O_L$ with $|\Gamma_c|\,\langle\!\langle j(F)_{[c]},\iota(x)_{[c]}\rangle\!\rangle=\iota(a_c)$.
Put $y=\sum_c|\Gamma_c|^{-1}a_c\in L$. Then $\iota(y)=\mathcal P_u(j^{\mathrm{pt}}_u(x),F)$.

The stabilizer $\Gamma_c$ injects into the stabilizer $\Gamma_{\tilde c}$ of a class
$\tilde c\in\mathrm{cls}_m$ of $C_u$ (Definition~\ref{pd-period-data}(4)), and $|\Gamma_{\tilde c}|$
divides $D_\Gamma(u)$ (Definition~\ref{th-column-chain}(2)). So $|\Gamma_c|$ divides $D_\Gamma(u)$.
Let $w$ be a finite place of $L$ and $c$ a point class with $a_c\neq0$. Then $v_w(a_c)\geq0$ and
$v_w(|\Gamma_c|)=e_wv_{p_w}(|\Gamma_c|)\leq e_wv_{p_w}(D_\Gamma(u))$, so
$v_w(|\Gamma_c|^{-1}a_c)\geq-e_wv_{p_w}(D_\Gamma(u))$. By the ultrametric inequality, $y=0$ or
$v_w(y)\geq-e_wv_{p_w}(D_\Gamma(u))$.

(2) Suppose that the stabilizer denominators of $D$ are defined. Let $u$ be a word, $\mu\neq\mu_0$ a
used tuple, $x$ a vector of the twisted point lattice of $(u,\mu)$ and $F$ an integral top form for
$\mu$. The comparison for $(u,\mu)$ describes the point classes of $(u,\mu)$ as the classes of a
level group whose level structure sits at the places over $g_0g_1$ and at the places $v$ of the
totally real subfield of $E$ over which the character $\xi'_\mu$ of the tie ramifies. The tie is
tied to the true columns, so by Lemma~\ref{pd-standard-comparison}(3) these $v$ are the places under
$Q_E(\mu)$, where the group $G$ of the point data of $\mu$ has its level subgroup
(Definition~\ref{pd-true-columns}(1)(g)). Hence the point classes of $(u,\mu)$ are those of the
point data of $\mu$, each matched with itself. The twisted forms are those of the true columns, so by
Lemma~\ref{pd-standard-comparison}(1) and (2) the twisted point lattice is that of the point data of
$\mu$ and $\mathcal P_{u,\mu}(j^{\mathrm{pt}}_{u,\mu}(x),F)=\langle\!\langle j(F),\iota(x)\rangle\!\rangle$.
The computation of (1) gives $y=\sum_c|\Gamma_c|^{-1}a_c$ with $a_c\in\mathcal O_L$ and
$\iota(y)=\mathcal P_{u,\mu}(j^{\mathrm{pt}}_{u,\mu}(x),F)$. Now $\Gamma_c$ injects into the
stabilizer of a class of level $m$ of $C_{u,\mu}$, whose order divides $D_\Gamma(u,\mu)$, and
$D_\Gamma(u,\mu)=D_\Gamma(u_0)=D_\Gamma(u)$ because the stabilizer denominators are defined. The
bound of (1) follows at every finite place $w$ of $L$, in particular at those with
$p_w\notin Q(\mu)$.
\end{proof}

\begin{remark}\label{pd-comparison-class-split}
A shorter route to the element $y$ of the proof of Theorem~\ref{pd-period-integrality-from-rows}
splits the period $P$ directly. By Lemma~\ref{pd-standard-comparison}(1),
$\mathcal P_{u,\mu}(j^{\mathrm{pt}}_{u,\mu}(x),F)=P(\iota(x),F)$, and $\iota(x)$ is the sum over the
point classes $c$ of the vectors $\iota(x)_{[c]}$. The period $P$ is linear in its first argument, so
this is $\sum_cP(\iota(x)_{[c]},F)=\sum_c|\Gamma_c|^{-1}\iota(a_c)$ with the $a_c$ of
Lemma~\ref{pd-class-term-integral}. This route uses neither the Hermitian product of
Lemma~\ref{pd-standard-comparison}(2) nor, for $\mu\neq\mu_0$, the tie, since
Definition~\ref{pd-true-columns}(4) makes the twisted point lattice that of the point data of $\mu$.
\end{remark}

\begin{definition}\label{pd-row-stage}
Let true columns with row weights $\sigma$ be given (Definition~\ref{pd-true-columns}), and let $\mu$
be a tuple. The space $\widetilde{\mathcal M}_\mu$ consists of the functions $F$ on the decorated
$\mathcal O_E$-lattices of rank $N(m)-1$ for $\mu$, with values in $V_\sigma(L)$, that are
equivariant under the level groups of the true columns. Such an $F$ is \emph{$\nabla$-integral} if
$F(\Lambda)\in\nabla_\sigma(N_\Lambda)$ for every decorated lattice $\Lambda$, where $N_\Lambda$ is its
underlying lattice (Definition~\ref{pd-weight-lattices}). A \emph{row stage} of the construction over
the true columns consists of later stages of the construction and, for every tuple $\mu$, a
$\mathbf C$-linear isomorphism $e_\mu$ from $\widetilde{\mathcal M}_\mu\otimes_{L,\iota}\mathbf C$ onto
the space $\mathcal A_\mu$ of top forms of the datum built from them, such that:
\begin{enumerate}
\item for every word $u$, every $F\in\widetilde{\mathcal M}_\mu\otimes_{L,\iota}\mathbf C$ and every
  point class $c$ of $\mu$, the row of $e_\mu(F)$ at $c$ is the value of $F$ at a decorated lattice
  attached to the representative of $c$, whose underlying lattice is $L_c$;
\item every element $\sum_{s\in S_\mu}\iota(c_s)s$ of $\mathcal A_\mu$ with $c_s\in\mathcal O$ is
  $e_\mu(F\otimes1)$ for a $\nabla$-integral $F\in\widetilde{\mathcal M}_\mu$.
\end{enumerate}
A definite-side datum is \emph{built at a row stage} if it is built from the true columns and from
later stages that refine those of a row stage.
\end{definition}
\leanref{Development/AbelianVarietyTheory/Definitions/AutomorphicMembership/NonzeroPlacesOverPrimesLifts.lean}{GlobalTheory.Mem.R3.Rw.Pre2}
\leanref{Formalism/LocalToGlobal/Columns/RowForms.lean}{Formalism.L2.Col.MtildeL}
\leanref{Development/AbelianVarietyTheory/Definitions/AutomorphicMembership/SameForm.lean}{GlobalTheory.Mem.Df.SameForm}

\begin{proposition}\label{pd-rows-integral-built}
Let $D$ be a definite-side datum built at a row stage (Definition~\ref{pd-row-stage}). Then the period
rows of $D$ are integral (Definition~\ref{pd-true-columns}).
\end{proposition}
\leanref{Development/AbelianVarietyTheory/Proofs/GlobalInterpolation/PeriodRowLatticesIntegralProvisionalRowsStatement.lean}{GlobalTheory.T.periodRowLatticesIntegral_provisionalRows}

\begin{proof}
Let $u$ be a word, $\mu$ a tuple, $G=\sum_{s\in S_\mu}\iota(c_s)s$ with $c_s\in\mathcal O$, and $c$ a
point class of $\mu$. By Definition~\ref{pd-row-stage}(2), $G=e_\mu(F\otimes1)$ for a $\nabla$-integral
$F\in\widetilde{\mathcal M}_\mu$. By (1), $R_c(G)$ is the value of $F\otimes1$ at a decorated lattice
$\Lambda_c$ with underlying lattice $L_c$, that is, $\iota(F(\Lambda_c))$. Since $F$ is
$\nabla$-integral, $F(\Lambda_c)\in\nabla_\sigma(L_c)$, so the row of $G$ at $c$ is integral.
\end{proof}

\begin{theorem}\label{pd-period-integrality-along-families}
Let $Q$ be a slot over true columns (Definition~\ref{pd-true-columns}) and $\Lambda$ link data
(Definition~\ref{is-slot-rows}). For a class $\mathcal A$ of families let $\mathcal W(\mathcal A)$, the
\emph{working subclass}, be the class of families in $\mathcal A$ that satisfy the conditions defining
$\mathcal C_N$ and $\mathcal C_{\mathrm{tw}}$ in Definition~\ref{asm-designed}, have chosen constants
$(3,C_H)$ for some $C_H$, and satisfy both size conditions there. Then the
following holds along $(Q,\Lambda)$: for every class $\mathcal A$ and every selection for
$\mathcal A$, eventually along $\mathcal W(\mathcal A)$, at all large scales
(Definition~\ref{pd-family}), every member $D_{x,i}$ satisfies
\begin{enumerate}
\item period integrality (Definition~\ref{pd-period-integrality}), and
\item twisted period integrality for the divisibility data and the twisted forms of the family
  (Definition~\ref{pd-twisted-period-integrality}).
\end{enumerate}
\end{theorem}
\leanref{Development/AbelianVarietyTheory/Proofs/GlobalInterpolation/ReferencePeriodIntegralBundlePointsOfOverColumnsStatement.lean}{GlobalTheory.T.referencePeriodIntegral_bundlePoints_of_overColumns}
\leanref{Development/AbelianVarietyTheory/Proofs/GlobalInterpolation/TwistedPeriodsIntegralAtBundlePointsStatement.lean}{GlobalTheory.T.twistedPeriods_integral_at_bundlePoints}
\leanref{Development/AbelianVarietyTheory/Definitions/GlobalTwisting/ValidForConstructedPointClassesEventuallyWorking.lean}{GlobalTheory.Fc.W3}

\begin{proof}
Let $\beta$ be a base point of $Q$. Since $Q$ is over true columns, the definite-side datum at
$\beta$ is built from true columns, its period rows are integral, its twisted forms are those of the
true columns, and its tuples have a tie tied to the true columns. By
Theorem~\ref{pd-period-integrality-from-rows}, period integrality holds at $\beta$, and twisted
period integrality holds at $\beta$ for every choice of divisibility data with these twisted forms.

Take for the condition on bundles of Definition~\ref{is-slot-rows} that these two properties hold at
the datum of the bundle; by the first paragraph every bundle at a base point of $Q$ satisfies it.
Let $\mathcal A$ be a class and take a selection for $\mathcal A$ whose bundles are built from
$(Q,\Lambda)$ and satisfy this condition. Let $\mathcal F\in\mathcal W(\mathcal A)$. Then
$\mathcal F\in\mathcal A$, so at every scale $x>x_1$ every member $D_{x,i}$ of $\mathcal F$ is the
datum of a selected bundle, whose divisibility data and twisted forms are those of $\mathcal F$.
Hence (1) and (2) hold at every member at every scale $x>x_1$ of every family of
$\mathcal W(\mathcal A)$. In particular they hold eventually along $\mathcal W(\mathcal A)$ at all
large scales, with $\varepsilon_1=1$, every largeness threshold and $x_0=0$.
\end{proof}

\begin{remark}\label{pd-integrality-inputs}
The construction of Chapter~\ref{chapter-assembly} uses the slots of good cores. These lie over true
columns: each good core is built at a row stage, so its period rows are integral by
Proposition~\ref{pd-rows-integral-built}, and its twisted forms, tie and jet conditions are those of
its true columns. By Theorem~\ref{pd-period-integrality-along-families}, period integrality and
twisted period integrality hold along these slots for the link data of the construction. The proof
of Theorem~\ref{asm-definite-side-exists} obtains in this way the period integrality and twisted
period integrality of item (10) of Definition~\ref{asm-definite-side-hypotheses}: for the class
$\mathcal A$ of the collection they hold eventually along $\operatorname{Des}(\mathcal A_N,3)$, which
is contained in $\mathcal W(\mathcal A)$. In the proof of
Theorem~\ref{asm-tied-from-green} they are hypotheses (3) and (4) of
Theorem~\ref{pd-multiplicity-from-integrality}.
\end{remark}
\leanref{Development/AbelianVarietyTheory/Proofs/GlobalInterpolation/GoodTwistCoresOverTrueColumnsStatement.lean}{GlobalTheory.T.goodTwistCores_over_trueColumns}
\leanref{Development/AutomorphicArithmetic/ThetaCorrespondence/ThetaPairLifts/ThirdPlaceSeries/SixRowsAtMemberStatement.lean}{GlobalTheory.T.period_entries_tie_and_level_valuations_at_built}

\section{Gains at inert primes}\label{section-pd-inert-gains}

\subsection{Gain data and the projection property}\label{subsection-pd-gain-data-projection}

\begin{definition}\label{pd-class-sum}
Let $D$ be a definite-side datum with size data, exponents $a_t$ and weights $w_{j,t,p,i}$
(Definition~\ref{pd-size-data}). For a lift $j$, a top class $t\in T$ and a prime $p$, the
\emph{class sum} is
\[
 \mathrm{VM}_{j,t}(p)=\sum_{\substack{0\leq i\leq\lfloor\log_pa_t\rfloor+1\\ i\ \mathrm{odd}}}
 \max\Big(0,\ \frac{a_t}{p^{i}}-w_{j,t,p,i}\Big).
\]
\end{definition}
\leanref{Formalism/LocalRingCardinality/DataParams₂.lean}{Formalism.BasicNotions.SizeData₂.vonMangoldtClassSum}

\begin{definition}\label{pd-gain-data}
Let $D$ be a definite-side datum with ranks $k,m,t,n$, Galois datum with group $A$ and generator
$a_0$ of $A_0$, host with top classes $T$, set $T_2$ and element $u_0\in A$
(Definition~\ref{hf-host}), and size data with scale $(x,\varepsilon,J)$, exponents $a_t,b_t$, width
$\eta$ and the derived quantities $b,a,b',w'',X',\delta_1,p_0$ of Definition~\ref{pd-size-data}. A
place $z$ of $D$ (Definition~\ref{pd-orbit-places}) is \emph{thick} if $z=(1,sa_0+u_0)$ for some
$s\in T_2$; then $\tau_z$ is a thick place of the host. Extend $t\mapsto a_t$ and $t\mapsto b_t$ by
$0$ to functions $\mathbf a,\mathbf b\colon A\to\mathbf Z$. For a prime $p$ let $h_p\in A$ be the
$A$-component of the Frobenius class of $p$ in $G_\Phi$ (which is $\operatorname{art}(p)$,
Definition~\ref{hf-framed-host-field}, when $p\nmid N$), let $n_p$ be its order, and let $s_p=2h_p$ if
$p$ is a signal prime and $s_p=h_p$ otherwise. \emph{Gain data} for $D$ consist of:
\begin{enumerate}
\item for each $h\in A$ a map $t\mapsto t^h$ on $A$ with $t\in t^h+\mathbf Nh$ and $(t+h)^h=t^h$,
  choosing a representative of each orbit of translation by $h$;
\item offsets $o_j(z)\in\mathbf R$ for integers $j\geq0$ and places $z\in\{0,1\}\times A$
  (Definition~\ref{pd-orbit-places}), used in the surpluses below and in the local gains of
  Definition~\ref{pd-served-selection}, such that for every $j$:
  $b'/x\leq o_j(z)\leq b'/x+\eta$ if $z$ is a top place; $b\leq o_j(z)\leq b+\eta+1$ if $z$ is
  thick; and $w''\leq o_j(z)\leq w''+\eta$ at every other place;
\item a function $\kappa_I(p,n',P,a',b')\in\mathbf R$ of primes $p$, integers $n'\geq1$, subsets
  $P\subset\mathbf Z/n'\mathbf Z$ and functions $a',b'\colon\mathbf Z/n'\mathbf Z\to\mathbf Z$;
\item a list $K_G$ of real constants, among them $K_\Lambda$, $K_\rho$, $K_\kappa$ and $K_N$;
\item for each prime $p$ an integer $e_p\geq0$, the \emph{Weyl exponent} at $p$.
\end{enumerate}
They determine the following.
\begin{enumerate}
\item[(a)] The \emph{representatives} at $p$ are the elements $t^{s_p}$, $t\in A$. For a
  representative $t_0$, $N(p,t_0)$ is the number of $t\in T$ with $t^{s_p}=t_0$, and the
  \emph{positions} of $t_0$ are the $i\in\{0,\dots,n_p-1\}$ with $t_0+ih_p\in T$; they form a set
  $P(p,t_0)\subset\mathbf Z/n_p\mathbf Z$.
\item[(b)] A representative $t_0$ is \emph{good} at $p$ if $P(p,t_0)$ is a tiled set of positions for
  $(x,\varepsilon,L_2)$ (Definition~\ref{hf-good-orbit}), where $L_2$ is twice the largest of the
  integers $\lceil\log_{p'}\big((kx^{26})^4\big)\rceil$ over the primes $p'$ of inert class.
\item[(c)] With $N_1$ and $r_\varepsilon$ as in Definition~\ref{hf-good-orbit}, put
  $\Lambda=K_\Lambda N_1^2\log a$, $N^*=2n\delta_1$ and
  $\rho=\big(4N^*x^\varepsilon\Lambda+K_\rho K_N\Lambda\log x\big)/k$. A prime $p$ is \emph{in band}
  if $p_0\leq p<X'$ and $X'(1-\rho)\leq p^e\leq X'(1+\rho)$ for some $e\geq1$. It is a
  \emph{candidate prime} if $p<X'$, $p$ is of inert class, $n_p\geq x^{\varepsilon/4}$, and $p$ is
  not in band.
\item[(d)] The \emph{local constant} is
  $\kappa(p;n',P,a',b')=\min\big(1,K_\kappa|P|r_\varepsilon^2k\,x^{-26}\big)$ if $p<p_0$, and
  $\kappa_I(p,n',P,a',b')$ if $p\geq p_0$. For a representative $t_0$ and
  $a',b'\colon A\to\mathbf Z$ put
  $\kappa_{p,t_0}(a',b')=\kappa\big(p;n_p,P(p,t_0),(a'(t_0+ih_p))_i,(b'(t_0+ih_p))_i\big)$.
\item[(e)] The \emph{deviation} at $p$ is $\Sigma(p)=m\sum_{t_0}\kappa_{p,t_0}(\mathbf a,\mathbf b)$,
  the sum over the good representatives at $p$. The set $V(p)$ of \emph{served representatives}
  is the set of good representatives at $p$ if $p$ is a candidate prime with
  $\Sigma(p)\leq x^{-\varepsilon}$, and is empty otherwise.
\item[(f)] For $t_0\in V(p)$, the \emph{surplus} $S(p,t_0)$ is a number computed from the offsets,
  the exponents $\mathbf a,\mathbf b$ and $K_G$ (Definition~\ref{pd-served-lattice}(3)); only the
  bound of Definition~\ref{pd-gain-link}(2) on it is used here.
\item[(g)] Let $W_a$ be the set of integers in $[a,\,a+a/(kx^{13})]$ and $W_b$ the set of integers in
  $[b',\,b'+N^*x^\varepsilon\Lambda b'/k]$. For $P\subset\mathbf Z/n'\mathbf Z$ and a set $W$ of
  integers let $B_P(W)$ be the set of functions $\mathbf Z/n'\mathbf Z\to\mathbf Z$ with values in
  $W$ on $P$ and $0$ off $P$. The \emph{mean} of a function $F$ over $P$ is its average over
  $(a',b')\in B_P(W_a)\times B_P(W_b)$.
\end{enumerate}
\end{definition}
\leanref{Formalism/SharpRestriction/GrowthBounds/IsResidueClassZero.lean}{Formalism.MainNotions.G.Inst}
\leanref{Formalism/SharpRestriction/GrowthBounds/IsResidueClassZero.lean}{Formalism.MainNotions.G.Consts}
\leanref{Formalism/SharpRestriction/GrowthBounds/IsResidueClassZero.lean}{Formalism.MainNotions.G.reps}
\leanref{Formalism/SharpRestriction/GrowthBounds/IsResidueClassZero.lean}{Formalism.MainNotions.G.Good}
\leanref{Formalism/SharpRestriction/GrowthBounds/IsResidueClassZero.lean}{Formalism.MainNotions.G.Cand}
\leanref{Formalism/SharpRestriction/GrowthBounds/IsResidueClassZero.lean}{Formalism.MainNotions.G.kapLoc}
\leanref{Formalism/SharpRestriction/GrowthBounds/IsResidueClassZero.lean}{Formalism.MainNotions.G.Sig}
\leanref{Formalism/SharpRestriction/GrowthBounds/IsResidueClassZero.lean}{Formalism.MainNotions.G.VV}
\leanref{Formalism/SharpRestriction/GrowthBounds/IsResidueClassZero.lean}{Formalism.MainNotions.G.meanLoc}

\begin{definition}\label{pd-gain-link}
Let $D$ be a definite-side datum with ranks $k,m,t,n$, $f=2n$ and $N(j)=k+j(t+1)$, and with size
data as in Definition~\ref{pd-size-data}. For $C,\beta\in\mathbf R$ and an inert prime $p$ put
\[
 \mathrm{fl}(p)=Cf\,N(m)^2\log_pa+Cf(\log_pa)^2,
\]
and $\mathrm{lo}(p)=C(S_c+k+\eta)\log_p(2a)+\mathrm{pr}(p)+\mathrm{or}(p)$, where
$\mathrm{pr}(p)=\mathrm{or}(p)=0$ for signal primes and, for primes of inert class,
\[
 \mathrm{pr}(p)=\begin{cases}
 C\,r(\varepsilon)\,b'\log x,&p<p_0,\\
 C\big(b'/x+\eta+S_c+(w''+t)\log x/\delta_1\big)+\beta,&p_0\leq p<X',\\
 0,&p\geq X',
 \end{cases}
\]
\[
 \mathrm{or}(p)=\begin{cases}
 C\log x\,\big(\Lambda_Ga/(pk)+a\,p^{-\ell(p)}\big),&p<p_0,\\
 C\log x\,\big(\Lambda_Ga/(pk)+a/p^2\big),&p_0\leq p<X',\\
 0,&p\geq X',
 \end{cases}
\]
with $r(\varepsilon)=\lceil8(4+8\varepsilon)/\varepsilon\rceil+1$ and
$\ell(p)=\lceil\log_p(k^2x^{78})\rceil+1$.

Let gain data for $D$ be given, with served representatives $V(p)$, surpluses $S(p,t_0)$ and
numbers $N(p,t_0)$ (Definition~\ref{pd-gain-data}). The \emph{gain link} holds for $D$ with
constants $C,\beta$ if:
\begin{enumerate}
\item for every word $u$ and every place $w$ with $p_w\neq T_u$: locally at $w$,
  $N^\circ_u\subset\mathfrak p_w^{r_w}M^{\mathrm{ram}}_u$ with $r_w=\max(0,\lfloor e_w\tau_w\rfloor)$,
  that is, for every $x\in N^\circ_u$ there is $\sigma\in\mathcal O_L\smallsetminus\mathfrak p_w$
  with $\sigma x\in\mathfrak p_w^{r_w}M^{\mathrm{ram}}_u$; moreover
  $\operatorname{len}_{u,w}\leq(\dim M_u)(\tau_w+\iota_w)$; and if
  $V(p_w)=\emptyset$ then $\tau_w=\iota_w=\operatorname{len}_{u,w}=0$;
\item for every prime $p$ and every $t_0\in V(p)$, $S(p,t_0)\leq2m\,N(p,t_0)\,a/p$;
\item for every place $w$ such that $p=p_w$ is inert and $p\leq X''$:
  \begin{enumerate}
  \item $0\leq\mathrm{gain}_j(w)\leq4|T|a/p$ for every lift $j$, and
  \item $\displaystyle\sum_{j=1}^{m}\sum_{t\in T_w}\big(\mathrm{VM}_{j,t}(p)-\mathrm{lo}(p)\big)
    -m\,\mathrm{fl}(p)\leq\sum_{j=1}^{m}\mathrm{gain}_j(w)-\tau_w$;
  \end{enumerate}
\item for every candidate prime $p$:
  \begin{enumerate}
  \item if $\Sigma(p)\leq x^{-\varepsilon}$, then $V(p)$ contains every good representative at $p$;
  \item for every good representative $t_0$ at $p$, the mean over $P(p,t_0)$ of
    $(a',b')\mapsto\kappa\big(p;n_p,P(p,t_0),a',b'\big)$ is at most $N(p,t_0)\,x^{-12}$; and if
    $p<p_0$, then $\kappa_{p,t_0}(a',b')\leq N(p,t_0)\,x^{-12}$ for all
    $a',b'\colon A\to\mathbf Z$.
  \end{enumerate}
\end{enumerate}
A family (Definition~\ref{pd-family}) satisfies the \emph{gain link} if there are $C$, $x_0$ and a
function $\delta$ with $\delta(x)\to0$ as $x\to\infty$ such that for every $x\geq x_0$ and every
member $i$ at $x$, the gain link holds for $D_{x,i}$ with the gain data of the family and constants
$C$ and $\beta=b^{\delta(x)}$, where $b$ is that of the size data $L_{x,i}$.
\end{definition}
\leanref{Formalism/SharpRestriction/GlobalTheoremInputs.lean}{Formalism.MainNotions.LinkG}
\leanref{Formalism/SharpRestriction/GlobalTheoremInputs.lean}{Formalism.MainNotions.LinkG_at}
\leanref{Formalism/SharpRestriction/GlobalTheoremInputs.lean}{Formalism.MainNotions.loG}
\leanref{Formalism/SharpRestriction/GlobalTheoremInputs.lean}{Formalism.MainNotions.flG}
\leanref{Formalism/SharpRestriction/GlobalTheoremInputs.lean}{Formalism.MainNotions.InertK}

Condition~(4a) holds for all gain data, by the definition of $V(p)$.
Condition~(3b) is the form in which the divisibility of the periods enters the inert-prime sum of
Chapter~\ref{chapter-inert-prime-balance}: summed over the places over $p$ with weight
$2\log p$, its left side is the inert-prime sum of the balance identity, up to the errors
$\mathrm{lo}$ and $\mathrm{fl}$.

\begin{definition}\label{pd-projection-property}
Let $p\geq5$ be a prime and $(W,Q,\sigma,\iota)$ a Hida setting of degree $n'$ at $p$ with rings
$\mathcal R\subset\mathcal K$ and $\varpi_0$ (Definition~\ref{pd-local-model}(1)). The setting is
\emph{unramified quadratic} if $\varpi_0$ is a uniformizer of $\mathcal R$, the residue field of
$\mathcal R$ has $p^{2n'}$ elements, and every integer is a square in $\mathcal R$. Let $k\geq1$ and
$r\geq0$, let $\lambda_0(s)\in\mathbf Z^k$, $s\in\mathbf Z/n'\mathbf Z$, be dominant weights, and
define $\lambda(s)$ from $\lambda_0(s)$ as in Definition~\ref{pd-local-model}. Let
$V\subset\mathcal K[x_{s,a,b}]$ be the span of the products $\prod_sf_s(x_s)$ with $f_s$ in the dual
Weyl module over $\mathcal K$ of the natural weight of $\lambda(s)$
(Definition~\ref{pd-weight-lattices}(1)), and $V^\vee$ its dual.
\begin{enumerate}
\item For a $W$-lattice $L\subset Q^k$, let $M_s$ be the $\mathcal R$-span of $\iota\sigma^s(L)$
  and $M_s^\vee=\{\phi\in\mathcal K^k:\phi\cdot v\in\mathcal R\text{ for all }v\in M_s\}$ its dual.
  Then $H(L)\subset V$ is the $\mathcal R$-span of the products $\prod_sf_s(x_s)$ with
  $f_s\in\nabla_{\lambda(s)}(M_s^\vee)$ (Definition~\ref{pd-weight-lattices}(5), formed over
  $\mathcal R$). Equivalently, $f_s\in\nabla_{\lambda(s)}(M_s^\vee)$ if and only if
  $\det(g)^{\lambda(s)_k}f_s(gx_s)$ lies in the dual Weyl module over $\mathcal R$ of the natural
  weight of $\lambda(s)$ for every $g$ whose columns form an $\mathcal R$-basis of $M_s$: the lattice
  is moved from the one of the standard basis by $f\mapsto\det(g)^{-\lambda(s)_k}f(g^{-1}x)$, which
  is the action of Definition~\ref{pd-local-model}(3) and is $\rho_{\lambda(s)}({}^tg^{-1})$ in the
  notation of Definition~\ref{pd-weight-lattices}(4). An
  \emph{edge} $e$ is a pair of $W$-lattices $L''\subset L$ with $pL\subset L''$ and
  $|L/L''|=p^{n'r}$, with \emph{source} $L''$ and \emph{target} $L$; $H(e)\subset V$ is defined as
  $H(L)$ with $M_s$ replaced by the $\mathcal R$-span of $\iota\sigma^s(L'')$ and
  $\varpi_0\,\iota\sigma^s(L)$. Put $A=\sum_s\sum_{i\leq r}\lambda(s)_i$,
  $R(e)=\varpi_0^{-A}H(e)$ and $R'(e)=\varpi_0^{-A}H(e)^\vee$, where $H^\vee\subset V^\vee$ consists of
  the forms that take $H$ into $\mathcal R$.
\item For $t\geq0$, $N_{\mathrm{src}}(L'')$ is $\varpi_0^{-t}H(L'')$ intersected with the $R(e)$ over
  the edges $e$ with source $L''$, and $N_{\mathrm{tgt}}(L)$ is $\varpi_0^{-t}H(L)^\vee$ intersected
  with the $R'(e)$ over the edges $e$ with target $L$. The \emph{edge pairing} is
  $\langle F,G\rangle=G(p^{A}F)$ for $F\in V$ and $G\in V^\vee$.
\item Let $t,D,c_*\geq0$ be integers and $\kappa\in\mathbf R$. The \emph{projection property}
  $\mathrm{LP}(r,\lambda_0;t,D,\kappa,c_*)$ holds if $4(D+1)\leq c_*$ and for every $W$-lattice
  $L\subset Q^k$ there are edges $e_1,\dots,e_M$ with target $L$ and vectors
  $F_i\in N_{\mathrm{src}}(\text{source of }e_i)$ such that $F_i\in\varpi_0^{c_*}R(e')$ for every edge
  $e'\neq e_i$ with the same source as $e_i$, and there are $n\geq(1-\kappa)\dim V$, an
  $\mathcal R$-linear map $u\colon\mathcal R^n\to N_{\mathrm{tgt}}(L)$ and an $n\times M$ matrix
  $(B_{ji})$ over $\mathcal R$ with $\sum_iB_{ji}\langle F_i,u(y)\rangle=p^Dy_j$ for all
  $y\in\mathcal R^n$ and $j\leq n$.
\end{enumerate}
\end{definition}
\leanref{Formalism/SharpRestriction/GrowthBounds/UnramifiedQuadraticSetting.lean}{Formalism.MainNotions.G.UnramifiedQuadraticSetting}
\leanref{Formalism/SharpRestriction/GrowthBounds/UnramifiedQuadraticSetting.lean}{Formalism.MainNotions.G.HL}
\leanref{Formalism/SharpRestriction/GrowthBounds/UnramifiedQuadraticSetting.lean}{Formalism.MainNotions.G.IsEdge}
\leanref{Formalism/SharpRestriction/GrowthBounds/UnramifiedQuadraticSetting.lean}{Formalism.MainNotions.G.Re}
\leanref{Formalism/SharpRestriction/GrowthBounds/UnramifiedQuadraticSetting.lean}{Formalism.MainNotions.G.Nsrc}
\leanref{Formalism/SharpRestriction/GrowthBounds/UnramifiedQuadraticSetting.lean}{Formalism.MainNotions.G.Ntgt}
\leanref{Formalism/SharpRestriction/GrowthBounds/UnramifiedQuadraticSetting.lean}{Formalism.MainNotions.G.GoodRankK}
\leanref{Formalism/SharpRestriction/GrowthBounds/UnramifiedQuadraticSetting.lean}{Formalism.MainNotions.G.LPprime}

\begin{lemma}\label{pd-extremal-projector}
Let $K$ be a field of characteristic $0$, $n\geq1$, and $V$ a finite-dimensional $K$-vector space
with endomorphisms $E_{ab}$, $1\leq a,b\leq n$, such that
$E_{ab}E_{cd}-E_{cd}E_{ab}=\delta_{bc}E_{ad}-\delta_{da}E_{cb}$ for all $a,b,c,d$ and $V$ is the sum
of the weight spaces $V_\mu=\{x\in V:E_{aa}x=\mu_ax\text{ for all }a\}$, $\mu\in\mathbf Z^n$. Let
$V^-\subset V$ be the sum of the images of the $E_{ab}$ with $a>b$, and put
$E^{(k)}_{ab}=E_{ab}^k/k!$. Let $\mu\in\mathbf Z^n$ satisfy $h_{ab}=\mu_a-\mu_b+b-a\geq1$ for all
$a<b$, let $T\geq\dim V$ be an integer, and for $a<b$ put
\[
 p_{ab}=\sum_{k=0}^{T}\frac{(-1)^k}{\binom{h_{ab}+k}{k}}\,E^{(k)}_{ba}E^{(k)}_{ab}.
\]
Let $(a_1,b_1),\dots,(a_N,b_N)$ be the pairs $a<b$, each listed once, in an order in which the pair
$(a,c)$ lies between $(a,b)$ and $(b,c)$ whenever $a<b<c$, and put
$p_\mu=p_{a_1b_1}p_{a_2b_2}\cdots p_{a_Nb_N}$. Then:
\begin{enumerate}
\item $p_\mu(V_\mu)\subset V_\mu$;
\item $p_\mu x=0$ for every $x\in V_\mu\cap V^-$;
\item $p_\mu x=x$ for every $x\in V_\mu$ with $E_{ab}x=0$ for all $a<b$.
\end{enumerate}
On $V_\mu$, $p_\mu$ is the extremal projector of Asherova, Smirnov and Tolstoy \cite{AST71}.
\end{lemma}
\leanref{Formalism/SharpRestriction/GrowthBounds/CovariantDegrees.lean}{Formalism.MainNotions.CovariantDegrees.adm5}
\leanref{Formalism/SharpRestriction/GrowthBounds/ObjPol.lean}{Formalism.MainNotions.G.ppFac}
\leanref{Formalism/SharpRestriction/GrowthBounds/ObjPol.lean}{Formalism.MainNotions.G.pp}
\leanref{Formalism/SharpRestriction/GrowthBounds/ObjPol.lean}{Formalism.MainNotions.G.NormalOrder}
\leanref{Formalism/SharpRestriction/GrowthBounds/TruncatedChains.lean}{Formalism.MainNotions.G.ExtremalProjector.extremal_projector_statement}

\subsection{Path generators}\label{subsection-pd-path-generators}

\begin{definition}\label{pd-symbols-path-generators}
Let $N\geq0$ and $n\geq1$ be integers. A \emph{lettering} of $\{0,\dots,N-1\}$ gives each
\emph{letter} $i$ a type, \emph{vector} or \emph{covector}. Let $R$ be a commutative ring, $R[z]$ the
polynomial ring in $z_{i,b}$, $0\leq i<N$, $1\leq b\leq n$, and $z_i=(z_{i,1},\dots,z_{i,n})$. A
polynomial has \emph{row degrees} $B\colon\{0,\dots,N-1\}\to\mathbf Z_{\geq0}$ if each of its
monomials has degree $B(i)$ in $z_i$ for every $i$.
\begin{enumerate}
\item Let $\mathfrak g$ be the $n\times n$ matrix of indeterminates over $R$, $\eta$ a further
  indeterminate, $D_a(\eta)$ the diagonal matrix with $\eta$ in place $a$ and $1$ elsewhere, and
  $\tau_a(\eta)=\mathfrak gD_a(\eta)\mathfrak g^{-1}$. For $S\subset\{0,\dots,N-1\}$ and $P\in R[z]$
  let $P_{S,a}$ be obtained from $P$ by replacing, for every $i\in S$, the row $z_i$ by
  $z_i\tau_a(\eta)$ if $i$ is a vector letter and by $z_i\,{}^t\tau_a(\eta^{-1})$ if it is a covector
  letter. Then $\operatorname{val}_S(P)$ is the least exponent of $\eta$ that occurs in $P_{S,n}$ and
  $\lambda_S(P)$ the largest that occurs in $P_{S,1}$, with $\operatorname{val}_S(0)=+\infty$ and
  $\lambda_S(0)=-\infty$.
\item \emph{Paths.} Let $B\colon\{0,\dots,N-1\}\to\mathbf Z_{\geq0}$. A \emph{path} of rank $n$ for
  $B$ is a sequence $\mu=(\mu_0,\dots,\mu_N)$ of vectors $\mu_i\in\mathbf Z^n$ with
  $\mu_i(1)\geq\dots\geq\mu_i(n)$ and $\mu_0=0$ such that for every letter $i$, with
  $(\alpha,\beta)=(\mu_i,\mu_{i+1})$ if $i$ is a vector letter and $(\alpha,\beta)=(\mu_{i+1},\mu_i)$ if
  it is a covector letter, $\alpha(a)\leq\beta(a)$ for all $a$, $\beta(a+1)\leq\alpha(a)$ for $a<n$, and
  $\sum_a(\beta(a)-\alpha(a))=B(i)$.
\item \emph{Column functions.} For a sequence $\mu=(\mu_0,\dots,\mu_N)$ in $\mathbf Z^n$ and
  $c\in\mathbf Z$ let $h_i(c)$ be the number of $a$ with $\mu_i(a)\geq c$, and put
  $T_\mu=\sum_i\sum_a|\mu_i(a)|$. For a row $x=(x_1,\dots,x_n)$ in a commutative ring $A$ and a
  function $F$ from the subsets of $\{1,\dots,n\}$ to $A$ put
  \[
   (F\wedge x)(J)=\sum_{m\in J}\epsilon_J(m)\,x_mF(J\smallsetminus\{m\}),\qquad
   (F\lrcorner x)(J)=\sum_{m\notin J}\epsilon_J(m)\,x_mF(J\cup\{m\}),
  \]
  with $\epsilon_J(m)=(-1)^{|\{l\in J:l>m\}|}$. The \emph{step} at the letter $i$ for $c$ sends $(F,x)$
  to $F\wedge x$ if $h_i(c)<h_{i+1}(c)$ and $c\geq1$ or $h_{i+1}(c)<h_i(c)$ and $c\leq0$, to
  $F\lrcorner x$ if $h_i(c)<h_{i+1}(c)$ and $c\leq0$ or $h_{i+1}(c)<h_i(c)$ and $c\geq1$, and to $F$ if
  $h_i(c)=h_{i+1}(c)$. The \emph{column function} $\mathrm{Col}_{\mu,c}$ is $F_N$, where $F_0(J)=1$ for
  $J=\emptyset$ and $F_0(J)=0$ otherwise, and $F_{i+1}$ is the step at $i$ for $c$ applied to
  $(F_i,z_i)$.
\item \emph{Path generators.} A \emph{standard filling} of $\mu$ is a function $I$ from $\mathbf Z$ to
  the subsets of $\{1,\dots,n\}$ such that $I(c)=\emptyset$ for $|c|>T_\mu$, $|I(c)|=h_N(c)$ for
  $c\geq1$ and $|I(c)|=n-h_N(c)$ for $c\leq0$, and, with $I^*(c)=I(c)$ for $c\geq1$ and $I^*(c)$ the
  complement of $I(c)$ for $c\leq0$, the set $I^*(c+1)$ has at most as many elements $\leq t$ as
  $I^*(c)$, for all $c\in\mathbf Z$ and $1\leq t\leq n$. For any function $I$ from $\mathbf Z$ to the
  subsets of $\{1,\dots,n\}$ put $\Psi(\mu,I)=\prod_{|c|\leq T_\mu}\mathrm{Col}_{\mu,c}(I(c))$; when $\mu$ is a
  path and $I$ a standard filling of it, $\Psi(\mu,I)$ is a \emph{path generator}.
\item \emph{Splitting.} Let $n=r+K$ with $r,K\geq1$. For $\mu$ as in (3) let $\mu^{\mathrm{pl}}_i$
  and $\mu^{\mathrm{tn}}_i$ be the first $r$ and the last $K$ entries of $\mu_i$, and for sequences
  $\nu$ in $\mathbf Z^r$ and $\kappa$ in $\mathbf Z^K$ of length $N+1$ let $(\nu,\kappa)$ be the
  sequence of the vectors $(\nu_i,\kappa_i)\in\mathbf Z^n$. Let $\iota_{\mathrm{pl}}$ and
  $\iota_{\mathrm{tn}}$ be the $R$-algebra maps to $R[z]$ from the polynomial rings in the $z_{i,b}$
  with $b\leq r$, respectively $b\leq K$, with $z_{i,b}\mapsto z_{i,b}$, respectively
  $z_{i,b}\mapsto z_{i,r+b}$. For $c\in\mathbf Z$ and $J\subset\{1,\dots,n\}$ put
  \[
   \mathrm{Col}^{\mathrm{pt}}_{\mu,c}(J)=\iota_{\mathrm{pl}}\big(\mathrm{Col}_{\mu^{\mathrm{pl}},c}(J\cap\{1,\dots,r\})\big)\,
   \iota_{\mathrm{tn}}\big(\mathrm{Col}_{\mu^{\mathrm{tn}},c}(\{b:r+b\in J\})\big).
  \]
  For $\omega\colon\{0,\dots,N-1\}\to\mathbf Z_{\geq0}$, the \emph{$\omega$-degree} of $z_{i,b}$ is
  $\omega(i)$ if $i$ is a vector letter and $b>r$ or a covector letter and $b\leq r$, and $0$
  otherwise, and that of a monomial is the sum over its variables, counted with multiplicity. A letter
  $i$ is \emph{used} for $c$ if $h_i(c)\neq h_{i+1}(c)$, and \emph{costly} for $c$ if
  $h_i(c)<h_{i+1}(c)$ and $h_i(c)\geq r$, or $h_{i+1}(c)<h_i(c)$ and $h_i(c)\leq r$. Let
  $o_\omega(\mu,c)$ be the sum of the $\omega(i)$ over the letters $i$ costly for $c$.
\item \emph{Lifted column functions.} Let $p\geq5$ be a prime and $(W,Q,\sigma,\iota)$ a Hida setting
  of degree $n'$ at $p$ with rings $\mathcal R\subset\mathcal K$ and $\varpi_0$
  (Definition~\ref{pd-local-model}), let each letter $i$ have a \emph{position} $\pi_i\in\mathbf Z$, and
  for $g\in GL_n(W)$ and $v\in\mathbf Z$ let $g_v=\iota\sigma^v(g)$. Put $\ell_i=z_ig_{\pi_i}$ if $i$
  is a vector letter and $\ell_i=z_i\,{}^tg_{\pi_i}^{-1}$ if it is a covector letter, rows of linear
  forms over $\mathcal K$. For $\mu$ as in (3), $c\in\mathbf Z$ and an integer $\zeta$, let
  $i_1<\dots<i_s$ be the letters used for $c$, let $F_0$ be as in (3), and for $1\leq t\leq s$ let
  $F_t$ be the step at $i_t$ for $c$ applied to $(F_{t-1}^{\,e_t},\ell_{i_t})$, where
  $F^{\,e}(J)=F(J)^e$, $e_1=1$ and $e_t=p^{\pi_{i_t}-\pi_{i_{t-1}}}$ for $t\geq2$. The \emph{lifted
  column function} is $\mathrm{Col}^g_{\mu,c,\zeta}=F_s^{\,p^{\zeta-\pi_{i_s}}}$ if $s\geq1$ and $F_0$ if $s=0$.
  An exponent of $p$ in this item that would be negative is replaced by $0$.
\item \emph{Substitutions.} In the setting of (6), with $n=r+K$ as in (5), let $\delta$ and $\delta'$
  be the diagonal matrices
  with entries $(1,\dots,1,\varpi_0,\dots,\varpi_0)$ and $(\varpi_0,\dots,\varpi_0,1,\dots,1)$, with $r$
  entries before the change. The \emph{lettered substitution} $\Sigma^{\mathrm L}_g$ is the
  $\mathcal K$-algebra endomorphism of $\mathcal K[z]$ with $z_i\mapsto z_i\,\delta\,g_{\pi_i}^{-1}$ for
  vector letters and $z_i\mapsto z_i\,\delta'\,{}^tg_{\pi_i}$ for covector letters, and for
  $\zeta\in\mathbf Z$, $\Pi_\zeta\colon\mathcal R[z]\to\mathcal K[z]$ is the $\mathcal R$-algebra map with
  $z_{i,b}\mapsto z_{i,b}^{\,p^{\zeta-\pi_i}}$. For $i\geq0$, a polynomial lies in $\varpi_0^i\mathcal
  R[z]$ if all its coefficients lie in $\varpi_0^i\mathcal R$.
\end{enumerate}
\end{definition}
\leanref{Formalism/SharpRestriction/GrowthBounds/ObjPol.lean}{Formalism.MainNotions.G.Letters}
\leanref{Formalism/SharpRestriction/GrowthBounds/ObjPol.lean}{Formalism.MainNotions.G.LiftLetters}
\leanref{Formalism/SharpRestriction/GrowthBounds/ObjPol.lean}{Formalism.MainNotions.G.Pol}
\leanref{Formalism/SharpRestriction/GrowthBounds/ObjPol.lean}{Formalism.MainNotions.G.valS}
\leanref{Formalism/SharpRestriction/GrowthBounds/ObjPol.lean}{Formalism.MainNotions.G.lam1S}
\leanref{Formalism/SharpRestriction/GrowthBounds/ObjPol.lean}{Formalism.MainNotions.G.HStrip}
\leanref{Formalism/SharpRestriction/GrowthBounds/ObjPol.lean}{Formalism.MainNotions.G.PPath}
\leanref{Formalism/SharpRestriction/GrowthBounds/ObjPol.lean}{Formalism.MainNotions.G.colH}
\leanref{Formalism/SharpRestriction/GrowthBounds/ObjPol.lean}{Formalism.MainNotions.G.wedgeR}
\leanref{Formalism/SharpRestriction/GrowthBounds/ObjPol.lean}{Formalism.MainNotions.G.contrR}
\leanref{Formalism/SharpRestriction/GrowthBounds/ObjPol.lean}{Formalism.MainNotions.G.colStep}
\leanref{Formalism/SharpRestriction/GrowthBounds/ObjPol.lean}{Formalism.MainNotions.G.Kcol}
\leanref{Formalism/SharpRestriction/GrowthBounds/ObjPol.lean}{Formalism.MainNotions.G.colRange}
\leanref{Formalism/SharpRestriction/GrowthBounds/ObjPol.lean}{Formalism.MainNotions.G.IsSt}
\leanref{Formalism/SharpRestriction/GrowthBounds/ObjPol.lean}{Formalism.MainNotions.G.Psi}
\leanref{Formalism/SharpRestriction/GrowthBounds/ObjPol.lean}{Formalism.MainNotions.G.joinPT}
\leanref{Formalism/SharpRestriction/GrowthBounds/ObjPol.lean}{Formalism.MainNotions.G.PsiPl}
\leanref{Formalism/SharpRestriction/GrowthBounds/ObjPol.lean}{Formalism.MainNotions.G.PsiTan}
\leanref{Formalism/SharpRestriction/GrowthBounds/ObjPol.lean}{Formalism.MainNotions.G.plTanCol}
\leanref{Formalism/SharpRestriction/GrowthBounds/ObjPol.lean}{Formalism.MainNotions.G.symWt}
\leanref{Formalism/SharpRestriction/GrowthBounds/ObjPol.lean}{Formalism.MainNotions.G.monWt}
\leanref{Formalism/SharpRestriction/GrowthBounds/ObjPol.lean}{Formalism.MainNotions.G.IsUsed}
\leanref{Formalism/SharpRestriction/GrowthBounds/ObjPol.lean}{Formalism.MainNotions.G.IsCostly}
\leanref{Formalism/SharpRestriction/GrowthBounds/ObjPol.lean}{Formalism.MainNotions.G.oc}
\leanref{Formalism/SharpRestriction/GrowthBounds/ObjPol.lean}{Formalism.MainNotions.G.letG}
\leanref{Formalism/SharpRestriction/GrowthBounds/ObjPol.lean}{Formalism.MainNotions.G.KKstate}
\leanref{Formalism/SharpRestriction/GrowthBounds/ObjPol.lean}{Formalism.MainNotions.G.KKc}
\leanref{Formalism/SharpRestriction/GrowthBounds/ObjPol.lean}{Formalism.MainNotions.G.liftSubst}
\leanref{Formalism/SharpRestriction/GrowthBounds/ObjPol.lean}{Formalism.MainNotions.G.powSubst}

\begin{definition}\label{pd-symbols-path-weights}
Let $N\geq0$ and $n\geq1$ be integers with a lettering of $\{0,\dots,N-1\}$, $R$ a commutative ring
and $S\subset\{0,\dots,N-1\}$; we use the notation of Definition~\ref{pd-symbols-path-generators}, and
write the monomials of $R[z]$ as $z^m=\prod_{i,b}z_{i,b}^{m(i,b)}$.
\begin{enumerate}
\item The \emph{$S$-weight} of $z^m$ is the vector $\mathrm{wt}_S(m)\in\mathbf Z^n$ whose entry at $a$
  is the sum of the $m(i,a)$ over the vector letters $i\in S$ minus the sum of the $m(i,a)$ over the
  covector letters $i\in S$.
\item For $a\neq b$ in $\{1,\dots,n\}$ and $k\geq0$, $E^{(k)}_{ab}$ is the $R$-linear endomorphism of
  $R[z]$ with $E^{(k)}_{ab}(z^m)=\sum_f\prod_{i\in S}c_i(f)\,z^{m_f}$, the sum over the functions
  $f\colon S\to\mathbf Z_{\geq0}$ with $\sum_if(i)=k$. For a vector letter $i$,
  $c_i(f)=\binom{m(i,b)}{f(i)}$ and $m_f$ moves $f(i)$ from the exponent of $z_{i,b}$ to that of
  $z_{i,a}$; for a covector letter $i$, $c_i(f)=(-1)^{f(i)}\binom{m(i,a)}{f(i)}$ and $m_f$ moves $f(i)$
  from the exponent of $z_{i,a}$ to that of $z_{i,b}$. Over $\mathbf Q$ it is $E_{ab}^k/k!$, where
  $E_{ab}=\sum_iz_{i,a}\,\partial/\partial z_{i,b}-\sum_{i'}z_{i',b}\,\partial/\partial z_{i',a}$, the
  sums over the vector letters $i$ and the covector letters $i'$ in $S$. For $a\leq n$ and $k\geq0$,
  $H^{(k)}_a$ is the $R$-linear endomorphism with $H^{(k)}_a(z^m)=\binom{\mathrm{wt}_S(m)(a)}{k}z^m$,
  where $\binom xk=x(x-1)\cdots(x-k+1)/k!$ for $x\in\mathbf Z$.
\item For $G\in R[z]$, $U_S(G)$ is the $R$-span of the $u(G)$, with $u$ a composite of finitely many of
  the operators $E^{(k)}_{ab}$ and $H^{(k)}_a$, the identity included. A \emph{weight} of $U_S(G)$ is a
  $\mu\in\mathbf Z^n$ such that $U_S(G)$ contains a nonzero element all of whose monomials have
  $S$-weight $\mu$.
\item For $\mu,\nu\in\mathbf Z^n$ write $\nu\preceq\mu$ if $\sum_{b\leq a}\nu(b)\leq\sum_{b\leq a}\mu(b)$
  for every $a\leq n$, with equality for $a=n$. A \emph{lower set} is a set $\pi$ of dominant vectors
  in $\mathbf Z^n$ that contains every dominant $\nu$ with $\nu\preceq\mu$ for some $\mu\in\pi$.
\end{enumerate}
\end{definition}
\leanref{Formalism/SharpRestriction/GrowthBounds/ObjPol.lean}{Formalism.MainNotions.G.wtS}
\leanref{Formalism/SharpRestriction/GrowthBounds/ObjPol.lean}{Formalism.MainNotions.G.divPow}
\leanref{Formalism/SharpRestriction/GrowthBounds/ObjPol.lean}{Formalism.MainNotions.G.binomPol}
\leanref{Formalism/SharpRestriction/GrowthBounds/ObjPol.lean}{Formalism.MainNotions.G.polUZgens}
\leanref{Formalism/SharpRestriction/GrowthBounds/ObjPol.lean}{Formalism.MainNotions.G.polGenSub}
\leanref{Formalism/SharpRestriction/GrowthBounds/ObjPol.lean}{Formalism.MainNotions.G.polWts}
\leanref{Formalism/SharpRestriction/GrowthBounds/StatementOfDegreeAndBound.lean}{Formalism.MainNotions.CovariantDegrees.DomLE}
\leanref{Formalism/SharpRestriction/GrowthBounds/StatementOfDegreeAndBound.lean}{Formalism.MainNotions.CovariantDegrees.DomLowerSet}

\begin{lemma}\label{pd-symbols-path-count}
Let $N\geq0$ and $n\geq1$ be integers with a lettering of $\{0,\dots,N-1\}$, and
$B\colon\{0,\dots,N-1\}\to\mathbf Z_{\geq0}$; we use the notation of
Definition~\ref{pd-symbols-path-generators}. Then there are finitely many pairs $(\mu,I)$ of a path $\mu$
of rank $n$ for $B$ and a standard filling $I$ of $\mu$, and their number is
$\prod_{i=0}^{N-1}\binom{n+B(i)-1}{B(i)}$, which is the number of monomials in the $z_{i,b}$ of row
degrees $B$.
\end{lemma}
\leanref{Formalism/ThetaLogCoordinates/JacobiPathWeights/FinalP.lean}{Formalism.InvariantFactorPaths.GelfandTsetlinBases.Top.countPAll}
\leanref{Formalism/ThetaLogCoordinates/JacobiPathWeights/PathBasisTheorem.lean}{Formalism.InvariantFactorPaths.GelfandTsetlinBases.Top.pathIdxFinite}
\leanref{Formalism/ThetaLogCoordinates/JacobiPathWeights/PathBasisTheorem.lean}{Formalism.InvariantFactorPaths.GelfandTsetlinBases.Top.polBRank}
\leanref{Formalism/ThetaLogCoordinates/JacobiPathWeights/JacobiMinorProofs/NormIndexCountBound.lean}{Formalism.InvariantFactorPaths.GelfandTsetlinBases.Count.norm_index_count_bound}

\begin{proof}
For a dominant $\nu\in\mathbf Z^n$ let $h_\nu(c)$ be the number of $a$ with $\nu(a)\geq c$. The
conditions of Definition~\ref{pd-symbols-path-generators}(4) on a standard filling $I$ of a path $\mu$
involve only $h_{\mu_N}$, apart from $I(c)=\emptyset$ for $|c|>T_\mu$, which they imply as
$T_\mu\geq|\mu_N(a)|$ for all $a$. We call such an $I$ a standard filling of $\mu_N$. Let $d_n(\nu)$ be
the number of Gelfand--Tsetlin patterns with top row $\nu$: $d_0=1$, and $d_n(\nu)=\sum_\kappa
d_{n-1}(\kappa)$ over the $\kappa\in\mathbf Z^{n-1}$ with $\nu(a+1)\leq\kappa(a)\leq\nu(a)$ for $a<n$. Two
facts are used.

(i) \emph{A dominant $\nu$ has $d_n(\nu)$ standard fillings.} For a standard filling $I$, the set
$I^*(c)$ has $h_\nu(c)$ elements, so it is $\{1,\dots,n\}$ for $c\leq\nu(n)$ and empty for $c>\nu(1)$,
and $I$ is determined by the $I^*(c)$ with $\nu(n)<c\leq\nu(1)$. Write the elements of such an $I^*(c)$
in increasing order down the column $c$ of the diagram of $\nu-\nu(n)(1,\dots,1)$. The condition that
$I^*(c+1)$ has at most as many elements $\leq t$ as $I^*(c)$, for every $t$, says that the $k$-th
element of $I^*(c)$ is at most the $k$-th element of $I^*(c+1)$. So the standard fillings of $\nu$
correspond to the semistandard tableaux of shape $\nu-\nu(n)(1,\dots,1)$ with entries in
$\{1,\dots,n\}$. Recording, for $k=n,n-1,\dots,0$, the shape formed by the entries at most $k$ identifies
these tableaux with the sequences of partitions in which each is obtained from the next by removing a
horizontal strip \cite[Chapter~I, \S1]{Mac95}, that is, with the Gelfand--Tsetlin patterns with top row
$\nu-\nu(n)(1,\dots,1)$; and $d_n$ does not change when a constant vector is added.

(ii) \emph{Pieri's formula.} For dominant $\nu$ and $b\geq0$, the sum of the $d_n(\nu')$ over the
$\nu'$ that are strips of size $b$ over $\nu$ (Definition~\ref{pd-symbols-path-generators}(2)) is
$d_n(\nu)\binom{n+b-1}b$, and so is the sum over the $\nu'$ over which $\nu$ is a strip of size $b$. By
(i), $d_n(\nu)$ is the number of semistandard tableaux of shape $\nu-\nu(n)(1,\dots,1)$ with entries in
$\{1,\dots,n\}$, which is the value at $(1,\dots,1)$ of the corresponding Schur polynomial in $n$
variables, and $\binom{n+b-1}b$ is the value there of the complete symmetric polynomial $h_b$. Pieri's
formula for the product of a Schur polynomial and $h_b$ \cite[Chapter~I, \S5]{Mac95}, evaluated at
$(1,\dots,1)$, gives the first sum, as the strips of size $b$ over $\nu$ are the $\nu'$ for which
$\nu'-\nu(n)(1,\dots,1)$ is obtained from $\nu-\nu(n)(1,\dots,1)$ by adding a horizontal strip of $b$
boxes in the first $n$ rows. The second sum is the first one for $\nu^*=(-\nu(n),\dots,-\nu(1))$: $\nu$
is a strip of size $b$ over $\nu'$ if and only if $\nu'^*$ is one over $\nu^*$, and
$d_n(\nu'^*)=d_n(\nu')$, by reversing and negating the rows of the patterns.

For $0\leq i\leq N$ and a dominant $\nu$ let $C_i(\nu)$ be the number of pairs of a sequence
$(\mu_i,\dots,\mu_N)$ with $\mu_i=\nu$ that satisfies the conditions of
Definition~\ref{pd-symbols-path-generators}(2) at the letters $i,\dots,N-1$, and a standard filling of
$\mu_N$. We show $C_i(\nu)=d_n(\nu)\prod_{l=i}^{N-1}\binom{n+B(l)-1}{B(l)}$ by descending induction on
$i$. For $i=N$ this is (i). For $i<N$, sorting by $\mu_{i+1}$ gives $C_i(\nu)=\sum_{\nu'}C_{i+1}(\nu')$, the
sum over the strips $\nu'$ of size $B(i)$ over $\nu$ if $i$ is a vector letter, and over the $\nu'$ over
which $\nu$ is such a strip if $i$ is a covector letter. These $\nu'$ are finite in number and dominant:
in the first case $\nu'(a+1)\leq\nu(a)\leq\nu'(a)$, and in the second $\nu'(a)\geq\nu(a+1)\geq\nu'(a+1)$.
The induction hypothesis and (ii) give the formula. The lemma is the case $i=0$, $\nu=0$, as $d_n(0)=1$;
and there are $\binom{n+B(i)-1}{B(i)}$ monomials of degree $B(i)$ in $z_{i,1},\dots,z_{i,n}$.
% Lean: countPAll_of countCont with countCont = countCont_of stripsMem gtPieri stdCountGT
%   (Count.norm_index_count_bound, by induction on N - i). Fact (i) is StdCountGT (stdCountGT, with dGT the
%   Gelfand-Tsetlin count by gtBox) and fact (ii) is GTPieri (gtPieri, with stripsUp, stripsDown and
%   StripsMem); the Lean proves both combinatorially. The monomial count is polynomial_b_rank.
\end{proof}

\begin{definition}\label{pd-symbols-path-canonical}
Let $n\geq1$ be an integer. For a dominant $\nu\in\mathbf Z^n$ and $c\in\mathbf Z$ let $h_\nu(c)$ be the
number of $a$ with $\nu(a)\geq c$.
\begin{enumerate}
\item A \emph{filling of shape $\nu$} is a function $I$ from $\mathbf Z$ to the subsets of
  $\{1,\dots,n\}$ with $|I(c)|=h_\nu(c)$ for $c\geq1$ and $|I(c)|=n-h_\nu(c)$ for $c\leq0$. As
  $h_\nu(c)=n$ for $c\leq\nu(n)$ and $h_\nu(c)=0$ for $c>\nu(1)$, such an $I$ has $I(c)=\emptyset$ for
  $c>\max(\nu(1),0)$ and for $c\leq\min(\nu(n),0)$, and $I(c)=\{1,\dots,n\}$ for $1\leq c\leq\nu(n)$ and
  for $\nu(1)<c\leq0$; in particular $I(c)=\emptyset$ whenever $|c|>|\nu(a)|$ for every $a$. Its
  \emph{weight} $\mathrm{wt}(I)\in\mathbf Z^n$
  has as entry at $m$ the number of $c\geq1$ with $m\in I(c)$ minus the number of $c\leq0$ with $m\in
  I(c)$. It is \emph{standard} if, with $I^*$ as in Definition~\ref{pd-symbols-path-generators}(4), the
  set $I^*(c+1)$ has at most as many elements $\leq t$ as $I^*(c)$ for all $c\in\mathbf Z$ and $1\leq
  t\leq n$. As $|I^*(c)|=h_\nu(c)$ for every $c$ and $h_\nu$ is nonincreasing, this says that for
  $k\leq h_\nu(c+1)$ the $k$-th smallest element of $I^*(c)$ is at most that of $I^*(c+1)$. By the proof
  of Lemma~\ref{pd-symbols-path-count}, the standard fillings of a path $\mu$
  (Definition~\ref{pd-symbols-path-generators}(4)) are the standard fillings of shape $\mu_N$.
\item The \emph{canonical filling} $\operatorname{can}_\nu$ of $\nu$ is given by
  $\operatorname{can}_\nu(c)=\{1,\dots,h_\nu(c)\}$ for $c\geq1$ and
  $\operatorname{can}_\nu(c)=\{h_\nu(c)+1,\dots,n\}$ for $c\leq0$. It is a standard filling of shape
  $\nu$, as $\operatorname{can}_\nu^*(c)=\{1,\dots,h_\nu(c)\}$ for all $c$ and $h_\nu$ is nonincreasing.
\item Let $A$ be a commutative ring and $h$ an $n\times n$ matrix over $A$. For
  $I,J\subset\{1,\dots,n\}$ with $|I|=|J|$, $\det h_{I,J}$ is the determinant of the submatrix of $h$ on
  the rows in $I$ and the columns in $J$, with $\det h_{\emptyset,\emptyset}=1$; for $|I|\neq|J|$ it is
  $0$. For a function $F$ from the subsets of $\{1,\dots,n\}$ to a commutative $A$-algebra put
  $(h\cdot F)(J)=\sum_IF(I)\det h_{I,J}$.
\item For $g\in GL_n(A)$ put $g_c=g$ for $c\geq1$ and $g_c={}^tg^{-1}$ for $c\leq0$. The
  \emph{bracket} of a filling $I$ of shape $\nu$ is
  $[\nu|I]_g=\prod_{c\in\mathbf Z}\det(g_c)_{\operatorname{can}_\nu(c),I(c)}\in A$; all but finitely
  many of the factors are $\det(g_c)_{\emptyset,\emptyset}=1$.
\end{enumerate}
% Lean: IsIndex fixes only the sizes |I c| (colHeight); the values empty and full in (1) follow from them. IsStd is
%   IsIndex with the order condition in the form nthElt (hatIdx I c) a <= nthElt (hatIdx I (c + 1)) a for
%   a < |hatIdx I (c + 1)|, the second form in (1).
\end{definition}
\leanref{Formalism/ThetaLogCoordinates/JacobiPathWeights/StatementA.lean}{Formalism.InvariantFactorPaths.GelfandTsetlinBases.IsIndex}
\leanref{Formalism/ThetaLogCoordinates/JacobiPathWeights/IndexWeightsDominance.lean}{Formalism.InvariantFactorPaths.GelfandTsetlinBases.idxWeight}
\leanref{Formalism/ThetaLogCoordinates/JacobiPathWeights/IndexWeightsDominance.lean}{Formalism.InvariantFactorPaths.GelfandTsetlinBases.hatIdx}
\leanref{Formalism/ThetaLogCoordinates/JacobiPathWeights/IndexWeightsDominance.lean}{Formalism.InvariantFactorPaths.GelfandTsetlinBases.IsStd}
\leanref{Formalism/ThetaLogCoordinates/JacobiPathWeights/StatementA.lean}{Formalism.InvariantFactorPaths.GelfandTsetlinBases.canIdx}
\leanref{Formalism/ThetaLogCoordinates/JacobiPathWeights/LiftWord.lean}{Formalism.InvariantFactorPaths.GelfandTsetlinBases.minorF}
\leanref{Formalism/ThetaLogCoordinates/JacobiPathWeights/LiftWord.lean}{Formalism.InvariantFactorPaths.GelfandTsetlinBases.actF}
\leanref{Formalism/ThetaLogCoordinates/JacobiPathWeights/IndexWeightsDominance.lean}{Formalism.InvariantFactorPaths.GelfandTsetlinBases.brk}

\begin{lemma}\label{pd-symbols-path-exterior}
Let $A$ be a commutative ring, $n\geq1$, $h$ an $n\times n$ matrix over $A$, $F$ a function from the
subsets of $\{1,\dots,n\}$ to a commutative $A$-algebra $A'$, and $x,y\in A'^n$ rows; we use the notation
of Definitions~\ref{pd-symbols-path-generators}(3) and~\ref{pd-symbols-path-canonical}(3). Let $F_0$ be
the function with $F_0(\emptyset)=1$ and $F_0(J)=0$ for $J\neq\emptyset$. Then $h\cdot F_0=F_0$ and
$h\cdot(F\wedge x)=(h\cdot F)\wedge(xh)$. If $h$ is invertible, then
$h\cdot(F\lrcorner y)=(h\cdot F)\lrcorner(y\,{}^th^{-1})$, and
$\det(h)\det({}^th^{-1})_{I,J}=\pm\det h_{I^{\mathrm c},J^{\mathrm c}}$ for all $I,J$ with $|I|=|J|$,
where ${}^{\mathrm c}$ denotes the complement in $\{1,\dots,n\}$.
\end{lemma}
\leanref{Formalism/ThetaLogCoordinates/JacobiPathWeights/JacobiMinorProofs/ActFormContract.lean}{Formalism.InvariantFactorPaths.GelfandTsetlinBases.actF_oneF}
\leanref{Formalism/ThetaLogCoordinates/JacobiPathWeights/JacobiMinorProofs/ActFormContract.lean}{Formalism.InvariantFactorPaths.GelfandTsetlinBases.actFormWedge}
\leanref{Formalism/ThetaLogCoordinates/JacobiPathWeights/JacobiMinorProofs/ActFormContract.lean}{Formalism.InvariantFactorPaths.GelfandTsetlinBases.actFormContract}
\leanref{Formalism/ThetaLogCoordinates/JacobiPathWeights/JacobiMinorProofs/JacobiMinorProof.lean}{Formalism.InvariantFactorPaths.GelfandTsetlinBases.JacobiMinor_proof'}

\begin{proof}
Let $\Lambda$ be the exterior algebra over $A'$ of $A'^n$, with $e_1,\dots,e_n$ the unit rows, and put
$e_J=e_{j_1}\wedge\dots\wedge e_{j_k}$ for $J=\{j_1<\dots<j_k\}$, $e_\emptyset=1$. To $F$ attach
$\omega_F=\sum_JF(J)e_J$. For $m\in J$, $e_{J\smallsetminus\{m\}}\wedge e_m=\epsilon_J(m)e_J$, so
$\omega_{F\wedge x}=\omega_F\wedge\sum_mx_me_m$. Let $\Lambda(h)$ be the algebra endomorphism of
$\Lambda$ with $e_m\mapsto e_mh$. Expanding $e_{i_1}h\wedge\dots\wedge e_{i_k}h$ gives
$\Lambda(h)e_I=\sum_J\det h_{I,J}e_J$, so $\omega_{h\cdot F}=\Lambda(h)\omega_F$. As $\Lambda(h)$ is an
algebra map, $\Lambda(h)1=1$ and $\Lambda(h)(\omega\wedge x)=\Lambda(h)\omega\wedge xh$; these are the
first two identities.

For a row $y$ let $\iota_y$ be the $A'$-linear map of $\Lambda$ with
$\iota_y(v_1\wedge\dots\wedge v_k)=\sum_l(-1)^{k-l}\langle v_l,y\rangle\,v_1\wedge\dots\wedge\widehat{v_l}\wedge
\dots\wedge v_k$, where $\langle v,y\rangle=\sum_mv_my_m$; the right side is alternating and multilinear in
the $v_l$, so $\iota_y$ is well defined. On $e_K$, removing $e_m$ from the place $l$ gives the sign
$(-1)^{k-l}=\epsilon_{K\smallsetminus\{m\}}(m)$, so $\omega_{F\lrcorner y}=\iota_y\omega_F$. If $h$ is
invertible and $y'=y\,{}^th^{-1}$, then $\langle vh,y'\rangle=\langle v,y\rangle$, so applying the formula
to the vectors $v_lh$ gives $\Lambda(h)\iota_y=\iota_{y'}\Lambda(h)$; this is the third identity.

For the last, let $(\star F)(J)=(-1)^{\sum_{j\in J}(n+j)}F(J^{\mathrm c})$. A check of signs gives
$\star(F\wedge x)=(\star F)\lrcorner x$. We show $h\cdot\star F=\det(h)\,\star({}^th^{-1}\cdot F)$ for invertible
$h$. Both sides are $A'$-linear in $F$, and every $F$ is a sum of multiples of iterated wedges
$F_0\wedge e_{i_1}\wedge\dots\wedge e_{i_l}$ of $F_0$ with unit rows. For $F_0$: $\star F_0$ is $\pm1$ at
$\{1,\dots,n\}$ and $0$ elsewhere, so $h\cdot\star F_0=\det(h)\star F_0$, while ${}^th^{-1}\cdot F_0=F_0$. If the
identity holds for $F$, then by the third identity for $h$ and the second for ${}^th^{-1}$,
\[
 h\cdot\star(F\wedge x)=(h\cdot\star F)\lrcorner(x\,{}^th^{-1})=\det(h)\big(\star({}^th^{-1}\cdot F)\big)\lrcorner
 (x\,{}^th^{-1})=\det(h)\star\big(({}^th^{-1}\cdot F)\wedge x\,{}^th^{-1}\big)=\det(h)\star\big({}^th^{-1}\cdot(F\wedge x)\big).
\]
Now take $F$ equal to $1$ at $I$ and $0$ elsewhere, and evaluate at $J^{\mathrm c}$: the left side is
$\pm\det h_{I^{\mathrm c},J^{\mathrm c}}$ and the right side is $\pm\det(h)\det({}^th^{-1})_{I,J}$.
% Lean: actF_oneF, actFormWedge (lap_row), actFormContract (by induction on F, with contractWedge), and
%   JacobiMinor_proof' = jacobi_gen through the Hodge star star (epsN J = sum over J of n + 1 + j, 0-based):
%   star_wedge and star_act (by form_induction), applied to basisF I at the complement of J; the sign is
%   (-1)^(sum I + sum J).
\end{proof}

\begin{lemma}\label{pd-symbols-path-action}
Let $A$ be a commutative ring, $N\geq0$ and $n\geq1$ integers with a lettering of $\{0,\dots,N-1\}$,
$g\in GL_n(A)$, and $\Sigma_g$ the $A$-algebra endomorphism of $A[z]$ that replaces each row $z_i$ by
$z_ig$ if $i$ is a vector letter and by $z_i\,{}^tg^{-1}$ if it is a covector letter; we use the notation
of Definitions~\ref{pd-symbols-path-generators} and~\ref{pd-symbols-path-canonical}. Let $\mu$ be a path
of rank $n$ for some $B$. Then $\Sigma_g(\mathrm{Col}_{\mu,c}(J))=(g_c\cdot\mathrm{Col}_{\mu,c})(J)$ for
all $c\in\mathbf Z$ and $J\subset\{1,\dots,n\}$, and for every function $I$ from $\mathbf Z$ to the
subsets of $\{1,\dots,n\}$
\[
 \Sigma_g(\Psi(\mu,I))=\sum_{I'}\Psi(\mu,I')\prod_{|c|\leq T_\mu}\det(g_c)_{I'(c),I(c)},
\]
the sum over the functions $I'$ from $\{c:|c|\leq T_\mu\}$ to the subsets of $\{1,\dots,n\}$.
\end{lemma}
\leanref{Formalism/ThetaLogCoordinates/JacobiPathWeights/IndexWeightsDominance.lean}{Formalism.InvariantFactorPaths.GelfandTsetlinBases.actAll}
\leanref{Formalism/ThetaLogCoordinates/JacobiPathWeights/JacobiMinorProofs/PathProductHighestWeight.lean}{Formalism.InvariantFactorPaths.GelfandTsetlinBases.A1hw.act_pathWord}
\leanref{Formalism/ThetaLogCoordinates/JacobiPathWeights/JacobiMinorProofs/LemmaA2indepOf.lean}{Formalism.InvariantFactorPaths.GelfandTsetlinBases.A3ev.col_act}
\leanref{Formalism/ThetaLogCoordinates/JacobiPathWeights/JacobiMinorProofs/LemmaA2indepOf.lean}{Formalism.InvariantFactorPaths.GelfandTsetlinBases.A3ev.act_psi}

\begin{proof}
Fix $c$, put $h=g_c$, and let $F_j$ be the function after the first $j$ steps of the definition of
$\mathrm{Col}_{\mu,c}$. We show $\Sigma_g\circ F_j=h\cdot F_j$ by induction on $j$. For $j=0$ this is
$h\cdot F_0=F_0$ (Lemma~\ref{pd-symbols-path-exterior}), as $\Sigma_g$ fixes constants. If
$h_j(c)=h_{j+1}(c)$, then $F_{j+1}=F_j$. Otherwise, as $\mu$ is a path, $h_j(c)<h_{j+1}(c)$ at a vector
letter and $h_{j+1}(c)<h_j(c)$ at a covector letter (proof of Lemma~\ref{pd-symbols-path-clauses}(3)), so
the step at $j$ is $F_j\wedge z_j$ for a vector letter with $c\geq1$ and for a covector letter with
$c\leq0$, and $F_j\lrcorner z_j$ in the other two cases. As $\Sigma_g$ is a ring homomorphism,
$\Sigma_g\circ F_{j+1}$ is the same step applied to $\Sigma_g\circ F_j=h\cdot F_j$ and the row
$\Sigma_g(z_j)$. If $c\geq1$, then $h=g$, and $\Sigma_g(z_j)$ is $z_jh$ at a vector letter, where the step
is a wedge, and $z_j\,{}^th^{-1}$ at a covector letter, where it is a contraction. If $c\leq0$, then
$h={}^tg^{-1}$, so $g={}^th^{-1}$, and $\Sigma_g(z_j)$ is $z_j\,{}^th^{-1}$ at a vector letter, where the
step is a contraction, and $z_jh$ at a covector letter, where it is a wedge. In all four cases
Lemma~\ref{pd-symbols-path-exterior} gives $\Sigma_g\circ F_{j+1}=h\cdot F_{j+1}$.

As $\Sigma_g$ is a ring homomorphism,
$\Sigma_g(\Psi(\mu,I))=\prod_{|c|\leq T_\mu}\sum_{J}\mathrm{Col}_{\mu,c}(J)\det(g_c)_{J,I(c)}$, and expanding
the product gives the formula.
% Lean: A1hw.act_pathWord with actFormWedge and actFormContract (by induction on j), col_act, and act_psi
%   (Finset.prod_sum over a finite set F of columns outside which the words are oneF).
\end{proof}

\begin{lemma}\label{pd-symbols-path-filling-weights}
Let $N\geq0$ and $n\geq1$ be integers with a lettering of $\{0,\dots,N-1\}$, $R$ a commutative ring and
$S=\{0,\dots,N-1\}$; we use the notation of Definitions~\ref{pd-symbols-path-generators},
\ref{pd-symbols-path-weights} and~\ref{pd-symbols-path-canonical}.
\begin{enumerate}
\item Let $\mu$ be a path of rank $n$ for some $B$ and $I$ a filling of shape $\mu_N$. Then every monomial
  $z^m$ of $\Psi(\mu,I)\in R[z]$ has $\mathrm{wt}_S(m)=\mathrm{wt}(I)$.
\item Let $\nu\in\mathbf Z^n$ be dominant and $I$ a filling of shape $\nu$. Then
  $\mathrm{wt}(\operatorname{can}_\nu)=\nu$ and $\mathrm{wt}(I)\preceq\nu$, and $\mathrm{wt}(I)=\nu$ only if
  $I=\operatorname{can}_\nu$.
\end{enumerate}
\end{lemma}
\leanref{Formalism/ThetaLogCoordinates/JacobiPathWeights/JacobiMinorProofs/IndexWeightsDominance.lean}{Formalism.InvariantFactorPaths.GelfandTsetlinBases.psi_weight}
\leanref{Formalism/ThetaLogCoordinates/JacobiPathWeights/JacobiMinorProofs/IndexWeightsDominance.lean}{Formalism.InvariantFactorPaths.GelfandTsetlinBases.index_weights_dominance}

\begin{proof}
(1) Put $s_c=1$ for $c\geq1$ and $s_c=-1$ for $c\leq0$, and let $1_J\in\mathbf Z^n$ be the indicator
vector of $J$ and $e_m=1_{\{m\}}$. We show that every monomial of $\mathrm{Col}_{\mu,c}(J)$ has
$S$-weight $s_c1_J$, by induction over the steps of its definition. The function $F_0$ is $1$ at
$\emptyset$ and $0$ elsewhere. As in the proof of Lemma~\ref{pd-symbols-path-action}, a wedge step occurs
at a vector letter with $c\geq1$ or a covector letter with $c\leq0$, where the variable $z_{j,m}$ has
$S$-weight $s_ce_m$; its terms are $\pm z_{j,m}F(J\smallsetminus\{m\})$ with $m\in J$, of weight
$s_ce_m+s_c1_{J\smallsetminus\{m\}}=s_c1_J$. A contraction step occurs at a covector letter with $c\geq1$
or a vector letter with $c\leq0$, where $z_{j,m}$ has weight $-s_ce_m$; its terms are
$\pm z_{j,m}F(J\cup\{m\})$ with $m\notin J$, of weight $s_c1_J$. The weight of a product of monomials is
the sum of their weights, so every monomial of $\Psi(\mu,I)$ has weight $\sum_{|c|\leq T_\mu}s_c1_{I(c)}$,
which is $\mathrm{wt}(I)$ as $I(c)=\emptyset$ for $|c|>T_\mu\geq|\mu_N(a)|$.

(2) Write $h=h_\nu$ and $s_c$ as in (1). As $\nu$ is dominant, $m\leq h(c)$ if and only if
$\nu(m)\geq c$; so the entry of $\mathrm{wt}(\operatorname{can}_\nu)$ at $m$ is the number of $c$ with
$1\leq c\leq\nu(m)$ minus the number with $\nu(m)<c\leq0$, which is $\nu(m)$. For $0\leq j\leq n$ let
$a_c(j)$ and $b_c(j)$ be the numbers of elements $\leq j$ of $I(c)$ and of $\operatorname{can}_\nu(c)$.
The sum of the first $j$ entries of $\mathrm{wt}(I)$ is $\sum_cs_ca_c(j)$, and that of $\nu$ is
$\sum_cs_cb_c(j)$. For $c\geq1$, $I(c)$ has $h(c)$ elements, so $a_c(j)\leq\min(j,h(c))=b_c(j)$. For
$c\leq0$, $I(c)$ has $n-h(c)$ elements, at most $n-j$ of them larger than $j$, so
$a_c(j)\geq\max(0,j-h(c))=b_c(j)$. Hence $s_ca_c(j)\leq s_cb_c(j)$ for all $c$ and $j$, with equality for
$j=n$; so $\mathrm{wt}(I)\preceq\nu$. If $\mathrm{wt}(I)=\nu$, then for every $j$ the terms
$s_c(a_c(j)-b_c(j))\leq0$ sum to $0$, so all vanish. With $j=h(c)$ this gives, for $c\geq1$,
$a_c(h(c))=h(c)=|I(c)|$, so $I(c)=\{1,\dots,h(c)\}$; and for $c\leq0$, $a_c(h(c))=0$, so $I(c)$ lies in
$\{h(c)+1,\dots,n\}$, which has $|I(c)|$ elements. So $I=\operatorname{can}_\nu$.
% Lean: (1) psi_weight (by induction over the steps, wedgeW and contrW); (2) index_weights_dominance
%   (sum_idxWeight, col_le, col_eq, idxWeight_can).
\end{proof}

\begin{lemma}\label{pd-symbols-path-top-weight}
Let $A$ be a commutative ring, $N\geq0$ and $n\geq1$ integers with a lettering of $\{0,\dots,N-1\}$,
$g\in GL_n(A)$ and $\Sigma_g$ as in Lemma~\ref{pd-symbols-path-action}, and $S=\{0,\dots,N-1\}$; we use
the notation of Definitions~\ref{pd-symbols-path-generators}, \ref{pd-symbols-path-weights}
and~\ref{pd-symbols-path-canonical}. For $P\in A[z]$ and $\nu\in\mathbf Z^n$ let $P_\nu$ be the sum of the
terms of $P$ whose monomials $z^m$ have $\mathrm{wt}_S(m)=\nu$. Let $\mu$ be a path of rank $n$ for some
$B$ and $I$ a filling of shape $\mu_N$. Then $\Sigma_g(\Psi(\mu,I))_\nu=0$ unless $\nu\preceq\mu_N$, and
$\Sigma_g(\Psi(\mu,I))_{\mu_N}=[\mu_N|I]_g\,\Psi(\mu,\operatorname{can}_{\mu_N})$.
\end{lemma}
\leanref{Formalism/ThetaLogCoordinates/JacobiPathWeights/IndexWeightsDominance.lean}{Formalism.InvariantFactorPaths.GelfandTsetlinBases.FactEV}
\leanref{Formalism/ThetaLogCoordinates/JacobiPathWeights/JacobiMinorProofs/LemmaA2indepOf.lean}{Formalism.InvariantFactorPaths.GelfandTsetlinBases.factEV_of}
\leanref{Formalism/ThetaLogCoordinates/JacobiPathWeights/PathBasisTheorem.lean}{Formalism.InvariantFactorPaths.GelfandTsetlinBases.Top.factEV}

\begin{proof}
By Lemma~\ref{pd-symbols-path-action},
$\Sigma_g(\Psi(\mu,I))=\sum_{I'}\Psi(\mu,I')\prod_{|c|\leq T_\mu}\det(g_c)_{I'(c),I(c)}$. A term with
$|I'(c)|\neq|I(c)|$ for some $c$ vanishes. Extended by $\emptyset$ to all $c$, the remaining $I'$ are
fillings of shape $\mu_N$, as $I(c)=\emptyset$ for $|c|>T_\mu$. By
Lemma~\ref{pd-symbols-path-filling-weights}(1), all monomials of $\Psi(\mu,I')$ have $S$-weight
$\mathrm{wt}(I')$, so $\Sigma_g(\Psi(\mu,I))_\nu$ is the sum of the terms with $\mathrm{wt}(I')=\nu$. By
Lemma~\ref{pd-symbols-path-filling-weights}(2), $\mathrm{wt}(I')\preceq\mu_N$, so this sum is empty unless
$\nu\preceq\mu_N$, and for $\nu=\mu_N$ its only term is the one with $I'=\operatorname{can}_{\mu_N}$. That
term is $\Psi(\mu,\operatorname{can}_{\mu_N})[\mu_N|I]_g$, since the factors of the bracket with
$|c|>T_\mu$ are $\det(g_c)_{\emptyset,\emptyset}=1$.
% Lean: factEV_of fact_index_weights psiWeight (act_psi, term_cases, term_hom, idxWeight_can).
\end{proof}

\begin{lemma}\label{pd-symbols-path-canonical-independent}
Let $F$ be a field, $N\geq0$ and $n\geq1$ integers with a lettering of $\{0,\dots,N-1\}$,
$B\colon\{0,\dots,N-1\}\to\mathbf Z_{\geq0}$ and $\nu\in\mathbf Z^n$ dominant; we use the notation of
Definitions~\ref{pd-symbols-path-generators} and~\ref{pd-symbols-path-canonical}. Then the polynomials
$\Psi(\mu,\operatorname{can}_\nu)\in F[z]$, $\mu$ a path of rank $n$ for $B$ with $\mu_N=\nu$, are linearly
independent over $F$.
\end{lemma}
\leanref{Formalism/ThetaLogCoordinates/JacobiPathWeights/StatementA.lean}{Formalism.InvariantFactorPaths.GelfandTsetlinBases.PathProductLinearIndependent}
\leanref{Formalism/ThetaLogCoordinates/JacobiPathWeights/JacobiMinorProofs/IsWeightedHomogeneousZeroStripPathWord.lean}{Formalism.InvariantFactorPaths.GelfandTsetlinBases.pathProductLinearIndependent}
\leanref{Formalism/ThetaLogCoordinates/JacobiPathWeights/JacobiMinorProofs/IsWeightedHomogeneousZeroStripPathWord.lean}{Formalism.InvariantFactorPaths.GelfandTsetlinBases.A1.psi_lead}

\begin{proof}
By induction on $N$. For $N=0$ the only path is $\mu=(0)$, which ends at $\nu$ only if $\nu=0$; then
$\operatorname{can}_0(c)=\emptyset$ for all $c$, every $\mathrm{Col}_{\mu,c}$ is $F_0$, and
$\Psi(\mu,\operatorname{can}_0)=1\neq0$.

Let $N\geq1$ and $\ell=N-1$. For a path $\mu$ for $B$ let $\mu'=(\mu_0,\dots,\mu_\ell)$, a path for the
restriction of $B$ to $\{0,\dots,\ell-1\}$, and $u_\mu=\prod_mz_{\ell,m}^{|\mu_N(m)-\mu_\ell(m)|}$. Order the
monomials in $z_{\ell,1},\dots,z_{\ell,n}$ lexicographically by their exponents, read at
$z_{\ell,n},z_{\ell,n-1},\dots,z_{\ell,1}$ in this order if $\ell$ is a vector letter and at
$z_{\ell,1},\dots,z_{\ell,n}$ if it is a covector letter; this is a total order compatible with
multiplication. Every monomial of $F[z]$ is the product of its \emph{part} in the row $z_\ell$ and a
monomial in the other rows; for $P\in F[z]$ and a monomial $u$ in the row $z_\ell$ let $P[u]$ be the sum
of the terms of $P$ whose part is $u$. We show:
\begin{enumerate}
\item[($*$)] every monomial of $\Psi(\mu,\operatorname{can}_{\mu_N})$ has part at most $u_\mu$, and
  $\Psi(\mu,\operatorname{can}_{\mu_N})[u_\mu]=\pm\Psi(\mu',\operatorname{can}_{\mu_\ell})\,u_\mu$.
\end{enumerate}
Fix $c$, write $h_j=h_j(c)$, and let $F_j$ be the functions of the definition of
$\mathrm{Col}_{\mu,c}$; then $F_\ell=\mathrm{Col}_{\mu',c}$ involves no variable of the row $z_\ell$. If
$h_\ell=h_N$, then $\operatorname{can}_{\mu_N}(c)=\operatorname{can}_{\mu_\ell}(c)$ and
$\mathrm{Col}_{\mu,c}(\operatorname{can}_{\mu_N}(c))=\mathrm{Col}_{\mu',c}(\operatorname{can}_{\mu_\ell}(c))$,
with part $1$. Otherwise let $t=\min(h_\ell,h_N)+1$. As in the proof of
Lemma~\ref{pd-symbols-path-clauses}(3), $h_N=h_\ell+1=t$ at a vector letter and $h_\ell=h_N+1=t$ at a
covector letter. In each of the four cases $\mathrm{Col}_{\mu,c}(\operatorname{can}_{\mu_N}(c))$ is a sum
of terms $\pm z_{\ell,m}F_\ell(J_m)$, over a set of $m$ that contains $t$, with
$J_t=\operatorname{can}_{\mu_\ell}(c)$:
\begin{itemize}
\item vector letter, $c\geq1$: a wedge at $\{1,\dots,t\}$, over the $m\leq t$;
\item vector letter, $c\leq0$: a contraction at $\{t+1,\dots,n\}$, over the $m\leq t$;
\item covector letter, $c\geq1$: a contraction at $\{1,\dots,t-1\}$, over the $m\geq t$;
\item covector letter, $c\leq0$: a wedge at $\{t,\dots,n\}$, over the $m\geq t$.
\end{itemize}
In each case $z_{\ell,m}<z_{\ell,t}$ for the other $m$. So every monomial of the column has part at most
$z_{\ell,t}$, and its terms with part $z_{\ell,t}$ are $\pm z_{\ell,t}\mathrm{Col}_{\mu',c}(
\operatorname{can}_{\mu_\ell}(c))$. As the order is total and compatible with multiplication, every monomial
of the product $\Psi(\mu,\operatorname{can}_{\mu_N})$ of these columns over $|c|\leq T_\mu$ has part at most
the product $u$ of the $z_{\ell,t(c)}$ over the $c$ with $h_\ell(c)\neq h_N(c)$, and the terms with part
$u$ are the product of the terms of the columns with the largest parts. The factors
$\mathrm{Col}_{\mu',c}(\operatorname{can}_{\mu_\ell}(c))$ with $T_{\mu'}<|c|\leq T_\mu$ are $1$, as in the
proof of Lemma~\ref{pd-symbols-path-clauses}(4); so this product is
$\pm\Psi(\mu',\operatorname{can}_{\mu_\ell})u$. Finally $u=u_\mu$. At a vector letter, $t(c)=m$ means
$h_N(c)=m$ and $h_\ell(c)=m-1$, that is, $\mu_\ell(m)<c\leq\mu_N(m)$, as $\mu_\ell$ and $\mu_N$ are dominant
and $h_N\leq h_\ell+1$; there are $\mu_N(m)-\mu_\ell(m)$ such $c$. At a covector letter they are the $c$
with $\mu_N(m)<c\leq\mu_\ell(m)$. This proves ($*$).

Suppose that $\sum_\mu a_\mu\Psi(\mu,\operatorname{can}_\nu)=0$ with $a_\mu\in F$ not all $0$, the sum
over the paths for $B$ with $\mu_N=\nu$. Among the $\mu$ with $a_\mu\neq0$ choose one with $u_\mu$
largest, and call it $u$. By ($*$), taking the terms with part $u$ gives
$\sum_{u_\mu=u}\pm a_\mu\Psi(\mu',\operatorname{can}_{\mu_\ell})u=0$. For these $\mu$, $u$ determines
$|\nu(m)-\mu_\ell(m)|$ for every $m$, and $\nu(m)-\mu_\ell(m)$ is $\geq0$ at a vector letter and $\leq0$
at a covector letter; so $\mu_\ell$ is one and the same $\nu'$ for all of them, and
$\mu\mapsto\mu'$ is injective on them. Dividing by $u$, the induction hypothesis for the paths $\mu'$
ending at $\nu'$ gives $a_\mu=0$ whenever $u_\mu=u$, a contradiction.
% Lean: pathProductLinearIndependent, by induction on N with li_of_lead for the weights rho n N (isVec N)
%   (lexicographic on the last row, reversed by rw_ for vector letters); (*) is psi_lead from col_lead (the
%   four cases, leading_wedge_lemma_12, leading_contract_lemma_13) and sum_Dc; end_determined.
\end{proof}

\begin{lemma}\label{pd-symbols-path-brackets}
Let $F$ be a field, $n\geq1$, $L$ the field of rational functions over $F$ in the entries of an
$n\times n$ matrix $\mathfrak g$ of indeterminates, and $\nu\in\mathbf Z^n$ dominant; we use the notation
of Definition~\ref{pd-symbols-path-canonical}. Then the brackets $[\nu|I]_{\mathfrak g}\in L$, $I$ a
standard filling of shape $\nu$, are linearly independent over $F$.
\end{lemma}
\leanref{Formalism/ThetaLogCoordinates/JacobiPathWeights/IndexWeightsDominance.lean}{Formalism.InvariantFactorPaths.GelfandTsetlinBases.FactBrk}
\leanref{Formalism/ThetaLogCoordinates/JacobiPathWeights/JacobiMinorProofs/JacobiMinorProof.lean}{Formalism.InvariantFactorPaths.GelfandTsetlinBases.fact_break}
\leanref{Formalism/ThetaLogCoordinates/JacobiPathWeights/PathBasisTheorem.lean}{Formalism.InvariantFactorPaths.GelfandTsetlinBases.Top.factBrk}
\leanref{Formalism/ThetaLogCoordinates/JacobiPathWeights/JacobiMinorProofs/LinearIndependentOfLead.lean}{Formalism.InvariantFactorPaths.GelfandTsetlinBases.path_basis_lemma}

\begin{proof}
Write $h=h_\nu$, put $l_0=\min(0,\nu(n))$ and $l_1=\max(0,\nu(1))$, and let $I$ be a filling of shape
$\nu$. Then $I(c)$ and $\operatorname{can}_\nu(c)$ are empty for $c\notin(l_0,l_1]$, so
$[\nu|I]_{\mathfrak g}=\prod_{l_0<c\leq l_1}\det(\mathfrak g_c)_{\operatorname{can}_\nu(c),I(c)}$. For
$c\geq1$ the factor is $\det\mathfrak g_{\{1,\dots,h(c)\},I^*(c)}$. For $c\leq0$,
Lemma~\ref{pd-symbols-path-exterior} gives
$\det(\mathfrak g)\det({}^t\mathfrak g^{-1})_{\operatorname{can}_\nu(c),I(c)}=\pm\det\mathfrak g_{\{1,\dots,h(c)\},I^*(c)}$,
as the complements of $\operatorname{can}_\nu(c)$ and $I(c)$ are $\{1,\dots,h(c)\}$ and $I^*(c)$. Hence
\[
 [\nu|I]_{\mathfrak g}=\pm\det(\mathfrak g)^{l_0}\,D(I),\qquad
 D(I)=\prod_{l_0<c\leq l_1}\det\mathfrak g_{\{1,\dots,h(c)\},I^*(c)}\in F[\mathfrak g].
\]
Call a family $T=(T(c))_{l_0<c\leq l_1}$ of subsets of $\{1,\dots,n\}$ \emph{semistandard} if for every
$c$ and $a$ such that $T(c+1)$ has at least $a$ elements, so has $T(c)$, and the $a$-th element of $T(c)$
is at most that of $T(c+1)$; put
$D_T=\prod_{l_0<c\leq l_1}\det\mathfrak g_{\{1,\dots,|T(c)|\},T(c)}$. Then $D(I)=D_T$ with $T(c)=I^*(c)$,
a set of $h(c)$ elements. If $I$ is standard, this $T$ is semistandard (proof of
Lemma~\ref{pd-symbols-path-count}(i)), and distinct $I$ give distinct $T$, since $I(c)$ is determined by
$I^*(c)$ and is empty outside $(l_0,l_1]$. As $F[\mathfrak g]\subset L$ and
$\det(\mathfrak g)^{l_0}\neq0$, it suffices to show that the $D_T$, over the semistandard $T$, are linearly
independent over $F$. We use leading terms.

Order the entries $\mathfrak g_{a,b}$ lexicographically by $(a,b)$, and for monomials $z^e$ and
$z^{e'}$ in the entries put $z^e<z^{e'}$ if $e(a,b)<e'(a,b)$ at the first entry $(a,b)$ at which $e$ and
$e'$ differ. This is a total order on monomials, and $z^e<z^{e'}$ implies $z^{e+f}<z^{e'+f}$. Say that
$P\in F[\mathfrak g]$ has \emph{leading monomial} $z^d$ if the coefficient of $z^d$ in $P$ is $1$ and all
other monomials of $P$ are smaller than $z^d$.

(a) \emph{Minors.} Let $b_1<\dots<b_h$. The minor $\det\mathfrak g_{\{1,\dots,h\},\{b_1,\dots,b_h\}}$ is the
sum over the permutations $\pi$ of $\{1,\dots,h\}$ of
$\operatorname{sgn}(\pi)\prod_{a\leq h}\mathfrak g_{a,b_{\pi(a)}}$, and these monomials are distinct. Let
$\pi\neq1$ and let $a$ be the least index with $\pi(a)\neq a$; then $\pi(a)>a$. The monomials for $\pi$ and
for $1$ agree at the entries in the rows before $a$, and neither involves an entry $(a,b)$ with $b<b_a$;
at $(a,b_a)$ the first has exponent $0$ and the second exponent $1$. So the minor has leading monomial
$\prod_{a\leq h}\mathfrak g_{a,b_a}$.

(b) \emph{Products.} If $P$ and $P'$ have leading monomials $z^d$ and $z^{d'}$, then $PP'$ has leading
monomial $z^{d+d'}$: every other product of a monomial of $P$ and one of $P'$ is smaller than $z^{d+d'}$,
by the compatibility of the order with multiplication. By (a), $D_T$ has leading monomial $z^{d(T)}$, where
$d(T)(a,b)$ is the number of $c$ such that the $a$-th element of $T(c)$ is $b$.

(c) \emph{Injectivity.} Let $T$ be semistandard. For $a\leq n$ let $\tau_a(c)$ be the $a$-th element of
$T(c)$, or $n+1$ if $T(c)$ has fewer than $a$ elements. By semistandardness, $\tau_a$ is nondecreasing in
$c$. For $b\leq n$, $d(T)(a,b)$ is the number of $c$ with $\tau_a(c)=b$. A nondecreasing function on the
interval $(l_0,l_1]$ with values in $\{1,\dots,n+1\}$ is determined by these numbers: $\tau_a(c)\leq v$ if
and only if $c$ is among the first $\sum_{b\leq v}d(T)(a,b)$ elements of the interval. So $d(T)$
determines every $\tau_a$, hence every $T(c)$, and distinct semistandard $T$ have distinct leading
monomials.

(d) \emph{Independence.} Suppose that $\sum_Ta_TD_T=0$ with $a_T\in F$ not all $0$. Among the $T$ with
$a_T\neq0$ choose the one with the largest leading monomial $z^{d(T_0)}$, unique by (c). Every monomial of
another $D_T$ with $a_T\neq0$ is at most $z^{d(T)}<z^{d(T_0)}$, so the coefficient of $z^{d(T_0)}$ in the
sum is $a_{T_0}\neq0$, a contradiction. The linear independence of the $D_T$ is also a case of the
standard basis theorem for bitableaux \cite{DEP80}.
% Lean: fact_break: brk_eq (factor, with JacobiMinor for c <= 0, and prod_Ioc_eq over (lo, hi]),
%   semi_colT, colT_inj, and LemmaA = path_basis_lemma: linearIndependent_of_Lead key_inj (d) with
%   Lead_tab, which is Lead_prod (b) of Lead_minor and Lead_det (mon_LL, (a)), and tab_inj (c) (semi_mono,
%   tabDiag_apply, mono_le_iff, mono_eq_of_counts). The order is LL key with key = toLex on (row, column);
%   LL_trans, LL_total and LL_add give the properties used.
\end{proof}

\begin{lemma}\label{pd-symbols-path-independent}
Let $F$ be a field, $N\geq0$ and $n\geq1$ integers with a lettering of $\{0,\dots,N-1\}$, and
$B\colon\{0,\dots,N-1\}\to\mathbf Z_{\geq0}$; we use the notation of
Definition~\ref{pd-symbols-path-generators}. Then the path generators $\Psi(\mu,I)\in F[z]$, indexed by
the pairs $(\mu,I)$ of a path $\mu$ of rank $n$ for $B$ and a standard filling $I$ of $\mu$, are
linearly independent over $F$.
\end{lemma}
\leanref{Formalism/ThetaLogCoordinates/JacobiPathWeights/PathBasisTheorem.lean}{Formalism.InvariantFactorPaths.GelfandTsetlinBases.Top.lemmaA2indep}
\leanref{Formalism/ThetaLogCoordinates/JacobiPathWeights/JacobiMinorProofs/LemmaA2indepOf.lean}{Formalism.InvariantFactorPaths.GelfandTsetlinBases.lemmaA2indep_of}

\begin{proof}
We use the notation of Definitions~\ref{pd-symbols-path-weights}(4) and~\ref{pd-symbols-path-canonical}.
For $\nu\in\mathbf Z^n$ put $\operatorname{ht}(\nu)=\sum_{j=1}^n\sum_{a\leq j}\nu(a)$. If $\nu\preceq\nu'$
and $\operatorname{ht}(\nu')\leq\operatorname{ht}(\nu)$, then $\nu=\nu'$: each partial sum of $\nu$ is at
most that of $\nu'$, and their totals agree, so all the partial sums agree.

Suppose that $\sum_xa_x\Psi_x=0$, over the pairs $x=(\mu,I)$ of the statement, with $\Psi_x=\Psi(\mu,I)$
and $a_x\in F$ almost all $0$ and not all $0$. Choose $x_0$ with $a_{x_0}\neq0$ whose path has
$\operatorname{ht}(\mu_N)$ largest, and let $\nu_0$ be the last vector $\mu_N$ of its path. Let $L$ and
$\mathfrak g$ be as in Lemma~\ref{pd-symbols-path-brackets}, $\Sigma_{\mathfrak g}$ as in
Lemma~\ref{pd-symbols-path-action} over $L$, and $S=\{0,\dots,N-1\}$; for $P\in L[z]$ let $P_{\nu_0}$ be
the sum of the terms of $P$ of $S$-weight $\nu_0$. The map $P\mapsto\Sigma_{\mathfrak g}(P)_{\nu_0}$ from
$F[z]$ to $L[z]$ is $F$-linear, so $\sum_xa_x\Sigma_{\mathfrak g}(\Psi_x)_{\nu_0}=0$.

Let $x=(\mu,I)$ with $a_x\neq0$ and $\Sigma_{\mathfrak g}(\Psi_x)_{\nu_0}\neq0$. By
Lemma~\ref{pd-symbols-path-top-weight}, $\nu_0\preceq\mu_N$. By the choice of $x_0$,
$\operatorname{ht}(\mu_N)\leq\operatorname{ht}(\nu_0)$. So $\mu_N=\nu_0$. For every $x=(\mu,I)$ with
$\mu_N=\nu_0$ the same lemma gives
$\Sigma_{\mathfrak g}(\Psi_x)_{\nu_0}=[\nu_0|I]_{\mathfrak g}\,\Psi(\mu,\operatorname{can}_{\nu_0})$. Hence
\[
 \sum_\mu\Big(\sum_Ia_{(\mu,I)}[\nu_0|I]_{\mathfrak g}\Big)\Psi(\mu,\operatorname{can}_{\nu_0})=0,
\]
the outer sum over the paths $\mu$ for $B$ with $\mu_N=\nu_0$ and the inner one over the standard
fillings $I$ of shape $\nu_0$. By Lemma~\ref{pd-symbols-path-canonical-independent} over the field $L$,
every inner sum vanishes. By Lemma~\ref{pd-symbols-path-brackets}, $a_{(\mu,I)}=0$ whenever $\mu_N=\nu_0$.
In particular $a_{x_0}=0$, a contradiction.
% Lean: lemmaA2indep_of factEV factBrk (height, domLe_eq, evC = weight-mu0 component of the generic
%   translate actAll (genMat k n), evC_psi, ev_gen), with fam_li = linearIndependent_smul of FactBrk and
%   pathProductLinearIndependent over GenField k n.
\end{proof}

\begin{lemma}\label{pd-symbols-path-basis}
Let $R$ be a commutative ring, $N\geq0$ and $n\geq1$ integers with a lettering of $\{0,\dots,N-1\}$,
and $B\colon\{0,\dots,N-1\}\to\mathbf Z_{\geq0}$; we use the notation of
Definition~\ref{pd-symbols-path-generators}. Then the path generators $\Psi(\mu,I)\in R[z]$, indexed by
the pairs $(\mu,I)$ of a path $\mu$ of rank $n$ for $B$ and a standard filling $I$ of $\mu$, are
linearly independent over $R$ and span the $R$-module of the polynomials of row degrees $B$.
\end{lemma}
\leanref{Formalism/SharpRestriction/GrowthBounds/StatementOfDegreeAndBound.lean}{Formalism.MainNotions.CovariantDegrees.PathBasisOverRing}
\leanref{Formalism/SharpRestriction/GrowthBounds/BinomialPolynomialIdentities/AdmissiblePolynomialBasisProv.lean}{Formalism.MainNotions.G.ColumnFamilies.admissible_standardPath_basis}

\begin{proof}
Let $\mathcal P$ be the set of the pairs $x=(\mu,I)$ of the statement, $\Psi_x=\Psi(\mu,I)$, and
$\mathcal M$ the set of the monomials of row degrees $B$. By Lemma~\ref{pd-symbols-path-count},
$\mathcal P$ and $\mathcal M$ are finite with the same number of elements.

Each $\Psi_x$ has row degrees $B$. Fix $c$ and a letter $i$. The step at $i$ for $c$ multiplies by an
entry of $z_i$ if $h_i(c)\neq h_{i+1}(c)$ and does nothing otherwise, so every monomial of
$\mathrm{Col}_{\mu,c}(J)$ has degree $1$ in $z_i$ in the first case and $0$ in the second. Hence every
monomial of $\Psi_x$ has degree in $z_i$ equal to the number of $c$ with $|c|\leq T_\mu$ and
$h_i(c)\neq h_{i+1}(c)$. If $i$ is a vector letter, then $\mu_i\leq\mu_{i+1}$ entrywise, so
$h_{i+1}(c)-h_i(c)$ is the number of $a$ with $\mu_i(a)<c\leq\mu_{i+1}(a)$; such $c$ satisfy
$|c|\leq T_\mu$, and the difference is $0$ or $1$, as $\mu_{i+1}(h+2)\leq\mu_i(h+1)<c$ for
$h=h_i(c)$. So
the number of these $c$ is $\sum_a(\mu_{i+1}(a)-\mu_i(a))=B(i)$; likewise for a covector letter, with
$\mu_i$ and $\mu_{i+1}$ exchanged.

The column functions are formed from the variables by sums of products with signs, so $\Psi_x\in R[z]$
is the image of $\Psi_x\in\mathbf Z[z]$. Write $\Psi_x=\sum_{m\in\mathcal M}C_{x,m}m$ in $\mathbf Z[z]$,
and let $C$ be the square integer matrix $(C_{x,m})$, after choosing a bijection of $\mathcal P$ with
$\mathcal M$. Let $\ell$ be a prime. By Lemma~\ref{pd-symbols-path-independent} over $\mathbf F_\ell$,
the images of the $\Psi_x$ in $\mathbf F_\ell[z]$ are linearly independent, so the rows of $C$ modulo
$\ell$ are linearly independent and $\ell$ does not divide $\det C$. Hence $\det C=\pm1$, and $C$ is
invertible over every commutative ring $R$. Over $R$, $\Psi_x=\sum_mC_{x,m}m$, and the monomials in
$\mathcal M$ form a basis of the $R$-module of the polynomials of row degrees $B$; as $C$ is invertible,
so do the $\Psi_x$.
% Lean: word_basis_of_path_restrictions transports GelfandTsetlinBases.Top.psiPath_basis along the renaming
%   rename (iv N n) and the index equivalence e17 (restPathTo, restPathFrom, restStd, restPsi, restPol);
%   psiPath_basis = path_basis_theorem = path_basis_of_count lemmaA2indep pathIdxFinite polBRank with
%   countPAll, through MinorDeterminantBasis.basis_of (det_unit, rows_indep); the row degrees are
%   pathProductDegrees (path_ncard), used in psi_support.
\end{proof}

\begin{lemma}\label{pd-symbols-path-truncate}
Let $A$ be a commutative ring, $N\geq0$ and $n\geq1$ integers with a lettering of $\{0,\dots,N-1\}$,
$g\in GL_n(A)$, $0\leq i\leq N$ and $S=\{0,\dots,i-1\}$, and let $\Sigma^S_g$ be the $A$-algebra
endomorphism of $A[z]$ that replaces each row $z_l$ with $l\in S$ by $z_lg$ if $l$ is a vector letter and
by $z_l\,{}^tg^{-1}$ if it is a covector letter, and fixes the other rows; we use the notation of
Definitions~\ref{pd-symbols-path-generators}, \ref{pd-symbols-path-weights}
and~\ref{pd-symbols-path-canonical}. Let $\mu$ be a path of rank $n$ for some $B$ and $I$ a filling of
shape $\mu_N$. Then every monomial $z^m$ of $\Sigma^S_g(\Psi(\mu,I))$ satisfies
$\mathrm{wt}_S(m)\preceq\mu_i$.
\end{lemma}
\leanref{Formalism/SharpRestriction/GrowthBounds/BinomialPolynomialIdentities/RestSplit.lean}{Formalism.MainNotions.G.ColumnFamilies.RestWord}
\leanref{Formalism/SharpRestriction/GrowthBounds/BinomialPolynomialIdentities/ExistsDominantTopWeightNonzero.lean}{Formalism.MainNotions.G.ColumnFamilies.restWord}
\leanref{Formalism/ThetaLogCoordinates/JacobiPathWeights/JacobiMinorProofs/ValuationProductsComponents.lean}{Formalism.InvariantFactorPaths.GelfandTsetlinBases.A0.trunc}

\begin{proof}
Let $\mu'=(\mu_0,\dots,\mu_i)$, a path for the restriction of $B$ to $S$, with the lettering of $S$. Fix
$c$ and let $F_j$ be the functions of the definition of $\mathrm{Col}_{\mu,c}$; then
$F_i=\mathrm{Col}_{\mu',c}$. Each step at a letter $l\geq i$ is linear in $F$, with coefficients that are
entries of the row $z_l$. So there are polynomials $Q_c(J',J)$ in the rows $z_l$, $l\geq i$, with
$\mathrm{Col}_{\mu,c}(J)=\sum_{J'}\mathrm{Col}_{\mu',c}(J')Q_c(J',J)$. Multiplying over $|c|\leq T_\mu$,
\[
 \Psi(\mu,I)=\sum_{I'}\Big(\prod_{|c|\leq T_\mu}\mathrm{Col}_{\mu',c}(I'(c))\Big)\prod_{|c|\leq T_\mu}Q_c(I'(c),I(c)),
\]
the sum over the functions $I'$ on $\{c:|c|\leq T_\mu\}$. By induction over the steps, $F_i(J)=0$ unless
$|J|=h_i(c)$ for $c\geq1$ and $|J|=n-h_i(c)$ for $c\leq0$, as in the proof of
Lemma~\ref{pd-symbols-path-word-bounds}. So only the $I'$ that, extended by $\emptyset$, are fillings of
shape $\mu_i$ contribute; for them $I'(c)=\emptyset$ for $T_{\mu'}<|c|\leq T_\mu$, where
$\mathrm{Col}_{\mu',c}(\emptyset)=1$, and the first product is $\Psi(\mu',I')$.

The map $\Sigma^S_g$ fixes the $Q_c$, which involve only rows outside $S$, and acts on $\Psi(\mu',I')$ as
the substitution $\Sigma_g$ of Lemma~\ref{pd-symbols-path-action} for the lettering of $S$. A monomial of
$\Sigma^S_g(\Psi(\mu,I))$ is therefore a product $z^{m_1}z^{m_2}$ of a monomial $z^{m_1}$ of
$\Sigma_g(\Psi(\mu',I'))$, for some such $I'$, and a monomial $z^{m_2}$ in the rows outside $S$, of
$S$-weight $0$. By Lemma~\ref{pd-symbols-path-top-weight} for the path $\mu'$, $\mathrm{wt}_S(m_1)\preceq
\mu_i$.
% Lean: restWord, with A0.trunc (the Q_c), A0.pathWord_card (the sizes), partialAct_eq_fullAct_of_vars_lt,
%   partialAct_fix_of_vars_ge, partialWeight_eq_zero_of_ge and fullAct_pathWord_support_dominated (FactEV).
\end{proof}

\begin{lemma}\label{pd-symbols-path-lower-top}
Let $F$ be a field, $N\geq0$ and $n\geq1$ integers with a lettering of $\{0,\dots,N-1\}$,
$B\colon\{0,\dots,N-1\}\to\mathbf Z_{\geq0}$, $\pi\subset\mathbf Z^n$ a lower set and $S=\{0,\dots,N-1\}$; we
use the notation of Definitions~\ref{pd-symbols-path-generators}, \ref{pd-symbols-path-weights}
and~\ref{pd-symbols-path-canonical}. Let $L$ and $\mathfrak g$ be as in
Lemma~\ref{pd-symbols-path-brackets} and $\Sigma_{\mathfrak g}$ as in Lemma~\ref{pd-symbols-path-action}
over $L$. Let $G\in F[z]$ have row degrees $B$, and suppose that every monomial $z^m$ of
$\Sigma_{\mathfrak g}(G)\in L[z]$ with $\mathrm{wt}_S(m)$ dominant has $\mathrm{wt}_S(m)\in\pi$. Then $G$
lies in the $F$-span of the path generators $\Psi(\mu,I)$ with $\mu$ a path of rank $n$ for $B$ such that
$\mu_N\in\pi$ and $I$ a standard filling of $\mu$.
\end{lemma}
\leanref{Formalism/SharpRestriction/GrowthBounds/BinomialPolynomialIdentities/RestSplit.lean}{Formalism.MainNotions.G.ColumnFamilies.RestTop}
\leanref{Formalism/SharpRestriction/GrowthBounds/BinomialPolynomialIdentities/ExistsDominantTopWeightNonzero.lean}{Formalism.MainNotions.G.ColumnFamilies.restTop}
\leanref{Formalism/SharpRestriction/GrowthBounds/BinomialPolynomialIdentities/ExistsDominantTopWeightNonzero.lean}{Formalism.MainNotions.G.ColumnFamilies.exists_dominant_top_weight_nonzero}

\begin{proof}
By Lemma~\ref{pd-symbols-path-basis}, $G=\sum_xc_x\Psi_x$ over the pairs $x=(\mu,I)$ of a path for $B$ and
a standard filling, with $c_x\in F$ almost all $0$. Let $x_1$ have $c_{x_1}\neq0$, and let $\nu_1$ be the
last vector of its path. Let $\operatorname{ht}$ be as in the proof of
Lemma~\ref{pd-symbols-path-independent}. Among the $x$ with $c_x\neq0$ whose path ends at some $\nu$ with
$\nu_1\preceq\nu$, choose $x_0$ with $\operatorname{ht}(\nu)$ largest, and let $\mu_0$ be the end of its
path; it is dominant, and $\nu_1\preceq\mu_0$. For $P\in L[z]$ let $P_{\mu_0}$ be the sum of its terms of
$S$-weight $\mu_0$.

Let $x=(\mu,I)$ with $c_x\neq0$ and $\Sigma_{\mathfrak g}(\Psi_x)_{\mu_0}\neq0$. By
Lemma~\ref{pd-symbols-path-top-weight}, $\mu_0\preceq\mu_N$, so $\nu_1\preceq\mu_N$, and by the choice of
$x_0$, $\operatorname{ht}(\mu_N)\leq\operatorname{ht}(\mu_0)$; as in the proof of
Lemma~\ref{pd-symbols-path-independent}, $\mu_N=\mu_0$. So, by the same lemma,
\[
 \Sigma_{\mathfrak g}(G)_{\mu_0}=\sum_\mu\Big(\sum_Ic_{(\mu,I)}[\mu_0|I]_{\mathfrak g}\Big)
 \Psi(\mu,\operatorname{can}_{\mu_0}),
\]
over the paths $\mu$ for $B$ with $\mu_N=\mu_0$ and the standard fillings $I$ of shape $\mu_0$. By
Lemmas~\ref{pd-symbols-path-canonical-independent} and~\ref{pd-symbols-path-brackets}, this is nonzero,
as $c_{x_0}\neq0$. So $\Sigma_{\mathfrak g}(G)$ has a monomial of $S$-weight $\mu_0$, and $\mu_0\in\pi$ by
the hypothesis. As $\pi$ is a lower set, $\nu_1\in\pi$. Hence every $x$ with $c_x\neq0$ has its path
ending in $\pi$.
% Lean: restTop, with Top.psiPath_basis and exists_dominant_top_weight_nonzero (the argument of
%   lemmaA2indep_of with the filter DomLe (x1 N) (x N)), fam_li fact_break.
\end{proof}

\begin{lemma}\label{pd-symbols-path-lower-set}
Let $F$ be a field, $N\geq0$ and $n\geq1$ integers with a lettering of $\{0,\dots,N-1\}$,
$B\colon\{0,\dots,N-1\}\to\mathbf Z_{\geq0}$, $0\leq i\leq N$ and $\pi\subset\mathbf Z^n$ a lower set
(Definition~\ref{pd-symbols-path-weights}(4)). Let $G\in F[z]$ have row degrees $B$, and put
$S=\{0,\dots,i-1\}$. Then $G$ lies in the $F$-span of the path generators $\Psi(\mu,I)$, with $\mu$ a
path of rank $n$ for $B$ such that $\mu_i\in\pi$ and $I$ a standard filling of $\mu$, if and only if
every dominant weight of $U_S(G)$ lies in $\pi$.
\end{lemma}
\leanref{Formalism/SharpRestriction/GrowthBounds/StatementOfDegreeAndBound.lean}{Formalism.MainNotions.CovariantDegrees.FirstLettersLowerSet}
\leanref{Formalism/SharpRestriction/GrowthBounds/BinomialPolynomialIdentities/AdmissibleWordBasis.lean}{Formalism.MainNotions.G.ColumnFamilies.admissible_word_basis}
\leanref{Formalism/SharpRestriction/GrowthBounds/BinomialPolynomialIdentities/TorusComponentMemDominatedWeightPolynomials.lean}{Formalism.MainNotions.G.ColumnFamilies.firstLettersLowerSet_of_restWord_restOne}
\leanref{Formalism/SharpRestriction/GrowthBounds/BinomialPolynomialIdentities/TorusComponentMemDominatedWeightPolynomials.lean}{Formalism.MainNotions.G.ColumnFamilies.restFwd}
\leanref{Formalism/SharpRestriction/GrowthBounds/BinomialPolynomialIdentities/CardPathIndexWeightMemEqMulChoose.lean}{Formalism.MainNotions.G.ColumnFamilies.restOne_of}
\leanref{Formalism/SharpRestriction/GrowthBounds/BinomialPolynomialIdentities/CardPathIndexWeightMemEqMulChoose.lean}{Formalism.MainNotions.G.ColumnFamilies.restSplit}
\leanref{Formalism/SharpRestriction/GrowthBounds/BinomialPolynomialIdentities/CardPathIndexWeightMemEqMulChoose.lean}{Formalism.MainNotions.G.ColumnFamilies.card_pathIndex_weight_mem_eq_mul_choose}

\begin{proof}
We use the notation of Definitions~\ref{pd-symbols-path-generators}, \ref{pd-symbols-path-weights}
and~\ref{pd-symbols-path-canonical}, and of Lemma~\ref{pd-symbols-path-translate} for the substitutions
$\Sigma^\varphi_g$, formed with $S$. Let $\mathcal P_i(\pi)$ be the set of the pairs $x=(\mu,I)$ of a path
$\mu$ for $B$ with $\mu_i\in\pi$ and a standard filling $I$ of $\mu$, $\Psi_x=\Psi(\mu,I)$, and $U$ the
$F$-span of these $\Psi_x$.

\emph{The span has its weights in $\pi$.} Let $\mathcal Q$ be the set of the $x\in F[z]$ such that for
every commutative ring $B'$, ring homomorphism $\varphi\colon F\to B'$ and $g\in GL_n(B')$, every monomial
$z^m$ of $\Sigma^\varphi_g(x)$ satisfies $\mathrm{wt}_S(m)\preceq\nu$ for some $\nu\in\pi$. It is an
$F$-subspace of $F[z]$. For $x\in\mathcal P_i(\pi)$, $\Sigma^\varphi_g(\Psi_x)$ is the substitution
$\Sigma^S_g$ of Lemma~\ref{pd-symbols-path-truncate} applied to the image of $\Psi_x$ in $B'[z]$, so its
monomials have $S$-weight $\preceq\mu_i\in\pi$, and $\Psi_x\in\mathcal Q$. Suppose $G\in U$. Then
$G\in\mathcal Q$. The second paragraph of the proof of Lemma~\ref{pd-symbols-path-translate}(3), with the
condition $\mathrm{wt}_S(m)(a)\in E$ replaced by this one, shows that $U_S(G)\subset\mathcal Q$. Let $\mu$
be a dominant weight of $U_S(G)$ and $x\in U_S(G)$ nonzero with all its monomials of $S$-weight $\mu$.
With $B'=F$, $\varphi$ the identity and $g=1$, a monomial of $x$ gives $\mu\preceq\nu$ for some
$\nu\in\pi$, so $\mu\in\pi$.

\emph{The converse.} Let $L$ and $\mathfrak g$ be as in Lemma~\ref{pd-symbols-path-brackets}, with
$\Sigma^S_{\mathfrak g}$ as in Lemma~\ref{pd-symbols-path-truncate} over $L$. Say that $P\in F[z]$
satisfies (C) if every monomial $z^m$ of $\Sigma^S_{\mathfrak g}(P)\in L[z]$ with $\mathrm{wt}_S(m)$
dominant has $\mathrm{wt}_S(m)\in\pi$. Suppose that every dominant weight of $U_S(G)$ lies in $\pi$. By
Lemma~\ref{pd-symbols-path-translate}(d), with $B=L$ and $g=\mathfrak g$, every monomial of
$\Sigma^S_{\mathfrak g}(G)$ is a monomial of the generic translate $\tilde G$, so by part (2) of that
lemma its $S$-weight is a weight of $U_S(G)$; hence $G$ satisfies (C). Each $\Psi_x$ with
$x\in\mathcal P_i(\pi)$ satisfies (C) by Lemma~\ref{pd-symbols-path-truncate}, as $\pi$ is a lower set.

Let $W$ be the $F$-span of the products $\Psi(\mu',I')z^\beta$, with $\mu'$ a path for the restriction of
$B$ to $S$ (with the lettering of $S$) such that $\mu'_i\in\pi$, $I'$ a standard filling of $\mu'$, and
$z^\beta$ a monomial in the rows $z_l$, $l\geq i$, of degree $B(l)$ in each of them. We show that every $P$
of row degrees $B$ satisfying (C) lies in $W$. Write $P=\sum_\beta P_\beta z^\beta$, with $z^\beta$ as
above and $P_\beta$ a polynomial in the rows in $S$ of row degrees $B$ there. The map
$\Sigma^S_{\mathfrak g}$ fixes $z^\beta$, which has $S$-weight $0$, and maps $P_\beta$ to a polynomial in
the rows in $S$; so the monomials of $\Sigma^S_{\mathfrak g}(P)$ are the products of those of the
$\Sigma^S_{\mathfrak g}(P_\beta)$ with the $z^\beta$, with the same $S$-weights, and each $P_\beta$
satisfies (C). On polynomials in the rows in $S$, $\Sigma^S_{\mathfrak g}$ is the substitution
$\Sigma_{\mathfrak g}$ of Lemma~\ref{pd-symbols-path-action} for the lettering of $S$. By
Lemma~\ref{pd-symbols-path-lower-top}, applied to $P_\beta$ with $N$ replaced by $i$, $P_\beta$ lies in
the span of the $\Psi(\mu',I')$ with $\mu'_i\in\pi$. So $P\in W$.

In particular $U\subset W$ and $G\in W$. By Lemma~\ref{pd-symbols-path-independent}, $\dim U$ is the number
of elements of $\mathcal P_i(\pi)$. In the notation of the proof of Lemma~\ref{pd-symbols-path-count}, an
element of $\mathcal P_i(\pi)$ is a path $\mu'$ for the restriction of $B$ to $S$ with $\mu'_i\in\pi$
together with one of the $C_i(\mu'_i)=d_n(\mu'_i)\prod_{l=i}^{N-1}\binom{n+B(l)-1}{B(l)}$ continuations
and fillings, and $d_n(\mu'_i)$ is the number of standard fillings of $\mu'_i$. So $\dim U$ is the number
of pairs $(\mu',I')$ times the number $\prod_{l=i}^{N-1}\binom{n+B(l)-1}{B(l)}$ of monomials $z^\beta$,
which is at least $\dim W$. Hence $U=W$, and $G\in U$.
% Lean: firstLettersLowerSet_of_restWord_restOne D restWord restOne. Forward: restFwd (the property
%   submodule dominatedWeightPolynomials, stable under divPow and binomPol by the Laurent-coefficient
%   argument) with standardWord_monomial_weight_dominated (RestWord). Converse: generic_monomials_specialise
%   and generic_translate_weights_mem give GoodC; restOne = restOne_of restWord restTop restSplit, where
%   U <= W is mem_span_wordMonomialProducts with standardWordFamily_satisfies_condition, the splitting is
%   restSplit (Top.factSplit, partialAction_monomial_later_letters), and the count is
%   card_pathIndex_weight_mem_eq_mul_choose (CountCont = countCont_of stripsMem gtPieri stdCountGT) with
%   ExponentSetCounting.expSet_card; then Submodule.eq_of_le_of_finrank_le.
\end{proof}

\begin{lemma}\label{pd-symbols-path-translate}
Let $F$ be a field, $N\geq0$ and $n\geq1$ integers with a lettering of $\{0,\dots,N-1\}$, and
$S\subset\{0,\dots,N-1\}$; we use the notation of Definitions~\ref{pd-symbols-path-generators}
and~\ref{pd-symbols-path-weights}. For a commutative ring $B$, a ring homomorphism
$\varphi\colon F\to B$ and $g\in GL_n(B)$, let $\Sigma^\varphi_g\colon F[z]\to B[z]$ be the ring
homomorphism that applies $\varphi$ to the coefficients and replaces, for every $i\in S$, the row $z_i$
by $z_ig$ if $i$ is a vector letter and by $z_i\,{}^tg^{-1}$ if it is a covector letter. Let
$A=F[\mathfrak g,\det(\mathfrak g)^{-1}]$, with $\mathfrak g$ the $n\times n$ matrix of indeterminates,
let $G\in F[z]$, and let $\tilde G=\Sigma_{\mathfrak g^{-1}}(G)\in A[z]$, formed with the inclusion
$F\to A$. Let $a\leq n$, and let $\mathcal E_a$ be the set of the $\mathrm{wt}_S(m)(a)$, $z^m$ a monomial
of $\tilde G$. Then:
\begin{enumerate}
\item $\mathcal E_a$ is the set of the exponents of $\eta$ that occur in $G_{S,a}$;
\item for every monomial $z^m$ of $\tilde G$, $\mathrm{wt}_S(m)$ is a weight of $U_S(G)$;
\item if a set $E\subset\mathbf Z$ contains $\mathcal E_a$, then $\mu(a)\in E$ for every weight $\mu$
  of $U_S(G)$.
\end{enumerate}
\end{lemma}
\leanref{Formalism/SharpRestriction/GrowthBounds/RootWeights/TwistSubEqMatrixSubst.lean}{Formalism.MainNotions.G.GLWeightTheory.matrixSubst}
\leanref{Formalism/SharpRestriction/GrowthBounds/RootWeights/BinomialStableContainsComponents.lean}{Formalism.MainNotions.G.GLWeightTheory.genericTranslate}
\leanref{Formalism/SharpRestriction/GrowthBounds/RootWeights/LaurentDegreesSubstitutionSubset.lean}{Formalism.MainNotions.G.GLWeightTheory.twist_exponents_are_weight_entries}
\leanref{Formalism/SharpRestriction/GrowthBounds/RootWeights/GeneratedModuleWeightLemmaOfInterfaces.lean}{Formalism.MainNotions.G.GLWeightTheory.generic_translate_weights_mem}
\leanref{Formalism/SharpRestriction/GrowthBounds/RootWeights/PropertySubmoduleWeightEntryComponent.lean}{Formalism.MainNotions.G.GLWeightTheory.weights_property_of_generic_translate}

\begin{proof}
We write $\Sigma_g$ for $\Sigma^\varphi_g$ when $\varphi$ is clear, and also for the same substitution
on $B[z]$ with $\varphi$ the identity. Then $\Sigma_{g_1}\circ\Sigma_{g_2}=\Sigma_{g_1g_2}$, as
$z_ig_1g_2$ and $z_i\,{}^tg_1^{-1}\,{}^tg_2^{-1}=z_i\,{}^t(g_1g_2)^{-1}$ are the composite
substitutions; and $G_{S,a}=\Sigma_{\tau_a(\eta)}(G)$ over $A[\eta,\eta^{-1}]$, as
$\tau_a(\eta)^{-1}=\tau_a(\eta^{-1})$. We use four facts.

(a) \emph{Diagonal matrices.} If $D$ is the diagonal matrix with entries $u_1,\dots,u_n\in B^\times$,
then $\Sigma^\varphi_D(cz^m)=\varphi(c)\prod_bu_b^{\mathrm{wt}_S(m)(b)}z^m$ for $c\in F$, since the
substitution sends $z_{i,b}$ to $u_bz_{i,b}$ for a vector letter $i\in S$ and to $u_b^{-1}z_{i,b}$ for
a covector letter $i\in S$, and fixes the other variables.

(b) \emph{Transvections.} For $b\neq c$ and $t\in B$ let $T_{bc}(t)=1+te_{bc}$, whose inverse is
$T_{bc}(-t)$. Then $\Sigma^\varphi_{T_{bc}(t)}(P)=\sum_{k\geq0}t^k\varphi(E^{(k)}_{bc}(P))$ for
$P\in F[z]$, a finite sum. Indeed $\Sigma_{T_{bc}(t)}$ replaces $z_{i,c}$ by $z_{i,c}+tz_{i,b}$ for a
vector letter $i\in S$ and $z_{i,b}$ by $z_{i,b}-tz_{i,c}$ for a covector letter $i\in S$, and expanding
the powers of these binomials in a monomial gives the coefficients and the exponents of
Definition~\ref{pd-symbols-path-weights}(2).

(c) \emph{Weight components.} For $x\in F[z]$ and $\nu\in\mathbf Z^n$ let $x_\nu$ be the sum of the
terms of $x$ whose monomials have $S$-weight $\nu$. If $W\subset F[z]$ is a subspace stable under the
$H^{(k)}_b$ and $x\in W$, then $x_\nu\in W$. Let $\mathcal B$ be the additive group of functions
$\mathbf Z\to\mathbf Z$ generated by the $x\mapsto\binom xk$, $k\geq0$. For $f\in\mathcal B$ and
$b\leq n$, the map $z^m\mapsto f(\mathrm{wt}_S(m)(b))z^m$ is an integral combination of the
$H^{(k)}_b$, so it preserves $W$. By Vandermonde's identity
$\binom{x+s}k=\sum_{i+l=k}\binom xi\binom sl$, the functions $x\mapsto\binom{x+s}k$, $s\in\mathbf Z$,
lie in $\mathcal B$. We show that for $w_0\in\mathbf Z$ and $M\geq0$ some $f\in\mathcal B$ agrees on
$[-M,M]$ with the indicator function $\delta_{w_0}$ of $w_0$. For $M=0$ take a multiple of
$\binom x0=1$. Given $f$ for $M$, put $g_1(x)=\binom{x+M}{2M+1}$ and $g_2(x)=\binom{x+M+1}{2M+2}$.
Both vanish on $[-M,M]$, $g_1(M+1)=g_2(M+1)=1$, $g_1(-M-1)=\binom{-1}{2M+1}=-1$ and $g_2(-M-1)=0$.
So $f+c_1g_2+c_2(g_2-g_1)$, with $c_1=\delta_{w_0}(M+1)-f(M+1)$ and
$c_2=\delta_{w_0}(-M-1)-f(-M-1)$, works for $M+1$. Taking $M\geq|\mathrm{wt}_S(m)(b)|$ for all the
monomials $z^m$ of $x$, we find that the sum of the terms of $x$ with $\mathrm{wt}_S(m)(b)=w_0$ lies in
$W$. Applying this for $b=1,\dots,n$ in turn gives $x_\nu\in W$.

(d) \emph{Specialization.} If $g\in GL_n(B)$, then every monomial of $\Sigma^\varphi_g(G)$ is a
monomial of $\tilde G$. Indeed the ring homomorphism $A\to B$ that extends $\varphi$ and sends
$\mathfrak g$ to $g^{-1}$ exists, as $\det(g^{-1})$ is a unit; it sends $\mathfrak g^{-1}$ to $g$,
so applied to the coefficients it maps $\tilde G$ to $\Sigma^\varphi_g(G)$.

(1) Let $D=D_a(\eta)$ over $B=A[\eta,\eta^{-1}]$. By (a), $\Sigma_D(\tilde G)$ is the sum of the
terms $\tilde G_m\eta^{\mathrm{wt}_S(m)(a)}z^m$, with $\tilde G_m\in A$ the coefficients of $\tilde G$,
so the exponents of $\eta$ that occur in it form the set $\mathcal E_a$. As
$\tau_a(\eta)=\mathfrak gD\mathfrak g^{-1}$, $G_{S,a}=\Sigma_{\mathfrak g}(\Sigma_D(\tilde G))$. If
$Q=\sum_k\eta^kQ_k$ with $Q_k\in A[z]$, then $\Sigma_{\mathfrak g}(Q)=\sum_k\eta^k\Sigma_{\mathfrak
g}(Q_k)$, so every exponent of $\eta$ that occurs in $\Sigma_{\mathfrak g}(Q)$ occurs in $Q$; as
$\Sigma_{\mathfrak g^{-1}}$ undoes $\Sigma_{\mathfrak g}$, the two sets of exponents are equal.

(2) Let $W=U_S(G)$; it is stable under the $E^{(k)}_{bc}$ and the $H^{(k)}_b$, and contains $G$. The
ring $A$ is a domain; let $K$ be its fraction field and $W_K\subset K[z]$ the $K$-span of $W$. By (b),
$\Sigma_{T_{bc}(t)}$ maps $W_K$ into itself for $t\in K$. If $D$ is diagonal with entries
$u_b\in K^\times$ and $x\in W$, then $\Sigma_D(x)=\sum_\nu\big(\prod_bu_b^{\nu(b)}\big)x_\nu$ by (a),
which lies in $W_K$ by (c); so $\Sigma_D$ maps $W_K$ into itself. By Gaussian elimination,
$\mathfrak g^{-1}\in GL_n(K)$ is a product of matrices $T_{bc}(t)$ and an invertible diagonal matrix, so
$\Sigma_{\mathfrak g^{-1}}$, the composite of the corresponding substitutions, maps $W_K$ into itself,
and $\tilde G\in W_K$. Let $z^m$ be a monomial of $\tilde G$. As $\tilde G$ is a $K$-linear combination of
elements of $W$, $z^m$ is a monomial of some $x\in W$. Then $x_{\mathrm{wt}_S(m)}$ lies in $W$ by (c), is
nonzero, and all its monomials have $S$-weight $\mathrm{wt}_S(m)$; so $\mathrm{wt}_S(m)$ is a weight of
$U_S(G)$.

(3) Let $\mathcal P$ be the set of the $x\in F[z]$ such that for every commutative ring $B$, every ring
homomorphism $\varphi\colon F\to B$ and every $g\in GL_n(B)$, each monomial $z^m$ of
$\Sigma^\varphi_g(x)$ has $\mathrm{wt}_S(m)(a)\in E$. It is an $F$-subspace of $F[z]$, and $G\in\mathcal P$
by (d), as $E\supset\mathcal E_a$.

Let $x\in\mathcal P$, and suppose that for every $B$ and $\varphi$ there are $h\in GL_n(B[T,T^{-1}])$ and
$x_k\in F[z]$, almost all $0$ and independent of $B$ and $\varphi$, with
$\Sigma^\varphi_h(x)=\sum_kT^k\varphi(x_k)$. Then each $x_k$ lies in $\mathcal P$: for $g\in GL_n(B)$,
$\Sigma_{gh}(x)=\Sigma_g(\Sigma_h(x))=\sum_kT^k\Sigma^\varphi_g(x_k)$; every monomial of this has
$\mathrm{wt}_S(m)(a)\in E$, as $x\in\mathcal P$; and a monomial of $\Sigma^\varphi_g(x_k)$ is a monomial of
the sum, since the coefficient of $T^k$ in its coefficient there is its coefficient in
$\Sigma^\varphi_g(x_k)$. By (b) with $h=T_{bc}(T)$, $E^{(k)}_{bc}(x)\in\mathcal P$ for all $b\neq c$ and
$k$. By (a) with $h$ the diagonal matrix with $T$ in place $b$ and $1$ elsewhere, the sum $x^{(w)}$ of
the terms of $x$ with $\mathrm{wt}_S(m)(b)=w$ lies in $\mathcal P$ for every $w$, and so does
$H^{(k)}_b(x)=\sum_w\binom wkx^{(w)}$. Hence $U_S(G)\subset\mathcal P$.

Let $\mu$ be a weight of $U_S(G)$, and $x\in U_S(G)$ nonzero with all its monomials of $S$-weight $\mu$.
With $B=F$, $\varphi$ the identity and $g=1$, a monomial $z^m$ of $x$ gives
$\mu(a)=\mathrm{wt}_S(m)(a)\in E$.
% Lean: (a) torus_acts_through_weight; (b) transvectionExpansion (matrixSubst_transvection_monomial);
%   (c) binomial_stable_contains_components (weightEntry_restriction_mem, binomialFunctions_window_indicator,
%   diagonalMultiplier_mem); (d) generic_monomials_specialise; (1) twist_exponents_are_weight_entries
%   (laurent_degrees_substitution_eq); (2) generic_translate_weights_mem (Matrix.Pivot.
%   exists_list_transvec_mul_diagonal_mul_list_transvec, scalarExtension_monomial_weights); (3)
%   weights_property_of_generic_translate (weightEntryPropertySubmodule, laurent_coefficients_mem_property,
%   propertySubmodule_dividedPower_stable, propertySubmodule_binomial_stable).
\end{proof}

\begin{lemma}\label{pd-symbols-path-valuation}
Let $F$ be a field, $N\geq0$ and $n\geq1$ integers with a lettering of $\{0,\dots,N-1\}$,
$S\subset\{0,\dots,N-1\}$, $G\in F[z]$ and $\theta,M\in\mathbf Z$, with the notation of
Definitions~\ref{pd-symbols-path-generators} and~\ref{pd-symbols-path-weights}. Then
$\theta\leq\operatorname{val}_S(G)$ if and only if $\theta\leq\mu(n)$ for every weight $\mu$ of $U_S(G)$,
and $\lambda_S(G)\leq M$ if and only if $\mu(1)\leq M$ for every weight $\mu$ of $U_S(G)$.
\end{lemma}
\leanref{Formalism/SharpRestriction/GrowthBounds/CovariantDegrees.lean}{Formalism.MainNotions.CovariantDegrees.ValuationsAdditive}
\leanref{Formalism/SharpRestriction/GrowthBounds/GLWeightTheory.lean}{Formalism.MainNotions.G.GLWeightTheory.valuation_field}

\begin{proof}
Let $\tilde G$ and the sets $\mathcal E_a$ be as in Lemma~\ref{pd-symbols-path-translate}. By its
part (1), $\operatorname{val}_S(G)$ is the least element of $\mathcal E_n$ and $\lambda_S(G)$ the largest
element of $\mathcal E_1$, with $\min\emptyset=+\infty$ and $\max\emptyset=-\infty$. Suppose
$\theta\leq\operatorname{val}_S(G)$. Then the set $E$ of the integers at least $\theta$ contains
$\mathcal E_n$, and part (3) gives $\theta\leq\mu(n)$ for every weight $\mu$ of $U_S(G)$. Conversely,
suppose that $\theta\leq\mu(n)$ for every weight $\mu$ of $U_S(G)$. By part (2), every element of
$\mathcal E_n$ is $\mu(n)$ for some weight $\mu$, hence at least $\theta$; so
$\theta\leq\operatorname{val}_S(G)$. The statement for $\lambda_S$ is proved in the same way, with $a=1$
and the set of the integers at most $M$.
% Lean: GLWeightTheory.generated_module_weight_lemma_of_interfaces (through valuationsAdditive_D7 and
%   generated_module_weight_lemma_part_b).
\end{proof}

\begin{lemma}\label{pd-symbols-path-word-bounds}
Let $N\geq0$ and $n\geq1$ be integers with a lettering of $\{0,\dots,N-1\}$,
$B\colon\{0,\dots,N-1\}\to\mathbf Z_{\geq0}$, $\mu$ a path of rank $n$ for $B$ and $I$ a function from
$\mathbf Z$ to the subsets of $\{1,\dots,n\}$, with the notation of
Definition~\ref{pd-symbols-path-generators}. Then for every letter $i$ the polynomial
$\Psi(\mu,I)\in\mathbf Z[z]$ satisfies $\mu_{i+1}(n)\leq\operatorname{val}_{\{0,\dots,i\}}(\Psi(\mu,I))$ and
$\lambda_{\{0,\dots,i\}}(\Psi(\mu,I))\leq\mu_{i+1}(1)$.
\end{lemma}
\leanref{Formalism/SharpRestriction/GrowthBounds/FrobeniusPaths/ComponentsInFiltrationContractionSingle.lean}{Formalism.MainNotions.LemmaBAlt.WordValuationBounds}
\leanref{Formalism/SharpRestriction/GrowthBounds/FrobeniusPaths/ComponentsInFiltrationContractionSingle.lean}{Formalism.MainNotions.LemmaBAlt.word_valuation_bounds}

\begin{proof}
We use the notation of Definition~\ref{pd-symbols-path-generators}, over the ring
$A=\mathbf Z[\mathfrak g,\det(\mathfrak g)^{-1}][\eta,\eta^{-1}]$. For $\varsigma=\pm1$ and
$s\in\mathbf Z$ let $F^\varsigma_s$ be the set of the polynomials in the $z$ over $A$ in which every
exponent $k$ of $\eta$ that occurs satisfies $\varsigma k\geq s$. Then
$F^\varsigma_sF^\varsigma_{s'}\subset F^\varsigma_{s+s'}$, each $F^\varsigma_s$ is closed under sums,
$F^\varsigma_s\subset F^\varsigma_{s'}$ for $s'\leq s$, the variables and the elements of
$\mathbf Z[\mathfrak g,\det(\mathfrak g)^{-1}]$ lie in $F^\varsigma_0$, and $\eta^k\in F^\varsigma_{\varsigma k}$.
For $P\in\mathbf Z[z]$, $s\leq\operatorname{val}_S(P)$ if and only if $P_{S,n}\in F^1_s$, and
$\lambda_S(P)\leq M$ if and only if $P_{S,1}\in F^{-1}_{-M}$.

Let $\mu'=\mu_{i+1}$, $S=\{0,\dots,i\}$ and $a\in\{1,n\}$. Let $u$ be the $a$-th column of $\mathfrak g$ and
$v$ the $a$-th row of $\mathfrak g^{-1}$, so that $\sum_mu_mv_m=1$ and
$\tau_a(\eta)=1+(\eta-1)uv$. Substituting $z_l\tau_a(\eta)$ for the row $z_l$ of a vector letter $l$
replaces $z_{l,m}$ by $z_{l,m}+(\eta-1)v_m\sum_bu_bz_{l,b}$, and substituting
$z_l\,{}^t\tau_a(\eta^{-1})$ for that of a covector letter replaces it by
$z_{l,m}+(\eta^{-1}-1)u_m\sum_bv_bz_{l,b}$. For functions $F$ from the subsets of $\{1,\dots,n\}$ to $A[z]$
and rows $x,y$ put $W_xF=F\wedge x$, $C_yF=F\lrcorner y$ and
$\Lambda_\theta(y,x)=1+(\theta-1)W_xC_y$. With $\delta_m$ the $m$-th unit row, the operators
$e_mF=F\wedge\delta_m$ and $\iota_mF=F\lrcorner\delta_m$ satisfy $e_me_b=-e_be_m$,
$\iota_m\iota_b=-\iota_b\iota_m$ and $\iota_me_b+e_b\iota_m=\delta_{mb}$, and $W_x=\sum_mx_me_m$,
$C_y=\sum_my_m\iota_m$. Hence $W_x^2=0$, $C_y^2=0$ and $C_yW_x+W_xC_y=\sum_my_mx_m$, and
from these, for $\Lambda=\Lambda_\theta(y,x)$ with $\sum_my_mx_m=1$ and $\theta\theta'=1$,
\[
 \Lambda e_b=(e_b+(\theta-1)y_bW_x)\Lambda,\qquad \Lambda\iota_b=(\iota_b+(\theta'-1)x_bC_y)\Lambda .
\]
Multiplying by the entries of a row $w$ and summing over $b$ gives $\Lambda(F\wedge w)=(\Lambda
F)\wedge w'$ with $w'_m=w_m+(\theta-1)x_m\sum_by_bw_b$, and $\Lambda(F\lrcorner w)=(\Lambda F)\lrcorner w''$
with $w''_m=w_m+(\theta'-1)y_m\sum_bx_bw_b$.

Fix $c$, put $\Lambda=\Lambda_\eta(u,v)$ if $c\geq1$ and $\Lambda=\Lambda_{\eta^{-1}}(v,u)$ if $c\leq0$, and
let $F_j$ be the function after $j$ steps of the definition of $\mathrm{Col}_{\mu,c}$, with the rows
$z_l$. Write $h_j=h_j(c)$. If $\beta$ is a strip over $\alpha$ as in
Definition~\ref{pd-symbols-path-generators}(2), both dominant, then $\alpha\leq\beta$ entrywise, and the
numbers $h\leq h'$ of their entries at least $c$ satisfy $h'\leq h+1$, as $\beta(h+2)\leq\alpha(h+1)<c$.
So $h_j\leq h_{j+1}\leq h_j+1$ at a vector letter $j$ and $h_{j+1}\leq h_j\leq h_{j+1}+1$ at a covector
letter. Hence a step at a letter of $S$ is a wedge with
$\Lambda=\Lambda_\eta(u,v)$ or a contraction with $\Lambda=\Lambda_{\eta^{-1}}(v,u)$ at a vector letter,
and a contraction with $\Lambda_\eta(u,v)$ or a wedge with $\Lambda_{\eta^{-1}}(v,u)$ at a covector letter;
in each case the twisted row is the $w'$ or $w''$ above. As $\Lambda F_0=F_0$, induction on $j$ gives
$(F_j(J))_{S,a}=(\Lambda F_j)(J)$ for $j\leq i+1$. The values of $F_j$ lie in $F^\varsigma_0$. By induction, $F_j(J)=0$
unless $|J|=h_j(c)$ for $c\geq1$, and unless $|J|=n-h_j(c)$ for $c\leq0$, as $F_0$ is supported at
$\emptyset$, $h_0(c)$ is $0$ or $n$, and the heights change by at most $1$ at each letter. If
$W_xF=0$, then $\Lambda_\theta(y,x)F=\theta F$, as $W_xC_y=1-C_yW_x$; if $C_yF=0$, then $\Lambda_\theta(y,x)F=F$;
and in general, if $\theta\in F^\varsigma_s$, then the values of $\Lambda_\theta(y,x)F_j$ lie in
$F^\varsigma_{\min(s,0)}$.

For $t\in\mathbf Z$ let $\varepsilon(t,c)$ be $1$ if $1\leq c\leq t$, $-1$ if $t<c\leq0$, and $0$
otherwise; then $\sum_{|c|\leq T}\varepsilon(t,c)=t$ for $|t|\leq T$.

Take $a=n$ and $\varsigma=1$. As $\mu'$ is dominant, $h_{i+1}(c)=n$ if and only if $c\leq\mu'(n)$. If $c\geq1$ and
$c\leq\mu'(n)$, then $F_{i+1}$ is supported on sets of $n$ elements, $W_vF_{i+1}=0$, and
$\Lambda F_{i+1}=\eta F_{i+1}$ has values in $F^1_1$. If $c\geq1$ and $c>\mu'(n)$, the values lie in
$F^1_0$. If $c\leq0$ and $c\leq\mu'(n)$, then $F_{i+1}$ is supported at $\emptyset$, $C_vF_{i+1}=0$, and
$\Lambda F_{i+1}=F_{i+1}$ has values in $F^1_0$. If $c\leq0$ and $c>\mu'(n)$, the values lie in $F^1_{-1}$, as
$\eta^{-1}\in F^1_{-1}$. In all cases $(F_{i+1}(J))_{S,n}\in F^1_{\varepsilon(\mu'(n),c)}$. The steps at the
letters after $i$ use untwisted rows, whose entries lie in $F^1_0$, so
$(\mathrm{Col}_{\mu,c}(J))_{S,n}\in F^1_{\varepsilon(\mu'(n),c)}$ for every $J$. As $P\mapsto P_{S,n}$ is a ring
homomorphism, $\Psi(\mu,I)_{S,n}\in F^1_s$ with $s=\sum_{|c|\leq T_\mu}\varepsilon(\mu'(n),c)=\mu'(n)$, so
$\mu'(n)\leq\operatorname{val}_S(\Psi(\mu,I))$.

Take $a=1$ and $\varsigma=-1$; now $h_{i+1}(c)=0$ if and only if $c>\mu'(1)$. If $c\geq1$ and $c\leq\mu'(1)$, the
values of $\Lambda F_{i+1}$ lie in $F^{-1}_{-1}$, as $\eta\in F^{-1}_{-1}$. If $c\geq1$ and $c>\mu'(1)$, then
$F_{i+1}$ is supported at $\emptyset$ and $\Lambda F_{i+1}=F_{i+1}$. If $c\leq0$ and $c\leq\mu'(1)$, the values
lie in $F^{-1}_0$, as $\eta^{-1}\in F^{-1}_1$. If $c\leq0$ and $c>\mu'(1)$, then $F_{i+1}$ is supported on sets
of $n$ elements and $\Lambda F_{i+1}=\eta^{-1}F_{i+1}$ has values in $F^{-1}_1$. So
$(\mathrm{Col}_{\mu,c}(J))_{S,1}\in F^{-1}_{-\varepsilon(\mu'(1),c)}$, and as before
$\Psi(\mu,I)_{S,1}\in F^{-1}_{-\mu'(1)}$, that is, $\lambda_S(\Psi(\mu,I))\leq\mu'(1)$.
\end{proof}

\begin{lemma}\label{pd-symbols-path-lift-twist}
Let $p\geq5$ be a prime, $(W,Q,\sigma,\iota)$ a Hida setting of degree $n'$ at $p$ with rings
$\mathcal R\subset\mathcal K$ and $\varpi_0$ (Definition~\ref{pd-local-model}), $n\geq1$ and
$M\in GL_n(W)$, and for $s\in\mathbf Z$ let $M_s=\iota\sigma^s(M)$; we use the notation of
Definition~\ref{pd-symbols-path-canonical}(3). Let $F$ be a function from the subsets of
$\{1,\dots,n\}$ to $\mathcal R[z]$, $o\geq0$ an integer, $s\leq s'$ integers and $e=s'-s$, and suppose
that every value of $M_s\cdot F$ lies in $\varpi_0^o\mathcal R[z]$. Let $F^{p^e}(J)=F(J)^{p^e}$. Then:
\begin{enumerate}
\item every value of $F$ lies in $\varpi_0^o\mathcal R[z]$;
\item every value of $M_{s'}\cdot F^{p^e}$ lies in $\varpi_0^{op^e}\mathcal R[z]$;
\item if $(M_s\cdot F)(J)\in\varpi_0^{o+1}\mathcal R[z]$, then $(M_{s'}\cdot F^{p^e})(J)$ lies in
  $\varpi_0^{op^e+1}\mathcal R[z]$, and in $\varpi_0^{op^e+2}\mathcal R[z]$ if $e\geq1$.
\end{enumerate}
\end{lemma}
\leanref{Formalism/SharpRestriction/GrowthBounds/FrobeniusPaths/SelfDualLatticesOneOrbit.lean}{Formalism.MainNotions.LemmaBAlt.step1}
\leanref{Formalism/SharpRestriction/GrowthBounds/FrobeniusPaths/SelfDualLatticesOneOrbit.lean}{Formalism.MainNotions.LemmaBAlt.sum_pow_dvd}
\leanref{Formalism/SharpRestriction/GrowthBounds/FrobeniusPaths/LiftedColumnOrder.lean}{Formalism.MainNotions.LemmaBAlt.FrobCong}
\leanref{Formalism/SharpRestriction/GrowthBounds/FrobeniusPaths/LatticeCountingFormula.lean}{Formalism.MainNotions.LemmaBAlt.frobCong}

\begin{proof}
Write $\varpi=\varpi_0$. By the proof of Lemma~\ref{pd-symbols-path-exterior}, $h\cdot(h'\cdot F)=(h'h)\cdot F$
for $n\times n$ matrices $h,h'$, as $\Lambda(h)\Lambda(h')=\Lambda(h'h)$. Since $\sigma$ is a ring
automorphism, $(M^{-1})_s=M_s^{-1}$, whose entries lie in $\mathcal R$. (1) follows from
$F=M_s^{-1}\cdot(M_s\cdot F)$.

For (2) and (3) write $M_s\cdot F=\varpi^oZ$, so that $F(I)=\varpi^o\sum_JZ(J)a_{JI}$ with
$a_{JI}=\det(M_s^{-1})_{J,I}$. Put $b_{JI}=\det(M_{s'}^{-1})_{J,I}$. The minors of $M_s^{-1}$ are the images
under $\iota\sigma^s$ of those of $M^{-1}\in GL_n(W)$. For $\beta\in W$, $\sigma^s(\beta)\in W$, and
iterating $\sigma(x)\equiv x^p\pmod{pW}$ gives $\sigma^s(\beta)^{p^e}\equiv\sigma^{s'}(\beta)\pmod{pW}$, as
the $p$-th power map is a ring endomorphism modulo $p$. Hence $a_{JI}^{p^e}\equiv b_{JI}\pmod{p\mathcal R}$.
Also $(\sum_lx_l)^{p^e}\equiv\sum_lx_l^{p^e}\pmod{p}$ in any commutative ring. So
$F(I)^{p^e}=\varpi^{op^e}\big(\sum_JZ(J)^{p^e}b_{JI}+p\rho(I)\big)$ with $\rho(I)\in\mathcal R[z]$, and
\[
 M_{s'}\cdot F^{p^e}=\varpi^{op^e}\big(Z^{p^e}+p\,(M_{s'}\cdot\rho)\big),
\]
as $M_{s'}\cdot(M_{s'}^{-1}\cdot Z^{p^e})=Z^{p^e}$. This gives (2). If $(M_s\cdot F)(J)\in\varpi^{o+1}\mathcal R[z]$,
then $\varpi$ divides $Z(J)$, so $Z(J)^{p^e}\in\varpi^{p^e}\mathcal R[z]$; as $p=\varpi^2$ and $p^e\geq1$, with
$p^e\geq2$ if $e\geq1$, this gives (3).
% Lean: step1 with L, L' the exterior powers abelian2 of M_s and its inverse, L2, L2' those at s', and hF
%   from phiS_cong (FrobCong = frobCong hf, through the constant coefficients of Witt vectors); sum_pow_dvd.
\end{proof}

\begin{lemma}\label{pd-symbols-path-lift-nominal}
Let $N,r,K\geq1$ be integers with $N+r\leq K$, $n=r+K$, and a lettering of $\{0,\dots,N-1\}$; we use
the notation of Definition~\ref{pd-symbols-path-generators}. Let $\mu$ be a path of rank $n$ for some $B$
and $c\in\mathbf Z$. For a letter $j$ let $P_j$ be the set of the $m>r$ if $j$ is a vector letter and of
the $m\leq r$ if it is a covector letter. For an integer $h$, a set $J\subset\{1,\dots,n\}$ is
\emph{nominal} for $h$ if, when $c\geq1$, $|J|=h$, $J\supset\{1,\dots,r\}$ if $h\geq r$ and
$J\subset\{1,\dots,r\}$ if $h\leq r$; and, when $c\leq0$, $J\subset\{r+1,\dots,n\}$. Let $j$ be a letter
used for $c$, $J\subset\{1,\dots,n\}$ and $m\in\{1,\dots,n\}$, and put $J_m=J\smallsetminus\{m\}$, with
$m\in J$, if the step at $j$ for $c$ is a wedge, and $J_m=J\cup\{m\}$, with $m\notin J$, if it is a
contraction. Then:
\begin{enumerate}
\item $\emptyset$ is nominal for $h_0(c)$;
\item if ($m\in P_j$ if and only if $j$ is costly for $c$) and $J_m$ is nominal for $h_j(c)$, then $J$ is
  nominal for $h_{j+1}(c)$;
\item if $j$ is costly for $c$ and $m\notin P_j$, then $J_m$ is not nominal for $h_j(c)$.
\end{enumerate}
\end{lemma}
\leanref{Formalism/SharpRestriction/GrowthBounds/FrobeniusPaths/NmHContrPos.lean}{Formalism.MainNotions.LemmaBAlt.NmH}
\leanref{Formalism/SharpRestriction/GrowthBounds/FrobeniusPaths/NmHContrPos.lean}{Formalism.MainNotions.LemmaBAlt.step_nom}
\leanref{Formalism/SharpRestriction/GrowthBounds/FrobeniusPaths/NmHContrPos.lean}{Formalism.MainNotions.LemmaBAlt.step_van}
\leanref{Formalism/SharpRestriction/GrowthBounds/FrobeniusPaths/ColOKOfPath.lean}{Formalism.MainNotions.LemmaBAlt.ColOK}
\leanref{Formalism/SharpRestriction/GrowthBounds/FrobeniusPaths/ColOKOfPath.lean}{Formalism.MainNotions.LemmaBAlt.colOK_of_path}

\begin{proof}
Write $h_i=h_i(c)$ and let $\mathrm{pl}=\{1,\dots,r\}$. As in the proof of
Lemma~\ref{pd-symbols-path-clauses}(3), at a used letter $j$ either $j$ is a vector letter and
$h_{j+1}=h_j+1$, or $j$ is a covector letter and $h_j=h_{j+1}+1$; the step is a wedge for a vector
letter with $c\geq1$ and for a covector letter with $c\leq0$. If $c\geq1$, then $h_0=0$, as $\mu_0=0$;
if $c\leq0$, then $h_0=n$ and $h_i\geq n-i\geq r+K-N>r$ for all $i$. (1) follows. For (2) and (3) there
are four cases.
\begin{itemize}
\item $c\geq1$, vector letter, wedge, $m\in J$. Here $j$ is costly if and only if $h_j\geq r$, and $P_j$
  consists of the $m>r$. (2): $|J|=h_j+1$; if $h_j\geq r$, then $\mathrm{pl}\subset J_m\subset J$; if
  $h_j<r$, then $m\leq r$ and $J_m\subset\mathrm{pl}$, so $J\subset\mathrm{pl}$, and if moreover
  $h_j+1=r$, then $J=\mathrm{pl}$. (3): $m\leq r$ and $h_j\geq r$, so a nominal $J_m$ would contain $m$.
\item $c\geq1$, covector letter, contraction, $m\notin J$. Here $j$ is costly if and only if $h_j\leq r$,
  and $P_j$ consists of the $m\leq r$. (2): $|J|=h_j-1$; if $h_j\leq r$, then
  $J\subset J_m\subset\mathrm{pl}$; if $h_j>r$, then $m>r$ and $\mathrm{pl}\subset J_m$, so
  $\mathrm{pl}\subset J$, and if moreover $h_j-1=r$, then $J=\mathrm{pl}$. (3): $m>r$ and $h_j\leq r$, so a
  nominal $J_m$ would lie in $\mathrm{pl}$.
\item $c\leq0$, covector letter, wedge, $m\in J$. As $h_j>r$, $j$ is not costly, so $m\notin P_j$, that is,
  $m>r$, and $J=J_m\cup\{m\}\subset\{r+1,\dots,n\}$ if $J_m$ is.
\item $c\leq0$, vector letter, contraction, $m\notin J$. Here $J\subset J_m$, which gives (2). As $h_j>r$,
  $j$ is costly; if $m\notin P_j$, then $m\leq r$, so $J_m\ni m$ is not nominal, which is (3).
\end{itemize}
% Lean: NmH c h J; (2) is step_nom (NmH_wedge_pos, NmH_contr_pos, NmH_wedge_neg, NmH_contr_neg), (3) is
%   step_van; the facts on heights are ColOK (rise, fall, start, dual), supplied by colOK_of_path.
\end{proof}

\begin{lemma}\label{pd-symbols-path-lift-order}
Let $p\geq5$ be a prime, $(W,Q,\sigma,\iota)$ a Hida setting of degree $n'$ at $p$ with rings
$\mathcal R\subset\mathcal K$ and $\varpi_0$, $N,r,K\geq1$ integers with $N+r\leq K$, $n=r+K$, and let the
letters of a lettering of $\{0,\dots,N-1\}$ have positions $\pi_0\leq\dots\leq\pi_{N-1}$; we use the
notation of Definitions~\ref{pd-symbols-path-generators} and~\ref{pd-symbols-path-canonical} and of
Lemma~\ref{pd-symbols-path-lift-nominal}, and write $A_s=\iota\sigma^s(A)$ for a matrix $A$ over $W$. Let
$\mu$ be a path of rank $n$ for some $B$ and $c\in\mathbf Z$ such that $\pi_i<\pi_{i'}$ whenever $i<i'$,
$i$ is used for $c$ and $i'$ is costly for $c$.
\begin{enumerate}
\item For rows $x_0,\dots,x_{N-1}$ over a commutative ring, $0\leq j\leq N$ and $s\in\mathbf Z$, let
  $Z_j(x;s)$ be formed as $\mathrm{Col}^g_{\mu,c,\zeta}$ in
  Definition~\ref{pd-symbols-path-generators}(6), with the rows $x_i$ in place of the $\ell_i$, the used
  letters $i<j$ in place of all used letters, and $s$ in place of $\zeta$. Put
  $o_j(s)=\sum_ip^{s-\pi_i}$, over the letters $i<j$ costly for $c$.
\item Let $y_i$ be the row with $y_{i,b}=\varpi_0z_{i,b}$ for $b\in P_i$ and $y_{i,b}=z_{i,b}$ otherwise.
  For $g,g'\in GL_n(W)$ let $x_i=y_i(g'^{-1}g)_{\pi_i}$ for a vector letter $i$ and
  $x_i=y_i\,{}^t(g^{-1}g')_{\pi_i}$ for a covector letter, and $h=g^{-1}g'$ if $c\geq1$ and
  $h={}^t(g'^{-1}g)$ if $c\leq0$.
\end{enumerate}
Then for $0\leq j\leq N$, every $s\in\mathbf Z$ with $s\geq\pi_i$ for all used letters $i<j$, and every
$J$: $(h_s\cdot Z_j(x;s))(J)\in\varpi_0^{o_j(s)}\mathcal R[z]$, and
$(h_s\cdot Z_j(x;s))(J)\in\varpi_0^{o_j(s)+1}\mathcal R[z]$ if $J$ is not nominal for $h_j(c)$.
\end{lemma}
\leanref{Formalism/SharpRestriction/GrowthBounds/FrobeniusPaths/ColumnKernelComponentsIntegral.lean}{Formalism.MainNotions.LemmaBAlt.KKs}
\leanref{Formalism/SharpRestriction/GrowthBounds/FrobeniusPaths/ColumnKernelComponentsIntegral.lean}{Formalism.MainNotions.LemmaBAlt.ZZ}
\leanref{Formalism/SharpRestriction/GrowthBounds/FrobeniusPaths/ColStepFactor.lean}{Formalism.MainNotions.LemmaBAlt.oS}
\leanref{Formalism/SharpRestriction/GrowthBounds/FrobeniusPaths/NmHContrPos.lean}{Formalism.MainNotions.LemmaBAlt.TS}
\leanref{Formalism/SharpRestriction/GrowthBounds/FrobeniusPaths/LiftedColumnOrder.lean}{Formalism.MainNotions.LemmaBAlt.mkTS}
\leanref{Formalism/SharpRestriction/GrowthBounds/FrobeniusPaths/NmHContrPos.lean}{Formalism.MainNotions.LemmaBAlt.Ord_main}

\begin{proof}
Write $\varpi=\varpi_0$, $h_i=h_i(c)$, and $V\equiv0\ (\varpi^a)$ for $V\in\varpi^a\mathcal R[z]$. From the
definition: $Z_{j+1}(x;s)=Z_j(x;s)$ if $j$ is not used; if $j$ is used,
$Z_{j+1}(x;s)$ is the step at $j$ for $c$ applied to $(Z_j(x;\pi_j),x_j)$, raised to the power
$p^{s-\pi_j}$; and $Z_j(x;s')=Z_j(x;s)^{p^{s'-s}}$ for $s\leq s'$ when $s\geq\pi_i$ for all used $i<j$,
while then $o_j(s')=p^{s'-s}o_j(s)$.

\emph{The twist.} Let $j$ be used and $s_0=\pi_j$. As $\sigma$ is a ring automorphism,
$(AA')_{s}=A_sA'_s$, $(A^{-1})_s=A_s^{-1}$ and $({}^tA)_s={}^tA_s$. If $c\geq1$, $h_{s_0}=(g^{-1}g')_{s_0}$:
at a vector letter the step is a wedge and $x_jh_{s_0}=y_j$, at a covector letter a contraction and
$x_j\,{}^th_{s_0}^{-1}=y_j$. If $c\leq0$, $h_{s_0}={}^t(g'^{-1}g)_{s_0}$: at a vector letter the step is a
contraction and $x_j\,{}^th_{s_0}^{-1}=y_j$, at a covector letter a wedge and $x_jh_{s_0}=y_j$. By
Lemma~\ref{pd-symbols-path-exterior}, in all cases $h_{s_0}\cdot\mathrm{step}(F,x_j)=\mathrm{step}(h_{s_0}\cdot
F,y_j)$, where $\mathrm{step}$ is the step at $j$ for $c$.

\emph{Induction on $j$.} For $j=0$, $Z_0(x;s)=F_0=h_s\cdot F_0$, $o_0=0$, and $F_0(J)=0$ unless
$J=\emptyset$, which is nominal for $h_0$ by Lemma~\ref{pd-symbols-path-lift-nominal}(1). Let the claim
hold for $j$. If $j$ is not used, nothing changes, as an unused letter is not costly and $h_{j+1}=h_j$.
Let $j$ be used, $s_0=\pi_j\leq s$, $U=h_{s_0}\cdot Z_j(x;s_0)$ and
$V=\mathrm{step}(Z_j(x;s_0),x_j)$, so that $h_{s_0}\cdot V=\mathrm{step}(U,y_j)$, whose value at $J$ is a sum of
terms $\pm y_{j,m}U(J_m)$, with $J_m$ as in Lemma~\ref{pd-symbols-path-lift-nominal}. We show (A):
$(h_{s_0}\cdot V)(J)\equiv0\ (\varpi^{o_{j+1}(s_0)})$, and $\equiv0\ (\varpi^{o_{j+1}(s_0)+1})$ if $J$ is not
nominal for $h_{j+1}$.

If $j$ is costly, then $o_{j+1}(s_0)=o_j(s_0)+1$, and every used $i<j$ has $\pi_i\leq s_0-1$. The claim
for $j$ at $s_0-1$ and Lemma~\ref{pd-symbols-path-lift-twist}(2),(3) with $e=1$ give $U(J')\equiv0\
(\varpi^{o_j(s_0)})$ for all $J'$, and $U(J')\equiv0\ (\varpi^{o_j(s_0)+2})$ if $J'$ is not nominal for $h_j$.
A term with $m\in P_j$ has $y_{j,m}=\varpi z_{j,m}$, so it is $\equiv0\ (\varpi^{o_j(s_0)+1})$; if $J$ is not
nominal for $h_{j+1}$, then $J_m$ is not nominal for $h_j$ by Lemma~\ref{pd-symbols-path-lift-nominal}(2),
and the term is $\equiv0\ (\varpi^{o_j(s_0)+2})$. For $m\notin P_j$, $J_m$ is not nominal for $h_j$ by
Lemma~\ref{pd-symbols-path-lift-nominal}(3), so the term is $\equiv0\ (\varpi^{o_j(s_0)+2})$.

If $j$ is not costly, then $o_{j+1}(s_0)=o_j(s_0)$, and the used $i<j$ have $\pi_i\leq s_0$; the claim for
$j$ at $s_0$ gives $U\equiv0\ (\varpi^{o_j(s_0)})$, so every term is. Let $J$ not be nominal for $h_{j+1}$.
A term with $m\in P_j$ has the factor $\varpi$. For $m\notin P_j$, $J_m$ is not nominal for $h_j$ by
Lemma~\ref{pd-symbols-path-lift-nominal}(2), so $U(J_m)\equiv0\ (\varpi^{o_j(s_0)+1})$.

This proves (A). Now $Z_{j+1}(x;s)=V^{p^{s-s_0}}$ and $o_{j+1}(s)=p^{s-s_0}o_{j+1}(s_0)$, as all costly
$i\leq j$ have $\pi_i\leq s_0$. Lemma~\ref{pd-symbols-path-lift-twist}(2),(3), with $M$ the matrix $h$ over $W$,
$s_0\leq s$ and (A), gives the claim for $j+1$ at $s$.
% Lean: Ord_main for ts = mkTS hF S hsq c mu pos zetaL yO M M' with T s Kv = vecMul Kv (abelian2 (M.map phiS s)),
%   M = g^-1 g' for c >= 1 and the transpose of g'^-1 g for c <= 0 (the four step identities by Lam_wedge and
%   Lam_contr); ZZ_unused, ZZ_used, ZZ_pow, oS_succ, oS_pow; colStep_dvd, yO_cond.
\end{proof}

\begin{lemma}\label{pd-symbols-path-lift-residue}
In the setting of Lemma~\ref{pd-symbols-path-lift-order}, with $y_i$ as there, let $\kappa$ be the residue
field of $\mathcal R$ and bars denote reduction modulo $\varpi_0$. For a letter $i$ let $n_i$ be the row
with $n_{i,b}=z_{i,b}$ if ($b\in P_i$ if and only if $i$ is costly for $c$) and $n_{i,b}=0$ otherwise, and
let $\mathcal N_j$ be formed as $\mathrm{Col}_{\mu,c}$ in Definition~\ref{pd-symbols-path-generators}(3), with
the rows $n_i$ and the first $j$ letters. For $s\in\mathbf Z$ let $\pi_s\colon\mathbf Z[z]\to\kappa[z]$ be
the ring map with $z_{i,b}\mapsto z_{i,b}^{p^{s-\pi_i}}$ for the letters with $\pi_i\leq s$. Then:
\begin{enumerate}
\item $\mathcal N_j(J)\neq0$ only if $J$ is nominal for $h_j(c)$; and if $j$ is costly for $c$, $m\notin P_j$
  and $J_m$ is as in Lemma~\ref{pd-symbols-path-lift-nominal}, then $\mathcal N_j(J_m)=0$;
\item for $0\leq j\leq N$, every $s$ with $s\geq\pi_i$ for all used letters $i<j$, and every $J$, there is
  $Y\in\mathcal R[z]$ with $Z_j(y;s)(J)=\varpi_0^{o_j(s)}Y$ and $\bar Y=\pi_s(\mathcal N_j(J))$;
\item $\mathcal N_N=\mathrm{Col}^{\mathrm{pt}}_{\mu,c}$.
\end{enumerate}
\end{lemma}
\leanref{Formalism/SharpRestriction/GrowthBounds/FrobeniusPaths/ColStepFactor.lean}{Formalism.MainNotions.LemmaBAlt.nomL}
\leanref{Formalism/SharpRestriction/GrowthBounds/FrobeniusPaths/ColStepFactor.lean}{Formalism.MainNotions.LemmaBAlt.SgW}
\leanref{Formalism/SharpRestriction/GrowthBounds/FrobeniusPaths/ColOKOfPath.lean}{Formalism.MainNotions.LemmaBAlt.SgW_supp}
\leanref{Formalism/SharpRestriction/GrowthBounds/FrobeniusPaths/ColOKOfPath.lean}{Formalism.MainNotions.LemmaBAlt.SgVanish_of}
\leanref{Formalism/SharpRestriction/GrowthBounds/FrobeniusPaths/ColStepFactor.lean}{Formalism.MainNotions.LemmaBAlt.Inv_main}
\leanref{Formalism/SharpRestriction/GrowthBounds/FrobeniusPaths/ColStepFactor.lean}{Formalism.MainNotions.LemmaBAlt.piS_frob}
\leanref{Formalism/SharpRestriction/GrowthBounds/FrobeniusPaths/WedgeSplitTan.lean}{Formalism.MainNotions.LemmaBAlt.SgW_eq}
\leanref{Formalism/SharpRestriction/GrowthBounds/FrobeniusPaths/WedgeSplitTan.lean}{Formalism.MainNotions.LemmaBAlt.plTanCol_eq_SgW}

\begin{proof}
Write $\varpi=\varpi_0$ and $h_i=h_i(c)$. (1) By induction on $j$: $\mathcal N_0=F_0$ is nonzero only at
$\emptyset$, which is nominal for $h_0$ (Lemma~\ref{pd-symbols-path-lift-nominal}(1)). At a used letter $j$,
$\mathcal N_{j+1}(J)$ is a sum of terms $\pm n_{j,m}\mathcal N_j(J_m)$, and $n_{j,m}\neq0$ only if ($m\in P_j$ if
and only if $j$ is costly); if $\mathcal N_{j+1}(J)\neq0$, some such $J_m$ is nominal for $h_j$, and $J$ is
nominal for $h_{j+1}$ by Lemma~\ref{pd-symbols-path-lift-nominal}(2). The second assertion follows from
the first and Lemma~\ref{pd-symbols-path-lift-nominal}(3).

(2) As $p=\varpi^2$, $\kappa$ has characteristic $p$, so $P\mapsto P^p$ is a ring endomorphism of
$\kappa[z]$ fixing the image of $\mathbf Z$; hence $\pi_s(P)^{p^{s'-s}}=\pi_{s'}(P)$ for $s\leq s'$ if $P$
involves only rows $z_i$ with $\pi_i\leq s$. The polynomial $\mathcal N_j(J)$ involves only the rows of the
used letters $i<j$. We argue by induction on $j$, with the recursion for $Z_j$ of the proof of
Lemma~\ref{pd-symbols-path-lift-order}. For $j=0$ take $Y=F_0(J)$. An unused letter changes nothing. Let $j$
be used, $s_0=\pi_j\leq s$, and $\mathrm{step}$ the step at $j$ for $c$.

If $j$ is costly, then every used $i<j$ has $\pi_i\leq s_0-1$, and the induction hypothesis at $s_0-1$ gives
$Z_j(y;s_0-1)(J')=\varpi^{o_j(s_0-1)}Y_1(J')$ with $\bar Y_1(J')=\pi_{s_0-1}(\mathcal N_j(J'))$. So
$Z_j(y;s_0)=\varpi^{o_j(s_0)}Y_1^p$ and $\overline{Y_1^p}(J')=\pi_{s_0}(\mathcal N_j(J'))$. The value of
$\mathrm{step}(\varpi^{o_j(s_0)}Y_1^p,y_j)$ at $J$ is $\varpi^{o_j(s_0)}$ times a sum of terms $\pm
y_{j,m}Y_1(J_m)^p$. For $m\in P_j$ the term is $\varpi z_{j,m}Y_1(J_m)^p$, and $z_{j,m}=n_{j,m}=\pi_{s_0}(n_{j,m})$.
For $m\notin P_j$, $\mathcal N_j(J_m)=0$ by (1), so $Y_1(J_m)=\varpi T$ with $T\in\mathcal R[z]$, and the term
is $\varpi\cdot z_{j,m}\varpi^{p-1}T^p$, whose quotient by $\varpi$ reduces to $0=\pi_{s_0}(n_{j,m})\pi_{s_0}(\mathcal
N_j(J_m))$. So the value is $\varpi^{o_j(s_0)+1}Y'$ with
$\bar Y'=\mathrm{step}(\pi_{s_0}\circ\mathcal N_j,\pi_{s_0}(n_j))(J)=\pi_{s_0}(\mathcal N_{j+1}(J))$, and
$o_j(s_0)+1=o_{j+1}(s_0)$.

If $j$ is not costly, the induction hypothesis at $s_0$ gives $Z_j(y;s_0)=\varpi^{o_j(s_0)}Y_1$, and the step
gives $\varpi^{o_j(s_0)}\mathrm{step}(Y_1,y_j)$. Modulo $\varpi$, $y_{j,b}$ is $0$ for $b\in P_j$ and $z_{j,b}$
otherwise, which is $\pi_{s_0}(n_{j,b})$; so $\mathrm{step}(Y_1,y_j)(J)$ reduces to
$\pi_{s_0}(\mathcal N_{j+1}(J))$, and $o_{j+1}(s_0)=o_j(s_0)$.

In both cases $Z_{j+1}(y;s)(J)=(\varpi^{o_{j+1}(s_0)}Y')^{p^{s-s_0}}=\varpi^{o_{j+1}(s)}Y'^{p^{s-s_0}}$, and
$\overline{Y'}^{\,p^{s-s_0}}=\pi_s(\mathcal N_{j+1}(J))$.

(3) In the notation of the proof of Lemma~\ref{pd-symbols-path-clauses}(3), we show
$\mathcal N_j(J)=P_j(J')T_j(J'')$ by induction on $j$. In each of the cases (i) to (vi) there, the variables
$z_{j,m}$ that are entries of $n_j$ are those whose terms that proof keeps: $m\leq r$ in cases (i) and (iv)
and $m>r$ in the others. With $\mathcal N_j=P_jT_j$ in place of the lowest parts, the computation there shows
that these terms sum to $P_{j+1}(J')T_{j+1}(J'')$. At $j=N$ this is
$\mathrm{Col}^{\mathrm{pt}}_{\mu,c}(J)$.
% Lean: (1) SgW_supp and SgVanish_of; (2) Inv_main (colStep_factor, sum_factor, piR_ker, piR_w, piS_frob,
%   ZZ_pow, ZZ_used); (3) plTanCol_eq_SgW = SgW_eq (wedge_split_tan, contr_split_tan, wedge_split_pl,
%   contr_split_pl, PLc_succ, TNc_succ, nomL_X, nomL_0).
\end{proof}

\begin{lemma}\label{pd-symbols-path-clauses}
Let $p$ be a prime, $\mathcal R$ a complete discrete valuation ring in which $p$ is neither $0$ nor a
unit, with residue field $\bar{\mathcal R}$, and $N\geq0$ an integer with a lettering of
$\{0,\dots,N-1\}$; we use the notation of Definition~\ref{pd-symbols-path-generators}.
\begin{enumerate}
\item Let $n\geq1$, $B\colon\{0,\dots,N-1\}\to\mathbf Z_{\geq0}$ and $\theta,M\colon
  \{0,\dots,N-1\}\to\mathbf Z$, and let $R$ be $\mathcal R$ or $\bar{\mathcal R}$. Then the set of the $G\in
  R[z]$ of row degrees $B$ with $\theta(i)\leq\operatorname{val}_{\{0,\dots,i\}}(G)$ and
  $\lambda_{\{0,\dots,i\}}(G)\leq M(i)$ for every letter $i$ is the $R$-span of the path generators
  $\Psi(\mu,I)$ in $R[z]$ with $\mu$ a path of rank $n$ for $B$ such that $\theta(i)\leq\mu_{i+1}(n)$ and
  $\mu_{i+1}(1)\leq M(i)$ for every letter $i$, and $I$ a standard filling of $\mu$.
\item Let $r,K\geq1$ with $N+r\leq K$, let $B\colon\{0,\dots,N-1\}\to\mathbf Z_{\geq0}$, and let
  $\nu$ and $\kappa$ be sequences of length $N+1$ in $\mathbf Z^r$ and $\mathbf Z^K$. Then
  $(\nu,\kappa)$ is a path of rank $r+K$ for $B$ if and only if there are $B',B''$ with $B'+B''=B$ such
  that $\nu$ is a path of rank $r$ for $B'$, $\kappa$ is a path of rank $K$ for $B''$, and for every
  letter $i$, $\kappa_{i+1}(1)\leq\nu_i(r)$ if $i$ is a vector letter and $\kappa_i(1)\leq\nu_{i+1}(r)$
  if it is a covector letter.
\item Let $r,K\geq1$ with $N+r\leq K$, $\omega\colon\{0,\dots,N-1\}\to\mathbf Z_{>0}$ strictly
  decreasing, $\mu$ a path of rank $r+K$ for some $B$, $c\in\mathbf Z$ and $J\subset\{1,\dots,r+K\}$.
  Then every monomial of $\mathrm{Col}_{\mu,c}(J)\in\mathcal R[z]$ has $\omega$-degree at least $o_\omega(\mu,c)$,
  and the sum of its terms of $\omega$-degree $o_\omega(\mu,c)$ is $\mathrm{Col}^{\mathrm{pt}}_{\mu,c}(J)$.
\item Let $r$, $K$ and $\omega$ be as in (3), let $(\nu,\kappa)$ be a path of rank $r+K$ for some
  $B$, and let $I'$ and $I''$ be standard fillings of $\nu$ and $\kappa$. Put $\mu=(\nu,\kappa)$ and
  $o=\sum_{|c|\leq T_\mu}o_\omega(\mu,c)$. Then there is a function $I$ from $\mathbf Z$ to the
  subsets of $\{1,\dots,r+K\}$ such that every monomial of $\Psi(\mu,I)\in\mathcal R[z]$ has
  $\omega$-degree at least $o$, and the sum of its terms of $\omega$-degree $o$ is
  $\iota_{\mathrm{pl}}(\Psi(\nu,I'))\,\iota_{\mathrm{tn}}(\Psi(\kappa,I''))$.
\item Suppose that $p\geq5$ and that $\mathcal R$, its fraction field $\mathcal K$ and an element
  $\varpi_0$ are the rings and the element of a Hida setting $(W,Q,\sigma,\iota)$ of degree $n'$ at
  $p$, with $\varpi_0$ a uniformizer of $\mathcal R$. Let
  $N,r,K\geq1$ with $N+r\leq K$, $n=r+K$, and let the letters have positions
  $\pi_0\leq\pi_1\leq\dots\leq\pi_{N-1}\leq\zeta$. Let $\mu$ be a path of rank $n$ for some $B$,
  $g\in GL_n(W)$ and $c\in\mathbf Z$ such that $\pi_i<\pi_{i'}$ whenever $i<i'$, $i$ is used for $c$ and
  $i'$ is costly for $c$. Put $\omega(i)=p^{\zeta-\pi_i}$ and $o=o_\omega(\mu,c)$, and let
  $J\subset\{1,\dots,n\}$. Then:
  \begin{enumerate}
  \item $\mathrm{Col}^g_{\mu,c,\zeta}(J)$ has its coefficients in $\mathcal R$;
  \item every monomial of $\mathrm{Col}^g_{\mu,c,\zeta}(J)$ has degree $p^{\zeta-\pi_i}$ in $z_i$ for every letter
    $i$ used for $c$;
  \item $\Sigma^{\mathrm L}_{g'}\big(\mathrm{Col}^g_{\mu,c,\zeta}(J)\big)\in\varpi_0^o\mathcal R[z]$ for every
    $g'\in GL_n(W)$;
  \item $\Sigma^{\mathrm L}_g\big(\mathrm{Col}^g_{\mu,c,\zeta}(J)\big)-\varpi_0^o\,\Pi_\zeta\big(\mathrm{Col}^{\mathrm{pt}}_{\mu,c}(J)\big)
    \in\varpi_0^{o+1}\mathcal R[z]$;
  \item $o\leq2p^{\zeta-\pi_0}$.
  \end{enumerate}
\item Let $F$ be a field, $r\geq1$, $S\subset\{0,\dots,N-1\}$, and $A\colon\{1,\dots,r\}\to S$ an
  injective map whose values are vector letters. Let $F[z]$ be the polynomial ring in the $z_{i,b}$,
  $0\leq i<N$, $1\leq b\leq r$, with $\operatorname{val}$ formed with $n=r$. Put
  $\Delta_A=\det(z_{A(l),b})_{1\leq l,b\leq r}$ and, for $d\colon\{0,\dots,N-1\}\to\mathbf Z_{\geq0}$, let
  $d_A(S)$ be the sum of the $d(i)$ over the vector letters $i\in S$ that are not values of $A$. Then:
  \begin{enumerate}
  \item if $G\in F[z]$ has row degrees $d$ and $\theta$ is an integer with
    $d_A(S)<\theta\leq\operatorname{val}_S(G)$, then $\Delta_A^{\,\theta-d_A(S)}$ divides $G$;
  \item if $G=\Delta_A^mH$ with $G,H\in F[z]$ and $m\geq0$, then
    $\operatorname{val}_{S'}(G)\leq\operatorname{val}_{S'}(H)+m$ for every $S'\subset\{0,\dots,N-1\}$.
  \end{enumerate}
\end{enumerate}
\end{lemma}
\leanref{Formalism/SharpRestriction/GrowthBounds/StatementOfDegreeAndBound.lean}{Formalism.MainNotions.CovariantDegrees.A3at}
\leanref{Formalism/SharpRestriction/GrowthBounds/StatementOfDegreeAndBound.lean}{Formalism.MainNotions.CovariantDegrees.DefiniteRankOne}
\leanref{Formalism/SharpRestriction/GrowthBounds/StatementOfDegreeAndBound.lean}{Formalism.MainNotions.CovariantDegrees.CondC}
\leanref{Formalism/SharpRestriction/GrowthBounds/StatementOfDegreeAndBound.lean}{Formalism.MainNotions.CovariantDegrees.RotatedSumBound}
\leanref{Formalism/SharpRestriction/GrowthBounds/StatementOfDegreeAndBound.lean}{Formalism.MainNotions.CovariantDegrees.PathWeightComponents}
\leanref{Formalism/SharpRestriction/GrowthBounds/StatementOfDegreeAndBound.lean}{Formalism.MainNotions.CovariantDegrees.ConclusionB3}
\leanref{Formalism/SharpRestriction/GrowthBounds/StatementOfDegreeAndBound.lean}{Formalism.MainNotions.CovariantDegrees.CanonicalTranslateJetIdentity}
\leanref{Formalism/SharpRestriction/GrowthBounds/StatementOfDegreeAndBound.lean}{Formalism.MainNotions.CovariantDegrees.Conclusion2}
\leanref{Formalism/SharpRestriction/GrowthBounds/TableLineDegreeFive.lean}{Formalism.MainNotions.tableLineDegreeFive.adm15}
\leanref{Formalism/SharpRestriction/GrowthBounds/TableLineDegreeFive.lean}{Formalism.MainNotions.tableLineDegreeFive.adm16}
\leanref{Formalism/SharpRestriction/GrowthBounds/TableLineDegreeFive.lean}{Formalism.MainNotions.tableLineDegreeFive.adm17}
\leanref{Formalism/SharpRestriction/GrowthBounds/TableLineDegreeFive.lean}{Formalism.MainNotions.tableLineDegreeFive.adm18}
\leanref{Formalism/SharpRestriction/GrowthBounds/TableLineDegreeFive.lean}{Formalism.MainNotions.tableLineDegreeFive.adm20}
\leanref{Formalism/SharpRestriction/GrowthBounds/BinomialPolynomialIdentities/SplittingThirdPartHolds.lean}{Formalism.MainNotions.G.ColumnFamilies.splitting_package_holds}
\leanref{Formalism/SharpRestriction/GrowthBounds/CovariantDegrees.lean}{Formalism.MainNotions.CovariantDegrees.ConclusionE}
\leanref{Formalism/SharpRestriction/GrowthBounds/TableLineDegreeFive.lean}{Formalism.MainNotions.tableLineDegreeFive.adm21}
\leanref{Formalism/SharpRestriction/GrowthBounds/WittFrameLocaliserOrders.lean}{Formalism.MainNotions.G.WittFrameLocaliserOrders.FieldMultidegreeDivisibility}
\leanref{Formalism/SharpRestriction/GrowthBounds/GLWeightTheory.lean}{Formalism.MainNotions.G.GLWeightTheory.support_coprimality_lemma_holds}
\leanref{Formalism/SharpRestriction/GrowthBounds/RootWeights/FrameChangePreservesMultidegree.lean}{Formalism.MainNotions.G.GLWeightTheory.height_bound_of_weights}
\leanref{Formalism/SharpRestriction/GrowthBounds/RootWeights/ValuationBoundDetPowerFactor.lean}{Formalism.MainNotions.G.GLWeightTheory.valuation_bound_det_power_factor}

\begin{proof}
We use the notation of Definitions~\ref{pd-symbols-path-generators} and~\ref{pd-symbols-path-weights}.
For $\alpha,\beta\in\mathbf Z^m$ and $b\geq0$, call $\beta$ a \emph{strip of size $b$ over $\alpha$} if
$\alpha(a)\leq\beta(a)$ for all $a$, $\beta(a+1)\leq\alpha(a)$ for $a<m$, and
$\sum_a(\beta(a)-\alpha(a))=b$. Thus $\mu$ is a path for $B$ if $\mu_0=0$, each $\mu_i$ is dominant, and
$\mu_{i+1}$ is a strip of size $B(i)$ over $\mu_i$ for a vector letter $i$, and $\mu_i$ one over
$\mu_{i+1}$ for a covector letter $i$.

(1) For $P\in R[z]$, a set $S$ of letters and $a\in\{1,n\}$, let $\mathcal E_{S,a}(P)$ be the set of
the exponents of $\eta$ that occur in $P_{S,a}$; so $\operatorname{val}_S(P)$ is the least element of
$\mathcal E_{S,n}(P)$ and $\lambda_S(P)$ the largest element of $\mathcal E_{S,1}(P)$. The map
$P\mapsto P_{S,a}$ is a ring homomorphism, it sends the constants to constants, and it commutes with
the maps induced by ring homomorphisms $R\to R'$. Hence $\mathcal E_{S,a}$ of the image of $P$ in
$R'[z]$ is contained in $\mathcal E_{S,a}(P)$, $\mathcal E_{S,a}(P+Q)\subset\mathcal E_{S,a}(P)\cup
\mathcal E_{S,a}(Q)$, and $\mathcal E_{S,a}(cP)\subset\mathcal E_{S,a}(P)$ for $c\in R$.

Let $\mathcal P$ be the set of pairs $x=(\mu,I)$ of a path $\mu$ of rank $n$ for $B$ and a standard
filling $I$ of $\mu$, write $\Psi_x=\Psi(\mu,I)$, and let $\mathcal P_0\subset\mathcal P$ consist of the
pairs with $\theta(i)\leq\mu_{i+1}(n)$ and $\mu_{i+1}(1)\leq M(i)$ for every letter $i$. For a ring $R$
let $A_R$ be the set of the $G$ of the statement; we show that $A_R$ is the $R$-span of the
$\Psi_x$ with $x\in\mathcal P_0$.

Let $x=(\mu,I)\in\mathcal P_0$. By Lemma~\ref{pd-symbols-path-basis}, $\Psi_x$ has row degrees $B$.
The column functions are built from the variables by sums of products with signs, so $\Psi_x\in R[z]$
is the image of $\Psi_x\in\mathbf Z[z]$. By Lemma~\ref{pd-symbols-path-word-bounds} and the first
property of $\mathcal E$, every element of $\mathcal E_{\{0,\dots,i\},n}(\Psi_x)$ is at least
$\mu_{i+1}(n)\geq\theta(i)$ and every element of $\mathcal E_{\{0,\dots,i\},1}(\Psi_x)$ is at most
$\mu_{i+1}(1)\leq M(i)$, for every letter $i$. So $\Psi_x\in A_R$. By the other two properties of
$\mathcal E$, $A_R$ is closed under sums and under multiplication by elements of $R$, and it contains
$0$; so it contains the span.

Conversely, let $G\in A_R$. Suppose first that $R$ is a field. By Lemma~\ref{pd-symbols-path-basis},
$G=\sum_{x\in\mathcal P}l_x\Psi_x$ for unique $l_x\in R$, almost all $0$. Let $i$ be a letter, put
$S=\{0,\dots,i\}$, and let $\pi$ be the set of the dominant $\nu\in\mathbf Z^n$ with
$\theta(i)\leq\nu(n)$ and $\nu(1)\leq M(i)$. It is a lower set: if $\nu$ is dominant and
$\nu\preceq\mu$ with $\mu\in\pi$, then $\nu(1)\leq\mu(1)\leq M(i)$, and $\nu(n)\geq\mu(n)\geq\theta(i)$,
by the equality of the sums of all entries and, for $n\geq2$, the inequality of the sums of the first
$n-1$. By Lemma~\ref{pd-symbols-path-valuation}, every weight $\mu$ of $U_S(G)$ satisfies
$\theta(i)\leq\mu(n)$ and $\mu(1)\leq M(i)$, so every dominant one lies in $\pi$. By
Lemma~\ref{pd-symbols-path-lower-set}, applied with $i+1$, $G$ is a linear combination of the
$\Psi_{(\mu,I)}$ with $\mu_{i+1}\in\pi$. By the uniqueness of the $l_x$, every $x=(\mu,I)$ with
$l_x\neq0$ has $\mu_{i+1}\in\pi$. As this holds for every letter $i$, $l_x\neq0$ implies
$x\in\mathcal P_0$, and $G$ lies in the span. This proves (1) for the residue field of $\mathcal R$.

Now let $R=\mathcal R$, with fraction field $\mathcal K$. By Lemma~\ref{pd-symbols-path-basis} over
$\mathcal R$, $G=\sum_xl_x\Psi_x$ with $l_x\in\mathcal R$. By the first property of $\mathcal E$, the
image $\sum_xl_x\Psi_x$ of $G$ in $\mathcal K[z]$ lies in $A_{\mathcal K}$. By the case of a field and
the linear independence of the $\Psi_x$ over $\mathcal K$, $l_x\neq0$ implies $x\in\mathcal P_0$, as
$\mathcal R\to\mathcal K$ is injective. So $G$ lies in the $\mathcal R$-span of the $\Psi_x$ with
$x\in\mathcal P_0$.

(2) Let $\alpha',\beta'\in\mathbf Z^r$ and $\alpha'',\beta''\in\mathbf Z^K$. Then $(\beta',\beta'')$ is
a strip of size $b$ over $(\alpha',\alpha'')$ if and only if there are $b',b''\geq0$ with $b'+b''=b$
such that $\beta'$ is a strip of size $b'$ over $\alpha'$, $\beta''$ is one of size $b''$ over
$\alpha''$, and $\beta''(1)\leq\alpha'(r)$. Indeed, the conditions at the entries split between the
two blocks except $\beta(r+1)\leq\alpha(r)$, which is $\beta''(1)\leq\alpha'(r)$, and the sizes are
$b'=\sum_a(\beta'(a)-\alpha'(a))$ and $b''=\sum_a(\beta''(a)-\alpha''(a))$, both nonnegative. Likewise
$(\nu_i,\kappa_i)$ is dominant if and only if $\nu_i$ and $\kappa_i$ are dominant and
$\kappa_i(1)\leq\nu_i(r)$.

Let $(\nu,\kappa)$ be a path for $B$. For a vector letter $i$ apply the criterion to the strip
$(\nu_{i+1},\kappa_{i+1})$ over $(\nu_i,\kappa_i)$, and for a covector letter to the strip
$(\nu_i,\kappa_i)$ over $(\nu_{i+1},\kappa_{i+1})$. This gives $B'(i)+B''(i)=B(i)$, the strip
conditions for $\nu$ with sizes $B'(i)$ and for $\kappa$ with sizes $B''(i)$, and
$\kappa_{i+1}(1)\leq\nu_i(r)$, respectively $\kappa_i(1)\leq\nu_{i+1}(r)$. As $(\nu_i,\kappa_i)$ is
dominant, so are $\nu_i$ and $\kappa_i$, and $\nu_0=0$, $\kappa_0=0$. So $\nu$ is a path for $B'$ and
$\kappa$ one for $B''$.

Conversely, let $B'$, $B''$, $\nu$ and $\kappa$ be as in (2). We show $\kappa_i(1)\leq\nu_i(r)$ for
$0\leq i\leq N$. For $i=0$ both sides are $0$. Let $i=j+1$. If $j$ is a vector letter, then
$\kappa_{j+1}(1)\leq\nu_j(r)\leq\nu_{j+1}(r)$, as $\nu_{j+1}$ is a strip over $\nu_j$. If $j$ is a
covector letter, then $\kappa_{j+1}(1)\leq\kappa_j(1)\leq\nu_{j+1}(r)$, as $\kappa_j$ is a strip over
$\kappa_{j+1}$. So each $(\nu_i,\kappa_i)$ is dominant, $(\nu_0,\kappa_0)=0$, and by the criterion,
with $b'=B'(i)$ and $b''=B''(i)$, $(\nu,\kappa)$ satisfies the strip conditions with sizes $B(i)$.

(3) Fix $\mu$ and $c$, write $h_j=h_j(c)$ for $0\leq j\leq N$ and $n=r+K$. As $\mu_j$ is dominant,
$\mu_j(a)\geq c$ if and only if $a\leq h_j$. If $\beta$ is a strip over $\alpha$, both dominant, then the
numbers $h,h'$ of entries at least $c$ of $\alpha$ and $\beta$ satisfy $h\leq h'\leq h+1$: the first as
$\alpha\leq\beta$, and if $\beta(h+2)\geq c$ then $\alpha(h+1)\geq\beta(h+2)\geq c$, which is false.
Hence $h_{j+1}\in\{h_j,h_j+1\}$ at a vector letter $j$ and $h_{j+1}\in\{h_j-1,h_j\}$ at a covector letter.
The numbers of entries at least $c$ of $\nu_j=\mu^{\mathrm{pl}}_j$ and $\kappa_j=\mu^{\mathrm{tn}}_j$ are
$\min(h_j,r)$ and $(h_j-r)^+$, where $x^+=\max(x,0)$. In the row $z_j$ the variables $z_{j,b}$ with
$b\leq r$ have $\omega$-degree $0$ and those with $b>r$ have $\omega$-degree $\omega(j)$ if $j$ is a
vector letter, and the other way round if $j$ is a covector letter.

For $J\subset\{1,\dots,n\}$ put $J'=J\cap\{1,\dots,r\}$ and $J''=\{b:r+b\in J\}$. Let $F_j$ be the
function after the first $j$ steps of the definition of $\mathrm{Col}_{\mu,c}$, and $P_j$ and $T_j$
those after $j$ steps for $\mathrm{Col}_{\nu,c}$ and $\mathrm{Col}_{\kappa,c}$, formed with the rows
$(z_{i,1},\dots,z_{i,r})$ and $(z_{i,r+1},\dots,z_{i,r+K})$; so $F_N=\mathrm{Col}_{\mu,c}$ and
$P_N(J')T_N(J'')=\mathrm{Col}^{\mathrm{pt}}_{\mu,c}(J)$. Let $o_j$ be the sum of the $\omega(i)$ over
the letters $i<j$ costly for $c$, so $o_N=o_\omega(\mu,c)$. Write $\operatorname{ord}P\geq t$ if every
monomial of $P$ has $\omega$-degree at least $t$, and $[P]_t$ for the sum of the terms of $P$ of
$\omega$-degree $t$. If $y$ is a variable of $\omega$-degree $e$ and $\operatorname{ord}P\geq t$, then
$\operatorname{ord}(\pm yP)\geq e+t$ and $[\pm yP]_{e+t}=\pm y[P]_t$. We show by induction on $j\leq N$:
\begin{enumerate}
\item[(A$_j$)] for every integer $w\geq0$ with $w\leq\omega(i)$ for all $i<j$, and every $J$,
  $\operatorname{ord}F_j(J)\geq o_j+(\min(h_j,r)-|J'|)^+w$ if $c\geq1$, and
  $\operatorname{ord}F_j(J)\geq o_j+|J'|w$ if $c\leq0$;
\item[(D$_j$)] $[F_j(J)]_{o_j}=P_j(J')T_j(J'')$ for every $J$.
\end{enumerate}
Then (3) is (A$_N$) with $w=0$, together with (D$_N$).

Two facts are used. If $c\geq1$, then $h_0=0$, and $F_j(J)\neq0$ only if $|J|=h_j$, as each wedge
step adds an element to $J$ and each contraction step removes one; likewise $T_j(J'')\neq0$ only if
$|J''|=(h_j-r)^+$. If $c\leq0$, then $h_0=n$, and $h_j\geq n-j\geq r$ for $j\leq N$, as $N\leq K$. As
for the signs, $\epsilon_J(r+b)=\epsilon_{J''}(b)$, and $\epsilon_J(m)=\epsilon_{J'}(m)$ for $m\leq r$
when $J''=\emptyset$.

For $j=0$, $F_0(J)=P_0(J')T_0(J'')$ is $1$ if $J=\emptyset$ and $0$ otherwise, and $o_0=0$. The bound
in (A$_0$) is $0$ when $J=\emptyset$ or $c\geq1$, and $F_0(J)=0$ otherwise. Let $j<N$ and suppose
(A$_j$) and (D$_j$). As $\omega$ is strictly decreasing, (A$_j$) holds with $w'=\omega(j)+1$; in
particular $\operatorname{ord}F_j(J)\geq o_j$. Let $w\geq0$ with $w\leq\omega(i)$ for $i\leq j$. If
$h_{j+1}=h_j$, then $F$, $P$, $T$ and $o$ do not change. Otherwise $F_{j+1}(J)$ is a sum of the terms
$\pm z_{j,m}F_j(J\smallsetminus\{m\})$ over $m\in J$ for a wedge step and $\pm z_{j,m}F_j(J\cup\{m\})$
over $m\notin J$ for a contraction step, and we bound each term by (A$_j$) with $w'$. The arithmetic is
$(d+1)w\leq\omega(j)+d\,w'$ for $d\geq0$, and $\omega(j)+dw\leq(d+1)w'$ and $\omega(j)+1\leq(d+1)w'$.
\begin{enumerate}
\item[(i)] $c\geq1$, $h_{j+1}=h_j+1\leq r$. The letter $j$ is a vector letter, the step is a wedge, $j$
  is not costly, $P_{j+1}=P_j\wedge(z_{j,1},\dots,z_{j,r})$ and $T_{j+1}=T_j$. A term with $m\leq r$
  has $\operatorname{ord}\geq o_j+(h_{j+1}-|J'|)^+w'$, and one with $m>r$ has
  $\operatorname{ord}\geq\omega(j)+o_j+(h_j-|J'|)^+w'$; both bounds are at least
  $o_{j+1}+(h_{j+1}-|J'|)^+w$. The terms with $m>r$ have $\operatorname{ord}>o_j$, so
  $[F_{j+1}(J)]_{o_j}$ is the sum of the $\pm z_{j,m}P_j(J'\smallsetminus\{m\})T_j(J'')$, $m\in J'$. As
  $(h_j-r)^+=0$, $T_j(J'')=0$ unless $J''=\emptyset$, and then the signs are those of
  $P_j\wedge(z_{j,1},\dots,z_{j,r})$ at $J'$.
\item[(ii)] $c\geq1$, $h_{j+1}=h_j+1$, $h_j\geq r$. A vector letter and a wedge step; $j$ is costly,
  $o_{j+1}=o_j+\omega(j)$, $P_{j+1}=P_j$ and $T_{j+1}=T_j\wedge(z_{j,r+1},\dots,z_{j,n})$. A term with
  $m\leq r$ has $\operatorname{ord}\geq o_j+(r-|J'|+1)w'$, which is at least
  $o_{j+1}+(r-|J'|)w$ and at least $o_{j+1}+1$. A term with $m>r$ has
  $\operatorname{ord}\geq o_{j+1}+(r-|J'|)w'$, and its part of $\omega$-degree $o_{j+1}$ is
  $\pm z_{j,m}P_j(J')T_j(J''\smallsetminus\{m-r\})$, with the sign $\epsilon_{J''}(m-r)$. The sum is
  $P_j(J')T_{j+1}(J'')$.
\item[(iii)] $c\geq1$, $h_{j+1}=h_j-1$, $h_j>r$. A covector letter and a contraction step; $j$ is not
  costly, $P_{j+1}=P_j$ and $T_{j+1}=T_j\lrcorner(z_{j,r+1},\dots,z_{j,n})$. A term with $m\leq r$ has
  $\operatorname{ord}\geq\omega(j)+o_j+(r-|J'|-1)^+w'$ and one with $m>r$ has
  $\operatorname{ord}\geq o_j+(r-|J'|)w'$; both are at least $o_j+(r-|J'|)w$. The terms with $m\leq r$
  have $\operatorname{ord}>o_j$, and those with $m>r$ contribute $\pm z_{j,m}P_j(J')T_j(J''\cup\{m-r\})$
  with the sign $\epsilon_{J''}(m-r)$, whose sum is $P_j(J')T_{j+1}(J'')$.
\item[(iv)] $c\geq1$, $h_{j+1}=h_j-1$, $h_j\leq r$. A covector letter and a contraction step; $j$ is
  costly, $o_{j+1}=o_j+\omega(j)$, $P_{j+1}=P_j\lrcorner(z_{j,1},\dots,z_{j,r})$ and $T_{j+1}=T_j$. A
  term with $m\leq r$ has $\operatorname{ord}\geq\omega(j)+o_j+(h_j-1-|J'|)^+w'$, at least
  $o_{j+1}+(h_{j+1}-|J'|)^+w$. If a term with $m>r$ is nonzero, then $|J|+1=h_j$, so $|J'|\leq h_j-1$,
  and the term has $\operatorname{ord}\geq o_j+(h_j-|J'|)w'$, which is at least
  $o_{j+1}+(h_j-1-|J'|)w$ and at least $o_{j+1}+1$. So $[F_{j+1}(J)]_{o_{j+1}}$ is the sum of the
  $\pm z_{j,m}P_j(J'\cup\{m\})T_j(J'')$, $m\leq r$, $m\notin J'$; as $(h_j-r)^+=0$, only $J''=\emptyset$
  contributes, with the signs of $P_j\lrcorner(z_{j,1},\dots,z_{j,r})$ at $J'$.
\item[(v)] $c\leq0$, $h_{j+1}=h_j+1$. A vector letter and a contraction step; $h_j\geq r$, so $j$ is
  costly, $o_{j+1}=o_j+\omega(j)$, $P_{j+1}=P_j$ and $T_{j+1}=T_j\lrcorner(z_{j,r+1},\dots,z_{j,n})$. A
  term with $m\leq r$ has $\operatorname{ord}\geq o_j+(|J'|+1)w'$, at least $o_{j+1}+|J'|w$ and at
  least $o_{j+1}+1$; one with $m>r$ has $\operatorname{ord}\geq o_{j+1}+|J'|w'$ and contributes
  $\pm z_{j,m}P_j(J')T_j(J''\cup\{m-r\})$ with the sign $\epsilon_{J''}(m-r)$.
\item[(vi)] $c\leq0$, $h_{j+1}=h_j-1$. A covector letter and a wedge step; $h_j>h_{j+1}\geq r$, so $j$
  is not costly, $P_{j+1}=P_j$ and $T_{j+1}=T_j\wedge(z_{j,r+1},\dots,z_{j,n})$. A term with $m\leq r$
  has $\operatorname{ord}\geq\omega(j)+o_j+(|J'|-1)w'\geq o_j+|J'|w$ and $\operatorname{ord}>o_j$; one with
  $m>r$ has $\operatorname{ord}\geq o_j+|J'|w'$ and contributes
  $\pm z_{j,m}P_j(J')T_j(J''\smallsetminus\{m-r\})$ with the sign $\epsilon_{J''}(m-r)$.
\end{enumerate}
In each case $P_{j+1}$ and $T_{j+1}$ are as stated, since the heights of $\nu$ and $\kappa$ are
$\min(h,r)$ and $(h-r)^+$ and the steps for $\nu$ and $\kappa$ are taken with the same $c$. This proves
(A$_{j+1}$) and (D$_{j+1}$).

(4) Put $I(c)=I'(c)\cup\{r+b:b\in I''(c)\}$, so that $I(c)\cap\{1,\dots,r\}=I'(c)$ and
$\{b:r+b\in I(c)\}=I''(c)$. By (3), each factor $\mathrm{Col}_{\mu,c}(I(c))$ of $\Psi(\mu,I)$ has all its
monomials of $\omega$-degree at least $o_\omega(\mu,c)$, with lowest part
$\iota_{\mathrm{pl}}(\mathrm{Col}_{\nu,c}(I'(c)))\,\iota_{\mathrm{tn}}(\mathrm{Col}_{\kappa,c}(I''(c)))$. For a
product of polynomials $Q_c$ whose monomials have $\omega$-degree at least $o_c$, the monomials have
$\omega$-degree at least $\sum_co_c$ and the part of degree $\sum_co_c$ is the product of the parts
of degree $o_c$. So every monomial of $\Psi(\mu,I)$ has $\omega$-degree at least $o$, and its part of
degree $o$ is the product over $|c|\leq T_\mu$ of the lowest parts above. As $T_\nu,T_\kappa\leq T_\mu$,
it remains to see that the factors with $T_\nu<|c|\leq T_\mu$ in the product for $\nu$ are $1$, and the
same for $\kappa$. For such $c$ all the entries of all the $\nu_i$ are less than $c$, or all are at
least $c$; so the heights do not change, $\mathrm{Col}_{\nu,c}=F_0$, $I'(c)=\emptyset$ as $I'$ is a
standard filling, and $\mathrm{Col}_{\nu,c}(\emptyset)=1$. The product is therefore
$\iota_{\mathrm{pl}}(\Psi(\nu,I'))\,\iota_{\mathrm{tn}}(\Psi(\kappa,I''))$.

(5)(a) As $g\in GL_n(W)$, the entries of $g_v=\iota\sigma^v(g)$ and of
$g_v^{-1}=\iota\sigma^v(g^{-1})$ lie in $\mathcal R$, so the entries of each row $\ell_i$ lie in
$\mathcal R[z]$. The function $F_0$ takes the values $0$ and $1$, and the steps and the powers in the
definition of $\mathrm{Col}^g_{\mu,c,\zeta}$ form sums of products with signs. So the values lie in
$\mathcal R[z]$.

(b) Let $i_1<\dots<i_s$ be the letters used for $c$. The entries of $\ell_i$ are linear forms in
$z_i$. We show by induction on $t$ that every monomial of each $F_t(J)$ has degree
$p^{\pi_{i_t}-\pi_{i_u}}$ in $z_{i_u}$ for $u\leq t$ and degree $0$ in the other rows. For $t=0$,
$F_0(J)$ is $0$ or $1$. Raising to the power $e_t=p^{\pi_{i_t}-\pi_{i_{t-1}}}$ multiplies the degrees by
$e_t$, and $p^{\pi_{i_t}-\pi_{i_{t-1}}}p^{\pi_{i_{t-1}}-\pi_{i_u}}=p^{\pi_{i_t}-\pi_{i_u}}$, as the positions
are nondecreasing; the step at $i_t$ multiplies by an entry of $\ell_{i_t}$ and adds degree $1$ in
$z_{i_t}$. The final power $p^{\zeta-\pi_{i_s}}$ gives the degrees $p^{\zeta-\pi_{i_u}}$.

(e) Put $E=\zeta-\pi_0$. For a costly letter $i$, $0\leq\zeta-\pi_i\leq E$, as the positions are
nondecreasing and at most $\zeta$. A costly letter is used, so if $i<i'$ are costly, then
$\pi_i<\pi_{i'}$; the exponents $\zeta-\pi_i$ of the costly letters are therefore distinct. Hence
$o=\sum_i p^{\zeta-\pi_i}$, over the costly $i$, is at most $\sum_{e=0}^Ep^e\leq2p^E$, since
$\sum_{e<E}p^e<p^E$ for $p\geq2$.

(c) Let $g'\in GL_n(W)$, and let $y_i$, $x_i$ and $h$ be as in Lemma~\ref{pd-symbols-path-lift-order}(2);
the set $P_i$ there consists of the $b$ for which $z_{i,b}$ has positive $\omega$-degree. The map
$\Sigma^{\mathrm L}_{g'}$ is a $\mathcal K$-algebra homomorphism, so it commutes with the steps and the
powers in the definition of $\mathrm{Col}^g_{\mu,c,\zeta}$, and
$\Sigma^{\mathrm L}_{g'}(\mathrm{Col}^g_{\mu,c,\zeta}(J))=Z_N(x';\zeta)(J)$ with
$x'_i=\Sigma^{\mathrm L}_{g'}(\ell_i)$, in the notation of Lemma~\ref{pd-symbols-path-lift-order}(1). For a
vector letter, $z_i\delta=y_i$ and
$\Sigma^{\mathrm L}_{g'}(z_ig_{\pi_i})=z_i\delta\,g'^{-1}_{\pi_i}g_{\pi_i}=y_i(g'^{-1}g)_{\pi_i}=x_i$; for a
covector letter, $z_i\delta'=y_i$ and $\Sigma^{\mathrm L}_{g'}(z_i\,{}^tg_{\pi_i}^{-1})=z_i\delta'\,{}^tg'_{\pi_i}\,
{}^tg_{\pi_i}^{-1}=y_i\,{}^t(g^{-1}g')_{\pi_i}=x_i$. So $x'=x$. By Lemma~\ref{pd-symbols-path-lift-order} with
$j=N$ and $s=\zeta$, every value of $h_\zeta\cdot Z_N(x;\zeta)$ lies in $\varpi_0^{o_N(\zeta)}\mathcal R[z]$, and
$o_N(\zeta)=o$. By Lemma~\ref{pd-symbols-path-lift-twist}(1), every value of $Z_N(x;\zeta)$ lies in
$\varpi_0^o\mathcal R[z]$.

(d) Take $g'=g$ in (c); then $x_i=y_i$, and $\Sigma^{\mathrm L}_g(\mathrm{Col}^g_{\mu,c,\zeta}(J))=Z_N(y;\zeta)(J)$.
By Lemma~\ref{pd-symbols-path-lift-residue}(2),(3), this is $\varpi_0^oY$ with $Y\in\mathcal R[z]$ whose reduction
modulo $\varpi_0$ is $\pi_\zeta(\mathrm{Col}^{\mathrm{pt}}_{\mu,c}(J))$. The polynomial
$\Pi_\zeta(\mathrm{Col}^{\mathrm{pt}}_{\mu,c}(J))\in\mathcal R[z]$ has the same reduction. So
$Y-\Pi_\zeta(\mathrm{Col}^{\mathrm{pt}}_{\mu,c}(J))\in\varpi_0\mathcal R[z]$, and the difference in (d) is
$\varpi_0^o$ times it.

(6) Let $Z_A=(z_{A(l),b})_{1\leq l,b\leq r}$, so that $\Delta_A=\det Z_A$, and let $Z_A^{\mathrm{ad}}$ be its
adjugate, so that $Z_AZ_A^{\mathrm{ad}}=\Delta_A\cdot1$; put $\gamma_b=(Z_A^{\mathrm{ad}})_{b,r}$ for $1\leq b\leq r$.
Expanding along the last row, $\gamma_b$ is the image of $\Delta_A$ under the substitution $\epsilon_b$ that
replaces the row $z_{A(r)}$ by the $b$-th unit vector and fixes the other variables; so $\gamma_b$ involves only
the $z_{A(l),c}$ with $l<r$ and $c\neq b$. Up to sign, $\gamma_b$ is the determinant of the matrix obtained from
$Z_A$ by deleting row $r$ and column $b$, whose entries are distinct variables, so it is not $0$: its value at a
suitable permutation matrix is $\pm1$.

(a) \emph{Step 1: $\Delta_A$ and $\gamma_r$ are coprime.} Let $h\in F[z]$ divide both. In a polynomial ring over a
field the degree of a product of nonzero polynomials in each variable is the sum of the degrees, so a divisor
of a nonzero polynomial involves only variables that the polynomial involves. Hence $h$ involves only the
$z_{A(l),c}$ with $l<r$ and $c<r$; so $\epsilon_b(h)=h$, and $h$ divides $\epsilon_b(\Delta_A)=\gamma_b$ for every
$b$. Then $h$ involves no variable $z_{A(l),c}$ with $c\neq b$, for every $b\leq r$; so $h$ is a constant, nonzero
as $\gamma_r\neq0$, hence a unit.

\emph{Step 2: a change of frame.} Let $F'=F[z][\gamma_r^{-1}]$, which contains $F[z]$, and let $g_0$ be the
$r\times r$ matrix over $F'$ that equals the identity except that its last column is ${}^t(\gamma_1,\dots,\gamma_r)$;
its determinant $\gamma_r$ is a unit of $F'$. Let $F'[u]$ be the polynomial ring over $F'$ in variables $u_{i,b}$,
$0\leq i<N$, $1\leq b\leq r$, with rows $u_i$. Let $\beta\colon F[z]\to F'[u]$ be the $F$-algebra map with
$z_i\mapsto u_ig_0^{-1}$ for the vector letters $i\in S$, $z_i\mapsto u_i\,{}^tg_0$ for the covector letters $i\in
S$ and $z_i\mapsto u_i$ for $i\notin S$, and $\epsilon\colon F'[u]\to F'$ the $F'$-algebra map with
$u_i\mapsto z_ig_0$, $u_i\mapsto z_i\,{}^tg_0^{-1}$ and $u_i\mapsto z_i$ in the three cases. Then
$\epsilon(\beta(P))=P$ for $P\in F[z]$. As $\beta$ maps the entries of each row $z_i$ to linear forms in $u_i$,
every monomial of $\beta(G)$ has degree $d(i)$ in $u_i$ for every $i$.

\emph{Step 3: weights.} For a monomial $u^m$ put $w(m)=\sum_im(i,r)-\sum_{i'}m(i',r)$, the sums over the vector
letters $i\in S$ and the covector letters $i'\in S$, and let $\beta(G)_m$ be the coefficient of $u^m$ in
$\beta(G)$. We show that $w(m)\geq\operatorname{val}_S(G)$ whenever $\beta(G)_m\neq0$. The polynomial $G_{S,r}$ of
Definition~\ref{pd-symbols-path-generators}(1) is a Laurent polynomial in $\eta$ with coefficients in
$F[\mathfrak g][\det\mathfrak g^{-1}][z]$. Apply to it the ring map that sends $\mathfrak g$ to $g_0$, which is
possible as $\det g_0$ is a unit of $F'$, sends the row $z_i$ to $\beta(z_i)$ for every $i$, and fixes $\eta$. It
sends $\tau_r(\eta)$ to $g_0D_r(\eta)g_0^{-1}$, so $z_i\tau_r(\eta)\mapsto u_iD_r(\eta)g_0^{-1}$ for the vector
letters $i\in S$ and $z_i\,{}^t\tau_r(\eta^{-1})\mapsto u_iD_r(\eta^{-1})\,{}^tg_0$ for the covector letters
$i\in S$. So the image of $G_{S,r}$ is $\beta(G)$ with $u_{i,r}$ replaced by $\eta u_{i,r}$ for the vector letters
and by $\eta^{-1}u_{i,r}$ for the covector letters in $S$, that is, $\sum_m\beta(G)_m\eta^{w(m)}u^m$. A ring map
that fixes $\eta$ creates no new exponents of $\eta$, so every $w(m)$ with $\beta(G)_m\neq0$ is an exponent of
$\eta$ in $G_{S,r}$, and is at least $\operatorname{val}_S(G)$.

\emph{Step 4: divisibility in $F'$.} For $l\leq r$, $\epsilon(u_{A(l),r})=\sum_bz_{A(l),b}\gamma_b$ is the entry
$(l,r)$ of $Z_AZ_A^{\mathrm{ad}}$, which is $\Delta_A$ for $l=r$ and $0$ for $l<r$. Put $e=\theta-d_A(S)\geq1$ and
let $u^m$ be a monomial of $\beta(G)$. If $m(A(l),r)>0$ for some $l<r$, then $\epsilon(u^m)=0$. Otherwise, in
$w(m)$ each vector letter $i\in S$ that is not a value of $A$ contributes $m(i,r)\leq d(i)$, the $A(l)$ with $l<r$
contribute $0$ and the covector letters at most $0$. By Step~3,
$\theta\leq w(m)\leq m(A(r),r)+d_A(S)$, so $\epsilon(u^m)$ is a multiple of
$\epsilon(u_{A(r),r})^{m(A(r),r)}=\Delta_A^{m(A(r),r)}$, hence of $\Delta_A^e$. By Step~2, $G=\epsilon(\beta(G))$ is
a multiple of $\Delta_A^e$ in $F'$.

\emph{Step 5.} So $G\gamma_r^s=\Delta_A^ew$ for some $s\geq0$ and $w\in F[z]$. The ring $F[z]$ is a unique
factorization domain, and by Step~1 no prime factor of $\Delta_A$ divides $\gamma_r$; hence $\Delta_A^e$ divides
$G$.

(b) The map $P\mapsto P_{S',r}$ is a ring homomorphism into the Laurent polynomials in $\eta$ over the domain
$F[\mathfrak g][\det\mathfrak g^{-1}][z]$, and there the lowest exponent of $\eta$ in a product of nonzero elements
is the sum of the lowest exponents. So $\operatorname{val}_{S'}(PQ)=\operatorname{val}_{S'}(P)+\operatorname{val}_{S'}(Q)$
for all $P,Q\in F[z]$. Moreover $\operatorname{val}_{S'}(\Delta_A)\leq1$: the ring map that sends $\mathfrak g$ to the
identity matrix, $z_{A(l),b}$ to $1$ for $b=l$ and to $0$ for $b\neq l$, the other variables to $0$, and fixes
$\eta$, sends $(\Delta_A)_{S',r}$ to the determinant of the diagonal matrix with entries $1,\dots,1,\eta$ if
$A(r)\in S'$ and $1,\dots,1,1$ otherwise, as the $A(l)$ are vector letters; so the exponent $1$ or $0$ occurs in
$(\Delta_A)_{S',r}$. Hence $\operatorname{val}_{S'}(G)=m\operatorname{val}_{S'}(\Delta_A)+\operatorname{val}_{S'}(H)\leq
\operatorname{val}_{S'}(H)+m$.
% Lean: (6) is CovariantDegrees.ConclusionE (= FieldMultidegreeDivisibility, the clause adm21 of
%   tableLineDegreeFive.admitted), a component of CovariantDegrees.Conclusion1 proved as
%   GLWeightTheory.support_coprimality_lemma_holds = prvE_of_gl_operators_form_representation (in
%   counting_statement_of_block_form; the extremal projector does not enter it). (a) is height_bound_of_weights
%   hSupp_proved hCop_proved: Step 1 is hCop_proved (adjugate_last_row_ne_zero, divisor_supported_of_supported,
%   cofactor_last_row_supported, lastRowSubstitution with substitution_det_eq_adjugate); Step 2 is
%   embed_GL2_standard_swap, frameMatrixInverse, eBeta, frameCoordinateEvaluation, frame_evaluation_comp_change and
%   frame_change_preserves_multidegree; Step 3 is hSupp_proved (mul_square, mul_Psi_support, mul_dl_coeff); Step 4
%   is hX with last_column_weight_bound; Step 5 is IsLocalization.surj and IsRelPrime.dvd_of_dvd_mul_right. (b) is
%   valuation_bound_det_power_factor (mul_valS_mul with mul_laurent_inf, mul_valS_detA_le).
% Lean: (1) is ColumnFamilies.splitting_third_part_holds = lemma_A3_over_domain and lemma_A3_over_field
%   (valuation_bounds_iff_mem_span_words, mem_span_words_of_valuation_bounds with bounded_dominant_lower_set
%   and support_subset_of_mem_span_image, valuation_bounds_of_mem_span_words with the sDegs facts
%   support_degrees_map_subset, _add_subset, _smul_subset). Lemma pd-symbols-path-word-bounds is
%   word_valuation_bounds, and Lemmas pd-symbols-path-basis and pd-symbols-path-valuation are proved in the
%   book (round nine). Lemma pd-symbols-path-independent = GelfandTsetlinBases.Top.lemmaA2indep and
%   FirstLettersLowerSet (Lemma pd-symbols-path-lower-set) are proved in the book (round ten). Earlier
%   notes: ValuationsAdditive (valuation_field = valuation_field_of_adm2 with weight_condition_from_integral_form).
% (2) is prvB1_of_D7 with hstrip_append, antitone_append, antitone_of_append (AntitoneOfAppend.lean).
% (3) is lemma_B2 = B2.kcol_low: mkDat, Dat.inv_pos (c >= 1) and Dat.inv_neg (c <= 0) with A' = ihA (om j + 1),
%   KG_card_pos, ordWedge/ordContr, wedge_pl, wedge_tan, contr_pl, contr_tan, ar0 to ar3, colH_strip, colH_pl,
%   colH_tan (EqEmptyIffSplit.lean).
% (4) is the second clause of lemma_B3, joint_type_lower_bound D (lemma_B2 D) with I = joinI I' I'',
%   low_prod, prod_plTan, psi_eq_prod_superset, kcol_of_notMem (JointTypeLatticeSaturated.lean).
% (5) is canonical_translate_jet_identity: (a) column_kernel_components_integral (letG_int, ZZ_mem);
%   (b) lifted_column_multidegree (ZZ_HomD, homD_KKs_and_slot_properties); (e) column_order_le_two_power.
%   (c) lifted_column_order (FrobCong from frobCong, mkTS, zetaL, lift_substitution_column_kernel,
%   Ord_main, TS.inv) and (d) lifted_column_congruence (lift_KKc, Inv_main, SgVanish_of, plTanCol_eq_SgW,
%   colOK_of_path, piR_ker), FrobeniusPaths/.
% Overall: CovariantDegrees.Conclusion2 D, proved as ColumnFamilies.splitting_package_holds D
%   (BinomialPolynomialIdentities/SplittingThirdPartHolds.lean): (1) is splitting_third_part_holds with
%   LemmaBAlt.lemma_A3_over_domain and lemma_A3_over_field (word_valuation_bounds, valuation_field,
%   admissible_standardPath_basis, admissible_word_basis); (2) is LemmaBAlt.prvB1_of_D7; (3) is
%   LemmaBAlt.lemma_B2; (4) is LemmaBAlt.nested_sets_lemma_of_admitted with adm_fourteen and
%   admissible_externalProduct; (5) is LemmaBAlt.canonical_translate_jet_identity (FrobeniusPaths/).
%   The extremal projector hypothesis ExtremalProjectorTools does not enter this proof.
\end{proof}

\subsection{Edge families}\label{subsection-pd-edge-families}

\begin{lemma}\label{pd-subspace-count}
Let $F$ be a finite field with $q$ elements and let $0\leq r\leq k$ be integers. Then $F^k$ has at
least $q^{r(k-r)}$ subspaces of dimension $r$.
\end{lemma}
\leanref{Formalism/SharpRestriction/GrowthBounds/PlaneCountTheoremWithConstants.lean}{Formalism.MainNotions.G.gauss_binomial_ge}
\leanref{Formalism/SharpRestriction/GrowthBounds/PlaneCountTheoremWithConstants.lean}{Formalism.MainNotions.G.rational_planes_count}
\leanref{Formalism/SharpRestriction/GrowthBounds/ValuationLowestDegreeRename.lean}{Formalism.MainNotions.SelectionLemmas.ColumnFamilyClosure}

\begin{proof}
Let $G(n,j)$ be the number of subspaces of dimension $j$ of $F^n$. Then $G(n,0)=1$,
$G(0,j+1)=0$ and $G(n+1,j+1)=G(n,j)+q^{j+1}G(n,j+1)$, the recursion of the Gaussian binomial
coefficients. We show $G(n,j)\geq q^{j(n-j)}$ for $j\leq n$ by induction on $n$. For $j=0$ both sides
are $1$. For $j+1=n+1$, $G(n+1,n+1)\geq G(n,n)\geq1$. For $j+1\leq n$,
$G(n+1,j+1)\geq q^{j+1}G(n,j+1)\geq q^{j+1}q^{(j+1)(n-j-1)}=q^{(j+1)(n-j)}$.
\end{proof}

\begin{lemma}\label{pd-edge-basis}
Let $p\geq5$ be a prime, $(W,Q,\sigma,\iota)$ a Hida setting of degree $n'$ at $p$ with rings
$\mathcal R\subset\mathcal K$ and $\varpi_0$, $k\geq1$ and $r\geq0$ integers, and
$\lambda_0(s)\in\mathbf Z^k$, $s\in\mathbf Z/n'\mathbf Z$, dominant weights. Let $V$, its dual
$V^\vee$, the lattices $H(L)$, $H(e)$, $R(e)$ and $R'(e)$, the integer $A$ and the edge pairing be
those of Definition~\ref{pd-projection-property}, and let $e$ be an edge with source $L''$ and
target $L$.
\begin{enumerate}
\item There are a $\mathcal K$-basis $(x_i)$ of $V$, with dual basis $(x_i^*)$ of $V^\vee$, and
  integers $\pi_i\leq A$ such that $H(L'')$, $H(L)$, $H(e)$, $R(e)$ and $R'(e)$ are the
  $\mathcal R$-spans of the $x_i$, of the $p^{\pi_i}x_i$, of the $\varpi_0^{\pi_i}x_i$, of the
  $\varpi_0^{\pi_i-A}x_i$ and of the $\varpi_0^{\pi_i-A}p^{-\pi_i}x_i^*$ respectively.
\item $H(L'')\subset R(e)$ and $H(L)^\vee\subset R'(e)$.
\item An element $G\in V^\vee$ lies in $R'(e)$ if and only if $\langle F,G\rangle\in\mathcal R$ for
  every $F\in R(e)$. An element $F\in V$ lies in $R(e)$ if and only if
  $\langle F,G\rangle\in\mathcal R$ for every $G\in R'(e)$.
\end{enumerate}
\end{lemma}
\leanref{Formalism/SharpRestriction/GrowthBounds/ObjBridge.lean}{Formalism.MainNotions.G.DualHeightBasis}
\leanref{Formalism/LocalToGlobal/GroupLattices/FockModels/ElementaryDivisorBasis.lean}{Formalism.L2.G.Model.dualHeightBasis}
\leanref{Formalism/LocalToGlobal/GroupLattices/FockModels/ElementaryDivisorBasis.lean}{Formalism.L2.G.Model.elementary_divisor_basis}
\leanref{Formalism/SharpRestriction/GrowthBounds/ObjBridge.lean}{Formalism.MainNotions.G.DualHeightInclusion}
\leanref{Formalism/LocalToGlobal/GroupLattices/FockModels/DualInclusionLattices.lean}{Formalism.L2.G.Model.dualInclusionLattices}
\leanref{Formalism/SharpRestriction/GrowthBounds/ObjBridge.lean}{Formalism.MainNotions.G.BePerfect}
\leanref{Formalism/LocalToGlobal/GroupLattices/FockModels/DualInclusionLattices.lean}{Formalism.L2.G.Model.bePerfect}

\begin{proof}
(1) The quotient $L/L''$ is killed by $p$ and has $p^{n'r}$ elements, and $W/pW$ has $p^{n'}$
elements. By the elementary divisor theorem $L''$ has a $W$-basis $b_1,\dots,b_k$ such that
$p^{-1}b_1,\dots,p^{-1}b_r,b_{r+1},\dots,b_k$ is a basis of $L$. For each $s$ let $g_s$ be the
matrix with columns $\iota\sigma^s(b_1),\dots,\iota\sigma^s(b_k)$, and for $c\in\mathcal K^\times$
let $d_c$ be the diagonal matrix with entries $c^{-1}$ in the first $r$ places and $1$ in the
others. The columns of $g_s$, $g_sd_{\varpi_0}$ and $g_sd_p=g_sd_{\varpi_0}^2$ are $\mathcal R$-bases
of the lattices from which $H(L'')$, $H(e)$ and $H(L)$ are formed at $s$. For a matrix $g$, the
lattice formed from the columns of $g$ is the image of the one formed from the standard basis under
the action of $g$ of Definition~\ref{pd-local-model}(3). Choose for each $s$ an $\mathcal R$-basis of
the latter lattice for $\lambda(s)$ consisting of polynomials $f$ with $f(tx)=t^\beta f$ for the
diagonal matrices $t$, for some $\beta\in\mathbf Z^k$, and put $\mu=\beta+\lambda(s)_k(1,\dots,1)$.
Then $d_c$ multiplies $f$ by $c^{\mu_1+\cdots+\mu_r}$, and $\mu$ is a weight of the representation of
highest weight $\lambda(s)$. Let the $x_i$ be the products over $s$ of the chosen polynomials moved by
the $g_s$, and $\pi_i$ the sum over $s$ of the corresponding $\mu_1+\cdots+\mu_r$. Then, for every $c\neq0$, the lattice formed from the columns of the $g_sd_c$ is the
$\mathcal R$-span of the $c^{\pi_i}x_i$, and $\pi_i\leq A$ because every weight $\mu$ of the
representation of highest weight $\lambda(s)$ satisfies
$\mu_1+\cdots+\mu_r\leq\lambda(s)_1+\cdots+\lambda(s)_r$. Taking $c=1,\varpi_0,p$ gives $H(L'')$,
$H(e)$ and $H(L)$. Then $R(e)=\varpi_0^{-A}H(e)$ is spanned by the
$\varpi_0^{\pi_i-A}x_i$. The lattice $H(e)^\vee$ is spanned by the $\varpi_0^{-\pi_i}x_i^*$, so
$R'(e)$ is spanned by the $\varpi_0^{-\pi_i-A}x_i^*=\varpi_0^{\pi_i-A}p^{-\pi_i}x_i^*$, since
$\varpi_0^2=p$.

(2) First, $\varpi_0^AH(L'')\subset H(e)$ and $\varpi_0^AH(e)\subset H(L)$. Let $s$ be a position, $g$
one of the matrices $g_s$ and $g_sd_{\varpi_0}$ of (1), and $f$ one of the chosen polynomials for
$\lambda(s)$, of weight $\mu$, moved by $g$. Moving $f$ by $gd_{\varpi_0}$ instead multiplies it by
$\varpi_0^{\mu_1+\cdots+\mu_r}$, so $\varpi_0^{\lambda(s)_1+\cdots+\lambda(s)_r}f=
\varpi_0^{\lambda(s)_1+\cdots+\lambda(s)_r-(\mu_1+\cdots+\mu_r)}\cdot\varpi_0^{\mu_1+\cdots+\mu_r}f$
lies in the lattice formed from the columns of $gd_{\varpi_0}$, as the exponent of the first factor
is $\geq0$. The lattice does not depend on the choice of the matrix whose columns form a basis, so
multiplication by $\varpi_0^{\lambda(s)_1+\cdots+\lambda(s)_r}$ carries the lattice formed from the
columns of $g$ into the one formed from those of $gd_{\varpi_0}$. Taking products over $s$ gives the
two inclusions. The first gives $H(L'')\subset\varpi_0^{-A}H(e)=R(e)$. For $G\in H(L)^\vee$ and $z\in
H(e)$, $(\varpi_0^AG)(z)=G(\varpi_0^Az)\in\mathcal R$ by the second; so $\varpi_0^AG\in H(e)^\vee$ and
$G\in\varpi_0^{-A}H(e)^\vee=R'(e)$.

(3) The bases $F_i=\varpi_0^{\pi_i-A}x_i$ of $R(e)$ and $G_j=\varpi_0^{\pi_j-A}p^{-\pi_j}x_j^*$ of
$R'(e)$ satisfy $\langle F_i,G_j\rangle=G_j(p^AF_i)=\delta_{ij}\,p^{A-\pi_i}\varpi_0^{2\pi_i-2A}=\delta_{ij}$.
So each of $R(e)$ and $R'(e)$ is the set of vectors whose pairings with the other lie in
$\mathcal R$.
\end{proof}

\begin{lemma}\label{pd-edge-transport}
Let $p\geq5$ be a prime, $(W,Q,\sigma,\iota)$ a Hida setting of degree $n'$ at $p$, $k\geq1$ and
$r,t,D,c_*\geq0$ integers, $\kappa\in\mathbf R$, and $\lambda_0(s)\in\mathbf Z^k$ dominant weights.
Suppose that $4(D+1)\leq c_*$ and that the condition of
Definition~\ref{pd-projection-property}(3) on edges, vectors, $n$, $u$ and $(B_{ji})$ holds for one
$W$-lattice $L_1\subset Q^k$. Then the projection property
$\mathrm{LP}(r,\lambda_0;t,D,\kappa,c_*)$ holds.
\end{lemma}
\leanref{Formalism/SharpRestriction/GrowthBounds/ObjBridge.lean}{Formalism.MainNotions.G.LPOneEvery}
\leanref{Formalism/LocalToGlobal/GroupLattices/FockModels/LoBlockActComp.lean}{Formalism.L2.G.Model.lPOneEvery}
\leanref{Formalism/LocalToGlobal/GroupLattices/FockModels/LoBlockActComp.lean}{Formalism.L2.G.Model.lo_at}

\begin{proof}
Let $L\subset Q^k$ be a $W$-lattice. Both $L_1$ and $L$ are spanned by bases of $Q^k$, so $L=\gamma L_1$
for some $\gamma\in GL_k(Q)$. Let $\gamma$ act on $\mathcal K[x_{s,a,b}]$ as in
Definition~\ref{pd-local-model}(3), with one Hida setting. This action preserves $V$ and carries
$H(M)$ onto $H(\gamma M)$ for every $W$-lattice $M$. The map $M\mapsto\gamma M$ preserves
inclusions, multiplication by $p$ and indices, so it carries an edge $e$ with source $L''$ and
target $M$ to the edge $\gamma e$ with source $\gamma L''$ and target $\gamma M$, and $\gamma$
carries $H(e)$ onto $H(\gamma e)$ and $R(e)$ onto $R(\gamma e)$. Let $\gamma$ act on $V^\vee$ by
$G\mapsto G\circ\gamma^{-1}$. This carries $H^\vee$ onto $(\gamma H)^\vee$ for every lattice $H$,
hence $R'(e)$ onto $R'(\gamma e)$, and
$\langle\gamma F,G\circ\gamma^{-1}\rangle=\langle F,G\rangle$. Hence $\gamma$ carries
$N_{\mathrm{src}}(L'')$ onto $N_{\mathrm{src}}(\gamma L'')$ and $N_{\mathrm{tgt}}(M)$ onto
$N_{\mathrm{tgt}}(\gamma M)$. Apply $\gamma$ to the edges $e_i$ and the vectors $F_i$ given at $L_1$,
replace $u$ by $y\mapsto u(y)\circ\gamma^{-1}$, and keep $n$ and $(B_{ji})$. These satisfy the
condition at $L$.
\end{proof}

\begin{definition}\label{pd-edge-family}
Let $p\geq5$ be a prime, $(W,Q,\sigma,\iota)$ a Hida setting of degree $n'$ at $p$ with rings
$\mathcal R\subset\mathcal K$ and $\varpi_0$, $k\geq1$ and $r\geq0$ integers, and
$\lambda_0(s)\in\mathbf Z^k$, $s\in\mathbf Z/n'\mathbf Z$, dominant weights. Let $V$, $V^\vee$, the
lattices $H(L)$, $R(e)$ and $R'(e)$ and the edge pairing be those of
Definition~\ref{pd-projection-property}. Let $e$ be an edge with source $L''$ and target $L$, and let
$t_S,t_T,D_1,c_*,n\geq0$ be integers. An \emph{edge family} at $e$ of size $n$, with \emph{levels}
$(t_S,t_T,D_1)$ and \emph{margin} $c_*$, is a family of vectors $F_1,\dots,F_n\in V$ such that:
\begin{enumerate}
\item $F_i\in\varpi_0^{-t_S}H(L'')\cap R(e)$ for every $i$;
\item $F_i\in\varpi_0^{c_*}R(e')$ for every $i$ and every edge $e'\neq e$ with source $L''$;
\item for every $c\in\mathcal R^n$ with some $c_i$ a unit there is
  $G\in\varpi_0^{-t_T}H(L)^\vee\cap R'(e)$, with $G\in\varpi_0^{c_*}R'(e')$ for every edge $e'\neq e$
  with target $L$, such that $\langle\sum_ic_iF_i,G\rangle$ lies in $\mathcal R$ and divides
  $p^{D_1}$ in $\mathcal R$.
\end{enumerate}
\end{definition}
\leanref{Formalism/SharpRestriction/GrowthBounds/ScaleEdges/ManyEdges.lean}{Formalism.MainNotions.Pf14.EdgeFam}
\leanref{Formalism/SharpRestriction/GrowthBounds/ScaleEdges/ManyEdges.lean}{Formalism.MainNotions.Pf14.ValLe}

\begin{lemma}\label{pd-edge-assembly}
Let $p\geq5$ be a prime, $(W,Q,\sigma,\iota)$ a Hida setting of degree $n'$ at $p$ with rings
$\mathcal R\subset\mathcal K$ and $\varpi_0$, $k\geq1$ and $r,t,D,c_*\geq0$ integers with
$2(D+1)\leq c_*$, $\kappa\in\mathbf R$, $\lambda_0(s)\in\mathbf Z^k$ dominant weights, and
$L\subset Q^k$ a $W$-lattice; let $N_{\mathrm{src}}$ and $N_{\mathrm{tgt}}$ be formed with $t$
(Definition~\ref{pd-projection-property}). Let $E$ be a finite set of edges with target $L$, and for
each $e\in E$ let an edge family at $e$ of size $n_e$, with levels $(t_S(e),t_T(e),D_1(e))$ and
margin $c_*$, be given (Definition~\ref{pd-edge-family}), where $t_S(e)\leq t$, $t_T(e)\leq t$ and
$D_1(e)\leq D$. Suppose that $\sum_{e\in E}n_e\geq(1-\kappa)\dim V$. Write the vectors of all the
families as $F_1,\dots,F_M$, $M=\sum_en_e$, and let $e_i$ be the edge whose family contains $F_i$.
Then $F_i\in N_{\mathrm{src}}(L''_i)$, where $L''_i$ is the source of $e_i$; $F_i\in\varpi_0^{c_*}R(e')$
for every edge $e'\neq e_i$ with source $L''_i$; and there is an $\mathcal R$-linear map
$u\colon\mathcal R^M\to N_{\mathrm{tgt}}(L)$ with $\langle F_i,u(y)\rangle=p^Dy_i$ for all
$y\in\mathcal R^M$ and $i\leq M$. In particular the condition of
Definition~\ref{pd-projection-property}(3) holds at $L$, with $n=M$ and $(B_{ji})$ the identity matrix.
\end{lemma}
\leanref{Formalism/SharpRestriction/GrowthBounds/ScaleEdges/ManyEdges.lean}{Formalism.MainNotions.Pf14.Asm}
\leanref{Formalism/SharpRestriction/GrowthBounds/ScaleEdges/ScaleNegMono.lean}{Formalism.MainNotions.Pf14.asm}
\leanref{Formalism/SharpRestriction/GrowthBounds/ScaleEdges/ManyEdges.lean}{Formalism.MainNotions.Pf14.S7}
\leanref{Formalism/SharpRestriction/GrowthBounds/ScaleEdges/GoodRankOfPrim.lean}{Formalism.MainNotions.Pf14.s7}
\leanref{Formalism/SharpRestriction/GrowthBounds/ScaleEdges/GoodRankOfPrim.lean}{Formalism.MainNotions.Pf14.goodRank_of_prim}
\leanref{Formalism/SharpRestriction/GrowthBounds/ScaleEdges/DvdOfAdd.lean}{Formalism.MainNotions.Pf14.smith_mem}

\begin{proof}
Let $F$ be a vector of the family at $e\in E$, and $L''$ the source of $e$. Then
$F\in\varpi_0^{-t_S(e)}H(L'')\subset\varpi_0^{-t}H(L'')$, $F\in R(e)$, and
$F\in\varpi_0^{c_*}R(e')\subset R(e')$ for the other edges $e'$ with source $L''$; so
$F\in N_{\mathrm{src}}(L'')$. In the same way each form $G$ given by
Definition~\ref{pd-edge-family}(3) for the family at $e$ lies in $N_{\mathrm{tgt}}(L)$, and the
pairing there divides $p^D$ because $D_1(e)\leq D$.

For $G\in N_{\mathrm{tgt}}(L)$ put $\phi(G)=(\langle F_i,G\rangle)_{i\leq M}$. Since
$F_i\in R(e_i)$ and $G\in R'(e_i)$, Lemma~\ref{pd-edge-basis}(3) gives $\phi(G)\in\mathcal R^M$, and
$S=\phi(N_{\mathrm{tgt}}(L))$ is an $\mathcal R$-submodule of $\mathcal R^M$. Let $c\in\mathcal R^M$
with $c_{i_0}$ a unit, and put $e_0=e_{i_0}$. Apply Definition~\ref{pd-edge-family}(3) to the family at
$e_0$ and the entries of $c$ indexed by it. This gives $G\in N_{\mathrm{tgt}}(L)$ such that the sum
$a$ of the $c_i\langle F_i,G\rangle$ over the $i$ with $e_i=e_0$ divides $p^D$, and
$G\in\varpi_0^{c_*}R'(e)$ for every $e\in E$ other than $e_0$. For $i$ with $e_i\neq e_0$,
Lemma~\ref{pd-edge-basis}(3) gives $\langle F_i,G\rangle\in\varpi_0^{c_*}\mathcal R\subset
p^{D+1}\mathcal R$, as $c_*\geq2(D+1)$ and $\varpi_0^2=p$. Hence $\sum_ic_i\phi(G)_i=a+p^{D+1}z$
with $z\in\mathcal R$. Its valuation is that of $a$, because $p$ is not a unit in $\mathcal R$, so it
divides $p^D$.

By the elementary divisor theorem over the discrete valuation ring $\mathcal R$, there are a basis
$v_1,\dots,v_M$ of $\mathcal R^M$ and $\alpha_1,\dots,\alpha_M\in\mathcal R$ such that the
$\alpha_jv_j$ span $S$. Let $c\in\mathcal R^M$ be the vector with $\sum_ic_iy_i$ equal to the $j$-th
coordinate of $y$ in this basis. Some entry of $c$ is a unit, since this coordinate of $v_j$ is $1$,
and $\sum_ic_is_i\in\alpha_j\mathcal R$ for every $s\in S$. By the previous paragraph $\alpha_j$
divides $p^D$. Hence $p^D\mathcal R^M\subset S$. Choose $G_i\in N_{\mathrm{tgt}}(L)$ with
$\phi(G_i)=p^D\epsilon_i$, where $\epsilon_i$ is the $i$-th standard basis vector, and put
$u(y)=\sum_iy_iG_i$.
\end{proof}

\begin{lemma}\label{pd-edge-count}
Let $p$ be a prime, $W$ a discrete valuation ring with uniformizer $p$ whose residue field has
$p^{n'}$ elements, $Q$ its fraction field, and $0\leq r\leq k$ integers. Then there are at least
$p^{n'r(k-r)}$ edges with target $W^k$, that is, $W$-lattices $L''\subset W^k$ with
$pW^k\subset L''$ and $|W^k/L''|=p^{n'r}$.
\end{lemma}
\leanref{Formalism/SharpRestriction/GrowthBounds/ScaleEdges/ManyEdges.lean}{Formalism.MainNotions.Pf14.ManyEdges}
\leanref{Formalism/SharpRestriction/GrowthBounds/ScaleEdges/ColSpanMap.lean}{Formalism.MainNotions.Pf14.manyEdges}

\begin{proof}
Let $A$ be an $r\times(k-r)$ matrix over $W/pW$. Let $\psi_A\colon W^k\to(W/pW)^r$ send $x$ to the
vector with entries $x_i+\sum_cA_{ic}x_{r+c}$ modulo $p$, $i\leq r$. It is $W$-linear, and it is
surjective already on the first $r$ coordinates. Its kernel $L_A$ contains $pW^k$, so it is a
lattice, and $W^k/L_A\cong(W/pW)^r$ has $p^{n'r}$ elements. The matrix $A$ is determined by $L_A$:
for each $c$, the residue of $a\in W^r$ is the $c$-th column of $A$ exactly when
$e_{r+c}-\sum_ia_ie_i\in L_A$. There are $p^{n'r(k-r)}$ matrices $A$, hence as many edges
$(L_A,W^k)$.
\end{proof}

\begin{remark}\label{pd-edge-count-subspaces}
The sources of the edges with target $W^k$ are the preimages in $W^k$ of the subspaces of dimension
$k-r$ of $(W/pW)^k$. The lattice $L_A$ in the proof of Lemma~\ref{pd-edge-count} is the preimage of
the graph of $-A$, a subspace of this kind. Lemma~\ref{pd-subspace-count}, applied to all subspaces
of dimension $k-r$, gives the same lower bound $p^{n'(k-r)r}$.
\end{remark}

\begin{lemma}\label{pd-edge-products}
Let $p\geq5$ be a prime, $(W,Q,\sigma,\iota)$ a Hida setting of degree $n'$ at $p$ with rings
$\mathcal R\subset\mathcal K$ and $\varpi_0$, $k\geq1$ and $r\geq0$ integers, and
$\lambda^\flat(s),\lambda^\sharp(s)\in\mathbf Z^k$, $s\in\mathbf Z/n'\mathbf Z$, dominant weights;
put $\lambda_0=\lambda^\flat+\lambda^\sharp$. Let $V^\flat$, $H^\flat$, $R^\flat$, respectively
$V^\sharp$, $H^\sharp$, $R^\sharp$ and $V$, $H$, $R$, be the space and lattices of
Definition~\ref{pd-projection-property} formed with $\lambda^\flat$, respectively $\lambda^\sharp$ and
$\lambda_0$. Let $F\in V^\flat$ and $F'\in V^\sharp$. Then $FF'\in V$. Moreover, for every $W$-lattice
$L$, every edge $e$ and all integers $i,j$: if $F\in\varpi_0^iH^\flat(L)$ and
$F'\in\varpi_0^jH^\sharp(L)$ then $FF'\in\varpi_0^{i+j}H(L)$; and if $F\in\varpi_0^iR^\flat(e)$ and
$F'\in\varpi_0^jR^\sharp(e)$ then $FF'\in\varpi_0^{i+j}R(e)$.
\end{lemma}
\leanref{Formalism/SharpRestriction/GrowthBounds/TableLineDegreeFive.lean}{Formalism.MainNotions.tableLineDegreeFive.LevelPrimesProviso}
\leanref{Formalism/LocalToGlobal/GroupLattices/FockModels/MulMemRiemannRochSpaceAdd.lean}{Formalism.L2.G.Model.products_sections_integral}
\leanref{Formalism/SharpRestriction/GrowthBounds/ScaleEdges/EdgeLatticeMulOfWeightSum.lean}{Formalism.MainNotions.Pf14.HL_mul}
\leanref{Formalism/SharpRestriction/GrowthBounds/ScaleEdges/EdgeLatticeMulOfWeightSum.lean}{Formalism.MainNotions.Pf14.Re_mul}

\begin{proof}
For each $s$, the weight $\lambda(s)$ formed from $\lambda_0(s)$ is the sum of those formed from
$\lambda^\flat(s)$ and $\lambda^\sharp(s)$, and so are the natural weights. A product of elements of
the dual Weyl modules of two natural weights has the sum of their column degrees and the invariance of
Definition~\ref{pd-weight-lattices}(1), so it lies in the dual Weyl module of the sum; this holds over
$\mathcal K$ and over $\mathcal R$. Moving by a matrix $g$ as in Definition~\ref{pd-local-model}(3) is
multiplicative, and the determinant factors multiply. Hence a product of a spanning product of
$H^\flat(L)$ and one of $H^\sharp(L)$ is a spanning product of $H(L)$, and the same holds for
$H^\flat(e)$, $H^\sharp(e)$ and $H(e)$, which are formed in the same way from other lattices at each
$s$. This gives $FF'\in V$ and the claim for $H$. The integer $A$ of $\lambda_0$ is the sum of those of
$\lambda^\flat$ and $\lambda^\sharp$, so the claim for $R=\varpi_0^{-A}H(e)$ follows from the one for
$H(e)$.
\end{proof}

\begin{lemma}\label{pd-edge-dimension-count}
Let $n'\geq1$ and $k\geq3$ be integers, $T\subset\mathbf Z/n'\mathbf Z$, $a>0$ a real number,
$c_{pp},c_0\geq0$ integers with $kc_0\leq a$, $a_T,b_T\colon\mathbf Z/n'\to\mathbf Z$, and
$\lambda_0(s)\in\mathbf Z^k$ dominant weights such that $\lambda_0(s)_1=a_T(s)$,
$\lambda_0(s)_k=-b_T(s)$ and $|\lambda_0(s)_i|\leq c_{pp}$ for $1<i<k$ when $s\in T$,
$|\lambda_0(s)_i|\leq c_{pp}$ for all $i$ when $s\notin T$, and $a+c_{pp}\leq a_T(s)\leq2a-c_{pp}$,
$b_T(s)\geq0$ and $b_T(s)+2c_{pp}\leq c_0$ for $s\in T$. For a dominant weight $\mu\in\mathbf Z^j$ put
$\dim_j(\mu)=\prod_{1\leq i<l\leq j}(\mu_i-\mu_l+l-i)/(l-i)$. Let $V$ be the space of
Definition~\ref{pd-projection-property} formed with $\lambda_0$, let $d_t=\lambda_0(t)_1-\lambda_0(t)_2$
for $t\in T$, and put
\[
 n_2=\prod_{t\in T}\dim_{k-1}\big(\lambda_0(t)_2,\dots,\lambda_0(t)_k\big)\prod_{s\notin T}\dim_k(\lambda_0(s)).
\]
Then $\dim V=\prod_s\dim_k(\lambda_0(s))$ and
\[
 \Big(1-\frac{kc_0}a\Big)^{|T|}\dim V\leq n_2\prod_{t\in T}\binom{d_t+k-1}{k-1}.
\]
\end{lemma}
\leanref{Formalism/SharpRestriction/GrowthBounds/TableLineDegreeFive.lean}{Formalism.MainNotions.tableLineDegreeFive.FewBadOrbits}
\leanref{Formalism/PeriodIntegrals/AdmissibleModel/WeylDim.lean}{Formalism.L2.G.Model.fewBadOrbits}
\leanref{Formalism/SharpRestriction/GrowthBounds/ScaleEdges/RemainingWeightVectorCount.lean}{Formalism.MainNotions.Pf14.dimV_eq}
\leanref{Formalism/SharpRestriction/GrowthBounds/ScaleEdges/RemainingWeightVectorCount.lean}{Formalism.MainNotions.Pf14.top_count}
\leanref{Formalism/SharpRestriction/GrowthBounds/ScaleEdges/RemainingWeightVectorCount.lean}{Formalism.MainNotions.Pf14.count_of_exact}
\leanref{Formalism/SharpRestriction/GrowthBounds/ScaleEdges/RemainingWeightVectorCount.lean}{Formalism.MainNotions.Pf14.remainingWeightVectorCount}

\begin{proof}
The space $V$ is the tensor product over $s$ of the dual Weyl modules of the natural weights of the
$\lambda(s)$, and $\dim_k(\lambda(s))=\dim_k(\lambda_0(s))$, because $\dim_k$ depends only on the
differences of the entries and is unchanged by $\mu\mapsto(-\mu_k,\dots,-\mu_1)$. The first claim is
the Weyl dimension formula \cite[Theorem~6.3]{FH91}. For the second, let $t\in T$ and
$\mu_i=\lambda_0(t)_{i+1}+b_T(t)$ for $i<k$. Then $\mu$ is dominant with $0\leq\mu_i\leq c_0$, and
$d_t\geq a$. Splitting off the factors with $i=1$,
\[
 \dim_k(\lambda_0(t))=\dim_{k-1}(\mu)\prod_{l=2}^{k}\frac{d_t+\mu_1-\mu_{l-1}+l-1}{l-1},\qquad
 \binom{d_t+k-1}{k-1}=\prod_{l=2}^{k}\frac{d_t+l-1}{l-1},
\]
and each factor of the first product is at most $1+c_0/a$ times the corresponding factor of the
second. Since $(1-kc_0/a)(1+c_0/a)^{k-1}\leq e^{-kc_0/a}e^{(k-1)c_0/a}\leq1$, we get
$(1-kc_0/a)\dim_k(\lambda_0(t))\leq\binom{d_t+k-1}{k-1}\dim_{k-1}(\lambda_0(t)_2,\dots,\lambda_0(t)_k)$.
Multiplying over $t\in T$, and by the factors with $s\notin T$, gives the claim.
\end{proof}

\begin{situation}\label{pd-edge-situation}
Let $C\geq0$ be a real number, $p\geq13$ a prime, $(W,Q,\sigma,\iota)$ an unramified quadratic Hida
setting of degree $n'$ at $p$ with rings $\mathcal R\subset\mathcal K$ and $\varpi_0$
(Definition~\ref{pd-projection-property}), $m\geq0$ and $r\geq2$ integers, $k=m+1$,
$T\subset\mathbf Z/n'\mathbf Z$, $L_1,L_2\geq0$ integers and $\delta,h,x,a,Y>0$ reals. Let $u$, $M_0$,
$G'$ and $N_1$ be the integers of Definition~\ref{hf-good-orbit} formed with $r$, $K=r-1$, $\delta$
and $h$ in place of $r_\varepsilon$, $K_\varepsilon$, $\delta_1$ and $h_\varepsilon$. Let
$\lambda_0(s)\in\mathbf Z^{k}$ be dominant weights and $a_T,b_T\colon\mathbf Z/n'\to\mathbf Z$,
$c_{pp},c_0,D_p,c_*\geq0$ integers such that:
\begin{enumerate}
\item $T$ satisfies the three conditions of a tiled set of positions for $(x,\cdot,L_2)$ with these
  integers, $2L_1\leq L_2<u$, $3G'\leq n'$ and $N_1^2\leq r|T|$;
\item $\lambda_0(s)_1=a_T(s)$, $\lambda_0(s)_k=-b_T(s)$ and $|\lambda_0(s)_i|\leq c_{pp}$ for
  $1<i<k$ when $s\in T$, $|\lambda_0(s)_i|\leq c_{pp}$ for all $i$ when $s\notin T$, and
  $a+c_{pp}\leq a_T(s)\leq2a-c_{pp}$, $b_T(s)\geq0$ and $b_T(s)+2c_{pp}\leq c_0$ for $s\in T$;
\item $r\leq Y$, $CN_1^2\log a\leq k$, $N_1+2r\leq k-r$, $256k^2|T|Y^2\leq p^{L_1-1}$,
  $C|T|k^2\log_pa\leq D_p$, $8(D_p+1)\leq c_*$ and
  $Ck^3r^2|T|\,p^{5G'}\big(c_*+r|T|c_0+2rn'c_{pp}\big)Y\leq a$.
\end{enumerate}
Let $V$, the lattices $H(L)$, $R(e)$ and $R'(e)$ and the edge pairing be those of
Definition~\ref{pd-projection-property} for $r$ and $\lambda_0$. Put $\omega=r|T|c_0+2rn'c_{pp}$ and,
for $C'\geq0$,
\[
 t_0(C')=C'|T|\Big(\frac{CN_1^2\log a\cdot a}{pk}+ap^{-L_1}+kp^{G'}\Big).
\]
\end{situation}
\leanref{Formalism/SharpRestriction/GrowthBounds/ScaleEdges/ManyEdges.lean}{Formalism.MainNotions.Pf14.Hyps}
\leanref{Formalism/SharpRestriction/GrowthBounds/ScaleEdges/ManyEdges.lean}{Formalism.MainNotions.Pf14.t0}
\leanref{Formalism/SharpRestriction/GrowthBounds/ScaleEdges/ManyEdges.lean}{Formalism.MainNotions.Pf14.om}

\begin{lemma}\label{pd-edge-numerics}
In Situation~\ref{pd-edge-situation}, let $C_t,C_D,C_r>0$ with $C\geq2$, $C_t\leq C$, $C_D\leq C$ and
$C_r+1\leq C$. Then:
\begin{enumerate}
\item $t_0(C_t)+|T|c_0+\omega\leq Cr|T|\big(CN_1^2\log a\cdot a/(pk)+ap^{-L_1}+c_0+kp^{G'}\big)
  +2rn'c_{pp}$;
\item $C_D|T|k^2\log_pa\leq D_p$;
\item $kc_0\leq a$ and $1-C|T|r^2k/Y\leq(1-C_r|T|r^2k/Y)(1-kc_0/a)^{|T|}$.
\end{enumerate}
\end{lemma}
\leanref{Formalism/SharpRestriction/GrowthBounds/ScaleEdges/ManyEdges.lean}{Formalism.MainNotions.Pf14.Num}
\leanref{Formalism/SharpRestriction/GrowthBounds/ScaleEdges/NumericProductInequality.lean}{Formalism.MainNotions.Pf14.num}
\leanref{Formalism/SharpRestriction/GrowthBounds/ScaleEdges/RankMulConductorDivLeOne.lean}{Formalism.MainNotions.Pf14.rank_mul_conductor_div_le_one}

\begin{proof}
Since $1\leq N_1$ and $N_1^2\leq r|T|$, we have $|T|\geq1$. As $C\geq2$, $r\geq2$, $Y\geq r$ and
$c_*\geq8$, the last inequality of (3) gives $a\geq1$ and $a\geq c_0r^2kY\geq kc_0$.

(1) Let $X=CN_1^2\log a\cdot a/(pk)+ap^{-L_1}+kp^{G'}\geq0$. The left side is
$C_t|T|X+(r+1)|T|c_0+2rn'c_{pp}$, and $C_t\leq C\leq Cr$ and $r+1\leq2r\leq Cr$.

(2) This follows from $C|T|k^2\log_pa\leq D_p$, since $C_D\leq C$ and $\log_pa\geq0$.

(3) The first claim was shown above. Put $A_r=C_r|T|r^2k/Y$ and $B=|T|kc_0/a$. By Bernoulli's
inequality $(1-kc_0/a)^{|T|}\geq1-B$, and this power is at most $1$. If $A_r\leq1$ the right side is
at least $(1-A_r)(1-B)\geq1-A_r-B$; if $A_r>1$ it is at least $1-A_r$. Finally
$B\leq|T|r^2k/Y\leq(C-C_r)|T|r^2k/Y$, because $c_0Y\leq r^2a$ and $C-C_r\geq1$.
\end{proof}

\subsection{Run systems and the count of symbols}\label{subsection-pd-symbols-count}

\begin{definition}\label{pd-symbols-run-system}
Let $n'\geq1$ and $r\geq1$ be integers and $T\subset\mathbf Z/n'\mathbf Z$; for $v\in\mathbf Z$ write
$\bar v$ for its class modulo $n'$. A \emph{run system} for $T$ of length $r$ consists of an integer
$J\geq1$ and integers $f(j,i)$ for $j\in\mathbf Z$ and $0\leq i<r$. For $j\in\mathbf Z$ put
\[
 \Lambda_j=\{v\in\mathbf Z:f(j-1,0)\leq v\leq f(j,r-1),\ \bar v\in T\},\qquad
 \Lambda'_j=\{v\in\mathbf Z:f(j-1,r-1)<v<f(j,0),\ \bar v\in T\}.
\]
An element $v\in\Lambda_j$ is \emph{new} if $v>f(j-1,r-1)$. For $v\in\mathbf Z$ put
$\operatorname{dist}_j(v)=\max(f(j,r-1)-v,0)$, and let $\operatorname{old}_j(v)$ be the $r$-th
largest even element of $\Lambda_j$ smaller than $v$, the smallest one if there are fewer than $r$,
and $0$ if there is none. For an integer $L\geq0$, $v$ is \emph{$L$-isolated in $\Lambda_j$} if
$v-y>L$ for the largest $y\in\Lambda_j$ with $y<v$, when there is one. An element of $\Lambda'_j$ is
\emph{extreme} if it is among the $r$ smallest or among the $r$ largest elements of $\Lambda'_j$. Let
$L_1,L_2,N_1,G'\geq0$ be integers. The run system is \emph{valid} for $(L_1,N_1)$ if, for all $j$
and $0\leq i<r$:
\begin{enumerate}
\item $f(j+J,i)=f(j,i)+n'$;
\item $f(j,0)<f(j,1)<\cdots<f(j,r-1)<f(j+1,0)$;
\item $f(j,i)$ is even and $\overline{f(j,i)}\in T$, and if $i+1<r$ no $v$ with
  $f(j,i)<v<f(j,i+1)$ has $\bar v\in T$;
\item $\Lambda'_j$ has at least $r$ elements, every extreme element of $\Lambda'_j$ is even and
  $L_1$-isolated in $\Lambda_j$, and $f(j,0)$ is $L_1$-isolated in $\Lambda_j$;
\item $\Lambda_j$ has at most $N_1$ elements;
\item $v-\operatorname{old}_j(v)>L_1$ for every new $v\in\Lambda_j$.
\end{enumerate}
It is \emph{spaced} for $(L_1,L_2,N_1,G')$ if it is valid for $(L_1,N_1)$,
$f(j,r-1)-f(j-1,0)<G'$ for every $j$, and every extreme element of $\Lambda'_j$ and $f(j,0)$ are
$L_2$-isolated in $\Lambda_j$ for every $j$.
\end{definition}
\leanref{Formalism/SharpRestriction/GrowthBounds/SegRing.lean}{Formalism.MainNotions.G.Frames}
\leanref{Formalism/SharpRestriction/GrowthBounds/SegRing.lean}{Formalism.MainNotions.G.Frames.Valid}
\leanref{Formalism/SharpRestriction/GrowthBounds/ScaleEdges/FramesOf.lean}{Formalism.MainNotions.Pf14.FramesOK}

\begin{lemma}\label{pd-symbols-run-system-exists}
Let $\delta,h,x$ be real numbers, $n'\geq2$ an even integer, $r\geq2$, $L_1\geq1$ and $L_2\geq0$
integers, and $T\subset\mathbf Z/n'\mathbf Z$. Let $u$, $M_0$, $G'$ and $N_1$ be the integers of
Definition~\ref{hf-good-orbit} formed with $r$, $K=r-1$, $\delta$ and $h$ in place of
$r_\varepsilon$, $K_\varepsilon$, $\delta_1$ and $h_\varepsilon$. Suppose that $T$ satisfies the
three conditions of a tiled set of positions for $(x,\cdot,L_2)$ with these integers, that
$2L_1\leq L_2<u$ and that $3G'\leq n'$. Then there is a run system for $T$ of length $r$ that is
spaced for $(L_1,L_2,N_1,G')$ (Definition~\ref{pd-symbols-run-system}).
\end{lemma}
\leanref{Formalism/SharpRestriction/GrowthBounds/ScaleEdges/FramesOf.lean}{Formalism.MainNotions.Pf14.frames_of}

\begin{proof}
Let $P\subset\{0,\dots,n'-1\}$ be the set of representatives of $T$, so that $\bar v\in T$ if and only
if the residue of $v$ modulo $n'$ lies in $P$; in this proof a \emph{point} is an integer $v$ with
$\bar v\in T$. Put $\ell=(6r+1)u$ and $S=\lfloor n'/\ell\rfloor$. Since $0\leq L_2<u$, $u\geq1$. As
$G'=(M_0+5)\ell$ and $3G'\leq n'$, $(M_0+2)\ell\leq n'$, so $S\geq M_0+2$; and $L_2<u\leq G'\leq n'$.
For $b\in\mathbf Z$ and $1\leq c\leq6r+1$ let $I_c(b)$ be the set of integers $v$ with
$b+(c-1)u\leq v<b+cu$. Say that $b$ \emph{starts a marked block} if for odd $c$ no element of $I_c(b)$
is a point, and for even $c$ exactly one element of $I_c(b)$ is a point and it is even; write
$\tau_b(c)$ for it.

\emph{Step 1: blocks.} Let $0\leq s<S$ be a marked tile (Definition~\ref{hf-good-orbit}) and
$q\in\mathbf Z$. Then $b=n'q+s\ell$ starts a marked block. Indeed $(s+1)\ell\leq S\ell\leq n'$, so for
$v\in I_c(b)$ the integer $v-n'q$ lies in the cell $I_{s,c}\subset\{0,\dots,n'-1\}$, and $v$ is a
point if and only if $v-n'q\in P$. The conditions on $b$ follow from those on the tile, and the point
of $I_c(b)$ for even $c$ is $n'q+y$ with $y\in P$ even, which is even as $n'$ is even.

\emph{Step 2: three facts.} Let $b$ start a marked block, and let $c,c'$ be even.
\begin{enumerate}
\item[(i)] If $c<c'$, then $\tau_b(c)<\tau_b(c')$, since $\tau_b(c)<b+cu\leq b+(c'-1)u\leq\tau_b(c')$.
\item[(ii)] Every point $v$ with $b\leq v<b+\ell$ is $\tau_b(c'')$ for some even $c''$: it lies in some
  $I_{c''}(b)$, and $c''$ is even because the odd cells contain no points.
\item[(iii)] If $v$ is a point with $v<\tau_b(c)$, then $v+u<\tau_b(c)$. Indeed $v\notin I_c(b)$, as
  $\tau_b(c)$ is the only point there, and $v\notin I_{c-1}(b)$, as $c-1$ is odd; so
  $v<b+(c-2)u$ and $v+u<b+(c-1)u\leq\tau_b(c)$.
\end{enumerate}

\emph{Step 3: the blocks of the run system.} Let $\mu_0<\mu_1<\dots<\mu_{J-1}$ be the marked tiles in
$\{0,\dots,S-1\}$. Condition~(2) of a tiled set, applied to $0$, gives one, so $J\geq1$. For
$j\in\mathbf Z$ write $j=Jq+i$ with $0\leq i<J$ and put $B_j=n'q+\mu_i\ell$. Then $B_{j+J}=B_j+n'$, and
each $B_j$ starts a marked block by Step~1. We show that $B_j+\ell\leq B_{j+1}$ and
$B_{j+1}-B_j<(M_0+2)\ell$. Condition~(2), applied to the residue of $\mu_i+1$ modulo $S$, gives
$0\leq i'\leq M_0$ such that the residue $\mu'$ of $\mu_i+1+i'$ modulo $S$ is a marked tile. As
$\mu_i+1+i'\leq S+M_0<2S$, $\mu'$ is $\mu_i+1+i'$ or $\mu_i+1+i'-S$. If $i+1<J$, then
$B_{j+1}-B_j=(\mu_{i+1}-\mu_i)\ell\geq\ell$, and $\mu_{i+1}\leq\mu_i+1+i'$: in the first case because
$\mu'>\mu_i$ is a marked tile, in the second because $\mu_{i+1}<S\leq\mu_i+1+i'$. So
$B_{j+1}-B_j\leq(M_0+1)\ell$. If $i+1=J$, then $B_{j+1}-B_j=n'+(\mu_0-\mu_{J-1})\ell$, which is at
least $\ell$ because $(\mu_{J-1}+1)\ell\leq S\ell\leq n'$. In this case $\mu'>\mu_{J-1}$ is impossible,
so $\mu'=\mu_{J-1}+1+i'-S\geq\mu_0$, and $B_{j+1}-B_j\leq n'+(M_0+1-S)\ell<(M_0+2)\ell$, because
$n'<(S+1)\ell$.

\emph{Step 4: the run system.} Put $f(j,i)=\tau_{B_j}(2r+2+2i)$ for $j\in\mathbf Z$ and $0\leq i<r$. We
check the conditions of Definition~\ref{pd-symbols-run-system}.

(1) Adding $n'$ to a point gives a point, so the point of $I_c(B_j+n')$ is $\tau_{B_j}(c)+n'$, and
$f(j+J,i)=f(j,i)+n'$.

(2) By (i), $f(j,0)<\dots<f(j,r-1)$. Also $f(j,r-1)<B_j+4ru<B_j+\ell\leq B_{j+1}<f(j+1,0)$.

(3) Each $f(j,i)$ is an even point. A point $v$ with $f(j,i)<v<f(j,i+1)$ satisfies
$B_j\leq v<B_j+\ell$, so by (ii) it is $\tau_{B_j}(c'')$ with $c''$ even, and by (i)
$2r+2+2i<c''<2r+4+2i$, which is impossible.

Spacing, first part: $f(j,r-1)<B_j+\ell$ and $f(j-1,0)\geq B_{j-1}$, so by Step~3
$f(j,r-1)-f(j-1,0)<(M_0+3)\ell\leq G'$.

(5) The set $\Lambda_j$ lies in the interval $[f(j-1,0),f(j-1,0)+G')$. As $G'\leq n'$, distinct
elements of $\Lambda_j$ have distinct residues, and these lie in $P\cap W(i_0,G')$, where $i_0$ is
the residue of $f(j-1,0)$. By condition~(3) of a tiled set, $|\Lambda_j|\leq N_1$.

(4) For $0\leq l<r$ the points $\tau_{B_{j-1}}(4r+2+2l)$ lie in $\Lambda'_j$: they exceed
$\tau_{B_{j-1}}(4r)=f(j-1,r-1)$, and they are less than $B_{j-1}+\ell\leq B_j<f(j,0)$. They are distinct
by (i), so $|\Lambda'_j|\geq r$. In the same way the points $\tau_{B_j}(2+2l)$, $0\leq l<r$, lie in
$\Lambda'_j$, since $f(j-1,r-1)<B_{j-1}+\ell\leq B_j$ and $\tau_{B_j}(2+2l)<\tau_{B_j}(2r+2)=f(j,0)$.
Let $x$ be an extreme element of $\Lambda'_j$. If $x$ is among the $r$ smallest, then
$x<B_{j-1}+\ell$, for otherwise the $r$ points $\tau_{B_{j-1}}(4r+2+2l)$ of $\Lambda'_j$ are all smaller
than $x$; and $x>f(j-1,r-1)\geq B_{j-1}$; so $x=\tau_{B_{j-1}}(c)$ for an even $c$ by (ii). If $x$ is
among the $r$ largest, then $x\geq B_j$, for otherwise the $r$ points $\tau_{B_j}(2+2l)$ of
$\Lambda'_j$ are all larger than $x$; and $x<f(j,0)<B_j+\ell$; so $x=\tau_{B_j}(c)$ for an even $c$. In
both cases $x$ is even, and by (iii) $x-y>u>L_2\geq L_1$ for every point $y<x$; so $x$ is
$L_2$-isolated and $L_1$-isolated in $\Lambda_j$. By (iii) the same holds for
$f(j,0)=\tau_{B_j}(2r+2)$. This gives (4) and the second part of the spacing.

(6) Let $v\in\Lambda_j$ be new, so $v>f(j-1,r-1)$. The even points $f(j-1,0),\dots,f(j-1,r-1)$ lie in
$\Lambda_j$, as $f(j-1,r-1)<f(j,0)\leq f(j,r-1)$, and are smaller than $v$. So $\operatorname{old}_j(v)$
is the $r$-th largest even element of $\Lambda_j$ smaller than $v$, and the $r$ largest such elements
lie in $[\operatorname{old}_j(v),v)$. Suppose $v-\operatorname{old}_j(v)\leq L_1$. Since $L_1\geq1$ and
$2L_1\leq L_2$, $L_1<L_2$, so these $r$ points lie in $[\operatorname{old}_j(v),\operatorname{old}_j(v)+L_2)$.
As $L_2\leq n'$, their residues are $r$ distinct elements of $P\cap W(i_1,L_2)$, where $i_1$ is the
residue of $\operatorname{old}_j(v)$. This contradicts condition~(1) of a tiled set, with
$K=r-1$. Hence $v-\operatorname{old}_j(v)>L_1$.
\end{proof}

\begin{definition}\label{pd-symbols-module}
Let $p$ be a prime, $\mathcal R$ a commutative ring, $r,K\geq1$ integers, $k=r+K$, $n'\geq2$ an
even integer, $T\subset\mathbf Z/n'\mathbf Z$, and $(J,f)$ a run system for $T$ of length $r$, with
the notation of Definition~\ref{pd-symbols-run-system}.
\begin{enumerate}
\item \emph{Weights.} Let $n\geq1$, $S\subset\mathbf Z$, and $A$ a polynomial over a commutative
  ring $R_0$ in variables $z_{v,c}$, $v\in\mathbf Z$, $1\leq c\leq n$. Let $\mathfrak g$ be the
  $n\times n$ matrix of indeterminates over $R_0$, $\eta$ a further indeterminate, and for
  $1\leq i\leq n$ let $\tau_i(\eta)=\mathfrak gD_i(\eta)\mathfrak g^{-1}$, where $D_i(\eta)$ is the
  diagonal matrix with $\eta$ in place $i$ and $1$ elsewhere. Let $A_{S,i}$ be obtained from $A$ by
  replacing, for every $v\in S$, the row $z_v=(z_{v,1},\dots,z_{v,n})$ by $z_v\tau_i(\eta)$ if $v$
  is even and by $z_v\,{}^t\tau_i(\eta^{-1})$ if $v$ is odd. Then $\operatorname{val}_S(A)$ is the
  least exponent of $\eta$ that occurs in $A_{S,n}$, and $\lambda_S(A)$ the largest that occurs in
  $A_{S,1}$, with $\operatorname{val}_S(0)=+\infty$ and $\lambda_S(0)=-\infty$.
\item \emph{Segment modules.} Let $\mathcal R[z]$ be the polynomial ring over $\mathcal R$ in
  $z_{v,c}$, $v\in\mathbf Z$, $1\leq c\leq k$; the $z_{v,c}$ with $c\leq r$ are the \emph{plane}
  variables and the others the \emph{normal} ones. For a monomial $\mu$ and $v\in\mathbf Z$ let
  $\operatorname{ord}_v(\mu)$ be the degree of $\mu$ in the normal variables $z_{v,c}$ if $v$ is
  even, and in the plane variables $z_{v,c}$ if $v$ is odd. Let $j\in\mathbf Z$,
  $m\colon\mathbf Z\to\mathbf Z$ and $\theta,o\colon\mathbf Z\to\mathbf R$. Consider the products
  $AG$ of a polynomial $A$ in the plane variables $z_{v,c}$ with $v\in\Lambda_j$ and a polynomial $G$
  in the normal variables $z_{v,c}$ with $v\in\Lambda_j$ new, for which there is
  $M\colon\mathbf Z\to\mathbf Z$ such that for every new $v\in\Lambda_j$:
  $\lambda_{\{y\leq v\}}(G)\leq M(v)$ and $M(v)+m(v)\leq\operatorname{val}_{\{y<v\}}(A)$ if $v$ is
  even; $\lambda_{\{y<v\}}(G)\leq M(v)$ and $M(v)+m(v)\leq\operatorname{val}_{\{y\leq v\}}(A)$ if
  $v$ is odd; and $\theta(v)\leq\operatorname{val}_{\{y<v\}}(A)$. Here the weights of $A$ are
  formed with $n=r$, and those of $G$ with $n=K$, numbering the variables $z_{v,r+1},\dots,z_{v,k}$
  as $1,\dots,K$. The \emph{segment module} $X_j(m,\theta,o)$ consists of the elements of the
  $\mathcal R$-span of these products all of whose monomials $\mu$ satisfy
  $\operatorname{ord}_v(\mu)\leq o(v)$ for every new $v\in\Lambda_j$.
\item \emph{Symbols.} Let $\mathcal R[X]$ be the polynomial ring over $\mathcal R$ in $X_{s,c}$,
  $s\in\mathbf Z/n'\mathbf Z$, $1\leq c\leq k$, and for $j\in\mathbf Z$ let
  $\phi_j\colon\mathcal R[z]\to\mathcal R[X]$ be the $\mathcal R$-algebra map with
  $\phi_j(z_{v,c})=X_{\bar v,c}^{\,p^{\operatorname{dist}_j(v)}}$. A monomial $X^\ell$ is a
  \emph{label} if $\ell(\bar v,c)<p^{\operatorname{dist}_j(v)}$ for all $j$, all new $v\in\Lambda_j$
  and all $c$. A polynomial has \emph{multidegree} $d\colon\mathbf Z/n'\mathbf Z\to\mathbf Z_{\geq0}$
  if each of its monomials has degree $d(s)$ in $X_{s,1},\dots,X_{s,k}$ for every $s$. Given
  segment modules $X_j\subset\mathcal R[z]$ for $j\in\mathbf Z$ and such a $d$, a \emph{symbol} is a
  polynomial of multidegree $d$ of the form $\phi_0(\xi_0)\phi_1(\xi_1)\cdots\phi_{J-1}(\xi_{J-1})X^\ell$
  with $\xi_j\in X_j$ and $X^\ell$ a label. The set of symbols is $\mathrm{Gen}(X_\bullet,d)$, and
  $N(X_\bullet,d)$ is its $\mathcal R$-span. The \emph{normal degree} of a monomial in the $X_{s,c}$ is
  its degree in the $X_{s,c}$ with $s$ even and $c>r$ together with those with $s$ odd and $c\leq r$.
\item \emph{Symbol spaces.} Let $C_{La},a,Y$ be real numbers with $a>0$, $L_1,N_s,m_*\geq0$
  integers and $\deg\colon\mathbf Z/n'\mathbf Z\to\mathbf Z_{\geq0}$. Put $L_a=C_{La}N_s^2\log a$
  and, for $j,v\in\mathbf Z$,
  $B_j(v)=\deg\big(\overline{\operatorname{old}_j(v)}\big)/p^{\operatorname{dist}_j(\operatorname{old}_j(v))}$,
  \[
   m_j(v)=m_*+\max\Big(\Big\lfloor\frac{B_j(v)}{k^2Y}\Big\rfloor,0\Big),\qquad
   \theta_j(v)=\frac{B_j(v)}2,\qquad
   o_j(v)=\frac{L_aa}{pk\,p^{\operatorname{dist}_j(v)}}+\epsilon_v\,a\,p^{-L_1-\operatorname{dist}_j(v)},
  \]
  with $\epsilon_v=10$ for $v$ even and $\epsilon_v=1$ for $v$ odd. The \emph{symbol spaces} for
  $(C_{La},\deg,L_1,N_s,m_*,a,Y)$ are the segment modules $X_j=X_j(m_j,\theta_j,o_j)$.
\end{enumerate}
\end{definition}
\leanref{Formalism/SharpRestriction/GrowthBounds/ObjPol.lean}{Formalism.MainNotions.G.valS}
\leanref{Formalism/SharpRestriction/GrowthBounds/ObjPol.lean}{Formalism.MainNotions.G.lam1S}
\leanref{Formalism/SharpRestriction/GrowthBounds/SegRing.lean}{Formalism.MainNotions.G.XX'}
\leanref{Formalism/SharpRestriction/GrowthBounds/SegRing.lean}{Formalism.MainNotions.G.phSeg}
\leanref{Formalism/SharpRestriction/GrowthBounds/SegRing.lean}{Formalism.MainNotions.G.IsLabel}
\leanref{Formalism/SharpRestriction/GrowthBounds/SegRing.lean}{Formalism.MainNotions.G.Gen}
\leanref{Formalism/SharpRestriction/GrowthBounds/SegRing.lean}{Formalism.MainNotions.G.vsDeg}
\leanref{Formalism/SharpRestriction/GrowthBounds/ScaleEdges/SymbolRealisableBound.lean}{Formalism.MainNotions.Pf14.symbolSpaces}

\begin{lemma}\label{pd-symbols-select}
Let $\mathcal R$ be a discrete valuation ring with residue field $\kappa$, let $\mathcal R[Z]$ be a
polynomial ring over $\mathcal R$ in any set of variables, and for $P\in\mathcal R[Z]$ let
$\bar P\in\kappa[Z]$ be its reduction. Let $s$ be a finite set of monomials and
$\mathcal G\subset\mathcal R[Z]$ a set of polynomials all of whose monomials lie in $s$. Then there is
an $\mathcal R$-submodule $S$ of the $\mathcal R$-span of $\mathcal G$ which is a direct summand of
$\mathcal R[Z]$, has a basis consisting of elements of $\mathcal G$, and has rank equal to the
dimension of the $\kappa$-span of the $\bar P$, $P\in\mathcal G$.
\end{lemma}
\leanref{Formalism/SharpRestriction/GrowthBounds/SegmentDegree/NonemptyDvrSettingOfStandard.lean}{Formalism.MainNotions.SelectionLemmas.Pf.RestSel}
\leanref{Formalism/SharpRestriction/GrowthBounds/SegmentDegree/SelIsComplRestrictSupport.lean}{Formalism.MainNotions.SelectionLemmas.Pf.restSel}

\begin{proof}
Let $\varpi$ be a uniformizer of $\mathcal R$. The $\bar P$, $P\in\mathcal G$, lie in the
finite-dimensional space of polynomials over $\kappa$ supported on $s$. Choose
$g_1,\dots,g_n\in\mathcal G$ such that $\bar g_1,\dots,\bar g_n$ form a basis of the $\kappa$-span of
the $\bar P$, and let $S$ be the $\mathcal R$-span of the $g_i$.

First, if $\sum_ic_ig_i=\varpi x$ with $c_i\in\mathcal R$ and $x\in\mathcal R[Z]$, then
$c_i=\varpi c'_i$ with $c'_i\in\mathcal R$ and $\sum_ic'_ig_i=x$. Indeed, reducing gives
$\sum_i\bar c_i\bar g_i=0$, so every $\bar c_i$ vanishes; then $\varpi(\sum_ic'_ig_i-x)=0$, and
$\mathcal R[Z]$ has no $\mathcal R$-torsion.

The $g_i$ are linearly independent over $\mathcal R$: if $\sum_ic_ig_i=0$, the first step applied
repeatedly shows that every $c_i$ is divisible by every power of $\varpi$, so $c_i=0$. Hence the $g_i$
form a basis of $S$, of rank $n$.

The module $S$ is saturated: if $a\in\mathcal R$ is nonzero and $ax\in S$, then $x\in S$. Write
$a=\epsilon\varpi^m$ with $\epsilon$ a unit; then $\varpi^mx\in S$, and the first step, applied $m$
times, gives $x\in S$.

The polynomials supported on $s$ form a free $\mathcal R$-module $N$ of finite rank, and $N$ is a
direct summand of $\mathcal R[Z]$, with complement the polynomials supported off $s$. As $S\subset N$
is saturated, $N/S$ is a finitely generated torsion-free module over the principal ideal domain
$\mathcal R$, hence free, so $S$ is a direct summand of $N$ and therefore of $\mathcal R[Z]$.
\end{proof}

\begin{lemma}\label{pd-symbols-normal-degree}
Let $p\geq13$ be a prime, $\mathcal R$ an integral domain, $n'\geq2$ an even integer, $r\geq2$ and
$K\geq1$ integers, $k=r+K$, $T\subset\mathbf Z/n'\mathbf Z$, $L_1,N_1\geq0$ integers, and $(J,f)$ a run
system for $T$ of length $r$ valid for $(L_1,N_1)$ (Definition~\ref{pd-symbols-run-system}). Let
$C_{La}\geq0$ and $C_k$ be real numbers, and let $\deg$, $E'$, $N_s$, $E_0$, $m_*$, $a$, $Y$ and $q$
satisfy the hypotheses (1) to (4) of Theorem~\ref{pd-symbols-count} with this $C_k$. Let $X_j$ be the
symbol spaces for $(C_{La},\deg,L_1,N_s,m_*,a,Y)$ over $\mathcal R$, $L_a=C_{La}N_s^2\log a$
(Definition~\ref{pd-symbols-module}), $d(t)=\max(\deg(t)-E'(t),0)$ for $t\in T$ and $d(s)=0$ for
$s\notin T$. Then every monomial of every element of $N(X_\bullet,d)$ has normal degree at most
\[
 |T|\Big(\frac{L_aa}{pk}+10ap^{-L_1}\Big)+|T|kq\leq10|T|\Big(\frac{L_aa}{pk}+ap^{-L_1}+kq\Big).
\]
\end{lemma}
\leanref{Formalism/SharpRestriction/GrowthBounds/SegmentDegree/NonemptyDvrSettingOfStandard.lean}{Formalism.MainNotions.SelectionLemmas.Pf.RestDeg}
\leanref{Formalism/SharpRestriction/GrowthBounds/SegmentDegree/StrictMonoIntExists.lean}{Formalism.MainNotions.SelectionLemmas.Pf.restDeg}
\leanref{Formalism/SharpRestriction/GrowthBounds/SegmentDegree/StrictMonoIntExists.lean}{Formalism.MainNotions.SelectionLemmas.Pf.gen_bound}

\begin{proof}
We use the notation of Definition~\ref{pd-symbols-run-system}. The element $f(0,r-1)$ lies in
$\Lambda_0$, since $f(-1,0)\leq f(-1,r-1)<f(0,0)\leq f(0,r-1)$, and
$\operatorname{dist}_0(f(0,r-1))=0$; so $N_s\geq|\Lambda_0|\geq1$ and $q\geq1$ by (2) and (3). Each
factor of $32pYk^2N_sq^2(E_0+m_*+1)$ is at least $1$, so $a\geq1$ by (4), and $L_a\geq0$. Put
$B=L_aa/(pk)+10ap^{-L_1}$.

\emph{Old elements are even.} Let $v\in\Lambda_j$ not be new, so $f(j-1,0)\leq v\leq f(j-1,r-1)$. By
condition~(3) of validity no point of $T$ lies strictly between consecutive $f(j-1,i)$, so
$v=f(j-1,i)$ for some $i$, and $v$ is even.

\emph{Segments.} Let $\xi\in X_j$, and let $\mu$ be a monomial of $\xi$. Since $\xi$ is a combination
of products $AG$ as in Definition~\ref{pd-symbols-module}(2), $\mu$ involves only variables $z_{v,c}$
with $v\in\Lambda_j$, and with $c\leq r$ unless $v$ is new. The map $\phi_j$ sends $\mu$ to a
monomial, and every monomial of $\phi_j(\xi)$ is such an image. As $n'$ is even, $\bar v$ is even
exactly when $v$ is. A variable $z_{v,c}$ with $v$ old has $v$ even and $c\leq r$, so its image is not
normal. For new $v$ the normal variables among the $z_{v,c}$ are those counted by
$\operatorname{ord}_v$. Hence, by the condition on the monomials of $X_j$,
\[
 \text{normal degree of }\phi_j(\mu)=\sum_{v\text{ new}}p^{\operatorname{dist}_j(v)}\operatorname{ord}_v(\mu)
 \leq\sum_{v\text{ new}}p^{\operatorname{dist}_j(v)}o_j(v)\leq n_jB,
\]
where $n_j$ is the number of new elements of $\Lambda_j$, because
$p^{\operatorname{dist}_j(v)}o_j(v)=L_aa/(pk)+\epsilon_vap^{-L_1}\leq B$.

\emph{Counting new elements.} The new elements of $\Lambda_j$ are the $v$ with $\bar v\in T$ and
$f(j-1,r-1)<v\leq f(j,r-1)$. For $0\leq j<J$ these intervals are disjoint, and their union is
$(f(-1,r-1),f(-1,r-1)+n']$, on which $v\mapsto\bar v$ is injective. So $\sum_{j<J}n_j\leq|T|$.

\emph{Labels.} Let $X^\ell$ be a label occurring in a symbol $P$ of multidegree $d$, and let $X^\mu$ be
a monomial of $P$; it is the product of $X^\ell$ and a monomial of
$\prod_{j<J}\phi_j(\xi_j)$. For $s\notin T$, $d(s)=0$, so $\mu$ and $\ell$ vanish at $(s,c)$. For
$t\in T$, the intervals $(f(j-1,r-1),f(j,r-1)]$, $j\in\mathbf Z$, cover $\mathbf Z$, as $f(j,r-1)$
increases with $j$ and $f(j+J,r-1)=f(j,r-1)+n'$; so $t=\bar v$ for a new element $v$ of some
$\Lambda_j$, and $\ell(t,c)<p^{\operatorname{dist}_j(v)}\leq q$ by (3). Hence the normal degree of
$X^\ell$ is at most $|T|kq$.

The normal degree of a monomial of a product is the sum of those of monomials of the factors, so every
monomial of a symbol has normal degree at most $\sum_{j<J}n_jB+|T|kq\leq|T|B+|T|kq$. A monomial of an
element of $N(X_\bullet,d)$ is a monomial of one of the symbols it is a combination of. The last
inequality of the statement holds because $L_aa/(pk)$, $ap^{-L_1}$ and $kq$ are nonnegative.
\end{proof}

\begin{definition}\label{pd-symbols-residue-coarse}
Let $p$ be a prime, $\mathcal R$ a commutative ring, $r,K\geq1$ integers, $k=r+K$, $n'\geq2$ an even
integer, $T\subset\mathbf Z/n'\mathbf Z$ and $(J,f)$ a run system for $T$ of length $r$, with the
notation of Definitions~\ref{pd-symbols-run-system} and~\ref{pd-symbols-module}. Let $C_{La}$, $a>0$
and $Y$ be real numbers, $L_1,N_s,m_*\geq0$ integers, $\deg,E'\colon\mathbf Z/n'\mathbf Z\to
\mathbf Z_{\geq0}$, $L_a$, $B_j(v)$ and the symbol spaces $X_j$ as in
Definition~\ref{pd-symbols-module}(4), and $d(t)=\max(\deg(t)-E'(t),0)$ for $t\in T$, $d(s)=0$ for
$s\notin T$.
\begin{enumerate}
\item \emph{The discriminant.} For $j\in\mathbf Z$ and a new $v\in\Lambda_j$ let
  $\rho_{j,v}(1)<\dots<\rho_{j,v}(r)$ be the $r$ largest even elements of $\Lambda_j$ smaller than $v$
  (for a valid run system there are at least $r$ of them, as the $f(j-1,i)$, $i<r$, are such
  elements), and $\sigma_j(v)=\max(\lfloor B_j(v)/(k^2Y)\rfloor,0)$. The \emph{discriminant} is
  \[
   D=\prod_{j=0}^{J-1}\phi_j\Big(\prod_{v}\det\big(z_{\rho_{j,v}(l),b}\big)_{1\leq l,b\leq r}^{\,\sigma_j(v)}\Big)
   \in\mathcal R[X],
  \]
  the inner product over the new $v\in\Lambda_j$. Let $E^D(t)$ be the sum, over $0\leq j<J$, the new
  $v\in\Lambda_j$ and the $l\leq r$ with $\overline{\rho_{j,v}(l)}=t$, of
  $\sigma_j(v)p^{\operatorname{dist}_j(\rho_{j,v}(l))}$; it is the degree of $D$ in
  $X_{t,1},\dots,X_{t,k}$, as $\phi_j$ sends $z_{y,b}$ to $X_{\bar y,b}^{\,p^{\operatorname{dist}_j(y)}}$
  and each determinant has degree $1$ in the variables of each of its rows. Put
  $d^D(t)=\max(d(t)-E^D(t),0)$.
\item \emph{Coarse twists.} For $j\in\mathbf Z$ let $R_j\geq0$ be the largest integer with
  $p^{R_j}\leq\max\big(1,\big\lfloor a/(8Yk^2N_sp^{\operatorname{dist}_j(f(j-1,0))})\big\rfloor\big)$,
  put $Q_j=p^{\max(0,\,f(j+1,r-1)-f(j,r-1)+R_{j+1}-R_j)}$ and $m_0=kN_s+1$, and let
  $\psi_j\colon\mathcal R[z]\to\mathcal R[X]$ be the $\mathcal R$-algebra map with
  $\psi_j(z_{v,c})=X_{\bar v,c}^{\,p^{\operatorname{dist}_j(v)+R_j}}$. A monomial $X^\ell$ is a
  \emph{coarse label} if $\ell(\bar v,c)<p^{\operatorname{dist}_j(v)+R_j}$ for all $j$, all new
  $v\in\Lambda_j$ and all $c$.
\item \emph{Coarse segment modules.} For $j,y\in\mathbf Z$ put
  $W_j(y)=\big\lfloor d^D(\bar y)/p^{\operatorname{dist}_j(y)+R_j}\big\rfloor-(k+r+1)$. For
  $b\colon\mathbf Z\to\mathbf Z_{\geq0}$ and $v\in\mathbf Z$ put
  $\theta'_{j,b}(v)=\frac23\big(W_j(\operatorname{old}_j(v))-m_0\big)$ if $\operatorname{old}_j(v)$ is
  not new and $\theta'_{j,b}(v)=\frac23b(\operatorname{old}_j(v))$ if it is, and
  \[
   o'_j(v)=\frac{L_aa}{pk\,p^{\operatorname{dist}_j(v)+R_j}}+\epsilon'_v\,a\,
   p^{-\operatorname{dist}_j(\operatorname{old}_j(v))-R_j},
  \]
  with $\epsilon'_v=4$ for $v$ even and $\epsilon'_v=0$ for $v$ odd. The module $M_j(b)$ consists of
  the elements of the segment module $X_j(m_0,\theta'_{j,b},o'_j)$, with $m_0$ the constant function,
  all of whose monomials $\mu$ have degree $b(v)$ in $z_{v,1},\dots,z_{v,k}$ for every $v$ and exponent
  less than $Q_j$ at $z_{f(j,i),c}$ for all $i<r$ and $c\leq r$.
\item The \emph{coarse generators} are the polynomials
  $u=\prod_{j=0}^{J-1}\psi_j(y_j)\cdot X^\ell$ of multidegree $d^D$ with $y_j\in M_j(b_j)$ for some
  $b_j$ and $X^\ell$ a coarse label. Their set is $\Lambda^{\mathrm c}$.
\end{enumerate}
\end{definition}
\leanref{Formalism/SharpRestriction/GrowthBounds/SegmentDegree/Coarse.lean}{Formalism.MainNotions.SelectionLemmas.Pf.Dat.Dl}
\leanref{Formalism/SharpRestriction/GrowthBounds/SegmentDegree/Coarse.lean}{Formalism.MainNotions.SelectionLemmas.Pf.Dat.DlSeg}
\leanref{Formalism/SharpRestriction/GrowthBounds/SegmentDegree/Coarse.lean}{Formalism.MainNotions.SelectionLemmas.Pf.Dat.phc}
\leanref{Formalism/SharpRestriction/GrowthBounds/SegmentDegree/Coarse.lean}{Formalism.MainNotions.SelectionLemmas.Pf.Dat.IsCLabel}
\leanref{Formalism/SharpRestriction/GrowthBounds/SegmentDegree/Coarse.lean}{Formalism.MainNotions.SelectionLemmas.Pf.Dat.Mj}
\leanref{Formalism/SharpRestriction/GrowthBounds/SegmentDegree/Coarse.lean}{Formalism.MainNotions.SelectionLemmas.Pf.Dat.lamGen}

\begin{lemma}\label{pd-symbols-residue-discriminant}
Let $C_{La}$ and $C_k$ be real numbers, $p\geq13$ a prime, $\mathcal R$ a discrete valuation ring with
residue field $\kappa$, $n'\geq2$ an even integer, $r\geq2$ and $K\geq1$ integers, $k=r+K$,
$T\subset\mathbf Z/n'\mathbf Z$, $L_1,N_1\geq0$ integers, and $(J,f)$ a run system for $T$ of length $r$
valid for $(L_1,N_1)$. Let $\deg$, $E'$, $N_s$, $E_0$, $m_*$, $a$, $Y$ and $q$ satisfy the hypotheses
(1) to (4) of Theorem~\ref{pd-symbols-count} with this $C_k$, and let $D$ be the discriminant of
Definition~\ref{pd-symbols-residue-coarse} formed with them. Then the reduction of $D$ modulo the
maximal ideal of $\mathcal R$ is nonzero.
\end{lemma}
\leanref{Formalism/SharpRestriction/GrowthBounds/SegmentDegree/Coarse.lean}{Formalism.MainNotions.SelectionLemmas.Pf.DiscriminantNonzeroModUniformiser}
\leanref{Formalism/SharpRestriction/GrowthBounds/SegmentDegree/DiscriminantNonzeroModUniformiser.lean}{Formalism.MainNotions.SelectionLemmas.Pf.discriminant_nonzero_mod_uniformiser}
\leanref{Formalism/SharpRestriction/GrowthBounds/SegmentDegree/DiscriminantNonzeroOfInjectiveTops.lean}{Formalism.MainNotions.SelectionLemmas.Pf.discriminant_nonzero_of_injective_tops}
\leanref{Formalism/SharpRestriction/GrowthBounds/SegmentDegree/LateVectorsDistinctTops.lean}{Formalism.MainNotions.SelectionLemmas.Pf.late_vectors_distinct_tops}

\begin{proof}
As $\kappa[X]$ is a domain, it suffices that the reduction of
$\phi_j\big(\det(z_{\rho_{j,v}(l),b})_{l,b}\big)=\det\big(X_{\overline{\rho_{j,v}(l)},b}^{\,q_l}\big)_{l,b}$,
with $q_l=p^{\operatorname{dist}_j(\rho_{j,v}(l))}$, is nonzero for $0\leq j<J$ and every new
$v\in\Lambda_j$. We show first that the residues $\overline{\rho_{j,v}(l)}$, $l\leq r$, are distinct.
Let $S$ be the set of even elements of $\Lambda_j$ smaller than $v$; then $\rho_{j,v}(1),\dots,
\rho_{j,v}(r)$ are the elements of $S$ that are at least $y=\rho_{j,v}(1)$. By conditions (1) to (3) of
validity, the $f(t,i)$ with $t\in\mathbf Z$ and $i<r$ are distinct even integers with residues in $T$,
and $f(t,i)<f(t',i')$ for $t<t'$. Suppose that $\rho_{j,v}(l)\geq y+n'$ for some $l$. For $i<r$ the
translate $w_i=f(t_i,i)$ of $f(0,i)$ by a multiple of $n'$ (condition (1)) with $y\leq w_i<y+n'$ lies in
$S$: $f(j-1,0)\leq y\leq w_i<y+n'\leq\rho_{j,v}(l)\leq f(j,r-1)$ and $w_i<\rho_{j,v}(l)<v$. The $w_i$ are
distinct and smaller than $\rho_{j,v}(l)$, so $S$ has at least $r+1$ elements that are at least $y$, a
contradiction. Hence the $\rho_{j,v}(l)$ lie in $[y,y+n')$, and their residues are distinct.

Now evaluate $X_{\overline{\rho_{j,v}(l)},l}$ at $1$ for every $l$ and all other variables at $0$. The
matrix becomes the identity, as the residues are distinct and $q_l\geq1$, so the determinant becomes
$1$ in $\kappa$, and its reduction is nonzero.
\end{proof}

\begin{lemma}\label{pd-symbols-residue-top-count}
Let $C_k\geq32$ be a real number, $p\geq13$ a prime, $n'\geq2$ an even integer, $r\geq2$ and $K\geq1$
integers, $k=r+K$, $T\subset\mathbf Z/n'\mathbf Z$, $L_1,N_1\geq0$ integers, and $(J,f)$ a run system
for $T$ of length $r$ valid for $(L_1,N_1)$. Let $\deg$, $E'$, $N_s$, $E_0$, $m_*$, $a$, $Y$ and $q$
satisfy the hypotheses (1) to (4) of Theorem~\ref{pd-symbols-count} with this $C_k$, let $d$ be as in
Lemma~\ref{pd-symbols-normal-degree}, and let $E^D$ and $d^D$ be the functions of
Definition~\ref{pd-symbols-residue-coarse}(1) formed with these data. Then
$E^D(t)\leq a/40\leq d(t)$ and $d^D(t)\geq\frac9{10}\deg(t)$ for $t\in T$, and $E^D(s)=0$ for
$s\notin T$.
\end{lemma}
\leanref{Formalism/SharpRestriction/GrowthBounds/SegmentDegree/DiscriminantNeZeroOfResidueNeZero.lean}{Formalism.MainNotions.SelectionLemmas.Pf.TopCountLowerBounds}
\leanref{Formalism/SharpRestriction/GrowthBounds/SegmentDegree/TopCountLowerBoundsOfOccurrence.lean}{Formalism.MainNotions.SelectionLemmas.Pf.top_count_lower_bounds_of_occurrence}
\leanref{Formalism/SharpRestriction/GrowthBounds/SegmentDegree/TopCountLowerBoundsOfOccurrence.lean}{Formalism.MainNotions.SelectionLemmas.Pf.Et_le}
\leanref{Formalism/SharpRestriction/GrowthBounds/SegmentDegree/TopCountLowerBoundsOfOccurrence.lean}{Formalism.MainNotions.SelectionLemmas.Pf.sg_mul_le}
\leanref{Formalism/SharpRestriction/GrowthBounds/SegmentDegree/TopCountLowerBoundsOfOccurrence.lean}{Formalism.MainNotions.SelectionLemmas.Pf.a_ge}
\leanref{Formalism/SharpRestriction/GrowthBounds/SegmentDegree/TopCountLowerBoundsOfOccurrence.lean}{Formalism.MainNotions.SelectionLemmas.Pf.k_sq_mul_Y_ge_Ns}
\leanref{Formalism/SharpRestriction/GrowthBounds/SegmentDegree/TopOccursInAtMostTwoSegments.lean}{Formalism.MainNotions.SelectionLemmas.Pf.top_occurs_in_at_most_two_segments}

\begin{proof}
We use the notation of Definitions~\ref{pd-symbols-run-system} and~\ref{pd-symbols-residue-coarse}.
As in the proof of Lemma~\ref{pd-symbols-normal-degree}, $N_s\geq1$ and $q\geq1$. Since $p$, $Y$,
$k^2$, $N_s$ and $q^2$ are at least $1$, the last inequality of (4) gives $32(E_0+m_*+1)\leq a$; so
$E_0\leq a/32$, $a\geq32$ and $\log a\geq1$. By (2), $K\geq N_s+r\geq3$, so $k\geq5$. By (4),
$32N_s\leq32N_s^2\log a\leq C_kN_s^2\log a\leq k$, hence $k^2Y\geq k^2\geq5k\geq160N_s$.

Let $0\leq j<J$, let $v\in\Lambda_j$ be new and $1\leq l\leq r$. The element $\rho_{j,v}(1)$ is
$\operatorname{old}_j(v)$, so $\rho_{j,v}(l)\geq\operatorname{old}_j(v)$ and
$\operatorname{dist}_j(\rho_{j,v}(l))\leq\operatorname{dist}_j(\operatorname{old}_j(v))$. As
$\sigma_j(v)\leq B_j(v)/(k^2Y)$, and $\operatorname{old}_j(v)\in\Lambda_j$ has its residue in $T$, where
$\deg\leq2a$ by (1),
\[
 \sigma_j(v)\,p^{\operatorname{dist}_j(\rho_{j,v}(l))}\leq
 \frac{\deg\big(\overline{\operatorname{old}_j(v)}\big)}{p^{\operatorname{dist}_j(\operatorname{old}_j(v))}k^2Y}
 \,p^{\operatorname{dist}_j(\operatorname{old}_j(v))}\leq\frac{2a}{k^2Y}.
\]
Let $t\in\mathbf Z/n'\mathbf Z$ and let $n_j(t)$ be the number of $y\in\Lambda_j$ with $\bar y=t$. The
$\rho_{j,v}(l)$, $l\leq r$, are distinct elements of $\Lambda_j$, so at most $n_j(t)$ of them have
residue $t$. As $\Lambda_j$ has at most $N_s$ elements by (2), and hence at most $N_s$ new ones,
\[
 E^D(t)\leq\sum_{j=0}^{J-1}N_s\cdot\frac{2a}{k^2Y}\,n_j(t).
\]
By condition~(2) of validity, $f(j-2,r-1)<f(j-1,0)$, so every element of $\Lambda_j$ lies in
$(f(j-2,r-1),f(j-1,r-1)]$ or in $(f(j-1,r-1),f(j,r-1)]$. For $a_0\in\{1,2\}$ the intervals
$(f(j-a_0,r-1),f(j-a_0+1,r-1)]$, $0\leq j<J$, are disjoint, and by conditions (1) and (2) their union
is an interval of length $n'$; so together they contain at most one integer with residue $t$. Hence
$\sum_{j<J}n_j(t)\leq2$, and $E^D(t)\leq4N_sa/(k^2Y)\leq a/40$.

Let $t\in T$. By (1), $\deg(t)\geq a$ and $E'(t)\leq E_0\leq a/32$, so $d(t)=\deg(t)-E'(t)\geq
a-a/32\geq a/40\geq E^D(t)$, and $d^D(t)=d(t)-E^D(t)\geq\deg(t)-a/32-a/40\geq\frac9{10}\deg(t)$, since
$a/32+a/40=9a/160\leq\deg(t)/10$. For $s\notin T$, every $\rho_{j,v}(l)$ has its residue in $T$, so
$E^D(s)=0$.
\end{proof}

\begin{definition}\label{pd-symbols-residue-pairs}
Let $\mathcal R$ be a commutative ring, $r,K\geq1$ integers, $n'\geq2$ an even integer,
$T\subset\mathbf Z/n'\mathbf Z$ and $(J,f)$ a run system for $T$ of length $r$, with the notation of
Definitions~\ref{pd-symbols-run-system} and~\ref{pd-symbols-module}. Let $j\in\mathbf Z$, let
$v_0<v_1<\dots<v_{N-1}$ be the elements of $\Lambda_j$, and letter $\{0,\dots,N-1\}$ so that $i$ is a
vector letter if $v_i$ is even and a covector letter if $v_i$ is odd. Let $m\colon\mathbf Z\to\mathbf Z$
and $\theta,o\colon\mathbf Z\to\mathbf R$.
\begin{enumerate}
\item A \emph{segment pair} for $(j,m,\theta,o)$ is a pair $(\nu,\kappa)$ of a path $\nu$ of rank $r$
  for some $B'$ and a path $\kappa$ of rank $K$ for some $B''$, for this lettering
  (Definition~\ref{pd-symbols-path-generators}(2)), such that for every letter $i$: if $v_i$ is not
  new, then $B''(i)=0$; if $v_i$ is new and even, then $\kappa_{i+1}(1)+m(v_i)\leq\nu_i(r)$ and
  $B''(i)\leq o(v_i)$; if $v_i$ is new and odd, then $\kappa_i(1)+m(v_i)\leq\nu_{i+1}(r)$ and
  $B'(i)\leq o(v_i)$; and if $v_i$ is new, then $\theta(v_i)\leq\nu_i(r)$.
\item Let $(\nu,\kappa)$ be a segment pair and $I'$, $I''$ standard fillings of $\nu$ and $\kappa$. The
  \emph{segment generator} $w(\nu,\kappa,I',I'')\in\mathcal R[z]$ is obtained from
  $\iota_{\mathrm{pl}}(\Psi(\nu,I'))\,\iota_{\mathrm{tn}}(\Psi(\kappa,I''))$, formed with $n=r+K$ as in
  Definition~\ref{pd-symbols-path-generators}(5), by renaming each variable $z_{i,c}$ to $z_{v_i,c}$.
\end{enumerate}
\end{definition}
\leanref{Formalism/SharpRestriction/GrowthBounds/ValuationLowestDegreeRename.lean}{Formalism.MainNotions.SelectionLemmas.segE}
\leanref{Formalism/SharpRestriction/GrowthBounds/ValuationLowestDegreeRename.lean}{Formalism.MainNotions.SelectionLemmas.segLL}
\leanref{Formalism/SharpRestriction/GrowthBounds/ValuationLowestDegreeRename.lean}{Formalism.MainNotions.SelectionLemmas.segRen}
\leanref{Formalism/SharpRestriction/GrowthBounds/ValuationLowestDegreeRename.lean}{Formalism.MainNotions.SelectionLemmas.segPairs}
\leanref{Formalism/SharpRestriction/GrowthBounds/StatementOfDegreeAndBound.lean}{Formalism.MainNotions.CovariantDegrees.XXZ}

\begin{lemma}\label{pd-symbols-residue-segment-span}
Let $p$ be a prime, $\mathcal R$ a complete discrete valuation ring in which $p$ is neither $0$ nor a
unit, $r,K\geq1$ integers, $n'\geq2$ an even integer, $T\subset\mathbf Z/n'\mathbf Z$, $L_1,N_1\geq0$
integers and $(J,f)$ a run system for $T$ of length $r$ valid for $(L_1,N_1)$. Let $j\in\mathbf Z$ with
$|\Lambda_j|+r\leq K$, $m\colon\mathbf Z\to\mathbf Z_{\geq0}$ and $\theta,o\colon\mathbf Z\to\mathbf R$.
Then the segment module $X_j(m,\theta,o)$ of Definition~\ref{pd-symbols-module}(2) is the
$\mathcal R$-span of the segment generators for $(j,m,\theta,o)$
(Definition~\ref{pd-symbols-residue-pairs}).
\end{lemma}
\leanref{Formalism/SharpRestriction/GrowthBounds/ValuationLowestDegreeRename.lean}{Formalism.MainNotions.SelectionLemmas.NestedTypeBound}
\leanref{Formalism/SharpRestriction/GrowthBounds/ValuationLowestDegreeMapFractionField.lean}{Formalism.MainNotions.G.Joints.space_statement_of_lemmas}

\begin{proof}
We use the notation of Definitions~\ref{pd-symbols-path-generators}, \ref{pd-symbols-module}
and~\ref{pd-symbols-residue-pairs}; $v_0<\dots<v_{N-1}$ are the elements of $\Lambda_j$, lettered as
there, and $S_i=\{0,\dots,i-1\}$ for $0\leq i\leq N$. Renaming $z_{i,c}$ to $z_{v_i,c}$ identifies the
polynomials in the $z_{i,c}$, $i<N$, with those in the $z_{v,c}$, $v\in\Lambda_j$. Under it the even $v_i$
correspond to the vector letters, and the rows $z_y$ with $y<v_i$, respectively $y\leq v_i$, that occur
are those of the letters in $S_i$, respectively $S_{i+1}$; so $\operatorname{val}_{\{y<v_i\}}$ and
$\lambda_{\{y<v_i\}}$ of a renamed polynomial are $\operatorname{val}_{S_i}$ and $\lambda_{S_i}$ of the
original, and likewise with $y\leq v_i$ and $S_{i+1}$. We use the renaming tacitly. If $\mu'$ is a
monomial in the plane variables of row degrees $B'$ and $\mu''$ one in the normal variables of row degrees
$B''$, as functions of the letters, then
\[
 \operatorname{ord}_{v_i}(\mu'\mu'')=B''(i)\ \text{if $v_i$ is even},\qquad
 \operatorname{ord}_{v_i}(\mu'\mu'')=B'(i)\ \text{if $v_i$ is odd}.\tag{$*$}
\]
Two facts on twisted degrees are used, for polynomials in $z_{i,b}$, $i<N$, $b\leq s$, over $\mathcal R$.

(a) If $\mu$ is a path of rank $s$ and $I$ a standard filling of $\mu$, then
$\mu_i(s)\leq\operatorname{val}_{S_i}(\Psi(\mu,I))$ and $\lambda_{S_i}(\Psi(\mu,I))\leq\mu_i(1)$ for
$0\leq i\leq N$. For $i\geq1$ this is Lemma~\ref{pd-symbols-path-word-bounds}, transferred from
$\mathbf Z$ to $\mathcal R$ as in the proof of Lemma~\ref{pd-symbols-path-clauses}(1). For $i=0$,
$P_{\emptyset,a}=P$, so only the exponent $0$ occurs, and $\mu_0=0$.

(b) Let $P\neq0$ have row degrees $B$, and let $E_v,E_l\subset\mathbf Z$ be finite sets. Then $P$ is an
$\mathcal R$-linear combination of path generators $\Psi(\mu,I)$ with $\mu$ a path of rank $s$ for $B$
such that for $0\leq i\leq N$: every $X\in E_v$ with $X\leq\operatorname{val}_{S_i}(P)$ satisfies
$X\leq\mu_i(s)$, and every $Y\in E_l$ with $\lambda_{S_i}(P)\leq Y$ satisfies $\mu_i(1)\leq Y$. Indeed, for
each letter $k$ choose an integer $\theta(k)\leq\operatorname{val}_{S_{k+1}}(P)$ with $X\leq\theta(k)$ for
every $X\in E_v$ with $X\leq\operatorname{val}_{S_{k+1}}(P)$: the twisted degree itself if it is finite,
and an upper bound of $E_v$ if it is $+\infty$. Likewise choose an integer
$M(k)\geq\lambda_{S_{k+1}}(P)$ with $M(k)\leq Y$ for every $Y\in E_l$ with $\lambda_{S_{k+1}}(P)\leq Y$.
Lemma~\ref{pd-symbols-path-clauses}(1) with these $\theta$ and $M$ gives a combination of the
$\Psi(\mu,I)$ with $\theta(k)\leq\mu_{k+1}(s)$ and $\mu_{k+1}(1)\leq M(k)$ for every letter $k$, which
gives the conditions at $i=k+1$. At $i=0$, $P_{\emptyset,a}=P\neq0$ has only the exponent $0$, so both
twisted degrees are $0=\mu_0(s)=\mu_0(1)$.

\emph{Segment generators lie in $X_j(m,\theta,o)$.} Let $(\nu,\kappa)$ be a segment pair, $\nu$ a path for
$B'$ and $\kappa$ one for $B''$, and $I'$, $I''$ standard fillings. Let $A$ be $\Psi(\nu,I')$ with
$z_{i,b}$ renamed $z_{v_i,b}$, and $G$ be $\Psi(\kappa,I'')$ with $z_{i,b}$ renamed $z_{v_i,r+b}$; so
$w(\nu,\kappa,I',I'')=AG$, and $A$ is a polynomial in the plane variables $z_{v,c}$, $v\in\Lambda_j$. By
Lemma~\ref{pd-symbols-path-basis}, $\Psi(\kappa,I'')$ has row degrees $B''$, and $B''(i)=0$ unless $v_i$
is new; so $G$ is a polynomial in the normal variables $z_{v,c}$ with $v\in\Lambda_j$ new. Put
$M(v_i)=\kappa_{i+1}(1)$ if $v_i$ is even and $M(v_i)=\kappa_i(1)$ if $v_i$ is odd. Let $v_i$ be new. By
(a), if $v_i$ is even, $\lambda_{\{y\leq v_i\}}(G)\leq\kappa_{i+1}(1)=M(v_i)$ and
$M(v_i)+m(v_i)\leq\nu_i(r)\leq\operatorname{val}_{\{y<v_i\}}(A)$; if $v_i$ is odd,
$\lambda_{\{y<v_i\}}(G)\leq\kappa_i(1)=M(v_i)$ and
$M(v_i)+m(v_i)\leq\nu_{i+1}(r)\leq\operatorname{val}_{\{y\leq v_i\}}(A)$; and
$\theta(v_i)\leq\nu_i(r)\leq\operatorname{val}_{\{y<v_i\}}(A)$. By $(*)$ and the bounds $B''(i)\leq o(v_i)$
and $B'(i)\leq o(v_i)$ of a segment pair, every monomial $\mu$ of $AG$ has $\operatorname{ord}_{v_i}(\mu)\leq
o(v_i)$. So $AG\in X_j(m,\theta,o)$.

\emph{Products lie in the span of generators.} Let $\mathcal G_0$ be the $\mathcal R$-span of the
$w(\nu,\kappa,I',I'')$ with $(\nu,\kappa)$ satisfying the conditions of
Definition~\ref{pd-symbols-residue-pairs}(1) except the two bounds by $o$. Let $AG$ be a product as in
Definition~\ref{pd-symbols-module}(2), with its function $M$. Write $A=\sum_{\beta'}A_{\beta'}$ and
$G=\sum_{\beta''}G_{\beta''}$ as sums of their parts of given row degrees. As in Step~4 of the proof of
Lemma~\ref{pd-symbols-residue-generators}, $\operatorname{val}_S(A_{\beta'})\geq\operatorname{val}_S(A)$
and $\lambda_S(G_{\beta''})\leq\lambda_S(G)$ for every $S$. Let $A_{\beta'}\neq0$ and $G_{\beta''}\neq0$. Let
$E_v$ be the finite set of the integers $M(v_i)+m(v_i)$ and $\lceil\theta(v_i)\rceil$, and $E_l$ that of
the $M(v_i)$, for $i<N$. By (b) with $E_v$ and $\emptyset$, $A_{\beta'}$ is a combination of the
$\Psi(\nu,I')$, renamed, with $\nu$ a path of rank $r$ for $\beta'$ such that $X\leq\nu_i(r)$ for every
$X\in E_v$ with $X\leq\operatorname{val}_{S_i}(A_{\beta'})$. By (b) with $\emptyset$ and $E_l$,
$G_{\beta''}$ is one of the $\Psi(\kappa,I'')$, renamed, with $\kappa$ a path of rank $K$ for $\beta''$
such that $\kappa_i(1)\leq Y$ for every $Y\in E_l$ with $\lambda_{S_i}(G_{\beta''})\leq Y$. Each such pair
$(\nu,\kappa)$ satisfies the conditions defining $\mathcal G_0$. First, $\beta''(i)=0$ unless $v_i$ is new,
as $G_{\beta''}\neq0$ involves no variable $z_{v_i,c}$ with $v_i$ not new. Let $v_i$ be new and even. Then
$\lambda_{S_{i+1}}(G_{\beta''})\leq\lambda_{\{y\leq v_i\}}(G)\leq M(v_i)$, so $\kappa_{i+1}(1)\leq M(v_i)$;
and $M(v_i)+m(v_i)\leq\operatorname{val}_{\{y<v_i\}}(A)\leq\operatorname{val}_{S_i}(A_{\beta'})$, so
$M(v_i)+m(v_i)\leq\nu_i(r)$; hence $\kappa_{i+1}(1)+m(v_i)\leq\nu_i(r)$. If $v_i$ is new and odd, then in the
same way, with $S_i$ and $S_{i+1}$ exchanged, $\kappa_i(1)+m(v_i)\leq\nu_{i+1}(r)$. Finally
$\theta(v_i)\leq\operatorname{val}_{\{y<v_i\}}(A)\leq\operatorname{val}_{S_i}(A_{\beta'})$, and as the twisted
degree is an integer or $+\infty$, $\lceil\theta(v_i)\rceil\leq\operatorname{val}_{S_i}(A_{\beta'})$; so
$\theta(v_i)\leq\lceil\theta(v_i)\rceil\leq\nu_i(r)$. So $A_{\beta'}G_{\beta''}\in\mathcal G_0$, and
$AG\in\mathcal G_0$.

\emph{The bound by $o$.} Let $\pi$ be the $\mathcal R$-linear map that keeps the terms whose monomials
$\mu$ satisfy $\operatorname{ord}_v(\mu)\leq o(v)$ for every new $v\in\Lambda_j$ and drops the others. For a
generator $w$ of $\mathcal G_0$ formed from $(\nu,\kappa)$ with row degrees $B'$ and $B''$, all monomials of
$w$ have the same $\operatorname{ord}_{v_i}$ by $(*)$; so $\pi(w)=w$ if $(\nu,\kappa)$ is a segment pair, and
$\pi(w)=0$ otherwise. Hence $\pi(\mathcal G_0)$ lies in the span of the segment generators. An element $x$
of $X_j(m,\theta,o)$ lies in $\mathcal G_0$ by the previous step and satisfies $\pi(x)=x$, so it lies in that
span.
% Lean: the first clause of NestedTypeBound is first_clause_segment_rename O h3 (with h3 the first component
%   of DefiniteRankOne, Lemma pd-symbols-path-clauses(1)), inside space_statement_at_ring_of_basis; it does
%   not use validity, |Lambda_j| + r <= K or m >= 0. Renaming: seg_val, segRen_mul. (a) is
%   standard_polynomial_mem_bounds; (b) is mem_span_standard_of_bounds (thresholds by exists_thr_top and
%   exists_thr_bot), used in crit_alt with Bv = {M + m} u {ceil th} and Bl = {M}. Products: crit_alt (exists_sum_polComp,
%   valS_le_polComp, lowest_degree_compose_le, pol_deg_zero, span_mul_mem); the bound by o: ord_proj
%   (word_ord, suppProj); generators: the second half of first_clause_segment_rename with
%   space_criterion_clause and pol_vars_deg.
\end{proof}

\begin{lemma}\label{pd-symbols-residue-top-degree}
Let $C_{La}$ and $C_k$ be real numbers, and let $p$, $\mathcal R$, $n'$, $r$, $K$, $k$, $T$, $L_1$,
$N_1$, $(J,f)$, $\deg$, $E'$, $N_s$, $E_0$, $m_*$, $a$, $Y$ and $q$ be as in
Lemma~\ref{pd-symbols-residue-rank}, except that $C_{La}$ and $C_k$ need not be at least $1088$; we use
the notation of Definition~\ref{pd-symbols-residue-coarse}. For $0\leq j<J$ let
$b_j\colon\mathbf Z\to\mathbf Z_{\geq0}$ and $y_j\in M_j(b_j)$, and let $X^\ell$ be a coarse label such
that $u=\prod_{j=0}^{J-1}\psi_j(y_j)\cdot X^\ell$ is nonzero and has multidegree $d^D$. Then
$d^D(\bar v)/p^{\operatorname{dist}_j(v)+R_j}-k<b_j(v)$ for $0\leq j<J$ and every $v\in\Lambda'_j$.
\end{lemma}
\leanref{Formalism/SharpRestriction/GrowthBounds/SegmentDegree/DiscriminantNeZeroOfResidueNeZero.lean}{Formalism.MainNotions.SelectionLemmas.Pf.TopDegreeComparison}
\leanref{Formalism/SharpRestriction/GrowthBounds/SegmentDegree/TopDegreeComparison.lean}{Formalism.MainNotions.SelectionLemmas.Pf.top_degree_comparison}

\begin{proof}
We use the notation of Definitions~\ref{pd-symbols-run-system}, \ref{pd-symbols-module}
and~\ref{pd-symbols-residue-coarse}. Fix $0\leq j<J$ and $v\in\Lambda'_j$, and put $t=\bar v$ and
$P=p^{\operatorname{dist}_j(v)+R_j}$. A polynomial in the $X_{s,c}$ has \emph{$t$-degree} $e$ if each of
its monomials has degree $e$ in $X_{t,1},\dots,X_{t,k}$. As
$f(j-1,0)\leq f(j-1,r-1)<v<f(j,0)\leq f(j,r-1)$, $v$ is a new element of $\Lambda_j$.

First, if $0\leq j'<J$, $x\in\Lambda_{j'}$ and $\bar x=t$, then $j'=j$ and $x=v$. Write $x=v-mn'$ with
$m\in\mathbf Z$. By condition~(2) of validity, $i\mapsto f(i,0)$ and $i\mapsto f(i,r-1)$ are strictly
increasing, and by condition~(1), $f(j-mJ,l)=f(j,l)-mn'$ for $l<r$. With $g=j-mJ$ this gives
$f(g-1,r-1)<x<f(g,0)$, while $f(j'-1,0)\leq x\leq f(j',r-1)$. If $j'<g$, then
$f(j',r-1)\leq f(g-1,r-1)<x$, and if $j'>g$, then $f(j'-1,0)\geq f(g,0)>x$; so $j'=g=j-mJ$, and $m=0$
as $0\leq j,j'<J$.

Next, let $0\leq j'<J$. By Definition~\ref{pd-symbols-module}(2) every monomial of $y_{j'}$ involves
only variables $z_{x,c}$ with $x\in\Lambda_{j'}$, and by Definition~\ref{pd-symbols-residue-coarse}(3) it
has degree $b_{j'}(x)$ in the row $z_x$. As $\psi_{j'}(z_{x,c})=X_{\bar
x,c}^{\,p^{\operatorname{dist}_{j'}(x)+R_{j'}}}$, the image of such a monomial has $t$-degree
$\sum_xp^{\operatorname{dist}_{j'}(x)+R_{j'}}b_{j'}(x)$, over the $x\in\Lambda_{j'}$ with $\bar x=t$. By the
first step, $\psi_{j'}(y_{j'})$ has $t$-degree $Pb_j(v)$ if $j'=j$ and $0$ otherwise. The monomial $X^\ell$
has $t$-degree $\sum_c\ell(t,c)$. Hence every monomial of $u$ has $t$-degree
$Pb_j(v)+\sum_c\ell(t,c)$, and as $u\neq0$ has multidegree $d^D$,
$d^D(t)=Pb_j(v)+\sum_c\ell(t,c)$. Since $X^\ell$ is a coarse label and $v$ is a new element of
$\Lambda_j$, $\ell(t,c)<P$ for $c\leq k$. So $d^D(t)\leq Pb_j(v)+k(P-1)<P(b_j(v)+k)$, which is the claim.
% Lean: top_degree_comparison; the first step is LL_unique, the t-degrees are DegAt with degAt_prod,
%   degAt_mul, degAt_monomial, multideg_imgMon and XX'_supported.
\end{proof}

\begin{lemma}\label{pd-symbols-residue-lifts-numeric}
Let $C_{La}\geq1088$ and $C_k\geq1088$ be real numbers, and let $p$, $\mathcal R$, $\kappa$, $n'$, $r$,
$K$, $k$, $T$, $L_1$, $N_1$, $(J,f)$, $\deg$, $E'$, $N_s$, $E_0$, $m_*$, $a$, $Y$ and $q$ be as in
Lemma~\ref{pd-symbols-residue-rank}; we use the notation of Definitions~\ref{pd-symbols-module}
and~\ref{pd-symbols-residue-coarse}. Let $j\in\mathbf Z$ and let $v$ be a new element of $\Lambda_j$.
Then:
\begin{enumerate}
\item $m_*<p^{R_j}$;
\item $p^{R_j}o'_j(v)+kp^{R_j}\leq o_j(v)$;
\item if $b\colon\mathbf Z\to\mathbf Z_{\geq0}$ satisfies
  $d^D(\bar x)/p^{\operatorname{dist}_j(x)+R_j}-k<b(x)$ for every $x\in\Lambda'_j$, then
  $\theta_j(v)\leq p^{R_j}\big(\theta'_{j,b}(v)-(r+K)N_s\big)$.
\end{enumerate}
\end{lemma}
\leanref{Formalism/SharpRestriction/GrowthBounds/SegmentDegree/ColumnConditionsOfNumericBounds.lean}{Formalism.MainNotions.SelectionLemmas.Pf.PowerExceedsMaxMultiplicity}
\leanref{Formalism/SharpRestriction/GrowthBounds/SegmentDegree/ReflectionMemGlevCastSucc.lean}{Formalism.MainNotions.SelectionLemmas.Pf.power_exceeds_max_multiplicity}
\leanref{Formalism/SharpRestriction/GrowthBounds/SegmentDegree/ReflectionMemGlevCastSucc.lean}{Formalism.MainNotions.SelectionLemmas.Pf.OrderNumericBound}
\leanref{Formalism/SharpRestriction/GrowthBounds/SegmentDegree/ThresholdPowerBoundOfCounts.lean}{Formalism.MainNotions.SelectionLemmas.Pf.order_numeric_bound_of_oldest}
\leanref{Formalism/SharpRestriction/GrowthBounds/SegmentDegree/ColumnConditionsOfNumericBounds.lean}{Formalism.MainNotions.SelectionLemmas.Pf.ThresholdPowerLowerBound}
\leanref{Formalism/SharpRestriction/GrowthBounds/SegmentDegree/ThresholdPowerBoundOfCounts.lean}{Formalism.MainNotions.SelectionLemmas.Pf.threshold_power_bound_of_counts}

\begin{proof}
We use the notation of Definitions~\ref{pd-symbols-run-system}, \ref{pd-symbols-module}
and~\ref{pd-symbols-residue-coarse}, and put $D_j=p^{\operatorname{dist}_j(f(j-1,0))}$ and
$T_j=a/(8Yk^2N_sD_j)$.

\emph{Step 1: the size of $p^{R_j}$.} As $f(j-1,0)\leq f(j-1,r-1)<f(j,0)\leq f(j,r-1)$ and the residue
of $f(j-1,0)$ lies in $T$ (conditions (2) and (3) of validity), $f(j-1,0)\in\Lambda_j$. As in the proof
of Lemma~\ref{pd-symbols-normal-degree}, $N_s\geq1$ and $q\geq1$; by (2), $K\geq N_s+r\geq3$, so
$k\geq5$; and $Y\geq1$ by (4). By (3), $D_j\leq q\leq q^2$, so by (4)
\[
 4p(E_0+m_*+1)\cdot8Yk^2N_sD_j\leq32pYk^2N_sq^2(E_0+m_*+1)\leq a,
\]
that is, $T_j\geq4p(E_0+m_*+1)\geq1$. So $\lfloor T_j\rfloor\geq1$, and $R_j$ is the largest integer
with $p^{R_j}\leq\lfloor T_j\rfloor$; hence
\[
 p^{R_j}\leq T_j<\lfloor T_j\rfloor+1\leq p^{R_j+1},\qquad a=8Yk^2N_sD_jT_j .
\]
In particular $a\geq0$.

\emph{Step 2: the oldest element.} Put $o=\operatorname{old}_j(v)$. The $f(j-1,i)$, $i<r$, are even
elements of $\Lambda_j$ smaller than $v$, as $v$ is new; so $o$ is the $r$-th largest even element of
$\Lambda_j$ smaller than $v$. Thus $o\in\Lambda_j$ and $f(j-1,0)\leq o<v$, and $v-o>L_1$ by
condition~(6) of validity. As $v\leq f(j,r-1)$, this gives
\[
 \operatorname{dist}_j(v)+L_1+1\leq\operatorname{dist}_j(o)\leq\operatorname{dist}_j(f(j-1,0)).
\]
Moreover, if $o$ is new, then $o\in\Lambda'_j$. Indeed, every $y\in\Lambda_j$ with $y\geq f(j,0)$ is one
of the $f(j,i)$, $i<r$: otherwise condition~(3) of validity, applied between consecutive $f(j,i)$, gives
$f(j,i)<y$ for every $i<r$, against $y\leq f(j,r-1)$. If $o\geq f(j,0)$, the $r$ even elements of
$\Lambda_j$ that are at least $o$ and smaller than $v$, together with $v$, would be $r+1$ distinct
elements among the $r$ elements $f(j,i)$. So $o<f(j,0)$, and if $o$ is new, then
$f(j-1,r-1)<o<f(j,0)$ and $o\in\Lambda'_j$.

(1) By Step 1, $4p(E_0+m_*+1)\leq T_j<p\cdot p^{R_j}$, so $m_*<4(E_0+m_*+1)<p^{R_j}$.

(2) Put $P=p^{L_1+\operatorname{dist}_j(v)}$. By Step 2, $pP\leq p^{\operatorname{dist}_j(o)}$ and
$P\leq D_j$. As $8YkN_s\geq1$, Step 1 gives
\[
 kp^{R_j}P\leq kT_jD_j\leq8YkN_s\cdot kT_jD_j=a,
\]
so $kp^{R_j}\leq a/P$. As $p\geq4$ and $a\geq0$, $4aP\leq apP\leq ap^{\operatorname{dist}_j(o)}$, so
$4ap^{-\operatorname{dist}_j(o)}\leq a/P$. By Definitions~\ref{pd-symbols-module}(4)
and~\ref{pd-symbols-residue-coarse}(3),
\[
 p^{R_j}o'_j(v)=\frac{L_aa}{pk\,p^{\operatorname{dist}_j(v)}}+\epsilon'_v\,a\,p^{-\operatorname{dist}_j(o)},\qquad
 o_j(v)=\frac{L_aa}{pk\,p^{\operatorname{dist}_j(v)}}+\epsilon_v\frac aP .
\]
If $v$ is even, then $\epsilon'_v=4$, $\epsilon_v=10$ and
$4ap^{-\operatorname{dist}_j(o)}+kp^{R_j}\leq2a/P\leq10a/P$. If $v$ is odd, then $\epsilon'_v=0$,
$\epsilon_v=1$ and $kp^{R_j}\leq a/P$. In both cases $p^{R_j}o'_j(v)+kp^{R_j}\leq o_j(v)$.

(3) Let $b$ be as in (3), $o=\operatorname{old}_j(v)$ and $P_o=p^{\operatorname{dist}_j(o)+R_j}$, and put
$c=\deg(\bar o)/P_o$ and $c'=d^D(\bar o)/P_o$. The residue $\bar o$ lies in $T$. By
Lemma~\ref{pd-symbols-residue-top-count}, which applies as $C_k\geq32$,
$d^D(\bar o)\geq\frac9{10}\deg(\bar o)$, so $c'\geq\frac9{10}c$; and $\deg(\bar o)\geq a$ by (1) of
Theorem~\ref{pd-symbols-count}. As $Y\geq1$ and $k\geq5$, $8Yk\geq40$, and Steps 1 and 2 give
\[
 40kN_sP_o\leq8Yk^2N_sD_jT_j=a\leq\deg(\bar o),
\]
so $c\geq40kN_s$, that is, $c/10\geq4kN_s$. Since $\theta_j(v)=B_j(v)/2=p^{R_j}c/2$, it suffices to show
that $c/2\leq\theta'_{j,b}(v)-kN_s$; recall that $k=r+K$.

If $o$ is not new, then $\theta'_{j,b}(v)=\frac23(W_j(o)-m_0)$, where
$W_j(o)=\lfloor c'\rfloor-(k+r+1)>c'-k-r-2$ and $m_0=kN_s+1$. As $r\leq k$ and $3\leq k\leq kN_s$,
$k+r+3+kN_s\leq4kN_s$, so $\theta'_{j,b}(v)>\frac23(c'-4kN_s)\geq\frac35c-\frac83kN_s$ and
$\theta'_{j,b}(v)-kN_s>\frac35c-4kN_s\geq\frac35c-\frac c{10}=\frac c2$.

If $o$ is new, then $o\in\Lambda'_j$ by Step 2, so $c'-k<b(o)$ by the hypothesis on $b$, and
$\theta'_{j,b}(v)=\frac23b(o)>\frac23(c'-k)\geq\frac35c-\frac23k$. As $\frac23k+kN_s\leq2kN_s\leq c/10$,
$\theta'_{j,b}(v)-kN_s>\frac35c-\frac c{10}=\frac c2$.
% Lean: Step 1 is basics, reflection_mem_Glev_castSucc, Tj (= T_j), Tj_ge (from hX.twist and P3),
%   pow_Rj_le and lt_pow_Rj; Step 2 is oldest_letter_in_segment (r_le_length, evL, the "frame" step) with
%   the validity bound L1 < x - oldest j x. (1) is power_exceeds_max_multiplicity; (2) is
%   order_numeric_bound_of_oldest (hA, hB, hLe, hT1); (3) is threshold_power_bound_of_counts C0
%   (TopCountLowerBounds, Lemma pd-symbols-residue-top-count; hc, hq, hsplit; natDiv_gt,
%   threshold_arith_w_case and threshold_arith_letter_case). All three are stated for every C0 (with
%   32 <= C0 for (3)); the book uses C0 = 1088.
\end{proof}

\begin{lemma}\label{pd-symbols-residue-lifts}
Let $C_{La}\geq1088$ and $C_k\geq1088$ be real numbers, and let $p$, $\mathcal R$, $\kappa$, $n'$, $r$,
$K$, $k$, $T$, $L_1$, $N_1$, $(J,f)$, $\deg$, $E'$, $N_s$, $E_0$, $m_*$, $a$, $Y$, $q$ and $X_j$ be as
in Lemma~\ref{pd-symbols-residue-rank}; we use the notation of
Definition~\ref{pd-symbols-residue-coarse}. For $j\in\mathbf Z$ let
$D^{\mathrm{seg}}_j=\prod_v\det(z_{\rho_{j,v}(l),b})_{1\leq l,b\leq r}^{\,\sigma_j(v)}$, the product over
the new $v\in\Lambda_j$, and let $\epsilon_j$ be the $\mathcal R$-algebra endomorphism of
$\mathcal R[z]$ with $\epsilon_j(z_{v,c})=z_{v,c}^{\,p^{R_j}}$. Let $j\in\mathbf Z$, let
$b\colon\mathbf Z\to\mathbf Z_{\geq0}$ satisfy $d^D(\bar v)/p^{\operatorname{dist}_j(v)+R_j}-k<b(v)$ for
every $v\in\Lambda'_j$, and let $z^e$ be a monomial with $e(v,c)<p^{R_j}$ for all $v$ and $c$, and
$e(v,c)=0$ unless $v$ is a new element of $\Lambda_j$. Let $w$ be a segment generator for
$(j,m_0,\theta'_{j,b},o'_j)$ (Definition~\ref{pd-symbols-residue-pairs}). Then there is $x\in X_j$ whose
reduction modulo the maximal ideal of $\mathcal R$ is that of $D^{\mathrm{seg}}_j\,\epsilon_j(w)\,z^e$.
\end{lemma}
\leanref{Formalism/SharpRestriction/GrowthBounds/SegmentDegree/SegmentReductionOfColumnWords.lean}{Formalism.MainNotions.SelectionLemmas.Pf.ColumnWordReductionLifts}
\leanref{Formalism/SharpRestriction/GrowthBounds/SegmentDegree/ColumnWordReductionOfFour.lean}{Formalism.MainNotions.SelectionLemmas.Pf.column_word_reduction_of_four}
\leanref{Formalism/SharpRestriction/GrowthBounds/SegmentDegree/FrobeniusReductionProduct.lean}{Formalism.MainNotions.SelectionLemmas.Pf.frobenius_reduction_product}
\leanref{Formalism/SharpRestriction/GrowthBounds/SegmentDegree/AllowedLettersOfLateRows.lean}{Formalism.MainNotions.SelectionLemmas.Pf.allowed_letters_of_late_rows}
\leanref{Formalism/SharpRestriction/GrowthBounds/SegmentDegree/ColumnConditionsOfNumericBounds.lean}{Formalism.MainNotions.SelectionLemmas.Pf.column_conditions_of_numeric_bounds}
\leanref{Formalism/SharpRestriction/GrowthBounds/SegmentDegree/WordOrderBoundOfNumericBound.lean}{Formalism.MainNotions.SelectionLemmas.Pf.word_order_bound_of_numeric_bound}
\leanref{Formalism/SharpRestriction/GrowthBounds/ValuationLowestDegreeMapFractionField.lean}{Formalism.MainNotions.G.Joints.StandardMonomialSpace.space_statement_sixth_clause}

\begin{proof}
We use the notation of Definitions~\ref{pd-symbols-run-system}, \ref{pd-symbols-module},
\ref{pd-symbols-residue-pairs} and~\ref{pd-symbols-residue-coarse}. In this proof $\kappa$ denotes the second
path of a segment pair, $\bar{\mathcal R}$ the residue field, and bars reductions modulo the maximal ideal;
$\bar{\mathcal R}$ has characteristic $p$. Put $Q=p^{R_j}$. Let $w$ be formed from the segment pair
$(\nu,\kappa)$, with $\nu$ a path for $B'$ and $\kappa$ one for $B''$, and standard fillings $I'$ and $I''$;
let $v_0<\dots<v_{N-1}$ be the elements of $\Lambda_j$, $A_0$ the polynomial $\Psi(\nu,I')$ with $z_{i,b}$
renamed $z_{v_i,b}$, and $G_0$ the polynomial $\Psi(\kappa,I'')$ with $z_{i,b}$ renamed $z_{v_i,r+b}$, so
that $w=A_0G_0$. Let $\ell_P$ and $\ell_N$ be the products of the $z_{v,c}^{e(v,c)}$ with $c\leq r$ and with
$c>r$, so that $z^e=\ell_P\ell_N$, and let $d_S$ be the sum of the $e(v,c)$ with $v$ odd and $c\leq r$, and
$d_V$ that of the $e(v,c)$ with $v$ even and $c>r$. As $e(v,c)<Q$, at most $N_s$ elements $v$ are new, and
there are $k=r+K$ values of $c$, $d_S+d_V\leq(r+K)N_sQ$.

\emph{Step 1: the reduction.} The polynomials $A_0$ and $G_0$ have integer coefficients, and over a ring of
characteristic $p$ a polynomial $P$ with coefficients in the prime field satisfies $P(z^Q)=P(z)^Q$. So
$D^{\mathrm{seg}}_j\,\epsilon_j(w)\,z^e$ reduces to $A_wG_w$, where
$A_w=\bar A_0^{\,Q}\,\bar\ell_P\,\bar D^{\mathrm{seg}}_j$ and $G_w=\bar G_0^{\,Q}\,\bar\ell_N$.

\emph{Step 2: the letters.} $A_w$ is a polynomial in the plane variables $z_{v,c}$ with $v\in\Lambda_j$, as the
$\rho_{j,x}(l)$ lie in $\Lambda_j$ and $e$ vanishes off the new elements. $G_w$ is a polynomial in the normal
variables $z_{v,c}$ with $v\in\Lambda_j$ new, as $B''(i)=0$ unless $v_i$ is new.

\emph{Step 3: the conditions of the segment module.} For every $S$ the substitution $P\mapsto P_{S,a}$ of
Definition~\ref{pd-symbols-module}(1) is a ring homomorphism, so $\operatorname{val}_S(PP')\geq
\operatorname{val}_S(P)+\operatorname{val}_S(P')$ and $\lambda_S(PP')\leq\lambda_S(P)+\lambda_S(P')$. As
$\tau_a(\eta)=1+(\eta-1)uv$ in the notation of the proof of Lemma~\ref{pd-symbols-path-word-bounds}, it
replaces a row $z_y$, $y\in S$, by linear forms in that row with coefficients of degree at most $1$ in $\eta$
if $y$ is even and in $\eta^{-1}$ if $y$ is odd. Hence a variable at an even $y$ has
$\operatorname{val}_S\geq0$ and $\lambda_S\leq1$, and one at an odd $y$ has $\operatorname{val}_S\geq-1$ and
$\lambda_S\leq0$; so $\operatorname{val}_S(\bar\ell_P)\geq-d_S$ and $\lambda_S(\bar\ell_N)\leq d_V$. Each
determinant $\det(z_{\rho_{j,x}(l),b})$ has even rows, so its $\operatorname{val}_S$ is $\geq0$; if all its
rows lie in $S$, the substitution multiplies it by $\det\tau_r(\eta)=\eta$, so its $\operatorname{val}_S$ is
$1$. Let $v=v_i$ be new. The rows $\rho_{j,v}(l)$ are even elements of $\Lambda_j$ smaller than $v$, so
$\operatorname{val}_S(\bar D^{\mathrm{seg}}_j)\geq\sigma_j(v)$ for $S=\{y<v\}$ and $S=\{y\leq v\}$. By fact
(a) of the proof of Lemma~\ref{pd-symbols-residue-segment-span}, which holds over $\bar{\mathcal R}$ in the
same way, $\operatorname{val}_{\{y<v\}}(\bar A_0)\geq\nu_i(r)$, $\operatorname{val}_{\{y\leq
v\}}(\bar A_0)\geq\nu_{i+1}(r)$, $\lambda_{\{y\leq v\}}(\bar G_0)\leq\kappa_{i+1}(1)$ and
$\lambda_{\{y<v\}}(\bar G_0)\leq\kappa_i(1)$.

Put $M_w(v_i)=Q\kappa_{i+1}(1)+d_V$ if $v_i$ is even and $M_w(v_i)=Q\kappa_i(1)+d_V$ if it is odd. Let $v=v_i$
be new and even. Then $\lambda_{\{y\leq v\}}(G_w)\leq Q\kappa_{i+1}(1)+d_V=M_w(v)$ and
$\operatorname{val}_{\{y<v\}}(A_w)\geq Q\nu_i(r)-d_S+\sigma_j(v)$. The segment pair gives
$\kappa_{i+1}(1)+m_0\leq\nu_i(r)$ with $m_0=(r+K)N_s+1$, and $m_*<Q$ by
Lemma~\ref{pd-symbols-residue-lifts-numeric}(1). As $m_j(v)=m_*+\sigma_j(v)$ and
$d_S+d_V\leq(r+K)N_sQ$,
\[
 M_w(v)+m_j(v)\leq Q\nu_i(r)-Q(r+K)N_s-Q+d_V+m_*+\sigma_j(v)\leq Q\nu_i(r)-d_S+\sigma_j(v)\leq
 \operatorname{val}_{\{y<v\}}(A_w).
\]
If $v$ is new and odd, the same argument with $\{y<v\}$ and $\{y\leq v\}$ exchanged and
$\kappa_i(1)+m_0\leq\nu_{i+1}(r)$ gives $\lambda_{\{y<v\}}(G_w)\leq M_w(v)$ and
$M_w(v)+m_j(v)\leq\operatorname{val}_{\{y\leq v\}}(A_w)$. Finally, for new $v$, $\nu_i(r)\geq\theta'_{j,b}(v)$,
$d_S\leq(r+K)N_sQ$ and $\sigma_j(v)\geq0$ give $\operatorname{val}_{\{y<v\}}(A_w)\geq
Q(\theta'_{j,b}(v)-(r+K)N_s)\geq\theta_j(v)$, by Lemma~\ref{pd-symbols-residue-lifts-numeric}(3). So
$A_wG_w$, with the function $M_w$, satisfies the conditions of Definition~\ref{pd-symbols-module}(2) for
$(j,m_j,\theta_j)$, over $\bar{\mathcal R}$.

\emph{Step 4: the bound by $o_j$.} Let $v=v_i$ be new and $\mu$ a monomial of $A_wG_w$. By
Lemma~\ref{pd-symbols-path-basis}, $\bar A_0$ and $\bar G_0$ have row degrees $B'$ and $B''$. If $v$ is even,
$\operatorname{ord}_v(\mu)$ is the degree of $\mu$ in the normal variables at $v$, which is
$QB''(i)+\sum_{c>r}e(v,c)\leq Qo'_j(v)+KQ$, as $B''(i)\leq o'_j(v)$ for a segment pair. If $v$ is odd, it is the
degree in the plane variables at $v$, which is $QB'(i)+\sum_{c\leq r}e(v,c)\leq Qo'_j(v)+rQ$, as
$\bar D^{\mathrm{seg}}_j$ involves only even rows. In both cases $\operatorname{ord}_v(\mu)\leq Qo'_j(v)+kQ\leq
o_j(v)$ by Lemma~\ref{pd-symbols-residue-lifts-numeric}(2).

\emph{Step 5: the lift.} By Steps 2 to 4, $A_wG_w$ lies in the segment module $X_j(m_j,\theta_j,o_j)$ formed
over $\bar{\mathcal R}$. The last two steps of the proof of Lemma~\ref{pd-symbols-residue-segment-span} apply
over $\bar{\mathcal R}$, as Lemma~\ref{pd-symbols-path-clauses}(1) holds there; so $A_wG_w=\sum_uc_u\bar u$ with
$c_u\in\bar{\mathcal R}$, over the segment generators $u$ for $(j,m_j,\theta_j,o_j)$, which have integer
coefficients. Lift each $c_u$ to $\tilde c_u\in\mathcal R$ and put $x=\sum_u\tilde c_uu$. As
$|\Lambda_j|+r\leq N_s+r\leq K$ by (2) of Theorem~\ref{pd-symbols-count} and $m_j\geq0$,
Lemma~\ref{pd-symbols-residue-segment-span} gives $x\in X_j(m_j,\theta_j,o_j)=X_j$, and $\bar x=A_wG_w$ is the
reduction of $D^{\mathrm{seg}}_j\,\epsilon_j(w)\,z^e$.
% Lean: column_word_reduction_of_four D hB12 1088 with h1 = frobenius_reduction_product (expand_pow_of_map,
%   ell_split; A_w = Aw, G_w = Gw, M_w = Mw), h2 = allowed_letters_of_late_rows (with late_row_injective),
%   h3 = column_conditions_of_numeric_bounds (vg_ellP, lg_ellN, vg_DlSeg with sDegs_det, vg_Aw, lg_Gw from the
%   additivity clauses of ValuationsAdditive; int_ineq; ell_degree_le_power_bound; the numeric inputs of
%   Lemma pd-symbols-residue-lifts-numeric), h4 = word_order_bound_of_numeric_bound (degAt_ellP, degAt_ellN,
%   ordDeg_prod), and the sixth clause of NestedTypeBound = space_statement_sixth_clause (crit_words over
%   FF O = crit_alt and ord_proj, XXZ_lift) with the first clause, Lemma pd-symbols-residue-segment-span.
\end{proof}

\begin{definition}\label{pd-symbols-residue-families}
Let $p$, $n'$, $r$, $K$, $k=r+K$, $T$, $(J,f)$ and the data of Definition~\ref{pd-symbols-residue-coarse}
be given, with that notation, and for $j,x\in\mathbf Z$ put $P_j(x)=p^{\operatorname{dist}_j(x)+R_j}$.
\begin{enumerate}
\item For $t\in\mathbf Z/n'\mathbf Z$, a \emph{remainder vector} at $t$ is an $\alpha\in\mathbf Z_{\geq0}^k$ with
  $\sum_c\alpha_c\leq d^D(t)$ such that for every $j$ and every new $x\in\Lambda_j$ with $\bar x=t$:
  $\alpha_c<P_j(x)$ for all $c$, and $\sum_c\alpha_c\equiv d^D(t)\pmod{P_j(x)}$. Let $\mathcal L$ be the set of
  the families $\alpha=(\alpha_t)_t$ of remainder vectors, and for $\alpha\in\mathcal L$ and $j,x\in\mathbf Z$
  let $n^\alpha_j(x)$ be the integer part of $(d^D(\bar x)-\sum_c\alpha_{\bar x,c})/P_j(x)$.
\item For $\alpha\in\mathcal L$, $E(\alpha)$ is the set of the families $e=(e_j)_{0\leq j<J}$ of vectors
  $e_j\in\mathbf Z_{\geq0}^r$ with $e_j(i)<rQ_j$ and $e_j(i)\equiv n^\alpha_j(f(j,i))\pmod{Q_j}$ for all $i<r$;
  for $j\in\mathbf Z$ write $e_j$ for $e_{j'}$ with $0\leq j'<J$ and $j'\equiv j\pmod J$.
\item For $\alpha\in\mathcal L$, $e\in E(\alpha)$ and $j\in\mathbf Z$, $b^{\alpha,e}_j\colon\mathbf Z\to\mathbf Z_{\geq0}$ is
  $0$ off $\Lambda_j$; at $x=f(j-1,i)$ it is the integer part of $(n^\alpha_{j-1}(x)-e_{j-1}(i))/Q_{j-1}$,
  with a negative difference replaced by $0$; at
  $x=f(j,i)$ it is $e_j(i)$; at the other $x\in\Lambda_j$ with $f(j,0)\leq x$ it is $0$; and at
  $x\in\Lambda'_j$ it is $n^\alpha_j(x)$.
\item For $j\in\mathbf Z$ and $b\colon\mathbf Z\to\mathbf Z_{\geq0}$: $b$ is \emph{balanced at $j$} if
  $d^D(\bar x)/P_j(x)-k<b(x)\leq d^D(\bar x)/P_j(x)$ for every $x\in\Lambda'_j$;
  \[
   S_j(b)=\big(1-a^{-N_s^2}\big)\prod_{x\in\Lambda'_j}\binom{b(x)+k-1}{k-1}
   \prod_{i<r}\binom{\max(0,W_j(f(j-1,i))-m_0)+k-1}{k-1}\Big(Q_j^{\,r-1}\max\Big(0,1-\frac r{2Y}\Big)\Big)^r;
  \]
  and for $\varepsilon\in\mathbf Z_{\geq0}^r$, $\tau_j(b,\varepsilon)$ agrees with $b$ except at the $x$ with
  $f(j,0)\leq x\leq f(j,r-1)$, where it is $\varepsilon(i)$ at $x=f(j,i)$ and $0$ elsewhere.
\end{enumerate}
\end{definition}
\leanref{Formalism/SharpRestriction/GrowthBounds/SegmentDegree/RestCountOf.lean}{Formalism.MainNotions.SelectionLemmas.Pf.Dat.Pz}
\leanref{Formalism/SharpRestriction/GrowthBounds/SegmentDegree/RestCountOf.lean}{Formalism.MainNotions.SelectionLemmas.Pf.Dat.laSet}
\leanref{Formalism/SharpRestriction/GrowthBounds/SegmentDegree/RestCountOf.lean}{Formalism.MainNotions.SelectionLemmas.Pf.Dat.LA}
\leanref{Formalism/SharpRestriction/GrowthBounds/SegmentDegree/RestCountOf.lean}{Formalism.MainNotions.SelectionLemmas.Pf.Dat.nz}
\leanref{Formalism/SharpRestriction/GrowthBounds/SegmentDegree/RestCountOf.lean}{Formalism.MainNotions.SelectionLemmas.Pf.Dat.EE}
\leanref{Formalism/SharpRestriction/GrowthBounds/SegmentDegree/RestCountOf.lean}{Formalism.MainNotions.SelectionLemmas.Pf.Dat.eZ}
\leanref{Formalism/SharpRestriction/GrowthBounds/SegmentDegree/RestCountOf.lean}{Formalism.MainNotions.SelectionLemmas.Pf.Dat.bfam}
\leanref{Formalism/SharpRestriction/GrowthBounds/SegmentDegree/RestCountOf.lean}{Formalism.MainNotions.SelectionLemmas.Pf.Dat.AdmLL}
\leanref{Formalism/SharpRestriction/GrowthBounds/SegmentDegree/RestCountOf.lean}{Formalism.MainNotions.SelectionLemmas.Pf.Dat.segBound}
\leanref{Formalism/SharpRestriction/GrowthBounds/SegmentDegree/RestCountOf.lean}{Formalism.MainNotions.SelectionLemmas.Pf.Dat.setTau}

\begin{lemma}\label{pd-symbols-residue-count-digits}
Let $k\geq1$, $n\geq0$ and $P\geq1$ be integers, and let $S$ be a finite set of vectors $v\in\mathbf Z_{\geq0}^k$ with
$\sum_cv_c\leq n$ that contains every such $v$ with $v_c<P$ for all $c$ and $\sum_cv_c\equiv n\pmod P$. Then:
\begin{enumerate}
\item $\binom{n+k-1}{k-1}\leq\sum_{v\in S}\binom{\lfloor(n-\sum_cv_c)/P\rfloor+k-1}{k-1}$;
\item if $kP\leq n$, then $S$ has at least $P^{k-1}$ elements.
\end{enumerate}
\end{lemma}
\leanref{Formalism/SharpRestriction/GrowthBounds/SegmentDegree/FactorialMulChooseEqProdRange.lean}{Formalism.MainNotions.SelectionLemmas.Pf.sum_digit_count}
\leanref{Formalism/SharpRestriction/GrowthBounds/SegmentDegree/FactorialMulChooseEqProdRange.lean}{Formalism.MainNotions.SelectionLemmas.Pf.sum_label_count}

\begin{proof}
(1) The left side is the number of $w\in\mathbf Z_{\geq0}^k$ with $\sum_cw_c=n$. Write $w_c=Pu_c+v_c$ with $0\leq v_c<P$.
Then $\sum_cv_c\leq n$ and $\sum_cv_c\equiv n\pmod P$, so $v\in S$, and the $w$ with a given $v$ correspond to the
$u\in\mathbf Z_{\geq0}^k$ with $\sum_cu_c=(n-\sum_cv_c)/P$, which number $\binom{(n-\sum_cv_c)/P+k-1}{k-1}$.

(2) For $u\in\{0,\dots,P-1\}^{k-1}$ let $v_1$ be the residue of $n-\sum_cu_c$ modulo $P$ in $\{0,\dots,P-1\}$, and
$v=(v_1,u)$. Then $v_c<P$ for all $c$, $\sum_cv_c\equiv n\pmod P$ and $\sum_cv_c<kP\leq n$, so $v\in S$; and $u\mapsto v$
is injective.
% Lean: (1) is sum_digit_count (sum_card_tuples, the fibres of w -> w mod P); (2) is sum_label_count (Fin.cons
%   of (n + (P - 1) sum u) mod P, which is the residue of n - sum u).
\end{proof}

\begin{lemma}\label{pd-symbols-residue-count-sum-facts}
Let $C_{La}\geq1088$ and $C_k\geq1088$ be real numbers, and let $p$, $\mathcal R$, $\kappa$, $n'$, $r$, $K$, $k$, $T$,
$L_1$, $N_1$, $(J,f)$, $\deg$, $E'$, $N_s$, $E_0$, $m_*$, $a$, $Y$ and $q$ be as in
Lemma~\ref{pd-symbols-residue-rank}; we use the notation of Definitions~\ref{pd-symbols-residue-coarse}
and~\ref{pd-symbols-residue-families}, and write $\mathcal N_j$ for the set of the new elements of $\Lambda_j$.
Then:
\begin{enumerate}
\item the residues of the elements of $\mathcal N_0,\dots,\mathcal N_{J-1}$ are distinct and form $T$, and $J\leq|T|$;
\item $\mathcal N_j$ is the disjoint union of $\Lambda'_j$ and $\{f(j,0),\dots,f(j,r-1)\}$, and the elements of
  $\Lambda'_j$ are smaller than $f(j,0)$;
\item if $x\in\mathcal N_j$ and $x'\in\mathcal N_{j'}$ have the same residue, then $P_j(x)=P_{j'}(x')$;
\item $\frac{36}5Yk^2N_sP_j(y)\leq d^D(\bar y)$ for $y\in\Lambda_j$, and $kP_j(x)\leq d^D(\bar x)$ for $x\in\mathcal N_j$;
\item $P_j(x)Q_j=P_{j+1}(x)$ for $x\leq f(j,r-1)$, $\prod_{j=0}^{J-1}Q_j=p^{n'}$, and
  $W_{j+J}(f(j+J-1,i))=W_j(f(j-1,i))$ for all $j$ and $i<r$;
\item $32Y\leq a$;
\item $\binom{\deg(t)+k-1}{k-1}\big(1-\frac5Y\big)\leq\binom{d^D(t)+k-1}{k-1}$ for $t\in T$;
\item for $j\in\mathbf Z$, $i<r$ and $x=f(j,i)$,
  $\binom{d^D(\bar x)+k-1}{k-1}\big(1-\frac1Y\big)\leq(P_j(x)Q_j)^{k-1}\binom{\max(0,W_{j+1}(x)-m_0)+k-1}{k-1}$.
\end{enumerate}
\end{lemma}
\leanref{Formalism/SharpRestriction/GrowthBounds/SegmentDegree/FamilyDegreesAdmissible.lean}{Formalism.MainNotions.SelectionLemmas.Pf.SumFacts}
\leanref{Formalism/SharpRestriction/GrowthBounds/SegmentDegree/FrameZeroStrictMono.lean}{Formalism.MainNotions.SelectionLemmas.Pf.sum_facts}
\leanref{Formalism/SharpRestriction/GrowthBounds/SegmentDegree/FrameZeroStrictMono.lean}{Formalism.MainNotions.SelectionLemmas.Pf.sum_kP}
\leanref{Formalism/SharpRestriction/GrowthBounds/SegmentDegree/FrameZeroStrictMono.lean}{Formalism.MainNotions.SelectionLemmas.Pf.sum_prodQ}
\leanref{Formalism/SharpRestriction/GrowthBounds/SegmentDegree/FrameZeroStrictMono.lean}{Formalism.MainNotions.SelectionLemmas.Pf.sum_compA}
\leanref{Formalism/SharpRestriction/GrowthBounds/SegmentDegree/FrameZeroStrictMono.lean}{Formalism.MainNotions.SelectionLemmas.Pf.sum_compW}

\begin{proof}
Put $D_j=p^{\operatorname{dist}_j(f(j-1,0))}$. By Step~1 of the proof of Lemma~\ref{pd-symbols-residue-lifts-numeric},
$p^{R_j}D_j\leq a/(8Yk^2N_s)<p^{R_{j+1}+1}D_{j+1}$ for every $j$. By condition~(1) of validity, $D_j$, $R_j$ and $Q_j$
depend only on $j$ modulo $J$, and $f(j+J,i)=f(j,i)+n'$, so $\Lambda_{j+J}=\Lambda_j+n'$ and
$\operatorname{dist}_{j+J}(y+n')=\operatorname{dist}_j(y)$.

(1) The sets $\mathcal N_j$, $0\leq j<J$, consist of the $x$ with $\bar x\in T$ in the consecutive intervals
$(f(j-1,r-1),f(j,r-1)]$, whose union $(f(-1,r-1),f(J-1,r-1)]$ has $n'$ elements. So every element of $T$ is the residue
of exactly one of these $x$. The $f(j,0)$, $0\leq j<J$, lie in different $\mathcal N_j$, so $J\leq|T|$.

(2) Let $x\in\mathcal N_j$. If $x<f(j,0)$, then $x\in\Lambda'_j$; otherwise $f(j,0)\leq x\leq f(j,r-1)$ and $x$ is one of
the $f(j,i)$ by condition~(3) of validity. Conversely $\Lambda'_j\subset\mathcal N_j$, and $f(j,i)\in\mathcal N_j$ as
$f(j-1,r-1)<f(j,0)$.

(3) By (1) and the periodicity, $x'=x+un'$ and $j'=j+uJ$ for an integer $u$; so
$\operatorname{dist}_{j'}(x')=\operatorname{dist}_j(x)$ and $R_{j'}=R_j$.

(4) Let $y\in\Lambda_j$. As $y\geq f(j-1,0)$, $P_j(y)\leq p^{R_j}D_j\leq a/(8Yk^2N_s)$. By
Lemma~\ref{pd-symbols-residue-top-count} and hypothesis (1) of Theorem~\ref{pd-symbols-count},
$d^D(\bar y)\geq\frac9{10}\deg(\bar y)\geq\frac9{10}a\geq\frac{36}5Yk^2N_sP_j(y)$. As $Y,N_s\geq1$, this is at least
$kP_j(y)$.

(5) The inequality $p^{R_j}D_j<p^{R_{j+1}+1}D_{j+1}$ gives $R_j+f(j,r-1)-f(j-1,0)\leq R_{j+1}+f(j+1,r-1)-f(j,0)$, and as
$f(j-1,0)<f(j,0)$ the exponent $e_j=f(j+1,r-1)-f(j,r-1)+R_{j+1}-R_j$ is positive; so $Q_j=p^{e_j}$, and
$P_j(x)Q_j=P_{j+1}(x)$ for $x\leq f(j,r-1)$. The sum of the $e_j$, $0\leq j<J$, telescopes to
$f(J,r-1)-f(0,r-1)+R_J-R_0=n'$. The last identity follows from the periodicity, as $d^D$ depends only on residues.

(6) With $j=0$, $8Yk^2N_s\leq p^{R_0}D_0\cdot8Yk^2N_s\leq a$, and $k\geq5$, $N_s\geq1$.

(7) This is a self-contained computation. Put $\theta=5/(kY)$. As $d=\max(\deg-E',0)$ and $d^D=\max(d-E^D,0)$,
$\deg(t)\leq d^D(t)+E'(t)+E^D(t)$, and $E'(t)\leq E_0$. By Step~1 of the proof of
Lemma~\ref{pd-symbols-residue-lifts-numeric}, $E_0kY\leq a$;
the proof of Lemma~\ref{pd-symbols-residue-top-count} gives $E^D(t)\leq4N_sa/(k^2Y)$, so $E^D(t)kY\leq4a$ as $N_s\leq k$;
and $a\leq\deg(t)$. So $E'(t)+E^D(t)\leq\theta\deg(t)$, and $(\deg(t)+1+l)(1-\theta)\leq d^D(t)+1+l$ for $l\geq0$. With
$M=k-1$, the product formula $\binom{y+M}M=\prod_{l<M}(y+1+l)/(l+1)$ and $(1-\theta)^M\geq1-M\theta\geq1-5/Y$ give
the claim.

(8) This is a self-contained computation. Put $P'=P_j(x)Q_j=P_{j+1}(x)$ by (5), $n=d^D(\bar x)$,
$B=\max(0,W_{j+1}(x)-m_0)$ and $\theta=1/(kY)$. As $x\in\Lambda_{j+1}$, (4) gives $\frac{36}5Yk^2N_sP'\leq n$. As
$W_{j+1}(x)=\lfloor n/P'\rfloor-(k+r+1)$ and $m_0=kN_s+1$, $P'B\geq n-(k+r+3+kN_s)P'\geq n-4kN_sP'\geq n-\theta n$.
Hence $(n+1+l)(1-\theta)\leq P'(B+1+l)$ for $l\geq0$, and as in (7), with $(1-\theta)^{k-1}\geq1-1/Y$,
$\binom{n+k-1}{k-1}(1-\frac1Y)\leq P'^{\,k-1}\binom{B+k-1}{k-1}$.
% Lean: sum_facts: (1) is prod_tiling and J_le_card_TT_of_valid; (2) is sum_split, sum_LLsub and sum_frlate; (3) is
%   sum_uniqP; (4) is sum_aP (top_count_lower_bounds_of_occurrence, l42b) and sum_kP; (5) is sum_PzQc
%   (Rj_add_first_frame_gap_le_last_frame_gap_add_Rj_succ, frame_zero_strict_mono), sum_prodQ and sum_perW; (6) is
%   thirty_two_Y_le_a (l42b at 0); (7) is sum_compA (Tj_ge, Et_le, choose_add_mul_one_sub_le_pow_mul_choose_add with
%   s = 1); (8) is sum_compW (natDiv_gt, the same comparison with s = P').
\end{proof}

\begin{lemma}\label{pd-symbols-residue-count-inputs-products}
Let $F$ be a field, $J\geq0$ an integer, and $F[z]$ and $F[X]$ the polynomial rings over $F$ in variables $z_v$,
$v\in\Sigma$, and $X_w$, $w\in\Omega$, for sets $\Sigma$ and $\Omega$. For $0\leq j<J$ let $g_j\colon\Sigma\to\Omega$
and $e_j\colon\Sigma\to\mathbf Z_{\geq0}$ be maps, and $\phi_j\colon F[z]\to F[X]$ the $F$-algebra map with
$\phi_j(z_v)=X_{g_j(v)}^{\,e_j(v)}$; it sends monomials to monomials. Let $\Phi$ be a set. For $f\in\Phi$ let
$\ell_f$ be a monomial in the $X_w$, and for $f\in\Phi$ and $0\leq j<J$ let $\mathcal S_{f,j}$ be a set of monomials
in the $z_v$ and $(u_{f,j,i})_{i\in I_{f,j}}$ a finite family of elements of $F[z]$, linearly independent over $F$,
all of whose monomials lie in $\mathcal S_{f,j}$. Suppose that for $f,f'\in\Phi$ and monomials $\mu_j\in\mathcal
S_{f,j}$ and $\mu'_j\in\mathcal S_{f',j}$, $0\leq j<J$, with
\[
 \ell_f\prod_{j<J}\phi_j(\mu_j)=\ell_{f'}\prod_{j<J}\phi_j(\mu'_j),
\]
we have $f=f'$ and $\mu_j=\mu'_j$ for every $j$. Then the polynomials $\ell_f\prod_{j<J}\phi_j(u_{f,j,i_j})$, indexed
by the pairs of an $f\in\Phi$ and a family $(i_j)_{j<J}$ with $i_j\in I_{f,j}$, are linearly independent over $F$.
\end{lemma}
\leanref{Formalism/SharpRestriction/GrowthBounds/SegmentDegree/LamProdMonomial.lean}{Formalism.MainNotions.SelectionLemmas.Pf.lam_li_all}
\leanref{Formalism/SharpRestriction/GrowthBounds/SegmentDegree/LamProdMonomial.lean}{Formalism.MainNotions.SelectionLemmas.Pf.lam_li_fixed}

\begin{proof}
For $u\in F[z]$ and a monomial $\mu$ write $u_\mu$ for the coefficient of $u$ at $\mu$. For $f\in\Phi$ let
$\mathcal M_f$ be the set of the monomials $\ell_f\prod_j\phi_j(\mu_j)$ with $\mu_j\in\mathcal S_{f,j}$. By the
hypothesis these sets are disjoint for distinct $f$, and each element of $\mathcal M_f$ determines the family
$(\mu_j)$.

\emph{Step 1: coefficients.} Let $f\in\Phi$ and let $u_j\in F[z]$, $0\leq j<J$, have all their monomials in
$\mathcal S_{f,j}$. Expanding the product,
\[
 \ell_f\prod_{j<J}\phi_j(u_j)=\sum_{(\mu_j)}\Big(\prod_{j<J}(u_j)_{\mu_j}\Big)\,\ell_f\prod_{j<J}\phi_j(\mu_j),
\]
the sum over the families of monomials $\mu_j$ of $u_j$. The monomials on the right lie in $\mathcal M_f$ and are
distinct. So every monomial of the product lies in $\mathcal M_f$, and for $\mu_j\in\mathcal S_{f,j}$ the
coefficient of the product at $\ell_f\prod_j\phi_j(\mu_j)$ is $\prod_j(u_j)_{\mu_j}$.

\emph{Step 2: one $f$.} Fix $f\in\Phi$, write $u_{j,i}=u_{f,j,i}$ and $I_j=I_{f,j}$, and let $c_a\in F$, for the
families $a=(a_j)$ with $a_j\in I_j$, satisfy $\sum_ac_a\,\ell_f\prod_j\phi_j(u_{j,a_j})=0$. By Step~1, for every
family $(\mu_j)$ of monomials
\[
 \sum_ac_a\prod_{j<J}(u_{j,a_j})_{\mu_j}=0;\tag{$*$}
\]
if some $\mu_j$ is not in $\mathcal S_{f,j}$, every product in this sum has the factor $(u_{j,a_j})_{\mu_j}=0$. Let
$b=(b_j)$ be a family. As $(u_{j,i})_{i\in I_j}$ is linearly independent, there is a linear form $\psi_j$ on $F[z]$
with $\psi_j(u_{j,b_j})=1$ and $\psi_j(u_{j,i})=0$ for $i\neq b_j$. Writing $\psi_j(u)=\sum_\mu u_\mu\psi_j(\mu)$,
\[
 c_b=\sum_ac_a\prod_{j<J}\psi_j(u_{j,a_j})=\sum_{(\mu_j)}\Big(\prod_{j<J}\psi_j(\mu_j)\Big)\sum_ac_a\prod_{j<J}
 (u_{j,a_j})_{\mu_j}=0
\]
by $(*)$, the outer sum over the families of monomials of the $u_{j,i}$. So the products with this $f$ are
linearly independent.

\emph{Step 3: all $f$.} Let a finite linear combination of the products of the statement vanish. By Step~1 the
products with a given $f$ have all their monomials in $\mathcal M_f$. As these sets are disjoint, the sum of the
terms with each $f$ vanishes, and by Step~2 all the coefficients are $0$.
% Lean: lam_li_all: Step 3 is li_graded with the sets M_f; Step 2 is lam_li_fixed, with the coefficient of Step 1
%   (lam_coeff through lam_expand and lam_prod_monomial) and the forms psi_j from lam_dual.
\end{proof}

\begin{lemma}\label{pd-symbols-residue-count-inputs-saturated}
Let $p$, $\mathcal R$, $r$, $K$, $n'$, $T$, $L_1$, $N_1$, $(J,f)$, $j$, $m$, $\theta$ and $o$ be as in
Lemma~\ref{pd-symbols-residue-segment-span}, so that $|\Lambda_j|+r\leq K$ and $m\geq0$. Let $P\in\mathcal R[z]$ and
let $c\in\mathcal R$ be nonzero. If $cP$ lies in the segment module $X_j(m,\theta,o)$, then so does $P$.
\end{lemma}
\leanref{Formalism/SharpRestriction/GrowthBounds/ValuationLowestDegreeMapFractionField.lean}{Formalism.MainNotions.G.Joints.StandardMonomialSpace.space_statement_saturated}
\leanref{Formalism/SharpRestriction/GrowthBounds/ValuationLowestDegreeMapFractionField.lean}{Formalism.MainNotions.G.Joints.StandardMonomialSpace.space_statement_linear_independent}

\begin{proof}
We use the notation of Definitions~\ref{pd-symbols-path-generators} and~\ref{pd-symbols-residue-pairs}; $v_0<\dots<
v_{N-1}$ are the elements of $\Lambda_j$, lettered as there. Let $\mathcal P$ be the set of the quadruples
$x=(\nu,\kappa,I',I'')$ of a path $\nu$ of rank $r$ for some $B'$, a path $\kappa$ of rank $K$ for some $B''$, and
standard fillings $I'$ of $\nu$ and $I''$ of $\kappa$, and let $w_x$ be
$\iota_{\mathrm{pl}}(\Psi(\nu,I'))\,\iota_{\mathrm{tn}}(\Psi(\kappa,I''))$, formed with $n=r+K$, with each $z_{i,c}$
renamed $z_{v_i,c}$. The segment generators for $(j,m,\theta,o)$ are the $w_x$ for which $(\nu,\kappa)$ is a segment
pair.

\emph{Step 1: a basis.} We show that the $w_x$, $x\in\mathcal P$, form a basis of the $\mathcal R$-module of the
polynomials in the $z_{v,c}$ with $v\in\Lambda_j$. By Lemma~\ref{pd-symbols-path-basis}, for each $B'$ the
$\Psi(\nu,I')$ with $\nu$ a path of rank $r$ for $B'$ form a basis of the polynomials of row degrees $B'$ in the
$z_{i,b}$, $b\leq r$. Every polynomial is the sum of its parts of given row degrees, so the $\Psi(\nu,I')$, over all
$B'$, form a basis of the polynomial ring in these variables; likewise the $\Psi(\kappa,I'')$ form a basis of the
polynomial ring in the $z_{i,b}$, $b\leq K$. Through $\iota_{\mathrm{pl}}$ and $\iota_{\mathrm{tn}}$ the polynomial
ring in the $z_{i,c}$, $c\leq r+K$, is the tensor product over $\mathcal R$ of these two rings, so the products
$\iota_{\mathrm{pl}}(\Psi(\nu,I'))\,\iota_{\mathrm{tn}}(\Psi(\kappa,I''))$ form a basis of it. The renaming is
injective on the variables, which gives the claim.

\emph{Step 2: saturation.} By Definition~\ref{pd-symbols-module}(2) every element of $X_j(m,\theta,o)$ is a
polynomial in the $z_{v,c}$ with $v\in\Lambda_j$. As $\mathcal R$ is a domain and $c\neq0$, $P$ has the same
monomials as $cP$, so $P$ is also such a polynomial, and by Step~1, $P=\sum_xa_xw_x$ with $a_x\in\mathcal R$. By
Lemma~\ref{pd-symbols-residue-segment-span}, $cP=\sum_xb_xw_x$ with $b_x\in\mathcal R$ and $b_x=0$ unless
$(\nu,\kappa)$ is a segment pair. By Step~1, $ca_x=b_x$ for every $x$; as $c\neq0$, $a_x=0$ unless $(\nu,\kappa)$ is a
segment pair. So $P$ lies in the span of the segment generators, which is $X_j(m,\theta,o)$ by
Lemma~\ref{pd-symbols-residue-segment-span}.
% Lean: the third clause of NestedTypeBound (Saturated), proved in space_statement_at_ring_of_basis as
%   space_statement_saturated: sat_range (the variables of P), sat_words_top with sat_psi_top and sat_prod_top (the
%   span of Step 1), and sat_li with space_statement_linear_independent (li_allDeg, li_ext and rename_injective, the
%   independence of Step 1). Both use hA2 = HyperplaneRank, the case of a DVR of Lemma pd-symbols-path-basis, and
%   the first clause h1 (Lemma pd-symbols-residue-segment-span).
\end{proof}

\begin{lemma}\label{pd-symbols-residue-count-inputs-reductions}
Let $C_{La}$ and $C_k$ be real numbers, and let $p$, $\mathcal R$, $\kappa$, $n'$, $r$, $K$, $k$, $T$, $L_1$, $N_1$,
$(J,f)$, $\deg$, $E'$, $N_s$, $E_0$, $m_*$, $a$, $Y$ and $q$ be as in Lemma~\ref{pd-symbols-residue-rank}, except that
$C_{La}$ and $C_k$ need not be at least $1088$; we use the notation of Definition~\ref{pd-symbols-residue-coarse}. Let
$j\in\mathbf Z$ and let $b\colon\mathbf Z\to\mathbf Z_{\geq0}$ vanish off $\Lambda_j$. Then the $\mathcal R$-module
$M_j(b)$ is free of finite rank, and the reductions modulo the maximal ideal of $\mathcal R$ of the elements of a
basis of $M_j(b)$ are linearly independent over $\kappa$.
\end{lemma}
\leanref{Formalism/SharpRestriction/GrowthBounds/SegmentDegree/RestLamOfRests.lean}{Formalism.MainNotions.SelectionLemmas.Pf.RestLamA}
\leanref{Formalism/SharpRestriction/GrowthBounds/SegmentDegree/LamSatRestrict.lean}{Formalism.MainNotions.SelectionLemmas.Pf.restLamA_of}
\leanref{Formalism/SharpRestriction/GrowthBounds/SegmentDegree/FiniteRestrictSupport.lean}{Formalism.MainNotions.SelectionLemmas.Pf.restFin}

\begin{proof}
We write bars for reductions modulo the maximal ideal, and let $\varpi$ be a uniformizer of $\mathcal R$.

\emph{Finite rank.} Every monomial of an element of $M_j(b)$ has degree $b(v)$ in $z_{v,1},\dots,z_{v,k}$ for every
$v$. As $b$ vanishes off the finite set $\Lambda_j$, there are finitely many such monomials, so $M_j(b)$ lies in a
free $\mathcal R$-module of finite rank. As $\mathcal R$ is Noetherian, $M_j(b)$ is finitely generated; it is
torsion-free, so it is free over the principal ideal domain $\mathcal R$.

\emph{Saturation.} Let $P\in\mathcal R[z]$ and $0\neq c\in\mathcal R$ with $cP\in M_j(b)$. As $\mathcal R$ is a domain,
$P$ and $cP$ have the same monomials, so the conditions on the monomials in
Definition~\ref{pd-symbols-residue-coarse}(3) hold for $P$. By (2) of Theorem~\ref{pd-symbols-count},
$|\Lambda_j|+r\leq N_s+r\leq K$, and $m_0\geq0$; so Lemma~\ref{pd-symbols-residue-count-inputs-saturated} gives
$P\in X_j(m_0,\theta'_{j,b},o'_j)$. Hence $P\in M_j(b)$.

\emph{Independence.} Let $y_1,\dots,y_n$ be a basis of $M_j(b)$, and let $c_i\in\mathcal R$ with $\sum_i\bar c_i\bar
y_i=0$. Then every coefficient of $\sum_ic_iy_i$ lies in $\varpi\mathcal R$, so $\sum_ic_iy_i=\varpi P$ with
$P\in\mathcal R[z]$, and $P\in M_j(b)$ by the saturation. Writing $P=\sum_id_iy_i$ with $d_i\in\mathcal R$, we get
$c_i=\varpi d_i$ for every $i$, so $\bar c_i=0$.
% Lean: restLamA_of: the finite rank is restFin (finite_restrictSupport' with segDegSet_finite) and
%   Module.free_of_finite_type_torsion_free'; the saturation is hsatX, the third clause of NestedTypeBound (Lemma
%   pd-symbols-residue-count-inputs-saturated), with lam_sat_restrict for the two support conditions; the
%   independence uses Module.finBasis and lam_dvd_of_res.
\end{proof}

\begin{lemma}\label{pd-symbols-residue-count-inputs-old}
Let $C_{La}\geq1088$ and $C_k\geq1088$ be real numbers, and let $p$, $\mathcal R$, $\kappa$, $n'$, $r$, $K$, $k$, $T$,
$L_1$, $N_1$, $(J,f)$, $\deg$, $E'$, $N_s$, $E_0$, $m_*$, $a$, $Y$ and $q$ be as in
Lemma~\ref{pd-symbols-residue-rank}; we use the notation of Definitions~\ref{pd-symbols-residue-coarse}
and~\ref{pd-symbols-residue-families}, and write $\mathcal N_j$ for the set of the new elements of $\Lambda_j$. Let
$0\leq j,j'<J$ and let $x\in\Lambda_j$ not be new. Then:
\begin{enumerate}
\item $x=f(j-1,i)$ for some $i<r$, and $x\in\mathcal N_{j-1}$;
\item if $x'\in\Lambda_{j'}$ is not new and $\bar x'=\bar x$, then $j'=j$ and $x'=x$;
\item if $x'\in\mathcal N_{j'}$ and $\bar x'=\bar x$, then $x'=f(j',i)$ for the $i$ of (1), and
  $P_j(x)=P_{j'}(x')Q_{j'}$.
\end{enumerate}
\end{lemma}
\leanref{Formalism/SharpRestriction/GrowthBounds/SegmentDegree/LamImgApply.lean}{Formalism.MainNotions.SelectionLemmas.Pf.LamGeo}
\leanref{Formalism/SharpRestriction/GrowthBounds/SegmentDegree/LamMZero.lean}{Formalism.MainNotions.SelectionLemmas.Pf.restLamGeo}
\leanref{Formalism/SharpRestriction/GrowthBounds/SegmentDegree/LamMZero.lean}{Formalism.MainNotions.SelectionLemmas.Pf.lam_W_late}

\begin{proof}
(1) As in the proof of Lemma~\ref{pd-symbols-normal-degree}, $x=f(j-1,i)$ for some $i<r$. Then
$f(j-2,r-1)<f(j-1,0)\leq x\leq f(j-1,r-1)$ and $\bar x\in T$, so $x\in\mathcal N_{j-1}$.

(2) By (1), $x\in\mathcal N_{j-1}$ and $x'\in\mathcal N_{j'-1}$. For $0\leq l<J$ the set $\mathcal N_{l-1}$ consists of
the integers with residue in $T$ in the interval $(f(l-2,r-1),f(l-1,r-1)]$. These intervals are consecutive for
$0\leq l<J$, and by condition~(1) of validity their union $(f(-2,r-1),f(J-2,r-1)]$ has $n'$ elements; so distinct
integers in it have distinct residues. Hence $x=x'$, and then $j=j'$, as the intervals are disjoint.

(3) Let $j_1=j-1$ and $\tilde x=x$ if $j\geq1$, and $j_1=J-1$ and $\tilde x=x+n'$ if $j=0$. Then $0\leq j_1<J$, and by
(1) and condition~(1) of validity, $\tilde x=f(j_1,i)\in\mathcal N_{j_1}$. As $\tilde x$ and $x'$ have the same
residue, Lemma~\ref{pd-symbols-residue-count-sum-facts}(1) gives $j'=j_1$ and $x'=\tilde x=f(j',i)$.

It remains to compute the exponents of $P_j(x)$ and $P_{j'}(x')Q_{j'}$. Put $s=x'-x$, so $s=0$ if $j\geq1$ and $s=n'$
if $j=0$; by condition~(1) of validity, $f(j'+\epsilon,t)=f(j-1+\epsilon,t)+s$ for $\epsilon\in\{-1,0,1\}$ and $t<r$.
Let $L\geq0$ be the largest integer with $p^L\leq\max\big(1,\lfloor a/(8Yk^2N_s)\rfloor\big)$, and for $l\in\mathbf Z$
put $d_l=\operatorname{dist}_l(f(l-1,0))$. For an integer $N\geq0$ and $d\geq0$, the largest $e\geq0$ with
$p^e\leq\max(1,\lfloor N/p^d\rfloor)$ is $\max(e_N-d,0)$, where $e_N$ is the one for $d=0$; and
$\lfloor\lfloor A\rfloor/p^d\rfloor=\lfloor A/p^d\rfloor$ for a real $A\geq0$. So $R_l=\max(L-d_l,0)$ for every $l$.
By validity, $f(j-2,0)<f(j-1,0)\leq f(j-1,r-1)<f(j,r-1)$, and $x\leq f(j-1,r-1)$ by (1). Hence
\[
 \operatorname{dist}_j(x)=f(j,r-1)-x,\quad d_j=f(j,r-1)-f(j-1,0),\quad
 \operatorname{dist}_{j-1}(x)=f(j-1,r-1)-x,\quad d_{j-1}=f(j-1,r-1)-f(j-2,0),
\]
and, shifting by $s$, $\operatorname{dist}_{j'}(x')=\operatorname{dist}_{j-1}(x)$, $d_{j'}=d_{j-1}$ and $d_{j'+1}=d_j$, so
$R_{j'}=R_{j-1}$ and $R_{j'+1}=R_j$. Put $\Delta=f(j,r-1)-f(j-1,r-1)>0$; then also $f(j'+1,r-1)-f(j',r-1)=\Delta$. As
$d_j-d_{j-1}=\Delta-\big(f(j-1,0)-f(j-2,0)\big)<\Delta$, we get $R_{j-1}\leq\max(L-d_j+\Delta,0)\leq R_j+\Delta$. So
the exponent of $Q_{j'}$ is $\Delta+R_j-R_{j-1}\geq0$, and the exponent of $P_{j'}(x')Q_{j'}$ is
\[
 \big(f(j-1,r-1)-x\big)+R_{j-1}+\Delta+R_j-R_{j-1}=f(j,r-1)-x+R_j ,
\]
the exponent of $P_j(x)$.
% Lean: restLamGeo: (1) is lam_W_late; (2) is the clause U2 (eq_of_mod_eq_of_mem_Ioc_fr_last with s = -1); (3) is
%   the clause LK (eq_of_mod_eq_of_mem_Ioc_fr_last, shifting by n' when j = 0, and Pz_eq_Pz_mul_Qc with s = 0 or n'
%   and the four facts hTT, hA, hB, hC of validity; it computes the exponents from Rj_eq_log_floor, which is
%   Nat.floor_div_natCast and nat_log_div_pow_eq_sub, and closes with omega). The clauses U1 and EX of LamGeo are (1)
%   of Lemma pd-symbols-residue-count-sum-facts. restLamGeo uses only the validity of the run system.
\end{proof}

\begin{lemma}\label{pd-symbols-residue-count-inputs-monomials}
Let $C_{La}\geq1088$ and $C_k\geq1088$ be real numbers, and let $p$, $\mathcal R$, $\kappa$, $n'$, $r$, $K$, $k$, $T$,
$L_1$, $N_1$, $(J,f)$, $\deg$, $E'$, $N_s$, $E_0$, $m_*$, $a$, $Y$ and $q$ be as in
Lemma~\ref{pd-symbols-residue-rank}; we use the notation of Definitions~\ref{pd-symbols-residue-coarse}
and~\ref{pd-symbols-residue-families}, write $\mathcal N_j$ for the set of the new elements of $\Lambda_j$, and
$\mu(v,c)$ for the exponent of $z_{v,c}$ in a monomial $\mu$ in the $z_{v,c}$. For $\alpha\in\mathcal L$,
$e\in E(\alpha)$ and $0\leq j<J$ let $\mathcal S_j(\alpha,e)$ be the set of the monomials $\mu$ in the variables
$z_{v,c}$ with $v\in\Lambda_j$, and $c\leq r$ unless $v$ is new, such that $\sum_c\mu(v,c)=b^{\alpha,e}_j(v)$ for every
$v$ and $\mu(f(j,i),c)<Q_j$ for all $i<r$ and $c\leq r$; and let $X^\alpha=\prod_{t,c}X_{t,c}^{\,\alpha_{t,c}}$. For a
family $\mu=(\mu_j)_{0\leq j<J}$ of monomials, $\mu_j$ in the $z_{v,c}$ with $v\in\Lambda_j$, and for
$t\in\mathbf Z/n'\mathbf Z$ and $c\leq k$, let $w_\mu(t,c)$ be the sum of the $\mu_j(x,c)$ over the $0\leq j<J$ and the
$x\in\Lambda_j\smallsetminus\mathcal N_j$ with $\bar x=t$. Then:
\begin{enumerate}
\item for such a family, $0\leq j_0<J$, $x_0\in\mathcal N_{j_0}$ and $c\leq k$, the exponent of $X_{\bar x_0,c}$ in
  $\prod_{j<J}\psi_j(\mu_j)$ is $P_{j_0}(x_0)\big(\mu_{j_0}(x_0,c)+Q_{j_0}w_\mu(\bar x_0,c)\big)$, and no $X_{s,c}$ with
  $s\notin T$ occurs in it;
\item let $\alpha,\alpha'\in\mathcal L$, $e\in E(\alpha)$, $e'\in E(\alpha')$, and $\mu_j\in\mathcal S_j(\alpha,e)$ and
  $\mu'_j\in\mathcal S_j(\alpha',e')$ for $0\leq j<J$; if $X^\alpha\prod_{j<J}\psi_j(\mu_j)=X^{\alpha'}\prod_{j<J}\psi_j(\mu'_j)$,
  then $\alpha=\alpha'$, $e=e'$ and $\mu_j=\mu'_j$ for every $j$.
\end{enumerate}
\end{lemma}
\leanref{Formalism/SharpRestriction/GrowthBounds/SegmentDegree/LamImgApply.lean}{Formalism.MainNotions.SelectionLemmas.Pf.lam_eval}
\leanref{Formalism/SharpRestriction/GrowthBounds/SegmentDegree/RestLamOfRests.lean}{Formalism.MainNotions.SelectionLemmas.Pf.RestLamInj}
\leanref{Formalism/SharpRestriction/GrowthBounds/SegmentDegree/LamImgApply.lean}{Formalism.MainNotions.SelectionLemmas.Pf.restLamInj_of}
\leanref{Formalism/SharpRestriction/GrowthBounds/SegmentDegree/LamMZero.lean}{Formalism.MainNotions.SelectionLemmas.Pf.restLamTau}

\begin{proof}
(1) As $\psi_j(z_{x,c})=X_{\bar x,c}^{\,P_j(x)}$ and $\mu_j$ involves only rows $z_x$ with $x\in\Lambda_j$, the exponent
of $X_{t,c}$ in $\prod_j\psi_j(\mu_j)$ is $\sum_{j<J}\sum_xP_j(x)\mu_j(x,c)$, the inner sum over the $x\in\Lambda_j$
with $\bar x=t$. All residues of elements of $\Lambda_j$ lie in $T$, which gives the second claim. Let
$t=\bar x_0$. By Lemma~\ref{pd-symbols-residue-count-sum-facts}(1), the only new $x$ in this sum is $x_0$, with
$j=j_0$. For $x\in\Lambda_j\smallsetminus\mathcal N_j$ with $\bar x=t$,
Lemma~\ref{pd-symbols-residue-count-inputs-old}(3) gives $P_j(x)=P_{j_0}(x_0)Q_{j_0}$. This gives the formula.

(2) \emph{Step 1: at a new element.} Let $0\leq j_0<J$, $x_0\in\mathcal N_{j_0}$, $t=\bar x_0$ and $c\leq k$. We show
that $\alpha_{t,c}=\alpha'_{t,c}$, $\mu_{j_0}(x_0,c)=\mu'_{j_0}(x_0,c)$ and $w_\mu(t,c)=w_{\mu'}(t,c)$. Put
$P=P_{j_0}(x_0)$ and $Q=Q_{j_0}$. Comparing the exponents of $X_{t,c}$, (1) gives
\[
 \alpha_{t,c}+P\big(\mu_{j_0}(x_0,c)+Qw_\mu(t,c)\big)=\alpha'_{t,c}+P\big(\mu'_{j_0}(x_0,c)+Qw_{\mu'}(t,c)\big).
\]
As $\alpha_t$ and $\alpha'_t$ are remainder vectors at $t$ and $x_0$ is new, $\alpha_{t,c}<P$ and $\alpha'_{t,c}<P$; so
$\alpha_{t,c}=\alpha'_{t,c}$ and $\mu_{j_0}(x_0,c)+Qw_\mu(t,c)=\mu'_{j_0}(x_0,c)+Qw_{\mu'}(t,c)$. If $x_0=f(j_0,i)$ for
some $i<r$ and $c\leq r$, then $\mu_{j_0}(x_0,c)$ and $\mu'_{j_0}(x_0,c)$ are less than $Q$ by the definition of
$\mathcal S_{j_0}$, so they agree, and then so do $w_\mu(t,c)$ and $w_{\mu'}(t,c)$. Otherwise
$w_\mu(t,c)=w_{\mu'}(t,c)=0$: if $x\in\Lambda_j\smallsetminus\mathcal N_j$ has $\bar x=t$, then $x_0=f(j_0,i)$ for some
$i$ by Lemma~\ref{pd-symbols-residue-count-inputs-old}(3), so $c>r$, and $\mu_j(x,c)=\mu'_j(x,c)=0$, as the monomials
in $\mathcal S_j$ involve no normal variable in a row that is not new. Then $\mu_{j_0}(x_0,c)=\mu'_{j_0}(x_0,c)$.

\emph{Step 2: $\alpha=\alpha'$.} Let $t\in\mathbf Z/n'\mathbf Z$. If $t\in T$, then $t$ is the residue of an
$x_0\in\mathcal N_{j_0}$ with $0\leq j_0<J$, by Lemma~\ref{pd-symbols-residue-count-sum-facts}(1), and Step~1 gives
$\alpha_t=\alpha'_t$. If $t\notin T$, then $d(t)=0$, so $d^D(t)=0$, and $\sum_c\alpha_{t,c}\leq d^D(t)$ gives $\alpha_t=0$;
likewise $\alpha'_t=0$.

\emph{Step 3: $\mu=\mu'$.} Let $0\leq j<J$, $v\in\Lambda_j$ and $c\leq k$. If $v$ is new, Step~1 gives
$\mu_j(v,c)=\mu'_j(v,c)$. If not, then $\bar v\in T$ is the residue of some $x_0\in\mathcal N_{j_0}$, $0\leq j_0<J$, by
Lemma~\ref{pd-symbols-residue-count-sum-facts}(1), and by Lemma~\ref{pd-symbols-residue-count-inputs-old}(2), $v$ is
the only element of the sets $\Lambda_{j'}\smallsetminus\mathcal N_{j'}$, $0\leq j'<J$, with residue $\bar v$; so
$w_\mu(\bar v,c)=\mu_j(v,c)$ and $w_{\mu'}(\bar v,c)=\mu'_j(v,c)$, and Step~1 gives $\mu_j(v,c)=\mu'_j(v,c)$. The
variables $z_{v,c}$ with $v\notin\Lambda_j$ occur in neither $\mu_j$ nor $\mu'_j$.

\emph{Step 4: $e=e'$.} Let $0\leq j<J$ and $i<r$. The element $f(j,i)$ of $\Lambda_j$ is new and satisfies
$f(j,0)\leq f(j,i)\leq f(j,r-1)$, so $b^{\alpha,e}_j(f(j,i))=e_j(i)$ and $b^{\alpha',e'}_j(f(j,i))=e'_j(i)$ by
Definition~\ref{pd-symbols-residue-families}(3). As $\mu_j=\mu'_j$ lies in $\mathcal S_j(\alpha,e)$ and in $\mathcal
S_j(\alpha',e')$, $e_j(i)=\sum_c\mu_j(f(j,i),c)=e'_j(i)$.
% Lean: (1) is lam_eval (lam_img_apply, with the clauses U1 and LK of LamGeo); (2) is restLamInj = restLamInj_of
%   (restLamGeo 1088) restLamTau: Step 1 is hdec (lam_digit, the bound of resSet, and hw0 with okVars), Step 2 is
%   hlaeq (the clause EX), Step 3 is hmeq (the clause U2) and Step 4 is restLamTau with segDegSet. The book's
%   S_j(alpha, e) is lamS, X^alpha is lamL la, w_mu is segmentWeightExponent and psi_j(mu) is imgMon (lamG j) (lamE j).
\end{proof}

\begin{definition}\label{pd-symbols-residue-count-inputs-index}
Let the data of Definition~\ref{pd-symbols-residue-coarse} be given, with that notation. Let $j\in\mathbf Z$, let
$v_0<\dots<v_{N-1}$ be the elements of $\Lambda_j$, lettered as in Definition~\ref{pd-symbols-residue-pairs}, and let
$b\colon\mathbf Z\to\mathbf Z_{\geq0}$. For $x\in\Lambda_j$ put $\beta_j(b,x)=\max(W_j(x)-m_0,0)$ if $x$ is not new and
$\beta_j(b,x)=b(x)$ if it is, and let $\beta_j(b)$ be the sum of the $\beta_j(b,x)$ over the $x\in\Lambda_j$ with
$x<f(j,0)$. We use the paths, standard fillings, $\Psi$, $\iota_{\mathrm{pl}}$, $\iota_{\mathrm{tn}}$, $\mu^{\mathrm{pl}}$
and $\mu^{\mathrm{tn}}$ of Definition~\ref{pd-symbols-path-generators} for this lettering, and the canonical fillings
of Definition~\ref{pd-symbols-path-canonical}(2).
\begin{enumerate}
\item A \emph{counting path} for $(j,b)$ is a path $\mu$ of rank $k$ for some $B$ such that: $B(i)=\beta_j(b,v_i)$ for
  every letter $i$ with $v_i<f(j,0)$; $\mu_{i+1}(m)=\beta_j(b)$ for $1\leq m\leq i'+1$ whenever $v_i=f(j,i')$ with
  $i'<r$; and for every letter $i$ with $v_i$ new, $\theta'_{j,b}(v_i)\leq\mu_i(r)$, and
  $\sum_{c=r+1}^k(\mu_{i+1}(c)-\mu_i(c))\leq o'_j(v_i)$ if $v_i$ is even, $\sum_{c=1}^r(\mu_i(c)-\mu_{i+1}(c))\leq
  o'_j(v_i)$ if $v_i$ is odd. Let $T_j(b)$ be the set of the pairs $(\mu,I'')$ of a counting path $\mu$ for $(j,b)$ and
  a standard filling $I''$ of $\mu^{\mathrm{tn}}$.
\item Let $\mathcal I_j(b)$ be the set of the triples $(\mu,I'',u)$ of a pair $(\mu,I'')\in T_j(b)$, with $\mu$ a path
  for $B$, and a function $u\colon\{0,\dots,r-1\}\times\{1,\dots,r\}\to\mathbf Z_{\geq0}$ such that
  $u(i',c)+2\beta_j(b)<Q_j$ for all $(i',c)$, and $B(i)+\sum_cu(i',c)=b(f(j,i'))$ whenever $v_i=f(j,i')$.
\item For $(\mu,I'',u)\in\mathcal I_j(b)$ let $y_j(b;\mu,I'',u)\in\mathcal R[z]$ be the product of
  \[
   D_j^{\,m_0}\prod_{i<r}z_{f(j-1,i),1}^{\,\max(b(f(j-1,i))-W_j(f(j-1,i)),0)}\prod_{i'<r,\,c\leq r}z_{f(j,i'),c}^{\,u(i',c)},
   \qquad D_j=\det\big(z_{f(j-1,a),c}\big)_{0\leq a<r,\,1\leq c\leq r},
  \]
  and of $\iota_{\mathrm{pl}}(\Psi(\mu^{\mathrm{pl}},\operatorname{can}_{\nu}))\,\iota_{\mathrm{tn}}(\Psi(\mu^{\mathrm{tn}},I''))$, with
  $\nu=\mu^{\mathrm{pl}}_N$ and each $z_{i,c}$ renamed $z_{v_i,c}$.
\end{enumerate}
For a set $S$ write $|S|$ for its number of elements if it is finite and $0$ otherwise.
\end{definition}
\leanref{Formalism/SharpRestriction/GrowthBounds/SegmentDegree/GoodPathWeight.lean}{Formalism.MainNotions.SelectionLemmas.Pf.Dat.Bdeg}
\leanref{Formalism/SharpRestriction/GrowthBounds/SegmentDegree/GoodPathWeight.lean}{Formalism.MainNotions.SelectionLemmas.Pf.Dat.beS}
\leanref{Formalism/SharpRestriction/GrowthBounds/SegmentDegree/GoodPathWeight.lean}{Formalism.MainNotions.SelectionLemmas.Pf.Dat.IsPATH}
\leanref{Formalism/SharpRestriction/GrowthBounds/SegmentDegree/GoodPathWeight.lean}{Formalism.MainNotions.SelectionLemmas.Pf.Dat.IsGood}
\leanref{Formalism/SharpRestriction/GrowthBounds/SegmentDegree/GoodPathWeight.lean}{Formalism.MainNotions.SelectionLemmas.Pf.Dat.Tset}
\leanref{Formalism/SharpRestriction/GrowthBounds/SegmentDegree/GoodPathWeight.lean}{Formalism.MainNotions.SelectionLemmas.Pf.Dat.Idx}
\leanref{Formalism/SharpRestriction/GrowthBounds/SegmentDegree/GoodPathWeight.lean}{Formalism.MainNotions.SelectionLemmas.Pf.Dat.Dw}
\leanref{Formalism/SharpRestriction/GrowthBounds/SegmentDegree/GoodPathWeight.lean}{Formalism.MainNotions.SelectionLemmas.Pf.Dat.famEl}
\leanref{Formalism/SharpRestriction/GrowthBounds/SegmentDegree/GoodPathWeight.lean}{Formalism.MainNotions.SelectionLemmas.Pf.Dat.restr}

\begin{lemma}\label{pd-symbols-residue-count-inputs-ratio}
Let $C_{La}\geq1088$ and $C_k\geq1088$ be real numbers, and let $p$, $\mathcal R$, $\kappa$, $n'$, $r$, $K$, $k$, $T$,
$L_1$, $N_1$, $(J,f)$, $\deg$, $E'$, $N_s$, $E_0$, $m_*$, $a$, $Y$ and $q$ be as in
Lemma~\ref{pd-symbols-residue-rank}; we use the notation of Definitions~\ref{pd-symbols-residue-coarse},
\ref{pd-symbols-residue-families} and~\ref{pd-symbols-residue-count-inputs-index}. Let $j\in\mathbf Z$ and let
$b\colon\mathbf Z\to\mathbf Z_{\geq0}$ be balanced at $j$. Then $4Y\beta_j(b)\leq Q_j$.
\end{lemma}
\leanref{Formalism/SharpRestriction/GrowthBounds/SegmentDegree/TauDegreeOnPathBound.lean}{Formalism.MainNotions.SelectionLemmas.Pf.BetaRatioBound}
\leanref{Formalism/SharpRestriction/GrowthBounds/SegmentDegree/BetaRatioBound.lean}{Formalism.MainNotions.SelectionLemmas.Pf.beta_ratio_bound}
\leanref{Formalism/ThetaLogCoordinates/ResidueClassEstimates/VandermondeRankEstimates/AdmissibleSmallRatio.lean}{Formalism.InvariantFactorPaths.BranchingDimensions.admissibleSmallRatio}
\leanref{Formalism/ThetaLogCoordinates/ResidueClassEstimates/VandermondeRankEstimates/AdmissibleSmallRatio.lean}{Formalism.InvariantFactorPaths.BranchingDimensions.geometricSumDistinct}

\begin{proof}
We use the notation of Definition~\ref{pd-symbols-run-system}. As in Step~1 of the proof of
Lemma~\ref{pd-symbols-residue-lifts-numeric}, $Y\geq1$, $k\geq5$, $N_s\geq1$ and $a>0$, and for every $j'$,
$T_{j'}<p^{R_{j'}+1}$ with $T_{j'}=a/(8Yk^2N_s\,p^{\operatorname{dist}_{j'}(f(j'-1,0))})$. By (4) of
Theorem~\ref{pd-symbols-count}, $128k^2N_sY^2\leq p^{L_1-1}$; the left side is at least $1$, so $L_1\geq1$. Put
$D_1=\operatorname{dist}_j(f(j,0))$.

\emph{Step 1: $\beta_j(b)\leq4a/p^{R_j+D_1+L_1+1}$.} Let $S$ be the set of the $x\in\Lambda_j$ with $x<f(j,0)$. By
condition~(4) of validity, $f(j,0)$ is $L_1$-isolated in $\Lambda_j$; so $\operatorname{dist}_j(x)\geq D_1+L_1+1$ for
$x\in S$, and distinct $x$ have distinct $\operatorname{dist}_j(x)$. Let $x\in S$ and $P=p^{\operatorname{dist}_j(x)+R_j}$.
If $x$ is not new, $\beta_j(b,x)=\max(W_j(x)-m_0,0)\leq d^D(\bar x)/P$. If $x$ is new, then $x\in\Lambda'_j$, and
$\beta_j(b,x)=b(x)\leq d^D(\bar x)/P$ as $b$ is balanced. Also $d^D(\bar x)\leq\deg(\bar x)\leq2a$ by (1) of
Theorem~\ref{pd-symbols-count}. Hence
\[
 \beta_j(b)\leq\frac{2a}{p^{R_j}}\sum_{x\in S}p^{-\operatorname{dist}_j(x)}\leq\frac{2a}{p^{R_j}}\cdot\frac2{p^{D_1+L_1+1}},
\]
as the $\operatorname{dist}_j(x)$ are distinct integers at least $D_1+L_1+1$, and $\sum_{e\geq e_0}p^{-e}\leq2p^{-e_0}$ for
$p\geq2$.

\emph{Step 2: comparison with $Q_j$.} Put $\delta=f(j+1,r-1)-f(j,r-1)\geq1$. Then
$\operatorname{dist}_{j+1}(f(j,0))=D_1+\delta$, so $T_{j+1}<p^{R_{j+1}+1}$ gives
$a<8pYk^2N_s\,p^{D_1+\delta+R_{j+1}}$. By (5) of Lemma~\ref{pd-symbols-residue-count-sum-facts} and its proof, the
exponent $\delta+R_{j+1}-R_j$ is positive and $Q_j=p^{\delta+R_{j+1}-R_j}$. With Step~1,
\[
 \frac{\beta_j(b)}{Q_j}\leq\frac{4a}{p^{D_1+L_1+1+\delta+R_{j+1}}}<\frac{32Yk^2N_s}{p^{L_1}}\leq\frac1{4pY}\leq\frac1{4Y},
\]
where the third inequality is $128k^2N_sY^2\,p\leq p^{L_1}$.
% Lean: beta_ratio_bound: hL1; hU1 (the gap at f(j, 0) from validity); hA1 (TE.Bdeg_bounds, TE.c_bounds,
%   geometricSumDistinct); hB1 (lt_pow_Rj at j + 1, hdist); hQ (with Rj_add_first_frame_gap_le_last_frame_gap_add_Rj_succ);
%   the inequalities of Step 2 are admissibleSmallRatio.
\end{proof}

\begin{lemma}\label{pd-symbols-residue-count-inputs-rows}
Let $C_{La}$ and $C_k$ be real numbers, and let $p$, $\mathcal R$, $\kappa$, $n'$, $r$, $K$, $k$, $T$, $L_1$, $N_1$,
$(J,f)$, $\deg$, $E'$, $N_s$, $E_0$, $m_*$, $a$, $Y$ and $q$ be as in Lemma~\ref{pd-symbols-residue-rank}, except that
$C_{La}$ and $C_k$ need not be at least $1088$; we use the notation of
Definition~\ref{pd-symbols-residue-count-inputs-index}. Let $j\in\mathbf Z$, $b\colon\mathbf Z\to\mathbf Z_{\geq0}$, and
let $\mu$ be a counting path for $(j,b)$, a path for $B$. Then $B(i)\leq2\beta_j(b)$ for every letter $i$ with
$v_i\geq f(j,0)$.
\end{lemma}
\leanref{Formalism/SharpRestriction/GrowthBounds/SegmentDegree/TauDegreeOnPathBound.lean}{Formalism.MainNotions.SelectionLemmas.Pf.TauDegreeOnPathBound}
\leanref{Formalism/SharpRestriction/GrowthBounds/SegmentDegree/DminusLeBeS.lean}{Formalism.MainNotions.SelectionLemmas.Pf.tau_degree_on_path_bound}
\leanref{Formalism/SharpRestriction/GrowthBounds/SegmentDegree/DminusLeBeS.lean}{Formalism.MainNotions.SelectionLemmas.Pf.dminus_le_beS}

\begin{proof}
Let $v_i\geq f(j,0)$. By Step~2 of the proof of Lemma~\ref{pd-symbols-residue-lifts-numeric}, the elements of
$\Lambda_j$ that are at least $f(j,0)$ are the $f(j,i')$, $i'<r$; so $v_i=f(j,i')$, which is even, and $i$ is a vector
letter. As $\mu$ is a path, $\mu_i(a)\leq\mu_{i+1}(a)$ for all $a$, $\mu_{i+1}(a+1)\leq\mu_i(a)$ for $a<k$, and
$\sum_a(\mu_{i+1}(a)-\mu_i(a))=B(i)$. Hence
\[
 B(i)=\mu_{i+1}(1)-\mu_i(k)+\sum_{a<k}\big(\mu_{i+1}(a+1)-\mu_i(a)\big)\leq\mu_{i+1}(1)-\mu_i(k),
\]
and $\mu_{i+1}(1)=\beta_j(b)$ by the definition of a counting path. For a letter $l$, $\mu_{l+1}(a)\geq\mu_l(a)$ if $l$ is a
vector letter, and $\mu_{l+1}(a)\geq\mu_l(a)-B(l)$ if it is a covector letter, as $\mu_l$ is then a strip of size $B(l)$
over $\mu_{l+1}$. As $\mu_0=0$, $\mu_i(k)$ is at least minus the sum of the $B(l)$ over the covector letters $l<i$. For
these, $v_l$ is odd, so $v_l<f(j,0)$ and $B(l)=\beta_j(b,v_l)$; so this sum is at most $\beta_j(b)$, and
$B(i)\leq2\beta_j(b)$.
% Lean: tau_degree_on_path_bound (segE_eq_fr_iff, step_vec, tele_le for the telescoping, head for mu_(i+1)(1),
%   neg_dminus_le_mu for the lower bound of mu_i(k) by minus dminus, and dminus_le_beS).
\end{proof}

\begin{lemma}\label{pd-symbols-residue-count-inputs-residues}
Let $r\geq2$ and $Q,\beta,c,\rho\geq0$ be integers with $2\beta<Q$. Then the number of the
$w\in\{0,\dots,Q-2\beta-1\}^r$ with $c+\sum_iw_i\equiv\rho\pmod Q$ is at least $Q^{r-1}(1-2r\beta/Q)$.
\end{lemma}
\leanref{Formalism/SharpRestriction/GrowthBounds/SegmentDegree/TauDegreeOnPathBound.lean}{Formalism.MainNotions.SelectionLemmas.Pf.ResidueClassCount}
\leanref{Formalism/SharpRestriction/GrowthBounds/SegmentDegree/DminusLeBeS.lean}{Formalism.MainNotions.SelectionLemmas.Pf.residue_class_count}
\leanref{Formalism/ThetaLogCoordinates/ResidueClassEstimates/VandermondeRankEstimates/LDigits.lean}{Formalism.InvariantFactorPaths.BranchingDimensions.cnt_ge}

\begin{proof}
This is a self-contained computation. Put $t=Q-2\beta$, and let $N_r(s)$ be the number of $w\in\{0,\dots,t-1\}^r$ with
$\sum_iw_i\equiv s\pmod Q$. For $r=2$: for each $w_1$ exactly one $w_2\in\{0,\dots,Q-1\}$ satisfies the congruence, and it
fails to be less than $t$ for at most $Q-t$ values of $w_1$, as distinct $w_1$ give distinct $w_2$; so $N_2(s)\geq
t-(Q-t)$. Summing over the first coordinate, $N_r(s)\geq t^{r-2}(t-(Q-t))=t^{r-2}(Q-4\beta)$. If $4\beta>Q$, the bound of the
statement is negative. Otherwise, induction on $m$ gives $q^m(q-(m+2)b)\leq(q-b)^m(q-2b)$ for $0\leq b\leq q$, and with
$q=Q$, $b=2\beta$ and $m=r-2$ this is $Q^{r-1}(1-2r\beta/Q)\leq t^{r-2}(Q-4\beta)$. The count of the statement is
$N_r(\rho-c)$.
% Lean: residue_class_count = BranchingDimensions.residueClassLowerBound (cnt_two, cnt_ge and
%   pow_mul_sub_natCast_add_two_mul_le_sub_pow_mul), with the shift of the residue by c.
\end{proof}

\begin{lemma}\label{pd-symbols-residue-count-inputs-segment-clauses}
For $n\geq1$ and a dominant $\nu\in\mathbf Z^n$ let $\operatorname{st}(\nu)$ be the number of standard fillings of shape
$\nu$ (Definition~\ref{pd-symbols-path-canonical}(1)). For $\nu,\nu'\in\mathbf Z^n$ and $b\geq0$, $\nu'$ is a \emph{strip
of size $b$ over} $\nu$ if $\nu(a)\leq\nu'(a)$ for all $a$, $\nu'(a+1)\leq\nu(a)$ for $a<n$, and
$\sum_a(\nu'(a)-\nu(a))=b$.
\begin{enumerate}
\item Let $R$ be a discrete valuation ring, $\Lambda$ a set each of whose elements is a vector letter or a covector
  letter, $n\geq1$ and $S\subset\Lambda$, and form $\operatorname{val}_S$ and $\lambda_S$ of the polynomials over $R$ in
  the $z_{v,c}$, $v\in\Lambda$, $1\leq c\leq n$, as in Definition~\ref{pd-symbols-path-generators}(1). Then for all such
  $F$ and $G$, with $+\infty+x=+\infty$ and $-\infty+x=-\infty$,
  \[
   \operatorname{val}_S(FG)=\operatorname{val}_S(F)+\operatorname{val}_S(G),\qquad
   \min\big(\operatorname{val}_S(F),\operatorname{val}_S(G)\big)\leq\operatorname{val}_S(F+G),
  \]
  \[
   \lambda_S(FG)=\lambda_S(F)+\lambda_S(G),\qquad\lambda_S(F+G)\leq\max\big(\lambda_S(F),\lambda_S(G)\big);
  \]
  for $v\in S$ and $c\leq n$, $0\leq\operatorname{val}_S(z_{v,c})$ and $\lambda_S(z_{v,c})\leq1$ if $v$ is a vector letter,
  and $-1\leq\operatorname{val}_S(z_{v,c})$ and $\lambda_S(z_{v,c})\leq0$ if it is a covector letter; and
  $\operatorname{val}_S\big(\det(z_{y_a,b})_{1\leq a,b\leq n}\big)=1$ for distinct vector letters $y_1,\dots,y_n\in S$.
\item Let $N\geq0$ and $n\geq1$ be integers and $\mu=(\mu_0,\dots,\mu_N)$ a sequence in $\mathbf Z^n$ with $\mu_N$
  dominant. Then the number of standard fillings of $\mu$ (Definition~\ref{pd-symbols-path-generators}(4)) is
  \[
   \prod_{1\leq a<b\leq n}\frac{\mu_N(a)-\mu_N(b)+b-a}{b-a}.
  \]
\item Let $m\geq0$ and let $\nu\in\mathbf Z^{m+1}$ be dominant. If $m=0$, then $\operatorname{st}(\nu)=1$. If $m\geq1$,
  then $\operatorname{st}(\nu)=\sum_\lambda\operatorname{st}(\lambda)$, the sum over the $\lambda\in\mathbf Z^m$ with
  $\nu(a+1)\leq\lambda(a)\leq\nu(a)$ for $1\leq a\leq m$.
\item Let $m,b\geq0$ and let $\nu\in\mathbf Z^{m+1}$ be dominant. Then $\operatorname{st}(\nu)\binom{b+m}m$ is the sum
  of the $\operatorname{st}(\nu')$ over the strips $\nu'$ of size $b$ over $\nu$, and it is also the sum of the
  $\operatorname{st}(\nu')$ over the $\nu'$ over which $\nu$ is a strip of size $b$.
\item Let $p$ be a prime, $\mathcal R$ a complete discrete valuation ring in which $p$ is neither $0$ nor a unit,
  $r,K\geq1$ integers, $n'\geq2$ an even integer, $T\subset\mathbf Z/n'\mathbf Z$, $L_1,N_1\geq0$ integers and
  $(J,f)$ a run system for $T$ of length $r$ valid for $(L_1,N_1)$. Let $j\in\mathbf Z$ with $|\Lambda_j|+r\leq K$, and
  let $v_0<\dots<v_{N-1}$ be the elements of $\Lambda_j$, lettered as in Definition~\ref{pd-symbols-residue-pairs}.
  Then the polynomials obtained from $\iota_{\mathrm{pl}}(\Psi(\nu,I'))\,\iota_{\mathrm{tn}}(\Psi(\kappa,I''))\in\mathcal
  R[z]$, formed with $n=r+K$ as in Definition~\ref{pd-symbols-path-generators}(5), by renaming each $z_{i,c}$ to
  $z_{v_i,c}$, indexed by the quadruples $(\nu,\kappa,I',I'')$ of a path $\nu$ of rank $r$ for some $B'$, a path
  $\kappa$ of rank $K$ for some $B''$ and standard fillings $I'$ of $\nu$ and $I''$ of $\kappa$, are linearly
  independent over $\mathcal R$.
\end{enumerate}
\end{lemma}
\leanref{Formalism/SharpRestriction/GrowthBounds/ValuationLowestDegreeRename.lean}{Formalism.MainNotions.SelectionLemmas.AdmB}
\leanref{Formalism/SharpRestriction/GrowthBounds/ValuationLowestDegreeRename.lean}{Formalism.MainNotions.SelectionLemmas.admB2}
\leanref{Formalism/SharpRestriction/GrowthBounds/ValuationLowestDegreeRename.lean}{Formalism.MainNotions.SelectionLemmas.stCount}
\leanref{Formalism/SharpRestriction/GrowthBounds/ValuationLowestDegreeRename.lean}{Formalism.MainNotions.SelectionLemmas.admB3}
\leanref{Formalism/SharpRestriction/GrowthBounds/ValuationLowestDegreeRename.lean}{Formalism.MainNotions.SelectionLemmas.ColumnFamilyLowerSet}
\leanref{Formalism/SharpRestriction/GrowthBounds/ValuationLowestDegreeRename.lean}{Formalism.MainNotions.SelectionLemmas.GenericProjectionCriterion}
\leanref{Formalism/SharpRestriction/GrowthBounds/ValuationLowestDegreeRename.lean}{Formalism.MainNotions.SelectionLemmas.NestedTypeBound}
\leanref{Formalism/SharpRestriction/GrowthBounds/ValuationLowestDegreeMapFractionField.lean}{Formalism.MainNotions.G.Joints.StandardMonomialSpace.valuation_transport_of_lemma}
\leanref{Formalism/SharpRestriction/GrowthBounds/JointsLit.lean}{Formalism.MainNotions.G.Joints.Lit.weyl_dimension_formula_index_sets}
\leanref{Formalism/SharpRestriction/GrowthBounds/JointsLit.lean}{Formalism.MainNotions.G.Joints.Lit.branching_rule_count_form}
\leanref{Formalism/SharpRestriction/GrowthBounds/JointsLit.lean}{Formalism.MainNotions.G.Joints.Lit.pieri_rule_count_form}

\begin{proof}
We use the notation of Definition~\ref{pd-symbols-path-generators}.

(1) Let $F$ and $G$ be such polynomials. They involve finitely many letters. Renaming the variables of these letters
into the polynomial ring of these letters alone, with the induced lettering and the preimage of $S$, changes neither
$\operatorname{val}_S$ nor $\lambda_S$, as the substitution $P\mapsto P_{S,a}$ commutes with this injective renaming.
Nor does the passage to the fraction field $L$ of $R$, as $R[z]\to L[z]$ is injective and commutes with the
substitution. Over $L$, $P\mapsto P_{S,a}$ is a ring homomorphism into the Laurent polynomials in $\eta$ over the domain
$L[\mathfrak g][\det\mathfrak g^{-1}][z]$. There every exponent of $\eta$ in a sum occurs in one of the summands, and
the least and the largest exponents of a product of nonzero elements are the sums of the least and of the largest
exponents of the factors. This gives the four rules when $F,G\neq0$, and they are clear when $F$ or $G$ is $0$.

Let $v\in S$ and $c\leq n$. If $v$ is a vector letter, $(z_{v,c})_{S,a}$ is the entry $c$ of $z_v\tau_a(\eta)$, and the
entries of $\tau_a(\eta)=\mathfrak gD_a(\eta)\mathfrak g^{-1}$ involve only the exponents $0$ and $1$ of $\eta$; so
$0\leq\operatorname{val}_S(z_{v,c})$ and $\lambda_S(z_{v,c})\leq1$. If $v$ is a covector letter, the entries of
${}^t\tau_a(\eta^{-1})$ involve only the exponents $0$ and $-1$, which gives the other two bounds. Let $y_1,\dots,y_n\in
S$ be distinct vector letters and $Z=(z_{y_a,b})$. Then $\det(Z)_{S,n}=\det(Z\tau_n(\eta))=\det(Z)\,\eta$, as
$\det\tau_n(\eta)=\det D_n(\eta)=\eta$, and $\det Z\neq0$ as its entries are distinct variables; so the only exponent
of $\eta$ that occurs is $1$.
% Lean: admB2 D, supplied in plane_count_theorem_with_constants as Joints.valuation_transport_of_lemma D hV.1
%   = StandardMonomialSpace.valuation_transport_of_lemma: the four rules by exists_finset_letters,
%   valuation_lowest_degree_rename and valuation_lowest_degree_map_fraction_field, then the first four clauses of
%   ValuationsAdditive over the fraction field (hV.1, hV = GLWeightTheory.valuation_field D from tool_statements_hold).
%   Its first component valuationsAdditive_D7 = valuations_additive_of_generated_weight_description D
%   generated_module_weight_lemma_part_b takes these four clauses from mul_valS_mul, mul_valS_add, mul_lam1S_mul and
%   mul_lam1S_add (RootWeights/MulLaurentSup.lean, ValuationBoundDetPowerFactor.lean: mul_valS_eq writes valS through
%   mul_Phi as the least exponent of a Laurent polynomial over MvPolynomial over GLCoord, a domain by mul_domGL, and
%   mul_laurent_inf, mul_laurent_sup), the argument above; the weight description gives only the last two clauses,
%   which admB2 drops. The variables are coord_clauses (sDegs_X_vec, sDegs_X_cov); the determinant is sDegs_det
%   (twistSub_det, detX_ne_zero).

(2) The argument of fact (i) in the proof of Lemma~\ref{pd-symbols-path-count} uses only that $\mu_N$ is dominant: the
conditions on a standard filling of $\mu$ involve only the numbers of entries of $\mu_N$ at least $c$, apart from
$I(c)=\emptyset$ for $|c|>T_\mu$, which they imply as $T_\mu\geq|\mu_N(a)|$ for all $a$. So the standard fillings of
$\mu$ correspond to the semistandard tableaux of shape $\lambda=\mu_N-\mu_N(n)(1,\dots,1)$ with entries in
$\{1,\dots,n\}$. Their number is the value at $(1,\dots,1)$ of the Schur polynomial $s_\lambda$ in $n$ variables
\cite[Chapter~I, (5.12)]{Mac95}, which is $\prod_{a<b}(\lambda_a-\lambda_b+b-a)/(b-a)$ \cite[Theorem~6.3]{FH91}. This
product depends only on the differences of the entries, so it is the product of the statement.
% Lean: admB3 D is rationalPlanesClauseThree of Loc6, supplied in local_model_results as
%   Joints.Lit.weyl_dimension_formula_index_sets: restStd identifies IsSt mu with IsStd (pEnd mu), stdCountGT counts
%   these by Gelfand-Tsetlin patterns (dGT), and Model.weylQ proves the product formula for dGT by branching. The Lean
%   proves the count combinatorially; the book cites it, as in Lemma pd-kappa-defects-digits-weyl.

(3) The standard fillings of shape $\nu$ are the standard fillings of the sequence $(\nu)$ of length one, as the
conditions on the cardinalities force $I(c)=\emptyset$ for $|c|>T_{(\nu)}$, since $T_{(\nu)}\geq|\nu(a)|$ for every $a$
(Definition~\ref{pd-symbols-path-canonical}(1)). By fact (i) in the proof of Lemma~\ref{pd-symbols-path-count},
$\operatorname{st}(\nu)=d_{m+1}(\nu)$, the number of Gelfand--Tsetlin patterns with top row $\nu$ defined there. For
$m=0$, $d_1(\nu)=d_0=1$. For $m\geq1$, the $\lambda$ of the sum are dominant, as $\lambda(a+1)\leq\nu(a+1)\leq\lambda(a)$,
so $\operatorname{st}(\lambda)=d_m(\lambda)$, and the claim is the recursion that defines $d_{m+1}$.
% Lean: ColumnFamilyLowerSet D is rationalPlanesClauseFour of Loc6, supplied as Joints.Lit.branching_rule_count_form:
%   stCount_eq_dGT for nu and for each mu (antitone_of_succ), and the unfolding of dGT through gtBox.

(4) By fact (i) as in (3), $\operatorname{st}=d_{m+1}$ on the dominant vectors of $\mathbf Z^{m+1}$. A strip over a
dominant vector is dominant, and so is a vector over which a dominant vector is a strip: if $\nu'$ is a strip over $\nu$,
then $\nu'(a+1)\leq\nu(a)\leq\nu'(a)$, and if $\nu$ is a strip over $\nu'$, then $\nu'(a+1)\leq\nu(a+1)\leq\nu'(a)$. So
the two identities are fact (ii) in the proof of Lemma~\ref{pd-symbols-path-count} with $n=m+1$, as
$\binom{n+b-1}b=\binom{b+m}m$.
% Lean: GenericProjectionCriterion D is rationalPlanesClauseFive of Loc6, supplied as Joints.Lit.pieri_rule_count_form:
%   GelfandTsetlinBases.Top.gtPieri (fact (ii)) with stCount_eq_dGT, hstrip_iff, strip_up_antitone and
%   strip_down_antitone, and Nat.choose_symm_add for the binomial coefficient.

(5) This is Step~1 of the proof of Lemma~\ref{pd-symbols-residue-count-inputs-saturated}, which shows that these
polynomials form a basis of the $\mathcal R$-module of the polynomials in the $z_{v,c}$ with $v\in\Lambda_j$.
% Lean: the second clause of NestedTypeBound D, supplied as Joints.space_statement_of_lemmas D hA2 hA3 hB1 hB3 and
%   proved in space_statement_at_ring_of_basis as space_statement_linear_independent O hA2 (li_allDeg, li_ext and
%   rename_injective), with hA2 = HyperplaneRank (Lemma pd-symbols-path-basis). It uses none of the hypotheses of
%   validity, |Lambda_j| + r <= K and m >= 0 of NestedTypeBound.
\end{proof}

\begin{lemma}\label{pd-symbols-residue-count-inputs-segment-plane}
Let $C_{La}$ and $C_k$ be real numbers, and let $p$, $\mathcal R$, $\kappa$, $n'$, $r$, $K$, $k$, $T$, $L_1$, $N_1$,
$(J,f)$, $\deg$, $E'$, $N_s$, $E_0$, $m_*$, $a$, $Y$ and $q$ be as in Lemma~\ref{pd-symbols-residue-rank}, except that
$C_{La}$ and $C_k$ need not be at least $1088$; we use the notation of
Definition~\ref{pd-symbols-residue-count-inputs-index}. Let $j\in\mathbf Z$ and $b\colon\mathbf Z\to\mathbf Z_{\geq0}$,
put $N=|\Lambda_j|$ and $\beta=\beta_j(b)$, and let $\mu$ be a path of rank $k$ for some $B$ such that $B(i)=\beta_j(b,v_i)$
for every letter $i$ with $v_i<f(j,0)$, and $\mu_{i+1}(m)=\beta$ for $1\leq m\leq i'+1$ whenever $v_i=f(j,i')$ with
$i'<r$. Put $\nu=\mu^{\mathrm{pl}}$, and let $h_i(c)$ be the number of $a\leq r$ with $\nu_i(a)\geq c$. Then $N\geq r$,
$v_{N-r+t}=f(j,t)$ for $0\leq t<r$, and:
\begin{enumerate}
\item $0\leq\nu_i(a)\leq\beta$ for $0\leq i\leq N$ and $1\leq a\leq r$;
\item $\nu_N(a)=\beta$ for $1\leq a\leq r$;
\item $h_{N-r+t}(c)\leq h_{N-r+t+1}(c)\leq h_{N-r+t}(c)+1$ for $0\leq t<r$ and $c\in\mathbf Z$;
\item the canonical filling $\operatorname{can}_{\nu_N}$ (Definition~\ref{pd-symbols-path-canonical}(2)) takes the value
  $\{1,\dots,r\}$ at $1\leq c\leq\beta$ and $\emptyset$ at the other $c$, and it is a standard filling of $\nu$.
\end{enumerate}
\end{lemma}
\leanref{Formalism/SharpRestriction/GrowthBounds/SegmentDegree/MemColRangeOfLeEntry.lean}{Formalism.MainNotions.SelectionLemmas.Pf.PlHyp}
\leanref{Formalism/SharpRestriction/GrowthBounds/SegmentDegree/DplusLeBeS.lean}{Formalism.MainNotions.SelectionLemmas.Pf.plHyp_of}
\leanref{Formalism/SharpRestriction/GrowthBounds/SegmentDegree/MemColRangeOfLeEntry.lean}{Formalism.MainNotions.SelectionLemmas.Pf.Ican_eq}
\leanref{Formalism/SharpRestriction/GrowthBounds/SegmentDegree/DiscPowerMonomialMemSupported.lean}{Formalism.MainNotions.SelectionLemmas.Pf.mem_isSt_Ican}

\begin{proof}
By Step~2 of the proof of Lemma~\ref{pd-symbols-residue-lifts-numeric}, the elements of $\Lambda_j$ that are at least
$f(j,0)$ are the $f(j,t)$, $t<r$, which are even and increase with $t$ by conditions~(2) and~(3) of validity. So
$N\geq r$, $v_{N-r+t}=f(j,t)$, the letters $N-r,\dots,N-1$ are vector letters, and the letters $i<n_0=N-r$ are those
with $v_i<f(j,0)$. By (2) of Theorem~\ref{pd-symbols-count}, $N+r\leq N_s+r\leq K$.

(1) We show $\mu_s(a)\geq0$ whenever $a+s\leq k$, by induction on $s$. For $s=0$, $\mu_0=0$. At a vector letter $s$,
$\mu_{s+1}(a)\geq\mu_s(a)$. At a covector letter $s$, $\mu_s$ is a strip over $\mu_{s+1}$, so
$\mu_{s+1}(a)\geq\mu_s(a+1)$, and $(a+1)+s\leq k$. As $i+a\leq N+r\leq k$ for $i\leq N$ and $a\leq r$, this gives
$\nu_i(a)\geq0$.

For the upper bound, let first $s\leq n_0$. At a vector letter $l$, $\mu_{l+1}(a)\leq\mu_l(a)+B(l)$, as the
nonnegative differences $\mu_{l+1}(a')-\mu_l(a')$ add up to $B(l)$; at a covector letter, $\mu_{l+1}(a)\leq\mu_l(a)$.
As $\mu_0=0$, $\mu_s(a)$ is at most the sum of the $B(l)$ over the vector letters $l<s$, which is at most
$\sum_{l<n_0}B(l)=\sum_{l<n_0}\beta_j(b,v_l)=\beta$. We show $\mu_{n_0+t}(a)\leq\beta$ for $t\leq r$ and all $a\leq k$, by
induction on $t$, the case $t=0$ being the bound just shown. The letter $l=n_0+t$ has $v_l=f(j,t)$, so
$\mu_{l+1}(1)=\beta$ by the hypothesis; and for $a\geq2$, $\mu_{l+1}(a)\leq\mu_l(a-1)\leq\beta$, as $\mu_{l+1}$ is a
strip over $\mu_l$.

(2) The letter $N-1$ has $v_{N-1}=f(j,r-1)$, so $\mu_N(a)=\beta$ for $a\leq r$ by the hypothesis.

(3) The letter $l=N-r+t$ is a vector letter, so $\nu_l(a)\leq\nu_{l+1}(a)$ for $a\leq r$ and $\nu_{l+1}(a+1)\leq\nu_l(a)$
for $a<r$, and $\nu_l$ and $\nu_{l+1}$ are dominant. The first inequalities give $h_l(c)\leq h_{l+1}(c)$. If
$\nu_{l+1}(a+1)\geq c$, then $\nu_l(a)\geq c$; so the entries of $\nu_{l+1}$ at least $c$, apart from the first, are
at most $h_l(c)$ in number, and $h_{l+1}(c)\leq h_l(c)+1$.

(4) By (1) and (2), $\beta\geq0$, and $h_N(c)=r$ for $c\leq\beta$ and $h_N(c)=0$ for $c>\beta$. So
$\operatorname{can}_{\nu_N}(c)=\{1,\dots,h_N(c)\}$ is $\{1,\dots,r\}$ for $1\leq c\leq\beta$ and $\emptyset$ for $c>\beta$,
and for $c\leq0$, $\operatorname{can}_{\nu_N}(c)=\{h_N(c)+1,\dots,r\}=\emptyset$. It is a standard filling of $\nu$. If
$1\leq c\leq\beta$, then $c\leq\nu_N(1)\leq T_\nu$; so $\operatorname{can}_{\nu_N}(c)=\emptyset$ for $|c|>T_\nu$. Its value at
$c$ has $h_N(c)$ elements for $c\geq1$, and $0=r-h_N(c)$ elements for $c\leq0$. Finally, with
$\operatorname{can}^*_{\nu_N}$ as in Definition~\ref{pd-symbols-path-generators}(4), $\operatorname{can}^*_{\nu_N}(c)$ is
$\{1,\dots,r\}$ for $c\leq\beta$ and $\emptyset$ for $c>\beta$, so $\operatorname{can}^*_{\nu_N}(c+1)\subset
\operatorname{can}^*_{\nu_N}(c)$ for every $c$.
% Lean: plHyp_of (k_r_le_segN, segE_eq_fr_iff, comb_vec with validity for the vector letters): rng is
%   T.mu_nonneg (the induction of (1)) with hlow (T.mu_le_dplus and dplus_le_beS, using hP.degs) and hup (the
%   induction on t with hP.head and the strip T.step_vec); fin is hP.head at the letter of f(j, r - 1); up is colH_step.
%   (4) is Ican_eq (colH_all, colH_none) and mem_isSt_Ican (mem_colRange_of_le_entry).
\end{proof}

\begin{lemma}\label{pd-symbols-residue-count-inputs-segment-members}
Let $C_{La}\geq1088$ and $C_k\geq1088$ be real numbers, and let $p$, $\mathcal R$, $\kappa$, $n'$, $r$, $K$, $k$, $T$,
$L_1$, $N_1$, $(J,f)$, $\deg$, $E'$, $N_s$, $E_0$, $m_*$, $a$, $Y$ and $q$ be as in
Lemma~\ref{pd-symbols-residue-rank}; we use the notation of Definitions~\ref{pd-symbols-residue-coarse},
\ref{pd-symbols-residue-families} and~\ref{pd-symbols-residue-count-inputs-index}. Let $j\in\mathbf Z$ and let
$b\colon\mathbf Z\to\mathbf Z_{\geq0}$ be balanced at $j$, vanish off $\Lambda_j$, and satisfy $W_j(x)\leq b(x)$ for every
$x\in\Lambda_j$ that is not new. Then $y_j(b;q)\in M_j(b)$ for every $q\in\mathcal I_j(b)$.
\end{lemma}
\leanref{Formalism/SharpRestriction/GrowthBounds/SegmentDegree/DiscPowerTimesMonomial.lean}{Formalism.MainNotions.SelectionLemmas.Pf.restMem_of_rests}
\leanref{Formalism/SharpRestriction/GrowthBounds/SegmentDegree/DiscPowerMonomialMemSupported.lean}{Formalism.MainNotions.SelectionLemmas.Pf.restMemA_of}
\leanref{Formalism/SharpRestriction/GrowthBounds/SegmentDegree/DiscPowerMonomialDegreeAtW.lean}{Formalism.MainNotions.SelectionLemmas.Pf.restMemB_of}
\leanref{Formalism/SharpRestriction/GrowthBounds/SegmentDegree/M0LeBwOfIsW.lean}{Formalism.MainNotions.SelectionLemmas.Pf.restMemC_of}
\leanref{Formalism/SharpRestriction/GrowthBounds/SegmentDegree/M0LeBwOfIsW.lean}{Formalism.MainNotions.SelectionLemmas.Pf.restMemD_of}
\leanref{Formalism/SharpRestriction/GrowthBounds/SegmentDegree/M0LeBwOfIsW.lean}{Formalism.MainNotions.SelectionLemmas.Pf.m0_le_Bw_of_isW}
\leanref{Formalism/SharpRestriction/GrowthBounds/SegmentDegree/DiscPowerMonomialMemSupported.lean}{Formalism.MainNotions.SelectionLemmas.Pf.exists_path_decomposition}

\begin{proof}
Let $N=|\Lambda_j|$, let $v_0<\dots<v_{N-1}$ be the elements of $\Lambda_j$, lettered as in
Definition~\ref{pd-symbols-residue-pairs}, and put $S_i=\{0,\dots,i-1\}$. We use the renaming of the proof of
Lemma~\ref{pd-symbols-residue-segment-span}, under which $\operatorname{val}_{\{y<v_i\}}$ and $\lambda_{\{y<v_i\}}$ of a
renamed polynomial are $\operatorname{val}_{S_i}$ and $\lambda_{S_i}$ of the original, and likewise with $y\leq v_i$ and
$S_{i+1}$. By Lemma~\ref{pd-symbols-division-rows}(1), the elements of $\Lambda_j$ that are not new are the even
$f(j-1,i)$, $i<r$. By Step~2 of the proof of Lemma~\ref{pd-symbols-residue-lifts-numeric}, the elements of $\Lambda_j$
that are at least $f(j,0)$ are the even $f(j,i')$, $i'<r$. The others are the elements of $\Lambda'_j$. So
$v_i=f(j-1,i)$ for $i<r$, $v_{N-r+i'}=f(j,i')$ for $i'<r$, and the $v_i$ with $r\leq i<N-r$ are the elements of
$\Lambda'_j$.

\emph{Step 1: $m_0\leq W_j(x)$ for every $x\in\Lambda_j$ that is not new.} Put $P=P_j(x)$. As $x$ is not new,
$x\leq f(j-1,r-1)<f(j,0)$, and $\beta_j(b,x)=\max(W_j(x)-m_0,0)$. By Step~2 of the proof of
Lemma~\ref{pd-symbols-residue-count-inputs-segment-atypical-dictionary}, $N_s\geq6$: the $f(0,i)$, the $f(1,i)$ and
the elements of $\Lambda'_1$ are at least $3r\geq6$ elements of $\Lambda_1$. With this, the lower bound of Step~3 of
that proof, in the case of an element that is not new, gives $\beta_j(b,x)\geq d^D(\bar x)/P-2kN_s$. By
Lemma~\ref{pd-symbols-residue-count-sum-facts}(4), $d^D(\bar x)/P\geq\frac{36}5Yk^2N_s$. As $Y\geq1$ and $k\geq1$,
\[
 \max(W_j(x)-m_0,0)\geq\frac{d^D(\bar x)}P-2kN_s\geq\Big(\frac{36}5Yk-2\Big)kN_s>0,
\]
so $W_j(x)-m_0>0$.

\emph{Step 2: the splitting of the path.} Let $q=(\mu,I'',u)\in\mathcal I_j(b)$, with $\mu$ a counting path for
$(j,b)$, a path of rank $k$ for $B$. Put $\nu=\mu^{\mathrm{pl}}$ and $\kappa=\mu^{\mathrm{tn}}$, and let $\operatorname{can}$
be the canonical filling of $\nu_N$. By (2) of Theorem~\ref{pd-symbols-count}, $N+r\leq N_s+r\leq K$, so
Lemma~\ref{pd-symbols-path-clauses}(2) gives $B'$ and $B''$ with $B'+B''=B$ such that $\nu$ is a path of rank $r$ for
$B'$, $\kappa$ is a path of rank $K$ for $B''$, and $\kappa_{i+1}(1)\leq\nu_i(r)$ at every vector letter $i$ and
$\kappa_i(1)\leq\nu_{i+1}(r)$ at every covector letter $i$. By
Lemma~\ref{pd-symbols-residue-count-inputs-segment-plane}(4), $\operatorname{can}$ is a standard filling of $\nu$; and
$I''$ is one of $\kappa$, by the definition of $T_j(b)$. Let $\pi$ be $\nu$, of rank $n=r$ for $B'$, or $\kappa$, of rank
$n=K$ for $B''$, and $I$ the standard filling $\operatorname{can}$ or $I''$ of it. By
Lemma~\ref{pd-symbols-path-clauses}(1) over $\mathcal R$, with $\theta(i)=\pi_{i+1}(n)$ and $M(i)=\pi_{i+1}(1)$, the
$\mathcal R$-span of the path generators $\Psi(\pi',I')$ with $\pi'$ a path for the same row degrees such that
$\pi_{i+1}(n)\leq\pi'_{i+1}(n)$ and $\pi'_{i+1}(1)\leq\pi_{i+1}(1)$ for every letter $i$ is the set of the polynomials $G$
of these row degrees with $\pi_{i+1}(n)\leq\operatorname{val}_{S_{i+1}}(G)$ and $\lambda_{S_{i+1}}(G)\leq\pi_{i+1}(1)$ for
every letter $i$. The generator $\Psi(\pi,I)$ is one of those spanning it, so it lies in this set. At $i=0$,
$S_0=\emptyset$ and $P_{\emptyset,a}=P$, so the only exponent of $\eta$ that can occur is $0=\pi_0(n)=\pi_0(1)$. Hence
$\Psi(\nu,\operatorname{can})$ and $\Psi(\kappa,I'')$ have row degrees $B'$ and $B''$, and for $0\leq i\leq N$
\[
 \nu_i(r)\leq\operatorname{val}_{S_i}\big(\Psi(\nu,\operatorname{can})\big),\qquad
 \lambda_{S_i}\big(\Psi(\kappa,I'')\big)\leq\kappa_i(1).
\]
Next, $\kappa_s=0$ for $s\leq r$. Indeed, the letters $0,\dots,r-1$ are vector letters, and at a vector letter $i$ a path has
$\mu_i\leq\mu_{i+1}$ entrywise; as $\mu_0=0$, $\kappa_s\geq0$ for $s\leq r$. Also $\mu_{s+1}(a+1)\leq\mu_s(a)$ at a vector
letter $s$ and $\mu_{s+1}\leq\mu_s$ at a covector letter, so induction on $s$ gives $\mu_s(a)\leq0$ for $a>s$; hence
$\kappa_s(c)=\mu_s(r+c)\leq0$ for $s\leq r$. So $B''(i)=\sum_c(\kappa_{i+1}(c)-\kappa_i(c))=0$ for $i<r$.

Write $y=y_j(b;q)=AG$, where $G$ is $\Psi(\kappa,I'')$ with each $z_{i,c}$ renamed $z_{v_i,r+c}$, and $A=A_0A_1A_2$ with
$A_2$ the polynomial $\Psi(\nu,\operatorname{can})$ with each $z_{i,c}$ renamed $z_{v_i,c}$,
\[
 A_0=D_j^{\,m_0}\prod_{i<r}z_{f(j-1,i),1}^{\,\max(b(f(j-1,i))-W_j(f(j-1,i)),0)},\qquad
 A_1=\prod_{i'<r,\,c\leq r}z_{f(j,i'),c}^{\,u(i',c)} .
\]

\emph{Step 3: twisted degrees of $A$.} Let $S\subset\mathbf Z$ contain the $f(j-1,i)$, $i<r$. Apply
Lemma~\ref{pd-symbols-residue-count-inputs-segment-clauses}(1) over $\mathcal R$ to the polynomials in the $z_{v,c}$,
$v\in\mathbf Z$, $c\leq r$, with the even $v$ as vector letters and the odd ones as covector letters, as in
Definition~\ref{pd-symbols-module}(1). It gives $\operatorname{val}_S(PP')=\operatorname{val}_S(P)+\operatorname{val}_S(P')$,
and $\operatorname{val}_S(D_j)=1$, as $D_j$ is the determinant of the rows of the distinct even elements $f(j-1,a)$ of
$S$. A variable $z_{v,c}$ with $v$ even has $\operatorname{val}_S(z_{v,c})\geq0$: by the same lemma if $v\in S$, and as
the substitution does not change it if $v\notin S$. So $\operatorname{val}_S(A)\geq m_0+\operatorname{val}_S(A_2)$. For a
new $v_i$ the $f(j-1,a)$
are smaller than $v_i$, so with Step~2
\[
 \operatorname{val}_{\{y<v_i\}}(A)\geq m_0+\nu_i(r),\qquad\operatorname{val}_{\{y\leq v_i\}}(A)\geq m_0+\nu_{i+1}(r).
\]

\emph{Step 4: the segment module.} $A$ is a polynomial in the plane variables $z_{v,c}$ with $v\in\Lambda_j$, as the
$f(j-1,i)$ and $f(j,i')$ lie in $\Lambda_j$. $G$ is a polynomial in the normal variables $z_{v_i,c}$ with $v_i$ new, as
$\Psi(\kappa,I'')$ has row degree $B''(i)=0$ at the letters $i<r$. Put $M(v_i)=\kappa_{i+1}(1)$ if $v_i$ is even,
$M(v_i)=\kappa_i(1)$ if $v_i$ is odd, and $M(v)=0$ for $v\notin\Lambda_j$. Let $v_i$ be new.
\begin{itemize}
\item If $v_i$ is even, $i$ is a vector letter. By Steps 2 and 3, $\lambda_{\{y\leq v_i\}}(G)\leq\kappa_{i+1}(1)=M(v_i)$ and
  $M(v_i)+m_0\leq\nu_i(r)+m_0\leq\operatorname{val}_{\{y<v_i\}}(A)$.
\item If $v_i$ is odd, $i$ is a covector letter, and in the same way $\lambda_{\{y<v_i\}}(G)\leq\kappa_i(1)=M(v_i)$ and
  $M(v_i)+m_0\leq\nu_{i+1}(r)+m_0\leq\operatorname{val}_{\{y\leq v_i\}}(A)$.
\item By the definition of a counting path, $\theta'_{j,b}(v_i)\leq\mu_i(r)=\nu_i(r)\leq\operatorname{val}_{\{y<v_i\}}(A)$, as
  $m_0\geq0$.
\end{itemize}
The factors $A_0$ and $A_1$ involve only even rows, and $A$ involves no normal variable. So for every monomial $\xi$ of
$y$, $\operatorname{ord}_{v_i}(\xi)$ is the degree $B''(i)$ of $G$ in $z_{v_i}$ if $v_i$ is even, and the degree $B'(i)$ of
$A_2$ in $z_{v_i}$ if $v_i$ is odd. At a vector letter $i$,
$B''(i)=\sum_{c=r+1}^k(\mu_{i+1}(c)-\mu_i(c))$, and at a covector letter $B'(i)=\sum_{c=1}^r(\mu_i(c)-\mu_{i+1}(c))$; by
the definition of a counting path both are at most $o'_j(v_i)$ when $v_i$ is new. Hence $y\in X_j(m_0,\theta'_{j,b},o'_j)$.

\emph{Step 5: the row degrees.} Each determinant $D_j$ has degree $1$ in each of its rows, and $W_j(x)\leq b(x)$ at the
$x=f(j-1,i)$; so $A_0$ has degree $m_0+b(x)-W_j(x)$ in $z_x$ at these $x$ and $0$ in the other rows. $A_1$ has degree
$\sum_cu(i',c)$ in $z_{f(j,i')}$ and $0$ in the other rows, and $A_2G$ has degree $B'(i)+B''(i)=B(i)$ in $z_{v_i}$ and $0$
in the rows off $\Lambda_j$.
\begin{itemize}
\item At $x=v_i=f(j-1,i)$, $i<r$: $x<f(j,0)$ and $x$ is not new, so $B(i)=\beta_j(b,x)=\max(W_j(x)-m_0,0)=W_j(x)-m_0$ by
  Step~1, and the degree of $y$ is $m_0+b(x)-W_j(x)+W_j(x)-m_0=b(x)$.
\item At $x=v_i\in\Lambda'_j$: $x<f(j,0)$ and $x$ is new, so the degree is $B(i)=\beta_j(b,x)=b(x)$.
\item At $x=f(j,i')=v_{N-r+i'}$: the degree is $B(N-r+i')+\sum_cu(i',c)=b(x)$, by the definition of $\mathcal I_j(b)$.
\item At $x\notin\Lambda_j$: the degree is $0=b(x)$.
\end{itemize}

\emph{Step 6: the exponents at the $f(j,i')$.} Let $i'<r$, $c\leq r$ and $i=N-r+i'$, so $v_i=f(j,i')$. The exponent of
$z_{v_i,c}$ in a monomial of $y$ is $u(i',c)$ plus its exponent in a monomial of $A_0A_2G$, which is at most the degree
$B(i)$ of $A_0A_2G$ in $z_{v_i}$. As $v_i\geq f(j,0)$, Lemma~\ref{pd-symbols-residue-count-inputs-rows} gives
$B(i)\leq2\beta_j(b)$, so the exponent is at most $u(i',c)+2\beta_j(b)<Q_j$ by the definition of $\mathcal I_j(b)$. With
Steps 4 and 5, $y\in M_j(b)$.
\end{proof}
% Lean: restMem_of D hA = restMem_of_rests 1088 (restMemA_of, restMemB_of, restMemC_of, restMemD_of) with
%   famEl_eq (y = plEmb memA * nmEmb memG, memA = discPowerTimesMonomial * memMu * memPl; A_0, A_1, A_2 here).
%   Step 1 is m0_le_Bw_of_isW (TE.c_bounds, the bound of Lemma pd-symbols-residue-count-sum-facts(4) from
%   top_count_lower_bounds; TE.Bdeg_bounds, its IsW case, with TE.six_le_Ns from LLset 1: then Bdeg > 0 and so
%   Bw - m0 > 0). Step 2 is exists_path_decomposition (the AdmB clause
%   BracketPolynomialSpan = RotatedSumBound, Lemma pd-symbols-path-clauses(2), with segN + r <= K from hX.card and
%   hX.NsK), mem_isSt_Ican with plHyp_of (Lemma pd-symbols-residue-count-inputs-segment-plane(4)),
%   standard_polynomial_mem_bounds (A3at with th i = mu (i+1) (lastIdx n) and M i = mu (i+1) 0, the clause
%   DefiniteRankOne, Lemma pd-symbols-path-clauses(1); pre_zero and sDegs_empty_subset at i = 0), and mem_tan_zero and
%   tangent_degree_zero_at_w (T.mu_nonpos, comb_nonneg_of_vec). Step 3 is mem_vg_A, from the AdmB clause admB2
%   (Lemma pd-symbols-residue-count-inputs-segment-clauses(1)): mem_vg_mul, mem_vg_pow and mem_vg_prod for products,
%   (hB2 O Z r LZ S).2.2 for the determinant, and mem_vg_X (sDegs_X_vec, sDegs_X_out) for the even variables, in S or
%   not. Step 4 is restMemA_of (seg_val,
%   CondC, IsGood) and restMemB_of (ordDeg_prod, T.step_vec, T.step_cov). Step 5 is restMemC_of
%   (disc_power_monomial_degree_at_w, mem_dMu_tau, hPATH.degs, hd). Step 6 is restMemD_of (monTau_apply,
%   tau_degree_on_path_bound = Lemma pd-symbols-residue-count-inputs-rows). The Lean indexes the family by
%   Idx j (restr j b); its conditions use b only at the new elements, where restr j b = b, so it is the book's
%   I_j(b).

\begin{lemma}\label{pd-symbols-residue-count-inputs-segment-kill}
Let $R$ be a commutative ring, $R[z]$ the polynomial ring over $R$ in a set $V$ of variables, and $Q$ a set; for $q\in Q$
let $m_q$ be a monomial and $w_q\in R[z]$.
\begin{enumerate}
\item Let $S\subset V$, and let $\chi_S$ be the $R$-algebra endomorphism of $R[z]$ that sends the variables in $S$ to $0$
  and fixes the others. If the $m_q\chi_S(w_q)$, $q\in Q$, are linearly independent over $R$, then so are the $m_qw_q$.
\item Let $R$ be a domain, $T\subset V$, and $w'_x\in R[z]$, $x\in X$, polynomials linearly independent over $R$ that
  involve no variable of $T$. For $q\in Q$ let $M_q$ be a monomial in the variables of $T$ and $x_q\in X$, such that
  $q\mapsto(M_q,x_q)$ is injective. Then the $M_qw'_{x_q}$, $q\in Q$, are linearly independent over $R$.
\end{enumerate}
\end{lemma}
\leanref{Formalism/SharpRestriction/GrowthBounds/SegmentDegree/ScaleSCoeffZero.lean}{Formalism.MainNotions.SelectionLemmas.Pf.kill}
\leanref{Formalism/SharpRestriction/GrowthBounds/SegmentDegree/ScaleSCoeffZero.lean}{Formalism.MainNotions.SelectionLemmas.Pf.li_of_kill}
\leanref{Formalism/SharpRestriction/GrowthBounds/SegmentDegree/KillOfEmpty.lean}{Formalism.MainNotions.SelectionLemmas.Pf.li_mono_words}

\begin{proof}
(1) Suppose that $\sum_{q\in Q_0}a_qm_qw_q=0$ for a finite $Q_0\subset Q$ and $a_q\in R$ not all $0$. For a monomial $m$
let $\deg_S(m)$ be its degree in the variables of $S$, and let $\delta$ be the least $\deg_S(m_q)$ over the $q\in Q_0$
with $a_q\neq0$. Let $\sigma$ be the $R$-algebra map from $R[z]$ to $R[z][t]$, with $t$ a further variable, that sends
each variable of $S$ to $t$ times itself and fixes the others. It sends a monomial $m$ to $t^{\deg_S(m)}m$, so the
constant term of $\sigma(P)$ in $t$ is $\chi_S(P)$, and $\sigma(m_qw_q)=t^{\deg_S(m_q)}m_q\sigma(w_q)$. The coefficient of
$t^\delta$ in $\sigma$ of the relation is therefore $\sum a_qm_q\chi_S(w_q)$, the sum over the $q\in Q_0$ with
$\deg_S(m_q)=\delta$, as the other $q$ with $a_q\neq0$ have $\deg_S(m_q)>\delta$. By hypothesis these $a_q$ are $0$,
which contradicts the choice of $\delta$.

(2) Let $\sum_{q\in Q_0}a_qM_qw'_{x_q}=0$ with $Q_0\subset Q$ finite, and let $q_0\in Q_0$. Let $\Phi$ be the $R$-linear
endomorphism of $R[z]$ that sends a monomial $m=m'm''$, with $m'$ a monomial in the variables of $T$ and $m''$ one in
the others, to $m''$ if $m'=M_{q_0}$ and to $0$ otherwise. As $w'_{x_q}$ involves no variable of $T$,
$\Phi(M_qw'_{x_q})$ is $w'_{x_q}$ if $M_q=M_{q_0}$ and $0$ otherwise. Applying $\Phi$ to the relation gives
$\sum a_qw'_{x_q}=0$, the sum over the $q\in Q_0$ with $M_q=M_{q_0}$. For these $q$ the $x_q$ are distinct, as
$q\mapsto(M_q,x_q)$ is injective, so all these $a_q$ are $0$, and $a_{q_0}$ among them.
% Lean: (1) is li_of_kill (scaleS, scaleS_monomial, scaleS_coeff_zero with offDeg, exists_min_image for delta);
%   (2) is li_mono_words (Phi built with basisMonomials.constr, hkey, hxinj and hfib).
\end{proof}

\begin{lemma}\label{pd-symbols-residue-count-inputs-segment-cut}
Let $r\geq1$ and $n_0\geq0$ be integers, $N=n_0+r$, fix a lettering of $\{0,\dots,N-1\}$, and let $\beta$ be an integer.
Let $\nu=(\nu_0,\dots,\nu_N)$ be a sequence of dominant vectors in $\mathbf Z^r$, let $h_i(c)$ be the number of $a$ with
$\nu_i(a)\geq c$, and suppose that $0\leq\nu_i(a)\leq\beta$ for all $i$ and $a$, $\nu_N(a)=\beta$ for all $a$, and
$h_{n_0+t}(c)\leq h_{n_0+t+1}(c)\leq h_{n_0+t}(c)+1$ for $0\leq t<r$ and $c\in\mathbf Z$. For $c\in\mathbf Z$ let $U(c)$
be the set of the $a\in\{1,\dots,r\}$ with $h_{n_0+r-a}(c)<h_{n_0+r-a+1}(c)$. Let $\nu^{\mathrm{cut}}$ be the sequence
with $\nu^{\mathrm{cut}}_i=\nu_{\min(i,n_0)}$, and $I^{\mathrm{cut}}$ the function with
$I^{\mathrm{cut}}(c)=\{1,\dots,r\}\smallsetminus U(c)$ for $1\leq c\leq\beta$ and $I^{\mathrm{cut}}(c)=\emptyset$
otherwise. Let $R$ be a commutative ring, form the polynomials of Definition~\ref{pd-symbols-path-generators} over $R$
with $n=r$, and let $\chi$ be the $R$-algebra endomorphism of $R[z]$ that sends $z_{n_0+t,c}$ to $0$ for $0\leq t<r$
and $c\neq r-t$ and fixes the other variables. Then:
\begin{enumerate}
\item there is $e\geq0$ with $\chi\big(\Psi(\nu,\operatorname{can}_{\nu_N})\big)=(-1)^e\prod_{c=1}^{\beta}\prod_{a\in U(c)}
  z_{n_0+r-a,a}\cdot\Psi(\nu^{\mathrm{cut}},I^{\mathrm{cut}})$;
\item for every function $I$ from $\mathbf Z$ to the subsets of $\{1,\dots,r\}$, $\Psi(\nu^{\mathrm{cut}},I)$ involves only
  the variables $z_{i,c}$ with $i<n_0$;
\item $I^{\mathrm{cut}}$ is a standard filling of $\nu^{\mathrm{cut}}$; and if $\nu$ is a path of rank $r$ for some $B$, then
  $\nu^{\mathrm{cut}}$ is a path of rank $r$ for the function equal to $B$ at the letters $i<n_0$ and to $0$ at the
  others;
\item if $\nu'$ is another sequence with these properties, for the same $n_0$ and $\beta$, such that
  $\nu'^{\mathrm{cut}}=\nu^{\mathrm{cut}}$ and $I'^{\mathrm{cut}}=I^{\mathrm{cut}}$, then $\nu'=\nu$.
\end{enumerate}
\end{lemma}
\leanref{Formalism/SharpRestriction/GrowthBounds/SegmentDegree/MemColRangeOfLeEntry.lean}{Formalism.MainNotions.SelectionLemmas.Pf.cutPath}
\leanref{Formalism/SharpRestriction/GrowthBounds/SegmentDegree/MemColRangeOfLeEntry.lean}{Formalism.MainNotions.SelectionLemmas.Pf.Icut}
\leanref{Formalism/SharpRestriction/GrowthBounds/SegmentDegree/KillKcolTau.lean}{Formalism.MainNotions.SelectionLemmas.Pf.riseA}
\leanref{Formalism/SharpRestriction/GrowthBounds/SegmentDegree/KillKcolTau.lean}{Formalism.MainNotions.SelectionLemmas.Pf.offN}
\leanref{Formalism/SharpRestriction/GrowthBounds/SegmentDegree/KillKcolTau.lean}{Formalism.MainNotions.SelectionLemmas.Pf.kill_Kcol_tau}
\leanref{Formalism/SharpRestriction/GrowthBounds/SegmentDegree/MemColRangeOfLeEntry.lean}{Formalism.MainNotions.SelectionLemmas.Pf.kill_Psi}
\leanref{Formalism/SharpRestriction/GrowthBounds/SegmentDegree/PsiCutSupported.lean}{Formalism.MainNotions.SelectionLemmas.Pf.Psi_cut_supported}
\leanref{Formalism/SharpRestriction/GrowthBounds/SegmentDegree/CardIcutFilter.lean}{Formalism.MainNotions.SelectionLemmas.Pf.isSt_cut}
\leanref{Formalism/SharpRestriction/GrowthBounds/SegmentDegree/DplusLeBeS.lean}{Formalism.MainNotions.SelectionLemmas.Pf.cut_ppath}
\leanref{Formalism/SharpRestriction/GrowthBounds/SegmentDegree/CardIcutFilter.lean}{Formalism.MainNotions.SelectionLemmas.Pf.path_eq_of_cut}

\begin{proof}
For $c\in\mathbf Z$ let $F^c_i$ be the function after the first $i$ steps of the definition of $\mathrm{Col}_{\nu,c}$
(Definition~\ref{pd-symbols-path-generators}(3)), so that $F^c_N=\mathrm{Col}_{\nu,c}$, and for $t\leq r$ let $U_t(c)$ be
the set of the $a\in U(c)$ with $r-a<t$, so that $U_r(c)=U(c)$. As $\nu_N(1)=\beta$ and $\nu_N(1)\geq0$, $\beta\geq0$.

\emph{Step 1: the rises.} We show $h_{n_0}(c)+|U_t(c)|=h_{n_0+t}(c)$ for $t\leq r$, by induction on $t$. The letter
$n_0+t$ is $n_0+r-a$ for $a=r-t$, and $U_{t+1}(c)$ is $U_t(c)$ together with $r-t$ if $h_{n_0+t}(c)<h_{n_0+t+1}(c)$, and
$U_t(c)$ otherwise. By the hypothesis on the $h_{n_0+t}$, the height grows by $1$ in the first case and stays in the
second. For $1\leq c\leq\beta$, $h_N(c)=r$, so $|U(c)|=r-h_{n_0}(c)$.

\emph{Step 2: the columns outside $[1,\beta]$.} If $c>\beta$, then $h_i(c)=0$ for all $i$; if $c\leq0$, then $h_i(c)=r$
for all $i$. In both cases the steps for $c$ do nothing, so $F^c_i=F_0$ for all $i$, and $F_0(\emptyset)=1$. The values of
$\operatorname{can}_{\nu_N}$ and of $I^{\mathrm{cut}}$ at these $c$ are $\emptyset$. As $1\leq c\leq\beta$ implies
$c\leq\nu_N(1)\leq T_\nu$, $\Psi(\nu,\operatorname{can}_{\nu_N})=\prod_{c=1}^\beta\mathrm{Col}_{\nu,c}(\{1,\dots,r\})$.

\emph{Step 3: the cut sequence.} The numbers of entries at least $c$ of $\nu^{\mathrm{cut}}_i$ are $h_i(c)$ for $i\leq n_0$
and $h_{n_0}(c)$ for $i\geq n_0$, so the steps of $\mathrm{Col}_{\nu^{\mathrm{cut}},c}$ after the first $n_0$ do nothing,
and $\mathrm{Col}_{\nu^{\mathrm{cut}},c}=F^c_{n_0}$. Each step multiplies by entries of the row of its letter, so the values
of $F^c_{n_0}$ involve only the rows $z_i$ with $i<n_0$; this proves (2). In $\Psi(\nu^{\mathrm{cut}},I^{\mathrm{cut}})$
the factors at the $c$ outside $[1,\beta]$ are $F_0(\emptyset)=1$ by Step 2. Let $1\leq c\leq\beta$ with
$c>T_{\nu^{\mathrm{cut}}}$. Then all entries of the $\nu_i$, $i\leq n_0$, are less than $c$, so $h_i(c)=0$ for $i\leq n_0$
and $F^c_{n_0}=F_0$; and $|U(c)|=r$ by Step 1, so $I^{\mathrm{cut}}(c)=\emptyset$ and $F^c_{n_0}(I^{\mathrm{cut}}(c))=1$.
Hence $\Psi(\nu^{\mathrm{cut}},I^{\mathrm{cut}})=\prod_{c=1}^\beta F^c_{n_0}(\{1,\dots,r\}\smallsetminus U(c))$.

\emph{Step 4: killing a column.} Let $c\geq1$. The map $\chi$ fixes the rows of the letters $i<n_0$, so
$\chi(F^c_i(J))=F^c_i(J)$ for $i\leq n_0$ and every $J$. We show by induction on $t\leq r$ that for every
$J\supset U_t(c)$ there is $e\geq0$ with
\[
 \chi\big(F^c_{n_0+t}(J)\big)=(-1)^e\prod_{a\in U_t(c)}z_{n_0+r-a,a}\cdot F^c_{n_0}\big(J\smallsetminus U_t(c)\big).
\]
For $t=0$ this was just shown. Let $l=n_0+t<N$. As $\chi$ is a ring homomorphism, $\chi(F^c_{l+1}(J))$ is the step at $l$
for $c$ applied to the function $\chi\circ F^c_l$ and the row $\chi(z_l)$, whose only nonzero entry is $z_{l,r-t}$. If
$h_l(c)=h_{l+1}(c)$, the step does nothing and $U_{t+1}(c)=U_t(c)$. Otherwise $h_l(c)<h_{l+1}(c)$, $U_{t+1}(c)$ is $U_t(c)$
together with $r-t\notin U_t(c)$, and, as $c\geq1$, the step is the wedge: $\chi(F^c_{l+1}(J))=\sum_{m\in J}\epsilon_J(m)\,
\chi(z_{l,m})\,\chi(F^c_l(J\smallsetminus\{m\}))$. Only $m=r-t$ remains, and it lies in $J$, so
$\chi(F^c_{l+1}(J))=\epsilon_J(r-t)\,z_{l,r-t}\,\chi(F^c_l(J\smallsetminus\{r-t\}))$. The induction hypothesis for
$J\smallsetminus\{r-t\}\supset U_t(c)$ gives the claim, as $(J\smallsetminus\{r-t\})\smallsetminus U_t(c)=J\smallsetminus
U_{t+1}(c)$ and $l=n_0+r-(r-t)$.

(1) Apply Step 4 with $t=r$ and $J=\{1,\dots,r\}$ for each $1\leq c\leq\beta$, take the product over $c$, and use Steps 2
and 3.

(3) The vector $\nu^{\mathrm{cut}}_N=\nu_{n_0}$ has $h_{n_0}(c)$ entries at least $c$. For $1\leq c\leq\beta$, Step 1 gives
$|I^{\mathrm{cut}}(c)|=r-|U(c)|=h_{n_0}(c)$; for $c>\beta$, $|I^{\mathrm{cut}}(c)|=0=h_{n_0}(c)$, as the entries are at most
$\beta$; for $c\leq0$, $|I^{\mathrm{cut}}(c)|=0=r-h_{n_0}(c)$, as the entries are at least $0$. If $1\leq c\leq\beta$ and
$c>T_{\nu^{\mathrm{cut}}}$, then $h_{n_0}(c)=0$ and $I^{\mathrm{cut}}(c)=\emptyset$; so $I^{\mathrm{cut}}(c)=\emptyset$ for
$|c|>T_{\nu^{\mathrm{cut}}}$. For the last condition of a standard filling, let $1\leq t'\leq r$. For $c\leq0$ the complement
of $I^{\mathrm{cut}}(c)=\emptyset$ is $\{1,\dots,r\}$, which has at least as many elements $\leq t'$ as any subset; and
$I^{\mathrm{cut}}(c+1)=\emptyset$ if $c+1>\beta$. Let $1\leq c<\beta$. As $a\leq t'$ if and only if $r-a\geq r-t'$, the
elements of $U(c)$ that are at most $t'$ form $U(c)\smallsetminus U_{r-t'}(c)$, whose number is
$(r-h_{n_0}(c))-(h_{n_0+r-t'}(c)-h_{n_0}(c))=r-h_{n_0+r-t'}(c)$ by Step 1. So $I^{\mathrm{cut}}(c)$ has
$t'-r+h_{n_0+r-t'}(c)$ elements at most $t'$, and this number does not increase with $c$, as $h_{n_0+r-t'}(c)$ does not.
If $\nu$ is a path for $B$, then $\nu^{\mathrm{cut}}_0=0$, its entries are dominant, its conditions at the letters
$i<n_0$ are those of $\nu$, and at the letters $i\geq n_0$, $\nu^{\mathrm{cut}}_{i+1}=\nu^{\mathrm{cut}}_i$ is a strip of size
$0$ over itself.

(4) For $i\leq n_0$, $\nu'_i=\nu'^{\mathrm{cut}}_i=\nu^{\mathrm{cut}}_i=\nu_i$. For $1\leq c\leq\beta$,
$U'(c)=\{1,\dots,r\}\smallsetminus I'^{\mathrm{cut}}(c)=U(c)$, so $U'_t(c)=U_t(c)$ and, by Step 1,
$h'_{n_0+t}(c)=h'_{n_0}(c)+|U'_t(c)|=h_{n_0+t}(c)$ for $t\leq r$. For the other $c$, $h'_{n_0+t}(c)$ and $h_{n_0+t}(c)$
are both $0$ or both $r$, as in Step 2. A dominant vector is determined by its numbers of entries at least $c$, as its
entry $a$ is at least $c$ if and only if $a$ is at most that number; so $\nu'_{n_0+t}=\nu_{n_0+t}$.
% Lean: Step 1 is card_riseA (riseA_succ_pos, riseA_succ_neg; the book's U(c) is riseA n0 nu c r with a = m + 1 and
%   tau_I m = n0 + r - a, U_t(c) is riseA n0 nu c t). Step 2 is Kcol_inactive and Psi_full (Ican_eq,
%   mem_colRange_of_le_entry). Step 3 is Kcol_cut, Kcol_supported and Psi_cut_supported for (2), and Psi_cut. Step 4 is
%   kill_Kcol_low and kill_Kcol_tau (kill_colStep, wedgeR, kill_X_not, kill_X_mem; offN is the set of the killed
%   variables). (1) is kill_Psi. (3) is isSt_cut (Icut_eq, pEnd_cut, card_Icut_filter with riseA_filter,
%   riseA_mono and colH_anti) and cut_ppath. (4) is path_eq_of_cut (riseA_eq_filter, eq_of_colH with le_iff_lt_colH).
\end{proof}

\begin{lemma}\label{pd-symbols-residue-count-inputs-segment-paths}
Let $C_{La}$ and $C_k$ be real numbers, and let $p$, $\mathcal R$, $\kappa$, $n'$, $r$, $K$, $k$, $T$, $L_1$, $N_1$,
$(J,f)$, $\deg$, $E'$, $N_s$, $E_0$, $m_*$, $a$, $Y$ and $q$ be as in Lemma~\ref{pd-symbols-residue-rank}, except that
$C_{La}$ and $C_k$ need not be at least $1088$; we use the notation of Definitions~\ref{pd-symbols-residue-coarse}
and~\ref{pd-symbols-residue-count-inputs-index}, and write $B(n)=\binom{n+k-1}{k-1}$. Let $j\in\mathbf Z$ and
$b\colon\mathbf Z\to\mathbf Z_{\geq0}$, and let $T^{\mathrm{all}}_j(b)$ be the set of the pairs $(\mu,I'')$ of a path $\mu$
of rank $k$ for some $B_\mu$ and a standard filling $I''$ of $\mu^{\mathrm{tn}}$ such that $B_\mu(i)=\beta_j(b,v_i)$ for
every letter $i$ with $v_i<f(j,0)$, and $\mu_{i+1}(m)=\beta_j(b)$ for $1\leq m\leq i'+1$ whenever $v_i=f(j,i')$ with
$i'<r$. Then $T^{\mathrm{all}}_j(b)$ is finite, it contains $T_j(b)$, and
\[
 |T^{\mathrm{all}}_j(b)|=\prod_{x\in\Lambda'_j}B(b(x))\prod_{i<r}B\big(\max(0,W_j(f(j-1,i))-m_0)\big).
\]
\end{lemma}
\leanref{Formalism/SharpRestriction/GrowthBounds/SegmentDegree/AtypicalPathSmallWeight.lean}{Formalism.MainNotions.SelectionLemmas.Pf.T.SAll}
\leanref{Formalism/SharpRestriction/GrowthBounds/SegmentDegree/AtypicalPathSmallWeight.lean}{Formalism.MainNotions.SelectionLemmas.Pf.T.rkA}
\leanref{Formalism/SharpRestriction/GrowthBounds/SegmentDegree/AtypicalPathSmallWeight.lean}{Formalism.MainNotions.SelectionLemmas.Pf.T.PathWeightEqRank}
\leanref{Formalism/SharpRestriction/GrowthBounds/SegmentDegree/ProdFromSelfEqOne.lean}{Formalism.MainNotions.SelectionLemmas.Pf.T.path_weight_eq_rank}
\leanref{Formalism/SharpRestriction/GrowthBounds/SegmentDegree/ProdFromSelfEqOne.lean}{Formalism.MainNotions.SelectionLemmas.Pf.T.pre_count}
\leanref{Formalism/SharpRestriction/GrowthBounds/SegmentDegree/DomOfStripDown.lean}{Formalism.MainNotions.SelectionLemmas.Pf.T.tau_count}
\leanref{Formalism/SharpRestriction/GrowthBounds/SegmentDegree/TailTGlueT.lean}{Formalism.MainNotions.SelectionLemmas.Pf.T.finite_SAll}
\leanref{Formalism/SharpRestriction/GrowthBounds/SegmentDegree/DistAddTwoLeOfSameVec.lean}{Formalism.MainNotions.SelectionLemmas.Pf.TE.letter_dictionary_bounds}

\begin{proof}
Let $N=|\Lambda_j|$ and $v_0<\dots<v_{N-1}$ the elements of $\Lambda_j$, lettered as in
Definition~\ref{pd-symbols-residue-pairs}, and put $n_0=N-r$, $B_0(i)=\beta_j(b,v_i)$ for $i<n_0$ and $\beta=\beta_j(b)$. By
Lemma~\ref{pd-symbols-residue-count-inputs-segment-plane}, $v_{n_0+t}=f(j,t)$ for $t<r$, the letters $n_0,\dots,N-1$ are
vector letters, and the letters $i<n_0$ are those with $v_i<f(j,0)$; so $\beta=\sum_{i<n_0}B_0(i)$. We use the function
$\operatorname{st}$ and the strips of Lemma~\ref{pd-symbols-residue-count-inputs-segment-clauses}.

For $0\leq s\leq N$ call a sequence $\mu=(\mu_0,\dots,\mu_N)$ in $\mathbf Z^k$ \emph{valid up to $s$} if $\mu_0=0$, $\mu_t$
is dominant for $t\leq s$, and for every letter $i<s$: if $i<n_0$, then $\mu_{i+1}$ is a strip of size $B_0(i)$ over
$\mu_i$ if $i$ is a vector letter, and $\mu_i$ is one over $\mu_{i+1}$ if it is a covector letter; if $i\geq n_0$, then
$\mu_{i+1}$ is a strip of some size over $\mu_i$ and $\mu_{i+1}(a)=\beta$ for $a\leq i-n_0+1$. As the letters $i\geq n_0$
are vector letters and $v_i=f(j,i-n_0)$, the paths of the pairs of $T^{\mathrm{all}}_j(b)$ are the sequences valid up
to $N$. For $\mu^0$ valid up to $s$ let $A(s,\mu^0)$ be the number of the pairs $(\mu,I'')\in T^{\mathrm{all}}_j(b)$ with
$\mu_t=\mu^0_t$ for $t\leq s$.

\emph{Step 1: finiteness.} Let $\mu$ be valid up to $s$. As in the proof of
Lemma~\ref{pd-symbols-residue-count-inputs-segment-plane}(1), $\mu_t(a)\leq\beta$ for $t\leq s$ and all $a$. At a letter
$l<n_0$, $\mu_{l+1}(a)\geq\mu_l(a)$ if $l$ is a vector letter and $\mu_{l+1}(a)\geq\mu_l(a)-B_0(l)$ if it is a covector
letter, and at the letters $l\geq n_0$, which are vector letters, $\mu_{l+1}(a)\geq\mu_l(a)$; so $\mu_t(a)\geq-\beta$.
Hence there are finitely many sequences valid up to $N$, and each $\mu^{\mathrm{tn}}$ has finitely many standard fillings,
as they are determined by their values at the $c$ with $|c|\leq T_{\mu^{\mathrm{tn}}}$. So $T^{\mathrm{all}}_j(b)$ is finite.
It contains $T_j(b)$, as a counting path satisfies the two conditions in its definition.

\emph{Step 2: the last $r$ letters.} We show $A(s,\mu^0)=\operatorname{st}(\tau_t(\mu^0_s))$ for $0\leq t\leq r$, $s=N-t$
and $\mu^0$ valid up to $s$, where $\tau_t(\lambda)=(\lambda(r-t+1),\dots,\lambda(k))\in\mathbf Z^{K+t}$ is the vector of the
last $K+t$ entries of $\lambda$. For $t=0$ the pairs counted are the $(\mu^0,I'')$ with $I''$ a standard filling of
$\mu^{0,\mathrm{tn}}$; by Lemma~\ref{pd-symbols-residue-count-inputs-segment-clauses}(2), applied to $\mu^{0,\mathrm{tn}}$
and to the sequence $(\mu^{0,\mathrm{tn}}_N)$ of length one, their number is $\operatorname{st}(\mu^{0,\mathrm{tn}}_N)$.

Let $t<r$, $s=N-t-1\geq n_0$, $\mu^0$ valid up to $s$ and $\lambda=\mu^0_s$. The letter $s$ is a vector letter. As in
Step~1, $\lambda(a)\leq\beta$ for all $a$; and $\lambda(a)=\beta$ for $a\leq r-t-1$, by the condition at the letter
$s-1$ if $s>n_0$, and trivially if $s=n_0$, as then $r-t-1=0$. Put $\nu=\tau_{t+1}(\lambda)=(\lambda(r-t),\dots,\lambda(k))$,
which is dominant, and for $\tau\in\mathbf Z^{K+t}$ let $g(\tau)\in\mathbf Z^k$ have the entries $\beta$ at
$a\leq r-t$ and $\tau(a-r+t)$ at $a>r-t$. We show that the $\lambda'$ for which $\mu^0$ with $\mu^0_{s+1}$ replaced by
$\lambda'$ is valid up to $s+1$ are the $g(\tau)$ with $\nu(l+1)\leq\tau(l)\leq\nu(l)$ for $1\leq l\leq K+t$. Such a
$\lambda'$ has $\lambda'(a)=\beta$ for $a\leq s-n_0+1=r-t$, so $\lambda'=g(\tau)$ with $\tau=\tau_t(\lambda')$, and the
strip conditions $\lambda(a)\leq\lambda'(a)$ and $\lambda'(a)\leq\lambda(a-1)$ at $a=r-t+l$ give
$\nu(l+1)\leq\tau(l)\leq\nu(l)$. Conversely, for such $\tau$, $g(\tau)$ is dominant, as $\tau(1)\leq\nu(1)\leq\beta$ and
$\tau(l+1)\leq\nu(l+1)\leq\tau(l)$; it is a strip over $\lambda$, as $\lambda(a)\leq\beta=g(\tau)(a)$ and, for $a<r-t$,
$g(\tau)(a+1)=\beta=\lambda(a)$, while the conditions at the other entries are those on $\tau$; and it has the entries
$\beta$ required at the letter $s$. As $\tau_t(g(\tau))=\tau$, the induction hypothesis gives
\[
 A(s,\mu^0)=\sum_\tau A(s+1,\mu^0\text{ with }\mu^0_{s+1}=g(\tau))=\sum_\tau\operatorname{st}(\tau)=\operatorname{st}(\nu),
\]
the sums over the $\tau$ above, the last equality by Lemma~\ref{pd-symbols-residue-count-inputs-segment-clauses}(3)
with $m=K+t\geq1$.

\emph{Step 3: the first $n_0$ letters.} We show $A(s,\mu^0)=\operatorname{st}(\mu^0_s)\prod_{s\leq i<n_0}B(B_0(i))$ for
$s\leq n_0$ and $\mu^0$ valid up to $s$, by descending induction on $s$. For $s=n_0$ this is Step~2 with $t=r$, as
$\tau_r(\lambda)=\lambda$. Let $s<n_0$ and $\lambda=\mu^0_s$. The $\lambda'$ that continue $\mu^0$ validly at $s+1$ are the
strips of size $B_0(s)$ over $\lambda$ if $s$ is a vector letter, and the $\lambda'$ over which $\lambda$ is a strip of size
$B_0(s)$ if it is a covector letter; they are dominant, as in the proof of
Lemma~\ref{pd-symbols-residue-count-inputs-segment-clauses}(4). By the induction hypothesis $A(s,\mu^0)$ is the sum of
the $\operatorname{st}(\lambda')\prod_{s<i<n_0}B(B_0(i))$ over these $\lambda'$, and by
Lemma~\ref{pd-symbols-residue-count-inputs-segment-clauses}(4) with $m=k-1$ and $b=B_0(s)$ this is
$\operatorname{st}(\lambda)\binom{B_0(s)+k-1}{k-1}\prod_{s<i<n_0}B(B_0(i))$.

\emph{Step 4: the count.} The sequence $0$ is valid up to $0$, so $|T^{\mathrm{all}}_j(b)|=A(0,0)=\operatorname{st}(0)
\prod_{i<n_0}B(B_0(i))$ by Step~3, and $\operatorname{st}(0)=1$ by
Lemma~\ref{pd-symbols-residue-count-inputs-segment-clauses}(2) for the sequence $(0)$ of length one. As in the proof of
Lemma~\ref{pd-symbols-residue-count-inputs-segment-members}, $v_i=f(j-1,i)$ for $i<r$, which is not new, so that
$B_0(i)=\max(W_j(f(j-1,i))-m_0,0)$; and the $v_i$ with $r\leq i<n_0$ are the elements of $\Lambda'_j$, which are new, so
that $B_0(i)=b(v_i)$. This gives the product of the statement.
% Lean: path_weight_eq_rank D admB3 ColumnFamilyLowerSet GenericProjectionCriterion, applied through the dictionary
%   TE.letter_dictionary_bounds (n0 = segN - r, B = Bdeg (restr j b) o segE, IsPATH iff IsPATHa by dict_early, dict_tau
%   and beS_eq; the product is its rank identity rkA = prod, with dict_w, dict_wE and LLset_eq_image). Valid up to s is
%   VP with StepOK, and sp_iff identifies SP with VP up to N; A(s, mu0) is the card of Af. Step 1 is finite_SAll
%   (path_le_beA, dminus_le_beA with neg_dminus_le_mu, finite_St). Step 2 is tau_count (wt_eq from admB3 at t = 0;
%   hheads, vp_le, hE, hterm with glueT and tailT_glueT, Af_card_sum, and ColumnFamilyLowerSet h4 at K + t). Step 3 is
%   pre_count (pieri_vec and pieri_cov from GenericProjectionCriterion, dom_of_strip_up, dom_of_strip_down, vp_update,
%   Af_step, Af_card_sum, prodFrom_succ). Step 4 is stCount_zero (admB3) and prodFrom.
\end{proof}

\begin{definition}\label{pd-symbols-residue-count-inputs-segment-typical}
Let the data of Definition~\ref{pd-symbols-residue-coarse} be given, with that notation, and let $j\in\mathbf Z$,
$b\colon\mathbf Z\to\mathbf Z_{\geq0}$, and $v_0<\dots<v_{N-1}$ the elements of $\Lambda_j$, lettered as in
Definition~\ref{pd-symbols-residue-pairs}; we use the notation of Definition~\ref{pd-symbols-residue-count-inputs-index}.
For a letter $i$ put $\Xi(i)=L_aa/(pk\,p^{\operatorname{dist}_j(v_i)+R_j})$, and for a letter $i\geq r$ let $o(i)$ be the
letter with $v_{o(i)}=\operatorname{old}_j(v_i)$. For $\lambda\in\mathbf Z^k$ put $\lambda^+=\sum_a\max(\lambda(a),0)$ and
$\lambda^-=\sum_a\max(-\lambda(a),0)$. A path $\mu$ of rank $k$ is \emph{typical} for $(j,b)$ if for every letter $i$:
\begin{enumerate}
\item $\mu_i^+-\mu_{i+1}^+\leq\Xi(i)$ if $i$ is a covector letter;
\item $\mu_i^--\mu_{i+1}^-\leq\Xi(i)$ if $i$ is a vector letter and $i\geq r$;
\item $\frac23\beta_j(b,v_{o(i)})\leq\mu_i(r)$ if $i\geq r$.
\end{enumerate}
\end{definition}
\leanref{Formalism/SharpRestriction/GrowthBounds/SegmentDegree/AtypicalPathSmallWeight.lean}{Formalism.MainNotions.SelectionLemmas.Pf.T.posS}
\leanref{Formalism/SharpRestriction/GrowthBounds/SegmentDegree/AtypicalPathSmallWeight.lean}{Formalism.MainNotions.SelectionLemmas.Pf.T.negS}
\leanref{Formalism/SharpRestriction/GrowthBounds/SegmentDegree/AtypicalPathSmallWeight.lean}{Formalism.MainNotions.SelectionLemmas.Pf.T.Xi}
\leanref{Formalism/SharpRestriction/GrowthBounds/SegmentDegree/AtypicalPathSmallWeight.lean}{Formalism.MainNotions.SelectionLemmas.Pf.T.Typ}
\leanref{Formalism/SharpRestriction/GrowthBounds/SegmentDegree/SortGetDSpec.lean}{Formalism.MainNotions.SelectionLemmas.Pf.TE.exists_od}

\begin{definition}\label{pd-symbols-residue-count-inputs-segment-atypical-setting}
Let $r,K\geq1$ and $n_0\geq0$ be integers, $k=r+K$ and $N=n_0+r$, fix a lettering of $\{0,\dots,N-1\}$, and let
$B\colon\{0,\dots,n_0-1\}\to\mathbf Z_{\geq0}$, $o\colon\{r,\dots,N-1\}\to\{0,\dots,N-1\}$ and
$\gamma\colon\{r,\dots,N-1\}\to\mathbf R$ be functions, $a$ and $L_a$ real numbers and $p$ an integer. Put
$\beta=\sum_{i<n_0}B(i)$, $\operatorname{rk}(B)=\prod_{i<n_0}\binom{B(i)+k-1}{k-1}$ and $\ell_0=\log(1+2aN)$.
\begin{enumerate}
\item These data form an \emph{atypical setting} if:
  \begin{enumerate}
  \item $k\leq B(i)\leq2a$ for $i<n_0$;
  \item $p^2B(i)\leq5B(i')$ for letters $i<i'<n_0$ of the same type;
  \item the letters $i<r$ and the letters $i\geq n_0-r$ are vector letters;
  \item $p\geq13$, $a\geq3$, $L_a>0$, $2r\leq n_0$, $N+r\leq K$, $N\leq a$ and $k\leq2(K-2N)$;
  \item for $r\leq i<N$, $o(i)<i$ is a vector letter, and the number of vector letters $l$ with $o(i)\leq l<i$ is $r$;
  \item for $r\leq i<N$, $\gamma(i)>0$, and $pB(l)\leq4\gamma(i)$ for every letter $l<\min(i,n_0)$.
  \end{enumerate}
  In an atypical setting $o(i)<n_0$ for $r\leq i<N$: otherwise the $r$ vector letters $l$ with $o(i)\leq l<i$ would lie
  in $\{n_0,\dots,i-1\}$, which has fewer than $r$ elements as $i<N$. By (c), every covector letter $i$ satisfies
  $r\leq i<n_0-r$.
\item The \emph{paths} of the setting are the paths $\mu$ of rank $k$ for some $B_\mu$
  (Definition~\ref{pd-symbols-path-generators}(2)) with $B_\mu(i)=B(i)$ for $i<n_0$ and $\mu_{i+1}(m)=\beta$ for
  $1\leq m\leq i-n_0+1$ whenever $i\geq n_0$. Their set is $\mathcal P$.
\item For $\lambda\in\mathbf Z^m$ and $u\geq0$ put $\lambda^+_{>u}=\sum_{u<c\leq m}\max(\lambda(c),0)$,
  $\lambda^-_{>u}=\sum_{u<c\leq m}\max(-\lambda(m+1-c),0)$ and $\lambda^\pm=\lambda^\pm_{>0}$. For $r\leq i<N$ put
  $\Xi(i)=L_a\gamma(i)/(pk)$. A path $\mu\in\mathcal P$ is \emph{typical} if for every letter $i$:
  $\mu_i^+-\mu_{i+1}^+\leq\Xi(i)$ if $i$ is a covector letter; $\mu_i^--\mu_{i+1}^-\leq\Xi(i)$ if $i$ is a vector letter
  and $i\geq r$; and $\frac23B(o(i))\leq\mu_i(r)$ if $i\geq r$. It is \emph{good} if for every letter $i\geq r$:
  $\frac23B(o(i))\leq\mu_i(r)$; $\sum_{c=r+1}^k(\mu_{i+1}(c)-\mu_i(c))\leq\Xi(i)+2B(o(i))$ if $i$ is a vector letter;
  and $\sum_{c=1}^r(\mu_i(c)-\mu_{i+1}(c))\leq\Xi(i)$ if $i$ is a covector letter.
\end{enumerate}
\end{definition}
\leanref{Formalism/SharpRestriction/GrowthBounds/SegmentDegree/AtypicalPathSmallWeight.lean}{Formalism.MainNotions.SelectionLemmas.Pf.T.HypT}
\leanref{Formalism/SharpRestriction/GrowthBounds/SegmentDegree/AtypicalPathSmallWeight.lean}{Formalism.MainNotions.SelectionLemmas.Pf.T.IsPATHa}
\leanref{Formalism/SharpRestriction/GrowthBounds/SegmentDegree/AtypicalPathSmallWeight.lean}{Formalism.MainNotions.SelectionLemmas.Pf.T.SP}
\leanref{Formalism/SharpRestriction/GrowthBounds/SegmentDegree/AtypicalPathSmallWeight.lean}{Formalism.MainNotions.SelectionLemmas.Pf.T.beA}
\leanref{Formalism/SharpRestriction/GrowthBounds/SegmentDegree/AtypicalPathSmallWeight.lean}{Formalism.MainNotions.SelectionLemmas.Pf.T.rkA}
\leanref{Formalism/SharpRestriction/GrowthBounds/SegmentDegree/AtypicalPathSmallWeight.lean}{Formalism.MainNotions.SelectionLemmas.Pf.T.ell0}
\leanref{Formalism/SharpRestriction/GrowthBounds/SegmentDegree/AtypicalPathSmallWeight.lean}{Formalism.MainNotions.SelectionLemmas.Pf.T.GoodA}
\leanref{Formalism/SharpRestriction/GrowthBounds/SegmentDegree/TypicalPathOrdersSmall.lean}{Formalism.MainNotions.SelectionLemmas.Pf.T.oldest_lt_n0}

\begin{lemma}\label{pd-symbols-residue-count-inputs-segment-atypical-rows}
Let $N\geq0$ and $n\geq1$ be integers with a lettering of $\{0,\dots,N-1\}$, $B\colon\{0,\dots,N-1\}\to\mathbf
Z_{\geq0}$, and $\mu$ a path of rank $n$ for $B$ (Definition~\ref{pd-symbols-path-generators}(2)). For $0\leq s\leq N$
let $d^+(s)$ and $d^-(s)$ be the sums of the $B(l)$ over the vector letters $l<s$ and over the covector letters $l<s$.
Then for $0\leq s\leq N$ and $1\leq m\leq n$:
\begin{enumerate}
\item $\mu_s(m)\geq0$ if $m+s\leq n$, and $\mu_s(m)\leq0$ if $m>s$;
\item $-d^-(s)\leq\mu_s(m)\leq d^+(s)$;
\item if $t\leq s$ and $v$ is the number of vector letters $l$ with $t\leq l<s$, then $\mu_s(m+v)\leq\mu_t(m)$ whenever
  $m+v\leq n$.
\end{enumerate}
\end{lemma}
\leanref{Formalism/SharpRestriction/GrowthBounds/SegmentDegree/NegDminusLeMu.lean}{Formalism.MainNotions.SelectionLemmas.Pf.T.mu_nonneg}
\leanref{Formalism/SharpRestriction/GrowthBounds/SegmentDegree/NegDminusLeMu.lean}{Formalism.MainNotions.SelectionLemmas.Pf.T.mu_nonpos}
\leanref{Formalism/SharpRestriction/GrowthBounds/SegmentDegree/NegDminusLeMu.lean}{Formalism.MainNotions.SelectionLemmas.Pf.T.dplus}
\leanref{Formalism/SharpRestriction/GrowthBounds/SegmentDegree/NegDminusLeMu.lean}{Formalism.MainNotions.SelectionLemmas.Pf.T.dminus}
\leanref{Formalism/SharpRestriction/GrowthBounds/SegmentDegree/NegDminusLeMu.lean}{Formalism.MainNotions.SelectionLemmas.Pf.T.mu_le_dplus}
\leanref{Formalism/SharpRestriction/GrowthBounds/SegmentDegree/NegDminusLeMu.lean}{Formalism.MainNotions.SelectionLemmas.Pf.T.neg_dminus_le_mu}
\leanref{Formalism/SharpRestriction/GrowthBounds/SegmentDegree/NegDminusLeMu.lean}{Formalism.MainNotions.SelectionLemmas.Pf.T.row_back}

\begin{proof}
At a vector letter $s$, $\mu_{s+1}$ is a strip of size $B(s)$ over $\mu_s$: $\mu_s(c)\leq\mu_{s+1}(c)$,
$\mu_{s+1}(c+1)\leq\mu_s(c)$, and, as the nonnegative differences $\mu_{s+1}(c)-\mu_s(c)$ add up to $B(s)$,
$\mu_{s+1}(c)\leq\mu_s(c)+B(s)$. At a covector letter the roles of $\mu_s$ and $\mu_{s+1}$ are exchanged. Each part is
proved by induction on $s$, starting from $\mu_0=0$, and for (3) from $s=t$.

(1) Let $s$ be a vector letter. If $m+s+1\leq n$, then $\mu_{s+1}(m)\geq\mu_s(m)\geq0$. If $m>s+1$, then
$\mu_{s+1}(m)\leq\mu_s(m-1)\leq0$. Let $s$ be a covector letter. If $m+s+1\leq n$, then $\mu_{s+1}(m)\geq\mu_s(m+1)\geq0$.
If $m>s+1$, then $\mu_{s+1}(m)\leq\mu_s(m)\leq0$.

(2) At a vector letter $s$, $\mu_s(m)\leq\mu_{s+1}(m)\leq\mu_s(m)+B(s)$; there $d^+$ grows by $B(s)$ and $d^-$ does not
change. At a covector letter $s$, $\mu_s(m)-B(s)\leq\mu_{s+1}(m)\leq\mu_s(m)$, and $d^-$ grows by $B(s)$.

(3) At a covector letter $s\geq t$, $v$ does not change and $\mu_{s+1}(m+v)\leq\mu_s(m+v)$. At a vector letter $s\geq t$,
$v$ grows by one and $\mu_{s+1}(m+v+1)\leq\mu_s(m+v)$.
% Lean: (1) is mu_nonneg (stated there for m + s < n, 0-based m) and mu_nonpos; (2) is mu_le_dplus (strip_le_add,
%   dsum_succ) and neg_dminus_le_mu; (3) is row_back (vcount, vcount_succ). The Lean states them for PPath.
\end{proof}

\begin{lemma}\label{pd-symbols-residue-count-inputs-segment-atypical-orders}
Let $r$, $K$, $n_0$, a lettering of $\{0,\dots,n_0+r-1\}$, $B$, $o$, $\gamma$, $a$, $L_a$ and $p$ form an atypical
setting (Definition~\ref{pd-symbols-residue-count-inputs-segment-atypical-setting}), with its notation. Then every
typical path in $\mathcal P$ is good.
\end{lemma}
\leanref{Formalism/SharpRestriction/GrowthBounds/SegmentDegree/AtypicalPathSmallWeight.lean}{Formalism.MainNotions.SelectionLemmas.Pf.T.TypicalPathOrdersSmall}
\leanref{Formalism/SharpRestriction/GrowthBounds/SegmentDegree/TypicalPathOrdersSmall.lean}{Formalism.MainNotions.SelectionLemmas.Pf.T.typical_path_orders_small}
\leanref{Formalism/SharpRestriction/GrowthBounds/SegmentDegree/TypicalPathOrdersSmall.lean}{Formalism.MainNotions.SelectionLemmas.Pf.T.nu_last_le}
\leanref{Formalism/SharpRestriction/GrowthBounds/SegmentDegree/TypicalPathOrdersSmall.lean}{Formalism.MainNotions.SelectionLemmas.Pf.T.dplus_le_geom}

\begin{proof}
Let $\mu\in\mathcal P$ be typical, a path for $B_\mu$, and let $i\geq r$ be a letter; we use $d^+$ of
Lemma~\ref{pd-symbols-residue-count-inputs-segment-atypical-rows}, formed with $B_\mu$. The first condition of a good
path is the third of a typical one. As $r+N\leq K<k$, part (1) of that lemma gives $\mu_s(c)\geq0$ for all $s\leq N$ and
$c\leq r$. Put $x=\mu_i$ and $y=\mu_{i+1}$.

\emph{Step 1: $\mu_i(r)\leq2B(o(i))$.} We show first that $d^+(m)\leq10p^{-2}B(l)$ for every vector letter $l<n_0$ and
$m\leq l$, by induction on $m$; as $B_\mu=B$ below $n_0$, $d^+(m)$ is then the sum of the $B(l')$ over the vector letters
$l'<m$. For $m=0$ the sum is empty. Let $m<l$. If $m$ is a covector letter, $d^+(m+1)=d^+(m)$. If $m$ is a vector letter,
the induction hypothesis for the letter $m$ and condition (b) give
\[
 d^+(m+1)=d^+(m)+B(m)\leq(1+10p^{-2})B(m)\leq(1+10p^{-2})\,5p^{-2}B(l)\leq10p^{-2}B(l),
\]
as $p\geq13$. Now $o(i)<n_0$ is a vector letter, and there are $r-1$ vector letters $l$ with $o(i)<l<i$. By parts (3) and
(2) of Lemma~\ref{pd-symbols-residue-count-inputs-segment-atypical-rows}, with $t=o(i)+1$, $s=i$ and $m=1$,
\[
 \mu_i(r)\leq\mu_{o(i)+1}(1)\leq d^+(o(i)+1)=d^+(o(i))+B(o(i))\leq(1+10p^{-2})B(o(i))\leq2B(o(i)).
\]

\emph{Step 2: a vector letter $i$.} Then $y$ is a strip over $x$. As the entries at $c\leq r$ are nonnegative, $x^-$ and
$y^-$ are the sums of $\max(-x(c),0)$ and $\max(-y(c),0)$ over $c>r$, and $z=\max(z,0)-\max(-z,0)$ gives
\[
 \sum_{c=r+1}^k\big(y(c)-x(c)\big)=\sum_{c=r+1}^k\max(y(c),0)-\sum_{c=r+1}^k\max(x(c),0)+x^--y^-.
\]
As $y(c+1)\leq x(c)$ for $c<k$, $\sum_{c=r+2}^k\max(y(c),0)\leq\sum_{c=r+1}^k\max(x(c),0)$, and
$\max(y(r+1),0)\leq x(r)$ as $x(r)\geq0$. So the left side is at most $x(r)+x^--y^-$, which is at most
$2B(o(i))+\Xi(i)$ by Step 1 and the second condition of a typical path.

\emph{Step 3: a covector letter $i$.} Then $x$ is a strip over $y$, so $y(c)\leq x(c)$ for all $c$. As the entries at
$c\leq r$ are nonnegative,
\[
 x^+-y^+=\sum_{c=1}^r\big(x(c)-y(c)\big)+\sum_{c=r+1}^k\big(\max(x(c),0)-\max(y(c),0)\big)\geq\sum_{c=1}^r\big(x(c)-y(c)\big),
\]
and $x^+-y^+\leq\Xi(i)$ by the first condition of a typical path.
% Lean: typical_path_orders_small: F1 is mu_nonneg; Step 1 is nu_last_le (oldest_lt_n0, vcount_od, row_back,
%   mu_le_dplus, dsum_succ, dplus_le_geom with H.H2 and p >= 13); Step 2 is the case hv = true (step_vec, negS_split,
%   sum_shift_le, int_pos_neg); Step 3 is the case hv = false (step_cov, posS_split).
\end{proof}

\begin{lemma}\label{pd-symbols-residue-count-inputs-segment-atypical-count}
Let $r$, $K$, $n_0$, a lettering of $\{0,\dots,n_0+r-1\}$, $B$, $o$, $\gamma$, $a$, $L_a$ and $p$ form an atypical
setting (Definition~\ref{pd-symbols-residue-count-inputs-segment-atypical-setting}), with its notation. Then
$\mathcal P$ is finite and $|\mathcal P|\leq\exp(2N^2\ell_0)$.
\end{lemma}
\leanref{Formalism/SharpRestriction/GrowthBounds/SegmentDegree/AtypicalPathSmallWeight.lean}{Formalism.MainNotions.SelectionLemmas.Pf.T.PathCountFinite}
\leanref{Formalism/SharpRestriction/GrowthBounds/SegmentDegree/WindowSmall.lean}{Formalism.MainNotions.SelectionLemmas.Pf.T.path_count_finite}
\leanref{Formalism/SharpRestriction/GrowthBounds/SegmentDegree/WindowSmall.lean}{Formalism.MainNotions.SelectionLemmas.Pf.T.path_le_beA}

\begin{proof}
By condition (d), $r+K=k\leq2(K-2N)$, so $K\geq r+4N$ and $k\geq4N$. Let $\mu\in\mathcal P$, a path for $B_\mu$, and
$1\leq s\leq N$; we use $d^\pm$ of Lemma~\ref{pd-symbols-residue-count-inputs-segment-atypical-rows}, formed with
$B_\mu$. By (c) every covector letter is smaller than $n_0$, where $B_\mu=B$; so $d^-(s)\leq\beta$, and
$\mu_s(m)\geq-\beta$ by part (2) of that lemma. Also $\mu_s(m)\leq\beta$: for $s\leq n_0$ by part (2), as
$d^+(s)\leq\beta$; for $s>n_0$ by induction on $s$, as the letter $s-1\geq n_0$ is a vector letter, $\mu_s(1)=\beta$ by
the definition of $\mathcal P$, and $\mu_s(m)\leq\mu_{s-1}(m-1)$ for $m\geq2$. By part (1), $\mu_s(m)\geq0$ for
$m\leq k-N$ and $\mu_s(m)\leq0$ for $m>N$. So $\mu_s(m)=0$ for $N<m\leq k-N$, the first $N$ entries of $\mu_s$ lie in
$\{0,\dots,\beta\}$ and its last $N$ entries in $\{-\beta,\dots,0\}$. As $\mu_0=0$, $\mu$ is determined by these $2N$
entries of $\mu_1,\dots,\mu_N$, so $\mathcal P$ is finite with at most $(\beta+1)^{2N^2}$ elements. By (a),
$\beta\leq2an_0\leq2aN$, and $(1+2aN)^{2N^2}=\exp(2N^2\ell_0)$.
% Lean: path_count_finite: two_N_lt_K (H.H4); hshape from mu_nonneg, mu_nonpos, path_le_beA (mu_le_dplus,
%   dplus_le_beA, the head at the late vector letters, step_vec) and neg_dminus_le_mu with dminus_le_beA; the
%   injective code into Fin N x Fin (2N) -> {0, ..., beA}; beA_le.
\end{proof}

\begin{lemma}\label{pd-symbols-residue-count-inputs-segment-atypical-tails}
Let $r$, $K$, $n_0$, a lettering of $\{0,\dots,n_0+r-1\}$, $B$, $o$, $\gamma$, $a$, $L_a$ and $p$ form an atypical
setting (Definition~\ref{pd-symbols-residue-count-inputs-segment-atypical-setting}), with its notation, and let
$\mu\in\mathcal P$. For $0\leq s\leq N$ put $h_s=\max(s-n_0,0)$ and
$\hat\mu_s=(\mu_s(h_s+1),\dots,\mu_s(k))\in\mathbf Z^{k-h_s}$; let $l^+_s(1)>\dots>l^+_s(q^+_s)$ be the vector letters
smaller than $\min(s,n_0)$, and put $L^+_s(t)=B(l^+_s(t))$ and $L^+_{s,>u}=\sum_{u<t\leq q^+_s}L^+_s(t)$ for $u\geq0$;
define $l^-_s(t)$, $q^-_s$, $L^-_s(t)$ and $L^-_{s,>u}$ in the same way with the covector letters. Then:
\begin{enumerate}
\item $(\hat\mu_s)^+_{>u}\leq L^+_{s,>u}$ and $(\hat\mu_s)^-_{>u}\leq L^-_{s,>u}$ for $0\leq s\leq N$ and $u\geq0$;
\item let $0\leq s_0\leq N$, $u\geq0$ and $X\geq0$ satisfy $(\hat\mu_{s_0})^+_{>u}\leq L^+_{s_0,>u}-X$, and let $c$ be
  the number of vector letters $l$ with $s_0\leq l<n_0$; then $(\hat\mu_N)^+_{>u+c}\leq L^+_{N,>u+c}-X$. The same holds
  with $-$ in place of $+$ and covector letters in place of vector letters.
\end{enumerate}
\end{lemma}
\leanref{Formalism/ThetaLogCoordinates/ResidueClassEstimates/StatementT.lean}{Formalism.InvariantFactorPaths.BranchingDimensions.PathTailBounds}
\leanref{Formalism/ThetaLogCoordinates/ResidueClassEstimates/VandermondeRankEstimates/N0LeNs.lean}{Formalism.InvariantFactorPaths.BranchingDimensions.pathTailBounds}
\leanref{Formalism/ThetaLogCoordinates/ResidueClassEstimates/VandermondeRankEstimates/N0LeNs.lean}{Formalism.InvariantFactorPaths.BranchingDimensions.abs_bound}
\leanref{Formalism/ThetaLogCoordinates/ResidueClassEstimates/VandermondeRankEstimates/N0LeNs.lean}{Formalism.InvariantFactorPaths.BranchingDimensions.deficit_persists_of_active_inactive_recursion}
\leanref{Formalism/ThetaLogCoordinates/ResidueClassEstimates/Letters.lean}{Formalism.InvariantFactorPaths.BranchingDimensions.tailPos}
\leanref{Formalism/ThetaLogCoordinates/ResidueClassEstimates/Letters.lean}{Formalism.InvariantFactorPaths.BranchingDimensions.tailNeg}
\leanref{Formalism/ThetaLogCoordinates/ResidueClassEstimates/Letters.lean}{Formalism.InvariantFactorPaths.BranchingDimensions.Letters.stage}
\leanref{Formalism/ThetaLogCoordinates/ResidueClassEstimates/Letters.lean}{Formalism.InvariantFactorPaths.BranchingDimensions.Letters.latestV}
\leanref{Formalism/ThetaLogCoordinates/ResidueClassEstimates/Letters.lean}{Formalism.InvariantFactorPaths.BranchingDimensions.Letters.SPos}

\begin{proof}
Call a letter $s$ \emph{active for $+$} if it is a vector letter and $s<n_0$, and \emph{active for $-$} if it is a
covector letter; by (c), covector letters are smaller than $n_0$. For a letter $s$ write $x=\hat\mu_s$ and
$y=\hat\mu_{s+1}$. We show, for $u\geq0$:
\begin{enumerate}
\item[(i)] if $s$ is active for $+$, then $y^+_{>u+1}\leq x^+_{>u}$, $L^+_{s+1,>u+1}=L^+_{s,>u}$ and
  $y^+-x^+\leq L^+_{s+1,>0}-L^+_{s,>0}$;
\item[(ii)] if $s$ is not active for $+$, then $y^+_{>u}\leq x^+_{>u}$ and $L^+_{s+1,>u}=L^+_{s,>u}$;
\end{enumerate}
and the same for $-$. For $\lambda\in\mathbf Z^m$ let $\lambda^\vee=(-\lambda(m),\dots,-\lambda(1))$, so that
$\lambda^-_{>u}=(\lambda^\vee)^+_{>u}$.

Let $s<n_0$; then $x=\mu_s$ and $y=\mu_{s+1}$. If $s$ is a vector letter, $y$ is a strip of size $B(s)$ over $x$, and
the list $l^+_{s+1}$ is $l^+_s$ with $s$ put in front, so $L^+_{s+1,>u+1}=L^+_{s,>u}$ and
$L^+_{s+1,>0}-L^+_{s,>0}=B(s)$. As $y(c+1)\leq x(c)$, $y^+_{>u+1}\leq x^+_{>u}$. As $x\leq y$ entrywise, $y^-\leq x^-$,
so $y^+-x^+=\sum_cy(c)-\sum_cx(c)+y^--x^-\leq B(s)$. This is (i) for $+$; for $-$, $s$ is not active, the list $l^-$
does not change, and $-y\leq-x$ gives $y^-_{>u}\leq x^-_{>u}$, which is (ii). If $s$ is a covector letter, $x$ is a strip
of size $B(s)$ over $y$, and then $y^\vee$ is a strip of size $B(s)$ over $x^\vee$. The same arguments give (i) for $-$
and (ii) for $+$.

Let $s\geq n_0$. Then $s$ is a vector letter by (c), active for neither sign, $\min(s,n_0)=\min(s+1,n_0)$, so the lists
do not change, and $h_{s+1}=h_s+1$. As $\mu_{s+1}$ is a strip over $\mu_s$,
\[
 y(c)=\mu_{s+1}(h_s+1+c)\leq\mu_s(h_s+c)=x(c),\qquad y(c)\geq\mu_s(h_s+1+c)=x(c+1),
\]
for $1\leq c\leq k-h_s-1$. The first inequality gives $y^+_{>u}\leq x^+_{>u}$. With $m=k-h_s-1$, the second gives
$\max(-y(m+1-c),0)\leq\max(-x(m+2-c),0)$, and the right side is the $c$-th term of $x^-_{>u}$; so
$y^-_{>u}\leq x^-_{>u}$. This is (ii) for both signs.

(1) For $s=0$, $\mu_0=0$ and the tails vanish. If $s$ is not active, (ii) passes the bound from $s$ to $s+1$. If $s$ is
active and $u\geq1$, then $y^+_{>u}\leq x^+_{>u-1}\leq L^+_{s,>u-1}=L^+_{s+1,>u}$; and for $u=0$,
$y^+\leq x^++L^+_{s+1,>0}-L^+_{s,>0}\leq L^+_{s+1,>0}$. The same for $-$.

(2) For $s_0\leq s\leq N$ let $c(s)$ be the number of active letters $l$ with $s_0\leq l<s$. We show
$(\hat\mu_s)^+_{>u+c(s)}\leq L^+_{s,>u+c(s)}-X$ by induction on $s$; for $s=s_0$ it is the hypothesis. If $s$ is active,
$c(s+1)=c(s)+1$ and (i) gives
$(\hat\mu_{s+1})^+_{>u+c(s)+1}\leq(\hat\mu_s)^+_{>u+c(s)}\leq L^+_{s,>u+c(s)}-X=L^+_{s+1,>u+c(s)+1}-X$. If not,
$c(s+1)=c(s)$ and (ii) gives the bound. As $c(N)=c$, this is (2). The same for $-$.
% Lean: pathTailBounds for BranchingDimensions.Letters with 1-based letters; the book's hat mu_s is stage mu s (of
%   length stLen s), its lists are latestV and latestC with sums SPos and SNeg, and the tails are tailPos and tailNeg
%   (tailNeg_eq through revW). (i) is pos_A and pos_active_zero (HStrip.tailPos_shift, HStrip.tailPos_zero through
%   wsize_eq_tails) and neg_A, neg_active_zero; (ii) is pos_I and neg_I, with frame_step (interl_of_path,
%   Interl.tailPos_le) for s >= n0. (1) is abs_bound and (2) is deficit_persists_of_active_inactive_recursion
%   (with card_actV, card_actC, latestV_frame).
\end{proof}

\begin{lemma}\label{pd-symbols-residue-count-inputs-segment-atypical-log}
For integers $k'\geq0$ and real $y\geq0$ put $f_{k'}(y)=\sum_{j=1}^{k'}\log((y+j)/j)$, and for $n\geq1$ and
$\lambda\in\mathbf Z^n$ put $W_n(\lambda)=\prod_{1\leq a<b\leq n}(\lambda(a)-\lambda(b)+b-a)/(b-a)$.
\begin{enumerate}
\item Let $k'\geq1$ and $L\geq0$ be integers, and $\alpha_1,\dots,\alpha_L\geq0$ and $B_1,\dots,B_L$ real numbers with
  $B_t\geq k'$ for all $t$, $4B_{t+1}\leq B_t$ for $t<L$, and $\sum_{t\geq u}\alpha_t\leq\sum_{t\geq u}B_t$ for
  $1\leq u\leq L$. Then $\sum_tf_{k'}(\alpha_t)\leq\sum_tf_{k'}(B_t)$. If moreover $1\leq u_0\leq L$ and $X\geq0$ satisfy
  $\sum_{t\geq u_0}\alpha_t\leq\sum_{t\geq u_0}B_t-X$, then $\sum_tf_{k'}(\alpha_t)-\sum_tf_{k'}(B_t)\leq-Xk'/(4B_{u_0})$.
\item Let $1\leq r\leq N$ and $K\geq2N$ be integers, $k=r+K$, $k'=K-2N$, $A_0\geq0$ a real number and $B_1,\dots,B_{N-r}$
  integers in $[0,A_0]$. Then
  \[
   0\leq\log\prod_{i=1}^{N-r}\binom{B_i+k-1}{k-1}-\sum_{i=1}^{N-r}f_{k'}(B_i)\leq N(r+2N)\log(1+A_0).
  \]
\item Let $r,K,N\geq0$ and $n\geq1$ be integers with $2N\leq K\leq n\leq r+K$, $k'=K-2N$, $\lambda\in\mathbf Z^n$ with
  $\lambda(1)\geq\dots\geq\lambda(n)$ and with at most $N$ positive and at most $N$ negative entries, and $A_0$ a real
  number with $\max(\lambda(1),0)+\max(-\lambda(n),0)\leq A_0$. Then $W_n(\lambda)>0$ and
  \[
   -2N^2\log(1+A_0)\leq\log W_n(\lambda)-\sum_{a=1}^nf_{k'}\big(\max(\lambda(a),0)\big)-\sum_{a=1}^nf_{k'}\big(\max(-\lambda(a),0)\big)
   \leq(2Nr+8N^2)\log(1+A_0).
  \]
\end{enumerate}
\end{lemma}
\leanref{Formalism/ThetaLogCoordinates/ResidueClassEstimates/Letters.lean}{Formalism.InvariantFactorPaths.BranchingDimensions.fT}
\leanref{Formalism/ThetaLogCoordinates/ResidueClassEstimates/Letters.lean}{Formalism.InvariantFactorPaths.BranchingDimensions.weylDim}
\leanref{Formalism/ThetaLogCoordinates/ResidueClassEstimates/StatementT.lean}{Formalism.InvariantFactorPaths.BranchingDimensions.LogBinomialDeficit}
\leanref{Formalism/ThetaLogCoordinates/ResidueClassEstimates/VandermondeRankEstimates/LogBinomialProperties.lean}{Formalism.InvariantFactorPaths.BranchingDimensions.logBinomialDeficit}
\leanref{Formalism/ThetaLogCoordinates/ResidueClassEstimates/StatementT.lean}{Formalism.InvariantFactorPaths.BranchingDimensions.LogRankEstimate}
\leanref{Formalism/ThetaLogCoordinates/ResidueClassEstimates/VandermondeRankEstimates/LogBinomialProperties.lean}{Formalism.InvariantFactorPaths.BranchingDimensions.logRankEstimate}
\leanref{Formalism/ThetaLogCoordinates/ResidueClassEstimates/StatementT.lean}{Formalism.InvariantFactorPaths.BranchingDimensions.LogWeylDimensionEstimate}
\leanref{Formalism/ThetaLogCoordinates/ResidueClassEstimates/VandermondeRankEstimates/LogWeylDimensionEstimate.lean}{Formalism.InvariantFactorPaths.BranchingDimensions.logWeylDimensionEstimate}

\begin{proof}
These are self-contained computations with $f'_{k'}(y)=\sum_{j=1}^{k'}1/(y+j)$, the derivative of $f_{k'}$ on
$[0,\infty)$.

(1) As $f_{k'}$ is concave, $f_{k'}(\alpha_t)-f_{k'}(B_t)\leq f'_{k'}(B_t)(\alpha_t-B_t)$. Summation by parts writes
$\sum_tf'_{k'}(B_t)(\alpha_t-B_t)$ as $\sum_ug_u\sum_{t\geq u}(\alpha_t-B_t)$, with $g_u=f'_{k'}(B_u)-f'_{k'}(B_{u-1})$ and
$f'_{k'}(B_0)=0$. As $f'_{k'}(y')\leq f'_{k'}(y)/2$ for $k'\leq y$ and $4y\leq y'$, $g_u\geq f'_{k'}(B_u)/2\geq0$; so
every term is at most $0$, and the term $u_0$ is at most $-Xf'_{k'}(B_{u_0})/2$. Finally $f'_{k'}(y)\geq k'/(2y)$ for
$y\geq k'$.

(2) For an integer $B\geq0$, $\log\binom{B+k-1}{k-1}=f_{k-1}(B)$, and $0\leq f_{k'+d}(y)-f_{k'}(y)\leq d\log(1+y)$ for
$d\geq0$; here $d=k-1-k'\leq r+2N$, and there are $N-r\leq N$ factors.

(3) $\log W_n(\lambda)$ is the sum over $a<b$ of $\log\big((\lambda(a)-\lambda(b)+b-a)/(b-a)\big)$, which lies between
$0$ and $\log(1+A_0)$. At most $(2N)^2$ pairs have $\lambda(a)\neq0\neq\lambda(b)$. For $\lambda(a)>0$, the $b>a$ with
$\lambda(b)=0$ form a run whose distances $b-a$ cover $\{1,\dots,k'\}$ up to at most $N$ missing and at most $r+2N$
extra values, so their terms add up to $f_{k'}(\lambda(a))$ up to $-N\log(1+A_0)$ and $+(r+2N)\log(1+A_0)$; the same
holds for $\lambda(b)<0$ and the $a<b$ with $\lambda(a)=0$. There are at most $N$ positive and $N$ negative entries.
% Lean: (1) is logBinomialDeficit (fT_concave, abel_tails, sum_incr, fT'_half, fT'_nonneg, fT'_lower2); (2) is
%   logRankEstimate (fT_nat, fT_add_bounds); (3) is logWeylDimensionEstimate (log_weylDim_lemma_4, G_split into
%   pa, pb and pc, sum_pa_le, pos_row and neg_row with sum_hh_le and le_sum_hh). The Lean states (2) for
%   BranchingDimensions.Letters (rankAA, n0 = Ns - r) and numbers the entries from 0. LogBinomialProperties, the
%   remaining input of pathWeightBound_of, is used only through fT k' 0 = 0.
\end{proof}

\begin{lemma}\label{pd-symbols-residue-count-inputs-segment-atypical-weight}
Let $r$, $K$, $n_0$, a lettering of $\{0,\dots,n_0+r-1\}$, $B$, $o$, $\gamma$, $a$, $L_a$ and $p$ form an atypical
setting (Definition~\ref{pd-symbols-residue-count-inputs-segment-atypical-setting}), with its notation, and let
$\mu\in\mathcal P$ be a path that is not typical. Then the number of standard fillings of $\mu^{\mathrm{tn}}$
(Definition~\ref{pd-symbols-path-generators}(4) and (5)) is at most
\[
 \exp\Big(9N^2\ell_0-\min\Big(\frac{L_a(K-2N)}{16k},\frac{K-2N}{12}\Big)\Big)\operatorname{rk}(B).
\]
\end{lemma}
\leanref{Formalism/SharpRestriction/GrowthBounds/SegmentDegree/AtypicalPathSmallWeight.lean}{Formalism.MainNotions.SelectionLemmas.Pf.T.AtypicalPathSmallWeight}
\leanref{Formalism/SharpRestriction/GrowthBounds/SegmentDegree/TypOfTypical.lean}{Formalism.MainNotions.SelectionLemmas.Pf.T.atypical_path_small_weight}
\leanref{Formalism/ThetaLogCoordinates/ResidueClassEstimates/StatementT.lean}{Formalism.InvariantFactorPaths.BranchingDimensions.PathWeightBound}
\leanref{Formalism/ThetaLogCoordinates/ResidueClassEstimates/FinalT.lean}{Formalism.InvariantFactorPaths.BranchingDimensions.Top.lemma_t_one}
\leanref{Formalism/ThetaLogCoordinates/ResidueClassEstimates/VandermondeRankEstimates/AtypicalWeightBoundOf.lean}{Formalism.InvariantFactorPaths.BranchingDimensions.pathWeightBound_of}
\leanref{Formalism/ThetaLogCoordinates/ResidueClassEstimates/VandermondeRankEstimates/N0LeNs.lean}{Formalism.InvariantFactorPaths.BranchingDimensions.pathEntryBounds}
\leanref{Formalism/ThetaLogCoordinates/ResidueClassEstimates/VandermondeRankEstimates/N0LeNs.lean}{Formalism.InvariantFactorPaths.BranchingDimensions.pathInterlacingCriterion}
\leanref{Formalism/SharpRestriction/GrowthBounds/SegmentDegree/ProdIccOneEqProdRange.lean}{Formalism.MainNotions.SelectionLemmas.Pf.T.toL_wt}
\leanref{Formalism/SharpRestriction/GrowthBounds/SegmentDegree/TypOfTypical.lean}{Formalism.MainNotions.SelectionLemmas.Pf.T.typ_of_Typical}

\begin{proof}
Put $k'=K-2N$; by (d), $k\leq2k'$, so $k'\geq1$. We use the notation of
Lemmas~\ref{pd-symbols-residue-count-inputs-segment-atypical-tails}
and~\ref{pd-symbols-residue-count-inputs-segment-atypical-log}, and write $f=f_{k'}$. Let $\mu$ be a path for $B_\mu$ and
$\lambda\in\mathbf Z^K$ the last vector of $\mu^{\mathrm{tn}}$; as $h_N=r$, $\lambda=\hat\mu_N$. By
Lemma~\ref{pd-symbols-residue-count-inputs-segment-clauses}(2), the number of standard fillings of $\mu^{\mathrm{tn}}$ is
$W_K(\lambda)$.

\emph{Step 1: the entries of $\lambda$.} The vector $\lambda$ is dominant. By
Lemma~\ref{pd-symbols-residue-count-inputs-segment-atypical-rows}(1), $\mu_N(m)\geq0$ for $m\leq k-N$ and
$\mu_N(m)\leq0$ for $m>N$; so $\lambda(c)=\mu_N(r+c)$ is positive only for $r+c\leq N$ and negative only for $c>K-N$, and
$\lambda$ has at most $N$ entries of each sign. Next, the entries of $\lambda$ lie between $\mu_{n_0}(k)$ and
$\mu_{n_0}(1)$: by the proof of Lemma~\ref{pd-symbols-residue-count-inputs-segment-atypical-tails}, for $n_0\leq s<N$ the
entries of $\hat\mu_{s+1}$ satisfy $\hat\mu_s(c+1)\leq\hat\mu_{s+1}(c)\leq\hat\mu_s(c)$, so by induction on $s$ the entries
of $\hat\mu_s$ lie between the last and the first entry of $\hat\mu_{n_0}=\mu_{n_0}$. By
Lemma~\ref{pd-symbols-residue-count-inputs-segment-atypical-rows}(2), $-\beta^-\leq\mu_{n_0}(k)$ and
$\mu_{n_0}(1)\leq\beta^+$, where $\beta^+$ and $\beta^-$ are the sums of the $B(l)$ over the vector and over the covector
letters $l<n_0$. Hence $\max(\lambda(1),0)+\max(-\lambda(K),0)\leq\beta^++\beta^-=\beta\leq2an_0\leq2aN$, by (a).

\emph{Step 2: the Weyl product and the rank.} By
Lemma~\ref{pd-symbols-residue-count-inputs-segment-atypical-log}(3), with $n=K$ and $A_0=2aN$, $W_K(\lambda)>0$ and
\[
 \log W_K(\lambda)\leq\sum_{c=1}^Kf\big(\max(\lambda(c),0)\big)+\sum_{c=1}^Kf\big(\max(-\lambda(c),0)\big)+9N^2\ell_0,
\]
as $2Nr+8N^2\leq9N^2$ by $2r\leq n_0\leq N$. By part (2) of that lemma, with $A_0=2a$ and the $B(i)$, $i<n_0$, as the
$B_i$, $\sum_{i<n_0}f(B(i))\leq\log\operatorname{rk}(B)$.

\emph{Step 3: comparison with the latest letters.} By (a) and (b), $L^+_N(t)\geq k\geq k'$ and
$4L^+_N(t+1)\leq p^2L^+_N(t+1)/5\leq L^+_N(t)$, as $p^2\geq20$. By
Lemma~\ref{pd-symbols-residue-count-inputs-segment-atypical-tails}(1) with $s=N$, $\lambda^+_{>u}\leq L^+_{N,>u}$ for all
$u$; for $u=q^+_N$ this gives $\lambda(c)\leq0$ for $c>q^+_N$. Let $\alpha_t=\max(\lambda(t),0)$ for $t\leq\min(K,q^+_N)$
and $\alpha_t=0$ for $K<t\leq q^+_N$. Then $\sum_{t>u}\alpha_t=\lambda^+_{>u}$ for $u\geq0$, and, as $f(0)=0$, the sum of
the $f(\max(\lambda(c),0))$ is $\sum_{t\leq q^+_N}f(\alpha_t)$. Put
$\Delta^+=\sum_{c=1}^Kf(\max(\lambda(c),0))-\sum_{t=1}^{q^+_N}f(L^+_N(t))$. By
Lemma~\ref{pd-symbols-residue-count-inputs-segment-atypical-log}(1), $\Delta^+\leq0$, and:
\begin{enumerate}
\item[($*$)] if $u\geq0$ and $X>0$ satisfy $\lambda^+_{>u}\leq L^+_{N,>u}-X$, then $u<q^+_N$, as otherwise the right side
  is $-X<0\leq\lambda^+_{>u}$, and $\Delta^+\leq-Xk'/(4L^+_N(u+1))$.
\end{enumerate}
The same holds for $\Delta^-$, formed with the $\max(-\lambda(c),0)$, $\lambda^-_{>u}$ and the $L^-_N(t)$. Every letter
$i<n_0$ is a vector or a covector letter, so $\sum_tf(L^+_N(t))+\sum_tf(L^-_N(t))=\sum_{i<n_0}f(B(i))$. With Step 2,
\[
 \log W_K(\lambda)\leq9N^2\ell_0+\log\operatorname{rk}(B)+\Delta^++\Delta^-,\qquad\Delta^+,\Delta^-\leq0.
\]

\emph{Step 4: a deficit.} We show $\Delta^++\Delta^-\leq-\min(L_ak'/(16k),k'/12)$; with Step 3 this gives the lemma
after taking exponentials. As $\mu$ is not typical, one of the three conditions fails at some letter $i$.

(a) Let $i$ be a covector letter with $\mu_i^+-\mu_{i+1}^+>\Xi(i)$. By (c), $r\leq i<n_0$, so $\hat\mu_i=\mu_i$ and
$\hat\mu_{i+1}=\mu_{i+1}$; by (d) and (f), $\Xi(i)>0$. As $i$ is a covector letter, $L^+_{i+1,>0}=L^+_{i,>0}$. By
Lemma~\ref{pd-symbols-residue-count-inputs-segment-atypical-tails}(1),
$\mu_{i+1}^+<\mu_i^+-\Xi(i)\leq L^+_{i+1,>0}-\Xi(i)$, and part (2) of that lemma with $s_0=i+1$ and $u=0$ gives
$\lambda^+_{>c}\leq L^+_{N,>c}-\Xi(i)$, with $c$ the number of vector letters $l$ with $i<l<n_0$. By ($*$), $c<q^+_N$, and
$L^+_N(c+1)=B(l)$ for the largest vector letter $l<i+1$, which is smaller than $i$. By (f), $pB(l)\leq4\gamma(i)$, and
$B(l)\geq k>0$ by (a); so
\[
 \Delta^+\leq-\frac{\Xi(i)k'}{4B(l)}\leq-\frac{L_a\gamma(i)}{pk}\cdot\frac{pk'}{16\gamma(i)}=-\frac{L_ak'}{16k}.
\]

(b) Let $i\geq r$ be a vector letter with $\mu_i^--\mu_{i+1}^->\Xi(i)$. For $s>n_0$ the entries of $\mu_s$ left out of
$\hat\mu_s$ equal $\beta\geq0$, by the definition of $\mathcal P$; so $(\hat\mu_s)^-=\mu_s^-$ for every $s$. As $i$ is a
vector letter, $L^-_{i+1,>0}=L^-_{i,>0}$. As in (a), Lemma~\ref{pd-symbols-residue-count-inputs-segment-atypical-tails}
gives $\lambda^-_{>c}\leq L^-_{N,>c}-\Xi(i)$ with $c$ the number of covector letters $l$ with $i<l<n_0$; by ($*$),
$c<q^-_N$ and $L^-_N(c+1)=B(l)$ for the largest covector letter $l<\min(i+1,n_0)$, which is smaller than $\min(i,n_0)$
as $i$ is a vector letter. As in (a), $\Delta^-\leq-L_ak'/(16k)$.

(c) Let $i\geq r$ with $\mu_i(r)<\frac23B(o(i))$. As $i<N$, $h_i\leq r-1$; put $u=r-1-h_i$, so that
$\hat\mu_i(u+1)=\mu_i(r)$, which is nonnegative by
Lemma~\ref{pd-symbols-residue-count-inputs-segment-atypical-rows}(1). The $h_i$ letters $l$ with $n_0\leq l<i$ are vector
letters by (c), and $o(i)<n_0$ is the $r$-th largest vector letter below $i$ by (e); so $o(i)=l^+_i(r-h_i)=l^+_i(u+1)$
and $L^+_{i,>u}=B(o(i))+L^+_{i,>u+1}$. By Lemma~\ref{pd-symbols-residue-count-inputs-segment-atypical-tails}(1),
\[
 (\hat\mu_i)^+_{>u}=\mu_i(r)+(\hat\mu_i)^+_{>u+1}\leq\tfrac23B(o(i))+L^+_{i,>u+1}=L^+_{i,>u}-\tfrac13B(o(i)),
\]
and part (2) of that lemma with $s_0=i$ gives $\lambda^+_{>u+c}\leq L^+_{N,>u+c}-\frac13B(o(i))$, with $c$ the number of
vector letters $l$ with $i\leq l<n_0$. The list $l^+_N$ begins with these $c$ letters and continues with $l^+_i$, so
$L^+_N(u+c+1)=B(o(i))$, which is positive by (a). By ($*$),
$\Delta^+\leq-\frac13B(o(i))k'/(4B(o(i)))=-k'/12$.
% Lean: atypical_path_small_weight D admB3 translates the setting into BranchingDimensions.Letters by toL (1-based
%   letters, Ns = N), with p = 13 and the b of TData equal to 13 gamma / p, so that XT 13 La b = Xi (toL_Wf, toL_Hyp,
%   toL_TData, toL_IsPATH, typ_of_Typical, rank_eq_prod_choose), and wt mu = the number of standard fillings by
%   toL_wt (admB3, Lemma pd-symbols-residue-count-inputs-segment-clauses(2)). It applies the second clause of
%   Top.lemma_t_one = pathWeightBound_of pathEntryBounds pathInterlacingCriterion logBinomialProperties
%   logWeylDimensionEstimate logRankEstimate pathTailBounds logBinomialDeficit. Step 1 is ent_sign (pathEntryBounds)
%   and hlab (pathInterlacingCriterion: pre_size, tails_range, interl_of_path) with be_le; Step 2 is hLW and hLR;
%   Step 3 is TTyp.bracket (applied to kaRows and to its reverse revW) with hLk, hL4v, hL4c, hTv, hTc and sum_lat;
%   Step 4 is the three cases hc1, hc2, hc3 with TTyp.key (lat_shift), heqv, heqc, TTyp.a_fin and, in (c), hh
%   (Bo i, oldest_threshold_eq), hsplitT and hds.
\end{proof}

\begin{lemma}\label{pd-symbols-residue-count-inputs-segment-atypical-dictionary}
Let $C_{La}\geq1088$ and $C_k\geq1088$ be real numbers, and let $p$, $\mathcal R$, $\kappa$, $n'$, $r$, $K$, $k$, $T$,
$L_1$, $N_1$, $(J,f)$, $\deg$, $E'$, $N_s$, $E_0$, $m_*$, $a$, $Y$ and $q$ be as in Lemma~\ref{pd-symbols-residue-rank};
we use the notation of Definitions~\ref{pd-symbols-residue-coarse}, \ref{pd-symbols-residue-families},
\ref{pd-symbols-residue-count-inputs-index} and~\ref{pd-symbols-residue-count-inputs-segment-typical} and of
Lemma~\ref{pd-symbols-residue-count-inputs-segment-paths}. Let $j\in\mathbf Z$, let $b\colon\mathbf Z\to\mathbf
Z_{\geq0}$ be balanced at $j$, put $N=|\Lambda_j|$, let $v_0<\dots<v_{N-1}$ be the elements of $\Lambda_j$, lettered as in
Definition~\ref{pd-symbols-residue-pairs}, and put $n_0=N-r$, $B(i)=\beta_j(b,v_i)$ for $i<n_0$, and
$\gamma(i)=a/p^{\operatorname{dist}_j(v_i)+R_j}$ and $o(i)$ as in
Definition~\ref{pd-symbols-residue-count-inputs-segment-typical} for $r\leq i<N$. Then $N\geq r$, and:
\begin{enumerate}
\item $r$, $K$, $n_0$, this lettering, $B$, $o$, $\gamma$, $a$, $L_a$ and $p$ form an atypical setting
  (Definition~\ref{pd-symbols-residue-count-inputs-segment-atypical-setting}), whose $\Xi(i)$ are those of
  Definition~\ref{pd-symbols-residue-count-inputs-segment-typical};
\item the paths of this setting are the paths of the pairs in $T^{\mathrm{all}}_j(b)$, its typical paths are the paths
  that are typical for $(j,b)$, and its good paths are counting paths for $(j,b)$;
\item $\operatorname{rk}(B)=|T^{\mathrm{all}}_j(b)|$ and $\exp(11N^2\ell_0-\min(L_a,k)/32)\leq a^{-N_s^2}$.
\end{enumerate}
\end{lemma}
\leanref{Formalism/SharpRestriction/GrowthBounds/SegmentDegree/GoodPathWeightOfParts.lean}{Formalism.MainNotions.SelectionLemmas.Pf.LetterDictionaryBounds}
\leanref{Formalism/SharpRestriction/GrowthBounds/SegmentDegree/DistAddTwoLeOfSameVec.lean}{Formalism.MainNotions.SelectionLemmas.Pf.TE.letter_dictionary_bounds}
\leanref{Formalism/SharpRestriction/GrowthBounds/SegmentDegree/DistAddTwoLeOfSameVec.lean}{Formalism.MainNotions.SelectionLemmas.Pf.TE.final_num}
\leanref{Formalism/SharpRestriction/GrowthBounds/SegmentDegree/FourHundredSixteenLeA.lean}{Formalism.MainNotions.SelectionLemmas.Pf.TE.Bdeg_bounds}
\leanref{Formalism/SharpRestriction/GrowthBounds/SegmentDegree/FourHundredSixteenLeA.lean}{Formalism.MainNotions.SelectionLemmas.Pf.TE.c_bounds}
\leanref{Formalism/SharpRestriction/GrowthBounds/SegmentDegree/FourHundredSixteenLeA.lean}{Formalism.MainNotions.SelectionLemmas.Pf.TE.a_big}
\leanref{Formalism/SharpRestriction/GrowthBounds/SegmentDegree/DistAddTwoLeOfSameVec.lean}{Formalism.MainNotions.SelectionLemmas.Pf.TE.dist_add_two_le_of_same_vec}
\leanref{Formalism/SharpRestriction/GrowthBounds/SegmentDegree/SortGetDSpec.lean}{Formalism.MainNotions.SelectionLemmas.Pf.TE.oldest_spec}

\begin{proof}
We use the notation of Definition~\ref{pd-symbols-run-system}, and $P_j(x)=p^{\operatorname{dist}_j(x)+R_j}$ as in
Definition~\ref{pd-symbols-residue-families}. By Step~1 of the proof of Lemma~\ref{pd-symbols-residue-lifts-numeric},
$Y\geq1$, $k\geq5$ and $q\geq1$, and by the proof of Lemma~\ref{pd-symbols-residue-top-count}, $\log a\geq1$. For
$x\in\Lambda_j$, $\operatorname{dist}_j(x)=f(j,r-1)-x$, so $P_j(x)\geq p^{x'-x}P_j(x')$ for $x\leq x'$ in $\Lambda_j$.

\emph{Step 1: the letters.} The elements of $\Lambda_j$ that are not new lie between $f(j-1,0)$ and $f(j-1,r-1)$ and have
their residues in $T$; by condition~(3) of validity they are the $f(j-1,i)$, $i<r$, so $v_i=f(j-1,i)$ for $i<r$. By the
first paragraph of the proof of Lemma~\ref{pd-symbols-residue-count-inputs-segment-plane}, $N\geq r$ and
$v_{n_0+t}=f(j,t)$ for $t<r$. By Lemma~\ref{pd-symbols-residue-count-sum-facts}(2), the letters $i$ with $r\leq i<n_0$
are the elements of $\Lambda'_j$, and the letters $i\geq r$ are those with $v_i$ new. As $\Lambda'_j$ has at least $r$
elements, $n_0\geq2r$ and $N\geq3r$; its $r$ largest elements $v_{n_0-r},\dots,v_{n_0-1}$ are extreme, hence even by
condition~(4) of validity, and the $f(j-1,i)$ and $f(j,t)$ are even. This is (c) of
Definition~\ref{pd-symbols-residue-count-inputs-segment-atypical-setting} and $2r\leq n_0$.

\emph{Step 2: the numbers.} By hypothesis (2) of Theorem~\ref{pd-symbols-count}, $N\leq N_s$ and $N+r\leq N_s+r\leq K$.
Also $N_s\geq6$. Indeed, by conditions (2) and (3) of validity the $f(0,i)$ and the $f(1,i)$, $i<r$, are $2r$ distinct
elements of $\Lambda_1$, and the elements of $\Lambda'_1\subset\Lambda_1$ lie strictly between $f(0,r-1)$ and $f(1,0)$;
by condition (4) there are at least $r$ of them. So $N_s\geq|\Lambda_1|\geq3r\geq6$, as $r\geq2$. By hypothesis (4) of
Theorem~\ref{pd-symbols-count}, as $p\geq13$ and $Y,k,q\geq1$,
$416N_s\leq32pYk^2N_sq^2(E_0+m_*+1)\leq a$; so $a\geq3$ and $N\leq a$. By (4) and $\log a\geq1$,
$1088N_s^2\leq C_kN_s^2\log a\leq k$, so $6N_s\leq k$ and $2(K-2N)\geq2(k-3N_s)\geq k$. Also $L_a=C_{La}N_s^2\log a>0$.
With $p\geq13$ and Step 1 this is (d).

\emph{Step 3: the sizes.} For $y\in\Lambda_j$: $\frac{36}5Yk^2N_sP_j(y)\leq d^D(\bar y)$ by
Lemma~\ref{pd-symbols-residue-count-sum-facts}(4), and $\frac9{10}a\leq d^D(\bar y)\leq2a$ by
Lemma~\ref{pd-symbols-residue-top-count} and hypothesis (1) of Theorem~\ref{pd-symbols-count}, as $\bar y\in T$ and
$d^D\leq d\leq\deg$. For $i<n_0$ we show
\[
 d^D(\bar v_i)/P_j(v_i)-2kN_s\leq B(i)\leq d^D(\bar v_i)/P_j(v_i).
\]
If $v_i$ is not new, $B(i)=\max(W_j(v_i)-m_0,0)$, and $W_j(v_i)-m_0=\lfloor d^D(\bar v_i)/P_j(v_i)\rfloor-(k+r+1)-(kN_s+1)$
lies between $d^D(\bar v_i)/P_j(v_i)-(k+r+3+kN_s)$ and $d^D(\bar v_i)/P_j(v_i)$; as $k+r+3\leq3k\leq kN_s$, this gives both
bounds. If $v_i$ is new, $v_i\in\Lambda'_j$ and $B(i)=b(v_i)$, and the bounds follow from the balance of $b$ at $j$, as
$k\leq2kN_s$. Hence, for $i<n_0$, $B(i)\leq2a/P_j(v_i)$, and $B(i)\geq\frac{36}5k^2N_s-2kN_s\geq k$, which is (a).

Let $i<i'<n_0$ be letters of the same type and $P'=P_j(v_{i'})$. Then $v_i$ and $v_{i'}$ have the same parity, so
$v_{i'}-v_i\geq2$ and $P_j(v_i)\geq p^2P'$, and $p^2B(i)\leq2a/P'$. As
$2kN_s\leq\frac1{18}\cdot\frac{36}5Yk^2N_s\leq d^D(\bar v_{i'})/(18P')$,
$B(i')\geq\frac{17}{18}d^D(\bar v_{i'})/P'\geq\frac{17}{18}\cdot\frac9{10}\cdot a/P'\geq\frac25a/P'$. So
$p^2B(i)\leq5B(i')$, which is (b).

\emph{Step 4: the oldest letters.} Let $r\leq i<N$. Then $v_i$ is new, so the $f(j-1,i')$, $i'<r$, are $r$ even
elements of $\Lambda_j$ smaller than $v_i$, and $\operatorname{old}_j(v_i)$ is the $r$-th largest even element of
$\Lambda_j$ smaller than $v_i$. So $o(i)<i$ is a vector letter, and the even elements $z$ of $\Lambda_j$ with
$\operatorname{old}_j(v_i)\leq z<v_i$, which are the $v_l$ for the vector letters $l$ with $o(i)\leq l<i$, are $r$ in
number. This is (e). For (f), $\gamma(i)=a/P_j(v_i)>0$; and for a letter $l<\min(i,n_0)$, $v_l<v_i$, so
$P_j(v_l)\geq pP_j(v_i)$ and $pB(l)\leq2a/P_j(v_i)\leq4\gamma(i)$. Finally
$\Xi(i)=L_a\gamma(i)/(pk)=L_aa/(pk\,p^{\operatorname{dist}_j(v_i)+R_j})$. This proves (1).

(2) For a letter $i$, $v_i<f(j,0)$ if and only if $i<n_0$, and $v_i=f(j,i')$ with $i'<r$ if and only if $i=n_0+i'$. As
$\beta_j(b)=\sum_{i<n_0}B(i)=\beta$, the conditions on a path of a pair in $T^{\mathrm{all}}_j(b)$ are those on a path of
the setting. The conditions of Definition~\ref{pd-symbols-residue-count-inputs-segment-typical} are those of a typical
path, as $B(o(i))=\beta_j(b,v_{o(i)})$ and the $\Xi(i)$ agree. Let $\mu$ be a good path, and $i$ a letter with $v_i$ new,
that is, $i\geq r$; put $o=\operatorname{old}_j(v_i)=v_{o(i)}$. If $o$ is not new, then
$\theta'_{j,b}(v_i)=\frac23(W_j(o)-m_0)\leq\frac23B(o(i))$; if $o$ is new, then
$\theta'_{j,b}(v_i)=\frac23b(o)=\frac23B(o(i))$. So $\theta'_{j,b}(v_i)\leq\mu_i(r)$. If $v_i$ is even, $i$ is a vector
letter, and by Step 3, $\Xi(i)+2B(o(i))\leq\Xi(i)+4a/P_j(o)=o'_j(v_i)$; if $v_i$ is odd, $\Xi(i)=o'_j(v_i)$. So $\mu$
satisfies the conditions of a counting path at $v_i$.

(3) By Step~4 of the proof of Lemma~\ref{pd-symbols-residue-count-inputs-segment-paths},
$|T^{\mathrm{all}}_j(b)|=\prod_{i<n_0}\binom{B(i)+k-1}{k-1}=\operatorname{rk}(B)$. By Step 2, $2aN\leq2a\cdot a/416$, so
$1+2aN\leq a^2$ and $\ell_0\leq2\log a$; hence $11N^2\ell_0\leq22N_s^2\log a$. As $C_{La},C_k\geq1088$,
$L_a\geq1088N_s^2\log a$ and $k\geq1088N_s^2\log a$, so $\min(L_a,k)/32\geq34N_s^2\log a$. The exponent is therefore at
most $-12N_s^2\log a\leq-N_s^2\log a$.
% Lean: letter_dictionary_bounds (TopCountLowerBounds 1088, Lemma pd-symbols-residue-top-count): n0 = segN - r and od
%   from exists_od (oldest_spec, sort_getD_spec); Step 1 is r_le_segN, two_r_le, dict_w, dict_early, dict_lateV and
%   comb_W_even; Step 2 is six_le_Ns (fr 0 i, LLset 1 and fr 1 i inside letters 1, with hX.card 1), a_big,
%   four_hundred_sixteen_le_a, log_a_ge, k_big and the fields N_add_r_le_K, H4 and N_le_a; Step 3 is Bdeg_bounds and
%   c_bounds (the fields H1 and H2, with dist_add_two_le_of_same_vec); Step 4 is the fields od_lt, od_vec, od_card,
%   bb_pos and bb_le (dist_lt). (2) is the clauses IsPATH <-> IsPATHa (dict_tau, beS_eq) and GoodA -> IsGood (the
%   definitions thc and oc of Coarse.lean); the book's typical paths for (j, b) are the Lean's Typ read through this
%   dictionary. (3) is the rank identity (hpt, prod_fin_lt, LLset_eq_image) and final_num.
\end{proof}

\begin{lemma}\label{pd-symbols-residue-count-inputs-segment-atypical}
Let $C_{La}\geq1088$ and $C_k\geq1088$ be real numbers, and let $p$, $\mathcal R$, $\kappa$, $n'$, $r$, $K$, $k$, $T$,
$L_1$, $N_1$, $(J,f)$, $\deg$, $E'$, $N_s$, $E_0$, $m_*$, $a$, $Y$ and $q$ be as in
Lemma~\ref{pd-symbols-residue-rank}; we use the notation of Definitions~\ref{pd-symbols-residue-coarse},
\ref{pd-symbols-residue-count-inputs-index} and~\ref{pd-symbols-residue-count-inputs-segment-typical} and of
Lemma~\ref{pd-symbols-residue-count-inputs-segment-paths}. Let $j\in\mathbf Z$ and let $b\colon\mathbf Z\to\mathbf
Z_{\geq0}$ be balanced at $j$. Then:
\begin{enumerate}
\item every pair $(\mu,I'')\in T^{\mathrm{all}}_j(b)$ with $\mu$ typical for $(j,b)$ lies in $T_j(b)$;
\item the number of the pairs $(\mu,I'')\in T^{\mathrm{all}}_j(b)$ with $\mu$ not typical for $(j,b)$ is at most
  $a^{-N_s^2}|T^{\mathrm{all}}_j(b)|$.
\end{enumerate}
\end{lemma}
\leanref{Formalism/SharpRestriction/GrowthBounds/SegmentDegree/AtypicalPathSmallWeight.lean}{Formalism.MainNotions.SelectionLemmas.Pf.T.HypT}
\leanref{Formalism/SharpRestriction/GrowthBounds/SegmentDegree/AtypicalPathSmallWeight.lean}{Formalism.MainNotions.SelectionLemmas.Pf.T.GoodA}
\leanref{Formalism/SharpRestriction/GrowthBounds/SegmentDegree/AtypicalPathSmallWeight.lean}{Formalism.MainNotions.SelectionLemmas.Pf.T.TypicalPathOrdersSmall}
\leanref{Formalism/SharpRestriction/GrowthBounds/SegmentDegree/AtypicalPathSmallWeight.lean}{Formalism.MainNotions.SelectionLemmas.Pf.T.PathCountFinite}
\leanref{Formalism/SharpRestriction/GrowthBounds/SegmentDegree/AtypicalPathSmallWeight.lean}{Formalism.MainNotions.SelectionLemmas.Pf.T.AtypicalPathSmallWeight}
\leanref{Formalism/SharpRestriction/GrowthBounds/SegmentDegree/AtypicalPathSmallWeight.lean}{Formalism.MainNotions.SelectionLemmas.Pf.T.lemT_of}
\leanref{Formalism/SharpRestriction/GrowthBounds/SegmentDegree/GoodPathWeightOfParts.lean}{Formalism.MainNotions.SelectionLemmas.Pf.LetterDictionaryBounds}

\begin{proof}
Form the atypical setting of Lemma~\ref{pd-symbols-residue-count-inputs-segment-atypical-dictionary}(1), and use the
notation of Definition~\ref{pd-symbols-residue-count-inputs-segment-atypical-setting}.

(1) Let $(\mu,I'')\in T^{\mathrm{all}}_j(b)$ with $\mu$ typical for $(j,b)$. By
Lemma~\ref{pd-symbols-residue-count-inputs-segment-atypical-dictionary}(2), $\mu$ is a typical path of the setting. By
Lemma~\ref{pd-symbols-residue-count-inputs-segment-atypical-orders} it is good, and by
Lemma~\ref{pd-symbols-residue-count-inputs-segment-atypical-dictionary}(2) again it is a counting path for $(j,b)$. So
$(\mu,I'')\in T_j(b)$.

(2) Let $\mathcal S$ be the set of the pairs $(\mu,I'')\in T^{\mathrm{all}}_j(b)$ with $\mu$ not typical for $(j,b)$. By
Lemma~\ref{pd-symbols-residue-count-inputs-segment-atypical-dictionary}(2), their paths are the paths in $\mathcal P$
that are not typical, and for each such $\mu$ the pairs $(\mu,I'')$ are as many as the standard fillings of
$\mu^{\mathrm{tn}}$. Put $m=\min\big(L_a(K-2N)/(16k),(K-2N)/12\big)$. By
Lemmas~\ref{pd-symbols-residue-count-inputs-segment-atypical-weight}
and~\ref{pd-symbols-residue-count-inputs-segment-atypical-count},
\[
 |\mathcal S|\leq\exp(2N^2\ell_0)\exp(9N^2\ell_0-m)\operatorname{rk}(B).
\]
By condition (d), $K-2N\geq k/2$, so $L_a(K-2N)/(16k)\geq L_a/32$ and $(K-2N)/12\geq k/24$; hence $m\geq\min(L_a,k)/32$.
By Lemma~\ref{pd-symbols-residue-count-inputs-segment-atypical-dictionary}(3),
\[
 |\mathcal S|\leq\exp\big(11N^2\ell_0-\min(L_a,k)/32\big)\operatorname{rk}(B)\leq a^{-N_s^2}|T^{\mathrm{all}}_j(b)|.
\]
% Lean: (1) is typical_path_orders_small (Typ -> GoodA, TypicalPathOrdersSmall.lean, 225 lines) with the part
%   GoodA -> IsGood of TE.letter_dictionary_bounds (TopCountLowerBounds, Lemma pd-symbols-residue-top-count; HypT from
%   exists_od, Bdeg_bounds, c_bounds and dist_add_two_le_of_same_vec; DistAddTwoLeOfSameVec.lean, SortGetDSpec.lean and
%   FourHundredSixteenLeA.lean, about 820 lines). (2) is the bound h3 of lemT_of: SB, the pairs with an atypical path,
%   has at most |SP| exp(9 N^2 ell0 - min(La (K - 2N) / (16k), (K - 2N) / 12)) rkA elements by
%   atypical_path_small_weight D admB3 (TypOfTypical.lean, CountFromEqFilterLength.lean, ProdIccOneEqProdRange.lean, about
%   650 lines, applying BranchingDimensions.Top.lemma_t_one = pathWeightBound_of in
%   ThetaLogCoordinates/ResidueClassEstimates/, about 3700 lines with VandermondeRankEstimates/), |SP| <= exp(2 N^2 ell0)
%   by path_count_finite (WindowSmall.lean, 180 lines), hmin (H.H4), and exp(11 N^2 ell0 - min(La, k) / 32) <= a^-(Ns^2)
%   by TE.final_num (the last clause of letter_dictionary_bounds). Here ell0 = log(1 + 2aN), N = |Lambda_j|, and
%   rkA = |T^all_j(b)| (Lemma pd-symbols-residue-count-inputs-segment-paths). In all, about 5600 lines.
\end{proof}

\begin{lemma}\label{pd-symbols-residue-count-inputs-segment}
Let $C_{La}\geq1088$ and $C_k\geq1088$ be real numbers, and let $p$, $\mathcal R$, $\kappa$, $n'$, $r$, $K$, $k$, $T$,
$L_1$, $N_1$, $(J,f)$, $\deg$, $E'$, $N_s$, $E_0$, $m_*$, $a$, $Y$ and $q$ be as in
Lemma~\ref{pd-symbols-residue-rank}; we use the notation of Definitions~\ref{pd-symbols-residue-coarse},
\ref{pd-symbols-residue-families} and~\ref{pd-symbols-residue-count-inputs-index}, and write
$B(n)=\binom{n+k-1}{k-1}$. Let $j\in\mathbf Z$ and let $b\colon\mathbf Z\to\mathbf Z_{\geq0}$ be balanced at $j$. Then:
\begin{enumerate}
\item if $b$ vanishes off $\Lambda_j$ and $W_j(x)\leq b(x)$ for every $x\in\Lambda_j$ that is not new, the elements
  $y_j(b;q)$, $q\in\mathcal I_j(b)$, lie in $M_j(b)$ and are linearly independent over $\mathcal R$;
\item $T_j(b)$ is finite, and
  \[
   \big(1-a^{-N_s^2}\big)\prod_{x\in\Lambda'_j}B(b(x))\prod_{i<r}B\big(\max(0,W_j(f(j-1,i))-m_0)\big)\leq|T_j(b)|;
  \]
\item if $T_j(b)$ is finite, then for every $\rho\in\mathbf Z_{\geq0}^r$,
  \[
   |T_j(b)|\Big(Q_j^{\,r-1}\max\Big(0,1-\frac r{2Y}\Big)\Big)^r\leq\sum_\varepsilon\big|\mathcal I_j(\tau_j(b,\varepsilon))\big|,
  \]
  the sum over the $\varepsilon\in\mathbf Z_{\geq0}^r$ with $\varepsilon(i)<rQ_j$ and $\varepsilon(i)\equiv\rho(i)\pmod{Q_j}$
  for all $i$.
\end{enumerate}
\end{lemma}
\leanref{Formalism/SharpRestriction/GrowthBounds/SegmentDegree/GoodPathWeight.lean}{Formalism.MainNotions.SelectionLemmas.Pf.RestMem}
\leanref{Formalism/SharpRestriction/GrowthBounds/SegmentDegree/GoodPathWeight.lean}{Formalism.MainNotions.SelectionLemmas.Pf.RestInd}
\leanref{Formalism/SharpRestriction/GrowthBounds/SegmentDegree/GoodPathWeight.lean}{Formalism.MainNotions.SelectionLemmas.Pf.GoodPathWeight}
\leanref{Formalism/SharpRestriction/GrowthBounds/SegmentDegree/GoodPathWeight.lean}{Formalism.MainNotions.SelectionLemmas.Pf.RestMon}
\leanref{Formalism/SharpRestriction/GrowthBounds/SegmentDegree/M0LeBwOfIsW.lean}{Formalism.MainNotions.SelectionLemmas.Pf.restMem_of}
\leanref{Formalism/SharpRestriction/GrowthBounds/SegmentDegree.lean}{Formalism.MainNotions.SelectionLemmas.Pf.restInd}
\leanref{Formalism/SharpRestriction/GrowthBounds/SegmentDegree/GoodPathWeightAlt.lean}{Formalism.MainNotions.SelectionLemmas.Pf.good_path_weight}
\leanref{Formalism/SharpRestriction/GrowthBounds/SegmentDegree.lean}{Formalism.MainNotions.SelectionLemmas.Pf.restMon}

\begin{proof}
We use the notation of Definition~\ref{pd-symbols-residue-count-inputs-index}.

(3) Let $T_j(b)$ be finite and $\rho\in\mathbf Z_{\geq0}^r$, and put $\beta=\beta_j(b)$ and $Q=Q_j$. If $1-\frac r{2Y}\leq0$,
the left side is $0$. Let $1-\frac r{2Y}>0$. By Lemma~\ref{pd-symbols-residue-count-inputs-ratio}, $4Y\beta\leq Q$; as
$Y\geq1$, $4\beta\leq Q$, so $t=Q-2\beta>0$. For $i'<r$ let $l(i')$ be the letter with $v_{l(i')}=f(j,i')$. For
$\tau=(\mu,I'')\in T_j(b)$, with $\mu$ a path for $B_\tau$, put $d_\tau(i')=B_\tau(l(i'))$; by
Lemma~\ref{pd-symbols-residue-count-inputs-rows}, $d_\tau(i')\leq2\beta$. Let $A$ be the set of the pairs $(\tau,u)$ of a
$\tau\in T_j(b)$ and a function $u\colon\{0,\dots,r-1\}\times\{1,\dots,r\}\to\{0,\dots,t-1\}$ with
$d_\tau(i')+\sum_cu(i',c)\equiv\rho(i')\pmod Q$ for every $i'$, and put $\varphi(\tau,u)(i')=d_\tau(i')+\sum_cu(i',c)$.

\emph{Step 1: the values of $\varphi$.} As $\varphi(\tau,u)(i')\leq2\beta+r(t-1)<r(2\beta+t)=rQ$, $\varphi(\tau,u)$ is one of
the $\varepsilon$ of the sum.

\emph{Step 2: the fibres.} Let $\varepsilon$ be one of these. The function $\tau_j(b,\varepsilon)$ agrees with $b$ below
$f(j,0)$, so $\beta_j(\tau_j(b,\varepsilon))=\beta$, and the counting paths for $(j,b)$ and for $(j,\tau_j(b,\varepsilon))$ are
the same: $\theta'_{j,b}(v)$ uses $b$ only at $\operatorname{old}_j(v)$ when it is new, and then
$\operatorname{old}_j(v)\in\Lambda'_j$ by Step~2 of the proof of Lemma~\ref{pd-symbols-residue-lifts-numeric}, so
$\operatorname{old}_j(v)<f(j,0)$. Let $(\tau,u)\in A$ with $\varphi(\tau,u)=\varepsilon$ and $\tau=(\mu,I'')$. Then
$u(i',c)+2\beta<t+2\beta=Q$ and $B_\tau(l(i'))+\sum_cu(i',c)=\varepsilon(i')=\tau_j(b,\varepsilon)(f(j,i'))$, so
$(\mu,I'',u)\in\mathcal I_j(\tau_j(b,\varepsilon))$. This map is injective. The set $\mathcal I_j(\tau_j(b,\varepsilon))$ is
finite, as $(\mu,I'',u)\mapsto((\mu,I''),u)$ maps it injectively to $T_j(b)\times\{0,\dots,Q-1\}^{r\times r}$. So the
number of $(\tau,u)\in A$ with $\varphi(\tau,u)=\varepsilon$ is at most $|\mathcal I_j(\tau_j(b,\varepsilon))|$.

\emph{Step 3: the count over $\tau$.} Fix $\tau\in T_j(b)$. The functions $u$ with $(\tau,u)\in A$ are the families of
rows $u(i',\cdot)\in\{0,\dots,t-1\}^r$ with $d_\tau(i')+\sum_cu(i',c)\equiv\rho(i')\pmod Q$. As $r\geq2$ and $2\beta<Q$,
Lemma~\ref{pd-symbols-residue-count-inputs-residues} gives at least $Q^{r-1}(1-2r\beta/Q)$ choices of each row, and
$2r\beta/Q\leq r/(2Y)$ as $4Y\beta\leq Q$. So there are at least $(Q^{r-1}(1-\frac r{2Y}))^r$ such $u$.

By Step~3, $|T_j(b)|(Q^{r-1}(1-\frac r{2Y}))^r\leq|A|$, and by Steps 1 and 2, $|A|$ is at most the sum of the statement.
% Lean: (3) is restMon = restMon_of 1088 beta_ratio_bound tau_degree_on_path_bound residue_class_count: the case
%   split on 1 - r / (2Y), hbeQ (Lemma pd-symbols-residue-count-inputs-ratio), hdtau (Lemma
%   pd-symbols-residue-count-inputs-rows), the set A with phi; Step 1 is hphiE, Step 2 hfib (restr_setTau_low,
%   beS_congr, IsPATH_congr, IsGood_congr with oldest_letter_in_segment, restr_setTau_tau, hfinI), Step 3 hcount
%   (residue_class_count, Lemma pd-symbols-residue-count-inputs-residues), and the two fibrewise counts of A.

(1) The elements $y_j(b;q)$ lie in $M_j(b)$ by Lemma~\ref{pd-symbols-residue-count-inputs-segment-members}. For their
linear independence let $N=|\Lambda_j|$, $n_0=N-r$ and $\beta=\beta_j(b)$, let $v_0<\dots<v_{N-1}$ be the elements of
$\Lambda_j$, and for $q=(\mu,I'',u)\in\mathcal I_j(b)$ put $\nu_q=\mu^{\mathrm{pl}}$ and $\kappa_q=\mu^{\mathrm{tn}}$. By
Lemma~\ref{pd-symbols-residue-count-inputs-segment-plane}, $v_{n_0+t}=f(j,t)$ for $t<r$, and $\nu_q$ satisfies the
hypotheses of Lemma~\ref{pd-symbols-residue-count-inputs-segment-cut} with this $n_0$ and $\beta$; we use the notation of
that lemma for $\nu_q$, with an index $q$. Write $y_j(b;q)=A_0z^{u}w_q$, where
\[
 A_0=D_j^{\,m_0}\prod_{i<r}z_{f(j-1,i),1}^{\,\max(b(f(j-1,i))-W_j(f(j-1,i)),0)},\qquad
 z^{u}=\prod_{i'<r,\,c\leq r}z_{f(j,i'),c}^{\,u(i',c)},
\]
and $w_q$ is $\iota_{\mathrm{pl}}(\Psi(\nu_q,\operatorname{can}_{\nu_{q,N}}))\,\iota_{\mathrm{tn}}(\Psi(\kappa_q,I''))$ with
each $z_{i,c}$ renamed $z_{v_i,c}$.

\emph{Step 4: the factor $A_0$.} The polynomial $A_0$ does not depend on $q$, and it is not $0$: $D_j$ is the
determinant of a matrix of distinct variables, as the $f(j-1,a)$, $a<r$, are distinct. As $\mathcal R[z]$ is a domain,
multiplication by $A_0$ is injective, so it suffices that the $z^{u}w_q$, $q\in\mathcal I_j(b)$, are linearly
independent.

\emph{Step 5: killing.} Let $S$ be the set of the plane variables $z_{f(j,t),c}$ with $t<r$, $c\leq r$ and $c\neq r-t$, and
$\chi_S$ as in Lemma~\ref{pd-symbols-residue-count-inputs-segment-kill}. By (1) of that lemma it suffices that the
$z^{u}\chi_S(w_q)$ are linearly independent. Under the renaming, the variables of $S$ come from the $z_{n_0+t,c}$ with
$c\leq r$ and $c\neq r-t$, and no normal variable lies in $S$. So $\chi_S(w_q)$ is the renaming of
$\chi\big(\iota_{\mathrm{pl}}(\Psi(\nu_q,\operatorname{can}_{\nu_{q,N}}))\big)\,\iota_{\mathrm{tn}}(\Psi(\kappa_q,I''))$, with
$\chi$ the map of Lemma~\ref{pd-symbols-residue-count-inputs-segment-cut}. By (1) of that lemma,
$\chi_S(w_q)=(-1)^{e_q}z^{d_q}w'_q$ for some $e_q\geq0$, where $z^{d_q}=\prod_{c=1}^\beta\prod_{a\in U_q(c)}z_{f(j,r-a),a}$
and $w'_q$ is $\iota_{\mathrm{pl}}(\Psi(\nu^{\mathrm{cut}}_q,I^{\mathrm{cut}}_q))\,\iota_{\mathrm{tn}}(\Psi(\kappa_q,I''))$ with
each $z_{i,c}$ renamed $z_{v_i,c}$; the variable $z_{n_0+r-a,a}$ is renamed $z_{f(j,r-a),a}$.

\emph{Step 6: independence.} Let $T$ be the set of the plane variables $z_{f(j,t),c}$, $t<r$, $c\leq r$, and put
$M_q=z^{u}z^{d_q}$, a monomial in the variables of $T$, and $x_q=(\nu^{\mathrm{cut}}_q,\kappa_q,I^{\mathrm{cut}}_q,I'')$.
By Lemma~\ref{pd-symbols-residue-count-inputs-segment-cut}(3), $\nu^{\mathrm{cut}}_q$ is a path of rank $r$ and
$I^{\mathrm{cut}}_q$ a standard filling of it; by Step~2 of the proof of
Lemma~\ref{pd-symbols-residue-count-inputs-segment-members}, $\kappa_q$ is a path of rank $K$, and $I''$ is a standard
filling of it. As $|\Lambda_j|+r\leq N_s+r\leq K$ by (2) of Theorem~\ref{pd-symbols-count},
Lemma~\ref{pd-symbols-residue-count-inputs-segment-clauses}(5) shows that the polynomials $w'_x$ formed in this way from
all such quadruples $x$ are linearly independent, and $w'_q=w'_{x_q}$. They involve no variable of $T$: the plane factor
involves only the rows $z_{v_i}$ with $i<n_0$, by Lemma~\ref{pd-symbols-residue-count-inputs-segment-cut}(2), and
$v_i<f(j,0)$ for these $i$; the other factor involves only normal variables. The map $q\mapsto(M_q,x_q)$ is injective.
Indeed, let $M_q=M_{q'}$ and $x_q=x_{q'}$ for $q=(\mu,I'',u)$ and $q'=(\mu',I''',u')$. Then $\kappa_q=\kappa_{q'}$ and
$I''=I'''$, and $\nu_q=\nu_{q'}$ by Lemma~\ref{pd-symbols-residue-count-inputs-segment-cut}(4), as both satisfy its
hypotheses with the same $n_0$ and $\beta$; so $\mu=\mu'$. Then $U_q(c)=U_{q'}(c)$ for every $c$, so
$z^{d_q}=z^{d_{q'}}$ and $z^{u}=z^{u'}$, and $u=u'$ as the $f(j,i')$, $i'<r$, are distinct. By
Lemma~\ref{pd-symbols-residue-count-inputs-segment-kill}(2), the $M_qw'_{x_q}$ are linearly independent, and so are the
$z^{u}\chi_S(w_q)=(-1)^{e_q}M_qw'_{x_q}$.
% Lean: the membership is restMem_of D hA (Lemma pd-symbols-residue-count-inputs-segment-members). The independence
%   is restInd D hA = restInd_of 1088 (restKey_of D hA 1088). Step 4 is restInd_of: c0 (A_0 here) with
%   StandardMonomialSpace.detX_ne_zero and fr_injective, mulLeft injective (plEmb_injective), and li_of_kill (Lemma
%   pd-symbols-residue-count-inputs-segment-kill(1)) with monTau (z^u) and wordEl (w_q). Step 5 is offSet (S), gp_preimage
%   and gt_preimage, kill_rename, kill_of_empty and kill_wordEl (kill_Psi under PlHyp from plHyp_of; dgT is d_q). Step 6
%   is restKey_of: tauSet (T), M = monTau + dgT, x (hx: cut_ppath, isSt_cut, and BracketPolynomialSpan for kappa), hw
%   (the second clause of NestedTypeBound, Lemma pd-symbols-residue-count-inputs-segment-clauses(5)), hM, hsup
%   (Psi_cut_supported, castAdd_ne_natAdd), hinj (path_eq_of_cut, joinPT_parts, monTau_apply), li_mono_words (Lemma
%   pd-symbols-residue-count-inputs-segment-kill(2)) and units_smul for the signs.

(2) By Lemma~\ref{pd-symbols-residue-count-inputs-segment-paths}, $T_j(b)$ is contained in the finite set
$T^{\mathrm{all}}_j(b)$, so it is finite. Let $T^{\mathrm{typ}}$ be the set of the pairs of $T^{\mathrm{all}}_j(b)$ whose
path is typical for $(j,b)$ (Definition~\ref{pd-symbols-residue-count-inputs-segment-typical}) and $T^{\mathrm{at}}$ the set
of the others. By Lemma~\ref{pd-symbols-residue-count-inputs-segment-atypical}, $T^{\mathrm{typ}}\subset T_j(b)$ and
$|T^{\mathrm{at}}|\leq a^{-N_s^2}|T^{\mathrm{all}}_j(b)|$. So
\[
 |T^{\mathrm{all}}_j(b)|=|T^{\mathrm{typ}}|+|T^{\mathrm{at}}|\leq|T_j(b)|+a^{-N_s^2}|T^{\mathrm{all}}_j(b)|,
\]
and $(1-a^{-N_s^2})|T^{\mathrm{all}}_j(b)|\leq|T_j(b)|$. Lemma~\ref{pd-symbols-residue-count-inputs-segment-paths} gives
$|T^{\mathrm{all}}_j(b)|$, which is the product of (2).
% Lean: (2) is good_path_weight D hA = good_path_weight_of_top_count D hA (top_count_lower_bounds 1088)
%   = good_path_weight_of_parts 1088 (path_weight_eq_rank D admB3 ColumnFamilyLowerSet GenericProjectionCriterion)
%   (lemT_of_adm D hA) (TE.letter_dictionary_bounds h41): the set S of pairs is Tset j (restr j b), hSA (S is in SAll,
%   Lemma pd-symbols-residue-count-inputs-segment-paths), hGS (SGood is in S, Lemma
%   pd-symbols-residue-count-inputs-segment-atypical(1)), hSfin, and hcard from LemT = lemT_of (rkA <= |SG| + |SB| with
%   SB the atypical pairs, and h3, Lemma pd-symbols-residue-count-inputs-segment-atypical(2)), with hrk (Lemma
%   pd-symbols-residue-count-inputs-segment-paths) and hexp. What remains open is marked in the proof of Lemma
%   pd-symbols-residue-count-inputs-segment-atypical.
\end{proof}

\begin{lemma}\label{pd-symbols-residue-count-inputs}
Let $C_{La}\geq1088$ and $C_k\geq1088$ be real numbers, and let $p$, $\mathcal R$, $\kappa$, $n'$, $r$,
$K$, $k$, $T$, $L_1$, $N_1$, $(J,f)$, $\deg$, $E'$, $N_s$, $E_0$, $m_*$, $a$, $Y$ and $q$ be as in
Lemma~\ref{pd-symbols-residue-rank}; we use the notation of Definitions~\ref{pd-symbols-residue-coarse}
and~\ref{pd-symbols-residue-families}, write $\operatorname{rk}M_j(b)$ for the rank of the
$\mathcal R$-module $M_j(b)$, and put
$\beta=p^{-n'rK}\big(1-|T|r^2k/Y\big)\prod_{t\in T}\binom{\deg(t)+k-1}{k-1}$. Then:
\begin{enumerate}
\item $\sum_{\alpha\in\mathcal L}\sum_{e\in E(\alpha)}\prod_{j=0}^{J-1}\operatorname{rk}M_j(b^{\alpha,e}_j)$ is at most the
  dimension of the $\kappa$-span of the reductions of the elements of $\Lambda^{\mathrm c}$;
\item for every $j$ there is $G_j\colon(\mathbf Z\to\mathbf Z_{\geq0})\to\mathbf R_{\geq0}$ such that $G_j(b)$ depends
  only on the values of $b$ at the new elements of $\Lambda_j$; $G_j(b)\leq\operatorname{rk}M_j(b)$ whenever $b$
  vanishes off $\Lambda_j$, $W_j(x)\leq b(x)$ for every $x\in\Lambda_j$ that is not new, and $b$ is balanced at
  $j$; and for every $b$ balanced at $j$ and every $\rho\in\mathbf Z_{\geq0}^r$,
  $S_j(b)\leq\sum_\varepsilon G_j(\tau_j(b,\varepsilon))$, the sum over the $\varepsilon\in\mathbf Z_{\geq0}^r$ with
  $\varepsilon(i)<rQ_j$ and $\varepsilon(i)\equiv\rho(i)\pmod{Q_j}$ for all $i$;
\item for $\alpha\in\mathcal L$, $e\in E(\alpha)$ and $j\in\mathbf Z$, $b^{\alpha,e}_j$ is balanced at $j$ and
  $W_j(x)\leq b^{\alpha,e}_j(x)$ for every $x\in\Lambda_j$ that is not new;
\item for all functions $G_j\colon(\mathbf Z\to\mathbf Z_{\geq0})\to\mathbf R_{\geq0}$, $j\in\mathbf Z$, with the first
  and the last property of (2), $\beta\leq\sum_{\alpha\in\mathcal L}\sum_{e\in E(\alpha)}\prod_{j=0}^{J-1}G_j(b^{\alpha,e}_j)$.
\end{enumerate}
\end{lemma}
\leanref{Formalism/SharpRestriction/GrowthBounds/SegmentDegree/RestCountOf.lean}{Formalism.MainNotions.SelectionLemmas.Pf.RestLam}
\leanref{Formalism/SharpRestriction/GrowthBounds/SegmentDegree/RestCountOf.lean}{Formalism.MainNotions.SelectionLemmas.Pf.RestSeg}
\leanref{Formalism/SharpRestriction/GrowthBounds/SegmentDegree/RestCountOf.lean}{Formalism.MainNotions.SelectionLemmas.Pf.Dat.SegOK}
\leanref{Formalism/SharpRestriction/GrowthBounds/SegmentDegree/RestCountOf.lean}{Formalism.MainNotions.SelectionLemmas.Pf.RestPad}
\leanref{Formalism/SharpRestriction/GrowthBounds/SegmentDegree/RestCountOf.lean}{Formalism.MainNotions.SelectionLemmas.Pf.RestSum}

\begin{proof}
We use the notation of Definitions~\ref{pd-symbols-run-system}, \ref{pd-symbols-residue-coarse}
and~\ref{pd-symbols-residue-families}.

(3) Let $\alpha\in\mathcal L$, $e\in E(\alpha)$, $j\in\mathbf Z$ and $b=b^{\alpha,e}_j$, and for $x\in\mathbf Z$
put $s(x)=\sum_c\alpha_{\bar x,c}$.

\emph{Step 1: remainders.} Let $j'\in\mathbf Z$ and let $x$ be a new element of $\Lambda_{j'}$. Each
$\alpha_{\bar x,c}$ is less than $P_{j'}(x)$, so $s(x)<kP_{j'}(x)$; and as $s(x)\leq d^D(\bar x)$ and
$s(x)\equiv d^D(\bar x)\pmod{P_{j'}(x)}$, $P_{j'}(x)$ divides $d^D(\bar x)-s(x)$ and
$n^\alpha_{j'}(x)=(d^D(\bar x)-s(x))/P_{j'}(x)$.

\emph{Step 2: the twists $Q_{j'}$.} Let $j'\in\mathbf Z$ and $D_{j'}=p^{\operatorname{dist}_{j'}(f(j'-1,0))}$.
By Step 1 of the proof of Lemma~\ref{pd-symbols-residue-lifts-numeric}, for $j'$ and for $j'+1$,
$p^{R_{j'}}D_{j'}\leq a/(8Yk^2N_s)<p^{R_{j'+1}+1}D_{j'+1}$; so
$R_{j'}+f(j',r-1)-f(j'-1,0)\leq R_{j'+1}+f(j'+1,r-1)-f(j',0)$. As $f(j'-1,0)<f(j',0)$, the exponent
$f(j'+1,r-1)-f(j',r-1)+R_{j'+1}-R_{j'}$ is positive and is the exponent of $Q_{j'}$. Hence
$P_{j'}(x)Q_{j'}=P_{j'+1}(x)$ for every $x\leq f(j',r-1)$. Moreover $R_{j'}$ and $Q_{j'}$ depend only on
the class of $j'$ modulo $J$, by condition~(1) of validity; so $e_{j'}(i)<rQ_{j'}$ for all $j'$ and $i<r$, by
the definition of $E(\alpha)$.

\emph{Step 3: balance.} Let $x\in\Lambda'_j$. Then $x$ is a new element of $\Lambda_j$ with $x<f(j,0)$, so
$b(x)=n^\alpha_j(x)=(d^D(\bar x)-s(x))/P_j(x)$ by Step 1, and $0\leq s(x)<kP_j(x)$ gives
$d^D(\bar x)/P_j(x)-k<b(x)\leq d^D(\bar x)/P_j(x)$. So $b$ is balanced at $j$.

\emph{Step 4: the old elements.} Let $x\in\Lambda_j$ not be new. As in the proof of
Lemma~\ref{pd-symbols-normal-degree}, $x=f(j-1,i)$ for one $i<r$, and $x$ is a new element of
$\Lambda_{j-1}$, as $f(j-2,r-1)<f(j-1,0)\leq x\leq f(j-1,r-1)$. Put $n=d^D(\bar x)$, $P=P_{j-1}(x)$,
$Q=Q_{j-1}$ and $\epsilon=e_{j-1}(i)$. By Step 2, $PQ=P_j(x)$ and $\epsilon<rQ$, and by Step 1, $s(x)<kP$
and $n^\alpha_{j-1}(x)=U$ with $U=\lfloor(n-s(x))/P\rfloor$. Let $A=\lfloor n/(PQ)\rfloor$, so that
$W_j(x)=A-(k+r+1)$, and $B=b(x)$, the integer part of $\max(U-\epsilon,0)/Q$. As $APQ\leq n$ and
$n-s(x)<(U+1)P$, $(AQ-k)P<(U+1)P$, so $AQ-k\leq U$. As $\max(U-\epsilon,0)<(B+1)Q$,
\[
 AQ<(B+1)Q+k+\epsilon<(B+1)Q+kQ+rQ,
\]
so $A<B+1+k+r$ and $W_j(x)\leq B$.

(4) Let the $G_j$ be as in (4). Write $B(n)=\binom{n+k-1}{k-1}$, which is at least $1$, and put
$\eta=(1-\frac5Y)(1-\frac1Y)(1-\frac r{2Y})$, $A=1-a^{-N_s^2}$,
$W'_j=\prod_{i<r}B(\max(0,W_j(f(j-1,i))-m_0))$ and $T'_j=\big(Q_j^{\,r-1}\max(0,1-\frac r{2Y})\big)^r$, so that
$S_j(b)=AW'_jT'_j\prod_{x\in\Lambda'_j}B(b(x))$. For $\alpha\in\mathcal L$ write $n^\alpha_j$ for the function
$x\mapsto n^\alpha_j(x)$. We use Lemma~\ref{pd-symbols-residue-count-sum-facts}, with $\mathcal N_j$ as there. If
$1-|T|r^2k/Y\leq0$, then $\beta\leq0$, and the right side of (4) is at least $0$; so let $1-|T|r^2k/Y>0$.

\emph{Step 5: the total loss.} By (6) of that lemma, $a\geq32Y\geq1$, and $N_s\geq1$; so $1-\frac1{32Y}\leq A\leq1$. Put
$u=1/Y$. As $|T|\geq1$, $r^2ku<1$; with $r\geq2$ and $k\geq5$ this gives $5u,u,ru/2\leq1$ and $7+r/2\leq r^2k$. Hence
$A\eta\geq(1-\frac u{32})(1-5u)(1-u)(1-\frac{ru}2)\geq1-(7+\frac r2)u\geq0$, and by Bernoulli's inequality and $J\leq|T|$
(part (1) there),
\[
 1-\frac{|T|r^2k}Y\leq\Big(1-\big(7+\tfrac r2\big)u\Big)^{|T|}\leq(A\eta)^{|T|}\leq A^J\eta^{|T|}.
\]

\emph{Step 6: one family.} Let $\alpha\in\mathcal L$ and $0\leq j<J$. As in Step~3, $n^\alpha_j$ is balanced at $j$.
By Definition~\ref{pd-symbols-residue-families}(3), $b^{\alpha,e}_j$ and $\tau_j(n^\alpha_j,e_j)$ agree at the elements
of $\mathcal N_j$, so $G_j(b^{\alpha,e}_j)=G_j(\tau_j(n^\alpha_j,e_j))$ by the first property in (2). The set $E(\alpha)$
consists of the families $(e_j)_{j<J}$ with $e_j$ in the set $E_j$ of the $\varepsilon\in\mathbf Z_{\geq0}^r$ with
$\varepsilon(i)<rQ_j$ and $\varepsilon(i)\equiv n^\alpha_j(f(j,i))\pmod{Q_j}$. So the last property in (2), with
$\rho(i)=n^\alpha_j(f(j,i))$, gives
\[
 \sum_{e\in E(\alpha)}\prod_{j=0}^{J-1}G_j(b^{\alpha,e}_j)=\prod_{j=0}^{J-1}\sum_{\varepsilon\in E_j}G_j(\tau_j(n^\alpha_j,\varepsilon))
 \geq\prod_{j=0}^{J-1}S_j(n^\alpha_j).
\]

\emph{Step 7: the sum over the families.} For $t\in\mathbf Z/n'\mathbf Z$ let $\mathcal L_t$ be the set of the remainder
vectors at $t$, so that $\mathcal L$ is the product of the $\mathcal L_t$, and put
\[
 \sigma_t=\sum_{v\in\mathcal L_t}\ \prod_{j=0}^{J-1}\ \prod_{x\in\Lambda'_j,\ \bar x=t}
 B\Big(\Big\lfloor\frac{d^D(t)-\sum_cv_c}{P_j(x)}\Big\rfloor\Big).
\]
Expanding the product over $t$ of these sums gives
$\sum_{\alpha\in\mathcal L}\prod_{j<J}S_j(n^\alpha_j)=\prod_{j<J}(AW'_jT'_j)\cdot\prod_t\sigma_t$.

\emph{Step 8: the factors $\sigma_t$.} Every summand of $\sigma_t$ is at least $1$. If $t\notin T$, no element of any
$\Lambda_j$ has residue $t$, so $0\in\mathcal L_t$ and $\sigma_t\geq1$. Let $0\leq j<J$, $x\in\Lambda'_j$ and $t=\bar x$.
By (3) of Lemma~\ref{pd-symbols-residue-count-sum-facts}, every $v$ with $v_c<P_j(x)$ for all $c$, $\sum_cv_c\leq d^D(t)$
and $\sum_cv_c\equiv d^D(t)\pmod{P_j(x)}$ lies in $\mathcal L_t$. The product in $\sigma_t$ has the factor for $(j,x)$,
and its other factors are at least $1$; so Lemma~\ref{pd-symbols-residue-count-digits}(1), with $n=d^D(t)$ and
$P=P_j(x)$, gives $\sigma_t\geq B(d^D(t))$. For $x=f(j,i)$, (4) of Lemma~\ref{pd-symbols-residue-count-sum-facts}
gives $kP_j(x)\leq d^D(t)$, and Lemma~\ref{pd-symbols-residue-count-digits}(2) gives $\sigma_t\geq|\mathcal L_t|\geq
P_j(x)^{k-1}$. By (1) and (2) of Lemma~\ref{pd-symbols-residue-count-sum-facts},
\[
 \prod_t\sigma_t\geq\prod_{t\in T}\sigma_t=\prod_{j=0}^{J-1}\Big(\prod_{x\in\Lambda'_j}\sigma_{\bar x}\cdot
 \prod_{i<r}\sigma_{\overline{f(j,i)}}\Big)\geq\prod_{j=0}^{J-1}\Big(\prod_{x\in\Lambda'_j}B(d^D(\bar x))\cdot
 \prod_{i<r}P_j(f(j,i))^{k-1}\Big).
\]

\emph{Step 9: comparison.} By (5) of Lemma~\ref{pd-symbols-residue-count-sum-facts},
$p^{n'rK}=\prod_{j<J}Q_j^{\,rK}$ and $\prod_{j<J}W'_j=\prod_{j<J}W'_{j+1}$. As $1-\frac r{2Y}\geq0$ and
$P^{k-1}Q^{r-1}Q^K=(PQ)^{k-1}$, Steps 7 and 8 give
\[
 p^{n'rK}\sum_{\alpha\in\mathcal L}\prod_{j<J}S_j(n^\alpha_j)\geq\prod_{j=0}^{J-1}A\prod_{x\in\Lambda'_j}B(d^D(\bar x))
 \prod_{i<r}\Big(1-\frac r{2Y}\Big)\big(P_j(x_i)Q_j\big)^{k-1}B\big(\max(0,W_{j+1}(x_i)-m_0)\big),
\]
with $x_i=f(j,i)$. By Step~5 and (1), (2) of the same lemma,
\[
 \Big(1-\frac{|T|r^2k}Y\Big)\prod_{t\in T}B(\deg(t))\leq\prod_{j=0}^{J-1}A\prod_{x\in\Lambda'_j}\eta B(\deg(\bar x))
 \prod_{i<r}\eta B(\deg(\bar x_i)).
\]
For $x\in\Lambda'_j$, $\eta B(\deg(\bar x))\leq(1-\frac5Y)B(\deg(\bar x))\leq B(d^D(\bar x))$ by (7) there. For $x_i$, by
(7) and (8) there, $\eta B(\deg(\bar x_i))\leq(1-\frac r{2Y})(1-\frac1Y)B(d^D(\bar x_i))\leq(1-\frac r{2Y})
(P_j(x_i)Q_j)^{k-1}B(\max(0,W_{j+1}(x_i)-m_0))$. Comparing the factors, $p^{n'rK}\beta$ is at most the left side of
the first display, and Step~6 gives (4).

(1) Let $\Phi$ be the set of the pairs $(\alpha,e)$ with $\alpha\in\mathcal L$ and $e\in E(\alpha)$, and let $\mathcal
S_j(\alpha,e)$, $X^\alpha$ and $\mu(v,c)$ be as in Lemma~\ref{pd-symbols-residue-count-inputs-monomials}, and
$\mathcal N_j$ the set of the new elements of $\Lambda_j$.

\emph{Step 10: the multidegree.} Let $(\alpha,e)\in\Phi$ and $\mu_j\in\mathcal S_j(\alpha,e)$ for $0\leq j<J$. We show
that $X^\alpha\prod_{j<J}\psi_j(\mu_j)$ has multidegree $d^D$. Write $b=b^{\alpha,e}$, and let $s\in\mathbf Z/n'\mathbf Z$.
If $s\notin T$, then $\alpha_s=0$ as in Step~2 of the proof of
Lemma~\ref{pd-symbols-residue-count-inputs-monomials}(2), and by part (1) of that lemma no $X_{s,c}$ occurs in
$\prod_j\psi_j(\mu_j)$; so the degree in $X_{s,1},\dots,X_{s,k}$ is $0=d^D(s)$. Let $t\in T$. By
Lemma~\ref{pd-symbols-residue-count-sum-facts}(1) there are $0\leq j_0<J$ and $x_0\in\mathcal N_{j_0}$ with $\bar x_0=t$.
Put $P=P_{j_0}(x_0)$, $Q=Q_{j_0}$, $n=n^\alpha_{j_0}(x_0)$ and $A=\sum_c\alpha_{t,c}$. By
Lemma~\ref{pd-symbols-residue-count-inputs-monomials}(1), and as $\mu_j$ has degree $b_j(x)$ in the row $z_x$, the degree
in $X_{t,1},\dots,X_{t,k}$ is $A+P\big(b_{j_0}(x_0)+QB\big)$, where $B$ is the sum of the $b_j(x)$ over the $0\leq j<J$
and the $x\in\Lambda_j\smallsetminus\mathcal N_j$ with $\bar x=t$. We show that $b_{j_0}(x_0)+QB=n$.

If $x_0<f(j_0,0)$, then $x_0\in\Lambda'_{j_0}$ and $b_{j_0}(x_0)=n$; and $B=0$, since an $x$ in the sum $B$ would give
$x_0=f(j_0,i)\geq f(j_0,0)$ by Lemma~\ref{pd-symbols-residue-count-inputs-old}(3).

If $x_0\geq f(j_0,0)$, then $x_0=f(j_0,i_0)$ for some $i_0<r$ by Lemma~\ref{pd-symbols-residue-count-sum-facts}(2), and
$b_{j_0}(x_0)=e_{j_0}(i_0)$. Let $j_1=j_0+1$ if $j_0+1<J$ and $j_1=0$ otherwise, and $x_1=f(j_1-1,i_0)$, which is $x_0$
or $x_0-n'$. Then $x_1\in\Lambda_{j_1}$ is not new and $\bar x_1=t$, and by
Lemma~\ref{pd-symbols-residue-count-inputs-old}(2) it is the only $x$ in the sum $B$; so $B=b_{j_1}(x_1)$. By
Definition~\ref{pd-symbols-residue-families}(3) and the periodicity of $n^\alpha$, $e$ and $Q$ in $j$, $b_{j_1}(x_1)$ is
the integer part of $\max(n-e_{j_0}(i_0),0)/Q$. By (3), $W_{j_1}(x_1)\leq b_{j_1}(x_1)$. Also $W_{j_1}(x_1)\geq1$: as
$Y,N_s\geq1$ and $r\leq k$, Lemma~\ref{pd-symbols-residue-count-sum-facts}(4) gives $d^D(t)\geq\frac{36}5k^2P_{j_1}(x_1)
\geq(k+r+2)P_{j_1}(x_1)$. So $n-e_{j_0}(i_0)\geq Q$. As $e\in E(\alpha)$, $e_{j_0}(i_0)\equiv n\pmod Q$, so $Q$ divides
$n-e_{j_0}(i_0)$, $B=(n-e_{j_0}(i_0))/Q$ and $b_{j_0}(x_0)+QB=n$.

In both cases the degree is $A+Pn$. As $\alpha_t$ is a remainder vector at $t$, $A\leq d^D(t)$ and $P$ divides
$d^D(t)-A$; so $Pn=d^D(t)-A$, and the degree is $d^D(t)$.

\emph{Step 11: bases.} Let $(\alpha,e)\in\Phi$ and $0\leq j<J$. As $b^{\alpha,e}_j$ vanishes off $\Lambda_j$,
Lemma~\ref{pd-symbols-residue-count-inputs-reductions} shows that $M_j(b^{\alpha,e}_j)$ is free of finite rank
$n_{\alpha,e,j}=\operatorname{rk}M_j(b^{\alpha,e}_j)$; let $y_{\alpha,e,j,1},\dots,y_{\alpha,e,j,n_{\alpha,e,j}}$ be a
basis, whose reductions are linearly independent over $\kappa$. Every monomial of an element of $M_j(b^{\alpha,e}_j)$
lies in $\mathcal S_j(\alpha,e)$: by Definition~\ref{pd-symbols-module}(2) it involves only the $z_{v,c}$ with
$v\in\Lambda_j$, and with $c\leq r$ unless $v$ is new, and the two other conditions are those of
Definition~\ref{pd-symbols-residue-coarse}(3). The same holds for the reductions.

\emph{Step 12: the count.} For $(\alpha,e)\in\Phi$ and a family $i=(i_j)_{j<J}$ with $1\leq i_j\leq n_{\alpha,e,j}$ put
$u_{\alpha,e,i}=\prod_{j<J}\psi_j(y_{\alpha,e,j,i_j})\cdot X^\alpha$. Its monomials are among the
$X^\alpha\prod_j\psi_j(\mu_j)$ with $\mu_j\in\mathcal S_j(\alpha,e)$, so it has multidegree $d^D$ by Step~10; and
$X^\alpha$ is a coarse label, since $\alpha_{\bar x,c}<P_j(x)$ for every $j$ and every new $x\in\Lambda_j$ by the
definition of a remainder vector. So $u_{\alpha,e,i}\in\Lambda^{\mathrm c}$, and its reduction is
$X^\alpha\prod_j\psi_j(\bar y_{\alpha,e,j,i_j})$. Apply Lemma~\ref{pd-symbols-residue-count-inputs-products} over
$\kappa$, with the variables $z_{v,c}$ and $X_{t,c}$, the maps $\psi_j$, the set $\Phi$, the monomials $X^\alpha$, the
sets $\mathcal S_j(\alpha,e)$ and the families $(\bar y_{\alpha,e,j,i})_i$ of Step~11; its hypothesis is
Lemma~\ref{pd-symbols-residue-count-inputs-monomials}(2). So the reductions of the $u_{\alpha,e,i}$ are linearly
independent over $\kappa$. There are $\sum_{(\alpha,e)\in\Phi}\prod_{j<J}n_{\alpha,e,j}$ of them, the left side of (1),
and they lie in the $\kappa$-span of the reductions of the elements of $\Lambda^{\mathrm c}$, which consists of
polynomials of multidegree $d^D$ and so has finite dimension. This gives (1).
% Lean: (1) is restLam_of D hA 1088 restLamDeg = restLam_of_rests 1088 (restLamA_of D hAdm) (restLamInj 1088)
%   restLamDeg. Step 10 is the second part of restLamDeg (lam_eval; hkey by cases on fr j0 0 <= x0: comb_frame_top,
%   bfam_at_w_letter_eq with lamd_nz_per, lamd_eZ_per and sum_Qc_per, the clause U2, restLamTau, restPad (part (3))
%   and lamd_Bw_pos from sum_aP; the clause LK otherwise; then hPnz). Step 11 is restLamA_of (Lemma
%   pd-symbols-residue-count-inputs-reductions) with lam_mem_lamS (XX'_supported). Step 12 is restLam_of_rests:
%   lam_li_all (Lemma pd-symbols-residue-count-inputs-products) with restLamInj (Lemma
%   pd-symbols-residue-count-inputs-monomials (2)), lam_map_phc, the first part of restLamDeg (IsCLabel), hmem
%   (membership in LamBar) and hfin (LamBar lies in the polynomials of multidegree a', finite_restrictSupport).

(2) Let $j\in\mathbf Z$, let $T_j(b)$, $\mathcal I_j(b)$ and $y_j(b;q)$ be as in
Definition~\ref{pd-symbols-residue-count-inputs-index}, and put $G_j(b)=|\mathcal I_j(b)|$.

\emph{Step 13: locality.} The sets $T_j(b)$ and $\mathcal I_j(b)$ depend only on the values of $b$ at the new elements
of $\Lambda_j$: $\beta_j(b,x)$ uses $b$ only at new $x$, and $\theta'_{j,b}(v)$ uses $b$ only at
$\operatorname{old}_j(v)$ when it is new. So $G_j(b)$ depends only on these values, and $G_j(b)\geq0$.

\emph{Step 14: the ranks.} Let $b$ vanish off $\Lambda_j$, satisfy $W_j(x)\leq b(x)$ for every $x\in\Lambda_j$ that is not
new, and be balanced at $j$. By Lemma~\ref{pd-symbols-residue-count-inputs-segment}(1), the $y_j(b;q)$, $q\in\mathcal
I_j(b)$, are linearly independent elements of $M_j(b)$, which is free of finite rank by
Lemma~\ref{pd-symbols-residue-count-inputs-reductions}. So $\mathcal I_j(b)$ is finite and $G_j(b)\leq\operatorname{rk}
M_j(b)$.

\emph{Step 15: the sum.} Let $b$ be balanced at $j$ and $\rho\in\mathbf Z_{\geq0}^r$. By the definition of $S_j(b)$, by
Lemma~\ref{pd-symbols-residue-count-inputs-segment}(2), multiplied by the nonnegative factor
$(Q_j^{\,r-1}\max(0,1-\frac r{2Y}))^r$, and by part (3) of that lemma, which applies as $T_j(b)$ is finite by its part
(2),
\[
 S_j(b)\leq|T_j(b)|\Big(Q_j^{\,r-1}\max\Big(0,1-\frac r{2Y}\Big)\Big)^r\leq\sum_\varepsilon G_j(\tau_j(b,\varepsilon)),
\]
with the sum of (2). This gives (2).
% Lean: (2) is restSeg_of 1088 (restMem_of D hA) (restInd D hA) restFin (good_path_weight D hA) restMon, with
%   G b = Nat.card (Idx j (restr j b)): Step 13 is the locality clause (restr), Step 14 is RestMem, RestInd and
%   restFin with LinearIndependent.fintype_card_le_finrank, and Step 15 is GoodPathWeight and RestMon through
%   Dat.segBound (Lemma pd-symbols-residue-count-inputs-segment).
% Lean: (4) is restSum = restSum_of (sum_facts, Ns_one_le): the case split on 1 - |T| r^2 k / Y, Step 5 is
%   sum_total_loss with first_sum_factor_bounds; Step 6 is step1 (sum_bfam_late, hGdep, Finset.prod_univ_sum over
%   Fintype.piFinset sumTset, hGsum with family_degrees_admissible); Step 7 is step2 (sum_LA, Dat.sum_segBound_eq;
%   the book's sigma_t is sumSt, with laSet the remainder vectors and sumhT the product); Step 8 is step3 (F.tiling,
%   F.split, sum_St_LL with sum_digit_count, sum_St_fr with sum_label_count, sum_St_nontop); Step 9 is hW (perW),
%   hQ (prodQ), regroup, hfin (tiling) and sum_fac_LL, sum_fac_fr (compA, compW).
% Lean: (3) is restPad: Step 1 is family_degrees_admissible (with sum_LLsub) and sum_la_lt; Step 2 is
%   sum_PzQc (Rj_add_first_frame_gap_le_last_frame_gap_add_Rj_succ, frame_zero_strict_mono), sum_Qc_per,
%   sum_Rj_per and eZ_lt_mul_Qc; Step 4 is comb_frame_top, sum_frlate and div_mul_sub_le_div_sub_div.
% Lean: (1) RestLam 1088 = restLam_of D hA 1088 restLamDeg; (2) RestSeg 1088 = restSeg_of 1088 (restMem_of D hA)
%   (restInd D hA) restFin (good_path_weight D hA) restMon; (3) RestPad 1088 = restPad; (4) RestSum 1088 1 =
%   restSum. restLam_of, restMem_of, restInd and good_path_weight use the AdmB clauses (the extremal
%   projectors of Lemma pd-extremal-projector and Lemma pd-subspace-count). The twelve SumFacts of (4) are proved in
%   FrameZeroStrictMono.lean. What remains open of (2) is marked in the proof of Lemma
%   pd-symbols-residue-count-inputs-segment-atypical.
\end{proof}

\begin{lemma}\label{pd-symbols-residue-generators}
Let $C_{La}\geq1088$ and $C_k\geq1088$ be real numbers, and let $p$, $\mathcal R$, $\kappa$, $n'$, $r$,
$K$, $k$, $T$, $L_1$, $N_1$, $(J,f)$, $\deg$, $E'$, $N_s$, $E_0$, $m_*$, $a$, $Y$, $q$, $X_j$ and $d$ be
as in Lemma~\ref{pd-symbols-residue-rank}. Let $D$ and $\Lambda^{\mathrm c}$ be the discriminant and
the coarse generators of Definition~\ref{pd-symbols-residue-coarse}. Then:
\begin{enumerate}
\item for every $u\in\Lambda^{\mathrm c}$ there is a symbol $P\in\mathrm{Gen}(X_\bullet,d)$ whose
  reduction modulo the maximal ideal is that of $Du$;
\item the $\kappa$-span of the reductions of the elements of $\Lambda^{\mathrm c}$ has dimension at least
  \[
   p^{-n'rK}\Big(1-\frac{|T|r^2k}{Y}\Big)\prod_{t\in T}\binom{\deg(t)+k-1}{k-1}.
  \]
\end{enumerate}
\end{lemma}
\leanref{Formalism/SharpRestriction/GrowthBounds/SegmentDegree/Coarse.lean}{Formalism.MainNotions.SelectionLemmas.Pf.GeneratorModuloUniformiser}
\leanref{Formalism/SharpRestriction/GrowthBounds/SegmentDegree/Coarse.lean}{Formalism.MainNotions.SelectionLemmas.Pf.RestCount}
\leanref{Formalism/SharpRestriction/GrowthBounds/SegmentDegree/DiscriminantNonzeroModUniformiser.lean}{Formalism.MainNotions.SelectionLemmas.Pf.generator_modulo_uniformiser_of_adm}
\leanref{Formalism/SharpRestriction/GrowthBounds/SegmentDegree/RestCountOf.lean}{Formalism.MainNotions.SelectionLemmas.Pf.restCount_of}
\leanref{Formalism/SharpRestriction/GrowthBounds/SegmentDegree/DiscriminantNeZeroOfResidueNeZero.lean}{Formalism.MainNotions.SelectionLemmas.Pf.generator_of_segment_reductions}
\leanref{Formalism/SharpRestriction/GrowthBounds/SegmentDegree/DiscriminantNeZeroOfResidueNeZero.lean}{Formalism.MainNotions.SelectionLemmas.Pf.Dat.zSeg}
\leanref{Formalism/SharpRestriction/GrowthBounds/SegmentDegree/DiscriminantNeZeroOfResidueNeZero.lean}{Formalism.MainNotions.SelectionLemmas.Pf.phc_eq}
\leanref{Formalism/SharpRestriction/GrowthBounds/SegmentDegree/CoarseLabelFactorisation.lean}{Formalism.MainNotions.SelectionLemmas.Pf.coarse_label_factorisation}
\leanref{Formalism/SharpRestriction/GrowthBounds/SegmentDegree/SegmentProductDegreeDetermined.lean}{Formalism.MainNotions.SelectionLemmas.Pf.segment_image_multihomogeneous}
\leanref{Formalism/SharpRestriction/GrowthBounds/SegmentDegree/SegmentProductDegreeDetermined.lean}{Formalism.MainNotions.SelectionLemmas.Pf.segment_product_degree_determined}
\leanref{Formalism/SharpRestriction/GrowthBounds/SegmentDegree/DiscriminantMultihomogeneous.lean}{Formalism.MainNotions.SelectionLemmas.Pf.discriminant_multihomogeneous}
\leanref{Formalism/SharpRestriction/GrowthBounds/SegmentDegree/DiscriminantNeZeroOfResidueNeZero.lean}{Formalism.MainNotions.SelectionLemmas.Pf.discriminant_ne_zero_of_residue_ne_zero}

\begin{proof}
We use the notation of Definitions~\ref{pd-symbols-run-system}, \ref{pd-symbols-module}
and~\ref{pd-symbols-residue-coarse}, and write bars for reductions modulo the maximal ideal of
$\mathcal R$. A polynomial $x$ in the $z_{v,c}$ has \emph{row degrees} $\beta\colon\mathbf Z\to
\mathbf Z_{\geq0}$ if each of its monomials has degree $\beta(v)$ in $z_{v,1},\dots,z_{v,k}$ for every
$v$.

(1) Let $u=\prod_{j=0}^{J-1}\psi_j(y_j)\cdot X^\ell\in\Lambda^{\mathrm c}$, with $y_j\in M_j(b_j)$ and
$X^\ell$ a coarse label. If $u=0$, then $P=0=\prod_{j<J}\phi_j(0)\cdot X^0$ is a symbol, since $X^0$ is a
label and $0$ has multidegree $d$. Let $u\neq0$.

\emph{Step 1: the label.} By the proof of Lemma~\ref{pd-symbols-normal-degree}, the new elements $v$ of
the $\Lambda_j$, $0\leq j<J$, have distinct residues. For such $j$ and $v$ and $c\leq k$ put
$e_j(v,c)=\lfloor\ell(\bar v,c)/p^{\operatorname{dist}_j(v)}\rfloor$, and $e_j(v,c)=0$ for the other $v$;
let $\ell_0(\bar v,c)=\ell(\bar v,c)-p^{\operatorname{dist}_j(v)}e_j(v,c)$ for these $j$ and $v$, and
$\ell_0(t,c)=\ell(t,c)$ at the other residues $t$. Then $X^\ell=X^{\ell_0}\prod_{j<J}\phi_j(z^{e_j})$ with
$z^{e_j}=\prod_{v,c}z_{v,c}^{e_j(v,c)}$, as $\phi_j(z_{v,c})=X_{\bar v,c}^{\,p^{\operatorname{dist}_j(v)}}$,
and $e_j(v,c)<p^{R_j}$ as $X^\ell$ is a coarse label. The monomial $X^{\ell_0}$ is a label. Let $i\in
\mathbf Z$, $w\in\Lambda_i$ new and $c\leq k$; write $i=j+mJ$ with $0\leq j<J$. By condition~(1) of a run
system $\Lambda_i=\Lambda_j+mn'$, so $v=w-mn'$ is new in $\Lambda_j$, $\bar v=\bar w$ and
$\operatorname{dist}_j(v)=\operatorname{dist}_i(w)$; and $\ell_0(\bar w,c)$ is the remainder of
$\ell(\bar v,c)$ modulo $p^{\operatorname{dist}_j(v)}$, which is less than $p^{\operatorname{dist}_i(w)}$.

\emph{Step 2: the factorization.} Let $\epsilon_j$ be the $\mathcal R$-algebra endomorphism of
$\mathcal R[z]$ with $\epsilon_j(z_{v,c})=z_{v,c}^{\,p^{R_j}}$; then $\psi_j=\phi_j\circ\epsilon_j$. Let
$D^{\mathrm{seg}}_j=\prod_v\det(z_{\rho_{j,v}(l),b})_{l,b}^{\,\sigma_j(v)}$, over the new $v\in\Lambda_j$,
so that $D=\prod_{j<J}\phi_j(D^{\mathrm{seg}}_j)$, and put
$z_j=D^{\mathrm{seg}}_j\,\epsilon_j(y_j)\,z^{e_j}\in\mathcal R[z]$. By Step 1,
\[
 Du=\prod_{j=0}^{J-1}\phi_j(z_j)\cdot X^{\ell_0}.
\]

\emph{Step 3: row degrees.} The polynomial $y_j$ has row degrees $b_j$, so $\epsilon_j(y_j)$ has row
degrees $p^{R_j}b_j$; each determinant $\det(z_{\rho_{j,v}(l),b})$ has degree $1$ in each of the
distinct rows $z_{\rho_{j,v}(l)}$ (the proof of Lemma~\ref{pd-symbols-residue-discriminant}); so $z_j$
has row degrees $\beta_j$ for some $\beta_j$. If $x\neq0$ has row degrees $\beta_j$, then $\phi_j(x)$
has the multidegree $t\mapsto\sum_{\bar v=t}p^{\operatorname{dist}_j(v)}\beta_j(v)$, which depends only on
$\beta_j$. Hence any nonzero product $\prod_{j<J}\phi_j(x_j)\cdot X^{\ell_0}$ with $x_j$ of row degrees
$\beta_j$ has one and the same multidegree.

\emph{Step 4: segment reductions.} We show that for each $0\leq j<J$ there is $x_j\in X_j$ with row
degrees $\beta_j$ and $\bar x_j=\bar z_j$. Fix such a $j$. As $u\neq0$ has multidegree $d^D$,
Lemma~\ref{pd-symbols-residue-top-degree} gives $d^D(\bar v)/p^{\operatorname{dist}_j(v)+R_j}-k<b_j(v)$
for every $v\in\Lambda'_j$. By Step 1, $e_j(v,c)<p^{R_j}$, and $e_j(v,c)=0$ unless $v$ is a new element of
$\Lambda_j$. By Definition~\ref{pd-symbols-residue-coarse}(3), $y_j$ lies in the segment module
$X_j(m_0,\theta'_{j,b_j},o'_j)$. As $|\Lambda_j|+r\leq N_s+r\leq K$ by (2) of
Theorem~\ref{pd-symbols-count} and $m_0\geq0$, Lemma~\ref{pd-symbols-residue-segment-span} writes $y_j$
as an $\mathcal R$-linear combination $\sum_wc_ww$ of segment generators for
$(j,m_0,\theta'_{j,b_j},o'_j)$. The map $y\mapsto D^{\mathrm{seg}}_j\epsilon_j(y)z^{e_j}$ is
$\mathcal R$-linear, and by Lemma~\ref{pd-symbols-residue-lifts} each
$D^{\mathrm{seg}}_j\epsilon_j(w)z^{e_j}$ has the same reduction as some $x_w\in X_j$. So
$x=\sum_wc_wx_w\in X_j$ and $\bar x=\bar z_j$.

Let $x_j$ be the sum of the terms of $x$ of row degrees $\beta_j$. Its reduction is the sum of the terms
of $\bar z_j$ of row degrees $\beta_j$, which is $\bar z_j$ by Step 3. It remains to see that
$x_j\in X_j$. Taking the part of given row degrees is $\mathcal R$-linear and keeps the bounds on
$\operatorname{ord}_v$ of Definition~\ref{pd-symbols-module}(2), so it suffices that the part of row
degrees $\beta_j$ of a product $AG$ as there is a sum of such products. Write $A=\sum_{\beta'}A_{\beta'}$
and $G=\sum_{\beta''}G_{\beta''}$ as sums of their parts of given row degrees; the part of $AG$ of row
degrees $\beta_j$ is the sum of the $A_{\beta'}G_{\beta''}$ with $\beta'+\beta''=\beta_j$. Each
$A_{\beta'}$ involves only variables of $A$, and each $G_{\beta''}$ only variables of $G$. The
substitution $P\mapsto P_{S,i}$ of Definition~\ref{pd-symbols-module}(1) replaces each row by linear forms
in the same row, so it commutes with taking the part of given row degrees; hence every exponent of
$\eta$ that occurs in $(A_{\beta'})_{S,i}$ occurs in $A_{S,i}$, and
$\operatorname{val}_S(A_{\beta'})\geq\operatorname{val}_S(A)$ and $\lambda_S(G_{\beta''})\leq\lambda_S(G)$.
So each $A_{\beta'}G_{\beta''}$ satisfies the conditions of Definition~\ref{pd-symbols-module}(2) with the
same $M$.

\emph{Step 5: the multidegree.} By Definition~\ref{pd-symbols-residue-coarse}(1), $D$ has multidegree
$E^D$, and $u$ has multidegree $d^D=\max(d-E^D,0)$. By Lemma~\ref{pd-symbols-residue-top-count}, which
applies as $C_k\geq32$, $E^D\leq d$; so $d^D=d-E^D$, and $Du$ has multidegree $d$. It is nonzero:
$\bar D\neq0$ by Lemma~\ref{pd-symbols-residue-discriminant}, and $\mathcal R[X]$ is a domain.

\emph{Step 6: the symbol.} Put $P=\prod_{j<J}\phi_j(x_j)\cdot X^{\ell_0}$. By Steps 2, 3 and 5, every
nonzero product as in Step 3 has multidegree $d$, so $P$ has multidegree $d$. As $x_j\in X_j$ and
$X^{\ell_0}$ is a label, $P\in\mathrm{Gen}(X_\bullet,d)$, and
$\bar P=\prod_j\phi_j(\bar x_j)\cdot X^{\ell_0}=\prod_j\phi_j(\bar z_j)\cdot X^{\ell_0}=\overline{Du}$.

(2) We use the notation of Definition~\ref{pd-symbols-residue-families} and
Lemma~\ref{pd-symbols-residue-count-inputs}. Let the $G_j$ be those of part (2) of that lemma. By its part
(4), $\beta\leq\sum_{\alpha,e}\prod_{j<J}G_j(b^{\alpha,e}_j)$. Each $b^{\alpha,e}_j$ vanishes off $\Lambda_j$, and
by part (3) it is balanced at $j$ and at least $W_j$ at the elements of $\Lambda_j$ that are not new; so
$0\leq G_j(b^{\alpha,e}_j)\leq\operatorname{rk}M_j(b^{\alpha,e}_j)$ by part (2), and the products and sums
compare termwise. By part (1), the sum of the products of the ranks is at most the dimension of the
$\kappa$-span of the reductions of the elements of $\Lambda^{\mathrm c}$, which is therefore at least $\beta$.
% Lean (restRank_of 1088 1, applied under AdmB D in SelectionLemmas.Pf.second_estimate_of_restrictions):
% Part (1) = GeneratorModuloUniformiser 1088 = generator_modulo_uniformiser_of_adm D hAdm 1088, which is
%   generator_of_segment_reductions; Steps 1 to 3, 5 and 6 above follow it (coarse_label_factorisation,
%   phc_eq, segment_image_multihomogeneous, discriminant_multihomogeneous, discriminant_ne_zero_of_residue_ne_zero,
%   segment_product_degree_determined); the bound E^D <= d of Step 5 is (TopCountLowerBounds 1088).2,
%   Lemma pd-symbols-residue-top-count. Step 4 = SegmentReductionLiftsMultihomogeneous 1088 =
%   segment_reduction_of_column_words D hB12 1088 hWord, with its hypothesis Step2 from top_degree_comparison
%   (Lemma pd-symbols-residue-top-degree): h1 is the first clause of NestedTypeBound (Lemma
%   pd-symbols-residue-segment-span), the span induction uses zSeg_add and zSeg_smul, hW is
%   ColumnWordReductionLifts 1088 (Lemma pd-symbols-residue-lifts), and the row-degree part is polComp with
%   XX'_polComp (exists_sum_polComp, valS_le_polComp, lowest_degree_compose_le, sDegs_polComp_subset).
%   The first clause of NestedTypeBound (Lemma pd-symbols-residue-segment-span) and top_degree_comparison
%   (Lemma pd-symbols-residue-top-degree) are proved in the book (round nine). hWord =
%   column_word_reduction_of_four D hB12 1088 is Lemma pd-symbols-residue-lifts, proved in the book (round
%   ten), with its numeric inputs, Lemma pd-symbols-residue-lifts-numeric (round eleven).
% Part (2) = RestCount 1088 1 = restCount_of 1088 1 (restLam_of D hA 1088 restLamDeg), which the book's
%   proof follows; its four inputs are Lemma pd-symbols-residue-count-inputs, of which (1), (2) and (4)
%   are open (round eleven):
%   (restLam_of D hA 1088 restLamDeg)
%   (restSeg_of 1088 (restMem_of D hA) (restInd D hA) restFin (good_path_weight D hA) restMon) restPad restSum.
% The AdmB clauses are supplied in plane_count_theorem_with_constants from tool_statements_hold
%   (ExtremalProjectorTools, the extremal projectors of Lemma pd-extremal-projector) and rational_planes_count
%   (Lemma pd-subspace-count).
\end{proof}

\begin{lemma}\label{pd-symbols-residue-rank}
Let $C_{La}\geq1088$ and $C_k\geq1088$ be real numbers. Let $p\geq13$ be a prime, $\mathcal R$ a
complete discrete valuation ring in which $p$ is neither $0$ nor a unit, with residue field $\kappa$,
$n'\geq2$ an even integer, $r\geq2$ and $K\geq1$ integers, $k=r+K$, $T\subset\mathbf Z/n'\mathbf Z$,
$L_1,N_1\geq0$ integers, and $(J,f)$ a run system for $T$ of length $r$ valid for $(L_1,N_1)$. Let
$\deg$, $E'$, $N_s$, $E_0$, $m_*$, $a$, $Y$ and $q$ satisfy the hypotheses (1) to (4) of
Theorem~\ref{pd-symbols-count} with this $C_k$. Let $X_j$ be the symbol spaces for
$(C_{La},\deg,L_1,N_s,m_*,a,Y)$ over $\mathcal R$, and $d$ as in
Lemma~\ref{pd-symbols-normal-degree}. Then the $\kappa$-span of the reductions $\bar P$ of the symbols
$P\in\mathrm{Gen}(X_\bullet,d)$ has dimension at least
\[
 p^{-n'rK}\Big(1-\frac{|T|r^2k}{Y}\Big)\prod_{t\in T}\binom{\deg(t)+k-1}{k-1}.
\]
\end{lemma}
\leanref{Formalism/SharpRestriction/GrowthBounds/SegmentDegree/NonemptyDvrSettingOfStandard.lean}{Formalism.MainNotions.SelectionLemmas.Pf.RestRank}
\leanref{Formalism/SharpRestriction/GrowthBounds/SegmentDegree/Coarse.lean}{Formalism.MainNotions.SelectionLemmas.Pf.restRank_of}

\begin{proof}
Let $D$ and $\Lambda^{\mathrm c}$ be the discriminant and the coarse generators of
Definition~\ref{pd-symbols-residue-coarse}, formed with the data of the statement, and write bars for
reductions modulo the maximal ideal of $\mathcal R$. By Lemma~\ref{pd-symbols-residue-discriminant},
$\bar D\neq0$; as $\kappa[X]$ is a domain, multiplication by $\bar D$ is an injective $\kappa$-linear
map of $\kappa[X]$ into itself. By Lemma~\ref{pd-symbols-residue-generators}(1) it sends $\bar u$, for
$u\in\Lambda^{\mathrm c}$, to the reduction of a symbol, so it maps the $\kappa$-span $\bar\Lambda$ of
these $\bar u$ into the $\kappa$-span $\bar N$ of the reductions of the symbols. The symbols have
multidegree $d$, so $\bar N$ lies in the finite-dimensional space of the polynomials of multidegree $d$
over $\kappa$. Hence $\dim\bar N\geq\dim\bar\Lambda$, which is at least the stated bound by
Lemma~\ref{pd-symbols-residue-generators}(2).
% Lean: restRank_of 1088 1; what is open in Lemma pd-symbols-residue-generators is marked in its proof.
\end{proof}

\begin{theorem}\label{pd-symbols-count}
There are real numbers $C_{La},C_k,C_{rk},C_{\deg}$ with the following property. Let $p\geq13$ be a
prime, $(W,Q,\sigma,\iota)$ an unramified quadratic Hida setting of degree $n'$ at $p$ with rings
$\mathcal R\subset\mathcal K$ (Definition~\ref{pd-projection-property}), $r\geq2$ and $K\geq1$
integers, $k=r+K$, $T\subset\mathbf Z/n'\mathbf Z$, $L_1,N_1\geq0$ integers, and $(J,f)$ a run
system for $T$ of length $r$ valid for $(L_1,N_1)$ (Definition~\ref{pd-symbols-run-system}). Let
$\deg,E'\colon\mathbf Z/n'\mathbf Z\to\mathbf Z_{\geq0}$ be functions, $N_s,E_0,m_*\geq0$ integers and
$a,Y,q$ real numbers such that:
\begin{enumerate}
\item $a\leq\deg(t)\leq2a$ and $E'(t)\leq E_0$ for $t\in T$;
\item $|\Lambda_j|\leq N_s$ for every $j$, $N_s\leq N_1$, $N_s\leq2|T|$ and $N_s+r\leq K$;
\item $p^{\operatorname{dist}_j(v)}\leq q$ for every $j$ and every $v\in\Lambda_j$;
\item $Y\geq1$, $C_kN_s^2\log a\leq k$, $128k^2N_sY^2\leq p^{L_1-1}$ and
  $32pYk^2N_sq^2(E_0+m_*+1)\leq a$.
\end{enumerate}
Let $X_j$ be the symbol spaces for $(C_{La},\deg,L_1,N_s,m_*,a,Y)$ over $\mathcal R$ and
$L_a=C_{La}N_s^2\log a$ (Definition~\ref{pd-symbols-module}). Let $d(t)=\max(\deg(t)-E'(t),0)$ for
$t\in T$ and $d(s)=0$ for $s\notin T$, $N=N(X_\bullet,d)$, and
\[
 \beta=\Big(1-\frac{C_{rk}|T|r^2k}{Y}\Big)\prod_{t\in T}\binom{\deg(t)+k-1}{k-1}.
\]
Then:
\begin{enumerate}
\item there is an $\mathcal R$-submodule $S\subset N$ which is a direct summand of $\mathcal R[X]$ and
  has a basis consisting of symbols in $\mathrm{Gen}(X_\bullet,d)$, whose rank $n_S$ satisfies
  $n_S\geq p^{-n'rK}\beta$ and $n_S\geq\beta/G$, where $G$ is the number of subspaces of dimension
  $r$ of $\mathbf F^k$ for a field $\mathbf F$ with $p^{n'}$ elements;
\item every monomial of every element of $N$ has normal degree at most
  $C_{\deg}|T|\big(L_aa/(pk)+ap^{-L_1}+kq\big)$.
\end{enumerate}
\end{theorem}
\leanref{Formalism/SharpRestriction/GrowthBounds/ValuationLowestDegreeRename.lean}{Formalism.MainNotions.SelectionLemmas.Conclusion2}
\leanref{Formalism/SharpRestriction/GrowthBounds/PlaneCountTheoremWithConstants.lean}{Formalism.MainNotions.G.plane_count_theorem_with_constants}
\leanref{Formalism/SharpRestriction/GrowthBounds/SegmentDegree/NonemptyDvrSettingOfStandard.lean}{Formalism.MainNotions.SelectionLemmas.Pf.second_estimate_of_rank_bound}

\begin{proof}
Take $C_{La}=C_k=1088$, $C_{rk}=1$ and $C_{\deg}=10$, and let the data be as in the statement. By
Definitions~\ref{pd-local-model} and~\ref{pd-projection-property}, $n'$ is even and $\mathcal R$ is a
complete discrete valuation ring in which $p=\varpi_0^2$ is neither $0$ nor a unit; let $\kappa$ be
its residue field.

(2) This is Lemma~\ref{pd-symbols-normal-degree}, since $C_{La}\geq0$.

(1) Every symbol has multidegree $d$, so its monomials lie in the finite set of monomials of
multidegree $d$. Lemma~\ref{pd-symbols-select}, applied to the set $\mathrm{Gen}(X_\bullet,d)$, gives
an $\mathcal R$-submodule $S\subset N$ which is a direct summand of $\mathcal R[X]$, has a basis of
symbols, and has rank $n_S$ equal to the dimension of the $\kappa$-span of the reductions of the
symbols. By Lemma~\ref{pd-symbols-residue-rank}, $n_S\geq p^{-n'rK}\beta$. Let $\mathbf F$ be a field
with $p^{n'}$ elements and $G$ the number of subspaces of dimension $r$ of $\mathbf F^k$. By
Lemma~\ref{pd-subspace-count}, $G\geq(p^{n'})^{r(k-r)}=p^{n'rK}>0$. If $\beta\leq0$, then
$\beta/G\leq0\leq n_S$; otherwise $\beta/G\leq p^{-n'rK}\beta\leq n_S$.
% Lean: the constants are those of SelectionLemmas.Pf.secondPartClosed = second_estimate_of_restrictions
%   = second_estimate_of_rank_bound D 1088 1 restDeg restSel (restRank_of ...); what is open in
%   Lemma pd-symbols-residue-rank is marked in its proof.
\end{proof}

\subsection{Twisted substitutions and symbol multipliers}\label{subsection-pd-symbol-multipliers}

\begin{lemma}\label{pd-symbols-space}
Let $C_{La},C_k,C_{rk},C_{\deg}$ be real numbers with the property of
Theorem~\ref{pd-symbols-count}. Let the data of Situation~\ref{pd-edge-situation} be given with
$C\geq3041+|C_k|$, and let $(J,f)$ be a run system for $T$ of length $r$ that is spaced for
$(L_1,L_2,N_1,G')$ (Definition~\ref{pd-symbols-run-system}). Put $K=k-r$, $N_s=\min(N_1,2|T|)$,
$q=p^{G'}$, $c'=c_*+\omega$ and $\deg(s)=\lambda_0(s)_1-\lambda_0(s)_2$ for
$s\in\mathbf Z/n'\mathbf Z$; for $j\in\mathbf Z$ put $q_j=p^{f(j+1,r-1)-f(j,r-1)}$, let
$\alpha_0=1+r(k-r)(q_0-1)$ and $\alpha_j=r(k-r)(q_j-1)$ for $j\neq0$, and let $\alpha_{\max}$ be the
largest $\alpha_j$ with $0\leq j<J$. Put
\[
 m_*=(N_1+1)\big((2q+3)c'+kq+c'\alpha_{\max}+3c'\big),\qquad E_0=15c'(N_1+rk)q^2,
\]
and let $E'\colon\mathbf Z/n'\mathbf Z\to\mathbf Z_{\geq0}$ satisfy $E'(t)\leq E_0$ for $t\in T$. Then
$\deg$, $E'$, $N_s$, $E_0$, $m_*$, $a$, $Y$ and $q$ satisfy the hypotheses (1) to (4) of
Theorem~\ref{pd-symbols-count}, so its conclusions (1) and (2) hold for them.
\end{lemma}
\leanref{Formalism/SharpRestriction/GrowthBounds/ScaleEdges/GainBoundAt.lean}{Formalism.MainNotions.Pf14.symbol_space_at}
\leanref{Formalism/SharpRestriction/GrowthBounds/ScaleEdges/GainBoundAt.lean}{Formalism.MainNotions.Pf14.GainBoundAt}
\leanref{Formalism/SharpRestriction/GrowthBounds/ScaleEdges/GainBoundAt.lean}{Formalism.MainNotions.Pf14.E0}
\leanref{Formalism/SharpRestriction/GrowthBounds/SourceCleaning.lean}{Formalism.MainNotions.SourceCleaning.mStar}
\leanref{Formalism/SharpRestriction/GrowthBounds/SourceCleaning.lean}{Formalism.MainNotions.SourceCleaning.cPrime}

\begin{proof}
As in the proof of Lemma~\ref{pd-edge-numerics}, $|T|\geq1$; since $r\geq2$, $Y\geq r\geq2$,
$c_*\geq8$ and $C\geq2$, the last inequality of (3) of the situation gives $a\geq2$. By (3) also
$k\geq N_1+3r\geq7$, and $c'\geq c_*\geq1$. A spaced run system is valid for $(L_1,N_1)$. We check
(1) to (4).

(1) Let $t\in T$. Then $\deg(t)=a_T(t)-\lambda_0(t)_2$, where $a+c_{pp}\leq a_T(t)\leq2a-c_{pp}$ and
$|\lambda_0(t)_2|\leq c_{pp}$ since $2<k$; so $a\leq\deg(t)\leq2a$. The bound on $E'$ is assumed.

(2) Let $v,v'\in\Lambda_j$ with $\bar v=\bar v'$. Then $|v-v'|\leq f(j,r-1)-f(j-1,0)<G'\leq n'$, so
$v=v'$. Hence $|\Lambda_j|\leq|T|$, and $|\Lambda_j|\leq N_1$ by validity; so
$|\Lambda_j|\leq N_s$. Also $N_s+r\leq N_1+r\leq k-2r\leq K$, since $N_1+2r\leq k-r$.

(3) For $v\in\Lambda_j$, $\operatorname{dist}_j(v)=f(j,r-1)-v\leq f(j,r-1)-f(j-1,0)<G'$.

(4) $Y\geq1$. Since $N_s\leq N_1$, $\log a\geq0$ and $C\geq|C_k|$,
$C_kN_s^2\log a\leq CN_1^2\log a\leq k$. Since $N_s\leq2|T|$,
$128k^2N_sY^2\leq256k^2|T|Y^2\leq p^{L_1-1}$. For the last inequality: by (2) of
Definition~\ref{pd-symbols-run-system}, $f(j+1,r-1)-f(j,r-1)\leq f(j+1,r-1)-f(j,0)<G'$, so
$q_j\leq q$ and $\alpha_{\max}\leq1+rkq$. As $c',q,r\geq1$, this gives $m_*\leq11(N_1+1)c'rkq$, and
$N_1\leq k\leq rk$ gives $E_0\leq30c'rkq^2$; so $E_0+m_*+1\leq c'rkq(30q+11N_1+12)$. With $p\leq q$
(as $G'\geq u>L_2\geq0$), $N_s\leq2|T|$ and $N_sN_1\leq N_1^2\leq r|T|$,
\[
 32pYk^2N_sq^2(E_0+m_*+1)\leq32k^3q^4c'rY\cdot N_s(30q+11N_1+12)\leq3040\,k^3r^2|T|q^5c'Y,
\]
using $N_s(30q+11N_1+12)\leq60|T|q+11r|T|+24|T|\leq95r|T|q$. Since $C\geq3040$, $q^5=p^{5G'}$ and
$c'=c_*+\omega$, the last inequality of (3) of the situation bounds this by $a$.
\end{proof}

\begin{lemma}\label{pd-symbols-edge-basis}
Let $p\geq5$ be a prime, $(W,Q,\sigma,\iota)$ a Hida setting of degree $n'$ at $p$, $k\geq1$ and
$r\geq0$ integers, and $\gamma\in GL_k(Q)$. Let $L''\subset Q^k$ be the $W$-lattice spanned by the
columns of $\gamma$, and $\delta_p$ the diagonal matrix with entries $p^{-1}$ in the first $r$
places and $1$ in the others. For $g\in GL_k(W)$ let $e_\gamma(g)$ be the pair of the $W$-lattices
spanned by the columns of $\gamma\,{}^tg^{-1}$ and of $\gamma\,{}^tg^{-1}\delta_p$, and let the
\emph{plane} $\Pi(g)\subset(W/pW)^k$ of $g$ be the span of the reductions modulo $p$ of the first
$r$ rows of $g^{-1}$. Then:
\begin{enumerate}
\item every edge with source $L''$ is $e_\gamma(g)$ for some $g\in GL_k(W)$;
\item if $r\leq k$, then $e_\gamma(g)$ is an edge with source $L''$ for every $g\in GL_k(W)$;
\item for $g,g'\in GL_k(W)$, $e_\gamma(g)=e_\gamma(g')$ if and only if $\Pi(g)=\Pi(g')$.
\end{enumerate}
\end{lemma}
\leanref{Formalism/SharpRestriction/GrowthBounds/ObjBridge.lean}{Formalism.MainNotions.G.EdgePlane}
\leanref{Formalism/SharpRestriction/GrowthBounds/ScaleEdges/EdgePlane.lean}{Formalism.MainNotions.Pf14.edgePlane}
\leanref{Formalism/SharpRestriction/GrowthBounds/UnramifiedQuadraticSetting.lean}{Formalism.MainNotions.G.IsEdgeBasis}
\leanref{Formalism/SharpRestriction/GrowthBounds/ObjPol.lean}{Formalism.MainNotions.G.ratPlane}

\begin{proof}
Let $g\in GL_k(W)$ and let $b_1,\dots,b_k$ be the columns of $\gamma\,{}^tg^{-1}$. Since
${}^tg^{-1}\in GL_k(W)$, they form a $W$-basis of $L''$, so the first lattice of $e_\gamma(g)$ is
$L''$, and the second is the lattice $L_g$ with basis $p^{-1}b_1,\dots,p^{-1}b_r,b_{r+1},\dots,b_k$.

(2) If $r\leq k$, then $L''\subset L_g$, $pL_g\subset L''$, and $L_g/L''\cong(W/pW)^r$ has $p^{n'r}$
elements; so $e_\gamma(g)$ is an edge.

(1) Let $e$ be an edge with source $L''$ and target $L$. As in the proof of
Lemma~\ref{pd-edge-basis}(1), the elementary divisor theorem gives a basis $b'_1,\dots,b'_k$ of
$L''$ such that $p^{-1}b'_1,\dots,p^{-1}b'_r,b'_{r+1},\dots,b'_k$ is a basis of $L$. The matrix with
columns $b'_i$ is $\gamma U$ for some $U\in GL_k(W)$, since both it and $\gamma$ have columns forming
a basis of $L''$. Put $g={}^tU^{-1}$. Then $\gamma\,{}^tg^{-1}=\gamma U$, and $e=e_\gamma(g)$.

(3) The map $w\mapsto p^{-1}\gamma w$ induces an isomorphism of $(W/pW)^k$ onto $p^{-1}L''/L''$.
Under it, $L_g/L''$ corresponds to the span of the reductions of the first $r$ columns of
${}^tg^{-1}$, which are the transposed first $r$ rows of $g^{-1}$; that is, to $\Pi(g)$. Since
$L''\subset L_g\subset p^{-1}L''$, $L_g$ is determined by $L_g/L''$. Hence $L_g=L_{g'}$ if and only
if $\Pi(g)=\Pi(g')$, and the first lattices of $e_\gamma(g)$ and $e_\gamma(g')$ are both $L''$.
\end{proof}

\begin{definition}\label{pd-symbols-twist}
Let $p\geq5$ be a prime, $(W,Q,\sigma,\iota)$ a Hida setting of degree $n'$ at $p$ with rings
$\mathcal R\subset\mathcal K$ and $\varpi_0$ (Definition~\ref{pd-local-model}(1)), and $1\leq r\leq k$
integers. Let $\mathcal K[X]\supset\mathcal R[X]$ be the polynomial rings in $X_{s,c}$,
$s\in\mathbf Z/n'\mathbf Z$, $1\leq c\leq k$, and $X_s=(X_{s,1},\dots,X_{s,k})$; write $s$ also for
its representative in $\{0,\dots,n'-1\}$. For $g\in GL_k(W)$ and $v\in\mathbf Z$ let
$g_v=\iota\sigma^v(g)$, applied to the entries; it depends on $v$ modulo $n'$.
\begin{enumerate}
\item The \emph{twisted substitution} $\Sigma_g$ is the $\mathcal K$-algebra endomorphism of
  $\mathcal K[X]$ with $X_s\mapsto X_s\,\delta\,g_s^{-1}$ for $s$ even and
  $X_s\mapsto X_s\,\delta'\,{}^tg_s$ for $s$ odd, where $\delta$ and $\delta'$ are the diagonal
  matrices with entries $(1,\dots,1,\varpi_0,\dots,\varpi_0)$ and
  $(\varpi_0,\dots,\varpi_0,1,\dots,1)$, with $r$ entries before the change. For $i\geq0$, a
  polynomial lies in $\varpi_0^i\mathcal R[X]$ if all its coefficients lie in $\varpi_0^i\mathcal R$.
\item Let $c',o\geq0$ be integers. A polynomial $\Phi\in\mathcal K[X]$ is \emph{realizable} at
  $(g,c',o)$ if there is $F\in\mathcal R[X]$, of multidegree $d$ whenever $\Phi$ has multidegree $d$
  (Definition~\ref{pd-symbols-module}(3)),
  such that $\Sigma_g(F)-\varpi_0^o\Phi\in\varpi_0^{o+c'}\mathcal R[X]$ and
  $\Sigma_{g'}(F)\in\varpi_0^{o+c'}\mathcal R[X]$ for every $g'\in GL_k(W)$ whose plane
  (Lemma~\ref{pd-symbols-edge-basis}) differs from that of $g$.
\item Let $T\subset\mathbf Z/n'\mathbf Z$ and $(J,f)$ a run system for $T$ of length $r$, with the
  notation of Definition~\ref{pd-symbols-run-system}; an exponent of $p$ below that would be negative
  is replaced by $0$. For $v\in\mathbf Z$ let $\ell_g(v)=X_{\bar v}\,g_v$, a row of $k$ linear forms.
  For $j\in\mathbf Z$ put $q_j=p^{f(j+1,r-1)-f(j,r-1)}$, $\alpha_0=1+r(k-r)(q_0-1)$ and
  $\alpha_j=r(k-r)(q_j-1)$ for $j\neq0$, and let $\Delta_j(g)$ be the determinant of the $r\times r$ matrix
  whose entry in row $l$, $0\leq l<r$, and column $b$, $1\leq b\leq r$, is
  $\ell_g(f(j,l))_b^{\,p^{f(j,r-1)-f(j,l)}}$. For $v\in\mathbf Z$ let $a_0(v)>a_1(v)>\cdots$ be the
  even integers $y$ with $v-rn'\leq y<v$ and $\bar y\in T$, put $a_l(v)=v$ when there are at most $l$
  of them, and let $\Delta^{\mathrm{top}}_v(g)$ be the determinant of the $r\times r$ matrix with
  entries $\ell_g(a_l(v))_b^{\,p^{v-a_l(v)}}$, $0\leq l<r$, $1\leq b\leq r$. For integers $c',q\geq0$
  put
  \[
   \mathrm{SP}_g=\prod_{j=0}^{J-1}\Delta_j(g)^{c'(\alpha_j+1)}\cdot
   \prod_{t\in T}\Delta^{\mathrm{top}}_t(g)^{(2q+3)c'+kq},
  \]
  with $t\in T$ taken in $\{0,\dots,n'-1\}$. The \emph{symbol multiplier} of $P\in\mathcal R[X]$ for
  $(g,c',q)$ is $P\cdot\Sigma_g(\mathrm{SP}_g)\in\mathcal K[X]$.
\end{enumerate}
\end{definition}
\leanref{Formalism/SharpRestriction/GrowthBounds/SourceCleaning.lean}{Formalism.MainNotions.SourceCleaning.subst}
\leanref{Formalism/SharpRestriction/GrowthBounds/ObjPol.lean}{Formalism.MainNotions.G.liftSubst}
\leanref{Formalism/SharpRestriction/GrowthBounds/ObjPol.lean}{Formalism.MainNotions.G.InWj}
\leanref{Formalism/SharpRestriction/GrowthBounds/SourceCleaning.lean}{Formalism.MainNotions.SourceCleaning.Realisable}
\leanref{Formalism/SharpRestriction/GrowthBounds/SourceCleaning.lean}{Formalism.MainNotions.SourceCleaning.Del}
\leanref{Formalism/SharpRestriction/GrowthBounds/SourceCleaning.lean}{Formalism.MainNotions.SourceCleaning.R0}
\leanref{Formalism/SharpRestriction/GrowthBounds/SourceCleaning.lean}{Formalism.MainNotions.SourceCleaning.DtopProd}
\leanref{Formalism/SharpRestriction/GrowthBounds/ScaleEdges/SymbolRealisableBound.lean}{Formalism.MainNotions.Pf14.SPtrue}
\leanref{Formalism/SharpRestriction/GrowthBounds/ScaleEdges/SymbolRealisableBound.lean}{Formalism.MainNotions.Pf14.symB}

\begin{definition}\label{pd-symbols-realize-keys}
Let $p\geq5$ be a prime, $(W,Q,\sigma,\iota)$ a Hida setting of degree $n'$ at $p$ with rings
$\mathcal R\subset\mathcal K$ and $\varpi_0$, $1\leq r\leq k$ integers, $T\subset\mathbf Z/n'\mathbf Z$,
$q\geq0$ an integer and $g\in GL_k(W)$, with $\Sigma_g$ as in Definition~\ref{pd-symbols-twist} and the
normal degree of Definition~\ref{pd-symbols-module}(3); a variable $X_{s,c}$ is \emph{normal} if $s$ is
even and $c>r$ or $s$ is odd and $c\leq r$. A polynomial \emph{lies on} $T$ if it involves only
variables $X_{s,c}$ with $s\in T$.
\begin{enumerate}
\item A \emph{key} of order $o\geq0$ at $g$ is a pair $(\kappa,\kappa_0)$ of polynomials in
  $\mathcal K[X]$ lying on $T$ such that $\kappa\in\mathcal R[X]$, $o\leq2q$, every monomial of
  $\kappa_0$ has normal degree $o$, $\kappa_0\in\varpi_0^o\mathcal R[X]$,
  $\Sigma_g(\kappa)-\kappa_0\in\varpi_0^{o+1}\mathcal R[X]$, and $\Sigma_{g'}(\kappa)\in
  \varpi_0^o\mathcal R[X]$ for every $g'\in GL_k(W)$.
\item A \emph{key product} of order $s$ is a product $\kappa_0^{(1)}\cdots\kappa_0^{(l)}$, $l\geq0$, of
  the second entries of keys of orders $o_1,\dots,o_l$ with $o_1+\dots+o_l=s$.
\item For $o\geq0$, $\mathrm{Sy}_g(o)$ is the set of the $P\in\mathcal K[X]$ all of whose monomials have
  normal degree $o$ such that $\varpi_0^oP$ is a finite sum of products $AH$ of a key product $A$ of
  order $o$ and a polynomial $H\in\mathcal R[X]$ that lies on $T$ and involves no normal variable. The
  \emph{symbol set} $\mathcal S_g$ is the set of finite sums of elements of the $\mathrm{Sy}_g(o)$,
  $o\geq0$.
\end{enumerate}
\end{definition}
\leanref{Formalism/SharpRestriction/GrowthBounds/SourceCleaning.lean}{Formalism.MainNotions.SourceCleaning.IsKey}
\leanref{Formalism/SharpRestriction/GrowthBounds/WittFrameLocaliserOrders.lean}{Formalism.MainNotions.G.WittFrameLocaliserOrders.bc_KP}
\leanref{Formalism/SharpRestriction/GrowthBounds/WittFrameLocaliserOrders.lean}{Formalism.MainNotions.G.WittFrameLocaliserOrders.bc_IT}
\leanref{Formalism/SharpRestriction/GrowthBounds/WittFrameLocaliserOrders.lean}{Formalism.MainNotions.G.WittFrameLocaliserOrders.bc_agen}
\leanref{Formalism/SharpRestriction/GrowthBounds/PathWeight/GeneratorImageIsSymbol.lean}{Formalism.MainNotions.G.symbolGeneratorData.Sy}
\leanref{Formalism/SharpRestriction/GrowthBounds/PathWeight/GeneratorImageIsSymbol.lean}{Formalism.MainNotions.G.symbolGeneratorData.SyS}

\begin{lemma}\label{pd-symbols-realize-closure}
Let the data and the symbol set $\mathcal S_g$ be as in Definition~\ref{pd-symbols-realize-keys}.
\begin{enumerate}
\item $\mathrm{Sy}_g(o)\cdot\mathrm{Sy}_g(o')\subset\mathrm{Sy}_g(o+o')$ for all $o,o'\geq0$, so
  $\mathcal S_g$ is closed under products.
\item Let $Z\in\mathcal S_g$, let $H_0\in\mathcal K[X]$ have all its monomials of normal degree $0$, and
  let $o\geq0$. Then the sum of the terms of $ZH_0$ of normal degree $o$ is $Z_oH_0$ for some
  $Z_o\in\mathcal S_g$.
\end{enumerate}
\end{lemma}
\leanref{Formalism/SharpRestriction/GrowthBounds/PathWeight/FixedSymbolFactorIntegral.lean}{Formalism.MainNotions.G.symbolGeneratorData.SyS_mul}
\leanref{Formalism/SharpRestriction/GrowthBounds/ScaleEdges/LeviInductionForms.lean}{Formalism.MainNotions.Pf14.SyS_compZ}

\begin{proof}
(1) Let $P\in\mathrm{Sy}_g(o)$ and $P'\in\mathrm{Sy}_g(o')$. The normal degree of a product of monomials
is the sum of theirs, so every monomial of $PP'$ has normal degree $o+o'$. Writing $\varpi_0^oP=\sum AH$
and $\varpi_0^{o'}P'=\sum A'H'$ as in the definition, $\varpi_0^{o+o'}PP'$ is the sum of the
$(AA')(HH')$, where $AA'$ is a key product of order $o+o'$ and $HH'\in\mathcal R[X]$ lies on $T$ and
involves no normal variable. Products of finite sums are finite sums of products.

(2) Write $Z=\sum_iP_i$ with $P_i\in\mathrm{Sy}_g(o_i)$. Every monomial of $P_iH_0$ has normal degree
$o_i$, so the sum of the terms of $ZH_0$ of normal degree $o$ is $Z_oH_0$ with $Z_o$ the sum of the
$P_i$ with $o_i=o$.
\end{proof}

\begin{lemma}\label{pd-symbols-realize-segments}
Let the data of Proposition~\ref{pd-symbols-realize} be given, let $g\in GL_k(W)$, and let
$\mathcal S_g$ be the symbol set of Definition~\ref{pd-symbols-realize-keys} for $T$, $q$ and $g$. Let
$j\in\mathbf Z$, $m\colon\mathbf Z\to\mathbf Z_{\geq0}$, $\theta,o\colon\mathbf Z\to\mathbf R$, and
$\xi$ an element of the segment module $X_j(m,\theta,o)$ over $\mathcal R$
(Definition~\ref{pd-symbols-module}(2)). Then $\phi_j(\xi)\in\mathcal S_g$.
\end{lemma}
\leanref{Formalism/SharpRestriction/GrowthBounds/ScaleEdges/ExistsPolFamilyOfMemGenOfVsDegLe.lean}{Formalism.MainNotions.Pf14.seg_SyS}
\leanref{Formalism/SharpRestriction/GrowthBounds/PathWeight/GeneratorImageIsSymbol.lean}{Formalism.MainNotions.G.symbolGeneratorData.GeneratorSpanInclusion}
\leanref{Formalism/SharpRestriction/GrowthBounds/PathWeight/GeneratorImageIsSymbol.lean}{Formalism.MainNotions.G.symbolGeneratorData.GeneratorImageIsSymbol}
\leanref{Formalism/SharpRestriction/GrowthBounds/ScaleEdges/LeviInductionForms.lean}{Formalism.MainNotions.Pf14.xx_mono}
\leanref{Formalism/SharpRestriction/GrowthBounds/ScaleEdges/ExistsPolFamilyOfMemGenOfVsDegLe.lean}{Formalism.MainNotions.Pf14.xx_span_MGen0}
\leanref{Formalism/SharpRestriction/GrowthBounds/PathWeight/GeneratorImageIsSymbol.lean}{Formalism.MainNotions.G.symbolGeneratorData.MGen0}
\leanref{Formalism/SharpRestriction/GrowthBounds/PathWeight/MdWeightMapDomain.lean}{Formalism.MainNotions.G.symbolGeneratorData.md_comp_gen}
\leanref{Formalism/SharpRestriction/GrowthBounds/PathWeight/MdWeightMapDomain.lean}{Formalism.MainNotions.G.symbolGeneratorData.md_val_comp}
\leanref{Formalism/SharpRestriction/GrowthBounds/PathWeight/MdWeightMapDomain.lean}{Formalism.MainNotions.G.symbolGeneratorData.md_lam_comp}
\leanref{Formalism/SharpRestriction/GrowthBounds/PathWeight/Lemma4OfRemainingStatements.lean}{Formalism.MainNotions.G.symbolGeneratorData.SyS_of_span}
\leanref{Formalism/SharpRestriction/GrowthBounds/PathWeight/GeneratorImageIsSymbol.lean}{Formalism.MainNotions.G.symbolGeneratorData.PGen}
\leanref{Formalism/SharpRestriction/GrowthBounds/PathWeight/CoeUntopDLeAlt.lean}{Formalism.MainNotions.G.symbolGeneratorData.generator_pair_conditions_of_admitted}
\leanref{Formalism/SharpRestriction/GrowthBounds/PathWeight/LetterWeightPos.lean}{Formalism.MainNotions.G.symbolGeneratorData.generator_path_homogeneous_holds}
\leanref{Formalism/SharpRestriction/GrowthBounds/PathWeight/LetterWeightPos.lean}{Formalism.MainNotions.G.symbolGeneratorData.key_column_leading_bound}
\leanref{Formalism/SharpRestriction/GrowthBounds/PathWeight/LetterWeightPos.lean}{Formalism.MainNotions.G.symbolGeneratorData.weighted_component_prod_lowest}

\begin{proof}
We use the notation of Definition~\ref{pd-symbols-module}(1) and (2). For a polynomial $P$ in the
variables $z_{v,c}$ and a function $d\colon\mathbf Z\to\mathbf Z_{\geq0}$ let $P_d$ be the sum of the
terms of $P$ whose degree in the row $z_v$ is $d(v)$ for every $v$; then $P=\sum_dP_d$.

\emph{Step 1: the weights of the parts.} The substitution that defines $P_{S,i}$ replaces each row
$z_v$ with $v\in S$ by $z_v$ times a matrix whose entries do not involve the $z_{v,c}$. So it carries a
polynomial whose degree in every row $z_v$ is $d(v)$ to another such polynomial, and $(P_d)_{S,i}$ is the
sum of the terms of $P_{S,i}$ of these degrees. Hence every exponent of $\eta$ that occurs in
$(P_d)_{S,i}$ occurs in $P_{S,i}$, and
\[
 \operatorname{val}_S(P_d)\geq\operatorname{val}_S(P),\qquad\lambda_S(P_d)\leq\lambda_S(P).
\]

\emph{Step 2: reduction to generators.} The element $\xi$ is an $\mathcal R$-linear combination of
products $AG$ as in Definition~\ref{pd-symbols-module}(2), each with a function $M$ satisfying the
inequalities there with $m$. As $m\geq0$, $M(v)+0\leq M(v)+m(v)$, so each $AG$ satisfies them with $m$
replaced by $0$ and the same $M$. Write $AG=\sum_{d',d''}A_{d'}G_{d''}$. By Step 1, each product
$A_{d'}G_{d''}$ satisfies the inequalities with $m$ replaced by $0$ and the same $M$: for example, for
an even new $v$, $\lambda_{\{y\leq v\}}(G_{d''})\leq\lambda_{\{y\leq v\}}(G)\leq M(v)$ and
$M(v)\leq\operatorname{val}_{\{y<v\}}(A)\leq\operatorname{val}_{\{y<v\}}(A_{d'})$. The polynomial
$A_{d'}$ involves the plane variables $z_{v,c}$ with $v\in\Lambda_j$ only, and $G_{d''}$ the normal
variables $z_{v,c}$ with $v\in\Lambda_j$ new only. So $\xi$ is an $\mathcal R$-linear combination of
\emph{generators}: products $AG$ of a polynomial $A$ in the plane variables $z_{v,c}$, $v\in\Lambda_j$,
and a polynomial $G$ in the normal variables $z_{v,c}$, $v\in\Lambda_j$ new, each with fixed degree
$d'(v)$, respectively $d''(v)$, in every row $z_v$, that satisfy the inequalities of
Definition~\ref{pd-symbols-module}(2) with $m=0$ for some $M$.

\emph{Step 3: linear combinations.} The set $\mathcal S_g$ is closed under sums. An element
$c\in\mathcal R$ lies in $\mathrm{Sy}_g(0)$, as $c=1\cdot c$ with the empty key product, so
$c\mathcal S_g\subset\mathcal S_g$ by Lemma~\ref{pd-symbols-realize-closure}(1). As $\phi_j$ is
$\mathcal R$-linear, $\phi_j(\xi)\in\mathcal S_g$ once $\phi_j(AG)\in\mathcal S_g$ for every generator
$AG$.

\emph{Step 4: the letters of the segment.} We use the notation of
Definitions~\ref{pd-symbols-run-system} and~\ref{pd-symbols-path-generators}. Let
$e_0<e_1<\dots<e_{N-1}$ be the elements of $\Lambda_j$, and give $\{0,\dots,N-1\}$ the lettering in
which $i$ is a vector letter if $e_i$ is even and a covector letter if $e_i$ is odd, with position
$\pi_i=e_i$. Let $\rho$ be the $\mathcal R$-algebra map from the polynomial ring in the $z_{i,b}$,
$0\leq i<N$, $1\leq b\leq k$, to the ring $\mathcal R[z]$ of Definition~\ref{pd-symbols-module} with
$z_{i,b}\mapsto z_{e_i,b}$, and $\rho'$ the $\mathcal K$-algebra map from $\mathcal K[z_{i,b}]$ to
$\mathcal K[X]$ with $z_{i,b}\mapsto X_{\bar e_i,b}$. Put $\zeta=f(j,r-1)$ and
$\omega(i)=p^{\zeta-e_i}=p^{\operatorname{dist}_j(e_i)}$; as $p>1$ and the $e_i$ increase, $\omega$ is
strictly decreasing with positive values. Three facts:
\begin{enumerate}
\item[(i)] $N\geq1$ and $N+r\leq K$: $f(j,r-1)\in\Lambda_j$, $N\leq N_1$ by condition~(5) of validity,
  and $N_1+2r\leq k-r=K$ by (3) of Situation~\ref{pd-edge-situation}.
\item[(ii)] If $e_i$ is not new, then $0,\dots,i$ are vector letters and $e_0,\dots,e_i$ are not new.
  Indeed the elements of $\Lambda_j$ that are not new are even (proof of
  Lemma~\ref{pd-symbols-normal-degree}), and $e_0$ is not new, as $f(j-1,0)\in\Lambda_j$ and
  $f(j-1,0)\leq f(j-1,r-1)$.
\item[(iii)] For a polynomial $A'$ in the $z_{i,b}$ with $b\leq r$ and $0\leq i<N$,
  $\operatorname{val}_{\{y\leq e_i\}}$ and $\lambda_{\{y\leq e_i\}}$ of
  $\rho(\iota_{\mathrm{pl}}(A'))$ (Definition~\ref{pd-symbols-module}(1), with $n=r$) are
  $\operatorname{val}_{\{0,\dots,i\}}(A')$ and $\lambda_{\{0,\dots,i\}}(A')$
  (Definition~\ref{pd-symbols-path-generators}(1)), and those for $\{y<e_i\}$ are those for
  $\{0,\dots,i-1\}$; the same holds for $\rho(\iota_{\mathrm{tn}}(G'))$ with $n=K$. Indeed the
  substitutions change only the rows $z_v$ with $v\in\Lambda_j$ in the set, and $v=e_{i'}$ is even
  exactly when $i'$ is a vector letter.
\end{enumerate}
A \emph{path generator of the segment} is a polynomial
$\rho\big(\iota_{\mathrm{pl}}(\Psi(\nu,I'))\,\iota_{\mathrm{tn}}(\Psi(\kappa,I''))\big)$ over
$\mathcal R$, where $(\nu,\kappa)$ is a path of rank $k$ for some $B$, and $I'$ and $I''$ are standard
fillings of $\nu$ and $\kappa$.

\emph{Step 5: generators are combinations of path generators.} Let $AG$ be a generator of Step 2,
with row degrees $d'$ and $d''$ and function $M$, and write $A=\rho(\iota_{\mathrm{pl}}(A'))$ and
$G=\rho(\iota_{\mathrm{tn}}(G'))$. Then $A'$ and $G'$ have row degrees $d'(e_i)$ and $d''(e_i)$ at the
letter $i$, and $d''(e_i)=0$ when $e_i$ is not new. Put $M_i=M(e_i)$, let $i^+=i+1$ if $i+1<N$ and
$i^+=i$ otherwise, and define
\[
 \theta_A(i)=\begin{cases}\operatorname{val}_{\{0,\dots,i\}}(A')&\text{if finite,}\\
 \max(M_i,M_{i^+})&\text{otherwise,}\end{cases}\qquad
 M_G(i)=\begin{cases}\lambda_{\{0,\dots,i\}}(G')&\text{if finite,}\\
 \min(M_i,M_{i^+})&\text{otherwise,}\end{cases}
\]
$M_A(i)=\lambda_{\{0,\dots,i\}}(A')$ if this is finite and $0$ otherwise, and
$\theta_G(i)=\operatorname{val}_{\{0,\dots,i\}}(G')$ if this is finite and $0$ otherwise. Then $A'$
satisfies the conditions of Lemma~\ref{pd-symbols-path-clauses}(1) with $n=r$, $\theta_A$ and $M_A$,
so it is an $\mathcal R$-linear combination of path generators $\Psi(\nu,I')$ with $\nu$ a path of
rank $r$ for $i\mapsto d'(e_i)$ and $\theta_A(i)\leq\nu_{i+1}(r)$ for every $i$. In the same way $G'$ is
a combination of path generators $\Psi(\kappa,I'')$ with $\kappa$ a path of rank $K$ for
$i\mapsto d''(e_i)$ and $\kappa_{i+1}(1)\leq M_G(i)$ for every $i$. So $AG$ is an $\mathcal R$-linear
combination of the polynomials $\rho(\iota_{\mathrm{pl}}(\Psi(\nu,I'))\,\iota_{\mathrm{tn}}(\Psi(\kappa,I'')))$
for such $\nu$ and $\kappa$, and it suffices that each $(\nu,\kappa)$ is a path of rank $k$. By (i)
and Lemma~\ref{pd-symbols-path-clauses}(2) it remains to check the condition there at each letter $i$.
Recall the inequalities of a generator (Definition~\ref{pd-symbols-module}(2) with $m=0$) and (iii).
\begin{itemize}
\item $e_i$ not new. By (ii) the letters $0,\dots,i$ are vector letters with $d''(e_{i'})=0$. At a
  vector letter $i'$ a path has $\kappa_{i'}\leq\kappa_{i'+1}$ entrywise with entries of equal sum, so
  here $\kappa_{i'+1}=\kappa_{i'}$, and $\kappa_{i+1}=\kappa_0=0$; and $\nu_{i'}\leq\nu_{i'+1}$
  entrywise, so $\nu_i\geq\nu_0=0$. Hence $\kappa_{i+1}(1)=0\leq\nu_i(r)$.
\item $e_i$ new and even. By (ii), $i\geq1$, and $(i-1)^+=i$. The generator gives
  $\lambda_{\{0,\dots,i\}}(G')\leq M_i\leq\operatorname{val}_{\{0,\dots,i-1\}}(A')$. In both cases of
  their definitions, $M_G(i)\leq M_i\leq\theta_A(i-1)$. So
  $\kappa_{i+1}(1)\leq M_G(i)\leq M_i\leq\theta_A(i-1)\leq\nu_i(r)$.
\item $e_i$ new and odd. Again $i\geq1$ and $(i-1)^+=i$. The generator gives
  $\lambda_{\{0,\dots,i-1\}}(G')\leq M_i\leq\operatorname{val}_{\{0,\dots,i\}}(A')$, so
  $M_G(i-1)\leq M_i\leq\theta_A(i)$, and $\kappa_i(1)\leq M_G(i-1)\leq M_i\leq\theta_A(i)\leq
  \nu_{i+1}(r)$.
\end{itemize}
Hence every generator $AG$ is an $\mathcal R$-linear combination of path generators of the segment,
and by Step 3 it suffices that $\phi_j(P)\in\mathcal S_g$ for every path generator $P$ of the segment.

\emph{Step 6: path generators are symbols.} Let
$P=\rho\big(\iota_{\mathrm{pl}}(\Psi(\nu,I'))\,\iota_{\mathrm{tn}}(\Psi(\kappa,I''))\big)$ with
$\mu=(\nu,\kappa)$ a path of rank $k$ for some $B$, and put $o_c=o_\omega(\mu,c)$ and
$o=\sum_{|c|\leq T_\mu}o_c$. If polynomials $f_1,\dots,f_m$ have all their monomials of
$\omega$-degree at least $o'_1,\dots,o'_m$, then all monomials of $f_1\cdots f_m$ have $\omega$-degree
at least $o'_1+\dots+o'_m$, and the sum of its terms of that $\omega$-degree is the product of the sums
of the terms of $\omega$-degree $o'_l$ of the $f_l$, since $\omega$-degrees add under multiplication. By
Lemma~\ref{pd-symbols-path-clauses}(4) there is $I$ such that the sum of the terms of $\omega$-degree $o$
of $\Psi(\mu,I)=\prod_{|c|\leq T_\mu}\mathrm{Col}_{\mu,c}(I(c))$ is
$\iota_{\mathrm{pl}}(\Psi(\nu,I'))\,\iota_{\mathrm{tn}}(\Psi(\kappa,I''))$; by (3) of the same lemma,
all monomials of $\mathrm{Col}_{\mu,c}(I(c))$ have $\omega$-degree at least $o_c$, and the sum of its terms of
$\omega$-degree $o_c$ is $\mathrm{Col}^{\mathrm{pt}}_{\mu,c}(I(c))$. Hence
\[
 \iota_{\mathrm{pl}}(\Psi(\nu,I'))\,\iota_{\mathrm{tn}}(\Psi(\kappa,I''))=\prod_{|c|\leq T_\mu}
 \mathrm{Col}^{\mathrm{pt}}_{\mu,c}(I(c)),
\]
and every monomial of $\mathrm{Col}^{\mathrm{pt}}_{\mu,c}(I(c))$ has $\omega$-degree $o_c$.

The map $\phi_j\circ\rho$ sends $z_{i,b}$ to $X_{\bar e_i,b}^{\,\omega(i)}$. As $n'$ is even, $\bar e_i$
is even exactly when $e_i$ is, so $X_{\bar e_i,b}$ is normal exactly when $z_{i,b}$ has
$\omega$-degree $\omega(i)$, and $\phi_j\circ\rho$ carries a monomial of $\omega$-degree $w$ to one of
normal degree $w$. So every monomial of $\phi_j(P)$ has normal degree $o$, and every monomial of
$\phi_j(\rho(\mathrm{Col}^{\mathrm{pt}}_{\mu,c}(I(c))))$ has normal degree $o_c$.

Fix $c$ with $|c|\leq T_\mu$, put $J=I(c)$,
$\mathfrak k_c=\rho'\big(\mathrm{Col}^g_{\mu,c,\zeta}(J)\big)$ and
$\mathfrak k_{0,c}=\varpi_0^{o_c}\phi_j\big(\rho(\mathrm{Col}^{\mathrm{pt}}_{\mu,c}(J))\big)$. We show that
$(\mathfrak k_c,\mathfrak k_{0,c})$ is a key of order $o_c$ at $g$ (Definition~\ref{pd-symbols-realize-keys}).
Lemma~\ref{pd-symbols-path-clauses}(5) applies: the setting is unramified quadratic, so $\varpi_0$ is a
uniformizer of $\mathcal R$; $p\geq13$; $N\geq1$ and $N+r\leq K$ by (i); the positions $e_i$ are
strictly increasing and at most $\zeta$; and $\omega(i)=p^{\zeta-e_i}$. For $g'\in GL_k(W)$,
$\Sigma_{g'}\circ\rho'=\rho'\circ\Sigma^{\mathrm L}_{g'}$, since $\bar e_i$ is even exactly when $e_i$
is and $g'_{\bar e_i}=g'_{e_i}$; and $\phi_j\circ\rho=\rho'\circ\Pi_\zeta$ on polynomials over
$\mathcal R$, as both send $z_{i,b}$ to $X_{\bar e_i,b}^{\,p^{\zeta-e_i}}$. Renaming variables
preserves $\varpi_0^i\mathcal R[\cdot]$. Now:
\begin{itemize}
\item $\mathfrak k_c$ and $\mathfrak k_{0,c}$ lie on $T$, as $\bar e_i\in T$; $\mathfrak k_c\in\mathcal R[X]$ by
  (5)(a); and $\mathfrak k_{0,c}\in\varpi_0^{o_c}\mathcal R[X]$, with all monomials of normal degree $o_c$.
\item $o_c\leq2p^{\zeta-e_0}\leq2p^{G'}=2q$ by (5)(e), since $e_0\geq f(j-1,0)$ and
  $f(j,r-1)-f(j-1,0)<G'$ by the spacing.
\item $\Sigma_g(\mathfrak k_c)-\mathfrak k_{0,c}=\rho'\big(\Sigma^{\mathrm L}_g(\mathrm{Col}^g_{\mu,c,\zeta}(J))-
  \varpi_0^{o_c}\Pi_\zeta(\mathrm{Col}^{\mathrm{pt}}_{\mu,c}(J))\big)\in\varpi_0^{o_c+1}\mathcal R[X]$ by (5)(d).
\item $\Sigma_{g'}(\mathfrak k_c)=\rho'\big(\Sigma^{\mathrm L}_{g'}(\mathrm{Col}^g_{\mu,c,\zeta}(J))\big)\in
  \varpi_0^{o_c}\mathcal R[X]$ for every $g'\in GL_k(W)$ by (5)(c).
\end{itemize}
As $\phi_j\circ\rho$ is a ring map,
$\varpi_0^o\phi_j(P)=\prod_{|c|\leq T_\mu}\mathfrak k_{0,c}$ is a key product of order $o$. It is the
product of this key product and $H=1$, which lies on $T$ and involves no normal variable, and every
monomial of $\phi_j(P)$ has normal degree $o$. So $\phi_j(P)\in\mathrm{Sy}_g(o)\subset\mathcal S_g$.
% Lean: the generators of Step 2 are MGen0 X.Fr j O K (xx_mono, xx_span_MGen0 with md_comp_gen).
%   Step 4: Mj, eL, LLs (vec = Even, pos = eL), plnm (= rho after plVar/tanVar), rho (= rho'),
%   segment_count_add_rank_le, segment_letter_count_pos, even_of_lower_letter, first_letter_is_lower.
%   Step 5 is generator_pair_conditions_of_admitted (CoeUntopDLeAlt.lean): thA, MA, thG, MG with nxt,
%   lattice_path_zero_prefix and lattice_path_nonneg_prefix for the letters that are not new, hpred for
%   i - 1, and span_mul_image_mem_span; its clauses are Lemma pd-symbols-path-clauses (1) and (2).
%   Step 6 is generator_path_homogeneous_holds (PathWeight/LetterWeightPos.lean): omL (strictAnti, pos),
%   weighted_component_prod_lowest, segment_substitution_weighted_homog with leading_part_normal_degree,
%   key_column_leading_bound (substitution_rename_lift, segment_substitution_eq_rename,
%   depth_filtration_rename, isKeyProduct_of_isKey) and key_polynomial_leading_part (isKeyProduct_mul),
%   with bc_IT_one; its clauses are Lemma pd-symbols-path-clauses (3), (4) and (5).
%   GeneratorSpanInclusion (every element of MGen0 is in the span of PGen, the path generators
%   plnm (PsiPl nu I' * PsiTan ka I'') of PPath with IsSt nu I', IsSt ka I'') is
%   generator_pair_conditions_of_admitted D a15 a16 (StandardMonomialSpanBis, PathPairCorrespondence);
%   GeneratorImageIsSymbol (phSeg of each element of PGen is in Sy o for some o) is
%   generator_path_homogeneous_holds D a17 a18 a20 (admB2, admB3, CoefficientValuationLemma), in
%   PathWeight/LetterWeightPos.lean. Here a15..a20 are adm15 = DefiniteRankOne, adm16 = RotatedSumBound,
%   adm17 = PathWeightComponents, adm18 = ConclusionB3, adm20 = CanonicalTranslateJetIdentity of
%   tableLineDegreeFive.admitted (taken apart in Pf14.sources_realisable_holds), supplied in
%   PlaneCountTheoremWithConstants.second_theorem_closed_everywhere through Joints.* from
%   tool_statements_hold (ExtremalProjectorTools, ColumnFamilies.splitting_package_holds). Step 3 is SyS_of_span.
\end{proof}

\begin{lemma}\label{pd-symbols-realize-points}
Let the data of Proposition~\ref{pd-symbols-realize} be given, let $g\in GL_k(W)$, and let
$\mathrm{Sy}_g(o)$ be as in Definition~\ref{pd-symbols-realize-keys} for $T$, $q$ and $g$. Let $s\in T$,
taken in $\{0,\dots,n'-1\}$, and $1\leq c\leq k$ be such that the variable $X_{s,c}$ is normal. Then
$X_{s,c}\,\Sigma_g\big(\Delta^{\mathrm{top}}_s(g)\big)\in\mathrm{Sy}_g(1)$.
\end{lemma}
\leanref{Formalism/SharpRestriction/GrowthBounds/SourceCleaning.lean}{Formalism.MainNotions.SourceCleaning.PointsLemma}
\leanref{Formalism/SharpRestriction/GrowthBounds/WittFrameLocaliserOrders.lean}{Formalism.MainNotions.G.WittFrameLocaliserOrders.keys_vanish_modulo_uniformiser}
\leanref{Formalism/SharpRestriction/GrowthBounds/WittFrameLocaliserOrders.lean}{Formalism.MainNotions.G.WittFrameLocaliserOrders.pairing_order_of_planes_and_rows}
\leanref{Formalism/SharpRestriction/GrowthBounds/WittFrameLocaliserOrders.lean}{Formalism.MainNotions.G.WittFrameLocaliserOrders.window_holds_even_tops}
\leanref{Formalism/SharpRestriction/GrowthBounds/WittFrameLocaliserOrders.lean}{Formalism.MainNotions.G.WittFrameLocaliserOrders.plane_pairings_of_first_key}
\leanref{Formalism/SharpRestriction/GrowthBounds/WittFrameLocaliserOrders.lean}{Formalism.MainNotions.G.WittFrameLocaliserOrders.explicit_minor_order_positive}
\leanref{Formalism/SharpRestriction/GrowthBounds/WittFrameLocaliserOrders.lean}{Formalism.MainNotions.G.WittFrameLocaliserOrders.bk_vanish1}
\leanref{Formalism/SharpRestriction/GrowthBounds/WittFrameLocaliserOrders.lean}{Formalism.MainNotions.G.WittFrameLocaliserOrders.third_key_residue_eq_zero}
\leanref{Formalism/PolynomialDimension/DetMulEqSumAdjugate.lean}{Formalism.CountingSumsFrames.Top.keyVanish}
\leanref{Formalism/SharpRestriction/GrowthBounds/PathWeight/LeadingDepthOneMemClosure.lean}{Formalism.MainNotions.G.symbolGeneratorData.leading_depth_one_mem_closure}

\begin{proof}
We use the notation of Definitions~\ref{pd-symbols-run-system} and~\ref{pd-symbols-twist}. The setting
is unramified quadratic, so $\varpi_0$ is a uniformizer of $\mathcal R$. Let $\bar{\mathcal R}$ be the
residue field of $\mathcal R$ and $\rho\colon\mathcal R[X]\to\bar{\mathcal R}[X]$ the reduction of the
coefficients; its kernel is $\varpi_0\mathcal R[X]$, and $\bar{\mathcal R}$ has characteristic $p$, as
$p=\varpi_0^2$. For $x\in W$ put $\hat x=\rho(\iota(x))$. For $v\in\mathbf Z$ let
$\ell'_g(v)=X_{\bar v}\,{}^tg_v^{-1}$, a second row of $k$ linear forms.

\emph{Step 1: the rows below $s$.} Put $a_l=a_l(s)$ and $e_l=s-a_l$ for $0\leq l<r$. By condition~(3)
of validity, $f(0,0)$ is even and its residue lies in $T$. Let $y$ be the integer congruent to $f(0,0)$
modulo $n'$ with $s-rn'\leq y<s-rn'+n'$. The $r$ integers $y+in'$, $0\leq i<r$, lie in $[s-rn',s)$, are
even, as $f(0,0)$ and $n'$ are, and have residue in $T$. So $a_0>a_1>\dots>a_{r-1}$ are the $r$
largest even integers in $[s-rn',s)$ with residues in $T$, and $e_l\geq1$. Let $\Xi$ be the $r\times r$ matrix with
entries $\Xi_{lb}=X_{\bar a_l,b}^{\,p^{e_l}}$, $0\leq l<r$, $1\leq b\leq r$. It lies on $T$, and it
involves only variables $X_{t,b}$ with $t$ even and $b\leq r$, none of which is normal.

\emph{Step 2: substitutions.} Let $g'\in GL_k(W)$ and $h=g'^{-1}g$. As $n'$ is even, $\Sigma_{g'}$
replaces $X_{\bar v}$ by $X_{\bar v}\,\delta\,g'^{-1}_v$ for even $v$ and by $X_{\bar v}\,\delta'\,{}^tg'_v$
for odd $v$. Hence $\Sigma_{g'}(\ell_g(v))=X_{\bar v}\,\delta\,h_v$ for even $v$, and
$\Sigma_{g'}(\ell'_g(v))=X_{\bar v}\,\delta'\,{}^t(h^{-1})_v$ for odd $v$. The entries of $h_v$ and of
$(h^{-1})_v$ lie in $\iota(W)\subset\mathcal R$. For $g'=g$ this gives $\Sigma_g(\ell_g(v))=X_{\bar v}\delta$
for even $v$ and $\Sigma_g(\ell'_g(v))=X_{\bar v}\delta'$ for odd $v$; in particular
$\Sigma_g(\Delta^{\mathrm{top}}_s(g))=\det\Xi$.

\emph{Step 3: reduction.} Let $v\in\mathbf Z$ and $m\geq0$ with $m\equiv v$ modulo $n'$. Then
$\sigma^v=\sigma^m$, as $\sigma^{n'}=1$. Since $\sigma(x)\equiv x^p$ modulo $pW$ for $x\in W$, induction
on $m$ gives $\sigma^m(x)\equiv x^{p^m}$ modulo $pW$, and $\rho(\iota(pW))=0$; so
$\rho(\iota\sigma^v(x))=\hat x^{\,p^m}$. The entries of $\delta$ in the places $i>r$ and those of
$\delta'$ in the places $i\leq r$ are $\varpi_0$. By Step 2, for even $v$ and odd $v'$,
\[
 \rho\big(\Sigma_{g'}(\ell_g(v))_b\big)=\sum_{i\leq r}X_{\bar v,i}\,\hat h_{ib}^{\,p^m},\qquad
 \rho\big(\Sigma_{g'}(\ell'_g(v'))_b\big)=\sum_{i>r}X_{\bar v',i}\,\widehat{(h^{-1})}_{bi}^{\,p^{m'}},
\]
with $m'\geq0$ and $m'\equiv v'$ modulo $n'$.

\emph{Step 4: $s$ even.} Then $c>r$. Let $\kappa$ be the determinant of the $(r+1)\times(r+1)$ matrix
whose row $0$ is $(\ell_g(s)_1,\dots,\ell_g(s)_r,\ell_g(s)_c)$ and whose row $l+1$, for $0\leq l<r$, is
$(\ell_g(a_l)_1^{p^{e_l}},\dots,\ell_g(a_l)_r^{p^{e_l}},\ell_g(a_l)_c^{p^{e_l}})$, and put
$\kappa_0=(-1)^r\varpi_0X_{s,c}\det\Xi$. Then $(\kappa,\kappa_0)$ is a key of order $1$ at $g$.
Both lie on $T$, by Step 1. The entries of $g_v$ lie in $\mathcal R$, so $\kappa\in\mathcal R[X]$. Also
$1\leq2q$, $\kappa_0\in\varpi_0\mathcal R[X]$, and every monomial of $\kappa_0$ has normal degree $1$,
as $X_{s,c}$ is normal and $\det\Xi$ involves no normal variable. By Step 2, $\Sigma_g(\kappa)$ is the
determinant of the matrix with row $0$ equal to $(X_{s,1},\dots,X_{s,r},\varpi_0X_{s,c})$ and row $l+1$
equal to $(\Xi_{l1},\dots,\Xi_{lr},(\varpi_0X_{\bar a_l,c})^{p^{e_l}})$. Expanding along the last column,
$\Sigma_g(\kappa)$ is $(-1)^r\varpi_0X_{s,c}\det\Xi=\kappa_0$ plus a sum of products of the
$(\varpi_0X_{\bar a_l,c})^{p^{e_l}}$ with cofactors in $\mathcal R[X]$. As $p^{e_l}\geq2$, this sum lies in
$\varpi_0^2\mathcal R[X]$.

Let $g'\in GL_k(W)$ and $h=g'^{-1}g$. Put $y_0=s$,
$E_0=0$, $y_{l+1}=a_l$ and $E_{l+1}=e_l$ for $l<r$, and $S=s+rn'$ and $N_l=y_l+rn'$ for $l\leq r$. Then
$N_l\geq0$, as $y_l\geq s-rn'$ and $s\geq0$; also $N_l\equiv y_l$ modulo $n'$ and $N_l+E_l=S$. Let
$b_1,\dots,b_{r+1}$ be $1,\dots,r,c$. By Steps 2 and 3, with $m=N_l$, and since $u\mapsto u^{p^{E_l}}$ is
additive on $\bar{\mathcal R}[X]$, the reduction of the entry of $\Sigma_{g'}(\kappa)$ in row $l$ and
column $b_\beta$ is
\[
 \Big(\sum_{i\leq r}X_{\bar y_l,i}\,\hat h_{ib_\beta}^{\,p^{N_l}}\Big)^{p^{E_l}}
 =\sum_{i\leq r}X_{\bar y_l,i}^{\,p^{E_l}}\,\hat h_{ib_\beta}^{\,p^S}.
\]
So the reduced matrix is the product of the $(r+1)\times r$ matrix
$(X_{\bar y_l,i}^{\,p^{E_l}})$ and the $r\times(r+1)$ matrix $(\hat h_{ib_\beta}^{\,p^S})$. It has rank at
most $r$ over the fraction field of $\bar{\mathcal R}[X]$, so $\rho(\Sigma_{g'}(\kappa))=0$, that is,
$\Sigma_{g'}(\kappa)\in\varpi_0\mathcal R[X]$. Hence $(\kappa,\kappa_0)$ is a key of order $1$. Finally
$\varpi_0X_{s,c}\Sigma_g(\Delta^{\mathrm{top}}_s(g))=\kappa_0\cdot(-1)^r$ by Step 2, and every monomial of
$X_{s,c}\Sigma_g(\Delta^{\mathrm{top}}_s(g))$ has normal degree $1$; so it lies in $\mathrm{Sy}_g(1)$.

\emph{Step 5: $s$ odd.} Then $c\leq r$. For $l<r$ put
\[
 \kappa_l=\sum_{b=1}^{k}\ell_g(a_l)_b^{\,p^{e_l}}\,\ell'_g(s)_b,\qquad
 \kappa_{0,l}=\varpi_0\sum_{b\leq r}\Xi_{lb}X_{s,b}.
\]
Then $(\kappa_l,\kappa_{0,l})$ is a key of order $1$ at $g$. Both lie on $T$, and
$\kappa_l\in\mathcal R[X]$, as the entries of $g_v$ and $g_v^{-1}$ lie in $\mathcal R$. Also $1\leq2q$,
$\kappa_{0,l}\in\varpi_0\mathcal R[X]$, and every monomial of $\kappa_{0,l}$ has normal degree $1$, as
$X_{s,b}$ is normal for $b\leq r$. By Step 2,
\[
 \Sigma_g(\kappa_l)=\kappa_{0,l}+\sum_{b>r}(\varpi_0X_{\bar a_l,b})^{p^{e_l}}X_{s,b},
\]
and the sum lies in $\varpi_0^2\mathcal R[X]$, as $p^{e_l}\geq2$. Let $g'\in GL_k(W)$ and $h=g'^{-1}g$, and
put $N=a_l+rn'\geq0$ and $S=N+e_l=s+rn'$. By Steps 2 and 3, with $m=N$ for $a_l$ and $m'=S$ for $s$,
\[
 \rho\big(\Sigma_{g'}(\kappa_l)\big)=\sum_{b=1}^k\Big(\sum_{i\leq r}X_{\bar a_l,i}^{\,p^{e_l}}\,
 \hat h_{ib}^{\,p^S}\Big)\Big(\sum_{i'>r}X_{s,i'}\,\widehat{(h^{-1})}_{bi'}^{\,p^S}\Big)
 =\sum_{i\leq r<i'}X_{\bar a_l,i}^{\,p^{e_l}}X_{s,i'}\Big(\sum_{b=1}^k\hat h_{ib}\,
 \widehat{(h^{-1})}_{bi'}\Big)^{p^S}.
\]
The inner sum is the entry $(i,i')$ of the reduction of $hh^{-1}=1$, which is $0$ for $i\neq i'$. So
$\Sigma_{g'}(\kappa_l)\in\varpi_0\mathcal R[X]$, and $(\kappa_l,\kappa_{0,l})$ is a key of order $1$.
Finally, the adjugate $\operatorname{adj}(\Xi)$ satisfies $\operatorname{adj}(\Xi)\,\Xi=\det(\Xi)\cdot1$;
applied to the column $(X_{s,1},\dots,X_{s,r})$ and read in row $c$, this gives
\[
 \varpi_0X_{s,c}\,\Sigma_g\big(\Delta^{\mathrm{top}}_s(g)\big)=\varpi_0X_{s,c}\det\Xi
 =\sum_{l<r}\kappa_{0,l}\operatorname{adj}(\Xi)_{cl}.
\]
The entries of $\operatorname{adj}(\Xi)$ lie in $\mathcal R[X]$, lie on $T$ and involve no normal
variable, by Step 1. Every monomial of $X_{s,c}\Sigma_g(\Delta^{\mathrm{top}}_s(g))$ has normal degree
$1$, so it lies in $\mathrm{Sy}_g(1)$.
\end{proof}

\begin{lemma}\label{pd-symbols-realize-labels}
Let the data of Proposition~\ref{pd-symbols-realize} be given, let $g\in GL_k(W)$, and let
$\mathrm{Sy}_g(o)$ be as in Definition~\ref{pd-symbols-realize-keys} for $T$, $q$ and $g$. Let $X^\ell$ be
a label (Definition~\ref{pd-symbols-module}(3)) that lies on $T$, of normal degree $o$. Then
$X^\ell\,\Sigma_g\big(\prod_{t\in T}\Delta^{\mathrm{top}}_t(g)^{kq}\big)\in\mathrm{Sy}_g(o)$.
\end{lemma}
\leanref{Formalism/SharpRestriction/GrowthBounds/PathWeight/GeneratorImageIsSymbol.lean}{Formalism.MainNotions.G.symbolGeneratorData.RestMon}
\leanref{Formalism/SharpRestriction/GrowthBounds/PathWeight/LeadingDepthOneMemClosure.lean}{Formalism.MainNotions.G.symbolGeneratorData.restMon}
\leanref{Formalism/SharpRestriction/GrowthBounds/PathWeight/LeadingDepthOneMemClosure.lean}{Formalism.MainNotions.G.symbolGeneratorData.mn_main}
\leanref{Formalism/SharpRestriction/GrowthBounds/PathWeight/LeadingDepthOneMemClosure.lean}{Formalism.MainNotions.G.symbolGeneratorData.mn_label_le}

\begin{proof}
We use the notation of Definitions~\ref{pd-symbols-run-system} and~\ref{pd-symbols-twist}, and put
$D_t=\Sigma_g(\Delta^{\mathrm{top}}_t(g))$ for $t\in T$. By Step 2 of the proof of
Lemma~\ref{pd-symbols-realize-points}, $D_t$ is the determinant of the matrix $\Xi$ formed there for $t$;
by Step 1 there, it lies in $\mathcal R[X]$, lies on $T$ and involves no normal variable.

\emph{The exponents of the label.} Let $t\in T$, let $v$ be its representative in $\{0,\dots,n'-1\}$, and
let $1\leq c\leq k$. As in the proof of Lemma~\ref{pd-symbols-normal-degree}, the intervals
$(f(j-1,r-1),f(j,r-1)]$, $j\in\mathbf Z$, cover $\mathbf Z$; let $j$ be such that
$f(j-1,r-1)<v\leq f(j,r-1)$. Then $v\in\Lambda_j$ is new, since $f(j-1,0)\leq f(j-1,r-1)<v$. As $X^\ell$
is a label, $\ell(t,c)<p^{\operatorname{dist}_j(v)}$. By the spacing,
$\operatorname{dist}_j(v)=f(j,r-1)-v\leq f(j,r-1)-f(j-1,0)<G'$, so $\ell(t,c)\leq p^{G'}=q$. Let $n_t$ be
the sum of the $\ell(t,c)$ over the $c$ for which $X_{t,c}$ is normal; then $n_t\leq kq$.

\emph{The factorization.} As $X^\ell$ lies on $T$, it is the product of
$\prod_{(t,c)}X_{t,c}^{\ell(t,c)}$, over the normal $X_{t,c}$ with $t\in T$, and of the product $N'$ of the
powers of the other variables in it. The normal degree of $X^\ell$ is $o=\sum_{t\in T}n_t$. Hence
\[
 X^\ell\,\Sigma_g\Big(\prod_{t\in T}\Delta^{\mathrm{top}}_t(g)^{kq}\Big)
 =\prod_{(t,c)}\big(X_{t,c}D_t\big)^{\ell(t,c)}\cdot N'\prod_{t\in T}D_t^{\,kq-n_t}.
\]
By Lemma~\ref{pd-symbols-realize-points} each $X_{t,c}D_t$ lies in $\mathrm{Sy}_g(1)$, so by
Lemma~\ref{pd-symbols-realize-closure}(1) the first product lies in $\mathrm{Sy}_g(o)$. The second factor
$H$ lies in $\mathcal R[X]$, lies on $T$ and involves no normal variable. So every monomial of $H$ has
normal degree $0$, and $H=1\cdot H$ with the empty key product $1$ of order $0$; that is,
$H\in\mathrm{Sy}_g(0)$. By Lemma~\ref{pd-symbols-realize-closure}(1) the product lies in
$\mathrm{Sy}_g(o)$.
\end{proof}

\begin{definition}\label{pd-symbols-cleaning-terms}
Let $p\geq5$ be a prime, $(W,Q,\sigma,\iota)$ a Hida setting of degree $n'$ at $p$ with rings
$\mathcal R\subset\mathcal K$ and $\varpi_0$, $1\leq r\leq k$ integers, $T\subset\mathbf Z/n'\mathbf Z$,
$(J,f)$ a run system for $T$ of length $r$, $q,c'\geq0$ integers and $g\in GL_k(W)$, with the notation
of Definitions~\ref{pd-symbols-run-system}, \ref{pd-symbols-twist} and~\ref{pd-symbols-realize-keys};
write $s$ also for the representative of $s\in\mathbf Z/n'\mathbf Z$ in $\{0,\dots,n'-1\}$.
\begin{enumerate}
\item Let $\mathcal I$ be the set of subsets $I\subset\{1,\dots,k\}$ with $r$ elements of which $r-1$
  lie in $\{1,\dots,r\}$; it has $r(k-r)$ elements. For such a set $I$ and $j\in\mathbf Z$ let
  $\Delta_{j,I}(g)$ be the determinant of the $r\times r$ matrix whose entry in row $l$, $0\leq l<r$,
  and column $b\in I$, the columns taken in increasing order, is
  $\ell_g(f(j,l))_b^{\,p^{f(j,r-1)-f(j,l)}}$; thus $\Delta_j(g)=\Delta_{j,\{1,\dots,r\}}(g)$. The
  \emph{localizer} of $g$ is
  \[
   \mathrm{Loc}_g=\Delta_0(g)\prod_{I\in\mathcal I}\Big(\prod_{j=0}^{J-1}\Delta_j(g)^{q_j-1}
   -\prod_{j=0}^{J-1}\Delta_{j,I}(g)^{q_j-1}\Big),
  \]
  and $\Delta^\alpha_g=\prod_{j=0}^{J-1}\Delta_j(g)^{\alpha_j}$. Put
  $\Delta^{\mathrm{top}}(g)=\prod_{t\in T}\Delta^{\mathrm{top}}_t(g)$.
\item A \emph{term} at $g$ is a triple $\tau=\big(i,(\kappa_l,\kappa_{0,l})_{l<n},h\big)$ of an integer
  $i\geq0$, keys $(\kappa_l,\kappa_{0,l})$ at $g$ for $T$ and $q$, of orders $o_l$, and a polynomial
  $h\in\mathcal R[X]$ that lies on $T$ and involves no normal variable. Its \emph{level} is
  $L(\tau)=i+\sum_lo_l$. Let $h^\natural$ be $h$ with each variable $X_{s,c}$ replaced by
  $\ell_g(s)_c$ if $s$ is even and by $\ell'_g(s)_c$ if $s$ is odd, where
  $\ell'_g(s)=X_s\,{}^tg_s^{-1}$. For integers $\mu_1,\mu_2\geq0$ put
  \[
   \mathrm{sym}_\tau(\mu_1,\mu_2)=\varpi_0^i\prod_l\kappa_{0,l}\cdot\Sigma_g(\Delta^\alpha_g)^{c'}
   \prod_{j=0}^{J-1}\Sigma_g(\Delta_j(g))^{\mu_2}\cdot\Sigma_g(\Delta^{\mathrm{top}}(g))^{\mu_1}\cdot h,
  \]
  \[
   \mathrm{raw}_\tau(\mu_1,\mu_2)=\varpi_0^i\prod_l\kappa_l\cdot\mathrm{Loc}_g^{\,c'}
   \prod_{j=0}^{J-1}\Delta_j(g)^{\mu_2}\cdot\Delta^{\mathrm{top}}(g)^{\mu_1}\cdot h^\natural.
  \]
\end{enumerate}
\end{definition}
\leanref{Formalism/SharpRestriction/GrowthBounds/SourceCleaning.lean}{Formalism.MainNotions.SourceCleaning.II}
\leanref{Formalism/SharpRestriction/GrowthBounds/SourceCleaning.lean}{Formalism.MainNotions.SourceCleaning.minorI}
\leanref{Formalism/SharpRestriction/GrowthBounds/SourceCleaning.lean}{Formalism.MainNotions.SourceCleaning.Rloc}
\leanref{Formalism/SharpRestriction/GrowthBounds/SourceCleaning.lean}{Formalism.MainNotions.SourceCleaning.hTrue}
\leanref{Formalism/SharpRestriction/GrowthBounds/SourceCleaning.lean}{Formalism.MainNotions.SourceCleaning.Term}
\leanref{Formalism/SharpRestriction/GrowthBounds/SourceCleaning.lean}{Formalism.MainNotions.SourceCleaning.Term.level}
\leanref{Formalism/SharpRestriction/GrowthBounds/SourceCleaning.lean}{Formalism.MainNotions.SourceCleaning.Term.sym}
\leanref{Formalism/SharpRestriction/GrowthBounds/SourceCleaning.lean}{Formalism.MainNotions.SourceCleaning.Term.raw}

\begin{lemma}\label{pd-symbols-cleaning-off-plane}
Let the data of Proposition~\ref{pd-symbols-realize} be given, and let $g,g'\in GL_k(W)$ have
different planes (Lemma~\ref{pd-symbols-edge-basis}). Let $\mathrm{Loc}_g$ be the localizer of $g$
for the run system of the data (Definition~\ref{pd-symbols-cleaning-terms}). Then
$\Sigma_{g'}(\mathrm{Loc}_g)\in\varpi_0\mathcal R[X]$.
\end{lemma}
\leanref{Formalism/SharpRestriction/GrowthBounds/SourceCleaning.lean}{Formalism.MainNotions.SourceCleaning.Conclusion1}
\leanref{Formalism/SharpRestriction/GrowthBounds/WittFrameLocaliserOrders.lean}{Formalism.MainNotions.G.WittFrameLocaliserOrders.localiser_vanishing_off_plane}
\leanref{Formalism/SharpRestriction/GrowthBounds/WittFrameLocaliserOrders.lean}{Formalism.MainNotions.G.WittFrameLocaliserOrders.localiser_vanishing_of_distinct_planes}
\leanref{Formalism/SharpRestriction/GrowthBounds/WittFrameLocaliserOrders.lean}{Formalism.MainNotions.G.WittFrameLocaliserOrders.localiser_nonzero_off_plane}
\leanref{Formalism/PolynomialDimension/DetMulEqSumAdjugate.lean}{Formalism.CountingSumsFrames.Top.localizerVanishes}
\leanref{Formalism/PolynomialDimension/DetMulEqSumAdjugate.lean}{Formalism.CountingSumsFrames.Top.graphMinor}
\leanref{Formalism/PolynomialDimension/MultinomialDepth/ExistsNonzeroExchangedMinorHolds.lean}{Formalism.CountingSumsFrames.piMinor_map}

\begin{proof}
We use the notation of Definitions~\ref{pd-symbols-run-system}, \ref{pd-symbols-twist}
and~\ref{pd-symbols-cleaning-terms}, and of the proof of Lemma~\ref{pd-symbols-realize-points}: $\rho$
is the reduction $\mathcal R[X]\to\bar{\mathcal R}[X]$ of the coefficients and $\hat x=\rho(\iota(x))$
for $x\in W$. The entries of $g_v$, $g_v^{-1}$, $g'_v$, $g'^{-1}_v$, $\delta$ and $\delta'$ lie in
$\mathcal R$, so $\Sigma_{g'}$ maps $\mathcal R[X]$ into itself and $\Sigma_{g'}(\mathrm{Loc}_g)$ lies in
$\mathcal R[X]$. Let $h=g'^{-1}g\in GL_k(W)$, and let $H$ be the $r\times k$ matrix with entries
$\hat h_{ib}$, $i\leq r$. Its entries lie in the image $F$ of $W/pW$ in $\bar{\mathcal R}$, a field with
$p^{n'}$ elements. For a set $I$ of $r$ columns let $H_I$ be the $r\times r$ matrix of the columns of
$H$ in $I$, in increasing order.

\emph{Step 1: an entry off the plane.} Some entry $\hat h_{ic}$ with $i\leq r<c$ is nonzero. Otherwise,
as $g'^{-1}=hg^{-1}$, row $i\leq r$ of $g'^{-1}$ is $\sum_dh_{id}$ times row $d$ of $g^{-1}$, and the
terms with $d>r$ vanish modulo $p$; so the reductions of the first $r$ rows of $g'^{-1}$ lie in the
plane $\Pi(g)$. The first $r$ rows of an invertible matrix over $W/pW$ are linearly independent, so
$\Pi(g)$ and $\Pi(g')$ both have dimension $r$, and $\Pi(g')=\Pi(g)$, which is false.

\emph{Step 2: the minors modulo $\varpi_0$.} By conditions (1) and (2) of a run system,
$f(j,0)<f(j+1,0)$ and $f(j,0)\leq f(j,l)$, so $f(j,l)\geq f(0,0)$ for $j\geq0$. Let $s_0\leq f(0,0)$ be a
multiple of $n'$, and put $m_{jl}=f(j,l)-s_0$ and $z_j=f(j,r-1)-s_0$, so that $0\leq m_{jl}\leq z_j$
and $m_{jl}\equiv f(j,l)$ modulo $n'$. The $f(j,l)$ are even by condition~(3) of validity; let $s_{jl}$
be the residue of $f(j,l)$. By Steps 2 and 3 of the proof of Lemma~\ref{pd-symbols-realize-points},
with $m=m_{jl}$,
\[
 \rho\big(\Sigma_{g'}(\ell_g(f(j,l))_b)\big)=\sum_{i\leq r}X_{s_{jl},i}\,\hat h_{ib}^{\,p^{m_{jl}}}.
\]
Raising to the power $p^{z_j-m_{jl}}$, which is additive on $\bar{\mathcal R}[X]$, the reduction of
$\Sigma_{g'}$ of the entry in row $l$ and column $b$ of the matrix defining $\Delta_{j,I}(g)$ is
$\sum_{i\leq r}X_{s_{jl},i}^{\,p^{z_j-m_{jl}}}\hat h_{ib}^{\,p^{z_j}}$. So the reduced matrix is the
product of the matrix $T_j=\big(X_{s_{jl},i}^{\,p^{z_j-m_{jl}}}\big)_{l,i}$ and the image of $H_I$ under
the ring map $y\mapsto y^{p^{z_j}}$, and
\[
 \rho\big(\Sigma_{g'}(\Delta_{j,I}(g))\big)=\det(T_j)\,\det(H_I)^{p^{z_j}}.
\]

\emph{Step 3: the exponents.} By condition~(1) of a run system $z_J=z_0+n'$, and $q_j=p^{z_{j+1}-z_j}$,
so $\sum_{j=0}^{J-1}p^{z_j}(q_j-1)=\sum_{j=0}^{J-1}(p^{z_{j+1}}-p^{z_j})=p^{z_0}(p^{n'}-1)$. Every nonzero
$y\in F$ satisfies $y^{p^{n'}-1}=1$. Hence, by Step 2, if $\det(H_I)\neq0$, then
\[
 \rho\Big(\Sigma_{g'}\Big(\prod_{j=0}^{J-1}\Delta_{j,I}(g)^{q_j-1}\Big)\Big)=\prod_{j=0}^{J-1}\det(T_j)^{q_j-1}.
\]

\emph{Step 4: vanishing.} If $\det(H_{\{1,\dots,r\}})=0$, then $\rho(\Sigma_{g'}(\Delta_0(g)))=0$ by
Step 2, so $\rho(\Sigma_{g'}(\mathrm{Loc}_g))=0$. Otherwise $H_0=H_{\{1,\dots,r\}}$ is invertible. Let
$\hat h_{ic}\neq0$ with $i\leq r<c$ be as in Step 1. Then column $c$ of $H$ is nonzero, so
$u=H_0^{-1}(\hat h_{1c},\dots,\hat h_{rc})$ has an entry $u_{i'}\neq0$. By Cramer's rule
$u_{i'}\det(H_0)$ is the determinant of $H_0$ with column $i'$ replaced by column $c$ of $H$. Up to the
order of the columns, this is $H_I$ for $I=(\{1,\dots,r\}\smallsetminus\{i'\})\cup\{c\}$, which lies in
$\mathcal I$; so $\det(H_I)\neq0$. By Step 3 for $\{1,\dots,r\}$ and for $I$, the factor of
$\mathrm{Loc}_g$ at $I$ reduces to
$\prod_j\det(T_j)^{q_j-1}-\prod_j\det(T_j)^{q_j-1}=0$. In both cases
$\rho(\Sigma_{g'}(\mathrm{Loc}_g))=0$. As $\varpi_0$ is a uniformizer of $\mathcal R$, the kernel of
$\rho$ is $\varpi_0\mathcal R[X]$.
\end{proof}

\begin{lemma}\label{pd-symbols-cleaning-expansion}
Let the data of Proposition~\ref{pd-symbols-realize} be given, let $g\in GL_k(W)$, and put
$D_j=\Sigma_g(\Delta_j(g))$ for $0\leq j<J$. Let $\mathrm{Loc}_g$ and $\Delta^\alpha_g$ be as in
Definition~\ref{pd-symbols-cleaning-terms} for the run system of the data. Then
$\Sigma_g(\mathrm{Loc}_g-\Delta^\alpha_g)$ is a finite sum of polynomials
\[
 \varpi_0^e\,X^\nu\prod_{j=0}^{J-1}D_j^{\,y_j}\cdot H,
\]
where $e\geq1$; $X^\nu$ is a monomial of degree at most $e$ in the normal variables $X_{s,c}$ with
$s\in T$; the $y_j\geq0$ are integers with $y_j\geq\alpha_j-e$; and $H\in\mathcal R[X]$ involves only
variables $X_{s,b}$ with $s\in T$ even and $b\leq r$.
\end{lemma}
\leanref{Formalism/SharpRestriction/GrowthBounds/SourceCleaning.lean}{Formalism.MainNotions.SourceCleaning.Conclusion2}
\leanref{Formalism/SharpRestriction/GrowthBounds/WittFrameLocaliserOrders.lean}{Formalism.MainNotions.G.WittFrameLocaliserOrders.localiser_expansion_lemma}
\leanref{Formalism/SharpRestriction/GrowthBounds/WittFrameLocaliserOrders.lean}{Formalism.MainNotions.G.WittFrameLocaliserOrders.complement_minor_order_pos}
\leanref{Formalism/SharpRestriction/GrowthBounds/WittFrameLocaliserOrders.lean}{Formalism.MainNotions.G.WittFrameLocaliserOrders.determinant_powers_lift_into_terms}
\leanref{Formalism/SharpRestriction/GrowthBounds/WittFrameLocaliserOrders.lean}{Formalism.MainNotions.G.WittFrameLocaliserOrders.card_complement_index_set}
\leanref{Formalism/PolynomialDimension/DetMulEqSumAdjugate.lean}{Formalism.CountingSumsFrames.Top.locExpand}

\begin{proof}
Say that a polynomial has \emph{order at least} $E$ if it is a finite sum of polynomials
$\varpi_0^eX^\nu H$ with $e\geq E$, $X^\nu$ a monomial of degree at most $e$ in the normal variables
$X_{s,c}$ with $s\in T$, and $H\in\mathcal R[X]$ in the variables $X_{s,b}$ with $s\in T$ even and
$b\leq r$. Sums and negatives of polynomials of order at least $E$ have order at least $E$, and a
product of polynomials of orders at least $E$ and $E'$ has order at least $E+E'$.

\emph{Step 1: the expansion.} Put $a=\prod_j\Delta_j(g)^{q_j-1}$ and
$b_I=\prod_j\Delta_{j,I}(g)^{q_j-1}$ for $I\in\mathcal I$. Expanding the product over $\mathcal I$, which
has $r(k-r)$ elements, and using the values of the $\alpha_j$ in Definition~\ref{pd-symbols-twist},
\[
 \mathrm{Loc}_g=\sum_{S\subset\mathcal I}(-1)^{|S|}\Delta_0(g)\,a^{r(k-r)-|S|}\prod_{I\in S}b_I
 =\sum_{S\subset\mathcal I}(-1)^{|S|}\prod_{j=0}^{J-1}\Delta_j(g)^{\alpha_j-|S|(q_j-1)}\prod_{I\in S}b_I.
\]
The term of $S=\emptyset$ is $\Delta^\alpha_g$.

\emph{Step 2: the other minors have order at least $1$.} Let $I\in\mathcal I$ and $j\in\mathbf Z$. As in
Step 4 of the proof of Proposition~\ref{pd-symbols-realize}, $\Sigma_g(\ell_g(v))=X_{\bar v}\,\delta$
for even $v$, and the $f(j,l)$ are even with residues $s_{jl}\in T$. So $\Sigma_g$ of the entry in row
$l$ and column $b$ of the matrix defining $\Delta_{j,I}(g)$ is $X_{s_{jl},b}^{\,P_l}$ if $b\leq r$ and
$(\varpi_0X_{s_{jl},b})^{P_l}$ if $b>r$, with $P_l=p^{f(j,r-1)-f(j,l)}\geq1$; here $X_{s_{jl},b}$ is
normal exactly when $b>r$. The first kind has order at least $0$, the second order at least $P_l\geq1$
($e=P_l$, $\nu$ of degree $P_l$). Each term of the expansion of the determinant is $\pm$ a product of
one entry from each column, and $I$ has a column $c>r$; so $\Sigma_g(\Delta_{j,I}(g))$ has order at
least $1$, and $\Sigma_g\big(\prod_{I\in S}b_I\big)$ has order at least
$E_S=|S|\sum_{j=0}^{J-1}(q_j-1)$.

\emph{Step 3: the terms.} Let $S\neq\emptyset$. Since $f(1,r-1)>f(0,r-1)$, $q_0\geq p\geq2$ and
$E_S\geq1$. The exponent $y_j=\alpha_j-|S|(q_j-1)$ is at least $0$ and at least $\alpha_j-E_S$.
Multiplying the polynomials $\varpi_0^eX^\nu H$, $e\geq E_S$, of the expansion of
$\pm\Sigma_g\big(\prod_{I\in S}b_I\big)$ by $\prod_jD_j^{y_j}$ gives polynomials of the stated form, with
$y_j\geq\alpha_j-E_S\geq\alpha_j-e$. Summing over $S\neq\emptyset$ gives
$\Sigma_g(\mathrm{Loc}_g-\Delta^\alpha_g)$ by Step 1.
\end{proof}

\begin{lemma}\label{pd-symbols-cleaning-difference}
Let the data of Proposition~\ref{pd-symbols-realize} be given, let $g\in GL_k(W)$ and $o\geq0$ an
integer, and form the terms at $g$ and the polynomials $\mathrm{sym}_\tau$ and $\mathrm{raw}_\tau$ of
Definition~\ref{pd-symbols-cleaning-terms} for the run system, $q$ and $c'$ of the data. Let $\tau$ be a
term at $g$ of level $L<o+c'$, and let $\mu_1,\mu_2\geq0$ be integers with
$\mu_1\geq(2q+1)(o+c'-1-L)$ and $\mu_2\geq o+c'-1-L$. Then
\[
 \Sigma_g\big(\mathrm{raw}_\tau(\mu_1,\mu_2)\big)-\mathrm{sym}_\tau(\mu_1,\mu_2)
 =\sum_y\mathrm{sym}_{\tau_y}(\mu_{1,y},\mu_{2,y})+\rho
\]
for a finite family of terms $\tau_y$ at $g$ with $L<L(\tau_y)<o+c'$ and integers
$\mu_{1,y}=\mu_1-(2q+1)(L(\tau_y)-L)\geq0$ and $\mu_{2,y}=\mu_2-(L(\tau_y)-L)\geq0$, and some
$\rho\in\varpi_0^{o+c'}\mathcal R[X]$.
\end{lemma}
\leanref{Formalism/SharpRestriction/GrowthBounds/WittFrameLocaliserOrders.lean}{Formalism.MainNotions.G.WittFrameLocaliserOrders.OrderClaimRemaining}
\leanref{Formalism/SharpRestriction/GrowthBounds/WittFrameLocaliserOrders.lean}{Formalism.MainNotions.G.WittFrameLocaliserOrders.order_claim_of_residue_form}
\leanref{Formalism/SharpRestriction/GrowthBounds/WittFrameLocaliserOrders.lean}{Formalism.MainNotions.G.WittFrameLocaliserOrders.DifferenceClaimInCoordinates}
\leanref{Formalism/SharpRestriction/GrowthBounds/WittFrameLocaliserOrders.lean}{Formalism.MainNotions.G.WittFrameLocaliserOrders.bc_claim'}
\leanref{Formalism/SharpRestriction/GrowthBounds/WittFrameLocaliserOrders.lean}{Formalism.MainNotions.G.WittFrameLocaliserOrders.bc_claim_of_tail}
\leanref{Formalism/SharpRestriction/GrowthBounds/WittFrameLocaliserOrders.lean}{Formalism.MainNotions.G.WittFrameLocaliserOrders.bc_gen}
\leanref{Formalism/SharpRestriction/GrowthBounds/WittFrameLocaliserOrders.lean}{Formalism.MainNotions.G.WittFrameLocaliserOrders.bc_gr_tel}
\leanref{Formalism/SharpRestriction/GrowthBounds/WittFrameLocaliserOrders.lean}{Formalism.MainNotions.G.WittFrameLocaliserOrders.bc_restTail}
\leanref{Formalism/SharpRestriction/GrowthBounds/WittFrameLocaliserOrders.lean}{Formalism.MainNotions.G.WittFrameLocaliserOrders.bx_coeff_normDeg}
\leanref{Formalism/SharpRestriction/GrowthBounds/WittFrameLocaliserOrders.lean}{Formalism.MainNotions.G.WittFrameLocaliserOrders.bc_absorb}
\leanref{Formalism/SharpRestriction/GrowthBounds/WittFrameLocaliserOrders.lean}{Formalism.MainNotions.G.WittFrameLocaliserOrders.bc_mono_absorb}
\leanref{Formalism/SharpRestriction/GrowthBounds/WittFrameLocaliserOrders.lean}{Formalism.MainNotions.G.WittFrameLocaliserOrders.subst_true_coordinates}

\begin{proof}
We use the notation of Definitions~\ref{pd-symbols-twist} and~\ref{pd-symbols-cleaning-terms}, and
put $D_j=\Sigma_g(\Delta_j(g))$ and $D^{\mathrm{top}}=\Sigma_g(\Delta^{\mathrm{top}}(g))$. Let
$\tau=(i,(\kappa_l,\kappa_{0,l})_l,h)$, with keys of orders $o_l$.

\emph{Step 0: three facts.} (a) By Step 4 of the proof of Proposition~\ref{pd-symbols-realize}, the
$D_j$, $D^{\mathrm{top}}$ and $\Sigma_g(\Delta^\alpha_g)$ lie in $\mathcal R[X]$, lie on $T$ and involve
no normal variable. (b) Let $P\in\mathcal R[X]$. The substitution $\Sigma_g$ replaces each row $X_s$ by
a row of linear forms in $X_s$ with coefficients in $\mathcal R$, in which the coefficients of the
normal variables lie in $\varpi_0\mathcal R$: these are the $X_{s,c}$ with $c>r$ for $s$ even and
$c\leq r$ for $s$ odd, the places where $\delta$, respectively $\delta'$, has the entry $\varpi_0$. So
$\Sigma_g(P)\in\mathcal R[X]$, the coefficient in $\Sigma_g(P)$ of a monomial of normal degree $N$ lies
in $\varpi_0^N\mathcal R$, and $\Sigma_g(P)$ lies on $T$ if $P$ does. (c) As $\Sigma_g(\ell_g(s))=X_s\delta$
for $s$ even and $\Sigma_g(\ell'_g(s))=X_s\delta'$ for $s$ odd, $\Sigma_g$ sends the replacement of a
variable $X_{s,c}$ that is not normal back to $X_{s,c}$; so $\Sigma_g(h^\natural)=h$, and
\[
 \Sigma_g\big(\mathrm{raw}_\tau(\mu_1,\mu_2)\big)-\mathrm{sym}_\tau(\mu_1,\mu_2)
 =\varpi_0^i\,\Psi\prod_{j=0}^{J-1}D_j^{\,\mu_2}\cdot(D^{\mathrm{top}})^{\mu_1}h,\qquad
 \Psi=\prod_l\Sigma_g(\kappa_l)\,\Sigma_g(\mathrm{Loc}_g)^{c'}-\prod_l\kappa_{0,l}\,
 \Sigma_g(\Delta^\alpha_g)^{c'}.
\]

\emph{Step 1: graded sets.} For an integer $a\geq0$, integers $\beta=(\beta_j)_{j<J}$ and $E\geq0$ let
$G(a,\beta;E)$ be the set of polynomials $\varpi_0^xA\,X^\nu\prod_{j<J}D_j^{\,y_j}\cdot H$ with $A$ a key
product of order $s$ (Definition~\ref{pd-symbols-realize-keys}), $X^\nu$ a monomial in the normal
variables $X_{t,c}$ with $t\in T$, integers $x,y_j\geq0$, and $H\in\mathcal R[X]$ lying on $T$ and
involving no normal variable, for which there is $e\geq E$ with
\[
 x+s=a+e,\qquad|\nu|\leq x,\qquad|\nu|\leq(2q+1)e,\qquad\beta_j\leq y_j+e\ \ (j<J).
\]
Let $\langle G(a,\beta;E)\rangle$ be the set of finite sums of its elements. Multiplying the data shows
$G(a,\beta;E)\,G(a',\beta';E')\subset G(a+a',\beta+\beta';E+E')$; also $1\in G(0,0;0)$ and
$G(a,\beta;1)\subset G(a,\beta;0)$. Hence, if $u_t\in\langle G(a_t,\beta_t;0)\rangle$ and
$v_t\in\langle G(a_t,\beta_t;1)\rangle$ for $t$ in a finite set, then
$\prod_t(u_t+v_t)-\prod_tu_t\in\langle G(\sum_ta_t,\sum_t\beta_t;1)\rangle$, by induction on the number
of factors and the identity $(u_0+v_0)P'-u_0P=(u_0+v_0)(P'-P)+v_0P$.

\emph{Step 2: the pieces.} (i) $\kappa_{0,l}\in G(o_l,0;0)$, with $x=0$, $A=\kappa_{0,l}$, $e=0$, $\nu=0$,
$y=0$ and $H=1$. (ii) $\Sigma_g(\kappa_l)-\kappa_{0,l}\in\langle G(o_l,0;1)\rangle$. Indeed let $P$ be this
difference and $uX^m$ a term of it. By the definition of a key $u\in\varpi_0^{o_l+1}\mathcal R$. Let $N$
be the normal degree of $X^m$. By Step 0(b) the coefficient of $X^m$ in $\Sigma_g(\kappa_l)$ lies in
$\varpi_0^N\mathcal R$, and that in $\kappa_{0,l}$ is $0$ unless $N=o_l$, when it lies in
$\varpi_0^{o_l}\mathcal R$; so $u\in\varpi_0^N\mathcal R$. With $x=\max(o_l+1,N)$ write $u=\varpi_0^xu'$
and $X^m=X^\nu X^{m'}$ with $X^\nu$ the part in the normal variables. As $P$ lies on $T$ (Step 0(b)),
$uX^m\in G(o_l,0;1)$ with this $x$, $A=1$, $e=x-o_l\geq1$, $y=0$ and $H=u'X^{m'}$: indeed $|\nu|=N\leq x$,
and $|\nu|\leq x=o_l+e\leq2q+e\leq(2q+1)e$, as $o_l\leq2q$ and $e\geq1$. (iii)
$\Sigma_g(\Delta^\alpha_g)\in G(0,\alpha;0)$, with $y=\alpha$ and everything else trivial. (iv)
$\Sigma_g(\mathrm{Loc}_g)-\Sigma_g(\Delta^\alpha_g)\in\langle G(0,\alpha;1)\rangle$ by
Lemma~\ref{pd-symbols-cleaning-expansion}: a polynomial $\varpi_0^eX^\nu\prod_jD_j^{y_j}H$ given there
lies in $G(0,\alpha;1)$ with $x=e$, $A=1$, $s=0$, as $|\nu|\leq e\leq(2q+1)e$, $\alpha_j\leq y_j+e$, and
$H$ lies on $T$ and involves no normal variable.

By Step 1, (i) and (ii) give $\prod_l\Sigma_g(\kappa_l)\in\langle G(L-i,0;0)\rangle$ and
$\prod_l\Sigma_g(\kappa_l)-\prod_l\kappa_{0,l}\in\langle G(L-i,0;1)\rangle$; (iii) and (iv) give
$\Sigma_g(\Delta^\alpha_g)^{c'}\in G(0,c'\alpha;0)$ and
$\Sigma_g(\mathrm{Loc}_g)^{c'}-\Sigma_g(\Delta^\alpha_g)^{c'}\in\langle G(0,c'\alpha;1)\rangle$. Since
\[
 \Psi=\prod_l\Sigma_g(\kappa_l)\big(\Sigma_g(\mathrm{Loc}_g)^{c'}-\Sigma_g(\Delta^\alpha_g)^{c'}\big)
 +\Big(\prod_l\Sigma_g(\kappa_l)-\prod_l\kappa_{0,l}\Big)\Sigma_g(\Delta^\alpha_g)^{c'},
\]
$\Psi\in\langle G(L-i,c'\alpha;1)\rangle$.

\emph{Step 3: absorption.} Let $\Phi$ be the set of polynomials of the form
$\sum_y\mathrm{sym}_{\tau_y}(\mu_{1,y},\mu_{2,y})+\rho$ described in the statement. It is closed under
addition and contains $\varpi_0^{o+c'}\mathcal R[X]$. By Steps 0(c) and 2 it suffices to show that
$\Gamma=\varpi_0^i\gamma\prod_jD_j^{\mu_2}\cdot(D^{\mathrm{top}})^{\mu_1}h$ lies in $\Phi$ for
$\gamma=\varpi_0^xA\,X^\nu\prod_jD_j^{y_j}H\in G(L-i,c'\alpha;1)$, with $s$ and $e\geq1$ as in Step 1. A
key product of order $s$ lies in $\varpi_0^s\mathcal R[X]$, and the other factors of $\Gamma$ lie in
$\mathcal R[X]$; so $\Gamma\in\varpi_0^{i+x+s}\mathcal R[X]=\varpi_0^{L+e}\mathcal R[X]$. If
$L+e\geq o+c'$, then $\Gamma\in\Phi$.

Let $L+e<o+c'$. Then $e\leq o+c'-1-L$, so $\mu_1=\mu'_1+(2q+1)e$ and $\mu_2=\mu'_2+e$ with
$\mu'_1,\mu'_2\geq0$; write also $x=|\nu|+x'$, $(2q+1)e=|\nu|+n_2$ and $y_j+e=c'\alpha_j+y'_j$, with
$x',n_2,y'_j\geq0$. For a normal variable $X_{t,c}$ with $t\in T$,
Lemma~\ref{pd-symbols-realize-points} writes $\varpi_0X_{t,c}\Sigma_g(\Delta^{\mathrm{top}}_t(g))$ as a
finite sum of products of a key product of order $1$ and a polynomial in $\mathcal R[X]$ that lies on
$T$ and involves no normal variable; multiplying by the $\Sigma_g(\Delta^{\mathrm{top}}_{t'}(g))$,
$t'\neq t$, which are of the second kind by Step 0(a), gives the same for
$\varpi_0X_{t,c}D^{\mathrm{top}}$. Taking products,
$\varpi_0^{|\nu|}X^\nu(D^{\mathrm{top}})^{|\nu|}$ is a finite sum of products $A_1H_1$ with $A_1$ a key
product of order $|\nu|$ and $H_1$ of the second kind. Then
\[
 \Gamma=\varpi_0^{i+x'}A\cdot\varpi_0^{|\nu|}X^\nu(D^{\mathrm{top}})^{|\nu|}\cdot
 \Sigma_g(\Delta^\alpha_g)^{c'}\prod_jD_j^{\mu'_2}\cdot(D^{\mathrm{top}})^{\mu'_1}\cdot
 \Big(\prod_jD_j^{y'_j}\cdot(D^{\mathrm{top}})^{n_2}H\,h\Big),
\]
as $\Sigma_g(\Delta^\alpha_g)^{c'}=\prod_jD_j^{c'\alpha_j}$. So $\Gamma$ is the sum, over the products
$A_1H_1$, of $\mathrm{sym}_{\tau'}(\mu'_1,\mu'_2)$ for the term $\tau'$ with integer $i+x'$, with the keys
whose second entries make up $A$ and $A_1$, and with $H_1\prod_jD_j^{y'_j}(D^{\mathrm{top}})^{n_2}Hh$,
which lies on $T$ and involves no normal variable. Its level is $i+x'+s+|\nu|=i+x+s=L+e$, so
$L<L(\tau')<o+c'$, $\mu'_1=\mu_1-(2q+1)(L(\tau')-L)$ and $\mu'_2=\mu_2-(L(\tau')-L)$. Hence
$\Gamma\in\Phi$.
\end{proof}

\begin{lemma}\label{pd-symbols-realize-cleaning}
Let the data of Proposition~\ref{pd-symbols-realize} be given, let $g\in GL_k(W)$, let $o\geq0$ be an
integer, and let
$F_{c'}=\prod_{j=0}^{J-1}\Delta_j(g)^{c'(\alpha_j+1)}\prod_{t\in T}\Delta^{\mathrm{top}}_t(g)^{(2q+3)c'}$
(Definition~\ref{pd-symbols-twist}). Let $A$ be a key product of order $o$ at $g$
(Definition~\ref{pd-symbols-realize-keys}, for $T$ and $q$), and let $H\in\mathcal R[X]$ lie on $T$ and
involve no normal variable. Then there is $F\in\mathcal R[X]$ that lies on $T$ such that
$\Sigma_g(F)-A\,\Sigma_g(F_{c'})\,H\in\varpi_0^{o+c'}\mathcal R[X]$ and
$\Sigma_{g'}(F)\in\varpi_0^{o+c'}\mathcal R[X]$ for every $g'\in GL_k(W)$ whose plane
(Lemma~\ref{pd-symbols-edge-basis}) differs from that of $g$.
\end{lemma}
\leanref{Formalism/SharpRestriction/GrowthBounds/WittFrameLocaliserOrders.lean}{Formalism.MainNotions.G.WittFrameLocaliserOrders.clean_symbol_order}
\leanref{Formalism/SharpRestriction/GrowthBounds/WittFrameLocaliserOrders.lean}{Formalism.MainNotions.G.WittFrameLocaliserOrders.bc_sym_of_KP}
\leanref{Formalism/PolynomialDimension/DetMulEqSumAdjugate.lean}{Formalism.CountingSumsFrames.Top.cleanInduction}
\leanref{Formalism/SharpRestriction/GrowthBounds/WittFrameLocaliserOrders.lean}{Formalism.MainNotions.G.WittFrameLocaliserOrders.highOrderTopsSubgroup}
\leanref{Formalism/SharpRestriction/GrowthBounds/WittFrameLocaliserOrders.lean}{Formalism.MainNotions.G.WittFrameLocaliserOrders.symbol_order_ge_level}
\leanref{Formalism/SharpRestriction/GrowthBounds/WittFrameLocaliserOrders.lean}{Formalism.MainNotions.G.WittFrameLocaliserOrders.raw_term_order_ge_off_plane}
\leanref{Formalism/SharpRestriction/GrowthBounds/WittFrameLocaliserOrders.lean}{Formalism.MainNotions.G.WittFrameLocaliserOrders.raw_symbol_order_other_planes}
\leanref{Formalism/SharpRestriction/GrowthBounds/WittFrameLocaliserOrders.lean}{Formalism.MainNotions.G.WittFrameLocaliserOrders.bs_restRawTops}

\begin{proof}
Form the terms at $g$ and $\mathrm{sym}_\tau$, $\mathrm{raw}_\tau$ of
Definition~\ref{pd-symbols-cleaning-terms} for the run system, $q$ and $c'$ of the data. For an integer
$N$ let $\mathcal H(N)$ be the set of $F\in\mathcal R[X]$ that lie on $T$ and satisfy
$\Sigma_{g'}(F)\in\varpi_0^N\mathcal R[X]$ for every $g'\in GL_k(W)$ whose plane differs from that of
$g$. It is an additive group, and $\mathcal H(N)\subset\mathcal H(N')$ for $N'\leq N$.

\emph{Step 1: symbols.} For every term $\tau$ and all $\mu_1,\mu_2$,
$\mathrm{sym}_\tau(\mu_1,\mu_2)\in\varpi_0^{L(\tau)}\mathcal R[X]$: the second entry of a key of order
$o_l$ lies in $\varpi_0^{o_l}\mathcal R[X]$, and the other factors lie in $\mathcal R[X]$ by Step 0(a)
of the proof of Lemma~\ref{pd-symbols-cleaning-difference}.

\emph{Step 2: raw polynomials.} For every term $\tau$ and all $\mu_1,\mu_2$,
$\mathrm{raw}_\tau(\mu_1,\mu_2)\in\mathcal H(L(\tau)+c')$. It lies in $\mathcal R[X]$, since the
$\kappa_l$ do and the entries of $\ell_g(v)$ and $\ell'_g(v)$ are linear forms with coefficients in
$\iota(W)\subset\mathcal R$. It lies on $T$: the $\kappa_l$ and $h$ do, and $\mathrm{Loc}_g$, the
$\Delta_j(g)$ and the $\Delta^{\mathrm{top}}_t(g)$ involve only rows $X_s$ with $s$ the residue of some
$f(j,l)$ or some $a_l(t)$, which lies in $T$ by condition~(3) of validity and by fact~(b) of Step 2 of
the proof of Proposition~\ref{pd-symbols-realize}. Let $g'$ have a plane different from that of $g$.
Then $\Sigma_{g'}$ maps $\mathcal R[X]$ into itself, $\Sigma_{g'}(\kappa_l)\in\varpi_0^{o_l}\mathcal R[X]$
by the definition of a key, and $\Sigma_{g'}(\mathrm{Loc}_g)\in\varpi_0\mathcal R[X]$ by
Lemma~\ref{pd-symbols-cleaning-off-plane}; so
$\Sigma_{g'}(\mathrm{raw}_\tau(\mu_1,\mu_2))\in\varpi_0^{i+\sum_lo_l+c'}\mathcal R[X]$.

\emph{Step 3: induction.} We show: for every term $\tau$ with $L(\tau)\geq o$ and all integers
$\mu_1,\mu_2\geq0$ with $\mu_1\geq(2q+1)(o+c'-1-L(\tau))$ and $\mu_2\geq o+c'-1-L(\tau)$, there is
$F\in\mathcal H(o+c')$ with $\Sigma_g(F)-\mathrm{sym}_\tau(\mu_1,\mu_2)\in\varpi_0^{o+c'}\mathcal R[X]$.
We argue by induction on $o+c'-L(\tau)$. If $L(\tau)\geq o+c'$, then $F=0$ works by Step 1. Otherwise
Lemma~\ref{pd-symbols-cleaning-difference} gives terms $\tau_y$, integers $\mu_{1,y},\mu_{2,y}$ and
$\rho$. For each $y$, $L(\tau_y)>L(\tau)\geq o$, and subtracting $(2q+1)(L(\tau_y)-L(\tau))$, respectively
$L(\tau_y)-L(\tau)$, from the inequalities for $\mu_1$ and $\mu_2$ gives
$\mu_{1,y}\geq(2q+1)(o+c'-1-L(\tau_y))$ and $\mu_{2,y}\geq o+c'-1-L(\tau_y)$. By induction there
are $F_y\in\mathcal H(o+c')$ with $\Sigma_g(F_y)-\mathrm{sym}_{\tau_y}(\mu_{1,y},\mu_{2,y})\in
\varpi_0^{o+c'}\mathcal R[X]$. Put $F=\mathrm{raw}_\tau(\mu_1,\mu_2)-\sum_yF_y$. By Step 2 and
$L(\tau)\geq o$, $F\in\mathcal H(o+c')$, and
\[
 \Sigma_g(F)-\mathrm{sym}_\tau(\mu_1,\mu_2)=\sum_y\big(\mathrm{sym}_{\tau_y}(\mu_{1,y},\mu_{2,y})
 -\Sigma_g(F_y)\big)+\rho\in\varpi_0^{o+c'}\mathcal R[X].
\]

\emph{Step 4: conclusion.} Let $(\kappa_l,\kappa_{0,l})_l$ be keys whose second entries have product $A$
and orders of sum $o$, and $\tau_0=(0,(\kappa_l,\kappa_{0,l})_l,H)$, a term of level $o$. The
conditions of Step 3 hold for $\mu_1=(2q+3)c'$ and $\mu_2=c'$, since $(2q+1)(c'-1)\leq(2q+3)c'$. As
$(\Delta^\alpha_g)^{c'}\prod_j\Delta_j(g)^{c'}=\prod_j\Delta_j(g)^{c'(\alpha_j+1)}$ and
$\Delta^{\mathrm{top}}(g)^{(2q+3)c'}=\prod_{t\in T}\Delta^{\mathrm{top}}_t(g)^{(2q+3)c'}$,
$\mathrm{sym}_{\tau_0}((2q+3)c',c')=A\,\Sigma_g(F_{c'})\,H$. The polynomial $F$ of Step 3 for $\tau_0$
has the stated properties.
\end{proof}

\begin{lemma}\label{pd-symbols-realize-core}
Let the data of Proposition~\ref{pd-symbols-realize} be given, let $g\in GL_k(W)$, and let
$\mathcal S_g$ be the symbol set of Definition~\ref{pd-symbols-realize-keys} for $T$, $q$ and $g$. Let
$o\geq0$ be an integer and
$F_{c'}=\prod_{j=0}^{J-1}\Delta_j(g)^{c'(\alpha_j+1)}\prod_{t\in T}\Delta^{\mathrm{top}}_t(g)^{(2q+3)c'}$
(Definition~\ref{pd-symbols-twist}). Let $Z\in\mathcal S_g$ and $Y'\in\mathcal R[X]$, and suppose that
every monomial of $M'=Z\,\Sigma_g(F_{c'})+p^{c'}Y'$ has normal degree $o$. Then $M'$ is realizable at
$(g,c',o)$.
\end{lemma}
\leanref{Formalism/SharpRestriction/GrowthBounds/PathWeight/FixedSymbolFactorIntegral.lean}{Formalism.MainNotions.G.symbolGeneratorData.core}
\leanref{Formalism/SharpRestriction/GrowthBounds/PathWeight/FixedSymbolFactorIntegral.lean}{Formalism.MainNotions.G.symbolGeneratorData.fixed_symbol_factor_integral}
\leanref{Formalism/SharpRestriction/GrowthBounds/PathWeight/FixedSymbolFactorIntegral.lean}{Formalism.MainNotions.G.symbolGeneratorData.SyS_comp}
\leanref{Formalism/SharpRestriction/GrowthBounds/PathWeight/FixedSymbolFactorIntegral.lean}{Formalism.MainNotions.G.symbolGeneratorData.terms_of_agen}

\begin{proof}
For a polynomial $P$ write $[P]_o$ for the sum of its terms of normal degree $o$. By Step 4 of the proof
of Proposition~\ref{pd-symbols-realize}, $\Sigma_g(\Delta_j(g))$ and $\Sigma_g(\Delta^{\mathrm{top}}_t(g))$
are determinants of matrices with entries $X_{\bar y,b}^{\,p^e}$, $b\leq r$, where $y$ is one of the
$f(j,l)$ or one of the $a_l(t)$, which are even and have residues in $T$ (by validity and by Step 1 of
the proof of Lemma~\ref{pd-symbols-realize-points}). So $\Sigma_g(F_{c'})$ lies in $\mathcal R[X]$, lies
on $T$ and involves no normal variable; every monomial of it has normal degree $0$.

Write $Z=\sum_iP_i$ with $P_i\in\mathrm{Sy}_g(o_i)$. Every monomial of $P_i\Sigma_g(F_{c'})$ has normal
degree $o_i$, so $[Z\Sigma_g(F_{c'})]_o=Z_o\Sigma_g(F_{c'})$, where $Z_o$ is the sum of the $P_i$ with
$o_i=o$. By the definition of $\mathrm{Sy}_g(o)$, $\varpi_0^oZ_o=\sum_mA_mH_m$, a finite sum with key
products $A_m$ of order $o$ and polynomials $H_m\in\mathcal R[X]$ that lie on $T$ and involve no normal
variable. As every monomial of $M'$ has normal degree $o$,
$M'=[M']_o=Z_o\Sigma_g(F_{c'})+p^{c'}[Y']_o$. Hence
\[
 \varpi_0^oM'-\sum_mA_m\,\Sigma_g(F_{c'})\,H_m=\varpi_0^op^{c'}[Y']_o=\varpi_0^{o+2c'}[Y']_o,
\]
which lies in $\varpi_0^{o+c'}\mathcal R[X]$, as $[Y']_o\in\mathcal R[X]$.

By Lemma~\ref{pd-symbols-realize-cleaning}, for each $m$ there is $F_m\in\mathcal R[X]$ with
$\Sigma_g(F_m)-A_m\Sigma_g(F_{c'})H_m\in\varpi_0^{o+c'}\mathcal R[X]$ and
$\Sigma_{g'}(F_m)\in\varpi_0^{o+c'}\mathcal R[X]$ for every $g'\in GL_k(W)$ whose plane differs from that
of $g$. Put $F^0=\sum_mF_m\in\mathcal R[X]$. Then $\Sigma_g(F^0)-\varpi_0^oM'\in\varpi_0^{o+c'}\mathcal
R[X]$, and $\Sigma_{g'}(F^0)\in\varpi_0^{o+c'}\mathcal R[X]$ for every such $g'$.

It remains to arrange the multidegree. If $M'=0$, then $F=0$ realizes $M'$ at $(g,c',o)$. If $M'$ has no
multidegree, $F=F^0$ does. Otherwise $M'\neq0$ has a multidegree $d$, which is unique. Let $F$ be the
sum of the terms of $F^0$ of multidegree $d$; it lies in $\mathcal R[X]$ and has multidegree $d$. For
every $g'\in GL_k(W)$, $\Sigma_{g'}$ substitutes for each $X_s$ a row of linear forms in $X_s$, so it
preserves multidegrees, and the sum of the terms of multidegree $d$ of $\Sigma_{g'}(P)$ is $\Sigma_{g'}$
of that of $P$. Taking the terms of multidegree $d$ preserves $\varpi_0^i\mathcal R[X]$, and
$\varpi_0^oM'$ has multidegree $d$. So $\Sigma_g(F)-\varpi_0^oM'$ and $\Sigma_{g'}(F)$ are the parts of
multidegree $d$ of $\Sigma_g(F^0)-\varpi_0^oM'$ and of $\Sigma_{g'}(F^0)$, and lie in
$\varpi_0^{o+c'}\mathcal R[X]$. Hence $F$ realizes $M'$ at $(g,c',o)$.
\end{proof}

\begin{proposition}\label{pd-symbols-realize}
Let $C_{La},C_k,C_{rk},C_{\deg}$ be real numbers with the property of
Theorem~\ref{pd-symbols-count}. There is $C_2>0$ such that the following holds for every $C\geq C_2$.
Let the data of Situation~\ref{pd-edge-situation} be given with this $C$, let $(J,f)$ be a run system
for $T$ of length $r$ spaced for $(L_1,L_2,N_1,G')$, form $K$, $N_s$, $q$, $c'$, $\deg$, $m_*$ and
$E_0$ as in Lemma~\ref{pd-symbols-space} and the symbol spaces $X_j$ for
$(C_{La},\deg,L_1,N_s,m_*,a,Y)$ over $\mathcal R$ (Definition~\ref{pd-symbols-module}), and let
$d_T(s)=\deg(s)$ for $s\in T$ and $d_T(s)=0$ for $s\notin T$. Then there is
$E'\colon\mathbf Z/n'\mathbf Z\to\mathbf Z_{\geq0}$ with $E'(t)\leq E_0$ and $E'(t)\leq\deg(t)$ for
$t\in T$ with the following property, where $d(t)=\deg(t)-E'(t)$ for $t\in T$ and $d(s)=0$ for
$s\notin T$. Let $g\in GL_k(W)$, $P\in\mathrm{Gen}(X_\bullet,d)$, and $t_1\geq0$ an integer at least
the normal degree of every monomial of $P$. Then there are polynomials $\Phi_o\in\mathcal K[X]$ of
multidegree $d_T$, for $o\geq0$, such that for every $o$ the polynomial $\Sigma_g(\Phi_o)$ has
multidegree $d_T$, lies in $\mathcal R[X]$ and is realizable at $(g,c',o)$
(Definition~\ref{pd-symbols-twist}), and $\sum_{o=0}^{t_1}\Sigma_g(\Phi_o)$ is the symbol multiplier
of $P$ for $(g,c',q)$.
\end{proposition}
\leanref{Formalism/SharpRestriction/GrowthBounds/ScaleEdges/SymbolRealisableBound.lean}{Formalism.MainNotions.Pf14.SourcesRealisable}
\leanref{Formalism/SharpRestriction/GrowthBounds/ScaleEdges/SymbolRealisableBound.lean}{Formalism.MainNotions.Pf14.SymbolRealisableBound}
\leanref{Formalism/SharpRestriction/GrowthBounds/ScaleEdges/SymbolRealisableBound.lean}{Formalism.MainNotions.Pf14.SymbolRealisableAt}
\leanref{Formalism/SharpRestriction/GrowthBounds/ScaleEdges/SymbolRealisableIffGeneralRank.lean}{Formalism.MainNotions.Pf14.sources_realisable_holds}
\leanref{Formalism/SharpRestriction/GrowthBounds/ScaleEdges/ExistsPolFamilyOfMemGenOfVsDegLe.lean}{Formalism.MainNotions.Pf14.exists_pol_family_of_mem_gen_of_vsDeg_le}

\begin{proof}
Put $C_2=30$, let $C\geq30$ and the data be as in the statement, and recall $c'=c_*+\omega$. As in the
proof of Lemma~\ref{pd-edge-numerics}, $|T|\geq1$. Since $G'\geq u>L_2\geq0$, $q=p^{G'}\geq p\geq13$.
By (1) and (3) of the situation, $G'\leq n'$, $N_1\leq k-3r\leq k$ and $Y\geq r\geq2$; and $N_1\geq1$
by Definition~\ref{hf-good-orbit}. We use the
notation of Definitions~\ref{pd-symbols-run-system} and~\ref{pd-symbols-twist}.

\emph{Step 1: the multidegree of $\mathrm{SP}_g$.} The entries of row $l$ of the matrix defining
$\Delta_j(g)$ are linear forms in $X_{\overline{f(j,l)}}$ raised to the power
$p^{f(j,r-1)-f(j,l)}$. So $\Delta_j(g)$ has multidegree $e_j$, where $e_j(s)$ is the sum of the
$p^{f(j,r-1)-f(j,l)}$ over the $l<r$ with $\overline{f(j,l)}=s$. In the same way
$\Delta^{\mathrm{top}}_v(g)$ has multidegree $e^{\mathrm{top}}_v$, where $e^{\mathrm{top}}_v(s)$ is the
sum of the $p^{v-a_l(v)}$ over the $l<r$ with $\overline{a_l(v)}=s$. Hence $\mathrm{SP}_g$ has
multidegree
\[
 E'(s)=\sum_{j=0}^{J-1}c'(\alpha_j+1)e_j(s)+\sum_{t\in T}\big((2q+3)c'+kq\big)e^{\mathrm{top}}_t(s),
\]
which does not depend on $g$. The $f(j,l)$ and the $a_l(t)$, $t\in T$, have residues in $T$, so
$E'(s)=0$ for $s\notin T$.

\emph{Step 2: $E'\leq E_0$.} Two facts on the run system:
\begin{enumerate}
\item[(a)] the $f(j,l)$ with $0\leq j<J$ and $0\leq l<r$ are distinct and lie in
  $[f(0,0),f(0,0)+n')$, by conditions (1) and (2) of a run system, so their residues are distinct;
\item[(b)] for every integer $t$ and $l<r$, $a_l(t)$ is even and $0<t-a_l(t)<G'$. Indeed, let $j_0$
  be the largest $j$ with $f(j,r-1)<t$. By the spacing,
  $t-f(j_0,0)\leq f(j_0+1,r-1)-f(j_0,0)<G'\leq n'$, so the $r$ even points $f(j_0,i)$, $i<r$, lie in
  $[t-rn',t)$. Hence $a_l(t)$ is the $(l+1)$-th largest even point of $T$ in $[t-rn',t)$, and
  $a_l(t)\geq f(j_0,0)>t-G'$.
\end{enumerate}
Let $s\in\mathbf Z/n'\mathbf Z$. By (a), at most one pair $(j_0,l_0)$ with $0\leq j_0<J$ has
$\overline{f(j_0,l_0)}=s$, so the first sum in $E'(s)$ is at most
$c'(\alpha_{j_0}+1)p^{f(j_0,r-1)-f(j_0,l_0)}$. Since
$f(j+1,r-1)-f(j,r-1)\leq f(j+1,r-1)-f(j,0)<G'$, $q_j\leq q$ and $\alpha_j\leq1+rkq$; and
$f(j_0,r-1)-f(j_0,l_0)\leq f(j_0,r-1)-f(j_0-1,0)<G'$. So the first sum is at most
$(c'(1+rkq)+c')q$. By (b), $e^{\mathrm{top}}_t(s)\leq rq$ for $t\in T$, and if
$e^{\mathrm{top}}_t(s)\neq0$ then $s=\overline{a_l(t)}$ for some $l$, so the representative of $t$
lies in the window $W(s,G')$. By condition~(3) of a tiled set this happens for at most $N_1$ elements
$t\in T$, and the second sum is at most $(N_1+1)((2q+3)c'+kq)rq$.

By (3) of the situation and $C\geq30$,
\[
 E_0=15c'(N_1+rk)q^2\leq30c'rkq^2\leq Ck^3r^2|T|q^5c'Y\leq a.
\]
As $c'\geq c_*\geq8$, $E_0\geq p$; so $a\geq p$, $\log_pa\geq1$, and
$k\leq C|T|k^2\log_pa\leq D_p<c_*\leq c'$. With $k\leq c'$, $q\geq3$, $N_1+1\leq2k$ and $rkq\geq2$ we
get $(2q+3)c'+kq\leq4qc'$ and $(c'(1+rkq)+c')q\leq2c'rkq^2$, so
\[
 E'(s)\leq2c'rkq^2+2k\cdot4qc'\cdot rq=10c'rkq^2\leq E_0.
\]

\emph{Step 3: $E'\leq\deg$ on $T$.} For $t\in T$, (2) of the situation gives
$\deg(t)=a_T(t)-\lambda_0(t)_2\geq a\geq E_0\geq E'(t)$. Hence $d(t)+E'(t)=d_T(t)$ for every $t$.

\emph{Step 4: the parts of the symbol multiplier.} Let $g$, $P$ and $t_1$ be as in the statement.
For even $v$ the representative of $\bar v$ is even, as $n'$ is even, and $g_v=g_{\bar v}$; so
$\Sigma_g(\ell_g(v))=\Sigma_g(X_{\bar v})\,g_v=X_{\bar v}\,\delta$, whose first $r$ entries are
$X_{\bar v,1},\dots,X_{\bar v,r}$. The $f(j,l)$ are even, and so are the $a_l(t)$ by (b). Hence
$\Sigma_g(\Delta_j(g))$ and $\Sigma_g(\Delta^{\mathrm{top}}_t(g))$ are the determinants of the
matrices with entries $X_{\overline{f(j,l)},b}^{\,p^{f(j,r-1)-f(j,l)}}$ and
$X_{\overline{a_l(t)},b}^{\,p^{t-a_l(t)}}$, $b\leq r$, and $\Sigma_g(\mathrm{SP}_g)$ is a polynomial with
integer coefficients in variables $X_{s,b}$ with $s$ even and $b\leq r$, none of them normal. The
symbol multiplier $M=P\cdot\Sigma_g(\mathrm{SP}_g)$ therefore lies in $\mathcal R[X]$, has multidegree
$d+E'=d_T$ by Steps 1 and 3, and each of its monomials has normal degree at most $t_1$. For $o\geq0$
let $M_o$ be the sum of the terms of $M$ of normal degree $o$. Then $M_o\in\mathcal R[X]$ has
multidegree $d_T$, and $M=\sum_{o=0}^{t_1}M_o$. Since $\delta$, $\delta'$ and the $g_s$ are invertible,
$\Sigma_g$ is an automorphism of $\mathcal K[X]$, and it preserves multidegrees. Put
$\Phi_o=\Sigma_g^{-1}(M_o)$. Then $\Phi_o$ has multidegree $d_T$, $\Sigma_g(\Phi_o)=M_o$ lies in
$\mathcal R[X]$ and has multidegree $d_T$, and $\sum_{o=0}^{t_1}\Sigma_g(\Phi_o)$ is the symbol
multiplier of $P$. It remains that each $M_o$ is realizable at $(g,c',o)$.

\emph{Step 5: realizability.} Let $\mathcal S_g$ be the symbol set of
Definition~\ref{pd-symbols-realize-keys} for $T$, $q$ and $g$, and write
$P=\prod_{j=0}^{J-1}\phi_j(\xi_j)\cdot X^\ell$ with $\xi_j\in X_j$ and $X^\ell$ a label. If $P=0$, then
$M_o=0$, which is realizable at $(g,c',o)$ with $F=0$. Otherwise $X^\ell$ lies on $T$: some monomial
of $P$ is the product of $X^\ell$ and a monomial of $\prod_j\phi_j(\xi_j)$, and it has degree
$d(s)=0$ in $X_{s,1},\dots,X_{s,k}$ for $s\notin T$. By Lemma~\ref{pd-symbols-realize-labels},
$X^\ell\,\Sigma_g\big(\prod_{t\in T}\Delta^{\mathrm{top}}_t(g)^{kq}\big)\in\mathcal S_g$; by
Lemma~\ref{pd-symbols-realize-segments}, each $\phi_j(\xi_j)\in\mathcal S_g$, as $m_j\geq0$. By
Lemma~\ref{pd-symbols-realize-closure}(1), the product $Z$ of these polynomials lies in $\mathcal S_g$.
With $F_{c'}$ as in Lemma~\ref{pd-symbols-realize-core},
$\mathrm{SP}_g=F_{c'}\prod_{t\in T}\Delta^{\mathrm{top}}_t(g)^{kq}$, so $M=Z\,\Sigma_g(F_{c'})$. As
$\Sigma_g(F_{c'})$ involves no normal variable, by the computation above,
Lemma~\ref{pd-symbols-realize-closure}(2) gives $M_o=Z_o\,\Sigma_g(F_{c'})$ with $Z_o\in\mathcal S_g$.
Lemma~\ref{pd-symbols-realize-core}, with $Y'=0$, shows that $M_o$ is realizable at $(g,c',o)$.
% Lean (exists_pol_family_of_mem_gen_of_vsDeg_le, the claim hreal; Step 5 follows it). What is still open
%   is marked in the proof of Lemma pd-symbols-realize-segments.
% Written out: Lemma pd-symbols-realize-labels (restMon, mn_main) with Lemma pd-symbols-realize-points
%   (PointsLemma = keys_vanish_modulo_uniformiser), Lemma pd-symbols-realize-core (core), and
%   Lemma pd-symbols-realize-cleaning (clean_symbol_order through cleanInduction, with
%   Lemmas pd-symbols-cleaning-off-plane, pd-symbols-cleaning-expansion and pd-symbols-cleaning-difference).
% Done above: eSP (E'), SPtrue_hom, del_bound, d_bound (tops_near, arow_lt_self_and_sub_lt_of_framesOK),
%   num_E, multidegree_bound_le_a, spdeg, eSP_off, sb_SPtrue, fixed_symbol_factor_integral, mn_DkQ_IT,
%   IT_hom, sb_surj, sb_multideg, sum_comp_range.
\end{proof}

\subsection{Flat models and test maps}\label{subsection-pd-flat-models-test-maps}

\begin{definition}\label{pd-symbols-flat-model}
Let $p\geq5$ be a prime, $(W,Q,\sigma,\iota)$ a Hida setting of degree $n'$ at $p$ with rings
$\mathcal R\subset\mathcal K$ and $\varpi_0$, $1\leq r\leq k$ integers,
$d\colon\mathbf Z/n'\mathbf Z\to\mathbf Z_{\geq0}$, and $\lambda^\flat(s)=(d(s),0,\dots,0)\in\mathbf Z^k$.
Let $V^\flat$, $H^\flat$ and $R^\flat$ be the space and lattices of
Definition~\ref{pd-projection-property} formed with $\lambda^\flat$, and $\mathcal K[X]_d$ the
polynomials of multidegree $d$ in $\mathcal K[X]$ (Definition~\ref{pd-symbols-twist}). Let
$\gamma\in GL_k(Q)$, with $L''$ and the pairs $e_\gamma(g)$ as in
Lemma~\ref{pd-symbols-edge-basis}. A \emph{flat model} of $V^\flat$ at $\gamma$ is an
$\mathcal R$-linear bijection $\Theta\colon\mathcal K[X]_d\to V^\flat$ such that:
\begin{enumerate}
\item an element $F\in\mathcal K[X]_d$ lies in $\mathcal R[X]$ if and only if
  $\Theta(F)\in H^\flat(L'')$;
\item for $g\in GL_k(W)$, $F\in\mathcal K[X]_d$ and $i\geq0$, $\Sigma_g(F)\in\varpi_0^i\mathcal R[X]$ if
  and only if $\Theta(F)\in\varpi_0^iR^\flat(e_\gamma(g))$.
\end{enumerate}
\end{definition}
\leanref{Formalism/SharpRestriction/GrowthBounds/ScaleEdges/FlatSrc.lean}{Formalism.MainNotions.Pf14.IsFlatIdent}

\begin{lemma}\label{pd-symbols-flat-letter}
Let $\mathcal R$ be a discrete valuation ring with fraction field $\mathcal K$, let $k\geq2$ and $d\geq0$
be integers, let $x=(x_{a,b})$ be a $k\times k$ matrix of variables and $Y=(Y_1,\dots,Y_k)$ a column
of variables, and let the \emph{parity} be even or odd. For even parity put
$\lambda=\nu=(d,0,\dots,0)\in\mathbf Z^k$; for odd parity put $\lambda=(0,\dots,0,-d)$ and
$\nu=(d,\dots,d,0)$. Let $\nabla_{\mathcal R}(\nu)$ and $\nabla_{\mathcal K}(\nu)$ be the dual Weyl
modules of Definition~\ref{pd-weight-lattices}(1), and $\mathcal R[Y]_d$ and $\mathcal K[Y]_d$ the
homogeneous polynomials of degree $d$ over $\mathcal R$ and $\mathcal K$. For $h\in GL_k(\mathcal K)$
let $\theta_h\colon\mathcal K[Y]\to\mathcal K[x]$ be the $\mathcal K$-algebra map with
$\theta_h(Y_c)=(h^{-1}x)_{c,1}$ for even parity and $\theta_h(Y_c)=(\operatorname{adj}(x)h)_{k,c}$ for
odd parity, where $\operatorname{adj}(x)$ is the adjugate of $x$.
\begin{enumerate}
\item $\theta_1$ maps $\mathcal R[Y]_d$ onto $\nabla_{\mathcal R}(\nu)$ and $\mathcal K[Y]_d$ onto
  $\nabla_{\mathcal K}(\nu)$.
\item Let $\psi\colon\mathcal K[x]\to\mathcal K[Y]$ be the $\mathcal K$-algebra map that substitutes
  for $x$ the matrix with first column $Y$ and all other entries $0$ for even parity, and the matrix
  $\Psi$ with $\Psi_{a,a}=Y_k$ for all $a$, $\Psi_{k,b}=-Y_b$ for $b<k$ and all other entries $0$ for
  odd parity. Then $\psi(\theta_1(Y_c))=qY_c$ for every $c$, where $q=1$ for even parity and
  $q=Y_k^{k-2}$ for odd parity.
\item For $h\in GL_k(\mathcal K)$ and $P\in\mathcal K[Y]_d$, the polynomial $\theta_h(P)(hx)$ obtained
  by substituting $hx$ for $x$ satisfies $\det(h)^{\lambda_k}\theta_h(P)(hx)=\theta_1(P)$. Hence a
  polynomial $f\in\mathcal K[x]$ satisfies $\det(h)^{\lambda_k}f(hx)\in\nabla_{\mathcal R}(\nu)$ if and
  only if $f=\theta_h(P)$ for some $P\in\mathcal R[Y]_d$.
\end{enumerate}
\end{lemma}
\leanref{Formalism/LocalToGlobal/GroupLattices/FockModels/FlatModels/LetterB.lean}{Formalism.L2.G.Model.theta_Hd}
\leanref{Formalism/LocalToGlobal/GroupLattices/FockModels/FlatModels/Letter.lean}{Formalism.L2.G.Model.psi_Nv}
\leanref{Formalism/LocalToGlobal/GroupLattices/FockModels/FlatModels/MemNablaLTwLetter.lean}{Formalism.L2.G.Model.twist_theta}
\leanref{Formalism/LocalToGlobal/GroupLattices/FockModels/FlatModels/MemNablaLTwLetter.lean}{Formalism.L2.G.Model.mem_nablaLTw_letter}

\begin{proof}
(1) For even parity, $\theta_1(Y_c)=x_{c,1}$ has degree $1$ in the first column of $x$ and $0$ in the
others, and the substitutions $x\mapsto x(1+yE_{lj})$ with $l<j$ of
Definition~\ref{pd-weight-lattices}(1) do not change the first column. For odd parity,
$\theta_1(Y_c)=\operatorname{adj}(x)_{k,c}$ is $\pm$ the minor of $x$ on the rows other than $c$ and
the first $k-1$ columns. It has degree $1$ in each of the first $k-1$ columns and $0$ in the last, and
$x\mapsto x(1+yE_{lj})$ adds $y$ times column $l$ to column $j$, which does not change the minor if
$j<k$ and does not touch its columns if $j=k$. So $\theta_1(Y_c)$ lies in the dual Weyl module of
$(1,0,\dots,0)$, respectively $(1,\dots,1,0)$. A product of elements of the dual Weyl modules of two
natural weights lies in the dual Weyl module of their sum, as in the proof of
Lemma~\ref{pd-edge-products}; so $\theta_1$ maps $\mathcal R[Y]_d$ into $\nabla_{\mathcal R}(\nu)$,
and $\mathcal K[Y]_d$ into $\nabla_{\mathcal K}(\nu)$. Conversely, by the standard basis theorem
\cite{DEP80}, over $\mathcal R$ and over $\mathcal K$ (Lemma~\ref{is-standard-basis}(1)),
$\nabla_{\mathcal R}(\nu)$ and $\nabla_{\mathcal K}(\nu)$ are spanned by the bideterminants of shape
$\nu$: products of minors of $x$, the $i$-th on the first $s_i$ columns, with $\#\{i:s_i\geq
t\}=\nu_t$ for every $t$. For $\nu=(d,0,\dots,0)$, $d$ of the $s_i$ are $1$ and the others $0$, so a
bideterminant is a product of $d$ entries of the first column, the image of a monomial of degree $d$.
For $\nu=(d,\dots,d,0)$, $d$ of the $s_i$ are $k-1$ and the others $0$, so a bideterminant is a
product of $d$ minors on the first $k-1$ columns, each of them $\pm\theta_1(Y_c)$ for some $c$; it is
the image of $\pm$ a monomial of degree $d$.

(2) For even parity $\psi(\theta_1(Y_c))=\psi(x_{c,1})=Y_c$. For odd parity
$\psi(\theta_1(Y_c))=\operatorname{adj}(\Psi)_{k,c}$. The matrix $\Psi$ is lower triangular with
diagonal entries $Y_k$, so $\operatorname{adj}(\Psi)\Psi=Y_k^k\cdot1$. In row $k$ this reads
$\operatorname{adj}(\Psi)_{k,k}Y_k=Y_k^k$ in column $k$, and
$\operatorname{adj}(\Psi)_{k,b}Y_k-\operatorname{adj}(\Psi)_{k,k}Y_b=0$ in column $b<k$. As
$\mathcal K[Y]$ is a domain and $k\geq2$, $\operatorname{adj}(\Psi)_{k,k}=Y_k^{k-1}$ and
$\operatorname{adj}(\Psi)_{k,b}=Y_k^{k-2}Y_b$.

(3) For even parity $\lambda_k=0$, as $k\geq2$, and $\theta_h(P)(hx)$ is $P$ evaluated at the first
column of $h^{-1}hx=x$, which is $\theta_1(P)$. For odd parity $\lambda_k=-d$; since
$\operatorname{adj}(hx)h=\operatorname{adj}(x)\operatorname{adj}(h)h=\det(h)\operatorname{adj}(x)$
and $P$ is homogeneous of degree $d$, $\theta_h(P)(hx)=\det(h)^d\theta_1(P)$. This gives the identity.
The map $f\mapsto\det(h)^{\lambda_k}f(hx)$ is a bijection of $\mathcal K[x]$ onto itself, with inverse
$F\mapsto\det(h)^{-\lambda_k}F(h^{-1}x)$, so the second claim follows from the identity and (1) over
$\mathcal R$.
\end{proof}

\begin{lemma}\label{pd-symbols-flat-exists}
Let $p\geq5$ be a prime, $(W,Q,\sigma,\iota)$ a Hida setting of degree $n'$ at $p$ with rings
$\mathcal R\subset\mathcal K$ and $\varpi_0$, $1\leq r<k$ integers,
$d\colon\mathbf Z/n'\mathbf Z\to\mathbf Z_{\geq0}$, and $V^\flat$ and $\mathcal K[X]_d$ as in
Definition~\ref{pd-symbols-flat-model}; write $s$ also for its representative in
$\{0,\dots,n'-1\}$. Let $\gamma\in GL_k(Q)$ and $\gamma_s=\iota\sigma^s(\gamma)$, applied to the
entries. Let $x_s$ be the $k\times k$ matrices of variables of Definition~\ref{pd-projection-property}
and $\operatorname{adj}(x_s)$ the adjugate of $x_s$. Let $\Theta_\gamma$ be the restriction to
$\mathcal K[X]_d$ of the $\mathcal K$-algebra map $\mathcal K[X]\to\mathcal K[x_{s,a,b}]$ that sends
$X_{s,c}$ to the entry $(\gamma_s^{-1}x_s)_{c,1}$ if $s$ is even and to the entry
$(\operatorname{adj}(x_s)\gamma_s)_{k,c}$ if $s$ is odd. Then $\Theta_\gamma$ is a flat model of
$V^\flat$ at $\gamma$.
\end{lemma}
\leanref{Formalism/SharpRestriction/GrowthBounds/ScaleEdges/SourcePrimitiveIntegrality.lean}{Formalism.MainNotions.Pf14.theta_exists}
\leanref{Formalism/SharpRestriction/GrowthBounds/ScaleEdges/SourcePrimitiveIntegrality.lean}{Formalism.MainNotions.Pf14.IsTrueTheta}
\leanref{Formalism/LocalToGlobal/GroupLattices/FockModels/FlatModels/LaTwTranspose.lean}{Formalism.L2.G.Model.thF}
\leanref{Formalism/LocalToGlobal/GroupLattices/FockModels/FlatModels/LaTwTranspose.lean}{Formalism.L2.G.Model.gen_equiv}
\leanref{Formalism/LocalToGlobal/GroupLattices/FockModels/FlatModels/LaTwTranspose.lean}{Formalism.L2.G.Model.gen_A}
\leanref{Formalism/LocalToGlobal/GroupLattices/FockModels/FlatModels/LaTwTranspose.lean}{Formalism.L2.G.Model.gen_R}

\begin{proof}
For a family $h=(h_s)$ of matrices in $GL_k(\mathcal K)$, one for each position $s$, let
$\Theta_h\colon\mathcal K[X]\to\mathcal K[x_{s,a,b}]$ be the $\mathcal K$-algebra map that is, on the
variables $X_s=(X_{s,1},\dots,X_{s,k})$, the map $\theta_{h_s}$ of
Lemma~\ref{pd-symbols-flat-letter} with values in the variables $x_s$, for the parity of $s$. Then
$\Theta_\gamma$ is $\Theta_h$ for $h_s=\gamma_s$. For $F\in\mathcal K[X]$ and matrices $N=(N_s)$
write $F^N$ for $F$ with the column $X_s$ replaced by $N_sX_s$ for every $s$; then
$(F^N)^M=F^{NM}$, with $(NM)_s=N_sM_s$, and $F^N$ has multidegree $d$ when $F$ has. Put
$a_s(h)=h_s^{-1}$ for $s$ even and $a_s(h)={}^th_s$ for $s$ odd. Then
$\theta_{h_s}(X_{s,c})=\theta_1\big((a_s(h)X_s)_c\big)$, so $\Theta_h(F)=\Theta_1(F^{a(h)})$. Write
$\mathcal R[X]_d=\mathcal R[X]\cap\mathcal K[X]_d$.

\emph{Step 1: products.} For each $s$ let $M_s$ be the image under $\theta_{h_s}$ of the homogeneous
polynomials of degree $d(s)$ in $X_s$ over $\mathcal R$. Then the $\mathcal R$-span of the products
$\prod_sf_s$ with $f_s\in M_s$ is $\Theta_h(\mathcal R[X]_d)$. Indeed
$\prod_s\theta_{h_s}(P_s)=\Theta_h(\prod_sP_s)$ with $\prod_sP_s\in\mathcal R[X]_d$; conversely
$\mathcal R[X]_d$ is spanned by the monomials of multidegree $d$, and each is a product over $s$ of
monomials of degree $d(s)$ in $X_s$. The same holds over $\mathcal K$.

\emph{Step 2: the weights.} For $\lambda^\flat(s)=(d(s),0,\dots,0)$, the weight $\lambda(s)$ of
Definition~\ref{pd-local-model} is $(d(s),0,\dots,0)$ for $s$ even and $(0,\dots,0,-d(s))$ for $s$
odd, with natural weight $\nu(s)=(d(s),0,\dots,0)$, respectively $(d(s),\dots,d(s),0)$. These are the
$\lambda$ and $\nu$ of Lemma~\ref{pd-symbols-flat-letter} for $d(s)$ and the parity of $s$. As
$1\leq r<k$, the integer of Definition~\ref{pd-projection-property}(1) for $\lambda^\flat$ is
$A=\sum_s\sum_{i\leq r}\lambda(s)_i=\sum_{s\text{ even}}d(s)$.

\emph{Step 3: a bijection.} By Lemma~\ref{pd-symbols-flat-letter}(1) over $\mathcal K$ and Step 1 for
$h=1$ over $\mathcal K$, $V^\flat=\Theta_1(\mathcal K[X]_d)$. Since $F\mapsto F^{a(\gamma)}$ is a
bijection of $\mathcal K[X]_d$ onto itself, $\Theta_\gamma(\mathcal K[X]_d)=V^\flat$. Let
$\Psi\colon\mathcal K[x_{s,a,b}]\to\mathcal K[X]$ be the $\mathcal K$-algebra map that is the map
$\psi$ of Lemma~\ref{pd-symbols-flat-letter}(2) at each position, from $x_s$ to $X_s$. By that lemma
$\Psi\Theta_1(X_{s,c})=q_sX_{s,c}$ with $q_s\neq0$, so $\Psi\Theta_1(F)=\prod_sq_s^{d(s)}\cdot F$ for
$F\in\mathcal K[X]_d$. As $\mathcal K[X]$ is a domain, $\Theta_1$ is injective on $\mathcal K[X]_d$,
and so is $\Theta_h$ for every $h$. So $\Theta_\gamma$ is a $\mathcal K$-linear bijection of
$\mathcal K[X]_d$ onto $V^\flat$.

\emph{Step 4: property (1).} The columns of $\gamma$ form a $W$-basis of $L''$, and $\iota\sigma^s$
maps $W$ into $\mathcal R$, so the columns of $\gamma_s$ form an $\mathcal R$-basis of the
$\mathcal R$-span of $\iota\sigma^s(L'')$. By Definition~\ref{pd-projection-property}(1),
$H^\flat(L'')$ is the $\mathcal R$-span of the products $\prod_sf_s(x_s)$ with
$\det(\gamma_s)^{\lambda(s)_k}f_s(\gamma_sx_s)\in\nabla_{\mathcal R}(\nu(s))$; the condition for this
basis gives it for every other basis $\gamma_su$, $u\in GL_k(\mathcal R)$, because
$\nabla_{\mathcal R}(\nu(s))$ is stable under $x\mapsto ux$ and $\det(u)$ is a unit. By
Lemma~\ref{pd-symbols-flat-letter}(3) these $f_s$ are the $\theta_{\gamma_s}(P_s)$ with $P_s$
homogeneous of degree $d(s)$ over $\mathcal R$, so $H^\flat(L'')=\Theta_\gamma(\mathcal R[X]_d)$ by
Step 1. Hence $\Theta_\gamma(F)\in H^\flat(L'')$ for $F\in\mathcal R[X]_d$. Conversely, if
$\Theta_\gamma(F)=\Theta_\gamma(G)$ with $G\in\mathcal R[X]_d$, then $F=G$ by Step 3.

\emph{Step 5: the lattice of an edge.} Let $g\in GL_k(W)$ and $e=e_\gamma(g)$, and put
$\gamma'=\gamma\,{}^tg^{-1}$, so that $\iota\sigma^s(\gamma')=\gamma_s\,{}^tg_s^{-1}$ with $g_s$ as
in Definition~\ref{pd-symbols-twist}. Let $D$ be the diagonal matrix with entries $\varpi_0^{-1}$ in the
first $r$ places and $1$ in the others, and $h_s=\gamma_s\,{}^tg_s^{-1}D$. By
Lemma~\ref{pd-symbols-edge-basis} the columns $b_i$ of $\gamma'$ form a basis of the source $L''$ of
$e$, and $p^{-1}b_1,\dots,p^{-1}b_r,b_{r+1},\dots,b_k$ a basis of its target $L$. The columns of
$h_s$ are $\varpi_0\,\iota\sigma^s(p^{-1}b_i)$ for $i\leq r$ and $\iota\sigma^s(b_i)$ for $i>r$, so
they form an $\mathcal R$-basis of the $\mathcal R$-span of $\iota\sigma^s(L'')$ and
$\varpi_0\,\iota\sigma^s(L)$, from which $H^\flat(e)$ is formed. As in Step 4,
$H^\flat(e)=\Theta_h(\mathcal R[X]_d)$, and $R^\flat(e)=\varpi_0^{-A}\Theta_h(\mathcal R[X]_d)$.

\emph{Step 6: property (2).} Let $N_s={}^tg_s^{-1}\delta$ for $s$ even and $N_s=g_s\delta'$ for $s$
odd, with $\delta$ and $\delta'$ as in Definition~\ref{pd-symbols-twist}(1). Since $\Sigma_g$ replaces
the row $X_s$ by $X_s\,{}^tN_s$, $\Sigma_g(F)=F^N$. Let $F\in\mathcal K[X]_d$ and
$G=F^{a(\gamma)a(h)^{-1}}\in\mathcal K[X]_d$. Then
$\Theta_h(G)=\Theta_1(G^{a(h)})=\Theta_1(F^{a(\gamma)})=\Theta_\gamma(F)$. Moreover
$a_s(\gamma)a_s(h)^{-1}=\gamma_s^{-1}h_s={}^tg_s^{-1}D$ for $s$ even, and
$a_s(\gamma)a_s(h)^{-1}={}^t\gamma_s\,{}^th_s^{-1}={}^t(h_s^{-1}\gamma_s)=g_sD^{-1}$ for $s$ odd. As
$\delta=\varpi_0D$ and $\delta'=D^{-1}$, $N_s$ is $\varpi_0a_s(\gamma)a_s(h)^{-1}$ for $s$ even and
$a_s(\gamma)a_s(h)^{-1}$ for $s$ odd. Since $F$ has degree $d(s)$ in $X_s$, Step 2 gives
$\Sigma_g(F)=\varpi_0^AG$. Hence $\Sigma_g(F)\in\varpi_0^i\mathcal R[X]$ if and only if
$G\in\varpi_0^{i-A}\mathcal R[X]_d$. As $\Theta_h$ is injective on $\mathcal K[X]_d$, this holds if and
only if $\Theta_\gamma(F)=\Theta_h(G)$ lies in $\varpi_0^{i-A}\Theta_h(\mathcal R[X]_d)$, which is
$\varpi_0^iR^\flat(e)$ by Step 5.
\end{proof}

\begin{lemma}\label{pd-symbols-primitive-basis}
Let $\mathcal R$ be a discrete valuation ring with fraction field $\mathcal K$, $n\geq1$ an integer and
$x=(x_{a,t})$ an $n\times n$ matrix of variables. Order the monomials
$x^e=\prod_{a,t}x_{a,t}^{e(a,t)}$ \emph{column-lexicographically}: $x^e>x^{e'}$ if $e(a,t)>e'(a,t)$ at
the first entry $(a,t)$, in the order $(1,1),(2,1),\dots,(n,1),(1,2),\dots$, at which $e$ and $e'$
differ; the \emph{leading monomial} of a nonzero polynomial is the largest monomial occurring in it.
For integers $\nu_1\geq\dots\geq\nu_n\geq0$ let $E_n(\nu)$ be the set of $n\times n$ matrices $e$ of
nonnegative integers with $\sum_ae(a,t)=\nu_t$ for every $t$ and
\[
 \sum_{b\leq a}e(b,t+1)\leq\sum_{b<a}e(b,t)\qquad(1\leq a\leq n,\ 1\leq t<n),
\]
and let $\dim_n(\nu)$ be as in Lemma~\ref{pd-edge-dimension-count}.
\begin{enumerate}
\item $\nabla_{\mathcal R}(\nu)$ (Definition~\ref{pd-weight-lattices}(1)) has an $\mathcal R$-basis
  $(b(e))_{e\in E_n(\nu)}$ of bideterminants of shape $\nu$ such that the leading monomial of $b(e)$ is
  $x^e$, with coefficient $1$.
\item $|E_n(\nu)|=\dim_n(\nu)$. If $n\geq2$, the matrices $e\in E_n(\nu)$ with $e(1,1)=\nu_1$ are
  $\dim_{n-1}(\nu_2,\dots,\nu_n)$ in number, and if moreover $\nu_n=0$, those whose first row is zero
  are $\dim_{n-1}(\nu_1,\dots,\nu_{n-1})$ in number.
\end{enumerate}
\end{lemma}
\leanref{Formalism/GradedWeylDivisibility/FinalD.lean}{Formalism.GradedWeylDivisibility.sSBasis}
\leanref{Formalism/GradedWeylDivisibility/FinalC.lean}{Formalism.GradedWeylDivisibility.colLexInitMinor}
\leanref{Formalism/SharpRestriction/GrowthBounds/DualFockLinear/TestSideCounting.lean}{Formalism.MainNotions.TestSideCounting.cntSpec}
\leanref{Formalism/SharpRestriction/GrowthBounds/DualFockLinear/TestSideCounting.lean}{Formalism.MainNotions.TestSideCounting.cntTop}

\begin{proof}
(1) By the standard basis theorem for bitableaux \cite{DEP80} (Lemma~\ref{is-standard-basis}(1)),
$\nabla_{\mathcal R}(\nu)$ has an $\mathcal R$-basis of the bideterminants of the semistandard tableaux
of shape $\nu$ with entries at most $n$: the product, over the columns of the tableau, of the minor of
$x$ on the rows given by the entries $I(1)<\dots<I(s)$ of the column and on the first $s$ columns. In
the expansion of such a minor the product $\prod_px_{I(p),p}$ of its diagonal entries is the leading
monomial, with coefficient $1$. Indeed, let $\pi\neq1$ be a permutation and $p$ the first index with
$\pi(p)\neq p$; then $\pi(p)>p$, the term $\prod_qx_{I(\pi(q)),q}$ agrees with the diagonal one at the
entries before $(I(p),p)$, and it has exponent $0$ at $(I(p),p)$, where the diagonal has $1$. The
leading monomial of a product is the product of the leading monomials, so the leading monomial of the
bideterminant is $x^e$, with coefficient $1$, where $e(a,t)$ is the number of entries $a$ in row $t$ of
the tableau. These matrices $e$ form $E_n(\nu)$: the column sums of $e$ are the row lengths $\nu_t$, a
row of the tableau is determined by its entries as it is weakly increasing, and the inequality at
$(a,t)$ for all $a$ says that the $i$-th entry of row $t+1$ exceeds the $i$-th entry of row $t$ for
every $i$, which is the strict increase of the columns.

(2) By (1) over $\mathcal K$, $|E_n(\nu)|=\dim_{\mathcal K}\nabla_{\mathcal K}(\nu)$, which is
$\dim_n(\nu)$ by the Weyl dimension formula \cite[Theorem~6.3]{FH91}. For every $e\in E_n(\nu)$ the
first row vanishes outside the first column: the inequality at $(1,t)$ gives $e(1,t+1)\leq0$.

Let $n\geq2$ and $e\in E_n(\nu)$ with $e(1,1)=\nu_1$. As the first column sums to $\nu_1$, it is
$(\nu_1,0,\dots,0)$. Let $e'(a,t)=e(a+1,t+1)$ for $a,t\leq n-1$. Its column sums are
$\nu_2,\dots,\nu_n$, and its inequality at $(a,t)$ is that of $e$ at $(a+1,t+1)$, since
$e(1,t+1)=e(1,t+2)=0$. Conversely, for $e'\in E_{n-1}(\nu_2,\dots,\nu_n)$ the matrix $e$ with
$e(1,1)=\nu_1$, $e(a+1,t+1)=e'(a,t)$ and all other entries of the first row and column $0$ lies in
$E_n(\nu)$: at $(1,t)$ the left side of its inequality is $0$; at $(a,t)$ with $a,t\geq2$ the inequality
is that of $e'$ at $(a-1,t-1)$; and at $(a,1)$ with $a\geq2$ the left side is at most $\nu_2$ and the
right side is $\nu_1$. So $e\mapsto e'$ is a bijection onto $E_{n-1}(\nu_2,\dots,\nu_n)$, which has
$\dim_{n-1}(\nu_2,\dots,\nu_n)$ elements by the first claim.

Let moreover $\nu_n=0$, so that the last column of every $e\in E_n(\nu)$ is zero, and let $e$ have zero
first row. Let $e'(a,t)=e(a+1,t)$ for $a,t\leq n-1$. Its column sums are $\nu_1,\dots,\nu_{n-1}$, and
its inequality at $(a,t)$ is that of $e$ at $(a+1,t)$, as the first row of $e$ is zero. Conversely, for
$e'\in E_{n-1}(\nu_1,\dots,\nu_{n-1})$ the matrix $e$ with zero first row and last column and
$e(a+1,t)=e'(a,t)$ lies in $E_n(\nu)$: at $(1,t)$ the left side of its inequality is $0$, at $(a+1,t)$
with $t<n-1$ it is the inequality of $e'$ at $(a,t)$, and at $(a,n-1)$ the left side is $0$. So these
$e$ are $\dim_{n-1}(\nu_1,\dots,\nu_{n-1})$ in number.
\end{proof}

\begin{lemma}\label{pd-symbols-primitive-independent}
Let $F$ be a field, $n'\geq2$ an even integer, $k\geq2$ an integer, $T\subset\mathbf Z/n'\mathbf Z$ and
$d\colon\mathbf Z/n'\mathbf Z\to\mathbf Z_{\geq0}$ with $d(s)=0$ for $s\notin T$; write $s$ also for its
representative in $\{0,\dots,n'-1\}$. For each $s$ let $x_s$ be a $k\times k$ matrix of variables,
$\nu(s)\in\mathbf Z^k$ with $\nu(s)_1\geq\dots\geq\nu(s)_k\geq0$, and $E'(s)\subset E_k(\nu(s))$
(Lemma~\ref{pd-symbols-primitive-basis}) the set $E_k(\nu(s))$ if $s\notin T$, the set of the $e$ with
$e(1,1)=\nu(s)_1$ if $s\in T$ is even, and the set of the $e$ with zero first row if $s\in T$ is odd. Let
$b_s(e)\in F[x_s]$, $e\in E'(s)$, be polynomials whose leading monomial for the column-lexicographic
order is $x_s^e$, with coefficient $1$. Let $\Theta_1\colon F[X]\to F[x_{s,a,b}]$ be the $F$-algebra
map with $\Theta_1(X_{s,c})=x_{s,c,1}$ for $s$ even and $\Theta_1(X_{s,c})=\operatorname{adj}(x_s)_{k,c}$
for $s$ odd. Let $J$ be a finite set of families $i=(i_s)$ with $i_s\in E'(s)$, and let $Q_i\in F[X]$,
$i\in J$, be polynomials of multidegree $d$ with
\[
 \sum_{i\in J}\Theta_1(Q_i)\prod_sb_s(i_s)=0.
\]
Then $Q_i=0$ for every $i\in J$.
\end{lemma}
\leanref{Formalism/SharpRestriction/GrowthBounds/DualFockLinear/ThAllMonomial.lean}{Formalism.MainNotions.TestSideCounting.bigInd}
\leanref{Formalism/SharpRestriction/GrowthBounds/DualFockLinear/DgFalseDiag.lean}{Formalism.MainNotions.TestSideCounting.posInd}
\leanref{Formalism/SharpRestriction/GrowthBounds/DualFockLinear/DgFalseDiag.lean}{Formalism.MainNotions.TestSideCounting.key_inj}

\begin{proof}
Fix a position $s$ and write $x=x_s$. Let $M(c)$ be $x_{c,1}$ if $s$ is even, and, if $s$ is odd, the
minor of $x$ on the rows other than $c$ and the first $k-1$ columns; then
$\Theta_1(X_{s,c})=\pm M(c)$. As in the proof of Lemma~\ref{pd-symbols-primitive-basis}(1), the leading
monomial of $M(c)$ is $x^{\delta(c)}$, with coefficient $1$, where $\delta(c)$ has the single nonzero
entry $1$ at $(c,1)$ if $s$ is even, and if $s$ is odd has entries $1$ at the places $(c_p,p)$, $p<k$,
for the rows $c_1<\dots<c_{k-1}$ other than $c$, and $0$ elsewhere. For $\alpha\in\mathbf Z_{\geq0}^k$
with $\sum_c\alpha_c=d(s)$ and $e\in E'(s)$, let $f_s(\alpha,e)$ be the product of
$\prod_c\Theta_1(X_{s,c})^{\alpha_c}$ and $b_s(e)$. It is $\pm$ a polynomial with leading monomial
$x^{\epsilon(\alpha,e)}$ and coefficient $1$, where $\epsilon(\alpha,e)=\sum_c\alpha_c\delta(c)+e$.

We show that $\epsilon$ is injective. Let $\epsilon(\alpha,e)=\epsilon(\alpha',e')$. If $s\notin T$,
then $d(s)=0$ and $\alpha=\alpha'=0$. If $s\in T$ is even, the first column of $e$ sums to
$e(1,1)=\nu(s)_1$, so it is $(\nu(s)_1,0,\dots,0)$, and the same holds for $e'$; the entries at $(c,1)$
give $\alpha_c=\alpha'_c$. If $s\in T$ is odd, then $e(a,t)=0$ for $a\leq t$, by induction on $t$: for
$t=1$ the first row of $e$ is zero, and for $t\geq2$ and $a\leq t$ the inequality of
Lemma~\ref{pd-symbols-primitive-basis} at $(a,t-1)$ gives $\sum_{b\leq a}e(b,t)\leq\sum_{b<a}e(b,t-1)=0$.
The same holds for $e'$. For $t<k$ the entry of $\delta(c)$ at $(t,t)$ is $1$ if $t<c$ and $0$
otherwise, so the entries at $(t,t)$ give $\sum_{c>t}\alpha_c=\sum_{c>t}\alpha'_c$ for $t<k$; with
$\sum_c\alpha_c=d(s)=\sum_c\alpha'_c$ this gives $\alpha=\alpha'$. In every case then $e=e'$.

Polynomials with distinct leading monomials are linearly independent: in a nontrivial linear relation,
the largest leading monomial among the terms with nonzero coefficient occurs in only one of them. So
the $f_s(\alpha,e)$ are linearly independent over $F$. The families for different positions lie in
disjoint sets of variables, so the products $\prod_sf_s(\alpha_s,e_s)$, over the families
$(\alpha_s,e_s)_s$, are linearly independent. Write
$Q_i=\sum_\alpha q_{i,\alpha}\prod_{s,c}X_{s,c}^{\alpha_{s,c}}$, the sum over the families
$\alpha=(\alpha_s)$ with $\sum_c\alpha_{s,c}=d(s)$. Then
\[
 0=\sum_{i\in J}\Theta_1(Q_i)\prod_sb_s(i_s)=\sum_{i\in J}\sum_\alpha q_{i,\alpha}\prod_sf_s(\alpha_s,i_s),
\]
and all the $q_{i,\alpha}$ vanish.
\end{proof}

\begin{lemma}\label{pd-symbols-primitive-transfer}
Let $p\geq5$ be a prime, $(W,Q,\sigma,\iota)$ a Hida setting of degree $n'$ at $p$ with rings
$\mathcal R\subset\mathcal K$ and $\varpi_0$, $1\leq r<k$ integers,
$d\colon\mathbf Z/n'\mathbf Z\to\mathbf Z_{\geq0}$, $\lambda^\flat(s)=(d(s),0,\dots,0)$ and
$\lambda^\sharp(s)\in\mathbf Z^k$ dominant weights, $s\in\mathbf Z/n'\mathbf Z$, and
$\lambda_0=\lambda^\flat+\lambda^\sharp$. Form $V^\flat$ and $\mathcal K[X]_d$ as in
Definition~\ref{pd-symbols-flat-model}, and the lattices $R$ and $R^\sharp$ of
Definition~\ref{pd-projection-property} with $\lambda_0$ and $\lambda^\sharp$; let $\nu(s)$ be the
natural weight of the weight $\lambda(s)$ formed from $\lambda^\sharp(s)$ as in
Definition~\ref{pd-local-model}. Let $\gamma\in GL_k(Q)$, $g\in GL_k(W)$ and $e=e_\gamma(g)$
(Lemma~\ref{pd-symbols-edge-basis}), let $\Theta_\gamma$ be the flat model of
Lemma~\ref{pd-symbols-flat-exists}, and let $\Theta_1\colon\mathcal K[X]\to\mathcal K[x_{s,a,b}]$ be
the $\mathcal K$-algebra map with $\Theta_1(X_{s,c})=x_{s,c,1}$ for $s$ even and
$\Theta_1(X_{s,c})=\operatorname{adj}(x_s)_{k,c}$ for $s$ odd. Let $I$ be a finite set and, for
$i\in I$, let $z_i$ be an $\mathcal R$-linear combination of products $\prod_sf_s(x_s)$ with
$f_s\in\nabla_{\mathcal R}(\nu(s))$. Then there are $y_i\in R^\sharp(e)$, $i\in I$, with the following
property. Let $P_i\in\mathcal K[X]_d$ and $Q_i\in\mathcal R[X]$ with $\Sigma_g(P_i)=Q_i$ for $i\in I$
(Definition~\ref{pd-symbols-twist}), and suppose that $\sum_i\Theta_\gamma(P_i)y_i\in\varpi_0R(e)$.
Then $\sum_i\Theta_1(Q_i)z_i\in\varpi_0\mathcal R[x_{s,a,b}]$.
\end{lemma}
\leanref{Formalism/SharpRestriction/GrowthBounds/DualFockLinear/DvdOfCompl.lean}{Formalism.MainNotions.TestSideCounting.TransStep}
\leanref{Formalism/SharpRestriction/GrowthBounds/DualFockLinear/TransStep.lean}{Formalism.MainNotions.TestSideCounting.transStep}
\leanref{Formalism/SharpRestriction/GrowthBounds/DualFockLinear/TransStep.lean}{Formalism.MainNotions.TestSideCounting.tr_theta}

\begin{proof}
Let $D$ and $h_s=\gamma_s\,{}^tg_s^{-1}D$ be as in Step 5 of the proof of
Lemma~\ref{pd-symbols-flat-exists}; the columns of $h_s$ form an $\mathcal R$-basis of the lattice from
which the lattices of $e$ are formed at $s$, for every weight. For a family $\mu$ of weights, with
weights $\mu'(s)$ formed from $\mu(s)$ as in Definition~\ref{pd-local-model}, put
$\tau^\mu(f)=\prod_s\det(h_s)^{\mu'(s)_k}f(hx)$ for $f\in\mathcal K[x_{s,a,b}]$, where $f(hx)$ is $f$
with $h_sx_s$ substituted for $x_s$. It is a $\mathcal K$-linear bijection of $\mathcal K[x_{s,a,b}]$
onto itself, and $\tau^{\mu+\mu''}(ff'')=\tau^\mu(f)\tau^{\mu''}(f'')$, as $\mu\mapsto\mu'$ is additive.
By Definition~\ref{pd-projection-property}(1), tested on the bases $h_s$ (the condition does not depend
on the basis, as in Step 4 of the proof of Lemma~\ref{pd-symbols-flat-exists}), $\tau^\mu$ maps the
lattice $H(e)$ formed with $\mu$ onto the $\mathcal R$-span of the products $\prod_sf_s$ with $f_s$ in
the dual Weyl module over $\mathcal R$ of the natural weight of $\mu'(s)$; in particular it maps it into
$\mathcal R[x_{s,a,b}]$. Let $A^\flat$, $A^\sharp$ and $A$ be the integers of
Definition~\ref{pd-projection-property}(1) for $\lambda^\flat$, $\lambda^\sharp$ and $\lambda_0$; then
$A=A^\flat+A^\sharp$.

Put $y_i=\varpi_0^{-A^\sharp}(\tau^{\lambda^\sharp})^{-1}(z_i)$. As $(\tau^{\lambda^\sharp})^{-1}(z_i)$
lies in the lattice $H^\sharp(e)$ formed with $\lambda^\sharp$, $y_i\in R^\sharp(e)$.

Let $P_i$ and $Q_i$ be as in the statement, and $v=\sum_i\Theta_\gamma(P_i)y_i$. By Step 6 of the proof
of Lemma~\ref{pd-symbols-flat-exists}, $\Theta_\gamma(P_i)=\Theta_h(G_i)$ with
$\varpi_0^{A^\flat}G_i=\Sigma_g(P_i)=Q_i$. By Lemma~\ref{pd-symbols-flat-letter}(3) at each position,
$\tau^{\lambda^\flat}(\Theta_h(G))=\Theta_1(G)$ for $G\in\mathcal K[X]_d$. Hence
\[
 \varpi_0^A\tau^{\lambda_0}(v)=\sum_i\varpi_0^{A^\flat}\tau^{\lambda^\flat}(\Theta_\gamma(P_i))\cdot
 \varpi_0^{A^\sharp}\tau^{\lambda^\sharp}(y_i)=\sum_i\Theta_1(Q_i)z_i .
\]
On the other hand $v\in\varpi_0R(e)=\varpi_0^{1-A}H(e)$, and $\tau^{\lambda_0}$ maps $H(e)$ into
$\mathcal R[x_{s,a,b}]$, so $\varpi_0^A\tau^{\lambda_0}(v)\in\varpi_0\mathcal R[x_{s,a,b}]$.
\end{proof}

\begin{lemma}\label{pd-symbols-primitive}
Let $C_{La},C_k,C_{rk},C_{\deg}$ be real numbers with the property of
Theorem~\ref{pd-symbols-count}. There is $C_4>0$ such that the following holds for every $C\geq C_4$.
Let the data of Situation~\ref{pd-edge-situation} with this $C$, and $(J,f)$, $K$, $N_s$, $q$, $c'$,
$\deg$, $m_*$, $E_0$, $X_j$, $d_T$, $E'$, $d$, $P_1,\dots,P_{n_b}$, $e$, $\gamma$, $g$,
$\Phi_1,\dots,\Phi_{n_b}$, $\lambda^\flat$, $\lambda^\sharp$, $V^\flat$ and $R^\sharp$ satisfy the
hypotheses of Proposition~\ref{pd-symbols-cloud}, and let $\Theta_\gamma$ be the flat model of
$V^\flat$ at $\gamma$ of Lemma~\ref{pd-symbols-flat-exists}, formed with $d_T$. Then there are vectors
$y_1,\dots,y_{n_2}\in R^\sharp(e)$, with $n_2$ the integer of Lemma~\ref{pd-edge-dimension-count},
such that $\sum_{b,i}c_{bi}\Theta_\gamma(\Phi_b)y_i\notin\varpi_0R(e)$ for every family $(c_{bi})$ in
$\mathcal R$ with some $c_{bi}$ a unit.
\end{lemma}
\leanref{Formalism/SharpRestriction/GrowthBounds/ScaleEdges/SourcePrimitiveIntegrality.lean}{Formalism.MainNotions.Pf14.SourcePrimitiveIntegrality}
\leanref{Formalism/SharpRestriction/GrowthBounds/ScaleEdges/SourcePrimitiveIntegrality.lean}{Formalism.MainNotions.Pf14.PrimStep}
\leanref{Formalism/SharpRestriction/GrowthBounds/ScaleEdges/SourcePrimitiveIntegrality.lean}{Formalism.MainNotions.Pf14.PrimAt}
\leanref{Formalism/SharpRestriction/GrowthBounds/DualFockLinear/PrimStep.lean}{Formalism.MainNotions.TestSideCounting.rest6a}
\leanref{Formalism/SharpRestriction/GrowthBounds/DualFockLinear/TheoremIIClosedOfTestBound.lean}{Formalism.MainNotions.TestSideCounting.primStep_of}
\leanref{Formalism/SharpRestriction/GrowthBounds/DualFockLinear/TestSideCounting.lean}{Formalism.MainNotions.TestSideCounting.card_all}
\leanref{Formalism/SharpRestriction/GrowthBounds/DualFockLinear/SubstSPtrue.lean}{Formalism.MainNotions.TestSideCounting.spStep}

\begin{proof}
Put $C_4=1$, let $C\geq1$ and let the data be as in the statement. By (3) of
Situation~\ref{pd-edge-situation}, $N_1+2r\leq k-r$, so $k\geq3$ and $1\leq r<k$. For $s\in T$,
$d_T(s)=\deg(s)=\lambda_0(s)_1-\lambda_0(s)_2$, so
$\lambda^\sharp(s)=(\lambda_0(s)_2,\lambda_0(s)_2,\lambda_0(s)_3,\dots,\lambda_0(s)_k)$; for $s\notin T$,
$\lambda^\sharp(s)=\lambda_0(s)$. These weights are dominant. Let $\nu(s)$ be the natural weight of the
weight $\lambda(s)$ formed from $\lambda^\sharp(s)$ as in Definition~\ref{pd-local-model}, and $\kappa$
the residue field of $\mathcal R$.

\emph{Step 1: the basis.} For each $s$ let $(b_s(e))_{e\in E_k(\nu(s))}$ be the basis of
$\nabla_{\mathcal R}(\nu(s))$ of Lemma~\ref{pd-symbols-primitive-basis}(1), and let
$E'(s)\subset E_k(\nu(s))$ be as in Lemma~\ref{pd-symbols-primitive-independent}: all of
$E_k(\nu(s))$ for $s\notin T$, the $e$ with $e(1,1)=\nu(s)_1$ for $s\in T$ even, and the $e$ with zero
first row for $s\in T$ odd. Let $I$ be the set of families $i=(i_s)$ with $i_s\in E'(s)$, and put
$z_i=\prod_sb_s(i_s)(x_s)$. We show $|I|=n_2$, using Lemma~\ref{pd-symbols-primitive-basis}(2) and that
$\dim_j$ depends only on the differences of the entries and is unchanged by
$\mu\mapsto(-\mu_j,\dots,-\mu_1)$. For $s\notin T$,
$|E'(s)|=\dim_k(\nu(s))=\dim_k(\lambda_0(s))$. For $s\in T$ even,
$\nu(s)=\lambda^\sharp(s)-\lambda_0(s)_k(1,\dots,1)$ and
$|E'(s)|=\dim_{k-1}(\nu(s)_2,\dots,\nu(s)_k)=\dim_{k-1}(\lambda_0(s)_2,\dots,\lambda_0(s)_k)$. For
$s\in T$ odd, $\lambda(s)=(-\lambda_0(s)_k,\dots,-\lambda_0(s)_3,-\lambda_0(s)_2,-\lambda_0(s)_2)$, so
$\nu(s)_k=0$, $(\nu(s)_1,\dots,\nu(s)_{k-1})=(\lambda_0(s)_2-\lambda_0(s)_k,\dots,
\lambda_0(s)_2-\lambda_0(s)_2)$ and $|E'(s)|=\dim_{k-1}(\lambda_0(s)_2,\dots,\lambda_0(s)_k)$. The
product of these numbers is $n_2$.

\emph{Step 2: the vectors.} Lemma~\ref{pd-symbols-primitive-transfer}, applied to $d_T$,
$\lambda^\sharp$, $\gamma$, $g$ and the $z_i$, gives $y_i\in R^\sharp(e)$ for $i\in I$; we number them
$y_1,\dots,y_{n_2}$.

\emph{Step 3: the multiplier modulo $\varpi_0$.} By Step 4 of the proof of
Proposition~\ref{pd-symbols-realize}, $\Sigma_g(\mathrm{SP}_g)$ is a product of powers of determinants of
$r\times r$ matrices whose row $l$ is $(X_{s_l,1}^{q_l},\dots,X_{s_l,r}^{q_l})$ for a position $s_l$ and
a power $q_l$ of $p$; so it lies in $\mathbf Z[X]$. In each determinant the positions $s_l$ are distinct:
for $\Delta_j(g)$ they are the residues of the $f(j,l)$, $l<r$, which lie in an interval of length
$f(j,r-1)-f(j,0)<G'\leq n'$; for $\Delta^{\mathrm{top}}_t(g)$ they are the residues of the distinct
$a_l(t)$, which lie in $(t-G',t)$ by fact (b) in Step 2 of that proof. Evaluating $X_{s_l,l}$ at $1$
for every $l$ and all other variables at $0$ sends such a determinant to $\det(1)=1$, in $\kappa$ as in
$\mathbf Z$; so its reduction modulo $\varpi_0$ is nonzero. As $\kappa[X]$ is a domain, the reduction of
$\Sigma_g(\mathrm{SP}_g)$ modulo $\varpi_0$ is nonzero.

\emph{Step 4: reduction.} Let $(c_{bi})$ be a family in $\mathcal R$ with $c_{b_0i_0}$ a unit, and
suppose that $\sum_{b,i}c_{bi}\Theta_\gamma(\Phi_b)y_i\in\varpi_0R(e)$. For $i\in I$ put
$P_i=\sum_bc_{bi}\Phi_b\in\mathcal K[X]_{d_T}$, $U_i=\sum_bc_{bi}P_b\in\mathcal R[X]$ and
$Q_i=U_i\,\Sigma_g(\mathrm{SP}_g)\in\mathcal R[X]$. As $\Sigma_g(\Phi_b)$ is the symbol multiplier
$P_b\,\Sigma_g(\mathrm{SP}_g)$, $\Sigma_g(P_i)=Q_i$; and
$\sum_{b,i}c_{bi}\Theta_\gamma(\Phi_b)y_i=\sum_i\Theta_\gamma(P_i)y_i$. By
Lemma~\ref{pd-symbols-primitive-transfer}, $\sum_i\Theta_1(Q_i)z_i\in\varpi_0\mathcal R[x_{s,a,b}]$. With
bars for reductions modulo $\varpi_0$, $\sum_i\Theta_1(\bar Q_i)\prod_s\bar b_s(i_s)=0$. The substitution
$\Sigma_g$ preserves multidegrees, so $\bar Q_i$ has multidegree $d_T$, and $d_T(s)=0$ for $s\notin T$;
the leading monomial of $\bar b_s(e)$ is still $x_s^e$, with coefficient $1$. By
Lemma~\ref{pd-symbols-primitive-independent} over $\kappa$, $\bar Q_i=0$ for every $i$. By Step 3,
$\bar U_{i_0}=0$, that is, $\sum_bc_{bi_0}P_b=\varpi_0y$ with $y\in\mathcal R[X]$.

\emph{Step 5: a contradiction.} Let $S$ be the span of $P_1,\dots,P_{n_b}$ and $S'$ a complement of $S$
in $\mathcal R[X]$, and write $y=\sum_bd_bP_b+y'$ with $d_b\in\mathcal R$ and $y'\in S'$. Then
$\sum_b(c_{bi_0}-\varpi_0d_b)P_b=\varpi_0y'$ lies in $S\cap S'=0$. As the $P_b$ are linearly independent,
$c_{bi_0}=\varpi_0d_b$ for every $b$, which contradicts that $c_{b_0i_0}$ is a unit.
\end{proof}

\begin{definition}\label{pd-symbols-test-operators}
Let $n'\geq1$, $k\geq1$ and $c_0\geq0$ be integers, $T\subset\mathbf Z/n'\mathbf Z$ and $\mathcal K$ a
field, and let $\mathcal K[X]$ be the polynomial ring in $X_{s,c}$, $s\in\mathbf Z/n'\mathbf Z$,
$1\leq c\leq k$. The \emph{operator closure} $\mathrm{Op}_{T,c_0}(M)$ of a $\mathcal K$-subspace
$M\subset\mathcal K[X]$ is the $\mathcal K$-span of the polynomials $D_1D_2\cdots D_lf$ with $f\in M$
and $l\geq0$, where each $D_i$ is an operator $X_{s,a}\,\partial/\partial X_{s,b}$ with $s\in T$ and
$1\leq a,b\leq k$, and each $s$ occurs in at most $c_0$ of the $D_i$.
\end{definition}
\leanref{Formalism/SharpRestriction/GrowthBounds/DualFockLinear/TestSideExplicitIdentification.lean}{Formalism.MainNotions.TestSideCounting.opWord}
\leanref{Formalism/SharpRestriction/GrowthBounds/DualFockLinear/TestSideExplicitIdentification.lean}{Formalism.MainNotions.TestSideCounting.OpCl}
\leanref{Formalism/SharpRestriction/GrowthBounds/DualFockLinear/TestSideExplicitIdentification.lean}{Formalism.MainNotions.TestSideCounting.SgOp}

\begin{lemma}\label{pd-symbols-test-size}
There are $C_K,C_1>0$ such that the following holds for every $C\geq C_1$. Let the data of
Situation~\ref{pd-edge-situation} be given with this $C$, and form $d_T$ and $\lambda^\flat$ as in
Proposition~\ref{pd-symbols-cloud}. Let $\lambda_0^T(s)=\lambda_0(s)$ for $s\in T$ and $\lambda_0^T(s)=0$
for $s\notin T$. For dominant weights $\mu_0(s)\in\mathbf Z^k$, $s\in\mathbf Z/n'\mathbf Z$, let $\mu(s)$ be
formed from $\mu_0(s)$ as in Definition~\ref{pd-local-model}, $\ell_s=\mu(s)_1-\mu(s)_k$, and
\[
 \kappa_1(\mu_0)=\sum_{s:\ \ell_s\neq0}\frac{k(k-1)}2\big\lfloor\log_p(\ell_s+k-1)\big\rfloor.
\]
Then $\kappa_1(\lambda_0^T)+\kappa_1(\lambda^\flat)\leq C_K|T|k^2\log_pa$.
\end{lemma}
\leanref{Formalism/SharpRestriction/GrowthBounds/DualFockLinear/RestCloudBOfThree.lean}{Formalism.MainNotions.TestSideCounting.SizeStep}
\leanref{Formalism/SharpRestriction/GrowthBounds/DualFockLinear/RestCloudBOfThree.lean}{Formalism.MainNotions.TestSideCounting.RestSize}
\leanref{Formalism/SharpRestriction/GrowthBounds/ObjBridge.lean}{Formalism.MainNotions.G.kap1}
\leanref{Formalism/SharpRestriction/GrowthBounds/DualFockLinear/RestCloudBOfTop.lean}{Formalism.MainNotions.TestSideCounting.topW}

\begin{proof}
Put $C_K=4$ and $C_1=2$, let $C\geq2$, and let the data be as in the statement. As in the proof of
Lemma~\ref{pd-edge-numerics}, $|T|\geq1$. By (3) of the situation, $k\geq N_1+3r\geq3$, $Y\geq r\geq2$
and $c_*\geq8$. As $p^{G'}\geq1$ and $c_0\geq0$, the last inequality of (3) gives
$a\geq2\cdot k\cdot4\cdot8\cdot2\geq8k$ and $a\geq k\cdot r|T|c_0\geq kc_0$. So $a\geq8$, $k\leq a/8$,
$c_0\leq a$ and $4a\leq a^2$.

\emph{A bound for one weight.} Let $\mu_0(s)\in\mathbf Z^k$ be dominant weights such that, with
$\ell_s$ as in the statement, $\ell_s=0$ for $s\notin T$ and $\ell_s+k\leq a^2$ for $s\in T$. Then
$\kappa_1(\mu_0)\leq2|T|k^2\log_pa$. Indeed, a term with $\ell_s\neq0$ has $s\in T$, and there
$k(k-1)/2\leq k^2$ and $\lfloor\log_p(\ell_s+k-1)\rfloor\leq\log_p(a^2)=2\log_pa$; and $\log_pa\geq0$.

\emph{The two weights.} Forming $\mu(s)$ from $\mu_0(s)$ as in Definition~\ref{pd-local-model} keeps
$\mu_1-\mu_k$, so $\ell_s=\mu_0(s)_1-\mu_0(s)_k$. For $s\notin T$, $\lambda_0^T(s)=0$ and
$\lambda^\flat(s)=0$, so both have $\ell_s=0$. Let $s\in T$. For $\lambda_0^T$, by (2) of the situation,
$\ell_s=a_T(s)+b_T(s)\leq(2a-c_{pp})+(c_0-2c_{pp})\leq2a+c_0\leq3a$. For $\lambda^\flat$,
$\ell_s=d_T(s)=a_T(s)-\lambda_0(s)_2\leq(2a-c_{pp})+c_{pp}=2a$, as $|\lambda_0(s)_2|\leq c_{pp}$ because
$2<k$. In both cases $\ell_s+k\leq3a+a/8\leq4a\leq a^2$. The bound for one weight, applied to
$\lambda_0^T$ and to $\lambda^\flat$, gives $\kappa_1(\lambda_0^T)+\kappa_1(\lambda^\flat)\leq4|T|k^2\log_pa$.
% Lean: restSize = sizeStep with CK = 4 and C1 = 2 (LogSizeBoundLe.lean): num_size2 (a >= 8k),
%   rank_mul_conductor_div_le_one (k c0 <= a), log_size_bound_le (the bound for one weight), applied to
%   topW TT lam0 (with lamAt_zero, ellWt_zero) and to flW (dT TT lam0) (ellWt_flat, flatDeg_bounds). kap1 p
%   lam0 is the sum over s with ellWt (lamAt lam0 s) /= 0 of N0 (m + 1) * deF p (m + 1) (ellWt (lamAt lam0 s)),
%   with N0 n = n (n - 1) / 2 and deF p n l = Nat.log p (l + n - 1).
\end{proof}

\begin{lemma}\label{pd-symbols-test-onto}
Let $p\geq5$ be a prime, $(W,Q,\sigma,\iota)$ a Hida setting of degree $n'$ at $p$ with rings
$\mathcal R\subset\mathcal K$ and $\varpi_0$, $k\geq1$ and $r\geq0$ integers,
$d\colon\mathbf Z/n'\mathbf Z\to\mathbf Z_{\geq0}$, $\lambda^\flat(s)=(d(s),0,\dots,0)$, and $\lambda'_0(s)\in
\mathbf Z^k$ dominant weights; put $\lambda_0=\lambda^\flat+\lambda'_0$. For each of $\lambda_0$,
$\lambda^\flat$ and $\lambda'_0$ form the space and the lattices $R(e)$ of
Definition~\ref{pd-projection-property}, written $R_{\lambda_0}(e)$, $R_{\lambda^\flat}(e)$ and
$R_{\lambda'_0}(e)$. Then for every edge $e$, $R_{\lambda_0}(e)$ is contained in the $\mathcal R$-span of
the products $y\beta$ in $\mathcal K[x_{s,a,b}]$ with $y\in R_{\lambda^\flat}(e)$ and
$\beta\in R_{\lambda'_0}(e)$.
\end{lemma}
\leanref{Formalism/SharpRestriction/GrowthBounds/DualFockLinear/RestCloudBOfThree.lean}{Formalism.MainNotions.TestSideCounting.OntoAt}
\leanref{Formalism/SharpRestriction/GrowthBounds/DualFockLinear/RestCloudBOfThree.lean}{Formalism.MainNotions.TestSideCounting.OntoAll}

\begin{proof}
Let $e$ be an edge with source $L''$ and target $L$. For $s\in\mathbf Z/n'\mathbf Z$ let $M_s$ be the
$\mathcal R$-span of $\iota\sigma^s(L'')$ and $\varpi_0\iota\sigma^s(L)$, as in
Definition~\ref{pd-projection-property}(1). It lies between the $\mathcal R$-spans of $\iota\sigma^s(pL)$
and $\iota\sigma^s(L)$, which are lattices in $\mathcal K^k$, so it is a lattice. For $\mu\in\mathbf Z^k$ write $\mu^{(s)}$
for the weight formed from $\mu$ at $s$ as in Definition~\ref{pd-local-model} ($\mu$ if $s$ is even,
$(-\mu_k,\dots,-\mu_1)$ if $s$ is odd); for a family $\mu_0$ of dominant weights write $H_{\mu_0}(e)$ for
the lattice $H(e)$ formed with $\mu_0$, and $A(\mu_0)$ for the integer $A$ of
Definition~\ref{pd-projection-property}(1). The weights $\lambda^\flat(s)$ are dominant, and
$\mu\mapsto\mu^{(s)}$ is additive and keeps dominance; so
$\lambda_0(s)^{(s)}=\lambda^\flat(s)^{(s)}+\lambda'_0(s)^{(s)}$ with both summands dominant, and
$A(\lambda_0)=A(\lambda^\flat)+A(\lambda'_0)$.

\emph{Step 1: one position.} By Lemma~\ref{pd-kappa-cloud-cartan}(2), applied to the lattice $M_s$ and
the dominant weights $\lambda^\flat(s)^{(s)}$ and $\lambda'_0(s)^{(s)}$, the lattice
$\nabla_{\lambda_0(s)^{(s)}}(M_s^\vee)$ is the $\mathcal R$-span of the products $ff'$ with
$f\in\nabla_{\lambda^\flat(s)^{(s)}}(M_s^\vee)$ and $f'\in\nabla_{\lambda'_0(s)^{(s)}}(M_s^\vee)$.

\emph{Step 2: all positions.} Let $\prod_sf_s(x_s)$ be one of the products that span
$H_{\lambda_0}(e)$, with $f_s\in\nabla_{\lambda_0(s)^{(s)}}(M_s^\vee)$. By Step 1,
$f_s=\sum_ic_{s,i}g_{s,i}h_{s,i}$ with $c_{s,i}\in\mathcal R$,
$g_{s,i}\in\nabla_{\lambda^\flat(s)^{(s)}}(M_s^\vee)$ and $h_{s,i}\in\nabla_{\lambda'_0(s)^{(s)}}(M_s^\vee)$.
Multiplying out,
\[
 \prod_sf_s(x_s)=\sum_{(i_s)}\Big(\prod_sc_{s,i_s}\Big)\prod_sg_{s,i_s}(x_s)\cdot\prod_sh_{s,i_s}(x_s),
\]
and the two products over $s$ lie in $H_{\lambda^\flat}(e)$ and $H_{\lambda'_0}(e)$. So $H_{\lambda_0}(e)$
lies in the $\mathcal R$-span of the products $PQ$ with $P\in H_{\lambda^\flat}(e)$ and
$Q\in H_{\lambda'_0}(e)$.

\emph{Step 3: the scaling.} Let $z\in R_{\lambda_0}(e)$, so $z=\varpi_0^{-A(\lambda_0)}z_0$ with
$z_0\in H_{\lambda_0}(e)$. By Step 2, $z_0=\sum_jc_jP_jQ_j$ with $c_j\in\mathcal R$, $P_j\in H_{\lambda^\flat}(e)$
and $Q_j\in H_{\lambda'_0}(e)$. As $A(\lambda_0)=A(\lambda^\flat)+A(\lambda'_0)$,
\[
 z=\sum_jc_j\big(\varpi_0^{-A(\lambda^\flat)}P_j\big)\big(\varpi_0^{-A(\lambda'_0)}Q_j\big),
\]
with $\varpi_0^{-A(\lambda^\flat)}P_j\in R_{\lambda^\flat}(e)$ and $\varpi_0^{-A(\lambda'_0)}Q_j\in
R_{\lambda'_0}(e)$.
% Lean: ontoAll. The lattices of e at s are colSpan O (gs * lb_D ...) (Model.edgeLat); additivity is
%   lamAt_add and Ar_add, dominance flW_dom and kp_lamAt_antitone. Step 1 is Model.cartanOnto_proved
%   (CartanOnto, Lemma pd-kappa-cloud-cartan(2)); Step 2 is tensor_onto
%   (UpperTheory.Reformulated.B.product_first_partition.tensor_onto); Step 3 is the span induction `key`
%   (with Hten_le_V and V_mul for membership in the spaces V).
\end{proof}

\begin{definition}\label{pd-symbols-test-fock}
Let $\mathcal K$ be a field and $\mathcal K[x]$ a polynomial ring in finitely many variables $x_v$. For a
monomial $x^m=\prod_vx_v^{m(v)}$ put $m!=\prod_vm(v)!$, and for $f\in\mathcal K[x]$ let $f_m$ be its
coefficient at $x^m$. The \emph{Fock pairing} of $f,g\in\mathcal K[x]$ is
$\langle f,g\rangle_{\mathrm F}=\sum_mf_mg_m\,m!$. If the variables are split into disjoint sets and
$f=\prod_sf_s$, $g=\prod_sg_s$ with $f_s,g_s$ in the variables of the $s$-th set, then
$\langle f,g\rangle_{\mathrm F}=\prod_s\langle f_s,g_s\rangle_{\mathrm F}$, as $m!$ is the product of the factorials
of the parts of $m$.
\end{definition}
\leanref{Formalism/LocalToGlobal/GroupLattices/FockModels/TransposeWords/FockR.lean}{Formalism.L2.G.Model.TS.fockR}
\leanref{Formalism/LocalToGlobal/GroupLattices/FockModels/TransposeWords/TransposeWordsProofs/ProdBlockEmbMonomial.lean}{Formalism.L2.G.Model.TS.fock_pairing_bilinear}

\begin{lemma}\label{pd-symbols-test-fock-position}
Let $p\geq5$ be a prime and $(W,Q,\sigma,\iota)$ a Hida setting of degree $n'$ at $p$ with rings
$\mathcal R\subset\mathcal K$. Let $k\geq1$, $x$ a $k\times k$ matrix of variables, $\mu\in\mathbf Z^k$
dominant, and $\nabla_{\mathcal K}(\mu)$, $\nabla_{\mathcal R}(\mu)$ and, for an $\mathcal R$-lattice
$\Lambda\subset\mathcal K^k$, $\nabla_\mu(\Lambda^\vee)$ as in Lemma~\ref{pd-kappa-cloud-cartan}. Put
$\ell=\mu_1-\mu_k$, and $N(\mu)=0$ if $\ell=0$ and
$N(\mu)=\frac{k(k-1)}2\lfloor\log_p(\ell+k-1)\rfloor$ otherwise. Then there is $c\in\mathcal K^\times$ such that
$B(P,Q)=c\langle P,Q\rangle_{\mathrm F}$ (Definition~\ref{pd-symbols-test-fock}) satisfies:
\begin{enumerate}
\item $B(P(gx),Q)=B(P,Q({}^tgx))$ for $g\in GL_k(\mathcal K)$ and $P,Q\in\mathcal K[x]$;
\item $B$ is nondegenerate on $\nabla_{\mathcal K}(\mu)$;
\item if $Q\in\nabla_{\mathcal K}(\mu)$ and $B(P,Q)\in\mathcal R$ for every $P\in\nabla_{\mathcal R}(\mu)$, then
  $Q\in\nabla_{\mathcal R}(\mu)$;
\item $B(P,p^{N(\mu)}Q)\in\mathcal R$ for $P,Q\in\nabla_{\mathcal R}(\mu)$;
\item for every $\mathcal R$-lattice $\Lambda\subset\mathcal K^k$: if $Q\in\nabla_{\mathcal K}(\mu)$ and
  $B(P,Q)\in\mathcal R$ for every $P\in\nabla_\mu(\Lambda^\vee)$, then $Q\in\nabla_\mu(\Lambda)$; and
  $B(P,p^{N(\mu)}Q)\in\mathcal R$ for $P\in\nabla_\mu(\Lambda^\vee)$ and $Q\in\nabla_\mu(\Lambda)$.
\end{enumerate}
\end{lemma}
\leanref{Formalism/SharpRestriction/GrowthBounds/DualFockLinear/FormCoreSKXOf.lean}{Formalism.L2.G.Model.TS.FormCoreSKX}
\leanref{Formalism/SharpRestriction/GrowthBounds/DualFockLinear/FormCoreSKXOf.lean}{Formalism.L2.G.Model.TS.formCoreSKX}
\leanref{Formalism/LocalToGlobal/GroupLattices/FockModels/TestSideLat.lean}{Formalism.L2.G.Model.PosCore}
\leanref{Formalism/SharpRestriction/GrowthBounds/DualFockLinear/JetMatrixAlgebra.lean}{Formalism.L2.G.Model.posCore_of_formX}

\begin{proof}
Let $\nu=\nu(\mu)=(\mu_1-\mu_k,\dots,\mu_k-\mu_k)$ be the natural weight of $\mu$, so that $\nu_1=\ell$,
$\nu_k=0$, $\nabla_{\mathcal K}(\mu)=\nabla_{\mathcal K}(\nu)$ and $\nabla_{\mathcal R}(\mu)=\nabla_{\mathcal R}(\nu)$,
and use the modules $\Delta_A(\nu)$ and the polynomial $h_\nu$ of Definition~\ref{pd-weight-lattices}(1)
for the matrix $x$. The field $\mathcal K$ contains $\iota(Q)$, so it has characteristic $0$, as $p$ is a
uniformizer of $W$ and so neither $0$ nor a unit there. The ring $\mathcal R$ is a local domain, and $p$ is
not a unit in $\mathcal R$: otherwise $\mathcal R=p\mathcal R$, and since $\mathcal R$ is a finitely generated
$W$-module and $p$ lies in the maximal ideal of $W$, Nakayama's lemma would give $\mathcal R=0$.

By Lemma~\ref{pd-kappa-gram-fock}(3), $N=\langle h_\nu,h_\nu\rangle_{\mathrm F}$ is a positive integer. Put
$c=N^{-1}$ and $B(P,Q)=c\langle P,Q\rangle_{\mathrm F}$; this is the form $B_0$ of
Lemma~\ref{pd-kappa-gram-form} with $r=k$. Step~1 and parts (2) and (3) of the proof of that lemma use of
$\mathcal R$ only that it is a domain whose fraction field $\mathcal K$ has characteristic $0$, so they hold
here: $B(P,Q)\in\mathcal R$ for $P\in\Delta_{\mathcal R}(\nu)$ and $Q\in\nabla_{\mathcal R}(\nu)$; a
$Q\in\nabla_{\mathcal K}(\nu)$ with $B(P,Q)=0$ for all $P\in\nabla_{\mathcal K}(\nu)$ is $0$; and a
$Q\in\nabla_{\mathcal K}(\nu)$ with $B(P,Q)\in\mathcal R$ for all $P\in\Delta_{\mathcal R}(\nu)$ lies in
$\nabla_{\mathcal R}(\nu)$. Moreover $\Delta_{\mathcal R}(\nu)\subset\nabla_{\mathcal R}(\nu)$, by
Lemma~\ref{pd-kappa-defects-sandwich} with $r=k$ and $l=\ell$, as every entry of $\nu$ is at most $\ell$.

(1) This is Lemma~\ref{pd-kappa-gram-fock}(1) over $\mathcal K$, multiplied by $c$.

(2) This is the second of the three facts above.

(3) Let $Q$ be as in (3). As $\Delta_{\mathcal R}(\nu)\subset\nabla_{\mathcal R}(\nu)$, $B(P,Q)\in\mathcal R$ for
every $P\in\Delta_{\mathcal R}(\nu)$, so $Q\in\nabla_{\mathcal R}(\nu)$ by the third fact.

(4) Let $P,Q\in\nabla_{\mathcal R}(\nu)$. If $\ell=0$, then $\nu=0$ and $h_\nu=1$, so $N=1$ and
$B(P,p^{N(\mu)}Q)=\langle P,Q\rangle_{\mathrm F}=\sum_mP_mQ_m\,m!$, which lies in $\mathcal R$ because $P$ and $Q$
have coefficients in $\mathcal R$. Let $\ell\neq0$; then $k\geq2$, and
Lemma~\ref{pd-kappa-defects-sandwich}, with $r=k$, $l=\ell$ and the local domain $\mathcal R$, gives
$p^{N(\mu)}Q\in\Delta_{\mathcal R}(\nu)$. The Fock pairing is
symmetric, so $B(P,p^{N(\mu)}Q)=B(p^{N(\mu)}Q,P)$, which lies in $\mathcal R$ by the first fact.
% Lean: (1) to (4) are formCoreSKX = formCoreSKX_of (formCoreGX_of Top.fockRBilin Top.fockRRowAct Top.hWNorm
%   Top.weylPair Top.dualWeylIn Top.nondegK): not_isUnit_p_of_hidaSetting (Module.Finite W O,
%   Submodule.eq_bot_of_le_smul_of_le_jacobson_bot) and the CharZero instances; B = (N : K)^-1 fockR with
%   N from hWNorm; (1) is hR = fockRRowAct; (2) is hG = nondegK (nondegK_of, part (2) of the proof of Lemma
%   pd-kappa-gram-form); (3) is hD = dualWeylIn (dualWeylIn_of, part (3) there) with hle from
%   Top.sandwichLocal .2; (4) is hbd with l = ellWt lam (nuOf_le_ellWt), sExp p (m + 1) (ellWt lam) = expN,
%   (Top.sandwichLocal .1, hsymm, hW = weylPair, which is weylPair_of, Step 1 there), and hzero (hw = 1,
%   N = 1, Top.fockRMap) when ellWt lam = 0.

(5) Let $g$ be a matrix whose columns form an $\mathcal R$-basis of $\Lambda$, and put $e=\mu_k$. The columns
of ${}^tg^{-1}$ form a basis of the dual lattice $\Lambda^\vee$, so $f\in\nabla_\mu(\Lambda)$ if and only if
$\det({}^tg^{-1})^ef({}^tg^{-1}x)\in\nabla_{\mathcal R}(\mu)$, and $f\in\nabla_\mu(\Lambda^\vee)$ if and only if
$\det(g)^ef(gx)\in\nabla_{\mathcal R}(\mu)$. By (1), as $\det(g)\det({}^tg^{-1})=1$,
\[
 B\big(\det(g)^eP(gx),\det({}^tg^{-1})^eQ({}^tg^{-1}x)\big)=B(P,Q)\qquad(P,Q\in\mathcal K[x]).
\]
Let $Q$ be as in the first claim and $P_0\in\nabla_{\mathcal R}(\mu)$. Then
$P=\det(g)^{-e}P_0(g^{-1}x)$ lies in $\nabla_\mu(\Lambda^\vee)$, and the identity gives
$B(P_0,\det({}^tg^{-1})^eQ({}^tg^{-1}x))=B(P,Q)\in\mathcal R$. As the substitution keeps
$\nabla_{\mathcal K}(\mu)$ (Lemma~\ref{pd-kappa-cloud-cartan}(1)), (3) gives
$\det({}^tg^{-1})^eQ({}^tg^{-1}x)\in\nabla_{\mathcal R}(\mu)$, that is, $Q\in\nabla_\mu(\Lambda)$. For the second
claim, put $P_0=\det(g)^eP(gx)$ and $Q_0=\det({}^tg^{-1})^eQ({}^tg^{-1}x)$, both in $\nabla_{\mathcal R}(\mu)$; then
$B(P,p^{N(\mu)}Q)=B(P_0,p^{N(\mu)}Q_0)\in\mathcal R$ by (4).
% Lean: (5) is posCore_of_formX (the identity is key, with ldlat_rowAct_inv_right and det_invT_zpow; up and
%   low are the two claims, through nablaLTwMem).
\end{proof}

\begin{lemma}\label{pd-symbols-test-fock-dual}
Let $p\geq5$ be a prime, $(W,Q,\sigma,\iota)$ a Hida setting of degree $n'$ at $p$ with rings
$\mathcal R\subset\mathcal K$ and $\varpi_0$, $k\geq1$ and $r\geq0$ integers, and
$\lambda_0(s)\in\mathbf Z^k$, $s\in\mathbf Z/n'\mathbf Z$, dominant weights. Form $V$, $H$, $R$ and $R'$ with
$\lambda_0$ as in Definition~\ref{pd-projection-property}, let $\kappa_1$ be as in
Lemma~\ref{pd-symbols-test-size}, and let $L^\vee$ and $e'^\vee$ be the duals of a lattice and of an edge as in
Lemma~\ref{pd-symbols-test-duality}. Then there are a $\mathcal K$-linear isomorphism $\Psi\colon V^\vee\to V$ and
$c\in\mathcal K^\times$ with $\Psi^{-1}(y)(u)=c\langle u,y\rangle_{\mathrm F}$ for all $u,y\in V$, such that for
every $W$-lattice $L$, $\Psi(H(L)^\vee)\subset H(L^\vee)$ and $p^{\kappa_1(\lambda_0)}H(L^\vee)\subset
\Psi(H(L)^\vee)$, and for every edge $e'$, $\Psi(R'(e'))\subset R(e'^\vee)$ and
$p^{\kappa_1(\lambda_0)}R(e'^\vee)\subset\Psi(R'(e'))$.
\end{lemma}
\leanref{Formalism/SharpRestriction/GrowthBounds/DualFockLinear/BaseMatrixAlgebra.lean}{Formalism.MainNotions.TestSideCounting.RestTSX}
\leanref{Formalism/SharpRestriction/GrowthBounds/DualFockLinear/JetMatrixAlgebra.lean}{Formalism.L2.G.Model.test_side_explicit_identification}
\leanref{Formalism/SharpRestriction/GrowthBounds/DualFockLinear/JetMatrixAlgebra.lean}{Formalism.L2.G.Model.ten_sideX}

\begin{proof}
For each $s$ let $\lambda(s)$ be formed from $\lambda_0(s)$ as in Definition~\ref{pd-local-model}, and let
$c_s$, $B_s$ and $N_s=N(\lambda(s))$ be the constant, the form and the integer of
Lemma~\ref{pd-symbols-test-fock-position} for $\lambda(s)$. As $\lambda(s)_1-\lambda(s)_k=\lambda_0(s)_1-\lambda_0(s)_k$,
$\sum_sN_s=\kappa_1(\lambda_0)$. Put $c=\prod_sc_s$. If $(b_{s,i})_i$ is a $\mathcal K$-basis of
$\nabla_{\mathcal K}(\lambda(s))$ for each $s$, the products $\prod_sb_{s,i_s}(x_s)$ form a $\mathcal K$-basis of $V$,
as the $x_s$ are disjoint sets of variables. By Definition~\ref{pd-symbols-test-fock},
\[
 c\Big\langle\prod_su_s(x_s),\prod_sy_s(x_s)\Big\rangle_{\mathrm F}=\prod_sB_s(u_s,y_s).
\]
So in such product bases the matrix of $(u,y)\mapsto c\langle u,y\rangle_{\mathrm F}$ on $V$ is the tensor product of
the matrices of the $B_s$, which are invertible by Lemma~\ref{pd-symbols-test-fock-position}(2). Hence this form
is nondegenerate on $V$, and there is a unique $\Psi$ with $\Psi^{-1}(y)=c\langle\,\cdot\,,y\rangle_{\mathrm F}$.

\emph{Lattices.} Let $M=(M_s)$ be a family of $\mathcal R$-lattices in $\mathcal K^k$, $M^\vee=(M_s^\vee)$, and $H(M)$
the $\mathcal R$-span of the products $\prod_sf_s(x_s)$ with $f_s\in\nabla_{\lambda(s)}(M_s^\vee)$. By
Lemma~\ref{pd-kappa-cloud-cartan}(1) each $\nabla_{\lambda(s)}(M_s^\vee)$ has an $\mathcal R$-basis $(Z_{s,i})_i$ that
is a $\mathcal K$-basis of $\nabla_{\mathcal K}(\lambda(s))$; let $(Z^*_{s,i})_i$ be the basis with
$B_s(Z_{s,i},Z^*_{s,j})$ equal to $1$ for $i=j$ and $0$ otherwise. By
Lemma~\ref{pd-symbols-test-fock-position}(5), every $Z^*_{s,j}$ lies in $\nabla_{\lambda(s)}(M_s)$, and for
$Q\in\nabla_{\lambda(s)}(M_s)$ the coordinates $B_s(Z_{s,i},p^{N_s}Q)$ of $p^{N_s}Q$ in the basis $(Z^*_{s,i})$ lie
in $\mathcal R$. The products $z_I=\prod_sZ_{s,i_s}(x_s)$ form an $\mathcal R$-basis of $H(M)$ and a
$\mathcal K$-basis of $V$, so $H(M)^\vee$ is spanned by the coordinate forms $z_I^*$; and
$\Psi(z_I^*)=\prod_sZ^*_{s,i_s}(x_s)$, since $c\langle z_J,\prod_sZ^*_{s,i_s}\rangle_{\mathrm F}=\prod_s
B_s(Z_{s,j_s},Z^*_{s,i_s})$. Hence $\Psi(H(M)^\vee)$ is the $\mathcal R$-span of these products, which lies in
$H(M^\vee)$ and contains $p^{\kappa_1(\lambda_0)}H(M^\vee)$.

\emph{Identification.} Let $L$ be a $W$-lattice with basis the columns of $g$, and $g_s=\iota\sigma^s(g)$. The
columns of $g_s$ form an $\mathcal R$-basis of the lattice $M_s$ from which $H(L)$ is formed, $L^\vee$ has basis the
columns of ${}^tg^{-1}$, and $\iota\sigma^s({}^tg^{-1})={}^tg_s^{-1}$; so the lattices for $L^\vee$ are the
$M_s^\vee$, and $H(L)=H(M)$, $H(L^\vee)=H(M^\vee)$. Let $e'$ be an edge with source $L''$ and target $L$. As in
the proof of Lemma~\ref{pd-edge-basis}(1), $L''$ has a basis $b_1,\dots,b_k$ with
$p^{-1}b_1,\dots,p^{-1}b_r,b_{r+1},\dots,b_k$ a basis of $L$; with $g$ the matrix of the $b_i$ and $D$ the diagonal
matrix with entries $\varpi_0^{-1}$ in the first $r$ places and $1$ in the others, the lattices from which $H(e')$
is formed are spanned by the columns of $g_sD$. The edge $e'^\vee$ has source $L^\vee$, with basis the columns of
${}^tg^{-1}\operatorname{diag}(p,\dots,p,1,\dots,1)$, and target $L''^\vee$, with basis the columns of ${}^tg^{-1}$;
as $p=\varpi_0^2$, the lattices from which $H(e'^\vee)$ is formed are spanned by the columns of
${}^tg_s^{-1}D^{-1}={}^t(g_sD)^{-1}$, the duals of those of $e'$. So $H(e'^\vee)=H(M^\vee)$ for the family $M$ of
$e'$, and $\Psi(H(e')^\vee)\subset H(e'^\vee)$ and $p^{\kappa_1(\lambda_0)}H(e'^\vee)\subset\Psi(H(e')^\vee)$.
Multiplying by $\varpi_0^{-A}$ gives the claims for $R'(e')=\varpi_0^{-A}H(e')^\vee$ and
$R(e'^\vee)=\varpi_0^{-A}H(e'^\vee)$.
% Lean: test_side_explicit_identification: per position TS.formCoreSKX with posCore_of_formX (Lemma
%   pd-symbols-test-fock-position), then ten_sideX (psi_exists, prodBasis, latBasis, tensorBasis, dualB,
%   TS.fock_pairing_block, dualLat_span, lb_tensor; up and low of PosCore); kap1 = sum of expN (hkappa).
%   Identification: e1 and e2 (extO_colSpan, latDual_colSpan, tw_invT) for lattices, and for edges
%   edgeBasis, leAt_of_basis, dual_edgeBasis, tw_dE, invT_mul_D, dE_sq_D, then scale_map and scale_comm.
\end{proof}

\begin{definition}\label{pd-symbols-fock-words-operators-symbol}
Let $\mathcal K$ be a field of characteristic $0$, $k\geq1$ and $d\geq0$ integers, $\mathcal K[Y]_d$ the space of
homogeneous polynomials of degree $d$ in $Y_1,\dots,Y_k$, $E_{ab}=Y_a\,\partial/\partial Y_b$ on $\mathcal K[Y]_d$ for
$1\leq a,b\leq k$, and $\mathcal K[P,Q]$ the polynomial ring in $P_1,\dots,P_k,Q_1,\dots,Q_k$. For
$g\in\mathbf Z_{\geq0}^k$ put $|g|=\sum_ig_i$, $g!=\prod_ig_i!$, $Y^g=\prod_iY_i^{g_i}$, and $P^g$, $Q^g$ likewise.
\begin{enumerate}
\item The \emph{symbol} of $A\in\operatorname{End}_{\mathcal K}(\mathcal K[Y]_d)$ is
  \[
   \kappa(A)=\sum_{|g|=d}\frac1{g!}\,A(Y^g)(P)\,Q^g\in\mathcal K[P,Q],
  \]
  where $F(P)$ is obtained from $F\in\mathcal K[Y]$ by replacing each $Y_i$ by $P_i$.
\item For $1\leq a,b\leq k$, $R_{ab}=P_a\,\partial/\partial P_b-Q_b\,\partial/\partial Q_a$, a derivation of $\mathcal
  K[P,Q]$; and $\tau=\sum_iP_iQ_i$.
\item A polynomial in $\mathcal K[P,Q]$ is \emph{bihomogeneous of bidegree $(n,n)$} if each of its monomials has degree
  $n$ in $P_1,\dots,P_k$ and degree $n$ in $Q_1,\dots,Q_k$. The \emph{weight at $a$} of a monomial is its degree in
  $P_a$ minus its degree in $Q_a$.
\item For $c\geq0$, $\mathcal W_c$ is the span in $\operatorname{End}_{\mathcal K}(\mathcal K[Y]_d)$ of the products
  $E_{a_1b_1}\cdots E_{a_lb_l}$ with $l\leq c$, the empty product being the identity.
\end{enumerate}
\end{definition}
\leanref{Formalism/PeriodIntegrals/UnitaryGroupMassFormula/PolynomialDifferentialWords/KappaRight.lean}{PI.OpWords.kappa}
\leanref{Formalism/PeriodIntegrals/UnitaryGroupMassFormula/PolynomialDifferentialWords/KappaRight.lean}{PI.OpWords.Rop}
\leanref{Formalism/PeriodIntegrals/UnitaryGroupMassFormula/PolynomialDifferentialWords/KappaRight.lean}{PI.OpWords.tt}
\leanref{Formalism/PeriodIntegrals/UnitaryGroupMassFormula/PolynomialDifferentialWords/KappaRight.lean}{PI.OpWords.Bihom}
\leanref{Formalism/PeriodIntegrals/UnitaryGroupMassFormula/PolynomialDifferentialWords/KappaRight.lean}{PI.OpWords.WtLe}
\leanref{Formalism/PeriodIntegrals/UnitaryGroupMassFormula/PolynomialDifferentialWords/KappaRight.lean}{PI.OpWords.Wc}

\begin{lemma}\label{pd-symbols-fock-words-operators-kappa}
Let $\mathcal K$ be a field of characteristic $0$, $k\geq1$ and $d\geq0$ integers; we use the notation of
Definition~\ref{pd-symbols-fock-words-operators-symbol}. Then:
\begin{enumerate}
\item for $A\in\operatorname{End}_{\mathcal K}(\mathcal K[Y]_d)$ and $m,g\in\mathbf Z_{\geq0}^k$, the coefficient of $P^mQ^g$
  in $\kappa(A)$ is the coefficient of $Y^m$ in $A(Y^g)$ divided by $g!$ if $|g|=d$, and $0$ otherwise; the map $\kappa$
  is linear and injective, and $\kappa(A)$ is bihomogeneous of bidegree $(d,d)$;
\item $\kappa(E_{ab}A)=P_a\,\partial\kappa(A)/\partial P_b$ and $\kappa(AE_{ab})=Q_b\,\partial\kappa(A)/\partial Q_a$; hence
  $\kappa(E_{ab}A-AE_{ab})=R_{ab}\kappa(A)$;
\item $\kappa(1)=\tau^d/d!$.
\end{enumerate}
\end{lemma}
\leanref{Formalism/PeriodIntegrals/UnitaryGroupMassFormula/PolynomialDifferentialWords/KappaCoeff.lean}{PI.OpWords.coeff_kappa}
\leanref{Formalism/PeriodIntegrals/UnitaryGroupMassFormula/PolynomialDifferentialWords/KappaMain.lean}{PI.OpWords.symbolMapInjective}
\leanref{Formalism/PeriodIntegrals/UnitaryGroupMassFormula/PolynomialDifferentialWords/KappaMain.lean}{PI.OpWords.symbolMapBihomogeneous}
\leanref{Formalism/PeriodIntegrals/UnitaryGroupMassFormula/PolynomialDifferentialWords/KappaMain.lean}{PI.OpWords.symbolMapMulLeft}
\leanref{Formalism/PeriodIntegrals/UnitaryGroupMassFormula/PolynomialDifferentialWords/KappaMain.lean}{PI.OpWords.symbolMapMulRight}
\leanref{Formalism/PeriodIntegrals/UnitaryGroupMassFormula/PolynomialDifferentialWords/KappaMain.lean}{PI.OpWords.symbolMapOne}

\begin{proof}
(1) The polynomial $A(Y^g)(P)Q^g$ has coefficient at $P^mQ^{g'}$ equal to the coefficient of $Y^m$ in $A(Y^g)$ if
$g'=g$, and $0$ otherwise; this gives the coefficients, and $\kappa$ is linear. If $\kappa(A)=\kappa(B)$, the coefficients
give $A(Y^g)=B(Y^g)$ for every $g$ with $|g|=d$; these monomials span $\mathcal K[Y]_d$, so $A=B$. If the coefficient of
$P^mQ^g$ is nonzero, then $|g|=d$, and $|m|=d$ as $A(Y^g)$ is homogeneous of degree $d$.

(2) The operator $P_a\,\partial/\partial P_b$ kills $Q^g$ and acts on $F(P)$ as $E_{ab}$ acts on $F$; this gives the first
identity. For the second, compare the coefficients at $P^mQ^g$ with $|g|=d$. As $E_{ab}Y^g=g_bY^{g'}$ with $g'=g-e_b+e_a$
(and $0$ if $g_b=0$), the coefficient in $\kappa(AE_{ab})$ is $g_b\,c_m(g')/g!$, where $c_m(h)$ is the coefficient of
$Y^m$ in $A(Y^h)$. In $Q_b\,\partial\kappa(A)/\partial Q_a$ it is $0$ if $g_b=0$, and otherwise $g'_a$ times the
coefficient of $P^mQ^{g'}$ in $\kappa(A)$, that is $g'_a\,c_m(g')/g'!$. As $g'!=(g-e_b)!\,g'_a$ and $g!=(g-e_b)!\,g_b$,
the two agree. The last identity is the difference of the first two.

(3) By the multinomial theorem the coefficient of $P^mQ^g$ in $\tau^d=(\sum_iP_iQ_i)^d$ is $d!/g!$ if $m=g$ and $|g|=d$,
and $0$ otherwise. By (1) the coefficient in $\kappa(1)$ is $1/g!$ if $m=g$ and $|g|=d$, and $0$ otherwise.
% Lean: (1) is coeff_kappa (kappa_apply, coeff_term, mono_val), symbolMapInjective (eq_sum_mono) and
%   symbolMapBihomogeneous; (2) is symbolMapMulLeft and symbolMapMulRight (Ed_mono, invFact_add_single), and the last
%   identity is kappa_ad in DivideAll.lean; (3) is symbolMapOne (coeff_sumElim_tt_pow, Nat.multinomial_spec).
\end{proof}

\begin{lemma}\label{pd-symbols-fock-words-operators-tau}
Let $\mathcal K$ be a field of characteristic $0$ and $k\geq1$ an integer; we use the notation of
Definition~\ref{pd-symbols-fock-words-operators-symbol}. Let $1\leq a,b\leq k$, $n\geq0$ and $c$ be integers. Then:
\begin{enumerate}
\item $R_{ab}(\tau g)=\tau R_{ab}(g)$ for $g\in\mathcal K[P,Q]$, and $R_{ab}$ maps bihomogeneous polynomials of bidegree
  $(n,n)$ to bihomogeneous polynomials of bidegree $(n,n)$;
\item if $f\in\mathcal K[P,Q]$ and $w\in\mathbf Z$ satisfy $R_{aa}f=wf$, every monomial of $f$ has weight $w$ at $a$;
\item if every monomial of $\tau g$ has weight at most $c$ at every $a'\leq k$, so does every monomial of $g$;
\item if $n\geq1$, $f$ is bihomogeneous of bidegree $(n,n)$ and $f(p,q)=0$ for all $p,q\in\mathcal K^k$ with
  $\sum_ip_iq_i=0$, then $f=\tau g$ with $g$ bihomogeneous of bidegree $(n-1,n-1)$.
\end{enumerate}
\end{lemma}
\leanref{Formalism/PeriodIntegrals/UnitaryGroupMassFormula/PolynomialDifferentialWords/PolarizationBihomogeneous.lean}{PI.OpWords.rttRelation}
\leanref{Formalism/PeriodIntegrals/UnitaryGroupMassFormula/PolynomialDifferentialWords/PolarizationBihomogeneous.lean}{PI.OpWords.polarizationBihomogeneous}
\leanref{Formalism/PeriodIntegrals/UnitaryGroupMassFormula/PolynomialDifferentialWords/PolarizationBihomogeneous.lean}{PI.OpWords.weightOfEigenvector}
\leanref{Formalism/PeriodIntegrals/UnitaryGroupMassFormula/PolynomialDifferentialWords/PolarizationBihomogeneous.lean}{PI.OpWords.pairingWeightCancel}
\leanref{Formalism/PeriodIntegrals/UnitaryGroupMassFormula/PolynomialDifferentialWords/IsWeightedHomogeneousPairingLeft.lean}{PI.OpWords.pairingDivisibility}

\begin{proof}
(1) As $\partial\tau/\partial P_b=Q_b$ and $\partial\tau/\partial Q_a=P_a$, $R_{ab}\tau=P_aQ_b-Q_bP_a=0$, and $R_{ab}$ is a
derivation. The operators $P_a\,\partial/\partial P_b$ and $Q_b\,\partial/\partial Q_a$ keep the degree of a monomial in the
$P_i$ and in the $Q_i$.

(2) For a monomial $\mu$, $P_a\,\partial\mu/\partial P_a$ and $Q_a\,\partial\mu/\partial Q_a$ are $\mu$ times its degrees in
$P_a$ and in $Q_a$; so $R_{aa}\mu$ is $\mu$ times its weight at $a$. Comparing coefficients in $R_{aa}f=wf$, each monomial
of $f$ has weight $w$ at $a$, as the characteristic is $0$.

(3) Fix $a'$, and for $n'\in\mathbf Z$ let $g_{n'}$ be the sum of the terms of $g$ whose monomials have weight $n'$ at $a'$.
Every monomial $P_iQ_i$ of $\tau$ has weight $0$ at $a'$, so the sum of the terms of $\tau g$ of weight $n'$ at $a'$ is
$\tau g_{n'}$. If $g$ had a monomial of weight $n'>c$ at $a'$, then $g_{n'}\neq0$, so $\tau g_{n'}\neq0$, as $\tau\neq0$ and
$\mathcal K[P,Q]$ is a domain, and $\tau g$ would have a monomial of weight $n'$ at $a'$.

(4) First, $\tau$ divides $f$. If $k=1$, every monomial of $f$ is a multiple of $P_1^nQ_1^n$, hence of $\tau=P_1Q_1$. Let
$k\geq2$, and fix $i_0\neq j$. Let $\chi$ and $\psi$ be the $\mathcal K$-algebra endomorphisms of $\mathcal K[P,Q]$ that fix
every $P_i$, with $\chi(Q_i)=P_{i_0}Q_i$ for all $i$, and $\psi(Q_i)=P_{i_0}Q_i$ for $i\neq i_0$ and
$\psi(Q_{i_0})=-\sum_{i\neq i_0}P_iQ_i$. They agree on every variable except $Q_{i_0}$, where
$\chi(Q_{i_0})-\psi(Q_{i_0})=\tau$; so $\chi(h)-\psi(h)\in\tau\mathcal K[P,Q]$ for every $h$. As $f$ has degree $n$ in the
$Q_i$, $\chi(f)=P_{i_0}^nf$. For $p,q\in\mathcal K^k$, $\psi(f)(p,q)=f(p,q')$ with $q'_i=p_{i_0}q_i$ for $i\neq i_0$ and
$q'_{i_0}=-\sum_{i\neq i_0}p_iq_i$, and $\sum_ip_iq'_i=0$; so $\psi(f)$ vanishes at every point, and $\psi(f)=0$ as
$\mathcal K$ is infinite. Hence $\tau$ divides $P_{i_0}^nf$. The prime element $P_{i_0}$ does not divide $\tau$, which does not
vanish at $p=e_j$, $q=(1,\dots,1)$; so $\tau$ divides $f$, one factor $P_{i_0}$ at a time.

Write $f=\tau g$. As $\tau$ has degree $1$ in the $P_i$ and $1$ in the $Q_i$, the part of $\tau g$ of degree $n'+1$ in the
$P_i$ is $\tau$ times the part of $g$ of degree $n'$, and likewise for the $Q_i$; as $\tau\neq0$, $g$ is bihomogeneous of
bidegree $(n-1,n-1)$.
% Lean: (1) is rttRelation and polarizationBihomogeneous (iwh_X_pderiv); (2) is weightOfEigenvector
%   (coeff_X_pderiv_self); (3) is pairingWeightCancel (comp_mul_zero with tt_va); (4) is pairingDivisibility = divT
%   (divT_one for k = 1; for k >= 2 chi_psi_dvd, psi_zero, chi_eq, dvd_pow_mul, not_X_dvd_tt and dvd_of_dvd_X_pow; the
%   degrees by homog_of_mul). The Lean also treats k = 0.
\end{proof}

\begin{lemma}\label{pd-symbols-fock-words-operators-exponential}
Let $\mathcal K$ be a field of characteristic $0$ and $k\geq1$ an integer; we use the notation of
Definition~\ref{pd-symbols-fock-words-operators-symbol}. Let $a\neq b$ and $s\in\mathcal K$, and for $p,q\in\mathcal K^k$ let
$\sigma_s(p,q)=(p',q')$ with $p'_b=p_b+sp_a$, $q'_a=q_a-sq_b$, and all other coordinates unchanged. Then for every
$g\in\mathcal K[P,Q]$ there is $N\geq0$ with $R_{ab}^ng=0$ for $n>N$, and
\[
 g(\sigma_s(p,q))=\sum_{n=0}^N\frac{s^n}{n!}\,(R_{ab}^ng)(p,q)\qquad\text{for all }p,q\in\mathcal K^k.
\]
\end{lemma}
\leanref{Formalism/PeriodIntegrals/UnitaryGroupMassFormula/PolynomialDifferentialWords/VanishExp.lean}{PI.OpWords.expFormula}

\begin{proof}
Write $R=R_{ab}$. For a variable $x$, $Rx$ is $P_a$ if $x=P_b$, $-Q_b$ if $x=Q_a$, and $0$ otherwise; as $a\neq b$,
$R(Rx)=0$, and $\sigma_s(p,q)$ is the point with coordinates $x(p,q)+s(Rx)(p,q)$. As $R$ is a derivation, induction on
$n$ gives $R^{n+1}(gx)=R^{n+1}(g)x+(n+1)R^n(g)R(x)$ for every $g$ and every variable $x$. The claim holds for constants,
with $N=0$, and for a sum when it holds for the summands, with the larger $N$. Let it hold for $g$ with $N$, and let $x$
be a variable. Then $R^n(gx)=0$ for $n>N+1$, and
\[
 (gx)(\sigma_s(p,q))=\sum_n\frac{s^n}{n!}(R^ng)(p,q)\big(x+sRx\big)(p,q)
 =\sum_n\frac{s^n}{n!}\big(R^n(g)x+nR^{n-1}(g)R(x)\big)(p,q)=\sum_n\frac{s^n}{n!}\big(R^n(gx)\big)(p,q),
\]
where the term with $R^{n-1}$ is $0$ for $n=0$. Every polynomial is built from constants by sums and products with
variables.
% Lean: expFormula (ExpFormula with shiftPt; rop_rop_X, rop_iter_mul_X and induction_on).
\end{proof}

\begin{lemma}\label{pd-symbols-fock-words-operators-quadric}
Let $\mathcal K$ be a field of characteristic $0$, and $k\geq1$ and $0\leq c<n$ integers; we use the notation of
Definition~\ref{pd-symbols-fock-words-operators-symbol}. Let $V$ be a subspace of the bihomogeneous polynomials of
bidegree $(n,n)$ in $\mathcal K[P,Q]$ with $R_{ab}V\subset V$ for all $a,b$, such that every monomial of every element of
$V$ has weight at most $c$ at every $a$. Then $f(p,q)=0$ for every $f\in V$ and all $p,q\in\mathcal K^k$ with
$\sum_ip_iq_i=0$.
\end{lemma}
\leanref{Formalism/PeriodIntegrals/UnitaryGroupMassFormula/PolynomialDifferentialWords/VanishMain.lean}{PI.OpWords.vanishingOnQuadric}
\leanref{Formalism/PeriodIntegrals/UnitaryGroupMassFormula/PolynomialDifferentialWords/VanishMain.lean}{PI.OpWords.vanish_at}

\begin{proof}
Let $f\in V$ and $\sum_ip_iq_i=0$. If $p=0$, then $f(p,q)=0$, since every monomial of $f$ has degree $n\geq1$ in the
$P_i$. Otherwise fix $i_0$ with $p_{i_0}\neq0$. We show that every $g\in V$ vanishes at every point $(p,q)$ of the quadric
with $p_{i_0}\neq0$, by induction on the number of $b\neq i_0$ with $p_b\neq0$.

If there is none, the quadric gives $p_{i_0}q_{i_0}=0$, so $q_{i_0}=0$. A monomial of $g$ that does not vanish at $(p,q)$
involves no $P_b$ with $b\neq i_0$ and not $Q_{i_0}$; as it has degree $n$ in the $P_i$, its weight at $i_0$ is $n>c$,
which is impossible. So $g(p,q)=0$.

Otherwise let $b\neq i_0$ with $p_b\neq0$, put $s=-p_b/p_{i_0}$, and let $(p',q')=\sigma_s(p,q)$ as in
Lemma~\ref{pd-symbols-fock-words-operators-exponential} for $(a,b)=(i_0,b)$. Then $p'_b=0$, $p'_i=p_i$ for $i\neq b$, and
\[
 \sum_ip'_iq'_i=\sum_ip_iq_i-sp_{i_0}q_b+sp_{i_0}q_b=0 ,
\]
so by induction every element of $V$ vanishes at $(p',q')$. Also $\sigma_{-s}(p',q')=(p,q)$. By
Lemma~\ref{pd-symbols-fock-words-operators-exponential}, $g(p,q)=\sum_n\frac{(-s)^n}{n!}(R_{i_0b}^ng)(p',q')=0$, as each
$R_{i_0b}^ng$ lies in $V$.
% Lean: vanishingOnQuadric (the case p = 0 through xdeg_of_bihom) with vanish_at, an induction on a finite set S
%   containing the b /= i0 with p_b /= 0: eval_base for the empty set, and shiftPt with expFormula, hquad and hback for
%   the step.
\end{proof}

\begin{lemma}\label{pd-symbols-fock-words-operators-words}
Let $\mathcal K$ be a field of characteristic $0$, and $k\geq1$ and $0\leq c\leq d$ integers; we use the notation of
Definition~\ref{pd-symbols-fock-words-operators-symbol}. If $g\in\mathcal K[P,Q]$ is bihomogeneous of bidegree $(c,c)$,
then $\tau^{d-c}g=\kappa(B)$ for some $B\in\mathcal W_c$.
\end{lemma}
\leanref{Formalism/PeriodIntegrals/UnitaryGroupMassFormula/PolynomialDifferentialWords/PairingPowMemSymbolRange.lean}{PI.OpWords.pairingPowMemSymbolRange}
\leanref{Formalism/PeriodIntegrals/UnitaryGroupMassFormula/PolynomialDifferentialWords/PairingPowMemSymbolRange.lean}{PI.OpWords.bihom_le_Sc}

\begin{proof}
By induction on $c$. For $c=0$, $g$ is a constant $\lambda$, and $\tau^dg=\kappa(\lambda d!\cdot1)$ by
Lemma~\ref{pd-symbols-fock-words-operators-kappa}(3). Let $c+1\leq d$ and put $n=d-c-1$. By linearity it suffices to treat a
monomial $g$ of bidegree $(c+1,c+1)$. Choose $a$ and $b$ with $P_a$ and $Q_b$ dividing $g$, and write $g=P_aQ_bg'$; then $g'$
and $P_a\,\partial g'/\partial P_b$ are bihomogeneous of bidegree $(c,c)$. By induction there are $B,B'\in\mathcal W_c$ with
$\kappa(B)=\tau^{n+1}g'$ and $\kappa(B')=\tau^{n+1}P_a\,\partial g'/\partial P_b$. By
Lemma~\ref{pd-symbols-fock-words-operators-kappa}(2), and as $\partial\tau/\partial P_b=Q_b$,
\[
 \kappa(E_{ab}B-B')=P_a\frac{\partial(\tau^{n+1}g')}{\partial P_b}-\tau^{n+1}P_a\frac{\partial g'}{\partial P_b}
 =(n+1)\tau^nP_aQ_bg'=(n+1)\tau^ng .
\]
So $\tau^{d-c-1}g=\kappa(B'')$ with $B''=(E_{ab}B-B')/(n+1)$, which lies in $\mathcal W_{c+1}$ since $E_{ab}\mathcal W_c\subset
\mathcal W_{c+1}$ and $\mathcal W_c\subset\mathcal W_{c+1}$.
% Lean: pairingPowMemSymbolRange = bihom_le_Sc (bihom_zero_const and symbolMapOne for c = 0; Ed_mul_Wc, Wc_mono,
%   symbolMapMulLeft and pderiv_inl_tt for the step).
\end{proof}

\begin{lemma}\label{pd-symbols-fock-words-operators}
Let $\mathcal K$ be a field of characteristic $0$, $k\geq1$ and $d,c\geq0$ integers, $\mathcal K[Y]_d$ the space of
homogeneous polynomials of degree $d$ in $Y_1,\dots,Y_k$, and $E_{ab}=Y_a\,\partial/\partial Y_b$ on $\mathcal K[Y]_d$
for $1\leq a,b\leq k$. Let $\mathcal W$ be a subspace of $\operatorname{End}_{\mathcal K}(\mathcal K[Y]_d)$ such that
$[E_{ab},A]\in\mathcal W$ for all $A\in\mathcal W$ and all $a,b$, and such that $\mathcal W$ is spanned by elements
$A$ for each of which there are integers $\omega(1),\dots,\omega(k)\leq c$ with $[E_{aa},A]=\omega(a)A$ for every
$a$. Then every element of $\mathcal W$ is a linear combination of products $E_{a_1b_1}\cdots E_{a_lb_l}$ with
$l\leq c$.
\end{lemma}
\leanref{Formalism/SharpRestriction/GrowthBounds/DualFockLinear/OpBound.lean}{Formalism.MainNotions.TestSideCounting.OpBound}
\leanref{Formalism/PeriodIntegrals/OrbitalBoundsConclusion/OpWords.lean}{operator_words_bound}
\leanref{Formalism/PeriodIntegrals/UnitaryGroupMassFormula/PolynomialDifferentialWords/Assemble.lean}{PI.OpWords.opBound_of_vanish}

\begin{proof}
We use the notation of Definition~\ref{pd-symbols-fock-words-operators-symbol}, and write $[E,A]=EA-AE$.

\emph{Step 1: the symbols.} Let $V=\kappa(\mathcal W)\subset\mathcal K[P,Q]$. By
Lemma~\ref{pd-symbols-fock-words-operators-kappa}(1), every element of $V$ is bihomogeneous of bidegree $(d,d)$, and by
part (2) of that lemma, $R_{ab}\kappa(A)=\kappa([E_{ab},A])\in V$ for $A\in\mathcal W$; so $R_{ab}V\subset V$. Let $A$ be one
of the spanning elements of the statement. Then $R_{aa}\kappa(A)=\kappa([E_{aa},A])=\omega(a)\kappa(A)$, so every monomial of
$\kappa(A)$ has weight $\omega(a)\leq c$ at $a$, by Lemma~\ref{pd-symbols-fock-words-operators-tau}(2). As $V$ is spanned by
these $\kappa(A)$, every monomial of every element of $V$ has weight at most $c$ at every $a$.

\emph{Step 2: division by $\tau$.} We show, by induction on $n\geq0$: if $V'$ is a subspace of the bihomogeneous
polynomials of bidegree $(c+n,c+n)$ with $R_{ab}V'\subset V'$ for all $a,b$, such that every monomial of every element of
$V'$ has weight at most $c$ at every $a$, then every $f\in V'$ is $\tau^ng$ with $g$ bihomogeneous of bidegree $(c,c)$. For
$n=0$ take $g=f$. For $n+1$, let $f\in V'$. As $c<c+n+1$, Lemma~\ref{pd-symbols-fock-words-operators-quadric} shows that
$f$ vanishes on the quadric $\sum_ip_iq_i=0$, and Lemma~\ref{pd-symbols-fock-words-operators-tau}(4) gives $f=\tau g$ with
$g$ bihomogeneous of bidegree $(c+n,c+n)$. Let $V''$ be the subspace of the $g$ with $\tau g\in V'$. As $\tau\neq0$ and
$\mathcal K[P,Q]$ is a domain, $V''$ consists of bihomogeneous polynomials of bidegree $(c+n,c+n)$; it satisfies
$R_{ab}V''\subset V''$ by part (1) of that lemma, and the bound on the weights by its part (3). By induction $g=\tau^ng''$
with $g''$ bihomogeneous of bidegree $(c,c)$, so $f=\tau^{n+1}g''$.

\emph{Step 3: the words.} Let $A\in\mathcal W$. If $c<d$, Steps 1 and 2 with $n=d-c$ give $\kappa(A)=\tau^{d-c}g$ with $g$
bihomogeneous of bidegree $(c,c)$; put $c'=c$. If $c\geq d$, put $c'=d$ and $g=\kappa(A)$, which is bihomogeneous of
bidegree $(d,d)$. In both cases $c'\leq\min(c,d)$ and $\kappa(A)=\tau^{d-c'}g$. By
Lemma~\ref{pd-symbols-fock-words-operators-words}, $\tau^{d-c'}g=\kappa(B)$ for some $B\in\mathcal W_{c'}$. As $\kappa$ is
injective, $A=B\in\mathcal W_{c'}\subset\mathcal W_c$.
% Lean: operator_words_bound = PI.OpWords.opBound_of_vanish PI.OpWords.vanishingOnQuadric = opWords_of
%   symbolMapInjective symbolMapMulLeft symbolMapMulRight symbolMapBihomogeneous weightOfEigenvector
%   polarizationBihomogeneous rttRelation pairingWeightCancel pairingDivisibility vanishingOnQuadric
%   pairingPowMemSymbolRange (DivideAll.lean). Step 1 is kappa_ad, hWst and hWwt (wtLe_span); Step 2 is divide_all,
%   with Vanish (Lemma pd-symbols-fock-words-operators-quadric), DivT, RTT and TTWt (Lemma
%   pd-symbols-fock-words-operators-tau); Step 3 is key in opWords_of, then Surj (Lemma
%   pd-symbols-fock-words-operators-words), KappaInj and Wc K k d c' <= Wc K k d c. The input RBihom of divide_all
%   is not used in its proof.
\end{proof}

\begin{lemma}\label{pd-symbols-fock-words-position}
Let $\mathcal K$ be a field of characteristic $0$, $k\geq2$ and $d,c\geq0$ integers, $x$ a $k\times k$ matrix of
variables, $\nu\in\mathbf Z_{\geq0}^k$ nonincreasing with $\nu_1\leq c$, and $\nabla_{\mathcal K}(\nu)\subset\mathcal K[x]$
the dual Weyl module of Definition~\ref{pd-weight-lattices}(1). Fix a parity, even or odd, let
$\theta_1\colon\mathcal K[Y]\to\mathcal K[x]$ be the map of Lemma~\ref{pd-symbols-flat-letter} for it, and let
$\mathcal K[Y]_d$ and the $E_{ab}$ be as in Lemma~\ref{pd-symbols-fock-words-operators}. Let
$\beta,s\in\nabla_{\mathcal K}(\nu)$. Then there are finitely many products $D_l$ of at most $c$ of the operators
$E_{ab}$, and $a_l\in\mathcal K$, such that
\[
 \langle\theta_1(P)s,\theta_1(Q)\beta\rangle_{\mathrm F}=\Big\langle\theta_1\Big(\sum_la_lD_lP\Big),\theta_1(Q)\Big\rangle_{\mathrm F}
 \qquad\text{for all }P,Q\in\mathcal K[Y]_d .
\]
\end{lemma}
\leanref{Formalism/SharpRestriction/GrowthBounds/DualFockLinear/OnePosDef.lean}{Formalism.MainNotions.TestSideCounting.OnePos}
\leanref{Formalism/SharpRestriction/GrowthBounds/DualFockLinear/OnePosOf.lean}{Formalism.MainNotions.TestSideCounting.onePos_of}
\leanref{Formalism/SharpRestriction/GrowthBounds/DualFockLinear/FRSmulRight.lean}{Formalism.MainNotions.TestSideCounting.exists_A}
\leanref{Formalism/SharpRestriction/GrowthBounds/DualFockLinear/LetterRowOperator.lean}{Formalism.MainNotions.TestSideCounting.comm_A}
\leanref{Formalism/PeriodIntegrals/OrbitalBoundsConclusion/RowLetter.lean}{row_letter_rule}

\begin{proof}
Write $\langle\ ,\ \rangle$ for the Fock pairing on $\mathcal K[x]$, which is symmetric, and for $1\leq a,b\leq k$ let
$R_{ab}=\sum_tx_{a,t}\,\partial/\partial x_{b,t}$, a derivation of $\mathcal K[x]$; it is the operator
$E^{(1)}_{ab}$ of Definition~\ref{pd-weight-lattices}(1). Let $\nu'$ be $(d,0,\dots,0)$ for even and
$(d,\dots,d,0)$ for odd parity. The proofs of Lemma~\ref{pd-symbols-flat-letter}(1) and (2) use of $\mathcal K$ only
that it is a field; so $\theta_1$ maps $\mathcal K[Y]_d$ onto $\nabla_{\mathcal K}(\nu')$, and it is injective on
$\mathcal K[Y]_d$, as $\psi(\theta_1(P))=q^dP$ with $q\neq0$. Let $L_{ab}$ be $E_{ab}$ for even parity and
$\delta_{ab}d-E_{ba}$ for odd parity, on $\mathcal K[Y]_d$. Five facts:
\begin{enumerate}
\item[(i)] If $P\in\mathcal K[Y]_d$ and $\langle\theta_1(P),\theta_1(Q)\rangle=0$ for every $Q\in\mathcal K[Y]_d$, then
  $P=0$; and every linear form on $\mathcal K[Y]_d$ is $Q\mapsto\langle\theta_1(P),\theta_1(Q)\rangle$ for some $P$.
  Indeed $\theta_1(P)\in\nabla_{\mathcal K}(\nu')$ pairs to $0$ with all of $\nabla_{\mathcal K}(\nu')$, so it is $0$ by
  part (2) of the proof of Lemma~\ref{pd-kappa-gram-form}, which uses of the field only that it has characteristic
  $0$ (as in the proof of Lemma~\ref{pd-symbols-test-fock-position}); and $P=0$ by injectivity. The second claim
  follows, as $\mathcal K[Y]_d$ has finite dimension.
\item[(ii)] $\langle R_{ab}f,f'\rangle=\langle f,R_{ba}f'\rangle$ for $f,f'\in\mathcal K[x]$: this is
  Lemma~\ref{pd-kappa-gram-fock}(2) with $k=1$.
\item[(iii)] $R_{ab}$ maps $\nabla_{\mathcal K}(\nu)$ into itself: it does not change the degree in any column, and
  $R_{ab}f$ is the coefficient of $y$ in $f((1+yE_{ba})x)$, so $R_{ab}$ commutes with the substitutions
  $x\mapsto x(1+yE_{lj})$, which multiply on the right.
\item[(iv)] $R_{ab}\theta_1(P)=\theta_1(L_{ab}P)$ for $P\in\mathcal K[Y]_d$. As $R_{ab}$ is a derivation,
  $R_{ab}\theta_1(P)=\sum_c\theta_1(\partial P/\partial Y_c)\,R_{ab}\theta_1(Y_c)$. For even parity
  $\theta_1(Y_c)=x_{c,1}$ and $R_{ab}x_{c,1}=\delta_{bc}x_{a,1}$, which gives $\theta_1(Y_a\partial P/\partial Y_b)$.
  For odd parity $\theta_1(Y_c)$ is the determinant $N_c$ of $x$ with row $c$ replaced by the last unit vector, and
  $R_{ab}N_c$ is $0$ if $b=c$ and otherwise the determinant with, in addition, row $b$ replaced by row $a$; so
  $R_{ab}N_c=\delta_{ab}N_c-\delta_{ac}N_b$ (for $a=c\neq b$ exchange the rows $b$ and $c$). With Euler's identity $\sum_cY_c\,\partial P/\partial
  Y_c=dP$ this gives $\theta_1(\delta_{ab}dP-Y_b\,\partial P/\partial Y_a)$.
\item[(v)] $\nabla_{\mathcal K}(\nu)$ is spanned by elements $b_i$ with $R_{aa}b_i=\mu_i(a)b_i$ and
  $0\leq\mu_i(a)\leq\nu_1$ for all $a$: by Lemma~\ref{is-standard-basis}(1) it is spanned by bideterminants, products
  of $\nu_1$ minors of $x$ on the first columns, and such a product is homogeneous of some degree $\mu(a)$ in the
  row $a$, which $R_{aa}$ multiplies by $\mu(a)$; each minor has degree at most $1$ in row $a$.
\end{enumerate}

\emph{Step 1: the operators $A(f,f')$.} By (i), for $f,f'\in\mathcal K[x]$ there is a unique endomorphism $A(f,f')$
of $\mathcal K[Y]_d$ with
\[
 \langle\theta_1(A(f,f')P),\theta_1(Q)\rangle=\langle\theta_1(P)f',\theta_1(Q)f\rangle\qquad(P,Q\in\mathcal K[Y]_d),
\]
and $A(f,f')$ is bilinear in $(f,f')$ by the uniqueness.

\emph{Step 2: commutators.} We show $[L_{ab},A(f,f')]=A(f,R_{ab}f')-A(R_{ba}f,f')$. By (iv), (ii) and the
derivation rule $(R_{ba}\theta_1(Q))f=R_{ba}(\theta_1(Q)f)-\theta_1(Q)R_{ba}f$,
\[
 \langle\theta_1(L_{ab}A(f,f')P),\theta_1(Q)\rangle=\langle\theta_1(A(f,f')P),R_{ba}\theta_1(Q)\rangle
 =\langle R_{ab}(\theta_1(P)f'),\theta_1(Q)f\rangle-\langle\theta_1(P)f',\theta_1(Q)R_{ba}f\rangle,
\]
where the middle term was rewritten with $R_{ba}\theta_1(Q)=\theta_1(L_{ba}Q)$ and Step~1. Also
$\langle\theta_1(A(f,f')L_{ab}P),\theta_1(Q)\rangle=\langle(R_{ab}\theta_1(P))f',\theta_1(Q)f\rangle$. As
$R_{ab}(\theta_1(P)f')-(R_{ab}\theta_1(P))f'=\theta_1(P)R_{ab}f'$, the difference of the two is
$\langle\theta_1((A(f,R_{ab}f')-A(R_{ba}f,f'))P),\theta_1(Q)\rangle$, and (i) gives the identity.

\emph{Step 3: the space $\mathcal W$.} Let $\mathcal W$ be the span of the $A(f,f')$ with
$f,f'\in\nabla_{\mathcal K}(\nu)$; by (v) and bilinearity it is spanned by the $A(b_i,b_j)$. For even parity
$E_{ab}=L_{ab}$, and for odd parity $[E_{ab},A]=-[L_{ba},A]$. So by Step~2 and (iii), $[E_{ab},A(f,f')]\in\mathcal W$
for $f,f'\in\nabla_{\mathcal K}(\nu)$, hence $[E_{ab},A]\in\mathcal W$ for every $A\in\mathcal W$. By Step~2 and (v),
$[E_{aa},A(b_i,b_j)]=\omega(a)A(b_i,b_j)$ with $\omega(a)=\mu_j(a)-\mu_i(a)$ for even and $\mu_i(a)-\mu_j(a)$ for odd
parity, and $\omega(a)\leq\nu_1\leq c$.

\emph{Step 4.} By Lemma~\ref{pd-symbols-fock-words-operators}, $A(\beta,s)=\sum_la_lD_l$ with products $D_l$ of at
most $c$ operators $E_{ab}$. By Step~1 with $(f,f')=(\beta,s)$,
$\langle\theta_1(P)s,\theta_1(Q)\beta\rangle=\langle\theta_1(\sum_la_lD_lP),\theta_1(Q)\rangle$.
% Lean: onePos_of_opBound = onePos_of hOB restAdj restStab letter_operators_holds dual_weyl_weights_bounded restRep
%   restRep2 (OnePosOf.lean). (i) is restRep (theta_inj_Hd with psi_Nv and scaleHd, Top.nondegK, theta_Hd) and
%   restRep2 (cForm bijective); (ii) is restAdj (Top.fockRRowDivAdj); (iii) is restStab (rowDivStable); (iv) is
%   RestTheta = row_letter_rule (rowLetter_of chainRule vecX keyCov hdHom, with euler_theta); (v) is RestWt =
%   rowWtSpan_of bidetMem bideterminant_span_holds (standardMonomialSpanDomain). Step 1 is exists_A and uniqV with
%   the bilinear A (LinearMap.mk2); Step 2 is comm_A (letter_operator_matches_row_operator, RE_mul) with
%   letterRowOperator = L and Ed_comm_eq; Step 3 is WW = map2, hWWspan, hstab, hwt; Step 4 is hOB with
%   Finsupp.mem_span_image_iff_linearCombination.
\end{proof}

\begin{lemma}\label{pd-symbols-test-fock-words}
Let $\mathcal K$ be a field of characteristic $0$, $n'\geq1$ and $k\geq2$ integers, $T\subset\mathbf Z/n'\mathbf Z$,
$d\colon\mathbf Z/n'\mathbf Z\to\mathbf Z_{\geq0}$, $c\geq0$ an integer, and $\nu(s)\in\mathbf Z_{\geq0}^k$
nonincreasing, with $\nu(t)_1\leq c$ for $t\in T$ and $\nu(s)=0$ for $s\notin T$. Let $\Theta_1$ be the map
$\Theta_\gamma$ of Lemma~\ref{pd-symbols-flat-exists} for $\gamma=1$, and let $\beta$ and $s$ lie in the span of the
products $\prod_tf_t(x_t)$ with $f_t\in\nabla_{\mathcal K}(\nu(t))$ (Definition~\ref{pd-weight-lattices}(1)).
Then there are finitely many products $D_l$ of operators $X_{t,a}\,\partial/\partial X_{t,b}$ with $t\in T$, in each
of which every $t$ occurs at most $c$ times, and $a_l\in\mathcal K$, such that for every $G\in\mathcal K[X]_d$ the
polynomial $E=\sum_la_lD_lG$ lies in $\mathcal K[X]_d$ and
\[
 \langle\Theta_1(G)s,\Theta_1(H)\beta\rangle_{\mathrm F}=\langle\Theta_1(E),\Theta_1(H)\rangle_{\mathrm F}
 \qquad\text{for every }H\in\mathcal K[X]_d .
\]
\end{lemma}
\leanref{Formalism/SharpRestriction/GrowthBounds/DualFockLinear/FirstWeightLeConductor.lean}{Formalism.MainNotions.TestSideCounting.TenOp}
\leanref{Formalism/SharpRestriction/GrowthBounds/TwistedSubstitution/OperatorWeightBoundHolds.lean}{Formalism.MainNotions.TestSideWords.tensor_operators_holds}

\begin{proof}
For a set $\mathcal L$ of products of operators $X_{t,a}\,\partial/\partial X_{t,b}$ and coefficients $a_l$, say that
$(\mathcal L,a)$ \emph{works} for $(\beta,s)$ if each $D_l\in\mathcal L$ involves only positions $t\in T$, each at most
$c$ times, and the identity of the statement holds for all $G,H\in\mathcal K[X]_d$ with $E=\sum_la_lD_lG$. Each
operator $X_{t,a}\,\partial/\partial X_{t,b}$ preserves the multidegree, so $E\in\mathcal K[X]_d$ in any case.

\emph{Step 1: products.} Let $\beta=\prod_tB_t(x_t)$ and $s=\prod_tS_t(x_t)$ with $B_t,S_t\in\nabla_{\mathcal K}(\nu(t))$.
For each position $t$ apply Lemma~\ref{pd-symbols-fock-words-position} with the parity of $t$, $d(t)$, $\nu(t)$, $B_t$
and $S_t$, and with $c$ for $t\in T$ and $0$ for $t\notin T$, where $\nu(t)=0$. It gives products $D_{t,l}$ of at most
$c$, respectively $0$, operators $E_{ab}$ and coefficients $a_{t,l}$. For a family $w=(l_t)_t$ let $D_w$ be the product
over $t$ of $D_{t,l_t}$ with $Y_a$ and $\partial/\partial Y_b$ replaced by $X_{t,a}$ and $\partial/\partial X_{t,b}$, and
$a_w=\prod_ta_{t,l_t}$. In $D_w$ each $t\in T$ occurs at most $c$ times and no $t\notin T$ occurs. The operators at
position $t$ act only on $X_t$, so for $G=\prod_tG_t(X_t)$ with $G_t$ homogeneous of degree $d(t)$,
$\sum_wa_wD_wG=\prod_t\big(\sum_la_{t,l}D_{t,l}G_t\big)(X_t)$. As $\Theta_1(G)=\prod_t\theta_1(G_t)(x_t)$, with $\theta_1$
of the parity of $t$, and the Fock pairing is multiplicative over the disjoint sets of variables $x_t$
(Definition~\ref{pd-symbols-test-fock}), for $H=\prod_tH_t(X_t)$ of the same kind
\[
 \langle\Theta_1(G)s,\Theta_1(H)\beta\rangle_{\mathrm F}=\prod_t\langle\theta_1(G_t)S_t,\theta_1(H_t)B_t\rangle_{\mathrm F}
 =\prod_t\Big\langle\theta_1\Big(\sum_la_{t,l}D_{t,l}G_t\Big),\theta_1(H_t)\Big\rangle_{\mathrm F}
 =\Big\langle\Theta_1\Big(\sum_wa_wD_wG\Big),\Theta_1(H)\Big\rangle_{\mathrm F}.
\]
Both sides are bilinear in $(G,H)$, and $\mathcal K[X]_d$ is spanned by the monomials of multidegree $d$, which are
such products. So the $D_w$ and $a_w$ work for $(\beta,s)$.

\emph{Step 2: the span.} Both sides of the identity are linear in $\beta$ and in $s$. If $(\mathcal L,a)$ works for
$(\beta,s)$ and $(\mathcal L',a')$ for $(\beta',s)$, then $\mathcal L\cup\mathcal L'$ with the sums of the coefficients
works for $(\beta+\beta',s)$, and $(\mathcal L,\lambda a)$ works for $(\lambda\beta,s)$; in the same way in $s$. By
Step~1, there is a family that works for every $(\beta,s)$ in the span of the products.
% Lean: tensor_operators_holds = tenOp_of_OB operator_weight_bound_holds = adjointTensorOfSingle_valid
%   TS.Top.fockRBlock (onePos_of_OB operator_words_bound), that is PI.TOW.tenOp_of_steps with stepDeg (the
%   multidegree), stepPure_of stepLift stepMonoFactor (Step 1: OnePos at each t with c or 0, concatW, embS,
%   Fintype.piFinset, fockRBlock and bilin_ext_restrict on the monomials) and stepSpan (Step 2). The case of one
%   position is Lemma pd-symbols-fock-words-position, and OpBound is Lemma pd-symbols-fock-words-operators.
\end{proof}

\begin{lemma}\label{pd-symbols-test-duality}
Let $C_{La},C_k,C_{rk},C_{\deg}$ be real numbers with the property of
Theorem~\ref{pd-symbols-count}. There is $C_1>0$ such that the following holds for every
$C\geq C_1$. Let the data of Situation~\ref{pd-edge-situation} be given with this $C$, form $d_T$,
$\lambda^\flat$ and $\lambda^\sharp$ as in Proposition~\ref{pd-symbols-cloud}, and let $\lambda_0^T$,
$\kappa_1$ be as in Lemma~\ref{pd-symbols-test-size} and $\lambda^{\sharp T}$ be formed from
$\lambda^\sharp$ as $\lambda_0^T$ from $\lambda_0$. Let $e$ be an edge with target $W^k$, $\gamma\in
GL_k(Q)$ and $g\in GL_k(W)$ with $e=e_\gamma(g)$ (Lemma~\ref{pd-symbols-edge-basis}), and $\Theta_\gamma$ the
flat model of $V^\flat$ at $\gamma$ of Lemma~\ref{pd-symbols-flat-exists}, formed with $d_T$. Form the
space $V$ and the lattices $H$, $R$, $R'$ of Definition~\ref{pd-projection-property} with
$\lambda_0^T$, the same objects $V^\flat$, $H^\flat$, $R^\flat$, $R'^\flat$ with $\lambda^\flat$, and the space
$V^\sharp$ with $\lambda^{\sharp T}$; let $A^\sharp$ be the integer $A$ of that definition for
$\lambda^{\sharp T}$. For a $W$-lattice $L$ let $L^\vee$ be its dual, and for an edge $e'$ with source
$L''$ and target $L$ let $e'^\vee$ be the edge with source $L^\vee$ and target $L''^\vee$. Then there are
$\mathcal K$-linear isomorphisms $\Psi_0\colon V^\vee\to V$ and $\Psi_f\colon(V^\flat)^\vee\to V^\flat$ such
that:
\begin{enumerate}
\item for $(\Psi,\kappa)$ equal to $(\Psi_0,\kappa_1(\lambda_0^T))$ with the lattices for $\lambda_0^T$, and to
  $(\Psi_f,\kappa_1(\lambda^\flat))$ with those for $\lambda^\flat$: for every $W$-lattice $L$,
  $\Psi(H(L)^\vee)\subset H(L^\vee)$ and $p^\kappa H(L^\vee)\subset\Psi(H(L)^\vee)$; and for every edge $e'$,
  $\Psi(R'(e'))\subset R(e'^\vee)$ and $p^\kappa R(e'^\vee)\subset\Psi(R'(e'))$;
\item for all $\beta,s\in V^\sharp$, $F\in\mathcal K[X]_{d_T}$ and $u\in V$ with $u=\Theta_\gamma(F)s$, there
  is $F'\in\mathcal K[X]_{d_T}$ in $\mathrm{Op}_{T,c_0}(\mathcal KF)$ (Definition~\ref{pd-symbols-test-operators})
  such that
  $p^{A^\sharp}p^{\kappa_1(\lambda_0^T)}\,\Psi_0^{-1}(z)(u)=T(\Theta_\gamma(F'))$ for every
  $T\in(V^\flat)^\vee$ and $z\in V$ with $z=\Psi_f(T)\beta$.
\end{enumerate}
\end{lemma}
\leanref{Formalism/SharpRestriction/GrowthBounds/DualFockLinear/RestCloudBOfThree.lean}{Formalism.MainNotions.TestSideCounting.TSProp}
\leanref{Formalism/SharpRestriction/GrowthBounds/DualFockLinear/RestCloudBOfThree.lean}{Formalism.MainNotions.TestSideCounting.TSOpAt}
\leanref{Formalism/SharpRestriction/GrowthBounds/DualFockLinear/RestCloudBOfThree.lean}{Formalism.MainNotions.TestSideCounting.TSOpStep}
\leanref{Formalism/SharpRestriction/GrowthBounds/DualFockLinear/RestCloudBOfThree.lean}{Formalism.MainNotions.TestSideCounting.RestTSOp}

\begin{proof}
Put $C_1=1$, let $C\geq1$ and let the data be as in the statement. By (3) of the situation,
$k\geq N_1+3r\geq3$. The weights $\lambda_0^T(s)$, $\lambda^\flat(s)$ and $\lambda^{\sharp T}(s)$ are dominant, and
$\lambda_0^T=\lambda^\flat+\lambda^{\sharp T}$.

(1) Lemma~\ref{pd-symbols-test-fock-dual}, applied to $\lambda_0^T$ and to $\lambda^\flat$, gives $\Psi_0$ and
$a_0\in\mathcal K^\times$ with $\Psi_0^{-1}(y)(u)=a_0\langle u,y\rangle_{\mathrm F}$ on $V$, and $\Psi_f$ and
$a_f\in\mathcal K^\times$ with $\Psi_f^{-1}(y)(u)=a_f\langle u,y\rangle_{\mathrm F}$ on $V^\flat$, with the lattice
properties (1).

(2) Let $\beta,s\in V^\sharp$, $F\in\mathcal K[X]_{d_T}$ and $u=\Theta_\gamma(F)s$. For each $t$ let $\nu(t)$ be the
natural weight of the weight formed from $\lambda^{\sharp T}(t)$ as in Definition~\ref{pd-local-model}; then
$V^\sharp$ is the span of the products $\prod_tf_t(x_t)$ with $f_t\in\nabla_{\mathcal K}(\nu(t))$. For $t\notin T$,
$\nu(t)=0$. For $t\in T$, $\nu(t)_1=\lambda^{\sharp T}(t)_1-\lambda^{\sharp T}(t)_k=\lambda_0(t)_2+b_T(t)\leq
c_{pp}+(c_0-2c_{pp})\leq c_0$, by (2) of the situation, as $2<k$. Lemma~\ref{pd-symbols-test-fock-words}, with
$c=c_0$, gives products $D_l$ of operators and $a_l\in\mathcal K$.

As in the proof of Lemma~\ref{pd-symbols-flat-exists}, $\Theta_\gamma(F)=\Theta_1(F^N)$, where $N_s=\gamma_s^{-1}$
for $s$ even and $N_s={}^t\gamma_s$ for $s$ odd, and $F^N$ is $F$ with $N_sX_s$ substituted for $X_s$. Put
$G=F^N$, $E=\sum_la_lD_lG\in\mathcal K[X]_{d_T}$,
$\kappa=p^{A^\sharp}p^{\kappa_1(\lambda_0^T)}a_0/a_f$ and $F'=\kappa E^{N^{-1}}$, of multidegree $d_T$; then
$\Theta_\gamma(F')=\kappa\Theta_1(E)$.

\emph{$F'$ lies in $\mathrm{Op}_{T,c_0}(\mathcal KF)$.} For a family $N'$ of invertible matrices, the chain rule gives
\[
 \Big(X_{t,a}\frac{\partial P}{\partial X_{t,b}}\Big)^{N'}=\sum_{a',b'}(N'_t)_{aa'}(N'^{-1}_t)_{b'b}\,
 X_{t,a'}\frac{\partial P^{N'}}{\partial X_{t,b'}} .
\]
By induction on the length, a product of such operators applied to $P$, transported by $N'$, is a combination of
products of operators at the same positions applied to $P^{N'}$. With $N'=N^{-1}$ and $P=G$, so that
$G^{N^{-1}}=F$, this gives $F'\in\mathrm{Op}_{T,c_0}(\mathcal KF)$.

\emph{The pairing.} Let $U\in(V^\flat)^\vee$ and $z=\Psi_f(U)\beta$. Then $U(y)=a_f\langle y,\Psi_f(U)\rangle_{\mathrm F}$
for $y\in V^\flat$, and $\Psi_f(U)=\Theta_1(H)$ for some $H\in\mathcal K[X]_{d_T}$, by Step~3 of the proof of
Lemma~\ref{pd-symbols-flat-exists}. As $\Theta_1(G)s=u$,
\[
 U(\Theta_\gamma(F'))=a_f\kappa\langle\Theta_1(E),\Theta_1(H)\rangle_{\mathrm F}
 =a_f\kappa\langle\Theta_1(G)s,\Theta_1(H)\beta\rangle_{\mathrm F}=a_f\kappa\langle u,z\rangle_{\mathrm F}
 =p^{A^\sharp}p^{\kappa_1(\lambda_0^T)}\Psi_0^{-1}(z)(u).
\]
% Lean: RestTSOp is restTSOp_of_ten = restTSOp_of restTSX (restOpF_of_ten hT) with C1 = 1. (1) is restTSX =
%   test_side_explicit_identification (Lemma pd-symbols-test-fock-dual) for lam = topW TT lam0 and flW d, inside
%   tsOpAt_of; (2) is tsOpAt_of (kappa = p^(Ar rest) (p^(kap1 lam) a0) / af, e1) with opFAt_of_ten:
%   first_weight_le_conductor, A x = Aa (vecE x) (la_tw x hm)^-1 and A' its inverse (the book's N), G = Lin A G0,
%   Lin_OpCl (Lin_word, Lin_X, Lin_pderiv, wspan_le_OpCl) and map_OpCl, mem_tens_dw for H, and hpair from
%   TenOp (Lemma pd-symbols-test-fock-words). The Lean asks part (1) for every lattice L and every edge e'.
\end{proof}

\begin{definition}\label{pd-symbols-test-maps}
Let $p\geq5$ be a prime, $(W,Q,\sigma,\iota)$ a Hida setting of degree $n'$ at $p$ with rings
$\mathcal R\subset\mathcal K$ and $\varpi_0$, $1\leq r<k$ integers, $T\subset\mathbf Z/n'\mathbf Z$,
$d\colon\mathbf Z/n'\mathbf Z\to\mathbf Z_{\geq0}$, $\lambda^\flat(s)=(d(s),0,\dots,0)$ and
$\lambda^\sharp(s)\in\mathbf Z^k$ dominant weights, and put $\lambda=\lambda^\flat+\lambda^\sharp$. Form
$V$, $H$, $R$ and $R'$ with $\lambda$, $V^\flat$, $H^\flat$ and $R^\flat$ with $\lambda^\flat$, and
$V^\sharp$ and its integer $A^\sharp$ (the integer $A$) with $\lambda^\sharp$, as in
Definition~\ref{pd-projection-property}. Let $\Theta\colon\mathcal K[X]_d\to V^\flat$ be an
$\mathcal R$-linear bijection (Definition~\ref{pd-symbols-flat-model}), $e$ an edge with target $L$, and
$c_0,\omega,\kappa\geq0$ integers. A family of \emph{test maps} for $(\lambda^\sharp,\Theta,e)$ with
constants $(c_0,\omega,\kappa)$ is a family $(\rho_\beta)_{\beta\in\mathcal B}$ of $\mathcal K$-linear maps
$\rho_\beta\colon V\to V^\flat$ such that:
\begin{enumerate}
\item $\rho_\beta(R(e))\subset R^\flat(e)$ for every $\beta$;
\item $\rho_\beta(R(e'))\subset\varpi_0^{-\omega}R^\flat(e')$ for every $\beta$ and every edge $e'$ with
  target $L$;
\item $\rho_\beta(H(L))\subset p^{A^\sharp}\varpi_0^{-\omega}H^\flat(L)$ for every $\beta$;
\item for every $\beta$, $F\in\mathcal K[X]_d$ and $s\in V^\sharp$ there is $F'\in\mathcal K[X]_d$ in
  $\mathrm{Op}_{T,c_0}(\mathcal KF)$ (Definition~\ref{pd-symbols-test-operators}) with
  $\rho_\beta(\Theta(F)s)=\Theta(F')$; here $\Theta(F)s\in V$ by Lemma~\ref{pd-edge-products};
\item every $u\in R(e)$ with $\rho_\beta(u)\in\varpi_0^{2\kappa+1}R^\flat(e)$ for all $\beta$ lies in
  $\varpi_0R(e)$.
\end{enumerate}
\end{definition}
\leanref{Formalism/SharpRestriction/GrowthBounds/DualFockLinear/VFlatSupported.lean}{Formalism.MainNotions.TestSideCounting.RhoAtG}
\leanref{Formalism/SharpRestriction/GrowthBounds/DualFockLinear/CloudBAtOfRho.lean}{Formalism.MainNotions.TestSideCounting.RhoAt}

\begin{lemma}\label{pd-symbols-test-maps-top}
Let $C_{La},C_k,C_{rk},C_{\deg}$ be real numbers with the property of
Theorem~\ref{pd-symbols-count}. There are $C_K,C_1>0$ such that the following holds for every
$C\geq C_1$. Let the data of Situation~\ref{pd-edge-situation} be given with this $C$, form $d_T$,
$\lambda^\flat$ and $\lambda^\sharp$ as in Proposition~\ref{pd-symbols-cloud}, and let $\lambda_0^T$,
$\lambda^{\sharp T}$ and $\kappa_1$ be as in Lemma~\ref{pd-symbols-test-duality}. Let $e$ be an edge with
target $W^k$, $\gamma\in GL_k(Q)$ and $g\in GL_k(W)$ with $e=e_\gamma(g)$
(Lemma~\ref{pd-symbols-edge-basis}), and $\Theta_\gamma$ the flat model of $V^\flat$ at $\gamma$ of
Lemma~\ref{pd-symbols-flat-exists}, formed with $d_T$. Then
$\kappa=\kappa_1(\lambda_0^T)+\kappa_1(\lambda^\flat)\leq C_K|T|k^2\log_pa$, and there is a family of test
maps for $(\lambda^{\sharp T},\Theta_\gamma,e)$ with constants $(c_0,r|T|c_0,\kappa)$
(Definition~\ref{pd-symbols-test-maps}).
\end{lemma}
\leanref{Formalism/SharpRestriction/GrowthBounds/DualFockLinear/RhoAtGOfPair.lean}{Formalism.MainNotions.TestSideCounting.PairAt}
\leanref{Formalism/SharpRestriction/GrowthBounds/DualFockLinear/RestCloudBOfThree.lean}{Formalism.MainNotions.TestSideCounting.pairAt_of_tsop}
\leanref{Formalism/SharpRestriction/GrowthBounds/DualFockLinear/RhoAtGOfPair.lean}{Formalism.MainNotions.TestSideCounting.rhoAtG_of_pair}
\leanref{Formalism/SharpRestriction/GrowthBounds/DualFockLinear/RhoTopStepOfPair.lean}{Formalism.MainNotions.TestSideCounting.rhoTopStep_of_pair}

\begin{proof}
Let $C_K$ and $C'_1$ be the constants of Lemma~\ref{pd-symbols-test-size}, $C''_1$ that of
Lemma~\ref{pd-symbols-test-duality}, and $C_1=\max(C'_1,C''_1)$. Let $C\geq C_1$ and the data be as in the
statement. By Lemma~\ref{pd-symbols-test-size}, $\kappa\leq C_K|T|k^2\log_pa$. The weights
$\lambda^{\sharp T}(s)$ are dominant, and $\lambda^\flat+\lambda^{\sharp T}=\lambda_0^T$. Form $V$, $H$, $R$,
$R'$ and $A$ with $\lambda_0^T$, the same objects $V^\flat,\dots,A^\flat$ with $\lambda^\flat$ and
$V^\sharp,\dots,A^\sharp$ with $\lambda^{\sharp T}$; then $A=A^\flat+A^\sharp$. Put $L=W^k$,
$\kappa_0=\kappa_1(\lambda_0^T)$ and $\kappa_f=\kappa_1(\lambda^\flat)$, write $\langle\,,\rangle$ and
$\langle\,,\rangle^\flat$ for the edge pairings of $\lambda_0^T$ and $\lambda^\flat$, and let $\Psi_0$ and
$\Psi_f$ be the isomorphisms of Lemma~\ref{pd-symbols-test-duality}. We use:
\begin{enumerate}
\item[(a)] Lemma~\ref{pd-symbols-test-onto}, with $d_T$ and $\lambda'_0=\lambda^{\sharp T}$, at the edge
  $e^\vee$: $R(e^\vee)$ lies in the $\mathcal R$-span of the products $y\beta$ with $y\in R^\flat(e^\vee)$
  and $\beta\in R^\sharp(e^\vee)$;
\item[(b)] the argument of the proof of Lemma~\ref{pd-symbols-rest-bound}, applied to $\lambda^{\sharp T}$,
  for which $\lambda^{\sharp T}(s)_1-\lambda^{\sharp T}(s)_k$ is $\lambda_0(s)_2+b_T(s)\leq c_0$ for $s\in T$
  and $0$ for $s\notin T$: $R^\sharp(e')\subset\varpi_0^{-r|T|c_0}H^\sharp(L''')$ for every edge $e'$ with
  source $L'''$;
\item[(c)] by Lemma~\ref{pd-edge-basis}(1), $H^\flat(L)$ has an $\mathcal R$-basis that is a
  $\mathcal K$-basis of $V^\flat$, so $y\in V^\flat$ lies in $H^\flat(L)$ if and only if
  $U(y)\in\mathcal R$ for every $U\in H^\flat(L)^\vee$.
\end{enumerate}

\emph{The maps.} Let $\mathcal B=R^\sharp(e^\vee)$. For $\beta\in V^\sharp$ and $U\in(V^\flat)^\vee$ put
$\Lambda_\beta(U)=p^{\kappa_0}\Psi_0^{-1}(\Psi_f(U)\beta)\in V^\vee$, which is defined as
$\Psi_f(U)\beta\in V$ by Lemma~\ref{pd-edge-products}, and let $\rho_\beta\colon V\to V^\flat$ be the
$\mathcal K$-linear map with $U(\rho_\beta(u))=p^{A^\sharp}\Lambda_\beta(U)(u)$ for all $U$. Then
\[
 \langle\rho_\beta(u),U\rangle^\flat=p^{A^\flat+A^\sharp}\Lambda_\beta(U)(u)=\langle u,\Lambda_\beta(U)\rangle .
 \tag{$*$}
\]
Let $e'$ be an edge with target $L$, $U\in R'^\flat(e')$ and $\beta'\in R^\sharp(e'^\vee)$. By
Lemma~\ref{pd-symbols-test-duality}(1) for $\Psi_f$, $\Psi_f(U)\in R^\flat(e'^\vee)$; by
Lemma~\ref{pd-edge-products}, $\Psi_f(U)\beta'\in R(e'^\vee)$; and by Lemma~\ref{pd-symbols-test-duality}(1)
for $\Psi_0$, $\Lambda_{\beta'}(U)\in R'(e')$. Together with ($*$) and Lemma~\ref{pd-edge-basis}(3), applied
to $\lambda_0^T$ and to $\lambda^\flat$, this gives: $\rho_{\beta'}(R(e'))\subset R^\flat(e')$ for every such
$e'$ and $\beta'\in R^\sharp(e'^\vee)$.

(1) This is the last statement with $e'=e$.

(2) Let $e'$ have target $L$ and $\beta\in\mathcal B$. By (b) for $e^\vee$, whose source is $L^\vee$,
$\beta=\varpi_0^{-r|T|c_0}\beta'$ with $\beta'\in H^\sharp(L^\vee)$, and $\beta'\in R^\sharp(e'^\vee)$ by
Lemma~\ref{pd-edge-basis}(2), as $e'^\vee$ has source $L^\vee$. As $\rho_\beta$ is linear in $\beta$,
$\rho_\beta(R(e'))=\varpi_0^{-r|T|c_0}\rho_{\beta'}(R(e'))\subset\varpi_0^{-r|T|c_0}R^\flat(e')$.

(3) Let $\beta=\varpi_0^{-r|T|c_0}\beta'$ be as in (2), $u\in H(L)$ and $U\in H^\flat(L)^\vee$. By
Lemma~\ref{pd-symbols-test-duality}(1), $\Psi_f(U)\in H^\flat(L^\vee)$; by Lemma~\ref{pd-edge-products},
$\Psi_f(U)\beta'\in H(L^\vee)$; and by Lemma~\ref{pd-symbols-test-duality}(1) again,
$\Lambda_{\beta'}(U)\in H(L)^\vee$. So $U(\rho_{\beta'}(u))=p^{A^\sharp}\Lambda_{\beta'}(U)(u)\in
p^{A^\sharp}\mathcal R$, and by (c) $\rho_{\beta'}(u)\in p^{A^\sharp}H^\flat(L)$. Hence
$\rho_\beta(u)\in p^{A^\sharp}\varpi_0^{-r|T|c_0}H^\flat(L)$.

(4) Let $F\in\mathcal K[X]_{d_T}$, $s\in V^\sharp$ and $u=\Theta_\gamma(F)s$.
Lemma~\ref{pd-symbols-test-duality}(2) gives $F'\in\mathcal K[X]_{d_T}$ in $\mathrm{Op}_{T,c_0}(\mathcal KF)$
with $p^{A^\sharp}p^{\kappa_0}\Psi_0^{-1}(\Psi_f(U)\beta)(u)=U(\Theta_\gamma(F'))$ for every
$U\in(V^\flat)^\vee$. The left side is $U(\rho_\beta(u))$, so $\rho_\beta(u)=\Theta_\gamma(F')$.

(5) Let $u\in R(e)$ with $\rho_\beta(u)\in\varpi_0^{2\kappa+1}R^\flat(e)$ for every $\beta\in\mathcal B$. By
Lemma~\ref{pd-edge-basis}(3) it suffices that $\langle u,G\rangle\in\varpi_0\mathcal R$ for every
$G\in R'(e)$. By Lemma~\ref{pd-symbols-test-duality}(1), $\Psi_0(G)\in R(e^\vee)$, so by (a) $\Psi_0(G)$ is an
$\mathcal R$-combination of products $y\beta$ with $y\in R^\flat(e^\vee)$ and $\beta\in\mathcal B$; so it
suffices that $\langle u,\Psi_0^{-1}(y\beta)\rangle\in\varpi_0\mathcal R$ for these. Put
$U=\Psi_f^{-1}(p^{\kappa_f}y)$, which lies in $R'^\flat(e)$ by Lemma~\ref{pd-symbols-test-duality}(1). Then
$\Lambda_\beta(U)=p^{\kappa}\Psi_0^{-1}(y\beta)$, and by ($*$)
$\langle u,\Lambda_\beta(U)\rangle=\langle\rho_\beta(u),U\rangle^\flat\in\varpi_0^{2\kappa+1}\mathcal R$, by
Lemma~\ref{pd-edge-basis}(3). As $p^\kappa=\varpi_0^{2\kappa}$, $\langle u,\Psi_0^{-1}(y\beta)\rangle\in
\varpi_0\mathcal R$.
% Lean: rhoTopStep_of_pair: pairAt_of_tsop (the three inputs repackaged as PairAt, with kappa0 = kap1 p lam,
%   kappaf = kap1 p (flW d), OntoAt at e.dual) and rhoAtG_of_pair with loss r |T| c0 (ordMax_top). In
%   rhoAtG_of_pair, Lambda beta = p^kappa0 Psi0^-1 o mulR beta o Psif, rho beta is its dual through
%   evalEquiv times p^(Ar rest), (*) is hI, the membership is hmem (edge_lattice_mul_of_weight_sum, b2), (b) is
%   Re_le_scale_HL with ordMax_top, (c) is HL_refl (hlSpans), and (1) to (5) are the five clauses in order
%   ((5) by span induction over OntoAt, with a3 and bePerfect).
\end{proof}

\begin{lemma}\label{pd-symbols-test-maps-peel}
Let $p\geq5$ be a prime, $(W,Q,\sigma,\iota)$ a Hida setting of degree $n'$ at $p$ with rings
$\mathcal R\subset\mathcal K$ and $\varpi_0$, $1\leq r<k$ integers, $T\subset\mathbf Z/n'\mathbf Z$,
$d\colon\mathbf Z/n'\mathbf Z\to\mathbf Z_{\geq0}$ with $d(s)=0$ for $s\notin T$, and
$\lambda^\flat(s)=(d(s),0,\dots,0)$. Let $\lambda^{\sharp T}(s),\mu(s)\in\mathbf Z^k$ be dominant weights with
$\lambda^{\sharp T}(s)=0$ for $s\notin T$ and $\mu(s)=0$ for $s\in T$, and put
$\lambda^\sharp=\lambda^{\sharp T}+\mu$. For $s\in\mathbf Z/n'\mathbf Z$ let $\mu(s)'$ be formed from $\mu(s)$
as in Definition~\ref{pd-local-model}, and let $M\geq0$ be an integer with
$\sum_s\sum_{l\leq r}(\mu(s)'_l-\mu(s)'_k)\leq M$. Let $\Theta\colon\mathcal K[X]_d\to V^\flat$ be an
$\mathcal R$-linear bijection, $e$ an edge and $c_0,\omega,\kappa\geq0$ integers. If there is a family of test
maps for $(\lambda^{\sharp T},\Theta,e)$ with constants $(c_0,\omega,\kappa)$, then there is one for
$(\lambda^\sharp,\Theta,e)$ with constants $(c_0,\omega+M,\kappa)$.
\end{lemma}
\leanref{Formalism/SharpRestriction/GrowthBounds/DualFockLinear/VFlatSupported.lean}{Formalism.MainNotions.TestSideCounting.rho_peel}
\leanref{Formalism/SharpRestriction/GrowthBounds/DualFockLinear/OneMemH0WeightZero.lean}{Formalism.MainNotions.TestSideCounting.Peel}
\leanref{Formalism/SharpRestriction/GrowthBounds/DualFockLinear/DualWeylZeroConst.lean}{Formalism.MainNotions.TestSideCounting.peelMap}
\leanref{Formalism/SharpRestriction/GrowthBounds/DualFockLinear/FilterOfSupported.lean}{Formalism.MainNotions.TestSideCounting.slice}

\begin{proof}
Put $\lambda^T=\lambda^\flat+\lambda^{\sharp T}$ and $\lambda=\lambda^T+\mu=\lambda^\flat+\lambda^\sharp$; then
$\lambda(s)=\lambda^T(s)$ and $\mu(s)=0$ for $s\in T$, and $\lambda^T(s)=0$ and $\lambda(s)=\mu(s)$ for
$s\notin T$, and the same holds for $\lambda^\sharp$, $\lambda^{\sharp T}$ and $\mu$. Form $V$, $H$, $R$, $R'$
and $A$ with $\lambda$, the same objects with superscript $T$ with $\lambda^T$ and with superscript $\mu$ with
$\mu$, and $V^\flat$, $H^\flat$, $R^\flat$ with $\lambda^\flat$; then $A=A^T+A^\mu$, and the integers of
$\lambda^\sharp$ and $\lambda^{\sharp T}$ differ by $A^\mu$. Let $L$ be the target of $e$ and
$(\rho^T_\beta)_{\beta\in\mathcal B}$ the given family; its maps go from $V^T$ to $V^\flat$.

\emph{Slices.} Write a monomial in the $x_{s,a,b}$ as $m_Tm_\circ$, with $m_T$ in the variables with $s\in T$
and $m_\circ$ in the others. For $G\in(V^\mu)^\vee$ extend $G$ to a $\mathcal K$-linear form $\tilde G$ on
$\mathcal K[x_{s,a,b}]$, and let $\sigma_G$ be the $\mathcal K$-linear map of $\mathcal K[x_{s,a,b}]$ with
$\sigma_G(m_Tm_\circ)=\tilde G(m_\circ)m_T$. If $P$ involves only variables with $s\in T$, then
$\sigma_G(PQ)=P\sigma_G(Q)$, and $\sigma_G(Q)=G(Q)$ if moreover $Q\in V^\mu$ involves only the others. The
dual Weyl module of the zero weight consists of the constants, so a spanning product $\prod_sf_s(x_s)$ of
$V$ is $PQ$ with $P=\prod_{s\in T}f_s$ a spanning product of $V^T$ and $Q=\prod_{s\notin T}f_s$ one of
$V^\mu$; and the same holds for the lattices $H(L')$ and $H(e')$ of a lattice $L'$ or an edge $e'$, with
those of $\lambda^T$ and $\mu$, as $1$ lies in the lattices of the zero weight. Hence:
\begin{enumerate}
\item[(i)] $\sigma_G(V)\subset V^T$, and if $c\in\mathcal K$ and $G(H^\mu(e'))\subset c\mathcal R$, then
  $\sigma_G(H(e'))\subset cH^T(e')$; the same holds for $H(L')$.
\end{enumerate}
Put $\pi_G=p^{A^\mu}\sigma_G\colon V\to V^T$. By (i):
\begin{enumerate}
\item[(ii)] if $N\in\mathbf Z$ and $\langle Q,G\rangle^\mu\in\varpi_0^N\mathcal R$ for every $Q\in R^\mu(e')$,
  then $\pi_G(R(e'))\subset\varpi_0^NR^T(e')$. Indeed, for $Q\in H^\mu(e')$,
  $\langle\varpi_0^{-A^\mu}Q,G\rangle^\mu=\varpi_0^{A^\mu}G(Q)$, so $G(H^\mu(e'))\subset\varpi_0^{N-A^\mu}\mathcal R$
  and $\pi_G(\varpi_0^{-A}H(e'))\subset\varpi_0^{2A^\mu-A+N-A^\mu}H^T(e')=\varpi_0^NR^T(e')$;
\item[(iii)] if $G(H^\mu(L))\subset c\mathcal R$, then $\pi_G(H(L))\subset p^{A^\mu}cH^T(L)$.
\end{enumerate}
Moreover, take the basis $(x_i)$, the dual basis $(x_i^*)$ and the integers $\pi_i$ of
Lemma~\ref{pd-edge-basis}(1) for $\mu$ and $e$. As in the proof of Lemma~\ref{pd-symbols-rest-bound},
$A^\mu-\pi_i\leq M$, so
\begin{enumerate}
\item[(iv)] $R'^\mu(e)\subset\varpi_0^{-M}H^\mu(L)^\vee$, as $R'^\mu(e)$ and $H^\mu(L)^\vee$ are spanned by the
  $\varpi_0^{\pi_i-A^\mu}p^{-\pi_i}x_i^*$ and by the $p^{-\pi_i}x_i^*$.
\end{enumerate}

\emph{The maps.} Index the new family by the pairs $(\beta,G)$ with $\beta\in\mathcal B$ and $G\in R'^\mu(e)$,
and put $\rho_{\beta,G}=\rho^T_\beta\circ\pi_G$. For $G\in R'^\mu(e)$ write $G=\varpi_0^{-M}G'$ with
$G'\in H^\mu(L)^\vee$ by (iv); then $G'\in R'^\mu(e')$ for every edge $e'$ with target $L$, by
Lemma~\ref{pd-edge-basis}(2).

(1) By Lemma~\ref{pd-edge-basis}(3), $\langle Q,G\rangle^\mu\in\mathcal R$ for $Q\in R^\mu(e)$, so
$\pi_G(R(e))\subset R^T(e)$ by (ii), and $\rho^T_\beta(R^T(e))\subset R^\flat(e)$.

(2) Let $e'$ have target $L$. By Lemma~\ref{pd-edge-basis}(3) for $e'$,
$\langle Q,G\rangle^\mu=\varpi_0^{-M}\langle Q,G'\rangle^\mu\in\varpi_0^{-M}\mathcal R$ for $Q\in R^\mu(e')$, so
$\pi_G(R(e'))\subset\varpi_0^{-M}R^T(e')$ by (ii), and $\rho^T_\beta$ maps this into
$\varpi_0^{-M-\omega}R^\flat(e')$.

(3) $G(H^\mu(L))\subset\varpi_0^{-M}\mathcal R$, so $\pi_G(H(L))\subset p^{A^\mu}\varpi_0^{-M}H^T(L)$ by (iii),
and $\rho^T_\beta$ maps this into $p^{A^\mu}p^{A^{\sharp T}}\varpi_0^{-M-\omega}H^\flat(L)$, where $A^{\sharp T}$
is the integer of $\lambda^{\sharp T}$ and $A^{\sharp T}+A^\mu$ that of $\lambda^\sharp$.

(4) Let $F\in\mathcal K[X]_d$ and let $z$ lie in the space of $\lambda^\sharp$. As $d$ vanishes off $T$,
$\Theta(F)\in V^\flat$ involves only the variables $x_{s,a,b}$ with $s\in T$, so
$\pi_G(\Theta(F)z)=\Theta(F)\pi'_G(z)$, where $\pi'_G(z)=p^{A^\mu}\sigma_G(z)$ lies in the space of
$\lambda^{\sharp T}$ by (i) for $\lambda^\sharp$. Property~(4) of $\rho^T_\beta$ gives $F'$.

(5) Let $u\in R(e)$ with $\rho_{\beta,G}(u)\in\varpi_0^{2\kappa+1}R^\flat(e)$ for all $\beta$ and $G$. For each
$G\in R'^\mu(e)$, $\pi_G(u)\in R^T(e)$ by (1), so $\pi_G(u)\in\varpi_0R^T(e)$ by property~(5) of the family
$\rho^T$. Put $y_i=\varpi_0^{\pi_i-A^\mu}x_i\in R^\mu(e)$ and
$G_i=\varpi_0^{\pi_i-A^\mu}p^{-\pi_i}x_i^*\in R'^\mu(e)$, so that $p^{A^\mu}G_i(y_j)$ is $1$ for $i=j$ and $0$
otherwise. We claim $u=\sum_i\pi_{G_i}(u)y_i$. Both sides are linear in $u$; for a spanning product $u=PQ$
as above, $\pi_{G_i}(u)=p^{A^\mu}G_i(Q)P$, and $\sum_ip^{A^\mu}G_i(Q)y_i=Q$. Now $\pi_{G_i}(u)=\varpi_0z_i$ with
$z_i\in R^T(e)$, and $z_iy_i\in R(e)$ by Lemma~\ref{pd-edge-products}, so $u\in\varpi_0R(e)$.
% Lean: rho_peel with Peel TT lam lamT sp and Peel TT rest restT sp. Slices: slice, extD (exists_ext),
%   slice_mul_left, slice_spec, slice_mul, prod_split, slice_tensor ((i)), slice_V, slice_Hten, peelMap,
%   peel_Hten; (ii) is peel_Re, (iii) is peel_HL, (iv) is Re'_le_scale (edge_lattice_adapted_basis); the
%   five clauses are G1, G2, G3 with L1, L2, L3, Op (V_flat_supported, slice_mul_left) and Nd with hexp.
\end{proof}

\begin{lemma}\label{pd-symbols-test-maps-cloud}
Let $p$, $(W,Q,\sigma,\iota)$, $\mathcal R\subset\mathcal K$, $\varpi_0$, $r$, $k$, $T$, $d$, $\lambda^\flat$,
$\lambda^\sharp$, $\lambda=\lambda^\flat+\lambda^\sharp$, the spaces and lattices and $\Theta$ be as in
Definition~\ref{pd-symbols-test-maps}, and let $R'^\flat$ and $R^\sharp$ be formed with $\lambda^\flat$ and
$\lambda^\sharp$. Let $e$ be an edge with target $L$, $c_0,\omega,\kappa\geq0$ integers, and
$(\rho_\beta)$ a family of test maps for $(\lambda^\sharp,\Theta,e)$ with constants $(c_0,\omega,\kappa)$.
Let $\Phi_1,\dots,\Phi_{n_b}\in\mathcal K[X]_d$ with $\Theta(\Phi_b)\in R^\flat(e)$, and $M$ their
$\mathcal K$-span. Let $y_1,\dots,y_n\in R^\sharp(e)$ and $(c_{bi})$ a family in $\mathcal R$ with
$v=\sum_{b,i}c_{bi}\Theta(\Phi_b)y_i\notin\varpi_0R(e)$. Then there are $F_1\in\mathrm{Op}_{T,c_0}(M)$ of
multidegree $d$ with $\Theta(F_1)\in R^\flat(e)$ and $\Theta(F_1)\notin\varpi_0R^\flat(e)$, a $\mathcal K$-linear
map $\Lambda\colon(V^\flat)^\vee\to V^\vee$ and $\xi\in\mathcal R$ such that, with the edge pairings of
$\lambda$ and $\lambda^\flat$:
\begin{enumerate}
\item $\xi\langle v,\Lambda U\rangle=p^\kappa\langle\Theta(F_1),U\rangle$ for every $U\in(V^\flat)^\vee$;
\item for $U\in R'^\flat(e)$, $\Lambda U\in R'(e)$ and $\langle v,\Lambda U\rangle\in\mathcal R$;
\item for every edge $e'$ with target $L$ and every integer $i$, $\Lambda$ maps $\varpi_0^iR'^\flat(e')$
  into $\varpi_0^{i-\omega}R'(e')$;
\item for every integer $i$, $\Lambda$ maps $\varpi_0^iH^\flat(L)^\vee$ into $\varpi_0^{i-\omega}H(L)^\vee$.
\end{enumerate}
\end{lemma}
\leanref{Formalism/SharpRestriction/GrowthBounds/DualFockLinear/CloudBAtOfRho.lean}{Formalism.MainNotions.TestSideCounting.cloudBAt_of_rho}

\begin{proof}
Let $A$, $A^\flat$ and $A^\sharp$ be the integers of $\lambda$, $\lambda^\flat$ and $\lambda^\sharp$; then
$A=A^\flat+A^\sharp$, and $p=\varpi_0^2$. We use that if a $\mathcal K$-linear map $f\colon V_0\to V_1$ maps an
$\mathcal R$-lattice $N_0$ into $cN_1$, then $U\mapsto U\circ f$ maps $tN_1^\vee$ into $tcN_0^\vee$.

\emph{The element $F_1$.} By Lemma~\ref{pd-edge-products}, each $\Theta(\Phi_b)y_i$ lies in $R(e)$, so $v\in R(e)$.
By property~(5) there is $\beta$ with $\rho_\beta(v)\notin\varpi_0^{2\kappa+1}R^\flat(e)$, and
$\rho_\beta(v)\in R^\flat(e)$ by property~(1). Let $a$ be the largest integer $\leq2\kappa$ with
$\rho_\beta(v)\in\varpi_0^aR^\flat(e)$, and $\rho_\beta(v)=\varpi_0^ay$; then $y\in R^\flat(e)$ and
$y\notin\varpi_0R^\flat(e)$. By property~(4), $\rho_\beta(\Theta(\Phi_b)y_i)=\Theta(F'_{bi})$ with
$F'_{bi}\in\mathrm{Op}_{T,c_0}(\mathcal K\Phi_b)\subset\mathrm{Op}_{T,c_0}(M)$ of multidegree $d$. Put
$F_1=\varpi_0^{-a}\sum_{b,i}c_{bi}F'_{bi}$. It lies in $\mathrm{Op}_{T,c_0}(M)$, has multidegree $d$, and
$\Theta(F_1)=\varpi_0^{-a}\rho_\beta(v)=y$.

\emph{The map.} Put $\Lambda U=\varpi_0^{2A^\flat-2A}\,U\circ\rho_\beta$ and $\xi=\varpi_0^{2\kappa-a}$.

(1) $\xi\langle v,\Lambda U\rangle=\varpi_0^{2\kappa-a+2A^\flat-2A}U(\rho_\beta(p^Av))=\varpi_0^{2\kappa+2A^\flat}U(y)
=p^\kappa\langle\Theta(F_1),U\rangle$.

(2) By property~(1), $\rho_\beta(H(e))=\varpi_0^A\rho_\beta(R(e))\subset\varpi_0^{A-A^\flat}H^\flat(e)$. So for
$U\in R'^\flat(e)=\varpi_0^{-A^\flat}H^\flat(e)^\vee$, $U\circ\rho_\beta\in\varpi_0^{A-2A^\flat}H(e)^\vee$ and
$\Lambda U\in\varpi_0^{-A}H(e)^\vee=R'(e)$. Then $\langle v,\Lambda U\rangle\in\mathcal R$ by
Lemma~\ref{pd-edge-basis}(3), as $v\in R(e)$.

(3) By property~(2), $\rho_\beta(H(e'))\subset\varpi_0^{A-A^\flat-\omega}H^\flat(e')$. For
$U\in\varpi_0^iR'^\flat(e')=\varpi_0^{i-A^\flat}H^\flat(e')^\vee$, this gives
$U\circ\rho_\beta\in\varpi_0^{i-2A^\flat+A-\omega}H(e')^\vee$ and $\Lambda U\in\varpi_0^{i-\omega-A}H(e')^\vee=
\varpi_0^{i-\omega}R'(e')$.

(4) By property~(3), $\rho_\beta(H(L))\subset\varpi_0^{2A^\sharp-\omega}H^\flat(L)$. For
$U\in\varpi_0^iH^\flat(L)^\vee$, $U\circ\rho_\beta\in\varpi_0^{i+2A^\sharp-\omega}H(L)^\vee$, and
$\Lambda U\in\varpi_0^{i-\omega}H(L)^\vee$, as $2A^\sharp+2A^\flat-2A=0$.
% Lean: cloudBAt_of_rho: hSre/hxre (v in Re, from hv), beta from h5, a and y from desc, F' from h4 (hG, hGmem
%   with OpCl_mono, hF1, hF1mem), Lambda = w^(2 Ar flW - 2 Ar lam0) dualMap (rho beta), cc = w^(2 kappa - a);
%   the general fact is dualMap_mem_scale_dualLat_of_maps_to_scale, and (1) to (4) are the four clauses.
\end{proof}

\begin{lemma}\label{pd-symbols-test-cloud}
Let $C_{La},C_k,C_{rk},C_{\deg}$ be real numbers with the property of
Theorem~\ref{pd-symbols-count}. There are $C_K,C_6>0$ such that the following holds for every
$C\geq C_6$. Let the data of Situation~\ref{pd-edge-situation} with this $C$, and $(J,f)$, $K$, $N_s$,
$q$, $c'$, $\deg$, $m_*$, $E_0$, $X_j$, $d_T$, $E'$, $d$, $P_1,\dots,P_{n_b}$, $e$, $\gamma$, $g$,
$\Phi_1,\dots,\Phi_{n_b}$, $\lambda^\flat$, $\lambda^\sharp$, $V^\flat$ and $R^\sharp$ satisfy the
hypotheses of Proposition~\ref{pd-symbols-cloud}; let $\Theta_\gamma$ be the flat model of $V^\flat$ at
$\gamma$ of Lemma~\ref{pd-symbols-flat-exists}, formed with $d_T$, and form the lattices $R^\flat$,
$R'^\flat$ and $H^\flat$ of Definition~\ref{pd-projection-property} with $\lambda^\flat$. Let $M$ be
the $\mathcal K$-span of $\Phi_1,\dots,\Phi_{n_b}$ and $L=W^k$ the target of $e$. Then there is an
integer $0\leq\kappa\leq C_K|T|k^2\log_pa$ with the following property. Let $y_1,\dots,y_n\in R^\sharp(e)$
and $(c_{bi})$ a family in $\mathcal R$ such that $v=\sum_{b,i}c_{bi}\Theta_\gamma(\Phi_b)y_i\notin
\varpi_0R(e)$. Then there are $F_1\in\mathrm{Op}_{T,c_0}(M)$ of multidegree $d_T$
(Definition~\ref{pd-symbols-test-operators}) with $\Theta_\gamma(F_1)\in R^\flat(e)$ and
$\Theta_\gamma(F_1)\notin\varpi_0R^\flat(e)$, a $\mathcal K$-linear map
$\Lambda\colon(V^\flat)^\vee\to V^\vee$ and $\xi\in\mathcal R$ such that, with the edge pairings of
$\lambda_0$ and $\lambda^\flat$:
\begin{enumerate}
\item $\xi\langle v,\Lambda U\rangle=p^\kappa\langle\Theta_\gamma(F_1),U\rangle$ for every
  $U\in(V^\flat)^\vee$;
\item for $U\in R'^\flat(e)$, $\Lambda U\in R'(e)$ and $\langle v,\Lambda U\rangle\in\mathcal R$;
\item for every edge $e'$ with target $L$ and every integer $i$, $\Lambda$ maps
  $\varpi_0^iR'^\flat(e')$ into $\varpi_0^{i-\omega}R'(e')$;
\item for every integer $i$, $\Lambda$ maps $\varpi_0^iH^\flat(L)^\vee$ into
  $\varpi_0^{i-\omega}H(L)^\vee$.
\end{enumerate}
\end{lemma}
\leanref{Formalism/SharpRestriction/GrowthBounds/DualFockLinear/TestSideExplicitIdentification.lean}{Formalism.MainNotions.TestSideCounting.RestCloudB}
\leanref{Formalism/SharpRestriction/GrowthBounds/DualFockLinear/TestSideExplicitIdentification.lean}{Formalism.MainNotions.TestSideCounting.CloudBStep}
\leanref{Formalism/SharpRestriction/GrowthBounds/DualFockLinear/TestSideExplicitIdentification.lean}{Formalism.MainNotions.TestSideCounting.CloudBAt}
\leanref{Formalism/SharpRestriction/GrowthBounds/DualFockLinear/LogSizeBoundLe.lean}{Formalism.MainNotions.TestSideCounting.restCloudB_of_op}
\leanref{Formalism/SharpRestriction/GrowthBounds/DualFockLinear/RestCloudBOfThree.lean}{Formalism.MainNotions.TestSideCounting.restCloudB_of_three}
\leanref{Formalism/SharpRestriction/GrowthBounds/DualFockLinear/RestCloudBOfThree.lean}{Formalism.MainNotions.TestSideCounting.pairAt_of_tsop}

\begin{proof}
Let $C_K$ and $C_1$ be the constants of Lemma~\ref{pd-symbols-test-maps-top}, and put $C_6=C_1$. Let
$C\geq C_6$ and the data be as in the statement, and let $\lambda_0^T$ and $\lambda^{\sharp T}$ be as in
Lemma~\ref{pd-symbols-test-duality}. By Lemma~\ref{pd-symbols-test-maps-top},
$\kappa=\kappa_1(\lambda_0^T)+\kappa_1(\lambda^\flat)$ satisfies $0\leq\kappa\leq C_K|T|k^2\log_pa$, and there is a
family of test maps for $(\lambda^{\sharp T},\Theta_\gamma,e)$ with constants $(c_0,r|T|c_0,\kappa)$.

Let $\mu(s)=\lambda_0(s)$ for $s\notin T$ and $\mu(s)=0$ for $s\in T$. These weights are dominant, and
$\lambda^\sharp=\lambda^{\sharp T}+\mu$, as $d_T(s)=0$ and $\lambda^\sharp(s)=\lambda_0(s)$ for $s\notin T$. For
$s\notin T$ the entries of $\lambda_0(s)$ have absolute value at most $c_{pp}$, so the sum of
Lemma~\ref{pd-symbols-test-maps-peel} is at most $2rn'c_{pp}$. That lemma, with $M=2rn'c_{pp}$, gives a family of
test maps for $(\lambda^\sharp,\Theta_\gamma,e)$ with constants $(c_0,\omega,\kappa)$, where
$\omega=r|T|c_0+2rn'c_{pp}$; here $\lambda^\flat+\lambda^\sharp=\lambda_0$.

For each $b$, $\Sigma_g(\Phi_b)$ is the symbol multiplier $P_b\Sigma_g(\mathrm{SP}_g)$, which lies in
$\mathcal R[X]$, as $P_b\in\mathcal R[X]$ and $\Sigma_g(\mathrm{SP}_g)\in\mathbf Z[X]$ (Step~3 of the proof of
Lemma~\ref{pd-symbols-primitive}). By property~(2) of the flat model with $i=0$,
$\Theta_\gamma(\Phi_b)\in R^\flat(e)$. Lemma~\ref{pd-symbols-test-maps-cloud}, applied to this family, the
$\Phi_b$, the $y_i$ and $(c_{bi})$, gives $F_1$, $\Lambda$ and $\xi$ with (1) to (4).
% Lean: restCloudB_of_op = restCloudB_of_two hop restSize = restCloudB_of_three hop restSize ontoAll =
%   restCloudB_of_pair (restPairTop_of hop hsz hon) = restCloudB_of_top (restRhoTop_of_pair ...) =
%   restCloudB_of_rho (restRho_of_top ...). The first paragraph is restPairTop_of and rhoTopStep_of_pair
%   (Lemma pd-symbols-test-maps-top); the second is rhoStep_of_top: rho_peel with specW TT lam0, top_add_spec,
%   specW_rest and ordMax_spec (Lemma pd-symbols-test-maps-peel), then rhoAt_of_G; the third is
%   cloudBStep_of_rho: hFs (subst_SPtrue and the flat model hTheta1) and cloudBAt_of_rho (Lemma
%   pd-symbols-test-maps-cloud) with edge_lattice_mul_of_weight_sum for hv.
\end{proof}

\subsection{One-sided elements and signed chain data}\label{subsection-pd-one-sided-chain-data}

\begin{definition}\label{pd-symbols-test-one-sided}
Let $p$ be a prime, $\mathcal R$ an integral domain, $n'\geq2$ an even integer, $r,K\geq1$ integers,
$k=r+K$, $T\subset\mathbf Z/n'\mathbf Z$, $(J,f)$ a run system for $T$ of length $r$, $N_1,G',c'\geq0$
integers, $\deg\colon\mathbf Z/n'\mathbf Z\to\mathbf Z_{\geq0}$ and $Y$ a real number; we use the notation of
Definitions~\ref{pd-symbols-run-system} and~\ref{pd-symbols-module}. Let $q=p^{G'}$, let $\alpha_{\max}$ be as in
Lemma~\ref{pd-symbols-space}, $m_*=(N_1+1)\big((2q+3)c'+kq+c'\alpha_{\max}+3c'\big)$, and for $j,x\in\mathbf Z$
let $B_j(x)$ be as in Definition~\ref{pd-symbols-module}(4) and
$\theta^1_j(x)=B_j(x)/2-\lfloor B_j(x)/(Yk^2)\rfloor$. A \emph{one-sided element} of normal degree $o\geq0$ is a
polynomial
\[
 \sigma=X^\ell\sum_{i=1}^n\prod_{j=0}^{J-1}\phi_j(\xi_{i,j})\in\mathcal R[X]
\]
such that $X^\ell$ is a label with $\ell(s,c)=0$ for $s\notin T$, every monomial of $\sigma$ has normal degree
$o$, and for each $0\leq j<J$ there are $o_j\colon\mathbf Z\to\mathbf R$ and $b_j\colon\mathbf Z\to\mathbf Z_{\geq0}$
such that every $\xi_{i,j}$ lies in the segment module $X_j(m_*,\theta^1_j,o_j)$, with $m_*$ the constant
function, and has all its monomials of degree $b_j(x)$ in $z_{x,1},\dots,z_{x,k}$ for $x\in\Lambda_j$, and
$b_j(x)\leq\frac{101}{100}\deg(\bar x)/p^{\operatorname{dist}_j(x)}$ for every $x\in\Lambda_j$ that is not new or
lies in $\Lambda'_j$.
\end{definition}
\leanref{Formalism/SharpRestriction/GrowthBounds/DualFockLinear/OneSidedLemma4Realisable.lean}{Formalism.MainNotions.TestSideCounting.OneSidedLemma4Element}
\leanref{Formalism/SharpRestriction/GrowthBounds/SourceCleaning.lean}{Formalism.MainNotions.SourceCleaning.monPh}
\leanref{Formalism/SharpRestriction/GrowthBounds/SourceCleaning.lean}{Formalism.MainNotions.SourceCleaning.mStar}
\leanref{Formalism/SharpRestriction/GrowthBounds/SourceCleaning.lean}{Formalism.MainNotions.SourceCleaning.levDeg}
\leanref{Formalism/SharpRestriction/GrowthBounds/SourceCleaning.lean}{Formalism.MainNotions.SourceCleaning.sl}

\begin{definition}\label{pd-symbols-test-chains}
Let $p$ be a prime, $\mathcal R$ a discrete valuation ring with fraction field $\mathcal K$ in which $p$ is not
a unit, $v$ the valuation of $\mathcal K$ with $v(p)=1$, and $r,J\geq1$ integers. A \emph{signed chain datum}
$\mathcal D$ of rank $r$ with $J$ levels consists of:
\begin{enumerate}
\item for each $j\in\mathbf Z/J\mathbf Z$ an integer $n_j\geq0$ and a type, vector or covector, for each $l<n_j$;
  the \emph{letters} $(j,l)$, $l<n_j$, and the \emph{links} $(j,m)$, $m<r$, are the \emph{positions}, and each
  position $y$ has a degree $d(y)\geq0$;
\item integers $g(j,m)\geq0$ for the links, with $P(j,m)=p^{g(j,m)}$;
\item for each $j$ a strictly increasing sequence of \emph{chains} $C_{j,1}\subsetneq\dots\subsetneq C_{j,n^c_j}$,
  sets of letters of level $j$ and indices $m<r$, each containing all the $m<r$, with real numbers $\theta_{j,i}$
  and $s_{j,i}$; and real numbers $\theta^L_{j,m}>0$ for $m+2\leq r$;
\item units $\gamma(\alpha)\in\mathcal R^\times$, for the exponent vectors $\alpha$ of the variables below, such
  that for every $k\geq1$ and $\epsilon=\pm1$ there are $\lambda_\mu\in\mathcal R$, for finitely many exponent
  vectors $\mu$ of total degree at most $6k$, with
  $v\big(\gamma(\alpha)^\epsilon-\sum_\mu\lambda_\mu\prod_y\binom{\alpha(y)}{\mu(y)}\big)\geq k$ for every $\alpha$;
\item a real number $c\geq1$.
\end{enumerate}
Let $\mathcal K[x]$ be the polynomial ring in $x_{y,a}$, $y$ a position and $a<r$, with the type of a letter and
the vector type for links. The \emph{model positions} are the letters $(j,l)$ and, for each link $(j,m)$, a
label, a lift and, if $m\geq1$, a top, all of vector type; let $\mathcal K[\tilde x]$ be the polynomial ring in
the $\tilde x_{y,a}$, $y$ a model position, and $\varphi\colon\mathcal K[\tilde x]\to\mathcal K[x]$ the
$\mathcal K$-algebra map sending the variables of a letter, a label or a top to those of the letter or link,
and those of the lift of $(j,m)$ to the $P(j,m)$-th powers of those of $(j,m)$. For residues $\rho(j,m,a)$ and
$F\in\mathcal K[x]$ let $F^\rho\in\mathcal K[\tilde x]$ be the sum of the terms $F_\mu x^\mu$ with
$\mu(j,m,a)\equiv\rho(j,m,a)\pmod{P(j,m)}$ for all links, each with the exponents of the link variables divided
by $P(j,m)$, rounded down, and moved to the lift. For a set $S$ of positions, $\operatorname{val}_S$ is formed as
in Definition~\ref{pd-symbols-module}(1), with vector positions in the role of the even $v$ and covector
positions in that of the odd $v$. The chain $C_{j,i}$ gives the set $\tilde C_{j,i}$ of model positions consisting
of its letters and the lifts of the $(j-1,m)$ with $m\in C_{j,i}$; for $m<r$ let $S_{j,m}$ be the set of the letters of level $j$
and the links $(j,m')$ with $m'>m$, and $\tilde S_{j,m}$ that of the lifts of level $j-1$, the letters of level
$j$ and the tops $(j,m')$ with $m'>m$. Put
\[
 e_P=r^2\Big(\log_p\Big(2\sum_yd(y)+r\Big)+1\Big),\quad D_m=e_P\sum_jn^c_j,\quad D_l=J(r-1)e_P,\quad
 \kappa_{\mathcal D}=\sum_yr\big(\log_p(d(y)+r)+1\big),
\]
$E(C_E)=D_m+D_l+C_E\kappa_{\mathcal D}$ and $n_0=v(\prod_yd(y)!)/2$; for $F,G\in\mathcal K[x]$ put
$B_{\mathcal D}(F,G)=\langle F,\sum_\mu\gamma(\mu)G_\mu x^\mu\rangle_{\mathrm F}$ and
$v_{\mathrm F}(F)=\min_\mu(v(F_\mu)+v(\mu!)/2)$. Let $C_H,C_E,C_6,C_{18}$ be real numbers.
\begin{itemize}
\item $\mathcal D$ satisfies \emph{(H) for $C_H$} if $rd_L(j)/P(j,m)+C_Hrc\leq s_{j+1,i}$ for all $j$, $m$ and
  $i$, where $d_L(j)$ is the sum of the degrees of the letters of level $j$.
\item $F\in\mathcal K[x]$ is \emph{admissible} if it has degree $d(y)$ in the row of each position $y$,
  $\operatorname{val}_{\tilde C_{j,i}}(F^\rho)\geq\theta_{j,i}$ for every chain $C_{j,i}$ and every $\rho$, and
  $\operatorname{val}_{S_{j,m}}(F)\geq\theta^L_{j,m}$ for $m+2\leq r$.
\item A \emph{projection pair} for $(\mathcal D,F)$ and $(C_E,C_6,C_{18})$ is a pair $(G',R')$ in
  $\mathcal K[\tilde x]$ such that the coefficients of $G'$ lie in $\mathcal R$;
  $\operatorname{val}_{\tilde C_{j,i}}(G')\geq\theta_{j,i}-s_{j,i}-C_6c$ for every chain and
  $\operatorname{val}_{\tilde S_{j,m}}(G')\geq\theta^L_{j,m}-C_{18}c$ for $m+2\leq r$; the coefficients of $R'$
  have valuation at least $c-E(C_E)$; $\varphi$ of every monomial of $G'$ and of $R'$ has degree $d(y)$ in
  each row; $v_{\mathrm F}(\varphi(G'+R'))\geq n_0-\kappa_{\mathcal D}$; and
  $v\big(B_{\mathcal D}(F,\varphi(G'+R'))\big)\leq E(C_E)+\kappa_{\mathcal D}+n_0$.
\end{itemize}
\end{definition}
\leanref{Formalism/SharpRestriction/GrowthBounds/Lemma5aAt.lean}{Formalism.MainNotions.PlaneGramData.SignedChainData}
\leanref{Formalism/SharpRestriction/GrowthBounds/Lemma5aAt.lean}{Formalism.MainNotions.PlaneGramData.SignedChainData.ph}
\leanref{Formalism/SharpRestriction/GrowthBounds/Lemma5aAt.lean}{Formalism.MainNotions.PlaneGramData.SignedChainData.Fdig}
\leanref{Formalism/SharpRestriction/GrowthBounds/Lemma5aAt.lean}{Formalism.MainNotions.PlaneGramData.SignedChainData.Bform}
\leanref{Formalism/SharpRestriction/GrowthBounds/Lemma5aAt.lean}{Formalism.MainNotions.PlaneGramData.SignedChainData.Hyp}
\leanref{Formalism/SharpRestriction/GrowthBounds/Lemma5aAt.lean}{Formalism.MainNotions.PlaneGramData.SignedChainData.XX}
\leanref{Formalism/SharpRestriction/GrowthBounds/Lemma5aAt.lean}{Formalism.MainNotions.PlaneGramData.SignedChainData.YY}

\begin{lemma}\label{pd-symbols-one-sided-real-fixed}
Let the data of Proposition~\ref{pd-symbols-realize} be given, let $g\in GL_k(W)$, and let $\alpha_{\max}$ be as in
Lemma~\ref{pd-symbols-space}; we use the notation of Definitions~\ref{pd-symbols-run-system},
\ref{pd-symbols-module} and~\ref{pd-symbols-twist}, with the integers $a_l(v)$ of
Definition~\ref{pd-symbols-twist}(3). Put $m_1=(2q+3)c'+kq+c'$ and $m_2=c'\alpha_{\max}+2c'$, and for $j\in\mathbf Z$
\[
 D^{\mathrm{one}}_j=\prod_v\det\big(z_{a_l(v),b}\big)_{0\leq l<r,\,1\leq b\leq r}^{\,m_1}\cdot
 \det\big(z_{f(j-1,l),b}\big)_{0\leq l<r,\,1\leq b\leq r}^{\,m_2}\in\mathcal R[z],
\]
the product over the new $v\in\Lambda_j$. Let $F_{c'}$ be as in Lemma~\ref{pd-symbols-realize-core}. Then there are
$H,Y\in\mathcal R[X]$, with $H$ lying on $T$ and involving no normal variable
(Definition~\ref{pd-symbols-realize-keys}), such that
\[
 \prod_{j=0}^{J-1}\phi_j\big(D^{\mathrm{one}}_j\big)=\Sigma_g(F_{c'})\cdot
 \Sigma_g\Big(\prod_{t\in T}\Delta^{\mathrm{top}}_t(g)^{kq}\Big)\cdot H+p^{c'}Y .
\]
\end{lemma}
\leanref{Formalism/SharpRestriction/GrowthBounds/PathWeight/GeneratorImageIsSymbol.lean}{Formalism.MainNotions.G.symbolGeneratorData.Dv}
\leanref{Formalism/SharpRestriction/GrowthBounds/PathWeight/GeneratorImageIsSymbol.lean}{Formalism.MainNotions.G.symbolGeneratorData.RestPh}
\leanref{Formalism/SharpRestriction/GrowthBounds/PathWeight/SymbolGeneratorData.lean}{Formalism.MainNotions.G.symbolGeneratorData.restPh}
\leanref{Formalism/SharpRestriction/GrowthBounds/PathWeight/SymbolGeneratorData.lean}{Formalism.MainNotions.G.symbolGeneratorData.ph_det_pow}
\leanref{Formalism/SharpRestriction/GrowthBounds/PathWeight/SymbolGeneratorData.lean}{Formalism.MainNotions.G.symbolGeneratorData.Spl_pow}
\leanref{Formalism/SharpRestriction/GrowthBounds/PathWeight/PhProdShift.lean}{Formalism.MainNotions.G.symbolGeneratorData.ph_prod_tops}

\begin{proof}
Let $S\subset\mathcal R[X]$ be the $\mathcal R$-subalgebra of the polynomials that lie on $T$ and involve no normal
variable; it contains $p$. By Step~4 of the proof of Proposition~\ref{pd-symbols-realize}, for every integer $t$
and every $j$, $D^{\mathrm{top}}_t=\Sigma_g(\Delta^{\mathrm{top}}_t(g))$ and $D_j=\Sigma_g(\Delta_j(g))$ are the
determinants of the matrices with entries $X_{\overline{a_l(t)},b}^{\,p^{t-a_l(t)}}$ and
$X_{\overline{f(j,l)},b}^{\,p^{f(j,r-1)-f(j,l)}}$, $0\leq l<r$, $1\leq b\leq r$; as the $a_l(t)$ and the $f(j,l)$ are
even with residues in $T$ (fact (b) of that proof and condition~(3) of validity), they lie in $S$.

\emph{Step 1: two facts.} (a) For a square matrix $N$ over a commutative ring and $e\geq0$, let $N^{[p^e]}$ have
the $p^e$-th powers of the entries of $N$. Then $\det(N^{[p^e]})=\det(N)^{p^e}+ph$ with $h$ a polynomial with integer
coefficients in the entries of $N$: in the polynomial ring over $\mathbf Z$ in the entries, the difference vanishes
modulo $p$, as raising to the $p^e$-th power is a ring endomorphism in characteristic $p$. (b) Let $D,h\in S$,
$P\geq1$ and $n,c\geq0$. Then $(D^P+ph)^n=D^{\max(n-c,0)}H'+p^cY'$ with $H',Y'\in S$. Indeed, in the binomial
expansion the terms $\binom ni D^{P(n-i)}(ph)^i$ with $i\geq c$ lie in $p^cS$, and those with $i<c$ have
$P(n-i)\geq n-i\geq\max(n-c,0)$. (c) If $f_i=a_iH_i+p^cY_i$ with $a_i,H_i,Y_i\in S$ for finitely many $i$, then
$\prod_if_i=(\prod_ia_i)H+p^cY$ with $H,Y\in S$, by induction from
$(aH_1+p^cY_1)(bH_2+p^cY_2)=abH_1H_2+p^c(aH_1Y_2+bY_1H_2+p^cY_1Y_2)$.

\emph{Step 2: the factors.} Let $j\in\mathbf Z$ and $v$ a new element of $\Lambda_j$, and write $v=s+n'u$ with
$0\leq s<n'$. The definition of the $a_l$ is invariant under translation by multiples of $n'$, so
$a_l(v)=a_l(s)+n'u$; and $a_l(v)<v\leq f(j,r-1)$ by fact (b), so
$\operatorname{dist}_j(a_l(v))=(v-a_l(v))+\operatorname{dist}_j(v)$. Hence
$\phi_j\big(\det(z_{a_l(v),b})\big)$ is the determinant of the matrix with entries
$\big(X_{\overline{a_l(s)},b}^{\,p^{s-a_l(s)}}\big)^{p^{\operatorname{dist}_j(v)}}$, which by (a) is
$(D^{\mathrm{top}}_s)^{p^{\operatorname{dist}_j(v)}}+ph_v$ with $h_v\in S$. In the same way, as
$f(j-1,l)\leq f(j-1,r-1)<f(j,r-1)$, $\phi_j\big(\det(z_{f(j-1,l),b})\big)=D_{j-1}^{\,q_{j-1}}+ph'_j$ with $h'_j\in S$
and $q_{j-1}=p^{f(j,r-1)-f(j-1,r-1)}$. As $m_1,m_2\geq c'$, (b) with $c=c'$ and (c) give
\[
 \prod_{j=0}^{J-1}\phi_j\big(D^{\mathrm{one}}_j\big)=\prod_{j=0}^{J-1}\prod_v\big(D^{\mathrm{top}}_{\bar v}\big)^{m_1-c'}
 \cdot\prod_{j=0}^{J-1}D_{j-1}^{\,m_2-c'}\cdot H_0+p^{c'}Y_0
\]
with $H_0,Y_0\in S$, where $D^{\mathrm{top}}_{\bar v}$ is formed with the representative of $\bar v$.

\emph{Step 3: the products.} The new elements of $\Lambda_j$, $0\leq j<J$, are the $v$ with
$f(-1,r-1)<v\leq f(J-1,r-1)$ and $\bar v\in T$, and $f(J-1,r-1)=f(-1,r-1)+n'$ by condition~(1) of validity; so their
residues run through $T$, each once, and the first double product is $\prod_{t\in T}(D^{\mathrm{top}}_t)^{m_1-c'}$. By
condition~(1) the entries of the matrix of $D_j$ depend only on $j$ modulo $J$, so $D_{-1}=D_{J-1}$ and the second
product is $\prod_{j=0}^{J-1}D_j^{\,m_2-c'}$. Now $m_1-c'=(2q+3)c'+kq$, and $m_2-c'=c'(\alpha_j+1)+c'(\alpha_{\max}-\alpha_j)$
with $\alpha_{\max}-\alpha_j\geq0$ for $0\leq j<J$. As $\Sigma_g$ is a ring homomorphism,
\[
 \prod_{t\in T}(D^{\mathrm{top}}_t)^{m_1-c'}\prod_{j=0}^{J-1}D_j^{\,m_2-c'}=\Sigma_g(F_{c'})\cdot
 \Sigma_g\Big(\prod_{t\in T}\Delta^{\mathrm{top}}_t(g)^{kq}\Big)\cdot\prod_{j=0}^{J-1}D_j^{\,c'(\alpha_{\max}-\alpha_j)} .
\]
So the statement holds with $H=\prod_{j<J}D_j^{\,c'(\alpha_{\max}-\alpha_j)}H_0\in S$ and $Y=Y_0$.
% Lean: restPh (for the Ctx X of the data, with Dv X.Fr j O K (X.m1 c') (X.m2 c')): Step 1 is ph_det_pow and
%   ph_det_pow_sub (a), Spl_pow = CountingSumsFrames.Top.powerSplit (b), Spl_mul and Spl_prod (c), in the subalgebra
%   SS = supported O (bc_St X.Fr (r + K)); Step 2 is ph_Dseg (topPairingMatrix, ph_Arow_shift, Frames.dist) and
%   ph_Wseg (Qj), with DsO and DlO mapping to sb (Dtop) and sb (Del) by ph_DsO_map and ph_DlO_map; Step 3 is
%   ph_prod_tops, ph_prod_shift with ph_DlO_period, and E1 to E3 (R0 = prod Del^alpha); H lies in bc_IT by
%   bc_IT_of_supp.
\end{proof}

\begin{lemma}\label{pd-symbols-division-rows}
Let $n'\geq2$ be an even integer, $r\geq1$ and $L_1,N_1\geq0$ integers, $T\subset\mathbf Z/n'\mathbf Z$, and $(J,f)$ a
run system for $T$ of length $r$ valid for $(L_1,N_1)$, with the notation of
Definition~\ref{pd-symbols-run-system}. For $v\in\mathbf Z$ let $a_0(v)>a_1(v)>\cdots$ be the even integers $y$ with
$v-rn'\leq y<v$ and $\bar y\in T$, and put $a_l(v)=v$ when there are at most $l$ of them, as in
Definition~\ref{pd-symbols-twist}(3). Let $j\in\mathbf Z$ and let $x\in\Lambda_j$ be new. Then:
\begin{enumerate}
\item the elements of $\Lambda_j$ that are not new are $f(j-1,0)<\dots<f(j-1,r-1)$, and they are even;
\item $a_0(x)>\dots>a_{r-1}(x)$ are even elements of $\Lambda_j$ smaller than $x$;
\item $\operatorname{old}_j(x)$ is an even element of $\Lambda_j$ smaller than $x$, and every even $y\in\Lambda_j$
  with $y<x$ that is not among $a_0(x),\dots,a_{r-1}(x)$ satisfies $y<\operatorname{old}_j(x)$ and is either not
  new or an element of $\Lambda'_j$.
\end{enumerate}
\end{lemma}
\leanref{Formalism/SharpRestriction/GrowthBounds/PathWeight/DvDegOtherMono.lean}{Formalism.MainNotions.G.symbolGeneratorData.dv_W_frame}
\leanref{Formalism/SharpRestriction/GrowthBounds/PathWeight/DvDegOtherMono.lean}{Formalism.MainNotions.G.symbolGeneratorData.dv_Arow_letter}
\leanref{Formalism/SharpRestriction/GrowthBounds/PathWeight/GeOldest.lean}{Formalism.MainNotions.G.symbolGeneratorData.ge_oldest}
\leanref{Formalism/SharpRestriction/GrowthBounds/PathWeight/GeOldest.lean}{Formalism.MainNotions.G.symbolGeneratorData.ge_older}

\begin{proof}
For integers $y<x$ let $N(y)$ be the number of even $t$ with $y<t<x$ and $\bar t\in T$. If $y_1<y_2<x$ and $y_2$ is
even with $\bar y_2\in T$, then $N(y_2)<N(y_1)$; so $N$ is injective on the even $y<x$ with $\bar y\in T$.

(1) By conditions (2) and (3) of validity, $f(j-1,0)<\dots<f(j-1,r-1)<f(j,0)\leq f(j,r-1)$, and the $f(j-1,l)$ are
even with residues in $T$; so they lie in $\Lambda_j$ and are not new. Conversely, let $y\in\Lambda_j$ with
$y\leq f(j-1,r-1)$. As $y\geq f(j-1,0)$, $\bar y\in T$ and no $v$ strictly between $f(j-1,i)$ and $f(j-1,i+1)$ has
$\bar v\in T$, $y$ is one of the $f(j-1,l)$.

(2) The integers congruent to $f(0,0)$ modulo $n'$ are even, as $n'$ is, and have residue in $T$; each interval
$[x-(i+1)n',x-in')$, $i<r$, contains one of them. So $a_0(x)>\dots>a_{r-1}(x)$ are even integers in $[x-rn',x)$ with
residues in $T$, and $N(a_l(x))=l$. Suppose $a_l(x)<f(j-1,0)$ for some $l<r$. As $x$ is new, each $f(j-1,i)$, $i<r$, lies in
$(a_l(x),x)$, so it is an even element of $[x-rn',x)$ with residue in $T$ larger than $a_l(x)$, that is, one of
$a_0(x),\dots,a_{l-1}(x)$; these are fewer than $r$, which is impossible. So $a_l(x)\geq f(j-1,0)$, and as
$a_l(x)<x\leq f(j,r-1)$, $a_l(x)\in\Lambda_j$.

(3) By (1), the $r$ elements $f(j-1,l)$ of $\Lambda_j$ are even and smaller than $x$. So $o=\operatorname{old}_j(x)$ is
the $r$-th largest even element of $\Lambda_j$ smaller than $x$. Every $t$ with $o<t<x$ and $\bar t\in T$ lies in
$\Lambda_j$, so $N(o)=r-1$. Let $y$ be as in (3). If $N(y)<r$, then $N(y)=N(a_{N(y)}(x))$, so $y=a_{N(y)}(x)$ by
injectivity, which is excluded. So $N(y)\geq r>N(o)$, and $y<o$. If $y<f(j,0)$, then either $y\leq f(j-1,r-1)$ and $y$
is not new, or $f(j-1,r-1)<y<f(j,0)$ and $y\in\Lambda'_j$. If $y\geq f(j,0)$, every even $t$ with $y<t<x$ and
$\bar t\in T$ satisfies $f(j,0)<t<f(j,r-1)$, so it is one of $f(j,1),\dots,f(j,r-2)$ by condition~(3) of validity;
then $N(y)\leq r-2$, a contradiction.
% Lean: (1) is dv_frame_letter and dv_W_frame; (2) is window_holds_even_tops (the r translates of fr 0 0),
%   dv_Arow_anti, dv_Arow_newer and dv_Arow_ge, giving dv_Arow_letter; (3) is ge_oldest (bd_vectors_before, ge_sort)
%   and ge_older (geN = N, ge_N_lt, ge_N_inj, ge_Arow, ge_frame).
\end{proof}

\begin{lemma}\label{pd-symbols-division-numeric}
There is $C_N\geq1$ such that the following holds for every $C\geq C_N$. Let the data of
Situation~\ref{pd-edge-situation} be given with this $C$, let $(J,f)$ be a run system for $T$ of length $r$ spaced
for $(L_1,L_2,N_1,G')$, put $K=k-r$, $c'=c_*+\omega$ and $\deg(s)=\lambda_0(s)_1-\lambda_0(s)_2$, let $q$,
$\alpha_{\max}$ and $m_*$ be as in Definition~\ref{pd-symbols-test-one-sided} and $B_j$ as in
Definition~\ref{pd-symbols-module}(4), and put $m_1=(2q+3)c'+kq+c'$ and $m_2=c'\alpha_{\max}+2c'$. Then for every
$j\in\mathbf Z$:
\begin{enumerate}
\item $|\Lambda_j|m_1+m_2\leq m_*$;
\item $m_*\leq\lfloor B_j(x)/(Yk^2)\rfloor\leq B_j(x)/100$ for every new $x\in\Lambda_j$.
\end{enumerate}
\end{lemma}
\leanref{Formalism/SharpRestriction/GrowthBounds/WittFrameLocaliserOrders.lean}{Formalism.MainNotions.G.WittFrameLocaliserOrders.segment_divisions_le_bound}
\leanref{Formalism/SharpRestriction/GrowthBounds/WittFrameLocaliserOrders.lean}{Formalism.MainNotions.G.WittFrameLocaliserOrders.numeric_symbol_bound_real}
\leanref{Formalism/SharpRestriction/GrowthBounds/WittFrameLocaliserOrders.lean}{Formalism.MainNotions.G.WittFrameLocaliserOrders.segment_length_bounds_real}
\leanref{Formalism/SharpRestriction/GrowthBounds/ScaleEdges/GainBoundAt.lean}{Formalism.MainNotions.Pf14.hypII_at}

\begin{proof}
(1) A spaced run system is valid, so $|\Lambda_j|\leq N_1$, and $m_*=(N_1+1)(m_1+m_2)$.

(2) This is a self-contained computation. Put $o=\operatorname{old}_j(x)$ and $\hat a=\lceil a\rceil$. By
Lemma~\ref{pd-symbols-division-rows}(3), $o\in\Lambda_j$, so $\bar o\in T$, and $\deg(\bar o)\geq a$ by (2) of the
situation, as in the proof of Lemma~\ref{pd-symbols-space}(1). As $\deg(\bar o)$ is an integer,
$\deg(\bar o)\geq\hat a$. And $\operatorname{dist}_j(o)\leq f(j,r-1)-f(j-1,0)<G'$. So $B_j(x)\geq\hat a/q$. The last
inequality of (3) of the situation gives $a\geq2$, so $a\leq\hat a\leq a^2$ and $\log\hat a\leq2\log a$. Let
$C_6\geq1$ be a real number and $C\geq1+2|C_6|+2|C_6+1|+|C_6|$. Then (2) and (3) of the situation give the same
conditions with $\hat a$ in place of $a$: $\hat a\leq\deg(t)\leq2\hat a$ for $t\in T$, $C_6N_1^2\log\hat a\leq k$,
$(C_6+1)|T|k^2\log_p\hat a\leq D_p$ and $C_6k^3r^2|T|\,p^{5G'}(c_*+\omega)Y\leq\hat a$, together with the other
conditions of (3). There is a $C_6$ for which these give, for every real $B\geq\hat a/q$, the inequalities
$m_*\leq\lfloor B/(Yk^2)\rfloor\leq B/100$. Take $C_N=1+2|C_6|+2|C_6+1|+|C_6|$ for this $C_6$.
% Lean: (1) is segment_divisions_le_bound (divisions_le_total_budget with the bound on the letters of validity).
%   (2) is the first part of numeric_symbol_bound_real = segment_length_bounds_real = slope_bounds_on_letters =
%   slope_bounds_alt (level_degree_ge_ratio for B >= ceil(a) / q, CountingSumsFrames.Top.paramsII0 =
%   secondParameterSizeBounds_holds through hypothesis_numeric_bounds), stated under HypII C6 (C6 + 1) C6 for the
%   constant C6 of numeric_symbol_bound_real, whose a is a natural number; hypII_at C6 (C6 + 1) C6 gives
%   C_N = 1 + 2|C6| + 2|C6 + 1| + |C6| and HypII from Hyps C with a replaced by its ceiling (hca1, hcsq, hlogc,
%   hlogbc; flatDeg_bounds and Nat.ceil_le for the degrees), c'' = 0 and cStar + 2 r n' cpp in place of cStar (so
%   that cPrime is c'). TwistBound p F q comes from the spacing (Lemma pd-symbols-space (4) and fact (b) in the proof of
%   Proposition pd-symbols-realize).
\end{proof}

\begin{lemma}\label{pd-symbols-division-domain}
Let $\mathcal R$ be an integral domain, $N\geq0$ an integer with a lettering of $\{0,\dots,N-1\}$, $r\geq1$, and
$\mathcal R[z]$ the polynomial ring in the $z_{i,b}$, $0\leq i<N$, $1\leq b\leq r$, with $\operatorname{val}$ formed with
$n=r$ (Definition~\ref{pd-symbols-path-generators}(1)), and $\Delta_A$ and $d_A(S)$ as in
Lemma~\ref{pd-symbols-path-clauses}(6). Let $(S_u,A_u,n_u)$, $1\leq u\leq t$, be triples of a set
$S_u\subset\{0,\dots,N-1\}$, an injective map $A_u\colon\{1,\dots,r\}\to S_u$ whose values are vector letters, and
an integer $n_u\geq0$, and put $n=n_1+\dots+n_t$. Let $G\in\mathcal R[z]$ have row degrees $d$, and suppose that
$\operatorname{val}_{S_u}(G)\geq d_{A_u}(S_u)+n$ for every $u$, with $d_{A_u}(S_u)$ formed with $d$. Then
$G=\prod_u\Delta_{A_u}^{n_u}\cdot Q$ for some $Q\in\mathcal R[z]$ with fixed row degrees, and
$\operatorname{val}_{S'}(Q)\geq\operatorname{val}_{S'}(G)-n$ for every $S'\subset\{0,\dots,N-1\}$.
\end{lemma}
\leanref{Formalism/SharpRestriction/GrowthBounds/PathWeight/DvDegOtherMono.lean}{Formalism.MainNotions.G.symbolGeneratorData.dv_div_one}
\leanref{Formalism/SharpRestriction/GrowthBounds/PathWeight/DvDegOtherMono.lean}{Formalism.MainNotions.G.symbolGeneratorData.dv_div_list}
\leanref{Formalism/SharpRestriction/GrowthBounds/WittFrameLocaliserOrders.lean}{Formalism.MainNotions.G.WittFrameLocaliserOrders.gauss_quotient_integral}
\leanref{Formalism/SharpRestriction/GrowthBounds/WittFrameLocaliserOrders.lean}{Formalism.MainNotions.G.WittFrameLocaliserOrders.valuation_fraction_field_extension}

\begin{proof}
For a map $A$ as in the statement let $e_A(i)$ be $1$ if $i$ is a value of $A$ and $0$ otherwise; $\Delta_A$ has row
degrees $e_A$. Let $\mathcal K$ be the fraction field of $\mathcal R$. The substitution that defines $P_{S,r}$
commutes with the inclusion $\mathcal R[z]\subset\mathcal K[z]$, which is injective on coefficients; so
$\operatorname{val}_S(P)$ is the same over $\mathcal R$ and over $\mathcal K$.

\emph{Step 1: one triple.} Let $S$, $A$ be as in the statement, $n_1\geq0$, and $G\in\mathcal R[z]$ of row degrees
$d$ with $\operatorname{val}_S(G)\geq d_A(S)+n_1$. We show that $G=\Delta_A^{n_1}Q$ with $Q\in\mathcal R[z]$ of row
degrees $d-n_1e_A$ and $\operatorname{val}_{S'}(Q)\geq\operatorname{val}_{S'}(G)-n_1$ for every $S'$. For $n_1=0$ take
$Q=G$. Let $n_1\geq1$. Lemma~\ref{pd-symbols-path-clauses}(6)(a) over $\mathcal K$, with $\theta=d_A(S)+n_1$, shows
that $\Delta_A^{n_1}$ divides $G$ in $\mathcal K[z]$. Order the monomials lexicographically, for some ordering of the
variables. The monomials $\prod_lz_{A(l),\sigma(l)}$ of $\Delta_A$, $\sigma$ a permutation, are distinct as $A$ is
injective, so its coefficients are $\pm1$; hence the leading coefficient of $\Delta_A^{n_1}$ is a unit of $\mathcal R$.
Division with remainder by $\Delta_A^{n_1}$, which inverts only this coefficient, gives $G=\Delta_A^{n_1}Q+R'$ with
$Q,R'\in\mathcal R[z]$ and no monomial of $R'$ divisible by the leading monomial of $\Delta_A^{n_1}$. If $R'\neq0$,
then $\Delta_A^{n_1}$ divides $R'$ in $\mathcal K[z]$, and the leading monomial of $R'$ is the product of the leading
monomials of $\Delta_A^{n_1}$ and of the quotient, which is impossible. So $G=\Delta_A^{n_1}Q$. As $\mathcal R$ is a
domain and $\Delta_A^{n_1}\neq0$ has row degrees $n_1e_A$, comparing the parts of given row degrees shows that $Q$
has row degrees $d-n_1e_A$. Finally, Lemma~\ref{pd-symbols-path-clauses}(6)(b) over $\mathcal K$ gives
$\operatorname{val}_{S'}(G)\leq\operatorname{val}_{S'}(Q)+n_1$.

\emph{Step 2: induction on $t$.} For $t=0$ take $Q=G$. Let $t\geq1$. As $n_1\leq n$, Step~1 applies to the first
triple and gives $G=\Delta_{A_1}^{n_1}Q_1$ with $Q_1$ of row degrees $d_1=d-n_1e_{A_1}\leq d$ and
$\operatorname{val}_{S'}(Q_1)\geq\operatorname{val}_{S'}(G)-n_1$. The sum $d_A(S)$ does not decrease when $d$ does
not, so for $u\geq2$, $\operatorname{val}_{S_u}(Q_1)\geq d_{A_u}(S_u)+n-n_1$ with $d_{A_u}(S_u)$ formed with $d_1$.
The induction hypothesis for the other $t-1$ triples gives $Q_1=\prod_{u\geq2}\Delta_{A_u}^{n_u}\cdot Q$ with
$\operatorname{val}_{S'}(Q)\geq\operatorname{val}_{S'}(Q_1)-(n-n_1)\geq\operatorname{val}_{S'}(G)-n$.
% Lean: Step 1 is dv_div_one (valuation_fraction_field_extension, the clause hE = FieldMultidegreeDivisibility
%   over FractionRing R0, gauss_quotient_integral = gauss_of_det_coefficients determinant_coefficients_units with
%   dvd_of_dvd_fraction_field_unit_lead for the lex order, dv_quot_hom with dv_detA_hom and dv_detA_ne); Step 2 is
%   dv_div_list (dv_degOther_mono).
\end{proof}

\begin{lemma}\label{pd-symbols-one-sided-real-division}
Let $C_{La},C_k,C_{rk},C_{\deg}$ be real numbers with the property of Theorem~\ref{pd-symbols-count}. There is
$C_9>0$ such that the following holds for every $C\geq C_9$. Let the data of Situation~\ref{pd-edge-situation} be
given with this $C$, let $(J,f)$ be a run system for $T$ of length $r$ spaced for $(L_1,L_2,N_1,G')$, put $K=k-r$,
$c'=c_*+\omega$ and $\deg(s)=\lambda_0(s)_1-\lambda_0(s)_2$, and let $q$, $\alpha_{\max}$, $m_*$ and $\theta^1_j$ be as
in Definition~\ref{pd-symbols-test-one-sided} for these data and $Y$, and $D^{\mathrm{one}}_j$ as in
Lemma~\ref{pd-symbols-one-sided-real-fixed}. Let $0\leq j<J$ and $b\colon\mathbf Z\to\mathbf Z_{\geq0}$ with
$b(x)\leq\frac{101}{100}\deg(\bar x)/p^{\operatorname{dist}_j(x)}$ for every $x\in\Lambda_j$ that is not new or lies
in $\Lambda'_j$. Let $AG$ be a product as in Definition~\ref{pd-symbols-module}(2) for $(j,m_*,\theta^1_j)$, with
$m_*$ the constant function, such that for every $x\in\Lambda_j$, $A$ has a fixed degree $d'(x)$ and $G$ a fixed
degree $d''(x)$ in the row $z_x$, with $d'(x)+d''(x)=b(x)$. Then $AG=D^{\mathrm{one}}_jA'G'$ for a polynomial $A'$ over
$\mathcal R$ in the plane variables $z_{v,c}$, $v\in\Lambda_j$, and a polynomial $G'$ over $\mathcal R$ in the normal
variables $z_{v,c}$, $v\in\Lambda_j$ new, each with a fixed degree in every row, such that $A'G'$ satisfies the
inequalities of Definition~\ref{pd-symbols-module}(2) with $m=0$ for some $M$.
\end{lemma}
\leanref{Formalism/SharpRestriction/GrowthBounds/DualFockLinear/OneSidedLemma4OfAdmittedTheoremII.lean}{Formalism.MainNotions.TestSideCounting.RestDivMU}
\leanref{Formalism/SharpRestriction/GrowthBounds/DualFockLinear/OneSidedLemma4OfAdmittedTheoremII.lean}{Formalism.MainNotions.TestSideCounting.restDivMU}
\leanref{Formalism/SharpRestriction/GrowthBounds/PathWeight/GeneratorImageIsSymbol.lean}{Formalism.MainNotions.G.symbolGeneratorData.MGen}

\begin{proof}
Let $C_9=C_N$ be the constant of Lemma~\ref{pd-symbols-division-numeric}. Let $C\geq C_9$ and let the data, $j$, $b$
and $AG$ be as in the statement, with the function $M$ of Definition~\ref{pd-symbols-module}(2) for $AG$. Let
$\mathcal N$ be the set of the new elements of $\Lambda_j$, let $a_l(x)$ be as in
Lemma~\ref{pd-symbols-division-rows}, put $m_1=(2q+3)c'+kq+c'$, $m_2=c'\alpha_{\max}+2c'$ and
$s_j(x)=\lfloor B_j(x)/(Yk^2)\rfloor$, so that $\theta^1_j(x)=B_j(x)/2-s_j(x)$, and for $x\in\mathcal N$ let $E_x$ be
the set of the even $y\in\Lambda_j$ with $y<x$ that are not among $a_0(x),\dots,a_{r-1}(x)$.

\emph{Step 1: the degrees below $x$.} Let $x\in\mathcal N$; we show $\sum_{y\in E_x}b(y)\leq B_j(x)/80$. Put
$o=\operatorname{old}_j(x)$ and $\hat a=\lceil a\rceil$. By (2) of the situation, $a\leq\deg(t)\leq2a$ for $t\in T$, as
in the proof of Lemma~\ref{pd-symbols-space}(1). As $\deg(t)$ is an integer and $a\leq\hat a$, this gives
$\hat a\leq\deg(t)\leq2\hat a$ for $t\in T$. By Lemma~\ref{pd-symbols-division-rows}(3), each $y\in E_x$ satisfies
$y<o$ and is not new or lies in $\Lambda'_j$, so $b(y)\leq\frac{101}{100}\deg(\bar y)/p^{\operatorname{dist}_j(y)}$,
and hence $b(y)\leq2.02\,\hat a/p^{\operatorname{dist}_j(y)}$. As $y<o\leq f(j,r-1)$ are both even,
$\operatorname{dist}_j(y)=\operatorname{dist}_j(o)+2m_y$ with an integer $m_y\geq1$ that determines $y$. As
$\bar o\in T$, $\hat a/p^{\operatorname{dist}_j(o)}\leq\deg(\bar o)/p^{\operatorname{dist}_j(o)}=B_j(x)$. Hence
\[
 \sum_{y\in E_x}b(y)\leq2.02\,B_j(x)\sum_{m\geq1}p^{-2m}=\frac{2.02\,B_j(x)}{p^2-1}\leq\frac{B_j(x)}{80},
\]
as $p\geq13$ gives $p^2-1\geq168$.

\emph{Step 2: the valuations of $A$.} Let $x\in\mathcal N$. By Lemma~\ref{pd-symbols-division-numeric}(2),
$m_*\leq s_j(x)\leq B_j(x)/100$, and $B_j(x)\geq0$. With Step~1 and $d'\leq b$,
\[
 \operatorname{val}_{\{y<x\}}(A)\geq\frac{B_j(x)}2-s_j(x)\geq\frac{B_j(x)}4+m_*+\frac{B_j(x)}{80}
 \geq\sum_{y\in E_x}d'(y)+m_* .\tag{$*$}
\]

\emph{Step 3: division.} Let $e_0<\dots<e_{N-1}$ be the elements of $\Lambda_j$, with the lettering of Step~4 of the
proof of Lemma~\ref{pd-symbols-realize-segments}, in which $i$ is a vector letter if $e_i$ is even. Renaming $z_{e_i,b}$
as $z_{i,b}$ makes $A$ a polynomial of row degrees $i\mapsto d'(e_i)$ in the $z_{i,b}$, $b\leq r$, and by fact (iii)
there the valuations for the sets $\{y<e_i\}$ and $\{y\leq e_i\}$ become those for $\{0,\dots,i-1\}$ and
$\{0,\dots,i\}$. For $x=e_i\in\mathcal N$ let $S_x=\{0,\dots,i-1\}$ and let $A_x(l)$ be the index of $a_{l-1}(x)$;
let $x_0=e_{i_0}$ be the smallest element of $\mathcal N$, which exists as $f(j,r-1)\in\mathcal N$, put
$S_0=\{0,\dots,i_0-1\}$, and let $A_0(l)$ be the index of $f(j-1,l-1)$. By Lemma~\ref{pd-symbols-division-rows}(1)
and (2), $A_x$ and $A_0$ are injective with values vector letters in $S_x$ and in $S_0$. The vector letters of $S_x$
that are not values of $A_x$ are the indices of the elements of $E_x$, so $d'_{A_x}(S_x)=\sum_{y\in E_x}d'(y)$; and the
even elements of $\Lambda_j$ below $x_0$ are not new, so they are values of $A_0$ by
Lemma~\ref{pd-symbols-division-rows}(1), and $d'_{A_0}(S_0)=0$. Take the triples $(S_x,A_x,m_1)$, $x\in\mathcal N$, and
$(S_0,A_0,m_2)$, with $n=|\mathcal N|m_1+m_2\leq|\Lambda_j|m_1+m_2\leq m_*$ by
Lemma~\ref{pd-symbols-division-numeric}(1). By $(*)$, $\operatorname{val}_{S_x}(A)\geq d'_{A_x}(S_x)+n$ and
$\operatorname{val}_{S_0}(A)\geq m_*\geq d'_{A_0}(S_0)+n$. Lemma~\ref{pd-symbols-division-domain} over the domain
$\mathcal R$ gives
\[
 A=\prod_{x\in\mathcal N}\Delta_{A_x}^{m_1}\cdot\Delta_{A_0}^{m_2}\cdot A'
\]
with $A'\in\mathcal R[z]$ of fixed row degrees and $\operatorname{val}_S(A')\geq\operatorname{val}_S(A)-m_*$ for
every $S$.

\emph{Step 4: the conclusion.} In the variables $z_{v,b}$, $\Delta_{A_x}=\det(z_{a_l(x),b})_{0\leq l<r,\,1\leq b\leq r}$
and $\Delta_{A_0}=\det(z_{f(j-1,l),b})_{0\leq l<r,\,1\leq b\leq r}$, so the product of the determinants is
$D^{\mathrm{one}}_j$, and $AG=D^{\mathrm{one}}_jA'G'$ with $G'=G$. The polynomial $A'$ involves only the plane
variables $z_{v,c}$ with $v\in\Lambda_j$, and $G'$ only the normal variables of the new rows, each with a fixed degree
in every row. For an even new $v$, $\lambda_{\{y\leq v\}}(G')\leq M(v)$ and
$M(v)\leq\operatorname{val}_{\{y<v\}}(A)-m_*\leq\operatorname{val}_{\{y<v\}}(A')$; for an odd new $v$ the same holds
with $\lambda_{\{y<v\}}(G')$ and $\operatorname{val}_{\{y\leq v\}}(A')$. So $A'G'$ satisfies the inequalities of
Definition~\ref{pd-symbols-module}(2) with $m=0$ and the same $M$.
% Lean: restDivMU hE restGeomU, with hE : FieldMultidegreeDivisibility (Lemma pd-symbols-path-clauses (6), the
%   clause h21 of tableLineDegreeFive.admitted; one_sided_lemma4_of_admitted_lemmas splits on it, and its negation
%   contradicts h21). The constant is that of numeric_symbol_bound_real, and HypII comes from hypII_at (Lemma
%   pd-symbols-division-numeric). Step 1 is restGeomU (ge_oldest, ge_older, hlev with deg <= 2 ceil(a), and
%   CountingSumsFrames.Top.lemma4Geom: 2.02 B / (p^2 - 1) <= B / 80); Step 2 is the second part of
%   numeric_symbol_bound_real (symbol_valuation_margin = Top.lemma4Arith) with dv_degOther_le; Step 3 is dv_segment
%   (the tasks ({y < i}, Ai i, m1) with eL (Ai i l) = Arow F (eL i) l and ({y < i0}, Aw, m2) with eL (Aw l) =
%   fr (j - 1) l, hms from segment_divisions_le_bound, the vanishing of otherVectorDegree for Aw, and dv_div_list);
%   Step 4 is hC' (CondM 0) and hDv (dv_plnm_det, dv_nonW_image). The book's A'G' is the Lean's P' in MGen0, and the
%   hypothesis on AG is membership in MGen with m = mStar, th = Bo / 2 - sl and dd = b.
\end{proof}

\begin{lemma}\label{pd-symbols-test-one-sided-real}
Let $C_{La},C_k,C_{rk},C_{\deg}$ be real numbers with the property of
Theorem~\ref{pd-symbols-count}. There is $C_8>0$ such that the following holds for every $C\geq C_8$. Let the
data of Situation~\ref{pd-edge-situation} be given with this $C$, let $(J,f)$ be a run system for $T$ of length
$r$ spaced for $(L_1,L_2,N_1,G')$, and put $K=k-r$, $c'=c_*+\omega$ and $\deg(s)=\lambda_0(s)_1-\lambda_0(s)_2$.
Let $\sigma$ be a one-sided element of normal degree $o$ for these data and $Y$
(Definition~\ref{pd-symbols-test-one-sided}). Then $\sigma$ is realizable at $(g,c',o)$
(Definition~\ref{pd-symbols-twist}) for every $g\in GL_k(W)$.
\end{lemma}
\leanref{Formalism/SharpRestriction/GrowthBounds/DualFockLinear/OneSidedLemma4Realisable.lean}{Formalism.MainNotions.TestSideCounting.one_sided_lemma4_realisable}
\leanref{Formalism/SharpRestriction/GrowthBounds/DualFockLinear/OneSidedLemma4OfAdmittedTheoremII.lean}{Formalism.MainNotions.TestSideCounting.one_sided_lemma4_of_admitted_theorem_II}

\begin{proof}
Let $C_9$ be the constant of Lemma~\ref{pd-symbols-one-sided-real-division} and $C_8=\max(30,C_9)$. Let $C\geq C_8$,
the data and $\sigma$ be as in the statement, and $g\in GL_k(W)$. As $C\geq30$, these data, with the integers of
Lemma~\ref{pd-symbols-space} and the symbol spaces formed from them, are data of
Proposition~\ref{pd-symbols-realize}; let $\mathcal S_g$ be the symbol set of
Definition~\ref{pd-symbols-realize-keys} for $T$, $q$ and $g$, and $D^{\mathrm{one}}_j$ and $F_{c'}$ as in
Lemma~\ref{pd-symbols-one-sided-real-fixed}. Write $\sigma=X^\ell\sum_{i=1}^n\prod_{j=0}^{J-1}\phi_j(\xi_{i,j})$ with
$o_j$ and $b_j$ as in Definition~\ref{pd-symbols-test-one-sided}.

\emph{Step 1: parts of given row degrees.} Fix $i$ and $0\leq j<J$. By Definition~\ref{pd-symbols-module}(2),
$\xi_{i,j}$ is an $\mathcal R$-linear combination of products $AG$ as there for $(j,m_*,\theta^1_j)$. For functions
$d',d''$ on $\Lambda_j$ let $A_{d'}$ and $G_{d''}$ be the sums of the terms of $A$ and of $G$ of degree $d'(x)$,
respectively $d''(x)$, in the row $z_x$ for every $x\in\Lambda_j$. By Step~1 of the proof of
Lemma~\ref{pd-symbols-realize-segments}, $\operatorname{val}_S(A_{d'})\geq\operatorname{val}_S(A)$ and
$\lambda_S(G_{d''})\leq\lambda_S(G)$ for every $S$, so each $A_{d'}G_{d''}$ satisfies the conditions of
Definition~\ref{pd-symbols-module}(2) for $(j,m_*,\theta^1_j)$ with the same $M$ as $AG$. As $A$ and $G$ involve only
rows $z_x$ with $x\in\Lambda_j$, the part of $A_{d'}G_{d''}$ of row degrees $b_j$ on $\Lambda_j$ is $A_{d'}G_{d''}$ if
$d'+d''=b_j$ and $0$ otherwise. Every monomial of $\xi_{i,j}$ has these row degrees, so $\xi_{i,j}$ is an
$\mathcal R$-linear combination of products $A_{d'}G_{d''}$ with $d'+d''=b_j$.

\emph{Step 2: division.} By Lemma~\ref{pd-symbols-one-sided-real-division}, each such product is
$D^{\mathrm{one}}_jA'G'$ with $A'G'$ a generator in the sense of Step~2 of the proof of
Lemma~\ref{pd-symbols-realize-segments}. By Steps 4 to 6 of that proof, $\phi_j(A'G')\in\mathcal S_g$, and by its
Step~3, $\mathcal S_g$ is closed under sums and under multiplication by elements of $\mathcal R$. Hence
$\xi_{i,j}=D^{\mathrm{one}}_j\xi'_{i,j}$ with $\xi'_{i,j}\in\mathcal R[z]$ and $\phi_j(\xi'_{i,j})\in\mathcal S_g$.

\emph{Step 3: the fixed factor.} Put $Z_1=\sum_i\prod_{j<J}\phi_j(\xi'_{i,j})$; it lies in $\mathcal R[X]$ and, by
Lemma~\ref{pd-symbols-realize-closure}(1), in $\mathcal S_g$. As the $\phi_j$ are ring homomorphisms,
$\sigma=X^\ell\prod_{j<J}\phi_j(D^{\mathrm{one}}_j)\cdot Z_1$, and Lemma~\ref{pd-symbols-one-sided-real-fixed} gives
$H$ and $Y$ with
\[
 \sigma=Z\,\Sigma_g(F_{c'})+p^{c'}Y',\qquad Z=X^\ell\,\Sigma_g\Big(\prod_{t\in T}\Delta^{\mathrm{top}}_t(g)^{kq}\Big)Z_1H,
 \qquad Y'=X^\ell Z_1Y\in\mathcal R[X].
\]

\emph{Step 4: realization.} The label $X^\ell$ has $\ell(s,c)=0$ for $s\notin T$, so it lies on $T$, and by
Lemma~\ref{pd-symbols-realize-labels}, $X^\ell\Sigma_g(\prod_{t\in T}\Delta^{\mathrm{top}}_t(g)^{kq})$ lies in
$\mathrm{Sy}_g(o')$, where $o'$ is the normal degree of $X^\ell$. The polynomial $H$ lies in $\mathcal R[X]$, lies on
$T$ and involves no normal variable, so $H\in\mathrm{Sy}_g(0)$, as $H=1\cdot H$ with the empty key product. By
Lemma~\ref{pd-symbols-realize-closure}(1), $Z\in\mathcal S_g$. Every monomial of $\sigma$ has normal degree $o$, so
Lemma~\ref{pd-symbols-realize-core}, applied to $Z$ and $Y'$, shows that $\sigma$ is realizable at $(g,c',o)$.
% Lean: one_sided_lemma4_realisable h6 with h6 = one_sided_lemma4_of_admitted_theorem_II (OneSidedLemma4 D C6U
%   (C6U + 1) C6U, from the clauses h15 to h21 of admitted through one_sided_lemma4_of_admitted_lemmas), under
%   HypII, which hypII_at C6U (C6U + 1) C6U supplies from Hyps C for C >= C4 (with cPrime = c').
%   OneSidedLemma4 is one_sided_lemma4_of_interfaces D C (C + 1) C restModel (restDivMU hE restGeomU)
%   (generator_pair_conditions_of_admitted D hA3 hB1) (generator_path_homogeneous_holds D hB2 hB3 hC1) (restMon D)
%   restPh: Step 1 is restModel (the parts of given row degrees, hwh); Step 2 is good_of_span with restDivMU (Lemma
%   pd-symbols-one-sided-real-division), GeneratorSpanInclusion and GeneratorImageIsSymbol (Steps 5 and 6 of the
%   proof of Lemma pd-symbols-realize-segments, through SyS_of_span) and Sy_IT (IT_C a) for scalars; Step 3 is hT
%   with restPh (Lemma pd-symbols-one-sided-real-fixed); Step 4 is core (Lemma pd-symbols-realize-core) with restMon
%   (Lemma pd-symbols-realize-labels), SyS_of_Sy (Sy_IT hH1) and inWPow_zero_mul.
\end{proof}

\begin{lemma}\label{pd-symbols-flat-real-multinomial}
Let $p$ be a prime, $\mathcal K$ a field of characteristic $0$ and $\mathcal R\subset\mathcal K$ a local subring
whose maximal ideal contains $p$. Let $n',k\geq1$ be integers, $d\colon\mathbf Z/n'\mathbf Z\to\mathbf Z_{\geq0}$,
and $\kappa_F=\sum_sk\lfloor\log_pd(s)\rfloor$, the sum over the $s$ with $d(s)\geq1$. Let $\mathcal K[X]$ be the
polynomial ring in variables $X_{s,c}$, $s\in\mathbf Z/n'\mathbf Z$, $1\leq c\leq k$, with the Fock pairing
$\langle\ ,\ \rangle_{\mathrm F}$ of Definition~\ref{pd-symbols-test-fock}, and for an exponent vector $\alpha$ put
$\alpha!=\prod_{s,c}\alpha(s,c)!$. Then:
\begin{enumerate}
\item if $\sum_c\alpha(s,c)=d(s)$ for every $s$, then $\prod_sd(s)!$ divides $p^{\kappa_F}\alpha!$ in $\mathcal R$;
\item if $G\in\mathcal R[X]$ has multidegree $d$ and $E\in\mathcal R[X]$, then
  $p^{\kappa_F}\langle G,E\rangle_{\mathrm F}/\prod_sd(s)!$ lies in $\mathcal R$.
\end{enumerate}
\end{lemma}
\leanref{Formalism/SharpRestriction/GrowthBounds/DualFockLinear/OneSidedLemma4Realisable.lean}{Formalism.MainNotions.TestSideCounting.MultDiv}
\leanref{Formalism/SharpRestriction/GrowthBounds/DualFockLinear/OneSidedLemma4Realisable.lean}{Formalism.MainNotions.TestSideCounting.kapFl}
\leanref{Formalism/SharpRestriction/GrowthBounds/DualFockLinear/FockRLinDual.lean}{Formalism.MainNotions.TestSideCounting.multDiv_kapFl}
\leanref{Formalism/PolynomialDimension/MultinomialDepth/FockGaussComparisonHolds.lean}{Formalism.CountingSumsFrames.fock_mult_val_le}
\leanref{Formalism/SharpRestriction/GrowthBounds/DualFockLinear/FockRAddRight.lean}{Formalism.MainNotions.TestSideCounting.int_pair}

\begin{proof}
(1) Write $v_p$ for the $p$-adic valuation of a positive integer. For integers $0\leq j\leq n$ with $n\geq1$,
Legendre's formula gives $v_p\binom nj=\sum_{i\geq1}\big(\lfloor n/p^i\rfloor-\lfloor j/p^i\rfloor-
\lfloor(n-j)/p^i\rfloor\big)$; each term is $0$ or $1$, and it is $0$ when $p^i>n$, so
$v_p\binom nj\leq\lfloor\log_pn\rfloor$. Fix $s$ and put $M_s=d(s)!/\prod_c\alpha(s,c)!$, the multinomial
coefficient. Writing it as $\binom{d(s)}{\alpha(s,1)}$ times the multinomial coefficient of
$\alpha(s,2),\dots,\alpha(s,k)$, and using that $\lfloor\log_p\rfloor$ is nondecreasing, induction on the number of
parts gives $v_p(M_s)\leq k\lfloor\log_pd(s)\rfloor$ if $d(s)\geq1$; if $d(s)=0$, then $M_s=1$. So
$M=\prod_sM_s$ satisfies $\prod_sd(s)!=M\alpha!$ and $M=p^vu$ with $v\leq\kappa_F$ and $p\nmid u$. Then
$au+bp=1$ for some integers $a,b$; as $p$ lies in the maximal ideal of the local ring $\mathcal R$, $au$ is a
unit, and so is $u$. Hence $p^{\kappa_F}\alpha!=\prod_sd(s)!\cdot p^{\kappa_F-v}u^{-1}$.

(2) The integer $\prod_sd(s)!$ is not $0$ in $\mathcal K$. By definition
$\langle G,E\rangle_{\mathrm F}=\sum_\alpha G_\alpha E_\alpha\,\alpha!$, the sum over the exponents $\alpha$ of the
monomials of $G$, which satisfy $\sum_c\alpha(s,c)=d(s)$ for every $s$. By (1),
$p^{\kappa_F}\alpha!=\prod_sd(s)!\cdot w_\alpha$ with $w_\alpha\in\mathcal R$, so
$p^{\kappa_F}\langle G,E\rangle_{\mathrm F}/\prod_sd(s)!=\sum_\alpha G_\alpha E_\alpha w_\alpha\in\mathcal R$.
% Lean: (1) is multDiv_kapFl (MultDiv O p (kapFl p k d) d): h1 to h3 (d s ! = multinomial * prod (alpha (s, c))!),
%   h4 from fock_val_prod and fock_mult_val_le = fock_mult_fact_le (induction with Nat.multinomial_insert and
%   Nat.factorization_choose_le_log), then huO (the unit, from Nat.gcd_eq_gcd_ab and the maximal ideal).
%   (2) is int_pair (fockR_eq_sum over G.support).
\end{proof}

\begin{lemma}\label{pd-symbols-flat-real-forms}
Let $p\geq5$ be a prime, $(W,Q,\sigma,\iota)$ a Hida setting of degree $n'$ at $p$ with rings
$\mathcal R\subset\mathcal K$ and $\varpi_0$, $1\leq r<k$ integers, $d\colon\mathbf Z/n'\mathbf Z\to\mathbf Z_{\geq0}$,
and $V^\flat$, $H^\flat$, $R^\flat$, $R'^\flat$ and $\mathcal K[X]_d$ as in Definition~\ref{pd-symbols-flat-model};
let $A$ be the integer of Definition~\ref{pd-projection-property}(1) for $\lambda^\flat$, and $\Theta_1$ the flat
model of $V^\flat$ at $1$ of Lemma~\ref{pd-symbols-flat-exists}. For a $W$-lattice $L\subset Q^k$ let $L^\vee$ be
its dual, and for an edge $e'$ with source $L''$ and target $L$ let $e'^\vee$ be the edge with source $L^\vee$ and
target $L''^\vee$. Put $\kappa_F=\sum_sk\lfloor\log_pd(s)\rfloor$, the sum over the $s$ with $d(s)\geq1$, and for
$y\in\mathcal K[X]_d$ let $U_y\in(V^\flat)^\vee$ be the form
\[
 U_y(u)=\frac{p^{\kappa_F}}{\prod_sd(s)!}\,\big\langle y,\Theta_1^{-1}(u)\big\rangle_{\mathrm F}\qquad(u\in V^\flat).
\]
For a family $N=(N_s)$ of $k\times k$ matrices over $\mathcal K$ and $F\in\mathcal K[X]$ let $F^N$ be $F$ with the
column $X_s={}^t(X_{s,1},\dots,X_{s,k})$ replaced by $N_sX_s$ for every $s$. Let $i\in\mathbf Z$ and
$y\in\mathcal K[X]_d$. Then:
\begin{enumerate}
\item $\langle F^N,G\rangle_{\mathrm F}=\langle F,G^{{}^tN}\rangle_{\mathrm F}$ for all $F,G\in\mathcal K[X]$, where
  ${}^tN=({}^tN_s)$;
\item if $L\subset Q^k$ is a $W$-lattice and $\Theta_1(y)\in\varpi_0^iH^\flat(L^\vee)$, then
  $U_y\in\varpi_0^iH^\flat(L)^\vee$;
\item if $e$ is an edge and $\Theta_1(y)\in\varpi_0^iR^\flat(e^\vee)$, then $U_y\in\varpi_0^iR'^\flat(e)$.
\end{enumerate}
\end{lemma}
\leanref{Formalism/SharpRestriction/GrowthBounds/DualFockLinear/FockRAddRight.lean}{Formalism.MainNotions.TestSideCounting.exists_thK}
\leanref{Formalism/SharpRestriction/GrowthBounds/DualFockLinear/FockRAddRight.lean}{Formalism.MainNotions.TestSideCounting.TY}
\leanref{Formalism/SharpRestriction/GrowthBounds/DualFockLinear/FockRLinDual.lean}{Formalism.MainNotions.TestSideCounting.fockR_Lin}
\leanref{Formalism/SharpRestriction/GrowthBounds/DualFockLinear/FockRLinDual.lean}{Formalism.MainNotions.TestSideCounting.fockR_Lin_dual}
\leanref{Formalism/SharpRestriction/GrowthBounds/DualFockLinear/FockRAddRight.lean}{Formalism.MainNotions.TestSideCounting.sandwich}
\leanref{Formalism/SharpRestriction/GrowthBounds/DualFockLinear/LaTwInvT.lean}{Formalism.MainNotions.TestSideCounting.FD_HL}
\leanref{Formalism/SharpRestriction/GrowthBounds/DualFockLinear/LaTwInvT.lean}{Formalism.MainNotions.TestSideCounting.FD_edge}

\begin{proof}
As in the proof of Lemma~\ref{pd-symbols-test-fock-position}, $\mathcal K$ has characteristic $0$ and $p$ lies in
the maximal ideal of the local ring $\mathcal R$, so Lemma~\ref{pd-symbols-flat-real-multinomial} applies. By
Lemma~\ref{pd-symbols-flat-exists} with $\gamma=1$, $\Theta_1$ is a $\mathcal K$-linear bijection of
$\mathcal K[X]_d$ onto $V^\flat$, so $U_y$ is defined. We use the notation of the proof of that lemma: the maps
$\Theta_h$ for families $h=(h_s)$ in $GL_k(\mathcal K)$, with $\Theta_h(F)=\Theta_1(F^{a(h)})$, where
$a_s(h)=h_s^{-1}$ for $s$ even and $a_s(h)={}^th_s$ for $s$ odd, and $(F^N)^M=F^{NM}$.

(1) The proof of Lemma~\ref{pd-kappa-gram-fock}(1) applies verbatim, with the columns $X_s$ in place of the columns
of $Z$ and the matrix $N_s$ depending on the column: $\langle FX_v,G\rangle_{\mathrm F}=\langle F,\partial_vG\rangle_{\mathrm F}$
for every variable $X_v$, and the chain rule gives
$\partial_{(s,b)}(G^{{}^tN})=\sum_a(N_s)_{ba}(\partial_{(s,a)}G)^{{}^tN}$.

\emph{A pairing of lattices.} Let $h=(h_s)$ be a family in $GL_k(\mathcal K)$, $h'=({}^th_s^{-1})$, and
$F,G\in\mathcal R[X]_d=\mathcal R[X]\cap\mathcal K[X]_d$. If $\Theta_1(y)=\Theta_{h'}(G)$, then
$U_y(\Theta_h(F))\in\mathcal R$. Indeed $\Theta_1(y)=\Theta_1(G^{a(h')})$ and
$\Theta_h(F)=\Theta_1(F^{a(h)})$, so $y=G^{a(h')}$ and $\Theta_1^{-1}(\Theta_h(F))=F^{a(h)}$, as $\Theta_1$ is
injective on $\mathcal K[X]_d$. For every $s$, $a_s(h')={}^ta_s(h)^{-1}$: for $s$ even both are ${}^th_s$, for $s$
odd both are $h_s^{-1}$. By (1),
\[
 \big\langle G^{a(h')},F^{a(h)}\big\rangle_{\mathrm F}=\big\langle G,F^{a(h)\,{}^ta(h')}\big\rangle_{\mathrm F}
 =\langle G,F\rangle_{\mathrm F},
\]
so $U_y(\Theta_h(F))=p^{\kappa_F}\langle G,F\rangle_{\mathrm F}/\prod_sd(s)!$, which lies in $\mathcal R$ by
Lemma~\ref{pd-symbols-flat-real-multinomial}(2), as $G$ has multidegree $d$.

(2) Let the columns of $\ell\in GL_k(Q)$ form a basis of $L$, and $\ell_s=\iota\sigma^s(\ell)$. The columns of
${}^t\ell^{-1}$ form a basis of $L^\vee$, and $\iota\sigma^s({}^t\ell^{-1})={}^t\ell_s^{-1}$. By Step~4 of the proof
of Lemma~\ref{pd-symbols-flat-exists}, applied to $\ell$ and to ${}^t\ell^{-1}$, $H^\flat(L)=\Theta_h(\mathcal R[X]_d)$
for $h=(\ell_s)$ and $H^\flat(L^\vee)=\Theta_{h'}(\mathcal R[X]_d)$. Write $\Theta_1(y)=\varpi_0^i\Theta_{h'}(G)$ with
$G\in\mathcal R[X]_d$, and $y=\varpi_0^iy'$; then $\Theta_1(y')=\Theta_{h'}(G)$ and $U_y=\varpi_0^iU_{y'}$. By the
pairing of lattices $U_{y'}$ takes $\Theta_h(\mathcal R[X]_d)=H^\flat(L)$ into $\mathcal R$, so
$U_y\in\varpi_0^iH^\flat(L)^\vee$.

(3) Let $e$ have source $L''$ and target $L$. As in the identification in the proof of
Lemma~\ref{pd-symbols-test-fock-dual}, there is a matrix $b$ whose columns form a basis of $L''$, such that with
$b_s=\iota\sigma^s(b)$ and $D$ the diagonal matrix with entries $\varpi_0^{-1}$ in the first $r$ places and $1$ in
the others, the lattices from which $H^\flat(e)$ is formed are spanned by the columns of the $b_sD$, and those from
which $H^\flat(e^\vee)$ is formed by the columns of the ${}^t(b_sD)^{-1}$. As in Step~5 of the proof of
Lemma~\ref{pd-symbols-flat-exists}, $H^\flat(e)=\Theta_h(\mathcal R[X]_d)$ for $h=(b_sD)$ and
$H^\flat(e^\vee)=\Theta_{h'}(\mathcal R[X]_d)$. As $R^\flat(e^\vee)=\varpi_0^{-A}H^\flat(e^\vee)$ and
$R'^\flat(e)=\varpi_0^{-A}H^\flat(e)^\vee$, the argument of (2), with $i-A$ in place of $i$, gives (3).
% Lean: U_y is TY thK cc y with thK from exists_thK (Th0, the flat model at 1 over Kf) and cc = p^kappa / prod d!.
%   (1) is fockR_Lin (Lin_X, Lin_chain); the pairing of lattices is sandwich (mem_tens_nabla, Th0_inj, TY_apply,
%   int_pair, fockR_Lin_dual with Aa_dual). (2) is FD_HL (HL_iff, la_tw_invT) and (3) is FD_edge (HLe_iff,
%   dual_edgeBasis, dualM).
\end{proof}

\begin{lemma}\label{pd-symbols-test-flat-real}
Let $C_{La},C_k,C_{rk},C_{\deg}$ be real numbers with the property of Theorem~\ref{pd-symbols-count}, and let the
data be as in Lemma~\ref{pd-symbols-test-flat} with $C\geq1$. Put $\kappa_F=\sum_sk\lfloor\log_pd_T(s)\rfloor$, the
sum over the $s$ with $d_T(s)\geq1$. Then there is $g_v\in GL_k(W)$ with the following property. Let $o\geq0$ and
$\sigma\in\mathcal R[X]$ of multidegree $d_T$ be realizable at $(g_v,c',o)$. Then there is
$U\in\varpi_0^{-o}H^\flat(L)^\vee\cap R'^\flat(e)$, with $U\in\varpi_0^{c'}R'^\flat(e')$ for every edge $e'\neq e$
with target $L$, such that for every $F\in\mathcal K[X]_{d_T}$ with $\Theta_\gamma(F)\in R^\flat(e)$
\[
 \langle\Theta_\gamma(F),U\rangle\in\frac{p^{\kappa_F}\langle\Sigma_g(F),\sigma\rangle_{\mathrm F}}{\prod_sd_T(s)!}
 +\varpi_0^{c'}\mathcal R .
\]
\end{lemma}
\leanref{Formalism/SharpRestriction/GrowthBounds/DualFockLinear/OneSidedLemma4Realisable.lean}{Formalism.MainNotions.TestSideCounting.RealAt}
\leanref{Formalism/SharpRestriction/GrowthBounds/DualFockLinear/ScaleZeroIff.lean}{Formalism.MainNotions.TestSideCounting.realAt_gen}
\leanref{Formalism/SharpRestriction/GrowthBounds/DualFockLinear/ScaleZeroIff.lean}{Formalism.MainNotions.TestSideCounting.restReal}

\begin{proof}
By (3) of the situation, $r\geq2$ and $k\geq N_1+3r$, so $1\leq r<k$. Let $A=\sum_{s\text{ even}}d_T(s)$ be the
integer of Definition~\ref{pd-projection-property}(1) for $\lambda^\flat$ (Step~2 of the proof of
Lemma~\ref{pd-symbols-flat-exists}), so that the edge pairing of $\lambda^\flat$ is $\langle u,U\rangle=U(p^Au)$.
Let $\Theta_1$ be the flat model of $V^\flat$ at $1$ of Lemma~\ref{pd-symbols-flat-exists}, and $U_y$, for
$y\in\mathcal K[X]_{d_T}$, the forms of Lemma~\ref{pd-symbols-flat-real-forms} for $d=d_T$, whose $\kappa_F$ is the
one of the statement; put $c_F=p^{\kappa_F}/\prod_sd_T(s)!$. We use the notation $F^N$ of that lemma, the matrices
$\delta$ and $\delta'$ of Definition~\ref{pd-symbols-twist}(1), and the diagonal matrix $\delta_p$ with entries
$p^{-1}$ in the first $r$ places and $1$ in the others.

\emph{Step 1: the element $g_v$.} Let $b=\gamma\,{}^tg^{-1}$. By Lemma~\ref{pd-symbols-edge-basis}, the columns of
$b$ form a basis of the source $L''$ of $e=e_\gamma(g)$, and those of $b\delta_p$ a basis of its target $L=W^k$. So
$g_v=b\delta_p$ lies in $GL_k(W)$. The columns of ${}^tg_v^{-1}={}^tb^{-1}\delta_p^{-1}$ form a basis of
$L^\vee=W^k$, and those of ${}^tg_v^{-1}\delta_p={}^tb^{-1}$ a basis of $L''^\vee$. So $e^\vee=e_1(g_v)$, in the
notation of Lemma~\ref{pd-symbols-edge-basis} with $\gamma=1$.

\emph{Step 2: the other edges.} Let $e'$ be an edge with target $L$. Its dual $e'^\vee$ has source $L^\vee=W^k$, the
lattice spanned by the columns of $1$, so $e'^\vee=e_1(g')$ for some $g'\in GL_k(W)$, by
Lemma~\ref{pd-symbols-edge-basis}(1). If $e'\neq e$, then $e'^\vee\neq e^\vee$, as $(e'^\vee)^\vee=e'$; so the plane
of $g'$ differs from that of $g_v$, by Lemma~\ref{pd-symbols-edge-basis}(3) and Step~1.

\emph{Step 3: the form.} Let $o$ and $\sigma$ be as in the statement, and let $F_v\in\mathcal R[X]$ of multidegree
$d_T$ realize $\sigma$ at $(g_v,c',o)$: $\Sigma_{g_v}(F_v)=\varpi_0^o\sigma+\varpi_0^{o+c'}E$ with
$E\in\mathcal R[X]$, and $\Sigma_{g'}(F_v)\in\varpi_0^{o+c'}\mathcal R[X]$ for every $g'\in GL_k(W)$ whose plane differs
from that of $g_v$. Put $U=\varpi_0^{-o}U_{F_v}$.

As $F_v\in\mathcal R[X]$, property~(1) of the flat model $\Theta_1$ gives $\Theta_1(F_v)\in H^\flat(W^k)$; as
$W^k=L$ is its own dual, Lemma~\ref{pd-symbols-flat-real-forms}(2) with $i=0$ gives $U_{F_v}\in H^\flat(L)^\vee$, so
$U\in\varpi_0^{-o}H^\flat(L)^\vee$. As $\sigma\in\mathcal R[X]$, $\Sigma_{g_v}(F_v)\in\varpi_0^o\mathcal R[X]$, and
property~(2) of $\Theta_1$ gives $\Theta_1(F_v)\in\varpi_0^oR^\flat(e_1(g_v))=\varpi_0^oR^\flat(e^\vee)$; so
$U_{F_v}\in\varpi_0^oR'^\flat(e)$ by Lemma~\ref{pd-symbols-flat-real-forms}(3), and $U\in R'^\flat(e)$. Let $e'\neq e$
be an edge with target $L$, and $g'$ as in Step~2. Then $\Sigma_{g'}(F_v)\in\varpi_0^{o+c'}\mathcal R[X]$, so
$\Theta_1(F_v)\in\varpi_0^{o+c'}R^\flat(e'^\vee)$ by property~(2) of $\Theta_1$, and
$U\in\varpi_0^{c'}R'^\flat(e')$ by Lemma~\ref{pd-symbols-flat-real-forms}(3).

\emph{Step 4: the pairing.} Let $F\in\mathcal K[X]_{d_T}$. We show that
$\langle\Theta_\gamma(F),U_{F_v}\rangle=c_F\langle\Sigma_g(F),\Sigma_{g_v}(F_v)\rangle_{\mathrm F}$. By the proof of
Lemma~\ref{pd-symbols-flat-exists}, $\Theta_\gamma(F)=\Theta_1(F^{a(\gamma)})$, where $a_s(\gamma)=\gamma_s^{-1}$ for
$s$ even and $a_s(\gamma)={}^t\gamma_s$ for $s$ odd; so
$\langle\Theta_\gamma(F),U_{F_v}\rangle=p^Ac_F\langle F_v,F^{a(\gamma)}\rangle_{\mathrm F}$. By Step~6 of that proof,
$\Sigma_g(F)=F^N$ and $\Sigma_{g_v}(F_v)=F_v^{N'}$, where $N_s={}^tg_s^{-1}\delta$ and $N'_s={}^tg_{v,s}^{-1}\delta$ for
$s$ even and $N_s=g_s\delta'$ and $N'_s=g_{v,s}\delta'$ for $s$ odd, with $g_{v,s}=\iota\sigma^s(g_v)$. Put
$b_s=\iota\sigma^s(b)=\gamma_s\,{}^tg_s^{-1}$, so that $g_{v,s}=b_s\delta_p$, and $M_s=a_s(\gamma)^{-1}N_s$, so that
$F^N=(F^{a(\gamma)})^M$; then $M_s=b_s\delta$ for $s$ even and $M_s={}^tb_s^{-1}\delta'$ for $s$ odd. As $p=\varpi_0^2$,
$\delta_p^{-1}\delta=p\delta^{-1}$ and $\delta_p\delta'=\delta'^{-1}$; so $N'_s={}^tb_s^{-1}\delta_p^{-1}\delta=
p\,{}^tM_s^{-1}$ for $s$ even and $N'_s=b_s\delta'^{-1}={}^tM_s^{-1}$ for $s$ odd. Since $F_v$ has degree $d_T(s)$ in
$X_s$, $F_v^{N'}=p^AF_v^{M^*}$ with $M^*_s={}^tM_s^{-1}$. By Lemma~\ref{pd-symbols-flat-real-forms}(1),
\[
 \langle\Sigma_g(F),\Sigma_{g_v}(F_v)\rangle_{\mathrm F}=p^A\big\langle(F^{a(\gamma)})^M,F_v^{M^*}\big\rangle_{\mathrm F}
 =p^A\big\langle F^{a(\gamma)},F_v^{M^*\,{}^tM}\big\rangle_{\mathrm F}=p^A\langle F^{a(\gamma)},F_v\rangle_{\mathrm F},
\]
and the Fock pairing is symmetric.

\emph{Step 5: the congruence.} Let $F\in\mathcal K[X]_{d_T}$ with $\Theta_\gamma(F)\in R^\flat(e)$. As
$e=e_\gamma(g)$, property~(2) of the flat model $\Theta_\gamma$ with $i=0$ gives $\Sigma_g(F)\in\mathcal R[X]$, of
multidegree $d_T$. By Step~4,
\[
 \langle\Theta_\gamma(F),U\rangle=\varpi_0^{-o}c_F\langle\Sigma_g(F),\varpi_0^o\sigma+\varpi_0^{o+c'}E\rangle_{\mathrm F}
 =c_F\langle\Sigma_g(F),\sigma\rangle_{\mathrm F}+\varpi_0^{c'}c_F\langle\Sigma_g(F),E\rangle_{\mathrm F},
\]
and $c_F\langle\Sigma_g(F),E\rangle_{\mathrm F}\in\mathcal R$ by Lemma~\ref{pd-symbols-flat-real-multinomial}(2), which
applies as in the proof of Lemma~\ref{pd-symbols-flat-real-forms}.
% Lean: restReal = realStep = realAt_gen with kappa = kapFl p (m + 1) dT and multDiv_kapFl (dvrOfStd p_mem).
%   Step 1 is exists_gv (b * lb_D with colSpan = colSpan 1), hrel (invT_mul_D) and hbv (dual_edgeBasis); Step 2
%   is edgePlane hE1 and hE3 with dual_edge and edgeBasis; Thetav is theta_exists at 1 (hv1, hv2), and Theta = Theta'
%   at hm by IsTrueTheta. Step 3: U = TY thK cc (w^-o Thetav F') for the realization (Ff, hA, hsym, hoth); the
%   three lattice clauses are FD_HL (with h1T), FD_edge with scale_shift and h1, and FD_edge for e' with hoth and
%   hpl. Step 4 is pair_formula (subst_Lin, fockR_Lin_dual with dualM and Aa_dual2, flat_Ar, p_eq). Step 5 is
%   inWj_iff, pol_of_map and int_pair (e1 to e4).
\end{proof}

\begin{definition}\label{pd-symbols-test-flat-chains-tied}
Let $p$, $n'$, $r$, $K$, $k=r+K$, $T$, $(J,f)$, $N_1$, $G'$, $c'$, $\deg$ and $Y$ be as in
Definition~\ref{pd-symbols-test-one-sided}, with $m_*$ and $B_j(x)$ as there, and let $\mathcal R$ be a discrete valuation
ring with fraction field $\mathcal K$ in which $p$ is not a unit; let $\mathcal K[X]$ be the polynomial ring in the
$X_{s,c}$, $s\in\mathbf Z/n'\mathbf Z$, $1\leq c\leq k$. Let $\mathcal D$ be a signed chain datum of rank $r$ with $J$
levels over $\mathcal R$ (Definition~\ref{pd-symbols-test-chains}), and for a level $j\in\mathbf Z/J\mathbf Z$ let
$\hat\jmath\in\{0,\dots,J-1\}$ be its representative.
\begin{enumerate}
\item \emph{Places.} The \emph{index} of the letter $(j,l)$ is $l$, and that of the link $(j,m)$ is $n_j+r-1-m$. The
  \emph{place} $x_j(u)\in\mathbf Z$ of the position of level $j$ with index $u$ is the $(u+1)$-st smallest element of
  $\Lambda'_{\hat\jmath}$ if $u<n_j$ ($0$ if there is none), and $f(\hat\jmath,u-n_j)$ if $n_j\leq u<n_j+r$. For a position
  $y$ of level $j$ with place $x$ put $\lambda(y)=p^{\operatorname{dist}_{\hat\jmath}(x)}$. For $0\leq u\leq n_j+r$ let
  $N_j(u)$ be the set of the positions of level $j$ with index less than $u$, and for $0\leq u<n_j+r$ put
  $s_j(u)=\lfloor B_{\hat\jmath}(x_j(u))/(Yk^2)\rfloor$.
\item \emph{Ties.} Let $d^{\mathrm{pl}}$ be a function from the positions to $\mathbf Z_{\geq0}$, and
  $\kappa(j,u)\in\mathbf Z^K$ for $j\in\mathbf Z/J\mathbf Z$ and $u\geq0$. For $0\leq u<n_j+r$ put
  $\theta^*_j(u)=\max(B/2,t_1,t_2)$, where $B=B_{\hat\jmath}(x_j(u))$, $t_1=\kappa(j,u+1)_1+m_*+s_j(u)$ if $x_j(u)$ is even
  and $t_1=B/2$ otherwise, and $t_2=\kappa(j,u-1)_1+m_*+s_j(u)$ if $u\geq1$ and $x_j(u-1)$ is odd and $t_2=B/2$
  otherwise. For real numbers $a_1,a_2$, the datum $\mathcal D$ is \emph{tied} to $(J,f)$ with $d^{\mathrm{pl}}$, $\kappa$ and
  the margins $(a_1,a_2)$ if $d=d^{\mathrm{pl}}$ and for every level $j$: $n_j=|\Lambda'_{\hat\jmath}|$; the letter $(j,l)$
  is a vector letter if and only if its place is even; $g(j,m)=f(\hat\jmath+1,r-1)-f(\hat\jmath,r-1)$ for every $m<r$;
  $n^c_j=n_j+1$, and for $0\leq i\leq n_j$ the chain $C_{j,i+1}$ consists of the indices $m<r$ and the letters $(j,l)$ with
  $l<i$, with $\theta_{j,i+1}=\theta^*_j(i)-a_1s_j(i)$; and $\theta^L_{j,m}=\theta^*_j(u)-a_2s_j(u)$ with $u=n_j+r-1-m$, for
  $m+2\leq r$. It is \emph{tied} to $(J,f)$ with $d^{\mathrm{pl}}$ and $\kappa$ if this holds with the margins
  $(\frac18,\frac38)$.
\item \emph{Level images.} Let $\mathcal K[x]$ be the polynomial ring of Definition~\ref{pd-symbols-test-chains} for
  $\mathcal D$, and $\mathcal K[w]$ the polynomial ring in variables $w_{y,b}$, $y$ a position and $1\leq b\leq K$. For an
  exponent vector $\mu$ of $\mathcal K[X]$, a position $y$ with place $x$ and $c\leq k$ write $\mu(y,c)=\mu(\bar x,c)$.
  For $\rho$ from the pairs of a position and an integer $1\leq c\leq k$ to $\mathbf Z_{\geq0}$, and $d^{\mathrm{pl}}$ and
  $d^{\mathrm{no}}$ from the positions to $\mathbf Z_{\geq0}$, the \emph{level image} of $G\in\mathcal K[X]$ is
  \[
   L_{\rho,d^{\mathrm{pl}},d^{\mathrm{no}}}(G)=\sum_\mu G_\mu\prod_{y}\prod_{b=1}^Kw_{y,b}^{\lfloor\mu(y,r+b)/\lambda(y)\rfloor}
   \cdot\prod_y\prod_{a=0}^{r-1}x_{y,a}^{\lfloor\mu(y,a+1)/\lambda(y)\rfloor}\in\mathcal K[w][x],
  \]
  the sum over the $\mu$ such that for every position $y$: $\mu(y,c)\equiv\rho(y,c)\pmod{\lambda(y)}$ for every $c\leq k$,
  $\sum_{c\leq r}\lfloor\mu(y,c)/\lambda(y)\rfloor=d^{\mathrm{pl}}(y)$ and $\sum_{c>r}\lfloor\mu(y,c)/\lambda(y)\rfloor=
  d^{\mathrm{no}}(y)$.
\item \emph{Joint projectors.} For a level $j$ and $0\leq u\leq n_j+r$ let $E^{j,u}_{ab}$, $1\leq a,b\leq K$, be the
  endomorphism $\sum_yw_{y,a}\,\partial/\partial w_{y,b}-\sum_{y'}w_{y',b}\,\partial/\partial w_{y',a}$ of $\mathcal K[w]$,
  the sums over the positions $y\in N_j(u)$ that are links or vector letters and the covector letters $y'\in N_j(u)$.
  For a family $\kappa'$ of $\kappa'(j,u)\in\mathbf Z^K$, $0\leq u\leq n_j+r$, let $V(\kappa')$ be the intersection over
  $(j,u)$ of the isotypic subspaces of weight $\kappa'(j,u)$ of the families $(E^{j,u}_{ab})_{a,b}$
  (Definition~\ref{pd-kappa-defects-segment-mean-operators}(1)). A \emph{joint isotypic projector} of weight $\kappa$ is
  a $\mathcal K$-linear endomorphism $e$ of $\mathcal K[w]$ that is the identity on $V(\kappa)$ and $0$ on $V(\kappa')$
  for every such family $\kappa'\neq\kappa$.
\item \emph{Membership.} For such $e$, a monomial $w^\nu$ and $L=\sum_\alpha L_\alpha x^\alpha\in\mathcal K[w][x]$, the
  \emph{plane coefficient} $[eL]_\nu\in\mathcal K[x]$ is $\sum_\alpha c_\alpha x^\alpha$, with $c_\alpha$ the coefficient of
  $w^\nu$ in $e(L_\alpha)$. A polynomial $G\in\mathcal K[X]$ satisfies the \emph{membership condition with the margins
  $(a_1,a_2)$} if for every signed chain datum $\mathcal D$ of rank $r$ with $J$ levels over $\mathcal R$, all
  $d^{\mathrm{pl}}$, $d^{\mathrm{no}}$ and $\kappa$ such that $\mathcal D$ is tied to $(J,f)$ with $d^{\mathrm{pl}}$,
  $\kappa$ and the margins $(a_1,a_2)$, every $\rho$, every joint isotypic projector $e$ of weight $\kappa$ and every
  monomial $w^\nu$, the plane coefficient $\big[e\,L_{\rho,d^{\mathrm{pl}},d^{\mathrm{no}}}(G)\big]_\nu$ is admissible
  for $\mathcal D$. It satisfies the \emph{membership condition} if it does with the margins $(\frac18,\frac38)$.
\end{enumerate}
\end{definition}
\leanref{Formalism/SharpRestriction/GrowthBounds/TwistedSubstitution/PlaneCoeff.lean}{Formalism.MainNotions.TestSideWords.posTL}
\leanref{Formalism/SharpRestriction/GrowthBounds/TwistedSubstitution/PlaneCoeff.lean}{Formalism.MainNotions.TestSideWords.thStar}
\leanref{Formalism/SharpRestriction/GrowthBounds/TwistedSubstitution/PlaneCoeff.lean}{Formalism.MainNotions.TestSideWords.slU}
\leanref{Formalism/SharpRestriction/GrowthBounds/TwistedSubstitution/PlaneCoeff.lean}{Formalism.MainNotions.TestSideWords.Tied}
\leanref{Formalism/SharpRestriction/GrowthBounds/TwistedSubstitution/PlaneCoeff.lean}{Formalism.MainNotions.TestSideWords.levImg}
\leanref{Formalism/SharpRestriction/GrowthBounds/TwistedSubstitution/PlaneCoeff.lean}{Formalism.MainNotions.TestSideWords.normE}
\leanref{Formalism/SharpRestriction/GrowthBounds/TwistedSubstitution/PlaneCoeff.lean}{Formalism.MainNotions.TestSideWords.jIso}
\leanref{Formalism/SharpRestriction/GrowthBounds/TwistedSubstitution/PlaneCoeff.lean}{Formalism.MainNotions.TestSideWords.IsJProj}
\leanref{Formalism/SharpRestriction/GrowthBounds/TwistedSubstitution/PlaneCoeff.lean}{Formalism.MainNotions.TestSideWords.oneT}
\leanref{Formalism/SharpRestriction/GrowthBounds/TwistedSubstitution/PlaneCoeff.lean}{Formalism.MainNotions.TestSideWords.planeCoeff}
\leanref{Formalism/SharpRestriction/GrowthBounds/TwistedSubstitution/PlaneCoeff.lean}{Formalism.MainNotions.TestSideWords.MemAt}
\leanref{Formalism/SharpRestriction/GrowthBounds/TwistedSubstitution/MemLin.lean}{Formalism.MainNotions.TestSideWords.TwistedSubstitutionGrowthBound}
\leanref{Formalism/SharpRestriction/GrowthBounds/TwistedSubstitution/MemLin.lean}{Formalism.MainNotions.TestSideWords.MemAtG}
\leanref{Formalism/SharpRestriction/GrowthBounds/TwistedSubstitution/MemLin.lean}{Formalism.MainNotions.TestSideWords.memAt_iff}

\begin{definition}\label{pd-symbols-test-flat-chains-membership-loss}
Let the notation be as in Definition~\ref{pd-symbols-test-flat-chains-tied}, with $\mathcal K[X]$, the run system
$(J,f)$ and its sets $\Lambda_j$ and $\Lambda'_j$ (Definition~\ref{pd-symbols-run-system}). For an integer $j$ and
$0\leq u<|\Lambda'_j|+r$ let $x(j,u)$ be the $(u+1)$-st smallest element of $\Lambda'_j$ if $u<|\Lambda'_j|$, and
$f(j,u-|\Lambda'_j|)$ otherwise, and put $s(j,u)=\lfloor B_j(x(j,u))/(Yk^2)\rfloor$; so the place $x_j(u)$ and the
number $s_j(u)$ of Definition~\ref{pd-symbols-test-flat-chains-tied}(1) are $x(\hat\jmath,u)$ and $s(\hat\jmath,u)$. A
\emph{loss} is a function $L\colon\mathbf Z\times\mathbf Z_{\geq0}\to\mathbf R$. Let $a_1,a_2$ be real numbers and $L$ a
loss.
\begin{enumerate}
\item A signed chain datum $\mathcal D$ of rank $r$ with $J$ levels over $\mathcal R$ is \emph{tied} to $(J,f)$ with
  $d^{\mathrm{pl}}$, $\kappa$, the margins $(a_1,a_2)$ and the loss $L$ if it satisfies the conditions of
  Definition~\ref{pd-symbols-test-flat-chains-tied}(2) for the margins $(a_1,a_2)$, except that the two equalities for
  the thresholds are replaced by the inequalities $\theta_{j,i+1}\leq\theta^*_j(i)-a_1s_j(i)-L(\hat\jmath,i)$ for
  $0\leq i\leq n_j$ and $\theta^L_{j,m}\leq\theta^*_j(u)-a_2s_j(u)-L(\hat\jmath,u)$, with $u=n_j+r-1-m$, for $m+2\leq r$.
\item A polynomial $G\in\mathcal K[X]$ satisfies the \emph{membership condition with the margins $(a_1,a_2)$ and the
  loss $L$} if it satisfies the condition of Definition~\ref{pd-symbols-test-flat-chains-tied}(5) with the data
  $\mathcal D$, $d^{\mathrm{pl}}$, $d^{\mathrm{no}}$ and $\kappa$ for which $\mathcal D$ is tied to $(J,f)$ with
  $d^{\mathrm{pl}}$, $\kappa$, the margins $(a_1,a_2)$ and the loss $L$.
\item For $s\in\mathbf Z/n'\mathbf Z$ the \emph{cost} $c_s$ is the loss with $c_s(j,u)=2$ if some $x\in\Lambda_j$ with
  $\bar x=s$ satisfies $x\leq x(j,u)$, and $c_s(j,u)=0$ otherwise (in particular when $u\geq|\Lambda'_j|+r$). The
  \emph{loss} of a product $D_1\cdots D_l$ of operators $D_i=X_{s_i,a_i}\,\partial/\partial X_{s_i,b_i}$
  (Definition~\ref{pd-symbols-test-operators}) is $c_{s_1}+\dots+c_{s_l}$.
\end{enumerate}
\end{definition}
\leanref{Formalism/SharpRestriction/GrowthBounds/TwistedSubstitution/MemAtDZeroIff.lean}{Formalism.MainNotions.TestSideWords.TwistedSubstitutionLossInequality}
\leanref{Formalism/SharpRestriction/GrowthBounds/TwistedSubstitution/MemAtDZeroIff.lean}{Formalism.MainNotions.TestSideWords.MemAtLe}
\leanref{Formalism/SharpRestriction/GrowthBounds/TwistedSubstitution/WordLossNonneg.lean}{Formalism.MainNotions.TestSideWords.touchT}
\leanref{Formalism/SharpRestriction/GrowthBounds/TwistedSubstitution/WordLossNonneg.lean}{Formalism.MainNotions.TestSideWords.opCost}
\leanref{Formalism/SharpRestriction/GrowthBounds/TwistedSubstitution/WordLossNonneg.lean}{Formalism.MainNotions.TestSideWords.wordLoss}

\begin{lemma}\label{pd-symbols-test-flat-chains-membership-one-position}
Let $p$ and $\mathcal R$ be as in Definition~\ref{pd-symbols-test-chains}, $n',r\geq1$ and $L_1,N_1\geq0$ integers,
$T\subset\mathbf Z/n'\mathbf Z$, and $(J,f)$ a run system for $T$ of length $r$ valid for $(L_1,N_1)$. Let $\mathcal D$
be a signed chain datum of rank $r$ with $J$ levels over $\mathcal R$ with $n_j=|\Lambda'_{\hat\jmath}|$ for every level
$j$, with the places, indices and sets $N_j(u)$ of Definition~\ref{pd-symbols-test-flat-chains-tied}(1), and let $c_s$ be
as in Definition~\ref{pd-symbols-test-flat-chains-membership-loss}(3).
\begin{enumerate}
\item For every $s\in T$ there is exactly one position whose place has residue $s$.
\item Let $y_0$ be a position whose place has residue $s$, $j$ a level and $0\leq u<n_j+r$. If $y_0\in N_j(u+1)$, or
  $y_0$ is a link of level $j-1$, then $c_s(\hat\jmath,u)=2$.
\end{enumerate}
\end{lemma}
\leanref{Formalism/SharpRestriction/GrowthBounds/DiagonalSupport/PosTLSurj.lean}{Formalism.MainNotions.Proposition5Numerics.posTL_surj}
\leanref{Formalism/SharpRestriction/GrowthBounds/PlaneCoefficients/GeoTouch.lean}{Formalism.MainNotions.PairMembershipBlocks.geoTop}
\leanref{Formalism/SharpRestriction/GrowthBounds/PlaneCoefficients/GeoTouch.lean}{Formalism.MainNotions.PairMembershipBlocks.geoTouch}

\begin{proof}
(1) Uniqueness is Lemma~\ref{pd-symbols-test-flat-chains-datum-numbers-places}(2). For existence, let $s\in T$ and
let $v$ be the integer with $\bar v=s$ and $f(-1,r-1)<v\leq f(-1,r-1)+n'=f(J-1,r-1)$, by (1) of validity. As
$f(i,r-1)$ increases with $i$, there is $0\leq\hat\jmath<J$ with $f(\hat\jmath-1,r-1)<v\leq f(\hat\jmath,r-1)$. If
$v<f(\hat\jmath,0)$, then $v\in\Lambda'_{\hat\jmath}$ is the place of a letter of level $\hat\jmath$. Otherwise
$f(\hat\jmath,0)\leq v\leq f(\hat\jmath,r-1)$ and $\bar v\in T$; by (3) of validity no integer with residue in $T$ lies
strictly between $f(\hat\jmath,i)$ and $f(\hat\jmath,i+1)$, so $v=f(\hat\jmath,i)$ for some $i<r$, the place of the link
$(\hat\jmath,r-1-i)$.

(2) The places of the positions of level $j$ increase with their indices: the letters come in the order of
$\Lambda'_{\hat\jmath}$, then the links with the places $f(\hat\jmath,0)<\dots<f(\hat\jmath,r-1)$. If $y_0\in N_j(u+1)$,
its index is at most $u$, so its place $x$ satisfies $x\leq x_j(u)$, and $x\in\Lambda_{\hat\jmath}$ by
Lemma~\ref{pd-symbols-test-flat-chains-datum-numbers-places}(1). If $y_0=(j-1,m)$, its place is congruent modulo $n'$,
by (1) of validity, to $x=f(\hat\jmath-1,r-1-m)$; then $\bar x=s$, $f(\hat\jmath-1,0)\leq x\leq f(\hat\jmath-1,r-1)<
f(\hat\jmath,r-1)$, so $x\in\Lambda_{\hat\jmath}$, and $x\leq f(\hat\jmath-1,r-1)<x_j(u)$, as $x_j(u)$ is the place of a
position of level $j$. In both cases $x\in\Lambda_{\hat\jmath}$, $\bar x=s$ and $x\leq x_j(u)=x(\hat\jmath,u)$.
% Lean: (1) is geoTop (posTL_surj: mn_cover, ph_mem_nonW, llPos_surj and ge_frame; posTL_inj). (2) is geoTouch
%   (posU_le_iff for the order of the places, LL_sub_letters and fr_mem_letters; w_mem_letters, DdG.shift and
%   w_lt_posU for a link of level j - 1).
\end{proof}

\begin{lemma}\label{pd-symbols-test-flat-chains-membership-one-transport}
Let the notation be as in Definition~\ref{pd-symbols-test-flat-chains-tied}. Let $y_0$ be a position whose place has a
residue $s$ that the place of no other position has, put $\lambda=\lambda(y_0)$, and let $1\leq c\leq k$, $\rho$,
$d^{\mathrm{pl}}$, $d^{\mathrm{no}}$ and $F\in\mathcal K[X]$ be given. If $c\leq r$, let $\xi=x_{y_0,c-1}$ and
$d^\xi=d^{\mathrm{pl}}$; if $c>r$, let $\xi=w_{y_0,c-r}$, acting on $\mathcal K[w][x]$ through the coefficients, and
$d^\xi=d^{\mathrm{no}}$. Write $L_\rho$ for $L_{\rho,d^{\mathrm{pl}},d^{\mathrm{no}}}$ and $L^\pm_\rho$ for the level image
with $d^\xi$ replaced by $d^\xi\pm\delta$, where $\delta$ is $1$ at $y_0$ and $0$ elsewhere; let $t=\rho(y_0,c)$, and let
$\rho[t']$ be $\rho$ with its value at $(y_0,c)$ replaced by $t'$.
\begin{enumerate}
\item If $t+1<\lambda$, then $L_\rho(\partial F/\partial X_{s,c})=(t+1)L_{\rho[t+1]}(F)+\lambda\,\xi\,\partial
  L_{\rho[t+1]}(F)/\partial\xi$.
\item If $t+1=\lambda$, then $L_\rho(\partial F/\partial X_{s,c})=\lambda\,\partial L^+_{\rho[0]}(F)/\partial\xi$.
\item If $0<t<\lambda$, then $L_\rho(X_{s,c}F)=L_{\rho[t-1]}(F)$.
\item If $t=0$, then $L_\rho(X_{s,c}F)=\xi\,L^-_{\rho[\lambda-1]}(F)$ if $d^\xi(y_0)\geq1$, and $L_\rho(X_{s,c}F)=0$ if
  $d^\xi(y_0)=0$.
\item If $t\geq\lambda$, then $L_\rho(\partial F/\partial X_{s,c})=L_\rho(X_{s,c}F)=0$.
\end{enumerate}
\end{lemma}
\leanref{Formalism/SharpRestriction/GrowthBounds/PlaneCoefficients/NoOpStep.lean}{Formalism.MainNotions.PairMembershipBlocks.TrDerPl}
\leanref{Formalism/SharpRestriction/GrowthBounds/PlaneCoefficients/NoOpStep.lean}{Formalism.MainNotions.PairMembershipBlocks.TrDerNo}
\leanref{Formalism/SharpRestriction/GrowthBounds/PlaneCoefficients/NoOpStep.lean}{Formalism.MainNotions.PairMembershipBlocks.TrMulPl}
\leanref{Formalism/SharpRestriction/GrowthBounds/PlaneCoefficients/NoOpStep.lean}{Formalism.MainNotions.PairMembershipBlocks.TrMulNo}
\leanref{Formalism/SharpRestriction/GrowthBounds/PlaneCoefficients/DoubleMonomialApply.lean}{Formalism.MainNotions.PairMembershipBlocks.trDerPl}
\leanref{Formalism/SharpRestriction/GrowthBounds/PlaneCoefficients/DoubleMonomialApply.lean}{Formalism.MainNotions.PairMembershipBlocks.trDerNo}
\leanref{Formalism/SharpRestriction/GrowthBounds/PlaneCoefficients/DoubleMonomialApply.lean}{Formalism.MainNotions.PairMembershipBlocks.trMulPl}
\leanref{Formalism/SharpRestriction/GrowthBounds/PlaneCoefficients/DoubleMonomialApply.lean}{Formalism.MainNotions.PairMembershipBlocks.trMulNo}

\begin{proof}
All the maps are $\mathcal K$-linear, so let $F=X^\mu$. As no other position has residue $s$, the operators change
only the exponent at $(y_0,c)$ among those that the level images read. Write $\mu(s,c)=\lambda q+t_0$ with
$0\leq t_0<\lambda$. Then $X^\mu$ contributes to $L_{\rho'}$ only if $\rho'(y_0,c)=t_0$, and its contribution has the
factor $\xi^q$. A residue at least $\lambda$ is never met, which gives (5).

(1), (2) $\partial X^\mu/\partial X_{s,c}=\mu(s,c)X^{\mu-e}$, $e$ the unit vector at $(s,c)$. If $t+1<\lambda$, then
$X^{\mu-e}$ has residue $t$ at $(y_0,c)$ exactly when $t_0=t+1$, with the same quotient $q$; its contribution is
$(t+1+\lambda q)$ times that of $X^\mu$ to $L_{\rho[t+1]}$, and $\xi\partial\xi^q/\partial\xi=q\xi^q$. If $t+1=\lambda$,
then $X^{\mu-e}$ has residue $t$ exactly when $t_0=0$ and $q\geq1$, with quotient $q-1$; so $X^\mu$ lies in the block
with $d^\xi+\delta$, and the contribution is $\lambda q\,\xi^{q-1}$ in place of $\xi^q$, that is
$\lambda\,\partial\xi^q/\partial\xi$. For $q=0$ both sides vanish.

(3), (4) $X_{s,c}X^\mu=X^{\mu+e}$. If $0<t<\lambda$, it has residue $t$ exactly when $t_0=t-1$, with the same quotient.
If $t=0$, it has residue $0$ exactly when $t_0=\lambda-1$, with quotient $q+1$; then $X^\mu$ lies in the block with
$d^\xi-\delta$, and the contribution is $\xi^{q+1}=\xi\cdot\xi^q$. If $d^\xi(y_0)=0$, no monomial with a quotient
$q+1\geq1$ at $(y_0,c)$ lies in the block, and the image is $0$.
% Lean: trDerPl and trDerNo are der_A_pl, der_B_pl and levImg_empty (and der_A_no, der_B_no), by the trichotomy of
%   rho + 1 and lowP; trMulPl and trMulNo are mul_A, mul_Bpl (mul_Bno), mul_Bpl_zero (mul_Bno_zero) and levImg_empty.
%   The computations are levImg_monomial with the digit identities nat_A, nat_A_div, nat_B and nat_B_div, and pl_euler,
%   pl_der, pl_mul, no_euler, no_der and no_mul for the effect on the double monomials (DoubleMonomialApply.lean).
\end{proof}

\begin{lemma}\label{pd-symbols-test-flat-chains-membership-one-twist}
Let $L$ be a commutative ring, $\Lambda$ a set each of whose elements is a vector letter or a covector letter, $n\geq1$,
and $L[z]$ the polynomial ring in the $z_{v,c}$, $v\in\Lambda$, $1\leq c\leq n$, with $\operatorname{val}_S$ and
$\lambda_S$ for $S\subset\Lambda$ as in Definition~\ref{pd-symbols-path-generators}(1). Let $S\subset\Lambda$, $v\in\Lambda$,
$c\leq n$ and $H\in L[z]$, and let $T$ be $\partial/\partial z_{v,c}$, the multiplication by $z_{v,c}$, or their composite
$z_{v,c}\,\partial/\partial z_{v,c}$. Then $\operatorname{val}_S(TH)\geq\operatorname{val}_S(H)-1$ and
$\lambda_S(TH)\leq\lambda_S(H)+1$, and both hold without the $1$ if $v\notin S$. Moreover
$\operatorname{val}_S(\partial H/\partial z_{v,c})\geq\operatorname{val}_S(H)$ if $v$ is a covector letter, and
$\operatorname{val}_S(z_{v,c}H)\geq\operatorname{val}_S(H)$ if $v$ is a vector letter.
\end{lemma}
\leanref{Formalism/SharpRestriction/GrowthBounds/TwistedSubstitution/TwistedMonomialFiltration/MaxTwistDegreePderivLeAddOne.lean}{Formalism.MainNotions.TestSideWords.TwistedFiltration.G2.valS_pderiv}
\leanref{Formalism/SharpRestriction/GrowthBounds/TwistedSubstitution/TwistedMonomialFiltration/MaxTwistDegreePderivLeAddOne.lean}{Formalism.MainNotions.TestSideWords.TwistedFiltration.G2.valS_pderiv_of}
\leanref{Formalism/SharpRestriction/GrowthBounds/TwistedSubstitution/TwistedMonomialFiltration/MaxTwistDegreePderivLeAddOne.lean}{Formalism.MainNotions.TestSideWords.TwistedFiltration.G2.valS_X_mul}
\leanref{Formalism/SharpRestriction/GrowthBounds/TwistedSubstitution/TwistedMonomialFiltration/MaxTwistDegreePderivLeAddOne.lean}{Formalism.MainNotions.TestSideWords.TwistedFiltration.G2.valS_X_mul_of}
\leanref{Formalism/SharpRestriction/GrowthBounds/MonomialSupport/Lam1SPderivAndXMulLeAddIndicator.lean}{Formalism.MainNotions.LetterSubstitution.lam1S_pderiv_and_X_mul_le_add_indicator}

\begin{proof}
Let $a$ be $n$ or $1$, and $P_{S,a}$ as in Definition~\ref{pd-symbols-path-generators}(1). If $v\notin S$, the
substitution does not touch the row $z_v$, so it commutes with $T$, and $T$ does not change the exponents of $\eta$.
Let $v\in S$. By the rules of Lemma~\ref{pd-symbols-residue-count-inputs-segment-clauses}(1),
$\operatorname{val}_S(z_{v,c}H)=\operatorname{val}_S(z_{v,c})+\operatorname{val}_S(H)$ and
$\lambda_S(z_{v,c}H)=\lambda_S(z_{v,c})+\lambda_S(H)$, with $\operatorname{val}_S(z_{v,c})\geq0$ and
$\lambda_S(z_{v,c})\leq1$ for a vector letter and $\operatorname{val}_S(z_{v,c})\geq-1$ and $\lambda_S(z_{v,c})\leq0$ for a
covector letter. For the derivative, let $M$ be $\tau_a(\eta)$ for a vector letter and ${}^t\tau_a(\eta^{-1})$ for a
covector letter, and $\nabla$ the column of the $\partial/\partial z_{v,c'}$, $c'\leq n$. The chain rule gives
$\nabla(H_{S,a})=M(\nabla H)_{S,a}$, so $(\nabla H)_{S,a}=M^{-1}\nabla(H_{S,a})$, where $M^{-1}$ is $\tau_a(\eta^{-1})$,
resp.\ ${}^t\tau_a(\eta)$. As $\tau_a(\eta^{\pm1})=1+(\eta^{\pm1}-1)\mathfrak gE_{aa}\mathfrak g^{-1}$, with $E_{aa}$ the
matrix unit, its entries are combinations of $1$ and $\eta^{\pm1}$; and $\nabla$ does not change the exponents of
$\eta$. So the exponents of $\eta$ in $(\partial H/\partial z_{v,c})_{S,a}$ lie between $\operatorname{val}_S(H)-1$ and
$\lambda_S(H)$ for a vector letter, and between $\operatorname{val}_S(H)$ and $\lambda_S(H)+1$ for a covector letter. The
composite $z_{v,c}\,\partial/\partial z_{v,c}$ lowers $\operatorname{val}_S$ by at most $1$ and raises $\lambda_S$ by at
most $1$, as in each case one of its two factors does neither.
% Lean: valS_pderiv, valS_pderiv_of, valS_X_mul and valS_X_mul_of (the lower bounds), maxTwistDegree_pderiv_le_add_one,
%   maxTwistDegree_pderiv_le_of, maxTwistDegree_X_mul_le_add_one and maxTwistDegree_X_mul_le_of (the upper bounds,
%   lam1S being maxTwistDegree), combined in lam1S_pderiv_and_X_mul_le_add_indicator (MaxTwistDegreePderivLeAddOne.lean
%   and the files of TwistedMonomialFiltration/).
\end{proof}

\begin{lemma}\label{pd-symbols-test-flat-chains-membership-one-plane}
Let $\mathcal D$ be a signed chain datum of rank $r$ with $J$ levels over $\mathcal R$
(Definition~\ref{pd-symbols-test-chains}). For $d'$ from the positions to $\mathbf Z_{\geq0}$ and real numbers
$\theta'_{j,i}$, one for each chain, and $\theta'^L_{j,m}$, for $m+2\leq r$, call $H\in\mathcal K[x]$ \emph{admissible for
$(d',\theta',\theta'^L)$} if it satisfies the conditions of admissibility for $\mathcal D$ with $d,\theta,\theta^L$
replaced by $d',\theta',\theta'^L$. Let $y_0$ be a position, $a<r$, $\xi=x_{y_0,a}$ and $\delta$ the function that is $1$
at $y_0$ and $0$ elsewhere. Suppose that $\theta_{j,i}\leq\theta'_{j,i}$ and $\theta^L_{j,m}\leq\theta'^L_{j,m}$ for all
$j,i,m$; that $\theta_{j,i}+1\leq\theta'_{j,i}$ whenever $y_0$ is a letter in $C_{j,i}$ or a link of level $j-1$; and that
$\theta^L_{j,m}+1\leq\theta'^L_{j,m}$ whenever $y_0\in S_{j,m}$. Let $H\in\mathcal K[x]$.
\begin{enumerate}
\item If $H$ is admissible for $(d+\delta,\theta',\theta'^L)$, then $\partial H/\partial\xi$ is admissible for $\mathcal D$.
\item If $H$ is admissible for $(d,\theta',\theta'^L)$, then $\xi\,\partial H/\partial\xi$ is admissible for $\mathcal D$.
\item If $d(y_0)\geq1$ and $H$ is admissible for $(d-\delta,\theta',\theta'^L)$, then $\xi H$ is admissible for
  $\mathcal D$.
\end{enumerate}
\end{lemma}
\leanref{Formalism/SharpRestriction/GrowthBounds/PlaneCoefficients/NoOpStep.lean}{Formalism.MainNotions.PairMembershipBlocks.ThDrop}
\leanref{Formalism/SharpRestriction/GrowthBounds/PlaneCoefficients/NoOpStep.lean}{Formalism.MainNotions.PairMembershipBlocks.PlDer}
\leanref{Formalism/SharpRestriction/GrowthBounds/PlaneCoefficients/NoOpStep.lean}{Formalism.MainNotions.PairMembershipBlocks.PlEul}
\leanref{Formalism/SharpRestriction/GrowthBounds/PlaneCoefficients/NoOpStep.lean}{Formalism.MainNotions.PairMembershipBlocks.PlMul}
\leanref{Formalism/SharpRestriction/GrowthBounds/PlaneCoefficients/XXgZeroP.lean}{Formalism.MainNotions.PairMembershipBlocks.plDer}
\leanref{Formalism/SharpRestriction/GrowthBounds/PlaneCoefficients/XXgZeroP.lean}{Formalism.MainNotions.PairMembershipBlocks.plEul}
\leanref{Formalism/SharpRestriction/GrowthBounds/PlaneCoefficients/XXgZeroP.lean}{Formalism.MainNotions.PairMembershipBlocks.plMul}
\leanref{Formalism/SharpRestriction/GrowthBounds/PlaneCoefficients/TouchCOfMem.lean}{Formalism.MainNotions.PairMembershipBlocks.Fdig_der_ex}
\leanref{Formalism/SharpRestriction/GrowthBounds/PlaneCoefficients/TouchCOfMem.lean}{Formalism.MainNotions.PairMembershipBlocks.Fdig_mul_ex}

\begin{proof}
Let $\tilde y_0$ be the model position of $y_0$: $y_0$ itself if it is a letter, its lift if it is a link; it has the
type of $y_0$, and $\tilde\xi=\tilde x_{\tilde y_0,a}$. As every chain contains all the $m<r$, $\tilde y_0\in\tilde C_{j,i}$
exactly when $y_0$ is a letter in $C_{j,i}$ or a link of level $j-1$.

\emph{Step 1: the residues.} If $y_0$ is a letter, then $(\partial H/\partial\xi)^\rho=\partial H^\rho/\partial\tilde\xi$
and $(\xi H)^\rho=\tilde\xi H^\rho$ for every $\rho$, as $F\mapsto F^\rho$ does not touch the letter variables. If $y_0$
is a link $(j_0,m_0)$, put $P=P(j_0,m_0)$ and $t=\rho(j_0,m_0,a)$. The computation of
Lemma~\ref{pd-symbols-test-flat-chains-membership-one-transport}, with $P$ in place of $\lambda$ and $\tilde\xi$ in place
of the level variable, gives: $(\partial H/\partial\xi)^\rho=(t+1)H^{\rho[t+1]}+P\tilde\xi\,\partial
H^{\rho[t+1]}/\partial\tilde\xi$ if $t+1<P$, $=P\,\partial H^{\rho[0]}/\partial\tilde\xi$ if $t+1=P$, and $=0$ if $t\geq P$;
and $(\xi H)^\rho=H^{\rho[t-1]}$ if $0<t<P$, $=\tilde\xi H^{\rho[P-1]}$ if $t=0$, and $=0$ if $t\geq P$. So for $T$ one of
$\partial/\partial\xi$, $\xi$ and $\xi\,\partial/\partial\xi$, each $(TH)^\rho$ is a linear combination of polynomials
$\tilde T(H^{\rho'})$, with $\tilde T$ a composite of at most one $\partial/\partial\tilde\xi$ and at most one
multiplication by $\tilde\xi$, and of polynomials $H^{\rho'}$.

\emph{Step 2: the valuations.} By Lemma~\ref{pd-symbols-test-flat-chains-membership-one-twist}, for a set $S$ of model
positions such a $\tilde T$ lowers $\operatorname{val}_S$ by at most $1$, and not at all if $\tilde y_0\notin S$: the
derivative lowers it only at a vector $\tilde y_0$, the multiplication only at a covector one. The same holds for the
operators $T$ on $\mathcal K[x]$ and the sets $S_{j,m}$. By Lemma~\ref{pd-symbols-residue-count-inputs-segment-clauses}(1),
$\operatorname{val}_S$ of a sum is at least the minimum, and scalars do not lower it.

\emph{Step 3.} Let $(T,d')$ be $(\partial/\partial\xi,d+\delta)$, $(\xi\,\partial/\partial\xi,d)$ or $(\xi,d-\delta)$, and
$H$ admissible for $(d',\theta',\theta'^L)$. Then $TH$ has degree $d(y)$ in the row of every $y$: in the second case, if
$d(y_0)=0$, then $H$ does not involve $\xi$ and $TH=0$. By Steps 1 and 2,
$\operatorname{val}_{\tilde C_{j,i}}((TH)^\rho)\geq\theta'_{j,i}-1\geq\theta_{j,i}$ if $\tilde y_0\in\tilde C_{j,i}$, and
$\geq\theta'_{j,i}\geq\theta_{j,i}$ otherwise; and $\operatorname{val}_{S_{j,m}}(TH)\geq\theta^L_{j,m}$ in the same way. So
$TH$ is admissible for $\mathcal D$.
% Lean: plDer = plDer' and plMul = plMul' (XXgZeroP.lean), with Fdig_der_ex and Fdig_mul_ex (TouchCOfMem.lean: Fdig_der_inl,
%   Fdig_der_lt, Fdig_der_eq, Fdig_mul_inl, Fdig_mul_zero, Fdig_mul_pos and Fdig_big) for Step 1, touchC_of_liftVar_mem_chainML
%   and modelLet_vec for the model position, valuation_ge_of_le_add_one, valS_pderiv(_of), valS_X_mul(_of),
%   flt_valS_add and valS_smul_ge for Step 2, and Pol_pderiv and Pol_X_mul for the degrees. plEul composes plDer' and
%   plMul' through d - delta, with the drop taken by the derivative at a vector y0 and by the multiplication at a
%   covector one, and pderiv_eq_zero_of_notMem_vars when d(y0) = 0.
\end{proof}

\begin{lemma}\label{pd-symbols-test-flat-chains-membership-one-normal}
Let $p$, $n'$, $\mathcal R$ and $\mathcal K$ be as in Lemma~\ref{pd-symbols-test-chains-projection-clauses}(1), $K\geq1$
an integer, and $\mathcal D$ a signed chain datum of rank $r$ with $J$ levels over $\mathcal R$. Form $\mathcal K[w]$,
the families $(E^{j,u}_{ab})$, the joint isotypic projectors and the plane coefficients as in
Definition~\ref{pd-symbols-test-flat-chains-tied}(3) to (5), with the indices and the sets $N_j(u)$ of its part (1), and
use the admissibility for $(d,\theta,\theta^L)$ of Lemma~\ref{pd-symbols-test-flat-chains-membership-one-plane}. Let
$y_0$ be a position, $1\leq b\leq K$, $\omega=w_{y_0,b}$, and $T$ one of $\partial/\partial\omega$, the multiplication by
$\omega$, and $\omega\,\partial/\partial\omega$, acting on $\mathcal K[w][x]$ through the coefficients. Let $G\in\mathcal
K[w][x]$, $\kappa'$ a family of weights $\kappa'(j,u)\in\mathbf Z^K$, $0\leq u\leq n_j+r$, and $e$ a joint isotypic
projector of weight $\kappa'$. Suppose that $[e_1G]_\nu$ is admissible for $(d,\theta,\theta^L)$ for every $\nu$, every
family $\kappa$ with $\kappa'(j,u)_1\leq\kappa(j,u)_1+1$ if $y_0\in N_j(u)$ and $\kappa'(j,u)_1\leq\kappa(j,u)_1$ otherwise,
and every joint isotypic projector $e_1$ of weight $\kappa$. Then $[eTG]_\nu$ is admissible for $(d,\theta,\theta^L)$ for
every $\nu$.
\end{lemma}
\leanref{Formalism/SharpRestriction/GrowthBounds/PlaneCoefficients/NoOpStep.lean}{Formalism.MainNotions.PairMembershipBlocks.NoOpStep}
\leanref{Formalism/SharpRestriction/GrowthBounds/MonomialSupport/Lam1SPderivAndXMulLeAddIndicator.lean}{Formalism.MainNotions.LetterSubstitution.noOpStep}
\leanref{Formalism/SharpRestriction/GrowthBounds/TwistedSubstitution/JointEIntList.lean}{Formalism.MainNotions.TestSideWords.eeN_comm}
\leanref{Formalism/SharpRestriction/GrowthBounds/TwistedSubstitution/JointEIntList.lean}{Formalism.MainNotions.TestSideWords.jointN_props}

\begin{proof}
Let $I$ be the set of the pairs $i=(j,u)$ with $0\leq u\leq n_j+r$, and put $S_i=N_j(u)$, with the links as vector
letters. As $\mathcal K$ has characteristic $0$, for every $i$ and $\mu\in\mathbf Z^K$ there is an isotypic projector
$e_{i,\mu}$ of weight $\mu$ for $S_i$, and every $z\in\mathcal K[w]$ is the sum of the finitely many nonzero $e_{i,\mu}z$
(Lemma~\ref{pd-symbols-test-chains-projection-ring-projectors}(1), with $K$ in place of $r$). For a family $\kappa$ put
$e_\kappa=\prod_ie_{i,\kappa(i)}$.

\emph{Step 1: the joint projectors.} Let $i,i'\in I$. If they have the same level, then $S_i\subset S_{i'}$ or
$S_{i'}\subset S_i$, and Lemma~\ref{pd-symbols-test-chains-projection-clauses}(1)(b) with $N=2$ shows that $e_{i,\mu}$ and
$e_{i',\mu'}$ commute. If not, $S_i\cap S_{i'}=\emptyset$, and as in the proof of
Lemma~\ref{pd-symbols-test-chains-projection-ring-projectors}(3) they commute. So $e_\kappa$ does not depend on the order
of the factors. It is the identity on $V(\kappa)$, as each factor is. If $\kappa''\neq\kappa$, say
$\kappa''(i)\neq\kappa(i)$, then $e_{i,\kappa(i)}$ vanishes on the isotypic subspace of weight $\kappa''(i)$ for $S_i$,
which contains $V(\kappa'')$, so $e_\kappa$ vanishes on $V(\kappa'')$. Hence $e_\kappa$ is a joint isotypic projector of
weight $\kappa$. Also $e_\kappa z$ lies in the isotypic subspace of weight $\kappa(i)$ for $S_i$, for every $i$, so in
$V(\kappa)$; and every $z$ is $\sum_\kappa e_\kappa z$ over a finite set of families, which may be enlarged.

\emph{Step 2: the families that are killed.} Let $\kappa$ be a family with $\kappa(i)_1+[y_0\in S_i]<\kappa'(i)_1$ for
some $i$, where $[y_0\in S_i]$ is $1$ or $0$. Let $z\in\mathcal K[w]$ and $z'=e_\kappa z$. By
Lemma~\ref{pd-symbols-test-chains-projection-clauses}(2)(b), $\lambda_{S_i}(z')\leq\kappa(i)_1$, as $e_{i,\mu}z'=0$ for
$\mu\neq\kappa(i)$. By Lemma~\ref{pd-symbols-test-flat-chains-membership-one-twist},
$\lambda_{S_i}(Tz')\leq\kappa(i)_1+[y_0\in S_i]<\kappa'(i)_1$, so $e_{i,\kappa'(i)}Tz'=0$ by the same part (2)(b). Write
$Tz'=\sum_{\kappa''}e_{\kappa''}Tz'$ over a finite set containing $\kappa'$. The term of $\kappa'$ has the factor
$e_{i,\kappa'(i)}$ and vanishes; the others lie in $V(\kappa'')$ with $\kappa''\neq\kappa'$, where $e$ vanishes. So
$eTz'=0$.

\emph{Step 3.} Let $A$ be a finite set of families with $G_\alpha=\sum_{\kappa\in A}e_\kappa G_\alpha$ for every
coefficient $G_\alpha\in\mathcal K[w]$ of $G$. Then $[eTG]_\nu=\sum_{\kappa\in A}[eTe_\kappa G]_\nu$. If $\kappa$ fails
the condition of the statement, then $\kappa'(i)_1>\kappa(i)_1+[y_0\in S_i]$ for some $i$, and the term vanishes by
Step 2. Otherwise $[eTe_\kappa G]_\nu=\sum_{\nu'}c_{\nu\nu'}[e_\kappa G]_{\nu'}$, a finite sum with $c_{\nu\nu'}$ the
coefficient of $w^\nu$ in $eTw^{\nu'}$, and each $[e_\kappa G]_{\nu'}$ is admissible for $(d,\theta,\theta^L)$ by the
hypothesis with $e_1=e_\kappa$. The polynomials admissible for $(d,\theta,\theta^L)$ form a $\mathcal K$-subspace, as
$\operatorname{val}_S$ of a sum is at least the minimum and $F\mapsto F^\rho$ is linear. So $[eTG]_\nu$ is admissible.
% Lean: noOpStep (MonomialSupport/Lam1SPderivAndXMulLeAddIndicator.lean): the projectors are eeN D K Dt i
%   (JointEIntList.lean, from polIsoDecompB_of), commuting by eeN_comm (filtered_projector_corollary_field, the clause
%   (1)(b), for one level, and comm_pres_isoSub with comm_isoProj for two), with jointE and jointN_props for Step 1 and
%   joint_dec for the decomposition. Step 2 is hkill (valuation_field, the clause (2)(b), and hTb from
%   lam1S_pderiv_and_X_mul_le_add_indicator); Step 3 is hsplit, planeCoeff_sum, XXg_sum and the case split hC, with
%   oneT_mul, planeCoeff_oneT and XXg_smul. MonomialSupport/ is 814 lines.
\end{proof}

\begin{lemma}\label{pd-symbols-test-flat-chains-membership-one}
Let $C_{La},C_k,C_{rk},C_{\deg}$ be real numbers with the property of Theorem~\ref{pd-symbols-count}, and $C$ a real
number. Let the data be as in Lemma~\ref{pd-symbols-test-flat} for this $C$, without $e$, $\gamma$, $g$ and the
$\Phi_b$, and use Definition~\ref{pd-symbols-test-flat-chains-membership-loss} for $\mathcal R$, $n'$, $r$, $K$, $T$,
$(J,f)$, $N_1$, $G'$, $c'$, $\deg$ and $Y$. Let $L$ be a loss and $F\in\mathcal K[X]$ a polynomial that satisfies the
membership condition with the margins $(0,\frac18)$ and the loss $L$. Then for every $s\in T$ and all $a,b\leq k$, the
polynomial $X_{s,a}\,\partial F/\partial X_{s,b}$ satisfies the membership condition with the margins $(0,\frac18)$ and
the loss $L+c_s$.
\end{lemma}
\leanref{Formalism/SharpRestriction/GrowthBounds/TwistedSubstitution/WordLossNonneg.lean}{Formalism.MainNotions.TestSideWords.OneOpStep}
\leanref{Formalism/SharpRestriction/GrowthBounds/PlaneCoefficients/RestOneOpOf.lean}{Formalism.MainNotions.PairMembershipBlocks.oneOpStep_of}
\leanref{Formalism/SharpRestriction/GrowthBounds/PlaneCoefficients/RestOneOpOf.lean}{Formalism.MainNotions.PairMembershipBlocks.OpSteps}

\begin{proof}
The ring $\mathcal R$ of the Hida setting is complete, finite over $W$ and has $\varpi_0^2=p$, so with $p\geq13$ it
satisfies the hypotheses of Lemma~\ref{pd-symbols-test-chains-projection-clauses}(1); the run system $(J,f)$ is valid for
$(L_1,N_1)$. For a loss $L'$ and $G\in\mathcal K[X]$ write $M_{L'}(G)$ for the membership condition with the margins
$(0,\frac18)$ and the loss $L'$, and say that a datum is tied with the loss $L'$ if it is tied to $(J,f)$ with some
$d^{\mathrm{pl}}$ and $\kappa$, the margins $(0,\frac18)$ and the loss $L'$. Put $c^1_s=c_s/2$, with the values $0$ and
$1$. We show, for every loss $L'$, every $F'$ with $M_{L'}(F')$ and every $c\leq k$:
\begin{enumerate}
\item[(A)] $M_{L'+c^1_s}(\partial F'/\partial X_{s,c})$;
\item[(B)] $M_{L'+c^1_s}(X_{s,c}F')$.
\end{enumerate}
Then (A) with $L'=L$ and $c=b$, and (B) with $L'=L+c^1_s$ and $c=a$, give the lemma, as $c_s=2c^1_s$.

\emph{Step 1: the datum.} Let $\mathcal D$ be tied with the loss $L'+c^1_s$, with $d^{\mathrm{pl}}$ and $\kappa$, and let
$d^{\mathrm{no}}$, $\rho$, a joint isotypic projector $e$ of weight $\kappa$ and $\nu$ be given. By
Lemma~\ref{pd-symbols-test-flat-chains-membership-one-position}(1) there is exactly one position $y_0$ whose place has
residue $s$. Let $\mathcal D^\uparrow$ be $\mathcal D$ with the thresholds $\theta^\uparrow_{j,i+1}=\theta_{j,i+1}+
c^1_s(\hat\jmath,i)$ and $\theta^{L\uparrow}_{j,m}=\theta^L_{j,m}+c^1_s(\hat\jmath,n_j+r-1-m)$. These are at least the old
ones, so $\mathcal D^\uparrow$ is a signed chain datum, tied with $d^{\mathrm{pl}}$, $\kappa$, the margins $(0,\frac18)$ and
the loss $L'$; so is $\mathcal D$, as $c^1_s\geq0$; and so are these data with $d$ replaced by $d^{\mathrm{pl}}\pm\delta$,
for $d^{\mathrm{pl}}\pm\delta$ in place of $d^{\mathrm{pl}}$, where $\delta$ is $1$ at $y_0$ and $0$ elsewhere. The
thresholds of $\mathcal D^\uparrow$ satisfy the hypotheses of Lemma~\ref{pd-symbols-test-flat-chains-membership-one-plane}
for $y_0$. Indeed, if $y_0$ is a letter in $C_{j,i+1}$, then $y_0\in N_j(i)\subset N_j(i+1)$; if $y_0\in S_{j,m}$, then
$y_0\in N_j(u+1)$ with $u=n_j+r-1-m$ the index of $(j,m)$, as the letters of level $j$ and the links $(j,m')$ with $m'>m$
have indices less than $u$; and in these cases, and when $y_0$ is a link of level $j-1$,
Lemma~\ref{pd-symbols-test-flat-chains-membership-one-position}(2) gives $c^1_s=1$ at the relevant index.

Let $\kappa_1$ be a family with $\kappa(j,u)_1\leq\kappa_1(j,u)_1+1$ if $y_0\in N_j(u)$ and $\kappa(j,u)_1\leq\kappa_1(j,u)_1$
otherwise. Then $\mathcal D$ is tied with $d^{\mathrm{pl}}$, $\kappa_1$, the margins $(0,\frac18)$ and the loss $L'$. For
$\theta^*_j(u)$ formed with $\kappa$ is at most $\theta^*_j(u)$ formed with $\kappa_1$ plus $c^1_s(\hat\jmath,u)$: it
involves $\kappa$ only through $\kappa(j,u+1)_1$ and $\kappa(j,u-1)_1$, and if $y_0$ lies in $N_j(u+1)$ or in
$N_j(u-1)\subset N_j(u+1)$, then $c^1_s(\hat\jmath,u)=1$ by
Lemma~\ref{pd-symbols-test-flat-chains-membership-one-position}(2).

The projector $e$ acts on the coefficients in $\mathcal K[w]$, so it commutes with $\partial/\partial x_{y,a}$ and with
the multiplication by $x_{y,a}$, and so does $G\mapsto[G]_\nu$. Write $L_{\rho'}$ for $L_{\rho',d^{\mathrm{pl}},d^{\mathrm{no}}}$.

\emph{Step 2: (A).} Let $c\leq r$ and $\xi=x_{y_0,c-1}$. By
Lemma~\ref{pd-symbols-test-flat-chains-membership-one-transport}(1), (2) and (5) there are $c_1,c_2,c_3\in\mathcal K$ and
$\rho_1,\rho_2$ with
\[
 \big[eL_\rho(\partial F'/\partial X_{s,c})\big]_\nu=c_1[eL_{\rho_1}(F')]_\nu+c_2\,\xi\frac{\partial}{\partial\xi}
 [eL_{\rho_1}(F')]_\nu+c_3\frac{\partial}{\partial\xi}[eL^+_{\rho_2}(F')]_\nu .
\]
By $M_{L'}(F')$ for $\mathcal D$, for $\mathcal D^\uparrow$ and for $\mathcal D^\uparrow$ with $d^{\mathrm{pl}}+\delta$, the
polynomial $[eL_{\rho_1}(F')]_\nu$ is admissible for $\mathcal D$ and for $(d,\theta^\uparrow,\theta^{L\uparrow})$, and
$[eL^+_{\rho_2}(F')]_\nu$ for $(d+\delta,\theta^\uparrow,\theta^{L\uparrow})$. By
Lemma~\ref{pd-symbols-test-flat-chains-membership-one-plane}(2) and (1) the three terms are admissible for $\mathcal D$, and
so is their sum. Let $c>r$ and $\omega=w_{y_0,c-r}$. The same parts of
Lemma~\ref{pd-symbols-test-flat-chains-membership-one-transport} give
$[eL_\rho(\partial F'/\partial X_{s,c})]_\nu=c_1[eL_{\rho_1}(F')]_\nu+c_2[e\,\omega\partial_\omega L_{\rho_1}(F')]_\nu+
c_3[e\,\partial_\omega L^+_{\rho_2}(F')]_\nu$, with $L^+$ formed with $d^{\mathrm{no}}+\delta$. The first term is admissible
for $\mathcal D$ as before. For the others apply Lemma~\ref{pd-symbols-test-flat-chains-membership-one-normal} to $\mathcal
D$, $G=L_{\rho_1}(F')$, resp.\ $L^+_{\rho_2}(F')$, $\kappa'=\kappa$ and $e$: its hypothesis is $M_{L'}(F')$ for $\mathcal
D$ tied with the families $\kappa_1$ of Step~1.

\emph{Step 3: (B).} By Lemma~\ref{pd-symbols-test-flat-chains-membership-one-transport}(3), (4) and (5) there are
$c_1,c_2\in\mathcal K$ and $\rho_1,\rho_2$ with $L_\rho(X_{s,c}F')=c_1L_{\rho_1}(F')+c_2\,\xi L^-_{\rho_2}(F')$, with
$c_2=0$ if $d^\xi(y_0)=0$; here $\xi$ and $d^\xi$ are as in that lemma. For $c\leq r$ the second term, when $c_2\neq0$, is
admissible for $\mathcal D$ by $M_{L'}(F')$ for $\mathcal D^\uparrow$ with $d^{\mathrm{pl}}-\delta$ and
Lemma~\ref{pd-symbols-test-flat-chains-membership-one-plane}(3). For $c>r$ it is $c_2[e\,\omega L^-_{\rho_2}(F')]_\nu$, and
Lemma~\ref{pd-symbols-test-flat-chains-membership-one-normal} applies as in Step~2.
% Lean: oneOpStep_of D (opSteps_of noOpStep) C CLa (RestOneOpOf.lean), for every C. geoTop gives the unique position
%   x0 of the levels with residue s, geoTouch the chains and links it touches (TouchAt), derLev the derivative
%   pderiv (s, b) with loss cost1 (from TrDerPl, TrDerNo, PlDer, PlEul and NoOpStep) and mulLev the multiplication by
%   X (s, a) with loss cost1 (from TrMulPl, TrMulNo, PlMul and NoOpStep), and opCost = 2 cost1 (opCost_eq).
%   memAtLe_iff (UpdDModelLet.lean) rewrites the condition through Base and MemLev, which vary d and the thresholds of a
%   datum; the book varies the datum itself. Step 1 is thLUp_pos, thLe_up, thLe_weak, thDrop_of (with TouchAt) and
%   thLe_ka with inN_le_cost (PlaneCoeffMonomial.lean); Step 2 is derLev (var_cases for c <= r or c > r) and Step 3 is
%   mulLev, with XXg_add and XXg_smul. The hypothesis DataII is used only for the admissible ring (hG1, hK) and the
%   validity of the run system.
\end{proof}

\begin{lemma}\label{pd-symbols-test-flat-chains-membership-bound}
Let $C_{La},C_k,C_{rk},C_{\deg}$ be real numbers with the property of Theorem~\ref{pd-symbols-count}, and $C\geq16$. Let
the data be as in Lemma~\ref{pd-symbols-test-flat} for this $C$, without $e$, $\gamma$, $g$ and the $\Phi_b$, and use
Definition~\ref{pd-symbols-test-flat-chains-membership-loss}. Let $D_1\cdots D_l$ be a product of operators as in
Definition~\ref{pd-symbols-test-operators}, with loss $L$. Then $0\leq L(j,u)\leq s(j,u)/8$ for every integer $j$ and
$0\leq u<|\Lambda'_j|+r$.
\end{lemma}
\leanref{Formalism/SharpRestriction/GrowthBounds/TwistedSubstitution/WordLossNonneg.lean}{Formalism.MainNotions.TestSideWords.WordBoundStep}
\leanref{Formalism/SharpRestriction/GrowthBounds/PlaneCoefficients/WBound.lean}{Formalism.MainNotions.PairMembershipBlocks.wordBoundStep}
\leanref{Formalism/SharpRestriction/GrowthBounds/PlaneCoefficients/WBound.lean}{Formalism.MainNotions.PairMembershipBlocks.wordLoss_le}

\begin{proof}
Each $c_s$ takes the values $0$ and $2$, so $L\geq0$. Let $D_i=X_{s_i,a_i}\,\partial/\partial X_{s_i,b_i}$, with
$s_i\in T$ and each $s$ equal to at most $c_0$ of the $s_i$. If $c_{s_i}(j,u)\neq0$, then $s_i$ is the residue of an
element of $\Lambda_j$. So $L(j,u)$ is at most $2$ times the number of the $i$ with $s_i$ among the residues of the
elements of $\Lambda_j$, which is at most $|\Lambda_j|c_0$. A spaced run system is valid, so $|\Lambda_j|\leq N_1$, and
$L(j,u)\leq2N_1c_0$.

Put $x=x(j,u)$, an element of $\Lambda_j$ larger than $f(j-1,r-1)$. As in the proof of
Lemma~\ref{pd-symbols-division-numeric}(2), $B_j(x)\geq a/p^{G'}$, since $\deg(\bar o)\geq a$ for the residue of
$o=\operatorname{old}_j(x)$ and $\operatorname{dist}_j(o)<G'$. By (1) of Situation~\ref{pd-edge-situation},
$N_1\leq N_1^2\leq r|T|$, as $N_1\geq1$ by its definition (Definition~\ref{hf-good-orbit}); so $|T|\geq1$ and
$N_1c_0\leq c_*+r|T|c_0+2rn'c_{pp}$. With $k^2\leq k^3$, $1\leq r^2$, $1\leq|T|$,
$p^{G'}\leq p^{5G'}$ and $16\leq C$, the last condition of (3) of the situation gives
\[
 16N_1c_0\,Yk^2p^{G'}\leq Ck^3r^2|T|\,p^{5G'}\big(c_*+r|T|c_0+2rn'c_{pp}\big)Y\leq a .
\]
So $16N_1c_0\leq a/(p^{G'}Yk^2)\leq B_j(x)/(Yk^2)$, and as $16N_1c_0$ is an integer, $16N_1c_0\leq s(j,u)$. Hence
$L(j,u)\leq2N_1c_0\leq s(j,u)/8$.
% Lean: wordBoundStep C CLa (16 <= C) (WBound.lean), with wordLoss_le (sum_ite_mem, the bound card (letters jz) <= N1
%   of validity), posU_mem, Bo_ge (oldest_mem_segment_letters, flatDeg_bounds), the bound N1 <= r |T| from N1^2 <= r |T|,
%   and size_arith with the last condition h12 of the situation; wordLoss_nonneg for L >= 0. WBound.lean, 195 lines.
\end{proof}

\begin{lemma}\label{pd-symbols-test-flat-chains-membership-words}
Let $C_{La},C_k,C_{rk},C_{\deg}$ be real numbers with the property of Theorem~\ref{pd-symbols-count}, and $C\geq16$. Let
the data be as in Lemma~\ref{pd-symbols-test-flat} for this $C$, without $e$, $\gamma$, $g$ and the $\Phi_b$. If
$F\in\mathcal K[X]$ satisfies the membership condition with the margins $(0,\frac18)$
(Definition~\ref{pd-symbols-test-flat-chains-tied}) and $D_1\cdots D_l$ is a product of operators as in
Definition~\ref{pd-symbols-test-operators}, then $D_1\cdots D_lF$ satisfies the membership condition with the margins
$(\frac18,\frac38)$.
\end{lemma}
\leanref{Formalism/SharpRestriction/GrowthBounds/TwistedSubstitution/MemRed.lean}{Formalism.MainNotions.TestSideWords.WordConditionAt}
\leanref{Formalism/SharpRestriction/GrowthBounds/TwistedSubstitution/WordLossNonneg.lean}{Formalism.MainNotions.TestSideWords.restWordStep_of}
\leanref{Formalism/SharpRestriction/GrowthBounds/TwistedSubstitution/WordLossNonneg.lean}{Formalism.MainNotions.TestSideWords.restWord_of}
\leanref{Formalism/SharpRestriction/GrowthBounds/PlaneCoefficients/WithThModelLet.lean}{Formalism.MainNotions.PairMembershipBlocks.restLeOfG}
\leanref{Formalism/SharpRestriction/GrowthBounds/TwistedSubstitution/MemAtDZeroIff.lean}{Formalism.MainNotions.TestSideWords.memAtLe_mono}
\leanref{Formalism/SharpRestriction/GrowthBounds/TwistedSubstitution/MemAtDZeroIff.lean}{Formalism.MainNotions.TestSideWords.memAtG_of_le}

\begin{proof}
We use Definition~\ref{pd-symbols-test-flat-chains-membership-loss}. Two remarks on the thresholds. (a) Let $G$ satisfy
the condition with the margins $(a_1,a_2)$ and the loss $L$, and let $(a'_1,a'_2)$ and $L'$ be such that
$a_1s(j,u)+L(j,u)\leq a'_1s(j,u)+L'(j,u)$ for $u\leq|\Lambda'_j|$ and $a_2s(j,u)+L(j,u)\leq a'_2s(j,u)+L'(j,u)$ for
$|\Lambda'_j|<u<|\Lambda'_j|+r$. A datum tied with the margins $(a'_1,a'_2)$ and the loss $L'$ is then tied with the
margins $(a_1,a_2)$ and the loss $L$, as its thresholds are bounded above by the smaller right-hand sides; so $G$
satisfies the condition with $(a'_1,a'_2)$ and $L'$. (b) A datum tied with the margins $(a_1,a_2)$ is tied with these
margins and the loss $0$.

\emph{Step 1: the loss $0$.} Let $F$ satisfy the membership condition with the margins $(0,\frac18)$, and let
$\mathcal D$ be tied with the margins $(0,\frac18)$ and the loss $0$, with $d^{\mathrm{pl}}$, $d^{\mathrm{no}}$ and
$\kappa$. Let $\mathcal D'$ be $\mathcal D$ with $\theta_{j,i+1}$ replaced by $\theta^*_j(i)$ and $\theta^L_{j,m}$ by
$\theta^*_j(u)-\frac18s_j(u)$, $u=n_j+r-1-m$. These are at least the old thresholds, so the $\theta^L_{j,m}$ of
$\mathcal D'$ are positive, and $\mathcal D'$ is a signed chain datum tied with the margins $(0,\frac18)$. The level
images, the joint isotypic projectors and the plane coefficients do not involve the thresholds, so each plane
coefficient of the condition is admissible for $\mathcal D'$, and hence for $\mathcal D$, whose thresholds are smaller.
So $F$ satisfies the condition with the margins $(0,\frac18)$ and the loss $0$.

\emph{Step 2: the word.} Let $D_i=X_{s_i,a_i}\,\partial/\partial X_{s_i,b_i}$. By descending induction on $i$ and
Lemma~\ref{pd-symbols-test-flat-chains-membership-one}, $D_i\cdots D_lF$ satisfies the condition with the margins
$(0,\frac18)$ and the loss $c_{s_i}+\dots+c_{s_l}$, starting from Step 1 for $i=l+1$. For $i=1$ the loss is the loss
$L$ of $D_1\cdots D_l$.

\emph{Step 3: the margins.} By Lemma~\ref{pd-symbols-test-flat-chains-membership-bound}, $0\leq L(j,u)\leq s(j,u)/8$ for
$u<|\Lambda'_j|+r$; so $0\cdot s+L\leq\frac18s+0$ and $\frac18s+L\leq\frac38s+0$ there, as $s(j,u)\geq8L(j,u)\geq0$. By
remark (a), $D_1\cdots D_lF$ satisfies the condition with the margins $(\frac18,\frac38)$ and the loss $0$, and by remark
(b) it satisfies the membership condition with the margins $(\frac18,\frac38)$.
% Lean: restWordStep_of C CLa restLeOfG (oneOpStep_of D h C CLa) (wordBoundStep C CLa) (WordLossNonneg.lean), so the
%   constant is max(1, 16) = 16 (restWord_of with RestOneOp at 1 and RestWordBound at 16). Remark (a) is memAtLe_mono and
%   remark (b) memAtG_of_le; Step 1 is restLeOfG (withTh, XX_withTh, withTh_normE, the positivity thL_pos); Step 2 is
%   hind (induction over the list, opWord = foldr); Step 3 is memAtLe_mono with wordLoss_nonneg, then memAtG_of_le.
\end{proof}

\begin{lemma}\label{pd-symbols-test-flat-chains-membership-monomials-pairs-even}
Let the data be as in Proposition~\ref{pd-symbols-realize}, with $n'$ even and a run system $(J,f)$ valid for some $(L_1,N_1)$, and
let $g\in GL_k(W)$. Then every monomial of $\Sigma_g(\mathrm{SP}_g)$ (Definition~\ref{pd-symbols-twist}) involves only variables
$X_{s,b}$ with $s$ even and $b\leq r$.
\end{lemma}
\leanref{Formalism/SharpRestriction/GrowthBounds/PlaneCoefficients/PlaneCoeffOneTTens.lean}{Formalism.MainNotions.PairMembershipBlocks.RestSP}
\leanref{Formalism/SharpRestriction/GrowthBounds/PlaneCoefficients/PmatDetMem.lean}{Formalism.MainNotions.PairMembershipBlocks.restSP}

\begin{proof}
By Step~4 of the proof of Proposition~\ref{pd-symbols-realize}, $\Sigma_g(\mathrm{SP}_g)$ is a product of powers of determinants of
$r\times r$ matrices whose entries are powers of variables $X_{\bar v,b}$ with $b\leq r$ and $v$ one of the even integers $f(j,l)$ and
$a_l(t)$. As $n'$ is even, $\bar v$ is even. A determinant is a sum of signed products of its entries, so every monomial of the product
involves only such variables.
% Lean: restSP (PmatDetMem.lean): SPR is the product of powers of the determinants of the pmat matrices; pmat_det_mem
%   (det_mem_sub, X_mem_supp) with hfr (validity) and hAr (window_holds_even_tops) and posOf_even. The Lean states it for SPR
%   over any commutative ring R and any c' and Q.
\end{proof}

\begin{lemma}\label{pd-symbols-test-flat-chains-membership-monomials-pairs-image}
Let $p$, $n'$, $r$, $K$, $k=r+K$, $T$ and a run system $(J,f)$ valid for some $(L_1,N_1)$ be as in
Definition~\ref{pd-symbols-test-flat-chains-tied}, and let $\mathcal D$ be a signed chain datum of rank $r$ with $J$ levels over
$\mathcal R$ with $n_j=|\Lambda'_{\hat\jmath}|$ for every level $j$, with the places and the numbers $\lambda(y)$ of
Definition~\ref{pd-symbols-test-flat-chains-tied}(1). For $0\leq j<J$ let $A_j$ be a polynomial over $\mathcal R$ in the plane
variables $z_{v,c}$ with $v\in\Lambda_j$, and $G_j$ one in the normal variables $z_{v,c}$ with $v\in\Lambda_j$ new
(Definition~\ref{pd-symbols-module}). Let $X^\ell$ be a label and $X^\beta$ a monomial in the variables $X_{s,b}$ with $s$ even and
$b\leq r$, and put
\[
 F=\prod_j\phi_j(A_jG_j)\cdot X^\ell X^\beta,\qquad P=\prod_j\phi_j(A_j)\cdot X^\ell X^\beta,\qquad N=\prod_j\phi_j(G_j).
\]
For $H\in\mathcal K[X]$ let $L^{\mathrm{pl}}(H)=\sum_\mu H_\mu\prod_y\prod_{a<r}x_{y,a}^{\lfloor\mu(y,a+1)/\lambda(y)\rfloor}$ and
$L^{\mathrm{no}}(H)=\sum_\mu H_\mu\prod_y\prod_{b\leq K}w_{y,b}^{\lfloor\mu(y,r+b)/\lambda(y)\rfloor}$, and for a function $d$ on the positions
let $[\,\cdot\,]_d$ be the sum of the terms of row degrees $d$. Let $\rho$, $d^{\mathrm{pl}}$ and $d^{\mathrm{no}}$ be as in
Definition~\ref{pd-symbols-test-flat-chains-tied}(3). If $\ell(\bar y,c)+\beta(\bar y,c)\equiv\rho(y,c)\pmod{\lambda(y)}$ for every
position $y$ and $c\leq k$, then, writing $[L^{\mathrm{pl}}(P)]_{d^{\mathrm{pl}}}=\sum_\alpha c_\alpha x^\alpha$,
\[
 L_{\rho,d^{\mathrm{pl}},d^{\mathrm{no}}}(F)=\sum_\alpha c_\alpha x^\alpha\cdot\big[L^{\mathrm{no}}(N)\big]_{d^{\mathrm{no}}};
\]
otherwise $L_{\rho,d^{\mathrm{pl}},d^{\mathrm{no}}}(F)=0$.
\end{lemma}
\leanref{Formalism/SharpRestriction/GrowthBounds/PlaneCoefficients/PlaneCoeffOneTTens.lean}{Formalism.MainNotions.PairMembershipBlocks.RestT}
\leanref{Formalism/SharpRestriction/GrowthBounds/PlaneCoefficients/RenameMemSupported.lean}{Formalism.MainNotions.PairMembershipBlocks.restT}

\begin{proof}
\openstep{The factorization of the level image is not yet written out: the exponents of $\prod_j\phi_j(A_jG_j)$ at the variables of a
position are multiples of $\lambda(y)$, so the residues are those of $X^\ell X^\beta$, and the plane and normal variables separate.}
% Lean: restT (RenameMemSupported.lean), with TensZeroRight, PlaneCoeffMonomial and ProdRangeZmod: levImg of symG is
%   tens (compD dPl (levP (zPl A l ga))) (compD dNo (levN (zNo G))) when InLow rho (l + ga), and 0 otherwise.
% Size: about 850 lines (lean_size of restT).
\end{proof}

\begin{lemma}\label{pd-symbols-test-flat-chains-membership-monomials-pairs-normal}
In the setting of Lemma~\ref{pd-symbols-test-flat-chains-membership-monomials-pairs-image}, suppose that $p\geq5$, that $n'$ is even,
that $\mathcal R$ is complete, its fraction field $\mathcal K$ has characteristic $0$ and its residue field contains a field of
$p^{n'}$ elements, and that $\mathcal D$ satisfies the conditions of a tie of Definition~\ref{pd-symbols-test-flat-chains-tied}(2)
other than those on $d$ and on the thresholds. Let $\kappa(j,u)\in\mathbf Z^K$, let
$e$ be a joint isotypic projector of weight $\kappa$, and suppose that $e\big([L^{\mathrm{no}}(N)]_{d^{\mathrm{no}}}\big)\neq0$. Then
for every level $j$, every $0\leq u<n_j+r$ and every integer $M$, with the weights of $G_{\hat\jmath}$ formed with $n=K$:
if $\lambda_{\{y<x_j(u)\}}(G_{\hat\jmath})\leq M$ then $\kappa(j,u)_1\leq M$, and if $\lambda_{\{y\leq x_j(u)\}}(G_{\hat\jmath})\leq M$ then
$\kappa(j,u+1)_1\leq M$.
\end{lemma}
\leanref{Formalism/SharpRestriction/GrowthBounds/PlaneCoefficients/PlaneCoeffOneTTens.lean}{Formalism.MainNotions.PairMembershipBlocks.RestN}
\leanref{Formalism/SharpRestriction/GrowthBounds/PlaneCoefficients/CompDEqHomComp.lean}{Formalism.MainNotions.PairMembershipBlocks.restN}

\begin{proof}
\openstep{The bounds on the weight $\kappa$ are not yet written out: the normal level image of $\prod_j\phi_j(G_j)$, restricted to the
row degrees $d^{\mathrm{no}}$, lies in the sum of the isotypic subspaces whose first weights are bounded by the twisted degrees
$\lambda$ of the $G_{\hat\jmath}$, and $e$ is $0$ on the others.}
% Lean: restN D (CompDEqHomComp.lean), with levN_zNo, lam_prod_le, typeBound_proved and jproj_ne; it takes O with a
%   uniformizer, CharZero (KK O), AdmissibleDVR p n' O and TieS Fr Dt.
% Size: about 1,580 lines (lean_size of restN).
\end{proof}

\begin{lemma}\label{pd-symbols-test-flat-chains-membership-monomials-pairs-twisted}
Let the notation be as in Definition~\ref{pd-symbols-test-flat-chains-tied}, with the run system $(J,f)$ valid for some $(L_1,N_1)$
and the numbers $m_*$ and $B_j(x)$, and let $n_j=|\Lambda'_{\hat\jmath}|$ for $j\in\mathbf Z/J\mathbf Z$. Let $j$ be a level, $A$ a
polynomial in the plane variables $z_{v,c}$ with $v\in\Lambda_{\hat\jmath}$, $G$ a polynomial in the normal variables, and
$M\colon\mathbf Z\to\mathbf Z$ such that $A$, $G$ and $M$ satisfy the inequalities of Definition~\ref{pd-symbols-module}(2) at the new
$v\in\Lambda_{\hat\jmath}$, for the functions $m_{\hat\jmath}$ and $\theta_{\hat\jmath}$ of the symbol spaces
(Definition~\ref{pd-symbols-module}(4), with $m_*$). Let
$\kappa(j,u)\in\mathbf Z^K$ satisfy, for $0\leq u<n_j+r$ and every integer $M'$, the two implications of
Lemma~\ref{pd-symbols-test-flat-chains-membership-monomials-pairs-normal} with $G$ in place of $G_{\hat\jmath}$. Then
$\operatorname{val}_{\{y<x_j(u)\}}(A)\geq\theta^*_j(u)$ for $0\leq u<n_j+r$.
\end{lemma}
\leanref{Formalism/SharpRestriction/GrowthBounds/PlaneCoefficients/PlaneCoeffOneTTens.lean}{Formalism.MainNotions.PairMembershipBlocks.RestTh}
\leanref{Formalism/SharpRestriction/GrowthBounds/PlaneCoefficients/ValSCongrSupp.lean}{Formalism.MainNotions.PairMembershipBlocks.restTh}

\begin{proof}
\openstep{The comparison of $\operatorname{val}_{\{y<x_j(u)\}}(A)$ with the three terms of $\theta^*_j(u)$ is not yet written out: the bound
$\theta_{\hat\jmath}$ at the oldest letter gives $B/2$, and the bounds $M(v)+m_{\hat\jmath}(v)$ at the even place $x_j(u)$ and at the
odd place $x_j(u-1)$ give the other two.}
% Lean: restTh (ValSCongrSupp.lean and OldestFar.lean); the condition on A at the letters that are not new is not stated, as in
%   Definition pd-symbols-module(2).
% Size: about 340 lines (lean_size of restTh).
\end{proof}

\begin{lemma}\label{pd-symbols-test-flat-chains-membership-monomials-pairs-plane}
In the setting of Lemma~\ref{pd-symbols-test-flat-chains-membership-monomials-pairs-image}, suppose that $n'$ is even, $J\geq2$,
$\mathcal K$ has characteristic $0$, and $\mathcal D$ satisfies the conditions of a tie of
Definition~\ref{pd-symbols-test-flat-chains-tied}(2) other than those on $d$ and on the thresholds. Let $t(j,u)$ be real numbers with
$\operatorname{val}_{\{y<x_j(u)\}}(A_{\hat\jmath})\geq t(j,u)$ for every level $j$ and $0\leq u<n_j+r$, let $d^{\mathrm{pl}}$ be a
function on the positions, and put $P'=[L^{\mathrm{pl}}(P)]_{d^{\mathrm{pl}}}$. For a level $j$ let $\delta(j-1,m')$ be the degree of
$L^{\mathrm{pl}}(X^\ell X^\beta)$ in the row of the link $(j-1,m')$, and put
$w_j=\sum_{m'<r}\big(d^{\mathrm{pl}}(j-1,m')\dot-\delta(j-1,m')\big)/P(j-1,m')$. Then, in the notation of
Definition~\ref{pd-symbols-test-chains}:
\begin{enumerate}
\item $\operatorname{val}_{\tilde C_{j,i+1}}(P'^{\rho'})\geq t(j,i)$ for every level $j$, $0\leq i\leq n_j$ and residues $\rho'$;
\item $\operatorname{val}_{S_{j,m}}(P')\geq t(j,n_j+r-1-m)-w_j$ for every level $j$ and $m+2\leq r$.
\end{enumerate}
\end{lemma}
\leanref{Formalism/SharpRestriction/GrowthBounds/PlaneCoefficients/PlaneCoeffOneTTens.lean}{Formalism.MainNotions.PairMembershipBlocks.RestPl}
\leanref{Formalism/SharpRestriction/GrowthBounds/PlaneCoefficients/PlaneCoeffOneTTens.lean}{Formalism.MainNotions.PairMembershipBlocks.dwB}
\leanref{Formalism/SharpRestriction/GrowthBounds/AevalValuation.lean}{Formalism.MainNotions.TraceSpectra.restChain}
\leanref{Formalism/SharpRestriction/GrowthBounds/PlaneCoefficients.lean}{Formalism.MainNotions.PairMembershipBlocks.restLink}

\begin{proof}
\openstep{The bounds at the chains and at the links of the plane level image are not yet written out: the twisted degrees of
$\prod_j\phi_j(A_j)\cdot X^\ell X^\beta$ at the chains and at the sets $S_{j,m}$ are bounded by those of the $A_{\hat\jmath}$ below the places,
with the loss $w_j$ at the links coming from the degrees of the links of level $j-1$ that $X^\ell X^\beta$ does not fill.}
% Lean: RestPl is the pair of RestChain and RestLink (restPl_of_restChain_of_restLink, PlaneCoefficients.lean): restChain
%   (AevalValuation.lean and AevalValuation/) for (1), and restLink = level_link_of_degree_letters linkL lkDeg lkVars
%   (restLink_of with valNN, homDec, restTr; MemTSetIff, MultidegSingle, SupportedOfSupport, RestPlOf, FrameZeroLe,
%   ValuationGeOfLeAddOne) for (2). The book's val at C-tilde of P'^rho' is valS Dt.modelLet (Dt.chainML j i) of Dt.Fdig
%   rho' P', and the one at S_{j,m} is valS Dt.trueLet (Dt.TSet j m).
% Size: about 4,200 lines (lean_size of restChain and restLink together).
\end{proof}

\begin{lemma}\label{pd-symbols-test-flat-chains-membership-monomials-pairs-numbers}
Let $C_{La},C_k,C_{rk},C_{\deg}$ be real numbers with the property of Theorem~\ref{pd-symbols-count}, and $C\geq1$. Let the data be
as in Lemma~\ref{pd-symbols-test-flat}, without $e$, $\gamma$, $g$ and the $\Phi_b$. Then $J\geq2$. Let moreover $E'$, $d$, $d'$, $X^\ell$
and $X^\beta$ be as in Lemma~\ref{pd-symbols-test-flat-chains-membership-monomials-pairs}, with $X^\beta$ a monomial in the variables
$X_{s,b}$ with $s$ even and $b\leq r$; let $\mathcal D$ satisfy the conditions of a tie other than those on $d$ and on the
thresholds, and $\rho$, $d^{\mathrm{pl}}$ and $d^{\mathrm{no}}$ the identities of that lemma and the congruences of
Lemma~\ref{pd-symbols-test-flat-chains-membership-monomials-pairs-image}. Then, with $w_j$ as in
Lemma~\ref{pd-symbols-test-flat-chains-membership-monomials-pairs-plane}, $w_j\leq\frac18s_j(n_j+r-1-m)$ for every level $j$ and
$m+2\leq r$.
\end{lemma}
\leanref{Formalism/SharpRestriction/GrowthBounds/PlaneCoefficients/PlaneCoeffOneTTens.lean}{Formalism.MainNotions.PairMembershipBlocks.RestArithStep}
\leanref{Formalism/SharpRestriction/GrowthBounds/PlaneCoefficients/FortyEightBoundArith.lean}{Formalism.MainNotions.PairMembershipBlocks.arithCondition_of_one_le}

\begin{proof}
\openstep{The bound of the loss at the links is not yet written out: the degree $d^{\mathrm{pl}}(j-1,m')$ of a link exceeds the degree
filled by $X^\ell X^\beta$ by at most the degree $d'$ there divided by $\lambda$, which the size conditions make small compared with
$P(j-1,m')s_j(u)$.}
% Lean: arithCondition_of_one_le C CLa hC (FortyEightBoundArith.lean), with the constant 1; 2 <= J is its first part.
% Size: about 860 lines (lean_size of arithCondition_of_one_le).
\end{proof}

\begin{lemma}\label{pd-symbols-test-flat-chains-membership-monomials-pairs}
Let $C_{La},C_k,C_{rk},C_{\deg}$ be real numbers with the property of Theorem~\ref{pd-symbols-count}, and $C\geq1$. Let
the data be as in Lemma~\ref{pd-symbols-test-flat}, without $e$, $\gamma$, $g$ and the $\Phi_b$, and let
$S=\Sigma_g(\mathrm{SP}_g)$ (Step~3 of the proof of Lemma~\ref{pd-symbols-primitive}); we use
Definitions~\ref{pd-symbols-module} and~\ref{pd-symbols-test-flat-chains-tied}, with the functions $m_j$ and $\theta_j$
of the symbol spaces. Let $E'\colon\mathbf Z/n'\mathbf Z\to\mathbf Z_{\geq0}$ be any function, not necessarily the one
of the data, and put $d(t)=\deg(t)-E'(t)$ for $t\in T$ (or
$0$ if this is negative) and $d(s)=0$ for $s\notin T$. For $0\leq j<J$ let $A_jG_j$ be one of the products of
Definition~\ref{pd-symbols-module}(2) that span the segment module $X_j(m_j,\theta_j,o)$, before the condition on the
orders: $A_j$ in the plane variables $z_{v,c}$ with $v\in\Lambda_j$, $G_j$ in the normal variables $z_{v,c}$ with
$v\in\Lambda_j$ new, and some $M_j$ with the conditions there for $m_j$ and $\theta_j$. Let $X^\ell$ be a label and
$X^\beta$ a monomial of $S$, put $F=\phi_0(A_0G_0)\cdots\phi_{J-1}(A_{J-1}G_{J-1})X^\ell X^\beta$ and
$d'(t)=d(t)+\sum_c\beta(t,c)$. Let $\mathcal D$ be tied to $(J,f)$ with $d^{\mathrm{pl}}$, $\kappa$ and the margins
$(0,\frac18)$ (Definition~\ref{pd-symbols-test-flat-chains-tied}(2)), and let $\rho$ and
$d^{\mathrm{no}}$ satisfy $\sum_c\rho(y,c)+\lambda(y)\big(d^{\mathrm{pl}}(y)+d^{\mathrm{no}}(y)\big)=d'(\bar y)$ for every
position $y$, with $\bar y$ the residue of its place. Then for every joint isotypic projector $e$ of weight $\kappa$ and
every monomial $w^\nu$, the plane coefficient $\big[e\,L_{\rho,d^{\mathrm{pl}},d^{\mathrm{no}}}(F)\big]_\nu$ is admissible
for $\mathcal D$.
\end{lemma}
\leanref{Formalism/SharpRestriction/GrowthBounds/PlaneCoefficients/PlanePairsOnBlocksAt.lean}{Formalism.MainNotions.PairMembershipBlocks.PlanePairsOnBlocksAt}
\leanref{Formalism/SharpRestriction/GrowthBounds/PlaneCoefficients/PlanePairsOnBlocksAt.lean}{Formalism.MainNotions.PairMembershipBlocks.MemAtB}
\leanref{Formalism/SharpRestriction/GrowthBounds/PlaneCoefficients/PlaneCoeffOneTTens.lean}{Formalism.MainNotions.PairMembershipBlocks.rest5aPairBStep_of}

\begin{proof}
Let the data, $F$, $\mathcal D$, $\rho$, $d^{\mathrm{no}}$, $e$ and $\nu$ be as in the statement, and let $M_j$ be the functions of the
products $A_jG_j$ in Definition~\ref{pd-symbols-module}(2). The run system is spaced, so valid, and $n'$ is even. Put
$P=\prod_j\phi_j(A_j)\cdot X^\ell X^\beta$, $N=\prod_j\phi_j(G_j)$, $P'=[L^{\mathrm{pl}}(P)]_{d^{\mathrm{pl}}}$ and
$N'=[L^{\mathrm{no}}(N)]_{d^{\mathrm{no}}}$, as in Lemma~\ref{pd-symbols-test-flat-chains-membership-monomials-pairs-image}.

\emph{Step 1: the monomial of $S$.} By Lemma~\ref{pd-symbols-test-flat-chains-membership-monomials-pairs-even}, $X^\beta$ involves only
variables $X_{s,b}$ with $s$ even and $b\leq r$. By Lemma~\ref{pd-symbols-test-flat-chains-membership-monomials-pairs-numbers}, $J\geq2$.

\emph{Step 2: the level image.} If the congruences of
Lemma~\ref{pd-symbols-test-flat-chains-membership-monomials-pairs-image} fail for $X^\ell X^\beta$ and $\rho$, that lemma gives
$L_{\rho,d^{\mathrm{pl}},d^{\mathrm{no}}}(F)=0$, and the plane coefficient is $0$, which is admissible. Otherwise
$L_{\rho,d^{\mathrm{pl}},d^{\mathrm{no}}}(F)=\sum_\alpha c_\alpha x^\alpha\cdot N'$ with $P'=\sum_\alpha c_\alpha x^\alpha$. As $e$ is
$\mathcal K$-linear, the plane coefficient is $\big[e\,L_{\rho,d^{\mathrm{pl}},d^{\mathrm{no}}}(F)\big]_\nu=aP'$, where $a$ is the coefficient
of $w^\nu$ in $e(N')$. If $a=0$ it is $0$. Otherwise it is admissible if $P'$ is, as a nonzero scalar changes neither the row degrees
nor the numbers $\operatorname{val}$. So let $a\neq0$; then $e(N')\neq0$.

\emph{Step 3: the normal part.} The setting is unramified quadratic with $p\geq13$, so $\mathcal R$ is complete, $\mathcal K$ has
characteristic $0$ and the residue field contains a field of $p^{n'}$ elements. As $e(N')\neq0$,
Lemma~\ref{pd-symbols-test-flat-chains-membership-monomials-pairs-normal} gives, for every level $j$, $0\leq u<n_j+r$ and integer $M$:
if $\lambda_{\{y<x_j(u)\}}(G_{\hat\jmath})\leq M$ then $\kappa(j,u)_1\leq M$, and if $\lambda_{\{y\leq x_j(u)\}}(G_{\hat\jmath})\leq M$ then
$\kappa(j,u+1)_1\leq M$.

\emph{Step 4: the plane factors.} For each level $j$, $A_{\hat\jmath}$, $G_{\hat\jmath}$ and $M_{\hat\jmath}$ satisfy the inequalities of
Definition~\ref{pd-symbols-module}(2) for $m_{\hat\jmath}$ and $\theta_{\hat\jmath}$, and Step 3 gives the remaining hypothesis of
Lemma~\ref{pd-symbols-test-flat-chains-membership-monomials-pairs-twisted}. So
$\operatorname{val}_{\{y<x_j(u)\}}(A_{\hat\jmath})\geq\theta^*_j(u)$ for $0\leq u<n_j+r$.

\emph{Step 5: $P'$ is admissible.} As $\mathcal D$ is tied with $d^{\mathrm{pl}}$, $d=d^{\mathrm{pl}}$, and $P'$ has degree $d(y)$ in the
row of every position $y$. By Step 4 and Lemma~\ref{pd-symbols-test-flat-chains-membership-monomials-pairs-plane}, with
$t(j,u)=\theta^*_j(u)$: for every chain, $\operatorname{val}_{\tilde C_{j,i+1}}(P'^{\rho'})\geq\theta^*_j(i)=\theta_{j,i+1}$, as the first
margin is $0$; and for $m+2\leq r$, with $u=n_j+r-1-m$,
$\operatorname{val}_{S_{j,m}}(P')\geq\theta^*_j(u)-w_j\geq\theta^*_j(u)-\frac18s_j(u)=\theta^L_{j,m}$, by
Lemma~\ref{pd-symbols-test-flat-chains-membership-monomials-pairs-numbers}. So $P'$ is admissible.
% Lean: rest5aPairBStep_of restT (restN D) restTh restPl restSP C CLa (arithCondition_of_one_le C CLa hC)
%   (PlaneCoeffOneTTens.lean). RestT (restT, RenameMemSupported.lean, with TensZeroRight, PlaneCoeffMonomial and
%   ProdRangeZmod): the level image of symG is tens (compD dPl (levP (zPl A l ga))) (compD dNo (levN (zNo G))) when the
%   residues of l + ga are rho, and 0 otherwise. RestN (restN, CompDEqHomComp.lean, with levN_zNo, lam_prod_le,
%   typeBound_proved and jproj_ne): e of the normal part nonzero bounds ka j u 0 and ka j (u + 1) 0 by lam1S of G_j below
%   and at the place x_j(u). RestTh (restTh, ValSCongrSupp.lean and OldestFar.lean): then valS of A_j below x_j(u) is at
%   least thStar. RestPl (restChain, AevalValuation/ 663 lines, and restLink through linkL, lkDeg, lkVars and restTr,
%   MemTSetIff, MultidegSingle, SupportedOfSupport, RestPlOf, FrameZeroLe, ValuationGeOfLeAddOne): the chain bounds
%   t(j, i) and the link bounds t(j, u) - dwB. RestSP (restSP, PmatDetMem.lean): the monomials of SPR involve only
%   plane variables at even residues. RestArithStep (arithCondition_of_one_le, FortyEightBoundArith.lean, C1 = 1):
%   2 <= J and dwB <= slU / 8 at the links. Step 1 is hgaP and hJ2; Step 2 is hsym, the rewriting by hT and the case split
%   on InLow (XX_zero), planeCoeff_oneT_tens and he0 (XX_zero, XX_smul); Step 3 is hNN; Step 4 is hval; Step 5 is hch and
%   hlk with compD_mem_Pol and t4, t8, t9 of the tie, and hAr'. The open parts are in Lemmas
%   pd-symbols-test-flat-chains-membership-monomials-pairs-image, -normal, -twisted, -plane and -numbers.
\end{proof}

\begin{lemma}\label{pd-symbols-test-flat-chains-membership-monomials}
Let $C_{La},C_k,C_{rk},C_{\deg}$ be real numbers with the property of Theorem~\ref{pd-symbols-count}. There is
$C_M>0$ such that the following holds for every $C\geq C_M$. Let the data be as in Lemma~\ref{pd-symbols-test-flat},
without $e$, $\gamma$, $g$ and the $\Phi_b$, and let $S=\Sigma_g(\mathrm{SP}_g)$ (Step~3 of the proof of
Lemma~\ref{pd-symbols-primitive}). Then for every $b$ and every monomial $X^\beta$ that occurs in $S$, the polynomial
$P_bX^\beta$ satisfies the membership condition with the margins $(0,\frac18)$
(Definition~\ref{pd-symbols-test-flat-chains-tied}).
\end{lemma}
\leanref{Formalism/SharpRestriction/GrowthBounds/TwistedSubstitution/Lemma5aMonomial.lean}{Formalism.MainNotions.TestSideWords.Lemma5aMonomialAt}
\leanref{Formalism/SharpRestriction/GrowthBounds/TwistedSubstitution/Lemma5aMonomial.lean}{Formalism.MainNotions.TestSideWords.Lemma5aMonomial}
\leanref{Formalism/SharpRestriction/GrowthBounds/PlaneCoefficients/PlaneCoeffOneTTens.lean}{Formalism.MainNotions.PairMembershipBlocks.rest5aMono_of_cut}

\begin{proof}
Put $C_M=1$, and let $C\geq1$, the data, $b$ and $X^\beta$ be as in the statement. By the hypotheses of
Proposition~\ref{pd-symbols-cloud}, $P_b\in\mathrm{Gen}(X_\bullet,d)$ for some $E'$ and the $d$ formed from it as in
Lemma~\ref{pd-symbols-test-flat-chains-membership-monomials-pairs}: $P_b=\phi_0(\xi_0)\cdots\phi_{J-1}(\xi_{J-1})X^\ell$
with $\xi_j\in X_j$ and $X^\ell$ a label, of multidegree $d$. So $P_bX^\beta$ has multidegree $d'$, with
$d'(t)=d(t)+\sum_c\beta(t,c)$.

\emph{Step 1: the blocks.} Let $X^\mu$ have multidegree $d'$ and contribute to $L_{\rho,d^{\mathrm{pl}},d^{\mathrm{no}}}$.
For a position $y$ and $c\leq k$, $\mu(\bar y,c)=\rho(y,c)+\lambda(y)\lfloor\mu(\bar y,c)/\lambda(y)\rfloor$, and summing
over $c$ gives $\sum_c\rho(y,c)+\lambda(y)(d^{\mathrm{pl}}(y)+d^{\mathrm{no}}(y))=d'(\bar y)$. So for $(\rho,d^{\mathrm{pl}},
d^{\mathrm{no}})$ without this property the level image of $P_bX^\beta$ is $0$, and its plane coefficients are admissible.

\emph{Step 2: linearity.} Let $Q$ be the set of the $F\in\mathcal K[X]$ that satisfy the conclusion of
Lemma~\ref{pd-symbols-test-flat-chains-membership-monomials-pairs} for every $\mathcal D$, $d^{\mathrm{pl}}$, $\kappa$,
$\rho$ and $d^{\mathrm{no}}$ as there, every $e$ and every $\nu$. As in the proof of
Lemma~\ref{pd-symbols-test-flat-chains-membership}, $Q$ is a $\mathcal K$-subspace.

\emph{Step 3: the pairs.} Each $X_j$ lies in the $\mathcal R$-span of the products $A_jG_j$ of
Lemma~\ref{pd-symbols-test-flat-chains-membership-monomials-pairs}, by Definition~\ref{pd-symbols-module}(2) and~(4). The
map $(\xi_0,\dots,\xi_{J-1})\mapsto\phi_0(\xi_0)\cdots\phi_{J-1}(\xi_{J-1})X^\ell X^\beta$ is $\mathcal R$-multilinear,
and by Lemma~\ref{pd-symbols-test-flat-chains-membership-monomials-pairs} it maps every tuple of products into $Q$.
Replacing one entry at a time by an element of the span, it maps every tuple from the spans into $Q$. So $P_bX^\beta\in
Q$, and with Step~1 it satisfies the membership condition with the margins $(0,\frac18)$.
% Lean: rest5aMono_of_cut gives C1 = 1, the constant of arithCondition_of_one_le passed as hAr (stmtE_of_pl), and is
%   rest5aMono_of_B, that is rest5aMonoStep_of_B (PlanePairsOnBlocksAt.lean): P_b is symG l 0 xs by hD.2.2.2.1 (the
%   symbols B b with E', xs and the label l), Step 1 is memAtG_of_B (blockDeg_of), Step 2 is memSubB, and Step 3 is
%   symbol_spaces_le_pair_span and symG_span (TwistedSubstitution/Lemma5aGeneratorProducts.lean) with
%   PlanePairsOnBlocksAt (Lemma pd-symbols-test-flat-chains-membership-monomials-pairs).
%   What remains open is in Lemmas pd-symbols-test-flat-chains-membership-monomials-pairs-image (about 850 lines),
%   -normal (about 1,580), -twisted (about 340), -plane (about 4,200) and -numbers (about 860), about 6,500 lines
%   together; the
%   reduction above is PlanePairsOnBlocksAt.lean and Lemma5aGeneratorProducts.lean, about 300 lines.
\end{proof}

\begin{lemma}\label{pd-symbols-test-flat-chains-membership}
Let $C_{La},C_k,C_{rk},C_{\deg}$ be real numbers with the property of Theorem~\ref{pd-symbols-count}. There is $C_1>0$
such that the following holds for every $C\geq C_1$. Let the data be as in Lemma~\ref{pd-symbols-test-flat}, without $e$,
$\gamma$, $g$ and the $\Phi_b$, and let $S=\Sigma_g(\mathrm{SP}_g)$, which is the same polynomial for every
$g\in GL_k(W)$ (Step~3 of the proof of Lemma~\ref{pd-symbols-primitive}). Then every $G\in\mathrm{Op}_{T,c_0}$ of the
$\mathcal K$-span of the $P_bS$ satisfies the membership condition of
Definition~\ref{pd-symbols-test-flat-chains-tied} for $\mathcal R$, $n'$, $r$, $K$, $T$, $(J,f)$, $N_1$, $G'$, $c'$,
$\deg$ and $Y$.
\end{lemma}
\leanref{Formalism/SharpRestriction/GrowthBounds/TwistedSubstitution/Proposition5Gram.lean}{Formalism.MainNotions.TestSideWords.MemStep}
\leanref{Formalism/SharpRestriction/GrowthBounds/TwistedSubstitution/Proposition5Gram.lean}{Formalism.MainNotions.TestSideWords.RestMEM}
\leanref{Formalism/SharpRestriction/GrowthBounds/TwistedSubstitution/MemRed.lean}{Formalism.MainNotions.TestSideWords.memCondition_of_theoremII}
\leanref{Formalism/SharpRestriction/GrowthBounds/TwistedSubstitution/MemRed.lean}{Formalism.MainNotions.TestSideWords.memAt_of_wordConditionAt}
\leanref{Formalism/SharpRestriction/GrowthBounds/TwistedSubstitution/MemRed.lean}{Formalism.MainNotions.TestSideWords.Lemma5aAt}
\leanref{Formalism/SharpRestriction/GrowthBounds/TwistedSubstitution/MemRed.lean}{Formalism.MainNotions.TestSideWords.WordConditionAt}

\begin{proof}
Let $M$ be the $\mathcal K$-span of the $P_bS$. For real $a_1,a_2$ the polynomials that satisfy the membership condition
with the margins $(a_1,a_2)$ form a $\mathcal K$-subspace of $\mathcal K[X]$: the level image, the projector $e$ and the
plane coefficient are $\mathcal K$-linear, and the admissible polynomials for $\mathcal D$ form a subspace, as
$\operatorname{val}_S(F+F')\geq\min(\operatorname{val}_S(F),\operatorname{val}_S(F'))$, $\operatorname{val}_S(cF)\geq
\operatorname{val}_S(F)$ and $(F+F')^\rho=F^\rho+F'^\rho$. Consider, for data as in the statement, the two assertions:
\begin{enumerate}
\item each $P_bS$ satisfies the membership condition with the margins $(0,\frac18)$;
\item if $f$ satisfies the membership condition with the margins $(0,\frac18)$ and $D_1\cdots D_l$ is a product of
  operators as in Definition~\ref{pd-symbols-test-operators}, then $D_1\cdots D_lf$ satisfies it with the margins
  $(\frac18,\frac38)$.
\end{enumerate}
Let $C_M$ be the constant of Lemma~\ref{pd-symbols-test-flat-chains-membership-monomials}. For $C\geq C_M$, (1) holds:
write $S=\sum_\beta S_\beta X^\beta$, the sum over the monomials of $S$; then $P_bS=\sum_\beta S_\beta\,P_bX^\beta$, and each
$P_bX^\beta$ satisfies the condition with the margins $(0,\frac18)$ by that lemma, so $P_bS$ does, by the subspace
property. For $C\geq16$, (2) holds by Lemma~\ref{pd-symbols-test-flat-chains-membership-words}.
Put $C_1=\max(C_M,16)$, and let $C\geq C_1$ and the data be as in the statement. By (1) every $f\in M$ satisfies the
condition with the margins $(0,\frac18)$. By (2) the generators $D_1\cdots D_lf$ of $\mathrm{Op}_{T,c_0}(M)$ satisfy it
with the margins $(\frac18,\frac38)$, which is the membership condition, and so does every element of their span.
% Lean: RestMEM D CLa Ck Crk Cdeg is restMEM_of_last (PlaneCoefficients.lean) = memCondition_of_theoremII D ...
%   (TwistedSubstitution/MemRed.lean): memAt_of_wordConditionAt, by span induction over the span of the symP (B b) and
%   over OpCl, with memAtG_zero, memAtG_add, memAtG_smul and memAt_iff (MemAt is MemAtG (1/8) (3/8), MemAtG sc sL being
%   the condition for TwistedSubstitutionGrowthBound with margins sc * slU and sL * slU), from Lemma5aTheoremII
%   (MemAtG 0 (1/8) for each symP (B b): lemma5a_of_monomial, whose conclusionAt_of_monomialAt is the reduction of (1)
%   to the monomials with symP_eq_sum and memSub, and Lemma pd-symbols-test-flat-chains-membership-monomials) and
%   RestWord (WordConditionAt: Lemma pd-symbols-test-flat-chains-membership-words, restWord_of_steps). The constant is
%   max(C1, C2) of memCondition_of_theoremII; the Lean's C2 for RestWord is 16. What remains open is in Lemmas
%   pd-symbols-test-flat-chains-membership-monomials-pairs-image, -normal, -twisted, -plane and -numbers (about 6,500
%   lines together).
\end{proof}

\begin{definition}\label{pd-symbols-test-flat-chains-datum-assembly}
Let the notation be as in Definition~\ref{pd-symbols-test-flat-chains-tied}, with $s_j(u)$, $\lambda(y)$ and the places,
let $d_T\colon\mathbf Z/n'\mathbf Z\to\mathbf Z_{\geq0}$, and let $\mathcal D$ be a signed chain datum of rank $r$ with
$J$ levels over $\mathcal R$, with its rings $\mathcal K[x]$ and $\mathcal K[w]$. For a position $y$ let $\bar y$ be the
residue modulo $n'$ of its place.
\begin{enumerate}
\item For $\rho$ from the pairs of a position and an integer $1\leq c\leq k$ to $\mathbf Z_{\geq0}$, $A\in\mathcal K[x]$ and
  $N\in\mathcal K[w]$, the \emph{assembly} $\operatorname{asm}_\rho(A,N)\in\mathcal K[X]$ is
  \[
   \prod_y\prod_{c=1}^kX_{\bar y,c}^{\rho(y,c)}\cdot A\big(x_{y,a}\mapsto X_{\bar y,a+1}^{\lambda(y)}\big)\cdot
   N\big(w_{y,b}\mapsto X_{\bar y,r+b}^{\lambda(y)}\big),
  \]
  the products over the positions $y$, with $0\leq a<r$ and $1\leq b\leq K$.
\item $\mathcal D$ is \emph{adapted} to $(J,f)$ with $d^{\mathrm{pl}}$ and $\kappa$ if it is tied to $(J,f)$ with
  $d^{\mathrm{pl}}$ and $\kappa$, its number $c$ is $c'$, and $s_{j,i+1}=\frac38s_j(i)$ for every level $j$ and
  $0\leq i\leq n_j$.
\item A triple $(\rho,d^{\mathrm{pl}},d^{\mathrm{no}})$, with $\rho$ as in (1) and $d^{\mathrm{pl}},d^{\mathrm{no}}$ from the
  positions to $\mathbf Z_{\geq0}$, is a \emph{block decomposition} of $d_T$ for $\mathcal D$ if $\rho(y,c)<\lambda(y)$ and
  $\sum_{c=1}^k\rho(y,c)+\lambda(y)\big(d^{\mathrm{pl}}(y)+d^{\mathrm{no}}(y)\big)=d_T(\bar y)$ for every position $y$ and
  every $c$.
\item For $d^{\mathrm{no}}$ and $\kappa$ as in Definition~\ref{pd-symbols-test-flat-chains-tied}(2), a \emph{normal test}
  is an $N_0\in\mathcal K[w]$ with integral coefficients and degree $d^{\mathrm{no}}(y)$ in the row of every position
  $y$, such that $\lambda_{N_j(u)}(N_0)\leq\kappa(j,u)_1+4c'$ for every level $j$ and $0\leq u\leq n_j+r$; here
  $\lambda_S$ is formed on $\mathcal K[w]$ as in Definition~\ref{pd-symbols-path-generators}(1), with $n=K$ and the
  links taken as vector positions.
\end{enumerate}
\end{definition}
\leanref{Formalism/SharpRestriction/GrowthBounds/DiagonalSupport/TwiceValuationFactorialProduct.lean}{Formalism.MainNotions.Proposition5Numerics.asm}
\leanref{Formalism/SharpRestriction/GrowthBounds/DiagonalSupport/TwiceValuationFactorialProduct.lean}{Formalism.MainNotions.Proposition5Numerics.labMon}
\leanref{Formalism/SharpRestriction/GrowthBounds/DiagonalSupport/TwiceValuationFactorialProduct.lean}{Formalism.MainNotions.Proposition5Numerics.upPl}
\leanref{Formalism/SharpRestriction/GrowthBounds/DiagonalSupport/TwiceValuationFactorialProduct.lean}{Formalism.MainNotions.Proposition5Numerics.upNo}
\leanref{Formalism/SharpRestriction/GrowthBounds/DiagonalSupport/TwiceValuationFactorialProduct.lean}{Formalism.MainNotions.Proposition5Numerics.SignedChainDegreeCondition}
\leanref{Formalism/SharpRestriction/GrowthBounds/DiagonalSupport/TwiceValuationFactorialProduct.lean}{Formalism.MainNotions.Proposition5Numerics.BlockOK}
\leanref{Formalism/SharpRestriction/GrowthBounds/DiagonalSupport/TwiceValuationFactorialProduct.lean}{Formalism.MainNotions.Proposition5Numerics.NormTest}

\begin{definition}\label{pd-symbols-test-flat-chains-datum-pick-weights}
Let the notation be as in Definition~\ref{pd-symbols-test-flat-chains-datum-assembly}, with a signed chain datum
$\mathcal D$, and let $\rho$ satisfy $\rho(y,c)<\lambda(y)$ for all $y$ and $c$. For integers $P\geq1$ and $t,b\geq0$ put
$\gamma(P,t,b)=(t+Pb)!/(t!\,b!\,(P!)^b)$.
\begin{enumerate}
\item For an exponent vector $\beta$ of $\mathcal K[x]$ put $\gamma^{\mathrm{pl}}_\rho(\beta)=\prod_{y}\prod_{a<r}
  \gamma(\lambda(y),\rho(y,a+1),\beta(y,a))$, and for one of $\mathcal K[w]$ put $\gamma^{\mathrm{no}}_\rho(\beta)=\prod_y
  \prod_{b=1}^K\gamma(\lambda(y),\rho(y,r+b),\beta(y,b))$. For $d^{\mathrm{pl}},d^{\mathrm{no}}$ put
  $c_\rho(d^{\mathrm{pl}},d^{\mathrm{no}})=\prod_{y,c}\rho(y,c)!\cdot\prod_y(\lambda(y)!)^{d^{\mathrm{pl}}(y)+d^{\mathrm{no}}(y)}$.
  For a function $g$ on exponent vectors and a polynomial $H$ put $\operatorname{diag}_gH=\sum_\beta g(\beta)H_\beta
  z^\beta$, in the variables of $H$.
\item For $d^{\mathrm{no}}$ put $D_{\mathrm{no}}(d^{\mathrm{no}})=\sum_{(j,u)}N_0(K)\,\ell_p\big(2\sum_{y\in
  N_j(u)}d^{\mathrm{no}}(y)+K-1\big)$, the sum over the pairs with $0\leq u\leq n_j+r$, where $N_0(K)$ is the integer of
  Lemma~\ref{pd-kappa-defects-segment-mean-letters-polynomials} and $\ell_p$ is as in
  Lemma~\ref{pd-symbols-test-chains-projection-clauses}.
\item For a monomial $X^\mu$ its \emph{residues} are $\rho_\mu(y,c)$, the remainder of $\mu(\bar y,c)$ modulo $\lambda(y)$;
  its \emph{blocks} are $d^{\mathrm{pl}}_\mu(y)=\sum_{c\leq r}\lfloor\mu(\bar y,c)/\lambda(y)\rfloor$ and
  $d^{\mathrm{no}}_\mu(y)=\sum_{c>r}\lfloor\mu(\bar y,c)/\lambda(y)\rfloor$; and its \emph{plane and normal exponents} are
  $\beta^{\mathrm{pl}}_\mu(y,a)=\lfloor\mu(\bar y,a+1)/\lambda(y)\rfloor$ and
  $\beta^{\mathrm{no}}_\mu(y,b)=\lfloor\mu(\bar y,r+b)/\lambda(y)\rfloor$.
\end{enumerate}
\end{definition}
\leanref{Formalism/SharpRestriction/GrowthBounds/DiagonalSupport/PickDefs.lean}{Formalism.MainNotions.Proposition5Numerics.gamR}
\leanref{Formalism/SharpRestriction/GrowthBounds/DiagonalSupport/PickDefs.lean}{Formalism.MainNotions.Proposition5Numerics.gamPl}
\leanref{Formalism/SharpRestriction/GrowthBounds/DiagonalSupport/PickDefs.lean}{Formalism.MainNotions.Proposition5Numerics.gamNo}
\leanref{Formalism/SharpRestriction/GrowthBounds/DiagonalSupport/PickDefs.lean}{Formalism.MainNotions.Proposition5Numerics.cstOf}
\leanref{Formalism/SharpRestriction/GrowthBounds/DiagonalSupport/PickDefs.lean}{Formalism.MainNotions.Proposition5Numerics.diagK}
\leanref{Formalism/SharpRestriction/GrowthBounds/DiagonalSupport/PickDefs.lean}{Formalism.MainNotions.Proposition5Numerics.DnoOf}
\leanref{Formalism/SharpRestriction/GrowthBounds/DiagonalSupport/PickDefs.lean}{Formalism.MainNotions.Proposition5Numerics.classOf}
\leanref{Formalism/SharpRestriction/GrowthBounds/DiagonalSupport/PickDefs.lean}{Formalism.MainNotions.Proposition5Numerics.blkPl}
\leanref{Formalism/SharpRestriction/GrowthBounds/DiagonalSupport/PickDefs.lean}{Formalism.MainNotions.Proposition5Numerics.blkNo}

\begin{lemma}\label{pd-symbols-test-flat-chains-datum-pick-inputs}
Let $C_{La},C_k,C_{rk},C_{\deg}$ be real numbers with the property of Theorem~\ref{pd-symbols-count}. There are
$C_{Dn},C_S>0$ such that the following holds for every $C\geq C_S$. Let the data be as in
Lemma~\ref{pd-symbols-test-flat-chains-membership}, put $\mathcal L=|T|k^2\log_pa$, and use
Definitions~\ref{pd-symbols-test-chains}, \ref{pd-symbols-test-flat-chains-datum-assembly}
and~\ref{pd-symbols-test-flat-chains-datum-pick-weights} for signed chain data of rank $r$ with $J$ levels over
$\mathcal R$ with $n_j=|\Lambda'_{\hat\jmath}|$ for every level $j$. Let $\rho$ satisfy $\rho(y,c)<\lambda(y)$.
\begin{enumerate}
\item Let $G\in\mathcal K[X]$ have multidegree $d_T$. For all $d^{\mathrm{pl}},d^{\mathrm{no}}$ and $A\in\mathcal K[x]$,
  $N\in\mathcal K[w]$ of degrees $d^{\mathrm{pl}}(y)$, resp.\ $d^{\mathrm{no}}(y)$, in the rows,
  \[
   \langle G,\operatorname{asm}_\rho(A,N)\rangle_{\mathrm F}=c_\rho(d^{\mathrm{pl}},d^{\mathrm{no}})\sum_{\beta,\beta'}
   L_{\beta,\beta'}\cdot\gamma^{\mathrm{pl}}_\rho(\beta)A_\beta\,\beta!\cdot\gamma^{\mathrm{no}}_\rho(\beta')N_{\beta'}\,\beta'!,
  \]
  where $L_{\beta,\beta'}$ is the coefficient of $x^\beta w^{\beta'}$ in $L_{\rho,d^{\mathrm{pl}},d^{\mathrm{no}}}(G)$; each
  coefficient in $\mathcal K[w]$ of $L_{\rho,d^{\mathrm{pl}},d^{\mathrm{no}}}(G)$ has degrees $d^{\mathrm{no}}$; and for
  every monomial $X^\mu$ of $G$, its residues and blocks form a block decomposition of $d_T$, the coefficient of
  $x^{\beta^{\mathrm{pl}}_\mu}w^{\beta^{\mathrm{no}}_\mu}$ in $L_{\rho_\mu,d^{\mathrm{pl}}_\mu,d^{\mathrm{no}}_\mu}(G)$ is
  $G_\mu$, $\beta^{\mathrm{pl}}_\mu$ and $\beta^{\mathrm{no}}_\mu$ have the degrees $d^{\mathrm{pl}}_\mu$ and
  $d^{\mathrm{no}}_\mu$, and $\mu!=c_{\rho_\mu}(d^{\mathrm{pl}}_\mu,d^{\mathrm{no}}_\mu)\,\beta^{\mathrm{pl}}_\mu!\,
  \beta^{\mathrm{no}}_\mu!\,\gamma^{\mathrm{pl}}_{\rho_\mu}(\beta^{\mathrm{pl}}_\mu)\,\gamma^{\mathrm{no}}_{\rho_\mu}
  (\beta^{\mathrm{no}}_\mu)$.
\item The values of $\gamma^{\mathrm{pl}}_\rho$ and $\gamma^{\mathrm{no}}_\rho$ are nonzero rational numbers with
  $v_p=0$, and $\gamma(\alpha)=\gamma^{\mathrm{pl}}_\rho(\alpha)$ defines units of $\mathcal R$ with the property of
  Definition~\ref{pd-symbols-test-chains}(4).
\item Let $\mathcal D$ be any signed chain datum of rank $r$ with $J$ levels over $\mathcal R$, $d^{\mathrm{no}}$ given and
  $y\in\mathcal K[w]$ of degrees $d^{\mathrm{no}}$. Then
  there are $\kappa$ and a joint isotypic projector $e$ of weight $\kappa$ for $\mathcal D$ such that
  $\langle eF,F'\rangle_{\mathrm F}=\langle F,eF'\rangle_{\mathrm F}$ for all $F,F'$; $e$ maps the polynomials of degrees
  $d^{\mathrm{no}}$ to such polynomials, and those $F$ with $v_{\mathrm F}(F)\geq0$ to polynomials with
  $v_{\mathrm F}(p^{D_{\mathrm{no}}(d^{\mathrm{no}})}eF)\geq0$; each $eF$ lies in the isotypic subspace of weight
  $\kappa(j,u)$ for $N_j(u)$, for every $(j,u)$; and $v_{\mathrm F}(ey)\leq v_{\mathrm F}(y)$.
\item Let $\mathcal D$, $d^{\mathrm{no}}$, $\kappa$ and $e$ be as in the conclusion of (3), and $\nu$ an exponent vector of
  degrees $d^{\mathrm{no}}$. Then there are a normal test $N_0$ for $d^{\mathrm{no}}$ and $\kappa$ and $c_Y\in\mathcal K$
  with $v(c_Y)\leq v(\prod_yd^{\mathrm{no}}(y)!)/2+D_{\mathrm{no}}(d^{\mathrm{no}})-v(\nu!)/2$, such that the coefficients
  of $N_0-\operatorname{diag}_{1/\gamma^{\mathrm{no}}_\rho}(c_Y\,e(w^\nu))$ have valuation at least $c'$.
\item $c'\geq1$; $\frac38s_j(u)<B_{\hat\jmath}(x_j(u))/2$ for every level $j$ and $u<n_j+r$; and for every block
  decomposition $(\rho,d^{\mathrm{pl}},d^{\mathrm{no}})$ of $d_T$ and exponent vectors $\beta,\beta'$ of degrees
  $d^{\mathrm{pl}}$, resp.\ $d^{\mathrm{no}}$,
  \[
   D_{\mathrm{no}}(d^{\mathrm{no}})+\frac{v(\prod_yd^{\mathrm{pl}}(y)!)-v(\beta!)}2+\frac{v(\prod_yd^{\mathrm{no}}(y)!)-
   v(\beta'!)}2\leq C_{Dn}\mathcal L .
  \]
\end{enumerate}
\end{lemma}
\leanref{Formalism/SharpRestriction/GrowthBounds/DiagonalSupport/PickDefs.lean}{Formalism.MainNotions.Proposition5Numerics.DigFacts}
\leanref{Formalism/SharpRestriction/GrowthBounds/DiagonalSupport/PickDefs.lean}{Formalism.MainNotions.Proposition5Numerics.GaAt}
\leanref{Formalism/SharpRestriction/GrowthBounds/DiagonalSupport/PickDefs.lean}{Formalism.MainNotions.Proposition5Numerics.JPAt}
\leanref{Formalism/SharpRestriction/GrowthBounds/DiagonalSupport/PickDefs.lean}{Formalism.MainNotions.Proposition5Numerics.NoAt}
\leanref{Formalism/SharpRestriction/GrowthBounds/DiagonalSupport/PickDefs.lean}{Formalism.MainNotions.Proposition5Numerics.SizeAt}
\leanref{Formalism/SharpRestriction/GrowthBounds/DiagonalSupport/ValuationOneEqZero.lean}{Formalism.MainNotions.Proposition5Numerics.digFacts}
\leanref{Formalism/SharpRestriction/GrowthBounds/DiagonalSupport/BinomCombFOnFinset.lean}{Formalism.MainNotions.Proposition5Numerics.gaAt}
\leanref{Formalism/SharpRestriction/GrowthBounds/TwistedSubstitution/JointEIntList.lean}{Formalism.MainNotions.TestSideWords.restJP}
\leanref{Formalism/SharpRestriction/GrowthBounds/DiagonalSupport/VBBMonomialOne.lean}{Formalism.MainNotions.Proposition5Numerics.noAt}
\leanref{Formalism/SharpRestriction/GrowthBounds/DiagonalSupport/SizeStep.lean}{Formalism.MainNotions.Proposition5Numerics.sizeStep}

\begin{proof}
\openstep{The five inputs are not yet written out: the Fock pairing of $G$ with an assembly in terms of its level images
and the factorization of $\mu!$ (1), the $p$-adic units $\gamma^{\mathrm{pl}}_\rho$ and $\gamma^{\mathrm{no}}_\rho$ and
their expansions (2), the choice of a joint isotypic projector that does not raise $v_{\mathrm F}$ (3), the normal test
approximating $\operatorname{diag}_{1/\gamma^{\mathrm{no}}_\rho}(c_Ye(w^\nu))$ (4), and the size bounds (5).}
% Lean: (1) digFacts (ValuationOneEqZero.lean, 275 lines, on LevImgCoeffBigM.lean and the files before it in
%   DiagonalSupport/); (2) gaAt (BinomCombFOnFinset.lean, 217 lines, with digAbs); (3) restJP
%   (TwistedSubstitution/JointEIntList.lean, 327 lines, from eeN and jointN_props); (4) noAt (VBBMonomialOne.lean, 109
%   lines, on LowestWeightCNonpos.lean and the files it uses); (5) sizeStep = sizeStep' (SizeStep.lean and
%   SizeStepAlt.lean, 175 lines, on LeCubeMul.lean and the counting files). The constants CDn and C1 are those of sizeStep.
% Size: about 3,700 lines together (the import closure of GramExponentArithmetic.lean in DiagonalSupport/ less
%   PickCore, PickDefs, PickLem and PickVal), shared in part with Lemmas pd-symbols-test-flat-chains-datum-membership
%   and -datum-numbers.
\end{proof}

\begin{lemma}\label{pd-symbols-test-flat-chains-datum-pick}
Let $C_{La},C_k,C_{rk},C_{\deg}$ be real numbers with the property of Theorem~\ref{pd-symbols-count}. There are
$C_{Dn},C_P>0$ such that the following holds for every $C\geq C_P$. Let the data and $S$ be as in
Lemma~\ref{pd-symbols-test-flat-chains-membership}, let $\kappa_F$ be as in Lemma~\ref{pd-symbols-test-flat-real}, and
put $\mathcal L=|T|k^2\log_pa$. Let $G\in\mathcal K[X]$ have multidegree $d_T$, lie in $\mathcal R[X]$ but not in
$\varpi_0\mathcal R[X]$, and satisfy the membership condition of Definition~\ref{pd-symbols-test-flat-chains-tied}. Then
there are a signed chain datum $\mathcal D$ of rank $r$ with $J$ levels over $\mathcal R$ that is adapted to $(J,f)$ with
some $d^{\mathrm{pl}}$ and $\kappa$, a block decomposition $(\rho,d^{\mathrm{pl}},d^{\mathrm{no}})$ of $d_T$ for $\mathcal D$, an
admissible $F_2$ with $v_{\mathrm F}(F_2)=0$ (Definition~\ref{pd-symbols-test-chains}), a normal test $N_0$ for $d^{\mathrm{no}}$
and $\kappa$, and $k_0\in\mathcal K$ with $v(k_0)\leq v(\prod_sd_T(s)!)-n_0+C_{Dn}\mathcal L$
(Definition~\ref{pd-symbols-test-flat-chains-datum-assembly}), such that
\[
 v\big(\langle G,\operatorname{asm}_\rho(A,N_0)\rangle_{\mathrm F}-k_0B_{\mathcal D}(F_2,A)\big)\geq c'+v\Big(\prod_sd_T(s)!\Big)-\kappa_F
\]
for every $A\in\mathcal K[x]$ with integral coefficients and degree $d^{\mathrm{pl}}(y)$ in the row of every position $y$.
\end{lemma}
\leanref{Formalism/SharpRestriction/GrowthBounds/DiagonalSupport/TwiceValuationFactorialProduct.lean}{Formalism.MainNotions.Proposition5Numerics.PickAt}
\leanref{Formalism/SharpRestriction/GrowthBounds/DiagonalSupport/TwiceValuationFactorialProduct.lean}{Formalism.MainNotions.Proposition5Numerics.PickStep}
\leanref{Formalism/SharpRestriction/GrowthBounds/DiagonalSupport/TwiceValuationFactorialProduct.lean}{Formalism.MainNotions.Proposition5Numerics.RestPick}
\leanref{Formalism/SharpRestriction/GrowthBounds/DiagonalSupport/GramExponentArithmetic.lean}{Formalism.MainNotions.Proposition5Numerics.restPick}

\begin{proof}
Let $C_{Dn}$ and $C_S$ be the constants of Lemma~\ref{pd-symbols-test-flat-chains-datum-pick-inputs}, and $C_P=C_S$. Let
$C\geq C_P$ and $G$ be as in the statement, and use the notation of
Definition~\ref{pd-symbols-test-flat-chains-datum-pick-weights}; put $P=\prod_sd_T(s)!$. We refer to the parts of
Lemma~\ref{pd-symbols-test-flat-chains-datum-pick-inputs} as (1) to (5).

\emph{Step 1: a monomial.} As $G\in\mathcal R[X]$ and $G\notin\varpi_0\mathcal R[X]$, some coefficient $G_{\mu_0}$ is a
unit. Let $\rho$, $d^{\mathrm{pl}}$, $d^{\mathrm{no}}$ be the residues and blocks of $X^{\mu_0}$, and $\beta_0$, $\beta'_0$
its plane and normal exponents. By (1), $(\rho,d^{\mathrm{pl}},d^{\mathrm{no}})$ is a block decomposition of $d_T$, the
coefficient of $x^{\beta_0}w^{\beta'_0}$ in $L_\rho=L_{\rho,d^{\mathrm{pl}},d^{\mathrm{no}}}(G)$ is $G_{\mu_0}$,
$\beta_0,\beta'_0$ have degrees $d^{\mathrm{pl}},d^{\mathrm{no}}$, and $v(\mu_0!)=v(c_\rho)+v(\beta_0!)+v(\beta'_0!)$, with
$c_\rho=c_\rho(d^{\mathrm{pl}},d^{\mathrm{no}})$, as the values of $\gamma^{\mathrm{pl}}_\rho,\gamma^{\mathrm{no}}_\rho$ have
valuation $0$ by (2). Let $y\in\mathcal K[w]$ be the coefficient of $x^{\beta_0}$ in $L_\rho$; it has degrees
$d^{\mathrm{no}}$ by (1), and $v_{\mathrm F}(y)\leq v(\beta'_0!)/2$.

\emph{Step 2: the datum.} Let $\mathcal D_0$ be the signed chain datum with $n_j=|\Lambda'_{\hat\jmath}|$, the letters
typed by the parity of their places, $d=d^{\mathrm{pl}}$, $g(j,m)=f(\hat\jmath+1,r-1)-f(\hat\jmath,r-1)$, $n^c_j=n_j+1$ and
the chains of Definition~\ref{pd-symbols-test-flat-chains-tied}(2), $s_{j,i+1}=\frac38s_j(i)$, $c=c'$ (which is at least
$1$ by (5)), the units $\gamma$ of (2) for $\rho$, and any positive thresholds. By (3) for $\mathcal D_0$ and $y$ there are
$\kappa$ and $e$ with the properties there and $v_{\mathrm F}(ey)\leq v_{\mathrm F}(y)$. Let $\mathcal D$ be $\mathcal D_0$
with $\theta_{j,i+1}=\theta^*_j(i)-\frac18s_j(i)$ and $\theta^L_{j,m}=\theta^*_j(u)-\frac38s_j(u)$, $u=n_j+r-1-m$, formed with
$\kappa$; these $\theta^L_{j,m}$ are positive, as $\theta^*_j(u)\geq B_{\hat\jmath}(x_j(u))/2>\frac38s_j(u)$ by (5). So
$\mathcal D$ is adapted to $(J,f)$ with $d^{\mathrm{pl}}$ and $\kappa$, and $e$ keeps the properties of (3) for $\mathcal
D$, as they involve only the types and the sets $N_j(u)$.

\emph{Step 3: $F_2$.} Let $q_1=v_{\mathrm F}(ey)\leq v(\beta'_0!)/2$, attained at a monomial $w^\nu$:
$q_1=q_2+v(\nu!)/2$ with $q_2=v((ey)_\nu)$. As $ey$ has degrees $d^{\mathrm{no}}$, so has $\nu$. Let
$P_c=[eL_\rho]_\nu$. Its coefficient at $x^{\beta_0}$ is $(ey)_\nu$, so $q_3=v_{\mathrm F}(P_c)\leq q_2+v(\beta_0!)/2$;
let $q_3$ be attained at $x^{m_3}$ and put $z=(P_c)_{m_3}\varpi_0^{v(m_3!)}$, so that $v(z)=q_3$ as $v(\varpi_0)=\frac12$.
Put $F_2=z^{-1}P_c$. Then $v_{\mathrm F}(F_2)=0$, and $F_2$ is admissible for $\mathcal D$ by the membership condition of
$G$.

\emph{Step 4: $N_0$ and $k_0$.} By (4) there are a normal test $N_0$ for $d^{\mathrm{no}}$ and $\kappa$ and $c_Y$ with
$v(c_Y)\leq v(\prod_yd^{\mathrm{no}}(y)!)/2+D_{\mathrm{no}}-v(\nu!)/2$, $D_{\mathrm{no}}=D_{\mathrm{no}}(d^{\mathrm{no}})$, and
$N_0-N_e$ with coefficients of valuation at least $c'$, where $N_e=\operatorname{diag}_{1/\gamma^{\mathrm{no}}_\rho}(c_Y
e(w^\nu))$. Put $k_0=c_\rho c_Y\nu!\,z$. As $X^{\mu_0}$ has multidegree $d_T$, $\mu_0!$ divides $P$, and
$n_0=v(\prod_yd^{\mathrm{pl}}(y)!)/2$. Hence
\[
 v(k_0)\leq v(c_\rho)+\frac{v(\prod_yd^{\mathrm{no}}(y)!)}2+D_{\mathrm{no}}+\frac{v(\nu!)}2+q_3
 \leq v(\mu_0!)+D_{\mathrm{no}}+\frac{v(\prod_yd^{\mathrm{no}}(y)!)-v(\beta'_0!)}2-\frac{v(\beta_0!)}2
 \leq v(P)-n_0+C_{Dn}\mathcal L,
\]
using $q_3\leq q_1-v(\nu!)/2+v(\beta_0!)/2\leq(v(\beta'_0!)-v(\nu!)+v(\beta_0!))/2$ and then (5).

\emph{Step 5: the congruence.} Let $A$ be integral of degrees $d^{\mathrm{pl}}$. By (1) with $N=N_e$, as
$\gamma^{\mathrm{no}}_\rho(\beta')(N_e)_{\beta'}=c_Y(e(w^\nu))_{\beta'}$ and $e$ is self-adjoint for the Fock pairing,
\[
 \langle G,\operatorname{asm}_\rho(A,N_e)\rangle_{\mathrm F}=c_\rho c_Y\sum_\beta\langle L_\beta,e(w^\nu)\rangle_{\mathrm F}
 \gamma^{\mathrm{pl}}_\rho(\beta)A_\beta\beta!=c_\rho c_Y\nu!\,\langle P_c,\operatorname{diag}_{\gamma^{\mathrm{pl}}_\rho}A\rangle_{
 \mathrm F}=k_0B_{\mathcal D}(F_2,A),
\]
with $L_\beta\in\mathcal K[w]$ the coefficient of $x^\beta$ in $L_\rho$. So $\langle G,\operatorname{asm}_\rho(A,N_0)\rangle_{
\mathrm F}-k_0B_{\mathcal D}(F_2,A)=\langle G,\operatorname{asm}_\rho(A,N_0-N_e)\rangle_{\mathrm F}$, and the assembly has
coefficients of valuation at least $c'$. As in Step~3 of the proof of Lemma~\ref{pd-symbols-test-flat-chains-datum}, this
pairing has valuation at least $c'+v(P)-\kappa_F$.
% Lean: restPick D CLa Ck Crk Cdeg (DiagonalSupport/GramExponentArithmetic.lean) = pickStep_all (the constants of
%   sizeStep) and pickAt_of: Step 1 is exists_unit_coeff and the third clause of digFacts (classOf, blkPl, blkNo, plMon,
%   noMon); Step 2 is mkDt, hJP = restJP, updTh with thF and thLF, thL_arith (with the clause hB of sizeStep) and
%   jpok_upd. Steps 3 to 5 are pick_core (PickCore.lean): vBB_le_coeff, vBB_attained, vK_w_pow and vBB_smul for Step 3,
%   hXX from MemAt; noAt and gram_exponent_arithmetic (with padicVal_expFact_le and hfact) for Step 4; value_exact
%   (PickVal.lean), asm_add_right and geQ_fock with asm_coefGe' for Step 5. What remains open is Lemma
%   pd-symbols-test-flat-chains-datum-pick-inputs; PickCore, PickDefs, PickLem and PickVal are about 550 lines.
\end{proof}

\begin{definition}\label{pd-symbols-test-flat-chains-datum-membership-levels}
Let the notation be as in Definitions~\ref{pd-symbols-test-chains}, \ref{pd-symbols-test-flat-chains-tied}
and~\ref{pd-symbols-test-flat-chains-datum-assembly}, and let $\mathcal D$ be a signed chain datum of rank $r$ with $J$
levels over $\mathcal R$ with $n_j=|\Lambda'_{\hat\jmath}|$ for every level $j$. Let $\mathcal R[z^{\mathrm{pl}}]$ and
$\mathcal R[z^{\mathrm{no}}]$ be the polynomial rings in the variables $z^{\mathrm{pl}}_{(j,x),a}$, $a<r$, and
$z^{\mathrm{no}}_{(j,x),b}$, $1\leq b\leq K$, for $j\in\mathbf Z/J\mathbf Z$ and $x\in\mathbf Z$.
\begin{enumerate}
\item For $Z\in\mathcal R[\tilde x]$ let $Z^{\mathrm{lv}}\in\mathcal R[z^{\mathrm{pl}}]$ be obtained by renaming
  $\tilde x_{\tilde y,a}$ to $z^{\mathrm{pl}}_{(j,x),a}$, where $(j,x)$ is: for the letter $(j,l)$, its level and its place;
  for the label and the top of the link $(j,m)$, $(j,f(\hat\jmath,r-1-m))$; for the lift of $(j,m)$,
  $(j+1,f(\hat\jmath'-1,r-1-m))$ with $\hat\jmath'$ the representative of $j+1$. For $N\in\mathcal R[w]$ let
  $N^{\mathrm{lv}}\in\mathcal R[z^{\mathrm{no}}]$ be obtained by renaming $w_{y,b}$ to $z^{\mathrm{no}}_{(j,x),b}$, with $j$
  the level and $x$ the place of $y$.
\item $\Phi^{\mathrm{pl}}$ and $\Phi^{\mathrm{no}}$ are the $\mathcal R$-algebra maps to $\mathcal R[X]$ with
  $\Phi^{\mathrm{pl}}(z^{\mathrm{pl}}_{(j,x),a})=X_{\bar x,a+1}^{\,p^{\operatorname{dist}_{\hat\jmath}(x)}}$ and
  $\Phi^{\mathrm{no}}(z^{\mathrm{no}}_{(j,x),b})=X_{\bar x,r+b}^{\,p^{\operatorname{dist}_{\hat\jmath}(x)}}$.
\item For $\rho$ from the pairs of a position and an integer $c\leq k$ to $\mathbf Z_{\geq0}$, the \emph{label of $\rho$}
  is $X^{\ell_\rho}=\prod_y\prod_{c\leq k}X_{\bar y,c}^{\rho(y,c)}$.
\item For $\delta$ from $(\mathbf Z/J\mathbf Z)\times\mathbf Z$ to $\mathbf Z_{\geq0}$ and $H\in\mathcal R[z^{\mathrm{pl}}]$,
  $H_\delta$ is the sum of the terms of $H$ whose monomials have degree $\delta(j,x)$ in the row $(j,x)$ for every
  $(j,x)$; $\delta$ \emph{occurs} in $H$ if some monomial of $H$ has these degrees.
\end{enumerate}
\end{definition}
\leanref{Formalism/SharpRestriction/GrowthBounds/DiagonalSupport/ExistsMapOfCoefGe.lean}{Formalism.MainNotions.Proposition5Numerics.mlPos}
\leanref{Formalism/SharpRestriction/GrowthBounds/DiagonalSupport/ExistsMapOfCoefGe.lean}{Formalism.MainNotions.Proposition5Numerics.tlPos}
\leanref{Formalism/SharpRestriction/GrowthBounds/DiagonalSupport/ExistsMapOfCoefGe.lean}{Formalism.MainNotions.Proposition5Numerics.xiPl}
\leanref{Formalism/SharpRestriction/GrowthBounds/DiagonalSupport/ExistsMapOfCoefGe.lean}{Formalism.MainNotions.Proposition5Numerics.xiNo}
\leanref{Formalism/SharpRestriction/GrowthBounds/DiagonalSupport/ExistsMapOfCoefGe.lean}{Formalism.MainNotions.Proposition5Numerics.PhiPl}
\leanref{Formalism/SharpRestriction/GrowthBounds/DiagonalSupport/ExistsMapOfCoefGe.lean}{Formalism.MainNotions.Proposition5Numerics.PhiNo}
\leanref{Formalism/SharpRestriction/GrowthBounds/DiagonalSupport/ExistsMapOfCoefGe.lean}{Formalism.MainNotions.Proposition5Numerics.labO}
\leanref{Formalism/SharpRestriction/GrowthBounds/DiagonalSupport/MembDefs.lean}{Formalism.MainNotions.Proposition5Numerics.homComp}
\leanref{Formalism/SharpRestriction/GrowthBounds/DiagonalSupport/ExistsMapOfCoefGe.lean}{Formalism.MainNotions.Proposition5Numerics.Occ}

\begin{lemma}\label{pd-symbols-test-flat-chains-datum-membership-parts}
Let $p$ be a prime, $\mathcal R$ a complete discrete valuation ring with fraction field $\mathcal K$ in which $p$ is
neither $0$ nor a unit, $n'\geq1$, $r\geq2$, $K\geq1$, $k=r+K$, $L_1,L_2,N_1,G'\geq0$ integers, $T\subset\mathbf Z/n'\mathbf Z$, $(J,f)$ a run
system for $T$ of length $r$ spaced for $(L_1,L_2,N_1,G')$, $\deg,d\colon\mathbf Z/n'\mathbf Z\to\mathbf Z_{\geq0}$, $Y$ a
real number and $c'\geq0$ an integer, with $d(s)=\deg(s)$ for $s\in T$, $d(s)=0$ for $s\notin T$, and
$128c'\leq\lfloor B_j(x)/(Yk^2)\rfloor$ for every $j$ and every new $x\in\Lambda_j$. Use
Definition~\ref{pd-symbols-test-flat-chains-datum-membership-levels}. Let $\mathcal D$ be a signed chain datum of rank $r$
with $J$ levels over $\mathcal R$ that is adapted to $(J,f)$ with $d^{\mathrm{pl}}$ and $\kappa$, $(\rho,d^{\mathrm{pl}},
d^{\mathrm{no}})$ a block decomposition of $d$ for $\mathcal D$, $N_o\in\mathcal R[w]$ whose image in $\mathcal K[w]$ is a
normal test for $d^{\mathrm{no}}$ and $\kappa$, and $Z_o\in\mathcal R[\tilde x]$ whose image $Z$ in $\mathcal K[\tilde x]$
satisfies the two bounds of Lemma~\ref{pd-symbols-test-flat-chains-datum-membership} and such that $\varphi$ of every
monomial of $Z_o$ has degree $d(y)$ in the row of every position $y$. Then:
\begin{enumerate}
\item $X^{\ell_\rho}\Phi^{\mathrm{pl}}(Z_o^{\mathrm{lv}})\Phi^{\mathrm{no}}(N_o^{\mathrm{lv}})=\operatorname{asm}_\rho(\varphi(Z),
  N_o)$ in $\mathcal K[X]$;
\item $X^{\ell_\rho}$ is a label, with $\ell_\rho(s,c)=0$ for $s\notin T$;
\item for every $\delta$ that occurs in $Z_o^{\mathrm{lv}}$, $X^{\ell_\rho}\Phi^{\mathrm{pl}}(Z^{\mathrm{lv}}_{o,\delta})
  \Phi^{\mathrm{no}}(N_o^{\mathrm{lv}})$ has multidegree $d$ and all its monomials have the same normal degree;
\item for every such $\delta$ there are $\xi_{i,j}\in\mathcal R[z]$, $1\leq i\leq n$, $0\leq j<J$, with
  $\Phi^{\mathrm{pl}}(Z^{\mathrm{lv}}_{o,\delta})\Phi^{\mathrm{no}}(N_o^{\mathrm{lv}})=\sum_i\prod_{j<J}\phi_j(\xi_{i,j})$, such
  that each $\xi_{i,j}$ and the functions $o_j$ and $b_j$ satisfy the conditions of
  Definition~\ref{pd-symbols-test-one-sided} on the $\xi_{i,j}$ of a one-sided element.
\end{enumerate}
\end{lemma}
\leanref{Formalism/SharpRestriction/GrowthBounds/DiagonalSupport/ExistsMapOfCoefGe.lean}{Formalism.MainNotions.Proposition5Numerics.MFH}
\leanref{Formalism/SharpRestriction/GrowthBounds/DiagonalSupport/ExistsMapOfCoefGe.lean}{Formalism.MainNotions.Proposition5Numerics.MH}
\leanref{Formalism/SharpRestriction/GrowthBounds/DiagonalSupport/MembTop.lean}{Formalism.MainNotions.Proposition5Numerics.RestAlgG}
\leanref{Formalism/SharpRestriction/GrowthBounds/DiagonalSupport/MembTop.lean}{Formalism.MainNotions.Proposition5Numerics.RestLabG}
\leanref{Formalism/SharpRestriction/GrowthBounds/DiagonalSupport/MembTop.lean}{Formalism.MainNotions.Proposition5Numerics.RestDegG}
\leanref{Formalism/SharpRestriction/GrowthBounds/DiagonalSupport/MembTop.lean}{Formalism.MainNotions.Proposition5Numerics.RestDecG}
\leanref{Formalism/SharpRestriction/GrowthBounds/DiagonalSupport/MembAlgF.lean}{Formalism.MainNotions.Proposition5Numerics.restAlgG}
\leanref{Formalism/SharpRestriction/GrowthBounds/DiagonalSupport/MembLabF.lean}{Formalism.MainNotions.Proposition5Numerics.restLabG}
\leanref{Formalism/SharpRestriction/GrowthBounds/DiagonalSupport/MembDegF.lean}{Formalism.MainNotions.Proposition5Numerics.restDegG}
\leanref{Formalism/SharpRestriction/GrowthBounds/DiagonalSupport/MembershipOfThreeParts.lean}{Formalism.MainNotions.Proposition5Numerics.decomposition_of_three_parts}

\begin{proof}
\openstep{The four parts are not yet written out: the algebraic identity between the assembly and the product of the
level maps (1), the label of $\rho$ (2), the multidegree and normal degree of the components (3), and the decomposition
of each component into products of elements of the segment modules $X_j(m_*,\theta^1_j,o_j)$ with the degree bounds,
from the bounds of $Z$ at the chains and at the links (4).}
% Lean: (1) restAlgG (MembAlgF.lean, 187 lines: plane_hom, normal_hom, label_hom, alg_identity); (2) restLabG
%   (MembLabF.lean, 153 lines); (3) restDegG (MembDegF.lean, 147 lines); (4) decomposition_of_three_parts restValBdG
%   restXXG_b restDdG_holds = restDecG_of atl valSlice valComp (MembershipOfThreeParts.lean), through DecIdentity,
%   ProjectionMemPartialSpan, DecompositionOfThreeParts, MembVB (263 lines), MembXXb (143) and MembDd (431), with
%   ThresholdGeOddPosition, PreChainLe, RestValBdG and the atl, valSlice and valComp of ImportsAlt.
% Size: about 6,700 lines (the import closure of MembFinal.lean in DiagonalSupport/ less MembTop.lean and the parts of
%   ExistsMapOfCoefGe.lean used for the top level), with (4) the largest part.
\end{proof}

\begin{lemma}\label{pd-symbols-test-flat-chains-datum-membership}
Let $C_{La},C_k,C_{rk},C_{\deg}$ be real numbers with the property of Theorem~\ref{pd-symbols-count}. There is $C_B>0$
such that the following holds for every $C\geq C_B$. Let the data be as in
Lemma~\ref{pd-symbols-test-flat-chains-membership}, let $\mathcal D$ be a signed chain datum of rank $r$ with $J$ levels
over $\mathcal R$ that is adapted to $(J,f)$ with $d^{\mathrm{pl}}$ and $\kappa$, $(\rho,d^{\mathrm{pl}},d^{\mathrm{no}})$ a
block decomposition of $d_T$ for $\mathcal D$, and $N_0$ a normal test for $d^{\mathrm{no}}$ and $\kappa$
(Definition~\ref{pd-symbols-test-flat-chains-datum-assembly}). Let $Z\in\mathcal K[\tilde x]$ have integral
coefficients, with $\operatorname{val}_{\tilde C_{j,i}}(Z)\geq\theta_{j,i}-s_{j,i}-24c$ for every chain and
$\operatorname{val}_{\tilde S_{j,m}}(Z)\geq\theta^L_{j,m}-60c$ for all $j$ and $m+2\leq r$, and such that $\varphi$ of every
monomial of $Z$ has degree $d(y)$ in the row of every position $y$ (Definition~\ref{pd-symbols-test-chains}). Then there
are one-sided elements $\sigma_1,\dots,\sigma_{n_q}$ of multidegree $d_T$ (Definition~\ref{pd-symbols-test-one-sided},
for $(J,f)$, $N_1$, $G'$, $c'$, $\deg$ and $Y$) with $\sum_q\sigma_q=\operatorname{asm}_\rho(\varphi(Z),N_0)$.
\end{lemma}
\leanref{Formalism/SharpRestriction/GrowthBounds/DiagonalSupport/TwiceValuationFactorialProduct.lean}{Formalism.MainNotions.Proposition5Numerics.MembAt}
\leanref{Formalism/SharpRestriction/GrowthBounds/DiagonalSupport/TwiceValuationFactorialProduct.lean}{Formalism.MainNotions.Proposition5Numerics.MembStep}
\leanref{Formalism/SharpRestriction/GrowthBounds/DiagonalSupport/TwiceValuationFactorialProduct.lean}{Formalism.MainNotions.Proposition5Numerics.RestMemb}
\leanref{Formalism/SharpRestriction/GrowthBounds/DiagonalSupport/MembFinal.lean}{Formalism.MainNotions.Proposition5Numerics.restMemb}

\begin{proof}
Put $C_B=465$. Let $C\geq C_B$ and the data be as in the statement, and use
Definition~\ref{pd-symbols-test-flat-chains-datum-membership-levels}.

\emph{Step 1: the hypotheses of Lemma~\ref{pd-symbols-test-flat-chains-datum-membership-parts}.} The prime $p$, $r\geq2$,
the spaced run system and $d_T$ are as there, with $d=d_T$. Let $x\in\Lambda_j$ be new. As in Step~5 of the proof of
Lemma~\ref{pd-symbols-test-flat-chains-datum-numbers}, $B_j(x)\geq a/p^{f(j,r-1)-f(j-1,0)}\geq a/p^{G'}$, so
Lemma~\ref{pd-symbols-test-flat-chains-datum-numbers-slack}(1), with $\alpha=0$, gives
$100r^2c'\leq\lfloor B_j(x)/(Yk^2)\rfloor$, as the other terms on its left are nonnegative; as $r\geq2$,
$128c'\leq400c'\leq100r^2c'$.

\emph{Step 2: integral models.} As $Z$ and $N_0$ have integral coefficients, they are the images of some
$Z_o\in\mathcal R[\tilde x]$ and $N_o\in\mathcal R[w]$, with the same monomials. So
Lemma~\ref{pd-symbols-test-flat-chains-datum-membership-parts} applies to $\mathcal D$, $(\rho,d^{\mathrm{pl}},
d^{\mathrm{no}})$, $N_o$ and $Z_o$.

\emph{Step 3: the decomposition.} Let $\delta_1,\dots,\delta_{n_q}$ be the $\delta$ that occur in $Z_o^{\mathrm{lv}}$, so
that $Z_o^{\mathrm{lv}}=\sum_qZ^{\mathrm{lv}}_{o,\delta_q}$, and let $\sigma_q$ be the image in $\mathcal K[X]$ of
$X^{\ell_\rho}\Phi^{\mathrm{pl}}(Z^{\mathrm{lv}}_{o,\delta_q})\Phi^{\mathrm{no}}(N_o^{\mathrm{lv}})$. By parts (2) to (4) of
Lemma~\ref{pd-symbols-test-flat-chains-datum-membership-parts}, $\sigma_q=X^{\ell_\rho}\sum_i\prod_{j<J}\phi_j(\xi_{i,j})$
is a one-sided element of multidegree $d_T$, of normal degree the common normal degree $o_q$ of its monomials. As
$\Phi^{\mathrm{pl}}$ is additive, part (1) gives $\sum_q\sigma_q=\operatorname{asm}_\rho(\varphi(Z),N_0)$.
% Lean: restMemb = membership_of_three_parts (MembershipOfThreeParts.lean) = membership_of_four_parts (MembTop.lean):
%   C1 = max C0 1 with C0 the constant of paramsII0 1 1, which is 465. Step 1 is mfh_of (Bo_ge, counting_hypotheses_of_hyps
%   with paramsII0 and margin_real); Step 2 is membAt_of (exists_map_of_coefGe); Step 3 is membAt_core
%   (ExistsMapOfCoefGe.lean), with the set S of the multidegrees of the monomials of xiPl Go, homComp and the four parts.
%   What remains open is Lemma pd-symbols-test-flat-chains-datum-membership-parts.
\end{proof}

\begin{lemma}\label{pd-symbols-test-flat-chains-datum-numbers-places}
Let $p$ and $\mathcal R$ be as in Definition~\ref{pd-symbols-test-chains}, $n',r\geq1$ and $L_1,N_1\geq0$ integers,
$T\subset\mathbf Z/n'\mathbf Z$, and $(J,f)$ a run system for $T$ of length $r$ valid for $(L_1,N_1)$
(Definition~\ref{pd-symbols-run-system}). Let $\mathcal D$ be a signed chain datum of rank $r$ with $J$ levels over
$\mathcal R$ with $n_j=|\Lambda'_{\hat\jmath}|$ for every level $j$, and let the places be as in
Definition~\ref{pd-symbols-test-flat-chains-tied}(1). Then:
\begin{enumerate}
\item the place $x$ of a position of level $j$ satisfies $f(\hat\jmath-1,r-1)<x\leq f(\hat\jmath,r-1)$ and $\bar x\in T$;
\item the places of distinct positions have distinct residues modulo $n'$, so $\mathcal D$ has at most $|T|$
  positions;
\item for every level $j$ and $0\leq u\leq n_j$, the place $x_j(u)$ is a new element of $\Lambda_{\hat\jmath}$.
\end{enumerate}
\end{lemma}
\leanref{Formalism/SharpRestriction/GrowthBounds/DiagonalSupport/SlackBoundReal.lean}{Formalism.MainNotions.Proposition5Numerics.posTL_seg}
\leanref{Formalism/SharpRestriction/GrowthBounds/DiagonalSupport/SlackBoundReal.lean}{Formalism.MainNotions.Proposition5Numerics.posTL_inj}
\leanref{Formalism/SharpRestriction/GrowthBounds/DiagonalSupport/SlackBoundReal.lean}{Formalism.MainNotions.Proposition5Numerics.card_TL_le}
\leanref{Formalism/SharpRestriction/GrowthBounds/DiagonalSupport/SlackBoundJoint.lean}{Formalism.MainNotions.Proposition5Numerics.posU_mem}

\begin{proof}
(1) Let $y$ be a position of level $j$ with place $x$. If $y$ is the letter $(j,l)$, then $l<n_j$, so $x$ is an element
of $\Lambda'_{\hat\jmath}$: $\bar x\in T$ and $f(\hat\jmath-1,r-1)<x<f(\hat\jmath,0)\leq f(\hat\jmath,r-1)$, by (2) of
validity. If $y$ is the link $(j,m)$, its index is $n_j+r-1-m$, so $x=f(\hat\jmath,r-1-m)$; then $\bar x\in T$ by (3) of
validity, and $f(\hat\jmath-1,r-1)<f(\hat\jmath,0)\leq x\leq f(\hat\jmath,r-1)$ by (2).

(2) By (2) of validity $f(i,r-1)$ increases with $i$, and by (1) of validity $f(J-1,r-1)=f(-1,r-1)+n'$. So by (1) the
places of the positions of the levels $j$, $0\leq\hat\jmath<J$, lie in the disjoint intervals
$(f(\hat\jmath-1,r-1),f(\hat\jmath,r-1)]$, all contained in $(f(-1,r-1),f(-1,r-1)+n']$. Within one level, the letters
have distinct places, the elements of $\Lambda'_{\hat\jmath}$ in increasing order; the links have the distinct places
$f(\hat\jmath,i)$, by (2) of validity; and the place of a letter is smaller than $f(\hat\jmath,0)$, so smaller than that of
a link. So distinct positions have distinct places in an interval of $n'$ integers, and their residues modulo $n'$ are
distinct. By (1) these residues lie in $T$.

(3) If $u<n_j$, then $x_j(u)\in\Lambda'_{\hat\jmath}$, and by (1) $f(\hat\jmath-1,0)\leq f(\hat\jmath-1,r-1)<x_j(u)\leq
f(\hat\jmath,r-1)$; so $x_j(u)$ is a new element of $\Lambda_{\hat\jmath}$. If $u=n_j$, then $x_j(u)=f(\hat\jmath,0)$, which
lies in $\Lambda_{\hat\jmath}$ and is new by (2) and (3) of validity.
% Lean: (1) is posTL_seg (llPos_mem, mem_LLset, fr_mono_of_valid); (2) is posTL_inj (posTL_win for the window of
%   length n', then Int.eq_zero_of_abs_lt_dvd) and card_TL_le; (3) is posU_mem (LL_sub_letters for u < nL, and
%   fr j 0 for u = nL).
\end{proof}

\begin{lemma}\label{pd-symbols-test-flat-chains-datum-numbers-letters}
Let $p\geq3$ be a prime, $\mathcal R$ as in Definition~\ref{pd-symbols-test-chains}, $n',r\geq1$ and
$L_1\leq L_2$, $N_1$, $G'\geq0$ integers, $T\subset\mathbf Z/n'\mathbf Z$, and $(J,f)$ a run system for $T$ of length $r$
spaced for $(L_1,L_2,N_1,G')$ (Definition~\ref{pd-symbols-run-system}). Let $\mathcal D$ be a signed chain datum of rank
$r$ with $J$ levels over $\mathcal R$ with $n_j=|\Lambda'_{\hat\jmath}|$ for every level $j$, let $\lambda(y)$ be as in
Definition~\ref{pd-symbols-test-flat-chains-tied}(1), and let $a\geq0$ be a real number with $\lambda(y)d(y)\leq2a$ for
every position $y$. Then for every level $j$ the sum $d_L(j)$ of the degrees of the letters of level $j$ satisfies
\[
 d_L(j)\leq3a\,p^{-D_j-L_1-1},\qquad D_j=f(\hat\jmath,r-1)-f(\hat\jmath,0).
\]
\end{lemma}
\leanref{Formalism/SharpRestriction/GrowthBounds/DiagonalSupport/SlackBoundJoint.lean}{Formalism.MainNotions.Proposition5Numerics.LL_gap}
\leanref{Formalism/SharpRestriction/GrowthBounds/DiagonalSupport/SlackBoundJoint.lean}{Formalism.MainNotions.Proposition5Numerics.dLL_le}

\begin{proof}
Let $y$ be a letter of level $j$, with place $x\in\Lambda'_{\hat\jmath}$. As $f(\hat\jmath,0)$ is $L_2$-isolated in
$\Lambda_{\hat\jmath}$ and $x$ is an element of $\Lambda_{\hat\jmath}$ smaller than $f(\hat\jmath,0)$, the largest such
element $x'$ satisfies $x\leq x'<f(\hat\jmath,0)-L_2$. Put $e(y)=f(\hat\jmath,r-1)-x=\operatorname{dist}_{\hat\jmath}(x)$;
then $\lambda(y)=p^{e(y)}$ and $e(y)\geq D_j+L_2+1$. Distinct letters of level $j$ have distinct places, so distinct
$e(y)$. Put $q=p^{-1}\leq\frac13$. Then $d(y)\leq2aq^{e(y)}$, and
\[
 d_L(j)\leq2a\sum_{e\geq D_j+L_2+1}q^e=\frac{2aq^{D_j+L_2+1}}{1-q}\leq3aq^{D_j+L_1+1},
\]
as $1/(1-q)\leq\frac32$ and $L_1\leq L_2$.
% Lean: dLL_le: LL_gap (the largest letter below fr j 0 is L2-isolated, from the spacing hF.2.2.2.2), the injectivity
%   llPos_inj, hterm (d <= 2 a q^e), then geom_sum_Ico_le_of_lt_one and hfrac.
\end{proof}

\begin{lemma}\label{pd-symbols-test-flat-chains-datum-numbers-sizes}
Let $p$ and $\mathcal R$ be as in Definition~\ref{pd-symbols-test-chains}, and $\mathcal D$ a signed chain datum of rank
$r$ with $J$ levels over $\mathcal R$. Let $N,k\geq0$ be integers and $a$ a real number with $r\leq k$, $a\geq2$,
$p\leq a$, $4N\leq a$ and $r\leq a$. Suppose that $\mathcal D$ has at most $N$ positions, that $n^c_j=n_j+1$ for every
level $j$, and that $d(y)\leq2a$ for every position $y$. Then $\kappa_{\mathcal D}\geq0$ and
\[
 D_m+D_l+3\kappa_{\mathcal D}\leq20Nk^2\log_pa .
\]
\end{lemma}
\leanref{Formalism/SharpRestriction/GrowthBounds/DiagonalSupport/SlackBoundReal.lean}{Formalism.MainNotions.Proposition5Numerics.sizes_le}
\leanref{Formalism/SharpRestriction/GrowthBounds/DiagonalSupport/SlackBoundReal.lean}{Formalism.MainNotions.Proposition5Numerics.lg_le_four}

\begin{proof}
Put $\ell=\log_pa$; then $\ell\geq1$, as $p\leq a$. If $n\geq0$ is an integer with $n+r\leq a^3$, then
$0\leq\log_p(n+r)+1\leq3\ell+\ell=4\ell$. As $a\geq2$, $4a\leq a^3$ and $a^2+a\leq a^3$. For a position $y$,
$d(y)+r\leq3a\leq a^3$, so each term $r(\log_p(d(y)+r)+1)$ of $\kappa_{\mathcal D}$ lies between $0$ and $4k\ell\leq4k^2\ell$;
hence $\kappa_{\mathcal D}\geq0$ and $\kappa_{\mathcal D}\leq4Nk^2\ell$. Next, $2\sum_yd(y)+r\leq4Na+a\leq a^2+a\leq a^3$,
so $0\leq e_P\leq4r^2\ell\leq4k^2\ell$. There are $\sum_j(n_j+r)$ positions, so $\sum_jn^c_j=\sum_j(n_j+1)\leq N$ and
$J(r-1)\leq\sum_j(n_j+r)\leq N$. Hence $D_m\leq4Nk^2\ell$, $D_l\leq4Nk^2\ell$ and $3\kappa_{\mathcal D}\leq12Nk^2\ell$.
% Lean: sizes_le: lg_le_four (lg p r n <= 4 log_p a for n + r <= a^3) and lg_nonneg; hkp0 and hkp for kappa, hlgS and
%   heP for e_P, card_TL_eq for the number of positions, hnch and hJr.
\end{proof}

\begin{lemma}\label{pd-symbols-test-flat-chains-datum-numbers-slack}
Let the data of Situation~\ref{pd-edge-situation} be given with $C\geq465$, and put $c'=c_*+\omega$, $q=p^{G'}$ and
$s(B)=\lfloor B/(Yk^2)\rfloor$ for real $B$. Then:
\begin{enumerate}
\item for every real $B\geq a/q$ and every real $\alpha$ with $0\leq\alpha\leq1+rkq$,
  \[
   100r^2c'+24(N_1+r)c_0+m_\alpha\leq s(B)\leq B/100,\qquad m_\alpha=(N_1+1)\big((2q+3)c'+kq+c'\alpha+3c'\big);
  \]
\item $rd/Q\leq3rp^{-L_1-1}B\leq s(B)/8$ for every integer $D\geq0$ and all reals $Q,d,B$ with $p^D\leq q$,
  $1\leq Q\leq q$, $d\leq3ap^{-D-L_1-1}$ and $B\geq a/(p^DQ)$.
\end{enumerate}
\end{lemma}
\leanref{Formalism/PolynomialDimension/MultinomialDepth/SecondParameterSizeBoundsHolds.lean}{Formalism.CountingSumsFrames.secondParameterSizeBounds_holds}
\leanref{Formalism/PolynomialDimension/MultinomialDepth/SecondParameterSlackBoundsHolds.lean}{Formalism.CountingSumsFrames.secondParameterSlackBounds_holds}
\leanref{Formalism/SharpRestriction/GrowthBounds/ScaleEdges/GainBoundAt.lean}{Formalism.MainNotions.Pf14.counting_hypotheses_of_hyps}

\begin{proof}
This is a self-contained computation. As in the proof of Lemma~\ref{pd-symbols-space}, $|T|\geq1$, $k\geq7$, $Y\geq r\geq2$
and $c'\geq c_*\geq1$; and $q\geq1$. Put $M=k^3r^2|T|q^5c'Y$; the last condition of (3) of the situation gives
$465M\leq a$.

(1) Let $u$ be as in Definition~\ref{hf-good-orbit}, formed with $\delta$. Then $u\delta\geq1$ and $G'\geq(6r+1)u$, so
$N_1>6eG'\delta\geq r$. Hence $k\geq N_1+3r\geq4r\geq8$ by (3) of the situation, and $r^2\leq N_1^2\leq r|T|$ gives
$r\leq|T|$ and $N_1+1\leq N_1^2\leq2r|T|$. As $r|T|\geq1$, $c_0\leq\omega\leq c'$. Let $0\leq\alpha\leq1+rkq$, and put
$S=(2q+3)c'+kq+c'\alpha+3c'$, so that $m_\alpha=(N_1+1)S$. As $rk\geq16$ and $c',q\geq1$,
$S\leq(2q+7)c'+kq+rkqc'\leq3rkc'q$. Put $W=Yk^2q$. Then:
\begin{itemize}
\item $r^2c'W\leq M$, as $k|T|q^4\geq1$;
\item $(N_1+r)c_0W\leq kc'W\leq M$, as $N_1+r\leq k$, $c_0\leq c'$ and $r^2|T|q^4\geq1$;
\item $(N_1+1)SW\leq2r|T|\cdot3rkc'q\cdot W=6k^3r^2|T|q^2c'Y\leq6M$;
\item $W\leq M$.
\end{itemize}
So, with $\Sigma=100r^2c'+24(N_1+r)c_0+m_\alpha$, $(\Sigma+1)W\leq131M\leq a$. For $B\geq a/q$ this gives
$B/(Yk^2)\geq\Sigma+1$, so $s(B)>B/(Yk^2)-1\geq\Sigma$. As $Y\geq2$ and $k\geq8$, $Yk^2\geq100$; as $B\geq0$,
$s(B)\leq B/(Yk^2)\leq B/100$.

(2) As $r^2|T|q^3\geq1$ and $kc'\geq1$, $2Yk^2q^2\leq100kc'q^2\cdot Yk^2\leq100M\leq a$. As $p^DQ\leq q^2$,
$B\geq a/(p^DQ)\geq2Yk^2$, so $s(B)>B/(Yk^2)-1\geq B/(2Yk^2)$. By (3) of the situation, $r\leq Y$ and $|T|\geq1$,
$48rYk^2\leq256k^2|T|Y^2\leq p^{L_1-1}\leq p^{L_1+1}$. Hence $3rp^{-L_1-1}B\leq B/(16Yk^2)\leq s(B)/8$. Finally
$rd/Q\leq3rp^{-L_1-1}\cdot a/(p^DQ)\leq3rp^{-L_1-1}B$.
% Lean: counting_hypotheses_of_hyps turns Hyps C into HypII C with N = |T|, N1 the orbit window count, c'' = cpp
%   (so cPrII is c') and J = 1. (1) is the second clause of ParamsII0 1 1 = secondParameterSizeBounds_holds, whose
%   constant is 64 + 400 + 1 = 465, with al = alpha: rank_le_orbitWindowCount (r <= N1), pII0_basic, pII0_om (c0 <= c'),
%   sizeBound_multiplicity_le (S <= 3 r k c' q), the four terms sizeBound_first_term to sizeBound_fourth_term of
%   pII0_slack (which states the bound as 155 M), and pII0_floor. (2) is the second clause of
%   ParamsII5 1 1 = secondParameterSlackBounds_holds, whose constant is 100: hdivX, slack_ge_half, h48r and
%   slackBound_coeff.
\end{proof}

\begin{lemma}\label{pd-symbols-test-flat-chains-datum-numbers}
Let $C_{La},C_k,C_{rk},C_{\deg}$ be real numbers with the property of Theorem~\ref{pd-symbols-count}. There are
$C_{Dm},C_N>0$ such that the following holds for every $C\geq C_N$. Let the data be as in
Lemma~\ref{pd-symbols-test-flat-chains-membership}, let $\kappa_F$ be as in Lemma~\ref{pd-symbols-test-flat-real}, and
put $\mathcal L=|T|k^2\log_pa$. Then $1\leq\mathcal L$ and $\kappa_F\leq C_{Dm}\mathcal L$, and for
every signed chain datum $\mathcal D$ of rank $r$ with $J$ levels over $\mathcal R$ that is adapted to $(J,f)$ with some
$d^{\mathrm{pl}}$ and $\kappa$, and every block decomposition $(\rho,d^{\mathrm{pl}},d^{\mathrm{no}})$ of $d_T$ for
$\mathcal D$ (Definition~\ref{pd-symbols-test-flat-chains-datum-assembly}): $\mathcal D$ satisfies (H) for $C_H=26$,
$4(D_m+D_l+3\kappa_{\mathcal D})<c$, $\kappa_{\mathcal D}\geq0$ and $D_m+D_l+3\kappa_{\mathcal D}\leq C_{Dm}\mathcal L$
(Definition~\ref{pd-symbols-test-chains}).
\end{lemma}
\leanref{Formalism/SharpRestriction/GrowthBounds/DiagonalSupport/TwiceValuationFactorialProduct.lean}{Formalism.MainNotions.Proposition5Numerics.NumAt}
\leanref{Formalism/SharpRestriction/GrowthBounds/DiagonalSupport/TwiceValuationFactorialProduct.lean}{Formalism.MainNotions.Proposition5Numerics.NumStep}
\leanref{Formalism/SharpRestriction/GrowthBounds/DiagonalSupport/TwiceValuationFactorialProduct.lean}{Formalism.MainNotions.Proposition5Numerics.RestNum}
\leanref{Formalism/SharpRestriction/GrowthBounds/DiagonalSupport/LeOfOneLeMul.lean}{Formalism.MainNotions.Proposition5Numerics.restNum}

\begin{proof}
Put $C_{Dm}=20$ and $C_N=465$. Let $C\geq C_N$ and the data be as in the statement, and put $q=p^{G'}$ and
$\ell=\log_pa$. Recall that $c'=c_*+\omega$, that $(J,f)$ is spaced for $(L_1,L_2,N_1,G')$, and that $d_T(s)=\deg(s)$ for
$s\in T$ and $d_T(s)=0$ otherwise.

\emph{Step 1: the size of $a$.} As in the proof of Lemma~\ref{pd-symbols-space}, $|T|\geq1$, $k\geq7$, $Y\geq r\geq2$,
$c'\geq1$, and $a\leq\deg(t)\leq2a$ for $t\in T$. By (2) of validity $f(j-1,0)<f(j,r-1)$, and the spacing gives
$f(j,r-1)-f(j-1,0)<G'$; so $G'\geq1$. As all the factors of the left side of the last condition of (3) of the
situation are at least $1$, it gives $4|T|\leq k^3|T|\leq a$, $r\leq r^2\leq a$ and $p\leq q^5\leq a$, and also $a\geq2$.
So $\ell\geq1$ and $\mathcal L=|T|k^2\ell\geq1$.

\emph{Step 2: $\kappa_F$.} The terms of $\kappa_F$ are those of the $s\in T$. For $s\in T$,
$d_T(s)\leq2a\leq a^2$, so $\lfloor\log_pd_T(s)\rfloor\leq2\ell$. Hence $\kappa_F\leq2|T|k\ell\leq20\mathcal L$.

\emph{Step 3: the size of $c'$.} By (3) of the situation, $C\mathcal L\leq D_p$ and $8(D_p+1)\leq c_*\leq c'$. As
$C\geq10$, $80\mathcal L\leq8C\mathcal L\leq8D_p<c'$.

\emph{Step 4: the sizes of $\mathcal D$.} Let $\mathcal D$ be adapted to $(J,f)$ with $d^{\mathrm{pl}}$ and $\kappa$, and
$(\rho,d^{\mathrm{pl}},d^{\mathrm{no}})$ a block decomposition of $d_T$ for $\mathcal D$. Then $d=d^{\mathrm{pl}}$, $c=c'$,
$n_j=|\Lambda'_{\hat\jmath}|$, $n^c_j=n_j+1$, $P(j,m)=p^{g_j}$ with $g_j=f(\hat\jmath+1,r-1)-f(\hat\jmath,r-1)$, and
$s_{j,i+1}=\frac38s_j(i)$ for $0\leq i\leq n_j$. Let $y$ be a position with place $x$. By
Lemma~\ref{pd-symbols-test-flat-chains-datum-numbers-places}(1), $\bar x\in T$, so the block decomposition gives
$\lambda(y)d(y)\leq\lambda(y)(d^{\mathrm{pl}}(y)+d^{\mathrm{no}}(y))\leq d_T(\bar x)\leq2a$; as $\lambda(y)\geq1$,
$d(y)\leq2a$. By Lemma~\ref{pd-symbols-test-flat-chains-datum-numbers-places}(2), $\mathcal D$ has at most $|T|$
positions. Lemma~\ref{pd-symbols-test-flat-chains-datum-numbers-sizes}, with $N=|T|$ and Step~1, gives
$\kappa_{\mathcal D}\geq0$ and $D_m+D_l+3\kappa_{\mathcal D}\leq20\mathcal L$; with Step~3,
$4(D_m+D_l+3\kappa_{\mathcal D})\leq80\mathcal L<c'=c$.

\emph{Step 5: condition (H) for $26$.} Let $j$ be a level, $m<r$, and $0\leq i\leq n_{j+1}$, so that $C_{j+1,i+1}$ is a
chain. Let $\hat\jmath'$ be the representative of $j+1$, $x=x_{j+1}(i)$, $B=B_{\hat\jmath'}(x)$ and $s=\lfloor
B/(Yk^2)\rfloor$, so that $s_{j+1,i+1}=\frac38s$. Put $D_j=f(\hat\jmath,r-1)-f(\hat\jmath,0)$ and
$E=f(\hat\jmath',r-1)-f(\hat\jmath'-1,0)$. Then $D_j+g_j=E$: if $\hat\jmath<J-1$ then $\hat\jmath'=\hat\jmath+1$, and if
$\hat\jmath=J-1$ then $\hat\jmath'=0$ and $f(J,i')=f(0,i')+n'$, $f(J-1,i')=f(-1,i')+n'$ by (1) of validity. By the
spacing, applied at $\hat\jmath$, $\hat\jmath+1$ and $\hat\jmath'$, and (2) of validity, $D_j$, $g_j$ and $E$ are less
than $G'$; so $p^{D_j},p^{g_j}\leq q$ and $p^E\leq q$.

By Lemma~\ref{pd-symbols-test-flat-chains-datum-numbers-places}(3), $x$ is a new element of $\Lambda_{\hat\jmath'}$. So
$o=\operatorname{old}_{\hat\jmath'}(x)$ lies in $\Lambda_{\hat\jmath'}$ (Lemma~\ref{pd-symbols-division-rows}(3)): $\bar o\in
T$, so $\deg(\bar o)\geq a$, and $\operatorname{dist}_{\hat\jmath'}(o)=f(\hat\jmath',r-1)-o\leq E$. Hence
$B\geq a/p^E=a/(p^{D_j}p^{g_j})\geq a/q$.

As $p\geq13$, $L_1\leq L_2$ by (1) of the situation and $\lambda(y)d(y)\leq2a$ by Step~4,
Lemma~\ref{pd-symbols-test-flat-chains-datum-numbers-letters} gives $d_L(j)\leq3ap^{-D_j-L_1-1}$. As $C\geq465$,
Lemma~\ref{pd-symbols-test-flat-chains-datum-numbers-slack}(2), with $D=D_j$, $Q=p^{g_j}$ and $d=d_L(j)$, gives
$rd_L(j)/P(j,m)\leq s/8$, and Lemma~\ref{pd-symbols-test-flat-chains-datum-numbers-slack}(1), with $\alpha=0$, gives
$100r^2c'\leq s$, as the other terms on its left are nonnegative. As
$r\geq2$, $26r\leq25r^2$, so $26rc\leq25r^2c'\leq s/4$, and
\[
 \frac{rd_L(j)}{P(j,m)}+26rc\leq\frac s8+\frac s4=\frac38s=s_{j+1,i+1}.
\]
% Lean: restNum D CLa Ck Crk Cdeg (DiagonalSupport/LeOfOneLeMul.lean) gives CDm = 20 and the C1 of numStep_of, which
%   is max (max C0 C5) 10 with C0 and C5 the constants of CountingSumsFrames.Top.paramsII0 1 1 and paramsII5 1 1; their
%   proofs give C0 = 465 and C5 = 100, so C1 = 465. numStep_of: Step 1 is hyps_basic and nb_of (NB: 2 <= a, p <= a,
%   4 |T| <= a, r <= a, 1 <= log_p a and 1 <= L, from factor_le_product_seven and the last condition h12; G' >= 1 is
%   hG1); Step 2 is kapFl_le (kap_fin); Step 3 is hcP (C >= 10, hDp, h11 and hcS). numAt_of: Step 4 is hd and hd2 (the
%   block decomposition hBk.2 with posTL_seg), hcard (card_TL_le) and sizes_le; Step 5 is slack_bound_joint (seg_shift
%   for D_j + g_j = E, Bo_ge with oldest_mem_segment_letters, posU_mem, dLL_le, hII5 and hII0, and slack_bound_real).
\end{proof}

\begin{lemma}\label{pd-symbols-test-flat-chains-datum}
Let $C_{La},C_k,C_{rk},C_{\deg}$ be real numbers with the property of Theorem~\ref{pd-symbols-count}. There are
$C_D,C_2>0$ such that the following holds for every $C\geq C_2$. Let the data and $S$ be as in
Lemma~\ref{pd-symbols-test-flat-chains-membership}, and let $\kappa_F$ be as in Lemma~\ref{pd-symbols-test-flat-real}.
Then there is an integer $0\leq D_X\leq C_D|T|k^2\log_pa$ with the following property. Let $G\in\mathcal K[X]$ have
multidegree $d_T$, lie in $\mathrm{Op}_{T,c_0}$ of the $\mathcal K$-span of the $P_bS$, lie in $\mathcal R[X]$ but not in
$\varpi_0\mathcal R[X]$, and satisfy the membership condition of Definition~\ref{pd-symbols-test-flat-chains-tied}.
Then there are a signed chain datum $\mathcal D$ of rank $r$ over $\mathcal R$ that satisfies (H) for $C_H=26$ and
$4(D_m+D_l+3\kappa_{\mathcal D})<c$, and an admissible $F_2$ with $v_{\mathrm F}(F_2)=0$
(Definition~\ref{pd-symbols-test-chains}), such that for every projection pair for $(\mathcal D,F_2)$ and
$(C_E,C_6,C_{18})=(2,24,60)$ there are one-sided elements $\sigma_1,\dots,\sigma_{n_q}$ of multidegree $d_T$
(Definition~\ref{pd-symbols-test-one-sided}, for $(J,f)$, $N_1$, $G'$, $c'$, $\deg$ and $Y$) for which
$p^{\kappa_F}\langle G,\sum_q\sigma_q\rangle_{\mathrm F}/\prod_sd_T(s)!$ lies in $\mathcal R$ and divides $p^{D_X}$.
\end{lemma}
\leanref{Formalism/SharpRestriction/GrowthBounds/TwistedSubstitution/Proposition5Gram.lean}{Formalism.MainNotions.TestSideWords.NSConcl}
\leanref{Formalism/SharpRestriction/GrowthBounds/TwistedSubstitution/Proposition5Gram.lean}{Formalism.MainNotions.TestSideWords.Proposition5At}
\leanref{Formalism/SharpRestriction/GrowthBounds/TwistedSubstitution/Proposition5Gram.lean}{Formalism.MainNotions.TestSideWords.Proposition5Gram}

\begin{proof}
The hypothesis that $G$ lies in $\mathrm{Op}_{T,c_0}$ of the $\mathcal K$-span of the $P_bS$ is not used.
Let $C_{Dn},C_P$, $C_B$ and $C_{Dm},C_N$ be the constants of Lemmas~\ref{pd-symbols-test-flat-chains-datum-pick},
\ref{pd-symbols-test-flat-chains-datum-membership} and~\ref{pd-symbols-test-flat-chains-datum-numbers}. Put
$C_D=C_{Dn}+2C_{Dm}+1$ and $C_2=\max(C_P,C_B,C_N,C_{Dn}+3C_{Dm})$. Let $C\geq C_2$ and the data be as in the statement,
and put $\mathcal L=|T|k^2\log_pa$, $D_X=\lceil(C_{Dn}+2C_{Dm})\mathcal L\rceil$ and $P=\prod_sd_T(s)!$. By
Lemma~\ref{pd-symbols-test-flat-chains-datum-numbers}, $\mathcal L\geq1$ and $\kappa_F\leq C_{Dm}\mathcal L$; so
$0\leq D_X<(C_{Dn}+2C_{Dm})\mathcal L+1\leq C_D\mathcal L$.

\emph{Step 1: the size of $c'$.} By (3) of Situation~\ref{pd-edge-situation}, $C\mathcal L\leq D_p$ and
$8(D_p+1)\leq c_*$, and $c'=c_*+\omega\geq c_*$. As $C\geq C_{Dn}+3C_{Dm}\geq0$,
$(C_{Dn}+3C_{Dm})\mathcal L\leq8C\mathcal L<c'$.

\emph{Step 2: the datum.} Let $G$ be as in the statement. Lemma~\ref{pd-symbols-test-flat-chains-datum-pick} gives a
signed chain datum $\mathcal D$ adapted to $(J,f)$ with $d^{\mathrm{pl}}$ and $\kappa$, a block decomposition
$(\rho,d^{\mathrm{pl}},d^{\mathrm{no}})$, an admissible $F_2$ with $v_{\mathrm F}(F_2)=0$, a normal test $N_0$ and
$k_0$ with the properties stated there; so $d=d^{\mathrm{pl}}$ and $c=c'$. By
Lemma~\ref{pd-symbols-test-flat-chains-datum-numbers}, $\mathcal D$ satisfies (H) for $26$,
$4(D_m+D_l+3\kappa_{\mathcal D})<c$, $\kappa_{\mathcal D}\geq0$ and $D_m+D_l+3\kappa_{\mathcal D}\leq C_{Dm}\mathcal L$.

\emph{Step 3: a Fock bound.} If $H\in\mathcal K[X]$ has coefficients of valuation at least $t$, then
$v(\langle G,H\rangle_{\mathrm F})\geq t+v(P)-\kappa_F$: by Lemma~\ref{pd-symbols-flat-real-multinomial}(1), $P$ divides
$p^{\kappa_F}\mu!$ for every exponent vector $\mu$ of multidegree $d_T$, and $G$ has integral coefficients and multidegree
$d_T$. If $A\in\mathcal K[x]$ has coefficients of valuation at least $t$ and $N\in\mathcal K[w]$ has integral
coefficients, so does $\operatorname{asm}_\rho(A,N)$; and $\varphi$ maps polynomials with coefficients of valuation at
least $t$ to such polynomials.

\emph{Step 4: a projection pair.} Let $(Z,R')$ be a projection pair for $(\mathcal D,F_2)$ and $(2,24,60)$. By
Lemma~\ref{pd-symbols-test-flat-chains-datum-membership}, applied to $Z$, there are one-sided elements
$\sigma_1,\dots,\sigma_{n_q}$ of multidegree $d_T$ with $\sum_q\sigma_q=\operatorname{asm}_\rho(\varphi(Z),N_0)$. Put
$x=\langle G,\operatorname{asm}_\rho(\varphi(Z),N_0)\rangle_{\mathrm F}$ and $E=E(2)=D_m+D_l+2\kappa_{\mathcal D}$; then
$E\leq D_m+D_l+3\kappa_{\mathcal D}\leq C_{Dm}\mathcal L$ and $E<c$. Let $A=\varphi(Z+R')$. It has degree $d(y)$ in the
row of every position $y$, and integral coefficients, as $Z$ is integral and the coefficients of $R'$ have valuation at
least $c-E\geq0$. By Lemma~\ref{pd-symbols-test-flat-chains-datum-pick},
$v(\langle G,\operatorname{asm}_\rho(A,N_0)\rangle_{\mathrm F}-k_0B_{\mathcal D}(F_2,A))\geq c'+v(P)-\kappa_F$, and by
Step 3, $v(\langle G,\operatorname{asm}_\rho(\varphi(R'),N_0)\rangle_{\mathrm F})\geq c-E+v(P)-\kappa_F$. As $(Z,R')$ is a
projection pair, $v(B_{\mathcal D}(F_2,A))\leq E+\kappa_{\mathcal D}+n_0$, so
\[
 v\big(k_0B_{\mathcal D}(F_2,A)\big)\leq v(P)-n_0+C_{Dn}\mathcal L+E+\kappa_{\mathcal D}+n_0\leq v(P)+(C_{Dn}+C_{Dm})\mathcal L .
\]
As $\operatorname{asm}_\rho$ is additive in its first argument and $\varphi$ is additive,
\[
 x=k_0B_{\mathcal D}(F_2,A)+\big(\langle G,\operatorname{asm}_\rho(A,N_0)\rangle_{\mathrm F}-k_0B_{\mathcal D}(F_2,A)\big)
 -\langle G,\operatorname{asm}_\rho(\varphi(R'),N_0)\rangle_{\mathrm F}.
\]
The last two terms have valuation at least $c'-C_{Dm}\mathcal L+v(P)-\kappa_F$, as $c=c'$, and this exceeds
$v(P)+(C_{Dn}+C_{Dm})\mathcal L$, since $c'>(C_{Dn}+3C_{Dm})\mathcal L\geq(C_{Dn}+2C_{Dm})\mathcal L+\kappa_F$ by Step 1.
So $v(x)\leq v(P)+(C_{Dn}+C_{Dm})\mathcal L$, and
$v(p^{\kappa_F}x/P)\leq\kappa_F+(C_{Dn}+C_{Dm})\mathcal L\leq(C_{Dn}+2C_{Dm})\mathcal L\leq D_X$. By Step 3 with $t=0$,
$v(x)\geq v(P)-\kappa_F$, so $v(p^{\kappa_F}x/P)\geq0$. Hence $p^{\kappa_F}x/P$ lies in $\mathcal R$ and divides $p^{D_X}$.
% Lean: Proposition5Gram D CLa Ck Crk Cdeg 26 2 24 60 is proposition5_holds (DiagonalSupport.lean) =
%   proposition5_of_three_parts D ... (restPick D ...) (restMemb D ...) (restNum D ...) (DiagonalSupport/
%   Proposition5OfThreeParts.lean): the constants CDn + 2 CDm + 1 and max (max C1 C2) (max C3 (CDn + 3 CDm)), then
%   proposition5_at_of_three_parts with DX = ceil((CDn + 2 CDm) large2). Step 1 is hcPge (hDp, hcS and cStar <= cPrime)
%   and hbig; Steps 2 and 4 are the uses of PickStep, NumStep and MembStep there, and the bound is value_core: hEE and
%   hRpos for E, hApol and hAint for A, e1 (hval), e2 and the Fock bound geQ_fock (geQ_expFact through multDiv_kapFl,
%   Lemma pd-symbols-flat-real-multinomial(1); asm_coefGe and ph_coefGe for the integrality), hmain (leQ_mul with hk0
%   and the last clause g6 of the pair), hsplit, herr and hx (leQ_add_of_lt), then leQ_scale, geQ_scale and valLe_of_vK.
%   The clause of the pair on v_F(ph(G + R)) is not used, and neither is the hypothesis G in OpCl of Proposition5At
%   (proposition5_at_of_three_parts introduces it as _). What remains open is in Lemmas
%   pd-symbols-test-flat-chains-datum-pick-inputs (about 3,700 lines) and -datum-membership-parts (about 6,700 lines),
%   which share files of DiagonalSupport/ (49 files, 8,240 lines); Lemma -datum-numbers is written out.
\end{proof}

\begin{lemma}\label{pd-symbols-test-flat-chains}
Let $C_{La},C_k,C_{rk},C_{\deg}$ be real numbers with the property of Theorem~\ref{pd-symbols-count}. There are
$C_D,C_1>0$ such that the following holds for every $C\geq C_1$. Let the data be as in
Lemma~\ref{pd-symbols-test-flat}, without $e$, $\gamma$, $g$ and the $\Phi_b$, and let $\kappa_F$ be as in
Lemma~\ref{pd-symbols-test-flat-real}. Let $S=\Sigma_g(\mathrm{SP}_g)$, which is the same polynomial for every
$g\in GL_k(W)$ (Step~3 of the proof of Lemma~\ref{pd-symbols-primitive}). Then there is an integer
$0\leq D_X\leq C_D|T|k^2\log_pa$ with the following property. Let $G\in\mathcal K[X]$ have multidegree $d_T$, lie
in $\mathrm{Op}_{T,c_0}$ of the $\mathcal K$-span of the $P_bS$, and lie in $\mathcal R[X]$ but not in
$\varpi_0\mathcal R[X]$. Then there are a signed chain datum $\mathcal D$ of rank $r$ over $\mathcal R$ that
satisfies (H) for $C_H=26$ and $4(D_m+D_l+3\kappa_{\mathcal D})<c$, and an admissible $F_2$ with
$v_{\mathrm F}(F_2)=0$ (Definition~\ref{pd-symbols-test-chains}), such that for every projection pair for
$(\mathcal D,F_2)$ and $(C_E,C_6,C_{18})=(2,24,60)$ there are one-sided elements $\sigma_1,\dots,\sigma_{n_q}$ of
multidegree $d_T$ (Definition~\ref{pd-symbols-test-one-sided}, for $(J,f)$, $N_1$, $G'$, $c'$, $\deg$ and $Y$)
for which $p^{\kappa_F}\langle G,\sum_q\sigma_q\rangle_{\mathrm F}/\prod_sd_T(s)!$ lies in $\mathcal R$ and divides
$p^{D_X}$.
\end{lemma}
\leanref{Formalism/SharpRestriction/GrowthBounds/DualFockLinear/RestXiPOfNS.lean}{Formalism.MainNotions.TestSideCounting.NSAt}
\leanref{Formalism/SharpRestriction/GrowthBounds/DualFockLinear/RestXiPOfNS.lean}{Formalism.MainNotions.TestSideCounting.NSStep}
\leanref{Formalism/SharpRestriction/GrowthBounds/DualFockLinear/RestXiPOfNS.lean}{Formalism.MainNotions.TestSideCounting.RestNS}
\leanref{Formalism/SharpRestriction/GrowthBounds/TwistedSubstitution/Proposition5Gram.lean}{Formalism.MainNotions.TestSideWords.restNS_of}
\leanref{Formalism/SharpRestriction/GrowthBounds/TwistedSubstitution/Proposition5Gram.lean}{Formalism.MainNotions.TestSideWords.nsStep_of}

\begin{proof}
Let $C'_1$ be the constant of Lemma~\ref{pd-symbols-test-flat-chains-membership}, and $C_D$ and $C_2$ those of
Lemma~\ref{pd-symbols-test-flat-chains-datum}; put $C_1=\max(C'_1,C_2)$. Let $C\geq C_1$ and the data be as in the
statement, and let $D_X$ be the integer of Lemma~\ref{pd-symbols-test-flat-chains-datum}, which satisfies
$0\leq D_X\leq C_D|T|k^2\log_pa$. Let $G$ be as in the statement. By
Lemma~\ref{pd-symbols-test-flat-chains-membership}, $G$ satisfies the membership condition of
Definition~\ref{pd-symbols-test-flat-chains-tied}, and Lemma~\ref{pd-symbols-test-flat-chains-datum} gives $\mathcal D$
and $F_2$ with the stated property.
% Sizes: what remains open is in Lemmas pd-symbols-test-flat-chains-membership-monomials-pairs-image, -normal, -twisted,
%   -plane and -numbers (about 6,500 lines together, counted by lean_size with the book's current citations), behind
%   Lemma pd-symbols-test-flat-chains-membership, and Lemmas pd-symbols-test-flat-chains-datum-pick-inputs (about 3,700
%   lines) and -datum-membership-parts (about 6,700 lines), behind Lemma pd-symbols-test-flat-chains-datum (they share
%   files of DiagonalSupport/, 49 files, 8,240 lines). The membership
%   condition is the Lean's MemAt, written out in Definition pd-symbols-test-flat-chains-tied.
% Lean: restNS_of takes C1 = max C1' C2 and the constant CDf of Proposition5Gram; nsStep_of is the rest of the proof
%   (DataII from the data, then nsAt_of with the DX of Proposition5At and MemStep for G).
% Lean: RestNS D CLa Ck Crk Cdeg 26 2 24 60 is restNS_of D ... (hM _ _ _ _) (hP _ _ _ _) in
%   TwistedSubstitution.lean (stmtE_of_TenMP). restNS_of is nsStep_of (TwistedSubstitution/Proposition5Gram.lean)
%   from two inputs: MemStep (every G of the operator closure satisfies MemAt: for every signed chain datum tied to
%   the run system, the plane coefficients of the projected level image of G are admissible), which is RestMEM =
%   restMEM_of_last / memCondition_of_theoremII (PlaneCoefficients.lean, TwistedSubstitution/MemRed.lean: Lemma5aTheoremII
%   through lemma5a_of_monomial, and RestWord through restWord_of_steps); and Proposition5At (from DataII, MemAt and
%   the conditions on G, the conclusion NSConcl), which is proposition5_holds (DiagonalSupport.lean) =
%   Proposition5OfThreeParts with RestMemb and RestNum.
\end{proof}

\subsection{Projection pairs}\label{subsection-pd-projection-pairs}

\begin{definition}\label{pd-symbols-test-chains-projection-clauses-mahler}
Let $p$ be a prime, $v_p$ the $p$-adic valuation of $\mathbf Q$, with $v_p(0)=+\infty$, and
$f\colon\mathbf Z_{\geq0}\to\mathbf Q$ a function, and put $(\Delta f)(m)=f(m+1)-f(m)$. The \emph{Mahler
coefficients} of $f$ are the numbers $a_j(f)=(\Delta^jf)(0)$, $j\geq0$. For real numbers $\theta$ and $s$, the
function $f$ has \emph{decay} $(\theta,s)$ at $p$ if $v_p(a_j(f))\geq\theta j+s$ for every $j\geq0$.
\end{definition}
\leanref{Formalism/SharpRestriction/GrowthBounds/RootWeights/PadicGeMulNat.lean}{Formalism.MainNotions.G.GLWeightTheory.mah}
\leanref{Formalism/SharpRestriction/GrowthBounds/RootWeights/PadicGeMulNat.lean}{Formalism.MainNotions.G.GLWeightTheory.Dec}
\leanref{Formalism/SharpRestriction/GrowthBounds/CovariantDegrees.lean}{Formalism.MainNotions.CovariantDegrees.PadicGe}

\begin{lemma}\label{pd-symbols-test-chains-projection-clauses-mahler-rules}
Let $p$ be a prime, $f,g,u\colon\mathbf Z_{\geq0}\to\mathbf Q$ functions and $\theta,s,s'$ real numbers; we use the
notation of Definition~\ref{pd-symbols-test-chains-projection-clauses-mahler}.
\begin{enumerate}
\item $f(m)=\sum_{j=0}^ma_j(f)\binom mj$ for every $m\geq0$. Each $a_j$ is $\mathbf Q$-linear in $f$. The Mahler
  coefficients of a constant $c$ are $a_0=c$ and $a_j=0$ for $j\geq1$, and those of $m\mapsto\binom ml$ are $a_l=1$ and
  $a_j=0$ for $j\neq l$.
\item If $\theta\geq0$ and $f$ has decay $(\theta,s)$, then $v_p(f(m))\geq s$ for every $m\geq0$.
\item If $v_p(f(m))\geq s$ for every $m\geq0$, then $v_p(a_j(f))\geq s$ for every $j\geq0$.
\item If $\theta\geq0$, $f$ has decay $(\theta,s)$ and $g$ has decay $(\theta,s')$, then $fg$ has decay
  $(\theta,s+s')$.
\item If $\theta\geq0$, $u$ has decay $(\theta,\theta)$, $g(0)=1$ and $g(m+1)=g(m)(1+u(m))$ for every $m\geq0$, then $g$
  has decay $(\theta,0)$.
\item If $\theta>0$ and $u$ has decay $(\theta,\theta)$, then $1+u(m)\neq0$ and $v_p(1+u(m))=0$ for every $m\geq0$, and
  the function $m\mapsto(1+u(m))^{-1}-1$ has decay $(\theta,\theta)$.
\end{enumerate}
\end{lemma}
\leanref{Formalism/SharpRestriction/GrowthBounds/RootWeights/MahPadicGeOfDvd.lean}{Formalism.MainNotions.G.GLWeightTheory.mah_newton}
\leanref{Formalism/SharpRestriction/GrowthBounds/RootWeights/MahPadicGeOfDvd.lean}{Formalism.MainNotions.G.GLWeightTheory.mah_add}
\leanref{Formalism/SharpRestriction/GrowthBounds/RootWeights/MahPadicGeOfDvd.lean}{Formalism.MainNotions.G.GLWeightTheory.mah_smul}
\leanref{Formalism/SharpRestriction/GrowthBounds/RootWeights/MahPadicGeOfDvd.lean}{Formalism.MainNotions.G.GLWeightTheory.mah_const}
\leanref{Formalism/SharpRestriction/GrowthBounds/RootWeights/MahPadicGeOfDvd.lean}{Formalism.MainNotions.G.GLWeightTheory.mah_choose}
\leanref{Formalism/SharpRestriction/GrowthBounds/RootWeights/MahPadicGeOfDvd.lean}{Formalism.MainNotions.G.GLWeightTheory.mah_val}
\leanref{Formalism/SharpRestriction/GrowthBounds/RootWeights/MahPadicGeOfDvd.lean}{Formalism.MainNotions.G.GLWeightTheory.mah_of_vals}
\leanref{Formalism/SharpRestriction/GrowthBounds/RootWeights/MahPadicGeOfDvd.lean}{Formalism.MainNotions.G.GLWeightTheory.hMul_proved}
\leanref{Formalism/SharpRestriction/GrowthBounds/RootWeights/MahPadicGeOfDvd.lean}{Formalism.MainNotions.G.GLWeightTheory.mah_prodLe}
\leanref{Formalism/SharpRestriction/GrowthBounds/RootWeights/MahPadicGeOfDvd.lean}{Formalism.MainNotions.G.GLWeightTheory.mah_running}
\leanref{Formalism/SharpRestriction/GrowthBounds/RootWeights/MahPadicGeOfDvd.lean}{Formalism.MainNotions.G.GLWeightTheory.mah_unit}
\leanref{Formalism/SharpRestriction/GrowthBounds/RootWeights/MahPadicGeOfDvd.lean}{Formalism.MainNotions.G.GLWeightTheory.mah_inv}

\begin{proof}
(1) The shift $f\mapsto f(\cdot+1)$ is $1+\Delta$, so $f(m)=((1+\Delta)^mf)(0)=\sum_j\binom mj(\Delta^jf)(0)$. The
operator $\Delta$ is linear and kills constants. By Pascal's rule, $\Delta$ maps $m\mapsto\binom ml$ to
$m\mapsto\binom m{l-1}$ if $l\geq1$ and to $0$ if $l=0$, and $\binom0l$ is $1$ for $l=0$ and $0$ otherwise; induction on
$j$ gives the Mahler coefficients of $m\mapsto\binom ml$.

(2) By (1), as the $\binom mj$ are integers and $\theta j\geq0$.

(3) By induction on $j$, for all $f$ at once: $a_0(f)=f(0)$, and $a_{j+1}(f)=a_j(\Delta f)$, where
$v_p((\Delta f)(m))\geq s$ for every $m$.

(4) Here $\Delta(fg)=(\Delta f)g+f(\Delta g)+(\Delta f)(\Delta g)$, and $a_j(\Delta f)=a_{j+1}(f)$. We show by
induction on $N$: for all real $s,s'$ and all $f,g$ with $v_p(a_j(f))\geq\theta j+s$ and $v_p(a_j(g))\geq\theta j+s'$
for $j\leq N$, $v_p(a_N(fg))\geq\theta N+s+s'$. For $N=0$, $a_0(fg)=f(0)g(0)$. For the step, $a_{N+1}(fg)=a_N(\Delta(fg))$
is the sum of $a_N((\Delta f)g)$, $a_N(f\Delta g)$ and $a_N((\Delta f)(\Delta g))$. For $j\leq N$, $\Delta f$ and $\Delta g$
satisfy the bounds with $s+\theta$ and $s'+\theta$ in place of $s$ and $s'$, and $f$ and $g$ satisfy them. So the three
terms have valuation at least $\theta N+s+s'+\theta$, as $\theta\geq0$, and $a_{N+1}(fg)$ as well.

(5) Here $\Delta g=gu$. We show by induction on $N$ that $v_p(a_j(g))\geq\theta j$ for $j\leq N$. For $N=0$,
$a_0(g)=1$. If it holds for $N$, then $a_{N+1}(g)=a_N(gu)$, and the claim in the proof of (4), with $s=0$ and
$s'=\theta$, gives $v_p(a_N(gu))\geq\theta N+\theta$.

(6) By (2), $v_p(u(m))\geq\theta>0$; so $1+u(m)\neq0$ and $v_p(1+u(m))=0$. Fix $N\geq0$. For $x\in\mathbf Q$ with
$1+x\neq0$,
\[
 (1+x)^{-1}-1=\sum_{i=0}^N(-x)^{i+1}+(-x)^{N+2}(1+x)^{-1}.
\]
The function $-u$ has decay $(\theta,\theta)$, so by (4) and induction on $i$ the function $(-u)^{i+1}$ has decay
$(\theta,(i+1)\theta)$, and $v_p(a_N((-u)^{i+1}))\geq\theta N+\theta$. The function
$w(m)=(-u(m))^{N+2}(1+u(m))^{-1}$ satisfies $v_p(w(m))\geq(N+2)\theta$, so $v_p(a_N(w))\geq(N+2)\theta\geq\theta N+\theta$
by (3). By the linearity of $a_N$, $v_p(a_N)$ of $m\mapsto(1+u(m))^{-1}-1$ is at least $\theta N+\theta$.
% Lean: (1) is mah_newton (shift_eq_sum_fwdDiff_iter), mah_add, mah_smul, mah_sum, mah_const and mah_choose;
%   (2) is mah_val, (3) is mah_of_vals, (4) is hMul_proved = mah_dec_mul through mah_prodLe (mah_leibniz,
%   mah_decLe_diff), (5) is mah_running and (6) is mah_unit with mah_inv (mah_geom, mah_dec_pow, mah_padicGe_pow).
%   PadicGe p x r is x = 0 or r <= v_p(x). MahPadicGeOfDvd.lean, 407 lines, and PadicGeMulNat.lean, 92 lines.
\end{proof}

\begin{lemma}\label{pd-symbols-test-chains-projection-clauses-mahler-factorials}
Let $p$ be a prime, $t\geq0$ an integer and $P=p^{t+1}$. For integers $m,\rho\geq0$ put $c(P,m)=(Pm)!/((P!)^m\,m!)$ and
$d(P,m,\rho)=\prod_{y=1}^\rho(1+Pm/y)$; we use the notation of
Definition~\ref{pd-symbols-test-chains-projection-clauses-mahler}.
\begin{enumerate}
\item Let $0\leq\rho<P$. Then there are integers $A_0,\dots,A_\rho$, with $p^l$ dividing $A_l$, such that
  $d(P,m,\rho)=\binom{\rho+Pm}\rho=\sum_{l=0}^\rho A_l\binom ml$ for every $m\geq0$.
\item $c(P,0)=1$, and $c(P,m+1)=c(P,m)\,d(P,m,P-1)$ for every $m\geq0$.
\item Let $0\leq\rho<P$. Then $c(P,m)$ and $d(P,m,\rho)$ are nonzero with $v_p(c(P,m))=v_p(d(P,m,\rho))=0$ for every
  $m\geq0$, and the four functions $m\mapsto c(P,m)$, $m\mapsto c(P,m)^{-1}$, $m\mapsto d(P,m,\rho)$ and
  $m\mapsto d(P,m,\rho)^{-1}$ have decay $(\frac12,0)$ and $(\frac14,0)$ at $p$.
\end{enumerate}
\end{lemma}
\leanref{Formalism/SharpRestriction/GrowthBounds/ObjPol.lean}{Formalism.MainNotions.G.cUnit}
\leanref{Formalism/SharpRestriction/GrowthBounds/ObjPol.lean}{Formalism.MainNotions.G.dUnit}
\leanref{Formalism/SharpRestriction/GrowthBounds/RootWeights/BinomialExpansionCoeffZeroOfDeg.lean}{Formalism.MainNotions.G.GLWeightTheory.unit_product_eq_choose}
\leanref{Formalism/SharpRestriction/GrowthBounds/RootWeights/BinomialExpansionCoeffZeroOfDeg.lean}{Formalism.MainNotions.G.GLWeightTheory.hdU_proved}
\leanref{Formalism/SharpRestriction/GrowthBounds/RootWeights/BinomialExpansionCoeffZeroOfDeg.lean}{Formalism.MainNotions.G.GLWeightTheory.coeff_binomial_power_divisibility}
\leanref{Formalism/SharpRestriction/GrowthBounds/RootWeights/MahPadicGeOfDvd.lean}{Formalism.MainNotions.G.GLWeightTheory.mah_cUnit_succ}
\leanref{Formalism/GrowthSums/FactorialBaseGrowth/Coarse.lean}{Formalism.GrowthSums.blockFactorialQuotientProd}
\leanref{Formalism/SharpRestriction/GrowthBounds/RootWeights/MahPadicGeOfDvd.lean}{Formalism.MainNotions.G.GLWeightTheory.mah_dU}
\leanref{Formalism/SharpRestriction/GrowthBounds/RootWeights/MahPadicGeOfDvd.lean}{Formalism.MainNotions.G.GLWeightTheory.mah_unitFn_half}
\leanref{Formalism/SharpRestriction/GrowthBounds/RootWeights/MahPadicGeOfDvd.lean}{Formalism.MainNotions.G.GLWeightTheory.hDec_of}
\leanref{Formalism/SharpRestriction/GrowthBounds/RootWeights/MahPadicGeOfDvd.lean}{Formalism.MainNotions.G.GLWeightTheory.hUnit_of}

\begin{proof}
(1) By induction on $\rho$, $d(P,m,\rho)=\prod_{y=1}^\rho(y+Pm)/y=\binom{\rho+Pm}\rho$. Let $A_l$ be the coefficient of
$X^\rho$ in $(1+X)^\rho((1+X)^P-1)^l\in\mathbf Z[X]$. Writing
$(1+X)^{\rho+Pm}=(1+X)^\rho\big(((1+X)^P-1)+1\big)^m=\sum_{l=0}^m\binom ml(1+X)^\rho((1+X)^P-1)^l$ and taking the
coefficient of $X^\rho$ gives $\binom{\rho+Pm}\rho=\sum_lA_l\binom ml$, and $A_l=0$ for $l>\rho$, as $X$ divides
$(1+X)^P-1$. As $p$ divides $\binom Pi$ for $0<i<P$, $(1+X)^P-1=X^P+ph$ with $h\in\mathbf Z[X]$. By induction on $l$,
using $f(X^P+ph)^{l+1}=X^P\cdot f(X^P+ph)^l+p\cdot fh(X^P+ph)^l$, the coefficient of $X^j$ in $f(X^P+ph)^l$ is
divisible by $p^{\max(l-\lfloor j/P\rfloor,0)}$ for every $f\in\mathbf Z[X]$ and $j\geq0$. For $j=\rho<P$ this is $p^l$.

(2) $c(P,0)=1$. Next, $(Pm+P)!=(Pm)!\prod_{y=1}^P(Pm+y)=(Pm)!\,P!\prod_{y=1}^P(1+Pm/y)$, and the factor for $y=P$ is
$m+1$. Dividing by $(P!)^{m+1}(m+1)!$ gives $c(P,m+1)=c(P,m)\prod_{y=1}^{P-1}(1+Pm/y)=c(P,m)\,d(P,m,P-1)$.

(3) Let $0\leq\rho<P$ and $u_\rho(m)=d(P,m,\rho)-1$. As $d(P,0,\rho)=1$, (1) and
Lemma~\ref{pd-symbols-test-chains-projection-clauses-mahler-rules}(1) give $a_0(u_\rho)=0$, $a_j(u_\rho)=A_j$ for
$1\leq j\leq\rho$ and $a_j(u_\rho)=0$ for $j>\rho$. As $v_p(A_j)\geq j\geq\frac j2+\frac12$ for $j\geq1$, $u_\rho$ has
decay $(\frac12,\frac12)$. By part (6) of that lemma, $d(P,m,\rho)=1+u_\rho(m)$ is nonzero with $v_p=0$, and
$m\mapsto d(P,m,\rho)^{-1}-1$ has decay $(\frac12,\frac12)$, hence $(\frac12,0)$. The constant $1$ has decay
$(\frac12,0)$, and by the linearity of the $a_j$ a sum of two functions with decay $(\frac12,0)$ has decay
$(\frac12,0)$; so $d(P,\cdot,\rho)=1+u_\rho$ and $d(P,\cdot,\rho)^{-1}=1+(d(P,\cdot,\rho)^{-1}-1)$ have decay
$(\frac12,0)$.

Apply this to $\rho=P-1$. By (2), $c(P,0)=1$ and $c(P,m+1)=c(P,m)(1+u_{P-1}(m))$, so by induction on $m$, $c(P,m)$ is
nonzero with $v_p=0$, and $c(P,\cdot)$ has decay $(\frac12,0)$ by part (5) of
Lemma~\ref{pd-symbols-test-chains-projection-clauses-mahler-rules}. Also $c(P,0)^{-1}=1$ and
$c(P,m+1)^{-1}=c(P,m)^{-1}\big(1+((1+u_{P-1}(m))^{-1}-1)\big)$, where $m\mapsto(1+u_{P-1}(m))^{-1}-1$ has decay
$(\frac12,\frac12)$; so $c(P,\cdot)^{-1}$ has decay $(\frac12,0)$ by part (5). Finally, decay $(\frac12,0)$ implies
decay $(\frac14,0)$, as $\frac j2\geq\frac j4$.
% Lean: (1) is unit_product_eq_choose and hdU_proved, with A_l = binomialExpansionCoeff P rho rho l
%   (choose_shift_binomial_expansion, binomial_expansion_coeff_zero_of_gt, binomial_expansion_coeff_padic_dvd through
%   coeff_binomial_power_divisibility; the clause HdU of PadicGeMulNat.lean). (2) is mah_cUnit_zero and mah_cUnit_succ
%   from Formalism.GrowthSums.blockFactorialQuotientProd (coarseUnit_succ, factorial_add_prod). (3) is mah_dU (the
%   decay of u_rho), mah_unitFn_half (the four cases of unitFn, with mah_one_add, mah_inv and mah_running), hDec_of (the
%   weakening to 1/4) and hUnit_of (the units, through mah_unit and mah_val, and induction for c).
%   BinomialExpansionCoeffZeroOfDeg.lean, 191 lines, shared with clause (5)(b).
\end{proof}

\begin{lemma}\label{pd-symbols-test-chains-projection-clauses}
For a finite set $\Lambda$ each of whose elements is a vector letter or a covector letter, an integer $n\geq1$ and a
field $L$, let $L[z]$ be the polynomial ring in the $z_{v,c}$, $v\in\Lambda$, $1\leq c\leq n$. For $S\subset\Lambda$ form
$\operatorname{val}_S$ and $\lambda_S$, the operators $E^{(k)}_{ab}$ and $H^{(k)}_a$, and the weights of $U_S(G)$ as in
Definitions~\ref{pd-symbols-path-generators}(1) and~\ref{pd-symbols-path-weights}, and let
$E^S_{ab}=\sum_vz_{v,a}\,\partial/\partial z_{v,b}-\sum_{v'}z_{v',b}\,\partial/\partial z_{v',a}$, the sums over the
vector letters $v$ and the covector letters $v'$ in $S$. For $\mu\in\mathbf Z^n$, a \emph{highest weight vector of
weight $\mu$ for $S$} is a nonzero $x\in L[z]$ with $E^S_{aa}x=\mu(a)x$ for every $a$ and $E^S_{ab}x=0$ for $a<b$, and the
\emph{isotypic subspace} $L[z]^{(\mu)}_S$ is the $L$-span of the $u(x)$, with $u$ a composite of finitely many of the
$E^S_{ab}$, the identity included, and $x$ a highest weight vector of weight $\mu$ for $S$. An \emph{isotypic projector of
weight $\lambda\in\mathbf Z^n$ for $S$} is an $L$-linear endomorphism $e$ of $L[z]$ that is the identity on
$L[z]^{(\lambda)}_S$ and $0$ on $L[z]^{(\mu)}_S$ for every $\mu\neq\lambda$. For a monomial $z^m$ put $m!=\prod_{v,c}m(v,c)!$; the
\emph{Fock pairing} is $\langle F,G\rangle_{\mathrm F}=\sum_mF_mG_m\,m!$, and the \emph{divided derivative} of $F$ by a
monomial $z^\beta$ is $D_\beta F=\sum_mF_m\prod_{v,c}\binom{m(v,c)}{\beta(v,c)}z^{m-\beta}$, the sum over the $m$ with
$m\geq\beta$ entrywise. For a prime $p$ let $v_p$ be the $p$-adic valuation of $\mathbf Q$, with $v_p(0)=+\infty$, and for
an integer $x\geq0$ let $\ell_p(x)$ be the largest integer $e\geq0$ with $p^e\leq x$ if $x\geq1$, and $\ell_p(0)=0$.
\begin{enumerate}
\item Let $p\geq5$ be a prime, $n'\geq2$ an even integer, $\mathcal R$ a complete discrete valuation ring in which $p$ is
  a nonzero element of the maximal ideal and whose residue field contains a field of $p^{n'}$ elements, $\mathcal K$ its
  fraction field and $v$ the valuation of $\mathcal K$ with $v(p)=1$; for $F\in\mathcal K[z]$ put
  $v_{\mathrm F}(F)=\min_m(v(F_m)+v(m!)/2)$. Let $n\geq1$ and $N\geq0$ be integers and
  $S_0\subset S_1\subset\dots\subset S_{N-1}\subset\Lambda$. Then:
  \begin{enumerate}
  \item for each $i<N$, the $E^{(k)}_{ab}$ with $a\neq b$ and the $H^{(k)}_a$ for $S_i$ carry the polynomials with
    coefficients in $\mathcal R$ to such polynomials, and the $F$ with $v_{\mathrm F}(F)\geq0$ to such $F$;
  \item let $e_{i,\lambda}$, for $i<N$ and $\lambda\in\mathbf Z^n$, be isotypic projectors of weight $\lambda$ for $S_i$
    on $\mathcal K[z]$. Then $e_{i,\lambda}$ commutes with the $E^{S_l}_{ab}$ for $l\geq i$ and with the $E^{S'}_{ab}$ for
    every $S'\subset\Lambda$ disjoint from $S_i$; the $e_{i,\lambda}$ commute with each other;
    $\langle e_{i,\lambda}F,G\rangle_{\mathrm F}=\langle F,e_{i,\lambda}G\rangle_{\mathrm F}$ for all $F,G$; and for every
    finite set $T$ of $N$-tuples $t=(t_0,\dots,t_{N-1})$ in $\mathbf Z^n$, every $d\colon\Lambda\to\mathbf Z_{\geq0}$ and
    every $F\in\mathcal K[z]$ of row degrees $d$, with
    \[
     \kappa=\sum_{i<N}\frac{n(n-1)}2\,\ell_p\Big(2\sum_{v\in S_i}d(v)+n-1\Big),\qquad
     e_T=\sum_{t\in T}e_{0,t_0}e_{1,t_1}\cdots e_{N-1,t_{N-1}},
    \]
    $p^\kappa e_TF$ has coefficients in $\mathcal R$ if $F$ has, and $v_{\mathrm F}(p^\kappa e_TF)\geq0$ if
    $v_{\mathrm F}(F)\geq0$.
  \end{enumerate}
\item Let $L$ be a field and $S\subset\Lambda$.
  \begin{enumerate}
  \item The four rules of Lemma~\ref{pd-symbols-residue-count-inputs-segment-clauses}(1) hold for $\operatorname{val}_S$
    and $\lambda_S$ on $L[z]$. For $G\in L[z]$ and $\theta,M\in\mathbf Z$, $\theta\leq\operatorname{val}_S(G)$ if and only
    if $\theta\leq\mu(n)$ for every weight $\mu$ of $U_S(G)$, and $\lambda_S(G)\leq M$ if and only if $\mu(1)\leq M$ for
    every weight $\mu$ of $U_S(G)$.
  \item Let $L$ have characteristic $0$, and let $e_\mu$, $\mu\in\mathbf Z^n$, be isotypic projectors of weight $\mu$ for
    $S$. Then for $G\in L[z]$ and $\theta,M\in\mathbf Z$, $\theta\leq\operatorname{val}_S(G)$ if and only if
    $\theta\leq\mu(n)$ for every $\mu$ with $e_\mu G\neq0$, and $\lambda_S(G)\leq M$ if and only if $\mu(1)\leq M$ for every
    $\mu$ with $e_\mu G\neq0$.
  \end{enumerate}
\item Let $\mathcal R'$ be a discrete valuation ring, $N\geq0$ an integer with a lettering of $\{0,\dots,N-1\}$, and
  $B\colon\{0,\dots,N-1\}\to\mathbf Z_{\geq0}$. Then the path generators $\Psi(\mu,I)\in\mathcal R'[z]$, for the paths
  $\mu$ of rank $n$ for $B$ and the standard fillings $I$ of $\mu$ (Definition~\ref{pd-symbols-path-generators}), are
  linearly independent over $\mathcal R'$ and span the $\mathcal R'$-module of the polynomials of row degrees $B$.
\item Let $p\geq5$ be a prime. For integers $P\geq1$ and $e,\rho\geq0$ put $c(P,e)=(Pe)!/((P!)^e\,e!)$ and
  $d(P,e,\rho)=\prod_{y=1}^\rho(1+Pe/y)$.
  \begin{enumerate}
  \item For integers $t,\rho,e\geq0$ with $\rho<p^{t+1}$,
    $(\rho+p^{t+1}e)!=\rho!\,e!\,(p^{t+1}!)^e\,c(p^{t+1},e)\,d(p^{t+1},e,\rho)$.
  \item Let $\sigma$ be a finite set and $t,\rho\colon\sigma\to\mathbf Z_{\geq0}$ with $\rho(k)<P_k=p^{t(k)+1}$ for every $k$,
    and for $e\colon\sigma\to\mathbf Z_{\geq0}$ put
    \[
     u(e)=\prod_kc(P_k,e(k))\,d(P_k,e(k),\rho(k)),\qquad f(e)=\prod_k\frac{(\rho(k)+P_ke(k))!}{\rho(k)!\,e(k)!\,(P_k!)^{e(k)}}.
    \]
    Then $u(e)\neq0$ and $v_p(u(e))=0$ for every $e$. For $\epsilon=\pm1$ and every integer $c\geq0$ there is a
    polynomial $Q$ over $\mathbf Q$ in variables $x_k$, $k\in\sigma$, of total degree at most $4c$ with
    $v_p(u(e)^\epsilon-Q(e))\geq c$ for every $e$; and there are finitely many $j\colon\sigma\to\mathbf Z_{\geq0}$ with
    $\sum_kj(k)\leq4c$ and $q_j\in\mathbf Q$ with $v_p(q_j)\geq0$ such that
    $v_p\big(f(e)^\epsilon-\sum_jq_j\prod_k\binom{e(k)}{j(k)}\big)\geq c$ for every $e$.
  \item Let $L$ be a field, $S\subset\Lambda$, $c\geq0$ an integer, $T$ a finite set of monomials of total degree at most
    $4c$, $q_m\in L$ for $m\in T$, $F\in L[z]$ and $G=\sum_{m\in T}q_mz^mD_mF$. Then
    $\operatorname{val}_S(F)\leq\operatorname{val}_S(G)+4c$ and $\lambda_S(G)\leq\lambda_S(F)+4c$.
  \end{enumerate}
\item
  \begin{enumerate}
  \item Let $L$ be a field, $S\subset\Lambda$, $F\in L[z]$ and $z^\alpha$, $z^\beta$ monomials, and put
    $G=z^\alpha D_\beta F$. Write $|\alpha|$ for the total degree of $z^\alpha$, and $\alpha_{\mathrm v}$ and
    $\alpha_{\mathrm c}$ for its degrees in the variables of the vector letters and of the covector letters in $S$, and
    likewise for $\beta$. Then
    \[
     \operatorname{val}_S(F)\leq\operatorname{val}_S(G)+|\alpha|+|\beta|,\qquad\lambda_S(G)\leq\lambda_S(F)+|\alpha|+|\beta|,
    \]
    \[
     \operatorname{val}_S(F)\leq\operatorname{val}_S(G)+\alpha_{\mathrm c}+\beta_{\mathrm v},\qquad
     \lambda_S(G)\leq\lambda_S(F)+\alpha_{\mathrm v}+\beta_{\mathrm c}.
    \]
  \item Let $p\geq5$ be a prime, $t\geq0$ an integer, $\sigma$ a finite set and $\rho,M\colon\sigma\to\mathbf Z_{\geq0}$ with
    $\rho(k)<p^t$ for every $k$. Then for every integer $c\geq0$ there are finitely many $l\colon\sigma\to\mathbf Z_{\geq0}$,
    each with $\sum_kl(k)\leq p^{-t}\sum_kM(k)+c$ and $M(k)\leq\rho(k)+p^tl(k)$ for every $k$, and integers $\gamma_l$,
    such that $p^c$ divides $\prod_k\binom{\rho(k)+p^te(k)}{M(k)}-\sum_l\gamma_l\prod_k\binom{e(k)}{l(k)}$ for every
    $e\colon\sigma\to\mathbf Z_{\geq0}$.
  \item Let $R$ be a commutative ring, $R[x]$ the polynomial ring in variables $x_k$, $k\in\sigma$, and
    $P\colon\sigma\to\mathbf Z_{\geq1}$, $\rho\colon\sigma\to\mathbf Z_{\geq0}$ and $k\in\sigma$. Suppose that for every
    monomial $x^m$ of $F\in R[x]$ and every $k'\in\sigma$ the remainder of $m(k')$ modulo $P(k')$ is $\rho(k')$. Then
    every monomial of $\partial F/\partial x_k$ has the same remainders at the $k'\neq k$, and at $k$ the remainder
    $\rho(k)-1$ if $\rho(k)\geq1$ and $P(k)-1$ if $\rho(k)=0$.
  \end{enumerate}
\end{enumerate}
\end{lemma}
\leanref{Formalism/SharpRestriction/GrowthBounds/TableLineDegreeFive.lean}{Formalism.MainNotions.tableLineDegreeFive.admitted}
\leanref{Formalism/SharpRestriction/GrowthBounds/TableLineDegreeFive.lean}{Formalism.MainNotions.tableLineDegreeFive.adm11}
\leanref{Formalism/SharpRestriction/GrowthBounds/TableLineDegreeFive.lean}{Formalism.MainNotions.tableLineDegreeFive.adm12}
\leanref{Formalism/SharpRestriction/GrowthBounds/TableLineDegreeFive.lean}{Formalism.MainNotions.tableLineDegreeFive.adm14}
\leanref{Formalism/SharpRestriction/GrowthBounds/TableLineDegreeFive.lean}{Formalism.MainNotions.tableLineDegreeFive.adm23}
\leanref{Formalism/SharpRestriction/GrowthBounds/TableLineDegreeFive.lean}{Formalism.MainNotions.tableLineDegreeFive.FrobeniusExpansionCongruence}
\leanref{Formalism/SharpRestriction/GrowthBounds/CovariantDegrees.lean}{Formalism.MainNotions.CovariantDegrees.FilteredProjectorCorollaryField}
\leanref{Formalism/SharpRestriction/GrowthBounds/CovariantDegrees.lean}{Formalism.MainNotions.CovariantDegrees.ValuationLemmaField}
\leanref{Formalism/SharpRestriction/GrowthBounds/CovariantDegrees.lean}{Formalism.MainNotions.CovariantDegrees.ShaLengthVersionTwo}
\leanref{Formalism/SharpRestriction/GrowthBounds/ConclusionA2.lean}{Formalism.MainNotions.CovariantDegrees.HyperplaneRank}
\leanref{Formalism/SharpRestriction/GrowthBounds/CovariantDegrees.lean}{Formalism.MainNotions.CovariantDegrees.UnitProductSecondLemma}
\leanref{Formalism/SharpRestriction/GrowthBounds/CovariantDegrees.lean}{Formalism.MainNotions.CovariantDegrees.FactorialUnitDecomposition}
\leanref{Formalism/SharpRestriction/GrowthBounds/CovariantDegrees.lean}{Formalism.MainNotions.CovariantDegrees.DegreeTruncationLemma}
\leanref{Formalism/SharpRestriction/GrowthBounds/CovariantDegrees.lean}{Formalism.MainNotions.CovariantDegrees.DerivativeLemma}
\leanref{Formalism/SharpRestriction/GrowthBounds/CovariantDegrees.lean}{Formalism.MainNotions.CovariantDegrees.DerivativeValuationBound}
\leanref{Formalism/SharpRestriction/GrowthBounds/CovariantDegrees.lean}{Formalism.MainNotions.CovariantDegrees.DerivativeDigitExpansion}
\leanref{Formalism/SharpRestriction/GrowthBounds/CovariantDegrees.lean}{Formalism.MainNotions.CovariantDegrees.TrueDerivativeDigitClass}
\leanref{Formalism/SharpRestriction/GrowthBounds/GLWeightTheory.lean}{Formalism.MainNotions.G.GLWeightTheory.true_derivative_digit_class}
\leanref{Formalism/SharpRestriction/GrowthBounds/ObjPol.lean}{Formalism.MainNotions.G.glOps}
\leanref{Formalism/SharpRestriction/GrowthBounds/ObjPol.lean}{Formalism.MainNotions.G.IsIsoProj}

\begin{proof}
(1) Number the elements of $\Lambda$ as $0,\dots,|\Lambda|-1$, and use
Lemmas~\ref{pd-kappa-defects-segment-mean-letters-operators}, \ref{pd-kappa-defects-segment-mean-letters-lattices}
and~\ref{pd-kappa-defects-segment-mean-letters-polynomials} with $\mathcal O=\mathcal R$, $K=\mathcal K$ and $r=n$.
Their hypotheses hold. A positive integer is $p^aq$ with $a\geq0$ and $q$ prime to $p$, and $q$ is a unit of $\mathcal R$:
otherwise $q$ and $p$ would lie in the maximal ideal, and so would $1=xp+yq$ for suitable integers $x,y$. As $p\neq0$
in $\mathcal R$, positive integers are nonzero in $\mathcal R$, and $\mathcal K$ has characteristic $0$. For the same
reason $v(m!)=v_p(m!)$ for every monomial $z^m$. So the $F$ with $v_{\mathrm F}(F)\geq0$ form the module
$\mathcal B_{\mathcal R}$ of Lemma~\ref{pd-kappa-defects-segment-mean-letters-lattices}: such an $F$ is the sum of its
terms $F_mz^m$, which lie in $\mathcal B_{\mathcal R}$, and the coefficient at $z^m$ of an element of
$\mathcal B_{\mathcal R}$ has valuation at least $-\frac12v_p(m!)$. For $i<N$ let $E^{(i)}=(E^{S_i}_{ab})_{a,b\leq n}$.
It is a $\mathfrak{gl}_n$-family on $\mathcal K[z]$ (Definition~\ref{pd-kappa-defects-segment-mean-letters-families}) by
Lemma~\ref{pd-kappa-defects-segment-mean-letters-operators}(1). Its highest weight vectors and isotypic subspaces
(Definition~\ref{pd-kappa-defects-segment-mean-operators}(1)) are those of the statement for $S_i$, and its isotypic
projectors for $\{\lambda\}$ are the isotypic projectors of weight $\lambda$ for $S_i$.

(1)(a) By Lemma~\ref{pd-kappa-defects-segment-mean-letters-operators}(1), $E^{S_i}_{aa}z^m=\mathrm{wt}_{S_i}(m)(a)z^m$,
so $H^{(k)}_a=\binom{E^{S_i}_{aa}}k$; and $E^{(k)}_{ab}=(E^{S_i}_{ab})^k/k!$ for $a\neq b$, as $\mathcal K$ has
characteristic $0$ (Definition~\ref{pd-symbols-path-weights}(2)). So (a) is
Lemma~\ref{pd-kappa-defects-segment-mean-letters-lattices}(1) for $S=S_i$.

(1)(b) Let $i\leq l<N$, $\lambda,\mu\in\mathbf Z^n$, $a,b\leq n$, and $S'\subset\Lambda$ disjoint from $S_i$. By
Lemma~\ref{pd-kappa-defects-segment-mean-letters-operators}(2), $E^{S'}_{ab}$ commutes with every $E^{S_i}_{cd}$; so
$E^{S'}_{ab}$, and also $E^{S_i}_{ab}$, map every isotypic subspace of $E^{(i)}$ into itself
(Lemma~\ref{pd-kappa-defects-segment-mean-letters-span}(3)). By part (5) of that lemma, $e_{i,\lambda}$ commutes with
$E^{S'}_{ab}$ and with $E^{S_i}_{ab}$. As $S_i\subset S_l$, $E^{S_l}_{ab}=E^{S_i}_{ab}+E^{S_l\smallsetminus S_i}_{ab}$
(Lemma~\ref{pd-kappa-defects-segment-mean-letters-operators}(2)), so $e_{i,\lambda}$ commutes with $E^{S_l}_{ab}$. Hence
$e_{i,\lambda}$ maps every isotypic subspace of $E^{(l)}$ into itself, by part (3), and commutes with $e_{l,\mu}$, by
part (5); for $i>l$ exchange the two. The Fock pairing of
Definition~\ref{pd-weight-lattices}(2) is $\langle\ ,\ \rangle_{\mathrm F}$, so
$\langle e_{i,\lambda}F,G\rangle_{\mathrm F}=\langle F,e_{i,\lambda}G\rangle_{\mathrm F}$ by
Lemma~\ref{pd-kappa-defects-segment-mean-letters-polynomials} for $S=S_i$ and $T=\{\lambda\}$.

For the last assertion let $M$ be $\mathcal R[z]$ or $\mathcal B_{\mathcal R}$, let $\mathrm{Pol}(d)$ be the space of
the polynomials of row degrees $d$, and let $F\in M\cap\mathrm{Pol}(d)$. For $i<N$ put
$c_i=\frac{n(n-1)}2\ell_p(2\sum_{v\in S_i}d(v)+n-1)$, so that $\kappa=\sum_ic_i$. By
Lemma~\ref{pd-kappa-defects-segment-mean-letters-polynomials} for $S=S_i$, $T=\{\lambda\}$ and $B=d$, each
$e_{i,\lambda}$ maps $\mathrm{Pol}(d)$ into itself, and $p^{c_i}e_{i,\lambda}(G)\in M$ for $G\in M\cap\mathrm{Pol}(d)$.
Let $t\in T$, put $G_N=F$ and $G_i=p^{c_i}e_{i,t_i}(G_{i+1})$ for $i=N-1,\dots,0$. By descending induction on $i$,
$G_i\in M\cap\mathrm{Pol}(d)$; and $G_0=p^\kappa e_{0,t_0}e_{1,t_1}\cdots e_{N-1,t_{N-1}}F$, as the $e_{i,\lambda}$ are
$\mathcal K$-linear. Summing over $t\in T$ gives $p^\kappa e_TF\in M$.
% Lean: (1) is filtered_projector_corollary_field = polynomial_isotypic_lattice_corollary D h2 h13 =
%   polynomial_isotypic_lattice_of_parts (PolynomialIsotypicLatticeOfParts.lean, 97 lines) with hCR =
%   complete_reducibility_of_semisimple (Lemma pd-kappa-defects-segment-mean-letters-span(4)), hOrth =
%   contravariant_orthogonality_holds, hDom = highest_weight_dominant_holds, hRes = subrep_restriction_of_reducibility,
%   hPol = polynomial_representation_holds (Lemma pd-kappa-defects-segment-mean-letters-operators), hLat =
%   integral_formula_lattice_holds (Lemma pd-kappa-defects-segment-mean-letters-lattices; hchar is
%   dvr_fraction_field_char_zero through natCast_isUnit_of_coprime, and InBB is coefficientwise by
%   mem_fock_lattice_iff_coeff) and h13 = isotypic_projector_theorem_field (Lemma
%   pd-kappa-defects-segment-mean-letters-integral). (a) is hL1. In (b), hdis and hle are isotypic_projector_commutes
%   with isotypic_stable_commuting, hP2 and hP3, hcomm' is isotypic_projector_commutes again, and the adjointness is
%   isotypic_projector_selfAdjoint with hOrth and hB2. The integrality is operator_list_lattice_scaling (the G_i) with
%   isotypic_projector_lattice_scaling for each e_(i, lambda) (IsotypicProjectorLatticeScaling.lean, 130 lines): Pol d
%   is stable, the restriction is an isotypic projector of the restricted family, its weight is dominant, ellOf and
%   ellWt are at most degIn d S (polynomial_weight_difference_le), and h13 is applied to latticeMultidegreePart L d.
%   This is the argument of CountingStatementReduction.subOps_int for T = {lambda} and phi = id, which Lemma
%   pd-kappa-defects-segment-mean-letters-polynomials follows. The Lean states (1) for n >= 0; RootWeights/ holds
%   about 8000 lines behind it, written out in the Lemmas pd-kappa-defects-segment-mean-letters-* and
%   pd-gain-weyl-invariant-denominator.

(2)(a) The argument for the four rules in the proof of
Lemma~\ref{pd-symbols-residue-count-inputs-segment-clauses}(1) applies over the field $L$ itself. The two equivalences
are Lemma~\ref{pd-symbols-path-valuation}, with the elements of $\Lambda$ numbered $0,\dots,|\Lambda|-1$.

(2)(b) Write $E_{ab}$ for $E^S_{ab}$. They satisfy $E_{ab}E_{cd}-E_{cd}E_{ab}=\delta_{bc}E_{ad}-\delta_{da}E_{cb}$, and
$E_{aa}$ multiplies a monomial $z^m$ by $\mathrm{wt}_S(m)(a)$. As $L$ has characteristic $0$, $E^{(k)}_{ab}=E_{ab}^k/k!$
for $a\neq b$ (Definition~\ref{pd-symbols-path-weights}(2)) and $H^{(k)}_a=E_{aa}(E_{aa}-1)\cdots(E_{aa}-k+1)/k!$, while
$E_{ab}=E^{(1)}_{ab}$ and $E_{aa}=H^{(1)}_a$. So $U_S(G)$ is the span $W_G$ of the $u(G)$, with $u$ a composite of the
$E_{ab}$. It is finite-dimensional, as the $E_{ab}$ preserve the row degrees, and spanned by weight vectors. By Weyl's
theorem \cite[\S6.3]{Hum72}, applied to $\mathfrak{sl}_n$, since the $E_{aa}$ act by integers on weight vectors, every
finite-dimensional subspace of $L[z]$ that is stable under the $E_{ab}$ and spanned by weight vectors is the sum of
the spans of the $u(v)$, with $v$ a highest weight vector in it and $u$ a composite of the $E_{ab}$. Applied to the
$W_x$, this shows that $L[z]$ is the sum of the isotypic subspaces $L[z]^{(\mu)}_S$. These are stable under the $E_{ab}$;
so each $e_\mu$, which is the identity on $L[z]^{(\mu)}_S$ and $0$ on the others, commutes with the $E_{ab}$, and
$e_\mu$ vanishes on $W_G$ when $e_\mu G=0$. Finally, for a highest weight vector $v$ of weight $\mu$ and a composite
$u$ of the $E_{ab}$, $u(v)$ is a sum of weight vectors of weights $\nu$ with $\nu\preceq\mu$, so with $\mu(n)\leq\nu(n)$
and $\nu(1)\leq\mu(1)$ \cite[\S\S20--21]{Hum72}.

By (2)(a), $\theta\leq\operatorname{val}_S(G)$ if and only if $\theta\leq\nu(n)$ for every weight $\nu$ of $W_G$. Suppose
this, and let $e_\mu G\neq0$. Write $G=\sum_{\mu'}G_{\mu'}$ with $G_{\mu'}$ in the span of the $u(v)$, $v\in W_G$ a highest
weight vector of weight $\mu'$. Then $e_\mu G=G_\mu\neq0$, so $W_G$ contains a highest weight vector of weight $\mu$,
and $\mu$ is a weight of $W_G$; hence $\theta\leq\mu(n)$. Conversely, suppose that $\theta\leq\mu(n)$ whenever
$e_\mu G\neq0$, let $\nu$ be a weight of $W_G$ and $x\in W_G$ a nonzero element all of whose monomials have $S$-weight
$\nu$. Write $x=\sum_\mu x_\mu$ with $x_\mu\in L[z]^{(\mu)}_S$. Then $x_\mu=e_\mu x$, which is $0$ if $e_\mu G=0$. So
every $x_\mu\neq0$ has $\theta\leq\mu(n)$, and its monomials have $S$-weights $\nu'$ with $\mu(n)\leq\nu'(n)$. A monomial
of $x$ occurs in some $x_\mu$, so $\theta\leq\nu(n)$. The statement for $\lambda_S$ is proved in the same way, with the
condition $\nu(1)\leq M$.

(3) This is Lemma~\ref{pd-symbols-path-basis} with $R=\mathcal R'$.

(4)(a) Put $P=p^{t+1}$. Then $(\rho+Pe)!=(Pe)!\prod_{y=1}^\rho(Pe+y)=(Pe)!\,\rho!\,d(P,e,\rho)$, and
$(Pe)!=e!\,(P!)^e\,c(P,e)$ by the definition of $c$.

(4)(b) We use Definition~\ref{pd-symbols-test-chains-projection-clauses-mahler}. Let $\epsilon=\pm1$, and for $k\in\sigma$
put $f_k(m)=c(P_k,m)^\epsilon d(P_k,m,\rho(k))^\epsilon$. By
Lemma~\ref{pd-symbols-test-chains-projection-clauses-mahler-factorials}(3), applied with $t(k)$ and $\rho(k)<P_k$, the
numbers $c(P_k,m)$ and $d(P_k,m,\rho(k))$ are nonzero with $v_p=0$; so $u(e)\neq0$ and $v_p(u(e))=0$. By the same
part and Lemma~\ref{pd-symbols-test-chains-projection-clauses-mahler-rules}(4), each $f_k$ has decay $(\frac14,0)$, and
$u(e)^\epsilon=\prod_kf_k(e(k))$.

Let $c\geq0$ be an integer, $J$ the finite set of the $j\colon\sigma\to\mathbf Z_{\geq0}$ with $\sum_kj(k)\leq4c$, and
$q_j=\prod_ka_{j(k)}(f_k)$ for every $j\colon\sigma\to\mathbf Z_{\geq0}$, so that $v_p(q_j)\geq\frac14\sum_kj(k)$. We
show that
\[
 v_p\Big(\prod_kf_k(x(k))-\sum_{j\in J}q_j\prod_k\binom{x(k)}{j(k)}\Big)\geq c
\]
for every $x\colon\sigma\to\mathbf Z_{\geq0}$. Choose an integer $B>4c$ with $B>x(k)$ for every $k$. By
Lemma~\ref{pd-symbols-test-chains-projection-clauses-mahler-rules}(1),
$f_k(x(k))=\sum_{j<B}a_j(f_k)\binom{x(k)}j$, as $\binom{x(k)}j=0$ for $j>x(k)$. Multiplying out,
$\prod_kf_k(x(k))=\sum_jq_j\prod_k\binom{x(k)}{j(k)}$, the sum over the $j$ with values in $\{0,\dots,B-1\}$, which
include those in $J$. A term with $j\notin J$ has $\sum_kj(k)>4c$, so $v_p(q_j)>c$, and its binomial coefficients are
integers. This proves the bound.

For $j\geq0$ let $\binom xj=x(x-1)\cdots(x-j+1)/j!$, a polynomial of degree $j$ over $\mathbf Q$, and put
$Q=\sum_{j\in J}q_j\prod_k\binom{x_k}{j(k)}$. It has total degree at most $4c$, and
$Q(e)=\sum_{j\in J}q_j\prod_k\binom{e(k)}{j(k)}$; so $v_p(u(e)^\epsilon-Q(e))\geq c$ for every $e$. By (4)(a), applied
for each $k$, $f(e)=u(e)$; so the $j\in J$, with the $q_j$, which satisfy $v_p(q_j)\geq0$, give the last assertion.
% Lean: (4)(a) is fact_lemma_2; (4)(b) is lemma_five_two_first_holds = unit_function_mahler_decay_of D (hDec_of with
%   hdU_proved, hMul_proved, hUnit_of with hdU_proved): the units are asm_unit_prod; f_k is asm_f, with decay by
%   asm_f_dec and u(e)^epsilon = prod_k f_k(e(k)) by asm_prod_eq; the bound is asm_multi (asm_newton_big,
%   boundedDegreeExponents = J, asm_q_ge); Q is built from asm_bin (asm_bin_deg, asm_bin_eval); f(e) = u(e) is asm_ratio.
%   MemBoundedDegreeExponentsIff.lean, 353 lines.

(5)(a) Let $a\leq n$, and let $M_v$ be the matrix by which the row $z_v$ is multiplied in $P\mapsto P_{S,a}$
(Definition~\ref{pd-symbols-path-generators}(1)): $\tau_a(\eta)$ if $v\in S$ is a vector letter, ${}^t\tau_a(\eta^{-1})$ if
$v\in S$ is a covector letter, and the identity if $v\notin S$. Let $w_{v,c}$ be further variables. The coefficient of
$w^\beta$ in $(z+w)^m$ is $\prod_{v,c}\binom{m(v,c)}{\beta(v,c)}z^{m-\beta}$, so $F(z+w)=\sum_\beta(D_\beta F)(z)\,w^\beta$.
Substituting $z_vM_v$ for $z_v$ in this identity gives
\[
 F_{S,a}\big((z_v+w_vM_v^{-1})_v\big)=F\big((z_vM_v+w_v)_v\big)=\sum_\beta(D_\beta F)_{S,a}\,w^\beta,
\]
where the first equality is the definition of $F_{S,a}$; so $(D_\beta F)_{S,a}$ is the coefficient of $w^\beta$ on the
left. The entries of $M_v^{-1}$ are those of
$\tau_a(\eta^{-1})=\mathfrak gD_a(\eta^{-1})\mathfrak g^{-1}$, which involve only the exponents $0$ and $-1$ of $\eta$, if
$v\in S$ is a vector letter, those of ${}^t\tau_a(\eta)$, with the exponents $0$ and $1$, if $v\in S$ is a covector
letter, and constants if $v\notin S$. Expanding the left side, every exponent of $\eta$ in $(D_\beta F)_{S,a}$ is
therefore at least the least exponent in $F_{S,a}$ minus $\beta_{\mathrm v}$, and at most the largest exponent in
$F_{S,a}$ plus $\beta_{\mathrm c}$. In the same way the exponents in $(z^\alpha)_{S,a}$ lie between
$-\alpha_{\mathrm c}$ and $\alpha_{\mathrm v}$, and the exponents of a product lie between the sums of the least and of the
largest exponents of the factors. With $a=n$ and $a=1$ this gives
$\operatorname{val}_S(G)\geq\operatorname{val}_S(F)-\alpha_{\mathrm c}-\beta_{\mathrm v}$ and
$\lambda_S(G)\leq\lambda_S(F)+\alpha_{\mathrm v}+\beta_{\mathrm c}$, the second pair of inequalities. The first pair
follows, as $\alpha_{\mathrm v},\alpha_{\mathrm c}\leq|\alpha|$ and $\beta_{\mathrm v},\beta_{\mathrm c}\leq|\beta|$.

(4)(c) For $m\in T$, the argument of (5)(a) gives that every exponent of $\eta$ in $(q_mz^mD_mF)_{S,n}$ is at least
$\operatorname{val}_S(F)-|m|\geq\operatorname{val}_S(F)-4c$, and every exponent in $(q_mz^mD_mF)_{S,1}$ is at most
$\lambda_S(F)+4c$. Every exponent of $\eta$ in a sum occurs in one of the summands.

(5)(b) Put $P=p^t$, and for integers $\rho',M',l\geq0$ let $A(\rho',M',l)$ be the coefficient of $X^{M'}$ in
$(1+X)^{\rho'}\big((1+X)^P-1\big)^l\in\mathbf Z[X]$. Writing $(1+X)^{\rho'+Pe'}=(1+X)^{\rho'}\big(((1+X)^P-1)+1\big)^{e'}$ and
taking the coefficient of $X^{M'}$ gives $\binom{\rho'+Pe'}{M'}=\sum_{l=0}^{M'}A(\rho',M',l)\binom{e'}l$, as $X$ divides
$(1+X)^P-1$, so that $A(\rho',M',l)=0$ for $l>M'$. By degrees, $A(\rho',M',l)=0$ if $\rho'+Pl<M'$. As $p$ divides
$\binom Pi$ for $0<i<P$, $(1+X)^P-1=X^P+ph$ with $h\in\mathbf Z[X]$, and by induction on $l$ the coefficient of $X^{M'}$
in $f\cdot(X^P+ph)^l$ is divisible by $p^{\max(l-\lfloor M'/P\rfloor,0)}$ for every $f\in\mathbf Z[X]$; so
$p^{\max(l-\lfloor M'/P\rfloor,0)}$ divides $A(\rho',M',l)$. Multiplying over $k\in\sigma$,
\[
 \prod_k\binom{\rho(k)+Pe(k)}{M(k)}=\sum_l\prod_kA\big(\rho(k),M(k),l(k)\big)\binom{e(k)}{l(k)},
\]
the sum over the $l$ with $l(k)\leq M(k)$ for all $k$. Take for the finite set of the statement the $l$ with
$M(k)\leq\rho(k)+Pl(k)$ for all $k$ and $\sum_k\max(l(k)-\lfloor M(k)/P\rfloor,0)<c$, and
$\gamma_l=\prod_kA(\rho(k),M(k),l(k))$. For each of the other $l$ in the sum, some factor $A$ vanishes or the product of
the $A$ is divisible by $p^c$. For the kept $l$, $\sum_kl(k)<c+\sum_k\lfloor M(k)/P\rfloor\leq p^{-t}\sum_kM(k)+c$.

(5)(c) It suffices to treat a monomial $x^m$ of $F$. If $m(k)=0$, its derivative is $0$. Otherwise the derivative is
$m(k)x^{m'}$, where $m'$ agrees with $m$ except that $m'(k)=m(k)-1$. If $\rho(k)\geq1$, then $m(k)-1$ has the remainder
$\rho(k)-1$ modulo $P(k)$. If $\rho(k)=0$, then $m(k)=P(k)q$ with $q\geq1$, and $m(k)-1=P(k)(q-1)+P(k)-1$ has the
remainder $P(k)-1$.
% Lean: (5)(c) is TrueDerivativeDigitClass, proved as GLWeightTheory.true_derivative_digit_class (pderiv_monomial,
%   sub_one_mod_eq_ite). (3) is HyperplaneRank, proved as BinomialPolynomialIdentities.HyperplaneRank =
%   admissible_polynomial_basis_prov. (2)(a) is ValuationsAdditive, proved as valuationsAdditive_D7 =
%   valuations_additive_of_generated_weight_description D generated_module_weight_lemma_part_b: the four rules are
%   mul_valS_mul, mul_valS_add, mul_lam1S_mul and mul_lam1S_add, the argument of Lemma
%   pd-symbols-residue-count-inputs-segment-clauses(1), and the last two clauses are the weight description
%   generated_module_weight_lemma_part_b, Lemma pd-symbols-path-valuation.
% Lean: (1) is filtered_projector_corollary_field = polynomial_isotypic_lattice_corollary D
%   (weight_condition_from_integral_form D) (isotypic_projector_theorem_field D), see the comment after (1).
%   (4)(b) is lemma_five_two_first_holds = unit_function_mahler_decay_of D with hDec_of, hMul_proved and hUnit_of
%   (Lemmas pd-symbols-test-chains-projection-clauses-mahler-rules and -mahler-factorials), with asm_multi and
%   asm_ratio (MemBoundedDegreeExponentsIff.lean, MahPadicGeOfDvd.lean, PadicGeMulNat.lean, about 850 lines).
%   (4)(a) is fact_lemma_2.
% Lean, the other parts: (2)(b) is ShaLengthVersionTwo = generated_weight_lemma_of_reducibility D h2
%   (valuationsAdditive_D7 D) = generated_weight_lemma_char_zero with StableSubspaceCompleteReducibility (Weyl),
%   HighestWeightModuleWeightBounds, PolynomialOperatorsGLRelations, MonomialsAreWeightVectors, PolynomialLocallyFinite,
%   DividedPowerFormulaCharZero and BinomialFormulaCharZero (polynomial_generated_eq_enveloping,
%   isotypic_component_weight_mem for the first direction, weight_type_of_isotypic_projections with
%   isotypic_projection_commutes_monoid for the second). (5)(a) is derivative_valuation_bound and (4)(c) is
%   degree_truncation (DerivativeValuationBounds.lean): the Lean's filtrations laurentCoeffFiltration of the twisted
%   polynomials, flt_Th_twist (the Taylor expansion twisted by the inverse matrices), flt_divDer_mem,
%   flt_twist_monomial_mem and flt_op_mem. (5)(b) is derivative_digit_expansion (binomialExpansionCoeff,
%   choose_shift_binomial_expansion, binomial_expansion_coeff_padic_dvd through coeff_binomial_power_divisibility,
%   binomial_expansion_coeff_zero_of_deg); it does not use rho(k) < p^t.
% Lean: in tableLineDegreeFive.admitted these are the clauses h11 = adm11 = FilteredProjectorCorollaryField (1),
%   h12 = adm12 = ValuationLemmaField = ValuationsAdditive and ShaLengthVersionTwo (2), h14 = adm14 = HyperplaneRank (3),
%   h23 = adm23 = UnitProductSecondLemma = FactorialUnitDecomposition and DegreeTruncationLemma (4), and
%   h24 = FrobeniusExpansionCongruence = DerivativeLemma = DerivativeValuationBound, DerivativeDigitExpansion and
%   TrueDerivativeDigitClass (5), taken apart in TestSideCounting.plane_gram_lemmas_of_admitted. They are discharged
%   through tool_statements_hold (ExtremalProjectorTools): (1) is GLWeightTheory.filtered_projector_corollary_field =
%   polynomial_isotypic_lattice_corollary with isotypic_projector_theorem_field; (2)(b) is
%   generated_weight_lemma_of_reducibility (adm2, weight_condition_from_integral_form); (4) is lemma_five_two_first_holds
%   (unit_function_mahler_decay_of with hDec_of, hMul_proved, hUnit_of) and degree_truncation
%   (DerivativeValuationBounds.lean); (5)(a) is derivative_valuation_bound (DerivativeValuationBounds.lean, 861 lines
%   with degree_truncation) and (5)(b) derivative_digit_expansion (BinomialExpansionCoeffZeroOfDeg.lean, 191 lines).
%   These proofs draw on RootWeights/ (126 files, about 18700 lines, shared with the other clauses of Conclusion1).
%   The Lean states (1) for every n >= 0 (n = 0 is trivial) and the parts over fields for any finite type of letters.
\end{proof}

\begin{definition}\label{pd-symbols-test-chains-projection-properties}
Let $p\geq5$ be a prime, $\mathcal K$ a field of characteristic $0$ and $v\colon\mathcal K\to\mathbf Q\cup\{+\infty\}$ a
valuation with $v(p)=1$. An element $x\in\mathcal K$ is \emph{integral} if $v(x)\geq0$, and a polynomial over $\mathcal K$
is integral if its coefficients are. For a finite set $\Lambda$ each of whose elements is a vector letter or a covector
letter and an integer $r\geq1$, let $\mathcal K[z]$ be the polynomial ring in the $z_{y,a}$, $y\in\Lambda$,
$1\leq a\leq r$, with $\operatorname{val}_S$ for $S\subset\Lambda$, the Fock pairing $\langle\ ,\ \rangle_{\mathrm F}$ and
the divided derivatives $D_\beta$ of Lemma~\ref{pd-symbols-test-chains-projection-clauses} for $n=r$. Put
$v_{\mathrm F}(F)=\min_\mu(v(F_\mu)+v(\mu!)/2)$, the minimum over the monomials $z^\mu$ of $F$, with
$v_{\mathrm F}(0)=+\infty$; for $d\colon\Lambda\to\mathbf Z_{\geq0}$ let $\mathrm{Pol}(d)$ be the space of the polynomials
of row degrees $d$; and put $\ell(x)=\log_p(x+r)+1$ for real $x\geq0$. For $S\subset\Lambda$, $\theta\in\mathbf Z$ and
$d$, a \emph{Fock projector} for $(S,\theta,d)$ is a $\mathcal K$-linear endomorphism $\pi$ of $\mathcal K[z]$ such that
for every $F\in\mathrm{Pol}(d)$: $\pi F\in\mathrm{Pol}(d)$ and $\operatorname{val}_S(\pi F)\geq\theta$; $\pi F=F$ if
$\operatorname{val}_S(F)\geq\theta$; and $\langle F-\pi F,G\rangle_{\mathrm F}=0$ for every $G\in\mathrm{Pol}(d)$ with
$\operatorname{val}_S(G)\geq\theta$. The field $\mathcal K$ with $v$ has the \emph{lifting properties} at $p$ if the
following five properties hold for all such $\Lambda$ and $r$.
\begin{itemize}
\item[(V)] Let $S\subset\Lambda$ and $F\in\mathcal K[z]$.
  \begin{enumerate}
  \item $\operatorname{val}_S(z_{y,a})\geq0$ if $y$ is a vector letter, and $\operatorname{val}_S(z_{y,a})\geq-1$ if it is
    a covector letter.
  \item Let $y\in\Lambda$ and let $z^\beta$ be a monomial in $z_{y,1},\dots,z_{y,r}$ of degree $|\beta|$. Then
    $\operatorname{val}_S(D_\beta F)\geq\operatorname{val}_S(F)-|\beta|$ if $y$ is a vector letter, and
    $\operatorname{val}_S(D_\beta F)\geq\operatorname{val}_S(F)$ if it is a covector letter.
  \item If $S\subset S'\subset\Lambda$ and $S'\smallsetminus S$ consists of vector letters, then
    $\operatorname{val}_S(F)\leq\operatorname{val}_{S'}(F)$.
  \item If $S'\subset S$, $S\smallsetminus S'$ consists of vector letters and $F\in\mathrm{Pol}(d)$, then
    $\operatorname{val}_{S'}(F)\geq\operatorname{val}_S(F)-\sum_{y\in S\smallsetminus S'}d(y)$.
  \end{enumerate}
\item[(S)] Let $\Lambda=\{0,\dots,N-1\}$ with a lettering, $d\colon\Lambda\to\mathbf Z_{\geq0}$, and
  $S_i=\{0,\dots,i-1\}$ for $0\leq i\leq N$. For integers $c$ and $h$ let $\omega(c,h)$ be $1$ if $c\geq1$ and $h=r$, $-1$
  if $c\leq0$ and $h<r$, and $0$ otherwise. Then there are finitely many integral $\Psi_b\in\mathrm{Pol}(d)$, linearly
  independent over $\mathcal K$, each a product $\Psi_b=\prod_kQ_{b,k}$ of finitely many polynomials with integers
  $c_{b,k}$, sets $O_{b,k}\subset\Lambda$ and integers $h_{b,k}(i)\geq0$ for $0\leq i\leq N$, such that, with
  $\pi_b(i)=\sum_k\omega(c_{b,k},h_{b,k}(i))$:
  \begin{enumerate}
  \item every integral $F\in\mathrm{Pol}(d)$ is $\sum_bc_b\Psi_b$ with integral $c_b$; if moreover
    $\theta\colon\{0,\dots,N\}\to\mathbf Z$ and $\operatorname{val}_{S_i}(F)\geq\theta(i)$ for every $i$, the $c_b$ can be
    chosen with $\theta(i)\leq\pi_b(i)$ for every $i$ whenever $c_b\neq0$;
  \item $Q_{b,k}$ has degree $1$ in the row of each letter of $O_{b,k}$ and $0$ in the other rows, and every letter $i$
    lies in $d(i)$ of the sets $O_{b,k}$, $k=1,2,\dots$;
  \item $\operatorname{val}_{S_i}(Q_{b,k})\geq\omega(c_{b,k},h_{b,k}(i))$, and if $c_{b,k}\geq1$, then $h_{b,k}(i)$ is at
    most the number of vector letters of $O_{b,k}$ smaller than $i$.
  \end{enumerate}
\item[(P)] Let $d\colon\Lambda\to\mathbf Z_{\geq0}$, $S\subset\Lambda$ and $\theta\in\mathbf Z$. There is a Fock projector
  for $(S,\theta,d)$, and every such $\pi$ satisfies, for every $F\in\mathrm{Pol}(d)$:
  \begin{enumerate}
  \item $\pi(z_{y,a}\,\partial F/\partial z_{y,a})=z_{y,a}\,\partial(\pi F)/\partial z_{y,a}$ for every vector letter
    $y\notin S$ and every $a\leq r$;
  \item $\pi\pi'F=\pi'\pi F$ for every Fock projector $\pi'$ for $(S',\theta',d)$ with $S\subset S'$ or $S\cap S'=\emptyset$;
  \item $v_{\mathrm F}(\pi F)\geq v_{\mathrm F}(F)-r^2\ell\big(2\sum_yd(y)\big)$.
  \end{enumerate}
\item[(D)] Let $\tau$ be a finite set, $r'\geq0$ an integer and $g\colon\tau\to\mathbf Z_{\geq0}$. There are elements
  $c_t(\rho,n)$ and $\gamma_t(\rho,\beta)$ of $\mathcal K$, for $t\in\tau$, $\rho,\beta\in\mathbf Z_{\geq0}^{r'}$ and
  $n\geq0$, with the following properties. If $\rho(a)<p^{g(t)}$ for every $a$, then $c_t(\rho,n)\neq0$, and for every
  $\beta$, with $|\beta|=\sum_a\beta(a)$, $v(\gamma_t(\rho,\beta))=0$ and
  \[
   \prod_a\big(\rho(a)+p^{g(t)}\beta(a)\big)!=c_t\big(\rho,|\beta|\big)\prod_a\beta(a)!\cdot\gamma_t(\rho,\beta).
  \]
  If $\rho=(\rho_t)_{t\in\tau}$ is a family with $\rho_t(a)<p^{g(t)}$ for all $t$ and $a$, $c\geq1$ is an integer and
  $\epsilon=\pm1$, there are finitely many $\mu\colon\tau\times\{1,\dots,r'\}\to\mathbf Z_{\geq0}$ with $\sum\mu\leq6c$ and
  integral $\lambda_\mu\in\mathcal K$ such that, for every $\beta\colon\tau\times\{1,\dots,r'\}\to\mathbf Z_{\geq0}$ and
  with $\beta_t=\beta(t,\cdot)$,
  \[
   v\Big(\Big(\prod_t\gamma_t(\rho_t,\beta_t)\Big)^\epsilon-\sum_\mu\lambda_\mu\prod_{t,a}\binom{\beta(t,a)}{\mu(t,a)}\Big)\geq c.
  \]
\item[(E)] Let $r',g\geq0$ and $c\geq1$ be integers and $\rho,M\in\mathbf Z_{\geq0}^{r'}$ with $\rho(a)<p^g$ for every $a$.
  In the polynomial ring $\mathcal K[x_1,\dots,x_{r'}]$, with the divided derivatives $D_\beta$ as above and $Q^{[P]}$ the
  polynomial $Q$ with each $x_a$ replaced by $x_a^P$, there are finitely many integral $e_t\in\mathcal K$ and
  $\rho'_t,\delta_t,\kappa_t\in\mathbf Z_{\geq0}^{r'}$ with $|\kappa_t|\leq|M|/p^g+c$ such that for every
  $\beta\in\mathbf Z_{\geq0}^{r'}$ every coefficient of
  \[
   D_M\big(x^{\rho+p^g\beta}\big)-\sum_te_t\,x^{\rho'_t}\big(x^{\delta_t}D_{\kappa_t}(x^\beta)\big)^{[p^g]}
  \]
  has valuation at least $c$.
\end{itemize}
\end{definition}
\leanref{Formalism/SharpRestriction/GrowthBounds/MonomialOrderSegments/LiftBoundOfHypotheses.lean}{Formalism.MainNotions.PlaneGram.V2}
\leanref{Formalism/SharpRestriction/GrowthBounds/MonomialOrderSegments/LiftBoundOfHypotheses.lean}{Formalism.MainNotions.PlaneGram.SPAN}
\leanref{Formalism/SharpRestriction/GrowthBounds/MonomialOrderSegments/LiftBoundOfHypotheses.lean}{Formalism.MainNotions.PlaneGram.IsFockProj}
\leanref{Formalism/SharpRestriction/GrowthBounds/MonomialOrderSegments/LiftBoundOfHypotheses.lean}{Formalism.MainNotions.PlaneGram.PROJ}
\leanref{Formalism/SharpRestriction/GrowthBounds/MonomialOrderSegments/LiftBoundOfHypotheses.lean}{Formalism.MainNotions.PlaneGram.DIG}
\leanref{Formalism/SharpRestriction/GrowthBounds/MonomialOrderSegments/LiftBoundOfHypotheses.lean}{Formalism.MainNotions.PlaneGram.DER}
\leanref{Formalism/SharpRestriction/GrowthBounds/MonomialOrderSegments/LiftBoundOfHypotheses.lean}{Formalism.MainNotions.PlaneGram.vBB}
\leanref{Formalism/SharpRestriction/GrowthBounds/MonomialOrderSegments/LiftBoundOfHypotheses.lean}{Formalism.MainNotions.PlaneGram.lg}

\begin{lemma}\label{pd-symbols-test-chains-projection-ring-columns}
Let $N\geq0$ and $r\geq1$ be integers with a lettering of $\{0,\dots,N-1\}$, $d\colon\{0,\dots,N-1\}\to\mathbf Z_{\geq0}$,
and $\mu$ a path of rank $r$ for $d$, with $h_i(c)$, $T_\mu$ and the column functions $\mathrm{Col}_{\mu,c}$ of
Definition~\ref{pd-symbols-path-generators}; a letter $i$ is \emph{used} for $c$ if $h_i(c)\neq h_{i+1}(c)$. For
integers $c$ and $h$ let $\omega(c,h)$ be $1$ if $c\geq1$ and $h=r$, $-1$ if $c\leq0$ and $h<r$, and $0$ otherwise, and
put $S_i=\{0,\dots,i-1\}$ for $0\leq i\leq N$. Let $c\in\mathbf Z$.
\begin{enumerate}
\item Over every commutative ring, for every $J\subset\{1,\dots,r\}$, every monomial of $\mathrm{Col}_{\mu,c}(J)$ has
  degree $1$ in the row $z_i$ of every letter $i$ used for $c$ and degree $0$ in the other rows.
\item If a letter is used for $c$, then $|c|\leq T_\mu$. Every letter $i$ is used for $d(i)$ of the integers $c'$ with
  $|c'|\leq T_\mu$.
\item If $c\geq1$, then for $0\leq i\leq N$, $h_i(c)$ is at most the number of the vector letters $j<i$ used for $c$.
\item $\mu_i(r)=\sum_{|c'|\leq T_\mu}\omega(c',h_i(c'))$ for $0\leq i\leq N$.
\item Let $p$ be a prime, $\mathcal R$ a complete discrete valuation ring in which $p$ is neither $0$ nor a unit, and
  $\mathcal K$ its fraction field. Let $J\subset\{1,\dots,r\}$ have $h_N(c)$ elements if $c\geq1$ and $r-h_N(c)$
  elements if $c\leq0$. Then $\operatorname{val}_{S_i}(\mathrm{Col}_{\mu,c}(J))\geq\omega(c,h_i(c))$ in $\mathcal K[z]$
  for $0\leq i\leq N$.
\end{enumerate}
\end{lemma}
\leanref{Formalism/SharpRestriction/GrowthBounds/ProjectionOfIsotypicDecomposition.lean}{Formalism.MainNotions.G.Joints.col_clause1}
\leanref{Formalism/SharpRestriction/GrowthBounds/ProjectionOfIsotypicDecomposition.lean}{Formalism.MainNotions.G.Joints.col_clause2}
\leanref{Formalism/SharpRestriction/GrowthBounds/ProjectionOfIsotypicDecomposition.lean}{Formalism.MainNotions.G.Joints.col_clause3}
\leanref{Formalism/SharpRestriction/GrowthBounds/ProjectionOfIsotypicDecomposition.lean}{Formalism.MainNotions.G.Joints.col_clause4}
\leanref{Formalism/SharpRestriction/GrowthBounds/ProjectionOfIsotypicDecomposition.lean}{Formalism.MainNotions.G.Joints.span_column_clauses_at}
\leanref{Formalism/SharpRestriction/GrowthBounds/MonomialOrderSegments/BrSpanG.lean}{Formalism.MainNotions.G.Joints.PlaneGramSpan.stage_identity}
\leanref{Formalism/SharpRestriction/GrowthBounds/MonomialOrderSegments/BrSpanG.lean}{Formalism.MainNotions.G.Joints.PlaneGramSpan.col_est}
\leanref{Formalism/SharpRestriction/GrowthBounds/MonomialOrderSegments/BrSpanG.lean}{Formalism.MainNotions.G.Joints.PlaneGramSpan.onePath}

\begin{proof}
For a letter $i$ let $(\alpha,\beta)=(\mu_i,\mu_{i+1})$ if $i$ is a vector letter and $(\alpha,\beta)=(\mu_{i+1},\mu_i)$
if it is a covector letter; so $\alpha(a)\leq\beta(a)$ for all $a$, $\beta(a+1)\leq\alpha(a)$ for $a<r$,
$\sum_a(\beta(a)-\alpha(a))=d(i)$, and $\beta$ is dominant. If $a<a'$, then $\beta(a')\leq\beta(a+1)\leq\alpha(a)$. As
$\alpha\leq\beta$ entrywise, the number of entries $\geq c$ of $\beta$ is that of $\alpha$ plus the number of the $a$ with
$\alpha(a)<c\leq\beta(a)$, and there is at most one such $a$: for two of them, $a<a'$, $c\leq\beta(a')\leq\alpha(a)<c$. So
$i$ is used for $c$ if and only if $\alpha(a)<c\leq\beta(a)$ for some $a$, and then the two counts differ by $1$.

(1) Let $F_j$ be the function after the first $j$ steps of the definition of $\mathrm{Col}_{\mu,c}$. By induction on
$j$, every monomial of every value of $F_j$ has degree $1$ in the row of each letter $i<j$ used for $c$ and $0$ in the
other rows: the values of $F_0$ are $1$ and $0$; the step at a letter not used for $c$ is the identity; and at a used
letter $i$ each value of $F\wedge z_i$ or $F\lrcorner z_i$ is a sum of terms $\pm z_{i,m}F(J')$.

(2) If $i$ is used for $c$, then $\alpha(a)<c\leq\beta(a)$ for some $a$, and
$-T_\mu\leq-|\alpha(a)|<c\leq|\beta(a)|\leq T_\mu$. The intervals $(\alpha(a),\beta(a)]$, $1\leq a\leq r$, are pairwise
disjoint, as $\beta(a')\leq\alpha(a)$ for $a<a'$, and they have $\sum_a(\beta(a)-\alpha(a))=d(i)$ integers in all, each of
absolute value at most $T_\mu$. So $i$ is used for $d(i)$ integers $c'$ with $|c'|\leq T_\mu$.

(3) By induction on $i$. As $\mu_0=0$ and $c\geq1$, $h_0(c)=0$. If $i$ is a vector letter, $h_{i+1}(c)$ is $h_i(c)$, or
$h_i(c)+1$ when $i$ is used for $c$, and then the count grows by $1$. If $i$ is a covector letter, $h_{i+1}(c)\leq h_i(c)$,
as $\mu_{i+1}\leq\mu_i$ entrywise.

(4) As $\mu_i$ is dominant, $h_i(c')=r$ if and only if $c'\leq\mu_i(r)$. So $\omega(c',h_i(c'))$ is $1$ if
$1\leq c'\leq\mu_i(r)$, $-1$ if $\mu_i(r)<c'\leq0$, and $0$ otherwise. As $|\mu_i(r)|\leq T_\mu$, the sum is $\mu_i(r)$.

(5) For $i=0$, $S_0=\emptyset$, and the substitution of Definition~\ref{pd-symbols-path-generators}(1) for the empty
set replaces no row; so $\operatorname{val}_\emptyset(G)$ is $0$ for $G\neq0$ and $+\infty$ for $G=0$. Also
$h_0(c)$ is $r$ if $c\leq0$ and $0$ if $c\geq1$, so $\omega(c,h_0(c))=0$. Let $i\geq1$.

If no letter is used for $c$, then $h_i(c)=h_0(c)$, so $\omega(c,h_i(c))=0$, and $\mathrm{Col}_{\mu,c}=F_0$ has the
values $1$ and $0$, with $\operatorname{val}_{S_i}\geq0$. Suppose that some letter is used for $c$. Let $e$ be $0$ if
$c\geq1$ and $1$ if $c\leq0$, put $c_0=1-e$, and let $\mu'_j\in\mathbf Z^r$, $0\leq j\leq N$, have the entries $1-e$ at
$a\leq h_j(c)$ and $-e$ at $a>h_j(c)$. Then $\mu'_0=0$, each $\mu'_j$ is dominant, and as $h_{j+1}(c)-h_j(c)$ is $0$ or
$1$ at a vector letter $j$ and $0$ or $-1$ at a covector letter, with $\pm1$ exactly when $j$ is used for $c$, $\mu'$ is a
path of rank $r$ for $d'$, where $d'(j)$ is $1$ if $j$ is used for $c$ and $0$ otherwise. For $c'\in\mathbf Z$, the
number of entries $\geq c'$ of $\mu'_j$ is $r$ if $c'+e\leq0$, $h_j(c)$ if $c'+e=1$, and $0$ otherwise. So the column
heights of $\mu'$ at $c_0$ are those of $\mu$ at $c$, no letter is used for $\mu'$ and any $c'\neq c_0$, and
$|c_0|\leq T_{\mu'}$ by (2) for $\mu'$. As $c_0\geq1$ if and only if $c\geq1$, and the steps depend only on the heights
and on whether $c\geq1$, $\mathrm{Col}_{\mu',c_0}=\mathrm{Col}_{\mu,c}$; and $\mathrm{Col}_{\mu',c'}=F_0$ for $c'\neq c_0$.
Let $I'(c_0)=J$ and $I'(c')=\emptyset$ for $c'\neq c_0$. This is a standard filling of $\mu'$: $|J|$ is $h'_N(c_0)$ or
$r-h'_N(c_0)$ as required; for $c'\neq c_0$, the height of $\mu'$ at $c'$ is $0$ if $c'\geq1$ and $r$ if $c'\leq0$; and the
sets $I'^*(c')$ are $\emptyset$ for $c'\geq1$ other than $c_0$, and $\{1,\dots,r\}$ for $c'\leq0$ other than $c_0$, so the
condition on the elements $\leq t$ holds. Hence $\Psi(\mu',I')=\mathrm{Col}_{\mu,c}(J)$, as $F_0(\emptyset)=1$.

By Lemma~\ref{pd-symbols-path-clauses}(1) over $\mathcal R$, for $d'$, $\theta(j)=\mu'_{j+1}(r)$ and
$M(j)=\mu'_{j+1}(1)$, the path generator $\Psi(\mu',I')$, which lies in the span described there, satisfies
$\operatorname{val}_{\{0,\dots,j\}}(\Psi(\mu',I'))\geq\mu'_{j+1}(r)$ for every letter $j$. The substitution that defines
$\operatorname{val}_S$ commutes with the inclusion $\mathcal R[z]\subset\mathcal K[z]$, which is injective on
coefficients, so the same holds in $\mathcal K[z]$. With $j=i-1$, $\mu'_i(r)$ is $1-e$ if $h_i(c)=r$ and $-e$ otherwise,
which is $\omega(c,h_i(c))$.
% Lean: (1) is col_clause1 (Kcol_hom, colStep_hom with wedgeR_hom and contrR_hom); (2) is col_clause3 (path_step,
%   mem_colRange) and col_clause2 (strip_count, with strip_lt, colH_strip, strip_le_one and colH_ne_iff); (3) is
%   col_clause4 = column_height_le_used_count; these are the clauses of ColumnDegreeBoundAtRing
%   (span_column_clauses_at), whose last clause col_clause5 the book does not use. (4) is stage_identity (colH_eq_iff,
%   sum_colVal). (5) is col_est: for i = 0 valS_empty_nonneg, the case without used letters Kcol_unused and colH_const,
%   and otherwise onePath P c (oneMu with e, hstrip_one, colH_onePath), the filling I', Psi_one and Kcol_congr, and the
%   second clause of PathSpanAtRing (third_degree_bound_gives_span_clause from A3at, Lemma pd-symbols-path-clauses(1)),
%   with StandardMonomialSpace.lift_Kcol_map and valS_map for the passage from O to its fraction field.
%   ProjectionOfIsotypicDecomposition.lean lines 452 to 830 and BrSpanG.lean, about 870 lines.
\end{proof}

\begin{lemma}\label{pd-symbols-test-chains-projection-ring-projectors}
Let $p$, $n'$, $\mathcal R$, $\mathcal K$ and $v$ be as in Lemma~\ref{pd-symbols-test-chains-projection-ring}, and let
$\Lambda$, $r\geq1$, $\mathcal K[z]$, $\operatorname{val}_S$, $\langle\ ,\ \rangle_{\mathrm F}$, $v_{\mathrm F}$ and
$\mathrm{Pol}(d)$ be as in Definition~\ref{pd-symbols-test-chains-projection-properties}. Let $S\subset\Lambda$, and for
$\mu\in\mathbf Z^r$ let $e_\mu$ be an isotypic projector of weight $\mu$ for $S$
(Lemma~\ref{pd-symbols-test-chains-projection-clauses}, with $n=r$). For $\theta\in\mathbf Z$ and $F\in\mathcal K[z]$
put $\pi_\theta F=\sum_\mu e_\mu F$, the sum over the $\mu$ with $\theta\leq\mu(r)$ and $e_\mu F\neq0$. Let
$d\colon\Lambda\to\mathbf Z_{\geq0}$ and $\theta\in\mathbf Z$.
\begin{enumerate}
\item The $e_\mu$ exist; for each $F$ only finitely many $e_\mu F$ are nonzero, $F=\sum_\mu e_\mu F$, and
  $\pi_\theta$ is a $\mathcal K$-linear endomorphism of $\mathcal K[z]$ and a Fock projector for $(S,\theta,d)$.
\item Every Fock projector $\pi$ for $(S,\theta,d)$ agrees with $\pi_\theta$ on $\mathrm{Pol}(d)$.
\item Let $S'\subset\Lambda$ with $S\subset S'$ or $S\cap S'=\emptyset$, and $\theta'\in\mathbf Z$, and let $\pi'_{\theta'}$
  be formed for $S'$ as $\pi_\theta$ is for $S$. Then $\pi_\theta$ commutes with the operators $E^{S'}_{ab}$ of
  Lemma~\ref{pd-symbols-test-chains-projection-clauses} and with $\pi'_{\theta'}$.
\item If $F\in\mathrm{Pol}(d)$ and $v_{\mathrm F}(F)\geq0$, then $v_{\mathrm F}(\pi_\theta F)\geq-\frac{r(r-1)}2\,
  \ell_p\big(2\sum_{y\in S}d(y)+r-1\big)$, with $\ell_p$ as in Lemma~\ref{pd-symbols-test-chains-projection-clauses}.
\end{enumerate}
\end{lemma}
\leanref{Formalism/SharpRestriction/GrowthBounds/ProjectionOfIsotypicDecomposition.lean}{Formalism.MainNotions.G.Joints.projection_at_of_isotypic}
\leanref{Formalism/SharpRestriction/GrowthBounds/ProjectionOfIsotypicDecomposition.lean}{Formalism.MainNotions.G.Joints.DecSys.projQ}
\leanref{Formalism/SharpRestriction/GrowthBounds/ProjectionOfIsotypicDecomposition.lean}{Formalism.MainNotions.G.Joints.DecSys.comm_projQ}
\leanref{Formalism/SharpRestriction/GrowthBounds/ProjectionOfIsotypicDecomposition.lean}{Formalism.MainNotions.G.Joints.e_mem_Pol}
\leanref{Formalism/SharpRestriction/GrowthBounds/ProjectionOfIsotypicDecomposition.lean}{Formalism.MainNotions.G.Joints.isFockProj_projQ}
\leanref{Formalism/SharpRestriction/GrowthBounds/ProjectionOfIsotypicDecomposition.lean}{Formalism.MainNotions.G.Joints.fock_nondeg}
\leanref{Formalism/SharpRestriction/GrowthBounds/ProjectionOfIsotypicDecomposition.lean}{Formalism.MainNotions.G.Joints.fockProj_unique}
\leanref{Formalism/SharpRestriction/GrowthBounds/ProjectionOfIsotypicDecomposition.lean}{Formalism.MainNotions.G.Joints.polIsoDecompB_of}
\leanref{Formalism/SharpRestriction/GrowthBounds/MonomialOrderSegments/SecondValuationOfAdmitted.lean}{Formalism.MainNotions.PlaneGram.br_proj_iff}

\begin{proof}
Number the elements of $\Lambda$ as $0,\dots,|\Lambda|-1$. As in the proof of
Lemma~\ref{pd-symbols-test-chains-projection-clauses}(1), $\mathcal K$ has characteristic $0$, the isotypic projectors
of weight $\mu$ for $S$ are the isotypic projectors of the $\mathfrak{gl}_r$-family $(E^S_{ab})$ for $\{\mu\}$, and the
Fock pairing is that of Lemma~\ref{pd-kappa-defects-segment-mean-letters-operators}(3).

(1) The $e_\mu$ exist by Lemma~\ref{pd-kappa-defects-segment-mean-letters-polynomials}. By
Lemma~\ref{pd-kappa-defects-segment-mean-letters-span}(4), $F=\sum_{\mu\in A}F_\mu$ for a finite set $A$ and
$F_\mu$ in the isotypic subspace of weight $\mu$; as $e_\mu$ is the identity on that subspace and $0$ on the others,
$e_\mu F=F_\mu$ for $\mu\in A$ and $e_\mu F=0$ for $\mu\notin A$. So $\pi_\theta F=\sum_{\mu\in A,\ \theta\leq\mu(r)}e_\mu F$
for every finite $A$ that contains the $\mu$ with $e_\mu F\neq0$; taking one such $A$ for two polynomials shows that
$\pi_\theta$ is $\mathcal K$-linear. Also $e_\nu e_\mu F=0$ for $\nu\neq\mu$ and $e_\mu e_\mu F=e_\mu F$, so
$e_\nu(\pi_\theta F)$ is $e_\nu F$ if $\theta\leq\nu(r)$ and $0$ otherwise.

Let $F\in\mathrm{Pol}(d)$. By Lemma~\ref{pd-kappa-defects-segment-mean-letters-polynomials}, each $e_\mu F$ lies in
$\mathrm{Pol}(d)$, and so does $\pi_\theta F$. By Lemma~\ref{pd-symbols-test-chains-projection-clauses}(2)(b), a
polynomial $G$ has $\operatorname{val}_S(G)\geq\theta$ if and only if $\theta\leq\nu(r)$ for every $\nu$ with
$e_\nu G\neq0$. Applied to $G=\pi_\theta F$, this gives $\operatorname{val}_S(\pi_\theta F)\geq\theta$. If
$\operatorname{val}_S(F)\geq\theta$, then every $\mu$ with $e_\mu F\neq0$ has $\theta\leq\mu(r)$, and
$\pi_\theta F=\sum_\mu e_\mu F=F$. By Lemma~\ref{pd-symbols-test-chains-projection-clauses}(1)(b), with $N=1$ and
$S_0=S$, $\langle e_\mu F,G\rangle_{\mathrm F}=\langle F,e_\mu G\rangle_{\mathrm F}$ for all $F,G$; summing over a finite
set of $\mu$ that serves for $F$ and for $G$ gives $\langle\pi_\theta F,G\rangle_{\mathrm F}=\langle F,\pi_\theta
G\rangle_{\mathrm F}$. So if $G\in\mathrm{Pol}(d)$ has $\operatorname{val}_S(G)\geq\theta$, then $\pi_\theta G=G$ and
$\langle F-\pi_\theta F,G\rangle_{\mathrm F}=\langle F,G\rangle_{\mathrm F}-\langle F,\pi_\theta G\rangle_{\mathrm F}=0$.
Hence $\pi_\theta$ is a Fock projector for $(S,\theta,d)$.

(2) Let $\pi$ be a Fock projector for $(S,\theta,d)$, and $F,G\in\mathrm{Pol}(d)$. As $\pi G$ lies in $\mathrm{Pol}(d)$
and has $\operatorname{val}_S(\pi G)\geq\theta$, $\langle F-\pi F,\pi G\rangle_{\mathrm F}=0$, so
$\langle\pi F,\pi G\rangle_{\mathrm F}=\langle F,\pi G\rangle_{\mathrm F}$; exchanging $F$ and $G$, and as the Fock
pairing is symmetric, $\langle\pi F,G\rangle_{\mathrm F}=\langle F,\pi G\rangle_{\mathrm F}$. Write $\pi_\theta$ as
$\pi_0$. As $\pi$ fixes $\pi_0F$ and $\pi_0$ fixes $\pi G$, both having $\operatorname{val}_S\geq\theta$,
\[
 \langle\pi_0F,G\rangle_{\mathrm F}=\langle\pi\pi_0F,G\rangle_{\mathrm F}=\langle\pi_0F,\pi G\rangle_{\mathrm F}
 =\langle F,\pi_0\pi G\rangle_{\mathrm F}=\langle F,\pi G\rangle_{\mathrm F}=\langle\pi F,G\rangle_{\mathrm F}.
\]
Take for $G$ a monomial $z^m$ of row degrees $d$: $\langle\pi_0F,z^m\rangle_{\mathrm F}=m!\,(\pi_0F)_m$, and $m!\neq0$ in
$\mathcal K$. So $\pi_0F$ and $\pi F$ have the same coefficients at the monomials of row degrees $d$, and both lie in
$\mathrm{Pol}(d)$; hence $\pi_0F=\pi F$.

(3) Let $e'_\nu$, $\nu\in\mathbf Z^r$, be the isotypic projectors for $S'$ from which $\pi'_{\theta'}$ is formed. If
$S\subset S'$, Lemma~\ref{pd-symbols-test-chains-projection-clauses}(1)(b) with $N=2$, $S_0=S$, $S_1=S'$, $e_{0,\mu}=e_\mu$
and $e_{1,\nu}=e'_\nu$ shows that each $e_\mu$ commutes with every $E^{S'}_{ab}$ and every $e'_\nu$. If
$S\cap S'=\emptyset$, the same part with $N=1$ shows that each $e_\mu$ commutes with every $E^{S'}_{ab}$; so $e_\mu$ maps
every isotypic subspace of $(E^{S'}_{ab})$ into itself, and commutes with every $e'_\nu$, by
Lemma~\ref{pd-kappa-defects-segment-mean-letters-span}(3) and~(5). Let $X$ be an endomorphism of $\mathcal K[z]$ that
commutes with every $e_\mu$, and $F\in\mathcal K[z]$ with the finite set $A$ of (1). Then
$XF=\sum_{\mu\in A}e_\mu XF$, so $A$ serves for $XF$, and $\pi_\theta XF=\sum_{\mu\in A,\theta\leq\mu(r)}Xe_\mu F=
X\pi_\theta F$. This applies to $X=E^{S'}_{ab}$ and to $X=e'_\nu$. In the same way $\pi'_{\theta'}$ commutes with every
endomorphism that commutes with all $e'_\nu$, which $\pi_\theta$ does.

(4) Let $T$ be the set of the $\mu$ with $\theta\leq\mu(r)$ and $e_\mu F\neq0$, a finite set by (1), so that
$\pi_\theta F=\sum_{\mu\in T}e_\mu F$. Lemma~\ref{pd-symbols-test-chains-projection-clauses}(1)(b) with $N=1$, $S_0=S$
and the set of the $1$-tuples $(\mu)$, $\mu\in T$, gives $v_{\mathrm F}(p^\kappa\pi_\theta F)\geq0$ with
$\kappa=\frac{r(r-1)}2\ell_p(2\sum_{y\in S}d(y)+r-1)$, and $v_{\mathrm F}(p^\kappa G)=\kappa+v_{\mathrm F}(G)$.
% Lean: projection_at_of_isotypic D hC hV (polIsoDecompB_of D h2) p O hS, through which
%   projection_of_integral_corollary gives FockProjectorAtRing p O. The Lean takes the e_mu from PolIsoDecompB (isoDecomp:
%   isoSub_indep and isoSub_top, with the complete reducibility adm2), and the book from Lemmas
%   pd-kappa-defects-segment-mean-letters-polynomials and -span(4), whose Lean proofs (subOps_exists through iso_indep and
%   subOps_top) run the same argument. (1) is decSys_of, DecSys.projQ (projFun_eq, cover2) and isFockProj_projQ, with
%   e_mem_Pol (polProj_glAct, isoSub_map_of_comm, comm_of_preserve), hV2 = ShaLengthVersionTwo (clause (2)(b)) and hself
%   (clause (1)(b), the self-adjointness); the book's orthogonality form of a Fock projector is br_proj_iff.
%   (2) is fockProj_unique with fock_nondeg. (3) is the key of projection_at_of_isotypic: for S in S' the clause hC with
%   N = 2 and the chain ![S, S'] (h1 and h3), for S and S' disjoint hcD and comm_of_preserve with isoSub_map_of_comm;
%   then DecSys.comm_projQ. (4) is hint (clause (1)(b) with N = 1) with vBB_pow_smul.
%   ProjectionOfIsotypicDecomposition.lean lines 1173 to 1700, about 530 lines.
\end{proof}

\begin{lemma}\label{pd-symbols-test-chains-projection-ring}
Let $p\geq5$ be a prime, $n'\geq2$ an even integer, $\mathcal R$ a complete discrete valuation ring of characteristic $0$
in which $p$ is not a unit and whose residue field contains a field of $p^{n'}$ elements, $\mathcal K$ its fraction field
and $v$ the valuation of $\mathcal K$ with $v(p)=1$. Then $\mathcal K$ with $v$ has the lifting properties at $p$
(Definition~\ref{pd-symbols-test-chains-projection-properties}).
\end{lemma}
\leanref{Formalism/SharpRestriction/GrowthBounds/DualFockLinear/Closed.lean}{Formalism.MainNotions.TestSideCounting.plane_gram_lemmas_of_admitted}
\leanref{Formalism/SharpRestriction/GrowthBounds/Lemma5aAt.lean}{Formalism.MainNotions.PlaneGramData.AdmAt}
\leanref{Formalism/SharpRestriction/GrowthBounds/MonomialOrderSegments/SecondValuationOfAdmitted.lean}{Formalism.MainNotions.PlaneGram.second_valuation_of_admitted}
\leanref{Formalism/SharpRestriction/GrowthBounds/MonomialOrderSegments/BrSpan.lean}{Formalism.MainNotions.PlaneGram.brSPAN}
\leanref{Formalism/SharpRestriction/GrowthBounds/MonomialOrderSegments/SecondValuationOfAdmitted.lean}{Formalism.MainNotions.PlaneGram.brPROJ}
\leanref{Formalism/SharpRestriction/GrowthBounds/MonomialOrderSegments/UnitLevelOne.lean}{Formalism.MainNotions.PlaneGram.brDIG}
\leanref{Formalism/SharpRestriction/GrowthBounds/MonomialOrderSegments/BrRenExpandMonomial.lean}{Formalism.MainNotions.PlaneGram.brDER}

\begin{proof}
A positive integer is $p^aq$ with $q$ prime to $p$, and $q$ is a unit of $\mathcal R$, as in the proof of
Lemma~\ref{pd-symbols-test-chains-projection-clauses}(1); so $v(x)=v_p(x)$ for $x\in\mathbf Q$. We use $c(P,e)$ and
$d(P,e,\rho)$ of Lemma~\ref{pd-symbols-test-chains-projection-clauses}(4). Then $c(P,0)=1$ and $d(P,0,\rho)=1$, and
$c(1,e)=e!/e!=1$ and $d(1,e,0)=1$, an empty product.

(D) Let $\tau$, $r'$ and $g$ be as in (D). For $t\in\tau$, $\rho,\beta\in\mathbf Z_{\geq0}^{r'}$ and $n\geq0$ put
$P_t=p^{g(t)}$,
\[
 c_t(\rho,n)=\prod_a\rho(a)!\cdot(P_t!)^n,\qquad\gamma_t(\rho,\beta)=\prod_ac\big(P_t,\beta(a)\big)\,d\big(P_t,\beta(a),\rho(a)\big).
\]
Let $\rho(a)<P_t$ for every $a$; then $c_t(\rho,n)\neq0$. If $g(t)=0$, then $P_t=1$ and $\rho=0$, so $\gamma_t(\rho,\beta)=1$
and the identity of (D) is $\prod_a\beta(a)!=\prod_a\beta(a)!$. If $g(t)\geq1$,
Lemma~\ref{pd-symbols-test-chains-projection-clauses}(4)(a), with $g(t)-1$ in place of $t$, gives
$(\rho(a)+P_t\beta(a))!=\rho(a)!\,\beta(a)!\,(P_t!)^{\beta(a)}c(P_t,\beta(a))\,d(P_t,\beta(a),\rho(a))$, and the product over
$a$ is the identity; and (4)(b), for the set $\{1,\dots,r'\}$ with the constant function $g(t)-1$, gives
$v(\gamma_t(\rho,\beta))=v_p(\gamma_t(\rho,\beta))=0$.

Let $\rho=(\rho_t)$ with $\rho_t(a)<P_t$, $c\geq1$ and $\epsilon=\pm1$. Apply (4)(b) of that lemma to
$\sigma=\tau\times\{1,\dots,r'\}$, with $t(t,a)=\max(g(t)-1,0)$ and $\rho(t,a)=\rho_t(a)$, which is less than
$p^{t(t,a)+1}$. It gives finitely many $j\colon\sigma\to\mathbf Z_{\geq0}$ with $\sum j\leq4c$ and $q_j\in\mathbf Q$ with
$v_p(q_j)\geq0$ such that $v_p(f(e)^\epsilon-\sum_jq_j\prod_k\binom{e(k)}{j(k)})\geq c$ for every $e\colon\sigma\to\mathbf
Z_{\geq0}$, where $f(e)=\prod_kc(p^{t(k)+1},e(k))\,d(p^{t(k)+1},e(k),\rho(k))$ by (4)(a). Let
$\beta\colon\sigma\to\mathbf Z_{\geq0}$, and let $\beta^0$ agree with $\beta$ at the $(t,a)$ with $g(t)\geq1$ and vanish at
the others. At the $(t,a)$ with $g(t)=0$ the factor of $f(\beta^0)$ is $c(p,0)d(p,0,0)=1$, as is the factor
$c(1,\beta(t,a))d(1,\beta(t,a),0)$ of $\gamma_t$; so $f(\beta^0)=\prod_t\gamma_t(\rho_t,\beta_t)$. As $\binom0{j(k)}=0$ for
$j(k)\geq1$, $\sum_jq_j\prod_k\binom{\beta^0(k)}{j(k)}=\sum_{j\in J_0}q_j\prod_k\binom{\beta(k)}{j(k)}$, where $J_0$ is the set
of the $j$ that vanish at the $(t,a)$ with $g(t)=0$. So the $\mu=j\in J_0$, with $\lambda_\mu=q_j$, which are integral,
satisfy (D), as $\sum\mu\leq4c\leq6c$.

(V) We use the substitution $F\mapsto F_{S,r}$ of Definition~\ref{pd-symbols-path-generators}(1), with $n=r$; the
entries of $\tau_r(\eta)=\mathfrak gD_r(\eta)\mathfrak g^{-1}$ are polynomials in $\eta$, and so are those of
$\eta\tau_r(\eta)^{-1}=\eta\,\tau_r(\eta^{-1})$. (a) The substitution fixes $1$, so $\operatorname{val}_S(1)=0$. The third
inequality of Lemma~\ref{pd-symbols-test-chains-projection-clauses}(5)(a), with $F=1$, $z^\alpha=z_{y,a}$ and $z^\beta=1$,
gives $0\leq\operatorname{val}_S(z_{y,a})+\alpha_{\mathrm c}$, where $\alpha_{\mathrm c}$ is $1$ if $y$ is a covector
letter in $S$ and $0$ otherwise. (b) The same inequality, with $z^\alpha=1$, gives $\operatorname{val}_S(F)\leq
\operatorname{val}_S(D_\beta F)+\beta_{\mathrm v}$, where $\beta_{\mathrm v}$ is at most $|\beta|$ if $y$ is a vector
letter and $0$ if it is a covector letter. (c) The substitution for $S$ does not touch the rows of $S'\smallsetminus S$,
so $F_{S',r}$ is obtained from $F_{S,r}$ by replacing each row $z_y$, $y\in S'\smallsetminus S$, by $z_y\tau_r(\eta)$. As
these entries are polynomials in $\eta$, the least exponent of $\eta$ does not decrease. (d) Let $V=S\smallsetminus S'$
and $D=\sum_{y\in V}d(y)$. Replacing in $F_{S,r}$ each row $z_y$, $y\in V$, by $\eta z_y\tau_r(\eta)^{-1}$ turns
$z_y\tau_r(\eta)$ into $\eta z_y$, so it gives $\eta^DF_{S',r}$, as $F$ has degree $d(y)$ in the row $z_y$. As the entries
of $\eta\tau_r(\eta)^{-1}$ are polynomials in $\eta$, $\operatorname{val}_{S'}(F)+D\geq\operatorname{val}_S(F)$.

(E) Let $r'$, $g$, $c$, $\rho$ and $M$ be as in (E), and $P=p^g$.
Lemma~\ref{pd-symbols-test-chains-projection-clauses}(5)(b), with $t=g$, $\sigma=\{1,\dots,r'\}$ and this $c$, gives
finitely many $l\in\mathbf Z_{\geq0}^{r'}$ with $|l|\leq|M|/P+c$ and $M\leq\rho+Pl$, and integers $\gamma_l$, such that
$p^c$ divides $\prod_a\binom{\rho(a)+P\beta(a)}{M(a)}-\sum_l\gamma_l\binom\beta l$ for every $\beta$, where
$\binom\beta l=\prod_a\binom{\beta(a)}{l(a)}$. For each $l$ write $\rho+Pl-M=\rho'_l+P\delta_l$ with $0\leq\rho'_l(a)<P$,
and take $e_l=\gamma_l$, $\rho'_l$, $\delta_l$ and $\kappa_l=l$. As $D_l(x^\beta)=\binom\beta lx^{\beta-l}$, vanishing
unless $l\leq\beta$, and $\rho'_l+P(\delta_l+\beta-l)=\rho+P\beta-M$ when $l\leq\beta$,
\[
 D_M\big(x^{\rho+P\beta}\big)-\sum_le_lx^{\rho'_l}\big(x^{\delta_l}D_l(x^\beta)\big)^{[P]}
 =\Big(\prod_a\binom{\rho(a)+P\beta(a)}{M(a)}-\sum_l\gamma_l\binom\beta l\Big)x^{\rho+P\beta-M},
\]
both sides being $0$ unless $M\leq\rho+P\beta$. The coefficient is $p^c$ times an integer, of valuation at least $c$.

(S) Let $\Lambda=\{0,\dots,N-1\}$ with a lettering, $d$, $S_i$ and $\omega$ be as in (S). Let $B$ be the set of the
pairs $b=(\mu,I)$ of a path $\mu$ of rank $r$ for $d$ and a standard filling $I$ of $\mu$
(Definition~\ref{pd-symbols-path-generators}). By Lemma~\ref{pd-symbols-test-chains-projection-clauses}(3), the path
generators $\Psi(\mu,I)\in\mathcal R[z]$, $(\mu,I)\in B$, are linearly independent over $\mathcal R$ and span the
$\mathcal R$-module of the polynomials of row degrees $d$. That module is generated by the finitely many monomials of row
degrees $d$, and $\mathcal R$ is noetherian, so $B$ is finite. Let $\Psi_b\in\mathcal K[z]$ be the image of
$\Psi(\mu,I)$; it is integral and lies in $\mathrm{Pol}(d)$, and the $\Psi_b$ are linearly independent over
$\mathcal K$, as a relation over $\mathcal K$ becomes one over $\mathcal R$ after multiplication by a common
denominator. For $b=(\mu,I)$ let $k$ run over the integers $c$ with $|c|\leq T_\mu$, and put
$Q_{b,c}=\mathrm{Col}_{\mu,c}(I(c))$, $c_{b,c}=c$, $h_{b,c}(i)=h_i(c)$ and $O_{b,c}$ the set of the letters used for $c$.
Then $\Psi_b=\prod_cQ_{b,c}$, and $\pi_b(i)=\sum_c\omega(c,h_i(c))=\mu_i(r)$ by
Lemma~\ref{pd-symbols-test-chains-projection-ring-columns}(4).

(1) Let $F\in\mathrm{Pol}(d)$ be integral. It is the image of some $F_0\in\mathcal R[z]$ of row degrees $d$, so
$F_0=\sum_bl_b\Psi(\mu,I)$ with $l_b\in\mathcal R$, and $F=\sum_bl_b\Psi_b$. Suppose moreover that
$\operatorname{val}_{S_i}(F)\geq\theta(i)$ for $0\leq i\leq N$. As $\mathcal R[z]\subset\mathcal K[z]$ is injective on
coefficients, $\operatorname{val}_S$ and $\lambda_S$ are the same for $F_0$ and $F$. For each letter $i$, $\lambda_{S_{i+1}}(F_0)$
is an integer or $-\infty$; let $M(i)$ be an integer at least $\lambda_{S_{i+1}}(F_0)$. By
Lemma~\ref{pd-symbols-path-clauses}(1) over $\mathcal R$, with $i\mapsto\theta(i+1)$ and $M$, $F_0$ is an
$\mathcal R$-combination of the $\Psi(\mu,I)$ with $\theta(i+1)\leq\mu_{i+1}(r)$ for every letter $i$. Let the $c_b$ be
its coefficients. If $c_b\neq0$, then $\theta(i)\leq\mu_i(r)=\pi_b(i)$ for $1\leq i\leq N$; and $F\neq0$, by the linear
independence, so $\operatorname{val}_{S_0}(F)=\operatorname{val}_\emptyset(F)=0$, as in the proof of
Lemma~\ref{pd-symbols-test-chains-projection-ring-columns}(5), and $\theta(0)\leq0=\mu_0(r)=\pi_b(0)$.

(2) This is Lemma~\ref{pd-symbols-test-chains-projection-ring-columns}(1) and~(2).

(3) The first assertion is Lemma~\ref{pd-symbols-test-chains-projection-ring-columns}(5) with $J=I(c)$, which has the
size required there as $I$ is a standard filling; the second is part (3) of that lemma.

(P) Let $\Lambda$, $r$, $d$, $S$ and $\theta$ be as in (P), and $\pi_\theta$ as in
Lemma~\ref{pd-symbols-test-chains-projection-ring-projectors}. By part (1) of that lemma, $\pi_\theta$ is a Fock
projector for $(S,\theta,d)$. Let $\pi$ be one, and $F\in\mathrm{Pol}(d)$; by part (2), $\pi F=\pi_\theta F$.

(a) Let $y\notin S$ be a vector letter and $a\leq r$. The operator $z_{y,a}\,\partial/\partial z_{y,a}$ is $E^{\{y\}}_{aa}$,
it maps $\mathrm{Pol}(d)$ into itself, and $\{y\}\cap S=\emptyset$. By part (3),
$\pi(E^{\{y\}}_{aa}F)=\pi_\theta E^{\{y\}}_{aa}F=E^{\{y\}}_{aa}\pi_\theta F=E^{\{y\}}_{aa}\pi F$.

(b) Let $\pi'$ be a Fock projector for $(S',\theta',d)$, with $S\subset S'$ or $S\cap S'=\emptyset$, and $\pi'_{\theta'}$
formed for $S'$. As $\pi$, $\pi'$, $\pi_\theta$ and $\pi'_{\theta'}$ map $\mathrm{Pol}(d)$ into itself, parts (2) and (3)
give $\pi\pi'F=\pi_\theta\pi'_{\theta'}F=\pi'_{\theta'}\pi_\theta F=\pi'\pi F$.

(c) If $F=0$ there is nothing to show. Otherwise $t=v_{\mathrm F}(F)$ is a rational number. Put $N=\lceil-t\rceil$ and
$F'=p^NF$, so that $v_{\mathrm F}(F')=t+N\geq0$. By part (4),
$v_{\mathrm F}(\pi F)=v_{\mathrm F}(\pi_\theta F')-N\geq-\kappa-N>t-\kappa-1$, where
$\kappa=\frac{r(r-1)}2\ell_p(2D_S+r-1)$ and $D_S=\sum_{y\in S}d(y)$. Let $D=\sum_yd(y)\geq D_S$. As $p^{\ell_p(x)}\leq x$
for $x\geq1$ and $\ell_p(0)=0$, $\ell_p(2D_S+r-1)\leq\log_p(2D+r)$, which is $\geq0$; and $\frac{r(r-1)}2\leq r^2$ and
$1\leq r^2$. So $\kappa+1\leq r^2(\log_p(2D+r)+1)=r^2\ell(2D)$, and $v_{\mathrm F}(\pi F)\geq v_{\mathrm F}(F)-r^2\ell(2D)$.
% Lean, (S): brSPAN S h3 (BrSpan.lean) with h3 = PolynomialSpanAtRing O, whose three parts are
%   isotypic_coefficient_of_degree_bound (h14, clause (3): PathColumnBasisAtRing), third_degree_bound_gives_span_clause
%   (h15, Lemma pd-symbols-path-clauses(1): PathSpanAtRing) and span_column_clauses (ColumnDegreeBoundAtRing, Lemma
%   pd-symbols-test-chains-projection-ring-columns(1) to (3)). The index set is Idx N r LL d, finite by idx_finite
%   (linear independence in a finitely generated module over a noetherian ring), with li_K for the independence over
%   the fraction field and psiK_mem; Psi b is Psi (KK O) mu I, col b k = Kcol mu (cI b k) N (I (cI b k)) for the
%   numbering cI of colRange mu, cidx = cI, hts b k i = h i (cI b k) and occ b k the used letters. (1) is span_all
%   (exists_lift, comb_of_span) and span_mem (with M from withBot_le_int and A3at), the case i = 0 being valS_empty_le
%   and b.zero; (2) is the first two clauses of h3c; (3) is col_est (with valS_empty_nonneg for i = 0) and the fourth
%   clause of h3c; the identity for pi_b is stage_identity.
% Lean, (P): brPROJ S h4 (SecondValuationOfAdmitted.lean) with h4 = FockProjectorAtRing p O =
%   projection_of_integral_corollary D h11 h12 (Lemma pd-symbols-test-chains-projection-ring-projectors): the existence
%   is hpi0 through br_proj_iff, (a) is the commutation with glAct for the disjoint set {i} (br_glAct_single), (b) is
%   the commutation of Fock projectors, and (c) is the bound of FockProjectorAtRing, scaled by c = p^ceil(-t)
%   (br_exists_v, br_le_vBB_smul, br_v_inv) and br_exp_le with proj_exponent_le.
% Lean, (V): second_valuation_of_admitted h2 (SecondValuationOfAdmitted.lean) with h2 = ValuationStabilityAtRing:
%   (a) and (b) are CoordinateMultiplicationAtRing = valuation_clauses_of_lattice_index D hL61.1 (Joints.lean,
%   valuation_of_letters_clauses: the third clause of DerivativeValuationBound with F = 1 or al = 0, covDeg and vecDeg);
%   (c) and (d) are EnlargingLettersAtRing = twisted_valuation_second_sentence_at (ProjectionOfIsotypicDecomposition.lean:
%   valS_le_union with the substitution g of twistMat 1, and valS_le_diff with T * twistMat (-1), twistMat_neg_mul and
%   hom_scale_sum; PGe_aeval for the nonnegative exponents). The Lean proves (c) and (d) over any commutative ring.
% Lean, (E): brDER S h6 (BrRenExpandMonomial.lean: brRen_ddiv, brRen_term) with h6 = derivation_of_lattice_index D
%   hL61.2.1 = derivation_valuation_at_prime (ProjectionOfIsotypicDecomposition.lean): derTriple P M rho l is
%   (rho'_l, delta_l, l), der_term the identity of exponents, divDer_monoL, and geQ_vK_pow_mul for the valuation of
%   p^c times an element of O.
% Lean, (D): brDIG D S hG h5 (UnitLevelOne.lean): const and gam are c_t and gamDig (p^g t) rho beta (the product of
%   cUnit and dUnit); h5 = digit_identity_of_decomposition D hL52.1 (ProjectionOfIsotypicDecomposition.lean) gives the
%   identity and the units (gamDig_one_zero for g = 0, the clause FactorialUnitDecomposition otherwise); the
%   approximation is lemma_five_two_first_holds at tt k = g - 1 (natural subtraction) and rh, evaluated at ze beta,
%   with unit_at_zero, unit_level_one and the filter of T.
% Lean: PlaneGram.stmt destructures AdmAt p O (Lemma5aAt.lean), the six properties IsotypicValuationAtRing,
%   ValuationStabilityAtRing (CoordinateMultiplicationAtRing and EnlargingLettersAtRing), PolynomialSpanAtRing,
%   FockProjectorAtRing, FactorialDigitSplitAtRing and DerivativeCongruenceAtRing of O, and turns the last five into the
%   properties of KK O with vQ p S: (V) = V2 is second_valuation_of_admitted h2, (S) = SPAN is brSPAN S h3 (BrSpanG.lean),
%   (P) = PROJ is brPROJ S h4, (D) = DIG is brDIG D S hG h5 and (E) = DER is brDER S h6. AdmAt comes from the clauses
%   through plane_gram_lemmas_of_admitted: isotypic_multiplicity_of_reducibility_alt (h12), valuation_clauses_of_lattice_index
%   (h24, its first part) with twisted_valuation_second_sentence, isotypic_coefficient_of_degree_bound (h14),
%   third_degree_bound_gives_span_clause (h15) with span_column_clauses, projection_of_integral_corollary (h11 and h12),
%   digit_identity_of_decomposition (h23) and derivation_of_lattice_index (h24, its second part); Joints.lean and
%   ProjectionOfIsotypicDecomposition.lean.
% The isotypic decomposition that projection_of_integral_corollary uses (polIsoDecompB_of, from the complete
%   reducibility adm2) is the one of Lemmas pd-kappa-defects-segment-mean-letters-span(4) and -polynomials.
\end{proof}

\begin{lemma}\label{pd-symbols-test-chains-projection-fock-product}
Let $p\geq5$ be a prime and $\mathcal K$ with $v$ a field that satisfies (P) of
Definition~\ref{pd-symbols-test-chains-projection-properties} at $p$, and use the notation of that definition. Let
$\Lambda$ and $r$ be as there, $d\colon\Lambda\to\mathbf Z_{\geq0}$, and $(S_x,\theta_x)$, $x\in X$, a finite family of
subsets $S_x\subset\Lambda$ and integers $\theta_x$ such that for all $x,y\in X$, $S_x\subset S_y$, $S_y\subset S_x$ or
$S_x\cap S_y=\emptyset$. Then there is a $\mathcal K$-linear endomorphism $\Pi$ of $\mathcal K[z]$ such that for all
$F,G\in\mathrm{Pol}(d)$:
\begin{enumerate}
\item $\Pi F\in\mathrm{Pol}(d)$, and $\operatorname{val}_{S_x}(\Pi F)\geq\theta_x$ for every $x\in X$;
\item if $\operatorname{val}_{S_x}(G)\geq\theta_x$ for every $x\in X$, then
  $\langle G,\Pi F\rangle_{\mathrm F}=\langle G,F\rangle_{\mathrm F}$;
\item $v_{\mathrm F}(\Pi F)\geq v_{\mathrm F}(F)-|X|\,r^2\ell\big(2\sum_yd(y)\big)$;
\item if $S''\subset\Lambda$ and $\theta''\in\mathbf Z$ are such that for every $x\in X$, $S_x\subset S''$ or
  $S_x\cap S''=\emptyset$, and $\operatorname{val}_{S''}(F)\geq\theta''$, then $\operatorname{val}_{S''}(\Pi F)\geq\theta''$;
\item $\Pi(z_{y,a}\,\partial F/\partial z_{y,a})=z_{y,a}\,\partial(\Pi F)/\partial z_{y,a}$ for every vector letter $y$ that
  lies in no $S_x$ and every $a\leq r$. Hence every monomial of $\Pi(z^\mu)$, for a monomial $z^\mu\in\mathrm{Pol}(d)$, has
  the exponents of $z^\mu$ in the variables $z_{y,a}$ of these letters.
\end{enumerate}
\end{lemma}
\leanref{Formalism/SharpRestriction/GrowthBounds/MonomialOrderSegments/ProjKeepsLetter.lean}{Formalism.MainNotions.PlaneGram.prodProj}
\leanref{Formalism/SharpRestriction/GrowthBounds/MonomialOrderSegments/ProjKeepsLetter.lean}{Formalism.MainNotions.PlaneGram.prodProj_entries}
\leanref{Formalism/SharpRestriction/GrowthBounds/MonomialOrderSegments/ProjKeepsLetter.lean}{Formalism.MainNotions.PlaneGram.proj_preserves}
\leanref{Formalism/SharpRestriction/GrowthBounds/MonomialOrderSegments/ProjKeepsLetter.lean}{Formalism.MainNotions.PlaneGram.fock_proj_right}
\leanref{Formalism/SharpRestriction/GrowthBounds/MonomialOrderSegments/ProjKeepsLetter.lean}{Formalism.MainNotions.PlaneGram.bge_proj}

\begin{proof}
\emph{Two facts.} (a) Let $\pi$ be a Fock projector for $(S,\theta,d)$, and let $S''\subset\Lambda$ and $\theta''\in\mathbf Z$
with $S\subset S''$, $S''\subset S$ or $S\cap S''=\emptyset$. If $F\in\mathrm{Pol}(d)$ and $\operatorname{val}_{S''}(F)\geq\theta''$,
then $\operatorname{val}_{S''}(\pi F)\geq\theta''$. Indeed, (P) gives a Fock projector $\pi''$ for $(S'',\theta'',d)$, and
$\pi''F=F$. By (P)(b), applied to $\pi$ if $S\subset S''$ or $S\cap S''=\emptyset$ and to $\pi''$ if $S''\subset S$,
$\pi\pi''F=\pi''\pi F$. So $\pi F=\pi''(\pi F)$, and $\operatorname{val}_{S''}(\pi''(\pi F))\geq\theta''$ as
$\pi F\in\mathrm{Pol}(d)$. (b) Let $\pi$ be a Fock projector for $(S,\theta,d)$ and $F,G\in\mathrm{Pol}(d)$ with
$\operatorname{val}_S(G)\geq\theta$. Then $\langle F-\pi F,G\rangle_{\mathrm F}=0$, and as the Fock pairing is symmetric,
$\langle G,\pi F\rangle_{\mathrm F}=\langle G,F\rangle_{\mathrm F}$.

\emph{Induction on $|X|$.} For $X=\emptyset$ take the identity. Let $X=X'\cup\{x\}$ with $x\notin X'$, let $\Pi'$ be given
for $X'$, let $\pi_x$ be a Fock projector for $(S_x,\theta_x,d)$, which exists by (P), and put $\Pi=\pi_x\Pi'$. Let
$F,G\in\mathrm{Pol}(d)$. (1) As $\Pi'F\in\mathrm{Pol}(d)$, $\Pi F\in\mathrm{Pol}(d)$ and
$\operatorname{val}_{S_x}(\Pi F)\geq\theta_x$. For $y\in X'$, $\operatorname{val}_{S_y}(\Pi'F)\geq\theta_y$, and fact (a)
for $\pi_x$ and $S''=S_y$ gives $\operatorname{val}_{S_y}(\Pi F)\geq\theta_y$. (2) By fact (b) for $\pi_x$ and $\Pi'F$, and
by (2) for $\Pi'$, $\langle G,\Pi F\rangle_{\mathrm F}=\langle G,\Pi'F\rangle_{\mathrm F}=\langle G,F\rangle_{\mathrm F}$.
(3) By (P)(c) for $\pi_x$ and (3) for $\Pi'$, $v_{\mathrm F}(\Pi F)\geq v_{\mathrm F}(\Pi'F)-r^2\ell(2\sum_yd(y))\geq
v_{\mathrm F}(F)-(|X'|+1)r^2\ell(2\sum_yd(y))$. (4) By (4) for $\Pi'$, $\operatorname{val}_{S''}(\Pi'F)\geq\theta''$, and
fact (a) for $\pi_x$ gives $\operatorname{val}_{S''}(\Pi F)\geq\theta''$. (5) By (5) for $\Pi'$ and (P)(a) for $\pi_x$,
applied to $\Pi'F\in\mathrm{Pol}(d)$,
$\Pi(z_{y,a}\partial F/\partial z_{y,a})=\pi_x(z_{y,a}\partial(\Pi'F)/\partial z_{y,a})=z_{y,a}\partial(\Pi F)/\partial
z_{y,a}$.

For the last assertion, $z_{y,a}\partial/\partial z_{y,a}$ multiplies each monomial $z^{\mu'}$ by $\mu'(y,a)$. So (5)
for $F=z^\mu$ says that $\mu(y,a)\,\Pi(z^\mu)_{\mu'}=\mu'(y,a)\,\Pi(z^\mu)_{\mu'}$ for every $\mu'$, and as $\mathcal K$ has
characteristic $0$, $\mu'(y,a)=\mu(y,a)$ whenever $\Pi(z^\mu)_{\mu'}\neq0$.
% Lean: prodProj by Finset.induction_on (pi_x composed with Pr); fact (a) is proj_preserves, fact (b) is fock_proj_right
%   (with fock_comm), (3) is bge_proj, (4) and (5) are its last two clauses (proj_preserves, and the commutation clause
%   of PROJ), and the last assertion is prodProj_entries with proj_keeps_letter (euler_eq_dg).
\end{proof}

\begin{lemma}\label{pd-symbols-test-chains-projection-diagonal}
Let $p\geq5$ be a prime and $\mathcal K$ with $v$ a field that satisfies (V) of
Definition~\ref{pd-symbols-test-chains-projection-properties} at $p$, and use the notation of that definition, with
$\Lambda$ and $r$ as there and $S\subset\Lambda$. For a function $g$ from the exponent vectors
$\mu\colon\Lambda\times\{1,\dots,r\}\to\mathbf Z_{\geq0}$ to $\mathcal K$ put $\operatorname{diag}_gF=\sum_\mu g(\mu)F_\mu z^\mu$.
Let $n\geq0$ be an integer, $M$ a finite set of exponent vectors $m$ with $\sum_{y,a}m(y,a)\leq n$, $\lambda_m\in\mathcal K$
for $m\in M$, and
\[
 g(\mu)=\sum_{m\in M}\lambda_m\prod_{y,a}\binom{\mu(y,a)}{m(y,a)}.
\]
Then $\operatorname{val}_S(\operatorname{diag}_gF)\geq\operatorname{val}_S(F)-n$ for every $F\in\mathcal K[z]$.
\end{lemma}
\leanref{Formalism/SharpRestriction/GrowthBounds/MonomialOrderSegments/SlopeOfBinomCombGen.lean}{Formalism.MainNotions.PlaneGram.approx_of_trunc}
\leanref{Formalism/SharpRestriction/GrowthBounds/MonomialOrderSegments/SlopeOfBinomCombGen.lean}{Formalism.MainNotions.PlaneGram.CostV}
\leanref{Formalism/SharpRestriction/GrowthBounds/MonomialOrderSegments/SlopeOfBinomCombGen.lean}{Formalism.MainNotions.PlaneGram.costV_ok}
\leanref{Formalism/SharpRestriction/GrowthBounds/MonomialOrderSegments/AddMemSupportOfMemSupportCoefOut.lean}{Formalism.MainNotions.PlaneGram.vala_dg_choose}
\leanref{Formalism/SharpRestriction/GrowthBounds/MonomialOrderSegments/AddMemSupportOfMemSupportCoefOut.lean}{Formalism.MainNotions.PlaneGram.dg_choose_eq}

\begin{proof}
The map $F\mapsto F_{S,r}$ of Definition~\ref{pd-symbols-path-generators}(1) is a ring homomorphism that fixes the
constants, so $\operatorname{val}_S(FG)\geq\operatorname{val}_S(F)+\operatorname{val}_S(G)$,
$\operatorname{val}_S(F+G)\geq\min(\operatorname{val}_S(F),\operatorname{val}_S(G))$ and
$\operatorname{val}_S(\lambda F)\geq\operatorname{val}_S(F)$ for $\lambda\in\mathcal K$. Say that a function $g$
\emph{costs} $k$ if $\operatorname{val}_S(\operatorname{diag}_gF)\geq\operatorname{val}_S(F)-k$ for every $F$. As
$\operatorname{diag}_{gh}=\operatorname{diag}_g\operatorname{diag}_h$, $\operatorname{diag}_{g+h}=\operatorname{diag}_g+
\operatorname{diag}_h$ and $\operatorname{diag}_{\lambda g}=\lambda\operatorname{diag}_g$: the constant $1$ costs $0$; if
$g$ costs $k$ and $h$ costs $l$, then $gh$ costs $k+l$; if $g$ and $h$ cost $k$, so do $g+h$ and $\lambda g$; and if $g$
costs $k$, it costs every $l\geq k$.

Let $(y,a)\in\Lambda\times\{1,\dots,r\}$, $k\geq0$, and $e$ the exponent vector with $e(y,a)=k$ and $0$ elsewhere. As
$D_e(z^\mu)=\binom{\mu(y,a)}kz^{\mu-e}$ if $\mu(y,a)\geq k$ and $D_e(z^\mu)=0$ otherwise, the function
$\mu\mapsto\binom{\mu(y,a)}k$ gives $\operatorname{diag}F=z^eD_eF$. By (V)(a), $\operatorname{val}_S(z^e)\geq0$ if $y$ is a
vector letter and $\operatorname{val}_S(z^e)\geq-k$ if it is a covector letter, and by (V)(b),
$\operatorname{val}_S(D_eF)\geq\operatorname{val}_S(F)-k$, respectively $\geq\operatorname{val}_S(F)$. So this function
costs $k$. For $m\in M$ the product of these functions over the $(y,a)$, with $k=m(y,a)$, costs $\sum_{y,a}m(y,a)\leq n$,
so it costs $n$, and so does $g$.
% Lean: approx_of_trunc (approx_of_trunc_gen: g is binomComb T lam) with the cost structure CostOK (CostV sg {S}) of
%   costV_ok (one, mul, add, smul, mono through dg_mul_fun and the like) and cost_prod_choose_gen; the cost of one
%   binomial is vala_dg_choose (dg_choose_eq, vala_ddiv_letter' and vala_X from V2, vala_mul from valS_mul_le).
\end{proof}

\begin{lemma}\label{pd-symbols-test-chains-projection-lift}
Let $p\geq5$ be a prime, $\mathcal K$ a field of characteristic $0$, $v\colon\mathcal K\to\mathbf Q\cup\{+\infty\}$ a
valuation with $v(p)=1$, and $\mathcal R$ the ring of the $x\in\mathcal K$ with $v(x)\geq0$. Signed chain data over
$\mathcal R$ and the notation of Definition~\ref{pd-symbols-test-chains} are defined as there for this $\mathcal R$,
$\mathcal K$ and $v$: the definition uses of $\mathcal R$ only that its elements are the $x\in\mathcal K$ with
$v(x)\geq0$. Suppose that $\mathcal K$ with $v$ satisfies (V) and (P) of
Definition~\ref{pd-symbols-test-chains-projection-properties} at $p$. Let $r,J\geq1$ be integers, $\mathcal D$ a
signed chain datum of rank $r$ with $J$ levels over $\mathcal R$ whose $c$ is an integer and satisfies
$4(D_m+D_l+3\kappa_{\mathcal D})<c$, and $F$ an admissible polynomial with $v_{\mathrm F}(F)=0$. Then there is
$\tilde G\in\mathcal R[x]$ of degree $d(y)$ in the row of every position $y$ such that $\operatorname{val}_{S_{j,m}}(\tilde G)\geq\theta^L_{j,m}-6c$ for all $j$ and all $m$ with
$m+2\leq r$, and $v\big(B_{\mathcal D}(F,\tilde G)\big)\leq n_0+D_l+1$.
\end{lemma}
\leanref{Formalism/SharpRestriction/GrowthBounds/MonomialOrderSegments/CollapseExponent.lean}{Formalism.MainNotions.PlaneGram.StepB}
\leanref{Formalism/SharpRestriction/GrowthBounds/MonomialOrderSegments/CardLinkIdxMulLe.lean}{Formalism.MainNotions.PlaneGram.stepB_of}

\begin{proof}
Use the notation of Definition~\ref{pd-symbols-test-chains-projection-properties} with $\Lambda$ the set of positions,
so that $\mathcal K[z]$ is $\mathcal K[x]$ and $\mathrm{Pol}(d)$ is the space of the polynomials of degree $d(y)$ in the
row of every position $y$; the integral elements of $\mathcal K$ are those of $\mathcal R$. As $v(1)=0$, $v(n)\geq0$ for
every integer $n$. If $q$ is an integer prime to $p$, then $aq+bp=1$ for some integers $a$ and $b$; as $v(bp)\geq1>v(1)$,
$v(aq)=0$, and as $v(a),v(q)\geq0$, $v(q)=0$. So $v(n)=v_p(n)$ for integers $n\geq1$. The numbers $e_P$,
$D_m$, $D_l$ and $\kappa_{\mathcal D}$ are $\geq0$, as $2\sum_yd(y)+r\geq1$; so $4(D_m+D_l+3\kappa_{\mathcal D})<c$ gives
$D_l<c$.

\emph{Step 1: a leading monomial.} As $v_{\mathrm F}(F)=0$, $v(F_\mu)\geq-v(\mu!)/2$ for every $\mu$, and some monomial
$x^{\alpha_0}$ of $F$ has $v(F_{\alpha_0})=-v(\alpha_0!)/2$. As $F\in\mathrm{Pol}(d)$, so is $x^{\alpha_0}$. For every
exponent vector $\mu$ of row degrees $d$, $\mu!$ divides $\prod_yd(y)!$, the quotient being a product of multinomial
coefficients; so $v(\mu!)\leq v(\prod_yd(y)!)=2n_0$.

\emph{Step 2: a projection.} Let $X$ be the set of the pairs $(j,m)$ with $m+2\leq r$; it has at most $J(r-1)$ elements.
For $(j,m),(j',m')\in X$, $S_{j,m'}\subset S_{j,m}$ if $m\leq m'$, and $S_{j,m}\cap S_{j',m'}=\emptyset$ if $j\neq j'$.
Lemma~\ref{pd-symbols-test-chains-projection-fock-product}, with the integers $\lceil\theta^L_{j,m}\rceil$, gives $\Pi$.
Put $Y=\Pi(x^{\alpha_0})$. Then $Y\in\mathrm{Pol}(d)$, $\operatorname{val}_{S_{j,m}}(Y)\geq\theta^L_{j,m}$ for
$(j,m)\in X$, and, as $v_{\mathrm F}(x^{\alpha_0})=v(\alpha_0!)/2$ and $J(r-1)r^2\ell(2\sum_yd(y))=J(r-1)e_P=D_l$,
\[
 v(Y_\mu)\geq\frac{v(\alpha_0!)}2-D_l-\frac{v(\mu!)}2\qquad\text{for every }\mu .
\]
As $F$ is admissible and $\operatorname{val}_{S_{j,m}}(F)$ is an integer or $+\infty$,
$\operatorname{val}_{S_{j,m}}(F)\geq\lceil\theta^L_{j,m}\rceil$ for $(j,m)\in X$. So by (2) of that lemma,
$\langle F,Y\rangle_{\mathrm F}=\langle F,x^{\alpha_0}\rangle_{\mathrm F}=F_{\alpha_0}\alpha_0!$, which has valuation
$v(\alpha_0!)/2$.

\emph{Step 3: the units $\gamma$.} By (4) of Definition~\ref{pd-symbols-test-chains}, with $k=c$ and $\epsilon=-1$, there
are finitely many exponent vectors $\mu'$ of total degree at most $6c$ and $\lambda_{\mu'}\in\mathcal R$ such that
\[
 g(\mu)=\sum_{\mu'}\lambda_{\mu'}\prod_{y,a}\binom{\mu(y,a)}{\mu'(y,a)}
\]
satisfies $v(\gamma(\mu)^{-1}-g(\mu))\geq c$ for every $\mu$. Each $g(\mu)$ lies in $\mathcal R$. By
Lemma~\ref{pd-symbols-test-chains-projection-diagonal},
$\operatorname{val}_S(\operatorname{diag}_gG)\geq\operatorname{val}_S(G)-6c$ for every $G$ and every set $S$ of positions.
As $\gamma(\mu)$ is a unit, $v(\gamma(\mu)g(\mu)-1)=v(g(\mu)-\gamma(\mu)^{-1})\geq c$.

\emph{Step 4: the pairing.} Put $Y'=\operatorname{diag}_gY$. By the definition of $B_{\mathcal D}$,
\[
 B_{\mathcal D}(F,Y')=\Big\langle F,\sum_\mu\gamma(\mu)g(\mu)Y_\mu x^\mu\Big\rangle_{\mathrm F}=\langle F,Y\rangle_{\mathrm F}
 +\langle F,E\rangle_{\mathrm F},\qquad E=\sum_\mu\big(\gamma(\mu)g(\mu)-1\big)Y_\mu x^\mu .
\]
By Steps 2 and 3, $v(E_\mu)\geq c+v(\alpha_0!)/2-D_l-v(\mu!)/2$, and $v(F_\mu)\geq-v(\mu!)/2$; so every term
$F_\mu E_\mu\,\mu!$ of $\langle F,E\rangle_{\mathrm F}$ has valuation at least $c+v(\alpha_0!)/2-D_l$, which is larger than
$v(\alpha_0!)/2$ as $D_l<c$. With Step 2, $v(B_{\mathcal D}(F,Y'))=v(\alpha_0!)/2$.

\emph{Step 5: the scaling.} Put $N=\lceil n_0+D_l-v(\alpha_0!)/2\rceil$, which is $\geq0$ by Step 1, and
$\tilde G=p^NY'$. As $\operatorname{diag}_g$ does not enlarge supports, $\tilde G\in\mathrm{Pol}(d)$. For a monomial $x^\mu$
of $\tilde G$, $v(g(\mu))\geq0$, Step 2 and Step 1 give
$v(\tilde G_\mu)\geq N+v(\alpha_0!)/2-D_l-n_0\geq0$; so $\tilde G\in\mathcal R[x]$. By Steps 2 and 3,
$\operatorname{val}_{S_{j,m}}(\tilde G)\geq\operatorname{val}_{S_{j,m}}(Y)-6c\geq\theta^L_{j,m}-6c$ for $(j,m)\in X$.
Finally $v(B_{\mathcal D}(F,\tilde G))=N+v(\alpha_0!)/2<n_0+D_l+1$.
% Lean: stepB_of. Step 1 is exists_of_vBB_zero and fact_upper; Step 2 is prodProj for linkIdx with the sets Tset and
%   thresholds ceil(thl) (Tset_pair, card_linkIdx_mul_le), p1 to p4, vala_iff_Aset, fock_monomial_right and vle_vN;
%   Step 3 is ga_dig through GaOK (eps = -1) and approx_of_trunc with costV_ok and vala_dg_choose (Lemma
%   pd-symbols-test-chains-projection-diagonal); Step 4 is hsplit, herr (vge_unit) and hpair (vle_add_of_lt, vge_fock,
%   bge_dg); Step 5 is the element p^N * dg g Y with N = ceil(n0 + Dl - v(alpha0!)/2), cge_of_bge and vala_smul.
\end{proof}

\begin{definition}\label{pd-symbols-test-chains-projection-good}
Let $p$, $\mathcal R$ and $\mathcal K$ be as in Definition~\ref{pd-symbols-test-chains} or as in
Lemma~\ref{pd-symbols-test-chains-projection-lift}, $r,J\geq1$ integers, $\mathcal D$ a signed chain datum of rank $r$
with $J$ levels over $\mathcal R$, and $\mathcal K[\tilde x]$ the polynomial ring of its model positions. For a level $j$
let $d_L(j)$ be the sum of the degrees of the letters of level $j$. A monomial of
$\mathcal K[\tilde x]$ is \emph{good} for $\mathcal D$ if its exponent of each variable of the label of a link $(j,m)$ is
less than $P(j,m)$; its degree in the row of each letter $y$ is $d(y)$; and for each link $(j,m)$, with
$e_{\mathrm{lab}}$, $e_{\mathrm{lift}}$ and $e_{\mathrm{top}}$ its degrees in the rows of the label, the lift and the top
of $(j,m)$ ($e_{\mathrm{top}}=0$ if $m=0$),
\[
 e_{\mathrm{lab}}+P(j,m)\,e_{\mathrm{lift}}+e_{\mathrm{top}}=d(j,m),\qquad e_{\mathrm{top}}\leq d_L(j).
\]
A polynomial is \emph{good} for $\mathcal D$ if its monomials are.
\end{definition}
\leanref{Formalism/SharpRestriction/GrowthBounds/MonomialOrderSegments/CollapseExponent.lean}{Formalism.MainNotions.PlaneGram.TwistData.GoodE}
\leanref{Formalism/SharpRestriction/GrowthBounds/MonomialOrderSegments/CollapseExponent.lean}{Formalism.MainNotions.PlaneGram.TwistData.Good}

\begin{lemma}\label{pd-symbols-test-chains-projection-model-level-val}
Let $\mathcal K$ be a field, $\Lambda$ a finite set each of whose elements is a vector letter or a covector letter, $r\geq1$
an integer, and $\mathcal K[z]$ and $\operatorname{val}_S$, for $S\subset\Lambda$, as in
Definition~\ref{pd-symbols-test-chains-projection-properties}. Let $S\subset\Lambda$ and $F,G\in\mathcal K[z]$.
\begin{enumerate}
\item $\operatorname{val}_S(F)\in\mathbf Z\cup\{+\infty\}$; $\operatorname{val}_S(FG)\geq\operatorname{val}_S(F)+
  \operatorname{val}_S(G)$, $\operatorname{val}_S(F+G)\geq\min(\operatorname{val}_S(F),\operatorname{val}_S(G))$, and
  $\operatorname{val}_S(\lambda F)\geq\operatorname{val}_S(F)$ for $\lambda\in\mathcal K$.
\item If no monomial of $F$ involves a variable $z_{y,a}$ with $y\in S$ a covector letter, then
  $\operatorname{val}_S(F)\geq0$. In particular $\operatorname{val}_S(F)\geq0$ if $F$ involves no variable $z_{y,a}$ with
  $y\in S$.
\item Let $\Lambda'$ be another such set, with its polynomial ring $\mathcal K[z']$, and $f\colon\Lambda'\to\Lambda$ an
  injective map that preserves the types. Let $f_*\colon\mathcal K[z']\to\mathcal K[z]$ be the $\mathcal K$-algebra map
  with $f_*(z'_{y,a})=z_{f(y),a}$. Then $\operatorname{val}_S(f_*F')=\operatorname{val}_{f^{-1}(S)}(F')$ for every
  $F'\in\mathcal K[z']$.
\item Let $O\subset\Lambda$ be disjoint from $S$, and write $F=\sum_\mu\mu F_\mu$, where $\mu$ runs over the monomials in
  the variables $z_{y,a}$ with $y\in O$ and the $F_\mu$ involve none of these variables. Then
  $\operatorname{val}_S(F_\mu)\geq\operatorname{val}_S(F)$ for every $\mu$.
\item Let $O\subset\Lambda$ be disjoint from $S$, and $\Theta$ a $\mathcal K$-linear endomorphism of $\mathcal K[z]$ such
  that $\Theta(\mu\nu)=\nu\,\Theta(\mu)$ and $\operatorname{val}_S(\Theta(\mu))\geq0$ for every monomial $\mu$ in the
  variables $z_{y,a}$ with $y\in O$ and every monomial $\nu$ in the other variables. Then
  $\operatorname{val}_S(\Theta F)\geq\operatorname{val}_S(F)$.
\item Let $d\colon\Lambda\to\mathbf Z_{\geq0}$ and $F_d$ the sum of the terms of $F$ of row degrees $d$. Then
  $\operatorname{val}_S(F_d)\geq\operatorname{val}_S(F)$.
\end{enumerate}
\end{lemma}
\leanref{Formalism/SharpRestriction/GrowthBounds/MonomialOrderSegments/AddMemSupportOfMemSupportCoefOut.lean}{Formalism.MainNotions.PlaneGram.valS_mul_le}
\leanref{Formalism/SharpRestriction/GrowthBounds/MonomialOrderSegments/AddMemSupportOfMemSupportCoefOut.lean}{Formalism.MainNotions.PlaneGram.valS_sum_le}
\leanref{Formalism/SharpRestriction/GrowthBounds/MonomialOrderSegments/AddMemSupportOfMemSupportCoefOut.lean}{Formalism.MainNotions.PlaneGram.valS_C_mul_le}
\leanref{Formalism/SharpRestriction/GrowthBounds/MonomialOrderSegments/AddMemSupportOfMemSupportCoefOut.lean}{Formalism.MainNotions.PlaneGram.valS_nonneg_of_vecS}
\leanref{Formalism/SharpRestriction/GrowthBounds/MonomialOrderSegments/AddMemSupportOfMemSupportCoefOut.lean}{Formalism.MainNotions.PlaneGram.valS_nonneg_of_free}
\leanref{Formalism/SharpRestriction/GrowthBounds/MonomialOrderSegments/AddMemSupportOfMemSupportCoefOut.lean}{Formalism.MainNotions.PlaneGram.valS_rename}
\leanref{Formalism/SharpRestriction/GrowthBounds/MonomialOrderSegments/AddMemSupportOfMemSupportCoefOut.lean}{Formalism.MainNotions.PlaneGram.valS_coefOut_le}
\leanref{Formalism/SharpRestriction/GrowthBounds/MonomialOrderSegments/AddMemSupportOfMemSupportCoefOut.lean}{Formalism.MainNotions.PlaneGram.valS_outside_linear}
\leanref{Formalism/SharpRestriction/GrowthBounds/MonomialOrderSegments/AddMemSupportOfMemSupportCoefOut.lean}{Formalism.MainNotions.PlaneGram.valS_hcomp_le}

\begin{proof}
We use the substitution $F\mapsto F_{S,r}$ of Definition~\ref{pd-symbols-path-generators}(1), with $n=r$, which replaces
the row $z_y$, $y\in S$, by $z_y\tau_r(\eta)$ for a vector letter and by $z_y\,{}^t\tau_r(\eta^{-1})$ for a covector
letter; $\operatorname{val}_S(F)$ is the least exponent of $\eta$ in its coefficients.

(1) The substitution is a homomorphism of $\mathcal K$-algebras, and the least exponent of $\eta$ in a product of two
Laurent polynomials in $\eta$ is at least the sum of the least exponents.

(2) The entries of $\tau_r(\eta)=\mathfrak gD_r(\eta)\mathfrak g^{-1}$ are polynomials in $\eta$. If $F$ involves no
variable of a covector letter in $S$, then $F_{S,r}$ is obtained by substituting polynomials in $\eta$, so only exponents
$\geq0$ occur.

(3) As $f$ preserves the types, $(f_*F')_{S,r}$ is $f_*$, applied to the variables, of $F'_{f^{-1}(S),r}$. As $f$ is
injective, $f_*$ maps distinct monomials to distinct monomials, so the two have the same coefficients.

(4) The substitution fixes the variables of the letters in $O$, so $F_{S,r}=\sum_\mu\mu(F_\mu)_{S,r}$, and
$(F_\mu)_{S,r}$ involves none of these variables. So the coefficient of $F_{S,r}$ at a monomial $\mu\nu$, with $\nu$ free
of them, is that of $(F_\mu)_{S,r}$ at $\nu$, and every exponent of $\eta$ that occurs in $(F_\mu)_{S,r}$ occurs in
$F_{S,r}$.

(5) Write $F=\sum_\mu\mu F_\mu$ as in (4). As $F_\mu$ is a combination of monomials $\nu$ in the other variables,
$\Theta F=\sum_\mu\Theta(\mu)F_\mu$. By (4) and (1), $\operatorname{val}_S(\Theta(\mu)F_\mu)\geq0+\operatorname{val}_S(F)$
for every $\mu$, and by (1) the sum has $\operatorname{val}_S\geq\operatorname{val}_S(F)$.

(6) The substitution replaces each row by linear combinations of the variables of the same row, so it maps
polynomials of row degrees $d$ to polynomials of row degrees $d$. Hence $(F_d)_{S,r}$ is the sum of the terms of
$F_{S,r}$ of row degrees $d$, and its coefficients are coefficients of $F_{S,r}$.
% Lean: (6) is valS_hcomp_le (coeff_subV_hcomp), with hcomp d F = F_d.
% Lean: (1) is valS_mul_le, valS_sum_le and valS_C_mul_le (valS takes values in WithTop Z); (2) is valS_nonneg_of_vecS
%   (substS_mem with hMat_mem: the entries of hMat lie in the image of the polynomials in T) and valS_nonneg_of_free;
%   (3) is valS_rename (substS_rename, coeff_rename_mapDomain); (4) is valS_coefOut_le with coefOut out mu F = F_mu
%   (coeff_substS_monomial); (5) is valS_outside_linear, with the letters of O as `out' and sum_mul_coefOut.
\end{proof}

\begin{lemma}\label{pd-symbols-test-chains-projection-model-level}
Let $p\geq5$, $\mathcal R$, $\mathcal K$ and $v$ be as in Lemma~\ref{pd-symbols-test-chains-projection-lift}, and suppose
that $\mathcal K$ with $v$ satisfies (S) of Definition~\ref{pd-symbols-test-chains-projection-properties} at $p$. Let
$r,J\geq1$ be integers, $\mathcal D$ a signed chain datum of rank $r$ with $J$ levels over $\mathcal R$
(Definition~\ref{pd-symbols-test-chains}), $j$ a level, and $d_L(j)$ the sum of the degrees of the letters of level $j$.
For $m<r$ let $T_{j,m}$ be the set of the letters of level $j$ and the labels of the links $(j,m')$ with $m'>m$, and
$Y_{j,m}$ that of the letters of level $j$ and the tops of the links $(j,m')$ with $m'>m$. A monomial of
$\mathcal K[\tilde x]$ is \emph{open at $j$} if its degree in the row of each letter $y$ of level $j$ is $d(y)$, its
degree in the row of the label of each link $(j,m)$ is $d(j,m)$, and it involves no top of level $j$. It is \emph{closed
at $j$} if its degree in the row of each letter $y$ of level $j$ is $d(y)$ and, for each link $(j,m)$, its degrees
$e_{\mathrm{lab}}$ and $e_{\mathrm{top}}$ in the rows of the label and the top of $(j,m)$ ($e_{\mathrm{top}}=0$ if $m=0$)
satisfy $e_{\mathrm{lab}}+e_{\mathrm{top}}=d(j,m)$ and $e_{\mathrm{top}}\leq d_L(j)$. Then there is a $\mathcal K$-linear
endomorphism $\Theta$ of $\mathcal K[\tilde x]$ with the following properties.
\begin{enumerate}
\item Let $\zeta$ be a monomial open at $j$. Then $\Theta(\zeta)\in\mathcal R[\tilde x]$, $\varphi(\Theta(\zeta))=\varphi(\zeta)$,
  and every monomial of $\Theta(\zeta)$ is closed at $j$ and has the exponents of $\zeta$ in the variables of the model
  positions other than the letters, labels and tops of level $j$.
\item Let $Z\in\mathcal R[\tilde x]$ have all its monomials open at $j$, $m+2\leq r$ and $\theta$ real. If
  $\operatorname{val}_{T_{j,m}}(Z)\geq\theta$, then $\operatorname{val}_{Y_{j,m}}(\Theta Z)\geq\theta$.
\item $\operatorname{val}_S(\Theta Z)\geq\operatorname{val}_S(Z)$ for every $Z\in\mathcal K[\tilde x]$ and every set $S$
  of model positions that contains no letter, label or top of level $j$.
\end{enumerate}
\end{lemma}
\leanref{Formalism/SharpRestriction/GrowthBounds/MonomialOrderSegments/SegDefs.lean}{Formalism.MainNotions.PlaneGram.SegStep}
\leanref{Formalism/SharpRestriction/GrowthBounds/MonomialOrderSegments/SegDefs.lean}{Formalism.MainNotions.PlaneGram.TwistData.PreE}
\leanref{Formalism/SharpRestriction/GrowthBounds/MonomialOrderSegments/SegDefs.lean}{Formalism.MainNotions.PlaneGram.TwistData.PostE}
\leanref{Formalism/SharpRestriction/GrowthBounds/MonomialOrderSegments/RenamePullLCoefOut.lean}{Formalism.MainNotions.PlaneGram.segStep_of}

\begin{proof}
We use Lemma~\ref{pd-symbols-test-chains-projection-model-level-val} for the model positions, the letters with their
types and the labels, lifts and tops as vector letters, and the notation $f_*$ of its part (3).

\emph{Step 1: the letters of the level.} Let $y_0,\dots,y_{n-1}$ be the letters of level $j$, and
$\Lambda=\{0,\dots,n+r-1\}$, lettered so that $l<n$ has the type of $y_l$ and $n,\dots,n+r-1$ are vector letters. For
$m<r$ put $l_m=n+r-1-m$, the \emph{index} of the link $(j,m)$, and let $d_\Lambda(l)=d(y_l)$ for $l<n$ and
$d_\Lambda(l_m)=d(j,m)$. Let $e_0$ and $e_1$ be the injective maps from $\Lambda$ to the model positions with
$e_0(l)=e_1(l)=y_l$ for $l<n$, $e_0(l_m)$ the label of $(j,m)$, and $e_1(l_m)$ the top of $(j,m)$ if $m\geq1$ and the
label of $(j,0)$ if $m=0$. They preserve the types. Let $\varepsilon$ map $l<n$ to $y_l$ and $l_m$ to the link $(j,m)$,
and $\varepsilon_*\colon\mathcal K[z]\to\mathcal K[x]$ the map with $z_{l,a}\mapsto x_{\varepsilon(l),a}$. As $\varphi$
maps the variables of a label or a top to those of its link, $\varphi\circ e_*=\varepsilon_*$ for $e=e_0,e_1$. With
$S_i=\{0,\dots,i-1\}$, the definitions give $e_0^{-1}(T_{j,m})=S_{l_m}=e_1^{-1}(Y_{j,m})$, and $e_0^{-1}(Y_{j,m})$
consists of indices $l<n$.

\emph{Step 2: the basis.} Property (S) for $\Lambda$ and $d_\Lambda$ gives integral $\Psi_b=\prod_kQ_{b,k}$ in
$\mathrm{Pol}(d_\Lambda)$, linearly independent, with $c_{b,k}$, $O_{b,k}$, $h_{b,k}$ and $\pi_b$. Call $Q_{b,k}$
\emph{marked} if $O_{b,k}$ contains some $l<n$, put $f_{b,k}=e_1$ if it is and $f_{b,k}=e_0$ if not, and
\[
 \Psi'_b=\prod_kf_{b,k*}(Q_{b,k})\in\mathcal K[\tilde x].
\]
(a) \emph{$\Psi'_b$ is integral.} For a polynomial $Q$ let $v_G(Q)$ be the least valuation of its coefficients. By
Gauss's lemma $v_G(QQ')=v_G(Q)+v_G(Q')$ \cite[\S5.1.2]{BGR84}, and $v_G(f_*Q)=v_G(Q)$ as $f_*$ maps distinct monomials to
distinct monomials. So $v_G(\Psi'_b)=\sum_kv_G(Q_{b,k})=v_G(\Psi_b)\geq0$.

(b) $\varphi(\Psi'_b)=\varepsilon_*(\Psi_b)$, by Step 1.

(c) \emph{Every monomial $\xi$ of $\Psi'_b$ is closed at $j$ and involves only variables of letters, labels and tops of
level $j$.} By (S)(2), $f_{b,k*}(Q_{b,k})$ has degree $1$ in the row of $f_{b,k}(l)$ for $l\in O_{b,k}$ and $0$ in the
other rows. So the degree of $\xi$ in the row of a model position $y$ is the number of $k$ with $y\in f_{b,k}(O_{b,k})$;
in particular $\xi$ involves only the positions $e_0(l)$ and $e_1(l)$, which are letters, labels and tops of level $j$.
For a letter $y_l$ this number is that of the $k$ with $l\in O_{b,k}$, which is $d(y_l)$ by (S)(2). For a link $(j,m)$,
$f_{b,k}(l)$ is the label or the top of $(j,m)$ exactly when $l=l_m$, so $e_{\mathrm{lab}}+e_{\mathrm{top}}$ is the number
of $k$ with $l_m\in O_{b,k}$, which is $d(j,m)$. The top of $(j,m)$ is $f_{b,k}(l_m)$ only for marked $Q_{b,k}$, and a
marked $O_{b,k}$ contains some $l<n$; so by (S)(2),
\[
 e_{\mathrm{top}}\leq\sum_k\big|O_{b,k}\cap\{0,\dots,n-1\}\big|=\sum_{l<n}d(y_l)=d_L(j).
\]
(d) \emph{$\operatorname{val}_{Y_{j,m}}(\Psi'_b)\geq\pi_b(l_m)$ for $m<r$.} By
Lemma~\ref{pd-symbols-test-chains-projection-model-level-val}(1) it is enough to show
$\operatorname{val}_{Y_{j,m}}(f_{b,k*}Q_{b,k})\geq\omega(c_{b,k},h_{b,k}(l_m))$ for every $k$. If $Q_{b,k}$ is marked, by
part (3) of that lemma and Step 1, $\operatorname{val}_{Y_{j,m}}(e_{1*}Q_{b,k})=\operatorname{val}_{S_{l_m}}(Q_{b,k})$,
which is at least $\omega(c_{b,k},h_{b,k}(l_m))$ by (S)(3). If not, $\operatorname{val}_{Y_{j,m}}(e_{0*}Q_{b,k})=
\operatorname{val}_{e_0^{-1}(Y_{j,m})}(Q_{b,k})\geq0$ by parts (3) and (2), as $Q_{b,k}$ has degree $0$ in the rows of the
$l<n$. And $\omega(c_{b,k},h_{b,k}(l_m))\leq0$: this is the definition of $\omega$ if $c_{b,k}\leq0$, and if
$c_{b,k}\geq1$, then by (S)(3) $h_{b,k}(l_m)$ is at most the number of elements of $O_{b,k}$ in $\{n,\dots,l_m-1\}$,
which is at most $r-1-m<r$.

\emph{Step 3: the map.} For a monomial $\lambda\in\mathrm{Pol}(d_\Lambda)$, (S)(1) gives integral $c_b(\lambda)\in\mathcal
K$ with $\lambda=\sum_bc_b(\lambda)\Psi_b$; we fix them, and put $c_b(\lambda)=0$ for the other monomials $\lambda$. Let
$\Theta'\colon\mathcal K[z]\to\mathcal K[\tilde x]$ be the $\mathcal K$-linear map with
$\Theta'(\lambda)=\sum_bc_b(\lambda)\Psi'_b$. If $Q\in\mathrm{Pol}(d_\Lambda)$ and $Q=\sum_ba_b\Psi_b$, then
$\Theta'(Q)=\sum_ba_b\Psi'_b$: the monomials of $Q$ lie in $\mathrm{Pol}(d_\Lambda)$, so
$Q=\sum_b\big(\sum_\lambda Q_\lambda c_b(\lambda)\big)\Psi_b$, and by linear independence
$\sum_\lambda Q_\lambda c_b(\lambda)=a_b$. For a monomial $\zeta$ of $\mathcal K[\tilde x]$ let $\zeta_{\mathrm{in}}$ be the
product of its powers of the variables of the positions $e_0(l)$, the letters of level $j$ and the labels of the links
$(j,m)$, put $\zeta_{\mathrm{out}}=\zeta/\zeta_{\mathrm{in}}$, and let $\lambda(\zeta)$ be the monomial of $\mathcal K[z]$
with $e_{0*}\lambda(\zeta)=\zeta_{\mathrm{in}}$. Let $\Theta$ be the $\mathcal K$-linear map with
$\Theta(\zeta)=\zeta_{\mathrm{out}}\,\Theta'(\lambda(\zeta))$.

\emph{Step 4: property (1).} Let $\zeta$ be open at $j$. The degree of $\lambda(\zeta)$ in the row of $l<n$ is that of
$\zeta$ in the row of $y_l$, and in the row of $l_m$ that of $\zeta$ in the row of the label of $(j,m)$; so
$\lambda(\zeta)\in\mathrm{Pol}(d_\Lambda)$, and the $c_b(\lambda(\zeta))$ are integral. With (a), $\Theta(\zeta)$ is
integral. By (b) and Step 1,
\[
 \varphi(\Theta(\zeta))=\varphi(\zeta_{\mathrm{out}})\sum_bc_b(\lambda(\zeta))\,\varepsilon_*(\Psi_b)
 =\varphi(\zeta_{\mathrm{out}})\,\varepsilon_*(\lambda(\zeta))=\varphi(\zeta_{\mathrm{out}})\,\varphi(\zeta_{\mathrm{in}})
 =\varphi(\zeta).
\]
A monomial of $\Theta(\zeta)$ is $\zeta_{\mathrm{out}}\xi$ with $\xi$ a monomial of some $\Psi'_b$. Now
$\zeta_{\mathrm{out}}$ involves no letter or label of level $j$, by definition, and no top of level $j$, as $\zeta$ is open
at $j$. So in the rows of the letters, labels and tops of level $j$ the monomial $\zeta_{\mathrm{out}}\xi$ has the degrees
of $\xi$, and it is closed at $j$ by (c); in the other variables $\xi$ has exponent $0$ by (c), so
$\zeta_{\mathrm{out}}\xi$ has the exponents of $\zeta$.

\emph{Step 5: property (2).} Let $Z$, $m$ and $\theta$ be as in (2). Write $Z=\sum_\mu\mu Z_\mu$, with $\mu$ running over
the monomials in the variables of the model positions outside $e_0(\Lambda)$ and $Z_\mu$ involving only variables of
positions in $e_0(\Lambda)$, and let $Q_\mu\in\mathcal K[z]$ with $e_{0*}Q_\mu=Z_\mu$. By the definition of $\Theta$,
$\Theta Z=\sum_\mu\mu\,\Theta'(Q_\mu)$. Fix $\mu$. It involves no letter of level $j$, and no top of level $j$ as the
monomials of $Z$ are open at $j$; so $\operatorname{val}_{Y_{j,m}}(\mu)\geq0$ by
Lemma~\ref{pd-symbols-test-chains-projection-model-level-val}(2). The monomials of $Q_\mu$ are monomials $\lambda(\zeta)$
with $\zeta$ a monomial of $Z$, so $Q_\mu\in\mathrm{Pol}(d_\Lambda)$ by Step 4, and its coefficients are coefficients of
$Z$, so it is integral. By parts (3) and (4) of that lemma, the latter for $O$ the positions outside $e_0(\Lambda)$, which
are disjoint from $T_{j,m}\subset e_0(\Lambda)$,
\[
 \operatorname{val}_{S_{l_m}}(Q_\mu)=\operatorname{val}_{T_{j,m}}(Z_\mu)\geq\operatorname{val}_{T_{j,m}}(Z)\geq\theta,
\]
so $\operatorname{val}_{S_{l_m}}(Q_\mu)\geq\lceil\theta\rceil$, as it is an integer. Property (S)(1) for $Q_\mu$, with
$\theta(l_m)=\lceil\theta\rceil$ and $\theta(i)=\operatorname{val}_{S_i}(Q_\mu)$ for the other $i$ (or $0$ if it is
$+\infty$), gives integral $a_b$ with $Q_\mu=\sum_ba_b\Psi_b$ and $\lceil\theta\rceil\leq\pi_b(l_m)$ whenever $a_b\neq0$.
By Step 3, $\Theta'(Q_\mu)=\sum_ba_b\Psi'_b$, and by (d) each term with $a_b\neq0$ has
$\operatorname{val}_{Y_{j,m}}\geq\pi_b(l_m)\geq\theta$. By part (1) of that lemma,
$\operatorname{val}_{Y_{j,m}}(\mu\,\Theta'(Q_\mu))\geq\theta$, and so $\operatorname{val}_{Y_{j,m}}(\Theta Z)\geq\theta$.

\emph{Step 6: property (3).} Let $S$ contain no letter, label or top of level $j$; then $S$ is disjoint from
$e_0(\Lambda)$. For monomials $\mu$ in the variables of the positions in $e_0(\Lambda)$ and $\nu$ in the others,
$\Theta(\mu\nu)=\nu\,\Theta'(\lambda(\mu))=\nu\,\Theta(\mu)$. By (c), $\Theta(\mu)=\sum_bc_b(\lambda(\mu))\Psi'_b$ involves
no variable of a position in $S$, so $\operatorname{val}_S(\Theta(\mu))\geq0$ by
Lemma~\ref{pd-symbols-test-chains-projection-model-level-val}(2). Part (5) of that lemma, with $O=e_0(\Lambda)$, gives
$\operatorname{val}_S(\Theta Z)\geq\operatorname{val}_S(Z)$.
% Lean: segStep_of hS (RenamePullLCoefOut.lean). Step 1: SPAN at N = nB j + r with sg = msg o e0 j and degrees
%   dO (eA j k); the letters are eB j and the index natAdd (nB j) m' is the link (j, rev m'), so l_m = nB j + r - 1 - m
%   = stg j m; e0, e1 and eA are eps (SecondEmbeddingInjective.lean), with ph_X_fk, TL_preimage, YP_preimage1 and
%   YP_preimage0. Step 2: Marked, fk and PsiP (SecondEmbeddingTangentIndex.lean); (a) is PsiP_integral =
%   integral_prod_rename (vG_prod, vG_mul by a monomial order, and vG_rename, ExistsMonomialOrder.lean), (b) PsiP_ph,
%   (c) PsiP_support (PsiP_HomD, cnt_mb, cnt_lab_add_tp, cnt_tp_le, cnt_tp_zero, cnt_not_segL), (d) PsiP_val
%   (valS_prod_le, valS_rename, card_unmarked). Step 3: cc lam (the choice from h5, 0 off PolO), TT = Theta', TT_expand,
%   outSeg, pullE (lambda(zeta)), Th and Th_monomial (MonomialMemPolKSg.lean). Step 4 is Th_integral (pullE_deg), Th_ph
%   and Th_post (Th_supp); Step 5 is Th_val (Th_eq, pullL_coefOut, valS_nonneg_of_free, valS_coefOut_le, h7 with
%   th = ceil(theta) at stg j m and untopD 0 elsewhere, vala_sum, vala_smul); Step 6 is Th_away (valS_outside_linear with
%   the predicate "not outSeg", valS_nonneg_of_vecS).
\end{proof}

\begin{lemma}\label{pd-symbols-test-chains-projection-model}
Let $p\geq5$, $\mathcal R$, $\mathcal K$ and $v$ be as in Lemma~\ref{pd-symbols-test-chains-projection-lift}, and suppose
that $\mathcal K$ with $v$ satisfies (V) and (S) of Definition~\ref{pd-symbols-test-chains-projection-properties} at $p$.
Let $r,J\geq1$ be integers, $\mathcal D$ a signed chain datum of rank $r$ with $J$ levels over $\mathcal R$
(Definition~\ref{pd-symbols-test-chains}), and $\theta_{j,m}$ real numbers for $j\in\mathbf Z/J\mathbf Z$ and $m+2\leq r$.
Let $\tilde G\in\mathcal R[x]$ have degree $d(y)$ in the row of every position $y$, and suppose that
$\operatorname{val}_{S_{j,m}}(\tilde G)\geq\theta_{j,m}$ for all $j$ and $m+2\leq r$. Then there is $Z\in\mathcal R[\tilde x]$,
good for $\mathcal D$ (Definition~\ref{pd-symbols-test-chains-projection-good}), with $\varphi(Z)=\tilde G$ and
$\operatorname{val}_{\tilde S_{j,m}}(Z)\geq\theta_{j,m}$ for all $j$ and $m+2\leq r$.
\end{lemma}
\leanref{Formalism/SharpRestriction/GrowthBounds/MonomialOrderSegments/CollapseExponent.lean}{Formalism.MainNotions.PlaneGram.StepC}
\leanref{Formalism/SharpRestriction/GrowthBounds/MonomialOrderSegments/ExistsOfMemSupportLin.lean}{Formalism.MainNotions.PlaneGram.stepC_of}

\begin{proof}
We use the notions open and closed at a level, and the sets $T_{j,m}$ and $Y_{j,m}$, of
Lemma~\ref{pd-symbols-test-chains-projection-model-level}. For a set $s$ of levels, $Z\in\mathcal R[\tilde x]$ is
\emph{prepared for $s$} if $\varphi(Z)=\tilde G$, its monomials involve no lift, are closed at the levels in $s$ and open
at the others, and for $m+2\leq r$, $\operatorname{val}_{Y_{j,m}}(Z)\geq\theta_{j,m}$ if $j\in s$ and
$\operatorname{val}_{T_{j,m}}(Z)\geq\theta_{j,m}$ if $j\notin s$.

\emph{Step 1: the start.} Let $Z_\emptyset$ be $\tilde G$ with each variable of a letter replaced by the corresponding
variable of the letter among the model positions, and each variable of a link $(j,m)$ by the corresponding variable of
its label. Then $Z_\emptyset\in\mathcal R[\tilde x]$, and $\varphi(Z_\emptyset)=\tilde G$. The monomials of $Z_\emptyset$
involve no lift and no top, and as $\tilde G$ has degree $d(y)$ in every row, they are open at every level. The renaming
keeps the types of the rows and maps $S_{j,m}$ onto $T_{j,m}$, and the substitution that defines $\operatorname{val}$
commutes with it; so $\operatorname{val}_{T_{j,m}}(Z_\emptyset)=\operatorname{val}_{S_{j,m}}(\tilde G)\geq\theta_{j,m}$.
So $Z_\emptyset$ is prepared for $\emptyset$.

\emph{Step 2: one more level.} Let $Z$ be prepared for $s$, let $j\notin s$, and let $\Theta$ be the map of
Lemma~\ref{pd-symbols-test-chains-projection-model-level} for $j$. We show that $\Theta Z$ is prepared for $s\cup\{j\}$.
The monomials $\zeta$ of $Z$ are open at $j$, so by (1) of that lemma $\Theta Z=\sum_\zeta Z_\zeta\Theta(\zeta)$ lies in
$\mathcal R[\tilde x]$ and $\varphi(\Theta Z)=\sum_\zeta Z_\zeta\varphi(\zeta)=\tilde G$. A monomial $\zeta'$ of $\Theta Z$ is a
monomial of some $\Theta(\zeta)$. It is closed at $j$, and it has the exponents of $\zeta$ in the variables of the lifts
and of the letters, labels and tops of the levels $j'\neq j$; so it involves no lift, and it is closed or open at $j'$
when $\zeta$ is. By (2), $\operatorname{val}_{Y_{j,m}}(\Theta Z)\geq\theta_{j,m}$. For $j'\neq j$ the sets $Y_{j',m}$ and
$T_{j',m}$ contain no letter, label or top of level $j$, so (3) gives the bounds at $j'$.

\emph{Step 3: the splitting.} By induction on the number of levels in $s$, Steps 1 and 2 give $Z_1$ prepared for the
set of all levels. For a monomial $\zeta$ let $\zeta^{\mathrm{sp}}$ have the exponents of $\zeta$, except that for each
link $(j,m)$ and $a<r$ the exponent $e$ of the variable of the label of $(j,m)$ in column $a$ is replaced by its
remainder modulo $P(j,m)$, and $\lfloor e/P(j,m)\rfloor$ is added to the exponent of the variable of the lift of $(j,m)$
in column $a$. Let $Z$ be the image of $Z_1$ under the linear map $\zeta\mapsto\zeta^{\mathrm{sp}}$. Its coefficients are
sums of coefficients of $Z_1$, so $Z\in\mathcal R[\tilde x]$. As $\varphi$ maps these two variables to $x_{(j,m),a}$ and
$x_{(j,m),a}^{P(j,m)}$, $\varphi(\zeta^{\mathrm{sp}})=\varphi(\zeta)$, and $\varphi(Z)=\varphi(Z_1)=\tilde G$.

\emph{Step 4: $Z$ is good.} Let $\zeta$ be a monomial of $Z_1$; it involves no lift and is closed at every level. In
$\zeta^{\mathrm{sp}}$ the label exponents are less than $P(j,m)$, and the rows of the letters and the tops are those of
$\zeta$. As $\zeta$ involves no lift, the degrees of $\zeta^{\mathrm{sp}}$ in the rows of the label and the lift of
$(j,m)$ satisfy $e_{\mathrm{lab}}+P(j,m)e_{\mathrm{lift}}=$ the degree of $\zeta$ in the row of the label. So
$e_{\mathrm{lab}}+P(j,m)e_{\mathrm{lift}}+e_{\mathrm{top}}=d(j,m)$ and $e_{\mathrm{top}}\leq d_L(j)$.

\emph{Step 5: the bounds.} Let $m+2\leq r$. The set $\tilde S_{j,m}$ is $Y_{j,m}$ together with the lifts of the links of
level $j-1$, which are vector positions; so (V)(c) gives $\operatorname{val}_{\tilde S_{j,m}}(Z_1)\geq
\operatorname{val}_{Y_{j,m}}(Z_1)\geq\theta_{j,m}$. Write $Z_1=\sum_\mu\lambda^\mu W_\mu$, where the $\lambda^\mu$ are
distinct monomials in the label variables and the $W_\mu$ involve no label. As $\tilde S_{j,m}$ contains no label, the
substitution that defines $\operatorname{val}_{\tilde S_{j,m}}$ fixes the label variables, so
$\operatorname{val}_{\tilde S_{j,m}}(W_\mu)\geq\operatorname{val}_{\tilde S_{j,m}}(Z_1)$ for every $\mu$. Now
$Z=\sum_\mu(\lambda^\mu)^{\mathrm{sp}}W_\mu$, and $(\lambda^\mu)^{\mathrm{sp}}$ is a monomial in the variables of labels and
lifts, which are vector positions; the substitution replaces each of its rows $z$ in $\tilde S_{j,m}$ by $z\tau_r(\eta)$,
whose entries involve only nonnegative powers of $\eta$, so $\operatorname{val}_{\tilde S_{j,m}}((\lambda^\mu)^{\mathrm{sp}})
\geq0$. Hence $\operatorname{val}_{\tilde S_{j,m}}(Z)\geq\min_\mu\operatorname{val}_{\tilde S_{j,m}}(W_\mu)\geq\theta_{j,m}$.
% Lean: stepC_of hV2 (segStep_of hSPAN) (ExistsOfMemSupportLin.lean). Prepared for s is Inv v p theta Gt s; Step 1 is
%   inv_start (rename by iotaL, ph_rename_iotaL, vala_rename_iff with preimage_TLset), Step 2 is inv_step
%   (PostE.of_agree, PreE.of_agree, notSegL_of_mem_YPset, notSegL_of_mem_TLset), and inv_all the induction. Step 3 is
%   Sp p (splitE, ph_Sp with phE_splitE, integral_Sp), Step 4 is goodE_splitE, and Step 5 is vala_superset_vec (V2(c))
%   and vala_Sp (coefOut for the labels labB, valS_coefOut_le, valS_monomial_nonneg_of_vec, valS_mul_le).
\end{proof}

\begin{definition}\label{pd-symbols-test-chains-projection-omega-blocks}
Let $p$, $\mathcal R$, $\mathcal K$ and $v$ be as in Lemma~\ref{pd-symbols-test-chains-projection-lift}, $r,J\geq1$
integers and $\mathcal D$ a signed chain datum of rank $r$ with $J$ levels over $\mathcal R$, with the notation of
Definition~\ref{pd-symbols-test-chains}; we identify a monomial with its exponent vector, and define $v_{\mathrm F}$ on
$\mathcal K[\tilde x]$ by the formula used on $\mathcal K[x]$. Let $\zeta$ be a monomial of $\mathcal K[\tilde x]$.
\begin{enumerate}
\item For a link $t$ let $\lambda_t(\zeta),w_t(\zeta),\tau_t(\zeta)\in\mathbf Z_{\geq0}^r$ be the exponents of $\zeta$ in
  the variables of the label, the lift and the top of $t$ ($\tau_t(\zeta)=0$ if $t=(j,0)$). Let $\zeta^{\mathrm c}$ and
  $\zeta^{\mathrm t}$ be the exponent vectors of the positions that agree with $\zeta$, respectively vanish, at the
  letters, and at a link $t$ are $\lambda_t(\zeta)+P(t)w_t(\zeta)$, respectively $\tau_t(\zeta)$. So
  $\varphi(\zeta)=x^{\zeta^{\mathrm c}+\zeta^{\mathrm t}}$.
\item For a family $\rho=(\rho_t)$ of $\rho_t\in\mathbf Z_{\geq0}^r$, indexed by the links, let
  $\operatorname{emb}_\rho(\zeta)$ be the exponent vector of the positions that agrees with $\zeta$ at the letters and is
  $\rho_t+P(t)w_t(\zeta)$ at a link $t$. So $\zeta^{\mathrm c}=\operatorname{emb}_{\lambda(\zeta)}(\zeta)$, with
  $\lambda(\zeta)=(\lambda_t(\zeta))_t$.
\item Monomials $\zeta$ and $\zeta'$ are \emph{of the same block} if they have the same degree in the row of every
  model position and $\lambda_t(\zeta)=\lambda_t(\zeta')$ and $\tau_t(\zeta)=\tau_t(\zeta')$ for every link $t$.
\item Suppose that $c$ is an integer, and let $c_t(\rho,n)$ and $\gamma_t(\rho,\beta)$, for the links $t$, be elements of
  $\mathcal K$ with the properties of (D) of Definition~\ref{pd-symbols-test-chains-projection-properties} for $\tau$ the
  set of links, $r'=r$ and the function $t\mapsto g(t)$, the second property for this $c$. Put
  \[
   \gamma^{\mathrm P}(\zeta)=\gamma(\zeta^{\mathrm c})\prod_t\gamma_t\big(\lambda_t(\zeta),w_t(\zeta)\big),\qquad
   \kappa_{\mathrm C}(\zeta)=\prod_tc_t\big(\lambda_t(\zeta),|w_t(\zeta)|\big)\Big/\prod_{t,a}\lambda_t(\zeta)(a)!,
  \]
  the products over the links $t$, with $|w|=\sum_aw(a)$.
\end{enumerate}
\end{definition}
\leanref{Formalism/SharpRestriction/GrowthBounds/MonomialOrderSegments/CollapseExponent.lean}{Formalism.MainNotions.PlaneGram.TwistData.collapseExponent}
\leanref{Formalism/SharpRestriction/GrowthBounds/MonomialOrderSegments/CollapseExponent.lean}{Formalism.MainNotions.PlaneGram.TwistData.tpE}
\leanref{Formalism/SharpRestriction/GrowthBounds/MonomialOrderSegments/CollapseExponent.lean}{Formalism.MainNotions.PlaneGram.TwistData.gaP}
\leanref{Formalism/SharpRestriction/GrowthBounds/MonomialOrderSegments/CollapseExponent.lean}{Formalism.MainNotions.PlaneGram.TwistData.DigOK}
\leanref{Formalism/SharpRestriction/GrowthBounds/MonomialOrderSegments/DdivOnLetterMonomial.lean}{Formalism.MainNotions.PlaneGram.TwistData.emb}
\leanref{Formalism/SharpRestriction/GrowthBounds/MonomialOrderSegments/GlobalTheory.lean}{Formalism.MainNotions.PlaneGram.TwistData.SameBlockAlt}
\leanref{Formalism/SharpRestriction/GrowthBounds/MonomialOrderSegments/GlobalTheory.lean}{Formalism.MainNotions.PlaneGram.TwistData.kapC}

\begin{lemma}\label{pd-symbols-test-chains-projection-omega-cost}
Let $\mathcal K$ be a field of characteristic $0$ and $v\colon\mathcal K\to\mathbf Q\cup\{+\infty\}$ a valuation, with
$\Lambda$, $r$, $\mathcal K[z]$, $\operatorname{val}_S$ and $\operatorname{diag}_g$ as in
Lemma~\ref{pd-symbols-test-chains-projection-diagonal}. Let $\sigma$ be a finite set. A \emph{cost relation} on the
functions $\mathbf Z_{\geq0}^\sigma\to\mathcal K$ is a relation ``$g$ costs $n$'' between such functions $g$ and integers
$n\geq0$ with the following four rules: the constant $1$ costs $0$; if $g$ costs $k$ and $h$ costs $l$, then $gh$ costs
$k+l$; if $g$ and $h$ cost $n$, so do $g+h$ and $\lambda g$ for every $\lambda\in\mathcal K$; and if $g$ costs $n$, it
costs every $n'\geq n$. For a set $\mathcal F$ of subsets of $\Lambda$, a function $g$ from the exponent vectors
$\mu\colon\Lambda\times\{1,\dots,r\}\to\mathbf Z_{\geq0}$ to $\mathcal K$ \emph{costs $n$ on $\mathcal F$} if for every
$S\in\mathcal F$, every real $\theta$ and every $F\in\mathcal K[z]$ with $\operatorname{val}_S(F)\geq\theta$,
$\operatorname{val}_S(\operatorname{diag}_gF)\geq\theta-n$.
\begin{enumerate}
\item For every set $\mathcal F$ of subsets of $\Lambda$, ``$g$ costs $n$ on $\mathcal F$'' is a cost relation.
\item Let $\rho,P,k\geq0$ be integers. There are integers $z_0,\dots,z_k\geq0$ with
  $\binom{\rho+P\beta}k=\sum_{i\leq k}z_i\binom\beta i$ for every integer $\beta\geq0$.
\item Let ``costs'' be a cost relation on the functions $\mathbf Z_{\geq0}^\sigma\to\mathcal K$, $\kappa$ a finite set,
  and $u_x\colon\mathbf Z_{\geq0}^\sigma\to\mathbf Z_{\geq0}$, for $x\in\kappa$, functions such that
  $\alpha\mapsto\binom{u_x(\alpha)}k$ costs $k$ for every $x\in\kappa$ and every $k\geq0$.
  \begin{enumerate}
  \item For all integers $\rho,P,m\geq0$ and $x\in\kappa$, $\alpha\mapsto\binom{\rho+Pu_x(\alpha)}m$ costs $m$.
  \item Let $n\geq0$ be an integer, $M$ a finite set of functions $m\colon\kappa\to\mathbf Z_{\geq0}$ with
    $\sum_xm(x)\leq n$, and $\lambda_m\in\mathcal K$ with $v(\lambda_m)\geq0$ for $m\in M$. Then the function
    $\alpha\mapsto\sum_{m\in M}\lambda_m\prod_{x\in\kappa}\binom{u_x(\alpha)}{m(x)}$ costs $n$, and its values have
    valuation at least $0$.
  \end{enumerate}
\item Suppose that $\mathcal K$ with $v$ satisfies (V) of Definition~\ref{pd-symbols-test-chains-projection-properties}
  at a prime $p\geq5$. Then for every $(y,a)\in\Lambda\times\{1,\dots,r\}$ and every integer $k\geq0$, the function
  $\mu\mapsto\binom{\mu(y,a)}k$ costs $k$ on the set of all subsets of $\Lambda$.
\end{enumerate}
\end{lemma}
\leanref{Formalism/SharpRestriction/GrowthBounds/MonomialOrderSegments/SlopeOfBinomCombGen.lean}{Formalism.MainNotions.PlaneGram.CostOK}
\leanref{Formalism/SharpRestriction/GrowthBounds/MonomialOrderSegments/SlopeOfBinomCombGen.lean}{Formalism.MainNotions.PlaneGram.CostV}
\leanref{Formalism/SharpRestriction/GrowthBounds/MonomialOrderSegments/SlopeOfBinomCombGen.lean}{Formalism.MainNotions.PlaneGram.costV_ok}
\leanref{Formalism/SharpRestriction/GrowthBounds/MonomialOrderSegments/SlopeOfBinomCombGen.lean}{Formalism.MainNotions.PlaneGram.CostOK.zero}
\leanref{Formalism/SharpRestriction/GrowthBounds/MonomialOrderSegments/SlopeOfBinomCombGen.lean}{Formalism.MainNotions.PlaneGram.CostOK.sum}
\leanref{Formalism/SharpRestriction/GrowthBounds/MonomialOrderSegments/SlopeOfBinomCombGen.lean}{Formalism.MainNotions.PlaneGram.choose_affine}
\leanref{Formalism/SharpRestriction/GrowthBounds/MonomialOrderSegments/SlopeOfBinomCombGen.lean}{Formalism.MainNotions.PlaneGram.cost_choose_affine}
\leanref{Formalism/SharpRestriction/GrowthBounds/MonomialOrderSegments/SlopeOfBinomCombGen.lean}{Formalism.MainNotions.PlaneGram.cost_prod_choose_gen}
\leanref{Formalism/SharpRestriction/GrowthBounds/MonomialOrderSegments/SlopeOfBinomCombGen.lean}{Formalism.MainNotions.PlaneGram.approx_of_trunc_gen}
\leanref{Formalism/SharpRestriction/GrowthBounds/MonomialOrderSegments/AddMemSupportOfMemSupportCoefOut.lean}{Formalism.MainNotions.PlaneGram.vala_dg_choose}
\leanref{Formalism/SharpRestriction/GrowthBounds/MonomialOrderSegments/PrimedEmbeddingInjective.lean}{Formalism.MainNotions.PlaneGram.costM_ok}
\leanref{Formalism/SharpRestriction/GrowthBounds/MonomialOrderSegments/PrimedEmbeddingInjective.lean}{Formalism.MainNotions.PlaneGram.costM_choose}

\begin{proof}
(1) For functions $g,h$, $\lambda\in\mathcal K$ and $F\in\mathcal K[z]$, $\operatorname{diag}_1F=F$,
$\operatorname{diag}_{gh}F=\operatorname{diag}_g(\operatorname{diag}_hF)$,
$\operatorname{diag}_{g+h}F=\operatorname{diag}_gF+\operatorname{diag}_hF$ and
$\operatorname{diag}_{\lambda g}F=\lambda\operatorname{diag}_gF$, by comparing coefficients. Let $S\in\mathcal F$,
$\theta$ real and $\operatorname{val}_S(F)\geq\theta$. The constant $1$ costs $0$, as $\operatorname{diag}_1F=F$. If $h$
costs $l$ and $g$ costs $k$, then $\operatorname{val}_S(\operatorname{diag}_hF)\geq\theta-l$, and so
$\operatorname{val}_S(\operatorname{diag}_g\operatorname{diag}_hF)\geq\theta-l-k$. If $g$ and $h$ cost $n$, then
$\operatorname{val}_S(\operatorname{diag}_gF+\operatorname{diag}_hF)\geq\theta-n$ and
$\operatorname{val}_S(\lambda\operatorname{diag}_gF)\geq\theta-n$, by
Lemma~\ref{pd-symbols-test-chains-projection-model-level-val}(1). If $g$ costs $n$ and $n'\geq n$, then
$\theta-n'\leq\theta-n$.

(2) Write $(1+X)^P=1+XQ$ with $Q\in\mathbf Z_{\geq0}[X]$. Then $\binom{\rho+P\beta}k$ is the coefficient of $X^k$ in
\[
 (1+X)^\rho(1+XQ)^\beta=\sum_i\binom\beta i(1+X)^\rho(XQ)^i ,
\]
and $(XQ)^i$ has no term of degree less than $i$. So $z_i$, the coefficient of $X^k$ in $(1+X)^\rho(XQ)^i$, will do.

(3) The zero function is $0\cdot1$, so it costs $0$, and hence every $n$; by induction, a finite sum of functions that
cost $n$ costs $n$, and a finite product of functions that cost $n_1,\dots,n_l$ costs $n_1+\dots+n_l$. (a) By (2),
$\binom{\rho+Pu_x(\alpha)}m=\sum_{i\leq m}z_i\binom{u_x(\alpha)}i$. Each $\binom{u_x(\alpha)}i$ costs $i$, hence $m$, and
so does $z_i\binom{u_x(\alpha)}i$; the sum costs $m$. (b) For $m\in M$ the product $\prod_x\binom{u_x(\alpha)}{m(x)}$
costs $\sum_xm(x)$, hence $n$; so does its multiple by $\lambda_m$, and so does the sum over $M$. Its values are sums
of products of the $\lambda_m$ with integers, and $v(n')\geq0$ for every integer $n'$, as $v(1)=0$.

(4) Let $e$ be the exponent vector with $e(y,a)=k$ and $0$ elsewhere, and $g(\mu)=\binom{\mu(y,a)}k$. As
$D_e(z^\mu)=\binom{\mu(y,a)}kz^{\mu-e}$ if $\mu(y,a)\geq k$ and $D_e(z^\mu)=0$ otherwise,
$\operatorname{diag}_gF=z^eD_eF$. Let $S\subset\Lambda$ and $\operatorname{val}_S(F)\geq\theta$. If $y$ is a vector
letter, then $\operatorname{val}_S(z^e)\geq0$ by (V)(a) and
Lemma~\ref{pd-symbols-test-chains-projection-model-level-val}(1), and $\operatorname{val}_S(D_eF)\geq\theta-k$ by
(V)(b). If $y$ is a covector letter, then $\operatorname{val}_S(z^e)\geq-k$ and $\operatorname{val}_S(D_eF)\geq\theta$.
In both cases $\operatorname{val}_S(z^eD_eF)\geq\theta-k$, by part (1) of that lemma.

For $\mathcal F=\{S\}$, (1), (3)(b) with the functions $u_{(y,a)}(\mu)=\mu(y,a)$, and (4) are the facts used in the
proof of Lemma~\ref{pd-symbols-test-chains-projection-diagonal}, where the cost is stated with $\theta=\operatorname{val}_S(F)$.
The two forms agree, as $\operatorname{val}_S$ takes values in $\mathbf Z\cup\{+\infty\}$ and
$\operatorname{diag}_g0=0$.
% Lean: a cost relation is CostOK T for a predicate T n g (its fields one, mul, add, smul, mono); "costs n on F" is
%   CostV sg F n g, and (1) is costV_ok, with vala_add and vala_smul (as costM_ok does for F = Set.univ). (2) is
%   choose_affine (one_add_X_pow_eq). (3): the zero function, sums and products are CostOK.zero, CostOK.sum and
%   CostOK.prod (cost_prod); (a) is cost_choose_affine and (b) is approx_of_trunc_gen, through cost_prod_choose_gen
%   (approx_of_trunc_gen also carries a function f within c of the combination, which the book does at each use).
%   (4) is vala_dg_choose (dg_choose_eq, vala_ddiv_letter' from V2, vala_X, vala_pow and vala_mul), giving costM_choose
%   for CostM D = CostV D.msg Set.univ. Lemma pd-symbols-test-chains-projection-diagonal is approx_of_trunc, the
%   case u = the coordinates; stepB_of applies it with costV_ok for Set.univ, which the book's per-S form covers.
\end{proof}

\begin{lemma}\label{pd-symbols-test-chains-projection-omega-units}
Let $p\geq5$, $\mathcal R$, $\mathcal K$ and $v$ be as in Lemma~\ref{pd-symbols-test-chains-projection-lift}, and suppose
that $\mathcal K$ with $v$ satisfies (V) of Definition~\ref{pd-symbols-test-chains-projection-properties} at $p$. Let
$r,J\geq1$ be integers, $\mathcal D$ a signed chain datum of rank $r$ with $J$ levels over $\mathcal R$ whose $c$ is an
integer, and $c_t$, $\gamma_t$ and $\gamma^{\mathrm P}$ as in
Definition~\ref{pd-symbols-test-chains-projection-omega-blocks}(4). Then there are functions $g_+$ and $g_-$ from the
monomials of $\mathcal K[\tilde x]$ to $\mathcal R$ such that:
\begin{enumerate}
\item if $\lambda_t(\zeta)(a)<P(t)$ for every link $t$ and every $a$, then $\gamma^{\mathrm P}(\zeta)$ is a unit of
  $\mathcal R$, $v(g_+(\zeta)-\gamma^{\mathrm P}(\zeta))\geq c$ and $v(g_-(\zeta)-\gamma^{\mathrm P}(\zeta)^{-1})\geq c$;
\item $\operatorname{val}_S(\operatorname{diag}_{g_\pm}Z)\geq\operatorname{val}_S(Z)-12c$ for every $Z\in\mathcal
  K[\tilde x]$ and every set $S$ of model positions that contains no label.
\end{enumerate}
\end{lemma}
\leanref{Formalism/SharpRestriction/GrowthBounds/MonomialOrderSegments/CollapseExponent.lean}{Formalism.MainNotions.PlaneGram.StepG}
\leanref{Formalism/SharpRestriction/GrowthBounds/MonomialOrderSegments/PrimedEmbeddingInjective.lean}{Formalism.MainNotions.PlaneGram.stepG_of}

\begin{proof}
We apply Lemmas~\ref{pd-symbols-test-chains-projection-model-level-val}
and~\ref{pd-symbols-test-chains-projection-omega-cost} to the model positions, the letters with their types and the
others as vector letters. In this proof ``$g$ costs $n$'' means that $g$ costs $n$ on the set of all sets of model
positions, a cost relation by Lemma~\ref{pd-symbols-test-chains-projection-omega-cost}(1); by part (4) of that lemma,
$\zeta\mapsto\binom{\zeta(y,a)}k$ costs $k$ for every model position $y$, every $a\leq r$ and every $k\geq0$.

\emph{Step 1: fixed label exponents.} Let $\rho=(\rho_t)$ be a family as in
Definition~\ref{pd-symbols-test-chains-projection-omega-blocks}(2) and $\epsilon=\pm1$. If $\rho_t(a)<P(t)$ for all
$t$ and $a$, we find $g_{\rho,\epsilon}$, from the monomials to $\mathcal R$, that costs $12c$ and satisfies
\[
 v\Big(g_{\rho,\epsilon}(\zeta)-\gamma(\operatorname{emb}_\rho\zeta)^\epsilon\prod_t\gamma_t\big(\rho_t,w_t(\zeta)\big)^\epsilon
 \Big)\geq c\qquad\text{for every }\zeta;
\]
otherwise we put $g_{\rho,\epsilon}=1$, which costs $0$, hence $12c$. For a position $y$ and $a\leq r$ put
$u_{(y,a)}(\zeta)=\operatorname{emb}_\rho(\zeta)(y,a)$, which is $\zeta(y,a)$ if $y$ is a letter and
$\rho_t(a)+P(t)w_t(\zeta)(a)$ if $y$ is a link $t$, with $w_t(\zeta)(a)$ the exponent of $\zeta$ at the variable
$\tilde x_{w,a}$ of the lift $w$ of $t$. By part (4) of that lemma, and by part (3)(a) for a link, each
$\zeta\mapsto\binom{u_{(y,a)}(\zeta)}k$ costs $k$. By (4) of Definition~\ref{pd-symbols-test-chains} with $k=c$
there are $\lambda_\mu\in\mathcal R$, for finitely many $\mu$ of total degree at most $6c$, with
$v(\gamma(\alpha)^\epsilon-\sum_\mu\lambda_\mu\prod_x\binom{\alpha(x)}{\mu(x)})\geq c$ for every $\alpha$. Put
$g_1(\zeta)=\sum_\mu\lambda_\mu\prod_x\binom{u_x(\zeta)}{\mu(x)}$, the value at $\alpha=\operatorname{emb}_\rho(\zeta)$
of this combination; it has values in $\mathcal R$ and costs $6c$, by part (3)(b) of that lemma. The second property
of (D) for $\rho$, $c$ and $\epsilon$ gives $\lambda'_{\mu'}\in\mathcal R$ for finitely many $\mu'$ with
$\sum\mu'\leq6c$, and $g_2(\zeta)=\sum_{\mu'}\lambda'_{\mu'}\prod_{t,a}\binom{w_t(\zeta)(a)}{\mu'(t,a)}$ satisfies
$v((\prod_t\gamma_t(\rho_t,w_t(\zeta)))^\epsilon-g_2(\zeta))\geq c$; it has values in $\mathcal R$ and costs $6c$, by
part (3)(b) with the functions $\zeta\mapsto w_t(\zeta)(a)$, which cost as required by part (4). Put
$g_{\rho,\epsilon}=g_1g_2$, which costs $6c+6c=12c$ by the product rule. With
$u=(\prod_t\gamma_t(\rho_t,w_t(\zeta)))^\epsilon$, of valuation $0$ by (D),
\[
 \gamma(\operatorname{emb}_\rho\zeta)^\epsilon u-g_1g_2=\big(\gamma(\operatorname{emb}_\rho\zeta)^\epsilon-g_1\big)u
 +g_1(u-g_2)
\]
has valuation at least $c$.

\emph{Step 2: the functions.} Put $g_+(\zeta)=g_{\lambda(\zeta),1}(\zeta)$ and $g_-(\zeta)=g_{\lambda(\zeta),-1}(\zeta)$.
(1) If $\lambda_t(\zeta)(a)<P(t)$ for all $t$ and $a$, then $\operatorname{emb}_{\lambda(\zeta)}(\zeta)=\zeta^{\mathrm c}$,
so the product in Step 1 is $\gamma^{\mathrm P}(\zeta)^{\pm1}$, and $\gamma^{\mathrm P}(\zeta)$ is a unit by (4) of
Definition~\ref{pd-symbols-test-chains} and by (D). (2) Let $S$ contain no label. Write $Z=\sum_\mu\mu Z_\mu$, with
$\mu$ running over the monomials in the variables of the labels and $Z_\mu$ free of them. The monomials of $\mu Z_\mu$
have the label exponents $\lambda(\mu)$, so $\operatorname{diag}_{g_+}(\mu Z_\mu)=\operatorname{diag}_{g_{\lambda(\mu),1}}
(\mu Z_\mu)$, of $\operatorname{val}_S$ at least $\operatorname{val}_S(\mu Z_\mu)-12c$. By
Lemma~\ref{pd-symbols-test-chains-projection-model-level-val}(1), (2) and (4), the last for $O$ the labels,
$\operatorname{val}_S(\mu Z_\mu)\geq0+\operatorname{val}_S(Z)$. Summing over $\mu$ gives (2) for $g_+$, and in the same
way for $g_-$.
% Lean: stepG_of hV2 (PrimedEmbeddingInjective.lean). The cost relation is CostM D = CostV D.msg Set.univ, with
%   costM_ok and costM_choose (from vala_dg_choose). Step 1 is exists_approx_gaR: gaR p gam rho eps is the product
%   of Step 1, embE p rho is emb_rho and embU p rho the functions u; g1 and g2 come from approx_of_trunc_gen with the
%   clause GaOK of Valid (hu from costM_choose and cost_choose_affine) and the clause of DigOK (hu from costM_choose),
%   and are multiplied (CostOK.mul, 6c + 6c = 12c); the other case is CostOK.mono with CostOK.one. Step 2: gp and
%   gm are gp (labF zeta) zeta and gm (labF zeta) zeta; (1) is gaR_one and gaR_neg_one (and v_gaP for the unit); (2) is
%   vala_dg_labels (sum_coefOut over the labels, dg_congr, valS_monomial_nonneg_of_free, vala_coefOut).
\end{proof}

\begin{lemma}\label{pd-symbols-test-chains-projection-omega-projection}
Let $p\geq5$, $\mathcal R$, $\mathcal K$ and $v$ be as in Lemma~\ref{pd-symbols-test-chains-projection-lift}, and suppose
that $\mathcal K$ with $v$ satisfies (P) of Definition~\ref{pd-symbols-test-chains-projection-properties} at $p$. Let
$r,J\geq1$ be integers and $\mathcal D$ a signed chain datum of rank $r$ with $J$ levels over $\mathcal R$, with the
notation of Definitions~\ref{pd-symbols-test-chains} and~\ref{pd-symbols-test-chains-projection-omega-blocks}. Then
there is a $\mathcal K$-linear endomorphism $\Pi$ of $\mathcal K[\tilde x]$ such that:
\begin{enumerate}
\item for every monomial $\zeta$, the monomials of $\Pi(\zeta)$ are of the same block as $\zeta$;
\item $v_{\mathrm F}(\Pi(\zeta))\geq v(\zeta!)/2-D_m$ for every good monomial $\zeta$
  (Definition~\ref{pd-symbols-test-chains-projection-good});
\item $\operatorname{val}_{\tilde C_{j,i}}(\Pi Z)\geq\theta_{j,i}-s_{j,i}$ for every $Z\in\mathcal K[\tilde x]$ and every
  chain $C_{j,i}$;
\item $\langle A,\Pi(\zeta)\rangle_{\mathrm F}=\langle A,\zeta\rangle_{\mathrm F}$ for every monomial $\zeta$ and every
  $A\in\mathcal K[\tilde x]$ with the row degrees of $\zeta$ and $\operatorname{val}_{\tilde C_{j,i}}(A)\geq
  \theta_{j,i}-s_{j,i}$ for every chain;
\item if $Z\in\mathcal K[\tilde x]$, $j$ is a level, $m<r$ and $\theta$ is real with
  $\operatorname{val}_{\tilde S_{j,m}}(Z)\geq\theta$, then $\operatorname{val}_{\tilde S_{j,m}}(\Pi Z)\geq\theta$.
\end{enumerate}
\end{lemma}
\leanref{Formalism/SharpRestriction/GrowthBounds/MonomialOrderSegments/CollapseExponent.lean}{Formalism.MainNotions.PlaneGram.StepP}
\leanref{Formalism/SharpRestriction/GrowthBounds/MonomialOrderSegments/SubstitutedExponentApplyBlock.lean}{Formalism.MainNotions.PlaneGram.stepP_of}

\begin{proof}
We use Lemma~\ref{pd-symbols-test-chains-projection-fock-product} for the model positions, the letters with their types
and the others as vector letters, with $X$ the set of the chains $C_{j,i}$, $S_x=\tilde C_{j,i}$ and
$\theta_x=\lceil\theta_{j,i}-s_{j,i}\rceil$. Its hypothesis holds: for $i\leq i'$, $C_{j,i}\subset C_{j,i'}$, so
$\tilde C_{j,i}\subset\tilde C_{j,i'}$; and for $j\neq j'$, $\tilde C_{j,i}$ and $\tilde C_{j',i'}$ are disjoint, as they
consist of letters of the levels $j$ and $j'$ and of lifts of the levels $j-1$ and $j'-1$. For each
$d\colon\{\text{model positions}\}\to\mathbf Z_{\geq0}$ that lemma gives $\Pi_d$. Let $\Pi$ be the $\mathcal K$-linear map
with $\Pi(\zeta)=\Pi_{d(\zeta)}(\zeta)$, $d(\zeta)$ the row degrees of $\zeta$; then $\Pi F=\Pi_dF$ for $F\in\mathrm{Pol}(d)$.
As $\operatorname{val}_S$ takes integer values, $\operatorname{val}_S\geq\theta$ is $\operatorname{val}_S\geq\lceil\theta\rceil$.

(1) By (1) of that lemma $\Pi(\zeta)$ has the row degrees of $\zeta$. The labels and the tops are vector positions in no
$\tilde C_{j,i}$, so by its (5) the monomials of $\Pi(\zeta)$ have the label and top exponents of $\zeta$.

(2) By (3) of that lemma, $v_{\mathrm F}(\Pi(\zeta))\geq v(\zeta!)/2-|X|\,r^2\ell(2\sum_yd(\zeta)(y))$, with
$|X|=\sum_jn^c_j$. For good $\zeta$ the sum of its row degrees is at most $\sum_yd(y)$, the sum over the positions: at a
link $t$, $e_{\mathrm{lab}}+e_{\mathrm{lift}}+e_{\mathrm{top}}\leq e_{\mathrm{lab}}+P(t)e_{\mathrm{lift}}+e_{\mathrm{top}}
=d(t)$. As $\ell$ is nondecreasing, $|X|\,r^2\ell(2\sum_yd(\zeta)(y))\leq e_P\sum_jn^c_j=D_m$.

(3) By (1) of that lemma each $\Pi(\zeta)$ has $\operatorname{val}_{\tilde C_{j,i}}\geq\theta_{j,i}-s_{j,i}$, and so does
$\Pi Z$, by Lemma~\ref{pd-symbols-test-chains-projection-model-level-val}(1).

(4) This is (2) of that lemma, with $G=A$ and $F=\zeta$.

(5) Write $Z=\sum_dZ_d$ with $Z_d$ of row degrees $d$. By
Lemma~\ref{pd-symbols-test-chains-projection-model-level-val}(6), $\operatorname{val}_{\tilde S_{j,m}}(Z_d)\geq\theta$.
Each $\tilde C_{j',i}$ is contained in $\tilde S_{j,m}$ (if $j'=j$: letters of level $j$ and lifts of level $j-1$) or
disjoint from it (if $j'\neq j$). So (4) of Lemma~\ref{pd-symbols-test-chains-projection-fock-product} gives
$\operatorname{val}_{\tilde S_{j,m}}(\Pi_dZ_d)\geq\lceil\theta\rceil$, and the sum $\Pi Z=\sum_d\Pi_dZ_d$ satisfies the
same bound.
% Lean: stepP_of v p C0 h1 hPROJ (SubstitutedExponentApplyBlock.lean): prodProj for every d with Sx = chainSet and
%   thx = ceil(th - sl), the pairs by chainSet_pair; Pj is the linear map with Pj (monomial zeta 1) = Pr (ldeg zeta)
%   (monomial zeta 1), and hhom. (1) is prodProj_entries with mlab_notMem_chainSet and mtp_notMem_chainSet; (2) is the
%   clause bge of prodProj with sum_ldeg_le and lg_mono'; (3) Aset_sum and Aset_smul; (4) the Fock clause; (5) sum_hcomp,
%   vala_hcomp, chainSet_Yset and the clause for S''.
\end{proof}

\begin{lemma}\label{pd-symbols-test-chains-projection-omega-expansion}
Let $p\geq5$, $\mathcal R$, $\mathcal K$ and $v$ be as in Lemma~\ref{pd-symbols-test-chains-projection-lift}, and suppose
that $\mathcal K$ with $v$ satisfies (V) and (E) of Definition~\ref{pd-symbols-test-chains-projection-properties} at
$p$. Let $r,J\geq1$ be integers, $\mathcal D$ a signed chain datum of rank $r$ with $J$ levels over $\mathcal R$ whose $c$
is an integer and that satisfies (H) for $13$, and $F$ an admissible polynomial with $v_{\mathrm F}(F)\geq0$. Let $M$ be an
exponent vector of the positions that vanishes at the letters and has degree at most $d_L(j)$ in the row of every link
$(j,m)$, let $D_M$ be the divided derivative and $\operatorname{diag}_g$ as in
Lemma~\ref{pd-symbols-test-chains-projection-diagonal}, and
\[
 G=\operatorname{diag}_{\gamma^{-1}}D_M\operatorname{diag}_\gamma F .
\]
Then for every family $\rho$ of residues as in Definition~\ref{pd-symbols-test-chains}, $G^\rho=Q_0+E$ with
$\operatorname{val}_{\tilde C_{j,i}}(Q_0)\geq\theta_{j,i}-s_{j,i}$ for every chain and every coefficient of $E$ of
valuation at least $c-n_0$.
\end{lemma}
\leanref{Formalism/SharpRestriction/GrowthBounds/MonomialOrderSegments/CollapseExponent.lean}{Formalism.MainNotions.PlaneGram.StepE}
\leanref{Formalism/SharpRestriction/GrowthBounds/MonomialOrderSegments/SumOnLetterApplyInl.lean}{Formalism.MainNotions.PlaneGram.stepE_of}

\begin{proof}
For a function $\theta'$ on the chains and a real number $a$, call $G'\in\mathcal K[x]$ \emph{split for $(\theta',a)$} if
$G'=G_0+E_0$, where the coefficients of $G_0$ have valuation at least $a$, those of $E_0$ at least $a+c$, and
$\operatorname{val}_{\tilde C_{j,i}}(G_0^\rho)\geq\theta'_{j,i}$ for every family $\rho$ and every chain. As $F\mapsto
F^\rho$ only selects and renames terms, $(G_0+E_0)^\rho=G_0^\rho+E_0^\rho$, and the coefficients of $E_0^\rho$ are
coefficients of $E_0$. We use Lemma~\ref{pd-symbols-test-chains-projection-model-level-val} for the model positions.

\emph{Step 1: the start.} The polynomial $F$ is split for $(\theta,-n_0)$, with $G_0=F$ and $E_0=0$: as
$v_{\mathrm F}(F)\geq0$, $v(F_\alpha)\geq-v(\alpha!)/2\geq-n_0$, since $\alpha!$ divides $\prod_yd(y)!$ for $\alpha$ of row
degrees $d$; and $\operatorname{val}_{\tilde C_{j,i}}(F^\rho)\geq\theta_{j,i}$ as $F$ is admissible.

\emph{Step 2: multiplication.} Let $G'$ be split for $(\theta',a)$, and $f$ a function on the exponent vectors of the
positions with $v(f)\geq0$ for which there are integral $\lambda_m$, for finitely many $m$ with $\sum m\leq6c$, with
$v(f(\alpha)-g(\alpha))\geq c$ for every $\alpha$, where $g(\alpha)=\sum_m\lambda_m\prod_x\binom{\alpha(x)}{m(x)}$. Then
$\operatorname{diag}_fG'$ is split for $(\theta'-6c,a)$. Indeed $\operatorname{diag}_fG'=\operatorname{diag}_gG_0+
(\operatorname{diag}_{f-g}G_0+\operatorname{diag}_fE_0)$; the first term has coefficients of valuation at least $a$, as
$g$ is integral, and the second at least $a+c$. For the bound, say that a function $h$ on the exponent vectors of the
positions \emph{costs $n$ after embedding} if for every family $\rho$ the function $\zeta\mapsto h(\operatorname{emb}_\rho
\zeta)$ on the monomials of $\mathcal K[\tilde x]$ costs $n$ on the set of all sets of model positions
(Lemma~\ref{pd-symbols-test-chains-projection-omega-cost}, for the model positions). This is a cost relation, as the
four rules of part (1) of that lemma hold for each $\rho$. For a position $y$ and $a\leq r$, $\alpha\mapsto
\binom{\alpha(y,a)}k$ costs $k$ after embedding: $\operatorname{emb}_\rho(\zeta)(y,a)$ is $\zeta(y,a)$ if $y$ is a letter,
and $\rho_t(a)+P(t)\zeta(w,a)$, $w$ the lift of $t$, if $y$ is a link $t$; use part (4) of that lemma, respectively
parts (3)(a) and (4). By part (3)(b), $g$ costs $6c$ after embedding. For a term $x^\alpha$ in the class of $\rho$, the
digit part is $\tilde x^{\alpha'}$ with $\operatorname{emb}_\rho(\alpha')=\alpha$; so $(\operatorname{diag}_gG_0)^\rho=
\operatorname{diag}_{g\circ\operatorname{emb}_\rho}(G_0^\rho)$. As $g\circ\operatorname{emb}_\rho$ costs $6c$ on the
set of all sets of model positions, the bound follows.

\emph{Step 3: one link.} Let $G'$ be split for $(\theta',a)$, $t$ a link of level $j_0$, $M_t\in\mathbf Z_{\geq0}^r$, and
$D$ the divided derivative by $M_t$ in the row of $t$. Then $DG'$ is split for $(\theta'',a)$, where
$\theta''_{j,i}=\theta'_{j,i}-|M_t|/P(t)-c$ if $j_0=j-1$ and $\theta''_{j,i}=\theta'_{j,i}$ otherwise. For each
$\rho_f\in\{0,\dots,P(t)-1\}^r$, (E) with $r'=r$, $g=g(t)$, $\rho_f$, $M_t$ and $c$ gives integral $e_s$ and
$\rho'_s,\delta_s,\kappa_s\in\mathbf Z_{\geq0}^r$ with $|\kappa_s|\leq|M_t|/P(t)+c$, for $s$ in a finite set $I(\rho_f)$,
such that, as $x^{\rho'_s}(x^{\delta_s}D_{\kappa_s}(x^\beta))^{[P(t)]}=\binom\beta{\kappa_s}x^{\rho'_s+P(t)(\delta_s+
\beta-\kappa_s)}$, every coefficient of
\[
 D_{M_t}\big(x^{\rho_f+P(t)\beta}\big)-\sum_{s\in I(\rho_f)}e_s\binom\beta{\kappa_s}x^{\rho'_s+P(t)(\delta_s+\beta-\kappa_s)}
\]
has valuation at least $c$, for every $\beta$; here $x$ is a single row, and $\binom\beta\kappa=\prod_a\binom{\beta(a)}
{\kappa(a)}$ is $0$ unless $\kappa\leq\beta$. For a monomial $x^\alpha$ of $\mathcal K[x]$ let $\alpha_t$ be its exponents in
the row of $t$, write $\alpha_t=\rho_f+P(t)\beta$ with $\rho_f(a)<P(t)$, and let $A(x^\alpha)$ be $x^\alpha$ with its row of
$t$ replaced by the sum above; extend $A$ linearly. As $D(x^\alpha)$ is $x^\alpha$ with the row of $t$ replaced by
$D_{M_t}(x^{\alpha_t})$, the coefficients of $D(x^\alpha)-A(x^\alpha)$ have valuation at least $c$. Put $G'_0=A(G_0)$ and
$E'_0=(DG_0-A(G_0))+DE_0$; then $DG'=G'_0+E'_0$, and the coefficients of $G'_0$ have valuation at least $a$ and those of
$E'_0$ at least $a+c$, as the $e_s$ and binomial coefficients are integral. For the digit parts, let $w$ be the lift of
$t$ and, for $s\in I(\rho_f)$, $\Theta_s$ the map $Q\mapsto e_s\tilde x_w^{\lfloor\rho'_s/P(t)\rfloor+\delta_s}
D_{\kappa_s,w}Q$, where $D_{\kappa_s,w}$ is the divided derivative by $\kappa_s$ in the row of $w$ and $\tilde x_w^\eta=
\prod_a\tilde x_{w,a}^{\eta(a)}$. Comparing the two sides on a monomial of $G_0$,
\[
 A(G_0)^\rho=\sum_{\rho_f}\ \sum_{s\in I(\rho_f),\ \rho'_s\equiv\rho_t}\Theta_s\big(G_0^{\rho[t\mapsto\rho_f]}\big),
\]
the congruence modulo $P(t)$ entrywise, where $\rho[t\mapsto\rho_f]$ is $\rho$ with $\rho_t$ replaced by $\rho_f$: a term
$x^\alpha$ of $G_0$ in the class of $\rho[t\mapsto\rho_f]$ has digit part $\tilde x^{\alpha'}$ with exponent
$\beta=\lfloor\alpha_t/P(t)\rfloor$ at $w$, and the digit part of its image under $A$ is $\sum_s\Theta_s(\tilde
x^{\alpha'})$, the sum over the $s$ with $\rho'_s\equiv\rho_t$. If $j_0=j-1$, then $w\in\tilde C_{j,i}$, and as $w$ is a
vector position, (V)(b) and Lemma~\ref{pd-symbols-test-chains-projection-model-level-val}(1) and (2) give
$\operatorname{val}_{\tilde C_{j,i}}(\Theta_sQ)\geq\operatorname{val}_{\tilde C_{j,i}}(Q)-|\kappa_s|$. If not,
$w\notin\tilde C_{j,i}$, and $\Theta_s(\mu\nu)=\nu\,\Theta_s(\mu)$ for monomials $\mu$ in the variables of $w$ and $\nu$
in the others, with $\operatorname{val}_{\tilde C_{j,i}}(\Theta_s\mu)\geq0$ by part (2) of that lemma; so part (5) gives
$\operatorname{val}_{\tilde C_{j,i}}(\Theta_sQ)\geq\operatorname{val}_{\tilde C_{j,i}}(Q)$. With part (1), $G'_0$
satisfies the bounds $\theta''$.

\emph{Step 4: all links.} Write $M=\sum_tM_t$, with $M_t$ supported in the row of the link $t$. Divided derivatives in
different rows compose, so $D_M$ is the composite of the $D$ of Step 3 for the links $t$, and by Step 3 $D_MG'$ is split
for $(\theta'',a)$ if $G'$ is split for $(\theta',a)$, where $\theta''_{j,i}=\theta'_{j,i}-\sum_{m<r}\big(|M_{(j-1,m)}|/
P(j-1,m)+c\big)$.

\emph{Step 5: the bound.} By (4) of Definition~\ref{pd-symbols-test-chains} with $k=c$ and $\epsilon=\pm1$, the functions
$\gamma$ and $\gamma^{-1}$ have the property of Step 2, and they are integral. By Steps 1, 2, 4 and 2, $G$ is split for
$(\theta'',-n_0)$ with $\theta''_{j,i}=\theta_{j,i}-12c-\sum_{m<r}(|M_{(j-1,m)}|/P(j-1,m)+c)$. Write $G=G_0+E_0$, and put
$Q_0=G_0^\rho$ and $E=E_0^\rho$; the coefficients of $E$ have valuation at least $c-n_0$. As $|M_{(j-1,m)}|\leq
d_L(j-1)$ and $\mathcal D$ satisfies (H) for $13$, $d_L(j-1)/P(j-1,m)\leq s_{j,i}/r-13c$, so
\[
 12c+\sum_{m<r}\Big(\frac{|M_{(j-1,m)}|}{P(j-1,m)}+c\Big)\leq12c+s_{j,i}-13rc+rc\leq s_{j,i},
\]
and $\operatorname{val}_{\tilde C_{j,i}}(Q_0)\geq\theta_{j,i}-s_{j,i}$.
% Lean: stepE_of v p hV2 hDER (SumOnLetterApplyInl.lean). Split for (theta', a) is ChOK v p theta' a; Step 1 is h0
%   (cge_of_bge and the clause of InXX); Step 2 is ChOK.dg (approx_of_trunc with the cost relation CostA D p, "costs n
%   after embedding", which is CostOK by costA_ok from costM_ok; costA_choose: emb_inl with costM_choose for a letter,
%   emb_inr with cost_choose_affine for a link; digitPart_dg, emb_liftE, dg_sub_fun); Step 3 is ChOK.ddiv (der' from DER, appMon and appL with
%   cge_appL and cge_ddiv_sub_appL through ddiv_onLetter_monomial and putL; digitPart_appL with opS = Theta_s, liftE_new
%   and inClass_new; vala_opS (vala_ddiv_letter from V2, valS_nonneg_of_vec) and vala_opS_out (vala_outside_linear),
%   mw_not_mem_chainSet); Step 4 is ChOK.ddiv_finset and ChOK.ddiv_all (ddiv_ddiv_of_disjoint); Step 5 is the rest of
%   stepE_of with GaOK, gaAct_eq, digitPart_add, cge_digitPart and sum_ldeg_div_add_le_sl (hypH', nxt_prv).
\end{proof}

\begin{lemma}\label{pd-symbols-test-chains-projection-omega}
Let $p\geq5$, $\mathcal R$, $\mathcal K$ and $v$ be as in Lemma~\ref{pd-symbols-test-chains-projection-lift}, and suppose
that $\mathcal K$ with $v$ satisfies (V), (P), (D) and (E) of
Definition~\ref{pd-symbols-test-chains-projection-properties} at $p$. Let $r,J\geq1$ be integers, $\mathcal D$ a signed
chain datum of rank $r$ with $J$ levels over $\mathcal R$ whose $c$ is an integer, that satisfies (H) for $13$ and
$4(D_m+D_l+3\kappa_{\mathcal D})<c$ (Definition~\ref{pd-symbols-test-chains}), and $F$ an admissible polynomial with
$v_{\mathrm F}(F)=0$. Then there are a $\mathcal K$-linear endomorphism $\Omega$ of $\mathcal K[\tilde x]$ and an integer
$0\leq N_2\leq D_m+\kappa_{\mathcal D}/2+1$ such that for every $Z\in\mathcal K[\tilde x]$ good for $\mathcal D$
(Definition~\ref{pd-symbols-test-chains-projection-good}) the polynomial $\Omega(Z)$ is good for $\mathcal D$ and:
\begin{enumerate}
\item $\operatorname{val}_{\tilde C_{j,i}}(\Omega(Z))\geq\theta_{j,i}-s_{j,i}-12c$ for every chain $C_{j,i}$;
\item $\operatorname{val}_{\tilde S_{j,m}}(\Omega(Z))\geq\operatorname{val}_{\tilde S_{j,m}}(Z)-24c$ for all $j$ and
  $m+2\leq r$;
\item if $Z\in\mathcal R[\tilde x]$, then $p^{N_2}\Omega(Z)\in\mathcal R[\tilde x]$ and
  $v\big(B_{\mathcal D}(F,\varphi(\Omega(Z)))-B_{\mathcal D}(F,\varphi(Z))\big)\geq c+n_0-2\kappa_{\mathcal D}-D_m$.
\end{enumerate}
\end{lemma}
\leanref{Formalism/SharpRestriction/GrowthBounds/MonomialOrderSegments/CollapseExponent.lean}{Formalism.MainNotions.PlaneGram.StepDEF}
\leanref{Formalism/SharpRestriction/GrowthBounds/MonomialOrderSegments/OmegaValue.lean}{Formalism.MainNotions.PlaneGram.stepDEF_of}

\begin{proof}
We use Definition~\ref{pd-symbols-test-chains-projection-omega-blocks} and
Lemma~\ref{pd-symbols-test-chains-projection-model-level-val} for the model positions, and write $\kappa$ for
$\kappa_{\mathcal D}$. As in the proof of Lemma~\ref{pd-symbols-test-chains-projection-lift}, $v(n)=v_p(n)$ for integers
$n\geq1$ and $D_m,\kappa\geq0$.

\emph{Step 1: the maps.} Property (D) for $\tau$ the set of links, $r'=r$ and $t\mapsto g(t)$ gives $c_t$ and $\gamma_t$
as in Definition~\ref{pd-symbols-test-chains-projection-omega-blocks}(4).
Lemma~\ref{pd-symbols-test-chains-projection-omega-units} gives $g_\pm$, and
Lemma~\ref{pd-symbols-test-chains-projection-omega-projection} gives $\Pi$. Put
\[
 \Omega=\operatorname{diag}_{g_-}\circ\Pi\circ\operatorname{diag}_{g_+},\qquad N_2=\lceil D_m+\kappa/2\rceil ,
\]
so that $0\leq N_2\leq D_m+\kappa/2+1$.

\emph{Step 2: factorials.} Let $\zeta$ be a good monomial. For an exponent vector $\alpha$ of row degrees $e$, the
quotient $\prod_ye(y)!/\alpha!$ is a product of multinomial coefficients with $r$ parts, and by the proof of
Lemma~\ref{pd-symbols-flat-real-multinomial}(1) the valuation of each, and of each $\binom{e(y)}k$, is at most
$r\lfloor\log_pe(y)\rfloor$ if $e(y)\geq1$.

(a) \emph{If $\zeta'$ is of the same block as $\zeta$, then $v(\zeta'!)-v(\zeta!)\leq\kappa$.} Write $\zeta=\zeta_1\zeta_2$
with $\zeta_1$ the part in the labels and tops and $\zeta_2$ that in the letters and lifts, and likewise $\zeta'$. Then
$\zeta'_1=\zeta_1$, and $\zeta'_2$ and $\zeta_2$ have the same row degrees $e$, with $e(y)=d(y)$ at a letter and
$e(w)\leq d(t)$ at the lift $w$ of a link $t$, as $P(t)e(w)\leq d(t)$. So $v(\zeta'_2!)\leq v(\prod e!)\leq
v(\zeta_2!)+\sum r\log_p(d(y)+r)\leq v(\zeta_2!)+\kappa$, the sum over the letters and links $y$ with $e\geq1$ at the
corresponding model position.

(b) \emph{If $\lambda_t(\zeta)(a)<P(t)$ for all $t$ and $a$, then $(\zeta^{\mathrm c})!(\zeta^{\mathrm t})!=
\kappa_{\mathrm C}(\zeta)\prod_t\gamma_t(\lambda_t(\zeta),w_t(\zeta))\,\zeta!$, and $v(\kappa_{\mathrm C}(\zeta))\geq
2n_0-2\kappa-v(\zeta!)$.} The factor of $(\zeta^{\mathrm c})!$ at a link $t$ is $\prod_a(\lambda_t(\zeta)(a)+P(t)w_t(\zeta)
(a))!$, which by (D) is $c_t(\lambda_t(\zeta),|w_t(\zeta)|)\prod_aw_t(\zeta)(a)!\,\gamma_t(\lambda_t(\zeta),w_t(\zeta))$;
and $\zeta!$ is the product of the factorials of the exponents of the letters, labels, lifts and tops, $(\zeta^{\mathrm
t})!$ that of the tops. This gives the identity. As $\zeta$ is good, $\zeta^{\mathrm c}+\zeta^{\mathrm t}$ has row degrees
$d$, and $(\zeta^{\mathrm c}+\zeta^{\mathrm t})!=\prod_x\binom{\zeta^{\mathrm c}(x)+\zeta^{\mathrm t}(x)}{\zeta^{\mathrm
t}(x)}(\zeta^{\mathrm c})!(\zeta^{\mathrm t})!$; so $v((\zeta^{\mathrm c})!)+v((\zeta^{\mathrm t})!)\geq
v(\prod_yd(y)!)-2\kappa=2n_0-2\kappa$. As the $\gamma_t$ are units, the bound for $\kappa_{\mathrm C}(\zeta)$ follows.

(c) In the same way $v((\zeta^{\mathrm c}+\zeta^{\mathrm t})!)\geq2n_0-\kappa$.

\emph{Step 3: goodness, (1) and (2).} Let $Z$ be good. The maps $\operatorname{diag}_g$ do not enlarge supports, and by
Lemma~\ref{pd-symbols-test-chains-projection-omega-projection}(1) the monomials of $\Pi(\zeta)$ are of the same block as
$\zeta$, hence good when $\zeta$ is, as goodness depends only on the row degrees and the label and top exponents. So
$\Omega(Z)$ is good. The sets $\tilde C_{j,i}$ and $\tilde S_{j,m}$ contain no label. By (3) of that lemma,
$\Pi(\operatorname{diag}_{g_+}Z)$ has $\operatorname{val}_{\tilde C_{j,i}}\geq\theta_{j,i}-s_{j,i}$, and
Lemma~\ref{pd-symbols-test-chains-projection-omega-units}(2) gives (1). By the same lemma, (5) of
Lemma~\ref{pd-symbols-test-chains-projection-omega-projection} and the same lemma again,
$\operatorname{val}_{\tilde S_{j,m}}(\Omega(Z))\geq\operatorname{val}_{\tilde S_{j,m}}(Z)-24c$, which is (2).

\emph{Step 4: integrality.} Let $Z$ be also integral. Then $\operatorname{diag}_{g_+}Z$ is integral. For a good monomial
$\zeta$ and a monomial $\zeta'$ of $\Pi(\zeta)$, by (2) of
Lemma~\ref{pd-symbols-test-chains-projection-omega-projection} and Step 2(a),
$v(\Pi(\zeta)_{\zeta'})\geq v(\zeta!)/2-D_m-v(\zeta'!)/2\geq-D_m-\kappa/2$. So the coefficients of $\Omega(Z)$ have
valuation at least $-D_m-\kappa/2\geq-N_2$, as $g_-$ is integral, and $p^{N_2}\Omega(Z)\in\mathcal R[\tilde x]$.

\emph{Step 5: the pairing.} Let $\Lambda(W)=B_{\mathcal D}(F,\varphi(W))$ for $W\in\mathcal K[\tilde x]$; it is
$\mathcal K$-linear, and for a monomial $\zeta$, with $\alpha=\zeta^{\mathrm c}+\zeta^{\mathrm t}$,
$\Lambda(\zeta)=F_\alpha\gamma(\alpha)\alpha!$. As $Z$ is integral and $\Omega$ and $\Lambda$ are linear, it is enough to
show $v(\Lambda(\Omega\zeta)-\Lambda(\zeta))\geq c+n_0-2\kappa-D_m$ for a good monomial $\zeta$. Now
$\Omega(\zeta)=g_+(\zeta)\operatorname{diag}_{g_-}(\Pi\zeta)$ and $\operatorname{diag}_{g_-}\zeta=g_-(\zeta)\zeta$, so
\[
 \Lambda(\Omega\zeta)-\Lambda(\zeta)=g_+(\zeta)\big(\Lambda(\operatorname{diag}_{g_-}\Pi\zeta)-
 \Lambda(\operatorname{diag}_{g_-}\zeta)\big)+\big(g_+(\zeta)g_-(\zeta)-1\big)\Lambda(\zeta).
\]
With $u=\gamma^{\mathrm P}(\zeta)$, a unit, $g_+g_--1=(g_+-u)g_-+u(g_--u^{-1})$ has valuation at least $c$ by
Lemma~\ref{pd-symbols-test-chains-projection-omega-units}(1). As $v_{\mathrm F}(F)\geq0$, $v(\Lambda(\zeta))\geq
v(\alpha!)/2\geq n_0-\kappa/2$ by Step 2(c). So the second term has valuation at least $c+n_0-\kappa/2$. For the first
term, as $g_+$ is integral, it remains to bound the difference in parentheses.

(i) \emph{An identity.} Put $M=\zeta^{\mathrm t}$; it vanishes at the letters, and its degree in the row of a link
$(j,m)$ is that of $\zeta$ in the row of the top, at most $d_L(j)$. Let $G=\operatorname{diag}_{\gamma^{-1}}D_M
\operatorname{diag}_\gamma F$, so that $G_\beta=\gamma(\beta)^{-1}\gamma(\beta+M)F_{\beta+M}\binom{\beta+M}M$ with
$\binom{\beta+M}M=\prod_x\binom{\beta(x)+M(x)}{M(x)}$. Let $\zeta'$ be of the same block as $\zeta$; then $\zeta'^{\mathrm
t}=M$, $\zeta'$ is good and $\kappa_{\mathrm C}(\zeta')=\kappa_{\mathrm C}(\zeta)$. As $(\zeta'^{\mathrm c}+M)!=
\binom{\zeta'^{\mathrm c}+M}M(\zeta'^{\mathrm c})!M!$, Step 2(b) gives
\[
 \Lambda(\zeta')=\kappa_{\mathrm C}(\zeta)\,\gamma^{\mathrm P}(\zeta')\,G_{\zeta'^{\mathrm c}}\,\zeta'! .
\]
(ii) \emph{The expansion.} Lemma~\ref{pd-symbols-test-chains-projection-omega-expansion}, with $M$ and
$\rho=\lambda(\zeta)$, gives $G^\rho=Q_0+E$. Let $\mu$ be the product of the powers of the label and top variables in
$\zeta$, and $A$ the sum of the terms of $\mu Q_0$ of the row degrees of $\zeta$. As $\mu$ involves only vector positions,
parts (2), (1) and (6) of Lemma~\ref{pd-symbols-test-chains-projection-model-level-val} give
$\operatorname{val}_{\tilde C_{j,i}}(A)\geq\theta_{j,i}-s_{j,i}$. Write $\zeta'=\mu\zeta'_2$, with $\zeta'_2$ in the
letters and lifts. As $(\mu Q_0)_{\mu\nu}=(Q_0)_\nu$ for every monomial $\nu$, $A_{\zeta'}=(Q_0)_{\zeta'_2}$; and the
term of $G$ at $\zeta'^{\mathrm c}$ is in the class of $\rho=\lambda(\zeta')$, with digit part $\tilde x^{\zeta'_2}$, so
$G_{\zeta'^{\mathrm c}}=(Q_0)_{\zeta'_2}+E_{\zeta'_2}$. Hence
\[
 e(\zeta')=g_-(\zeta')\gamma^{\mathrm P}(\zeta')G_{\zeta'^{\mathrm c}}-A_{\zeta'}
 =\big(g_-(\zeta')-\gamma^{\mathrm P}(\zeta')^{-1}\big)\gamma^{\mathrm P}(\zeta')G_{\zeta'^{\mathrm c}}+E_{\zeta'_2}
\]
has valuation at least $c-n_0$, as $v(G_\beta)\geq v(F_{\beta+M})\geq-n_0$ (Step 1 of the proof of
Lemma~\ref{pd-symbols-test-chains-projection-omega-expansion}).

(iii) \emph{The difference.} If every monomial of $W$ is of the same block as $\zeta$, then by (i)
\[
 \Lambda(\operatorname{diag}_{g_-}W)=\sum_{\zeta'}W_{\zeta'}g_-(\zeta')\Lambda(\zeta')=\kappa_{\mathrm C}(\zeta)\Big(
 \langle A,W\rangle_{\mathrm F}+\sum_{\zeta'}W_{\zeta'}e(\zeta')\,\zeta'!\Big).
\]
This applies to $W=\Pi\zeta$, by (1) of Lemma~\ref{pd-symbols-test-chains-projection-omega-projection}, and to
$W=\zeta$; and by (4) of that lemma $\langle A,\Pi\zeta\rangle_{\mathrm F}=\langle A,\zeta\rangle_{\mathrm F}$. So the
difference in parentheses is a sum of terms $\kappa_{\mathrm C}(\zeta)\,w\,e(\zeta')\,\zeta'!$ with $\zeta'$ of the same
block as $\zeta$ and $w=\Pi(\zeta)_{\zeta'}$ or $w=-1$, $\zeta'=\zeta$; in both cases $v(w)\geq v(\zeta!)/2-D_m-
v(\zeta'!)/2$, by (2) of that lemma and as $D_m\geq0$. By Step 2(b) for $\zeta$ and for $\zeta'$,
$v(\kappa_{\mathrm C}(\zeta))\geq2n_0-2\kappa-\min(v(\zeta!),v(\zeta'!))$. So the term has valuation at least
\[
 2n_0-2\kappa-\min\big(v(\zeta!),v(\zeta'!)\big)+\frac{v(\zeta!)-v(\zeta'!)}2-D_m+c-n_0+v(\zeta'!)\geq
 c+n_0-2\kappa-D_m .
\]
% Lean: stepDEF_of h1 hDIG (stepG_of hV2) (stepP_of v p C0 h1 hPROJ) (stepE_of v p hV2 hDER) (OmegaValue.lean):
%   Omega = dg gm o Pj o dg gp and N2 = ceil(Dm + kapD / 2). Step 1 is digOK_exists (GlobalTheory.lean). Step 2: (a) is
%   blk_fact (fact_eq_out_mul_in with outP, fact_upper, fact_lower, kapL_mono with dUp, kapL_le_kap'); (b) is
%   fact_eq_kapC, vge_kapC and wt_lower (fact_add_lower); (c) is fact_phE_lower. Steps 3 and 4 are omega_clauses
%   (MlabNotMemYset.lean: good_Pj, Good.dg, mlab_notMem_chainSet, mlab_notMem_Yset, cge_Pj_monomial, cge_dg).
%   Step 5 is omega_value (Lam_monomial, the split hexp, v_gaP, hsplit, vge_Lam_monomial) with
%   TwistData.value_monomial (SubstitutedExponentApplyLabel.lean) for (i) to (iii): GM F M is G, Lam_monomial_blk and
%   kappa_constant_same_block for (i); A = hcomp (ldeg zeta) (monomial mu 1 * Q0), coeff_digitPart_eIn (class_iff),
%   vge_coeff_GM and herr for (ii); key, hP4, hterm for (iii).
\end{proof}

\begin{lemma}\label{pd-symbols-test-chains-projection}
Let $p\geq5$ be a prime, $n'\geq2$ an even integer, and $\mathcal R$ a complete discrete valuation ring of
characteristic $0$ in which $p$ is not a unit and whose residue field contains a field of $p^{n'}$ elements; we
use the notation of Definition~\ref{pd-symbols-test-chains} for $\mathcal R$. Let $r,J\geq1$ be integers,
$\mathcal D$ a signed chain datum of rank $r$ with $J$ levels over $\mathcal R$ that satisfies (H) for $C_H=26$
and $4(D_m+D_l+3\kappa_{\mathcal D})<c$, and $F$ an admissible polynomial with $v_{\mathrm F}(F)=0$. Then there
is a projection pair for $(\mathcal D,F)$ and $(2,24,60)$.
\end{lemma}
\leanref{Formalism/SharpRestriction/GrowthBounds/Lemma5aAt.lean}{Formalism.MainNotions.PlaneGramData.ProjectionDerivativeEstimate}
\leanref{Formalism/SharpRestriction/GrowthBounds/Lemma5aAt.lean}{Formalism.MainNotions.PlaneGramData.aux₃}
\leanref{Formalism/SharpRestriction/GrowthBounds/Lemma5aAt.lean}{Formalism.MainNotions.PlaneGramData.HidaSetting}
\leanref{Formalism/SharpRestriction/GrowthBounds/DualFockLinear/Closed.lean}{Formalism.MainNotions.TestSideCounting.plane_gram_of_admitted}
\leanref{Formalism/SharpRestriction/GrowthBounds/MonomialOrderSegments.lean}{Formalism.MainNotions.PlaneGram.stmt}
\leanref{Formalism/SharpRestriction/GrowthBounds/MonomialOrderSegments/BrData.lean}{Formalism.MainNotions.PlaneGram.Br.plane_gram_of_theorem_at}
\leanref{Formalism/SharpRestriction/GrowthBounds/MonomialOrderSegments/LiftBoundOfHypotheses.lean}{Formalism.MainNotions.PlaneGram.LiftBoundUniform}
\leanref{Formalism/SharpRestriction/GrowthBounds/MonomialOrderSegments/PlaneGramTheoremOfLemmas.lean}{Formalism.MainNotions.PlaneGram.plane_gram_theorem_of_five_lemmas}
\leanref{Formalism/SharpRestriction/GrowthBounds/MonomialOrderSegments/TestBoundOfRhoAndNormSets.lean}{Formalism.MainNotions.PlaneGram.plane_gram_theorem_of_combined_lemmas}
\leanref{Formalism/SharpRestriction/GrowthBounds/MonomialOrderSegments/TestBoundOfRhoAndNormSets.lean}{Formalism.MainNotions.PlaneGram.stepBC_of}

\begin{proof}
\emph{Step 1: the field.} Let $\mathcal K$ be the fraction field of $\mathcal R$, of characteristic $0$, and $v$ its
valuation with $v(p)=1$, with values in $\mathbf Q\cup\{+\infty\}$; the elements of $\mathcal R$ are the $x\in\mathcal K$
with $v(x)\geq0$, so Lemmas~\ref{pd-symbols-test-chains-projection-lift} to~\ref{pd-symbols-test-chains-projection-omega}
apply to $\mathcal K$ and $v$ with this $\mathcal R$. By
Lemma~\ref{pd-symbols-test-chains-projection-ring}, $\mathcal K$ with $v$ has the lifting properties at $p$. As in the
proof of Lemma~\ref{pd-symbols-test-chains-projection-lift}, $v(n)=v_p(n)$ for integers $n\geq1$, and $D_m,D_l\geq0$.
Moreover $\kappa_{\mathcal D}\geq1$: its sum runs over the positions, among which are the $Jr\geq1$ links, and each term
$r(\log_p(d(y)+r)+1)$ is at least $1$.

\emph{Step 2: an integral $c$.} Let $c'=\lceil c\rceil$, and let $\mathcal D'$ be $\mathcal D$ with $c'$ in place of $c$;
it is a signed chain datum, as $c'\geq c\geq1$, and $c'<c+1\leq2c$. The numbers $e_P$, $D_m$, $D_l$, $\kappa_{\mathcal D}$,
$n_0$ and $E(C_E)$, the form $B_{\mathcal D}$, the sets $\tilde C_{j,i}$ and $\tilde S_{j,m}$, admissibility and goodness
(Definition~\ref{pd-symbols-test-chains-projection-good}) do not involve $c$, so they are the same for $\mathcal D'$. The
datum $\mathcal D'$ satisfies (H) for $13$, as
\[
 \frac{rd_L(j)}{P(j,m)}+13rc'\leq\frac{rd_L(j)}{P(j,m)}+26rc\leq s_{j+1,i},
\]
and $4(D_m+D_l+3\kappa_{\mathcal D})<c\leq c'$.

\emph{Step 3: three lemmas.} Lemma~\ref{pd-symbols-test-chains-projection-lift}, applied to $\mathcal D'$ and $F$, gives
$\tilde G\in\mathcal R[x]$ of degree $d(y)$ in the row of every position $y$ with
$\operatorname{val}_{S_{j,m}}(\tilde G)\geq\theta^L_{j,m}-6c'$ for $m+2\leq r$ and
$v(B_{\mathcal D}(F,\tilde G))\leq n_0+D_l+1$. Lemma~\ref{pd-symbols-test-chains-projection-model}, with
$\theta_{j,m}=\theta^L_{j,m}-6c'$, gives a good $Z\in\mathcal R[\tilde x]$ with $\varphi(Z)=\tilde G$ and
$\operatorname{val}_{\tilde S_{j,m}}(Z)\geq\theta^L_{j,m}-6c'$ for $m+2\leq r$.
Lemma~\ref{pd-symbols-test-chains-projection-omega}, applied to $\mathcal D'$ and $F$, gives $\Omega$ and $N_2$; then
$Z'=\Omega(Z)$ is good and satisfies (1) to (3) of that lemma.

\emph{Step 4: the pair.} Put $G'=p^{N_2}Z'$ and $R'=0$. We check the conditions of a projection pair for
$(\mathcal D,F)$ and $(2,24,60)$.
\begin{itemize}
\item The coefficients of $G'$ lie in $\mathcal R$, by (3) of Lemma~\ref{pd-symbols-test-chains-projection-omega}.
\item As $p^{N_2}$ is a nonzero constant, $\operatorname{val}_S(G')=\operatorname{val}_S(Z')$ for every set $S$ of model
  positions. By (1), $\operatorname{val}_{\tilde C_{j,i}}(G')\geq\theta_{j,i}-s_{j,i}-12c'\geq\theta_{j,i}-s_{j,i}-24c$ for
  every chain. By (2), $\operatorname{val}_{\tilde S_{j,m}}(G')\geq\theta^L_{j,m}-6c'-24c'\geq\theta^L_{j,m}-60c$ for
  $m+2\leq r$.
\item The coefficients of $R'=0$ have valuation $+\infty\geq c-E(2)$.
\item Let $\zeta$ be a monomial of $G'$; it is good. As $\varphi$ maps the variables of a letter to those of the letter,
  and those of the label, the lift and the top of a link $(j,m)$ to those of $(j,m)$, the second ones raised to the
  power $P(j,m)$, the monomial $\varphi(\zeta)$ has degree $d(y)$ in the row of each letter $y$ and degree
  $e_{\mathrm{lab}}+P(j,m)e_{\mathrm{lift}}+e_{\mathrm{top}}=d(j,m)$ in the row of each link $(j,m)$.
\item So $\varphi(G'+R')=\varphi(G')$ lies in $\mathcal R[x]$ and has degree $d(y)$ in every row. For an exponent vector
  $\mu$ of these row degrees, $v(\prod_yd(y)!)-v(\mu!)$ is the sum over the positions $y$ of the $p$-adic valuations of
  the multinomial coefficients $d(y)!/\prod_a\mu(y,a)!$. By the proof of
  Lemma~\ref{pd-symbols-flat-real-multinomial}(1), with $r$ parts in place of $k$, each is at most
  $r\lfloor\log_pd(y)\rfloor\leq r(\log_p(d(y)+r)+1)$ if $d(y)\geq1$, and $0$ if $d(y)=0$. So
  $v(\mu!)\geq2n_0-\kappa_{\mathcal D}$, and $v_{\mathrm F}(\varphi(G'+R'))\geq n_0-\kappa_{\mathcal D}/2\geq
  n_0-\kappa_{\mathcal D}$.
\item As $B_{\mathcal D}$ and $\varphi$ are linear, $B_{\mathcal D}(F,\varphi(G'))=p^{N_2}B_{\mathcal D}(F,\varphi(Z'))$,
  and
  \[
   B_{\mathcal D}(F,\varphi(Z'))=B_{\mathcal D}(F,\tilde G)+\big(B_{\mathcal D}(F,\varphi(Z'))-B_{\mathcal D}(F,\varphi(Z))\big).
  \]
  The first term has valuation at most $n_0+D_l+1$. By (3), the second has valuation at least
  $c'+n_0-2\kappa_{\mathcal D}-D_m$, which is larger, as $c'>4(D_m+D_l+3\kappa_{\mathcal D})\geq
  D_m+D_l+2\kappa_{\mathcal D}+1$ by Step~1. So $v(B_{\mathcal D}(F,\varphi(Z')))\leq n_0+D_l+1$, and as
  $N_2\leq D_m+\kappa_{\mathcal D}/2+1$, $E(2)=D_m+D_l+2\kappa_{\mathcal D}$ and $\kappa_{\mathcal D}\geq1$,
  \[
   v\big(B_{\mathcal D}(F,\varphi(G'+R'))\big)\leq N_2+n_0+D_l+1\leq D_m+D_l+\frac{\kappa_{\mathcal D}}2+2+n_0
   \leq E(2)+\kappa_{\mathcal D}+n_0 .
  \]
\end{itemize}
% Lean: Step 1 is PlaneGram.stmt (the AdmAt clauses, Lemma pd-symbols-test-chains-projection-ring) and kapD_ge_one.
%   Step 2 is Br.plane_gram_of_theorem_at: toData Dt (c := ceil c, positions renamed through eO, gO and fM, Fin J for
%   ZMod J), valid_toData, hypH_toData (natCeil_le_two_mul_of_one_le), hc_toData and inXX_toData, then the translation
%   back (exists_pre, phDeg_mem, valS_transport, Dm_eq, Dl_eq, kapD_eq, normShift_eq). The Lean's LiftBoundUniform is
%   LiftBoundWith 13 12 30 2. Step 3 is plane_gram_theorem_of_five_lemmas = plane_gram_theorem_of_combined_lemmas with
%   stepBC_of (stepB_of and stepC_of at theta = thl - 6c) and stepDEF_of. Step 4 is the rest of
%   plane_gram_theorem_of_combined_lemmas: G' = p^N2 * Omega Z, R = 0; ph_mem_PolK, bge_of_cge (fact_lower,
%   vN_multinomial_le, natLog_le_term) with kapL_le_kap', Lam_smul and vle_add_of_lt, vle_pow_p.
% Lean: ProjectionDerivativeEstimate p O 26 2 24 60 is the hypothesis h14b = PlaneGramData.aux3 D 26 2 24 60,
%   which follows from admitted through plane_gram_of_admitted (DualFockLinear/Closed.lean): HidaSetting p O is
%   Nonempty (UniformizedDVR p O) and some AdmissibleDVR p n' O, the hypotheses of the statement. The proof is
%   PlaneGram.stmt D (MonomialOrderSegments.lean) with the six ring properties AdmAt p O of
%   plane_gram_lemmas_of_admitted (Joints.isotypic_multiplicity_of_reducibility_alt,
%   valuation_clauses_of_lattice_index, twisted_valuation_second_sentence, isotypic_coefficient_of_degree_bound,
%   third_degree_bound_gives_span_clause, span_column_clauses, projection_of_integral_corollary,
%   digit_identity_of_decomposition, derivation_of_lattice_index, from the clauses h11, h12, h14, h15, h23 and h24
%   of admitted): Br.plane_gram_of_theorem_at S hG (the translation of the signed chain datum into toData Dt, with
%   c replaced by its ceiling) applied to LiftBoundUniform (KK O) (vQ p S) p 1, which is
%   plane_gram_theorem_of_five_lemmas from second_valuation_of_admitted, brSPAN, brPROJ, brDIG and brDER (stepB_of,
%   stepC_of, segStep_of, stepDEF_of, stepG_of, stepP_of, stepE_of).
% Sizes: MonomialOrderSegments/ is about 10100 lines. What remains open is in Lemma
%   pd-symbols-test-chains-projection-ring: (S) and (P) from the clauses, about 4,500 lines. The clause h15 is Lemma
%   pd-symbols-path-clauses(1); the other five are stated in Lemma pd-symbols-test-chains-projection-clauses.
\end{proof}

\subsection{The projection property below $p_0$}\label{subsection-pd-projection-below-p0}

\begin{lemma}\label{pd-symbols-test-flat-xi}
Let $C_{La},C_k,C_{rk},C_{\deg}$ be real numbers with the property of Theorem~\ref{pd-symbols-count}. There are
$C_{t0},C_D,C_1>0$ such that the following holds for every $C\geq C_1$. Let the data be as in
Lemma~\ref{pd-symbols-test-flat}, and $\kappa_F$ as in Lemma~\ref{pd-symbols-test-flat-real}. Then there are
integers $t_F,D_X\geq0$ with $t_F\leq t_0(C_{t0})+|T|c_0$ and $D_X\leq C_D|T|k^2\log_pa$ such that for every
$F_1\in\mathrm{Op}_{T,c_0}(M)$ of multidegree $d_T$ with $\Theta_\gamma(F_1)\in R^\flat(e)$ and
$\Theta_\gamma(F_1)\notin\varpi_0R^\flat(e)$ there are one-sided elements $\sigma_1,\dots,\sigma_{n_q}$ of
multidegree $d_T$ and normal degrees $o_q\leq t_F$, for $(J,f)$, $N_1$, $G'$, $c'$, $\deg$ and $Y$, such that
$p^{\kappa_F}\langle\Sigma_g(F_1),\sum_q\sigma_q\rangle_{\mathrm F}/\prod_sd_T(s)!$ lies in $\mathcal R$ and divides
$p^{D_X}$.
\end{lemma}
\leanref{Formalism/SharpRestriction/GrowthBounds/DualFockLinear/OneSidedLemma4Realisable.lean}{Formalism.MainNotions.TestSideCounting.XiAt}
\leanref{Formalism/SharpRestriction/GrowthBounds/DualFockLinear/RestXiPOfNS.lean}{Formalism.MainNotions.TestSideCounting.restXi_of_NS}
\leanref{Formalism/SharpRestriction/GrowthBounds/DualFockLinear/XiStepNOfP.lean}{Formalism.MainNotions.TestSideCounting.xiStepN_of_P}
\leanref{Formalism/SharpRestriction/GrowthBounds/DualFockLinear/Lemma4ElementZero.lean}{Formalism.MainNotions.TestSideCounting.xiAt_of}
\leanref{Formalism/SharpRestriction/GrowthBounds/DualFockLinear/RestNDOfAdm.lean}{Formalism.MainNotions.TestSideCounting.restND_of_adm}

\begin{proof}
Let $C_D$ and $C'_1$ be the constants of Lemma~\ref{pd-symbols-test-flat-chains}, $C_{t0}=\max(1,|C_{\deg}|)$, and
$C_1=\max(C'_1,3041+|C_k|,|C_{La}|)$, so that Lemma~\ref{pd-symbols-space} applies. Let $C\geq C_1$ and
the data be as in the statement, and let $D_X$ be the integer of Lemma~\ref{pd-symbols-test-flat-chains}.

\emph{Step 1: normal degrees.} Put $L_a=C_{La}N_s^2\log a$ and $X_b=C_{\deg}|T|\big(L_aa/(pk)+ap^{-L_1}+kq\big)$. By
Lemma~\ref{pd-symbols-space} and Theorem~\ref{pd-symbols-count}(2), every monomial of every $P_b$ has normal degree
at most $X_b$. The polynomial $\Sigma_g(\mathrm{SP}_g)$ is a product of determinants whose entries are powers of
variables $X_{s,b}$ with $s$ even and $b\leq r$ (Step~3 of the proof of Lemma~\ref{pd-symbols-primitive}), which
are not normal; so its monomials have normal degree $0$, and the monomials of the symbol multiplier
$\Sigma_g(\Phi_b)=P_b\Sigma_g(\mathrm{SP}_g)$ have normal degree at most $\lfloor X_b\rfloor$. The substitution
$\Sigma_g$ is of the form $F\mapsto F^N$ of the proof of Lemma~\ref{pd-symbols-test-duality}, so by the chain rule
there, $\Sigma_g$ maps $\mathrm{Op}_{T,c_0}(M)$ into $\mathrm{Op}_{T,c_0}$ of the span of the
$\Sigma_g(\Phi_b)$. An operator $X_{t,a}\,\partial/\partial X_{t,b}$ changes the normal degree of a monomial by at
most $1$, and each $t\in T$ occurs at most $c_0$ times in a product. Hence every monomial of $\Sigma_g(F_1)$, for
$F_1\in\mathrm{Op}_{T,c_0}(M)$, has normal degree at most $t_F=\lfloor X_b\rfloor+|T|c_0$. As $N_s\leq N_1$,
$\log a\geq0$ and $|C_{La}|\leq C$, $X_b\leq t_0(C_{t0})$, so $t_F\leq t_0(C_{t0})+|T|c_0$.

\emph{Step 2: the one-sided elements.} Let $F_1$ be as in the statement and $G=\Sigma_g(F_1)$. By Step 1, $G$ lies
in $\mathrm{Op}_{T,c_0}$ of the span of the $\Sigma_g(\Phi_b)=P_b\Sigma_g(\mathrm{SP}_g)$, and it has multidegree
$d_T$. By property~(2) of the flat model with $i=0$ and $i=1$,
$G\in\mathcal R[X]$ and $G\notin\varpi_0\mathcal R[X]$. Lemma~\ref{pd-symbols-test-flat-chains} gives a signed chain
datum $\mathcal D$ and $F_2$. The ring $\mathcal R$ is a complete discrete valuation ring of characteristic $0$ in
which $p=\varpi_0^2$ is not a unit, and its residue field contains a field of $p^{n'}$ elements: the map from $W$ to
the residue field of $\mathcal R$ kills $p$, so it induces an embedding of the residue field of $W$. So
Lemma~\ref{pd-symbols-test-chains-projection} gives a projection pair for $(\mathcal D,F_2)$, and
Lemma~\ref{pd-symbols-test-flat-chains} then gives one-sided elements $\sigma_q$, of normal degrees $o_q$, with the
divisibility for $G$.

\emph{Step 3: the normal degrees.} The Fock pairing pairs only equal monomials. If $o_q>t_F$, no monomial of
$\sigma_q$ is a monomial of $G$, by Step 1, so $\langle G,\sigma_q\rangle_{\mathrm F}=0$. Replacing these $\sigma_q$ by
$0$, a one-sided element of normal degree $0$ (with $n=0$), changes neither $\langle G,\sum_q\sigma_q\rangle_{\mathrm F}$
nor the other conditions, and afterwards $o_q\leq t_F$ for all $q$.
% Lean: restXi_of_NS = restXi_of_P (restXiP_of_NS: xiAtP_of_NS, Lemmas pd-symbols-test-flat-chains and
%   pd-symbols-test-chains-projection, the latter at O with HidaSetting p O from dvrOfStd hSet and
%   admissible_dvr_of_unramified_quadratic_setting, CharZero from dvr_fraction_field_charZero) = restXi_of_N (xiStepN_of_P: Step 2, with subst_eq_Lin, Lin_OpCl,
%   subst_SPtrue, substPol and the flat model hTheta1) = restXi_of (xiAt_of: Step 3, with lemma4_element_zero and
%   lemma4_element_normal_degree; NDAt from restND_of_adm: Step 1, with symbol_space_at,
%   fixed_symbol_normal_degree_zero, nd_OpCl and level_lower_bound_terms).
\end{proof}

\begin{lemma}\label{pd-symbols-test-flat}
Let $C_{La},C_k,C_{rk},C_{\deg}$ be real numbers with the property of
Theorem~\ref{pd-symbols-count}. There are $C_{t0},C_D,C_7>0$ such that the following holds for every
$C\geq C_7$. Let the data of Situation~\ref{pd-edge-situation} with this $C$, and $(J,f)$, $K$, $N_s$,
$q$, $c'$, $\deg$, $m_*$, $E_0$, $X_j$, $d_T$, $E'$, $d$, $P_1,\dots,P_{n_b}$, $e$, $\gamma$, $g$,
$\Phi_1,\dots,\Phi_{n_b}$, $\lambda^\flat$, $\lambda^\sharp$, $V^\flat$ and $R^\sharp$ satisfy the
hypotheses of Proposition~\ref{pd-symbols-cloud}; let $\Theta_\gamma$, $R^\flat$, $R'^\flat$,
$H^\flat$, $M$ and $L=W^k$ be as in Lemma~\ref{pd-symbols-test-cloud}. Then there are integers
$t_F,D_F\geq0$ with
\[
 t_F\leq t_0(C_{t0})+|T|c_0,\qquad D_F\leq C_D|T|k^2\log_pa,
\]
such that for every $F_1\in\mathrm{Op}_{T,c_0}(M)$ of multidegree $d_T$ with
$\Theta_\gamma(F_1)\in R^\flat(e)$ and $\Theta_\gamma(F_1)\notin\varpi_0R^\flat(e)$ there is
$U\in\varpi_0^{-t_F}H^\flat(L)^\vee\cap R'^\flat(e)$, with $U\in\varpi_0^{c'}R'^\flat(e')$ for every
edge $e'\neq e$ with target $L$, such that $\langle\Theta_\gamma(F_1),U\rangle$ lies in $\mathcal R$
and divides $p^{D_F}$.
\end{lemma}
\leanref{Formalism/SharpRestriction/GrowthBounds/DualFockLinear/TestSideExplicitIdentification.lean}{Formalism.MainNotions.TestSideCounting.RestFlat}
\leanref{Formalism/SharpRestriction/GrowthBounds/DualFockLinear/TestSideExplicitIdentification.lean}{Formalism.MainNotions.TestSideCounting.FlatStep}
\leanref{Formalism/SharpRestriction/GrowthBounds/DualFockLinear/TestSideExplicitIdentification.lean}{Formalism.MainNotions.TestSideCounting.FlatAt}
\leanref{Formalism/SharpRestriction/GrowthBounds/DualFockLinear/ScaleZeroIff.lean}{Formalism.MainNotions.TestSideCounting.restFlat_of_Xi}

\begin{proof}
Let $C_{t0}$, $C_D$ and $C'_1$ be the constants of Lemma~\ref{pd-symbols-test-flat-xi}, $C_8$ that of
Lemma~\ref{pd-symbols-test-one-sided-real}, and $C_7=\max(C'_1,C_8,C_D,1)$. Let $C\geq C_7$ and the data be as in
the statement, and $\kappa_F$ as in Lemma~\ref{pd-symbols-test-flat-real}. Lemma~\ref{pd-symbols-test-flat-xi}
gives $t_F$ and $D_F=D_X$ with the stated bounds, and Lemma~\ref{pd-symbols-test-flat-real} gives $g_v$. As
$C_D\leq C$ and $\log_pa\geq0$, (3) of the situation gives $D_F\leq C|T|k^2\log_pa\leq D_p$, and
$8(D_p+1)\leq c_*\leq c'$; so $2D_F<c'$.

Let $F_1$ be as in the statement. Lemma~\ref{pd-symbols-test-flat-xi} gives one-sided elements $\sigma_q$ of
multidegree $d_T$ and normal degrees $o_q\leq t_F$, and $a=p^{\kappa_F}\langle\Sigma_g(F_1),\sum_q\sigma_q\rangle_{\mathrm F}
/\prod_sd_T(s)!\in\mathcal R$ and $b\in\mathcal R$ with $ab=p^{D_F}$. Each $\sigma_q$ lies in $\mathcal R[X]$, and by
Lemma~\ref{pd-symbols-test-one-sided-real} it is realizable at $(g_v,c',o_q)$. So
Lemma~\ref{pd-symbols-test-flat-real} gives $U_q\in\varpi_0^{-o_q}H^\flat(L)^\vee\cap R'^\flat(e)$, with
$U_q\in\varpi_0^{c'}R'^\flat(e')$ for $e'\neq e$ with target $L$, and $z_q\in\mathcal R$ with
$\langle\Theta_\gamma(F_1),U_q\rangle=p^{\kappa_F}\langle\Sigma_g(F_1),\sigma_q\rangle_{\mathrm F}/\prod_sd_T(s)!+
\varpi_0^{c'}z_q$. As $o_q\leq t_F$, $U_q\in\varpi_0^{-t_F}H^\flat(L)^\vee$. Put $U=\sum_qU_q$; it has the stated
lattice properties, and by linearity $\langle\Theta_\gamma(F_1),U\rangle=a+\varpi_0^{c'}z$ with $z=\sum_qz_q$. As
$c'>2D_F$ and $p=\varpi_0^2$,
\[
 a+\varpi_0^{c'}z=a\big(1+b\varpi_0^{c'-2D_F}z\big),
\]
and $1+b\varpi_0^{c'-2D_F}z$ is a unit of $\mathcal R$. So $\langle\Theta_\gamma(F_1),U\rangle$ lies in $\mathcal R$
and divides $ab=p^{D_F}$.
% Lean: restFlat_of_Xi = restFlat_of hXi (restReal D CLa Ck Crk Cdeg): RestXi gives XiAt (a bound tF for the
%   orders o q, a bound DX, and for every F1 as in the statement a family sig q of OneSidedLemma4Element with
%   ValLe p DX of p^kappa * fock (subst ... F1) (sum sig) / prod (dT s)!), and restReal gives RealAt (forms
%   T in w^(-o) (HL flW e.tgt)^vee and in Re' (flW) e, in w^cP Re' e' for e' /= e, realizing the Fock
%   pairing modulo w^cP). RestXi = restXi_of_NS from RestNS (NSAt: SignedChainData, through
%   ProjectionDerivativeEstimate in xiAtP_of_NS), Lemma pd-symbols-test-flat-xi. In restFlat_of the
%   realisability of the sig q is one_sided_lemma4_realisable (Lemma pd-symbols-test-one-sided-real), under
%   hypII_at; the constants are max (max C2 C3) (max C4 CDf); the last step is valLe_perturb (h2D from h10 and
%   h11 of Hyps), with fock_sum_right for the linearity.
\end{proof}

\begin{lemma}\label{pd-symbols-test}
Let $C_{La},C_k,C_{rk},C_{\deg}$ be real numbers with the property of
Theorem~\ref{pd-symbols-count}. There are $C_{t0},C_{D1},C_5>0$ such that the following holds for
every $C\geq C_5$. Let the data of Situation~\ref{pd-edge-situation} with this $C$, and $(J,f)$, $K$,
$N_s$, $q$, $c'$, $\deg$, $m_*$, $E_0$, $X_j$, $d_T$, $E'$, $d$, $P_1,\dots,P_{n_b}$, $e$, $\gamma$,
$g$, $\Phi_1,\dots,\Phi_{n_b}$, $\lambda^\flat$, $\lambda^\sharp$, $V^\flat$ and $R^\sharp$ satisfy the
hypotheses of Proposition~\ref{pd-symbols-cloud}, and let $\Theta_\gamma$ be the flat model of
$V^\flat$ at $\gamma$ of Lemma~\ref{pd-symbols-flat-exists}, formed with $d_T$. Then there are
integers $t_T,D_1\geq0$ with
\[
 t_T\leq t_0(C_{t0})+|T|c_0+\omega,\qquad D_1\leq C_{D1}|T|k^2\log_pa,
\]
such that for all $y_1,\dots,y_n\in R^\sharp(e)$ and every family $(c_{bi})$ in $\mathcal R$ with
$\sum_{b,i}c_{bi}\Theta_\gamma(\Phi_b)y_i\notin\varpi_0R(e)$ there is
$G\in\varpi_0^{-t_T}H(W^k)^\vee\cap R'(e)$, with $G\in\varpi_0^{c_*}R'(e')$ for every edge $e'\neq e$
with target $W^k$, such that $\langle\sum_{b,i}c_{bi}\Theta_\gamma(\Phi_b)y_i,G\rangle$ lies in
$\mathcal R$ and divides $p^{D_1}$.
\end{lemma}
\leanref{Formalism/SharpRestriction/GrowthBounds/ScaleEdges/SourcePrimitiveIntegrality.lean}{Formalism.MainNotions.Pf14.RemainingInputsAfter}
\leanref{Formalism/SharpRestriction/GrowthBounds/ScaleEdges/SourcePrimitiveIntegrality.lean}{Formalism.MainNotions.Pf14.TestStep}
\leanref{Formalism/SharpRestriction/GrowthBounds/ScaleEdges/SourcePrimitiveIntegrality.lean}{Formalism.MainNotions.Pf14.TestAt}
\leanref{Formalism/SharpRestriction/GrowthBounds/DualFockLinear/TestSideExplicitIdentification.lean}{Formalism.MainNotions.TestSideCounting.test_bound_of_two_halves}

\begin{proof}
Let $C_K,C_6$ be the constants of Lemma~\ref{pd-symbols-test-cloud} and $C_{t0},C_D,C_7$ those of
Lemma~\ref{pd-symbols-test-flat}, and put $C_{D1}=C_K+C_D$ and $C_5=\max(C_6,C_7)$. Let $C\geq C_5$
and the data be as in the statement, and let $\kappa$ and $t_F,D_F$ be the integers given by the two
lemmas. Put $t_T=t_F+\omega$ and $D_1=\kappa+D_F$; they satisfy the stated bounds.

Let $y_1,\dots,y_n\in R^\sharp(e)$ and $(c_{bi})$ be a family in $\mathcal R$ with
$v=\sum_{b,i}c_{bi}\Theta_\gamma(\Phi_b)y_i\notin\varpi_0R(e)$. Lemma~\ref{pd-symbols-test-cloud} gives
$F_1$, $\Lambda$ and $\xi$, and Lemma~\ref{pd-symbols-test-flat}, applied to $F_1$, gives $U$. Put
$G=\Lambda U$. By (4) of Lemma~\ref{pd-symbols-test-cloud} with $i=-t_F$,
$G\in\varpi_0^{-t_F-\omega}H(W^k)^\vee=\varpi_0^{-t_T}H(W^k)^\vee$. By (2), $G\in R'(e)$ and
$\langle v,G\rangle\in\mathcal R$. Let $e'\neq e$ be an edge with target $W^k$. Then
$U\in\varpi_0^{c'}R'^\flat(e')$, so by (3) $G\in\varpi_0^{c'-\omega}R'(e')=\varpi_0^{c_*}R'(e')$, as
$c'=c_*+\omega$. Finally $u=\langle\Theta_\gamma(F_1),U\rangle$ satisfies $ub=p^{D_F}$ for some
$b\in\mathcal R$, and by (1) $\xi\langle v,G\rangle=p^\kappa u$; so
$\langle v,G\rangle\cdot\xi b=p^{\kappa+D_F}=p^{D_1}$.
% Round eleven: Lemmas pd-symbols-test-cloud and pd-symbols-test-flat are proved from the lemmas before
%   them. Round twelve: Lemmas pd-symbols-test-fock-position, pd-symbols-test-flat-real and
%   pd-symbols-test-one-sided-real are proved, the last from Lemma pd-symbols-one-sided-real-division. Round
%   thirteen: the division and Lemma pd-symbols-test-fock-words are proved; what remains open is marked in Lemmas
%   pd-symbols-fock-words-operators, pd-symbols-test-flat-chains and pd-symbols-test-chains-projection.
% Lean: Lemma pd-symbols-test-cloud is RestCloudB, proved as restCloudB_of_op from RestTSOp
%   (restTSOp_of_ten with TenOp; CloudBAtOfRho, RhoAtGOfPair, RestCloudBOfTop/OfTwo/OfThree).
% Lemma pd-symbols-test-flat is RestFlat, proved as restFlat_of_Xi from restXi_of_NS with RestNS
%   (ScaleZeroIff, RestXiPOfNS, XiStepNOfP). Both hold under tableLineDegreeFive.admitted, RestFlat
%   also under SourceCleaning.aux3 and PlaneGramData.aux3 (cleaning_of_admitted, plane_gram_of_admitted
%   in DualFockLinear/Closed.lean). Assembled in TestSideCounting.stmtE_of_OpNS.
% The chain for RestCloudB, read down to its inputs: restCloudB_of_op = restCloudB_of_two hop restSize =
%   restCloudB_of_three hop restSize ontoAll = restCloudB_of_pair (restPairTop_of hop hsz hon), and
%   restCloudB_of_pair = restCloudB_of_top (restRhoTop_of_pair ...) = restCloudB_of_rho (restRho_of_top ...).
%   restPairTop_of is pairAt_of_tsop with three inputs: SizeStep (restSize = sizeStep with CK = 4, C1 = 2:
%   kap1 p (topW lam0) + kap1 p (flW dT) <= CK |T| k^2 log_p a), OntoAll (ontoAll, through tensor_onto and
%   cartanOnto_proved: a scaled element of Re for lam = flW d + rest lies in the span of the products of
%   elements of Re for flW d and for rest) and TSOpStep (TSOpAt: isomorphisms Psi0, Psif from the duals
%   with TSProp, and for each product Theta(F) s an F' in OpCl TT c0 (span F) with the pairing identity).
%   These three, the
%   steps restRhoTop_of_pair, restRho_of_top and restCloudB_of_rho, and the RestFlat chain are not yet in
%   the book.
\end{proof}

\begin{proposition}\label{pd-symbols-cloud}
Let $C_{La},C_k,C_{rk},C_{\deg}$ be real numbers with the property of
Theorem~\ref{pd-symbols-count}. There are $C_{t0},C_{D1},C_3>0$ such that the following holds for
every $C\geq C_3$. Let the data of Situation~\ref{pd-edge-situation} be given with this $C$, let
$(J,f)$ be a run system for $T$ of length $r$ spaced for $(L_1,L_2,N_1,G')$, and form $K$, $N_s$,
$q$, $c'$, $\deg$, $m_*$, $E_0$, the symbol spaces $X_j$ and $d_T$ as in
Proposition~\ref{pd-symbols-realize}. Let $E'\colon\mathbf Z/n'\mathbf Z\to\mathbf Z_{\geq0}$ satisfy
$E'(t)\leq E_0$ and $E'(t)\leq\deg(t)$ for $t\in T$, and let $d$ be formed from it as there. Let
$P_1,\dots,P_{n_b}\in\mathrm{Gen}(X_\bullet,d)$ be linearly independent over $\mathcal R$ and span a
direct summand of $\mathcal R[X]$. Let $e$ be an edge with target $W^k$, and $\gamma\in GL_k(Q)$ and
$g\in GL_k(W)$ with $e=e_\gamma(g)$ (Lemma~\ref{pd-symbols-edge-basis}). Let
$\Phi_1,\dots,\Phi_{n_b}\in\mathcal K[X]$ have multidegree $d_T$, with $\Sigma_g(\Phi_b)$ the symbol
multiplier of $P_b$ for $(g,c',q)$ (Definition~\ref{pd-symbols-twist}). Let
$\lambda^\flat(s)=(d_T(s),0,\dots,0)$ and $\lambda^\sharp=\lambda_0-\lambda^\flat$, and form $V^\flat$
and the lattices $R^\sharp$ of Definition~\ref{pd-projection-property} with them. Then there are a flat
model $\Theta$ of $V^\flat$ at $\gamma$ (Definition~\ref{pd-symbols-flat-model}), integers
$t_T,D_1\geq0$ with
\[
 t_T\leq t_0(C_{t0})+|T|c_0+\omega,\qquad D_1\leq C_{D1}|T|k^2\log_pa,
\]
and vectors $y_1,\dots,y_{n_2}\in R^\sharp(e)$, with $n_2$ the integer of
Lemma~\ref{pd-edge-dimension-count}, such that for every family $(c_{bi})$ in $\mathcal R$ with some
$c_{bi}$ a unit there is $G\in\varpi_0^{-t_T}H(W^k)^\vee\cap R'(e)$, with $G\in\varpi_0^{c_*}R'(e')$
for every edge $e'\neq e$ with target $W^k$, such that
$\langle\sum_{b,i}c_{bi}\Theta(\Phi_b)y_i,G\rangle$ lies in $\mathcal R$ and divides $p^{D_1}$.
\end{proposition}
\leanref{Formalism/SharpRestriction/GrowthBounds/ScaleEdges/RestCloudXOfY.lean}{Formalism.MainNotions.Pf14.RestCloudY}
\leanref{Formalism/SharpRestriction/GrowthBounds/ScaleEdges/RestCloudXOfY.lean}{Formalism.MainNotions.Pf14.CloudStepY}
\leanref{Formalism/SharpRestriction/GrowthBounds/ScaleEdges/RestCloudXOfY.lean}{Formalism.MainNotions.Pf14.CloudAtY}
\leanref{Formalism/SharpRestriction/GrowthBounds/ScaleEdges/SourcePrimitiveIntegrality.lean}{Formalism.MainNotions.Pf14.restCloudY_of}

\begin{proof}
Let $C_4$ be the constant of Lemma~\ref{pd-symbols-primitive}, let $C_{t0},C_{D1},C_5$ be those of
Lemma~\ref{pd-symbols-test}, and put $C_3=\max(C_4,C_5)$. Let $C\geq C_3$ and the data be as in the
statement. Since $K\geq1$, $1\leq r<k$, so Lemma~\ref{pd-symbols-flat-exists}, applied with $d_T$ in
place of $d$, gives the flat model $\Theta=\Theta_\gamma$ of $V^\flat$ at $\gamma$; here $\lambda^\flat$
is the weight $(d_T(s),0,\dots,0)$ of Definition~\ref{pd-symbols-flat-model}. As $C\geq C_4$,
Lemma~\ref{pd-symbols-primitive} gives $y_1,\dots,y_{n_2}\in R^\sharp(e)$ such that
$\sum_{b,i}c_{bi}\Theta(\Phi_b)y_i\notin\varpi_0R(e)$ whenever some $c_{bi}$ is a unit. As $C\geq C_5$,
Lemma~\ref{pd-symbols-test} gives integers $t_T$ and $D_1$ with the stated bounds. Let $(c_{bi})$ be a
family in $\mathcal R$ with some $c_{bi}$ a unit. Then
$\sum_{b,i}c_{bi}\Theta(\Phi_b)y_i\notin\varpi_0R(e)$, and Lemma~\ref{pd-symbols-test}, applied to
$y_1,\dots,y_{n_2}$ and $(c_{bi})$, gives the form $G$.
\end{proof}

\begin{lemma}\label{pd-symbols-flat-source}
Let $p\geq5$ be a prime, $(W,Q,\sigma,\iota)$ a Hida setting of degree $n'$ at $p$ with rings
$\mathcal R\subset\mathcal K$ and $\varpi_0$, $1\leq r\leq k$ integers,
$d\colon\mathbf Z/n'\mathbf Z\to\mathbf Z_{\geq0}$, $\lambda^\flat(s)=(d(s),0,\dots,0)$, and $V^\flat$,
$H^\flat$, $R^\flat$ and $\mathcal K[X]_d$ as in Definition~\ref{pd-symbols-flat-model}. Let
$\gamma\in GL_k(Q)$ with lattice $L''$ (Lemma~\ref{pd-symbols-edge-basis}), $\Theta$ a flat model of
$V^\flat$ at $\gamma$, $g\in GL_k(W)$, $e=e_\gamma(g)$, and $c',o\geq0$ integers. Let
$\Phi\in\mathcal K[X]_d$ be such that $\Sigma_g(\Phi)$ has multidegree $d$ and is realizable at
$(g,c',o)$ (Definition~\ref{pd-symbols-twist}). Then there is $F^v\in\varpi_0^{-o}H^\flat(L'')$ with
$F^v-\Theta(\Phi)\in\varpi_0^{c'}R^\flat(e)$ and $F^v\in\varpi_0^{c'}R^\flat(e')$ for every edge
$e'\neq e$ with source $L''$.
\end{lemma}
\leanref{Formalism/SharpRestriction/GrowthBounds/ScaleEdges/FlatSrc.lean}{Formalism.MainNotions.Pf14.flatSrc}

\begin{proof}
Let $F\in\mathcal R[X]$ realize $\Sigma_g(\Phi)$ at $(g,c',o)$. Since $\Sigma_g(\Phi)$ has
multidegree $d$, so does $F$, and $F\in\mathcal K[X]_d$. Put $F^v=\varpi_0^{-o}\Theta(F)$. By
property (1) of the flat model, $\Theta(F)\in H^\flat(L'')$, so $F^v\in\varpi_0^{-o}H^\flat(L'')$.
The polynomial $F-\varpi_0^o\Phi$ lies in $\mathcal K[X]_d$, and
$\Sigma_g(F-\varpi_0^o\Phi)=\Sigma_g(F)-\varpi_0^o\Sigma_g(\Phi)\in\varpi_0^{o+c'}\mathcal R[X]$. By
property (2), $\Theta(F)-\varpi_0^o\Theta(\Phi)\in\varpi_0^{o+c'}R^\flat(e)$; multiplying by
$\varpi_0^{-o}$ gives $F^v-\Theta(\Phi)\in\varpi_0^{c'}R^\flat(e)$. Let $e'\neq e$ be an edge with
source $L''$. By Lemma~\ref{pd-symbols-edge-basis}(1), $e'=e_\gamma(g')$ for some $g'\in GL_k(W)$,
and by Lemma~\ref{pd-symbols-edge-basis}(3) the planes of $g'$ and $g$ differ. So
$\Sigma_{g'}(F)\in\varpi_0^{o+c'}\mathcal R[X]$, hence $\Theta(F)\in\varpi_0^{o+c'}R^\flat(e')$ by
property (2), and $F^v\in\varpi_0^{c'}R^\flat(e')$.
\end{proof}

\begin{lemma}\label{pd-symbols-rest-bound}
Let $p\geq5$ be a prime, $(W,Q,\sigma,\iota)$ a Hida setting of degree $n'$ at $p$ with rings
$\mathcal R\subset\mathcal K$ and $\varpi_0$, $k\geq3$ and $r\geq0$ integers,
$T\subset\mathbf Z/n'\mathbf Z$, $a$ a real number, $c_{pp},c_0\geq0$ integers and
$a_T,b_T\colon\mathbf Z/n'\mathbf Z\to\mathbf Z$. Let $\lambda_0(s)\in\mathbf Z^k$ be dominant weights
such that $\lambda_0(s)_1=a_T(s)$, $\lambda_0(s)_k=-b_T(s)$ and $|\lambda_0(s)_i|\leq c_{pp}$ for
$1<i<k$ when $s\in T$, $|\lambda_0(s)_i|\leq c_{pp}$ for all $i$ when $s\notin T$, and
$a+c_{pp}\leq a_T(s)\leq2a-c_{pp}$, $b_T(s)\geq0$ and $b_T(s)+2c_{pp}\leq c_0$ for $s\in T$. Let
$\lambda^\sharp(s)=(\lambda_0(s)_2,\lambda_0(s)_2,\lambda_0(s)_3,\dots,\lambda_0(s)_k)$ for $s\in T$
and $\lambda^\sharp(s)=\lambda_0(s)$ for $s\notin T$, and let $H^\sharp$ and $R^\sharp$ be the
lattices of Definition~\ref{pd-projection-property} formed with $\lambda^\sharp$. Then
$R^\sharp(e)\subset\varpi_0^{-\omega}H^\sharp(L'')$ for every edge $e$ with source $L''$, where
$\omega=r|T|c_0+2rn'c_{pp}$.
\end{lemma}
\leanref{Formalism/SharpRestriction/GrowthBounds/ScaleEdges/EdgeLatticeAdaptedBasis.lean}{Formalism.MainNotions.Pf14.Re_le_scale_HL}
\leanref{Formalism/SharpRestriction/GrowthBounds/ScaleEdges/OrdMaxRest.lean}{Formalism.MainNotions.Pf14.ordMax_rest}

\begin{proof}
The weights $\lambda^\sharp(s)$ are dominant. Let $\lambda(s)$ be formed from $\lambda^\sharp(s)$ as
in Definition~\ref{pd-local-model}, and take the basis $(x_i)$ and the integers $\pi_i$ of
Lemma~\ref{pd-edge-basis}(1) for $\lambda^\sharp$ and $e$. In its proof,
$\pi_i=\sum_s(\mu_1+\cdots+\mu_r)$ with $\mu=\beta+\lambda(s)_k(1,\dots,1)$, where $\beta$ is the
weight of a polynomial under the diagonal matrices, so $\beta\geq0$ and $\mu_l\geq\lambda(s)_k$. Hence
\[
 A-\pi_i\leq\sum_s\sum_{l\leq r}\big(\lambda(s)_l-\lambda(s)_k\big)
 \leq r\sum_s\big(\lambda^\sharp(s)_1-\lambda^\sharp(s)_k\big),
\]
as $\lambda(s)_1-\lambda(s)_k=\lambda^\sharp(s)_1-\lambda^\sharp(s)_k$ for odd $s$ too. For $s\in T$,
$\lambda^\sharp(s)_1-\lambda^\sharp(s)_k=\lambda_0(s)_2+b_T(s)\leq c_{pp}+b_T(s)\leq c_0$, and for
$s\notin T$ it is at most $2c_{pp}$. So $A-\pi_i\leq\omega$, and each basis vector
$\varpi_0^{\pi_i-A}x_i=\varpi_0^{\omega-(A-\pi_i)}\cdot\varpi_0^{-\omega}x_i$ of $R^\sharp(e)$ lies in
$\varpi_0^{-\omega}H^\sharp(L'')$.
\end{proof}

\begin{proposition}\label{pd-edge-step}
There are $C_1,C_t,C_D,C_r>0$ with the following property. In Situation~\ref{pd-edge-situation} with
$C\geq C_1$, let $e$ be an edge with target $W^k$. Then there are integers $n,t_S,t_T,D_1\geq0$ with
\[
 t_S\leq t_0(C_t)+\omega,\qquad t_T\leq t_0(C_t)+|T|c_0+\omega,\qquad D_1\leq C_D|T|k^2\log_pa,
\]
\[
 n\geq\Big(1-\frac{C_r|T|r^2k}{Y}\Big)\Big(1-\frac{kc_0}a\Big)^{|T|}\frac{\dim V}{p^{n'r(k-r)}},
\]
and an edge family at $e$ (Definition~\ref{pd-edge-family}) of size $n$, with levels $(t_S,t_T,D_1)$
and margin $c_*$.
\end{proposition}
\leanref{Formalism/SharpRestriction/GrowthBounds/ScaleEdges/ManyEdges.lean}{Formalism.MainNotions.Pf14.EdgeStep}
\leanref{Formalism/SharpRestriction/GrowthBounds/ScaleEdges/ManyEdges.lean}{Formalism.MainNotions.Pf14.RestEdge}
\leanref{Formalism/SharpRestriction/GrowthBounds/ScaleEdges/InitialLevelMonoWeights.lean}{Formalism.MainNotions.Pf14.edgeStep_of}
\leanref{Formalism/SharpRestriction/GrowthBounds/ScaleEdges/InitialLevelMonoWeights.lean}{Formalism.MainNotions.Pf14.restEdge_of}

\begin{proof}
Let $C_{La},C_k,C_{rk},C_{\deg}$ be real numbers with the property of
Theorem~\ref{pd-symbols-count}, $C_2$ the constant of Proposition~\ref{pd-symbols-realize} and
$C_{t0},C_{D1},C_3$ those of Proposition~\ref{pd-symbols-cloud}. Put $C_t=\max(C_{t0},|C_{\deg}|)$,
$C_D=C_{D1}$, $C_r=\max(C_{rk},1)$ and
$C_1=\max(3041+|C_k|,C_2,C_3,2,|C_{La}|,C_D)$. Since $C\geq2$ and $C\geq C_D$,
Lemma~\ref{pd-edge-numerics} applies with $1$, $C_D$ and $1$ in place of $C_t$, $C_D$ and $C_r$; it
gives $C_D|T|k^2\log_pa\leq D_p$ and $kc_0\leq a$. Also $k\geq3r+1\geq7$ by (3). Put $c'=c_*+\omega$. For $s\in T$ let
$d_s=\lambda_0(s)_1-\lambda_0(s)_2$ and $\lambda^\flat(s)=(d_s,0,\dots,0)$, and for $s\notin T$ let
$\lambda^\flat(s)=0$. Put $\lambda^\sharp=\lambda_0-\lambda^\flat$. Both are dominant, since
$\lambda^\sharp(s)=(\lambda_0(s)_2,\lambda_0(s)_2,\lambda_0(s)_3,\dots,\lambda_0(s)_k)$ for $s\in T$.
Write $V^\flat$, $H^\flat$, $R^\flat$ and $V^\sharp$, $H^\sharp$, $R^\sharp$ for the objects of
Definition~\ref{pd-projection-property} formed with $\lambda^\flat$ and $\lambda^\sharp$, and let
$L''$ be the source of $e$ and $L=W^k$ its target. We first construct the following.
\begin{enumerate}
\item[(a)] An integer $t_1\leq t_0(C_t)$ and vectors $F_1,\dots,F_{n_1}\in\varpi_0^{-t_1}H^\flat(L'')$
  and $S_1,\dots,S_{n_1}\in R^\flat(e)$ such that $F_b-S_b\in\varpi_0^{c'}R^\flat(e)$ and
  $F_b\in\varpi_0^{c'}R^\flat(e')$ for every edge $e'\neq e$ with source $L''$, where
  \[
   n_1\geq p^{-n'r(k-r)}\Big(1-\frac{C_r|T|r^2k}Y\Big)\prod_{t\in T}\binom{d_t+k-1}{k-1}.
  \]
\item[(b)] An integer $o_S\leq\omega$ and vectors
  $y_1,\dots,y_{n_2}\in R^\sharp(e)\cap\varpi_0^{-o_S}H^\sharp(L'')$, with $n_2$ the integer of
  Lemma~\ref{pd-edge-dimension-count}.
\item[(c)] Integers $t_T\leq t_0(C_t)+|T|c_0+\omega$ and $D_1\leq C_D|T|k^2\log_pa$ such that, for every
  family $(c_{bi})$ in $\mathcal R$ with some $c_{bi}$ a unit, there is
  $G\in\varpi_0^{-t_T}H(L)^\vee\cap R'(e)$, with $G\in\varpi_0^{c_*}R'(e')$ for every edge $e'\neq e$
  with target $L$, such that $\langle\sum_{b,i}c_{bi}S_by_i,G\rangle$ lies in $\mathcal R$ and
  divides $p^{D_1}$.
\end{enumerate}
Since $256k^2|T|Y^2\leq p^{L_1-1}$, we have $L_1\geq1$. By (1) of the situation and as $n'$ is even,
Lemma~\ref{pd-symbols-run-system-exists} gives a run system $(J,f)$ for $T$ of length $r$ spaced for
$(L_1,L_2,N_1,G')$. Form $K=k-r$, $N_s$, $q$, $\deg$, $m_*$ and $E_0$ as in
Lemma~\ref{pd-symbols-space}, the symbol spaces $X_j$ for $(C_{La},\deg,L_1,N_s,m_*,a,Y)$ over
$\mathcal R$, and $d_T$ as in Proposition~\ref{pd-symbols-realize}. Then $\deg(t)=d_t$ for $t\in T$,
and $\lambda^\flat(s)=(d_T(s),0,\dots,0)$ for every $s$. As $C\geq C_2$,
Proposition~\ref{pd-symbols-realize} gives $E'$ with $E'(t)\leq E_0$ and $E'(t)\leq d_t$ for $t\in T$;
let $d$ be formed from it as there.

As $C\geq3041+|C_k|$, Lemma~\ref{pd-symbols-space} and Theorem~\ref{pd-symbols-count}(1) give a
submodule $S\subset N(X_\bullet,d)$, a direct summand of $\mathcal R[X]$, with a basis of symbols
$P_1,\dots,P_{n_1}$, where
\[
 n_1\geq p^{-n'rK}\Big(1-\frac{C_{rk}|T|r^2k}Y\Big)\prod_{t\in T}\binom{d_t+k-1}{k-1}
 \geq p^{-n'r(k-r)}\Big(1-\frac{C_r|T|r^2k}Y\Big)\prod_{t\in T}\binom{d_t+k-1}{k-1},
\]
because $C_r\geq C_{rk}$ and the product is nonnegative. By Theorem~\ref{pd-symbols-count}(2), every
monomial of every $P_b$ has normal degree at most
\[
 \mathcal D=C_{\deg}|T|\Big(\frac{C_{La}N_s^2\log a\cdot a}{pk}+ap^{-L_1}+kp^{G'}\Big).
\]
Let $t_1=\lfloor\mathcal D\rfloor$ if $\mathcal D\geq0$ and $t_1=0$ otherwise. Normal degrees are
integers, so they are at most $t_1$. Since $|C_{\deg}|\leq C_t$, $|C_{La}|\leq C$, $N_s\leq N_1$ and
$\log a\geq0$, the absolute value of the bracket in $\mathcal D$ is at most the bracket in $t_0(C_t)$,
so $\mathcal D\leq t_0(C_t)$; and $t_0(C_t)\geq0$. Hence $t_1\leq t_0(C_t)$.

Let $\gamma\in GL_k(Q)$ have columns forming a $W$-basis of $L''$. By
Lemma~\ref{pd-symbols-edge-basis}(1) there is $g\in GL_k(W)$ with $e=e_\gamma(g)$. For each $b$,
Proposition~\ref{pd-symbols-realize}, applied to $g$, $P_b$ and $t_1$, gives polynomials
$\Phi_{b,o}$ of multidegree $d_T$, $o\geq0$, such that $\Sigma_g(\Phi_{b,o})$ has multidegree $d_T$,
lies in $\mathcal R[X]$ and is realizable at $(g,c',o)$, and $\sum_{o\leq t_1}\Sigma_g(\Phi_{b,o})$ is
the symbol multiplier of $P_b$ for $(g,c',q)$. Put $\Phi_b=\sum_{o=0}^{t_1}\Phi_{b,o}$; it has
multidegree $d_T$, and $\Sigma_g(\Phi_b)$ is the symbol multiplier of $P_b$. The $P_b$ are linearly
independent and span the direct summand $S$. As $C\geq C_3$, Proposition~\ref{pd-symbols-cloud},
applied to the $P_b$, $e$, $\gamma$, $g$ and the $\Phi_b$, gives a flat model $\Theta$ of $V^\flat$
at $\gamma$, integers $t_T$ and $D_1$, and vectors $y_1,\dots,y_{n_2}\in R^\sharp(e)$.

For (a), put $S_b=\Theta(\Phi_b)$. For each $b$ and $o\leq t_1$, Lemma~\ref{pd-symbols-flat-source}
gives $F^v_{b,o}\in\varpi_0^{-o}H^\flat(L'')$ with $F^v_{b,o}-\Theta(\Phi_{b,o})\in\varpi_0^{c'}R^\flat(e)$
and $F^v_{b,o}\in\varpi_0^{c'}R^\flat(e')$ for every edge $e'\neq e$ with source $L''$. Put
$F_b=\sum_{o=0}^{t_1}F^v_{b,o}$. Then $F_b\in\varpi_0^{-t_1}H^\flat(L'')$,
$F_b-S_b=\sum_o\big(F^v_{b,o}-\Theta(\Phi_{b,o})\big)\in\varpi_0^{c'}R^\flat(e)$ and
$F_b\in\varpi_0^{c'}R^\flat(e')$. Since $\Sigma_g(\Phi_{b,o})\in\mathcal R[X]$, property (2) of the
flat model with $i=0$ gives $\Theta(\Phi_{b,o})\in R^\flat(e)$, so $S_b\in R^\flat(e)$. With the
bounds on $n_1$ and $t_1$ this is (a).

For (b), take $o_S=\omega$: by Lemma~\ref{pd-symbols-rest-bound}, $y_i\in R^\sharp(e)\subset
\varpi_0^{-\omega}H^\sharp(L'')$. For (c), Proposition~\ref{pd-symbols-cloud} gives the forms $G$,
with $D_1\leq C_D|T|k^2\log_pa$ and $t_T\leq t_0(C_{t0})+|T|c_0+\omega\leq t_0(C_t)+|T|c_0+\omega$,
since $t_0(C')$ increases with $C'$.
% Lean (edgeStep_of): basis_fin, edgePlane, flatSrc, level_lower_bound_terms, initial_level_mono_weights,
% restCloudX_of_Y (Re_le_scale_HL, ordMax_rest), restCloud_of_X (count_of_exact), norm_index_count_bound.
% The four statements used above carry their proofs, and what remains of them is marked there.
Given (a), (b) and (c), the $n_1n_2$ products $F_by_i$ form an edge family at $e$ with levels
$(t_1+o_S,t_T,D_1)$ and margin $c_*$. By Lemma~\ref{pd-edge-products},
$F_by_i\in\varpi_0^{-t_1-o_S}H(L'')$, $S_by_i\in R(e)$ and
$F_by_i-S_by_i=(F_b-S_b)y_i\in\varpi_0^{c'}R(e)$. Hence $F_by_i\in R(e)$, which is
condition~(1) of Definition~\ref{pd-edge-family}. Let $e'\neq e$ be an edge with source $L''$. By
Lemma~\ref{pd-edge-basis}(2) for $\lambda^\sharp$ and $e'$,
$y_i\in\varpi_0^{-o_S}H^\sharp(L'')\subset\varpi_0^{-o_S}R^\sharp(e')$, so
Lemma~\ref{pd-edge-products} gives $F_by_i\in\varpi_0^{c'-o_S}R(e')\subset\varpi_0^{c_*}R(e')$, as
$c'-o_S\geq c_*$; this is condition~(2). For condition~(3), let $(c_{bi})$ have a unit entry and take
$G$ as in (c). Then
\[
 \Big\langle\sum_{b,i}c_{bi}F_by_i,G\Big\rangle=\Big\langle\sum_{b,i}c_{bi}S_by_i,G\Big\rangle
 +\Big\langle\sum_{b,i}c_{bi}(F_b-S_b)y_i,G\Big\rangle.
\]
The first term divides $p^{D_1}$. In the second, the vector lies in $\varpi_0^{c'}R(e)$ and
$G\in R'(e)$, so by Lemma~\ref{pd-edge-basis}(3) the term lies in $\varpi_0^{c'}\mathcal R\subset
p^{D_1+1}\mathcal R$; here $c'\geq c_*\geq8(D_p+1)\geq2(D_1+1)$ by Lemma~\ref{pd-edge-numerics}(2).
So the sum has the valuation of the first term and divides $p^{D_1}$.

Finally $t_1+o_S\leq t_0(C_t)+\omega$. If $1-C_r|T|r^2k/Y\geq0$, the bound on $n_1$ in (a) and
Lemma~\ref{pd-edge-dimension-count} give
\[
 n_1n_2\geq\Big(1-\frac{C_r|T|r^2k}{Y}\Big)\Big(1-\frac{kc_0}a\Big)^{|T|}\frac{\dim V}{p^{n'r(k-r)}};
\]
otherwise the right side is negative. Take $n=n_1n_2$, $t_S=t_1+o_S$ and the $t_T$, $D_1$ of (c).
\end{proof}

\begin{theorem}\label{pd-gain-theorem-two}
There is $C_0\geq0$ with the following property. Let $p\geq13$ be a prime, $(W,Q,\sigma,\iota)$ an
unramified quadratic Hida setting of degree $n'$ at $p$ (Definition~\ref{pd-projection-property}),
$m\geq0$ and $r\geq2$ integers, $T\subset\mathbf Z/n'\mathbf Z$, $L_1,L_2\geq0$ integers and
$\delta,h,x,a,Y>0$ reals. Let $u$, $M_0$, $G'$ and $N_1$ be the integers of
Definition~\ref{hf-good-orbit} formed with $r$, $K=r-1$, $\delta$ and $h$ in place of $r_\varepsilon$,
$K_\varepsilon$, $\delta_1$ and $h_\varepsilon$; thus $u=\lceil1/\delta\rceil$ and
$M_0=\lceil(h+3)10^{6r+1}\log x\rceil$. Let
$\lambda_0(s)\in\mathbf Z^{m+1}$ be dominant weights and $a_T,b_T\colon\mathbf Z/n'\to\mathbf Z$,
$c_{pp},c_0,D_p,c_*\geq0$ integers. Suppose that:
\begin{enumerate}
\item $T$ satisfies the three conditions of a tiled set of positions for $(x,\cdot,L_2)$ with these
  integers, $2L_1\leq L_2<u$, $3G'\leq n'$ and $N_1^2\leq r|T|$;
\item $\lambda_0(s)_1=a_T(s)$, $\lambda_0(s)_{m+1}=-b_T(s)$ and $|\lambda_0(s)_i|\leq c_{pp}$ for
  $1<i<m+1$ when $s\in T$, $|\lambda_0(s)_i|\leq c_{pp}$ for all $i$ when $s\notin T$, and
  $a+c_{pp}\leq a_T(s)\leq2a-c_{pp}$, $b_T(s)\geq0$ and $b_T(s)+2c_{pp}\leq c_0$ for $s\in T$;
\item $r\leq Y$, $C_0N_1^2\log a\leq m+1$, $N_1+2r\leq m+1-r$,
  $256(m+1)^2|T|Y^2\leq p^{L_1-1}$, $C_0|T|(m+1)^2\log_pa\leq D_p$, $8(D_p+1)\leq c_*$ and
  \[
   C_0(m+1)^3r^2|T|\,p^{5G'}\big(c_*+r|T|c_0+2rn'c_{pp}\big)Y\leq a .
  \]
\end{enumerate}
Then for every integer
\[
 t\geq C_0r|T|\Big(\frac{C_0N_1^2\log a\cdot a}{p(m+1)}+ap^{-L_1}+c_0+(m+1)p^{G'}\Big)+2rn'c_{pp}
\]
the projection property $\mathrm{LP}\big(r,\lambda_0;t,D_p,C_0|T|r^2(m+1)/Y,c_*\big)$ holds.
\end{theorem}
\leanref{Formalism/SharpRestriction/GrowthBounds/TableLineDegreeFive.lean}{Formalism.MainNotions.tableLineDegreeFive.aux₃}
\leanref{Formalism/SharpRestriction/GrowthBounds/TableLineDegreeFive.lean}{Formalism.MainNotions.tableLineDegreeFive.OrbitHyp}
\leanref{Formalism/SharpRestriction/GrowthBounds/TableLineDegreeFive.lean}{Formalism.MainNotions.tableLineDegreeFive.WeightHyp}
\leanref{Formalism/SharpRestriction/GrowthBounds/TableLineDegreeFive.lean}{Formalism.MainNotions.tableLineDegreeFive.tBound}
\leanref{Formalism/SharpRestriction/GrowthBounds/PlaneCountTheoremWithConstants.lean}{Formalism.MainNotions.G.second_theorem_absolute_constant}
\leanref{Formalism/SharpRestriction/GrowthBounds/PlaneCountTheoremWithConstants.lean}{Formalism.MainNotions.G.SecondTheoremReduced}

\begin{proof}
Let $C_1,C_t,C_D,C_r$ be the constants of Proposition~\ref{pd-edge-step}, and put
$C_0=\max(C_1,2,C_t,C_D,C_r+1)$. Let the data of the theorem be given, and put $k=m+1$; they form
Situation~\ref{pd-edge-situation} with $C=C_0$. Let $t$ be an integer at least the bound of the
theorem, and put $\kappa=C_0|T|r^2k/Y$ and
$\rho=(1-C_r|T|r^2k/Y)(1-kc_0/a)^{|T|}$.

By (3), $4(D_p+1)\leq c_*$. So, by Lemma~\ref{pd-edge-transport}, it suffices to show the condition
of Definition~\ref{pd-projection-property}(3) at the lattice $W^k$. By Lemma~\ref{pd-edge-count}
there is a set $E$ of at least $p^{n'r(k-r)}$ edges with target $W^k$. By
Proposition~\ref{pd-edge-step}, each $e\in E$ carries an edge family of size
$n_e\geq\rho\dim V/p^{n'r(k-r)}$, with levels $(t_S(e),t_T(e),D_1(e))$ and margin $c_*$, where
$t_S(e)$ and $t_T(e)$ are at most $t_0(C_t)+|T|c_0+\omega$ and $D_1(e)\leq C_D|T|k^2\log_pa$. By
Lemma~\ref{pd-edge-numerics}(1), $t_0(C_t)+|T|c_0+\omega$ is at most the bound of the theorem, hence
at most $t$; by Lemma~\ref{pd-edge-numerics}(2), $D_1(e)\leq D_p$. If $\rho\geq0$, then
\[
 \sum_{e\in E}n_e\geq|E|\,\frac{\rho\dim V}{p^{n'r(k-r)}}\geq\rho\dim V;
\]
if $\rho<0$, then $\sum_en_e\geq0>\rho\dim V$. By Lemma~\ref{pd-edge-numerics}(3),
$1-\kappa\leq\rho$, so $\sum_{e\in E}n_e\geq(1-\kappa)\dim V$. Since $2(D_p+1)\leq c_*$,
Lemma~\ref{pd-edge-assembly} gives the condition at $W^k$.
\end{proof}
% Lean: second_theorem_absolute_constant takes C0 from second_theorem_closed_everywhere, i.e. from
% the closed form TraceSpectra.stmtE (tableLineDegreeFive.ClosedFormOfAdmitted, through
% PairMembershipBlocks.stmtE_of_chain, Proposition5Numerics.stmtE_of_mem, TestSideWords.stmtE_of_MP,
% TestSideCounting.stmtE_of_OpNS and theorem_II_closed_of_test_bound) applied to the 37 clauses of
% tableLineDegreeFive.admitted supplied by ExtremalProjectorTools and Loc6. The argument above is
% Pf14.theorem_II_of_split_cloud_bounds, which reduces through theorem_II_of_cloud_bound_Y and
% theorem_II_of_cloud_bound_X to theorem_II_of_edge_bound and
% prvE_of_gl_operators_form_representation with manyEdges, s7, asm and num.

\subsection{Split families and the edge rank condition}\label{subsection-pd-edge-rank-condition}

\begin{definition}\label{pd-kappa-split-rank}
Let $p\geq5$ be a prime and $(W,Q,\sigma,\iota)$ an unramified quadratic Hida setting of degree $n'$
at $p$ with rings $\mathcal R\subset\mathcal K$ and $\varpi_0$
(Definition~\ref{pd-projection-property}), and let $v$ be the valuation of $\mathcal K$ with
$v(p)=1$. Let $k\geq1$, let $\lambda_0(s)\in\mathbf Z^k$, $s\in\mathbf Z/n'\mathbf Z$, be dominant
weights, and let $V$, $V^\vee$, the lattices $H(L)$, the edges, $R(e)$, $R'(e)$ and the edge pairing
$\langle\ ,\ \rangle$ be those of Definition~\ref{pd-projection-property} with $r=1$.
\begin{enumerate}
\item For integers $t,c\geq0$ and an edge $e$ with source $L''$ and target $L$, the \emph{cores} at
  $e$ are
  \[
   \mathcal C_{t,c}(e)=\varpi_0^{-t}H(L'')\cap R(e)\cap\bigcap_{e'}\varpi_0^{c}R(e'),\qquad
   \mathcal C'_{t,c}(e)=\varpi_0^{-t}H(L)^\vee\cap R'(e)\cap\bigcap_{e''}\varpi_0^{c}R'(e''),
  \]
  where $e'$ runs over the edges other than $e$ with source $L''$ and $e''$ over the edges other
  than $e$ with target $L$.
\item Let $\mathcal C\subset V$ and $\mathcal C'\subset V^\vee$ be $\mathcal R$-submodules and
  $D\in\mathbf R$. A \emph{split family} of depth $D$ and size $n\geq1$ on $(\mathcal C,\mathcal C')$
  consists of $x_1,\dots,x_n\in\mathcal C$, $y_1,\dots,y_n\in\mathcal C'$, a nonzero
  $z\in\mathcal R$ with $v(z)\leq D$ and an $n\times n$ matrix $(M_{lj})$ over $\mathcal R$ such that
  $\sum_l\langle x_i,y_l\rangle M_{lj}=z\delta_{ij}$ for all $i,j$. The empty family is a split
  family of size $0$ and of every depth.
\item For integers $t,c\geq0$ and reals $D$ and $X$, the \emph{edge rank condition}
  $\mathrm{ER}(\lambda_0;t,c,D,X)$ holds if for every edge $e$ there is a split family of depth $D$
  and of some size $n\geq X$ on $(\mathcal C_{t,c}(e),\mathcal C'_{t,c}(e))$.
\item Put $[k]=(p^{n'k}-1)/(p^{n'}-1)$ and $\mu(\lambda_0)=\dim V/[k]$.
\item Let $k\geq2$ and $T\subset\mathbf Z/n'\mathbf Z$. For $\alpha,\beta\colon T\to\mathbf Z_{\geq0}$
  the \emph{flat weights} $\phi(\alpha,\beta)$ are $\phi(\alpha,\beta)(s)=(\alpha(s),0,\dots,0,-\beta(s))$
  for $s\in T$ and $\phi(\alpha,\beta)(s)=0$ for $s\notin T$. The \emph{top part} $\lambda_0^T$ of
  $\lambda_0$ is $\lambda_0^T(s)=\lambda_0(s)$ for $s\in T$ and $\lambda_0^T(s)=0$ for $s\notin T$.
\end{enumerate}
\end{definition}
\leanref{Formalism/SharpRestriction/GrowthBounds/ColumnSpanDensity/RhoLowerBoundLargeConductor.lean}{Formalism.MainNotions.flipOneThree.Pf.CCL}
\leanref{Formalism/SharpRestriction/GrowthBounds/ColumnSpanDensity/RhoLowerBoundLargeConductor.lean}{Formalism.MainNotions.flipOneThree.Pf.CCL'}
\leanref{Formalism/SharpRestriction/GrowthBounds/ObjSeg.lean}{Formalism.MainNotions.G.vLEO}
\leanref{Formalism/SharpRestriction/GrowthBounds/ObjSeg.lean}{Formalism.MainNotions.G.FactorsLE}
\leanref{Formalism/SharpRestriction/GrowthBounds/ObjSeg.lean}{Formalism.MainNotions.G.RhoGE}
\leanref{Formalism/SharpRestriction/GrowthBounds/ColumnSpanDensity/RhoLowerBoundLargeConductor.lean}{Formalism.MainNotions.flipOneThree.Pf.RhoL}
\leanref{Formalism/SharpRestriction/GrowthBounds/ObjSeg.lean}{Formalism.MainNotions.G.gaussianNumber}
\leanref{Formalism/SharpRestriction/GrowthBounds/ColumnSpanDensity/RhoLowerBoundLargeConductor.lean}{Formalism.MainNotions.flipOneThree.Pf.dL}
\leanref{Formalism/SharpRestriction/GrowthBounds/ColumnSpanDensity/RhoLowerBoundLargeConductor.lean}{Formalism.MainNotions.flipOneThree.Pf.flatW}
\leanref{Formalism/SharpRestriction/GrowthBounds/ColumnSpanDensity/RhoLowerBoundLargeConductor.lean}{Formalism.MainNotions.flipOneThree.Pf.topW}

\begin{lemma}\label{pd-kappa-edge-count}
Let $p\geq5$ be a prime, $(W,Q,\sigma,\iota)$ an unramified quadratic Hida setting of degree $n'$
at $p$, $k\geq1$, and $L\subset Q^k$ a $W$-lattice. Among the edges with $r=1$
(Definition~\ref{pd-projection-property}), at least $[k]=(p^{n'k}-1)/(p^{n'}-1)$ have target $L$ and
at least $[k]$ have source $L$.
\end{lemma}
\leanref{Formalism/SharpRestriction/GrowthBounds/ColumnSpanDensity/RhoLowerBoundOfEdgeCount.lean}{Formalism.MainNotions.flipOneThree.Pf.RestEdgeCount}
\leanref{Formalism/SharpRestriction/GrowthBounds/ColumnSpanDensity/RhoUpperBoundSmallDepth.lean}{Formalism.MainNotions.flipOneThree.Pf.restEdgeCount}
\leanref{Formalism/SharpRestriction/GrowthBounds/ColumnSpanDensity/RhoUpperBoundOfEdgeCount.lean}{Formalism.MainNotions.flipOneThree.Pf.RestEdgeCountSrc}
\leanref{Formalism/SharpRestriction/GrowthBounds/ColumnSpanDensity/RhoUpperBoundSmallDepth.lean}{Formalism.MainNotions.flipOneThree.Pf.restEdgeCountSrc}

\begin{proof}
The residue field $F=W/pW$ has $p^{n'}$ elements. An edge with target $L$ is a lattice $L''$ with
$pL\subset L''\subset L$ and $|L/L''|=p^{n'}$, that is, the inverse image of a hyperplane of the
$F$-vector space $L/pL\cong F^k$. An edge with source $L$ is a lattice $L'$ with
$L\subset L'\subset p^{-1}L$ and $|L'/L|=p^{n'}$, that is, the inverse image of a line of
$p^{-1}L/L\cong F^k$. The space $F^k$ has $[k]$ lines and $[k]$ hyperplanes.
\end{proof}

\begin{lemma}\label{pd-kappa-rank-monotone}
Let $p$, $(W,Q,\sigma,\iota)$, $k$ and $\lambda_0$ be as in Definition~\ref{pd-kappa-split-rank},
and let $t,t',c\geq0$ be integers.
\begin{enumerate}
\item If $\mathrm{ER}(\lambda_0;t,c,D,X)$ holds, $t\leq t'$, $D\leq D'$ and $X'\leq X$, then
  $\mathrm{ER}(\lambda_0;t',c,D',X')$ holds.
\item If $X\leq0$, then $\mathrm{ER}(\lambda_0;t,c,D,X)$ holds for every $D$.
\item Let $D\in\mathbf R$. The set of $\kappa\geq0$ such that
  $\mathrm{ER}(\lambda_0;t,c,D,(1-\kappa)\mu(\lambda_0))$ holds contains $1$ and contains its
  infimum.
\item Let $r,D,c_*\geq0$ be integers with $4(D+1)\leq c_*$, and let $\kappa\geq1$. Then
  $\mathrm{LP}(r,\lambda_0;t,D,\kappa,c_*)$ holds (Definition~\ref{pd-projection-property}).
\end{enumerate}
\end{lemma}
\leanref{Formalism/SharpRestriction/GrowthBounds/ColumnSpanDensity/LatticeCountPrimeMonoLevel.lean}{Formalism.MainNotions.flipOneThree.Pf.RhoL_mono}
\leanref{Formalism/SharpRestriction/GrowthBounds/ColumnSpanDensity/LatticeCountPrimeMonoLevel.lean}{Formalism.MainNotions.flipOneThree.Pf.RhoL_zero}
\leanref{Formalism/SharpRestriction/GrowthBounds/ColumnSpanDensity/LatticeCountPrimeMonoLevel.lean}{Formalism.MainNotions.flipOneThree.Pf.rho_inf_mem}
\leanref{Formalism/SharpRestriction/GrowthBounds/ColumnSpanDensity/LatticeCountPrimeMonoLevel.lean}{Formalism.MainNotions.flipOneThree.Pf.lpprime_of_one_le}

\begin{proof}
(1) For an $\mathcal R$-module $H$, $\varpi_0^{-t}H\subset\varpi_0^{-t'}H$, so the cores for $t$ are
contained in those for $t'$. A split family of depth $D$ and size $n\geq X$ is one of depth $D'$ and
size $n\geq X'$.

(2) Take the empty family at every edge.

(3) The value $\kappa=1$ belongs to the set by (2). Let $\kappa_0$ be the infimum and
$X_0=(1-\kappa_0)\mu(\lambda_0)$. If $X_0\leq0$, then $\kappa_0$ belongs to the set by (2). Otherwise
$\mu(\lambda_0)>0$, and there is $\kappa$ in the set with
$(1-\kappa)\mu(\lambda_0)>\lceil X_0\rceil-1$. At every edge there is then a split family of size
$n>\lceil X_0\rceil-1$. Since $n$ is an integer, $n\geq\lceil X_0\rceil\geq X_0$.

(4) Take no edges and $n=0$ in Definition~\ref{pd-projection-property}(3); then
$n\geq(1-\kappa)\dim V$.
\end{proof}

\begin{lemma}\label{pd-kappa-rank-criterion}
Let $p$, $(W,Q,\sigma,\iota)$, $k$ and $\lambda_0$ be as in Definition~\ref{pd-kappa-split-rank},
let $t,c,D\geq0$ be integers with $8D<c$ and $4(D+1)\leq c$, and let $\kappa\in\mathbf R$. If
$\mathrm{ER}(\lambda_0;t,c,D,(1-\kappa)\mu(\lambda_0))$ holds, then the projection property
$\mathrm{LP}(1,\lambda_0;t,D,\kappa,c)$ holds (Definition~\ref{pd-projection-property}).
\end{lemma}
\leanref{Formalism/SharpRestriction/GrowthBounds/ColumnSpanDensity/RhoLowerBoundLargeConductor.lean}{Formalism.MainNotions.flipOneThree.Pf.RhoLowerBoundLargeConductor}
\leanref{Formalism/SharpRestriction/GrowthBounds/ColumnSpanDensity/RhoLowerBoundOfEdgeCount.lean}{Formalism.MainNotions.flipOneThree.Pf.rho_lower_bound_of_edge_count}
\leanref{Formalism/SharpRestriction/GrowthBounds/ColumnSpanDensity/RhoUpperBoundSmallDepth.lean}{Formalism.MainNotions.flipOneThree.Pf.rho_lower_bound_large_conductor}
\leanref{Formalism/SharpRestriction/GrowthBounds/ColumnSpanDensity/IsUnitDetOfOneAdd.lean}{Formalism.MainNotions.flipOneThree.Pf.goodRank_core}

\begin{proof}
Let $L\subset Q^k$ be a $W$-lattice. By Lemma~\ref{pd-kappa-edge-count} there are distinct edges
$e_1,\dots,e_E$ with target $L$ and $E\geq[k]$. For each $j$ take a split family
$(x_{j,i})_i$, $(y_{j,i})_i$, $z_j$, $M_j$ of depth $D$ and size $n_j\geq(1-\kappa)\mu(\lambda_0)$
on $(\mathcal C_{t,c}(e_j),\mathcal C'_{t,c}(e_j))$, with $z_j=1$ if $n_j=0$. Since $D$ is an
integer and $v(z_j)\leq D$, $p^D=z_jq_j$ with $q_j\in\mathcal R$. Put
$y'_{j,l}=\sum_{l'}(M_j)_{l'l}\,y_{j,l'}$. Then $y'_{j,l}\in\mathcal C'_{t,c}(e_j)$ and
$\langle x_{j,i},y'_{j,l}\rangle=z_j\delta_{il}$.

Let $j\neq j'$. Then $x_{j,i}\in R(e_j)$, and $y'_{j',l}\in\varpi_0^cR'(e_j)$ because $e_j$ is an
edge other than $e_{j'}$ with target $L$. By Lemma~\ref{pd-edge-basis}(3),
$\langle x_{j,i},y'_{j',l}\rangle\in\varpi_0^c\mathcal R$. Number the pairs $(j,i)$,
$1\leq i\leq n_j$, as $1,\dots,n$ with $n=\sum_jn_j$, and write $x_\rho$, $y'_\rho$, $z_\rho$, $q_\rho$
and $e_\rho$ for the corresponding $x_{j,i}$, $y'_{j,i}$, $z_j$, $q_j$ and $e_j$. Then
$\langle x_\rho,y'_\tau\rangle=z_\tau\delta_{\rho\tau}+\varpi_0^cb_{\rho\tau}$ with
$b_{\rho\tau}\in\mathcal R$. Since $\varpi_0^2=p$ and $c\geq2D+1$, the matrix
$P_{\rho\tau}=\delta_{\rho\tau}+q_\tau\varpi_0^{c-2D}b_{\rho\tau}$ satisfies
$\langle x_\rho,y'_\tau\rangle=P_{\rho\tau}z_\tau$, and it is congruent to the identity modulo
$\varpi_0$, so it is invertible over $\mathcal R$. Let $(B_{j\rho})=P^{-1}$.

The cores $\mathcal C_{t,c}(e_\rho)$ lie in $\varpi_0^{-t}H(L''_\rho)$, in $R(e_\rho)$ and in
$\varpi_0^cR(e')\subset R(e')$ for the other edges $e'$ with the source $L''_\rho$ of $e_\rho$. So
$x_\rho\in N_{\mathrm{src}}(L''_\rho)$ and $x_\rho\in\varpi_0^cR(e')$ for every edge $e'\neq e_\rho$
with source $L''_\rho$. In the same way $y'_\tau\in\mathcal C'_{t,c}(e_\tau)\subset
N_{\mathrm{tgt}}(L)$. Let $u(y)=\sum_\tau q_\tau y_\tau y'_\tau$ for $y\in\mathcal R^n$; it lies in
$N_{\mathrm{tgt}}(L)$, and
\[
 \langle x_\rho,u(y)\rangle=\sum_\tau P_{\rho\tau}z_\tau q_\tau y_\tau=p^D(Py)_\rho,\qquad
 \sum_\rho B_{j\rho}\langle x_\rho,u(y)\rangle=p^Dy_j .
\]
Finally $n\geq E(1-\kappa)\mu(\lambda_0)\geq(1-\kappa)\dim V$ if $\kappa\leq1$, because $E\geq[k]$,
and $n\geq0\geq(1-\kappa)\dim V$ otherwise. With $4(D+1)\leq c$, this is the condition of
Definition~\ref{pd-projection-property}(3) at $L$, with the edges $e_\rho$ and the vectors
$F_\rho=x_\rho$.
\end{proof}

\begin{lemma}\label{pd-kappa-rank-upper}
Let $p$, $(W,Q,\sigma,\iota)$, $k$ and $\lambda_0$ be as in Definition~\ref{pd-kappa-split-rank},
let $t,c,n\geq0$ be integers and $D\geq0$ a real with $2D<c$. If for every edge $e$ there is a split
family of depth $D$ and size $n$ on $(\mathcal C_{t,c}(e),\mathcal C'_{t,c}(e))$, then
$n\leq\mu(\lambda_0)$.
\end{lemma}
\leanref{Formalism/SharpRestriction/GrowthBounds/ColumnSpanDensity/RhoLowerBoundLargeConductor.lean}{Formalism.MainNotions.flipOneThree.Pf.RhoUpperBoundSmallDepth}
\leanref{Formalism/SharpRestriction/GrowthBounds/ColumnSpanDensity/RhoUpperBoundOfEdgeCount.lean}{Formalism.MainNotions.flipOneThree.Pf.rho_upper_bound_of_edge_count}
\leanref{Formalism/SharpRestriction/GrowthBounds/ColumnSpanDensity/RhoUpperBoundSmallDepth.lean}{Formalism.MainNotions.flipOneThree.Pf.rho_upper_bound_small_depth}

\begin{proof}
We may assume $n\geq1$. By Lemma~\ref{pd-kappa-edge-count} there are distinct edges
$e_1,\dots,e_E$ with source $W^k$ and $E\geq[k]$. For each $j$ take a split family
$(x_{j,a})_a$, $(y_{j,b})_b$, $z_j$, $M_j$ of depth $D$ and size $n$ on the cores at $e_j$, and put
$y'_{j,b}=\sum_l(M_j)_{lb}\,y_{j,l}\in\mathcal C'_{t,c}(e_j)$, so that
$\langle x_{j,a},y'_{j,b}\rangle=z_j\delta_{ab}$. The valuation of $z_j$ in powers of $\varpi_0$ is
$2v(z_j)\leq2D<c$, so $\varpi_0^{c-1}=z_jz'_j$ with $z'_j\in\mathcal R$. For $j\neq l$,
$x_{j,a}\in\varpi_0^cR(e_l)$, because $e_l$ is an edge other than $e_j$ with source $W^k$, and
$y'_{l,b}\in R'(e_l)$. By Lemma~\ref{pd-edge-basis}(3),
$\langle x_{j,a},y'_{l,b}\rangle\in\varpi_0^c\mathcal R=\varpi_0z_lz'_l\mathcal R$. Hence
$\langle x_{j,a},y'_{l,b}\rangle=(\delta_{(j,a),(l,b)}+\beta_{(j,a),(l,b)})z_l$ with all
$\beta_{(j,a),(l,b)}\in\varpi_0\mathcal R$, and the matrix $P=1+\beta$ is invertible over
$\mathcal R$. If $\sum_{j,a}g_{j,a}x_{j,a}=0$ with $g_{j,a}\in\mathcal K$, pairing with $y'_{l,b}$ and
dividing by $z_l$ gives $gP=0$, so $g=0$. Thus the $En$ vectors $x_{j,a}$ are linearly independent
in $V$, and $n[k]\leq nE\leq\dim V$.
\end{proof}

\begin{lemma}\label{pd-kappa-edge-move}
Let $p$, $(W,Q,\sigma,\iota)$, $k$ and $\lambda_0$ be as in Definition~\ref{pd-kappa-split-rank},
let $t,c,n\geq0$ be integers, $D\in\mathbf R$, and $e$ and $e'$ edges. If there is a split family of
depth $D$ and size $n$ on $(\mathcal C_{t,c}(e),\mathcal C'_{t,c}(e))$, then there is one on
$(\mathcal C_{t,c}(e'),\mathcal C'_{t,c}(e'))$.
\end{lemma}
\leanref{Formalism/SharpRestriction/GrowthBounds/ColumnSpanDensity/RhoLowerBoundLargeConductor.lean}{Formalism.MainNotions.flipOneThree.Pf.RestEdge}
\leanref{Formalism/SharpRestriction/GrowthBounds/ColumnSpanDensity/EdgeIndep.lean}{Formalism.MainNotions.flipOneThree.Pf.restEdge_of}
\leanref{Formalism/SharpRestriction/GrowthBounds/ColumnSpanDensity/EdgeIndep.lean}{Formalism.MainNotions.flipOneThree.Pf.RestFunct}
\leanref{Formalism/SharpRestriction/GrowthBounds/ColumnSpanDensity/InvariantSectionsTransportGL.lean}{Formalism.MainNotions.flipOneThree.Pf.restFunct}

\begin{proof}
By the elementary divisor theorem, as in the proof of Lemma~\ref{pd-edge-basis}(1), the source of
$e$ has a $W$-basis $b_1,\dots,b_k$ such that $p^{-1}b_1,b_2,\dots,b_k$ is a basis of its target,
and the same holds for $e'$. So $e'=\gamma e$ for some $\gamma\in GL_k(Q)$. Let $\gamma$ act on $V$
and, by $G\mapsto G\circ\gamma^{-1}$, on $V^\vee$, as in the proof of
Lemma~\ref{pd-edge-transport}. This action carries $H(M)$ onto $H(\gamma M)$ for every $W$-lattice
$M$, and $R(f)$ onto $R(\gamma f)$ and $R'(f)$ onto $R'(\gamma f)$ for every edge $f$, and it preserves
the edge pairing. The map $f\mapsto\gamma f$ is a bijection from the edges with the source of $e$
onto the edges with the source of $e'$, and from the edges with the target of $e$ onto the edges with
the target of $e'$. Hence $\gamma$ carries $\mathcal C_{t,c}(e)$ onto $\mathcal C_{t,c}(e')$ and
$\mathcal C'_{t,c}(e)$ onto $\mathcal C'_{t,c}(e')$, and it carries a split family on the first
pair to a split family on the second with the same $z$ and $(M_{lj})$.
\end{proof}

\begin{lemma}\label{pd-kappa-setting-change}
Let $p\geq5$ be a prime, $n'\geq2$ even, $k\geq1$ and $r,t,D,c_*\geq0$ integers,
$\kappa\in\mathbf R$, and $\lambda_0(s)\in\mathbf Z^k$, $s\in\mathbf Z/n'\mathbf Z$, dominant
weights. If the projection property $\mathrm{LP}(r,\lambda_0;t,D,\kappa,c_*)$ holds for one
unramified quadratic Hida setting of degree $n'$ at $p$ (Definition~\ref{pd-projection-property}),
it holds for every other.
\end{lemma}
\leanref{Formalism/SharpRestriction/GrowthBounds/ColumnSpanDensity/RhoLowerBoundLargeConductor.lean}{Formalism.MainNotions.flipOneThree.Pf.LPRingIndep}
\leanref{Formalism/SharpRestriction/GrowthBounds/ColumnSpanDensity/LpRingIndepOf.lean}{Formalism.MainNotions.flipOneThree.Pf.lpRingIndep_of''}
\leanref{Formalism/SharpRestriction/GrowthBounds/ColumnSpanDensity/LpRingIndepOf.lean}{Formalism.MainNotions.flipOneThree.Pf.restRSIso_of}
\leanref{Formalism/SharpRestriction/GrowthBounds/ColumnSpanDensity/RestOIso.lean}{Formalism.MainNotions.flipOneThree.Pf.restOIso}
\leanref{Formalism/SharpRestriction/GrowthBounds/ObjBridge.lean}{Formalism.MainNotions.G.WittIdent}

\begin{proof}
Let $(W_i,Q_i,\sigma_i,\iota_i)$, with rings $\mathcal R_i\subset\mathcal K_i$ and $\varpi_i$, for
$i=1,2$, be two such settings. Each $W_i$ is a complete discrete valuation ring with uniformizer $p$
and a residue field with $p^{n'}$ elements, so it is the ring of Witt vectors of that field, and a
ring homomorphism from it to a complete ring of this kind is determined by its reduction
\cite[Chapter~II, \S5]{Ser79}. Hence there is an isomorphism $\phi\colon W_1\to W_2$, and
$\phi\sigma_1=\sigma_2\phi$ on $W_1$, because both sides reduce to $x\mapsto\phi(x)^p$. Each
$\mathcal R_i$ is complete, finite over $\iota_i(W_i)$, with uniformizer $\varpi_i$, $\varpi_i^2=p$
and a residue field with $p^{2n'}$ elements. So it contains the unramified extension of degree $2$
of $\iota_i(W_i)$ and is generated over it by $\varpi_i$, a root of $X^2-p$
\cite[Chapter~III, \S5]{Ser79}. Therefore $\phi$ extends to an isomorphism
$\mathcal R_1\to\mathcal R_2$ compatible with $\iota_1$ and $\iota_2$ that carries $\varpi_1$ to
$\varpi_2$, and to the fraction fields. The space $V$, the lattices $H(L)$ and $H(e)$, the edges, the
edge pairing and the projection property are defined from these data, so the isomorphisms carry the
projection property in the first setting to the projection property in the second.
\end{proof}

\begin{lemma}\label{pd-kappa-weight-tensor}
Let $p$, $(W,Q,\sigma,\iota)$, $k$ and $\lambda_0$ be as in Definition~\ref{pd-kappa-split-rank},
and let $T\subset\mathbf Z/n'\mathbf Z$. Let $\lambda_0^T$ be the top part of $\lambda_0$ and
$\lambda_0^O=\lambda_0-\lambda_0^T$, so that $\lambda_0^O(s)=\lambda_0(s)$ for $s\notin T$ and
$\lambda_0^O(s)=0$ for $s\in T$. Write $V^T$, $H^T$, $R^T$, $R'^T$, $A^T$ and
$\langle\ ,\ \rangle^T$, respectively $V^O$, $H^O$, $R^O$, $R'^O$, $A^O$ and
$\langle\ ,\ \rangle^O$, for the space, lattices, integer and edge pairing of
Definition~\ref{pd-projection-property} with $r=1$ formed with $\lambda_0^T$, respectively
$\lambda_0^O$, and $V$, $H$, $R$, $R'$, $A$ and $\langle\ ,\ \rangle$ for those formed with
$\lambda_0$.
\begin{enumerate}
\item For $F\in V^T$ and $F'\in V^O$ the product $FF'$ lies in $V$, and $F\otimes F'\mapsto FF'$ is
  an isomorphism $V^T\otimes_{\mathcal K}V^O\to V$. In particular $\dim V=\dim V^T\dim V^O$.
  Moreover $A=A^T+A^O$.
\item For every $W$-lattice $L$, $H(L)$ is the $\mathcal R$-span of the products $FF'$ with
  $F\in H^T(L)$ and $F'\in H^O(L)$. For every edge $e$, $H(e)$ is the $\mathcal R$-span of the
  products $FF'$ with $F\in H^T(e)$ and $F'\in H^O(e)$.
\item For $G\in(V^T)^\vee$ and $G'\in(V^O)^\vee$ let $G\otimes G'\in V^\vee$ be the form with
  $(G\otimes G')(FF')=G(F)G'(F')$. Then
  $\langle FF',G\otimes G'\rangle=\langle F,G\rangle^T\langle F',G'\rangle^O$. Let $i,j$ be integers.
  If $G\in\varpi_0^iH^T(L)^\vee$ and $G'\in\varpi_0^jH^O(L)^\vee$ for a $W$-lattice $L$, then
  $G\otimes G'\in\varpi_0^{i+j}H(L)^\vee$. If $G\in\varpi_0^iR'^T(e)$ and
  $G'\in\varpi_0^jR'^O(e)$ for an edge $e$, then $G\otimes G'\in\varpi_0^{i+j}R'(e)$.
\end{enumerate}
\end{lemma}
\leanref{Formalism/SharpRestriction/GrowthBounds/ColumnSpanDensity/FactorsLEOfIndex.lean}{Formalism.MainNotions.flipOneThree.Pf.RestTensor}
\leanref{Formalism/SharpRestriction/GrowthBounds/ColumnSpanDensity/RestTensor.lean}{Formalism.MainNotions.flipOneThree.Pf.restTensor}
\leanref{Formalism/SharpRestriction/GrowthBounds/ColumnSpanDensity/Tensor.lean}{Formalism.MainNotions.flipOneThree.Pf.restTensor_of}
\leanref{Formalism/SharpRestriction/GrowthBounds/ColumnSpanDensity/FactorsLEOfIndex.lean}{Formalism.MainNotions.flipOneThree.Pf.Ar_add}
\leanref{Formalism/SharpRestriction/GrowthBounds/ColumnSpanDensity/FactorsLEOfIndex.lean}{Formalism.MainNotions.flipOneThree.Pf.nu_scale}

\begin{proof}
At each $s$ one of $\lambda_0^T(s)$ and $\lambda_0^O(s)$ is $\lambda_0(s)$ and the other is $0$. The
weight $\lambda(s)$ formed from $0$ is $0$, and the dual Weyl module of the natural weight $0$, over
$\mathcal K$ or over $\mathcal R$, consists of the constants, since its elements have degree $0$ in
every column. A constant is unchanged by the moves of Definition~\ref{pd-projection-property}(1),
because the determinant factor is $\det(g)^0=1$. Hence $V^T$ is the span of the products
$\prod_{s\in T}f_s(x_s)$ and $V^O$ that of the products $\prod_{s\notin T}f_s(x_s)$, with $f_s$ in the
dual Weyl module over $\mathcal K$ of the natural weight of $\lambda(s)$, and $H^T(L)$ and $H^O(L)$ are
the $\mathcal R$-spans of the same products with $f_s\in\nabla_{\lambda(s)}(M_s^\vee)$; the same
holds for $H^T(e)$ and $H^O(e)$.

The spaces $V^T$ and $V^O$ consist of polynomials in the disjoint sets of variables $x_s$, $s\in T$,
and $x_s$, $s\notin T$. Multiplication is an isomorphism from the tensor product of the polynomial
rings in these two sets of variables onto the polynomial ring in all the $x_s$. It carries
$V^T\otimes V^O$ onto the span of the products $\prod_sf_s(x_s)$, which is $V$. Since $H(L)$ and
$H(e)$ are the $\mathcal R$-spans of the products $\prod_sf_s(x_s)$ with $f_s$ in the lattices at $s$,
this gives (1) and (2), except the last claim of (1). With $r=1$, $A=\sum_s\lambda(s)_1$, and each
$\lambda(s)_1$ is the sum of the corresponding entries for $\lambda_0^T$ and $\lambda_0^O$, one of
which is $0$.

(3) The form $G\otimes G'$ exists by (1). By (1),
$\langle FF',G\otimes G'\rangle=(G\otimes G')(p^{A^T}F\cdot p^{A^O}F')=\langle F,G\rangle^T\langle
F',G'\rangle^O$. Let $G=\varpi_0^iG_0$ and $G'=\varpi_0^jG'_0$ with $G_0(H^T(L))\subset\mathcal R$
and $G'_0(H^O(L))\subset\mathcal R$. Then $(G_0\otimes G'_0)(FF')\in\mathcal R$ for $F\in H^T(L)$ and
$F'\in H^O(L)$, so $G_0\otimes G'_0$ takes $H(L)$ into $\mathcal R$ by (2), and
$G\otimes G'=\varpi_0^{i+j}G_0\otimes G'_0\in\varpi_0^{i+j}H(L)^\vee$. The same argument with $H(e)$
shows that $G_0\otimes G'_0\in H(e)^\vee$ when $G_0\in H^T(e)^\vee$ and $G'_0\in H^O(e)^\vee$. Since
$R'=\varpi_0^{-A}H(e)^\vee$ and $A=A^T+A^O$, the claim for $R'(e)$ follows.
\end{proof}

\begin{lemma}\label{pd-kappa-weight-order}
Let $p$, $(W,Q,\sigma,\iota)$, $k$ and $\lambda_0$ be as in Definition~\ref{pd-kappa-split-rank},
and let $c\geq0$ be an integer with $|\lambda_0(s)_i|\leq c$ for all $s$ and $i$. Let $e$ be an edge
with source $L''$ and target $L$. Then $R(e)\subset\varpi_0^{-2n'c}H(L'')$ and
$R'(e)\subset\varpi_0^{-2n'c}H(L)^\vee$.
\end{lemma}
\leanref{Formalism/SharpRestriction/GrowthBounds/ColumnSpanDensity/FactorsLEOfIndex.lean}{Formalism.MainNotions.flipOneThree.Pf.RestOrdBound}
\leanref{Formalism/SharpRestriction/GrowthBounds/ColumnSpanDensity/ResidueClassCountOfTensorBasisInclusion.lean}{Formalism.MainNotions.flipOneThree.Pf.restOrdBound}
\leanref{Formalism/SharpRestriction/GrowthBounds/ColumnSpanDensity/ResidueClassCountOfTensorBasisInclusion.lean}{Formalism.MainNotions.flipOneThree.Pf.steps_rev}
\leanref{Formalism/SharpRestriction/GrowthBounds/ColumnSpanDensity/ResidueClassCountOfTensorBasisInclusion.lean}{Formalism.MainNotions.flipOneThree.Pf.Ar_sub_le}
\leanref{Formalism/SharpRestriction/GrowthBounds/ColumnSpanDensity/EdgeLatRS.lean}{Formalism.MainNotions.flipOneThree.Pf.P_rev}

\begin{proof}
Put $E=\sum_s\lambda(s)_k$. We first show $\varpi_0^{-E}H(e)\subset H(L'')$ and
$\varpi_0^{-E}H(L)\subset H(e)$. Let $x_i$ and $\pi_i$ be the basis and the integers constructed in
the proof of Lemma~\ref{pd-edge-basis}(1) with $r=1$, so that $H(L'')$, $H(e)$ and $H(L)$ are the
$\mathcal R$-spans of the $x_i$, of the $\varpi_0^{\pi_i}x_i$ and of the $\varpi_0^{2\pi_i}x_i$.
There $\pi_i=\sum_s\mu^{(s)}_1$ with $\mu^{(s)}=\beta+\lambda(s)_k(1,\dots,1)$, where $f(tx)=t^\beta f$
for the diagonal matrices $t$ and a polynomial $f$ in the dual Weyl module of the natural weight of
$\lambda(s)$. The entry $\beta_1$ is the degree of $f$ in the first row, so $\beta_1\geq0$ and
$\mu^{(s)}_1\geq\lambda(s)_k$. Hence $\pi_i\geq E$, and $\varpi_0^{-E}\varpi_0^{\pi_i}x_i$ and
$\varpi_0^{-E}\varpi_0^{2\pi_i}x_i=\varpi_0^{\pi_i-E}\varpi_0^{\pi_i}x_i$ lie in $H(L'')$ and in
$H(e)$ respectively.

The entries of $\lambda(s)$ are those of $\lambda_0(s)$ or their negatives, so
$A-E=\sum_s(\lambda(s)_1-\lambda(s)_k)\leq2n'c$. Therefore
$R(e)=\varpi_0^{-A}H(e)=\varpi_0^{E-A}\varpi_0^{-E}H(e)\subset\varpi_0^{E-A}H(L'')\subset
\varpi_0^{-2n'c}H(L'')$. Let $G\in R'(e)$, so $G=\varpi_0^{-A}G_0$ with $G_0(H(e))\subset\mathcal R$.
For $h\in H(L)$ we have $\varpi_0^{-E}h\in H(e)$, so $G_0(h)\in\varpi_0^E\mathcal R$ and
$\varpi_0^{2n'c}G(h)\in\varpi_0^{2n'c-A+E}\mathcal R\subset\mathcal R$. Hence
$\varpi_0^{2n'c}G\in H(L)^\vee$.
\end{proof}

\begin{lemma}\label{pd-kappa-cloud-dimensions}
For a dominant weight $\mu\in\mathbf Z^j$ put
$\dim_j(\mu)=\prod_{1\leq i<l\leq j}(\mu_i-\mu_l+l-i)/(l-i)$. Let $m\geq2$ and $P,Q,c\geq0$ be
integers, put $h(P,Q)=(P,0,\dots,0,-Q)\in\mathbf Z^{m+1}$, let $\lambda\in\mathbf Z^{m+1}$ be dominant
with $\lambda_1=P+c$, $\lambda_{m+1}=-Q-c$ and $|\lambda_i|\leq c$ for $1<i<m+1$, and put
$\lambda^{\mathrm{cl}}=(\lambda_2,\dots,\lambda_m)$. Let $t\geq0$ be a real number with $2c\leq tP$
and $2c\leq tQ$. Then $\dim_{m-1}(\lambda^{\mathrm{cl}})>0$,
\[
 \dim_{m-1}(\lambda^{\mathrm{cl}})\dim_{m+1}(h(P,Q))\leq\dim_{m+1}(\lambda)
 \leq(1+t)^{2m-1}\dim_{m-1}(\lambda^{\mathrm{cl}})\dim_{m+1}(h(P,Q)),
\]
and $\dim_{m+1}(h(P+2c,Q+2c))\leq(1+t)^{2m-1}\dim_{m+1}(h(P,Q))$.
\end{lemma}
\leanref{Formalism/SharpRestriction/GrowthBounds/ColumnSpanDensity/CloudBoundOfPropAndWeylDimension.lean}{Formalism.MainNotions.flipOneThree.Pf.top_bounds}
\leanref{Formalism/SharpRestriction/GrowthBounds/ColumnSpanDensity/CloudBoundOfPropAndWeylDimension.lean}{Formalism.MainNotions.flipOneThree.Pf.cwt_lower}
\leanref{Formalism/PolynomialDimension/MultinomialDepth/WeylDimPosOfAntitone.lean}{Formalism.CountingSumsFrames.cwt_bound}
\leanref{Formalism/PolynomialDimension/MultinomialDepth/WeylDimPosOfAntitone.lean}{Formalism.CountingSumsFrames.hwt_bound}
\leanref{Formalism/PolynomialDimension/MultinomialDepth/WeylDimPosOfAntitone.lean}{Formalism.CountingSumsFrames.weylDim_cwt}

\begin{proof}
The number $\dim_j(\mu)$ is unchanged when a constant is added to all entries of $\mu$, and it is
positive for dominant $\mu$. Put $\mu_i=\lambda_{i+1}+c$ for $i<m$; then $\mu$ is dominant,
$0\leq\mu_i\leq2c$, $\dim_{m-1}(\mu)=\dim_{m-1}(\lambda^{\mathrm{cl}})>0$ and
$\lambda+c(1,\dots,1)=(P+2c,\mu_1,\dots,\mu_{m-1},-Q)$. Separating in Weyl's product the factors
with $i=1$ or $l=m+1$ gives
\[
 \dim_{m+1}(\lambda)=\dim_{m-1}(\mu)\prod_{j=1}^{m-1}\frac{P+2c-\mu_j+j}{j}
 \prod_{j=1}^{m-1}\frac{\mu_j+Q+m-j}{m-j}\cdot\frac{P+Q+2c+m}{m},
\]
and with $c=0$ and $\mu=0$ it gives the same formula for $\dim_{m+1}(h(P,Q))$, without the factor
$\dim_{m-1}(\mu)$. Since $0\leq\mu_j\leq2c$, each of the $2m-1$ factors of the first formula is at
least the corresponding factor for $h(P,Q)$, and at most $1+t$ times it, because
$P+2c+j\leq(1+t)(P+j)$, $Q+2c+m-j\leq(1+t)(Q+m-j)$ and $P+Q+2c+m\leq(1+t)(P+Q+m)$. This gives the
two inequalities for $\dim_{m+1}(\lambda)$. The factors of $\dim_{m+1}(h(P+2c,Q+2c))$ are
$(P+2c+j)/j$, $(Q+2c+m-j)/(m-j)$ and $(P+Q+4c+m)/m$, each at most $1+t$ times the corresponding
factor for $h(P,Q)$, since $4c\leq t(P+Q)$.
\end{proof}

\begin{lemma}\label{pd-kappa-cloud-cartan}
Let $\mathcal O$ be a discrete valuation ring with fraction field $K$ and uniformizer $\varpi$, let
$k\geq1$, and let $x$ be a $k\times k$ matrix of variables. For a dominant $\mu\in\mathbf Z^k$ write
$\nabla_K(\mu)$ and $\nabla_{\mathcal O}(\mu)$ for the dual Weyl modules over $K$ and over
$\mathcal O$ of the natural weight $\nu(\mu)=(\mu_1-\mu_k,\dots,\mu_k-\mu_k)$ of $\mu$
(Definition~\ref{pd-weight-lattices}(1)). For an $\mathcal O$-lattice $\Lambda\subset K^k$ let
$\nabla_\mu(\Lambda^\vee)\subset K[x]$ be the set of $f$ with
$\det(g)^{\mu_k}f(gx)\in\nabla_{\mathcal O}(\mu)$ for every $g\in GL_k(K)$ whose columns form an
$\mathcal O$-basis of $\Lambda$, as in Definition~\ref{pd-projection-property}(1). Let
$\mu,\mu'\in\mathbf Z^k$ be dominant and $\Lambda\subset K^k$ an $\mathcal O$-lattice.
\begin{enumerate}
\item If $\det(g)^{\mu_k}f(gx)\in\nabla_{\mathcal O}(\mu)$ for one $g$ whose columns form a basis of
  $\Lambda$, then $f\in\nabla_\mu(\Lambda^\vee)$. Moreover $\nabla_\mu(\Lambda^\vee)\subset
  \nabla_K(\mu)$, it is the $\mathcal O$-span of a $K$-basis of $\nabla_K(\mu)$, and for every
  $f\in\nabla_K(\mu)$ there is $n\geq0$ with $\varpi^nf\in\nabla_\mu(\Lambda^\vee)$.
\item $\nabla_{\mu+\mu'}(\Lambda^\vee)$ is the $\mathcal O$-span of the products $ff'$ with
  $f\in\nabla_\mu(\Lambda^\vee)$ and $f'\in\nabla_{\mu'}(\Lambda^\vee)$.
\item If $f\in\nabla_\mu(\Lambda^\vee)$, $f'\in\nabla_{\mu'}(\Lambda^\vee)$ and
  $ff'\in\varpi\nabla_{\mu+\mu'}(\Lambda^\vee)$, then $f\in\varpi\nabla_\mu(\Lambda^\vee)$ or
  $f'\in\varpi\nabla_{\mu'}(\Lambda^\vee)$.
\end{enumerate}
\end{lemma}
\leanref{Formalism/SharpRestriction/GrowthBounds/HermitianTensor/MemInvariantSectionsIffIntegral.lean}{Formalism.MainNotions.flipOneThree.Pf.mem_invariant_sections_iff_integral}
\leanref{Formalism/SharpRestriction/GrowthBounds/HermitianTensor/CloudIso.lean}{Formalism.MainNotions.flipOneThree.Pf.exists_pow_mem}
\leanref{Formalism/SharpRestriction/GrowthBounds/ProjectionLemmaBaseLevel.lean}{Formalism.MainNotions.flipOneThree.CartanOnto}
\leanref{Formalism/LocalToGlobal/GroupLattices/FockModels/Cartan.lean}{Formalism.L2.G.Model.cartanOnto_proved}
\leanref{Formalism/LocalToGlobal/GroupLattices/FockModels/Cartan.lean}{Formalism.L2.G.Model.dualWeyl_le_prod}
\leanref{Formalism/SharpRestriction/GrowthBounds/ProjectionLemmaBaseLevel.lean}{Formalism.MainNotions.flipOneThree.CartanDomain}
\leanref{Formalism/LocalToGlobal/GroupLattices/FockModels/Cartan.lean}{Formalism.L2.G.Model.cartanDomain_proved}

\begin{proof}
(1) Let $g$ and $g'$ be matrices whose columns form bases of $\Lambda$. Then $g'=gu$ with
$u\in GL_k(\mathcal O)$. The substitution $x\mapsto ux$ maps $\mathcal O[x]$ into itself, preserves
the degree in each column and commutes with the substitutions $x\mapsto x(1+yE_{lj})$ of
Definition~\ref{pd-weight-lattices}(1), so it maps $\nabla_{\mathcal O}(\mu)$ into itself. Since
$\det(g')^{\mu_k}f(g'x)$ is $\det(u)^{\mu_k}$ times the image of $\det(g)^{\mu_k}f(gx)$ under this
substitution, and $\det(u)$ is a unit, the first claim follows. The same argument over $K$ shows that
$f\mapsto\det(g)^{\mu_k}f(gx)$ maps $\nabla_K(\mu)$ onto itself, with inverse
$\tau_g\colon F\mapsto\det(g)^{-\mu_k}F(g^{-1}x)$. The module $\nabla_{\mathcal O}(\mu)$ consists of
the elements of $\nabla_K(\mu)$ with coefficients in $\mathcal O$, because both are cut out by the
same degree and invariance conditions. So $\nabla_\mu(\Lambda^\vee)=\tau_g(\nabla_{\mathcal O}(\mu))
\subset\nabla_K(\mu)$. By the standard basis theorem \cite{DEP80}, over $\mathcal O$ and over $K$
(Lemma~\ref{is-standard-basis}(1)), the bideterminants of shape $\nu(\mu)$ form an $\mathcal O$-basis
of $\nabla_{\mathcal O}(\mu)$ and a $K$-basis of $\nabla_K(\mu)$; their images under $\tau_g$ give the
basis claim. For $f\in\nabla_K(\mu)$ some multiple $\varpi^n\tau_g^{-1}(f)$ has coefficients in
$\mathcal O$, so it lies in $\nabla_{\mathcal O}(\mu)$, and $\varpi^nf\in\nabla_\mu(\Lambda^\vee)$.

(2) The natural weight of $\mu+\mu'$ is $\nu(\mu)+\nu(\mu')$, and $\tau_g(FF')=\tau_g(F)\tau_g(F')$,
where the three maps $\tau_g$ are formed with $\mu$, $\mu'$ and $\mu+\mu'$. A product of elements of
$\nabla_{\mathcal O}(\mu)$ and $\nabla_{\mathcal O}(\mu')$ lies in $\nabla_{\mathcal O}(\mu+\mu')$, as
in the proof of Lemma~\ref{pd-edge-products}. By (1) it therefore suffices to show that
$\nabla_{\mathcal O}(\mu+\mu')$ is spanned by such products. By the standard basis theorem it is
spanned by bideterminants: products of minors of $x$, the $i$-th on the first $s_i$ columns, where
$\#\{i:s_i\geq t\}=\nu(\mu)_t+\nu(\mu')_t$ for every $t$. For a non-increasing $\nu$ the column
lengths $s$ of the diagram of $\nu$ satisfy $\#\{s\geq t\}=\nu_t$. So the multiset of the $s_i$ is
the union of the column lengths of $\nu(\mu)$ and of $\nu(\mu')$, and the bideterminant is the product
of a bideterminant of shape $\nu(\mu)$ and one of shape $\nu(\mu')$. These lie in
$\nabla_{\mathcal O}(\mu)$ and $\nabla_{\mathcal O}(\mu')$, since a minor on the first $s$ columns
has degree $1$ in each of them and is unchanged by adding a multiple of an earlier column to a later
one.

(3) Put $F=\tau_g^{-1}(f)\in\nabla_{\mathcal O}(\mu)$ and $F'=\tau_g^{-1}(f')\in\nabla_{\mathcal O}(\mu')$.
By the multiplicativity in (2), $FF'=\varpi X$ with $X\in\nabla_{\mathcal O}(\mu+\mu')$. The ring
$\mathcal O[x]/\varpi\mathcal O[x]$ is a polynomial ring over the residue field, hence a domain, so
$\varpi$ divides $F$ or $F'$ in $\mathcal O[x]$. Say $F=\varpi F_1$ with $F_1\in\mathcal O[x]$. Then
$F_1\in\nabla_K(\mu)$ has coefficients in $\mathcal O$, so $F_1\in\nabla_{\mathcal O}(\mu)$ and
$f=\varpi\tau_g(F_1)\in\varpi\nabla_\mu(\Lambda^\vee)$. The other case is the same.
\end{proof}

\begin{lemma}\label{pd-kappa-cloud-primitive}
Let $\mathcal O$, $K$, $\varpi$, $k$, $x$, $\nabla_K(\mu)$ and $\nabla_\mu(\Lambda^\vee)$ be as in
Lemma~\ref{pd-kappa-cloud-cartan}. Let $\mu,\zeta\in\mathbf Z^k$ be dominant,
$\Lambda\subset K^k$ an $\mathcal O$-lattice, and $z\in\nabla_\zeta(\Lambda^\vee)$ with
$z\notin\varpi\nabla_\zeta(\Lambda^\vee)$.
\begin{enumerate}
\item For $F\in\nabla_K(\mu)$, $zF\in\nabla_{\mu+\zeta}(\Lambda^\vee)$ if and only if
  $F\in\nabla_\mu(\Lambda^\vee)$.
\item There is a $K$-linear map $\rho\colon K[x]\to K[x]$ with
  $\rho(\nabla_{\mu+\zeta}(\Lambda^\vee))\subset\nabla_\mu(\Lambda^\vee)$ and $\rho(zF)=F$ for every
  $F\in\nabla_K(\mu)$.
\end{enumerate}
\end{lemma}
\leanref{Formalism/SharpRestriction/GrowthBounds/HermitianTensor/CloudIso.lean}{Formalism.MainNotions.flipOneThree.Pf.iso_one}
\leanref{Formalism/SharpRestriction/GrowthBounds/HermitianTensor/CloudIso.lean}{Formalism.MainNotions.flipOneThree.Pf.retract_one}

\begin{proof}
(1) If $F\in\nabla_\mu(\Lambda^\vee)$, then $zF\in\nabla_{\mu+\zeta}(\Lambda^\vee)$ by
Lemma~\ref{pd-kappa-cloud-cartan}(2). Conversely, let $zF\in\nabla_{\mu+\zeta}(\Lambda^\vee)$. By
Lemma~\ref{pd-kappa-cloud-cartan}(1) there is $n\geq0$ with $\varpi^nF\in\nabla_\mu(\Lambda^\vee)$. We
show by induction on $n$ that this implies $F\in\nabla_\mu(\Lambda^\vee)$. For $n=0$ there is nothing
to show. If $\varpi^{n+1}F\in\nabla_\mu(\Lambda^\vee)$, then $(\varpi^{n+1}F)z=\varpi(\varpi^nzF)$ with
$\varpi^nzF\in\nabla_{\mu+\zeta}(\Lambda^\vee)$. Since $z\notin\varpi\nabla_\zeta(\Lambda^\vee)$,
Lemma~\ref{pd-kappa-cloud-cartan}(3) gives $\varpi^{n+1}F=\varpi F'$ with
$F'\in\nabla_\mu(\Lambda^\vee)$. So $\varpi^nF\in\nabla_\mu(\Lambda^\vee)$, and the induction
hypothesis gives $F\in\nabla_\mu(\Lambda^\vee)$.

(2) By Lemma~\ref{pd-kappa-cloud-cartan}(1), $z\in\nabla_K(\zeta)$, so multiplication by $z$ maps
$\nabla_K(\mu)$ into $\nabla_K(\mu+\zeta)$, and $\nabla_\mu(\Lambda^\vee)$ and
$\nabla_{\mu+\zeta}(\Lambda^\vee)$ are the $\mathcal O$-spans of $K$-bases of $\nabla_K(\mu)$ and
$\nabla_K(\mu+\zeta)$. By (1), for $F\in\nabla_K(\mu)$, $zF$ lies in the second lattice if and only if
$F$ lies in the first. Lemma~\ref{pd-kappa-cloud-retraction} gives a $K$-linear map
$\rho_0\colon\nabla_K(\mu+\zeta)\to\nabla_K(\mu)$ with $\rho_0(zF)=F$ and
$\rho_0(\nabla_{\mu+\zeta}(\Lambda^\vee))\subset\nabla_\mu(\Lambda^\vee)$. Let $\rho$ be any
$K$-linear extension of $\rho_0$ to $K[x]$.
\end{proof}

\begin{lemma}\label{pd-kappa-cloud-transpose}
Let $K$ be a field, $\mathcal O\subset K$ a subring of characteristic $0$, $r\geq1$, and $Z$ an
$r\times r$ matrix of variables. Dual Weyl modules are those of Definition~\ref{pd-weight-lattices}(1),
indexed by natural weights, over $K$ or over $\mathcal O$. Let $N,e\geq0$ be integers and
$\alpha\in\mathbf Z^r$ non-increasing with $e\leq\alpha_i\leq N$ for all $i$ and $\alpha_r=e$, and put
$\alpha^*=(N-\alpha_r,\dots,N-\alpha_1)$ and $\alpha-e=(\alpha_1-e,\dots,\alpha_r-e)$. Then there is a
$K$-linear map $\tau\colon K[Z]\to K[Z]$ such that:
\begin{enumerate}
\item $\tau$ maps the dual Weyl module over $K$ of $\alpha^*$ into that of $\alpha-e$, and the dual
  Weyl module over $\mathcal O$ of $\alpha^*$ onto that of $\alpha-e$;
\item for $l\in GL_r(K)$ and $f$ in the dual Weyl module over $K$ of $\alpha^*$, if $f_l(Z)=f(lZ)$,
  then $\tau(f_l)({}^tl\,Z)=\det(l)^{N-e}\,\tau(f)(Z)$.
\end{enumerate}
\end{lemma}
\leanref{Formalism/SharpRestriction/GrowthBounds/HermitianTensor/DetFactsPAlt.lean}{Formalism.MainNotions.flipOneThree.Pf.TauP}
\leanref{Formalism/SharpRestriction/GrowthBounds/HermitianTensor/TauP.lean}{Formalism.MainNotions.flipOneThree.Pf.tauP}

\begin{proof}
Write $\operatorname{adj}(M)$ for the adjugate of a square matrix $M$, so that
$M\operatorname{adj}(M)=\operatorname{adj}(M)M=\det(M)$ and $\operatorname{adj}(MM')=
\operatorname{adj}(M')\operatorname{adj}(M)$. Put $\bar k=r+1-k$ for $1\leq k\leq r$, let $P_0$ be the
permutation matrix with $(P_0)_{kb}=1$ if $k=\bar b$ and $0$ otherwise, and $\varepsilon=\det P_0=\pm1$.
For $A$ equal to $K$ or $\mathcal O$, let $S$ be the $A$-algebra endomorphism of $A[Z]$ with
$S(f)(Z)=f({}^t\operatorname{adj}(Z)P_0)$; it replaces $Z_{ij}$ by $\operatorname{adj}(Z)_{\bar ji}$, and
it commutes with the inclusion $\mathcal O[Z]\subset K[Z]$.

\emph{Step 1.} Let $\mu\in\mathbf Z^r$ be non-increasing with $\mu_r\geq0$, and $f$ in the dual Weyl
module over $A$ of $\mu$. We show that $S(f)$ lies in the dual Weyl module over $A$ of $\psi$, where
$\psi_c=\sum_{j\neq\bar c}\mu_j$; this $\psi$ is non-increasing with $\psi_r\geq0$. The entry
$\operatorname{adj}(Z)_{ai}$ is, up to sign, the minor of $Z$ without the $i$-th row and the $a$-th
column, so it has degree $1$ in each column of $Z$ other than the $a$-th and degree $0$ in the $a$-th.
Hence $S$ replaces each variable of the $j$-th column by a polynomial of degree $1$ in every column
except the $\bar j$-th, and carries a monomial of degree $\mu_j$ in the $j$-th column, for each $j$, to
a polynomial of degree $\sum_{j:\bar j\neq c}\mu_j=\psi_c$ in the $c$-th column. Let $l<j$ and let $y$
be a variable. Since $\operatorname{adj}(1+yE_{lj})=1-yE_{lj}$, $P_0^2=1$ and
$P_0E_{jl}P_0=E_{\bar j\bar l}$,
\[
 {}^t\operatorname{adj}\big(Z(1+yE_{lj})\big)P_0={}^t\operatorname{adj}(Z)(1-yE_{jl})P_0
 ={}^t\operatorname{adj}(Z)P_0\,(1-yE_{\bar j\bar l}).
\]
As $\bar j<\bar l$, $f$ is unchanged by $Z\mapsto Z(1-yE_{\bar j\bar l})$, so $S(f)$ is unchanged by
$Z\mapsto Z(1+yE_{lj})$.

\emph{Step 2.} Put $D=\sum_i(N-\alpha_i)$ and $Q=D-(N-e)$; then $Q\geq0$, because $N-e=N-\alpha_r$ is
one of the terms of $D$. Let $f$ lie in the dual Weyl module over $A$ of $\alpha^*$. Step 1 with
$\mu=\alpha^*$, so $\mu_j=N-\alpha_{\bar j}$, gives $\psi_c=D-(N-\alpha_c)=(\alpha_c-e)+Q$. By
Lemma~\ref{pd-kappa-cloud-lift}(1) below, applied over $A$ with $\beta=\alpha-e$, which is non-increasing
with $\beta_r=0$, $S(f)=\det(Z)^Qu$ with $u$ in the dual Weyl module over $A$ of $\alpha-e$.

\emph{Step 3: the map $\tau$.} Multiplication by $\det(Z)^Q$ is injective on $K[Z]$; let $\sigma$ be a
$K$-linear left inverse of it, and put $\tau=\sigma\circ S$. By Step 2, $\tau(f)=u$ lies in the dual
Weyl module of $\alpha-e$ over $K$, respectively over $\mathcal O$, when $f$ lies in that of $\alpha^*$
over $K$, respectively over $\mathcal O$. This is the first part of (1), and $\tau$ maps the module
over $\mathcal O$ of $\alpha^*$ into that of $\alpha-e$.

\emph{Step 4: onto.} Let $g$ lie in the dual Weyl module over $\mathcal O$ of $\beta=\alpha-e$, and put
$D'=\sum_j\beta_j$ and $Q'=D'-\beta_1$. Step 1 with $\mu=\beta$ puts $S(g)$ in the module over
$\mathcal O$ of $\psi'$, with $\psi'_c=D'-\beta_{\bar c}=(\beta_1-\beta_{\bar c})+Q'$. The weight
$\gamma_c=\beta_1-\beta_{\bar c}$ is non-increasing with $\gamma_r=0$, so
Lemma~\ref{pd-kappa-cloud-lift}(1) gives $S(g)=\det(Z)^{Q'}h$ with $h$ in the module over $\mathcal O$
of $\gamma$. Put $c=N-\alpha_1$ and $f=\varepsilon^{c+D'+Q'}\det(Z)^ch$. The polynomial $\det(Z)^c$ lies
in the module over $\mathcal O$ of $(c,\dots,c)$, being a product of minors on all $r$ columns, and a
product lies in the module of the sum of the weights, as in the proof of
Lemma~\ref{pd-edge-products}. Since $c+\gamma_i=N-\alpha_{\bar i}$, $f$ lies in the module over
$\mathcal O$ of $\alpha^*$. We show $S(f)=\det(Z)^Qg$; then $\tau(f)=g$.

First, $\det\operatorname{adj}(Z)=\det(Z)^{r-1}$, so $S(\det Z)=E$ with $E=\varepsilon\det(Z)^{r-1}$.
Second, $S$ applied to the entries of $\operatorname{adj}(Z)$ gives
$\operatorname{adj}({}^t\operatorname{adj}(Z)P_0)=\varepsilon P_0\,{}^t\operatorname{adj}(\operatorname{adj}(Z))$,
and $\det(Z)\operatorname{adj}(\operatorname{adj}(Z))=\det(Z)^{r-1}Z$. Hence
$\det(Z)\,S(S(Z_{ij}))=\det(Z)\,S(\operatorname{adj}(Z)_{\bar ji})=E\,Z_{ij}$. The polynomial $g$ is
homogeneous of total degree $D'$, so $\det(Z)^{D'}S(S(g))=E^{D'}g$, and
\[
 \det(Z)^{D'}E^{Q'}S(f)=\varepsilon^{c+D'+Q'}E^c\det(Z)^{D'}S(\det(Z)^{Q'}h)
 =\varepsilon^{c+D'+Q'}E^c\det(Z)^{D'}S(S(g))=\varepsilon^{c+D'+Q'}E^{c+D'}g .
\]
Now $c+\beta_1=N-e$ gives $c+D'=(N-e)+Q'$, and $D+D'=r(N-e)$ gives $D'+Q=(r-1)(N-e)$. With
$\varepsilon^2=1$ this yields $\varepsilon^{c+D'+Q'}E^{c+D'}=\varepsilon^{Q'}\det(Z)^{(r-1)(N-e)+(r-1)Q'}
=E^{Q'}\det(Z)^{D'}\det(Z)^Q$. Since $\mathcal O[Z]$ is a domain and $\det(Z)^{D'}E^{Q'}\neq0$,
$S(f)=\det(Z)^Qg$.

\emph{Step 5: part (2).} Let $l\in GL_r(K)$ and $f$ in the module over $K$ of $\alpha^*$; then $f_l$
lies in it too, by Lemma~\ref{pd-gain-holds-hecke-rows}(2). Since
${}^t\operatorname{adj}({}^tlZ)={}^t\operatorname{adj}({}^tl)\,{}^t\operatorname{adj}(Z)=
\operatorname{adj}(l)\,{}^t\operatorname{adj}(Z)$, we have
$l\,{}^t\operatorname{adj}({}^tlZ)P_0=\det(l)\,{}^t\operatorname{adj}(Z)P_0$. So
$S(f_l)({}^tlZ)$ is $f$ with $Z_{ij}$ replaced by $\det(l)\operatorname{adj}(Z)_{\bar ji}$, and as $f$ is
homogeneous of total degree $D$, $S(f_l)({}^tlZ)=\det(l)^DS(f)$. By Step 2,
$S(f)=\det(Z)^Q\tau(f)$ and $S(f_l)=\det(Z)^Q\tau(f_l)$, so
\[
 \det(l)^Q\det(Z)^Q\,\tau(f_l)({}^tlZ)=\det(l)^{D}\det(Z)^Q\tau(f)
 =\det(l)^{Q+N-e}\det(Z)^Q\tau(f),
\]
and cancelling $\det(l)^Q\det(Z)^Q$ gives (2).
\end{proof}

\begin{lemma}\label{pd-kappa-gram-fock}
Let $A$ be a commutative ring, $r\geq1$, $Z$ an $r\times r$ matrix of variables, and $[\ ,\ ]$ the Fock
pairing on $A[Z]$ (Definition~\ref{pd-weight-lattices}(2)). Let $\nabla_A(\nu)$, $\Delta_A(\nu)$,
$h_\nu$ and the operators $E^{(k)}_{ab}$ be as in Definition~\ref{pd-weight-lattices}(1), with $Z$ in
place of $X$, and let $\nu_1\geq\cdots\geq\nu_r\geq0$ be integers.
\begin{enumerate}
\item $[f(gZ),f']=[f,f'({}^tgZ)]$ for all $f,f'\in A[Z]$ and all $r\times r$ matrices $g$ over $A$.
\item $[E^{(k)}_{ab}f,f']=[f,E^{(k)}_{ba}f']$ for all $f,f'\in A[Z]$, $1\leq a,b\leq r$ and $k\geq0$.
\item There is an integer $N_\nu>0$ such that $[h_\nu,h_\nu]$ is the image of $N_\nu$ in $A$ for
  every $A$.
\item If $A$ is a field of characteristic $0$ and $F\in\nabla_A(\nu)$, then
  $[h_\nu,F]=F(1)\,[h_\nu,h_\nu]$, where $F(1)$ is the value of $F$ at the identity matrix.
\item Let $F\in A[Z]$ have degree $\nu_s$ in the entries of the $s$-th column for each $s$, and let
  $Z^e=\prod_{a,s}Z_{as}^{e_{as}}$ be a monomial with $e_{as}=0$ for $a<s$ and $\sum_ae_{as}=\nu_s$
  for each $s$. Let $w_e$ be the composite of the operators $E^{(e_{as})}_{sa}$, $a>s$, applying
  first those with $s=1$, then those with $s=2$, and so on. Then the coefficient of $Z^e$ in $F$ is
  $(w_eF)(1)$.
\item Let $A$ be a domain with fraction field $K$ and $Q\in\nabla_K(\nu)$. If the coefficient of $Q$
  at every monomial $Z^e$ with $e_{as}=0$ for $a<s$ lies in $A$, then $Q\in\nabla_A(\nu)$; if all
  these coefficients are $0$, then $Q=0$.
\end{enumerate}
\end{lemma}
\leanref{Formalism/LocalToGlobal/GroupLattices/FockModels/TransposeWords/FockR.lean}{Formalism.L2.G.Model.TS.fockR}
\leanref{Formalism/LocalToGlobal/GroupLattices/FockModels/TransposeWords/TransposeWordsProofs/FockPairingRowAction.lean}{Formalism.L2.G.Model.TS.fock_pairing_row_action}
\leanref{Formalism/LocalToGlobal/GroupLattices/FockModels/TransposeWords/TransposeWordsProofs/FockPairingDividedAdjoint.lean}{Formalism.L2.G.Model.TS.fock_pairing_divided_adjoint}
\leanref{Formalism/GradedWeylDivisibility/FockSpaceTableaux/FockQRowDivAdj.lean}{Formalism.GradedWeylDivisibility.fock_space_row_divided_adjoint}
\leanref{Formalism/LocalToGlobal/GroupLattices/FockModels/TransposeWords/TransposeWordsProofs/ProdBlockEmbMonomial.lean}{Formalism.L2.G.Model.TS.fockR_hw_norm}
\leanref{Formalism/LocalToGlobal/GroupLattices/FockModels/TransposeWords/TransposeWordsProofs/FockPairingHighestWeightProof.lean}{Formalism.L2.G.Model.TS.fock_pairing_highest_weight_proof}
\leanref{Formalism/GradedWeylDivisibility/FinalB.lean}{Formalism.GradedWeylDivisibility.hwLine}
\leanref{Formalism/LocalToGlobal/GroupLattices/FockModels/TransposeWords/TransposeWordsProofs/CwlApplyWordAppendSingle.lean}{Formalism.L2.G.Model.TS.coeffWord'}
\leanref{Formalism/LocalToGlobal/GroupLattices/FockModels/TransposeWords/TransposeWordsProofs/LowerTriangularOverField.lean}{Formalism.L2.G.Model.TS.lower_triangular_over_field}
\leanref{Formalism/GradedWeylDivisibility/FinalC.lean}{Formalism.GradedWeylDivisibility.leadInterlaced}
\leanref{Formalism/GradedWeylDivisibility/FinalC.lean}{Formalism.GradedWeylDivisibility.sSYTExists}
\leanref{Formalism/GradedWeylDivisibility/FinalC.lean}{Formalism.GradedWeylDivisibility.column_lex_bideterminant}

\begin{proof}
(1) Write $\partial_v$ for the derivative in the variable $Z_v$ and $m!$ for the product of the
factorials of the entries of an exponent $m$. If $m$ has entry $m_v$ at $v$, then the exponent
$m+\epsilon_v$ of $Z^mZ_v$ has $(m+\epsilon_v)!=(m_v+1)\,m!$, and $\partial_vZ^{m+\epsilon_v}=(m_v+1)Z^m$.
Comparing coefficients gives $[fZ_v,h]=[f,\partial_vh]$ for all $f,h$. The chain rule gives
$\partial_{(b,s)}\big(h(gZ)\big)=\sum_ag_{ab}\,(\partial_{(a,s)}h)(gZ)$. We prove (1) by induction on
$f$. For a constant $f$ both sides are $f$ times the constant term of $f'$. If (1) holds for $f$, then
$(fZ_{as})(gZ)=\sum_bg_{ab}f(gZ)Z_{bs}$, and
\[
 [(fZ_{as})(gZ),f']=\sum_bg_{ab}[f(gZ),\partial_{(b,s)}f']=\sum_bg_{ab}[f,(\partial_{(b,s)}f')({}^tgZ)]
 =[f,\partial_{(a,s)}(f'({}^tgZ))]=[fZ_{as},f'({}^tgZ)] .
\]
Both sides of (1) are additive in $f$, so (1) holds for all $f$.

(2) Over $\mathbf C$, take $g=1+tE_{ba}$ with $t\in\mathbf C$ in (1). The substitution $Z\mapsto gZ$
replaces the $b$-th row by its sum with $t$ times the $a$-th row, so
$f(gZ)=\sum_kt^kE^{(k)}_{ab}f$; likewise ${}^tg=1+tE_{ab}$ and $f'({}^tgZ)=\sum_kt^kE^{(k)}_{ba}f'$. Both
sides of (1) are then polynomials in $t$, and comparing the coefficients of $t^k$ gives (2) over
$\mathbf C$, hence for polynomials with rational coefficients. Both sides of (2) are bilinear, and
$E^{(k)}_{ab}$ and the Fock pairing carry polynomials with integer coefficients to such and commute
with ring homomorphisms applied to the coefficients. So (2) for monomials with integer coefficients
gives (2) over every $A$.

(3) The polynomial $h_\nu$ has integer coefficients and $h_\nu(1)=1$, since every leading principal
minor is $1$ at the identity matrix. Over $\mathbf Z$, $[h_\nu,h_\nu]=\sum_m(h_\nu)_m^2\,m!$ is a sum of
nonnegative terms, not all $0$; call it $N_\nu$. The Fock pairing commutes with ring homomorphisms
applied to the coefficients, so $[h_\nu,h_\nu]=N_\nu$ in every $A$.

(4) Let $G$ be the sum of the terms of $F$ of degree $\nu_a$ in the $a$-th row for every $a$. The
substitutions $Z\mapsto Z(1+yE_{lj})$ do not change the degree in each row, so $G\in\nabla_A(\nu)$.
By Lemma~\ref{pd-gain-holds-hecke-rows}(1), every monomial of an element of $\nabla_A(\nu)$ has degree
at most $\nu_1+\cdots+\nu_a$ in the first $a$ rows. Let $a<b$ and $k\geq1$. The polynomial
$E^{(k)}_{ab}G$ lies in $\nabla_A(\nu)$, since the operators $E^{(k)}_{ab}$ preserve it (as in the proof
of Lemma~\ref{pd-gain-weyl-comparison}), and it moves $k$ units of degree from the $b$-th row to the
$a$-th, so its monomials have degree $\nu_1+\cdots+\nu_a+k$ in the first $a$ rows. Hence
$E^{(k)}_{ab}G=0$; in particular $G$ is killed by the raising operators $E^{(1)}_{ab}$, $a<b$. The
representation $\nabla_A(\nu)$ is irreducible with highest weight vector $h_\nu$
\cite[Lecture~15]{FH91}, and a vector killed by the raising operators is a multiple of it; so
$G=ch_\nu$ for some $c\in A$. The monomials of $h_\nu$ have degree $\nu_a$ in
the $a$-th row, and the Fock pairing pairs only equal monomials, so $[h_\nu,F]=[h_\nu,G]=c[h_\nu,h_\nu]$.
A monomial of degree $\nu_s$ in each column is nonzero at the identity matrix only if it is
$m_\nu=\prod_sZ_{ss}^{\nu_s}$, which has degree $\nu_a$ in the $a$-th row. So $F(1)$ is the coefficient
of $m_\nu$ in $F$, which is its coefficient in $G=ch_\nu$, that is, $c\,h_\nu(1)=c$.

(5) Let $y_{as}$, $a>s$, be variables, and let $L_s$ be the matrix $1+\sum_{a>s}y_{as}E_{as}$. The
substitution $Z\mapsto L_sZ$ adds $y_{as}$ times the $s$-th row to the $a$-th row for each $a>s$, and
these operations commute for fixed $s$. So $w_eF$ is the coefficient of
$y^e=\prod_{a>s}y_{as}^{e_{as}}$ in $F(L_1L_2\cdots L_rZ)$. The product $L=L_1\cdots L_r$ is the lower
unitriangular matrix with entry $y_{as}$ at $(a,s)$ for $a>s$. A monomial $Z^m$ takes the value
$\prod_{a>s}y_{as}^{m_{as}}$ at $L$ if $m_{as}=0$ for all $a<s$, and $0$ otherwise. It contributes to the
coefficient of $y^e$ only if moreover $m_{as}=e_{as}$ for $a>s$. If $Z^m$ occurs in $F$, its column
sums are $\nu_s=\sum_ae_{as}$, so then also $m_{ss}=e_{ss}$ and $m=e$. Hence $(w_eF)(1)$ is the
coefficient of $Z^e$ in $F$.

(6) Order the monomials by their exponents at the entries $(1,1),(2,1),\dots,(r,1),(1,2),\dots$ in
turn: a monomial is larger if it has the larger exponent at the first entry where the two differ.
The leading monomial of a nonzero polynomial is its largest monomial with nonzero coefficient. In
this order the leading monomial of a bideterminant whose minors have increasing rows is the product
of the diagonal terms of its minors, with coefficient $1$, and the standard monomial theory of
bideterminants \cite{DEP80} gives: (i) the leading monomial $Z^m$ of a nonzero element of
$\nabla_K(\nu)$ satisfies $\sum_{b\leq a}m_{b,s+1}\leq\sum_{b<a}m_{bs}$ for all $a$ and $s<r$, and
these inequalities force $m_{as}=0$ for $a<s$, by induction on $s$; (ii) a monomial with column sums
$\nu$ satisfying them is the leading monomial of a bideterminant $b$ of shape $\nu$ whose minors have
increasing rows. Such a $b$ has integer coefficients and lies in $\nabla_A(\nu)$, as in the proof of
Lemma~\ref{pd-kappa-cloud-cartan}(2).

Let $Q$ be as in the first assertion. We argue by induction on the leading monomial, the order being
a well-order. If $Q\neq0$, its leading monomial $Z^m$ has column sums $\nu$ and is lower triangular by
(i), so its coefficient $c$ lies in $A$. Let $b$ be as in (ii). Then $Q-cb\in\nabla_K(\nu)$ is $0$ or has
a smaller leading monomial, and its coefficients at the lower triangular monomials lie in $A$. By
induction $Q-cb\in\nabla_A(\nu)$, so $Q\in\nabla_A(\nu)$. For the second assertion, if $Q\neq0$ the
coefficient at its leading monomial, which is lower triangular, is not $0$.
\end{proof}

\begin{definition}\label{pd-kappa-defects-gl-modules}
Let $K$ be a field of characteristic $0$ and $n\geq1$ an integer, and let $\epsilon_a$ be the $a$-th
unit vector of $\mathbf Z^n$.
\begin{enumerate}
\item A \emph{weight module} for $\mathfrak{gl}_n$ over $K$ is a $K$-vector space $V$ with
  endomorphisms $E_{ab}$ ($1\leq a,b\leq n$) such that
  $E_{ab}E_{cd}-E_{cd}E_{ab}=\delta_{bc}E_{ad}-\delta_{da}E_{cb}$ for all $a,b,c,d$, together with a
  decomposition $V=\bigoplus_{\mu\in\mathbf Z^n}V_\mu$ such that $E_{cc}$ acts on $V_\mu$ as
  multiplication by $\mu_c$ and $E_{ab}V_\mu\subset V_{\mu+\epsilon_a-\epsilon_b}$. The elements of
  $V_\mu$ are the \emph{weight vectors} of weight $\mu$, and $x_\mu$ is the component of $x\in V$ in
  $V_\mu$.
\item A \emph{subrepresentation} of $V$ is a subspace stable under every $E_{ab}$, and
  $\langle S\rangle$ is the smallest subrepresentation containing a set $S\subset V$. A \emph{weight}
  of a subspace $U$ is a $\mu$ with $U\cap V_\mu\neq0$. A \emph{highest weight vector} of weight
  $\lambda$ is an $x\in V_\lambda$ with $E_{ab}x=0$ for all $a<b$. For $\mu,\lambda\in\mathbf Z^n$,
  $\mu\preceq\lambda$ means that $\mu_1+\cdots+\mu_k\leq\lambda_1+\cdots+\lambda_k$ for
  $1\leq k\leq n$, with equality for $k=n$.
\item Let $K=\mathbf Q_p$ for a prime $p$. For $a\neq b$ and $k\geq0$ put $E^{(k)}_{ab}=E_{ab}^k/k!$.
  The \emph{divided-power hull} $\mathcal D(S)$ of a set $S\subset V$ is the smallest
  $\mathbf Z_p$-submodule of $V$ that contains $S$ and is stable under the $E^{(k)}_{ab}$ with $a\neq b$
  and $k\geq0$.
\end{enumerate}
\end{definition}
\leanref{Formalism/MahlerMeasure/MahlerMeasureProofs/WeightSpaceDimensionMonotone.lean}{Formalism.GM.DimensionMonotonicity.GlW}
\leanref{Formalism/MahlerMeasure/MahlerMeasureProofs/WeightSpaceDimensionMonotone.lean}{Formalism.GM.DimensionMonotonicity.GlW.gS}
\leanref{Formalism/MahlerMeasure/MahlerMeasureProofs/WeightSpaceDimensionMonotone.lean}{Formalism.GM.DimensionMonotonicity.GlW.IsHW}
\leanref{Formalism/MahlerMeasure/MahlerMeasureProofs/WeightSpaceDimensionMonotone.lean}{Formalism.GM.DimensionMonotonicity.Dle}
\leanref{Formalism/ThetaLogCoordinates/TraceIsometryWeights/DivStable.lean}{Formalism.InvariantFactorPaths.PW.DivStable}
\leanref{Formalism/ThetaLogCoordinates/TraceIsometryWeights/DivStable.lean}{Formalism.InvariantFactorPaths.PW.ZpClosed}
\leanref{Formalism/ThetaLogCoordinates/TraceIsometryWeights/HasLatBasis.lean}{Formalism.InvariantFactorPaths.PW.DeltaLat}

\begin{lemma}\label{pd-kappa-defects-gl-highest}
Let $V$ be a weight module for $\mathfrak{gl}_n$ over a field $K$ of characteristic $0$
(Definition~\ref{pd-kappa-defects-gl-modules}).
\begin{enumerate}
\item If $U\subset V$ is a subrepresentation, then $x_\mu\in U$ for every $x\in U$ and every $\mu$. A
  finite-dimensional subrepresentation has finitely many weights.
\item Let $h\in V$ be a highest weight vector of weight $\lambda$. Every weight $\mu$ of $\langle
  h\rangle$ satisfies $\mu\preceq\lambda$, and $x_\lambda\in Kh$ for every $x\in\langle h\rangle$.
\item Let $W\subset V$ be a finite-dimensional subrepresentation and $(h_i)_{i\in I}$ nonzero highest
  weight vectors in $W$ such that every highest weight vector in $W$ lies in the span of the $h_i$.
  Then $W\subset\sum_{i\in I}\langle h_i\rangle$.
\item Let $W\subset V$ be a finite-dimensional subrepresentation and $h,h'\in W$ highest weight
  vectors of the same weight with $h\neq0$. Then there is a $K$-linear map
  $\theta\colon\langle h\rangle\to V$ that commutes with every $E_{ab}$ and maps $\langle h\rangle$
  onto $\langle h'\rangle$.
\end{enumerate}
\end{lemma}
\leanref{Formalism/MahlerMeasure/MahlerMeasureProofs/WeightSpaceDimensionMonotone.lean}{Formalism.GM.DimensionMonotonicity.GlW.graded_of}
\leanref{Formalism/MahlerMeasure/MahlerMeasureProofs/WeightSpaceDimensionMonotone.lean}{Formalism.GM.DimensionMonotonicity.GlW.wt_le_of_mem}
\leanref{Formalism/MahlerMeasure/MahlerMeasureProofs/WeightSpaceDimensionMonotone.lean}{Formalism.GM.DimensionMonotonicity.GlW.comp_top_of_mem}
\leanref{Formalism/PolynomialDimension/MultinomialDepth/MonoCount.lean}{Formalism.CountingSumsFrames.GlX.span_hw}
\leanref{Formalism/PolynomialDimension/MultinomialDepth/MonoCount.lean}{Formalism.CountingSumsFrames.GlX.ker_le}
\leanref{Formalism/ThetaLogCoordinates/TraceIsometryWeights/SchurTraceProofs/WeylInt.lean}{Formalism.InvariantFactorPaths.PW.wl_transport}

\begin{proof}
(1) Let $x\in U$, and let $\mu^1,\dots,\mu^s$ be the weights $\mu$ with $x_\mu\neq0$. For each $i$,
the product over $j\neq i$ of the operators $(E_{cc}-\mu^j_c)/(\mu^i_c-\mu^j_c)$, with $c$ chosen so
that $\mu^i_c\neq\mu^j_c$, maps $x$ to $x_{\mu^i}$; it is a polynomial in the $E_{cc}$, so
$x_{\mu^i}\in U$. Weight vectors of distinct weights are linearly independent, so a
finite-dimensional $U$ has finitely many weights.

Parts (2) to (4) are standard facts of highest weight theory \cite[\S\S20--21]{Hum72}, applied to
$\mathfrak{sl}_n$, since the $E_{aa}$ act on every $V_\mu$ by integers. For (2), $\langle h\rangle$ is
spanned by the vectors obtained from $h$ by applying the $E_{ab}$ with $a>b$ \cite[\S20.2]{Hum72}, and
such an $E_{ab}$ maps $V_\mu$ to $V_{\mu+\epsilon_a-\epsilon_b}$ with
$\mu+\epsilon_a-\epsilon_b\preceq\mu$. For (3), $W$ is the direct sum of irreducible
subrepresentations by Weyl's theorem \cite[\S6.3]{Hum72}, and each of them is generated by a highest
weight vector of $W$, which lies in the span of the $h_i$ and hence in $\sum_i\langle h_i\rangle$. For
(4), $\langle h\rangle$ is irreducible, and $\langle h'\rangle$ is either $0$ or irreducible with the same
highest weight, so the two are isomorphic or $\langle h'\rangle=0$ \cite[\S20.3]{Hum72}.
\end{proof}

\begin{definition}\label{pd-kappa-defects-projector}
Let $K$ be a field of characteristic $0$, $W$ a $K$-vector space, $\iota$ a set and $E_{ab}$
($a,b\in\iota$) endomorphisms of $W$ such that
$E_{ab}E_{cd}-E_{cd}E_{ab}=\delta_{bc}E_{ad}-\delta_{da}E_{cb}$ for all $a,b,c,d\in\iota$. Let $l\geq0$
be an integer.
\begin{enumerate}
\item For a finite set $S\subset\iota$ and $z\in\iota$ put $\Omega_{S,z,0}=1$ and
  $\Omega_{S,z,q+1}=\sum_{i\in S}E_{zi}\,\Omega_{S,z,q}\,E_{iz}$ for $q\geq0$, and
  \[
   P^l_{S,z}=\sum_{q=0}^lc_{|S|}(q)\,\Omega_{S,z,q},\qquad c_m(q)=\frac{(-1)^q\,m!}{(q+m)!\,q!}.
  \]
\item For a list $L=(z_1,\dots,z_k)$ of elements of $\iota$ put $\Pi^l_L=1$ if $k=0$, and
  $\Pi^l_L=P^l_{\{z_2,\dots,z_k\},z_1}\Pi^l_{(z_2,\dots,z_k)}$ if $k\geq1$.
\item If $E'_{ab}$ ($a,b\in\iota$) are endomorphisms of a second $K$-vector space $W'$ with the same
  relations, $\operatorname{ad}E_{ab}$ is the endomorphism $\phi\mapsto E'_{ab}\phi-\phi E_{ab}$ of
  $\operatorname{Hom}_K(W,W')$.
\end{enumerate}
\end{definition}
\leanref{Formalism/ThetaLogCoordinates/TraceIsometryWeights/DivStable.lean}{Formalism.InvariantFactorPaths.PW.GlRel}
\leanref{Formalism/ThetaLogCoordinates/TraceIsometryWeights/DivStable.lean}{Formalism.InvariantFactorPaths.PW.Om}
\leanref{Formalism/ThetaLogCoordinates/TraceIsometryWeights/DivStable.lean}{Formalism.InvariantFactorPaths.PW.Pstep}
\leanref{Formalism/ThetaLogCoordinates/TraceIsometryWeights/DivStable.lean}{Formalism.InvariantFactorPaths.PW.InvOp}
\leanref{Formalism/ThetaLogCoordinates/TraceIsometryWeights/DivStable.lean}{Formalism.InvariantFactorPaths.PW.homAct}

\begin{lemma}\label{pd-kappa-defects-projector-invariant}
Let $K$, $W$, $\iota$, the $E_{ab}$ and $l$ be as in Definition~\ref{pd-kappa-defects-projector}, and
let $L$ be a list of distinct elements of $\iota$.
\begin{enumerate}
\item Let $E'_{ab}$ be endomorphisms of a $K$-vector space $W'$ with the same relations. Then the
  $\operatorname{ad}E_{ab}$ on $\operatorname{Hom}_K(W,W')$ satisfy these relations too. If
  $\theta\colon W\to W'$ is $K$-linear with $\theta E_{ab}=E'_{ab}\theta$ for all $a,b$, then
  $\theta\Pi^l_L=\Pi'^l_L\theta$, where $\Pi'^l_L$ is formed with the $E'_{ab}$.
\item Let $N\subset W$ be a subspace stable under every $E_{ab}$, such that $x=0$ whenever $x\in N$,
  $z\in\iota$ and $E_{zz}x=kx$ for an integer $k$ with $|k|>l$. Let $u\in N$ with $E_{aa}u=0$ for every
  entry $a$ of $L$. Then $\Pi^l_Lu\in N$, $E_{ab}\Pi^l_Lu=0$ for all entries $a,b$ of $L$, and
  $\Pi^l_Lu-u$ lies in the sum of the images of the $E_{ab}$ with $a\neq b$.
\end{enumerate}
\end{lemma}
\leanref{Formalism/ThetaLogCoordinates/TraceIsometryWeights/StatementA.lean}{Formalism.InvariantFactorPaths.PW.InvInv}
\leanref{Formalism/ThetaLogCoordinates/TraceIsometryWeights/StatementA.lean}{Formalism.InvariantFactorPaths.PW.InvNat}
\leanref{Formalism/ThetaLogCoordinates/TraceIsometryWeights/StatementA.lean}{Formalism.InvariantFactorPaths.PW.HomRel}
\leanref{Formalism/ThetaLogCoordinates/TraceIsometryWeights/SchurTraceProofs/InvInv.lean}{Formalism.InvariantFactorPaths.PW.invariant_inverse_holds}
\leanref{Formalism/ThetaLogCoordinates/TraceIsometryWeights/SchurTraceProofs/InvInv.lean}{Formalism.InvariantFactorPaths.PW.invariant_natural_holds}
\leanref{Formalism/ThetaLogCoordinates/TraceIsometryWeights/SchurTraceProofs/HomRelationHolds.lean}{Formalism.InvariantFactorPaths.PW.hom_relation_holds}
\leanref{Formalism/ThetaLogCoordinates/TraceIsometryWeights/SchurTraceProofs/InvInv.lean}{Formalism.InvariantFactorPaths.PW.Om_E}

\begin{proof}
(1) For the first assertion, expand $\operatorname{ad}E_{ab}\operatorname{ad}E_{cd}\phi
-\operatorname{ad}E_{cd}\operatorname{ad}E_{ab}\phi$; the terms $E'_{ab}\phi E_{cd}$ and
$E'_{cd}\phi E_{ab}$ cancel, and the relations for the $E'$ and for the $E$ give the rest. The second
holds because $\Pi^l_L$ is a polynomial in the $E_{ab}$ and $\Pi'^l_L$ is the same polynomial in the
$E'_{ab}$.

(2) We first treat one factor. Let $S\subset\iota$ be finite, $z\in\iota\smallsetminus S$, and $u\in N$
with $E_{ab}u=0$ for all $a,b\in S$ and $E_{zz}u=0$. Write $\Omega_q=\Omega_{S,z,q}$ and $c(q)=c_{|S|}(q)$,
and put $y=P^l_{S,z}u$. We show that $y\in N$, that $E_{ab}y=0$ for all $a,b\in S\cup\{z\}$, and that
$y-u$ lies in the sum of the images of the $E_{zi}$, $i\in S$.

(a) For $a,b\in S$, for $a=b=z$ and for $a=b\notin S\cup\{z\}$, the relations give
\[
 E_{ab}\Omega_{q+1}-\Omega_{q+1}E_{ab}=\sum_{i\in S}E_{zi}(E_{ab}\Omega_q-\Omega_qE_{ab})E_{iz}:
\]
the remaining terms are $\sum_iE_{zi}\Omega_q[E_{ab},E_{iz}]+\sum_i[E_{ab},E_{zi}]\Omega_qE_{iz}$, which
is $E_{zb}\Omega_qE_{az}-E_{zb}\Omega_qE_{az}$ in the first case, $-\Omega_{q+1}+\Omega_{q+1}$ in the
second, and $0$ in the third. By induction on $q$, these $E_{ab}$ commute with every $\Omega_q$. Hence
$E_{ab}y=0$ for $a,b\in S$, and $E_{zz}y=0$.

(b) For $j\in S$ put $C_j=\sum_{i\in S}E_{iz}E_{ji}-E_{jz}E_{zz}+(|S|+1)E_{jz}$. A computation with the
relations, by induction on $q$, gives
\[
 E_{jz}\Omega_{q+1}=\Omega_{q+1}E_{jz}+(q+1)\,\Omega_q\,(C_j+qE_{jz}).
\]
Since $E_{ji}u=0$ for $i\in S$ and $E_{zz}u=0$, $C_ju=(|S|+1)E_{jz}u$, so
$E_{jz}\Omega_{q+1}u=\Omega_{q+1}E_{jz}u+(q+1)(q+1+|S|)\,\Omega_qE_{jz}u$. As
$c(q)+c(q+1)(q+1)(q+1+|S|)=0$, induction on $l'$ gives
$E_{jz}\sum_{q\leq l'}c(q)\Omega_qu=c(l')\,\Omega_{l'}E_{jz}u$, and for $l'=l$,
$E_{jz}y=c(l)\,\Omega_lE_{jz}u$.

(c) Let $x\in N$ with $E_{zz}x=-kx$ for an integer $k\geq0$. For $i\in S$,
$E_{zz}E_{iz}x=-(k+1)E_{iz}x$, since $E_{zz}E_{iz}-E_{iz}E_{zz}=-E_{iz}$. As
$\Omega_{q+1}x=\sum_{i\in S}E_{zi}\Omega_q(E_{iz}x)$ and $N$ is stable, induction on $q$ shows that
$\Omega_qx=0$ whenever $k+q>l$: for $q=0$ the hypothesis on $N$ gives $x=0$. For $j\in S$ the vector
$x=E_{jz}u$ satisfies $E_{zz}x=-x$, so $\Omega_lE_{jz}u=0$, and $E_{jz}y=0$ by (b).

(d) Let $b\in S$. Then $y\in N$, and $E_{bz}y=0$ by (c), $E_{bb}y=0$ and $E_{zz}y=0$ by (a). Put
$y_k=E_{zb}^ky\in N$. The relations give $E_{zz}y_k=ky_k$, $E_{bb}y_k=-ky_k$ and
$E_{bz}y_{k+1}=-(k+1)k\,y_k$. By the hypothesis on $N$, $y_k=0$ for $k>l$; and if $y_{k+1}=0$ with
$k\geq1$, then $y_k=0$. Descending induction from $k=l+1$ gives $E_{zb}y=y_1=0$. With (a) and (c) this
gives $E_{ab}y=0$ for all $a,b\in S\cup\{z\}$.

(e) The space $N$ is stable, so $y\in N$. Finally $y-u=\sum_{q=1}^lc(q)\Omega_qu$, as $c(0)=1$, and for
$q\geq1$ every term of $\Omega_q$ begins with some $E_{zi}$, $i\in S$.

Now we prove (2) by induction on the length of $L$. For $L$ empty, $\Pi^l_Lu=u$. Let $L=(z,L')$, let
$S$ be the set of entries of $L'$, and put $u'=\Pi^l_{L'}u$. By induction $u'\in N$, $E_{ab}u'=0$ for
$a,b\in S$, and $u'-u$ lies in the sum of the images. Since $z\notin S$, $E_{zz}$ commutes with every
factor $P^l_{S',z'}$ of $\Pi^l_{L'}$ by (a), as $z\notin S'\cup\{z'\}$; so $E_{zz}u'=\Pi^l_{L'}E_{zz}u=0$.
The case of one factor, applied to $u'$, gives the assertions for $\Pi^l_Lu=P^l_{S,z}u'$, and
$\Pi^l_Lu-u=(P^l_{S,z}u'-u')+(u'-u)$.
\end{proof}
% Lean: (1) is hom_relation_holds and invariant_natural_holds (Om_nat, Pstep_nat). (2) is
% invariant_inverse_holds; the one-factor case is Pstep_inv: (a) Om_comm_S, diag_comm_omega and
% Om_comm_out (used through InvOp_comm_out for E_zz u' = 0), (b) Om_E with Cj, E_Om_apply, cc_rec and
% E_Pstep_apply, (c) Om_vanish, (d) lower_kill, (e) Pstep_mem and Pstep_sub_mem.

\begin{lemma}\label{pd-kappa-defects-projector-integral}
Let $p$ be a prime and let $W$, $\iota$, the $E_{ab}$ and $l$ be as in
Definition~\ref{pd-kappa-defects-projector}, with $K=\mathbf Q_p$. Let $M\subset W$ be a
$\mathbf Z_p$-submodule with $E_{ab}^k(M)\subset k!\,M$ for all $a\neq b$ in $\iota$ and all $k\geq0$.
Let $L$ be a list of $k$ distinct elements of $\iota$, and put
$w_L=\binom k2\lfloor\log_p(l+k-1)\rfloor$, and $w_L=0$ if $k\leq1$. Then
$p^{w_L}\Pi^l_L(M)\subset M$.
\end{lemma}
\leanref{Formalism/ThetaLogCoordinates/TraceIsometryWeights/StatementA.lean}{Formalism.InvariantFactorPaths.PW.InvInt}
\leanref{Formalism/ThetaLogCoordinates/TraceIsometryWeights/SchurTraceProofs/InvInt.lean}{Formalism.InvariantFactorPaths.PW.invariant_integral_holds}
\leanref{Formalism/ThetaLogCoordinates/TraceIsometryWeights/SchurTraceProofs/InvInt.lean}{Formalism.InvariantFactorPaths.PW.Om_insert}
\leanref{Formalism/ThetaLogCoordinates/TraceIsometryWeights/SchurTraceProofs/InvInt.lean}{Formalism.InvariantFactorPaths.PW.Om_int}
\leanref{Formalism/ThetaLogCoordinates/TraceIsometryWeights/SchurTraceProofs/InvInt.lean}{Formalism.InvariantFactorPaths.PW.Pstep_int}

\begin{proof}
Write $E^{(k)}_{ab}=E_{ab}^k/k!$, and put $\lambda_N=\lfloor\log_pN\rfloor$ for $N\geq1$ and
$\lambda_0=0$. For $0\leq j\leq m$, $v_p\binom mj\leq\lambda_m$, by Legendre's formula as in the proof
of Lemma~\ref{pd-gain-weyl-column-projector}(2); so $p^{\lambda_N}/\binom mj\in\mathbf Z_p$ when
$m\leq N$.

\emph{Step 1.} Let $S\subset\iota$ be finite, $i\in\iota\smallsetminus S$ and $z\notin S\cup\{i\}$. The
$E_{zj}$ with $j\neq z$ commute with each other, and so do the $E_{jz}$, by the relations. So the maps
$T\mapsto E_{zi}TE_{iz}$ and $T\mapsto\sum_{j\in S}E_{zj}TE_{jz}$ of $\operatorname{End}(W)$ commute, and
$\Omega_{S,z,q}$ is the $q$-th power of the second applied to $1$. The binomial theorem gives
\[
 \Omega_{S\cup\{i\},z,q}=\sum_{k=0}^q\binom qk\,E_{zi}^k\,\Omega_{S,z,q-k}\,E_{iz}^k .
\]

\emph{Step 2.} We show: if $z\notin S$, $c,q\geq0$ are integers, $N\geq q+c$ and $u\in M$, then
$p^{|S|\lambda_N}\frac{c!}{(q+c)!\,q!}\Omega_{S,z,q}u\in M$. For $S=\emptyset$, $\Omega_{S,z,0}=1$ and
$\Omega_{S,z,q}=0$ for $q\geq1$. For $S\cup\{i\}$ with $i\notin S$ and $z\neq i$, Step 1 writes the
element as
\[
 \sum_{k=0}^q\frac{p^{\lambda_N}}{\binom{c+k}k}\,E^{(k)}_{zi}\Big(p^{|S|\lambda_N}
 \frac{(c+k)!}{(q+c)!\,(q-k)!}\,\Omega_{S,z,q-k}\,E^{(k)}_{iz}u\Big).
\]
Here $E^{(k)}_{iz}u\in M$, the bracket lies in $M$ by induction applied with $c+k$ and $q-k$, its image
under $E^{(k)}_{zi}$ lies in $M$, and $p^{\lambda_N}/\binom{c+k}k\in\mathbf Z_p$ as $c+k\leq N$.

\emph{Step 3.} For $z\notin S$ and $u\in M$, Step 2 with $c=|S|$, $N=l+|S|$ and $q\leq l$ gives
$p^{|S|\lambda_{l+|S|}}P^l_{S,z}u\in M$.

\emph{Step 4.} We show $p^{\binom k2\lambda_N}\Pi^l_Lu\in M$ for $u\in M$ and $N\geq l+k-1$, by induction
on $k$; for $N=l+k-1$ this is the lemma. For $k=0$, $\Pi^l_Lu=u$. Let $L=(z,L')$ with $L'$ of length
$k-1$ and set of entries $S$. By induction $u'=p^{\binom{k-1}2\lambda_N}\Pi^l_{L'}u\in M$, and Step 3
gives $p^{(k-1)\lambda_{l+k-1}}P^l_{S,z}u'\in M$. As $\lambda_{l+k-1}\leq\lambda_N$ and
$\binom{k-1}2+(k-1)=\binom k2$, $p^{\binom k2\lambda_N}\Pi^l_Lu\in M$.
\end{proof}

\begin{lemma}\label{pd-kappa-defects-schur}
Let $V$ be a weight module for $\mathfrak{gl}_n$ over a field $K$ of characteristic $0$, $l\geq0$ an
integer, and $W\subset V$ a finite-dimensional subrepresentation with $|\mu'_c-\mu_c|\leq l$ for all
weights $\mu,\mu'$ of $W$ and all $c$ (Definition~\ref{pd-kappa-defects-gl-modules}). Let $h\in W$ be a
nonzero highest weight vector, $U=\langle h\rangle$, and $\phi$ an endomorphism of $U$ that commutes
with every $E_{cc}$. Let $\Pi^l_{(1,\dots,n)}$ be formed with the operators $\operatorname{ad}E_{ab}$ on
$\operatorname{End}_K(U)$ (Definition~\ref{pd-kappa-defects-projector}). Then
$\Pi^l_{(1,\dots,n)}\phi=(\operatorname{tr}\phi/\dim U)\cdot1_U$.
\end{lemma}
\leanref{Formalism/ThetaLogCoordinates/TraceIsometryWeights/StatementB.lean}{Formalism.InvariantFactorPaths.PW.SchurTrace}
\leanref{Formalism/ThetaLogCoordinates/TraceIsometryWeights/StatementB.lean}{Formalism.InvariantFactorPaths.PW.HomInv}
\leanref{Formalism/ThetaLogCoordinates/TraceIsometryWeights/SchurTraceProofs/SchurTraceOfHomInv.lean}{Formalism.InvariantFactorPaths.PW.schurTrace_of}
\leanref{Formalism/ThetaLogCoordinates/TraceIsometryWeights/SchurTraceProofs/SchurTraceOfHomInv.lean}{Formalism.InvariantFactorPaths.PW.homInv_of}
\leanref{Formalism/ThetaLogCoordinates/TraceIsometryWeights/SchurTraceProofs/SchurTraceOfHomInv.lean}{Formalism.InvariantFactorPaths.PW.hom_eig_zero}

\begin{proof}
Let $\lambda$ be the weight of $h$. The space $U\subset W$ is finite-dimensional, and it is the sum of the
$U\cap V_\mu$ by Lemma~\ref{pd-kappa-defects-gl-highest}(1). The $\operatorname{ad}E_{ab}$ satisfy the
relations by Lemma~\ref{pd-kappa-defects-projector-invariant}(1).

\emph{Step 1.} Let $\psi\in\operatorname{End}_K(U)$ and $z\in\{1,\dots,n\}$ with
$(\operatorname{ad}E_{zz})\psi=k\psi$ for an integer $k$ with $|k|>l$. For $y\in U\cap V_\mu$,
$E_{zz}\psi(y)=\psi(E_{zz}y)+k\psi(y)=(\mu_z+k)\psi(y)$, so every nonzero component of $\psi(y)$ has a
weight $\nu$ with $\nu_z=\mu_z+k$. If $y\neq0$, then $|\nu_z-\mu_z|\leq l<|k|$ for such a $\nu$, so
$\psi(y)=0$. Hence $\psi=0$. So Lemma~\ref{pd-kappa-defects-projector-invariant}(2) applies to the
operators $\operatorname{ad}E_{ab}$, with $N=\operatorname{End}_K(U)$, $L=(1,\dots,n)$ and $u=\phi$, as
$(\operatorname{ad}E_{cc})\phi=0$ for every $c$. Thus $\psi_0=\Pi^l_{(1,\dots,n)}\phi$ commutes with
every $E_{ab}$, and $\psi_0-\phi$ is a sum of endomorphisms $E_{ab}\chi-\chi E_{ab}$.

\emph{Step 2.} As $\psi_0$ commutes with the $E_{cc}$, $\psi_0(h)\in U\cap V_\lambda$, so
$\psi_0(h)=ch$ for some $c\in K$ by Lemma~\ref{pd-kappa-defects-gl-highest}(2). The kernel of
$\psi_0-c\cdot1_U$ is a subrepresentation containing $h$, so it is $U$, and $\psi_0=c\cdot1_U$.

\emph{Step 3.} The trace of $E_{ab}\chi-\chi E_{ab}$ is $0$, so
$c\dim U=\operatorname{tr}\psi_0=\operatorname{tr}\phi$. As $U\neq0$, $c=\operatorname{tr}\phi/\dim U$.
\end{proof}
% Lean: schurTrace_of = schurTrace_of_homInv (homInv_of invariant_inverse_holds). Step 1 is homInv_of
% with mapsInto U top and hom_eig_zero; Step 2 is isWt_of_eig, comp_top_of_mem and the gS_le argument
% (hle); Step 3 is hom_trace_holds (trace_mul_comm).

\begin{lemma}\label{pd-kappa-defects-weyl-integral}
Let $p$ be a prime, $V$ a weight module for $\mathfrak{gl}_n$ over $\mathbf Q_p$, $l\geq0$ an integer,
and $W\subset V$ a finite-dimensional subrepresentation with $|\mu'_c-\mu_c|\leq l$ for all weights
$\mu,\mu'$ of $W$ and all $c$. Let $M\subset V$ be an additive subgroup with $E^{(k)}_{ab}(M)\subset M$
for $a\neq b$ and $k\geq0$, and with $x_\mu\in M$ for every $x\in M$ and every $\mu$. Let
$\lambda\in\mathbf Z^n$, let $H$ be the set of highest weight vectors of weight $\lambda$ that lie in
$W\cap M$, and let $I$ be the subrepresentation generated by the highest weight vectors of weight
$\lambda$ in $W$. Put $w=\binom n2\lfloor\log_p(l+n-1)\rfloor$, and $w=0$ if $n=1$. Then
$p^w(I\cap M)\subset\mathcal D(H)$ (Definition~\ref{pd-kappa-defects-gl-modules}).
\end{lemma}
\leanref{Formalism/ThetaLogCoordinates/TraceIsometryWeights/StatementB.lean}{Formalism.InvariantFactorPaths.PW.WeylInt}
\leanref{Formalism/ThetaLogCoordinates/TraceIsometryWeights/HasLatBasis.lean}{Formalism.InvariantFactorPaths.PW.LatData}
\leanref{Formalism/ThetaLogCoordinates/TraceIsometryWeights/HasLatBasis.lean}{Formalism.InvariantFactorPaths.PW.IsoW}
\leanref{Formalism/ThetaLogCoordinates/TraceIsometryWeights/DischargeHeightBound.lean}{Formalism.InvariantFactorPaths.PW.Top.weylInt}
\leanref{Formalism/ThetaLogCoordinates/TraceIsometryWeights/SchurTraceProofs/WeylInt.lean}{Formalism.InvariantFactorPaths.PW.weylInt_of}
\leanref{Formalism/ThetaLogCoordinates/TraceIsometryWeights/SchurTraceProofs/WeylInt.lean}{Formalism.InvariantFactorPaths.PW.wl_key}

\begin{proof}
Let $x\in I\cap M$. If $W$ contains no nonzero highest weight vector of weight $\lambda$, then $I=0$
and $x=0$. Otherwise fix one, $h$, and put $U=\langle h\rangle\subset W$ and $d=\dim U\geq1$. For
$z\in I$ let $\tau(z)=z_\lambda$; so $\tau\in\operatorname{Hom}(I,V)$, where $I$ carries the
restrictions of the $E_{ab}$. Let $\Pi=\Pi^l_{(1,\dots,n)}$ be formed with the operators
$\operatorname{ad}E_{ab}$ on $\operatorname{Hom}(I,V)$ (Definition~\ref{pd-kappa-defects-projector}).

\emph{Step 1: weights.} The subrepresentation $I$ lies in $W$ and in the sum of the $\langle v\rangle$
over the highest weight vectors $v$ of weight $\lambda$ in $W$. By
Lemma~\ref{pd-kappa-defects-gl-highest}(2), every weight $\mu$ of $I$ satisfies $\mu\preceq\lambda$.
Let $y\in I$ and $a<b$. Then $E_{ab}(y_\lambda)=(E_{ab}y)_{\lambda+\epsilon_a-\epsilon_b}$, and
$\lambda+\epsilon_a-\epsilon_b\not\preceq\lambda$, so $E_{ab}(y_\lambda)=0$: $y_\lambda$ is a highest
weight vector of weight $\lambda$. It lies in $W$ by Lemma~\ref{pd-kappa-defects-gl-highest}(1).

\emph{Step 2: $\Pi\tau$ is $d^{-1}$ times the inclusion.} Let $v$ be a highest weight vector of weight $\lambda$ in $W$
and $z\in\langle v\rangle$. Lemma~\ref{pd-kappa-defects-gl-highest}(4) gives a linear map
$\theta\colon U\to V$ commuting with the $E_{ab}$, with $\theta(U)=\langle v\rangle\subset I$, and
$z=\theta(y)$ for some $y\in U$. Let $\tau_U\in\operatorname{End}(U)$ be $y\mapsto y_\lambda$. Since
$\theta$ commutes with the $E_{cc}$, it maps $U\cap V_\mu$ to $V_\mu$, so $\tau\circ\theta=\theta\circ\tau_U$.
Composition with $\theta$ on the right, $\operatorname{Hom}(I,V)\to\operatorname{Hom}(U,V)$, and on the left,
$\operatorname{End}(U)\to\operatorname{Hom}(U,V)$, commute with the operators $\operatorname{ad}E_{ab}$; by
Lemma~\ref{pd-kappa-defects-projector-invariant}(1),
$(\Pi\tau)\circ\theta=\Pi(\tau\circ\theta)=\Pi(\theta\circ\tau_U)=\theta\circ\Pi\tau_U$, with $\Pi$ formed
on each space. The map $\tau_U$ commutes with the $E_{cc}$, and by
Lemma~\ref{pd-kappa-defects-gl-highest}(2) it is the projection of $U$ onto $Kh$ along the other
weight spaces, so $\operatorname{tr}\tau_U=1$. Lemma~\ref{pd-kappa-defects-schur} gives
$\Pi\tau_U=d^{-1}\cdot1_U$, so $(\Pi\tau)(z)=\theta(d^{-1}y)=d^{-1}z$. As $I$ is spanned by such $z$,
$(\Pi\tau)(z)=d^{-1}z$ for every $z\in I$.

\emph{Step 3: integrality.} Let $\mathcal H$ be the set of $\phi\in\operatorname{Hom}(I,V)$ with
$\phi(I\cap M)\subset\mathcal D(H)$. It is a $\mathbf Z_p$-submodule, since $\mathcal D(H)$ is one. For
$a\neq b$ and $k\geq0$,
\[
 \frac1{k!}(\operatorname{ad}E_{ab})^k\phi=\sum_{i+j=k}(-1)^jE^{(i)}_{ab}\,\phi\,E^{(j)}_{ab},
\]
and $I\cap M$ is stable under $E^{(j)}_{ab}$ and $\mathcal D(H)$ under $E^{(i)}_{ab}$; so
$(\operatorname{ad}E_{ab})^k\mathcal H\subset k!\,\mathcal H$. The map $\tau$ lies in $\mathcal H$: for
$y\in I\cap M$, $y_\lambda$ is a highest weight vector of weight $\lambda$ in $W$ by Step 1, and lies in
$M$ by hypothesis, so $y_\lambda\in H$. The operators $\operatorname{ad}E_{ab}$ satisfy the relations by
Lemma~\ref{pd-kappa-defects-projector-invariant}(1), so
Lemma~\ref{pd-kappa-defects-projector-integral}, with the list $(1,\dots,n)$, gives
$p^w\Pi\tau\in\mathcal H$. By Step 2, $p^wd^{-1}x\in\mathcal D(H)$, and multiplying by the integer $d$
gives $p^wx\in\mathcal D(H)$.
\end{proof}
% Lean: weylInt_of. The case without a nonzero highest weight vector is hbot; Step 1 is wl_iso_le_iSup,
% wl_iso_wt and wl_comp_hw; Step 2 is hid, from wl_transport (GlX.ker_le, GlX.tr), wl_key
% (invariant_natural_holds with homAct_precomp and homAct_postcomp, wl_comp_comm, wl_trace and the
% SchurTrace hypothesis); Step 3 is hom_divided_holds (homAct_pow_apply), hom_relation_holds and
% invariant_integral_holds, applied to topC = tau, with deltaLat_closed and deltaLat_div. WeylInt
% assumes LatData p G W M l; weylInt_of uses its fields stable, fd, finW, bound, div and graded, and
% the book assumes only these.

\begin{lemma}\label{pd-kappa-defects-sandwich-core}
Let $p$ be a prime, $r\geq1$ and $l\geq0$ integers, and $\nu_1\geq\cdots\geq\nu_r\geq0$ integers,
and use the modules of Definition~\ref{pd-weight-lattices}(1) for an $r\times r$ matrix $X$ of
variables. Suppose that for all $f,g\in\nabla_{\mathbf Q_p}(\nu)$, every monomial of $f$, every
monomial of $g$ and every $1\leq a\leq r$, the degree of the first monomial in the $a$-th row of $X$
is at most $l$ plus that of the second. Put $w=\frac{r(r-1)}2\lfloor\log_p(l+r-1)\rfloor$, and
$w=0$ if $r=1$. Then $p^w\nabla_{\mathbf Z_p}(\nu)\subset\Delta_{\mathbf Z_p}(\nu)$.
\end{lemma}
\leanref{Formalism/ThetaLogCoordinates/SandwichModelDomains/SpreadSquare.lean}{Formalism.InvariantFactorPaths.INST.CoreQp}
\leanref{Formalism/ThetaLogCoordinates/SandwichModelDomains/SpreadSquare.lean}{Formalism.InvariantFactorPaths.INST.Spread}
\leanref{Formalism/ThetaLogCoordinates/SandwichModelDomains/StatementW.lean}{Formalism.InvariantFactorPaths.INST.sExp}
\leanref{Formalism/ThetaLogCoordinates/SandwichModelDomains/SandwichModelProofs/SwappedHighestWeightNeZero.lean}{Formalism.InvariantFactorPaths.INST.coreQp_of}

\begin{proof}
\emph{The weight module.} Let $V=\mathbf Q_p[X]$, and for $1\leq a,b\leq r$ let $E_{ab}$ be the
derivation $\sum_sX_{as}\,\partial/\partial X_{bs}$. These satisfy the relations of
Definition~\ref{pd-kappa-defects-gl-modules}(1). Let $V_\mu$ be the span of the monomials of degree
$\mu_a$ in the $a$-th row of $X$ for every $a$. Then $E_{cc}$ multiplies such a monomial by $\mu_c$, and
$E_{ab}$ moves one unit of degree from the $b$-th row to the $a$-th, so $V$ is a weight module for
$\mathfrak{gl}_r$ over $\mathbf Q_p$, and $x\in V_\mu$ if and only if every monomial of $x$ has degree
$\mu_a$ in the $a$-th row for every $a$. For $a\neq b$ and $k\geq0$, $E_{ab}^k/k!$ is the operator
$E^{(k)}_{ab}$ of Definition~\ref{pd-weight-lattices}(1), as in the proof of
Lemma~\ref{pd-gain-weyl-comparison}.

\emph{The subrepresentation.} Let $W=\nabla_{\mathbf Q_p}(\nu)$. It lies in the finite-dimensional
space of polynomials of degree $\nu_s$ in the $s$-th column for every $s$, and it is stable under the
$E_{ab}$, which act on the rows of $X$ and so commute with the substitutions $X\mapsto X(1+yE_{lj})$.
Let $\mu,\mu'$ be weights of $W$, with nonzero $f\in W\cap V_\mu$ and $g\in W\cap V_{\mu'}$, and let
$m$ and $m'$ be monomials of $f$ and of $g$. Their degrees in the $a$-th row are $\mu_a$ and $\mu'_a$,
and the hypothesis, applied to $(f,g,m,m')$ and to $(g,f,m',m)$, gives $|\mu'_a-\mu_a|\leq l$.

\emph{The lattice.} Let $M=\mathbf Z_p[X]$. The operators $E^{(k)}_{ab}$ with $a\neq b$ map $M$ into
itself, as they are defined by a substitution with integer coefficients, and the components $x_\mu$ of
an $x\in M$ lie in $M$, as they are sums of some of its terms.

\emph{The highest weight vectors.} The polynomial $h_\nu$ is nonzero and lies in $W$, every monomial of
$h_\nu$ has degree $\nu_a$ in the $a$-th row, and $E_{ab}h_\nu=0$ for $a<b$, by the proof of
Lemma~\ref{pd-gain-weyl-comparison}; so $h_\nu$ is a highest weight vector of weight $\nu$. Let $f\in W$
be a highest weight vector. Then $E^{(k)}_{ab}f=E_{ab}^kf/k!=0$ for $a<b$ and $k\geq1$, so
$f\in\mathbf Q_ph_\nu$ by the highest weight theory \cite[Lecture~15]{FH91}, as in the proof of
Lemma~\ref{pd-kappa-gram-fock}(4). By Lemma~\ref{pd-kappa-defects-gl-highest}(3), applied to the single
vector $h_\nu$, $W\subset\langle h_\nu\rangle$. So $W$ lies in the subrepresentation $I$ generated by the
highest weight vectors of weight $\nu$ in $W$.

\emph{Conclusion.} Let $x\in\nabla_{\mathbf Z_p}(\nu)$. The conditions defining the dual Weyl modules
are polynomial identities, so $x\in W\subset I$, and $x\in M$. Let $H$ be the set of highest weight
vectors of weight $\nu$ in $W\cap M$. Lemma~\ref{pd-kappa-defects-weyl-integral}, with $n=r$ and
$\lambda=\nu$, gives $p^wx\in\mathcal D(H)$, with $w=\binom r2\lfloor\log_p(l+r-1)\rfloor$ as in the
statement. Let $f\in H$. Then $f=ch_\nu$ with $c\in\mathbf Q_p$; at the identity matrix every leading
principal minor is $1$, so $c=f(1)$, which lies in $\mathbf Z_p$ because $f$ has coefficients in
$\mathbf Z_p$. Hence $f\in\Delta_{\mathbf Z_p}(\nu)$. By its definition, $\Delta_{\mathbf Z_p}(\nu)$ is a
$\mathbf Z_p$-submodule stable under the $E^{(k)}_{ab}$. So $\mathcal D(H)\subset\Delta_{\mathbf Z_p}(\nu)$,
and $p^wx\in\Delta_{\mathbf Z_p}(\nu)$.
\end{proof}
% Lean: coreQp_of rowGlOne rowGlDiv swpWt swpPol swpInt. The Lean works with the transpose swp
% (cor_emb), on polynomials in n x r variables, where GR p r nu = GAe acts on the second index; the
% weight module is (GR, gl_one_first_line_action, gl_first_line_divided_powers, swap_weight_compat).
% The subrepresentation and the spread are cor_latData (cor_stable, cor_W_le_pol, cor_bound, with
% polLatData_of for the lattice: intPol_zpClosed, divided_integral_holds, intPol_whc); the highest
% weight vectors are swapped_highest_weight_isHW (hwRowWeight, hwRaise), cor_line (hwLine) and
% cor_W_le_iso (GlX.span_hw); the conclusion is PW.Top.weylInt with dualWeylMap, cor_delta_le (cor_gen
% with cor_eval_hw, cor_L_div with weylRowDivStable) and cor_emb_inj.

\begin{lemma}\label{pd-kappa-defects-sandwich}
Let $p$ be a prime, $r\geq1$ and $l\geq0$ integers, and $\nu_1\geq\cdots\geq\nu_r\geq0$ integers
with $\nu_1\leq l$ or $\nu_1\leq\nu_r+l$, and put $w=\frac{r(r-1)}2\lfloor\log_p(l+r-1)\rfloor$, and
$w=0$ if $r=1$. Let $\mathcal O$ be a local domain in which $p$ is not a unit. Then
$p^w\nabla_{\mathcal O}(\nu)\subset\Delta_{\mathcal O}(\nu)\subset\nabla_{\mathcal O}(\nu)$
(Definition~\ref{pd-weight-lattices}(1), for an $r\times r$ matrix of variables).
\end{lemma}
\leanref{Formalism/ThetaLogCoordinates/SandwichModelDomains/StatementW.lean}{Formalism.InvariantFactorPaths.INST.SandwichLocal}
\leanref{Formalism/ThetaLogCoordinates/SandwichModelDomains/FinalW.lean}{Formalism.InvariantFactorPaths.INST.Top.sandwichLocal}
\leanref{Formalism/ThetaLogCoordinates/SandwichModelDomains/StatementW.lean}{Formalism.InvariantFactorPaths.INST.LBound}
\leanref{Formalism/ThetaLogCoordinates/SandwichModelDomains/SandwichModelProofs/Spread.lean}{Formalism.InvariantFactorPaths.INST.spread_top_holds}
\leanref{Formalism/ThetaLogCoordinates/SandwichModelDomains/SandwichModelProofs/Spread.lean}{Formalism.InvariantFactorPaths.INST.spread_square_holds}
\leanref{Formalism/ThetaLogCoordinates/SandwichModelDomains/SandwichModelProofs/Descent.lean}{Formalism.InvariantFactorPaths.INST.descent_over_integers}
\leanref{Formalism/ThetaLogCoordinates/SandwichModelDomains/SandwichModelProofs/Ascent.lean}{Formalism.InvariantFactorPaths.INST.ascent_domain_holds}
\leanref{Formalism/ThetaLogCoordinates/SandwichModelDomains/SandwichModelProofs/Ascent.lean}{Formalism.InvariantFactorPaths.INST.units_of_local_ring}

\begin{proof}
A bideterminant of shape $\nu$ is a product $\prod_k\det X_k$, where $X_k$ is the submatrix of $X$ on
some rows $i_{k,1}<\cdots<i_{k,s_k}$ and the first $s_k$ columns, and $\nu_t$ is the number of $k$
with $s_k\geq t$. It has integer coefficients and lies in $\nabla_{\mathbf Z}(\nu)$. By the standard
monomial theory \cite{DEP80}, for every domain $A$ these bideterminants span
$\nabla_A(\nu)$ over $A$, and for $A=\mathbf Z$ the standard ones among them form a finite basis.

\emph{Step 0.} For every commutative ring $A$, $\Delta_A(\nu)\subset\nabla_A(\nu)$: the operators
$E^{(k)}_{ab}$ preserve $\nabla_A(\nu)$ and $h_\nu\in\nabla_A(\nu)$, by the argument in the proof of
Lemma~\ref{pd-gain-weyl-comparison}, which uses nothing about the ring. This is the second inclusion.

\emph{Step 1: the spread.} Let $A$ be a domain. In a monomial of a bideterminant as above, the degree in
the $a$-th row is the number of $k$ such that $a$ is one of the rows of $X_k$. It is at most the number
of $k$ with $s_k\geq1$, which is $\nu_1$, and at least the number of $k$ with $s_k=r$, which is
$\nu_r$, since then $X_k$ uses every row. As these bideterminants span $\nabla_A(\nu)$, every monomial
of every element of $\nabla_A(\nu)$ has all its row degrees between $\nu_r$ and $\nu_1$. If
$\nu_1\leq l$, a row degree of one monomial is at most $\nu_1\leq l$, which is at most $l$ plus a row
degree of another; if $\nu_1\leq\nu_r+l$, it is at most $\nu_r+l$, which is again at most $l$ plus a
row degree of another. So the hypothesis of Lemma~\ref{pd-kappa-defects-sandwich-core} holds, with
$A=\mathbf Q_p$.

\emph{Step 2: over $\mathbf Z_p$.} By Lemma~\ref{pd-kappa-defects-sandwich-core},
$p^w\nabla_{\mathbf Z_p}(\nu)\subset\Delta_{\mathbf Z_p}(\nu)$.

\emph{Step 3: over $\mathbf Z$.} We show that for every $f\in\nabla_{\mathbf Z}(\nu)$ there is an
integer $c$ prime to $p$ with $cp^wf\in\Delta_{\mathbf Z}(\nu)$. Over $\mathbf Q$ the representation
$\nabla_{\mathbf Q}(\nu)$ is irreducible with highest weight vector $h_\nu$ \cite[Lecture~15]{FH91},
so it is spanned by the polynomials obtained from $h_\nu$ by the operators $E^{(k)}_{ab}$, which have
integer coefficients. So each element of the finite basis of $\nabla_{\mathbf Z}(\nu)$ has a nonzero
integer multiple in $\Delta_{\mathbf Z}(\nu)$, and the product $N$ of these multipliers satisfies
$N\nabla_{\mathbf Z}(\nu)\subset\Delta_{\mathbf Z}(\nu)$. Write $N=p^kc$ with $p\nmid c$. Let
$f\in\nabla_{\mathbf Z}(\nu)$. The conditions defining the dual Weyl modules are polynomial
identities with integer coefficients, so $f\in\nabla_{\mathbf Z_p}(\nu)$, and $p^wf\in
\Delta_{\mathbf Z_p}(\nu)$ by Step 2. The module $\Delta_{\mathbf Z_p}(\nu)$ is the
$\mathbf Z_p$-span of polynomials with integer coefficients; approximating the coefficients of a
$\mathbf Z_p$-combination of them by integers modulo $p^k$ gives $d\in\Delta_{\mathbf Z}(\nu)$ and a
polynomial $z$ over $\mathbf Z_p$ with $p^wf=d+p^kz$. The coefficients of the integer polynomial
$p^wf-d$ are divisible by $p^k$ in $\mathbf Z_p$, hence in $\mathbf Z$, so $p^wf-d=p^kg$ with $g$ over
$\mathbf Z$. By Step 0, $p^kg\in\nabla_{\mathbf Z}(\nu)$, and then $g\in\nabla_{\mathbf Z}(\nu)$,
because the conditions defining $\nabla_{\mathbf Z}(\nu)$ are linear and $\mathbf Z[X,y]$ has no
torsion. Hence $cp^wf=cd+Ng\in\Delta_{\mathbf Z}(\nu)$.

\emph{Step 4: over a domain.} Let $\mathcal O'$ be a domain in which every integer prime to $p$ is a
unit. The map $\mathbf Z[X]\to\mathcal O'[X]$ commutes with the operators $E^{(k)}_{ab}$ and fixes
$h_\nu$, so it carries $\Delta_{\mathbf Z}(\nu)$ into $\Delta_{\mathcal O'}(\nu)$. Let $b$ be a
bideterminant of shape $\nu$. By Step 3 there is $c$ prime to $p$ with $cp^wb\in\Delta_{\mathbf
Z}(\nu)$, so $cp^wb\in\Delta_{\mathcal O'}(\nu)$, and $p^wb\in\Delta_{\mathcal O'}(\nu)$ because $c$
is a unit. These $b$ span $\nabla_{\mathcal O'}(\nu)$, so
$p^w\nabla_{\mathcal O'}(\nu)\subset\Delta_{\mathcal O'}(\nu)$.

\emph{Step 5.} In $\mathcal O$ every integer $c$ prime to $p$ is a unit. Indeed $up+vc=1$ for some
integers $u,v$; in the local ring $\mathcal O$ one of $up$ and $vc$ is a unit, and $up$ is not, since
$p$ is not a unit. So Step 4 applies to $\mathcal O$ and gives the first inclusion.
\end{proof}
% Lean: sandwichLocal = sandwichDomain (sandwichDomain_of sandwichZ ascentDomain, with weylLeDual for
% the second inclusion) and unitsOfLocal; sandwichZ = descent_over_integers (dsc_denominator from
% sSBasis and weylRational, dsc_split, dualWeylMap, dsc_approx, dsc_divide, dualWeylSat) applied to
% sandwichZp = coreQp with spread_top_holds or spread_square_holds (spr_upper, spr_lower);
% ascent_domain_holds uses standardMonomialSpanDomain, bidetMem and asc_map_weyl.

\begin{lemma}\label{pd-kappa-gram-form}
Let $p$ and $(W,Q,\sigma,\iota)$ be as in Definition~\ref{pd-kappa-split-rank}, with rings
$\mathcal R\subset\mathcal K$, let $r\geq1$, $Z$ an $r\times r$ matrix of variables and
$\nu_1\geq\cdots\geq\nu_r=0$ integers. Let $[\ ,\ ]$ be the Fock pairing on $\mathcal K[Z]$ and put
$B_0(f,f')=[f,f']/[h_\nu,h_\nu]$ (Lemma~\ref{pd-kappa-gram-fock}). Let $w=0$ if $\nu_1=0$ and
$w=\frac{r(r-1)}2\lfloor\log_p(\nu_1+r-1)\rfloor$ otherwise. Then:
\begin{enumerate}
\item $B_0(f(gZ),f')=B_0(f,f'({}^tgZ))$ for all $f,f'\in\mathcal K[Z]$ and all $r\times r$ matrices
  $g$ over $\mathcal K$;
\item if $Q\in\nabla_{\mathcal K}(\nu)$ and $B_0(P,Q)=0$ for every $P\in\nabla_{\mathcal K}(\nu)$, then
  $Q=0$;
\item if $Q\in\nabla_{\mathcal K}(\nu)$ and $B_0(P,Q)\in\mathcal R$ for every
  $P\in\nabla_{\mathcal R}(\nu)$, then $Q\in\nabla_{\mathcal R}(\nu)$;
\item $B_0(P,p^wQ)\in\mathcal R$ for all $P,Q\in\nabla_{\mathcal R}(\nu)$.
\end{enumerate}
\end{lemma}
\leanref{Formalism/LocalToGlobal/GroupLattices/FockModels/TestSideFinal.lean}{Formalism.L2.G.Model.formCore}
\leanref{Formalism/LocalToGlobal/GroupLattices/FockModels/TransposeWords/FormCoreSK.lean}{Formalism.L2.G.Model.TS.Top.formCoreSK}
\leanref{Formalism/LocalToGlobal/GroupLattices/FockModels/TransposeWords/TransposeWordsProofs/FormSK.lean}{Formalism.L2.G.Model.TS.formCoreSK_of}
\leanref{Formalism/LocalToGlobal/GroupLattices/FockModels/TransposeWords/TransposeWordsProofs/WeylFormNondegenerateGlobal.lean}{Formalism.L2.G.Model.TS.formCoreG_of}
\leanref{Formalism/LocalToGlobal/GroupLattices/FockModels/TransposeWords/TransposeWordsProofs/WeylFormNondegenerateGlobal.lean}{Formalism.L2.G.Model.TS.weylPair_of}
\leanref{Formalism/LocalToGlobal/GroupLattices/FockModels/TransposeWords/TransposeWordsProofs/WeylFormNondegenerateGlobal.lean}{Formalism.L2.G.Model.TS.dualWeylIn_of}
\leanref{Formalism/LocalToGlobal/GroupLattices/FockModels/TransposeWords/TransposeWordsProofs/WeylFormNondegenerateGlobal.lean}{Formalism.L2.G.Model.TS.nondegK_of}

\begin{proof}
The field $\mathcal K$ contains $\iota(Q)$, so it has characteristic $0$, and
$[h_\nu,h_\nu]=N_\nu\neq0$ by Lemma~\ref{pd-kappa-gram-fock}(3). Part (1) is
Lemma~\ref{pd-kappa-gram-fock}(1). For a composite $u$ of operators $E^{(k)}_{ab}$ let $u^t$ be the
composite of the corresponding $E^{(k)}_{ba}$ in the reverse order; by Lemma~\ref{pd-kappa-gram-fock}(2),
$[uf,f']=[f,u^tf']$. The operators $E^{(k)}_{ab}$ preserve $\nabla_{\mathcal R}(\nu)$ and
$\nabla_{\mathcal K}(\nu)$, and $h_\nu\in\nabla_{\mathcal R}(\nu)$, as in the proof of
Lemma~\ref{pd-gain-weyl-comparison}; so $\Delta_{\mathcal R}(\nu)\subset\nabla_{\mathcal R}(\nu)$ and
$\Delta_{\mathcal K}(\nu)\subset\nabla_{\mathcal K}(\nu)$.

\emph{Step 1.} If $P\in\Delta_{\mathcal R}(\nu)$ and $Q\in\nabla_{\mathcal R}(\nu)$, then
$B_0(P,Q)\in\mathcal R$. It suffices to take $P=uh_\nu$ for a composite $u$. Then
$u^tQ\in\nabla_{\mathcal R}(\nu)$, and by Lemma~\ref{pd-kappa-gram-fock}(4),
$B_0(uh_\nu,Q)=[h_\nu,u^tQ]/[h_\nu,h_\nu]=(u^tQ)(1)$, which lies in $\mathcal R$.

\emph{Step 2.} Let $Q\in\nabla_{\mathcal K}(\nu)$ and $Z^e$ a monomial with $e_{as}=0$ for $a<s$. If
$\sum_ae_{as}\neq\nu_s$ for some $s$, the coefficient $Q_e$ of $Z^e$ in $Q$ is $0$, as $Q$ has degree
$\nu_s$ in the $s$-th column. Otherwise, with $w_e$ as in Lemma~\ref{pd-kappa-gram-fock}(5), parts (5),
(4) and (2) of that lemma give $Q_e=(w_eQ)(1)=[h_\nu,w_eQ]/[h_\nu,h_\nu]=B_0(w_e^th_\nu,Q)$.

(2) For $Q$ as in (2), $w_e^th_\nu\in\Delta_{\mathcal K}(\nu)\subset\nabla_{\mathcal K}(\nu)$, so all
the $Q_e$ of Step 2 are $0$, and $Q=0$ by Lemma~\ref{pd-kappa-gram-fock}(6).

(3) For $Q$ as in (3), $w_e^th_\nu\in\Delta_{\mathcal R}(\nu)\subset\nabla_{\mathcal R}(\nu)$, so all
the $Q_e$ of Step 2 lie in $\mathcal R$, and $Q\in\nabla_{\mathcal R}(\nu)$ by
Lemma~\ref{pd-kappa-gram-fock}(6), applied with $A=\mathcal R$.

(4) The Fock pairing is symmetric. If $\nu_1=0$, then $\nu=0$, an element of $\nabla_{\mathcal R}(0)$
has degree $0$ in every column and so is a constant, and $h_0=1$; so
$\nabla_{\mathcal R}(0)=\mathcal R=\Delta_{\mathcal R}(0)$, and Step 1 gives (4) with $w=0$. Otherwise,
the ring $\mathcal R$ is a local domain in which $p=\varpi_0^2$ is not a unit, and $\nu_t\leq\nu_1$ for
every $t$, so Lemma~\ref{pd-kappa-defects-sandwich} with $l=\nu_1$ gives
$p^w\nabla_{\mathcal R}(\nu)\subset\Delta_{\mathcal R}(\nu)$. So $p^wQ\in\Delta_{\mathcal R}(\nu)$ and
$B_0(P,p^wQ)=B_0(p^wQ,P)\in\mathcal R$ by Step 1.
\end{proof}

\begin{lemma}\label{pd-kappa-gram-divisors}
Let $\mathcal O$ be a discrete valuation ring with fraction field $K$ and uniformizer $\varpi$, put
$\pi=\varpi^2$, and let $N\geq0$ be an integer. Let $V$ be a $K$-vector space with a basis
$b_1,\dots,b_n$, let $\Lambda$ be the $\mathcal O$-span of the $b_i$, and let $B_0$ be a $K$-bilinear
form on $V$ such that:
\begin{enumerate}
\item if $Q\in V$ and $B_0(P,Q)=0$ for all $P\in V$, then $Q=0$;
\item if $Q\in V$ and $B_0(P,Q)\in\mathcal O$ for all $P\in\Lambda$, then $Q\in\Lambda$;
\item $B_0(P,\pi^NQ)\in\mathcal O$ for all $P,Q\in\Lambda$.
\end{enumerate}
Then there are $x_1,\dots,x_n,x'_1,\dots,x'_n\in\Lambda$ and integers $0\leq\kappa_i\leq N$ with
$\pi^NB_0(x_i,x'_j)=\delta_{ij}\pi^{\kappa_i}$ for all $i$ and $j$.
\end{lemma}
\leanref{Formalism/SharpRestriction/GrowthBounds/HermitianTensor/GrowthOfBasic.lean}{Formalism.MainNotions.flipOneThree.Pf.pb_abs}

\begin{proof}
By (3) the matrix $M$ with entries $M_{kl}=\pi^NB_0(b_k,b_l)$ has entries in $\mathcal O$. For
$d\in\mathcal O^n$ put $\operatorname{vec}(d)=\sum_kd_kb_k$; then
$\pi^NB_0(\operatorname{vec}(d),\operatorname{vec}(d'))=(dM)\cdot d'$. Let $N'\subset\mathcal O^n$ be the
submodule of the rows $dM$. By the elementary divisor theorem over the principal ideal domain
$\mathcal O$, there are a basis $\beta_1,\dots,\beta_n$ of $\mathcal O^n$, an integer $n'\leq n$, an
injective map $f$ from $\{1,\dots,n'\}$ to $\{1,\dots,n\}$ and nonzero $a_1,\dots,a_{n'}\in\mathcal O$
such that the $a_i\beta_{f(i)}$ form a basis of $N'$. Let $\gamma_j\in\mathcal O^n$ be such that the
coordinate of $v\in\mathcal O^n$ along $\beta_j$ is $v\cdot\gamma_j$. Then
$\pi^NB_0(\operatorname{vec}(d),\operatorname{vec}(\gamma_j))$ is the coordinate of $dM$ along
$\beta_j$, and every element of $N'$ has its coordinate along $\beta_{f(i)}$ in $a_i\mathcal O$ and its
coordinates along the $\beta_j$ with $j$ not in the image of $f$ equal to $0$.

The map $f$ is onto. Otherwise, for $j$ not in its image, $B_0(b_k,\operatorname{vec}(\gamma_j))=0$
for all $k$, so $\operatorname{vec}(\gamma_j)=0$ by (1), and $\gamma_j=0$ because the $b_k$ are linearly
independent; but $\beta_j\cdot\gamma_j=1$. So $n'=n$, and we may number so that $f$ is the identity.
Choose $d_i$ with $d_iM=a_i\beta_i$. Then
$\pi^NB_0(\operatorname{vec}(d_i),\operatorname{vec}(\gamma_j))=\delta_{ij}a_i$.

Next, $a_i$ divides $\pi^N$. Put $Q=(\pi^N/a_i)\operatorname{vec}(\gamma_i)$. For every $k$,
$\pi^NB_0(b_k,\operatorname{vec}(\gamma_i))$ is the coordinate along $\beta_i$ of the $k$-th row of
$M$, which lies in $a_i\mathcal O$; so $B_0(b_k,Q)\in\mathcal O$, hence $B_0(P,Q)\in\mathcal O$ for all
$P\in\Lambda$, and $Q\in\Lambda$ by (2). Write $Q=\operatorname{vec}(q)$ with $q\in\mathcal O^n$. By
linear independence $q=(\pi^N/a_i)\gamma_i$, so $\pi^N/a_i=(\pi^N/a_i)(\beta_i\cdot\gamma_i)=
\beta_i\cdot q\in\mathcal O$.

Hence $a_i=u_i\varpi^{m_i}$ with $u_i$ a unit and $0\leq m_i\leq2N$. Put
$\kappa_i=\lceil m_i/2\rceil\leq N$, $x_i=\operatorname{vec}(d_i)$ and
$x'_i=u_i^{-1}\varpi^{2\kappa_i-m_i}\operatorname{vec}(\gamma_i)$. These lie in $\Lambda$, and
$\pi^NB_0(x_i,x'_j)=\delta_{ij}u_i^{-1}\varpi^{2\kappa_i-m_i}a_i=\delta_{ij}\pi^{\kappa_i}$.
\end{proof}

\begin{lemma}\label{pd-kappa-cloud-gram}
Let $p$ and $(W,Q,\sigma,\iota)$ be as in Definition~\ref{pd-kappa-split-rank}, with rings
$\mathcal R\subset\mathcal K$, let $r\geq1$, and let $Z$ be an $r\times r$ matrix of variables. Let
$\nu\in\mathbf Z^r$ be non-increasing with $\nu_r=0$, and put $n=\dim_r(\nu)$
(Lemma~\ref{pd-kappa-cloud-dimensions}). Then there are a $\mathcal K$-bilinear form $B$ on
$\mathcal K[Z]$, elements $x_b$ and $x'_b$, $1\leq b\leq n$, of the dual Weyl module over $\mathcal R$
of $\nu$ (Definition~\ref{pd-weight-lattices}(1)), and integers
$0\leq\kappa_b\leq\frac{r(r-1)}2\lfloor\log_p(\nu_1+r-1)\rfloor$ (which is $0$ if $r=1$), such that:
\begin{enumerate}
\item $B(f(lZ),f')=B(f,f'({}^tl\,Z))$ for all $r\times r$ matrices $l$ over $\mathcal K$ and all
  $f,f'\in\mathcal K[Z]$;
\item $B(f,f')\in\mathcal R$ for $f$ and $f'$ in the dual Weyl module over $\mathcal R$ of $\nu$;
\item $B(x_b,x'_{b'})=\delta_{bb'}\,p^{\kappa_b}$ for all $b$ and $b'$.
\end{enumerate}
\end{lemma}
\leanref{Formalism/SharpRestriction/GrowthBounds/HermitianTensor/DetFactsPAlt.lean}{Formalism.MainNotions.flipOneThree.Pf.SmithP}
\leanref{Formalism/SharpRestriction/GrowthBounds/HermitianTensor/PbSmithLam.lean}{Formalism.MainNotions.flipOneThree.Pf.smithP}
\leanref{Formalism/SharpRestriction/GrowthBounds/HermitianTensor/PbSmithLam.lean}{Formalism.MainNotions.flipOneThree.Pf.pb_smith_lam}
\leanref{Formalism/SharpRestriction/GrowthBounds/HermitianTensor/PbSmithLam.lean}{Formalism.MainNotions.flipOneThree.Pf.pb_finrank}
\leanref{Formalism/SharpRestriction/GrowthBounds/LatticeSystems.lean}{Formalism.MainNotions.G.LatticeSystems.la_basis}

\begin{proof}
Let $B_0$ and $w$ be as in Lemma~\ref{pd-kappa-gram-form}, and put $B=p^wB_0$. By the standard basis
theorem \cite{DEP80}, over $\mathcal R$ and over $\mathcal K$ (Lemma~\ref{is-standard-basis}(1)), the
bideterminants of shape $\nu$ form an $\mathcal R$-basis of $\nabla_{\mathcal R}(\nu)$ and a
$\mathcal K$-basis $b_1,\dots,b_{n'}$ of $\nabla_{\mathcal K}(\nu)$, and $n'=\dim_r(\nu)=n$ by the Weyl
dimension formula \cite[Theorem~6.3]{FH91}. The ring $\mathcal R$ is a discrete valuation ring with
uniformizer $\varpi_0$ and $\varpi_0^2=p$. Parts (2), (3) and (4) of Lemma~\ref{pd-kappa-gram-form} are
hypotheses (1), (2) and (3) of Lemma~\ref{pd-kappa-gram-divisors} for $V=\nabla_{\mathcal K}(\nu)$, the
$b_i$, the restriction of $B_0$, $\varpi=\varpi_0$, $\pi=p$ and $N=w$. That lemma gives
$x_b,x'_b\in\nabla_{\mathcal R}(\nu)$ and integers $0\leq\kappa_b\leq w$ with
$B(x_b,x'_{b'})=p^wB_0(x_b,x'_{b'})=\delta_{bb'}p^{\kappa_b}$, which is (3). Part (1) is
Lemma~\ref{pd-kappa-gram-form}(1), and (2) holds because $B(f,f')=B_0(f,p^wf')$ lies in $\mathcal R$ by
Lemma~\ref{pd-kappa-gram-form}(4). Finally
$w\leq\frac{r(r-1)}2\lfloor\log_p(\nu_1+r-1)\rfloor$, and $w=0$ if $r=1$.
\end{proof}

\begin{lemma}\label{pd-kappa-cloud-lift}
Let $A$ be a domain, $r\geq1$ and $k=r+2$, and let $Y$ be a $k\times k$ and $Z$ an $r\times r$ matrix
of variables. Let $\hat Z$ be the block diagonal matrix with diagonal blocks $1$, $Z$ and $1$. Dual
Weyl modules are those of Definition~\ref{pd-weight-lattices}(1) over $A$, indexed by natural weights.
Let $\beta\in\mathbf Z^r$ be non-increasing with $\beta_r\geq0$.
\begin{enumerate}
\item Let $Q\geq0$. Every element of the dual Weyl module of $(\beta_1+Q,\dots,\beta_r+Q)$ is
  $\det(Z)^Qu_0$ with $u_0$ in the dual Weyl module of $\beta$.
\item Let $c\geq\beta_1$. For every $u$ in the dual Weyl module of $\beta$ there is $s$ in the dual
  Weyl module of $(c,\beta_1,\dots,\beta_r,0)$ such that every monomial of $s$ has degree $c$ in the
  first row of $Y$ and $0$ in the last row, and $s(\hat Z)=u$.
\end{enumerate}
\end{lemma}
\leanref{Formalism/SharpRestriction/GrowthBounds/HermitianTensor/DerivativeRepresentation.lean}{Formalism.MainNotions.flipOneThree.Pf.detDivP}
\leanref{Formalism/SharpRestriction/GrowthBounds/HermitianTensor/InitialMinorLiftFirstEntry.lean}{Formalism.MainNotions.flipOneThree.Pf.liftP}
\leanref{Formalism/SharpRestriction/GrowthBounds/HermitianTensor/InitialMinorLiftFirstEntry.lean}{Formalism.MainNotions.flipOneThree.Pf.qc_lift_bidet}

\begin{proof}
By the standard basis theorem over $A$ \cite{DEP80} (Lemma~\ref{is-standard-basis}(1)), the dual Weyl
module of a natural weight $\mu$ is spanned by the bideterminants of shape $\mu$: products of minors
of $Z$, the $l$-th on rows $I_l(1)<\cdots<I_l(s_l)$ and the first $s_l$ columns, where
$\#\{l:s_l\geq t\}=\mu_t$ for every $t$. A minor on the first $q$ columns lies in the module of
$(1,\dots,1,0,\dots,0)$, with $q$ entries $1$, and products lie in the module of the sum of the
weights, as in the proof of Lemma~\ref{pd-kappa-cloud-cartan}(2).

(1) For $\mu=(\beta_1+Q,\dots,\beta_r+Q)$ the number of $l$ with $s_l=r$ is $\beta_r+Q\geq Q$, and a minor
on $r$ increasing rows and all $r$ columns is $\det(Z)$. Removing $Q$ of these factors from a
bideterminant of shape $\mu$ leaves one of shape $\beta$. So each spanning element, and hence each
element, is $\det(Z)^Q$ times an element of the module of $\beta$.

(2) It suffices to treat a bideterminant $u=\prod_lM_l$ of shape $\beta$, since sums and multiples of
solutions are solutions. For a minor $M$ of $Z$ on rows $i_1<\cdots<i_q$ and the first $q$ columns,
let $\tilde M$ be the minor of $Y$ on the rows $1,i_1+1,\dots,i_q+1$ and the first $q+1$ columns. It
lies in the module of $(1,\dots,1,0,\dots,0)$ with $q+1$ entries $1$, each of its monomials has degree
$1$ in the first row and $0$ in the last row, and $\tilde M(\hat Z)=M$, since the rows
$1,i_1+1,\dots,i_q+1$ and the first $q+1$ columns of $\hat Z$ form the block diagonal matrix with
blocks $1$ and the rows $i_1,\dots,i_q$ and first $q$ columns of $Z$. The number $d$ of factors $M_l$
with $s_l\geq1$ is $\beta_1\leq c$, and the factors with $s_l=0$ are $1$. Put
$s=Y_{11}^{c-d}\prod_{l:s_l\geq1}\tilde M_l$. It lies in the module of
$(c-d)(1,0,\dots,0)+\sum_{l:s_l\geq1}(1,\dots,1,0,\dots,0)=(c,\beta_1,\dots,\beta_r,0)$, its monomials
have degree $c$ in the first row and $0$ in the last, and $s(\hat Z)=u$.
\end{proof}

\begin{lemma}\label{pd-kappa-cloud-pairing}
Let $p$ and $(W,Q,\sigma,\iota)$ be as in Definition~\ref{pd-kappa-split-rank}, with rings
$\mathcal R\subset\mathcal K$, let $r\geq1$, and let $Z$ and $Z'$ be $r\times r$ matrices of
variables. Dual Weyl modules are those of Definition~\ref{pd-weight-lattices}(1), indexed by natural
weights. Let $a,c,N\geq0$ be integers with $c\leq8a$, and $\alpha\in\mathbf Z^r$ non-increasing with
$0\leq\alpha_i\leq N$ and $\alpha_i\leq\alpha_j+c$ for all $i,j$; put
$\alpha^*=(N-\alpha_r,\dots,N-\alpha_1)$ and $n=\dim_r(\alpha)$
(Lemma~\ref{pd-kappa-cloud-dimensions}). Then there are a $\mathcal K$-linear map
$L\colon\mathcal K[Z,Z']\to\mathcal K$, elements $u_b$ of the dual Weyl module over $\mathcal R$ of
$\alpha$ and $u'_b$ of that of $\alpha^*$, for $1\leq b\leq n$, and integers
$0\leq\kappa_b\leq\lceil(r+1)^2(\log_p(8a+r+1)+1)\rceil$, such that:
\begin{enumerate}
\item $L(f(lZ)f'(lZ'))=\det(l)^NL(f(Z)f'(Z'))$ for $l\in GL_r(\mathcal K)$, $f$ in the dual Weyl module
  over $\mathcal K$ of $\alpha$ and $f'$ in that of $\alpha^*$;
\item $L(f(Z)f'(Z'))\in\mathcal R$ for $f$ in the dual Weyl module over $\mathcal R$ of $\alpha$ and
  $f'$ in that of $\alpha^*$;
\item $L(u_b(Z)u'_{b'}(Z'))=\delta_{bb'}\,p^{\kappa_b}$ for all $b$ and $b'$.
\end{enumerate}
\end{lemma}
\leanref{Formalism/SharpRestriction/GrowthBounds/HermitianTensor/CoeffSumMonomialLinear.lean}{Formalism.MainNotions.flipOneThree.Pf.QPh}
\leanref{Formalism/SharpRestriction/GrowthBounds/HermitianTensor/CoeffSumMonomialLinear.lean}{Formalism.MainNotions.flipOneThree.Pf.PhInv}
\leanref{Formalism/SharpRestriction/GrowthBounds/HermitianTensor.lean}{Formalism.MainNotions.flipOneThree.Pf.qPh_proved}
\leanref{Formalism/SharpRestriction/GrowthBounds/HermitianTensor/DetYNeZero.lean}{Formalism.MainNotions.flipOneThree.Pf.qPh_of}
\leanref{Formalism/SharpRestriction/GrowthBounds/HermitianTensor/DetFactsP.lean}{Formalism.MainNotions.flipOneThree.Pf.numP}

\begin{proof}
Put $e=\alpha_r$, so that $e\leq\alpha_i\leq N$ for all $i$. Let $\tau$ be the map of
Lemma~\ref{pd-kappa-cloud-transpose} for $N$, $e$ and $\alpha$, with $\mathcal O=\mathcal R$, and let
$B$, $x_b$, $x'_b$ and $\kappa_b$ be given by Lemma~\ref{pd-kappa-cloud-gram} for
$\nu=\alpha-e$, which is non-increasing with $\nu_r=0$. The number of the $x_b$ is
$\dim_r(\alpha-e)=\dim_r(\alpha)=n$, since $\dim_r$ is unchanged when a constant is added to all
entries. By Lemma~\ref{pd-kappa-cloud-lift}(1) over $\mathcal K$ and over $\mathcal R$, every $f$ in
the dual Weyl module of $\alpha$ is $\det(Z)^eg$ with $g$ in that of $\alpha-e$, over the same ring.
Multiplication by $\det(Z)^e$ is injective on $\mathcal K[Z]$; let $\sigma$ be a $\mathcal K$-linear
left inverse of it. Let $L$ be the $\mathcal K$-linear map with
$L(f(Z)f'(Z'))=B(\sigma(f),\tau(f'))$ for $f,f'\in\mathcal K[Z]$; it exists because
$\mathcal K[Z,Z']$ is the tensor product of $\mathcal K[Z]$ and $\mathcal K[Z']$.

(1) Let $f=\det(Z)^eg$ and $f'$ be as in (1), and $l\in GL_r(\mathcal K)$. Then
$f(lZ)=\det(l)^e\det(Z)^eg(lZ)$, so $\sigma(f(lZ))=\det(l)^eg(lZ)$. With $f'_l(Z)=f'(lZ)$,
Lemma~\ref{pd-kappa-cloud-gram}(1) and Lemma~\ref{pd-kappa-cloud-transpose}(2) give
\[
 L(f(lZ)f'(lZ'))=\det(l)^eB(g(lZ),\tau(f'_l))=\det(l)^eB\big(g,\tau(f'_l)({}^tl\,Z)\big)
 =\det(l)^NB(g,\tau(f'))=\det(l)^NL(f(Z)f'(Z')).
\]

(2) Let $f=\det(Z)^eg$ with $g$ over $\mathcal R$, and $f'$ over $\mathcal R$. By
Lemma~\ref{pd-kappa-cloud-transpose}(1), $\tau(f')$ lies in the dual Weyl module over $\mathcal R$ of
$\alpha-e$, so $L(f(Z)f'(Z'))=B(g,\tau(f'))\in\mathcal R$ by Lemma~\ref{pd-kappa-cloud-gram}(2).

(3) Put $u_b=\det(Z)^ex_b$. The polynomial $\det(Z)^e$ lies in the dual Weyl module over $\mathcal R$
of $(e,\dots,e)$, being a product of minors on all $r$ columns, and a product lies in the module of
the sum of the weights, as in the proof of Lemma~\ref{pd-edge-products}; so $u_b$ lies in the module
over $\mathcal R$ of $\alpha$. By Lemma~\ref{pd-kappa-cloud-transpose}(1) choose $u'_b$ in the module
over $\mathcal R$ of $\alpha^*$ with $\tau(u'_b)=x'_b$. Then $L(u_b(Z)u'_{b'}(Z'))=B(x_b,x'_{b'})=
\delta_{bb'}p^{\kappa_b}$. Finally $\nu_1=\alpha_1-e\leq c\leq8a$, $r(r-1)/2\leq(r+1)^2$ and
$\lfloor\log_p(\nu_1+r-1)\rfloor\leq\log_p(8a+r+1)$, so
$\kappa_b\leq\lceil(r+1)^2(\log_p(8a+r+1)+1)\rceil$.
\end{proof}

\begin{lemma}\label{pd-kappa-cloud-substitution}
Let $p$ and $(W,Q,\sigma,\iota)$ be as in Definition~\ref{pd-kappa-split-rank}, with rings
$\mathcal R\subset\mathcal K$, let $m\geq2$, $r=m-1$ and $k=m+1$, and let $Y$ be a $k\times k$ and $Z$
an $r\times r$ matrix of variables. Let $\hat Z$ be the block diagonal matrix with diagonal blocks
$1$, $Z$ and $1$. Dual Weyl modules are those of Definition~\ref{pd-weight-lattices}(1), indexed by
natural weights, and $h_\nu$ is the polynomial defined there. Let $A$ be $\mathcal R$ or $\mathcal K$,
and regard $f(Y\hat Z)$, for $f\in A[Y]$, as a polynomial in $Y$ with coefficients in $A[Z]$.
\begin{enumerate}
\item If $f$ lies in the dual Weyl module over $A$ of $\nu\in\mathbf Z^k$, then every coefficient of
  $f(Y\hat Z)$ lies in the dual Weyl module over $A$ of $(\nu_2,\dots,\nu_{k-1})$.
\item If every monomial of $f$ has degree $d$ in the entries of the $i$-th row of $Y$, then so does
  every monomial of $f(Y\hat Z)$.
\item Let $P,Q\geq0$. If $F$ lies in the dual Weyl module over $\mathcal K$ of
  $(P+Q,Q,\dots,Q,0)$, the natural weight of $(P,0,\dots,0,-Q)$, then $F(Y\hat Z)=\det(Z)^QF(Y)$.
\item Let $a,b\geq0$ and $\nu=(a+b,b,\dots,b,0)$, the natural weight of $(a,0,\dots,0,-b)$. If $G$
  lies in the dual Weyl module over $\mathcal K$ of $\nu$ and every monomial of $G$ has degree $a+b$
  in the first row and $0$ in the last row, then $G=G(1)\,h_\nu$, where $G(1)$ is the value at the
  identity matrix.
\end{enumerate}
\end{lemma}
\leanref{Formalism/SharpRestriction/GrowthBounds/HermitianTensor/CoeffSumMonomialLinear.lean}{Formalism.MainNotions.flipOneThree.Pf.subY}
\leanref{Formalism/SharpRestriction/GrowthBounds/HermitianTensor/SubCoeffP.lean}{Formalism.MainNotions.flipOneThree.Pf.subCoeffP}
\leanref{Formalism/SharpRestriction/GrowthBounds/HermitianTensor/QaWeightRowInd.lean}{Formalism.MainNotions.flipOneThree.Pf.subRowP}
\leanref{Formalism/SharpRestriction/GrowthBounds/HermitianTensor/SubYDPsiMonomial.lean}{Formalism.MainNotions.flipOneThree.Pf.subMuP_of}
\leanref{Formalism/SharpRestriction/GrowthBounds/HermitianTensor/TopLinePOf.lean}{Formalism.MainNotions.flipOneThree.Pf.topLineP_of}

\begin{proof}
(1) Scaling the $j$-th column of $Z$ by $t$ scales the column $j+1$ of $Y\hat Z$ by $t$ and fixes its
other columns, and $f$ is homogeneous of degree $\nu_{j+1}$ in that column; so every coefficient of
$f(Y\hat Z)$ is homogeneous of degree $\nu_{j+1}$ in the $j$-th column of $Z$. For $i<j$ the block
diagonal matrix formed from $Z(1+yE_{ij})$ is $\hat Z(1+yE_{i+1,j+1})$, and
$f(X(1+yE_{i+1,j+1}))=f(X)$; so $f(Y\hat Z)$, and with it each coefficient, is unchanged when $Z$ is
replaced by $Z(1+yE_{ij})$.

(2) The substitution replaces each entry of $Y$ by a linear combination of entries of the same row.

(3) By Lemma~\ref{pd-kappa-harmonic-model}(1), $F$ is a linear combination of polynomials
$a(Y_{11},\dots,Y_{k1})\,b(c(Y))$ with $b$ homogeneous of degree $Q$, where $c(Y)$ is the vector of
signed minors of Lemma~\ref{pd-kappa-harmonic-minors}. The substitution fixes the first column, and
the first $m$ columns of $Y\hat Z$ are the first $m$ columns of $Y$ multiplied by the block diagonal
matrix with blocks $1$ and $Z$. So $c(Y\hat Z)=\det(Z)c(Y)$, and $b(c(Y\hat Z))=\det(Z)^Qb(c(Y))$.

(4) By Lemma~\ref{pd-kappa-harmonic-model}(1), $G=\delta\Psi(f)$ for some
$f=\sum f_{u,v}x^uy^v\in\mathrm{Pol}(a,b)$, and $\delta\Psi(x^uy^v)=\prod_iY_{i1}^{u_i}c_i(Y)^{v_i}$.
Since $Y_{i1}$ has degree $1$ in the $i$-th row and $c_i(Y)$ has degree $1$ in each row other than the
$i$-th, every monomial of $\delta\Psi(x^uy^v)$ has degree $u_1+b-v_1$ in the first row and $u_k+b-v_k$
in the last. These are $a+b$ and $0$ only for $x^uy^v=x_1^ay_k^b$, because $u_1\leq a$ and $v_k\leq b$.
By hypothesis $G$ has no monomials of other degrees, so comparing monomials gives
$G=f_{a,b}\,Y_{11}^ac_k(Y)^b$, where $f_{a,b}$ is the coefficient of $x_1^ay_k^b$. The polynomial
$h_\nu=\delta_1^a\delta_m^b$ is a product of minors on the first columns, so it lies in the dual Weyl
module of $\nu$, as in the proof of Lemma~\ref{pd-kappa-cloud-cartan}(2), and it has degree $a+b$ in
the first row and $0$ in the last. So $h_\nu=f'\,Y_{11}^ac_k(Y)^b$ for some $f'\in\mathcal K$, and
$f'\neq0$ because $h_\nu(1)=1$. Hence $G$ is a multiple of $h_\nu$, and evaluating at the identity
matrix gives $G=G(1)h_\nu$.
\end{proof}

\begin{lemma}\label{pd-kappa-cloud-contraction}
Let $p$ and $(W,Q,\sigma,\iota)$ be as in Definition~\ref{pd-kappa-split-rank}, with rings
$\mathcal R\subset\mathcal K$, and let $\nabla_{\mathcal K}(\mu)$ and $\nabla_{\mathcal R}(\mu)$ be as in
Lemma~\ref{pd-kappa-cloud-cartan} with $\mathcal O=\mathcal R$. Let $m\geq2$ and $k=m+1$, let
$a,c,P,Q\geq0$ be integers with $2c\leq8a$, put $\ell=\log_p(8a+m)+1$ and
$h(P,Q)=(P,0,\dots,0,-Q)\in\mathbf Z^k$, and let $\lambda\in\mathbf Z^k$ be dominant with
$\lambda_1=P+c$, $\lambda_k=-Q-c$ and $|\lambda_i|\leq c$ for $1<i<k$. Put $\omega=\lambda-h(P,Q)$
and $\omega^*=(-\omega_k,\dots,-\omega_1)$; both are dominant, with first entry $c$ and last entry
$-c$. Let $x$, $Y$ and $Y'$ be $k\times k$ matrices of variables, and let $h_\nu\in\mathcal R[x]$ be
the polynomial of Definition~\ref{pd-weight-lattices} for the natural weight $\nu$ of $h(2c,2c)$.
Then there is a $\mathcal K$-linear map $\Gamma\colon\mathcal K[Y,Y']\to\mathcal K[x]$ such that, for
all $f\in\nabla_{\mathcal K}(\lambda)$ and $f'\in\nabla_{\mathcal K}(\omega^*)$:
\begin{enumerate}
\item $\Gamma(f(Y)f'(Y'))\in\nabla_{\mathcal K}(h(P+2c,Q+2c))$;
\item $\Gamma(f(gY)f'(gY'))=\Gamma(f(Y)f'(Y'))(gx)$ for every $g\in GL_k(\mathcal K)$;
\item if $f\in\nabla_{\mathcal R}(\lambda)$ and $f'\in\nabla_{\mathcal R}(\omega^*)$, then
  $\Gamma(f(Y)f'(Y'))\in\nabla_{\mathcal R}(h(P+2c,Q+2c))$;
\item there are $s_b\in\nabla_{\mathcal R}(\omega)$, $s'_b\in\nabla_{\mathcal R}(\omega^*)$ and
  integers $0\leq h_b\leq\lceil m^2\ell\rceil$, for $1\leq b\leq n$ with
  $n=\dim_{m-1}(\lambda_2,\dots,\lambda_m)$ (Lemma~\ref{pd-kappa-cloud-dimensions}), such that every
  monomial of $s_b$ and of $s'_b$ has degree $2c$ in the entries of the first row, and
  \[
   \Gamma\big((Fs_b)(Y)\,s'_{b'}(Y')\big)=\delta_{bb'}\,p^{h_b}\,h_\nu F
   \qquad(F\in\nabla_{\mathcal K}(h(P,Q)),\ 1\leq b,b'\leq n).
  \]
\end{enumerate}
\end{lemma}
\leanref{Formalism/SharpRestriction/GrowthBounds/HermitianTensor/CloudADef.lean}{Formalism.MainNotions.flipOneThree.Pf.ContractionStandardLattice}
\leanref{Formalism/SharpRestriction/GrowthBounds/HermitianTensor.lean}{Formalism.MainNotions.flipOneThree.Pf.contraction_standard_lattice_holds}
\leanref{Formalism/SharpRestriction/GrowthBounds/HermitianTensor/ModelIdentificationHolds.lean}{Formalism.MainNotions.flipOneThree.Pf.contraction_of_pairing_GL}

\begin{proof}
Put $r=m-1\geq1$, let $Z$ and $Z'$ be $r\times r$ matrices of variables, and let $\hat Z$ and
$\hat Z'$ be as in Lemma~\ref{pd-kappa-cloud-substitution}. In this proof dual Weyl modules are
indexed by natural weights, as there. The natural weights of $\lambda$, $\omega$ and $\omega^*$ are
\[
 \nu^\lambda=(P+Q+2c,\lambda_2+Q+c,\dots,\lambda_m+Q+c,0),\quad
 \nu^\omega=(2c,\lambda_2+c,\dots,\lambda_m+c,0),\quad
 \nu^*=(2c,c-\lambda_m,\dots,c-\lambda_2,0).
\]
Put $N=Q+2c$ and $\alpha=(\lambda_2+Q+c,\dots,\lambda_m+Q+c)\in\mathbf Z^r$, the middle part of
$\nu^\lambda$. Since $|\lambda_i|\leq c$, $Q\leq\alpha_i\leq N$ and $\alpha_i\leq\alpha_j+2c$, and
$2c\leq8a$. The middle part of $\nu^*$ is $\alpha^*=(N-\alpha_r,\dots,N-\alpha_1)$, and that of
$\nu^\omega$ is $\beta=\alpha-Q$. Lemma~\ref{pd-kappa-cloud-pairing}, applied with $r$, $a$, $2c$, $N$
and $\alpha$, gives $L$, $u_b$, $u'_b$ and $\kappa_b\leq\lceil m^2\ell\rceil$ for $1\leq b\leq n$,
where $n=\dim_r(\alpha)=\dim_{m-1}(\lambda_2,\dots,\lambda_m)$.

\emph{The map $\Gamma$.} For $G\in\mathcal K[Y,Y']$ substitute $x\hat Z$ for $Y$ and $x\hat Z'$ for
$Y'$, regard the result as a polynomial in $x$ with coefficients in $\mathcal K[Z,Z']$, and apply $L$
to each coefficient; the result is $\Gamma(G)$. This is $\mathcal K$-linear, and directly from the
definition:
\begin{enumerate}
\item[(a)] $\Gamma(G)(gx)=\Gamma(G(gY,gY'))$ for $g\in GL_k(\mathcal K)$, since $(gx)\hat Z=g(x\hat Z)$;
\item[(b)] for $G=v(Y)t(Y')$, the coefficient of $x^u$ in $\Gamma(G)$ is the sum over $u=u'+u''$ of
  $L(v_{u'}(Z)t_{u''}(Z'))$, where $v_{u'}$ and $t_{u''}$ are the coefficients of $Y^{u'}$ in
  $v(Y\hat Z)$ and of $Y^{u''}$ in $t(Y\hat Z)$;
\item[(c)] the value of $\Gamma(G)$ at the identity matrix is $L(G(\hat Z,\hat Z'))$;
\item[(d)] if $F\in\mathcal K[Y]$ satisfies $F(Y\hat Z)=\det(Z)^QF(Y)$, then
  $\Gamma(F(Y)G)=F(x)\,\Gamma_Q(G)$, where $\Gamma_Q(G)$ is defined as $\Gamma(G)$ with the map
  $f\mapsto L(\det(Z)^Qf)$ in place of $L$.
\end{enumerate}

\emph{Step 1.} Let $\mu,\mu'\in\mathbf Z^k$ be natural weights with last entries $0$ and middle parts
$\alpha$ and $\alpha^*$, and let $v$ and $t$ lie in the dual Weyl modules over $\mathcal K$ of $\mu$
and $\mu'$. We show that $F=\Gamma(v(Y)t(Y'))$ lies in the dual Weyl module of
$\mu^+=(\mu_1+\mu'_1,N,\dots,N,0)$. Let $l\in GL_r(\mathcal K)$. Since
$(x\hat l)\hat Z=x\widehat{lZ}$, the polynomial $F(x\hat l)$ is obtained by applying $L$ to the
coefficients of $v(x\widehat{lZ})\,t(x\widehat{lZ'})$. By Lemma~\ref{pd-kappa-cloud-substitution}(1),
the coefficients of $v(x\hat Z)$ and $t(x\hat Z')$ lie in the dual Weyl modules over $\mathcal K$ of
$\alpha$ and $\alpha^*$, so Lemma~\ref{pd-kappa-cloud-pairing}(1) gives $F(x\hat l)=\det(l)^NF(x)$.
Since $\mathcal K$ is infinite, taking $l$ diagonal shows that $F$ is homogeneous of degree $N$ in each
of the columns $2,\dots,m$, and taking $l=1+yE_{ij}$ with $i<j$ and $y\in\mathcal K$ shows that $F$ is
unchanged when $y$ times the column $i+1$ is added to the column $j+1$, also for a variable $y$. The substitution fixes the first and the last columns, in which $v$
and $t$ have degrees $\mu_1$ and $\mu'_1$, and $0$; so $F$ has degree $\mu_1+\mu'_1$ in the first column
and $0$ in the last. For $2\leq j\leq m$, adding $y$ times the first column of $x$ to its $j$-th
column replaces $x\hat Z$ by $x\hat Z\prod_i(1+yZ_{j-1,i}E_{1,i+1})$, a product of additions of
multiples of the first column to later ones, and the same holds for $Z'$; so $v$, $t$ and $F$ are
unchanged. Adding a multiple of a column to the last column does not change $F$, which does not
involve the last column. Hence $F$ lies in the dual Weyl module of $\mu^+$.

\emph{Step 2: parts (1), (2) and (3).} Part (1) is Step 1 with $\mu=\nu^\lambda$ and $\mu'=\nu^*$,
since $(P+Q+4c,N,\dots,N,0)$ is the natural weight of $h(P+2c,Q+2c)$, and part (2) is (a). For (3),
by Lemma~\ref{pd-kappa-cloud-substitution}(1) over $\mathcal R$ the coefficients of $f(Y\hat Z)$ and
$f'(Y\hat Z')$ lie in the dual Weyl modules over $\mathcal R$ of $\alpha$ and $\alpha^*$, so by (b) and
Lemma~\ref{pd-kappa-cloud-pairing}(2) the polynomial $\Gamma(f(Y)f'(Y'))$ has coefficients in
$\mathcal R$. By (1) it lies in $\nabla_{\mathcal K}(h(P+2c,Q+2c))$, and the elements of this module
with coefficients in $\mathcal R$ form $\nabla_{\mathcal R}(h(P+2c,Q+2c))$, as in the proof of
Lemma~\ref{pd-kappa-cloud-cartan}(1).

\emph{Step 3: the elements.} By Lemma~\ref{pd-kappa-cloud-lift}(1) over $\mathcal R$,
$u_b=\det(Z)^Qw_b$ with $w_b$ in the dual Weyl module over $\mathcal R$ of $\beta$. Since
$\beta_1=\lambda_2+c\leq2c$, Lemma~\ref{pd-kappa-cloud-lift}(2) gives $s_b$ in the module over
$\mathcal R$ of $(2c,\beta,0)=\nu^\omega$, of degree $2c$ in the first row and $0$ in the last, with
$s_b(\hat Z)=w_b$. Since $\alpha^*_1=c-\lambda_m\leq2c$, it also gives $s'_b$ in the module over
$\mathcal R$ of $(2c,\alpha^*,0)=\nu^*$, of degree $2c$ in the first row and $0$ in the last, with
$s'_b(\hat Z)=u'_b$.

\emph{Step 4: the identity in (4).} Let $\zeta$ and $\eta$ be the natural weights of $h(2c,2c)$ and
$h(0,Q)$, so that $\zeta+\eta$ is that of $h(2c,Q+2c)$; then $h_\nu=h_\zeta$. From the definition of
these polynomials as products of leading principal minors, $h_{\zeta+\eta}=h_\zeta h_\eta$ and
$h_\eta(1)=1$. By Lemma~\ref{pd-kappa-cloud-substitution}(3), every $F\in\nabla_{\mathcal K}(h(P,Q))$
satisfies $F(Y\hat Z)=\det(Z)^QF(Y)$, and so does $F_0=h_\eta$, which lies in
$\nabla_{\mathcal K}(h(0,Q))$ as a product of minors on the first columns. By (d),
$\Gamma\big((Fs_b)(Y)s'_{b'}(Y')\big)=F\,\Phi_{bb'}$ for all such $F$, where
$\Phi_{bb'}=\Gamma_Q(s_b(Y)s'_{b'}(Y'))$ does not depend on $F$. Put
$G_{bb'}=\Gamma\big((F_0s_b)(Y)s'_{b'}(Y')\big)=F_0\Phi_{bb'}$. The product $F_0s_b$ lies in the dual
Weyl module of $\eta+\nu^\omega$, whose middle part is $Q+\beta=\alpha$, so by Step 1, $G_{bb'}$ lies in
that of $(Q+4c,N,\dots,N,0)=\zeta+\eta$. The polynomial $F_0=\delta_m^Q$ has degree $Q$ in the first row
and $0$ in the last, and $s_b$ and $s'_{b'}$ have degrees $2c$ and $0$. By
Lemma~\ref{pd-kappa-cloud-substitution}(2), and because $L$ acts on the coefficients in $Z$ and $Z'$,
every monomial of $G_{bb'}$ has degree $Q+4c$ in the first row and $0$ in the last.
Lemma~\ref{pd-kappa-cloud-substitution}(4), with $a=2c$ and $b=Q+2c$, gives
$G_{bb'}=G_{bb'}(1)h_{\zeta+\eta}$. By (c), Lemma~\ref{pd-kappa-cloud-substitution}(3) at $Y=1$ and
Lemma~\ref{pd-kappa-cloud-pairing}(3),
\[
 G_{bb'}(1)=L\big(F_0(\hat Z)s_b(\hat Z)s'_{b'}(\hat Z')\big)=L\big(\det(Z)^Qw_b(Z)u'_{b'}(Z')\big)
 =L\big(u_b(Z)u'_{b'}(Z')\big)=\delta_{bb'}p^{\kappa_b}.
\]
So $h_\eta\Phi_{bb'}=\delta_{bb'}p^{\kappa_b}h_\zeta h_\eta$, and since $h_\eta\neq0$,
$\Phi_{bb'}=\delta_{bb'}p^{\kappa_b}h_\zeta$. Hence
$\Gamma\big((Fs_b)(Y)s'_{b'}(Y')\big)=\delta_{bb'}p^{\kappa_b}h_\nu F$ for
$F\in\nabla_{\mathcal K}(h(P,Q))$, which is (4) with $h_b=\kappa_b$.
% Lean: ContractionStandardLattice = contraction_of_pairing_GL qPh_proved CartanOnto
%   = contraction_of_pairing_lemmas (HermitianTensor/ContractionOfPairingLemmas.lean), with Gamma =
%   lamOf L (lcoef L after sub2, the substitution subY on each letter). (a) to (d) are lamEquivP,
%   lamCoeffP, lamEvalP and lamMulP; Step 1 is lamMemP; Step 2 uses subCoeffP, subMapP and intCapP;
%   Step 3 is detDivP and liftP (with shapeC_lo and shapeC_lo'); Step 4 is subMuP_of, rowA_lam,
%   topLineP_of and hwMulP.
\end{proof}

\begin{lemma}\label{pd-kappa-cloud-position}
Let $p$, $(W,Q,\sigma,\iota)$, $\mathcal R\subset\mathcal K$, $m\geq2$, $k=m+1$, $a,c,P,Q$, $\ell$,
$h(P,Q)$, $\lambda$, $\omega$, $\omega^*$, $x$, $Y$, $Y'$ and $n$ be as in
Lemma~\ref{pd-kappa-cloud-contraction}, let $\varpi_0$ be the uniformizer of the setting, and put
$\zeta=h(2c,2c)$. Let $\nabla_{\mathcal K}(\mu)$ and $\nabla_\mu(\Lambda^\vee)$ be as in
Lemma~\ref{pd-kappa-cloud-cartan} with $\mathcal O=\mathcal R$. For $u\in\mathcal K^\times$ let $d_u$
be the diagonal matrix with entries $u^{-1},1,\dots,1$. Then there is a $\mathcal K$-linear map
$\Gamma\colon\mathcal K[Y,Y']\to\mathcal K[x]$ with the following properties.
\begin{enumerate}
\item For $f\in\nabla_{\mathcal K}(\lambda)$ and $f'\in\nabla_{\mathcal K}(\omega^*)$,
  $\Gamma(f(Y)f'(Y'))\in\nabla_{\mathcal K}(h(P+2c,Q+2c))$. If moreover $f\in\nabla_\lambda(\Lambda^\vee)$
  and $f'\in\nabla_{\omega^*}(\Lambda^\vee)$ for an $\mathcal R$-lattice $\Lambda\subset\mathcal K^k$,
  then $\Gamma(f(Y)f'(Y'))\in\nabla_{h(P+2c,Q+2c)}(\Lambda^\vee)$.
\item For every $g\in GL_k(\mathcal K)$ there are $z_0,s_b,s'_b\in\mathcal K[x]$ and integers
  $0\leq h_b\leq\lceil m^2\ell\rceil$, for $1\leq b\leq n$, such that
  $\Gamma\big((Fs_b)(Y)s'_{b'}(Y')\big)=\delta_{bb'}p^{h_b}z_0F$ for all
  $F\in\nabla_{\mathcal K}(h(P,Q))$ and all $b,b'$, and such that for every $u\in\mathcal K^\times$,
  with $\Lambda_u$ the lattice spanned by the columns of $gd_u$:
  $u^{2c}z_0\in\nabla_\zeta(\Lambda_u^\vee)$, $u^{2c}z_0\notin\varpi_0\nabla_\zeta(\Lambda_u^\vee)$,
  $u^cs_b\in\nabla_\omega(\Lambda_u^\vee)$ and $u^cs'_b\in\nabla_{\omega^*}(\Lambda_u^\vee)$.
\end{enumerate}
\end{lemma}
\leanref{Formalism/SharpRestriction/GrowthBounds/HermitianTensor/CloudDefs.lean}{Formalism.MainNotions.flipOneThree.Pf.CloudOne}
\leanref{Formalism/SharpRestriction/GrowthBounds/HermitianTensor/CloudA.lean}{Formalism.MainNotions.flipOneThree.Pf.cloudOne_of}
\leanref{Formalism/SharpRestriction/GrowthBounds/HermitianTensor/MemInvariantSectionsIffIntegral.lean}{Formalism.MainNotions.flipOneThree.Pf.core_id}
\leanref{Formalism/SharpRestriction/GrowthBounds/HermitianTensor/MemInvariantSectionsIffIntegral.lean}{Formalism.MainNotions.flipOneThree.Pf.hw_prim}

\begin{proof}
Let $\Gamma$ be the map of Lemma~\ref{pd-kappa-cloud-contraction}. The first claim of (1) is part
(1) of that lemma. Let $\Lambda$ be spanned by the columns of $g$. By
Lemma~\ref{pd-kappa-cloud-cartan}(1), $\det(g)^{\lambda_k}f(gx)\in\nabla_{\mathcal R}(\lambda)$ and
$\det(g)^{\omega^*_k}f'(gx)\in\nabla_{\mathcal R}(\omega^*)$. By parts (3) and (2) of
Lemma~\ref{pd-kappa-cloud-contraction}, $\Gamma$ of the product of these two, which is
$\det(g)^{\lambda_k+\omega^*_k}\Gamma(f(Y)f'(Y'))(gx)$, lies in
$\nabla_{\mathcal R}(h(P+2c,Q+2c))$. Since $\lambda_k+\omega^*_k=-Q-2c$ is the last entry of
$h(P+2c,Q+2c)$, Lemma~\ref{pd-kappa-cloud-cartan}(1) gives the second claim of (1).

(2) Let $s^0_b$, $s'^0_b$ and $h_b$ be given by Lemma~\ref{pd-kappa-cloud-contraction}(4), and put
\[
 z_0=\det(g)^{2c}h_\nu(g^{-1}x),\qquad s_b=\det(g)^cs^0_b(g^{-1}x),\qquad
 s'_b=\det(g)^cs'^0_b(g^{-1}x).
\]
Let $F\in\nabla_{\mathcal K}(h(P,Q))$. Then $F(gx)\in\nabla_{\mathcal K}(h(P,Q))$, as in the proof of
Lemma~\ref{pd-kappa-cloud-cartan}(1), and $F(gx)s^0_b\in\nabla_{\mathcal K}(\lambda)$ because
$h(P,Q)+\omega=\lambda$. Since $Fs_b=\det(g)^c(F(gy)s^0_b(y))|_{y=g^{-1}x}$, parts (2) and (4) of
Lemma~\ref{pd-kappa-cloud-contraction}, the first with $g^{-1}$, give
\[
 \Gamma\big((Fs_b)(Y)s'_{b'}(Y')\big)=\det(g)^{2c}\,\Gamma\big((F(gY)s^0_b(Y))s'^0_{b'}(Y')\big)(g^{-1}x)
 =\delta_{bb'}p^{h_b}z_0F .
\]
Let $u\in\mathcal K^\times$. If every monomial of $f\in\mathcal K[x]$ has degree $D$ in the first row,
then $f(d_ux)=u^{-D}f(x)$; and $\det(gd_u)=u^{-1}\det(g)$. The polynomial $h_\nu$ is a product of
$\nu_1=4c$ leading principal minors, each of degree $1$ in the first row, and the last entry of $\zeta$
is $-2c$. Hence
\[
 \det(gd_u)^{-2c}(u^{2c}z_0)(gd_ux)=u^{2c}\det(g)^{-2c}\cdot u^{2c}\det(g)^{2c}h_\nu(d_ux)
 =u^{4c}h_\nu(d_ux)=h_\nu .
\]
Now $h_\nu$ is a product of minors on the first columns, so it lies in $\nabla_{\mathcal R}(\zeta)$, as
in the proof of Lemma~\ref{pd-kappa-cloud-cartan}(2), and Lemma~\ref{pd-kappa-cloud-cartan}(1) gives
$u^{2c}z_0\in\nabla_\zeta(\Lambda_u^\vee)$. If $u^{2c}z_0=\varpi_0z'$ with
$z'\in\nabla_\zeta(\Lambda_u^\vee)$, the same computation gives $h_\nu=\varpi_0y$ with
$y\in\nabla_{\mathcal R}(\zeta)\subset\mathcal R[x]$; evaluating at the identity matrix gives
$1\in\varpi_0\mathcal R$, which is false. The last entries of $\omega$ and $\omega^*$ are $-c$, and
$s^0_b$ and $s'^0_b$ have degree $2c$ in the first row, so in the same way
$\det(gd_u)^{-c}(u^cs_b)(gd_ux)=s^0_b\in\nabla_{\mathcal R}(\omega)$ and
$\det(gd_u)^{-c}(u^cs'_b)(gd_ux)=s'^0_b\in\nabla_{\mathcal R}(\omega^*)$. This gives the remaining
claims, by Lemma~\ref{pd-kappa-cloud-cartan}(1).
\end{proof}

\begin{lemma}\label{pd-kappa-cloud-maps}
Let $p$ and $(W,Q,\sigma,\iota)$ be as in Definition~\ref{pd-kappa-split-rank}, $m\geq2$ and
$k=m+1$, $T\subset\mathbf Z/n'\mathbf Z$, $a,c_{cl}\geq0$ integers with $2c_{cl}\leq8a$,
$a_w,b_w\colon\mathbf Z/n'\to\mathbf Z$, and $\lambda_0(s)\in\mathbf Z^k$ dominant weights with
$\lambda_0(s)_1=a_w(s)$, $\lambda_0(s)_k=-b_w(s)$ and $|\lambda_0(s)_i|\leq c_{cl}$ for $1<i<k$
when $s\in T$, and $\lambda_0(s)=0$ when $s\notin T$. Suppose that $A=a_w-c_{cl}$ and
$B=b_w-c_{cl}$ are nonnegative on $T$, and put $\phi=\phi(A,B)$ and
$\phi^+=\phi(A+2c_{cl},B+2c_{cl})$. Write $V^\phi$, $R^\phi$, $\mathcal C^\phi_{t,c}$,
$\mathcal C'^\phi_{t,c}$ and $A^\phi$ for the space, lattice, cores and integer of
Definitions~\ref{pd-projection-property} and~\ref{pd-kappa-split-rank} formed with $\phi$, the same
with $\phi^+$ for $\phi^+$, and $V$, $\mathcal C_{t,c}$, $\mathcal C'_{t,c}$ and $A$ for those formed
with $\lambda_0$. Put
$n_b=\prod_{s\in T}\dim_{m-1}(\lambda(s)_2,\dots,\lambda(s)_m)$, with $\dim_{m-1}$ as in
Lemma~\ref{pd-kappa-cloud-dimensions}, and $\ell=\log_p(8a+m)+1$. Let $e$ be an edge. Then there are
$\mathcal K$-linear maps $\theta\colon V^\phi\to V^{\phi^+}$, $U_b\colon V^\phi\to V$ and
$L_b\colon V\to V^{\phi^+}$ and integers $0\leq h_b\leq|T|\lceil m^2\ell\rceil$, for
$1\leq b\leq n_b$, such that:
\begin{enumerate}
\item $L_{b'}U_b=\delta_{bb'}p^{h_b}\theta$ for all $b$ and $b'$;
\item for $F\in V^\phi$, $\theta(F)\in R^{\phi^+}(e)$ if and only if $F\in R^\phi(e)$;
\item for all integers $t,c\geq0$ with $2|T|c_{cl}\leq c$, $\theta$ carries $\mathcal C^\phi_{t,c}(e)$
  into $\mathcal C^{\phi^+}_{t,c}(e)$, each $U_b$ carries $\mathcal C^\phi_{t,c}(e)$ into
  $\mathcal C_{t,c}(e)$, and $p^{A^{\phi^+}-A}(G\circ L_b)\in\mathcal C'_{t,c-2|T|c_{cl}}(e)$ for every
  $b$ and every $G\in\mathcal C'^{\phi^+}_{t,c}(e)$.
\end{enumerate}
\end{lemma}
\leanref{Formalism/SharpRestriction/GrowthBounds/ColumnSpanDensity/CloudPropOfMapsRetractionBasis.lean}{Formalism.MainNotions.flipOneThree.Pf.RestCloudMaps}
\leanref{Formalism/SharpRestriction/GrowthBounds/HermitianTensor/TheoremIOfTwoLetterClausesContraction.lean}{Formalism.MainNotions.flipOneThree.Pf.cloud_maps_of_contraction}
\leanref{Formalism/SharpRestriction/GrowthBounds/HermitianTensor/CloudC.lean}{Formalism.MainNotions.flipOneThree.Pf.restCloudMaps_of}
\leanref{Formalism/SharpRestriction/GrowthBounds/HermitianTensor/CloudDefs.lean}{Formalism.MainNotions.flipOneThree.Pf.RestCloudLat}
\leanref{Formalism/SharpRestriction/GrowthBounds/HermitianTensor/CloudB.lean}{Formalism.MainNotions.flipOneThree.Pf.restCloudLat_of}
\leanref{Formalism/SharpRestriction/GrowthBounds/HermitianTensor/ProdMemHten.lean}{Formalism.MainNotions.flipOneThree.Pf.cloudTen}
\leanref{Formalism/SharpRestriction/GrowthBounds/HermitianTensor/LastNeZero.lean}{Formalism.MainNotions.flipOneThree.Pf.posOK_at_top}
\leanref{Formalism/SharpRestriction/GrowthBounds/HermitianTensor/CloudIso.lean}{Formalism.MainNotions.flipOneThree.Pf.cloudIso}
\leanref{Formalism/SharpRestriction/GrowthBounds/HermitianTensor/CcDualWtAnti.lean}{Formalism.MainNotions.flipOneThree.Pf.cc_CCL}
\leanref{Formalism/SharpRestriction/GrowthBounds/HermitianTensor/CcDualWtAnti.lean}{Formalism.MainNotions.flipOneThree.Pf.cc_CCL'}

\begin{proof}
For a family $\mu=(\mu(s))_{s\in\mathbf Z/n'}$ of dominant weights in $\mathbf Z^k$ let $\bar\mu(s)$
be formed from $\mu(s)$ as $\lambda(s)$ is formed from $\lambda_0(s)$ in
Definition~\ref{pd-local-model}, and write $V^\mu$, $H^\mu$, $R^\mu$, $R'^\mu$ and
$A^\mu=\sum_s\bar\mu(s)_1$ for the space, lattices and integer of
Definition~\ref{pd-projection-property} with $r=1$ formed with $\mu$. For a family $M=(M_s)$ of
$\mathcal R$-lattices in $\mathcal K^k$ let $H^\mu(M)$ be the $\mathcal R$-span of the products
$\prod_sf_s(x_s)$ with $f_s\in\nabla_{\bar\mu(s)}(M_s^\vee)$, in the notation of
Lemma~\ref{pd-kappa-cloud-cartan} with $\mathcal O=\mathcal R$. By
Definition~\ref{pd-projection-property}(1), $H^\mu(L)$ and $H^\mu(e)$ are of this form.

Put $\omega=\lambda_0-\phi$, $\omega^*(s)=(-\omega(s)_k,\dots,-\omega(s)_1)$ and
$\zeta=\phi(2c_{cl},2c_{cl})$, formed with the constant functions $2c_{cl}$ on $T$. For $s\in T$,
$\omega(s)$ has first entry $a_w(s)-A(s)=c_{cl}$, last entry $B(s)-b_w(s)=-c_{cl}$ and the middle
entries of $\lambda_0(s)$; for $s\notin T$, $\omega(s)=0$. So $\omega$, $\omega^*$ and $\zeta$ are
dominant, $\phi+\omega=\lambda_0$ and $\phi+\zeta=\phi^+$. For $s\in T$, $\bar\omega(s)$ and
$\bar\omega^*(s)=(-\bar\omega(s)_k,\dots,-\bar\omega(s)_1)$ have first entry $c_{cl}$ and last entry
$-c_{cl}$, and $\bar\zeta(s)=h(2c_{cl},2c_{cl})$ in the notation of
Lemma~\ref{pd-kappa-cloud-dimensions}. Hence, with $d=|T|c_{cl}$,
\[
 A^\zeta=2d,\qquad A^\omega=A^{\omega^*}=d,\qquad A^{\phi^+}=A^\phi+2d,\qquad A=A^\phi+d,\qquad
 A^{\phi^+}-A=d,
\]
and the last entries of the $\bar\omega^*(s)$ sum to $-d$.

\emph{Step 1: one position.} For $s\in T$ put $(P_s,Q_s)=(A(s),B(s))$ if $s$ is even and
$(P_s,Q_s)=(B(s),A(s))$ if $s$ is odd. Then $\lambda(s)=\bar\lambda_0(s)$ is dominant with first
entry $P_s+c_{cl}$, last entry $-Q_s-c_{cl}$ and middle entries of absolute value at most $c_{cl}$;
moreover $\bar\phi(s)=h(P_s,Q_s)$, $\bar\phi^+(s)=h(P_s+2c_{cl},Q_s+2c_{cl})$ and
$\bar\omega(s)=\lambda(s)-h(P_s,Q_s)$. Lemma~\ref{pd-kappa-cloud-position}, applied with $c=c_{cl}$
to $\lambda(s)$, gives a map $\Gamma_s\colon\mathcal K[Y,Y']\to\mathcal K[x]$ with its properties (1)
and (2), where $n=n_s=\dim_{m-1}(\lambda(s)_2,\dots,\lambda(s)_m)$. For $s\notin T$ all the weights
$\bar\mu(s)$ above are $0$, the dual Weyl modules of the weight $0$ consist of the constants, and
$\nabla_0(\Lambda^\vee)=\mathcal R$ for every lattice $\Lambda$. Let
$\Gamma_s(f(Y)f'(Y'))=f(x)f'(x)$, $n_s=1$, $z_0=s_1=s'_1=1$ and $h_1=0$; then (1) and (2) of
Lemma~\ref{pd-kappa-cloud-position} hold at $s$ as well.

\emph{Step 2: products over the positions.} As in the proof of Lemma~\ref{pd-kappa-weight-tensor},
multiplication identifies $V^\mu$ with the tensor product over $s$ of the dual Weyl modules over
$\mathcal K$ of the natural weights of the $\bar\mu(s)$. So, by Step 1, there is a bilinear map
$\Gamma\colon V\times V^{\omega^*}\to V^{\phi^+}$ with
$\Gamma(\prod_sf_s(x_s),\prod_sf'_s(x_s))=\prod_s\Gamma_s(f_s(Y)f'_s(Y'))(x_s)$, and multiplication
gives bilinear maps $V^\phi\times V^\zeta\to V^{\phi^+}$ and $V^\phi\times V^\omega\to V$
(Lemma~\ref{pd-edge-products}). Let $M$ be a family of lattices. By
Lemma~\ref{pd-kappa-cloud-cartan}(2) at each $s$, $H^\phi(M)H^\zeta(M)\subset H^{\phi^+}(M)$ and
$H^\phi(M)H^\omega(M)\subset H^{\lambda_0}(M)$, and by Lemma~\ref{pd-kappa-cloud-position}(1) at
each $s$, $\Gamma$ maps $H^{\lambda_0}(M)\times H^{\omega^*}(M)$ into $H^{\phi^+}(M)$; it suffices to
check these on spanning products. Hence, for integers $i,j$, products and $\Gamma$ map
$\varpi_0^iH^{\mu}(M)\times\varpi_0^jH^{\mu'}(M)$ into $\varpi_0^{i+j}H^{\mu''}(M)$ for each of the
three triples $(\mu,\mu',\mu'')$.

Let $g=(g_s)$ be a family of matrices in $GL_k(\mathcal K)$. For $s\in T$ let $z_{0,s}$,
$s_{\beta,s}$, $s'_{\beta,s}$ and $h_{\beta,s}$, $\beta\leq n_s$, be given by
Lemma~\ref{pd-kappa-cloud-position}(2) for $g_s$, and for $s\notin T$ take those of Step 1. Number the
families $(\beta_s)_s$ with $1\leq\beta_s\leq n_s$ as $b=1,\dots,n_b$, where $n_b=\prod_sn_s$ is the
number of the statement, and put
\[
 z_0=\prod_sz_{0,s}(x_s),\quad s_b=\prod_ss_{\beta_s,s}(x_s),\quad s'_b=\prod_ss'_{\beta_s,s}(x_s),\quad
 h_b=\sum_sh_{\beta_s,s}\leq|T|\lceil m^2\ell\rceil .
\]
For $F=\prod_sF_s(x_s)$ with each $F_s$ in the dual Weyl module of $\bar\phi(s)$, Step 1 gives
$\Gamma(Fs_b,s'_{b'})=\prod_s\Gamma_s\big((F_ss_{\beta_s,s})(Y)s'_{\beta'_s,s}(Y')\big)(x_s)$, which
is $p^{h_b}Fz_0$ if $b=b'$ and $0$ otherwise, since then some factor vanishes. By linearity,
$\Gamma(Fs_b,s'_{b'})=\delta_{bb'}p^{h_b}Fz_0$ for all $F\in V^\phi$. For $u\in\mathcal K^\times$ let
$M(u)$ be the family of the lattices spanned by the columns of the $g_sd_u$, with $d_u$ as in
Lemma~\ref{pd-kappa-cloud-position}. Taking products over $s$ in
Lemma~\ref{pd-kappa-cloud-position}(2) gives
\[
 u^{A^\zeta}z_0\in H^\zeta(M(u)),\qquad u^{A^\omega}s_b\in H^\omega(M(u)),\qquad
 u^{A^{\omega^*}}s'_b\in H^{\omega^*}(M(u)). \tag{a}
\]
Moreover:
\[
 \text{if }F\in V^\phi\text{ and }F\cdot u^{A^\zeta}z_0\in H^{\phi^+}(M(u)),\text{ then }
 F\in H^\phi(M(u)). \tag{b}
\]
Indeed, for each $s$ the element $z_s=u^{\bar\zeta(s)_1}z_{0,s}$ lies in
$\nabla_{\bar\zeta(s)}(M(u)_s^\vee)$ and not in $\varpi_0\nabla_{\bar\zeta(s)}(M(u)_s^\vee)$, so
Lemma~\ref{pd-kappa-cloud-primitive}(2) gives $\rho_s$ with
$\rho_s(\nabla_{\bar\phi^+(s)}(M(u)_s^\vee))\subset\nabla_{\bar\phi(s)}(M(u)_s^\vee)$ and
$\rho_s(z_sF_s)=F_s$ for $F_s$ in the dual Weyl module of $\bar\phi(s)$. The linear map $\rho$ with
$\rho(\prod_sf_s(x_s))=\prod_s\rho_s(f_s)(x_s)$ maps $H^{\phi^+}(M(u))$ into $H^\phi(M(u))$, and
$\rho(u^{A^\zeta}z_0F)=F$ for $F\in V^\phi$, first for products and then by linearity, because
$u^{A^\zeta}z_0=\prod_sz_s(x_s)$.

\emph{Step 3: the edge $e$.} Let $L''$ and $L$ be the source and target of $e$. As in the proof of
Lemma~\ref{pd-edge-basis}(1) with $r=1$, there are $g_s\in GL_k(\mathcal K)$ such that
$H^\mu(L'')=H^\mu(M(1))$, $H^\mu(e)=H^\mu(M(\varpi_0))$ and $H^\mu(L)=H^\mu(M(p))$ for every family
$\mu$. Take $z_0$, $s_b$, $s'_b$ and $h_b$ of Step 2 for these $g_s$, and put
\[
 \theta(F)=Fz_0,\qquad U_b(F)=Fs_b,\qquad L_b(G)=\Gamma(G,s'_b).
\]
Then (1) is the identity of Step 2. By (a) with $u=1$, $\varpi_0$ and $p$, and since
$R^\mu(e)=\varpi_0^{-A^\mu}H^\mu(e)$ and $p^{A^{\omega^*}}=\varpi_0^{2d}$,
\[
 z_0\in H^\zeta(L'')\cap R^\zeta(e),\qquad s_b\in H^\omega(L'')\cap R^\omega(e),\qquad
 s'_b\in R^{\omega^*}(e)\cap\varpi_0^{-2d}H^{\omega^*}(L).
\]

(2) If $F\in R^\phi(e)$, then $\theta(F)\in\varpi_0^{-A^\phi-A^\zeta}H^{\phi^+}(e)=R^{\phi^+}(e)$ by
Step 2. Conversely, if $\theta(F)\in R^{\phi^+}(e)$, then
$(\varpi_0^{A^\phi}F)\cdot\varpi_0^{A^\zeta}z_0=\varpi_0^{A^{\phi^+}}\theta(F)\in H^{\phi^+}(M(\varpi_0))$,
so $\varpi_0^{A^\phi}F\in H^\phi(e)$ by (b), and $F\in R^\phi(e)$.

(3) Let $t,c\geq0$ be integers with $2d\leq c$, and let $F\in\mathcal C^\phi_{t,c}(e)$: so
$F\in\varpi_0^{-t}H^\phi(L'')$, $F\in R^\phi(e)$, and $F\in\varpi_0^cR^\phi(e')$ for every edge
$e'\neq e$ with source $L''$. By Step 2, $z_0\in H^\zeta(L'')$ gives
$\theta(F)\in\varpi_0^{-t}H^{\phi^+}(L'')$, and $z_0\in R^\zeta(e)$ gives
$\theta(F)\in R^{\phi^+}(e)$. For $e'$ as above, $z_0\in H^\zeta(L'')\subset R^\zeta(e')$ by
Lemma~\ref{pd-edge-basis}(2), which gives $\theta(F)\in\varpi_0^cR^{\phi^+}(e')$. So
$\theta(F)\in\mathcal C^{\phi^+}_{t,c}(e)$. The same argument with
$s_b\in H^\omega(L'')\cap R^\omega(e)$ and $A=A^\phi+A^\omega$ gives
$U_b(F)\in\mathcal C_{t,c}(e)$.

Let $G\in\mathcal C'^{\phi^+}_{t,c}(e)$, and put $G_b=p^{A^{\phi^+}-A}(G\circ L_b)=\varpi_0^{2d}(G\circ
L_b)$. For $v\in H(L)$, Step 2 and $\varpi_0^{2d}s'_b\in H^{\omega^*}(L)$ give
$\Gamma(v,s'_b)\in\varpi_0^{-2d}H^{\phi^+}(L)$; as $G\in\varpi_0^{-t}H^{\phi^+}(L)^\vee$,
$G_b(v)\in\varpi_0^{-t}\mathcal R$. So $G_b\in\varpi_0^{-t}H(L)^\vee$. For $v\in H(e)$,
$s'_b\in R^{\omega^*}(e)=\varpi_0^{-d}H^{\omega^*}(e)$ gives
$\Gamma(v,s'_b)\in\varpi_0^{-d}H^{\phi^+}(e)$; as
$G\in R'^{\phi^+}(e)=\varpi_0^{-A^{\phi^+}}H^{\phi^+}(e)^\vee$,
$G_b(v)\in\varpi_0^{2d-d-A^{\phi^+}}\mathcal R=\varpi_0^{-A}\mathcal R$. So
$G_b\in\varpi_0^{-A}H(e)^\vee=R'(e)$. Let $e'\neq e$ be an edge with target $L$. The proof of
Lemma~\ref{pd-kappa-weight-order} shows that $\varpi_0^{-E}H^{\omega^*}(L)\subset H^{\omega^*}(e')$,
where $E=-d$ is the sum of the last entries of the $\bar\omega^*(s)$. Hence
$s'_b\in\varpi_0^{-2d}H^{\omega^*}(L)\subset\varpi_0^{-3d}H^{\omega^*}(e')$, and
$\Gamma(v,s'_b)\in\varpi_0^{-3d}H^{\phi^+}(e')$ for $v\in H(e')$. As
$G\in\varpi_0^cR'^{\phi^+}(e')=\varpi_0^{c-A^{\phi^+}}H^{\phi^+}(e')^\vee$,
$G_b(v)\in\varpi_0^{2d-3d+c-A^{\phi^+}}\mathcal R=\varpi_0^{c-2d-A}\mathcal R$, so
$G_b\in\varpi_0^{c-2d}R'(e')$. Therefore $G_b\in\mathcal C'_{t,c-2d}(e)$, and $2d=2|T|c_{cl}$.
\end{proof}

\begin{lemma}\label{pd-kappa-cloud-retraction}
Let $\mathcal O$ be a discrete valuation ring with fraction field $K$, let $V$ and $V'$ be
finite-dimensional $K$-vector spaces, and let $M\subset V$ and $M'\subset V'$ be the
$\mathcal O$-spans of $K$-bases of $V$ and $V'$. Let $\theta\colon V\to V'$ be a $K$-linear map such
that $\theta(F)\in M'$ if and only if $F\in M$, for every $F\in V$. Then there is a $K$-linear map
$\rho\colon V'\to V$ with $\rho\theta=\mathrm{id}_V$ and $\rho(M')\subset M$.
\end{lemma}
\leanref{Formalism/SharpRestriction/GrowthBounds/ColumnSpanDensity/FactorsLEOfFintype.lean}{Formalism.MainNotions.flipOneThree.Pf.LatRetract}
\leanref{Formalism/SharpRestriction/GrowthBounds/ColumnSpanDensity/CloudPropOfMapsBasis.lean}{Formalism.MainNotions.flipOneThree.Pf.latRetract}

\begin{proof}
The map $\theta$ is injective. Indeed, let $F\neq0$ with $\theta(F)=0$, let $x$ be a nonzero coordinate
of $F$ in the basis that spans $M$, and $\varpi$ a uniformizer. Then $(\varpi x)^{-1}F\notin M$,
while $\theta((\varpi x)^{-1}F)=0\in M'$. So $\theta$ maps $M$ isomorphically onto $\theta(M)\subset M'$.
The quotient $M'/\theta(M)$ is torsion-free: if $y\in M'$ and $ay=\theta(x)$ with $x\in M$ and
$a\in\mathcal O$ nonzero, then $\theta(a^{-1}x)=y\in M'$, so $a^{-1}x\in M$ and $y\in\theta(M)$. A
finitely generated torsion-free module over $\mathcal O$ is free, so $M'=\theta(M)\oplus N$ for an
$\mathcal O$-submodule $N$. Let $\rho_0\colon M'\to M$ be the projection onto $\theta(M)$ followed by
the inverse of $\theta$, and $\rho$ its $K$-linear extension, which exists because $M'$ is spanned by a
$K$-basis of $V'$. Then $\rho\theta$ is the identity on $M$, hence on $V$, and $\rho(M')\subset M$.
\end{proof}

\begin{lemma}\label{pd-kappa-cloud-smith}
Let $\mathcal O$ be a discrete valuation ring with maximal ideal $\mathfrak m$, let
$z_1,z_2\in\mathcal O$ be nonzero, $\sigma,\rho\geq0$ integers, $N$ a $\sigma\times\rho$ matrix and $B$
a $\rho\times\sigma$ matrix over $\mathcal O$. Then there are integers $r,s\geq0$ with $r+s=\sigma$
and matrices $P$, $Q$, $P'$, $P''$ and $M$ over $\mathcal O$, of sizes $r\times\sigma$, $\rho\times r$,
$s\times\sigma$, $\sigma\times s$ and $(s+\rho)\times(s+\rho)$, such that $PNQ=z_1z_2\mathbf 1_r$,
$P'P''=\mathbf 1_s$ and
\[
 \begin{pmatrix}z_1\mathbf 1_s&P'N\\BP''&z_2\mathbf 1_\rho\end{pmatrix}M=z_1z_2\mathbf 1_{s+\rho}.
\]
\end{lemma}
\leanref{Formalism/SharpRestriction/GrowthBounds/ColumnSpanDensity/FactorsLEOfFintype.lean}{Formalism.MainNotions.flipOneThree.Pf.LemmaSMat}
\leanref{Formalism/SharpRestriction/GrowthBounds/ColumnSpanDensity/LemmaSMat.lean}{Formalism.MainNotions.flipOneThree.Pf.KP.lemmaSMat}

\begin{proof}
Write $z=z_1z_2$. We first show, by induction on $\sigma$, that there are $r+s=\sigma$ and $P$, $Q$,
$P'$, $P''$ with $PNQ=z\mathbf 1_r$, $P'P''=\mathbf 1_s$ and all entries of $P'N$ in $z\mathfrak m$. If
all entries of $N$ lie in $z\mathfrak m$, take $r=0$, $s=\sigma$ and $P'=P''=\mathbf 1_\sigma$.
Otherwise let $N_{i_0j_0}$ be an entry of least valuation. It divides every entry of $N$, and it
divides $z$, since otherwise all entries would lie in $z\mathfrak m$. Write $N_{ij_0}=f_iN_{i_0j_0}$,
$N_{i_0j}=d_jN_{i_0j_0}$ and $z=uN_{i_0j_0}$, so that $d_{j_0}=1$ and the row $i_0$ of $N$ is
$N_{i_0j_0}d$ for the row vector $d=(d_j)$. Let $S$ be the $(\sigma-1)\times\sigma$ matrix such that
the rows of $SN$ are the $N_i-f_iN_{i_0}$, $i\neq i_0$, where $N_i$ is the $i$-th row of $N$, and $E$
the $\sigma\times(\sigma-1)$ matrix that puts a vector into the coordinates $i\neq i_0$, so that
$SE=\mathbf 1$. The column $j_0$ of $N_1=SN$ is $0$. Apply the claim to $N_1$, with
$r_1+s_1=\sigma-1$ and $P_1$, $Q_1$, $P'_1$, $P''_1$. Let $\epsilon$ be the $j_0$-th standard basis
column of $\mathcal O^\rho$ and $\epsilon'$ the $i_0$-th one of $\mathcal O^\sigma$, and put
\[
 P=\begin{pmatrix}u\,\epsilon'^t\\P_1S\end{pmatrix},\quad
 Q=\begin{pmatrix}\epsilon&Q_1-\epsilon dQ_1\end{pmatrix},\quad P'=P'_1S,\quad P''=EP''_1 .
\]
Since $u\epsilon'^tN=zd$, $d\epsilon=1$ and $N_1\epsilon=0$, the matrix $PNQ$ has the blocks $z$,
$z(dQ_1-dQ_1)=0$, $0$ and $P_1N_1Q_1=z\mathbf 1$.
Moreover $P'P''=P'_1P''_1=\mathbf 1$ and $P'N=P'_1N_1$ has entries in $z\mathfrak m$.

Now write $P'N=zX$, with the entries of $X$ in $\mathfrak m$, and $C=BP''$. The matrices
$\mathbf 1-XC$ and $\mathbf 1-CX$ are congruent to $\mathbf 1$ modulo $\mathfrak m$, so they are
invertible over $\mathcal O$, and multiplying out shows that
\[
 M=\begin{pmatrix}z_2(\mathbf 1-XC)^{-1}&-zX(\mathbf 1-CX)^{-1}\\
 -C(\mathbf 1-XC)^{-1}&z_1(\mathbf 1-CX)^{-1}\end{pmatrix}
\]
has the required property.
\end{proof}

\begin{lemma}\label{pd-kappa-cloud-upper}
Let $p$, $(W,Q,\sigma,\iota)$, $k$ and $\lambda_0$ be as in Definition~\ref{pd-kappa-split-rank},
$t,c,n\geq0$ integers and $D\geq0$ a real with $2D<c$. If for some edge $e$ there is a split family
of depth $D$ and size $n$ on $(\mathcal C_{t,c}(e),R'(e))$, then $n\leq\mu(\lambda_0)$.
\end{lemma}
\leanref{Formalism/SharpRestriction/GrowthBounds/ColumnSpanDensity/NLeDAll.lean}{Formalism.MainNotions.flipOneThree.Pf.nLeD}
\leanref{Formalism/SharpRestriction/GrowthBounds/ColumnSpanDensity/NLeDAll.lean}{Formalism.MainNotions.flipOneThree.Pf.nLeD_all}
\leanref{Formalism/SharpRestriction/GrowthBounds/ColumnSpanDensity/NLeDAll.lean}{Formalism.MainNotions.flipOneThree.Pf.nEdge}

\begin{proof}
Let $e'$ be an edge, and $\gamma\in GL_k(Q)$ with $e'=\gamma e$, as in the proof of
Lemma~\ref{pd-kappa-edge-move}. The action of $\gamma$ described there carries
$\mathcal C_{t,c}(e)$ onto $\mathcal C_{t,c}(e')$ and $R'(e)$ onto $R'(e')$ and preserves the edge
pairing, so there is such a family at every edge. The proof of Lemma~\ref{pd-kappa-rank-upper} uses
of the forms $y_{j,b}$ only that they lie in $R'(e_j)$, so it applies to these families and gives
$n\leq\mu(\lambda_0)$.
\end{proof}

\begin{lemma}\label{pd-kappa-cloud-count}
Let $p$, $(W,Q,\sigma,\iota)$ and $k$ be as in Definition~\ref{pd-kappa-split-rank}, and let
$\lambda^1$, $\lambda^2$ and $\lambda_0$ be families of dominant weights in $\mathbf Z^k$, indexed
by $s\in\mathbf Z/n'\mathbf Z$. Write $V^i$, $R^i$, $R'^i$, $A^i$, $\langle\ ,\ \rangle^i$,
$\mathcal C^i_{t,c}$ and $\mathcal C'^i_{t,c}$ for the space, lattices, integer, edge pairing and
cores of Definitions~\ref{pd-projection-property} and~\ref{pd-kappa-split-rank} formed with
$\lambda^i$, and $V$, $A$, $\mathcal C_{t,c}$ and $\mathcal C'_{t,c}$ for those formed with
$\lambda_0$. Let $e$ be an edge, $t,c,c'\geq0$ integers with $c'\leq c$, $D_1\geq0$ a real with
$4D_1<c$, and $n_b,h\geq0$ integers. Let $\theta\colon V^1\to V^2$, $U_b\colon V^1\to V$ and
$L_b\colon V\to V^2$, $1\leq b\leq n_b$, be $\mathcal K$-linear maps and $0\leq h_b\leq h$ integers
such that:
\begin{enumerate}
\item[(a)] $L_{b'}U_b=\delta_{bb'}p^{h_b}\theta$ for all $b$ and $b'$;
\item[(b)] for $F\in V^1$, $\theta(F)\in R^2(e)$ if and only if $F\in R^1(e)$;
\item[(c)] $\theta(\mathcal C^1_{t,c}(e))\subset\mathcal C^2_{t,c}(e)$ and
  $U_b(\mathcal C^1_{t,c}(e))\subset\mathcal C_{t,c}(e)$ for every $b$;
\item[(d)] $p^{A^2-A}(G\circ L_b)\in\mathcal C'_{t,c'}(e)$ for every $b$ and every
  $G\in\mathcal C'^2_{t,c}(e)$.
\end{enumerate}
Suppose that there are split families of depth $D_1$ and sizes $n_1$ on
$(\mathcal C^1_{t,c}(e),\mathcal C'^1_{t,c}(e))$ and $n_2$ on
$(\mathcal C^2_{t,c}(e),\mathcal C'^2_{t,c}(e))$. Then there are an integer
$r\geq n_1+n_2-\mu(\lambda^2)$ and a split family of depth $2D_1+h$ and size $n_br$ on
$(\mathcal C_{t,c'}(e),\mathcal C'_{t,c'}(e))$.
\end{lemma}
\leanref{Formalism/SharpRestriction/GrowthBounds/ColumnSpanDensity/CloudCount.lean}{Formalism.MainNotions.flipOneThree.Pf.cloud_count}
\leanref{Formalism/SharpRestriction/GrowthBounds/ColumnSpanDensity/CloudCount.lean}{Formalism.MainNotions.flipOneThree.Pf.be_pull}
\leanref{Formalism/SharpRestriction/GrowthBounds/ColumnSpanDensity/FactorsLEOfFintype.lean}{Formalism.MainNotions.flipOneThree.Pf.wit_of_factors}

\begin{proof}
Replacing the forms $y_j$ of a split family by $\sum_lM_{lj}y_l$, which lie in the same
$\mathcal R$-module, we may take $M$ to be the identity. So let $x^1_1,\dots,x^1_{n_1}$,
$y^1_1,\dots,y^1_{n_1}$ and $z_1$, and $x^2_1,\dots,x^2_{n_2}$, $y^2_1,\dots,y^2_{n_2}$ and $z_2$, be
the two families, with $\langle x^i_a,y^i_b\rangle^i=z_i\delta_{ab}$ and $v(z_i)\leq D_1$, where
$z_i=1$ if $n_i=0$. For $F\in V^1$ and $G\in(V^2)^\vee$,
\[
 \langle U_bF,p^{A^2-A}(G\circ L_{b'})\rangle=G(p^{A^2}L_{b'}U_bF)=\langle L_{b'}U_bF,G\rangle^2 .
\]

The lattices $R^1(e)$ and $R^2(e)$ are spanned by bases (Lemma~\ref{pd-edge-basis}(1)), so by (b)
and Lemma~\ref{pd-kappa-cloud-retraction} there is a $\mathcal K$-linear $\rho\colon V^2\to V^1$
with $\rho\theta=\mathrm{id}$ and $\rho(R^2(e))\subset R^1(e)$. Put
$\tau_l=p^{A^1-A^2}(y^1_l\circ\rho)$. Then $\langle F,\tau_l\rangle^2=\langle\rho F,y^1_l\rangle^1$ for
$F\in V^2$. For $F\in R^2(e)$ this lies in $\mathcal R$, by Lemma~\ref{pd-edge-basis}(3) for
$\lambda^1$, so $\tau_l\in R'^2(e)$ by Lemma~\ref{pd-edge-basis}(3) for $\lambda^2$. Moreover
$\langle\theta(x^1_a),\tau_l\rangle^2=z_1\delta_{al}$.

Let $N$ be the $n_1\times n_2$ matrix with entries $\langle\theta(x^1_a),y^2_b\rangle^2$ and $B$ the
$n_2\times n_1$ matrix with entries $\langle x^2_b,\tau_l\rangle^2$. They lie in $\mathcal R$ by
Lemma~\ref{pd-edge-basis}(3), since $\theta(x^1_a)$ and $x^2_b$ lie in $R^2(e)$ by (c) and $y^2_b$ and
$\tau_l$ in $R'^2(e)$. Lemma~\ref{pd-kappa-cloud-smith} gives $r+s=n_1$ and $P$, $Q$, $P'$, $P''$,
$M$. The vectors $X_i=\sum_aP'_{ia}\theta(x^1_a)$, $i\leq s$, and $x^2_1,\dots,x^2_{n_2}$, and the
forms $Y_i=\sum_aP''_{ai}\tau_a$, $i\leq s$, and $y^2_1,\dots,y^2_{n_2}$, have the matrix of pairings
\[
 \begin{pmatrix}z_1P'P''&P'N\\BP''&z_2\mathbf 1\end{pmatrix}
 =\begin{pmatrix}z_1\mathbf 1&P'N\\BP''&z_2\mathbf 1\end{pmatrix}.
\]
With $z_1z_2$ and $M$ they form a split family of depth $2D_1$ and size $s+n_2$ on
$(\mathcal C^2_{t,c}(e),R'^2(e))$. Since $4D_1<c$, Lemma~\ref{pd-kappa-cloud-upper} gives
$s+n_2\leq\mu(\lambda^2)$, so $r=n_1-s\geq n_1+n_2-\mu(\lambda^2)$.

Put $F_i=\sum_aP_{ia}x^1_a\in\mathcal C^1_{t,c}(e)$ and $g_j=\sum_bQ_{bj}y^2_b\in\mathcal
C'^2_{t,c}(e)$ for $i,j\leq r$. Then $\langle\theta(F_i),g_j\rangle^2=(PNQ)_{ij}=z_1z_2\delta_{ij}$.
The vectors $U_b(F_i)$ lie in $\mathcal C_{t,c}(e)\subset\mathcal C_{t,c'}(e)$ by (c), and the forms
$p^{h-h_{b'}}p^{A^2-A}(g_j\circ L_{b'})$ lie in $\mathcal C'_{t,c'}(e)$ by (d). By the identity above
and (a), the pairing of $U_b(F_i)$ with $p^{h-h_{b'}}p^{A^2-A}(g_j\circ L_{b'})$ is
$\delta_{bb'}p^{h}\langle\theta(F_i),g_j\rangle^2=\delta_{bb'}\delta_{ij}p^hz_1z_2$. With $p^hz_1z_2$,
of valuation at most $h+2D_1$, and the identity matrix, these $n_br$ vectors and forms are a split
family of the required kind.
\end{proof}

\begin{lemma}\label{pd-kappa-weight-change}
Let $p$ and $(W,Q,\sigma,\iota)$ be as in Definition~\ref{pd-kappa-split-rank}, $m\geq2$ and
$k=m+1$, $T\subset\mathbf Z/n'\mathbf Z$, $c_{cl}\geq0$ an integer and
$a_w,b_w\colon\mathbf Z/n'\to\mathbf Z$. Let $\lambda_0(s)\in\mathbf Z^k$ be dominant weights with
$\lambda_0(s)_1=a_w(s)$, $\lambda_0(s)_k=-b_w(s)$ and $|\lambda_0(s)_i|\leq c_{cl}$ for $1<i<k$ when
$s\in T$, and $|\lambda_0(s)_i|\leq c_{cl}$ for all $i$ when $s\notin T$.
\begin{enumerate}
\item Let $t,c_3\geq0$ be integers with $2n'c_{cl}\leq c_3$, and $D,\kappa\in\mathbf R$. If
  $\mathrm{ER}(\lambda_0^T;t,c_3,D,(1-\kappa)\mu(\lambda_0^T))$ holds, then
  $\mathrm{ER}(\lambda_0;t+2n'c_{cl},c_3-2n'c_{cl},D,(1-\kappa)\mu(\lambda_0))$ holds.
\item Suppose that $\lambda_0(s)=0$ for $s\notin T$. Let $a\geq0$ be an integer and $b'>0$ a real,
  and suppose that $A=a_w-c_{cl}$ and $B=b_w-c_{cl}$ take values at least $b'/3$ on $T$, that
  $2c_{cl}\leq8a$ and that $12|T|kc_{cl}/b'\leq1/9$. Let $t,c_2\geq0$ be integers and
  $D_1,\kappa_1,\kappa_2$ reals with $0\leq D_1$, $4D_1<c_2$, $2|T|c_{cl}\leq c_2$,
  $0\leq\kappa_1\leq1$ and $0\leq\kappa_2$. If
  $\mathrm{ER}(\phi(A,B);t,c_2,D_1,(1-\kappa_1)\mu(\phi(A,B)))$ and
  $\mathrm{ER}(\phi(A+2c_{cl},B+2c_{cl});t,c_2,D_1,(1-\kappa_2)\mu(\phi(A+2c_{cl},B+2c_{cl})))$ hold,
  then
  \[
   \mathrm{ER}\Big(\lambda_0;t,c_2-2|T|c_{cl},2D_1+|T|\lceil m^2\ell\rceil,
   \Big(1-\kappa_1-2\kappa_2-\frac{12|T|kc_{cl}}{b'}\Big)\mu(\lambda_0)\Big)
  \]
  holds, where $\ell=\log_p(8a+m)+1$.
\end{enumerate}
\end{lemma}
\leanref{Formalism/SharpRestriction/GrowthBounds/ColumnSpanDensity/RhoLowerBoundLargeConductor.lean}{Formalism.MainNotions.flipOneThree.Pf.ResidueClassCount}
\leanref{Formalism/SharpRestriction/GrowthBounds/ColumnSpanDensity/TensorBoundOfOrderBound.lean}{Formalism.MainNotions.flipOneThree.Pf.tensor_bound_of_order_bound}
\leanref{Formalism/SharpRestriction/GrowthBounds/ColumnSpanDensity/RhoLowerBoundLargeConductor.lean}{Formalism.MainNotions.flipOneThree.Pf.CloudWeightLowerBound}
\leanref{Formalism/SharpRestriction/GrowthBounds/ColumnSpanDensity/CloudPropOfMapsBasis.lean}{Formalism.MainNotions.flipOneThree.Pf.cloud_bound_of_maps}
\leanref{Formalism/SharpRestriction/GrowthBounds/ColumnSpanDensity/CloudPropOfMapsRetractionBasis.lean}{Formalism.MainNotions.flipOneThree.Pf.cloud_prop_of_maps_retraction_basis}
\leanref{Formalism/SharpRestriction/GrowthBounds/ColumnSpanDensity/CloudBoundOfPropAndWeylDimension.lean}{Formalism.MainNotions.flipOneThree.Pf.cloud_bound_of_prop_and_weyl_dimension}
\leanref{Formalism/SharpRestriction/GrowthBounds/ColumnSpanDensity/CloudBoundOfPropAndWeylDimension.lean}{Formalism.MainNotions.flipOneThree.Pf.share_abs}

\begin{proof}
(1) Let $\lambda_0^O=\lambda_0-\lambda_0^T$ and use the notation of
Lemma~\ref{pd-kappa-weight-tensor}. The weights $\lambda_0^T$ and $\lambda_0^O$ are dominant, and
every entry of $\lambda_0^O$ has absolute value at most $c_{cl}$. Put $c=2n'c_{cl}$ and
$d_O=\dim V^O$. Let $e$ be an edge with source $L''$ and target $L$, and let $x_1,\dots,x_n$,
$y_1,\dots,y_n$, $z$ and $(M_{lj})$ be a split family of depth $D$ and size
$n\geq(1-\kappa)\mu(\lambda_0^T)$ on the cores of $\lambda_0^T$ at $e$ formed with $t$ and $c_3$. Let
$x^O_i$ and $\pi_i$ be given by Lemma~\ref{pd-edge-basis}(1) for $\lambda_0^O$ and $e$, and put
$r_i=\varpi_0^{\pi_i-A^O}x^O_i$ and $\rho_i=\varpi_0^{\pi_i-A^O}p^{-\pi_i}(x^O_i)^*$ for $i\leq d_O$.
Then $r_i\in R^O(e)$, $\rho_i\in R'^O(e)$ and
$\langle r_i,\rho_j\rangle^O=\delta_{ij}p^{A^O}\varpi_0^{2\pi_i-2A^O}p^{-\pi_i}=\delta_{ij}$. By
Lemma~\ref{pd-kappa-weight-order}, $r_i\in\varpi_0^{-c}H^O(L'')$ and
$\rho_i\in\varpi_0^{-c}H^O(L)^\vee$.

The products $x_ar_i$ lie in the core $\mathcal C_{t+c,c_3-c}(e)$ of $\lambda_0$. Indeed,
Lemma~\ref{pd-edge-products} with $\lambda^\flat=\lambda_0^T$ and $\lambda^\sharp=\lambda_0^O$ gives
$x_ar_i\in\varpi_0^{-t-c}H(L'')$ and $x_ar_i\in R(e)$. For an edge $e'\neq e$ with source $L''$,
$x_a\in\varpi_0^{c_3}R^T(e')$ and $r_i\in\varpi_0^{-c}H^O(L'')\subset\varpi_0^{-c}R^O(e')$ by
Lemma~\ref{pd-edge-basis}(2), so $x_ar_i\in\varpi_0^{c_3-c}R(e')$. In the same way, with
Lemma~\ref{pd-kappa-weight-tensor}(3) in place of Lemma~\ref{pd-edge-products} and with
$H^O(L)^\vee\subset R'^O(e'')$ for the edges $e''$ with target $L$, the forms $y_l\otimes\rho_j$ lie
in $\mathcal C'_{t+c,c_3-c}(e)$. By Lemma~\ref{pd-kappa-weight-tensor}(3),
\[
 \sum_{l,j}\langle x_ar_i,y_l\otimes\rho_j\rangle M_{lb}\delta_{ji'}
 =\sum_l\langle x_a,y_l\rangle^TM_{lb}\,\delta_{ii'}=z\,\delta_{ab}\delta_{ii'} .
\]
So the $nd_O$ vectors $x_ar_i$ and forms $y_l\otimes\rho_j$, with $z$ and the matrix with entries
$M_{lb}\delta_{ji'}$, form a split family of depth $D$ on the cores of $\lambda_0$ at $e$ formed with
$t+c$ and $c_3-c$. By Lemma~\ref{pd-kappa-weight-tensor}(1), $\mu(\lambda_0)=\mu(\lambda_0^T)d_O$, so
its size satisfies $nd_O\geq(1-\kappa)\mu(\lambda_0)$.

(2) Let $e$ be an edge, and put $\phi=\phi(A,B)$, $\phi^+=\phi(A+2c_{cl},B+2c_{cl})$,
$h=|T|\lceil m^2\ell\rceil$ and $\epsilon=12|T|kc_{cl}/b'$. Since $A$ and $B$ are at least $b'/3>0$ on
$T$, Lemma~\ref{pd-kappa-cloud-maps} applies and gives $n_b$, $\theta$, the $U_b$ and $L_b$, and
integers $h_b\leq h$. With $c=c_2$ and $c'=c_2-2|T|c_{cl}$ they satisfy conditions (a) to (d) of
Lemma~\ref{pd-kappa-cloud-count} for $\lambda^1=\phi$ and $\lambda^2=\phi^+$. The two edge rank
conditions give split families at $e$ of depth $D_1$ and sizes $n_1\geq(1-\kappa_1)\mu(\phi)$ and
$n_2\geq(1-\kappa_2)\mu(\phi^+)$ on the cores formed with $t$ and $c_2$. Since $0\leq D_1$ and
$4D_1<c_2$, Lemma~\ref{pd-kappa-cloud-count} gives a split family of depth $2D_1+h$ and size $n_br$
on the cores of $\lambda_0$ at $e$ formed with $t$ and $c_2-2|T|c_{cl}$, where
$r\geq n_1+n_2-\mu(\phi^+)\geq(1-\kappa_1)\mu(\phi)-\kappa_2\mu(\phi^+)$. It remains to show
\[
 n_b\big((1-\kappa_1)\mu(\phi)-\kappa_2\mu(\phi^+)\big)\geq(1-\kappa_1-2\kappa_2-\epsilon)\mu(\lambda_0).
\]

Put $t_0=6c_{cl}/b'$, so that $2c_{cl}\leq t_0A(s)$ and $2c_{cl}\leq t_0B(s)$ for $s\in T$. For
$s\in T$ the weight $\lambda(s)$ has first entry $A(s)+c_{cl}$ and last entry $-B(s)-c_{cl}$ if $s$ is
even, first entry $B(s)+c_{cl}$ and last entry $-A(s)-c_{cl}$ if $s$ is odd, and middle entries of
absolute value at most $c_{cl}$. By the Weyl dimension formula \cite[Theorem~6.3]{FH91},
$\dim V=\prod_s\dim_k(\lambda(s))$, where the factors with $s\notin T$ are $1$, and
$\dim V^\phi=\prod_{s\in T}\dim_k(h(A(s),B(s)))$ in the notation of
Lemma~\ref{pd-kappa-cloud-dimensions}, since $\dim_k$ is unchanged by
$\mu\mapsto(-\mu_k,\dots,-\mu_1)$; the same holds for $\phi^+$. That lemma, applied at each $s\in T$
with $t=t_0$ and multiplied over $T$, gives, with $E=(1+t_0)^{(2m-1)|T|}$,
\[
 n_b\dim V^\phi\leq\dim V\leq En_b\dim V^\phi,\qquad\dim V^{\phi^+}\leq E\dim V^\phi .
\]
Now $(2m-1)|T|t_0\leq\epsilon$, so $E\leq e^{\epsilon}$ and $(1-\epsilon)E\leq(1-\epsilon)e^\epsilon\leq1$;
as $\epsilon\leq1/9$, also $E\leq2$. Dividing by $[k]$ gives
$(1-\epsilon)\mu(\lambda_0)\leq n_b\mu(\phi)\leq\mu(\lambda_0)$ and $\mu(\phi^+)\leq2\mu(\phi)$. Since
$0\leq\kappa_1\leq1$, $\kappa_2\geq0$ and $\epsilon\geq0$,
\[
 n_b\big((1-\kappa_1)\mu(\phi)-\kappa_2\mu(\phi^+)\big)\geq(1-\kappa_1)n_b\mu(\phi)-2\kappa_2n_b\mu(\phi)
 \geq(1-\kappa_1)(1-\epsilon)\mu(\lambda_0)-2\kappa_2\mu(\lambda_0),
\]
which is at least $(1-\kappa_1-2\kappa_2-\epsilon)\mu(\lambda_0)$. Since $n_b\geq0$, the size $n_br$
is at least this number, at every edge.
\end{proof}

\subsection{Product spaces and the flat defect}\label{subsection-pd-product-spaces-flat-defect}

\begin{definition}\label{pd-kappa-product-space}
Let $p\geq5$ be a prime, $(W,Q,\sigma,\iota)$ an unramified quadratic Hida setting of degree $n'$ at
$p$ with rings $\mathcal R\subset\mathcal K$ and $\varpi_0$, $m\geq1$ and $k=m+1$,
$T\subset\mathbf Z/n'\mathbf Z$, and $A,B\colon T\to\mathbf Z_{\geq0}$. For $s\in T$ put
$(d_s,d'_s)=(A(s),B(s))$ if $s$ is even and $(d_s,d'_s)=(B(s),A(s))$ if $s$ is odd. The
\emph{product space} $V^\times(A,B)$ is the $\mathcal K$-span, in the polynomial ring over
$\mathcal K$ on one $k\times k$ matrix $x_{s,\epsilon}$ of variables for each $s\in T$ and
$\epsilon\in\{0,1\}$, of the products $\prod_{s\in T}f_s(x_{s,1})g_s(x_{s,0})$ with $f_s$ in the dual
Weyl module over $\mathcal K$ of the natural weight of $(d_s,0,\dots,0)$ and $g_s$ in that of
$(0,\dots,0,-d'_s)$ (Definition~\ref{pd-weight-lattices}(1)). Its lattices $H^\times(L)$ and
$H^\times(e)$, the modules $R^\times(e)=\varpi_0^{-A^\times}H^\times(e)$ and
$R'^\times(e)=\varpi_0^{-A^\times}H^\times(e)^\vee$ with $A^\times=\sum_{s\in T}d_s$, the pairing
$\langle F,G\rangle^\times=G(p^{A^\times}F)$, the cores $\mathcal C^\times_{t,c}(e)$ and
$\mathcal C'^\times_{t,c}(e)$ and split families on them are defined as in
Definitions~\ref{pd-projection-property} and~\ref{pd-kappa-split-rank}, with $r=1$ and the $2|T|$
positions $(s,1)$ and $(s,0)$, both of which use $\iota\sigma^s$.
\end{definition}
\leanref{Formalism/SharpRestriction/GrowthBounds/ObjSeg.lean}{Formalism.MainNotions.G.hhDeg}
\leanref{Formalism/SharpRestriction/GrowthBounds/ObjSeg.lean}{Formalism.MainNotions.G.letterWt}
\leanref{Formalism/SharpRestriction/GrowthBounds/ObjSeg.lean}{Formalism.MainNotions.G.hhWt}
\leanref{Formalism/SharpRestriction/GrowthBounds/ObjSeg.lean}{Formalism.MainNotions.G.VHH}
\leanref{Formalism/SharpRestriction/GrowthBounds/ObjSeg.lean}{Formalism.MainNotions.G.HHL}
\leanref{Formalism/SharpRestriction/GrowthBounds/ObjSeg.lean}{Formalism.MainNotions.G.HHLe}
\leanref{Formalism/SharpRestriction/GrowthBounds/ObjSeg.lean}{Formalism.MainNotions.G.AHH}
\leanref{Formalism/SharpRestriction/GrowthBounds/ObjSeg.lean}{Formalism.MainNotions.G.ReHH}
\leanref{Formalism/SharpRestriction/GrowthBounds/ObjSeg.lean}{Formalism.MainNotions.G.beHH}
\leanref{Formalism/SharpRestriction/GrowthBounds/ObjSeg.lean}{Formalism.MainNotions.G.CCe}
\leanref{Formalism/SharpRestriction/GrowthBounds/ObjSeg.lean}{Formalism.MainNotions.G.CCe'}

\begin{lemma}\label{pd-kappa-dimensions}
Let $p$, $(W,Q,\sigma,\iota)$, $m\geq1$, $k=m+1$ and $T$ be as in
Definition~\ref{pd-kappa-product-space}. For integers $d\geq0$ put $P_m(d)=\binom{d+m}{m}$, and for
integers $\alpha,\beta\geq0$ put $\operatorname{Har}_m(\alpha,\beta)=P_m(\alpha)P_m(\beta)
-P_m(\alpha-1)P_m(\beta-1)$, where the subtracted product is $0$ if $\alpha=0$ or $\beta=0$. Let
$A,B\colon T\to\mathbf Z_{\geq0}$.
\begin{enumerate}
\item $\dim V^\times(A,B)=\prod_{s\in T}P_m(A(s))P_m(B(s))$.
\item $\dim V=\prod_{s\in T}\operatorname{Har}_m(A(s),B(s))$ for the space $V$ of the flat weights
  $\phi(A,B)$ (Definition~\ref{pd-kappa-split-rank}).
\item $\sum_{j=0}^{\min(\alpha,\beta)}\operatorname{Har}_m(\alpha-j,\beta-j)=P_m(\alpha)P_m(\beta)$
  for all integers $\alpha,\beta\geq0$.
\end{enumerate}
\end{lemma}
\leanref{Formalism/SharpRestriction/GrowthBounds/ColumnSpanDensity/MeanDefs.lean}{Formalism.MainNotions.flipOneThree.Pf.RestDims}
\leanref{Formalism/SharpRestriction/GrowthBounds/ColumnSpanDensity/FinrankLetter.lean}{Formalism.MainNotions.flipOneThree.Pf.restDims_of}
\leanref{Formalism/SharpRestriction/GrowthBounds/ProjectionLemmaBaseLevel.lean}{Formalism.MainNotions.flipOneThree.WeylV}
\leanref{Formalism/PolynomialDimension/DigitClassSize.lean}{Formalism.CountingSumsFrames.dimPol}
\leanref{Formalism/PolynomialDimension/DigitClassSize.lean}{Formalism.CountingSumsFrames.dimHar}
\leanref{Formalism/PolynomialDimension/DigitClassCount.lean}{Formalism.CountingSumsFrames.Top.levelCount}

\begin{proof}
The dual Weyl module over $\mathcal K$ of a dominant $\nu$ has the dimension given by Weyl's formula
$\prod_{i<j}(\nu_i-\nu_j+j-i)/(j-i)$ \cite[Theorem~6.3]{FH91}. For $(d,0,\dots,0)$ and for
$(d,\dots,d,0)$, the natural weights of $(d,0,\dots,0)$ and $(0,\dots,0,-d)$, this is $P_m(d)$. A
product of polynomials in disjoint sets of variables spans a tensor product, which gives (1). For
$s\in T$ the weight $\lambda(s)$ is $(A(s),0,\dots,0,-B(s))$ or $(B(s),0,\dots,0,-A(s))$, and Weyl's
formula for $(\alpha,0,\dots,0,-\beta)$ is $\operatorname{Har}_m(\alpha,\beta)$, a short
computation; positions outside $T$ carry the weight $0$ and contribute the factor $1$. This gives
(2). In (3) the sum telescopes: its terms are
$P_m(\alpha-j)P_m(\beta-j)-P_m(\alpha-j-1)P_m(\beta-j-1)$, and the last subtracted product is $0$.
\end{proof}

\begin{lemma}\label{pd-kappa-harmonic-minors}
Let $K$ be a field, $m\geq1$ and $k=m+1$, and let $Y$ be a $k\times k$ matrix of variables. For
$1\leq i\leq k$ let $c_i(Y)$ be $(-1)^{i-1}$ times the minor of $Y$ on the rows other than $i$ and the
first $m$ columns, and write $c(Y)$ for the column vector of the $c_i(Y)$. Dual Weyl modules are those
of Definition~\ref{pd-weight-lattices}(1), over $K$.
\begin{enumerate}
\item For $P\geq0$, the dual Weyl module of $(P,0,\dots,0)$ consists of the polynomials
  $a(Y_{11},\dots,Y_{k1})$ with $a\in K[x_1,\dots,x_k]$ homogeneous of degree $P$.
\item For $Q\geq0$, the dual Weyl module of $(Q,\dots,Q,0)$, the natural weight of
  $(0,\dots,0,-Q)$, consists of the polynomials $b(c(Y))$ with $b\in K[y_1,\dots,y_k]$ homogeneous
  of degree $Q$.
\item The polynomials $c_1(Y),\dots,c_k(Y)$ are algebraically independent over $K$.
\item $\sum_iY_{i1}c_i(Y)=0$.
\item For $g\in GL_k(K)$, $c(gY)=\det(g)\,{}^tg^{-1}c(Y)$.
\end{enumerate}
\end{lemma}
\leanref{Formalism/SharpRestriction/GrowthBounds/HermitianTensor/ShapeCovering/CovWeyl.lean}{Formalism.MainNotions.flipOneThree.Pf.ShapeCovering.first_column_polynomials_characterize}
\leanref{Formalism/SharpRestriction/GrowthBounds/HermitianTensor/ShapeCovering/CovWeyl.lean}{Formalism.MainNotions.flipOneThree.Pf.ShapeCovering.signed_minor_polynomials_characterize}
\leanref{Formalism/SharpRestriction/GrowthBounds/HermitianTensor/ShapeCovering/CovDet.lean}{Formalism.MainNotions.flipOneThree.Pf.ShapeCovering.signed_minors_alg_independent}
\leanref{Formalism/SharpRestriction/GrowthBounds/ColumnSpanDensity/DeltaPsiH.lean}{Formalism.MainNotions.flipOneThree.Pf.sum_covY}
\leanref{Formalism/SharpRestriction/GrowthBounds/HermitianTensor/ShapeCovering/CovDet.lean}{Formalism.MainNotions.flipOneThree.Pf.ShapeCovering.signed_minors_cofactor_transform}

\begin{proof}
By the standard basis theorem \cite{DEP80} (Lemma~\ref{is-standard-basis}(1)), the dual Weyl module of
a natural weight $\nu$ is spanned by the bideterminants of shape $\nu$: products of minors of $Y$,
the $l$-th on the first $s_l$ columns, where $\#\{l:s_l\geq t\}=\nu_t$ for every $t$. A minor on the
first $s$ columns lies in the dual Weyl module of $(1,\dots,1,0,\dots,0)$, with $s$ entries $1$, as in
the proof of Lemma~\ref{pd-kappa-cloud-cartan}(2), and a product of elements of the dual Weyl modules
of $\nu$ and $\nu'$ lies in that of $\nu+\nu'$, as in the proof of Lemma~\ref{pd-edge-products}.

(1) For $\nu=(P,0,\dots,0)$ the conditions say that $P$ of the $s_l$ are $1$ and the others $0$. So
the bideterminants are the products of $P$ entries of the first column, which span the polynomials
$a(Y_{11},\dots,Y_{k1})$ with $a$ homogeneous of degree $P$; conversely these lie in the module, since
each $Y_{i1}$ is a minor on the first column.

(2) For $\nu=(Q,\dots,Q,0)$ the conditions say that $Q$ of the $s_l$ are $m$ and the others $0$. A
minor on the first $m$ columns is $0$ if two of its rows coincide, and otherwise it is $\pm c_i(Y)$,
where $i$ is the row it omits. So the bideterminants are the products of $Q$ of the $\pm c_i(Y)$, and
they span the polynomials $b(c(Y))$ with $b$ homogeneous of degree $Q$; conversely these lie in the
module, since each $c_i(Y)$ is up to sign a minor on the first $m$ columns.

(3) Let $z_1,\dots,z_k$ be variables and $L$ the fraction field of $K[z_1,\dots,z_k]$. Let $N\in
GL_k(L)$ be the upper triangular matrix with first row $(z_1,\dots,z_k)$, entry $z_1^{-1}$ in place
$(2,2)$, entries $1$ in the places $(i,i)$ with $i\geq3$, and zeros elsewhere. Then $\det N=1$, so
$Q=N^{-1}$ has determinant $1$ and adjugate $N$. Substitute for the first $m$ columns of $Y$ the
columns $2,\dots,k$ of $Q$. Then $c_i(Y)$ becomes $(-1)^{i+1}$ times the minor of $Q$ on the rows other
than $i$ and the columns other than $1$, which is the entry $(1,i)$ of the adjugate of $Q$, that is
$z_i$. Since $z_1,\dots,z_k$ are algebraically independent over $K$, so are the $c_i(Y)$.

(4) The $k\times k$ matrix whose first column is the first column of $Y$ and whose other columns are
the first $m$ columns of $Y$ has two equal columns, so its determinant is $0$. Expanding it along the
first column gives $\sum_iY_{i1}c_i(Y)$.

(5) Let $\tilde Y$ be the matrix with first column $0$ whose columns $2,\dots,k$ are the first $m$
columns of $Y$. As in (3), $c_i(Y)$ is the entry $(1,i)$ of the adjugate $\operatorname{adj}(\tilde
Y)$, and the matrix formed in the same way from $gY$ is $g\tilde Y$. Since $\operatorname{adj}(g\tilde
Y)=\operatorname{adj}(\tilde Y)\operatorname{adj}(g)$ and $\operatorname{adj}(g)=\det(g)g^{-1}$,
$c_i(gY)=\det(g)\sum_jc_j(Y)(g^{-1})_{ji}$.
\end{proof}

\begin{lemma}\label{pd-kappa-harmonic-decomposition}
Let $K$ be a field of characteristic $0$ and $k\geq1$, and let $x=(x_1,\dots,x_k)$ and
$y=(y_1,\dots,y_k)$ be variables. For integers $A,B$ let $\mathrm{Pol}(A,B)\subset K[x,y]$ be the
space of polynomials homogeneous of degree $A$ in $x$ and $B$ in $y$, which is $0$ if $A<0$ or
$B<0$. Put $h=\sum_ix_iy_i$ and $\Delta=\sum_i\partial^2/\partial x_i\partial y_i$, and let
$\mathrm{Har}(A,B)$ be the kernel of $\Delta$ on $\mathrm{Pol}(A,B)$. For $A,B\geq0$ and
$0\leq j\leq\min(A,B)$ let $\pi_{A,B}$ and $\pi_{A,B,j}$ be the linear maps
\[
 \pi_{A,B}=\sum_{l=0}^{\min(A,B)}\frac{(-1)^l}{\binom{A+B+k-2}{l}}\,\frac{h^l}{l!}\,\frac{\Delta^l}{l!},
 \qquad
 \pi_{A,B,j}=\binom{A+B+k-j-1}{j}^{-1}\frac{h^j}{j!}\,\pi_{A-j,B-j}\,\frac{\Delta^j}{j!},
\]
where $h^l$ acts by multiplication.
\begin{enumerate}
\item $\Delta(hf)=h\Delta f+(A+B+k)f$ for $f\in\mathrm{Pol}(A,B)$.
\item Let $f\in\mathrm{Pol}(A,B)$. Then $\pi_{A,B}f\in\mathrm{Har}(A,B)$ and
  $f-\pi_{A,B}f\in h\,\mathrm{Pol}(A-1,B-1)$. If $f$ is harmonic, $\pi_{A,B}f=f$. If $A,B\geq1$, then
  $\pi_{A,B}(hG)=0$ for $G\in\mathrm{Pol}(A-1,B-1)$.
\item Let $f\in\mathrm{Pol}(A,B)$. Then $f=\sum_{j=0}^{\min(A,B)}\pi_{A,B,j}f$, and
  $\pi_{A,B,j}f=h^jw_j$ with $w_j\in\mathrm{Har}(A-j,B-j)$. For $i,j\leq\min(A,B)$ and
  $w\in\mathrm{Har}(A-i,B-i)$, $\pi_{A,B,j}(h^iw)$ is $h^iw$ if $i=j$ and $0$ otherwise.
\end{enumerate}
\end{lemma}
\leanref{Formalism/PolynomialDimension/StatementHar.lean}{Formalism.CountingSumsFrames.harProj}
\leanref{Formalism/PolynomialDimension/StatementHar.lean}{Formalism.CountingSumsFrames.levelProj}
\leanref{Formalism/PolynomialDimension/StatementHar.lean}{Formalism.CountingSumsFrames.SlTwo}
\leanref{Formalism/PolynomialDimension/StatementHar.lean}{Formalism.CountingSumsFrames.HarProj}
\leanref{Formalism/PolynomialDimension/StatementHar.lean}{Formalism.CountingSumsFrames.LevelProj}
\leanref{Formalism/PolynomialDimension/MultinomialDepth/HarmonicDecompositionHolds.lean}{Formalism.CountingSumsFrames.Har2.ctr_h}
\leanref{Formalism/PolynomialDimension/MultinomialDepth/HarmonicDecompositionHolds.lean}{Formalism.CountingSumsFrames.harmonicProjection_holds}
\leanref{Formalism/PolynomialDimension/MultinomialDepth/HarmonicDecompositionHolds.lean}{Formalism.CountingSumsFrames.isotypicProjection_holds}

\begin{proof}
(1) By the Leibniz rule,
$\partial^2(hf)/\partial x_i\partial y_i=f+x_i\,\partial f/\partial x_i+y_i\,\partial f/\partial y_i
+h\,\partial^2f/\partial x_i\partial y_i$. Summing over $i$ and using Euler's identities
$\sum_ix_i\,\partial f/\partial x_i=Af$ and $\sum_iy_i\,\partial f/\partial y_i=Bf$ gives (1).

(2) and (3). From (1), by induction on $l$, $\Delta(h^{l+1}G)=h^{l+1}\Delta G+(l+1)(a+b+k+l)h^lG$
for $G\in\mathrm{Pol}(a,b)$. Hence, for $w\in\mathrm{Har}(a,b)$ and $n\geq0$,
\[
 \Delta^n(h^iw)=i(i-1)\cdots(i-n+1)\cdot(a+b+k+i-1)(a+b+k+i-2)\cdots(a+b+k+i-n)\,h^{i-n}w,
 \tag{a}
\]
which is $0$ if $n>i$. Put $M=A+B+k-2$ and $c_l=(-1)^l/(\binom Ml\,l!^2)$, so that
$\pi_{A,B}=\sum_lc_lh^l\Delta^l$ and $c_l+(l+1)(M-l)c_{l+1}=0$. For $f\in\mathrm{Pol}(A,B)$, the
formula for $\Delta(h^{l+1}G)$ applied to $G=\Delta^{l+1}f$ and this recursion give, by induction on
$N$, $\Delta\big(\sum_{l\leq N}c_lh^l\Delta^lf\big)=c_Nh^N\Delta^{N+1}f$. For $N=\min(A,B)$,
$\Delta^{N+1}f=0$, so $\pi_{A,B}f$ is harmonic. Its term $l=0$ is $f$ and the others are multiples of
$h$ by elements of $\mathrm{Pol}(A-1,B-1)$, which gives $f-\pi_{A,B}f\in h\,\mathrm{Pol}(A-1,B-1)$, and
$\pi_{A,B}f=f$ if $\Delta f=0$. By induction on $A$, writing $f=\pi_{A,B}f+hG$, every
$f\in\mathrm{Pol}(A,B)$ is a sum $\sum_jh^jw_j$ with $w_j\in\mathrm{Har}(A-j,B-j)$.

Such a sum determines the $w_j$: if $\sum_{j<n}h^jw_j=0$ with $w_j\in\mathrm{Har}(a_j,b_j)$, apply
$\Delta^{n-1}$; by (a) the terms $j<n-1$ vanish and the last becomes a nonzero integer multiple of
$w_{n-1}$, so $w_{n-1}=0$, and induction gives the others. In particular a harmonic multiple of $h$ by
an element of some $\mathrm{Pol}(a,b)$ is $0$. For $G\in\mathrm{Pol}(A-1,B-1)$, $\pi_{A,B}(hG)$ is
harmonic and, by the above, equal to $hG-hG'$ with $G'\in\mathrm{Pol}(A-1,B-1)$; so it is $0$. This
proves (2).

For (3), $\Delta^j$ maps $\mathrm{Pol}(A,B)$ into $\mathrm{Pol}(A-j,B-j)$, so $\pi_{A,B,j}f=h^jw_j$ with
$w_j$ harmonic by (2). Let $w\in\mathrm{Har}(A-i,B-i)$. If $i<j$, $\Delta^j(h^iw)=0$ by (a). If $i>j$,
$\Delta^j(h^iw)$ is a multiple of $h\cdot h^{i-j-1}w$ with $h^{i-j-1}w\in\mathrm{Pol}(A-j-1,B-j-1)$,
which $\pi_{A-j,B-j}$ kills by (2). If $i=j$, (a) gives $\Delta^j(h^jw)/j!=j!\binom{A+B+k-j-1}jw$, and
$\pi_{A-j,B-j}w=w$, so $\pi_{A,B,j}(h^jw)=h^jw$. Applying this to the terms of $f=\sum_ih^iw_i$ gives
$\sum_j\pi_{A,B,j}f=\sum_ih^iw_i=f$.
\end{proof}

\begin{lemma}\label{pd-kappa-harmonic-fock}
Let $K$ be a field of characteristic $0$, $p$ a prime, $\mathcal O\subset K$ a subring in which every
integer prime to $p$ is invertible, $k\geq1$, and $x$, $y$, $\mathrm{Pol}(A,B)$, $h$ and $\Delta$ as in
Lemma~\ref{pd-kappa-harmonic-decomposition}. For the monomial $x^uy^v$ put
$u!=\prod_iu_i!$. Let $F\subset K[x,y]$ be the set of polynomials $\sum c_{u,v}x^uy^v$ with
$u!\,c_{u,v}\in\mathcal O$ for all $u,v$. For $g\in GL_k(K)$ let $F_g$ and $G_g$ be the sets of the
polynomials $f(gx,{}^tg^{-1}y)$ with $f\in F$, respectively $f\in\mathcal O[x,y]$. Write
$\lfloor\log_p0\rfloor=0$. Let $g\in GL_k(K)$ and $A,B\geq0$.
\begin{enumerate}
\item Let $f\in\mathrm{Pol}(A,B)$. If $f\in F_g$, then $A!\,f\in G_g$. If $f\in G_g$, then
  $p^{(k-1)\lfloor\log_pA\rfloor}f/A!\in F_g$.
\item $F_g$ is an $\mathcal O$-module, and $h^lf/l!$ and $\Delta^lf/l!$ lie in $F_g$ for every
  $f\in F_g$ and $l\geq0$.
\item If $0\leq n\leq N$ are integers with $N\geq1$, and $l\geq0$ with $l\leq n$, then
  $p^{\lfloor\log_pN\rfloor}/\binom nl\in\mathcal O$.
\end{enumerate}
\end{lemma}
\leanref{Formalism/PolynomialDimension/StatementHar.lean}{Formalism.CountingSumsFrames.FockLat}
\leanref{Formalism/PolynomialDimension/StatementHar.lean}{Formalism.CountingSumsFrames.FockGaussXY}
\leanref{Formalism/PolynomialDimension/MultinomialDepth/SqrtFockGaussLatticeComparisonHolds.lean}{Formalism.CountingSumsFrames.Har4.fockGaussLatticeComparison_holds}
\leanref{Formalism/PolynomialDimension/MultinomialDepth/DiagLaplacianDivPowMemSqrtFockLattice.lean}{Formalism.CountingSumsFrames.divPow_mem_fockLattice}
\leanref{Formalism/PolynomialDimension/MultinomialDepth/DiagLaplacianDivPowMemSqrtFockLattice.lean}{Formalism.CountingSumsFrames.laplacian_linearSubst}
\leanref{Formalism/PolynomialDimension/MultinomialDepth/DiagLaplacianDivPowMemSqrtFockLattice.lean}{Formalism.CountingSumsFrames.pow_log_div_choose_mem}

\begin{proof}
(3) By Legendre's formula $v_p(n!)=\sum_{i\geq1}\lfloor n/p^i\rfloor$, the $p$-adic valuation of
$\binom nl$ is $\sum_{i\geq1}\big(\lfloor n/p^i\rfloor-\lfloor l/p^i\rfloor-\lfloor(n-l)/p^i\rfloor\big)$.
Each term is $0$ or $1$, and it is $0$ when $p^i>n$. So $\binom nl=p^eu$ with $e\leq\lfloor\log_pN\rfloor$
and $u$ prime to $p$, and $p^{\lfloor\log_pN\rfloor}/\binom nl=p^{\lfloor\log_pN\rfloor-e}u^{-1}\in\mathcal O$.

(1) Write $f=f_0(gx,{}^tg^{-1}y)$. The substitution preserves the bidegree, so $f_0\in\mathrm{Pol}(A,B)$,
and $f\in F_g$, respectively $f\in G_g$, if and only if $f_0\in F$, respectively
$f_0\in\mathcal O[x,y]$. Let $c$ be the coefficient of a monomial $x^uy^v$ of $f_0$, so $\sum_iu_i=A$.
The multinomial coefficient $A!/u!$ is the product of the $k-1$ binomial coefficients
$\binom{u_1+\cdots+u_l}{u_l}$, $2\leq l\leq k$, so by (3) its $p$-adic valuation is at most
$(k-1)\lfloor\log_pA\rfloor$. If $u!\,c\in\mathcal O$, then $A!\,c=(A!/u!)\,u!\,c\in\mathcal O$. If
$c\in\mathcal O$, then $u!\,p^{(k-1)\lfloor\log_pA\rfloor}c/A!=p^{(k-1)\lfloor\log_pA\rfloor}(A!/u!)^{-1}c$
lies in $\mathcal O$.

(2) The set $F$ is an $\mathcal O$-module. Let $e_i$ be multiplication by $x_iy_i$ and
$d_i=\partial^2/\partial x_i\partial y_i$, and let $\epsilon_i\in\mathbf Z^k$ be the $i$-th unit vector. The
coefficient of $x^{u+n\epsilon_i}y^{v+n\epsilon_i}$ in $e_i^nf/n!$ is $c_{u,v}/n!$, and
$(u+n\epsilon_i)!\,c_{u,v}/n!=\binom{u_i+n}n\,u!\,c_{u,v}$. The coefficient of $x^uy^v$ in $d_i^nf/n!$ is
$c'=\frac{(u_i+n)!\,(v_i+n)!}{u_i!\,v_i!\,n!}\,c_{u+n\epsilon_i,v+n\epsilon_i}$, and
$u!\,c'=\binom{v_i+n}n(u+n\epsilon_i)!\,c_{u+n\epsilon_i,v+n\epsilon_i}$. So the operators $e_i^n/n!$ and
$d_i^n/n!$ map $F$ into itself. The $e_i$ commute with each other, and so do the $d_i$, so
$h^l/l!=(\sum_ie_i)^l/l!$ is the sum of the products $\prod_ie_i^{n_i}/n_i!$ over $n_1+\cdots+n_k=l$,
and similarly for $\Delta^l/l!$; hence both map $F$ into itself. Finally let
$\tau_g(f)=f(gx,{}^tg^{-1}y)$. Then $\tau_g(h)=h$, and $\Delta\tau_g=\tau_g\Delta$ by the chain rule,
since the substitution $x\mapsto gx$, $y\mapsto{}^tg^{-1}y$ preserves the pairing of the $x$-derivatives
with the $y$-derivatives. So $h^l/l!$ and $\Delta^l/l!$ also map $F_g=\tau_g(F)$ into itself.
\end{proof}

\begin{lemma}\label{pd-kappa-harmonic-model}
Let $p$ and $(W,Q,\sigma,\iota)$ be as in Definition~\ref{pd-kappa-split-rank}, with rings
$\mathcal R\subset\mathcal K$, let $m\geq2$ and $k=m+1$, and let $x=(x_1,\dots,x_k)$ and
$y=(y_1,\dots,y_k)$ be variables and $Y$ and $Y'$ be $k\times k$ matrices of variables. For integers
$A,B\geq0$ let $\mathrm{Pol}(A,B)\subset\mathcal K[x,y]$ be the polynomials homogeneous of degree $A$
in $x$ and $B$ in $y$, put $h=\sum_ix_iy_i$, and let $\mathrm{Har}(A,B)$ be the set of
$f\in\mathrm{Pol}(A,B)$ with $\sum_i\partial^2f/\partial x_i\partial y_i=0$. For $g\in GL_k(\mathcal K)$
let $G_g\subset\mathcal K[x,y]$ be the set of the polynomials $f(gx,{}^tg^{-1}y)$ with
$f\in\mathcal R[x,y]$. Let $\Psi\colon\mathcal K[x,y]\to\mathcal K[Y,Y']$ be the
$\mathcal K$-algebra map with $\Psi(x_i)=Y_{i1}$ and $\Psi(y_i)$ equal to $(-1)^{i-1}$ times the minor
of $Y'$ on the rows other than $i$ and the first $m$ columns, and let
$\delta\colon\mathcal K[Y,Y']\to\mathcal K[Y]$ substitute $Y$ for $Y'$. Let $V_{P,Q}$,
$V_{P,Q}(\Lambda)$, $\nabla_{\mathcal K}(\lambda)$ and $\nabla_\lambda(\Lambda)$ be as in
Lemma~\ref{pd-kappa-harmonic-single}, and $h(a,b)=(a,0,\dots,0,-b)\in\mathbf Z^k$.
\begin{enumerate}
\item $\Psi$ is injective, $\Psi(\mathrm{Pol}(P,Q))=V_{P,Q}$ for all $P,Q\geq0$, and $\delta\Psi$
  maps $\mathrm{Pol}(a,b)$ onto $\nabla_{\mathcal K}(h(a,b))$, with kernel $h\,\mathrm{Pol}(a-1,b-1)$
  if $a,b\geq1$ and $0$ otherwise.
\item For every $\mathcal R$-lattice $\Lambda\subset\mathcal K^k$ there is $g\in GL_k(\mathcal K)$ such
  that $V_{P,Q}(\Lambda)=\Psi(\mathrm{Pol}(P,Q)\cap G_g)$ and
  $\nabla_{h(a,b)}(\Lambda)=\delta\Psi(\mathrm{Pol}(a,b)\cap G_g)$ for all $P,Q,a,b\geq0$.
\item Every $f\in\mathrm{Pol}(A,B)$ can be written uniquely as
  $f=\sum_{j=0}^{\min(A,B)}h^jD_j(f)$ with $D_j(f)\in\mathrm{Har}(A-j,B-j)$.
\item If $f\in\mathrm{Pol}(A,B)\cap G_g$ for some $g\in GL_k(\mathcal K)$, then
  $p^Eh^jD_j(f)\in G_g$ for every $j\leq\min(A,B)$, where $E=(k+1)\lfloor\log_p(A+B+k)\rfloor$.
\end{enumerate}
\end{lemma}
\leanref{Formalism/SharpRestriction/GrowthBounds/ColumnSpanDensity/ModelIdentificationOfExplicitMap.lean}{Formalism.MainNotions.flipOneThree.Pf.TiePsi}
\leanref{Formalism/SharpRestriction/GrowthBounds/HermitianTensor/NuOfWtPQCast.lean}{Formalism.MainNotions.flipOneThree.Pf.tiePsi_of}
\leanref{Formalism/SharpRestriction/GrowthBounds/ColumnSpanDensity/SingleWeightOfModelIdentification.lean}{Formalism.MainNotions.flipOneThree.Pf.TwoLetterModelIdentification}
\leanref{Formalism/SharpRestriction/GrowthBounds/ColumnSpanDensity/ModelIdentificationOfExplicitMap.lean}{Formalism.MainNotions.flipOneThree.Pf.model_identification_of_explicit_map}
\leanref{Formalism/PolynomialDimension/StatementHar.lean}{Formalism.CountingSumsFrames.LevelProj}
\leanref{Formalism/PolynomialDimension/FinalHar.lean}{Formalism.CountingSumsFrames.Top.levelProj}
\leanref{Formalism/PolynomialDimension/StatementHar.lean}{Formalism.CountingSumsFrames.GaussProj}
\leanref{Formalism/PolynomialDimension/FinalHar.lean}{Formalism.CountingSumsFrames.Top.gaussProj}

\begin{proof}
Let $c(Y')$ be the column vector of the signed minors $c_i(Y')$ of
Lemma~\ref{pd-kappa-harmonic-minors}, so that $\Psi(x_i)=Y_{i1}$ and $\Psi(y_i)=c_i(Y')$. The natural
weights of $(P,0,\dots,0)$ and $(0,\dots,0,-Q)$ are $(P,0,\dots,0)$ and $(Q,\dots,Q,0)$.

(1) The map $\Psi$ is the composite of the map $\mathcal K[x,y]\to\mathcal K[x,Y']$ that fixes the
$x_i$ and sends $y_i$ to $c_i(Y')$, which is injective because the $c_i(Y')$ are algebraically
independent (Lemma~\ref{pd-kappa-harmonic-minors}(3)), and the renaming of each $x_i$ as $Y_{i1}$,
which is injective because the $Y_{i1}$ are distinct variables not among those of $Y'$. So $\Psi$ is
injective. The space $\mathrm{Pol}(P,Q)$ is spanned by the products $a(x)b(y)$ with $a$ and $b$
homogeneous of degrees $P$ and $Q$, and $\Psi(a(x)b(y))=a(Y_{11},\dots,Y_{k1})\,b(c(Y'))$. By
Lemma~\ref{pd-kappa-harmonic-minors}(1) and~(2) these are the products $f(Y)g(Y')$ that span
$V_{P,Q}$; so $\Psi(\mathrm{Pol}(P,Q))=V_{P,Q}$. The map $\delta$ sends $f(Y)g(Y')$ to $f(Y)g(Y)$, and
the product of elements of $\nabla_{\mathcal K}(a,0,\dots,0)$ and $\nabla_{\mathcal K}(0,\dots,0,-b)$
lies in $\nabla_{\mathcal K}(h(a,b))$, as in the proof of Lemma~\ref{pd-edge-products}, since the
natural weights add. So $\delta\Psi$ maps $\mathrm{Pol}(a,b)$ into $\nabla_{\mathcal K}(h(a,b))$. The
image is all of it: by Lemma~\ref{pd-kappa-cloud-cartan}(1), $\nabla_{\mathcal K}(h(a,b))$ is spanned
by the lattice $\nabla_{h(a,b)}(\mathcal R^k)$, and by the proof of (2) below with
$\Lambda=\mathcal R^k$, every element of this lattice is $\delta\Psi(f)$ with $f\in\mathrm{Pol}(a,b)$.
By Lemma~\ref{pd-kappa-harmonic-minors}(4), $\delta\Psi(h)=\sum_iY_{i1}c_i(Y)=0$, so
$h\,\mathrm{Pol}(a-1,b-1)$ lies in the kernel. In the notation of Lemma~\ref{pd-kappa-dimensions},
$\mathrm{Pol}(a,b)$ has dimension $P_m(a)P_m(b)$, and $\nabla_{\mathcal K}(h(a,b))$ has dimension
$\operatorname{Har}_m(a,b)$ by Weyl's formula, as in the proof of that lemma. So the kernel has
dimension $P_m(a-1)P_m(b-1)$ if $a,b\geq1$ and $0$ otherwise. Multiplication by $h$ is injective, so
$h\,\mathrm{Pol}(a-1,b-1)$ has the same dimension and is the kernel.

(2) Let $\Lambda^\vee=\{\phi\in\mathcal K^k:\phi\cdot v\in\mathcal R\text{ for all }v\in\Lambda\}$, let
$g_0\in GL_k(\mathcal K)$ be a matrix whose columns form an $\mathcal R$-basis of $\Lambda^\vee$, and put
$g=g_0^{-1}$. By Definition~\ref{pd-projection-property}(1) and Lemma~\ref{pd-kappa-cloud-cartan}(1),
applied to $\Lambda^\vee$, a polynomial $F$ lies in $\nabla_\lambda(\Lambda)$ if and only if
$\det(g_0)^{\lambda_k}F(g_0Y)$ lies in the dual Weyl module over $\mathcal R$ of the natural weight of
$\lambda$; and that module consists of the elements of $\nabla_{\mathcal K}(\lambda)$ with
coefficients in $\mathcal R$, as in the proof of Lemma~\ref{pd-kappa-cloud-cartan}(1). We use three
facts.

First, $F\in\nabla_{(P,0,\dots,0)}(\Lambda)$ if and only if
$F(Y)=a\big((g_0^{-1}Y)_{11},\dots,(g_0^{-1}Y)_{k1}\big)$ with $a\in\mathcal R[x]$ homogeneous of
degree $P$. Indeed, the condition is that $F(g_0Y)$ lies in $\nabla_{\mathcal K}(P,0,\dots,0)$ and has
coefficients in $\mathcal R$. By Lemma~\ref{pd-kappa-harmonic-minors}(1), $F(g_0Y)=a(Y_{11},\dots,Y_{k1})$
with $a$ homogeneous of degree $P$, and its coefficients are those of $a$.

Second, $G\in\nabla_{(0,\dots,0,-Q)}(\Lambda)$ if and only if $G(Y)=\det(g_0)^Qb(c(g_0^{-1}Y))$ with
$b\in\mathcal R[y]$ homogeneous of degree $Q$. Indeed, the condition is that
$\det(g_0)^{-Q}G(g_0Y)$ lies in $\nabla_{\mathcal K}(0,\dots,0,-Q)$ and has coefficients in
$\mathcal R$. By Lemma~\ref{pd-kappa-harmonic-minors}(2), $\det(g_0)^{-Q}G(g_0Y)=b(c(Y))$ with $b$
homogeneous of degree $Q$. If $b$ has coefficients in $\mathcal R$, so does $b(c(Y))$, because the
$c_i$ have integer coefficients. Conversely, suppose that $b(c(Y))$ has coefficients in $\mathcal R$,
and let $n\geq0$ be least with $\beta=\varpi_0^nb\in\mathcal R[y]$. If $n\geq1$, then $\beta(c(Y))\in
\varpi_0\mathcal R[Y]$, so the reduction $\bar\beta$ of $\beta$ modulo $\varpi_0$ vanishes on the
signed minors over the residue field. By Lemma~\ref{pd-kappa-harmonic-minors}(3) over the residue
field, $\bar\beta=0$, so $\varpi_0^{n-1}b\in\mathcal R[y]$, against the choice of $n$. So $n=0$.

Third, let $a\in\mathcal K[x]$, and $b\in\mathcal K[y]$ homogeneous of degree $Q$.
Lemma~\ref{pd-kappa-harmonic-minors}(5) with $g_0^{-1}$ gives ${}^tg_0\,c(Y')=\det(g_0)\,c(g_0^{-1}Y')$,
so
\[
 \Psi\big(a(gx)\,b({}^tg^{-1}y)\big)=a\big((g_0^{-1}Y)_{11},\dots,(g_0^{-1}Y)_{k1}\big)\cdot
 \det(g_0)^Q\,b\big(c(g_0^{-1}Y')\big).
\]

Let $f\in\mathrm{Pol}(P,Q)\cap G_g$, and write $f=f_0(gx,{}^tg^{-1}y)$ with $f_0\in\mathcal R[x,y]$;
then $f_0\in\mathrm{Pol}(P,Q)$, because the substitution preserves the bidegree. Expanding $f_0$ in
monomials $x^uy^v$ with coefficients in $\mathcal R$, the three facts show that $\Psi(f)\in
V_{P,Q}(\Lambda)$. Conversely, by the first two facts and the third, every product $F(Y)G(Y')$ that
spans $V_{P,Q}(\Lambda)$ is $\Psi(a(gx)b({}^tg^{-1}y))$ with $a\in\mathcal R[x]$ and $b\in\mathcal R[y]$
homogeneous of degrees $P$ and $Q$. So every element of $V_{P,Q}(\Lambda)$ is $\Psi(f_1)$ with
$f_1\in\mathrm{Pol}(P,Q)\cap G_g$, and $V_{P,Q}(\Lambda)=\Psi(\mathrm{Pol}(P,Q)\cap G_g)$. By
Lemma~\ref{pd-kappa-cloud-cartan}(2), applied to $\Lambda^\vee$ and the weights $(a,0,\dots,0)$ and
$(0,\dots,0,-b)$, whose sum is $h(a,b)$, the lattice $\nabla_{h(a,b)}(\Lambda)$ is the $\mathcal R$-span
of the products $F(Y)G(Y)$ with $F\in\nabla_{(a,0,\dots,0)}(\Lambda)$ and
$G\in\nabla_{(0,\dots,0,-b)}(\Lambda)$. These are the images under $\delta$ of the products
$F(Y)G(Y')$ that span $V_{a,b}(\Lambda)$. Hence
$\nabla_{h(a,b)}(\Lambda)=\delta(V_{a,b}(\Lambda))=\delta\Psi(\mathrm{Pol}(a,b)\cap G_g)$.

(3) Let $\pi_{A,B}$ and $\pi_{A,B,j}$ be the maps of Lemma~\ref{pd-kappa-harmonic-decomposition}. By part
(3) of that lemma, $f=\sum_j\pi_{A,B,j}f$ with $\pi_{A,B,j}f=h^jw_j$ and $w_j\in\mathrm{Har}(A-j,B-j)$.
If also $f=\sum_jh^jw'_j$ with $w'_j\in\mathrm{Har}(A-j,B-j)$, applying $\pi_{A,B,i}$ and the last
claim of that part gives $h^iw_i=h^iw'_i$, so $w_i=w'_i$. Thus $D_j(f)=w_j$ and
$h^jD_j(f)=\pi_{A,B,j}f$; in particular $D_0(f)=\pi_{A,B}f$, since $\pi_{A,B,0}=\pi_{A,B}$.

(4) The ring $\mathcal R$ is a discrete valuation ring with residue field of characteristic $p$, so
every integer prime to $p$ is invertible in $\mathcal R$. Let $F_g$ be as in
Lemma~\ref{pd-kappa-harmonic-fock} with $\mathcal O=\mathcal R$, let $f\in\mathrm{Pol}(A,B)\cap G_g$,
and put $N=A+B+k$ and $L=\lfloor\log_pN\rfloor$. By part (1) of that lemma,
$f'=p^{(k-1)\lfloor\log_pA\rfloor}f/A!$ lies in $F_g$. For $a,b\geq0$ with $a+b+k\leq N$, each term of
$p^L\pi_{a,b}$ is $(-1)^lp^L/\binom{a+b+k-2}l$ times the composite of $\Delta^l/l!$ and $h^l/l!$; by
parts (2) and (3) of that lemma, $p^L\pi_{a,b}$ maps $F_g$ into itself. The map $p^{2L}\pi_{A,B,j}$ is
$p^L/\binom{A+B+k-j-1}j$ times the composite of $\Delta^j/j!$, $p^L\pi_{A-j,B-j}$ and $h^j/j!$, so it
also maps $F_g$ into itself. Hence $p^{2L}\pi_{A,B,j}f'\in F_g\cap\mathrm{Pol}(A,B)$, and part (1) of
that lemma gives
$A!\,p^{2L}\pi_{A,B,j}f'=p^{(k-1)\lfloor\log_pA\rfloor+2L}\,h^jD_j(f)\in G_g$. Since
$(k-1)\lfloor\log_pA\rfloor+2L\leq(k+1)L=E$ and $G_g$ is an $\mathcal R$-module,
$p^Eh^jD_j(f)\in G_g$.
\end{proof}
% Lean: (1) and (2) are TiePsi = tiePsi_of D explicit_map_injective
%   explicit_map_onto_two_letter_space two_letter_lattice_gauss_inverse CartanOnto
%   (ColumnSystems.adm3 D), through tie3, tie4 and tie5b (HermitianTensor/NuOfWtPQCast.lean), with
%   g = g0^{-1} for Lambda the column span of g0; model_identification_of_explicit_map repackages
%   them as TwoLetterModelIdentification (alpha = Psi^{-1}, beta = deltaB o PsiXY). (3) is LevelProj
%   (levelProj = h^j D_j, uniqueness from its fourth clause) with harProj = levelProj at j = 0
%   (HarProj); (4) is GaussProj = gaussLatticeProjection_holds = Har4.gaussProj_of
%   fockLatticeProjection_holds, through gauss_dual_of_fock, pow_log_smul_harmonicProj_mem and
%   pow_log_smul_isotypicProj_mem.

\begin{lemma}\label{pd-kappa-harmonic-gauss}
Let $K$ be a field, $\mathcal O\subset K$ a subring and $k\geq1$, let $x=(x_1,\dots,x_k)$ and
$y=(y_1,\dots,y_k)$ be variables, and put $h=\sum_ix_iy_i$. For $g\in GL_k(K)$ let $G_g\subset K[x,y]$
be the set of the polynomials $f(gx,{}^tg^{-1}y)$ with $f\in\mathcal O[x,y]$. Let $g\in GL_k(K)$ and
$j\geq0$.
\begin{enumerate}
\item If $G\in K[x,y]$ and $h^jG\in G_g$, then $G\in G_g$.
\item Let $a,b\geq0$, and let $w\in K[x,y]$ be homogeneous of degree $a$ in $x$ and $b$ in $y$ with
  $\sum_i\partial^2w/\partial x_i\partial y_i=0$. Then $h^jw\in G_g$ if and only if $w\in G_g$.
\end{enumerate}
\end{lemma}
\leanref{Formalism/PolynomialDimension/StatementHar.lean}{Formalism.CountingSumsFrames.GaussLemmaH}
\leanref{Formalism/PolynomialDimension/StatementHar.lean}{Formalism.CountingSumsFrames.LevelLattice}
\leanref{Formalism/PolynomialDimension/MultinomialDepth/SqrtFockGaussLatticeComparisonHolds.lean}{Formalism.CountingSumsFrames.Har4.gaussLemmaPairing_holds}
\leanref{Formalism/PolynomialDimension/MultinomialDepth/SqrtFockGaussLatticeComparisonHolds.lean}{Formalism.CountingSumsFrames.Har4.isotypicLattice_holds}
\leanref{Formalism/PolynomialDimension/MultinomialDepth/SqrtFockGaussLatticeComparisonHolds.lean}{Formalism.CountingSumsFrames.Har4.coeff_mem_of_h_mul}

\begin{proof}
Let $\tau_g(f)=f(gx,{}^tg^{-1}y)$. It is a ring automorphism of $K[x,y]$ with inverse $\tau_{g^{-1}}$,
and $\tau_g(h)=h$, since $\sum_i(gx)_i({}^tg^{-1}y)_i={}^tx\,{}^tg\,{}^tg^{-1}y=h$. So $F\in G_g$ if and
only if $\tau_{g^{-1}}(F)\in\mathcal O[x,y]$, and $\tau_{g^{-1}}(h^jF)=h^j\tau_{g^{-1}}(F)$.

(1) By the above it suffices to show that $G\in\mathcal O[x,y]$ when $h^jG\in\mathcal O[x,y]$, and by
induction on $j$, applied to $h^{j-1}(hG)$, it suffices to treat $j=1$. Suppose that
$hG\in\mathcal O[x,y]$ and that some coefficient of $G$ is not in $\mathcal O$. Among the monomials
whose coefficient in $G$ is not in $\mathcal O$, choose one, $\mu$, with the largest exponent of
$x_1$, and put $M=x_1y_1\mu$. The coefficient of $M$ in $hG$ is the sum over $i$ of the coefficients
of $M$ in $x_iy_iG$. For $i=1$ this is the coefficient of $\mu$ in $G$. For $i\neq1$ it is $0$ if
$x_iy_i$ does not divide $M$, and otherwise it is the coefficient in $G$ of $M/(x_iy_i)$, whose
exponent of $x_1$ is one more than that of $\mu$; by the choice of $\mu$, it lies in $\mathcal O$.
Since the coefficient of $M$ in $hG$ lies in $\mathcal O$, so does that of $\mu$ in $G$, which is a
contradiction.

(2) If $h^jw\in G_g$, then $w\in G_g$ by (1). Conversely, if $w\in G_g$, then
$\tau_{g^{-1}}(h^jw)=h^j\tau_{g^{-1}}(w)$ is a product of polynomials with coefficients in
$\mathcal O$, so $h^jw\in G_g$.
\end{proof}

\begin{lemma}\label{pd-kappa-harmonic-single}
Let $p$ and $(W,Q,\sigma,\iota)$ be as in Definition~\ref{pd-kappa-split-rank}, with rings
$\mathcal R\subset\mathcal K$, let $m\geq2$ and $a,P,Q\geq0$ be integers with $P+Q+1\leq8a$, and put
$u_0=\lceil(m+2)(\log_p(8a+m)+1)\rceil$. Let $Y$ and $Y'$ be $(m+1)\times(m+1)$ matrices of
variables. For a dominant weight $\lambda\in\mathbf Z^{m+1}$ let $\nabla_{\mathcal K}(\lambda)$ be the
dual Weyl module over $\mathcal K$ of the natural weight of $\lambda$, and for an
$\mathcal R$-lattice $\Lambda\subset\mathcal K^{m+1}$ let $\nabla_\lambda(\Lambda)$ be its lattice of
Definition~\ref{pd-weight-lattices}(5), formed over $\mathcal R$. Let $V_{P,Q}\subset\mathcal K[Y,Y']$
be the span of the products $f(Y)g(Y')$ with $f\in\nabla_{\mathcal K}(P,0,\dots,0)$ and
$g\in\nabla_{\mathcal K}(0,\dots,0,-Q)$, and $V_{P,Q}(\Lambda)$ the $\mathcal R$-span of those with
$f\in\nabla_{(P,0,\dots,0)}(\Lambda)$ and $g\in\nabla_{(0,\dots,0,-Q)}(\Lambda)$. For
$0\leq i\leq\min(P,Q)$ put $\lambda_i=(P-i,0,\dots,0,-(Q-i))$, a weight in the variables $Y$. Then
there are $\mathcal K$-linear maps $\eta_i\colon\nabla_{\mathcal K}(\lambda_i)\to V_{P,Q}$ and
$\pi_i\colon V_{P,Q}\to\nabla_{\mathcal K}(\lambda_i)$, $0\leq i\leq\min(P,Q)$, with
$\sum_i\eta_i\pi_i=\mathrm{id}$, such that
$p^{u_0}\pi_i(V_{P,Q}(\Lambda))\subset\nabla_{\lambda_i}(\Lambda)$ and
$p^{u_0}\eta_i(\nabla_{\lambda_i}(\Lambda))\subset V_{P,Q}(\Lambda)$ for every $\mathcal R$-lattice
$\Lambda\subset\mathcal K^{m+1}$ and every $i$.
\end{lemma}
\leanref{Formalism/SharpRestriction/GrowthBounds/ColumnSpanDensity/SingleWeightLatticeComparison.lean}{Formalism.MainNotions.flipOneThree.Pf.SingleWeightLatticeComparison}
\leanref{Formalism/SharpRestriction/GrowthBounds/ColumnSpanDensity/SingleWeightOfModelIdentification.lean}{Formalism.MainNotions.flipOneThree.Pf.single_weight_of_model_identification}
\leanref{Formalism/SharpRestriction/GrowthBounds/ColumnSpanDensity/HarProjEqZeroIff.lean}{Formalism.MainNotions.flipOneThree.Pf.oneTop_abs}
\leanref{Formalism/SharpRestriction/GrowthBounds/ColumnSpanDensity/SingleWeightOfModelIdentification.lean}{Formalism.MainNotions.flipOneThree.Pf.expo_le}

\begin{proof}
Let $k=m+1$, and let $x$, $y$, $\mathrm{Pol}(A,B)$, $h$, $\mathrm{Har}(A,B)$, $G_g$, $\Psi$, $\delta$
and $D_j$ be as in Lemma~\ref{pd-kappa-harmonic-model}. Put
$E=(m+2)\lfloor\log_p(P+Q+m+1)\rfloor$. Since $P+Q+m+1\leq8a+m$, $E\leq u_0$. By
Lemma~\ref{pd-kappa-harmonic-model}(1), $\Psi$ restricts to an isomorphism from
$\mathrm{Pol}(P,Q)$ onto $V_{P,Q}$; let $\alpha$ be its inverse. For $i\leq\min(P,Q)$ let $\beta_i$
be the restriction of $\delta\Psi$ to $\mathrm{Pol}(P-i,Q-i)$; it maps onto
$\nabla_{\mathcal K}(\lambda_i)$, with kernel $h\,\mathrm{Pol}(P-i-1,Q-i-1)$ or $0$.

For $x'\in V_{P,Q}$ put $\pi_i(x')=\beta_i(D_i(\alpha(x')))$. Choose a linear section $\sigma_i$ of
$\beta_i$, and put $\eta_i(y')=\Psi(h^iD_0(\sigma_i(y')))$ for $y'\in\nabla_{\mathcal K}(\lambda_i)$;
it lies in $V_{P,Q}$ because $h^iD_0(\sigma_i(y'))\in\mathrm{Pol}(P,Q)$. If
$f\in\mathrm{Pol}(P-i,Q-i)$ and $\beta_i(f)=0$, then $f=hG$ with $G$ of bidegree $(P-i-1,Q-i-1)$, and
writing $G=\sum_jh^jD_j(G)$ shows, by the uniqueness in Lemma~\ref{pd-kappa-harmonic-model}(3), that
$D_0(f)=0$. Hence $\eta_i(\beta_i(f))=\Psi(h^iD_0(f))$ for every $f\in\mathrm{Pol}(P-i,Q-i)$.

Let $x'\in V_{P,Q}$ and $f=\alpha(x')$. Since $D_i(f)$ is harmonic, uniqueness gives
$D_0(D_i(f))=D_i(f)$, so $\eta_i(\pi_i(x'))=\Psi(h^iD_i(f))$. Summing over $i$,
$\sum_i\eta_i\pi_i(x')=\Psi\big(\sum_ih^iD_i(f)\big)=\Psi(f)=x'$.

Let $\Lambda$ be an $\mathcal R$-lattice and $g$ the matrix of Lemma~\ref{pd-kappa-harmonic-model}(2).
The set $G_g$ is an $\mathcal R$-module; it is the set $G_g$ of Lemma~\ref{pd-kappa-harmonic-gauss}
with $\mathcal O=\mathcal R$.

Let $x'\in V_{P,Q}(\Lambda)$. By Lemma~\ref{pd-kappa-harmonic-model}(2), $f=\alpha(x')\in G_g$, and by
part (4) of that lemma, $p^Eh^iD_i(f)\in G_g$. By Lemma~\ref{pd-kappa-harmonic-gauss}(1),
$p^ED_i(f)\in G_g$, so $p^ED_i(f)\in G_g\cap\mathrm{Pol}(P-i,Q-i)$, and
$p^E\pi_i(x')=\beta_i(p^ED_i(f))\in\nabla_{\lambda_i}(\Lambda)$ by part (2). As $E\leq u_0$,
$p^{u_0}\pi_i(x')\in\nabla_{\lambda_i}(\Lambda)$.

Let $y'\in\nabla_{\lambda_i}(\Lambda)$. By part (2), $y'=\beta_i(f)$ with
$f\in\mathrm{Pol}(P-i,Q-i)\cap G_g$. Part (4) with $j=0$ gives $p^{E_i}D_0(f)\in G_g$, where
$E_i=(m+2)\lfloor\log_p(P+Q-2i+m+1)\rfloor\leq E$; so $p^ED_0(f)\in G_g$. Since $p^ED_0(f)$ is
harmonic, Lemma~\ref{pd-kappa-harmonic-gauss}(2) gives $p^Eh^iD_0(f)\in G_g$, and
$p^Eh^iD_0(f)\in\mathrm{Pol}(P,Q)$. By part (2),
$p^E\eta_i(y')=\Psi(p^Eh^iD_0(f))\in V_{P,Q}(\Lambda)$, and so
$p^{u_0}\eta_i(y')\in V_{P,Q}(\Lambda)$.
\end{proof}

\begin{lemma}\label{pd-kappa-harmonic-maps}
Let $p$ and $(W,Q,\sigma,\iota)$ be as in Definition~\ref{pd-kappa-split-rank}, $m\geq2$ and
$k=m+1$, $T\subset\mathbf Z/n'\mathbf Z$, $a\geq0$ an integer and $A,B\colon T\to\mathbf Z_{\geq0}$
with $A(s)+B(s)+1\leq8a$ for $s\in T$. Put $u_0=\lceil(m+2)(\log_p(8a+m)+1)\rceil$ and $u=|T|u_0$.
Let $J$ be the set of functions $j\colon T\to\mathbf Z_{\geq0}$ with $j\leq\min(A,B)$, and for
$j\in J$ write $V_j$, $A_j$, $\mathcal C_{j,t,c}$ and $\mathcal C'_{j,t,c}$ for the space, integer and
cores of Definitions~\ref{pd-projection-property} and~\ref{pd-kappa-split-rank} formed with the flat
weights $\phi(A-j,B-j)$. Let $V^\times=V^\times(A,B)$, $A^\times$ and the cores
$\mathcal C^\times_{t,c}$ and $\mathcal C'^\times_{t,c}$ be as in
Definition~\ref{pd-kappa-product-space}. Then there are $\mathcal K$-linear maps
$\eta_j\colon V_j\to V^\times$ and $\pi_j\colon V^\times\to V_j$, $j\in J$, such that
$\sum_{j\in J}\eta_j\pi_j$ is the identity of $V^\times$ and, for all integers $t,c\geq0$, every edge $e$
and every $j\in J$:
\begin{enumerate}
\item $p^u\pi_j(x)\in\mathcal C_{j,t,c}(e)$ for every $x\in\mathcal C^\times_{t,c}(e)$;
\item $p^{2u+A^\times-A_j}(\Phi\circ\eta_j)\in\mathcal C'_{j,t,c}(e)$ for every
  $\Phi\in\mathcal C'^\times_{t,c}(e)$.
\end{enumerate}
\end{lemma}
\leanref{Formalism/SharpRestriction/GrowthBounds/ColumnSpanDensity/RhoSumBoundOfMapsSplitting.lean}{Formalism.MainNotions.flipOneThree.Pf.RhoSumLatticeMaps}
\leanref{Formalism/SharpRestriction/GrowthBounds/ColumnSpanDensity/LatticeMapsOfFunctorComparison.lean}{Formalism.MainNotions.flipOneThree.Pf.lattice_maps_of_functor_comparison}
\leanref{Formalism/SharpRestriction/GrowthBounds/ColumnSpanDensity/LatticeMapsOfFunctorComparison.lean}{Formalism.MainNotions.flipOneThree.Pf.LatticeFunctorComparison}
\leanref{Formalism/SharpRestriction/GrowthBounds/ColumnSpanDensity/LatticeMapsOfFunctorComparison.lean}{Formalism.MainNotions.flipOneThree.Pf.Ar_le_AHH}
\leanref{Formalism/SharpRestriction/GrowthBounds/ColumnSpanDensity/FunctorComparisonOfSingleWeight.lean}{Formalism.MainNotions.flipOneThree.Pf.functor_comparison_of_single_weight}
\leanref{Formalism/SharpRestriction/GrowthBounds/ColumnSpanDensity/FunctorComparisonOfSingleWeightFacts.lean}{Formalism.MainNotions.flipOneThree.Pf.functor_comparison_of_single_weight_facts}

\begin{proof}
For $s\in T$ let $(d_s,d'_s)$ be as in Definition~\ref{pd-kappa-product-space}. Then
$d_s+d'_s+1\leq8a$, $\min(d_s,d'_s)=\min(A(s),B(s))$, and the weight $\lambda_j(s)$ formed from
$\phi(A-j,B-j)$ at $s$ is $(d_s-j(s),0,\dots,0,-(d'_s-j(s)))$; for $s\notin T$ it is $0$. Take the maps
$\eta^{(s)}_i$ and $\pi^{(s)}_i$ of Lemma~\ref{pd-kappa-harmonic-single} for $P=d_s$ and $Q=d'_s$,
with $Y=x_{s,1}$ and $Y'=x_{s,0}$ on the side of $V^\times$ and $Y=x_s$ on the side of $V_j$.

As in the proof of Lemma~\ref{pd-kappa-weight-tensor}, multiplication identifies $V^\times$ with the
tensor product over $s\in T$ of the spaces $V_{d_s,d'_s}$, and $V_j$ with the tensor product over
$s\in T$ of the $\nabla_{\mathcal K}(\lambda_j(s))$, the positions outside $T$ contributing the
constants. Let $\eta'_j$ and $\pi'_j$ be the tensor products of the $\eta^{(s)}_{j(s)}$ and of the
$\pi^{(s)}_{j(s)}$. Then $\sum_{j\in J}\eta'_j\pi'_j$ is the tensor product of the identities
$\sum_i\eta^{(s)}_i\pi^{(s)}_i$, the identity. For a family $N=(N_s)$ of $\mathcal R$-lattices in
$\mathcal K^k$ let $H_j(N)$ be the $\mathcal R$-span of the products $\prod_sf_s(x_s)$ with
$f_s\in\nabla_{\lambda_j(s)}(N_s)$, and $H^\times(N)$ that of the products of elements of the
$V_{d_s,d'_s}(N_s)$, $s\in T$. Applying Lemma~\ref{pd-kappa-harmonic-single} factor by factor gives
\[
 p^u\pi'_j(H^\times(N))\subset H_j(N),\qquad p^u\eta'_j(H_j(N))\subset H^\times(N),
\]
since $u=|T|u_0$. Write $H_j(L)$, $H_j(e)$, $R_j(e)$ and $R'_j(e)$ for the lattices of
Definition~\ref{pd-projection-property} formed with $\phi(A-j,B-j)$. By
Definitions~\ref{pd-projection-property}(1) and~\ref{pd-kappa-product-space}, $H_j(L)$ and
$H^\times(L)$ are $H_j(N)$ and $H^\times(N)$ for the family $N$ formed from the $W$-lattice $L$, and
the same holds at an edge.

Put $d_j=A^\times-A_j$. Since $A_j=\sum_{s\in T}(d_s-j(s))$ and $A^\times=\sum_{s\in T}d_s$, we have
$d_j\geq0$. Let $\eta_j=\varpi_0^{-d_j}\eta'_j$ and $\pi_j=\varpi_0^{d_j}\pi'_j$; then
$\sum_j\eta_j\pi_j$ is the identity. Let $N$ be a family of lattices and $E$ an integer. If
$x\in\varpi_0^EH^\times(N)$, then $p^u\pi_j(x)\in\varpi_0^{E+d_j}H_j(N)$. If
$\Phi\in\varpi_0^EH^\times(N)^\vee$ and $y\in H_j(N)$, then
\[
 p^{2u+A^\times-A_j}\Phi(\eta_j(y))=\varpi_0^{d_j}p^u\,\Phi\big(p^u\eta'_j(y)\big)
 \in\varpi_0^{E+d_j}\mathcal R,
\]
since $p^{A^\times-A_j}=\varpi_0^{2d_j}$; so $p^{2u+A^\times-A_j}(\Phi\circ\eta_j)\in\varpi_0^{E+d_j}H_j(N)^\vee$.

Now let $x\in\mathcal C^\times_{t,c}(e)$, where $e$ has source $L''$ and target $L$. From
$x\in\varpi_0^{-t}H^\times(L'')$ we get $p^u\pi_j(x)\in\varpi_0^{-t+d_j}H_j(L'')\subset
\varpi_0^{-t}H_j(L'')$; from $x\in R^\times(e)=\varpi_0^{-A^\times}H^\times(e)$ we get
$p^u\pi_j(x)\in\varpi_0^{-A_j}H_j(e)=R_j(e)$; and from $x\in\varpi_0^{c-A^\times}H^\times(e')$ for an
edge $e'\neq e$ with source $L''$ we get $p^u\pi_j(x)\in\varpi_0^{c-A_j}H_j(e')=\varpi_0^cR_j(e')$.
This is (1). The same three steps with the forms, using $H(L)^\vee$ and $R'=\varpi_0^{-A}H(e)^\vee$,
give (2).
\end{proof}

\begin{lemma}\label{pd-kappa-harmonic-block}
Let $\mathcal O$ be a discrete valuation ring with maximal ideal $\mathfrak m$, $J$ a finite set,
$n\geq0$ an integer, $z\in\mathcal O$ nonzero, and $A_j$, $j\in J$, $n\times n$ matrices over
$\mathcal O$ with $\sum_jA_j=z\mathbf 1_n$. Then there are integers $\nu_j\geq0$ with
$\sum_j\nu_j\geq n$ and matrices $P_j$ and $Q_j$ over $\mathcal O$, of sizes $\nu_j\times n$ and
$n\times\nu_j$, with $P_jA_jQ_j=z\mathbf 1_{\nu_j}$ for every $j$.
\end{lemma}
\leanref{Formalism/SharpRestriction/GrowthBounds/ColumnSpanDensity/RhoSumBoundOfMapsSplitting.lean}{Formalism.MainNotions.flipOneThree.Pf.BlockSplit}
\leanref{Formalism/SharpRestriction/GrowthBounds/ColumnSpanDensity/BlockSplit.lean}{Formalism.MainNotions.flipOneThree.Pf.KP.blockSplit}

\begin{proof}
Fix $j$. By the elementary divisor theorem over $\mathcal O$ there are a basis $v_1,\dots,v_n$ of
$\mathcal O^n$ and nonzero $a_1,\dots,a_d\in\mathcal O$, $d\leq n$, such that $a_1v_1,\dots,a_dv_d$
is a basis of the image $A_j\mathcal O^n$. Let $I_j$ be the set of $i\leq d$ with $a_i\mid z$, and
$\nu_j=|I_j|$. For $i\in I_j$ choose $w_i$ with $A_jw_i=a_iv_i$, and put $w'_i=(z/a_i)w_i$, so that
$A_jw'_i=zv_i$. Let $Q_j$ have the columns $w'_i$, $i\in I_j$, and $P_j$ the rows given by the
coordinates along the $v_i$, $i\in I_j$, in the basis $(v_i)$. Then $P_jA_jQ_j=z\mathbf 1_{\nu_j}$.
For $i\notin I_j$ with $i\leq d$, $a_i$ has larger valuation than $z$, so $a_iv_i\in\mathfrak mz\mathcal O^n$.
Hence $A_j\mathcal O^n\subset W_j+\mathfrak mz\mathcal O^n$, where $W_j$ is spanned by the $v_i$,
$i\in I_j$.

Summing over $j$, $z\mathcal O^n=(\sum_jA_j)\mathcal O^n\subset W+\mathfrak mz\mathcal O^n$, where
$W=\sum_jW_j$ is generated by $\sum_j\nu_j$ elements. By Nakayama's lemma $z\mathcal O^n\subset W$.
Since $z\mathcal O^n$ is free of rank $n$, tensoring with the fraction field gives
$n\leq\sum_j\nu_j$.
\end{proof}

\begin{lemma}\label{pd-kappa-harmonic-split}
Let $p$ and $(W,Q,\sigma,\iota)$ be as in Definition~\ref{pd-kappa-split-rank}, $m\geq2$ and
$k=m+1$, $T\subset\mathbf Z/n'\mathbf Z$, $a\geq0$ an integer and $A,B\colon T\to\mathbf Z_{\geq0}$
with $A(s)+B(s)+1\leq8a$ for $s\in T$. Let $t,c\geq0$ be integers, $D\geq0$ and $X$ reals, and $e$ an
edge, and put $u_0=\lceil(m+2)(\log_p(8a+m)+1)\rceil$. Suppose that there is a split family of depth
$D$ and size $n\geq X$ on $(\mathcal C^\times_{t,c}(e),\mathcal C'^\times_{t,c}(e))$
(Definition~\ref{pd-kappa-product-space}). Then there are integers $n(j)\geq0$, for the functions
$j\colon T\to\mathbf Z_{\geq0}$ with $j\leq\min(A,B)$, such that $X\leq\sum_jn(j)$ and, for each
such $j$, there is a split family of depth $D+4|T|u_0$ and size $n(j)$ on the cores
$(\mathcal C_{t,c}(e),\mathcal C'_{t,c}(e))$ of the flat weights $\phi(A-j,B-j)$.
\end{lemma}
\leanref{Formalism/SharpRestriction/GrowthBounds/ColumnSpanDensity/RhoLowerBoundLargeConductor.lean}{Formalism.MainNotions.flipOneThree.Pf.RhoSumLowerBoundEdge}
\leanref{Formalism/SharpRestriction/GrowthBounds/ColumnSpanDensity/RhoSumBoundOfMapsSplitting.lean}{Formalism.MainNotions.flipOneThree.Pf.rho_sum_bound_of_maps_splitting}

\begin{proof}
Let $x_1,\dots,x_n$, $y_1,\dots,y_n$, $z$ and the matrix $M$ be the split family, so that
$(\langle x_i,y_l\rangle^\times)_{il}M=z\mathbf 1_n$. If $n=0$, take $n(j)=0$
for every $j$. Let $n\geq1$, and let $J$, $u$, $V_j$, $A_j$, $\eta_j$ and $\pi_j$ be as in
Lemma~\ref{pd-kappa-harmonic-maps}; write $\langle\ ,\ \rangle_j$ and $\langle\ ,\ \rangle^\times$
for the pairings on $V_j$ and on $V^\times(A,B)$. For $j\in J$ put
\[
 x_{j,i}=p^u\pi_j(x_i)\in\mathcal C_{j,t,c}(e),\qquad
 y_{j,l}=p^{2u+A^\times-A_j}(y_l\circ\eta_j)\in\mathcal C'_{j,t,c}(e).
\]
Then $\langle x_{j,i},y_{j,l}\rangle_j=p^{3u}\langle\eta_j\pi_j(x_i),y_l\rangle^\times$, and this lies
in $\mathcal R$ by Lemma~\ref{pd-edge-basis}(3), since the cores lie in $R(e)$ and $R'(e)$. Let $G_j$
be the $n\times n$ matrix with these entries and $\Gamma_j=G_jM$. Since
$\sum_j\eta_j\pi_j$ is the identity, $\sum_jG_j=p^{3u}(\langle x_i,y_l\rangle^\times)_{il}$, so
$\sum_j\Gamma_j=p^{3u}z\mathbf 1_n$. Lemma~\ref{pd-kappa-harmonic-block}, applied to the $\Gamma_j$
with $p^{3u}z\neq0$, gives integers $n(j)=\nu_j$ with $\sum_jn(j)\geq n\geq X$ and matrices $P_j$,
$Q_j$ with $P_j\Gamma_jQ_j=p^{3u}z\mathbf 1$. Put $x'_{j,a}=\sum_i(P_j)_{ai}x_{j,i}$ and
$y'_{j,b}=\sum_l(MQ_j)_{lb}y_{j,l}$. They lie in the cores, which are $\mathcal R$-modules, and
$\langle x'_{j,a},y'_{j,b}\rangle_j=(P_jG_jMQ_j)_{ab}=p^{3u}z\delta_{ab}$. With $p^{3u}z$, of valuation
at most $D+3u\leq D+4|T|u_0$, and the identity matrix, they form a split family of size $n(j)$.
\end{proof}

\begin{definition}\label{pd-kappa-function}
Let $p$ be a prime, $n',m,N_1\geq1$ and $G',c_{cl},D_p,c_*\geq0$ integers,
$T\subset\mathbf Z/n'\mathbf Z$, $Y\geq1$, $a\geq1$, $b'>0$ and $b_l\leq b_u$ reals, and
$C_0,C_A,C_t>0$ constants. Call $s\in\mathbf Z/n'$ even if its representative in
$\{0,\dots,n'-1\}$ is even, suppose that $T$ has $q\geq2$ even elements, and for odd $y\in T$ let
$d(y)\geq1$ be least with $y-d(y)$ an even element of $T$. Put $k=m+1$, $Q=p^{G'}$ and
\[
 \bar a=\lceil a\rceil,\quad u_a=\Big\lfloor a+\frac a{kY}\Big\rfloor,\quad l_b=\lceil b_l\rceil,\quad
 u_b=\lfloor b_u\rfloor,\quad\bar\Lambda=C_0N_1^2\log\bar a,\quad J_0=\frac{\bar\Lambda b'}k,\quad
 J=\lfloor J_0\rfloor,
\]
\[
 \ell=\log_p(8\bar a+m)+1,\quad u_0=\lceil(m+2)\ell\rceil,\quad
 c_\flat=c_*+2|T|c_{cl}+2n'c_{cl},
\]
\[
 D_0=q\big(C_AY^2N_1^2k^2\ell+4|T|k\log_p\bar a\big)+8|T|k\log_p\bar a,\qquad D_1=D_0+4|T|u_0 .
\]
For $A,B\colon T\to\mathbf Z_{\geq0}$, extended by $0$ to $\mathbf Z/n'$, put
\[
 t_0(A,B)=\sum_{y\in T\ \mathrm{odd}}\min\Big(B(y),\frac{A(y-d(y))}{p^{d(y)}}\Big)
 +C_t|T|\Big(\frac{\bar a}{p^2}+\frac{b'}p+\frac{\bar\Lambda\bar a}{pk}+\frac{\bar\Lambda b'}k
 +kQ+Qc_\flat\Big)
\]
and $t^*(A,B)=\lceil t_0(A,B)+|T|J_0\rceil$. Let $\mathcal S$ be an unramified quadratic Hida setting
of degree $n'$ at $p$. For $\alpha,\beta\colon T\to\mathbf Z_{\geq0}$, the \emph{flat defect}
$\kappa_\flat(\alpha,\beta)$ is the least $\kappa\geq0$ such that
$\mathrm{ER}\big(\phi(\alpha,\beta);t^*(\alpha,\beta),c_\flat,D_1,(1-\kappa)\mu(\phi(\alpha,\beta))\big)$
holds in $\mathcal S$ (Definition~\ref{pd-kappa-split-rank}); it exists and lies in $[0,1]$ by
Lemma~\ref{pd-kappa-rank-monotone}(3). For $a_w,b_w\colon T\to\mathbf Z$ put, with
$x^+=\max(x,0)$,
\[
 \kappa_{\mathcal S}(a_w,b_w)=\kappa_-+2\kappa_\flat\big((a_w+c_{cl})^+,(b_w+c_{cl})^+\big)
 +\frac{12|T|kc_{cl}}{b'},
\]
where $\kappa_-=\kappa_\flat\big((a_w-c_{cl})^+,(b_w-c_{cl})^+\big)$ if $a_w(s)-c_{cl}\geq\bar a$
and $b_w(s)-c_{cl}\geq b_l$ for every $s\in T$, and $\kappa_-=1$ otherwise.
\end{definition}
\leanref{Formalism/SharpRestriction/GrowthBounds/ColumnSpanDensity/RhoLowerBoundLargeConductor.lean}{Formalism.MainNotions.flipOneThree.Pf.Dat}
\leanref{Formalism/SharpRestriction/GrowthBounds/ColumnSpanDensity/MeanDefs.lean}{Formalism.MainNotions.flipOneThree.Pf.Dat.ua}
\leanref{Formalism/SharpRestriction/GrowthBounds/ColumnSpanDensity/RhoLowerBoundLargeConductor.lean}{Formalism.MainNotions.flipOneThree.Pf.Dat.c4}
\leanref{Formalism/SharpRestriction/GrowthBounds/ColumnSpanDensity/RhoLowerBoundLargeConductor.lean}{Formalism.MainNotions.flipOneThree.Pf.Dat.J0}
\leanref{Formalism/SharpRestriction/GrowthBounds/ProjectionLemmaBaseLevel.lean}{Formalism.MainNotions.flipOneThree.totalDegreeBudget}
\leanref{Formalism/SharpRestriction/GrowthBounds/FlipOneThree.lean}{Formalism.MainNotions.flipOneThree.depthD}
\leanref{Formalism/SharpRestriction/GrowthBounds/ColumnSpanDensity/RhoLowerBoundLargeConductor.lean}{Formalism.MainNotions.flipOneThree.Pf.Dat.D1}
\leanref{Formalism/SharpRestriction/GrowthBounds/ObjSeg.lean}{Formalism.MainNotions.G.t0G4}
\leanref{Formalism/SharpRestriction/GrowthBounds/ColumnSpanDensity/RhoLowerBoundLargeConductor.lean}{Formalism.MainNotions.flipOneThree.Pf.Dat.tStar}
\leanref{Formalism/SharpRestriction/GrowthBounds/ColumnSpanDensity/RhoLowerBoundLargeConductor.lean}{Formalism.MainNotions.flipOneThree.Pf.Dat.kF}
\leanref{Formalism/SharpRestriction/GrowthBounds/ColumnSpanDensity/RhoLowerBoundLargeConductor.lean}{Formalism.MainNotions.flipOneThree.Pf.Dat.Inside}
\leanref{Formalism/SharpRestriction/GrowthBounds/ColumnSpanDensity/RhoLowerBoundLargeConductor.lean}{Formalism.MainNotions.flipOneThree.Pf.Dat.kText}

\subsection{Segment lattices and flat theta maps}\label{subsection-pd-segment-lattices}

\begin{definition}\label{pd-kappa-defects-shares}
Let $p$ be a prime, $\mathcal O$ a discrete valuation ring with fraction field $\mathcal K$ in which $p$
is neither $0$ nor a unit, and $v$ the valuation of $\mathcal K$ with $v(p)=1$. Let $V$ be a
torsion-free $\mathcal O$-module, $\beta\colon V\times V\to\mathcal K$ a map, $L\subset V$ a submodule
and $D$ a real number. For a submodule $L'\subset V$ let $\operatorname{rk}L'$ be the largest number
of $\mathcal O$-linearly independent elements of $L'$, and $0$ if there is no largest.
\begin{enumerate}
\item An element $u$ of a submodule $L_1$ is \emph{primitive} in $L_1$ if $u=cx$ with $c\in\mathcal O$
  and $x\in L_1$ implies that $c$ is a unit.
\item A \emph{good sublattice} of $L$ of depth $D$ is a submodule $L_1\subset L$ such that $x\in L_1$
  whenever $x\in L$ and $cx\in L_1$ for a nonzero $c\in\mathcal O$, and such that for every primitive
  $u$ in $L_1$ there is $u'\in L$ with $\beta(u,u')\neq0$ and $v(\beta(u,u'))\leq D$.
\item For a real $\kappa$, $L$ has a \emph{good share} of depth $D$ and defect $\kappa$ if it has a good
  sublattice $L_1$ of depth $D$ with $\operatorname{rk}L_1\geq(1-\kappa)\operatorname{rk}L$.
\item The \emph{bad share} $\operatorname{bad}_D(L,\beta)$ is $0$ if $\operatorname{rk}L=0$, and
  otherwise $(\operatorname{rk}L-n)/\operatorname{rk}L$, where $n$ is the largest rank of a good
  sublattice of $L$ of depth $D$.
\end{enumerate}
\end{definition}
\leanref{Formalism/SharpRestriction/GrowthBounds/ObjSeg.lean}{Formalism.MainNotions.G.SatIn}
\leanref{Formalism/SharpRestriction/GrowthBounds/ObjSeg.lean}{Formalism.MainNotions.G.PrimIn}
\leanref{Formalism/SharpRestriction/GrowthBounds/ObjSeg.lean}{Formalism.MainNotions.G.vLE}
\leanref{Formalism/SharpRestriction/GrowthBounds/ObjSeg.lean}{Formalism.MainNotions.G.GoodSub}
\leanref{Formalism/SharpRestriction/GrowthBounds/ObjSeg.lean}{Formalism.MainNotions.G.GoodShare}
\leanref{Formalism/SharpRestriction/GrowthBounds/ObjSeg.lean}{Formalism.MainNotions.G.badShare}

\begin{lemma}\label{pd-kappa-defects-bad-share}
Let $p$, $\mathcal O$, $V$, $\beta$, $L$ and $D$ be as in Definition~\ref{pd-kappa-defects-shares}. Then
$\operatorname{bad}_D(L,\beta)\geq0$, and $L$ has a good share of depth $D$ and defect
$\operatorname{bad}_D(L,\beta)$.
\end{lemma}
\leanref{Formalism/SharpRestriction/GrowthBounds/Joints.lean}{Formalism.MainNotions.G.Joints.badShare_spec}

\begin{proof}
The submodule $0$ is a good sublattice of $L$ of every depth: if $x\in L$ and $cx=0$ with $c\neq0$,
then $x=0$, as $V$ is torsion-free; and $0$ has no primitive element, since $0=0\cdot0$ and $0$ is not
a unit. If $\operatorname{rk}L=0$, the bad share is $0$, and $L_1=0$ is a good share of defect $0$.
Otherwise $\operatorname{rk}L$ is finite and positive, and every submodule $L_1\subset L$ has
$\operatorname{rk}L_1\leq\operatorname{rk}L$. So the largest rank $n$ of a good sublattice exists and
$0\leq n\leq\operatorname{rk}L$, whence the bad share is $\geq0$; a good sublattice of rank $n$ has
$n=(1-\operatorname{bad}_D(L,\beta))\operatorname{rk}L$.
\end{proof}

\begin{definition}\label{pd-kappa-defects-segment-lattice}
Let $p$, $\mathcal O\subset\mathcal K$ and $v$ be as in Definition~\ref{pd-kappa-defects-shares}, let
$r\geq1$ and $s\geq0$ be integers, and write $x\dot-y=\max(x-y,0)$. A \emph{segment datum} $S$ of
length $s$ consists of integers $\eta,G,a,b_+\geq0$ and $d_i,b_i\geq0$ ($1\leq i\leq s$) with
$\eta\leq s$ and $d_i\leq G$, and of reals $Y$, $\Lambda$ and $E_0$. Put $P_i=p^{G-d_i}$ and
$\hat a=\lfloor(a-E_0)/p^G\rfloor$. Let $\alpha_1,\dots,\alpha_s,\alpha_+\geq0$ be integers and $\mu$
a real number.
\begin{enumerate}
\item A \emph{block} $b$ consists of vectors $\rho_i,\rho'_i\in\mathbf Z_{\geq0}^{r+1}$ and integers
  $o_i,w_i,\varsigma_i\geq0$ for $1\leq i\leq s$, and an integer $o_+\geq0$. Let $|\rho|$ be the sum of
  the entries of $\rho$, and put $\hat b_i=\lfloor(b_i\dot-|\rho_i|)/P_i\rfloor$,
  $\hat\alpha_i=\lfloor(\alpha_i\dot-|\rho'_i|)/P_i\rfloor$, $o^\flat_i=w_i+\varsigma_i$ for $i\leq\eta$,
  $o^\flat_i=o_i$ for $i>\eta$,
  \[
   \operatorname{cap}(l)=\hat a+\sum_{i=1}^l(\hat b_i-o^\flat_i)\quad(0\leq l\leq s),\qquad
   O_+=\operatorname{cap}(s)+\Big(\frac\Lambda{r+1}+\frac4{p^2}\Big)\sum_{i=1}^s\frac{2a}{P_i}
   +\frac{\Lambda b_+}{r+1},
  \]
  and $\operatorname{sp}(b)=\prod_{i\leq\eta}(\hat b_i\dot-w_i)!\,\varsigma_i!$.
\item The block $b$ is \emph{kept} if: for every $i$, the entries of $\rho_i$ and of $\rho'_i$ are less
  than $P_i$, $|\rho_i|\leq b_i$, $|\rho'_i|\leq\alpha_i$, and $P_i$ divides $b_i-|\rho_i|$ and
  $\alpha_i-|\rho'_i|$; for $i>\eta$, $o_i\leq\hat b_i$, $o_i\leq\operatorname{cap}(i-1)+\Lambda\hat
  b_i/(r+1)$ and $w_i=\varsigma_i=0$; for $i\leq\eta$, $o_i=0$, $w_i\leq\hat b_i$,
  $\varsigma_i\leq\hat\alpha_i$, $w_i+\varsigma_i\leq(1+4/p^2)\hat b_i+\Lambda\operatorname{cap}(i-1)/(r+1)$
  and $\varsigma_i\leq2\Lambda\hat a/(r+1)$; and $o_+\leq O_+$. The entries of a kept block are
  bounded, so there are finitely many kept blocks.
\item Let $\mathcal P$ be the polynomial ring over $\mathcal K$ in the variables $x_{i,c}$, $y_{i,c}$
  ($1\leq i\leq s$), $x_{+,c}$ and $\xi_c$ ($1\leq c\leq r$), and $z_l$, $z'_l$ ($1\leq l\leq s-\eta$, and
  $l=+$). The \emph{letters} are $x_1,y_1,\dots,x_s,y_s,x_+$, in this order, where $x_i=(x_{i,c})_c$,
  $y_i=(y_{i,c})_c$ and $x_+=(x_{+,c})_c$; the $x_i$ and $x_+$ are \emph{vector} letters and the $y_i$
  are \emph{covector} letters. For a block $b$ give the letters the degrees $\delta(x_i)=w_i$ and
  $\delta(y_i)=\hat\alpha_i\dot-\varsigma_i$ for $i\leq\eta$, $\delta(x_i)=o_i$ and
  $\delta(y_i)=\hat\alpha_i$ for $i>\eta$, and $\delta(x_+)=o_+$.
\item For a letter $\lambda$ write $\lambda_c$ for its $c$-th variable. For an initial segment
  $\mathcal L$ of the letters and $1\leq c,d\leq r$, let
  $E^{\mathcal L}_{cd}=\sum_\lambda\lambda_c\,\partial/\partial\lambda_d
  -\sum_{\lambda'}\lambda'_d\,\partial/\partial\lambda'_c$, the first sum over the vector letters
  $\lambda$ in $\mathcal L$ and the second over the covector letters $\lambda'$ in $\mathcal L$. The
  \emph{weight} of a monomial for $\mathcal L$ is the sum of its exponents at the variables $\lambda_1$,
  for the vector letters $\lambda$ in $\mathcal L$, minus the sum of its exponents at the variables
  $\lambda'_1$, for the covector letters $\lambda'$ in $\mathcal L$. For an integer $M$, let $Y(\mathcal L,M)$ be the largest subspace of the span of the monomials in the
  letters of weight at most $M$ for $\mathcal L$ that is stable under all the $E^{\mathcal L}_{cd}$.
\item For $1\leq l\leq s-\eta$ put $u_l=\eta+l$, let $\mathcal L_l$ consist of the letters up to and
  including $x_{u_l}$, and put $M_l=\lfloor\operatorname{cap}(u_l-1)-\mu\rfloor$; let $\mathcal L_+$
  consist of all the letters and put $M_+=\lfloor\operatorname{cap}(s)-\mu\rfloor$. Put
  $\hat E_l=z_lz'_l+\sum_cx_{u_l,c}y_{u_l,c}$ and $\hat E_+=z_+z'_++\sum_cx_{+,c}\xi_c$. Let
  $j_1,\dots,j_{s-\eta},j_+\geq0$ be integers with $j_l\leq\min(\delta(x_{u_l}),\delta(y_{u_l}))$ and
  $j_+\leq\min(\delta(x_+),b_+)$, and let $\delta'$ be $\delta$ lowered by $j_l$ at $x_{u_l}$ and at
  $y_{u_l}$ and by $j_+$ at $x_+$. The \emph{generators} of level $(j,j_+)$ are the products
  \[
   \prod_l\hat E_l^{j_l}\cdot\hat E_+^{j_+}\cdot\prod_lz_l^{\hat b_{u_l}\dot-o_{u_l}}\cdot
   z_+^{\alpha_+\dot-o_+}\cdot f\cdot g,
  \]
  where $f$ is a polynomial in the letters with coefficients in $\mathcal O$, of total degree
  $\delta'(\lambda)$ in the variables of each letter $\lambda$, that lies in $Y(\mathcal L_l,M_l)$ for
  every $l$ and in $Y(\mathcal L_+,M_+)$, and $g$ is a polynomial in the $\xi_c$ with coefficients in
  $\mathcal O$, homogeneous of degree $b_+-j_+$.
\item The \emph{block lattice} $L_b$ is the set of elements of the $\mathcal K$-span of the generators
  of all levels whose coefficient $\gamma$ at each monomial, with exponent vector $\nu$, satisfies
  $v(\gamma)+\frac12v_p(\nu!)+\frac12v_p(\operatorname{sp}(b))\geq0$, where $\nu!$ is as in
  Definition~\ref{pd-weight-lattices}(2).
\item Let $V_S$ be the $\mathcal O$-module of the families $(f_b)$ of elements of $\mathcal P$, indexed
  by the blocks, with finitely many nonzero entries. The \emph{segment lattice}
  $L_S(\alpha,\alpha_+;\mu)\subset V_S$ consists of the families with $f_b\in L_b$ for kept $b$ and
  $f_b=0$ otherwise. The \emph{segment pairing} on $V_S$ is
  $\beta_S(f,f')=\sum_b\operatorname{sp}(b)[f_b,f'_b]$, with the Fock pairing of
  Definition~\ref{pd-weight-lattices}(2).
\item For a real $D$, the \emph{segment defect} is
  $\kappa_S(\mu,D;\alpha,\alpha_+)=\operatorname{bad}_D(L_S(\alpha,\alpha_+;\mu),\beta_S)$
  (Definition~\ref{pd-kappa-defects-shares}).
\end{enumerate}
\end{definition}
\leanref{Formalism/SharpRestriction/GrowthBounds/ObjSeg.lean}{Formalism.MainNotions.G.SegData}
\leanref{Formalism/SharpRestriction/GrowthBounds/ObjSeg.lean}{Formalism.MainNotions.G.SegBlock}
\leanref{Formalism/SharpRestriction/GrowthBounds/ObjSeg.lean}{Formalism.MainNotions.G.coarseA}
\leanref{Formalism/SharpRestriction/GrowthBounds/ObjSeg.lean}{Formalism.MainNotions.G.coarseC}
\leanref{Formalism/SharpRestriction/GrowthBounds/ObjSeg.lean}{Formalism.MainNotions.G.oEff}
\leanref{Formalism/SharpRestriction/GrowthBounds/ObjSeg.lean}{Formalism.MainNotions.G.cap}
\leanref{Formalism/SharpRestriction/GrowthBounds/ObjSeg.lean}{Formalism.MainNotions.G.Oplus}
\leanref{Formalism/SharpRestriction/GrowthBounds/ObjSeg.lean}{Formalism.MainNotions.G.DigitsOK}
\leanref{Formalism/SharpRestriction/GrowthBounds/ObjSeg.lean}{Formalism.MainNotions.G.Kept}
\leanref{Formalism/SharpRestriction/GrowthBounds/ObjSeg.lean}{Formalism.MainNotions.G.segChain}
\leanref{Formalism/SharpRestriction/GrowthBounds/ObjPol.lean}{Formalism.MainNotions.G.glAct}
\leanref{Formalism/SharpRestriction/GrowthBounds/ObjSeg.lean}{Formalism.MainNotions.G.chWt}
\leanref{Formalism/SharpRestriction/GrowthBounds/ObjSeg.lean}{Formalism.MainNotions.G.YCond}
\leanref{Formalism/SharpRestriction/GrowthBounds/ObjSeg.lean}{Formalism.MainNotions.G.YjtZ}
\leanref{Formalism/SharpRestriction/GrowthBounds/ObjSeg.lean}{Formalism.MainNotions.G.segTies}
\leanref{Formalism/SharpRestriction/GrowthBounds/ObjSeg.lean}{Formalism.MainNotions.G.blockData}
\leanref{Formalism/SharpRestriction/GrowthBounds/ObjSeg.lean}{Formalism.MainNotions.G.BlockData.gens}
\leanref{Formalism/SharpRestriction/GrowthBounds/ObjSeg.lean}{Formalism.MainNotions.G.BlockData.SSZ}
\leanref{Formalism/SharpRestriction/GrowthBounds/ObjSeg.lean}{Formalism.MainNotions.G.bbSpec}
\leanref{Formalism/SharpRestriction/GrowthBounds/ObjSeg.lean}{Formalism.MainNotions.G.specConst}
\leanref{Formalism/SharpRestriction/GrowthBounds/ObjSeg.lean}{Formalism.MainNotions.G.blockLat}
\leanref{Formalism/SharpRestriction/GrowthBounds/ObjSeg.lean}{Formalism.MainNotions.G.segLat}
\leanref{Formalism/SharpRestriction/GrowthBounds/ObjSeg.lean}{Formalism.MainNotions.G.segForm}
\leanref{Formalism/SharpRestriction/GrowthBounds/ObjSeg.lean}{Formalism.MainNotions.G.kapT}

\begin{definition}\label{pd-kappa-defects-numbers}
Let $p\geq5$ be a prime, $n'\geq2$ even, $m\geq1$, $G'$, $N_1$, $c_{cl}$, $D_p$, $c_*\geq0$ and $a\geq1$
integers, $T\subset\mathbf Z/n'\mathbf Z$, and $Y,b'>0$ reals; let $s_1<\cdots<s_q$ be the even
elements of $T$ and $g_j$ the cyclic distance from $s_j$ to $s_{j+1}$, as in
Theorem~\ref{pd-gain-theorem-one}. Let $C_0,C_A,C_t,C_{E1},C_{E2},C_m,C_\kappa>0$ be constants, $c\geq0$
an integer and $a_T,b_T\colon\mathbf Z/n'\to\mathbf Z_{\geq0}$. Put $k=m+1$, $N=|T|$, $Q=p^{G'}$,
$\Lambda=C_0N_1^2\log a$, $J_0=\Lambda b'/k$, $\ell=\log_p(8a+m)+1$,
\[
 E_0=C_{E1}\Big(\frac a{p^2}+\frac{b'}p+\frac{\Lambda N_1a}{pk}+\frac{\Lambda b'}k\Big)+C_{E2}k^2N_1Q^2c,
 \qquad m_0=C_mkN_1Qc+12c,
\]
\[
 \kappa_0=2Nk\log_pa,\qquad D_A=C_AY^2N_1^2k^2\ell+2\kappa_0,\qquad
 \kappa_P(\mu)=C_\kappa N\Big(e^{-\Lambda/8}+\frac{kQ(\mu+N_1k)}{b'}+\frac1Y+\frac{kE_0}a\Big)
\]
for $\mu\in\mathbf R$, and $\kappa_P=\kappa_P(m_0)$. Let $t_0$ be formed as $t_0(A,B)$ in
Definition~\ref{pd-kappa-function}, with $a$, $\Lambda$, $c$, $a_T|_T$ and $b_T|_T$ in place of $\bar a$,
$\bar\Lambda$, $c_\flat$, $A$ and $B$. For $1\leq j\leq q$ let $y_{j,1},\dots,y_{j,S_j}$ be the odd
elements of $T$ strictly between $s_j$ and $s_{j+1}$ (indices modulo $q$), at cyclic distances
$d_{j,1}<\cdots<d_{j,S_j}$ from $s_j$. The \emph{segment datum} $S_j(b_T)$ of length $S_j$ has
$\eta_j=|\{i:p^{d_{j,i}}<a/b'\}|$, $G=g_j$, $d_i=d_{j,i}$, $b_i=b_T(y_{j,i})$, $b_+=b_T(s_{j+1})$,
and $a$, $Y$, $\Lambda$, $E_0$ (Definition~\ref{pd-kappa-defects-segment-lattice}). For an unramified
quadratic Hida setting of degree $n'$ at $p$ with ring $\mathcal R$, the \emph{segment lattice} of the
$j$-th segment is $L_{S_j(b_T)}(\alpha,\alpha_+;m_0)$ with its pairing, formed over $\mathcal R$ with
$r=m$, where $\alpha_i=a_T(y_{j,i})$ and $\alpha_+=a_T(s_{j+1})$, and its \emph{segment defect} is
$\kappa_j(a_T,b_T)=\kappa_{S_j(b_T)}(m_0,D_A;\alpha,\alpha_+)$.
\end{definition}
\leanref{Formalism/SharpRestriction/GrowthBounds/ColumnSystems.lean}{Formalism.MainNotions.ColumnSystems.Loc}
\leanref{Formalism/SharpRestriction/GrowthBounds/ColumnSystems.lean}{Formalism.MainNotions.ColumnSystems.Loc.seg}
\leanref{Formalism/SharpRestriction/GrowthBounds/ObjSeg.lean}{Formalism.MainNotions.G.segData}
\leanref{Formalism/SharpRestriction/GrowthBounds/ObjSeg.lean}{Formalism.MainNotions.G.segATop}
\leanref{Formalism/SharpRestriction/GrowthBounds/ObjSeg.lean}{Formalism.MainNotions.G.segAPlus}
\leanref{Formalism/SharpRestriction/GrowthBounds/ObjSeg.lean}{Formalism.MainNotions.G.errorTermBound}
\leanref{Formalism/SharpRestriction/GrowthBounds/ObjSeg.lean}{Formalism.MainNotions.G.m0G4}
\leanref{Formalism/SharpRestriction/GrowthBounds/ObjSeg.lean}{Formalism.MainNotions.G.kap0G4}
\leanref{Formalism/SharpRestriction/GrowthBounds/ObjSeg.lean}{Formalism.MainNotions.G.kappaMarginBound}
\leanref{Formalism/SharpRestriction/GrowthBounds/ObjSeg.lean}{Formalism.MainNotions.G.t0G4}
\leanref{Formalism/SharpRestriction/GrowthBounds/FlipOneThree.lean}{Formalism.MainNotions.flipOneThree.depthD}
\leanref{Formalism/SharpRestriction/GrowthBounds/ProjectionLemmaBaseLevel.lean}{Formalism.MainNotions.flipOneThree.kapPI}
\leanref{Formalism/SharpRestriction/GrowthBounds/ProjectionLemmaBaseLevel.lean}{Formalism.MainNotions.flipOneThree.kapTI}

\begin{lemma}\label{pd-kappa-defects-margin}
Let $C_c\geq0$, $C_0\geq1$, $C_{E1},C_{E2},C_m\geq1$, $C_\kappa>0$ and
$C_H\geq1000(C_{E1}+C_{E2}+1)(C_c+1)$. Let $p\geq5$ be a prime and $n'\geq2$ an even integer, and let
$m$, $G'$, $N_1$, $c_{cl}$, $D_p$, $c_*$, $T$, $Y$, $Y_b$, $b'$, $b_l$, $b_u$ and an integer $a$ satisfy hypotheses
(1) to (3) of Theorem~\ref{pd-gain-theorem-one} with $C_0$ and $C_H$. Let $c\geq0$ be an integer with
$c\leq C_c(c_*+|T|k\log_pa+|T|c_{cl}+n'c_{cl})$, and form the margin term $\kappa_P$ of
Definition~\ref{pd-kappa-defects-numbers} with these constants. Then
$\kappa_P\leq C_\kappa(C_m+4C_{E1}+C_{E2}+15)|T|N_1/Y$.
\end{lemma}
\leanref{Formalism/SharpRestriction/GrowthBounds/BinomialRatio/KapCore.lean}{Formalism.MainNotions.ColumnSystems.Pf.rest2}
\leanref{Formalism/SharpRestriction/GrowthBounds/BinomialRatio/KapCore.lean}{Formalism.MainNotions.ColumnSystems.Pf.kap_core}
\leanref{Formalism/SharpRestriction/GrowthBounds/BinomialRatio/ThresholdLowerBounds.lean}{Formalism.MainNotions.ColumnSystems.Pf.std_of}
\leanref{Formalism/SharpRestriction/GrowthBounds/BinomialRatio/ThresholdLowerBounds.lean}{Formalism.MainNotions.ColumnSystems.Pf.T2}
\leanref{Formalism/SharpRestriction/GrowthBounds/ColumnSystems.lean}{Formalism.MainNotions.ColumnSystems.ExtendedKappaBound}

\begin{proof}
Write $N=|T|$, $F=c_*+Nc_{cl}+n'c_{cl}$, $\beta=b'/Y$ and $\alpha=a/Y$. We first collect facts from the
hypotheses.
\begin{enumerate}
\item[(a)] $k\geq2N_1+3\geq5$, as $2N_1+2\leq m$ and $N_1\geq1$; $Q\geq1$, as $p\geq5$; $N\geq q\geq2$;
  and $kY\leq k^2Y^2\leq p$, so $Y\leq p$.
\item[(b)] By (3), $D_p\geq0$ and $c_*\geq8(D_p+1)\geq8$, so $F\geq8$; as $C_H\geq1$, the last inequality
  of (3) gives $b'\geq C_Hk^3NQ^4FY\geq1$.
\item[(c)] By (2), $a>pb'\geq5$, so $\log a\geq\log4\geq1$, and $\Lambda=C_0N_1^2\log a\geq1$.
\item[(d)] By (2), $2Y_bN\Lambda<k$ with $Y_b\geq C_H\geq1$, so $2N\Lambda\leq k$ and $\Lambda\leq k$.
\item[(e)] By (2), $kY^2N\Lambda\log a\leq kY^2Y_bN\Lambda\log a\leq a/b'$. As $YN\log a\geq1$, this gives
  $kY\Lambda b'\leq a$; as $k\geq1$ and $Y\Lambda\geq1$, also $\Lambda Yb'\leq a$ and $kb'\leq a$, and
  $b'\leq a$.
\item[(f)] By (3), $C_HNY^2(N_1^5+k^3)\log_pa\leq D_p$ with $C_HY^2\geq1$ and $k\leq k^3$, so
  $Nk\log_pa\leq D_p\leq c_*$. Hence $c\leq C_c(2c_*+Nc_{cl}+n'c_{cl})\leq2C_cF\leq C_HF$, since
  $2C_c\leq C_H$.
\item[(g)] So $k^3NQ^4cY\leq C_Hk^3NQ^4FY\leq b'$, and as $NQ^2\geq1$, $k^3Q^2cY\leq b'$. As $C_HF\geq1$,
  also $k^3NQ^4Y\leq b'$, and so $k^3QY\leq b'$.
\end{enumerate}
We bound the four terms in the bracket of $\kappa_P$. By hypothesis (1),
$e^{-\Lambda/8}\leq e^{-\log(NY)}=1/(NY)\leq N_1/Y$, and $1/Y\leq N_1/Y$. Next
\[
 kQ(m_0+N_1k)=C_mN_1\cdot k^2Q^2c+12\cdot kQc+N_1\cdot k^2Q,
\]
and $k^2Q^2c$, $kQc$ and $k^2Q$ are at most $\beta$ by (g), so
$kQ(m_0+N_1k)/b'\leq(C_mN_1+12+N_1)\beta/b'\leq(C_m+13)N_1/Y$. Finally
\[
 kE_0=C_{E1}\Big(\frac{ka}{p^2}+\frac{kb'}p+\frac{\Lambda N_1a}p+\Lambda b'\Big)+C_{E2}N_1k^3Q^2c .
\]
Here $ka/p^2\leq\alpha$ as $kY\leq p^2$; $kb'/p\leq\alpha$ as $kb'\leq a$ and $Y\leq p$;
$\Lambda N_1a/p\leq N_1\alpha$ as $\Lambda Y\leq kY\leq p$; $\Lambda b'\leq\alpha$ as $\Lambda Yb'\leq a$;
and $N_1k^3Q^2c\leq N_1\beta\leq N_1\alpha$ by (g) and $b'\leq a$. So
$kE_0/a\leq(4C_{E1}+C_{E2})N_1/Y$. Adding, the bracket is at most $(C_m+4C_{E1}+C_{E2}+15)N_1/Y$, and
multiplying by $C_\kappa N$ gives the claim.
\end{proof}
% Lean: rest2 (FamilyFourthBound, TransportedAlt.of_explicit with thresholds 1, 1 and
% T2 CE1 CE2 Cc 1 = 1000(CE1 + CE2 + 1)(Cc + 1), and Cs = Ckap (Cm + 4 CE1 + CE2 + 15)): the facts
% (a) to (g) are those of std_of with H = 1 (k5, Q1, N2, hb1, hloga, hLa1, hLaN, LaLeK, hab, kYLab, kb, ba,
% hNkL, hcF, hHc', hHcF, hHbF, Hc2, Hb1), and the computation is kap_core. ExtendedKappaBound also
% assumes the windows InWin for a_T and b_T; std_of does not use them for the facts that kap_core
% needs, and the book omits them.

\begin{definition}\label{pd-kappa-defects-windows-blocks}
Let $p$, $\mathcal O\subset\mathcal K$, $v$, $r$, a segment datum $S$, integers
$\alpha_1,\dots,\alpha_s,\alpha_+\geq0$ and a real number $\mu$ be as in
Definition~\ref{pd-kappa-defects-segment-lattice}, and put $L=L_S(\alpha,\alpha_+;\mu)$.
\begin{enumerate}
\item For a block $b$, $V_b\subset\mathcal P$ is the $\mathcal K$-span of the generators of all levels
  of $b$; so $L_b\subset V_b$.
\item Let $g\geq0$ be an integer and $D$ a real number. A \emph{block-diagonal family} of size $g$ and
  depth $D$ in $L$ consists of $x_1,\dots,x_g,y_1,\dots,y_g\in L$, blocks $b_1,\dots,b_g$ and nonzero
  $z_1,\dots,z_g\in\mathcal O$ with $v(z_i)\leq D$, such that the entries of $x_i$ and of $y_i$ at every
  block other than $b_i$ are $0$ and $\beta_S(x_i,y_l)=\delta_{il}z_i$ for all $i$ and $l$.
\end{enumerate}
\end{definition}
\leanref{Formalism/SharpRestriction/GrowthBounds/BinomialRatio/TotalDegreeEplus.lean}{Formalism.MainNotions.ColumnSystems.Pf.Vb}
\leanref{Formalism/SharpRestriction/GrowthBounds/ObjSeg.lean}{Formalism.MainNotions.G.blockRank}
\leanref{Formalism/SharpRestriction/GrowthBounds/ObjSeg.lean}{Formalism.MainNotions.G.segRank}
\leanref{Formalism/SharpRestriction/GrowthBounds/BinomialRatio/ExtendedStatementOfParts.lean}{Formalism.MainNotions.ColumnSystems.Pf.BlockDiag}

\begin{lemma}\label{pd-kappa-defects-windows-ranks}
Let $p$, $\mathcal O\subset\mathcal K$, $r$, $S$, $\alpha$, $\alpha_+$ and $\mu$ be as in
Definition~\ref{pd-kappa-defects-windows-blocks}. Then $\operatorname{rk}L_S(\alpha,\alpha_+;\mu)$
(Definition~\ref{pd-kappa-defects-shares}) is the sum of $\dim_{\mathcal K}V_b$ over the kept blocks
$b$, and each of these $V_b$ is finite-dimensional.
\end{lemma}
\leanref{Formalism/SharpRestriction/GrowthBounds/BinomialRatio/ExtendedStatementOfParts.lean}{Formalism.MainNotions.ColumnSystems.Pf.RkAt}
\leanref{Formalism/SharpRestriction/GrowthBounds/BinomialRatio/TotalDegreeEplus.lean}{Formalism.MainNotions.ColumnSystems.Pf.rkAt_all}
\leanref{Formalism/SharpRestriction/GrowthBounds/BinomialRatio/TotalDegreeEplus.lean}{Formalism.MainNotions.ColumnSystems.Pf.finrank_segLat}
\leanref{Formalism/SharpRestriction/GrowthBounds/BinomialRatio/TotalDegreeEplus.lean}{Formalism.MainNotions.ColumnSystems.Pf.finrank_lattice}

\begin{proof}
Put $L=L_S(\alpha,\alpha_+;\mu)$ and let $B$ be the finite set of kept blocks. The generators of a block
have bounded total degree, so each $V_b$ is finite-dimensional. Let $V$ be the space of the families
$(f_b)$ with $f_b\in V_b$ for $b\in B$ and $f_b=0$ otherwise; it has dimension
$\sum_{b\in B}\dim V_b$, and $L\subset V$. Every $f\in V$ has a nonzero multiple in $L$: choose a
nonzero $c\in\mathcal O$ with $v(c\gamma)\geq0$ for the finitely many coefficients $\gamma$ of the
entries of $f$; then $cf\in L$, since $v_p(\nu!)\geq0$ and $\operatorname{sp}(b)$ is a positive integer.
Elements of $L$ that are linearly independent over $\mathcal O$ are linearly independent over
$\mathcal K$, as $\mathcal K$ is the fraction field of $\mathcal O$; and nonzero multiples in $L$ of a
basis of $V$ are linearly independent over $\mathcal O$. So $\operatorname{rk}L=\dim V$.
\end{proof}
% Lean: rkAt_all = finrank_segLat (finrank_lattice, smul_mem_bbSpec, spanSSZ_finite via spanSSZ_le and
% restrictTotalDegree) with kf_kept_finite.

\begin{lemma}\label{pd-kappa-defects-windows-partners}
Let $p$, $\mathcal O\subset\mathcal K$, $v$, $r$, $S$, $\alpha$, $\alpha_+$ and $\mu$ be as in
Definition~\ref{pd-kappa-defects-windows-blocks}, and $D\geq0$ a real number. If $L_1$ is a good
sublattice of $L_S(\alpha,\alpha_+;\mu)$ of depth $D$ for the segment pairing $\beta_S$
(Definition~\ref{pd-kappa-defects-shares}), then $L_S(\alpha,\alpha_+;\mu)$ has a block-diagonal
family of size $\operatorname{rk}L_1$ and depth $D$.
\end{lemma}
\leanref{Formalism/SharpRestriction/GrowthBounds/BinomialRatio/ExtendedStatementOfParts.lean}{Formalism.MainNotions.ColumnSystems.Pf.PartnersToFactorsAt}
\leanref{Formalism/SharpRestriction/GrowthBounds/BinomialRatio/BlockFG.lean}{Formalism.MainNotions.ColumnSystems.Pf.partners_to_factors_holds}
\leanref{Formalism/SharpRestriction/GrowthBounds/BinomialRatio/PartnersFactorsOfKernelSmith.lean}{Formalism.MainNotions.ColumnSystems.Pf.SmithData}
\leanref{Formalism/SharpRestriction/GrowthBounds/BinomialRatio/FockZeroRight.lean}{Formalism.MainNotions.ColumnSystems.Pf.smithData_seg}
\leanref{Formalism/SharpRestriction/GrowthBounds/BinomialRatio/PartnersFactorsOfKernelSmith.lean}{Formalism.MainNotions.ColumnSystems.Pf.ker_core}
\leanref{Formalism/SharpRestriction/GrowthBounds/BinomialRatio/PartnersFactorsOfKernelSmith.lean}{Formalism.MainNotions.ColumnSystems.Pf.blockDiag_of_smith}

\begin{proof}
Put $L=L_S(\alpha,\alpha_+;\mu)$, let $B$ be the finite set of kept blocks and $\pi$ a uniformizer of
$\mathcal O$. The elements of $L$ are the families $(f_b)$ with $f_b\in L_b$ for $b\in B$ and $f_b=0$
otherwise.

\emph{Step 1.} For $b\in B$, $L_b$ is free of finite rank over $\mathcal O$. Indeed $V_b$ is
finite-dimensional (Lemma~\ref{pd-kappa-defects-windows-ranks}), so only finitely many monomials occur
in its elements, and the condition defining $L_b$ bounds below the valuations of the coefficients at
these monomials. So taking coefficients embeds $L_b$ in a finitely generated $\mathcal O$-module. As
$\mathcal O$ is noetherian, $L_b$ is finitely generated; it is torsion-free, hence free.

\emph{Step 2.} For $b\in B$ there are bases $(\alpha^b_i)$ and $(\beta^b_i)$ of $L_b$ with
$[\alpha^b_i,\beta^b_l]=0$ for $i\neq l$. Let $(e_a)$ be a basis of $L_b$ and $G$ the matrix
$([e_a,e_c])$ over $\mathcal K$. By the elementary divisor theorem over $\mathcal O$, applied to a
nonzero multiple of $G$ with entries in $\mathcal O$, there are matrices $P$ and $Q$ over $\mathcal O$
with unit determinants such that $PGQ$ is diagonal; put $\alpha^b_i=\sum_aP_{ia}e_a$ and
$\beta^b_l=\sum_cQ_{cl}e_c$.

\emph{Step 3.} Let $\alpha_1,\dots,\alpha_n$ be the families whose only nonzero entry is one of the
$\alpha^b_i$, at the block $b$, for $b\in B$ and all $i$, and $\beta_1,\dots,\beta_n$ the families
formed in the same way from the $\beta^b_i$, in the same order; let $b(i)$ be the block of $\alpha_i$
and $\beta_i$. The $\alpha_i$ form a basis of $L$, and the $\beta_i$ span $L$. If $i\neq l$, then
$\beta_S(\alpha_i,\beta_l)=0$: the terms of $\beta_S$ at blocks other than $b(i)$ vanish, and the term at
$b(i)$ is $0$ if $b(l)\neq b(i)$ and $\operatorname{sp}(b(i))[\alpha^{b(i)}_{i'},\beta^{b(i)}_{l'}]=0$
otherwise. By bilinearity,
$\beta_S\big(\sum_ia_i\alpha_i,\sum_l\eta_l\beta_l\big)=\sum_ia_i\eta_i\beta_S(\alpha_i,\beta_i)$ for
$a,\eta\in\mathcal O^n$.

\emph{Step 4.} Let $I_D$ be the set of $i$ with $\beta_S(\alpha_i,\beta_i)\neq0$ and
$v(\beta_S(\alpha_i,\beta_i))\leq D$. We show: if $a\in\mathcal O^n$, $\sum_ia_i\alpha_i\in L_1$ and
$a_i=0$ for $i\in I_D$, then $a=0$. Otherwise write $a=\pi^ka'$ with $a'\in\mathcal O^n$ and $a'_{i_0}$ a
unit for some $i_0$. As $L_1$ is saturated in $L$, $u=\sum_ia'_i\alpha_i$ lies in $L_1$. It is primitive
in $L_1$: if $u=cx$ with $x=\sum_i\xi_i\alpha_i\in L_1$, then $a'_{i_0}=c\xi_{i_0}$, so $c$ is a unit. As
$L_1$ is good, there is $y\in L$ with $\beta_S(u,y)\neq0$ and $v(\beta_S(u,y))\leq D$. Write
$y=\sum_l\eta_l\beta_l$; by Step 3, $\beta_S(u,y)=\sum_ia'_i\eta_i\beta_S(\alpha_i,\beta_i)$. For
$i\in I_D$, $a'_i=0$, as $\pi^ka'_i=a_i=0$; for $i\notin I_D$, $\beta_S(\alpha_i,\beta_i)$ is $0$ or has
valuation $>D$, and then so does $a'_i\eta_i\beta_S(\alpha_i,\beta_i)$. So the sum is $0$ or has
valuation $>D$, a contradiction.

\emph{Step 5.} The map $L_1\to\mathcal O^{I_D}$ that sends $\sum_ia_i\alpha_i$ to $(a_i)_{i\in I_D}$ is
$\mathcal O$-linear, and injective by Step 4. So $\operatorname{rk}L_1\leq|I_D|$.

\emph{Step 6.} Let $g=\operatorname{rk}L_1$, choose distinct $i_1,\dots,i_g\in I_D$, and put
$\kappa_t=\beta_S(\alpha_{i_t},\beta_{i_t})$, so $\kappa_t\neq0$ and $v(\kappa_t)\leq D$. If
$\kappa_t\in\mathcal O$, put $c_t=1$ and $z_t=\kappa_t$; otherwise put $c_t=\kappa_t^{-1}\in\mathcal O$
and $z_t=1$, so $v(z_t)=0\leq D$. Then $x_t=\alpha_{i_t}$, $y_t=c_t\beta_{i_t}$, the blocks $b(i_t)$ and
the $z_t$ form a block-diagonal family of size $g$ and depth $D$: $\beta_S(x_t,y_t)=c_t\kappa_t=z_t$,
and $\beta_S(x_t,y_{t'})=0$ for $t\neq t'$ by Step 3.
\end{proof}
% Lean: partners_to_factors_holds = partners_factors_of_smith_bases restSmith, through
% partners_factors_of_kernel_smith and partners_factors_at_of_smith. Step 1 is blockFG_all (fg_inter,
% bb_coeff, spanSSZ_finite); Step 2 is smith_block (Module.free_of_finite_type_torsion_free',
% KP.smithPairing); Step 3 is smithData_seg (span_single, kf_kept_finite) and bilinAt_all
% (segForm_bilin, fock_bilin); Step 4 is ker_core (prim_div, bad_sum, bad_mul); Step 5 is rank_le_of_ker
% with countSmithAt_of; Step 6 is blockDiag_of_smith (scale_to_O, segForm_smul_right).

\begin{lemma}\label{pd-kappa-defects-windows-model-segment}
For every $C_c\geq0$ there is $T_0>0$ such that for every $C_0\geq T_0$ there is $T_1>0$ such that for
all $C_{E1},C_{E2},C_m\geq T_1$ there is $T_2>0$ with the following property for every $C_H\geq T_2$.
Let $p\geq5$ be a prime and $n'\geq2$ an even integer, and let $m$, $G'$, $N_1$, $c_{cl}$, $D_p$, $c_*$,
$T$, $Y$, $Y_b$, $b'$, $b_l$, $b_u$ and an integer $a$ satisfy hypotheses (1) to (3) of
Theorem~\ref{pd-gain-theorem-one} with $C_0$ and $C_H$. Let $c\geq0$ be an integer with
$c\leq C_c(c_*+|T|k\log_pa+|T|c_{cl}+n'c_{cl})$, and let $a_T,b_T\colon\mathbf Z/n'\to\mathbf Z_{\geq0}$
satisfy $a\leq a_T(s)\leq a+a/(kY)+2J_0$ and $b_l\leq b_T(s)\leq b_u+2J_0$ for $s\in T$, where
$J_0=\Lambda b'/k$. Form the numbers of Definition~\ref{pd-kappa-defects-numbers} with these
constants; the segment data do not involve $C_m$, $C_A$, $C_t$ or $C_\kappa$. For an integer $d$ put
$P_m(d)=\binom{d+m}m$ if $d\geq0$ and $P_m(d)=0$ if $d<0$. Let an unramified quadratic Hida setting of
degree $n'$ at $p$ with ring $\mathcal R$ be given, an index $1\leq j\leq q$ and a real number
$\mu\geq0$, and put $\hat a_j=\lfloor(a-E_0)/p^{g_j}\rfloor$ and $n=\lceil\mu\rceil$. If
$n\leq\hat a_jN_1\Lambda/k$, then
\[
 \big(1-e^{-\Lambda/8}\big)\big(1-e^{-\Lambda/3}\big)\Big(1-\frac{mn}{\hat a_j}\Big)P_m(\hat a_j)
 \prod_{i=1}^{S_j}P_m\big(b_T(y_{j,i})\big)P_m\big(a_T(y_{j,i})\big)\cdot P_m\big(b_T(s_{j+1})\big)
 \leq r_j(\mu),
\]
where $r_j(\mu)$ is the sum of $\dim V_b$ over the kept blocks $b$ of the segment datum $S_j(b_T)$, for
$\alpha_i=a_T(y_{j,i})$, $\alpha_+=a_T(s_{j+1})$ and the real number $\mu$, formed over $\mathcal R$ with
$r=m$ (Definition~\ref{pd-kappa-defects-windows-blocks}).
\end{lemma}
\leanref{Formalism/SharpRestriction/GrowthBounds/BinomialRatio/RankDefs.lean}{Formalism.MainNotions.ColumnSystems.Pf.RestSegLow}
\leanref{Formalism/SharpRestriction/GrowthBounds/BinomialRatio/RankDefs.lean}{Formalism.MainNotions.ColumnSystems.Pf.SegLowAt}
\leanref{Formalism/SharpRestriction/GrowthBounds/BinomialRatio/RankDefs.lean}{Formalism.MainNotions.ColumnSystems.Pf.lowFac}
\leanref{Formalism/SharpRestriction/GrowthBounds/BinomialRatio/RankDefs.lean}{Formalism.MainNotions.ColumnSystems.Pf.sd}
\leanref{Formalism/SharpRestriction/GrowthBounds/ColumnSystems.lean}{Formalism.MainNotions.ColumnSystems.symDim}
\leanref{Formalism/SharpRestriction/GrowthBounds/BinomialRatio/SegLow.lean}{Formalism.MainNotions.ColumnSystems.Pf.restSegLow_of}

\begin{proof}
Let $T_0$, $T_1$ and $T_2$ be the numbers of Lemma~\ref{pd-kappa-defects-digits-lower}, and let the data
be as in the statement. Write $b_i=b_T(y_{j,i})$, $\alpha_i=a_T(y_{j,i})$ and $P_i=p^{g_j-d_{j,i}}$ for
$1\leq i\leq S_j$, and let $\mathcal D$ be the set of the digits $(\rho,\rho')$ for $j$
(Definition~\ref{pd-kappa-defects-segment-mean-conditions}). It is finite, as the entries of digits are
less than the $P_i$.

\emph{Step 1: the sum of the $d_j$.} For integers $P\geq1$ and $g\geq0$ let $\mathcal D(P,g)$ be the set
of the $\rho\in\mathbf Z^{m+1}_{\geq0}$ with entries less than $P$, $|\rho|\leq g$ and $P$ dividing
$g-|\rho|$. Then
\[
 \sum_{\rho\in\mathcal D(P,g)}P_m\big((g-|\rho|)/P\big)=P_m(g).
\]
Indeed $P_m(g)=\binom{g+m}m$ is the number of the $e\in\mathbf Z^{m+1}_{\geq0}$ with $|e|=g$. Let $\rho$ be
the vector of the remainders of the entries of such an $e$ modulo $P$. Then
$|e|=|\rho|+P\sum_c\lfloor e_c/P\rfloor$, so $\rho\in\mathcal D(P,g)$. The $e$ with a given $\rho$ are the
$\rho+Pq$ with $q\in\mathbf Z^{m+1}_{\geq0}$ and $|q|=(g-|\rho|)/P$, and there are $P_m((g-|\rho|)/P)$ of
them. The digits for $j$ are the pairs with $\rho_i\in\mathcal D(P_i,b_i)$ and
$\rho'_i\in\mathcal D(P_i,\alpha_i)$ for every $i$, and $d_j(\rho,\rho')$ is
$P_m(\hat a_j)P_m(b_T(s_{j+1}))$ times the product over $i$ of $P_m((b_i-|\rho_i|)/P_i)$ and
$P_m((\alpha_i-|\rho'_i|)/P_i)$. Expanding the product of the sums by the identity gives
\[
 \sum_{(\rho,\rho')\in\mathcal D}d_j(\rho,\rho')
 =P_m(\hat a_j)\prod_{i=1}^{S_j}P_m(b_i)P_m(\alpha_i)\cdot P_m\big(b_T(s_{j+1})\big).
\]

\emph{Step 2: the sum of the $r_j$.} By definition, $r_j(\mu;\rho,\rho')$ is the sum of $\dim V_b$ over the
kept blocks $b$ of $S_j(b_T)$ whose vectors are $\rho$ and $\rho'$. A block has one pair of vectors, so
every kept block occurs in at most one of these sums, and
$\sum_{(\rho,\rho')\in\mathcal D}r_j(\mu;\rho,\rho')\leq r_j(\mu)$.

\emph{Step 3.} Let $F$ be the factor $(1-e^{-\Lambda/8})(1-e^{-\Lambda/3})(1-mn/\hat a_j)$. As
$n\leq\hat a_jN_1\Lambda/k$, Lemma~\ref{pd-kappa-defects-digits-lower} gives
$Fd_j(\rho,\rho')\leq r_j(\mu;\rho,\rho')$ for every $(\rho,\rho')\in\mathcal D$. Summing over $\mathcal D$,
Steps 1 and 2 give the claim.
\end{proof}
% Lean: RestSegLow = restSegLow_of hL = TransportedAlt.mono of RestLow (Lemma
%   pd-kappa-defects-digits-lower) with the same thresholds, through segLow_of. Step 1 is
%   sum_prod_sd_digits_eq_prod_sd with digit_sum (card_tuple, the fibres of the reduction modulo P),
%   symDim_nat and the form of mmB; the set of digits is the product of the dset, each of whose members
%   satisfies DigitsOK (hdig). Step 2 is finsum_classes with the key b |-> (b.rz, b.ry) and
%   kf_kept_finite. Step 3 is the digit-wise bound LowAt summed with Finset.sum_le_sum.

\begin{lemma}\label{pd-kappa-defects-windows-model-step}
Let $C_c\geq0$, $C_0\geq1$, $C_{E1},C_{E2}\geq1$ and $C_H\geq1000(C_{E1}+C_{E2}+1)(C_c+1)$. Let $p\geq5$
be a prime and $n'\geq2$ an even integer, and let $m$, $G'$, $N_1$, $c_{cl}$, $D_p$, $c_*$, $T$, $Y$,
$Y_b$, $b'$, $b_l$, $b_u$ and an integer $a$ satisfy hypotheses (1) to (3) of
Theorem~\ref{pd-gain-theorem-one} with $C_0$ and $C_H$. Let $c\geq0$ be an integer with
$c\leq C_c(c_*+|T|k\log_pa+|T|c_{cl}+n'c_{cl})$, and let $a_T,b_T\colon\mathbf Z/n'\to\mathbf Z_{\geq0}$
with $a\leq a_T(s)\leq a+a/(kY)+2J_0$ for $s\in T$, where $J_0=\Lambda b'/k$. Form the numbers of
Definition~\ref{pd-kappa-defects-numbers} with these constants; those used here do not involve $C_m$,
$C_A$, $C_t$ or $C_\kappa$. Let $P_m$ be as in Lemma~\ref{pd-kappa-defects-windows-model-segment}, and
let $\mu\geq0$ be a real number with
\[
 \beta=e^{-\Lambda/8}+\frac{kQ(\mu+N_1k)}{b'}+\frac1Y+\frac{kE_0}a<\frac16 .
\]
Let $1\leq j\leq q$, put $\hat a_j=\lfloor(a-E_0)/p^{g_j}\rfloor$, $n=\lceil\mu\rceil$ and
$\Pi_j=\prod_{i=1}^{S_j}P_m(b_T(y_{j,i}))P_m(a_T(y_{j,i}))\cdot P_m(b_T(s_{j+1}))$, and let $R$ be a
real number such that
$(1-e^{-\Lambda/8})(1-e^{-\Lambda/3})(1-mn/\hat a_j)P_m(\hat a_j)\Pi_j\leq R$ if
$n\leq\hat a_jN_1\Lambda/k$. Then
\[
 (1-3\beta)P_m\big(a_T(s_j)\big)\Pi_j\leq p^{g_jm}R .
\]
\end{lemma}
\leanref{Formalism/SharpRestriction/GrowthBounds/BinomialRatio/FirstCoefficientFacts.lean}{Formalism.MainNotions.ColumnSystems.Pf.seg_step}
\leanref{Formalism/SharpRestriction/GrowthBounds/BinomialRatio/FirstCoefficientFacts.lean}{Formalism.MainNotions.ColumnSystems.Pf.brk}
\leanref{Formalism/SharpRestriction/GrowthBounds/BinomialRatio/FirstCoefficientFacts.lean}{Formalism.MainNotions.ColumnSystems.Pf.first_coefficient_facts}
\leanref{Formalism/SharpRestriction/GrowthBounds/BinomialRatio/FirstCoefficientFacts.lean}{Formalism.MainNotions.ColumnSystems.Pf.seg_num}
\leanref{Formalism/SharpRestriction/GrowthBounds/BinomialRatio/FirstCoefficientFacts.lean}{Formalism.MainNotions.ColumnSystems.Pf.ratio_num}
\leanref{Formalism/SharpRestriction/GrowthBounds/BinomialRatio/FirstCoefficientFacts.lean}{Formalism.MainNotions.ColumnSystems.Pf.choose_ratio}
\leanref{Formalism/SharpRestriction/GrowthBounds/BinomialRatio/FirstCoefficientFacts.lean}{Formalism.MainNotions.ColumnSystems.Pf.four_fac}
\leanref{Formalism/SharpRestriction/GrowthBounds/BinomialRatio/FirstCoefficientFacts.lean}{Formalism.MainNotions.ColumnSystems.Pf.chain_num}
\leanref{Formalism/SharpRestriction/GrowthBounds/BinomialRatio/ThresholdLowerBounds.lean}{Formalism.MainNotions.ColumnSystems.Pf.std_of}
\leanref{Formalism/SharpRestriction/GrowthBounds/BinomialRatio/ThresholdLowerBounds.lean}{Formalism.MainNotions.ColumnSystems.Pf.error_ratio_le_epsilon}

\begin{proof}
Write $N=|T|$, $P=p^{g_j}$, $\hat a=\hat a_j$, $e_1=e^{-\Lambda/8}$, $e_2=e^{-\Lambda/3}$,
$D=2a/(kY)$ and $\rho=(D+E_0+Q)/a$.

\emph{Facts.} The facts (a) to (g) in the proof of Lemma~\ref{pd-kappa-defects-margin} hold. They give
$k\geq5$, so $m=k-1\geq1$; $Q\geq1$, $Y\geq1$, $N_1\geq1$, $N\geq2$, $\Lambda\geq1$, $b'>0$, $kb'\leq a$ and
$b'\leq a$; and, as $Y\geq1$, $k^3Q\leq b'$ and $k^3Q^2c\leq b'$. Also $k^2\leq k^2Y^2\leq p$ by (1), and
$E_0\geq0$ as $c\geq0$. Three more facts follow.
\begin{enumerate}
\item[(h)] $E_0+Q\leq aN_1\Lambda/k$. Put $A=N_1\Lambda/k^2$. Each of $a/p^2$, $b'/p$,
  $\Lambda N_1a/(pk)$, $\Lambda b'/k$, $k^2N_1Q^2c$ and $Q$ is at most $aA$: for the first three use
  $k\leq k^2\leq p$, $kb'\leq a$ and $N_1\Lambda\geq1$; for the fourth $kb'\leq a$; for the fifth
  $k^4Q^2c\leq kb'\leq a$ and $\Lambda\geq1$; for the last $k^2Q\leq k^3Q\leq b'\leq a$. So
  $E_0+Q\leq(4C_{E1}+C_{E2}+1)aA$. Now $4C_{E1}+C_{E2}+1\leq C_H\leq C_HN\Lambda^2\leq k$ by (2), so
  $E_0+Q\leq kaA=aN_1\Lambda/k$.
\item[(i)] $E_0+Q\leq a/240$. As $C_H\geq1000$ and $b_l\leq b_u$, the number $\delta$ of
  Theorem~\ref{pd-gain-theorem-one} is at least $40N_1\Lambda/k$, and $\delta<1/6$ by (2); use (h).
\item[(j)] $a\leq a_T(s)\leq a+D$ for $s\in T$. Indeed $2J_0\leq a/(kY)$, since
  $2Y\Lambda b'\leq kY\Lambda b'\leq a$ by (e).
\end{enumerate}

\emph{Step 1.} As $p\geq1$ and $g_j\leq G'$ by (1), $1\leq P\leq Q$. By the definition of the floor,
$a-E_0<P(\hat a+1)$. By (i), $a/2+Q\leq a-E_0$, so $a/(2Q)+1\leq(a-E_0)/Q\leq(a-E_0)/P<\hat a+1$, and
$\hat a\geq a/(2Q)>0$.

\emph{Step 2.} The four terms of $\beta$ are $\geq0$, so each is $<1/6$. As $N_1k\geq5$,
$kQ\mu+5kQ\leq kQ(\mu+N_1k)<b'/6\leq a/6$, and $kE_0<a/6$.
(i) By Step 1, $2Qk(\mu+1)\leq2(kQ\mu+5kQ)<a/3<2Q\hat a$, so $k(\mu+1)<\hat a$. As $n<\mu+1$ and $m\leq k$,
$mn\leq k(\mu+1)<\hat a$, so $0\leq mn/\hat a\leq1$; and $nk<\hat a\leq\hat aN_1\Lambda$, so
$n\leq\hat aN_1\Lambda/k$.
(ii) As $kY\geq5$, $D\leq2a/5$; as $k\geq5$, $E_0<a/30$ and $Q<a/150$. So $D+E_0+Q\leq a$ and
$0\leq\rho\leq1$.
(iii) Put $S=kQ/a$. As $a\leq2Q\hat a$ by Step 1, $k\leq2S\hat a$, so $mn/\hat a\leq k(\mu+1)/\hat a\leq
2S(\mu+1)$. As $b'\leq a$ and $N_1k\geq5$, $S(\mu+5)\leq kQ(\mu+N_1k)/b'$. Also $m\rho\leq k\rho=2/Y+kE_0/a+S$.
Since $e_2\leq e_1$,
\[
 e_1+e_2+\frac{mn}{\hat a}+m\rho\leq2e_1+2S(\mu+1)+S+\frac2Y+\frac{kE_0}a\leq3\beta .
\]

\emph{Step 3.} For every integer $t\geq1$, $(1-\rho)(a_T(s_j)+t)\leq P(\hat a+t)$. Indeed $1-\rho\geq0$ and
$a_T(s_j)\leq a+D$ by (j), so
\[
 (1-\rho)\big(a_T(s_j)+t\big)\leq(1-\rho)(a+D+t)=a+D+t-(D+E_0+Q)-\rho(D+t)\leq a-E_0-Q+t,
\]
while $P(\hat a+t)=P(\hat a+1)+P(t-1)>a-E_0+t-1\geq a-E_0-Q+t$, as $P\geq1$ and $Q\geq1$.

\emph{Step 4.} Let $\theta,P'\geq0$ be reals and $x,A\geq0$ integers with $\theta(x+t)\leq P'(A+t)$ for
every integer $t\geq1$. Then $\theta^lP_l(x)\leq P'^lP_l(A)$ for every $l\geq0$, by induction on $l$:
$P_0=1$, and $(l+1)P_{l+1}(y)=(y+l+1)P_l(y)$, so multiplying the inequality for $l$ by the hypothesis at
$t=l+1$ gives the one for $l+1$. By Step 3 this applies to $\theta=1-\rho$, $x=a_T(s_j)$, $A=\hat a$ and
$P'=P$, and gives $(1-\rho)^mP_m(a_T(s_j))\leq P^mP_m(\hat a)$.

\emph{Step 5.} Put $x=mn/\hat a$, $z=m\rho$, $w=(1-\rho)^m$ and $\ell=(1-e_1)(1-e_2)(1-x)$, so
$0\leq\ell\leq1$ by Step 2(i). Bernoulli's inequality gives $w\geq1-z$, as $\rho\leq1$. Since
$\ell\geq1-e_1-e_2-x$ and $z\geq0$,
\[
 1-3\beta\leq1-e_1-e_2-x-z\leq\ell-z\leq\ell(1-z)\leq\ell w,
\]
using Step 2(iii) for the first inequality.

\emph{Step 6.} By Step 2(i) the hypothesis on $R$ applies, and gives $\ell P_m(\hat a)\Pi_j\leq R$. With
Steps 5 and 4, as $\Pi_j$, $P_m(a_T(s_j))$ and $\ell$ are $\geq0$,
\[
 (1-3\beta)P_m\big(a_T(s_j)\big)\Pi_j\leq\ell wP_m\big(a_T(s_j)\big)\Pi_j\leq\ell P^mP_m(\hat a)\Pi_j
 \leq P^mR .
\]
\end{proof}
% Lean: seg_step with Std C0 Cde CE1 CE2 1 X from std_of (thresholds 1, 1 and T2 CE1 CE2 Cc 1). The
% facts are std_of: (h) is N5 = error_ratio_le_epsilon (hk2' from chordCoeffCubic_inversionDenom_inv
% and Ck N La^2 <= k), (i) is the combination of N5, LaK and de6 in first_coefficient_facts, (j) is
% winA. Step 1 is first_coefficient_facts, Step 2 is seg_num, Step 3 is ratio_num, Step 4 is
% choose_ratio, Step 5 is one_add_mul_le_pow and four_fac, Step 6 is chain_num (with symDim_toNat).

\begin{lemma}\label{pd-kappa-defects-windows-model}
For every $C_c\geq0$ there is $T_0>0$ such that for every $C_0\geq T_0$ there is $T_1>0$ such that for
all $C_{E1},C_{E2},C_m\geq T_1$ there is $T_2>0$ such that for every $C_H\geq T_2$ there is
$C^0_\kappa>0$ with the following property for every $C_\kappa\geq C^0_\kappa$. Let $p\geq5$ be a prime
and $n'\geq2$ an even integer, and let $m$, $G'$, $N_1$, $c_{cl}$, $D_p$, $c_*$, $T$, $Y$, $Y_b$, $b'$,
$b_l$, $b_u$ and an integer $a$ satisfy hypotheses (1) to (3) of Theorem~\ref{pd-gain-theorem-one}
with $C_0$ and $C_H$. Let $c\geq0$ be an integer with $c\leq C_c(c_*+|T|k\log_pa+|T|c_{cl}+n'c_{cl})$,
and let $a_T,b_T\colon\mathbf Z/n'\to\mathbf Z_{\geq0}$ satisfy $a\leq a_T(s)\leq a+a/(kY)+2J_0$ and
$b_l\leq b_T(s)\leq b_u+2J_0$ for $s\in T$, where $J_0=\Lambda b'/k$. Form the numbers of
Definition~\ref{pd-kappa-defects-numbers} with these constants. Let an unramified quadratic Hida
setting of degree $n'$ at $p$ with ring $\mathcal R$ be given, and a real number $\mu\geq0$. For
$1\leq j\leq q$ let $r_j(\mu)$ be the sum of $\dim V_b$ over the kept blocks $b$ of the segment datum
$S_j(b_T)$, for $\alpha_i=a_T(y_{j,i})$, $\alpha_+=a_T(s_{j+1})$ and the real number $\mu$, formed over
$\mathcal R$ with $r=m$ (Definition~\ref{pd-kappa-defects-windows-blocks}). Then
\[
 (1-\kappa_P(\mu))\dim V^\times(a_T|_T,b_T|_T)\leq[k]\prod_{j=1}^qr_j(\mu),
\]
where $[k]=(p^{n'k}-1)/(p^{n'}-1)$ (Definitions~\ref{pd-kappa-split-rank}
and~\ref{pd-kappa-product-space}).
\end{lemma}
\leanref{Formalism/SharpRestriction/GrowthBounds/ColumnSystems.lean}{Formalism.MainNotions.ColumnSystems.ExtendedModelRank}
\leanref{Formalism/SharpRestriction/GrowthBounds/BinomialRatio/FamilyFourthBound.lean}{Formalism.MainNotions.ColumnSystems.Pf.RemainderSegmentBoundSecond}
\leanref{Formalism/SharpRestriction/GrowthBounds/BinomialRatio/FirstCoefficientFacts.lean}{Formalism.MainNotions.ColumnSystems.Pf.restAlt2_of}

\begin{proof}
\emph{The constants.} Let $C_c\geq0$, and let $T^S_0$, $T^S_1$ and $T^S_2$ be the thresholds of
Lemma~\ref{pd-kappa-defects-windows-model-segment}. Put $T_0=\max(T^S_0,1)$; for $C_0\geq T_0$ put
$T_1=\max(T^S_1,1)$; for $C_{E1},C_{E2},C_m\geq T_1$ put
$T_2=\max\big(T^S_2,1000(C_{E1}+C_{E2}+1)(C_c+1)\big)$; and put $C^0_\kappa=3$. Let $C_H\geq T_2$ and
$C_\kappa\geq3$, and let the data, the setting and $\mu$ be as in the statement. Write $N=|T|$ and
$\beta$ for the number of Lemma~\ref{pd-kappa-defects-windows-model-step}, so that
$\kappa_P(\mu)=C_\kappa N\beta$, and $\Pi_j$ for its product, so that
$r'_j=P_m(a_T(s_j))\Pi_j$ is the product of the factors of the $j$-th segment.

\emph{Signs.} The facts (a) to (g) in the proof of Lemma~\ref{pd-kappa-defects-margin} hold; they give
$N\geq q\geq2$, $k,Q,Y,\Lambda\geq1$, $b'>0$ and $a\geq5$, and $E_0\geq0$ as $c\geq0$. So $\beta\geq0$. Also
$p^{n'}\geq2$, so $[k]>0$, and the $r_j(\mu)$ and $r'_j$ are $\geq0$.

\emph{The dimension.} The field $\mathcal K$ has characteristic $0$, since $p$ is neither $0$ nor a unit
in $\mathcal R$. By Lemma~\ref{pd-kappa-dimensions}(1),
$\dim V^\times(a_T|_T,b_T|_T)=\prod_{s\in T}P_m(a_T(s))P_m(b_T(s))$. Every element of $T$ is exactly one
of the $s_j$ and the $y_{j,i}$, and $j\mapsto s_{j+1}$ is a bijection onto the even elements of $T$; so
$\prod_jr'_j=\dim V^\times(a_T|_T,b_T|_T)$.

\emph{The bound.} If $C_\kappa N\beta\geq1$, the left side of the claim is $\leq0$, and the right side is
$\geq0$. Otherwise, as $C_\kappa\geq3$ and $N\geq2$, $6\beta\leq C_\kappa N\beta<1$, so $\beta<1/6$. For
each $j$, Lemma~\ref{pd-kappa-defects-windows-model-segment} shows that $R=r_j(\mu)$ satisfies the
hypothesis of Lemma~\ref{pd-kappa-defects-windows-model-step}, which gives
$(1-3\beta)r'_j\leq p^{g_jm}r_j(\mu)$. All factors are $\geq0$, as $1-3\beta>0$; multiplying over $j$, and
using $\sum_jg_j=n'$,
\[
 (1-3\beta)^q\dim V^\times(a_T|_T,b_T|_T)\leq p^{n'm}\prod_{j=1}^qr_j(\mu).
\]
By Bernoulli's inequality and $q\leq N$, $(1-3\beta)^q\geq1-3q\beta\geq1-C_\kappa N\beta=1-\kappa_P(\mu)$.
Finally $p^{n'm}\leq[k]$: with $x=p^{n'}\geq2$, $x^m(x-1)=x^{m+1}-x^m\leq x^{m+1}-1$. As the dimension is
$\geq0$, the claim follows.
\end{proof}
% Lean: restAlt2_of (restSegLow_of hL) (dimHH_of h3) charZeroAt in closed_extendedAllConclusions. The
% constants are TransportedAlt.and of threshold_standard_one (std_of with thresholds 1, 1 and
% T2 CE1 CE2 Cc 1) and RestSegLow (Lemma pd-kappa-defects-windows-model-segment), with Ck0 = 3. The
% argument is rank_alternative_lower: kap_eq, the case 1 <= Ckap N brk, hB : brk < 1/6, seg_step for
% each j (Lemma pd-kappa-defects-windows-model-step), prod_letters, sum_segG, one_add_mul_le_pow with
% numSeg_le_card, and Ne_ge. The dimension is DimHHAt = dimHH_of h3 (finrank_tensor and Weyl's formula
% adm3), the book's Lemma pd-kappa-dimensions(1), with CharZeroAt = charZeroAt. The Lean's Loc data have
% real ccl, DLP and cs.

\begin{definition}\label{pd-kappa-defects-windows-lift-pairs-model}
Let $p\geq5$ be a prime, $(W,Q,\sigma,\iota)$ a Hida setting of degree $n'$ at $p$ with rings
$\mathcal R\subset\mathcal K$ and $\varpi_0$ (Definition~\ref{pd-local-model}(1)), $m\geq1$ and $k=m+1$,
$T\subset\mathbf Z/n'\mathbf Z$ and $a_T,b_T\colon\mathbf Z/n'\mathbf Z\to\mathbf Z_{\geq0}$. Form the product space
$V^\times=V^\times(a_T|_T,b_T|_T)$, its lattices $H^\times(L)$, the modules $R^\times(e)$, the cores
$\mathcal C^\times_{t,c}(e)$ and $\mathcal C'^\times_{t,c}(e)$ and the pairing $\langle\ ,\ \rangle^\times$, with the
degrees $d_s$ and $d'_s$, as in Definition~\ref{pd-kappa-product-space}.
\begin{enumerate}
\item \emph{Letters.} The letters are the pairs $(s,1)$ and $(s,0)$ with $s\in T$; $(s,1)$ is a vector letter of degree
  $d_s$ and $(s,0)$ a covector letter of degree $d'_s$, both of position $s\in\{0,\dots,n'-1\}$. Let $\mathcal K[z]$ be the
  polynomial ring in the $z_{x,c}$, for the letters $x$ and $1\leq c\leq k$, and $\mathrm{Pol}_T\subset\mathcal K[z]$ the
  polynomials of these row degrees (Definition~\ref{pd-symbols-path-generators}). For $g\in GL_k(W)$, $\Sigma_g$ is the
  lettered substitution $\Sigma^{\mathrm L}_g$ of Definition~\ref{pd-symbols-path-generators}(7) for these letters, with
  $n=k$ and $r=1$.
\item \emph{The scalar.} Put $z_T=p^{\kappa_Z}/\prod_xd_x!$, where $x$ runs over the letters, $d_x$ is the degree of $x$,
  and $\kappa_Z=k\sum_x\lfloor\log_pd_x\rfloor$ with $\lfloor\log_p0\rfloor=0$. On $\mathcal K[z]$, $[\ ,\ ]$ is the Fock
  pairing $[f,g]=\sum_\kappa\kappa!\,f_\kappa g_\kappa$ of Definition~\ref{pd-weight-lattices}(2), in the variables $z_{x,c}$.
\item \emph{Edges.} For a $W$-lattice $L\subset Q^k$ put $L^\vee=\{x\in Q^k:x\cdot y\in W\text{ for all }y\in L\}$, and for a
  pair $e$ of lattices with source $L''$ and target $L$ let $e^\vee$ be the pair with source $L^\vee$ and target
  $L''^\vee$. For a $k\times k$ matrix $h$ over $Q$ let $L_h$ be the $W$-span of its columns, and for $g\in GL_k(W)$ with
  $h\,{}^tg^{-1}$ invertible let $e_h(g)$ be the pair of the $W$-spans of the columns of $h\,{}^tg^{-1}$ and of
  $h\,{}^tg^{-1}\delta_p$, with $\delta_p$ the diagonal matrix with entries $p^{-1},1,\dots,1$, as in
  Lemma~\ref{pd-symbols-edge-basis} with $r=1$.
\item \emph{Flat theta maps.} Let $h$ be a $k\times k$ matrix over $Q$. A \emph{flat theta map} for $h$ is an
  $\mathcal R$-linear bijection $\Theta\colon\mathrm{Pol}_T\to V^\times$ such that: (a) for $F\in\mathrm{Pol}_T$,
  $F\in\mathcal R[z]$ if and only if $\Theta F\in H^\times(L_h)$; (b) for every $g\in GL_k(W)$ with $h\,{}^tg^{-1}$
  invertible, every $F\in\mathrm{Pol}_T$ and every integer $i\geq0$, $\Sigma_g(F)\in\varpi_0^i\mathcal R[z]$ if and only if
  $\Theta F\in\varpi_0^iR^\times(e_h(g))$ (Definition~\ref{pd-symbols-path-generators}(7)); and (c) every edge with source
  $L_h$ is $e_h(g)$ for some $g\in GL_k(W)$ with $h\,{}^tg^{-1}$ invertible.
\item \emph{Symbol families.} Let the data of Definition~\ref{pd-kappa-defects-numbers} be given, and for $1\leq j\leq q$
  let $\mathcal P_j$ be the ring $\mathcal P$ of Definition~\ref{pd-kappa-defects-segment-lattice}(3) for the segment datum
  $S_j(b_T)$ and $r=m$, with $s=S_j$, $\eta=\eta_j$, $P_i=p^{g_j-d_{j,i}}$ and $u_l=\eta+l$. Let
  $\iota_j\colon\mathcal P_j\to\mathcal K[z]$ be the $\mathcal K$-algebra map with
  \[
   x_{i,c}\mapsto z_{(y_{j,i},1),c+1}^{P_i},\quad y_{i,c}\mapsto z_{(y_{j,i},0),c+1}^{P_i},\quad
   x_{+,c}\mapsto z_{(s_{j+1},1),c+1},\quad \xi_c\mapsto z_{(s_{j+1},0),c+1},
  \]
  \[
   z_l\mapsto z_{(y_{j,u_l},1),1}^{P_{u_l}},\quad z'_l\mapsto z_{(y_{j,u_l},0),1}^{P_{u_l}},\quad
   z_+\mapsto z_{(s_{j+1},1),1},\quad z'_+\mapsto z_{(s_{j+1},0),1}.
  \]
  For a block $b$ of $S_j(b_T)$, with its numbers $\hat b_i$ for $\alpha_i=a_T(y_{j,i})$
  (Definition~\ref{pd-kappa-defects-segment-lattice}(1)), and $f\in\mathcal P_j$ put
  \[
   \iota_j(b,f)=\prod_{i=1}^{s}\prod_{c=1}^kz_{(y_{j,i},1),c}^{\rho_i(c)}z_{(y_{j,i},0),c}^{\rho'_i(c)}\cdot
   \prod_{i=1}^{\eta}z_{(y_{j,i},1),1}^{P_i(\hat b_i\dot-w_i)}z_{(y_{j,i},0),1}^{P_i\varsigma_i}\cdot\iota_j(f).
  \]
  For a tuple $\beta=(\beta_j)$ of blocks, $\beta_j$ a block of $S_j(b_T)$, and a tuple $f=(f_j)$ with $f_j\in\mathcal P_j$, the
  \emph{symbol family} is $\Phi(\beta,f)=\prod_{j=1}^q\iota_j(\beta_j,f_j)$.
\item \emph{Symbol modules.} For a block $b$ of $S_j(b_T)$ and a real number $\mu$, $Z_j(b;\mu)\subset\mathcal P_j$ is the
  $\mathcal R$-span of the generators of all levels of $b$ (Definition~\ref{pd-kappa-defects-segment-lattice}(5)), formed
  over $\mathcal R$ with $\alpha_i=a_T(y_{j,i})$, $\alpha_+=a_T(s_{j+1})$, $\mu$ and $r=m$.
\end{enumerate}
\end{definition}
\leanref{Formalism/SharpRestriction/GrowthBounds/ColumnSystems.lean}{Formalism.MainNotions.ColumnSystems.hhLetters}
\leanref{Formalism/SharpRestriction/GrowthBounds/ColumnSystems.lean}{Formalism.MainNotions.ColumnSystems.hhDegs}
\leanref{Formalism/SharpRestriction/GrowthBounds/CoarseMultiplicity/RestBlockForm.lean}{Formalism.MainNotions.ColumnSystems.Pf.kapZ}
\leanref{Formalism/SharpRestriction/GrowthBounds/CoarseMultiplicity/RestBlockForm.lean}{Formalism.MainNotions.ColumnSystems.Pf.zTen}
\leanref{Formalism/SharpRestriction/GrowthBounds/ObjPol.lean}{Formalism.MainNotions.G.fock}
\leanref{Formalism/SharpRestriction/GrowthBounds/UnramifiedQuadraticSetting.lean}{Formalism.MainNotions.G.latDual}
\leanref{Formalism/SharpRestriction/GrowthBounds/UnramifiedQuadraticSetting.lean}{Formalism.MainNotions.G.Edge.dual}
\leanref{Formalism/SharpRestriction/GrowthBounds/UnramifiedQuadraticSetting.lean}{Formalism.MainNotions.G.colSpan}
\leanref{Formalism/SharpRestriction/GrowthBounds/BinomialRatio/FinalTheorems.lean}{Formalism.MainNotions.ColumnSystems.Pf.IsFlatTheta}
\leanref{Formalism/SharpRestriction/GrowthBounds/BinomialRatio/FinalTheorems.lean}{Formalism.MainNotions.ColumnSystems.Pf.symVar}
\leanref{Formalism/SharpRestriction/GrowthBounds/BinomialRatio/FinalTheorems.lean}{Formalism.MainNotions.ColumnSystems.Pf.symLev}
\leanref{Formalism/SharpRestriction/GrowthBounds/BinomialRatio/FinalTheorems.lean}{Formalism.MainNotions.ColumnSystems.Pf.digitMon}
\leanref{Formalism/SharpRestriction/GrowthBounds/BinomialRatio/FinalTheorems.lean}{Formalism.MainNotions.ColumnSystems.Pf.sTopMon}
\leanref{Formalism/SharpRestriction/GrowthBounds/BinomialRatio/FinalTheorems.lean}{Formalism.MainNotions.ColumnSystems.Pf.ioBlock}
\leanref{Formalism/SharpRestriction/GrowthBounds/BinomialRatio/FinalTheorems.lean}{Formalism.MainNotions.ColumnSystems.Pf.symFam}
\leanref{Formalism/SharpRestriction/GrowthBounds/CoarseMultiplicity/RestBlockForm.lean}{Formalism.MainNotions.ColumnSystems.Pf.InSSZ}
% Lean: the letters are hhLetters TT (vec = the Bool, pos = the value of s) with degrees hhDegs TT aT bT = hhDeg, and
%   Pol_T is Pol (KK O) (m + 1) (hhDegs TT aT bT), taken as an O-module. Sigma_g is liftSubst p n' O w 1 (hhLetters
%   TT) g; the Lean's g is a matrix over the Witt vectors WF p n', moved to W by a Witt identification phi (WittIdent,
%   matQ phi g and algOfIdent), and the book writes the matrix over W. z_T is zTen with kapZ (Nat.log). e^vee is
%   Edge.dual = (latDual tgt, latDual src); e_h(g) is IsEdgeBasis p 1 e (h * ((matQ phi g)^-1)^T). A flat theta map
%   is IsFlatTheta (with InA for (a) and InWjRg for (b)). iota_j(b, f) is ioBlock: digitMon, sTopMon (the factors
%   with i < sigma) and the substitution aeval of symVar raised to symLev; Phi is symFam; Z_j(b; mu) is
%   (blockData p (seg j) (segATop j) (segAPlus j) mu b).SSZ O, the O-span of BlockData.gens over the levels LevelOK.

\begin{lemma}\label{pd-kappa-defects-windows-lift-pairs-flat-dual}
Let $p\geq5$ be a prime, $(W,Q,\sigma,\iota)$ a Hida setting of degree $n'$ at $p$ and $k\geq1$, and use the notation $L^\vee$,
$e^\vee$, $L_h$, $e_h(g)$ and $\delta_p$ of Definition~\ref{pd-kappa-defects-windows-lift-pairs-model}(3). Let $e$ be an edge
(Definition~\ref{pd-projection-property}(1) with $r=1$).
\begin{enumerate}
\item If $\ell\in GL_k(Q)$, then $L_\ell^\vee=L_{{}^t\ell^{-1}}$ and $(L_\ell^\vee)^\vee=L_\ell$.
\item There is $h\in GL_k(Q)$ with $e=e_h(1)$.
\item If $h\in GL_k(Q)$ and $e=e_h(1)$, then $e^\vee=e_{h'}(1)$ for $h'={}^th^{-1}\delta_p^{-1}$, $e^\vee$ is an edge, and
  $(e^\vee)^\vee=e$.
\end{enumerate}
\end{lemma}
\leanref{Formalism/LocalToGlobal/GroupLattices/FockModels/TestSideEdge.lean}{Formalism.L2.G.Model.latDual_colSpan}
\leanref{Formalism/LocalToGlobal/GroupLattices/FockModels/TestSideEdge.lean}{Formalism.L2.G.Model.latDual_latDual}
\leanref{Formalism/LocalToGlobal/GroupLattices/FockModels/TestSideEdge.lean}{Formalism.L2.G.Model.edgeBasis}
\leanref{Formalism/LocalToGlobal/GroupLattices/FockModels/TestSideEdge.lean}{Formalism.L2.G.Model.dual_edgeBasis}
\leanref{Formalism/LocalToGlobal/GroupLattices/FockModels/TestSideEdge.lean}{Formalism.L2.G.Model.dual_edge}

\begin{proof}
(1) Let $x\in Q^k$. For $y=\ell a$ with $a\in W^k$, $x\cdot y=({}^t\ell x)\cdot a$. So $x\in L_\ell^\vee$ if and only if
${}^t\ell x\in W^k$, that is, $x={}^t\ell^{-1}b$ with $b\in W^k$; hence $L_\ell^\vee=L_{{}^t\ell^{-1}}$. Applying this to
${}^t\ell^{-1}$, whose inverse transpose is $\ell$, gives $(L_\ell^\vee)^\vee=L_\ell$.

(2) Let $L''$ and $L$ be the source and target of $e$. As in the proof of Lemma~\ref{pd-edge-basis}(1), the elementary divisor
theorem gives a $W$-basis $b_1,\dots,b_k$ of $L''$ such that $p^{-1}b_1,b_2,\dots,b_k$ is a basis of $L$. The matrix $h$ with
columns $b_1,\dots,b_k$ is invertible, $L_h=L''$ and $L_{h\delta_p}=L$; so $e=e_h(1)$.

(3) Let $L''=L_h$ and $L=L_{h\delta_p}$ be the source and target of $e$. By (1), the source $L^\vee$ of $e^\vee$ is
$L_{{}^th^{-1}\delta_p^{-1}}=L_{h'}$, and its target $L''^\vee$ is $L_{{}^th^{-1}}=L_{h'\delta_p}$; so $e^\vee=e_{h'}(1)$. It is an
edge by Lemma~\ref{pd-symbols-edge-basis}(2), with $\gamma=h'$, $r=1$ and $g=1$. By (1), $(L^\vee)^\vee=L$ and
$(L''^\vee)^\vee=L''$, so $(e^\vee)^\vee=e$.
\end{proof}
% Lean: (1) is latDual_colSpan and latDual_latDual; (2) is edgeBasis (restEdgeBasis = lc_main, from the Smith normal form
%   lc_snf); (3) is dual_edgeBasis (the basis (b^-1)^T * dE (m + 1) 1 p, with lb_D_dE and invT_mul_D) and dual_edge
%   (isEdge_of_basis and latDual_latDual).

\begin{lemma}\label{pd-kappa-defects-windows-lift-pairs-flat-theta}
Let $p\geq5$ be a prime, $(W,Q,\sigma,\iota)$ a Hida setting of degree $n'$ at $p$ with rings $\mathcal R\subset\mathcal K$ and
$\varpi_0$, $m\geq1$, $T\subset\mathbf Z/n'\mathbf Z$ and $a_T,b_T\colon\mathbf Z/n'\mathbf Z\to\mathbf Z_{\geq0}$, with the notation
of Definition~\ref{pd-kappa-defects-windows-lift-pairs-model}; write $s$ also for its representative in $\{0,\dots,n'-1\}$. Let
$h\in GL_k(Q)$ and $h_s=\iota\sigma^s(h)$, applied to the entries, and let $\operatorname{adj}(x_{s,0})$ be the adjugate of the
matrix $x_{s,0}$ of variables of Definition~\ref{pd-kappa-product-space}. Let $\Theta_h$ be the restriction to
$\mathrm{Pol}_T$ of the $\mathcal K$-algebra map $\mathcal K[z]\to\mathcal K[x_{s,\epsilon}]$ with
\[
 z_{(s,1),c}\mapsto(h_s^{-1}x_{s,1})_{c,1},\qquad z_{(s,0),c}\mapsto\big(\operatorname{adj}(x_{s,0})\,h_s\big)_{k,c}.
\]
Then $\Theta_h$ is a flat theta map for $h$.
\end{lemma}
\leanref{Formalism/SharpRestriction/GrowthBounds/CoarseMultiplicity/MulMatQOne.lean}{Formalism.MainNotions.ColumnSystems.Pf.TF.flatTheta}
\leanref{Formalism/SharpRestriction/GrowthBounds/CoarseMultiplicity/MulMatQOne.lean}{Formalism.MainNotions.ColumnSystems.Pf.TF.Avec_eq}
\leanref{Formalism/LocalToGlobal/GroupLattices/FockModels/FlatModels/LaTwTranspose.lean}{Formalism.L2.G.Model.thF}
\leanref{Formalism/LocalToGlobal/GroupLattices/FockModels/FlatModels/LaTwTranspose.lean}{Formalism.L2.G.Model.gen_equiv}
\leanref{Formalism/LocalToGlobal/GroupLattices/FockModels/FlatModels/LaTwTranspose.lean}{Formalism.L2.G.Model.gen_A}
\leanref{Formalism/LocalToGlobal/GroupLattices/FockModels/FlatModels/LaTwTranspose.lean}{Formalism.L2.G.Model.gen_R}
\leanref{Formalism/LocalToGlobal/GroupLattices/FockModels/FlatModels/LetterWtAnti.lean}{Formalism.L2.G.Model.tensor_lattice_twist_eq}
\leanref{Formalism/LocalToGlobal/GroupLattices/FockModels/FlatModels/LetterWtAnti.lean}{Formalism.L2.G.Model.tens_Le}

\begin{proof}
Give the vector letter $(s,1)$ even parity and the covector letter $(s,0)$ odd parity. The proof of
Lemma~\ref{pd-symbols-flat-exists} applies with the letters in place of the positions: the variables $z_x$ in place of
$X_s$, the matrices $x_{s,1}$ and $x_{s,0}$ in place of $x_s$, the degree $d_x$ of the letter in place of $d(s)$, the parity
of the letter in place of that of the position, $\gamma=h$ and $r=1$, so that $1\leq r<k$. Both letters at $s$ use
$\iota\sigma^s$ (Definition~\ref{pd-kappa-product-space}), so $\Theta_h$ is the map $\Theta_\gamma$ of that lemma for this
lettering. We list what each step uses.
\begin{itemize}
\item Step 2. The weights of $V^\times$ at $(s,1)$ and $(s,0)$ are $(d_s,0,\dots,0)$ and $(0,\dots,0,-d'_s)$, which are the
  weights $\lambda$ of Lemma~\ref{pd-symbols-flat-letter} for $d_s$ with even parity and for $d'_s$ with odd parity. As
  $k\geq2$, the integer $A$ of Definition~\ref{pd-projection-property}(1) with $r=1$ is the sum of the first entries of these
  weights, $\sum_{s\in T}d_s=A^\times$.
\item Steps 1 and 3 show that $\Theta_h$ is a $\mathcal K$-linear bijection of $\mathrm{Pol}_T$ onto $V^\times$, and Step 4,
  with $L''=L_h$, gives property (a) of a flat theta map.
\item Steps 5 and 6, for $g\in GL_k(W)$ and the pair $e_h(g)$, give property (b). There $\Sigma_g$ is the lettered
  substitution, which replaces $z_x$ by $z_x\,\delta\,g_s^{-1}$ for $x=(s,1)$ and by $z_x\,\delta'\,{}^tg_s$ for $x=(s,0)$, with
  the matrices $\delta$ and $\delta'$ of Definition~\ref{pd-symbols-twist}(1) for $r=1$.
\end{itemize}
Property (c) is Lemma~\ref{pd-symbols-edge-basis}(1) with $\gamma=h$ and $r=1$; for $g\in GL_k(W)$ the matrix
$h\,{}^tg^{-1}$ is invertible, as $h$ is.
\end{proof}
% Lean: flatTheta hS phi hphi m TT aT bT hm h hh: Theta = gen_equiv (thF hS m vecH (la_tw h), with the lettering vecH TT =
%   the Bool of the letter) composed with VS = VHH; (a) is gen_A with tensor_lattice_twist_eq, (b) is gen_R (lift_eq,
%   Bm_eq, key_iff) with tens_Le and mem_scale2, and Avec_eq for A^x = AHH; (c) is (edgePlane ...).1. These are the
%   lemmas, general in the lettering, that Lemma pd-symbols-flat-exists follows.

\begin{lemma}\label{pd-kappa-defects-windows-lift-pairs-flat-forms}
Let $p\geq5$ be a prime, $(W,Q,\sigma,\iota)$ a Hida setting of degree $n'$ at $p$ with rings $\mathcal R\subset\mathcal K$ and
$\varpi_0$, $m\geq1$, $T\subset\mathbf Z/n'\mathbf Z$ and $a_T,b_T\colon\mathbf Z/n'\mathbf Z\to\mathbf Z_{\geq0}$, with the notation
of Definition~\ref{pd-kappa-defects-windows-lift-pairs-model}, and let $\Theta_1$ be the flat theta map for the identity
matrix of Lemma~\ref{pd-kappa-defects-windows-lift-pairs-flat-theta}. Let $\Xi\colon V^\times\to(V^\times)^\vee$ be the
$\mathcal K$-linear map with
\[
 \Xi(v)(u)=z_T\,\big[\Theta_1^{-1}(u),\Theta_1^{-1}(v)\big]\qquad(u,v\in V^\times).
\]
For a family $N=(N_x)$ of $k\times k$ matrices over $\mathcal K$, one for each letter $x$, and $F\in\mathcal K[z]$, let $F^N$ be
$F$ with the column ${}^t(z_{x,1},\dots,z_{x,k})$ replaced by $N_x\,{}^t(z_{x,1},\dots,z_{x,k})$ for every letter $x$. Let
$i\in\mathbf Z$ and $v\in V^\times$.
\begin{enumerate}
\item $[F^N,G]=[F,G^{{}^tN}]$ for all $F,G\in\mathcal K[z]$, where ${}^tN=({}^tN_x)$.
\item If $L\subset Q^k$ is a $W$-lattice and $v\in\varpi_0^iH^\times(L^\vee)$, then $\Xi(v)\in\varpi_0^iH^\times(L)^\vee$.
\item If $e$ is an edge and $v\in\varpi_0^iR^\times(e^\vee)$, then $\Xi(v)\in\varpi_0^iR'^\times(e)$.
\end{enumerate}
\end{lemma}
\leanref{Formalism/SharpRestriction/GrowthBounds/CoarseMultiplicity/LaTwInvT.lean}{Formalism.MainNotions.ColumnSystems.Pf.TF.XiH}
\leanref{Formalism/SharpRestriction/GrowthBounds/CoarseMultiplicity/FockRAddRight.lean}{Formalism.MainNotions.ColumnSystems.Pf.TF.exists_thK}
\leanref{Formalism/SharpRestriction/GrowthBounds/CoarseMultiplicity/FockRAddRight.lean}{Formalism.MainNotions.ColumnSystems.Pf.TF.TY}
\leanref{Formalism/SharpRestriction/GrowthBounds/CoarseMultiplicity/FockRAddRight.lean}{Formalism.MainNotions.ColumnSystems.Pf.TF.fockR_Lin}
\leanref{Formalism/SharpRestriction/GrowthBounds/CoarseMultiplicity/FockRAddRight.lean}{Formalism.MainNotions.ColumnSystems.Pf.TF.multDiv_kapFl}
\leanref{Formalism/SharpRestriction/GrowthBounds/CoarseMultiplicity/FockRAddRight.lean}{Formalism.MainNotions.ColumnSystems.Pf.TF.sandwich}
\leanref{Formalism/SharpRestriction/GrowthBounds/CoarseMultiplicity/LaTwInvT.lean}{Formalism.MainNotions.ColumnSystems.Pf.TF.FD_HL}
\leanref{Formalism/SharpRestriction/GrowthBounds/CoarseMultiplicity/LaTwInvT.lean}{Formalism.MainNotions.ColumnSystems.Pf.TF.FD_edge}

\begin{proof}
As in the proof of Lemma~\ref{pd-symbols-test-fock-position}, $\mathcal K$ has characteristic $0$ and $p$ lies in the maximal
ideal of the local ring $\mathcal R$. So Lemma~\ref{pd-symbols-flat-real-multinomial} applies with the letters in place of the
elements of $\mathbf Z/n'\mathbf Z$, as its statement and proof use only that the indices form a finite set; for the degrees
$d_x$ of the letters its number $p^{\kappa_F}/\prod_xd_x!$ is $z_T$. By
Lemma~\ref{pd-kappa-defects-windows-lift-pairs-flat-theta}, $\Theta_1$ is a $\mathcal K$-linear bijection of $\mathrm{Pol}_T$ onto
$V^\times$, so $\Xi$ is defined. For a family $h=(h_x)$ in $GL_k(\mathcal K)$, one for each letter, put $a_x(h)=h_x^{-1}$ for a
vector letter and $a_x(h)={}^th_x$ for a covector letter, and $\Theta_h(F)=\Theta_1(F^{a(h)})$ for $F\in\mathrm{Pol}_T$. For
$h\in GL_k(Q)$ and $h_x=\iota\sigma^s(h)$ at $x=(s,\epsilon)$ this is the map of
Lemma~\ref{pd-kappa-defects-windows-lift-pairs-flat-theta}, since $(h_s^{-1}x_{s,1})_{c,1}=\sum_b(h_s^{-1})_{cb}(x_{s,1})_{b,1}$
and $(\operatorname{adj}(x_{s,0})h_s)_{k,c}=\sum_b({}^th_s)_{cb}\operatorname{adj}(x_{s,0})_{k,b}$. Also $(F^N)^M=F^{NM}$, with
$(NM)_x=N_xM_x$, and $F^N$ has the row degrees of $F$.

(1) The proof of Lemma~\ref{pd-kappa-gram-fock}(1) applies, with the columns of the $z_x$ in place of those of $Z$ and the
matrix $N_x$ depending on the column: $[Fz_v,G]=[F,\partial_vG]$ for every variable $z_v$, and the chain rule gives
$\partial_{(x,b)}(G^{{}^tN})=\sum_a(N_x)_{ba}(\partial_{(x,a)}G)^{{}^tN}$.

\emph{A pairing of lattices.} Let $h=(h_x)$ be a family in $GL_k(\mathcal K)$, $h'=({}^th_x^{-1})$, and
$F,G\in\mathcal R[z]\cap\mathrm{Pol}_T$. Then $\Xi(\Theta_{h'}(G))(\Theta_h(F))\in\mathcal R$. Indeed, for every letter
$a_x(h')={}^ta_x(h)^{-1}$: for a vector letter both are ${}^th_x$, for a covector letter both are $h_x^{-1}$. By (1),
\[
 \Xi\big(\Theta_{h'}(G)\big)\big(\Theta_h(F)\big)=z_T\big[F^{a(h)},G^{a(h')}\big]=z_T\big[F,G^{a(h')\,{}^ta(h)}\big]=z_T[F,G],
\]
which lies in $\mathcal R$ by Lemma~\ref{pd-symbols-flat-real-multinomial}(2), as $G$ has the row degrees $d_x$.

(2) Let $\ell\in GL_k(Q)$ with $L=L_\ell$, and $\ell_s=\iota\sigma^s(\ell)$. By
Lemma~\ref{pd-kappa-defects-windows-lift-pairs-flat-dual}(1), $L^\vee=L_{{}^t\ell^{-1}}$, and
$\iota\sigma^s({}^t\ell^{-1})={}^t\ell_s^{-1}$. By Step~4 of the proof of Lemma~\ref{pd-symbols-flat-exists}, in the form used
in the proof of Lemma~\ref{pd-kappa-defects-windows-lift-pairs-flat-theta}, $H^\times(L)=\Theta_h(\mathcal R[z]\cap\mathrm{Pol}_T)$
for the family $h_x=\ell_s$ at $x=(s,\epsilon)$, and $H^\times(L^\vee)=\Theta_{h'}(\mathcal R[z]\cap\mathrm{Pol}_T)$. Write
$v=\varpi_0^i\Theta_{h'}(G)$ with $G\in\mathcal R[z]\cap\mathrm{Pol}_T$. By the pairing of lattices $\Xi(\Theta_{h'}(G))$ takes
$H^\times(L)$ into $\mathcal R$, so $\Xi(v)\in\varpi_0^iH^\times(L)^\vee$.

(3) By Lemma~\ref{pd-kappa-defects-windows-lift-pairs-flat-dual}, $e=e_h(1)$ for some $h\in GL_k(Q)$, and $e^\vee=e_{h'}(1)$ for
$h'={}^th^{-1}\delta_p^{-1}$. Let $D$ be the diagonal matrix with entries $\varpi_0^{-1},1,\dots,1$, and put $h_s=\iota\sigma^s(h)$
and $h'_s=\iota\sigma^s(h')={}^th_s^{-1}\delta_p^{-1}$. As in Step~5 of the proof of Lemma~\ref{pd-symbols-flat-exists},
$H^\times(e)=\Theta_{\hat h}(\mathcal R[z]\cap\mathrm{Pol}_T)$ for the family $\hat h_x=h_sD$, and
$H^\times(e^\vee)=\Theta_{\hat h'}(\mathcal R[z]\cap\mathrm{Pol}_T)$ for $\hat h'_x=h'_sD$. As $p=\varpi_0^2$,
$\delta_p^{-1}D=D^{-1}$, so $\hat h'_x={}^th_s^{-1}D^{-1}={}^t\hat h_x^{-1}$. Since $R^\times(e^\vee)=\varpi_0^{-A^\times}H^\times(e^\vee)$
and $R'^\times(e)=\varpi_0^{-A^\times}H^\times(e)^\vee$, the argument of (2), with $i-A^\times$ in place of $i$, gives (3).
\end{proof}
% Lean: Xi is XiH thK zK = TY thK zK, with thK = exists_thK (Th0 over Kf, the map Theta_1) and zK = zTen; TY_apply gives
%   the formula, with fockR symmetric (tsact_symm). (1) is fockR_Lin (Lin_X, Lin_chain); the pairing of lattices is
%   sandwich (mem_tens_nabla, TY_apply, fockR_Lin_dual with Aa_dual, int_pair with MultDiv from multDiv_kapFl and the
%   maximal ideal from settingCZ). (2) is FD_HL (HHL_iff, la_tw_invT) and (3) is FD_edge (HHLe_iff with edgeLat_of,
%   dual_edgeBasis, dualM), both stated for every scalar c in Kf.

\begin{lemma}\label{pd-kappa-defects-windows-lift-pairs-flat}
Let $p\geq5$ be a prime, $(W,Q,\sigma,\iota)$ a Hida setting of degree $n'$ at $p$ with rings
$\mathcal R\subset\mathcal K$ and $\varpi_0$, $m\geq1$, $T\subset\mathbf Z/n'\mathbf Z$ and
$a_T,b_T\colon\mathbf Z/n'\mathbf Z\to\mathbf Z_{\geq0}$, with the notation of
Definition~\ref{pd-kappa-defects-windows-lift-pairs-model}, and let $e$ be an edge
(Definition~\ref{pd-projection-property}(1) with $r=1$). Then there are $k\times k$ matrices $h$ and $h'$ over $Q$, a flat
theta map $\Theta$ for $h$, a flat theta map $\Theta'$ for $h'$ and a $\mathcal K$-linear map
$\Xi\colon V^\times\to(V^\times)^\vee$ such that:
\begin{enumerate}
\item $e=e_h(1)$, and $e^\vee$ is an edge with $e^\vee=e_{h'}(1)$;
\item $\Xi(\mathcal C^\times_{t,c}(e^\vee))\subset\mathcal C'^\times_{t,c}(e)$ for all integers $t,c\geq0$;
\item $\Sigma_1(\mathrm{Pol}_T)\subset\mathrm{Pol}_T$;
\item if $F,F'\in\mathrm{Pol}_T$ and $y\in\mathcal R$ satisfy $z_T[\Sigma_1F,\Sigma_1F']=y$, then
  $\langle\Theta F,\Xi(\Theta'F')\rangle^\times=y$.
\end{enumerate}
\end{lemma}
\leanref{Formalism/SharpRestriction/GrowthBounds/CoarseMultiplicity/PairedLiftAtOfIdentity.lean}{Formalism.MainNotions.ColumnSystems.Pf.TestFlatAt1}
\leanref{Formalism/SharpRestriction/GrowthBounds/CoarseMultiplicity/PairedLiftAtOfIdentity.lean}{Formalism.MainNotions.ColumnSystems.Pf.testFlatAt1}

\begin{proof}
Let $L''$ and $L$ be the source and target of $e$. By Lemma~\ref{pd-kappa-defects-windows-lift-pairs-flat-dual}(2) there is
$h\in GL_k(Q)$ with $e=e_h(1)$, and by its part (3), with $h'={}^th^{-1}\delta_p^{-1}$, $e^\vee$ is an edge with
$e^\vee=e_{h'}(1)$. This is (1). Let $\Theta=\Theta_h$ and $\Theta'=\Theta_{h'}$ be the flat theta maps of
Lemma~\ref{pd-kappa-defects-windows-lift-pairs-flat-theta}, and $\Xi$ the map of
Lemma~\ref{pd-kappa-defects-windows-lift-pairs-flat-forms}. We use the notation $F^N$, $a(h)$ and
$\Theta_h(F)=\Theta_1(F^{a(h)})$ of the proof of that lemma, with $h_x=\iota\sigma^s(h)$ at $x=(s,\epsilon)$, and let $D$ be the
diagonal matrix with entries $\varpi_0^{-1},1,\dots,1$.

(2) Let $t,c\geq0$ and $v\in\mathcal C^\times_{t,c}(e^\vee)$ (Definition~\ref{pd-kappa-split-rank}(1)). The source of $e^\vee$ is
$L^\vee$, so $v\in\varpi_0^{-t}H^\times(L^\vee)$, and Lemma~\ref{pd-kappa-defects-windows-lift-pairs-flat-forms}(2) gives
$\Xi(v)\in\varpi_0^{-t}H^\times(L)^\vee$. As $v\in R^\times(e^\vee)$, part (3) of that lemma with $i=0$ gives
$\Xi(v)\in R'^\times(e)$. Let $e''\neq e$ be an edge with target $L$. By
Lemma~\ref{pd-kappa-defects-windows-lift-pairs-flat-dual}(2) and~(3), $e''^\vee$ is an edge, with source $L^\vee$, and
$e''^\vee\neq e^\vee$, since $(e''^\vee)^\vee=e''$ and $(e^\vee)^\vee=e$. So $v\in\varpi_0^cR^\times(e''^\vee)$, and part (3) of
Lemma~\ref{pd-kappa-defects-windows-lift-pairs-flat-forms} gives $\Xi(v)\in\varpi_0^cR'^\times(e'')$. Hence
$\Xi(v)\in\mathcal C'^\times_{t,c}(e)$.

(3) The substitution $\Sigma_1$ replaces the row $z_x$ by $z_x\delta$ for a vector letter and by $z_x\delta'$ for a covector
letter, where $\delta$ and $\delta'$ are the diagonal matrices with entries $1,\varpi_0,\dots,\varpi_0$ and $\varpi_0,1,\dots,1$.
So $\Sigma_1F=F^N$ with $N_x=\delta$, respectively $N_x=\delta'$, and $\Sigma_1F$ has the row degrees of $F$.

(4) Let $F,F'\in\mathrm{Pol}_T$, and put $P=F^{a(h)}$ and $P'=F'^{a(h')}$, so that $\Theta F=\Theta_1(P)$ and
$\Theta'F'=\Theta_1(P')$. As $\delta=\varpi_0D$ and $\delta'=D^{-1}$, and $F$ has degree $d_x$ in $z_x$,
$\Sigma_1F=\varpi_0^{A^\times}F^{N'}$ with $N'_x=D$ for $x=(s,1)$ and $N'_x=D^{-1}$ for $x=(s,0)$; here
$A^\times=\sum_{s\in T}d_s$ is the sum of the degrees of the vector letters. Then $F^{N'}=P^M$ with $M_x=a_x(h)^{-1}N'_x$, that
is, $M_x=h_sD$ for $x=(s,1)$ and $M_x={}^th_s^{-1}D^{-1}$ for $x=(s,0)$. In the same way
$\Sigma_1F'=\varpi_0^{A^\times}P'^{M'}$ with $M'_x=h'_sD$ for $x=(s,1)$ and $M'_x={}^th'^{-1}_sD^{-1}$ for $x=(s,0)$, where
$h'_s={}^th_s^{-1}\delta_p^{-1}$. As $p=\varpi_0^2$, $\delta_p^{-1}D=D^{-1}$ and $\delta_pD^{-1}=D$; so
$M'_x={}^th_s^{-1}D^{-1}$ for $x=(s,1)$ and $M'_x=h_s\delta_pD^{-1}=h_sD$ for $x=(s,0)$, and in both cases $M'_x={}^tM_x^{-1}$.
By Lemma~\ref{pd-kappa-defects-windows-lift-pairs-flat-forms}(1),
\[
 [\Sigma_1F,\Sigma_1F']=\varpi_0^{2A^\times}\big[P^M,P'^{{}^tM^{-1}}\big]=p^{A^\times}\big[P,P'^{{}^tM^{-1}\,{}^tM}\big]
 =p^{A^\times}[P,P'] .
\]
By the definition of $\Xi$ and of the pairing (Definition~\ref{pd-kappa-product-space}),
$\langle\Theta F,\Xi(\Theta'F')\rangle^\times=\Xi(\Theta_1P')(p^{A^\times}\Theta_1P)=p^{A^\times}z_T[P,P']$. Hence
$\langle\Theta F,\Xi(\Theta'F')\rangle^\times=z_T[\Sigma_1F,\Sigma_1F']$, which gives (4).
\end{proof}
% Lean: testFlatAt1: edgeBasis and dual_edgeBasis for h = b and h' (with mul_matQ_one at g = 1), dual_edge for (1);
%   flatTheta for Theta and Theta'; exists_thK and multDiv_kapFl (with settingCZ for p in the maximal ideal) for
%   Xi = XiH thK (zK ...); (2) is Xi_CC (FD_HL for the source part, FD_edge at scale 1 for R'(e), and FD_edge with
%   dual_edge for the other edges); (3) is lift_Pol (lift_eq and Lin_Pol); (4) is pairing: lift_one (Sigma_1 F =
%   w^A Lin (Aa (b D)) (Lin (Aa b^-1) F), through lift_eq and Bm_eq), dualM and Aa_dual' for M' = M^-T, fockR_Lin_dual,
%   XiH_apply and p^A = w^A w^A (Avec_eq, p_eq).

\begin{lemma}\label{pd-kappa-defects-windows-lift-pairs-source-symbols}
Let $p\geq5$ be a prime, $(W,Q,\sigma,\iota)$ a Hida setting of degree $n'$ at $p$ with rings $\mathcal R\subset\mathcal K$ and
$\varpi_0$, $m\geq1$, $T\subset\mathbf Z/n'\mathbf Z$ with $q\geq2$ even elements, and
$a_T,b_T\colon\mathbf Z/n'\mathbf Z\to\mathbf Z_{\geq0}$; let the data and constants of
Definition~\ref{pd-kappa-defects-numbers} be given, and use the notation of
Definition~\ref{pd-kappa-defects-windows-lift-pairs-model}. Let $1\leq j\leq q$, let $b$ and $b'$ be blocks of $S_j(b_T)$,
$\mu$ and $\mu'$ real numbers, $f\in Z_j(b;\mu)$ and $f'\in Z_j(b';\mu')$, and write $o_+$ for the entry $o_+$ of $b$
(Definition~\ref{pd-kappa-defects-segment-lattice}(1)).
\begin{enumerate}
\item $\iota_j(b,f)\in\mathcal R[z]$, and $\iota_j(b,f)$ involves only the variables of the letters $(y_{j,i},\epsilon)$,
  $1\leq i\leq S_j$, and $(s_{j+1},\epsilon)$.
\item If $b$ is kept and $o_+\leq a_T(s_{j+1})$, then $\iota_j(b,f)$ has degree $d_x$ in $z_x$ for each letter $x$ of (1).
\item If $b$ and $b'$ are kept, $b\neq b'$ and $\mu=\mu'$, then no monomial occurs in both $\iota_j(b,f)$ and
  $\iota_j(b',f')$, and $[\iota_j(b,f),\iota_j(b',f')]=0$.
\item Let $\beta=(\beta_j)$ and $\beta'=(\beta'_j)$ be tuples of blocks, $\beta_j$ and $\beta'_j$ blocks of $S_j(b_T)$, and
  $f=(f_j)$ and $f'=(f'_j)$ tuples with $f_j\in Z_j(\beta_j;\mu)$ and $f'_j\in Z_j(\beta'_j;\mu')$ for every $j$. Then
  $\Phi(\beta,f)\in\mathcal R[z]$ and
  \[
   [\Phi(\beta,f),\Phi(\beta',f')]=\prod_{j=1}^q\big[\iota_j(\beta_j,f_j),\iota_j(\beta'_j,f'_j)\big].
  \]
  If every $\beta_j$ is kept and its entry $o_+$ is at most $a_T(s_{j+1})$, then $\Phi(\beta,f)\in\mathrm{Pol}_T$.
\end{enumerate}
\end{lemma}
\leanref{Formalism/SharpRestriction/GrowthBounds/BinomialRatio/SymInt.lean}{Formalism.MainNotions.ColumnSystems.Pf.symIntAt}
\leanref{Formalism/SharpRestriction/GrowthBounds/BinomialRatio/SymInt.lean}{Formalism.MainNotions.ColumnSystems.Pf.ssz_le_intAlg}
\leanref{Formalism/SharpRestriction/GrowthBounds/CoarseMultiplicity/IoBlockSupported.lean}{Formalism.MainNotions.ColumnSystems.Pf.ioBlock_supported}
\leanref{Formalism/SharpRestriction/GrowthBounds/CoarseMultiplicity/WeightWLet.lean}{Formalism.MainNotions.ColumnSystems.Pf.TD_wLet}
\leanref{Formalism/SharpRestriction/GrowthBounds/CoarseMultiplicity/SymbolFormOfOrthogonality.lean}{Formalism.MainNotions.ColumnSystems.Pf.symPolAt_of}
\leanref{Formalism/SharpRestriction/GrowthBounds/CoarseMultiplicity/IoBlockEqSum.lean}{Formalism.MainNotions.ColumnSystems.Pf.sep}
\leanref{Formalism/SharpRestriction/GrowthBounds/CoarseMultiplicity/IoBlockEqSum.lean}{Formalism.MainNotions.ColumnSystems.Pf.orthAt}
\leanref{Formalism/SharpRestriction/GrowthBounds/CoarseMultiplicity/IoBlockSupported.lean}{Formalism.MainNotions.ColumnSystems.Pf.prodAt}

\begin{proof}
Write $s=S_j$, $\eta=\eta_j$, $P_i=p^{g_j-d_{j,i}}$, $\alpha_i=a_T(y_{j,i})$, $b_i=b_T(y_{j,i})$, $\alpha_+=a_T(s_{j+1})$ and
$b_+=b_T(s_{j+1})$, and write $\rho_i,\rho'_i,o_i,w_i,\varsigma_i,o_+$ for the entries of $b$ and $\hat b_i$, $\hat\alpha_i$ for its
numbers (Definition~\ref{pd-kappa-defects-segment-lattice}(1)). The map $\iota_j$ sends each variable of $\mathcal P_j$ to a
power with positive exponent of a variable $z_{x,c}$ with $x$ one of the letters of (1), and distinct variables to powers of
distinct variables; so it sends distinct monomials to distinct monomials. The variables sent to the column $c=1$ are $z_l$,
$z'_l$, $z_+$ and $z'_+$, of the letters $(y_{j,u_l},\epsilon)$ with $u_l>\eta$ and $(s_{j+1},\epsilon)$; every variable sent
to a letter $(y_{j,i},\epsilon)$ is sent to a $P_i$-th power. Call the polynomials named $f$ and $g$ in the generators of
Definition~\ref{pd-kappa-defects-segment-lattice}(5) $f_0$ and $g_0$. As $Z_j(b;\mu)$ is the $\mathcal R$-span of the
generators of all levels of $b$, a monomial of $f$ is a monomial of some generator.

(1) The factors $\hat E_l$ and $\hat E_+$ and the powers of $z_l$ and $z_+$ have integer coefficients, and $f_0$ and $g_0$
have coefficients in $\mathcal R$; so every generator, and $f$, has coefficients in $\mathcal R$. So has $\iota_j(f)$, and
$\iota_j(b,f)$ is a monomial in the variables of the letters of (1) times $\iota_j(f)$.

(2) As $\iota_j(b,f)$ is linear in $f$, it suffices to take for $f$ a generator $G$ of a level $(j_1,\dots,j_{s-\eta},j_+)$. The
conditions on a kept block give $w_i\leq\hat b_i$ and $\varsigma_i\leq\hat\alpha_i$ for $i\leq\eta$ and $o_i\leq\hat b_i$ for
$i>\eta$, and $o_+\leq\alpha_+$ holds by hypothesis; so the truncated differences in the degrees $\delta$, in the factors
$z_l^{\hat b_{u_l}\dot-o_{u_l}}$ and $z_+^{\alpha_+\dot-o_+}$ and in $\iota_j(b,\cdot)$ are differences. Each factor
$\hat E_l$ has degree one in the variables of $x_{u_l}$ and $z_l$ together and one in those of $y_{u_l}$ and $z'_l$ together,
and likewise $\hat E_+$ for $x_+$ and $z_+$, and for $\xi$ and $z'_+$. Counting the degrees of $\iota_j(b,G)$ letter by
letter:
\begin{itemize}
\item at $(y_{j,i},1)$ with $i\leq\eta$: $|\rho_i|+P_i(\hat b_i-w_i)$ from the monomial factor and $P_i\delta'(x_i)=P_iw_i$
  from $\iota_j(G)$, in all $|\rho_i|+P_i\hat b_i$; at $(y_{j,i},0)$: $|\rho'_i|+P_i\varsigma_i+P_i(\hat\alpha_i-\varsigma_i)
  =|\rho'_i|+P_i\hat\alpha_i$;
\item at $(y_{j,u_l},1)$: $|\rho_{u_l}|$ plus $P_{u_l}$ times $(\hat b_{u_l}-o_{u_l})+(o_{u_l}-j_l)+j_l$, from
  $z_l^{\hat b_{u_l}-o_{u_l}}$, from $f_0$ and from $\hat E_l^{j_l}$, in all $|\rho_{u_l}|+P_{u_l}\hat b_{u_l}$; at
  $(y_{j,u_l},0)$: $|\rho'_{u_l}|+P_{u_l}\big((\hat\alpha_{u_l}-j_l)+j_l\big)=|\rho'_{u_l}|+P_{u_l}\hat\alpha_{u_l}$;
\item at $(s_{j+1},1)$: $(\alpha_+-o_+)+(o_+-j_+)+j_+=\alpha_+$, from $z_+^{\alpha_+-o_+}$, from $f_0$ and from
  $\hat E_+^{j_+}$; at $(s_{j+1},0)$: $(b_+-j_+)+j_+=b_+$, from $g_0$ and from $\hat E_+^{j_+}$.
\end{itemize}
For a kept block $|\rho_i|\leq b_i$ and $P_i$ divides $b_i-|\rho_i|$, so $|\rho_i|+P_i\hat b_i=b_i$; likewise
$|\rho'_i|+P_i\hat\alpha_i=\alpha_i$. These are the degrees $d_x$: by Definition~\ref{pd-kappa-product-space}, for odd $y$
the letter $(y,1)$ has degree $b_T(y)$ and $(y,0)$ has degree $a_T(y)$, and for even $s$ the letter $(s,1)$ has degree
$a_T(s)$ and $(s,0)$ has degree $b_T(s)$.

(3) Every monomial $\nu$ of $f$ has degree $o_{u_l}$ in the variables of $x_{u_l}$ and $z'_l$ together, for every $l$, and
degree $o_+$ in those of $x_+$ and $z'_+$ together: in a generator of level $(j_1,\dots,j_+)$, $f_0$ contributes
$\delta'(x_{u_l})=o_{u_l}-j_l$ and each factor $\hat E_l$ contributes one, and the other factors do not involve these
variables; likewise for $x_+$. Use tildes for the entries of $b'$. Suppose that a monomial $M$ occurs in both, so that
$M=M_b\,\iota_j(\nu)=M_{b'}\,\iota_j(\nu')$ with $M_b$ and $M_{b'}$ the monomial factors of $\iota_j(b,\cdot)$ and
$\iota_j(b',\cdot)$ and $\nu$, $\nu'$ monomials of $f$ and $f'$.
\begin{itemize}
\item In $M$, the exponent of $z_{(y_{j,i},1),c}$ is $\rho_i(c)$ plus a multiple of $P_i$. As the entries of $\rho_i$ and
  $\tilde\rho_i$ are less than $P_i$, $\rho_i(c)=\tilde\rho_i(c)$. Likewise $\rho'_i=\tilde\rho'_i$, and so $\hat b_i$ and
  $\hat\alpha_i$ are the same for $b$ and $b'$.
\item For $i\leq\eta$, $\iota_j(\nu)$ does not involve $z_{(y_{j,i},\epsilon),1}$, so the exponents of
  $z_{(y_{j,i},1),1}$ and $z_{(y_{j,i},0),1}$ in $M$ are $\rho_i(1)+P_i(\hat b_i-w_i)$ and $\rho'_i(1)+P_i\varsigma_i$. Hence
  $w_i=\tilde w_i$ and $\varsigma_i=\tilde\varsigma_i$.
\item So $M_b=M_{b'}$, hence $\iota_j(\nu)=\iota_j(\nu')$ and $\nu=\nu'$. By the degrees above, $o_{u_l}=\tilde o_{u_l}$ for
  every $l$ and $o_+=\tilde o_+$. For a kept block $o_i=0$ for $i\leq\eta$ and $w_i=\varsigma_i=0$ for $i>\eta$.
\end{itemize}
So all entries of $b$ and $b'$ agree, and $b=b'$, a contradiction. The Fock pairing $\sum_\kappa\kappa!\,F_\kappa G_\kappa$ of
two polynomials without a common monomial is $0$.

(4) Every element of $T$ is exactly one of the $s_{j+1}$ and the $y_{j,i}$. So by (1) the factors $\iota_j(\beta_j,f_j)$ of
$\Phi(\beta,f)$ lie in $\mathcal R[z]$ and involve disjoint sets of variables for distinct $j$, and the same holds for
$\beta'$ and $f'$. The Fock pairing of products of polynomials in disjoint sets of variables is the product of the pairings
(Definition~\ref{pd-symbols-test-fock}). If every $\beta_j$ is kept with $o_+\leq a_T(s_{j+1})$, then by (2) the factor
$\iota_j(\beta_j,f_j)$ has degree $d_x$ in every letter $x$ of the $j$-th segment and degree $0$ in the others, so
$\Phi(\beta,f)$ has degree $d_x$ in every letter.
\end{proof}
% Lean: (1) is symIntAt (ssz_le_intAlg, aeval_mem_intAlg) with ioBlock_supported (symVar_mem); (2) is TD_wLet with
%   ioBlock_wh (weighted homogeneity for the weight wLet x of a letter), and Phi in Pol is symPolAt_of (natPartSum),
%   which takes o_+ <= a_T(s_{j+1}) from OPlusAt; (3) is orthAt through ioBlock_eq_sum and sep (ioExp_mod for the
%   digits, coarse_spec and baseExp for w and sigma, symLev_pos and coarse_eval for mu = mu', omB_wo and omB_wp with the
%   balanced weights of hom_of_ssz for o and o_+); (4) is prodAt (fock_prod_disj, segVars_disj). The Lean indexes the
%   letters i and the columns c from 0; its sigma is eta.

\begin{lemma}\label{pd-kappa-defects-windows-lift-pairs-source-order-segment}
Let $p\geq5$ be a prime, $r\geq1$ and $s\geq0$ integers, and $S$ a segment datum of length $s$, with the notation of
Definition~\ref{pd-kappa-defects-segment-lattice}, such that $G\geq2$, $d_1<\dots<d_s$ are odd with $0<d_i<G$, $\Lambda\geq0$
and $E_0\geq0$. Let $b'>0$, $\delta$, $\beta_e$, $a_S$ and $Q$ be real numbers such that:
\begin{enumerate}
\item for $1\leq i\leq s$, $i\leq\eta$ if and only if $p^{d_i}<a/b'$;
\item $b'(1-\beta_e)\leq b_i\leq b'(1+\beta_e)$ for every $i$, and $b_+\leq b'(1+\beta_e)$;
\item $0\leq\beta_e\leq\delta/2$ and $\delta<1/6$;
\item no odd power $p^e$ lies in $[(a/b')(1-2\delta),(a/b')(1+2\delta)]$;
\item $a\leq a_S$, $s\leq p$ and $p^G\leq Q$.
\end{enumerate}
Let $\alpha_1,\dots,\alpha_s,\alpha_+\geq0$ be integers. For a block $b$ put
\[
 \operatorname{ord}(b)=\sum_{i=1}^s\Big(P_io^\flat_i+\sum_{c=2}^{r+1}\rho_i(c)+\rho'_i(1)\Big)+o_+,
\]
and put $R=a/p^2+b'/p+\Lambda a/(p(r+1))+\Lambda b'/(r+1)+(r+1)Q$. Then:
\begin{enumerate}
\item for every block $b$,
  $O_+\leq\dfrac a{p^2}+\dfrac{2sb'}p+\Big(\dfrac\Lambda{r+1}+\dfrac4{p^2}\Big)\dfrac{2sa}p+\dfrac{2\Lambda b'}{r+1}$;
\item for every kept block $b$, $\operatorname{ord}(b)\leq\sum_{i=1}^s\min\big(b_i,a_S/p^{d_i}\big)+30(s+1)R$.
\end{enumerate}
\end{lemma}
\leanref{Formalism/SharpRestriction/GrowthBounds/BinomialRatio/TopDegreeLeTwoB.lean}{Formalism.MainNotions.ColumnSystems.Pf.SegHyp}
\leanref{Formalism/SharpRestriction/GrowthBounds/BinomialRatio/TopDegreeLeTwoB.lean}{Formalism.MainNotions.ColumnSystems.Pf.SegHyp.Pcap}
\leanref{Formalism/SharpRestriction/GrowthBounds/BinomialRatio/TopDegreeLeTwoB.lean}{Formalism.MainNotions.ColumnSystems.Pf.SegHyp.top_ord}
\leanref{Formalism/SharpRestriction/GrowthBounds/BinomialRatio/TopDegreeLeTwoB.lean}{Formalism.MainNotions.ColumnSystems.Pf.SegHyp.dig}
\leanref{Formalism/SharpRestriction/GrowthBounds/BinomialRatio/TopDegreeLeTwoB.lean}{Formalism.MainNotions.ColumnSystems.Pf.SegHyp.oplus}
\leanref{Formalism/SharpRestriction/GrowthBounds/BinomialRatio/TopDegreeLeTwoB.lean}{Formalism.MainNotions.ColumnSystems.Pf.SegHyp.seg_ord}
\leanref{Formalism/SharpRestriction/GrowthBounds/BinomialRatio/TopDegreeLeTwoB.lean}{Formalism.MainNotions.ColumnSystems.Pf.Rerr}

\begin{proof}
Let $b$ be a block, with the notation of Definition~\ref{pd-kappa-defects-segment-lattice}(1).

\emph{Step 1: four bounds.}
\begin{enumerate}
\item[(a)] $\hat a\leq a/p^G$, as $E_0\geq0$.
\item[(b)] $P_i\hat b_i\leq b_i\dot-|\rho_i|\leq b_i$ for every $i$, by the definition of $\hat b_i$.
\item[(c)] $b_i\leq2b'$ for every $i$, and $b_+\leq2b'$: by (3), $\beta_e\leq\delta/2<1$, and (2) gives the bounds.
\item[(d)] $p\sum_{i'<i}p^{d_{i'}}\leq p^{d_i}$ for $1\leq i\leq s$. By induction on $i$: for $i=1$ the sum is empty. As the
  $d_i$ are odd and increasing, $d_{i+1}\geq d_i+2$; so, using the bound for $i$,
  $p\big(\sum_{i'<i}p^{d_{i'}}+p^{d_i}\big)\leq p^{d_i}(1+p)\leq p^{d_i+2}\leq p^{d_{i+1}}$.
\end{enumerate}

\emph{Step 2: the caps.} We show $P_i\operatorname{cap}(i-1)\leq a/p^{d_i}+2b'/p$ for every $i$. As the $o^\flat_{i'}$ are
$\geq0$, $\operatorname{cap}(i-1)\leq\hat a+\sum_{i'<i}\hat b_{i'}$. By (a), $P_i\hat a\leq p^{G-d_i}a/p^G=a/p^{d_i}$. For
$i'<i$, $P_i\hat b_{i'}=p^{d_{i'}-d_i}P_{i'}\hat b_{i'}\leq2b'p^{d_{i'}}/p^{d_i}$ by (b) and (c). By (d) the sum of these over
$i'<i$ is at most $(2b'/p^{d_i})(p^{d_i}/p)=2b'/p$.

\emph{Step 3: one letter.} Let $b$ be kept, let $1\leq i\leq s$, and put
\[
 e=\frac{2b'}p+\frac{2\Lambda b'}{r+1}+\frac{8b'}{p^2}+\frac{\Lambda a}{p(r+1)}+\frac{2\Lambda b'}{p(r+1)} .
\]
We show $P_io^\flat_i\leq\min(b_i,a_S/p^{d_i})+e$. As $d_i\geq1$, $p\leq p^{d_i}$, so $a/p^{d_i}\leq a/p$; and
$\delta\geq0$ by (3).
\begin{itemize}
\item Let $i\leq\eta$. Then $p^{d_i}<a/b'$ by (1), and as $d_i$ is odd, (4) gives $p^{d_i}<(a/b')(1-2\delta)$. By (2),
  $b_ip^{d_i}\leq(1+\beta_e)b'p^{d_i}\leq(1+\beta_e)(1-2\delta)a\leq a\leq a_S$, as
  $(1+\beta_e)(1-2\delta)\leq1+\delta/2-2\delta\leq1$ by (3). So the minimum is $b_i$. As $b$ is kept,
  $o^\flat_i=w_i+\varsigma_i\leq(1+4/p^2)\hat b_i+\Lambda\operatorname{cap}(i-1)/(r+1)$. Multiplying by $P_i$ and using (b),
  (c) and Step~2,
  \[
   P_io^\flat_i\leq b_i\Big(1+\frac4{p^2}\Big)+\frac\Lambda{r+1}\Big(\frac ap+\frac{2b'}p\Big)
   \leq b_i+\frac{8b'}{p^2}+\frac{\Lambda a}{p(r+1)}+\frac{2\Lambda b'}{p(r+1)} .
  \]
\item Let $i>\eta$. Then $p^{d_i}\geq a/b'$ by (1), and (4) gives $p^{d_i}>(a/b')(1+2\delta)$. By (2),
  $b_ip^{d_i}\geq(1-\beta_e)b'p^{d_i}\geq(1-\beta_e)(1+2\delta)a\geq a$, as
  $(1-\beta_e)(1+2\delta)\geq1+2\delta-\delta/2-\delta^2\geq1$ by (3). With $a\leq a_S$, this gives
  $a/p^{d_i}\leq\min(b_i,a_S/p^{d_i})$. As $b$ is kept, $o^\flat_i=o_i\leq\operatorname{cap}(i-1)+\Lambda\hat b_i/(r+1)$.
  Multiplying by $P_i$ and using Step~2, (b) and (c),
  $P_io^\flat_i\leq a/p^{d_i}+2b'/p+2\Lambda b'/(r+1)$.
\end{itemize}
In both cases the claim follows, as the terms of $e$ are $\geq0$.

\emph{Step 4: the digits.} If $b$ is kept, every entry of $\rho_i$ and $\rho'_i$ is less than $P_i\leq p^G\leq Q$, so
$\sum_{c=2}^{r+1}\rho_i(c)+\rho'_i(1)\leq rQ+Q\leq2(r+1)Q$.

(1) As $d_i<G$, $P_i\geq p$; with (b) and (c), $\hat b_i\leq2b'/p$. As $o^\flat_i\geq0$, (a) and $G\geq2$ give
$\operatorname{cap}(s)\leq\hat a+\sum_i\hat b_i\leq a/p^2+2sb'/p$. Also $\sum_i2a/P_i\leq2sa/p$, and
$\Lambda b_+/(r+1)\leq2\Lambda b'/(r+1)$ by (c). Adding gives (1), as $\Lambda/(r+1)+4/p^2\geq0$.

(2) Let $b$ be kept. As $1/p\leq1/5$, $8b'/p^2\leq2b'/p$ and $2\Lambda b'/(p(r+1))\leq2\Lambda b'/(r+1)$, so
$e\leq4b'/p+\Lambda a/(p(r+1))+4\Lambda b'/(r+1)$. Likewise
$(\Lambda/(r+1)+4/p^2)\,2sa/p=s\big(2\Lambda a/(p(r+1))+8a/p^3\big)\leq s\big(2\Lambda a/(p(r+1))+2a/p^2\big)$. As $b$ is kept,
$o_+\leq O_+$. Summing Steps 3 and 4 over $i$ and adding (1),
\[
 \operatorname{ord}(b)\leq\sum_{i=1}^s\min\Big(b_i,\frac{a_S}{p^{d_i}}\Big)
 +s\Big(\frac{6b'}p+\frac{3\Lambda a}{p(r+1)}+\frac{4\Lambda b'}{r+1}+2(r+1)Q+\frac{2a}{p^2}\Big)+\frac a{p^2}
 +\frac{2\Lambda b'}{r+1} .
\]
The five terms of $R$ are $\geq0$; so the term with the factor $s$ is at most $6sR$, and $a/p^2+2\Lambda b'/(r+1)\leq2R$.
As $6sR+2R\leq30(s+1)R$, this gives (2).
\end{proof}
% Lean: SegHyp p S b' de be aS Q bundles the hypotheses (sig is (1), bT and bP are (2), be0, beLe and de6 are (3), band
%   is (4), aS, sp and Q are (5)); sp is not used here. Step 1 is base_exponent_le_a, coarseA_mul,
%   SegHyp.top_degree_le_two_b and mul_sum_pow_lt_le_pow_of_segHyp; Step 2 is SegHyp.Pcap (cap_le); Step 3 is
%   SegHyp.top_ord, with e = errTop (the case split on i < S.sigma, with h3 in each case for the product of the two
%   factors); Step 4 is SegHyp.dig;
%   (1) is SegHyp.oplus and (2) is SegHyp.seg_ord (hEt, i3, k3, linarith). The Lean's cap p S b i, with i from 0, is
%   the book's cap(i - 1); oEff is o^flat and the Lean's sigma is eta.

\begin{lemma}\label{pd-kappa-defects-windows-lift-pairs-source-order}
Let $C_c\geq0$, $C_0\geq1$, $C_{E1},C_{E2}\geq1$ and $C_H\geq1000(C_{E1}+C_{E2}+1)(C_c+1)$. Let $p\geq5$ be a prime and
$n'\geq2$ an even integer, and let $m$, $G'$, $N_1$, $c_{cl}$, $D_p$, $c_*$, $T$, $Y$, $Y_b$, $b'$, $b_l$, $b_u$ and an integer
$a$ satisfy hypotheses (1) to (3) of Theorem~\ref{pd-gain-theorem-one} with $C_0$ and $C_H$. Let $c\geq0$ be an integer with
$c\leq C_c(c_*+|T|k\log_pa+|T|c_{cl}+n'c_{cl})$, and let $a_T,b_T\colon\mathbf Z/n'\to\mathbf Z_{\geq0}$ satisfy
$a\leq a_T(s)\leq a+a/(kY)+2J_0$ and $b_l\leq b_T(s)\leq b_u+2J_0$ for $s\in T$, where $J_0=\Lambda b'/k$. Form the numbers of
Definition~\ref{pd-kappa-defects-numbers} with these constants and any $C_m,C_A,C_t,C_\kappa>0$. For a block $b$ of
$S_j(b_T)$, with $\alpha_i=a_T(y_{j,i})$ and $r=m$, put
\[
 \operatorname{ord}(b)=\sum_{i=1}^{S_j}\Big(P_io^\flat_i+\sum_{c=2}^k\rho_i(c)+\rho'_i(1)\Big)+o_+,
\]
in the notation of Definition~\ref{pd-kappa-defects-segment-lattice}(1), and put
$R_0=a/p^2+b'/p+\Lambda a/(pk)+\Lambda b'/k+kQ$. Then:
\begin{enumerate}
\item for every $j$ and every kept block $b$ of $S_j(b_T)$,
  $\operatorname{ord}(b)\leq\sum_{i=1}^{S_j}\min\big(b_T(y_{j,i}),a_T(s_j)/p^{d_{j,i}}\big)+30(S_j+1)R_0$;
\item for every $j$ and every kept block $b$ of $S_j(b_T)$, its entry $o_+$ is at most $a_T(s_{j+1})$;
\item if $C_t\geq30$, then $\sum_j\operatorname{ord}(\beta_j)\leq t_0$ for every tuple $\beta=(\beta_j)$ with $\beta_j$ a kept
  block of $S_j(b_T)$.
\end{enumerate}
\end{lemma}
\leanref{Formalism/SharpRestriction/GrowthBounds/ObjSeg.lean}{Formalism.MainNotions.G.blockOrd}
\leanref{Formalism/SharpRestriction/GrowthBounds/BinomialRatio/TopDegreeLeTwoB.lean}{Formalism.MainNotions.ColumnSystems.Pf.SegHyp.seg_ord}
\leanref{Formalism/SharpRestriction/GrowthBounds/BinomialRatio/TopDegreeLeTwoB.lean}{Formalism.MainNotions.ColumnSystems.Pf.segHyp_of}
\leanref{Formalism/SharpRestriction/GrowthBounds/CoarseMultiplicity/OPlusAtOf.lean}{Formalism.MainNotions.ColumnSystems.Pf.oPlusAt_of}
\leanref{Formalism/SharpRestriction/GrowthBounds/BinomialRatio/TopDegreeLeTwoB.lean}{Formalism.MainNotions.ColumnSystems.Pf.ord_all}

\begin{proof}
The facts (a) to (g) in the proof of Lemma~\ref{pd-kappa-defects-margin} hold. Let $\delta$ be the number of
Theorem~\ref{pd-gain-theorem-one}, and put $\beta_e=(b_u-b_l)/b'+2\Lambda/k$.

\emph{The windows for $b_T$.} Let $s\in T$. As $b_l\leq b'\leq b_u$ by hypothesis (2) of
Theorem~\ref{pd-gain-theorem-one}, $b'(1-\beta_e)\leq b'-(b_u-b_l)\leq b_l\leq b_T(s)$ and
$b_T(s)\leq b_u+2J_0\leq b'+(b_u-b_l)+2\Lambda b'/k=b'(1+\beta_e)$. Also $\beta_e\geq0$, and as $C_H\geq1000$ and
$N_1\geq1$, $\beta_e+18N_1\Lambda/k\leq(b_u-b_l)/b'+20N_1\Lambda/k\leq\delta/2$; so $0\leq\beta_e\leq\delta/2$.

\emph{The segments.} Let $1\leq j\leq q$. The segment datum $S_j(b_T)$, with $r=m$, satisfies the hypotheses of
Lemma~\ref{pd-kappa-defects-windows-lift-pairs-source-order-segment} with $b'$, $\delta$, $\beta_e$, $a_S=a_T(s_j)$ and $Q$.
Indeed $n'$ is even, so $g_j$ is even and positive, hence $g_j\geq2$, and the distances $d_{j,i}$ from the even $s_j$ to the
odd $y_{j,i}$ are odd, increasing and strictly between $0$ and $g_j$. By fact (c), $\Lambda\geq1$; and $E_0\geq0$ as
$c\geq0$. Condition (1) is the definition of $\eta_j$, as the $d_{j,i}$ increase. Condition (2) is the window for $b_T$ at
the $y_{j,i}$ and at $s_{j+1}$. Condition (3) was shown above, with $\delta<1/6$ from hypothesis (2) of the theorem, and
(4) is also part of that hypothesis. For (5): $a\leq a_T(s_j)$ by hypothesis; $S_j\leq N_1\leq k\leq k^2Y^2\leq p$ and
$p^{g_j}\leq p^{G'}=Q$ by hypothesis (1). The kept blocks of $S_j(b_T)$ are formed with $\alpha_i=a_T(y_{j,i})$, the number
$\operatorname{ord}$ of the statement is that of the lemma with $r=m$, so $r+1=k$, and $R=R_0$.

(1) This is Lemma~\ref{pd-kappa-defects-windows-lift-pairs-source-order-segment}(2) for $S_j(b_T)$, as $b_i=b_T(y_{j,i})$.

(2) Let $b$ be a kept block of $S_j(b_T)$, and write $s=S_j$. Then $o_+\leq O_+$, and
Lemma~\ref{pd-kappa-defects-windows-lift-pairs-source-order-segment}(1) bounds $O_+$ by a sum of four terms. By fact (a),
$k\geq5$ and $|T|\geq2$, and $k^2\leq k^2Y^2\leq p$, so $5k\leq p$; by hypothesis (1), $2s\leq2N_1\leq k$; by facts (c)
and (d), $\Lambda\geq1$ and $4\Lambda\leq2|T|\Lambda\leq k$; and $pb'<a$ by hypothesis (2). Hence:
\begin{itemize}
\item $a/p^2\leq a/25$;
\item $2sb'/p\leq a/25$, as $50sb'\leq25kb'\leq5pb'<5a\leq pa$;
\item $(\Lambda/k+4/p^2)\,2sa/p\leq2\Lambda sa/p+8sa/p^3\leq a/4+4a/125$, as $8\Lambda s\leq k^2\leq p$, $8s\leq4k\leq4p/5$ and
  $p^2\geq25$;
\item $2\Lambda b'/k\leq2\Lambda b'\leq a/10$, as $20\Lambda b'\leq5kb'\leq pb'<a$.
\end{itemize}
The sum of these bounds is less than $a$, and $a\leq a_T(s_{j+1})$ by hypothesis; so $o_+\leq a_T(s_{j+1})$.

(3) The first sum in $t_0$ runs over the odd $y\in T$, and these are the $y_{j,i}$, with $y_{j,i}-d(y_{j,i})=s_j$ and
$d(y_{j,i})=d_{j,i}$. So summing (1) over $j$ gives
\[
 \sum_j\operatorname{ord}(\beta_j)\leq\sum_{y\in T\ \mathrm{odd}}\min\Big(b_T(y),\frac{a_T(y-d(y))}{p^{d(y)}}\Big)
 +30R_0\sum_j(S_j+1).
\]
Every element of $T$ is exactly one of the $s_{j+1}$ and the $y_{j,i}$, so $\sum_j(S_j+1)=|T|$. As $C_t\geq30$, $R_0\geq0$ and
$Qc\geq0$, $30R_0\leq C_t(R_0+Qc)$, and $t_0$ is the first sum plus $C_t|T|(R_0+Qc)$.
\end{proof}
% Lean: the facts are Std C0 Cde CE1 CE2 1 X from std_of under the thresholds of thr_std (those of the statement):
%   the windows for b_T are hwinB (eB1, eB2) with be = beta_e, and 0 <= be <= de/2 is hbe0 and hbeLe (hCv). The
%   segments are segHyp_of St j (segG_even, dist_pos_lt, dist_odd, dist_strictMono, sigma_iff, St.winB, St.band,
%   St.winA, St.segB, St.pk). (1) is (segHyp_of St j).seg_ord. (2) is oPlusAt_of: (segHyp_of St j).oplus, then e1 to e4
%   from St.pab, St.pk, St.N1k, St.N2, St.LaN and St.segB, and St.winA at segNext. (3) is ord_all St hCt bs hK: hseg is
%   (segHyp_of St j).seg_ord (hK j), summed with sum_segS for the count of the letters, and a3 for 30 R <= Ct (R + Q c).

\begin{definition}\label{pd-kappa-defects-windows-lift-pairs-source-lift-relaxed}
Let the data and constants of Definition~\ref{pd-kappa-defects-numbers} be given, together with reals $b_l\leq b_u$ and
$C_H>0$; let $\delta=C_HN_1\Lambda/k+2(b_u-b_l)/b'$ be the number of Theorem~\ref{pd-gain-theorem-one}, and put
$\kappa_1=13kN_1Q$.
\begin{enumerate}
\item Let $E\geq0$ be a real number and $1\leq j\leq q$. A block $b$ of $S_j(b_T)$, for $\alpha_i=a_T(y_{j,i})$ and $r=m$, with
  the notation of Definition~\ref{pd-kappa-defects-segment-lattice}(1), is \emph{relaxed of slack} $E$ if:
  \begin{enumerate}
  \item for every $i$, the entries of $\rho_i$ and of $\rho'_i$ are less than $P_i$, $|\rho_i|\leq b_i$, $|\rho'_i|\leq\alpha_i$,
    and $P_i$ divides $b_i-|\rho_i|$ and $\alpha_i-|\rho'_i|$;
  \item for $i>\eta_j$, $o_i\leq\hat b_i$, $o_i\leq\hat b_i(1-\delta/4)+E$ and $w_i=\varsigma_i=0$;
  \item for $i\leq\eta_j$, $o_i=0$, $w_i\leq\hat b_i$, $\varsigma_i\leq\hat\alpha_i$ and $\varsigma_i\leq2\Lambda\hat a/k+E$;
  \item $o_+\leq a_T(s_{j+1})$ and $o_+\leq a/p^2+3b'/p+(\Lambda/k+4/p^2)\,4a/p+2\Lambda b'/k+E$.
  \end{enumerate}
\item Let an unramified quadratic Hida setting of degree $n'$ at $p$ with rings $\mathcal R\subset\mathcal K$ and
  $\varpi_0$ be given, with the notation of Definition~\ref{pd-kappa-defects-windows-lift-pairs-model}. For an integer
  $o\geq0$, a polynomial $F\in\mathrm{Pol}_T$ is \emph{clean of order} $o$ if $\varpi_0^oF\in\mathcal R[z]$ and
  $\Sigma_gF\in\varpi_0^c\mathcal R[z]$ for every $g\in GL_k(W)$ whose plane differs from that of $1$
  (Lemma~\ref{pd-symbols-edge-basis}, with $r=1$).
\item In the setting of (2), for real numbers $\mu$ and $E\geq0$ and an integer $o\geq0$, $S(\mu,E,o)\subset\mathcal K[z]$ is the
  $\mathcal R$-span of the symbol families $\Phi(\beta,f)$ with $\beta=(\beta_j)$, $\beta_j$ a block of $S_j(b_T)$ relaxed of slack
  $E$, $\sum_j\operatorname{ord}(\beta_j)\leq o$ (Lemma~\ref{pd-kappa-defects-windows-lift-pairs-source-order}), and $f=(f_j)$,
  $f_j$ a generator of some level of $\beta_j$ (Definition~\ref{pd-kappa-defects-segment-lattice}(5)), formed over
  $\mathcal R$ with $\alpha_i=a_T(y_{j,i})$, $\alpha_+=a_T(s_{j+1})$, $\mu$ and $r=m$.
\end{enumerate}
\end{definition}
\leanref{Formalism/SharpRestriction/GrowthBounds/CoarseMultiplicity/CoarseMultiplicityAtIdentity.lean}{Formalism.MainNotions.ColumnSystems.Pf.ka1}
\leanref{Formalism/SharpRestriction/GrowthBounds/CoarseMultiplicity/CoarseMultiplicityAtIdentity.lean}{Formalism.MainNotions.ColumnSystems.Pf.localOrderCap}
\leanref{Formalism/SharpRestriction/GrowthBounds/CoarseMultiplicity/CoarseMultiplicityAtIdentity.lean}{Formalism.MainNotions.ColumnSystems.Pf.BLk}
\leanref{Formalism/SharpRestriction/GrowthBounds/CoarseMultiplicity/CoarseMultiplicityAtIdentity.lean}{Formalism.MainNotions.ColumnSystems.Pf.IsCleanAwayFromPlane}
\leanref{Formalism/SharpRestriction/GrowthBounds/CoarseMultiplicity/CoarseMultiplicityAtIdentity.lean}{Formalism.MainNotions.ColumnSystems.Pf.genSet}
\leanref{Formalism/SharpRestriction/GrowthBounds/CoarseMultiplicity/CoarseMultiplicityAtIdentity.lean}{Formalism.MainNotions.ColumnSystems.Pf.ordFam}
\leanref{Formalism/SharpRestriction/GrowthBounds/CoarseMultiplicity/CoarseMultiplicityAtIdentity.lean}{Formalism.MainNotions.ColumnSystems.Pf.SSg}
% Lean: ka1 is kappa_1 and localOrderCap the bound of (d) without E; BLk C0 Cde CE1 CE2 X E j b is (1) (DigitsOK is (a),
%   with de = X.de C0 Cde, La = X.La C0 and A0 = a-hat); IsCleanAwayFromPlane X O w 1 o F is (2) (InA, InWjRg at the g'
%   with IsGLW g' and ratPlane g' != ratPlane 1); SSg O C0 Cde CE1 CE2 X mu E o is (3), with genSet the union of the gens
%   over the levels LevelOK and ordFam the sum of blockOrd. The Lean defines (2) for any g in place of 1.

\begin{lemma}\label{pd-kappa-defects-windows-lift-pairs-source-lift-kept}
Let $C_c\geq0$, $C_0\geq1$, $C_{E1},C_{E2}\geq1$ and $C_H\geq1000(C_{E1}+C_{E2}+1)(C_c+1)$. Let $p$, $n'$, $m$, $G'$, $N_1$,
$c_{cl}$, $D_p$, $c_*$, $T$, $Y$, $Y_b$, $b'$, $b_l$, $b_u$, $a$, $c\geq0$, $a_T$ and $b_T$ satisfy the hypotheses of
Lemma~\ref{pd-kappa-defects-windows-lift-pairs-source-order} with these constants, and form the numbers of
Definition~\ref{pd-kappa-defects-numbers} with these constants and any $C_m,C_A,C_t,C_\kappa>0$. Then for every $j$, every
kept block of $S_j(b_T)$, for $\alpha_i=a_T(y_{j,i})$ and $r=m$, is relaxed of slack $0$
(Definition~\ref{pd-kappa-defects-windows-lift-pairs-source-lift-relaxed}(1)).
\end{lemma}
\leanref{Formalism/SharpRestriction/GrowthBounds/CoarseMultiplicity/CoarseMultiplicityAtIdentity.lean}{Formalism.MainNotions.ColumnSystems.Pf.KeptBLAt}
\leanref{Formalism/SharpRestriction/GrowthBounds/CoarseMultiplicity/PolyUpperBoundDelta.lean}{Formalism.MainNotions.ColumnSystems.Pf.restKeptBL}
\leanref{Formalism/SharpRestriction/GrowthBounds/CoarseMultiplicity/PolyUpperBoundDelta.lean}{Formalism.MainNotions.ColumnSystems.Pf.keptBL_of_std}
\leanref{Formalism/SharpRestriction/GrowthBounds/CoarseMultiplicity/PolyUpperBoundDelta.lean}{Formalism.MainNotions.ColumnSystems.Pf.kept_upper_level_bound}
\leanref{Formalism/SharpRestriction/GrowthBounds/CoarseMultiplicity/PolyUpperBoundDelta.lean}{Formalism.MainNotions.ColumnSystems.Pf.upper_real_bound}
\leanref{Formalism/SharpRestriction/GrowthBounds/CoarseMultiplicity/PolyUpperBoundDelta.lean}{Formalism.MainNotions.ColumnSystems.Pf.extra_order_le_maximal}
\leanref{Formalism/SharpRestriction/GrowthBounds/CoarseMultiplicity/PolyUpperBoundDelta.lean}{Formalism.MainNotions.ColumnSystems.Pf.maximal_order_le_a}

\begin{proof}
The facts (a) to (g) in the proof of Lemma~\ref{pd-kappa-defects-margin} hold. Put $\beta_e=(b_u-b_l)/b'+2\Lambda/k$ and
$\varepsilon_0=N_1\Lambda/k$. As in the proof of Lemma~\ref{pd-kappa-defects-windows-lift-pairs-source-order},
$b'(1-\beta_e)\leq b_T(s)\leq b'(1+\beta_e)$ for $s\in T$, and $S_j(b_T)$ satisfies the hypotheses of
Lemma~\ref{pd-kappa-defects-windows-lift-pairs-source-order-segment} with $b'$, $\delta$, $\beta_e$, $a_S=a_T(s_j)$ and $Q$.
Moreover:
\begin{enumerate}
\item[(i)] $0\leq\delta<1/6$, $40\varepsilon_0\leq\delta$, $\beta_e\leq\delta/2-18\varepsilon_0$ and $\Lambda/k\leq\varepsilon_0$. Indeed
  $\delta<1/6$ by hypothesis (2) of Theorem~\ref{pd-gain-theorem-one}; as $b_l\leq b_u$, $C_H\geq1000$ and $N_1\geq1$,
  $\delta\geq C_HN_1\Lambda/k\geq40\varepsilon_0\geq0$ and $\beta_e+18\varepsilon_0\leq(b_u-b_l)/b'+20N_1\Lambda/k\leq\delta/2$.
\item[(ii)] $2b'/p\leq\varepsilon_0b'$ and $kQ\leq\varepsilon_0b'$. Indeed, by facts (a) and (c), $2k\leq k^2\leq p$ and
  $N_1\Lambda\geq1$, so $2/p\leq1/k\leq\varepsilon_0$; and by fact (g), $kQ\leq k^3QY/k^2\leq b'/k\leq\varepsilon_0b'$.
\end{enumerate}
Let $b$ be a kept block of $S_j(b_T)$, with the notation of Definition~\ref{pd-kappa-defects-segment-lattice}(1), and write
$s=S_j$, $\eta=\eta_j$ and $P_i=p^{g_j-d_{j,i}}$. Conditions (a) and (c) of a block relaxed of slack $0$ are conditions of a
kept block, and so are $o_i\leq\hat b_i$ and $w_i=\varsigma_i=0$ for $i>\eta$.

\emph{The orders $o_i$ with $i>\eta$.} Let $i>\eta$, and put $P=P_i$, $A=\hat b_i$ and $t=a/p^{d_{j,i}}$. As $b$ is kept,
$o_i\leq\operatorname{cap}(i-1)+\Lambda A/k$; we show that the right side is at most $A(1-\delta/4)$.
\begin{itemize}
\item By Step~2 of the proof of Lemma~\ref{pd-kappa-defects-windows-lift-pairs-source-order-segment} and (ii),
  $P\operatorname{cap}(i-1)\leq t+2b'/p\leq t+\varepsilon_0b'$.
\item As in Step~3 of that proof, $p^{d_{j,i}}>(a/b')(1+2\delta)$, so $t(1+2\delta)<b'$.
\item The $k$ entries of $\rho_i$ are less than $P$, and $P$ divides $b_i-|\rho_i|$, so $AP=b_i-|\rho_i|>b_i-kP$. As
  $P\leq p^{g_j}\leq Q$, (ii) gives $kP\leq\varepsilon_0b'$; with $b_i\geq b'(1-\beta_e)$ and (i),
  $AP>b'(1-\beta_e)-\varepsilon_0b'\geq b'(1-\delta/2+17\varepsilon_0)$.
\item As $0\leq\delta<1/6$ and $0\leq40\varepsilon_0\leq\delta$,
  $(1-\delta/2)(1-\delta/4)\leq(1-\delta/2+17\varepsilon_0)(1-\delta/4-\varepsilon_0)-\varepsilon_0$ and
  $1\leq(1+2\delta)(1-\delta/2)(1-\delta/4)$: the difference of the two sides of the first is
  $\varepsilon_0(15-15\delta/4-17\varepsilon_0)\geq0$, and the right side of the second is
  $1+\delta(5/4-11\delta/8+\delta^2/4)\geq1$.
\end{itemize}
So $t<b'/(1+2\delta)\leq b'\big((1-\delta/2+17\varepsilon_0)(1-\delta/4-\varepsilon_0)-\varepsilon_0\big)$, and
\[
 P\operatorname{cap}(i-1)\leq t+\varepsilon_0b'<b'(1-\delta/2+17\varepsilon_0)(1-\delta/4-\varepsilon_0)\leq AP(1-\delta/4-\varepsilon_0),
\]
as $1-\delta/4-\varepsilon_0>0$. Also $P\Lambda A/k\leq\varepsilon_0AP$ by (i). Adding,
$P(\operatorname{cap}(i-1)+\Lambda A/k)\leq AP(1-\delta/4)$, and dividing by $P>0$ gives the claim.

\emph{The order $o_+$.} As $b$ is kept, $o_+\leq O_+$. By Step~1 of the proof of
Lemma~\ref{pd-kappa-defects-windows-lift-pairs-source-order-segment} and $g_j\geq2$, $\hat a\leq a/p^2$ and
$\hat b_i\leq b_i/P_i\leq2b'/P_i$; and $\operatorname{cap}(s)\leq\hat a+\sum_i\hat b_i$, as the $o^\flat_i$ are $\geq0$. The $P_i$ are
distinct powers $p^e$ with $1\leq e<g_j$, so $\sum_i1/P_i\leq\sum_{e\geq1}p^{-e}=1/(p-1)\leq5/(4p)$. Hence
$\sum_i\hat b_i\leq3b'/p$ and $\sum_i2a/P_i\leq4a/p$, and $\Lambda b_+/k\leq2\Lambda b'/k$, as $b_+\leq b'(1+\beta_e)\leq2b'$. So
$o_+\leq L$ with $L=a/p^2+3b'/p+(\Lambda/k+4/p^2)\,4a/p+2\Lambda b'/k$. Finally, by fact (d) and $|T|\geq2$,
$4\Lambda\leq2|T|\Lambda\leq k$, so $\Lambda/k\leq1/4$; by hypothesis (2), $b'<a/p\leq a/5$; and $p^2\geq25$. So
$L\leq a/25+3a/25+(1/4+4/25)\cdot4a/5+a/10<a\leq a_T(s_{j+1})$.
\end{proof}
% Lean: restKeptBL = TransportedAlt.mono with thr_std (thresholds 1, 1 and 1000(CE1 + CE2 + 1)(Cc + 1)) of keptBL_of_std:
%   the digits and the conditions for i <= eta are those of Kept; for i > eta, kept_upper_level_bound (segHyp_of,
%   coarse_bounds, SegHyp.Pcap, band and sig, Std.twoP, Std.kQeps, Std.be18, Std.forty_epsilon_le_delta, Std.Laeps)
%   with upper_real_bound (poly_upper_bound_delta and poly_lower_bound_delta); o_+ is extra_order_le_maximal
%   (cap_le, base_exponent_le_a, coarseA_mul, top_degree_le_two_b, SegHyp.geo) and maximal_order_le_a with St.winA. The
%   Lean's eps is epsilon_0 and be is beta_e.

\subsection{Lift modules and factorization systems}\label{subsection-pd-factorization-systems}

\begin{definition}\label{pd-kappa-defects-windows-lift-pairs-source-lift-tail-systems}
Let the data and constants of Definition~\ref{pd-kappa-defects-numbers} be given, together with reals $b_l\leq b_u$ and
$C_H>0$, and let an unramified quadratic Hida setting of degree $n'$ at $p$ with rings $\mathcal R\subset\mathcal K$ and
$\varpi_0$ be given, with the notation of Definitions~\ref{pd-kappa-defects-windows-lift-pairs-model}
and~\ref{pd-kappa-defects-windows-lift-pairs-source-lift-relaxed}.
\begin{enumerate}
\item The \emph{$\varpi_0$-degree} $d_\varpi(M)$ of a monomial $z^M$ of $\mathcal K[z]$ is the sum of its exponents at the
  variables $z_{x,b}$ with $x$ a vector letter and $b\geq2$, or $x$ a covector letter and $b=1$.
\item Let $\mu$ be a real number and $0\leq n\leq c$ and $o\geq0$ integers, and put
  $S_\epsilon=S(\mu+(n-\epsilon)\kappa_1,(c-n+\epsilon)\kappa_1,o+\epsilon)$ for $1\leq\epsilon<n$. The \emph{lift module}
  $\mathcal L(\mu,n,o)$ is the set of the $\sigma\in\mathcal K[z]$ for which there are $F\in\mathrm{Pol}_T$, clean of order
  $o$, and $G_\epsilon\in S_\epsilon$ for $1\leq\epsilon<n$ with
  $\Sigma_1F-\sigma-\sum_{\epsilon=1}^{n-1}\varpi_0^\epsilon G_\epsilon\in\varpi_0^n\mathcal R[z]$.
\item With $\mu$, $n$, $o$ and the $S_\epsilon$ of (2), a \emph{factorization system} for $\sigma\in\mathcal K[z]$ consists of a
  finite set $A$; for each $a\in A$ an integer $m_a\geq0$, polynomials $f_a,\sigma_a\in\mathcal K[z]$, an integer $o_a\geq0$
  and a map $d_a$ from the letters of Definition~\ref{pd-kappa-defects-windows-lift-pairs-model}(1) to $\mathbf Z_{\geq0}$;
  and polynomials $f_R,\sigma_R\in\mathcal K[z]$ and an integer $o_R\geq0$; such that:
  \begin{enumerate}
  \item for every $a\in A$: $f_a$ and $\sigma_a$ lie in $\mathcal R[z]$; $\Sigma_gf_a\in\varpi_0^{o_a}\mathcal R[z]$ for every
    $g\in GL_k(W)$; $\Sigma_1f_a-\varpi_0^{o_a}\sigma_a\in\varpi_0^{o_a+1}\mathcal R[z]$; every monomial of $\sigma_a$ has
    $\varpi_0$-degree $o_a$; and $\Sigma_1f_a$ and $\sigma_a$ have row degrees $d_a$;
  \item $f_R$ and $\sigma_R$ lie in $\mathcal R[z]$, and $\Sigma_1f_R=\varpi_0^{o_R}\sigma_R$;
  \item the product $F_A=f_R\prod_{a\in A}f_a^{m_a}$ lies in $\mathrm{Pol}_T$, and $\Sigma_gF_A\in\varpi_0^{o+c}\mathcal R[z]$
    for every $g\in GL_k(W)$ whose plane differs from that of $1$ (Lemma~\ref{pd-symbols-edge-basis}, with $r=1$);
  \item $o_R+\sum_am_ao_a=o$ and $\sigma_R\prod_a\sigma_a^{m_a}=\sigma$;
  \item for every $u\colon A\to\mathbf Z_{\geq0}$ with $u_a\leq m_a$ for all $a$, every integer $e$ with $1\leq e<n$ and
    $\sum_au_a\leq e$, and every monomial $z^M$ of row degrees $\sum_au_ad_a$ with $d_\varpi(M)\leq\sum_au_ao_a+e$,
    \[
     \sigma_R\prod_a\sigma_a^{m_a-u_a}\cdot z^M\in S_e .
    \]
  \end{enumerate}
  The bounds $\Sigma_gf_a\in\varpi_0^{o_a}\mathcal R[z]$ of (a), for all $g\in GL_k(W)$, are not used: the proof of
  Lemma~\ref{pd-kappa-defects-windows-lift-pairs-source-lift-tail-factorization}(2) uses only the other conditions. At $g=1$
  the bound also follows from the next condition of (a), as $\sigma_a\in\mathcal R[z]$.
\item Let $1\leq j\leq q$, $b$ a block of $S_j(b_T)$ and $\mu'$ a real number, and form the generators of $b$ over
  $\mathcal R$ with $\alpha_i=a_T(y_{j,i})$, $\alpha_+=a_T(s_{j+1})$, $\mu'$ and $r=m$
  (Definition~\ref{pd-kappa-defects-segment-lattice}(5)), with the notation there and $s=S_j$, $\eta=\eta_j$. Number the
  letters $x_1,y_1,\dots,x_s,y_s,x_+$ of $\mathcal P_j$ as $0,1,\dots,2s$; the $x_i$ and $x_+$ are vector letters and the
  $y_i$ covector letters. A \emph{column generator} of level $(j_1,\dots,j_{s-\eta},j_+)$ of $b$ is the product of that
  definition with $f=\Psi(\nu,I)$ and $g=\prod_c\xi_c^{\gamma_c}$, where $\nu$ is a path of rank $m$ for the degrees
  $\delta'$ of the level with $\nu_{2u_l-1}(1)\leq M_l$ for every $l$ and $\nu_{2s+1}(1)\leq M_+$, $I$ is a standard filling of
  $\nu$, $\Psi(\nu,I)$ is formed with the variable $z_{i,c}$ of Definition~\ref{pd-symbols-path-generators} read as the
  $c$-th variable of the letter numbered $i$, and $\gamma\in\mathbf Z_{\geq0}^m$ has $\sum_c\gamma_c=b_+-j_+$.
\end{enumerate}
\end{definition}
\leanref{Formalism/SharpRestriction/GrowthBounds/StandardPolynomials/MultidegAdd.lean}{Formalism.MainNotions.ColumnSystems.Pf.ndM}
\leanref{Formalism/SharpRestriction/GrowthBounds/StandardPolynomials/MultidegAdd.lean}{Formalism.MainNotions.ColumnSystems.Pf.tailSub}
\leanref{Formalism/SharpRestriction/GrowthBounds/StandardPolynomials/MultidegAdd.lean}{Formalism.MainNotions.ColumnSystems.Pf.liftSub}
\leanref{Formalism/SharpRestriction/GrowthBounds/StandardPolynomials/MultidegAdd.lean}{Formalism.MainNotions.ColumnSystems.Pf.FacSys}
\leanref{Formalism/SharpRestriction/GrowthBounds/StandardPolynomials/ColGen.lean}{Formalism.MainNotions.ColumnSystems.Pf.colGens}
% Lean: (1) is ndM X M = degUS (hhLetters TT) 1 M. (2) is liftSub C0 Cde CE1 CE2 X O w mm n o: raw with
%   IsCleanAwayFromPlane X O w 1 o, and lift1 raw - sigma in tailSub, whose tau epsilon lie in SSe epsilon = SSg at
%   mm + (n - epsilon) ka1, (c - (n - epsilon)) ka1 and o + epsilon. (3) is the structure FacSys: kappa, mult, f, sg, od,
%   dg, fR, sgR, oR are A, m_a, f_a, sigma_a, o_a, d_a, f_R, sigma_R, o_R; (a) is fac (IsFac: inA, ord, sym, sgInt),
%   sgHom, liftPol and sgPol; (b) is rigA, rigLift and rigInt; (c) is total and clean; (d) is ordEq and symEq; (e) is
%   core. The clause ord of IsFac is for every g with IsGLW g, as stated in (a); FacSys.liftable uses none of it, and
%   IsFac.mul uses it only at g = 1. (4) is colGens O (blockData ...): BlockData.gens with y = psiR (KK O) (lettersAt j jp) x for x in sAll, the
%   pairs (Tx) of a path and a standard filling with mu (J k) 0 <= thr k for every k (sK), and g = prod_c X (xiVar c)^gamma c.
%   The Lean's J k = 2 uTop k + 1 numbers the letters from 0; it is the book's 2 u_l - 1 with u_l from 1.

\begin{lemma}\label{pd-kappa-defects-windows-lift-pairs-source-lift-tail-divisible}
Let $p\geq5$ be a prime, $(W,Q,\sigma,\iota)$ a Hida setting of degree $n'$ at $p$ with rings $\mathcal R\subset\mathcal K$ and
$\varpi_0$, $m\geq1$, $k=m+1$, $T\subset\mathbf Z/n'\mathbf Z$ and $a_T,b_T\colon\mathbf Z/n'\mathbf Z\to\mathbf Z_{\geq0}$, and use
the notation of Definition~\ref{pd-kappa-defects-windows-lift-pairs-model}. For a monomial $z^M$ let $d_\varpi(M)$ be the sum of
its exponents at the variables $z_{x,b}$ with $x$ a vector letter and $b\geq2$, or $x$ a covector letter and $b=1$
(Definition~\ref{pd-kappa-defects-windows-lift-pairs-source-lift-tail-systems}(1)). Let $F\in\mathcal R[z]$ and
$g\in GL_k(W)$. Then for every monomial $z^M$ the coefficient of $z^M$ in $\Sigma_gF$ lies in
$\varpi_0^{d_\varpi(M)}\mathcal R$.
\end{lemma}
\leanref{Formalism/SharpRestriction/GrowthBounds/FrobeniusPaths/LiftSubstitutionCoeffDivisible.lean}{Formalism.MainNotions.LemmaBAlt.usSub}
\leanref{Formalism/SharpRestriction/GrowthBounds/FrobeniusPaths/LiftSubstitutionCoeffDivisible.lean}{Formalism.MainNotions.LemmaBAlt.liftSubst_X_mem}
\leanref{Formalism/SharpRestriction/GrowthBounds/FrobeniusPaths/LiftSubstitutionCoeffDivisible.lean}{Formalism.MainNotions.LemmaBAlt.lift_substitution_coeff_divisible}

\begin{proof}
Let $U\subset\mathcal K[z]$ be the set of the polynomials whose coefficient at every monomial $z^M$ lies in
$\varpi_0^{d_\varpi(M)}\mathcal R$. As $d_\varpi(M+M')=d_\varpi(M)+d_\varpi(M')$ and $d_\varpi(0)=0$, $U$ is a subring of
$\mathcal K[z]$ containing $\mathcal R$, and an $\mathcal R$-module. Let $x$ be a letter, of position $\pi$, and $1\leq b\leq k$. By
Definition~\ref{pd-symbols-path-generators}(7) with $r=1$, if $x$ is a vector letter, $\Sigma_g$ sends $z_{x,b}$ to
$\sum_a\delta_a(g_\pi^{-1})_{ab}z_{x,a}$ with $\delta_1=1$ and $\delta_a=\varpi_0$ for $a\geq2$; if $x$ is a covector letter, it
sends $z_{x,b}$ to $\sum_a\delta'_a(g_\pi)_{ba}z_{x,a}$ with $\delta'_1=\varpi_0$ and $\delta'_a=1$ for $a\geq2$. As
$g\in GL_k(W)$, the entries of $g_\pi=\iota\sigma^\pi(g)$ and of $g_\pi^{-1}=\iota\sigma^\pi(g^{-1})$ lie in $\mathcal R$. In
both cases the factor $\delta_a$, respectively $\delta'_a$, is $\varpi_0^{d_\varpi(z_{x,a})}$, so every term lies in $U$, and
$\Sigma_g(z_{x,b})\in U$. As $\Sigma_g$ is a $\mathcal K$-algebra homomorphism and $F$ is an $\mathcal R$-linear combination of
monomials, $\Sigma_gF\in U$.
\end{proof}
% Lean: lift_substitution_coeff_divisible at r = 1 and LL = hhLetters TT: usSub = goodSub w (degUS LL r) is U (a
%   subalgebra by degUS_zero and degUS_add), liftSubst_X_mem gives the images of the variables (wX_mem_goodSub,
%   X_mem_goodSub, sigMat_inv for g^-1, C_mem_sub), and aeval_mem_sub concludes for F with InA F. The Lean states it for
%   any r, any lifted lettering and any k.

\begin{lemma}\label{pd-kappa-defects-windows-lift-pairs-source-lift-tail-factorization}
Let the data, the setting, $\mu$, $0\leq n\leq c$ and $o\geq0$ be as in
Definition~\ref{pd-kappa-defects-windows-lift-pairs-source-lift-tail-systems}(2).
\begin{enumerate}
\item The lift module $\mathcal L(\mu,n,o)$ is an $\mathcal R$-submodule of $\mathcal K[z]$.
\item If $\sigma\in\mathcal K[z]$ has a factorization system with $\mu$, $n$ and $o$
  (Definition~\ref{pd-kappa-defects-windows-lift-pairs-source-lift-tail-systems}(3)), then $\sigma\in\mathcal L(\mu,n,o)$.
\end{enumerate}
\end{lemma}
\leanref{Formalism/SharpRestriction/GrowthBounds/StandardPolynomials/MultidegAdd.lean}{Formalism.MainNotions.ColumnSystems.Pf.TCls}
\leanref{Formalism/SharpRestriction/GrowthBounds/StandardPolynomials/MultidegAdd.lean}{Formalism.MainNotions.ColumnSystems.Pf.TCls.mul}
\leanref{Formalism/SharpRestriction/GrowthBounds/StandardPolynomials/MultidegAdd.lean}{Formalism.MainNotions.ColumnSystems.Pf.TCls.prod}
\leanref{Formalism/SharpRestriction/GrowthBounds/StandardPolynomials/MultidegAdd.lean}{Formalism.MainNotions.ColumnSystems.Pf.tail_cls}
\leanref{Formalism/SharpRestriction/GrowthBounds/StandardPolynomials/MultidegAdd.lean}{Formalism.MainNotions.ColumnSystems.Pf.tail_of_ssg}
\leanref{Formalism/SharpRestriction/GrowthBounds/StandardPolynomials/MultidegAdd.lean}{Formalism.MainNotions.ColumnSystems.Pf.tail_of_wj}
\leanref{Formalism/SharpRestriction/GrowthBounds/StandardPolynomials/MultidegAdd.lean}{Formalism.MainNotions.ColumnSystems.Pf.FacSys.liftable}

\begin{proof}
Write $\varpi=\varpi_0$ and $S_\epsilon$ as in Definition~\ref{pd-kappa-defects-windows-lift-pairs-source-lift-tail-systems}(2),
and let $\mathcal T$ be the set of the $P\in\mathcal K[z]$ for which there are $G_\epsilon\in S_\epsilon$, $1\leq\epsilon<n$, with
$P-\sum_\epsilon\varpi^\epsilon G_\epsilon\in\varpi^n\mathcal R[z]$. Then $\sigma\in\mathcal L(\mu,n,o)$ if and only if
$\Sigma_1F-\sigma\in\mathcal T$ for some $F\in\mathrm{Pol}_T$ clean of order $o$.

(1) The $S_\epsilon$ and $\varpi^n\mathcal R[z]$ are $\mathcal R$-modules, so $\mathcal T$ is one: sums and $\mathcal R$-multiples of
the $G_\epsilon$ serve for sums and multiples, and $G_\epsilon=0$ for $0$. The polynomials clean of order $o$ form an
$\mathcal R$-module (Step~2 of the proof of Lemma~\ref{pd-kappa-defects-windows-lift-pairs-source-lift}), and $\Sigma_1$ is
$\mathcal K$-linear. So if $F$ serves for $\sigma$ and $F'$ for $\sigma'$, then $F+F'$ serves for $\sigma+\sigma'$, $aF$ for
$a\sigma$ with $a\in\mathcal R$, and $0$ for $0$.

(2) Let a factorization system for $\sigma$ be given, with the notation of
Definition~\ref{pd-kappa-defects-windows-lift-pairs-source-lift-tail-systems}(3). For $u\colon A\to\mathbf Z_{\geq0}$ put
$s_u=\sum_au_a$, $d_u=\sum_au_ad_a$ and $o_u=\sum_au_ao_a$. For a map $d$ from the letters to $\mathbf Z_{\geq0}$ and integers
$o',c'\geq0$, say that $P\in\mathcal K[z]$ is \emph{bounded by} $(d,o',c')$ if every monomial occurring in $P$ has row degrees
$d$ and, for every monomial $z^M$, the coefficient of $z^M$ in $P$ lies in $\varpi^{\max(c',d_\varpi(M)\dot-o')}\mathcal R$.

\emph{Step 1: the polynomial.} Put $F=\varpi^{-o}F_A$. It lies in $\mathrm{Pol}_T$ with $F_A$, and $\varpi^oF=F_A\in\mathcal R[z]$, as
$f_R$ and the $f_a$ lie in $\mathcal R[z]$. For $g\in GL_k(W)$ whose plane differs from that of $1$,
$\Sigma_gF=\varpi^{-o}\Sigma_gF_A\in\varpi^c\mathcal R[z]$. So $F$ is clean of order $o$.

\emph{Step 2: its substitution.} For $a\in A$ put $t_a=\varpi^{-o_a}\Sigma_1f_a-\sigma_a$, so that
$\Sigma_1f_a=\varpi^{o_a}(t_a+\sigma_a)$. As $\Sigma_1$ is a ring homomorphism, $\Sigma_1f_R=\varpi^{o_R}\sigma_R$ and
$o_R+\sum_am_ao_a=o$,
\[
 \Sigma_1F=\varpi^{-o}\,\Sigma_1f_R\prod_a(\Sigma_1f_a)^{m_a}=\sigma_R\prod_a(t_a+\sigma_a)^{m_a}.
\]

\emph{Step 3: each $t_a$ is bounded by $(d_a,o_a,1)$.} As $\Sigma_1f_a$ and $\sigma_a$ have row degrees $d_a$, so has $t_a$. Let
$z^M$ be a monomial. If $d_\varpi(M)\dot-o_a\leq1$, its coefficient in $t_a$ lies in $\varpi\mathcal R$, as
$t_a=\varpi^{-o_a}(\Sigma_1f_a-\varpi^{o_a}\sigma_a)\in\varpi^{-o_a}\varpi^{o_a+1}\mathcal R[z]$. Otherwise $d_\varpi(M)>o_a+1$. Then
$z^M$ does not occur in $\sigma_a$, whose monomials have $\varpi_0$-degree $o_a$, and, as $f_a\in\mathcal R[z]$, the coefficient
of $z^M$ in $\Sigma_1f_a$ lies in $\varpi^{d_\varpi(M)}\mathcal R$ by
Lemma~\ref{pd-kappa-defects-windows-lift-pairs-source-lift-tail-divisible}. So the coefficient of $z^M$ in $t_a$ lies in
$\varpi^{d_\varpi(M)-o_a}\mathcal R$.

\emph{Step 4: products.} If $P$ is bounded by $(d,o',c')$ and $P'$ by $(d'',o'',c'')$, then $PP'$ is bounded by
$(d+d'',o'+o'',c'+c'')$. Indeed, a monomial occurring in $PP'$ is a product $z^{M_1}z^{M_2}$ with $z^{M_1}$ occurring in $P$
and $z^{M_2}$ in $P'$, so it has row degrees $d+d''$. The coefficient of $z^M$ in $PP'$ is the sum, over $M_1+M_2=M$, of the
products of the coefficients of $z^{M_1}$ in $P$ and of $z^{M_2}$ in $P'$. As $d_\varpi(M)=d_\varpi(M_1)+d_\varpi(M_2)$,
\[
 \max\big(c'+c'',d_\varpi(M)\dot-(o'+o'')\big)\leq\max\big(c',d_\varpi(M_1)\dot-o'\big)+\max\big(c'',d_\varpi(M_2)\dot-o''\big),
\]
so every such product lies in $\varpi^{\max(c'+c'',d_\varpi(M)\dot-(o'+o''))}\mathcal R$. As $1$ is bounded by $(0,0,0)$, Step 3
and induction show that $P_u=\prod_at_a^{u_a}$ is bounded by $(d_u,o_u,s_u)$ for every $u\colon A\to\mathbf Z_{\geq0}$.

\emph{Step 5: the expansion.} Put $U_u=\sigma_R\prod_a\sigma_a^{m_a-u_a}$ and $N_u=\prod_a\binom{m_a}{u_a}$ for the $u$ with
$u_a\leq m_a$ for all $a$. By the binomial formula and Step 2, $\Sigma_1F=\sum_uN_uU_uP_u$, the sum over these $u$. The term
$u=0$ is $\sigma_R\prod_a\sigma_a^{m_a}=\sigma$. Hence $\Sigma_1F-\sigma=\sum_{u\neq0}N_uU_uP_u$.

\emph{Step 6: the terms.} Let $u\neq0$, so that $s_u\geq1$, and write $P_u=\sum_M\pi_Mz^M$. Let $z^M$ occur in $P_u$, and put
$e=\max(s_u,d_\varpi(M)\dot-o_u)$. By Step 4, $z^M$ has row degrees $d_u$, and $\pi_M=\varpi^e\gamma$ with $\gamma\in\mathcal R$.
\begin{itemize}
\item If $e<n$, then $1\leq s_u\leq e$ and $d_\varpi(M)\leq o_u+e$, so condition (e) gives $U_uz^M\in S_e$. Then
  $\pi_MU_uz^M=\varpi^e(\gamma U_uz^M)$ with $\gamma U_uz^M\in S_e$, which lies in $\mathcal T$.
\item If $e\geq n$, then $U_uz^M\in\mathcal R[z]$, as $\sigma_R$ and the $\sigma_a$ lie in $\mathcal R[z]$, so
  $\pi_MU_uz^M\in\varpi^n\mathcal R[z]\subset\mathcal T$.
\end{itemize}

\emph{Step 7.} As $\mathcal T$ is an $\mathcal R$-module, Step 6 gives $N_uU_uP_u=N_u\sum_M\pi_MU_uz^M\in\mathcal T$ for every
$u\neq0$, and Step 5 gives $\Sigma_1F-\sigma\in\mathcal T$. With Step 1, $\sigma\in\mathcal L(\mu,n,o)$.
\end{proof}
% Lean: (1) is the submodule structure of liftSub (add_mem', zero_mem', smul_mem', with cleanSub) and of tailSub (the
%   set T). (2) is FacSys.liftable: raw = C (w^o)^-1 * Ftot and its cleanness (Step 1, inWj_div at the planes); h2 and h3
%   (Step 2); hcls from tail_cls (Step 3: the case split on ndM M - od <= 1, and lift_substitution_coeff_divisible at
%   g = 1); TCls.one, TCls.mul, TCls.pow and TCls.prod (Step 4; TCls w d o c P is "bounded by (d, o, c)"); h4 and h5
%   (Step 5, add_pow and prod_univ_sum); e6, e7 and the case split on e < n with tail_of_ssg and F.core, or tail_of_wj
%   with wj_smul_pow and wj_mono (Step 6). The clause ord of IsFac is not used here, at any g: tail_cls uses inA and
%   sym, and the proof uses inA and sgInt.

\begin{lemma}\label{pd-kappa-defects-windows-lift-pairs-source-lift-tail-columns}
Let the data and the setting be as in Definition~\ref{pd-kappa-defects-windows-lift-pairs-source-lift-tail-systems}, let
$1\leq j\leq q$, $b$ a block of $S_j(b_T)$ and $\mu'$ a real number, and form the generators of $b$ and its column generators
over $\mathcal R$ with $\alpha_i=a_T(y_{j,i})$, $\alpha_+=a_T(s_{j+1})$, $\mu'$ and $r=m$
(Definition~\ref{pd-kappa-defects-windows-lift-pairs-source-lift-tail-systems}(4)). Then every generator of a level of $b$ lies
in the $\mathcal R$-span of the column generators of that level.
\end{lemma}
\leanref{Formalism/SharpRestriction/GrowthBounds/StandardPolynomials/LatticeLeSpanPsiImage.lean}{Formalism.MainNotions.ColumnSystems.Pf.lattice_le_span_psi_image}
\leanref{Formalism/SharpRestriction/GrowthBounds/StandardPolynomials/LatticeLeSpanPsiImage.lean}{Formalism.MainNotions.ColumnSystems.Pf.xiSpan}
\leanref{Formalism/SharpRestriction/GrowthBounds/StandardPolynomials/LatticeLeSpanPsiImage.lean}{Formalism.MainNotions.ColumnSystems.Pf.genSet_le}

\begin{proof}
Use the notation of Definition~\ref{pd-kappa-defects-segment-lattice}(5) with $s=S_j$ and $\eta=\eta_j$, and number the
letters as in Definition~\ref{pd-kappa-defects-windows-lift-pairs-source-lift-tail-systems}(4). Let $G$ be a generator of level
$\ell=(j_1,\dots,j_{s-\eta},j_+)$, the product of
$H=\prod_l\hat E_l^{j_l}\cdot\hat E_+^{j_+}\cdot\prod_lz_l^{\hat b_{u_l}\dot-o_{u_l}}\cdot z_+^{\alpha_+\dot-o_+}$ with $f$ and
$g$. Let $X'$ be the set of the pairs $(\nu,I)$ of a path $\nu$ of rank $m$ for $\delta'$ with $\nu_{2u_l-1}(1)\leq M_l$ for every
$l$ and $\nu_{2s+1}(1)\leq M_+$ and a standard filling $I$ of $\nu$.

\emph{The polynomial $f$.} The set $\mathcal L_l$ consists of the first $2u_l-1$ letters, and $\mathcal L_+$ of all $2s+1$
letters. The ring $\mathcal R$ is a complete discrete valuation ring with $p=\varpi_0^2$ a nonzero element of its maximal ideal
(Definition~\ref{pd-local-model}(1)). We use Steps 1 and 2 of the proof of
Lemma~\ref{pd-kappa-defects-digits-chain-dimension}, with $\mathcal R'=\mathcal R$, $N=2s+1$, $r=m$, $B=\delta'$ and the
conditions $(J,M)$ equal to $(2u_l-1,M_l)$ for $1\leq l\leq s-\eta$ and $(2s+1,M_+)$. The polynomial $f$ has coefficients in
$\mathcal R$ and row degrees $\delta'$, and lies in every $Y(\mathcal L_l,M_l)$ and in $Y(\mathcal L_+,M_+)$. So by Step 2, for
each condition $(J,M)$, $f$ lies in the $\mathcal R$-span of the $\Psi(\nu,I)$, with $\nu$ a path of rank $m$ for $\delta'$ and
$I$ a standard filling of $\nu$, that have $\nu_J(1)\leq M$. By Step 1 all these $\Psi(\nu,I)$ are linearly independent over
$\mathcal R$, so the coefficients of $f$ in them are unique, and they vanish unless $\nu$ satisfies every condition. Hence
$f=\sum_{(\nu,I)\in X'}c_{\nu,I}\Psi(\nu,I)$ with $c_{\nu,I}\in\mathcal R$.

\emph{The polynomial $g$.} It is homogeneous of degree $b_+-j_+$ in the $\xi_c$, with coefficients in $\mathcal R$, so
$g=\sum_\gamma c_\gamma\prod_c\xi_c^{\gamma_c}$ with $c_\gamma\in\mathcal R$, the sum over the $\gamma\in\mathbf Z_{\geq0}^m$ with
$\sum_c\gamma_c=b_+-j_+$.

Hence $G=Hfg=\sum c_{\nu,I}c_\gamma\,H\Psi(\nu,I)\prod_c\xi_c^{\gamma_c}$, and each $H\Psi(\nu,I)\prod_c\xi_c^{\gamma_c}$ is a column
generator of level $\ell$.
\end{proof}
% Lean: genSet_le hA2 hA3 S B (blockData_J_le) for B = blockData of the segment at the margin mu': for a generator in
%   BlockData.gens, lattice_le_span_psi_image gives y in the O-span of the psiR x, x in sAll (hrange, one_cond for Step 2
%   with hA3, mem_span_iInter with the linear independence hli for Step 1, psiR_map), xiSpan gives g in the O-span of the
%   monomials, and key (span_induction in y and g) expands the product. hA2 = HyperplaneRank and hA3 = DefiniteRankOne
%   are clauses 15 and 6 of admitted D: Lemma pd-symbols-path-basis over a discrete valuation ring and Lemma
%   pd-symbols-path-clauses(1). The Lean states it for every BlockData whose J k are at most the number of letters.

\begin{situation}\label{pd-kappa-defects-windows-lift-pairs-source-lift-tail-existence-situation}
Let $C_c\geq0$, $C_0$, $C_{E1}$, $C_{E2}$ and $C_H$ be real numbers. Let $p\geq5$ be a prime and $n'\geq2$ an even integer, and
let $m$, $G'$, $N_1$, $c_{cl}$, $D_p$, $c_*$, $T$, $Y$, $Y_b$, $b'$, $b_l$, $b_u$ and an integer $a$ satisfy hypotheses (1) to (3)
of Theorem~\ref{pd-gain-theorem-one} with $C_0$ and $C_H$. Let $c\geq0$ be an integer with
$c\leq C_c(c_*+|T|k\log_pa+|T|c_{cl}+n'c_{cl})$, and let $a_T,b_T\colon\mathbf Z/n'\to\mathbf Z_{\geq0}$ satisfy
$a\leq a_T(s)\leq a+a/(kY)+2J_0$ and $b_l\leq b_T(s)\leq b_u+2J_0$ for $s\in T$, where $J_0=\Lambda b'/k$. Form the numbers of
Definition~\ref{pd-kappa-defects-numbers} with these constants and any $C_m,C_A,C_t,C_\kappa>0$, and let $\mu=0$ or
$\mu=12c$. Let an unramified quadratic Hida setting of degree $n'$ at $p$ with rings $\mathcal R\subset\mathcal K$ and
$\varpi_0$ be given, with the notation of Definitions~\ref{pd-kappa-defects-windows-lift-pairs-model},
\ref{pd-kappa-defects-windows-lift-pairs-source-lift-relaxed}
and~\ref{pd-kappa-defects-windows-lift-pairs-source-lift-tail-systems}. Let $1\leq n\leq c$, let $\beta=(\beta_j)$ with
$\beta_j$ a block of $S_j(b_T)$ relaxed of slack $(c-n)\kappa_1$, and put $o=\sum_j\operatorname{ord}(\beta_j)$ and
$S_e=S(\mu+(n-e)\kappa_1,(c-n+e)\kappa_1,o+e)$ for $1\leq e<n$. Form the generators of each $\beta_j$ with $\mu+n\kappa_1$, as
in Definition~\ref{pd-kappa-defects-windows-lift-pairs-source-lift-tail-systems}(4). For $1\leq j\leq q$ write $s=S_j$ and
$\eta=\eta_j$, use the notation of Definition~\ref{pd-kappa-defects-segment-lattice} for $\beta_j$ and its generators, with
$\alpha_+=a_T(s_{j+1})$ and $b_+=b_T(s_{j+1})$, number the letters of $\mathcal P_j$ as $0,1,\dots,2s$ as in
Definition~\ref{pd-kappa-defects-windows-lift-pairs-source-lift-tail-systems}(4), and let $\mathcal P'_j\subset\mathcal P_j$ be the
subring of the polynomials in the variables of the letters.
\end{situation}
% Lean: the hypotheses shared by ColsAt, PreAt, TCoreAt and FacSysAt: X : Loc with NeZero X.m and Fact X.p.Prime, a
%   discrete valuation ring O with Algebra (WF p n') O and S : UniformizedDVR p O with S.SqP, D1num X, HIext C0 CYb Cde
%   Ck CD Cb X and CBound Cc X (hypotheses (1) to (3) and the bound on c, with CYb = Cde = Ck = CD = Cb = C_H),
%   mm = 0 or 12 c, 1 <= n <= X.c, and BLk C0 Cde CE1 CE2 X ((c - n) ka1) j (beta j) for every j. Here o is
%   ordFam C0 CE1 CE2 X beta, S_e is SSe C0 Cde CE1 CE2 X O mm n o e, the generators are those of
%   BDk C0 CE1 CE2 (mm + n ka1) X beta j, and P'_j is MvPolynomial (Fin (2 s + 1) x Fin m) (KK O), sent into P_j by
%   rename (chainEmb ...). The book states these in an unramified quadratic Hida setting, as for Lemma
%   pd-kappa-defects-windows-lift-pairs-source-lift-tail.

\begin{definition}\label{pd-kappa-defects-windows-lift-pairs-source-lift-tail-existence-levels}
In Situation~\ref{pd-kappa-defects-windows-lift-pairs-source-lift-tail-existence-situation}:
\begin{enumerate}
\item For $\sigma\in\mathcal K[z]$ and an integer $o'\geq0$, a \emph{lifted factor} of $\sigma$ of order $o'$ is an
  $f\in\mathcal R[z]$ such that $\sigma\in\mathcal R[z]$, $\Sigma_gf\in\varpi_0^{o'}\mathcal R[z]$ for every $g\in GL_k(W)$, and
  $\Sigma_1f-\varpi_0^{o'}\sigma\in\varpi_0^{o'+1}\mathcal R[z]$. The bound at $g=1$ also follows from the last condition, as
  $\sigma\in\mathcal R[z]$. Condition (a) of
  Definition~\ref{pd-kappa-defects-windows-lift-pairs-source-lift-tail-systems}(3) asks that $f_a$ be a lifted factor of
  $\sigma_a$ of order $o_a$.
\item A \emph{level choice} $\ell$ for $\beta$ consists, for each $j$, of a level $\ell_j=(\ell_{j,1},\dots,\ell_{j,s-\eta},\ell_{j,+})$
  of $\beta_j$ (Definition~\ref{pd-kappa-defects-segment-lattice}(5)) and of $\gamma_j\in\mathbf Z_{\geq0}^m$ with
  $\sum_c\gamma_{j,c}=b_+-\ell_{j,+}$. Write $\delta'_j$ for the degrees $\delta'$ of the letters of $\mathcal P_j$ at the level
  $\ell_j$.
\item For a level choice $\ell$, $1\leq j\leq q$ and $G\in\mathcal P'_j$ put
  \[
   W_j(G)=\prod_l\hat E_l^{\ell_{j,l}}\cdot\hat E_+^{\ell_{j,+}}\cdot\prod_lz_l^{\hat b_{u_l}\dot-o_{u_l}}\cdot
   z_+^{\alpha_+\dot-o_+}\cdot G\cdot\prod_c\xi_c^{\gamma_{j,c}},
  \]
  and, for $u\in\mathbf Z_{\geq0}^{s-\eta}$, $H_j(u,G)=\iota_j\big(\prod_l\hat E_l^{u_l}\cdot\hat E_+^{\ell_{j,+}}\cdot G\big)$, all formed
  from $\beta_j$. The \emph{base monomial} of $\ell$ is
  \[
   z^{A_\ell}=\prod_{j=1}^q\iota_j\Big(\beta_j,\prod_lz_l^{\hat b_{u_l}\dot-o_{u_l}}\cdot z_+^{\alpha_+\dot-o_+}\cdot
   \prod_c\xi_c^{\gamma_{j,c}}\Big),
  \]
  a monomial, as $\iota_j$ sends monomials to monomials. For a monomial $z^A$ of $\mathcal K[z]$ write $A(x,c)$ for its exponent at
  $z_{x,c}$.
\end{enumerate}
\end{definition}
\leanref{Formalism/SharpRestriction/GrowthBounds/CoarseMultiplicity/RestRawTail1Of.lean}{Formalism.MainNotions.ColumnSystems.Pf.IsFac}
\leanref{Formalism/SharpRestriction/GrowthBounds/StandardPolynomials/TCoreDef.lean}{Formalism.MainNotions.ColumnSystems.Pf.LvDat}
\leanref{Formalism/SharpRestriction/GrowthBounds/StandardPolynomials/TCoreDef.lean}{Formalism.MainNotions.ColumnSystems.Pf.LvDat.LL}
\leanref{Formalism/SharpRestriction/GrowthBounds/StandardPolynomials/TCoreDef.lean}{Formalism.MainNotions.ColumnSystems.Pf.LvDat.word}
\leanref{Formalism/SharpRestriction/GrowthBounds/StandardPolynomials/TCoreDef.lean}{Formalism.MainNotions.ColumnSystems.Pf.LvDat.head}
\leanref{Formalism/SharpRestriction/GrowthBounds/StandardPolynomials/TCoreDef.lean}{Formalism.MainNotions.ColumnSystems.Pf.LvDat.stdPoly}
\leanref{Formalism/SharpRestriction/GrowthBounds/StandardPolynomials/RestAssumptionsOfThreeInterfaces.lean}{Formalism.MainNotions.ColumnSystems.Pf.LvDat.exists_std}
% Lean: (1) is IsFac w f o sg: inA, ord (for every g with IsGLW g), sym and sgInt. IsFac.mul (RestRawTail1Of.lean) uses
%   ord only at g = 1. (2) is LvDat C0 CE1 CE2 (mm + n ka1) X beta: lv, jp, ok (LevelOK), gamma and hgamma; delta'_j is
%   (L.LL j).deg, the degrees of lettersAt (lv j) (jp j). (3): W_j is L.word O j, H_j(u, G) is L.head O j u G (with ioD =
%   ioA, the map iota_j), and A_ell is the Std of LvDat.exists_std, with monomial Std 1 = prod_j L.stdPoly O j, where
%   L.stdPoly O j = digitMon * sTopMon * ioD of the monomial (LvDat.stdPoly_mono).

\begin{definition}\label{pd-kappa-defects-windows-lift-pairs-source-lift-tail-existence-column-systems}
In Situation~\ref{pd-kappa-defects-windows-lift-pairs-source-lift-tail-existence-situation}, let $\ell$ be a level choice for
$\beta$, let $1\leq j\leq q$ and $G\in\mathcal P'_j$, and write $\varsigma_i,o_i,w_i$ for the entries of $\beta_j$ and $\hat b_i$, $\hat a$
and $P_i=p^{g_j-d_{j,i}}$ for its numbers (Definition~\ref{pd-kappa-defects-segment-lattice}(1)). A \emph{column system} for
$G$ at $j$ consists of a finite set $I$, an integer $a_0\geq0$ and, for each $a\in I$, a polynomial $K_a\in\mathcal P'_j$, a
map $d_a$ from the letters of $\mathcal P_j$ to $\mathbf Z_{\geq0}$, a monomial $z^{\pi_a}$ of $\mathcal K[z]$, a polynomial
$f_a\in\mathcal K[z]$, an integer $o_a\geq0$ and a map $e_a$ from the letters of
Definition~\ref{pd-kappa-defects-windows-lift-pairs-model}(1) to $\mathbf Z_{\geq0}$, such that, with
$\sigma_a=z^{\pi_a}\iota_j(K_a)$:
\begin{enumerate}
\item each $K_a$ has coefficients in $\mathcal R$ and row degrees $d_a$, with $d_a(\lambda)\leq1$ for every letter $\lambda$;
  $\sum_ad_a=\delta'_j$; and $\prod_aK_a=G$;
\item each $f_a$ is a lifted factor of $\sigma_a$ of order $o_a$
  (Definition~\ref{pd-kappa-defects-windows-lift-pairs-source-lift-tail-existence-levels}(1)), every monomial of $\sigma_a$ has
  $\varpi_0$-degree $o_a$, and $f_a$ and $\sigma_a$ have row degrees $e_a$;
\item the row degrees of $\pi_a$ and the values of $e_a$ are at most $Q$, and $e_a$ vanishes at the letters $(s_{j'},0)$,
  $1\leq j'\leq q$;
\item $a_0+\sum_{i\leq\eta}\varsigma_i\leq\hat a$; the exponent vector $\sum_a\pi_a$ is $p^{g_j}a_0$ at $z_{(s_j,1),1}$,
  $P_i(\hat b_i\dot-w_i)$ at $z_{(y_{j,i},1),1}$ for $i\leq\eta$, $P_i(\hat b_i\dot-o_i)$ at $z_{(y_{j,i},1),1}$ for $i>\eta$, and $0$
  at every other variable; and $\sum_ao_a=\sum_{i=1}^sP_i\delta'_j(x_i)+\delta'_j(x_+)$.
\end{enumerate}
Of the bounds $\Sigma_gf_a\in\varpi_0^{o_a}\mathcal R[z]$ in (2), the proof of
Lemma~\ref{pd-kappa-defects-windows-lift-pairs-source-lift-tail-existence-prefactorization} uses those at the $g$ whose
plane differs from that of $1$ (Lemma~\ref{pd-symbols-edge-basis}, with $r=1$), and no others.
\end{definition}
\leanref{Formalism/SharpRestriction/GrowthBounds/StandardPolynomials/ColsDef.lean}{Formalism.MainNotions.ColumnSystems.Pf.ColSys}
% Lean: ColSys C0 CE1 CE2 X O w (mm + n ka1) beta L j G: iota, fin, Kc, dK, pl, fC, oC, dgC and a0 are I, K_a, d_a, pi_a,
%   f_a, o_a, e_a and a_0. (1) is KInt (IntPol), KPol, digit_le_one, dKsum and Kprod; (2) is facC, homC, fPol and polC;
%   (3) is plQ, dgCQ and dgCY, where topEnd X j' is s_{j'+1}; (4) is a0le, plSum and oSum, with Qj = p^g_j, A0 = a-hat,
%   coarseA = b-hat_i, sigma = eta (the Lean's i < sigma, from 0, is the book's i <= eta, from 1), topStart X j = s_j,
%   topOdd X j i = y_{j,i}, and the letters 2 i and 2 s of L.LL j the x_i and x_+. The bounds facC.ord used are those in
%   cleaned_product_order_family_eq (RigFactorWeightedHomogeneousLetter.lean), at g' with IsGLW g' and ratPlane g' different
%   from ratPlane 1; nothing else reads a ColSys's facC.ord.

\begin{definition}\label{pd-kappa-defects-windows-lift-pairs-source-lift-tail-existence-columns-data}
In Situation~\ref{pd-kappa-defects-windows-lift-pairs-source-lift-tail-existence-situation}, let $\ell$ be a level choice for
$\beta$ and $1\leq j\leq q$, let $\nu$ be a path of rank $m$ for $\delta'_j$ with $\nu_{2u_l-1}(1)\leq M_l$ for every $l$ and
$\nu_{2s+1}(1)\leq M_+$, and let $I$ be a standard filling of $\nu$. Use the notation of the situation for $\beta_j$, with
$P_i=p^{g_j-d_{j,i}}$, put $\mu'=\mu+n\kappa_1$, and let $o'_i$ be $w_i$ for $i\leq\eta$ and $o_i$ for $i>\eta$.
\begin{enumerate}
\item The \emph{augmented letters} are $0,1,\dots,2s+1$, with $s_j$ taken in $\{0,\dots,n'-1\}$: the letter $0$ is a vector
  letter of position $s_j$; for $1\leq i\leq s$
  the letters $2i-1$ and $2i$ are a vector and a covector letter of position $s_j+d_{j,i}$; and $2s+1$ is a vector letter of
  position $\zeta=s_j+g_j$. For $t\geq1$, the augmented letter $t$ is the letter of $\mathcal P_j$ numbered $t-1$. Let
  $\lambda_0=(s_j,1)$, $\lambda_{2i-1}=(y_{j,i},1)$, $\lambda_{2i}=(y_{j,i},0)$ and $\lambda_{2s+1}=(s_{j+1},1)$, letters of
  Definition~\ref{pd-kappa-defects-windows-lift-pairs-model}(1), and $P^S_t=p^{\zeta-\pi_t}$, where $\pi_t$ is the position of $t$;
  so $P^S_0=p^{g_j}$, $P^S_{2i-1}=P^S_{2i}=P_i$ and $P^S_{2s+1}=1$. Let $\phi$ be the $\mathcal K$-algebra map from the polynomial ring
  in the $z_{t,b}$, $0\leq t\leq2s+1$, $1\leq b\leq k$, to $\mathcal K[z]$ with $z_{t,b}\mapsto z_{\lambda_t,b}$.
\item Put $a_0=\lfloor\hat a-\mu'\rfloor-\sum_{i\leq\eta}\varsigma_i$, $\theta(0)=0$, and
  $\theta(T)=a_0+\sum_{i\colon2i\leq T}(\hat b_i\dot-o'_i)$ for $1\leq T\leq2s+2$. The \emph{augmented sequence}
  $\tilde\nu=(\tilde\nu_0,\dots,\tilde\nu_{2s+2})$ in $\mathbf Z^k$ has $\tilde\nu_0=0$, and $\tilde\nu_{t+1}$ has first entry
  $\theta(t+1)$ and last $m$ entries those of $\nu_t$.
\item For $c\in\mathbf Z$ with $|c|\leq T_\nu$ or $1\leq c\leq\theta(2s+2)$ put, in the notation of
  Definition~\ref{pd-symbols-path-generators}:
  \begin{itemize}
  \item $K_c=\mathrm{Col}_{\nu,c}(I(c))\in\mathcal P'_j$, and $d_c(\lambda)=1$ if the letter $\lambda$ of $\mathcal P_j$ is used for
    $c$ in $\nu$ and $d_c(\lambda)=0$ otherwise;
  \item $z^{\pi_c}=\prod_tz_{\lambda_t,1}^{P^S_t}$, the product over the $t$ with $\theta(t)<c\leq\theta(t+1)$;
  \item $J_c\subset\{1,\dots,k\}$ the set of the $b+1$ with $b\in I(c)$, together with $1$ if $1\leq c\leq\theta(2s+2)$;
  \item $f_c=\phi\big(\mathrm{Col}^1_{\tilde\nu,c,\zeta}(J_c)\big)$, the lifted column function of item (6) there for the augmented
    letters, with $n=k$ and $g=1$;
  \item $o_c=o_\omega(\tilde\nu,c)$, with the splitting $k=1+m$ of item (5) there and $\omega(t)=P^S_t$;
  \item $e_c(x)$ the sum of the $P^S_t$ over the augmented letters $t$ with $\lambda_t=x$ that are used for $c$ in $\tilde\nu$.
  \end{itemize}
\end{enumerate}
\end{definition}
\leanref{Formalism/SharpRestriction/GrowthBounds/StandardPolynomials/ColDef.lean}{Formalism.MainNotions.ColumnSystems.Pf.lamS}
\leanref{Formalism/SharpRestriction/GrowthBounds/StandardPolynomials/ColDef.lean}{Formalism.MainNotions.ColumnSystems.Pf.posS}
\leanref{Formalism/SharpRestriction/GrowthBounds/StandardPolynomials/ColDef.lean}{Formalism.MainNotions.ColumnSystems.Pf.LLseg}
\leanref{Formalism/SharpRestriction/GrowthBounds/StandardPolynomials/ColDef.lean}{Formalism.MainNotions.ColumnSystems.Pf.zS}
\leanref{Formalism/SharpRestriction/GrowthBounds/StandardPolynomials/ColDef.lean}{Formalism.MainNotions.ColumnSystems.Pf.phiS}
\leanref{Formalism/SharpRestriction/GrowthBounds/StandardPolynomials/ColsS.lean}{Formalism.MainNotions.ColumnSystems.Pf.PS}
\leanref{Formalism/SharpRestriction/GrowthBounds/StandardPolynomials/ColsS.lean}{Formalism.MainNotions.ColumnSystems.Pf.a0Z}
\leanref{Formalism/SharpRestriction/GrowthBounds/StandardPolynomials/ColsS.lean}{Formalism.MainNotions.ColumnSystems.Pf.nuS}
\leanref{Formalism/SharpRestriction/GrowthBounds/StandardPolynomials/NotCostlyZero.lean}{Formalism.MainNotions.ColumnSystems.Pf.Col.muP}
\leanref{Formalism/SharpRestriction/GrowthBounds/StandardPolynomials/ColsS.lean}{Formalism.MainNotions.ColumnSystems.Pf.colsS}
\leanref{Formalism/SharpRestriction/GrowthBounds/StandardPolynomials/ColsS.lean}{Formalism.MainNotions.ColumnSystems.Pf.dKS}
\leanref{Formalism/SharpRestriction/GrowthBounds/StandardPolynomials/ColsS.lean}{Formalism.MainNotions.ColumnSystems.Pf.plS}
\leanref{Formalism/SharpRestriction/GrowthBounds/StandardPolynomials/SymbolEq.lean}{Formalism.MainNotions.ColumnSystems.Pf.Col.Jc}
\leanref{Formalism/SharpRestriction/GrowthBounds/StandardPolynomials/ColsS.lean}{Formalism.MainNotions.ColumnSystems.Pf.fCS}
\leanref{Formalism/SharpRestriction/GrowthBounds/StandardPolynomials/ColsS.lean}{Formalism.MainNotions.ColumnSystems.Pf.oCS}
\leanref{Formalism/SharpRestriction/GrowthBounds/StandardPolynomials/ColsS.lean}{Formalism.MainNotions.ColumnSystems.Pf.dgCS}
% Lean: the augmented letters are Fin (2 S_j + 2), from 0, with lamS (the letter of K[z]), posS (the position, as an
%   integer s_j + d, not reduced modulo n'), LLseg (types and positions) and zS = zeta; PS t = p^(zeta - pi_t) and phiS
%   is phi. a0Z is a_0 (an integer; the Lean's sums over l < sigma, from 0, are the book's i <= eta), nuS is theta, with
%   the sum over the l from 0 with 2 l + 2 <= T, and muP nuS ka (joinPT of the first row and kaS ka = Fin.cons 0 ka) is the
%   augmented sequence, where ka = y.1.1.mu is nu. colsS is the set of c, colRange ka being |c| <= T_nu; Kcol (KK O) ka c
%   (2 s + 1) (I'' c) is K_c, dKS is d_c, plS is pi_c, Jc is J_c (joinI with the plane part {0}), fCS is f_c = rename
%   (phiS) of KKc (LLseg) 1 zS (muP ..) c (Jc ..), oCS is oc (PS) (muP ..) c and dgCS is e_c (dgS with all letters).

\begin{lemma}\label{pd-kappa-defects-windows-lift-pairs-source-lift-tail-existence-columns-numerics}
For every $C_c\geq0$ there is $T_0>0$ such that for every $C_0\geq T_0$ there is $T_1>0$ such that for all
$C_{E1},C_{E2}\geq T_1$ there is $T_2>0$ with the following property for every $C_H\geq T_2$. In
Situation~\ref{pd-kappa-defects-windows-lift-pairs-source-lift-tail-existence-situation} with these constants, let
$1\leq j\leq q$, use the notation of the situation for $\beta_j$, and put $\mu'=\mu+n\kappa_1$. Then
$\sum_{i\leq\eta}\varsigma_i\leq\lfloor\hat a-\mu'\rfloor$, and $\sum_{l\leq i}w_l+\sum_{l\leq\eta}\varsigma_l\leq\lfloor\hat a-\mu'\rfloor$
for every $1\leq i\leq\eta$.
\end{lemma}
\leanref{Formalism/SharpRestriction/GrowthBounds/StandardPolynomials/ColNumDef.lean}{Formalism.MainNotions.ColumnSystems.Pf.ColumnNumericsAt}
\leanref{Formalism/SharpRestriction/GrowthBounds/StandardPolynomials/ColNumDef.lean}{Formalism.MainNotions.ColumnSystems.Pf.ColumnNumericsBeyond}
\leanref{Formalism/SharpRestriction/GrowthBounds/StandardPolynomials/ColXNum.lean}{Formalism.MainNotions.ColumnSystems.Pf.GenusGluing.m1num_of_std}
\leanref{Formalism/SharpRestriction/GrowthBounds/StandardPolynomials/ColXNum.lean}{Formalism.MainNotions.ColumnSystems.Pf.GenusGluing.column_numerics_holds}

\begin{proof}
\emph{The constants.} Let $T_0=T_1=1$ and $T_2=1000(C_{E1}+C_{E2}+1)(C_c+1)$, and let the data be as in the statement.
Write $s=S_j$, $\eta=\eta_j$, $P=p^{g_j}$, $P_i=p^{g_j-d_{j,i}}$ and $B_i=b_T(y_{j,i})$, put $\varepsilon=N_1\Lambda/k$, and let
$\delta$ be the number of Theorem~\ref{pd-gain-theorem-one} and $\beta_e=(b_u-b_l)/b'+2\Lambda/k$. The block $\beta_j$ is relaxed of
slack $E=(c-n)\kappa_1$, where $\kappa_1=13kN_1Q$.

\emph{Facts.} The facts (a) to (g) in the proof of Lemma~\ref{pd-kappa-defects-margin} and the fact (h) in the proof of
Lemma~\ref{pd-kappa-defects-windows-model-step} hold. They give $k\geq2N_1+3\geq5$, $Q\geq1$, $Y\geq1$, $N_1\geq1$, $N\geq2$,
$\Lambda\geq1$, $kb'\leq a$, $k^3NQ^4cY\leq b'$, $k^3Q\leq b'$ and $E_0+Q\leq\varepsilon a$; and $k^2\leq k^2Y^2\leq p$ by
hypothesis (1) of Theorem~\ref{pd-gain-theorem-one}. As $C_H\geq1000$ and $b_l\leq b_u$, $\delta\geq C_HN_1\Lambda/k\geq40\varepsilon$,
and $\delta<1/6$ by hypothesis (2) of the theorem; so $\varepsilon\leq1/240$, and $1/k\leq\varepsilon$ as $N_1\Lambda\geq1$. As in the
proof of Lemma~\ref{pd-kappa-defects-windows-lift-pairs-source-order}, $b'(1-\beta_e)\leq b_T(t)\leq b'(1+\beta_e)$ for $t\in T$, and
$0\leq\beta_e\leq\delta/2<1/12$. By hypothesis (1), $s\leq N_1$ and $g_j\leq G'$, so $P_i\leq P\leq Q$ for every $i$.

\emph{Step 1: the size of $\hat a$.} By the definition of the floor, $\hat aP>a-E_0-P\geq a-(E_0+Q)\geq(1-\varepsilon)a\geq a/2$; so
$\hat a>0$. As $N\geq2$ and $Y\geq1$, $2k^3Q^4c\leq b'$, so $k^4Q^4c\leq kb'/2\leq a/2<\hat aP\leq\hat aQ$, and $k^4Q^3c<\hat a$. Hence
$k^3Q^3c\leq\hat a/k\leq\varepsilon\hat a$. As $2N_1\leq k$, $k^3\geq125$ and $Q\geq1$,
\[
 12c+13kN_1^2Qc\leq\tfrac{12}{125}k^3Q^3c+\tfrac{13}4k^3Q^3c\leq4\varepsilon\hat a .
\]

\emph{Step 2: the margin.} As $\mu\leq12c$, $n\leq c$ and $N_1\geq1$,
$\mu'+N_1E=\mu+n\kappa_1+N_1(c-n)\kappa_1\leq12c+N_1c\kappa_1=12c+13kN_1^2Qc\leq4\varepsilon\hat a$, by Step 1.

\emph{Step 3: the sum of the $\varsigma_i$.} By condition (c) of a relaxed block, $\varsigma_i\leq2\Lambda\hat a/k+E$ for $i\leq\eta$, and
this bound is $\geq0$. As $\eta\leq s\leq N_1$, $\sum_{i\leq\eta}\varsigma_i\leq N_1(2\Lambda\hat a/k+E)=2\varepsilon\hat a+N_1E$.

\emph{The first inequality.} By Steps 2 and 3, $\sum_{i\leq\eta}\varsigma_i+\mu'\leq6\varepsilon\hat a\leq\hat a$. The sum is an integer,
so it is at most $\lfloor\hat a-\mu'\rfloor$.

\emph{The second inequality.} Let $1\leq i\leq\eta$. By condition (a) of a relaxed block, for every $l$ the $k$ entries of $\rho_l$ are
less than $P_l$ and $P_l$ divides $B_l-|\rho_l|$, so $\hat b_lP_l=B_l-|\rho_l|$ and $B_l-kP_l<\hat b_lP_l\leq B_l$. As
$2kQ\leq k^3Q\leq b'$ and $\beta_e<1/12$, $\hat b_iP_i>B_i-kQ\geq b'(1-\beta_e)-b'/2\geq5b'/12$. Let $l<i$. As $d_{j,l}<d_{j,i}$,
$pP_i\leq P_l$, so $p\hat b_lP_i\leq\hat b_lP_l\leq B_l\leq b'(1+\beta_e)\leq13b'/12<3\hat b_iP_i$, and $p\hat b_l\leq3\hat b_i$. By
condition (c), $w_l\leq\hat b_l\leq3\hat b_i/p$, and $w_i\leq\hat b_i$. There are fewer than $s\leq N_1$ indices $l<i$, and
$3N_1\leq kN_1\leq\varepsilon k^2\leq\varepsilon p$, as $N_1\leq N_1\Lambda=\varepsilon k$; so
\[
 \sum_{l\leq i}w_l\leq\hat b_i+3N_1\hat b_i/p\leq(1+\varepsilon)\hat b_i .
\]
Condition (a) of a relaxed block is the first condition of a kept block, so $\rho,\rho'$ are digits for $j$, and
Lemma~\ref{pd-kappa-defects-digits-count-coarse}(2) gives $\hat b_i(1+\delta)<\hat a$. As $\delta<1/6$ and $8\varepsilon\leq\delta$,
\[
 \big((1+\varepsilon)\hat b_i+6\varepsilon\hat a\big)(1+\delta)\leq(1+\varepsilon)\hat a+6\varepsilon\hat a+\varepsilon\hat a\leq(1+\delta)\hat a,
\]
so $(1+\varepsilon)\hat b_i+6\varepsilon\hat a\leq\hat a$. With Steps 2 and 3,
$\sum_{l\leq i}w_l+\sum_{l\leq\eta}\varsigma_l+\mu'\leq(1+\varepsilon)\hat b_i+2\varepsilon\hat a+4\varepsilon\hat a\leq\hat a$, and as the two sums are
integers, their total is at most $\lfloor\hat a-\mu'\rfloor$.
% Lean: column_numerics_holds = TransportedAlt.mono of m1num_of_std (ColXNum.lean) over threshold_standard_one (thresholds
%   1, 1 and T2 CE1 CE2 Cc 1), which gives Std C0 Cde CE1 CE2 1 X (the facts; St.N5 is (h), St.forty_epsilon_le_delta,
%   St.de6, St.eps_ge, St.be18 and St.two_rank_mul_Q_le). Step 1 is column_first_coefficient_bound (hA0lt from
%   Int.sub_one_lt_floor); Step 2 is column_margin_bound; Step 3 is column_partial_sum_S with the bound (c) on sigma_i
%   of BLk; the first inequality is column_first_goal with Int.le_floor. The second is coarse_bounds (for l and i),
%   column_lower_ratio (hlow, with H.bT, H.dmono and dist_pos_lt), column_partial_sum_W, three_level_le_eps_prime,
%   coarse_degree_alternative (first part) and column_second_goal. The statement does not use the Hida setting, mu
%   beyond mu <= 12 c, or the generators.
\end{proof}

\begin{lemma}\label{pd-kappa-defects-windows-lift-pairs-source-lift-tail-existence-columns-shape}
In Situation~\ref{pd-kappa-defects-windows-lift-pairs-source-lift-tail-existence-situation}, let $\ell$, $j$, $\nu$ and $I$ be as
in Definition~\ref{pd-kappa-defects-windows-lift-pairs-source-lift-tail-existence-columns-data}, with its notation, and suppose
that the two inequalities of Lemma~\ref{pd-kappa-defects-windows-lift-pairs-source-lift-tail-existence-columns-numerics} hold
for $\beta_j$. Number the letters of $\mathcal P_j$ as $0,\dots,2s$. Then:
\begin{enumerate}
\item $a_0\geq0$, $\theta(0)=0$, $a_0\leq\theta(T)$ for $T\geq1$, and $\theta(T)\leq\theta(T+1)$ for $0\leq T\leq2s+1$;
\item $\theta(i+2)=\theta(i+1)$ for every covector letter $i$ of $\mathcal P_j$;
\item $\nu_{i+1}(b)\leq\theta(i+1)$ for every vector letter $i$ of $\mathcal P_j$ and $1\leq b\leq m$;
\item $2s+1<m$.
\end{enumerate}
\end{lemma}
\leanref{Formalism/SharpRestriction/GrowthBounds/StandardPolynomials/NotCostlyZero.lean}{Formalism.MainNotions.ColumnSystems.Pf.Col.PH}
\leanref{Formalism/SharpRestriction/GrowthBounds/StandardPolynomials/ColXPH.lean}{Formalism.MainNotions.ColumnSystems.Pf.GenusGluing.floor_cap}
\leanref{Formalism/SharpRestriction/GrowthBounds/StandardPolynomials/ColXPH.lean}{Formalism.MainNotions.ColumnSystems.Pf.GenusGluing.first_row_le}
\leanref{Formalism/SharpRestriction/GrowthBounds/StandardPolynomials/ColXPH.lean}{Formalism.MainNotions.ColumnSystems.Pf.GenusGluing.sum_deg_le}
\leanref{Formalism/SharpRestriction/GrowthBounds/StandardPolynomials/ColXPH.lean}{Formalism.MainNotions.ColumnSystems.Pf.GenusGluing.ph_seg}

\begin{proof}
(1) The first inequality of Lemma~\ref{pd-kappa-defects-windows-lift-pairs-source-lift-tail-existence-columns-numerics} is
$a_0\geq0$. By conditions (b) and (c) of a relaxed block, $o'_i\leq\hat b_i$, so the terms $\hat b_i\dot-o'_i=\hat b_i-o'_i$ of
$\theta$ are $\geq0$, and $\theta(T)$ grows with $T\geq1$, from $\theta(1)=a_0\geq0=\theta(0)$.

(2) The letter $i$ is a covector letter if and only if $i=2l-1$ with $1\leq l\leq s$. Then the sums defining $\theta(2l)$ and
$\theta(2l+1)$ run over the same indices $i'\leq l$.

(3) Let $L\geq\eta$ be an integer with $L\leq s$. By Definition~\ref{pd-kappa-defects-segment-lattice}(1),
$\operatorname{cap}(L)=\hat a+\sum_{i\leq L}(\hat b_i-o^\flat_i)=\hat a-\sum_{i\leq\eta}\varsigma_i+\sum_{i\leq L}(\hat b_i-o'_i)$, and the
two sums are integers; so $\lfloor\operatorname{cap}(L)-\mu'\rfloor=\theta(2L+1)$. Let $i$ be a vector letter of $\mathcal P_j$.
As $\nu_{i+1}$ is dominant, $\nu_{i+1}(b)\leq\nu_{i+1}(1)$.
\begin{itemize}
\item If $i=2s$, the letter $x_+$, then $\nu_{2s+1}(1)\leq M_+=\lfloor\operatorname{cap}(s)-\mu'\rfloor=\theta(2s+1)$
  (Definition~\ref{pd-kappa-defects-segment-lattice}(5)).
\item If $i=2u_l-2$, the letter $x_{u_l}$ with $u_l>\eta$, then
  $\nu_{2u_l-1}(1)\leq M_l=\lfloor\operatorname{cap}(u_l-1)-\mu'\rfloor=\theta(2u_l-1)$, as $u_l-1\geq\eta$.
\item If $i=2l-2$ with $l\leq\eta$, the letter $x_l$: at a vector letter $i'$ the first entry of a path grows by at most the
  degree $\delta'_j(i')$, and at a covector letter it does not grow (Definition~\ref{pd-symbols-path-generators}(2)); so
  $\nu_{2l-1}(1)\leq\sum_{l'\leq l}\delta'_j(x_{l'})=\sum_{l'\leq l}w_{l'}$, as the level lowers only the degrees of the $x_{u_l}$,
  $y_{u_l}$ and $x_+$, and $\delta(x_{l'})=w_{l'}$ for $l'\leq\eta$. By the second inequality of
  Lemma~\ref{pd-kappa-defects-windows-lift-pairs-source-lift-tail-existence-columns-numerics}, this is at most $a_0$, which is at
  most $\theta(2l-1)$ by (1).
\end{itemize}

(4) By hypothesis (1) of Theorem~\ref{pd-gain-theorem-one}, $s=S_j\leq N_1$ and $2N_1+2\leq m$.
\end{proof}
% Lean: ph_seg L j y hy hB hnum1 hnum2 hm (ColXPH.lean) gives Col.PH (L.LL j).vec (L.LL j).deg mu (nuS ..): ka0, kadom and
%   kastep are those of the path y.1.1; nu0 is h0; numono is (1) with hge and ha0; nucov is (2); C is (3), with floor_cap
%   (hfl, for L0 >= sigma), bd_J_last and bd_thr_last for x_+, bd_J_lt and bd_thr_lt for x_{u_l}, and first_row_le with
%   sum_deg_le and hnum2 for l < sigma; hK is (4), passed in colsAt_of as hm from HIext (segS j <= N1 and 2 N1 + 2 <= m).

\begin{lemma}\label{pd-kappa-defects-windows-lift-pairs-source-lift-tail-existence-columns-lifted}
Let $p\geq5$ be a prime, $(W,Q,\sigma,\iota)$ a Hida setting of degree $n'$ at $p$ with rings $\mathcal R\subset\mathcal K$ and
$\varpi_0$, where $\varpi_0$ is a uniformizer of $\mathcal R$, and $N,r,K\geq1$ integers with $n=r+K$. Fix a lettering of
$\{0,\dots,N-1\}$ whose letters have positions $\pi_0\leq\pi_1\leq\dots\leq\pi_{N-1}$, an integer $\zeta\geq\pi_{N-1}$, and use the
notation of Definition~\ref{pd-symbols-path-generators}, with $\omega(i)=p^{\zeta-\pi_i}$. Let $\mu=(\mu_0,\dots,\mu_N)$ be a
sequence in $\mathbf Z^n$ and $c\in\mathbf Z$, and write $h_i=h_i(c)$. Say that $\mu$ satisfies the \emph{column conditions} for $c$
if, for $0\leq i<N$:
(C1) if $h_i<h_{i+1}$, then $i$ is a vector letter and $h_{i+1}=h_i+1$;
(C2) if $h_{i+1}<h_i$, then $i$ is a covector letter and $h_i=h_{i+1}+1$;
(C3) if $c\geq1$, then $h_0=0$;
(C4) if $c\leq0$, then $h_{i'}>r$ for $0\leq i'\leq N$.
\begin{enumerate}
\item If $N+r\leq K$, every path of rank $n$ for some $B$ satisfies the column conditions for every $c$.
\item For every $g\in GL_n(W)$ and $J\subset\{1,\dots,n\}$, $\mathrm{Col}^g_{\mu,c,\zeta}(J)$ has its coefficients in $\mathcal R$.
\item Suppose that $\mu$ satisfies the column conditions for $c$ and that $\pi_i<\pi_{i'}$ whenever $i<i'$, $i$ is used for $c$
  and $i'$ is costly for $c$. Put $o=o_\omega(\mu,c)$, and let $g\in GL_n(W)$ and $J\subset\{1,\dots,n\}$. Then
  $\Sigma^{\mathrm L}_{g'}(\mathrm{Col}^g_{\mu,c,\zeta}(J))\in\varpi_0^o\mathcal R[z]$ for every $g'\in GL_n(W)$.
\item If moreover every $\mu_i$ is dominant, then
  $\Sigma^{\mathrm L}_g(\mathrm{Col}^g_{\mu,c,\zeta}(J))-\varpi_0^o\,\Pi_\zeta(\mathrm{Col}^{\mathrm{pt}}_{\mu,c}(J))\in\varpi_0^{o+1}\mathcal R[z]$.
\end{enumerate}
\end{lemma}
\leanref{Formalism/SharpRestriction/GrowthBounds/FrobeniusPaths/ColOKOfPath.lean}{Formalism.MainNotions.LemmaBAlt.ColOK}
\leanref{Formalism/SharpRestriction/GrowthBounds/FrobeniusPaths/ColOKOfPath.lean}{Formalism.MainNotions.LemmaBAlt.colOK_of_path}
\leanref{Formalism/SharpRestriction/GrowthBounds/FrobeniusPaths/ColumnKernelComponentsIntegral.lean}{Formalism.MainNotions.LemmaBAlt.column_kernel_components_integral}
\leanref{Formalism/SharpRestriction/GrowthBounds/FrobeniusPaths/LiftedColumnOrder.lean}{Formalism.MainNotions.LemmaBAlt.lifted_column_order}
\leanref{Formalism/SharpRestriction/GrowthBounds/FrobeniusPaths/LiftedColumnCongruence.lean}{Formalism.MainNotions.LemmaBAlt.lifted_column_congruence}

\begin{proof}
(1) Let $\mu$ be a path. As in the proof of Lemma~\ref{pd-symbols-path-clauses}(3), $h_i\leq h_{i+1}\leq h_i+1$ at a vector letter
$i$, as $\mu_{i+1}$ is a strip over $\mu_i$, and $h_{i+1}\leq h_i\leq h_{i+1}+1$ at a covector letter; this gives (C1) and (C2). As
$\mu_0=0$, $h_0=0$ for $c\geq1$, which is (C3), and $h_0=n$ for $c\leq0$. In the latter case $h$ drops by at most one at each letter,
so $h_{i'}\geq n-i'\geq r+K-N\geq2r>r$, which is (C4).

(2) The proof of Lemma~\ref{pd-symbols-path-clauses}(5)(a) uses nothing about $\mu$.

(3) and (4). Lemmas~\ref{pd-symbols-path-lift-nominal}, \ref{pd-symbols-path-lift-order} and~\ref{pd-symbols-path-lift-residue}
are stated for a path with $N+r\leq K$, but their proofs use these only through (C1) to (C4). In the proof of
Lemma~\ref{pd-symbols-path-lift-nominal}, the facts on the heights derived from the path at its start are (C1) to (C4). The proofs of
Lemma~\ref{pd-symbols-path-lift-order} and of Lemma~\ref{pd-symbols-path-lift-residue}(1) and (2) use the path only through
Lemma~\ref{pd-symbols-path-lift-nominal}, the definitions of used and costly letters and the positions. The proof of
Lemma~\ref{pd-symbols-path-lift-residue}(3) runs through the cases (i) to (vi) of the proof of
Lemma~\ref{pd-symbols-path-clauses}(3), which use of the path only (C1) to (C4) (the changes of the heights, $h_0=0$ for $c\geq1$ and
$h_j\geq r$ for $c\leq0$) and the dominance of the $\mu_j$ (the numbers of entries at least $c$ of $\mu^{\mathrm{pl}}_j$ and
$\mu^{\mathrm{tn}}_j$ are $\min(h_j,r)$ and $(h_j-r)^+$). So these three lemmas hold for $\mu$, with (3) of the last only when the
$\mu_j$ are dominant, and the proofs of Lemma~\ref{pd-symbols-path-clauses}(5)(c) and (d) give (3) and (4).
\end{proof}
% Lean: (1) is colOK_of_path (colH_strip); (2) is column_kernel_components_integral (no hypothesis on mu); (3) is
%   lifted_column_order, with H : ColOK LL mu c, Monotone LL.pos, LL.pos i <= z and the position condition hposc; (4) is
%   lifted_column_congruence with SgVanish_of LL mu c H and plTanCol_eq_SgW O LL mu c H hdom. The Lean's lemmas
%   Ord_main, Inv_main, SgW_supp, SgW_eq and step_nom take ColOK (and Dominant for SgW_eq) as hypotheses; the book's
%   Lemmas pd-symbols-path-lift-nominal, -order and -residue state them for a path with N + r <= K, which
%   colOK_of_path turns into ColOK. canonical_translate_jet_identity (Lemma pd-symbols-path-clauses(5)) is (1) to (4)
%   for a path, with P.dom.

\begin{lemma}\label{pd-kappa-defects-windows-lift-pairs-source-lift-tail-existence-columns-augmented}
In Situation~\ref{pd-kappa-defects-windows-lift-pairs-source-lift-tail-existence-situation}, let $\ell$, $j$, $\nu$ and $I$ be as
in Definition~\ref{pd-kappa-defects-windows-lift-pairs-source-lift-tail-existence-columns-data}, with its notation, and suppose that
the conclusions (1) to (4) of Lemma~\ref{pd-kappa-defects-windows-lift-pairs-source-lift-tail-existence-columns-shape} hold. Give the
augmented letters their types and their positions $\pi_t$, and use the splitting $k=1+m$ (Definition~\ref{pd-symbols-path-generators}(5)
with $r=1$ and $K=m$). Then:
\begin{enumerate}
\item every $\tilde\nu_t$ is dominant;
\item $\tilde\nu$ satisfies the column conditions for every $c\in\mathbf Z$
  (Lemma~\ref{pd-kappa-defects-windows-lift-pairs-source-lift-tail-existence-columns-lifted}, with $N=2s+2$, $r=1$ and $K=m$);
\item for every $c$, $\pi_t<\pi_{t'}$ whenever $t<t'$, $t$ is used for $c$ and $t'$ is costly for $c$.
\end{enumerate}
\end{lemma}
\leanref{Formalism/SharpRestriction/GrowthBounds/StandardPolynomials/NotCostlyZero.lean}{Formalism.MainNotions.ColumnSystems.Pf.Col.PH.dom}
\leanref{Formalism/SharpRestriction/GrowthBounds/StandardPolynomials/NotCostlyZero.lean}{Formalism.MainNotions.ColumnSystems.Pf.Col.PH.colOK}
\leanref{Formalism/SharpRestriction/GrowthBounds/StandardPolynomials/NotCostlyZero.lean}{Formalism.MainNotions.ColumnSystems.Pf.Col.PH.posc}

\begin{proof}
Write $h_i=h_i(c)$ for $\nu$, of rank $m$, and $H_t=H_t(c)$ for $\tilde\nu$, of rank $k$, and number the letters of $\mathcal P_j$ as
$0,\dots,2s$, so that the augmented letter $t\geq1$ is the letter $t-1$. By the definition of $\tilde\nu$, $H_0=k$ if $c\leq0$ and
$H_0=0$ if $c\geq1$, and $H_t=[c\leq\theta(t)]+h_{t-1}$ for $t\geq1$, where $[\cdot]$ is $1$ or $0$. By conclusion (1) of
Lemma~\ref{pd-kappa-defects-windows-lift-pairs-source-lift-tail-existence-columns-shape}, $\theta\geq0$ and $\theta$ is nondecreasing.
We first show $\nu_i(b)\leq\theta(i+1)$ for $0\leq i\leq2s+1$ and $1\leq b\leq m$, by induction on $i$: $\nu_0=0$; at a vector letter $i$,
$\nu_{i+1}(b)\leq\theta(i+1)\leq\theta(i+2)$ by conclusion (3) of that lemma; at a covector letter $i$,
$\nu_{i+1}(b)\leq\nu_i(b)\leq\theta(i+1)\leq\theta(i+2)$, as
$\nu_i$ is a strip over $\nu_{i+1}$.

(1) $\tilde\nu_0=0$, and $\tilde\nu_{t+1}$ is $\nu_t$, which is dominant, preceded by $\theta(t+1)\geq\nu_t(1)$.

(2) At the augmented letter $0$, a vector letter: if $c\leq0$, then $H_0=k=1+h_0=H_1$, as $\theta(1)\geq0\geq c$; if $c\geq1$, then $H_0=0$
and $H_1=[c\leq\theta(1)]\leq1$. At the augmented letter $i+1$: if $i$ is a covector letter, then $\theta(i+2)=\theta(i+1)$ by
conclusion (2) of Lemma~\ref{pd-kappa-defects-windows-lift-pairs-source-lift-tail-existence-columns-shape} and
$h_{i+1}\leq h_i\leq h_{i+1}+1$, so $H_{i+2}\leq H_{i+1}\leq H_{i+2}+1$. If $i$ is a vector letter, then $h_i\leq h_{i+1}\leq h_i+1$; if
$h_i<h_{i+1}$, some $\nu_{i+1}(b)\geq c$, so $c\leq\theta(i+1)\leq\theta(i+2)$ by conclusion (3) of that lemma, and $H_{i+2}=H_{i+1}+1$;
otherwise $H_{i+2}-H_{i+1}=[c\leq\theta(i+2)]-[c\leq\theta(i+1)]\in\{0,1\}$. These give (C1) and (C2), and $H_0=0$ for $c\geq1$ is (C3). For
(C4), let $c\leq0$. Then $H_0=k\geq2$, and $h_0=m$ drops by at most one at each letter, so for $t\geq1$, $h_{t-1}\geq m-(2s+1)\geq1$ by
conclusion (4) of that lemma, and $H_t=1+h_{t-1}\geq2>1$.

(3) Let $t'$ be costly for $c$. The augmented letter $0$ is not: $H_0=H_1$ if $c\leq0$, and $H_0=0\leq H_1$ if $c\geq1$. So $t'=i+1$. If
$i$ were a covector letter, then $H_{i+2}<H_{i+1}\leq1$, so $c\geq1$ by (C4), $H_{i+1}=1$ and $H_{i+2}=0$; as $\theta(i+2)=\theta(i+1)$,
this gives $c>\theta(i+1)$, $h_{i+1}=0$ and $h_i=1$, so some $\nu_i(b)\geq c>\theta(i+1)$, which is false. So $t'$ is a vector letter,
$x_{i'}$ or $x_+$, of position $s_j+d_{j,i'}$ or $\zeta=s_j+g_j$, and every augmented letter $t<t'$ has position $s_j$ or
$s_j+d_{j,i''}$ with $i''<i'$, respectively $i''\leq s$. As $0<d_{j,1}<\dots<d_{j,s}<g_j$, $\pi_t<\pi_{t'}$.
\end{proof}
% Lean: Col.PH is the conclusion of Lemma ..-columns-shape (ph_seg); muP nu ka is the augmented sequence, with colH_muP,
%   H0 and HS for the heights. The bound nu_i(b) <= theta(i + 1) is PH.ka_le (PH.nonneg); (1) is PH.dom; (2) is PH.colOK
%   (colH_muP_succ_ne_castSucc_iff_vec with PH.stepV, PH.stepC and PH.pos_le; one_lt_colH_muP_of_PH_of_lt_one with
%   PH.dual1 and PH.dualLow for (iv)); (3) is PH.posc with PH.not_costly_zero and PH.costly_succ, and posS_lt for the
%   positions (ColDef.lean).

\begin{lemma}\label{pd-kappa-defects-windows-lift-pairs-source-lift-tail-existence-columns-symbol}
In Situation~\ref{pd-kappa-defects-windows-lift-pairs-source-lift-tail-existence-situation}, let $\ell$, $j$, $\nu$ and $I$ be as
in Definition~\ref{pd-kappa-defects-windows-lift-pairs-source-lift-tail-existence-columns-data}, with its notation, suppose that the
conclusions (1) to (4) of Lemma~\ref{pd-kappa-defects-windows-lift-pairs-source-lift-tail-existence-columns-shape} hold, and let
$\Pi_\zeta$ be formed for the augmented letters, with their positions. Then for every $c$ in the range of
Definition~\ref{pd-kappa-defects-windows-lift-pairs-source-lift-tail-existence-columns-data}(3),
\[
 \phi\big(\Pi_\zeta(\mathrm{Col}^{\mathrm{pt}}_{\tilde\nu,c}(J_c))\big)=z^{\pi_c}\,\iota_j(K_c),
\]
with $\mathrm{Col}^{\mathrm{pt}}$ formed with the splitting $k=1+m$ (Definition~\ref{pd-symbols-path-generators}(5)).
\end{lemma}
\leanref{Formalism/SharpRestriction/GrowthBounds/StandardPolynomials/MonomialSumEqProd.lean}{Formalism.MainNotions.ColumnSystems.Pf.col_sym_eq}
\leanref{Formalism/SharpRestriction/GrowthBounds/StandardPolynomials/SymbolEq.lean}{Formalism.MainNotions.ColumnSystems.Pf.Col.symbol_eq}

\begin{proof}
\openstep{The identity is not yet written out: the plane factor of the column function of the augmented sequence at $c$ is the
product of the variables $z_{\lambda_t,1}^{P^S_t}$ over the $t$ with $\theta(t)<c\leq\theta(t+1)$, and its tangent factor is the image
under $\iota_j$ of the column function $K_c$ of $\nu$.}
% Lean: col_sym_eq (MonomialSumEqProd.lean): symbol_eq h O (colPsiS X O j) c (I'' c) (SymbolEq.lean), col_tan_eq
%   (Kcol_eq_KG, map_KG, KG_congr, phiS_succ, lev_succ; KGShift.lean), monomial_sum_eq_prod for z^{pi_c}.
% Size: about 480 lines (lean_size of col_sym_eq).
\end{proof}

\begin{lemma}\label{pd-kappa-defects-windows-lift-pairs-source-lift-tail-existence-columns-counts}
In Situation~\ref{pd-kappa-defects-windows-lift-pairs-source-lift-tail-existence-situation}, let $\ell$, $j$, $\nu$ and $I$ be as
in Definition~\ref{pd-kappa-defects-windows-lift-pairs-source-lift-tail-existence-columns-data}, with its notation, and let $C$ be the
set of the $c\in\mathbf Z$ with $|c|\leq T_\nu$ or $1\leq c\leq\theta(2s+2)$. Then $\prod_{c\in C}K_c=\Psi(\nu,I)$. If the conclusions (1)
to (4) of Lemma~\ref{pd-kappa-defects-windows-lift-pairs-source-lift-tail-existence-columns-shape} hold, then moreover
$\sum_{c\in C}d_c=\delta'_j$, $\sum_{c\in C}o_c=\sum_{i=1}^sP_i\delta'_j(x_i)+\delta'_j(x_+)$, and the exponent vector
$\sum_{c\in C}\pi_c$ is $p^{g_j}a_0$ at $z_{(s_j,1),1}$, $P_i(\hat b_i\dot-o'_i)$ at $z_{(y_{j,i},1),1}$ for $1\leq i\leq s$, and $0$ at
every other variable.
\end{lemma}
\leanref{Formalism/SharpRestriction/GrowthBounds/StandardPolynomials/ColXSum.lean}{Formalism.MainNotions.ColumnSystems.Pf.GenusGluing.Kprod_seg}
\leanref{Formalism/SharpRestriction/GrowthBounds/StandardPolynomials/ColXSum.lean}{Formalism.MainNotions.ColumnSystems.Pf.GenusGluing.dKsum_seg}
\leanref{Formalism/SharpRestriction/GrowthBounds/StandardPolynomials/ColXSum.lean}{Formalism.MainNotions.ColumnSystems.Pf.GenusGluing.plSum_seg}
\leanref{Formalism/SharpRestriction/GrowthBounds/StandardPolynomials/ColXSum.lean}{Formalism.MainNotions.ColumnSystems.Pf.GenusGluing.oSum_seg}

\begin{proof}
\openstep{The four counts over the columns are not yet written out: the factors $K_c$ with $c\notin[-T_\nu,T_\nu]$ are $1$; a letter
$\lambda$ is used for $\delta'_j(\lambda)$ values of $c$; the augmented letter $t$ contributes $z_{\lambda_t,1}^{P^S_t}$ to the
$\theta(t+1)-\theta(t)$ values $\theta(t)<c\leq\theta(t+1)$; and the costly letters give the orders.}
% Lean: Kprod_seg (the factors outside colRange are Kcol .. = 1, as I is a standard filling), dKsum_seg, plSum_seg and
%   oSum_seg (ColXSum.lean, with ColXSeg.lean, ColXCnt.lean, PH.oc_eq and PH.costly_succ).
% Size: about 500 lines (lean_size of Kprod_seg, dKsum_seg, plSum_seg and oSum_seg together).
\end{proof}

\begin{lemma}\label{pd-kappa-defects-windows-lift-pairs-source-lift-tail-existence-columns-degrees}
In Situation~\ref{pd-kappa-defects-windows-lift-pairs-source-lift-tail-existence-situation}, let $\ell$, $j$, $\nu$ and $I$ be as
in Definition~\ref{pd-kappa-defects-windows-lift-pairs-source-lift-tail-existence-columns-data}, with its notation, and let $c$ be in the
range of its part (3). Then:
\begin{enumerate}
\item $K_c$ has coefficients in $\mathcal R$ and row degrees $d_c$; $f_c$ has row degrees $e_c$; the row degrees of $\pi_c$ and the values
  of $e_c$ are at most $Q$; and $e_c$ vanishes at the letters $(s_{j'},0)$, $1\leq j'\leq q$;
\item if the conclusions (1) to (4) of Lemma~\ref{pd-kappa-defects-windows-lift-pairs-source-lift-tail-existence-columns-shape} hold,
  then $z^{\pi_c}\iota_j(K_c)$ has row degrees $e_c$, and each of its monomials has $\varpi_0$-degree $o_c$.
\end{enumerate}
\end{lemma}
\leanref{Formalism/SharpRestriction/GrowthBounds/StandardPolynomials/HomGMonomial.lean}{Formalism.MainNotions.ColumnSystems.Pf.ColF.colKInt}
\leanref{Formalism/SharpRestriction/GrowthBounds/StandardPolynomials/HomGMonomial.lean}{Formalism.MainNotions.ColumnSystems.Pf.ColF.colKPol}
\leanref{Formalism/SharpRestriction/GrowthBounds/StandardPolynomials/HomGMonomial.lean}{Formalism.MainNotions.ColumnSystems.Pf.ColF.colFPol}
\leanref{Formalism/SharpRestriction/GrowthBounds/StandardPolynomials/HomGMonomial.lean}{Formalism.MainNotions.ColumnSystems.Pf.ColF.colPlQ}
\leanref{Formalism/SharpRestriction/GrowthBounds/StandardPolynomials/HomGMonomial.lean}{Formalism.MainNotions.ColumnSystems.Pf.ColF.colDgQ}
\leanref{Formalism/SharpRestriction/GrowthBounds/StandardPolynomials/HomGMonomial.lean}{Formalism.MainNotions.ColumnSystems.Pf.ColF.colDgY}
\leanref{Formalism/SharpRestriction/GrowthBounds/StandardPolynomials/IoDChain.lean}{Formalism.MainNotions.ColumnSystems.Pf.ColF.colHomC}
\leanref{Formalism/SharpRestriction/GrowthBounds/StandardPolynomials/IoDChain.lean}{Formalism.MainNotions.ColumnSystems.Pf.ColF.colPolC}

\begin{proof}
\openstep{The degrees of the column data are not yet written out: the row degrees of $K_c$, of $f_c$ and of
$z^{\pi_c}\iota_j(K_c)$, the bounds by $Q$, and the $\varpi_0$-degree $o_c$ of the monomials of $z^{\pi_c}\iota_j(K_c)$.}
% Lean: colKInt, colKPol (kcol_pol), colFPol, colPlQ, colDgQ (sum_lam_le, dgS_le) and colDgY (HomGMonomial.lean) for (1);
%   colPolC and colHomC (IoDChain.lean, with h : Col.PH) for (2). Only colHomC and colPolC use the conclusions of Lemma
%   ..-columns-shape; colFPol, colPlQ and colDgQ use 2 <= q, p > 0 and g_j <= G'.
% Size: about 1,450 lines (lean_size of the eight lemmas together; colHomC, colPolC and colFPol share most of it).
\end{proof}

\begin{lemma}\label{pd-kappa-defects-windows-lift-pairs-source-lift-tail-existence-columns-clauses}
In Situation~\ref{pd-kappa-defects-windows-lift-pairs-source-lift-tail-existence-situation}, let $\ell$, $j$, $\nu$ and $I$ be as
in Definition~\ref{pd-kappa-defects-windows-lift-pairs-source-lift-tail-existence-columns-data}, with its notation. Suppose that
the conclusions (1) to (4) of Lemma~\ref{pd-kappa-defects-windows-lift-pairs-source-lift-tail-existence-columns-shape} hold, and
that $a_0+\sum_{i\leq\eta}\varsigma_i\leq\hat a$. Then the set of the $c\in\mathbf Z$ with $|c|\leq T_\nu$ or $1\leq c\leq\theta(2s+2)$,
the integer $a_0$, and the $K_c$, $d_c$, $z^{\pi_c}$, $f_c$, $o_c$ and $e_c$ form a column system for $\Psi(\nu,I)$ at $j$
(Definition~\ref{pd-kappa-defects-windows-lift-pairs-source-lift-tail-existence-column-systems}).
\end{lemma}
\leanref{Formalism/SharpRestriction/GrowthBounds/StandardPolynomials/ColXSum.lean}{Formalism.MainNotions.ColumnSystems.Pf.GenusGluing.dKsum_seg}
\leanref{Formalism/SharpRestriction/GrowthBounds/StandardPolynomials/ColXSum.lean}{Formalism.MainNotions.ColumnSystems.Pf.GenusGluing.Kprod_seg}
\leanref{Formalism/SharpRestriction/GrowthBounds/StandardPolynomials/ColXSum.lean}{Formalism.MainNotions.ColumnSystems.Pf.GenusGluing.plSum_seg}
\leanref{Formalism/SharpRestriction/GrowthBounds/StandardPolynomials/ColXSum.lean}{Formalism.MainNotions.ColumnSystems.Pf.GenusGluing.oSum_seg}
\leanref{Formalism/SharpRestriction/GrowthBounds/StandardPolynomials/MonomialSumEqProd.lean}{Formalism.MainNotions.ColumnSystems.Pf.colFac}
\leanref{Formalism/SharpRestriction/GrowthBounds/StandardPolynomials/IoDChain.lean}{Formalism.MainNotions.ColumnSystems.Pf.ColF.colHomC}
\leanref{Formalism/SharpRestriction/GrowthBounds/StandardPolynomials/IoDChain.lean}{Formalism.MainNotions.ColumnSystems.Pf.ColF.colPolC}
\leanref{Formalism/SharpRestriction/GrowthBounds/StandardPolynomials/HomGMonomial.lean}{Formalism.MainNotions.ColumnSystems.Pf.ColF.colFPol}
\leanref{Formalism/SharpRestriction/GrowthBounds/StandardPolynomials/HomGMonomial.lean}{Formalism.MainNotions.ColumnSystems.Pf.ColF.colKInt}
\leanref{Formalism/SharpRestriction/GrowthBounds/StandardPolynomials/HomGMonomial.lean}{Formalism.MainNotions.ColumnSystems.Pf.ColF.colKPol}
\leanref{Formalism/SharpRestriction/GrowthBounds/StandardPolynomials/HomGMonomial.lean}{Formalism.MainNotions.ColumnSystems.Pf.ColF.colPlQ}
\leanref{Formalism/SharpRestriction/GrowthBounds/StandardPolynomials/HomGMonomial.lean}{Formalism.MainNotions.ColumnSystems.Pf.ColF.colDgQ}
\leanref{Formalism/SharpRestriction/GrowthBounds/StandardPolynomials/HomGMonomial.lean}{Formalism.MainNotions.ColumnSystems.Pf.ColF.colDgY}
\leanref{Formalism/SharpRestriction/GrowthBounds/StandardPolynomials/ColXRen.lean}{Formalism.MainNotions.ColumnSystems.Pf.GenusGluing.liftSubst_rename}
\leanref{Formalism/SharpRestriction/GrowthBounds/StandardPolynomials/ColXRen.lean}{Formalism.MainNotions.ColumnSystems.Pf.GenusGluing.inWjRg_rename}

\begin{proof}
Let $C$ be the set of the $c$ of the statement. Give the augmented letters their types and positions $\pi_t$
(Definition~\ref{pd-kappa-defects-windows-lift-pairs-source-lift-tail-existence-columns-data}(1)), put $\omega(t)=P^S_t=p^{\zeta-\pi_t}$,
and use the splitting $k=1+m$; so $o_c=o_\omega(\tilde\nu,c)$. The positions satisfy $\pi_0\leq\pi_1\leq\dots\leq\pi_{2s+1}=\zeta$, as
$0<d_{j,1}<\dots<d_{j,s}<g_j$.

\emph{Step 1: the augmented sequence.} By
Lemma~\ref{pd-kappa-defects-windows-lift-pairs-source-lift-tail-existence-columns-augmented}, the $\tilde\nu_t$ are dominant, $\tilde\nu$
satisfies the column conditions for every $c$, and $\pi_t<\pi_{t'}$ whenever $t<t'$, $t$ is used for $c$ and $t'$ is costly for $c$.
So Lemma~\ref{pd-kappa-defects-windows-lift-pairs-source-lift-tail-existence-columns-lifted}(2) to (4) applies to $\tilde\nu$, with
$N=2s+2$, $r=1$, $K=m$ and $n=k$; the setting is unramified quadratic, so $\varpi_0$ is a uniformizer of $\mathcal R$.

\emph{Step 2: renaming.} The map $\phi$ sends the variables of the augmented letter $t$ to those of $\lambda_t$. The letters $\lambda_t$
are distinct, as $q\geq2$ gives $s_{j+1}\neq s_j$; $\lambda_t$ has the type of $t$; and its position is the residue of $\pi_t$ modulo
$n'$. As $\sigma^{n'}=1$, $g_\pi$ depends only on $\pi$ modulo $n'$, so $\Sigma_g\circ\phi=\phi\circ\Sigma^{\mathrm L}_g$ for
$g\in GL_k(W)$, with $\Sigma^{\mathrm L}_g$ formed for the augmented letters (Definition~\ref{pd-symbols-path-generators}(7)). As $\phi$
maps distinct monomials to distinct monomials, $\phi(F)\in\varpi_0^i\mathcal R[z]$ if and only if $F\in\varpi_0^i\mathcal R[z]$.

\emph{Step 3: the conditions.} Let $c\in C$ and $\sigma_c=z^{\pi_c}\iota_j(K_c)$.
\begin{itemize}
\item (1) $K_c$ has coefficients in $\mathcal R$ and row degrees $d_c$ by
  Lemma~\ref{pd-kappa-defects-windows-lift-pairs-source-lift-tail-existence-columns-degrees}(1), and $d_c\leq1$ by definition. By
  Lemma~\ref{pd-kappa-defects-windows-lift-pairs-source-lift-tail-existence-columns-counts}, $\sum_cd_c=\delta'_j$ and
  $\prod_cK_c=\Psi(\nu,I)$.
\item (2) By Lemma~\ref{pd-kappa-defects-windows-lift-pairs-source-lift-tail-existence-columns-symbol},
  $\sigma_c=\phi(\Pi_\zeta(\mathrm{Col}^{\mathrm{pt}}_{\tilde\nu,c}(J_c)))$, which lies in $\mathcal R[z]$, as $\Pi_\zeta$ maps $\mathcal R[z]$
  into itself and the column functions have integer coefficients. The polynomial $f_c=\phi(\mathrm{Col}^1_{\tilde\nu,c,\zeta}(J_c))$ lies in
  $\mathcal R[z]$ by Lemma~\ref{pd-kappa-defects-windows-lift-pairs-source-lift-tail-existence-columns-lifted}(2) and Step 2. For
  $g\in GL_k(W)$, Step 2 and part (3) of that lemma, with $g'=g$ and the matrix $1$ in the column function, give
  $\Sigma_gf_c=\phi(\Sigma^{\mathrm L}_g(\mathrm{Col}^1_{\tilde\nu,c,\zeta}(J_c)))\in\varpi_0^{o_c}\mathcal R[z]$. By part (4) and Step 2,
  $\Sigma_1f_c-\varpi_0^{o_c}\sigma_c=\phi\big(\Sigma^{\mathrm L}_1(\mathrm{Col}^1_{\tilde\nu,c,\zeta}(J_c))-\varpi_0^{o_c}\Pi_\zeta(\mathrm{Col}^{\mathrm{pt}}_{\tilde\nu,c}(J_c))\big)
  \in\varpi_0^{o_c+1}\mathcal R[z]$. So $f_c$ is a lifted factor of $\sigma_c$ of order $o_c$. By
  Lemma~\ref{pd-kappa-defects-windows-lift-pairs-source-lift-tail-existence-columns-degrees}, every monomial of $\sigma_c$ has
  $\varpi_0$-degree $o_c$, and $f_c$ and $\sigma_c$ have row degrees $e_c$.
\item (3) This is Lemma~\ref{pd-kappa-defects-windows-lift-pairs-source-lift-tail-existence-columns-degrees}(1).
\item (4) The bound on $a_0$ is assumed, and the sums of the $\pi_c$ and of the $o_c$ are those of
  Lemma~\ref{pd-kappa-defects-windows-lift-pairs-source-lift-tail-existence-columns-counts}.
\end{itemize}
% Lean: the fields of the ColSys in colsAt_of (ColTop.lean), with h = ph_seg (Lemma ..-columns-shape): KInt, KPol, fPol, plQ,
%   dgCQ, dgCY, homC and polC are colKInt, colKPol, colFPol, colPlQ, colDgQ, colDgY, colHomC and colPolC (Lemma
%   ..-columns-degrees); digit_le_one is dKS <= 1; dKsum, Kprod, plSum and oSum are dKsum_seg, Kprod_seg, plSum_seg and
%   oSum_seg (Lemma ..-columns-counts); facC is colFac (MonomialSumEqProd.lean): hOK = h.colOK and hposc = h.posc (Lemma
%   ..-columns-augmented, Step 1), col_sym_eq (Lemma ..-columns-symbol), column_kernel_components_integral with inA_rename,
%   lifted_column_order with inWjRg_rename (lamS_injective, pos_congr), lifted_column_congruence with SgVanish_of and
%   plTanCol_eq_SgW (h.dom) and liftSubst_rename (Lemma ..-columns-lifted and Step 2), and col_powSubst_int for sigma_c.
%   colFac applies lifted_column_order and lifted_column_congruence to the augmented sequence through ColOK, with only
%   2 s + 1 < m; Lemma pd-symbols-path-clauses(5), which assumes a path with N + r <= K, is not used here.
\end{proof}

\begin{lemma}\label{pd-kappa-defects-windows-lift-pairs-source-lift-tail-existence-columns}
For every $C_c\geq0$ there is $T_0>0$ such that for every $C_0\geq T_0$ there is $T_1>0$ such that for all
$C_{E1},C_{E2}\geq T_1$ there is $T_2>0$ with the following property for every $C_H\geq T_2$. In
Situation~\ref{pd-kappa-defects-windows-lift-pairs-source-lift-tail-existence-situation} with these constants, let $\ell$ be a
level choice for $\beta$ and $1\leq j\leq q$, let $\nu$ be a path of rank $m$ for $\delta'_j$ with $\nu_{2u_l-1}(1)\leq M_l$ for
every $l$ and $\nu_{2s+1}(1)\leq M_+$, and let $I$ be a standard filling of $\nu$. Then $\Psi(\nu,I)\in\mathcal P'_j$, formed as in
Definition~\ref{pd-kappa-defects-windows-lift-pairs-source-lift-tail-systems}(4), has a column system at $j$.
\end{lemma}
\leanref{Formalism/SharpRestriction/GrowthBounds/StandardPolynomials/ColsDef.lean}{Formalism.MainNotions.ColumnSystems.Pf.ColsAt}
\leanref{Formalism/SharpRestriction/GrowthBounds/StandardPolynomials/ColsDef.lean}{Formalism.MainNotions.ColumnSystems.Pf.RestCols}
\leanref{Formalism/SharpRestriction/GrowthBounds/StandardPolynomials/ColTop.lean}{Formalism.MainNotions.ColumnSystems.Pf.colsAt_of}
\leanref{Formalism/SharpRestriction/GrowthBounds/StandardPolynomials/ColTop.lean}{Formalism.MainNotions.ColumnSystems.Pf.restCols}

\begin{proof}
\emph{The constants.} Let $C_c\geq0$, and let $T_0$, $T_1$ and $T_2$ be the thresholds of
Lemma~\ref{pd-kappa-defects-windows-lift-pairs-source-lift-tail-existence-columns-numerics}. Let the constants and the data be as
in the statement, with $C_0\geq T_0$, $C_{E1},C_{E2}\geq T_1$ and $C_H\geq T_2$, and use the notation of
Definition~\ref{pd-kappa-defects-windows-lift-pairs-source-lift-tail-existence-columns-data}.

\emph{Step 1: the shape.} By Lemma~\ref{pd-kappa-defects-windows-lift-pairs-source-lift-tail-existence-columns-numerics}, the
two inequalities there hold for $\beta_j$, and
Lemma~\ref{pd-kappa-defects-windows-lift-pairs-source-lift-tail-existence-columns-shape} gives its conclusions (1) to (4).

\emph{Step 2: the integer $a_0$.} By Step 1, $a_0\geq0$. As $\mu\in\{0,12c\}$ and $\kappa_1\geq0$, $\mu'\geq0$, so
$a_0+\sum_{i\leq\eta}\varsigma_i=\lfloor\hat a-\mu'\rfloor\leq\hat a$, as $\hat a$ is an integer.

\emph{Step 3.} By Steps 1 and 2,
Lemma~\ref{pd-kappa-defects-windows-lift-pairs-source-lift-tail-existence-columns-clauses} applies, and the data of
Definition~\ref{pd-kappa-defects-windows-lift-pairs-source-lift-tail-existence-columns-data}(3), with $a_0$, form a column
system for $\Psi(\nu,I)$ at $j$.
\end{proof}
% Lean: ColsAt C0 CYb Cde Ck CD Cb CE1 CE2 Cc in its TransportedAlt form RestCols, proved as restCols =
%   TransportedAlt.mono colsAt_of column_numerics_holds, with the thresholds of ColumnNumericsBeyond. colsAt_of (ColTop.lean):
%   hNum gives hnum1 and hnum2 (Step 1 with Lemma ..-columns-numerics), h = ph_seg (Lemma ..-columns-shape), and the ColSys
%   has I = colsS, Kc = Kcol, dK = dKS, pl = plS, fC = fCS, oC = oCS, dgC = dgCS and a0 = toNat (a0Z), with a0le from
%   hnum1 and Int.floor_le_floor (Step 2, column_kappa_nonneg for mu' >= 0) and the other fields as in Lemma
%   ..-columns-clauses (Step 3).
% Size of the parts left open: Lemmas pd-kappa-defects-windows-lift-pairs-source-lift-tail-existence-columns-symbol
%   (about 480 lines), -columns-counts (about 500) and -columns-degrees (about 1,450), in the proof of Lemma
%   -columns-clauses; Lemma -columns-numerics is written out.

\begin{definition}\label{pd-kappa-defects-windows-lift-pairs-source-lift-tail-existence-prefactorizations}
In Situation~\ref{pd-kappa-defects-windows-lift-pairs-source-lift-tail-existence-situation}, let $\ell$ be a level choice for
$\beta$, with base monomial $z^{A_\ell}$ (Definition~\ref{pd-kappa-defects-windows-lift-pairs-source-lift-tail-existence-levels}).
\begin{enumerate}
\item Let $G=(G_j)$ with $G_j\in\mathcal P'_j$. A \emph{prefactorization} of $G$ for $\ell$ consists of: for each $j$, a finite set
  $I_j$ and, for $a\in I_j$, a polynomial $K_{j,a}\in\mathcal P'_j$, a map $d_{j,a}$, a monomial $z^{\pi_{j,a}}$, a polynomial
  $f_{j,a}\in\mathcal K[z]$, an integer $o_{j,a}\geq0$ and a map $e_{j,a}$, as in
  Definition~\ref{pd-kappa-defects-windows-lift-pairs-source-lift-tail-existence-column-systems}; for each $j$ and
  $1\leq l\leq S_j-\eta_j$, a polynomial $f^H_{j,l}\in\mathcal K[z]$, an integer $o^H_{j,l}\geq0$ and a map $e^H_{j,l}$ from the
  letters to $\mathbf Z_{\geq0}$; a finite set $I_M$ and, for $a\in I_M$, an integer $m_a\geq0$, a polynomial
  $f^M_a\in\mathcal K[z]$, a monomial $z^{\kappa_a}$, an integer $o^M_a\geq0$ and a map $e^M_a$; and a polynomial
  $f_R\in\mathcal K[z]$, a monomial $z^\rho$, an integer $o_R\geq0$ and a map $e_R$; such that, with
  $\sigma_{j,a}=z^{\pi_{j,a}}\iota_j(K_{j,a})$, $\sigma^H_{j,l}=\iota_j(\hat E_l)$ and
  $\sigma_R=\prod_j\iota_j(\hat E_+)^{\ell_{j,+}}\cdot z^\rho$, where $\hat E_l$ and $\hat E_+$ are formed from $\beta_j$:
  \begin{enumerate}
  \item for every $j$, the $K_{j,a}$ and $d_{j,a}$ satisfy condition (1) of
    Definition~\ref{pd-kappa-defects-windows-lift-pairs-source-lift-tail-existence-column-systems} for $G_j$; each $f_{j,a}$ is a
    lifted factor of $\sigma_{j,a}$ of order $o_{j,a}$, every monomial of $\sigma_{j,a}$ has $\varpi_0$-degree $o_{j,a}$, and
    $\Sigma_1f_{j,a}$ and $\sigma_{j,a}$ have row degrees $e_{j,a}$; and the row degrees of $\pi_{j,a}$ are at most $Q$, and $0$
    at the letters $(s_{j'},0)$, $1\leq j'\leq q$;
  \item each $f^H_{j,l}$ is a lifted factor of $\sigma^H_{j,l}$ of order $o^H_{j,l}$, every monomial of $\sigma^H_{j,l}$ has
    $\varpi_0$-degree $o^H_{j,l}$, and $\Sigma_1f^H_{j,l}$ and $\sigma^H_{j,l}$ have row degrees $e^H_{j,l}$;
  \item each $f^M_a$ is a lifted factor of $z^{\kappa_a}$ of order $o^M_a=d_\varpi(\kappa_a)$, and $\Sigma_1f^M_a$ and
    $z^{\kappa_a}$ have row degrees $e^M_a$;
  \item $f_R$ and $\sigma_R$ lie in $\mathcal R[z]$, $\Sigma_1f_R=\varpi_0^{o_R}\sigma_R$, every monomial of $\sigma_R$ has
    $\varpi_0$-degree $o_R$, and $\sigma_R$ has row degrees $e_R$;
  \item the values of the $e_{j,a}$, $e^H_{j,l}$ and $e^M_a$ are at most $Q$, and these maps vanish at the letters $(s_{j'},0)$;
  \item $e_R+\sum_{j,a}e_{j,a}+\sum_{j,l}\ell_{j,l}e^H_{j,l}+\sum_am_ae^M_a$ is the map $x\mapsto d_x$ of the row degrees of
    $\mathrm{Pol}_T$; $A_\ell=\rho+\sum_{j,a}\pi_{j,a}+\sum_am_a\kappa_a$; and
    $o_R+\sum_{j,a}o_{j,a}+\sum_{j,l}\ell_{j,l}o^H_{j,l}+\sum_am_ao^M_a=o$;
  \item the product $F=f_R\prod_{j,a}f_{j,a}\prod_{j,l}(f^H_{j,l})^{\ell_{j,l}}\prod_a(f^M_a)^{m_a}$ lies in $\mathrm{Pol}_T$, and
    $\Sigma_gF\in\varpi_0^{o+c}\mathcal R[z]$ for every $g\in GL_k(W)$ whose plane differs from that of $1$
    (Lemma~\ref{pd-symbols-edge-basis}, with $r=1$).
  \end{enumerate}
  The bounds on the substitutions $\Sigma_g$ of the lifted factors in (a) to (c) are not used: the proof of
  Lemma~\ref{pd-kappa-defects-windows-lift-pairs-source-lift-tail-existence-assembly} passes these lifted factors to
  condition (a) of a factorization system, where they are not used either.
\item The level choice $\ell$ satisfies the \emph{tail condition} if the following holds for every integer $e$ with
  $1\leq e<n$. Let $u=(u_j)$ with $u_j\in\mathbf Z_{\geq0}^{S_j-\eta_j}$ and $u_{j,l}\leq\ell_{j,l}$ for all $l$. Let $G'=(G'_j)$
  with $G'_j\in\mathcal P'_j$ having coefficients in $\mathcal R$ and row degrees $d'_j\leq\delta'_j$, and lying in
  $Y(\mathcal L_l,M_l+N_1e)$ for every $l$ and in $Y(\mathcal L_+,M_++N_1e)$, with the $\mathcal L_l$, $M_l$, $\mathcal L_+$ and
  $M_+$ of the generators of $\beta_j$. Let $z^{A'}$ be a monomial with
  \[
   \sum_{c=1}^k\big(A_\ell(x,c)\dot-A'(x,c)\big)\leq Qe\quad\text{and}\quad\sum_{c=1}^k\big(A'(x,c)\dot-A_\ell(x,c)\big)\leq Qe
  \]
  for every letter $x$, and with $A'(x,c)=A_\ell(x,c)$ for the letters $x=(s_{j'},0)$. If
  $P=\prod_jH_j(u_j,G'_j)\cdot z^{A'}$ lies in $\mathrm{Pol}_T$ and every monomial of $P$ has $\varpi_0$-degree at most $o+e$,
  then $P\in S_e$.
\end{enumerate}
\end{definition}
\leanref{Formalism/SharpRestriction/GrowthBounds/StandardPolynomials/PreSys.lean}{Formalism.MainNotions.ColumnSystems.Pf.PreSys}
\leanref{Formalism/SharpRestriction/GrowthBounds/StandardPolynomials/MonomialSumEqProdTensor.lean}{Formalism.MainNotions.ColumnSystems.Pf.PreSys.sgR}
\leanref{Formalism/SharpRestriction/GrowthBounds/StandardPolynomials/MonomialSumEqProdTensor.lean}{Formalism.MainNotions.ColumnSystems.Pf.TCoreBody}
% Lean: (1) is PreSys C0 CE1 CE2 X O w (mm + n ka1) beta L G Std, with Std = A_ell: iotaC, Kc, dK, pl, fC, oC and dgC;
%   fH, oH and dgH; iotaM, multM, fM, aM, oM and dgM; fR, R, oR and dR. (a) is KInt, KPol, digit_le_one, dKsum, Kprod,
%   facC, homC, liftC, polC, plQ and plY; (b) is facH, homH, liftH and polH; (c) is facM, homM, liftM and polM; (d) is
%   rigA, rigLift, rigHom, rigPol and rigInt, sigma_R being PreSys.sgR; (e) is dgCQ, dgHQ, dgMQ, dgCY, dgHY and dgMY; (f)
%   is degEq, stdEq and ordEq; (g) is total and clean. PreSys.toFacSys passes facC, facH and facM to FacSys.fac, which
%   FacSys.liftable does not read. (2) is TCoreBody C0 Cde CE1 CE2 X O mm n beta L Std, with ju, Gu, du and A for u, G',
%   d'_j and A'; the conditions of the generators are T.J k and T.thr k of BDk C0 CE1 CE2 (mm + n ka1) X beta j.

\begin{definition}\label{pd-kappa-defects-windows-lift-pairs-source-lift-tail-existence-prefactorization-keys}
In Situation~\ref{pd-kappa-defects-windows-lift-pairs-source-lift-tail-existence-situation}, write the indices $j$ of the
$s_j$ modulo $q$.
\begin{enumerate}
\item For $x\in T$ and an integer $P\geq1$ put
  \[
   h_x=\sum_{c=1}^kz_{(x,1),c}z_{(x,0),c},\qquad\hat E_{x,P}=\sum_{c=1}^kz_{(x,1),c}^Pz_{(x,0),c}^P,
  \]
  and let $\mathbf 1_x$ be the map that sends the letters $(x,1)$ and $(x,0)$ to $1$ and the other letters to $0$.
\item For $1\leq j\leq q$, $1\leq i\leq S_j$ and $2\leq t\leq k$, with $d=d_{j,i}$, put
  \[
   \kappa^S_{j,i}=\sum_{c=1}^kz_{(s_j,1),c}^{p^d}z_{(y_{j,i},0),c},\qquad
   \kappa^W_{j,i,t}=z_{(y_{j,i},1),t}z_{(s_j,1),1}^{p^d}-z_{(y_{j,i},1),1}z_{(s_j,1),t}^{p^d}.
  \]
\item For $2\leq t\leq k$ and $1\leq j_0\leq q$ put
  \[
   D_t=\prod_{j=1}^qz_{(s_j,1),1}^{p^{g_j}-1}-\prod_{j=1}^qz_{(s_j,1),t}^{p^{g_j}-1},\qquad
   L_{j_0}=z_{(s_{j_0},1),1}\prod_{t=2}^kD_t.
  \]
\end{enumerate}
\end{definition}
\leanref{Formalism/SharpRestriction/GrowthBounds/StandardPolynomials/TopKeyFormSelfSoft.lean}{Formalism.MainNotions.ColumnSystems.Pf.hTop}
\leanref{Formalism/SharpRestriction/GrowthBounds/CoarseMultiplicity/HookKeyOrderGeOne.lean}{Formalism.MainNotions.ColumnSystems.Pf.hKey}
\leanref{Formalism/SharpRestriction/GrowthBounds/StandardPolynomials/TopKeyFormSelfSoft.lean}{Formalism.MainNotions.ColumnSystems.Pf.EhatT}
\leanref{Formalism/SharpRestriction/GrowthBounds/StandardPolynomials/TopKeyFormSelfSoft.lean}{Formalism.MainNotions.ColumnSystems.Pf.dgH}
\leanref{Formalism/SharpRestriction/GrowthBounds/CoarseMultiplicity/RestRawTail1Of.lean}{Formalism.MainNotions.ColumnSystems.Pf.kSkey}
\leanref{Formalism/SharpRestriction/GrowthBounds/CoarseMultiplicity/RestRawTail1Of.lean}{Formalism.MainNotions.ColumnSystems.Pf.kWkey}
\leanref{Formalism/SharpRestriction/GrowthBounds/CoarseMultiplicity/RestRawTail1Of.lean}{Formalism.MainNotions.ColumnSystems.Pf.locC}
\leanref{Formalism/SharpRestriction/GrowthBounds/CoarseMultiplicity/RestRawTail1Of.lean}{Formalism.MainNotions.ColumnSystems.Pf.locR}
% Lean: h_x is hTop O X x = hKey O (m + 1) (x, true) (x, false), E_{x,P} is EhatT O X x P and 1_x is dgH x. kappa^S_{j,i}
%   is kSkey O X j i, kappa^W_{j,i,t} is kWkey O X j i t with t = c.succ (columns from 0), D_t is locC O X t and L_{j0} is
%   locR O X j0. topStart X j is s_j, topOdd X j i is y_{j,i} and topEnd X j is s_{j+1}.

\begin{lemma}\label{pd-kappa-defects-windows-lift-pairs-source-lift-tail-existence-prefactorization-products}
In Situation~\ref{pd-kappa-defects-windows-lift-pairs-source-lift-tail-existence-situation}, let $f,f',\sigma,\sigma'\in\mathcal
K[z]$ and let $o,o'\geq0$ be integers. Use the lifted factors of
Definition~\ref{pd-kappa-defects-windows-lift-pairs-source-lift-tail-existence-levels}(1).
\begin{enumerate}
\item If $f$ is a lifted factor of $\sigma$ of order $o$ and $f'$ one of $\sigma'$ of order $o'$, then $ff'$ is a lifted factor
  of $\sigma\sigma'$ of order $o+o'$, and $f^u$ is one of $\sigma^u$ of order $uo$ for every integer $u\geq0$.
\item Let $e\geq1$ be an integer and $r\in\mathcal R[z]$. If $f,\sigma\in\mathcal R[z]$, $\Sigma_gf\in\varpi_0^o\mathcal R[z]$
  for every $g\in GL_k(W)$, and $\Sigma_1f=\varpi_0^o(\sigma+\varpi_0^er)$, then $f$ is a lifted factor of $\sigma$ of order
  $o$.
\item If $f$ is a lifted factor of $\sigma$ of order $o$, $\sigma'\in\mathcal R[z]$ and $\sigma-\sigma'\in\varpi_0\mathcal
  R[z]$, then $f$ is a lifted factor of $\sigma'$ of order $o$.
\end{enumerate}
\end{lemma}
\leanref{Formalism/SharpRestriction/GrowthBounds/CoarseMultiplicity/RestRawTail1Of.lean}{Formalism.MainNotions.ColumnSystems.Pf.IsFac.mul}
\leanref{Formalism/SharpRestriction/GrowthBounds/StandardPolynomials/IdentityLiftWeightedHomog.lean}{Formalism.MainNotions.ColumnSystems.Pf.IsFac.one}
\leanref{Formalism/SharpRestriction/GrowthBounds/StandardPolynomials/IdentityLiftWeightedHomog.lean}{Formalism.MainNotions.ColumnSystems.Pf.IsFac.pow}
\leanref{Formalism/SharpRestriction/GrowthBounds/StandardPolynomials/IdentityLiftWeightedHomog.lean}{Formalism.MainNotions.ColumnSystems.Pf.IsFac.of_exact}
\leanref{Formalism/SharpRestriction/GrowthBounds/StandardPolynomials/IdentityLiftWeightedHomog.lean}{Formalism.MainNotions.ColumnSystems.Pf.IsFac.congr_sg}

\begin{proof}
Write $\varpi=\varpi_0$. (1) The polynomials $ff'$ and $\sigma\sigma'$ lie in $\mathcal R[z]$. As $\Sigma_g$ is a ring
homomorphism, $\Sigma_g(ff')=\Sigma_gf\cdot\Sigma_gf'\in\varpi^{o+o'}\mathcal R[z]$ for every $g$. Next
\[
 \Sigma_1(ff')-\varpi^{o+o'}\sigma\sigma'=(\Sigma_1f-\varpi^o\sigma)\,\Sigma_1f'+\varpi^o\sigma\,(\Sigma_1f'-\varpi^{o'}\sigma').
\]
The first term lies in $\varpi^{o+1}\varpi^{o'}\mathcal R[z]$, as $\Sigma_1f'\in\varpi^{o'}\mathcal R[z]$, and the second in
$\varpi^o\varpi^{o'+1}\mathcal R[z]$, as $\sigma\in\mathcal R[z]$. The polynomial $1$ is a lifted factor of $1$ of order $0$, as
$\Sigma_g1=1$; induction on $u$ gives the powers.

(2) Here $\Sigma_1f-\varpi^o\sigma=\varpi^{o+e}r\in\varpi^{o+1}\mathcal R[z]$, as $e\geq1$.

(3) $\Sigma_1f-\varpi^o\sigma'=(\Sigma_1f-\varpi^o\sigma)+\varpi^o(\sigma-\sigma')$, and both terms lie in
$\varpi^{o+1}\mathcal R[z]$.
\end{proof}
% Lean: (1) is IsFac.mul (inWj_mul for inA, ord and sgInt; sym through the displayed identity, with h'.ord at g = 1) and
%   IsFac.pow (IsFac.one, induction); (2) is IsFac.of_exact (wj_smul_pow and wj_mono); (3) is IsFac.congr_sg.

\begin{lemma}\label{pd-kappa-defects-windows-lift-pairs-source-lift-tail-existence-prefactorization-hooks}
In Situation~\ref{pd-kappa-defects-windows-lift-pairs-source-lift-tail-existence-situation}, use the notation of
Definition~\ref{pd-kappa-defects-windows-lift-pairs-source-lift-tail-existence-prefactorization-keys} and the
$\varpi_0$-degree of Definition~\ref{pd-kappa-defects-windows-lift-pairs-source-lift-tail-systems}(1), and let $x\in T$.
\begin{enumerate}
\item $\Sigma_gh_x=\varpi_0h_x$ for every $g\in GL_k(W)$, and every monomial of $h_x$ has $\varpi_0$-degree $1$.
\item Let $e\geq0$ be an integer and $P=p^e$. Then $h_x^P$ is a lifted factor of $\hat E_{x,P}$ of order $P$
  (Definition~\ref{pd-kappa-defects-windows-lift-pairs-source-lift-tail-existence-levels}(1)), every monomial of
  $\hat E_{x,P}$ has $\varpi_0$-degree $P$, and $\Sigma_1(h_x^P)$ and $\hat E_{x,P}$ have row degrees $P\mathbf 1_x$.
\item $h_x\neq0$, and $\hat E_{x,P}\neq0$ for every integer $P\geq1$.
\end{enumerate}
\end{lemma}
\leanref{Formalism/SharpRestriction/GrowthBounds/CoarseMultiplicity/HookKeyOrderGeOne.lean}{Formalism.MainNotions.ColumnSystems.Pf.hook_key_lift_exact}
\leanref{Formalism/SharpRestriction/GrowthBounds/CoarseMultiplicity/HookKeyOrderGeOne.lean}{Formalism.MainNotions.ColumnSystems.Pf.hook_key_order_ge_one}
\leanref{Formalism/SharpRestriction/GrowthBounds/StandardPolynomials/TopKeyFormSelfSoft.lean}{Formalism.MainNotions.ColumnSystems.Pf.hTop_nd}
\leanref{Formalism/SharpRestriction/GrowthBounds/StandardPolynomials/TopKeyFormSelfSoft.lean}{Formalism.MainNotions.ColumnSystems.Pf.hTop_ord}
\leanref{Formalism/SharpRestriction/GrowthBounds/StandardPolynomials/TopKeyFormSelfSoft.lean}{Formalism.MainNotions.ColumnSystems.Pf.softOK_hTop}
\leanref{Formalism/SharpRestriction/GrowthBounds/StandardPolynomials/TopKeyFormSelfSoft.lean}{Formalism.MainNotions.ColumnSystems.Pf.hpow_congr}
\leanref{Formalism/SharpRestriction/GrowthBounds/StandardPolynomials/TopKeyFormSelfSoft.lean}{Formalism.MainNotions.ColumnSystems.Pf.softH}
\leanref{Formalism/SharpRestriction/GrowthBounds/StandardPolynomials/MonomialWeightedHomogeneous.lean}{Formalism.MainNotions.ColumnSystems.Pf.top_letter_ne_zero}
\leanref{Formalism/SharpRestriction/GrowthBounds/StandardPolynomials/MonomialWeightedHomogeneous.lean}{Formalism.MainNotions.ColumnSystems.Pf.top_letter_power_ne_zero}

\begin{proof}
Write $\varpi=\varpi_0$, and let $\delta_1=1$, $\delta_c=\varpi$ for $c\geq2$, $\delta'_1=\varpi$ and $\delta'_c=1$ for $c\geq2$.
Then $\delta_c=\varpi^{d_\varpi(z_{(x,1),c})}$, $\delta'_c=\varpi^{d_\varpi(z_{(x,0),c})}$ and $\delta_c\delta'_c=\varpi$ for every $c$.
The letters $(x,1)$ and $(x,0)$ both have position $x$; write $g_x=\iota\sigma^x(g)$.

(1) By Definition~\ref{pd-symbols-path-generators}(7) with $r=1$, $\Sigma_g$ sends $z_{(x,1),b}$ to
$\sum_a\delta_a(g_x^{-1})_{ab}z_{(x,1),a}$ and $z_{(x,0),b}$ to $\sum_{a'}\delta'_{a'}(g_x)_{ba'}z_{(x,0),a'}$. As
$\sum_b(g_x^{-1})_{ab}(g_x)_{ba'}$ is $1$ if $a=a'$ and $0$ otherwise,
\[
 \Sigma_gh_x=\sum_{a,a'}\delta_a\delta'_{a'}\Big(\sum_b(g_x^{-1})_{ab}(g_x)_{ba'}\Big)z_{(x,1),a}z_{(x,0),a'}
 =\sum_a\delta_a\delta'_az_{(x,1),a}z_{(x,0),a}=\varpi h_x .
\]
The monomial $z_{(x,1),c}z_{(x,0),c}$ has $\varpi_0$-degree $d_\varpi(z_{(x,1),c})+d_\varpi(z_{(x,0),c})$, which is $0+1$ for $c=1$
and $1+0$ for $c\geq2$.

(2) By (1), $\Sigma_g(h_x^P)=\varpi^Ph_x^P\in\varpi^P\mathcal R[z]$ for every $g$, as $h_x\in\mathcal R[z]$. For elements
$u_1,\dots,u_k$ of a commutative ring, $(\sum_cu_c)^{p^e}-\sum_cu_c^{p^e}$ is a multiple of $p$, as the $p$-th power map is a
ring endomorphism modulo $p$. With $u_c=z_{(x,1),c}z_{(x,0),c}$ in $\mathbf Z[z]$ this gives $h_x^P-\hat E_{x,P}\in p\mathcal
R[z]\subset\varpi\mathcal R[z]$. So $\Sigma_1(h_x^P)-\varpi^P\hat E_{x,P}=\varpi^P(h_x^P-\hat E_{x,P})\in\varpi^{P+1}\mathcal R[z]$,
and $\hat E_{x,P}\in\mathcal R[z]$; hence $h_x^P$ is a lifted factor of $\hat E_{x,P}$ of order $P$. Each monomial
$z_{(x,1),c}^Pz_{(x,0),c}^P$ of $\hat E_{x,P}$ has $\varpi_0$-degree $P$, by the computation in (1), and degree $P$ in the
variables of each of the letters $(x,1)$ and $(x,0)$ and $0$ in the others. The monomials of $\Sigma_1(h_x^P)=\varpi^Ph_x^P$ are
products of $P$ monomials $z_{(x,1),c}z_{(x,0),c}$, so they have the same row degrees.

(3) Setting $z_{(x,1),1}=z_{(x,0),1}=1$ and all other variables to $0$ sends $h_x$ and $\hat E_{x,P}$ to $1$.
\end{proof}
% Lean: (1) is hook_key_lift_exact at r = 1 for the letters (x, true) and (x, false) (hinv from g^-1 g = 1 and hd: dW l b *
%   dW l' b = w), giving hTop_ord, and hTop_nd (wN_pair). (2) is softH S hS x e: softOK_hTop (IsFac.of_exact with e = 1,
%   r = 0) raised to the power p^e (SoftOK.pow), then IsFac.congr_sg with hpow_congr (sum_pow_char_pow in the residue
%   field of O, where p = 0, and residue_zero_mem_order_one), EhatT_wh with wN_pair, and identity_lift_mem_Pol of hTop_pol.
%   (3) is top_letter_ne_zero and top_letter_power_ne_zero, by evaluation at the first column.

\begin{lemma}\label{pd-kappa-defects-windows-lift-pairs-source-lift-tail-existence-prefactorization-key-factors}
In Situation~\ref{pd-kappa-defects-windows-lift-pairs-source-lift-tail-existence-situation}, use the notation of
Definition~\ref{pd-kappa-defects-windows-lift-pairs-source-lift-tail-existence-prefactorization-keys}. Let $1\leq j\leq q$,
$1\leq i\leq S_j$ and $2\leq t\leq k$, put $d=d_{j,i}$, and let $P\geq0$ be an integer.
\begin{enumerate}
\item $\Sigma_g\kappa^S_{j,i}$ and $\Sigma_g\kappa^W_{j,i,t}$ lie in $\varpi_0\mathcal R[z]$ for every $g\in GL_k(W)$.
\item $(\kappa^S_{j,i})^P$ is a lifted factor of $(z_{(s_j,1),1}^{p^d}z_{(y_{j,i},0),1})^P$ of order $P$, $\kappa^W_{j,i,t}$ is
  one of $z_{(y_{j,i},1),t}z_{(s_j,1),1}^{p^d}$ of order $1$, and $D_t$ is one of $\prod_{j'}z_{(s_{j'},1),1}^{p^{g_{j'}}-1}$ of
  order $0$ (Definition~\ref{pd-kappa-defects-windows-lift-pairs-source-lift-tail-existence-levels}(1)).
\item Each of $(\kappa^S_{j,i})^P$, $\kappa^W_{j,i,t}$ and $D_t$ has the row degrees of the monomial of which (2) makes it a lifted
  factor, and so has its image under $\Sigma_1$.
\end{enumerate}
\end{lemma}
\leanref{Formalism/SharpRestriction/GrowthBounds/CoarseMultiplicity/KeyMinorDvd.lean}{Formalism.MainNotions.ColumnSystems.Pf.key_pow_dvd}
\leanref{Formalism/SharpRestriction/GrowthBounds/CoarseMultiplicity/KeyMinorDvd.lean}{Formalism.MainNotions.ColumnSystems.Pf.kS_ord}
\leanref{Formalism/SharpRestriction/GrowthBounds/CoarseMultiplicity/KeyMinorDvd.lean}{Formalism.MainNotions.ColumnSystems.Pf.key_minor_dvd}
\leanref{Formalism/SharpRestriction/GrowthBounds/CoarseMultiplicity/KeyMinorDvd.lean}{Formalism.MainNotions.ColumnSystems.Pf.kW_ord}
\leanref{Formalism/SharpRestriction/GrowthBounds/CoarseMultiplicity/KeyMinorDvd.lean}{Formalism.MainNotions.ColumnSystems.Pf.key3Ord}
\leanref{Formalism/SharpRestriction/GrowthBounds/CoarseMultiplicity/KeyOne.lean}{Formalism.MainNotions.ColumnSystems.Pf.key3One_S}
\leanref{Formalism/SharpRestriction/GrowthBounds/CoarseMultiplicity/KeyOne.lean}{Formalism.MainNotions.ColumnSystems.Pf.key_lift_identity_lower}
\leanref{Formalism/SharpRestriction/GrowthBounds/CoarseMultiplicity/KeyOne.lean}{Formalism.MainNotions.ColumnSystems.Pf.loc4One}
\leanref{Formalism/SharpRestriction/GrowthBounds/StandardPolynomials/NormalWeightCovectorZero.lean}{Formalism.MainNotions.ColumnSystems.Pf.softKS}
\leanref{Formalism/SharpRestriction/GrowthBounds/StandardPolynomials/NormalWeightCovectorZero.lean}{Formalism.MainNotions.ColumnSystems.Pf.softKSP}
\leanref{Formalism/SharpRestriction/GrowthBounds/StandardPolynomials/NormalWeightCovectorZero.lean}{Formalism.MainNotions.ColumnSystems.Pf.softKW}
\leanref{Formalism/SharpRestriction/GrowthBounds/StandardPolynomials/NormalWeightCovectorZero.lean}{Formalism.MainNotions.ColumnSystems.Pf.softLoc}

\begin{proof}
Write $\varpi=\varpi_0$ and $y=y_{j,i}$, and let $\delta_c$ and $\delta'_c$ be as in the proof of
Lemma~\ref{pd-kappa-defects-windows-lift-pairs-source-lift-tail-existence-prefactorization-hooks}. As $0<d<g_j$, $p^d-1\geq1$,
and the position $y$ is $s_j+d$ modulo $n'$.

(1) Let $g\in GL_k(W)$, and put $M=\iota\sigma^{s_j}(g^{-1})$, $N'=\iota\sigma^{y}(g^{-1})$ and $N=\iota\sigma^{y}(g)$,
matrices over $\mathcal R$ with $N'N=1$. By Definition~\ref{pd-symbols-path-generators}(7) with $r=1$, $\Sigma_g$ sends
$z_{(s_j,1),c}$ to $u_c=\sum_b\delta_bM_{bc}z_{(s_j,1),b}$, $z_{(y,1),c}$ to $u'_c=\sum_b\delta_bN'_{bc}z_{(y,1),b}$ and
$z_{(y,0),c}$ to $v_c=\sum_b\delta'_bN_{cb}z_{(y,0),b}$. For $\beta\in W$, iterating $\sigma(x)\equiv x^p\pmod{pW}$ gives
$\sigma^{s_j}(\beta)^{p^d}\equiv\sigma^{s_j+d}(\beta)=\sigma^y(\beta)\pmod{pW}$, as $\sigma^{n'}=1$; so
$M_{bc}^{p^d}\equiv N'_{bc}$ modulo $p\mathcal R\subset\varpi\mathcal R$. Work modulo $\varpi\mathcal R[z]$. As
$\delta_b\in\varpi\mathcal R$ for $b\geq2$, $u_c\equiv M_{1c}z_{(s_j,1),1}$ and $u'_c\equiv N'_{1c}z_{(y,1),1}$, so
$u_c^{p^d}\equiv N'_{1c}z_{(s_j,1),1}^{p^d}$. Hence
\[
 \Sigma_g\kappa^S_{j,i}=\sum_cu_c^{p^d}v_c\equiv z_{(s_j,1),1}^{p^d}\sum_b\Big(\sum_cN'_{1c}N_{cb}\Big)\delta'_bz_{(y,0),b}
 =\delta'_1z_{(s_j,1),1}^{p^d}z_{(y,0),1}\equiv0,
\]
\[
 \Sigma_g\kappa^W_{j,i,t}=u'_tu_1^{p^d}-u'_1u_t^{p^d}\equiv N'_{1t}N'_{11}z_{(y,1),1}z_{(s_j,1),1}^{p^d}
 -N'_{11}N'_{1t}z_{(y,1),1}z_{(s_j,1),1}^{p^d}=0 .
\]
Both polynomials have coefficients in $\mathcal R$, so they lie in $\varpi\mathcal R[z]$.

(2) The substitution $\Sigma_1$ sends $z_{x,c}$ to $\delta_cz_{x,c}$ for a vector letter $x$ and to $\delta'_cz_{x,c}$ for a
covector letter. So
\[
 \Sigma_1\kappa^S_{j,i}=\varpi\Big(z_{(s_j,1),1}^{p^d}z_{(y,0),1}+\varpi^{p^d-1}\sum_{c\geq2}z_{(s_j,1),c}^{p^d}z_{(y,0),c}\Big),\qquad
 \Sigma_1\kappa^W_{j,i,t}=\varpi\big(z_{(y,1),t}z_{(s_j,1),1}^{p^d}-\varpi^{p^d-1}z_{(y,1),1}z_{(s_j,1),t}^{p^d}\big),
\]
\[
 \Sigma_1D_t=\prod_{j'}z_{(s_{j'},1),1}^{p^{g_{j'}}-1}-\varpi^E\prod_{j'}z_{(s_{j'},1),t}^{p^{g_{j'}}-1},\qquad
 E=\sum_{j'}(p^{g_{j'}}-1)\geq1,
\]
as $g_1\geq1$. The bounds on $\Sigma_g$ are (1) for $\kappa^S_{j,i}$ and $\kappa^W_{j,i,t}$, and for $D_t$ they are
$\Sigma_gD_t\in\mathcal R[z]$, by Lemma~\ref{pd-kappa-defects-windows-lift-pairs-source-lift-tail-divisible} as
$D_t\in\mathcal R[z]$. So
Lemma~\ref{pd-kappa-defects-windows-lift-pairs-source-lift-tail-existence-prefactorization-products}(2) shows that
$\kappa^S_{j,i}$, $\kappa^W_{j,i,t}$ and $D_t$ are lifted factors of the stated monomials, for $\kappa^S_{j,i}$ with $P=1$.
Lemma~\ref{pd-kappa-defects-windows-lift-pairs-source-lift-tail-existence-prefactorization-products}(1) gives the powers.

(3) Each monomial of $\kappa^S_{j,i}$ has degree $p^d$ in the variables of $(s_j,1)$ and $1$ in those of $(y,0)$; each of the
two monomials of $\kappa^W_{j,i,t}$ has degree $1$ in those of $(y,1)$ and $p^d$ in those of $(s_j,1)$; and each of the two of
$D_t$ has degree $p^{g_{j'}}-1$ in those of $(s_{j'},1)$ for every $j'$, as the $s_{j'}$ are distinct. These are the row degrees
of the monomials in (2), and $P$-th powers multiply them by $P$. The substitution $\Sigma_1$ multiplies each monomial by a power
of $\varpi$, so it keeps the row degrees.
\end{proof}
% Lean: (1) is key3Ord, from kS_ord and kW_ord: hlift writes the substitution as the image of a polynomial over O in u, v,
%   M, N' and N (liftSubst_X_vec, liftSubst_X_cov), hMN is the Frobenius congruence frobCong with segPos_val for
%   y = s_j + d, and key_pow_dvd and key_minor_dvd compute modulo C w as in the proof. (2) is softKS (IsFac.of_exact with
%   e = p^d - 1 >= 1, two_le_pd, key3One_S), softKSP (SoftOK.pow), softKW (key_lift_identity_lower) and softLoc (loc4One,
%   one_le_Mp, and inWjRg_zero_of_inA from lift_substitution_coeff_divisible for the order 0). (3) is kS_pol, kW_pol and
%   loc_pol with identity_lift_mem_Pol.

\begin{lemma}\label{pd-kappa-defects-windows-lift-pairs-source-lift-tail-existence-prefactorization-localizer}
In Situation~\ref{pd-kappa-defects-windows-lift-pairs-source-lift-tail-existence-situation}, use the notation of
Definition~\ref{pd-kappa-defects-windows-lift-pairs-source-lift-tail-existence-prefactorization-keys}, and let $1\leq j_0\leq q$.
Then $\Sigma_gL_{j_0}\in\varpi_0\mathcal R[z]$ for every $g\in GL_k(W)$ whose plane differs from that of $1$
(Lemma~\ref{pd-symbols-edge-basis}, with $r=1$).
\end{lemma}
\leanref{Formalism/SharpRestriction/GrowthBounds/CoarseMultiplicity/LiftSubstStart.lean}{Formalism.MainNotions.ColumnSystems.Pf.locR_lift}
\leanref{Formalism/SharpRestriction/GrowthBounds/CoarseMultiplicity/LiftSubstStart.lean}{Formalism.MainNotions.ColumnSystems.Pf.loc4Ord_of}
\leanref{Formalism/SharpRestriction/GrowthBounds/CoarseMultiplicity/Loc4DivOf.lean}{Formalism.MainNotions.ColumnSystems.Pf.map_xvL}
\leanref{Formalism/SharpRestriction/GrowthBounds/CoarseMultiplicity/Loc4DivOf.lean}{Formalism.MainNotions.ColumnSystems.Pf.loc4Div_of}
\leanref{Formalism/SharpRestriction/GrowthBounds/CoarseMultiplicity/Loc4ResOfLine.lean}{Formalism.MainNotions.ColumnSystems.Pf.loc4Res_of}
\leanref{Formalism/SharpRestriction/GrowthBounds/CoarseMultiplicity/Loc4ResOfLine.lean}{Formalism.MainNotions.ColumnSystems.Pf.resPowAt}
\leanref{Formalism/SharpRestriction/GrowthBounds/CoarseMultiplicity/Loc4ResOfLine.lean}{Formalism.MainNotions.ColumnSystems.Pf.resTelAt}
\leanref{Formalism/SharpRestriction/GrowthBounds/CoarseMultiplicity/PsiRes0Ne.lean}{Formalism.MainNotions.ColumnSystems.Pf.resLineAt}
\leanref{Formalism/SharpRestriction/GrowthBounds/CoarseMultiplicity/PsiRes0Ne.lean}{Formalism.MainNotions.ColumnSystems.Pf.loc4Res}
\leanref{Formalism/SharpRestriction/GrowthBounds/StandardPolynomials/NormalWeightCovectorZero.lean}{Formalism.MainNotions.ColumnSystems.Pf.locPow_ord}

\begin{proof}
Write $\varpi=\varpi_0$, and let $F=\mathcal R/\varpi\mathcal R$, a field as $\varpi$ is a uniformizer
(Definition~\ref{pd-projection-property}). Let $g\in GL_k(W)$ have plane different from that of $1$. For $1\leq c\leq k$ let
$\psi(c)\in F$ be the residue of $\iota((g^{-1})_{1c})$, and for $1\leq j\leq q$ let $\psi_j(c)$ be that of
$\iota\sigma^{s_j}((g^{-1})_{1c})$, with $s_j\in\{0,\dots,n'-1\}$. As $\sigma(x)\equiv x^p\pmod{pW}$ and $\sigma^{n'}=1$,
$\psi_j(c)=\psi(c)^{p^{s_j}}$ and $\psi(c)^{p^{n'}}=\psi(c)$.

\emph{Step 1: the residue.} By Definition~\ref{pd-symbols-path-generators}(7) with $r=1$, $\Sigma_g$ sends $z_{(s_j,1),c}$ to
$\sum_b\delta_b\,\iota\sigma^{s_j}(g^{-1})_{bc}z_{(s_j,1),b}$, with $\delta_1=1$ and $\delta_b=\varpi$ for $b\geq2$. So
$\Sigma_gL_{j_0}$ has coefficients in $\mathcal R$, and its image in $F[z]$ is
\[
 \psi_{j_0}(1)\,z_{(s_{j_0},1),1}\prod_{t=2}^k\Big(\prod_jz_{(s_j,1),1}^{p^{g_j}-1}\Big)
 \Big(\prod_j\psi_j(1)^{p^{g_j}-1}-\prod_j\psi_j(t)^{p^{g_j}-1}\Big).
\]
It remains to show that this vanishes.

\emph{Step 2: a telescoping product.} Let $y\in F$ be nonzero with $y^{p^{n'}}=y$, and put $A_j=y^{p^{s_j}}$. As $s_j+g_j$ is
$s_{j+1}$ for $j<q$ and $s_1+n'$ for $j=q$, $A_j^{p^{g_j}}=A_{j+1}$ with $A_{q+1}=A_1$. So
$\prod_jA_j^{p^{g_j}-1}\cdot\prod_jA_j=\prod_jA_{j+1}=\prod_jA_j$, and as the $A_j$ are nonzero,
$\prod_j(y^{p^{s_j}})^{p^{g_j}-1}=1$.

\emph{Step 3: the line.} If $\psi(1)=0$, then $\psi_{j_0}(1)=0$ and the residue vanishes. Otherwise some $\psi(t)$ with
$t\geq2$ is nonzero. Indeed, an element of $W$ outside $pW$ is a unit, and its image under $\iota$ is a unit of $\mathcal R$, so
its residue is nonzero. If $\psi(t)=0$ for all $t\geq2$, then $(g^{-1})_{1t}\in pW$ for $t\geq2$, and $(g^{-1})_{11}\notin pW$, as
$\psi(1)\neq0$; so the reduction modulo $p$ of the first row of $g^{-1}$ is a nonzero multiple of that of the first row of $1$,
and the plane of $g$ is that of $1$, a contradiction. Then Step 2 with $y=\psi(1)$ and with $y=\psi(t)$ gives
$\prod_j\psi_j(1)^{p^{g_j}-1}=1=\prod_j\psi_j(t)^{p^{g_j}-1}$, so the factor for $t$ vanishes, and so does the residue.
\end{proof}
% Lean: locPow_ord is inWjRg_pow of loc4Ord_of (loc4Div_of loc4Res): locR_lift and liftSubst_start write the substitution
%   as the image of locExpr over O; loc4Div_of reduces modulo w (map_xvL, C_dvd_iff_map_hom_eq_zero) to LocalResidueBound,
%   which loc4Res_of proves from resPowAt (Step 1's psi_j(c) = psi(c)^(p^s_j), through frobCong), resTelAt (Step 2) and
%   resLineAt (Step 3: psiRes0_ne for the units, and the plane ratPlane as the span of the constant coefficients of the
%   first row of g^-1). The Lean's LocalOrderBound assumes 2 <= q, which the situation gives.

\begin{lemma}\label{pd-kappa-defects-windows-lift-pairs-source-lift-tail-existence-prefactorization-room}
In Situation~\ref{pd-kappa-defects-windows-lift-pairs-source-lift-tail-existence-situation} with $C_0\geq1$,
$C_{E1},C_{E2}\geq17$ and $C_H\geq1000(C_{E1}+C_{E2}+1)(C_c+1)$, write $\rho_{j,i}$, $\rho'_{j,i}$, $\varsigma_{j,i}$ and $o_{j,+}$
for the entries of $\beta_j$, and $\hat a_j$ for the number $\hat a$ of $S_j(b_T)$
(Definition~\ref{pd-kappa-defects-segment-lattice}(1)). Let $a_{0,1},\dots,a_{0,q}\geq0$ be integers with
$a_{0,j}+\sum_{i\leq\eta_j}\varsigma_{j,i}\leq\hat a_j$ for every $j$, and let $1\leq j_0\leq q$. For $1\leq j\leq q$ put
\[
 U_j=p^{g_j}a_{0,j}+p^{g_j}\sum_{i=1}^{S_j}\varsigma_{j,i}+\sum_{i=1}^{S_j}p^{d_{j,i}}\Big(\sum_{t=2}^k\rho_{j,i}(t)+\rho'_{j,i}(1)\Big)
 +cm(p^{g_j}-1)+\epsilon_j,
\]
with $\epsilon_{j_0}=c$ and $\epsilon_j=0$ for $j\neq j_0$, and put $U_{q+1}=U_1$. Then $U_{j+1}+o_{j,+}\leq a_T(s_{j+1})$ for every
$j$.
\end{lemma}
\leanref{Formalism/SharpRestriction/GrowthBounds/StandardPolynomials/RestRoom.lean}{Formalism.MainNotions.ColumnSystems.Pf.usedX}
\leanref{Formalism/SharpRestriction/GrowthBounds/StandardPolynomials/RestRoom.lean}{Formalism.MainNotions.ColumnSystems.Pf.RoomAt}
\leanref{Formalism/SharpRestriction/GrowthBounds/StandardPolynomials/RestRoom.lean}{Formalism.MainNotions.ColumnSystems.Pf.RestRoom}
\leanref{Formalism/SharpRestriction/GrowthBounds/StandardPolynomials/ThresholdStandardSeventeen.lean}{Formalism.MainNotions.ColumnSystems.Pf.error_exponent_lower_bound}
\leanref{Formalism/SharpRestriction/GrowthBounds/StandardPolynomials/ThresholdStandardSeventeen.lean}{Formalism.MainNotions.ColumnSystems.Pf.room_of_std}
\leanref{Formalism/SharpRestriction/GrowthBounds/StandardPolynomials/ThresholdStandardSeventeen.lean}{Formalism.MainNotions.ColumnSystems.Pf.threshold_standard_seventeen}
\leanref{Formalism/SharpRestriction/GrowthBounds/StandardPolynomials/ThresholdStandardSeventeen.lean}{Formalism.MainNotions.ColumnSystems.Pf.restRoom}

\begin{proof}
The facts (a) to (g) in the proof of Lemma~\ref{pd-kappa-defects-margin} hold; we use from them $N_1\geq1$, $k\geq5$, $Q\geq1$,
$\Lambda\geq1$ and $b'>0$. By hypothesis (1) of Theorem~\ref{pd-gain-theorem-one}, $g_j\leq G'$ and $S_j\leq N_1$ for every $j$.
Also $c\geq n\geq1$, and $\kappa_1\geq0$. Let $1\leq j\leq q$ and put $j'=j+1$, so that $s_{j'}=s_{j+1}$.

\emph{The digits.} By condition (a) of a relaxed block, $\rho_{j',i}(t)<P_{j',i}$ and $\rho'_{j',i}(1)<P_{j',i}$, where
$P_{j',i}=p^{g_{j'}-d_{j',i}}$. As $p^{d_{j',i}}P_{j',i}=p^{g_{j'}}\leq p^{G'}=Q$, each term of the sum over $i$ in $U_{j'}$ is at
most $p^{d_{j',i}}kP_{j',i}\leq kQ$, and the sum is at most $N_1kQ$. Also
$cm(p^{g_{j'}}-1)+\epsilon_{j'}\leq c\big(m(p^{g_{j'}}-1)+1\big)\leq ckp^{g_{j'}}\leq ckQ$.

\emph{The first two terms.} By condition (b) of a relaxed block, $\varsigma_{j',i}=0$ for $i>\eta_{j'}$. So
$a_{0,j'}+\sum_i\varsigma_{j',i}\leq\hat a_{j'}\leq(a-E_0)/p^{g_{j'}}$, and the first two terms of $U_{j'}$ add up to at most
$a-E_0$. Hence $U_{j'}\leq a-E_0+N_1kQ+ckQ$.

\emph{The order $o_{j,+}$.} By condition (d) of a block relaxed of slack $(c-n)\kappa_1$, $o_{j,+}\leq L_0+(c-n)\kappa_1\leq
L_0+c\kappa_1$, with $L_0=a/p^2+3b'/p+(\Lambda/k+4/p^2)\,4a/p+2\Lambda b'/k$.

\emph{The error exponent.} We show $L_0+c\kappa_1+N_1kQ+ckQ\leq E_0$. As $p\geq5$, $(4/p^2)(4a/p)=(16/p)(a/p^2)\leq4a/p^2$, and
as $N_1\geq1$, $(\Lambda/k)(4a/p)\leq4\Lambda N_1a/(pk)$. So, as $C_{E1}\geq17$ and all the terms are nonnegative,
\[
 L_0\leq\frac{5a}{p^2}+\frac{3b'}p+\frac{4\Lambda N_1a}{pk}+\frac{2\Lambda b'}k
 \leq C_{E1}\Big(\frac a{p^2}+\frac{b'}p+\frac{\Lambda N_1a}{pk}+\frac{\Lambda b'}k\Big).
\]
As $kQ\geq1$, $c\geq1$ and $N_1\geq1$, each of $kN_1Qc$, $N_1kQ$ and $ckQ$ is at most $k^2N_1Q^2c$, so
$c\kappa_1+N_1kQ+ckQ=13kN_1Qc+N_1kQ+ckQ\leq15k^2N_1Q^2c\leq C_{E2}k^2N_1Q^2c$. Adding gives the claim, by the definition of $E_0$
(Definition~\ref{pd-kappa-defects-numbers}).

\emph{Conclusion.} $U_{j'}+o_{j,+}\leq a-E_0+N_1kQ+ckQ+L_0+c\kappa_1\leq a\leq a_T(s_{j+1})$, the last by the window for $a_T$.
\end{proof}
% Lean: restRoom = TransportedAlt.mono of room_of_std over threshold_standard_seventeen, whose thresholds are 1, 17 and
%   T2 CE1 CE2 Cc 1 = 1000(CE1 + CE2 + 1)(Cc + 1), with Std C0 Cde CE1 CE2 1 X from std_of (the facts). room_of_std: hN2
%   (the digits, (hbeta j').1.1, pd_mul_P, St.segB), hN3, hsum (sg i = 0 for i >= sigma), hA and hfl (the floor A0), hoP
%   (BLk's bound on oPlus with (c - n) ka1 <= c ka1), haP (St.winA) and error_exponent_lower_bound (h1, h2: 17 <= CE1,
%   CE2). U_j is usedX X beta a0 j0 j and j + 1 is rot j; the Lean's hypothesis on a0 is ColSys.a0le, over the integers.

\begin{lemma}\label{pd-kappa-defects-windows-lift-pairs-source-lift-tail-existence-prefactorization-family}
In Situation~\ref{pd-kappa-defects-windows-lift-pairs-source-lift-tail-existence-situation}, let $\ell$ be a level choice for
$\beta$ with base monomial $z^{A_\ell}$, and let $G=(G_j)$ with $G_j\in\mathcal P'_j$ nonzero of row degrees $\delta'_j$
(Definition~\ref{pd-kappa-defects-windows-lift-pairs-source-lift-tail-existence-levels}). Put $W(G)=(W_j(G_j))$. Then:
\begin{enumerate}
\item $\Phi(\beta,W(G))=\prod_jH_j(\ell_j,G_j)\cdot z^{A_\ell}$, and $\Phi(\beta,W(G))\neq0$;
\item $\Phi(\beta,W(G))\in\mathrm{Pol}_T$;
\item every monomial of $\Phi(\beta,W(G))$ has $\varpi_0$-degree $o$
  (Definition~\ref{pd-kappa-defects-windows-lift-pairs-source-lift-tail-systems}(1)).
\end{enumerate}
\end{lemma}
\leanref{Formalism/SharpRestriction/GrowthBounds/StandardPolynomials/SymFamilyWordEqHeadMulMonomial.lean}{Formalism.MainNotions.ColumnSystems.Pf.sym_family_word_eq_head_mul_monomial}
\leanref{Formalism/SharpRestriction/GrowthBounds/StandardPolynomials/SymFamilyWordEqHeadMulMonomial.lean}{Formalism.MainNotions.ColumnSystems.Pf.sym_family_word_ne_zero}
\leanref{Formalism/SharpRestriction/GrowthBounds/StandardPolynomials/IoBlockNeZero.lean}{Formalism.MainNotions.ColumnSystems.Pf.ioD_ne_zero}
\leanref{Formalism/SharpRestriction/GrowthBounds/StandardPolynomials/SymFamilyWordEqHeadMulMonomial.lean}{Formalism.MainNotions.ColumnSystems.Pf.balanced_word_homogeneous}
\leanref{Formalism/SharpRestriction/GrowthBounds/StandardPolynomials/GenWh.lean}{Formalism.MainNotions.ColumnSystems.Pf.gen_wh}
\leanref{Formalism/SharpRestriction/GrowthBounds/StandardPolynomials/IoBlockNeZero.lean}{Formalism.MainNotions.ColumnSystems.Pf.TD_wLetBL}
\leanref{Formalism/SharpRestriction/GrowthBounds/StandardPolynomials/IoBlockNeZero.lean}{Formalism.MainNotions.ColumnSystems.Pf.symPolBL}
\leanref{Formalism/SharpRestriction/GrowthBounds/StandardPolynomials/IoBlockNeZero.lean}{Formalism.MainNotions.ColumnSystems.Pf.TD_wN}
\leanref{Formalism/SharpRestriction/GrowthBounds/StandardPolynomials/IoBlockNeZero.lean}{Formalism.MainNotions.ColumnSystems.Pf.symNdBL}

\begin{proof}
(1) By Definition~\ref{pd-kappa-defects-windows-lift-pairs-model}(5), $\iota_j(\beta_j,\cdot)$ is the algebra map $\iota_j$
followed by multiplication by a monomial. As $W_j(G_j)$ is the product of
$\prod_l\hat E_l^{\ell_{j,l}}\hat E_+^{\ell_{j,+}}G_j$ with $\prod_lz_l^{\hat b_{u_l}\dot-o_{u_l}}z_+^{\alpha_+\dot-o_+}\prod_c\xi_c^{\gamma_{j,c}}$,
$\iota_j(\beta_j,W_j(G_j))$ is $H_j(\ell_j,G_j)$ times the $j$-th factor of the product that defines $z^{A_\ell}$; taking the
product over $j$ gives the identity. By Definition~\ref{pd-kappa-defects-windows-lift-pairs-model}(5), $\iota_j(\hat E_l)$ and
$\iota_j(\hat E_+)$ are the polynomials $\hat E_{x,P}$ and $h_x$ of
Definition~\ref{pd-kappa-defects-windows-lift-pairs-source-lift-tail-existence-prefactorization-keys} for $x=y_{j,u_l}$,
$P=P_{u_l}$ and $x=s_{j+1}$, which are nonzero by
Lemma~\ref{pd-kappa-defects-windows-lift-pairs-source-lift-tail-existence-prefactorization-hooks}(3); and $\iota_j(G_j)\neq0$, as
$\iota_j$ sends distinct monomials to distinct monomials (proof of
Lemma~\ref{pd-kappa-defects-windows-lift-pairs-source-symbols}). As $\mathcal K[z]$ is a domain, $\Phi(\beta,W(G))\neq0$.

For (2) and (3), fix $j$, and use the notation of Definition~\ref{pd-kappa-defects-windows-lift-pairs-form-block-data} for the
block $\beta_j$; so $\iota_j(\beta_j,X^\nu)=z^{\varepsilon(\nu)}$ with $\varepsilon=\varepsilon_{\beta_j}$
(Lemma~\ref{pd-kappa-defects-windows-lift-pairs-form-block-monomials}(2)). The computation in the proof of
Lemma~\ref{pd-kappa-defects-windows-lift-pairs-form-block-homogeneous}(1) applies to $W_j(G_j)$, the product of
Definition~\ref{pd-kappa-defects-segment-lattice}(5) for the level $\ell_j$ with $f=G_j$ and $g=\prod_c\xi_c^{\gamma_{j,c}}$: it
uses only that $f$ has degree $\delta'_j(\lambda)$ in the variables of each letter $\lambda\neq\xi$, that $g$ has degree
$b_+-\ell_{j,+}$, and the bounds on a level. So every $X^\nu$ that occurs in $W_j(G_j)$ has
$\sum_v\nu_v\omega(v)=\Omega_{\beta_j}(\omega,\omega_L)$ for every balanced weight $(\omega,\omega_L)$.

(2) Let $x$ be a letter, and let $\omega(v)=P_v$ if $x(v)=x$ and $\omega(v)=0$ otherwise, and $\omega_L(\lambda)=\omega(\lambda_1)$.
As in the proof of (a) in Lemma~\ref{pd-kappa-defects-windows-lift-pairs-form-block-valuations}, this is a balanced weight, and the
degree of $z^{\varepsilon(\nu)}$ in the variables of $x$ is $\sum_c\beta^{\beta_j}_{x,c}+\sum_v\nu_v\omega(v)$. The computation of
$t_x$ there evaluates $\sum_c\beta^{\beta_j}_{x,c}+\Omega_{\beta_j}(\omega,\omega_L)$ to $d_x$ for the letters at the $y_{j,i}$ and at
$s_{j+1}$, and to $0$ at the others. It uses the digit conditions, $w_i\leq\hat b_i$ and $\varsigma_i\leq\hat\alpha_i$ for
$i\leq\eta$, $o_i\leq\hat b_i$ for $i>\eta$, and $o_+\leq a_T(s_{j+1})$; these are conditions (a) to (d) of a relaxed block
(Definition~\ref{pd-kappa-defects-windows-lift-pairs-source-lift-relaxed}(1)). So every monomial of $\iota_j(\beta_j,W_j(G_j))$ has
degree $d_x$ in the letters $x$ at the $y_{j,i}$ and at $s_{j+1}$, and $0$ in the others. Every element of $T$ is exactly one of
the $s_{j+1}$ and the $y_{j,i}$, so every monomial of $\Phi(\beta,W(G))=\prod_j\iota_j(\beta_j,W_j(G_j))$ has degree $d_x$ in every
letter $x$.

(3) Let $\omega(v)=P_v\,d_\varpi(z_{x(v),c(v)})$ and $\omega_L(\lambda)=\omega(\lambda_1)$. By
Definition~\ref{pd-kappa-defects-windows-lift-pairs-model}(5), $\omega$ is $P_i$ at the $x_{i,c}$, $0$ at the $y_{i,c}$, $1$ at
the $x_{+,c}$, $0$ at the $\xi_c$, $0$ at $z_l$ and $z_+$, $P_{u_l}$ at $z'_l$ and $1$ at $z'_+$. So $\omega(\lambda_c)=\omega_L(\lambda)$,
$\omega(z_l)+\omega(z'_l)=P_{u_l}=\omega_L(x_{u_l})+\omega_L(y_{u_l})$ and $\omega(z_+)+\omega(z'_+)=1=\omega_L(x_+)+\omega_L(\xi)$,
and $(\omega,\omega_L)$ is a balanced weight, with
$\Omega_{\beta_j}(\omega,\omega_L)=\sum_iP_i\delta(x_i)+\delta(x_+)=\sum_{i\leq\eta}P_iw_i+\sum_{i>\eta}P_io_i+o_+$. As $d_\varpi$ is
additive, the $\varpi_0$-degree of $z^{\varepsilon(\nu)}$ is $d_\varpi(\beta^{\beta_j})+\sum_v\nu_v\omega(v)$, and
$d_\varpi(\beta^{\beta_j})=\sum_i\big(\sum_{c\geq2}\rho_i(c)+\rho'_i(1)\big)+\sum_{i\leq\eta}P_i\varsigma_i$. Adding, every monomial of
$\iota_j(\beta_j,W_j(G_j))$ has $\varpi_0$-degree $\sum_i\big(P_io^\flat_i+\sum_{c\geq2}\rho_i(c)+\rho'_i(1)\big)+o_+=\operatorname{ord}(\beta_j)$,
and every monomial of the product has $\varpi_0$-degree $\sum_j\operatorname{ord}(\beta_j)=o$.
\end{proof}
% Lean: (1) is sym_family_word_eq_head_mul_monomial (LvDat.ioBlock_word) and sym_family_word_ne_zero (word_head_formula,
%   top_letter_power_ne_zero, top_letter_ne_zero, ioD_ne_zero through ioBlock_coeff). (2) and (3) are symPolBL and
%   symNdBL applied to balanced_word_homogeneous (gen_wh with y = G_j of row degrees delta'_j, Pol from
%   column_word_mem_polynomials_of_degree in the use below): ioBlock_wh' for the weights wLet x and wN, then TD_wLetBL
%   (DigitsOK, BLk's o_+ <= a_T(s_{j+1}) and the conditions on w, sigma and o) and natPartSum for (2), TD_wN for (3).

\begin{lemma}\label{pd-kappa-defects-windows-lift-pairs-source-lift-tail-existence-prefactorization}
For every $C_c\geq0$ there is $T_0>0$ such that for every $C_0\geq T_0$ there is $T_1>0$ such that for all
$C_{E1},C_{E2}\geq T_1$ there is $T_2>0$ with the following property for every $C_H\geq T_2$. In
Situation~\ref{pd-kappa-defects-windows-lift-pairs-source-lift-tail-existence-situation} with these constants, let $\ell$ be a
level choice for $\beta$, and let $G=(G_j)$ with $G_j\in\mathcal P'_j$ nonzero and having a column system at $j$, for every $j$.
Then $G$ has a prefactorization for $\ell$.
\end{lemma}
\leanref{Formalism/SharpRestriction/GrowthBounds/StandardPolynomials/PreDef.lean}{Formalism.MainNotions.ColumnSystems.Pf.PreAt}
\leanref{Formalism/SharpRestriction/GrowthBounds/StandardPolynomials/PreDef.lean}{Formalism.MainNotions.ColumnSystems.Pf.RestPre}
\leanref{Formalism/SharpRestriction/GrowthBounds/StandardPolynomials/ColumnPolynomialCovectorZero.lean}{Formalism.MainNotions.ColumnSystems.Pf.preAt_of}
\leanref{Formalism/SharpRestriction/GrowthBounds/StandardPolynomials/ColumnPolynomialCovectorZero.lean}{Formalism.MainNotions.ColumnSystems.Pf.restPre}

\begin{proof}
\emph{The constants.} Let $C_c\geq0$, and take $T_0=1$, $T_1=17$ and $T_2=1000(C_{E1}+C_{E2}+1)(C_c+1)$, the thresholds of
Lemma~\ref{pd-kappa-defects-windows-lift-pairs-source-lift-tail-existence-prefactorization-room}. Let $C_0\geq T_0$,
$C_{E1},C_{E2}\geq T_1$ and $C_H\geq T_2$, and let the data of
Situation~\ref{pd-kappa-defects-windows-lift-pairs-source-lift-tail-existence-situation}, $\ell$ and $G$ be as in the
statement. Write $\varpi=\varpi_0$, use the notation of
Definition~\ref{pd-kappa-defects-windows-lift-pairs-source-lift-tail-existence-prefactorization-keys}, and put $j_0=1$. For
$1\leq j\leq q$ write $\rho_{j,i}$, $\rho'_{j,i}$, $o_{j,i}$, $w_{j,i}$, $\varsigma_{j,i}$ and $o_{j,+}$ for the entries of
$\beta_j$, $\hat b_{j,i}$ and $\hat a_j$ for its numbers, $P_{j,i}=p^{g_j-d_{j,i}}$ and $u_{j,l}=\eta_j+l$
(Definition~\ref{pd-kappa-defects-segment-lattice}), and $\ell_{j,l}$, $\ell_{j,+}$ and $\gamma_{j,c}$ for the parts of $\ell$.
For each $j$ fix a column system of $G_j$ at $j$, with index set $I_j$, integer $a_{0,j}$ and data $K_{j,a}$, $d_{j,a}$,
$\pi_{j,a}$, $f_{j,a}$, $o_{j,a}$ and $e_{j,a}$
(Definition~\ref{pd-kappa-defects-windows-lift-pairs-source-lift-tail-existence-column-systems}), and put
$\sigma_{j,a}=z^{\pi_{j,a}}\iota_j(K_{j,a})$. By its condition (1), $G_j=\prod_aK_{j,a}$ has row degrees $\sum_ad_{j,a}=\delta'_j$.

\emph{Step 1: the room.} By condition (4) of the column systems, $a_{0,j}+\sum_{i\leq\eta_j}\varsigma_{j,i}\leq\hat a_j$ for every
$j$. So Lemma~\ref{pd-kappa-defects-windows-lift-pairs-source-lift-tail-existence-prefactorization-room} applies to the
$a_{0,j}$ and $j_0$: with its numbers $U_j$ and $\epsilon_j$, $U_{j+1}+o_{j,+}\leq a_T(s_{j+1})$ for every $j$, where
$U_{q+1}=U_1$.

\emph{Step 2: the data.} The column part consists of the $I_j$ and the data of the column systems. For each $j$ and
$1\leq l\leq S_j-\eta_j$ put $x=y_{j,u_{j,l}}$, $P=P_{j,u_{j,l}}$, $f^H_{j,l}=h_x^P$, $o^H_{j,l}=P$ and $e^H_{j,l}=P\mathbf 1_x$.
Let $I_M$ be the disjoint union of four sets: the pairs $(j,i)$ with $1\leq i\leq S_j$, of the first kind; the triples
$(j,i,t)$ with $2\leq t\leq k$; the pairs $(j,i)$ again, of the third kind; and the integers $t$ with $2\leq t\leq k$. To
$a\in I_M$ assign $m_a$, $f^M_a$, $z^{\kappa_a}$ and $o^M_a$ as follows, and let $e^M_a$ be the row degrees of $z^{\kappa_a}$:
\begin{itemize}
\item to $(j,i)$ of the first kind, $\varsigma_{j,i}$, $(\kappa^S_{j,i})^{P_{j,i}}$,
  $z_{(s_j,1),1}^{p^{g_j}}z_{(y_{j,i},0),1}^{P_{j,i}}$ and $P_{j,i}$;
\item to $(j,i,t)$, $\rho_{j,i}(t)$, $\kappa^W_{j,i,t}$, $z_{(y_{j,i},1),t}z_{(s_j,1),1}^{p^{d_{j,i}}}$ and $1$;
\item to $(j,i)$ of the third kind, $\rho'_{j,i}(1)$, $\kappa^S_{j,i}$, $z_{(s_j,1),1}^{p^{d_{j,i}}}z_{(y_{j,i},0),1}$ and $1$;
\item to $t$, $c$, $D_t$, $\prod_jz_{(s_j,1),1}^{p^{g_j}-1}$ and $0$.
\end{itemize}
Put $o_R=\sum_j\ell_{j,+}$, and let $\rho$ be the exponent vector with the entries $\rho_{j,i}(1)$ at $z_{(y_{j,i},1),1}$,
$\rho'_{j,i}(t)$ at $z_{(y_{j,i},0),t}$ for $2\leq t\leq k$, $\gamma_{j,c}$ at $z_{(s_{j+1},0),c+1}$ for $1\leq c\leq m$, and
$a_T(s_{j+1})-o_{j,+}-U_{j+1}+\epsilon_{j+1}$ at $z_{(s_{j+1},1),1}$, and $0$ at the other variables; by Step 1 these entries are
$\geq0$. Put $f_R=\prod_jh_{s_{j+1}}^{\ell_{j,+}}\cdot z^\rho$, and let $e_R$ be $\sum_j\ell_{j,+}\mathbf 1_{s_{j+1}}$ plus the row
degrees of $z^\rho$.

\emph{Step 3: the images.} By Definition~\ref{pd-kappa-defects-windows-lift-pairs-model}(5),
$\iota_j(\hat E_l)=\hat E_{y_{j,u_{j,l}},P_{j,u_{j,l}}}$ and $\iota_j(\hat E_+)=h_{s_{j+1}}$. So $\sigma^H_{j,l}=\hat E_{x,P}$ in
the notation of Step 2, and $\sigma_R=f_R$.

\emph{Step 4: the base monomial.} We show $A_\ell=\rho+\sum_{j,a}\pi_{j,a}+\sum_am_a\kappa_a$. The truncated differences in the
definition of $A_\ell$ are differences, by conditions (b) to (d) of a relaxed block, and $\varsigma_{j,i}=0$ for $i>\eta_j$, by
condition (b). So, by Definitions~\ref{pd-kappa-defects-windows-lift-pairs-source-lift-tail-existence-levels}(3)
and~\ref{pd-kappa-defects-windows-lift-pairs-model}(5), $A_\ell$ has the entries $\rho_{j,i}(1)+P_{j,i}(\hat b_{j,i}-w_{j,i})$ for
$i\leq\eta_j$ and $\rho_{j,i}(1)+P_{j,i}(\hat b_{j,i}-o_{j,i})$ for $i>\eta_j$ at $z_{(y_{j,i},1),1}$; $\rho_{j,i}(t)$ at
$z_{(y_{j,i},1),t}$ and $\rho'_{j,i}(t)$ at $z_{(y_{j,i},0),t}$ for $t\geq2$; $\rho'_{j,i}(1)+P_{j,i}\varsigma_{j,i}$ at
$z_{(y_{j,i},0),1}$; $a_T(s_{j+1})-o_{j,+}$ at $z_{(s_{j+1},1),1}$; $\gamma_{j,c}$ at $z_{(s_{j+1},0),c+1}$; and $0$ at the other
variables. By condition (4) of the column systems, $\sum_a\pi_{j,a}$ has the entries $p^{g_j}a_{0,j}$ at $z_{(s_j,1),1}$ and
$P_{j,i}(\hat b_{j,i}-w_{j,i})$, respectively $P_{j,i}(\hat b_{j,i}-o_{j,i})$, at $z_{(y_{j,i},1),1}$, and $0$ elsewhere. As
$p^{d_{j,i}}P_{j,i}=p^{g_j}$ and the $k-1=m$ integers $t$ each contribute $c(p^{g_j}-1)$, the sum $\sum_am_a\kappa_a$ has the
entries $\rho_{j,i}(t)$ at $z_{(y_{j,i},1),t}$, $P_{j,i}\varsigma_{j,i}+\rho'_{j,i}(1)$ at $z_{(y_{j,i},0),1}$, and
$U_j-p^{g_j}a_{0,j}-\epsilon_j$ at $z_{(s_j,1),1}$, and $0$ elsewhere. Adding the three vectors, the entry at $z_{(s_{j+1},1),1}$
is $(a_T(s_{j+1})-o_{j,+}-U_{j+1}+\epsilon_{j+1})+p^{g_{j+1}}a_{0,j+1}+(U_{j+1}-p^{g_{j+1}}a_{0,j+1}-\epsilon_{j+1})=a_T(s_{j+1})-o_{j,+}$,
and the other entries are those of $A_\ell$.

\emph{Step 5: the product of the images.} As $\iota_j$ is a $\mathcal K$-algebra map, Steps 3 and 4 and
Lemma~\ref{pd-kappa-defects-windows-lift-pairs-source-lift-tail-existence-prefactorization-family}(1) give
\[
 \sigma_R\prod_{j,a}\sigma_{j,a}\prod_{j,l}(\sigma^H_{j,l})^{\ell_{j,l}}\prod_a(z^{\kappa_a})^{m_a}
 =\prod_j\iota_j\Big(\prod_l\hat E_l^{\ell_{j,l}}\cdot\hat E_+^{\ell_{j,+}}\cdot\prod_aK_{j,a}\Big)\cdot
 z^{\rho+\sum_{j,a}\pi_{j,a}+\sum_am_a\kappa_a}=\prod_jH_j(\ell_j,G_j)\cdot z^{A_\ell}=\Phi(\beta,W(G)).
\]

\emph{Step 6: conditions (a) to (e).}
\begin{itemize}
\item (a) The conditions on the $K_{j,a}$, $d_{j,a}$, $f_{j,a}$ and $\sigma_{j,a}$ are conditions (1) and (2) of the column systems,
  except that $\Sigma_1f_{j,a}$ has row degrees $e_{j,a}$: this holds as $f_{j,a}$ has them and $\Sigma_1$ multiplies each
  monomial $z^M$ by $\varpi^{d_\varpi(M)}$. The row degrees of the $\pi_{j,a}$ are at most $Q$ by condition (3); by condition (4),
  $\sum_a\pi_{j,a}$ vanishes at the variables of the letters $(s_{j'},0)$, and so does each $\pi_{j,a}$.
\item (b) This is Lemma~\ref{pd-kappa-defects-windows-lift-pairs-source-lift-tail-existence-prefactorization-hooks}(2), with Step 3,
  as $P=P_{j,u_{j,l}}$ is a power of $p$.
\item (c) This is Lemma~\ref{pd-kappa-defects-windows-lift-pairs-source-lift-tail-existence-prefactorization-key-factors}(2)
  and~(3), with $(z_{(s_j,1),1}^{p^{d_{j,i}}}z_{(y_{j,i},0),1})^{P_{j,i}}=z^{\kappa_a}$ for the first kind. The $\varpi_0$-degree of
  $z^{\kappa_a}$ is $P_{j,i}$, from $z_{(y_{j,i},0),1}$, for the first kind, $1$, from $z_{(y_{j,i},1),t}$ or $z_{(y_{j,i},0),1}$, for
  the second and third, and $0$ for the fourth; these are the $o^M_a$.
\item (d) $f_R=\sigma_R$ lies in $\mathcal R[z]$. By
  Lemma~\ref{pd-kappa-defects-windows-lift-pairs-source-lift-tail-existence-prefactorization-hooks}(1) every monomial of $h_x$ has
  $\varpi_0$-degree $1$, and $d_\varpi(\rho)=0$, as $\rho$ vanishes at the variables $z_{x,b}$ with $x$ a vector letter and
  $b\geq2$ or $x$ a covector letter and $b=1$. So every monomial of $\sigma_R$ has $\varpi_0$-degree $o_R$, and
  $\Sigma_1f_R=\varpi^{o_R}\sigma_R$. By Lemma~\ref{pd-kappa-defects-windows-lift-pairs-source-lift-tail-existence-prefactorization-hooks}(2)
  with $P=1$, $h_{s_{j+1}}$ has row degrees $\mathbf 1_{s_{j+1}}$, so $\sigma_R$ has row degrees $e_R$.
\item (e) The bounds on $e_{j,a}$ are condition (3) of the column systems. The values of $e^H_{j,l}$ are $0$ and
  $P_{j,u_{j,l}}\leq p^{g_j}\leq p^{G'}=Q$, by hypothesis (1) of Theorem~\ref{pd-gain-theorem-one}, and $y_{j,u_{j,l}}$ is odd, so
  $e^H_{j,l}$ vanishes at the letters $(s_{j'},0)$. The values of the $e^M_a$ are $0$, $P_{j,i}p^{d_{j,i}}=p^{g_j}$, $P_{j,i}$, $1$,
  $p^{d_{j,i}}$ and $p^{g_j}-1$, at most $Q$, as the $s_j$ are distinct; and the $z^{\kappa_a}$ involve only the letters $(s_j,1)$,
  $(y_{j,i},0)$ and $(y_{j,i},1)$.
\end{itemize}

\emph{Step 7: condition (f).} The identity of exponents is Step 4. By (a) to (d), every monomial of $\sigma_R$, $\sigma_{j,a}$,
$\sigma^H_{j,l}$ and $z^{\kappa_a}$ has $\varpi_0$-degree $o_R$, $o_{j,a}$, $o^H_{j,l}$ and $o^M_a$, and row degrees $e_R$, $e_{j,a}$,
$e^H_{j,l}$ and $e^M_a$. As both are additive, every monomial of the product of Step 5 has $\varpi_0$-degree
$o'=o_R+\sum_{j,a}o_{j,a}+\sum_{j,l}\ell_{j,l}o^H_{j,l}+\sum_am_ao^M_a$ and row degrees
$e'=e_R+\sum_{j,a}e_{j,a}+\sum_{j,l}\ell_{j,l}e^H_{j,l}+\sum_am_ae^M_a$. By Step 5 and
Lemma~\ref{pd-kappa-defects-windows-lift-pairs-source-lift-tail-existence-prefactorization-family}, this product is
$\Phi(\beta,W(G))$, which is nonzero, lies in $\mathrm{Pol}_T$ and has all its monomials of $\varpi_0$-degree $o$. A monomial
that occurs in it gives $o'=o$ and $e'=(x\mapsto d_x)$.

\emph{Step 8: condition (g).} Let $F$ be the product of (g). Its factors $f_R$, $f_{j,a}$, $f^H_{j,l}$ and $f^M_a$ have row degrees
$e_R$, $e_{j,a}$ (condition (2) of the column systems), $e^H_{j,l}$ and $e^M_a$
(Lemmas~\ref{pd-kappa-defects-windows-lift-pairs-source-lift-tail-existence-prefactorization-hooks}(2)
and~\ref{pd-kappa-defects-windows-lift-pairs-source-lift-tail-existence-prefactorization-key-factors}(3), as $\Sigma_1$ keeps row
degrees). By (f), $F$ has row degrees $x\mapsto d_x$, so $F\in\mathrm{Pol}_T$.

Let $g\in GL_k(W)$ have plane different from that of $1$. The entry of $\rho$ at $z_{(s_{j_0},1),1}$ is at least $\epsilon_{j_0}=c$,
so $z^\rho=z_{(s_{j_0},1),1}^c\,z^{\rho_0}$ with $z^{\rho_0}$ a monomial. The factors $D_t$ of the fourth kind have multiplicity
$c$, and $z_{(s_{j_0},1),1}^c\prod_tD_t^c=L_{j_0}^c$. So
\[
 F=L_{j_0}^c\cdot\prod_jh_{s_{j+1}}^{\ell_{j,+}}z^{\rho_0}\cdot\prod_{j,a}f_{j,a}\cdot\prod_{j,l}(f^H_{j,l})^{\ell_{j,l}}\cdot
 {\prod_a}'(f^M_a)^{m_a},
\]
with $\prod'$ over the $a$ of the first three kinds. The substitution $\Sigma_g$ is a ring homomorphism. By
Lemma~\ref{pd-kappa-defects-windows-lift-pairs-source-lift-tail-existence-prefactorization-localizer},
$\Sigma_gL_{j_0}^c\in\varpi^c\mathcal R[z]$; by
Lemma~\ref{pd-kappa-defects-windows-lift-pairs-source-lift-tail-existence-prefactorization-hooks}(1),
$\Sigma_gh_x\in\varpi\mathcal R[z]$; by Lemma~\ref{pd-kappa-defects-windows-lift-pairs-source-lift-tail-divisible},
$\Sigma_gz^{\rho_0}\in\mathcal R[z]$; and as $f_{j,a}$, $f^H_{j,l}$ and $f^M_a$ are lifted factors of orders $o_{j,a}$, $o^H_{j,l}$ and
$o^M_a$, their images under $\Sigma_g$ lie in $\varpi^{o_{j,a}}\mathcal R[z]$, $\varpi^{o^H_{j,l}}\mathcal R[z]$ and
$\varpi^{o^M_a}\mathcal R[z]$. So $\Sigma_gF\in\varpi^N\mathcal R[z]$ with
$N=c+o_R+\sum_{j,a}o_{j,a}+\sum_{j,l}\ell_{j,l}o^H_{j,l}+\sum_a'm_ao^M_a$, which is $c+o$ by (f), as $o^M_a=0$ for the fourth kind.
The bounds on $\Sigma_g$ of the lifted factors are used only here, at the $g$ whose plane differs from that of $1$.
\end{proof}
% Lean: restPre = TransportedAlt.mono preAt_of restRoom, with the thresholds of restRoom. preAt_of: hrm is RoomAt at
%   j0 = 0 with the a0le of the ColSys (Step 1), and lemmaRData builds the PreSys (Step 2): fH j i = hTop (l2x ..) ^ l2P ..,
%   multM = multMk, fM = fMk, aM = aMk, oM = oMk and dgM = dgMk on IM X (the four kinds are inl, inr inl, inr inr inl and
%   inr inr inr), fR = l2sR with R = rig = rig0 + c at ((topStart j0, true), 0), oR = sum_j jp j and dR as e_R. Step 3 is
%   lifted_block_generator (ioD_Ehat) and rigid_factor_product (ioD_Eplus). Step 4 is std_eq (stdPoly_eq, bMk_eq, sum_aMk,
%   and hX, hrot for the shift j + 1 = rot j) with hsg from BLk. Step 5 is sym_family_eq_sigma_product. Step 6 is the
%   fields KInt to plY (identity_lift_mem_Pol, column_polynomial_covector_zero), facH to polH
%   (soft_height_bound_segment_power), facM to polM (softMk), rigA to rigInt (identity_lift_normal_homog with
%   rig_factor_weighted_homogeneous_total, nd_rig) and dgCQ to dgMY (level_degree_le_prime_pow,
%   level_degree_covector_zero, dgMk_le, dgMk_Y). Step 7 is stdEq, multidegree_sum_eq_target_degrees and
%   order_sum_eq_order_family: wh_total, then IsWeightedHomogeneous.inj_right against symPolBL and symNdBL, with
%   sym_family_word_ne_zero. Step 8 is total_product_mem_polynomials and cleaned_product_order_family_eq
%   (product_split_constant with multMk0, locPow_ord, hTop_ord, inWjRg_zero_of_inA, order_total_bound and
%   sum_split_constant).

\begin{lemma}\label{pd-kappa-defects-windows-lift-pairs-source-lift-tail-existence-segment-numbers}
For every $C_c\geq0$ there is $T_0>0$ such that for every $C_0\geq T_0$ there is $T_1>0$ such that for all
$C_{E1},C_{E2}\geq T_1$ there is $T_2>0$ with the following property for every $C_H\geq T_2$. In
Situation~\ref{pd-kappa-defects-windows-lift-pairs-source-lift-tail-existence-situation} with these constants, let $1\leq e<n$
and $1\leq j\leq q$, write $s=S_j$ and $\eta=\eta_j$, and let $b$ be a block of $S_j(b_T)$ relaxed of slack $(c-n)\kappa_1$
(Definition~\ref{pd-kappa-defects-windows-lift-pairs-source-lift-relaxed}(1)), with entries $o_i,w_i,\varsigma_i,o_+$ and numbers
$\hat b_i$ and $\operatorname{cap}_b(l)$ for $\alpha_i=a_T(y_{j,i})$ (Definition~\ref{pd-kappa-defects-segment-lattice}(1)). Then:
\begin{enumerate}
\item $1+2Qe+o_i\leq\hat b_i$ for $\eta<i\leq s$;
\item let $b'$ be a block of $S_j(b_T)$, with entries $o'_i,w'_i,\varsigma'_i,o'_+$ and numbers $\hat b'_i$, $\hat\alpha'_i$ and
  $\operatorname{cap}_{b'}(l)$, whose vectors satisfy condition (a) of a relaxed block, such that $\hat b_i\leq\hat b'_i+k$ for every
  $i$; $o'_i\leq\hat b'_i$, $o'_i\leq o_i+2(k+Qe)$ and $w'_i=\varsigma'_i=0$ for $i>\eta$; $o'_i=0$, $w'_i\leq\hat b'_i$,
  $\varsigma'_i\leq\hat\alpha'_i$, $\varsigma'_i\leq\varsigma_i+k+Qe$ and $w'_i\leq w_i+k+Qe$ for $i\leq\eta$; and
  $o'_+\leq a_T(s_{j+1})$ and $o'_+\leq o_++Qe$. Then $b'$ is relaxed of slack $(c-n+e)\kappa_1$, and for $0\leq l\leq s$ and
  every integer $V\leq(N_1+1)(k+Qe)$,
  \[
   \big\lfloor\operatorname{cap}_b(l)-(\mu+n\kappa_1)\big\rfloor+N_1e+V\leq\big\lfloor\operatorname{cap}_{b'}(l)-(\mu+(n-e)\kappa_1)\big\rfloor .
  \]
\end{enumerate}
\end{lemma}
\leanref{Formalism/SharpRestriction/GrowthBounds/StandardPolynomials/TCNumDef.lean}{Formalism.MainNotions.ColumnSystems.Pf.NumCoreAt}
\leanref{Formalism/SharpRestriction/GrowthBounds/StandardPolynomials/TCNumDef.lean}{Formalism.MainNotions.ColumnSystems.Pf.RestNumCore}
\leanref{Development/AbelianVarietyTheory/Proofs/GlobalInterpolation/TensorOfPairLift.lean}{GlobalTheory.EX.RangeOne.numCore_std}
\leanref{Development/AbelianVarietyTheory/Proofs/GlobalInterpolation/TensorOfPairLift.lean}{GlobalTheory.EX.RangeOne.numCore_thr}
\leanref{Development/AbelianVarietyTheory/Proofs/GlobalInterpolation/TensorOfPairLift.lean}{GlobalTheory.EX.RangeOne.core_roomR}

\begin{proof}
\emph{The constants.} Let $T_0=T_1=1$ and $T_2=1000(C_{E1}+C_{E2}+1)(C_c+1)$, and let the data be as in the statement. Write
$P_i=p^{g_j-d_{j,i}}$, $B_i=b_T(y_{j,i})$, $\varepsilon=N_1\Lambda/k$, $E=(c-n)\kappa_1$ and $E'=(c-n+e)\kappa_1=E+e\kappa_1$, and let
$\delta$ and $\beta_e$ be as in the proof of
Lemma~\ref{pd-kappa-defects-windows-lift-pairs-source-lift-tail-existence-columns-numerics}. The facts listed there hold: $k\geq5$,
$Q\geq1$, $Y\geq1$, $N\geq2$, $k^3NQ^4cY\leq b'$, $40\varepsilon\leq\delta<1/6$, $1/k\leq\varepsilon$, $2N_1+3\leq k$,
$b'(1-\beta_e)\leq b_T(t)$ for $t\in T$, $0\leq\beta_e\leq1/12$ and $P_i\leq Q$. Also $c\geq n\geq1$ and $e\geq1$, so $Qe\geq1$ and
$kQ\geq1$.

\emph{Step 1: two numerical bounds.} As $e\kappa_1=13N_1kQe$, and $N_1e$, $N_1k$, $N_1Qe$, $k$ and $Qe$ are each at most $N_1kQe$,
\[
 \text{(A)}\quad 3k+2Qe\leq e\kappa_1,\qquad
 \text{(B)}\quad N_1e+(N_1+1)(k+Qe)+N_1(3k+2Qe)\leq e\kappa_1 ;
\]
the left side of (B) is $N_1e+4N_1k+3N_1Qe+k+Qe\leq10N_1kQe$.

\emph{Step 2: room.} Let $b''$ be a block of $S_j(b_T)$ whose vectors satisfy condition (a) of a relaxed block, with numbers
$\hat b''_i$. We show $c(\kappa_1+2Q+3k)+4\leq\frac\delta4\hat b''_i$ for every $i$. As in the proof of
Lemma~\ref{pd-kappa-defects-windows-lift-pairs-source-lift-tail-existence-columns-numerics},
$\hat b''_iP_i>B_i-kP_i\geq\frac{11}{12}b'-kQ$. As $N\geq2$, $Y\geq1$ and $Q^2\leq Q^4$, $2k^3Q^2c\leq b'$; so
$\hat b''_iQ\geq\hat b''_iP_i>\frac{11}6k^3Q^2c-kQ$, and $\hat b''_i\geq\frac{11}6k^3Qc-k$. As $\delta k\geq40\varepsilon k\geq40$,
$\frac\delta4\hat b''_i\geq\frac{10}k\hat b''_i\geq\frac{55}3k^2Qc-10$. On the other side, $N_1\leq k$ gives $\kappa_1\leq13k^2Q$, so
$c(\kappa_1+2Q+3k)\leq13k^2Qc+2Qc+3kQc$. As $k\geq5$ and $Qc\geq1$, $2Qc+3kQc+14\leq(\frac2{25}+\frac35+\frac{14}{25})k^2Qc\leq\frac{16}3k^2Qc$,
which gives the claim.

(1) Let $i>\eta$. By condition (b) of a relaxed block, $o_i\leq\hat b_i(1-\delta/4)+E$, and $E\leq c\kappa_1$. By Step 2 for $b$,
$\frac\delta4\hat b_i\geq c\kappa_1+2Qc+3kc+4$. So $o_i\leq\hat b_i-2Qc-3kc-4\leq\hat b_i-2Qe-1$, as $e\leq c$.

(2) \emph{The relaxed conditions.} Condition (a) for $b'$ is assumed. Let $i>\eta$. Then $o'_i\leq\hat b'_i$ and
$w'_i=\varsigma'_i=0$, and, by condition (b) for $b$ and as $1-\delta/4\in[0,1]$,
\[
 o'_i\leq o_i+2(k+Qe)\leq\hat b_i(1-\delta/4)+E+2(k+Qe)\leq\hat b'_i(1-\delta/4)+k+E+2k+2Qe\leq\hat b'_i(1-\delta/4)+E',
\]
by (A). Let $i\leq\eta$. Then $o'_i=0$, $w'_i\leq\hat b'_i$ and $\varsigma'_i\leq\hat\alpha'_i$, and by condition (c) for $b$ and (A),
$\varsigma'_i\leq\varsigma_i+k+Qe\leq2\Lambda\hat a/k+E+k+Qe\leq2\Lambda\hat a/k+E'$; the number $\hat a$ of the segment datum does not
depend on the block. Finally $o'_+\leq a_T(s_{j+1})$, and by condition (d) for $b$, $o'_+\leq o_++Qe$ is at most the bound there
with $E+Qe\leq E'$. So $b'$ is relaxed of slack $E'$.

\emph{The capacities.} Write $o^\flat_i$ and $o'^\flat_i$ for the numbers of $b$ and $b'$ of
Definition~\ref{pd-kappa-defects-segment-lattice}(1). For every $i$,
$(\hat b_i-o^\flat_i)-(\hat b'_i-o'^\flat_i)\leq3k+2Qe$: for $i>\eta$ this follows from $\hat b_i\leq\hat b'_i+k$ and
$o'_i\leq o_i+2(k+Qe)$, and for $i\leq\eta$, where $o^\flat_i=w_i+\varsigma_i$, from $\hat b_i\leq\hat b'_i+k$ and the bounds on
$w'_i$ and $\varsigma'_i$. Summing over $i\leq l$, with $l\leq s\leq N_1$, $\operatorname{cap}_b(l)-\operatorname{cap}_{b'}(l)\leq N_1(3k+2Qe)$.
So, with (B) and $V\leq(N_1+1)(k+Qe)$,
\[
 \operatorname{cap}_b(l)-(\mu+n\kappa_1)+N_1e+V\leq\operatorname{cap}_{b'}(l)-(\mu+n\kappa_1)+e\kappa_1=\operatorname{cap}_{b'}(l)-(\mu+(n-e)\kappa_1).
\]
As $N_1e+V$ is an integer, $\lfloor x\rfloor+N_1e+V\leq\lfloor y\rfloor$ whenever $x+N_1e+V\leq y$, and the claim follows.
\end{proof}
% Lean: numCore_thr = TransportedAlt.mono of numCore_std over thr_std (thresholds 1, 1 and T2 CE1 CE2 Cc 1, giving Std C0
%   Cde CE1 CE2 1 X). Step 1 is b1 to b5, hnumA and hnumB; Step 2 is core_room = core_roomR (coarse_bounds, segHyp_of,
%   St.levels, St.be18, St.forty_epsilon_le_delta, St.eps_ge, St.Hc); (1) is goalR (hEk, hQec). In (2), the relaxed
%   conditions are the four refine goals (hE', h2R, hAR, h4R, hPR), and the capacities are hterm with core_termS and
%   core_termU, core_capdiff and core_floor. The Lean states the comparison for every natural number l; for l > s,
%   cap(l) = cap(s).

\begin{definition}\label{pd-kappa-defects-windows-lift-pairs-source-lift-tail-existence-segment-block}
In Situation~\ref{pd-kappa-defects-windows-lift-pairs-source-lift-tail-existence-situation}, let $1\leq j\leq q$, write $s=S_j$,
$\eta=\eta_j$ and $P_i=p^{g_j-d_{j,i}}$, and number the letters of $\mathcal P_j$ as in
Definition~\ref{pd-kappa-defects-windows-lift-pairs-source-lift-tail-systems}(4). Let $u\in\mathbf Z_{\geq0}^{s-\eta}$, put
$\bar u_{\eta+l}=u_l$ for $1\leq l\leq s-\eta$ and $\bar u_i=0$ for $i\leq\eta$, let $j_+\geq0$ be an integer, $d'$ a map from the
letters of $\mathcal P_j$ to $\mathbf Z_{\geq0}$, and $z^{A'}$ a monomial of $\mathcal K[z]$, with exponent $A'(x,t)$ at $z_{x,t}$. For
$1\leq i\leq s$ and $1\leq t\leq k$, with $y=y_{j,i}$, let $q_i(t)$ and $r_i(t)$ be the quotient and the remainder of $A'((y,1),t)$
by $P_i$, and $q'_i(t)$ and $r'_i(t)$ those of $A'((y,0),t)$. Put $\tau_i=\sum_{t=2}^kq_i(t)$, $\tau'_i=\sum_{t=2}^kq'_i(t)$ and
$\tau_+=\sum_{t=2}^kA'((s_{j+1},1),t)$.
\begin{enumerate}
\item The \emph{block of $z^{A'}$} is the block $b(u,j_+,d',A')$ of $S_j(b_T)$ with the vectors $\rho_i=(r_i(t))_t$ and
  $\rho'_i=(r'_i(t))_t$; with $o_i=\bar u_i+d'(x_i)+q'_i(1)+\tau_i$ and $w_i=\varsigma_i=0$ for $i>\eta$; with $o_i=0$,
  $w_i=d'(x_i)+\tau_i$ and $\varsigma_i=q'_i(1)$ for $i\leq\eta$; and with $o_+=j_++d'(x_+)+\tau_+$.
\item Its \emph{quotient monomial} is
  $m(A')=\prod_{i=1}^s\prod_{c=1}^mx_{i,c}^{q_i(c+1)}y_{i,c}^{q'_i(c+1)}\cdot\prod_{c=1}^mx_{+,c}^{A'((s_{j+1},1),c+1)}\in\mathcal P'_j$.
\end{enumerate}
\end{definition}
\leanref{Formalism/SharpRestriction/GrowthBounds/StandardPolynomials/SegmentBlockOfMonomial.lean}{Formalism.MainNotions.ColumnSystems.Pf.segmentBlockOfMonomial}
\leanref{Formalism/SharpRestriction/GrowthBounds/StandardPolynomials/SegmentBlockOfMonomial.lean}{Formalism.MainNotions.ColumnSystems.Pf.mEps}
\leanref{Formalism/SharpRestriction/GrowthBounds/StandardPolynomials/SegmentBlockOfMonomial.lean}{Formalism.MainNotions.ColumnSystems.Pf.ezA}
\leanref{Formalism/SharpRestriction/GrowthBounds/StandardPolynomials/SegmentBlockOfMonomial.lean}{Formalism.MainNotions.ColumnSystems.Pf.y0A}
% Lean: segmentBlockOfMonomial C0 CE1 CE2 X j ju jp du A, with ju = u, jp = j_+, du = d' and A = A'; jE S ju i is u-bar_i.
%   z0A, y0A, ezA and eyA are q_i(1), q'_i(1), q_i(c + 1) and q'_i(c + 1), and axA the exponents A'((s_{j+1}, 1), c + 1);
%   mEps (ezA) (eyA) (axA) is m(A').

\begin{lemma}\label{pd-kappa-defects-windows-lift-pairs-source-lift-tail-existence-segment-block-numbers}
In Situation~\ref{pd-kappa-defects-windows-lift-pairs-source-lift-tail-existence-situation}, let $j$, $u$, $\bar u$, $j_+$, $d'$ and
$z^{A'}$ be as in Definition~\ref{pd-kappa-defects-windows-lift-pairs-source-lift-tail-existence-segment-block}, with its notation,
and let $b=b(u,j_+,d',A')$, with numbers $\hat b_i$ and $\hat\alpha_i$ for $\alpha_i=a_T(y_{j,i})$
(Definition~\ref{pd-kappa-defects-segment-lattice}(1)).
\begin{enumerate}
\item $\operatorname{ord}(b)=\sum_{i=1}^s\big(P_i(\bar u_i+d'(x_i))+\sum_{t=2}^kA'((y_{j,i},1),t)+A'((y_{j,i},0),1)\big)+j_++d'(x_+)+\tau_+$
  (Lemma~\ref{pd-kappa-defects-windows-lift-pairs-source-order}).
\item For every block $b''$ of $S_j(b_T)$, with numbers $\hat b''_i$, $\hat b''_i\leq\hat b_i+k$ for every $i$.
\item Let $1\leq i\leq s$ and $y=y_{j,i}$. If $P_i(\bar u_i+d'(x_i))+\sum_tA'((y,1),t)=b_T(y)$, then
  $\hat b_i=\bar u_i+d'(x_i)+q_i(1)+\tau_i$; if $P_i(\bar u_i+d'(y_i))+\sum_tA'((y,0),t)=a_T(y)$, then
  $\hat\alpha_i=\bar u_i+d'(y_i)+q'_i(1)+\tau'_i$. If both identities hold for every $i$, the vectors of $b$ satisfy condition (a) of a
  relaxed block (Definition~\ref{pd-kappa-defects-windows-lift-pairs-source-lift-relaxed}(1)).
\item Let $b''$ be a block of $S_j(b_T)$ whose vectors $\varrho_i,\varrho'_i$ satisfy condition (a) of a relaxed block, $Q''\geq0$ an
  integer, $1\leq i\leq s$ and $y=y_{j,i}$.
  \begin{enumerate}
  \item If $\sum_{t=2}^k\big(A'((y,1),t)\dot-\varrho_i(t)\big)\leq Q''$, then $\tau_i\leq Q''$.
  \item If $Q''\geq1$ (a bound that is not used) and $A'((y,0),1)\leq\varrho'_i(1)+P_it'+Q''$ for an integer $t'\geq0$, then
    $q'_i(1)\leq t'+Q''$.
  \item If $Q''\geq1$ (a bound that is not used), $D\geq1+2Q''$ is an integer, $\varrho_i(1)+P_iD\leq A'((y,1),1)+Q''$ and
    $A'((y,0),1)\leq\varrho'_i(1)+Q''$, then $q'_i(1)\leq q_i(1)$.
  \end{enumerate}
\end{enumerate}
\end{lemma}
\leanref{Formalism/SharpRestriction/GrowthBounds/StandardPolynomials/BlockYZeroLeOfBound.lean}{Formalism.MainNotions.ColumnSystems.Pf.nb_ord}
\leanref{Formalism/SharpRestriction/GrowthBounds/StandardPolynomials/BlockYZeroLeOfBound.lean}{Formalism.MainNotions.ColumnSystems.Pf.nb_coarseA_ge}
\leanref{Formalism/SharpRestriction/GrowthBounds/StandardPolynomials/BlockYZeroLeOfBound.lean}{Formalism.MainNotions.ColumnSystems.Pf.nb_coarseA}
\leanref{Formalism/SharpRestriction/GrowthBounds/StandardPolynomials/BlockYZeroLeOfBound.lean}{Formalism.MainNotions.ColumnSystems.Pf.nb_coarseC}
\leanref{Formalism/SharpRestriction/GrowthBounds/StandardPolynomials/BlockYZeroLeOfBound.lean}{Formalism.MainNotions.ColumnSystems.Pf.nb_digits}
\leanref{Formalism/SharpRestriction/GrowthBounds/StandardPolynomials/BlockYZeroLeOfBound.lean}{Formalism.MainNotions.ColumnSystems.Pf.nb_ez_le}
\leanref{Formalism/SharpRestriction/GrowthBounds/StandardPolynomials/BlockYZeroLeOfBound.lean}{Formalism.MainNotions.ColumnSystems.Pf.block_y_zero_le_of_bound}
\leanref{Formalism/SharpRestriction/GrowthBounds/StandardPolynomials/BlockYZeroLeOfBound.lean}{Formalism.MainNotions.ColumnSystems.Pf.block_y_zero_le_z_zero}

\begin{proof}
Write $\rho_i,\rho'_i$ and $o_i,w_i,\varsigma_i,o_+$ for the entries of $b$.

(1) For every $i$ the number $o^\flat_i$ of $b$ is $\bar u_i+d'(x_i)+q'_i(1)+\tau_i$: for $i>\eta$ it is $o_i$, and for $i\leq\eta$ it is
$w_i+\varsigma_i$, with $\bar u_i=0$. As $P_iq_i(t)+r_i(t)=A'((y_{j,i},1),t)$ and $P_iq'_i(1)+r'_i(1)=A'((y_{j,i},0),1)$,
\[
 P_io^\flat_i+\sum_{t=2}^k\rho_i(t)+\rho'_i(1)=P_i(\bar u_i+d'(x_i))+\sum_{t=2}^kA'((y_{j,i},1),t)+A'((y_{j,i},0),1),
\]
and $o_+=j_++d'(x_+)+\tau_+$. Summing gives $\operatorname{ord}(b)$.

(2) Write $B=b_T(y_{j,i})$. The $k$ entries of $\rho_i$ are remainders by $P_i$, so $|\rho_i|\leq k(P_i-1)$. Hence
$\hat b''_i=\lfloor(B\dot-|\rho''_i|)/P_i\rfloor\leq\lfloor B/P_i\rfloor\leq\lfloor(B\dot-|\rho_i|)/P_i\rfloor+k=\hat b_i+k$, with $\rho''_i$
the vector of $b''$.

(3) As $\sum_tA'((y,1),t)=P_i\sum_tq_i(t)+|\rho_i|$, the first identity gives $b_T(y)-|\rho_i|=P_i(\bar u_i+d'(x_i)+q_i(1)+\tau_i)$, so
$\hat b_i$ is as stated; in the same way the second identity gives $\hat\alpha_i$. Then $|\rho_i|\leq b_T(y)$, $P_i$ divides
$b_T(y)-|\rho_i|$, and the entries of $\rho_i$ are less than $P_i$; the same holds for $\rho'_i$ and $a_T(y)$. This is condition (a).

(4) (a) For $t\geq2$, $P_iq_i(t)\leq A'((y,1),t)$ and $\varrho_i(t)<P_i$, so $q_i(t)\leq A'((y,1),t)\dot-\varrho_i(t)$: if $q_i(t)\geq1$, then
$A'((y,1),t)-\varrho_i(t)>P_i(q_i(t)-1)\geq q_i(t)-1$. Summing over $t\geq2$ gives $\tau_i\leq Q''$.
(b) As $\varrho'_i(1)<P_i$ and $Q''\leq Q''P_i$, $A'((y,0),1)<P_i(t'+1)+Q''P_i$, so $q'_i(1)\leq t'+Q''$.
(c) By (b) with $t'=0$, $q'_i(1)\leq Q''$. As $D\geq1+2Q''$, $A'((y,1),1)\geq P_i(1+2Q'')-Q''\geq P_iQ''$, so $q_i(1)\geq Q''$.
\end{proof}
% Lean: (1) is nb_ord (nb_oEff, nb_sum_mod_div); (2) is nb_coarseA_ge; (3) is nb_coarseA, nb_coarseC and nb_digits (nb_key);
%   (4) is nb_ez_le (nb_q_le), block_y_zero_le_of_bound (nb_div_le) and block_y_zero_le_z_zero, whose hypotheses hQ : 1 <= Qe
%   are not used.

\begin{lemma}\label{pd-kappa-defects-windows-lift-pairs-source-lift-tail-existence-segment-weights}
Let $K$ be a field of characteristic $0$ and $N\geq0$, $r\geq1$ integers, fix a lettering of $\{0,\dots,N-1\}$, and use the spaces
$Y(\mathcal L,M)$ of Definition~\ref{pd-kappa-defects-segment-lattice}(4) in $K[z]=K[z_{i,c}:0\leq i<N,\ 1\leq c\leq r]$, with
$Y(\mathcal L,+\infty)=K[z]$. Let $J\geq0$ be an integer and $\mathcal L$ the set of the letters $i<J$.
\begin{enumerate}
\item If $M\leq M'$ in $\mathbf Z\cup\{+\infty\}$, then $Y(\mathcal L,M)\subset Y(\mathcal L,M')$. If $F\in Y(\mathcal L,M)$ and
  $F'\in Y(\mathcal L,M')$, then $FF'\in Y(\mathcal L,M+M')$.
\item For every letter $a$ and $c\leq r$, $z_{a,c}\in Y(\mathcal L,v(a))$, where $v(a)=1$ if $a$ is a vector letter in $\mathcal L$ and
  $v(a)=0$ otherwise.
\item For a vector letter $a$ and a covector letter $a'$, $\sum_cz_{a,c}z_{a',c}\in Y(\mathcal L,\pi(a,a'))$, where $\pi(a,a')=1$ if
  $a\in\mathcal L$ and $a'\notin\mathcal L$, and $\pi(a,a')=0$ otherwise.
\end{enumerate}
\end{lemma}
\leanref{Formalism/SharpRestriction/GrowthBounds/StandardPolynomials/RecursionWithDensities.lean}{Formalism.MainNotions.ColumnSystems.Pf.ycond_of_stable}
\leanref{Formalism/SharpRestriction/GrowthBounds/StandardPolynomials/RecursionWithDensities.lean}{Formalism.MainNotions.ColumnSystems.Pf.ycond_mono}
\leanref{Formalism/SharpRestriction/GrowthBounds/StandardPolynomials/RecursionWithDensities.lean}{Formalism.MainNotions.ColumnSystems.Pf.ycond_mul}
\leanref{Formalism/SharpRestriction/GrowthBounds/StandardPolynomials/RecursionWithDensities.lean}{Formalism.MainNotions.ColumnSystems.Pf.ycond_X}
\leanref{Formalism/SharpRestriction/GrowthBounds/StandardPolynomials/RecursionWithDensities.lean}{Formalism.MainNotions.ColumnSystems.Pf.ycond_pair}

\begin{proof}
Write $E_{cd}=E^{\mathcal L}_{cd}$. By definition, $Y(\mathcal L,M)$ contains every subspace that is stable under the $E_{cd}$ and
contained in the span of the monomials of weight at most $M$ for $\mathcal L$, and it is itself such a subspace.

(1) The first claim holds as $Y(\mathcal L,M)$ is stable and lies in the span of the monomials of weight at most $M'$. For the
second, let $V$ be the span of the products $FF'$ with $F\in Y(\mathcal L,M)$ and $F'\in Y(\mathcal L,M')$. A monomial of $FF'$ is
a product of a monomial of $F$ and one of $F'$, and the weight is additive; so $V$ lies in the span of the monomials of weight
at most $M+M'$. Each $E_{cd}$ is a derivation, so $E_{cd}(FF')=E_{cd}(F)F'+FE_{cd}(F')\in V$. Hence $V\subset Y(\mathcal L,M+M')$.

(2) Let $U$ be the span of the $z_{a,c'}$, $c'\leq r$. If $a\in\mathcal L$ is a vector letter, $E_{cd}z_{a,c'}=\delta_{dc'}z_{a,c}$; if
it is a covector letter, $E_{cd}z_{a,c'}=-\delta_{cc'}z_{a,d}$; and if $a\notin\mathcal L$, $E_{cd}z_{a,c'}=0$. So $U$ is stable. The
weight of $z_{a,c'}$ is $1$ or $0$ for a vector letter in $\mathcal L$, $-1$ or $0$ for a covector letter in $\mathcal L$, and $0$ for
$a\notin\mathcal L$; so it is at most $v(a)$, and $U\subset Y(\mathcal L,v(a))$.

(3) If $a,a'\in\mathcal L$, then $E_{cd}\sum_ez_{a,e}z_{a',e}=z_{a,c}z_{a',d}-z_{a,c}z_{a',d}=0$, and the weight of each $z_{a,e}z_{a',e}$
is $0$; so the line it spans lies in $Y(\mathcal L,0)$. Otherwise, the span of the products $z_{a,e}z_{a',e'}$ is stable, by (2) and
the Leibniz rule, and the weight of each product is at most $v(a)+0$, which is $\pi(a,a')$ in these cases, as a covector letter
has weight at most $0$.
\end{proof}
% Lean: ycond_of_stable (the maximality), ycond_mono and ycond_mul (Submodule.mul_le, chWt_add, glAct_mul'), ycond_X
%   (glAct_X') and ycond_pair (the case split hboth), all for YCond K LL j M with M in WithTop Z, over a field K of
%   characteristic 0.

\begin{lemma}\label{pd-kappa-defects-windows-lift-pairs-source-lift-tail-existence-segment-quotient}
Let $p$, $\mathcal O\subset\mathcal K$, $r\geq1$, a segment datum $S$ of length $s$, integers $\alpha_1,\dots,\alpha_s,\alpha_+\geq0$ and a
real number $\mu$ be as in Definition~\ref{pd-kappa-defects-segment-lattice}, and let $b'$ be a block, with entries $o'_i$ and
$o'_+$, numbers $\hat b'_i$, degrees $\delta$ of the letters (Definition~\ref{pd-kappa-defects-segment-lattice}(3)) and the
$\mathcal L_l$, $M_l$, $\mathcal L_+$, $M_+$ and $u_l=\eta+l$ of its generators (Definition~\ref{pd-kappa-defects-segment-lattice}(5)).
Let $u,\zeta,\zeta'\in\mathbf Z_{\geq0}^{s-\eta}$, $j_+,\zeta_+\geq0$ integers, $\gamma\in\mathbf Z_{\geq0}^r$, and $G$ a polynomial in the
letters with coefficients in $\mathcal O$ and row degrees $d_G$. Put $R_l=\sum_cx_{u_l,c}y_{u_l,c}$, and suppose:
\begin{enumerate}
\item $\zeta'_l\leq\zeta_l$ and $\hat b'_{u_l}\dot-o'_{u_l}=\zeta_l-\zeta'_l$ for every $l$, and $\alpha_+\dot-o'_+=\zeta_+$;
\item $\delta(x_{u_l})=d_G(x_{u_l})+u_l+\zeta'_l$ and $\delta(y_{u_l})=d_G(y_{u_l})+u_l+\zeta'_l$ for every $l$,
  $\delta(x_+)=d_G(x_+)+j_+$, and $\delta(\lambda)=d_G(\lambda)$ for the other letters $\lambda$;
\item $b_+=j_++\sum_c\gamma_c$;
\item for every $m\in\mathbf Z_{\geq0}^{s-\eta}$ with $m\leq\zeta'$, the polynomial $G\prod_lR_l^{\zeta'_l-m_l}$ lies in
  $Y(\mathcal L_l,M_l)$ for every $l$ and in $Y(\mathcal L_+,M_+)$.
\end{enumerate}
Then $\prod_l\hat E_l^{u_l}\cdot\hat E_+^{j_+}\cdot G\cdot\prod_lz_l^{\zeta_l}z_l'^{\zeta'_l}\cdot z_+^{\zeta_+}\cdot\prod_c\xi_c^{\gamma_c}$
lies in the $\mathcal O$-span of the generators of all levels of $b'$.
\end{lemma}
\leanref{Formalism/SharpRestriction/GrowthBounds/StandardPolynomials/DensitySumsAndDomination.lean}{Formalism.MainNotions.ColumnSystems.Pf.uTot}
\leanref{Formalism/SharpRestriction/GrowthBounds/StandardPolynomials/DensitySumsAndDomination.lean}{Formalism.MainNotions.ColumnSystems.Pf.uTot_mem}

\begin{proof}
As $z_lz'_l=\hat E_l-R_l$, the binomial formula gives
\[
 z_l^{\zeta_l}z_l'^{\zeta'_l}=z_l^{\zeta_l-\zeta'_l}(\hat E_l-R_l)^{\zeta'_l}
 =\sum_{m_l=0}^{\zeta'_l}(-1)^{\zeta'_l-m_l}\binom{\zeta'_l}{m_l}z_l^{\zeta_l-\zeta'_l}\hat E_l^{m_l}R_l^{\zeta'_l-m_l}.
\]
Multiplying over $l$, the product of the statement is the sum, over the $m\in\mathbf Z_{\geq0}^{s-\eta}$ with $m\leq\zeta'$, of
$\pm\prod_l\binom{\zeta'_l}{m_l}$ times
\[
 g_m=\prod_l\hat E_l^{u_l+m_l}\cdot\hat E_+^{j_+}\cdot\prod_lz_l^{\zeta_l-\zeta'_l}\cdot z_+^{\zeta_+}\cdot f_m\cdot\prod_c\xi_c^{\gamma_c},\qquad
 f_m=G\prod_lR_l^{\zeta'_l-m_l}.
\]
We show that $g_m$ is a generator of $b'$ of level $(u+m,j_+)$. By (1), the exponents of $z_l$ and $z_+$ are
$\hat b'_{u_l}\dot-o'_{u_l}$ and $\alpha_+\dot-o'_+$. As $R_l$ has degree $1$ in the variables of each of $x_{u_l}$ and $y_{u_l}$, (2)
shows that $f_m$ has degree $\delta(x_{u_l})-(u_l+m_l)$ in those of $x_{u_l}$ and of $y_{u_l}$, $\delta(x_+)-j_+$ in those of $x_+$,
and $\delta(\lambda)$ in those of every other letter $\lambda$. So $u_l+m_l\leq\min(\delta(x_{u_l}),\delta(y_{u_l}))$, $j_+\leq\delta(x_+)$,
$j_+\leq b_+$ by (3), and these are the degrees $\delta'$ of the level $(u+m,j_+)$. The polynomial $f_m$ has coefficients in
$\mathcal O$ and lies in the spaces $Y$ of the level by (4); and $\prod_c\xi_c^{\gamma_c}$ is homogeneous of degree
$\sum_c\gamma_c=b_+-j_+$ by (3). So $g_m$ is a generator, and the product lies in the $\mathcal O$-span of the generators.
\end{proof}
% Lean: uTot_mem (DensitySumsAndDomination.lean): hZS and hexp (EhS_eq, add_pow) expand the product; eterm writes each term
%   as an integer multiple cf of a generator; hLv, from hdeg, are the level conditions, and the generator's fields are hz,
%   hx, hdeg, hGI, hGP, hY and hB (xi_mono_mem). The Lean's hdeg is (2), with jE S (ju + sU) the numbers u_l + zeta'_l.

\begin{lemma}\label{pd-kappa-defects-windows-lift-pairs-source-lift-tail-existence-tail-condition-segment}
For every $C_c\geq0$ there is $T_0>0$ such that for every $C_0\geq T_0$ there is $T_1>0$ such that for all
$C_{E1},C_{E2}\geq T_1$ there is $T_2>0$ with the following property for every $C_H\geq T_2$. In
Situation~\ref{pd-kappa-defects-windows-lift-pairs-source-lift-tail-existence-situation} with these constants, let $\ell$ be a
level choice for $\beta$, let $1\leq e<n$ and $1\leq j\leq q$, use the notation of the situation for $\beta_j$, with
$P_i=p^{g_j-d_{j,i}}$, and write $\ell_{j,l}$, $\ell_{j,+}$ and $\gamma_{j,c}$ for the parts of $\ell$
(Definition~\ref{pd-kappa-defects-windows-lift-pairs-source-lift-tail-existence-levels}). Let $u\in\mathbf Z_{\geq0}^{s-\eta}$ with
$u_l\leq\ell_{j,l}$ for every $l$, and put $\bar u_{u_l}=u_l$ for $1\leq l\leq s-\eta$ and $\bar u_i=0$ for $i\leq\eta$. Let
$G'\in\mathcal P'_j$ have coefficients in $\mathcal R$ and row degrees $d'\leq\delta'_j$, and lie in $Y(\mathcal L_l,M_l+N_1e)$ for
every $l$ and in $Y(\mathcal L_+,M_++N_1e)$, with the $\mathcal L_l$, $M_l$, $\mathcal L_+$ and $M_+$ of the generators of
$\beta_j$. Let $z^{A'}$ be a monomial, with exponent $A'(x,t)$ at $z_{x,t}$, such that for $1\leq i\leq s$, with $y=y_{j,i}$:
\begin{enumerate}
\item $\rho_i(1)+P_i(\hat b_i\dot-w_i)\leq A'((y,1),1)+Qe$ if $i\leq\eta$, and $\rho_i(1)+P_i(\hat b_i\dot-o_i)\leq A'((y,1),1)+Qe$
  if $i>\eta$;
\item $\sum_{t=2}^k\big(A'((y,1),t)\dot-\rho_i(t)\big)\leq Qe$;
\item $A'((y,0),1)\leq\rho'_i(1)+P_i\varsigma_i+Qe$ if $i\leq\eta$, and $A'((y,0),1)\leq\rho'_i(1)+Qe$ if $i>\eta$;
\item $\sum_{t=2}^kA'((s_{j+1},1),t)\leq Qe$, $A'((s_{j+1},0),1)=0$, and $A'((s_{j+1},0),c+1)=\gamma_{j,c}$ for $1\leq c\leq m$;
\item $P_i(\bar u_i+d'(x_i))+\sum_tA'((y,1),t)=b_T(y)$, $P_i(\bar u_i+d'(y_i))+\sum_tA'((y,0),t)=a_T(y)$, and
  $\ell_{j,+}+d'(x_+)+\sum_tA'((s_{j+1},1),t)=a_T(s_{j+1})$.
\end{enumerate}
Let $z^{A'_j}$ be the product of the factors of $z^{A'}$ in the variables of the letters at the $y_{j,i}$ and at $s_{j+1}$, and
put
\[
 \nu=\sum_{i=1}^s\Big(P_i\big(\bar u_i+d'(x_i)\big)+\sum_{t=2}^kA'((y_{j,i},1),t)+A'((y_{j,i},0),1)\Big)+\ell_{j,+}+d'(x_+)
 +\sum_{t=2}^kA'((s_{j+1},1),t).
\]
Then there are a block $b'$ of $S_j(b_T)$ relaxed of slack $(c-n+e)\kappa_1$
(Definition~\ref{pd-kappa-defects-windows-lift-pairs-source-lift-relaxed}(1)) with $\operatorname{ord}(b')=\nu$, and
$u'\in Z_j(b';\mu+(n-e)\kappa_1)$ (Definition~\ref{pd-kappa-defects-windows-lift-pairs-model}(6)), with
$H_j(u,G')\,z^{A'_j}=\iota_j(b',u')$.
\end{lemma}
\leanref{Formalism/SharpRestriction/GrowthBounds/StandardPolynomials/RestSegCore.lean}{Formalism.MainNotions.ColumnSystems.Pf.SegCoreAt}
\leanref{Formalism/SharpRestriction/GrowthBounds/StandardPolynomials/RestSegCore.lean}{Formalism.MainNotions.ColumnSystems.Pf.RestSegCore}
\leanref{Formalism/SharpRestriction/GrowthBounds/StandardPolynomials/RestSegCore.lean}{Formalism.MainNotions.ColumnSystems.Pf.segMon}
\leanref{Formalism/SharpRestriction/GrowthBounds/StandardPolynomials/RestSegCore.lean}{Formalism.MainNotions.ColumnSystems.Pf.nuSeg}
\leanref{Formalism/SharpRestriction/GrowthBounds/StandardPolynomials/RestSegCoreOfNum.lean}{Formalism.MainNotions.ColumnSystems.Pf.segCoreAt_of_num}
\leanref{Formalism/SharpRestriction/GrowthBounds/StandardPolynomials/RestSegCoreOfNum.lean}{Formalism.MainNotions.ColumnSystems.Pf.restSegCore_of_num}
\leanref{Formalism/SharpRestriction/GrowthBounds/StandardPolynomials/TCNumDef.lean}{Formalism.MainNotions.ColumnSystems.Pf.NumCoreAt}

\begin{proof}
\emph{The constants.} Let $C_c\geq0$, and take the thresholds $T_0=T_1=1$ and $T_2=1000(C_{E1}+C_{E2}+1)(C_c+1)$ of
Lemma~\ref{pd-kappa-defects-windows-lift-pairs-source-lift-tail-existence-segment-numbers}. Let the constants and the data be as
in the statement; so $Qe\geq1$, and $\beta_j$ is relaxed of slack $(c-n)\kappa_1$. Let $b'=b(u,\ell_{j,+},d',A')$, with the notation
$q_i,q'_i,r_i,r'_i,\tau_i,\tau'_i,\tau_+$ of Definition~\ref{pd-kappa-defects-windows-lift-pairs-source-lift-tail-existence-segment-block};
write $o'_i,w'_i,\varsigma'_i,o'_+$ for its entries other than its vectors $r_i,r'_i$, and $\hat b'_i$, $\hat\alpha'_i$ for its numbers.

\emph{Step 1: the numbers of $b'$.} By (5) and
Lemma~\ref{pd-kappa-defects-windows-lift-pairs-source-lift-tail-existence-segment-block-numbers}(3), the vectors of $b'$ satisfy
condition (a) of a relaxed block, $\hat b'_i=\bar u_i+d'(x_i)+q_i(1)+\tau_i$ and $\hat\alpha'_i=\bar u_i+d'(y_i)+q'_i(1)+\tau'_i$. Apply
part (4) of that lemma with $b''=\beta_j$ and $Q''=Qe$. By (2) and (4)(a), $\tau_i\leq Qe$. By (3) and (4)(b), with $t'=\varsigma_i$
for $i\leq\eta$ and $t'=0$ for $i>\eta$, $q'_i(1)\leq\varsigma_i+Qe$ for $i\leq\eta$ and $q'_i(1)\leq Qe$ for $i>\eta$. Let $i>\eta$. By
Lemma~\ref{pd-kappa-defects-windows-lift-pairs-source-lift-tail-existence-segment-numbers}(1), $D=\hat b_i-o_i\geq1+2Qe$, and (1) and
(3) are the remaining hypotheses of (4)(c); so $q'_i(1)\leq q_i(1)$. By (4), $\tau_+\leq Qe$.

\emph{Step 2: the level.} At the level $\ell_j$ of $\beta_j$ (Definition~\ref{pd-kappa-defects-segment-lattice}(5)),
$\delta'_j(x_i)=w_i$ for $i\leq\eta$, $\delta'_j(x_{u_l})=o_{u_l}-\ell_{j,l}$ and $\delta'_j(x_+)=o_+-\ell_{j,+}$, with $\ell_{j,l}\leq o_{u_l}$
and $\ell_{j,+}\leq o_+$. As $d'\leq\delta'_j$ and $u_l\leq\ell_{j,l}$: $d'(x_i)\leq w_i$ for $i\leq\eta$, $\bar u_i+d'(x_i)\leq o_i$ for
$i>\eta$, and $\ell_{j,+}+d'(x_+)\leq o_+$.

\emph{Step 3: the block $b'$.} We check the hypotheses of
Lemma~\ref{pd-kappa-defects-windows-lift-pairs-source-lift-tail-existence-segment-numbers}(2) for $b=\beta_j$ and $b'$. Condition
(a) for $b'$ is Step 1, and $\hat b_i\leq\hat b'_i+k$ is
Lemma~\ref{pd-kappa-defects-windows-lift-pairs-source-lift-tail-existence-segment-block-numbers}(2). Let $i>\eta$. Then
$o'_i=\bar u_i+d'(x_i)+q'_i(1)+\tau_i\leq\hat b'_i$, as $q'_i(1)\leq q_i(1)$; $o'_i\leq o_i+2Qe$ by Steps 1 and 2; and
$w'_i=\varsigma'_i=0$. Let $i\leq\eta$. Then $o'_i=0$; $w'_i=d'(x_i)+\tau_i\leq\hat b'_i$, as $\bar u_i=0$; $\varsigma'_i=q'_i(1)\leq\hat\alpha'_i$;
and $\varsigma'_i\leq\varsigma_i+Qe$ and $w'_i\leq w_i+Qe$ by Steps 1 and 2. By the third identity of (5),
$o'_+=\ell_{j,+}+d'(x_+)+\tau_+=a_T(s_{j+1})-A'((s_{j+1},1),1)\leq a_T(s_{j+1})$, and $o'_+\leq o_++Qe$ by Steps 1 and 2. So $b'$ is
relaxed of slack $(c-n+e)\kappa_1$, and the comparison of capacities of that lemma holds. By
Lemma~\ref{pd-kappa-defects-windows-lift-pairs-source-lift-tail-existence-segment-block-numbers}(1), $\operatorname{ord}(b')=\nu$.

\emph{Step 4: the polynomial $u'$.} Let $G''=G'\,m(A')$, of row degrees $d''$ equal to $d'$ plus $\tau_i$ at $x_i$, $\tau'_i$ at $y_i$
and $\tau_+$ at $x_+$. Put $\zeta_l=q_{u_l}(1)$, $\zeta'_l=q'_{u_l}(1)$ and $\zeta_+=A'((s_{j+1},1),1)$, and
\[
 u'=\prod_l\hat E_l^{u_l}\cdot\hat E_+^{\ell_{j,+}}\cdot G''\cdot\prod_lz_l^{\zeta_l}z_l'^{\zeta'_l}\cdot z_+^{\zeta_+}\cdot\prod_c\xi_c^{\gamma_{j,c}} .
\]
We apply Lemma~\ref{pd-kappa-defects-windows-lift-pairs-source-lift-tail-existence-segment-quotient} to $b'$ and
$\mu+(n-e)\kappa_1$, with $j_+=\ell_{j,+}$, $\gamma=\gamma_j$ and $G=G''$, which has coefficients in $\mathcal R$.
\begin{itemize}
\item (1) $\zeta'_l\leq\zeta_l$ and $\hat b'_{u_l}-o'_{u_l}=\zeta_l-\zeta'_l$ by Step 1, and $a_T(s_{j+1})-o'_+=\zeta_+$ by Step 3.
\item (2) By Definition~\ref{pd-kappa-defects-segment-lattice}(3) and Step 1, the degrees of the letters for $b'$ are
  $w'_i=d''(x_i)$ and $\hat\alpha'_i-\varsigma'_i=d''(y_i)$ for $i\leq\eta$, $o'_{u_l}=d''(x_{u_l})+u_l+\zeta'_l$ and
  $\hat\alpha'_{u_l}=d''(y_{u_l})+u_l+\zeta'_l$, and $o'_+=d''(x_+)+\ell_{j,+}$.
\item (3) $\sum_c\gamma_{j,c}=b_+-\ell_{j,+}$ by
  Definition~\ref{pd-kappa-defects-windows-lift-pairs-source-lift-tail-existence-levels}(2).
\item (4) Let $m\leq\zeta'$, number the letters of $\mathcal P_j$ as $0,\dots,2s$, and let $(\mathcal L,M)$ be one of the pairs
  $(\mathcal L_l,M_l)$ and $(\mathcal L_+,M_+)$ for $b'$ and $\mu+(n-e)\kappa_1$, so that
  $M=\lfloor\operatorname{cap}_{b'}(L)-(\mu+(n-e)\kappa_1)\rfloor$ with $L=u_l-1$ or $L=s$, and $\mathcal L$ is the set of the letters
  $i<J$ for some $J$. With the same $\mathcal L$, $M^\beta=\lfloor\operatorname{cap}_{\beta_j}(L)-(\mu+n\kappa_1)\rfloor$ is the bound of the
  corresponding pair for $\beta_j$, and by hypothesis $G'\in Y(\mathcal L,M^\beta+N_1e)$. Let $v$ and $\pi$ be as in
  Lemma~\ref{pd-kappa-defects-windows-lift-pairs-source-lift-tail-existence-segment-weights}, which applies as $\mathcal K$ has
  characteristic $0$. By its parts (1) and (2), as $v$ vanishes at the covector letters,
  $m(A')\in Y(\mathcal L,\sum_i\tau_iv(x_i)+\tau_+v(x_+))$; by its parts (1) and (3),
  $\prod_lR_l^{\zeta'_l-m_l}\in Y(\mathcal L,\sum_l(\zeta'_l-m_l)\pi(x_{u_l},y_{u_l}))$. The letters $x_{u_l}$ and $y_{u_l}$ are numbered
  $2u_l-2$ and $2u_l-1$, so $\pi(x_{u_l},y_{u_l})=1$ only if $J=2u_l-1$, for at most one $l$, and then $x_+\notin\mathcal L$; so
  $\sum_l\pi(x_{u_l},y_{u_l})+v(x_+)\leq1$. As $\tau_i$, $\tau_+$ and $\zeta'_l$ are at most $Qe$ by Step 1, the sum $V$ of the two
  bounds is at most $sQe+Qe\leq(N_1+1)(k+Qe)$. By part (1) of that lemma and Step 3,
  $G''\prod_lR_l^{\zeta'_l-m_l}\in Y(\mathcal L,M^\beta+N_1e+V)\subset Y(\mathcal L,M)$.
\end{itemize}
So $u'\in Z_j(b';\mu+(n-e)\kappa_1)$.

\emph{Step 5: the identity.} By Definition~\ref{pd-kappa-defects-windows-lift-pairs-model}(5), $\iota_j(b',u')$ is the product of
$\prod_i\prod_tz_{(y_{j,i},1),t}^{r_i(t)}z_{(y_{j,i},0),t}^{r'_i(t)}$, of
$\prod_{i\leq\eta}z_{(y_{j,i},1),1}^{P_iq_i(1)}z_{(y_{j,i},0),1}^{P_iq'_i(1)}$, as $\hat b'_i-w'_i=q_i(1)$ and $\varsigma'_i=q'_i(1)$ for $i\leq\eta$,
and of
\[
 \iota_j(u')=H_j(u,G')\cdot\iota_j(m(A'))\cdot\prod_lz_{(y_{j,u_l},1),1}^{P_{u_l}\zeta_l}z_{(y_{j,u_l},0),1}^{P_{u_l}\zeta'_l}\cdot
 z_{(s_{j+1},1),1}^{\zeta_+}\cdot\prod_cz_{(s_{j+1},0),c+1}^{\gamma_{j,c}},
\]
where
$\iota_j(m(A'))=\prod_i\prod_{t\geq2}z_{(y_{j,i},1),t}^{P_iq_i(t)}z_{(y_{j,i},0),t}^{P_iq'_i(t)}\cdot\prod_{t\geq2}z_{(s_{j+1},1),t}^{A'((s_{j+1},1),t)}$.
Collecting, the exponent at $z_{(y_{j,i},1),t}$ is $r_i(t)+P_iq_i(t)=A'((y_{j,i},1),t)$ and that at $z_{(y_{j,i},0),t}$ is
$A'((y_{j,i},0),t)$, for every $i$ and $t$; those at the variables of $(s_{j+1},1)$ are those of $A'$; and those at the variables of
$(s_{j+1},0)$ are $0$ and the $\gamma_{j,c}$, which are those of $A'$ by (4). So $\iota_j(b',u')=H_j(u,G')\,z^{A'_j}$.
\end{proof}
% Lean: restSegCore_of_num numCore_thr, with the thresholds of numCore_thr (Lemma ..-segment-numbers). segCoreAt_of_num
%   (RestSegCoreOfNum.lean): b' = segmentBlockOfMonomial; Step 1 is hA', hC', hdig' (nb_coarseA, nb_coarseC, nb_digits), hez
%   (nb_ez_le), hy0U and hy0S (block_y_zero_le_of_bound), hroom and hyz (block_y_zero_le_z_zero); Step 2 is hlb (level_bounds),
%   lev_deg, hduW and hduU; Step 3 is hnum (NumCoreAt (ii)) with nb_coarseA_ge, and nb_ord; Step 4 is uTot_mem, with hdeg
%   (chain_deg_U, segment_chain_degree_even, chain_deg_E, dEps), hB (L.hgamma) and hY (ycond_mEps, ycond_pairs, ycond_mul,
%   rowInc_le with pWt_sum_le, ycond_congr, ycond_mono, thr_eq and hthr at capIdx k); Step 5 is segMon_eq and ioBlock_eq.
%   The Lean states SegCoreAt over any discrete valuation ring O with a uniformizer whose square is p, as for the statement.

\begin{lemma}\label{pd-kappa-defects-windows-lift-pairs-source-lift-tail-existence-tail-condition}
For every $C_c\geq0$ there is $T_0>0$ such that for every $C_0\geq T_0$ there is $T_1>0$ such that for all
$C_{E1},C_{E2}\geq T_1$ there is $T_2>0$ with the following property for every $C_H\geq T_2$. In
Situation~\ref{pd-kappa-defects-windows-lift-pairs-source-lift-tail-existence-situation} with these constants, every level
choice for $\beta$ satisfies the tail condition
(Definition~\ref{pd-kappa-defects-windows-lift-pairs-source-lift-tail-existence-prefactorizations}(2)).
\end{lemma}
\leanref{Formalism/SharpRestriction/GrowthBounds/StandardPolynomials/TCoreDef.lean}{Formalism.MainNotions.ColumnSystems.Pf.TCoreAt}
\leanref{Formalism/SharpRestriction/GrowthBounds/StandardPolynomials/TCoreDef.lean}{Formalism.MainNotions.ColumnSystems.Pf.RestTCore}
\leanref{Formalism/SharpRestriction/GrowthBounds/StandardPolynomials/TCGlue.lean}{Formalism.MainNotions.ColumnSystems.Pf.restTCore_of_seg}

\begin{proof}
\emph{The constants.} Let $C_c\geq0$, and let $T_0$, $T_1$ and $T_2$ be the thresholds of
Lemma~\ref{pd-kappa-defects-windows-lift-pairs-source-lift-tail-existence-tail-condition-segment}. Let the constants and the
data of the situation be given, with $C_0\geq T_0$, $C_{E1},C_{E2}\geq T_1$ and $C_H\geq T_2$, let $\ell$ be a level choice for
$\beta$ with base monomial $z^{A_\ell}$, and let $e$, $u=(u_j)$, $G'=(G'_j)$ with row degrees $d'_j$, and $z^{A'}$ be as in
Definition~\ref{pd-kappa-defects-windows-lift-pairs-source-lift-tail-existence-prefactorizations}(2), with
$P=\prod_jH_j(u_j,G'_j)\cdot z^{A'}\in\mathrm{Pol}_T$ and every monomial of $P$ of $\varpi_0$-degree at most $o+e$. We show
$P\in S_e$. If $P=0$, this holds; let $P\neq0$. Write $P_{j,i}=p^{g_j-d_{j,i}}$, $u_{j,l}=\eta_j+l$, and $\bar u_{j,i}$ for the
numbers $\bar u_i$ of Lemma~\ref{pd-kappa-defects-windows-lift-pairs-source-lift-tail-existence-tail-condition-segment} for $j$
and $u_j$.

\emph{Step 1: the segments.} Every element of $T$ is exactly one of the $s_{j+1}$ and the $y_{j,i}$. So
$z^{A'}=\prod_jz^{A'_j}$, with $z^{A'_j}$ as in
Lemma~\ref{pd-kappa-defects-windows-lift-pairs-source-lift-tail-existence-tail-condition-segment}, and
$P=\prod_jH_j(u_j,G'_j)\,z^{A'_j}$.

\emph{Step 2: the row degrees.} By Definition~\ref{pd-kappa-defects-windows-lift-pairs-model}(5), $\iota_j(\hat E_l)$ has degree
$P_{j,u_{j,l}}$ in the variables of each of the letters $(y_{j,u_{j,l}},1)$ and $(y_{j,u_{j,l}},0)$, $\iota_j(\hat E_+)$ has degree
$1$ in those of each of $(s_{j+1},1)$ and $(s_{j+1},0)$, and every monomial of $\iota_j(G'_j)$ has degree $P_{j,i}d'_j(x_i)$ in the
variables of $(y_{j,i},1)$, $P_{j,i}d'_j(y_i)$ in those of $(y_{j,i},0)$, $d'_j(x_+)$ in those of $(s_{j+1},1)$, and $0$ in those of
$(s_{j+1},0)$. So every monomial of $P$ has degree $P_{j,i}(\bar u_{j,i}+d'_j(x_i))+\sum_tA'((y_{j,i},1),t)$ in the variables of
$(y_{j,i},1)$, $P_{j,i}(\bar u_{j,i}+d'_j(y_i))+\sum_tA'((y_{j,i},0),t)$ in those of $(y_{j,i},0)$, and
$\ell_{j,+}+d'_j(x_+)+\sum_tA'((s_{j+1},1),t)$ in those of $(s_{j+1},1)$. As $P\neq0$ lies in $\mathrm{Pol}_T$, these are the
degrees $b_T(y_{j,i})$, $a_T(y_{j,i})$ and $a_T(s_{j+1})$ of these letters (Definition~\ref{pd-kappa-product-space}). These are
the identities (5) of Lemma~\ref{pd-kappa-defects-windows-lift-pairs-source-lift-tail-existence-tail-condition-segment} for
every $j$.

\emph{Step 3: the bounds (1) to (4).} The entries of $A_\ell$ are listed in Step 4 of the proof of
Lemma~\ref{pd-kappa-defects-windows-lift-pairs-source-lift-tail-existence-prefactorization}. Let $1\leq j\leq q$ and
$1\leq i\leq S_j$, and put $y=y_{j,i}$. The hypothesis $\sum_t(A_\ell((y,1),t)\dot-A'((y,1),t))\leq Qe$ gives (1), from its term
$t=1$. The hypothesis $\sum_t(A'((y,1),t)\dot-A_\ell((y,1),t))\leq Qe$ gives (2), from its terms $t\geq2$, where
$A_\ell((y,1),t)=\rho_{j,i}(t)$. The same hypothesis at $(y,0)$ gives (3), from $t=1$, and at $(s_{j+1},1)$, where $A_\ell$
vanishes for $t\geq2$, the first part of (4). The rest of (4) is the hypothesis $A'=A_\ell$ at the letters $(s_{j'},0)$, with the
entries $0$ and $\gamma_{j,c}$ of $A_\ell$ there.

\emph{Step 4: the order.} By Lemma~\ref{pd-kappa-defects-windows-lift-pairs-source-lift-tail-existence-prefactorization-hooks},
every monomial of $\iota_j(\hat E_l)$ has $\varpi_0$-degree $P_{j,u_{j,l}}$ and every monomial of $\iota_j(\hat E_+)$ has
$\varpi_0$-degree $1$. The map $\iota_j$ sends $x_{i,c}$ to a $P_{j,i}$-th power of a variable of $\varpi_0$-degree $1$, $x_{+,c}$ to
a variable of $\varpi_0$-degree $1$, and $y_{i,c}$ to a power of a variable of $\varpi_0$-degree $0$; so every monomial of
$\iota_j(G'_j)$ has $\varpi_0$-degree $\sum_iP_{j,i}d'_j(x_i)+d'_j(x_+)$. Adding the $\varpi_0$-degree of $z^{A'}$, every monomial of
$P$ has $\varpi_0$-degree $\sum_j\big(\nu_j+A'((s_{j+1},0),1)\big)$, where $\nu_j$ is the number $\nu$ of
Lemma~\ref{pd-kappa-defects-windows-lift-pairs-source-lift-tail-existence-tail-condition-segment} for $j$. By Step 3,
$A'((s_{j+1},0),1)=0$. As $P\neq0$, $\sum_j\nu_j\leq o+e$.

\emph{Step 5: the blocks.} By Steps 2 and 3,
Lemma~\ref{pd-kappa-defects-windows-lift-pairs-source-lift-tail-existence-tail-condition-segment} applies to each $j$, with $u_j$,
$G'_j$ and $z^{A'}$. It gives a block $b'_j$ of $S_j(b_T)$ relaxed of slack $(c-n+e)\kappa_1$ with $\operatorname{ord}(b'_j)=\nu_j$,
and $u'_j\in Z_j(b'_j;\mu+(n-e)\kappa_1)$ with $H_j(u_j,G'_j)\,z^{A'_j}=\iota_j(b'_j,u'_j)$. By Step 1, $P=\Phi(b',u')$ for
$b'=(b'_j)$ and $u'=(u'_j)$, and $\sum_j\operatorname{ord}(b'_j)\leq o+e$ by Step 4.

\emph{Step 6.} Each $u'_j$ is an $\mathcal R$-linear combination of generators of $b'_j$ formed with $\mu+(n-e)\kappa_1$, and
$\Phi(b',\cdot)$ is $\mathcal R$-linear in each entry, as $\iota_j(b'_j,\cdot)$ is the algebra map $\iota_j$ followed by
multiplication by a monomial. Expanding, $P$ is an $\mathcal R$-linear combination of symbol families $\Phi(b',f')$ with each
$f'_j$ such a generator. These lie in $S_e=S(\mu+(n-e)\kappa_1,(c-n+e)\kappa_1,o+e)$
(Definition~\ref{pd-kappa-defects-windows-lift-pairs-source-lift-relaxed}(3)), so $P\in S_e$.
\end{proof}
% Lean: restTCore_of_seg (restSegCore_of_num numCore_thr), with the thresholds of RestSegCore. tcoreAt_of_seg (TCGlue.lean):
%   hT0 is the case P = 0; hsplit is Step 1 (tg_mono_split, tgTop_bij in TgProdPlaces.lean); hwh with tg_head_wh and
%   tg_segMon_wh computes weighted degrees for any weight constant on the columns t >= 2 and balanced at each position;
%   hdeg with wLet x, tg_eval_odd and tg_eval_end is Step 2; the five bullets of key, with hS = tg_std_eq and tgStdE for
%   the entries of A_ell, are Step 3; hnu with tgWN and tg_nd_eval is Step 4 (hA0 for A'((s_{j+1}, 0), 1) = 0); key applies
%   SegCoreAt (Step 5); and symFam_span into the span SSe is Step 6.

\begin{lemma}\label{pd-kappa-defects-windows-lift-pairs-source-lift-tail-existence-quotient}
Let $K$ be a field of characteristic $0$ and $N\geq0$, $r\geq1$ integers, fix a lettering of $\{0,\dots,N-1\}$, and use the
notation of Definition~\ref{pd-symbols-path-generators} and the spaces $Y(\mathcal L,M)$ of
Definition~\ref{pd-kappa-defects-segment-lattice}(4) in $K[z]=K[z_{i,c}:0\leq i<N,\ 1\leq c\leq r]$, with
$Y(\mathcal L,+\infty)=K[z]$. Let $0\leq J\leq N$, let $\mathcal L$ be the set of the letters $i<J$, and let
$M\in\mathbf Z\cup\{+\infty\}$. Let $G_u,G_r\in K[z]$ with $G_r\neq0$ of row degrees $d_r$ and $G_uG_r\in Y(\mathcal L,M)$, and let
$c'\geq0$ be an integer with $\sum_id_r(i)\leq c'$, the sum over the covector letters $i$. Then $G_u\in Y(\mathcal L,M+c')$.
\end{lemma}
\leanref{Formalism/SharpRestriction/GrowthBounds/StandardPolynomials/YcondFactor.lean}{Formalism.MainNotions.ColumnSystems.Pf.ycond_factor}

\begin{proof}
If $M=+\infty$ there is nothing to prove. Let $M\in\mathbf Z$, and write $\lambda=\lambda_{\{0,\dots,J-1\}}$
(Definition~\ref{pd-symbols-path-generators}(1)). By Lemma~\ref{pd-kappa-defects-paths-isotypic}(3), a polynomial lies in
$Y(\mathcal L,M')$ if and only if its value of $\lambda$ is at most $M'$, for every integer $M'$. The rule
$\lambda(FF')=\lambda(F)+\lambda(F')$ holds in $K[z]$: the proof of
Lemma~\ref{pd-symbols-residue-count-inputs-segment-clauses}(1) proves it over a field.

We show $\lambda(G_r)\geq-c'$. Otherwise $\lambda(G_r)\leq-c'-1$, so $G_r\in Y(\mathcal L,-c'-1)$, which lies in the span of the
monomials of weight at most $-c'-1$ for $\mathcal L$. As $G_r\neq0$, it has a monomial $z^m$. Its exponents at the variables of
each letter $i$ add up to $d_r(i)$, so its weight for $\mathcal L$ is at least minus the sum of its exponents at the variables
$z_{i,1}$ of the covector letters $i\in\mathcal L$, hence at least $-\sum_id_r(i)\geq-c'$, the sum over the covector letters.
This is a contradiction.

So $\lambda(G_r)$ is an integer $v\geq-c'$. As $G_uG_r\in Y(\mathcal L,M)$, $\lambda(G_u)+v\leq M$, so
$\lambda(G_u)\leq M-v\leq M+c'$, also when $\lambda(G_u)=-\infty$. Hence $G_u\in Y(\mathcal L,M+c')$.
\end{proof}
% Lean: ycond_factor D LL jj M Gu Gr dr hr hne h cov hcov: the case top is ycond_top; otherwise mem_ycond_iff D and the
%   clause lam1S (F * G) = lam1S F + lam1S G of valuationsAdditive_D7 D, the ValuationsAdditive of Lemma
%   pd-symbols-residue-count-inputs-segment-clauses(1) over a field; hlow is the bound lambda(G_r) >= -c' (ycond_le, chWt
%   and the row degrees dr). The Lean states it for every integer jj >= 0, the letters i < jj; jj > N gives the case J = N.

\begin{lemma}\label{pd-kappa-defects-windows-lift-pairs-source-lift-tail-existence-assembly}
In Situation~\ref{pd-kappa-defects-windows-lift-pairs-source-lift-tail-existence-situation}, let $\ell$ be a level choice for
$\beta$, and let $G=(G_j)$ with $G_j\in\mathcal P'_j$ nonzero and lying in $Y(\mathcal L_l,M_l)$ for every $l$ and in
$Y(\mathcal L_+,M_+)$, with the $\mathcal L_l$, $M_l$, $\mathcal L_+$ and $M_+$ of the generators of $\beta_j$. Suppose that $G$
has a prefactorization for $\ell$ and that $\ell$ satisfies the tail condition
(Definition~\ref{pd-kappa-defects-windows-lift-pairs-source-lift-tail-existence-prefactorizations}). Then the symbol family
$\Phi(\beta,W(G))$, with $W(G)=(W_j(G_j))$
(Definition~\ref{pd-kappa-defects-windows-lift-pairs-source-lift-tail-existence-levels}), has a factorization system with
$\mu$, $n$ and $o$.
\end{lemma}
\leanref{Formalism/SharpRestriction/GrowthBounds/StandardPolynomials/ZeroOfMultideg.lean}{Formalism.MainNotions.ColumnSystems.Pf.PreSys.toFacSys}
\leanref{Formalism/SharpRestriction/GrowthBounds/StandardPolynomials/ZeroOfMultideg.lean}{Formalism.MainNotions.ColumnSystems.Pf.PreSys.core}
\leanref{Formalism/SharpRestriction/GrowthBounds/StandardPolynomials/MonomialSumEqProdTensor.lean}{Formalism.MainNotions.ColumnSystems.Pf.PreSys.ident}
\leanref{Formalism/SharpRestriction/GrowthBounds/StandardPolynomials/TCoreDef.lean}{Formalism.MainNotions.ColumnSystems.Pf.LvDat.ioBlock_word}
\leanref{Formalism/SharpRestriction/GrowthBounds/StandardPolynomials/YcondFactor.lean}{Formalism.MainNotions.ColumnSystems.Pf.hom_total}

\begin{proof}
Use the notation of Definition~\ref{pd-kappa-defects-windows-lift-pairs-source-lift-tail-existence-prefactorizations}(1) for
the prefactorization. The field $\mathcal K$ has characteristic $0$, as in the proof of
Lemma~\ref{pd-kappa-defects-digits-chain-dimension}.

\emph{Step 1: the identity.} Let $v$ be a map on the set of the pairs $(j,a)$ with $a\in I_j$, $v^H$ one on the set of the pairs
$(j,l)$, and $v^M$ one on $I_M$, all with values in $\mathbf Z_{\geq0}$, and let $z^M$ be a monomial. As $\iota_j$ is a
$\mathcal K$-algebra map and $H_j(u,G)=\prod_l\iota_j(\hat E_l)^{u_l}\cdot\iota_j(\hat E_+)^{\ell_{j,+}}\cdot\iota_j(G)$,
\[
 \sigma_R\prod_{j,a}\sigma_{j,a}^{v_{j,a}}\prod_{j,l}(\sigma^H_{j,l})^{v^H_{j,l}}\prod_a(z^{\kappa_a})^{v^M_a}\cdot z^M
 =\prod_jH_j\Big(v^H_j,\prod_aK_{j,a}^{v_{j,a}}\Big)\cdot z^{\rho+\sum_{j,a}v_{j,a}\pi_{j,a}+\sum_av^M_a\kappa_a+M}.
\]

\emph{Step 2: the system.} Let $A$ be the disjoint union of the sets of pairs $(j,a)$ with $a\in I_j$, of pairs $(j,l)$, and of
$I_M$. To $(j,a)$ assign the multiplicity $1$, the polynomials $f_{j,a}$ and $\sigma_{j,a}$, the order $o_{j,a}$ and the row
degrees $e_{j,a}$; to $(j,l)$ the multiplicity $\ell_{j,l}$, $f^H_{j,l}$, $\sigma^H_{j,l}$, $o^H_{j,l}$ and $e^H_{j,l}$; to
$a\in I_M$ the multiplicity $m_a$, $f^M_a$, $z^{\kappa_a}$, $o^M_a$ and $e^M_a$; and take $f_R$, $\sigma_R$ and $o_R$. Write $m_a$,
$f_a$, $\sigma_a$, $o_a$ and $d_a$ for the assigned data at $a\in A$. We check the conditions of
Definition~\ref{pd-kappa-defects-windows-lift-pairs-source-lift-tail-systems}(3) for $\sigma=\Phi(\beta,W(G))$.
\begin{itemize}
\item Condition (a) is (a) to (c) of the prefactorization; the only monomial of $z^{\kappa_a}$ has $\varpi_0$-degree
  $d_\varpi(\kappa_a)=o^M_a$.
\item Condition (b) is (d), and condition (c) is (g), as $F_A=f_R\prod_af_a^{m_a}$ is the product $F$ of (g).
\item The first identity of (d) is the last identity of (f). For the second, the definition of $\iota_j(\beta_j,\cdot)$ as a
  monomial times $\iota_j$ gives
  $\iota_j(\beta_j,W_j(G_j))=H_j(\ell_j,G_j)\cdot\iota_j\big(\beta_j,\prod_lz_l^{\hat b_{u_l}\dot-o_{u_l}}z_+^{\alpha_+\dot-o_+}
  \prod_c\xi_c^{\gamma_{j,c}}\big)$, so $\Phi(\beta,W(G))=\prod_jH_j(\ell_j,G_j)\cdot z^{A_\ell}$. Step 1 with $v\equiv1$,
  $v^H_{j,l}=\ell_{j,l}$, $v^M=m$ and $M=0$, together with $\prod_aK_{j,a}=G_j$ and $A_\ell=\rho+\sum_{j,a}\pi_{j,a}+\sum_am_a\kappa_a$,
  gives $\sigma_R\prod_a\sigma_a^{m_a}=\Phi(\beta,W(G))$.
\end{itemize}

\emph{Step 3: condition (e).} Let $u\colon A\to\mathbf Z_{\geq0}$ with $u_a\leq m_a$ for all $a$, let $e$ be an integer with
$1\leq e<n$ and $\sum_au_a\leq e$, and let $z^M$ be a monomial of row degrees $\sum_au_ad_a$ with
$d_\varpi(M)\leq\sum_au_ao_a+e$. Write $u_{j,a}$, $u^H_{j,l}$ and $u^M_a$ for the values of $u$ on the three parts of $A$; so
$u_{j,a}\leq1$. Put $G'_j=\prod_aK_{j,a}^{1-u_{j,a}}$, $G''_j=\prod_aK_{j,a}^{u_{j,a}}$, $u'_{j,l}=\ell_{j,l}-u^H_{j,l}$,
\[
 B=\sum_{j,a}(1-u_{j,a})\pi_{j,a}+\sum_a(m_a-u^M_a)\kappa_a,\qquad B'=\sum_{j,a}u_{j,a}\pi_{j,a}+\sum_au^M_a\kappa_a,
\]
and $A'=\rho+B+M$. By Step 1, $P=\sigma_R\prod_a\sigma_a^{m_a-u_a}\cdot z^M$ equals $\prod_jH_j(u'_j,G'_j)\cdot z^{A'}$, and by (f),
$A_\ell=\rho+B+B'$. We check the hypotheses of the tail condition for $e$, $u'$, $G'=(G'_j)$ and $z^{A'}$; it then gives
$P\in S_e$.
\begin{itemize}
\item $G'_j$ has coefficients in $\mathcal R$, as the $K_{j,a}$ have, and row degrees $\sum_a(1-u_{j,a})d_{j,a}\leq\sum_ad_{j,a}=\delta'_j$.
\item $G'_jG''_j=G_j$ lies in $Y(\mathcal L_l,M_l)$ for every $l$ and in $Y(\mathcal L_+,M_+)$, and $G''_j\neq0$ as $G_j\neq0$.
  $G''_j$ has row degrees $d''_j=\sum_au_{j,a}d_{j,a}$, so $d''_j(\lambda)\leq\sum_au_{j,a}\leq e$ at every letter $\lambda$, as
  $d_{j,a}(\lambda)\leq1$. The covector letters of $\mathcal P_j$ are $y_1,\dots,y_s$, and $s=S_j\leq N_1$ by hypothesis (1) of
  Theorem~\ref{pd-gain-theorem-one}; so the sum of $d''_j$ over the covector letters is at most $N_1e$.
  Lemma~\ref{pd-kappa-defects-windows-lift-pairs-source-lift-tail-existence-quotient} gives
  $G'_j\in Y(\mathcal L_l,M_l+N_1e)$ for every $l$ and $G'_j\in Y(\mathcal L_+,M_++N_1e)$.
\item Let $x$ be a letter, and write $|N|_x=\sum_cN(x,c)$ for the row degree at $x$ of a monomial $z^N$. As $A_\ell-A'=B'-M$,
  $\sum_c(A_\ell(x,c)\dot-A'(x,c))\leq|B'|_x=\sum_{j,a}u_{j,a}|\pi_{j,a}|_x+\sum_au^M_ae^M_a(x)$, using (c) for $|\kappa_a|_x$; by (a)
  and (e) this is at most $(\sum_{j,a}u_{j,a}+\sum_au^M_a)Q\leq Qe$. Likewise
  $\sum_c(A'(x,c)\dot-A_\ell(x,c))\leq|M|_x=\sum_au_ad_a(x)$, which is at most $Qe$ by (e).
\item At a letter $x=(s_{j'},0)$ the $d_a$ vanish by (e), and so do the row degrees of the $\pi_{j,a}$ and of the $\kappa_a$, by
  (a), (c) and (e). So the row degrees of $M$ and of $B'$ at $x$ are $0$, their exponents at the variables of $x$ vanish, and
  $A'(x,c)=A_\ell(x,c)$.
\item Every monomial of $\sigma_R$ has row degrees $e_R$, every monomial of $\sigma_a$ has row degrees $d_a$, and $z^M$ has row
  degrees $\sum_au_ad_a$. So every monomial of $P$ has row degrees $e_R+\sum_a(m_a-u_a)d_a+\sum_au_ad_a=e_R+\sum_am_ad_a$, which
  is $x\mapsto d_x$ by (f); so $P\in\mathrm{Pol}_T$.
\item As $d_\varpi$ is additive, every monomial of $P$ has $\varpi_0$-degree
  $o_R+\sum_a(m_a-u_a)o_a+d_\varpi(M)\leq o_R+\sum_am_ao_a+e=o+e$, by (a) to (d) and (f).
\end{itemize}
\end{proof}
% Lean: PreSys.toFacSys D P hStd hT hG0 hGY hcov. The set A is PreSys.kappa, the sum of (Sigma j, iotaC j), (Sigma j, Fin t)
%   and iotaM, with mult, f, sg, od and dg (SumMulLe.lean); Step 1 is PreSys.ident. Step 2: fac, sgHom (PreSys.sgNd),
%   liftPol, sgPol (PreSys.sgMd), rigA, rigLift, rigInt, clean and total (PreSys.fprod), ordEq (PreSys.ordkappa), and
%   symEq = PreSys.symkappa hStd (LvDat.ioBlock_word and PreSys.ident at 1, lv, multM and 0). Step 3 is PreSys.core: key
%   is PreSys.ident, hSB the decomposition Std = (R + B) + Brem, and hT applied with ju = lv - uH, Gu = prod Kc^(1 - uC),
%   du and A = R + B + M: glIntProd, chain_pol_prod, ycond_factor with Gr = prod Kc^uC and cov = N1 e (hcov,
%   digit_le_one), the two sums (plQ, dgMQ with polM, dgQ, sum_mul_le), zero_of_multideg with dgY, plY and dgMY, and
%   hom_total for mdHom l (degEq, PreSys.degkappa, sum_split) and for ndHom (ordkappa, sum_split). hcov is supplied in
%   facSysAt_of as cov_count with segS j <= N1 from HIext. The Lean states it for any n, mm and beta, with a CharZero
%   instance on KK O.

\begin{lemma}\label{pd-kappa-defects-windows-lift-pairs-source-lift-tail-existence}
For every $C_c\geq0$ there is $T_0>0$ such that for every $C_0\geq T_0$ there is $T_1>0$ such that for all
$C_{E1},C_{E2}\geq T_1$ there is $T_2>0$ with the following property for every $C_H\geq T_2$. Let $p\geq5$ be a prime and
$n'\geq2$ an even integer, and let $m$, $G'$, $N_1$, $c_{cl}$, $D_p$, $c_*$, $T$, $Y$, $Y_b$, $b'$, $b_l$, $b_u$ and an integer
$a$ satisfy hypotheses (1) to (3) of Theorem~\ref{pd-gain-theorem-one} with $C_0$ and $C_H$. Let $c\geq0$ be an integer with
$c\leq C_c(c_*+|T|k\log_pa+|T|c_{cl}+n'c_{cl})$, and let $a_T,b_T\colon\mathbf Z/n'\to\mathbf Z_{\geq0}$ satisfy
$a\leq a_T(s)\leq a+a/(kY)+2J_0$ and $b_l\leq b_T(s)\leq b_u+2J_0$ for $s\in T$, where $J_0=\Lambda b'/k$. Form the numbers of
Definition~\ref{pd-kappa-defects-numbers} with these constants and any $C_m,C_A,C_t,C_\kappa>0$, and let $\mu=0$ or
$\mu=12c$. Let an unramified quadratic Hida setting of degree $n'$ at $p$ with rings $\mathcal R\subset\mathcal K$ and
$\varpi_0$ be given, with the notation of Definitions~\ref{pd-kappa-defects-windows-lift-pairs-model},
\ref{pd-kappa-defects-windows-lift-pairs-source-lift-relaxed}
and~\ref{pd-kappa-defects-windows-lift-pairs-source-lift-tail-systems}. Let $1\leq n\leq c$, let $\beta=(\beta_j)$ with
$\beta_j$ a block of $S_j(b_T)$ relaxed of slack $(c-n)\kappa_1$, and let $f=(f_j)$ with $f_j$ a column generator of some
level of $\beta_j$ formed with $\mu+n\kappa_1$. Then the symbol family $\Phi(\beta,f)$ has a factorization system with $\mu$,
$n$ and $o=\sum_j\operatorname{ord}(\beta_j)$.
\end{lemma}
\leanref{Formalism/SharpRestriction/GrowthBounds/StandardPolynomials/ColGen.lean}{Formalism.MainNotions.ColumnSystems.Pf.FacSysAt}
\leanref{Formalism/SharpRestriction/GrowthBounds/StandardPolynomials/ColGen.lean}{Formalism.MainNotions.ColumnSystems.Pf.FacSysCondition}
\leanref{Formalism/SharpRestriction/GrowthBounds/StandardPolynomials/RestAssumptionsOfThreeInterfaces.lean}{Formalism.MainNotions.ColumnSystems.Pf.facSysAt_of}
\leanref{Formalism/SharpRestriction/GrowthBounds/StandardPolynomials/RestAssumptionsOfThreeInterfaces.lean}{Formalism.MainNotions.ColumnSystems.Pf.restFacSys_of}

\begin{proof}
\emph{The constants.} Let $C_c\geq0$. Let $T_0$ be the largest of the thresholds $T_0$ of
Lemmas~\ref{pd-kappa-defects-windows-lift-pairs-source-lift-tail-existence-columns},
\ref{pd-kappa-defects-windows-lift-pairs-source-lift-tail-existence-prefactorization}
and~\ref{pd-kappa-defects-windows-lift-pairs-source-lift-tail-existence-tail-condition} for $C_c$; for $C_0\geq T_0$ let $T_1$ be
the largest of their thresholds $T_1$; and for $C_{E1},C_{E2}\geq T_1$ let $T_2$ be the largest of their thresholds $T_2$. Let
$C_H\geq T_2$, and let the data, $\mu$, the setting, $n$, $\beta$ and $f$ be as in the statement. They form
Situation~\ref{pd-kappa-defects-windows-lift-pairs-source-lift-tail-existence-situation}, in which the three lemmas apply.

\emph{Step 1: the level choice.} By
Definition~\ref{pd-kappa-defects-windows-lift-pairs-source-lift-tail-systems}(4), for each $j$ there are a level $\ell_j$ of
$\beta_j$, a path $\nu_j$ of rank $m$ for $\delta'_j$ with $(\nu_j)_{2u_l-1}(1)\leq M_l$ for every $l$ and
$(\nu_j)_{2s+1}(1)\leq M_+$, a standard filling $I_j$ of $\nu_j$, and $\gamma_j\in\mathbf Z_{\geq0}^m$ with
$\sum_c\gamma_{j,c}=b_+-\ell_{j,+}$, such that $f_j=W_j(G_j)$ for $G_j=\Psi(\nu_j,I_j)\in\mathcal P'_j$, where $W_j$ is formed with
the level choice $\ell=(\ell_j,\gamma_j)_j$ (Definition~\ref{pd-kappa-defects-windows-lift-pairs-source-lift-tail-existence-levels}).
So $\Phi(\beta,f)=\Phi(\beta,W(G))$ for $G=(G_j)$.

\emph{Step 2: the path generators.} Let $1\leq j\leq q$. The ring $\mathcal R$ is a complete discrete valuation ring with
$p=\varpi_0^2$ a nonzero element of its maximal ideal (Definition~\ref{pd-local-model}(1)). We use Steps 1 and 2 of the proof
of Lemma~\ref{pd-kappa-defects-digits-chain-dimension}, with $\mathcal R'=\mathcal R$, $N=2s+1$, $r=m$, $B=\delta'_j$ and the
conditions $(J,M)$ equal to $(2u_l-1,M_l)$ for $1\leq l\leq s-\eta$ and $(2s+1,M_+)$. By Step 1 there, the path generators of
rank $m$ for $\delta'_j$ are linearly independent over $\mathcal K$, so $G_j\neq0$. For each condition $(J,M)$, $(\nu_j)_J(1)\leq M$,
so $G_j$ lies in the $\mathcal R$-span of the path generators of that step, and by Step 2 there it lies in $Y(\mathcal L,M)$ for the
set $\mathcal L$ of the first $J$ letters. So $G_j\in Y(\mathcal L_l,M_l)$ for every $l$ and $G_j\in Y(\mathcal L_+,M_+)$.

\emph{Step 3: the three inputs.} By Lemma~\ref{pd-kappa-defects-windows-lift-pairs-source-lift-tail-existence-columns}, each
$G_j$ has a column system at $j$. As the $G_j$ are nonzero,
Lemma~\ref{pd-kappa-defects-windows-lift-pairs-source-lift-tail-existence-prefactorization} gives a prefactorization of $G$
for $\ell$. By Lemma~\ref{pd-kappa-defects-windows-lift-pairs-source-lift-tail-existence-tail-condition}, $\ell$ satisfies the
tail condition.

\emph{Step 4.} By Steps 2 and 3, Lemma~\ref{pd-kappa-defects-windows-lift-pairs-source-lift-tail-existence-assembly} applies,
and $\Phi(\beta,f)=\Phi(\beta,W(G))$ has a factorization system with $\mu$, $n$ and $o$.
\end{proof}
% Lean: FacSysAt C0 CYb Cde Ck CD Cb CE1 CE2 Cc in its TransportedAlt form FacSysCondition, supplied as restFacSys_of hA
%   restCols restPre (restTCore_of_seg (restSegCore_of_num numCore_thr)), which is TransportedAlt.mono of facSysAt_of on
%   TransportedAlt.and of RestCols, RestPre and RestTCore (the constants), with hA2 = HyperplaneRank and hA3 =
%   DefiniteRankOne from admitted D. facSysAt_of hA2 hA3 hC hP hT: colGens unfolded by choose gives lv, jp, hOK, y, hy,
%   gamma and hx (Step 1, with L : LvDat and hs'); LvDat.exists_std gives Std; psi_facts (psiK_li, one_cond with hA3,
%   psiK_mem, dvr_fraction_field_charZero) is Step 2; hC (ColsAt), hP (PreAt) and hT (TCoreAt) are Step 3; and
%   PreSys.toFacSys D hStd hT (psi_facts) (cov_count with segS j <= N1, hHI.1.2.2.2.2.1) is Step 4.
% Size of the parts left open: Lemmas pd-kappa-defects-windows-lift-pairs-source-lift-tail-existence-columns-symbol
%   (about 480 lines), -columns-counts (about 500) and -columns-degrees (about 1,450), in the proof of Lemma -columns;
%   Lemmas -columns-numerics and -tail-condition-segment are written out; facSysAt_of with PreSys.toFacSys, about
%   460 lines, is written out here, and Lemma -prefactorization is written out in full.

\begin{lemma}\label{pd-kappa-defects-windows-lift-pairs-source-lift-tail}
For every $C_c\geq0$ there is $T_0>0$ such that for every $C_0\geq T_0$ there is $T_1>0$ such that for all
$C_{E1},C_{E2}\geq T_1$ there is $T_2>0$ with the following property for every $C_H\geq T_2$. Let $p\geq5$ be a prime and
$n'\geq2$ an even integer, and let $m$, $G'$, $N_1$, $c_{cl}$, $D_p$, $c_*$, $T$, $Y$, $Y_b$, $b'$, $b_l$, $b_u$ and an integer
$a$ satisfy hypotheses (1) to (3) of Theorem~\ref{pd-gain-theorem-one} with $C_0$ and $C_H$. Let $c\geq0$ be an integer with
$c\leq C_c(c_*+|T|k\log_pa+|T|c_{cl}+n'c_{cl})$, and let $a_T,b_T\colon\mathbf Z/n'\to\mathbf Z_{\geq0}$ satisfy
$a\leq a_T(s)\leq a+a/(kY)+2J_0$ and $b_l\leq b_T(s)\leq b_u+2J_0$ for $s\in T$, where $J_0=\Lambda b'/k$. Form the numbers of
Definition~\ref{pd-kappa-defects-numbers} with these constants and any $C_m,C_A,C_t,C_\kappa>0$, and let $\mu=0$ or
$\mu=12c$. Let an unramified quadratic Hida setting of degree $n'$ at $p$ with rings $\mathcal R\subset\mathcal K$ and
$\varpi_0$ be given, with the notation of Definitions~\ref{pd-kappa-defects-windows-lift-pairs-model}
and~\ref{pd-kappa-defects-windows-lift-pairs-source-lift-relaxed}. Let $1\leq n\leq c$, let $\beta=(\beta_j)$ with $\beta_j$ a
block of $S_j(b_T)$ relaxed of slack $(c-n)\kappa_1$, let $f=(f_j)$ with $f_j$ a generator of some level of $\beta_j$ formed with
$\mu+n\kappa_1$, and put $o=\sum_j\operatorname{ord}(\beta_j)$. Then there are $F\in\mathrm{Pol}_T$, clean of order $o$, and
$G_\epsilon\in S(\mu+(n-\epsilon)\kappa_1,(c-n+\epsilon)\kappa_1,o+\epsilon)$ for $1\leq\epsilon<n$, such that
\[
 \Sigma_1F-\Phi(\beta,f)-\sum_{\epsilon=1}^{n-1}\varpi_0^\epsilon G_\epsilon\in\varpi_0^n\mathcal R[z].
\]
\end{lemma}
\leanref{Formalism/SharpRestriction/GrowthBounds/CoarseMultiplicity/CoarseMultiplicityAtIdentity.lean}{Formalism.MainNotions.ColumnSystems.Pf.RawTailBody}
\leanref{Formalism/SharpRestriction/GrowthBounds/CoarseMultiplicity/CoarseMultiplicityAtIdentity.lean}{Formalism.MainNotions.ColumnSystems.Pf.RawTailAt1}
\leanref{Formalism/SharpRestriction/GrowthBounds/CoarseMultiplicity/CoarseMultiplicityAtIdentity.lean}{Formalism.MainNotions.ColumnSystems.Pf.RestRawTail1}
\leanref{Formalism/SharpRestriction/GrowthBounds/CoarseMultiplicity/RestRawTail1Of.lean}{Formalism.MainNotions.ColumnSystems.Pf.restRawTail1_of}
\leanref{Formalism/SharpRestriction/GrowthBounds/CoarseMultiplicity/RestRawTail1Of.lean}{Formalism.MainNotions.ColumnSystems.Pf.AsmRTAt}
\leanref{Formalism/SharpRestriction/GrowthBounds/StandardPolynomials/LatticeLeSpanPsiImage.lean}{Formalism.MainNotions.ColumnSystems.Pf.asmRTAt_of_facSys}
\leanref{Formalism/SharpRestriction/GrowthBounds/StandardPolynomials/LatticeLeSpanPsiImage.lean}{Formalism.MainNotions.ColumnSystems.Pf.restAsmRT_of_facSys}

Alpöge proves this lemma from a version with two further hypotheses, which hold and which the proof below, Alpöge's, does not
use. Let $1\leq j\leq q$, $1\leq i\leq S_j$, $d=d_{j,i}$ and $2\leq c\leq k$. The first hypothesis concerns the polynomials
$\sum_{c'=1}^kz_{(s_j,1),c'}^{p^d}z_{(y_{j,i},0),c'}$ and $z_{(y_{j,i},1),c}z_{(s_j,1),1}^{p^d}-z_{(y_{j,i},1),1}z_{(s_j,1),c}^{p^d}$:
their substitutions $\Sigma_g$ lie in $\varpi_0\mathcal R[z]$ for every $g\in GL_k(W)$, and their substitutions $\Sigma_1$ are
given by formulas that follow from the definition of $\Sigma_1$. The second concerns
$D_c=\prod_{j'}z_{(s_{j'},1),1}^{p^{g_{j'}}-1}-\prod_{j'}z_{(s_{j'},1),c}^{p^{g_{j'}}-1}$: the substitution $\Sigma_1(D_c)$ is given by
such a formula, and $\Sigma_g\big(z_{(s_{j_0},1),1}\prod_{c'=2}^kD_{c'}\big)\in\varpi_0\mathcal R[z]$ for every $j_0$ and every
$g\in GL_k(W)$ whose plane differs from that of $1$.
% Lean: AsmRTAt = CoarseMultiplicityBoundAtIdentity -> LocalBoundAtIdentity -> RawTailAt1; the polynomials are kSkey,
%   kWkey (c = c.succ), locC and locR j0 (RestRawTail1Of.lean). restRawTail1_of h3 h4 hA is TransportedAlt.mono
%   (fun ... h => h h3 h4) hA, and asmRTAt_of_facSys ignores both hypotheses.

\begin{proof}
\emph{The constants.} Let $C_c\geq0$, and let $T_0$, $T_1$ and $T_2$ be the thresholds of
Lemma~\ref{pd-kappa-defects-windows-lift-pairs-source-lift-tail-existence}. Let $C_H\geq T_2$, and let the data, $\mu$, the
setting, $n$, $\beta$ and $f$ be as in the statement. By
Definition~\ref{pd-kappa-defects-windows-lift-pairs-source-lift-tail-systems}(2), the claim is that $\Phi(\beta,f)$ lies in the
lift module $\mathcal L(\mu,n,o)$, which is an $\mathcal R$-module by
Lemma~\ref{pd-kappa-defects-windows-lift-pairs-source-lift-tail-factorization}(1).

\emph{Step 1: column generators.} By Lemma~\ref{pd-kappa-defects-windows-lift-pairs-source-lift-tail-columns} with
$\mu'=\mu+n\kappa_1$, each $f_j$ is an $\mathcal R$-linear combination of column generators of $\beta_j$ formed with
$\mu+n\kappa_1$. The symbol family $\Phi(\beta,f)=\prod_j\iota_j(\beta_j,f_j)$ is $\mathcal R$-linear in each $f_j$, as
$\iota_j(\beta_j,\cdot)$ is the algebra map $\iota_j$ followed by multiplication by a monomial. Expanding, $\Phi(\beta,f)$ is an
$\mathcal R$-linear combination of symbol families $\Phi(\beta,f')$ with each $f'_j$ a column generator of $\beta_j$ formed with
$\mu+n\kappa_1$. So it suffices to show $\Phi(\beta,f')\in\mathcal L(\mu,n,o)$ for these.

\emph{Step 2: factorization.} For such $f'$, Lemma~\ref{pd-kappa-defects-windows-lift-pairs-source-lift-tail-existence} gives a
factorization system for $\Phi(\beta,f')$ with $\mu$, $n$ and $o$, and
Lemma~\ref{pd-kappa-defects-windows-lift-pairs-source-lift-tail-factorization}(2) gives $\Phi(\beta,f')\in\mathcal L(\mu,n,o)$.
\end{proof}
% Lean: RawTailAt1 C0 CYb Cde Ck CD Cb CE1 CE2 Cc, that is RawTailBody C0 Cde CE1 CE2 X O S.w mm 1, for any discrete
%   valuation ring O with a uniformizer w, w^2 = p, and any algebra structure of WF p n' on O. RestRawTail1 is
%   restRawTail1_of h3 h4 hA (RestRawTail1Of.lean) with hA = rest_assumptions_of_three_interfaces hD restCols restPre
%   (restTCore_of_seg (restSegCore_of_num numCore_thr)) = restAsmRT_of_facSys hD (restFacSys_of ...), as assembled in
%   tensor_closed (Development/AbelianVarietyTheory/Proofs/GlobalInterpolation/TensorOfPairLift.lean). The proof is
%   asmRTAt_of_facSys hA2 hA3 (LatticeLeSpanPsiImage.lean): hmem is symFam_span with G j = colGens (Step 1, with
%   genSet_le for each s j, Lemma pd-kappa-defects-windows-lift-pairs-source-lift-tail-columns) and, for symbols in
%   colGens, FacSys.liftable (Lemma pd-kappa-defects-windows-lift-pairs-source-lift-tail-factorization(2)) of the FacSys
%   given by FacSysAt (Lemma pd-kappa-defects-windows-lift-pairs-source-lift-tail-existence); the membership in liftSub
%   is the conclusion. The thresholds are those of FacSysCondition. AsmRTAt also assumes CoarseMultiplicityBoundAtIdentity
%   and LocalBoundAtIdentity, which restRawTail1_of supplies as h3 = third_key_of_order key3Ord (the substitutions of the
%   keys kSkey and kWkey; KeyOne.lean, KeyMinorDvd.lean) and h4 = local_fourth_of_order (loc4Ord_of (loc4Div_of loc4Res))
%   (locC and locR; Loc4*.lean, PsiRes0Ne.lean); asmRTAt_of_facSys does not use them (about 4,200 lines).
% Size of the part left open: Lemmas pd-kappa-defects-windows-lift-pairs-source-lift-tail-existence-columns-numerics,
%   -columns-clauses and -tail-condition-segment, about 5,400 lines, under the proof of Lemma
%   pd-kappa-defects-windows-lift-pairs-source-lift-tail-existence.

\begin{lemma}\label{pd-kappa-defects-windows-lift-pairs-source-lift}
For every $C_c\geq0$ there is $T_0>0$ such that for every $C_0\geq T_0$ there is $T_1>0$ such that for all
$C_{E1},C_{E2}\geq T_1$ there is $T_2>0$ with the following property for every $C_H\geq T_2$. Let $p\geq5$ be a prime and
$n'\geq2$ an even integer, and let $m$, $G'$, $N_1$, $c_{cl}$, $D_p$, $c_*$, $T$, $Y$, $Y_b$, $b'$, $b_l$, $b_u$ and an integer
$a$ satisfy hypotheses (1) to (3) of Theorem~\ref{pd-gain-theorem-one} with $C_0$ and $C_H$. Let $c\geq0$ be an integer with
$c\leq C_c(c_*+|T|k\log_pa+|T|c_{cl}+n'c_{cl})$, and let $a_T,b_T\colon\mathbf Z/n'\to\mathbf Z_{\geq0}$ satisfy
$a\leq a_T(s)\leq a+a/(kY)+2J_0$ and $b_l\leq b_T(s)\leq b_u+2J_0$ for $s\in T$, where $J_0=\Lambda b'/k$. Form the numbers of
Definition~\ref{pd-kappa-defects-numbers} with these constants and any $C_m,C_A,C_t,C_\kappa>0$, and let $\mu=0$ or
$\mu=12c$. Let an unramified quadratic Hida setting of degree $n'$ at $p$ with rings $\mathcal R\subset\mathcal K$ and
$\varpi_0$ be given, with the notation of Definition~\ref{pd-kappa-defects-windows-lift-pairs-model}. Let $\beta=(\beta_j)$ be a
tuple with $\beta_j$ a kept block of $S_j(b_T)$, $f=(f_j)$ with $f_j\in Z_j(\beta_j;13kN_1Qc+\mu)$, and
$o=\sum_j\operatorname{ord}(\beta_j)$ (Lemma~\ref{pd-kappa-defects-windows-lift-pairs-source-order}). Then there is
$F\in\mathrm{Pol}_T$ such that:
\begin{enumerate}
\item $\varpi_0^oF\in\mathcal R[z]$;
\item $\Sigma_1F-\Phi(\beta,f)\in\varpi_0^c\mathcal R[z]$;
\item $\Sigma_gF\in\varpi_0^c\mathcal R[z]$ for every $g\in GL_k(W)$ whose plane differs from that of $1$
  (Lemma~\ref{pd-symbols-edge-basis}, with $r=1$).
\end{enumerate}
\end{lemma}
\leanref{Formalism/SharpRestriction/GrowthBounds/CoarseMultiplicity/CoarseMultiplicityAtIdentity.lean}{Formalism.MainNotions.ColumnSystems.Pf.CoarseMultiplicityAtIdentity}
\leanref{Formalism/SharpRestriction/GrowthBounds/CoarseMultiplicity/Lemma44IdentityOfRawTail.lean}{Formalism.MainNotions.ColumnSystems.Pf.lemma44_identity_of_raw_tail}
\leanref{Formalism/SharpRestriction/GrowthBounds/CoarseMultiplicity/Lemma44IdentityOfRawTail.lean}{Formalism.MainNotions.ColumnSystems.Pf.lemma44_at_body}
\leanref{Formalism/SharpRestriction/GrowthBounds/CoarseMultiplicity/Lemma44IdentityOfRawTail.lean}{Formalism.MainNotions.ColumnSystems.Pf.t44}
\leanref{Formalism/SharpRestriction/GrowthBounds/CoarseMultiplicity/Lemma44IdentityOfRawTail.lean}{Formalism.MainNotions.ColumnSystems.Pf.symFam_span}

\begin{proof}
\emph{The constants.} Let $C_c\geq0$, and let $T_0$, $T_1$ and $T_2$ be the largest of the thresholds of
Lemma~\ref{pd-kappa-defects-windows-lift-pairs-source-lift-tail} and of $1$, $1$ and $1000(C_{E1}+C_{E2}+1)(C_c+1)$, those
of Lemma~\ref{pd-kappa-defects-windows-lift-pairs-source-lift-kept}. Let $C_H\geq T_2$, and let the data, $\mu$, the setting,
$\beta$ and $f$ be as in the statement. Use the notation of
Definition~\ref{pd-kappa-defects-windows-lift-pairs-source-lift-relaxed}, with $\kappa_1=13kN_1Q$, so that
$13kN_1Qc+\mu=\mu+c\kappa_1$.

\emph{Step 1: the symbol family.} By Lemma~\ref{pd-kappa-defects-windows-lift-pairs-source-lift-kept}, every $\beta_j$ is
relaxed of slack $0$. Each $f_j$ is an $\mathcal R$-linear combination of generators of $\beta_j$ formed with $\mu+c\kappa_1$,
and $\Phi(\beta,f)=\prod_j\iota_j(\beta_j,f_j)$ is $\mathcal R$-linear in each $f_j$, as $\iota_j(\beta_j,\cdot)$ is the algebra map
$\iota_j$ followed by multiplication by a monomial. Expanding, $\Phi(\beta,f)$ is an $\mathcal R$-linear combination of symbol
families $\Phi(\beta,f')$ with each $f'_j$ a generator of $\beta_j$ formed with $\mu+c\kappa_1$. So
$\Phi(\beta,f)\in S(\mu+c\kappa_1,0,o)$.

\emph{Step 2: the induction.} Claim: for $0\leq n\leq c$, every integer $o'\geq0$ and every
$\sigma\in S(\mu+n\kappa_1,(c-n)\kappa_1,o')$, there is $F$ clean of order $o'$ with $\Sigma_1F-\sigma\in\varpi_0^n\mathcal R[z]$.
The proof is by induction on $n$. The polynomials clean of order $o'$ form an $\mathcal R$-module, as $\mathrm{Pol}_T$,
$\mathcal R[z]$ and $\varpi_0^c\mathcal R[z]$ are $\mathcal R$-modules and the $\Sigma_g$ are $\mathcal K$-linear. So the $\sigma$
for which such an $F$ exists form an $\mathcal R$-module, and it suffices to treat the spanning families
$\sigma=\Phi(\beta',f')$ of Definition~\ref{pd-kappa-defects-windows-lift-pairs-source-lift-relaxed}(3).
\begin{itemize}
\item Let $n=0$, and take $F=0$. Each $f'_j$ is a generator, so lies in $Z_j(\beta'_j;\mu)$, and
  Lemma~\ref{pd-kappa-defects-windows-lift-pairs-source-symbols}(4) gives $\sigma\in\mathcal R[z]$; so
  $\Sigma_10-\sigma\in\varpi_0^0\mathcal R[z]$.
\item Let $1\leq n\leq c$. Here each $\beta'_j$ is relaxed of slack $(c-n)\kappa_1$, $o_1=\sum_j\operatorname{ord}(\beta'_j)\leq o'$,
  and each $f'_j$ is a generator formed with $\mu+n\kappa_1$. Lemma~\ref{pd-kappa-defects-windows-lift-pairs-source-lift-tail}
  gives $F_0$ clean of order $o_1$ and $G_\epsilon\in S(\mu+(n-\epsilon)\kappa_1,(c-(n-\epsilon))\kappa_1,o_1+\epsilon)$ for
  $1\leq\epsilon<n$ with $\Sigma_1F_0-\sigma-\sum_\epsilon\varpi_0^\epsilon G_\epsilon\in\varpi_0^n\mathcal R[z]$. As $n-\epsilon<n$, the
  induction hypothesis gives $F_\epsilon$ clean of order $o_1+\epsilon$ with
  $\Sigma_1F_\epsilon-G_\epsilon\in\varpi_0^{n-\epsilon}\mathcal R[z]$. Then $\varpi_0^\epsilon F_\epsilon$ is clean of order $o_1$, as
  $\varpi_0^{o_1}\varpi_0^\epsilon F_\epsilon\in\mathcal R[z]$ and $\Sigma_g(\varpi_0^\epsilon F_\epsilon)=\varpi_0^\epsilon\Sigma_gF_\epsilon$. Put
  $F=F_0-\sum_\epsilon\varpi_0^\epsilon F_\epsilon$. It is clean of order $o_1$, hence of order $o'$, as
  $\varpi_0^{o'}F=\varpi_0^{o'-o_1}\varpi_0^{o_1}F$. Finally
  \[
   \Sigma_1F-\sigma=\Big(\Sigma_1F_0-\sigma-\sum_\epsilon\varpi_0^\epsilon G_\epsilon\Big)
   -\sum_\epsilon\varpi_0^\epsilon\big(\Sigma_1F_\epsilon-G_\epsilon\big)\in\varpi_0^n\mathcal R[z].
  \]
\end{itemize}

\emph{Step 3.} Step 2 with $n=c$, $o'=o$ and $\sigma=\Phi(\beta,f)$, which lies in $S(\mu+c\kappa_1,0,o)$ by Step 1, gives
$F$ clean of order $o$ with $\Sigma_1F-\Phi(\beta,f)\in\varpi_0^c\mathcal R[z]$. As $F$ is clean of order $o$, it lies in
$\mathrm{Pol}_T$ and satisfies (1) and (3), and the congruence is (2).
\end{proof}
% Lean: lemma44_identity_of_raw_tail hR hK = lemma44_at_body hK X O S hD hH hc mm 1 hR, with hR = RawTailAt1 (Lemma
%   pd-kappa-defects-windows-lift-pairs-source-lift-tail, from RestRawTail1) and hK = KeptBLAt (Lemma
%   pd-kappa-defects-windows-lift-pairs-source-lift-kept, from restKeptBL); the thresholds are those of RestRawTail1 and
%   restKeptBL. Step 1 is hmar and symFam_span (UpperTheory.Reformulated.B.SymmetricFamilySpan.symFam_span) into SSg;
%   Step 2 is t44 (strong induction on n; the submodule M with cleanSub and wjSub, ssg_int with symIntAt for n = 0, and
%   clean_mono, clean_smul_pow and wj_smul_pow for n >= 1); Step 3 is the end of lemma44_at_body. The Lean states it for
%   any discrete valuation ring O with a uniformizer w, w^2 = p, and any algebra structure of the Witt vectors on O; it
%   is applied with the one of the setting.
% Size of the part left open: Lemmas pd-kappa-defects-windows-lift-pairs-source-lift-tail-existence-columns-numerics,
%   -columns-clauses and -tail-condition-segment, about 5,400 lines, under the proof of Lemma
%   pd-kappa-defects-windows-lift-pairs-source-lift-tail-existence.

\begin{lemma}\label{pd-kappa-defects-windows-lift-pairs-source}
For every $C_c\geq0$ there is $T_0>0$ such that for every $C_0\geq T_0$ there is $T_1>0$ such that for
all $C_{E1},C_{E2}\geq T_1$ there is $T_2>0$ with the following property for every $C_H\geq T_2$ and every $C_t\geq30$.
Let $p\geq5$ be a prime and $n'\geq2$ an even integer, and let $m$, $G'$, $N_1$, $c_{cl}$, $D_p$, $c_*$, $T$, $Y$, $Y_b$, $b'$,
$b_l$, $b_u$ and an integer $a$ satisfy hypotheses (1) to (3) of Theorem~\ref{pd-gain-theorem-one} with $C_0$ and $C_H$.
Let $c\geq1$ be an integer with $c\leq C_c(c_*+|T|k\log_pa+|T|c_{cl}+n'c_{cl})$, and let
$a_T,b_T\colon\mathbf Z/n'\to\mathbf Z_{\geq0}$ satisfy $a\leq a_T(s)\leq a+a/(kY)+2J_0$ and
$b_l\leq b_T(s)\leq b_u+2J_0$ for $s\in T$, where $J_0=\Lambda b'/k$. Form the numbers of
Definition~\ref{pd-kappa-defects-numbers} with these constants and any $C_m,C_A,C_\kappa>0$, and put $t=\lfloor t_0\rfloor$ and
$\mu_0=13kN_1Qc+12c$. Let an unramified quadratic Hida setting of degree $n'$ at $p$ with rings
$\mathcal R\subset\mathcal K$ and $\varpi_0$ be given, with the notation of
Definition~\ref{pd-kappa-defects-windows-lift-pairs-model}. Let $h$ be a $k\times k$ matrix over $Q$, $\Theta$ a flat theta
map for $h$, and $e$ an edge with $e=e_h(1)$. Let $\beta=(\beta_j)$ be a tuple with $\beta_j$ a kept block of $S_j(b_T)$, for
$\alpha_i=a_T(y_{j,i})$ and $r=m$ (Definition~\ref{pd-kappa-defects-segment-lattice}), and let $f=(f_j)$ with
$f_j\in Z_j(\beta_j;\mu_0)$. Then there is $F\in\mathrm{Pol}_T$ with $\Theta F\in\mathcal C^\times_{t,c}(e)$ and
$\Sigma_1F-\Phi(\beta,f)\in\varpi_0^c\mathcal R[z]$.
\end{lemma}
\leanref{Formalism/SharpRestriction/GrowthBounds/CoarseMultiplicity/PairedLiftAtOfIdentity.lean}{Formalism.MainNotions.ColumnSystems.Pf.SrcHalfAt1}
\leanref{Formalism/SharpRestriction/GrowthBounds/CoarseMultiplicity/PairedLiftAtOfIdentity.lean}{Formalism.MainNotions.ColumnSystems.Pf.srcHalf1_of}

\begin{proof}
\emph{The constants.} Let $C_c\geq0$, and let $T_0$, $T_1$ and $T_2$ be the largest of the thresholds of
Lemma~\ref{pd-kappa-defects-windows-lift-pairs-source-lift} and of
Lemma~\ref{pd-kappa-defects-windows-lift-pairs-source-order}, the latter being $1$, $1$ and
$1000(C_{E1}+C_{E2}+1)(C_c+1)$. Let $C_H\geq T_2$ and $C_t\geq30$, and let the data, the setting, $h$, $\Theta$, $e$, $\beta$ and
$f$ be as in the statement. Put $o=\sum_j\operatorname{ord}(\beta_j)$, and let $L''=L_h$ be the source of $e$.

\emph{The polynomial.} Lemma~\ref{pd-kappa-defects-windows-lift-pairs-source-lift} with $\mu=12c$, for which
$13kN_1Qc+\mu=\mu_0$, gives $F\in\mathrm{Pol}_T$ with $\varpi_0^oF\in\mathcal R[z]$,
$\Sigma_1F-\Phi(\beta,f)\in\varpi_0^c\mathcal R[z]$, and $\Sigma_gF\in\varpi_0^c\mathcal R[z]$ for every $g\in GL_k(W)$ whose
plane differs from that of $1$. The second property is the congruence of the statement. It remains to show
$\Theta F\in\mathcal C^\times_{t,c}(e)$, that is, the three conditions of Definition~\ref{pd-kappa-split-rank}(1).
\begin{itemize}
\item As $\varpi_0^oF\in\mathcal R[z]$, property (a) of the flat theta map $\Theta$ gives
  $\Theta(\varpi_0^oF)\in H^\times(L'')$, so $\Theta F\in\varpi_0^{-o}H^\times(L'')$. By
  Lemma~\ref{pd-kappa-defects-windows-lift-pairs-source-order}(3), $o\leq t_0$; as $o$ is an integer, $o\leq t$. So
  $\Theta F\in\varpi_0^{-t}H^\times(L'')$.
\item By Lemma~\ref{pd-kappa-defects-windows-lift-pairs-source-symbols}(4), $\Phi(\beta,f)\in\mathcal R[z]$, so
  $\Sigma_1F\in\mathcal R[z]$. Property (b) with $g=1$ and $i=0$ gives $\Theta F\in R^\times(e_h(1))=R^\times(e)$.
\item Let $e'\neq e$ be an edge with source $L''$. By property (c), $e'=e_h(g)$ for some $g\in GL_k(W)$. By
  Lemma~\ref{pd-symbols-edge-basis}(3), with $\gamma=h$ and $r=1$, the plane of $g$ differs from that of $1$, as
  $e_h(g)\neq e_h(1)$. So $\Sigma_gF\in\varpi_0^c\mathcal R[z]$, and property (b) with $i=c$ gives
  $\Theta F\in\varpi_0^cR^\times(e')$.
\end{itemize}
\end{proof}
% Lean: paired_lifts_of_identity uses SrcHalfAt1 C0 ... CE1 CE2 13 Ct Cc = srcHalf1_of (le_refl 13) hCt hstd
%   (lemma44_identity_of_raw_tail hR hK) monoAt planeAt (symIntAt C0 CE1 CE2), with hstd = thr_std and hR, hK
%   from RestRawTail1 and restKeptBL; the thresholds are those of these three. srcHalf1_of: hmar (mu_0 <= m0 13, an
%   equality for Cm = 13) and monoAt (the symbol modules shrink as mu grows) for the symbols; h44 =
%   CoarseMultiplicityAtIdentity at mm = 12 c for F (Lemma pd-kappa-defects-windows-lift-pairs-source-lift); hot from
%   ord_all; the first clause from (a) of IsFlatTheta (h1); the second from hint (symIntAt and hcong) and (b) at i = 0
%   (h2, scale_one_mem); the third from (c) (h3), planeAt and (b) at i = c.

\subsection{Split families from segment lattices}\label{subsection-pd-split-families-segments}

\begin{lemma}\label{pd-kappa-defects-windows-lift-pairs-form-numbers}
Let $p\geq5$ be a prime, $(W,Q,\sigma,\iota)$ a Hida setting of degree $n'$ at $p$ with rings $\mathcal R\subset\mathcal K$ and
$\varpi_0$, $m\geq1$, $T\subset\mathbf Z/n'\mathbf Z$ with $q\geq2$ even elements, and
$a_T,b_T\colon\mathbf Z/n'\mathbf Z\to\mathbf Z_{\geq0}$, with the notation of
Definition~\ref{pd-kappa-defects-windows-lift-pairs-model} and of Definition~\ref{pd-kappa-defects-numbers}, and
$\lfloor\log_p0\rfloor=0$. For $1\leq j\leq q$ put
\[
 \lambda_j=\sum_{i=1}^{S_j}\big(\lfloor\log_pa_T(y_{j,i})\rfloor+\lfloor\log_pb_T(y_{j,i})\rfloor\big)
 +\lfloor\log_pa_T(s_{j+1})\rfloor+\lfloor\log_pb_T(s_{j+1})\rfloor,
\]
$\varphi_j=a_T(s_{j+1})!\,b_T(s_{j+1})!\prod_{i=1}^{S_j}a_T(y_{j,i})!\,b_T(y_{j,i})!$ and $z_j=p^{k\lambda_j}/\varphi_j$.
\begin{enumerate}
\item $z_T=\prod_{j=1}^qz_j$.
\item If $i\geq0$ is an integer, $P\in\varpi_0^i\mathcal R[z]\cap\mathrm{Pol}_T$ and $P'\in\mathcal R[z]$, then
  $z_T[P,P']\in\varpi_0^i\mathcal R$.
\item If $a\geq1$ is a real number with $a_T(s)\leq a^2$ and $b_T(s)\leq a^2$ for every $s\in T$, then
  $\sum_j(m+2)\lambda_j\leq4\kappa_0$, where $\kappa_0=2|T|k\log_pa$.
\end{enumerate}
\end{lemma}
\leanref{Formalism/SharpRestriction/GrowthBounds/CoarseMultiplicity/RestBlockForm.lean}{Formalism.MainNotions.ColumnSystems.Pf.lgSeg}
\leanref{Formalism/SharpRestriction/GrowthBounds/CoarseMultiplicity/RestBlockForm.lean}{Formalism.MainNotions.ColumnSystems.Pf.facSeg}
\leanref{Formalism/SharpRestriction/GrowthBounds/CoarseMultiplicity/RestBlockForm.lean}{Formalism.MainNotions.ColumnSystems.Pf.zSeg}
\leanref{Formalism/SharpRestriction/GrowthBounds/CoarseMultiplicity/LogLeOfLeSq.lean}{Formalism.MainNotions.ColumnSystems.Pf.zProdAt}
\leanref{Formalism/SharpRestriction/GrowthBounds/CoarseMultiplicity/FactorialExponentEqProd.lean}{Formalism.MainNotions.ColumnSystems.Pf.intAt}
\leanref{Formalism/SharpRestriction/GrowthBounds/CoarseMultiplicity/LogLeOfLeSq.lean}{Formalism.MainNotions.ColumnSystems.Pf.numAt}

\begin{proof}
Every element of $T$ is exactly one of the $s_{j+1}$ and the $y_{j,i}$, and the two letters $(s,1)$ and $(s,0)$ at $s\in T$ have
the degrees $a_T(s)$ and $b_T(s)$, in some order (Definition~\ref{pd-kappa-product-space}).

(1) Hence $\kappa_Z=k\sum_x\lfloor\log_pd_x\rfloor=k\sum_{s\in T}\big(\lfloor\log_pa_T(s)\rfloor+\lfloor\log_pb_T(s)\rfloor\big)
=k\sum_j\lambda_j$ and $\prod_xd_x!=\prod_{s\in T}a_T(s)!\,b_T(s)!=\prod_j\varphi_j$, so $z_T=\prod_jz_j$.

(2) As in the proof of Lemma~\ref{pd-kappa-defects-windows-lift-pairs-flat-forms},
Lemma~\ref{pd-symbols-flat-real-multinomial} applies with the letters as indices, and its number $p^{\kappa_F}/\prod_xd_x!$ is
$z_T$. Write $P=\varpi_0^iP_0$ with $P_0\in\mathcal R[z]$; then $P_0$ has the row degrees $d_x$, as $P$ has. By part (2) of
that lemma $z_T[P_0,P']\in\mathcal R$, so $z_T[P,P']=\varpi_0^iz_T[P_0,P']\in\varpi_0^i\mathcal R$.

(3) For an integer $0\leq d\leq a^2$, $\lfloor\log_pd\rfloor\leq2\log_pa$: for $d=0$ the left side is $0\leq2\log_pa$, and
otherwise $\lfloor\log_pd\rfloor\leq\log_pd\leq\log_pa^2$. By the partition used in (1),
$\sum_j\lambda_j=\sum_{s\in T}\big(\lfloor\log_pa_T(s)\rfloor+\lfloor\log_pb_T(s)\rfloor\big)\leq4|T|\log_pa$. As $m+2\leq2k$
and $\log_pa\geq0$, $\sum_j(m+2)\lambda_j\leq2k\cdot4|T|\log_pa=4\kappa_0$.
\end{proof}
% Lean: lgSeg, facSeg and zSeg are lambda_j, phi_j and z_j. (1) is zProdAt (kapZ_sum, facSeg_prod, natPart and
%   natPartSum for the partition of TT into the segments, ft_prod_letters); (2) is intAt (tensor_factorial_integral with
%   fact_multinomial_le and pow_div_mem, the Lean's form of Lemma pd-symbols-flat-real-multinomial with the factor w^i);
%   (3) is numAt with log_le_of_le_sq, where the Lean takes a_T(s), b_T(s) <= a^2 and a >= 3 from Std (winA, winB,
%   loga).

\begin{definition}\label{pd-kappa-defects-windows-lift-pairs-form-block-data}
Let $p\geq5$ be a prime, $(W,Q,\sigma,\iota)$ a Hida setting of degree $n'$ at $p$ with rings $\mathcal R\subset\mathcal K$ and
$\varpi_0$, $m\geq1$, $T\subset\mathbf Z/n'\mathbf Z$ with $q\geq2$ even elements, and
$a_T,b_T\colon\mathbf Z/n'\mathbf Z\to\mathbf Z_{\geq0}$; let the data and constants of
Definition~\ref{pd-kappa-defects-numbers} be given, and use the notation of
Definition~\ref{pd-kappa-defects-windows-lift-pairs-model}. Let $1\leq j\leq q$ and let $b$ be a block of $S_j(b_T)$, for
$\alpha_i=a_T(y_{j,i})$ and $r=m$. Write $s=S_j$, $\eta=\eta_j$, $P_i=p^{g_j-d_{j,i}}$, $\alpha_+=a_T(s_{j+1})$ and
$b_+=b_T(s_{j+1})$, and use the entries and numbers of $b$ of Definition~\ref{pd-kappa-defects-segment-lattice}(1) and the
degrees $\delta$ of its part~(3). Let $V$ be the set of variables of $\mathcal P_j$. For an exponent vector
$\nu\colon V\to\mathbf Z_{\geq0}$ write $X^\nu=\prod_{v\in V}v^{\nu_v}$, $|\nu|=\sum_v\nu_v$ and $\nu!=\prod_v\nu_v!$, and write
$u_\nu$ for the coefficient of $X^\nu$ in $u\in\mathcal P_j$.
\begin{enumerate}
\item \emph{Operators.} For exponent vectors $\nu$ and $\tau$ put $\binom\nu\tau=\prod_v\binom{\nu_v}{\tau_v}$, and let $D_\tau$
  be the $\mathcal K$-linear map of $\mathcal P_j$ with $D_\tau X^\nu=\binom\nu\tau X^{\nu-\tau}$ if $\tau_v\leq\nu_v$ for every
  $v$, and $D_\tau X^\nu=0$ otherwise. For a finite set $\Delta$ of exponent vectors and a family
  $\theta=(\theta_\tau)_{\tau\in\Delta}$ in $\mathcal R$ put
  \[
   G_{\Delta,\theta}(u)=\sum_{\tau\in\Delta}\theta_\tau X^\tau D_\tau u\quad(u\in\mathcal P_j),\qquad
   e_{\Delta,\theta}(\nu)=\sum_{\tau\in\Delta}\theta_\tau\binom\nu\tau\in\mathcal R .
  \]
\item \emph{Exponents.} The map $\iota_j$ sends each $v\in V$ to a power $z_{x(v),c(v)}^{P_v}$ of a variable of
  $\mathcal K[z]$, with $P_v=P_i$ for $x_{i,c}$ and $y_{i,c}$, $P_v=P_{u_l}$ for $z_l$ and $z'_l$, and $P_v=1$ for $x_{+,c}$,
  $\xi_c$, $z_+$ and $z'_+$. Let $\beta^b$ be the exponent vector of the monomial
  \[
   \prod_{i=1}^s\prod_{c=1}^kz_{(y_{j,i},1),c}^{\rho_i(c)}z_{(y_{j,i},0),c}^{\rho'_i(c)}\cdot
   \prod_{i=1}^{\eta}z_{(y_{j,i},1),1}^{P_i(\hat b_i\dot-w_i)}z_{(y_{j,i},0),1}^{P_i\varsigma_i}
  \]
  that multiplies $\iota_j(f)$ in $\iota_j(b,f)$. For an exponent vector $\nu$ on $V$ put
  $\varepsilon_b(\nu)=\beta^b+\sum_vP_v\nu_v\,e_{x(v),c(v)}$, where $e_{x,c}$ is the exponent vector of $z_{x,c}$. Put
  $\beta_v=\beta^b_{x(v),c(v)}$, and let $R_b$ be the product of the $\beta^b_{x,c}!$ over the pairs $(x,c)$ that are not of
  the form $(x(v),c(v))$.
\item \emph{Balanced weights.} Call $x_1,y_1,\dots,x_s,y_s,x_+$ and $\xi=(\xi_c)_c$ the letters of $\mathcal P_j$. A
  \emph{balanced weight} is a pair of maps $\omega\colon V\to\mathbf Z_{\geq0}$ and $\omega_L$ from the letters of
  $\mathcal P_j$ to $\mathbf Z_{\geq0}$ with $\omega(\lambda_c)=\omega_L(\lambda)$ for every letter $\lambda$ and every $c$,
  $\omega(z_l)+\omega(z'_l)=\omega_L(x_{u_l})+\omega_L(y_{u_l})$ for every $l$, and
  $\omega(z_+)+\omega(z'_+)=\omega_L(x_+)+\omega_L(\xi)$. Put
  \[
   \Omega_b(\omega,\omega_L)=\sum_{l=1}^{s-\eta}(\hat b_{u_l}\dot-o_{u_l})\,\omega(z_l)+(\alpha_+\dot-o_+)\,\omega(z_+)
   +\sum_{\lambda\neq\xi}\delta(\lambda)\,\omega_L(\lambda)+b_+\,\omega_L(\xi).
  \]
  An exponent vector $\nu$ on $V$ is \emph{homogeneous} for $b$ if $\sum_v\nu_v\,\omega(v)=\Omega_b(\omega,\omega_L)$ for every
  balanced weight $(\omega,\omega_L)$.
\item \emph{Letter degrees.} For a letter $x$ of Definition~\ref{pd-kappa-defects-windows-lift-pairs-model}(1) let
  $\omega_x(v)=1$ if $x(v)=x$ and $\omega_x(v)=0$ otherwise, and $\omega_{L,x}(\lambda)=1$ if $x(\lambda_1)=x$ and
  $\omega_{L,x}(\lambda)=0$ otherwise. Put $c_x(b)=\Omega_b(\omega_x,\omega_{L,x})$ and
  $N_b=\prod_xc_x(b)!\cdot\operatorname{sp}(b)$, the product over the letters $x$.
\item \emph{Odd variables.} The variables $x_{i,c}$, $y_{i,c}$, $z_l$ and $z'_l$ are \emph{odd}; the others are the $x_{+,c}$,
  $\xi_c$, $z_+$ and $z'_+$.
\item \emph{The leading exponent} $\nu^0_b$ is the exponent vector on $V$ with $\nu^0_b(\lambda_1)=\delta(\lambda)$ for the
  letters $\lambda\neq\xi$, $\nu^0_b(\xi_1)=b_+$, $\nu^0_b(z_l)=\hat b_{u_l}\dot-o_{u_l}$, $\nu^0_b(z_+)=\alpha_+\dot-o_+$, and
  $0$ at the other variables. Put $\Phi_b=R_b\prod_v\beta_v!\,(P_v!)^{\nu^0_b(v)}$.
\end{enumerate}
\end{definition}
\leanref{Formalism/SharpRestriction/GrowthBounds/ObjPol.lean}{Formalism.MainNotions.G.divDer}
\leanref{Formalism/SharpRestriction/GrowthBounds/CoarseMultiplicity/BFDefs.lean}{Formalism.MainNotions.ColumnSystems.Pf.gOp}
\leanref{Formalism/SharpRestriction/GrowthBounds/CoarseMultiplicity/BFDefs.lean}{Formalism.MainNotions.ColumnSystems.Pf.gEv}
\leanref{Formalism/SharpRestriction/GrowthBounds/CoarseMultiplicity/BFDefs.lean}{Formalism.MainNotions.ColumnSystems.Pf.baseExp}
\leanref{Formalism/SharpRestriction/GrowthBounds/CoarseMultiplicity/BFDefs.lean}{Formalism.MainNotions.ColumnSystems.Pf.ioExp}
\leanref{Formalism/SharpRestriction/GrowthBounds/CoarseMultiplicity/ExpFactIoExp.lean}{Formalism.MainNotions.ColumnSystems.Pf.Rblk}
\leanref{Formalism/SharpRestriction/GrowthBounds/CoarseMultiplicity/BFDefs.lean}{Formalism.MainNotions.ColumnSystems.Pf.Balanced}
\leanref{Formalism/SharpRestriction/GrowthBounds/CoarseMultiplicity/DegSplit.lean}{Formalism.MainNotions.ColumnSystems.Pf.OmB}
\leanref{Formalism/SharpRestriction/GrowthBounds/CoarseMultiplicity/BFDefs.lean}{Formalism.MainNotions.ColumnSystems.Pf.HomExp}
\leanref{Formalism/SharpRestriction/GrowthBounds/CoarseMultiplicity/BFDefs.lean}{Formalism.MainNotions.ColumnSystems.Pf.lamW}
\leanref{Formalism/SharpRestriction/GrowthBounds/CoarseMultiplicity/BFDefs.lean}{Formalism.MainNotions.ColumnSystems.Pf.lamL}
\leanref{Formalism/SharpRestriction/GrowthBounds/CoarseMultiplicity/BFDefs.lean}{Formalism.MainNotions.ColumnSystems.Pf.cdeg}
\leanref{Formalism/SharpRestriction/GrowthBounds/CoarseMultiplicity/BFDefs.lean}{Formalism.MainNotions.ColumnSystems.Pf.Nblk}
\leanref{Formalism/SharpRestriction/GrowthBounds/CoarseMultiplicity/BFDefs.lean}{Formalism.MainNotions.ColumnSystems.Pf.IsOddVar}
\leanref{Formalism/SharpRestriction/GrowthBounds/CoarseMultiplicity/LeadingExponentHomogeneous.lean}{Formalism.MainNotions.ColumnSystems.Pf.mu0}
\leanref{Formalism/SharpRestriction/GrowthBounds/CoarseMultiplicity/LeadingExponentHomogeneous.lean}{Formalism.MainNotions.ColumnSystems.Pf.PhiB}
% Lean: V is BV C0 CE1 CE2 X j = BVar (2 S_j + 1) m t: inl (a, c) is the variable of column c (from 0) of the a-th letter
%   of x_1, y_1, ..., x_s, y_s, x_+, xi (a from 0, xi is a = 2 S_j + 1), and inr (l, false), inr (l, true) are z_l, z'_l
%   (l = t for z_+, z'_+). D_tau is divDer tau, G is gOp T q, e is gEv T q. (x(v), c(v)) is symVar v, P_v is symLev v,
%   beta^b is baseExp (the exponents of digitMon and sTopMon), epsilon_b is ioExp and R_b is Rblk. Balanced, OmB and
%   HomExp are (3), with omega_L on Fin (2 S_j + 2); lamW and lamL are omega_x and omega_{L,x}, cdeg is c_x(b), Nblk is
%   N_b, IsOddVar is (5), mu0 is nu^0_b and PhiB is Phi_b.

\begin{lemma}\label{pd-kappa-defects-windows-lift-pairs-form-block-monomials}
Let the data of Definition~\ref{pd-kappa-defects-windows-lift-pairs-form-block-data} be given, with $1\leq j\leq q$ and a
block $b$ of $S_j(b_T)$.
\begin{enumerate}
\item For every finite set $\Delta$ of exponent vectors, family $\theta$ in $\mathcal R$ and $u\in\mathcal P_j$, the
  coefficient of $X^\nu$ in $G_{\Delta,\theta}(u)$ is $e_{\Delta,\theta}(\nu)u_\nu$.
\item $\iota_j(b,X^\nu)=z^{\varepsilon_b(\nu)}$ for every $\nu$, the map $\nu\mapsto\varepsilon_b(\nu)$ is injective, and
  $[\iota_j(b,u),\iota_j(b,u')]=\sum_\nu u_\nu u'_\nu\,\varepsilon_b(\nu)!$ for all $u,u'\in\mathcal P_j$.
\item $\varepsilon_b(\nu)!=R_b\prod_v(\beta_v+P_v\nu_v)!$ for every $\nu$.
\item If $b$ is kept, then for every odd variable $v$ there is an integer $t_v\geq0$ with $P_v=p^{t_v+1}$, and
  $\beta_v<P_v$; for every other variable $v$, $P_v=1$ and $\beta_v=0$.
\end{enumerate}
\end{lemma}
\leanref{Formalism/SharpRestriction/GrowthBounds/CoarseMultiplicity/TopOddNeEnd.lean}{Formalism.MainNotions.ColumnSystems.Pf.eigAt}
\leanref{Formalism/SharpRestriction/GrowthBounds/CoarseMultiplicity/TopOddNeEnd.lean}{Formalism.MainNotions.ColumnSystems.Pf.eig_one}
\leanref{Formalism/SharpRestriction/GrowthBounds/CoarseMultiplicity/IoBlockMonomial.lean}{Formalism.MainNotions.ColumnSystems.Pf.ioBlock_monomial}
\leanref{Formalism/SharpRestriction/GrowthBounds/CoarseMultiplicity/IoBlockMonomial.lean}{Formalism.MainNotions.ColumnSystems.Pf.ioExp_inj}
\leanref{Formalism/SharpRestriction/GrowthBounds/CoarseMultiplicity/IoBlockMonomial.lean}{Formalism.MainNotions.ColumnSystems.Pf.ioFockAt}
\leanref{Formalism/SharpRestriction/GrowthBounds/CoarseMultiplicity/ExpFactIoExp.lean}{Formalism.MainNotions.ColumnSystems.Pf.expFact_ioExp}
\leanref{Formalism/SharpRestriction/GrowthBounds/CoarseMultiplicity/CharZeroOfW.lean}{Formalism.MainNotions.ColumnSystems.Pf.oddVar_data}
\leanref{Formalism/SharpRestriction/GrowthBounds/CoarseMultiplicity/CharZeroOfW.lean}{Formalism.MainNotions.ColumnSystems.Pf.endVar_data}

\begin{proof}
(1) Let $\tau\in\Delta$. If $\tau_v\leq\nu_v$ for every $v$, then $X^\tau D_\tau X^\nu=\binom\nu\tau X^\nu$. Otherwise
$X^\tau D_\tau X^\nu=0=\binom\nu\tau X^\nu$, as some factor $\binom{\nu_v}{\tau_v}$ vanishes. So
$G_{\Delta,\theta}(X^\nu)=e_{\Delta,\theta}(\nu)X^\nu$, and (1) follows by linearity.

(2) By Definition~\ref{pd-kappa-defects-windows-lift-pairs-model}(5), $\iota_j(b,\cdot)$ is the $\mathcal K$-algebra map
$\iota_j$ followed by multiplication by $z^{\beta^b}$, and $\iota_j(X^\nu)=\prod_vz_{x(v),c(v)}^{P_v\nu_v}$; so
$\iota_j(b,X^\nu)=z^{\varepsilon_b(\nu)}$. By the list in Definition~\ref{pd-kappa-defects-windows-lift-pairs-model}(5),
distinct variables $v$ give distinct pairs $(x(v),c(v))$, as the $y_{j,i}$ are distinct and differ from $s_{j+1}$. So the
entry of $\varepsilon_b(\nu)$ at $(x(v),c(v))$ is $\beta_v+P_v\nu_v$, and as $P_v\geq1$, $\varepsilon_b(\nu)$ determines $\nu$.
Hence $\iota_j(b,u)=\sum_\nu u_\nu z^{\varepsilon_b(\nu)}$, with distinct monomials for distinct $\nu$, and the Fock pairing
$[F,G]=\sum_\kappa\kappa!\,F_\kappa G_\kappa$ gives the formula.

(3) By (2), the entries of $\varepsilon_b(\nu)$ are $\beta_v+P_v\nu_v$ at the pairs $(x(v),c(v))$ and $\beta^b_{x,c}$ at the
other pairs.

(4) Let $v$ be odd. Then $x(v)$ is a letter at some $y_{j,i}$, and $P_v=p^{g_j-d_{j,i}}$ with $0<d_{j,i}<g_j$. If $v$ is
$x_{i,c}$ or $y_{i,c}$, then $c(v)=c+1\geq2$, and $\beta_v$ is $\rho_i(c+1)$ or $\rho'_i(c+1)$. If $v$ is $z_l$ or $z'_l$, then
$c(v)=1$ at $y_{j,u_l}$ with $u_l>\eta$, and $\beta_v$ is $\rho_{u_l}(1)$ or $\rho'_{u_l}(1)$. The entries of $\rho_i$ and
$\rho'_i$ of a kept block are less than $P_i$. If $v$ is not odd, then $x(v)$ is a letter at $s_{j+1}$, where $\beta^b$
vanishes, and $P_v=1$.
\end{proof}
% Lean: (1) is eigAt (eig_one); (2) is ioBlock_monomial, ioExp_inj (symVar_inj, coarse_eval, symLev_pos) and ioFockAt
%   (fock_monoMap); (3) is expFact_ioExp; (4) is oddVar_data (lev_pow, baseExp_odd and the digits of Kept) and
%   endVar_data (baseExp_end).

\begin{lemma}\label{pd-kappa-defects-windows-lift-pairs-form-block-homogeneous}
Let the data of Definition~\ref{pd-kappa-defects-windows-lift-pairs-form-block-data} be given, with $1\leq j\leq q$ and a
block $b$ of $S_j(b_T)$.
\begin{enumerate}
\item Let $b$ be kept, and let $\mu$ be a real number. If $X^\nu$ occurs in an element of the $\mathcal K$-span of the
  generators of all levels of $b$ formed with $\mu$ (Definition~\ref{pd-kappa-defects-segment-lattice}(5)), then $\nu$ is
  homogeneous for $b$.
\item For every letter $x$ of Definition~\ref{pd-kappa-defects-windows-lift-pairs-model}(1), $(\omega_x,\omega_{L,x})$ is a
  balanced weight. If $\nu$ is homogeneous for $b$, then $\sum_{x(v)=x}\nu_v=c_x(b)$ for every letter $x$. If $b$ is kept
  and $\nu$ is homogeneous for $b$, then $\nu!\operatorname{sp}(b)$ divides $N_b$.
\item $\nu^0_b$ is homogeneous for $b$. If $\nu$ and $\nu'$ are homogeneous for $b$, then
  $\prod_vf(P_v)^{\nu_v}=\prod_vf(P_v)^{\nu'_v}$ for every map $f$ from the positive integers to $\mathbf Q$.
\end{enumerate}
\end{lemma}
\leanref{Formalism/SharpRestriction/GrowthBounds/CoarseMultiplicity/SymVarInlFst.lean}{Formalism.MainNotions.ColumnSystems.Pf.homAt}
\leanref{Formalism/SharpRestriction/GrowthBounds/CoarseMultiplicity/SymVarInlFst.lean}{Formalism.MainNotions.ColumnSystems.Pf.lam_balanced}
\leanref{Formalism/SharpRestriction/GrowthBounds/CoarseMultiplicity/SymVarInlFst.lean}{Formalism.MainNotions.ColumnSystems.Pf.sum_fiber}
\leanref{Formalism/SharpRestriction/GrowthBounds/CoarseMultiplicity/SymVarInlFst.lean}{Formalism.MainNotions.ColumnSystems.Pf.dvdAt}
\leanref{Formalism/SharpRestriction/GrowthBounds/CoarseMultiplicity/LeadingExponentHomogeneous.lean}{Formalism.MainNotions.ColumnSystems.Pf.leading_exponent_homogeneous}
\leanref{Formalism/SharpRestriction/GrowthBounds/CoarseMultiplicity/LeadingExponentHomogeneous.lean}{Formalism.MainNotions.ColumnSystems.Pf.lev_balanced}
\leanref{Formalism/SharpRestriction/GrowthBounds/CoarseMultiplicity/LeadingExponentHomogeneous.lean}{Formalism.MainNotions.ColumnSystems.Pf.lev_const}

\begin{proof}
(1) The argument does not use that $b$ is kept. Let $(\omega,\omega_L)$ be a balanced weight. Say that a polynomial has
$\omega$-degree $d$ if every monomial $X^\nu$ in it has $\sum_v\nu_v\omega(v)=d$; the $\omega$-degree of a product is the sum of those of the factors. In a generator of level
$(j_1,\dots,j_{s-\eta},j_+)$:
\begin{itemize}
\item $\hat E_l$ has $\omega$-degree $\omega(z_l)+\omega(z'_l)=\omega_L(x_{u_l})+\omega_L(y_{u_l})$, as each
  $x_{u_l,c}y_{u_l,c}$ has the latter, and likewise $\hat E_+$ has $\omega$-degree $\omega_L(x_+)+\omega_L(\xi)$;
\item $z_l^{\hat b_{u_l}\dot-o_{u_l}}$ and $z_+^{\alpha_+\dot-o_+}$ have $\omega$-degrees $(\hat b_{u_l}\dot-o_{u_l})\omega(z_l)$
  and $(\alpha_+\dot-o_+)\omega(z_+)$;
\item $f$ has degree $\delta'(\lambda)$ in the variables of each letter $\lambda\neq\xi$, so $\omega$-degree
  $\sum_\lambda\delta'(\lambda)\omega_L(\lambda)$, and $g$ has $\omega$-degree $(b_+-j_+)\omega_L(\xi)$.
\end{itemize}
The level bounds give $j_l\leq\delta(x_{u_l})$, $j_l\leq\delta(y_{u_l})$, $j_+\leq\delta(x_+)$ and $j_+\leq b_+$, and $\delta'$ is
$\delta$ lowered by $j_l$ at $x_{u_l}$ and $y_{u_l}$ and by $j_+$ at $x_+$. So the sum of these $\omega$-degrees is
$\Omega_b(\omega,\omega_L)$. Every generator, and so every element of their $\mathcal K$-span, has $\omega$-degree
$\Omega_b(\omega,\omega_L)$.

(2) The map $\iota_j$ sends the variables $\lambda_c$ of a letter $\lambda$ of $\mathcal P_j$ to variables of the one letter
$x(\lambda_1)$, so $\omega_x(\lambda_c)=\omega_{L,x}(\lambda)$. It sends $z_l$ and $z'_l$ to variables of the letters
$x(x_{u_l,1})$ and $x(y_{u_l,1})$, and $z_+$ and $z'_+$ to variables of $x(x_{+,1})$ and $x(\xi_1)$; this gives the other two
conditions. For homogeneous $\nu$, $\sum_{x(v)=x}\nu_v=\sum_v\nu_v\omega_x(v)=c_x(b)$. A product of factorials divides the
factorial of the sum of the arguments, as multinomial coefficients are integers; so $\prod_{x(v)=x}\nu_v!$ divides
$c_x(b)!$. As every $v$ has one letter $x(v)$, multiplying over $x$ shows that $\nu!$ divides $\prod_xc_x(b)!$, and so
$\nu!\operatorname{sp}(b)$ divides $N_b$. This argument does not use that $b$ is kept.

(3) Let $(\omega,\omega_L)$ be a balanced weight. As $\omega(\lambda_1)=\omega_L(\lambda)$ for every letter $\lambda$,
\[
 \sum_v\nu^0_b(v)\omega(v)=\sum_{\lambda\neq\xi}\delta(\lambda)\omega_L(\lambda)+b_+\omega_L(\xi)
 +\sum_l(\hat b_{u_l}\dot-o_{u_l})\omega(z_l)+(\alpha_+\dot-o_+)\omega(z_+)=\Omega_b(\omega,\omega_L).
\]
For an integer $P\geq1$ let $\omega^P(v)=1$ if $P_v=P$ and $\omega^P(v)=0$ otherwise, and $\omega^P_L(\lambda)=1$ if
$P_{\lambda_1}=P$ and $\omega^P_L(\lambda)=0$ otherwise. As $P_{\lambda_c}=P_{\lambda_1}$,
$P_{z_l}=P_{z'_l}=P_{x_{u_l,1}}=P_{y_{u_l,1}}$ and $P_{z_+}=P_{z'_+}=P_{x_{+,1}}=P_{\xi_1}=1$, the pair $(\omega^P,\omega^P_L)$ is a
balanced weight. So for homogeneous $\nu$ the number $\sum_{P_v=P}\nu_v=\Omega_b(\omega^P,\omega^P_L)$ does not depend on $\nu$,
and neither does $\prod_vf(P_v)^{\nu_v}=\prod_Pf(P)^{\sum_{P_v=P}\nu_v}$.
\end{proof}
% Lean: (1) is homAt (hgen: IsWeightedHomogeneous of each factor, with
%   sum_deg_sub_mul_add_eq_sum_deg_mul for the level bounds hok, then the weighted homogeneous submodule); the Lean
%   states it for kept blocks and does not use keptness. (2) is lam_balanced, sum_fiber and dvdAt
%   (Nat.prod_factorial_dvd_factorial_sum, prod_fiberwise); DvdAt assumes Kept, which dvdAt does not use (it is the
%   discarded binder in its intro). (3) is leading_exponent_homogeneous, and lev_const with
%   lev_balanced (the weights lvW and lvL, through sl_Z, symbol_level_lower_letter, sl_E, symbol_level_last, sl_U and
%   sl_xi); the Lean states lev_const for maps f into the natural numbers.

\begin{lemma}\label{pd-kappa-defects-windows-lift-pairs-form-block-saturated}
Let the data of Definition~\ref{pd-kappa-defects-windows-lift-pairs-form-block-data} be given, with $1\leq j\leq q$ and a
block $b$ of $S_j(b_T)$, and let $\mu$ be a real number. If $u$ lies in the $\mathcal K$-span of $Z_j(b;\mu)$ and has
coefficients in $\mathcal R$, then $u\in Z_j(b;\mu)$.
\end{lemma}
\leanref{Formalism/SharpRestriction/GrowthBounds/CoarseMultiplicity/SatBlock.lean}{Formalism.MainNotions.ColumnSystems.Pf.satAt}
\leanref{Formalism/SharpRestriction/GrowthBounds/CoarseMultiplicity/SatBlock.lean}{Formalism.MainNotions.ColumnSystems.Pf.sat_block}

\begin{proof}
Use the notation of the proof of Lemma~\ref{pd-kappa-defects-digits-block-rank} for the segment datum $S_j(b_T)$, with
$\mathcal O=\mathcal R$, $r=m$, $\alpha_i=a_T(y_{j,i})$, $\alpha_+=a_T(s_{j+1})$ and $\mu$; the $\mathcal K$-span of
$Z_j(b;\mu)$ is $V_b$. That proof shows $u=\sum_{\ell,\beta}\pi_\ell\,y_{\ell,\beta}\,\xi^\beta$ with
$y_{\ell,\beta}\in Y_b(\ell)$, and $u=\tau(B)$ with $B=\sum_{\ell,\beta}\sigma_\ell\,y_{\ell,\beta}\,\xi^\beta$. Let $\zeta$ be the
endomorphism of the $\mathcal K$-algebra $\mathcal P_j$ that sends the $z_l$ and $z_+$ to $1$ and fixes the other variables,
and $\gamma$ the one that sends $z'_l$ to $z'_l-q_l$ and $z'_+$ to $z'_+-q_+$ and fixes the other variables. Then
$\gamma(\zeta(\tau(w)))=\zeta(w)$ for every variable $w$: $\tau(z'_l)=z_lz'_l+q_l$, $\zeta$ sends it to $z'_l+q_l$ and
$\gamma$ to $z'_l$, as $\zeta$ and $\gamma$ fix $q_l$; likewise for $z'_+$; and the three maps fix the other variables,
except that $\zeta$ sends $z_l$ and $z_+$ to $1$, which $\gamma$ fixes. So $\zeta(B)=\gamma(\zeta(u))$, which has
coefficients in $\mathcal R$, as $u$ has and $\zeta$ and $\gamma$ send variables to polynomials with coefficients in
$\mathcal R$. The homomorphism $\psi$ of that proof sends $z_l$ and $z_+$ to $1$, so $\psi(\zeta(B))=\psi(B)$, which is the
polynomial with coefficient $y_{\ell,\beta}$ at $\prod_lS_l^{j_l}S_+^{j_+}\prod_c\Xi_c^{\beta_c}$. As $\psi$ sends each variable to
a variable or to $1$, and $\zeta(B)$ has coefficients in $\mathcal R$, every $y_{\ell,\beta}$ has coefficients in
$\mathcal R$. Then $\pi_\ell\,y_{\ell,\beta}\,\xi^\beta$ is a generator of level $\ell$, with $f=y_{\ell,\beta}$ and $g=\xi^\beta$,
and $u\in Z_j(b;\mu)$.
\end{proof}
% Lean: satAt = sat_block for every BlockData D: span_eq_range and Ltot_apply give v with u = Ltot v = tau B (ell_eq_tau);
%   beF and giF are epsilon and gamma, gi_be_tau (with qq_fix) gives gamma(epsilon(tau P)) = epsilon(P), and h1 the
%   integrality; ppi_be, ppi_int and ppi_base with eExp_inj extract the coefficients y; the generators are those of
%   (v i, xiMon) with xiMon_mem.

\begin{lemma}\label{pd-kappa-defects-windows-lift-pairs-form-block-shift-operators}
Let $K$ be a field and $V$ a finite set, let $K[V]$ be the polynomial ring in the variables $v\in V$, and write $X^\nu$ for the
monomial with exponent vector $\nu\colon V\to\mathbf Z_{\geq0}$. For exponent vectors $\nu$ and $\tau$ let $D_\tau$ be the
$K$-linear map of $K[V]$ with $D_\tau X^\nu=\binom\nu\tau X^{\nu-\tau}$ if $\tau_w\leq\nu_w$ for every $w$ and $D_\tau X^\nu=0$
otherwise, as in Definition~\ref{pd-kappa-defects-windows-lift-pairs-form-block-data}(1). For $v\in V$ let $e_v$ be its exponent
vector, and for an integer $n\geq0$ put $C_{v,n}=X^{ne_v}D_{ne_v}$. Say that $v$ \emph{does not occur} in $F\in K[V]$ if
$\nu_v=0$ for every monomial $X^\nu$ of $F$. Let $v\in V$, $n\geq0$ and $F,G\in K[V]$.
\begin{enumerate}
\item $C_{v,n}X^\nu=\binom{\nu_v}nX^\nu$ for every $\nu$. So $C_{v,0}F=F$, every monomial of $C_{v,n}F$ is a monomial of $F$,
  and if the coefficients of $F$ lie in a subring of $K$, so do those of $C_{v,n}F$.
\item $C_{v,n}(FG)=\sum_{a=0}^nC_{v,a}(F)\,C_{v,n-a}(G)$.
\item If $v$ does not occur in $F$, then $C_{v,n}F=0$ for $n\geq1$, and $C_{v,n}(FG)=F\,C_{v,n}(G)$.
\item $C_{v,n}(v^z)=\binom znv^z$ for every integer $z\geq0$.
\item Let $E=vm+E'$, where $v$ does not occur in $m,E'\in K[V]$. Then for all integers $J,l\geq0$,
  $C_{v,l}(E^J)=\binom JlE^{J-l}(vm)^l$ if $l\leq J$, and $C_{v,l}(E^J)=0$ if $l>J$.
\end{enumerate}
\end{lemma}
\leanref{Formalism/SharpRestriction/GrowthBounds/CoarseMultiplicity/CthMonomial.lean}{Formalism.MainNotions.ColumnSystems.Pf.cth}
\leanref{Formalism/SharpRestriction/GrowthBounds/CoarseMultiplicity/CthMonomial.lean}{Formalism.MainNotions.ColumnSystems.Pf.coeff_cth}
\leanref{Formalism/SharpRestriction/GrowthBounds/CoarseMultiplicity/CthMonomial.lean}{Formalism.MainNotions.ColumnSystems.Pf.cth_mul}
\leanref{Formalism/SharpRestriction/GrowthBounds/CoarseMultiplicity/CthMonomial.lean}{Formalism.MainNotions.ColumnSystems.Pf.cth_free}
\leanref{Formalism/SharpRestriction/GrowthBounds/CoarseMultiplicity/CthMonomial.lean}{Formalism.MainNotions.ColumnSystems.Pf.cth_mul_free}
\leanref{Formalism/SharpRestriction/GrowthBounds/CoarseMultiplicity/ChainEmbInj.lean}{Formalism.MainNotions.ColumnSystems.Pf.cth_X_pow}
\leanref{Formalism/SharpRestriction/GrowthBounds/CoarseMultiplicity/LettersAtUpdate.lean}{Formalism.MainNotions.ColumnSystems.Pf.cth_lin}
\leanref{Formalism/SharpRestriction/GrowthBounds/CoarseMultiplicity/CthMonomial.lean}{Formalism.MainNotions.ColumnSystems.Pf.cth_pow}

\begin{proof}
(1) If $n\leq\nu_v$, then $D_{ne_v}X^\nu=\binom{\nu_v}nX^{\nu-ne_v}$, and multiplying by $X^{ne_v}$ gives $\binom{\nu_v}nX^\nu$.
Otherwise $C_{v,n}X^\nu=0=\binom{\nu_v}nX^\nu$. The rest follows by linearity, as $\binom{\nu_v}0=1$ and binomial coefficients
are integers.

(2) Both sides are bilinear in $F$ and $G$, so it suffices to take $F=X^\nu$ and $G=X^{\nu'}$. By (1) the two sides are then
$\binom{\nu_v+\nu'_v}nX^{\nu+\nu'}$ and $\sum_a\binom{\nu_v}a\binom{\nu'_v}{n-a}X^{\nu+\nu'}$, which agree by Vandermonde's
identity.

(3) By (1), $C_{v,n}$ multiplies each monomial of $F$ by $\binom0n=0$ if $n\geq1$. So in (2) only the term with $a=0$ remains,
and $C_{v,0}F=F$.

(4) This is (1) for $\nu=ze_v$.

(5) By (3), $C_{v,a}E'=0$ for $a\geq1$, and $C_{v,a}(vm)=m\,C_{v,a}(v)$, which is $\binom1a\,vm$ by (4). So $C_{v,1}E=vm$ and
$C_{v,a}E=0$ for $a\geq2$. We induct on $J$. As $v$ does not occur in $E^0=1$, $C_{v,l}(1)$ is $1$ for $l=0$ and $0$ for
$l\geq1$, by (1) and (3). For the step, (2) and the values of $C_{v,a}E$ give
$C_{v,l}(E^{J+1})=C_{v,l}(E^J)\,E+C_{v,l-1}(E^J)\,vm$, where the second term is absent for $l=0$. If $l\leq J$, induction and
Pascal's rule $\binom Jl+\binom J{l-1}=\binom{J+1}l$ give $\binom{J+1}lE^{J+1-l}(vm)^l$. If $l=J+1$, only the second term
remains, and it is $(vm)^{J+1}$. If $l>J+1$, both terms are $0$.
\end{proof}
% Lean: C_{v,n} is cth v n = monomial (single v n) 1 * divDer (single v n), over any field K and any type of variables.
%   (1) is coeff_cth (eig_one), with intPol_cth for the coefficients; (2) is cth_mul (Nat.add_choose_eq, by induction on
%   the monomials of F and G); (3) is cth_free and cth_mul_free; (4) is cth_X_pow; (5) is cth_lin (the values of C_{v,a}E)
%   and cth_pow (induction on J through cth_mul and Nat.choose_succ_succ').

\begin{lemma}\label{pd-kappa-defects-windows-lift-pairs-form-block-shift-spaces}
Let $p$ and $\mathcal O\subset\mathcal K$ be as in Definition~\ref{pd-kappa-defects-shares}, and $r\geq1$ and $s\geq0$ integers;
form the letters $x_1,y_1,\dots,x_s,y_s,x_+$, their variables and the spaces $Y(\mathcal L,M)$ as in
Definition~\ref{pd-kappa-defects-segment-lattice}(3) and~(4), and let $C_{w,n}$ be the operators of
Lemma~\ref{pd-kappa-defects-windows-lift-pairs-form-block-shift-operators} on the polynomial ring over $\mathcal K$ in the
variables of the letters. Let $\mathcal L$ be an initial segment of the letters, $M$ an integer and $f\in Y(\mathcal L,M)$.
\begin{enumerate}
\item $f\in Y(\mathcal L,M')$ for every integer $M'\geq M$.
\item $C_{w,n}f\in Y(\mathcal L,M+n)$ for every variable $w$ of a letter and every integer $n\geq0$.
\item $\lambda_c\lambda'_{c'}f\in Y(\mathcal L,M+1)$ for every letter $\lambda$, every covector letter $\lambda'$ and all
  $1\leq c,c'\leq r$.
\end{enumerate}
\end{lemma}
\leanref{Formalism/SharpRestriction/GrowthBounds/CoarseMultiplicity/OtherEFree.lean}{Formalism.MainNotions.ColumnSystems.Pf.yjtZ_mono}
\leanref{Formalism/SharpRestriction/GrowthBounds/CoarseMultiplicity/DivDerZero.lean}{Formalism.MainNotions.ColumnSystems.Pf.yjtZ_cth}
\leanref{Formalism/SharpRestriction/GrowthBounds/CoarseMultiplicity/DivDerZero.lean}{Formalism.MainNotions.ColumnSystems.Pf.yjtZ_X}
\leanref{Formalism/SharpRestriction/GrowthBounds/CoarseMultiplicity/BlockWithThresholds.lean}{Formalism.MainNotions.ColumnSystems.Pf.isMult_XX}
\leanref{Formalism/SharpRestriction/GrowthBounds/CoarseMultiplicity/DivDerZero.lean}{Formalism.MainNotions.ColumnSystems.Pf.ycond_op}

\begin{proof}
(1) The space $Y(\mathcal L,M)$ is stable under the $E^{\mathcal L}_{cd}$ and lies in the span of the monomials of weight at
most $M$ for $\mathcal L$, hence in the span of those of weight at most $M'$. So it lies in the largest such stable subspace,
$Y(\mathcal L,M')$.

For (2) and (3), number the letters $x_1,y_1,\dots,x_s,y_s,x_+$ as $0,\dots,2s$. With the $x_i$ and $x_+$ as the vector
letters this is a lettering of $\{0,\dots,2s\}$ (Definition~\ref{pd-symbols-path-generators}), and $\mathcal L=\{0,\dots,i-1\}$
for some $i$. The field $\mathcal K$ has characteristic $0$: if its characteristic were a prime $\ell$, the integer $p$ would be
$0$ or a unit of $\mathcal O$. So by Lemma~\ref{pd-kappa-defects-paths-isotypic}(3), a polynomial $G$ in the variables of the
letters lies in $Y(\mathcal L,M'')$ if and only if $\lambda_{\mathcal L}(G)\leq M''$, for every integer $M''$. In particular
$\lambda_{\mathcal L}(f)\leq M$. Lemma~\ref{pd-symbols-test-chains-projection-clauses}(5)(a) applies over $\mathcal K$, with the
set of the letters as $\Lambda$, $n=r$ and $S=\mathcal L$.

(2) Take $F=f$ and $z^\alpha=z^\beta=w^n$, so that $G=C_{w,n}f$. If $w$ is a variable of a vector letter in $\mathcal L$, then
$\alpha_{\mathrm v}=n$ and $\beta_{\mathrm c}=0$; if it is a variable of a covector letter in $\mathcal L$, then
$\alpha_{\mathrm v}=0$ and $\beta_{\mathrm c}=n$; otherwise both are $0$. So $\lambda_{\mathcal L}(C_{w,n}f)\leq M+n$.

(3) Take $z^\beta=1$. With $F=f$ and $z^\alpha=\lambda'_{c'}$, $\alpha_{\mathrm v}=0$, as $\lambda'$ is a covector letter; so
$\lambda_{\mathcal L}(\lambda'_{c'}f)\leq M$. With $F=\lambda'_{c'}f$ and $z^\alpha=\lambda_c$, $\alpha_{\mathrm v}\leq1$; so
$\lambda_{\mathcal L}(\lambda_c\lambda'_{c'}f)\leq M+1$.
\end{proof}
% Lean: the Lean's YjtZ O LL J thr is the intersection of the spaces Y(L_k, thr k) with the polynomials of the row degrees
%   LL.deg (chPol) and with integral coefficients (IntPol), for thresholds thr k in WithTop Z. (1) is yjt_mono' (through
%   ycond_mono'), as yjtZ_mono; (2) is yjtZ_cth, which also keeps the row degrees and the integrality (coeff_cth,
%   intPol_cth); (3) is isMult_XX, two applications of yjtZ_X, first at the covector letter with n = 0 and then at the
%   other letter with n = 1. All three go through ycond_op: PolBridge.mem_ycond_iff (Lemma
%   pd-kappa-defects-paths-isotypic(3)) and flt_op_lam1S (RootWeights/DerivativeValuationBounds.lean, the second bound
%   of DerivativeValuationBound, Lemma pd-symbols-test-chains-projection-clauses(5)(a)), with vecDeg al + covDeg be <= n.

\begin{lemma}\label{pd-kappa-defects-windows-lift-pairs-form-block-shift-generators}
Let the data of Definition~\ref{pd-kappa-defects-windows-lift-pairs-form-block-data} be given, with $1\leq j\leq q$ and a block
$b$ of $S_j(b_T)$, and let $C_{v,n}$ be the operators of
Lemma~\ref{pd-kappa-defects-windows-lift-pairs-form-block-shift-operators} on $\mathcal P_j$. Let $v$ be an odd variable,
$n\geq0$ an integer and $\mu'$ a real number. Then $C_{v,n}u\in Z_j(b;\mu'-n)$ for every generator $u$ of a level of $b$ formed
with $\mu'$, as in the definition of $Z_j(b;\mu')$ (Definition~\ref{pd-kappa-defects-windows-lift-pairs-model}(6)).
\end{lemma}
\leanref{Formalism/SharpRestriction/GrowthBounds/CoarseMultiplicity/OtherEFree.lean}{Formalism.MainNotions.ColumnSystems.Pf.cth_gen}
\leanref{Formalism/SharpRestriction/GrowthBounds/CoarseMultiplicity/LettersAtUpdate.lean}{Formalism.MainNotions.ColumnSystems.Pf.shift_mem}
\leanref{Formalism/SharpRestriction/GrowthBounds/CoarseMultiplicity/LettersAtUpdate.lean}{Formalism.MainNotions.ColumnSystems.Pf.tie_mem}
\leanref{Formalism/SharpRestriction/GrowthBounds/CoarseMultiplicity/LettersAtUpdate.lean}{Formalism.MainNotions.ColumnSystems.Pf.lettersAt_update}
\leanref{Formalism/SharpRestriction/GrowthBounds/CoarseMultiplicity/LettersAtUpdate.lean}{Formalism.MainNotions.ColumnSystems.Pf.levelOK_update}
\leanref{Formalism/SharpRestriction/GrowthBounds/CoarseMultiplicity/ChainEmbInj.lean}{Formalism.MainNotions.ColumnSystems.Pf.tail_chain}
\leanref{Formalism/SharpRestriction/GrowthBounds/CoarseMultiplicity/BlockWithThresholds.lean}{Formalism.MainNotions.ColumnSystems.Pf.gen_mem}
\leanref{Formalism/SharpRestriction/GrowthBounds/CoarseMultiplicity/GTheta.lean}{Formalism.MainNotions.ColumnSystems.Pf.tieOK_blockBase}

\begin{proof}
Write $s=S_j$ and $\eta=\eta_j$, and use the notation of Definition~\ref{pd-kappa-defects-segment-lattice}(5), formed over
$\mathcal R$ with $\alpha_i=a_T(y_{j,i})$, $\alpha_+=a_T(s_{j+1})$ and $r=m$; the bounds $M_l$ and $M_+$ are those formed with
$\mu'$. For an integer $a$, the bounds formed with $\mu'-a$ are $M_l+a$ and $M_++a$. By
Lemma~\ref{pd-kappa-defects-windows-lift-pairs-form-block-shift-spaces}(1), $Z_j(b;\mu'')\subset Z_j(b;\mu''-a)$ for every real
$\mu''$ and every integer $a\geq0$. Put $\pi_l=\sum_cx_{u_l,c}y_{u_l,c}$, so that $\hat E_l=z_lz'_l+\pi_l$. Let
\[
 u=\prod_l\hat E_l^{j_l}\cdot T(f),\qquad
 T(f)=\hat E_+^{j_+}\prod_lz_l^{\hat b_{u_l}\dot-o_{u_l}}\cdot z_+^{\alpha_+\dot-o_+}\cdot f\cdot g,
\]
be a generator of level $(j_1,\dots,j_{s-\eta},j_+)$ formed with $\mu'$, and write $T(f')$ for the same product with a
polynomial $f'$ in the letters in place of $f$. For an integer $a\geq0$, say that such an $f'$ is \emph{admissible with} $a$
if it has coefficients in $\mathcal R$, has the degree $\delta'(\lambda)$ of this level in the variables of each letter
$\lambda$, lies in $Y(\mathcal L_l,M_l+a)$ for every $l$, and lies in $Y(\mathcal L_+,M_++a)$. Then $\prod_l\hat E_l^{j_l}T(f')$
is a generator of the same level formed with $\mu'-a$. The polynomial $f$ is admissible with $0$, and so with every
$a\geq0$, by Lemma~\ref{pd-kappa-defects-windows-lift-pairs-form-block-shift-spaces}(1).

\emph{Step 1: lowering the level.} Let $1\leq i\leq s-\eta$, $0\leq l\leq j_i$ and $a\geq0$ be integers, let $f'$ be
admissible with $a$, and let $q$ be $x_{u_i,c}y_{u_i,c}$ for some $c$, or $-\pi_i$. Then
\[
 \hat E_i^{j_i-l}q^l\prod_{k\neq i}\hat E_k^{j_k}\cdot T(f')\in Z_j(b;\mu'-a-l).
\]
Indeed, $y_{u_i}$ is a covector letter, so by Lemma~\ref{pd-kappa-defects-windows-lift-pairs-form-block-shift-spaces}(3) every
product $x_{u_i,c}y_{u_i,c'}$ maps $Y(\mathcal L,M'')$ into $Y(\mathcal L,M''+1)$, for every initial segment $\mathcal L$ and
every integer $M''$. As these spaces are subspaces, $q$ does too, and $q^l$ maps $Y(\mathcal L,M'')$ into
$Y(\mathcal L,M''+l)$. Also $q$ has integer coefficients, and multiplying by $q$ raises by $1$ the degrees in the variables of
$x_{u_i}$ and of $y_{u_i}$ and keeps the others. So $q^lf'$ has coefficients in $\mathcal R$, lies in $Y(\mathcal L_k,M_k+a+l)$
for every $k$ and in $Y(\mathcal L_+,M_++a+l)$, and has the degrees $\delta'$ raised by $l$ at $x_{u_i}$ and $y_{u_i}$. As
$j_i\leq\min(\delta(x_{u_i}),\delta(y_{u_i}))$, these are the degrees of the level with $j_i$ lowered to $j_i-l$, which is again
a level. So the product is a generator of that level formed with $\mu'-a-l$, with $q^lf'$ in place of $f$ and the same $g$.

\emph{Step 2: a variable of one $\hat E_i$.} Let $1\leq i\leq s-\eta$, and suppose that $\hat E_i=vm+E'$, where $v$ does not
occur in $m$, in $E'$ or in the $\hat E_k$ with $k\neq i$, and that $vm=q$ or $vm=\hat E_i+q$ with $q$ as in Step~1. Suppose
also that for every $0\leq b'\leq n$ there are an integer $c_{b'}\geq0$ and an $f_{b'}$ admissible with $b'$ such that
$C_{v,b'}T(f)=c_{b'}T(f_{b'})$. Then $C_{v,n}u\in Z_j(b;\mu'-n)$. Indeed, $v$ does not occur in $A=\prod_{k\neq i}\hat E_k^{j_k}$,
so by Lemma~\ref{pd-kappa-defects-windows-lift-pairs-form-block-shift-operators}(3), (2) and~(5),
\[
 C_{v,n}u=A\,C_{v,n}\big(\hat E_i^{j_i}T(f)\big)=\sum_l\binom{j_i}l\,c_{n-l}\,\hat E_i^{j_i-l}(vm)^l\,A\,T(f_{n-l}),
\]
the sum over the integers $0\leq l\leq\min(j_i,n)$. If $vm=q$, the term at $l$ lies in $Z_j(b;\mu'-(n-l)-l)=Z_j(b;\mu'-n)$ by
Step~1. If $vm=\hat E_i+q$, then $\hat E_i^{j_i-l}(vm)^l=\sum_{a'=0}^l\binom l{a'}\hat E_i^{j_i-(l-a')}q^{l-a'}$, and by Step~1,
with $l-a'\leq j_i$, the summand at $a'$ times $A\,T(f_{n-l})$ lies in $Z_j(b;\mu'-n+a')$, which lies in $Z_j(b;\mu'-n)$ as
$a'\geq0$. As $Z_j(b;\mu'-n)$ is an $\mathcal R$-module, $C_{v,n}u$ lies in it.

\emph{Step 3: the odd variables.} These are the $x_{i,c}$, $y_{i,c}$, $z_l$ and $z'_l$
(Definition~\ref{pd-kappa-defects-windows-lift-pairs-form-block-data}(5)). The polynomial $\hat E_+$ involves only the
variables of $x_+$, the $\xi_c$, $z_+$ and $z'_+$, and $g$ only the $\xi_c$.
\begin{itemize}
\item $v=x_{u_i,c}$ or $v=y_{u_i,c}$, for some $1\leq i\leq s-\eta$ and some $c$. Then $\hat E_i=x_{u_i,c}y_{u_i,c}+E'$ with
  $E'=z_iz'_i+\sum_{c'\neq c}x_{u_i,c'}y_{u_i,c'}$; take $m=y_{u_i,c}$, respectively $m=x_{u_i,c}$, and $q=vm=x_{u_i,c}y_{u_i,c}$.
  The variable $v$ does not occur in $m$, in $E'$, or in the $\hat E_k$ with $k\neq i$, whose letters $x_{u_k}$ and $y_{u_k}$
  differ from $x_{u_i}$ and $y_{u_i}$. In $T(f)$ it occurs only in $f$, so $C_{v,b'}T(f)=T(C_{v,b'}f)$ by
  Lemma~\ref{pd-kappa-defects-windows-lift-pairs-form-block-shift-operators}(3). By part (1) of that lemma and
  Lemma~\ref{pd-kappa-defects-windows-lift-pairs-form-block-shift-spaces}(2), $C_{v,b'}f$ is admissible with $b'$. Step~2
  applies with $c_{b'}=1$ and $f_{b'}=C_{v,b'}f$.
\item $v=x_{i,c}$ or $v=y_{i,c}$ with $i\leq\eta$. Then $v$ occurs in no $\hat E_k$, and in $T(f)$ only in $f$. By
  Lemma~\ref{pd-kappa-defects-windows-lift-pairs-form-block-shift-operators}(3), $C_{v,n}u=\prod_k\hat E_k^{j_k}T(C_{v,n}f)$,
  and $C_{v,n}f$ is admissible with $n$, as in the first case. So $C_{v,n}u$ is a generator formed with $\mu'-n$.
\item $v=z_i$ for some $1\leq i\leq s-\eta$. Then $\hat E_i=z_iz'_i+\pi_i$; take $m=z'_i$, $E'=\pi_i$ and $q=-\pi_i$, so that
  $vm=\hat E_i+q$. The variable $z_i$ occurs in no $\hat E_k$ with $k\neq i$, and $T(f)=z_i^eT'$ with $e=\hat b_{u_i}\dot-o_{u_i}$
  and $T'$ a product in which $z_i$ does not occur. By
  Lemma~\ref{pd-kappa-defects-windows-lift-pairs-form-block-shift-operators}(3) and~(4), $C_{z_i,b'}T(f)=\binom e{b'}T(f)$.
  Step~2 applies with $c_{b'}=\binom e{b'}$ and $f_{b'}=f$.
\item $v=z'_i$ for some $1\leq i\leq s-\eta$. Then $\hat E_i=z'_iz_i+\pi_i$; take $m=z_i$, $E'=\pi_i$ and $q=-\pi_i$, so that
  $vm=\hat E_i+q$. The variable $z'_i$ occurs in no $\hat E_k$ with $k\neq i$ and not in $T(f)$, so by
  Lemma~\ref{pd-kappa-defects-windows-lift-pairs-form-block-shift-operators}(1) and~(3), $C_{z'_i,0}T(f)=T(f)$ and
  $C_{z'_i,b'}T(f)=0$ for $b'\geq1$. Step~2 applies with $c_0=1$, $c_{b'}=0$ for $b'\geq1$, and $f_{b'}=f$.
\end{itemize}
\end{proof}
% Lean: cth_gen D7 hT n hthr v hv hL hy hg, for any BlockData D with TieOK D (tieOK_blockBase for the block data of a
%   segment) and thresholds with thr k + n <= thr' k; the book's thresholds are the bounds formed with mu' and mu' - n.
%   T(f) is tailP O r D jp y g, "admissible with a" is y in YjtZ O (lettersAt j jp) J (thr + a), and the generators are
%   gen_mem. Step 1 is shift_mem (IsMult.pow, lettersAt_update, levelOK_update, gen_mem, SSt_mono), with IsMult for
%   x_{u_i,c} y_{u_i,c} (isMult_XX, using TieOK.vecE) and for -pi_i (IsMult.sum and IsMult.neg of pairP). Step 2 is tie_mem
%   (cth_lin, prodE_split, cth_mul_free with otherE_free, cth_mul, cth_pow, and for vm = E_i + q add_pow). The cases of
%   Step 3 are those of cth_gen on BVar: inl (a, c) with a = tieW or tieE (Ehat_eq, hXX, and Hc from tail_chain and
%   yjtZ_cth), inl (a, c) with a at no tie position (cth_mul_free, tail_chain, gen_mem with yjtZ_cth), and inr (i, false)
%   and inr (i, true), the z_i and z'_i (Ehat_eq', hPP, cth_X_pow with yjtZ_mono, respectively cth_free); OddV excludes
%   x_+, xi, z_+ and z'_+.

\begin{lemma}\label{pd-kappa-defects-windows-lift-pairs-form-block-shift}
Let the data of Definition~\ref{pd-kappa-defects-windows-lift-pairs-form-block-data} be given, with $1\leq j\leq q$ and a
kept block $b$ of $S_j(b_T)$. Let $c\geq0$ be an integer, $\Delta$ a finite set of exponent vectors $\tau$ on $V$ with
$|\tau|\leq4c$ and $\tau_v=0$ at every variable $v$ that is not odd, $\theta$ a family in $\mathcal R$, and $\mu$ a real
number. Then $G_{\Delta,\theta}(u)\in Z_j(b;\mu-4c)$ for every $u\in Z_j(b;\mu)$.
\end{lemma}
\leanref{Formalism/SharpRestriction/GrowthBounds/CoarseMultiplicity/GTheta.lean}{Formalism.MainNotions.ColumnSystems.Pf.gThetaAt}
\leanref{Formalism/SharpRestriction/GrowthBounds/CoarseMultiplicity/GTheta.lean}{Formalism.MainNotions.ColumnSystems.Pf.op_SSt}
\leanref{Formalism/SharpRestriction/GrowthBounds/CoarseMultiplicity/OtherEFree.lean}{Formalism.MainNotions.ColumnSystems.Pf.cth_SSt}

\begin{proof}
The module $Z_j(b;\mu')$ shrinks as $\mu'$ grows, since the bounds $M_l$ and $M_+$ of
Definition~\ref{pd-kappa-defects-segment-lattice}(5) decrease and the spaces $Y(\mathcal L,M)$ grow with $M$. As
$G_{\Delta,\theta}$ is a sum of the $\theta_\tau X^\tau D_\tau$, it suffices to show that $X^\tau D_\tau$ maps $Z_j(b;\mu)$ into
$Z_j(b;\mu-|\tau|)$, which lies in $Z_j(b;\mu-4c)$. For a variable $v$ write $e_v$ for its exponent vector, and put
$C_{v,n}=X^{ne_v}D_{ne_v}$. If $\tau_v=0$ and $n\geq1$, then $X^{\tau+ne_v}D_{\tau+ne_v}=C_{v,n}\circ X^\tau D_\tau$, as both
multiply $X^\nu$ by $\binom{\nu_v}n\binom\nu\tau$ (Lemma~\ref{pd-kappa-defects-windows-lift-pairs-form-block-monomials}(1)). By
induction on the number of variables at which $\tau$ is nonzero, it therefore suffices that $C_{v,n}$ maps $Z_j(b;\mu')$
into $Z_j(b;\mu'-n)$ for every odd $v$, every $n\geq0$ and every real $\mu'$. As $C_{v,n}$ is $\mathcal K$-linear and
$Z_j(b;\mu'-n)$ is an $\mathcal R$-module, it suffices to check this on the generators, which is
Lemma~\ref{pd-kappa-defects-windows-lift-pairs-form-block-shift-generators}. The argument does not use that $b$ is kept.
% Lean: gThetaAt (tieOK_blockBase, the thresholds through thr_shift) and op_SSt (induction on tau with op_single_add)
%   reduce to cth_SSt: cth v n = X^{n e_v} D_{n e_v} maps SSt O r D thr into SSt O r D thr' when thr + n <= thr', which is
%   cth_gen on each generator (Lemma pd-kappa-defects-windows-lift-pairs-form-block-shift-generators). GThetaAt assumes
%   that b is kept; gThetaAt does not use it. The RootWeights/ files behind yjtZ_cth (DerivativeValuationBounds and the
%   weight lemmas, about 2,400 lines) are cited through Lemmas pd-kappa-defects-paths-isotypic(3) and
%   pd-symbols-test-chains-projection-clauses(5)(a).
\end{proof}

\begin{lemma}\label{pd-kappa-defects-windows-lift-pairs-form-block-integral}
Let $p$ be a prime, $\mathcal R$ a discrete valuation ring with fraction field $\mathcal K$ and a uniformizer $\varpi_0$ with
$\varpi_0^2=p$, and $v$ the valuation of $\mathcal K$ with $v(p)=1$. Let $V$ be a finite set, $s\geq1$ an integer, and $u$
and $u'$ polynomials over $\mathcal K$ in the variables $V$ whose coefficients $u_\nu$ and $u'_\nu$ at every exponent vector
$\nu$ satisfy $v(u_\nu)+\frac12v_p(\nu!)+\frac12v_p(s)\geq0$ and $v(u'_\nu)+\frac12v_p(\nu!)+\frac12v_p(s)\geq0$.
\begin{enumerate}
\item $u_\nu u'_\nu\,\nu!\,s\in\mathcal R$ for every $\nu$.
\item If $N\geq1$ is an integer and $\nu!\,s$ divides $N$, then $\varpi_0^{v_p(N)}u_\nu\in\mathcal R$.
\end{enumerate}
\end{lemma}
\leanref{Formalism/SharpRestriction/GrowthBounds/CoarseMultiplicity/MemIntPolOfCoeff.lean}{Formalism.MainNotions.ColumnSystems.Pf.bb_pair}

\begin{proof}
As $p=\varpi_0^2$ lies in the maximal ideal of $\mathcal R$, every integer prime to $p$ is a unit of $\mathcal R$; so
$v(n)=v_p(n)$ for every integer $n\geq1$.

(1) Adding the two conditions, $v(u_\nu u'_\nu\nu!\,s)=v(u_\nu)+v(u'_\nu)+v_p(\nu!)+v_p(s)\geq0$.

(2) Write $N=p^eN'$ with $e=v_p(N)$ and $N'$ prime to $p$, and $N=\nu!\,s\,d$ with $d$ an integer, and put
$z=\varpi_0^eu_\nu$. Then $z^2N'=p^eN'u_\nu^2=u_\nu^2\,\nu!\,s\,d$, which lies in $\mathcal R$ by (1) with $u'=u$. As $N'$ is a
unit, $z^2\in\mathcal R$, so $v(z)=v(z^2)/2\geq0$ and $z\in\mathcal R$.
\end{proof}
% Lean: (1) is bb_pair (vK_coeff_add_half_padicVal_nonneg_of_bbSpec for the span bbSpec, padicValNat_le_vK_natCast,
%   which gives v(n) >= v_p(n)); (2) is the private exists_algebraMap_eq_of_w_pow_padicVal_mul in MemIntPolOfCoeff.lean
%   (natCast_isUnit_of_coprime, and integrality of z from z^2 through IsIntegrallyClosed).

\begin{lemma}\label{pd-kappa-defects-windows-lift-pairs-form-block-valuations}
Let $C_c\geq0$, $C_0\geq1$, $C_{E1},C_{E2}\geq1$ and $C_H\geq1000(C_{E1}+C_{E2}+1)(C_c+1)$. Let $p$, $n'$, $m$, $G'$, $N_1$,
$c_{cl}$, $D_p$, $c_*$, $T$, $Y$, $Y_b$, $b'$, $b_l$, $b_u$, $a$, $c\geq0$, $a_T$ and $b_T$ satisfy the hypotheses of
Lemma~\ref{pd-kappa-defects-windows-lift-pairs-source-order} with these constants, and form the numbers of
Definition~\ref{pd-kappa-defects-numbers} with these constants and any $C_m,C_A,C_t,C_\kappa>0$. Let $1\leq j\leq q$, $b$ a
kept block of $S_j(b_T)$, $\lambda_j$ and $\varphi_j$ as in Lemma~\ref{pd-kappa-defects-windows-lift-pairs-form-numbers}, and
$c_x(b)$, $\nu^0_b$ and $\Phi_b$ as in Definition~\ref{pd-kappa-defects-windows-lift-pairs-form-block-data}, which do not
depend on the Hida setting. Put $C_b=\prod_xc_x(b)!$. Then
\[
 v_p(\varphi_j)\leq k\lambda_j+v_p(C_b)+v_p(\Phi_b),\qquad v_p(C_b)+v_p(\Phi_b)\leq\lambda_j+v_p(\varphi_j).
\]
\end{lemma}
\leanref{Formalism/SharpRestriction/GrowthBounds/CoarseMultiplicity/LeadingExponentHomogeneous.lean}{Formalism.MainNotions.ColumnSystems.Pf.KapNumAt}
\leanref{Formalism/SharpRestriction/GrowthBounds/CoarseMultiplicity/CoarseMultAt.lean}{Formalism.MainNotions.ColumnSystems.Pf.kapNumAt_of}
\leanref{Formalism/SharpRestriction/GrowthBounds/CoarseMultiplicity/CoarseMultAt.lean}{Formalism.MainNotions.ColumnSystems.Pf.val_digit}
\leanref{Formalism/SharpRestriction/GrowthBounds/CoarseMultiplicity/Kummer.lean}{Formalism.MainNotions.ColumnSystems.Pf.expFact_mul}
\leanref{Formalism/SharpRestriction/GrowthBounds/CoarseMultiplicity/Kummer.lean}{Formalism.MainNotions.ColumnSystems.Pf.multinomial_val}
\leanref{Formalism/SharpRestriction/GrowthBounds/CoarseMultiplicity/CoarseMultAt.lean}{Formalism.MainNotions.ColumnSystems.Pf.TrueDegAt}
\leanref{Formalism/SharpRestriction/GrowthBounds/CoarseMultiplicity/CoarseMultAt.lean}{Formalism.MainNotions.ColumnSystems.Pf.CoarseMultAt}
\leanref{Formalism/SharpRestriction/GrowthBounds/CoarseMultiplicity/BFTrueDeg.lean}{Formalism.MainNotions.ColumnSystems.Pf.trueDegAt_of}
\leanref{Formalism/SharpRestriction/GrowthBounds/CoarseMultiplicity/BFTrueDeg.lean}{Formalism.MainNotions.ColumnSystems.Pf.tdeg_eq}
\leanref{Formalism/SharpRestriction/GrowthBounds/CoarseMultiplicity/WeightWLet.lean}{Formalism.MainNotions.ColumnSystems.Pf.TD_wLet}
\leanref{Formalism/SharpRestriction/GrowthBounds/CoarseMultiplicity/BFCoarse.lean}{Formalism.MainNotions.ColumnSystems.Pf.coarseMultAt_of}
\leanref{Formalism/SharpRestriction/GrowthBounds/CoarseMultiplicity/BFCoarse.lean}{Formalism.MainNotions.ColumnSystems.Pf.card_supp_le}
\leanref{Formalism/SharpRestriction/GrowthBounds/CoarseMultiplicity/BFCoarse.lean}{Formalism.MainNotions.ColumnSystems.Pf.cdeg_le_tdeg}

\begin{proof}
Write $\nu^0=\nu^0_b$, and for a letter $x$ let $t_x$ be the degree of the monomial $z^{\varepsilon_b(\nu^0)}$ in the variables
of $x$. The proof uses two facts:
\begin{enumerate}
\item[(a)] $\prod_xt_x!=\varphi_j$ and $\sum_x\lfloor\log_pt_x\rfloor=\lambda_j$, with $\lfloor\log_p0\rfloor=0$;
\item[(b)] $v_p(C_b)\leq v_p(\nu^0!)+\lambda_j$.
\end{enumerate}

\emph{Proof of (a).} By Definition~\ref{pd-kappa-defects-windows-lift-pairs-form-block-data}(2), the degree of
$z^{\varepsilon_b(\nu)}$ in the variables of a letter $x$ is $\sum_c\beta^b_{x,c}+\sum_{x(v)=x}P_v\nu_v$, for every exponent vector
$\nu$. Fix a letter $x$, put $\omega(v)=P_v$ if $x(v)=x$ and $\omega(v)=0$ otherwise, for $v\in V$, and
$\omega_L(\lambda)=\omega(\lambda_1)$ for the letters $\lambda$ of $\mathcal P_j$. Then $(\omega,\omega_L)$ is a balanced weight.
Indeed, $\iota_j$ sends the variables $\lambda_c$ of a letter $\lambda$ of $\mathcal P_j$ to variables of the one letter
$x(\lambda_1)$, with $P_{\lambda_c}=P_{\lambda_1}$, so $\omega(\lambda_c)=\omega_L(\lambda)$. It sends $z_l$ and $z'_l$ to variables
of the letters $x(x_{u_l,1})$ and $x(y_{u_l,1})$, the two letters at $y_{j,u_l}$, with
$P_{z_l}=P_{z'_l}=P_{x_{u_l,1}}=P_{y_{u_l,1}}=P_{u_l}$; so $\omega(z_l)+\omega(z'_l)$ and $\omega_L(x_{u_l})+\omega_L(y_{u_l})$ are
both $P_{u_l}$ if $x$ is one of these two letters, and both $0$ otherwise. Likewise $\omega(z_+)+\omega(z'_+)$ and
$\omega_L(x_+)+\omega_L(\xi)$ are both $1$ if $x$ is a letter at $s_{j+1}$, and both $0$ otherwise. As $\nu^0$ is homogeneous for
$b$ (Lemma~\ref{pd-kappa-defects-windows-lift-pairs-form-block-homogeneous}(3)),
\[
 t_x=\sum_c\beta^b_{x,c}+\sum_v\nu^0_v\,\omega(v)=\sum_c\beta^b_{x,c}+\Omega_b(\omega,\omega_L).
\]
Write $b_i=b_T(y_{j,i})$ and $\alpha_i=a_T(y_{j,i})$. The entries $\beta^b_{x,c}$ vanish at the letters $x$ that are not at one
of the $y_{j,i}$.
\begin{itemize}
\item $x=(y_{j,i},1)$. Here $\omega_L$ is $P_i$ at $x_i$ and $0$ at the other letters, $\omega(z_l)$ is $P_i$ if $u_l=i$ and $0$
  otherwise, and $\omega(z_+)=0$. So $t_x=|\rho_i|+P_i(\hat b_i\dot-w_i)+P_iw_i$ if $i\leq\eta$, and
  $t_x=|\rho_i|+P_i(\hat b_i\dot-o_i)+P_io_i$ if $i>\eta$. As $b$ is kept, $w_i\leq\hat b_i$, respectively $o_i\leq\hat b_i$, and
  $|\rho_i|\leq b_i$ with $P_i$ dividing $b_i-|\rho_i|$, so that $P_i\hat b_i=b_i-|\rho_i|$. Hence $t_x=b_T(y_{j,i})$.
\item $x=(y_{j,i},0)$. Here $\omega_L$ is $P_i$ at $y_i$ and $0$ at the other letters, and $\omega(z_l)=\omega(z_+)=0$. So
  $t_x=|\rho'_i|+P_i\varsigma_i+P_i(\hat\alpha_i\dot-\varsigma_i)$ if $i\leq\eta$, and $t_x=|\rho'_i|+P_i\hat\alpha_i$ if $i>\eta$. As
  $b$ is kept, $\varsigma_i\leq\hat\alpha_i$, and $|\rho'_i|\leq\alpha_i$ with $P_i$ dividing $\alpha_i-|\rho'_i|$. Hence
  $t_x=a_T(y_{j,i})$.
\item $x=(s_{j+1},1)$. Here $\omega_L$ is $1$ at $x_+$ and $0$ at the other letters, $\omega(z_+)=1$ and $\omega(z_l)=0$. So
  $t_x=(\alpha_+\dot-o_+)+o_+=a_T(s_{j+1})$, as $o_+\leq a_T(s_{j+1})=\alpha_+$ by
  Lemma~\ref{pd-kappa-defects-windows-lift-pairs-source-order}(2).
\item $x=(s_{j+1},0)$. Here $\omega_L$ is $1$ at $\xi$ and $0$ at the other letters, and $\omega(z_l)=\omega(z_+)=0$. So
  $t_x=b_+=b_T(s_{j+1})$.
\item At every other letter $x$, $\omega=0$, and so $t_x=0$.
\end{itemize}
The positions $y_{j,1},\dots,y_{j,S_j}$ and $s_{j+1}$ are distinct, and $0!=1$; so (a) follows from the definitions of
$\varphi_j$ and $\lambda_j$ in Lemma~\ref{pd-kappa-defects-windows-lift-pairs-form-numbers}.

\emph{Proof of (b).} For a letter $x$ let $V_x$ be the set of the $v\in V$ with $x(v)=x$. By
Definition~\ref{pd-kappa-defects-windows-lift-pairs-form-block-data}(6), $\nu^0$ vanishes except at the variables $\lambda_1$ of
the letters and at the $z_l$ and $z_+$, and $\iota_j$ sends the $\lambda_1$ to variables of the second column and the $z_l$ and
$z_+$ to variables of the first column. As distinct variables have distinct pairs $(x(v),c(v))$ (proof of
Lemma~\ref{pd-kappa-defects-windows-lift-pairs-form-block-monomials}(2)), at most two $v\in V_x$ have $\nu^0_v\neq0$. By
Lemma~\ref{pd-kappa-defects-windows-lift-pairs-form-block-homogeneous}(2), $c_x(b)=\sum_{v\in V_x}\nu^0_v$. So
$c_x(b)!=M_x\prod_{v\in V_x}\nu^0_v!$, where $M_x$ is $1$ or a binomial coefficient $\binom{c_x(b)}e$. By Legendre's formula,
$v_p\binom te=\sum_{i\geq1}\big(\lfloor t/p^i\rfloor-\lfloor e/p^i\rfloor-\lfloor(t-e)/p^i\rfloor\big)$, whose terms are $0$ or $1$
and vanish for $p^i>t$; so $v_p(M_x)\leq\lfloor\log_pc_x(b)\rfloor$. Also $c_x(b)\leq\sum_{v\in V_x}P_v\nu^0_v\leq t_x$, as $P_v\geq1$
and by the degree formula at the start of the proof of (a); so $v_p(M_x)\leq\lfloor\log_pt_x\rfloor$. Every $v\in V$ lies in one
$V_x$, so $C_b=\nu^0!\prod_xM_x$, and by (a), $v_p(C_b)\leq v_p(\nu^0!)+\sum_x\lfloor\log_pt_x\rfloor=v_p(\nu^0!)+\lambda_j$.
% Lean: (a) is TrueDegAt = trueDegAt_of restOPlus. Its tdeg_eq applies ioBlock_wh' to the monomial at mu0
%   (leading_exponent_weighted_homog) with the weight wLet x on the z_{x,c}; the pullback symLev * wLet x is the book's
%   omega, h1 to h3 of ioBlock_wh' are the three balance conditions, and weight_wLet reads off t_x. TD_wLet evaluates the
%   degree letter by letter (the digit conditions of Kept, w, sigma and o at most the coarse numbers, and
%   o_+ <= a_T(s_{j+1}) from OPlusAt = oPlusAt_of, Lemma pd-kappa-defects-windows-lift-pairs-source-order(2)); prod_tdeg
%   and its additive form sum_tdeg give (a). (b) is CoarseMultAt = coarseMultAt_of restOPlus: sum_fiber at mu0,
%   symbol_variable_first_columns and card_supp_le (with symVar_inj), Nat.multinomial_spec, multinomial_val (by
%   induction through Nat.factorization_choose_le_log, the bound that the book takes from Legendre's formula),
%   cdeg_le_tdeg, and sum_tdeg.

\emph{The factorials of the variables.} Let $\nu$ be an exponent vector on $V$. By
Lemma~\ref{pd-kappa-defects-windows-lift-pairs-form-block-monomials}(4), $P_v=p^{t_v+1}$ and $\beta_v<P_v$ at the odd variables
$v$, and $P_v=1$ and $\beta_v=0$ at the others. Apply Lemma~\ref{pd-symbols-test-chains-projection-clauses}(4) to the set
$V$, with $t'_v=t_v$, $\rho(v)=\beta_v$ and $e(v)=\nu_v$ at the odd $v$, and $t'_v=\rho(v)=e(v)=0$ at the others; then
$\rho(v)<p^{t'_v+1}$ for every $v$. By the first assertion of (4)(b), $u(e)\neq0$ and $v_p(u(e))=0$. At an odd $v$, (4)(a) gives
$(\beta_v+P_v\nu_v)!=\beta_v!\,\nu_v!\,(P_v!)^{\nu_v}c(P_v,\nu_v)\,d(P_v,\nu_v,\beta_v)$, where $c(P_v,\nu_v)\,d(P_v,\nu_v,\beta_v)$ is the
factor of $u(e)$ at $v$. At a $v$ that is not odd, (4)(a) with $t=\rho=e=0$ gives $c(p,0)\,d(p,0,0)=1$, which is the factor of
$u(e)$ at $v$, and $(\beta_v+P_v\nu_v)!=\nu_v!=\beta_v!\,\nu_v!\,(P_v!)^{\nu_v}$. Multiplying over $v$,
\[
 \prod_v(\beta_v+P_v\nu_v)!=u(e)\prod_v\beta_v!\,\nu_v!\,(P_v!)^{\nu_v}\quad\text{with }v_p(u(e))=0 .
\]
With $\nu=\nu^0$ and Lemma~\ref{pd-kappa-defects-windows-lift-pairs-form-block-monomials}(3), this gives
$v_p(\varepsilon_b(\nu^0)!)=v_p(\Phi_b)+v_p(\nu^0!)$.

\emph{The multinomial factor.} Each letter $x$ has the $k$ variables $z_{x,1},\dots,z_{x,k}$, so
$\prod_xt_x!=\varepsilon_b(\nu^0)!\cdot M$, where $M$ is the product over $x$ of the multinomial coefficients of $t_x$ with the
$k$ parts $\varepsilon_b(\nu^0)_{x,c}$. By Legendre's formula, $v_p$ of a multinomial coefficient of $t$ with $k$ parts
$t_1,\dots,t_k$ is $\sum_{i\geq1}\big(\lfloor t/p^i\rfloor-\sum_c\lfloor t_c/p^i\rfloor\big)$, whose terms lie between $0$ and
$k-1$ and vanish for $p^i>t$. So $0\leq v_p(M)\leq(k-1)\sum_x\lfloor\log_pt_x\rfloor=m\lambda_j$ by (a), and
$v_p(\varphi_j)=v_p(\Phi_b)+v_p(\nu^0!)+v_p(M)$.

\emph{The two bounds.} By Lemma~\ref{pd-kappa-defects-windows-lift-pairs-form-block-homogeneous}(3) and~(2), $\nu^0$ is
homogeneous for $b$ and $\nu^0!$ divides $C_b$; so $v_p(\nu^0!)\leq v_p(C_b)$. Hence
$v_p(\varphi_j)\leq v_p(\Phi_b)+v_p(C_b)+m\lambda_j$, which gives the first bound as $m\leq k$. By (b),
$v_p(C_b)+v_p(\Phi_b)\leq v_p(\nu^0!)+v_p(\Phi_b)+\lambda_j=v_p(\varepsilon_b(\nu^0)!)+\lambda_j\leq v_p(\varphi_j)+\lambda_j$, as
$v_p(M)\geq0$.
\end{proof}
% Lean: kapNumAt_of h52 hT hM, with RestKapNum = restKapNum_of h52 restOPlus; the thresholds are those of restOPlus
%   (oPlusAt_of under thr_std), which are those of Lemma pd-kappa-defects-windows-lift-pairs-source-order. hvd is
%   val_digit (the identity of (4)(a) and the unit of (4)(b), taken over all variables with t, rho and the exponent zero
%   at those that are not odd), hvE the first paragraph through expFact_ioExp, expFact_mul (multinomial_val) the second, hdv the
%   divisibility from dvdAt at mu0 (leading_exponent_homogeneous), and the two bounds are omega.

\begin{lemma}\label{pd-kappa-defects-windows-lift-pairs-form-block-factorial}
Let $C_c\geq0$, $C_0\geq1$, $C_{E1},C_{E2}\geq1$ and $C_H\geq1000(C_{E1}+C_{E2}+1)(C_c+1)$. Let $p$, $n'$, $m$, $G'$, $N_1$,
$c_{cl}$, $D_p$, $c_*$, $T$, $Y$, $Y_b$, $b'$, $b_l$, $b_u$, $a$, $c\geq0$, $a_T$ and $b_T$ satisfy the hypotheses of
Lemma~\ref{pd-kappa-defects-windows-lift-pairs-source-order} with these constants, form the numbers of
Definition~\ref{pd-kappa-defects-numbers} with these constants and any $C_m,C_A,C_t,C_\kappa>0$, and let $\lambda_j$ and
$z_j$ be as in Lemma~\ref{pd-kappa-defects-windows-lift-pairs-form-numbers}. Let an unramified quadratic Hida setting of
degree $n'$ at $p$ with rings $\mathcal R\subset\mathcal K$ and $\varpi_0$ be given, with the notation of
Definition~\ref{pd-kappa-defects-windows-lift-pairs-form-block-data}, and $v$ the valuation of $\mathcal K$ with $v(p)=1$.
Let $1\leq j\leq q$ and $b$ a kept block of $S_j(b_T)$. Then there are a nonzero $\kappa\in\mathcal R$ with
$v(\kappa)\leq(m+2)\lambda_j$, a finite set $\Delta$ of exponent vectors $\tau$ on $V$ with $|\tau|\leq4c$ and $\tau_v=0$ at
every variable $v$ that is not odd, and a family $\theta$ in $\mathcal R$, such that for every exponent vector $\nu$ homogeneous
for $b$ there is $h\in\mathcal R$ with
\[
 z_j\,p^{v_p(N_b)}\,\varepsilon_b(\nu)!\,e_{\Delta,\theta}(\nu)=\kappa\,\nu!\operatorname{sp}(b)\,(1+\varpi_0^ch).
\]
\end{lemma}
\leanref{Formalism/SharpRestriction/GrowthBounds/CoarseMultiplicity/BFDefs.lean}{Formalism.MainNotions.ColumnSystems.Pf.FactAt}
\leanref{Formalism/SharpRestriction/GrowthBounds/CoarseMultiplicity/CharZeroOfW.lean}{Formalism.MainNotions.ColumnSystems.Pf.factAt_of}
\leanref{Formalism/SharpRestriction/GrowthBounds/CoarseMultiplicity/CharZeroOfW.lean}{Formalism.MainNotions.ColumnSystems.Pf.KapAt}
\leanref{Formalism/SharpRestriction/GrowthBounds/CoarseMultiplicity/LeadingExponentHomogeneous.lean}{Formalism.MainNotions.ColumnSystems.Pf.kapAt_of}
\leanref{Formalism/SharpRestriction/GrowthBounds/CoarseMultiplicity/LeadingExponentHomogeneous.lean}{Formalism.MainNotions.ColumnSystems.Pf.kap_transport}
\leanref{Formalism/SharpRestriction/GrowthBounds/CoarseMultiplicity/FactAbs.lean}{Formalism.MainNotions.ColumnSystems.Pf.fact_abs}

\begin{proof}
Write $\varphi_j$ as in Lemma~\ref{pd-kappa-defects-windows-lift-pairs-form-numbers}, so that $z_j=p^{k\lambda_j}/\varphi_j$, and
$C_b=\prod_xc_x(b)!$, so that $N_b=C_b\operatorname{sp}(b)$. As $p=\varpi_0^2$ lies in the maximal ideal of $\mathcal R$, the
integers prime to $p$ are units of $\mathcal R$.

\emph{The scalar.} Put $A=\Phi_b$, $B=\varphi_j\operatorname{sp}(b)$, $E=k\lambda_j+v_p(N_b)$ and $D=(m+2)\lambda_j$. As
$v_p(N_b)=v_p(C_b)+v_p(\operatorname{sp}(b))$ and $m+2=k+1$, the two bounds of
Lemma~\ref{pd-kappa-defects-windows-lift-pairs-form-block-valuations} give $v_p(B)\leq E+v_p(A)\leq D+v_p(B)$. Write
$A=p^{v_p(A)}A'$ and $B=p^{v_p(B)}B'$ with $A'$ and $B'$ prime to $p$, hence units of $\mathcal R$. Put $n=E+v_p(A)-v_p(B)$, so
$0\leq n\leq D$, and $\kappa=p^nA'B'^{-1}$. Then $\kappa\neq0$, $v(\kappa)=n\leq(m+2)\lambda_j$, and $p^EA=\kappa B$. Let $\nu$ be
homogeneous for $b$. By Lemma~\ref{pd-kappa-defects-windows-lift-pairs-form-block-homogeneous}(3), applied to $\nu$ and
$\nu^0_b$ with $f(P)=P!$, $R_b\prod_v\beta_v!\,(P_v!)^{\nu_v}=\Phi_b$. Hence
\[
 z_j\,p^{v_p(N_b)}R_b\prod_v\beta_v!\,(P_v!)^{\nu_v}=\frac{p^EA}{\varphi_j}=\kappa\operatorname{sp}(b).\qquad(*)
\]

\emph{The operator.} By Lemma~\ref{pd-kappa-defects-windows-lift-pairs-form-block-monomials}(4), $P_v=p^{t_v+1}$ and
$\beta_v<P_v$ at the odd variables $v$, and $P_v=1$ and $\beta_v=0$ at the others. Apply
Lemma~\ref{pd-symbols-test-chains-projection-clauses}(4)(b) to the set $V$, with $t'_v=t_v$ and $\rho(v)=\beta_v$ at the odd
$v$ and $t'_v=\rho(v)=0$ at the others, the exponent $-1$ and $c$. It gives a finite set of $\tau\colon V\to\mathbf Z_{\geq0}$ with
$|\tau|\leq4c$ and $q_\tau\in\mathbf Q$ with $v_p(q_\tau)\geq0$ such that
$v_p\big(f(e)^{-1}-\sum_\tau q_\tau\binom e\tau\big)\geq c$ for every $e\colon V\to\mathbf Z_{\geq0}$, and $v_p(f(e))=0$, as
$f(e)=u(e)$ by (4)(a). Let $\Delta$ be the set of these $\tau$ with $\tau_v=0$ at every $v$ that is not odd, and
$\theta_\tau=q_\tau$, which lies in $\mathcal R$. Let $\nu$ be an exponent vector, and $e$ the vector equal to $\nu$ at the odd
variables and $0$ at the others. A $\tau$ not in $\Delta$ has $\binom e\tau=0$, and for $\tau\in\Delta$, $\binom e\tau=\binom\nu\tau$; so
$g=\sum_\tau q_\tau\binom e\tau=e_{\Delta,\theta}(\nu)$. Put $h_0=f(e)g-1=-f(e)\big(f(e)^{-1}-g\big)$; then $v_p(h_0)\geq c$. The
factors of $f(e)$ at the variables that are not odd are $1$, and $(\beta_v+P_v\nu_v)!=\beta_v!\,\nu_v!\,(P_v!)^{\nu_v}$ there. So
\[
 \prod_v(\beta_v+P_v\nu_v)!\cdot e_{\Delta,\theta}(\nu)=\prod_v\beta_v!\,(P_v!)^{\nu_v}\cdot\nu!\,(1+h_0).\qquad(**)
\]
As $v_p(h_0)\geq c$ and $p^c=\varpi_0^{2c}$, $h_0=\varpi_0^ch$ with $h=\varpi_0^cp^{-c}h_0\in\mathcal R$.

\emph{Conclusion.} Let $\nu$ be homogeneous for $b$. By
Lemma~\ref{pd-kappa-defects-windows-lift-pairs-form-block-monomials}(3), $(**)$ and $(*)$,
\[
 z_jp^{v_p(N_b)}\varepsilon_b(\nu)!\,e_{\Delta,\theta}(\nu)=z_jp^{v_p(N_b)}R_b\prod_v\beta_v!\,(P_v!)^{\nu_v}\cdot\nu!\,(1+\varpi_0^ch)
 =\kappa\,\nu!\operatorname{sp}(b)\,(1+\varpi_0^ch).
\]
\end{proof}
% Lean: factAt_of h52 (kapAt_of hKa). kapAt_of is the scalar: kap_transport (nat_split for A' and B') with A = PhiB,
%   B = facSeg * specConst, E = (m + 1) lgSeg + v_p(Nblk), D = (m + 2) lgSeg and KapNumAt (Lemma
%   pd-kappa-defects-windows-lift-pairs-form-block-valuations), and lev_const at mu0 for (*). factAt_of is the operator:
%   oddVar_data and endVar_data, fact_abs h52 (FactorialUnitDecomposition, with inv = true, filtered to the tau that
%   vanish off the odd variables), rat_int for q_tau and h in O, and expFact_ioExp; the h of the statement is w^c y.

\begin{lemma}\label{pd-kappa-defects-windows-lift-pairs-form-block}
For every $C_c\geq0$ there is $T_0>0$ such that for every $C_0\geq T_0$ there is $T_1>0$ such that for all
$C_{E1},C_{E2},C_m\geq T_1$ there is $T_2>0$ with the following property for every $C_H\geq T_2$. Let $p$, $n'$, $m$, $G'$,
$N_1$, $c_{cl}$, $D_p$, $c_*$, $T$, $Y$, $Y_b$, $b'$, $b_l$, $b_u$, $a$, $c\geq1$, $a_T$ and $b_T$ satisfy the hypotheses of
Lemma~\ref{pd-kappa-defects-windows-lift-pairs-source} with these constants, form the numbers of
Definition~\ref{pd-kappa-defects-numbers} with these constants and any $C_A,C_t,C_\kappa>0$, put $\mu_0=13kN_1Qc+12c$, and let
$\lambda_j$ and $z_j$ be as in Lemma~\ref{pd-kappa-defects-windows-lift-pairs-form-numbers}. Let an unramified quadratic Hida
setting of degree $n'$ at $p$ with rings $\mathcal R\subset\mathcal K$ and $\varpi_0$ be given, with the notation of
Definition~\ref{pd-kappa-defects-windows-lift-pairs-model}, and $v$ the valuation of $\mathcal K$ with $v(p)=1$. Let
$1\leq j\leq q$ and $b$ a kept block of $S_j(b_T)$, and let $L_b\subset\mathcal P_j$ be its block lattice and
$\operatorname{sp}(b)$ its number, formed over $\mathcal R$ with $\alpha_i=a_T(y_{j,i})$, $\alpha_+=a_T(s_{j+1})$, $\mu=m_0$ and
$r=m$ (Definition~\ref{pd-kappa-defects-segment-lattice}(1) and~(6)). Then there are a nonzero $\kappa\in\mathcal R$ with
$v(\kappa)\leq(m+2)\lambda_j$ and maps $\psi,\psi'\colon\mathcal P_j\to\mathcal P_j$ with $\psi(u),\psi'(u)\in Z_j(b;\mu_0)$
for $u\in L_b$, such that for all $u,u'\in L_b$ and $a\in\mathcal R$ with $\operatorname{sp}(b)[u,u']=a$ there is
$x\in\mathcal R$ with
\[
 z_j\big[\iota_j(b,\psi(u)),\iota_j(b,\psi'(u'))\big]=\kappa\,(a+\varpi_0^cx).
\]
\end{lemma}
\leanref{Formalism/SharpRestriction/GrowthBounds/CoarseMultiplicity/RestBlockForm.lean}{Formalism.MainNotions.ColumnSystems.Pf.BlockFormAt}
\leanref{Formalism/SharpRestriction/GrowthBounds/CoarseMultiplicity/RestBlockForm.lean}{Formalism.MainNotions.ColumnSystems.Pf.RestBlockForm}
\leanref{Formalism/SharpRestriction/GrowthBounds/CoarseMultiplicity/BlockForm.lean}{Formalism.MainNotions.ColumnSystems.Pf.restBlockForm_of}
\leanref{Formalism/SharpRestriction/GrowthBounds/CoarseMultiplicity/MemIntPolOfCoeff.lean}{Formalism.MainNotions.ColumnSystems.Pf.blockFormAt_of}

\begin{proof}
\emph{The constants.} Let $C_c\geq0$, and put $T_0=1$, $T_1=17$ and $T_2=1000(C_{E1}+C_{E2}+1)(C_c+1)$. Let $C_0\geq T_0$,
$C_{E1},C_{E2},C_m\geq T_1$ and $C_H\geq T_2$; then the constants satisfy the hypotheses of
Lemmas~\ref{pd-kappa-defects-windows-lift-pairs-form-block-valuations}
and~\ref{pd-kappa-defects-windows-lift-pairs-form-block-factorial}. Let the data, the setting, $j$ and $b$ be as in the
statement, and use the notation of Definition~\ref{pd-kappa-defects-windows-lift-pairs-form-block-data}. By hypothesis (1) of
Theorem~\ref{pd-gain-theorem-one}, $N_1\geq1$; also $k\geq1$, $Q=p^{G'}\geq1$ and $c\geq1$. As $C_m\geq17$,
$m_0-\mu_0=(C_m-13)kN_1Qc\geq4kN_1Qc\geq4c$, so $m_0-4c\geq\mu_0$.

\emph{The maps.} Lemma~\ref{pd-kappa-defects-windows-lift-pairs-form-block-factorial} gives a nonzero $\kappa\in\mathcal R$
with $v(\kappa)\leq(m+2)\lambda_j$, $\Delta$ and $\theta$. Put $e=v_p(N_b)$, where $N_b\geq1$, and for $u\in\mathcal P_j$ put
$\psi(u)=\varpi_0^eG_{\Delta,\theta}(u)$ and $\psi'(u)=\varpi_0^eu$.

\emph{Integrality.} Let $u\in L_b$, and let $\nu$ be an exponent vector with $u_\nu\neq0$. As $u$ lies in the $\mathcal K$-span
of the generators of $b$ formed with $m_0$, $\nu$ is homogeneous for $b$ by
Lemma~\ref{pd-kappa-defects-windows-lift-pairs-form-block-homogeneous}(1), and $\nu!\operatorname{sp}(b)$ divides $N_b$ by
part (2) of that lemma. The coefficients of $u$ satisfy the condition of
Definition~\ref{pd-kappa-defects-segment-lattice}(6), so
Lemma~\ref{pd-kappa-defects-windows-lift-pairs-form-block-integral}(2), with $s=\operatorname{sp}(b)$ and $N=N_b$, gives
$\varpi_0^eu_\nu\in\mathcal R$. So $\varpi_0^eu_\nu\in\mathcal R$ for every $\nu$.

\emph{The images.} The module $Z_j(b;\mu)$ shrinks as $\mu$ grows, since the bounds $M_l$ and $M_+$ of
Definition~\ref{pd-kappa-defects-segment-lattice}(5) decrease and the spaces $Y(\mathcal L,M)$ grow with $M$. By
Lemma~\ref{pd-kappa-defects-windows-lift-pairs-form-block-shift} with $\mu=m_0$, $G_{\Delta,\theta}$ maps $Z_j(b;m_0)$ into
$Z_j(b;m_0-4c)\subset Z_j(b;\mu_0)$, and $Z_j(b;m_0)\subset Z_j(b;\mu_0)$. Let $u\in L_b$; it lies in the $\mathcal K$-span of
$Z_j(b;m_0)$. As $G_{\Delta,\theta}$ is $\mathcal K$-linear, $\psi(u)$ and $\psi'(u)$ lie in the $\mathcal K$-span of
$Z_j(b;\mu_0)$. By Lemma~\ref{pd-kappa-defects-windows-lift-pairs-form-block-monomials}(1), their coefficients at $X^\nu$ are
$e_{\Delta,\theta}(\nu)\varpi_0^eu_\nu$ and $\varpi_0^eu_\nu$, which lie in $\mathcal R$ by the integrality. So $\psi(u)$ and
$\psi'(u)$ lie in $Z_j(b;\mu_0)$ by Lemma~\ref{pd-kappa-defects-windows-lift-pairs-form-block-saturated}.

\emph{The pairing.} Let $u,u'\in L_b$ and $a\in\mathcal R$ with $\operatorname{sp}(b)[u,u']=a$. By
Lemma~\ref{pd-kappa-defects-windows-lift-pairs-form-block-monomials}(1), every monomial of $\psi(u)$ is a monomial of $u$; so by
part (2) of that lemma, as $\varpi_0^{2e}=p^e$,
\[
 z_j\big[\iota_j(b,\psi(u)),\iota_j(b,\psi'(u'))\big]=\sum_\nu z_j\,p^e\,e_{\Delta,\theta}(\nu)\,u_\nu u'_\nu\,\varepsilon_b(\nu)!,
\]
the sum over the $\nu$ with $u_\nu\neq0$. For such $\nu$, $\nu$ is homogeneous for $b$;
Lemma~\ref{pd-kappa-defects-windows-lift-pairs-form-block-factorial} gives $h_\nu\in\mathcal R$ with
$z_jp^e\varepsilon_b(\nu)!\,e_{\Delta,\theta}(\nu)=\kappa\,\nu!\operatorname{sp}(b)(1+\varpi_0^ch_\nu)$, and
Lemma~\ref{pd-kappa-defects-windows-lift-pairs-form-block-integral}(1) gives
$y_\nu=u_\nu u'_\nu\,\nu!\operatorname{sp}(b)\in\mathcal R$. So the term at $\nu$ is $\kappa y_\nu+\kappa\varpi_0^ch_\nu y_\nu$. As
$[u,u']=\sum_\nu\nu!\,u_\nu u'_\nu$ (Definition~\ref{pd-weight-lattices}(2)), $\sum_\nu y_\nu=a$, and the sum is
$\kappa(a+\varpi_0^cx)$ with $x=\sum_\nu h_\nu y_\nu\in\mathcal R$.
\end{proof}
% Lean: restBlockForm_of h52 hK = blockFormAt_of h17 hstd ioFockAt eigAt homAt dvdAt satAt (gThetaAt D)
%   (factAt_of h52 (kapAt_of hKa)), with the thresholds of thr_std, of h17 (1, 17 and 1, for Cm >= 17) and of
%   RestKapNum = restKapNum_of h52 restOPlus. In blockFormAt_of: hN1 from Std, hmle is m_0 - 4c >= mu_0, hint is the
%   integrality (homAt, dvdAt, bb_pair and exists_algebraMap_eq_of_w_pow_padicVal_mul), the images use
%   gOp_mem_span_of_mem_span with gThetaAt, ssz_antitone_of_le, mem_intPol_of_coeff and satAt, and the pairing is key
%   (hfac from FactAt, bb_pair) summed through ioFockAt over the support of u (hsub). The Lean's s and s' are
%   sc . gOp T q and sc ., with sc = w^v, v = v_p(Nblk).
% Size: no part of this proof is left open. Lemma pd-kappa-defects-windows-lift-pairs-form-block-shift (about 4,000 lines,
%   of which about 2,400 are RootWeights/ files cited through Lemmas pd-kappa-defects-paths-isotypic(3) and
%   pd-symbols-test-chains-projection-clauses(5)(a)) and facts (a) and (b) in the proof of Lemma
%   pd-kappa-defects-windows-lift-pairs-form-block-valuations (about 1,000 lines) are written out above.

\begin{lemma}\label{pd-kappa-defects-windows-lift-pairs-form}
For every $C_c\geq0$ there is $T_0>0$ such that for every $C_0\geq T_0$ there is $T_1>0$ such that for
all $C_{E1},C_{E2},C_m\geq T_1$ there is $T_2>0$ with the following property for every $C_H\geq T_2$.
Let $p$, $n'$, $m$, $G'$, $N_1$, $c_{cl}$, $D_p$, $c_*$, $T$, $Y$, $Y_b$, $b'$, $b_l$, $b_u$, $a$, $c\geq1$, $a_T$ and $b_T$ satisfy
the hypotheses of Lemma~\ref{pd-kappa-defects-windows-lift-pairs-source} with these constants, form the numbers of
Definition~\ref{pd-kappa-defects-numbers} with these constants and any $C_A,C_t,C_\kappa>0$, and put
$\mu_0=13kN_1Qc+12c$. Let an unramified quadratic Hida setting of degree $n'$ at $p$ with rings
$\mathcal R\subset\mathcal K$ and $\varpi_0$ be given, with the notation of
Definition~\ref{pd-kappa-defects-windows-lift-pairs-model}. For $1\leq j\leq q$ let $L_j$ be the segment lattice of the
$j$-th segment and $\beta_j$ its pairing (Definition~\ref{pd-kappa-defects-numbers}). For a tuple $b=(b_1,\dots,b_q)$, with
$b_j$ a block of $S_j(b_T)$, a tuple $u=(u_1,\dots,u_q)$ with $u_j\in L_j$ is \emph{carried by} $b$ if, for each $j$, the
entries of $u_j$ at the blocks other than $b_j$ are $0$. Then there are, for every tuple $b$ of blocks, a nonzero
$\kappa_b\in\mathcal R$ with $v(\kappa_b)\leq4\kappa_0$, and, for every tuple $u$ carried by $b$, tuples $f(b,u)$ and
$f'(b,u)$ of elements of the $\mathcal P_j$, with the following properties.
\begin{enumerate}
\item If every $b_j$ is kept, then $f(b,u)_j$ and $f'(b,u)_j$ lie in $Z_j(b_j;\mu_0)$ for every $j$.
\item Let every $b_j$ and every $b'_j$ be kept, let $u$ be carried by $b$ and $u'$ by $b'$, with
  $\beta_j(u_j,u'_j)=a_j\in\mathcal R$ for every $j$, and let $P,P'\in\mathrm{Pol}_T$ satisfy
  $P-\Phi(b,f(b,u))\in\varpi_0^c\mathcal R[z]$ and $P'-\Phi(b',f'(b',u'))\in\varpi_0^c\mathcal R[z]$. Then
  \[
   z_T[P,P']-\delta_{bb'}\kappa_b\prod_{j=1}^qa_j\in\varpi_0^c\mathcal R .
  \]
\end{enumerate}
\end{lemma}
\leanref{Formalism/SharpRestriction/GrowthBounds/CoarseMultiplicity/RestBlockForm.lean}{Formalism.MainNotions.ColumnSystems.Pf.FormAt}
\leanref{Formalism/SharpRestriction/GrowthBounds/CoarseMultiplicity/RestBlockForm.lean}{Formalism.MainNotions.ColumnSystems.Pf.RestForm}
\leanref{Formalism/SharpRestriction/GrowthBounds/CoarseMultiplicity/BFFinal.lean}{Formalism.MainNotions.ColumnSystems.Pf.restForm_of_oplus}
\leanref{Formalism/SharpRestriction/GrowthBounds/CoarseMultiplicity/RestBlockForm.lean}{Formalism.MainNotions.ColumnSystems.Pf.KeptAll}

\begin{proof}
\emph{The constants.} Let $C_c\geq0$, and let $T_0$, $T_1$ and $T_2$ be the largest of the thresholds of
Lemma~\ref{pd-kappa-defects-windows-lift-pairs-form-block} and of $1$, $1$ and $1000(C_{E1}+C_{E2}+1)(C_c+1)$, which are those
of Lemmas~\ref{pd-kappa-defects-margin} and~\ref{pd-kappa-defects-windows-lift-pairs-source-order}. Let $C_H\geq T_2$, and let
the data and the setting be as in the statement. The facts (a) to (g) in the proof of Lemma~\ref{pd-kappa-defects-margin}
hold; they give $q\geq2$ and $\log a\geq1$, so $a\geq e$ and, as $a$ is an integer, $a\geq3$. Let $s\in T$.
\begin{itemize}
\item As in fact (j) of the proof of Lemma~\ref{pd-kappa-defects-windows-model-step}, $a_T(s)\leq a(1+2/(kY))$. As
  $kY\geq1$ and $a\geq3$, $a_T(s)\leq3a\leq a^2$.
\item Put $\beta_e=(b_u-b_l)/b'+2\Lambda/k$. As $b_l\leq b'$ by hypothesis (2) of Theorem~\ref{pd-gain-theorem-one},
  $b_T(s)\leq b_u+2J_0\leq b'+(b_u-b_l)+2\Lambda b'/k=b'(1+\beta_e)$. As $C_H\geq1000$ and $N_1\geq1$,
  $\beta_e\leq\delta/2-18N_1\Lambda/k\leq\delta/2$ for the number $\delta$ of that theorem, and $\delta<1/6$ by (2); so
  $\beta_e\leq1$ and $b_T(s)\leq2b'$. By (2), $pb'<a$, and $p\geq5$; so $2b'\leq a$, and $b_T(s)\leq a\leq a^2$.
\end{itemize}
So Lemma~\ref{pd-kappa-defects-windows-lift-pairs-form-numbers}(3) gives $\sum_j(m+2)\lambda_j\leq4\kappa_0$.

\emph{The scalars and the symbols.} For each $j$ and each block $b$ of $S_j(b_T)$: if $b$ is kept, let $\kappa_{j,b}$,
$\psi_{j,b}$ and $\psi'_{j,b}$ be given by Lemma~\ref{pd-kappa-defects-windows-lift-pairs-form-block}; otherwise put
$\kappa_{j,b}=1$ and let $\psi_{j,b}$ and $\psi'_{j,b}$ be the identity. For a tuple $b=(b_j)$ put
$\kappa_b=\prod_j\kappa_{j,b_j}$. It is nonzero, as $\mathcal R$ is a domain, and
$v(\kappa_b)\leq\sum_j(m+2)\lambda_j\leq4\kappa_0$, as $v(\kappa_{j,b})\leq(m+2)\lambda_j$ in both cases. For $u$ carried by $b$
let $u_j(b_j)\in\mathcal P_j$ be the entry of $u_j$ at $b_j$, and put $f(b,u)_j=\psi_{j,b_j}(u_j(b_j))$ and
$f'(b,u)_j=\psi'_{j,b_j}(u_j(b_j))$.

(1) If $b_j$ is kept, the entry $u_j(b_j)$ of $u_j\in L_j$ lies in the block lattice $L_{b_j}$
(Definition~\ref{pd-kappa-defects-segment-lattice}(7)), so $f(b,u)_j$ and $f'(b,u)_j$ lie in $Z_j(b_j;\mu_0)$ by
Lemma~\ref{pd-kappa-defects-windows-lift-pairs-form-block}.

(2) Let $b$, $b'$, $u$, $u'$, the $a_j$, $P$ and $P'$ be as in (2), and put $\sigma=\Phi(b,f(b,u))$ and
$\sigma'=\Phi(b',f'(b',u'))$. By (1), Lemma~\ref{pd-kappa-defects-windows-lift-pairs-source-order}(2) and
Lemma~\ref{pd-kappa-defects-windows-lift-pairs-source-symbols}(4), $\sigma$ and $\sigma'$ lie in
$\mathcal R[z]\cap\mathrm{Pol}_T$. So $P-\sigma$ and $P'-\sigma'$ lie in $\varpi_0^c\mathcal R[z]\cap\mathrm{Pol}_T$, and
$P\in\mathcal R[z]$. The Fock pairing is bilinear and symmetric, so
\[
 [P,P']=[\sigma,\sigma']+[P-\sigma,\sigma']+[P'-\sigma',P],
\]
and by Lemma~\ref{pd-kappa-defects-windows-lift-pairs-form-numbers}(2) the last two terms, multiplied by $z_T$, lie in
$\varpi_0^c\mathcal R$. By Lemma~\ref{pd-kappa-defects-windows-lift-pairs-form-numbers}(1) and
Lemma~\ref{pd-kappa-defects-windows-lift-pairs-source-symbols}(4),
\[
 z_T[\sigma,\sigma']=\prod_{j=1}^qz_j\big[\iota_j(b_j,f(b,u)_j),\iota_j(b'_j,f'(b',u')_j)\big].
\]
\begin{itemize}
\item If $b=b'$, then for each $j$ the entries of $u_j$ and $u'_j$ at the blocks other than $b_j$ are $0$, so
  $a_j=\beta_j(u_j,u'_j)=\operatorname{sp}(b_j)[u_j(b_j),u'_j(b_j)]$, with $u_j(b_j),u'_j(b_j)\in L_{b_j}$.
  Lemma~\ref{pd-kappa-defects-windows-lift-pairs-form-block} gives $x_j\in\mathcal R$ with $j$-th factor
  $\kappa_{j,b_j}(a_j+\varpi_0^cx_j)$. Expanding, $\prod_j(a_j+\varpi_0^cx_j)=\prod_ja_j+\varpi_0^cy$ with $y\in\mathcal R$,
  so $z_T[\sigma,\sigma']=\kappa_b\prod_ja_j+\varpi_0^c\kappa_by$.
\item If $b\neq b'$, then $b_{j_0}\neq b'_{j_0}$ for some $j_0$, both kept, and the $j_0$-th factor is $0$ by
  Lemma~\ref{pd-kappa-defects-windows-lift-pairs-source-symbols}(3); so $z_T[\sigma,\sigma']=0$.
\end{itemize}
In both cases $z_T[P,P']-\delta_{bb'}\kappa_b\prod_ja_j\in\varpi_0^c\mathcal R$.
\end{proof}
% Lean: RestForm = restForm_of_oplus (lemma_five_two_first_holds D) restOPlus = symbol_form_of_orthogonality
%   (restBlockForm_of h52 hK) orthAt restOPlus = restForm_of' hB orthAt prodAt (restSymPol_of restOPlus) intAt numAt
%   zProdAt, whose thresholds are those of thr_std, RestBlockForm and RestSymPol (restOPlus, under thr_std). The proof is
%   formAt_of: hex (kappa = 1 and the identity at blocks that are not kept), vLEO_prod with NumAt (h5), famOK_mem for
%   (1); for (2) the symbols are integral (symIntAt) and in Pol (SymPolAt from OPlusAt), hdec splits fock P Q, IntAt
%   (h4) bounds the two error terms, and hmain uses ProdAt (h2) and ZProdAt (h6), with segForm_single and
%   prod_add_mul when beta = beta', and OrthAt (h1) at j0 otherwise. The constants paragraph follows numAt, which takes
%   a_T, b_T <= a^2 from Std (std_of with H = 1): loga gives a >= 3 (Real.exp_one_gt_d9), winA (= hwinA of std_of, the
%   book's fact (j)) gives a_T <= 3a, winB (hwinB, from b_l <= b' and eB1) with be <= 1 (beLe, de6) and pab give
%   b_T <= 2 b' <= a; be is the book's beta_e.

\begin{lemma}\label{pd-kappa-defects-windows-lift-pairs}
For every $C_c\geq0$ there is $T_0>0$ such that for every $C_0\geq T_0$ there is $T_1>0$ such that for
all $C_{E1},C_{E2},C_m\geq T_1$ there is $T_2>0$ such that for every $C_H\geq T_2$ there is $C^0_t>0$
with the following property for every $C_t\geq C^0_t$. Let $p\geq5$ be a prime and $n'\geq2$ an even
integer, and let $m$, $G'$, $N_1$, $c_{cl}$, $D_p$, $c_*$, $T$, $Y$, $Y_b$, $b'$, $b_l$, $b_u$ and an
integer $a$ satisfy hypotheses (1) to (3) of Theorem~\ref{pd-gain-theorem-one} with $C_0$ and $C_H$.
Let $c\geq1$ be an integer with $c\leq C_c(c_*+|T|k\log_pa+|T|c_{cl}+n'c_{cl})$, and let
$a_T,b_T\colon\mathbf Z/n'\to\mathbf Z_{\geq0}$ satisfy $a\leq a_T(s)\leq a+a/(kY)+2J_0$ and
$b_l\leq b_T(s)\leq b_u+2J_0$ for $s\in T$, where $J_0=\Lambda b'/k$. Form the numbers of
Definition~\ref{pd-kappa-defects-numbers} with these constants and any $C_A,C_\kappa>0$, and put
$t=\lfloor t_0\rfloor$. Let an unramified quadratic Hida setting of degree $n'$ at $p$ with rings
$\mathcal R\subset\mathcal K$ and $\varpi_0$ be given, and an edge $e$. For $1\leq j\leq q$ let $L_j$ be
the segment lattice of the $j$-th segment and $\beta_j$ its pairing. For a tuple $b=(b_1,\dots,b_q)$,
with $b_j$ a block of the segment datum $S_j(b_T)$, a tuple $u=(u_1,\dots,u_q)$ with $u_j\in L_j$ is
\emph{carried by} $b$ if, for each $j$, the entries of $u_j$ at the blocks other than $b_j$ are $0$.
Then there are, for every tuple $b$ of blocks, a nonzero $\kappa_b\in\mathcal R$ with
$v(\kappa_b)\leq4\kappa_0$, and, for every tuple $u$ carried by $b$, elements
$F(b,u)\in\mathcal C^\times_{t,c}(e)$ and $G(b,u)\in\mathcal C'^\times_{t,c}(e)$, with the following
property. If $u$ is carried by $b$, $u'$ is carried by $b'$, and $\beta_j(u_j,u'_j)=a_j\in\mathcal R$
for every $j$, then
\[
 \langle F(b,u),G(b',u')\rangle^\times-\delta_{bb'}\kappa_b\prod_{j=1}^qa_j\in\varpi_0^c\mathcal R
\]
(Definition~\ref{pd-kappa-product-space}).
\end{lemma}
\leanref{Formalism/SharpRestriction/GrowthBounds/BinomialRatio/PairLift.lean}{Formalism.MainNotions.ColumnSystems.Pf.RestPairLift}
\leanref{Formalism/SharpRestriction/GrowthBounds/BinomialRatio/PairLift.lean}{Formalism.MainNotions.ColumnSystems.Pf.PairLiftAt}
\leanref{Formalism/SharpRestriction/GrowthBounds/BinomialRatio/PairLift.lean}{Formalism.MainNotions.ColumnSystems.Pf.FamOK}
\leanref{Formalism/SharpRestriction/GrowthBounds/CoarseMultiplicity/PairedLiftAtOfIdentity.lean}{Formalism.MainNotions.ColumnSystems.Pf.paired_lifts_of_identity}

\begin{proof}
\emph{The constants.} Let $C_c\geq0$, and let $T_0$, $T_1$ and $T_2$ be the largest of the thresholds of
Lemmas~\ref{pd-kappa-defects-windows-lift-pairs-source} and~\ref{pd-kappa-defects-windows-lift-pairs-form}; for
$C_H\geq T_2$ put $C^0_t=30$. Let $C_t\geq30$, let the data, the setting and the edge $e$ be as in the statement, use the
notation of Definition~\ref{pd-kappa-defects-windows-lift-pairs-model}, and put $\mu_0=13kN_1Qc+12c$.

\emph{The lifts.} Lemma~\ref{pd-kappa-defects-windows-lift-pairs-flat} gives $h$, $h'$, $\Theta$, $\Theta'$ and $\Xi$ for $e$,
and Lemma~\ref{pd-kappa-defects-windows-lift-pairs-form} gives the $\kappa_b$ and, for $u$ carried by $b$, the tuples
$f(b,u)$ and $f'(b,u)$. Let $b$ be a tuple of blocks and $u$ carried by $b$. If some $b_j$ is not kept, put $F(b,u)=0$ and
$G(b,u)=0$; they lie in the cores, which are $\mathcal R$-modules. Otherwise $f(b,u)_j$ and $f'(b,u)_j$ lie in
$Z_j(b_j;\mu_0)$ for every $j$, by Lemma~\ref{pd-kappa-defects-windows-lift-pairs-form}(1).
Lemma~\ref{pd-kappa-defects-windows-lift-pairs-source}, applied to $h$, $\Theta$ and the edge $e=e_h(1)$, gives
$P_{b,u}\in\mathrm{Pol}_T$ with $\Theta P_{b,u}\in\mathcal C^\times_{t,c}(e)$ and
$\Sigma_1P_{b,u}-\Phi(b,f(b,u))\in\varpi_0^c\mathcal R[z]$. Applied to $h'$, $\Theta'$ and the edge $e^\vee=e_{h'}(1)$, it
gives $P'_{b,u}\in\mathrm{Pol}_T$ with $\Theta'P'_{b,u}\in\mathcal C^\times_{t,c}(e^\vee)$ and
$\Sigma_1P'_{b,u}-\Phi(b,f'(b,u))\in\varpi_0^c\mathcal R[z]$. Put $F(b,u)=\Theta P_{b,u}$ and
$G(b,u)=\Xi(\Theta'P'_{b,u})$. Then $F(b,u)\in\mathcal C^\times_{t,c}(e)$, and $G(b,u)\in\mathcal C'^\times_{t,c}(e)$ by
Lemma~\ref{pd-kappa-defects-windows-lift-pairs-flat}(2).

\emph{The pairings.} Let $u$ be carried by $b$ and $u'$ by $b'$, with $\beta_j(u_j,u'_j)=a_j\in\mathcal R$ for every $j$.
\begin{itemize}
\item If every $b_j$ and every $b'_j$ is kept, then $\Sigma_1P_{b,u}$ and $\Sigma_1P'_{b',u'}$ lie in $\mathrm{Pol}_T$ by
  Lemma~\ref{pd-kappa-defects-windows-lift-pairs-flat}(3), so Lemma~\ref{pd-kappa-defects-windows-lift-pairs-form}(2) gives
  $\nu\in\mathcal R$ with $z_T[\Sigma_1P_{b,u},\Sigma_1P'_{b',u'}]=\delta_{bb'}\kappa_b\prod_ja_j+\varpi_0^c\nu$. By
  Lemma~\ref{pd-kappa-defects-windows-lift-pairs-flat}(4), $\langle F(b,u),G(b',u')\rangle^\times$ is the same element.
\item If every $b_j$ is kept and some $b'_j$ is not, then $G(b',u')=0$ and $b\neq b'$, so the pairing and
  $\delta_{bb'}\kappa_b\prod_ja_j$ are both $0$.
\item If some $b_{j_0}$ is not kept, then $F(b,u)=0$. The entry of $u_{j_0}\in L_{j_0}$ at $b_{j_0}$ is $0$, as $b_{j_0}$ is
  not kept (Definition~\ref{pd-kappa-defects-segment-lattice}(7)), and its other entries are $0$, as $u$ is carried by $b$.
  So $u_{j_0}=0$, $a_{j_0}=\beta_{j_0}(0,u'_{j_0})=0$, and the pairing and $\prod_ja_j$ are both $0$.
\end{itemize}
\end{proof}
% Size: the open parts are Lemmas pd-kappa-defects-windows-lift-pairs-source-lift-tail-existence-columns-symbol,
%   -columns-counts and -columns-degrees (about 2,400 lines), under the proof of Lemma
%   pd-kappa-defects-windows-lift-pairs-source-lift-tail-existence. Lemma pd-kappa-defects-windows-lift-pairs-form-block and
%   parts (1) and (2) of Lemma pd-kappa-defects-windows-lift-pairs-source-order have no open part. The Witt
%   identification WittIdent (WittPI.wittIdentity_holds, about 400 lines) only moves the Lean's matrices from the Witt
%   vectors to W, which the book's notation does not need.
% Lean: paired_lifts_of_identity hW hRT hForm, with Ct0 = 30: paired_lift_at_of_identity hW testFlatAt1
%   (srcHalf1_of (le_refl 13) hCt hstd (lemma44_identity_of_raw_tail hR hK) monoAt planeAt symIntAt) hF.
%   In tensor_closed, hW is the field lWitt of Loc8All, hRT = restRawTail1_of (third_key_of_order
%   key3Ord) (local_fourth_of_order (loc4Ord_of (loc4Div_of loc4Res))) (rest_assumptions_of_three_interfaces
%   hA restCols restPre (restTCore_of_seg (restSegCore_of_num numCore_thr))), and hForm =
%   restForm_of_oplus (lemma_five_two_first_holds D) restOPlus. The Lean's Loc data have real ccl, DLP
%   and cs. The thresholds are TransportedAlt.and of RestRawTail1, restKeptBL, RestForm and thr_std. The
%   proof follows paired_lift_at_of_identity: hF and hG build the lifts (0 unless KeptAll and FamOK), and the
%   pairing splits into the three cases above (hform with hPol and hpair; if_neg; famOK_zero with
%   segForm_zero_left).

\begin{lemma}\label{pd-kappa-defects-windows-lift}
For every $C_c\geq0$ there is $T_0>0$ such that for every $C_0\geq T_0$ there is $T_1>0$ such that for
all $C_{E1},C_{E2},C_m\geq T_1$ there is $T_2>0$ such that for every $C_H\geq T_2$ there is $C^0_t>0$
with the following property for every $C_t\geq C^0_t$. Let $p\geq5$ be a prime and $n'\geq2$ an even
integer, and let $m$, $G'$, $N_1$, $c_{cl}$, $D_p$, $c_*$, $T$, $Y$, $Y_b$, $b'$, $b_l$, $b_u$ and an
integer $a$ satisfy hypotheses (1) to (3) of Theorem~\ref{pd-gain-theorem-one} with $C_0$ and $C_H$.
Let $c\geq0$ be an integer with $c\leq C_c(c_*+|T|k\log_pa+|T|c_{cl}+n'c_{cl})$, and let
$a_T,b_T\colon\mathbf Z/n'\to\mathbf Z_{\geq0}$ satisfy $a\leq a_T(s)\leq a+a/(kY)+2J_0$ and
$b_l\leq b_T(s)\leq b_u+2J_0$ for $s\in T$, where $J_0=\Lambda b'/k$. Form the numbers of
Definition~\ref{pd-kappa-defects-numbers} with these constants and any $C_A,C_\kappa>0$; none of the
numbers used here involves $C_A$ or $C_\kappa$. Let an unramified quadratic Hida setting of degree $n'$
at $p$ with rings $\mathcal R\subset\mathcal K$ and $\varpi_0$ be given, and reals $D_j\geq0$ and integers
$g_j\geq0$ ($1\leq j\leq q$) such that, for each $j$, the segment lattice of the $j$-th segment has a
block-diagonal family of size $g_j$ and depth $D_j$ (Definition~\ref{pd-kappa-defects-windows-blocks}).
Put $D=\sum_jD_j+4\kappa_0$ and $t=\lfloor t_0\rfloor$. If $4D<c$, then for every edge $e$ there are
$F_I\in\mathcal C^\times_{t,c}(e)$, $G_I\in\mathcal C'^\times_{t,c}(e)$ and nonzero $\delta_I\in\mathcal R$
with $v(\delta_I)\leq D$, for $I$ in the product of the sets $\{1,\dots,g_j\}$, such that
$\langle F_I,G_J\rangle^\times-\delta_I\delta_{IJ}\in\varpi_0^c\mathcal R$ for all $I$ and $J$
(Definition~\ref{pd-kappa-product-space}).
\end{lemma}
\leanref{Formalism/SharpRestriction/GrowthBounds/BinomialRatio/ExtendedStatementOfParts.lean}{Formalism.MainNotions.ColumnSystems.Pf.LiftVecAt}
\leanref{Formalism/SharpRestriction/GrowthBounds/BinomialRatio/ExtendedStatementOfParts.lean}{Formalism.MainNotions.ColumnSystems.Pf.RestLiftVec}
\leanref{Formalism/SharpRestriction/GrowthBounds/BinomialRatio/PairLift.lean}{Formalism.MainNotions.ColumnSystems.Pf.restLiftVec_of}

\begin{proof}
\emph{The constants.} Let $C_c\geq0$, and let $T^P_0$, $T^P_1$, $T^P_2$ be the thresholds of
Lemma~\ref{pd-kappa-defects-windows-lift-pairs}. Put $T_0=\max(T^P_0,1)$; for $C_0\geq T_0$ put
$T_1=\max(T^P_1,1)$; for $C_{E1},C_{E2},C_m\geq T_1$ put
$T_2=\max\big(T^P_2,1000(C_{E1}+C_{E2}+1)(C_c+1)\big)$; and for $C_H\geq T_2$ let $C^0_t$ be the constant
of Lemma~\ref{pd-kappa-defects-windows-lift-pairs}.

\emph{The lift.} Let $C_t\geq C^0_t$, and let the data, the setting, the $D_j$ and $g_j$ and an edge $e$
be given, with $4D<c$. The facts (a) to (g) in the proof of Lemma~\ref{pd-kappa-defects-margin} hold;
they give $k\geq1$ and $a\geq5$, so $\kappa_0=2|T|k\log_pa\geq0$. As the $D_j$ are $\geq0$, $c>4D\geq0$, so
$c\geq1$. Let $\kappa_b$, $F$ and $G$ be as in Lemma~\ref{pd-kappa-defects-windows-lift-pairs} at $e$. For
each $j$ let $x^j_i,y^j_i\in L_j$, the blocks $b^j_i$ and $z^j_i\in\mathcal R$ ($1\leq i\leq g_j$) form a
block-diagonal family of size $g_j$ and depth $D_j$ in $L_j$. For $I$ in the product of the sets
$\{1,\dots,g_j\}$ let $b(I)=(b^j_{I_j})_j$, $x(I)=(x^j_{I_j})_j$ and $y(I)=(y^j_{I_j})_j$; then $x(I)$ and
$y(I)$ are carried by $b(I)$. Put
\[
 F_I=F\big(b(I),x(I)\big),\qquad G_I=G\big(b(I),y(I)\big),\qquad
 \delta_I=\kappa_{b(I)}\prod_{j=1}^qz^j_{I_j}.
\]
Then $F_I\in\mathcal C^\times_{t,c}(e)$ and $G_I\in\mathcal C'^\times_{t,c}(e)$. As $\mathcal R$ is a domain,
$\delta_I\neq0$, and $v(\delta_I)=v(\kappa_{b(I)})+\sum_jv(z^j_{I_j})\leq4\kappa_0+\sum_jD_j=D$. For all
$I$ and $J$, $\beta_j(x^j_{I_j},y^j_{J_j})=\delta_{I_jJ_j}z^j_{I_j}\in\mathcal R$, so
Lemma~\ref{pd-kappa-defects-windows-lift-pairs} gives
\[
 \langle F_I,G_J\rangle^\times-\delta_{b(I)b(J)}\kappa_{b(I)}\prod_{j=1}^q\delta_{I_jJ_j}z^j_{I_j}
 \in\varpi_0^c\mathcal R .
\]
If $I=J$, the subtracted term is $\delta_I$. If $I\neq J$, then $I_j\neq J_j$ for some $j$, and the
subtracted term is $0=\delta_I\delta_{IJ}$.
\end{proof}
% Lean: restLiftVec_of hPL = TransportedAlt.mono of TransportedAlt.and thr_std hPL (thresholds 1, 1 and
% T2 CE1 CE2 Cc 1 for std_of), with the Ct0 of RestPairLift (Lemma pd-kappa-defects-windows-lift-pairs).
% The argument is liftVecAt_of: hkap (kap0 >= 0 from p5 and loga), hc1, FamOK for the families of
% BlockDiag (hfx, hfy), vLEO_prod and vLEO_mul for the valuation, and the product of the
% if-then-else pairings for I /= J.

\begin{lemma}\label{pd-kappa-defects-windows-perturb}
Let $\mathcal R$ be a discrete valuation ring with fraction field $\mathcal K$ and a uniformizer
$\varpi_0$ with $\varpi_0^2=p$ for a prime $p$, and let $v$ be the valuation of $\mathcal K$ with
$v(p)=1$. Let $\mathcal C$ and $\mathcal C'$ be sets, $\langle\ ,\ \rangle\colon\mathcal C\times\mathcal C'\to
\mathcal K$ a map, $n\geq1$ and $c\geq0$ integers and $D$ a real number with $4D<c$. Let
$x_1,\dots,x_n\in\mathcal C$, $y_1,\dots,y_n\in\mathcal C'$, nonzero $\delta_1,\dots,\delta_n\in\mathcal R$
with $v(\delta_i)\leq D$, and $\nu_{ij}\in\mathcal R$ satisfy
$\langle x_i,y_j\rangle=\delta_i\delta_{ij}+\varpi_0^c\nu_{ij}$ for all $i,j$. Then there are a nonzero
$z\in\mathcal R$ with $v(z)\leq D$ and an $n\times n$ matrix $M$ over $\mathcal R$ with
$\sum_l\langle x_i,y_l\rangle M_{lj}=z\delta_{ij}$ for all $i,j$.
\end{lemma}
\leanref{Formalism/SharpRestriction/GrowthBounds/BinomialRatio/FactorsOfPerturb.lean}{Formalism.MainNotions.ColumnSystems.Pf.factors_of_perturb}
\leanref{Formalism/SharpRestriction/GrowthBounds/BinomialRatio/FactorsOfPerturb.lean}{Formalism.MainNotions.ColumnSystems.Pf.isUnit_one_add}
\leanref{Formalism/SharpRestriction/GrowthBounds/BinomialRatio/FactorsOfPerturb.lean}{Formalism.MainNotions.ColumnSystems.Pf.liftAt_of}

\begin{proof}
Write $\delta_i=u_i\varpi_0^{k_i}$ with $u_i$ a unit and $k_i\geq0$. Then $k_i/2=v(\delta_i)\leq D$, so
$k_i\leq2D<c$. Let $i_0$ be an index with $k_{i_0}$ largest. Put
$N_{ij}=u_i^{-1}\varpi_0^{c-k_i-1}\nu_{ij}\in\mathcal R$; then
$\delta_i\delta_{ij}+\varpi_0^c\nu_{ij}=\delta_i(1+\varpi_0N)_{ij}$. The determinant of $1+\varpi_0N$ is
$1$ modulo $\varpi_0$, so it is a unit and $1+\varpi_0N$ has an inverse $Q$ over $\mathcal R$. Put
$z_l=u_{i_0}u_l^{-1}\varpi_0^{k_{i_0}-k_l}\in\mathcal R$, so that $\delta_lz_l=\delta_{i_0}$, and
$M_{lj}=Q_{lj}z_j$. Then
\[
 \sum_l\langle x_i,y_l\rangle M_{lj}=\delta_i\sum_l(1+\varpi_0N)_{il}Q_{lj}z_j=\delta_i\delta_{ij}z_j
 =\delta_{i_0}\delta_{ij},
\]
and $z=\delta_{i_0}$ is nonzero with $v(z)\leq D$.
\end{proof}
% Lean: factors_of_perturb (eq_unit_pow, vLEO_iff with ramE_two, isUnit_one_add, the matrix
% Q * diagonal zl); liftAt_of passes from LiftVecAt to LiftAt with it.

\begin{lemma}\label{pd-kappa-defects-windows}
For every $C_c\geq0$ there is $T_0>0$ such that for every $C_0\geq T_0$ there is $T_1>0$ such that for
all $C_{E1},C_{E2},C_m\geq T_1$ there is $T_2>0$ such that for every $C_H\geq T_2$ there are
$C_\kappa,C_t>0$ with the following property. Let $p\geq5$ be a prime and $n'\geq2$ an even integer,
and let $m$, $G'$, $N_1$, $c_{cl}$, $D_p$, $c_*$, $T$, $Y$, $Y_b$, $b'$, $b_l$, $b_u$ and an integer $a$ satisfy
hypotheses (1) to (3) of Theorem~\ref{pd-gain-theorem-one} with $C_0$ and $C_H$. Let $c\geq0$ be an
integer with $c\leq C_c(c_*+|T|k\log_pa+|T|c_{cl}+n'c_{cl})$, and let
$a_T,b_T\colon\mathbf Z/n'\to\mathbf Z_{\geq0}$ satisfy $a\leq a_T(s)\leq a+a/(kY)+2J_0$ and
$b_l\leq b_T(s)\leq b_u+2J_0$ for $s\in T$, where $J_0=\Lambda b'/k$. Form the numbers of
Definition~\ref{pd-kappa-defects-numbers} with these constants and any $C_A>0$; none of the numbers
used here involves $C_A$. Let an unramified quadratic Hida setting of
degree $n'$ at $p$ be given, and reals $D_j,\kappa_j\geq0$ ($1\leq j\leq q$) such that, for each $j$,
the segment lattice of the $j$-th segment has a good share of depth $D_j$ and defect $\kappa_j$ for its
pairing. If $4(\sum_jD_j+4\kappa_0)<c$, then for every edge $e$ there is a split family of depth
$\sum_jD_j+4\kappa_0$ and size at least
$(1-\kappa_P-\sum_j\kappa_j)\dim V^\times(a_T|_T,b_T|_T)/[k]$ on
$(\mathcal C^\times_{t,c}(e),\mathcal C'^\times_{t,c}(e))$, where $t=\lfloor t_0\rfloor$
(Definition~\ref{pd-kappa-product-space}).
\end{lemma}
\leanref{Formalism/SharpRestriction/GrowthBounds/ColumnSystems.lean}{Formalism.MainNotions.ColumnSystems.HerrExtensionClass}
\leanref{Formalism/SharpRestriction/GrowthBounds/ColumnSystems.lean}{Formalism.MainNotions.ColumnSystems.DjHyp}
\leanref{Formalism/SharpRestriction/GrowthBounds/ColumnSystems.lean}{Formalism.MainNotions.ColumnSystems.HIext}
\leanref{Formalism/SharpRestriction/GrowthBounds/ColumnSystems.lean}{Formalism.MainNotions.ColumnSystems.ExtendedAllConclusions}
\leanref{Formalism/SharpRestriction/GrowthBounds/PlaneCountTheoremWithConstants.lean}{Formalism.MainNotions.G.extended_windows_theorem}

\begin{proof}
\emph{The constants.} Let $C_c\geq0$. Let $T^M_0$ and $T^L_0$ be the numbers $T_0$ of
Lemmas~\ref{pd-kappa-defects-windows-model} and~\ref{pd-kappa-defects-windows-lift} for $C_c$, and put
$T_0=\max(T^M_0,T^L_0,1)$. For $C_0\geq T_0$ let $T^M_1$ and $T^L_1$ be their numbers $T_1$, and put
$T_1=\max(T^M_1,T^L_1,1)$. For $C_{E1},C_{E2},C_m\geq T_1$ let $T^M_2$ and $T^L_2$ be their numbers
$T_2$, and put $T_2=\max\big(T^M_2,T^L_2,1000(C_{E1}+C_{E2}+1)(C_c+1)\big)$. For $C_H\geq T_2$ let
$C_\kappa=C^0_\kappa$ and $C_t=C^0_t$ be the constants of the two lemmas.

\emph{Signs.} Let the data, the setting, with rings $\mathcal R\subset\mathcal K$ and $\varpi_0$, the
$D_j$ and $\kappa_j$, and an edge $e$ be given, and put $D=\sum_jD_j+4\kappa_0$, so $4D<c$. As
$C_0\geq1$, $C_{E1},C_{E2},C_m\geq1$ and $C_H\geq1000(C_{E1}+C_{E2}+1)(C_c+1)$, the facts (a) to (g) in
the proof of Lemma~\ref{pd-kappa-defects-margin} hold. With hypothesis (1) they give $Y\geq1$, $b'>0$,
$k,Q\geq1$, $\Lambda\geq1$ and $a\geq5$; so $E_0\geq0$ and $\kappa_P\geq0$, and $m_0\geq0$ as $c\geq0$.
Also $[k]>0$.

\emph{The split family.} For each $j$ let $L_j$ be the segment lattice of the $j$-th segment, $L^j_1$ a
good sublattice of $L_j$ of depth $D_j$ with $\operatorname{rk}L^j_1\geq(1-\kappa_j)\operatorname{rk}L_j$,
and $g_j=\operatorname{rk}L^j_1$. By Lemma~\ref{pd-kappa-defects-windows-partners}, $L_j$ has a
block-diagonal family of size $g_j$ and depth $D_j$. Lemma~\ref{pd-kappa-defects-windows-lift} gives
$F_I$, $G_I$ and $\delta_I$ as there, for the $g=\prod_jg_j$ elements $I$ of the product of the sets
$\{1,\dots,g_j\}$. If $g=0$, the empty family is a split family of depth $D$ and size $0$. Otherwise
Lemma~\ref{pd-kappa-defects-windows-perturb}, applied to the $F_I$ and $G_I$ in some order, gives a
nonzero $z\in\mathcal R$ with $v(z)\leq D$ and a matrix $M$ over $\mathcal R$ that make them a split
family of depth $D$ and size $g$ on $(\mathcal C^\times_{t,c}(e),\mathcal C'^\times_{t,c}(e))$.

\emph{The size.} It remains to show $(1-\kappa_P-\sum_j\kappa_j)\dim V^\times(a_T|_T,b_T|_T)\leq[k]g$.
If $\sum_j\kappa_j>1$, the left side is at most $0$, as $\kappa_P\geq0$. Otherwise
$0\leq\kappa_j\leq1$ for every $j$. Let $r_j=r_j(m_0)$ be as in
Lemma~\ref{pd-kappa-defects-windows-model}; by Lemma~\ref{pd-kappa-defects-windows-ranks},
$r_j=\operatorname{rk}L_j$, so $g_j\geq(1-\kappa_j)r_j\geq0$. Since
$(1-s)(1-\kappa)=1-s-\kappa+s\kappa\geq1-s-\kappa$ for $s,\kappa\geq0$, induction on the number of
factors gives $\prod_j(1-\kappa_j)\geq1-\sum_j\kappa_j$, and so
\[
 g=\prod_jg_j\geq\prod_j(1-\kappa_j)r_j\geq\Big(1-\sum_j\kappa_j\Big)\prod_jr_j .
\]
Lemma~\ref{pd-kappa-defects-windows-model} with $\mu=m_0$ gives
$(1-\kappa_P)\dim V^\times(a_T|_T,b_T|_T)\leq[k]\prod_jr_j$. As $\kappa_P$, $\sum_j\kappa_j$, the
dimension and $1-\sum_j\kappa_j$ are nonnegative,
\[
 \Big(1-\kappa_P-\sum_j\kappa_j\Big)\dim V^\times\leq\Big(1-\sum_j\kappa_j\Big)(1-\kappa_P)\dim V^\times
 \leq\Big(1-\sum_j\kappa_j\Big)[k]\prod_jr_j\leq[k]g .
\]
\end{proof}
% Lean: the lemma is the clause HerrExtensionClass of Closed (ExtendedAllConclusions D), with CYb = Cde =
%   Ck = CD = Cb = CH, as first_theorem_admitted_closed takes it from extended_windows_theorem (hExt);
%   in closed_of_rests it is RestExt = restExt_of restRk hL1 hLift hA2 (ExtendedStatementOfParts.lean).
%   The constants follow restExt_of: TransportedAlt.and of thr_std (thresholds 1, 1 and
%   T2 CE1 CE2 Cc 1), RestRk, PartnersToFactors, RestLift and RemainderSegmentBoundSecond, with Ckap >= Ck0
%   and Ct >= Ct0. The proof follows extended_statement_of_parts: the signs are kapP_nonneg (from std_of)
%   and hm0; the split family is PartnersToFactorsAt (Lemma pd-kappa-defects-windows-partners) and
%   LiftAt = liftAt_of (LiftVecAt, Lemma pd-kappa-defects-windows-lift, and factors_of_perturb, Lemma
%   pd-kappa-defects-windows-perturb), with RhoGE as FactorsLE of size prod g_j; the size is
%   Top.shareProd (share_product_bound), RkAt (Lemma pd-kappa-defects-windows-ranks) and
%   ExtendedModelRank at m0 (Lemma pd-kappa-defects-windows-model). The Lean states the lemma for Loc data
%   with real ccl, DLP and cs; the book needs integer data.

\subsection{Instances and the block bound}\label{subsection-pd-instances-block-bound}

\begin{definition}\label{pd-kappa-defects-segment-mean-operators}
Let $K$ be a field, $V$ a $K$-vector space, $n\geq1$ an integer and $E=(E_{ab})_{1\leq a,b\leq n}$ a
family of endomorphisms of $V$. Let $\mathcal U$ be the monoid of endomorphisms generated by the
$E_{ab}$.
\begin{enumerate}
\item For $\mu\in\mathbf Z^n$ the \emph{weight space} $V_\mu$ is the set of $x\in V$ with $E_{aa}x=\mu_ax$
  for every $a$. A \emph{highest weight vector} of weight $\lambda$ is a nonzero $v\in V_\lambda$ with
  $E_{ab}v=0$ for all $a<b$. The \emph{isotypic subspace} $V^{(\lambda)}$ is the span of the $u(v)$ with
  $u\in\mathcal U$ and $v$ a highest weight vector of weight $\lambda$. For a set $S\subset\mathbf Z^n$,
  an \emph{isotypic projector} for $S$ is an endomorphism $e$ of $V$ that is the identity on
  $V^{(\lambda)}$ for $\lambda\in S$ and $0$ on $V^{(\mu)}$ for $\mu\notin S$.
\item $E$ is an \emph{irreducible highest weight module} of weight $\lambda$ generated by $v\in V$ if
  $E_{ab}E_{cd}-E_{cd}E_{ab}=\delta_{bc}E_{ad}-\delta_{da}E_{cb}$ for all $a,b,c,d$, every $x\in V$ lies
  in a finite-dimensional subspace stable under the $E_{ab}$, $V$ is the sum of its weight spaces,
  $\lambda_1\geq\cdots\geq\lambda_n$, $v$ is a highest weight vector of weight $\lambda$, and the $u(v)$
  with $u\in\mathcal U$ span $V$.
\item Let $K$ be the fraction field of a discrete valuation ring $\mathcal R$. Let $\mathcal U_{\mathbf Z}$
  be the monoid generated by the $E_{ab}^k/k!$ with $a\neq b$ and $k\geq0$ and the
  $E_{aa}(E_{aa}-1)\cdots(E_{aa}-k+1)/k!$ with $k\geq0$. For $v\in V$ let $\Delta(v)$ be the
  $\mathcal R$-span of the $u(v)$ with $u\in\mathcal U_{\mathbf Z}$, and for $\lambda\in\mathbf Z^n$ let
  $\nabla(\lambda,v)$ be the set of $x\in V$ such that for every $u\in\mathcal U_{\mathbf Z}$ there is
  $c\in\mathcal R$ with $u(x)-cv\in\sum_{\mu\neq\lambda}V_\mu$. A \emph{lattice} for $E$ is a finitely
  generated $\mathcal R$-submodule of $V$ stable under the generators of $\mathcal U_{\mathbf Z}$.
\item Let $\mathcal R$ be a commutative ring, $r\geq0$, and $A_{r+1}=\mathcal R[x_{ab}:1\leq a,b\leq
  r+1][\det(x)^{-1}]$, where $x=(x_{ab})$. For an $\mathcal R$-algebra $A$ and $g\in GL_{r+1}(A)$ let
  $f(g)$ be the image of $f\in A_{r+1}$ under $x\mapsto g$. For $c\in\mathbf Z$ and $\kappa\in\mathbf Z^r$,
  $H_{r+1}(c,\kappa)\subset A_{r+1}$ is the set of the $f$ with
  $f(bg)=u_1^cu_2^{\kappa_1}\cdots u_{r+1}^{\kappa_r}f(g)$ for every $\mathcal R$-algebra $A$, every
  $g\in GL_{r+1}(A)$ and every lower triangular $b\in GL_{r+1}(A)$ with diagonal entries
  $u_1,\dots,u_{r+1}$. Put $g_0=x_{1,r+1}\in A_{r+1}$.
\end{enumerate}
\end{definition}
\leanref{Formalism/SharpRestriction/GrowthBounds/ObjPol.lean}{Formalism.MainNotions.G.GLOps}
\leanref{Formalism/SharpRestriction/GrowthBounds/ObjPol.lean}{Formalism.MainNotions.G.weightSp}
\leanref{Formalism/SharpRestriction/GrowthBounds/ObjPol.lean}{Formalism.MainNotions.G.IsHWV}
\leanref{Formalism/SharpRestriction/GrowthBounds/ObjPol.lean}{Formalism.MainNotions.G.Umon}
\leanref{Formalism/SharpRestriction/GrowthBounds/ObjPol.lean}{Formalism.MainNotions.G.isoSub}
\leanref{Formalism/SharpRestriction/GrowthBounds/ObjPol.lean}{Formalism.MainNotions.G.IsIsoProjSet}
\leanref{Formalism/SharpRestriction/GrowthBounds/ObjPol.lean}{Formalism.MainNotions.G.IsGLRep}
\leanref{Formalism/SharpRestriction/GrowthBounds/ObjPol.lean}{Formalism.MainNotions.G.Dominant}
\leanref{Formalism/SharpRestriction/GrowthBounds/ObjPol.lean}{Formalism.MainNotions.G.IsIrrHW}
\leanref{Formalism/SharpRestriction/GrowthBounds/ObjPol.lean}{Formalism.MainNotions.G.divPowE}
\leanref{Formalism/SharpRestriction/GrowthBounds/ObjPol.lean}{Formalism.MainNotions.G.binomE}
\leanref{Formalism/SharpRestriction/GrowthBounds/ObjPol.lean}{Formalism.MainNotions.G.UZgens}
\leanref{Formalism/SharpRestriction/GrowthBounds/ObjPol.lean}{Formalism.MainNotions.G.UZmon}
\leanref{Formalism/SharpRestriction/GrowthBounds/ObjPol.lean}{Formalism.MainNotions.G.Delta}
\leanref{Formalism/SharpRestriction/GrowthBounds/ObjPol.lean}{Formalism.MainNotions.G.otherWts}
\leanref{Formalism/SharpRestriction/GrowthBounds/ObjPol.lean}{Formalism.MainNotions.G.Nabla}
\leanref{Formalism/SharpRestriction/GrowthBounds/ObjPol.lean}{Formalism.MainNotions.G.UZStable}
\leanref{Formalism/SharpRestriction/GrowthBounds/ObjPol.lean}{Formalism.MainNotions.G.UZLattice}
\leanref{Formalism/SharpRestriction/GrowthBounds/ObjPol.lean}{Formalism.MainNotions.G.GLCoord}
\leanref{Formalism/SharpRestriction/GrowthBounds/FlipOneThree.lean}{Formalism.MainNotions.flipOneThree.evalG}
\leanref{Formalism/SharpRestriction/GrowthBounds/FlipOneThree.lean}{Formalism.MainNotions.flipOneThree.chiB}
\leanref{Formalism/SharpRestriction/GrowthBounds/FlipOneThree.lean}{Formalism.MainNotions.flipOneThree.HB}
\leanref{Formalism/SharpRestriction/GrowthBounds/FlipOneThree.lean}{Formalism.MainNotions.flipOneThree.g0r}

\begin{definition}\label{pd-kappa-defects-segment-mean-conditions}
Let $C_0,C_H,C_{E1},C_{E2},C_c,C_\kappa,C_s$ be reals. An \emph{instance} consists of a prime $p\geq5$,
an even integer $n'\geq2$, data $m\geq1$, $G'$, $N_1$, $c_{cl}$, $D_p$, $c_*$, $T$, $Y$, $Y_b$, $b'$, $b_l$, $b_u$
and an integer $a$ as in Theorem~\ref{pd-gain-theorem-one}, an integer $c\geq0$, reals $x_a,x_b$ and
$\mu$, a function $b_T\colon\mathbf Z/n'\to\mathbf Z_{\geq0}$, an unramified quadratic Hida setting of
degree $n'$ at $p$ with ring $\mathcal R$ and fraction field $\mathcal K$, and an index $1\leq j\leq q$. Form
the numbers of Definition~\ref{pd-kappa-defects-numbers} with these constants; those used here do not
involve $C_A$, $C_t$ or $C_m$. Put
$W=[a,a+a/(kY)+x_a]\cap\mathbf Z$, $s=S_j$ and $\kappa_1=\lceil m^2\ell\rceil$. The instance
\emph{satisfies its hypotheses} if hypotheses (1) to (3) of Theorem~\ref{pd-gain-theorem-one} hold
with $C_0$ and $C_H$, $kY^2\log a\leq a/b'$, $c\leq C_c(c_*+|T|k\log_pa+|T|c_{cl}+n'c_{cl})$,
$0\leq x_a,x_b\leq2J_0$, $b_l\leq b_T(s')\leq b_u+x_b$ for $s'\in T$, and $0\leq\mu\leq b'/(kQY)$. A
\emph{window function} is an $a_T\colon\mathbf Z/n'\to\mathbf Z_{\geq0}$ with
$a\leq a_T(s')\leq a+a/(kY)+x_a$ for $s'\in T$; for it, $r_{j'}(\mu')$ is as in
Lemma~\ref{pd-kappa-defects-windows-model}, and for $1\leq j'\leq q$, $\rho,\rho'$ and $\mu'\geq0$,
$r_{j'}(\mu';\rho,\rho')$ is the sum of $\dim V_b$ over the kept blocks $b$ of $S_{j'}(b_T)$, formed with
$\alpha_i=a_T(y_{j',i})$, $\alpha_+=a_T(s_{j'+1})$, $\mu'$ and $r=m$, whose vectors $\rho_i$ and
$\rho'_i$ are the $\rho_i$ and $\rho'_i$. \emph{Digits} for $j'$ are tuples
$\rho_i,\rho'_i\in\mathbf Z_{\geq0}^{m+1}$ ($1\leq i\leq S_{j'}$) that satisfy the first condition for a
kept block in Definition~\ref{pd-kappa-defects-segment-lattice}(2) for these $\alpha_i$; for them put
\[
 d_{j'}(\rho,\rho')=P_m(\hat a)\prod_{i=1}^{S_{j'}}P_m\big(\lfloor(b_i-|\rho_i|)/P_i\rfloor\big)
 P_m\big(\lfloor(\alpha_i-|\rho'_i|)/P_i\rfloor\big)\cdot P_m(b_+),
\]
with $\hat a$, $b_i$, $P_i$, $b_+$ those of $S_{j'}(b_T)$ and $P_m$ as in
Lemma~\ref{pd-kappa-defects-windows-model-segment}. A \emph{block lettering} is the lettering of the
letters $x_1,y_1,\dots,x_s,y_s,x_+$ of Definition~\ref{pd-kappa-defects-segment-lattice}(3), the $x_i$
and $x_+$ of vector type and the $y_i$ of covector type, with the degrees $\delta'$ of a level of a
block kept for $S_j(b_T)$ and some $\alpha\in W^s$ and $\alpha_+\in W$, formed with $\mu$ and $r=m$
(Definition~\ref{pd-kappa-defects-segment-lattice}(5)). For a lettering of $N$ letters with degrees
$B$, let $\mathrm{Pol}(B)\subset\mathcal K[z]$ be the polynomials in the $z_{i,c}$ ($0\leq i<N$,
$1\leq c\leq m$) of row degrees $B$ (Definition~\ref{pd-symbols-path-generators}); for $0\leq l\leq N$
let $E^{(l)}=(E^{(l)}_{ab})_{a,b\leq m}$ be the operators $E^{\mathcal L}_{ab}$ of
Definition~\ref{pd-kappa-defects-segment-lattice}(4) for the first $l$ letters, acting on
$\mathcal K[z]$, and $E'^{(l)}=(E^{(l)}_{ab})_{a,b\leq m-1}$; let $\mathcal B$ be the $\mathcal R$-span
of the monomials $\gamma z^\nu$ with $v(\gamma)+\frac12v_p(\nu!)\geq0$. The notions of
Definition~\ref{pd-kappa-defects-segment-mean-operators} are used with these operators, and $[\ ,\ ]$
is the Fock pairing. The \emph{conditions} on the instance are:
\begin{enumerate}
\item If it satisfies its hypotheses, then for every block lettering, every $0\leq l\leq2s+1$ and each
  of $E=E^{(l)}$ and $E=E'^{(l)}$, every set of weights has an isotypic projector, and every
  isotypic projector $e$ satisfies $[eF,G]=[F,eG]$ for $F,G\in\mathrm{Pol}(\delta')$,
  $p^{\kappa_1}e(F)\in\mathcal R[z]$ for $F\in\mathrm{Pol}(\delta')\cap\mathcal R[z]$, and
  $p^{\kappa_1}e(F)\in\mathcal B$ for $F\in\mathrm{Pol}(\delta')\cap\mathcal B$. For $l\leq l'\leq2s+1$,
  each isotypic projector of $E^{(l)}$ commutes with each of $E^{(l')}$ and of $E'^{(l')}$, and each of
  $E'^{(l)}$ with each of $E'^{(l')}$.
\item For every $1\leq h\leq m$, every finite-dimensional $\mathcal K$-vector space $V$, viewed as an
  $\mathcal R$-module, and every family $(E_{ab})_{a,b\leq h}$ on $V$ that is an irreducible highest
  weight module of weight $\lambda$ generated by $v$ with $\lambda_1-\lambda_h\leq8a$:
  $p^{\kappa_1}\nabla(\lambda,v)\subset\Delta(v)$, and every lattice $M$ with $v\in M$ and
  $M\cap V_\lambda=\mathcal Rv$ satisfies $\Delta(v)\subset M\subset\nabla(\lambda,v)$.
\item If it satisfies its hypotheses, then for every block lettering, every $0\leq l\leq2s+1$ and every
  $\lambda\in\mathbf Z^m$ with $\lambda_1\geq\cdots\geq\lambda_{m-1}$, there is a $\mathcal K$-linear
  endomorphism $P$ of $\mathcal K[z]$ with $P^2=P$ such that, with $V_\lambda$ the weight space of $E^{(l)}$: the image of $P$
  is the set of $x\in V_\lambda$ with $E^{(l)}_{ab}x=0$ for $a<b\leq m-1$; the intersection of the kernel
  of $P$ with $V_\lambda$ is $V_\lambda\cap\sum_{b<a\leq m-1}E^{(l)}_{ab}(\mathcal K[z])$; $[PF,G]=[F,PG]$
  for $F,G\in V_\lambda$; and $p^{\kappa_1}P(F)\in\mathcal R[z]$ for
  $F\in\mathrm{Pol}(\delta')\cap\mathcal R[z]\cap V_\lambda$.
\item For every discrete valuation ring $\mathcal R'$, $r\geq1$, $c'\in\mathbf Z$ and
  $\kappa\in\mathbf Z^r$, formed over $\mathcal R'$: if $f\in H_{r+1}(c',\kappa)$ and $g_0$ divides $f$,
  then $f=g_0f_1$ with $f_1\in H_{r+1}(c'-1,\kappa)$.
\item For $\mathcal R'$, $r$, $c'$ and $\kappa$ as in (4): if $f\in H_{r+1}(c',\kappa)$, $h\in A_{r+1}$,
  $0\neq\gamma\in\mathcal R'$ and $g_0$ divides $f-\gamma h$, then $g_0$ divides $h-f'$ for some
  $f'\in H_{r+1}(c',\kappa)$.
\item For every field $K'$ of characteristic $0$, $r\geq1$, $c'\in\mathbf Z$ and $\kappa\in\mathbf Z^r$
  with $\kappa_1\leq c'$, formed over $K'$, with $\langle S\rangle$ the $K'$-span of a set $S$:
  $\dim\langle\overline{H_{r+1}(c',\kappa)}\rangle+\dim\langle H_{r+1}(c'-1,\kappa)\rangle
  =\dim\langle H_{r+1}(c',\kappa)\rangle$, where the bar is the image in $A_{r+1}/(g_0)$, and
  $\dim\langle\overline{H_{r+1}(c',\kappa)}\rangle$ is the sum of $\dim\langle H_r(c',\kappa')\rangle$
  over the $\kappa'\in\mathbf Z^{r-1}$ with $\kappa_{i+1}\leq\kappa'_i\leq\kappa_i$ for every $i$.
\item For every discrete valuation ring $\mathcal R'$, $r\geq0$, $c'\in\mathbf Z$ and
  $\kappa\in\mathbf Z^{r+1}$, formed over $\mathcal R'$: let $\varepsilon\colon A_{r+2}\to
  A_{r+1}[X_1,\dots,X_{r+1}]$ be the evaluation at the matrix with first row $(1,X_1,\dots,X_{r+1})$ and,
  below it, the column $0$ followed by the matrix $(x_{ab})$ of $A_{r+1}$. Its determinant is a unit;
  $\varepsilon$ is injective on $H_{r+2}(c',\kappa)$; for $f\in H_{r+2}(c',\kappa)$ the coefficients of
  $\varepsilon(f)$ lie in $H_{r+1}(\kappa_1,(\kappa_2,\dots,\kappa_{r+1}))$; if moreover
  $\varepsilon(f)=\gamma F$ with $0\neq\gamma\in\mathcal R'$ and the coefficients of $F$ in
  $H_{r+1}(\kappa_1,(\kappa_2,\dots,\kappa_{r+1}))$, then $F=\varepsilon(f')$ for some
  $f'\in H_{r+2}(c',\kappa)$; and for $e\geq0$, $g_0^e$ divides $f$ if and only if $X_{r+1}^e$ divides
  $\varepsilon(f)$.
\item For every lettering of $N$ letters with degrees $B$, every $q'\geq0$ and all $J_1,\dots,J_{q'}\in
  \{0,\dots,N\}$ and $M_1,\dots,M_{q'}\in\mathbf Z\cup\{+\infty\}$: the set of the
  $F\in\mathrm{Pol}(B)\cap\mathcal R[z]$ that lie in $Y(\mathcal L_k,M_k)$ for every $k$, formed for
  this lettering as in Definition~\ref{pd-kappa-defects-segment-lattice}(4) with $\mathcal L_k$ the set
  of the first $J_k$ letters (for $M_k=+\infty$ there is no bound on the weight), is the $\mathcal R$-span of the path generators $\Psi(\mu',I)$ over $\mathcal K$ of rank $m$
  (Definition~\ref{pd-symbols-path-generators}) with $\mu'$ a path for $B$, $I$ a standard filling of
  $\mu'$, and $\mu'_{J_k}(1)\leq M_k$ for every $k$.
\item If it satisfies its hypotheses, then for every window function $a_T$ and every real
  $\mu'\geq0$, $(1-\kappa_P(\mu'))\dim V^\times(a_T|_T,b_T|_T)\leq[k]\prod_{j'=1}^qr_{j'}(\mu')$, and
  $\kappa_P(\mu')\leq C_s|T|N_1/Y$ if moreover $\mu'\leq b'/(kQY)$.
\item If it satisfies its hypotheses, then for every window function, every $1\leq j'\leq q$ and all
  digits $\rho,\rho'$ for $j'$, $r_{j'}(0;\rho,\rho')\leq d_{j'}(\rho,\rho')$.
\item If it satisfies its hypotheses, then for every window function, every real
  $0\leq\mu'\leq b'/(kQY)$ and all digits $\rho,\rho'$ for $j$,
  $d_j(\rho,\rho')/2\leq r_j(\mu';\rho,\rho')\leq d_j(\rho,\rho')$.
\end{enumerate}
\end{definition}
\leanref{Formalism/SharpRestriction/GrowthBounds/FlipOneThree.lean}{Formalism.MainNotions.flipOneThree.Inst}
\leanref{Formalism/SharpRestriction/GrowthBounds/FlipOneThree.lean}{Formalism.MainNotions.flipOneThree.HI}
\leanref{Formalism/SharpRestriction/GrowthBounds/FlipOneThree.lean}{Formalism.MainNotions.flipOneThree.Inst.winA}
\leanref{Formalism/SharpRestriction/GrowthBounds/FlipOneThree.lean}{Formalism.MainNotions.flipOneThree.Inst.InWinA}
\leanref{Formalism/SharpRestriction/GrowthBounds/FlipOneThree.lean}{Formalism.MainNotions.flipOneThree.kappaOneExponent}
\leanref{Formalism/SharpRestriction/GrowthBounds/FlipOneThree.lean}{Formalism.MainNotions.flipOneThree.Inst.BlockLetters}
\leanref{Formalism/SharpRestriction/GrowthBounds/FlipOneThree.lean}{Formalism.MainNotions.flipOneThree.chOps}
\leanref{Formalism/SharpRestriction/GrowthBounds/FlipOneThree.lean}{Formalism.MainNotions.flipOneThree.chOps'}
\leanref{Formalism/SharpRestriction/GrowthBounds/FlipOneThree.lean}{Formalism.MainNotions.flipOneThree.PTool}
\leanref{Formalism/SharpRestriction/GrowthBounds/FlipOneThree.lean}{Formalism.MainNotions.flipOneThree.resMat}
\leanref{Formalism/SharpRestriction/GrowthBounds/FlipOneThree.lean}{Formalism.MainNotions.flipOneThree.admitted}
\leanref{Formalism/SharpRestriction/GrowthBounds/FlipOneThree.lean}{Formalism.MainNotions.flipOneThree.adm1}
\leanref{Formalism/SharpRestriction/GrowthBounds/FlipOneThree.lean}{Formalism.MainNotions.flipOneThree.adm2}
\leanref{Formalism/SharpRestriction/GrowthBounds/FlipOneThree.lean}{Formalism.MainNotions.flipOneThree.adm3}
\leanref{Formalism/SharpRestriction/GrowthBounds/FlipOneThree.lean}{Formalism.MainNotions.flipOneThree.ArchimedeanSignatureMatching}
\leanref{Formalism/SharpRestriction/GrowthBounds/FlipOneThree.lean}{Formalism.MainNotions.flipOneThree.adm5}
\leanref{Formalism/SharpRestriction/GrowthBounds/FlipOneThree.lean}{Formalism.MainNotions.flipOneThree.adm6}
\leanref{Formalism/SharpRestriction/GrowthBounds/FlipOneThree.lean}{Formalism.MainNotions.flipOneThree.adm7}
\leanref{Formalism/SharpRestriction/GrowthBounds/FlipOneThree.lean}{Formalism.MainNotions.flipOneThree.adm8}
\leanref{Formalism/SharpRestriction/GrowthBounds/FlipOneThree.lean}{Formalism.MainNotions.flipOneThree.adm9}
\leanref{Formalism/SharpRestriction/GrowthBounds/FlipOneThree.lean}{Formalism.MainNotions.flipOneThree.adm10}
\leanref{Formalism/SharpRestriction/GrowthBounds/FlipOneThree.lean}{Formalism.MainNotions.flipOneThree.adm11}
\leanref{Formalism/SharpRestriction/GrowthBounds/ObjSeg.lean}{Formalism.MainNotions.G.chPol}
\leanref{Formalism/SharpRestriction/GrowthBounds/ObjPol.lean}{Formalism.MainNotions.G.IntPol}
\leanref{Formalism/SharpRestriction/GrowthBounds/ObjPol.lean}{Formalism.MainNotions.G.fock}
\leanref{Formalism/SharpRestriction/GrowthBounds/ObjSeg.lean}{Formalism.MainNotions.G.Yjt}
\leanref{Formalism/SharpRestriction/GrowthBounds/ColumnSystems.lean}{Formalism.MainNotions.ColumnSystems.rankOnDigits}
\leanref{Formalism/SharpRestriction/GrowthBounds/ColumnSystems.lean}{Formalism.MainNotions.ColumnSystems.dimMM}
\leanref{Formalism/SharpRestriction/GrowthBounds/ColumnSystems.lean}{Formalism.MainNotions.ColumnSystems.digitBlock}
\leanref{Formalism/SharpRestriction/GrowthBounds/ObjSeg.lean}{Formalism.MainNotions.G.coarseA}
\leanref{Formalism/SharpRestriction/GrowthBounds/ObjSeg.lean}{Formalism.MainNotions.G.coarseC}

\begin{definition}\label{pd-kappa-defects-segment-mean-theorem-a-classes}
Let an instance be given (Definition~\ref{pd-kappa-defects-segment-mean-conditions}), with $W$ as there,
$s=S_j$, $t=S_j-\eta_j$, $P_i=p^{g_j-d_{j,i}}$ and $b_i=b_T(y_{j,i})$ for $1\leq i\leq s$, and $P_m$ as in
Lemma~\ref{pd-kappa-defects-windows-model-segment}; write $S=S_j(b_T)$.
\begin{enumerate}
\item For $\alpha\in W^s$ and $\alpha_+\in W$, let $a_T^{\alpha,\alpha_+}\colon\mathbf Z/n'\to\mathbf Z_{\geq0}$ take
  the value $\alpha_i$ at $y_{j,i}$ ($1\leq i\leq s$), $\alpha_+$ at $s_{j+1}$, and $a$ elsewhere; these points
  are distinct.
\item A \emph{digit class} is a pair $\delta=(\rho,\rho')$ of tuples $\rho_i,\rho'_i\in\mathbf Z^{m+1}_{\geq0}$
  ($1\leq i\leq s$) whose entries are less than $P_i$; there are finitely many. It is \emph{compatible}
  with $\alpha\in W^s$ if $|\rho_i|\leq b_i$, $|\rho'_i|\leq\alpha_i$, and $P_i$ divides $b_i-|\rho_i|$ and
  $\alpha_i-|\rho'_i|$, for every $i$. Let $A(\delta)$ be the set of these $\alpha$, and for
  $\alpha\in A(\delta)$ put
  \[
   d(\delta;\alpha)=P_m(\hat a)\prod_{i=1}^sP_m\Big(\frac{b_i-|\rho_i|}{P_i}\Big)
   P_m\Big(\frac{\alpha_i-|\rho'_i|}{P_i}\Big)\cdot P_m(b_+),
  \]
  with $\hat a$ and $b_+$ those of $S$. Let $M(\alpha)$ be the sum of the $d(\delta;\alpha)$ over the digit
  classes $\delta$ compatible with $\alpha$.
\item For $\alpha\in W^s$, $\alpha_+\in W$ and a block $b$, the block lattice $L_b$ and the space $V_b$ are
  those of $L_S(\alpha,\alpha_+;\mu)$, formed over $\mathcal R$ with $r=m$
  (Definitions~\ref{pd-kappa-defects-segment-lattice} and~\ref{pd-kappa-defects-windows-blocks}), and
  $\beta_b(f,f')=\operatorname{sp}(b)[f,f']$ is the \emph{block pairing} on $\mathcal P$. Whether $b$ is kept
  depends only on $\alpha$; the vectors $\rho_i,\rho'_i$ of a block kept for $\alpha$ form a digit class
  compatible with $\alpha$, the \emph{class} of the block.
\end{enumerate}
\end{definition}
\leanref{Formalism/SharpRestriction/GrowthBounds/CountingStatementReduction.lean}{Formalism.MainNotions.G.CountingStatementReduction.realW}
\leanref{Formalism/SharpRestriction/GrowthBounds/CountingStatementReduction.lean}{Formalism.MainNotions.G.CountingStatementReduction.DCl}
\leanref{Formalism/SharpRestriction/GrowthBounds/CountingStatementReduction.lean}{Formalism.MainNotions.G.CountingStatementReduction.digSet}
\leanref{Formalism/SharpRestriction/GrowthBounds/CountingStatementReduction.lean}{Formalism.MainNotions.G.CountingStatementReduction.clSet}
\leanref{Formalism/SharpRestriction/GrowthBounds/CountingStatementReduction.lean}{Formalism.MainNotions.G.CountingStatementReduction.Compat}
\leanref{Formalism/SharpRestriction/GrowthBounds/CountingStatementReduction.lean}{Formalism.MainNotions.G.CountingStatementReduction.mCl}
\leanref{Formalism/SharpRestriction/GrowthBounds/CountingStatementReduction.lean}{Formalism.MainNotions.G.CountingStatementReduction.MCl}
\leanref{Formalism/SharpRestriction/GrowthBounds/CountingStatementReduction.lean}{Formalism.MainNotions.G.CountingStatementReduction.clSum}
\leanref{Formalism/SharpRestriction/GrowthBounds/CountingStatementReduction.lean}{Formalism.MainNotions.G.CountingStatementReduction.blockForm}

\begin{lemma}\label{pd-kappa-defects-segment-mean-theorem-a-average}
Let $A$ and $\Delta$ be finite sets, $R\subset\Delta\times A$ a relation, written $\delta\mathrel Ra$, and
$n,r,d\colon\Delta\times A\to\mathbf R$ functions with $n,d\geq0$; let $\varepsilon,K\geq0$ be real numbers.
Suppose that:
\begin{enumerate}
\item $d(\delta,a)/2\leq r(\delta,a)$ whenever $\delta\mathrel Ra$;
\item for every $\delta\in\Delta$, the sum of $n(\delta,a)$ over the $a$ with $\delta\mathrel Ra$ is at most
  $\varepsilon$ times the sum of $d(\delta,a)$ over these $a$;
\item for all $a,a'\in A$, the sum of $d(\delta,a)$ over the $\delta$ with $\delta\mathrel Ra$ is at most $K$
  times the sum of $d(\delta,a')$ over the $\delta$ with $\delta\mathrel Ra'$.
\end{enumerate}
For $a\in A$ let $N(a)$ and $\rho(a)$ be the sums of $n(\delta,a)$ and of $r(\delta,a)$ over the $\delta$
with $\delta\mathrel Ra$, and read $N(a)/\rho(a)$ as $0$ if $\rho(a)=0$. Then
$\sum_{a\in A}N(a)/\rho(a)\leq2\varepsilon K|A|$.
\end{lemma}
\leanref{Formalism/SharpRestriction/GrowthBounds/CountingStatementReduction.lean}{Formalism.MainNotions.G.CountingStatementReduction.double_count}

\begin{proof}
Let $M(a)$ be the sum of $d(\delta,a)$ over the $\delta$ with $\delta\mathrel Ra$, and $T=\sum_{a\in A}M(a)$.
The $N(a)$ and $M(a)$ are $\geq0$, and $M(a)/2\leq\rho(a)$ by (1). Exchanging the order of summation, (2)
gives $\sum_aN(a)\leq\varepsilon T$. If $T=0$, then every $N(a)$ is $0$, and the sum in question is
$0\leq2\varepsilon K|A|$. Let $T>0$. By (3), $T\leq K|A|M(a)$ for every $a$; so $K|A|>0$ and $M(a)>0$,
whence $\rho(a)\geq M(a)/2>0$ and $T\leq2K|A|\rho(a)$. So $N(a)/\rho(a)\leq N(a)\cdot2K|A|/T$, and summing
over $a$ gives $\sum_aN(a)/\rho(a)\leq\varepsilon T\cdot2K|A|/T=2\varepsilon K|A|$.
\end{proof}

\begin{lemma}\label{pd-kappa-defects-segment-mean-theorem-a-flat}
Let $C_H\geq1$, and let an instance satisfy its hypotheses
(Definition~\ref{pd-kappa-defects-segment-mean-conditions}), with $S_j\leq Y$. Then
$M(\alpha)\leq e^3M(\alpha')$ for all $\alpha,\alpha'\in W^{S_j}$
(Definition~\ref{pd-kappa-defects-segment-mean-theorem-a-classes}).
\end{lemma}
\leanref{Formalism/SharpRestriction/GrowthBounds/CountingStatementReduction.lean}{Formalism.MainNotions.G.CountingStatementReduction.RestFlat}
\leanref{Formalism/SharpRestriction/GrowthBounds/CountingStatementReduction.lean}{Formalism.MainNotions.G.CountingStatementReduction.restFlat_exp}
\leanref{Formalism/SharpRestriction/GrowthBounds/CountingStatementReduction.lean}{Formalism.MainNotions.G.CountingStatementReduction.restFlat_proved}
\leanref{Formalism/SharpRestriction/GrowthBounds/CountingStatementReduction.lean}{Formalism.MainNotions.G.CountingStatementReduction.MCl_eq}
\leanref{Formalism/SharpRestriction/GrowthBounds/CountingStatementReduction.lean}{Formalism.MainNotions.G.CountingStatementReduction.choose_ratio}
\leanref{Formalism/SharpRestriction/GrowthBounds/CountingStatementReduction.lean}{Formalism.MainNotions.G.CountingStatementReduction.Kflat_num}

\begin{proof}
Write $s=S_j$, $N=|T|$ and $u=\lfloor a+a/(kY)+x_a\rfloor$, so that $W\subset[a,u]$.

\emph{The product form.} In the notation of Step~1 of the proof of
Lemma~\ref{pd-kappa-defects-windows-model-segment}, the digit classes compatible with $\alpha$ are the pairs
with $\rho_i\in\mathcal D(P_i,b_i)$ and $\rho'_i\in\mathcal D(P_i,\alpha_i)$ for every $i$, and
$d(\delta;\alpha)$ is the product of a factor that depends only on the $\rho_i$ and the factors
$P_m((\alpha_i-|\rho'_i|)/P_i)$. Summing over the $\rho'_i$ with the identity of that Step~1 gives
$M(\alpha)=c\prod_{i=1}^sP_m(\alpha_i)$, where $c\geq0$ does not depend on $\alpha$.

\emph{One entry.} Let $1\leq x\leq y$ be integers. Then $x^lP_l(y)\leq y^lP_l(x)$ for every $l\geq0$, by
induction on $l$: $(l+1)P_{l+1}(z)=(z+l+1)P_l(z)$, and $x(y+l+1)\leq y(x+l+1)$. The hypotheses give
$a>pb'>0$ (Theorem~\ref{pd-gain-theorem-one}(2)), so $a\geq1$. As $P_m$ is increasing on the integers
$\geq0$ and $a\leq\alpha_i,\alpha'_i\leq u$,
\[
 P_m(\alpha_i)\leq P_m(u)\leq(u/a)^mP_m(a)\leq(u/a)^mP_m(\alpha'_i).
\]
Hence $M(\alpha)\leq((u/a)^m)^sM(\alpha')$.

\emph{The factor.} The hypotheses give $1\leq Y$, $k^2Y^2\leq p$, $N\geq1$ and $\Lambda\geq0$
(Theorem~\ref{pd-gain-theorem-one}(1)), $b'>0$, $C_H\leq Y_b$ and $2Y_bN\Lambda<k$ (its (2)), and
$0\leq x_a\leq2J_0$. As $2Y_bN\geq1$, $\Lambda\leq k$, so $x_a\leq2\Lambda b'/k\leq2b'$. Put
$\theta=1/(kY)+x_a/a$. Then $u/a\leq1+\theta\leq e^\theta$, so $((u/a)^m)^s\leq e^{sm\theta}$. As $m<k$ and
$s\leq Y$, $sm\leq kY$, so $sm/(kY)\leq1$; and $smx_a\leq2kYb'\leq2k^2Y^2b'\leq2pb'<2a$, so $smx_a/a<2$.
So $sm\theta\leq3$.
\end{proof}
% Lean: restFlat_exp hCYb hH hs = restFlat_proved (MCl_eq: mCl_factor and symDim_digits, the identity of
%   Step 1 of Lemma pd-kappa-defects-windows-model-segment; choose_ratio for one entry) with Kflat_le =
%   Kflat_num (hLk, hxb, hsr, t1, t2). The hypothesis s <= Y comes from remaining_of_segment_bound.

\begin{lemma}\label{pd-kappa-defects-segment-mean-theorem-a-blocks}
Let $p$, $\mathcal O\subset\mathcal K$, $v$, $r$, a segment datum $S$, integers
$\alpha_1,\dots,\alpha_s,\alpha_+\geq0$ and a real number $\mu$ be as in
Definition~\ref{pd-kappa-defects-segment-lattice}, put $L=L_S(\alpha,\alpha_+;\mu)$, and for a block $b$ let
$\beta_b(f,f')=\operatorname{sp}(b)[f,f']$. Let $D$ be a real number.
\begin{enumerate}
\item For every kept block $b$, $\operatorname{rk}L_b=\dim_{\mathcal K}V_b$
  (Definition~\ref{pd-kappa-defects-windows-blocks}).
\item For each kept block $b$ let $L_2(b)$ be an $\mathcal O$-submodule of $L_b$, and let $L_2$ be the set
  of the families $(f_b)$ with $f_b\in L_2(b)$ for kept $b$ and $f_b=0$ otherwise. If every $L_2(b)$ is a
  good sublattice of $L_b$ of depth $D$ for $\beta_b$ (Definition~\ref{pd-kappa-defects-shares}), then
  $L_2$ is a good sublattice of $L$ of depth $D$ for $\beta_S$. If every $L_2(b)$ has finite rank, then
  $\operatorname{rk}L_2=\sum_b\operatorname{rk}L_2(b)$, the sum over the kept blocks; if some $L_2(b)$ has
  arbitrarily many linearly independent elements, then $\operatorname{rk}L_2=0$.
\item If $L$ has a good share of depth $D$ and defect $\kappa\geq0$ for $\beta_S$, then
  $\kappa_S(\mu,D;\alpha,\alpha_+)\leq\kappa$. Always $\kappa_S(\mu,D;\alpha,\alpha_+)\leq1$.
\end{enumerate}
\end{lemma}
\leanref{Formalism/SharpRestriction/GrowthBounds/CountingStatementReduction.lean}{Formalism.MainNotions.G.CountingStatementReduction.restRk_proved}
\leanref{Formalism/SharpRestriction/GrowthBounds/CountingStatementReduction.lean}{Formalism.MainNotions.G.CountingStatementReduction.finrank_blockLat}
\leanref{Formalism/SharpRestriction/GrowthBounds/CountingStatementReduction.lean}{Formalism.MainNotions.G.CountingStatementReduction.segSub}
\leanref{Formalism/SharpRestriction/GrowthBounds/CountingStatementReduction.lean}{Formalism.MainNotions.G.CountingStatementReduction.segSub_good}
\leanref{Formalism/SharpRestriction/GrowthBounds/CountingStatementReduction.lean}{Formalism.MainNotions.G.CountingStatementReduction.finrank_segSub}
\leanref{Formalism/SharpRestriction/GrowthBounds/CountingStatementReduction.lean}{Formalism.MainNotions.G.CountingStatementReduction.badShare_le_of_goodShare}

\begin{proof}
(1) As in the proof of Lemma~\ref{pd-kappa-defects-windows-ranks}, every element of $V_b$ has a nonzero
multiple in $L_b$, elements of $L_b$ that are linearly independent over $\mathcal O$ are linearly
independent over $\mathcal K$, and nonzero multiples in $L_b$ of a basis of $V_b$ are linearly independent
over $\mathcal O$. So $\operatorname{rk}L_b=\dim V_b$, which is finite.

(2) As $L_2(b)\subset L_b$, $L_2\subset L$. Let $x\in L$ and $c\neq0$ with $cx\in L_2$. For kept $b$,
$x_b\in L_b$ and $cx_b\in L_2(b)$, so $x_b\in L_2(b)$ as $L_2(b)$ is good; for the other $b$, $x_b=0$. So
$x\in L_2$. Let $u$ be primitive in $L_2$, and $\pi$ a uniformizer of $\mathcal O$. Some kept $b$ has $u_b$
primitive in $L_2(b)$: otherwise, for each kept $b$, $u_b=c_by_b$ with $y_b\in L_2(b)$ and $c_b$ not a unit,
so $c_b\in\pi\mathcal O$ and $u_b=\pi y'_b$ with $y'_b\in L_2(b)$; then $u=\pi y'$ with $y'\in L_2$, and
$\pi$ is not a unit. As $L_2(b)$ is good, there is $u'\in L_b$ with $\beta_b(u_b,u')\neq0$ and
$v(\beta_b(u_b,u'))\leq D$. The family $f$ with $f_b=u'$ and all other entries $0$ lies in $L$, and
$\beta_S(u,f)=\beta_b(u_b,u')$. So $L_2$ is good. Finally, $L_2$ is the direct sum of the $L_2(b)$ over
the finitely many kept blocks, which gives the ranks.

(3) If $\operatorname{rk}L=0$, the segment defect is $0$. Otherwise let $n$ be the largest rank of a good
sublattice of $L$ of depth $D$ (Lemma~\ref{pd-kappa-defects-bad-share}). A good sublattice $L_1$ with
$\operatorname{rk}L_1\geq(1-\kappa)\operatorname{rk}L$ has $\operatorname{rk}L_1\leq n$, so
$\kappa_S(\mu,D;\alpha,\alpha_+)=(\operatorname{rk}L-n)/\operatorname{rk}L\leq\kappa$; and as $n\geq0$, it is
at most $1$.
\end{proof}
% Lean: (1) is restRk_proved = finrank_blockLat (finrank_eq_finrank_span, span_inf_bbSpec with
%   exists_smul_mem_intPol and intPol_le_bbSpec); (2) is segSub_good, finrank_segSub, finrank_segSub_zero
%   and segLat_eq_segSub; (3) is badShare_le_of_goodShare and badShare_le_one.

\begin{definition}\label{pd-kappa-defects-segment-mean-theorem-a-levels}
Let an instance be given, with $s$, $t$, $W$, $S$, the $P_i$, the digit classes and $A(\delta)$ as in
Definition~\ref{pd-kappa-defects-segment-mean-theorem-a-classes}, and put $u_l=\eta_j+l$ for $1\leq l\leq t$.
For $\alpha\in W^s$, $\alpha_+\in W$ and a block $b$, form the numbers, letters and degrees $\delta$ of
Definition~\ref{pd-kappa-defects-segment-lattice} for $S$, $\alpha$, $\alpha_+$, $\mu$ and $r=m$, with
polynomials over $\mathbf Q_p$ in place of $\mathcal K$.
\begin{enumerate}
\item A \emph{level} of $b$ is a pair $\omega=(j,j_+)$ as in
  Definition~\ref{pd-kappa-defects-segment-lattice}(5), and $\delta'_\omega$ is the lowered degree $\delta'$
  formed there. The levels depend on $\alpha$ and not on $\alpha_+$.
\item Let $Y_\omega$ be the $\mathbf Q_p$-vector space of the polynomials in the letters, with coefficients in
  $\mathbf Q_p$ and of total degree $\delta'_\omega(\lambda)$ in the variables of each letter $\lambda$, that
  lie in $Y(\mathcal L_l,M_l)$ for $1\leq l\leq t$ and in $Y(\mathcal L_+,M_+)$
  (Definition~\ref{pd-kappa-defects-segment-lattice}(4) and~(5), over $\mathbf Q_p$), and let
  $Y^{\mathbf Z}_\omega$ be the set of its elements with coefficients in $\mathbf Z_p$. The \emph{level
  lattice} $\Lambda_\omega$ is the $\mathbf Z_p$-span of the products $fg$ with $f\in Y^{\mathbf Z}_\omega$
  and $g$ a homogeneous polynomial of degree $b_+-j_+$ in $\xi_1,\dots,\xi_m$ with coefficients in
  $\mathbf Z_p$. For $n\geq0$ put $N_\xi(n)=\binom{n+m-1}{m-1}$, the number of monomials of degree $n$ in
  $\xi_1,\dots,\xi_m$.
\item For vectors of variables $w=(w_c)_{c\leq m}$ and $e=(e_c)_{c\leq m}$ and integers $o,n,z\geq0$, let
  $E_{we}$ be the multiplication by $\sum_cw_ce_c$, $F_{we}=\sum_c\partial^2/\partial w_c\partial e_c$, and
  \[
   \phi(w,e;o,n,z)=\sum_{i=0}^{o}\frac{(-1)^i}{(i!)^2\binom{o+n+m+z-1}{i}}\,E_{we}^iF_{we}^i .
  \]
  Put $\Phi^Y_\omega=\prod_{l=1}^t\phi\big(x_{u_l},y_{u_l};\delta'_\omega(x_{u_l}),\delta'_\omega(y_{u_l}),
  \hat b_{u_l}\dot-o_{u_l}\big)$ and $\Phi_\omega=\phi\big(x_+,\xi;o_+-j_+,b_+-j_+,\alpha_+\dot-o_+\big)\Phi^Y_\omega$;
  the factors act on disjoint sets of variables and commute.
\item With the Fock pairing $[\ ,\ ]$ over $\mathbf Q_p$ (Definition~\ref{pd-weight-lattices}(2)), the
  \emph{level pairings} are
  \[
   \beta_\omega(f,f')=\frac{[\Phi_\omega f,f']}{(b_+-j_+)!\prod_\lambda\delta'_\omega(\lambda)!}\quad\text{on }
   \Lambda_\omega,\qquad
   \beta^Y_\omega(f,f')=\frac{[\Phi^Y_\omega f,f']}{\prod_\lambda\delta'_\omega(\lambda)!}\quad\text{on }
   Y^{\mathbf Z}_\omega,
  \]
  the products over the letters $\lambda$; they are linear in $f$.
\item For a $\mathbf Z_p$-submodule $\Lambda$ of a finite-dimensional $\mathbf Q_p$-vector space, a map
  $\beta\colon\Lambda\times\Lambda\to\mathbf Q_p$ that is linear in the first variable, and an integer $D$, let
  $\operatorname{good}_D(\beta,\Lambda)$ be the largest $n$ for which there are
  $x_1,\dots,x_n,y_1,\dots,y_n\in\Lambda$ with $\beta(x_i,y_l)=p^D\delta_{il}$ for all $i,l$ (such $x_i$ are
  linearly independent, so $n\leq\operatorname{rk}\Lambda$), and
  $\operatorname{bad}_D(\beta,\Lambda)=\operatorname{rk}\Lambda-\operatorname{good}_D(\beta,\Lambda)$.
\item Put $\operatorname{bad}_\omega(\alpha,\alpha_+,b;D)=\operatorname{bad}_D(\beta_\omega,\Lambda_\omega)$,
  $\operatorname{bad}^Y_\omega(\alpha,b;D)=\operatorname{bad}_D(\beta^Y_\omega,Y^{\mathbf Z}_\omega)$, and let the
  \emph{badness} $g_D(\alpha,\alpha_+,b)$ be the sum of the $\operatorname{bad}_\omega(\alpha,\alpha_+,b;D)$ over
  the levels $\omega$ of $b$.
\item Let $\delta=(\rho,\rho')$ be a digit class. For $1\leq i\leq t$ let $\tau(\delta,i)$ be the largest
  $y\in W$ with $|\rho'_{u_i}|\leq y$ and $P_{u_i}\mid y-|\rho'_{u_i}|$, and $0$ if there is none, and for
  $\alpha\in W^s$ let $\alpha^{(i)}$ be $\alpha$ with its $u_i$-th entry replaced by $\tau(\delta,i)$. Let
  $I(\delta,\alpha)$ be the set of the pairs $(b,\omega)$ of a block $b$ kept for $\alpha$ of class $\delta$ and
  a level $\omega$ of $b$.
\item Put $c_S=3(s+2)m^2\ell$, with $\ell$ and $\kappa_0$ as in Definition~\ref{pd-kappa-defects-numbers};
  for a real number $C_M$ put $D_1(C_M)=t(Yc_S+C_Mc_S)+C_Mc_S$; and put
  $D_{\mathrm{lv}}=10(s+1)(m+1)\ell+\kappa_0$.
\end{enumerate}
\end{definition}
\leanref{Formalism/SharpRestriction/GrowthBounds/MonomialFactorials/BlockBoundUniformiser.lean}{Formalism.MainNotions.G.SymmetricFormBlocks.gB}
\leanref{Formalism/SharpRestriction/GrowthBounds/MonomialFactorials/BlockBoundUniformiser.lean}{Formalism.MainNotions.G.SymmetricFormBlocks.levels}
\leanref{Formalism/SharpRestriction/GrowthBounds/MonomialFactorials/BlockBoundUniformiser.lean}{Formalism.MainNotions.G.SymmetricFormBlocks.lv}
\leanref{Formalism/SharpRestriction/GrowthBounds/MonomialFactorials/BlockBoundUniformiser.lean}{Formalism.MainNotions.G.SymmetricFormBlocks.Ljj}
\leanref{Formalism/MahlerMeasure/GoodCount.lean}{Formalism.GM.M.Y}
\leanref{Formalism/MahlerMeasure/GoodCount.lean}{Formalism.GM.M.YZ}
\leanref{Formalism/MahlerMeasure/WithMargin.lean}{Formalism.GM.M.XiPolZ}
\leanref{Formalism/SharpRestriction/GrowthBounds/MonomialFactorials/BlockBoundUniformiser.lean}{Formalism.MainNotions.G.SymmetricFormBlocks.dimY}
\leanref{Formalism/SharpRestriction/GrowthBounds/MonomialFactorials/BlockBoundUniformiser.lean}{Formalism.MainNotions.G.SymmetricFormBlocks.dXi}
\leanref{Formalism/SharpRestriction/GrowthBounds/MonomialFactorials/BlockBoundUniformiser.lean}{Formalism.MainNotions.G.SymmetricFormBlocks.rL}
\leanref{Formalism/MahlerMeasure/LevelSum.lean}{Formalism.GM.tieE}
\leanref{Formalism/MahlerMeasure/LevelSum.lean}{Formalism.GM.tieF}
\leanref{Formalism/MahlerMeasure/LevelSum.lean}{Formalism.GM.tiePhi}
\leanref{Formalism/MahlerMeasure/LevelLattice.lean}{Formalism.GM.M.BlockData.phiAll}
\leanref{Formalism/MahlerMeasure/StatementTM.lean}{Formalism.GM.M.TieData.phis}
\leanref{Formalism/SharpRestriction/GrowthBounds/MonomialFactorials/BlockBoundUniformiser.lean}{Formalism.MainNotions.G.SymmetricFormBlocks.degFact}
\leanref{Formalism/SharpRestriction/GrowthBounds/MonomialFactorials/BlockBoundUniformiser.lean}{Formalism.MainNotions.G.SymmetricFormBlocks.bejj}
\leanref{Formalism/MahlerMeasure/GoodCount.lean}{Formalism.GM.M.fockN}
\leanref{Formalism/MahlerMeasure/GoodCount.lean}{Formalism.GM.M.goodCount}
\leanref{Formalism/MahlerMeasure/GoodCount.lean}{Formalism.GM.M.bad}
\leanref{Formalism/SharpRestriction/GrowthBounds/MonomialFactorials/BlockBoundUniformiser.lean}{Formalism.MainNotions.G.SymmetricFormBlocks.badLvl}
\leanref{Formalism/SharpRestriction/GrowthBounds/MonomialFactorials/BlockBoundUniformiser.lean}{Formalism.MainNotions.G.SymmetricFormBlocks.badY}
\leanref{Formalism/SharpRestriction/GrowthBounds/MonomialFactorials/BlockBoundUniformiser.lean}{Formalism.MainNotions.G.SymmetricFormBlocks.bLvl}
\leanref{Formalism/SharpRestriction/GrowthBounds/MonomialFactorials/BlockBoundUniformiser.lean}{Formalism.MainNotions.G.SymmetricFormBlocks.bY}
\leanref{Formalism/SharpRestriction/GrowthBounds/MonomialFactorials/BlockBoundUniformiser.lean}{Formalism.MainNotions.G.SymmetricFormBlocks.blockBadness}
\leanref{Formalism/SharpRestriction/GrowthBounds/MonomialFactorials/BlockBoundUniformiser.lean}{Formalism.MainNotions.G.SymmetricFormBlocks.cbox}
\leanref{Formalism/SharpRestriction/GrowthBounds/MonomialFactorials/BlockBoundUniformiser.lean}{Formalism.MainNotions.G.SymmetricFormBlocks.topA}
\leanref{Formalism/SharpRestriction/GrowthBounds/MonomialFactorials/BlockBoundUniformiser.lean}{Formalism.MainNotions.G.SymmetricFormBlocks.upd}
\leanref{Formalism/SharpRestriction/GrowthBounds/MonomialFactorials/BlockBoundUniformiser.lean}{Formalism.MainNotions.G.SymmetricFormBlocks.idx}
\leanref{Formalism/SharpRestriction/GrowthBounds/MonomialFactorials/BlockBoundUniformiser.lean}{Formalism.MainNotions.G.SymmetricFormBlocks.kapS}
\leanref{Formalism/SharpRestriction/GrowthBounds/MonomialFactorials/BlockBoundUniformiser.lean}{Formalism.MainNotions.G.SymmetricFormBlocks.D1r}
\leanref{Formalism/SharpRestriction/GrowthBounds/MonomialFactorials/BlockBoundUniformiser.lean}{Formalism.MainNotions.G.SymmetricFormBlocks.levelPassageLoss}

\begin{lemma}\label{pd-kappa-defects-segment-mean-theorem-a-free}
Let $C_H\geq2$ and $C_{E1},C_{E2}\geq0$, let an instance satisfy its hypotheses
(Definition~\ref{pd-kappa-defects-segment-mean-conditions}), and let $\delta$ be a digit class. Then
$I(\delta,\alpha)=I(\delta,\alpha')$ for all $\alpha,\alpha'\in A(\delta)$
(Definition~\ref{pd-kappa-defects-segment-mean-theorem-a-levels}).
\end{lemma}
\leanref{Formalism/SharpRestriction/GrowthBounds/MonomialFactorials/BlockBoundUniformiser.lean}{Formalism.MainNotions.G.SymmetricFormBlocks.RestFree}
\leanref{Formalism/SharpRestriction/GrowthBounds/MonomialFactorials/DigitsOKOfCompat.lean}{Formalism.MainNotions.G.SymmetricFormBlocks.restFree_proved}
\leanref{Formalism/SharpRestriction/GrowthBounds/MonomialFactorials/DigitsOKOfCompat.lean}{Formalism.MainNotions.G.SymmetricFormBlocks.idx_sub}
\leanref{Formalism/SharpRestriction/GrowthBounds/MonomialFactorials/CoarseCLower.lean}{Formalism.MainNotions.G.SymmetricFormBlocks.segFacts}

\begin{proof}
\emph{Facts.} Write $\delta=(\rho,\rho')$. As in the proof of
Lemma~\ref{pd-kappa-defects-segment-mean-letters-degrees}, the hypotheses give $(m+1)^2\leq p$, $m\geq4$,
$a>pb'\geq25b'$ and $b_i\leq5b'/2$ for every $i$. By hypothesis (3) of Theorem~\ref{pd-gain-theorem-one},
as $C_H$, $|T|$, $Y$ and $c_*+|T|c_{cl}+n'c_{cl}$ are $\geq1$, $kQ\leq b'$; and $P_i\leq p^{g_j}\leq Q$, so
$(m+1)P_i\leq b'$. By hypothesis (2), $4\Lambda\leq2Y_b|T|\Lambda<k$, as $Y_b\geq C_H\geq2$ and
$\Lambda\geq0$. As $C_{E1},C_{E2}\geq0$, $E_0\geq0$, so $\hat a\leq a/p^{g_j}$.

By symmetry it suffices to show $I(\delta,\alpha)\subset I(\delta,\alpha')$. For $\alpha''\in A(\delta)$ write
$\hat\alpha''_i=(\alpha''_i-|\rho'_i|)/P_i$. As the entries of $\rho'_i$ are less than $P_i$,
$|\rho'_i|<(m+1)P_i\leq b'$, so $\hat\alpha'_iP_i>a-b'$.

Let $(b,\omega)\in I(\delta,\alpha)$. Whether $b$ is kept depends on $\alpha$ only through its first
condition, which holds for $\alpha'$ as $\alpha'\in A(\delta)$, and through $\varsigma_i\leq\hat\alpha_i$ for
$i\leq\eta_j$ (Definition~\ref{pd-kappa-defects-segment-lattice}(2)). For $i\leq\eta_j$,
$\varsigma_i\leq2\Lambda\hat a/(m+1)$, so $\varsigma_iP_i\leq2\Lambda(a/p^{g_j})p^{g_j}/(m+1)<a/2<a-b'<
\hat\alpha'_iP_i$, and $\varsigma_i\leq\hat\alpha'_i$. So $b$ is kept for $\alpha'$, of class $\delta$. The level
conditions on $\omega=(j,j_+)$ involve $\alpha$ only through $j_l\leq\delta(y_{u_l})=\hat\alpha_{u_l}$. As $b$ is
kept and $u_l>\eta_j$, $j_l\leq o_{u_l}\leq\hat b_{u_l}$; and $b_{u_l}-|\rho_{u_l}|\leq5b'/2<a-b'<
\alpha'_{u_l}-|\rho'_{u_l}|$, so $\hat b_{u_l}\leq\hat\alpha'_{u_l}$. So $\omega$ is a level of $b$ for $\alpha'$.
\end{proof}
% Lean: restFree_proved hS = antisymmetry of idx_sub hS, with hS = segFacts (2 <= CYb, 1 <= Cb, 0 <= CD,
%   0 <= CE1, CE2): kP, P1, PQ, La4, E0, bTop and 25 b' < a. The kept conditions are hk with hSG for the
%   condition on sg; the level is mem_lv with hAC (coarseA <= coarseC for aTop').

\begin{lemma}\label{pd-kappa-defects-segment-mean-theorem-a-ratio}
Let $C_H\geq2$ and $C_{E1},C_{E2}\geq0$, let an instance satisfy its hypotheses, and let $\delta$ be a digit
class, $\alpha\in A(\delta)$ and $1\leq i\leq t$. Then $\alpha^{(i)}\in A(\delta)$ and
$d(\delta;\alpha^{(i)})\leq e^3d(\delta;\alpha)$
(Definitions~\ref{pd-kappa-defects-segment-mean-theorem-a-classes}
and~\ref{pd-kappa-defects-segment-mean-theorem-a-levels}).
\end{lemma}
\leanref{Formalism/SharpRestriction/GrowthBounds/MonomialFactorials/BlockBoundUniformiser.lean}{Formalism.MainNotions.G.SymmetricFormBlocks.RestRatio}
\leanref{Formalism/SharpRestriction/GrowthBounds/MonomialFactorials/RatioConditionOfSegFacts.lean}{Formalism.MainNotions.G.SymmetricFormBlocks.ratioCondition_of_segFacts}
\leanref{Formalism/SharpRestriction/GrowthBounds/MonomialFactorials/RatioConditionOfSegFacts.lean}{Formalism.MainNotions.G.SymmetricFormBlocks.ratio_num}

\begin{proof}
Write $\delta=(\rho,\rho')$, $u=u_i$, $P=P_u$, $c=|\rho'_u|$ and $\tau=\tau(\delta,i)$. As $\alpha\in A(\delta)$,
$\alpha_u\in W$, $c\leq\alpha_u$ and $P\mid\alpha_u-c$; so $\tau$ is the largest element of $W$ with these
properties, and $\alpha_u\leq\tau$. The other entries of $\alpha^{(i)}$ are those of $\alpha$, so
$\alpha^{(i)}\in A(\delta)$. Put $x=(\alpha_u-c)/P$ and $y=(\tau-c)/P$, integers with $x\leq y$; then
$d(\delta;\alpha^{(i)})$ is $d(\delta;\alpha)$ with the factor $P_m(x)$ replaced by $P_m(y)$.

The facts in the proof of Lemma~\ref{pd-kappa-defects-segment-mean-theorem-a-free} give $(m+1)P\leq b'$
and $4\Lambda<m+1$, and $4(m+1)\leq(m+1)^2\leq p$, so $(m+1)b'\leq pb'/4<a/4$; also $Y\geq1$ and
$0\leq x_a\leq2J_0$. As $c<(m+1)P$ and $\alpha_u\geq a$, $xP=\alpha_u-c\geq a-(m+1)P\geq3a/4$; as
$P\leq b'<3a/4$, $x\geq1$. Next $(y-x)P=\tau-\alpha_u\leq a/(kY)+x_a$, $ma/(kY)\leq a$, and
$mx_a\leq(m+1)x_a\leq2\Lambda b'\leq(m+1)b'/2\leq a/8$. So $m(y-x)P\leq9a/8\leq3xP$, and $m(y-x)\leq3x$. By
the inequality $x^mP_m(y)\leq y^mP_m(x)$ of the proof of
Lemma~\ref{pd-kappa-defects-segment-mean-theorem-a-flat} and $y/x\leq e^{(y-x)/x}$,
\[
 P_m(y)\leq(y/x)^mP_m(x)\leq e^{m(y-x)/x}P_m(x)\leq e^3P_m(x).
\]
As the other factors of $d(\delta;\alpha)$ are $\geq0$, $d(\delta;\alpha^{(i)})\leq e^3d(\delta;\alpha)$.
\end{proof}
% Lean: ratioCondition_of_segFacts hS: upd_mem and topA_spec for the membership; coarseC_lower (hxP2, hx1),
%   htop, hrk and hJ for m(y - x) <= 3x (hnum); ratio_num (choose_ratio and add_one_le_exp); prod_ratio for
%   the other factors.

\begin{lemma}\label{pd-kappa-defects-segment-mean-theorem-a-depth}
Let $C_H\geq1$, $C_M\geq0$ and $C\geq18(1+C_M)+20$, and let an instance satisfy its hypotheses. Then
$D_1(C_M)\geq0$ and $YD_1(C_M)+D_{\mathrm{lv}}\leq D_A$, where $D_A$ is formed with $C_A=C$
(Definitions~\ref{pd-kappa-defects-numbers} and~\ref{pd-kappa-defects-segment-mean-theorem-a-levels}).
\end{lemma}
\leanref{Formalism/SharpRestriction/GrowthBounds/MonomialFactorials/BlockBoundUniformiser.lean}{Formalism.MainNotions.G.SymmetricFormBlocks.RestDepth}
\leanref{Formalism/SharpRestriction/GrowthBounds/MonomialFactorials/KappaZeroNonneg.lean}{Formalism.MainNotions.G.SymmetricFormBlocks.restDepth_proved}
\leanref{Formalism/SharpRestriction/GrowthBounds/MonomialFactorials/KappaZeroNonneg.lean}{Formalism.MainNotions.G.SymmetricFormBlocks.depth_num}

\begin{proof}
Write $s=S_j$ and $N=|T|$. The hypotheses give $Y\geq1$, $N_1\geq1$ and $s\leq N_1$
(Theorem~\ref{pd-gain-theorem-one}(1)), and $a\geq1$; also $t\leq s$, $\ell\geq1$ and
$\kappa_0=2Nk\log_pa\geq0$. So $c_S\geq0$ and $D_1(C_M)\geq0$. Put $B=N_1k^2\ell$. As $s+2\leq3N_1$ and
$m^2\leq k^2$, $c_S\leq9B$ and $tc_S\leq9N_1B$; and $10(s+1)k\ell\leq20B$. Using $Y\leq Y^2$ and $1\leq N_1$,
\[
 YD_1(C_M)+10(s+1)k\ell=Y^2tc_S+C_MYtc_S+C_MYc_S+10(s+1)k\ell\leq(29+18C_M)Y^2N_1B\leq CY^2N_1^2k^2\ell .
\]
Adding $\kappa_0\leq2\kappa_0$ gives the claim, as $D_A=CY^2N_1^2k^2\ell+2\kappa_0$.
\end{proof}
% Lean: restDepth_proved CM C hCYb hCM hC hH: hiFacts, lgA_ge_one_facts, kap0A_nonneg_of_hiFacts and
%   depth_num.

\begin{lemma}\label{pd-kappa-defects-segment-mean-theorem-a-block-bound-base}
Let an instance be given (Definition~\ref{pd-kappa-defects-segment-mean-conditions}). As $\mathcal R$ is
complete and $p$ lies in its maximal ideal, the map $\mathbf Z\to\mathcal R$ extends to a homomorphism
$\mathbf Z_p\to\mathcal R$, and so to an embedding $\iota_p\colon\mathbf Q_p\to\mathcal K$. For a polynomial $P$
over $\mathbf Q_p$ let $P^\iota$ be the polynomial over $\mathcal K$ obtained by applying $\iota_p$ to its
coefficients. Let $\alpha\in W^{S_j}$, $\alpha_+\in W$ and $b$ a block, and form the generators of
Definition~\ref{pd-kappa-defects-segment-lattice}(5) over $\mathbf Q_p$, as in
Definition~\ref{pd-kappa-defects-segment-mean-theorem-a-levels}.
\begin{enumerate}
\item If $P$ is a generator of level $\omega$ over $\mathbf Q_p$, then $P^\iota$ is a generator of level $\omega$
  over $\mathcal R$. So $P\mapsto P^\iota$ maps the $\mathbf Q_p$-span of the generators over $\mathbf Q_p$ of all
  levels of $b$ into $V_b$ (Definition~\ref{pd-kappa-defects-windows-blocks}).
\item If $P_1,\dots,P_n$ are polynomials over $\mathbf Q_p$ that are linearly independent over $\mathbf Q_p$,
  then $P_1^\iota,\dots,P_n^\iota$ are linearly independent over $\mathcal K$.
\item For polynomials $P$ and $Q$ over $\mathbf Q_p$, $[P^\iota,Q^\iota]=\iota_p([P,Q])$, where $[\ ,\ ]$ is the
  Fock pairing (Definition~\ref{pd-weight-lattices}(2)).
\end{enumerate}
\end{lemma}
\leanref{Formalism/SharpRestriction/GrowthBounds/MonomialFactorials/PadicIntegersEmbeddingMapPrime.lean}{Formalism.MainNotions.G.SymmetricFormBlocks.iotaX}
\leanref{Formalism/SharpRestriction/GrowthBounds/MonomialFactorials/PadicIntegersEmbeddingMapPrime.lean}{Formalism.MainNotions.G.SymmetricFormBlocks.iotaX_int}
\leanref{Formalism/SharpRestriction/GrowthBounds/MonomialFactorials/MemIntPolIff.lean}{Formalism.MainNotions.G.SymmetricFormBlocks.map_gens}
\leanref{Formalism/SharpRestriction/GrowthBounds/MonomialFactorials/MemIntPolIff.lean}{Formalism.MainNotions.G.SymmetricFormBlocks.map_SSZ}
\leanref{Formalism/SharpRestriction/GrowthBounds/MonomialFactorials/StructureSheafH2BaseChange.lean}{Formalism.MainNotions.G.SymmetricFormBlocks.map_ycond}
\leanref{Formalism/SharpRestriction/GrowthBounds/MonomialFactorials/LiMap.lean}{Formalism.MainNotions.G.SymmetricFormBlocks.li_map}
\leanref{Formalism/SharpRestriction/GrowthBounds/MonomialFactorials/BlockBoundOfDualFamilies.lean}{Formalism.MainNotions.G.SymmetricFormBlocks.fock_map}

\begin{proof}
(1) The factors $\hat E_l$ and $\hat E_+$ and the powers of the $z_l$ and $z_+$ have integer coefficients and are
unchanged. A polynomial $g$ in the $\xi_c$ with coefficients in $\mathbf Z_p$ goes to one with coefficients in
$\mathcal R$ of the same degree. A polynomial $f\in Y^{\mathbf Z}_\omega$ goes to a polynomial with coefficients in
$\mathcal R$ of the same degrees in the letters, and $f^\iota$ lies in $Y(\mathcal L_l,M_l)$ over $\mathcal K$ for
every $l$ and in $Y(\mathcal L_+,M_+)$: for an initial segment $\mathcal L$ and an integer $M$, the
$\mathcal K$-span of the $h^\iota$ with $h$ in the space $Y(\mathcal L,M)$ over $\mathbf Q_p$ consists of
polynomials whose monomials have weight at most $M$, and it is stable under the $E^{\mathcal L}_{cd}$, which
have integer coefficients and so commute with $P\mapsto P^\iota$; so it lies in the largest such subspace,
$Y(\mathcal L,M)$ over $\mathcal K$. So $P^\iota$ is a generator of level $\omega$ over $\mathcal R$. As
$(cP)^\iota=\iota_p(c)P^\iota$ and $V_b$ is the $\mathcal K$-span of these generators, the second statement
follows.

(2) Extend $P_1,\dots,P_n$ to a $\mathbf Q_p$-basis of the polynomial ring over $\mathbf Q_p$. The polynomial ring
over $\mathcal K$ is its base change along $\iota_p$, so the images of this basis form a $\mathcal K$-basis of
it; in particular the $P_i^\iota$ are linearly independent.

(3) The coefficients of $P^\iota$ and $Q^\iota$ are the images under $\iota_p$ of those of $P$ and $Q$, and the
factorials $\kappa!$ in the Fock pairing are integers.
\end{proof}
% Lean: iotaX X = iotaQ (IsFractionRing.lift of iota0, injective by padic_integers_embedding_injective), with
%   iotaX_int; (1) is map_gens (map_Ehat, map_Eplus, map_XiPolZ, map_YZ = map_Y with map_ycond and
%   map_mem_restrictSupport, and map_int) and map_SSZ with map_smul_Q; (2) is li_map (Basis.extend and
%   Algebra.TensorProduct.basis with algebraTensorAlgEquiv); (3) is fock_map.

\begin{lemma}\label{pd-kappa-defects-segment-mean-theorem-a-block-bound-rank}
Let an instance be given, $(\alpha,\alpha_+)\in W^{S_j}\times W$ and $b$ a block kept for $\alpha$. Then
$\operatorname{rk}L_b$ is at most the sum of the ranks $\operatorname{rk}\Lambda_\omega$ of the level lattices over
the levels $\omega$ of $b$ (Definitions~\ref{pd-kappa-defects-segment-mean-theorem-a-classes}
and~\ref{pd-kappa-defects-segment-mean-theorem-a-levels}).
\end{lemma}
\leanref{Formalism/SharpRestriction/GrowthBounds/MonomialFactorials/BlockBoundUniformiser.lean}{Formalism.MainNotions.G.SymmetricFormBlocks.RestBCdim}
\leanref{Formalism/SharpRestriction/GrowthBounds/MonomialFactorials/MemLevelsIff.lean}{Formalism.MainNotions.G.SymmetricFormBlocks.restBCdim_of}
\leanref{Formalism/SharpRestriction/GrowthBounds/MonomialFactorials/RestBCdimOfYbc.lean}{Formalism.MainNotions.G.SymmetricFormBlocks.restBClvl_of}
\leanref{Formalism/SharpRestriction/GrowthBounds/MonomialFactorials/RestDecompProved.lean}{Formalism.MainNotions.G.SymmetricFormBlocks.restYbc_of}
\leanref{Formalism/SharpRestriction/GrowthBounds/MonomialFactorials/CmapMemRestrictSupport.lean}{Formalism.MainNotions.G.SymmetricFormBlocks.restFunY_proved}

\begin{proof}
By Lemma~\ref{pd-kappa-defects-segment-mean-theorem-a-blocks}(1), $\operatorname{rk}L_b=\dim V_b$. For a level
$\omega$ let $V^{\mathcal K}_\omega$ be the $\mathcal K$-span of the generators of level $\omega$ over $\mathcal R$.
Then $V_b$ is the sum of the $V^{\mathcal K}_\omega$ over the finitely many levels, and they lie in $V_b$, which
is finite-dimensional (Lemma~\ref{pd-kappa-defects-windows-ranks}). So $\dim V_b$ is at most the sum of the
$\dim V^{\mathcal K}_\omega$, and it suffices to show $\dim V^{\mathcal K}_\omega\leq\operatorname{rk}\Lambda_\omega$ for
each level $\omega=(j,j_+)$.

Let $W_\omega$ be the $\mathbf Q_p$-span of $\Lambda_\omega$. It is finite-dimensional, as $\Lambda_\omega$ lies in the
span of the products $y\xi^\beta$ of the elements $y$ of a basis of $Y_\omega$ with the monomials $\xi^\beta$ of
degree $b_+-j_+$; and $\operatorname{rk}\Lambda_\omega=\dim W_\omega$, as $\mathbf Q_p$ is the fraction field of
$\mathbf Z_p$. Put $F'_\omega=\prod_l\hat E_l^{j_l}\cdot\hat E_+^{j_+}\cdot\prod_lz_l^{\hat b_{u_l}\dot-o_{u_l}}\cdot
z_+^{\alpha_+\dot-o_+}$, a polynomial with integer coefficients, and let $\iota_p$ and $P\mapsto P^\iota$ be as in
Lemma~\ref{pd-kappa-defects-segment-mean-theorem-a-block-bound-base}. The map $h\mapsto F'_\omega h^\iota$ on
$W_\omega$ is additive and satisfies $F'_\omega(ch)^\iota=\iota_p(c)F'_\omega h^\iota$; so the $\mathcal K$-span $T$ of
its image is spanned by the images of a basis of $W_\omega$, and $\dim T\leq\dim W_\omega$.

We show $V^{\mathcal K}_\omega\subset T$. A generator of level $\omega$ over $\mathcal R$ is $F'_\omega fg$, with $f$ a
polynomial over $\mathcal R$ as in Definition~\ref{pd-kappa-defects-segment-lattice}(5) and $g$ a homogeneous
polynomial of degree $b_+-j_+$ in the $\xi_c$ over $\mathcal R$. The polynomial $g$ is an $\mathcal R$-combination
of monomials $\xi^\beta=(\xi^\beta)^\iota$. For $f$, view $\mathcal K$ as a $\mathbf Q_p$-vector space through
$\iota_p$, let $U$ be the finite-dimensional subspace spanned by the coefficients of $f$, $u_1,\dots,u_n$ a
basis of $U$, and $\varphi_1,\dots,\varphi_n\colon\mathcal K\to\mathbf Q_p$ $\mathbf Q_p$-linear maps that restrict to
the coordinate functions of this basis on $U$. For a polynomial $h$ over $\mathcal K$ let $\varphi_k(h)$ be the
polynomial over $\mathbf Q_p$ obtained by applying $\varphi_k$ to its coefficients. Then
$f=\sum_ku_k\varphi_k(f)^\iota$, comparing coefficients. Each $\varphi_k(f)$ lies in $Y_\omega$: its monomials are
among those of $f$, so it has the degrees $\delta'_\omega$ and the weight bounds; and for an initial segment
$\mathcal L$ and an integer $M$, the $\mathbf Q_p$-span of the $\varphi_k(h)$ with $h$ in $Y(\mathcal L,M)$ over
$\mathcal K$ consists of polynomials of weight at most $M$ and is stable under the $E^{\mathcal L}_{cd}$, which
commute with $\varphi_k$ as they have integer coefficients, so it lies in $Y(\mathcal L,M)$ over $\mathbf Q_p$.
For $e_k\geq0$ large, $p^{e_k}\varphi_k(f)$ has coefficients in $\mathbf Z_p$, so it lies in $Y^{\mathbf Z}_\omega$,
and $f=\sum_k\iota_p(p^{-e_k})u_k(p^{e_k}\varphi_k(f))^\iota$. So $F'_\omega fg$ is a $\mathcal K$-combination of
the $F'_\omega(h\xi^\beta)^\iota$ with $h\in Y^{\mathbf Z}_\omega$, and $h\xi^\beta\in\Lambda_\omega$. Hence
$F'_\omega fg\in T$, and $\dim V^{\mathcal K}_\omega\leq\dim T\leq\dim W_\omega=\operatorname{rk}\Lambda_\omega$.
\end{proof}
% Lean: restBCdim_of (restBClvl_of (restYbc_of_funY restFunY_proved)): finrank_blockLat, the inclusion hU
%   of span SSZ in the iSup of the level spans (mem_levels_iff'), span_SSZ_finite and finrank_biSup_le;
%   restBClvl_of is finrank_span_blockData_gens with Ljj_span_fd, finrank_span_image_le, xi_span and
%   finrank_Ljj_span (finrank_eq_span_gen); RestYbc is restYbc_of: restDecomp_proved (the basis bV of the
%   span of the coefficients and the coordinate maps through a left inverse rV), restFunY_proved
%   (cmap_mem_restrictSupport, cmap_ycond with cmap_glAct) and restClear_proved (p^e clears the
%   denominators). The kept condition is not used.

\begin{lemma}\label{pd-kappa-defects-segment-mean-theorem-a-block-bound-dual}
Let an instance be given, $(\alpha,\alpha_+)\in W^{S_j}\times W$, $b$ a block and $D$ an integer, and let
$L^{\mathbf Z}_b$ be the $\mathbf Z_p$-span of the generators over $\mathbf Q_p$ of all levels of $b$
(Definition~\ref{pd-kappa-defects-segment-mean-theorem-a-levels}). Let $n\geq0$, let
$x_1,\dots,x_n,y_1,\dots,y_n\in L^{\mathbf Z}_b$ and let $e$ be an integer such that $[x_i,y_l]=p^e\delta_{il}$
for all $i,l$, with the Fock pairing over $\mathbf Q_p$, and $e\leq v_p(\nu!)+D+D_{\mathrm{lv}}$ for every
monomial $\nu$ with a nonzero coefficient in an element of $L^{\mathbf Z}_b$. Then $L_b$ has a good sublattice
$L_2$ of depth $D+D_{\mathrm{lv}}$ for $\beta_b$ with $n\leq\operatorname{rk}L_2$
(Definitions~\ref{pd-kappa-defects-shares} and~\ref{pd-kappa-defects-segment-mean-theorem-a-classes}).
\end{lemma}
\leanref{Formalism/SharpRestriction/GrowthBounds/MonomialFactorials/BlockBoundOfDualFamilies.lean}{Formalism.MainNotions.G.SymmetricFormBlocks.block_dual_family_bound}
\leanref{Formalism/SharpRestriction/GrowthBounds/MonomialFactorials/StructureSheafH1BaseChange.lean}{Formalism.MainNotions.G.SymmetricFormBlocks.goodSub_of_dual}

\begin{proof}
As the setting is unramified quadratic, $\varpi_0$ is a uniformizer of $\mathcal R$ with $\varpi_0^2=p$, so
$v(\varpi_0)=1/2$ and $v$ takes values in $\frac12\mathbf Z$ on $\mathcal K^\times$. Let $\iota_p$ and
$P\mapsto P^\iota$ be as in Lemma~\ref{pd-kappa-defects-segment-mean-theorem-a-block-bound-base}. Put
$F=\lfloor D+D_{\mathrm{lv}}\rfloor$, $k=F-e-v_p(\operatorname{sp}(b))$ and $\lambda=\varpi_0^k$, so
$v(\lambda)=k/2$.

\emph{Into $L_b$.} Let $z\in L^{\mathbf Z}_b$. Then $\lambda z^\iota\in L_b$. It lies in $V_b$ by
Lemma~\ref{pd-kappa-defects-segment-mean-theorem-a-block-bound-base}(1). The generators over $\mathbf Q_p$ have
coefficients in $\mathbf Z_p$, so $z$ does, and $\iota_p(\mathbf Z_p)\subset\mathcal R$. Let $\nu$ be a monomial at
which $\lambda z^\iota$ has a nonzero coefficient $\gamma$. Then $v(\gamma)\geq k/2$, and $z$ has a nonzero
coefficient at $\nu$, so $e-v_p(\nu!)\leq D+D_{\mathrm{lv}}$; as the left side is an integer, it is at most $F$.
Hence $v(\gamma)+\frac12v_p(\nu!)+\frac12v_p(\operatorname{sp}(b))\geq\frac12(F-e+v_p(\nu!))\geq0$, which is the
condition of Definition~\ref{pd-kappa-defects-segment-lattice}(6).

\emph{The dual family.} Put $x'_i=\lambda x_i^\iota$ and $y'_i=\lambda y_i^\iota$, elements of $L_b$. By
Lemma~\ref{pd-kappa-defects-segment-mean-theorem-a-block-bound-base}(3),
$\beta_b(x'_i,y'_l)=\operatorname{sp}(b)\lambda^2\iota_p(p^e)\delta_{il}=c\,\delta_{il}$ with
$c=\operatorname{sp}(b)\lambda^2p^e\neq0$, and $v(c)=v_p(\operatorname{sp}(b))+k+e=F\leq D+D_{\mathrm{lv}}$.

\emph{The sublattice.} Let $L_2$ be the set of the elements of $L_b$ that lie in the $\mathcal K$-span of the
$x'_i$. If $u\in L_b$ and $ou\in L_2$ with $0\neq o\in\mathcal R$, then $u=o^{-1}(ou)$ lies in that span, so
$u\in L_2$. Let $u=\sum_it_ix'_i$ be primitive in $L_2$, with $t_i\in\mathcal K$. If every $t_i/\varpi_0$ lay in
$\mathcal R$, then $u=\varpi_0u'$ with $u'=\sum_i(t_i/\varpi_0)x'_i\in L_2$, and $\varpi_0$ is not a unit; so some
$t_l/\varpi_0$ is not in $\mathcal R$, that is $v(t_l)<1/2$, and so $v(t_l)\leq0$. As the block pairing is
linear in the first variable, $\beta_b(u,y'_l)=t_lc\neq0$, and $v(t_lc)\leq F\leq D+D_{\mathrm{lv}}$. So $L_2$ is
a good sublattice of depth $D+D_{\mathrm{lv}}$.

\emph{The rank.} The identity $\beta_b(\sum_it_ix'_i,y'_l)=t_lc$ shows that the $x'_i$ are linearly independent
over $\mathcal K$, hence over $\mathcal R$, and they lie in $L_2$. The map $u\mapsto(\beta_b(u,y'_l))_{l\leq n}$
from $L_2$ to $\mathcal K^n$ is $\mathcal R$-linear, and it is injective: if $u=\sum_it_ix'_i$ maps to $0$, then
$t_lc=0$, so $t_l=0$, for every $l$. Any $n+1$ elements of $\mathcal K^n$ are linearly dependent over $\mathcal K$,
hence over $\mathcal R$ after clearing denominators; so any $n+1$ elements of $L_2$ are linearly dependent over
$\mathcal R$. So $n$ is the largest number of linearly independent elements of $L_2$, that is
$\operatorname{rk}L_2=n$; in particular $n\leq\operatorname{rk}L_2$. The upper bound is needed for this: without
a largest number, $\operatorname{rk}L_2$ would be $0$ (Definition~\ref{pd-kappa-defects-shares}).
\end{proof}
% Lean: block_dual_family_bound C0 CE1 CE2 X aTop aPlus b D n x y e hx hy hxy hmon: F, k and lam = w^k
%   (vK_zpow, vK_w with ramE_eq_two); hmemL (map_SSZ, SSZ_int, coeff_of_int, iotaX_int); hform (fock_smul_both,
%   fock_map); hvc (vK_p, valuation_natCast_padicVal, specConst_ne_zero); then goodSub_of_dual with
%   L2 = L meet the K-span of the x i (vK_le_zero_of_not_int for v(t_l) <= 0). Its last paragraph proves
%   rank L2 = n (hlow from the independence of the x i, hup from the injective map phi into Fin n -> K and
%   IsFractionRing.rank_right_eq) and closes with finrank_eq_of_rank_eq; the statements of goodSub_of_dual
%   and block_dual_family_bound record only n <= finrank L2, which is what the book states. The kept
%   condition and the hypotheses of the instance are not used.

\begin{definition}\label{pd-kappa-defects-segment-mean-theorem-a-block-bound-families-tie}
Let an instance be given, $(\alpha,\alpha_+)\in W^{S_j}\times W$, $b$ a block and $\omega=(j,j_+)$ a level of $b$, and use
the notation of Definitions~\ref{pd-kappa-defects-segment-lattice}
and~\ref{pd-kappa-defects-segment-mean-theorem-a-levels}, over $\mathbf Q_p$. Put $\zeta_l=\hat b_{u_l}\dot-o_{u_l}$ for
$1\leq l\leq t$ and $\zeta_+=\alpha_+\dot-o_+$.
\begin{enumerate}
\item For vectors of variables $w=(w_c)_{c\leq m}$ and $e=(e_c)_{c\leq m}$, two further variables $z,z'$ and integers
  $o,n,\zeta\geq0$, let $\hat E_{we}$ be the multiplication by $zz'+\sum_cw_ce_c$, let $F_{we}$ be as in
  Definition~\ref{pd-kappa-defects-segment-mean-theorem-a-levels}(3), and put
  \[
   \pi(w,e,z,z';o,n,\zeta)F=\sum_{i=0}^{o}\frac{(-1)^i}{(i!)^2\binom{o+n+m+\zeta-1}{i}}\,\hat E_{we}^iF_{we}^i\big(z^\zeta F\big).
  \]
\item The \emph{tie projection} of $\omega$ is $\Pi_\omega=\pi_+\circ\pi_1\circ\cdots\circ\pi_t$, where
  $\pi_l=\pi\big(x_{u_l},y_{u_l},z_l,z'_l;\delta'_\omega(x_{u_l}),\delta'_\omega(y_{u_l}),\zeta_l\big)$ and
  $\pi_+=\pi\big(x_+,\xi,z_+,z'_+;\delta'_\omega(x_+),b_+-j_+,\zeta_+\big)$. Put
  $\Theta_\omega F=\prod_l\hat E_l^{j_l}\cdot\hat E_+^{j_+}\cdot\Pi_\omega F$, a $\mathbf Q_p$-linear map.
\item Put $d_\omega=(b_+-j_+)!\prod_\lambda\delta'_\omega(\lambda)!$, the denominator of $\beta_\omega$, and
  \[
   c_\omega=\Big(\prod_lj_l!\cdot j_+!\Big)^2\prod_{l=1}^t\binom{\delta'_\omega(x_{u_l})+\delta'_\omega(y_{u_l})+m+\zeta_l+j_l}{j_l}
   \zeta_l!\cdot\binom{\delta'_\omega(x_+)+b_+-j_++m+\zeta_++j_+}{j_+}\zeta_+!,
  \]
  positive integers.
\item Put $A=(3s+2)a$ and $U=(t+1)\big\lceil(m+1)\big(\log_p(8A+m)+1\big)\big\rceil$.
\end{enumerate}
\end{definition}
\leanref{Formalism/MahlerMeasure/TieVars.lean}{Formalism.GM.M.tieP}
\leanref{Formalism/MahlerMeasure/WithMargin.lean}{Formalism.GM.M.BlockData.Pall}
\leanref{Formalism/MahlerMeasure/WithMargin.lean}{Formalism.GM.M.BlockData.PSZ}
\leanref{Formalism/SharpRestriction/GrowthBounds/MonomialFactorials/GoodCountMem.lean}{Formalism.MainNotions.G.SymmetricFormBlocks.Th}
\leanref{Formalism/SharpRestriction/GrowthBounds/MonomialFactorials/GoodCountMem.lean}{Formalism.MainNotions.G.SymmetricFormBlocks.lvC}
\leanref{Formalism/SharpRestriction/GrowthBounds/MonomialFactorials/BlockBoundUniformiser.lean}{Formalism.MainNotions.G.SymmetricFormBlocks.degFact}
\leanref{Formalism/SharpRestriction/GrowthBounds/MonomialFactorials/DegFactNeZero.lean}{Formalism.MainNotions.G.SymmetricFormBlocks.expU}
\leanref{Formalism/SharpRestriction/GrowthBounds/MonomialFactorials/BlockValuationZpOfLatticeInclusion.lean}{Formalism.MainNotions.G.SymmetricFormBlocks.aBig}

\begin{lemma}\label{pd-kappa-defects-segment-mean-theorem-a-block-bound-families-fock-raising}
Let $k$ be a field of characteristic $0$, $[\ ,\ ]$ the Fock pairing on a polynomial ring over $k$
(Definition~\ref{pd-weight-lattices}(2)), $q\geq1$, and $w_1,\dots,w_q,e_1,\dots,e_q$ distinct variables. Let $E$ be
the multiplication by $\sum_cw_ce_c$, $F=\sum_c\partial^2/\partial w_c\partial e_c$, $D_w=\sum_cw_c\,\partial/\partial w_c$
and $D_e=\sum_ce_c\,\partial/\partial e_c$. For integers $x$ and $i\geq0$ write $x^{\underline i}=x(x-1)\cdots(x-i+1)$.
\begin{enumerate}
\item The pairing is bilinear and symmetric, and $[vf,g]=[f,\partial g/\partial v]$ for every variable $v$ and all
  $f,g$. So $[Ef,g]=[f,Fg]$.
\item Let $g$ be a polynomial with $D_wg=\alpha g$ and $D_eg=\beta g$, where $\alpha,\beta\in k$. Then
  $D_w(Eg)=(\alpha+1)Eg$, $D_e(Eg)=(\beta+1)Eg$, and $F(E^{j+1}g)=E^{j+1}Fg+(j+1)(\alpha+\beta+q+j)E^jg$ for $j\geq0$.
\item Let moreover $\alpha=a$ and $\beta=b$ be integers $\geq0$ and $Fg=0$, and put $N_0=a+b+q-1$. Then
  $F^iE^lg=l^{\underline i}(N_0+l)^{\underline i}E^{l-i}g$ for all $i,l\geq0$, where the right side is $0$ for $i>l$,
  as $l^{\underline i}=0$ then.
\item Let $g$ be as in (3) and $u$ a polynomial with $Fu=0$. Then $[E^iu,E^lg]=0$ for $i\neq l$, and
  $[E^iu,E^ig]=i!\,(N_0+i)^{\underline i}\,[u,g]$.
\end{enumerate}
\end{lemma}
\leanref{Formalism/MahlerMeasure/MahlerMeasureProofs/RemoveS.lean}{Formalism.GM.ProjS.fockK_symm}
\leanref{Formalism/MahlerMeasure/MahlerMeasureProofs/RemoveS.lean}{Formalism.GM.ProjS.fockK_X_mul}
\leanref{Formalism/MahlerMeasure/MahlerMeasureProofs/RemoveS.lean}{Formalism.GM.ProjS.fockK_tieE}
\leanref{Formalism/MahlerMeasure/MahlerMeasureProofs/RemoveS.lean}{Formalism.GM.ProjS.F_E_pow_comm}
\leanref{Formalism/MahlerMeasure/MahlerMeasureProofs/LevelAct.lean}{Formalism.GM.F_E_pow}
\leanref{Formalism/MahlerMeasure/MahlerMeasureProofs/LevelAct.lean}{Formalism.GM.F_pow_E_pow}

\begin{proof}
(1) Bilinearity and symmetry are clear from the definition. For exponent vectors $\kappa,\kappa'$ and the unit
vector $\varepsilon_v$ at $v$, $[vx^\kappa,x^{\kappa'}]$ and $[x^\kappa,\partial x^{\kappa'}/\partial v]$ are both
$(\kappa_v+1)\kappa!$ if $\kappa'=\kappa+\varepsilon_v$, and both $0$ otherwise. Applying this to $e_c$ and to $w_c$ and
summing over $c$ gives $[Ef,g]=[f,Fg]$.

(2) By the product rule, $\partial(Eg)/\partial e_c=w_cg+E\,\partial g/\partial e_c$ and
\[
 \frac{\partial^2(Eg)}{\partial w_c\partial e_c}=g+w_c\frac{\partial g}{\partial w_c}+e_c\frac{\partial g}{\partial e_c}
 +E\frac{\partial^2g}{\partial w_c\partial e_c};
\]
summing over $c$, $F(Eg)=EFg+(\alpha+\beta+q)g$. Likewise $D_w(Eg)=Eg+ED_wg$ and $D_e(Eg)=Eg+ED_eg$. The formula for
$F(E^{j+1}g)$ follows by induction on $j$, applying the case $j=0$ to $E^{j+1}g$, whose eigenvalues are $\alpha+j+1$
and $\beta+j+1$, and using $(j+1)(X+j)+X+2j+2=(j+2)(X+j+1)$ with $X=\alpha+\beta+q$.

(3) By (2), $F(E^{l+1}g)=(l+1)(N_0+l+1)E^lg$. Induction on $i$: $F^{i+1}E^0g=F^iFg=0$, and
$F^{i+1}E^{l+1}g=(l+1)(N_0+l+1)F^iE^lg=(l+1)^{\underline{i+1}}(N_0+l+1)^{\underline{i+1}}E^{l-i}g$.

(4) By (1) and (3), $[E^iu,E^lg]=[u,F^iE^lg]=l^{\underline i}(N_0+l)^{\underline i}[u,E^{l-i}g]$. For $i>l$ this is $0$.
For $i<l$, $[u,E^{l-i}g]=[E^{l-i}g,u]=[E^{l-i-1}g,Fu]=0$. For $i=l$, $l^{\underline l}=l!$.
\end{proof}
% Lean: (1) is fockK_symm, fockK_add_left, fockK_smul_left, fockK_X_mul (mfac_add_single) and fockK_tieE; (2) is
%   F_S_mul_k, eulL_w_S_mul, eulL_e_S_mul, eulL_tieE_pow and F_E_pow_comm (with eigenvalues in k), and F_S_mul,
%   euler_pow, F_E_pow (with natural eigenvalues); (3) is F_pow_E_pow (descFactorial); (4) is the private
%   fockK_tieE_pow_tieE_pow_eq_ite (for i <= l) and the argument of step_orth (F_pow_E_pow with
%   Nat.descFactorial_eq_zero_iff_lt for i > l).

\begin{lemma}\label{pd-kappa-defects-segment-mean-theorem-a-block-bound-families-fock-tie}
Let $k$ be a field of characteristic $0$, $m\geq1$, and $w=(w_c)_{c\leq m}$, $e=(e_c)_{c\leq m}$, $z$ and $z'$ distinct
variables of a polynomial ring over $k$, with the Fock pairing. Let $E$, $F$, $D_w$ and $D_e$ be the operators of
Lemma~\ref{pd-kappa-defects-segment-mean-theorem-a-block-bound-families-fock-raising} for the $m$ pairs
$(w_c,e_c)$, and $\hat E$, $\hat F$, $D_{\hat w}$ and $D_{\hat e}$ those for the $m+1$ pairs $(z,z'),(w_1,e_1),\dots,
(w_m,e_m)$; so $\hat E$ is the multiplication by $zz'+\sum_cw_ce_c$. For integers $o,n,\zeta\geq0$ put
$\pi_{o,n}=\pi(w,e,z,z';o,n,\zeta)$ and $\phi_{o,n}=\phi(w,e;o,n,\zeta)$
(Definitions~\ref{pd-kappa-defects-segment-mean-theorem-a-block-bound-families-tie}(1)
and~\ref{pd-kappa-defects-segment-mean-theorem-a-levels}(3)). A polynomial $y$ is \emph{pure of bidegree} $(o,n)$
if $D_wy=oy$, $D_ey=ny$ and $\partial y/\partial z=\partial y/\partial z'=0$. Let $y$ and $y'$ be pure of bidegree
$(o,n)$.
\begin{enumerate}
\item $\hat F\pi_{o,n}(y)=0$, $D_{\hat w}\pi_{o,n}(y)=(o+\zeta)\pi_{o,n}(y)$ and $D_{\hat e}\pi_{o,n}(y)=n\pi_{o,n}(y)$.
\item $[\pi_{o,n}(y),\pi_{o,n}(y')]=\zeta!\,[y,\phi_{o,n}(y')]$.
\item $[\hat E^j\pi_{o,n}(y),\hat E^j\pi_{o,n}(y')]=(j!)^2\binom{o+n+m+\zeta+j}{j}\zeta!\,[y,\phi_{o,n}(y')]$ for $j\geq0$.
\item If $y''$ is pure of bidegree $(o',n')$ and $j\neq j'$, then $[\hat E^j\pi_{o,n}(y),\hat E^{j'}\pi_{o',n'}(y'')]=0$.
\end{enumerate}
\end{lemma}
\leanref{Formalism/MahlerMeasure/MahlerMeasureProofs/RemoveS.lean}{Formalism.GM.ProjS.fh_P}
\leanref{Formalism/MahlerMeasure/MahlerMeasureProofs/RemoveS.lean}{Formalism.GM.ProjS.eig_P}
\leanref{Formalism/MahlerMeasure/MahlerMeasureProofs/RemoveS.lean}{Formalism.GM.ProjS.fockForm_augLift_augLift}
\leanref{Formalism/MahlerMeasure/MahlerMeasureProofs/RemoveS.lean}{Formalism.GM.ProjS.fockForm_raiseAug_pow_augLift}
\leanref{Formalism/MahlerMeasure/MahlerMeasureProofs/FockBlocks.lean}{Formalism.GM.FockB.step_same}
\leanref{Formalism/MahlerMeasure/MahlerMeasureProofs/FockBlocks.lean}{Formalism.GM.FockB.step_orth}

\begin{proof}
Put $A=o+n+m+\zeta-1$, so $A\geq o$ as $m\geq1$; put $c_i=(-1)^i/((i!)^2\binom Ai)$ for $0\leq i\leq o$ and
$u_i=F^i(z^\zeta y)$ for $i\geq0$, so that $\pi_{o,n}(y)=\sum_{i=0}^oc_i\hat E^iu_i$.

\emph{The polynomials $u_i$.} The operator $F$ involves neither $z$ nor $z'$, and each of its terms lowers the
degree in the $w_c$ and the degree in the $e_c$ by one. So $D_wu_i=(o-i)u_i$, $D_eu_i=(n-i)u_i$,
$D_{\hat w}u_i=(o+\zeta-i)u_i$, $D_{\hat e}u_i=(n-i)u_i$ and $\partial u_i/\partial z'=0$. Hence
$\hat Fu_i=\partial^2u_i/\partial z\partial z'+Fu_i=u_{i+1}$, and $\hat F^i(z^\zeta y)=u_i$. As the eigenvalues of $D_w$
on monomials are $\geq0$ and $D_wu_{o+1}=-u_{o+1}$, $u_{o+1}=0$.

(1) By Lemma~\ref{pd-kappa-defects-segment-mean-theorem-a-block-bound-families-fock-raising}(2) for the $m+1$
pairs, $D_{\hat w}\hat E^iu_i=(o+\zeta)\hat E^iu_i$ and $D_{\hat e}\hat E^iu_i=n\hat E^iu_i$, which gives the second and
third claims. For the first, $\hat Fu_0=u_1$, and for $1\leq i\leq o$ the same lemma, with $q=m+1$, $\alpha=o+\zeta-i$
and $\beta=n-i$, gives $\hat F(\hat E^iu_i)=\hat E^iu_{i+1}+i(A+1-i)\hat E^{i-1}u_i$. So
\[
 \hat F\pi_{o,n}(y)=\sum_{i=0}^{o-1}\big(c_i+(i+1)(A-i)c_{i+1}\big)\hat E^iu_{i+1}+c_o\hat E^ou_{o+1}.
\]
Here $u_{o+1}=0$, and $(i+1)(A-i)c_{i+1}=-c_i$, as $(i+1)\binom A{i+1}=(A-i)\binom Ai$ with $A-i>0$ for $i<o$.

(2) Let $Q$ be a polynomial with $\hat FQ=0$. For $i\geq1$,
$[\hat E^iu_i,Q]=[u_i,\hat F^iQ]=0$ by Lemma~\ref{pd-kappa-defects-segment-mean-theorem-a-block-bound-families-fock-raising}(1);
so $[\pi_{o,n}(y),Q]=[z^\zeta y,Q]$. By (1) this applies to $Q=\pi_{o,n}(y')$. Write $u'_i=F^i(z^\zeta y')$. By the same
lemma, for both sets of pairs, and by $\hat F^i(z^\zeta y)=F^i(z^\zeta y)$,
\[
 [z^\zeta y,\hat E^iu'_i]=[\hat F^i(z^\zeta y),u'_i]=[F^i(z^\zeta y),u'_i]=[z^\zeta y,E^iF^i(z^\zeta y')];
\]
summing with the $c_i$, $[z^\zeta y,\pi_{o,n}(y')]=[z^\zeta y,\phi_{o,n}(z^\zeta y')]$. As $E$ and $F$ involve no $z$,
$\phi_{o,n}(z^\zeta y')=z^\zeta\phi_{o,n}(y')$ and $\partial\phi_{o,n}(y')/\partial z=0$. Applying
Lemma~\ref{pd-kappa-defects-segment-mean-theorem-a-block-bound-families-fock-raising}(1) $\zeta$ times with $v=z$,
$[z^\zeta y,z^\zeta\phi_{o,n}(y')]=[y,\partial^\zeta(z^\zeta\phi_{o,n}(y'))/\partial z^\zeta]=\zeta!\,[y,\phi_{o,n}(y')]$.

(3) By (1),
Lemma~\ref{pd-kappa-defects-segment-mean-theorem-a-block-bound-families-fock-raising}(4) applies to $u=\pi_{o,n}(y)$
and $g=\pi_{o,n}(y')$, with $q=m+1$, $a=o+\zeta$, $b=n$ and $N_0=o+n+m+\zeta$. It gives
$[\hat E^j\pi_{o,n}(y),\hat E^j\pi_{o,n}(y')]=j!\,(N_0+j)^{\underline j}[\pi_{o,n}(y),\pi_{o,n}(y')]$, and
$(N_0+j)^{\underline j}=j!\binom{N_0+j}j$. Now use (2).

(4) By (1), $g=\pi_{o',n'}(y'')$ satisfies the hypotheses of
Lemma~\ref{pd-kappa-defects-segment-mean-theorem-a-block-bound-families-fock-raising}(3) with $q=m+1$, $a=o'+\zeta$
and $b=n'$, and $\hat F\pi_{o,n}(y)=0$. Part (4) of that lemma gives the claim.
\end{proof}
% Lean: the polynomials u_i are eig_u, eig_v (eulL_tieF_pow), Fh_v, Fh_pow_u and v_top (eq_zero_of_eulL). (1) is
%   eig_P (eulL_tieE_pow) and fh_P (F_E_pow_comm, coef_cancel, and the telescoping
%   sum_range_smul_add_smul_eq_zero_of_recurrence); (2) is fockK_P_left and fockForm_augLift_augLift (step1 with
%   fockK_tieE_pow_right and Fh_pow_u; hPX, hPd, fockK_X_pow_mul and pderiv_pow_X_pow_mul); (3) is
%   fockForm_raiseAug_pow_augLift with step_same (Nat.descFactorial_eq_factorial_mul_choose); (4) is step_orth
%   (fh_P, eig_P, F_pow_E_pow, fockK_tieE_pow and fockK_tieE_pow_right). The Lean's TieVars is the
%   distinctness of the variables, and its ZSFree and BiHom make up purity (biHom_of_good, zsfree_of_good).

\begin{lemma}\label{pd-kappa-defects-segment-mean-theorem-a-block-bound-families-fock}
Let an instance be given, $(\alpha,\alpha_+)\in W^{S_j}\times W$, $b$ a block and $\omega,\omega'$ levels of $b$
(Definitions~\ref{pd-kappa-defects-segment-mean-theorem-a-levels}
and~\ref{pd-kappa-defects-segment-mean-theorem-a-block-bound-families-tie}). If $\omega\neq\omega'$, then
$[\Theta_\omega f,\Theta_{\omega'}f']=0$ for all $f\in\Lambda_\omega$ and $f'\in\Lambda_{\omega'}$. For all
$f,f'\in\Lambda_\omega$, $[\Theta_\omega f,\Theta_\omega f']=c_\omega[f,\Phi_\omega f']$.
\end{lemma}
\leanref{Formalism/SharpRestriction/GrowthBounds/MonomialFactorials/GoodCountMem.lean}{Formalism.MainNotions.G.SymmetricFormBlocks.Th_fock}
\leanref{Formalism/SharpRestriction/GrowthBounds/MonomialFactorials/GoodCountMem.lean}{Formalism.MainNotions.G.SymmetricFormBlocks.Th_orth}
\leanref{Formalism/MahlerMeasure/LevelLattice.lean}{Formalism.GM.M.FockBlocks}
\leanref{Formalism/MahlerMeasure/MahlerMeasureProofs/FockBlocks.lean}{Formalism.GM.fock_block_decomposition}

\begin{proof}
Both sides of each identity are additive in $f$ and in $f'$ and $\mathbf Z_p$-homogeneous, as $\Theta_\omega$,
$\Theta_{\omega'}$ and $\Phi_\omega$ are linear and the Fock pairing is bilinear. So it suffices to take
$f=yg$ and $f'=y'g'$, with $y\in Y^{\mathbf Z}_\omega$, $y'$ in $Y^{\mathbf Z}_\omega$ or $Y^{\mathbf Z}_{\omega'}$, and
$g,g'$ homogeneous polynomials in the $\xi_c$ with coefficients in $\mathbf Z_p$ of the degrees of the levels.

\emph{The ties.} Index the ties by $a\in\{+,1,\dots,t\}$, with the letters $w_l=x_{u_l}$, $e_l=y_{u_l}$, $w_+=x_+$
and $e_+=\xi$, and the variables $z_a,z'_a$. For a level $\omega=(j,j_+)$ write $j_a$ for $j_l$ or $j_+$, and put
$\delta(\xi)=b_+$ and $\delta'_\omega(\xi)=b_+-j_+$. Then $\delta'_\omega(w_a)=\delta(w_a)-j_a$ and
$\delta'_\omega(e_a)=\delta(e_a)-j_a$, as in the proof of
Lemma~\ref{pd-kappa-defects-segment-mean-theorem-a-block-bound-families-valuation-degrees}; so these numbers depend
on $\omega$ only through $j_a$. Let $\pi^\omega_a$ and $\phi^\omega_a$ be the operators $\pi_{o,n}$ and $\phi_{o,n}$ of
Lemma~\ref{pd-kappa-defects-segment-mean-theorem-a-block-bound-families-fock-tie} for the variables of the tie
$a$, with $(o,n)=(\delta'_\omega(w_a),\delta'_\omega(e_a))$ and $\zeta=\zeta_a$, and let $\hat E_a$ also denote the
polynomial $z_az'_a+\sum_cw_{a,c}e_{a,c}$. By
Definitions~\ref{pd-kappa-defects-segment-mean-theorem-a-levels}(3)
and~\ref{pd-kappa-defects-segment-mean-theorem-a-block-bound-families-tie}(2),
$\Pi_\omega=\pi^\omega_+\circ\pi^\omega_1\circ\cdots\circ\pi^\omega_t$, $\Phi_\omega=\phi^\omega_+\circ\phi^\omega_1\circ\cdots
\circ\phi^\omega_t$ and $\Theta_\omega f=\prod_a\hat E_a^{j_a}\cdot\Pi_\omega f$, and $c_\omega$ is the product over the ties
of $c^\omega_a=(j_a!)^2\binom{\delta'_\omega(w_a)+\delta'_\omega(e_a)+m+\zeta_a+j_a}{j_a}\zeta_a!$.

\emph{Commuting operators.} The sets $V_a$ of the variables of $w_a$, $e_a$, $z_a$ and $z'_a$ are pairwise
disjoint. Let $\mathcal A_a$ be the algebra of operators generated by the multiplications by the variables in $V_a$
and the derivatives in them. Multiplications by and derivatives in distinct variables commute, so elements of
$\mathcal A_a$ and of $\mathcal A_{a'}$ commute for $a\neq a'$. The operators $\pi^\omega_a$, $\phi^\omega_a$, the
multiplication by $\hat E_a$, and $D_{w_a}$, $D_{e_a}$, $\partial/\partial z_a$ and $\partial/\partial z'_a$ lie in
$\mathcal A_a$. So an element of $\mathcal A_{a'}$, $a'\neq a$, maps polynomials pure of bidegree $(o,n)$ for the
tie $a$ to polynomials pure of bidegree $(o,n)$ for the tie $a$. The polynomial $yg$ is pure of bidegree
$(\delta'_\omega(w_a),\delta'_\omega(e_a))$ for every tie $a$: $y$ has total degree $\delta'_\omega(\lambda)$ in the
variables of each letter $\lambda$, $g$ has degree $b_+-j_+$ in the $\xi_c$, and neither involves the $z_a,z'_a$. The
same holds for $y'g'$ and its level.

\emph{Induction over the ties.} Order the ties as $a_0=+$, $a_1=1$, \dots, $a_t=t$. For a level $\omega$ and
$0\leq i\leq t+1$ let $\Pi^\omega_i=\pi^\omega_{a_i}\circ\cdots\circ\pi^\omega_{a_t}$, $\Phi^\omega_i=\phi^\omega_{a_i}\circ
\cdots\circ\phi^\omega_{a_t}$, $\Theta^\omega_i=\prod_{i'\geq i}\hat E_{a_{i'}}^{j_{a_{i'}}}\cdot\Pi^\omega_i$ and
$C^\omega_i=\prod_{i'\geq i}c^\omega_{a_{i'}}$, so that $\Theta^\omega_0=\Theta_\omega$, $\Phi^\omega_0=\Phi_\omega$ and
$C^\omega_0=c_\omega$. Let $\omega'=(j',j'_+)$ be a level. We show, by descending induction on $i$, that for all
polynomials $Y$ and $Y'$ that are pure of bidegree $(\delta'_\omega(w_a),\delta'_\omega(e_a))$, respectively
$(\delta'_{\omega'}(w_a),\delta'_{\omega'}(e_a))$, for every tie $a=a_{i'}$ with $i'\geq i$:
\begin{enumerate}
\item[(i)] if $\omega'=\omega$, then $[\Theta^\omega_iY,\Theta^\omega_iY']=C^\omega_i[Y,\Phi^\omega_iY']$;
\item[(ii)] if $j_{a_{i'}}\neq j'_{a_{i'}}$ for some $i'\geq i$, then $[\Theta^\omega_iY,\Theta^{\omega'}_iY']=0$.
\end{enumerate}
For $i=t+1$, (i) holds and (ii) is vacuous. Let $i\leq t$, $a=a_i$, $Z=\Theta^\omega_{i+1}Y$ and
$Z'=\Theta^{\omega'}_{i+1}Y'$. The operators $\Theta^\omega_{i+1}$ and $\Theta^{\omega'}_{i+1}$ are composites of elements
of the $\mathcal A_{a_{i'}}$, $i'>i$. So $\pi^\omega_a$ commutes with the multiplication by
$\prod_{i'>i}\hat E_{a_{i'}}^{j_{a_{i'}}}$, whence $\Theta^\omega_iY=\hat E_a^{j_a}\pi^\omega_a(Z)$ and likewise
$\Theta^{\omega'}_iY'=\hat E_a^{j'_a}\pi^{\omega'}_a(Z')$; and $Z$ and $Z'$ are pure of the bidegrees of $Y$ and $Y'$ for
the tie $a$. If $j_a\neq j'_a$,
Lemma~\ref{pd-kappa-defects-segment-mean-theorem-a-block-bound-families-fock-tie}(4) gives
$[\Theta^\omega_iY,\Theta^{\omega'}_iY']=0$. If $j_a=j'_a$, then $\pi^{\omega'}_a=\pi^\omega_a$ and
$\phi^{\omega'}_a=\phi^\omega_a$, and part (3) of that lemma gives
\[
 [\Theta^\omega_iY,\Theta^{\omega'}_iY']=c^\omega_a[Z,\phi^\omega_aZ']=c^\omega_a[\Theta^\omega_{i+1}Y,\Theta^{\omega'}_{i+1}(\phi^\omega_aY')],
\]
as $\phi^\omega_a\in\mathcal A_a$ commutes with $\Theta^{\omega'}_{i+1}$. The polynomial $\phi^\omega_aY'$ is pure of the
bidegrees of $Y'$ for the ties $a_{i'}$, $i'>i$. In case (ii) some $i'>i$ has $j_{a_{i'}}\neq j'_{a_{i'}}$, and the
induction gives $0$. In case (i) it gives $c^\omega_aC^\omega_{i+1}[Y,\Phi^\omega_{i+1}\phi^\omega_aY']=C^\omega_i[Y,\Phi^\omega_iY']$,
as $\phi^\omega_a$ commutes with $\Phi^\omega_{i+1}$.

\emph{Conclusion.} Take $i=0$, $Y=yg$ and $Y'=y'g'$. If $\omega'\neq\omega$, then $j_a\neq j'_a$ for some tie $a$, and
(ii) gives the first identity; (i) gives the second.
\end{proof}
% Lean: Th_fock and Th_orth reduce by Submodule.span_induction to the generators and apply
%   (fock_block_decomposition p N r G hB hr).2 and .1, which prove FockBlocks for every block datum with
%   BlockOK and r >= 1 (gB_blockOK). fock_block_decomposition = fock_blocks_of_weights: the ties are the list
%   LA D = none :: some i (wA, eA, zA, sA, zzA, oA, nA, jA), with pall_eq, phi_eq and ehat_eq; deg_fst, deg_snd
%   and degU give the lowered degrees; tvA and disjA (pairwiseA) the distinct and disjoint variables; goodA the
%   purity of y g. The algebras A_a are Ops (Va ...), with ops_commute, good_of_ops, exists_op_eq_chainP,
%   chainE_mem and chainPhi_mem; the induction is chain_main (case (i), with chain_rw, good_chain, phi_chain and
%   step_same) and chain_orth (case (ii), with step_orth). FockBlocks states the same-level identity for Theta
%   divided by prod j_i! jp!, so its constant omits the factor (prod j_i! jp!)^2, which chain_main's cst carries
%   and fock_blocks_of_weights and Th_fock clear and restore with field_simp. FockBlocks is stated for y in Y
%   over Q_p, which contains Y^Z.

\begin{lemma}\label{pd-kappa-defects-segment-mean-theorem-a-block-bound-families-lattice-lowering-weights}
Let $p$ be a prime and $r\geq1$, $s\geq0$ integers, and consider polynomials over $\mathbf Q_p$ in the variables of the
letters $x_1,y_1,\dots,x_s,y_s,x_+$ of Definition~\ref{pd-kappa-defects-segment-lattice}(3), with $r$ variables each.
Let $\mathcal L$ be an initial segment of the letters and $M$ an integer, let $E^{\mathcal L}_{cd}$, the weight for
$\mathcal L$ and $Y(\mathcal L,M)$ be as in Definition~\ref{pd-kappa-defects-segment-lattice}(4), over $\mathbf Q_p$, and
let $P_M$ be the span of the monomials of weight at most $M$. For a letter $\mu$ and $1\leq e\leq r$ write
$\partial_{\mu,e}=\partial/\partial\mu_e$.
\begin{enumerate}
\item If a $\mathbf Q_p$-linear map $T$ of polynomials commutes with every $E^{\mathcal L}_{cd}$ and maps $P_M$ into
  itself, then $T$ maps $Y(\mathcal L,M)$ into itself.
\item For a letter $\mu$ and $1\leq c,d,e\leq r$, $\partial_{\mu,e}E^{\mathcal L}_{cd}=E^{\mathcal L}_{cd}\partial_{\mu,e}+\epsilon$,
  where $\epsilon=0$ if $\mu\notin\mathcal L$, $\epsilon=\delta_{ec}\partial_{\mu,d}$ if $\mu\in\mathcal L$ is a vector
  letter, and $\epsilon=-\delta_{ed}\partial_{\mu,c}$ if $\mu\in\mathcal L$ is a covector letter.
\item If $\mu$ is a vector letter or does not lie in $\mathcal L$, then $\partial_{\mu,e}$ maps $Y(\mathcal L,M)$ into
  itself for every $e$.
\item If $\lambda\in\mathcal L$ is a vector letter and $\lambda'\in\mathcal L$ a covector letter, then
  $F=\sum_e\partial_{\lambda,e}\partial_{\lambda',e}$ commutes with every $E^{\mathcal L}_{cd}$ and maps $Y(\mathcal L,M)$
  into itself.
\end{enumerate}
\end{lemma}
\leanref{Formalism/PolynomialDimension/MultinomialDepth/PartnerM.lean}{Formalism.CountingSumsFrames.ycond_keep}
\leanref{Formalism/SharpRestriction/GrowthBounds/MonomialFactorials/DegSubSingle.lean}{Formalism.MainNotions.G.SymmetricFormBlocks.pderiv_glAct}
\leanref{Formalism/SharpRestriction/GrowthBounds/MonomialFactorials/DegSubSingle.lean}{Formalism.MainNotions.G.SymmetricFormBlocks.wt_sub_single}
\leanref{Formalism/SharpRestriction/GrowthBounds/MonomialFactorials/DegSubSingle.lean}{Formalism.MainNotions.G.SymmetricFormBlocks.pderiv_ycond}
\leanref{Formalism/SharpRestriction/GrowthBounds/MonomialFactorials/DegSubSingle.lean}{Formalism.MainNotions.G.SymmetricFormBlocks.tieF_ycond}

\begin{proof}
By definition $Y(\mathcal L,M)$ lies in $P_M$ and is stable under the $E^{\mathcal L}_{cd}$, and it contains every
subspace of $P_M$ with this property.

(1) The subspace $T(Y(\mathcal L,M))$ lies in $P_M$, and it is stable under the $E^{\mathcal L}_{cd}$, as
$E^{\mathcal L}_{cd}Ty=TE^{\mathcal L}_{cd}y$ and $E^{\mathcal L}_{cd}y\in Y(\mathcal L,M)$. So it lies in $Y(\mathcal L,M)$.

(2) As partial derivatives commute, for a vector letter $\lambda$ and a polynomial $f$,
$\partial_{\mu,e}(\lambda_c\partial_{\lambda,d}f)=\delta_{\lambda\mu}\delta_{ec}\partial_{\lambda,d}f+\lambda_c\partial_{\lambda,d}\partial_{\mu,e}f$,
and for a covector letter $\lambda'$,
$\partial_{\mu,e}(\lambda'_d\partial_{\lambda',c}f)=\delta_{\lambda'\mu}\delta_{ed}\partial_{\lambda',c}f+\lambda'_d\partial_{\lambda',c}\partial_{\mu,e}f$.
Summing over the letters of $\mathcal L$ with the signs of $E^{\mathcal L}_{cd}$ gives the formula.

(3) A derivative $\partial_{\mu,e}$ maps a monomial to a multiple of a monomial whose weight is lowered by $1$ if
$\mu\in\mathcal L$ is a vector letter and $e=1$, raised by $1$ if $\mu\in\mathcal L$ is a covector letter and $e=1$, and
unchanged otherwise. Let $S'$ be the span of the $\partial_{\mu,e'}y$ with $1\leq e'\leq r$ and $y\in Y(\mathcal L,M)$. By
the hypothesis on $\mu$, the $\partial_{\mu,e'}$ do not raise weights, so $S'\subset P_M$. By (2),
$E^{\mathcal L}_{cd}\partial_{\mu,e'}y=\partial_{\mu,e'}(E^{\mathcal L}_{cd}y)-\epsilon y$, where $\epsilon y$ is $0$ or
$\delta_{e'c}\partial_{\mu,d}y$; both terms lie in $S'$, as $E^{\mathcal L}_{cd}y\in Y(\mathcal L,M)$. So $S'$ is stable under
the $E^{\mathcal L}_{cd}$, and $S'\subset Y(\mathcal L,M)$.

(4) By (2) for $\lambda'$ and then for $\lambda$,
\[
 \partial_{\lambda,e}\partial_{\lambda',e}E^{\mathcal L}_{cd}=E^{\mathcal L}_{cd}\partial_{\lambda,e}\partial_{\lambda',e}
 +\delta_{ec}\partial_{\lambda,d}\partial_{\lambda',e}-\delta_{ed}\partial_{\lambda,e}\partial_{\lambda',c}.
\]
Summed over $e$, the last two terms give $\partial_{\lambda,d}\partial_{\lambda',c}-\partial_{\lambda,d}\partial_{\lambda',c}=0$;
so $F$ commutes with $E^{\mathcal L}_{cd}$. By the weight changes in (3), for $e=1$ the derivative $\partial_{\lambda',1}$
raises the weight by $1$ and $\partial_{\lambda,1}$ lowers it by $1$, and for $e\neq1$ neither changes it; so $F$ maps
$P_M$ into itself. Now apply (1).
\end{proof}
% Lean: (1) is CountingSumsFrames.ycond_keep (the image of YCond under T is in the sSup), with ycond_le_supp
%   and ycond_glStable; (2) is pderiv_glAct (glAct is E^L_cd, the variable (i, c) is the book's i-th letter's c-th
%   variable); the weight change is wt_sub_single, with c = 0 the book's first variable; (3) is pderiv_ycond
%   (hypothesis: a < J -> LL.vec a = true), through the span S' and pderiv_keep; (4) is tieF_ycond, through
%   ycond_keep, with the commutation computed by pderiv_glAct and the weight by pderiv_keep twice.

\begin{definition}\label{pd-kappa-defects-segment-mean-theorem-a-block-bound-families-lattice-lowering-pairs}
Let an instance be given, $(\alpha,\alpha_+)\in W^{S_j}\times W$ and $b$ a block, with the degrees $\delta$ of
Definition~\ref{pd-kappa-defects-segment-lattice}(3) and the $u_l$ of
Definition~\ref{pd-kappa-defects-segment-mean-theorem-a-levels}, and let $j_1,\dots,j_t,j_+\geq0$ be integers. Put
$\delta'_{(j,j_+)}(x_{u_l})=\delta(x_{u_l})\dot-j_l$ and $\delta'_{(j,j_+)}(y_{u_l})=\delta(y_{u_l})\dot-j_l$ for
$1\leq l\leq t$, $\delta'_{(j,j_+)}(x_+)=\delta(x_+)\dot-j_+$, and $\delta'_{(j,j_+)}(\lambda)=\delta(\lambda)$ for the other
letters $\lambda$. Let $Y_{(j,j_+)}$, $Y^{\mathbf Z}_{(j,j_+)}$ and $\Lambda_{(j,j_+)}$ be formed as $Y_\omega$,
$Y^{\mathbf Z}_\omega$ and $\Lambda_\omega$ in Definition~\ref{pd-kappa-defects-segment-mean-theorem-a-levels}(2), with
$\delta'_{(j,j_+)}$ in place of $\delta'_\omega$ and $b_+\dot-j_+$ in place of $b_+-j_+$. If $(j,j_+)$ is a level $\omega$ of
$b$, they are $Y_\omega$, $Y^{\mathbf Z}_\omega$ and $\Lambda_\omega$.
\end{definition}
\leanref{Formalism/MahlerMeasure/WithMargin.lean}{Formalism.GM.M.BlockData.lettersAt}
\leanref{Formalism/ThetaLogCoordinates/BihomogeneousPairs/LetterChainEmbeddings/ChainEmbInj.lean}{Formalism.InvariantFactorPaths.KP.LL.Lat}

\begin{lemma}\label{pd-kappa-defects-segment-mean-theorem-a-block-bound-families-lattice-lowering-spaces}
Let an instance be given, $(\alpha,\alpha_+)\in W^{S_j}\times W$, $b$ a block and $j_1,\dots,j_t,j_+\geq0$ integers, with
$Y_{(j,j_+)}$ as in Definition~\ref{pd-kappa-defects-segment-mean-theorem-a-block-bound-families-lattice-lowering-pairs}.
Let $1\leq l\leq t$ and $1\leq c\leq m$, and let $j^{(l)}$ be $j$ with $j_l$ raised by one.
\begin{enumerate}
\item The operator $F_l=\sum_{e=1}^m\partial^2/\partial x_{u_l,e}\partial y_{u_l,e}$ maps $Y_{(j,j_+)}$ into
  $Y_{(j^{(l)},j_+)}$.
\item The derivative $\partial/\partial x_{+,c}$ maps $Y_{(j,j_+)}$ into $Y_{(j,j_++1)}$.
\end{enumerate}
\end{lemma}
\leanref{Formalism/ThetaLogCoordinates/BihomogeneousPairs/LetterChainEmbeddings/PairStep.lean}{Formalism.InvariantFactorPaths.KP.LL.HypJoin}
\leanref{Formalism/ThetaLogCoordinates/BihomogeneousPairs/LetterChainEmbeddings/PairStep.lean}{Formalism.InvariantFactorPaths.KP.LL.HypU}
\leanref{Formalism/SharpRestriction/GrowthBounds/MonomialFactorials/BlockValuationZpOfValuationBound.lean}{Formalism.MainNotions.G.SymmetricFormBlocks.hypTie_proved}
\leanref{Formalism/SharpRestriction/GrowthBounds/MonomialFactorials/BlockValuationZpOfValuationBound.lean}{Formalism.MainNotions.G.SymmetricFormBlocks.hypU_proved}

\begin{proof}
\emph{Degrees.} Let $f\in Y_{(j,j_+)}$. Each monomial of $f$ has total degree $\delta'_{(j,j_+)}(\lambda)$ in the
variables of each letter $\lambda$. A derivative in a variable of a letter $\lambda$ maps a monomial to $0$ or to a
multiple of a monomial whose degree in $\lambda$ is one less, with the degrees in the other letters unchanged; it
gives $0$ if the degree in $\lambda$ is $0$. As $(d\dot-k)-1=d\dot-(k+1)$ when $d\dot-k\geq1$, the monomials of $F_lf$
have the degrees $\delta'_{(j^{(l)},j_+)}$, and those of $\partial f/\partial x_{+,c}$ the degrees $\delta'_{(j,j_++1)}$.

\emph{Initial segments.} The spaces $Y(\mathcal L_{l'},M_{l'})$ and $Y(\mathcal L_+,M_+)$ in the definition of
$Y_{(j,j_+)}$ do not depend on $(j,j_+)$. Let $\mathcal L$ be one of these segments, $M$ its integer and
$f\in Y(\mathcal L,M)$. For (2): $x_+$ is a vector letter, so $\partial f/\partial x_{+,c}\in Y(\mathcal L,M)$ by
Lemma~\ref{pd-kappa-defects-segment-mean-theorem-a-block-bound-families-lattice-lowering-weights}(3). For (1): $x_{u_l}$
is a vector letter, and $y_{u_l}$, which follows it, a covector letter. If $y_{u_l}\in\mathcal L$, then
$x_{u_l}\in\mathcal L$, and $F_lf\in Y(\mathcal L,M)$ by part~(4) of that lemma. Otherwise each $\partial/\partial y_{u_l,e}$
maps $Y(\mathcal L,M)$ into itself by part~(3), and so does each $\partial/\partial x_{u_l,e}$; so $F_lf\in Y(\mathcal L,M)$.
\end{proof}
% Lean: hypTie_proved (HypJoin) and hypU_proved (HypU), for every block datum with BlockOK. The degrees are hpol
%   (Pol p r LL into Pol p r LL' for LL' lowered by one at one letter, through pderiv_monomial) with hdeg and
%   sum_update_ite; ycond_vec says that YCond depends only on which letters are vector letters. In hypTie_proved
%   the case E < J k is tieF_ycond, and the other case is pderiv_ycond at E (not in the segment) and at W (a
%   vector letter, by the admissibility hB.1); hypU_proved is pderiv_ycond at uPos, a vector letter (hB.2.2.1).

\begin{lemma}\label{pd-kappa-defects-segment-mean-theorem-a-block-bound-families-lattice-lowering-sum}
Let $p$ be a prime, $V$ a $\mathbf Q_p$-vector space, $S$ a finite set and $(D_c)_{c\in S}$ pairwise commuting
$\mathbf Q_p$-linear maps $V\to V$. Let $(\Lambda(l))_{l\geq0}$ be $\mathbf Z_p$-submodules of $V$ such that
$(k!)^{-2}D_c^kx\in\Lambda(l+k)$ for all $c\in S$, $k,l\geq0$ and $x\in\Lambda(l)$. Let $L,i\geq0$ be integers with
$i<p^{L+1}$. Then $p^{|S|L}(i!)^{-2}\big(\sum_{c\in S}D_c\big)^ix\in\Lambda(l+i)$ for all $l\geq0$ and $x\in\Lambda(l)$.
\end{lemma}
\leanref{Formalism/ThetaLogCoordinates/BihomogeneousPairs/LetterChainEmbeddings/PlusOp.lean}{Formalism.InvariantFactorPaths.KP.LL.po_Fsum}

\begin{proof}
By induction on $|S|$, for all $i<p^{L+1}$ at once. If $S$ is empty, $(\sum_{c\in S}D_c)^i$ is the identity for
$i=0$ and $0$ for $i\geq1$. Otherwise let $S=S'\cup\{a\}$ with $a\notin S'$, and put $D'=\sum_{c\in S'}D_c$, which
commutes with $D_a$. By the binomial theorem and
$\binom ik(i!)^{-2}=\binom ik^{-1}(k!)^{-2}((i-k)!)^{-2}$,
\[
 p^{|S|L}(i!)^{-2}(D_a+D')^ix=\sum_{k=0}^ip^L\binom ik^{-1}\,(k!)^{-2}D_a^k\Big(p^{|S'|L}((i-k)!)^{-2}D'^{\,i-k}x\Big).
\]
By induction the element in brackets lies in $\Lambda(l+i-k)$, so by hypothesis its image under $(k!)^{-2}D_a^k$ lies in
$\Lambda(l+i)$. As $k\leq i<p^{L+1}$, $v_p\binom ik\leq L$ by
Lemma~\ref{pd-kappa-defects-segment-mean-theorem-a-block-bound-families-valuation-factorials}(1), so
$p^L\binom ik^{-1}\in\mathbf Z_p$, and each term lies in $\Lambda(l+i)$.
\end{proof}
% Lean: po_Fsum (Finset induction, Commute.add_pow, the identity key, and smul_mem_of_norm with
%   norm_choose_inv_le for |p^L / binom(i, k)| <= 1). It is stated for all j <= n with Nat.log p n <= L.

\begin{lemma}\label{pd-kappa-defects-segment-mean-theorem-a-block-bound-families-lattice-lowering}
Let an instance be given, $(\alpha,\alpha_+)\in W^{S_j}\times W$ and $b$ a block, and use the ties $a$, their letters
$w_a,e_a$ and the notation of the proof of
Lemma~\ref{pd-kappa-defects-segment-mean-theorem-a-block-bound-families-fock}; let
$F_a=\sum_c\partial^2/\partial w_{a,c}\partial e_{a,c}$. Let $j_1,\dots,j_t,j_+\geq0$ be integers, $a$ a tie, $i,L\geq0$
integers with $i<p^{L+1}$, and $(j',j'_+)$ the pair $(j,j_+)$ with $j_a$ raised by $i$. Then
$p^{mL}(i!)^{-2}F_a^if\in\Lambda_{(j',j'_+)}$ for every $f\in\Lambda_{(j,j_+)}$
(Definition~\ref{pd-kappa-defects-segment-mean-theorem-a-block-bound-families-lattice-lowering-pairs}).
\end{lemma}
\leanref{Formalism/ThetaLogCoordinates/BihomogeneousPairs/LetterChainEmbeddings/PairStep.lean}{Formalism.InvariantFactorPaths.KP.LL.LatTie}
\leanref{Formalism/ThetaLogCoordinates/BihomogeneousPairs/LetterChainEmbeddings/PairStep.lean}{Formalism.InvariantFactorPaths.KP.LL.LatPlus}
\leanref{Formalism/ThetaLogCoordinates/BihomogeneousPairs/LetterChainEmbeddings/LatticeEqualHolds.lean}{Formalism.InvariantFactorPaths.KP.LL.lattice_equal_holds}
\leanref{Formalism/ThetaLogCoordinates/BihomogeneousPairs/LetterChainEmbeddings/PlusOp.lean}{Formalism.InvariantFactorPaths.KP.LL.lattice_plus_holds}
\leanref{Formalism/ThetaLogCoordinates/BihomogeneousPairs/LetterChainEmbeddings/PairStep.lean}{Formalism.InvariantFactorPaths.KP.LL.HypJoin}
\leanref{Formalism/ThetaLogCoordinates/BihomogeneousPairs/LetterChainEmbeddings/PairStep.lean}{Formalism.InvariantFactorPaths.KP.LL.HypU}
\leanref{Formalism/SharpRestriction/GrowthBounds/MonomialFactorials/BlockValuationZpOfValuationBound.lean}{Formalism.MainNotions.G.SymmetricFormBlocks.hypTie_proved}
\leanref{Formalism/SharpRestriction/GrowthBounds/MonomialFactorials/BlockValuationZpOfValuationBound.lean}{Formalism.MainNotions.G.SymmetricFormBlocks.hypU_proved}

\begin{proof}
\emph{A tie $a=l$.} Here $w_l=x_{u_l}$, $e_l=y_{u_l}$ and $j'_+=j_+$. The map $T=p^{mL}(i!)^{-2}F_l^i$ is
$\mathbf Q_p$-linear, and $\Lambda_{(j,j_+)}$ is the $\mathbf Z_p$-span of the products $yg$ with
$y\in Y^{\mathbf Z}_{(j,j_+)}$ and $g$ a homogeneous polynomial of degree $b_+\dot-j_+$ in the $\xi_c$ with coefficients
in $\mathbf Z_p$; so it suffices to show $T(yg)\in\Lambda_{(j',j_+)}$ for these products. As $g$ involves no
variable of $x_{u_l}$ or $y_{u_l}$, $F_l(yg)=(F_ly)g$, and so $T(yg)=(Ty)g$. By
Lemma~\ref{pd-kappa-defects-segment-mean-theorem-a-block-bound-families-lattice-lowering-spaces}(1), applied $i$
times, $F_l^iy\in Y_{(j',j_+)}$, and so $Ty\in Y_{(j',j_+)}$. The polynomial $y$ has coefficients in $\mathbf Z_p$,
and $(i!)^{-2}F_l^i$ multiplies bounds by at most $p^{mL}$, by the step on powers of $F$ in the proof of
Lemma~\ref{pd-kappa-defects-segment-mean-theorem-a-plus-window-level-bounds}, for the $m$ pairs
$(x_{u_l,c},y_{u_l,c})$ and $i<p^{L+1}$. So $Ty$ has coefficients in $\mathbf Z_p$, and $Ty\in Y^{\mathbf Z}_{(j',j_+)}$.
As the degree $b_+\dot-j_+$ is the same for $(j',j_+)$, $(Ty)g$ is one of the products spanning $\Lambda_{(j',j_+)}$.

\emph{The tie $a=+$.} Here $w_+=x_+$ and $e_+=\xi$. For $1\leq c\leq m$ put $D_c=\partial^2/\partial x_{+,c}\partial\xi_c$,
so that $F_+=\sum_cD_c$; the $D_c$ commute. For $l\geq0$ put $\Lambda(l)=\Lambda_{(j,l)}$. We check the hypothesis
of Lemma~\ref{pd-kappa-defects-segment-mean-theorem-a-block-bound-families-lattice-lowering-sum}. Let $k,l\geq0$, let
$y\in Y^{\mathbf Z}_{(j,l)}$, and let $g$ be homogeneous of degree $b_+\dot-l$ in the $\xi_c$ with coefficients in
$\mathbf Z_p$. As $y$ involves no $\xi_c$ and $g$ no $x_{+,c}$,
\[
 (k!)^{-2}D_c^k(yg)=\Big(\frac1{k!}\frac{\partial^ky}{\partial x_{+,c}^k}\Big)\Big(\frac1{k!}\frac{\partial^kg}{\partial\xi_c^k}\Big).
\]
For a variable $v$, $(1/k!)\partial^k/\partial v^k$ maps $x^\kappa$ to $\binom{\kappa_v}kx^{\kappa-k\varepsilon_v}$, so it
preserves coefficients in $\mathbf Z_p$. By
Lemma~\ref{pd-kappa-defects-segment-mean-theorem-a-block-bound-families-lattice-lowering-spaces}(2), applied $k$
times, the first factor lies in $Y_{(j,l+k)}$, hence in $Y^{\mathbf Z}_{(j,l+k)}$. The second has coefficients in
$\mathbf Z_p$, and its monomials have degree $(b_+\dot-l)-k=b_+\dot-(l+k)$ in the $\xi_c$; a monomial of degree
less than $k$ in $\xi_c$ is killed. So $(k!)^{-2}D_c^k$ maps the products spanning $\Lambda(l)$ to products spanning
$\Lambda(l+k)$, and, being linear, it maps $\Lambda(l)$ into $\Lambda(l+k)$. Lemma~\ref{pd-kappa-defects-segment-mean-theorem-a-block-bound-families-lattice-lowering-sum},
with $S=\{1,\dots,m\}$, gives $p^{mL}(i!)^{-2}F_+^if\in\Lambda(j_++i)=\Lambda_{(j,j_++i)}$ for every
$f\in\Lambda(j_+)$.
\end{proof}
% Lean: LatTie and LatPlus are stated for every block datum with BlockOK under HypJoin (tieF at a tie maps Y of
%   (j, jp) into Y of (update j i (j i + 1), jp)) and HypU (pderiv at (uPos, c) maps Y of (j, jp) into Y of
%   (j, jp + 1)), for all (j, jp), with Nat.log p m <= L (the book's i < p^(L+1)); hypTie_proved and hypU_proved
%   (Lemma -lattice-lowering-spaces) prove HypJoin and HypU for every block datum with BlockOK (gB_blockOK).
%   LatTie is lattice_equal_holds = lat_tie: span_map_mem reduces to the generators rename chainEmb y * g of
%   Lat; tieF_pow_gen (tieF_gen, with xi_free, wv_ne_xi and ev_ne_xi) moves F^i past g; Y_iter (HypJoin
%   iterated, update_single) and int_Fpow (contr_Fsum with Mc = p^L, contr_D) give y in YZ (yz_mk). LatPlus is
%   lattice_plus_holds = po_Fsum (Lemma -lattice-lowering-sum) with Lam l = Lat p D j l and po_coord:
%   po_pow_right, po_pow_left and commute_pderiv split the operator, po_Y_iter (HypU iterated) and po_int
%   (contr_divDeriv) give the first factor, po_xi_div (po_R_pow, po_R_step) the second.

\begin{lemma}\label{pd-kappa-defects-segment-mean-theorem-a-block-bound-families-lattice-step}
Let an instance be given, $(\alpha,\alpha_+)\in W^{S_j}\times W$ and $b$ a block, and use the ties $a$, the operators
$\pi^\omega_a$ and the polynomials $\hat E_a$ of the proof of
Lemma~\ref{pd-kappa-defects-segment-mean-theorem-a-block-bound-families-fock}. For $0\leq k\leq t$ let $P^{(k)}_b$ be the
$\mathbf Z_p$-span of the products
\[
 \prod_l\hat E_l^{j_l}\cdot\hat E_+^{j_+}\cdot\prod_{l\leq k}z_l^{\zeta_l}\cdot z_+^{\zeta_+}\cdot
 (\pi^\omega_{k+1}\circ\cdots\circ\pi^\omega_t)(f)
\]
over the levels $\omega=(j,j_+)$ of $b$ and the $f\in\Lambda_\omega$, and let $P^{(+)}_b$ be the $\mathbf Z_p$-span of the
$\Theta_\omega(f)$ over the same $\omega$ and $f$. Let $A\geq1$ be an integer with $\delta(\lambda)\leq2A$ for every letter
$\lambda$, $b_+\leq2A$, $\zeta_l\leq2A$ for every $l$ and $\zeta_+\leq2A$, and let $u\geq m\lfloor\log_p(2A)\rfloor+
\lfloor\log_p(6A+m)\rfloor$ be an integer. Then $p^uP^{(+)}_b\subset P^{(0)}_b$, and $p^uP^{(k)}_b\subset P^{(k+1)}_b$ for
$0\leq k<t$.
\end{lemma}
\leanref{Formalism/ThetaLogCoordinates/BihomogeneousPairs/LetterChainEmbeddings/PairStep.lean}{Formalism.InvariantFactorPaths.KP.LL.Mid}
\leanref{Formalism/ThetaLogCoordinates/BihomogeneousPairs/LetterChainEmbeddings/PairStep.lean}{Formalism.InvariantFactorPaths.KP.LL.MidP}
\leanref{Formalism/ThetaLogCoordinates/BihomogeneousPairs/LetterChainEmbeddings/PairStep.lean}{Formalism.InvariantFactorPaths.KP.LL.PairStep}
\leanref{Formalism/ThetaLogCoordinates/BihomogeneousPairs/LetterChainEmbeddings/PstScalMem.lean}{Formalism.InvariantFactorPaths.KP.LL.pair_induction_holds}
\leanref{Formalism/ThetaLogCoordinates/BihomogeneousPairs/LetterChainEmbeddings/TieFPowMulFree.lean}{Formalism.InvariantFactorPaths.KP.LL.harmonicProjection_expand}
\leanref{Formalism/ThetaLogCoordinates/BihomogeneousPairs/LetterChainEmbeddings/PairStep.lean}{Formalism.InvariantFactorPaths.KP.LL.JoinCondition}
\leanref{Formalism/ThetaLogCoordinates/BihomogeneousPairs/LetterChainEmbeddings/PairStep.lean}{Formalism.InvariantFactorPaths.KP.LL.PlusCondition}
\leanref{Formalism/ThetaLogCoordinates/BihomogeneousPairs/LetterChainEmbeddings/StepTie.lean}{Formalism.InvariantFactorPaths.KP.LL.StepTie_of}
\leanref{Formalism/ThetaLogCoordinates/BihomogeneousPairs/LetterChainEmbeddings/StepPlus.lean}{Formalism.InvariantFactorPaths.KP.LL.StepPlus_of}

\begin{proof}
For a tie $a$ write $F_a$ as in
Lemma~\ref{pd-kappa-defects-segment-mean-theorem-a-block-bound-families-lattice-lowering}, put $\delta(\xi)=b_+$ and
$L=\lfloor\log_p(2A)\rfloor$. Let $\omega=(j,j_+)$ be a level of $b$ and $f\in\Lambda_\omega$.

\emph{Expansion at one tie.} Let $a$ be a tie, $o=\delta(w_a)-j_a$ and $n=\delta(e_a)-j_a$, so that $\pi^\omega_a$ is formed
with $(o,n)$ and $\zeta_a$, and put $A_a=o+n+m+\zeta_a-1$. As $F_a$ involves no $z_a$, Definition~\ref{pd-kappa-defects-segment-mean-theorem-a-block-bound-families-tie}(1)
gives, for every polynomial $Y$,
\[
 \pi^\omega_a(Y)=\sum_{i=0}^o\frac{(-1)^i}{(i!)^2\binom{A_a}i}\,\hat E_a^iz_a^{\zeta_a}F_a^iY .
\]
Let $Y=\Pi(f)$, where $\Pi$ is a composite of operators $\pi^\omega_{a'}$ for ties $a'\neq a$. As in the proof of
Lemma~\ref{pd-kappa-defects-segment-mean-theorem-a-block-bound-families-fock}, $F_a$ commutes with $\Pi$, so
$F_a^iY=\Pi(F_a^if)$. The polynomial $f$ is pure of bidegree $(o,n)$ for the tie $a$, so $F_a^if=0$ for
$i>\min(o,n)$. Let $i\leq\min(o,n)$. Then $\omega$ with $j_a$ raised by $i$ is a level $\omega_i$, and $\Pi$ is also the
composite of the $\pi^{\omega_i}_{a'}$, which depend on the level only through $j_{a'}$. As $i\leq2A<p^{L+1}$,
Lemma~\ref{pd-kappa-defects-segment-mean-theorem-a-block-bound-families-lattice-lowering} gives
$f_i=p^{mL}(i!)^{-2}F_a^if\in\Lambda_{\omega_i}$. As $i\leq A_a<6A+m$, $v_p\binom{A_a}i\leq\lfloor\log_p(6A+m)\rfloor\leq u-mL$
by Lemma~\ref{pd-kappa-defects-segment-mean-theorem-a-block-bound-families-valuation-factorials}(1); so
$g_i=(-1)^ip^{u-mL}/\binom{A_a}i$ lies in $\mathbf Z_p$, and
\[
 p^u\hat E_a^{j_a}\pi^\omega_a(\Pi f)=\sum_{i=0}^{\min(o,n)}g_i\,\hat E_a^{j_a+i}z_a^{\zeta_a}\,\Pi(f_i). \tag{$*$}
\]

\emph{The inclusions.} For $P^{(+)}_b$ take $a=+$ and $\Pi=\pi^\omega_1\circ\cdots\circ\pi^\omega_t$: multiplying $(*)$ by
$\prod_l\hat E_l^{j_l}$ writes $p^u\Theta_\omega(f)$ as a $\mathbf Z_p$-combination of the products
$\prod_l\hat E_l^{j_l}\cdot\hat E_+^{j_++i}\cdot z_+^{\zeta_+}\cdot(\pi^{\omega_i}_1\circ\cdots\circ\pi^{\omega_i}_t)(f_i)$, which
lie in $P^{(0)}_b$. For $P^{(k)}_b$ with $k<t$ take $a=k+1$ and $\Pi=\pi^\omega_{k+2}\circ\cdots\circ\pi^\omega_t$: multiplying
$(*)$ by $\prod_{l\neq k+1}\hat E_l^{j_l}\cdot\hat E_+^{j_+}\cdot\prod_{l\leq k}z_l^{\zeta_l}\cdot z_+^{\zeta_+}$ writes $p^u$ times
a spanning element of $P^{(k)}_b$ as a $\mathbf Z_p$-combination of spanning elements of $P^{(k+1)}_b$, for the levels
$\omega_i$. As $p^u$ times each spanning element lies in the target, so does $p^u$ times the span.
\end{proof}
% Lean: P^(k) is Mid p D k (midGen, with the projections PL over (finRange t).drop k and Lat p D j jp, the
%   level lattice) and P^(+) is MidP; the steps are the second halves of JoinCondition (StepTie_of, st_gen) and
%   PlusCondition (StepPlus_of), with u >= r Nat.log p (2a) + Nat.log p (6a + r). Both apply PairStep =
%   pair_induction_holds, whose second half is the expansion (*) (harmonicProjection_expand, pst_scal_mem for
%   the integrality of g_i through norm_choose_inv_le), to Lam l = the image under the later projections of
%   Lat at the level with j_a = l. The commutations are CommPL (comm_lattice_holds) and PLLin
%   (lattice_linear_holds), the independence of the later projections from j_a is ParamPL
%   (param_lattice_holds), F_a^i f = 0 for i > min(o, n) is tieF_pow_apply_eq_zero_of_min_lt with biHom_tie
%   and biHom_plus, and the lowering is LatTie and LatPlus (Lemma -lattice-lowering). The first halves of
%   PairStep, JoinCondition and PlusCondition (the reverse inclusions, through decomp) are not used here.

\begin{lemma}\label{pd-kappa-defects-segment-mean-theorem-a-block-bound-families-lattice}
Let $C_H\geq1$ and $C_{E1},C_{E2}\geq0$, and let an instance satisfy its hypotheses. Let
$(\alpha,\alpha_+)\in W^{S_j}\times W$ and $b$ a block kept for $\alpha$, and let $L^{\mathbf Z}_b$ be as in
Lemma~\ref{pd-kappa-defects-segment-mean-theorem-a-block-bound-dual}. Then $p^U\Theta_\omega f\in L^{\mathbf Z}_b$ for
every level $\omega$ of $b$ and every $f\in\Lambda_\omega$
(Definitions~\ref{pd-kappa-defects-segment-mean-theorem-a-levels}
and~\ref{pd-kappa-defects-segment-mean-theorem-a-block-bound-families-tie}).
\end{lemma}
\leanref{Formalism/SharpRestriction/GrowthBounds/MonomialFactorials/BlockValuationZpOfLatticeInclusion.lean}{Formalism.MainNotions.G.SymmetricFormBlocks.RestLL}
\leanref{Formalism/SharpRestriction/GrowthBounds/MonomialFactorials/BlockValuationZpOfValuationBound.lean}{Formalism.MainNotions.G.SymmetricFormBlocks.restLL_proved}
\leanref{Formalism/SharpRestriction/GrowthBounds/MonomialFactorials/GoodCountMem.lean}{Formalism.MainNotions.G.SymmetricFormBlocks.Th_mem_PSZ}

\begin{proof}
Let $P^{\mathbf Z}_b$ be the $\mathbf Z_p$-span of the $\Theta_{\omega'}(yg)$ over the levels $\omega'$ of $b$, with
$y\in Y^{\mathbf Z}_{\omega'}$ and $g$ a homogeneous polynomial in the $\xi_c$ with coefficients in $\mathbf Z_p$ of degree
$b_+-j'_+$. As $\Theta_\omega$ is linear and $\Lambda_\omega$ is the $\mathbf Z_p$-span of the products $yg$ for $\omega$,
$\Theta_\omega(\Lambda_\omega)\subset P^{\mathbf Z}_b$. It remains to show $p^UP^{\mathbf Z}_b\subset L^{\mathbf Z}_b$.

Write $s=S_j$ and $A=(3s+2)a$, and put $u=\lceil(m+1)(\log_p(8A+m)+1)\rceil$, so that $U=(t+1)u$
(Definition~\ref{pd-kappa-defects-segment-mean-theorem-a-block-bound-families-tie}(4)). As in the proof of
Lemma~\ref{pd-kappa-defects-segment-mean-theorem-a-block-bound-families-valuation}, $A\geq1$, and $\delta(\lambda)$ for
every letter $\lambda$, $b_+$, the $\zeta_l$ and $\zeta_+$ are at most $2A$. As $\lfloor\log_p(2A)\rfloor$ and
$\lfloor\log_p(6A+m)\rfloor$ are at most $\log_p(8A+m)$, $u\geq m\lfloor\log_p(2A)\rfloor+\lfloor\log_p(6A+m)\rfloor$. So
Lemma~\ref{pd-kappa-defects-segment-mean-theorem-a-block-bound-families-lattice-step} applies; in its notation
$P^{(+)}_b=P^{\mathbf Z}_b$, as $\Lambda_\omega$ is spanned by the products $yg$, and $P^{(t)}_b=L^{\mathbf Z}_b$, as for $k=t$ and
$f=yg$ its products are the generators of level $\omega$ over $\mathbf Q_p$
(Definition~\ref{pd-kappa-defects-segment-lattice}(5)). Its inclusions give
$p^{(k+1)u}P^{(+)}_b\subset P^{(k)}_b$ for $0\leq k\leq t$, by induction on $k$; for $k=t$ this is
$p^UP^{\mathbf Z}_b\subset L^{\mathbf Z}_b$.
\end{proof}
% Lean: Th_mem_PSZ for the first step. restLL_proved = restLL_of with hypTie_proved and hypU_proved:
%   restLL_of applies InvariantFactorPaths.KP.Top.levelLattice to gB aTop aPlus b with A = aBig = (3s + 2)a,
%   whose degree hypotheses are gB_deg_le, and takes its third part, p^expU PSZ in SSZ. levelLattice =
%   LevelLattice_of midTop midPEq stepTie stepPlus: asm_num is the bound for u, middle_top_holds (Mid t = SSZ)
%   and middle_p_eq_holds (MidP = PSZ) the identifications, and asm_chain composes the steps. levelLattice also
%   gives the equality of the Q_p-spans and p^expU SSZ in PSZ, which this book does not use.

\begin{lemma}\label{pd-kappa-defects-segment-mean-theorem-a-block-bound-families-valuation-factorials}
Let $p$ be a prime and $L\geq0$ an integer.
\begin{enumerate}
\item If $0\leq k\leq n<p^{L+1}$ are integers, then $v_p\binom nk\leq L$.
\item If $f_1,\dots,f_k\geq0$ are integers, then $f_1!\cdots f_k!$ divides $(f_1+\cdots+f_k)!$.
\item If moreover $f_1+\cdots+f_k<p^{L+1}$, then $v_p\big((f_1+\cdots+f_k)!\big)\leq\sum_{i=1}^kv_p(f_i!)+kL$.
\end{enumerate}
\end{lemma}
\leanref{Formalism/SharpRestriction/GrowthBounds/MonomialFactorials/ValFactSumLe.lean}{Formalism.MainNotions.G.SymmetricFormBlocks.val_fact_sum_le}
\leanref{Formalism/SharpRestriction/GrowthBounds/MonomialFactorials/ValFactSumLe.lean}{Formalism.MainNotions.G.SymmetricFormBlocks.fd}

\begin{proof}
(1) By Legendre's formula $v_p(n!)=\sum_{i\geq1}\lfloor n/p^i\rfloor$,
$v_p\binom nk=\sum_{i\geq1}\big(\lfloor n/p^i\rfloor-\lfloor k/p^i\rfloor-\lfloor(n-k)/p^i\rfloor\big)$. Each term is
$0$ or $1$, as $\lfloor x+y\rfloor-\lfloor x\rfloor-\lfloor y\rfloor$ is, and the terms with $i\geq L+1$ are $0$, as
$n<p^i$.

(2) As $\binom{f_1+f_2}{f_1}f_1!f_2!=(f_1+f_2)!$, $f_1!f_2!$ divides $(f_1+f_2)!$; the claim follows by induction on
$k$.

(3) By induction on $k$; for $k=0$ both sides are $0$. Put $T=f_2+\cdots+f_k$. Then
$(f_1+T)!=\binom{f_1+T}{f_1}f_1!\,T!$, and $f_1+T<p^{L+1}$, so by (1)
$v_p((f_1+T)!)\leq L+v_p(f_1!)+v_p(T!)$. As $T<p^{L+1}$, induction bounds $v_p(T!)$ by
$\sum_{i=2}^kv_p(f_i!)+(k-1)L$.
\end{proof}
% Lean: (1) is Mathlib's Nat.factorization_choose_le_log (with Nat.log p n <= L); (2) is
%   Nat.factorial_mul_factorial_dvd_factorial_add and Nat.prod_factorial_dvd_factorial_sum, used through fd
%   and fL_dvd; (3) is val_fact_sum_le (Finset induction, Nat.choose_mul_factorial_mul_factorial).

\begin{lemma}\label{pd-kappa-defects-segment-mean-theorem-a-block-bound-families-valuation-degrees}
Let an instance be given, $(\alpha,\alpha_+)\in W^{S_j}\times W$ and $b$ a block, use the notation of
Definitions~\ref{pd-kappa-defects-segment-lattice}, \ref{pd-kappa-defects-segment-mean-theorem-a-levels}
and~\ref{pd-kappa-defects-segment-mean-theorem-a-block-bound-families-tie} over $\mathbf Q_p$, and let
$L^{\mathbf Z}_b$ be as in Lemma~\ref{pd-kappa-defects-segment-mean-theorem-a-block-bound-dual}. Sort the
variables into groups: the group of a letter $\lambda$ consists of its variables $\lambda_1,\dots,\lambda_m$, together
with $z_l$ if $\lambda=x_{u_l}$, with $z'_l$ if $\lambda=y_{u_l}$, and with $z_+$ if $\lambda=x_+$; the last group $\Xi$
consists of $\xi_1,\dots,\xi_m$ and $z'_+$. Put $d(x_{u_l})=\delta(x_{u_l})+\zeta_l$ for $1\leq l\leq t$,
$d(x_+)=\delta(x_+)+\zeta_+$, $d(\Xi)=b_+$ and $d(\lambda)=\delta(\lambda)$ for the other letters $\lambda$. Then every
monomial with a nonzero coefficient in an element of $L^{\mathbf Z}_b$ has total degree $d(\gamma)$ in the
variables of each group $\gamma$.
\end{lemma}
\leanref{Formalism/SharpRestriction/GrowthBounds/MonomialFactorials/SLetCastSucc.lean}{Formalism.MainNotions.G.SymmetricFormBlocks.SSZ_support}
\leanref{Formalism/SharpRestriction/GrowthBounds/MonomialFactorials/SLetCastSucc.lean}{Formalism.MainNotions.G.SymmetricFormBlocks.hom_gens}
\leanref{Formalism/SharpRestriction/GrowthBounds/MonomialFactorials/SLetCastSucc.lean}{Formalism.MainNotions.G.SymmetricFormBlocks.wL}
\leanref{Formalism/SharpRestriction/GrowthBounds/MonomialFactorials/SLetCastSucc.lean}{Formalism.MainNotions.G.SymmetricFormBlocks.dLvl}
\leanref{Formalism/SharpRestriction/GrowthBounds/MonomialFactorials/FinPosEqCast.lean}{Formalism.MainNotions.G.SymmetricFormBlocks.dTot}
\leanref{Formalism/SharpRestriction/GrowthBounds/MonomialFactorials/FinPosEqCast.lean}{Formalism.MainNotions.G.SymmetricFormBlocks.dLvl_eq}

\begin{proof}
Let $\omega=(j,j_+)$ be a level of $b$, and consider a generator
$\prod_l\hat E_l^{j_l}\cdot\hat E_+^{j_+}\cdot\prod_lz_l^{\zeta_l}\cdot z_+^{\zeta_+}\cdot f\cdot g$ of level $\omega$ over
$\mathbf Q_p$. Each factor is homogeneous in the variables of each group. The factor
$\hat E_l=z_lz'_l+\sum_cx_{u_l,c}y_{u_l,c}$ has degree $1$ in the groups of $x_{u_l}$ and of $y_{u_l}$ and $0$ in the
others; $\hat E_+=z_+z'_++\sum_cx_{+,c}\xi_c$ has degree $1$ in the groups of $x_+$ and $\Xi$; $z_l$ and $z_+$ have
degree $1$ in the groups of $x_{u_l}$ and $x_+$; $f$ has degree $\delta'_\omega(\lambda)$ in the group of each letter
$\lambda$; and $g$ has degree $b_+-j_+$ in $\Xi$. So the generator has degree $j_l+\zeta_l+\delta'_\omega(x_{u_l})$ in the
group of $x_{u_l}$, $j_l+\delta'_\omega(y_{u_l})$ in that of $y_{u_l}$, $j_++\zeta_++\delta'_\omega(x_+)$ in that of $x_+$,
$j_++(b_+-j_+)$ in $\Xi$, and $\delta'_\omega(\lambda)$ in the group of every other letter. The letters
$x_{u_1},y_{u_1},\dots,x_{u_t},y_{u_t},x_+$ are distinct, $j_l\leq\min(\delta(x_{u_l}),\delta(y_{u_l}))$ and
$j_+\leq\min(\delta(x_+),b_+)$; so $\delta'_\omega$ is $\delta-j_l$ at $x_{u_l}$ and at $y_{u_l}$, $\delta-j_+$ at $x_+$, and
$\delta$ at the other letters, and these degrees are the $d(\gamma)$, for every level. A monomial with a nonzero
coefficient in a $\mathbf Z_p$-combination of generators has a nonzero coefficient in one of them.
\end{proof}
% Lean: SSZ_support (span induction) with hom_gens: each generator is weighted homogeneous of degree
%   dLvl G j jp a for the weight wL G a of the group a (zLet and sLet place z_l, z'_l, z_+ and z'_+), by
%   hom_Ehat, hom_Eplus, hom_chain and hom_xi. dLvl_eq (level_room, from BlockOK, which holds for the block
%   data of an instance by gB_blockOK) shows dLvl G j jp = dTot G, which does not depend on the level.

\begin{lemma}\label{pd-kappa-defects-segment-mean-theorem-a-block-bound-families-valuation}
Let $C_H\geq1$ and $C_{E1},C_{E2}\geq0$, and let an instance satisfy its hypotheses. Let
$(\alpha,\alpha_+)\in W^{S_j}\times W$ and $b$ a block kept for $\alpha$, and let $L^{\mathbf Z}_b$ be as in
Lemma~\ref{pd-kappa-defects-segment-mean-theorem-a-block-bound-dual}. Then
$v_p(c_\omega d_\omega)+2U\leq v_p(\nu!)+D_{\mathrm{lv}}$ for every level $\omega$ of $b$ and every monomial $\nu$ with a
nonzero coefficient in an element of $L^{\mathbf Z}_b$
(Definitions~\ref{pd-kappa-defects-segment-mean-theorem-a-levels}
and~\ref{pd-kappa-defects-segment-mean-theorem-a-block-bound-families-tie}).
\end{lemma}
\leanref{Formalism/SharpRestriction/GrowthBounds/MonomialFactorials/BlockValuationZpOfLatticeInclusion.lean}{Formalism.MainNotions.G.SymmetricFormBlocks.RestVal}
\leanref{Formalism/SharpRestriction/GrowthBounds/MonomialFactorials/BlockValuationZpOfNumericBound.lean}{Formalism.MainNotions.G.SymmetricFormBlocks.restVal_of_num}
\leanref{Formalism/SharpRestriction/GrowthBounds/MonomialFactorials/BlockValuationOverZpHolds.lean}{Formalism.MainNotions.G.SymmetricFormBlocks.restValNum_proved}

\begin{proof}
Write $s=S_j$, $A=(3s+2)a$ and $L=\lfloor\log_p(8A+m)\rfloor$, and let $\ell$ and $\kappa_0$ be as in
Definition~\ref{pd-kappa-defects-numbers}. Let $\omega=(j,j_+)$ be a level of $b$, $\nu$ a monomial with a nonzero
coefficient in an element of $L^{\mathbf Z}_b$, and the groups $\gamma$ and the degrees $d(\gamma)$ as in
Lemma~\ref{pd-kappa-defects-segment-mean-theorem-a-block-bound-families-valuation-degrees}.

\emph{Facts.} As in the proof of Lemma~\ref{pd-kappa-defects-segment-mean-letters-degrees}, the hypotheses give
$m\geq4$, $2s+2\leq m$, $p\geq(m+1)^2$ and $a\geq1$; so $3s+2\leq p$. As $b$ is kept for $\alpha$, the degrees satisfy
$\delta(x_i)\leq a/10$, $\delta(y_i)\leq2a$ and $\delta(x_+)\leq a+sa/10+2as+a/10$, and $b_+\leq a/10$ and
$\zeta_l\leq\hat b_{u_l}\leq a/10$; and $W$ lies in $[a,2a]$, so $\zeta_+\leq\alpha_+\leq2a$. All of these are at most
$2A=(6s+4)a$. Also $j_l\leq\delta(x_{u_l})$ and $j_+\leq\delta(x_+)$, and $\delta'_\omega\leq\delta$.

\emph{The factorials.} For each group $\gamma$ let $F_\gamma$ be the product of the factorials of the summands of
$d(\gamma)$ in the proof of
Lemma~\ref{pd-kappa-defects-segment-mean-theorem-a-block-bound-families-valuation-degrees}: $F_\gamma=j_l!\,\zeta_l!\,
\delta'_\omega(x_{u_l})!$ for the group of $x_{u_l}$, $j_l!\,\delta'_\omega(y_{u_l})!$ for that of $y_{u_l}$,
$j_+!\,\zeta_+!\,\delta'_\omega(x_+)!$ for that of $x_+$, $j_+!\,(b_+-j_+)!$ for $\Xi$, and $\delta'_\omega(\lambda)!$ for
every other letter $\lambda$. The product of the $F_\gamma$ is
$(\prod_lj_l!\cdot j_+!)^2\prod_l\zeta_l!\cdot\zeta_+!\cdot\prod_\lambda\delta'_\omega(\lambda)!\cdot(b_+-j_+)!$; so, by
Definition~\ref{pd-kappa-defects-segment-mean-theorem-a-block-bound-families-tie}(3),
\[
 c_\omega d_\omega=\prod_\gamma F_\gamma\cdot\prod_{l=1}^tB_l\cdot B_+,\qquad
 B_l=\binom{\delta'_\omega(x_{u_l})+\delta'_\omega(y_{u_l})+m+\zeta_l+j_l}{j_l},\quad
 B_+=\binom{\delta'_\omega(x_+)+b_+-j_++m+\zeta_++j_+}{j_+}.
\]
By Lemma~\ref{pd-kappa-defects-segment-mean-theorem-a-block-bound-families-valuation-factorials}(2), $F_\gamma$
divides $d(\gamma)!$ for every $\gamma$, so $v_p(\prod_\gamma F_\gamma)\leq\sum_\gamma v_p(d(\gamma)!)$. The top entries of
the $B_l$ and of $B_+$ are at most $4\cdot2A+m$, which is less than $p^{L+1}$; so by
Lemma~\ref{pd-kappa-defects-segment-mean-theorem-a-block-bound-families-valuation-factorials}(1),
$v_p(c_\omega d_\omega)\leq\sum_\gamma v_p(d(\gamma)!)+(t+1)L$.

\emph{The monomial.} By
Lemma~\ref{pd-kappa-defects-segment-mean-theorem-a-block-bound-families-valuation-degrees}, the exponents of $\nu$
at the variables of a group $\gamma$ add up to $d(\gamma)$, and $d(\gamma)\leq4A<p^{L+1}$. By
Lemma~\ref{pd-kappa-defects-segment-mean-theorem-a-block-bound-families-valuation-factorials}(3),
$v_p(d(\gamma)!)$ is at most the sum of the $v_p(\nu_v!)$ over the variables $v$ of $\gamma$ plus $L$ times their
number. The groups partition the $(2s+2)m+2(t+1)$ variables, and $v_p(\nu!)$ is the sum of all the $v_p(\nu_v!)$. So
\[
 v_p(c_\omega d_\omega)\leq v_p(\nu!)+\big((2s+2)m+3(t+1)\big)L .
\]

\emph{The numbers.} As $8A+m\leq(3s+2)(8a+m)$ and $3s+2\leq p$, $\log_p(8A+m)\leq1+\log_p(8a+m)=\ell$; so
$0\leq L\leq\ell$, and $\ell\geq1$. As $\lceil x\rceil<x+1$ and $t\leq s$,
\[
 U\leq(t+1)\big((m+1)(\ell+1)+1\big)\leq(s+1)\big((m+1)\ell+m+2\big)\leq(s+1)(2m+3)\ell ,
\]
and $((2s+2)m+3(t+1))L\leq(s+1)(2m+3)\ell$. So $((2s+2)m+3(t+1))L+2U\leq(s+1)(6m+9)\ell\leq10(s+1)(m+1)\ell$, which is
at most $D_{\mathrm{lv}}=10(s+1)(m+1)\ell+\kappa_0$, as $\kappa_0=2|T|(m+1)\log_pa\geq0$. With the last paragraph,
$v_p(c_\omega d_\omega)+2U\leq v_p(\nu!)+D_{\mathrm{lv}}$.
\end{proof}
% Lean: restVal_of_num with SSZ_support (Lemma -valuation-degrees) and restValNum_proved. The facts are hiFacts
%   and gB_deg_le (degrees at most 2A with A = aBig; segChain_deg_le, coarseA_le, and aPlus in winA for zPlus).
%   val_gen is the next two paragraphs: lvC_eq and prod_fL write lvC degFact as the product of the fL G j jp a
%   (the book's F_gamma) and the binT and binP; val_lvC_le bounds it by fL_dvd and Nat.factorization_choose_le_log
%   (the top entries bounded through dTot_le); val_letters_le is the monomial paragraph (weight_eq, letOf,
%   val_fact_sum_le, card of BVar = (N + 1) r + (t + 1) 2 with N = 2s + 1). The numbers are hbig, hL1
%   (Real.natLog_le_logb), hU (Nat.ceil_lt_add_one), hk0 and loss_num.

\begin{lemma}\label{pd-kappa-defects-segment-mean-theorem-a-block-bound-families}
Let $C_H\geq1$ and $C_{E1},C_{E2}\geq0$, and let an instance satisfy its hypotheses. Let
$(\alpha,\alpha_+)\in W^{S_j}\times W$, let $b$ be a block kept for $\alpha$, let $D\geq0$ be an integer, and let
$L^{\mathbf Z}_b$ be as in Lemma~\ref{pd-kappa-defects-segment-mean-theorem-a-block-bound-dual}. Then there are
$n\geq0$, $x_1,\dots,x_n,y_1,\dots,y_n\in L^{\mathbf Z}_b$ and an integer $e$ such that $[x_i,y_l]=p^e\delta_{il}$
for all $i,l$, $e\leq v_p(\nu!)+D+D_{\mathrm{lv}}$ for every monomial $\nu$ with a nonzero coefficient in an
element of $L^{\mathbf Z}_b$, and the sum of the $\operatorname{rk}\Lambda_\omega$ over the levels $\omega$ of $b$ is
at most $n+g_D(\alpha,\alpha_+,b)$ (Definition~\ref{pd-kappa-defects-segment-mean-theorem-a-levels}).
\end{lemma}
\leanref{Formalism/SharpRestriction/GrowthBounds/MonomialFactorials/BlockBoundUniformiser.lean}{Formalism.MainNotions.G.SymmetricFormBlocks.BlockBoundIntegers}
\leanref{Formalism/SharpRestriction/GrowthBounds/MonomialFactorials/BlockValuationOverZpHolds.lean}{Formalism.MainNotions.G.SymmetricFormBlocks.block_valuation_over_Zp_holds}

\begin{proof}
Use the notation of Definition~\ref{pd-kappa-defects-segment-mean-theorem-a-block-bound-families-tie}, and let
$\Omega$ be the set of the levels of $b$, a finite set that contains $(0,0)$.

\emph{The level families.} For $\omega\in\Omega$ let $n_\omega=\operatorname{good}_D(\beta_\omega,\Lambda_\omega)$, and
choose $u^\omega_1,\dots,u^\omega_{n_\omega},v^\omega_1,\dots,v^\omega_{n_\omega}\in\Lambda_\omega$ with
$\beta_\omega(u^\omega_{i'},v^\omega_i)=p^D\delta_{ii'}$, that is $[\Phi_\omega u^\omega_{i'},v^\omega_i]=p^Dd_\omega\delta_{ii'}$. Let
$M$ be the largest of the integers $v_p(c_\omega d_\omega)$, $\omega\in\Omega$. As $v_p(c_\omega d_\omega)\leq M$, there is
$s_\omega\in\mathbf Z_p$ with $s_\omega c_\omega d_\omega=p^M$.

\emph{The family.} For $\omega\in\Omega$ and $1\leq i,i'\leq n_\omega$ put $x_{\omega,i}=p^U\Theta_\omega v^\omega_i$ and
$y_{\omega,i'}=s_\omega p^U\Theta_\omega u^\omega_{i'}$. By
Lemma~\ref{pd-kappa-defects-segment-mean-theorem-a-block-bound-families-lattice}, these lie in $L^{\mathbf Z}_b$. By
Lemma~\ref{pd-kappa-defects-segment-mean-theorem-a-block-bound-families-fock}, $[x_{\omega,i},y_{\omega',i'}]=0$ for
$\omega\neq\omega'$, and, as the Fock pairing is symmetric,
\[
 [x_{\omega,i},y_{\omega,i'}]=p^{2U}s_\omega c_\omega[v^\omega_i,\Phi_\omega u^\omega_{i'}]
 =p^{2U}s_\omega c_\omega[\Phi_\omega u^\omega_{i'},v^\omega_i]=p^{2U}s_\omega c_\omega d_\omega p^D\delta_{ii'}
 =p^{D+2U+M}\delta_{ii'}.
\]
Number the pairs $(\omega,i)$ as $1,\dots,n$, with $n=\sum_\omega n_\omega$. This gives
$x_1,\dots,x_n,y_1,\dots,y_n\in L^{\mathbf Z}_b$ with $[x_i,y_l]=p^e\delta_{il}$, where $e=D+2U+M$.

\emph{The ranks.} By Definition~\ref{pd-kappa-defects-segment-mean-theorem-a-levels}(5) and~(6),
$\operatorname{rk}\Lambda_\omega=n_\omega+\operatorname{bad}_\omega(\alpha,\alpha_+,b;D)$; summing over $\Omega$, the sum of
the $\operatorname{rk}\Lambda_\omega$ is $n+g_D(\alpha,\alpha_+,b)$.

\emph{The valuation.} Choose $\omega_0\in\Omega$ with $v_p(c_{\omega_0}d_{\omega_0})=M$. By
Lemma~\ref{pd-kappa-defects-segment-mean-theorem-a-block-bound-families-valuation} at $\omega_0$,
$M+2U\leq v_p(\nu!)+D_{\mathrm{lv}}$ for every monomial $\nu$ with a nonzero coefficient in an element of
$L^{\mathbf Z}_b$; so $e\leq v_p(\nu!)+D+D_{\mathrm{lv}}$.
\end{proof}
% Lean: block_valuation_over_Zp_holds hCYb hCE1 hCE2 hH = block_valuation_Zp_of_numeric_bound
%   (restValNum_proved) = block_valuation_Zp_of_valuation_bound (restVal_of_num) =
%   block_valuation_Zp_of_lattice_inclusion with restLL_proved. That is blkZ_gen with U = expU (aBig X) and
%   Mx = expM, the largest v_p(lvC degFact) over the levels (le_expM): nLv_mem gives the level families,
%   exists_scal the scalars sc, Th_mem_PSZ with RestLL the membership, Th_fock and Th_orth (with fockK_symm)
%   the pairings, and the count is Fintype.card_sigma with le_add_tsub (badLvl = bad = rank - goodCount);
%   the valuation is exists_expM with RestVal. The Lean's sum is over lv aTop b, the levels of gB aTop 0 b,
%   and blkZ_gen's over the levels of gB aTop aPlus b, which are the same (lv_eq). levels_nonempty gives
%   (0, 0) in the levels.

\begin{lemma}\label{pd-kappa-defects-segment-mean-theorem-a-block-bound}
Let $C_H\geq1$ and $C_{E1},C_{E2}\geq0$, and let an instance satisfy its hypotheses. Let
$(\alpha,\alpha_+)\in W^{S_j}\times W$, let $b$ be a block kept for $\alpha$, and let $D\geq0$ be an integer.
Then $L_b$ has a good sublattice $L_2$ of depth $D+D_{\mathrm{lv}}$ for $\beta_b$ with
$\operatorname{rk}L_b\leq\operatorname{rk}L_2+g_D(\alpha,\alpha_+,b)$
(Definitions~\ref{pd-kappa-defects-segment-mean-theorem-a-classes}
and~\ref{pd-kappa-defects-segment-mean-theorem-a-levels}).
\end{lemma}
\leanref{Formalism/SharpRestriction/GrowthBounds/MonomialFactorials/BlockBoundUniformiser.lean}{Formalism.MainNotions.G.SymmetricFormBlocks.BlockBoundUniformiser}
\leanref{Formalism/SharpRestriction/GrowthBounds/MonomialFactorials/BlockBoundOfDualFamilies.lean}{Formalism.MainNotions.G.SymmetricFormBlocks.block_bound_of_dual_families}

\begin{proof}
By Lemma~\ref{pd-kappa-defects-segment-mean-theorem-a-block-bound-families} there are $n\geq0$, elements
$x_1,\dots,x_n,y_1,\dots,y_n$ of $L^{\mathbf Z}_b$ and an integer $e$ with $[x_i,y_l]=p^e\delta_{il}$ and
$e\leq v_p(\nu!)+D+D_{\mathrm{lv}}$ for every monomial $\nu$ occurring in an element of $L^{\mathbf Z}_b$, such that
the sum of the $\operatorname{rk}\Lambda_\omega$ over the levels $\omega$ of $b$ is at most
$n+g_D(\alpha,\alpha_+,b)$. By Lemma~\ref{pd-kappa-defects-segment-mean-theorem-a-block-bound-dual}, $L_b$ has a
good sublattice $L_2$ of depth $D+D_{\mathrm{lv}}$ for $\beta_b$ with $n\leq\operatorname{rk}L_2$. By
Lemma~\ref{pd-kappa-defects-segment-mean-theorem-a-block-bound-rank},
\[
 \operatorname{rk}L_b\leq\sum_\omega\operatorname{rk}\Lambda_\omega\leq n+g_D(\alpha,\alpha_+,b)\leq
 \operatorname{rk}L_2+g_D(\alpha,\alpha_+,b).
\]
\end{proof}
% Lean: openE_of_open4 uses block_bound_of_dual_families (block_valuation_over_Zp_holds hCYb hCE1 hCE2 hH)
%   RestBCdim: hZ gives the family (Lemma pd-kappa-defects-segment-mean-theorem-a-block-bound-families),
%   block_dual_family_bound the good sublattice (-dual), and RestBCdim = restBCdim_of (restBClvl_of
%   (restYbc_of_funY restFunY_proved)) the rank bound (-rank).

\subsection{Level pairings and top models}\label{subsection-pd-level-pairings-top-models}

\begin{definition}\label{pd-kappa-defects-segment-mean-theorem-a-plus-pairing}
Let an instance be given, $\alpha\in W^{S_j}$, $b$ a block and $\omega=(j,j_+)$ a level of $b$, and use the
notation of Definition~\ref{pd-kappa-defects-segment-mean-theorem-a-levels}. The \emph{level pairing without
the factor at $x_+$} is
\[
 \beta^-_\omega(f,f')=\frac{[\Phi^Y_\omega f,f']}{(b_+-j_+)!\prod_\lambda\delta'_\omega(\lambda)!}\quad\text{on }\Lambda_\omega,
\]
and for an integer $D$ put $\operatorname{bad}^-_\omega(\alpha,b;D)=\operatorname{bad}_D(\beta^-_\omega,\Lambda_\omega)$.
\end{definition}
\leanref{Formalism/SharpRestriction/GrowthBounds/MonomialFactorials/RestPlusOf.lean}{Formalism.MainNotions.G.SymmetricFormBlocks.phisB}
\leanref{Formalism/SharpRestriction/GrowthBounds/MonomialFactorials/RestPlusOf.lean}{Formalism.MainNotions.G.SymmetricFormBlocks.bePhi}

\begin{lemma}\label{pd-kappa-defects-segment-mean-theorem-a-plus-rank}
Let an instance be given, $\alpha\in W^{S_j}$, $b$ a block and $\omega=(j,j_+)$ a level of $b$. Then
$\operatorname{rk}\Lambda_\omega\leq\dim Y_\omega\cdot N_\xi(b_+-j_+)$
(Definition~\ref{pd-kappa-defects-segment-mean-theorem-a-levels}).
\end{lemma}
\leanref{Formalism/SharpRestriction/GrowthBounds/MonomialFactorials/RestPlusOf.lean}{Formalism.MainNotions.G.SymmetricFormBlocks.RestRankL}
\leanref{Formalism/SharpRestriction/GrowthBounds/MonomialFactorials/RestRankLProved.lean}{Formalism.MainNotions.G.SymmetricFormBlocks.finrank_Ljj_le}

\begin{proof}
Let $y_1,\dots,y_d$ be a basis of $Y_\omega$, $d=\dim Y_\omega$, and let $\xi^\beta$ run over the $N_\xi(b_+-j_+)$
monomials of degree $b_+-j_+$ in the $\xi_c$. Every $f\in Y^{\mathbf Z}_\omega$ is a $\mathbf Q_p$-combination of the
$y_a$, and every homogeneous $g$ of degree $b_+-j_+$ in the $\xi_c$ is one of the $\xi^\beta$; so $\Lambda_\omega$ lies in
the $\mathbf Q_p$-span of the $dN_\xi(b_+-j_+)$ products $y_a\xi^\beta$. The rank of a $\mathbf Z_p$-submodule of a
$\mathbf Q_p$-vector space is at most the dimension of any subspace containing it.
\end{proof}
% Lean: restRankL_proved = finrank_Ljj_le (xi_mem_span, Submodule.mul_mem_mul, finrank_zp_le,
%   finrank_range_le_card and card_xiMons).

\begin{lemma}\label{pd-kappa-defects-segment-mean-theorem-a-plus-phi}
Let an instance be given, $\alpha\in W^{S_j}$, $b$ a block, $\omega=(j,j_+)$ a level of $b$ and $D$ an integer. Then
$\operatorname{bad}^-_\omega(\alpha,b;D)\leq\operatorname{bad}^Y_\omega(\alpha,b;D)\cdot N_\xi(b_+-j_+)$
(Definitions~\ref{pd-kappa-defects-segment-mean-theorem-a-levels}
and~\ref{pd-kappa-defects-segment-mean-theorem-a-plus-pairing}).
\end{lemma}
\leanref{Formalism/SharpRestriction/GrowthBounds/MonomialFactorials/RestPlusOf.lean}{Formalism.MainNotions.G.SymmetricFormBlocks.RestPlusA}
\leanref{Formalism/SharpRestriction/GrowthBounds/MonomialFactorials/PlusAWICondition.lean}{Formalism.MainNotions.G.SymmetricFormBlocks.plusAW_of_blockOK}
\leanref{Formalism/SharpRestriction/GrowthBounds/MonomialFactorials/RestPlusAOf.lean}{Formalism.MainNotions.G.SymmetricFormBlocks.plusA_of}

\begin{proof}
Write $n=b_+-j_+$, $g=\operatorname{good}_D(\beta^Y_\omega,Y^{\mathbf Z}_\omega)$ and $\Delta'=\prod_\lambda\delta'_\omega(\lambda)!$,
and let $x_1,\dots,x_g,y_1,\dots,y_g\in Y^{\mathbf Z}_\omega$ with $\beta^Y_\omega(x_i,y_{i'})=p^D\delta_{ii'}$. For a
monomial $\xi^\beta$ of degree $n$ let $c_\beta=n!/\beta!$, an integer. We show that the $x_i\xi^\beta$ and the
$c_\beta y_i\xi^\beta$, which lie in $\Lambda_\omega$, satisfy
$\beta^-_\omega(x_i\xi^\beta,c_{\beta'}y_{i'}\xi^{\beta'})=p^D\delta_{ii'}\delta_{\beta\beta'}$. The operator $\Phi^Y_\omega$
involves only the variables of the letters, and differentiation in these variables kills $\xi^\beta$; so
$\Phi^Y_\omega(x_i\xi^\beta)=(\Phi^Y_\omega x_i)\xi^\beta$. For polynomials $F,G$ in the variables of the letters,
$[F\xi^\beta,G\xi^{\beta'}]$ is $\beta![F,G]$ if $\beta=\beta'$ and $0$ otherwise. Hence
\[
 \beta^-_\omega(x_i\xi^\beta,c_{\beta'}y_{i'}\xi^{\beta'})=\delta_{\beta\beta'}\,\frac{c_\beta\,\beta!\,[\Phi^Y_\omega x_i,y_{i'}]}
 {n!\,\Delta'}=\delta_{\beta\beta'}\beta^Y_\omega(x_i,y_{i'})=p^D\delta_{ii'}\delta_{\beta\beta'}.
\]
So $\operatorname{good}_D(\beta^-_\omega,\Lambda_\omega)\geq gN_\xi(n)$. By
Lemma~\ref{pd-kappa-defects-segment-mean-theorem-a-plus-rank}, $\operatorname{rk}\Lambda_\omega\leq\dim Y_\omega\cdot
N_\xi(n)$, and $\operatorname{rk}Y^{\mathbf Z}_\omega=\dim Y_\omega$, as every element of $Y_\omega$ has a multiple in
$Y^{\mathbf Z}_\omega$. So $\operatorname{bad}^-_\omega(\alpha,b;D)\leq(\dim Y_\omega-g)N_\xi(n)=
\operatorname{bad}^Y_\omega(\alpha,b;D)\cdot N_\xi(n)$.
\end{proof}
% Lean: plusACondition = restPlusA_of (plusAWICondition) = plusA_of (finrank_Ljj_le, finrank_lattice) with
%   plusAW_of_blockOK: goodCount_mem, cfn = (B - jp)!/m!, mono_mem_xiPolZ, bePhi_smul_right,
%   phisB_rename_mul (with pderiv_chain_mono), fockK_rename_mul, cf_mul, and gcL_bdd for the bound.

\begin{lemma}\label{pd-kappa-defects-segment-mean-theorem-a-plus-lattice}
Let $p$ be an odd prime, $V$ a finite-dimensional $\mathbf Q_p$-vector space, $L\subset V$ a finitely generated
$\mathbf Z_p$-submodule, $B\colon V\times V\to\mathbf Q_p$ a map that is $\mathbf Q_p$-linear in the first variable
and satisfies $B(u,v)=B(v,u)$ for $u,v\in L$, and $D$ an integer. Then there are $n_1\geq0$ and
$e_1,\dots,e_{n_1}\in L$, linearly independent over $\mathbf Z_p$, with $B(e_a,e_b)=0$ for $a\neq b$, $B(e_a,e_a)\neq0$
and $v_p(B(e_a,e_a))\leq D$ for every $a$, and $n_1\leq\operatorname{rk}L\leq n_1+\operatorname{bad}_D(B,L)$
(Definition~\ref{pd-kappa-defects-segment-mean-theorem-a-levels}(5)).
\end{lemma}
\leanref{Formalism/SharpRestriction/GrowthBounds/MonomialFactorials/BePhiAddLeft.lean}{Formalism.MainNotions.G.SymmetricFormBlocks.OrthLat}
\leanref{Formalism/SharpRestriction/GrowthBounds/MonomialFactorials/OrthCountOf.lean}{Formalism.MainNotions.G.SymmetricFormBlocks.orthLat_of}
\leanref{Formalism/SharpRestriction/GrowthBounds/MonomialFactorials/OrthBasisLatOf.lean}{Formalism.MainNotions.G.SymmetricFormBlocks.orthBasisLat_of}
\leanref{Formalism/SharpRestriction/GrowthBounds/MonomialFactorials/OrthSplitOfMin.lean}{Formalism.MainNotions.G.SymmetricFormBlocks.orthSplit_of}
\leanref{Formalism/PeriodIntegrals/UnitaryGroupMassFormula/MinimalLatticeNorm/NormTwoEqOne.lean}{PI.LatMinWork.main}
\leanref{Formalism/PeriodIntegrals/UnitaryGroupMassFormula/DiagonalRank/DetOneSubPNe.lean}{PI.UnitaryGroupMassFormula.DiagRank.diagRank}

\begin{proof}
\emph{Step 1: a vector of minimal valuation.} Suppose $B(u,v)\neq0$ for some $u,v\in L$. Then there is $x\in L$
with $B(x,x)\neq0$ and $B(u,x)\in\mathbf Z_pB(x,x)$ for every $u\in L$. Indeed, let $S$ be a finite generating set of
$L$ and choose $a,b\in S$ with $|B(a,b)|_p$ largest among the pairs in $S$. By linearity in each variable and
the ultrametric inequality, $|B(u,v)|_p\leq|B(a,b)|_p$ for all $u,v\in L$; so $|B(a,b)|_p>0$. If
$|B(a,a)|_p=|B(a,b)|_p$, take $x=a$; if $|B(b,b)|_p=|B(a,b)|_p$, take $x=b$. Otherwise take $x=a+b$: then
$B(x,x)=2B(a,b)+B(a,a)+B(b,b)$, where $|2B(a,b)|_p=|B(a,b)|_p$ as $p$ is odd and
$|B(a,a)+B(b,b)|_p<|B(a,b)|_p$, so $|B(x,x)|_p=|B(a,b)|_p$. In each case $|B(u,x)/B(x,x)|_p\leq1$ for $u\in L$.

\emph{Step 2: an orthogonal basis.} We show that $L$ has a basis $f_1,\dots,f_m$ over $\mathbf Z_p$ with
$B(f_a,f_b)=0$ for $a\neq b$, by induction on $\operatorname{rk}L$. If $B$ vanishes on $L\times L$, any basis will
do; $L$ is free, being finitely generated and torsion-free over $\mathbf Z_p$. Otherwise take $x$ as in Step 1, and
let $L'$ be the set of $u\in L$ with $B(u,x)=0$. Every $u\in L$ is $cx+(u-cx)$ with $c=B(u,x)/B(x,x)\in\mathbf Z_p$
and $u-cx\in L'$; and if $cx+y=0$ with $y\in L'$, then $cB(x,x)=0$, so $c=0$. So $L=\mathbf Z_px\oplus L'$ and
$\operatorname{rk}L'=\operatorname{rk}L-1$. By induction $L'$ has an orthogonal basis, and adding $x$, which is
orthogonal to $L'$ by symmetry, gives one of $L$.

\emph{Step 3: counting.} Let $f_1,\dots,f_m$ be as in Step 2, so $m=\operatorname{rk}L$, put $c_a=B(f_a,f_a)$, and let
$S$ be the set of $a$ with $c_a\neq0$ and $v_p(c_a)\leq D$. We show $\operatorname{good}_D(B,L)\leq|S|$. Let
$x_1,\dots,x_g,y_1,\dots,y_g\in L$ with $B(x_i,y_l)=p^D\delta_{il}$, and write $x_i=\sum_aX_{ia}f_a$ and
$y_l=\sum_aY_{la}f_a$ with $X_{ia},Y_{la}\in\mathbf Z_p$. By orthogonality and symmetry,
$\sum_aX_{ia}c_aY_{la}=p^D\delta_{il}$. Put $s_a=c_a/p^D$ for $a\in S$ and $s_a=0$ otherwise; for $a\notin S$, $c_a=0$
or $v_p(c_a)\geq D+1$, so $c_a=p^Ds_a+p^{D+1}u_a$ with $u_a\in\mathbf Z_p$ ($u_a=0$ for $a\in S$). Then
$X\operatorname{diag}(s)Y^t=1-pN$ with $N=X\operatorname{diag}(u)Y^t$ integral, and $\det(1-pN)\equiv1\pmod p$ is
nonzero. So $g$ is the rank of $X\operatorname{diag}(s)Y^t$, which is at most the rank of $\operatorname{diag}(s)$, the
number of $a$ with $s_a\neq0$, which is at most $|S|$.

\emph{Conclusion.} Let $e_1,\dots,e_{n_1}$ be the $f_a$ with $a\in S$. They are linearly independent, orthogonal,
with $B(e_a,e_a)\neq0$ of valuation at most $D$; $n_1=|S|\leq m=\operatorname{rk}L$; and
$\operatorname{bad}_D(B,L)=m-\operatorname{good}_D(B,L)\geq m-n_1$ by Step 3.
\end{proof}
% Lean: OrthLat is orthLat_of' (orthBasisLat_of freeLat_proved (orthSplit_of_min PI.exists_minimalNormVector))
%   (goodLeSmall_of diagonal_rank_bound). Step 1 is PI.exists_minimalNormVector = PI.LatMinWork.main
%   (exists_max_image over S x S, bound_span, norm_two_eq_one); Step 2 is orthBasisLat_of (strong induction
%   on the rank) with orthSplit_of (the submodule K = {u in L : B u x = 0}) and splitRank_proved, and
%   freeLat_proved when B vanishes; Step 3 is goodLeSmall_of (coordinates X, Y of a dual family, goodCount_mem)
%   with diagonal_rank_bound = DiagRank.diagRank (exists_int_mul, det_one_sub_p_ne, Matrix.rank_diagonal); the
%   conclusion is orthCount_of and orthLat_of. The Lean states it for submodules of polynomial rings over Q_p
%   inside a finite-dimensional subspace W.

\begin{lemma}\label{pd-kappa-defects-segment-mean-theorem-a-plus-orth}
Let an instance be given, let $C_M$ be a real number and put $D_1=D_1(C_M)$. Let $\delta$ be a digit class,
$\alpha\in A(\delta)$ and $(b,\omega)\in I(\delta,\alpha)$ (Definition~\ref{pd-kappa-defects-segment-mean-theorem-a-levels}).
Then there are $n_1\geq0$ and $e_1,\dots,e_{n_1}\in\Lambda_\omega$, linearly independent over $\mathbf Z_p$, with
$\beta^-_\omega(e_a,e_b)=0$ for $a\neq b$, $\beta^-_\omega(e_a,e_a)\neq0$ and
$v_p(\beta^-_\omega(e_a,e_a))\leq\lfloor D_1\rfloor$ for every $a$, and
$n_1\leq\operatorname{rk}\Lambda_\omega\leq n_1+\operatorname{bad}^-_\omega(\alpha,b;\lfloor D_1\rfloor)$
(Definition~\ref{pd-kappa-defects-segment-mean-theorem-a-plus-pairing}).
\end{lemma}
\leanref{Formalism/SharpRestriction/GrowthBounds/MonomialFactorials/BejjSmulLeft.lean}{Formalism.MainNotions.G.SymmetricFormBlocks.RestOrth}
\leanref{Formalism/SharpRestriction/GrowthBounds/MonomialFactorials/BePhiAddLeft.lean}{Formalism.MainNotions.G.SymmetricFormBlocks.restOrth_of}
\leanref{Formalism/SharpRestriction/GrowthBounds/MonomialFactorials/PhisBSelfAdj.lean}{Formalism.MainNotions.G.SymmetricFormBlocks.phisB_selfAdj}

\begin{proof}
Apply Lemma~\ref{pd-kappa-defects-segment-mean-theorem-a-plus-lattice} to $L=\Lambda_\omega$, inside the
finite-dimensional $\mathbf Q_p$-span of the products $y\xi^\beta$ of the proof of
Lemma~\ref{pd-kappa-defects-segment-mean-theorem-a-plus-rank}, $B=\beta^-_\omega$ and the integer
$D=\lfloor D_1\rfloor$; here $p\geq5$. The module
$\Lambda_\omega$ is finitely generated, as it lies in the finitely generated $\mathbf Z_p$-module of the polynomials of
bounded degree with coefficients in $\mathbf Z_p$. The map $\beta^-_\omega$ is linear in the first variable, and it is
symmetric: $\Phi^Y_\omega$ is the product of the operators $\phi$ of the ties, which act on disjoint sets of
variables and so commute, and each of them is self-adjoint for the Fock pairing, since multiplication by
$\sum_cw_ce_c$ and the operator $\sum_c\partial^2/\partial w_c\partial e_c$ are adjoint for it; a product of commuting
self-adjoint operators is self-adjoint, and the Fock pairing is symmetric.
\end{proof}
% Lean: restOrth_of CM hp2 OrthLat symmCondition restLjjFin_proved, with p != 2 from hiFacts (the book's
%   instance has p >= 5). symmCondition = restSymm_of commCondition: phisB_selfAdj (foldr_selfAdj with
%   GammaExpansion.fock_phi_symm, through fock_E_pow, and fock_comm) and commCondition (commute_tiePhi, the
%   ties at distinct positions use distinct letters); restLjjFin_proved = ljjFin_gen (Ljj_fg, Ljj_span_fd);
%   bePhi_smul_left and bePhi_add_left for linearity. RestOrth is stated for D = floor(D1r CM), for the
%   digit classes in clSet, aTop in cbox dl and (b, level) in idx dl aTop, with the pairing of gB aTop 0 b.

\begin{remark}\label{pd-kappa-defects-segment-mean-theorem-a-plus-orth-general}
The proof of Lemma~\ref{pd-kappa-defects-segment-mean-theorem-a-plus-orth} uses neither the digit class nor
the value of the integer $\lfloor D_1\rfloor$. For every $\alpha\in W^{S_j}$, every block $b$, every level $\omega$
of $b$ and every integer $D$, the same argument gives $n_1\geq0$ and $e_1,\dots,e_{n_1}\in\Lambda_\omega$ with the
properties of that lemma, with $D$ in place of $\lfloor D_1\rfloor$.
\end{remark}

\begin{lemma}\label{pd-kappa-defects-segment-mean-theorem-a-plus-window-level-bounds}
Let $p$ be a prime, $r\geq1$, and $w_1,\dots,w_r,e_1,\dots,e_r$ distinct variables of a polynomial ring over
$\mathbf Q_p$, with $E$ and $F$ as in
Lemma~\ref{pd-kappa-defects-segment-mean-theorem-a-block-bound-families-fock-raising} for $q=r$. A polynomial is
\emph{bounded by} a real number $B\geq0$ if all its coefficients have absolute value at most $B$. Let $u$ be bounded
by $B$, and let $L\geq0$ be an integer.
\begin{enumerate}
\item If $0\leq l\leq J<p^{L+1}$, then $\big((l!)^2\binom Jl\big)^{-1}E^lF^lu$ is bounded by $Bp^{(r+1)L}$.
\item Let $B=p^k$ with an integer $k\geq0$, and let $o,n,\zeta\geq0$ be integers with $o+n+r+\zeta<p^L$. Then
  $\phi(w,e;o,n,\zeta)u$ (Definition~\ref{pd-kappa-defects-segment-mean-theorem-a-levels}(3)) is bounded by
  $p^{k+(r+1)L}$.
\end{enumerate}
\end{lemma}
\leanref{Formalism/MahlerMeasure/MahlerMeasureProofs/PhiInt.lean}{Formalism.GM.M.PhiIntegrality.Bnd}
\leanref{Formalism/MahlerMeasure/IsIntegral.lean}{Formalism.GM.M.PhiInt}
\leanref{Formalism/MahlerMeasure/MahlerMeasureProofs/PhiInt.lean}{Formalism.GM.M.PhiIntegrality.contr_divDeriv}
\leanref{Formalism/MahlerMeasure/MahlerMeasureProofs/PhiInt.lean}{Formalism.GM.M.PhiIntegrality.contr_D}
\leanref{Formalism/MahlerMeasure/MahlerMeasureProofs/PhiInt.lean}{Formalism.GM.M.PhiIntegrality.contr_E}
\leanref{Formalism/MahlerMeasure/MahlerMeasureProofs/PhiInt.lean}{Formalism.GM.M.PhiIntegrality.contr_Fsum}
\leanref{Formalism/MahlerMeasure/MahlerMeasureProofs/PhiInt.lean}{Formalism.GM.projectorIntegral}
\leanref{Formalism/SharpRestriction/GrowthBounds/FockGramWindow/NormFockLe.lean}{Formalism.MainNotions.G.GramWindowAlt.bnd_Xl}
\leanref{Formalism/SharpRestriction/GrowthBounds/FockGramWindow/NormFockLe.lean}{Formalism.MainNotions.G.GramWindowAlt.bnd_phi}

\begin{proof}
For $j\geq k\geq0$ with $j<p^{L+1}$, $|\binom jk^{-1}|_p\leq p^L$, by
Lemma~\ref{pd-kappa-defects-segment-mean-theorem-a-block-bound-families-valuation-factorials}(1). Bounds are
preserved by sums, by the ultrametric inequality, and multiplied by $|a|_p$ under multiplication by a scalar $a$.

\emph{Elementary operators.} For a variable $v$ and $s\geq0$, $(1/s!)\partial^s/\partial v^s$ maps $x^\kappa$ to
$\binom{\kappa_v}sx^{\kappa-s\varepsilon_v}$, so it preserves the bound $B$. Hence so does
$(k!)^{-2}(\partial^2/\partial w_c\partial e_c)^k$, which is the composite of these operators for $w_c$ and for $e_c$, with
$s=k$. The operator $E$ multiplies by a polynomial with integer coefficients, so it preserves bounds.

\emph{Powers of $F$.} Let $j<p^{L+1}$. We show that $(j!)^{-2}F_S^j$ multiplies bounds by at most $p^{|S|L}$, where
$F_S=\sum_{c\in S}\partial^2/\partial w_c\partial e_c$ for a set $S$ of indices; this is clear for $S$ empty. For
$S=\{c\}\cup S'$ with $c\notin S'$, the two parts commute, and
\[
 \frac{F_S^j}{(j!)^2}=\sum_{m=0}^j\binom jm^{-1}\,\frac{(\partial^2/\partial w_c\partial e_c)^m}{(m!)^2}
 \cdot\frac{F_{S'}^{j-m}}{((j-m)!)^2},
\]
as $\binom jm/(j!)^2=\binom jm^{-1}(m!)^{-2}((j-m)!)^{-2}$. By induction each term multiplies bounds by at most
$p^L\cdot1\cdot p^{|S'|L}$. With $S$ all of $\{1,\dots,r\}$, $(j!)^{-2}F^j$ multiplies bounds by at most $p^{rL}$.

(1) Write $\big((l!)^2\binom Jl\big)^{-1}E^lF^l=\binom Jl^{-1}E^l\circ(l!)^{-2}F^l$, and use $l\leq J<p^{L+1}$.

(2) Put $A=o+n+r+\zeta-1$ and $L_0=\lfloor\log_p(o+n+r+\zeta)\rfloor$, so that $A<p^{L_0+1}$; as
$o+n+r+\zeta<p^L$, $L_0+1\leq L$. By Definition~\ref{pd-kappa-defects-segment-mean-theorem-a-levels}(3),
$\phi(w,e;o,n,\zeta)=\sum_{j=0}^o(-1)^j\binom Aj^{-1}E^j\circ(j!)^{-2}F^j$, with $j\leq o\leq A$. By the bounds above
with $L_0$ in place of $L$, each term multiplies bounds by at most $p^{L_0}\cdot1\cdot p^{rL_0}$, and so does
$\phi(w,e;o,n,\zeta)$. So if $u'$ has coefficients in $\mathbf Z_p$, then $p^{(r+1)(L_0+1)}\phi(w,e;o,n,\zeta)u'$
has coefficients in $\mathbf Z_p$. As $u$ is bounded by $p^k$, $u'=p^ku$ has coefficients in $\mathbf Z_p$, and
$p^{(r+1)(L_0+1)}\phi(w,e;o,n,\zeta)u'=p^{k+(r+1)(L_0+1)}\phi(w,e;o,n,\zeta)u$. Hence $\phi(w,e;o,n,\zeta)u$ is bounded
by $p^{k+(r+1)(L_0+1)}\leq p^{k+(r+1)L}$.
\end{proof}
% Lean: norm_choose_inv_le (Nat.factorization_choose_le_log), bnd_sum, bnd_smul and bnd_linear; contr_divDeriv
%   (pderiv_pow_monomial), contr_D and contr_E; contr_Fsum (Finset induction, Commute.add_pow,
%   inv_factorial_sq_mul_choose_eq); (1) is bnd_Xl (with Mc = p^(Nat.log p J)); (2) is bnd_phi, for B = p^k and
%   Lam = L >= Nat.log p (o + n + r + z) + 1: it applies PhiInt = projectorIntegral = phi_integrality to p^k u
%   (bnd_pow_iff). phi_integrality shows Contr (p^((r+1) L0)) for tiePhi with L0 = Nat.log p (o + n + r + z) (the
%   book's multiplication of bounds, Remark -bounds-sharp) and exports the integrality of p^((r+1)(L0+1)) tiePhi u
%   for integral u. The hypotheses that w, e are injective with distinct values and r >= 1 are not used.

\begin{remark}\label{pd-kappa-defects-segment-mean-theorem-a-plus-window-level-bounds-sharp}
Let $p$, $r$, the variables $w_c,e_c$ and $E$, $F$ be as in
Lemma~\ref{pd-kappa-defects-segment-mean-theorem-a-plus-window-level-bounds}, let $L,o,n,\zeta\geq0$ be integers with
$o+n+r+\zeta<p^{L+1}$, and let $u$ be bounded by a real number $B\geq0$. The proof of part~(2) of that lemma, with $L$
in place of $L_0$, shows that $\phi(w,e;o,n,\zeta)u$ is bounded by $Bp^{(r+1)L}$. For $B=p^k$ and
$L=\lfloor\log_p(o+n+r+\zeta)\rfloor$ this is smaller by the factor $p^{r+1}$ than the bound of part~(2) at the least
admissible exponent, $L+1$.
\end{remark}
\leanref{Formalism/MahlerMeasure/MahlerMeasureProofs/PhiInt.lean}{Formalism.GM.M.PhiIntegrality.phi_integrality}
% Lean: the internal step hPhi of phi_integrality, Contr (Mc^(r+1)) for tiePhi with Mc = p^(Nat.log p (o + n + r + z));
%   it is not exported, and the window bound uses part (2).

\begin{lemma}\label{pd-kappa-defects-segment-mean-theorem-a-plus-window-level-pair}
Let $p$, $r$, the variables $w_c,e_c$, $E$ and $F$ be as in
Lemma~\ref{pd-kappa-defects-segment-mean-theorem-a-plus-window-level-bounds}, with $D_w$ and $D_e$ as in
Lemma~\ref{pd-kappa-defects-segment-mean-theorem-a-block-bound-families-fock-raising}. Let $o,n\geq0$ be integers,
$J=\min(o,n)$ and $q=o+n+r-1$, and let $v$ be a polynomial with $D_wv=ov$ and $D_ev=nv$. For a polynomial $u$ and
$0\leq l\leq J$ put $\mu_l(u)=[E^lF^lu,v]/\big((l!)^2\binom Jl\big)$, and for $0\leq i\leq J$ put
$c(i)=\sum_{l=0}^{J-i}(-1)^l\binom{q-l}{J-l-i}\mu_l(u)$. Then $c(J)=[u,v]$ and, for every integer $\zeta\geq0$,
\[
 \binom{q+\zeta}J[\phi(w,e;o,n,\zeta)u,v]=\sum_{l=0}^J(-1)^l\binom{q+\zeta-l}{J-l}\mu_l(u)=\sum_{i=0}^Jc(i)\binom\zeta i .
\]
\end{lemma}
\leanref{Formalism/SharpRestriction/GrowthBounds/FockGramWindow/PhiPair.lean}{Formalism.MainNotions.G.GramWindowAlt.mX}
\leanref{Formalism/SharpRestriction/GrowthBounds/FockGramWindow/PhiPair.lean}{Formalism.MainNotions.G.GramWindowAlt.phi_pair}
\leanref{Formalism/SharpRestriction/GrowthBounds/FockGramWindow/PhiPair.lean}{Formalism.MainNotions.G.GramWindowAlt.cW}
\leanref{Formalism/SharpRestriction/GrowthBounds/FockGramWindow/PhiPair.lean}{Formalism.MainNotions.G.GramWindowAlt.cW_sum}
\leanref{Formalism/SharpRestriction/GrowthBounds/FockGramWindow/PhiPair.lean}{Formalism.MainNotions.G.GramWindowAlt.cW_top}

\begin{proof}
Put $A=q+\zeta$, so $A\geq J$ as $r\geq1$. By
Definition~\ref{pd-kappa-defects-segment-mean-theorem-a-levels}(3), $[\phi(w,e;o,n,\zeta)u,v]$ is the sum over
$0\leq l\leq o$ of $(-1)^l\big((l!)^2\binom Al\big)^{-1}[E^lF^lu,v]$. Each term of $F$ lowers the degrees in the $w_c$
and in the $e_c$ by one, so $F^lv=0$ for $l>J$, and then $[E^lF^lu,v]=[F^lu,F^lv]=0$ by
Lemma~\ref{pd-kappa-defects-segment-mean-theorem-a-block-bound-families-fock-raising}(1). For $l\leq J$,
$\binom AJ\binom Jl=\binom Al\binom{A-l}{J-l}$, which gives the first identity. For the second, Vandermonde's
identity $\binom{q+\zeta-l}{J-l}=\sum_{i=0}^{J-l}\binom\zeta i\binom{q-l}{J-l-i}$ and an exchange of the sums give
$\sum_ic(i)\binom\zeta i$. For $i=J$ only $l=0$ contributes, and $\mu_0(u)=[u,v]$.
\end{proof}
% Lean: phi_pair (with InvariantFactorPaths.KP.LL.tieF_pow_apply_eq_zero_of_min_lt for F^l v = 0 and Nat.choose_mul),
%   cW_sum (Nat.add_choose_eq) and cW_top. The book's mu_l(u) is mX w e J l u v.

\begin{lemma}\label{pd-kappa-defects-segment-mean-theorem-a-plus-window-level-lead}
Let $p$ be a prime, $n_1,J\geq0$ integers and $R$ a real number. For $1\leq a,b\leq n_1$ and $0\leq i\leq J$ let
$c_{ab}(i)\in\mathbf Q_p$, with $c_{ab}(J)=\delta_{ab}d_a$, where $d_a\neq0$ and $v_p(d_a)\leq R$ for every $a$. Let $M$
be the $n_1\times n_1$ matrix with entries $M_{ab}(X)=\sum_{i=0}^Jc_{ab}(i)\binom Xi$ in $\mathbf Q_p[X]$, where
$\binom Xi=X(X-1)\cdots(X-i+1)/i!$. Then $\det M\neq0$; it has degree $n_1J$ and leading coefficient
$\lambda=\prod_ad_a/(J!)^{n_1}$; and $v_p(\lambda)+n_1J/(p-1)\leq n_1(R+\Lambda)$ for every real
$\Lambda\geq J/(p-1)-v_p(J!)$. If $J\geq1$, then $J/(p-1)-v_p(J!)\leq\lfloor\log_pJ\rfloor+1$.
\end{lemma}
\leanref{Formalism/SharpRestriction/GrowthBounds/FockGramWindow/BinMatNatDegreeLe.lean}{Formalism.MainNotions.G.GramWindowAlt.lead_bound}
\leanref{Formalism/PolynomialDimension/ChangeLatticeGoodRank.lean}{Formalism.CountingSumsFrames.Top.detLeadDiag}
\leanref{Formalism/MahlerMeasure/MahlerMeasureProofs/MeanGam.lean}{Formalism.GM.MeanGammaLemmas.legendre'}

\begin{proof}
The polynomial $\binom Xi$ has degree $i$ and leading coefficient $1/i!$, so every $M_{ab}$ has degree at most $J$, with
coefficient $\delta_{ab}d_a/J!$ at $X^J$. In the expansion of $\det M$ over permutations, each product of $n_1$
entries has degree at most $n_1J$, and its coefficient at $X^{n_1J}$ is the product of the coefficients at $X^J$,
which is $0$ unless the permutation is the identity. So the coefficient of $\det M$ at $X^{n_1J}$ is
$\lambda=\prod_ad_a/(J!)^{n_1}\neq0$, which gives the first two claims. Then
$v_p(\lambda)+n_1J/(p-1)=\sum_av_p(d_a)+n_1\big(J/(p-1)-v_p(J!)\big)\leq n_1(R+\Lambda)$. For the last claim, by Legendre's
formula $(p-1)v_p(J!)=J-s$, where $s$ is the sum of the digits of $J$ in base $p$; there are $\lfloor\log_pJ\rfloor+1$
digits, each at most $p-1$, so $J/(p-1)-v_p(J!)=s/(p-1)\leq\lfloor\log_pJ\rfloor+1$.
\end{proof}
% Lean: lead_bound with binK_natDegree, binK_coeff_top, binMat_natDegree_le, binMat_coeff_top and detLeadDiag =
%   det_leading_diagonal (Matrix.det_apply', Polynomial.coeff_prod_of_natDegree_le); legendre' uses
%   sub_one_mul_padicValNat_factorial, Nat.digit_sum_le and Nat.digits_len. The Lean states the bound with
%   -log_p of the norm of the leading coefficient, which is its valuation.

\begin{lemma}\label{pd-kappa-defects-segment-mean-theorem-a-plus-window-level-smith-gauss}
Let $K$ be a field with an ultrametric absolute value $|\ |$, and $\rho\geq0$ a real number. For
$P=\sum_ka_kX^k\in K[X]$ let $|P|_\rho=\max_k|a_k|\rho^k$ be its \emph{Gauss norm} of radius $\rho$, with $|0|_\rho=0$.
\begin{enumerate}
\item For $P,Q\in K[X]$ and $a\in K$, $|P+Q|_\rho\leq\max(|P|_\rho,|Q|_\rho)$, $|PQ|_\rho\leq|P|_\rho|Q|_\rho$ and
  $|a|_\rho=|a|$.
\item If $\rho>0$, then $|PQ|_\rho=|P|_\rho|Q|_\rho$ for all $P,Q\in K[X]$.
\item Let $A$ be an $n\times n$ matrix over $K[X]$ and $b_1,\dots,b_n\geq0$ real numbers with $|A_{ab}|_\rho\leq b_a$
  for all $a,b$. Then $|\det A|_\rho\leq b_1\cdots b_n$.
\item If $\rho\leq1$ and $|P|_1\leq1$, then $|P|_\rho\leq\max(|P(0)|,\rho)$.
\item For $q,a\in K$, $|qX+a|_\rho=\max(|a|,|q|\rho)$.
\item Let $U$ be an $n\times n$ matrix over $K$ whose entries have absolute value at most $1$, with $|\det U|=1$. Then
  $U$ is invertible, and the entries of $U^{-1}$ have absolute value at most $1$.
\end{enumerate}
\end{lemma}
\leanref{Formalism/MahlerMeasure/StatementD.lean}{Formalism.GM.gaussR}
\leanref{Formalism/MahlerMeasure/StatementD.lean}{Formalism.GM.GaussMulRad}
\leanref{Formalism/MahlerMeasure/MahlerMeasureProofs/GaussRad.lean}{Formalism.GM.gauss_norm_multiplicative}
\leanref{Formalism/MahlerMeasure/MahlerMeasureProofs/LieMat.lean}{Formalism.GM.DetBound.gaussR_det_le}
\leanref{Formalism/MahlerMeasure/MahlerMeasureProofs/LieMat.lean}{Formalism.GM.DetBound.gaussR_le_max_coeff_zero}
\leanref{Formalism/MahlerMeasure/MahlerMeasureProofs/LieRoots.lean}{Formalism.GM.GaussNormRoots.gaussR_lin}
\leanref{Formalism/MahlerMeasure/MahlerMeasureProofs/LieMat.lean}{Formalism.GM.DetBound.norm_inv_entry_le}

\begin{proof}
Write $Q=\sum_kb_kX^k$.

(1) The coefficients of $P+Q$ are the sums of those of $P$ and $Q$. The coefficient of $X^k$ in $PQ$ is
$\sum_{i+j=k}a_ib_j$, and $|a_ib_j|\rho^k=|a_i|\rho^i\cdot|b_j|\rho^j\leq|P|_\rho|Q|_\rho$. The ultrametric inequality
gives both bounds. A constant $a$ has the single coefficient $a$, at $X^0$.

(2) By (1) it suffices to show $|PQ|_\rho\geq|P|_\rho|Q|_\rho$, and we may assume $P,Q\neq0$, so that
$|P|_\rho,|Q|_\rho>0$. Let $i_0$ be the least index with $|a_{i_0}|\rho^{i_0}=|P|_\rho$, and $j_0$ the least with
$|b_{j_0}|\rho^{j_0}=|Q|_\rho$. The coefficient of $X^{i_0+j_0}$ in $PQ$ is $a_{i_0}b_{j_0}$ plus the $a_ib_j$ with
$i+j=i_0+j_0$ and $(i,j)\neq(i_0,j_0)$. For these $i<i_0$ or $j<j_0$, so
$|a_ib_j|\rho^{i_0+j_0}<|P|_\rho|Q|_\rho=|a_{i_0}b_{j_0}|\rho^{i_0+j_0}$. By the ultrametric inequality the coefficient has
absolute value $|a_{i_0}b_{j_0}|$, and so $|PQ|_\rho\geq|P|_\rho|Q|_\rho$. For $\rho=1$ this is \cite[\S5.1.2]{BGR84}.

(3) The determinant is $\sum_\sigma\operatorname{sgn}(\sigma)\prod_aA_{\sigma(a)a}$ over the permutations $\sigma$, and by
(1) each term has Gauss norm at most $\prod_ab_{\sigma(a)}=\prod_ab_a$.

(4) The coefficient at $X^0$ contributes $|P(0)|$. For $k\geq1$, $|a_k|\leq|P|_1\leq1$ and $\rho^k\leq\rho$.

(5) The polynomial has the coefficients $a$ at $X^0$ and $q$ at $X^1$.

(6) As $\det U\neq0$, $U^{-1}=(\det U)^{-1}\operatorname{adj}(U)$. The entry of $\operatorname{adj}(U)$ at $(a,b)$ is the
determinant of $U$ with its $b$-th row replaced by the $a$-th unit vector, a matrix whose entries have absolute value
at most $1$; by (3), for constant polynomials with $b_a=1$, it has absolute value at most $1$. As $|\det U|=1$, so do
the entries of $U^{-1}$.
\end{proof}
% Lean: (1) is gaussR_sum_le (nnnorm_sum_mul_le), gaussR_mul_le and gaussR_C; (2) is gauss_norm_multiplicative =
%   le_antisymm gaussR_mul_le le_gaussR_mul, with gaussR_attained_least for i0, j0 and
%   IsUltrametricDist.nnnorm_add_eq_max_of_nnnorm_ne_nnnorm; (3) is gaussR_det_le (Matrix.det_apply,
%   gaussR_prod_le, gaussR_neg, Equiv.prod_comp); (4) is gaussR_le_max_coeff_zero; (5) is gaussR_lin; (6) is
%   norm_inv_entry_le (Matrix.inv_def, Matrix.adjugate_apply, norm_det_le_one). The Gauss norm is gaussR rho,
%   with rho in NNReal.

\begin{lemma}\label{pd-kappa-defects-segment-mean-theorem-a-plus-window-level-smith-binomial}
Let $p$ be a prime, $h\geq0$ and $0\leq k\leq p^h$ integers, and $t\in\mathbf Z$. Then the coefficients of
\[
 \binom{p^hX+t}k=\frac1{k!}\prod_{i=0}^{k-1}(p^hX+t-i)\in\mathbf Q[X]
\]
have $p$-adic absolute value at most $1$.
\end{lemma}
\leanref{Formalism/MahlerMeasure/StatementB.lean}{Formalism.GM.MahlerDisc}
\leanref{Formalism/MahlerMeasure/MahlerMeasureProofs/Mahler.lean}{Formalism.GM.mahler_disc_bound}
\leanref{Formalism/MahlerMeasure/MahlerMeasureProofs/Mahler.lean}{Formalism.GM.bB_mem}
\leanref{Formalism/MahlerMeasure/StatementD.lean}{Formalism.GM.MahlerK}
\leanref{Formalism/MahlerMeasure/MahlerMeasureProofs/MahlerK.lean}{Formalism.GM.mahler_binomial_unit_ball}

\begin{proof}
Write $B_k(t)$ for this polynomial, and call a rational polynomial integral if its coefficients have $p$-adic absolute
value at most $1$; integral polynomials form a ring. For a polynomial $Q$,
$\binom{Q+1}{k+1}=\binom Q{k+1}+\binom Qk$ and $\binom Q{k+1}=\frac Q{k+1}\binom{Q-1}k$. So
$B_{k+1}(t+1)=B_{k+1}(t)+B_k(t)$ and $B_{k+1}(0)=\frac{p^h}{k+1}X\cdot B_k(-1)$. We argue by induction on $k$, the
case $k=0$ being $B_0(t)=1$. Let $k+1\leq p^h$, and suppose that $B_k(t)$ is integral for every $t$. As
$p^{v_p(k+1)}\leq k+1\leq p^h$, $p^h/(k+1)$ has absolute value at most $1$, so $B_{k+1}(0)$ is integral. The identities
$B_{k+1}(t+1)=B_{k+1}(t)+B_k(t)$ and $B_{k+1}(t-1)=B_{k+1}(t)-B_k(t-1)$ then give, by induction on $|t|$, that
$B_{k+1}(t)$ is integral for every $t\in\mathbf Z$.
\end{proof}
% Lean: mahler_disc_bound = bB_mem, with bB c i t = C (1/i!) * (descPochhammer i).comp (C c * X + C t), bB_zero,
%   bB_pascal (descPochhammer_comp_add_one), bB_succ_zero and div_mem_zpS, by Int.induction_on in t; zpT is the
%   ring of integral polynomials. MahlerK (mahler_binomial_unit_ball) transfers the bound to binK K k composed with
%   p^h X + t0 over a normed algebra K over Q_p (norm_algebraMap').

\begin{lemma}\label{pd-kappa-defects-segment-mean-theorem-a-plus-window-level-smith-point}
Let $p$ be a prime and $K$ a field containing $\mathbf Q_p$, with an ultrametric absolute value $|\ |$ that extends the
$p$-adic absolute value; for an integer $h\geq0$ and $x\in K$ put $\nu_h(x)=-\log_p\max(|x|,p^{-h})$. Let $n,J,h\geq0$ be
integers with $J\leq p^h$, $E\geq0$ a real number, and $c_{ab}(i)\in K$ with $|c_{ab}(i)|\leq1$ for $1\leq a,b\leq n$ and
$0\leq i\leq J$, and let $M$ be the $n\times n$ matrix over $K[X]$ with entries $M_{ab}(X)=\sum_{i=0}^Jc_{ab}(i)\binom Xi$.
Suppose that $\det M=\lambda\prod_{\zeta\in R}(X-\zeta)$, with $\lambda\in K$ nonzero and $R$ a finite multiset of elements
of $K$. Let $t\in\mathbf Z$, let $U$ and $V$ be $n\times n$ matrices over $K$ whose entries have absolute value at most
$1$, with $|\det U|=|\det V|=1$, and let $d_1,\dots,d_n\in K$ with $M(t)=U\operatorname{diag}(d_1,\dots,d_n)V$. Let $\beta$
be the number of $a$ with $|d_a|\leq p^{-E}$. Then, with the sum and the count over $R$ taken with multiplicities,
\[
 \beta E\leq-\log_p|\lambda|+\sum_{\zeta\in R}\nu_h(t-\zeta)+E\,\big|\{\zeta\in R:|t-\zeta|\leq p^{-h}\}\big|.
\]
\end{lemma}
\leanref{Formalism/MahlerMeasure/Statement.lean}{Formalism.GM.vmin}
\leanref{Formalism/MahlerMeasure/StatementD.lean}{Formalism.GM.binMatK}
\leanref{Formalism/MahlerMeasure/StatementD.lean}{Formalism.GM.ValuationWindowLowerBound}
\leanref{Formalism/MahlerMeasure/FinalD.lean}{Formalism.GM.Top.valuation_window_lower_bound}
\leanref{Formalism/MahlerMeasure/LieRoots.lean}{Formalism.GM.LieMat}
\leanref{Formalism/MahlerMeasure/MahlerMeasureProofs/LieMat.lean}{Formalism.GM.lieMat_of}
\leanref{Formalism/MahlerMeasure/LieRoots.lean}{Formalism.GM.LieRoots}
\leanref{Formalism/MahlerMeasure/MahlerMeasureProofs/LieRoots.lean}{Formalism.GM.lieRoots_of}

\begin{proof}
Put $\rho=p^{-E}$, so $\rho\leq1$, and $L=p^hX+t$. Gauss norms are those of
Lemma~\ref{pd-kappa-defects-segment-mean-theorem-a-plus-window-level-smith-gauss}.

\emph{An upper bound.} We show $|(\det M)(L)|_\rho\leq\rho^\beta$. Let $B$ be the matrix with entries $M_{ab}(L)$, so that
$\det B=(\det M)(L)$. Its entries are the sums of the $c_{ab}(i)\binom{p^hX+t}i$, and $i\leq J\leq p^h$; by
Lemma~\ref{pd-kappa-defects-segment-mean-theorem-a-plus-window-level-smith-binomial}, whose coefficients keep their
absolute values in $K$, and by part~(1) of the Gauss norm lemma, $|B_{ab}|_1\leq1$. The constant terms are
$B(0)=M(t)=U\operatorname{diag}(d)V$. By part~(6), $U$ and $V$ are invertible with $U^{-1}$ and $V^{-1}$ of entries of
absolute value at most $1$. Put $B'=U^{-1}BV^{-1}$. Then $B'(0)=\operatorname{diag}(d)$, and $|B'_{ab}|_1\leq1$ by
part~(1). By part~(4), $|B'_{ab}|_\rho\leq\max(\delta_{ab}|d_a|,\rho)$. If $|d_a|\leq\rho$ this is at most $\rho$, and
otherwise it is at most $1$, as $|d_a|=|B'_{aa}(0)|\leq|B'_{aa}|_1\leq1$. By part~(3), $|\det B'|_\rho\leq\rho^\beta$. As
$\det B=\det U\det B'\det V$ and $|\det U|=|\det V|=1$, part~(1) gives $|\det B|_\rho\leq\rho^\beta$.

\emph{The roots.} Composing with $L$, $(\det M)(L)=\lambda\prod_{\zeta\in R}\big(p^hX+(t-\zeta)\big)$. As $|p^h|=p^{-h}$ in
$K$, parts (2) and (5) give
\[
 |(\det M)(L)|_\rho=|\lambda|\prod_{\zeta\in R}\max\big(|t-\zeta|,p^{-h-E}\big)\leq p^{-E\beta}.
\]
All factors are positive; taking $\log_p$,
$\log_p|\lambda|+\sum_\zeta\log_p\max(|t-\zeta|,p^{-h-E})\leq-E\beta$. For a real $u\geq0$,
$-\log_p\max(u,p^{-h-E})\leq-\log_p\max(u,p^{-h})+E$ if $u\leq p^{-h}$, as the left side is at most $h+E$ and the
right side is $h+E$; and the two sides are equal without the $E$ if $u>p^{-h}$, as then both maxima are $u$. Summing
over $\zeta$ with $u=|t-\zeta|$ gives the claim.
\end{proof}
% Lean: valuation_window_lower_bound = lieRoots (lieMat ...) with lieMat = lieMat_of mahler_binomial_unit_ball and
%   lieRoots = lieRoots_of gauss_norm_multiplicative. The upper bound is lieMat_of: hBint (binMatK, MahlerK),
%   hB0, B' = U^-1 B V^-1 (hB'0, hB'int with norm_inv_entry_le), hrow (gaussR_le_max_coeff_zero), hdetB'
%   (gaussR_det_le with b i = rho or 1), hBB' and hdetBB. The roots are lieRoots_of: hcomp
%   (multiset_prod_comp), hreal (GaussMulRad, gaussR_multiset_prod, gaussR_lin, hnph), hlog
%   (logb_multiset_prod) and vmin_step, summed with multiset_sum_le_sum.

\begin{lemma}\label{pd-kappa-defects-segment-mean-theorem-a-plus-window-level-smith-count}
Let $p$ be a prime and $K$ a field containing $\mathbf Q_p$, with an ultrametric absolute value that extends the $p$-adic
absolute value, and let $\nu_h$ be as in
Lemma~\ref{pd-kappa-defects-segment-mean-theorem-a-plus-window-level-smith-point}. Let $t_0,d\in\mathbf Z$ with $d$ prime
to $p$, put $t_j=t_0+dj$, and let $Y_0,h\geq0$ be integers.
\begin{enumerate}
\item For $\zeta\in K$, $\sum_{j=0}^{Y_0-1}\nu_h(t_j-\zeta)\leq\sum_{i=1}^h\big(Y_0/p^i+1\big)$.
\item If $Y_0\leq p^h$ and $R$ is a finite multiset of elements of $K$, then
  $\sum_{j=0}^{Y_0-1}\big|\{\zeta\in R:|t_j-\zeta|\leq p^{-h}\}\big|\leq|R|$, the counts taken with multiplicities.
\end{enumerate}
\end{lemma}
\leanref{Formalism/MahlerMeasure/QuotientUpperBoundsOfPowerBounds.lean}{Formalism.GM.WindowCountAP}
\leanref{Formalism/MahlerMeasure/MahlerMeasureProofs/WindowAP.lean}{Formalism.GM.window_count_progression}
\leanref{Formalism/MahlerMeasure/QuotientUpperBoundsOfPowerBounds.lean}{Formalism.GM.DiscsDisjointAP}
\leanref{Formalism/MahlerMeasure/MahlerMeasureProofs/WindowAP.lean}{Formalism.GM.root_discs_disjoint_progression}

\begin{proof}
\emph{Integers close to each other.} Let $i\geq0$ and $j,j'\geq0$ with $|t_j-t_{j'}|<p^{1-i}$. The absolute value in
$K$ of the integer $t_j-t_{j'}=d(j-j')$ is its $p$-adic absolute value, which lies in $p^{\mathbf Z}\cup\{0\}$, so it is at
most $p^{-i}$. So $p^i$ divides $d(j-j')$, and, as $d$ is prime to $p$, $p^i$ divides $j-j'$.

(1) First, $\nu_h(x)\leq\big|\{1\leq i\leq h:|x|<p^{1-i}\}\big|$ for $x\in K$. Indeed $\nu_h(x)\leq h$, as
$\max(|x|,p^{-h})\geq p^{-h}$. For a real $v$, $\min(v,h)\leq|\{1\leq i\leq h:i-1<v\}|$, by induction on $h$. And if
$i-1<\nu_h(x)$, then $|x|\leq\max(|x|,p^{-h})=p^{-\nu_h(x)}<p^{1-i}$. Summing over $j$ and exchanging the sums, it
suffices to show that, for $1\leq i\leq h$, at most $Y_0/p^i+1$ of the $j<Y_0$ have $|t_j-\zeta|<p^{1-i}$. For two
such $j,j'$, $|t_j-t_{j'}|\leq\max(|t_j-\zeta|,|t_{j'}-\zeta|)<p^{1-i}$, so $p^i$ divides $j-j'$ by the first paragraph.
So these $j$ lie in one residue class modulo $p^i$, and $j\mapsto\lfloor j/p^i\rfloor$ is injective on them, with
values in $\{0,\dots,\lfloor Y_0/p^i\rfloor\}$; there are at most $\lfloor Y_0/p^i\rfloor+1\leq Y_0/p^i+1$ of them.

(2) For $\zeta\in K$ there is at most one $j<Y_0$ with $|t_j-\zeta|\leq p^{-h}$. For two such $j,j'$,
$|t_j-t_{j'}|\leq p^{-h}<p^{1-h}$, so $p^h$ divides $j-j'$ by the first paragraph; as $|j-j'|<Y_0\leq p^h$, $j=j'$. So
each element of $R$ is counted for at most one $j$, and the sum is at most $|R|$, by induction on $R$.
\end{proof}
% Lean: (1) is window_count_progression: vmin_le (count, by induction on h) for each j, Finset.sum_comm, and card_le
%   (card_le_div_add_one_of_subset_range_of_mod_eq, with norm_sub_le_max, progr_sub, dvd_of_norm_lt and
%   dvd_cancel). (2) is root_discs_disjoint_progression: key (induction on the multiset) with huniq.

\begin{lemma}\label{pd-kappa-defects-segment-mean-theorem-a-plus-window-level-smith-mean}
Let $p$ be a prime and $K$ a field containing $\mathbf Q_p$, with an ultrametric absolute value that extends the $p$-adic
absolute value, and let $\nu_h$ be as in
Lemma~\ref{pd-kappa-defects-segment-mean-theorem-a-plus-window-level-smith-point}. Let $R$ be a finite multiset of
elements of $K$, $\lambda$, $\Lambda_0$ and $E>0$ real numbers, and $n_0,J\geq0$ integers with $|R|\leq Jn_0$ and
$\lambda+|R|/(p-1)\leq n_0\Lambda_0$. Let $t_0,d\in\mathbf Z$ with $d$ prime to $p$, put $t_j=t_0+dj$, and let $Y_0\geq1$
and $h\geq0$ be integers with $Y_0\leq p^h$. Let $\beta_0,\dots,\beta_{Y_0-1}\geq0$ be integers such that, for every
$j<Y_0$, with the sums and counts over $R$ taken with multiplicities,
\[
 \beta_jE\leq\lambda+\sum_{\zeta\in R}\nu_h(t_j-\zeta)+E\,\big|\{\zeta\in R:|t_j-\zeta|\leq p^{-h}\}\big|.
\]
Then $\frac1{Y_0}\sum_{j=0}^{Y_0-1}\beta_j\leq n_0\big(J/Y_0+(\Lambda_0+Jh/Y_0)/E\big)$.
\end{lemma}
\leanref{Formalism/MahlerMeasure/QuotientUpperBoundsOfPowerBounds.lean}{Formalism.GM.WindowMeanAP}
\leanref{Formalism/MahlerMeasure/MahlerMeasureProofs/WindowMeanAP.lean}{Formalism.GM.windowMeanAP_of}
\leanref{Formalism/MahlerMeasure/CoarseLevelSlackInequalities.lean}{Formalism.GM.Top.windowMeanAP}

\begin{proof}
Summing the hypothesis over $j$ and exchanging the sums over $j$ and over $R$,
\[
 E\sum_j\beta_j\leq Y_0\lambda+\sum_{\zeta\in R}\sum_j\nu_h(t_j-\zeta)+E\sum_j\big|\{\zeta\in R:|t_j-\zeta|\leq p^{-h}\}\big|.
\]
By Lemma~\ref{pd-kappa-defects-segment-mean-theorem-a-plus-window-level-smith-count}(1) and
$\sum_{i=1}^hp^{-i}=(1-p^{-h})/(p-1)\leq1/(p-1)$, each inner sum over $j$ in the middle term is at most
$Y_0/(p-1)+h$. By part~(2) of that lemma the last term is at most $E|R|$. So
\[
 E\sum_j\beta_j\leq Y_0\Big(\lambda+\frac{|R|}{p-1}\Big)+|R|h+E|R|\leq Y_0n_0\Lambda_0+Jn_0h+EJn_0,
\]
using $Y_0>0$, $h\geq0$, $E>0$ and the hypotheses on $|R|$. Dividing by $EY_0$ gives the claim.
\end{proof}
% Lean: windowMeanAP_of windowCountAP discsDisjointAP: hsum (Finset.sum_le_sum of the hypothesis), swap, hV
%   (window_count_progression with sumS = geom), hDD (root_discs_disjoint_progression), and the final
%   field_simp. The Lean takes beta : Z -> N and the hypothesis at beta (progr t0 d j); the book's beta_j is
%   beta (t_j).

\begin{lemma}\label{pd-kappa-defects-segment-mean-theorem-a-plus-window-level-smith}
Let $p$ be a prime, $n_1,J,h\geq0$, $t_0\geq0$ and $Y_0\geq1$ integers with $Y_0\leq p^h$ and $J+1\leq p^h$, and $E>0$
and $\Lambda_0$ real numbers. For $1\leq a,b\leq n_1$ and $0\leq i\leq J$ let $c_{ab}(i)\in\mathbf Z_p$, and let $M$ be
the matrix of polynomials $M_{ab}(X)=\sum_{i=0}^Jc_{ab}(i)\binom Xi$, as in
Lemma~\ref{pd-kappa-defects-segment-mean-theorem-a-plus-window-level-lead}. Suppose that $\det M\neq0$ and that its
leading coefficient $\lambda$ satisfies $v_p(\lambda)+\deg(\det M)/(p-1)\leq n_1\Lambda_0$. For each integer $t\geq0$
the matrix $M(t)$ has entries in $\mathbf Z_p$; let $U_t,V_t\in GL_{n_1}(\mathbf Z_p)$ and $d_1(t),\dots,d_{n_1}(t)\in\mathbf Z_p$
with $M(t)=U_t\operatorname{diag}(d_1(t),\dots,d_{n_1}(t))V_t$. Then
\[
 \sum_{t=t_0}^{t_0+Y_0-1}\big|\{a:|d_a(t)|_p\leq p^{-E}\}\big|\leq Y_0n_1\Big(\frac J{Y_0}+\frac{\Lambda_0+Jh/Y_0}E\Big).
\]
\end{lemma}
\leanref{Formalism/MahlerMeasure/StatementD.lean}{Formalism.GM.WindowFull}
\leanref{Formalism/MahlerMeasure/QuotientUpperBoundsOfPowerBounds.lean}{Formalism.GM.WindowFullAP}
\leanref{Formalism/MahlerMeasure/CoarseLevelSlackInequalities.lean}{Formalism.GM.Top.windowFullAP}
\leanref{Formalism/SharpRestriction/GrowthBounds/FockGramWindow/GramWindow2Proved.lean}{Formalism.MainNotions.G.GramWindowAlt.gramWindow2_proved}

\begin{proof}
Let $\bar{\mathbf Q}_p$ be an algebraic closure of $\mathbf Q_p$, with the absolute value extending the $p$-adic one
\cite[Chapter~II, Theorem~4.8]{Neu99}; it is ultrametric. Regard the $c_{ab}(i)$ and $M$ over $\bar{\mathbf Q}_p$. As
$\det M\neq0$, it splits: $\det M=\lambda\prod_{\zeta\in R}(X-\zeta)$ with $R$ a multiset of $\deg(\det M)$ elements of
$\bar{\mathbf Q}_p$, and $-\log_p|\lambda|=v_p(\lambda)$. Each entry of $M$ has degree at most $J$, so, expanding the
determinant, $|R|=\deg(\det M)\leq n_1J$. The hypothesis on $\lambda$ reads
$-\log_p|\lambda|+|R|/(p-1)\leq n_1\Lambda_0$.

Let $0\leq j<Y_0$ and $t=t_0+j$. The entries of $U_t$ and $V_t$ lie in $\mathbf Z_p$, so they have absolute value at
most $1$, and $\det U_t$ and $\det V_t$ are units of $\mathbf Z_p$, of absolute value $1$. Let $\beta_j$ be the number
of $a$ with $|d_a(t)|_p\leq p^{-E}$. As $J\leq p^h$ and $E\geq0$,
Lemma~\ref{pd-kappa-defects-segment-mean-theorem-a-plus-window-level-smith-point} over $\bar{\mathbf Q}_p$ gives
\[
 \beta_jE\leq-\log_p|\lambda|+\sum_{\zeta\in R}\nu_h(t_0+j-\zeta)+E\,\big|\{\zeta\in R:|t_0+j-\zeta|\leq p^{-h}\}\big|.
\]
Lemma~\ref{pd-kappa-defects-segment-mean-theorem-a-plus-window-level-smith-mean}, with $d=1$, $n_0=n_1$ and
$-\log_p|\lambda|$ in place of $\lambda$, applies, as $Y_0\geq1$, $Y_0\leq p^h$ and $E>0$. It gives
$\frac1{Y_0}\sum_j\beta_j\leq n_1\big(J/Y_0+(\Lambda_0+Jh/Y_0)/E\big)$. Multiplying by $Y_0$ gives the claim.
\end{proof}
% Lean: the second half of gramWindow2_proved: M(t) = binMatK at t = t0 + j (binMatK_eval_nat, hM); the passage to
%   K = PadicAlgCl p, where det M splits with as many roots as its degree (split_det, splits_form, binMatK_map,
%   with norm_algebraMap' for the norms), U', V', dg' the images of U, V, d (hsm: PadicInt.norm_le_one and
%   PadicInt.isUnit_iff); then Formalism.GM.Top.windowFullAP = windowFullAP_of valuation_window_lower_bound
%   windowMeanAP, applied with d = 1 (isCoprime_one_left), E = k + 1 and La0. windowFullAP_of gets |R| <= n J
%   from binMatK_det_natDegree_le (natDegree_C_mul and natDegree_multiset_prod_X_sub_C_eq_card) and applies
%   valuation_window_lower_bound (Lemma -smith-point) at each t0 + j inside windowMeanAP (Lemma -smith-mean).

\begin{lemma}\label{pd-kappa-defects-segment-mean-theorem-a-plus-window-level-gram}
Let $p$ be a prime, $n_1,J,h\geq0$, $t_0\geq0$, $D'\geq0$ and $W'\geq1$ integers with $W'\leq p^h$ and $J+1\leq p^h$, and
$\Lambda_0$ a real number. Let $V$ be a $\mathbf Q_p$-vector space, $e_1,\dots,e_{n_1}\in V$ linearly independent over
$\mathbf Z_p$ with $\mathbf Z_p$-span $\Lambda_e$, and $\Psi_0,\dots,\Psi_{W'-1}$ bilinear forms on $V$. Let $M$ be a
matrix of polynomials as in Lemma~\ref{pd-kappa-defects-segment-mean-theorem-a-plus-window-level-smith}, with
$c_{ab}(i)\in\mathbf Z_p$, $\det M\neq0$ and $v_p(\lambda)+\deg(\det M)/(p-1)\leq n_1\Lambda_0$. Suppose that for each
$0\leq i<W'$ there is a nonzero $s_i\in\mathbf Z_p$ with $s_i\Psi_i(e_a,e_b)=M_{ab}(t_0+i)$ for all $a,b$. Then, with
$\operatorname{bad}$ as in Definition~\ref{pd-kappa-defects-segment-mean-theorem-a-levels}(5),
\[
 \sum_{i=0}^{W'-1}\operatorname{bad}_{D'}(\Psi_i,\Lambda_e)\leq W'n_1\Big(\frac J{W'}+\frac{\Lambda_0+Jh/W'}{D'+1}\Big).
\]
\end{lemma}
\leanref{Formalism/SharpRestriction/GrowthBounds/FockGramWindow/GramWindow2.lean}{Formalism.MainNotions.G.GramWindowAlt.GramWindow2}
\leanref{Formalism/SharpRestriction/GrowthBounds/FockGramWindow/GramWindow2Proved.lean}{Formalism.MainNotions.G.GramWindowAlt.gramWindow2_proved}
\leanref{Formalism/SharpRestriction/GrowthBounds/FockGramWindow/GramWindow2Proved.lean}{Formalism.MainNotions.G.GramWindowAlt.bad_le_smith}
\leanref{Formalism/MahlerMeasure/FinalD.lean}{Formalism.GM.Top.smithZp}

\begin{proof}
Fix $i$ and put $t=t_0+i$. The entries of $M(t)$ lie in $\mathbf Z_p$, as the $\binom ti$ are integers. By the elementary
divisor theorem over $\mathbf Z_p$, $M(t)=U\operatorname{diag}(d_1,\dots,d_{n_1})V$ with $U,V\in GL_{n_1}(\mathbf Z_p)$ and
$d_a\in\mathbf Z_p$. We show $\operatorname{bad}_{D'}(\Psi_i,\Lambda_e)\leq|\{a:|d_a|_p\leq p^{-(D'+1)}\}|$; then
Lemma~\ref{pd-kappa-defects-segment-mean-theorem-a-plus-window-level-smith} with $Y_0=W'$ and $E=D'+1$ gives the
claim.

Let $P=U^{-1}$ and $Q=V^{-1}$, matrices over $\mathbf Z_p$, and put $x_a=\sum_bP_{ab}e_b$ and $y_a=\sum_bQ_{ba}e_b$, which
lie in $\Lambda_e$. By bilinearity, $s_i\Psi_i(x_a,y_{a'})=(PM(t)Q)_{aa'}=\delta_{aa'}d_a$. Let $S$ be the set of the $a$
with $|d_a|_p>p^{-(D'+1)}$, that is $d_a\neq0$ and $v_p(d_a)\leq D'$. For $a\in S$, $g_a=p^{D'}s_i/d_a$ lies in
$\mathbf Z_p$, as $v_p(s_i)\geq0$; and $\Psi_i(x_a,g_{a'}y_{a'})=p^{D'}\delta_{aa'}$ for $a,a'\in S$. So
$\operatorname{good}_{D'}(\Psi_i,\Lambda_e)\geq|S|$, and, as $\operatorname{rk}\Lambda_e=n_1$,
$\operatorname{bad}_{D'}(\Psi_i,\Lambda_e)\leq n_1-|S|$.
\end{proof}
% Lean: gramWindow2_proved: M j = the Z_p-matrix of binMatK at t0 + j (binMatK_eval_nat, hM), U V d from
%   smithZp = smith_form_p_adic_integers, hpt = bad_le_smith (the dual family x, cc y; hrank from
%   finrank_span_eq_card), and the count by Lemma -plus-window-level-smith. GramWindow2 takes c in Q_p with norm at
%   most 1 and s in Q_p nonzero with norm at most 1, which are the book's elements of Z_p, and linear maps in each
%   variable.

\begin{lemma}\label{pd-kappa-defects-segment-mean-theorem-a-plus-window-level}
Let an instance be given, $\alpha\in W^{S_j}$, $b$ a block and $\omega=(j,j_+)$ a level of $b$, with the degrees
$\delta$ and $\delta'_\omega$ formed with $\alpha$ (Definitions~\ref{pd-kappa-defects-segment-mean-theorem-a-levels}
and~\ref{pd-kappa-defects-segment-mean-theorem-a-plus-pairing}). Let $A\geq1$ be an integer with $\delta(\lambda)\leq2A$
for every letter $\lambda$, $b_+\leq2A$ and $\hat b_{u_l}\dot-o_{u_l}\leq2A$ for $1\leq l\leq t$. Let $D$ be an integer,
let $e_1,\dots,e_{n_1}\in\Lambda_\omega$ be linearly independent over $\mathbf Z_p$, with $\beta^-_\omega(e_a,e_b)=0$ for
$a\neq b$, and $\beta^-_\omega(e_a,e_a)\neq0$ and $v_p(\beta^-_\omega(e_a,e_a))\leq D$ for every $a$, and let $\Lambda_e$
be their $\mathbf Z_p$-span. For an integer $\zeta\geq0$ let $\beta_\omega[\zeta]$ be the level pairing $\beta_\omega$
with $\zeta$ in place of $\alpha_+\dot-o_+$. Put $J=\min(\delta'_\omega(x_+),b_+-j_+)$,
$\ell_A=\lfloor\log_p(8A+m)\rfloor+1$ and $x_A=\big((t+1)(m+1)+(2s+2)m+2(t+1)\big)\ell_A$. Let $z_0,h,D'\geq0$ and
$W'\geq1$ be integers with $W'\leq p^h$ and $J+1\leq p^h$. Then
\[
 \sum_{i=0}^{W'-1}\operatorname{bad}_{D'}\big(\beta_\omega[z_0+i],\Lambda_e\big)\leq
 W'n_1\Big(\frac J{W'}+\frac{x_A+D+\ell_A+Jh/W'}{D'+1}\Big).
\]
\end{lemma}
\leanref{Formalism/SharpRestriction/GrowthBounds/FockGramWindow/LvlDef.lean}{Formalism.MainNotions.G.GramWindowAlt.WinLvl}
\leanref{Formalism/SharpRestriction/GrowthBounds/FockGramWindow/WinLvl.lean}{Formalism.MainNotions.G.GramWindowAlt.winLvl}
\leanref{Formalism/SharpRestriction/GrowthBounds/FockGramWindow/LvlDef.lean}{Formalism.MainNotions.G.GramWindowAlt.JN}
\leanref{Formalism/SharpRestriction/GrowthBounds/FockGramWindow/LvlDef.lean}{Formalism.MainNotions.G.GramWindowAlt.lgN}
\leanref{Formalism/SharpRestriction/GrowthBounds/FockGramWindow/LvlDef.lean}{Formalism.MainNotions.G.GramWindowAlt.x0N}

\begin{proof}
Write $o=\delta'_\omega(x_+)$ and $n=b_+-j_+$, so $J=\min(o,n)$, and put $q=o+n+m-1$. Let $d_\omega$ be as in
Definition~\ref{pd-kappa-defects-segment-mean-theorem-a-block-bound-families-tie}(3), the denominator of
$\beta_\omega$ and of $\beta^-_\omega$. As $\Phi_\omega=\phi(x_+,\xi;o,n,\alpha_+\dot-o_+)\circ\Phi^Y_\omega$,
\[
 \beta_\omega[\zeta](u,v)=[\phi(x_+,\xi;o,n,\zeta)\Phi^Y_\omega u,v]/d_\omega,\qquad\beta^-_\omega(u,v)=[\Phi^Y_\omega u,v]/d_\omega .
\]
For $u\in\Lambda_\omega$ and $0\leq l\leq J$ let $\mu_l(u,v)$ be the number $\mu_l(u)$ of
Lemma~\ref{pd-kappa-defects-segment-mean-theorem-a-plus-window-level-pair} for the variables $x_{+,c},\xi_c$ and
the polynomial $v$, with $r=m$. Note $8A+m<p^{\ell_A}$.

\emph{Integrality.} We show $|p^{x_A}\mu_l(\Phi^Y_\omega u,v)/d_\omega|_p\leq1$ for $u,v\in\Lambda_\omega$ and $0\leq l\leq J$.
The elements of $\Lambda_\omega$ have coefficients in $\mathbf Z_p$. The operator $\Phi^Y_\omega$ is the composite of the
$t$ operators $\phi(x_{u_l},y_{u_l};\delta'_\omega(x_{u_l}),\delta'_\omega(y_{u_l}),\hat b_{u_l}\dot-o_{u_l})$, and
$\delta'_\omega(x_{u_l})+\delta'_\omega(y_{u_l})+m+(\hat b_{u_l}\dot-o_{u_l})\leq6A+m<p^{\ell_A}$; so by
Lemma~\ref{pd-kappa-defects-segment-mean-theorem-a-plus-window-level-bounds}(2), $\Phi^Y_\omega u$ is bounded by
$p^{t(m+1)\ell_A}$. As $J\leq b_+\leq2A<p^{\ell_A}$, part (1) of that lemma shows that
$U=\big((l!)^2\binom Jl\big)^{-1}E^lF^l\Phi^Y_\omega u$ is bounded by $p^{(t+1)(m+1)\ell_A}$, where $E$ and $F$ are formed
with the $x_{+,c}$ and the $\xi_c$. Next, $v$ is a $\mathbf Z_p$-combination of products $yg$ with $y\in Y^{\mathbf Z}_\omega$
and $g$ homogeneous of degree $b_+-j_+$ in the $\xi_c$. With the groups of
Lemma~\ref{pd-kappa-defects-segment-mean-theorem-a-block-bound-families-valuation-degrees}, every monomial $\nu$ of
$v$ has degree $\delta'_\omega(\lambda)$ in the group of each letter $\lambda$ and $b_+-j_+$ in $\Xi$, and
$d_\omega$ is the product of the factorials of these degrees, each at most $2A<p^{\ell_A}$. By
Lemma~\ref{pd-kappa-defects-segment-mean-theorem-a-block-bound-families-valuation-factorials}(3), applied to each
group with $L=\ell_A$, $v_p(d_\omega)\leq v_p(\nu!)+\big((2s+2)m+2(t+1)\big)\ell_A$. Hence each term
$U_\nu v_\nu\nu!/d_\omega$ of $[U,v]/d_\omega$ has absolute value at most
$p^{(t+1)(m+1)\ell_A+((2s+2)m+2(t+1))\ell_A}=p^{x_A}$, and so does their sum.

\emph{The matrix.} For $1\leq a,b\leq n_1$ and $0\leq l\leq J$ put $m_{ab}(l)=p^{x_A}\mu_l(\Phi^Y_\omega e_a,e_b)/d_\omega$, of
absolute value at most $1$, and let $c_{ab}(i)=\sum_{l=0}^{J-i}(-1)^l\binom{q-l}{J-l-i}m_{ab}(l)$ for $0\leq i\leq J$; these
lie in $\mathbf Z_p$. The polynomial $e_b$ has degree $o$ in the $x_{+,c}$ and $n$ in the $\xi_c$, as $y$ has degree
$\delta'_\omega(x_+)$ in $x_+$ and $g$ has degree $b_+-j_+$. So
Lemma~\ref{pd-kappa-defects-segment-mean-theorem-a-plus-window-level-pair}, multiplied by $p^{x_A}/d_\omega$, gives, for
every $\zeta\geq0$,
\[
 p^{x_A}\binom{q+\zeta}J\beta_\omega[\zeta](e_a,e_b)=\sum_{i=0}^Jc_{ab}(i)\binom\zeta i=M_{ab}(\zeta),
\]
with $M$ as in Lemma~\ref{pd-kappa-defects-segment-mean-theorem-a-plus-window-level-lead}, and
$c_{ab}(J)=p^{x_A}\beta^-_\omega(e_a,e_b)=\delta_{ab}d_a$, where $d_a=p^{x_A}\beta^-_\omega(e_a,e_a)\neq0$ has
$v_p(d_a)\leq x_A+D$. If $J\geq1$, $\lfloor\log_pJ\rfloor+1\leq\ell_A$, and if $J=0$, $J/(p-1)-v_p(J!)=0$; so
Lemma~\ref{pd-kappa-defects-segment-mean-theorem-a-plus-window-level-lead}, with $R=x_A+D$ and $\Lambda=\ell_A$, gives
$\det M\neq0$ and $v_p(\lambda)+\deg(\det M)/(p-1)\leq n_1(x_A+D+\ell_A)$ for its leading coefficient $\lambda$.

\emph{The window.} Apply Lemma~\ref{pd-kappa-defects-segment-mean-theorem-a-plus-window-level-gram} to the
polynomial ring, the family $e_1,\dots,e_{n_1}$, $t_0=z_0$, $\Psi_i=\beta_\omega[z_0+i]$, which is bilinear, the nonzero
$s_i=p^{x_A}\binom{q+z_0+i}J\in\mathbf Z_p$ (nonzero as $q\geq J$), and $\Lambda_0=x_A+D+\ell_A$. It gives the claim.
\end{proof}
% Lean: winLvl = win_abs for the tie of x_+ and xi (w = the x_+ variables, e = xiVar, o = LL.deg uPos - jp,
%   n = B - jp, T = phisB, dF = degFact) with gramWindow2_proved (Lemma -plus-window-level-gram), Ljj_biHom
%   (biHom_plus, lettersAt_uPos), and hV = legendre' with log_le_lgN. Integrality is int_lvl: Ljj_int, bnd_phisB
%   (bnd_foldr over the t ties with bnd_phi), bnd_Xl, and norm_fock_le with Ljj_val (Ljj_hom, prod_dJ, dJ_le and
%   val_letters_le); the exponent is x0N. The matrix is win_abs: c = cW (o + n + r - 1) J of p^x0N mX / dF, with
%   cW_norm, cW_top (with horth and mX_zero for the diagonal top coefficients), lead_bound with neglog_pow_mul, and
%   phi_pair with cW_sum and binMatK_eval_nat for the identity. WinLvl holds for every block datum G with
%   BlockOK; the book applies it to G = gB aTop 0 b, where setZ G z is gB aTop aPlus b for z = aPlus - oPlus.

\begin{lemma}\label{pd-kappa-defects-segment-mean-theorem-a-plus-window-numbers}
Let $C_H\geq2$, $C_{E1},C_{E2}\geq0$ and $C_M\geq1$, and let an instance satisfy its hypotheses. Let
$\alpha\in W^{S_j}$, let $b$ be a block kept for $\alpha$ and $J$ an integer with $0\leq J\leq b_+$. Put
$D_1=D_1(C_M)$ and $A=(3s+2)a$, and let $\ell_A$ and $x_A$ be as in
Lemma~\ref{pd-kappa-defects-segment-mean-theorem-a-plus-window-level}. Then $o_+\leq a$, $YD_1\geq0$, and there is an
integer $h\geq0$ with $|W|\leq p^h$, $J+1\leq p^h$ and
\[
 \frac J{|W|}+\frac{x_A+\lfloor D_1\rfloor+\ell_A+Jh/|W|}{\lfloor YD_1\rfloor+1}\leq\frac5Y.
\]
\end{lemma}
\leanref{Formalism/SharpRestriction/GrowthBounds/FockGramWindow/RestWinOfLvl.lean}{Formalism.MainNotions.G.GramWindowAlt.RestWinNum}
\leanref{Formalism/SharpRestriction/GrowthBounds/FockGramWindow/InitialPointNumericBound.lean}{Formalism.MainNotions.G.GramWindowAlt.restWinNum_proved}

\begin{proof}
Write $s=S_j$, $\ell=\log_p(8a+m)+1$ and $c_S=3(s+2)m^2\ell$
(Definitions~\ref{pd-kappa-defects-numbers} and~\ref{pd-kappa-defects-segment-mean-theorem-a-levels}).

\emph{Facts.} As in the proof of Lemma~\ref{pd-kappa-defects-segment-mean-theorem-a-top-sizes}, the hypotheses give
$m\geq4$, $2s+2\leq m$, $(m+1)^2\leq p$, so $p\geq25$, $a\geq1$, $b'>0$, $pb'<a$, $0\leq\Lambda$, $4\Lambda<m+1$,
$E_0\geq0$, $b_v\leq\frac52b'$ for every $v$, $b_+\leq\frac52b'$, $(m+1)b'\leq a/4$ and $pP_v\leq P_{v'}$ for $v'<v$.
By hypothesis (1) of Theorem~\ref{pd-gain-theorem-one}, $Y\geq1$ and $k^2Y^2\leq p$, where $k=m+1\geq5$. Also
$x_a\geq0$, and $W$ is the set of the integers from $a$ to $u_W=\lfloor a+a/(kY)+x_a\rfloor\leq2a$; so $a\in W$,
$|W|=u_W-a+1\leq a+1$, and $|W|\geq a/(kY)+x_a\geq a/(kY)$.

\emph{The degree $o_+$.} As $b$ is kept, $o_+\leq O_+$. First, $\operatorname{cap}(s)\leq\hat a+s\cdot\frac52b'$, as each
$\hat b_v-o^\flat_v$ is at most $b_v$. As $g_j\geq1$ and $E_0\geq0$, $\hat a\leq a/p\leq a/25$; and as $s\leq(m+1)/2$,
$sb'\leq a/8$. Next, $\Lambda/(m+1)\leq1/4$ and $4/p^2\leq4/625$. The numbers $f_v=1/P_v$ lie in $(0,1]$ and satisfy
$pf_{v'}\leq f_v$ for $v'<v$, so $\sum_vf_v\leq f_s(1+p^{-1}+p^{-2}+\cdots)\leq p/(p-1)\leq25/24$, and
$\sum_v2a/P_v\leq2a\cdot\frac{25}{24}$. Finally $\Lambda b_+/(m+1)\leq b_+/4\leq\frac58b'$, and $b'\leq a/20$. So
\[
 o_+\leq\frac a{25}+\frac{5a}{16}+\Big(\frac14+\frac4{625}\Big)\frac{25a}{12}+\frac a{32}\leq a .
\]

\emph{The depth.} As $p>1$ and $8a+m\geq1$, $\ell\geq1$, so $c_S\geq\ell\geq1$. As $t\geq0$, $Y\geq1$ and
$C_M\geq1$, $D_1=t(Yc_S+C_Mc_S)+C_Mc_S\geq c_S>0$; so $YD_1\geq0$.

\emph{The exponent $h$.} Let $h\geq0$ be the least integer with $a+1\leq p^h$. Then $|W|\leq a+1\leq p^h$; and
$J\leq b_+\leq\frac52b'\leq a$, so $J+1\leq p^h$. As $a+1\geq2$, $p^{h-1}<a+1$, so
$h\leq\log_p(a+1)+1\leq\ell\leq c_S$.

\emph{The quotient $q=J/|W|$.} As $k^2Y^2b'\leq pb'<a$, $kYb'<a/(kY)\leq|W|$. As $J\leq\frac52b'$ and $k\geq5$,
$2YJ\leq5Yb'\leq kYb'<|W|$; so $0\leq q\leq1/(2Y)\leq1$.

\emph{The term $x_A+\ell_A$.} As $3s+2\leq p$, $\log_p(8A+m)\leq\log_p\big((3s+2)(8a+m)\big)\leq1+\log_p(8a+m)$, so
$\ell_A\leq\ell+1\leq2\ell$. With $t\leq s$ and $m\geq4$, the factor $C=(t+1)(m+1)+(2s+2)m+2(t+1)+1$ satisfies
$C\leq3(s+1)(m+1)+1$ and $2C\leq12(s+2)m\leq3(s+2)m^2$. So $x_A+\ell_A=C\ell_A\leq2C\ell\leq c_S$.

\emph{The bound.} Put $E=\lfloor YD_1\rfloor+1\geq YD_1>0$. By the last three paragraphs, $qh\leq h\leq c_S\leq D_1$ and
$x_A+\ell_A\leq c_S\leq D_1$, and $\lfloor D_1\rfloor\leq D_1$; so the numerator of the second term is at most $3D_1$,
and the second term is at most $3D_1/(YD_1)=3/Y$. With $q\leq1/(2Y)$, the left side is at most $4/Y\leq5/Y$.
\end{proof}
% Lean: restWinNum_proved hS hpk CM hCM with hS = segFacts. The degree o_+ is oplus_le_a (cap_le' with
%   B = 5b'/2, first_coefficient_le_div for A0 <= a/25, geo_lemma with P_step and hgeo for the geometric
%   sum); the depth is kap_facts (lgA <= kapS, 0 < kapS <= D1r); h is Nat.clog p (a + 1), with win_facts for
%   the size of winA and clog_le_lg with hlk for h <= lgA <= kapS; q is q_num; x_A + l_A is hx_inst
%   (initial_point_nat_cast, lgN_le, initial_point_numeric_bound); the bound is eps_num.

\begin{lemma}\label{pd-kappa-defects-segment-mean-theorem-a-plus-window}
Let $C_0\geq1$, $C_{E1},C_{E2}\geq1$, $C_H\geq2$ and $C_M\geq1$, and let an instance satisfy its hypotheses.
Let $\delta$ be a digit class, $\alpha\in A(\delta)$ and $(b,\omega)\in I(\delta,\alpha)$, and put $D_1=D_1(C_M)$
(Definitions~\ref{pd-kappa-defects-segment-mean-theorem-a-levels}
and~\ref{pd-kappa-defects-segment-mean-theorem-a-plus-pairing}). Let $e_1,\dots,e_{n_1}\in\Lambda_\omega$ be linearly
independent over $\mathbf Z_p$, with $\beta^-_\omega(e_a,e_b)=0$ for $a\neq b$, and $\beta^-_\omega(e_a,e_a)\neq0$ and
$v_p(\beta^-_\omega(e_a,e_a))\leq\lfloor D_1\rfloor$ for every $a$. Let $\Lambda_e$ be the $\mathbf Z_p$-span of the $e_a$, and for $\alpha_+\in W$ let $\beta_\omega(\alpha_+)$ be the level
pairing $\beta_\omega$ formed with $\alpha_+$ (Definition~\ref{pd-kappa-defects-segment-mean-theorem-a-levels}). Then
\[
 \sum_{\alpha_+\in W}\operatorname{bad}_{\lfloor YD_1\rfloor}(\beta_\omega(\alpha_+),\Lambda_e)\leq|W|\cdot\frac{5n_1}Y.
\]
\end{lemma}
\leanref{Formalism/SharpRestriction/GrowthBounds/MonomialFactorials/BejjSmulLeft.lean}{Formalism.MainNotions.G.SymmetricFormBlocks.RestWin}
\leanref{Formalism/SharpRestriction/GrowthBounds/FockGramWindow.lean}{Formalism.MainNotions.G.GramWindowAlt.restWinE_proved}

\begin{proof}
Write $s=S_j$, $D=\lfloor D_1\rfloor$, $D'=\lfloor YD_1\rfloor$ and $A=(3s+2)a$, and let $\ell_A$ and $x_A$ be as in
Lemma~\ref{pd-kappa-defects-segment-mean-theorem-a-plus-window-level}. As $(b,\omega)\in I(\delta,\alpha)$, $b$ is kept for
$\alpha\in W^s$ and $\omega$ is a level of $b$. Put $J=\min(\delta'_\omega(x_+),b_+-j_+)$, so $0\leq J\leq b_+$.

\emph{The numbers.} By Lemma~\ref{pd-kappa-defects-segment-mean-theorem-a-plus-window-numbers}, $o_+\leq a$,
$YD_1\geq0$, so $D'\geq0$, and there is $h\geq0$ with $|W|\leq p^h$, $J+1\leq p^h$ and
\[
 \frac J{|W|}+\frac{x_A+D+\ell_A+Jh/|W|}{D'+1}\leq\frac5Y .
\]

\emph{The degrees.} As in the proof of Lemma~\ref{pd-kappa-defects-segment-mean-letters-degrees}, as $b$ is kept for
$\alpha$, the degrees $\delta$ formed with $\alpha$ satisfy $\delta(x_v)\leq a/10$, $\delta(y_v)\leq2a$ and
$\delta(x_+)\leq a+sa/10+2as+a/10$, and $b_+\leq a/10$; also $\hat b_{u_l}\dot-o_{u_l}\leq\hat b_{u_l}\leq a/10$. All of
these are at most $2A=(6s+4)a$, and $A\geq1$ as $a\geq1$.

\emph{The window.} As $a/(kY)\geq0$ and $x_a\geq0$, $W$ is the set of the integers $a+i$ with $0\leq i<|W|$, and
$|W|\geq1$. The level pairing $\beta_\omega(\alpha_+)$ depends on $\alpha_+$ only through the exponent $\alpha_+\dot-o_+$
of the factor $\phi(x_+,\xi;\cdot)$ of $\Phi_\omega$, and $\Lambda_\omega$ and $\beta^-_\omega$ do not depend on $\alpha_+$. As
$o_+\leq a$, $(a+i)\dot-o_+=z_0+i$ with $z_0=a-o_+\geq0$; so $\beta_\omega(a+i)=\beta_\omega[z_0+i]$ in the notation of
Lemma~\ref{pd-kappa-defects-segment-mean-theorem-a-plus-window-level}.

\emph{The bound.} By the last two paragraphs, Lemma~\ref{pd-kappa-defects-segment-mean-theorem-a-plus-window-level}
applies with $A$, $D$, the family $e_1,\dots,e_{n_1}$, $z_0$, $W'=|W|$, $h$ and $D'$. With the first paragraph,
\[
 \sum_{\alpha_+\in W}\operatorname{bad}_{D'}(\beta_\omega(\alpha_+),\Lambda_e)=\sum_{i=0}^{|W|-1}
 \operatorname{bad}_{D'}(\beta_\omega[z_0+i],\Lambda_e)\leq|W|n_1\Big(\frac J{|W|}+\frac{x_A+D+\ell_A+Jh/|W|}{D'+1}\Big)
 \leq|W|\cdot\frac{5n_1}Y .
\]
\end{proof}
% Lean: restWinE_proved = restWin_of_lvl (thresholds 1, 1 and 2) with winLvl (Lemma
%   pd-kappa-defects-segment-mean-theorem-a-plus-window-level, open), restWinNum_proved hS hpk CM hCM (Lemma
%   pd-kappa-defects-segment-mean-theorem-a-plus-window-numbers) and a_mem_win hS. restWin_of_lvl: haTop', hK
%   and hlv from cbox and idx, hJ (JN <= bPlus), gB_deg_le at aPlus = a for the degree bounds with A = aBig,
%   hmain = winLvl at G = gB aTop 0 b with z0 = a - oPlus, W = card winA and k = floor(Y D1r), hbe (bejj of
%   gB aTop aPlus b is that of setZ G (aPlus - oPlus)), winA_sum, and the numeric bound. The Lean states
%   RestWin for every family with the properties of Lemma pd-kappa-defects-segment-mean-theorem-a-plus-orth
%   at floor(D1r CM), and ignores the conditions AdmNoK. gB_deg_le proves the degree bounds through
%   segChain_deg_le, as in Lemma pd-kappa-defects-segment-mean-letters-degrees.

\begin{lemma}\label{pd-kappa-defects-segment-mean-theorem-a-plus}
Let $C_0\geq1$, $C_{E1},C_{E2}\geq1$, $C_H\geq2$ and $C_M\geq1$, and let an instance satisfy its hypotheses.
Let $\delta$ be a digit class, $\alpha\in A(\delta)$ and $(b,\omega)\in I(\delta,\alpha)$, with $\omega=(j,j_+)$,
and put $D_1=D_1(C_M)$ (Definition~\ref{pd-kappa-defects-segment-mean-theorem-a-levels}). Then
\[
 \sum_{\alpha_+\in W}\operatorname{bad}_\omega(\alpha,\alpha_+,b;\lfloor YD_1\rfloor)\leq
 |W|\Big(\operatorname{bad}^Y_\omega(\alpha,b;\lfloor D_1\rfloor)+\frac5Y\dim Y_\omega\Big)N_\xi(b_+-j_+).
\]
\end{lemma}
\leanref{Formalism/SharpRestriction/GrowthBounds/MonomialFactorials/BlockBoundUniformiser.lean}{Formalism.MainNotions.G.SymmetricFormBlocks.RestPlus}
\leanref{Formalism/SharpRestriction/GrowthBounds/MonomialFactorials/RestPlusOf.lean}{Formalism.MainNotions.G.SymmetricFormBlocks.restPlus_of}

\begin{proof}
Write $n=b_+-j_+$, $D=\lfloor D_1\rfloor$ and $D'=\lfloor YD_1\rfloor$. By
Lemma~\ref{pd-kappa-defects-segment-mean-theorem-a-plus-orth} there are $n_1\geq0$ and $e_1,\dots,e_{n_1}\in\Lambda_\omega$
as there; let $\Lambda_e\subset\Lambda_\omega$ be their $\mathbf Z_p$-span, of rank $n_1$. Fix $\alpha_+\in W$. A dual family
in $\Lambda_e$ for $\beta_\omega(\alpha_+)$ with pairings $p^{D'}$ is one in $\Lambda_\omega$, so
$\operatorname{good}_{D'}(\beta_\omega(\alpha_+),\Lambda_e)\leq\operatorname{good}_{D'}(\beta_\omega(\alpha_+),\Lambda_\omega)$,
and the former is at most $n_1$. Hence
\[
 \operatorname{bad}_\omega(\alpha,\alpha_+,b;D')\leq(\operatorname{rk}\Lambda_\omega-n_1)
 +\operatorname{bad}_{D'}(\beta_\omega(\alpha_+),\Lambda_e).
\]
Summing over $\alpha_+\in W$ and using Lemma~\ref{pd-kappa-defects-segment-mean-theorem-a-plus-window},
\[
 \sum_{\alpha_+\in W}\operatorname{bad}_\omega(\alpha,\alpha_+,b;D')\leq|W|\Big((\operatorname{rk}\Lambda_\omega-n_1)
 +\frac{5n_1}Y\Big).
\]
By Lemmas~\ref{pd-kappa-defects-segment-mean-theorem-a-plus-orth}
and~\ref{pd-kappa-defects-segment-mean-theorem-a-plus-phi},
$\operatorname{rk}\Lambda_\omega-n_1\leq\operatorname{bad}^-_\omega(\alpha,b;D)\leq\operatorname{bad}^Y_\omega(\alpha,b;D)N_\xi(n)$;
and by Lemmas~\ref{pd-kappa-defects-segment-mean-theorem-a-plus-orth}
and~\ref{pd-kappa-defects-segment-mean-theorem-a-plus-rank}, $n_1\leq\operatorname{rk}\Lambda_\omega\leq\dim
Y_\omega\cdot N_\xi(n)$. As $Y>0$, this gives the claim.
\end{proof}
% Lean: openE_of_open4 uses restPlus_of CM hY (plusACondition CM) restRankL_proved b, where b : RestPlusB is
%   restPlusB_of CM hY RestOrth RestWin from restricted_bounds_weaken_from_seven_to_six; restPlusB_of uses
%   bad_le_sub with bejj_bdd (the good counts are bounded) and finrank_span_eq_card for the rank of the span.
%   RestOrth is restOrth_of (Lemma pd-kappa-defects-segment-mean-theorem-a-plus-orth) and RestWin is
%   restWinE_proved (Lemma pd-kappa-defects-segment-mean-theorem-a-plus-window, open).

\begin{definition}\label{pd-kappa-defects-segment-mean-theorem-a-top-model}
Let $p$ be a prime and $s\geq0$, $r\geq1$ integers, and put $N=2s+1$. A \emph{model instance} consists of letters
$\lambda_0,\dots,\lambda_{N-1}$, each of vector or covector type, with variables $\lambda_{g,c}$ ($1\leq c\leq r$) and
degrees $\deg\lambda_g\geq0$; integers $t\geq0$, $\pi_1,\dots,\pi_t\geq0$ and $z_1,\dots,z_t\geq0$ (the \emph{ties});
finitely many pairs $(J_k,M_k)$ of an integer $0\leq J_k\leq N$ and $M_k\in\mathbf Z\cup\{+\infty\}$ (the
\emph{conditions}); real numbers $a$ and $E$; integers $n^0_i,h_i,R_i\geq0$ for $1\leq i\leq t$; and an
\emph{index} $\varrho$ with $0\leq\varrho\leq r-1$.
\begin{enumerate}
\item For $n\in\mathbf Z_{\geq0}^t$, the lettering $\lambda(n)$ has the same letters, with the degree of
  $\lambda_{\pi_i+1}$ replaced by $n_i$ for every $i$. For a lettering $\lambda'$ with degrees $\deg'$, let
  $Y(\lambda')$ be the $\mathbf Q_p$-vector space of the polynomials over $\mathbf Q_p$ in the $\lambda_{g,c}$, of total
  degree $\deg'\lambda_g$ in the variables of each letter, that lie in $Y(\mathcal L_k,M_k)$ for every $k$, where
  $\mathcal L_k$ consists of the first $J_k$ letters (Definition~\ref{pd-kappa-defects-segment-lattice}(4) over
  $\mathbf Q_p$, with no bound on the weight if $M_k=+\infty$); let $Y^{\mathbf Z}(\lambda')$ be the set of its elements
  with coefficients in $\mathbf Z_p$; and put
  \[
   \beta(\lambda')(F,G)=\frac{[\Phi(\lambda')F,G]}{\prod_g\deg'\lambda_g!},\qquad
   \Phi(\lambda')=\phi_1\circ\cdots\circ\phi_t,\quad
   \phi_i=\phi\big(\lambda_{\pi_i},\lambda_{\pi_i+1};\deg'\lambda_{\pi_i},\deg'\lambda_{\pi_i+1},z_i\big),
  \]
  with $\phi$ as in Definition~\ref{pd-kappa-defects-segment-mean-theorem-a-levels}(3) and the Fock pairing
  $[\ ,\ ]$ over $\mathbf Q_p$.
\item Put $\ell=\log_p(8a+r)+1$, $\kappa=3(s+2)r^2\ell$ and $D_1=t(E+\kappa)+\kappa$. The \emph{progression} $P$ is the
  set of the $n\in\mathbf Z^t$ with $n_i\in\{n^0_i+h_ij:0\leq j<R_i\}$ for every $i$; put $n^1_i=n^0_i+h_i(R_i-1)$, and
  for $n\in\mathbf Z^t$ let $n[i]$ be $n$ with its $i$-th entry replaced by $n^1_i$. Let $S_i$ be the sum of the
  degrees in $\lambda(n^1)$ of the letters $\lambda_g$ with $g\leq\pi_i$, and
  \[
   \varepsilon_i=\frac{2S_i}{R_i}+\frac1E\Big(2\kappa+\frac{2S_i\ell}{R_i}\Big).
  \]
\item The model instance \emph{satisfies the model hypotheses} if: $p\neq2$, $r\geq2$ and $\varrho=r-1$; the $\pi_i$ are
  strictly increasing, and for every $i$, $\pi_i+1<N$, $\lambda_{\pi_i}$ is of vector type, $\lambda_{\pi_i+1}$ is of
  covector type and some $J_k$ equals $\pi_i+1$; $\lambda_g$ is of vector type exactly when $g$ is even; every $J_k$
  is $N$ or one of the $\pi_i+1$; $a\geq1$ and $E\geq1$; for every $i$, $R_i\geq1$, $h_i>0$, $p\nmid h_i$ and
  $S_i\leq n^0_i$; and, with $\deg^1$ the degrees of $\lambda(n^1)$, $\sum_g\deg^1\lambda_g\leq8a+r$ and, for every
  $i$, each of $\deg^1\lambda_{\pi_i}+\deg^1\lambda_{\pi_i+1}+r+z_i+1$, $R_i$ and $2S_i+1$ is at most $8a+r$.
\end{enumerate}
\end{definition}
\leanref{Formalism/SharpRestriction/GrowthBounds/FlipOneThree.lean}{Formalism.MainNotions.flipOneThree.InstM}
\leanref{Formalism/SharpRestriction/GrowthBounds/FlipOneThree.lean}{Formalism.MainNotions.flipOneThree.InstM.prog}
\leanref{Formalism/SharpRestriction/GrowthBounds/FlipOneThree.lean}{Formalism.MainNotions.flipOneThree.InstM.top}
\leanref{Formalism/SharpRestriction/GrowthBounds/FlipOneThree.lean}{Formalism.MainNotions.flipOneThree.InstM.eps}
\leanref{Formalism/SharpRestriction/GrowthBounds/FlipOneThree.lean}{Formalism.MainNotions.flipOneThree.InstM.D1}
\leanref{Formalism/SharpRestriction/GrowthBounds/FlipOneThree.lean}{Formalism.MainNotions.flipOneThree.InstM.Hyp}
\leanref{Formalism/MahlerMeasure/StatementTM.lean}{Formalism.GM.M.Letters.setTies}
\leanref{Formalism/MahlerMeasure/StatementTM.lean}{Formalism.GM.M.TieData.phis}
\leanref{Formalism/MahlerMeasure/GoodCount.lean}{Formalism.GM.M.TieData}
\leanref{Formalism/MahlerMeasure/GoodCount.lean}{Formalism.GM.M.TieData.S}

\begin{lemma}\label{pd-kappa-defects-segment-mean-theorem-a-top-mean}
Let a model instance satisfy the model hypotheses
(Definition~\ref{pd-kappa-defects-segment-mean-theorem-a-top-model}). Then
\[
 \frac1{|P|}\sum_{n\in P}\operatorname{bad}_{\lfloor D_1\rfloor}\big(\beta(\lambda(n)),Y^{\mathbf Z}(\lambda(n))\big)\leq
 \sum_{i=1}^t\varepsilon_i\cdot\frac1{|P|}\sum_{n\in P}\dim Y\big(\lambda(n[i])\big),
\]
with $\operatorname{bad}$ as in Definition~\ref{pd-kappa-defects-segment-mean-theorem-a-levels}(5).
\end{lemma}
\leanref{Formalism/SharpRestriction/GrowthBounds/FlipOneThree.lean}{Formalism.MainNotions.flipOneThree.ModelLatticeRank}
\leanref{Formalism/SharpRestriction/GrowthBounds/CountingStatementReduction.lean}{Formalism.MainNotions.G.CountingStatementReduction.alternative_mean_value_theorem}
\leanref{Formalism/ThetaLogCoordinates/PartnerMeanCoordinates/LatticeTheoremEps28Valid.lean}{Formalism.InvariantFactorPaths.MC.Top.theoremMV7}

\begin{proof}
\openstep{The mean value estimate over the progression for the badness of the polynomial parts of a model
instance is not yet written out.}
% Size: most of the about 39,400 lines of the old open step of Lemma pd-kappa-defects-segment-mean-theorem-a-top;
%   about 27,300 lines of its closure lie under Formalism (MahlerMeasureProofs about 9,700,
%   ThetaLogCoordinates/FockResolutionOnto about 7,700, PolynomialDimension/MultinomialDepth about 5,200,
%   ThetaLogCoordinates/TraceIsometryWeights about 2,300).
% Lean: alternative_mean_value_theorem D X = MC.Top.theoremMV7 = MC.theoremMV7 (TheoremMV7) applied to the
%   eleven parts of X.Hyp.
\end{proof}

\begin{lemma}\label{pd-kappa-defects-segment-mean-theorem-a-top-box}
Let $C_H\geq2$ and $C_{E1},C_{E2}\geq0$, and let an instance satisfy its hypotheses; use the notation of
Definitions~\ref{pd-kappa-defects-segment-mean-theorem-a-classes}
and~\ref{pd-kappa-defects-segment-mean-theorem-a-levels}. Let $\delta=(\rho,\rho')$ be a digit class, and for
$1\leq k\leq s$ let $W_k(\delta)$ be the set of the $y\in W$ with $|\rho'_k|\leq y$ and $P_k\mid y-|\rho'_k|$.
\begin{enumerate}
\item If $A(\delta)$ is nonempty, then $A(\delta)=W_1(\delta)\times\cdots\times W_s(\delta)$.
\item Let $\alpha_0\in A(\delta)$ and $(b,\omega)\in I(\delta,\alpha_0)$, with $\omega=(j,j_+)$. Then
  $j_i\leq(y-|\rho'_{u_i}|)/P_{u_i}$ for every $1\leq i\leq t$ and every $y\in W_{u_i}(\delta)$.
\end{enumerate}
\end{lemma}
\leanref{Formalism/SharpRestriction/GrowthBounds/MonomialFactorials/DimYOvr.lean}{Formalism.MainNotions.G.SymmetricFormBlocks.cbox_eq}
\leanref{Formalism/SharpRestriction/GrowthBounds/MonomialFactorials/ExponentMatrixFirstDepthEq.lean}{Formalism.MainNotions.G.SymmetricFormBlocks.lvlRoom_of}

\begin{proof}
(1) Compatibility of $\delta$ with $\alpha\in W^s$ consists of conditions on $\rho$ alone, which hold as $A(\delta)$
is nonempty, and of the conditions $\alpha_k\in W_k(\delta)$.

(2) Let $y\in W_{u_i}(\delta)$, and let $\alpha$ be $\alpha_0$ with its $u_i$-th entry replaced by $y$. By (1),
$\alpha\in A(\delta)$, and by Lemma~\ref{pd-kappa-defects-segment-mean-theorem-a-free}, $(b,\omega)\in I(\delta,\alpha)$.
So $\omega$ is a level of $b$ for $\alpha$, and as $u_i>\eta_j$, $j_i\leq\delta(y_{u_i})=\hat\alpha_{u_i}=
(y-|\rho'_{u_i}|)/P_{u_i}$ (Definition~\ref{pd-kappa-defects-segment-lattice}(3) and~(5)).
\end{proof}
% Lean: cbox_eq (mem_cbox_iff_mem_piFinset_and_compat, mem_Wl_iff_mem_winA_and_dvd_sub; W_k(delta) is Wl) and
%   lvlRoom_of (LvlRoom, through hF = restFree_proved, mem_levels, chain_eta and coarseC).

\begin{definition}\label{pd-kappa-defects-segment-mean-theorem-a-top-attached}
Let an instance be given, with the notation of
Definitions~\ref{pd-kappa-defects-segment-mean-theorem-a-classes}
and~\ref{pd-kappa-defects-segment-mean-theorem-a-levels}, let $\delta=(\rho,\rho')$ be a digit class and $W_k(\delta)$
as in Lemma~\ref{pd-kappa-defects-segment-mean-theorem-a-top-box}. Let $\alpha\in W^s$, $b$ a block and $\omega=(j,j_+)$
a level of $b$ for $\alpha$, and suppose that every $W_{u_i}(\delta)$ is nonempty and $j_i\leq(y-|\rho'_{u_i}|)/P_{u_i}$
for every $i$ and $y\in W_{u_i}(\delta)$. Let $N_i(\delta,\omega)$ be the set of the numbers
$(y-|\rho'_{u_i}|)/P_{u_i}-j_i$ with $y\in W_{u_i}(\delta)$. The \emph{model instance attached to}
$(\delta,\alpha,b,\omega)$ (Definition~\ref{pd-kappa-defects-segment-mean-theorem-a-top-model}) has the prime $p$, the
integers $s$ and $r=m$; the letters $x_1,y_1,\dots,x_s,y_s,x_+$, in this order, as $\lambda_0,\dots,\lambda_{2s}$, with
their types and the degrees $\delta'_\omega$ formed with $\alpha$; the ties $\pi_i=2u_i-2$, the position of $x_{u_i}$,
and $z_i=\hat b_{u_i}\dot-o_{u_i}$; the conditions $(2u_l-1,M_l)$ for $1\leq l\leq t$ and $(2s+1,M_+)$ of
Definition~\ref{pd-kappa-defects-segment-lattice}(5); the numbers $a$ and $E=Yc_S$; $n^0_i=\min N_i(\delta,\omega)$,
$h_i=1$ and $R_i=\max N_i(\delta,\omega)-\min N_i(\delta,\omega)+1$; and the index $\varrho=m-1$.
\end{definition}
\leanref{Formalism/SharpRestriction/GrowthBounds/MonomialFactorials/DimYOvr.lean}{Formalism.MainNotions.G.SymmetricFormBlocks.XMof}
\leanref{Formalism/SharpRestriction/GrowthBounds/MonomialFactorials/DimYOvr.lean}{Formalism.MainNotions.G.SymmetricFormBlocks.Wl}
\leanref{Formalism/SharpRestriction/GrowthBounds/MonomialFactorials/DimYOvr.lean}{Formalism.MainNotions.G.SymmetricFormBlocks.Nl}
\leanref{Formalism/SharpRestriction/GrowthBounds/MonomialFactorials/DimYOvr.lean}{Formalism.MainNotions.G.SymmetricFormBlocks.cm}
\leanref{Formalism/SharpRestriction/GrowthBounds/ObjSeg.lean}{Formalism.MainNotions.G.segTies}

\begin{lemma}\label{pd-kappa-defects-segment-mean-theorem-a-top-sizes}
Let $C_H\geq2$ and $C_{E1},C_{E2}\geq0$, and let an instance satisfy its hypotheses. Let $\delta$ be a digit class,
$\alpha\in A(\delta)$ and $(b,\omega)\in I(\delta,\alpha)$, so that the model instance attached to
$(\delta,\alpha,b,\omega)$ is defined (Lemma~\ref{pd-kappa-defects-segment-mean-theorem-a-top-box}). Then, in the
notation of Definition~\ref{pd-kappa-defects-segment-mean-theorem-a-top-model} for this model instance,
$S_i\leq n^0_i$ and $z_i\leq2a$ for every $i$, $\deg^1\lambda_g\leq2a$ for every $g$, $\sum_g\deg^1\lambda_g\leq8a$, and
$a\geq1$.
\end{lemma}
\leanref{Formalism/SharpRestriction/GrowthBounds/MonomialFactorials/RestTMHypOf.lean}{Formalism.MainNotions.G.SymmetricFormBlocks.RestTMSize}
\leanref{Formalism/SharpRestriction/GrowthBounds/MonomialFactorials/LowerBoundInitialMultiplicity.lean}{Formalism.MainNotions.G.SymmetricFormBlocks.restTMSize_proved}

\begin{proof}
Write $\delta=(\rho,\rho')$ and $s=S_j$, and let $\mathcal M$ be the attached model instance.

\emph{Facts.} As in the proofs of Lemmas~\ref{pd-kappa-defects-segment-mean-letters-degrees}
and~\ref{pd-kappa-defects-segment-mean-theorem-a-free}, the hypotheses give $m\geq4$, $2s+2\leq m$, $(m+1)^2\leq p$,
$a\geq1$, $pb'<a$, $0\leq\Lambda$, $4\Lambda<m+1$ and $E_0\geq0$, so $\hat a\leq a$; $b_v\leq5b'/2$ and
$(m+1)P_v\leq b'$ for $1\leq v\leq s$, and $b_+\leq5b'/2$; and every element of $W$ is at most
$a+a/(kY)+x_a\leq2a$. As $4(m+1)\leq(m+1)^2\leq p$, $4(m+1)b'\leq pb'<a$, so $(m+1)b'\leq a/4$, and as $m+1\geq5$,
$\frac52b'\leq a/8$. For $1\leq v\leq s$, $P_v=p^{g_j-d_{j,v}}\geq p$, as $d_{j,v}<g_j$; and $pP_v\leq P_{v'}$ for
$v'<v$, as $d_{j,v'}<d_{j,v}$.

\emph{The top lettering.} Let $\alpha^1$ be $\alpha$ with each entry $\alpha_{u_l}$ replaced by $\tau(\delta,l)$, the
largest element of $W_{u_l}(\delta)$. By Lemma~\ref{pd-kappa-defects-segment-mean-theorem-a-top-box}(1),
$\alpha^1\in A(\delta)$, and by Lemma~\ref{pd-kappa-defects-segment-mean-theorem-a-free},
$(b,\omega)\in I(\delta,\alpha^1)$: $b$ is kept for $\alpha^1$ and $\omega$ is a level of $b$ for $\alpha^1$. As in the
proof of Lemma~\ref{pd-kappa-defects-segment-mean-theorem-a-top-fibre}, $\lambda(n^1)$ is the lettering of the
letters $x_1,y_1,\dots,x_s,y_s,x_+$ with the degrees $\delta'_\omega$ formed with $\alpha^1$. Below, $\delta$ and
$\delta'$ are the degrees formed with $\alpha^1$; so $\deg^1=\delta'\leq\delta$ letterwise.

\emph{One letter.} Let $1\leq v\leq s$ and let $\lambda$ be $x_v$ or $y_v$. Then $\delta'(\lambda)P_v\leq2a$. Indeed,
by Definition~\ref{pd-kappa-defects-segment-lattice}(3) and the conditions for $b$ to be kept, $\delta(x_v)$ is
$w_v\leq\hat b_v$ if $v\leq\eta_j$ and $o_v\leq\hat b_v$ if $v>\eta_j$, and $\delta(y_v)$ is at most $\hat\alpha^1_v$.
Here $\hat b_vP_v\leq b_v\leq\frac52b'\leq2a$ and $\hat\alpha^1_vP_v\leq\alpha^1_v\leq2a$.

\emph{The letter $x_+$.} As $b$ is kept, $\delta'(x_+)\leq\delta(x_+)=o_+\leq O_+$. We bound the three terms of $O_+$.
First, $\operatorname{cap}(s)\leq\hat a+\sum_v\hat b_v\leq a+s\cdot\frac52b'$, and as $s\leq(m+1)/2$,
$s\cdot\frac52b'\leq\frac54(m+1)b'\leq a/2$. Next, $\Lambda/(m+1)\leq1/4$ and $4/p^2\leq1/4$; and
$\sum_v2a/P_v\leq s\cdot2a/p\leq a/5$, as $10s\leq5m\leq(m+1)^2\leq p$; so the second term is at most $a/10$.
Finally $\Lambda b_+/(m+1)\leq\frac14\cdot\frac52b'\leq a/32$. So $\delta'(x_+)\leq a+a/2+a/10+a/32\leq2a$.

\emph{The degrees and the $z_i$.} By the last two paragraphs and $P_v\geq1$, $\deg^1\lambda_g\leq2a$ for every $g$.
Also $z_i=\hat b_{u_i}\dot-o_{u_i}\leq\hat b_{u_i}\leq b_{u_i}\leq\frac52b'\leq2a$.

\emph{The total.} For $g<2s$, $\lambda_g$ is $x_v$ or $y_v$ for some $v$, and $P_v\geq p$; so
$p\,\delta'(\lambda_g)\leq2a$. Summing over these $2s$ letters, $p\sum_{g<2s}\delta'(\lambda_g)\leq2s\cdot2a\leq2ap$, as
$2s\leq m\leq p$. So $\sum_g\deg^1\lambda_g\leq2a+\delta'(x_+)\leq4a\leq8a$.

\emph{The sums $S_i$.} Fix $i$, and put $u=u_i$ and $P=P_u$. The letters $\lambda_g$ with $g\leq\pi_i$ are
$x_1,y_1,\dots,x_{u-1},y_{u-1},x_u$. As $u>\eta_j$, $\delta'(x_u)\leq o_u\leq\hat b_u$, so $\delta'(x_u)P\leq
b_u\leq\frac52b'$. For $v<u$ and $\lambda$ equal to $x_v$ or $y_v$, $pP\leq P_v$, so
$\delta'(\lambda)P\leq\delta'(\lambda)P_v/p\leq2a/p$. There are fewer than $2s+1$ such letters, so
\[
 S_iP\leq\tfrac52b'+(2s+1)\frac{2a}p .
\]
On the other side, $n^0_i=\min N_i(\delta,\omega)$ is $(y-|\rho'_u|)/P-j_i$ for some $y\in W_u(\delta)$. Here $y\geq a$;
$|\rho'_u|\leq(m+1)P$, as the $m+1$ entries of $\rho'_u$ are less than $P$; and, $\omega$ being a level of $b$ for
$\alpha$ and $u>\eta_j$, $j_i\leq\delta(x_u)=o_u\leq\hat b_u$, so $j_iP\leq b_u\leq\frac52b'$. Hence
\[
 n^0_iP=y-|\rho'_u|-j_iP\geq a-(m+1)P-\tfrac52b'.
\]
Now $(m+1)P\leq b'$ and $b'\leq a/p$, so $5b'+(m+1)P\leq6a/p$; and $4s+8\leq2m+4\leq(m+1)^2\leq p$. So
$5b'+(m+1)P+(2s+1)\cdot2a/p\leq(4s+8)a/p\leq a$, which gives $S_iP\leq n^0_iP$, and $S_i\leq n^0_i$ as $P>0$.

Finally $a\geq1$ by the facts above.
\end{proof}
% Lean: restTMSize_proved hS hFree = restTMSize_of hS hFree (restSI_proved hS hFree), with hS = segFacts
%   (2 <= CYb, 1 <= Cb, 0 <= CD, 0 <= CE1, CE2) and hFree = restFree_proved. The facts are SegFacts (kP, P1,
%   La4, E0, bTop, kb) and HIFacts (m4, s_le, p_ge, a1, La0, b0, b_lt, bT_le, win_le), with r5, b_small,
%   p_le_P and P_step (segDists_lt_G, segDists_lt). The top lettering is setTies_top (top_eq, lLL_ovr) at
%   aTopPt, with aTopPt_mem, hFree and digits_of_idx. One letter is chain_wt and
%   lLL_deg_mul_segP_le_two_mul_a_of_kept (window_entry_le_two_mul_a, coarseA_upper); x_+ is oplus_le
%   (cap_le with B = 5b'/2 and base_exponent_le_a); the degrees are lLL_deg_le; z_i is coarseA_upper with
%   b_small; the total is the third clause of restTMSize_of. S_i <= n0 is restSI_proved: S_bound
%   (segChain_deg_even) for the upper bound, lower_bound_initial_multiplicity (cast_sub_ge, mem_lv) for
%   n0 P, and hsum with h48 (4s + 8 <= p) and hb'p (b' <= a/p).

\begin{lemma}\label{pd-kappa-defects-segment-mean-theorem-a-top-hyp}
Let $C_H\geq2$ and $C_{E1},C_{E2}\geq0$, and let an instance satisfy its hypotheses. Let $\delta$ be a digit class,
$\alpha\in A(\delta)$ and $(b,\omega)\in I(\delta,\alpha)$. Then the model instance attached to $(\delta,\alpha,b,\omega)$
satisfies the model hypotheses (Definitions~\ref{pd-kappa-defects-segment-mean-theorem-a-top-model}
and~\ref{pd-kappa-defects-segment-mean-theorem-a-top-attached}).
\end{lemma}
\leanref{Formalism/SharpRestriction/GrowthBounds/MonomialFactorials/ExponentMatrixFirstDepthEq.lean}{Formalism.MainNotions.G.SymmetricFormBlocks.RestTMHyp}
\leanref{Formalism/SharpRestriction/GrowthBounds/MonomialFactorials/RestTMHypOf.lean}{Formalism.MainNotions.G.SymmetricFormBlocks.restTMHyp_of}

\begin{proof}
As $p\geq5$, $p\neq2$, and $r=m\geq2N_1+2\geq2$ by hypothesis (1) of Theorem~\ref{pd-gain-theorem-one}; the
index is $\varrho=m-1=r-1$. The ties
$\pi_i=2u_i-2$ increase with $i$; $\pi_i+1=2u_i-1<2s+1$; $\lambda_{\pi_i}=x_{u_i}$ is a vector letter and
$\lambda_{\pi_i+1}=y_{u_i}$ a covector letter; and the condition $(2u_i-1,M_i)$ has $J=\pi_i+1$. The letters $x_k$ and
$x_+$ sit at the even positions and the $y_k$ at the odd ones, and every $J$ is $2s+1$ or some $\pi_i+1$. By
Lemma~\ref{pd-kappa-defects-segment-mean-theorem-a-top-sizes}, $a\geq1$. As $Y\geq1$, $s\geq0$, $m\geq2$ and $\ell\geq1$,
$E=Yc_S=3Y(s+2)m^2\ell\geq1$. As $W_{u_i}(\delta)$ is nonempty, $R_i\geq1$, and $h_i=1$ is positive and prime to $p$. By
Lemma~\ref{pd-kappa-defects-segment-mean-theorem-a-top-sizes}, $S_i\leq n^0_i$, and the sum of the $\deg^1\lambda_g$ is
at most $8a\leq8a+r$. For every $i$, $\deg^1\lambda_{\pi_i}+\deg^1\lambda_{\pi_i+1}+r+z_i+1\leq6a+r+1\leq8a+r$. As
$\lambda_{\pi_i+1}$ is the tied letter, $\deg^1\lambda_{\pi_i+1}=n^1_i=n^0_i+R_i-1\leq2a$; so $R_i\leq2a+1\leq8a+r$, and
$2S_i+1\leq2n^0_i+1\leq4a+1\leq8a+r$.
\end{proof}
% Lean: restTMHyp_of hH (restTMSize_proved ...): p_ne_two, r_ge_two, the admissibility from hpos, hlt, hmono and
%   hvec, Conds from segTies.J, 1 <= a from RestTMSize, 1 <= Y kapS (lgA_ge_one), hYl1 and hdd for the
%   progression, and the size bounds through key, hd, htop (setTies_deg_succ_pos_eq_of_strictMono) and htop'.

\begin{lemma}\label{pd-kappa-defects-segment-mean-theorem-a-top-eps}
Let $C_0\geq1$, $C_{E1},C_{E2}\geq0$ and $C_H\geq2$, and let an instance satisfy its hypotheses. Let $\delta$ be a digit
class, $\alpha\in A(\delta)$ and $(b,\omega)\in I(\delta,\alpha)$. Then the numbers $\varepsilon_i$ of the model instance
attached to $(\delta,\alpha,b,\omega)$ are at most $31/Y$ (Definitions~\ref{pd-kappa-defects-segment-mean-theorem-a-top-model}
and~\ref{pd-kappa-defects-segment-mean-theorem-a-top-attached}).
\end{lemma}
\leanref{Formalism/SharpRestriction/GrowthBounds/MonomialFactorials/ExponentMatrixFirstDepthEq.lean}{Formalism.MainNotions.G.SymmetricFormBlocks.RestTMEps}
\leanref{Formalism/SharpRestriction/GrowthBounds/MonomialFactorials/EpsilonBoundMConstant.lean}{Formalism.MainNotions.G.SymmetricFormBlocks.restTMEps_proved}

\begin{proof}
Write $\delta=(\rho,\rho')$, fix $i$, and put $u=u_i$, $P=P_u$, $\sigma=|\rho'_u|$ and
$u_W=\lfloor a+a/(kY)+x_a\rfloor$, so $W=[a,u_W]\cap\mathbf Z$. Let $\mathcal M$ be the attached model instance, and
let $\alpha^1$ be $\alpha$ with each entry $\alpha_{u_l}$ replaced by $\tau(\delta,l)$; as in the proof of
Lemma~\ref{pd-kappa-defects-segment-mean-theorem-a-top-fibre}, $\lambda(n^1)$ is the lettering of the degrees
$\delta'_\omega$ formed with $\alpha^1$, and $b$ is kept for $\alpha^1\in A(\delta)$.

\emph{Facts.} As in the proof of Lemma~\ref{pd-kappa-defects-segment-mean-theorem-a-free}, $(m+1)P_v\leq b'$ and
$b_v\leq5b'/2$ for every $v$, $m\geq4$, $(m+1)^2\leq p$ and $a>pb'$, so $(m+1)b'\leq a/4$; and $2s+2\leq m$, as
$s\leq N_1$ and $2N_1+2\leq m$. By hypothesis (1) of Theorem~\ref{pd-gain-theorem-one}, $Y\geq1$ and
$k^2Y^2\leq p$. By its hypothesis (2), $kY^2Y_b|T|\Lambda\log a\leq a/b'$; as $Y_b\geq C_H\geq2$ and $|T|$,
$\Lambda=C_0N_1^2\log a$ and $\log a$ are at least $1$ ($a\geq100$), $kY^2b'\leq a/2$. Hence $4Yb'\leq kY^2b'\leq
a/2$. As in the proof of Lemma~\ref{pd-kappa-defects-segment-mean-theorem-a-ratio}, every element of $W$ is at most
$2a$.

\emph{The length.} We show $R_iP\geq a/(2kY)$. First, $(1+P)(m+1)\leq2b'$, so
$(1+P)\cdot2kY\leq4Yb'\leq a$, that is $1+P\leq a/(2kY)$; in the same way $(2P+1)k\leq3b'$ and $3b'Y\leq a$ give
$a+2P\leq a+a/(kY)-1<u_W$. As $\sigma<(m+1)P\leq b'<a$, there are $y_-\in[a,a+P)$ and $y_+\in(u_W-P,u_W]$ with
$y_\pm\equiv\sigma\pmod P$; both lie in $W_u(\delta)$, and $y_-<y_+$. Their images under the map $\theta_i$ of the
proof of Lemma~\ref{pd-kappa-defects-segment-mean-theorem-a-top-fibre} lie in $N_i(\delta,\omega)$ and differ by
$(y_+-y_-)/P$, so $R_iP\geq y_+-y_-+P\geq u_W-a-P>a/(kY)-1-P\geq a/(2kY)$.

\emph{The sum $S_i$.} The letters $\lambda_g$ with $g\leq\pi_i$ are $x_1,y_1,\dots,x_{u-1},y_{u-1},x_u$, with the
degrees $\delta'_\omega\leq\delta$ formed with $\alpha^1$. For $x_u$: as $u>\eta_j$ and $b$ is kept,
$\delta(x_u)=o_u\leq\hat b_u\leq b_u/P$, so its degree times $P$ is at most $b_u\leq5b'/2$. For $x_v$ or $y_v$ with
$v<u$: by Definition~\ref{pd-kappa-defects-segment-lattice}(3) and the kept conditions, the degree is at most
$\hat b_v\leq b_v/P_v$ or at most $\hat\alpha^1_v\leq\alpha^1_v/P_v$, and $b_v,\alpha^1_v\leq2a$; and
$P_v=p^{g_j-d_{j,v}}\geq pP$ as $d_{j,v}<d_{j,u}$. So its degree times $P$ is at most $2a/p$. Hence
$S_iP\leq\frac52b'+(2s+1)\frac{2a}p$.

\emph{The ratio.} As $2s+1\leq k$ and $k^2Y^2\leq p$, $(2s+1)\frac{2a}p\cdot2kY^2\leq4a$; and
$\frac52b'\cdot2kY^2\leq\frac52a$. So $(\frac52b'+(2s+1)\frac{2a}p)Y\leq\frac{13a}{4kY}\leq7\cdot\frac a{2kY}$, and by
the last two paragraphs $S_iPY\leq7R_iP$, that is $S_i/R_i\leq7/Y$.

\emph{The bound.} For $\mathcal M$, $\kappa=c_S=3(s+2)m^2\ell\geq24\ell$, as $m\geq2$, and $E=Y\kappa$. So
\[
 \varepsilon_i=2\frac{S_i}{R_i}+\frac2Y+\frac2Y\cdot\frac{S_i}{R_i}\cdot\frac\ell\kappa\leq\frac{14}Y+\frac2Y
 +\frac2Y\cdot7\cdot\frac1{24}\leq\frac{17}Y\leq\frac{31}Y.
\]
\end{proof}
% Lean: restTMEps_proved hCYb hC0 hH hS hFree 1 le_rfl = restTMEps_of hS (epsFacts hCYb hC0 hH hS) hFree, which
%   gives (30 + CM)/Y for every CM >= 1 (the book takes CM = 1). The facts are SegFacts (kP, P1, kb, Y1, xa0, bTop)
%   and EpsFacts (pk, and kYb from hlog with Yb N La log a >= 2); the length is hroomR, hwin, hroom and Yl_bound
%   (exists_low, exists_high, lvlRoom_of) with hYP; the sum is S_bound at aTopPt (setTies_top, digits_of_idx;
%   lLL_deg_le_chain, segChain_deg_even, lLL_deg_mul_segP_le_two_mul_a_of_kept, P_step); the ratio is hfin and
%   hSY; the bound is epsilon_bound_M_constant = eps_num with CM = 2, which gives (15 + 2)/Y.

\begin{lemma}\label{pd-kappa-defects-segment-mean-theorem-a-top-fibre}
Let $C_0\geq1$, $C_{E1},C_{E2}\geq0$ and $C_H\geq2$, and let an instance satisfy its hypotheses. Let $\delta$ be a digit
class, $\alpha_0\in A(\delta)$ and $(b,\omega)\in I(\delta,\alpha_0)$, and put $D_1=D_1(1)$
(Definition~\ref{pd-kappa-defects-segment-mean-theorem-a-levels}). For $\alpha\in A(\delta)$ let $F(\alpha)$ be the set
of the $\alpha'\in A(\delta)$ with $\alpha'_k=\alpha_k$ for every $k$ other than $u_1,\dots,u_t$. Then for every
$\alpha\in A(\delta)$,
\[
 \sum_{\alpha'\in F(\alpha)}\operatorname{bad}^Y_\omega(\alpha',b;\lfloor D_1\rfloor)\leq\frac{31}Y\sum_{i=1}^t
 \sum_{\alpha'\in F(\alpha)}\dim Y_\omega(\alpha'^{(i)}),
\]
with $\operatorname{bad}^Y_\omega(\alpha',b;\cdot)$ formed with $\alpha'$ and $Y_\omega(\alpha'^{(i)})$ with $\alpha'^{(i)}$.
\end{lemma}
\leanref{Formalism/SharpRestriction/GrowthBounds/MonomialFactorials/ExponentMatrixFirstDepthEq.lean}{Formalism.MainNotions.G.SymmetricFormBlocks.fibre_ineq}
\leanref{Formalism/SharpRestriction/GrowthBounds/MonomialFactorials/ExponentMatrixFirstDepthEq.lean}{Formalism.MainNotions.G.SymmetricFormBlocks.prog_eq}
\leanref{Formalism/SharpRestriction/GrowthBounds/MonomialFactorials/DimYOvr.lean}{Formalism.MainNotions.G.SymmetricFormBlocks.bY_ovr}
\leanref{Formalism/SharpRestriction/GrowthBounds/MonomialFactorials/DimYOvr.lean}{Formalism.MainNotions.G.SymmetricFormBlocks.dimY_ovr}

\begin{proof}
Write $\delta=(\rho,\rho')$. By Lemmas~\ref{pd-kappa-defects-segment-mean-theorem-a-top-box}
and~\ref{pd-kappa-defects-segment-mean-theorem-a-free}, $A(\delta)=\prod_kW_k(\delta)$, $(b,\omega)\in I(\delta,\alpha)$,
and the model instance $\mathcal M$ attached to $(\delta,\alpha,b,\omega)$ is defined. For $1\leq i\leq t$ and
$y\in W_{u_i}(\delta)$ put $\theta_i(y)=(y-|\rho'_{u_i}|)/P_{u_i}-j_i$.

\emph{The progression.} The set $W_{u_i}(\delta)$ consists of the terms of an arithmetic progression with difference
$P_{u_i}$ that lie in the interval $W$, and $\theta_i(y)\geq0$ by
Lemma~\ref{pd-kappa-defects-segment-mean-theorem-a-top-box}(2). So $\theta_i$ is increasing and injective on
$W_{u_i}(\delta)$, and its image $N_i(\delta,\omega)$ consists of consecutive integers; it is
$\{n^0_i,\dots,n^0_i+R_i-1\}$ with $h_i=1$. As $F(\alpha)$ is the set of the $\alpha'$ obtained from $\alpha$ by
replacing the entries $\alpha_{u_i}$ by elements of $W_{u_i}(\delta)$, the map
$\theta\colon\alpha'\mapsto(\theta_i(\alpha'_{u_i}))_i$ is a bijection from $F(\alpha)$ onto the progression $P$ of
$\mathcal M$. As $\tau(\delta,i)$ is the largest element of $W_{u_i}(\delta)$, $\theta(\alpha'^{(i)})=\theta(\alpha')[i]$.

\emph{The letters.} Let $\alpha'\in F(\alpha)$ and $n=\theta(\alpha')$. Of the degrees $\delta'_\omega$ formed with
$\alpha'$, only those of the $y_{u_i}$ involve the entries $\alpha'_{u_i}$, and as $u_i>\eta_j$,
$\delta'_\omega(y_{u_i})=\hat\alpha'_{u_i}-j_i=n_i$. The ties, the numbers $z_i$ and the conditions do not involve
$\alpha'$. So the letters with the degrees $\delta'_\omega$ formed with $\alpha'$ are the lettering $\lambda(n)$ of
$\mathcal M$; hence $Y_\omega$, $Y^{\mathbf Z}_\omega$ and $\beta^Y_\omega$ formed with $\alpha'$ are $Y(\lambda(n))$,
$Y^{\mathbf Z}(\lambda(n))$ and $\beta(\lambda(n))$, and $\operatorname{bad}^Y_\omega(\alpha',b;D)=
\operatorname{bad}_D(\beta(\lambda(n)),Y^{\mathbf Z}(\lambda(n)))$ for every $D$. Likewise $Y_\omega(\alpha'^{(i)})=
Y(\lambda(n[i]))$.

\emph{The estimate.} For $\mathcal M$, $\kappa=3(s+2)m^2\ell=c_S$, so $D_1(\mathcal M)=t(Yc_S+c_S)+c_S=D_1(1)$. By
Lemma~\ref{pd-kappa-defects-segment-mean-theorem-a-top-hyp}, $\mathcal M$ satisfies the model hypotheses, and $P$ is
nonempty. Lemma~\ref{pd-kappa-defects-segment-mean-theorem-a-top-mean}, multiplied by $|P|$, and
Lemma~\ref{pd-kappa-defects-segment-mean-theorem-a-top-eps} give
\[
 \sum_{n\in P}\operatorname{bad}_{\lfloor D_1\rfloor}\big(\beta(\lambda(n)),Y^{\mathbf Z}(\lambda(n))\big)\leq
 \sum_{i=1}^t\varepsilon_i\sum_{n\in P}\dim Y(\lambda(n[i]))\leq\frac{31}Y\sum_{i=1}^t\sum_{n\in P}\dim Y(\lambda(n[i])).
\]
Transported to $F(\alpha)$ by the bijection $\theta$, this is the claim.
\end{proof}
% Lean: fibre_ineq D eps a b hb hr hne (hM _) (hH ...) (hE ...): h1 = hM hH, prog_eq (interval_eq, Nl_conv) and
%   exponent_matrix_first_depth_eq for D1; reindex (cm_inj) moves the sums from the piFinset of the Wl to that
%   of the Nl; bY_ovr and dimY_ovr (lLL_ovr) identify the badness and the dimensions; ovr_update and top_eq
%   identify alpha'^(i) with the update at the top. The fibre is ovr (uTop S) a n for n in the piFinset of the
%   Wl (uTop i).

\begin{lemma}\label{pd-kappa-defects-segment-mean-theorem-a-top}
There is $C_M\geq1$ with the following property. Let $C_0\geq1$, $C_{E1},C_{E2}\geq0$ and $C_H\geq2$, and let
an instance satisfy its hypotheses. Let $\delta$ be a digit class, $\alpha_0\in A(\delta)$ and
$(b,\omega)\in I(\delta,\alpha_0)$, so that $(b,\omega)\in I(\delta,\alpha)$ for every $\alpha\in A(\delta)$
(Lemma~\ref{pd-kappa-defects-segment-mean-theorem-a-free}), and put $D_1=D_1(C_M)$. Then
\[
 \sum_{\alpha\in A(\delta)}\operatorname{bad}^Y_\omega(\alpha,b;\lfloor D_1\rfloor)\leq\frac{30+C_M}Y
 \sum_{i=1}^t\sum_{\alpha\in A(\delta)}\dim Y_\omega(\alpha^{(i)}),
\]
where $\operatorname{bad}^Y_\omega(\alpha,b;\cdot)$ is formed with $\alpha$ and $Y_\omega(\alpha^{(i)})$ with
$\alpha^{(i)}$ (Definition~\ref{pd-kappa-defects-segment-mean-theorem-a-levels}).
\end{lemma}
\leanref{Formalism/SharpRestriction/GrowthBounds/MonomialFactorials/BlockBoundUniformiser.lean}{Formalism.MainNotions.G.SymmetricFormBlocks.RestTM}
\leanref{Formalism/SharpRestriction/GrowthBounds/MonomialFactorials/EpsilonBoundMConstant.lean}{Formalism.MainNotions.G.SymmetricFormBlocks.restTM_closed}
\leanref{Formalism/SharpRestriction/GrowthBounds/MonomialFactorials/RestTMExists.lean}{Formalism.MainNotions.G.SymmetricFormBlocks.restTM_exists}

\begin{proof}
Take $C_M=1$, so $D_1=D_1(1)$ and $(30+C_M)/Y=31/Y$. Let the constants, the instance, $\delta$, $\alpha_0$ and
$(b,\omega)$ be as in the statement. By Lemma~\ref{pd-kappa-defects-segment-mean-theorem-a-top-box}(1),
$A(\delta)=\prod_kW_k(\delta)$. For $\alpha\in A(\delta)$ let $F(\alpha)$ be as in
Lemma~\ref{pd-kappa-defects-segment-mean-theorem-a-top-fibre}. Every $F(\alpha)$ has
$f=\prod_i|W_{u_i}(\delta)|\geq1$ elements, and for $\alpha'\in A(\delta)$ the $\alpha\in A(\delta)$ with $\alpha'\in F(\alpha)$
are the $f$ elements that agree with $\alpha'$ outside the entries $u_1,\dots,u_t$. So for every function $G$ on
$A(\delta)$, $\sum_{\alpha\in A(\delta)}\sum_{\alpha'\in F(\alpha)}G(\alpha')=f\sum_{\alpha'\in A(\delta)}G(\alpha')$. Summing
the inequality of Lemma~\ref{pd-kappa-defects-segment-mean-theorem-a-top-fibre} over $\alpha\in A(\delta)$ and applying
this to $G(\alpha')=\operatorname{bad}^Y_\omega(\alpha',b;\lfloor D_1\rfloor)$ and to $G(\alpha')=\dim Y_\omega(\alpha'^{(i)})$
gives the claim multiplied by $f$; dividing by $f$ gives the claim.
\end{proof}
% Lean: restTM_closed = restTM_exists (CM = 1: restTM_of with alternative_mean_value_theorem = MC.Top.theoremMV7)
%   applied to restFree_proved (segFacts), restTMSize_proved and restTMEps_proved, under 2 <= CYb, 1 <= Cb,
%   0 <= CD, 1 <= C0 and 0 <= CE1, CE2; restTM_of = restTM_of_rests hM hF (restTMHyp_of hH hS) hE, which is
%   box_of_fibres (sum_ovr) over cbox = piFinset Wl (cbox_eq) with fibre_ineq (Lemma
%   pd-kappa-defects-segment-mean-theorem-a-top-fibre) for each a in the box, and lvlRoom_of.

\begin{lemma}\label{pd-kappa-defects-segment-mean-theorem-a-directions-independent}
Let an instance be given, $\alpha\in W^{S_j}$, $\alpha_+\in W$ and $b$ a block, and form the generators of
Definition~\ref{pd-kappa-defects-segment-lattice}(5) over $\mathbf Q_p$, as in
Definition~\ref{pd-kappa-defects-segment-mean-theorem-a-levels}. For a level $\omega$ of $b$ let
$V^{\mathbf Q}_\omega$ be the $\mathbf Q_p$-span of the generators of level $\omega$. Then the sum of the
$V^{\mathbf Q}_\omega$ is direct: if $v_\omega\in V^{\mathbf Q}_\omega$ for the levels $\omega$ in a finite set
$\Omega$ and $\sum_{\omega\in\Omega}v_\omega=0$, then every $v_\omega$ is $0$.
\end{lemma}
\leanref{Formalism/MahlerMeasure/MahlerMeasureProofs/SSDirect.lean}{Formalism.GM.M.ss_direct}
\leanref{Formalism/MahlerMeasure/MahlerMeasureProofs/SSDirect.lean}{Formalism.GM.M.SSD.absDirect}
\leanref{Formalism/MahlerMeasure/MahlerMeasureProofs/SSDirect.lean}{Formalism.GM.M.topV}

\begin{proof}
Write $t=S_j-\eta_j$. For a monomial $n$ in the variables of $\mathcal P$ let $\deg'(n)\in\mathbf Z_{\geq0}^{t+1}$
be the vector of its exponents at $z'_1,\dots,z'_t,z'_+$; order such vectors entrywise, and let $|e|$ be the
sum of the entries of $e$. Identify a level $\omega=(j,j_+)$ with the vector $(j_1,\dots,j_t,j_+)$.

\emph{Step 1: one level.} Let $\omega=(j,j_+)$ be a level and $F_\omega=\prod_l\hat E_l^{j_l}\cdot\hat E_+^{j_+}$.
A generator of level $\omega$ is $F_\omega Q$, where $Q=\prod_lz_l^{\hat b_{u_l}\dot-o_{u_l}}\cdot
z_+^{\alpha_+\dot-o_+}fg$ contains none of the variables $z'_l$, $z'_+$. So every $v\in V^{\mathbf Q}_\omega$ is
$F_\omega Q$ with $Q$ in the span of such products, again free of these variables. Now
$\hat E_l=z_lz'_l+\sum_cx_{u_l,c}y_{u_l,c}$ and $\hat E_+=z_+z'_++\sum_cx_{+,c}\xi_c$, where the sums are free of the
$z'$. Expanding the product, $F_\omega=m_\omega+R_\omega$, where
$m_\omega=\prod_l(z_lz'_l)^{j_l}\cdot(z_+z'_+)^{j_+}$ has $\deg'(m_\omega)=\omega$, and every monomial of $R_\omega$
has $\deg'\leq\omega$ and $\deg'\neq\omega$. As $Q$ is free of the $z'$, it follows that every monomial of $v$
has $\deg'\leq\omega$, and that for every monomial $n$ free of the $z'$ the coefficient of $v$ at $m_\omega n$
is the coefficient of $Q$ at $n$. So if the coefficients of $v$ at all monomials $n'$ with
$\deg'(n')=\omega$ vanish, then $Q=0$ and $v=0$.

\emph{Step 2: the sum.} We argue by induction on $|\Omega|$; for $\Omega$ empty there is nothing to show.
Choose $\omega\in\Omega$ with $|\omega|$ largest, and let $n'$ be a monomial with $\deg'(n')=\omega$. If
$v_{\omega'}$, with $\omega'\in\Omega$, has a nonzero coefficient at $n'$, then $\omega\leq\omega'$ by Step 1;
as $|\omega'|\leq|\omega|$, $\omega'=\omega$. So the coefficient of $\sum_{\omega'\in\Omega}v_{\omega'}=0$ at $n'$ is
that of $v_\omega$, and by Step 1, $v_\omega=0$. The induction hypothesis for $\Omega\setminus\{\omega\}$ gives
the rest.
\end{proof}
% Lean: ss_direct D = absDirect (scalS) with a = lv (lv_inj): Step 1 is topV (memV and gens_sub write
%   v = pre * Q with Q in the span of Bset; good_pre from good_tie, Good.mul, Good.pow, Good.prod;
%   sfree_span; Good.mul_sfree) and Q_of_top; Step 2 is absDirect with topcoeff and eq_of_le_sum.
%   ss_direct holds for every block datum over Q_p; the kept condition and BlockOK are not used.

\begin{lemma}\label{pd-kappa-defects-segment-mean-theorem-a-directions}
Let an instance be given, $(\alpha,\alpha_+)\in W^{S_j}\times W$ and $b$ a block kept for $\alpha$. Then the sum
of $\dim Y_\omega\cdot N_\xi(b_+-j_+)$ over the levels $\omega=(j,j_+)$ of $b$ is at most $\dim_{\mathcal K}V_b$
(Definitions~\ref{pd-kappa-defects-segment-mean-theorem-a-classes}
and~\ref{pd-kappa-defects-segment-mean-theorem-a-levels}).
\end{lemma}
\leanref{Formalism/SharpRestriction/GrowthBounds/MonomialFactorials/BlockBoundUniformiser.lean}{Formalism.MainNotions.G.SymmetricFormBlocks.RestDir}
\leanref{Formalism/SharpRestriction/GrowthBounds/MonomialFactorials/RestDirOf.lean}{Formalism.MainNotions.G.SymmetricFormBlocks.restDir_of}
\leanref{Formalism/SharpRestriction/GrowthBounds/MonomialFactorials/LvlDimGen.lean}{Formalism.MainNotions.G.SymmetricFormBlocks.lvlDim_gen}
\leanref{Formalism/SharpRestriction/GrowthBounds/MonomialFactorials/SumFinrankLeOfIndep.lean}{Formalism.MainNotions.G.SymmetricFormBlocks.restSumSpan_proved}
\leanref{Formalism/SharpRestriction/GrowthBounds/MonomialFactorials/RestDimBCProved.lean}{Formalism.MainNotions.G.SymmetricFormBlocks.restDimBC_proved}

\begin{proof}
Form the generators over $\mathbf Q_p$ and the spaces $V^{\mathbf Q}_\omega$ as in
Lemma~\ref{pd-kappa-defects-segment-mean-theorem-a-directions-independent}, and let $V^{\mathbf Q}_b$ be the
$\mathbf Q_p$-span of the generators over $\mathbf Q_p$ of all levels of $b$.

\emph{Step 1: comparison with $V_b$.} Let $\iota_p$ and $P\mapsto P^\iota$ be as in
Lemma~\ref{pd-kappa-defects-segment-mean-theorem-a-block-bound-base}. By parts (1) and (2) of that lemma,
$P\mapsto P^\iota$ maps $V^{\mathbf Q}_b$ into $V_b$ and maps $\mathbf Q_p$-linearly independent elements of
$V^{\mathbf Q}_b$ to $\mathcal K$-linearly independent elements of $V_b$. As $V_b$ is finite-dimensional
(Lemma~\ref{pd-kappa-defects-windows-ranks}), so is $V^{\mathbf Q}_b$, and
$\dim_{\mathbf Q_p}V^{\mathbf Q}_b\leq\dim_{\mathcal K}V_b$.

\emph{Step 2: one level.} Let $\omega=(j,j_+)$ be a level, $d=\dim Y_\omega$ and $y_1,\dots,y_d$ a basis of the
finite-dimensional space $Y_\omega$. Multiplying each $y_a$ by a power of $p$, we may assume that its
coefficients lie in $\mathbf Z_p$, so $y_a\in Y^{\mathbf Z}_\omega$. Let $F'_\omega=F_\omega\prod_l
z_l^{\hat b_{u_l}\dot-o_{u_l}}\cdot z_+^{\alpha_+\dot-o_+}$, with $F_\omega$ as in the proof of
Lemma~\ref{pd-kappa-defects-segment-mean-theorem-a-directions-independent}, and let $\xi^\beta$ run over the
$N_\xi(b_+-j_+)$ monomials of degree $b_+-j_+$ in $\xi_1,\dots,\xi_m$. The $F'_\omega y_a\xi^\beta$ are
generators of level $\omega$, and they are linearly independent. Indeed, the $y_a$ involve only the variables
of the letters, and the $\xi_c$ are not among them; so for each $\beta$ and each monomial $n$ in the variables
of the letters, the coefficient of $\sum_{a,\beta}c_{a\beta}y_a\xi^\beta$ at $n\xi^\beta$ is the coefficient of
$\sum_ac_{a\beta}y_a$ at $n$, and the $y_a\xi^\beta$ are linearly independent. Multiplication by $F'_\omega$ is
injective, as $F'_\omega\neq0$: by Step 1 of the proof of
Lemma~\ref{pd-kappa-defects-segment-mean-theorem-a-directions-independent}, $F_\omega$ has the coefficient $1$ at
$m_\omega$. So $\dim Y_\omega\cdot N_\xi(b_+-j_+)\leq\dim V^{\mathbf Q}_\omega$.

\emph{Step 3: the sum.} Each $V^{\mathbf Q}_\omega$ lies in $V^{\mathbf Q}_b$, which is finite-dimensional by
Step 1, and by Lemma~\ref{pd-kappa-defects-segment-mean-theorem-a-directions-independent} their sum is direct.
So the sum of the $\dim V^{\mathbf Q}_\omega$ over the levels of $b$ is the dimension of a subspace of
$V^{\mathbf Q}_b$. With Steps 2 and 1, the sum of the $\dim Y_\omega\cdot N_\xi(b_+-j_+)$ is at most
$\dim V^{\mathbf Q}_b\leq\dim_{\mathcal K}V_b$.
\end{proof}
% Lean: restDir_of (restDirQ_of lvlDimCondition restSumSpan_proved) restDimBC_proved. Step 1 is
%   restDimBC_proved and span_SSZ_fin, with iotaX, map_SSZ and li_map (Lemma
%   pd-kappa-defects-segment-mean-theorem-a-block-bound-base). Step 2 is lvlDimCondition = lvlDim_gen with den_Y
%   (exists_pow_smul_intPol), preF_ne_zero (good_pre), xiMons_free, xiMons_mem and card_xiMons_eq, through
%   li_rename_mon (decode, decode_gen) and li_mul_left. Step 3 is restSumSpan_proved = sum_gens_le with
%   semisimple_direct_sum (ss_direct) and sum_finrank_le_of_indep. The kept condition is not used; the
%   levels lv aTop b are those of gB aTop 0 b, equal to those of gB aTop aPlus b (lv_eq).

\begin{lemma}\label{pd-kappa-defects-segment-mean-theorem-a-segment}
There is $C\geq1$ with the following property. Let $C_c\geq0$, $C_0\geq1$, $C_{E1},C_{E2}\geq1$ and $C_H\geq2$
be real numbers, $C_\kappa$ and $C_s$ any real numbers, and let an instance with these constants satisfy its
hypotheses and condition (11) of Definition~\ref{pd-kappa-defects-segment-mean-conditions}. Put
$t=S_j-\eta_j$, suppose that $C(t+1)N_1<Y$, and form $D_A$ with $C_A=C$
(Definition~\ref{pd-kappa-defects-numbers}). Then there are integers $g(\alpha,\alpha_+,b)\geq0$, for
$(\alpha,\alpha_+)\in W^{S_j}\times W$ and blocks $b$, and a real number $\varepsilon\geq0$ with
$2\varepsilon e^3\leq C(t+1)N_1/Y$, such that, in the notation of
Definition~\ref{pd-kappa-defects-segment-mean-theorem-a-classes}:
\begin{enumerate}
\item for all $(\alpha,\alpha_+)\in W^{S_j}\times W$ and every block $b$ kept for $\alpha$, $L_b$ has a good
  sublattice $L_2$ of depth $D_A$ for $\beta_b$ with $\operatorname{rk}L_b\leq\operatorname{rk}L_2+g(\alpha,\alpha_+,b)$;
\item for every digit class $\delta$,
  \[
   \sum_{\alpha\in A(\delta)}\sum_{\alpha_+\in W}\sum_bg(\alpha,\alpha_+,b)\leq
   \varepsilon\sum_{\alpha\in A(\delta)}\sum_{\alpha_+\in W}d(\delta;\alpha),
  \]
  the inner sum on the left over the blocks $b$ kept for $\alpha$ of class $\delta$.
\end{enumerate}
\end{lemma}
\leanref{Formalism/SharpRestriction/GrowthBounds/CountingStatementReduction.lean}{Formalism.MainNotions.G.CountingStatementReduction.SegmentWeightedCountingBound}
\leanref{Formalism/SharpRestriction/GrowthBounds/CountingStatementReduction.lean}{Formalism.MainNotions.G.CountingStatementReduction.DigitClassRemainingUniform}
\leanref{Formalism/SharpRestriction/GrowthBounds/MonomialFactorials/SegmentWeightedCountingBoundOfAllRemainingBounds.lean}{Formalism.MainNotions.G.SymmetricFormBlocks.rest1DAE_of}
\leanref{Formalism/SharpRestriction/GrowthBounds/MonomialFactorials/ClosedOfOpenConditions.lean}{Formalism.MainNotions.G.SymmetricFormBlocks.restAllE_of_open}
\leanref{Formalism/SharpRestriction/GrowthBounds/MonomialFactorials/SegmentWeightedCountingBoundOfAllRemainingBounds.lean}{Formalism.MainNotions.G.SymmetricFormBlocks.segment_weighted_counting_bound_of_all_remaining_bounds}
\leanref{Formalism/SharpRestriction/GrowthBounds/MonomialFactorials/SegmentWeightedCountingBoundOfAllRemainingBounds.lean}{Formalism.MainNotions.G.SymmetricFormBlocks.digit_class_dimension_comparison_of_remaining_estimates}
\leanref{Formalism/SharpRestriction/GrowthBounds/MonomialFactorials/Imports.lean}{Formalism.MainNotions.G.SymmetricFormBlocks.assemble}

\begin{proof}
\emph{The constants.} Let $C_M\geq1$ be the constant of
Lemma~\ref{pd-kappa-defects-segment-mean-theorem-a-top}, and put
$C=\max\big(2((30+C_M)e^3+5)e^3,\,18(1+C_M)+20\big)$. Let the constants and the instance be as in the
statement, write $s=S_j$, and use the notation of
Definitions~\ref{pd-kappa-defects-segment-mean-theorem-a-classes}
and~\ref{pd-kappa-defects-segment-mean-theorem-a-levels}. Put $D_1=D_1(C_M)$, $D'=\lfloor YD_1\rfloor$,
\[
 g(\alpha,\alpha_+,b)=g_{D'}(\alpha,\alpha_+,b),\qquad\varepsilon=\frac{(30+C_M)te^3}Y+\frac5Y .
\]
By Lemma~\ref{pd-kappa-defects-segment-mean-theorem-a-depth}, $D_1\geq0$ and $YD_1+D_{\mathrm{lv}}\leq D_A$; so
$D'$ is an integer with $0\leq D'\leq YD_1$. As $Y\geq1$, $\varepsilon\geq0$; and as
$(30+C_M)te^3+5\leq((30+C_M)e^3+5)(t+1)$ and $N_1\geq1$,
\[
 2\varepsilon e^3=\frac{2((30+C_M)te^3+5)e^3}Y\leq\frac{2((30+C_M)e^3+5)e^3(t+1)}Y\leq\frac{C(t+1)N_1}Y .
\]

\emph{Property (1).} Lemma~\ref{pd-kappa-defects-segment-mean-theorem-a-block-bound} with $D=D'$ gives a good
sublattice $L_2$ of $L_b$ of depth $D'+D_{\mathrm{lv}}$ with $\operatorname{rk}L_b\leq\operatorname{rk}L_2+
g(\alpha,\alpha_+,b)$. As $D'+D_{\mathrm{lv}}\leq D_A$, $L_2$ is also good of depth $D_A$.

\emph{Property (2).} Fix a digit class $\delta=(\rho,\rho')$ and write $A=A(\delta)$. If $A$ is empty, both
sides are $0$. Otherwise fix $\alpha_0\in A$ and put $I=I(\delta,\alpha_0)$; by
Lemma~\ref{pd-kappa-defects-segment-mean-theorem-a-free}, $I(\delta,\alpha)=I$ for every $\alpha\in A$. So for
$\alpha\in A$ and $\alpha_+\in W$ the inner sum on the left is the sum of the
$\operatorname{bad}_\omega(\alpha,\alpha_+,b;D')$ over $(b,\omega)\in I$. For $(b,\omega)\in I$, with
$\omega=(j,j_+)$, and $\alpha\in A$, put $e(b,\omega)=N_\xi(b_+-j_+)\geq0$, and let $y_\alpha(b,\omega)=\dim Y_\omega$
and $q_\alpha(b,\omega)=\operatorname{bad}^Y_\omega(\alpha,b;\lfloor D_1\rfloor)$, formed with $\alpha$.

(a) For $\alpha\in A$, $\sum_{(b,\omega)\in I}y_\alpha(b,\omega)e(b,\omega)\leq d(\delta;\alpha)$. Indeed, fix
$\alpha_+\in W$. By Lemma~\ref{pd-kappa-defects-segment-mean-theorem-a-directions}, for each block $b$ kept for
$\alpha$ of class $\delta$ the sum over its levels is at most $\dim V_b$. The sum of $\dim V_b$ over these
blocks is $r_j(\mu;\rho,\rho')$ for the window function $a_T^{\alpha,\alpha_+}$
(Definition~\ref{pd-kappa-defects-segment-mean-conditions}), for which $\rho,\rho'$ are digits for $j$ and
$d_j(\rho,\rho')=d(\delta;\alpha)$; as $0\leq\mu\leq b'/(kQY)$, condition (11) bounds it by $d(\delta;\alpha)$.

(b) By Lemma~\ref{pd-kappa-defects-segment-mean-theorem-a-plus}, for $\alpha\in A$,
\[
 \sum_{\alpha_+\in W}\sum_{(b,\omega)\in I}\operatorname{bad}_\omega(\alpha,\alpha_+,b;D')\leq
 |W|\sum_{(b,\omega)\in I}\Big(q_\alpha(b,\omega)+\frac5Yy_\alpha(b,\omega)\Big)e(b,\omega).
\]

(c) By Lemma~\ref{pd-kappa-defects-segment-mean-theorem-a-top}, for $(b,\omega)\in I$,
$\sum_{\alpha\in A}q_\alpha(b,\omega)\leq\frac{30+C_M}Y\sum_{i=1}^t\sum_{\alpha\in A}y_{\alpha^{(i)}}(b,\omega)$; and by
Lemma~\ref{pd-kappa-defects-segment-mean-theorem-a-ratio}, $\alpha^{(i)}\in A$ and
$d(\delta;\alpha^{(i)})\leq e^3d(\delta;\alpha)$.

By (c) and then (a) at the $\alpha^{(i)}$,
\[
 \sum_{\alpha\in A}\sum_{(b,\omega)\in I}q_\alpha(b,\omega)e(b,\omega)\leq\frac{30+C_M}Y\sum_{i=1}^t\sum_{\alpha\in A}
 \sum_{(b,\omega)\in I}y_{\alpha^{(i)}}(b,\omega)e(b,\omega)\leq\frac{30+C_M}Y\sum_{i=1}^t\sum_{\alpha\in A}d(\delta;\alpha^{(i)}),
\]
which is at most $\frac{(30+C_M)te^3}Y\sum_{\alpha\in A}d(\delta;\alpha)$; and by (a),
$\frac5Y\sum_{\alpha\in A}\sum_{(b,\omega)\in I}y_\alpha(b,\omega)e(b,\omega)\leq\frac5Y\sum_{\alpha\in A}d(\delta;\alpha)$.
Summing (b) over $\alpha\in A$, the left side of (2) is at most $|W|\varepsilon\sum_{\alpha\in A}d(\delta;\alpha)$,
which is its right side.
\end{proof}
% Lean: DigitClassRemainingUniform is rest1DAE_of (restAllE_of_open (openE_of_open4 ...)), the end of the
%   chain of theorem_A_closed_form_holds; its thresholds are those of restWinE_proved (1, 1 and 2), kept by
%   thr_mono, thr_mono_ge, thresholds_mono_normalized and thr_mono_seg, and C is chosen in restAllE_of_open
%   with CM from restTM_closed (CM = 1). The Lean states it for every Inst with AdmNoK; of the conditions
%   only adm11 is used, the others being passed down to RestWinE, which ignores them. The argument is
%   segment_weighted_counting_bound_of_all_remaining_bounds with RestAll: the constants and property (1) are
%   e1, hBlk, GoodSub_mono and RestDepth; property (2) is digit_class_dimension_comparison_of_remaining_estimates,
%   with hFree, clSum_levels, claim2 (RestDir and rank_le, with realise and adm11) for (a), RestPlus for (b),
%   RestTM and RestRatio for (c), and the inequality assemble.

\begin{lemma}\label{pd-kappa-defects-segment-mean-theorem-a}
For every $C_c\geq0$ there is $T_0>0$ such that for every $C_0\geq T_0$ there is $T_1>0$ such that for
all $C_{E1},C_{E2}\geq T_1$ there is $T_2>0$ such that for every $C_H\geq T_2$ and all reals
$C_\kappa,C_s$ there is $C_A>0$ with the following property. Let an instance satisfy conditions (1) to
(11) of Definition~\ref{pd-kappa-defects-segment-mean-conditions} with these constants, and its
hypotheses; conditions (1) to (10) are not used in the proof. Form $D_A$ with $C_A$
(Definition~\ref{pd-kappa-defects-numbers}). Then the mean of
$\kappa_{S_j(b_T)}(\mu,D_A;\alpha,\alpha_+)$, formed over $\mathcal R$ with $r=m$, over
$(\alpha,\alpha_+)\in W^{S_j}\times W$ is at most $C_A(S_j-\eta_j+1)N_1/Y$.
\end{lemma}
\leanref{Formalism/SharpRestriction/GrowthBounds/FlipOneThree.lean}{Formalism.MainNotions.flipOneThree.IntegralityFromMetricForLargeConstants}
\leanref{Formalism/SharpRestriction/GrowthBounds/FlipOneThree.lean}{Formalism.MainNotions.flipOneThree.IntegralityFromMetric}
\leanref{Formalism/SharpRestriction/GrowthBounds/FockGramWindow.lean}{Formalism.MainNotions.G.GramWindowAlt.theorem_A_closed_form_holds}

\begin{proof}
\emph{The constants.} Take $T_0=T_1=1$ and $T_2=2$. Let $C_c\geq0$, $C_0\geq1$, $C_{E1},C_{E2}\geq1$ and
$C_H\geq2$, let $C_\kappa$ and $C_s$ be real numbers, and let $C_A=C\geq1$ be the constant of
Lemma~\ref{pd-kappa-defects-segment-mean-theorem-a-segment}. Let an instance satisfy the conditions and its
hypotheses; of the conditions, only (11) is used below. Write $s=S_j$, $t=S_j-\eta_j$, $S=S_j(b_T)$ and $D=D_A$,
use the notation of Definition~\ref{pd-kappa-defects-segment-mean-theorem-a-classes}, and for
$(\alpha,\alpha_+)\in W^s\times W$ write $L(\alpha,\alpha_+)=L_S(\alpha,\alpha_+;\mu)$ and
$\kappa(\alpha,\alpha_+)=\kappa_S(\mu,D;\alpha,\alpha_+)$, formed over $\mathcal R$ with $r=m$. As $x_a\geq0$,
$a\in W$.

\emph{Small range.} If $Y\leq C_A(t+1)N_1$, then, as $\kappa(\alpha,\alpha_+)\leq1$
(Lemma~\ref{pd-kappa-defects-segment-mean-theorem-a-blocks}(3)) and $Y\geq1$, the mean is at most
$1\leq C_A(t+1)N_1/Y$.

\emph{The segment bound.} Let $C_A(t+1)N_1<Y$. Lemma~\ref{pd-kappa-defects-segment-mean-theorem-a-segment}
gives integers $g(\alpha,\alpha_+,b)\geq0$ and $\varepsilon\geq0$ with $2\varepsilon e^3\leq C_A(t+1)N_1/Y$ and
its properties (1) and (2). As $s\leq N_1$ (hypothesis (1) of Theorem~\ref{pd-gain-theorem-one}) and
$C_A\geq1$, $s\leq C_A(t+1)N_1<Y$; so Lemma~\ref{pd-kappa-defects-segment-mean-theorem-a-flat} gives
$M(\alpha)\leq e^3M(\alpha')$ for all $\alpha,\alpha'\in W^s$.

\emph{One pair.} Let $(\alpha,\alpha_+)\in W^s\times W$, and let $G$ and $R$ be the sums of
$g(\alpha,\alpha_+,b)$ and of $\operatorname{rk}L_b$ over the blocks $b$ kept for $\alpha$. We show
$\kappa(\alpha,\alpha_+)\leq G/R$, reading $G/R$ as $0$ if $R=0$. For each kept $b$ let $L_2(b)$ be a good
sublattice of $L_b$ of depth $D$ as in property (1), and let $L_2$ be the family module of
Lemma~\ref{pd-kappa-defects-segment-mean-theorem-a-blocks}(2), a good sublattice of $L(\alpha,\alpha_+)$ of depth
$D$. If some kept $L_b$ has arbitrarily many linearly independent elements, then by that lemma, applied to
the $L_b$ themselves, $\operatorname{rk}L(\alpha,\alpha_+)=0$, and $\kappa(\alpha,\alpha_+)=0\leq G/R$. Otherwise
the $L_2(b)\subset L_b$ have finite rank, and by that lemma $\operatorname{rk}L(\alpha,\alpha_+)=R$ and
$\operatorname{rk}L_2=\sum_b\operatorname{rk}L_2(b)\geq R-G$. If $R=0$, then $\kappa(\alpha,\alpha_+)=0$. If $R>0$,
then $\operatorname{rk}L_2\geq(1-G/R)R$, so $L(\alpha,\alpha_+)$ has a good share of depth $D$ and defect
$G/R\geq0$, and $\kappa(\alpha,\alpha_+)\leq G/R$ by part (3) of that lemma.

\emph{Digit classes.} By Lemma~\ref{pd-kappa-defects-segment-mean-theorem-a-blocks}(1),
$\operatorname{rk}L_b=\dim V_b$. Every kept block has a class compatible with $\alpha$; grouping the kept
blocks by their class, $G$ and $R$ are the sums, over the digit classes $\delta$ compatible with $\alpha$, of
the numbers $n(\delta;\alpha,\alpha_+)$ and $r(\delta;\alpha,\alpha_+)$, the sums of $g(\alpha,\alpha_+,b)$ and of
$\dim V_b$ over the blocks $b$ kept for $\alpha$ of class $\delta$. The function $a_T^{\alpha,\alpha_+}$ is a
window function, as its values $\alpha_i$, $\alpha_+$ and $a$ lie in $W$. For it, the digit classes
compatible with $\alpha$ are the digits for $j$, $d(\delta;\alpha)=d_j(\rho,\rho')$ and
$r(\delta;\alpha,\alpha_+)=r_j(\mu;\rho,\rho')$, where $\delta=(\rho,\rho')$. As $0\leq\mu\leq b'/(kQY)$, condition
(11) gives $d(\delta;\alpha)/2\leq r(\delta;\alpha,\alpha_+)$.

\emph{The mean.} Apply Lemma~\ref{pd-kappa-defects-segment-mean-theorem-a-average} to $A=W^s\times W$, the
set $\Delta$ of digit classes, the relation ``$\delta$ is compatible with $\alpha$'', the functions $n$, $r$ and
$d(\delta;\alpha)$ (extended by $0$), $\varepsilon$ and $K=e^3$. Its hypothesis (1) is the last paragraph, (2) is
property (2) of Lemma~\ref{pd-kappa-defects-segment-mean-theorem-a-segment}, and (3) is the bound
$M(\alpha)\leq e^3M(\alpha')$, as the sum of the $d(\delta;\alpha)$ over the classes compatible with $\alpha$ is
$M(\alpha)$. So the sum of the $G/R$ over $W^s\times W$ is at most $2\varepsilon e^3|W|^{s+1}$, and the mean of
$\kappa(\alpha,\alpha_+)$ over $W^s\times W$ is at most $2\varepsilon e^3\leq C_A(t+1)N_1/Y$.
\end{proof}
% Lean: Loc8All.p13A = GramWindowAlt.theorem_A_closed_form_holds = closed_of_minimal_vector D
%   PI.exists_minimalNormVector diagonal_rank_bound (restWinE_proved D). closed_of_minimal_vector is a chain
%   of reductions (closed_of_splitting_matrix, closed_of_orthogonal_splitting,
%   integralityFromMetric_of_orthBasisLat, closed_of_orthogonal_basis, closed_of_orthogonal_lattice,
%   closed_of_symmetric_dims, closed_of_symmetric_windows, closed_of_coefficient_function,
%   closed_of_symmetric_finite, integrality_from_metric_of_orthogonality_window,
%   theorem_A_closed_of_nine_open_bounds, ..., closed_of_open_directions, closed_of_open_conditions,
%   theorem_A_closed_of_all_remaining_bounds), each supplying one input of Lemma
%   pd-kappa-defects-segment-mean-theorem-a-segment (DigitClassRemainingUniform), and ends in
%   counting_statement_uniform_of_digit_class, which is the proof above: thr_mono_ge (T2 = max T2 1),
%   counting_statement_of_digit_class_rank (restRk_proved), counting_statement_of_digit_class_rank_form,
%   remaining_of_segment_bound (s <= Y and restFlat_exp), counting_theorem_of_remaining_sum =
%   counting_statement_of_block_form (counting_statement_small_range with meanKapT_le_one; meanBox_mono
%   with badShare_le_of_goodShare and segSub_good) after block_form_of_segment_form (shareB, finrank_segSub,
%   finrank_segSub_zero) and segment_form_of_scaled_form (shareB_eq, mean_of_classes: keptSum_eq_classes,
%   rank_ge with realise and adm11, double_count, box_card). The conditions other than adm11 enter only
%   through AdmNoK, which the chain passes down to RestWinE unused.

\subsection{Isotypic projectors and restriction}\label{subsection-pd-isotypic-projectors-restriction}

\begin{definition}\label{pd-kappa-defects-segment-mean-letters-families}
Let $K$ be a field of characteristic $0$ and $n\geq1$ an integer, use the notions of
Definition~\ref{pd-kappa-defects-segment-mean-operators}(1), and let $\epsilon_a$ be the $a$-th unit vector
of $\mathbf Z^n$.
\begin{enumerate}
\item A \emph{$\mathfrak{gl}_n$-family} on a $K$-vector space $V$ is a family $E=(E_{ab})_{1\leq a,b\leq n}$
  of endomorphisms of $V$ with $E_{ab}E_{cd}-E_{cd}E_{ab}=\delta_{bc}E_{ad}-\delta_{da}E_{cb}$ for all
  $a,b,c,d$, such that every $x\in V$ lies in a finite-dimensional subspace stable under the $E_{ab}$ and
  $V$ is the sum of its weight spaces $V_\mu$. A \emph{subrepresentation} is a subspace $W$ stable under
  every $E_{ab}$, and $E|_W$ is the family of the restrictions of the $E_{ab}$ to $W$.
\item Let $E$ and $E'$ be $\mathfrak{gl}_n$-families on $V$ and $V'$. A $K$-linear map $T\colon V\to V'$
  is \emph{equivariant} if $TE_{ab}=E'_{ab}T$ for all $a,b$. The family $\operatorname{Hom}(E,E')$ on
  $\operatorname{Hom}_K(V,V')$ consists of the maps $\phi\mapsto E'_{ab}\phi-\phi E_{ab}$. An
  \emph{invariant projector} of $E$ is an isotypic projector for $\{0\}$.
\item Let $V$ be finite-dimensional. Its \emph{root length} $\operatorname{rl}(V)$ is the largest integer
  $k\geq0$ such that $\mu-\nu=k(\epsilon_a-\epsilon_b)$ for some $a\neq b$ and some $\mu,\nu$ with
  $V_\mu\neq0$ and $V_\nu\neq0$, and $0$ if there is none. The \emph{spread} of $\lambda\in\mathbf Z^n$
  with $\lambda_1\geq\lambda_n$ is $\operatorname{spr}(\lambda)=\lambda_1-\lambda_n$. Put $N_0(n)=n(n-1)/2$,
  and for a prime $p$ and an integer $x\geq0$ let $\lfloor\log_px\rfloor$ be the largest $d\geq0$ with
  $p^d\leq x$ if $x\geq1$, and $0$ if $x=0$.
\end{enumerate}
\end{definition}
\leanref{Formalism/SharpRestriction/GrowthBounds/ObjPol.lean}{Formalism.MainNotions.G.IsGLRep}
\leanref{Formalism/SharpRestriction/GrowthBounds/RootWeights/OperatorsKillInvariantProjection.lean}{Formalism.MainNotions.G.GLWeightTheory.IsSubrepresentation}
\leanref{Formalism/SharpRestriction/GrowthBounds/RootWeights/OperatorsKillInvariantProjection.lean}{Formalism.MainNotions.G.GLWeightTheory.restrictRepresentation}
\leanref{Formalism/SharpRestriction/GrowthBounds/RootWeights/OperatorsKillInvariantProjection.lean}{Formalism.MainNotions.G.GLWeightTheory.IsEquivariantMap}
\leanref{Formalism/SharpRestriction/GrowthBounds/RootWeights/OperatorsKillInvariantProjection.lean}{Formalism.MainNotions.G.GLWeightTheory.homRepresentation}
\leanref{Formalism/SharpRestriction/GrowthBounds/ObjPol.lean}{Formalism.MainNotions.G.IsInvProj}
\leanref{Formalism/SharpRestriction/GrowthBounds/ObjPol.lean}{Formalism.MainNotions.G.ellOf}
\leanref{Formalism/SharpRestriction/GrowthBounds/ObjPol.lean}{Formalism.MainNotions.G.ellWt}
\leanref{Formalism/SharpRestriction/GrowthBounds/ObjPol.lean}{Formalism.MainNotions.G.N0}
\leanref{Formalism/SharpRestriction/GrowthBounds/ObjPol.lean}{Formalism.MainNotions.G.deF}

\begin{lemma}\label{pd-kappa-defects-segment-mean-letters-span}
Let $K$ be a field of characteristic $0$, $n\geq1$, and $E$ a $\mathfrak{gl}_n$-family on a $K$-vector
space $V$ (Definition~\ref{pd-kappa-defects-segment-mean-letters-families}), with isotypic subspaces
$V^{(\lambda)}$, $\lambda\in\mathbf Z^n$ (Definition~\ref{pd-kappa-defects-segment-mean-operators}(1)).
\begin{enumerate}
\item For a subrepresentation $W$, $E|_W$ is a $\mathfrak{gl}_n$-family on $W$. Its weight spaces are the
  $W\cap V_\mu$, and its highest weight vectors are the highest weight vectors of $E$ that lie in $W$.
\item If $V$ is finite-dimensional and $V\neq0$, then $V$ contains a highest weight vector.
\item Each $V^{(\lambda)}$ is stable under every $E_{ab}$ and under every endomorphism of $V$ that
  commutes with all $E_{ab}$.
\item $V$ is the sum of the $V^{(\lambda)}$.
\item Let $T\subset\mathbf Z^n$ and let $e$ be an isotypic projector for $T$. Then $e$ commutes with every
  endomorphism of $V$ that maps each $V^{(\lambda)}$ into itself. If $T=\{\lambda\}$ and $W$ is a
  subrepresentation, then $e(W)\subset W$, the restriction of $e$ to $W$ is an isotypic projector of
  $E|_W$ for $\{\lambda\}$, and $e$ vanishes on $W$ if $W$ contains no highest weight vector of weight
  $\lambda$.
\end{enumerate}
\end{lemma}
\leanref{Formalism/SharpRestriction/GrowthBounds/RootWeights/ExistsHighestWeightVectorFindim.lean}{Formalism.MainNotions.G.GLWeightTheory.exists_highest_weight_vector_findim}
\leanref{Formalism/SharpRestriction/GrowthBounds/RootWeights/ExistsHighestWeightVectorFindim.lean}{Formalism.MainNotions.G.GLWeightTheory.complete_reducibility_findim}
\leanref{Formalism/SharpRestriction/GrowthBounds/RootWeights/ExistsHighestWeightVectorFindim.lean}{Formalism.MainNotions.G.GLWeightTheory.complete_reducibility_of_semisimple}
\leanref{Formalism/SharpRestriction/GrowthBounds/RootWeights/ExistsHighestWeightVectorFindim.lean}{Formalism.MainNotions.G.GLWeightTheory.subrep_restriction_of_reducibility}
\leanref{Formalism/SharpRestriction/GrowthBounds/RootWeights/IsotypicProjectorSelfAdjoint.lean}{Formalism.MainNotions.G.GLWeightTheory.isotypic_stable_gl_operator}
\leanref{Formalism/SharpRestriction/GrowthBounds/RootWeights/IsotypicProjectorSelfAdjoint.lean}{Formalism.MainNotions.G.GLWeightTheory.isotypic_stable_commuting}
\leanref{Formalism/SharpRestriction/GrowthBounds/RootWeights/IsotypicProjectorSelfAdjoint.lean}{Formalism.MainNotions.G.GLWeightTheory.isotypic_projector_commutes}
\leanref{Formalism/SharpRestriction/GrowthBounds/CountingStatementReduction.lean}{Formalism.MainNotions.G.CountingStatementReduction.isoProjSet_comm_abs}

\begin{proof}
Let $\mathcal U$ be the monoid generated by the $E_{ab}$. The relations give
$E_{cc}E_{ab}-E_{ab}E_{cc}=(\delta_{ca}-\delta_{cb})E_{ab}$, so $E_{ab}$ maps $V_\mu$ into
$V_{\mu+\epsilon_a-\epsilon_b}$, and $V$ with the $V_\mu$ is a weight module for $\mathfrak{gl}_n$
(Definition~\ref{pd-kappa-defects-gl-modules}(1)).

(1) The relations hold on $W$, and if $x\in W$ lies in a finite-dimensional subspace $U$ stable under the
$E_{ab}$, it lies in $U\cap W$. By Lemma~\ref{pd-kappa-defects-gl-highest}(1), the weight components of an
element of $W$ lie in $W$, so $W$ is the sum of the $W\cap V_\mu$, which are the weight spaces of $E|_W$.
The last assertion holds as the operators of $E|_W$ are restrictions.

(2) Put $h(\mu)=-\sum_cc\mu_c$. By Lemma~\ref{pd-kappa-defects-gl-highest}(1), $V$ has finitely many
weights, and it has one, as $V\neq0$. Let $\mu$ be a weight with $h(\mu)$ largest and $0\neq v\in V_\mu$.
For $a<b$, $E_{ab}v$ lies in $V_{\mu+\epsilon_a-\epsilon_b}$, and
$h(\mu+\epsilon_a-\epsilon_b)=h(\mu)+b-a>h(\mu)$, so $E_{ab}v=0$. Hence $v$ is a highest weight vector.

(3) $V^{(\lambda)}$ is spanned by the $u(v)$ with $u\in\mathcal U$ and $v$ a highest weight vector of weight
$\lambda$, and $E_{ab}u\in\mathcal U$. If $X$ commutes with all $E_{ab}$, then $Xu(v)=u(Xv)$, and $Xv$ is
$0$ or a highest weight vector of weight $\lambda$.

(4) Let $U$ be the sum of the $V^{(\lambda)}$, a subrepresentation by (3). Suppose first that $V$ is
finite-dimensional. By Lemma~\ref{pd-kappa-defects-paths-reducible}, $U$ has a complement $U'$ that is a
subrepresentation. If $U'\neq0$, then by (1) and (2) it contains a highest weight vector of $E$, of some
weight $\lambda$, which lies in $V^{(\lambda)}\subset U$; this contradicts $U\cap U'=0$. So $U=V$. In
general, let $x\in V$, and let $W$ be a finite-dimensional subspace stable under the $E_{ab}$ that
contains $x$. By (1) and the finite-dimensional case, $x$ lies in the sum of the isotypic subspaces
$W^{(\lambda)}$ of $E|_W$. The monoid generated by the operators of $E|_W$ consists of the restrictions of
the elements of $\mathcal U$, so $W^{(\lambda)}\subset V^{(\lambda)}$ by (1), and $x\in U$.

(5) Let $x\in V$; by (4), $x=\sum_\mu x_\mu$ with $x_\mu\in V^{(\mu)}$. If $X$ maps each $V^{(\mu)}$ into
itself, then $eXx_\mu$ and $Xex_\mu$ are both $Xx_\mu$ if $\mu\in T$ and both $0$ otherwise, so $eX=Xe$.
Let $T=\{\lambda\}$ and let $W$ be a subrepresentation. By (1) and (4), applied to $E|_W$, $W$ is the sum of
the isotypic subspaces $W^{(\mu)}$ of $E|_W$, and $W^{(\mu)}\subset V^{(\mu)}$ as in (4). So $e$ is the
identity on $W^{(\lambda)}$ and $0$ on the $W^{(\mu)}$ with $\mu\neq\lambda$. Hence $e(W)\subset W$, and the
restriction is an isotypic projector of $E|_W$ for $\{\lambda\}$. If $W$ contains no highest weight vector
of weight $\lambda$, then $W^{(\lambda)}=0$ by (1), and $e$ vanishes on $W$.
\end{proof}
% Lean: (1) is the private isGLRep_restrictStableOperators, with weight_components_mem_of_sum_mem; (2) is
%   exists_highest_weight_vector_findim (h is raisingFunctional); (4) is complete_reducibility_findim
%   (from adm2) and complete_reducibility_of_semisimple; (5) is isoProjSet_comm_abs and
%   subrep_restriction_of_reducibility.

\begin{lemma}\label{pd-kappa-defects-segment-mean-letters-forms}
Let $K$, $n$, $E$, $V$ and the $V^{(\lambda)}$ be as in
Lemma~\ref{pd-kappa-defects-segment-mean-letters-span}, and let $\beta$ be a $K$-bilinear form on $V$ with
$\beta(E_{ab}x,y)=\beta(x,E_{ba}y)$ for all $a,b$ and $x,y\in V$.
\begin{enumerate}
\item Let $\lambda\neq\mu$. Then $\beta(x,y)=0$ if $x\in V_\lambda$ and $y\in V_\mu$, and if
  $x\in V^{(\lambda)}$ and $y\in V^{(\mu)}$.
\item Let $T\subset\mathbf Z^n$ and let $e$ be an isotypic projector for $T$. Then
  $\beta(ex,y)=\beta(x,ey)$ for all $x,y\in V$.
\item Suppose that $x=0$ whenever $\beta(x,y)=0$ for all $y\in V$. Then $V$ is the direct sum of the
  $V^{(\lambda)}$, and every $T\subset\mathbf Z^n$ has an isotypic projector.
\end{enumerate}
\end{lemma}
\leanref{Formalism/SharpRestriction/GrowthBounds/CountingStatementReduction.lean}{Formalism.MainNotions.G.CountingStatementReduction.wt_orth}
\leanref{Formalism/SharpRestriction/GrowthBounds/RootWeights/ContravariantOrthogonalityHolds.lean}{Formalism.MainNotions.G.GLWeightTheory.contravariant_orthogonality_holds}
\leanref{Formalism/SharpRestriction/GrowthBounds/RootWeights/ContravariantOrthogonalityHolds.lean}{Formalism.MainNotions.G.GLWeightTheory.isotypic_le_lowering_span}
\leanref{Formalism/SharpRestriction/GrowthBounds/RootWeights/ContravariantOrthogonalityHolds.lean}{Formalism.MainNotions.G.GLWeightTheory.hwv_orthogonal_lowering_span}
\leanref{Formalism/SharpRestriction/GrowthBounds/RootWeights/IsotypicProjectorSelfAdjoint.lean}{Formalism.MainNotions.G.GLWeightTheory.isotypic_projector_selfAdjoint}
\leanref{Formalism/SharpRestriction/GrowthBounds/CountingStatementReduction.lean}{Formalism.MainNotions.G.CountingStatementReduction.isoProjSet_selfadj_abs}
\leanref{Formalism/SharpRestriction/GrowthBounds/CountingStatementReduction.lean}{Formalism.MainNotions.G.CountingStatementReduction.iso_indep_abs}
\leanref{Formalism/SharpRestriction/GrowthBounds/CountingStatementReduction.lean}{Formalism.MainNotions.G.CountingStatementReduction.isoProjSet_exists_abs}

\begin{proof}
(1) Let $x\in V_\lambda$, $y\in V_\mu$, and choose $a$ with $\lambda_a\neq\mu_a$. Then
$\lambda_a\beta(x,y)=\beta(E_{aa}x,y)=\beta(x,E_{aa}y)=\mu_a\beta(x,y)$, so $\beta(x,y)=0$.

For the second assertion, let $H_\mu$ be the span of the vectors $u(w)$, with $w$ a highest weight vector of
weight $\mu$ and $u$ a product of operators $E_{cd}$ with $c>d$, the empty product included. Each $E_{cd}$
with $c>d$ maps $H_\mu$ into itself. By the relations,
$E_{ab}E_{cd}=E_{cd}E_{ab}+\delta_{bc}E_{ad}-\delta_{da}E_{cb}$; induction on the length of $u$, starting from
$E_{ab}w=0$ for $a<b$, $E_{aa}w=\mu_aw$ and $E_{ab}w\in H_\mu$ for $a>b$, shows that $E_{ab}u(w)\in H_\mu$ for
all $a,b$. So $H_\mu$ is a subrepresentation, and it contains $V^{(\mu)}$. Let $v$ be a highest weight vector
of weight $\lambda$. If $u$ is empty, $\beta(v,u(w))=0$ by the first assertion; otherwise $u=E_{cd}u'$ with
$c>d$, and $\beta(v,u(w))=\beta(E_{dc}v,u'(w))=0$. So $\beta(v,y)=0$ for $y\in V^{(\mu)}$. For
$u=E_{a_1b_1}\cdots E_{a_kb_k}$, $\beta(u(v),y)=\beta(v,E_{b_ka_k}\cdots E_{b_1a_1}y)$, and $V^{(\mu)}$ is
stable under the $E_{ab}$ (Lemma~\ref{pd-kappa-defects-segment-mean-letters-span}(3)). So
$\beta(u(v),y)=0$, and the assertion follows by linearity.

(2) By Lemma~\ref{pd-kappa-defects-segment-mean-letters-span}(4) it suffices to take $x\in V^{(\lambda)}$ and
$y\in V^{(\mu)}$. Then $ex$ is $x$ or $0$, and $ey$ is $y$ or $0$. If $\lambda=\mu$, both sides are
$\beta(x,y)$ if $\lambda\in T$ and $0$ otherwise; if $\lambda\neq\mu$, both vanish by (1).

(3) Let $x\in V^{(\lambda)}$ lie in the sum of the $V^{(\mu)}$ with $\mu\neq\lambda$, and let
$y=\sum_\mu y_\mu$ with $y_\mu\in V^{(\mu)}$ (Lemma~\ref{pd-kappa-defects-segment-mean-letters-span}(4)). By
(1), $\beta(x,y_\mu)=0$ for $\mu\neq\lambda$, as $x\in V^{(\lambda)}$, and $\beta(x,y_\lambda)=0$, as $x$ is a
sum of elements of the $V^{(\mu)}$ with $\mu\neq\lambda$. So $\beta(x,y)=0$ for all $y$, and $x=0$. Hence the
sum of the $V^{(\lambda)}$ is direct, and the projection of $V$ onto $\bigoplus_{\lambda\in T}V^{(\lambda)}$
along $\bigoplus_{\lambda\notin T}V^{(\lambda)}$ is an isotypic projector for $T$.
\end{proof}

\begin{lemma}\label{pd-kappa-defects-segment-mean-letters-trace}
Let $K$ be a field of characteristic $0$, $n\geq1$, and $E$ a $\mathfrak{gl}_n$-family on a
finite-dimensional $K$-vector space $V$, with isotypic subspaces $V^{(\lambda)}$
(Definition~\ref{pd-kappa-defects-segment-mean-letters-families}); let $\mathcal U$ be the monoid generated
by the $E_{ab}$.
\begin{enumerate}
\item $E$ has an invariant projector, and every invariant projector of $E$ maps $V$ into the space of the
  $x$ with $E_{ab}x=0$ for all $a,b$.
\item Let $E'$ be a $\mathfrak{gl}_n$-family on a finite-dimensional $K$-vector space $V'$. Then
  $\operatorname{Hom}(E,E')$ is a $\mathfrak{gl}_n$-family. If $T\colon V\to V'$ is equivariant and $P$,
  $P'$ are invariant projectors of $E$ and $E'$, then $TP=P'T$.
\item Let $\lambda\in\mathbf Z^n$, $e$ an isotypic projector of $E$ for $\{\lambda\}$, $v_0$ a highest
  weight vector of weight $\lambda$, and $L$ the span of the $u(v_0)$ with $u\in\mathcal U$, a
  subrepresentation. Let $P$ be an invariant projector of $\operatorname{Hom}(E|_L,E)$, $b_1,\dots,b_d$ a
  basis of $L$ and $b_1^*,\dots,b_d^*$ its dual basis, and write $\xi\otimes x$ for the map
  $y\mapsto\xi(y)x$. Then
  \[
   e(x)=d\sum_{i=1}^dP(b_i^*\otimes x)(b_i)\qquad(x\in V).
  \]
\end{enumerate}
\end{lemma}
\leanref{Formalism/SharpRestriction/GrowthBounds/RootWeights/InvariantProjectionRangeTrivial.lean}{Formalism.MainNotions.G.GLWeightTheory.invariant_projection_exists}
\leanref{Formalism/SharpRestriction/GrowthBounds/RootWeights/InvariantProjectionRangeTrivial.lean}{Formalism.MainNotions.G.GLWeightTheory.invariant_projection_kills_operators}
\leanref{Formalism/SharpRestriction/GrowthBounds/RootWeights/InvariantProjectionRangeTrivial.lean}{Formalism.MainNotions.G.GLWeightTheory.isotypic_zero_killed_by_operators}
\leanref{Formalism/SharpRestriction/GrowthBounds/RootWeights/WeightProjectionDiagonalApplyEqSmul.lean}{Formalism.MainNotions.G.GLWeightTheory.homRepresentation_isRep}
\leanref{Formalism/SharpRestriction/GrowthBounds/RootWeights/InvariantProjectionRangeTrivial.lean}{Formalism.MainNotions.G.GLWeightTheory.equivariant_commutes_isotypic}
\leanref{Formalism/SharpRestriction/GrowthBounds/RootWeights/IsotypicProjectorTraceFormula.lean}{Formalism.MainNotions.G.GLWeightTheory.isotypic_projector_trace_formula}
\leanref{Formalism/SharpRestriction/GrowthBounds/RootWeights/SchurInvariantProjectionTrace.lean}{Formalism.MainNotions.G.GLWeightTheory.schur_invariant_projection_trace}
\leanref{Formalism/SharpRestriction/GrowthBounds/RootWeights/SchurInvariantProjectionTrace.lean}{Formalism.MainNotions.G.GLWeightTheory.exists_equivariant_cyclic_to_hwv}

\begin{proof}
(1) Let $I$ be the space of the $x$ with $E_{ab}x=0$ for all $a,b$. It is a subrepresentation, so by
Lemma~\ref{pd-kappa-defects-paths-reducible} it has a complement $I'$ that is a subrepresentation; let
$P_0$ be the projection onto $I$ along $I'$. For $y\in V$, $E_{cd}y=E_{cd}(y-P_0y)\in I'$, so
$P_0E_{cd}=0$, and $P_0u=0$ for $u\neq1$ in $\mathcal U$. A highest weight vector $w$ of weight $0$ lies
in $I$: $E_{ab}w=0$ for $a\leq b$, and for $a>b$ because $w$ is a vector of weight $0$ killed by $E_{ba}$
for the copy of $\mathfrak{sl}_2$ spanned by $E_{ba}$, $E_{ab}$ and $E_{bb}-E_{aa}$ \cite[\S7.2]{Hum72}.
So $u(w)$ is $w$ or $0$ for $u\in\mathcal U$, and $V^{(0)}\subset I$, on which $P_0$ is the identity. A
highest weight vector $v$ of weight $\mu\neq0$ satisfies $v=\mu_c^{-1}E_{cc}v$ for some $c$ with
$\mu_c\neq0$, so $P_0v=0$, and $P_0$ vanishes on $V^{(\mu)}$. So $P_0$ is an invariant projector. Every
invariant projector is the identity on $V^{(0)}$ and $0$ on the other $V^{(\mu)}$, so by
Lemma~\ref{pd-kappa-defects-segment-mean-letters-span}(4) it maps $V$ into $V^{(0)}\subset I$.

(2) The maps $\phi\mapsto E'_{ab}\phi$ and $\phi\mapsto\phi E_{cd}$ commute, so the relations for $E$ and
$E'$ give those for $\operatorname{Hom}(E,E')$, as in Lemma~\ref{pd-kappa-defects-projector-invariant}(1).
Let $\pi_\mu$ and $\pi'_\mu$ be the projections of $V$ and $V'$ onto their weight spaces of weight $\mu$
along the others. Every $\phi$ is the finite sum of the $\pi'_{\mu'}\phi\pi_\mu$, and
$\pi'_{\mu'}\phi\pi_\mu$ is a weight vector of weight $\mu'-\mu$. As $\operatorname{Hom}_K(V,V')$ is
finite-dimensional, $\operatorname{Hom}(E,E')$ is a $\mathfrak{gl}_n$-family. Next, for $u\in\mathcal U$
and a highest weight vector $v$, $Tu(v)=u'(Tv)$, where $u'$ is the same product of the $E'_{ab}$, and $Tv$
is $0$ or a highest weight vector of the weight of $v$; so $T(V^{(\mu)})\subset V'^{(\mu)}$. For $x\in
V^{(0)}$, $TPx=Tx=P'Tx$; for $x\in V^{(\mu)}$ with $\mu\neq0$, $TPx=0=P'Tx$. By
Lemma~\ref{pd-kappa-defects-segment-mean-letters-span}(4), $TP=P'T$.

(3) $L$ is stable under the $E_{ab}$, and $E|_L$ is a $\mathfrak{gl}_n$-family
(Lemma~\ref{pd-kappa-defects-segment-mean-letters-span}(1)); so $F=\operatorname{Hom}(E|_L,E)$ and
$F_L=\operatorname{Hom}(E|_L,E|_L)$ are $\mathfrak{gl}_n$-families by (2). By (1), every map in the image of
an invariant projector of $F$ or $F_L$ is killed by every operator of the family, that is, it is
equivariant.

\emph{Step 1.} Let $P_L$ be an invariant projector of $F_L$ and $g\in\operatorname{End}_K(L)$. Then
$P_L(g)$ is equivariant, so it maps the weight space $L_\lambda$ into itself, and $L_\lambda=Kv_0$ by
Lemma~\ref{pd-kappa-defects-gl-highest}(2). So $P_L(g)v_0=cv_0$ for some $c\in K$; the equivariant map
$P_L(g)-c$ kills $v_0$, hence every $u(v_0)$, and $P_L(g)=c$. The trace is an equivariant map from $F_L$
to the family $0$ on $K$, as $\operatorname{tr}(E_{ab}g-gE_{ab})=0$, and the identity is an invariant
projector of that family. By (2), $\operatorname{tr}(P_L(g))=\operatorname{tr}(g)$, that is,
$dP_L(g)=\operatorname{tr}(g)$.

\emph{Step 2.} Let $v$ be a highest weight vector of weight $\lambda$ and $u\in\mathcal U$; then
$e(u(v))=u(v)$. The spans of the $u'(v_0)$ and of the $u'(v)$, $u'\in\mathcal U$, are finite-dimensional
representations generated by highest weight vectors of the same weight, so they are irreducible and
isomorphic \cite[\S\S6.3, 20.2, 20.3]{Hum72}; an isomorphism maps $L_\lambda=Kv_0$ onto the line $Kv$. So
there is an equivariant $\varphi\colon L\to V$ with $\varphi(v_0)=v$. Put $l=u(v_0)$, so that
$\varphi(l)=u(v)$. Composition with $\varphi$ is an equivariant map from $F_L$ to $F$, so by (2) and
Step 1,
\[
 dP(b_i^*\otimes u(v))=dP(\varphi\circ(b_i^*\otimes l))=\varphi\circ dP_L(b_i^*\otimes l)=b_i^*(l)\,\varphi .
\]
So $d\sum_iP(b_i^*\otimes u(v))(b_i)=\varphi\big(\sum_ib_i^*(l)b_i\big)=\varphi(l)=u(v)$.

\emph{Step 3.} Let $\mu\neq\lambda$, $y\in V^{(\mu)}$ and $\xi$ a linear form on $L$. The space
$V^{(\mu)}$ is a subrepresentation, and its inclusion into $V$ is equivariant; by (2), applied to
composition with it, $P(\xi\otimes y)$ takes values in $V^{(\mu)}$. It is equivariant. If
$P(\xi\otimes y)(v_0)\neq0$, it is a highest weight vector of weight $\lambda$ in $V^{(\mu)}$, so it lies
in $V^{(\lambda)}\cap V^{(\mu)}$, on which $e$ is both the identity and $0$; this is impossible. So
$P(\xi\otimes y)$ kills $v_0$, hence $L$, and both sides of the formula vanish at $x=y$, as $e(y)=0$.

By Lemma~\ref{pd-kappa-defects-segment-mean-letters-span}(4), the vectors of Steps 2 and 3 span $V$, and
both sides of the formula are linear in $x$.
\end{proof}
% Lean: (1) is invariant_projection_exists (from adm2) and invariant_projection_kills_operators (through
%   invariant_projection_range_trivial and isotypic_zero_killed_by_operators; its sl_2 step is the private
%   E_apply_eq_zero_of_isHWV_zero). (2) is homRepresentation_isRep and equivariant_commutes_isotypic (with
%   equivariant_maps_isotypic). (3) is isotypic_projector_trace_formula, with invariant_projection_equivariant,
%   schur_invariant_projection_trace (Step 1; cyclic_highest_weightSpace_line, trivialRep_invariant_projection),
%   exists_equivariant_cyclic_to_hwv and equivariant_lifts_monoid (Step 2), and equivariant_image_highest_vector,
%   highest_vector_mem_isotypic and equivariant_kills_cyclic_of_zero (Step 3).

\begin{lemma}\label{pd-kappa-defects-segment-mean-letters-separators}
Let $\mathcal O$ be a discrete valuation ring whose fraction field $K$ has characteristic $0$, $n\geq1$,
$(E_{ab})_{a,b\leq n}$ endomorphisms of a $K$-vector space $V$ with weight spaces $V_\mu$
(Definition~\ref{pd-kappa-defects-segment-mean-operators}(1)), and $M\subset V$ an $\mathcal O$-submodule
stable under the $\binom{E_{aa}}k=E_{aa}(E_{aa}-1)\cdots(E_{aa}-k+1)/k!$, $a\leq n$, $k\geq0$.
\begin{enumerate}
\item For $\lambda\neq\nu$ in $\mathbf Z^n$ there are an endomorphism $T$ of $V$ with $T(M)\subset M$ and a
  function $g\colon\mathbf Z^n\to\mathbf Z$ with $g(\lambda)=1$, $g(\nu)=0$ and $Tx=g(\mu)x$ for every
  $\mu$ and $x\in V_\mu$.
\item If $x\in M$ is a sum of elements $x_\mu\in V_\mu$, with $\mu$ in a finite set $F$, then $x_\mu\in M$
  for every $\mu\in F$.
\end{enumerate}
\end{lemma}
\leanref{Formalism/SharpRestriction/GrowthBounds/RootWeights/IntegralEnvelopingPreservesSubmodule.lean}{Formalism.MainNotions.G.GLWeightTheory.lattice_stable_weight_separator}
\leanref{Formalism/SharpRestriction/GrowthBounds/RootWeights/IntegralEnvelopingPreservesSubmodule.lean}{Formalism.MainNotions.G.GLWeightTheory.lattice_weight_components_mem}
\leanref{Formalism/SharpRestriction/GrowthBounds/RootWeights/WeightProjectionPreservesLattice.lean}{Formalism.MainNotions.G.GLWeightTheory.weightProjection_preserves_lattice}

\begin{proof}
(1) On $V_\mu$, $\binom{E_{aa}}k$ is multiplication by $\binom{\mu_a}k$, where
$\binom xk=x(x-1)\cdots(x-k+1)/k!\in\mathbf Z$ for $x\in\mathbf Z$. For $c\in\mathbf Z$ and $\delta\geq0$, the
operator $T_{a,c,\delta}=\sum_{i+j=\delta}\binom cj\binom{E_{aa}}i$ maps $M$ into $M$, and by Vandermonde's
identity it acts on $V_\mu$ as $\sum_{i+j=\delta}\binom cj\binom{\mu_a}i=\binom{\mu_a+c}\delta$. Choose $a$
with $\lambda_a\neq\nu_a$. If $\nu_a<\lambda_a$, take $T=T_{a,-\nu_a,\lambda_a-\nu_a}$ and
$g(\mu)=\binom{\mu_a-\nu_a}{\lambda_a-\nu_a}$. If $\nu_a>\lambda_a$, take $T=1-T_{a,-\lambda_a,\nu_a-\lambda_a}$
and $g(\mu)=1-\binom{\mu_a-\lambda_a}{\nu_a-\lambda_a}$. In both cases $g(\lambda)=1$ and $g(\nu)=0$.

(2) Let $\mu\in F$, and for $\nu\in F\smallsetminus\{\mu\}$ let $T_\nu$ and $g_\nu$ be as in (1) for the pair
$(\mu,\nu)$. The product of the $T_\nu$ maps $M$ into $M$ and $x$ to
$\sum_{\mu'\in F}\prod_\nu g_\nu(\mu')x_{\mu'}=x_\mu$.
\end{proof}

\begin{lemma}\label{pd-kappa-defects-segment-mean-letters-integral}
Let $p$ be a prime, $\mathcal O$ a discrete valuation ring whose fraction field $K$ has characteristic
$0$, with $p$ in the maximal ideal of $\mathcal O$, $n\geq1$, and $E$ a $\mathfrak{gl}_n$-family on a
finite-dimensional $K$-vector space $V$ (Definition~\ref{pd-kappa-defects-segment-mean-letters-families}).
Let $M\subset V$ be a finitely generated $\mathcal O$-submodule that spans $V$ over $K$ and is stable under
the $E_{ab}^k/k!$ with $a\neq b$ and the $\binom{E_{aa}}k$, $k\geq0$. Let $\Sigma\subset\mathbf Z^n$ be a
finite set of dominant weights and, for $\lambda\in\Sigma$, $e_\lambda$ an isotypic projector of $E$ for
$\{\lambda\}$. Let $c$ be the largest of the numbers
$c_\lambda=N_0(n)\lfloor\log_p(\operatorname{rl}(V)+\operatorname{spr}(\lambda)+n-1)\rfloor$, $\lambda\in\Sigma$,
and $c=0$ if $\Sigma$ is empty. Then $p^c\sum_{\lambda\in\Sigma}e_\lambda(x)\in M$ for every $x\in M$.
\end{lemma}
\leanref{Formalism/SharpRestriction/GrowthBounds/CovariantDegrees.lean}{Formalism.MainNotions.CovariantDegrees.IsotypicProjectorTheoremField}
\leanref{Formalism/SharpRestriction/GrowthBounds/RootWeights/LatticeTheoremSingleRepLinf.lean}{Formalism.MainNotions.G.GLWeightTheory.isotypicProjectorTheoremField_of_linfty_bound}
\leanref{Formalism/SharpRestriction/GrowthBounds/RootWeights/LatticeTheoremSingleRepLinf.lean}{Formalism.MainNotions.G.GLWeightTheory.lattice_theorem_single_rep_linf}
\leanref{Formalism/SharpRestriction/GrowthBounds/RootWeights/ExistsClearDenominators.lean}{Formalism.MainNotions.G.GLWeightTheory.lattice_adapted_basis_exists}
\leanref{Formalism/SharpRestriction/GrowthBounds/RootWeights/HomLatticeDividedPowerStable.lean}{Formalism.MainNotions.G.GLWeightTheory.hom_lattice_dividedPower_stable}
\leanref{Formalism/SharpRestriction/GrowthBounds/RootWeights/WeightMemWeightsOfNeZero.lean}{Formalism.MainNotions.G.GLWeightTheory.linf_homRep_le_ell_add}

\begin{proof}
As $c\geq c_\lambda$ for $\lambda\in\Sigma$ and $M$ is an $\mathcal O$-module, it suffices to show
$p^{c_\lambda}e_\lambda(x)\in M$ for each $\lambda\in\Sigma$ and $x\in M$. Fix $\lambda$ and write
$e=e_\lambda$. If $V$ contains no highest weight vector of weight $\lambda$, then $V^{(\lambda)}=0$, $e$
vanishes on every isotypic subspace, and $e=0$ by
Lemma~\ref{pd-kappa-defects-segment-mean-letters-span}(4). Otherwise let $v_0$ be a highest weight vector
of weight $\lambda$, $L$ the span of the $u(v_0)$ with $u$ in the monoid generated by the $E_{ab}$,
$d=\dim L$, and $F=\operatorname{Hom}(E|_L,E)$, a $\mathfrak{gl}_n$-family on $\operatorname{Hom}_K(L,V)$
(Lemma~\ref{pd-kappa-defects-segment-mean-letters-trace}(2)).

\emph{A basis.} The $\mathcal O$-module $M\cap L$ is finitely generated, as $\mathcal O$ is noetherian, and
torsion-free, so it is free. It spans $L$, as every $y\in L$ has a nonzero multiple in $M$. So it has a
basis $b_1,\dots,b_d$ that is a $K$-basis of $L$, and $b_i^*(y)\in\mathcal O$ for $y\in M\cap L$.

\emph{A lattice for $F$.} Let $H$ be the set of the $\phi\in\operatorname{Hom}_K(L,V)$ with
$\phi(M\cap L)\subset M$, that is, with $\phi(b_i)\in M$ for all $i$. For $x\in M$, $b_i^*\otimes x\in H$.
Let $a\neq b$. The maps $\phi\mapsto E_{ab}\phi$ and $\phi\mapsto\phi E_{ab}$ commute, so
\[
 \frac1{k!}F_{ab}^k\phi=\sum_{i+j=k}(-1)^j\,\frac{E_{ab}^i}{i!}\,\phi\,\frac{E_{ab}^j}{j!}\Big|_L ,
\]
and $H$ is stable under $F_{ab}^k/k!$, as $M\cap L$ is stable under the $E_{ab}^j/j!$. Let $\pi_\mu$ be the
projection of $V$ onto $V_\mu$ along the other weight spaces. It maps $M$ into $M$ by
Lemma~\ref{pd-kappa-defects-segment-mean-letters-separators}(2), and $L$ into $L$ by
Lemma~\ref{pd-kappa-defects-gl-highest}(1). So for $\phi\in H$ each term of
$\phi=\sum_{\mu,\mu'}\pi_{\mu'}\phi\pi_\mu$ lies in $H$; it is a weight vector of $F$ of weight $\mu'-\mu$,
and $\binom{F_{aa}}k$ multiplies it by the integer $\binom{\mu'_a-\mu_a}k$. So $H$ is stable under the
$\binom{F_{aa}}k$.

\emph{The weights of $F$.} Let $\kappa$ be a weight of $F$ with $\sum_a\kappa_a=0$ and $\phi\neq0$ of
weight $\kappa$. As $L$ is the sum of its weight spaces, $\phi(y)\neq0$ for some $y\in L_\nu$, and then
$\phi(y)\in V_{\kappa+\nu}$. The representation $L$ is irreducible with highest weight $\lambda$, so the
entries of its weight $\nu$ lie between $\lambda_n$ and $\lambda_1$ \cite[\S\S20--21]{Hum72}. The weights
of $V$ are permuted by the permutations of the entries \cite[Theorem~21.2]{Hum72}; for a weight $w$ of $V$
and $a\neq b$, $w$ minus $w$ with its $a$-th and $b$-th entries exchanged is
$(w_a-w_b)(\epsilon_a-\epsilon_b)$, so $w_a-w_b\leq\operatorname{rl}(V)$. Choose $c_1,c_2$ with
$\kappa_{c_1}\leq0\leq\kappa_{c_2}$. For every $a$,
\[
 \kappa_a=(\kappa+\nu)_a-\nu_a\leq\operatorname{rl}(V)+(\kappa+\nu)_{c_1}-\nu_a
 \leq\operatorname{rl}(V)+\nu_{c_1}-\nu_a\leq\operatorname{rl}(V)+\operatorname{spr}(\lambda),
\]
and in the same way $-\kappa_a\leq\operatorname{rl}(V)+\nu_a-\nu_{c_2}-\kappa_{c_2}\leq
\operatorname{rl}(V)+\operatorname{spr}(\lambda)$.

\emph{The bound.} Let $P$ be an invariant projector of $F$
(Lemma~\ref{pd-kappa-defects-segment-mean-letters-trace}(1)), and $I$ the space of the $\phi$ killed by
every $F_{ab}$. By that lemma, $P$ maps into $I$, and a nonzero element of $I$ is a highest weight vector
of weight $0$, on which $P$ is the identity; so $P$ is a projection onto $I$. By
Lemma~\ref{pd-kappa-defects-segment-mean-letters-span}(4), its kernel is the sum of the isotypic subspaces
$F^{(\mu)}$ with $\mu\neq0$. This is the sum of the irreducible subrepresentations on which some $F_{ab}$
is nonzero: $F^{(\mu)}$ is the sum of the subrepresentations generated by its highest weight vectors,
which are irreducible \cite[\S\S6.3, 20.2]{Hum72}, and on which some $F_{cc}$ acts as $\mu_c\neq0$; and an
irreducible subrepresentation on which some $F_{ab}$ is nonzero is generated by a highest weight vector
(Lemma~\ref{pd-kappa-defects-segment-mean-letters-span}(2)) whose weight is not $0$, since highest weight
vectors of weight $0$ lie in $I$. So $P$ is the projection $\operatorname{Inv}$ of
Lemma~\ref{pd-gain-weyl-invariant-denominator} for $F$. That lemma, applied to $F$, the module $H$ and the
bound $\operatorname{rl}(V)+\operatorname{spr}(\lambda)$, gives $p^{c_\lambda}P(H)\subset H$. So
$p^{c_\lambda}P(b_i^*\otimes x)(b_i)\in M$ for $x\in M$, and by
Lemma~\ref{pd-kappa-defects-segment-mean-letters-trace}(3),
$p^{c_\lambda}e(x)=d\sum_ip^{c_\lambda}P(b_i^*\otimes x)(b_i)\in M$.
\end{proof}

\begin{lemma}\label{pd-kappa-defects-segment-mean-letters-operators}
Let $K$ be a field of characteristic $0$, $N\geq0$ and $1\leq n\leq r$ integers, a lettering of
$\{0,\dots,N-1\}$ (Definition~\ref{pd-symbols-path-generators}) and $S\subset\{0,\dots,N-1\}$. On
$K[z]=K[z_{i,c}:0\leq i<N,\ 1\leq c\leq r]$ let
$E^S_{ab}=\sum_iz_{i,a}\,\partial/\partial z_{i,b}-\sum_{i'}z_{i',b}\,\partial/\partial z_{i',a}$ for
$1\leq a,b\leq r$, the sums over the vector letters $i$ and the covector letters $i'$ in $S$, and put
$E^{S,n}=(E^S_{ab})_{a,b\leq n}$. Let $[\ ,\ ]$ be the Fock pairing
(Definition~\ref{pd-weight-lattices}(2)).
\begin{enumerate}
\item $E^{S,n}$ is a $\mathfrak{gl}_n$-family
  (Definition~\ref{pd-kappa-defects-segment-mean-letters-families}). A monomial $z^\nu$ is a weight vector
  of weight $(\mathrm{wt}_S(\nu)_a)_{a\leq n}$ (Definition~\ref{pd-symbols-path-weights}(1)), every weight
  vector of weight $\mu$ is a combination of monomials of weight $\mu$, and the $E^S_{ab}$ preserve the
  space $\mathrm{Pol}(B)$ of the polynomials of row degrees $B$, which is finite-dimensional, for every
  $B\colon\{0,\dots,N-1\}\to\mathbf Z_{\geq0}$.
\item If $S\cap S'=\emptyset$, every $E^S_{ab}$ commutes with every $E^{S'}_{cd}$. If $S\subset S'$, then
  $E^{S'}_{ab}=E^S_{ab}+E^{S'\smallsetminus S}_{ab}$.
\item $[E^S_{ab}F,G]=[F,E^S_{ba}G]$ for all $F,G\in K[z]$, and $F=0$ if $[F,G]=0$ for all $G$.
\item Let $F,F'\in\mathrm{Pol}(B)$, and let $z^\nu$ and $z^{\nu'}$ be monomials of $F$ and $F'$. Then
  $\mathrm{wt}_S(\nu)_a-\mathrm{wt}_S(\nu')_b\leq\sum_{i\in S}B(i)$ for all $a,b\leq r$.
\end{enumerate}
\end{lemma}
\leanref{Formalism/SharpRestriction/GrowthBounds/CountingStatementReduction.lean}{Formalism.MainNotions.G.CountingStatementReduction.subOps}
\leanref{Formalism/SharpRestriction/GrowthBounds/CountingStatementReduction.lean}{Formalism.MainNotions.G.CountingStatementReduction.subOps_rep}
\leanref{Formalism/SharpRestriction/GrowthBounds/RootWeights/PolynomialMemWeightSpaceIff.lean}{Formalism.MainNotions.G.GLWeightTheory.polynomial_isGLRep}
\leanref{Formalism/SharpRestriction/GrowthBounds/RootWeights/GlOperatorPreservesMultidegree.lean}{Formalism.MainNotions.G.GLWeightTheory.gl_operator_preserves_multidegree}
\leanref{Formalism/SharpRestriction/GrowthBounds/RootWeights/GlOperatorPreservesMultidegree.lean}{Formalism.MainNotions.G.GLWeightTheory.disjoint_letters_commute}
\leanref{Formalism/SharpRestriction/GrowthBounds/RootWeights/GlOperatorPreservesMultidegree.lean}{Formalism.MainNotions.G.GLWeightTheory.gl_operator_union_add}
\leanref{Formalism/SharpRestriction/GrowthBounds/RootWeights/ExponentFactorialAddSingleOne.lean}{Formalism.MainNotions.G.GLWeightTheory.fock_form_contravariant}
\leanref{Formalism/SharpRestriction/GrowthBounds/CountingStatementReduction.lean}{Formalism.MainNotions.G.CountingStatementReduction.fock_nondeg'}
\leanref{Formalism/SharpRestriction/GrowthBounds/RootWeights/IsotypicProjectorLatticeScaling.lean}{Formalism.MainNotions.G.GLWeightTheory.polynomial_weight_difference_le}

\begin{proof}
(1) For a letter $i$ and $a,b\leq r$ put $D^i_{ab}=z_{i,a}\,\partial/\partial z_{i,b}$. Then
$D^i_{ab}D^i_{cd}-D^i_{cd}D^i_{ab}=\delta_{bc}D^i_{ad}-\delta_{da}D^i_{cb}$, the operators $-D^i_{ba}$ satisfy the
same relations, and $D^i_{ab}$ commutes with $D^{i'}_{cd}$ for $i\neq i'$. As $E^S_{ab}$ is the sum of the
$D^i_{ab}$ over the vector letters and of the $-D^{i'}_{ba}$ over the covector letters in $S$, the relations
hold. Next, $D^i_{aa}z^\nu=\nu(i,a)z^\nu$, so $E^S_{aa}z^\nu=\mathrm{wt}_S(\nu)_az^\nu$; monomials are weight
vectors, and weight vectors of distinct weights are linearly independent, so a weight vector of weight
$\mu$ is a combination of monomials of weight $\mu$ and $K[z]$ is the sum of its weight spaces. The operator
$D^i_{ab}$ moves one unit of exponent from $z_{i,b}$ to $z_{i,a}$, so it preserves row degrees. The space
$\mathrm{Pol}(B)$ has the finite basis of the monomials of row degrees $B$, and every polynomial lies in the
sum of finitely many such spaces. The relations, the weights and the finite-dimensional stable subspaces
for $a,b\leq n$ are among these.

(2) For disjoint $S$ and $S'$ the two operators are sums of $D^i_{ab}$ or $-D^i_{ba}$ for distinct letters.
The second assertion holds by definition.

(3) For a variable $z_t$ and exponent vectors $\kappa,\kappa'$, $[z_tz^\kappa,z^{\kappa'}]$ and
$[z^\kappa,\partial z^{\kappa'}/\partial z_t]$ are both $(\kappa_t+1)\kappa!$ if $\kappa'$ is $\kappa$ plus
the unit vector at $t$, and both $0$ otherwise. So multiplication by $z_t$ and $\partial/\partial z_t$ are
adjoint, and $[D^i_{ab}F,G]=[\partial F/\partial z_{i,b},\partial G/\partial z_{i,a}]=[F,D^i_{ba}G]$; summing
over the letters in $S$ gives the first assertion. As $[F,z^\kappa]=\kappa!F_\kappa$ and $\kappa!\neq0$ in $K$,
the second follows.

(4) For a vector letter $i\in S$ the contribution of $i$ to $\mathrm{wt}_S(\nu)_a-\mathrm{wt}_S(\nu')_b$ is
$\nu(i,a)-\nu'(i,b)\leq\nu(i,a)\leq B(i)$, and for a covector letter it is $\nu'(i,b)-\nu(i,a)\leq B(i)$.
\end{proof}

\begin{lemma}\label{pd-kappa-defects-segment-mean-letters-lattices}
Let $p$ be a prime and $\mathcal O$ a discrete valuation ring whose fraction field $\mathcal K$ has
characteristic $0$, with $p$ in the maximal ideal of $\mathcal O$, and let $v$ be the valuation of
$\mathcal K$ with $v(p)=1$. Let $N$, $r$, the lettering, $S$ and $\mathrm{Pol}(B)$ be as in
Lemma~\ref{pd-kappa-defects-segment-mean-letters-operators} with $K=\mathcal K$, and let
$\mathcal B_{\mathcal O}$ be the $\mathcal O$-span of the monomials $\gamma z^\nu$ with
$v(\gamma)+\frac12v_p(\nu!)\geq0$.
\begin{enumerate}
\item $\mathcal O[z]$ and $\mathcal B_{\mathcal O}$ are stable under the $(E^S_{ab})^k/k!$ with $a\neq b$ and
  the $\binom{E^S_{aa}}k$, $k\geq0$.
\item For every $B\colon\{0,\dots,N-1\}\to\mathbf Z_{\geq0}$, the $\mathcal O$-modules
  $\mathcal O[z]\cap\mathrm{Pol}(B)$ and $\mathcal B_{\mathcal O}\cap\mathrm{Pol}(B)$ are finitely generated
  and span $\mathrm{Pol}(B)$ over $\mathcal K$.
\end{enumerate}
\end{lemma}
\leanref{Formalism/SharpRestriction/GrowthBounds/RootWeights/IntegralPolynomialsEnvelopingStable.lean}{Formalism.MainNotions.G.GLWeightTheory.integral_formula_lattice_holds}
\leanref{Formalism/SharpRestriction/GrowthBounds/RootWeights/DividedPowerMonomialExpansion.lean}{Formalism.MainNotions.G.GLWeightTheory.integral_formula_stability_holds}
\leanref{Formalism/SharpRestriction/GrowthBounds/RootWeights/DividedPowerPadicValuationBound.lean}{Formalism.MainNotions.G.GLWeightTheory.divided_power_padic_valuation_bound}
\leanref{Formalism/SharpRestriction/GrowthBounds/RootWeights/MultidegreePartFinitelyGenerated.lean}{Formalism.MainNotions.G.GLWeightTheory.fock_lattice_finite_spanning_holds}
\leanref{Formalism/SharpRestriction/GrowthBounds/CountingStatementReduction.lean}{Formalism.MainNotions.G.CountingStatementReduction.lat_facts}

\begin{proof}
(1) By Definition~\ref{pd-symbols-path-weights}(2), $(E^S_{ab})^k/k!$ for $a\neq b$ is the operator
$E^{(k)}_{ab}$ there, and $\binom{E^S_{aa}}k$ multiplies $z^\nu$ by $\binom{\mathrm{wt}_S(\nu)_a}k$. Both map
monomials to integral combinations of monomials, so $\mathcal O[z]$ is stable. Let
$v(\gamma)+\frac12v_p(\nu!)\geq0$. The binomials multiply $\gamma z^\nu$ by an integer. A term of
$E^{(k)}_{ab}(\gamma z^\nu)$ is $\gamma c\,z^{\nu_f}$, with $c=\prod_ic_i(f)$; suppose $c\neq0$. For each
letter $i$ let $(s_i,t_i)$ be $(b,a)$ if $i$ is a vector letter and $(a,b)$ if it is a covector letter.
Then $\nu_f$ arises from $\nu$ by moving $f(i)$ from $(i,s_i)$ to $(i,t_i)$, $|c|=\prod_i\binom{\nu(i,s_i)}{f(i)}$,
and comparing the factorials of the two changed entries of each letter gives
\[
 \nu!\prod_i\binom{\nu(i,t_i)+f(i)}{f(i)}=\nu_f!\prod_i\binom{\nu(i,s_i)}{f(i)} .
\]
So $v_p(\nu!)\leq v_p(\nu_f!)+v_p(c)$. Integers prime to $p$ are units of $\mathcal O$, so
$v(\gamma c)=v(\gamma)+v_p(c)$, and
$v(\gamma c)+\frac12v_p(\nu_f!)\geq v(\gamma)+\frac12v_p(\nu!)+\frac12v_p(c)\geq0$. Hence
$\mathcal B_{\mathcal O}$ is stable.

(2) The monomials of row degrees $B$ form a finite basis of $\mathrm{Pol}(B)$, and
$\mathcal O[z]\cap\mathrm{Pol}(B)$ is their $\mathcal O$-span. The coefficient $\gamma$ of an element of
$\mathcal B_{\mathcal O}$ at $z^\nu$ satisfies $v(\gamma)\geq-\frac12v_p(\nu!)$. So
$\mathcal B_{\mathcal O}\cap\mathrm{Pol}(B)$ lies between $\mathcal O[z]\cap\mathrm{Pol}(B)$ and
$p^{-e}(\mathcal O[z]\cap\mathrm{Pol}(B))$, where $e$ is the largest $v_p(\nu!)$ over these monomials. As
$\mathcal O$ is noetherian, both modules are finitely generated, and both span $\mathrm{Pol}(B)$.
\end{proof}
% Lean: the stability of IntPol is integral_formula_stability_holds (divided_power_monomial_expansion);
%   that of B = bbSpec O p _ 1 goes through mem_bbSpec_one_iff' (InBB, coefficientwise) and
%   divided_power_padic_valuation_bound, which concludes v_p(nu!) <= 2 v_p(c) + v_p(nu_f!); the
%   finiteness and spanning are fock_lattice_finite_spanning_holds (multidegree_part_finitely_generated,
%   multidegree_part_spans).

\begin{lemma}\label{pd-kappa-defects-segment-mean-letters-polynomials}
Let $p$, $\mathcal O\subset\mathcal K$, $v$, $N$, $r$, the lettering, $S$ and $\mathcal B_{\mathcal O}$ be as in
Lemma~\ref{pd-kappa-defects-segment-mean-letters-lattices}, let $1\leq n\leq r$, $E=E^{S,n}$
(Lemma~\ref{pd-kappa-defects-segment-mean-letters-operators}) and $T\subset\mathbf Z^n$. Then $E$ has an
isotypic projector for $T$, and every isotypic projector $e$ of $E$ for $T$ satisfies $[eF,G]=[F,eG]$ for
all $F,G\in\mathcal K[z]$. Let $B\colon\{0,\dots,N-1\}\to\mathbf Z_{\geq0}$, $D=\sum_{i\in S}B(i)$ and
$c=N_0(n)\lfloor\log_p(2D+n-1)\rfloor$. Then $e(\mathrm{Pol}(B))\subset\mathrm{Pol}(B)$, and for
$F\in\mathrm{Pol}(B)$: if $F\in\mathcal O[z]$, then $p^ce(F)\in\mathcal O[z]$, and if
$F\in\mathcal B_{\mathcal O}$, then $p^ce(F)\in\mathcal B_{\mathcal O}$.
\end{lemma}
\leanref{Formalism/SharpRestriction/GrowthBounds/CountingStatementReduction.lean}{Formalism.MainNotions.G.CountingStatementReduction.ptool_sub}
\leanref{Formalism/SharpRestriction/GrowthBounds/CountingStatementReduction.lean}{Formalism.MainNotions.G.CountingStatementReduction.subOps_exists}
\leanref{Formalism/SharpRestriction/GrowthBounds/CountingStatementReduction.lean}{Formalism.MainNotions.G.CountingStatementReduction.subOps_selfadj}
\leanref{Formalism/SharpRestriction/GrowthBounds/CountingStatementReduction.lean}{Formalism.MainNotions.G.CountingStatementReduction.subOps_int}
\leanref{Formalism/SharpRestriction/GrowthBounds/CountingStatementReduction.lean}{Formalism.MainNotions.G.CountingStatementReduction.subOps_int_lat}

\begin{proof}
By Lemma~\ref{pd-kappa-defects-segment-mean-letters-operators}(1) and (3), $E$ is a $\mathfrak{gl}_n$-family and
the Fock pairing satisfies the hypotheses of Lemma~\ref{pd-kappa-defects-segment-mean-letters-forms}(3) for
$E$; that lemma and its part (2) give the first two assertions. For $\mu\in\mathbf Z^n$ let $e_\mu$ be an
isotypic projector for $\{\mu\}$. Let $W=\mathrm{Pol}(B)$, a finite-dimensional subrepresentation, and
$F\in W$. By Lemma~\ref{pd-kappa-defects-segment-mean-letters-span}(4), $F=\sum_{\mu\in A}F_\mu$ for a
finite set $A$ and $F_\mu\in V^{(\mu)}$, so $e(F)=\sum_{\mu\in A\cap T}F_\mu=\sum_{\mu\in A\cap T}e_\mu(F)$. Let
$\Sigma$ be the set of the $\mu\in A\cap T$ such that $W$ contains a highest weight vector of weight $\mu$. By
Lemma~\ref{pd-kappa-defects-segment-mean-letters-span}(5), each $e_\mu$ maps $W$ into $W$, restricts to an
isotypic projector of $E|_W$, and vanishes on $W$ if $\mu\notin\Sigma$. So $e(F)=\sum_{\mu\in\Sigma}e_\mu(F)\in W$.
The weights in $\Sigma$ are dominant, as weights of highest weight vectors in a finite-dimensional
representation \cite[\S\S20--21]{Hum72}.

Let $M$ be $\mathcal O[z]\cap W$ or $\mathcal B_{\mathcal O}\cap W$. By
Lemma~\ref{pd-kappa-defects-segment-mean-letters-lattices}, $M$ is finitely generated, spans $W$, and is
stable under the divided powers and the binomials of $E|_W$. By
Lemma~\ref{pd-kappa-defects-segment-mean-letters-integral}, applied to $E|_W$, $M$, $\Sigma$ and the
restrictions of the $e_\mu$, $p^{c'}e(F)\in M$ for $F\in M$, where $c'$ is the largest of the numbers
$N_0(n)\lfloor\log_p(\operatorname{rl}(W)+\operatorname{spr}(\mu)+n-1)\rfloor$, $\mu\in\Sigma$, or $0$. By
Lemma~\ref{pd-kappa-defects-segment-mean-letters-operators}(1), a weight vector of weight $\mu$ is a
combination of monomials $z^\nu$ with $\mathrm{wt}_S(\nu)_a=\mu_a$ for $a\leq n$. If $\mu-\mu'=k(\epsilon_a-\epsilon_b)$
for weights $\mu,\mu'$ of $W$ and $a\neq b$, then $k=\mu_a-\mu'_a\leq D$ by part (4) of that lemma; so
$\operatorname{rl}(W)\leq D$. For $\mu\in\Sigma$, $\operatorname{spr}(\mu)=\mu_1-\mu_n\leq D$ by the same part.
So $c'\leq c$, and $p^ce(F)=p^{c-c'}p^{c'}e(F)\in M$.
\end{proof}
% Lean: ptool_sub; subOps_exists and subOps_selfadj (iso_indep, subOps_top, subOps_orth with
%   fock_form_contravariant); subOps_int_lat and subOps_int, which writes e F through setProj_eq_sum,
%   restricts to Pol d with subrep_restriction_of_reducibility (hR1 to hR6), uses
%   highest_weight_dominant_holds for the dominance and applies isotypic_projector_theorem_field to
%   latticeMultidegreePart L d; the bounds rl(W) <= D and spr(mu) <= D are hb1 and hb2 there.

\begin{lemma}\label{pd-kappa-defects-segment-mean-letters-degrees}
Let $C_0,C_H,C_{E1},C_{E2},C_c,C_\kappa,C_s$ be reals with $C_H\geq1$ and $C_{E1},C_{E2}\geq0$, and let an
instance satisfy its hypotheses (Definition~\ref{pd-kappa-defects-segment-mean-conditions}). Let $\delta'$ be
the degrees of a block lettering, $S$ a set of its letters, $D=\sum_{x\in S}\delta'(x)$ and $1\leq n\leq m$.
Then $m\geq4$, $2D+m\leq p^2(8a+m)^2$ and $N_0(n)\lfloor\log_p(2D+n-1)\rfloor\leq\kappa_1$.
\end{lemma}
\leanref{Formalism/SharpRestriction/GrowthBounds/CountingStatementReduction.lean}{Formalism.MainNotions.G.CountingStatementReduction.kap_num}
\leanref{Formalism/SharpRestriction/GrowthBounds/CountingStatementReduction.lean}{Formalism.MainNotions.G.CountingStatementReduction.block_size}
\leanref{Formalism/SharpRestriction/GrowthBounds/CountingStatementReduction.lean}{Formalism.MainNotions.G.CountingStatementReduction.hiFacts}
\leanref{Formalism/SharpRestriction/GrowthBounds/CountingStatementReduction.lean}{Formalism.MainNotions.G.CountingStatementReduction.exp_le}
\leanref{Formalism/ThetaLogCoordinates/TraceIsometryWeights/SchurTraceProofs/Discharge.lean}{Formalism.InvariantFactorPaths.PW.kapC_le_real}

\begin{proof}
\emph{The hypotheses.} By hypothesis (1) of Theorem~\ref{pd-gain-theorem-one}, $1\leq N_1$ and
$2N_1+2\leq m$, so $m\geq4$, and $2s+2\leq m$, as $s=S_j\leq N_1$; also $k^2\leq k^2Y^2\leq p$, as $Y\geq1$, so
$p\geq(m+1)^2\geq25$. As $q\geq2$, $T\neq\emptyset$, so $|T|Y\geq1$ and $\Lambda\geq8\log(|T|Y)\geq0$. By
hypothesis (2), $1\leq C_H\leq Y_b$ and $2Y_b|T|\Lambda<k$, so $\Lambda<k/2$; and $0<b'$ and $p<a/b'$, so
$a>pb'\geq25b'>0$ and, $a$ being an integer, $a\geq1$. Also $b_l\leq b'\leq b_u$ and $b_u-b_l\leq b'/2$, so
$b_u\leq3b'/2$, and $J_0=\Lambda b'/k\leq b'/2$, so $x_a,x_b\leq2J_0\leq b'$. Hence
$b_T(s')\leq b_u+x_b\leq5b'/2\leq a/10$ for $s'\in T$, and, as $kY\geq5$,
$a+a/(kY)+x_a\leq a+a/5+a/25\leq2a$.

\emph{The degrees.} The block lettering comes from a kept block of $S_j(b_T)$ with $\alpha\in W^s$ and
$\alpha_+\in W$ (Definition~\ref{pd-kappa-defects-segment-lattice}), and $\delta'\leq\delta$ letterwise. Here
$b_i=b_T(y_{j,i})$ and $b_+=b_T(s_{j+1})$ with $y_{j,i},s_{j+1}\in T$, so $b_i,b_+\leq a/10$, and
$\alpha_i\leq2a$. As $C_{E1},C_{E2}\geq0$, $E_0\geq0$, so $\hat a\leq a$; also $\hat b_i\leq b_i$ and
$\hat\alpha_i\leq\alpha_i$. So $\delta(x_i)\leq\hat b_i\leq a/10$ (it is $w_i$ or $o_i$) and
$\delta(y_i)\leq\hat\alpha_i\leq2a$. Further $\operatorname{cap}(s)\leq\hat a+\sum_i\hat b_i\leq a+sa/10$,
$\Lambda/(m+1)<1/2$, $4/p^2<1/2$ and $P_i\geq1$, so
$\delta(x_+)=o_+\leq O_+\leq a+sa/10+2as+b_+\leq a+sa/10+2as+a/10$. Every degree is therefore at most
$2a+sa/10+4as+a/10\leq3am$, and as there are $2s+1\leq m$ letters, $D\leq3am^2$.

\emph{The bounds.} $2D+m\leq6am^2+m\leq64am^2\leq p^2(8a+m)^2$, as $p\geq m$ and $(8a+m)^2\geq64a$. If
$2D+n-1\geq1$, then $\lfloor\log_p(2D+n-1)\rfloor\leq\log_p(p^2(8a+m)^2)=2\ell$, with $\ell$ as in
Definition~\ref{pd-kappa-defects-numbers}. As $2N_0(n)\leq m^2$,
$N_0(n)\lfloor\log_p(2D+n-1)\rfloor\leq m^2\ell\leq\kappa_1$.
\end{proof}
% Lean: kap_num = exp_le with block_size; hiFacts derives the consequences of HI (with CYb >= 1; its
%   proof of 2s+2 <= m uses segS <= N1), segChain_deg_le (cap_le, oPlus_le, base_exponent_le_a) bounds
%   the degrees with A = 2a and B = a/10, degIn_le_sum sums over the 2s+1 letters, final_num gives
%   2D + m <= p^2(8a+m)^2, and kapC_le_real the last bound through Nat.log p (2D + m - 1); r_ge is m >= 4.

\begin{lemma}\label{pd-kappa-defects-segment-mean-letters-highest}
Let $K$ be a field of characteristic $0$, $r\geq2$, $E$ a $\mathfrak{gl}_r$-family on a $K$-vector space $V$
(Definition~\ref{pd-kappa-defects-segment-mean-letters-families}) and $E'=(E_{ab})_{a,b\leq r-1}$. Let $\beta$
be a bilinear form on $V$ with $\beta(E_{ab}x,y)=\beta(x,E_{ba}y)$ for all $a,b,x,y$, such that $x=0$
whenever $\beta(x,y)=0$ for all $y$. Let $\lambda\in\mathbf Z^r$, $\lambda'=(\lambda_1,\dots,\lambda_{r-1})$,
$e'$ an isotypic projector of $E'$ for $\{\lambda'\}$, and $V_\lambda$ the weight space of $E$. Then there is
an endomorphism $P$ of $V$ with $P^2=P$ such that the image of $P$ is the set of the $x\in V_\lambda$ with
$E_{ab}x=0$ for $a<b\leq r-1$, the intersection of the kernel of $P$ with $V_\lambda$ is
$V_\lambda\cap\sum_{b<a\leq r-1}E_{ab}(V)$, $\beta(PF,G)=\beta(F,PG)$ for $F,G\in V_\lambda$, and $PF=e'(F)$
for $F\in V_\lambda$.
\end{lemma}
\leanref{Formalism/SharpRestriction/GrowthBounds/CountingStatementReduction.lean}{Formalism.MainNotions.G.CountingStatementReduction.ext_abs}
\leanref{Formalism/SharpRestriction/GrowthBounds/CountingStatementReduction.lean}{Formalism.MainNotions.G.CountingStatementReduction.ext_poly}
\leanref{Formalism/SharpRestriction/GrowthBounds/CountingStatementReduction.lean}{Formalism.MainNotions.G.CountingStatementReduction.weight_vector_own_isotypic_highest}
\leanref{Formalism/SharpRestriction/GrowthBounds/CountingStatementReduction.lean}{Formalism.MainNotions.G.CountingStatementReduction.weight_vector_other_isotypic_in_lowering_image}

\begin{proof}
The family $E'$ is a $\mathfrak{gl}_{r-1}$-family, as $V_\mu$ lies in the weight space of $E'$ of weight
$(\mu_1,\dots,\mu_{r-1})$, and $\beta$ satisfies the hypotheses of
Lemma~\ref{pd-kappa-defects-segment-mean-letters-forms} for $E'$. Write $V'_\nu$ and $V'^{(\nu)}$,
$\nu\in\mathbf Z^{r-1}$, for the weight spaces and isotypic subspaces of $E'$. Let $\pi$ be the projection of
$V$ onto $V_\lambda$ along the other weight spaces of $E$, and put $P=e'\pi$. Let $I$ be the space of the
$x$ with $E_{ab}x=0$ for $a<b\leq r-1$, and $Z=\sum_{b<a\leq r-1}E_{ab}(V)$.

(a) \emph{$e'$ commutes with $\pi$.} For $c\leq r-1$, $E_{cc}$ maps each $V'^{(\nu)}$ into itself
(Lemma~\ref{pd-kappa-defects-segment-mean-letters-span}(3)); so does $E_{rr}$, which commutes with every
$E_{ab}$, $a,b\leq r-1$, by the relations. By Lemma~\ref{pd-kappa-defects-segment-mean-letters-span}(5),
$e'$ commutes with every $E_{cc}$, so it maps each weight space of $E$ into itself and commutes with $\pi$.

(b) \emph{Two facts on $E'$.} Let $\nu\in\mathbf Z^{r-1}$ and $x\in V'^{(\nu)}$. As in the proof of
Lemma~\ref{pd-kappa-defects-segment-mean-letters-forms}(1), $V'^{(\nu)}$ lies in the span of the $u(w)$
with $w$ a highest weight vector of $E'$ of weight $\nu$ and $u$ a product of operators $E_{cd}$ with
$c>d$, $c\leq r-1$. For $u$ empty, $u(w)\in V'_\nu\cap I$. Otherwise $u(w)\in Z$, and $u(w)$ has a weight
$\nu'\neq\nu$, since $h(\nu')<h(\nu)$ for the function $h$ of the proof of
Lemma~\ref{pd-kappa-defects-segment-mean-letters-span}(2). So $x=y+z$ with $y\in V'_\nu\cap I$ and $z\in Z$
a sum of weight vectors of weights other than $\nu$. If $x\in V'_\nu$, then $z=x-y\in V'_\nu$, so $z=0$ and
$x\in I$. If $x\in V'_{\nu''}$ with $\nu''\neq\nu$: $Z$ is stable under the $E_{cc}$, $c\leq r-1$, by the
relations, so the weight components of $z=x-y$ lie in $Z$, by the proof of
Lemma~\ref{pd-kappa-defects-gl-highest}(1), which uses only the $E_{cc}$; the component of weight $\nu''$
is $x$, so $x\in Z$.

(c) \emph{The image.} By (a), $Px=e'(\pi x)\in V_\lambda$, and $e'(V)\subset V'^{(\lambda')}$ by
Lemma~\ref{pd-kappa-defects-segment-mean-letters-span}(4). As $V_\lambda\subset V'_{\lambda'}$, (b) gives
$Px\in V_\lambda\cap I$. Conversely, a nonzero $x\in V_\lambda\cap I$ is a highest weight vector of $E'$ of
weight $\lambda'$, so $e'x=x$ and $Px=x$. Hence $P^2=P$, with image $V_\lambda\cap I$.

(d) \emph{The kernel.} Let $x\in V_\lambda$ with $Px=e'x=0$. By
Lemma~\ref{pd-kappa-defects-segment-mean-letters-span}(4), $x=x-e'x$ is a sum of elements $x_\nu\in V'^{(\nu)}$
with $\nu\neq\lambda'$. The projection $\pi'$ of $V$ onto $V'_{\lambda'}$ along the other weight spaces of
$E'$ maps each $V'^{(\nu)}$ into itself (Lemma~\ref{pd-kappa-defects-gl-highest}(1)), and $\pi'x=x$; so
$x=\sum_\nu\pi'x_\nu$ with $\pi'x_\nu\in V'^{(\nu)}\cap V'_{\lambda'}\subset Z$ by (b). Conversely, let
$x\in V_\lambda\cap Z$ and $y\in V$. By Lemma~\ref{pd-kappa-defects-segment-mean-letters-forms}(1) for $E$,
$\beta(x,y')=\beta(x,\pi y')$ for all $y'$; by part (2) of that lemma for $E'$, and (a),
$\beta(e'x,y)=\beta(x,e'y)=\beta(x,\pi e'y)=\beta(x,Py)$. As $Py\in I$, $\beta(E_{ab}z',Py)=\beta(z',E_{ba}Py)=0$
for $b<a\leq r-1$ and $z'\in V$, so $\beta(x,Py)=0$. Hence $e'x=0$, and $Px=e'x=0$.

(e) For $F,G\in V_\lambda$, $PF=e'F$, and $\beta(PF,G)=\beta(e'F,G)=\beta(F,e'G)=\beta(F,PG)$ by
Lemma~\ref{pd-kappa-defects-segment-mean-letters-forms}(2) for $E'$.
\end{proof}
% Lean: ext_abs with phi = Fin.castLE (the Lean allows any injective phi), applied by ext_poly with
%   polynomial_isGLRep, ext_rep, complete_reducibility_of_semisimple, fock_form_contravariant and
%   ext_fock_nondeg. P = e' * wtProj; (a) is hdiag, hcd and hcomm (isotypic_projector_commutes); (b) is
%   ext_Z_le, weight_vector_own_isotypic_highest and weight_vector_other_isotypic_in_lowering_image (lowIm_diag,
%   weight_components_mem_invariant); (c) is hmem and hfix; (d) is hoth, hlow (wtProj_stable) and, for the
%   converse, wt_orth, hadj and hcomm; (e) is hPwt and hadj.

\begin{lemma}\label{pd-kappa-defects-segment-mean-letters}
For every $C_c\geq0$ there is $T_0>0$ such that for every $C_0\geq T_0$ there is $T_1>0$ such that for
all $C_{E1},C_{E2}\geq T_1$ there is $T_2>0$ such that for every $C_H\geq T_2$ and all reals
$C_\kappa,C_s$, every instance satisfies conditions (1) and (3) of
Definition~\ref{pd-kappa-defects-segment-mean-conditions} with these constants.
\end{lemma}
\leanref{Formalism/SharpRestriction/GrowthBounds/CountingStatementReduction.lean}{Formalism.MainNotions.G.CountingStatementReduction.positivity_condition_eventually}
\leanref{Formalism/SharpRestriction/GrowthBounds/CountingStatementReduction.lean}{Formalism.MainNotions.G.CountingStatementReduction.exponent_bound_eventually}
\leanref{Formalism/SharpRestriction/GrowthBounds/CountingStatementReduction.lean}{Formalism.MainNotions.G.CountingStatementReduction.positivity_condition_holds}
\leanref{Formalism/SharpRestriction/GrowthBounds/CountingStatementReduction.lean}{Formalism.MainNotions.G.CountingStatementReduction.exponent_condition_holds}

\begin{proof}
Take $T_0=T_1=T_2=1$. Let $C_c\geq0$, $C_0\geq1$, $C_{E1},C_{E2}\geq1$, $C_H\geq1$ and $C_\kappa,C_s$ be
reals, and let an instance satisfy its hypotheses, which conditions (1) and (3) assume. Its ring
$\mathcal R$ is a discrete valuation ring whose fraction field $\mathcal K$ contains $\iota(Q)$, so it has
characteristic $0$, and $p=\varpi_0^2$ lies in the maximal ideal of $\mathcal R$
(Definitions~\ref{pd-local-model}(1) and~\ref{pd-projection-property}). With $v$ the valuation of
$\mathcal K$ with $v(p)=1$, the module $\mathcal B$ is the module $\mathcal B_{\mathcal R}$ of
Lemma~\ref{pd-kappa-defects-segment-mean-letters-lattices}. By
Lemma~\ref{pd-kappa-defects-segment-mean-letters-degrees}, $m\geq4$. Fix a block lettering, with $N=2s+1$
letters and degrees $\delta'$, and $0\leq l\leq2s+1$; let $S$ be the set of the first $l$ letters and
$D=\sum_{x\in S}\delta'(x)$. In the notation of Lemma~\ref{pd-kappa-defects-segment-mean-letters-operators}
with $r=m$, $E^{(l)}=E^{S,m}$ and $E'^{(l)}=E^{S,m-1}$.

\emph{Condition (1): the projectors.} Let $E$ be $E^{S,n}$ with $n\in\{m,m-1\}$, and $T$ a set of weights.
By Lemma~\ref{pd-kappa-defects-segment-mean-letters-polynomials}, $E$ has an isotypic projector for $T$, and
every isotypic projector $e$ for $T$ satisfies $[eF,G]=[F,eG]$ for all $F,G$, and $p^ce(F)\in\mathcal R[z]$
for $F\in\mathrm{Pol}(\delta')\cap\mathcal R[z]$ and $p^ce(F)\in\mathcal B$ for
$F\in\mathrm{Pol}(\delta')\cap\mathcal B$, where $c=N_0(n)\lfloor\log_p(2D+n-1)\rfloor$. By
Lemma~\ref{pd-kappa-defects-segment-mean-letters-degrees}, $c\leq\kappa_1$, and
$p^{\kappa_1}e(F)=p^{\kappa_1-c}p^ce(F)$.

\emph{Condition (1): the commutation.} Let $l\leq l'\leq2s+1$, let $S\subset S'$ be the sets of the first
$l$ and $l'$ letters, and let $(n,n')$ be one of $(m,m)$, $(m,m-1)$ and $(m-1,m-1)$. Let $e$ be an isotypic
projector of $E^{S,n}$ and $e'$ one of $E^{S',n'}$. For $a,b\leq n'$,
$E^{S'}_{ab}=E^S_{ab}+E^{S'\smallsetminus S}_{ab}$
(Lemma~\ref{pd-kappa-defects-segment-mean-letters-operators}(2)). Here $E^S_{ab}$ maps each isotypic subspace
of $E^{S,n}$ into itself, as $n'\leq n$, and so does $E^{S'\smallsetminus S}_{ab}$, which commutes with every
$E^S_{cd}$ (Lemma~\ref{pd-kappa-defects-segment-mean-letters-span}(3)). By
Lemma~\ref{pd-kappa-defects-segment-mean-letters-span}(5), $e$ commutes with every $E^{S'}_{ab}$,
$a,b\leq n'$. So $e$ maps each isotypic subspace of $E^{S',n'}$ into itself, by the same part (3), and $e'$
commutes with $e$ by part (5). The three pairs $(n,n')$ give the three commutations of condition (1).

\emph{Condition (3).} Let $\lambda\in\mathbf Z^m$ and $\lambda'=(\lambda_1,\dots,\lambda_{m-1})$. The
hypothesis $\lambda_1\geq\cdots\geq\lambda_{m-1}$ of condition (3) is not used: the argument below, and
Lemma~\ref{pd-kappa-defects-segment-mean-letters-highest}, hold for every $\lambda$. By the first step there is an isotypic
projector $e'$ of $E'^{(l)}$ for $\{\lambda'\}$. Lemma~\ref{pd-kappa-defects-segment-mean-letters-highest},
applied to $V=\mathcal K[z]$, $r=m$, $E=E^{(l)}$ and the Fock pairing, which satisfies its hypotheses by
Lemma~\ref{pd-kappa-defects-segment-mean-letters-operators}(3), gives an endomorphism $P$ with $P^2=P$ and the
three properties of condition (3) on its image, its kernel and the Fock pairing, and with $PF=e'(F)$ for
$F\in V_\lambda$. For $F\in\mathrm{Pol}(\delta')\cap\mathcal R[z]\cap V_\lambda$,
$p^{\kappa_1}P(F)=p^{\kappa_1}e'(F)\in\mathcal R[z]$ by the first step.
\end{proof}
% Lean: thr_of_ge (thresholds 1; the proofs use CYb >= 1 and CE1, CE2 >= 0) with positivity_condition_holds
%   and exponent_condition_holds. The first: r_ge, chOps_eq and chOps'_eq (subOps with id and Fin.castLE),
%   ptool_sub with kap_num for the projectors, and chOps_comm3 (subOps_comm, subOps_comm_op) for the
%   commutation. The second is exponent_bound_of_positivity, through ext_poly. The ring facts are
%   dvr_fraction_field_char_zero with X.setting (UniformizedDVR) and X.instance_ring_standard
%   (AdmissibleDVR).

\begin{lemma}\label{pd-kappa-defects-segment-mean-weyl-lattices}
Let $\mathcal O$ be a discrete valuation ring whose fraction field $K$ has characteristic $0$, $n\geq1$,
and $(E_{ab})_{a,b\leq n}$ endomorphisms of a $K$-vector space $V$ such that $V$ is the sum of its weight
spaces $V_\mu$ (Definition~\ref{pd-kappa-defects-segment-mean-operators}(1)). Let $\lambda\in\mathbf Z^n$,
$v\in V$, and let $M\subset V$ be an $\mathcal O$-submodule that spans $V$ over $K$, is stable under the
generators of $\mathcal U_{\mathbf Z}$ and has $M\cap V_\lambda=\mathcal Ov$. Then $\Delta(v)\subset M\subset\nabla(\lambda,v)$,
where $\mathcal U_{\mathbf Z}$, $\Delta(v)$ and $\nabla(\lambda,v)$ are formed as in
Definition~\ref{pd-kappa-defects-segment-mean-operators}(3) with $\mathcal O$ in place of $\mathcal R$.
\end{lemma}
\leanref{Formalism/SharpRestriction/GrowthBounds/RootWeights/IntegralEnvelopingPreservesSubmodule.lean}{Formalism.MainNotions.G.GLWeightTheory.highest_weight_lattice_bounds}
\leanref{Formalism/SharpRestriction/GrowthBounds/RootWeights/HighestWeightLatticeOfFirstPart.lean}{Formalism.MainNotions.G.GLWeightTheory.highest_weight_lattice_of_first_part}

\begin{proof}
The hypothesis that $M$ spans $V$ is not used. As $M$ is stable under the generators of
$\mathcal U_{\mathbf Z}$, it is stable under every
$u\in\mathcal U_{\mathbf Z}$. Since $v\in M$, every $u(v)$ lies in $M$, so $\Delta(v)\subset M$. Let $x\in M$
and $u\in\mathcal U_{\mathbf Z}$. Then $u(x)\in M$; write $u(x)=y+z$ with $y\in V_\lambda$ and $z$ a sum of
elements of the $V_\mu$ with $\mu\neq\lambda$. By
Lemma~\ref{pd-kappa-defects-segment-mean-letters-separators}(2), $y\in M$, so $y\in M\cap V_\lambda$ and
$y=cv$ with $c\in\mathcal O$. Then $u(x)-cv=z$, and $x\in\nabla(\lambda,v)$.
\end{proof}
% Lean: highest_weight_lattice_bounds (stated with the equivalence: x in M and in the weight space iff
%   x = c v for some c in O); integral_enveloping_preserves_submodule, weight_component_decomposition,
%   lattice_weight_components_mem and sum_other_weights_mem.

\begin{lemma}\label{pd-kappa-defects-segment-mean-weyl}
Every instance satisfies condition (2) of Definition~\ref{pd-kappa-defects-segment-mean-conditions}.
\end{lemma}
\leanref{Formalism/SharpRestriction/GrowthBounds/Joints.lean}{Formalism.MainNotions.G.Joints.weight_condition_of_integral_form}
\leanref{Formalism/SharpRestriction/GrowthBounds/CovariantDegrees.lean}{Formalism.MainNotions.CovariantDegrees.NablaDeltaComparisonField}
\leanref{Formalism/SharpRestriction/GrowthBounds/GLWeightTheory.lean}{Formalism.MainNotions.G.GLWeightTheory.nabla_delta_comparison_field}
\leanref{Formalism/SharpRestriction/GrowthBounds/Joints.lean}{Formalism.MainNotions.G.Joints.integral_form_exponent_le_ceil}

\begin{proof}
Let $1\leq h\leq m$, $V$, $E=(E_{ab})_{a,b\leq h}$, $\lambda$ and $v$ be as in condition (2). The ring
$\mathcal R$ is a discrete valuation ring whose fraction field $\mathcal K$ contains $\iota(Q)$, so it has
characteristic $0$, and $p=\varpi_0^2$ lies in the maximal ideal of $\mathcal R$, as $\varpi_0$ is a
uniformizer (Definitions~\ref{pd-local-model}(1) and~\ref{pd-projection-property}); so $\mathcal R$ has
residue characteristic $p$.

\emph{The first inclusion.} The $E_{ab}$ satisfy the commutation relations, $V$ is finite-dimensional and
the sum of its weight spaces, $\lambda$ is dominant, $v\in V_\lambda$ is nonzero with $E_{ab}v=0$ for $a<b$,
and $V$ is spanned by the $u(v)$ with $u$ in the monoid generated by the $E_{ab}$. So
Lemma~\ref{pd-gain-weyl-highest-lattices} applies, with $n=h$ and $\mathcal O=\mathcal R$. Its module
$\Delta$ is $\Delta(v)$, and its module $\nabla$ is $\nabla(\lambda,v)$, since $V$ is the direct sum of its
weight spaces and $u(x)-cv\in\sum_{\mu\neq\lambda}V_\mu$ says that the component of $u(x)$ in $V_\lambda$ is
$cv$. That lemma gives $p^{N_0d}\nabla(\lambda,v)\subset\Delta(v)$, where $N_0=h(h-1)/2$ and
$d=\lfloor\log_p(\lambda_1-\lambda_h+h-1)\rfloor$, with $d=0$ if $\lambda_1-\lambda_h+h-1=0$. As $h\leq m$
and $\lambda_1-\lambda_h\leq8a$, we have $N_0\leq m^2$ and $d\leq\log_p(8a+m)\leq\ell$, with $\ell$ as in
Definition~\ref{pd-kappa-defects-numbers}; so $N_0d\leq m^2\ell\leq\kappa_1$. As $\Delta(v)$ is an
$\mathcal R$-module, $p^{\kappa_1}x=p^{\kappa_1-N_0d}\,p^{N_0d}x\in\Delta(v)$ for $x\in\nabla(\lambda,v)$.

\emph{The second.} Let $M$ be a lattice for $E$ with $v\in M$ and $M\cap V_\lambda=\mathcal Rv$. Then $M$
spans $V$ over $\mathcal K$: the $\mathcal K$-span of $M$ is stable under each $E_{ab}$, which is the
generator $E_{ab}^1/1!$ of $\mathcal U_{\mathbf Z}$ if $a\neq b$ and $\binom{E_{aa}}1$ if $a=b$; so it is
stable under the monoid generated by the $E_{ab}$, and as it contains $v$, it contains the $u(v)$, which
span $V$. By Lemma~\ref{pd-kappa-defects-segment-mean-weyl-lattices},
$\Delta(v)\subset M\subset\nabla(\lambda,v)$.
\end{proof}
% Lean: Joints.weight_condition_of_integral_form D X hT12, where hT12 : NablaDeltaComparisonField D comes
%   from tool_statements_hold (GLWeightTheory.problem_holds of CovariantDegrees.admitted); the field is the
%   theorem GLWeightTheory.nabla_delta_comparison_field = integral_form_of_linf_lemma =
%   highest_weight_lattice_of_first_part (restA_of_linf_integrality ...), which does not depend on
%   ExtremalProjectorTools. The first half is LatticeNablaDeltaBound (Lemma pd-gain-weyl-highest-lattices),
%   with integral_form_exponent_le_ceil bounding N0 (n+1) deF p (n+1) (ellWt lam) by kappa_1; the second is
%   highest_weight_lattice_bounds. The Lean passes the spanning of V by M (mq_top_of_irr), which the
%   second half does not use. The ring facts are X.setting (UniformizedDVR) and X.instance_ring_standard.

\begin{lemma}\label{pd-kappa-defects-restriction-universal}
Let $R$ be an integral domain, $r\geq0$, $c\in\mathbf Z$ and $\kappa\in\mathbf Z^r$, and form $A_{r+1}$ and
$H_{r+1}(c,\kappa)$ over $R$ (Definition~\ref{pd-kappa-defects-segment-mean-operators}(4)). For units
$u_1,\dots,u_{r+1}$ of a ring put $\chi_{c,\kappa}(u)=u_1^cu_2^{\kappa_1}\cdots u_{r+1}^{\kappa_r}$. Let $U$ be the
localization of the polynomial ring over $R$ in variables $\beta_{ij}$ ($1\leq j\leq i\leq r+1$) and
$\gamma_{ij}$ ($1\leq i,j\leq r+1$) at the product of the $\beta_{ii}$ and $\det(\gamma_{ij})$. Let $B_0$ be the
lower triangular matrix with entries $\beta_{ij}$ for $j\leq i$, and $G_0=(\gamma_{ij})$. Then $U$ is an integral
domain, $B_0,G_0\in GL_{r+1}(U)$, and every $f\in A_{r+1}$ with
$f(B_0G_0)=\chi_{c,\kappa}(\beta_{11},\dots,\beta_{r+1,r+1})f(G_0)$ lies in $H_{r+1}(c,\kappa)$.
\end{lemma}
\leanref{Formalism/SharpRestriction/GrowthBounds/CountingStatementReduction.lean}{Formalism.MainNotions.G.CountingStatementReduction.UA}
\leanref{Formalism/SharpRestriction/GrowthBounds/CountingStatementReduction.lean}{Formalism.MainNotions.G.CountingStatementReduction.B0}
\leanref{Formalism/SharpRestriction/GrowthBounds/CountingStatementReduction.lean}{Formalism.MainNotions.G.CountingStatementReduction.G0}
\leanref{Formalism/SharpRestriction/GrowthBounds/CountingStatementReduction.lean}{Formalism.MainNotions.G.CountingStatementReduction.HB_of_universal}

\begin{proof}
$U$ is the localization of a polynomial ring over an integral domain at a nonzero element, so it is an
integral domain, and $\det B_0=\prod_i\beta_{ii}$ and $\det G_0$ are units of $U$. Let $A$ be an
$R$-algebra, $g\in GL_{r+1}(A)$, and $b\in GL_{r+1}(A)$ lower triangular with diagonal entries
$u_1,\dots,u_{r+1}$. The $u_i$ are units, since their product is $\det b$. So there is a homomorphism of
$R$-algebras $\psi\colon U\to A$ with $\psi(\beta_{ij})=b_{ij}$ and $\psi(\gamma_{ij})=g_{ij}$; it carries
$B_0$ to $b$ and $G_0$ to $g$. For $h\in GL_{r+1}(U)$ the maps $f\mapsto\psi(f(h))$ and
$f\mapsto f(\psi(h))$ are homomorphisms of $R$-algebras $A_{r+1}\to A$ that agree on the $x_{ab}$, so they
are equal.
Applying $\psi$ to the identity for $f$ gives $f(bg)=\chi_{c,\kappa}(u)f(g)$.
\end{proof}

\begin{lemma}\label{pd-kappa-defects-restriction-basic}
Let $R$ be a commutative ring, $r\geq0$, $c,c'\in\mathbf Z$ and $\kappa,\kappa'\in\mathbf Z^r$, and form $A_{r+1}$,
$H_{r+1}$ and $g_0$ over $R$ (Definition~\ref{pd-kappa-defects-segment-mean-operators}(4)).
\begin{enumerate}
\item $H_{r+1}(c,\kappa)$ is an $R$-submodule of $A_{r+1}$, and
  $H_{r+1}(c,\kappa)H_{r+1}(c',\kappa')\subset H_{r+1}(c+c',\kappa+\kappa')$.
\item $\det(x)^N\in H_{r+1}(N,(N,\dots,N))$ for every $N\in\mathbf Z$, and $g_0\in H_{r+1}(1,(0,\dots,0))$.
\item Let $\varphi\colon R\to S$ be a homomorphism of commutative rings, and let $\varphi_*$ be the homomorphism
  from $A_{r+1}$ to $A_{r+1}$ formed over $S$ that applies $\varphi$ to the coefficients. Then $\varphi_*$ maps
  $H_{r+1}(c,\kappa)$ into $H_{r+1}(c,\kappa)$ formed over $S$, and $\varphi_*(\det(x)^N)=\det(x)^N$.
\item If $R$ is an integral domain, $0\neq a\in R$, $f\in A_{r+1}$ and $af\in H_{r+1}(c,\kappa)$, then
  $f\in H_{r+1}(c,\kappa)$.
\end{enumerate}
\end{lemma}
\leanref{Formalism/SharpRestriction/GrowthBounds/InitialMinorTransposeMemHighestWeight.lean}{Formalism.MainNotions.G.Joints.KP.hbLin}
\leanref{Formalism/SharpRestriction/GrowthBounds/InitialMinorTransposeMemHighestWeight.lean}{Formalism.MainNotions.G.Joints.KP.hbMul}
\leanref{Formalism/SharpRestriction/GrowthBounds/InitialMinorTransposeMemHighestWeight.lean}{Formalism.MainNotions.G.Joints.KP.hbDet}
\leanref{Formalism/SharpRestriction/GrowthBounds/InitialMinorTransposeMemHighestWeight.lean}{Formalism.MainNotions.G.Joints.KP.redBasic}
\leanref{Formalism/SharpRestriction/GrowthBounds/InitialMinorTransposeMemHighestWeight.lean}{Formalism.MainNotions.G.Joints.KP.redHB}
\leanref{Formalism/SharpRestriction/GrowthBounds/CountingStatementReduction.lean}{Formalism.MainNotions.G.CountingStatementReduction.HB_of_smul}
\leanref{Formalism/SharpRestriction/GrowthBounds/InitialMinorTransposeMemHighestWeight.lean}{Formalism.MainNotions.G.Joints.KP.kr_HB_of_smul}

\begin{proof}
Let $A$, $g$, $b$ and $u$ be as in Definition~\ref{pd-kappa-defects-segment-mean-operators}(4), and write
$\chi_{c,\kappa}$ as in Lemma~\ref{pd-kappa-defects-restriction-universal}.

(1) Evaluation at $g$ and at $bg$ are homomorphisms of $R$-algebras $A_{r+1}\to A$, so the $f$ with
$f(bg)=\chi_{c,\kappa}(u)f(g)$ form an $R$-submodule. The second claim follows from
$\chi_{c+c',\kappa+\kappa'}=\chi_{c,\kappa}\chi_{c',\kappa'}$.

(2) $\det(bg)=\det(b)\det(g)=u_1\cdots u_{r+1}\det(g)$, and $\chi_{N,(N,\dots,N)}(u)=(u_1\cdots u_{r+1})^N$. As
$b_{1j}=0$ for $j>1$, the entry $(1,r+1)$ of $bg$ is $u_1g_{1,r+1}$, and $\chi_{1,0}(u)=u_1$.

(3) Let $A$ be an $S$-algebra, viewed as an $R$-algebra through $\varphi$, and $g\in GL_{r+1}(A)$. Then
$(\varphi_*f)(g)=f(g)$, as both sides are homomorphisms in $f$ that agree on $R$ and on the $x_{ab}$. So
the identity for $\varphi_*f$ is the identity for $f$ over the same ring. The map $\varphi_*$ fixes
$\det(x)$, a polynomial with integer coefficients.

(4) Let $U$, $B_0$ and $G_0$ be as in Lemma~\ref{pd-kappa-defects-restriction-universal}, and
$\beta=(\beta_{11},\dots,\beta_{r+1,r+1})$. The identity for $af$ at $B_0$ and $G_0$ is
$af(B_0G_0)=\chi_{c,\kappa}(\beta)af(G_0)$. As $a\neq0$ in the integral domain $U$,
$f(B_0G_0)=\chi_{c,\kappa}(\beta)f(G_0)$, and Lemma~\ref{pd-kappa-defects-restriction-universal} applies.
\end{proof}

\begin{lemma}\label{pd-kappa-defects-restriction-divide}
Let $R$ be an integral domain, $r\geq0$, $c\in\mathbf Z$ and $\kappa\in\mathbf Z^r$, and form $A_{r+1}$, $H_{r+1}$
and $g_0$ over $R$ (Definition~\ref{pd-kappa-defects-segment-mean-operators}(4)). If $f_1\in A_{r+1}$ and
$g_0f_1\in H_{r+1}(c,\kappa)$, then $f_1\in H_{r+1}(c-1,\kappa)$.
\end{lemma}
\leanref{Formalism/SharpRestriction/GrowthBounds/InitialMinorTransposeMemHighestWeight.lean}{Formalism.MainNotions.G.Joints.KP.kp1Dom}

\begin{proof}
Let $U$, $B_0$ and $G_0$ be as in Lemma~\ref{pd-kappa-defects-restriction-universal}, and
$\beta=(\beta_{11},\dots,\beta_{r+1,r+1})$. As $B_0$ is lower triangular, the entry $(1,r+1)$ of $B_0G_0$ is
$\beta_{11}\gamma_{1,r+1}$. So the identity for $g_0f_1$ at $B_0$ and $G_0$ reads
\[
 \beta_{11}\gamma_{1,r+1}f_1(B_0G_0)=\chi_{c,\kappa}(\beta)\gamma_{1,r+1}f_1(G_0)
 =\beta_{11}\gamma_{1,r+1}\chi_{c-1,\kappa}(\beta)f_1(G_0).
\]
As $\beta_{11}\gamma_{1,r+1}\neq0$ in the integral domain $U$,
$f_1(B_0G_0)=\chi_{c-1,\kappa}(\beta)f_1(G_0)$, and Lemma~\ref{pd-kappa-defects-restriction-universal}
gives $f_1\in H_{r+1}(c-1,\kappa)$.
\end{proof}

\begin{lemma}\label{pd-kappa-defects-restriction-prime}
Let $R$ be an integral domain, $n\geq1$ and $A_n=R[x_{ab}:1\leq a,b\leq n][\det(x)^{-1}]$.
\begin{enumerate}
\item Let $q\in R[x_{ab}]$ be a prime element whose value $q(1)$ at the identity matrix is not a unit of $R$,
  and let $e\geq0$. Then $q$ is a prime element of $A_n$. If $f\in A_n$ and $q^e$ divides $f$ in
  $A_n[x_{11}^{-1}]$, then $q^e$ divides $f$ in $A_n$.
\item If $n\geq2$, then $g_0=x_{1n}$ is a prime element of $A_n$; and if $0\neq a\in R$, $y\in A_n$ and $g_0$
  divides $ay$, then $g_0$ divides $y$.
\item Let $R$ be a discrete valuation ring with uniformizer $\varpi$ and residue map $\varphi$, and $\varphi_*$ as in
  Lemma~\ref{pd-kappa-defects-restriction-basic}(3). If $y\in A_n$ and $\varphi_*(y)=0$, then
  $y\in\varpi A_n$.
\end{enumerate}
\end{lemma}
\leanref{Formalism/SharpRestriction/GrowthBounds/CountingStatementReduction.lean}{Formalism.MainNotions.G.CountingStatementReduction.away_pow_dvd}
\leanref{Formalism/SharpRestriction/GrowthBounds/CountingStatementReduction.lean}{Formalism.MainNotions.G.CountingStatementReduction.away_prime}
\leanref{Formalism/SharpRestriction/GrowthBounds/CountingStatementReduction.lean}{Formalism.MainNotions.G.CountingStatementReduction.cell_pow_dvd}
\leanref{Formalism/SharpRestriction/GrowthBounds/InitialMinorTransposeMemHighestWeight.lean}{Formalism.MainNotions.G.Joints.KP.top_right_not_dvd_det}
\leanref{Formalism/SharpRestriction/GrowthBounds/InitialMinorTransposeMemHighestWeight.lean}{Formalism.MainNotions.G.Joints.KP.g0rCancel}
\leanref{Formalism/SharpRestriction/GrowthBounds/InitialMinorTransposeMemHighestWeight.lean}{Formalism.MainNotions.G.Joints.KP.redKer}

\begin{proof}
(1) Let $D$ be an integral domain, $q\in D$ prime and $s\in D$ not divisible by $q$. Then $q$ is prime in
$D[s^{-1}]$, and if $q^e$ divides $z\in D$ in $D[s^{-1}]$, then $zs^k=q^ew$ for some $k\geq0$ and $w\in D$;
as $q$ does not divide $s^k$, $q^e$ divides $z$ in $D$. Now $q$ does not divide $\det(x)$ in $R[x_{ab}]$:
evaluating at the identity matrix, $q(1)$ would divide $1$. So $q$ is prime in $A_n$. If $q$ divided $x_{11}$
in $A_n$, it would divide $x_{11}$ in $R[x_{ab}]$, by the descent with $D=R[x_{ab}]$ and $s=\det(x)$, and
$q(1)$ would divide $1$. So $q$ does not divide
$x_{11}$ in $A_n$, and the descent of divisibility, with $D=A_n$ and $s=x_{11}$, gives the second claim.

(2) The variable $x_{1n}$ is prime in $R[x_{ab}]$, and its value at the identity matrix is $0$ as $n\geq2$. By
(1), $g_0$ is prime in $A_n$. If $g_0$ divided $a$ in $A_n$, then $x_{1n}$ would divide $a$ in $R[x_{ab}]$, by
the descent in (1) with $s=\det(x)$, and comparing constant terms would give $a=0$. So $g_0$ divides $y$.

(3) Write $y=P\det(x)^{-t}$ with $P\in R[x_{ab}]$ and $t\geq0$. Then $\varphi_*(P)=\varphi_*(y)\det(x)^t=0$,
and the polynomial ring over the residue field embeds in its localization, so every coefficient of $P$ lies
in the kernel $\varpi R$ of $\varphi$. Hence $P=\varpi Q$ with $Q\in R[x_{ab}]$, and $y=\varpi Q\det(x)^{-t}$.
\end{proof}

\begin{lemma}\label{pd-kappa-defects-restriction-polynomials}
Let $R$ be a commutative ring, $n\geq1$ and $\nu\in\mathbf Z^n_{\geq0}$, and form $A_n$ and the $H_n$ over $R$
(Definition~\ref{pd-kappa-defects-segment-mean-operators}(4)); write $H_n(\nu)=H_n(\nu_1,(\nu_2,\dots,\nu_n))$.
Regard $R[x_{ab}]$ as a subring of $A_n$, and for $P\in R[x_{ab}]$ let $P^{\mathsf T}(x)=P({}^tx)$. Let
$\nabla_R(\nu)$ be the set of $P\in R[x_{ab}]$ that are homogeneous of degree $\nu_t$ in the entries of the
$t$-th column of $x$ for every $t$ and satisfy $P(x(1+yE_{lj}))=P(x)$ in $R[x_{ab},y]$ for all $l<j$, where
$E_{lj}$ is the matrix unit; for non-increasing $\nu$ this is the dual Weyl module of
Definition~\ref{pd-weight-lattices}(1).
\begin{enumerate}
\item Let $s_1,\dots,s_K\in\{0,\dots,n\}$ with $\#\{k:s_k\geq t\}=\nu_t$ for every $t$, let
  $I_k(1),\dots,I_k(s_k)\in\{1,\dots,n\}$, and let $D$ be the product over $k$ of the minors of $x$ on the
  rows $I_k(1),\dots,I_k(s_k)$ and the first $s_k$ columns. Such a $D$ is a \emph{bideterminant of shape
  $\nu$}. Then $D^{\mathsf T}\in H_n(\nu)$.
\item If $R$ is an integral domain and $P\in R[x_{ab}]$ lies in $H_n(\nu)$, then $P^{\mathsf T}\in\nabla_R(\nu)$.
\end{enumerate}
\end{lemma}
\leanref{Formalism/SharpRestriction/GrowthBounds/InitialMinorTransposeMemHighestWeight.lean}{Formalism.MainNotions.G.Joints.KP.bidetHB}
\leanref{Formalism/SharpRestriction/GrowthBounds/InitialMinorTransposeMemHighestWeight.lean}{Formalism.MainNotions.G.Joints.KP.initial_minor_transpose_mem_highest_weight}
\leanref{Formalism/SharpRestriction/GrowthBounds/InitialMinorTransposeMemHighestWeight.lean}{Formalism.MainNotions.G.Joints.KP.polyHB}
\leanref{Formalism/GradedWeylDivisibility/StatementA.lean}{Formalism.GradedWeylDivisibility.bideterminant}

\begin{proof}
(1) Let $M$ be the minor of $x$ on the rows $I(1),\dots,I(s)$ and the first $s$ columns. Then $M^{\mathsf T}(g)$
is the minor of $g$ on the first $s$ rows and the columns $I(1),\dots,I(s)$. Let $A$, $g$, $b$ and $u$ be as in
Definition~\ref{pd-kappa-defects-segment-mean-operators}(4). As $b_{ij}=0$ for $j>i$, the first $s$ rows of
$bg$ are $b'$ times the first $s$ rows of $g$, where $b'$ is the upper left $s\times s$ block of $b$, which
is lower triangular with diagonal $u_1,\dots,u_s$. So $M^{\mathsf T}(bg)=u_1\cdots u_sM^{\mathsf T}(g)$, that
is, $M^{\mathsf T}\in H_n(w)$ with $w_t=1$ for $t\leq s$ and $w_t=0$ otherwise. By
Lemma~\ref{pd-kappa-defects-restriction-basic}(1), $D^{\mathsf T}$ is in $H_n(\nu')$, where $\nu'_t$ is the
number of $k$ with $s_k\geq t$, which is $\nu_t$.

(2) Let $F$ be the fraction field of $R[x_{ab},y,y_1,\dots,y_n]$, and put $g={}^tx\in GL_n(F)$. With
$b=\operatorname{diag}(y_1,\dots,y_n)$, the identity for $P$ reads
$P(\operatorname{diag}(y)\,{}^tx)=y_1^{\nu_1}\cdots y_n^{\nu_n}P({}^tx)$, that is,
$P^{\mathsf T}(x\operatorname{diag}(y))=y_1^{\nu_1}\cdots y_n^{\nu_n}P^{\mathsf T}(x)$. Multiplying the $t$-th
column of $x$ by $y_t$ thus multiplies $P^{\mathsf T}$ by $\prod_ty_t^{\nu_t}$, so $P^{\mathsf T}$ is homogeneous
of degree $\nu_t$ in the entries of the $t$-th column. For $l<j$ the matrix $b=1+yE_{jl}$ is lower triangular
with diagonal entries $1$, and the identity reads $P((1+yE_{jl})\,{}^tx)=P({}^tx)$, that is,
$P^{\mathsf T}(x(1+yE_{lj}))=P^{\mathsf T}(x)$.
\end{proof}

\begin{lemma}\label{pd-kappa-defects-restriction-standard}
Let $R$ be an integral domain, $n\geq1$ and $\nu\in\mathbf Z^n_{\geq0}$, and let $\nabla_R(\nu)$ be as in
Lemma~\ref{pd-kappa-defects-restriction-polynomials}. Let $E_n(\nu)$ be the set of $n\times n$ matrices $e$ of
nonnegative integers with $\sum_ae(a,t)=\nu_t$ for every $t$ and
$\sum_{b\leq a}e(b,t+1)\leq\sum_{b<a}e(b,t)$ for $1\leq a\leq n$ and $1\leq t<n$.
\begin{enumerate}
\item $\nabla_R(\nu)$ is the $R$-span of the bideterminants of shape $\nu$
  (Lemma~\ref{pd-kappa-defects-restriction-polynomials}(1)).
\item $\nabla_R(\nu)$ has an $R$-basis indexed by $E_n(\nu)$.
\end{enumerate}
\end{lemma}
\leanref{Formalism/GradedWeylDivisibility/FinalC.lean}{Formalism.GradedWeylDivisibility.standardMonomialSpanDomain}
\leanref{Formalism/GradedWeylDivisibility/FinalD.lean}{Formalism.GradedWeylDivisibility.sSBasis}
\leanref{Formalism/GradedWeylDivisibility/StatementD.lean}{Formalism.GradedWeylDivisibility.ssIdx}

If $\nu$ is not non-increasing, there is no bideterminant of shape $\nu$, as $\#\{k:s_k\geq t\}$ does not
increase with $t$, and $E_n(\nu)$ is empty, as the inequality at $a=n$ gives $\nu_{t+1}\leq\nu_t$; so both
parts say $\nabla_R(\nu)=0$. For non-increasing $\nu$ the lemma is the standard basis theorem for
bitableaux, which Lemma~\ref{pd-symbols-primitive-basis}(1) states over a discrete valuation ring
\cite{DEP80}.

\begin{lemma}\label{pd-kappa-defects-restriction-lift}
Let $R$ be a discrete valuation ring with residue field $k$ and residue map $\varphi$, $r\geq0$, $c\in\mathbf Z$
and $\kappa\in\mathbf Z^r$, and let $\varphi_*$ be as in Lemma~\ref{pd-kappa-defects-restriction-basic}(3). For
every $h\in H_{r+1}(c,\kappa)$ formed over $k$ there is $f\in H_{r+1}(c,\kappa)$ formed over $R$ with
$\varphi_*(f)=h$ (Definition~\ref{pd-kappa-defects-segment-mean-operators}(4)).
\end{lemma}
\leanref{Formalism/SharpRestriction/GrowthBounds/InitialMinorTransposeMemHighestWeight.lean}{Formalism.MainNotions.G.Joints.kp_lift}
\leanref{Formalism/SharpRestriction/GrowthBounds/InitialMinorTransposeMemHighestWeight.lean}{Formalism.MainNotions.G.Joints.KP.lift_of}
\leanref{Formalism/SharpRestriction/GrowthBounds/InitialMinorTransposeMemHighestWeight.lean}{Formalism.MainNotions.G.Joints.KP.residue_span_lift_bideterminant}

\begin{proof}
Write $h=P\det(x)^{-N_0}$ with $P\in k[x_{ab}]$ and $N_0\geq0$. Choose $M\geq0$ with $c+M\geq0$ and $\kappa_i+M\geq0$
for all $i$, and put $N=N_0+M$ and $\nu=(c+N,\kappa_1+N,\dots,\kappa_r+N)\in\mathbf Z^{r+1}_{\geq0}$. The polynomial
$P'=P\det(x)^M$ equals $h\det(x)^N$, so it lies in $H_{r+1}(\nu_1,(\nu_2,\dots,\nu_{r+1}))$ over $k$ by
Lemma~\ref{pd-kappa-defects-restriction-basic}(1) and~(2). By
Lemmas~\ref{pd-kappa-defects-restriction-polynomials}(2) and~\ref{pd-kappa-defects-restriction-standard}(1),
$P'^{\mathsf T}$ is a $k$-linear combination of bideterminants of shape $\nu$ over $k$. Each of these is the
image under $\varphi$ of the bideterminant over $R$ on the same rows, whose transpose lies in
$H_{r+1}(\nu_1,(\nu_2,\dots,\nu_{r+1}))$ over $R$ by
Lemma~\ref{pd-kappa-defects-restriction-polynomials}(1). As $\varphi$ is onto,
Lemma~\ref{pd-kappa-defects-restriction-basic}(1) gives $Q\in R[x_{ab}]$ with $Q^{\mathsf T}$ in this space and
$\varphi_*(Q)=P'^{\mathsf T}$. Put $f=Q^{\mathsf T}\det(x)^{-N}$, which lies in $H_{r+1}(c,\kappa)$ by
Lemma~\ref{pd-kappa-defects-restriction-basic}(1) and~(2). As $\varphi_*$ commutes with transposition and fixes
$\det(x)$, $\varphi_*(f)=P'\det(x)^{-N}=h$.
\end{proof}

\begin{lemma}\label{pd-kappa-defects-restriction-count}
For $n\geq1$ and $\nu\in\mathbf Z^n_{\geq0}$ let $e_n(\nu)$ be the number of elements of $E_n(\nu)$
(Lemma~\ref{pd-kappa-defects-restriction-standard}), and let $e_0$ of the empty sequence be $1$.
\begin{enumerate}
\item For $r\geq0$ and $\nu\in\mathbf Z^{r+1}_{\geq0}$, $e_{r+1}(\nu)=\sum_\mu e_r(\mu)$, the sum over the
  $\mu\in\mathbf Z^r$ with $\nu_{t+1}\leq\mu_t\leq\nu_t$ for $1\leq t\leq r$.
\item For $n\geq0$, $\nu\in\mathbf Z^n_{\geq0}$ and $M\geq0$, $e_n(\nu_1+M,\dots,\nu_n+M)=e_n(\nu)$.
\item For $r\geq0$, an integer $c\geq0$ and $\kappa\in\mathbf Z^{r+1}_{\geq0}$ with $\kappa_1\leq c+1$,
  \[
   e_{r+2}(c+1,\kappa_1,\dots,\kappa_{r+1})=e_{r+2}(c,\kappa_1,\dots,\kappa_{r+1})+\sum_\mu e_{r+1}(c+1,\mu_1,\dots,\mu_r),
  \]
  the sum over the $\mu\in\mathbf Z^r$ with $\kappa_{t+1}\leq\mu_t\leq\kappa_t$ for $1\leq t\leq r$.
\end{enumerate}
\end{lemma}
\leanref{Formalism/SharpRestriction/GrowthBounds/InitialMinorTransposeMemHighestWeight.lean}{Formalism.MainNotions.G.Joints.KP.br}
\leanref{Formalism/SharpRestriction/GrowthBounds/InitialMinorTransposeMemHighestWeight.lean}{Formalism.MainNotions.G.Joints.KP.shift_of}
\leanref{Formalism/SharpRestriction/GrowthBounds/InitialMinorTransposeMemHighestWeight.lean}{Formalism.MainNotions.G.Joints.KP.diff_of}

\begin{proof}
(1) Let $e\in E_{r+1}(\nu)$. First, $\sum_{b\leq a}e(b,t)=0$ for $a<t$, by induction on $a$: the inequality at
$(a,t-1)$ bounds this sum by $\sum_{b\leq a-1}e(b,t-1)$, which is $0$ by induction, or empty for $a=1$. In
particular $e(a,r+1)=0$ for $a\leq r$. Put $\mu_t=\sum_{a\leq r}e(a,t)$ and $e'(a,t)=e(a,t)$ for $a,t\leq r$.
Then $\mu_t\leq\nu_t$, the inequality at $(r+1,t)$ reads $\nu_{t+1}\leq\mu_t$, and $e'\in E_r(\mu)$, as its
inequalities are among those of $e$. Conversely, for such $\mu$ and $e'\in E_r(\mu)$ let $e(a,t)=e'(a,t)$ for
$a,t\leq r$, $e(a,r+1)=0$ for $a\leq r$, $e(r+1,t)=\nu_t-\mu_t$ for $t\leq r$ and $e(r+1,r+1)=\nu_{r+1}$. Its
column sums are the $\nu_t$; its inequality at $(r+1,t)$ is $\nu_{t+1}\leq\mu_t$, that at $(a,r)$ with $a\leq r$
has left side $0$, and the others are those of $e'$. So $e\in E_{r+1}(\nu)$. The two maps are inverse to each
other, by the vanishing of $e(a,r+1)$ for $a\leq r$ and the column sums.

(2) It suffices to treat $M=1$. For $n=0$ there is nothing to prove; for $n=r+1$ use induction on $n$. By (1)
both sides are sums over $\mu$, and $\mu\mapsto(\mu_1+1,\dots,\mu_r+1)$ is a bijection from the range for
$\nu$ onto that for $(\nu_1+1,\dots,\nu_n+1)$, with $e_r(\mu_1+1,\dots,\mu_r+1)=e_r(\mu)$ by induction.

(3) By (1), $e_{r+2}(c',\kappa_1,\dots,\kappa_{r+1})$ is the sum of $e_{r+1}(\mu_1,\mu')$ over the
$\mu_1$ with $\kappa_1\leq\mu_1\leq c'$ and the $\mu'\in\mathbf Z^r$ with $\kappa_{t+1}\leq\mu'_t\leq\kappa_t$. As
$\kappa_1\leq c+1$, the range of $\mu_1$ for $c'=c+1$ is that for $c'=c$ together with $\mu_1=c+1$.
\end{proof}

\begin{lemma}\label{pd-kappa-defects-restriction-dimension}
Let $K$ be a field, $r\geq0$, $c\in\mathbf Z$, $\kappa\in\mathbf Z^r$ and $N\in\mathbf Z$ with $c+N\geq0$ and
$\kappa_i+N\geq0$ for all $i$, and put $\nu=(c+N,\kappa_1+N,\dots,\kappa_r+N)$. Then $H_{r+1}(c,\kappa)$, formed
over $K$ (Definition~\ref{pd-kappa-defects-segment-mean-operators}(4)), is a $K$-vector space of dimension
$e_{r+1}(\nu)$ (Lemma~\ref{pd-kappa-defects-restriction-count}).
\end{lemma}
\leanref{Formalism/SharpRestriction/GrowthBounds/InitialMinorTransposeMemHighestWeight.lean}{Formalism.MainNotions.G.Joints.KP.finrankHB_of}
\leanref{Formalism/SharpRestriction/GrowthBounds/InitialMinorTransposeMemHighestWeight.lean}{Formalism.MainNotions.G.Joints.KP.dualWeylImage_det_shift}
\leanref{Formalism/SharpRestriction/GrowthBounds/InitialMinorTransposeMemHighestWeight.lean}{Formalism.MainNotions.G.Joints.KP.det_pow_mul_mem_dualWeyl}

\begin{proof}
Put $n=r+1$ and $H_n(\nu')=H_n(\nu'_1,(\nu'_2,\dots,\nu'_n))$, and for $\nu'\in\mathbf Z^n_{\geq0}$ let
$W(\nu')\subset A_n$ be the image of $\nabla_K(\nu')$ under $P\mapsto P^{\mathsf T}$
(Lemma~\ref{pd-kappa-defects-restriction-polynomials}). Transposition is injective and $K[x_{ab}]$ embeds in
$A_n$, so $\dim W(\nu')=e_n(\nu')$ by Lemma~\ref{pd-kappa-defects-restriction-standard}(2). By
Lemmas~\ref{pd-kappa-defects-restriction-standard}(1), \ref{pd-kappa-defects-restriction-polynomials}(1)
and~\ref{pd-kappa-defects-restriction-basic}(1), $W(\nu')\subset H_n(\nu')$.

Let $M\geq0$. For $P\in\nabla_K(\nu')$, $\det(x)^MP^{\mathsf T}$ is the transpose of the polynomial $\det(x)^MP$, and
lies in $H_n(\nu'_1+M,\dots,\nu'_n+M)$ by Lemma~\ref{pd-kappa-defects-restriction-basic}(1) and~(2); by
Lemma~\ref{pd-kappa-defects-restriction-polynomials}(2) it lies in $W(\nu'_1+M,\dots,\nu'_n+M)$. Both spaces
have dimension $e_n(\nu')$ by Lemma~\ref{pd-kappa-defects-restriction-count}(2), and multiplication by
$\det(x)^M$ is injective, so $\det(x)^MW(\nu')=W(\nu'_1+M,\dots,\nu'_n+M)$.

Let $f\in H_n(c,\kappa)$, and write $f=P\det(x)^{-M'}$ with $P\in K[x_{ab}]$ and $M'\geq0$. Put $M=M'+|N|$ and
$e=N+|N|\geq0$. Then $\det(x)^M\det(x)^Nf=P\det(x)^e$ is a polynomial, and it lies in
$H_n(\nu_1+M,\dots,\nu_n+M)$ by Lemma~\ref{pd-kappa-defects-restriction-basic}(1) and~(2). By
Lemma~\ref{pd-kappa-defects-restriction-polynomials}(2) it lies in $W(\nu_1+M,\dots,\nu_n+M)=\det(x)^MW(\nu)$,
so $\det(x)^Nf\in W(\nu)$. Conversely $\det(x)^{-N}W(\nu)\subset H_n(c,\kappa)$ by
Lemma~\ref{pd-kappa-defects-restriction-basic}(1) and~(2). So $H_n(c,\kappa)=\det(x)^{-N}W(\nu)$, of dimension
$e_n(\nu)$.
\end{proof}

\begin{lemma}\label{pd-kappa-defects-restriction-cell}
Let $R$ be an integral domain, $r\geq0$, $c\in\mathbf Z$ and $\kappa\in\mathbf Z^{r+1}$, and form $A_{r+1}$, $A_{r+2}$
and $H_{r+2}(c,\kappa)$ over $R$ (Definition~\ref{pd-kappa-defects-segment-mean-operators}(4)). Let
$\varepsilon\colon A_{r+2}\to A_{r+1}[X_1,\dots,X_{r+1}]$ be the evaluation at the matrix with first row
$(1,X_1,\dots,X_{r+1})$ and, below it, the column $0$ followed by the matrix $(x_{ab})$ of $A_{r+1}$, whose
determinant $\det(x)$ is a unit. Let $C=A_{r+2}[x_{11}^{-1}]$, and in $C$ put $\theta=x_{11}$,
$W_j=\theta^{-1}x_{1,j+1}$ and $D_{ij}=x_{i+1,j+1}-x_{i+1,1}W_j$ for $1\leq i,j\leq r+1$. Then $C$ is an integral
domain containing $A_{r+2}$, $D\in GL_{r+1}(C)$, and every $f\in H_{r+2}(c,\kappa)$ satisfies
$f=\theta^c\,\varepsilon(f)(W,D)$ in $C$, where $\varepsilon(f)(W,D)$ is the image of $\varepsilon(f)$ under the
homomorphism $A_{r+1}[X]\to C$ that evaluates the coefficients at $D$ and sends $X_j$ to $W_j$.
\end{lemma}
\leanref{Formalism/SharpRestriction/GrowthBounds/FlipOneThree.lean}{Formalism.MainNotions.flipOneThree.resMat}
\leanref{Formalism/SharpRestriction/GrowthBounds/CountingStatementReduction.lean}{Formalism.MainNotions.G.CountingStatementReduction.cell_mul}
\leanref{Formalism/SharpRestriction/GrowthBounds/CountingStatementReduction.lean}{Formalism.MainNotions.G.CountingStatementReduction.bigcell}
\leanref{Formalism/SharpRestriction/GrowthBounds/CountingStatementReduction.lean}{Formalism.MainNotions.G.CountingStatementReduction.RES_eval}
\leanref{Formalism/SharpRestriction/GrowthBounds/CountingStatementReduction.lean}{Formalism.MainNotions.G.CountingStatementReduction.cellFormula}

\begin{proof}
$C$ is an integral domain containing $A_{r+2}$, as $x_{11}\neq0$ in the integral domain $A_{r+2}$. Let $L$ be the
lower triangular matrix over $C$ with first column $(\theta,x_{21},\dots,x_{r+2,1})$, the identity in its lower
right $(r+1)\times(r+1)$ block, and $0$ elsewhere; let $M'$ be the matrix with first row $(1,W_1,\dots,W_{r+1})$
and, below it, the column $0$ followed by $D$. Then $x=LM'$: the first row of $LM'$ is
$(\theta,\theta W_1,\dots)=(x_{11},x_{12},\dots)$, and its entry $(i+1,j+1)$ is $x_{i+1,1}W_j+D_{ij}=x_{i+1,j+1}$.
So $\det(x)=\theta\det D$, and $D$ is invertible over $C$. As $L$ is lower triangular with diagonal entries
$\theta,1,\dots,1$, the identity for $f$ gives $f(x)=f(LM')=\theta^cf(M')$. Here $f(x)=f$, evaluation at the
matrix $x$ being the inclusion $A_{r+2}\subset C$. The matrix $M'$ is the image of the matrix of $\varepsilon$
under the homomorphism $A_{r+1}[X]\to C$ of the statement, and evaluation commutes with homomorphisms, so
$f(M')=\varepsilon(f)(W,D)$.
\end{proof}

\begin{lemma}\label{pd-kappa-defects-segment-mean-restriction}
Every instance satisfies conditions (4) to (7) of
Definition~\ref{pd-kappa-defects-segment-mean-conditions}. The two equations of condition (6) hold for
every field $K'$, of any characteristic.
\end{lemma}
\leanref{Formalism/SharpRestriction/GrowthBounds/CountingStatementReduction.lean}{Formalism.MainNotions.G.CountingStatementReduction.key_positivity_condition}
\leanref{Formalism/SharpRestriction/GrowthBounds/InitialMinorTransposeMemHighestWeight.lean}{Formalism.MainNotions.G.Joints.saturated_restriction_holds}
\leanref{Formalism/SharpRestriction/GrowthBounds/InitialMinorTransposeMemHighestWeight.lean}{Formalism.MainNotions.G.Joints.restriction_dimension_holds}
\leanref{Formalism/SharpRestriction/GrowthBounds/CountingStatementReduction.lean}{Formalism.MainNotions.G.CountingStatementReduction.residue_division_transfer}

\begin{proof}
The four conditions concern rings and the spaces $H_{r+1}(c',\kappa)$ only, and not the instance.

\emph{Condition (4).} A discrete valuation ring is an integral domain, so this is
Lemma~\ref{pd-kappa-defects-restriction-divide}.

\emph{Condition (5).} Let $\varpi$ be a uniformizer of $\mathcal R'$, $\varphi$ the residue map onto the residue
field $k$, and $\varphi_*$ as in Lemma~\ref{pd-kappa-defects-restriction-basic}(3); let $r\geq1$, $c'$ and $\kappa$
be given. We show by induction on $m\geq0$: if $f\in H_{r+1}(c',\kappa)$, $h\in A_{r+1}$ and $g_0$ divides
$f-\varpi^mh$, then $g_0$ divides $h-f'$ for some $f'\in H_{r+1}(c',\kappa)$. For $m=0$ take $f'=f$. For the step,
let $f-\varpi^{m+1}h=g_0k$ with $k\in A_{r+1}$. Applying $\varphi_*$ gives $\varphi_*(f)=g_0\varphi_*(k)$. By
Lemma~\ref{pd-kappa-defects-restriction-basic}(3), $\varphi_*(f)$ lies in $H_{r+1}(c',\kappa)$ over $k$, so
Lemma~\ref{pd-kappa-defects-restriction-divide} over $k$ gives $\varphi_*(k)\in H_{r+1}(c'-1,\kappa)$ over $k$,
and Lemma~\ref{pd-kappa-defects-restriction-lift} gives $f_1\in H_{r+1}(c'-1,\kappa)$ with
$\varphi_*(f_1)=\varphi_*(k)$. By Lemma~\ref{pd-kappa-defects-restriction-basic}(1) and~(2), $g_0f_1$ and
$f-g_0f_1$ lie in $H_{r+1}(c',\kappa)$, and $\varphi_*(f-g_0f_1)=0$. By
Lemma~\ref{pd-kappa-defects-restriction-prime}(3), $f-g_0f_1=\varpi y$ with $y\in A_{r+1}$, and
$y\in H_{r+1}(c',\kappa)$ by Lemma~\ref{pd-kappa-defects-restriction-basic}(4). Now
$\varpi(y-\varpi^mh)=f-\varpi^{m+1}h-g_0f_1=g_0(k-f_1)$, so $g_0$ divides $y-\varpi^mh$ by
Lemma~\ref{pd-kappa-defects-restriction-prime}(2). The induction hypothesis for $y$ and $h$ gives $f'$.

Now let $f$, $h$ and $\gamma$ be as in condition (5), and write $\gamma=u\varpi^m$ with $u$ a unit. Then $g_0$ divides
$f-\varpi^m(uh)$, so $g_0$ divides $uh-f'$ for some $f'\in H_{r+1}(c',\kappa)$. So $u^{-1}f'\in H_{r+1}(c',\kappa)$ by
Lemma~\ref{pd-kappa-defects-restriction-basic}(1), and $g_0$ divides $h-u^{-1}f'=u^{-1}(uh-f')$.

\emph{Condition (6).} Let $K'$ be a field of any characteristic, and let $r\geq1$, $c'$ and $\kappa$ with
$\kappa_1\leq c'$ be given. By
Lemma~\ref{pd-kappa-defects-restriction-basic}(1) each $\langle H_{r+1}(c',\kappa)\rangle$ is
$H_{r+1}(c',\kappa)$, and $\langle\overline{H_{r+1}(c',\kappa)}\rangle$ is the image of $H_{r+1}(c',\kappa)$ in
$A_{r+1}/(g_0)$. Choose $N\geq0$ with $c'-1+N\geq0$ and $\kappa_i+N\geq0$ for all $i$, and put $c_0=c'-1+N$ and
$\kappa'_i=\kappa_i+N$. By Lemma~\ref{pd-kappa-defects-restriction-dimension}, $H_{r+1}(c',\kappa)$ and
$H_{r+1}(c'-1,\kappa)$ have the finite dimensions $e_{r+1}(c_0+1,\kappa')$ and $e_{r+1}(c_0,\kappa')$.

The kernel of the map $H_{r+1}(c',\kappa)\to A_{r+1}/(g_0)$ is $g_0H_{r+1}(c'-1,\kappa)$: an element divisible by
$g_0$ is $g_0f_1$ with $f_1\in H_{r+1}(c'-1,\kappa)$ by Lemma~\ref{pd-kappa-defects-restriction-divide}, and
$g_0H_{r+1}(c'-1,\kappa)\subset H_{r+1}(c',\kappa)$ by Lemma~\ref{pd-kappa-defects-restriction-basic}(1)
and~(2). Multiplication by $g_0\neq0$ is injective on the integral domain $A_{r+1}$, so the kernel has dimension
$\dim H_{r+1}(c'-1,\kappa)$, and the first equation follows.

For the second, Lemma~\ref{pd-kappa-defects-restriction-count}(3), with $r-1$, $c_0$ and $\kappa'$ in place of $r$,
$c$ and $\kappa$ (here $\kappa'_1\leq c_0+1$ as $\kappa_1\leq c'$), gives
$e_{r+1}(c_0+1,\kappa')-e_{r+1}(c_0,\kappa')=\sum_\mu e_r(c_0+1,\mu)$, over the $\mu\in\mathbf Z^{r-1}$ with
$\kappa'_{t+1}\leq\mu_t\leq\kappa'_t$. The map $\kappa''\mapsto(\kappa''_i+N)_i$ is a bijection from the
$\kappa''\in\mathbf Z^{r-1}$ with $\kappa_{i+1}\leq\kappa''_i\leq\kappa_i$ onto these $\mu$, and
$\dim H_r(c',\kappa'')=e_r(c_0+1,(\kappa''_i+N)_i)$ by Lemma~\ref{pd-kappa-defects-restriction-dimension} with
$r-1$ and $N$, as $c'+N=c_0+1$ and $\kappa''_i+N\geq\kappa_{i+1}+N\geq0$. So the sum of the
$\dim H_r(c',\kappa'')$ is $\dim H_{r+1}(c',\kappa)-\dim H_{r+1}(c'-1,\kappa)$, which is
$\dim\langle\overline{H_{r+1}(c',\kappa)}\rangle$ by the first equation.

\emph{Condition (7).} Let $r\geq0$, $c'$ and $\kappa\in\mathbf Z^{r+1}$ be given, and let $C$, $\theta$, $W$ and $D$
be as in Lemma~\ref{pd-kappa-defects-restriction-cell} over $\mathcal R'$. The matrix of $\varepsilon$ is block upper
triangular, so its determinant is $\det(x)$, a unit.

\emph{Injectivity.} If $f,f'\in H_{r+2}(c',\kappa)$ and $\varepsilon(f)=\varepsilon(f')$, then
$f=\theta^{c'}\varepsilon(f)(W,D)=f'$ in $C$ by Lemma~\ref{pd-kappa-defects-restriction-cell}, and $A_{r+2}\subset C$.

\emph{Coefficients.} Let $f\in H_{r+2}(c',\kappa)$, $A$ an $\mathcal R'$-algebra, $g\in GL_{r+1}(A)$ and
$b\in GL_{r+1}(A)$ lower triangular with diagonal $u_1,\dots,u_{r+1}$. Over $A[X_1,\dots,X_{r+1}]$, the block
diagonal matrix with blocks $1$ and $b$ times the matrix with first row $(1,X)$ and lower blocks $0$ and $g$ is
the matrix with first row $(1,X)$ and lower blocks $0$ and $bg$. The first factor is lower triangular with
diagonal $1,u_1,\dots,u_{r+1}$, and the second has determinant $\det g$, a unit. So the identity for $f$ says that
$f$ at the third matrix is $u_1^{\kappa_1}\cdots u_{r+1}^{\kappa_{r+1}}$ times $f$ at the second. As evaluation
commutes with homomorphisms, $f$ at the matrix with lower blocks $0$ and $g'$ is $\varepsilon(f)$ with its
coefficients evaluated at $g'$. Comparing the coefficients of each monomial in $X$, every coefficient $F$ of
$\varepsilon(f)$ satisfies $F(bg)=u_1^{\kappa_1}u_2^{\kappa_2}\cdots u_{r+1}^{\kappa_{r+1}}F(g)$, that is,
$F\in H_{r+1}(\kappa_1,(\kappa_2,\dots,\kappa_{r+1}))$.

\emph{Saturation.} Let $f\in H_{r+2}(c',\kappa)$, $0\neq\gamma\in\mathcal R'$ and $F$ with $\varepsilon(f)=\gamma F$,
and write $\gamma=u\varpi^k$ with $u$ a unit and $\varpi$ a uniformizer. By
Lemma~\ref{pd-kappa-defects-restriction-cell}, $f=\varpi^k\cdot u\theta^{c'}F(W,D)$ in $C$. The constant $\varpi$ is
prime in $\mathcal R'[x_{ab}]$ and is not a unit of $\mathcal R'$, so
Lemma~\ref{pd-kappa-defects-restriction-prime}(1) gives $f=\varpi^kf_0$ with $f_0\in A_{r+2}$. Put
$f'=u^{-1}f_0$. Then $f=\gamma f'$, so $f'\in H_{r+2}(c',\kappa)$ by
Lemma~\ref{pd-kappa-defects-restriction-basic}(4), and $\gamma\varepsilon(f')=\gamma F$ gives $\varepsilon(f')=F$, as
$A_{r+1}[X]$ is an integral domain. The hypothesis on the coefficients of $F$ is not used.

\emph{Divisibility.} The entry $(1,r+2)$ of the matrix of $\varepsilon$ is $X_{r+1}$, so $\varepsilon(g_0)=X_{r+1}$,
and if $g_0^e$ divides $f$ then $X_{r+1}^e$ divides $\varepsilon(f)$. Conversely, let $f\in H_{r+2}(c',\kappa)$ and
$\varepsilon(f)=X_{r+1}^eF_1$. As $W_{r+1}=\theta^{-1}g_0$, Lemma~\ref{pd-kappa-defects-restriction-cell} gives
$f=g_0^e\cdot\theta^{c'-e}F_1(W,D)$ in $C$. The variable $x_{1,r+2}$ is prime in $\mathcal R'[x_{ab}]$ and its value at the
identity matrix is $0$, so Lemma~\ref{pd-kappa-defects-restriction-prime}(1) shows that $g_0^e$ divides $f$ in
$A_{r+2}$.
\end{proof}
% Lean: tKP1 = CountingStatementReduction.key_positivity_condition, the proof of KP.kp1Dom (Lemma
%   pd-kappa-defects-restriction-divide). tKP2 = Joints.saturated_restriction_holds =
%   KP.sizes_bound_of_log_estimates with kp_lift (Lemma pd-kappa-defects-restriction-lift), hbLin, hbMul, hbDet,
%   redBasic, redHB, redKer, kp1Dom and g0rCancel; the induction is KP.top_right_division_lift. tKP3 =
%   Joints.restriction_dimension_holds = KP.lattice_count_of_highest_weight = restriction_dimension_of_division
%   (it introduces the CharZero instance of adm6 and never uses it; FinrankHB, finrank_quotient_add_eq,
%   sum_finrank_shifted_eq and Diff are stated over every field)
%   with finrankHB_of (shift_of br) and diff_of br; the first equation is finrank_quotient_add_eq and the sum is
%   sum_finrank_shifted_eq. tRES = CountingStatementReduction.residue_division_transfer: res_unit, res_inj,
%   res_coeff (through HB_low and resEv_const), res_sat (cell_pow_dvd with C w, HB_of_smul) and res_div
%   (residue_matrix_generic_corner_entry, cW_last and cell_pow_dvd with X (0, last)).

\begin{lemma}\label{pd-kappa-defects-paths-reducible}
Let $K$ be a field of characteristic $0$, $n\geq0$, $V$ a finite-dimensional $K$-vector space and
$(E_{ab})_{a,b\leq n}$ a family of endomorphisms of $V$ with
$E_{ab}E_{cd}-E_{cd}E_{ab}=\delta_{bc}E_{ad}-\delta_{da}E_{cb}$ for all $a,b,c,d$, such that $V$ is the sum of
its weight spaces (Definition~\ref{pd-kappa-defects-segment-mean-operators}(1)). Then every subspace of
$V$ stable under the $E_{ab}$ has a complement stable under the $E_{ab}$.
\end{lemma}
\leanref{Formalism/SharpRestriction/GrowthBounds/CovariantDegrees.lean}{Formalism.MainNotions.CovariantDegrees.adm2}
\leanref{Formalism/SharpRestriction/GrowthBounds/RootWeights/WeightConditionFromIntegralForm.lean}{Formalism.MainNotions.G.GLWeightTheory.weight_condition_from_integral_form}

The $E_{ab}$ make $V$ a representation of $\mathfrak{gl}_n(K)$. The central element $\sum_aE_{aa}$ acts on the
weight space of weight $\mu$ by $\sum_a\mu_a$, so $V$ is the direct sum of its eigenspaces, which are
stable; on each of them the stable subspaces are those of $\mathfrak{sl}_n(K)$, which is semisimple. So
the lemma is Weyl's theorem on complete reducibility, whose proof by the Casimir element works over
every field of characteristic $0$ \cite[\S6.3]{Hum72}.

\begin{lemma}\label{pd-kappa-defects-paths-isotypic}
Let $K$ be a field of characteristic $0$, $N\geq0$ and $r\geq1$ integers, and fix a lettering of
$\{0,\dots,N-1\}$ (Definition~\ref{pd-symbols-path-generators}). For $S\subset\{0,\dots,N-1\}$ let
$E^S=(E^S_{cd})_{c,d\leq r}$ be the operators
$E^S_{cd}=\sum_lz_{l,c}\,\partial/\partial z_{l,d}-\sum_{l'}z_{l',d}\,\partial/\partial z_{l',c}$ on
$K[z]=K[z_{l,c}:0\leq l<N,\ 1\leq c\leq r]$, the sums over the vector letters $l$ and the covector letters
$l'$ in $S$. For $0\leq i\leq N$ put $E^{(i)}=E^{\{0,\dots,i-1\}}$; these are the operators $E^{\mathcal L}_{cd}$
of Definition~\ref{pd-kappa-defects-segment-lattice}(4) for the set $\mathcal L$ of the letters $0,\dots,i-1$.
Fix $S$, and let $V^{(\mu)}$, for $\mu\in\mathbf Z^r$, be the isotypic subspaces of $E^S$
(Definition~\ref{pd-kappa-defects-segment-mean-operators}(1)).
\begin{enumerate}
\item There are endomorphisms $e_\mu$ ($\mu\in\mathbf Z^r$) of $K[z]$, each an isotypic projector for
  $\{\mu\}$, such that for every $G\in K[z]$ only finitely many $e_\mu(G)$ are nonzero, $G$ is their sum,
  and $e_\mu(G)\in V^{(\mu)}$ for every $\mu$.
\item $V^{(\mu)}=0$ unless $\mu_1\geq\cdots\geq\mu_r$.
\item Let $S=\mathcal L=\{0,\dots,i-1\}$ with $0\leq i\leq N$, so that $E^S=E^{(i)}$, and let $e_\mu$ be as in
  (1), $M\in\mathbf Z$ and $G\in K[z]$. Then $G$ lies in $Y(\mathcal L,M)$
  (Definition~\ref{pd-kappa-defects-segment-lattice}(4)) if and only if
  $\lambda_{\{0,\dots,i-1\}}(G)\leq M$ (Definition~\ref{pd-symbols-path-generators}(1)), and this holds
  if and only if $\mu_1\leq M$ for every $\mu$ with $e_\mu(G)\neq0$.
\end{enumerate}
\end{lemma}
\leanref{Formalism/SharpRestriction/GrowthBounds/ProjectionOfIsotypicDecomposition.lean}{Formalism.MainNotions.G.Joints.PolIsoDecompB}
\leanref{Formalism/SharpRestriction/GrowthBounds/ProjectionOfIsotypicDecomposition.lean}{Formalism.MainNotions.G.Joints.polIsoDecompB_of}
\leanref{Formalism/SharpRestriction/GrowthBounds/PolBridge.lean}{Formalism.MainNotions.G.PolBridge.isoSub_eq_bot}
\leanref{Formalism/SharpRestriction/GrowthBounds/PolBridge.lean}{Formalism.MainNotions.G.PolBridge.mem_ycond_iff}
\leanref{Formalism/SharpRestriction/GrowthBounds/PolBridge.lean}{Formalism.MainNotions.G.PolBridge.lowest_degree_le_iff_types}

\begin{proof}
Write $E=E^S$ and $V=K[z]$. \emph{The module.} The $E_{cd}$ satisfy the
relations of $\mathfrak{gl}_r$; a monomial $z^m$ is a weight vector of weight $\mathrm{wt}_S(m)$
(Definition~\ref{pd-symbols-path-weights}(1)), and $E_{cd}$ maps weight vectors of weight $\nu$ to weight
vectors of weight $\nu+\epsilon_c-\epsilon_d$. So $V$ is a weight module
(Definition~\ref{pd-kappa-defects-gl-modules}). The $E_{cd}$ preserve row degrees, so every $x\in V$ lies
in a finite-dimensional subrepresentation, the polynomials whose row degrees are among those of the
monomials of $x$. By Lemma~\ref{pd-kappa-defects-paths-reducible}, every subrepresentation of a
finite-dimensional subrepresentation $P$ has a complement in $P$ that is a subrepresentation. We use three
facts of highest weight theory \cite[\S\S20--21]{Hum72}, applied as in the proof of
Lemma~\ref{pd-kappa-defects-gl-highest}: (a) a nonzero finite-dimensional subrepresentation contains a
nonzero highest weight vector; (b) a nonzero highest weight vector in the span of the $u(v)$, with $u$ in
the monoid generated by the $E_{cd}$ and $v$ highest weight vectors of weights in a set $\Pi$, has its
weight in $\Pi$; (c) the weight of a highest weight vector in a finite-dimensional weight module is
dominant.

(1) The $V^{(\mu)}$ are subrepresentations. They are independent: otherwise some
$V^{(\mu)}\cap\sum_{\nu\neq\mu}V^{(\nu)}$ contains a nonzero $x$, which lies in a finite-dimensional
subrepresentation $P$; the subrepresentation $P\cap V^{(\mu)}\cap\sum_{\nu\neq\mu}V^{(\nu)}$ contains a nonzero
highest weight vector by (a), whose weight is $\mu$ and differs from $\mu$ by (b). They span $V$: let
$x\in V$ lie in the finite-dimensional subrepresentation $P$, and let $P'$ be a complement of
$P\cap\sum_\mu V^{(\mu)}$ in $P$ that is a subrepresentation. If $P'\neq0$, it contains a nonzero highest
weight vector $v$ by (a), and $v\in V^{(\mu)}$ for its weight $\mu$, which contradicts
$P'\cap\sum_\mu V^{(\mu)}=0$. So $P'=0$, and $x\in\sum_\mu V^{(\mu)}$. Hence $V=\bigoplus_\mu V^{(\mu)}$, and the
projections $e_\mu$ onto the summands have the properties of (1). As the summands are
subrepresentations, each $e_\mu$ commutes with the $E_{cd}$.

(2) If $V^{(\mu)}\neq0$, there is a nonzero highest weight vector $v$ of weight $\mu$. It lies in a
finite-dimensional subrepresentation, so $\mu$ is dominant by (c).

(3) \emph{The first equivalence.} By Lemma~\ref{pd-symbols-path-valuation} over $K$,
$\lambda_{\mathcal L}(G)\leq M$ if and only if $\nu(1)\leq M$ for every weight $\nu$ of $U_{\mathcal L}(G)$. As $K$ has
characteristic $0$, $E^{(k)}_{cd}=E_{cd}^k/k!$ and $H^{(k)}_c$ is $\binom{E_{cc}}k$ on weight vectors; so a
subspace is stable under the operators of Definition~\ref{pd-symbols-path-weights}(2) if and only if it is a
subrepresentation, and $U_{\mathcal L}(G)$ is the smallest subrepresentation containing $G$. The weight of
a monomial for $\mathcal L$ (Definition~\ref{pd-kappa-defects-segment-lattice}(4)) is the first entry of its
weight. Suppose $G\in Y(\mathcal L,M)$. Then $U_{\mathcal L}(G)\subset Y(\mathcal L,M)$, as the latter is a
subrepresentation, so a nonzero weight vector of weight $\nu$ in $U_{\mathcal L}(G)$ has a monomial of weight
$\nu(1)\leq M$ for $\mathcal L$. Conversely, suppose $\nu(1)\leq M$ for every weight $\nu$ of $U_{\mathcal L}(G)$. By
Lemma~\ref{pd-kappa-defects-gl-highest}(1), the weight components of an element of $U_{\mathcal L}(G)$ lie in
it, so the weight of each monomial of such an element is a weight of $U_{\mathcal L}(G)$. Hence
$U_{\mathcal L}(G)$ is a subrepresentation in the span of the monomials of weight at most $M$, so it lies in
$Y(\mathcal L,M)$, and so does $G$.

\emph{The second equivalence.} By the first step, it suffices to show: if $e_\mu(G)\neq0$, then $\mu$ is a
weight of $U=U_{\mathcal L}(G)$; and every weight $\nu$ of $U$ has $\nu(1)\leq\mu(1)$ for some $\mu$ with
$e_\mu(G)\neq0$. The space $U$ is a finite-dimensional subrepresentation, so by the argument of (1) applied
inside $U$, $G=\sum_\nu f_\nu$ with $f_\nu$ in the span of the $u(v)$ with $v\in U$ a highest weight vector of
weight $\nu$. Then $f_\nu\in V^{(\nu)}$, so $e_\mu(G)=f_\mu$; if $\mu$ is not a weight of $U$, there is no nonzero
highest weight vector of weight $\mu$ in $U$, and $e_\mu(G)=f_\mu=0$. Next let $x\in U$ be a nonzero weight vector
of weight $\nu$, and write $x=\sum_\mu e_\mu(x)$. As $e_\mu$ commutes with the $E_{cd}$ and $x$ is a combination of
the $u(G)$, $e_\mu(x)=0$ whenever $e_\mu(G)=0$. Each $V^{(\mu)}$ is a sum of the $\langle v\rangle$ with $v$ a highest
weight vector of weight $\mu$, whose weights $\nu'$ satisfy $\nu'\preceq\mu$, so $\nu'(1)\leq\mu(1)$, by
Lemma~\ref{pd-kappa-defects-gl-highest}(2). Taking the weight-$\nu$ component of $x$, some $\mu$ with
$e_\mu(G)\neq0$ has $\nu(1)\leq\mu(1)$.
\end{proof}
% Lean: (1) is polIsoDecompB_of D h2 = isoDecomp (decomp_of_indep_top with isoSub_indep and isoSub_top,
%   through exists_hwv, hwv_wt_mem and cr_of from h2 = adm2, Lemma pd-kappa-defects-paths-reducible), with
%   isGLRep_glOps (the relations, the weight decomposition and polynomial_locally_finite); the Lean states it
%   for every set S of letters. (2) is isoSub_eq_bot (hwv_dominant: polynomial_locally_finite,
%   subrep_restriction_of_reducibility, highest_weight_dominant_holds). (3): mem_ycond_iff, through the
%   sixth clause of valuationsAdditive_D7 (Lemma pd-symbols-path-valuation), polGen_le_ycond, ycond_le,
%   scalarExtension_monomial_weights (binomial_stable_contains_components) and le_sSup; and
%   lowest_degree_le_iff_types = generated_weight_lemma_of_reducibility (generated_weight_lemma_char_zero:
%   polynomial_generated_eq_enveloping from DividedPowerFormulaCharZero and BinomialFormulaCharZero,
%   isotypic_component_weight_mem with StableSubspaceCompleteReducibility, and
%   weight_type_of_isotypic_projections with isotypic_projection_commutes_monoid and
%   HighestWeightModuleWeightBounds). Facts (a) to (c) are exists_hwv, hwv_wt_mem and
%   highest_weight_dominant_holds, which the Lean proves directly; the book cites them.

\begin{lemma}\label{pd-kappa-defects-paths-lower-types}
Let $K$ be a field of characteristic $0$, $N\geq0$ and $n\geq1$ integers, a lettering of $\{0,\dots,N-1\}$, and
$S\subset\{0,\dots,N-1\}$. Let $V^{(\lambda)}$ be the isotypic subspaces of $E^S$ in $K[z]$, as
in Lemma~\ref{pd-kappa-defects-paths-isotypic} with $r=n$, and let $\pi$ be a lower set in $\mathbf Z^n$
(Definition~\ref{pd-symbols-path-weights}(4)). Then $G\in K[z]$ lies in $\sum_{\lambda\in\pi}V^{(\lambda)}$ if and
only if every dominant weight of $U_S(G)$ (Definition~\ref{pd-symbols-path-weights}(3)) lies in $\pi$.
\end{lemma}
\leanref{Formalism/SharpRestriction/GrowthBounds/BinomialPolynomialIdentities/LowerSetTypesSpan.lean}{Formalism.MainNotions.G.ColumnFamilies.RestTypes}
\leanref{Formalism/SharpRestriction/GrowthBounds/BinomialPolynomialIdentities/LiePolarizationDerEqZeroOfDisjoint.lean}{Formalism.MainNotions.G.ColumnFamilies.restTypes}
\leanref{Formalism/SharpRestriction/GrowthBounds/BinomialPolynomialIdentities/LiePolarizationDerEqZeroOfDisjoint.lean}{Formalism.MainNotions.G.ColumnFamilies.mem_iSup_isotypic_iff_projections}
\leanref{Formalism/SharpRestriction/GrowthBounds/BinomialPolynomialIdentities/LiePolarizationDerEqZeroOfDisjoint.lean}{Formalism.MainNotions.G.ColumnFamilies.weights_generated_dominated}
\leanref{Formalism/SharpRestriction/GrowthBounds/BinomialPolynomialIdentities/LiePolarizationDerEqZeroOfDisjoint.lean}{Formalism.MainNotions.G.ColumnFamilies.isotypic_le_restrictSupport_dominated}
\leanref{Formalism/SharpRestriction/GrowthBounds/RootWeights/IsotypicProjectionCommutesOperators.lean}{Formalism.MainNotions.G.GLWeightTheory.isotypic_component_weight_mem}

\begin{proof}
Write $E=E^S$, let $e_\mu$ ($\mu\in\mathbf Z^n$) be the endomorphisms of
Lemma~\ref{pd-kappa-defects-paths-isotypic}(1), and let $\mathcal U$ be the monoid generated by the $E_{cd}$. A
monomial $z^m$ is a weight vector of weight $\mathrm{wt}_S(m)$ for $E$, so a weight vector of weight $\nu$ is a
polynomial all of whose monomials have $S$-weight $\nu$. We use five facts.

(a) \emph{$G\in\sum_{\lambda\in\pi}V^{(\lambda)}$ if and only if $\mu\in\pi$ for every $\mu$ with $e_\mu(G)\neq0$.} If
$G$ lies in the sum and $\mu\notin\pi$, then $e_\mu(G)=0$, as $e_\mu$ vanishes on $V^{(\lambda)}$ for $\lambda\neq\mu$.
Conversely, $G$ is the sum of the $e_\mu(G)\in V^{(\mu)}$, and the nonzero terms have $\mu\in\pi$.

(b) \emph{Each $e_\mu$ commutes with every $u\in\mathcal U$.} It suffices to take $u=E_{cd}$. Each $V^{(\nu)}$ is
stable under $E_{cd}$, as it is spanned by the $u'(v)$ with $u'\in\mathcal U$. Write $x=\sum_\nu x_\nu$ with
$x_\nu\in V^{(\nu)}$, by Lemma~\ref{pd-kappa-defects-paths-isotypic}(1). As $e_\mu$ is the identity on $V^{(\mu)}$
and $0$ on the other $V^{(\nu)}$, $e_\mu(E_{cd}x)=E_{cd}x_\mu=E_{cd}e_\mu(x)$.

(c) \emph{$U_S(G)$ is the span of the $u(G)$ with $u\in\mathcal U$.} As $K$ has characteristic $0$, the operators
of Definition~\ref{pd-symbols-path-weights}(2) are $E^{(k)}_{cd}=E_{cd}^k/k!$ for $c\neq d$ and
$H^{(k)}_c=\binom{E_{cc}}k$, as $E_{cc}$ multiplies $z^m$ by $\mathrm{wt}_S(m)(c)$; so each is a polynomial in the
$E_{cd}$. Conversely $E_{cd}=E^{(1)}_{cd}$ for $c\neq d$ and $E_{cc}=H^{(1)}_c$.

(d) \emph{Every monomial of an element of $V^{(\mu)}$ has $S$-weight $\nu$ with $\nu\preceq\mu$.} As in the proof of
Lemma~\ref{pd-kappa-defects-gl-highest}(2) \cite[\S20.2]{Hum72}, the span of the vectors $u(v)$, with $v$ a highest
weight vector of weight $\mu$ and $u$ a product of operators $E_{cd}$ with $c>d$, is stable under every $E_{ab}$; so it
contains $V^{(\mu)}$. Such an $E_{cd}$ maps a weight vector of weight $\nu$ to one of weight
$\nu+\epsilon_c-\epsilon_d\preceq\nu$, so each $u(v)$ is a weight vector of some weight $\nu\preceq\mu$.

(e) \emph{If $e_\mu(G)\neq0$, then $\mu$ is a weight of $U_S(G)$.} The argument of the proof of
Lemma~\ref{pd-kappa-defects-paths-isotypic}(3) applies to every $S$: the finite-dimensional subrepresentation $U=U_S(G)$ is the sum of
the spans of the $u(v)$ with $u\in\mathcal U$ and $v\in U$ a highest weight vector, by
Lemma~\ref{pd-kappa-defects-paths-reducible} and the argument of part (1) of that lemma; so $G=\sum_\nu f_\nu$
with $f_\nu\in V^{(\nu)}$ spanned by such $u(v)$ with $v$ of weight $\nu$, and $e_\mu(G)=f_\mu$ is $0$ unless $U$
contains a highest weight vector of weight $\mu$.

\emph{Suppose $G\in\sum_{\lambda\in\pi}V^{(\lambda)}$.} Let $\nu$ be a dominant weight of $U_S(G)$, and $x\in U_S(G)$
nonzero with all its monomials of $S$-weight $\nu$. By (c), $x$ is a combination of the $u(G)$ with $u\in\mathcal U$,
and by (b), $e_\mu(u(G))=u(e_\mu(G))$; so $e_\mu(x)=0$ whenever $e_\mu(G)=0$. Write $x=\sum_\mu e_\mu(x)$. For each $\mu$
with $e_\mu(x)\neq0$ we have $e_\mu(G)\neq0$, so $\mu\in\pi$ by (a), and every monomial of $e_\mu(x)\in V^{(\mu)}$ has
$S$-weight $\preceq\mu$ by (d). Hence every monomial of $x$ has $S$-weight $\preceq\lambda$ for some $\lambda\in\pi$. A
monomial of $x$ has $S$-weight $\nu$, so $\nu\preceq\lambda$ for some $\lambda\in\pi$, and $\nu\in\pi$, as $\nu$ is
dominant and $\pi$ is a lower set.

\emph{Suppose that every dominant weight of $U_S(G)$ lies in $\pi$.} Let $e_\mu(G)\neq0$. By (e), $\mu$ is a weight
of $U_S(G)$. As $e_\mu(G)\in V^{(\mu)}$, $V^{(\mu)}\neq0$, so $\mu$ is dominant by
Lemma~\ref{pd-kappa-defects-paths-isotypic}(2). Hence $\mu\in\pi$, and $G\in\sum_{\lambda\in\pi}V^{(\lambda)}$ by (a).
\end{proof}
% Lean: restTypes D (with e and sys from isoDecomp, Lemma pd-kappa-defects-paths-isotypic(1)), for every set S
%   of letters, as in the book; the Lean also allows n = 0. (a) is mem_iSup_isotypic_iff_projections (sys.zero_on, sys.cover, sys.mem);
%   (b) is isotypic_projection_commutes_monoid (isotypic_projection_commutes_operators with
%   isotypicSubmodule_stable and isoSub_top); (c) is polynomial_generated_eq_enveloping; (d) is
%   isotypic_le_restrictSupport_dominated (isotypic_le_lowering_span with hwv_lowering_span_stable,
%   lowering_weight_dominated with dominanceOrder_add_negative_root, monomials_are_weight_vectors); (e) is
%   isotypic_component_weight_mem (stable_subspace_complete_reducibility, isotypic_projection_of_sum). The first
%   direction is weights_generated_dominated with hpi.2; the second uses isoSub_glOps_eq_bot_of_not_dominant
%   (dominant_of_isHighestWeightVector, through hwv_dominant).

\begin{lemma}\label{pd-kappa-defects-paths-lower}
Let $K$ be a field of characteristic $0$, $N\geq0$ and $n\geq1$ integers, a lettering of $\{0,\dots,N-1\}$ and
$B\colon\{0,\dots,N-1\}\to\mathbf Z_{\geq0}$, and let $0\leq i\leq N$. Let $V^{(\lambda)}$ be the isotypic
subspaces of $E^{(i)}$, as in Lemma~\ref{pd-kappa-defects-paths-isotypic} with $r=n$, and let $\pi$ be a lower
set in $\mathbf Z^n$ (Definition~\ref{pd-symbols-path-weights}(4)). Then the set of the
$F\in\sum_{\lambda\in\pi}V^{(\lambda)}$ of row degrees $B$ is the $K$-span of the path generators $\Psi(\mu,I)$
over $K$ of rank $n$ for $B$ with $\mu_i\in\pi$ (Definition~\ref{pd-symbols-path-generators}).
\end{lemma}
\leanref{Formalism/SharpRestriction/GrowthBounds/BinomialPolynomialIdentities/LowerSetTypesSpan.lean}{Formalism.MainNotions.G.ColumnFamilies.LowerSetTypesSpan}
\leanref{Formalism/SharpRestriction/GrowthBounds/BinomialPolynomialIdentities/LiePolarizationDerEqZeroOfDisjoint.lean}{Formalism.MainNotions.G.ColumnFamilies.lower_set_types_span}

\begin{proof}
Put $S=\{0,\dots,i-1\}$, so that $E^S=E^{(i)}$. Let $F\in K[z]$ have row degrees $B$. By
Lemma~\ref{pd-kappa-defects-paths-lower-types}, $F\in\sum_{\lambda\in\pi}V^{(\lambda)}$ if and only if every
dominant weight of $U_S(F)$ lies in $\pi$, and by Lemma~\ref{pd-symbols-path-lower-set} this holds if and only if
$F$ lies in the $K$-span of the $\Psi(\mu,I)$ with $\mu_i\in\pi$. So the set of the lemma is contained in that
span. Conversely, each path generator has row degrees $B$ (proof of Lemma~\ref{pd-symbols-path-basis}), so every
element of the span has row degrees $B$, and by the two equivalences it lies in
$\sum_{\lambda\in\pi}V^{(\lambda)}$.
\end{proof}
% Lean: lower_set_types_span D combines restTypes D (Lemma pd-kappa-defects-paths-lower-types) with
%   admissible_word_basis D (Lemma pd-symbols-path-lower-set), and standardWord_mem_multihomogeneous
%   (pathProductDegrees) for the converse inclusion in Pol.

\begin{lemma}\label{pd-kappa-defects-paths-nested}
Let $\mathcal R'$ be a discrete valuation ring of characteristic $0$ with fraction field $K$, $N\geq0$ and
$n\geq1$ integers, a lettering of $\{0,\dots,N-1\}$ and $B\colon\{0,\dots,N-1\}\to\mathbf Z_{\geq0}$, and let
$E^{(i)}$ and $V^{(\lambda)}$ be as in Lemma~\ref{pd-kappa-defects-paths-isotypic} with $r=n$. For
$0\leq i\leq N$ let $\pi_i$ be a set of $\lambda\in\mathbf Z^n$ with $\lambda_1\geq\cdots\geq\lambda_n$ such that every
such $\nu$ with $\sum_{b\leq a}\nu_b\leq\sum_{b\leq a}\lambda_b$ for all $a$, with equality for $a=n$, for some
$\lambda\in\pi_i$, lies in $\pi_i$. Let $\mathcal I\subset\{0,\dots,N\}$. Then the set of the
$F\in\mathcal R'[z]$ of row degrees $B$ such that $F\in\sum_{\lambda\in\pi_i}V^{(\lambda)}$, for $E^{(i)}$, for
every $i\in\mathcal I$ is the $\mathcal R'$-span of the path generators $\Psi(\mu,I)$ over $\mathcal R'$ of
rank $n$ for $B$ (Definition~\ref{pd-symbols-path-generators}) with $\mu_i\in\pi_i$ for every
$i\in\mathcal I$.
\end{lemma}
\leanref{Formalism/SharpRestriction/GrowthBounds/StatementOfDegreeAndBound.lean}{Formalism.MainNotions.CovariantDegrees.SaturatedLowerSetSpan}
\leanref{Formalism/SharpRestriction/GrowthBounds/StatementOfDegreeAndBound.lean}{Formalism.MainNotions.CovariantDegrees.nestLat}
\leanref{Formalism/SharpRestriction/GrowthBounds/StatementOfDegreeAndBound.lean}{Formalism.MainNotions.CovariantDegrees.nestSpan}
\leanref{Formalism/SharpRestriction/GrowthBounds/StatementOfDegreeAndBound.lean}{Formalism.MainNotions.CovariantDegrees.DomLowerSet}
\leanref{Formalism/SharpRestriction/GrowthBounds/BinomialPolynomialIdentities.lean}{Formalism.MainNotions.G.ColumnFamilies.Top.b15}

\begin{proof}
Let $X$ be the set of the pairs $(\mu,I)$ of a path $\mu$ of rank $n$ for $B$ and a standard filling $I$ of $\mu$,
and write $\Psi_{\mathcal R'}(x)$ and $\Psi_K(x)$ for the path generator of $x\in X$ over $\mathcal R'$ and over $K$.
The column functions are given by the same formulas over both rings, so $\Psi_K(x)$ is the image of
$\Psi_{\mathcal R'}(x)$ under the inclusion $\mathcal R'[z]\subset K[z]$, which we write $F\mapsto\bar F$. By
Lemma~\ref{pd-symbols-path-basis} over $\mathcal R'$ and over $K$, each family is linearly independent and
spans the polynomials of row degrees $B$. The field $K$ has characteristic $0$, and each $\pi_i$ is a lower set
(Definition~\ref{pd-symbols-path-weights}(4)).

Let $F\in\mathcal R'[z]$ have row degrees $B$, with $\bar F\in\sum_{\lambda\in\pi_i}V^{(\lambda)}$ for every
$i\in\mathcal I$. Write $F=\sum_xc_x\Psi_{\mathcal R'}(x)$ with $c_x\in\mathcal R'$; then
$\bar F=\sum_xc_x\Psi_K(x)$. Let $i\in\mathcal I$. By Lemma~\ref{pd-kappa-defects-paths-lower}, $\bar F$ is a
$K$-combination of the $\Psi_K(x)$ with $\mu_i\in\pi_i$. By the linear independence of the $\Psi_K(x)$, the two
expressions agree, so $c_x=0$ unless $\mu_i\in\pi_i$. Hence $F$ lies in the $\mathcal R'$-span of the
$\Psi_{\mathcal R'}(x)$ with $\mu_i\in\pi_i$ for every $i\in\mathcal I$.

Conversely, each such $\Psi_{\mathcal R'}(x)$ has row degrees $B$, and for $i\in\mathcal I$ its image $\Psi_K(x)$ lies
in $\sum_{\lambda\in\pi_i}V^{(\lambda)}$ by Lemma~\ref{pd-kappa-defects-paths-lower}. Both conditions are
preserved by sums and by multiplication by elements of $\mathcal R'$.
\end{proof}
% Lean: ColumnFamilies.Top.b15 D : SaturatedLowerSetSpan D = admissible_lowerSet_span =
%   saturatedLowerSetSpan_of_lowerSetTypesSpan D (lower_set_types_span D), whose first part is the lemma
%   (nestLat = nestSpan); the second part, on intersections modulo p^E, is not used here. The proof
%   follows the first part: standardPath_basis_any_ring over O and over KK O (Lemma
%   pd-symbols-path-basis), map_Psi', span_words_eq_polynomials and span_standard_words_eq, hT from
%   LowerSetTypesSpan (Lemma pd-kappa-defects-paths-lower), and hLIK for the uniqueness of the coefficients.

\begin{lemma}\label{pd-kappa-defects-segment-mean-paths}
Every instance satisfies condition (8) of Definition~\ref{pd-kappa-defects-segment-mean-conditions}.
\end{lemma}
\leanref{Formalism/SharpRestriction/GrowthBounds/PolBridge.lean}{Formalism.MainNotions.G.PolBridge.combined_weight_of_nested_sets}

\begin{proof}
Let a lettering of $N$ letters with degrees $B$, $q'\geq0$, $J_1,\dots,J_{q'}$ and $M_1,\dots,M_{q'}$ be as in
condition (8). The ring $\mathcal R$ of the setting is a discrete valuation ring whose fraction field
$\mathcal K$ has characteristic $0$. Put $\mathcal I=\{J_1,\dots,J_{q'}\}$, and for $0\leq i\leq N$ let $\pi_i$
be the set of the $\lambda\in\mathbf Z^m$ with $\lambda_1\geq\cdots\geq\lambda_m$ and $\lambda_1\leq M_k$ for every $k$ with
$J_k=i$. If $\nu$ is such that its partial sums are at most those of some $\lambda\in\pi_i$, then
$\nu_1\leq\lambda_1$; so $\pi_i$ satisfies the hypothesis of Lemma~\ref{pd-kappa-defects-paths-nested}.

Let $0\leq i\leq N$, let $e_\mu$ be as in Lemma~\ref{pd-kappa-defects-paths-isotypic}(1) for $E^{(i)}$, with
$r=m$, and let $G\in\mathcal K[z]$. We show that $G\in\sum_{\lambda\in\pi_i}V^{(\lambda)}$ if and only if
$G\in Y(\mathcal L_k,M_k)$ for every $k$ with $J_k=i$. First, $G\in\sum_{\lambda\in\pi_i}V^{(\lambda)}$ if and only
if $\mu\in\pi_i$ for every $\mu$ with $e_\mu(G)\neq0$: if $G$ lies in the sum, then $e_\mu(G)=0$ for
$\mu\notin\pi_i$, as $e_\mu$ vanishes on $V^{(\lambda)}$ for $\lambda\neq\mu$; conversely $G$ is the sum of the
$e_\mu(G)\in V^{(\mu)}$. If $e_\mu(G)\neq0$, then $V^{(\mu)}\neq0$, so $\mu_1\geq\cdots\geq\mu_m$ by
Lemma~\ref{pd-kappa-defects-paths-isotypic}(2). So the condition says that $\mu_1\leq M_k$ whenever
$e_\mu(G)\neq0$ and $J_k=i$. For $M_k\in\mathbf Z$ this is $G\in Y(\mathcal L_k,M_k)$ by
Lemma~\ref{pd-kappa-defects-paths-isotypic}(3), and for $M_k=+\infty$ there is no condition.

Hence an $F\in\mathrm{Pol}(B)\cap\mathcal R[z]$ lies in $Y(\mathcal L_k,M_k)$ for every $k$ if and only if
$F\in\sum_{\lambda\in\pi_i}V^{(\lambda)}$ for every $i\in\mathcal I$. By Lemma~\ref{pd-kappa-defects-paths-nested},
these $F$ form the $\mathcal R$-span of the path generators $\Psi(\mu,I)$ over $\mathcal R$ with
$\mu_i\in\pi_i$ for $i\in\mathcal I$. As the entries of a path are non-increasing, this condition says that
$\mu_{J_k}(1)\leq M_k$ for every $k$. The column functions are defined by the same formulas over
$\mathcal R$ and over $\mathcal K$, so $\Psi(\mu,I)$ over $\mathcal R$ is $\Psi(\mu,I)$ over $\mathcal K$.
\end{proof}
% Lean: instanceClauseEight of Loc8All = PolBridge.combined_weight_of_nested_sets D X
%   (GLWeightTheory.weight_condition_from_integral_form D) (ColumnFamilies.Top.b15 D). The proof follows it:
%   Ix and pi (hpi from DomLE at the first entry), hcond (mem_iSup_iso_iff, isoSub_eq_bot, ycond_top,
%   mem_ycond_iff and lowest_degree_le_iff_types with polIsoDecompB_of), hmemL and hY for YjtZ, nestLat =
%   nestSpan from b15, and hgen with StandardMonomialSpace.Psi_map.

\begin{lemma}\label{pd-kappa-defects-segment-mean-ratio}
For every $C_c\geq0$ there is $T_0>0$ such that for every $C_0\geq T_0$ there is $T_1>0$ such that for
all $C_{E1},C_{E2},C_m\geq T_1$ there is $T_2>0$ such that for every $C_H\geq T_2$ there is
$C^0_\kappa>0$ with the following property for every $C_\kappa\geq C^0_\kappa$: every instance satisfies
condition (9) of Definition~\ref{pd-kappa-defects-segment-mean-conditions} with these constants and
$C_s=C_\kappa(C_m+4C_{E1}+C_{E2}+16)$.
\end{lemma}
\leanref{Formalism/SharpRestriction/GrowthBounds/Joints.lean}{Formalism.MainNotions.G.Joints.extension_gives_ratio_bound}
\leanref{Formalism/SharpRestriction/GrowthBounds/Joints.lean}{Formalism.MainNotions.G.Joints.kap_margin}
\leanref{Formalism/SharpRestriction/GrowthBounds/Joints.lean}{Formalism.MainNotions.G.Joints.extended_hypothesis_of_window}

\begin{proof}
Let $C_c\geq0$, and let $T_0$, $T_1$, $T_2$ and $C^0_\kappa$ be the numbers of
Lemma~\ref{pd-kappa-defects-windows-model}, replaced by $\max(T_0,1)$, $\max(T_1,1)$,
$\max\big(T_2,1000(C_{E1}+C_{E2}+1)(C_c+1)\big)$ and $\max(C^0_\kappa,1)$. Let $C_\kappa\geq C^0_\kappa$,
let an instance satisfy its hypotheses, and let $a_T$ be a window function. As $x_a,x_b\leq2J_0$,
$a_T$ and $b_T$ satisfy $a\leq a_T(s')\leq a+a/(kY)+2J_0$ and $b_l\leq b_T(s')\leq b_u+2J_0$ for $s'\in T$.
So Lemma~\ref{pd-kappa-defects-windows-model} applies and gives the first claim for every $\mu'\geq0$.
The constants also satisfy the hypotheses of Lemma~\ref{pd-kappa-defects-margin-share}, whose
constant $C_s$ is the one in the statement; it gives the second claim.
\end{proof}
% Lean: extension_gives_ratio_bound with hAlt2 = ExtendedModelRank (Lemma pd-kappa-defects-windows-model)
% and h2Ext = ExtendedKappaBound (Lemma pd-kappa-defects-margin, rest2 with Cs = Ckap (Cm + 4 CE1 + CE2
% + 15)); extended_hypothesis_of_window passes from the windows x_a, x_b to 2 J0; the second claim is
% kap_margin with loc_kap_facts (m0 >= 0), whose constant is Cs + Ckap.

\subsection{Digits and the mean of the segment defect}\label{subsection-pd-digits-segment-mean}

\begin{lemma}\label{pd-kappa-defects-digits-half}
Let $k\geq5$, $Q\geq1$, $Y\geq1$, $N_1\geq1$, $\Lambda\geq24$, $b'>0$, $a$, $\hat a$, $\mu$ and $M$ be real numbers
with $kY\Lambda b'\leq a$, $k^3Q\leq b'$, $a/(2Q)\leq\hat a$, $0\leq\mu\leq b'/(kQY)$ and $0\leq M<\mu+1$. Then
$M\leq\hat aN_1\Lambda/k$ and
\[
 \big(1-e^{-\Lambda/8}\big)\big(1-e^{-\Lambda/3}\big)\Big(1-\frac{(k-1)M}{\hat a}\Big)\geq\frac12 .
\]
\end{lemma}
\leanref{Formalism/SharpRestriction/GrowthBounds/BinomialRatio/LeAndHalfLeProdOneSubExpNegOfBounds.lean}{Formalism.MainNotions.ColumnSystems.Pf.le_and_half_le_prod_one_sub_exp_neg_of_bounds}

\begin{proof}
As $Y\geq1$, $\mu kQ\leq\mu kQY\leq b'$, and $kQ\leq k^3Q\leq b'$; so $MkQ\leq(\mu+1)kQ\leq2b'$. As $kY\geq5$ and
$\Lambda\geq24$, $120b'\leq kY\Lambda b'\leq a\leq2Q\hat a$. So $\hat a>0$ and $kMQ\leq2b'\leq Q\hat a/30$, that is,
$kM\leq\hat a/30$. As $N_1\Lambda\geq24$, this gives $kM\leq\hat aN_1\Lambda$, the first claim, and
$0\leq(k-1)M/\hat a\leq1/30$. As $e^x\geq1+x$, $e^3\geq4$ and $e^8\geq9$, so $e^{-\Lambda/8}\leq e^{-3}\leq1/4$ and
$e^{-\Lambda/3}\leq e^{-8}\leq1/9$. The product is therefore at least
$\frac34\cdot\frac89\cdot\frac{29}{30}=\frac{29}{45}\geq\frac12$.
\end{proof}

\begin{definition}\label{pd-kappa-defects-digits-chains}
Let $r\geq1$ and $N\geq0$ be integers, and fix a lettering of $\{0,\dots,N-1\}$ and
$B\colon\{0,\dots,N-1\}\to\mathbf Z_{\geq0}$ (Definition~\ref{pd-symbols-path-generators}).
\begin{enumerate}
\item The \emph{Weyl dimension} of $\lambda\in\mathbf Z^r$ is
  $\operatorname{wd}(\lambda)=\prod_{1\leq a<b\leq r}(\lambda_a-\lambda_b+b-a)/(b-a)$.
\item Let $q'\geq0$, $J_1,\dots,J_{q'}\in\{0,\dots,N\}$ and $M_1,\dots,M_{q'}\in\mathbf Z\cup\{+\infty\}$. The
  \emph{chain sum} $\chi_B(J,M)$ is the sum of $\operatorname{wd}(\mu_N)$ over the paths $\mu$ of rank $r$ for $B$
  (Definition~\ref{pd-symbols-path-generators}(2)) with $\mu_{J_k}(1)\leq M_k$ for every $k$.
\item Let a segment datum $S$ of length $s$, integers $\alpha_1,\dots,\alpha_s,\alpha_+\geq0$ and a real number $\mu$
  be as in Definition~\ref{pd-kappa-defects-segment-lattice}; the blocks, the kept blocks, their levels and
  the degrees $\delta'$ do not involve $\mathcal O$ or $v$. For a block $b$ and a level $\ell$ of $b$, with last
  entry $j_+$, let $\chi_b(\ell)$ be the chain sum of rank $r$ for the letters $x_1,y_1,\dots,x_s,y_s,x_+$ with the
  degrees $\delta'$ of $\ell$, where $J$ lists the numbers of letters of $\mathcal L_1,\dots,\mathcal L_{s-\eta}$ and
  $\mathcal L_+$, and $M$ lists $M_1,\dots,M_{s-\eta}$ and $M_+$. The \emph{chain count} of $b$ is
  $c(b)=\sum_\ell\chi_b(\ell)\binom{b_+-j_++r-1}{r-1}$, the sum over the levels $\ell$ of $b$.
\item In the situation of Definition~\ref{pd-kappa-defects-segment-mean-conditions}, for $a_T$ and $b_T$, an
  index $1\leq j'\leq q$, a real number $\mu'\geq0$ and digits $\rho,\rho'$ for $j'$, the \emph{chain count}
  $c_{j'}(\mu';\rho,\rho')$ is the sum of the $c(b)$ over the kept blocks $b$ of $S_{j'}(b_T)$, formed with
  $\alpha_i=a_T(y_{j',i})$, $\alpha_+=a_T(s_{j'+1})$, $\mu'$ and $r=m$, whose vectors $\rho_i$ and $\rho'_i$ are the
  $\rho_i$ and $\rho'_i$.
\end{enumerate}
\end{definition}
\leanref{Formalism/SharpRestriction/GrowthBounds/BinomialRatio/CountDefs.lean}{Formalism.MainNotions.ColumnSystems.Pf.wdim}
\leanref{Formalism/SharpRestriction/GrowthBounds/BinomialRatio/CountDefs.lean}{Formalism.MainNotions.ColumnSystems.Pf.Strip}
\leanref{Formalism/SharpRestriction/GrowthBounds/BinomialRatio/CountDefs.lean}{Formalism.MainNotions.ColumnSystems.Pf.IsChain}
\leanref{Formalism/SharpRestriction/GrowthBounds/BinomialRatio/CountDefs.lean}{Formalism.MainNotions.ColumnSystems.Pf.ChainOK}
\leanref{Formalism/SharpRestriction/GrowthBounds/BinomialRatio/CountDefs.lean}{Formalism.MainNotions.ColumnSystems.Pf.chainSum}
\leanref{Formalism/SharpRestriction/GrowthBounds/BinomialRatio/CountDefs.lean}{Formalism.MainNotions.ColumnSystems.Pf.sdXi}
\leanref{Formalism/SharpRestriction/GrowthBounds/BinomialRatio/CountDefs.lean}{Formalism.MainNotions.ColumnSystems.Pf.chainCount}

\begin{lemma}\label{pd-kappa-defects-digits-block-rank}
Let $p$, $\mathcal O\subset\mathcal K$, $v$, $r\geq1$, a segment datum $S$, integers
$\alpha_1,\dots,\alpha_s,\alpha_+\geq0$ and a real number $\mu$ be as in
Definition~\ref{pd-kappa-defects-segment-lattice}, and let $b$ be a block. For a level $\ell$ of $b$ let
$Y_b(\ell)$ be the $\mathcal K$-space of the polynomials in the letters, of total degree $\delta'(\lambda)$ in the
variables of each letter $\lambda$, that lie in $Y(\mathcal L_l,M_l)$ for every $l$ and in $Y(\mathcal L_+,M_+)$.
Then, with $V_b$ as in Definition~\ref{pd-kappa-defects-windows-blocks},
\[
 \dim_{\mathcal K}V_b=\sum_\ell\dim_{\mathcal K}Y_b(\ell)\binom{b_+-j_++r-1}{r-1},
\]
the sum over the levels $\ell$ of $b$, with $j_+$ the last entry of $\ell$.
\end{lemma}
\leanref{Formalism/SharpRestriction/GrowthBounds/BinomialRatio/CountDefs.lean}{Formalism.MainNotions.ColumnSystems.Pf.BlockRankAt}
\leanref{Formalism/SharpRestriction/GrowthBounds/BinomialRatio/BlockInj.lean}{Formalism.MainNotions.ColumnSystems.Pf.blockRankAt}

\begin{proof}
The levels of $b$ form a finite set, as their entries are bounded by degrees of $b$. For a level
$\ell=(j,j_+)$ put
\[
 \pi_\ell=\prod_l\hat E_l^{j_l}\cdot\hat E_+^{j_+}\cdot\prod_lz_l^{\hat b_{u_l}\dot-o_{u_l}}\cdot
 z_+^{\alpha_+\dot-o_+},
\]
and let $B_\ell$ be the set of the $\beta\in\mathbf Z^r_{\geq0}$ with $|\beta|=b_+-j_+$; it has
$\binom{b_+-j_++r-1}{r-1}$ elements. Write $\xi^\beta=\prod_c\xi_c^{\beta_c}$. Each $Y_b(\ell)$ is
finite-dimensional, as a subspace of the polynomials of given degrees in the letters. Let $\Phi$ be the
$\mathcal K$-linear map from $\bigoplus_{\ell,\beta\in B_\ell}Y_b(\ell)$ to $\mathcal P$ sending
$(y_{\ell,\beta})$ to $\sum_{\ell,\beta}\pi_\ell\,y_{\ell,\beta}\,\xi^\beta$. Its domain has dimension
$\sum_\ell\dim Y_b(\ell)\binom{b_+-j_++r-1}{r-1}$. So it suffices to show that the image of $\Phi$ is
$V_b$ and that $\Phi$ is injective.

\emph{The image.} A generator of level $\ell$ is $\pi_\ell fg$, with $f$ and $g$ as in
Definition~\ref{pd-kappa-defects-segment-lattice}(5). Then $f\in Y_b(\ell)$, and $g=\sum_\beta
g_\beta\xi^\beta$ with $\beta\in B_\ell$ and $g_\beta\in\mathcal O$. So $\pi_\ell fg=\sum_\beta\pi_\ell(g_\beta
f)\xi^\beta$ lies in the image of $\Phi$, and hence so does $V_b$. Conversely, let $y\in Y_b(\ell)$ and
$\beta\in B_\ell$. As $y$ has finitely many coefficients, there is a nonzero $c\in\mathcal O$ such that $cy$
has coefficients in $\mathcal O$. Then $cy\in Y_b(\ell)$, as $Y_b(\ell)$ is a $\mathcal K$-space, and
$\xi^\beta$ is homogeneous of degree $b_+-j_+$ with coefficient $1$. So $\pi_\ell(cy)\xi^\beta$ is a
generator, and $\pi_\ell y\xi^\beta=c^{-1}\pi_\ell(cy)\xi^\beta$ lies in $V_b$.

\emph{Injectivity.} Put $q_l=\sum_cx_{u_l,c}y_{u_l,c}$ and $q_+=\sum_cx_{+,c}\xi_c$. Let $\tau$ be the
endomorphism of the $\mathcal K$-algebra $\mathcal P$ with $\tau(z'_l)=z_lz'_l+q_l=\hat E_l$ and
$\tau(z'_+)=z_+z'_++q_+=\hat E_+$ that fixes the other variables. It is injective: the homomorphism $\theta$
from $\mathcal P$ to its fraction field with $\theta(z'_l)=(z'_l-q_l)/z_l$ and $\theta(z'_+)=(z'_+-q_+)/z_+$ that
fixes the other variables also fixes $q_l$ and $q_+$, so $\theta(\tau(w))=w$ for every variable $w$, and
$\theta\circ\tau$ is the inclusion of $\mathcal P$ in its fraction field. As $\tau$ fixes the variables of the letters, the $\xi_c$
and the $z_l$ and $z_+$,
\[
 \pi_\ell\,y\,\xi^\beta=\tau\big(\sigma_\ell\,y\,\xi^\beta\big),\qquad
 \sigma_\ell=\prod_lz_l'^{\,j_l}\cdot z_+'^{\,j_+}\cdot\prod_lz_l^{\hat b_{u_l}\dot-o_{u_l}}\cdot z_+^{\alpha_+\dot-o_+}.
\]
Now let $\Phi((y_{\ell,\beta}))=0$. As $\tau$ is injective, $\sum_{\ell,\beta}\sigma_\ell\,y_{\ell,\beta}\,\xi^\beta=0$.
Let $\mathcal P_0\subset\mathcal P$ be the ring of the polynomials in the letters, and let $\psi$ be the
homomorphism of $\mathcal K$-algebras from $\mathcal P$ to the polynomial ring over $\mathcal P_0$ in
variables $S_l$, $S_+$ and $\Xi_c$ that sends each variable of a letter to itself in the coefficient ring, $\xi_c$ to $\Xi_c$, $z'_l$ to $S_l$, $z'_+$ to $S_+$, and $z_l$ and $z_+$ to $1$. It sends
$\sigma_\ell y\xi^\beta$ to the monomial $\prod_lS_l^{j_l}S_+^{j_+}\prod_c\Xi_c^{\beta_c}$ with coefficient $y$, and
distinct pairs $(\ell,\beta)$ give distinct monomials. So the coefficients $y_{\ell,\beta}$ of the image of the
sum are all $0$, and $\Phi$ is injective.
\end{proof}
% Lean: blockRankAt = blockRankAt_of restInj, for every BlockData D; the book states it for the block data
%   blockData p S aTop aPlus m b of a segment. Yjt (KK O) (D.lettersAt j jp) D.T.J D.T.thr is Y_b(l), and
%   sdXi r (D.B - jp) is the binomial coefficient. The map is Ltot (sum of ell x beta = mulLeft pre, mulRight
%   xiMon, rename chainEmb) over IF (the levels Lv and antidiagonalTuple r (B - jp)); finrank_dom with
%   card_xi; span_eq_range (mem_gens_span_le_range_of_levelOK, IF_range_le_span with
%   exist_integer_multiples_of_finset and xiMon_mem); ltot_inj through tau (tauF), rho (rhoF, rho_tau,
%   tau_inj), ell_eq_tau with base, and ppi (piF, ppi_base, eExp_inj); then finrank_range_of_inj.

\begin{lemma}\label{pd-kappa-defects-digits-weyl}
Let $N\geq0$ and $n\geq1$ be integers, fix a lettering of $\{0,\dots,N-1\}$ and
$B\colon\{0,\dots,N-1\}\to\mathbf Z_{\geq0}$, and let $\mu$ be a path of rank $n$ for $B$
(Definition~\ref{pd-symbols-path-generators}). Then $\mu$ has $\operatorname{wd}(\mu_N)$ standard fillings
(Definition~\ref{pd-kappa-defects-digits-chains}(1)).
\end{lemma}
\leanref{Formalism/SharpRestriction/GrowthBounds/BinomialRatio/CountDefs.lean}{Formalism.MainNotions.ColumnSystems.Pf.StWeylAt}
\leanref{Formalism/SharpRestriction/GrowthBounds/BinomialRatio/StWeyl.lean}{Formalism.MainNotions.ColumnSystems.Pf.stWeylAt}
\leanref{Formalism/SharpRestriction/GrowthBounds/JointsLit.lean}{Formalism.MainNotions.G.Joints.Lit.weyl_dimension_formula_index_sets}

\begin{proof}
Put $\lambda=\mu_N-\mu_N(n)(1,\dots,1)$, a partition with at most $n$ parts. By fact (i) in the proof of
Lemma~\ref{pd-symbols-path-count}, the standard fillings of $\mu$ correspond to the semistandard tableaux
of shape $\lambda$ with entries in $\{1,\dots,n\}$. Their number is the value at $(1,\dots,1)$ of the
Schur polynomial $s_\lambda$ in $n$ variables, which is the sum of the monomials of these tableaux
\cite[Chapter~I, (5.12)]{Mac95}, and that value is
$\prod_{1\leq a<b\leq n}(\lambda_a-\lambda_b+b-a)/(b-a)$ \cite[Theorem~6.3]{FH91}. This product is
$\operatorname{wd}(\lambda)=\operatorname{wd}(\mu_N)$, as $\operatorname{wd}$ depends only on the differences
of the entries.
\end{proof}
% Lean: stWeylAt D = Joints.Lit.weyl_dimension_formula_index_sets D (admB3) at pEnd mu, dominant by
%   P.dom: restStd identifies IsSt mu with IsStd (pEnd mu), stdCountGT counts these by Gelfand-Tsetlin
%   patterns (dGT), and weylQ proves the product formula for dGT by branching (weylGT). The Lean proves
%   the count combinatorially; the book cites it.

\begin{lemma}\label{pd-kappa-defects-digits-chain-dimension}
Let $p$ be a prime and $\mathcal R'$ a complete discrete valuation ring in whose maximal ideal $p$ is a nonzero
element, with fraction field $K$. Let $N\geq0$ and $r\geq1$ be integers, fix a lettering of $\{0,\dots,N-1\}$
and $B\colon\{0,\dots,N-1\}\to\mathbf Z_{\geq0}$, and let $q'\geq0$, $J_1,\dots,J_{q'}\in\{0,\dots,N\}$ and
$M_1,\dots,M_{q'}\in\mathbf Z\cup\{+\infty\}$. Then the $K$-space of the $F\in K[z]$ of row degrees $B$ that lie in
$Y(\mathcal L_k,M_k)$ for every $k$, with $\mathcal L_k$ the first $J_k$ letters (as in condition (8) of
Definition~\ref{pd-kappa-defects-segment-mean-conditions}), has dimension $\chi_B(J,M)$, the chain sum of
rank $r$ (Definition~\ref{pd-kappa-defects-digits-chains}).
\end{lemma}
\leanref{Formalism/SharpRestriction/GrowthBounds/BinomialRatio/CountDefs.lean}{Formalism.MainNotions.ColumnSystems.Pf.DimYAt}
\leanref{Formalism/SharpRestriction/GrowthBounds/BinomialRatio/CardEqChainSum.lean}{Formalism.MainNotions.ColumnSystems.Pf.dimY_of}

\begin{proof}
For a ring $R$ let $\mathrm{Pol}_R(B)\subset R[z]$ be the polynomials of row degrees $B$, and let $X$ be the
set of the pairs $(\mu,I)$ of a path $\mu$ of rank $r$ for $B$ and a standard filling $I$ of $\mu$. Write
$\Psi_{\mathcal R'}(\mu,I)$ and $\Psi_K(\mu,I)$ for the path generator over $\mathcal R'$ and over $K$. The
column functions are given by the same formulas over both rings, so $\Psi_K(\mu,I)$ is the image of
$\Psi_{\mathcal R'}(\mu,I)$ under the inclusion $\mathcal R'[z]\subset K[z]$, which we write $G\mapsto\bar G$.
The field $K$ has characteristic $0$: if it had characteristic $l>0$, then $l\neq p$ as $p\neq0$ in
$\mathcal R'$, so $p$ would be a unit of $\mathbf F_l\subset\mathcal R'$.

\emph{Step 1: the basis.} By Lemma~\ref{pd-symbols-path-basis} with $R=\mathcal R'$, the
$\Psi_{\mathcal R'}(\mu,I)$ with $(\mu,I)\in X$ are linearly independent over $\mathcal R'$ and span
$\mathrm{Pol}_{\mathcal R'}(B)$. Clearing denominators, the $\Psi_K(\mu,I)$ are linearly independent over $K$.
They lie in the finite-dimensional space $\mathrm{Pol}_K(B)$, so $X$ is finite, and there is an integer
$d_0$ with $|\mu_i(a)|\leq d_0$ for all $(\mu,I)\in X$, $0\leq i\leq N$ and $1\leq a\leq r$.

\emph{Step 2: one condition.} Fix $k$. We show that for $G\in\mathrm{Pol}_{\mathcal R'}(B)$, $\bar G$ lies in
$Y(\mathcal L_k,M_k)$ if and only if $G$ lies in the $\mathcal R'$-span of the $\Psi_{\mathcal R'}(\mu,I)$
with $(\mu,I)\in X$ and $\mu_{J_k}(1)\leq M_k$. If $M_k=+\infty$, both conditions hold for every $G$, the
second by Step 1. If $J_k=0$, then $\mathcal L_k$ is empty, the operators $E^{\mathcal L_k}_{cd}$ are $0$ and
every monomial has weight $0$ for $\mathcal L_k$; so $Y(\mathcal L_k,M_k)$ is $K[z]$ if $M_k\geq0$ and $0$
otherwise. As $\mu_0=0$, the condition $\mu_0(1)\leq M_k$ holds for all pairs if $M_k\geq0$ and for none
otherwise, and $\bar G=0$ only for $G=0$; so the two conditions agree.

Let now $J_k=i+1\geq1$ and $M_k\in\mathbf Z$, so that $\mathcal L_k=\{0,\dots,i\}$. By
Lemma~\ref{pd-kappa-defects-paths-isotypic}(3), $\bar G\in Y(\mathcal L_k,M_k)$ if and only if
$\lambda_{\{0,\dots,i\}}(\bar G)\leq M_k$. The inclusion $\mathcal R'[z]\subset K[z]$ is injective and commutes
with $P\mapsto P_{S,1}$, so $\lambda_{\{0,\dots,i\}}(\bar G)=\lambda_{\{0,\dots,i\}}(G)$. We apply
Lemma~\ref{pd-symbols-path-clauses}(1) over $\mathcal R'$. For $\theta\equiv-d_0$ and $M\equiv d_0$ every
$(\mu,I)\in X$ satisfies its conditions on the path, so by Step 1 every $G\in\mathrm{Pol}_{\mathcal R'}(B)$ has
$-d_0\leq\operatorname{val}_{\{0,\dots,i'\}}(G)$ and $\lambda_{\{0,\dots,i'\}}(G)\leq d_0$ for every letter $i'$.
Let $M'$ be $M_k$ at $i$ and $d_0$ at the other letters. If $\lambda_{\{0,\dots,i\}}(G)\leq M_k$, then $G$
satisfies the conditions of that lemma for $\theta\equiv-d_0$ and $M'$, so it lies in the span of the
$\Psi_{\mathcal R'}(\mu,I)$ with $\mu_{i+1}(1)\leq M_k$. Conversely, these $\Psi_{\mathcal R'}(\mu,I)$ satisfy
the conditions on the path for $\theta\equiv-d_0$ and $M'$, so by the lemma every $G$ in their span has
$\lambda_{\{0,\dots,i\}}(G)\leq M_k$.

\emph{Step 3: the space.} Let $X'$ be the set of the $(\mu,I)\in X$ with $\mu_{J_k}(1)\leq M_k$ for every $k$,
and $Y$ the space of the lemma. We show that $Y$ is the $K$-span of the $\Psi_K(\mu,I)$ with
$(\mu,I)\in X'$. Each of these lies in $\mathrm{Pol}_K(B)$ and, by Step 2, in every $Y(\mathcal L_k,M_k)$.
Conversely, let $F\in Y$. There is a nonzero $d\in\mathcal R'$ with $dF=\bar G$ for some $G\in\mathcal R'[z]$,
and then $G\in\mathrm{Pol}_{\mathcal R'}(B)$ and $\bar G\in Y(\mathcal L_k,M_k)$ for every $k$. By Step 2, for
each $k$, $G$ lies in the span of the $\Psi_{\mathcal R'}(\mu,I)$ with $\mu_{J_k}(1)\leq M_k$. By Step 1 the
coefficients of $G$ in the $\Psi_{\mathcal R'}(\mu,I)$ are unique, so they vanish outside $X'$. Hence $dF$,
and with it $F$, lies in the $K$-span of the $\Psi_K(\mu,I)$ with $(\mu,I)\in X'$.

\emph{Step 4: the count.} By Steps 1 and 3, $\dim_KY$ is the number of elements of $X'$. Sorting them by
$\mu$, this is the sum, over the paths $\mu$ of rank $r$ for $B$ with $\mu_{J_k}(1)\leq M_k$ for every $k$,
of the number of standard fillings of $\mu$. By Lemma~\ref{pd-kappa-defects-digits-weyl} that number is
$\operatorname{wd}(\mu_N)$, so the sum is $\chi_B(J,M)$.
\end{proof}
% Lean: dimYAt_of_adm hA = dimY_of hA2 hA3 (stWeylAt D), with hA2 = HyperplaneRank (the case R = O of
%   Lemma pd-symbols-path-basis, admissible_polynomial_basis) and hA3 = DefiniteRankOne (Lemma
%   pd-symbols-path-clauses(1)), clauses adm15 and adm6 of ColumnSystems.admitted. The characteristic is
%   dvr_fraction_field_charZero. Step 1 is psiK_li, psiK_mem (Psi_map), tx_finite and exists_bound.
%   Step 2 is one_cond: the case top (ycond_top), the case J k = 0 (ycond_zero_top, ycond_zero_bot,
%   chWt_zero), and the case J k = i + 1 with mem_ycond_iff (Lemma pd-kappa-defects-paths-isotypic(3)),
%   lowest_degree_map_of_injective and single_cond (A3at with th = -Bd and M = Bd, then M updated at i).
%   Step 3 is span_eq_Yjt (clear_den, map_mem_Pol_iff, mem_span_iInter, map_span). Step 4 is
%   finrank_eq_nat_card_basis and card_eq_chainSum (pathOf, chain_of_path, fiber_card, StWeylAt).

\begin{lemma}\label{pd-kappa-defects-digits-count}
Let $p\geq5$ be a prime and $n'\geq2$ an even integer, let the integers, reals, constants, $c\geq0$, $a_T$ and
$b_T$ of Definition~\ref{pd-kappa-defects-numbers} be given, and form the numbers there. Let an unramified
quadratic Hida setting of degree $n'$ at $p$ with ring $\mathcal R$ be given, an index $1\leq j\leq q$, a real
number $\mu\geq0$ and digits $\rho,\rho'$ for $j$, and form $r_j(\mu;\rho,\rho')$ with $a_T$ and $b_T$ as in
Definition~\ref{pd-kappa-defects-segment-mean-conditions}. Then $r_j(\mu;\rho,\rho')=c_j(\mu;\rho,\rho')$
(Definition~\ref{pd-kappa-defects-digits-chains}).
\end{lemma}
\leanref{Formalism/SharpRestriction/GrowthBounds/BinomialRatio/RankOnDigitsCast.lean}{Formalism.MainNotions.ColumnSystems.Pf.rankOnDigits_cast}
\leanref{Formalism/SharpRestriction/GrowthBounds/BinomialRatio/RankOnDigitsCast.lean}{Formalism.MainNotions.ColumnSystems.Pf.blockRank_cast}

\begin{proof}
Both sides are sums over the same kept blocks $b$, of $\dim V_b$ and of $c(b)$. By
Lemma~\ref{pd-kappa-defects-digits-block-rank}, $\dim V_b$ is the sum over the levels $\ell$ of $b$ of
$\dim Y_b(\ell)$ times the binomial coefficient. The space $Y_b(\ell)$ is the space of
Lemma~\ref{pd-kappa-defects-digits-chain-dimension} for the letters of $b$ with the degrees $\delta'$ of $\ell$, and
the conditions $(\mathcal L_l,M_l)$ and $(\mathcal L_+,M_+)$, whose sets of letters are initial. The ring
$\mathcal R$ is a complete discrete valuation ring with $p=\varpi_0^2$ a nonzero element of its maximal ideal
(Definition~\ref{pd-local-model}(1)), so that lemma gives $\dim Y_b(\ell)=\chi_b(\ell)$. Hence
$\dim V_b=c(b)$.
\end{proof}

\begin{lemma}\label{pd-kappa-defects-digits-branching}
Let $r\geq0$ and $\nu\in\mathbf Z^{r+1}$ with $\nu_1\geq\cdots\geq\nu_{r+1}$. Then
$\operatorname{wd}(\nu)=\sum_\lambda\operatorname{wd}(\lambda)$, the sum over the $\lambda\in\mathbf Z^r$ with
$\nu_{a+1}\leq\lambda_a\leq\nu_a$ for $1\leq a\leq r$ (Definition~\ref{pd-kappa-defects-digits-chains}(1)).
\end{lemma}
\leanref{Formalism/SharpRestriction/GrowthBounds/BinomialRatio/CountDefs.lean}{Formalism.MainNotions.ColumnSystems.Pf.NumBranchAt}
\leanref{Formalism/SharpRestriction/GrowthBounds/BinomialRatio/StWeyl.lean}{Formalism.MainNotions.ColumnSystems.Pf.numBranch_of}

\begin{proof}
Adding $-\nu_{r+1}$ to all entries of $\nu$ and of the $\lambda$ changes no Weyl dimension and no condition,
so we may assume $\nu_{r+1}=0$; then $\nu$ and the $\lambda$ are partitions. The $\lambda$ are those for which
$\nu/\lambda$ is a horizontal strip with $\lambda$ of at most $r$ parts, and the branching formula
$s_\nu(x_1,\dots,x_{r+1})=\sum_\lambda s_\lambda(x_1,\dots,x_r)x_{r+1}^{|\nu|-|\lambda|}$
\cite[Chapter~I, (5.11)]{Mac95}, evaluated at $(1,\dots,1)$, gives the claim, as
$s_\nu(1,\dots,1)=\operatorname{wd}(\nu)$ \cite[Theorem~6.3]{FH91}.
\end{proof}
% Lean: numBranch_of (adm2, adm3 of ColumnSystems.admitted): after the shift by mu (last) (wdim_shift), adm3
%   identifies wdim with the dimension of the dual Weyl module over Q and adm2 is its branching to GL_r;
%   the sum is reindexed by the shift (antitone_of_interlace). The book cites the branching rule.

\begin{lemma}\label{pd-kappa-defects-digits-top-rows}
Let $C_c\geq0$, $C_0\geq1$, $C_{E1},C_{E2}\geq1$ and $C_H\geq1000(C_{E1}+C_{E2}+1)(C_c+1)$. Let $p\geq5$
be a prime and $n'\geq2$ an even integer, and let $m$, $G'$, $N_1$, $c_{cl}$, $D_p$, $c_*$, $T$, $Y$,
$Y_b$, $b'$, $b_l$, $b_u$ and an integer $a$ satisfy hypotheses (1) to (3) of
Theorem~\ref{pd-gain-theorem-one} with $C_0$ and $C_H$. Let $c\geq0$ be an integer with
$c\leq C_c(c_*+|T|k\log_pa+|T|c_{cl}+n'c_{cl})$, and let $a_T,b_T\colon\mathbf Z/n'\to\mathbf Z_{\geq0}$
satisfy $a\leq a_T(s)\leq a+a/(kY)+2J_0$ and $b_l\leq b_T(s)\leq b_u+2J_0$ for $s\in T$, where
$J_0=\Lambda b'/k$. Form the numbers of Definition~\ref{pd-kappa-defects-numbers} with these
constants. Let $1\leq j\leq q$, and let $b$ be a block kept for the segment datum $S_j(b_T)$, with
$\alpha_i=a_T(y_{j,i})$ and $r=m$ (Definition~\ref{pd-kappa-defects-segment-lattice}). Then for
$1\leq i\leq\eta_j$,
\[
 \sum_{l=1}^iw_l\leq\operatorname{cap}(i-1)\quad\text{and}\quad\sum_{l=1}^iw_l\leq\operatorname{cap}(i).
\]
\end{lemma}
\leanref{Formalism/SharpRestriction/GrowthBounds/BinomialRatio/UpperCountOfTopRows.lean}{Formalism.MainNotions.ColumnSystems.Pf.SecondKernelAtLevel}
\leanref{Formalism/SharpRestriction/GrowthBounds/BinomialRatio/UpperCountOfTopRows.lean}{Formalism.MainNotions.ColumnSystems.Pf.LeadingRowsBeyondThreshold}
\leanref{Formalism/SharpRestriction/GrowthBounds/BinomialRatio/AiAddMulLeA0.lean}{Formalism.MainNotions.ColumnSystems.Pf.top_rows_bound_holds}

\begin{proof}
Write $s=S_j$, $\eta=\eta_j$, $\hat a=\hat a_j$, $\varepsilon=N_1\Lambda/k$ and
$\beta_e=(b_u-b_l)/b'+2\Lambda/k$, and let $\delta$ be the number of Theorem~\ref{pd-gain-theorem-one}. For
$1\leq l\leq s$ put $B_l=b_T(y_{j,l})$, $P_l=p^{g_j-d_{j,l}}$ and $\hat b_l=\lfloor(B_l-|\rho_l|)/P_l\rfloor$, with
$\rho_l$ the vector of $b$.

\emph{Facts.} (1) The facts (a) to (g) in the proof of Lemma~\ref{pd-kappa-defects-margin} hold, and so do
the facts (h) and (i) and Step~1 in the proof of Lemma~\ref{pd-kappa-defects-windows-model-step}, which use
only the hypotheses here. They give $k\geq5$, $\Lambda\geq1$, $N_1\geq1$, $b'>0$, $k^2\leq k^2Y^2\leq p$,
$b'\leq a$, $k^3Q\leq b'$, $E_0+Q\leq\varepsilon a$, $40\varepsilon\leq\delta<1/6$ and $\hat a\geq a/(2Q)$. So
$2kQ\leq k^3Q\leq b'\leq a\leq2Q\hat a$, and $k\leq\hat a$.

(2) $0\leq\beta_e\leq\delta/2-18\varepsilon$, and $b'(1-\beta_e)\leq b_T(t)\leq b'(1+\beta_e)$ for $t\in T$. Indeed
$\delta/2-18\varepsilon-\beta_e=(C_H/2-18)N_1\Lambda/k-2\Lambda/k\geq0$ as $C_H\geq40$. As $b_l\leq b'\leq b_u$,
$b_T(t)\geq b_l\geq b'-(b_u-b_l)$ and $b_T(t)\leq b_u+2\Lambda b'/k\leq b'+(b_u-b_l)+2\Lambda b'/k$. In particular
$\beta_e\leq1/12$, as $\delta<1/6$.

(3) $\hat b_lP_l\leq B_l$ and $B_l-kP_l<\hat b_lP_l$. Indeed the entries of $\rho_l$ are less than $P_l$, so
$|\rho_l|<(m+1)P_l=kP_l$, and $P_l$ divides $B_l-|\rho_l|\geq0$, so $\hat b_lP_l=B_l-|\rho_l|$.

(4) $\hat b_l(1+\delta)<\hat a$ for $l\leq\eta$. The number $d_{j,l}$ is odd, as it is the cyclic distance from the
even $s_j$ to the odd $y_{j,l}$ and $n'$ is even. As the $d_{j,\cdot}$ increase, the $l$ with
$p^{d_{j,l}}<a/b'$ are $1,\dots,\eta$, so $p^{d_{j,l}}<a/b'$. As no odd power of $p$ lies in
$[(a/b')(1-2\delta),(a/b')(1+2\delta)]$ by hypothesis (2), $p^{d_{j,l}}<(a/b')(1-2\delta)$. Now
$P_lp^{d_{j,l}}=p^{g_j}$, and $B_l\leq b'(1+\beta_e)$ by (2), so by (3)
\[
 \hat b_lp^{g_j}\leq B_lp^{d_{j,l}}<a(1-2\delta)(1+\beta_e).
\]
By (1) and (2), $(1-2\delta)(1+\beta_e)(1+\delta)\leq(1+\delta/2)(1-2\delta)(1+\delta)=1-\delta/2-5\delta^2/2-\delta^3\leq
1-\varepsilon$, so $\hat b_l(1+\delta)p^{g_j}<a(1-\varepsilon)$. On the other hand, by the definition of the floor,
$g_j\leq G'$ and (1), $\hat ap^{g_j}>a-E_0-p^{g_j}\geq a-E_0-Q\geq a(1-\varepsilon)$.

(5) $\hat b_l\leq2(\hat b_i+k)/p$ for $l<i\leq\eta$. As $d_{j,l}<d_{j,i}$, $pP_i\leq P_l$. By (3) and (2),
$\hat b_lpP_i\leq\hat b_lP_l\leq B_l\leq b'(1+\beta_e)$ and $b'(1-\beta_e)\leq B_i<(\hat b_i+k)P_i$; and
$1+\beta_e\leq2(1-\beta_e)$ as $\beta_e\leq1/12$.

\emph{The sum.} Let $1\leq i\leq\eta$, and put $c_1=2(\hat b_i+k)/p$ and $c_2=2\Lambda\hat a/k$. As $b$ is kept and
$r+1=k$, $\varsigma_l\leq c_2$ for $l\leq\eta$; with (5), $\hat b_l+\varsigma_l\leq c_1+c_2$ for $l<i$, and
$\hat b_i+\varsigma_i\leq\hat b_i+c_2$. As $i\leq S_j\leq N_1$ by hypothesis (1),
\[
 \sum_{l=1}^i(\hat b_l+\varsigma_l)\leq\hat b_i+N_1(c_1+c_2).
\]
By (4) and (1), $\hat b_i\leq\hat a$ and $k\leq\hat a$, so $c_1\leq4\hat a/k^2\leq\Lambda\hat a/k$, as $p\geq k^2$ and
$4\leq\Lambda k$. So $N_1(c_1+c_2)\leq3\varepsilon\hat a$. By (4), $(\hat a-\hat b_i)(1+\delta)>\delta\hat a\geq40\varepsilon\hat a$, and
$1+\delta\leq7/6$, so $\hat a-\hat b_i\geq34\varepsilon\hat a$. Hence $\sum_{l=1}^i(\hat b_l+\varsigma_l)\leq\hat a$.

\emph{The bound.} Let $0\leq n\leq i$. For $l\leq i\leq\eta$, $o^\flat_l=w_l+\varsigma_l$ and $w_l\leq\hat b_l$, as $b$ is kept.
So $w_l\leq(\hat b_l-o^\flat_l)+(\hat b_l+\varsigma_l)$ for $l\leq n$, and $w_l\leq\hat b_l+\varsigma_l$ for $n<l\leq i$.
Summing and using the sum above,
\[
 \sum_{l=1}^iw_l\leq\sum_{l=1}^n(\hat b_l-o^\flat_l)+\hat a=\operatorname{cap}(n).
\]
The cases $n=i-1$ and $n=i$ are the claims.
\end{proof}
% Lean: top_rows_bound_holds = TransportedAlt.of_explicit with thresholds 1, 1 and T2 CE1 CE2 Cc 1, through
%   std_of and partial_width_le_cap at n = i and n = i + 1 (0-based). Facts (1) and (2) are fields of Std:
%   k5, La1, N1pos, bpos, pk, ba, kQ, N5, LaK, de6, half_ratio_le_segment_base_of_std, and be, be0, beLe,
%   winB (be X C0 = (bh - bl)/b' + 2 La/k, with the window bT <= bh + 2 J0). (3) is coarse_bounds. (4) is
%   coarse_degree_alternative (first part): dist_odd, sigma_iff, h.band, dich_S with poly_S. (5) is
%   coarse_ratio_small with dist_strictMono and levels. The sum is hterm and hsumR with Ai_add_mul_le_A0
%   (c1, c2, h1 to h4); the bound is hcmp, summed.

\begin{lemma}\label{pd-kappa-defects-digits-count-upper}
For every $C_c\geq0$ there is $T_0>0$ such that for every $C_0\geq T_0$ there is $T_1>0$ such that for
all $C_{E1},C_{E2},C_m\geq T_1$ there is $T_2>0$ with the following property for every $C_H\geq T_2$.
Let $p\geq5$ be a prime and $n'\geq2$ an even integer, and let $m$, $G'$, $N_1$, $c_{cl}$, $D_p$, $c_*$,
$T$, $Y$, $Y_b$, $b'$, $b_l$, $b_u$ and an integer $a$ satisfy hypotheses (1) to (3) of
Theorem~\ref{pd-gain-theorem-one} with $C_0$ and $C_H$. Let $c\geq0$ be an integer with
$c\leq C_c(c_*+|T|k\log_pa+|T|c_{cl}+n'c_{cl})$, and let $a_T,b_T\colon\mathbf Z/n'\to\mathbf Z_{\geq0}$
satisfy $a\leq a_T(s)\leq a+a/(kY)+2J_0$ and $b_l\leq b_T(s)\leq b_u+2J_0$ for $s\in T$, where
$J_0=\Lambda b'/k$. Form the numbers of Definition~\ref{pd-kappa-defects-numbers} with these
constants; the segment data do not involve $C_m$, $C_A$, $C_t$ or $C_\kappa$. Let $1\leq j\leq q$, $\mu\geq0$
and digits $\rho,\rho'$ for $j$ be given. Then $c_j(\mu;\rho,\rho')\leq d_j(\rho,\rho')$
(Definitions~\ref{pd-kappa-defects-digits-chains} and~\ref{pd-kappa-defects-segment-mean-conditions}).
\end{lemma}
\leanref{Formalism/SharpRestriction/GrowthBounds/BinomialRatio/CountDefs.lean}{Formalism.MainNotions.ColumnSystems.Pf.CountUpAt}
\leanref{Formalism/SharpRestriction/GrowthBounds/BinomialRatio/CountDefs.lean}{Formalism.MainNotions.ColumnSystems.Pf.RestCountUp}

\begin{proof}
\emph{The constants.} Let $T_0=T_1=1$ and $T_2=1000(C_{E1}+C_{E2}+1)(C_c+1)$, and let the data be as in the
statement. Write $s=S_j$, $\eta=\eta_j$, $\hat a=\hat a_j$, and for $1\leq i\leq s$ let $\hat b_i$ and $\hat\alpha_i$ be
the numbers of Definition~\ref{pd-kappa-defects-segment-lattice}(1) for the digits $\rho,\rho'$; they do not depend
on the rest of a block. Number the letters $x_1,y_1,\dots,x_s,y_s,x_+$ of the segment $0,\dots,2s$, and for a vector
$v\in\mathbf Z^m$ and $e\in\mathbf Z$ let $(e,v)\in\mathbf Z^{m+1}$ be $v$ preceded by $e$. The constants satisfy the
hypotheses of Lemmas~\ref{pd-kappa-defects-windows-model-step} and~\ref{pd-kappa-defects-digits-top-rows}, and
Step~1 in the proof of the first gives $\hat a\geq a/(2Q)>0$.

\emph{The sources.} Let $Z$ be the set of the triples $z=(b,\ell,\kappa)$ of a kept block $b$ of $S_j(b_T)$ with
vectors $\rho,\rho'$, a level $\ell=(j_1,\dots,j_{s-\eta},j_+)$ of $b$, and a path $\kappa=(\kappa_0,\dots,\kappa_{2s+1})$ of
rank $m$ for the degrees $\delta'$ of $\ell$ with $\kappa_{J}(1)\leq M$ for each of the pairs $(J,M)$ of
Definition~\ref{pd-kappa-defects-digits-chains}(3). By that definition,
\[
 c_j(\mu;\rho,\rho')=\sum_{z\in Z}\operatorname{wd}(\kappa_{2s+1})\binom{b_+-j_++m-1}{m-1}.
\]

\emph{The targets.} Give the letters $0,\dots,2s$ new types and degrees: $0$ is a vector letter of degree $\hat a$,
$2i-1$ a vector letter of degree $\hat b_i$, and $2i$ a covector letter of degree $\hat\alpha_i$, for $1\leq i\leq s$. Let
$\Theta$ be the set of the triples $(\nu,\lambda,j_+)$ of a path $\nu$ of rank $m+1$ for these degrees, a
$\lambda\in\mathbf Z^m$ with $\nu_{2s+1}(a+1)\leq\lambda_a\leq\nu_{2s+1}(a)$ for $1\leq a\leq m$, and an integer
$0\leq j_+\leq b_+$, and give $(\nu,\lambda,j_+)$ the weight $\operatorname{wd}(\lambda)\binom{b_+-j_++m-1}{m-1}>0$. The
sum over $j_+$ of the binomial coefficients is $\binom{b_++m}m=P_m(b_+)$. By
Lemma~\ref{pd-kappa-defects-digits-branching}, the sum over $\lambda$ of $\operatorname{wd}(\lambda)$ is
$\operatorname{wd}(\nu_{2s+1})$. By Lemma~\ref{pd-kappa-defects-digits-chain-dimension} with $q'=0$ (over
$\mathbf Z_5$, say), the sum over $\nu$ of $\operatorname{wd}(\nu_{2s+1})$ is the dimension of the space of
polynomials of these row degrees in $m+1$ variables per letter, which is $P_m(\hat a)\prod_iP_m(\hat b_i)P_m(\hat\alpha_i)$.
So the total weight of $\Theta$ is $d_j(\rho,\rho')$.

\emph{The facts.} Let $z=(b,\ell,\kappa)\in Z$ and $c_l=\operatorname{cap}(l)$ for $b$. Then:
\begin{enumerate}
\item[(F1)] $c_0=\hat a$;
\item[(F2)] $\delta'(x_i)\leq\hat b_i$ and $c_i\leq c_{i-1}+\hat b_i-\delta'(x_i)$;
\item[(F3)] $c_{i-1}+\hat b_i-\delta'(x_i)-c_i+\delta'(y_i)=\hat\alpha_i$;
\item[(F4)] $\kappa_{2i-1}(1)\leq c_{i-1}$ and $\kappa_{2i-1}(1)\leq c_i$, for $1\leq i\leq s$;
\item[(F5)] $\kappa_{2s+1}(1)\leq c_s$.
\end{enumerate}
(F1) is the definition. Recall $c_i=c_{i-1}+\hat b_i-o^\flat_i$. For $i\leq\eta$, $\delta'(x_i)=w_i\leq\hat b_i$,
$\delta'(y_i)=\hat\alpha_i-\varsigma_i$ and $o^\flat_i=w_i+\varsigma_i$, with $\varsigma_i\leq\hat\alpha_i$ as $b$ is kept. For
$i=u_l>\eta$, $\delta'(x_i)=o_i-j_l$, $\delta'(y_i)=\hat\alpha_i-j_l$ and $o^\flat_i=o_i\leq\hat b_i$, with
$j_l\leq\min(o_i,\hat\alpha_i)$. In both cases $\delta'(x_i)\leq o^\flat_i$ and $o^\flat_i-\delta'(x_i)+\delta'(y_i)=\hat\alpha_i$,
which give (F2) and (F3). For (F4) with $i=u_l>\eta$: the set $\mathcal L_l$ consists of the first $2i-1$ letters, so
$\kappa_{2i-1}(1)\leq M_l=\lfloor c_{i-1}-\mu\rfloor\leq c_{i-1}$, and $c_i\geq c_{i-1}$ as $o_i\leq\hat b_i$. For (F4) with
$i\leq\eta$: at a vector letter the first entry of a path grows by at most the degree, as
$\kappa_{l+1}(1)-\kappa_l(1)\leq\sum_a(\kappa_{l+1}(a)-\kappa_l(a))$, and at a covector letter it does not grow. As
$\kappa_0=0$ and the vector letters among the first $2i-1$ are $x_1,\dots,x_i$, with $\delta'(x_l)=w_l$,
$\kappa_{2i-1}(1)\leq\sum_{l\leq i}w_l$, and Lemma~\ref{pd-kappa-defects-digits-top-rows} bounds this by $c_{i-1}$
and by $c_i$. For (F5), $\mathcal L_+$ consists of all the letters, so
$\kappa_{2s+1}(1)\leq M_+=\lfloor c_s-\mu\rfloor\leq c_s$.

\emph{The map.} For $z\in Z$ let $\nu=\nu(z)$ be the sequence with $\nu_0=0$,
$\nu_{2i-1}=(c_{i-1},\kappa_{2i-2})$ and $\nu_{2i}=(c_{i-1}+\hat b_i-\delta'(x_i),\kappa_{2i-1})$ for $1\leq i\leq s$, and
$\nu_{2s+1}=(c_s,\kappa_{2s})$, and put $\Phi(z)=(\nu,\kappa_{2s+1},j_+)$. We check that $\Phi(z)\in\Theta$. At the
letter $0$, $\nu_1=(\hat a,0)$ by (F1), with $\hat a\geq0$. At the letter $2i-1$ compare $\nu_{2i-1}$ and $\nu_{2i}$:
their tails form the strip $(\kappa_{2i-2},\kappa_{2i-1})$ of size $\delta'(x_i)$ at the vector letter $x_i$; the first
entries satisfy $c_{i-1}\leq c_{i-1}+\hat b_i-\delta'(x_i)$ by (F2); the second entry
$\kappa_{2i-1}(1)$ of $\nu_{2i}$ is at most the first entry $c_{i-1}$ of $\nu_{2i-1}$ by (F4); and the sizes add to
$\hat b_i$. At the covector letter $2i$ compare $\nu_{2i+1}$ and $\nu_{2i}$ in the same way: the tails form the strip of
size $\delta'(y_i)$ at the covector letter $y_i$, $c_i\leq c_{i-1}+\hat b_i-\delta'(x_i)$ by (F2),
$\kappa_{2i-1}(1)\leq c_i$ by (F4), and the size is $\hat\alpha_i$ by (F3). The conditions at a letter force both of
its vectors to be non-increasing, so $\nu$ is a path for the new degrees. As
$x_+$ is a vector letter, $\kappa_{2s}(a)\leq\kappa_{2s+1}(a)$ and $\kappa_{2s+1}(a+1)\leq\kappa_{2s}(a)$; with (F5) this
says that $\lambda=\kappa_{2s+1}$ satisfies $\nu_{2s+1}(a+1)\leq\lambda_a\leq\nu_{2s+1}(a)$. Finally
$j_+\leq\min(\delta(x_+),b_+)$, and the weight of $\Phi(z)$ is the term of $z$ in the sum for $c_j(\mu;\rho,\rho')$.

\emph{Injectivity.} Let $\Phi(z)=\Phi(z')$. Then $\kappa=\kappa'$, as $\kappa_l$ is $\nu_{l+1}$ without its first entry for
$l\leq2s$ and $\kappa_{2s+1}=\lambda$, and $j_+=j'_+$. The path $\kappa$ determines the degrees $\delta'$, the sizes of its
strips. With $\hat b_i$ and $\hat\alpha_i$ fixed by the digits, these determine the block and the level: for $i\leq\eta$,
$w_i=\delta'(x_i)$, $\varsigma_i=\hat\alpha_i-\delta'(y_i)$ and $o_i=0$; for $i=u_l>\eta$, $j_l=\hat\alpha_i-\delta'(y_i)$,
$o_i=\delta'(x_i)+j_l$ and $w_i=\varsigma_i=0$; and $o_+=\delta'(x_+)+j_+$. So $z=z'$.

\emph{The bound.} The weights of $\Theta$ are positive and finitely many are nonzero, so the sum over $Z$, which is
the sum of the weights of the distinct targets $\Phi(z)$, is at most the total weight $d_j(\rho,\rho')$ of $\Theta$.
\end{proof}
% Lean: RestCountUp = upper_count_of_top_rows (chainDim_of hY) (numBranch_of ...) top_rows_bound_holds =
%   restCountUp_of_rows through restFirstRows_of (restSTopRows_of h) and restMatchUp_of; the thresholds are
%   those of std_of and of top_rows_bound_holds (1, 1 and T2 CE1 CE2 Cc 1). countUp_of: chainCount_eq
%   (the sources, Fs), Gt and Gt_sum (hockey_fin, NumBranchAt = Lemma pd-kappa-defects-digits-branching,
%   ChainDimAt = chainDim_of, Lemma pd-kappa-defects-digits-chain-dimension at Z_5 with q = 0 and
%   finrank_Pol), prod_kLetters with kLetters (the new types and degrees) and mmB, then finsum_le_of_inj.
%   The facts are LevFacts (levFactsAt_of with cap_zero, levW, levE, cap_succ, level_bounds); (F4) for
%   i > eta and (F5) are thr_U and thr_last (firstRows_of), (F4) for i <= eta is row_odd (strip_first,
%   row_even) with SecondKernelAtLevel (Lemma pd-kappa-defects-digits-top-rows). The map is muOf with
%   fstOf, checked by chain_of_facts (strip_cons); injectivity is muOf_inj and srcInjAt.

\begin{definition}\label{pd-kappa-defects-digits-count-pairs}
Let $p\geq5$ be a prime, $\Lambda$ a real number, and $m\geq1$, $s\geq0$, $A\geq0$ and
$\hat b_1,\dots,\hat b_s,\hat\alpha_1,\dots,\hat\alpha_s\geq0$ integers. Give the letters $0,1,\dots,2s$ types
and degrees: $0$ is a vector letter of degree $A$, and for $1\leq i\leq s$, $2i-1$ is a vector letter of degree
$\hat b_i$ and $2i$ a covector letter of degree $\hat\alpha_i$. For a vector $v$ with integer entries write $|v|$ for the
sum of its entries; for vectors $\alpha,\beta$ of the same length, $\beta$ is a \emph{strip of size $g$ over} $\alpha$ if
$\alpha(a)\leq\beta(a)$ and $\beta(a+1)\leq\alpha(a)$ for all $a$ and $|\beta|-|\alpha|=g$, as in
Definition~\ref{pd-symbols-path-generators}(2).
\begin{enumerate}
\item A \emph{pair} is a pair $(\nu,\lambda)$ of a path $\nu=(\nu_0,\dots,\nu_{2s+1})$ of rank $m+1$ for these degrees
  (Definition~\ref{pd-symbols-path-generators}(2)) and a $\lambda\in\mathbf Z^m$ with
  $\nu_{2s+1}(a+1)\leq\lambda_a\leq\nu_{2s+1}(a)$ for $1\leq a\leq m$. There are finitely many pairs, as each entry of
  $\nu_l$ has absolute value at most the sum of the degrees.
\item For a pair put $c_l=\nu_{2l+1}(1)$ for $0\leq l\leq s$, and $t_i=\nu_{2i}(2)$, $\bar o_i=c_{i-1}+\hat b_i-c_i$ and
  $n_i=-\lambda_{m'+i}$ for $1\leq i\leq s$, where $m'=\max(m-s,0)$ and $\lambda_a=0$ for $a>m$; so
  $n_i=-\lambda_{m-s+i}$ when $s\leq m$. The pair is \emph{kept} if for every $1\leq i\leq s$
  \[
   \bar o_i-t_i+\frac{4(\hat b_i+m)}{p^2}\leq\frac{\Lambda\max(c_{i-1},\hat b_i)}{m+1}
   \quad\text{and}\quad
   \hat\alpha_i\Big(1-\frac\Lambda{m+1}\Big)\leq n_i.
  \]
\item The \emph{kept sum} $K_A$ is the sum of $\operatorname{wd}(\lambda)$
  (Definition~\ref{pd-kappa-defects-digits-chains}(1)) over the kept pairs $(\nu,\lambda)$.
\end{enumerate}
\end{definition}
\leanref{Formalism/SharpRestriction/GrowthBounds/BinomialRatio/CountLow.lean}{Formalism.MainNotions.ColumnSystems.Pf.IsPair}
\leanref{Formalism/SharpRestriction/GrowthBounds/BinomialRatio/CountUp.lean}{Formalism.MainNotions.ColumnSystems.Pf.kLetters}
\leanref{Formalism/SharpRestriction/GrowthBounds/BinomialRatio/CountLow.lean}{Formalism.MainNotions.ColumnSystems.Pf.pairC}
\leanref{Formalism/SharpRestriction/GrowthBounds/BinomialRatio/CountLow.lean}{Formalism.MainNotions.ColumnSystems.Pf.pairRho}
\leanref{Formalism/SharpRestriction/GrowthBounds/BinomialRatio/CountLow.lean}{Formalism.MainNotions.ColumnSystems.Pf.pairO}
\leanref{Formalism/SharpRestriction/GrowthBounds/BinomialRatio/CountLow.lean}{Formalism.MainNotions.ColumnSystems.Pf.pairN}
\leanref{Formalism/SharpRestriction/GrowthBounds/BinomialRatio/CountLow.lean}{Formalism.MainNotions.ColumnSystems.Pf.KeptPair}
\leanref{Formalism/SharpRestriction/GrowthBounds/BinomialRatio/CountLow.lean}{Formalism.MainNotions.ColumnSystems.Pf.keptPairSum}

\begin{lemma}\label{pd-kappa-defects-digits-count-binomial}
Let $m\geq1$ be an integer, and for integers $d\geq0$ put $P_m(d)=\binom{d+m}m$, as in
Lemma~\ref{pd-kappa-defects-windows-model-segment}.
\begin{enumerate}
\item If $\hat a\geq1$ and $0\leq n\leq\hat a$ are integers, then $(1-mn/\hat a)P_m(\hat a)\leq P_m(\hat a-n)$.
\item If $m\geq4$, $B\geq m+1$ is an integer and $\Lambda$ is a real number with $0\leq4\Lambda\leq m+1$, then
  \[
   \big(1-e^{-\Lambda/3}\big)P_m(B)\leq\sum_{j=0}^J\binom{B-j+m-1}{m-1},\qquad
   J=\Big\lfloor\frac{\Lambda B}{m+1}\Big\rfloor .
  \]
\end{enumerate}
\end{lemma}
\leanref{Formalism/SharpRestriction/GrowthBounds/BinomialRatio/CountLow.lean}{Formalism.MainNotions.ColumnSystems.Pf.cl_sd_margin}
\leanref{Formalism/SharpRestriction/GrowthBounds/BinomialRatio/CountLow.lean}{Formalism.MainNotions.ColumnSystems.Pf.cl_hockey_low}

\begin{proof}
For integers $y,l\geq0$, $(y+l+1)P_l(y)=(l+1)P_{l+1}(y)$, as $(y+l+1)\binom{y+l}l=(l+1)\binom{y+l+1}{l+1}$. Hence
$P_l(y)=\prod_{t=1}^l(y+t)/t$.

(1) Put $q=n/\hat a\in[0,1]$. We show $(1-q)^lP_l(\hat a)\leq P_l(\hat a-n)$ by induction on $l$; for $l=0$ both sides
are $1$. As $0\leq(1-q)(\hat a+l+1)=\hat a-n+(1-q)(l+1)\leq\hat a-n+l+1$, the identity and the induction hypothesis
give
\[
 (l+1)(1-q)^{l+1}P_{l+1}(\hat a)=(1-q)(\hat a+l+1)\cdot(1-q)^lP_l(\hat a)\leq(\hat a-n+l+1)P_l(\hat a-n)
 =(l+1)P_{l+1}(\hat a-n).
\]
By Bernoulli's inequality, $1-mq\leq(1-q)^m$, as $-q\geq-2$; this gives (1).

(2) Pascal's rule $P_m(g+1)=P_m(g)+\binom{g+1+m-1}{m-1}$, applied $l$ times, gives
$\sum_{j=0}^{l-1}\binom{B-j+m-1}{m-1}+P_m(B-l)=P_m(B)$ for $0\leq l\leq B$. As $\Lambda B/(m+1)\leq B/4$,
$J+1\leq B/4+1\leq B$. So the sum in the claim is $P_m(B)-P_m(B-J-1)$, and it suffices to show
$P_m(B-J-1)\leq e^{-\Lambda/3}P_m(B)$. Put $\theta=1-(J+1)/(B+m)$, which lies in $[0,1]$. For $1\leq t\leq m$,
$(J+1)(B+t)/(B+m)\leq J+1$, so $B-J-1+t\leq\theta(B+t)$; by the product formula,
$P_m(B-J-1)\leq\theta^mP_m(B)$. As $1-x\leq e^{-x}$, $\theta^m\leq e^{-m(J+1)/(B+m)}$. Finally, $\Lambda B<(J+1)(m+1)$ by
the definition of $J$, and $3mB-(m+1)(B+m)=(B-m-1)(2m-1)+(m+1)(m-1)\geq0$; so
$\Lambda(B+m)(m+1)\leq3m\Lambda B\leq3m(J+1)(m+1)$, that is, $m(J+1)/(B+m)\geq\Lambda/3$.
\end{proof}
% Lean: (1) is cl_sd_margin (one_sub_div_pow_mul_sd_le_sd_sub by induction on m with cl_sd_step, and
%   one_add_mul_le_pow); (2) is cl_hockey_low (cl_hockey, cl_sd_ratio_le with theta, Real.add_one_le_exp). Here
%   P_m(d) is sd m d and the binomial coefficient is sdXi m (B - j).

\begin{lemma}\label{pd-kappa-defects-digits-count-pair-number}
Let $p$, $\Lambda$, $m$, $s$, $A$, the $\hat b_i$ and the $\hat\alpha_i$ be as in
Definition~\ref{pd-kappa-defects-digits-count-pairs}, and let $M\geq0$ be an integer with $\hat b_i\leq M$ and
$\hat\alpha_i\leq M$ for all $i$.
\begin{enumerate}
\item Let $\nu$ be a path of rank $m+1$ for the degrees of that definition, let $0\leq l\leq2s+1$, and let $u_l$ and $v_l$ be
  the numbers of vector and of covector letters among $0,\dots,l-1$. Then $\nu_l(a)\leq0$ for $a>u_l$, $\nu_l(a)\geq0$
  for $a\leq m+1-v_l$, and $|\nu_l(a)|\leq A+lM$.
\item There are at most $\big(2(A+(2s+1)M)+1\big)^{(2s+3)(2s+1)}$ pairs.
\end{enumerate}
\end{lemma}
\leanref{Formalism/SharpRestriction/GrowthBounds/BinomialRatio/PairCount.lean}{Formalism.MainNotions.ColumnSystems.Pf.pc_shape}
\leanref{Formalism/SharpRestriction/GrowthBounds/BinomialRatio/PairCount.lean}{Formalism.MainNotions.ColumnSystems.Pf.pc_count}

\begin{proof}
(1) By induction on $l$; for $l=0$, $\nu_0=0$. Let $l\leq2s$. If $l$ is a vector letter, of degree $g$, then
$u_{l+1}=u_l+1$, $v_{l+1}=v_l$, and $\nu_{l+1}$ is a strip of size $g$ over $\nu_l$: $\nu_l(a)\leq\nu_{l+1}(a)$,
$\nu_{l+1}(a+1)\leq\nu_l(a)$, and $\nu_{l+1}(a)\leq\nu_l(a)+g$, as the rises of the entries add up to $g$. If
$a>u_{l+1}$, then $a-1>u_l$, so $\nu_{l+1}(a)\leq\nu_l(a-1)\leq0$. If $a\leq m+1-v_{l+1}$, then
$\nu_{l+1}(a)\geq\nu_l(a)\geq0$. If $l$ is a covector letter, of degree $g$, then $u_{l+1}=u_l$, $v_{l+1}=v_l+1$, and
$\nu_l$ is a strip of size $g$ over $\nu_{l+1}$; the same argument, with the roles exchanged, gives the sign
conditions. In both cases $|\nu_{l+1}(a)|\leq|\nu_l(a)|+g$, and $g\leq A$ for $l=0$ and $g\leq M$ otherwise.

(2) Put $T=A+(2s+1)M$. For $l\leq2s+1$, $u_l\leq s+1$ and $v_l\leq s$. So by (1), the entries $\nu_l(a)$ with
$s+2\leq a\leq m+1-s$ vanish, and the others, at most $2s+1$ of them, lie in $[-T,T]$. As
$\nu_{2s+1}(a+1)\leq\lambda_a\leq\nu_{2s+1}(a)$, the entries $\lambda_a$ with $s+2\leq a\leq m-s$ vanish, and the others,
at most $2s+1$ of them, lie in $[-T,T]$. So a pair is determined by at most $(2s+2)(2s+1)+(2s+1)=(2s+3)(2s+1)$
integers in $[-T,T]$.
\end{proof}
% Lean: (1) is pc_shape, with pcU and pcL for u_l and v_l (pcU_vec, pcU_cov, pcL_vec, pcL_cov) and pc_strip_le;
%   the Lean bound is 0 at l = 0. (2) is pc_count, through the injection pcEnc into a product of copies of
%   Icc (-T) T (pcEnc_inj, pcEnc_abs), with Set.ncard_le_ncard_of_injOn.

\begin{lemma}\label{pd-kappa-defects-digits-count-pair-facts}
Let $p$, $\Lambda$, $m$, $s$, $A$, the $\hat b_i$ and the $\hat\alpha_i$ be as in
Definition~\ref{pd-kappa-defects-digits-count-pairs}, let $(\nu,\lambda)$ be a pair, and let $c_l$, $t_i$ and $\bar o_i$
be as there. Put $\kappa_l=(\nu_{l+1}(2),\dots,\nu_{l+1}(m+1))\in\mathbf Z^m$ for $0\leq l\leq2s$ and
$\kappa_{2s+1}=\lambda$, and $\bar w_i=|\kappa_{2i-1}|-|\kappa_{2i-2}|$ and $\bar\varsigma_i=\bar o_i-\bar w_i$ for
$1\leq i\leq s$.
\begin{enumerate}
\item $\kappa_0=0$ and $c_0=A$.
\item Let $1\leq i\leq s$. Then $\kappa_{2i-1}$ is a strip of size $\bar w_i$ over $\kappa_{2i-2}$ and a strip of size
  $\hat\alpha_i-\bar\varsigma_i$ over $\kappa_{2i}$; $0\leq\bar w_i\leq\hat b_i$; $\bar\varsigma_i=\nu_{2i}(1)-\nu_{2i+1}(1)$
  and $0\leq\bar\varsigma_i\leq\hat\alpha_i$; and $t_i\leq c_{i-1}$ and $t_i\leq c_i$.
\item $\lambda$ is a strip of size $|\lambda|-|\kappa_{2s}|\geq0$ over $\kappa_{2s}$, and $\lambda_1\leq c_s$.
\item $\nu_{2l+1}(2)\leq\sum_{i\leq l}\bar w_i$ for $0\leq l\leq s$, and $t_i\leq\nu_{2i-1}(2)+\bar w_i$ for
  $1\leq i\leq s$.
\item $c_l\leq A+\sum_{i\leq l}\hat b_i$ for $0\leq l\leq s$.
\item If $s<m$, then $|\lambda|-|\kappa_{2s}|\leq c_s+\sum_{i=1}^s\hat\alpha_i-\sum_{i=1}^s(-\lambda_{m-s+i})$.
\end{enumerate}
\end{lemma}
\leanref{Formalism/SharpRestriction/GrowthBounds/BinomialRatio/CoarseDegreeRatioConsecutive.lean}{Formalism.MainNotions.ColumnSystems.Pf.pair_top}
\leanref{Formalism/SharpRestriction/GrowthBounds/BinomialRatio/CoarseDegreeRatioConsecutive.lean}{Formalism.MainNotions.ColumnSystems.Pf.pairC_zero}
\leanref{Formalism/SharpRestriction/GrowthBounds/BinomialRatio/PairFacts.lean}{Formalism.MainNotions.ColumnSystems.Pf.pf_prev_le}
\leanref{Formalism/SharpRestriction/GrowthBounds/BinomialRatio/PairFacts.lean}{Formalism.MainNotions.ColumnSystems.Pf.pf_rho_le}
\leanref{Formalism/SharpRestriction/GrowthBounds/BinomialRatio/KeptBlocksOfSLine.lean}{Formalism.MainNotions.ColumnSystems.Pf.pairC_le_sum}
\leanref{Formalism/SharpRestriction/GrowthBounds/BinomialRatio/PairFacts.lean}{Formalism.MainNotions.ColumnSystems.Pf.pf_d_le}

\begin{proof}
Let $\beta$ be a strip of size $g$ over $\alpha$, both in $\mathbf Z^{r+1}$ with $r\geq1$, and let $\alpha',\beta'$ be
$\alpha,\beta$ without their first entries. Then $\beta'$ is a strip of size $g'=|\beta'|-|\alpha'|\geq0$ over
$\alpha'$, $\alpha(1)\leq\beta(1)$, $\beta(1)-\alpha(1)+g'=g$ and $\beta(2)\leq\alpha(1)$; and, as all entries rise and
their rises add up to $g$, $\beta(a)\leq\alpha(a)+g$ for every $a$. The path $\nu$ starts at $\nu_0=0$; at the vector
letter $2i-1$, $\nu_{2i}$ is a strip of size $\hat b_i$ over $\nu_{2i-1}$, and at the covector letter $2i$, $\nu_{2i}$ is
a strip of size $\hat\alpha_i$ over $\nu_{2i+1}$.

(1) At the vector letter $0$, $\nu_1$ is a strip of size $A$ over $0$, so $\nu_1(a)\geq0$ and $\nu_1(a+1)\leq0$ for
all $a$. So $\kappa_0=0$ and $c_0=\nu_1(1)=|\nu_1|=A$.

(2) By the remark, $\kappa_{2i-1}$ is a strip of size $\bar w_i\geq0$ over $\kappa_{2i-2}$, with
$\nu_{2i}(1)-\nu_{2i-1}(1)+\bar w_i=\hat b_i$ and $\nu_{2i}(1)\geq\nu_{2i-1}(1)$, so $\bar w_i\leq\hat b_i$. Likewise
$\kappa_{2i-1}$ is a strip of size $g'=|\kappa_{2i-1}|-|\kappa_{2i}|\geq0$ over $\kappa_{2i}$, with
$\nu_{2i}(1)-\nu_{2i+1}(1)+g'=\hat\alpha_i$ and $\nu_{2i}(1)\geq\nu_{2i+1}(1)$. Then
$\bar o_i=\nu_{2i-1}(1)+\hat b_i-\nu_{2i+1}(1)=\bar w_i+\nu_{2i}(1)-\nu_{2i+1}(1)$, so
$\bar\varsigma_i=\nu_{2i}(1)-\nu_{2i+1}(1)$, which lies in $[0,\hat\alpha_i]$, and $g'=\hat\alpha_i-\bar\varsigma_i$. Finally
$t_i=\nu_{2i}(2)\leq\nu_{2i-1}(1)=c_{i-1}$ and $t_i\leq\nu_{2i+1}(1)=c_i$, by the second condition on the two strips.

(3) The conditions $\nu_{2s+1}(a+1)\leq\lambda_a$ and $\lambda_{a+1}\leq\nu_{2s+1}(a+1)$ say that
$\kappa_{2s}(a)\leq\lambda_a$ and $\lambda_{a+1}\leq\kappa_{2s}(a)$; so $\lambda$ is a strip over $\kappa_{2s}$, of size
$|\lambda|-|\kappa_{2s}|\geq0$. And $\lambda_1\leq\nu_{2s+1}(1)=c_s$.

(4) By (1), $\nu_1(2)=0$. Let $1\leq l\leq s$. By (2), $\nu_{2l}(2)=\kappa_{2l-1}(1)\leq\kappa_{2l-2}(1)+\bar w_l=
\nu_{2l-1}(2)+\bar w_l$, which is the second claim for $i=l$, and $\nu_{2l+1}(2)=\kappa_{2l}(1)\leq\kappa_{2l-1}(1)=
\nu_{2l}(2)$. The first claim follows by induction on $l$.

(5) By (2), $\bar o_i=\bar w_i+\bar\varsigma_i\geq0$, so $c_i=c_{i-1}+\hat b_i-\bar o_i\leq c_{i-1}+\hat b_i$; and $c_0=A$
by (1).

(6) For a vector $v$ put $\pi(v)=\sum_a\max(0,-v(a))$. If $\beta$ is a strip of size $g$ over $\alpha$, then
$\pi(\beta)\leq\pi(\alpha)$, as the entries rise, and $\pi(\alpha)\leq\pi(\beta)+g$, as the entries of $\beta$ exceed those
of $\alpha$ by $g$ in all. By (1) and (2) and induction on $l$, $\pi(\kappa_{2l})\leq\sum_{i\leq l}\hat\alpha_i$.

Write $m=q+1+s$ with $q\geq0$ and $\kappa=\kappa_{2s}$, so that $|\lambda|-|\kappa|=\sum_{a=1}^m(\lambda_a-\kappa(a))$. As
$\lambda_{a+1}\leq\kappa(a)$, the terms with $a\leq q+1$ telescope:
$\sum_{a=1}^{q+1}(\lambda_a-\kappa(a))\leq\lambda_1-\kappa(q+1)$. Here $\lambda_1\leq c_s$ by (3), and
$\kappa(q+1)=\nu_{2s+1}(q+2)\geq0$ by Lemma~\ref{pd-kappa-defects-digits-count-pair-number}(1), as at most $s$
covector letters precede $2s+1$ and $q+2\leq m+1-s$. The terms with $a=q+1+i=m-s+i$, $1\leq i\leq s$, give
$\sum_i\lambda_{m-s+i}-\sum_i\kappa(m-s+i)$, and $-\sum_i\kappa(m-s+i)\leq\pi(\kappa)\leq\sum_i\hat\alpha_i$.
\end{proof}
% Lean: (1) chain_first_tail and pairC_zero; (2) pair_top (strip_tail for both letters; t3 to t6 and the last two
%   clauses); (3) inter_strip; (4) pf_prev_le and pf_rho_le (strip_first); (5) pairC_le_sum; (6) pf_d_le, with
%   pf_tele, pf_mass (pfMass_up, pfMass_down) and pc_shape for the sign of K q. The book's 1-based indices are
%   the Lean's plus one: c_l is pairC mu l, t_i is pairRho mu (i - 1), bar o_i is pairO mu Ai (i - 1), bar w_i is
%   pairW mu (i - 1), bar varsigma_i is pairSg mu Ai (i - 1), and -lambda_{m-s+i} is pairN lam (i - 1).

\begin{lemma}\label{pd-kappa-defects-digits-count-coarse}
Let $C_c\geq0$, $C_0\geq1$, $C_{E1},C_{E2}\geq1$ and $C_H\geq1000(C_{E1}+C_{E2}+1)(C_c+1)$. Let $p$, $n'$, $m$, $G'$,
$N_1$, $c_{cl}$, $D_p$, $c_*$, $T$, $Y$, $Y_b$, $b'$, $b_l$, $b_u$, $a$, $c$, $a_T$ and $b_T$ satisfy the hypotheses of
Lemma~\ref{pd-kappa-defects-digits-count-upper} with these constants, form the numbers of
Definition~\ref{pd-kappa-defects-numbers}, and let $\delta$ be the number of Theorem~\ref{pd-gain-theorem-one}. Let
$1\leq j\leq q$ and digits $\rho,\rho'$ for $j$ be given; write $s=S_j$, $\eta=\eta_j$, $\hat a=\hat a_j$ and
$P_i=p^{g_j-d_{j,i}}$, and let $\hat b_i$ and $\hat\alpha_i$ ($1\leq i\leq s$) be the numbers of
Definition~\ref{pd-kappa-defects-segment-lattice}(1) for the digits, with $\alpha_i=a_T(y_{j,i})$; by hypothesis (1) of
Theorem~\ref{pd-gain-theorem-one}, $s\leq N_1<m$. Then:
\begin{enumerate}
\item $p\leq P_i$, $\hat b_iP_i\leq2a$ and $\hat\alpha_iP_i\leq2a$ for $1\leq i\leq s$;
\item $\hat b_i(1+\delta)<\hat a$ for $1\leq i\leq\eta$, and $\hat a<\hat b_i(1-\delta)$ for $\eta<i\leq s$;
\item $p^2\hat b_l\leq2\hat b_{l+1}$ for $1\leq l<s$, and $\sum_{l<i}\hat b_l\leq4\hat b_i/p^2$ for $1\leq i\leq s$.
\end{enumerate}
\end{lemma}
\leanref{Formalism/SharpRestriction/GrowthBounds/BinomialRatio/CoarseDegreeAlternative.lean}{Formalism.MainNotions.ColumnSystems.Pf.coarse_degree_alternative}
\leanref{Formalism/SharpRestriction/GrowthBounds/BinomialRatio/CoarseDegreeAlternative.lean}{Formalism.MainNotions.ColumnSystems.Pf.dich_U}
\leanref{Formalism/SharpRestriction/GrowthBounds/BinomialRatio/CoarseDegreeRatioConsecutive.lean}{Formalism.MainNotions.ColumnSystems.Pf.coarse_degree_ratio_consecutive}
\leanref{Formalism/SharpRestriction/GrowthBounds/BinomialRatio/CoarseDegreeRatioConsecutive.lean}{Formalism.MainNotions.ColumnSystems.Pf.coarse_degree_prefix_sum}

\begin{proof}
Write $\varepsilon=N_1\Lambda/k$ and $B_i=b_T(y_{j,i})$, and let $\beta_e$ be as in the proof of
Lemma~\ref{pd-kappa-defects-digits-top-rows}. The facts (1) to (4) in that proof hold, with the digits in place of the
vectors of a kept block, as they use only the hypotheses here and the first condition for a kept block; so does the
fact (j) in the proof of Lemma~\ref{pd-kappa-defects-windows-model-step}.

(1) As $y_{j,i}$ lies strictly between $s_j$ and $s_{j+1}$, $d_{j,i}<g_j$, so $P_i\geq p$. By the facts (2) and (3),
$\hat b_iP_i\leq B_i\leq b'(1+\beta_e)\leq2b'\leq2a$. The entries of $\rho'_i$ are less than $P_i$, $|\rho'_i|\leq\alpha_i$
and $P_i$ divides $\alpha_i-|\rho'_i|$, so $\hat\alpha_iP_i=\alpha_i-|\rho'_i|\leq\alpha_i$; and $\alpha_i\leq a+2a/(kY)\leq2a$
by the fact (j), as $kY\geq5$.

(2) The first claim is the fact (4). Let $\eta<i\leq s$. As the $d_{j,i'}$ increase with $i'$ and $\eta$ is the number of
the $i'$ with $p^{d_{j,i'}}<a/b'$, $p^{d_{j,i}}\geq a/b'$. The number $d_{j,i}$ is odd, as in the fact (4), so
hypothesis (2) of Theorem~\ref{pd-gain-theorem-one} gives $p^{d_{j,i}}>(a/b')(1+2\delta)$. By the facts (2) and (3),
$\hat b_iP_i>B_i-kP_i\geq b'(1-\beta_e)-kP_i$, and $kP_i\leq kQ\leq\varepsilon b'$, as $k^3Q\leq b'$ and
$\varepsilon k=N_1\Lambda\geq1$. So $\hat b_iP_i>b'(1-\beta_e-\varepsilon)$, where $1-\beta_e-\varepsilon\geq1-\delta/2>0$, as
$\beta_e\leq\delta/2-18\varepsilon$. Multiplying by $p^{d_{j,i}}$, with $P_ip^{d_{j,i}}=p^{g_j}$,
\[
 \hat b_ip^{g_j}>(1-\beta_e-\varepsilon)(1+2\delta)a .
\]
As $\delta<1/6$, $(1-\beta_e-\varepsilon)(1-\delta)(1+2\delta)\geq(1-\delta/2)(1-\delta)(1+2\delta)=
1+\delta(1/2-5\delta/2+\delta^2)\geq1$. Hence $\hat b_i(1-\delta)p^{g_j}>a\geq a-E_0\geq\hat ap^{g_j}$, as $E_0\geq0$.

(3) As $d_{j,l}<d_{j,l+1}$ are both odd, $d_{j,l}+2\leq d_{j,l+1}$, so $p^2P_{l+1}\leq P_l$. By the facts (2) and (3),
$p^2\hat b_lP_{l+1}\leq\hat b_lP_l\leq B_l\leq b'(1+\beta_e)$ and $\hat b_{l+1}P_{l+1}>b'(1-\beta_e)-kP_{l+1}$. Here
$kP_{l+1}\leq kQ\leq b'/12$, as $k^3Q\leq b'$ and $k^2\geq12$, and $\beta_e\leq1/12$; so
$2\hat b_{l+1}P_{l+1}>\frac53b'\geq b'(1+\beta_e)\geq p^2\hat b_lP_{l+1}$. For the sum, argue by induction on $i$; for $i=1$
it is $0$. If the bound holds for some $i<s$, then
$\sum_{l\leq i}\hat b_l\leq\hat b_i(1+4/p^2)\leq(2\hat b_{i+1}/p^2)(1+4/25)\leq4\hat b_{i+1}/p^2$, as $p\geq5$.
\end{proof}
% Lean: (1) is coarse_bounds with winB and winA (as in pairCount_of_std and the hC step of
%   kept_pair_of_three_lines). (2) is coarse_degree_alternative: dich_S (Lemma pd-kappa-defects-digits-top-rows,
%   fact (4)) and dich_U with poly_U, sigma_iff, h.band and Std.kQeps. (3) is coarse_degree_ratio_consecutive
%   (dist_strictMono, dist_odd, coarse_bounds, winB, beLe, kQ) and coarse_degree_prefix_sum. Here Std is std_of
%   with H = 1 and the thresholds 1, 1 and T2 CE1 CE2 Cc 1.

\begin{lemma}\label{pd-kappa-defects-digits-count-pair-bound}
Let $C_c\geq0$, $C_0\geq1$, $C_{E1},C_{E2}\geq1$ and $C_H\geq1000(C_{E1}+C_{E2}+1)(C_c+1)$. Let the data, $j$, the digits,
$s$, $\hat a$, the $\hat b_i$ and the $\hat\alpha_i$ be as in Lemma~\ref{pd-kappa-defects-digits-count-coarse}. Let
$A\geq0$ be an integer with
$\hat a(1-N_1\Lambda/k)\leq A\leq\hat a$. Then there are at most $(8a)^{4(N_1+1)^2}$ pairs for $p$, $\Lambda$, $m$, $s$,
$A$ and the $\hat b_i$ and $\hat\alpha_i$ (Definition~\ref{pd-kappa-defects-digits-count-pairs}).
\end{lemma}
\leanref{Formalism/SharpRestriction/GrowthBounds/BinomialRatio/ThresholdStandardLarge.lean}{Formalism.MainNotions.ColumnSystems.Pf.PairCountAt}
\leanref{Formalism/SharpRestriction/GrowthBounds/BinomialRatio/PairCount.lean}{Formalism.MainNotions.ColumnSystems.Pf.pairCount_of_std}
\leanref{Formalism/SharpRestriction/GrowthBounds/BinomialRatio/PairCount.lean}{Formalism.MainNotions.ColumnSystems.Pf.restPairCount}

\begin{proof}
The facts in the proof of Lemma~\ref{pd-kappa-defects-digits-count-coarse} hold. Put $M=\lfloor2a/p\rfloor$. By
Lemma~\ref{pd-kappa-defects-digits-count-coarse}(1), $\hat b_i\leq2a/P_i\leq2a/p$, and likewise for $\hat\alpha_i$; so
$\hat b_i,\hat\alpha_i\leq M$. By Lemma~\ref{pd-kappa-defects-digits-count-pair-number}(2), there are at most
$\big(2(A+(2s+1)M)+1\big)^{(2s+3)(2s+1)}$ pairs. Here $A\leq\hat a\leq(a-E_0)/p^{g_j}\leq a$, and
$(2s+1)M\leq k\cdot2a/p\leq a$, as $2s+1\leq2N_1+3\leq k$ and $2k\leq k^2\leq p$; as $a\geq1$,
$2(A+(2s+1)M)+1\leq4a+1\leq8a$. Finally $(2s+3)(2s+1)\leq(2s+2)^2\leq4(N_1+1)^2$, and $8a\geq1$.
\end{proof}
% Lean: restPairCount = TransportedAlt.of_explicit with thresholds 1, 1 and T2 CE1 CE2 Cc 1, through std_of and
%   pairCount_of_std: pc_count with M = floor(2a/p), coarseA_mul and pc_coarseC_mul with winB and winA,
%   base_exponent_le_a for A <= a, and he for the exponent. PairCountAt is stated for D1num, HIext and CBound and
%   both bounds on A; only the upper bound is used.

\begin{lemma}\label{pd-kappa-defects-digits-count-weyl-bound}
Let $r\geq1$ be an integer and $\lambda\in\mathbf Z^r$ with $\lambda_1\geq\cdots\geq\lambda_r$. Then
$\operatorname{wd}(\lambda)\leq\prod_{a=1}^r\binom{|\lambda_a|+r-1}{r-1}$
(Definition~\ref{pd-kappa-defects-digits-chains}(1)).
\end{lemma}
\leanref{Formalism/SharpRestriction/GrowthBounds/BinomialRatio/WdimUp.lean}{Formalism.MainNotions.ColumnSystems.Pf.WdimUpAt}
\leanref{Formalism/SharpRestriction/GrowthBounds/BinomialRatio/WdimUp.lean}{Formalism.MainNotions.ColumnSystems.Pf.wdimUp_of}

\begin{proof}
Give the letters $0,\dots,2r$ types and degrees: for $1\leq a\leq r$, the letter $a-1$ is a vector letter of degree
$\max(\lambda_a,0)$ and the letter $2r-a$ a covector letter of degree $\max(-\lambda_a,0)$, and the letter $2r$ is a vector
letter of degree $0$. For $0\leq n\leq2r+1$ let $\nu_n\in\mathbf Z^r$ be given by $\nu_n(a)=\max(\lambda_a,0)$ for $a\leq n$
and $\nu_n(a)=0$ for $a>n$ if $n\leq r$, and by $\nu_n(a)=\lambda_a$ for $a\geq2r+1-n$ and $\nu_n(a)=\max(\lambda_a,0)$ for
$a<2r+1-n$ if $n>r$. Then $\nu_0=0$ and $\nu_{2r+1}=\lambda$. At the vector letter $n<r$, $\nu_{n+1}$ differs from $\nu_n$
only at the entry $n+1$, raised from $0$ to $\max(\lambda_{n+1},0)$; at the covector letter $n$ with $r\leq n<2r$, $\nu_{n+1}$
differs from $\nu_n$ only at the entry $a=2r-n$, lowered from $\max(\lambda_a,0)$ to $\lambda_a$; and $\nu_{2r+1}=\nu_{2r}$. As
$\lambda$ is non-increasing, so are the $\nu_n$, and at each letter the two vectors satisfy the strip conditions of
Definition~\ref{pd-symbols-path-generators}(2) with the given degree. So $\nu$ is a path of rank $r$ for these degrees.

By Lemma~\ref{pd-kappa-defects-digits-chain-dimension} with $q'=0$, as in the proof of
Lemma~\ref{pd-kappa-defects-digits-count-upper}, the sum of $\operatorname{wd}(\mu_{2r+1})$ over the paths $\mu$ of rank $r$
for these degrees is the number of monomials of these row degrees, $\prod_l\binom{\deg(l)+r-1}{r-1}$, and every term is
$\geq0$, as $\mu_{2r+1}$ is non-increasing. So $\operatorname{wd}(\lambda)$ is at most this product. For each $a$, one of
$\max(\lambda_a,0)$ and $\max(-\lambda_a,0)$ is $0$ and the other is $|\lambda_a|$, and a letter of degree $0$ contributes
$1$; so the product is $\prod_a\binom{|\lambda_a|+r-1}{r-1}$.
\end{proof}
% Lean: wdimUp_of hCD: wu_single (the term of the path in chainSum with q = 0, ChainDimAt, le_finsum_of_nonneg),
%   the letters wuLL with degrees wuDeg, the path wuSt (wu_chain, wu_step, wu_sum) and wu_prod.

\begin{lemma}\label{pd-kappa-defects-digits-count-weight}
Let $p$, $\Lambda$, $m$, $s$, $A$, the $\hat b_i$ and the $\hat\alpha_i$ be as in
Definition~\ref{pd-kappa-defects-digits-count-pairs}, with $2s+1\leq m$, let $(\nu,\lambda)$ be a pair, and let $t_i$ be as
there. Write $x^+=\max(x,0)$ and $P_{m-1}(d)=\binom{d+m-1}{m-1}$ for integers $d\geq0$. Then
\[
 \operatorname{wd}(\lambda)\leq P_{m-1}(\lambda_1^+)\prod_{i=1}^sP_{m-1}(t_i^+)\prod_{i=1}^sP_{m-1}\big((-\lambda_{m-s+i})^+\big).
\]
\end{lemma}
\leanref{Formalism/SharpRestriction/GrowthBounds/BinomialRatio/NotKept.lean}{Formalism.MainNotions.ColumnSystems.Pf.WtBoundAt}
\leanref{Formalism/SharpRestriction/GrowthBounds/BinomialRatio/WtBound.lean}{Formalism.MainNotions.ColumnSystems.Pf.wtBound_of}
\leanref{Formalism/SharpRestriction/GrowthBounds/BinomialRatio/WtBound.lean}{Formalism.MainNotions.ColumnSystems.Pf.wb_shape}
\leanref{Formalism/SharpRestriction/GrowthBounds/BinomialRatio/WtBound.lean}{Formalism.MainNotions.ColumnSystems.Pf.wb_follow}

\begin{proof}
Let $\kappa_l$ be the vectors of Lemma~\ref{pd-kappa-defects-digits-count-pair-facts}. By parts (1) and (2) of that lemma,
$\kappa_0=0$; at the vector letter $2i+1$, $\kappa_{2i+1}(a)\geq\kappa_{2i}(a)$ and $\kappa_{2i+1}(a+1)\leq\kappa_{2i}(a)$; and
at the covector letter $2i+2$, $\kappa_{2i+2}(a)\leq\kappa_{2i+1}(a)$ and $\kappa_{2i+1}(a+1)\leq\kappa_{2i+2}(a)$.

(a) \emph{For $0\leq i\leq s$: $\kappa_{2i}(a)\geq0$ for $a\leq i$, $\kappa_{2i}(a)=0$ for $i<a\leq m-i$, and
$\kappa_{2i}(a)\leq0$ for $a>m-i$.} By induction on $i$; for $i=0$, $\kappa_0=0$. Suppose it holds for some $i<s$; as
$2s+1\leq m$, $i+1<m-i$. The rules at the vector letter give $\kappa_{2i+1}(a)\geq0$ for $a\leq i+1$,
$\kappa_{2i+1}(a)=0$ for $i+1<a\leq m-i$, and $\kappa_{2i+1}(a)\leq0$ for $a>m-i$; then the rules at the covector letter
give the claim for $i+1$.

(b) \emph{For $0\leq i_0\leq i<s$: $\kappa_{2i+2}(i-i_0+1)\leq t_{i_0+1}$.} Here $t_{i_0+1}=\nu_{2i_0+2}(2)=\kappa_{2i_0+1}(1)$. For
$i=i_0$, $\kappa_{2i_0+2}(1)\leq\kappa_{2i_0+1}(1)$. If it holds for $i$, then
$\kappa_{2i+4}(i-i_0+2)\leq\kappa_{2i+3}(i-i_0+2)\leq\kappa_{2i+2}(i-i_0+1)\leq t_{i_0+1}$.

Put $\kappa=\kappa_{2s}$. By Lemma~\ref{pd-kappa-defects-digits-count-pair-facts}(3), $\kappa(a)\leq\lambda_a$ and
$\lambda_{a+1}\leq\kappa(a)$. By (a) with $i=s$: $\lambda_1\geq\kappa(1)\geq0$. For $2\leq a\leq s+1$, $\lambda_a\geq\kappa(a)\geq0$
and, by (b) with $i=s-1$ and $i_0=s+1-a$, $\lambda_a\leq\kappa(a-1)\leq t_{s+2-a}$; so $|\lambda_a|\leq t_{s+2-a}^+$. For
$s+2\leq a\leq m-s$, $0=\kappa(a)\leq\lambda_a\leq\kappa(a-1)=0$. For $a>m-s$, $\lambda_a\leq\kappa(a-1)\leq0$, so
$|\lambda_a|=(-\lambda_a)^+$. The sequence $\lambda$ is non-increasing, so
Lemma~\ref{pd-kappa-defects-digits-count-weyl-bound} with $r=m$, and the monotonicity of $P_{m-1}$, give the claim; the
entries $a=m-s+i$ give the last product.
\end{proof}
% Lean: wtBound_of hCD: wdimUp_of (Lemma pd-kappa-defects-digits-count-weyl-bound), wb_shape ((a), through wb_one,
%   wb_vstep and wb_cstep, for wbK mu (2 i + 1), the tail of the Lean's mu (2 i + 1)), wb_follow ((b)), wb_last and
%   the comparison g with wb_sdXi_mono; the factors are sdXi r of the toNat of lam 0, pairRho mu i and pairN lam i.

\begin{lemma}\label{pd-kappa-defects-digits-count-bracket-ratio}
Let $C_c\geq0$, $C_0\geq1$, $C_{E1},C_{E2}\geq1$ and $C_H\geq1000(C_{E1}+C_{E2}+1)(C_c+1)$. Let the data, $j$, the digits,
$s$ and the $\hat\alpha_i$ be as in Lemma~\ref{pd-kappa-defects-digits-count-coarse}. Then $p^2\hat\alpha_l\leq2\hat\alpha_{l+1}$
for $1\leq l<s$, and $\sum_{l<i}\hat\alpha_l\leq4\hat\alpha_i/p^2$ for $1\leq i\leq s$.
\end{lemma}
\leanref{Formalism/SharpRestriction/GrowthBounds/BinomialRatio/CoarseThresholdRatioConsecutive.lean}{Formalism.MainNotions.ColumnSystems.Pf.coarse_threshold_ratio_consecutive}
\leanref{Formalism/SharpRestriction/GrowthBounds/BinomialRatio/CoarseThresholdRatioConsecutive.lean}{Formalism.MainNotions.ColumnSystems.Pf.coarse_threshold_prefix_sum}

\begin{proof}
Write $\alpha_i=a_T(y_{j,i})$ and $P_i=p^{g_j-d_{j,i}}$. The facts in the proof of
Lemma~\ref{pd-kappa-defects-digits-count-coarse} hold. They give $k\geq5$, $Y\geq1$, $k^3Q\leq b'\leq a$, $P_i\leq Q$ and
$a\leq\alpha_i\leq a+2a/(kY)$. Let $1\leq l<s$. As in the proof of Lemma~\ref{pd-kappa-defects-digits-count-coarse}(3),
$p^2P_{l+1}\leq P_l$. For each $i$ the entries of $\rho'_i$ are less than $P_i$, $|\rho'_i|\leq\alpha_i$ and $P_i$ divides
$\alpha_i-|\rho'_i|$, so $\hat\alpha_iP_i=\alpha_i-|\rho'_i|$, and $\alpha_i-kP_i<\hat\alpha_iP_i\leq\alpha_i$. As $kY\geq5$,
$\alpha_l\leq7a/5$, so $p^2\hat\alpha_lP_{l+1}\leq\hat\alpha_lP_l\leq7a/5$. As $k^2\geq25$,
$kP_{l+1}\leq kQ\leq b'/25\leq a/25$, so $\hat\alpha_{l+1}P_{l+1}>a-kP_{l+1}\geq24a/25$. Hence
$2\hat\alpha_{l+1}P_{l+1}>48a/25\geq p^2\hat\alpha_lP_{l+1}$, which gives the first claim. The second follows from the
first by induction on $i$, as in the proof of Lemma~\ref{pd-kappa-defects-digits-count-coarse}(3).
\end{proof}
% Lean: coarse_threshold_ratio_consecutive (dist_strictMono, dist_odd, levels, coarse_bounds for ry, winA, Std.kQ,
%   pab; hwin is 1 + 2/(kY) <= 7/5 and hkP is k P' <= b'/25) and coarse_threshold_prefix_sum (induction on i), for
%   Std with H >= 1; they are used with H = 1, from std_of with thresholds 1, 1 and T2 CE1 CE2 Cc 1.

\begin{lemma}\label{pd-kappa-defects-digits-count-bracket-facts}
Let $C_c\geq0$, $C_0\geq1$, $C_{E1},C_{E2}\geq1$ and $C_H\geq1000(C_{E1}+C_{E2}+1)(C_c+1)$. Let the data, $j$, the digits,
$s$, $\hat a$, the $\hat b_i$ and the $\hat\alpha_i$ be as in Lemma~\ref{pd-kappa-defects-digits-count-coarse}. Let $A\geq0$ be
an integer with $\hat a(1-N_1\Lambda/k)\leq A\leq\hat a$, let $(\nu,\lambda)$ be a pair for $p$, $\Lambda$, $m$, $s$, $A$ and
the $\hat b_i$ and $\hat\alpha_i$, and let $c_l$, $t_i$, $\bar o_i$ and $n_i$ be as in
Definition~\ref{pd-kappa-defects-digits-count-pairs}. Then:
\begin{enumerate}
\item $c_0=A$, and $c_l\geq0$ for $0\leq l\leq s$;
\item $0\leq t_i\leq c_{i-1}$, $t_i\leq\hat b_i(1+4/p^2)$ and $\bar o_i-t_i\geq-4\hat b_i/p^2$ for $1\leq i\leq s$;
\item $0\leq\lambda_1\leq c_s$;
\item $n_i=-\lambda_{m-s+i}$ and $0\leq n_i\leq\hat\alpha_i(1+4/p^2)$ for $1\leq i\leq s$.
\end{enumerate}
The bounds on $A$ are not used.
\end{lemma}
\leanref{Formalism/SharpRestriction/GrowthBounds/BinomialRatio/CoarseDegreesLargeBeyond.lean}{Formalism.MainNotions.ColumnSystems.Pf.PairFactsAt}
\leanref{Formalism/SharpRestriction/GrowthBounds/BinomialRatio/CoarseDegreesLargeBeyond.lean}{Formalism.MainNotions.ColumnSystems.Pf.RestPairFacts}
\leanref{Formalism/SharpRestriction/GrowthBounds/BinomialRatio/CoarseDegreesLargeOfStd.lean}{Formalism.MainNotions.ColumnSystems.Pf.restPairFacts_of}
\leanref{Formalism/SharpRestriction/GrowthBounds/BinomialRatio/CoarseDegreesLargeOfStd.lean}{Formalism.MainNotions.ColumnSystems.Pf.NUpAt}
\leanref{Formalism/SharpRestriction/GrowthBounds/BinomialRatio/CoarseThresholdRatioConsecutive.lean}{Formalism.MainNotions.ColumnSystems.Pf.nUp_of_std}
\leanref{Formalism/SharpRestriction/GrowthBounds/BinomialRatio/CoarseThresholdRatioConsecutive.lean}{Formalism.MainNotions.ColumnSystems.Pf.nu_le_last}
\leanref{Formalism/SharpRestriction/GrowthBounds/BinomialRatio/CoarseThresholdRatioConsecutive.lean}{Formalism.MainNotions.ColumnSystems.Pf.last_ge_mass}

\begin{proof}
The facts in the proof of Lemma~\ref{pd-kappa-defects-digits-count-coarse} hold. By hypothesis (1) of
Theorem~\ref{pd-gain-theorem-one}, $s\leq N_1$ and $2N_1+2\leq m$, so $2s+1\leq m$. Let $\kappa_l$, $\bar w_i$ and
$\bar\varsigma_i$ be as in Lemma~\ref{pd-kappa-defects-digits-count-pair-facts}. As $2s+1\leq m$, the claim (a) in the proof
of Lemma~\ref{pd-kappa-defects-digits-count-weight} holds, and so do the rules at the letters stated there: at the vector
letter $2i+1$, $\kappa_{2i}(a)\leq\kappa_{2i+1}(a)$, and at the covector letter $2i+2$, $\kappa_{2i+1}(a+1)\leq\kappa_{2i+2}(a)$.

(1) $c_0=A$ by Lemma~\ref{pd-kappa-defects-digits-count-pair-facts}(1). Let $1\leq i\leq s$. Then
$t_i=\nu_{2i}(2)=\kappa_{2i-1}(1)\geq\kappa_{2i-2}(1)$, by the rule at the vector letter $2i-1$, and $\kappa_{2i-2}(1)\geq0$, by
(a) if $i\geq2$ and as $\kappa_0=0$ if $i=1$. So $t_i\geq0$, and $c_i\geq t_i\geq0$ by
Lemma~\ref{pd-kappa-defects-digits-count-pair-facts}(2).

(2) The bounds $0\leq t_i\leq c_{i-1}$ are (1) and Lemma~\ref{pd-kappa-defects-digits-count-pair-facts}(2). By parts (4) and
(2) of that lemma and Lemma~\ref{pd-kappa-defects-digits-count-coarse}(3),
$\nu_{2i-1}(2)\leq\sum_{l<i}\bar w_l\leq\sum_{l<i}\hat b_l\leq4\hat b_i/p^2$. By the same parts,
$t_i\leq\nu_{2i-1}(2)+\bar w_i\leq\hat b_i(1+4/p^2)$; and as $\bar o_i=\bar w_i+\bar\varsigma_i$ with $\bar\varsigma_i\geq0$,
$\bar o_i-t_i\geq\bar\varsigma_i-\nu_{2i-1}(2)\geq-4\hat b_i/p^2$.

(3) By Lemma~\ref{pd-kappa-defects-digits-count-pair-facts}(3), $\kappa_{2s}(1)\leq\lambda_1\leq c_s$; and
$\kappa_{2s}(1)\geq0$, by (a) if $s\geq1$ and as $\kappa_0=0$ if $s=0$.

(4) As $s\leq m$, $n_i=-\lambda_{m-s+i}$. By Lemma~\ref{pd-kappa-defects-digits-count-pair-facts}(3),
$\lambda_{m-s+i}\leq\kappa_{2s}(m-s+i-1)$, which is $0$ for $i=1$, as $s<m-s$, and $\leq0$ for $i\geq2$, by (a). So $n_i\geq0$.
For $i\leq l<s$ and $2\leq a\leq m$, the rules at the letters $2l+1$ and $2l+2$ give
$\kappa_{2l}(a)\leq\kappa_{2l+1}(a)\leq\kappa_{2l+2}(a-1)$. As $s<m$, induction on $l$ gives
$\kappa_{2i}(m)\leq\kappa_{2s}(m-s+i)$, and $\kappa_{2s}(m-s+i)\leq\lambda_{m-s+i}$ by
Lemma~\ref{pd-kappa-defects-digits-count-pair-facts}(3). So $n_i\leq-\kappa_{2i}(m)\leq\pi(\kappa_{2i})$, where
$\pi(v)=\sum_a\max(0,-v(a))$; and $\pi(\kappa_{2i})\leq\sum_{l\leq i}\hat\alpha_l$, as shown in the proof of
Lemma~\ref{pd-kappa-defects-digits-count-pair-facts}(6). By Lemma~\ref{pd-kappa-defects-digits-count-bracket-ratio},
$\sum_{l<i}\hat\alpha_l\leq4\hat\alpha_i/p^2$, so $n_i\leq\hat\alpha_i(1+4/p^2)$.
\end{proof}
% Lean: restPairFacts_of restNUp, with the thresholds 1, 1 and T2 CE1 CE2 Cc 1 of std_of (H = 1) and of restNUp;
%   hr is 2s + 1 <= m, from segB and N1k. (1) is pairC_zero and pfa_c_nonneg (pf_rho_le_c, pfa_rho_nonneg with
%   wb_vstep, and pfa_odd_nonneg from wb_one and wb_shape). (2) is pf_rho_le_c, pfa_rho_up and pfa_u_low (pf_rho_le,
%   pf_u_ge, pfa_prev_real from pf_prev_le and pair_top, and hNq from coarse_degree_prefix_sum). (3) is
%   pair_first_weight_bounds (wb_last, pfa_odd_nonneg). (4) is pfa_n_nonneg (wb_shape, wb_last) and NUpAt = nUp_of_std:
%   nu_le_last (nu_follow, with wb_vstep and wb_cstep), last_ge_mass (pf_mass) and coarse_threshold_prefix_sum.

\begin{lemma}\label{pd-kappa-defects-digits-count-bracket-large}
Let $C_c\geq0$, $C_0\geq1$, $C_{E1},C_{E2}\geq1$ and $C_H\geq4000(C_{E1}+C_{E2}+1)(C_c+1)$. Let the data, $j$, the digits,
$s$, $\hat a$, the $\hat b_i$ and the $\hat\alpha_i$ be as in Lemma~\ref{pd-kappa-defects-digits-count-coarse}, and let $A\geq0$ be
an integer with $\hat a(1-N_1\Lambda/k)\leq A$. Then $k^3\leq A$, $k^3\leq\hat b_i$ and $k^3\leq\hat\alpha_i$ for $1\leq i\leq s$,
$k^2\leq p$, $k\geq10$, $\Lambda\geq1$ and $4\Lambda\leq k$.
\end{lemma}
\leanref{Formalism/SharpRestriction/GrowthBounds/BinomialRatio/CoarseDegreesLargeBeyond.lean}{Formalism.MainNotions.ColumnSystems.Pf.CoarseDegreesLargeAt}
\leanref{Formalism/SharpRestriction/GrowthBounds/BinomialRatio/CoarseDegreesLargeBeyond.lean}{Formalism.MainNotions.ColumnSystems.Pf.CoarseDegreesLargeBeyond}
\leanref{Formalism/SharpRestriction/GrowthBounds/BinomialRatio/CoarseDegreesLargeOfStd.lean}{Formalism.MainNotions.ColumnSystems.Pf.coarse_degrees_large_of_std}
\leanref{Formalism/SharpRestriction/GrowthBounds/BinomialRatio/CoarseDegreesLargeOfStd.lean}{Formalism.MainNotions.ColumnSystems.Pf.coarse_degrees_large_holds}

\begin{proof}
Write $\varepsilon=N_1\Lambda/k$, and let $\delta$ be the number of Theorem~\ref{pd-gain-theorem-one} and $\beta_e$ the number
of the proof of Lemma~\ref{pd-kappa-defects-digits-top-rows}. The facts in the proof of
Lemma~\ref{pd-kappa-defects-digits-count-coarse} hold. They give $k^2\leq p$, $\Lambda\geq1$, $N_1\geq1$, $b'\leq a$,
$40\varepsilon\leq\delta<1/6$, $\hat a\geq a/(2Q)$ and $\beta_e\leq1/12$. So $240N_1\Lambda<k$; as $N_1\Lambda\geq1$, this gives
$k\geq10$, $4\Lambda\leq k$ and $\varepsilon\leq1/2$. By the fact (b) in the proof of Lemma~\ref{pd-kappa-defects-margin},
$b'\geq C_Hk^3NQ^4FY$ with $F\geq8$; as $C_H\geq1$ and $N,Q,Y\geq1$, $4k^3Q\leq b'$. Let $1\leq i\leq s$, and put
$B_i=b_T(y_{j,i})$, $\alpha_i=a_T(y_{j,i})$ and $P_i=p^{g_j-d_{j,i}}\leq Q$.

\emph{The degree $A$.} $A\geq\hat a(1-\varepsilon)\geq\hat a/2\geq a/(4Q)\geq b'/(4Q)\geq k^3$.

\emph{The degree $\hat b_i$.} By the facts (2) and (3) in the proof of Lemma~\ref{pd-kappa-defects-digits-top-rows},
$B_i\geq b'(1-\beta_e)\geq b'/2$ and $\hat b_iP_i>B_i-kP_i$. So $\hat b_i>B_i/P_i-k\geq b'/(2Q)-k\geq2k^3-k\geq k^3$.

\emph{The degree $\hat\alpha_i$.} As in the proof of Lemma~\ref{pd-kappa-defects-digits-count-bracket-ratio},
$\hat\alpha_iP_i>\alpha_i-kP_i$, and $\alpha_i\geq a$. So $\hat\alpha_i>a/Q-k\geq2k^3-k\geq k^3$, as $4k^3Q\leq b'\leq a$.
\end{proof}
% Lean: coarse_degrees_large_holds = TransportedAlt.of_explicit with thresholds 1, 1 and T2 CE1 CE2 Cc 4 =
%   4000 (CE1 + CE2 + 1)(Cc + 1), through coarse_degrees_large_of_std (Std with H = 4). The facts are k5, Y1, Q1,
%   N2, one_le_level_alt, La1, bpos, ba, de6 and LaK (hK240, hk10, hLa4), pk (hkp), and Hb with H = 4 (hb4 :
%   4 k^3 Q <= b'). A uses half_ratio_le_segment_base_of_std and hfrac; hat b_i and hat alpha_i use
%   coarse_degree_gt_quotient (coarse_bounds), levels, winB with beLe and be0 (be <= 1/2), and winA. The book reads
%   4 k^3 Q <= b' off the fact (b) of Lemma pd-kappa-defects-margin, where std_of derives Hb from H <= C_H F.

\begin{lemma}\label{pd-kappa-defects-digits-count-bracket-core}
Let $s\geq0$ and $m\geq1$ be integers with $2s+1\leq m$, and put $k=m+1$. Let $\Lambda$, $p$, $A$ and $l_0$ be real numbers,
and let $\hat b_i$, $\hat\alpha_i$, $t_i$, $u_i$, $n_i$ ($1\leq i\leq s$) and $c_l$ ($0\leq l\leq s$) be real numbers such that:
\begin{enumerate}
\item $k\geq10$, $\Lambda\geq1$, $4\Lambda\leq k$, $k^2\leq p$, $k^3\leq A$, and $k^3\leq\hat b_i$ and $k^3\leq\hat\alpha_i$ for
  all $i$;
\item $c_0=A$, $c_l\geq0$ for all $l$, and $c_i=c_{i-1}+\hat b_i-u_i-t_i$ for $1\leq i\leq s$;
\item $0\leq t_i\leq c_{i-1}$, $t_i\leq\hat b_i(1+4/p^2)$ and $u_i\geq-4\hat b_i/p^2$ for all $i$;
\item $0\leq l_0\leq c_s$, and $0\leq n_i\leq\hat\alpha_i(1+4/p^2)$ for all $i$;
\item for some $i_0$, $u_{i_0}+4(\hat b_{i_0}+m)/p^2>\Lambda\max(c_{i_0-1},\hat b_{i_0})/k$ or
  $\hat\alpha_{i_0}(1-\Lambda/k)>n_{i_0}$.
\end{enumerate}
Then there is a real number $\theta\geq0$ with $\theta^{m-1}\leq e^{2-\Lambda/5}$ and
\[
 (l_0+m)\prod_{i=1}^s(t_i+m)\prod_{i=1}^s(n_i+m)\leq\theta\,A\prod_{i=1}^s\hat b_i\prod_{i=1}^s\hat\alpha_i .
\]
\end{lemma}
\leanref{Formalism/SharpRestriction/GrowthBounds/BinomialRatio/Bracket.lean}{Formalism.MainNotions.ColumnSystems.Pf.brk_core}
\leanref{Formalism/SharpRestriction/GrowthBounds/BinomialRatio/Bracket.lean}{Formalism.MainNotions.ColumnSystems.Pf.brk_hat}
\leanref{Formalism/SharpRestriction/GrowthBounds/BinomialRatio/NkExponent.lean}{Formalism.MainNotions.ColumnSystems.Pf.nk_iv}
\leanref{Formalism/SharpRestriction/GrowthBounds/BinomialRatio/NkExponent.lean}{Formalism.MainNotions.ColumnSystems.Pf.nk_vi}
\leanref{Formalism/SharpRestriction/GrowthBounds/BinomialRatio/NkExponent.lean}{Formalism.MainNotions.ColumnSystems.Pf.nk_tele}
\leanref{Formalism/SharpRestriction/GrowthBounds/BinomialRatio/NkExponent.lean}{Formalism.MainNotions.ColumnSystems.Pf.nk_bracket}
\leanref{Formalism/SharpRestriction/GrowthBounds/BinomialRatio/NkExponent.lean}{Formalism.MainNotions.ColumnSystems.Pf.nk_exponent}

\begin{proof}
Put $\eta=1/k^2$, $\tau=\Lambda/(4k)$ and $\delta=4/p^2$.

\emph{Numbers.} As $k\geq10$, $\eta\leq1/100$ and $m\geq9$. As $p^2\geq k^4\geq4k^2$, $\delta\leq\eta$. For $y\geq k^3$,
$m\leq k\leq\eta y$. As $4\Lambda\leq k$, $0<\tau\leq1/16$. For $y\geq k^3$ this gives
\[
 (y+m)(1+\delta)\leq(y+\eta y)(1+\eta)=y(1+\eta)^2\leq e^{2\eta}y .
\]
We call this bound (i).

\emph{Shifted numbers.} Put $c'_l=c_l+m\geq m>0$, and for $1\leq i\leq s$ put $x_i=t_i+m$, $h_i=(\hat b_i+m)(1+\delta)$ and
$v_i=u_i+\delta(\hat b_i+m)$. By (2), $c'_i=c'_{i-1}+h_i-v_i-x_i$. By (3), $0\leq x_i\leq c'_{i-1}$; $x_i\leq h_i$, as
$t_i+m\leq\hat b_i(1+\delta)+m(1+\delta)$; and $v_i\geq0$, as $u_i\geq-\delta\hat b_i$.

\emph{Two inequalities.} Let $x,c,h,v$ be real numbers with $0\leq x\leq c$, $x\leq h$ and $v\geq0$. Then
\[
 x(c+h-v-x)=ch-(c-x)(h-x)-xv\leq ch .
\]
If moreover $c,h>0$, $0<\tau<1/4$ and $2\tau\max(c,h)<v$, then $x(c+h-v-x)\leq(1-\tau)ch$. Indeed, let $c\leq h$; the
other case is the same with $c$ and $h$ exchanged. If $x\geq c/2$, then $xv\geq(c/2)(2\tau h)=\tau ch$. If $x<c/2$, then
$(c-x)(h-x)\geq(c/2)(h/2)\geq\tau ch$.

\emph{The product.} Let $\kappa_i,\mu_i\geq0$ ($1\leq i\leq s$) satisfy $x_ic'_i\leq\kappa_ie^{2\eta}\hat b_ic'_{i-1}$ and
$n_i+m\leq\mu_ie^{2\eta}\hat\alpha_i$ for every $i$, and $\prod_i\kappa_i\prod_i\mu_i\leq e^{-\tau}$. We show that
$\theta=e^\eta e^{4s\eta}\prod_i\kappa_i\prod_i\mu_i$ has the two properties. By induction on $l$,
$c'_l\prod_{i\leq l}x_i\leq c'_0\prod_{i\leq l}\kappa_ie^{2\eta}\hat b_i$. As $l_0+m\leq c'_s$ by (4), and
$c'_0=A+m\leq(1+\eta)A\leq e^\eta A$ by (1) and (2),
\[
 (l_0+m)\prod_ix_i\prod_i(n_i+m)\leq e^\eta A\prod_i\kappa_ie^{2\eta}\hat b_i\prod_i\mu_ie^{2\eta}\hat\alpha_i
 =\theta\,A\prod_i\hat b_i\prod_i\hat\alpha_i .
\]
Next, $0\leq\theta\leq e^{(4s+1)\eta-\tau}$, so $\theta^{m-1}\leq\exp\big((k-2)(4s+1)/k^2-(k-2)\Lambda/(4k)\big)$. As
$4s+1\leq2m-1=2k-3$, $(k-2)(4s+1)\leq(k-2)(2k-3)\leq2k^2$; and $(k-2)/(4k)\geq1/5$, as $k\geq10$. So
$\theta^{m-1}\leq e^{2-\Lambda/5}$.

\emph{The choice.} For every $i$, the first of the two inequalities, with $(x,c,h,v)=(x_i,c'_{i-1},h_i,v_i)$, and (i) give
$x_ic'_i\leq c'_{i-1}h_i\leq e^{2\eta}\hat b_ic'_{i-1}$; and by (4) and (1),
$n_i+m\leq\hat\alpha_i(1+\delta)+\eta\hat\alpha_i\leq(1+2\eta)\hat\alpha_i\leq e^{2\eta}\hat\alpha_i$. So $\kappa_i=\mu_i=1$
satisfy the conditions at $i$, and it remains to choose $\kappa_{i_0}$ and $\mu_{i_0}$.

If $u_{i_0}+4(\hat b_{i_0}+m)/p^2\leq\Lambda\max(c_{i_0-1},\hat b_{i_0})/k$, then $n_{i_0}<\hat\alpha_{i_0}(1-\Lambda/k)$ by (5).
Take $\kappa_{i_0}=1$ and $\mu_{i_0}=1-\Lambda/k+\eta$, which is $\geq0$ as $\Lambda/k\leq1/4$. Then
$n_{i_0}+m\leq\hat\alpha_{i_0}(1-\Lambda/k)+\eta\hat\alpha_{i_0}=\mu_{i_0}\hat\alpha_{i_0}\leq\mu_{i_0}e^{2\eta}\hat\alpha_{i_0}$, and
the product of the $\kappa_i$ and $\mu_i$ is $\mu_{i_0}\leq e^{\eta-\Lambda/k}\leq e^{-\tau}$, as $\eta\leq3\Lambda/(4k)$.

Otherwise $v_{i_0}=u_{i_0}+4(\hat b_{i_0}+m)/p^2>\Lambda\max(c_{i_0-1},\hat b_{i_0})/k$. Take $\mu_{i_0}=1$ and
$\kappa_{i_0}=1-\tau$. As $m\leq\hat b_{i_0}$, $c'_{i_0-1}\leq2\max(c_{i_0-1},\hat b_{i_0})$; by (i) and $e^{1/50}\leq2$,
$h_{i_0}\leq2\hat b_{i_0}$. So $2\tau\max(c'_{i_0-1},h_{i_0})\leq\Lambda\max(c_{i_0-1},\hat b_{i_0})/k<v_{i_0}$, and the second of
the two inequalities, with (i), gives $x_{i_0}c'_{i_0}\leq(1-\tau)c'_{i_0-1}h_{i_0}\leq(1-\tau)e^{2\eta}\hat b_{i_0}c'_{i_0-1}$.
The product of the $\kappa_i$ and $\mu_i$ is $1-\tau\leq e^{-\tau}$.
\end{proof}
% Lean: brk_core, with k = m + 1, eta = 1/k^2, tau = La/(4k) and delta = 4/p^2: the numbers are heta1, hdelta_eta,
%   hm_eta, htau0, htau1 and hr9, and (i) is brk_hat. The shifted numbers are c, hstep0 and hfacts. The two
%   inequalities are nk_iv and nk_vi. The product is fin: nk_bracket (through nk_tele), hcA and nk_exponent, with theta
%   = exp(eta) exp(2 eta)^s exp(2 eta)^s times the product of the kappa_i and mu_i. The choice is the case split hP:
%   mu0 = 1 - La/k + eta with eta <= 3 La/(4k), and kappa_{i0} = 1 - tau with hmax (exp(1/50) <= 2).

\begin{lemma}\label{pd-kappa-defects-digits-count-bracket}
For every $C_c\geq0$ there is $T_0>0$ such that for every $C_0\geq T_0$ there is $T_1>0$ such that for
all $C_{E1},C_{E2},C_m\geq T_1$ there is $T_2>0$ with the following property for every $C_H\geq T_2$.
Let $p$, $n'$, $m$, $G'$, $N_1$, $c_{cl}$, $D_p$, $c_*$, $T$, $Y$, $Y_b$, $b'$, $b_l$, $b_u$, $a$, $c$, $a_T$ and $b_T$
satisfy the hypotheses of Lemma~\ref{pd-kappa-defects-digits-count-upper} with these constants, and form the numbers
of Definition~\ref{pd-kappa-defects-numbers}. Let $1\leq j\leq q$ and digits $\rho,\rho'$ for $j$ be given, and let
$\hat b_i$ and $\hat\alpha_i$ ($1\leq i\leq s=S_j$) be the numbers of Definition~\ref{pd-kappa-defects-segment-lattice}(1)
for them, with $\alpha_i=a_T(y_{j,i})$; by hypothesis (1) of Theorem~\ref{pd-gain-theorem-one}, $s\leq N_1<m$. Let
$A\geq0$ be an integer with $\hat a_j(1-N_1\Lambda/k)\leq A\leq\hat a_j$, let $(\nu,\lambda)$ be a pair for $p$, $\Lambda$,
$m$, $s$, $A$ and the $\hat b_i$ and $\hat\alpha_i$ that is not kept, with $t_i$ as in
Definition~\ref{pd-kappa-defects-digits-count-pairs}, and write $x^+=\max(x,0)$. Then there is a real number $\theta\geq0$
with $\theta^{m-1}\leq e^{2-\Lambda/5}$ and
\[
 (\lambda_1^++m)\prod_{i=1}^s(t_i^++m)\prod_{i=1}^s\big((-\lambda_{m-s+i})^++m\big)\leq
 \theta\,A\prod_{i=1}^s\hat b_i\prod_{i=1}^s\hat\alpha_i .
\]
\end{lemma}
\leanref{Formalism/SharpRestriction/GrowthBounds/BinomialRatio/NotKept.lean}{Formalism.MainNotions.ColumnSystems.Pf.BracketAt}
\leanref{Formalism/SharpRestriction/GrowthBounds/BinomialRatio/NotKept.lean}{Formalism.MainNotions.ColumnSystems.Pf.RestBracket}
\leanref{Formalism/SharpRestriction/GrowthBounds/BinomialRatio/Bracket.lean}{Formalism.MainNotions.ColumnSystems.Pf.restBracket_of}
\leanref{Formalism/SharpRestriction/GrowthBounds/BinomialRatio/Bracket.lean}{Formalism.MainNotions.ColumnSystems.Pf.bracket_of}

\begin{proof}
\emph{The constants.} Let $T_0=T_1=1$, and for $C_{E1},C_{E2}\geq1$ let $T_2=4000(C_{E1}+C_{E2}+1)(C_c+1)$. Let
$C_H\geq T_2$, let the data and the pair be as in the statement, and write $\hat a=\hat a_j$. The constants satisfy the
hypotheses of Lemmas~\ref{pd-kappa-defects-digits-count-bracket-facts} and~\ref{pd-kappa-defects-digits-count-bracket-large}.
Let $c_l$, $\bar o_i$ and $n_i$ be as in Definition~\ref{pd-kappa-defects-digits-count-pairs}. By hypothesis (1) of
Theorem~\ref{pd-gain-theorem-one}, $2s+1\leq2N_1+2\leq m$, and by
Lemma~\ref{pd-kappa-defects-digits-count-bracket-facts}(4), $n_i=-\lambda_{m-s+i}$.

\emph{The bound.} We apply Lemma~\ref{pd-kappa-defects-digits-count-bracket-core} with $l_0=\lambda_1$ and
$u_i=\bar o_i-t_i$, and with $\Lambda$, $p$, $A$, the $\hat b_i$, $\hat\alpha_i$, $t_i$, $n_i$ and $c_l$. Its hypothesis (1) is
Lemma~\ref{pd-kappa-defects-digits-count-bracket-large}. In (2), $c_0=A$ and $c_l\geq0$ by
Lemma~\ref{pd-kappa-defects-digits-count-bracket-facts}(1), and $c_i=c_{i-1}+\hat b_i-\bar o_i$ is the definition of
$\bar o_i$. Hypotheses (3) and (4) are parts (2), (3) and (4) of
Lemma~\ref{pd-kappa-defects-digits-count-bracket-facts}. As the pair is not kept, one of the two conditions of
Definition~\ref{pd-kappa-defects-digits-count-pairs}(2) fails at some $i_0$; with $k=m+1$, this is (5). The lemma gives
$\theta$. Finally $\lambda_1$, the $t_i$ and the $n_i$ are $\geq0$, so they are $\lambda_1^+$, $t_i^+$ and
$(-\lambda_{m-s+i})^+$.
\end{proof}
% Lean: RestBracket = restBracket_of (restPairFacts_of restNUp) coarse_degrees_large_holds, TransportedAlt.mono of
%   the conjunction of RestPairFacts (thresholds 1, 1 and T2 CE1 CE2 Cc 1) and CoarseDegreesLargeBeyond (1, 1 and
%   T2 CE1 CE2 Cc 4), whose maxima are the thresholds above. bracket_of: hF (Lemma
%   pd-kappa-defects-digits-count-bracket-facts), hN (Lemma pd-kappa-defects-digits-count-bracket-large), hroom
%   from HIext, hk' from not KeptPair (not_forall), and brk_core (Lemma pd-kappa-defects-digits-count-bracket-core)
%   with rho = pairRho, u = pairO - pairRho, nn = pairN and cc = pairC; nk_toNat_cast for the positive parts.

\begin{lemma}\label{pd-kappa-defects-digits-count-not-kept}
For every $C_c\geq0$ there is $T_0>0$ such that for every $C_0\geq T_0$ there is $T_1>0$ such that for
all $C_{E1},C_{E2},C_m\geq T_1$ there is $T_2>0$ with the following property for every $C_H\geq T_2$.
Let $p$, $n'$, $m$, $G'$, $N_1$, $c_{cl}$, $D_p$, $c_*$, $T$, $Y$, $Y_b$, $b'$, $b_l$, $b_u$, $a$, $c$, $a_T$ and $b_T$
satisfy the hypotheses of Lemma~\ref{pd-kappa-defects-digits-count-upper} with these constants, and form the numbers
of Definition~\ref{pd-kappa-defects-numbers}. Let $1\leq j\leq q$ and digits $\rho,\rho'$ for $j$ be given, and let
$\hat b_i$ and $\hat\alpha_i$ ($1\leq i\leq S_j$) be the numbers of Definition~\ref{pd-kappa-defects-segment-lattice}(1)
for them, with $\alpha_i=a_T(y_{j,i})$; by hypothesis (1) of Theorem~\ref{pd-gain-theorem-one}, $S_j\leq N_1<m$. Let
$A\geq0$ be an integer with $\hat a_j(1-N_1\Lambda/k)\leq A\leq\hat a_j$, and let $(\nu,\lambda)$ be a pair for $p$,
$\Lambda$, $m$, $S_j$, $A$ and the $\hat b_i$ and $\hat\alpha_i$ that is not kept
(Definition~\ref{pd-kappa-defects-digits-count-pairs}). Then, with $P_m$ as in
Lemma~\ref{pd-kappa-defects-windows-model-segment},
\[
 \operatorname{wd}(\lambda)\leq e^{2-\Lambda/5}(8a+k)^{(2N_1+2)^2}P_m(A)\prod_{i=1}^{S_j}P_m(\hat b_i)P_m(\hat\alpha_i).
\]
The factor $(8a+k)^{(2N_1+2)^2}$ is not used: it is at least $1$, and the proof gives the bound without it.
\end{lemma}
\leanref{Formalism/SharpRestriction/GrowthBounds/BinomialRatio/ThresholdStandardLarge.lean}{Formalism.MainNotions.ColumnSystems.Pf.NotKeptWtAt}
\leanref{Formalism/SharpRestriction/GrowthBounds/BinomialRatio/ThresholdStandardLarge.lean}{Formalism.MainNotions.ColumnSystems.Pf.RestNotKeptWt}
\leanref{Formalism/SharpRestriction/GrowthBounds/BinomialRatio/NotKept.lean}{Formalism.MainNotions.ColumnSystems.Pf.restNotKeptWt_of}

\begin{proof}
Let $T_0$, $T_1$ and $T_2$ be the numbers of Lemma~\ref{pd-kappa-defects-digits-count-bracket}, let the data and the pair
be as in the statement, write $s=S_j$, and let $t_i$ be as in Definition~\ref{pd-kappa-defects-digits-count-pairs},
$x^+=\max(x,0)$ and $P_{m-1}(d)=\binom{d+m-1}{m-1}$. Lemma~\ref{pd-kappa-defects-digits-count-bracket} gives
$\theta\geq0$ with $\theta^{m-1}\leq e^{2-\Lambda/5}$ and
\[
 (\lambda_1^++m)\prod_i(t_i^++m)\prod_i\big((-\lambda_{m-s+i})^++m\big)\leq\theta\,A\prod_i\hat b_i\prod_i\hat\alpha_i .
\]
As $2s+1\leq2N_1+2\leq m$ by hypothesis (1) of Theorem~\ref{pd-gain-theorem-one},
Lemma~\ref{pd-kappa-defects-digits-count-weight} bounds $\operatorname{wd}(\lambda)$ by the product $\Pi'$ of $P_{m-1}$ at
$\lambda_1^+$, the $t_i^+$ and the $(-\lambda_{m-s+i})^+$. For integers $y\geq0$, $P_{m-1}(y)=\prod_{t=1}^{m-1}(y+t)/t$, so
$y^{m-1}/(m-1)!\leq P_{m-1}(y)\leq(y+m)^{m-1}/(m-1)!$. Applying the upper bound to the $2s+1$ factors of $\Pi'$, raising the
displayed inequality to the power $m-1$, and applying the lower bound to $A$, the $\hat b_i$ and the $\hat\alpha_i$,
\[
 \operatorname{wd}(\lambda)\leq\Pi'\leq\theta^{m-1}P_{m-1}(A)\prod_iP_{m-1}(\hat b_i)P_{m-1}(\hat\alpha_i).
\]
Finally $P_{m-1}(y)\leq P_m(y)$, as $\binom{y+m}m=\binom{y+m-1}{m-1}+\binom{y+m-1}m$; $\theta^{m-1}\leq e^{2-\Lambda/5}$; and
$(8a+k)^{(2N_1+2)^2}\geq1$.
\end{proof}
% Lean: RestNotKeptWt = restNotKeptWt_of (wtBound_of hCD) (restBracket_of (restPairFacts_of restNUp)
%   coarse_degrees_large_holds), as in stmtEExt_of, with the thresholds of RestBracket. notKeptWt_of: WtBoundAt
%   (Lemma pd-kappa-defects-digits-count-weight, with hroom from HIext), BracketAt (Lemma
%   pd-kappa-defects-digits-count-bracket), nk_core (nk_b_up and nk_b_low, from Nat.choose_le_pow_div and
%   Nat.pow_le_choose), nk_prod_le_dimMM (nk_sdXi_le_sd) and hbig. The factor (8a + k)^((2 N1 + 2)^2) is not used.

\begin{lemma}\label{pd-kappa-defects-digits-count-tail}
For every $C_c\geq0$ there is $T_0>0$ such that for every $C_0\geq T_0$ there is $T_1>0$ such that for
all $C_{E1},C_{E2},C_m\geq T_1$ there is $T_2>0$ with the following property for every $C_H\geq T_2$.
Let the data, $j$, the digits and the $\hat b_i$ and $\hat\alpha_i$ be as in
Lemma~\ref{pd-kappa-defects-digits-count-not-kept} with these constants, and let $K_A$ be the kept sums for $p$,
$\Lambda$, $m$, $S_j$ and the $\hat b_i$ and $\hat\alpha_i$ (Definition~\ref{pd-kappa-defects-digits-count-pairs}). If
$A\geq0$ is an integer with $\hat a_j(1-N_1\Lambda/k)\leq A\leq\hat a_j$, then
\[
 \big(1-e^{-\Lambda/8}\big)P_m(A)\prod_{i=1}^{S_j}P_m(\hat b_i)P_m(\hat\alpha_i)\leq K_A .
\]
\end{lemma}
\leanref{Formalism/SharpRestriction/GrowthBounds/BinomialRatio/CountLow.lean}{Formalism.MainNotions.ColumnSystems.Pf.TailAt}
\leanref{Formalism/SharpRestriction/GrowthBounds/BinomialRatio/CountLow.lean}{Formalism.MainNotions.ColumnSystems.Pf.RestTail}
\leanref{Formalism/SharpRestriction/GrowthBounds/BinomialRatio/ThresholdStandardLarge.lean}{Formalism.MainNotions.ColumnSystems.Pf.restTail_of}
\leanref{Formalism/SharpRestriction/GrowthBounds/BinomialRatio/ThresholdStandardLarge.lean}{Formalism.MainNotions.ColumnSystems.Pf.tail_of}
\leanref{Formalism/SharpRestriction/GrowthBounds/BinomialRatio/ThresholdStandardLarge.lean}{Formalism.MainNotions.ColumnSystems.Pf.tail_core}
\leanref{Formalism/SharpRestriction/GrowthBounds/BinomialRatio/ThresholdStandardLarge.lean}{Formalism.MainNotions.ColumnSystems.Pf.tail_num}

\begin{proof}
\emph{The constants.} Let $T_0$, $T_1$ and $T_2$ be the largest of the numbers $T_0$, $T_1$ and $T_2$ of
Lemma~\ref{pd-kappa-defects-digits-count-not-kept} and of $880$, $1$ and $1000(C_{E1}+C_{E2}+1)(C_c+1)$, respectively,
and let the data and $A$ be as in the statement. Write $s=S_j$, $\Pi=\prod_iP_m(\hat b_i)P_m(\hat\alpha_i)$,
$M=P_m(A)\Pi$, $W=e^{2-\Lambda/5}(8a+k)^{(2N_1+2)^2}$ and $N_p=(8a)^{4(N_1+1)^2}$, and let $F$ be the finite set of
the pairs.

\emph{The total.} The sum of $\operatorname{wd}(\lambda)$ over $F$ is $M$. Indeed, for a path $\nu$, the sum of
$\operatorname{wd}(\lambda)$ over the $\lambda$ with $(\nu,\lambda)\in F$ is $\operatorname{wd}(\nu_{2s+1})$ by
Lemma~\ref{pd-kappa-defects-digits-branching}; and by Lemma~\ref{pd-kappa-defects-digits-chain-dimension} with $q'=0$,
as in the proof of Lemma~\ref{pd-kappa-defects-digits-count-upper}, the sum of $\operatorname{wd}(\nu_{2s+1})$ over the
paths $\nu$ is the dimension of the space of polynomials of these row degrees in $m+1$ variables per letter, which is
$P_m(A)\Pi$.

\emph{The bound.} For a kept pair, $\operatorname{wd}(\lambda)-WM\leq\operatorname{wd}(\lambda)$, as $W,M\geq0$; for a
pair that is not kept, $\operatorname{wd}(\lambda)-WM\leq0$ by Lemma~\ref{pd-kappa-defects-digits-count-not-kept}.
Summing over $F$ gives $M-|F|WM\leq K_A$. By Lemma~\ref{pd-kappa-defects-digits-count-pair-bound}, whose hypotheses
on the constants hold, $|F|\leq N_p$. So it suffices to show $N_pW\leq e^{-\Lambda/8}$.

\emph{The numbers.} The facts (a) to (g) in the proof of Lemma~\ref{pd-kappa-defects-margin} hold. They give $k\geq5$,
$N_1\geq1$, $\log a\geq1$, $b'\geq k^3Q\geq125$, $b'\leq a$ and $k\leq kb'\leq a$; so $a\geq9$. Put $L=\log a$, so that
$e^{2L}=a^2$. As $a\geq9$ and $k\leq a$, $8a\leq8a+k\leq9a\leq e^{2L}$; and $4(N_1+1)^2=(2N_1+2)^2\leq16N_1^2$. So
$(8a)^{4(N_1+1)^2}$ and $(8a+k)^{(2N_1+2)^2}$ are at most $e^{32N_1^2L}$, and $N_pW\leq e^{64N_1^2L+2-\Lambda/5}$, with
$\Lambda=C_0N_1^2L$. As $C_0\geq880$ and $N_1^2L\geq1$, $64N_1^2L+2\leq66N_1^2L\leq\frac3{40}\Lambda=\frac\Lambda5-\frac\Lambda8$.
Hence $N_pW\leq e^{-\Lambda/8}$.
\end{proof}
% Lean: RestTail = restTail_of hCD hNB restPairCount (restNotKeptWt_of ...), as in stmtEExt_of: the thresholds
%   are those of threshold_standard_large (880, 1 and T2 CE1 CE2 Cc 1), RestPairCount and RestNotKeptWt.
%   tail_of: tail_core with Np and W as above (pair_sum from ChainDimAt and NumBranchAt, kept_sum), and tail_num
%   (hb9, hka from kQ, kb and ba; loga, N1pos). The sets of the proof are pairSet and its toFinset.

\begin{lemma}\label{pd-kappa-defects-digits-count-kept}
Let $C_c\geq0$, $C_0\geq1$, $C_{E1},C_{E2}\geq1$ and $C_H\geq1000(C_{E1}+C_{E2}+1)(C_c+1)$. Let the data, $j$, the digits,
$s$, $\eta$, $\hat a$, the $\hat b_i$ and the $\hat\alpha_i$ be as in Lemma~\ref{pd-kappa-defects-digits-count-coarse}, and
put $b_+=b_T(s_{j+1})$. Let $\mu\geq0$ be a real number with $n=\lceil\mu\rceil\leq\hat aN_1\Lambda/k$.
\begin{enumerate}
\item $n\leq\hat a$, and $A=\hat a-n$ satisfies $\hat a(1-N_1\Lambda/k)\leq A\leq\hat a$.
\item Let $(\nu,\lambda)$ be a kept pair for $p$, $\Lambda$, $m$, $s$, $A$ and the $\hat b_i$ and $\hat\alpha_i$
  (Definition~\ref{pd-kappa-defects-digits-count-pairs}), and let $j_+$ be an integer with $0\leq j_+\leq\Lambda b_+/k$.
  Let $c_l$, $\bar o_i$, $\kappa_l$, $\bar w_i$ and $\bar\varsigma_i$ be as in
  Lemma~\ref{pd-kappa-defects-digits-count-pair-facts}. Let $b$ be the block with the vectors $\rho_i$ and $\rho'_i$, with
  $(o_i,w_i,\varsigma_i)=(0,\bar w_i,\bar\varsigma_i)$ for $i\leq\eta$ and $(o_i,w_i,\varsigma_i)=(\bar o_i,0,0)$ for $i>\eta$,
  and with $o_+=|\lambda|-|\kappa_{2s}|+j_+$; these entries are $\geq0$. Let $\ell$ be the level
  $(\bar\varsigma_{u_1},\dots,\bar\varsigma_{u_{s-\eta}},j_+)$, with $u_l=\eta+l$. Then $b$ is a kept block of $S_j(b_T)$, for
  $\alpha_i=a_T(y_{j,i})$, $\mu$ and $r=m$ (Definition~\ref{pd-kappa-defects-segment-lattice}); $\ell$ is a level of $b$;
  $\kappa=(\kappa_0,\dots,\kappa_{2s+1})$ is a path of rank $m$ for the letters $x_1,y_1,\dots,x_s,y_s,x_+$ with the
  degrees $\delta'$ of $\ell$; and $\kappa_J(1)\leq M$ for each of the pairs $(J,M)$ of
  Definition~\ref{pd-kappa-defects-digits-chains}(3).
\end{enumerate}
\end{lemma}
\leanref{Formalism/SharpRestriction/GrowthBounds/BinomialRatio/CountLow.lean}{Formalism.MainNotions.ColumnSystems.Pf.KeptBlockAt}
\leanref{Formalism/SharpRestriction/GrowthBounds/BinomialRatio/CountLow.lean}{Formalism.MainNotions.ColumnSystems.Pf.KeptBlocksBeyondThreshold}
\leanref{Formalism/SharpRestriction/GrowthBounds/BinomialRatio/CountLow.lean}{Formalism.MainNotions.ColumnSystems.Pf.pairMap}
\leanref{Formalism/SharpRestriction/GrowthBounds/BinomialRatio/KeptBlocksHolds.lean}{Formalism.MainNotions.ColumnSystems.Pf.kept_blocks_holds}
\leanref{Formalism/SharpRestriction/GrowthBounds/BinomialRatio/KeptBlocksOfSLine.lean}{Formalism.MainNotions.ColumnSystems.Pf.pairU_at}
\leanref{Formalism/SharpRestriction/GrowthBounds/BinomialRatio/KeptBlocksHolds.lean}{Formalism.MainNotions.ColumnSystems.Pf.pairS_at}
\leanref{Formalism/SharpRestriction/GrowthBounds/BinomialRatio/PairFacts.lean}{Formalism.MainNotions.ColumnSystems.Pf.pf_X}
\leanref{Formalism/SharpRestriction/GrowthBounds/BinomialRatio/CoarseDegreeRatioConsecutive.lean}{Formalism.MainNotions.ColumnSystems.Pf.kept_pair_of_three_lines}
\leanref{Formalism/SharpRestriction/GrowthBounds/BinomialRatio/CoarseDegreeRatioConsecutive.lean}{Formalism.MainNotions.ColumnSystems.Pf.pair_level}
\leanref{Formalism/SharpRestriction/GrowthBounds/BinomialRatio/CoarseDegreeRatioConsecutive.lean}{Formalism.MainNotions.ColumnSystems.Pf.pair_chain}
\leanref{Formalism/SharpRestriction/GrowthBounds/BinomialRatio/CoarseDegreeRatioConsecutive.lean}{Formalism.MainNotions.ColumnSystems.Pf.pair_thr}

\begin{proof}
Write $\varepsilon=N_1\Lambda/k$ and let $\delta$ be as in Lemma~\ref{pd-kappa-defects-digits-count-coarse}.

\emph{Facts.} The facts in the proof of Lemma~\ref{pd-kappa-defects-digits-count-coarse} hold. They give $k\geq5$,
$N_1\geq1$, $2N_1+3\leq k$, $s\leq N_1$, $\Lambda\geq1$, $k^2\leq p$, $40\varepsilon\leq\delta<1/6$ and $\hat a>0$; and
$\Lambda\leq k$, by the fact (d) in the proof of Lemma~\ref{pd-kappa-defects-margin}. Hence $\varepsilon\leq1$, $s<m$, and
$\Lambda/k\leq20\varepsilon\leq\delta/2$; and $4/p^2\leq4/k^4\leq20/k\leq\delta/2$ and
$12N_1/p^2\leq12/k^3\leq1/(2k)\leq\delta/80$, as $40/k\leq40\varepsilon\leq\delta$ and $N_1\leq k$.

(1) $n\leq\hat a\varepsilon\leq\hat a$, so $A=\hat a-n\geq\hat a(1-\varepsilon)$ and $A\leq\hat a$.

(2) Let $t_i$ be as in Definition~\ref{pd-kappa-defects-digits-count-pairs}. By
Lemma~\ref{pd-kappa-defects-digits-count-pair-facts}(2) and (3), $\bar w_i\geq0$, $\bar\varsigma_i\geq0$,
$\bar o_i=\bar w_i+\bar\varsigma_i\geq0$ and $|\lambda|-|\kappa_{2s}|\geq0$, so the entries of $b$ are $\geq0$.

\emph{Step 1: the indices $i>\eta$.} Let $\eta<i\leq s$. We show $\bar o_i\leq\hat b_i$ and
$\bar o_i\leq c_{i-1}+\Lambda\hat b_i/k$. By Lemma~\ref{pd-kappa-defects-digits-count-pair-facts}(5) and
Lemma~\ref{pd-kappa-defects-digits-count-coarse}(3) and (2),
$c_{i-1}\leq A+\sum_{l<i}\hat b_l\leq\hat a+4\hat b_i/p^2\leq\hat b_i(1-\delta)+\hat b_i\delta/2=\hat b_i(1-\delta/2)$. So
$\max(c_{i-1},\hat b_i)=\hat b_i$, and the first condition for a kept pair, with $t_i\leq c_{i-1}$
(Lemma~\ref{pd-kappa-defects-digits-count-pair-facts}(2)), gives $\bar o_i\leq t_i+\Lambda\hat b_i/k\leq
c_{i-1}+\Lambda\hat b_i/k$. As $\Lambda/k\leq\delta/2$, also $\bar o_i\leq\hat b_i(1-\delta/2)+\hat b_i\delta/2=\hat b_i$.

\emph{Step 2: the indices $i\leq\eta$.} Let $1\leq i\leq\eta$. We show $\bar o_i\leq\hat b_i(1+4/p^2)+\Lambda c_{i-1}/k$ and
$\bar\varsigma_i\leq2\Lambda\hat a/k$. Put $\delta_0=3\hat a(4/p^2+\Lambda/k)$.
\begin{enumerate}
\item[(a)] $\hat b_l\leq\hat a$ for $l\leq i$, by Lemma~\ref{pd-kappa-defects-digits-count-coarse}(2).
\item[(b)] $\nu_{2l-1}(2)\leq4\hat b_l/p^2$ for $1\leq l\leq s$: by
  Lemma~\ref{pd-kappa-defects-digits-count-pair-facts}(4) and (2),
  $\nu_{2l-1}(2)\leq\sum_{i'<l}\bar w_{i'}\leq\sum_{i'<l}\hat b_{i'}$, and
  Lemma~\ref{pd-kappa-defects-digits-count-coarse}(3) applies.
\item[(c)] $t_l\leq\hat b_l(1+4/p^2)$ for $1\leq l\leq s$, as $t_l\leq\nu_{2l-1}(2)+\bar w_l$
  (Lemma~\ref{pd-kappa-defects-digits-count-pair-facts}(4)) and $\bar w_l\leq\hat b_l$.
\item[(d)] $c_l\leq2\hat a$ for $0\leq l\leq i-1$: by Lemma~\ref{pd-kappa-defects-digits-count-pair-facts}(5),
  Lemma~\ref{pd-kappa-defects-digits-count-coarse}(3) and (a), $c_l\leq A+\sum_{i'\leq l}\hat b_{i'}\leq\hat a+4\hat b_{l+1}/p^2
  \leq2\hat a$.
\item[(e)] $A-l\delta_0\leq c_l$ for $0\leq l\leq i-1$. For $l=0$ this is $c_0=A$. If it holds for $l<i-1$, then the first
  condition for a kept pair at $l+1$, with (c), (d) and (a), gives
  $\bar o_{l+1}\leq t_{l+1}+2\Lambda\hat a/k\leq\hat b_{l+1}+4\hat b_{l+1}/p^2+2\Lambda\hat a/k\leq\hat b_{l+1}+\delta_0$; so
  $c_{l+1}=c_l+\hat b_{l+1}-\bar o_{l+1}\geq c_l-\delta_0$.
\item[(f)] $\hat b_i\leq c_{i-1}$. As $i-1<s\leq N_1$ and $\delta_0\geq0$, (e) gives
  $c_{i-1}\geq A-N_1\delta_0\geq\hat a(1-\varepsilon')$ with $\varepsilon'=4\varepsilon+12N_1/p^2$. By the facts,
  $\varepsilon'\leq\delta/10+\delta/80$, so $\varepsilon'(1+\delta)\leq\delta$ and $(1+\delta)(1-\varepsilon')\geq1$. As
  $\hat b_i(1+\delta)<\hat a$, $\hat b_i\leq\hat b_i(1+\delta)(1-\varepsilon')\leq\hat a(1-\varepsilon')\leq c_{i-1}$.
\end{enumerate}
So $\max(c_{i-1},\hat b_i)=c_{i-1}$, and the first condition for a kept pair gives $\bar o_i-t_i\leq\Lambda c_{i-1}/k-
4(\hat b_i+m)/p^2$. With (c), $\bar o_i\leq\hat b_i(1+4/p^2)+\Lambda c_{i-1}/k$. By
Lemma~\ref{pd-kappa-defects-digits-count-pair-facts}(4), $\bar\varsigma_i-\nu_{2i-1}(2)=\bar o_i-\bar w_i-\nu_{2i-1}(2)\leq
\bar o_i-t_i$; with (b), $\bar\varsigma_i\leq\Lambda c_{i-1}/k\leq2\Lambda\hat a/k$, by (d).

\emph{Step 3: the last letter.} As $s<m$, $n_i=-\lambda_{m-s+i}$, and
Lemma~\ref{pd-kappa-defects-digits-count-pair-facts}(6) and the second condition for a kept pair give
$|\lambda|-|\kappa_{2s}|\leq c_s+\sum_i\hat\alpha_i-(1-\Lambda/k)\sum_i\hat\alpha_i=c_s+(\Lambda/k)\sum_i\hat\alpha_i$.

\emph{Step 4: the capacities.} For every $i$ the number $o^\flat_i$ of $b$ is $\bar o_i$: for $i\leq\eta$ it is
$w_i+\varsigma_i=\bar w_i+\bar\varsigma_i$, and for $i>\eta$ it is $o_i$. As $c_i-c_{i-1}=\hat b_i-\bar o_i$ and $c_0=A$, the
capacities of $b$ are $\operatorname{cap}(l)=\hat a+c_l-A=c_l+n\geq c_l$.

\emph{The block is kept.} Its vectors satisfy the first condition, as $\rho,\rho'$ are digits. For $i>\eta$, $o_i=\bar o_i\leq
\hat b_i$, $o_i\leq\operatorname{cap}(i-1)+\Lambda\hat b_i/k$ by Steps 1 and 4, and $w_i=\varsigma_i=0$. For $i\leq\eta$, $o_i=0$,
$w_i=\bar w_i\leq\hat b_i$ and $\varsigma_i=\bar\varsigma_i\leq\hat\alpha_i$
(Lemma~\ref{pd-kappa-defects-digits-count-pair-facts}(2)), $w_i+\varsigma_i=\bar o_i\leq(1+4/p^2)\hat b_i+
\Lambda\operatorname{cap}(i-1)/k$ and $\varsigma_i\leq2\Lambda\hat a/k$ by Steps 2 and 4. Finally, by Step 3,
$j_+\leq\Lambda b_+/k$, Step 4 and Lemma~\ref{pd-kappa-defects-digits-count-coarse}(1),
\[
 o_+\leq c_s+\frac\Lambda k\sum_i\hat\alpha_i+\frac{\Lambda b_+}k\leq\operatorname{cap}(s)+\Big(\frac\Lambda k+\frac4{p^2}\Big)
 \sum_i\frac{2a}{P_i}+\frac{\Lambda b_+}k=O_+ .
\]

\emph{The level.} The degrees of $b$ are $\delta(x_i)=\bar w_i$ and $\delta(y_i)=\hat\alpha_i-\bar\varsigma_i$ for $i\leq\eta$,
$\delta(x_i)=\bar o_i$ and $\delta(y_i)=\hat\alpha_i$ for $i>\eta$, and $\delta(x_+)=o_+$. For $i=u_l$,
$\bar\varsigma_i\leq\bar o_i$ and $\bar\varsigma_i\leq\hat\alpha_i$; and $j_+\leq o_+$, while $j_+\leq\Lambda b_+/k\leq b_+$ as
$\Lambda\leq k$. So $\ell$ is a level of $b$, and its degrees are $\delta'(x_i)=\bar w_i$ and
$\delta'(y_i)=\hat\alpha_i-\bar\varsigma_i$ for every $i$, and $\delta'(x_+)=|\lambda|-|\kappa_{2s}|$.

\emph{The path.} $\kappa_0=0$, and by Lemma~\ref{pd-kappa-defects-digits-count-pair-facts}(2) and (3), at the vector
letter $x_i$ the vector $\kappa_{2i-1}$ is a strip of size $\bar w_i$ over $\kappa_{2i-2}$, at the covector letter $y_i$ the
vector $\kappa_{2i-1}$ is a strip of size $\hat\alpha_i-\bar\varsigma_i$ over $\kappa_{2i}$, and at the vector letter $x_+$ the
vector $\lambda=\kappa_{2s+1}$ is a strip of size $|\lambda|-|\kappa_{2s}|$ over $\kappa_{2s}$. The strip conditions make
each $\kappa_l$ non-increasing. So $\kappa$ is a path for the degrees $\delta'$.

\emph{The thresholds.} The pairs $(J,M)$ are $(2u_l-1,\lfloor\operatorname{cap}(u_l-1)-\mu\rfloor)$, as
$\mathcal L_l$ has $2u_l-1$ letters, and $(2s+1,\lfloor\operatorname{cap}(s)-\mu\rfloor)$. By Step 4, $\operatorname{cap}(l)-\mu
=c_l+n-\mu\geq c_l$, and $c_l$ is an integer, so $c_l\leq\lfloor\operatorname{cap}(l)-\mu\rfloor$. Now
$\kappa_{2u_l-1}(1)=t_{u_l}\leq c_{u_l-1}$ and $\kappa_{2s+1}(1)=\lambda_1\leq c_s$, by
Lemma~\ref{pd-kappa-defects-digits-count-pair-facts}(2) and (3).
\end{proof}
% Lean: KeptBlocksBeyondThreshold = kept_blocks_holds = kept_blocks_of_S_line restPairS =
%   kept_blocks_of_three_lines restPairU restPairS restPairX = kept_blocks_of_parts (kept_pairs_of_three_lines ...)
%   kept_levels_hold kept_thresholds_hold, with kept_chain_holds; all thresholds are 1, 1 and T2 CE1 CE2 Cc 1
%   (std_of with H = 1). The block, level and path are pairMap (pairBlock, pairLevel, pairChain) of the target
%   (mu, lam, jp), with the toNat of the entries. (1) is the hle and hA1, hA2 steps of kept_pair_of_three_lines
%   (LaK, de6). Step 1 is pairU_at (coarse_degree_alternative, second part; coarse_degree_prefix_sum;
%   pairC_le_sum; hp2, hLk). Step 2 is pairS_at: (a) hAl, (b) hkap, (c) hrho, (d) hcup, (e) hlow, (f) hci_ge (hq1,
%   hq2, heps), and the conclusion with pf_u_ge. Step 3 is pf_X through restPairX (s + 1 <= m from N1k and segB).
%   Step 4 is oEff_pair and cap_pair. The block is kept by kept_pair_of_three_lines (hcap; the bound for o_+ with
%   hC and hq). The level is pair_level (segment_chain_degree_even, chain_deg_E, chain_deg_U), the path is
%   pair_chain (levU, levW, levE, chain_first_tail, inter_strip), and the thresholds are pair_thr (cap_pair,
%   pair_top, inter_strip).

\begin{lemma}\label{pd-kappa-defects-digits-count-match}
Let $C_c\geq0$, $C_0\geq1$, $C_{E1},C_{E2}\geq1$ and $C_H\geq1000(C_{E1}+C_{E2}+1)(C_c+1)$. Let the data, $j$, the digits,
$s$, $\hat a$, the $\hat b_i$ and the $\hat\alpha_i$ be as in Lemma~\ref{pd-kappa-defects-digits-count-coarse}, put
$b_+=b_T(s_{j+1})$ and $J=\lfloor\Lambda b_+/k\rfloor$, and let $\mu\geq0$ be a real number with
$n=\lceil\mu\rceil\leq\hat aN_1\Lambda/k$. Let $Z$ be the set of the triples $(b,\ell,\kappa)$ of a kept block $b$ of
$S_j(b_T)$, formed with $\alpha_i=a_T(y_{j,i})$, $\alpha_+=a_T(s_{j+1})$, $\mu$ and $r=m$
(Definition~\ref{pd-kappa-defects-segment-lattice}), whose vectors are the $\rho_i$ and $\rho'_i$, a level $\ell$ of $b$,
and a path $\kappa=(\kappa_0,\dots,\kappa_{2s+1})$ of rank $m$ for the letters $x_1,y_1,\dots,x_s,y_s,x_+$ with the degrees
$\delta'$ of $\ell$ such that $\kappa_{J'}(1)\leq M'$ for each of the pairs $(J',M')$ of
Definition~\ref{pd-kappa-defects-digits-chains}(3). Suppose that there is an injective map $\Phi$ from $Z$ to the set of
the triples $(\nu,\lambda,j_+)$ of a pair $(\nu,\lambda)$ for $p$, $\Lambda$, $m$, $s$, $\hat a$ and the $\hat b_i$ and
$\hat\alpha_i$ (Definition~\ref{pd-kappa-defects-digits-count-pairs}) and an integer $j_+\geq0$, such that
$\Phi(b,\ell,\kappa)$ has $\lambda=\kappa_{2s+1}$ and $j_+$ equal to the last entry of $\ell$. Let $K$ be the kept sum
$K_{\hat a-n}$ for $p$, $\Lambda$, $m$, $s$ and the $\hat b_i$ and $\hat\alpha_i$. Then
\[
 K\sum_{j_+=0}^J\binom{b_+-j_++m-1}{m-1}\leq c_j(\mu;\rho,\rho')
\]
(Definition~\ref{pd-kappa-defects-digits-chains}(4)). The hypothesis on $\Phi$ is used only to show that $Z$ is finite.
\end{lemma}
\leanref{Formalism/SharpRestriction/GrowthBounds/BinomialRatio/CountLow.lean}{Formalism.MainNotions.ColumnSystems.Pf.MatchAt}
\leanref{Formalism/SharpRestriction/GrowthBounds/BinomialRatio/CountLow.lean}{Formalism.MainNotions.ColumnSystems.Pf.MatchLowAt}
\leanref{Formalism/SharpRestriction/GrowthBounds/BinomialRatio/CountLow.lean}{Formalism.MainNotions.ColumnSystems.Pf.match_of}
\leanref{Formalism/SharpRestriction/GrowthBounds/BinomialRatio/CountLow.lean}{Formalism.MainNotions.ColumnSystems.Pf.matchLow_of}
\leanref{Formalism/SharpRestriction/GrowthBounds/BinomialRatio/CountLow.lean}{Formalism.MainNotions.ColumnSystems.Pf.cl_chain_ext}
\leanref{Formalism/SharpRestriction/GrowthBounds/BinomialRatio/CountLow.lean}{Formalism.MainNotions.ColumnSystems.Pf.cl_prod_le}

\begin{proof}
By Lemma~\ref{pd-kappa-defects-digits-count-kept}(1), $A=\hat a-n\geq0$; and $J\leq b_+$, as $\Lambda\leq k$ by the fact
(d) in the proof of Lemma~\ref{pd-kappa-defects-margin}.

\emph{The sources.} For $z=(b,\ell,\kappa)\in Z$ let $j_+$ be the last entry of $\ell$. By
Definition~\ref{pd-kappa-defects-digits-chains}(3) and (4),
\[
 c_j(\mu;\rho,\rho')=\sum_{z\in Z}\operatorname{wd}(\kappa_{2s+1})\binom{b_+-j_++m-1}{m-1},
\]
a sum of terms $\geq0$, as $\kappa_{2s+1}$ is non-increasing. The set $Z$ is finite: $\Phi$ maps it injectively to
triples $(\nu,\lambda,j_+)$ with $j_+\leq b_+$, as a level has $j_+\leq b_+$, and there are finitely many pairs
(Definition~\ref{pd-kappa-defects-digits-count-pairs}(1)). This is the only use of $\Phi$; the finiteness of $Z$ also
follows from that of the kept blocks, their levels and the paths.

\emph{The map.} Let $\mathcal T$ be the set of the triples $t=(\nu,\lambda,j_+)$ of a kept pair $(\nu,\lambda)$ for $A$ and an
integer $0\leq j_+\leq J$. For $t\in\mathcal T$ let $\Psi(t)=(b,\ell,\kappa)$ be the block, the level and the path of
Lemma~\ref{pd-kappa-defects-digits-count-kept}(2). By that lemma, $\Psi(t)\in Z$, and its term in the sum is
$\operatorname{wd}(\lambda)\binom{b_+-j_++m-1}{m-1}$, as $\kappa_{2s+1}=\lambda$ and the last entry of $\ell$ is $j_+$.

\emph{Injectivity.} Let $\Psi(t)=\Psi(t')$ with $t=(\nu,\lambda,j_+)$ and $t'=(\nu',\lambda',j'_+)$. The last entries of the
levels give $j_+=j'_+$, and the paths give $\lambda=\lambda'$ and that $\nu_l$ and $\nu'_l$ agree outside their first entries
for $1\leq l\leq2s+1$. The sum of the entries of $\nu_l$ is the same for all paths for the given degrees: it is $0$ for
$l=0$, and from $\nu_l$ to $\nu_{l+1}$ it changes by the degree of the letter $l$, with the sign $+$ at a vector letter and
$-$ at a covector letter. So the first entries agree too, and $t=t'$.

\emph{The bound.} Expanding the product,
\[
 K\sum_{j_+=0}^J\binom{b_+-j_++m-1}{m-1}=\sum_{t\in\mathcal T}\operatorname{wd}(\lambda)\binom{b_+-j_++m-1}{m-1},
\]
which is the sum of the terms of the distinct elements $\Psi(t)$ of $Z$, and hence at most $c_j(\mu;\rho,\rho')$.
\end{proof}
% Lean: RestMatch = restMatch_of hCD hMU (restMatchLow_of kept_blocks_holds), as in stmtEExt_of. match_of:
%   chainCount_eq (the sources, Fs), cl_prod_le with the injection t |-> pairMap t (MatchLowAt, from matchLow_of
%   with KeptBlockAt, Lemma pd-kappa-defects-digits-count-kept), wdim_pos_of_antitone with antitone_lam_of_inter,
%   and cl_Fs_nonneg. Injectivity is cl_chain_ext (pairChain, cl_chain_total) with the level. The hypothesis on
%   Phi is match_of's MatchUpAt (IsChain (mmLetters ...), Inter, lam = the last vector, the same jp), which
%   RestMatchUp = restMatchUp_of supplies from the map of Lemma pd-kappa-defects-digits-count-upper; match_of
%   uses it only in cl_Fs_finite (finsum_le_of_inj into the finite support Gt_support_finite).

\begin{lemma}\label{pd-kappa-defects-digits-count-lower}
For every $C_c\geq0$ there is $T_0>0$ such that for every $C_0\geq T_0$ there is $T_1>0$ such that for
all $C_{E1},C_{E2},C_m\geq T_1$ there is $T_2>0$ with the following property for every $C_H\geq T_2$.
Let $p$, $n'$, $m$, $G'$, $N_1$, $c_{cl}$, $D_p$, $c_*$, $T$, $Y$, $Y_b$, $b'$, $b_l$, $b_u$, $a$, $c$, $a_T$ and $b_T$
satisfy the hypotheses of Lemma~\ref{pd-kappa-defects-digits-count-upper} with these constants, and form
the numbers of Definition~\ref{pd-kappa-defects-numbers}. Let $1\leq j\leq q$, $\mu\geq0$ and digits $\rho,\rho'$
for $j$ be given, and put $\hat a_j=\lfloor(a-E_0)/p^{g_j}\rfloor$ and $n=\lceil\mu\rceil$. If
$n\leq\hat a_jN_1\Lambda/k$, then
\[
 \big(1-e^{-\Lambda/8}\big)\big(1-e^{-\Lambda/3}\big)\Big(1-\frac{mn}{\hat a_j}\Big)d_j(\rho,\rho')
 \leq c_j(\mu;\rho,\rho').
\]
\end{lemma}
\leanref{Formalism/SharpRestriction/GrowthBounds/BinomialRatio/CountDefs.lean}{Formalism.MainNotions.ColumnSystems.Pf.CountLowAt}
\leanref{Formalism/SharpRestriction/GrowthBounds/BinomialRatio/CountDefs.lean}{Formalism.MainNotions.ColumnSystems.Pf.RestCountLow}

\begin{proof}
\emph{The constants.} Let $T_0$, $T_1$ and $T_2$ be the largest of the numbers $T_0$, $T_1$ and $T_2$ of
Lemma~\ref{pd-kappa-defects-digits-count-tail} and of $1$, $1$ and $1000(C_{E1}+C_{E2}+1)(C_c+1)$, respectively. The
constants then satisfy the hypotheses of Lemmas~\ref{pd-kappa-defects-digits-count-tail},
\ref{pd-kappa-defects-digits-count-kept} and~\ref{pd-kappa-defects-digits-count-match}. Let the data be as in the
statement; write $s=S_j$, $\hat a=\hat a_j$ and $b_+=b_T(s_{j+1})$, let $\hat b_i$ and $\hat\alpha_i$ be the numbers of
Definition~\ref{pd-kappa-defects-segment-lattice}(1) for the digits, with $\alpha_i=a_T(y_{j,i})$, and put
$\Pi=\prod_{i=1}^sP_m(\hat b_i)P_m(\hat\alpha_i)$, with $P_m$ as in
Lemma~\ref{pd-kappa-defects-windows-model-segment}. By Definition~\ref{pd-kappa-defects-segment-mean-conditions},
$d_j(\rho,\rho')=P_m(\hat a)\,\Pi\,P_m(b_+)$.

\emph{Facts.} The facts (a) to (g) in the proof of Lemma~\ref{pd-kappa-defects-margin} hold, and so do the fact (2) in
the proof of Lemma~\ref{pd-kappa-defects-digits-top-rows} and Step~1 in the proof of
Lemma~\ref{pd-kappa-defects-windows-model-step}. They give $k\geq5$, so $m\geq4$; $\Lambda\geq1$; $\hat a\geq a/(2Q)>0$, so
$\hat a\geq1$; $4\Lambda\leq2N\Lambda\leq k=m+1$, by (d), as $N=|T|\geq2$; and $b_+\geq b'(1-\beta_e)\geq b'/2\geq k^3Q/2\geq
m+1$, as $s_{j+1}\in T$, $\beta_e\leq1/12$ and $k^3Q\leq b'$. By Lemma~\ref{pd-kappa-defects-digits-count-kept}(1),
$0\leq n\leq\hat a$, and $A=\hat a-n$ satisfies $\hat a(1-N_1\Lambda/k)\leq A\leq\hat a$.

\emph{The bound.} Put $e_1=1-e^{-\Lambda/8}$ and $e_2=1-e^{-\Lambda/3}$, which are $\geq0$, $J=\lfloor\Lambda b_+/k\rfloor$
and $H=\sum_{j_+=0}^J\binom{b_+-j_++m-1}{m-1}$, and let $K=K_A$ be the kept sum for $p$, $\Lambda$, $m$, $s$ and the
$\hat b_i$ and $\hat\alpha_i$ (Definition~\ref{pd-kappa-defects-digits-count-pairs}). Then:
\begin{itemize}
\item $(1-mn/\hat a)P_m(\hat a)\leq P_m(A)$, by Lemma~\ref{pd-kappa-defects-digits-count-binomial}(1);
\item $e_1P_m(A)\,\Pi\leq K$, by Lemma~\ref{pd-kappa-defects-digits-count-tail}; in particular $K\geq0$;
\item $e_2P_m(b_+)\leq H$, by Lemma~\ref{pd-kappa-defects-digits-count-binomial}(2) with $B=b_+$;
\item $KH\leq c_j(\mu;\rho,\rho')$, by Lemma~\ref{pd-kappa-defects-digits-count-match}. Its map $\Phi$ is the map
  $z\mapsto\Phi(z)$ of the proof of Lemma~\ref{pd-kappa-defects-digits-count-upper}, whose constants $1$, $1$ and
  $1000(C_{E1}+C_{E2}+1)(C_c+1)$ are at most these: its targets are pairs for $\hat a$ with an integer $j_+\geq0$, it
  keeps $\kappa_{2s+1}$ and $j_+$, and it is injective (\emph{The map} and \emph{Injectivity} in that proof).
\end{itemize}
As $e_1\Pi e_2P_m(b_+)\geq0$,
\[
 e_1e_2\Big(1-\frac{mn}{\hat a}\Big)d_j(\rho,\rho')\leq e_1\Pi e_2P_m(b_+)P_m(A)\leq Ke_2P_m(b_+)\leq KH\leq
 c_j(\mu;\rho,\rho').
\]
\end{proof}
% Lean: RestCountLow = restCountLow_of'' (chainDim_of hY) hT hMU kept_blocks_holds = restCountLow_of hT
%   (restMatch_of hCD hMU (restMatchLow_of kept_blocks_holds)), as in stmtEExt_of, with
%   hT = restTail_of hCD hNB restPairCount (restNotKeptWt_of ...) (Lemma pd-kappa-defects-digits-count-tail). The
%   thresholds are those of threshold_standard_one (1, 1 and T2 CE1 CE2 Cc 1), RestTail and RestMatch, which
%   includes those of RestMatchUp = restMatchUp_of (the larger of those of RestFirstRows and of 1, 1 and
%   T2 CE1 CE2 Cc 1). The proof
%   is countLow_of: A0n and hA0pos (half_ratio_le_segment_base_of_std), heps (LaK, de6), hMle, hB (winB, beLe,
%   kQ), h4La (LaN, N2), t3 = cl_sd_margin, t1 = TailAt, t2 = cl_hockey_low, t0 = MatchAt, eB (mmB, dimMM,
%   symDim) and the final calc.

\begin{lemma}\label{pd-kappa-defects-digits-upper}
For every $C_c\geq0$ there is $T_0>0$ such that for every $C_0\geq T_0$ there is $T_1>0$ such that for
all $C_{E1},C_{E2},C_m\geq T_1$ there is $T_2>0$ with the following property for every $C_H\geq T_2$.
Let $p\geq5$ be a prime and $n'\geq2$ an even integer, and let $m$, $G'$, $N_1$, $c_{cl}$, $D_p$, $c_*$,
$T$, $Y$, $Y_b$, $b'$, $b_l$, $b_u$ and an integer $a$ satisfy hypotheses (1) to (3) of
Theorem~\ref{pd-gain-theorem-one} with $C_0$ and $C_H$. Let $c\geq0$ be an integer with
$c\leq C_c(c_*+|T|k\log_pa+|T|c_{cl}+n'c_{cl})$, and let $a_T,b_T\colon\mathbf Z/n'\to\mathbf Z_{\geq0}$
satisfy $a\leq a_T(s)\leq a+a/(kY)+2J_0$ and $b_l\leq b_T(s)\leq b_u+2J_0$ for $s\in T$, where
$J_0=\Lambda b'/k$. Form the numbers of Definition~\ref{pd-kappa-defects-numbers} with these
constants; the segment data do not involve $C_m$, $C_A$, $C_t$ or $C_\kappa$. Let an unramified
quadratic Hida setting of degree $n'$ at $p$ with ring $\mathcal R$ be given, an index $1\leq j\leq q$, a
real number $\mu\geq0$ and digits $\rho,\rho'$ for $j$, and form $r_j(\mu;\rho,\rho')$ and $d_j(\rho,\rho')$ with
$a_T$ and $b_T$ as in Definition~\ref{pd-kappa-defects-segment-mean-conditions}. Then
$r_j(\mu;\rho,\rho')\leq d_j(\rho,\rho')$.
\end{lemma}
\leanref{Formalism/SharpRestriction/GrowthBounds/BinomialRatio/RankDefs.lean}{Formalism.MainNotions.ColumnSystems.Pf.RestUp}
\leanref{Formalism/SharpRestriction/GrowthBounds/BinomialRatio/RankDefs.lean}{Formalism.MainNotions.ColumnSystems.Pf.UpAt}
\leanref{Formalism/SharpRestriction/GrowthBounds/BinomialRatio/RankDefs.lean}{Formalism.MainNotions.ColumnSystems.Pf.mmB}
\leanref{Formalism/SharpRestriction/GrowthBounds/BinomialRatio/RankOnDigitsCast.lean}{Formalism.MainNotions.ColumnSystems.Pf.restUp_of}

\begin{proof}
Let $T_0$, $T_1$ and $T_2$ be the numbers of Lemma~\ref{pd-kappa-defects-digits-count-upper}. By
Lemma~\ref{pd-kappa-defects-digits-count}, $r_j(\mu;\rho,\rho')=c_j(\mu;\rho,\rho')$, and
Lemma~\ref{pd-kappa-defects-digits-count-upper} gives $c_j(\mu;\rho,\rho')\leq d_j(\rho,\rho')$.
\end{proof}
% Lean: RestUp = restUp_of blockRankAt hY charZeroAt hCU (restUp_of_adm), with hY = dimYAt_of_adm hA and
%   hCU = upper_count_of_top_rows (chainDim_of hY) (numBranch_of ...) top_rows_bound_holds; restUp_of is
%   TransportedAlt.mono of hCU (the same thresholds) with rankOnDigits_cast (Lemma
%   pd-kappa-defects-digits-count), the ring being a UniformizedDVR by dvrOfStd.

\begin{lemma}\label{pd-kappa-defects-digits-lower}
For every $C_c\geq0$ there is $T_0>0$ such that for every $C_0\geq T_0$ there is $T_1>0$ such that for
all $C_{E1},C_{E2},C_m\geq T_1$ there is $T_2>0$ with the following property for every $C_H\geq T_2$.
Let $p$, $n'$, $m$, $G'$, $N_1$, $c_{cl}$, $D_p$, $c_*$, $T$, $Y$, $Y_b$, $b'$, $b_l$, $b_u$, $a$, $c$, $a_T$ and $b_T$
satisfy the hypotheses of Lemma~\ref{pd-kappa-defects-digits-upper} with these constants, and form the
numbers of Definition~\ref{pd-kappa-defects-numbers}. Let an unramified quadratic Hida setting of degree
$n'$ at $p$ with ring $\mathcal R$ be given, an index $1\leq j\leq q$, a real number $\mu\geq0$ and digits
$\rho,\rho'$ for $j$, and form $r_j(\mu;\rho,\rho')$ and $d_j(\rho,\rho')$ with $a_T$ and $b_T$ as in
Definition~\ref{pd-kappa-defects-segment-mean-conditions}. Put $\hat a_j=\lfloor(a-E_0)/p^{g_j}\rfloor$ and
$n=\lceil\mu\rceil$. If $n\leq\hat a_jN_1\Lambda/k$, then
\[
 \big(1-e^{-\Lambda/8}\big)\big(1-e^{-\Lambda/3}\big)\Big(1-\frac{mn}{\hat a_j}\Big)d_j(\rho,\rho')
 \leq r_j(\mu;\rho,\rho').
\]
\end{lemma}
\leanref{Formalism/SharpRestriction/GrowthBounds/BinomialRatio/RankDefs.lean}{Formalism.MainNotions.ColumnSystems.Pf.RestLow}
\leanref{Formalism/SharpRestriction/GrowthBounds/BinomialRatio/RankDefs.lean}{Formalism.MainNotions.ColumnSystems.Pf.LowAt}
\leanref{Formalism/SharpRestriction/GrowthBounds/BinomialRatio/RankOnDigitsCast.lean}{Formalism.MainNotions.ColumnSystems.Pf.restLow_of}

\begin{proof}
Let $T_0$, $T_1$ and $T_2$ be the numbers of Lemma~\ref{pd-kappa-defects-digits-count-lower}. By
Lemma~\ref{pd-kappa-defects-digits-count}, $r_j(\mu;\rho,\rho')=c_j(\mu;\rho,\rho')$, and
Lemma~\ref{pd-kappa-defects-digits-count-lower} gives the bound for $c_j(\mu;\rho,\rho')$.
\end{proof}
% Lean: RestLow = restLow_of blockRankAt hY charZeroAt hCL, with hY as for the upper bound and
%   hCL = restCountLow_of'' (chainDim_of hY) hT hMU kept_blocks_holds, as in stmtEExt_of; restLow_of is
%   TransportedAlt.mono of hCL with rankOnDigits_cast (Lemma pd-kappa-defects-digits-count). The same
%   bound, summed over the digits by segLow_of, gives Lemma pd-kappa-defects-windows-model-segment.

\begin{lemma}\label{pd-kappa-defects-segment-mean-digits}
For every $C_c\geq0$ there is $T_0>0$ such that for every $C_0\geq T_0$ there is $T_1>0$ such that for
all $C_{E1},C_{E2},C_m\geq T_1$ there is $T_2>0$ such that for every $C_H\geq T_2$ and all reals
$C_\kappa,C_s$, every instance satisfies conditions (10) and (11) of
Definition~\ref{pd-kappa-defects-segment-mean-conditions} with these constants.
\end{lemma}
\leanref{Formalism/SharpRestriction/GrowthBounds/ColumnSystems.lean}{Formalism.MainNotions.ColumnSystems.ExtendedDigitRankBound}
\leanref{Formalism/SharpRestriction/GrowthBounds/ColumnSystems.lean}{Formalism.MainNotions.ColumnSystems.RootMultiplicityTwo}
\leanref{Formalism/SharpRestriction/GrowthBounds/ColumnSystems.lean}{Formalism.MainNotions.ColumnSystems.SegmentMarginBound}
\leanref{Formalism/SharpRestriction/GrowthBounds/Joints.lean}{Formalism.MainNotions.G.Joints.fourth_extension_gives_ratio_bound}
\leanref{Formalism/SharpRestriction/GrowthBounds/Joints.lean}{Formalism.MainNotions.G.Joints.residue_bound_of_extension}
\leanref{Formalism/SharpRestriction/GrowthBounds/BinomialRatio/LeAndHalfLeProdOneSubExpNegOfBounds.lean}{Formalism.MainNotions.ColumnSystems.Pf.one_segment_of_upper_lower}

\begin{proof}
\emph{The constants.} Let $C_c\geq0$. Let $T_0$ be the largest of $24$ and the numbers $T_0$ of
Lemmas~\ref{pd-kappa-defects-digits-upper} and~\ref{pd-kappa-defects-digits-lower} for $C_c$; for
$C_0\geq T_0$ let $T_1$ be the largest of $1$ and their numbers $T_1$; for $C_{E1},C_{E2},C_m\geq T_1$ let
$T_2$ be the largest of $1000(C_{E1}+C_{E2}+1)(C_c+1)$ and their numbers $T_2$. Let $C_H\geq T_2$ and let
$C_\kappa$, $C_s$ be reals.

\emph{The windows.} Let an instance satisfy its hypotheses, and let $a_T$ be a window function. As
$x_a,x_b\leq2J_0$, $a\leq a_T(s')\leq a+a/(kY)+2J_0$ and $b_l\leq b_T(s')\leq b_u+2J_0$ for $s'\in T$. So
Lemmas~\ref{pd-kappa-defects-digits-upper} and~\ref{pd-kappa-defects-digits-lower} apply to the data of
the instance with $a_T$ and $b_T$.

\emph{Facts.} As $C_0,C_{E1},C_{E2},C_m\geq1$ and $C_H\geq1000(C_{E1}+C_{E2}+1)(C_c+1)$, the facts (a)
to (g) in the proof of Lemma~\ref{pd-kappa-defects-margin} hold, and so do the facts (h) and (i) in the
proof of Lemma~\ref{pd-kappa-defects-windows-model-step} and its Step~1 for every index $1\leq j'\leq q$,
which use only these hypotheses. They give $k\geq5$, $Q\geq1$, $Y\geq1$, $N_1\geq1$, $b'>0$,
$kY\Lambda b'\leq a$, $k^3Q\leq b'$ and $\hat a_{j'}\geq a/(2Q)$, where
$\hat a_{j'}=\lfloor(a-E_0)/p^{g_{j'}}\rfloor$; and $\Lambda=C_0N_1^2\log a\geq C_0\geq24$, as $\log a\geq1$
by (c).

\emph{The bound for each segment.} Let $1\leq j'\leq q$, let $0\leq\mu'\leq b'/(kQY)$ and let $\rho,\rho'$ be
digits for $j'$. We show that
\[
 d_{j'}(\rho,\rho')/2\leq r_{j'}(\mu';\rho,\rho')\leq d_{j'}(\rho,\rho').
\]
The upper bound is Lemma~\ref{pd-kappa-defects-digits-upper}. For the lower bound,
Lemma~\ref{pd-kappa-defects-digits-half}, with $M=\lceil\mu'\rceil$, $\hat a=\hat a_{j'}$ and $k-1=m$, gives
$\lceil\mu'\rceil\leq\hat a_{j'}N_1\Lambda/k$ and shows that the factor in front of $d_{j'}(\rho,\rho')$ in
Lemma~\ref{pd-kappa-defects-digits-lower} is at least $1/2$. As $d_{j'}(\rho,\rho')\geq0$, that lemma gives
$d_{j'}(\rho,\rho')/2\leq r_{j'}(\mu';\rho,\rho')$.

\emph{Condition (11).} This is the bound for $j'=j$.

\emph{Condition (10).} As $b'>0$ and $k,Q,Y>0$, the margin $\mu'=0$ satisfies $0\leq\mu'\leq b'/(kQY)$. So
the upper half of the bound with $\mu'=0$ gives $r_{j'}(0;\rho,\rho')\leq d_{j'}(\rho,\rho')$ for every
$1\leq j'\leq q$ and all digits $\rho,\rho'$ for $j'$.
\end{proof}
% Lean: ExtendedDigitRankBound and RootMultiplicityTwo are r2.toExtendedDigitRankBound and r2 in
%   closed_of_rests, with r2 from RootMultiplicityTwoBeyond = one_segment_of_upper_lower hU hL
%   (restUp_of blockRankAt hY charZeroAt (upper_count_of_top_rows ...) and restLow_of ..., as in
%   stmtEExt_of); Joints.fourth_extension_gives_ratio_bound and Joints.residue_bound_of_extension pass
%   to conditions (10) and (11) through extended_hypothesis_of_window. The thresholds are those of
%   one_segment_of_upper_lower: TransportedAlt.of_explicit with 24, 1 and T2 CE1 CE2 Cc 1 for Std (std_of),
%   and those of hU and hL. The lower bound uses Std.La_ge, Std.half_ratio_le_segment_base_of_std and
%   le_and_half_le_prod_one_sub_exp_neg_of_bounds (Lemma pd-kappa-defects-digits-half). The bound for
%   each segment is RootMultiplicityTwo = SegmentMarginBound with M X = b'/(kQY), for every segment j.
%   Condition (10) is RootMultiplicityTwo.toExtendedDigitRankBound: its second part at m = 0, with
%   0 <= b'/(kQY) from g10 (b' > 0) and the positivity of k, Q and Y.

\begin{lemma}\label{pd-kappa-defects-segment-mean}
For every $C_c\geq0$ there is $T_0>0$ such that for every $C_0\geq T_0$ there is $T_1>0$ such that for
all $C_{E1},C_{E2},C_m\geq T_1$ there is $T_2>0$ such that for every $C_H\geq T_2$ there is $C_A>0$
with the following property. Let $p\geq5$ be a prime and $n'\geq2$ an even integer, and let $m$, $G'$,
$N_1$, $c_{cl}$, $D_p$, $c_*$, $T$, $Y$, $Y_b$, $b'$, $b_l$, $b_u$ and an integer $a$ satisfy hypotheses (1) to
(3) of Theorem~\ref{pd-gain-theorem-one} with $C_0$ and $C_H$, and $kY^2\log a\leq a/b'$. Let $c\geq0$
be an integer with $c\leq C_c(c_*+|T|k\log_pa+|T|c_{cl}+n'c_{cl})$, let $x_a,x_b$ be reals with
$0\leq x_a,x_b\leq2J_0$, where $J_0=\Lambda b'/k$, let $b_T\colon\mathbf Z/n'\to\mathbf Z_{\geq0}$
satisfy $b_l\leq b_T(s)\leq b_u+x_b$ for $s\in T$, and let $\mu$ be a real with
$0\leq\mu\leq b'/(kQY)$, where $Q=p^{G'}$. Form the numbers of
Definition~\ref{pd-kappa-defects-numbers} with these constants; the segment data and $D_A$ do not
involve $C_t$, $C_m$, $C_\kappa$ or $a_T$. Put $W=[a,a+a/(kY)+x_a]\cap\mathbf Z$, and fix
an unramified quadratic Hida setting of degree $n'$ at $p$ with ring $\mathcal R$ and an index
$1\leq j\leq q$. Then the mean of $\kappa_{S_j(b_T)}(\mu,D_A;\alpha,\alpha_+)$, formed over $\mathcal R$
with $r=m$, over $(\alpha,\alpha_+)\in W^{S_j}\times W$ is at most $C_A(S_j-\eta_j+1)N_1/Y$.
\end{lemma}
\leanref{Formalism/SharpRestriction/GrowthBounds/ProjectionLemmaBaseLevel.lean}{Formalism.MainNotions.flipOneThree.PathPairCorrespondence}
\leanref{Formalism/SharpRestriction/GrowthBounds/FlipOneThree.lean}{Formalism.MainNotions.flipOneThree.Conclusion1}
\leanref{Formalism/SharpRestriction/GrowthBounds/FlipOneThree.lean}{Formalism.MainNotions.flipOneThree.HI}
\leanref{Formalism/SharpRestriction/GrowthBounds/FlipOneThree.lean}{Formalism.MainNotions.flipOneThree.meanKapT}
\leanref{Formalism/SharpRestriction/GrowthBounds/FlipOneThree.lean}{Formalism.MainNotions.flipOneThree.Inst}
\leanref{Formalism/SharpRestriction/GrowthBounds/Joints.lean}{Formalism.MainNotions.G.Joints.bad_share_bound_of_extension}

\begin{proof}
\emph{The constants.} Let $C_c\geq0$. Let $T_0$ be the largest of the numbers $T_0$ of
Lemmas~\ref{pd-kappa-defects-segment-mean-theorem-a}, \ref{pd-kappa-defects-segment-mean-letters},
\ref{pd-kappa-defects-segment-mean-ratio} and~\ref{pd-kappa-defects-segment-mean-digits} for $C_c$; for
$C_0\geq T_0$ let $T_1$ be the largest of their numbers $T_1$; for $C_{E1},C_{E2},C_m\geq T_1$ let $T_2$
be the largest of their numbers $T_2$. For $C_H\geq T_2$ let $C_\kappa$ be the number $C^0_\kappa$ of
Lemma~\ref{pd-kappa-defects-segment-mean-ratio}, put $C_s=C_\kappa(C_m+4C_{E1}+C_{E2}+16)$, and let $C_A$
be the constant of Lemma~\ref{pd-kappa-defects-segment-mean-theorem-a} for $C_\kappa$ and $C_s$.

\emph{The bound.} Let the data be as in the statement. With $\mu$, $x_a$, $x_b$, $b_T$, the setting and
$j$ they form an instance (Definition~\ref{pd-kappa-defects-segment-mean-conditions}), since
$m\geq2N_1+2\geq1$ by (1), and the instance satisfies its hypotheses. By
Lemmas~\ref{pd-kappa-defects-segment-mean-letters}, \ref{pd-kappa-defects-segment-mean-weyl},
\ref{pd-kappa-defects-segment-mean-restriction}, \ref{pd-kappa-defects-segment-mean-paths},
\ref{pd-kappa-defects-segment-mean-ratio} and~\ref{pd-kappa-defects-segment-mean-digits} it satisfies
conditions (1) to (11) with $C_0$, $C_H$, $C_{E1}$, $C_{E2}$, $C_c$, $C_\kappa$ and $C_s$. So
Lemma~\ref{pd-kappa-defects-segment-mean-theorem-a} bounds the mean by $C_A(S_j-\eta_j+1)N_1/Y$.
\end{proof}
% Lean: in first_theorem_admitted_closed, PathPairCorrespondence is bad_share_bound_of_extension hconds1:
%   thresholds_and (h.p13A D) (thresholds_and (h.thresholdClauseOne D) (h.thresholdClauseThree D)) gives
%   hS1, and CA comes from hS1 Ckap (Cs + Ckap) with Ckap, Cs from the extended windows theorem; then
%   hSt X applies to the eleven conditions ha1 X (thresholdClauseOne =
%   CountingStatementReduction.positivity_condition_eventually), Joints.weight_condition_of_integral_form
%   (hT12 of tool_statements_hold), ha3 X (exponent_bound_eventually), h.tKP1, h.tKP2, h.tKP3, h.tRES,
%   h.instanceClauseEight, Joints.extension_gives_ratio_bound (ExtendedModelRank and ExtendedKappaBound),
%   Joints.fourth_extension_gives_ratio_bound (ExtendedDigitRankBound) and
%   Joints.residue_bound_of_extension (RootMultiplicityTwo). The thresholds are the maxima of those of
%   the extended windows theorem and of theorem A. Conclusion1 is stated for an Inst X: Loc data with
%   real ccl, DLP, cs, a setting, a segment X.j, the margin X.m and the window extensions X.xa, X.xb;
%   HI is WindowHypotheses, CBound Cc, the bounds on xa, xb and bT, logl and the bounds on X.m; the
%   segment datum is X.S C0 CE1 CE2 (with c = X.c) and the depth is depthD CA, with the same CA.

\subsection{The projection property between $p_0$ and $X'$}\label{subsection-pd-projection-p0-xprime}

\begin{lemma}\label{pd-kappa-defects-numerical}
Let $C_0,C_A,C_m,C^*_H>0$ and $T_M=\max(108,8C_A+93,1/C_0^2,4C^*_H+1,2C_m+24)$. Let $C_H\geq T_M$, and
let $p$, $n'$, $m$, $G'$, $N_1$, $c_{cl}$, $D_p$, $c_*$, $T$, $Y$, $Y_b$, $a$, $b'$, $b_l$, $b_u$
satisfy hypotheses (1) to (3) of Theorem~\ref{pd-gain-theorem-one} with $C_0$ and $C_H$. Form the
numbers of Definition~\ref{pd-kappa-function} with $C_0$ and $C_A$, and put
$m_0=C_mkN_1Qc_\flat+12c_\flat$ with $Q=p^{G'}$. Then:
\begin{enumerate}
\item hypotheses (1) to (3) of Theorem~\ref{pd-gain-theorem-one} hold with $\bar a$, $C^*_H$ and
  $Y_b/2$ in place of $a$, $C_H$ and $Y_b$;
\item $c_\flat\leq4(c_*+|T|k\log_p\bar a+|T|c_{cl}+n'c_{cl})$;
\item $0\leq D_0$ and $4D_0<c_\flat$;
\item $J_0\geq0$, $\bar a+\bar a/(kY)\leq u_a+J+c_{cl}\leq\bar a+\bar a/(kY)+2J_0$ and
  $u_b+J+c_{cl}\leq b_u+2J_0$;
\item $kY^2\log\bar a\leq\bar a/b'$;
\item $0\leq m_0\leq b'/(kQY)$.
\end{enumerate}
\end{lemma}
\leanref{Formalism/SharpRestriction/GrowthBounds/ColumnSpanDensity/FactorizationBoundMonoLattices.lean}{Formalism.MainNotions.flipOneThree.Pf.FixedWeightNumericalFacts}
\leanref{Formalism/SharpRestriction/GrowthBounds/ColumnSpanDensity/LocalHypothesisOfGlobal.lean}{Formalism.MainNotions.flipOneThree.Pf.MeanKappaNumericalFacts}
\leanref{Formalism/SharpRestriction/GrowthBounds/ColumnSpanDensity/ThresholdFixedWeightFacts.lean}{Formalism.MainNotions.flipOneThree.Pf.threshold_fixed_weight_facts}
\leanref{Formalism/SharpRestriction/GrowthBounds/ColumnSpanDensity/ThresholdFixedWeightFacts.lean}{Formalism.MainNotions.flipOneThree.Pf.threshold_mean_kappa_facts}
\leanref{Formalism/SharpRestriction/GrowthBounds/ColumnSpanDensity/TwoHundredLeRadius.lean}{Formalism.MainNotions.flipOneThree.Pf.Cst.TM}
\leanref{Formalism/SharpRestriction/GrowthBounds/ColumnSpanDensity/NumBase.lean}{Formalism.MainNotions.flipOneThree.Pf.Cst.TLP}

\begin{lemma}\label{pd-kappa-defects-margin-share}
Let $C_c\geq0$, $C_0\geq1$, $C_{E1},C_{E2},C_m\geq1$, $C_\kappa>0$ and
$C_H\geq1000(C_{E1}+C_{E2}+1)(C_c+1)$, and put $C_s=C_\kappa(C_m+4C_{E1}+C_{E2}+16)$. Let $p\geq5$ be a
prime and $n'\geq2$ an even integer, and let $m$, $G'$, $N_1$, $c_{cl}$, $D_p$, $c_*$, $T$, $Y$, $Y_b$,
$b'$, $b_l$, $b_u$ and an integer $a$ satisfy hypotheses (1) to (3) of
Theorem~\ref{pd-gain-theorem-one} with $C_0$ and $C_H$. Let $c\geq0$ be an integer with
$c\leq C_c(c_*+|T|k\log_pa+|T|c_{cl}+n'c_{cl})$, and let $a_T,b_T\colon\mathbf Z/n'\to\mathbf Z_{\geq0}$
satisfy $a\leq a_T(s)\leq a+a/(kY)+2J_0$ and $b_l\leq b_T(s)\leq b_u+2J_0$ for $s\in T$, where
$J_0=\Lambda b'/k$. Form the numbers of Definition~\ref{pd-kappa-defects-numbers} with these
constants. Then $\kappa_P(\mu)\leq C_s|T|N_1/Y$ for every real $\mu$ with $0\leq\mu\leq b'/(kQY)$.
\end{lemma}
\leanref{Formalism/SharpRestriction/GrowthBounds/ProjectionLemmaBaseLevel.lean}{Formalism.MainNotions.flipOneThree.admB3}
\leanref{Formalism/SharpRestriction/GrowthBounds/Joints.lean}{Formalism.MainNotions.G.Joints.kappa_share_bound_of_extension}
\leanref{Formalism/SharpRestriction/GrowthBounds/Joints.lean}{Formalism.MainNotions.G.Joints.kappa_share_of_extension_bounds}
\leanref{Formalism/SharpRestriction/GrowthBounds/Joints.lean}{Formalism.MainNotions.G.Joints.kap_margin}

\begin{proof}
Put $C'_s=C_\kappa(C_m+4C_{E1}+C_{E2}+15)$, so that $C_s=C'_s+C_\kappa$. By
Lemma~\ref{pd-kappa-defects-margin}, $\kappa_P=\kappa_P(m_0)\leq C'_s|T|N_1/Y$. By hypotheses (1) and
(2), $Y\geq1$, $N_1\geq1$ and $b'>0$; also $k\geq1$ and $Q\geq1$, and $m_0=C_mkN_1Qc+12c\geq0$. Let
$0\leq\mu\leq b'/(kQY)$. Then $kQ\mu/b'\leq1/Y$, so
\[
 \frac{kQ(\mu+N_1k)}{b'}\leq\frac{kQ\mu}{b'}+\frac{kQ(m_0+N_1k)}{b'}\leq\frac{kQ(m_0+N_1k)}{b'}+\frac1Y .
\]
The other three terms in the bracket of $\kappa_P(\mu)$ do not involve $\mu$, so
$\kappa_P(\mu)\leq\kappa_P(m_0)+C_\kappa|T|/Y\leq C'_s|T|N_1/Y+C_\kappa|T|N_1/Y=C_s|T|N_1/Y$.
\end{proof}
% Lean: the second clause of admB3 with Cs + Ckap, from kappa_share_bound_of_extension:
% kappa_share_of_extension_bounds with h2 = ExtendedKappaBound (rest2, Lemma pd-kappa-defects-margin),
% loc_kap_facts for the signs and kap_margin for the step from m0 to mu.

\begin{lemma}\label{pd-kappa-product-defects}
There is $T_0>0$ such that for every $C_0\geq T_0$ there are constants $C_A,C_t,C_s>0$ and $T_3\geq0$
with the following property. Let $C_H\geq T_3$, and let $p$, $n'$, $m$, $G'$, $N_1$, $c_{cl}$,
$D_p$, $c_*$, $T$, $Y$, $Y_b$, $a$, $b'$, $b_l$, $b_u$ satisfy hypotheses (1) to (3) of
Theorem~\ref{pd-gain-theorem-one}. Let $s_1<\cdots<s_q$ be the even elements of $T$, and for each $j$
(indices modulo $q$) let $\Sigma_j\subset T$ consist of the $S_j$ odd elements of $T$ strictly between
$s_j$ and $s_{j+1}$, as in Theorem~\ref{pd-gain-theorem-one}, together with $s_{j+1}$. Form the
numbers of Definition~\ref{pd-kappa-function} with $C_0$, $C_A$ and $C_t$, and put
$I'_a=[\bar a,u_a+J+c_{cl}]\cap\mathbf Z$ and $I'_b=[l_b,u_b+J+c_{cl}]\cap\mathbf Z$. Then for every
unramified quadratic Hida setting of degree $n'$ at $p$ there are a real number $\kappa_P$ and
functions $\kappa_1,\dots,\kappa_q$ on the pairs $(A,B)$ of functions $T\to\mathbf Z_{\geq0}$ such
that:
\begin{enumerate}
\item for $A\in I_a'^T$, $B\in I_b'^T$ and every edge $e$, there is a split family of depth $D_0$
  and size at least $\big(1-\kappa_P-\sum_j\kappa_j(A,B)\big)\dim V^\times(A,B)/[k]$ on
  $(\mathcal C^\times_{t,c_\flat}(e),\mathcal C'^\times_{t,c_\flat}(e))$, where
  $t=\lfloor t_0(A,B)\rfloor$ (Definition~\ref{pd-kappa-product-space});
\item if $I'_a$ and $I'_b$ are not empty, then $\kappa_P\leq C_s|T|N_1/Y$;
\item $\kappa_j(A,B)$ depends on $A$ only through the restriction of $A$ to $\Sigma_j$;
\item for every $j$ and every $B\in I_b'^T$, the sum of $\kappa_j(A,B)$ over the restrictions
  $A\in I_a'^{\Sigma_j}$ is at most $|I'_a|^{S_j+1}C_A(S_j+1)N_1/Y$.
\end{enumerate}
\end{lemma}
\leanref{Formalism/SharpRestriction/GrowthBounds/ProjectionLemmaBaseLevel.lean}{Formalism.MainNotions.flipOneThree.admB2}
\leanref{Formalism/SharpRestriction/GrowthBounds/ProjectionLemmaBaseLevel.lean}{Formalism.MainNotions.flipOneThree.admB3}
\leanref{Formalism/SharpRestriction/GrowthBounds/ProjectionLemmaBaseLevel.lean}{Formalism.MainNotions.flipOneThree.PathPairCorrespondence}
\leanref{Formalism/SharpRestriction/GrowthBounds/ProjectionLemmaBaseLevel.lean}{Formalism.MainNotions.flipOneThree.kapPI}
\leanref{Formalism/SharpRestriction/GrowthBounds/ProjectionLemmaBaseLevel.lean}{Formalism.MainNotions.flipOneThree.kapTI}
\leanref{Formalism/SharpRestriction/GrowthBounds/ColumnSpanDensity/LocalHypothesisOfGlobal.lean}{Formalism.MainNotions.flipOneThree.Pf.segment_mean_kappa_bound}
\leanref{Formalism/SharpRestriction/GrowthBounds/ObjSeg.lean}{Formalism.MainNotions.G.segATop}
\leanref{Formalism/SharpRestriction/GrowthBounds/ObjSeg.lean}{Formalism.MainNotions.G.segAPlus}

\begin{proof}
\emph{The constants.} Let $T_0^W$ and $T_0^S$ be the numbers $T_0$ of
Lemmas~\ref{pd-kappa-defects-windows} and~\ref{pd-kappa-defects-segment-mean} for $C_c=4$, and put
$T_0=\max(T_0^W,T_0^S,1)$. Let $C_0\geq T_0$, let $T_1^W$ and $T_1^S$ be the numbers $T_1$ of the two
lemmas for $C_0$, and put $C_E=\max(T_1^W,T_1^S,1)$. Let $T_2^W$ and $T_2^S$ be their numbers $T_2$ for
$C_{E1}=C_{E2}=C_m=C_E$, and put $C^*_H=\max(T_2^W,T_2^S,5000(2C_E+1))$. For $C_H=C^*_H$,
Lemma~\ref{pd-kappa-defects-windows} gives constants $C_\kappa,C_t>0$ and
Lemma~\ref{pd-kappa-defects-segment-mean} a constant $C_A>0$. Put $C'_s=C_\kappa(6C_E+15)$ and
$C_s=C'_s+C_\kappa$, and let $T_3=\max(108,8C_A+93,1/C_0^2,4C^*_H+1,2C_E+24)$ be the number $T_M$ of
Lemma~\ref{pd-kappa-defects-numerical} for $C_0$, $C_A$, $C_m=C_E$ and $C^*_H$.

\emph{The data.} Let $C_H\geq T_3$, let the data satisfy hypotheses (1) to (3) of
Theorem~\ref{pd-gain-theorem-one}, and form the numbers of Definition~\ref{pd-kappa-function}. By
Lemma~\ref{pd-kappa-defects-numerical}, its conclusions (1) to (6) hold. By (1), the data with $\bar a$
and $Y_b/2$ in place of $a$ and $Y_b$ satisfy hypotheses (1) to (3) of
Theorem~\ref{pd-gain-theorem-one} with $C_0$ and $C^*_H$; by (2), $c=c_\flat$ satisfies the bound on
$c$ in Lemmas~\ref{pd-kappa-defects-margin}, \ref{pd-kappa-defects-windows}
and~\ref{pd-kappa-defects-segment-mean} with $C_c=4$ and $\bar a$ in place of $a$. Form the numbers of
Definition~\ref{pd-kappa-defects-numbers} with $\bar a$, $c=c_\flat$ and the constants $C_0$, $C_A$,
$C_t$, $C_{E1}=C_{E2}=C_m=C_E$ and $C_\kappa$. Then $\Lambda=\bar\Lambda$, $J_0$ and $t_0$ are those of
Definition~\ref{pd-kappa-function}, $m_0$ is that of Lemma~\ref{pd-kappa-defects-numerical}, and
$D_0=qD_A+4\kappa_0$.

Fix an unramified quadratic Hida setting of degree $n'$ at $p$, with ring $\mathcal R$; so $p\geq5$
and $n'\geq2$ is even (Definition~\ref{pd-local-model}). For $A,B\colon T\to\mathbf Z_{\geq0}$ let
$A^0$ and $B^0$ be their extensions by $0$ to $\mathbf Z/n'$, and put $\kappa_j(A,B)=\kappa_j(A^0,B^0)$,
the segment defects of Definition~\ref{pd-kappa-defects-numbers} over $\mathcal R$; and let $\kappa_P$
be the number of that definition.

(3) The segment defect $\kappa_j(A,B)$ is formed from the segment datum $S_j(B^0)$, which does not
involve $A$, and from the values of $A$ at $y_{j,1},\dots,y_{j,S_j}$ and $s_{j+1}$, which are the
elements of $\Sigma_j$.

(2) Suppose that $I'_a$ and $I'_b$ are not empty, and choose $a_0\in I'_a$ and $b_0\in I'_b$. Let
$a_T,b_T\colon\mathbf Z/n'\to\mathbf Z_{\geq0}$ be the functions with the values $a_0$ and $b_0$ on $T$
and $0$ elsewhere. By (4) of Lemma~\ref{pd-kappa-defects-numerical}, for $s\in T$,
$\bar a\leq a_T(s)\leq u_a+J+c_{cl}\leq\bar a+\bar a/(kY)+2J_0$ and
$b_l\leq l_b\leq b_T(s)\leq u_b+J+c_{cl}\leq b_u+2J_0$; and by (6), $0\leq m_0\leq b'/(kQY)$. So
Lemma~\ref{pd-kappa-defects-margin-share} applies with $C_c=4$, $C_0\geq1$,
$C_{E1}=C_{E2}=C_m=C_E\geq1$, $C_\kappa$ and $C^*_H\geq5000(2C_E+1)=1000(2C_E+1)(4+1)$, to the data with
$\bar a$ and $Y_b/2$, to $c=c_\flat$, to these $a_T$ and $b_T$, and to $\mu=m_0$. Its constant is
$C_\kappa(6C_E+16)=C_s$, so $\kappa_P=\kappa_P(m_0)\leq C_s|T|N_1/Y$.

(1) Let $A\in I_a'^T$, $B\in I_b'^T$ and $e$ an edge. By (4) of
Lemma~\ref{pd-kappa-defects-numerical}, for $s\in T$,
\[
 \bar a\leq A(s)\leq u_a+J+c_{cl}\leq\bar a+\frac{\bar a}{kY}+2J_0,\qquad
 b_l\leq l_b\leq B(s)\leq u_b+J+c_{cl}\leq b_u+2J_0 .
\]
Since $p>1$, $m\geq1$, $C_A>0$ and $\bar a\geq1$ (as $a>pb'>0$ by hypothesis (2)), $D_A\geq0$. For
each $j$,
Lemma~\ref{pd-kappa-defects-bad-share} shows that $\kappa_j(A,B)\geq0$ and that the segment lattice of
the $j$-th segment has a good share of depth $D_A$ and defect $\kappa_j(A,B)$. By (3) of
Lemma~\ref{pd-kappa-defects-numerical}, $4(qD_A+4\kappa_0)=4D_0<c_\flat$. So
Lemma~\ref{pd-kappa-defects-windows}, with $C^*_H$, $a_T=A^0$, $b_T=B^0$, $D_j=D_A$ and
$\kappa_j=\kappa_j(A,B)$, gives at $e$ a split family of depth $D_0$ and size at least
$(1-\kappa_P-\sum_j\kappa_j(A,B))\dim V^\times(A,B)/[k]$ on the cores
$(\mathcal C^\times_{t,c_\flat}(e),\mathcal C'^\times_{t,c_\flat}(e))$ with $t=\lfloor t_0(A,B)\rfloor$.

(4) Fix $j$ and $B\in I_b'^T$. Put $x_a=u_a+J+c_{cl}-\bar a-\bar a/(kY)$ and $x_b=2J_0$. By (4) of
Lemma~\ref{pd-kappa-defects-numerical}, $0\leq x_a\leq2J_0$ and $0\leq x_b$, and
$b_l\leq B(s)\leq b_u+x_b$ for $s\in T$, as in (1). By (5) and (6) of that lemma,
$kY^2\log\bar a\leq\bar a/b'$ and $0\leq m_0\leq b'/(kQY)$. So
Lemma~\ref{pd-kappa-defects-segment-mean} applies with $C^*_H$, $b_T=B^0$ and $\mu=m_0$. Its window is
$W=[\bar a,u_a+J+c_{cl}]\cap\mathbf Z=I'_a$, since $u_a+J+c_{cl}$ is an integer, and its segment defect
at $(\alpha,\alpha_+)$ is $\kappa_j(A,B)$ for every $A$ with $A(y_{j,i})=\alpha_i$ and
$A(s_{j+1})=\alpha_+$. As $(\alpha,\alpha_+)$ runs over $W^{S_j}\times W$, these restrictions of $A$
run once over $I_a'^{\Sigma_j}$. Hence the sum of $\kappa_j(A,B)$ over the restrictions
$A\in I_a'^{\Sigma_j}$ is at most $|I'_a|^{S_j+1}C_A(S_j-\eta_j+1)N_1/Y\leq|I'_a|^{S_j+1}C_A(S_j+1)N_1/Y$;
if $I'_a$ is empty, the sum is $0$.
\end{proof}
% Lean: the constants follow theorem_I_of_intermediate_bounds (Cc = 4, CE1 = CE2 = Cm = T1, CH = T2,
% the record K) and first_theorem_admitted_closed (maxima of the thresholds of extended_windows_theorem
% and of the threshold clauses; CA from hS1 Ckap (Cs + Ckap); the AdmB constant Cs + Ckap is C_s here).
% T3 = Cst.TM is the threshold of threshold_fixed_weight_facts and threshold_mean_kappa_facts.
% (1) is admB2 as applied in keyIneq_of (hF.hyp', hWin from winA, winB and Nat.le_ceil, hF.cb, hF.d0):
% bad_share_of_extension_proposition = depth_bound_of_extension with HerrExtensionClass and
% badShare_spec (depthD_nonneg). (2) is the second clause of admB3 (Lemma pd-kappa-defects-margin-share)
% as applied in meanX_of: after the cases hca and hcb of an empty window, at the constant functions
% fun _ => a0 and fun _ => b0 with a0 in WaP and b0 in WbP (hWin from winA, winB and Nat.le_ceil), with
% mm = m0I (hF.m0nn, hF.m0le). (3) is the definition of kapTI through
% segATop and segAPlus. (4) is segment_mean_kappa_bound: Conclusion1 at instOf K X R jj B (aT = 0,
% m = m0I, xa, xb = 2 J0), with local_hypothesis_of_global, instance_window_eq and (S).t <= segS.

\begin{lemma}\label{pd-kappa-product-segments}
There is $T_0>0$ such that for every $C_0\geq T_0$ there are constants $C_A,C_t,C_s>0$ and $T_3\geq0$
with the following property. Let $C_H\geq T_3$, and let $p$, $n'$, $m$, $G'$, $N_1$, $c_{cl}$,
$D_p$, $c_*$, $T$, $Y$, $Y_b$, $a$, $b'$, $b_l$, $b_u$ satisfy hypotheses (1) to (3) of
Theorem~\ref{pd-gain-theorem-one}. Let $s_1<\cdots<s_q$ be the even elements of $T$ and $S_j$ the
number of odd elements of $T$ strictly between $s_j$ and $s_{j+1}$ (indices modulo $q$), as there.
Form the numbers of Definition~\ref{pd-kappa-function} with $C_0$, $C_A$ and $C_t$, and put
$I'_a=[\bar a,u_a+J+c_{cl}]\cap\mathbf Z$ and $I'_b=[l_b,u_b+J+c_{cl}]\cap\mathbf Z$. Then for every
unramified quadratic Hida setting of degree $n'$ at $p$ there are a real number $\kappa_P$ and
functions $\kappa_1,\dots,\kappa_q$ on the pairs $(A,B)$ of functions $T\to\mathbf Z_{\geq0}$ such
that:
\begin{enumerate}
\item for $A\in I_a'^T$, $B\in I_b'^T$ and every edge $e$, there is a split family of depth $D_0$
  and size at least $\big(1-\kappa_P-\sum_j\kappa_j(A,B)\big)\dim V^\times(A,B)/[k]$ on
  $(\mathcal C^\times_{t,c_\flat}(e),\mathcal C'^\times_{t,c_\flat}(e))$, where
  $t=\lfloor t_0(A,B)\rfloor$ (Definition~\ref{pd-kappa-product-space});
\item if $I'_a$ and $I'_b$ are not empty, then $\kappa_P\leq C_s|T|N_1/Y$;
\item for every $j$ and every $B\in I_b'^T$,
  $\sum_{A\in I_a'^T}\kappa_j(A,B)\leq|I'_a|^{|T|}\,C_A(S_j+1)N_1/Y$.
\end{enumerate}
\end{lemma}
\leanref{Formalism/SharpRestriction/GrowthBounds/ProjectionLemmaBaseLevel.lean}{Formalism.MainNotions.flipOneThree.admB2}
\leanref{Formalism/SharpRestriction/GrowthBounds/ProjectionLemmaBaseLevel.lean}{Formalism.MainNotions.flipOneThree.admB3}
\leanref{Formalism/SharpRestriction/GrowthBounds/ProjectionLemmaBaseLevel.lean}{Formalism.MainNotions.flipOneThree.PathPairCorrespondence}
\leanref{Formalism/SharpRestriction/GrowthBounds/ProjectionLemmaBaseLevel.lean}{Formalism.MainNotions.flipOneThree.kapPI}
\leanref{Formalism/SharpRestriction/GrowthBounds/ProjectionLemmaBaseLevel.lean}{Formalism.MainNotions.flipOneThree.kapTI}
\leanref{Formalism/SharpRestriction/GrowthBounds/ColumnSpanDensity/LocalHypothesisOfGlobal.lean}{Formalism.MainNotions.flipOneThree.Pf.segment_sum_kappa_bound}
\leanref{Formalism/SharpRestriction/GrowthBounds/ColumnSpanDensity/FactorizationBoundMonoLattices.lean}{Formalism.MainNotions.flipOneThree.Pf.keyIneq_of}

\begin{proof}
Let $T_0$ and, for $C_0\geq T_0$, the constants $C_A$, $C_t$, $C_s$ and $T_3$ be those of
Lemma~\ref{pd-kappa-product-defects}. Let the data and a setting be given, and let $\kappa_P$ and
$\kappa_1,\dots,\kappa_q$ be given by that lemma. Parts (1) and (2) are parts (1) and (2) of that
lemma. For (3), fix $j$ and $B\in I_b'^T$, let $\Sigma_j$ be as there, and write $N=|T|$. By part (3)
of that lemma, $\kappa_j(A,B)$ depends only on the restriction $A'$ of $A$ to $\Sigma_j$, and each
$A'\in I_a'^{\Sigma_j}$ is the restriction of $|I'_a|^{N-S_j-1}$ functions $A\in I_a'^T$, since
$\Sigma_j$ has $S_j+1$ elements. By part (4),
\[
 \sum_{A\in I_a'^T}\kappa_j(A,B)=|I'_a|^{N-S_j-1}\sum_{A'\in I_a'^{\Sigma_j}}\kappa_j(A',B)
 \leq|I'_a|^{N-S_j-1}\,|I'_a|^{S_j+1}\,\frac{C_A(S_j+1)N_1}Y=|I'_a|^{N}\,\frac{C_A(S_j+1)N_1}Y .
\]

% Lean: (3) is segment_sum_kappa_bound, through sum_pi_marginal_injective (sgTop_inj) and
%   sum_pi_option_split, from segment_mean_kappa_bound, which is part (4) of
%   pd-kappa-product-defects. kappa_P = kapPI and kappa_j = kapTI jj, so that
%   Dat.kx = kapPI + sum_jj kapTI.
\end{proof}

\begin{lemma}\label{pd-kappa-product-rank}
There is $T_0>0$ such that for every $C_0\geq T_0$ there are constants $C_A,C_t,C_s>0$ and $T_3\geq0$
with the following property. Let $C_H\geq T_3$, and let $p$, $n'$, $m$, $G'$, $N_1$, $c_{cl}$,
$D_p$, $c_*$, $T$, $Y$, $Y_b$, $a$, $b'$, $b_l$, $b_u$ satisfy hypotheses (1) to (3) of
Theorem~\ref{pd-gain-theorem-one}. Form the numbers of Definition~\ref{pd-kappa-function} with
$C_0$, $C_A$ and $C_t$, and put $I'_a=[\bar a,u_a+J+c_{cl}]\cap\mathbf Z$ and
$I'_b=[l_b,u_b+J+c_{cl}]\cap\mathbf Z$. Then for every unramified quadratic Hida setting of degree
$n'$ at $p$ there is a function $\kappa_\times$ on the pairs $(A,B)$ of functions
$T\to\mathbf Z_{\geq0}$ such that:
\begin{enumerate}
\item for $A\in I_a'^T$, $B\in I_b'^T$ and every edge $e$, there is a split family of depth $D_0$
  and size at least $(1-\kappa_\times(A,B))\dim V^\times(A,B)/[k]$ on
  $(\mathcal C^\times_{t,c_\flat}(e),\mathcal C'^\times_{t,c_\flat}(e))$, where
  $t=\lfloor t_0(A,B)\rfloor$ (Definition~\ref{pd-kappa-product-space});
\item the mean of $\kappa_\times$ over $I_a'^T\times I_b'^T$ is at most $(C_A+C_s)|T|N_1/Y$.
\end{enumerate}
\end{lemma}
\leanref{Formalism/SharpRestriction/GrowthBounds/ProjectionLemmaBaseLevel.lean}{Formalism.MainNotions.flipOneThree.AdmBE}
\leanref{Formalism/SharpRestriction/GrowthBounds/ProjectionLemmaBaseLevel.lean}{Formalism.MainNotions.flipOneThree.AdmB}
\leanref{Formalism/SharpRestriction/GrowthBounds/ProjectionLemmaBaseLevel.lean}{Formalism.MainNotions.flipOneThree.PathPairCorrespondence}
\leanref{Formalism/SharpRestriction/GrowthBounds/ProjectionLemmaBaseLevel.lean}{Formalism.MainNotions.flipOneThree.admB2}
\leanref{Formalism/SharpRestriction/GrowthBounds/ProjectionLemmaBaseLevel.lean}{Formalism.MainNotions.flipOneThree.admB3}
\leanref{Formalism/SharpRestriction/GrowthBounds/PlaneCountTheoremWithConstants.lean}{Formalism.MainNotions.G.first_theorem_admitted_closed}
\leanref{Formalism/SharpRestriction/GrowthBounds/ColumnSpanDensity/RhoLowerBoundLargeConductor.lean}{Formalism.MainNotions.flipOneThree.Pf.Dat.kx}
\leanref{Formalism/SharpRestriction/GrowthBounds/ColumnSpanDensity/LocalHypothesisOfGlobal.lean}{Formalism.MainNotions.flipOneThree.Pf.meanX_of}

\begin{proof}
Let $T_0$ and, for $C_0\geq T_0$, the constants $C_A$, $C_t$, $C_s$ and $T_3$ be those of
Lemma~\ref{pd-kappa-product-segments}. Let the data and a setting be given, take $\kappa_P$ and
$\kappa_1,\dots,\kappa_q$ of that lemma, and put $\kappa_\times(A,B)=\kappa_P+\sum_{j=1}^q\kappa_j(A,B)$.
Then (1) is part (1) of Lemma~\ref{pd-kappa-product-segments}.

For (2), write $N=|T|$. If $I'_a$ or $I'_b$ is empty there is nothing to prove, the mean over the empty
set being $0$. Otherwise, by parts (2) and (3) of Lemma~\ref{pd-kappa-product-segments},
\[
 \sum_{A\in I_a'^T}\sum_{B\in I_b'^T}\kappa_\times(A,B)
 =|I'_a|^N|I'_b|^N\kappa_P+\sum_{j=1}^q\sum_{B\in I_b'^T}\sum_{A\in I_a'^T}\kappa_j(A,B)
 \leq|I'_a|^N|I'_b|^N\Big(\frac{C_sNN_1}Y+\sum_{j=1}^q\frac{C_A(S_j+1)N_1}Y\Big).
\]
The even elements $s_1,\dots,s_q$ and the odd elements of $T$ strictly between consecutive ones are
distinct elements of $T$, so $\sum_j(S_j+1)\leq N$. Dividing by $|I'_a|^N|I'_b|^N$, the mean of
$\kappa_\times$ is at most $(C_A+C_s)NN_1/Y$.
\end{proof}

\begin{lemma}\label{pd-kappa-numerical}
Let $T_0$ be as in Lemma~\ref{pd-kappa-product-rank}, let $C_0\geq T_0$, and let $C_A$ and $C_t$ be
the constants given there. There are $T_4\geq0$ and $C_n>0$ such that the following hold for all
$C_H\geq T_4$, all $C\geq C_n$ and all data satisfying hypotheses (1) to (3) of
Theorem~\ref{pd-gain-theorem-one}, with the numbers of Definition~\ref{pd-kappa-function},
$I_a(\tau)=[\bar a+\tau,u_a+\tau]\cap\mathbf Z$ and $I_b(\tau)=[l_b+\tau,u_b+\tau]\cap\mathbf Z$ for
$\tau\geq0$, and $I'_a$ and $I'_b$ as in Lemma~\ref{pd-kappa-product-rank}.
\begin{enumerate}
\item $m\geq2$, $8(D_p+1)\leq c_*$, $D_0\geq0$, $4D_1<c_\flat$ and
  $2D_1+|T|\lceil m^2\ell\rceil\leq D_p$.
\item $2c_{cl}\leq8\bar a$, $b'/3\leq\bar a$, $b'/3\leq b_l$, $b_u\geq0$ and
  $12|T|kc_{cl}/b'\leq\min(1/9,1/Y)$.
\item $J\leq\bar a$, $J\leq l_b$ and $(u_a+J+c_{cl})+(u_b+J+c_{cl})+1\leq8\bar a$.
\item $I_a(0)$ and $I_b(0)$ are not empty, and there is $\delta\geq0$ such that
  \[
   \sum_{j=0}^{J}\frac{\operatorname{Har}_m(\alpha,\beta)}{P_m(\alpha+j)P_m(\beta+j)}\geq\delta
  \]
  for all $0\leq\tau\leq c_{cl}$, $\alpha\in I_a(\tau)$ and $\beta\in I_b(\tau)$, and
  $|I'_a|^{|T|}|I'_b|^{|T|}\leq2\delta^{|T|}|I_a(0)|^{|T|}|I_b(0)|^{|T|}$
  (Lemma~\ref{pd-kappa-dimensions}).
\item $|I_a(0)|^{|T|}|I_b(0)|^{|T|}-|I_a(c_{cl})\cap I_a(0)|^{|T|}|I_b(c_{cl})\cap I_b(0)|^{|T|}
  \leq|I_a(0)|^{|T|}|I_b(0)|^{|T|}/Y$.
\item For all $a_w,b_w\colon T\to\mathbf Z$ with values in $I_a(0)$ and $I_b(0)$,
  $t^*(a_w+c_{cl},b_w+c_{cl})+2n'c_{cl}\leq t_I$, where $t_I$ is formed with $C$ as in
  Theorem~\ref{pd-gain-theorem-one}.
\end{enumerate}
\end{lemma}
\leanref{Formalism/SharpRestriction/GrowthBounds/ColumnSpanDensity/RhoLowerBoundLargeConductor.lean}{Formalism.MainNotions.flipOneThree.Pf.FactsLP}
\leanref{Formalism/SharpRestriction/GrowthBounds/ColumnSpanDensity/ThrLPOf.lean}{Formalism.MainNotions.flipOneThree.Pf.thrLP}
\leanref{Formalism/SharpRestriction/GrowthBounds/ColumnSpanDensity/MidMeanOf.lean}{Formalism.MainNotions.flipOneThree.Pf.FM}
\leanref{Formalism/SharpRestriction/GrowthBounds/ColumnSpanDensity/FactorizationBoundMonoLattices.lean}{Formalism.MainNotions.flipOneThree.Pf.FixedWeightNumericalFacts}
\leanref{Formalism/SharpRestriction/GrowthBounds/ColumnSpanDensity/DensityNumericalFacts.lean}{Formalism.MainNotions.flipOneThree.Pf.DensityNumericalFacts}
\leanref{Formalism/SharpRestriction/GrowthBounds/ColumnSpanDensity/LocalHypothesisOfGlobal.lean}{Formalism.MainNotions.flipOneThree.Pf.MeanKappaNumericalFacts}
\leanref{Formalism/SharpRestriction/GrowthBounds/ColumnSpanDensity/ColumnSpanNumericalFacts.lean}{Formalism.MainNotions.flipOneThree.Pf.ColumnSpanNumericalFacts}
\leanref{Formalism/SharpRestriction/GrowthBounds/ColumnSpanDensity/TheoremIOfLatticeFunctorComparison.lean}{Formalism.MainNotions.flipOneThree.Pf.thrFM}

The lemma lists the facts that the book uses. Alpöge's records of them contain further facts, among
them $b'>0$ and $p>0$, which are not stated here.
% Departs from the Lean: FactsLP also has bpos and ppos; the parts of FM also carry HypI (with
% Yb/2), CBoundI, the bounds m0nn and m0le on m0I, windowLowerBound, logl and the nonnegativity
% facts of ColumnSpanNumericalFacts; not yet stated.

\begin{proposition}\label{pd-kappa-projection}
Let $C_0,C_A,C_t,C>0$, and let $p$, $n'$, $m$, $G'$, $N_1$, $c_{cl}$, $D_p$, $c_*$, $T$, $Y$, $a$,
$b'$, $b_l$, $b_u$ be as in Definition~\ref{pd-kappa-function} and satisfy conditions (1), (2) and
(6) of Lemma~\ref{pd-kappa-numerical}, with the numbers of Definition~\ref{pd-kappa-function} and
$t_I$ formed with $C$. Let $\mathcal S$ be an unramified quadratic Hida setting of degree $n'$ at $p$,
let $a_w,b_w\colon\mathbf Z/n'\to\mathbf Z$ and dominant weights $\lambda_0(s)\in\mathbf Z^k$ satisfy
$\lambda_0(s)_1=a_w(s)$, $\lambda_0(s)_k=-b_w(s)$, $|\lambda_0(s)_i|\leq c_{cl}$ for $1<i<k$,
$a_w(s)\in I_a$ and $b_w(s)\in I_b$ when $s\in T$, and $|\lambda_0(s)_i|\leq c_{cl}$ for all $i$
when $s\notin T$, where $I_a=[\bar a,u_a]\cap\mathbf Z$ and $I_b=[l_b,u_b]\cap\mathbf Z$, and let
$t\geq t_I$ be an integer. Then
$\mathrm{LP}\big(1,\lambda_0;t,D_p,\kappa_{\mathcal S}(a_w|_T,b_w|_T),c_*\big)$ holds in $\mathcal S$.
\end{proposition}
\leanref{Formalism/SharpRestriction/GrowthBounds/ColumnSpanDensity/RhoLowerBoundLargeConductor.lean}{Formalism.MainNotions.flipOneThree.Pf.MidLP}
\leanref{Formalism/SharpRestriction/GrowthBounds/ColumnSpanDensity/LatticeCountPrimeMonoLevel.lean}{Formalism.MainNotions.flipOneThree.Pf.midLP_of}

\begin{proof}
Write $\kappa=\kappa_{\mathcal S}(a_w|_T,b_w|_T)$. Since $8(D_p+1)\leq c_*$, we have $8D_p<c_*$ and
$4(D_p+1)\leq c_*$. If $\kappa\geq1$, the projection property holds by
Lemma~\ref{pd-kappa-rank-monotone}(4).

Let $\kappa<1$. The other two terms of $\kappa$ are nonnegative, so $\kappa_-<1$, and
$a_w(s)-c_{cl}\geq\bar a$ and $b_w(s)-c_{cl}\geq b_l$ for $s\in T$. Put $A=a_w-c_{cl}$ and
$B=b_w-c_{cl}$ on $T$. They take values at least $\bar a\geq b'/3$ and $b_l\geq b'/3$, by condition
(2), so $(a_w-c_{cl})^+=A$, $(a_w+c_{cl})^+=A+2c_{cl}$, and the same for $B$. Thus
\[
 \kappa=\kappa_1+2\kappa_2+\frac{12|T|kc_{cl}}{b'},\qquad\kappa_1=\kappa_\flat(A,B),\quad
 \kappa_2=\kappa_\flat(A+2c_{cl},B+2c_{cl}).
\]
Put $t^*=t^*(A+2c_{cl},B+2c_{cl})$. Each term $\min(B(y),A(y-d(y))/p^{d(y)})$ of $t_0$ is
nondecreasing in $A$ and $B$, so $t^*(A,B)\leq t^*$. By the definition of the flat defects and
Lemma~\ref{pd-kappa-rank-monotone}(1),
\[
 \mathrm{ER}\big(\phi(A,B);t^*,c_\flat,D_1,(1-\kappa_1)\mu(\phi(A,B))\big),\qquad
 \mathrm{ER}\big(\phi(A+2c_{cl},B+2c_{cl});t^*,c_\flat,D_1,(1-\kappa_2)\mu(\phi(A+2c_{cl},B+2c_{cl}))\big)
\]
hold, with $0\leq\kappa_1\leq1$ and $\kappa_2\geq0$.

The top part $\lambda_0^T$ is dominant, satisfies the conditions on $\lambda_0$ in
Lemma~\ref{pd-kappa-weight-change}, and vanishes outside $T$. By conditions (1) and (2),
$D_1\geq D_0\geq0$, $4D_1<c_\flat$, $2|T|c_{cl}\leq c_\flat$, $2c_{cl}\leq8\bar a$ and
$12|T|kc_{cl}/b'\leq1/9$. So Lemma~\ref{pd-kappa-weight-change}(2), with the integer $\bar a$ in
place of $a$, gives
$\mathrm{ER}\big(\lambda_0^T;t^*,c_\flat-2|T|c_{cl},2D_1+|T|\lceil m^2\ell\rceil,(1-\kappa)\mu(\lambda_0^T)\big)$.
Since $c_\flat-2|T|c_{cl}=c_*+2n'c_{cl}$, Lemma~\ref{pd-kappa-weight-change}(1) gives
$\mathrm{ER}\big(\lambda_0;t^*+2n'c_{cl},c_*,2D_1+|T|\lceil m^2\ell\rceil,(1-\kappa)\mu(\lambda_0)\big)$.
By condition (6), $t^*+2n'c_{cl}\leq t_I\leq t$, and by condition (1),
$2D_1+|T|\lceil m^2\ell\rceil\leq D_p$. So $\mathrm{ER}(\lambda_0;t,c_*,D_p,(1-\kappa)\mu(\lambda_0))$
holds by Lemma~\ref{pd-kappa-rank-monotone}(1), and Lemma~\ref{pd-kappa-rank-criterion} with
$c=c_*$ and $D=D_p$ gives the projection property.
\end{proof}

\begin{proposition}\label{pd-kappa-mean}
Let $C_0,C_A,C_t,C_s>0$, and let $p$, $n'$, $m$, $G'$, $N_1$, $c_{cl}$, $D_p$, $c_*$, $T$, $Y$, $a$,
$b'$, $b_l$, $b_u$ be as in Definition~\ref{pd-kappa-function} and satisfy conditions (1) to (5) of
Lemma~\ref{pd-kappa-numerical}, with the numbers of Definition~\ref{pd-kappa-function} and the sets
$I_a(\tau)$, $I_b(\tau)$, $I'_a$, $I'_b$ of Lemmas~\ref{pd-kappa-numerical}
and~\ref{pd-kappa-product-rank}. Let $\mathcal S$ be an unramified quadratic Hida setting of degree
$n'$ at $p$, and suppose that in $\mathcal S$ there is a function $\kappa_\times$ with properties (1)
and (2) of Lemma~\ref{pd-kappa-product-rank}. Then the mean of $\kappa_{\mathcal S}$ over
$I_a(0)^T\times I_b(0)^T$ is at most $\big(6(C_A+C_s)+8\big)|T|N_1/Y$.
\end{proposition}
\leanref{Formalism/SharpRestriction/GrowthBounds/ColumnSpanDensity/MeanDefs.lean}{Formalism.MainNotions.flipOneThree.Pf.MidMean}
\leanref{Formalism/SharpRestriction/GrowthBounds/ColumnSpanDensity/MidMeanOf.lean}{Formalism.MainNotions.flipOneThree.Pf.midMean_of}
\leanref{Formalism/SharpRestriction/GrowthBounds/ColumnSpanDensity/FactorizationBoundMonoLattices.lean}{Formalism.MainNotions.flipOneThree.Pf.keyIneq_of}
\leanref{Formalism/SharpRestriction/GrowthBounds/ColumnSpanDensity/DensityNumericalFacts.lean}{Formalism.MainNotions.flipOneThree.Pf.avg_of}
\leanref{Formalism/SharpRestriction/GrowthBounds/ColumnSpanDensity/ColumnSpanNumericalFacts.lean}{Formalism.MainNotions.flipOneThree.Pf.mean_of}

\begin{proof}
Write $N=|T|$ and, for $j\colon T\to\mathbf Z_{\geq0}$, $\mu_j(A,B)=\mu(\phi(A-j,B-j))$.

\emph{Step 1: for $A\in I_a'^T$ and $B\in I_b'^T$,}
\[
 \sum_{j\in\{0,\dots,J\}^T}\kappa_\flat(A-j,B-j)\,\mu_j(A,B)\leq\kappa_\times(A,B)\,
 \frac{\dim V^\times(A,B)}{[k]}.\tag{$\ast$}
\]
By Lemma~\ref{pd-kappa-edge-count} there is an edge $e_0$. Let $t_1=\lfloor t_0(A,B)\rfloor$. By
property (1) of $\kappa_\times$ there is a split family of depth $D_0$ and size at least
$(1-\kappa_\times(A,B))\dim V^\times(A,B)/[k]$ on the cores of $V^\times(A,B)$ at $e_0$ with $t_1$ and
$c_\flat$. By condition (3), $A(s)+B(s)+1\leq8\bar a$, and $D_0\geq0$ by condition (1). So
Lemma~\ref{pd-kappa-harmonic-split} gives integers $n(j)\geq0$, for $j\leq\min(A,B)$, with
\[
 \sum_{j\leq\min(A,B)}n(j)\geq(1-\kappa_\times(A,B))\frac{\dim V^\times(A,B)}{[k]},
\]
and split families of depth $D_0+4Nu_0=D_1$ and size $n(j)$ on the cores of $\phi(A-j,B-j)$ at $e_0$
with $t_1$ and $c_\flat$. By Lemma~\ref{pd-kappa-edge-move} there are such families at every edge.
Since $2D_1<c_\flat$, Lemma~\ref{pd-kappa-rank-upper} gives $n(j)\leq\mu_j(A,B)$.

Let $j\in\{0,\dots,J\}^T$. By condition (3), $J\leq\bar a\leq A(s)$ and $J\leq l_b\leq B(s)$, so
$j\leq\min(A,B)$. Replacing $A$ and $B$ by $A-j$ and $B-j$ lowers each of the at most $N$ terms
$\min(B(y),A(y-d(y))/p^{d(y)})$ by at most $J\leq J_0$, so $t_0(A,B)\leq t_0(A-j,B-j)+NJ_0$ and
$t_1\leq t^*(A-j,B-j)$. If $n(j)>0$, Lemma~\ref{pd-kappa-rank-monotone}(1) gives
$\mathrm{ER}(\phi(A-j,B-j);t^*(A-j,B-j),c_\flat,D_1,n(j))$, which is the edge rank condition with
$(1-\kappa)\mu_j(A,B)$ for $\kappa=1-n(j)/\mu_j(A,B)\in[0,1)$; so
$\kappa_\flat(A-j,B-j)\leq1-n(j)/\mu_j(A,B)$. In all cases
$n(j)\leq(1-\kappa_\flat(A-j,B-j))\mu_j(A,B)$, because $\kappa_\flat\leq1$. Summing over
$j\leq\min(A,B)$ and using $\sum_{j\leq\min(A,B)}\mu_j(A,B)=\dim V^\times(A,B)/[k]$, which follows
from Lemma~\ref{pd-kappa-dimensions}, gives
\[
 (1-\kappa_\times(A,B))\frac{\dim V^\times(A,B)}{[k]}\leq\frac{\dim V^\times(A,B)}{[k]}
 -\sum_{j\in\{0,\dots,J\}^T}\kappa_\flat(A-j,B-j)\,\mu_j(A,B),
\]
which is ($\ast$).

\emph{Step 2: for $\tau\in\{0,c_{cl}\}$, the mean of $\kappa_\flat$ over
$I_a(\tau)^T\times I_b(\tau)^T$ is at most $2(C_A+C_s)NN_1/Y$.} By Lemma~\ref{pd-kappa-dimensions},
$\mu_j(A,B)[k]/\dim V^\times(A,B)=\prod_{s\in T}\rho(A(s),B(s),j(s))$ with
$\rho(\alpha,\beta,i)=\operatorname{Har}_m(\alpha-i,\beta-i)/(P_m(\alpha)P_m(\beta))$. Dividing
($\ast$) by $\dim V^\times(A,B)/[k]>0$ gives
\[
 \sum_{j\in\{0,\dots,J\}^T}\kappa_\flat(A-j,B-j)\prod_{s\in T}\rho(A(s),B(s),j(s))\leq\kappa_\times(A,B).
\]
For $\alpha\in I_a(\tau)^T$, $\beta\in I_b(\tau)^T$ and $j\in\{0,\dots,J\}^T$ we have
$\alpha+j\in I_a'^T$ and $\beta+j\in I_b'^T$, because $\tau\leq c_{cl}$. Sum the last inequality over
$(A,B)\in I_a'^T\times I_b'^T$ and keep on the left only the terms with $A=\alpha+j$ and $B=\beta+j$;
all terms are nonnegative. This gives
\[
 \sum_{\alpha,\beta}\kappa_\flat(\alpha,\beta)\prod_{s\in T}\sum_{i=0}^{J}
 \frac{\operatorname{Har}_m(\alpha(s),\beta(s))}{P_m(\alpha(s)+i)P_m(\beta(s)+i)}
 \leq\sum_{A,B}\kappa_\times(A,B),
\]
the first sum over $I_a(\tau)^T\times I_b(\tau)^T$ and the second over $I_a'^T\times I_b'^T$. By
condition (4) each inner sum is at least $\delta$, and $\delta>0$ because $I'_a$ and $I'_b$ are not
empty. With property (2) of $\kappa_\times$ and condition (4),
\[
 \delta^N\sum_{\alpha,\beta}\kappa_\flat(\alpha,\beta)\leq|I'_a|^N|I'_b|^N\frac{(C_A+C_s)NN_1}Y
 \leq2\delta^N|I_a(0)|^N|I_b(0)|^N\frac{(C_A+C_s)NN_1}Y .
\]
Since $|I_a(\tau)|=|I_a(0)|$ and $|I_b(\tau)|=|I_b(0)|$, this is the claim.

\emph{Step 3.} Let $U=|I_a(0)|^N|I_b(0)|^N$ and sum $\kappa_{\mathcal S}$ over
$I_a(0)^T\times I_b(0)^T$. The pairs with $a_w-c_{cl}\geq\bar a$ and $b_w-c_{cl}\geq b_l$ on $T$
contribute $\kappa_-=\kappa_\flat(a_w-c_{cl},b_w-c_{cl})$ at distinct points of
$I_a(0)^T\times I_b(0)^T$, so in total at most $2U(C_A+C_s)NN_1/Y$ by Step 2 with $\tau=0$. These are
the pairs in $(I_a(c_{cl})\cap I_a(0))^T\times(I_b(c_{cl})\cap I_b(0))^T$, so by condition (5) the
other pairs number at most $U/Y$, and they contribute $\kappa_-=1$ each. The terms
$2\kappa_\flat(a_w+c_{cl},b_w+c_{cl})$ sum to twice the sum of $\kappa_\flat$ over
$I_a(c_{cl})^T\times I_b(c_{cl})^T$, at most $4U(C_A+C_s)NN_1/Y$ by Step 2. The constant term is at
most $1/Y$ by condition (2). Since $NN_1\geq1$, the sum is at most
\[
 U\Big(\frac{6(C_A+C_s)NN_1}Y+\frac2Y\Big)\leq U\big(6(C_A+C_s)+8\big)\frac{NN_1}Y .
\]
\end{proof}

\begin{theorem}\label{pd-gain-theorem-one}
There are $T_0>0$ and functions $T_2,C\colon[T_0,\infty)\to\mathbf R_{>0}$ with the following
property. Let $C_0\geq T_0$ and $C_H\geq T_2(C_0)$, put $C=C(C_0)$, and let $p$ be a prime,
$n',m,G',N_1,c_{cl},D_p,c_*\geq0$ integers, $T\subset\mathbf Z/n'\mathbf Z$ and $Y,Y_b,a,b',b_l,b_u$
reals. Call $s\in\mathbf Z/n'$ even if its representative in $\{0,\dots,n'-1\}$ is even; let the even
elements of $T$ be $s_1<\cdots<s_q$, let $g_j$ be the cyclic distance from $s_j$ to $s_{j+1}$ (indices
modulo $q$; $g_j=n'$ if $q=1$) and $S_j$ the number of odd elements of $T$ strictly between them; for an
odd $y\in T$ let $d(y)\geq1$ be least with $y-d(y)$ an even element of $T$. Put
$\Lambda=C_0N_1^2\log a$, $k=m+1$, $\delta=C_HN_1\Lambda/k+2(b_u-b_l)/b'$,
$I_a=[a,a+a/(kY)]\cap\mathbf Z$ and $I_b=[b_l,b_u]\cap\mathbf Z$. Suppose:
\begin{enumerate}
\item $1\leq Y$, $k^2Y^2\leq p$, $q\geq2$, $g_j\leq G'$ and $S_j\leq N_1$ for every $j$, $1\leq N_1$,
  $2N_1+2\leq m$, and $8\log(|T|Y)\leq\Lambda$;
\item $0<b'$, $b_l\leq b'\leq b_u$, $Y_b|T|\Lambda b'/k\leq b_u-b_l\leq b'/2$, $C_H\leq Y_b$,
  $2Y_b|T|\Lambda<k$, $kY^2Y_b|T|\Lambda\log a\leq a/b'$, $\delta<1/6$, $p<a/b'$,
  $C_H|T|\Lambda^2\leq k$, and no odd power $p^{e}$ lies in $[(a/b')(1-2\delta),(a/b')(1+2\delta)]$;
\item $c_{cl}\leq\Lambda b'/k$, $C_H|T|Y^2(N_1^5+k^3)\log_pa\leq D_p$, $8(D_p+1)\leq c_*$ and
  $C_Hk^3|T|p^{4G'}(c_*+|T|c_{cl}+n'c_{cl})Y\leq b'$.
\end{enumerate}
Then there is a function $\kappa\geq0$ on the pairs $(a_w,b_w)$ of functions $T\to\mathbf Z$ whose mean
over $(a_w,b_w)\in I_a^T\times I_b^T$ is at most $C|T|N_1/Y$, such that for every unramified quadratic
Hida setting of degree $n'$ at $p$, all $a_w,b_w\colon\mathbf Z/n'\to\mathbf Z$, all dominant
weights $\lambda_0(s)\in\mathbf Z^{k}$ with $\lambda_0(s)_1=a_w(s)$, $\lambda_0(s)_k=-b_w(s)$,
$|\lambda_0(s)_i|\leq c_{cl}$ for $1<i<k$, $a_w(s)\in I_a$ and $b_w(s)\in I_b$ when $s\in T$, and
$|\lambda_0(s)_i|\leq c_{cl}$ for all $i$ when $s\notin T$, and every integer $t\geq t_I$, the
projection property
$\mathrm{LP}\big(1,\lambda_0;t,D_p,\kappa(a_w|_T,b_w|_T),c_*\big)$ holds. Here
\[
 t_I=\sum_{y\in T\ \mathrm{odd}}\min\Big(b_w(y),\frac{a_w(y-d(y))}{p^{d(y)}}\Big)+C|T|N_1\Big(\frac a{p^2}
 +\frac{b'}p+\frac{\Lambda a}{pk}+\frac{\Lambda b'}k\Big)+C|T|\big(Qc+kQ+k\log_pa+c_{cl}\big)+2n'c_{cl},
\]
with $Q=p^{G'}$ and $c=c_*+C|T|k\log_pa+4|T|c_{cl}+2n'c_{cl}$.
\end{theorem}
\leanref{Formalism/SharpRestriction/GrowthBounds/ProjectionLemmaBaseLevel.lean}{Formalism.MainNotions.flipOneThree.Conclusion2}
\leanref{Formalism/SharpRestriction/GrowthBounds/ProjectionLemmaBaseLevel.lean}{Formalism.MainNotions.flipOneThree.HypI}
\leanref{Formalism/SharpRestriction/GrowthBounds/ProjectionLemmaBaseLevel.lean}{Formalism.MainNotions.flipOneThree.WeightI}
\leanref{Formalism/SharpRestriction/GrowthBounds/ProjectionLemmaBaseLevel.lean}{Formalism.MainNotions.flipOneThree.tI}
\leanref{Formalism/SharpRestriction/GrowthBounds/PlaneCountTheoremWithConstants.lean}{Formalism.MainNotions.G.first_theorem_threshold_functions}

\begin{proof}
Let $T_0$ be as in Lemma~\ref{pd-kappa-product-rank}. For $C_0\geq T_0$ let $C_A$, $C_t$, $C_s$
and $T_3$ be as there and $T_4$ and $C_n$ as in Lemma~\ref{pd-kappa-numerical}, and put
$T_2(C_0)=\max(T_3,T_4,1)$ and $C(C_0)=\max\big(C_n,6(C_A+C_s)+8\big)$. Let $C_H\geq T_2(C_0)$ and
let the data satisfy (1) to (3). Form the numbers of Definition~\ref{pd-kappa-function} with $C_0$,
$C_A$ and $C_t$; then $I_a=[\bar a,u_a]\cap\mathbf Z$ and $I_b=[l_b,u_b]\cap\mathbf Z$, and
conditions (1) to (6) of Lemma~\ref{pd-kappa-numerical} hold with $C=C(C_0)$.

If there is no unramified quadratic Hida setting of degree $n'$ at $p$, the function $\kappa=0$
has the required properties. Otherwise fix one, $\mathcal S_0$, and put $\kappa=\kappa_{\mathcal S_0}$
(Definition~\ref{pd-kappa-function}). The flat defects are nonnegative, $\kappa_-$ is a flat defect
or $1$, and $b'>0$, so $\kappa\geq0$. By Lemma~\ref{pd-kappa-product-rank} there is a function
$\kappa_\times$ in $\mathcal S_0$ with its properties (1) and (2), so
Proposition~\ref{pd-kappa-mean} bounds the mean of $\kappa$ over $I_a^T\times I_b^T$ by
$(6(C_A+C_s)+8)|T|N_1/Y\leq C|T|N_1/Y$. Let a setting $\mathcal S$, functions $a_w,b_w$, weights
$\lambda_0$ and an integer $t\geq t_I$ be as in the statement. Proposition~\ref{pd-kappa-projection}
gives $\mathrm{LP}(1,\lambda_0;t,D_p,\kappa(a_w|_T,b_w|_T),c_*)$ in $\mathcal S_0$, and
Lemma~\ref{pd-kappa-setting-change} carries it to $\mathcal S$.
% Lean: first_theorem_threshold_functions = first_theorem_uniform_thresholds and
% first_theorem_with_constants = Loc8All.p13B applied to first_theorem_admitted_closed. Here p13B is
% flipOneThree.Pf.theorem_I_closed_form_holds = theorem_I_of_pairing_GL -> theorem_I_of_contraction
% -> theorem_I_of_two_letter_clauses_contraction -> theorem_I_of_explicit_minor_map ->
% theorem_I_of_model_identification -> theorem_I_of_lattice_functor_comparison ->
% theorem_I_of_finer_open_statements -> second_estimate_of_restrictions ->
% theorem_I_of_intermediate_bounds, which is the assembly above: kappa = Dat.kText at a fixed setting
% R0, hR = lpRingIndep_of'', hLP = midLP_of, hMean = midMean_of, hN1 = thrLP, hN2 = thrFM.
% Lemmas pd-kappa-weight-change, pd-kappa-harmonic-split, pd-kappa-cloud-maps,
% pd-kappa-harmonic-single and pd-kappa-product-rank are proved above, and so are
% pd-kappa-cloud-contraction, pd-kappa-harmonic-model and pd-kappa-product-segments. What remains open
% is Lemmas pd-kappa-cloud-transpose and pd-kappa-cloud-gram (TauP and SmithP, the open step in the
% proof of pd-kappa-cloud-pairing) and, below Lemma pd-kappa-product-defects, Lemmas
% pd-kappa-defects-windows-model (ExtendedModelRank) and pd-kappa-defects-windows-lift (LiftVecAt),
% behind pd-kappa-defects-windows (HerrExtensionClass), and pd-kappa-defects-segment-mean
% (PathPairCorrespondence). Lemma pd-kappa-defects-sandwich-core (CoreQp) is proved.
\end{proof}

\subsection{Group coherence and the gain law}\label{subsection-pd-gain-law}

\begin{definition}\label{pd-gain-coherence}
Let $D$ be a definite-side datum with ranks $k,m,t,n$, host with CM field $E$, totally real subfield
$F$, complex conjugation $\sigma$ and level primes $g_0\neq g_1$, and with local models $G(u,w)$, a
tie of them and a comparison $\operatorname{cmp}$ (Definition~\ref{pd-lattice-instance}). Let
$J\in GL_k(E)$ be the Gram matrix of the space $Y_0$ of $D$, $U_J$ the unitary group of $J$ over $F$,
and $\mathbf A_f$ the finite adeles of $F$. \emph{Group coherence} for $D$ consists of the following
data, subject to the compatibilities stated with them.
\begin{enumerate}
\item \emph{Levels in degree $0$.} For every word $u$, an open compact subgroup
  $K_u=\prod_{w'}K_{u,w'}$ of $U_J(\mathbf A_f)$ that stabilizes an $\mathcal O_E$-lattice of
  $E^k$ at every place $w'$ of $F$ and whose component at some place over $g_i$ lies, after
  conjugation, in the integral matrices that are unipotent upper triangular modulo that place
  ($i=0,1$). For every place $w$ admissible for $u$, with $p=p_w$, a \emph{tie} of $K_u$ to
  $G(u,w)$: the rank of $G(u,w)$ is $k$; a bijection $v\mapsto w'_v$ from the orbits of $G(u,w)$ to
  the places of $F$ over $p$, with isomorphisms $Q_v\cong F_{w'_v}$ that carry $W_v$ onto the
  integers; splittings $E\otimes_FF_{w'}\cong F_{w'}\times F_{w'}$ that exchange the factors under
  $\sigma$; matrices $\mathrm{fr}_{w'}\in GL_k(F_{w'})$; for every component $c$ a point
  $y_c\in U_J(\mathbf A_f)$ with trivial components over $p$, such that $c\mapsto y_c$ is a bijection
  onto $U_J(F)\backslash U_J(\mathbf A_f)/K^{(p)}_u$, where $K^{(p)}_u$ is the group of
  $x\in U_J(\mathbf A_f)$ with $x_{w'}\in K_{u,w'}$ for all $w'\nmid p$; and isomorphisms of the
  groups $\Gamma_c$ of the levels of $G(u,w)$ with $U_J(F)\cap y_cK^{(p)}_uy_c^{-1}$, under which
  $\gamma\in\Gamma_c$ acts on $Q_v^k$ by $\mathrm{fr}_{w'_v}^{-1}\gamma_1\mathrm{fr}_{w'_v}$, where
  $\gamma_1$ is the first factor of $\gamma$ under the splitting. The classes
  $U_J(F)\backslash U_J(\mathbf A_f)/K_u$ correspond bijectively to the finite set $\mathrm{cls}_0$
  of classes in degree $0$ of $C_u$.
\item \emph{Forms in degree $0$.} For every word $u$, an identification of the forms in degree $0$
  of $u$ with functions on the classes of the level $K_u$, compatible with the Hecke data at $T_u$,
  with the values and norms of the forms and with the weights $\lambda_{u,0}$.
\item \emph{Rational comparison.} A matrix $B\in GL_k(E)$ with ${}^t\sigma(B)JB=1$ whose columns
  are an $\mathcal O_E$-basis of the self-dual lattice in degree $0$, such that every
  $\mathrm{fr}_{w'}$ is the first factor of the image of $B$; and, for every $(u,w)$, a tie of
  $\operatorname{cmp}_{u,w}$ to the forms of (2): for $x\in V_{u,0}$, a component $c$ and a lattice
  tuple $L$ there is a polynomial $\operatorname{val}(x,c,L)$ with coefficients in the coefficient
  field $L_D$ of $D$ with $\operatorname{cmp}_{u,w}(x)(c,L)=\epsilon_w(\operatorname{val}(x,c,L))$,
  whose values at the archimedean points of the level are the values of $x$ times the character of
  the level.
\item \emph{Standard rings.} For every $(u,w)$ and every orbit $v$ of $G(u,w)$, with $n'=n'_v$:
  isomorphisms of $W_v$ with the Witt vectors $W(\mathbf F_{p^{n'}})$, of $Q_v$ with their fraction
  field, and of $\mathcal R\subset\mathcal K$ with the valuation ring and field of a finite
  extension, determined by $w$ and $n'$, of the ring $\mathcal R_0=W(\mathbf F_{p^{2n'}})[\omega]$,
  $\omega^2=p$; these carry $\iota_v$ and the inclusion of $W_v$ in $\mathcal R$ to the standard maps,
  and $\omega$ to $\varpi_0$.
\item \emph{Local models of the levels.} Every $G(u,w)$ is the local model at $p_w$ built from the
  level $K_u$: its orbits, components and groups are those of the tie (1), and for every orbit $v$,
  an element $x\in U_J(F_{w'_v})$ lies in $K_{u,w'_v}$ if and only if
  $\mathrm{fr}_{w'_v}^{-1}x_1\mathrm{fr}_{w'_v}$ and its inverse have integral entries.
\end{enumerate}
\end{definition}
\leanref{Formalism/LocalToGlobal/GroupLattices/StdSubAtOfRing.lean}{Formalism.L2.G.GroupCompatibilityAll}
\leanref{Formalism/LocalToGlobal/HermitianColumns/BaseModelLevelEq.lean}{Formalism.L2.Herm.Herm0}
\leanref{Formalism/LocalToGlobal/HermitianColumns/CompletionIntegersIff.lean}{Formalism.L2.Herm.GLevTie}
\leanref{Formalism/LocalToGlobal/GroupLattices/DownNaturalLevelZero.lean}{Formalism.L2.Forms.Forms0}
\leanref{Formalism/LocalToGlobal/GroupLattices/ComparisonMaps/TiedRec.lean}{Formalism.L2.G.Cmp.CmpTied}
\leanref{Formalism/LocalToGlobal/GroupLattices/StdSubAtOfRing.lean}{Formalism.L2.G.RingCompatibleAt}
\leanref{Formalism/LocalToGlobal/GroupLattices/StdSubAtOfRing.lean}{Formalism.L2.G.GatIs}

Item (2) is stated in summary: Alpöge's data in degree $0$ have further fields (the classes and
values of the forms, their naturality under the Hecke data and their norms), which this chapter
uses only through the tie of (3).
% Departs from the Lean: item (2) summarizes Formalism.L2.Forms.Forms0 (atT, cls0, val0, gl, glev0,
% lv0, fMe0, fv0, the up and down naturality ties, norm0, bw0); not yet written out field by field.

\begin{lemma}\label{pd-gain-projection-monotone}
Let $p\geq5$ be a prime, $(W,Q,\sigma,\iota)$ a Hida setting of degree $n'$ at $p$, $r\geq0$,
$\lambda_0(s)\in\mathbf Z^k$ dominant weights, $t,D,c_*\geq0$ integers and $\kappa\leq\kappa'$ real
numbers (Definition~\ref{pd-projection-property}). If $\mathrm{LP}(r,\lambda_0;t,D,\kappa,c_*)$
holds, then so do $\mathrm{LP}(r,\lambda_0;t,D,\kappa',c_*)$ and
$\mathrm{LP}(r,\lambda_0;t,D,\min(1,\kappa),c_*)$.
\end{lemma}
\leanref{Formalism/SharpRestriction/GrowthBounds/Joints.lean}{Formalism.MainNotions.G.Joints.LPprime_mono_kappa}
\leanref{Formalism/SharpRestriction/GrowthBounds/Joints.lean}{Formalism.MainNotions.G.Joints.LPprime_min_one}

\begin{proof}
The number $\kappa$ enters the projection property only through the condition
$n\geq(1-\kappa)\dim V$, which becomes weaker as $\kappa$ grows; this gives the first statement, and
the second when $\kappa\leq1$. If $\kappa>1$, the same edges, vectors, $n$, $u$ and $(B_{ji})$
satisfy $n\geq0=(1-\min(1,\kappa))\dim V$.
\end{proof}

\begin{lemma}\label{pd-gain-good-segments}
Let $x>1$ and $0<\varepsilon<1/52$, with $u$, $G'$ and $N_1$ as in Definition~\ref{hf-good-orbit},
and let $n'\geq1$ and $L_2<u$ be integers with $3G'\leq n'$. Let $P\subset\{0,\dots,n'-1\}$ be a good
set of positions for $(x,\varepsilon,L_2)$, viewed in $\mathbf Z/n'\mathbf Z$. Then $P$ has at least
two even elements; and if $s_1<\dots<s_q$ are its even elements, with $g_j$ and $S_j$ as in
Theorem~\ref{pd-gain-theorem-one}, then $g_j\leq G'$ and $S_j\leq N_1$ for every $j$.
\end{lemma}
\leanref{Formalism/SharpRestriction/GrowthBounds/TopChoiceBounds.lean}{Formalism.MainNotions.G.TopChoiceBounds.orbit_segments}
\leanref{Formalism/SharpRestriction/GrowthBounds/TopChoiceBounds.lean}{Formalism.MainNotions.G.TopChoiceBounds.bs_segs}

\begin{proof}
Since $n'\geq3G'\geq(M_0+1)(6r_\varepsilon+1)u$, there are at least $M_0+1$ tiles, and they cover
$\{0,\dots,n'-1\}$ except fewer than $(6r_\varepsilon+1)u$ positions at the end. Let $0\leq i<n'$.
Since $G'=(M_0+5)(6r_\varepsilon+1)u$, the positions $i+1,\dots,i+G'$, taken modulo $n'$, contain
$M_0+1$ tiles that are consecutive modulo the number of tiles. By condition (2) of
Definition~\ref{hf-good-orbit} one of them is marked, and its second cell contains an even element
of $P$. So for every $i$ there is $1\leq D\leq G'$ such that $i+D$ is, modulo $n'$, an even element
of $P$. In particular $P$ has an even element. For an even element $s\in P$, let $D\geq1$ be least
with $s+D$ an even element of $P$; then $D\leq G'<n'$, so $s+D\neq s$, $P$ has at least two even
elements, and the cyclic distance from $s$ to the next even element is $D\leq G'$. The elements of
$P$ strictly between them lie in the window $W(s+1,G')$, so there are at most $N_1$ of them by
condition (3) of Definition~\ref{hf-good-orbit}.
\end{proof}

\begin{lemma}\label{pd-gain-second-hypotheses}
Let $D$ be a definite-side datum with ranks $k,m,t,n$, size data with scale $(x,\varepsilon,J)$,
cloud size $S_c$ and derived quantities $a,b',\delta_1,p_0$ (Definition~\ref{pd-size-data}), and gain
data with constants $K_G$ (Definition~\ref{pd-gain-data}). Let $K_\Lambda$ and $K_\kappa$ be the
constants of $K_G$ so named there and $K_t,K_D,K_s,K'_\kappa$ four further constants of $K_G$, and
suppose that the constant $C_0$ of Theorem~\ref{pd-gain-theorem-two} is at most each of
$K_\kappa,K_t,K_\Lambda,K_D,K_s,K'_\kappa$. Put $c_G=4S_c$ and suppose $c_G\geq0$ and $a-c_G\geq1$.
Let $r=r_\varepsilon$, $u$, $G'$ and $N_1$ be as in Definition~\ref{hf-good-orbit}, $L_2$ as in
Definition~\ref{pd-gain-data}(b), $\Lambda=K_\Lambda N_1^2\log a$ and $c_0=2b'+2c_G$. For a prime
$p$ and integers $n'\geq1$ and $N\geq0$ put $L_1(p)=\lceil\log_p\big((kx^{26})^4\big)\rceil$,
$Q=p^{G'}$ and
\[
 D_{\mathrm{II}}(p,n')=\Big\lceil K_D\cdot40n'\delta_1\,(N_1^3+r^5+k^2)\log_pa+2rn'c_G
 +40rn'\delta_1c_0\Big\rceil,
\]
\[
 t_{\mathrm{II}}(p,n',N)=K_trN\Big(\frac{\Lambda a}{pk}+ap^{-L_1(p)}+c_0+kQ\Big)+2rn'c_G,\qquad
 \kappa_{\mathrm{II}}(N)=\min\big(1,K_\kappa Nr^2k\,x^{-26}\big).
\]
Let $p$ be a prime with $13\leq p<p_0$, $(W,Q,\sigma,\iota)$ an unramified quadratic Hida setting of
degree $n'$ at $p$ (Definition~\ref{pd-projection-property}), and $T\subset\mathbf Z/n'\mathbf Z$ a
tiled set of positions for $(x,\varepsilon,L_2)$ with $2L_1(p)\leq L_2<u$, $3G'\leq n'$ and
$|T|\leq40n'\delta_1$. Let $a_w,b_w\colon\mathbf Z/n'\mathbf Z\to\mathbf Z$ take values in the
windows $W_a$ and $W_b$ of Definition~\ref{pd-gain-data}(g) on $T$ and satisfy
\[
 a-c_G+\lfloor c_G\rfloor\leq a_w(s)\leq2(a-c_G)-\lfloor c_G\rfloor,\qquad
 b_w(s)+2\lfloor c_G\rfloor\leq c_0\qquad(s\in T),
\]
and let $\lambda_0(s)\in\mathbf Z^k$ be dominant weights with $\lambda_0(s)_1=a_w(s)$,
$\lambda_0(s)_k=-b_w(s)$ and $|\lambda_0(s)_i|\leq\lfloor c_G\rfloor$ for $1<i<k$ when $s\in T$, and
$|\lambda_0(s)_i|\leq\lfloor c_G\rfloor$ for all $i$ when $s\notin T$. Suppose that $r\geq2$,
$N_1^2\leq r|T|$, $r\leq x^{26}$, $k^2x^{78}\leq p^{L_1(p)-1}$, $256k^2|T|x^{52}\leq k^2x^{78}$,
$N_1+3r\leq k$, $K'_\kappa N_1^2\log a\leq k$ and
\[
 K_sk^3r^2|T|\,Q^5\big(8(D_{\mathrm{II}}(p,n')+1)+r|T|c_0+2rn'c_G\big)x^{26}\leq a-c_G .
\]
Then $\mathrm{LP}\big(r,\lambda_0;\lceil t_{\mathrm{II}}(p,n',|T|)\rceil,D_{\mathrm{II}}(p,n'),
\kappa_{\mathrm{II}}(|T|),8(D_{\mathrm{II}}(p,n')+1)\big)$ holds.
\end{lemma}
\leanref{Formalism/SharpRestriction/GrowthBounds/Joints.lean}{Formalism.MainNotions.G.Joints.type_two_hypothesis_of_counts}
\leanref{Formalism/SharpRestriction/GrowthBounds/SmallOrderPrimes.lean}{Formalism.MainNotions.AuxBlock8.adm6}
\leanref{Formalism/SharpRestriction/GrowthBounds/IsResidueClassZero.lean}{Formalism.MainNotions.G.DIIf}
\leanref{Formalism/SharpRestriction/GrowthBounds/IsResidueClassZero.lean}{Formalism.MainNotions.G.tIIf}
\leanref{Formalism/SharpRestriction/GrowthBounds/IsResidueClassZero.lean}{Formalism.MainNotions.G.kapIIf}

\begin{proof}
Apply Theorem~\ref{pd-gain-theorem-two} with $m+1=k$, $r$, the set $T$, $L_1=L_1(p)$, $L_2$,
$\delta=\delta_1$, $h=h_\varepsilon$ and $x$, so that $K=r-1=K_\varepsilon$ and $u,M_0,G',N_1$ are
those of the statement; with $a_T=a_w$, $b_T=b_w$, $a-c_G$ in place of $a$, $Y=x^{26}$,
$c_{pp}=\lfloor c_G\rfloor$, $\lfloor c_0\rfloor$ in place of $c_0$, $D_p=D_{\mathrm{II}}(p,n')$ and
$c_*=8(D_{\mathrm{II}}(p,n')+1)$. Its hypotheses hold.
\begin{enumerate}
\item[(1)] These are hypotheses of the lemma.
\item[(2)] The conditions on $\lambda_0$ and on $a_w$ are hypotheses. On $T$,
  $b_w(s)\geq\lceil b'\rceil>0$ because $b_w(s)\in W_b$, and $b_w(s)+2\lfloor c_G\rfloor\leq
  \lfloor c_0\rfloor$ because the left side is an integer at most $c_0$.
\item[(3)] $r\leq Y$ is a hypothesis. Since $1\leq a-c_G\leq a$ and $C_0\leq K'_\kappa$,
  $C_0N_1^2\log(a-c_G)\leq K'_\kappa N_1^2\log a\leq k$. Next $N_1+2r\leq k-r$, and
  $256k^2|T|Y^2\leq k^2x^{78}\leq p^{L_1-1}$. Since $C_0\leq K_D$, $|T|\leq40n'\delta_1$,
  $k^2\leq N_1^3+r^5+k^2$ and the remaining terms of $D_{\mathrm{II}}$ are nonnegative,
  $C_0|T|k^2\log_p(a-c_G)\leq D_{\mathrm{II}}(p,n')$. The condition $8(D_p+1)\leq c_*$ is an
  equality. The last inequality follows from the displayed hypothesis, since $C_0\leq K_s$,
  $\lfloor c_0\rfloor\leq c_0$ and $\lfloor c_G\rfloor\leq c_G$.
\end{enumerate}
Theorem~\ref{pd-gain-theorem-two} gives
$\mathrm{LP}\big(r,\lambda_0;t,D_{\mathrm{II}},C_0|T|r^2kx^{-26},8(D_{\mathrm{II}}+1)\big)$ for every
integer
\[
 t\geq C_0r|T|\Big(\frac{C_0N_1^2\log(a-c_G)\,(a-c_G)}{pk}+(a-c_G)p^{-L_1}+\lfloor c_0\rfloor+kQ\Big)
 +2rn'\lfloor c_G\rfloor .
\]
The right side is at most $t_{\mathrm{II}}(p,n',|T|)$, because $C_0\leq K_t$, $C_0\leq K_\Lambda$,
$a-c_G\leq a$, $\lfloor c_0\rfloor\leq c_0$ and $\lfloor c_G\rfloor\leq c_G$; so
$t=\lceil t_{\mathrm{II}}\rceil$ is allowed. Finally
$\min(1,C_0|T|r^2kx^{-26})\leq\kappa_{\mathrm{II}}(|T|)$ because $C_0\leq K_\kappa$, and
Lemma~\ref{pd-gain-projection-monotone} replaces $C_0|T|r^2kx^{-26}$ first by its minimum with $1$
and then by $\kappa_{\mathrm{II}}(|T|)$.
\end{proof}

\begin{lemma}\label{pd-gain-first-hypotheses}
Let $D$, the size data, the gain data, $K_G$, $K_\Lambda$, $c_G$, $r$, $u$, $G'$, $N_1$, $L_2$,
$\Lambda$ and $Q$ be as in Lemma~\ref{pd-gain-second-hypotheses}, with $c_G\geq0$, let $K_\rho$ be the
constant of $K_G$ so named in Definition~\ref{pd-gain-data}, $W_a,W_b$ the windows there, and
$N^*=2n\delta_1$. Let $K_{\mathrm I},K_{D,\mathrm I},K_1,K_2$ be further constants of $K_G$ and
$C_H,C\geq0$ and $C_m$ real numbers such that the conclusion of Theorem~\ref{pd-gain-theorem-one}
holds with $C_0=K_\Lambda$, this $C_H$ and $C(C_0)=C$, and
\[
 C_H\leq K_{\mathrm I},\quad C_H\leq K_{D,\mathrm I},\quad2C_H\leq K_\rho,\quad2C\leq K_1,\quad
 C\leq K_2,\quad C\leq C_m .
\]
Put $Y=x^{13}$, $Y_b=x^{\varepsilon}/20$, $I_b=N^*x^{\varepsilon}\Lambda b'/k$ and
$\rho_0=(4N^*x^{\varepsilon}\Lambda+K_\rho N_1\Lambda)/k$, and call a prime $p$ \emph{in the narrow
band} if $p_0\leq p<X'$ and $X'(1-\rho_0)\leq p^e\leq X'(1+\rho_0)$ for some $e\geq1$. For a prime
$p$, integers $n'\geq1$, $N\geq0$, a set $T\subset\mathbf Z/n'\mathbf Z$ and
$a',b'_w\colon\mathbf Z/n'\mathbf Z\to\mathbf Z$ put
\[
 D_{\mathrm I}(p,n')=\big\lceil K_{D,\mathrm I}\cdot40n'\delta_1\,Y^2(N_1^5+k^3)\log_pa\big\rceil,
\]
\[
 c_6=8(D_{\mathrm I}(p,n')+1)+40K_2n'\delta_1k\log_pa+160n'\delta_1c_G+2n'c_G,
\]
\[
 t_{\mathrm I}(p,n',T,a',b'_w)=\sum_{y}\min\Big(b'_w(y),\frac{a'(y-d(y))}{p^{d(y)}}\Big)
 +K_1|T|N_1\Big(\frac a{p^2}+\frac{\Lambda a}{pk}\Big)+K_2|T|\big(Qc_6+kQ+k\log_pa+c_G\big)
 +2n'c_G,
\]
where $y$ runs over the odd elements of $T$ and $d(y)$ is as in Theorem~\ref{pd-gain-theorem-one}
(the sum is $0$ if $T$ has no even element). Say that \emph{Hypothesis I} holds at $(p,n',N)$ if
\[
 k^2Y^2\leq p<a/b',\quad\rho_0<1/6,\quad2N_1+2\leq k-1,\quad8\log(NY)\leq\Lambda,\quad
 Y_bN\Lambda b'/k\leq I_b,\quad2I_b\leq b',\quad K_{\mathrm I}\leq Y_b,
\]
\[
 kY^2Y_bN\Lambda\log a\leq X',\quad K_{\mathrm I}N\Lambda^2\leq k,\quad2Y_bN\Lambda<k,\quad
 c_G\leq\Lambda b'/k,\quad N\leq40n'\delta_1,
\]
\[
 K_{\mathrm I}k^3NQ^4\big(8(D_{\mathrm I}(p,n')+1)+Nc_G+n'c_G\big)Y\leq b' .
\]
Suppose that the local constants $\kappa_I$ of the gain data are least: $\kappa_I(p,n',T,a',b'_w)$
is the infimum of the $\kappa\geq0$ such that
$\mathrm{LP}\big(1,\lambda_0;\lceil t_{\mathrm I}(p,n',T,a',b'_w)\rceil,D_{\mathrm I}(p,n'),\kappa,
8(D_{\mathrm I}(p,n')+1)\big)$ holds for every unramified quadratic Hida setting of degree $n'$ at
$p$ and all dominant weights $\lambda_0$ with $\lambda_0(s)_1=a'(s)$, $\lambda_0(s)_k=-b'_w(s)$ and
$|\lambda_0(s)_i|\leq\lfloor c_G\rfloor$ for $1<i<k$ when $s\in T$, and
$|\lambda_0(s)_i|\leq\lfloor c_G\rfloor$ for all $i$ when $s\notin T$.

Let $p$ be a prime with $p_0\leq p<X'$, not in the narrow band, $n'\geq1$, and
$T\subset\mathbf Z/n'\mathbf Z$ a tiled set of positions for $(x,\varepsilon,L_2)$ with $L_2<u$ and
$3G'\leq n'$ such that Hypothesis I holds at $(p,n',|T|)$. Then $\kappa_I(p,n',T,a',b'_w)\geq0$ for
all $a',b'_w$, and the mean over $T$ of $(a',b'_w)\mapsto\kappa_I(p,n',T,a',b'_w)$
(Definition~\ref{pd-gain-data}(g)) is at most $C_m|T|N_1/x^{13}$.
\end{lemma}
\leanref{Formalism/SharpRestriction/GrowthBounds/Joints.lean}{Formalism.MainNotions.G.Joints.type_one_hypothesis_of_bounds}
\leanref{Formalism/SharpRestriction/GrowthBounds/Joints.lean}{Formalism.MainNotions.G.Joints.type_one_hypothesis_of_segments}
\leanref{Formalism/SharpRestriction/GrowthBounds/Joints.lean}{Formalism.MainNotions.G.Joints.hypI_69_to_13}
\leanref{Formalism/SharpRestriction/GrowthBounds/Joints.lean}{Formalism.MainNotions.G.Joints.tI_le_tIf}
\leanref{Formalism/SharpRestriction/GrowthBounds/SmallOrderPrimes.lean}{Formalism.MainNotions.AuxBlock8.HypI}
\leanref{Formalism/SharpRestriction/GrowthBounds/SmallOrderPrimes.lean}{Formalism.MainNotions.AuxBlock8.adm8}
\leanref{Formalism/SharpRestriction/GrowthBounds/DefsInst.lean}{Formalism.MainNotions.G.KapIDef}
\leanref{Formalism/SharpRestriction/GrowthBounds/IsResidueClassZero.lean}{Formalism.MainNotions.G.DIf}
\leanref{Formalism/SharpRestriction/GrowthBounds/IsResidueClassZero.lean}{Formalism.MainNotions.G.tIf}
\leanref{Formalism/SharpRestriction/GrowthBounds/IsResidueClassZero.lean}{Formalism.MainNotions.G.InBand0}

\begin{proof}
The set of $\kappa$ in the definition of $\kappa_I$ consists of nonnegative numbers and contains
$1$: with no edges and $n=0$ the projection property holds for $\kappa=1$, since
$4(D_{\mathrm I}+1)\leq8(D_{\mathrm I}+1)$. So $0\leq\kappa_I\leq1$.

By Lemma~\ref{pd-gain-good-segments}, $T$ has at least two even elements, and $g_j\leq G'$ and
$S_j\leq N_1$ for every $j$. Apply Theorem~\ref{pd-gain-theorem-one} with $C_0=K_\Lambda$, $C_H$,
$m=k-1$, $Y=x^{13}$, $Y_b=x^\varepsilon/20$, $b_l=b'$, $b_u=b'+I_b$,
$c_{cl}=\lfloor c_G\rfloor$, $D_p=D_{\mathrm I}(p,n')$ and $c_*=8(D_{\mathrm I}(p,n')+1)$; its
$\Lambda=K_\Lambda N_1^2\log a$ is the $\Lambda$ here, and its $\delta$ satisfies
$2\delta=2C_HN_1\Lambda/k+4N^*x^\varepsilon\Lambda/k\leq\rho_0$ because $2C_H\leq K_\rho$. Its
hypotheses hold.
\begin{enumerate}
\item[(1)] $1\leq Y$ since $x>1$. The conditions $k^2Y^2\leq p$, $2N_1+2\leq m$ and
  $8\log(|T|Y)\leq\Lambda$ are in Hypothesis I; $q\geq2$, $g_j\leq G'$ and $S_j\leq N_1$ were shown
  above; and $N_1\geq1$ by its definition.
\item[(2)] $0<b'=b_l\leq b'\leq b_u$. The conditions $Y_b|T|\Lambda b'/k\leq b_u-b_l=I_b$,
  $2I_b\leq b'$, $2Y_b|T|\Lambda<k$ and $p<a/b'$ are in Hypothesis I. Next $C_H\leq K_{\mathrm I}\leq
  Y_b$ and $C_H|T|\Lambda^2\leq K_{\mathrm I}|T|\Lambda^2\leq k$. Since $a/b'=X'$,
  $kY^2Y_b|T|\Lambda\log a\leq a/b'$ is in Hypothesis I. From $2\delta\leq\rho_0<1/6$ we get
  $\delta<1/6$. An odd power $p^e$ in $[X'(1-2\delta),X'(1+2\delta)]$ would lie in
  $[X'(1-\rho_0),X'(1+\rho_0)]$ and put $p$ in the narrow band.
\item[(3)] $c_{cl}\leq c_G\leq\Lambda b'/k$. Since $C_H\leq K_{D,\mathrm I}$ and
  $|T|\leq40n'\delta_1$, $C_H|T|Y^2(N_1^5+k^3)\log_pa\leq D_{\mathrm I}(p,n')$. The condition
  $8(D_p+1)\leq c_*$ is an equality, and the last inequality follows from the last condition of
  Hypothesis I, since $C_H\leq K_{\mathrm I}$ and $c_{cl}\leq c_G$.
\end{enumerate}
Theorem~\ref{pd-gain-theorem-one} gives a function $\kappa\geq0$ on $I_a^T\times I_b^T$ with mean
at most $C|T|N_1/Y$, where $I_a=W_a$ and $I_b=[b',b'+I_b]\cap\mathbf Z=W_b$, such that
$\mathrm{LP}(1,\lambda_0;t,D_{\mathrm I},\kappa(a'|_T,b'_w|_T),8(D_{\mathrm I}+1))$ holds for every
$t\geq t_I$, with $t_I$ the threshold of that theorem.

Let $a',b'_w$ take values in $W_a$ and $W_b$ on $T$. Then $t_I\leq t_{\mathrm
I}(p,n',T,a',b'_w)$. The sums over the odd $y\in T$ agree. Since $p<a/b'$, $b'\leq a/p$, so
$b'/p\leq a/p^2$ and $\Lambda b'/k\leq\Lambda a/(pk)$; hence
$C|T|N_1(a/p^2+b'/p+\Lambda a/(pk)+\Lambda b'/k)\leq2C|T|N_1(a/p^2+\Lambda a/(pk))$, and $2C\leq K_1$.
In the remaining terms $C\leq K_2$, $c_{cl}\leq c_G$, and the number
$c=c_*+C|T|k\log_pa+4|T|c_{cl}+2n'c_{cl}$ of Theorem~\ref{pd-gain-theorem-one} is at most $c_6$
because $|T|\leq40n'\delta_1$. So $\kappa(a'|_T,b'_w|_T)$ belongs to the set whose infimum is
$\kappa_I(p,n',T,a',b'_w)$, and $\kappa_I(p,n',T,a',b'_w)\leq\kappa(a'|_T,b'_w|_T)$.

Restriction to $T$ is a bijection from $B_T(W_a)\times B_T(W_b)$ onto $W_a^T\times W_b^T$, so the
mean of $\kappa_I(p,n',T,\cdot,\cdot)$ over $T$ is at most the mean of $\kappa$, which is at most
$C|T|N_1/x^{13}\leq C_m|T|N_1/x^{13}$.
\end{proof}

\begin{lemma}\label{pd-gain-weyl-column-projector}
Let $K$ be a field of characteristic $0$, $n\geq1$, and $V$ a finite-dimensional $K$-vector space
with endomorphisms $E_{ab}$, $1\leq a,b\leq n$, such that
$E_{ab}E_{cd}-E_{cd}E_{ab}=\delta_{bc}E_{ad}-\delta_{da}E_{cb}$ for all $a,b,c,d$ and $V$ is the sum
of the weight spaces $V_\mu=\{x\in V:E_{aa}x=\mu_ax\text{ for all }a\}$, $\mu\in\mathbf Z^n$. Let
$L\geq0$ be an integer such that $|\mu_a|\leq L$ for every $a$ and every $\mu$ with $V_\mu\neq0$ and
$\mu_1+\cdots+\mu_n=0$. For $1\leq c\leq n$ let $C_c$ be the endomorphism
$T\mapsto\sum_{i<c}E_{ci}TE_{ic}$ of $\operatorname{End}_K(V)$, and put
\[
 P_c=\sum_{q=0}^{L}\frac{(-1)^q}{q!\,c(c+1)\cdots(c+q-1)}\,C_c^q(1).
\]
\begin{enumerate}
\item Let $u\in V_0$ with $E_{ab}u=0$ for all $a,b<c$. Then $P_cu\in V_0$, $E_{ab}P_cu=0$ for all
  $a,b<c$, $E_{jc}P_cu=0$ for all $j<c$, and $P_cu-u$ lies in the sum of the images of the
  $E_{ci}$, $i<c$.
\item Let $\mathcal O\subset K$ be a discrete valuation ring with fraction field $K$ and residue
  characteristic $p$, and $M\subset V$ an $\mathcal O$-submodule with $\frac1{k!}E_{ab}^k(M)\subset M$
  for all $a\neq b$ and $k\geq0$. Then $p^{(c-1)d}P_c(M)\subset M$, where
  $d=\lfloor\log_p(L+n-1)\rfloor$, and $d=0$ if $L+n-1=0$.
\end{enumerate}
\end{lemma}
\leanref{Formalism/SharpRestriction/GrowthBounds/RootWeights/DegreeZeroWeightsBounded.lean}{Formalism.MainNotions.G.GLWeightTheory.extremalProjector}
\leanref{Formalism/SharpRestriction/GrowthBounds/RootWeights/CommutatorTripleProductLeibniz.lean}{Formalism.MainNotions.G.GLWeightTheory.extremal_projector_properties}
\leanref{Formalism/SharpRestriction/GrowthBounds/RootWeights/CommutatorTripleProductLeibniz.lean}{Formalism.MainNotions.G.GLWeightTheory.raising_casimirPower_commutator}
\leanref{Formalism/SharpRestriction/GrowthBounds/RootWeights/CasimirPowerScaledInsertDividedPower.lean}{Formalism.MainNotions.G.GLWeightTheory.casimirPowerScaled_insert_dividedPower}
\leanref{Formalism/SharpRestriction/GrowthBounds/RootWeights/CasimirPowerScaledInsertDividedPower.lean}{Formalism.MainNotions.G.GLWeightTheory.scaled_casimirPower_integral_bound}
\leanref{Formalism/SharpRestriction/GrowthBounds/RootWeights/CasimirPowerScaledInsertDividedPower.lean}{Formalism.MainNotions.G.GLWeightTheory.extremalProjector_denominator_bound}

\begin{proof}
(1) By the commutation relations the $E_{ci}$ with $i<c$ commute with each other, and so do the
$E_{ic}$. Hence $C_c^q(1)$ is the sum of the composites
$E_{ci_1}\cdots E_{ci_q}E_{i_qc}\cdots E_{i_1c}$ over $i_1,\dots,i_q<c$, and each composite maps
every $V_\mu$ into itself. For $a,b<c$ the commutation relations give
$E_{ab}C_c(T)-C_c(T)E_{ab}=C_c(E_{ab}T-TE_{ab})$, so $E_{ab}$ commutes with every $C_c^q(1)$. This
gives the first two assertions. Let $j<c$ and $y=E_{jc}u$. A computation with the commutation
relations, using $u\in V_0$ and $E_{ab}u=0$ for $a,b<c$, gives
\[
 E_{jc}C_c^q(1)u=C_c^q(1)y+q(c-1+q)\,C_c^{q-1}(1)y\qquad(q\geq1).
\]
Let $\alpha_q$ be the coefficient of $C_c^q(1)$ in $P_c$. Then $\alpha_{q+1}(q+1)(c+q)=-\alpha_q$,
so the sum telescopes and $E_{jc}P_cu=\alpha_LC_c^L(1)y$. The vector $y$ has weight
$\epsilon_j-\epsilon_c$, where $\epsilon_a$ is the $a$-th unit vector. Each composite in $C_c^L(1)$
first applies $L$ operators $E_{ic}$ to $y$, which gives a vector of a weight with entry sum $0$ and
$c$-th entry $-1-L$; this vector is $0$ by the choice of $L$. So $E_{jc}P_cu=0$. Finally
$P_cu-u=\sum_{q\geq1}\alpha_qC_c^q(1)u$, and each composite in $C_c^q(1)$ with $q\geq1$ begins with
some $E_{ci}$, $i<c$.

(2) Write $E^{(k)}_{ab}=\frac1{k!}E_{ab}^k$. For $S\subset\{1,\dots,c-1\}$ and integers $b,q\geq0$
let $C_S(T)=\sum_{i\in S}E_{ci}TE_{ic}$ and $Q_S(b,q)=\frac{b!}{(q+b)!\,q!}C_S^q(1)$, so that
$P_c=\sum_{q\leq L}(-1)^qQ_S(c-1,q)$ for $S=\{1,\dots,c-1\}$. Let $i<c$ with $i\notin S$. Then
$C_{S\cup\{i\}}=C_S+C_{\{i\}}$, the maps $C_S$ and $C_{\{i\}}$ commute because the $E_{ci'}$ commute
with each other and so do the $E_{i'c}$, and $C_{\{i\}}^k(T)=(k!)^2E^{(k)}_{ci}TE^{(k)}_{ic}$. The
binomial theorem gives
\[
 Q_{S\cup\{i\}}(b,q)=\sum_{k=0}^{q}\binom{q+b}{k}^{-1}E^{(k)}_{ci}\,Q_S(b,q-k)\,E^{(k)}_{ic}.
\]
By Legendre's formula, $v_p\binom Nk=\sum_{l\geq1}\big(\lfloor N/p^l\rfloor-\lfloor k/p^l\rfloor
-\lfloor(N-k)/p^l\rfloor\big)$ for $0\leq k\leq N$, a sum of terms $0$ or $1$ that vanish for
$p^l>N$; so $v_p\binom Nk\leq\lfloor\log_pN\rfloor$. Integers prime to $p$ are units in $\mathcal O$.
By induction on $|S|$ it follows that $p^{|S|d}Q_S(b,q)(M)\subset M$ whenever $q+b\leq L+n-1$: the
map $Q_\emptyset(b,q)$ is $1$ if $q=0$ and $0$ otherwise, and the induction step uses the display
with $N=q+b$. For $S=\{1,\dots,c-1\}$, $b=c-1$ and $q\leq L$ this gives (2).
\end{proof}

\begin{lemma}\label{pd-gain-weyl-invariant-denominator}
Let $K$, $n$, $V$, the $E_{ab}$ and $L$ be as in Lemma~\ref{pd-gain-weyl-column-projector}, let
$\mathcal O\subset K$ be a discrete valuation ring with fraction field $K$ and residue
characteristic $p$, and let $M\subset V$ be an $\mathcal O$-submodule with
$\frac1{k!}E_{ab}^k(M)\subset M$ for all $a\neq b$ and $k\geq0$ and $\binom{E_{aa}}k(M)\subset M$ for
all $a$ and $k\geq0$, where $\binom{E_{aa}}k=\frac1{k!}E_{aa}(E_{aa}-1)\cdots(E_{aa}-k+1)$. Let
$V^{\mathrm{inv}}$ be the space of $x\in V$ with $E_{ab}x=0$ for all $a,b$, and let $\operatorname{Inv}$
be the projection of $V$ onto $V^{\mathrm{inv}}$ whose kernel is the sum of the irreducible
subrepresentations of $V$ on which some $E_{ab}$ is nonzero. Then
$p^{N_0d}\operatorname{Inv}(M)\subset M$, where $N_0=n(n-1)/2$ and $d$ is as in
Lemma~\ref{pd-gain-weyl-column-projector}(2).
\end{lemma}
\leanref{Formalism/SharpRestriction/GrowthBounds/RootWeights/InvariantDenominatorLinftyBound.lean}{Formalism.MainNotions.G.GLWeightTheory.InvariantDenominatorLinftyBound}
\leanref{Formalism/SharpRestriction/GrowthBounds/RootWeights/DenominatorLemmaWithoutProjector.lean}{Formalism.MainNotions.G.GLWeightTheory.linf_lemma_without_projector}
\leanref{Formalism/SharpRestriction/GrowthBounds/RootWeights/InvariantProjectionDenominatorBound.lean}{Formalism.MainNotions.G.GLWeightTheory.invariant_projection_denominator_bound}
\leanref{Formalism/SharpRestriction/GrowthBounds/RootWeights/InvariantProjectionDenominatorBound.lean}{Formalism.MainNotions.G.GLWeightTheory.denominator_chain_bound}
\leanref{Formalism/SharpRestriction/GrowthBounds/RootWeights/MemWeightSpaceIffDiagonalEigen.lean}{Formalism.MainNotions.G.GLWeightTheory.weight_projections_integral_findim}

\begin{proof}
The $E_{ab}$ give a representation of $\mathfrak{gl}_n$ on $V$. It is semisimple by Weyl's theorem
\cite[\S6.3]{Hum72}, applied to $\mathfrak{sl}_n$, since the $E_{aa}$ act semisimply; so
$\operatorname{Inv}$ is defined and commutes with every $E_{ab}$. Hence
$\operatorname{Inv}E_{ab}=E_{ab}\operatorname{Inv}=0$, and $\operatorname{Inv}$ vanishes on $V_\mu$ for
$\mu\neq0$, since $V^{\mathrm{inv}}\subset V_0$.

Let $x\in M$. Choose an integer $N$ with $|\mu_a|\leq N$ whenever $V_\mu\neq0$, and put
$f(t)=\binom{N+t}N\binom{N-t}N$. Then $f(0)=1$ and $f(m)=0$ for the integers $0<|m|\leq N$, so the
projection $\pi_0$ of $V$ onto $V_0$ along the other weight spaces is $\prod_af(E_{aa})$. The
polynomial $f$ takes integer values at the integers, so it is an integral linear combination of the
$\binom tk$, and $\pi_0(M)\subset M$. Put $y_0=\pi_0x$; then $\operatorname{Inv}y_0=\operatorname{Inv}x$.

For $c=1,\dots,n$ put $y_c=P_cy_{c-1}$, with $P_c$ as in Lemma~\ref{pd-gain-weyl-column-projector}.
By induction on $c$: $y_c\in V_0$, $E_{ab}y_c=0$ for all $a,b\leq c$,
$\operatorname{Inv}y_c=\operatorname{Inv}x$ and $p^{(0+1+\cdots+(c-1))d}y_c\in M$. Indeed, part (1) of
Lemma~\ref{pd-gain-weyl-column-projector} gives $y_c\in V_0$, $E_{ab}y_c=0$ for $a,b<c$ and
$E_{jc}y_c=0$ for $j<c$. For $j<c$, the vector $y_c$ has weight $0$ for the copy of
$\mathfrak{sl}_2$ spanned by $E_{jc}$, $E_{cj}$ and $E_{jj}-E_{cc}$ and is killed by $E_{jc}$, so
$E_{cj}y_c=0$; and $E_{cc}y_c=0$ as $y_c\in V_0$. Since $y_c-y_{c-1}$ lies in the images of the
$E_{ci}$, $\operatorname{Inv}y_c=\operatorname{Inv}y_{c-1}$. Part (2) gives the last assertion. For
$c=n$, $y_n\in V^{\mathrm{inv}}$, so $\operatorname{Inv}x=y_n$ and
$p^{N_0d}\operatorname{Inv}x=p^{N_0d}y_n\in M$.
\end{proof}

\begin{lemma}\label{pd-gain-weyl-highest-lattices}
Let $K$, $n$, $V$, the $E_{ab}$ and $\mathcal O$ be as in
Lemma~\ref{pd-gain-weyl-invariant-denominator}. Let $\lambda\in\mathbf Z^n$ with
$\lambda_1\geq\cdots\geq\lambda_n$ and $v\in V_\lambda$ nonzero with $E_{ab}v=0$ for all $a<b$, and
suppose that $V$ is spanned by the vectors obtained from $v$ by repeatedly applying the $E_{ab}$.
Let $\mathcal U$ be the set of composites of the operators $\frac1{k!}E_{ab}^k$ ($a\neq b$) and
$\binom{E_{aa}}k$ ($k\geq0$), let $\Delta$ be the $\mathcal O$-span of the $uv$ with $u\in\mathcal U$,
and let $\nabla$ be the set of $x\in V$ such that for every $u\in\mathcal U$ the component of $ux$ in
$V_\lambda$ lies in $\mathcal Ov$. Then $p^{N_0d}\nabla\subset\Delta$, where $N_0=n(n-1)/2$ and
$d=\lfloor\log_p(\lambda_1-\lambda_n+n-1)\rfloor$, and $d=0$ if $\lambda_1-\lambda_n+n-1=0$.
\end{lemma}
\leanref{Formalism/SharpRestriction/GrowthBounds/RootWeights/HighestWeightLatticeOfFirstPart.lean}{Formalism.MainNotions.G.GLWeightTheory.LatticeNablaDeltaBound}
\leanref{Formalism/SharpRestriction/GrowthBounds/RootWeights/LinfAdjointLeEllWtHighest.lean}{Formalism.MainNotions.G.GLWeightTheory.restA_of_linf_integrality}
\leanref{Formalism/SharpRestriction/GrowthBounds/RootWeights/LinfAdjointLeEllWtHighest.lean}{Formalism.MainNotions.G.GLWeightTheory.linf_adjoint_le_ellWt_highest}
\leanref{Formalism/SharpRestriction/GrowthBounds/RootWeights/ExistsScalarMultipleMemDelta.lean}{Formalism.MainNotions.G.GLWeightTheory.nabla_finitely_generated}
\leanref{Formalism/SharpRestriction/GrowthBounds/RootWeights/ExistsScalarMultipleMemDelta.lean}{Formalism.MainNotions.G.GLWeightTheory.nabla_to_delta_maps_span_top}
\leanref{Formalism/SharpRestriction/GrowthBounds/RootWeights/DeltaIntegralEnvelopingStable.lean}{Formalism.MainNotions.G.GLWeightTheory.nabla_to_delta_maps_stable}
\leanref{Formalism/SharpRestriction/GrowthBounds/RootWeights/TraceProjectionIsInvProj.lean}{Formalism.MainNotions.G.GLWeightTheory.trace_projection_isInvProj}

\begin{proof}
As in the proof of Lemma~\ref{pd-gain-weyl-invariant-denominator}, $V$ is semisimple; it is
generated by the highest weight vector $v$, so it is irreducible, $V_\lambda=Kv$, and the entries of
its weights lie between $\lambda_n$ and $\lambda_1$ \cite[\S\S20--21]{Hum72}. Let $\xi$ be the
linear form on $V$ with $\xi(v)=1$ that vanishes on the weight spaces other than $V_\lambda$.

The sets $\Delta$ and $\nabla$ are $\mathcal O$-submodules of $V$ stable under $\mathcal U$, and
$\Delta$ spans $V$. The weight space of $\Delta$ for $\lambda$ is $\mathcal Ov$, by Kostant's theorem
on the integral form of the enveloping algebra \cite[\S27.1]{Hum72} (the $\binom{E_{aa}}k$ act on
weight vectors by integers). Each $uv$ with $u\in\mathcal U$ is a weight vector in $\Delta$, so
$v\in\nabla$; and $\Delta\subset\nabla$, because $\mathcal U$ is closed under composition. If
$\xi(ux)=0$ for all
$u\in\mathcal U$, the subrepresentation generated by $x$ meets $V_\lambda$ in $0$, so $x=0$ by
irreducibility. So finitely many of the forms $\xi\circ u$, $u\in\mathcal U$, form a basis of the
dual of $V$; they take values in $\mathcal O$ on $\nabla$, so $\nabla$ is finitely generated.

Let $H$ be the set of $\phi\in\operatorname{End}_K(V)$ with $\phi(\nabla)\subset\Delta$. The space
$\operatorname{End}_K(V)$ is a representation with $E_{ab}\cdot\phi=E_{ab}\phi-\phi E_{ab}$, and
\[
 \frac1{k!}(E_{ab}\cdot)^k\phi=\sum_{i+j=k}(-1)^j\tfrac1{i!}E_{ab}^i\,\phi\,\tfrac1{j!}E_{ab}^j,\qquad
 \Big(\binom{E_{aa}\cdot}k\phi\Big)(y)=\sum_{i+j=k}\binom{-\mu_a}j\binom{E_{aa}}i\phi(y)
 \quad(y\in V_\mu).
\]
The weight components of an element of $\nabla$ lie in $\nabla$, as in the proof of
Lemma~\ref{pd-gain-weyl-invariant-denominator}, so $H$ is stable under these operators. It is
finitely generated, as $\nabla$ and $\Delta\subset\nabla$ are. It spans $\operatorname{End}_K(V)$:
for $\phi\in\operatorname{End}_K(V)$ and generators $g_1,\dots,g_r$ of $\nabla$ there is a nonzero
$c\in\mathcal O$ with $c\phi(g_i)\in\Delta$ for all $i$, since $\Delta$ spans $V$. The weights of
$\operatorname{End}_K(V)$ are differences of weights of $V$, so their entries have absolute value at
most $\lambda_1-\lambda_n$. Lemma~\ref{pd-gain-weyl-invariant-denominator} therefore applies to
$\operatorname{End}_K(V)$, the module $H$ and $L=\lambda_1-\lambda_n$.

The map $\phi_0=\xi(\cdot)v$ lies in $H$: for $x\in\nabla$, $\xi(x)v$ is the component of $x$ in
$V_\lambda$, which lies in $\mathcal Ov\subset\Delta$. Since $V$ is irreducible and
$V_\lambda=Kv$, the endomorphisms commuting with all $E_{ab}$ are the scalars. The trace vanishes on
every $E_{ab}\cdot\phi$, and these span the sum of the nontrivial irreducible subrepresentations of
$\operatorname{End}_K(V)$, which therefore is the kernel of the trace. So the projection
$\operatorname{Inv}$ of $\operatorname{End}_K(V)$ is $\phi\mapsto(\operatorname{tr}\phi/\dim V)\cdot1$,
and $\operatorname{Inv}\phi_0=(\dim V)^{-1}\cdot1$. By
Lemma~\ref{pd-gain-weyl-invariant-denominator}, $p^{N_0d}(\dim V)^{-1}\cdot1\in H$; multiplying by
the integer $\dim V$ gives $p^{N_0d}\cdot1\in H$, that is, $p^{N_0d}\nabla\subset\Delta$.
\end{proof}

\begin{lemma}\label{pd-gain-weyl-comparison}
Let $p\geq5$ be a prime and $\mathcal R$ the ring of a Hida setting at $p$
(Definition~\ref{pd-local-model}(1)). Let $\lambda\in\mathbf Z^k$ be dominant with
$\nu=\lambda-\lambda_k$, and put
$w(\lambda)=\frac{k(k-1)}2\big\lfloor\log_p(\lambda_1-\lambda_k+k-1)\big\rfloor$. Then
$p^{w(\lambda)}\nabla_{\mathcal R}(\nu)\subset\Delta_{\mathcal R}(\nu)\subset\nabla_{\mathcal R}(\nu)$
(Definition~\ref{pd-weight-lattices}(1)).
\end{lemma}
\leanref{Formalism/LocalToGlobal/GroupLattices/FockModels/Sandwich.lean}{Formalism.L2.G.Model.sandwich}
\leanref{Formalism/LocalToGlobal/GroupLattices/FockModels/Sandwich.lean}{Formalism.L2.G.Model.sandwichO}
\leanref{Formalism/SharpRestriction/GrowthBounds/GLWeightTheory.lean}{Formalism.MainNotions.G.GLWeightTheory.nabla_delta_comparison_field}

\begin{proof}
Let $\mathcal K$ be the fraction field of $\mathcal R$. It contains $\iota(Q)$, so it has
characteristic $0$, and $p=\varpi_0^2$ lies in the maximal ideal of $\mathcal R$, so $\mathcal R$ has
residue characteristic $p$. Let $E_{ab}$ act on $\mathcal K[X]$ as the derivation
$\sum_cx_{ac}\,\partial/\partial x_{bc}$. Then $\frac1{k!}E_{ab}^kf$ is the coefficient of $y^k$ in
$f$ with the $b$-th row of $X$ replaced by its sum with $y$ times the $a$-th row, that is, the
operator $E^{(k)}_{ab}$ of Definition~\ref{pd-weight-lattices}(1) when $a\neq b$; for $a=b$ that
operator is $\binom{E_{aa}}k$, and $E_{aa}$ multiplies a polynomial homogeneous of degree $d_a$ in
the $a$-th row by $d_a$. The $E_{ab}$ satisfy the commutation relations of
Lemma~\ref{pd-gain-weyl-column-projector}.

These operators act on the rows of $X$, so they commute with the substitutions $X\mapsto
X(1+yE_{lj})$ and preserve the degree in each column; hence they preserve $\nabla_{\mathcal R}(\nu)$
and $\nabla_{\mathcal K}(\nu)$. The polynomial $h_\nu$ lies in $\nabla_{\mathcal R}(\nu)$: right
multiplication by $1+yE_{lj}$ with $l<j$ fixes the leading principal minors, and $\delta_s$ has
degree $1$ in each of the first $s$ columns, so $h_\nu$ has degree $\nu_a$ in the $a$-th column.
So $\Delta_{\mathcal R}(\nu)\subset\nabla_{\mathcal R}(\nu)$.

For the first inclusion take $V=\nabla_{\mathcal K}(\nu)$ and $v=h_\nu$. Adding a multiple of the
$a$-th row to the $b$-th with $a<b$ fixes the leading principal minors, so $E_{ab}h_\nu=0$ for
$a<b$; and $h_\nu$ has degree $\nu_a$ in the $a$-th row, so $h_\nu\in V_\nu$. The representation
$\nabla_{\mathcal K}(\nu)$ of $GL_k$ is irreducible with highest weight $\nu$ and highest weight
vector $h_\nu$ \cite[Lecture~15]{FH91}, so it is spanned by the vectors obtained from $h_\nu$ by
applying the $E_{ab}$ repeatedly, and $V_\nu=\mathcal Kh_\nu$. In the notation of
Lemma~\ref{pd-gain-weyl-highest-lattices}, with $n=k$ and $\lambda=\nu$, the module $\Delta$ is
$\Delta_{\mathcal R}(\nu)$. Let $f\in\nabla_{\mathcal R}(\nu)$ and $u\in\mathcal U$. Then
$uf\in\nabla_{\mathcal R}(\nu)$, and its component in $V_\nu$ is the sum of its monomials of degree
$\nu_a$ in the $a$-th row for every $a$. This component has coefficients in $\mathcal R$ and equals
$ch_\nu$ for some $c\in\mathcal K$. At the identity matrix every $\delta_s$ is $1$, so $c$ is the
value there of a polynomial with coefficients in $\mathcal R$, and $c\in\mathcal R$. So $f$ lies
in the module $\nabla$ of Lemma~\ref{pd-gain-weyl-highest-lattices}. That lemma gives
$p^{w}\nabla_{\mathcal R}(\nu)\subset\Delta_{\mathcal R}(\nu)$ with
$w=\frac{k(k-1)}2\lfloor\log_p(\nu_1-\nu_k+k-1)\rfloor=w(\lambda)$.
\end{proof}

\begin{lemma}\label{pd-gain-weyl-sandwich}
Let $D$ be a definite-side datum with ranks $k,m,t,n$, $f=2n$, size data with $a\geq4$ and
$5\leq k\leq a$, and local models $G(u,w)$ at the places admissible for each word $u$, with a tie
(Definition~\ref{pd-lattice-instance}), such that $\lambda(v,s)_1-\lambda(v,s)_k\leq3a$ for every
position of every $G(u,w)$. Let $C_e\geq1$ and $e_p=\lceil C_efk^2\log_pa\rceil$ for every prime $p$.
For a local model let $M^\nabla$ be the module defined as its integral module $M$
(Definition~\ref{pd-local-model}(3)) with $H^\nabla(L)$ in place of $H^\Delta(L)$. Then:
\begin{enumerate}
\item $e_p\leq(C_e+1)fk^2\log_pa$ for every prime $p\leq a$;
\item for every word $u$ and place $w$ admissible for $u$, with $p=p_w$, the modules $M$ and
  $M^\nabla$ of $G(u,w)$ satisfy $p^{e_p}M^\nabla\subset M\subset M^\nabla$.
\end{enumerate}
\end{lemma}
\leanref{Formalism/SharpRestriction/GrowthBounds/Joints.lean}{Formalism.MainNotions.G.Joints.weyl_sandwich_of_factors}
\leanref{Formalism/SharpRestriction/GrowthBounds/Joints.lean}{Formalism.MainNotions.G.Joints.weylSandwich_of_model}
\leanref{Formalism/SharpRestriction/GrowthBounds/Joints.lean}{Formalism.MainNotions.G.Joints.expBound_of}
\leanref{Formalism/SharpRestriction/GrowthBounds/SmallOrderPrimes.lean}{Formalism.MainNotions.AuxBlock8.adm9}
\leanref{Formalism/SharpRestriction/GrowthBounds/DefsInst.lean}{Formalism.MainNotions.G.EpDef}
\leanref{Formalism/SharpRestriction/GrowthBounds/TopChoiceBounds.lean}{Formalism.MainNotions.G.TopChoiceBounds.sandwich_of_exponent_clause}

\begin{proof}
For (1): $p\leq a$ gives $\log_pa\geq1$, and $f,k\geq1$, so $fk^2\log_pa\geq1$ and
$e_p\leq C_efk^2\log_pa+1\leq(C_e+1)fk^2\log_pa$.

For (2), let $\lambda=\lambda(v,s)$ be the weight at a position. Since $\lambda_1-\lambda_k\leq3a$,
$k\leq a$ and $a\geq4$, $\lambda_1-\lambda_k+k-1\leq a^2$, so
$\lfloor\log_p(\lambda_1-\lambda_k+k-1)\rfloor\leq2\log_pa$ and
$w(\lambda)\leq\frac{k^2}2\cdot2\log_pa=k^2\log_pa$. The positions of $G(u,w)$ correspond under the tie
to the $f$ embeddings in $\mathcal T$, so the sum $W$ of the $w(\lambda(v,s))$ over all positions is at
most $fk^2\log_pa\leq e_p$. For a lattice tuple $L$, each factor of a product spanning $H^\nabla(L)$
is some $f_{v,s}\in\nabla_{\lambda(v,s)}(M^\vee_{v,s})$; after a change of basis of $M_{v,s}$ this is
the case of the standard lattices, and Lemma~\ref{pd-gain-weyl-comparison} gives
$p^{w(\lambda(v,s))}f_{v,s}\in\Delta_{\lambda(v,s)}(M^\vee_{v,s})$ and
$\Delta_{\lambda(v,s)}(M^\vee_{v,s})\subset\nabla_{\lambda(v,s)}(M^\vee_{v,s})$. Taking products,
$p^WH^\nabla(L)\subset H^\Delta(L)\subset H^\nabla(L)$. The conditions defining $M$ and $M^\nabla$
other than membership in these lattices are linear and the same for both, so
$p^WM^\nabla\subset M\subset M^\nabla$, and $p^{e_p}M^\nabla=p^{e_p-W}p^WM^\nabla\subset M$.
\end{proof}

\begin{definition}\label{pd-gain-quantities}
Let $D$ be a definite-side datum with ranks $k,m,t,n$, size data and gain data with constants $K_G$
(Definitions~\ref{pd-size-data} and~\ref{pd-gain-data}), with the notation there and in
Definition~\ref{pd-orbit-places}; let $\mathbf a,\mathbf b$ be the exponents extended by $0$, and
$e_p$ the Weyl exponents of the gain data (Definition~\ref{pd-gain-data}(5)). Put $c_G=4S_c$, let $D_{\mathrm{II}},t_{\mathrm{II}}$
and $D_{\mathrm I},t_{\mathrm I}$ be as in Lemmas~\ref{pd-gain-second-hypotheses}
and~\ref{pd-gain-first-hypotheses}, with $L_1(p)$ and $\Lambda=K_\Lambda N_1^2\log a$ as there, and
for a prime $p$, with orbit length $\ell_p$, and a representative $t_0$ at $p$ put
$D(p)=D_{\mathrm{II}}(p,\ell_p)$ and $t(p,t_0)=t_{\mathrm{II}}(p,\ell_p,N(p,t_0))$ if $p<p_0$, and
$D(p)=D_{\mathrm I}(p,\ell_p)$ and
$t(p,t_0)=t_{\mathrm I}\big(p,n_p,P(p,t_0),(\mathbf a(t_0+ih_p))_i,(\mathbf b(t_0+ih_p))_i\big)$
otherwise. The gain data determine the following \emph{gain quantities}.
\begin{enumerate}
\item For a lift $j$ and a place $z$, $\mathrm{fill}_j(z)$ is as in
  Definition~\ref{pd-lattice-instance}. For $s\in T$, the \emph{class sum with fills} is
  \[
   \mathrm{VM}^G_{j,s}(p)=\sum_{e\geq1\ \mathrm{odd}}\max\Big(0,\frac{a_s}{p^e}-\mathrm{fill}_j
   \big(\operatorname{pos}_p((0,s),e)\big)\Big),
  \]
  and $\mathrm{VM}^0_{j,s}(p)$ is the same sum with $\mathrm{fill}_j+\kappa_j$ in place of
  $\mathrm{fill}_j$.
\item The \emph{counted classes} at $p$ are $T(p)=T$ if $p$ is a signal prime or $p\geq X'$, and
  otherwise the $t\in T$ with $t^{s_p}\in V(p)$.
\item The \emph{depth} at $p$ is $\tau(p)=0$ if $V(p)=\emptyset$ and otherwise
  $\tau(p)=e_p+m\sum_{t_0\in V(p)}\big(\lceil t(p,t_0)\rceil+D(p)+1\big)$; the \emph{index excess} is
  $\iota(p)=2\Sigma(p)\max_{t_0\in V(p)}S(p,t_0)$ if $V(p)\neq\emptyset$ and $0$ otherwise.
\item For a lift $j$, the \emph{gain} $g_j(p)$ is $0$ if $p>X''$ or if $p$ is of inert class with
  $h_p=0$; otherwise it is
  \[
   g_j(p)=\sum_{t_0\in V(p)}\max\big(g_U(p,j,t_0)-\mathrm{Pr}(p,j,t_0),g_0(p,j,t_0),0\big)
   +\sum_{t_0\notin V(p)}\max\big(g_0(p,j,t_0),0\big),
  \]
  the second sum over the representatives at $p$ not in $V(p)$, where $g_0$ is the local gain of
  Definition~\ref{pd-served-selection} and $g_U(p,j,t_0)$ and $\mathrm{Pr}(p,j,t_0)$, the \emph{upper
  gain} and its \emph{correction}, are numbers computed from the offsets, the exponents and $K_G$.
  The surplus is $S(p,t_0)=\big\lceil\sum_{j=1}^m\max\big(0,g_U-\mathrm{Pr}-\max(g_0,0)\big)
  \big\rceil$, with the arguments $(p,j,t_0)$.
\item The \emph{lower term} is $\mathrm{lo}^G(p)=C(c_G+\eta)\log_p(2a)+\mathrm{pr}^G(p)+
  \mathrm{or}^G(p)$, where for primes of inert class $\mathrm{pr}^G(p)$ is $C\,r(\varepsilon)b'\log
  x$ if $p<p_0$, $b'/x+\eta+Cw''\log x/\delta_1+\beta$ if $p_0\leq p<X'$, and $0$ if $p\geq X'$,
  and $\mathrm{or}^G(p)=C\log x\,(\Lambda a/(pk)+ap^{-L_1(p)})$ if $p<p_0$,
  $C\log x\,(\Lambda a/(pk)+a/p^2)$ if $p_0\leq p<X'$, and $0$ if $p\geq X'$; both are $0$ for
  signal primes. Here $C,\beta$ are real numbers.
\item The size data are \emph{filled} by the gain data at a place $w$, with $p=p_w$, if
  $\mathrm{fill}_j(z)+\kappa_j=w_j(z)$ for every lift $j$ and place $z$, the weights of the size data
  are $w_{j,t,p,e}=\mathrm{fill}_j(\operatorname{pos}_p((0,t),e))+\kappa_j$, $T_w=T(p)$,
  $\mathrm{gain}_j(w)=g_j(p)$ for every lift $j$ when $w$ is in the support and $p$ is inert,
  $\tau_w=\tau(p)$, $\iota_w=\iota(p)$, and $\Lambda_G=\Lambda$.
\item Let $D$ carry local models and a lattice instance (Definition~\ref{pd-lattice-instance}), with
  maps $t\mapsto v(t)$ and submodules $V_{\mathrm{ram}}$. Let $w$ be admissible for a word $u$,
  $p=p_w$, and $M$ the integral module of $G(u,w)$. For $V'\subset V(p)$ let $y_{V'}$ be the
  composite of the $m$-th powers of the Hecke operators $y_{v(t_0)}$, $t_0\in V'$
  (Definition~\ref{pd-local-model}(4)), and $s(V')=\sum_{t_0\in V(p)\smallsetminus V'}\lceil
  S(p,t_0)\rceil$. The \emph{served local lattice} $N_P(u,w)$ is the $\mathcal R$-span of the
  $y_{V'}(p^{s(V')}F)$, for $V'\subset V(p)$ and $F\in M$, intersected with $p^{\lceil\tau(p)\rceil}M$
  and with $V_{\mathrm{ram}}$.
\end{enumerate}
\end{definition}
\leanref{Formalism/SharpRestriction/GrowthBounds/IsResidueClassZero.lean}{Formalism.MainNotions.G.askG}
\leanref{Formalism/SharpRestriction/GrowthBounds/IsResidueClassZero.lean}{Formalism.MainNotions.G.askP}
\leanref{Formalism/SharpRestriction/GrowthBounds/IsResidueClassZero.lean}{Formalism.MainNotions.G.TTP}
\leanref{Formalism/SharpRestriction/GrowthBounds/IsResidueClassZero.lean}{Formalism.MainNotions.G.tauP}
\leanref{Formalism/SharpRestriction/GrowthBounds/IsResidueClassZero.lean}{Formalism.MainNotions.G.ioP}
\leanref{Formalism/SharpRestriction/GrowthBounds/IsResidueClassZero.lean}{Formalism.MainNotions.G.gainP}
\leanref{Formalism/SharpRestriction/GrowthBounds/IsResidueClassZero.lean}{Formalism.MainNotions.G.Sv}
\leanref{Formalism/SharpRestriction/GrowthBounds/IsResidueClassZero.lean}{Formalism.MainNotions.G.lo}
\leanref{Formalism/SharpRestriction/GrowthBounds/IsResidueClassZero.lean}{Formalism.MainNotions.G.ParametersFill}
\leanref{Formalism/SharpRestriction/GrowthBounds/IsResidueClassZero.lean}{Formalism.MainNotions.G.yProd}
\leanref{Formalism/SharpRestriction/GrowthBounds/IsResidueClassZero.lean}{Formalism.MainNotions.G.NP}

\begin{lemma}\label{pd-gain-coherence-consequences}
Let $D$ be a definite-side datum with coefficient field $L_D$, local models $G(u,w)$, a tie of
them, a lattice instance with a comparison, and group coherence
(Definition~\ref{pd-gain-coherence}). Then the following hold at every word $u$ and every place $w$
admissible for $u$.
\begin{enumerate}
\item \emph{Trivial stabilizers and finite classes.} For every component $c$ of $G(u,w)$, the only
  element of $\Gamma_c$ that fixes a lattice tuple is $1$, and the lattice tuples fall into finitely
  many orbits under $\Gamma_c$.
\item \emph{Standard rings.} For every orbit $v$ of $G(u,w)$ there are an unramified quadratic Hida
  setting $(W_v,Q_v,\sigma_v,\iota'_v)$ of degree $n'_v$ at $p_w$ with ring $\mathcal R'$
  (Definition~\ref{pd-projection-property}) and a finite homomorphism $\mathcal R'\to\mathcal R$,
  extending to the fraction fields, that carries $\iota'_v$, the inclusion of $W_v$ and the element
  $\varpi'_0$ of the setting to $\iota_v$, the inclusion of $W_v$ and $\varpi_0$.
\item \emph{Rationality.} For every $x\in V_{u,0}$, component $c$ and lattice tuple $L$, the
  polynomial $\operatorname{cmp}_{u,w}(x)(c,L)$ is the image under $\epsilon_w$ of a polynomial with
  coefficients in $L_D$.
\item \emph{Descent.} Suppose that the submodule $V_{\mathrm{ram}}$ of the lattice instance at
  $(u,w)$ is the $\mathcal K$-span of $\operatorname{cmp}_{u,w}(M_{u,\mathrm{ram}})$. Then for every
  orbit $v$ there is an $L_D$-linear map $y'_v$ of $V_{u,0}$ with
  $y_v\circ\operatorname{cmp}_{u,w}=\operatorname{cmp}_{u,w}\circ y'_v$; and, for any gain data with
  $V(p_w)\neq\emptyset$, the served local lattice $N_P(u,w)$
  (Definition~\ref{pd-gain-quantities}(7)) is contained in the $\mathcal R$-span of the
  $\operatorname{cmp}_{u,w}(y)$ with $y\in M_{u,\mathrm{ram}}$ and
  $\operatorname{cmp}_{u,w}(y)\in N_P(u,w)$.
\end{enumerate}
\end{lemma}
\leanref{Formalism/LocalToGlobal/GroupLattices/StdSubAtOfRing.lean}{Formalism.L2.G.gatE_of_gatIs}
\leanref{Formalism/LocalToGlobal/GroupLattices/StdSubAtOfRing.lean}{Formalism.L2.G.stdSubAt_of_ring}
\leanref{Formalism/LocalToGlobal/GroupLattices/ComparisonMaps/NpDescFOfValRatF.lean}{Formalism.L2.G.Cmp.CmpTied.valRat}
\leanref{Formalism/SharpRestriction/GrowthBounds/PlaneCountTheoremWithConstants.lean}{Formalism.MainNotions.G.TrivialStabiliserFiniteClasses}
\leanref{Formalism/SharpRestriction/GrowthBounds/Joints.lean}{Formalism.MainNotions.G.Joints.StdSubAt}
\leanref{Formalism/LocalToGlobal/GroupLattices/TrivialTransfer/GenCoprime.lean}{Formalism.L2.G.Triv.gen_triv}
\leanref{Formalism/LocalToGlobal/GroupLattices/ComparisonMaps/YvRatOfValRat.lean}{Formalism.L2.G.Cmp.yvRat_of_valRat}
\leanref{Formalism/LocalToGlobal/GroupLattices/ComparisonMaps/NpDescOfValRat.lean}{Formalism.L2.G.Cmp.npDesc_of_valRat}

\begin{proof}
For (1), let $\gamma\in\Gamma_c$ fix a lattice tuple $L$. By (1) and (5) of
Definition~\ref{pd-gain-coherence}, $\gamma$ is an element of $U_J(F)$ that lies in
$y_cK^{(p)}_uy_c^{-1}$ and stabilizes, at each place $w'$ over $p$, the lattice that corresponds to
$L$; choosing $x_{w'}$ that carries the lattice of $K_{u,w'}$ to it, $\gamma$ lies in $xK_ux^{-1}$ for
some $x\in U_J(\mathbf A_f)$. The Hermitian form $J$ is totally definite, so $U_J(F)\cap xK_ux^{-1}$
is finite, and $\gamma$ has finite order $N$. Write $N=g_0^aN'$ with $g_0\nmid N'$. The element
$\gamma^{g_0^a}$ has order prime to $g_0$ and lies, after conjugation, in the group of integral
matrices unipotent upper triangular modulo a place over $g_0$, which is a pro-$g_0$ group; so
$\gamma^{g_0^a}=1$. The same argument at $g_1$ shows that the order of $\gamma$ is a power of $g_1$.
Since $g_0\neq g_1$, $\gamma=1$. The orbits of the lattice tuples under $\Gamma_c$ correspond to
classes of $U_J(F)\backslash U_J(\mathbf A_f)/K_u$, which correspond to the finite set
$\mathrm{cls}_0$ by Definition~\ref{pd-gain-coherence}(1).

For (2), transport the standard unramified quadratic setting on $W(\mathbf F_{p^{n'}})$ and
$\mathcal R_0=W(\mathbf F_{p^{2n'}})[\omega]$ along the isomorphisms of
Definition~\ref{pd-gain-coherence}(4): $\mathcal R_0$ has uniformizer $\omega$ with $\omega^2=p$,
residue field with $p^{2n'}$ elements, and every integer is a square in it, and the finite extension
of (4) gives the map to $\mathcal R$.

For (3), Definition~\ref{pd-gain-coherence}(3) gives
$\operatorname{cmp}_{u,w}(x)(c,L)=\epsilon_w(\operatorname{val}(x,c,L))$.

For (4), the map $\operatorname{cmp}_{u,w}$ is injective and $\epsilon_w$-semilinear and carries a
basis of $\Lambda_{u,0}$ to linearly independent vectors, so a function in the $\mathcal K$-span of
$\operatorname{cmp}_{u,w}(V_{u,0})$ lies in $\operatorname{cmp}_{u,w}(V_{u,0})$ if and only if its
values are images of polynomials over $L_D$: its coordinates in that basis solve a linear system whose
coefficients are such values. The Hecke operator $y_v$ preserves the $\mathcal K$-span and, by its
formula, carries functions with values over $L_D$ to such functions. By (3) it therefore carries
$\operatorname{cmp}_{u,w}(V_{u,0})$ into itself, and
$y'_v=\operatorname{cmp}_{u,w}^{-1}\circ y_v\circ\operatorname{cmp}_{u,w}$ is $L_D$-linear. Since
$M$ is the $\mathcal R$-span of $\operatorname{cmp}_{u,w}(\Lambda_{u,0})$, the span of the
$y_{V'}(p^{s(V')}F)$ and the module $p^{\lceil\tau(p)\rceil}M$ are $\mathcal R$-spans of images of
vectors of $V_{u,0}$, and $V_{\mathrm{ram}}$ is the $\mathcal K$-span of the image of
$M_{u,\mathrm{ram}}$. Each of the three is the base change to $\mathcal R$ of the corresponding module
over the valuation ring of $L_D$ at $w$, along the flat map $\epsilon_w$, and base change commutes
with finite intersections; so $N_P(u,w)$ is spanned by images of vectors of $M_{u,\mathrm{ram}}$
whose images it contains.
\end{proof}

\begin{lemma}\label{pd-gain-base-change}
Let $p\geq5$ be a prime, $(W,Q,\sigma,\iota)$ a Hida setting of degree $n'$ at $p$ with rings
$\mathcal R\subset\mathcal K$ and element $\varpi_0$, and $\phi\colon\mathcal R\to\mathcal R'$ a finite
homomorphism to another complete discrete valuation ring, extending to
$\phi\colon\mathcal K\to\mathcal K'$, such that $(W,Q,\sigma,\phi\circ\iota)$ with $\mathcal R'$ and
$\phi(\varpi_0)$ is again a Hida setting. If $\mathrm{LP}(r,\lambda_0;t,D,\kappa,c_*)$ holds for the
first setting, it holds for the second (Definition~\ref{pd-projection-property}).
\end{lemma}
\leanref{Formalism/LocalToGlobal/GroupLattices/FockModels/LPBase.lean}{Formalism.L2.G.Model.lpBaseChange}
\leanref{Formalism/SharpRestriction/GrowthBounds/UnramifiedQuadraticSetting.lean}{Formalism.MainNotions.G.LPBaseChange}

\begin{proof}
If $\lambda_0$ is not dominant, $V=0$ and the property holds trivially. Otherwise the $W$-lattices,
the edges and the exponents $A$ are the same for both settings. The lattices $H(L)$, $H(e)$, $R(e)$,
$R'(e)$, $N_{\mathrm{src}}$ and $N_{\mathrm{tgt}}$ of the second setting are the base changes along
$\phi$ of those of the first: they are built from the spans of the $\iota\sigma^s(L)$ by Weyl-module
constructions, duals and finite intersections, all of which commute with the flat base change
$\mathcal R\to\mathcal R'$. The edge pairing is compatible with $\phi$, and $\phi(\varpi_0)$ replaces
$\varpi_0$. So the images under $\phi$ of the vectors $F_i$, the base change of $u$ and the matrix
$(\phi(B_{ji}))$ witness the projection property for the second setting, with the same $n$.
\end{proof}

\begin{definition}\label{pd-gain-admissibility}
Let $D$ be a definite-side datum with ranks $k,m,t,n$, $f=2n$, size data and gain data with Weyl
exponents $e_p$ and constants $K_G$, with the gain quantities of
Definition~\ref{pd-gain-quantities}, and with local models and a lattice instance
(Definition~\ref{pd-lattice-instance}). Use the notation of Definitions~\ref{pd-gain-data}
and~\ref{hf-good-orbit} and Lemmas~\ref{pd-gain-second-hypotheses} and~\ref{pd-gain-first-hypotheses},
with $r=r_\varepsilon$; let $\rho$ be as in Definition~\ref{pd-gain-data}(c) and $\rho_0$ as in
Lemma~\ref{pd-gain-first-hypotheses}; and call $t\in A$ \emph{even at} $p$ if $t=t^{s_p}+ih_p$ with
$i$ even and $0\leq i<\ell_p$. Let $C_{31}>0$, $C_e,C_m\geq0$ and $\delta>0$. The gain data are
\emph{admissible} with constants $(C_{31},C_e,C_m,\delta)$ if the following hold.
\begin{enumerate}
\item \emph{Orbits.} For every prime $p$ of inert class with $n_p\geq x^{\varepsilon/4}$, if
  $L_2<u$ and $3G'\leq n_p$, then for every good representative $t_0$ at $p$:
  $n_p\delta_1/(C_{31}\log x)\leq N(p,t_0)\leq40n_p\delta_1$; for every top class $s$ with
  representative $t_0$ that is even at $p$ there is $1\leq d\leq G'$ such that $s+dh_p\in T$ is even
  at $p$ and at most $N_1$ of the $s+eh_p$ with $0<e<d$ lie in $T$; and for every top class $s$ with
  representative $t_0$, fewer than $r$ of the integers $1\leq d\leq L_1(p)$ have $s-dh_p\in T$ even at
  $p$.
\item \emph{Tops.} $n\delta_1/2\leq|T|\leq2n\delta_1$ and $B\leq|T|x^{-2+\varepsilon}$, where $B$
  is the sum, over the primes $p<X'$ of inert class with $n_p\geq x^{\varepsilon/4}$, of $\log p/p$
  times the number of top classes $t$ whose representative $t^{s_p}$ is not good at $p$.
\item \emph{Small orders.} The sum of $\log p/p$ over the primes $p<X''$ of inert class with
  $n_p<x^{\varepsilon/4}$ is at most $x^{-h_\varepsilon+\varepsilon}$.
\item \emph{Bands.} $\rho_0\leq\rho\leq x^{-2}$, and the sums of $\log p/p$ over the primes in the
  narrow band and over the primes in band are each at most $x^{-1}$.
\item \emph{Sizes of $\Lambda$ and $Q$.} $x^{\varepsilon-\delta}\leq\Lambda\leq
  x^{\varepsilon+\delta}$, and for every prime $p$, with $Q=p^{G'}$: $\log Q\leq
  x^{\varepsilon/8+\delta}$ if $p<p_0$, and $\log Q\leq x^{11\varepsilon/24+\delta}$ if $p<X'$.
\item \emph{Second family of bounds.} The conclusion of Lemma~\ref{pd-gain-second-hypotheses} holds
  for all $p$, settings, $T$, $a_w,b_w$ and $\lambda_0$ that satisfy its hypotheses.
\item \emph{Local constants.} For every prime $p$ with $p_0\leq p<X'$ not in the narrow band, every
  unramified quadratic Hida setting of degree $n'$ at $p$, every tiled set of positions
  $T\subset\mathbf Z/n'\mathbf Z$ for $(x,\varepsilon,L_2)$ with $L_2<u$, $3G'\leq n'$ and
  Hypothesis~I at $(p,n',|T|)$, all $a',b'_w$ with values in $W_a$ and $W_b$ on $T$, and all dominant
  $\lambda_0$ with $\lambda_0(s)_1=a'(s)$, $\lambda_0(s)_k=-b'_w(s)$ and
  $|\lambda_0(s)_i|\leq\lfloor c_G\rfloor$ for $1<i<k$ when $s\in T$ and for all $i$ when
  $s\notin T$, the projection property
  $\mathrm{LP}\big(1,\lambda_0;\lceil t_{\mathrm I}(p,n',T,a',b'_w)\rceil,D_{\mathrm I}(p,n'),
  \kappa_I(p,n',T,a',b'_w),8(D_{\mathrm I}(p,n')+1)\big)$ holds.
\item \emph{First family of bounds.} The conclusion of Lemma~\ref{pd-gain-first-hypotheses} holds
  with $C_m$, for all $p$, $n'$ and $T$ that satisfy its hypotheses.
\item \emph{Weyl sandwich.} $e_p\leq(C_e+1)fk^2\log_pa$ for every prime $p\leq a$, and, if the local
  models have a tie, $p_w^{e_{p_w}}M^\nabla\subset M\subset M^\nabla$ for every word $u$ and place $w$
  admissible for $u$, in the notation of Lemma~\ref{pd-gain-weyl-sandwich}.
\end{enumerate}
For the next two conditions, let a tie be given under which the lattice instance is tied to the gain
data, let $w$ be admissible for a word $u$, $p=p_w$, $M$ the integral module of $G(u,w)$,
$M_{\mathrm{ram}}=M\cap V_{\mathrm{ram}}$, and $K^{(m)}$ the composite of the first $m$ maps of the
chain of the lattice instance at $(u,w)$; write $g_U-\mathrm{Pr}$ and $g_0$ for their values at
$(p,j,t_0)$.
\begin{enumerate}
\setcounter{enumi}{9}
\item \emph{Units at the twisted places.} For every set $V''$ of representatives at $p$ and every
  $F\in M$, $K^{(m)}(y_{V''}F)$ has valuation at least
  \[
   \sum_{j=1}^m\Big(\sum_{t_0\in V''}\max(g_U-\mathrm{Pr},g_0,0)+\sum_{t_0\notin V''}\max(g_0,0)
   \Big)-m\,v_p(D_\Gamma(u))
  \]
  (Definition~\ref{pd-local-model}(7)), where the second sum is over the representatives at $p$
  not in $V''$.
\item \emph{Served lattices.} $N_P(u,w)\subset p^{\lceil\tau(p)\rceil}M_{\mathrm{ram}}$; the length
  of $M_{\mathrm{ram}}/N_P(u,w)$ as an $\mathcal R$-module is at most
  $v_{\mathcal R}(p)\dim M_u\,(\tau(p)+\iota(p))$; and for every $F\in N_P(u,w)$, $K^{(m)}F$ has
  valuation at least $\sum_{j=1}^m\big(g_j(p)+F_j(w)-E_{u,j}(w)\big)$.
\item \emph{Size bounds.} The size data satisfy the size bounds chosen with the family
  (Definition~\ref{pd-family}(4)), $4S_c\leq\Lambda b'/k$, and $N(m)+1\leq w''$.
\end{enumerate}
A family whose members carry gain data with Weyl exponents, one list $K_G$, local models and lattice
instances is \emph{admissible} if there are $C_{31}>0$ and $C_e,C_m\geq0$ such that for every
$\delta>0$ there is $x_0>1$ with the gain data of every member at every scale $x\geq x_0$ admissible
with $(C_{31},C_e,C_m,\delta)$.
\end{definition}
\leanref{Formalism/SharpRestriction/GrowthBounds/SmallOrderPrimes.lean}{Formalism.MainNotions.AuxBlock8.admitted}
\leanref{Formalism/SharpRestriction/GrowthBounds/SmallOrderPrimes.lean}{Formalism.MainNotions.AuxBlock8.AdmF}
\leanref{Formalism/SharpRestriction/GrowthBounds/SmallOrderPrimes.lean}{Formalism.MainNotions.AuxBlock8.adm1}
\leanref{Formalism/SharpRestriction/GrowthBounds/SmallOrderPrimes.lean}{Formalism.MainNotions.AuxBlock8.adm2}
\leanref{Formalism/SharpRestriction/GrowthBounds/SmallOrderPrimes.lean}{Formalism.MainNotions.AuxBlock8.adm3}
\leanref{Formalism/SharpRestriction/GrowthBounds/SmallOrderPrimes.lean}{Formalism.MainNotions.AuxBlock8.ArchimedeanSignatureMatching}
\leanref{Formalism/SharpRestriction/GrowthBounds/SmallOrderPrimes.lean}{Formalism.MainNotions.AuxBlock8.adm5}
\leanref{Formalism/SharpRestriction/GrowthBounds/SmallOrderPrimes.lean}{Formalism.MainNotions.AuxBlock8.adm7}
\leanref{Formalism/SharpRestriction/GrowthBounds/SmallOrderPrimes.lean}{Formalism.MainNotions.AuxBlock8.unitBoundAtTwistedPlaces}
\leanref{Formalism/SharpRestriction/GrowthBounds/SmallOrderPrimes.lean}{Formalism.MainNotions.AuxBlock8.admG1}
\leanref{Formalism/SharpRestriction/GrowthBounds/SmallOrderPrimes.lean}{Formalism.MainNotions.AuxBlock8.TypePrimeSizeConds}

\begin{lemma}\label{pd-gain-orbit-bounds}
Let $D$ be a definite-side datum with size data, gain data and the notation of
Definition~\ref{pd-gain-admissibility}, with local models, a lattice instance and group coherence
(Definition~\ref{pd-gain-coherence}). Suppose that conditions (6) and (7) of
Definition~\ref{pd-gain-admissibility} hold, that $\rho_0\leq\rho$, and that at every candidate
prime $p$ and good representative $t_0$ at $p$, with $P=P(p,t_0)$ and $a_w,b_w$ the functions
$i\mapsto\mathbf a(t_0+ih_p)$ and $i\mapsto\mathbf b(t_0+ih_p)$:
\begin{enumerate}
\item[(I)] if $p\geq p_0$, then $L_2<u$, $3G'\leq n_p$, Hypothesis~I holds at $(p,n_p,|P|)$, and
  $a_w,b_w$ take values in $W_a$ and $W_b$ on $P$;
\item[(II)] if $p<p_0$, then $T=P$, $n'=n_p$, $a_w$ and $b_w$ satisfy the hypotheses of
  Lemma~\ref{pd-gain-second-hypotheses} other than those on the setting and on $\lambda_0$.
\end{enumerate}
Then for every tie under which the lattice instance is tied to the gain data, every word $u$, place
$w$ admissible for $u$ with $p=p_w$, and $t_0\in V(p)$, the projection property
$\mathrm{LP}\big(r_w,\lambda_0(v,\cdot);\lceil t(p,t_0)\rceil,D(p),\kappa_{p,t_0}(\mathbf a,\mathbf b),
8(D(p)+1)\big)$ holds at the orbit $v=v(t_0)$ of $G(u,w)$, where $r_w$ is the integer $r$ of the
local model.
\end{lemma}
\leanref{Formalism/SharpRestriction/GrowthBounds/Joints.lean}{Formalism.MainNotions.G.Joints.linear_programming_input_of_hypotheses}
\leanref{Formalism/SharpRestriction/GrowthBounds/CornerGLContinuity.lean}{Formalism.MainNotions.CornerGLContinuity.OrbitLatticeBounds}
\leanref{Formalism/SharpRestriction/GrowthBounds/Joints.lean}{Formalism.MainNotions.G.Joints.HypIAt}
\leanref{Formalism/SharpRestriction/GrowthBounds/Joints.lean}{Formalism.MainNotions.G.Joints.HypIIAt}

\begin{proof}
Let $t_0\in V(p)$. Then $p$ is a candidate prime and $t_0$ is good at $p$. Since the lattice
instance is tied to the gain data (Definition~\ref{pd-lattice-instance}(4)), the orbit $v=v(t_0)$
has $n'_v=n_p$, its top positions are $P=P(p,t_0)$, so $|P|=N(p,t_0)$, its base weights at the top
positions are $a_w$ and $b_w$, its constant $c$ is at most $c_G$, hence at most $\lfloor c_G\rfloor$,
and $r_w=r_\varepsilon$ if $p<p_0$ and $r_w=1$ otherwise. By
Lemma~\ref{pd-gain-coherence-consequences}(2) the Hida setting of $v$ comes from an unramified
quadratic setting by a finite homomorphism $\mathcal R'\to\mathcal R$. If $p<p_0$, the hypotheses
of Lemma~\ref{pd-gain-second-hypotheses} hold for this setting by (II), and condition (6) gives the
projection property over $\mathcal R'$ with $\kappa_{\mathrm{II}}(N(p,t_0))=\kappa_{p,t_0}(\mathbf
a,\mathbf b)$. If $p\geq p_0$, $p$ is not in band, and since $\rho_0\leq\rho$ it is not in the
narrow band; with (I), condition (7) gives the projection property over $\mathcal R'$ with
$\kappa_I=\kappa_{p,t_0}(\mathbf a,\mathbf b)$. In both cases $t(p,t_0)$ and $D(p)$ are the
thresholds used there. Lemma~\ref{pd-gain-base-change} carries the property to $\mathcal R$.
\end{proof}

\begin{lemma}\label{pd-gain-holds-series}
Let $p\geq2$ and $y$ be real numbers, $(w_e)_{e\geq0}$ and $(w'_e)_{e\geq0}$ sequences of real
numbers and $Q$ a set of odd positive integers. For a real number $c$ put
$\Sigma_w(c)=\sum_{e\in Q}\max(0,y/p^e-w_e-c)$.
\begin{enumerate}
\item If $y\geq0$ and $w_e\geq0$ for every $e$, then $\Sigma_w(0)\leq py/(p^2-1)$.
\item If $c\geq0$, $2c\leq y\leq A$ for a real number $A\geq1$, $w_e\geq1$ and $0\leq w'_e\leq w_e+c$
  for every $e$, then
  $\Sigma_w(0)-3c\log_pA\leq\sum_{e\in Q}\max\big(0,(y-2c)/p^e-w'_e\big)$.
\item If $y\geq0$, $\kappa\geq0$ and $w_e\geq0$ for every $e$, then $\Sigma_w(\kappa)\leq\Sigma_w(0)$.
\item If $0\leq y\leq A$ for a real number $A\geq1$, $\kappa\geq0$ and $w_e\geq1$ for every $e$, then
  $\Sigma_w(0)\leq\Sigma_w(\kappa)+\kappa\log_pA$.
\end{enumerate}
\end{lemma}
\leanref{Formalism/SharpRestriction/GrowthBounds/LogarithmicSums/TsumTruncatedShiftLe.lean}{Formalism.MainNotions.AuxBlock8.Pf.tsum_pos_le}
\leanref{Formalism/SharpRestriction/GrowthBounds/LogarithmicSums/TsumTruncatedShiftLe.lean}{Formalism.MainNotions.AuxBlock8.Pf.tsum_truncated_shift_le}
\leanref{Formalism/SharpRestriction/GrowthBounds/LogarithmicSums/GainBoundAllAddOneOfGainBoundAllRequired.lean}{Formalism.MainNotions.AuxBlock8.Pf.As.shift_ge}
\leanref{Formalism/SharpRestriction/GrowthBounds/LogarithmicSums/GainBoundAllAddOneOfGainBoundAllRequired.lean}{Formalism.MainNotions.AuxBlock8.Pf.As.shift_le}

\begin{proof}
In each part the terms of the series are at most $|y|p^{-e}$, so they converge. (1) Each term is at
most $y/p^e$, and $\sum_{e\ \mathrm{odd}}y/p^e=py/(p^2-1)$. (3) Term by term. (2) If
$y/p^e-w_e>0$, then $y/p^e>1$, so $p^e<y\leq A$ and $1\leq e\leq\lfloor\log_pA\rfloor$; for such
$e$ the difference $(y/p^e-w_e)-((y-2c)/p^e-w'_e)=2c/p^e+w'_e-w_e$ is at most $3c$. At the other $e$
the term on the left is $0$, and at most $\log_pA$ integers lie in $[1,\lfloor\log_pA\rfloor]$.
(4) As in (2), a nonzero term of $\Sigma_w(0)$ has $p^e<A$, and it exceeds the term of
$\Sigma_w(\kappa)$ by at most $\kappa$. The odd $e\geq1$ with $p^e<A$ are at most $\log_pA$ in
number.
\end{proof}

\begin{lemma}\label{pd-gain-holds-numerics}
Let $0<\varepsilon<1/52$, let $K_G$ be a list of nonnegative constants as in
Definition~\ref{pd-gain-data}, and let $C_m\geq0$. There is $x_0$ such that the following holds for
every $x\geq x_0$. Let $D$ be a definite-side datum with ranks $k,m,t,n$ and top classes $T$, with size
data with scale $(x,\varepsilon,J)$ that satisfy the size bounds chosen with families
(Definition~\ref{pd-family}(4)), and with gain data with constants $K_G$ that satisfy condition (5) of
Definition~\ref{pd-gain-admissibility} with $\delta=\varepsilon/48$. Use the notation of
Definitions~\ref{hf-good-orbit} and~\ref{pd-gain-data}, with $K_\kappa$ the constant of
Definition~\ref{pd-gain-data}(d). Then:
\begin{enumerate}
\item $K_\kappa r_\varepsilon^2k\,x^{-26}\leq x^{-12}$ and $C_mN_1x^{-13}\leq x^{-12}$;
\item $L_2<u$ and $3G'\leq x^{\varepsilon/4}$;
\item $a/(kx^{13})\geq1$ and $N^*x^\varepsilon\Lambda b'/k\geq1$, so that $W_a$ and $W_b$ are not
  empty;
\item with $y=x^{\varepsilon/3}$,
  $(x^\varepsilon+1)\,m\,|T|\,x^{-12}\,y(2+y)\leq x^{-9}/4$.
\end{enumerate}
\end{lemma}
\leanref{Formalism/SharpRestriction/GrowthBounds/LogarithmicSums/MeanNum.lean}{Formalism.MainNotions.AuxBlock8.Pf.numMean_ev}
\leanref{Formalism/SharpRestriction/GrowthBounds/LogarithmicSums/SlopeEstimateOfLemmaBSlopeLine.lean}{Formalism.MainNotions.AuxBlock8.Pf.NumMean}

\begin{lemma}\label{pd-gain-holds-hypothesis-one}
Let $0<\varepsilon<1/52$, let $K_G$ be a list of nonnegative constants, and let $C_m\geq0$ and
$C_{31}>0$. There is $x_0$ such that the following holds for every $x\geq x_0$. Let $D$ be a
definite-side datum with size data with scale $(x,\varepsilon,J)$ and gain data with constants $K_G$,
and suppose that the gain data satisfy conditions (1) with $C_{31}$, (4), (5) with
$\delta=\varepsilon/48$ and (12) of Definition~\ref{pd-gain-admissibility}, and the conclusions of
Lemma~\ref{pd-gain-holds-numerics} with $C_m$. Then Hypothesis I of
Lemma~\ref{pd-gain-first-hypotheses} holds at $(p,n_p,N(p,t_0))$ for every candidate prime
$p\geq p_0$ and every good representative $t_0$ at $p$.
\end{lemma}
\leanref{Formalism/SharpRestriction/GrowthBounds/LogarithmicSums/MeanNum.lean}{Formalism.MainNotions.AuxBlock8.Pf.restHypI_ev}
\leanref{Formalism/SharpRestriction/GrowthBounds/LogarithmicSums/SlopeEstimateOfLemmaBSlopeLine.lean}{Formalism.MainNotions.AuxBlock8.Pf.RestHypI}

\begin{lemma}\label{pd-gain-holds-star-numerics}
Let $0<\varepsilon<1/52$, let $K_G$ be a list of nonnegative constants as in
Definition~\ref{pd-gain-data}, with $K_t,K_D$ as in Lemma~\ref{pd-gain-second-hypotheses},
$K_{D,\mathrm I},K_1,K_2$ as in Lemma~\ref{pd-gain-first-hypotheses} and $K_{\mathrm{Pr}}$ as in
Definition~\ref{pd-served-lattice}, and let $C_{31}>0$. There are $C_*\geq0$ and $x_0$ such that the
following holds for every $C\geq C_*$ and every $x\geq x_0$. Let $D$ be a definite-side datum with
ranks $k,m,t,n$, size data with scale $(x,\varepsilon,J)$ that satisfy the size bounds chosen with
families and have parameter choices (Definition~\ref{pd-family}(4)), and gain data with constants
$K_G$ that satisfy condition (5) of Definition~\ref{pd-gain-admissibility} with
$\delta=\varepsilon/48$. Use the notation of Definitions~\ref{hf-good-orbit} and~\ref{pd-gain-data}
and Lemmas~\ref{pd-gain-second-hypotheses} and~\ref{pd-gain-first-hypotheses}, with $r=r_\varepsilon$,
$Q=p^{G'}$ for a prime $p$, $Y=x^{13}$, $G_x=C_{31}\log x/\delta_1$ and $\beta=b^{x^{-\varepsilon/2}}$.
Then:
\begin{enumerate}
\item $L_2<u$, $3G'\leq x^{\varepsilon/4}$ and $\log x\geq1$;
\item $G_x(w''+\eta+c_G+N(m))\leq b'$ and $c_0\leq3b'$;
\item $kQ\leq b'$ for every prime $p<p_0$, and $K_D(N_1^3+r^5+k^2)\log_pa\leq b'$ for every prime
  $p$;
\item $3\leq C$, $K_tr\leq C$, $13+4K_t+160C_{31}\leq C$, $K_{\mathrm{Pr}}C_{31}\leq C$ and
  $K_1N_1\leq C\log x$;
\item for every prime $p$ with $p_0\leq p<X'$ and every integer $0\leq n'\leq n$,
  \[
   K_{\mathrm{Pr}}G_x(\eta+c_G+N(m))+K_2\big(Qc_6+kQ+k\log_pa+c_G\big)+2G_xc_G
   +40C_{31}\log x\cdot K_{D,\mathrm I}Y^2(N_1^5+k^3)\log_pa+3\leq\beta,
  \]
  where $c_6=c_6(p,n')$ is as in Lemma~\ref{pd-gain-first-hypotheses}.
\end{enumerate}
\end{lemma}
\leanref{Formalism/SharpRestriction/GrowthBounds/LogarithmicSums/StarNum.lean}{Formalism.MainNotions.AuxBlock8.Pf.St.numStar_ev}
\leanref{Formalism/SharpRestriction/GrowthBounds/LogarithmicSums/StarNum.lean}{Formalism.MainNotions.AuxBlock8.Pf.St.CStar}
\leanref{Formalism/SharpRestriction/GrowthBounds/LogarithmicSums/PartnerCountOfCentral.lean}{Formalism.MainNotions.AuxBlock8.Pf.St.NumStar}
\leanref{Formalism/SharpRestriction/GrowthBounds/LogarithmicSums/StarNumB.lean}{Formalism.MainNotions.AuxBlock8.Pf.St.B.facts_ev}
\leanref{Formalism/SharpRestriction/GrowthBounds/LogarithmicSums/StarNumB.lean}{Formalism.MainNotions.AuxBlock8.Pf.St.B.Facts}

\begin{lemma}\label{pd-gain-holds-star}
Let $\mathcal F$ be a family of definite-side data with exponent $\varepsilon$ whose members carry gain
data with Weyl exponents, one list $K_G$ of nonnegative constants, local models and lattice
instances, and suppose that $\mathcal F$ is admissible (Definition~\ref{pd-gain-admissibility}) with
constants $C_{31}>0$ and $C_e,C_m\geq0$. There is $C_0\geq0$ such that for every $C\geq C_0$ the
following holds at all large scales $x$, for every member $D$ at $x$, with $\beta=b^{x^{-\varepsilon/2}}$.
Let $\mathbf a',\mathbf b'\colon A\to\mathbf Z$ satisfy $a\leq\mathbf a'(t)\leq a+a/(kx^{13})$ and
$\mathbf b'(t)\geq b'$ for $t\in T$, and form the gain quantities of
Definition~\ref{pd-gain-quantities} with $\mathbf a',\mathbf b'$ in place of $\mathbf a,\mathbf b$ and
with $C$ and $\beta$. Then for every candidate prime $p$, every good representative $t_0$ at $p$ and
every lift $j$,
\[
 \sum_{s\in T(p,t_0)}\big(\mathrm{VM}^0_{j,s}(p)-\mathrm{lo}^G(p)\big)\leq g_U(p,j,t_0)+\delta(p,t_0)
 -\mathrm{Pr}(p,j,t_0)-\max(\lceil t(p,t_0)\rceil,0)-D(p)-1,
\]
with $T(p,t_0)$, $g_U$, $\mathrm{Pr}$ and $\delta(p,t_0)$ as in Definitions~\ref{pd-served-lattice}
and~\ref{pd-served-selection}.
\end{lemma}
\leanref{Formalism/SharpRestriction/GrowthBounds/LogarithmicSums/StarFam.lean}{Formalism.MainNotions.AuxBlock8.Pf.St.star_fam}
\leanref{Formalism/SharpRestriction/GrowthBounds/LogarithmicSums/Conclusion1OfGainBoundAll.lean}{Formalism.MainNotions.AuxBlock8.Pf.StarAt}
\leanref{Formalism/SharpRestriction/GrowthBounds/LogarithmicSums/Conclusion1OfGainBoundAll.lean}{Formalism.MainNotions.AuxBlock8.Pf.StarAll}
\leanref{Formalism/SharpRestriction/GrowthBounds/LogarithmicSums/PartnerCountOfCentral.lean}{Formalism.MainNotions.AuxBlock8.Pf.St.starAll_of}

\begin{proof}
Let $C_*$ and $x_0$ be as in Lemma~\ref{pd-gain-holds-star-numerics} for $\varepsilon$, $K_G$ and
$C_{31}$, and $C_0=C_*$. Let $C\geq C_0$. By admissibility with $\delta=\varepsilon/48$, at all large
scales every member satisfies conditions (1) and (5) of Definition~\ref{pd-gain-admissibility} with
$(C_{31},C_e,C_m,\varepsilon/48)$; at all large scales every member also has parameter choices for its
size data and satisfies the size bounds chosen with the family. So at all large scales $x\geq x_0$
the conclusions of Lemma~\ref{pd-gain-holds-star-numerics} hold for every member. Fix such a member
$D$, $\mathbf a',\mathbf b'$ as in the statement, a candidate prime $p$, a good representative $t_0$
at $p$ and a lift $j$, and write $a'_s=\mathbf a'(s)$, $b'_s=\mathbf b'(s)$, $N=N(p,t_0)$,
$r=r_\varepsilon\geq1$ and $G_x=C_{31}\log x/\delta_1$.

\emph{Basic facts.} As in the proof of Lemma~\ref{pd-gain-holds-gain-bound}, the parameter choices
give $0\leq c_G\leq w''\leq b'$, $w''\geq500$, $2c_G\leq a$ and $0\leq\kappa_j\leq N(m)$; so $2c_G\leq
a\leq a'_s\leq2a$ for $s\in T$. Put $\varphi_j(z)=\mathrm{fill}_j(z)+\kappa_j$, formed with
$\mathbf b'$; then $\varphi_j(z)\geq1$ at every place, as there. Since $p<X'$ and $a/X'=b'\geq1$,
$p\leq a$, so $\log_p(2a)\geq1$. Put $\gamma_p=p/(p^2-1)$, so that $0\leq\gamma_p\leq1$. As $p$ is
of inert class, $\ell_p=n_p$, $s_p=h_p$, $\operatorname{pos}_p((0,s),e)=(0,s+eh_p)$, and $T(p,t_0)$
consists of the $t_0+ih_p\in T$ with $0\leq i<n_p$, so $N=|T(p,t_0)|$.

\emph{Orbit sizes.} By Lemma~\ref{pd-gain-holds-star-numerics}(1), $L_2<u$ and $3G'\leq
x^{\varepsilon/4}\leq n_p$. So condition (1) applies at $t_0$ and gives
$n_p\delta_1/(C_{31}\log x)\leq N$. Hence $N\geq1$, $n_p\delta_1\leq NC_{31}\log x$ and
$n_p\leq NG_x$.

\emph{A lower bound for the unit gain.} For $s\in T(p,t_0)$ put
$m_j(s)=\sum_{e\ \mathrm{odd}}\min\big(a'_s/p^e,\varphi_j((0,s+eh_p))\big)\geq0$. Since
$\max(0,y-w)+\min(y,w)=y$ and $\sum_{e\ \mathrm{odd}}p^{-e}=\gamma_p$,
$\mathrm{VM}^0_{j,s}(p)+m_j(s)=\gamma_pa'_s$. If $s$ is at an even position, this gives
$\mathrm{VM}^0_{j,s}(p)+m_j(s)-3c_G\log_p(2a)\leq\gamma_p(a'_s-2c_G)$, because $2\gamma_pc_G\leq
3c_G\log_p(2a)$. If $s$ is at an odd position, put $w'_e=\mathrm{fill}_j((0,s+eh_p))+c_G$, so that
$0\leq w'_e\leq\varphi_j((0,s+eh_p))+c_G$. Lemma~\ref{pd-gain-holds-series}(2), with $y=a'_s$,
$c=c_G$, $A=2a$, $w_e=\varphi_j((0,s+eh_p))$ and $Q$ the odd positive integers, together with the
same identity for $y=a'_s-2c_G$ and the weights $w'_e$, gives
\[
 \mathrm{VM}^0_{j,s}(p)-3c_G\log_p(2a)\leq\gamma_p(a'_s-2c_G)-\sum_{e\ \mathrm{odd}}
 \min\Big(\frac{a'_s-2c_G}{p^e},w'_e\Big).
\]
Summing over $T(p,t_0)$ and comparing with Definition~\ref{pd-served-lattice}(1),
\[
 \sum_{s\in T(p,t_0)}\mathrm{VM}^0_{j,s}(p)+\sum_{s\ \mathrm{even}}m_j(s)-3Nc_G\log_p(2a)\leq
 g_U(p,j,t_0)+\delta(p,t_0),\tag{$*$}
\]
the second sum over the $s\in T(p,t_0)$ at even positions. As $C\geq3$ by
Lemma~\ref{pd-gain-holds-star-numerics}(4), $3c_G\log_p(2a)\leq C(c_G+\eta)\log_p(2a)$, the first
term of $\mathrm{lo}^G(p)$. By $(*)$ it therefore suffices to show
\[
 \mathrm{Pr}(p,j,t_0)+\max(\lceil t(p,t_0)\rceil,0)+D(p)+1\leq\sum_{s\ \mathrm{even}}m_j(s)
 +N\big(\mathrm{lo}^G(p)-C(c_G+\eta)\log_p(2a)\big).\tag{$\dagger$}
\]
Put $u_0=w''+\eta+c_G+N(m)$. By Lemma~\ref{pd-gain-holds-star-numerics}(2) and $n_p\leq NG_x$,
$n_pu_0\leq NG_xu_0\leq Nb'$, and so $n_pc_G\leq Nb'$.

\emph{The case $p<p_0$.} Here $\mathrm{Pr}(p,j,t_0)=r\big(N(2b'+c_G)+n_pu_0\big)\leq
3r(Nb'+n_pu_0)\leq6rNb'$. Put $S_1=\Lambda a/(pk)+ap^{-L_1(p)}\geq0$. Then
$t(p,t_0)=t_{\mathrm{II}}(p,n_p,N)=K_trN(S_1+c_0+kQ)+2rn_pc_G\geq0$, so $\max(\lceil
t\rceil,0)\leq t+1$, and $D(p)=D_{\mathrm{II}}(p,n_p)\leq40n_p\delta_1D_w+2rn_pc_G+40rn_p\delta_1c_0+1$
with $D_w=K_D(N_1^3+r^5+k^2)\log_pa$. Also
$\mathrm{lo}^G(p)-C(c_G+\eta)\log_p(2a)=Crb'\log x+C\log x\,S_1$. Put $Z=Nb'r\log x\geq1$. By
Lemma~\ref{pd-gain-holds-star-numerics}(2) and~(3) and $n_p\delta_1\leq NC_{31}\log x$:
$K_trN(c_0+kQ)\leq4K_tZ$; $2rn_pc_G\leq2Z$, twice; $40n_p\delta_1D_w\leq40C_{31}Z$; and
$40rn_p\delta_1c_0\leq120C_{31}Z$. With the three terms $1$, each at most $Z$, the left side of
$(\dagger)$ is at most $(13+4K_t+160C_{31})Z+K_trNS_1$, which is at most $CZ+CN\log x\,S_1$ by
Lemma~\ref{pd-gain-holds-star-numerics}(4). This is $N\big(\mathrm{lo}^G(p)-C(c_G+\eta)\log_p(2a)\big)$,
and $m_j\geq0$, so $(\dagger)$ holds.

\emph{The case $p\geq p_0$.} The order $n_p$ is even: by a condition of the size data the index $o$
of the Galois datum is at most $x^{\varepsilon/100}<x^{\varepsilon/4}\leq n_p$, and an element of $A$
of odd order has order at most $o$, as in the proof of Lemma~\ref{pd-gain-link-admissible-places}.
Here
\[
 \mathrm{Pr}(p,j,t_0)=\sum_{s\ \mathrm{odd}}\big(\mathrm{fill}_j((0,s))-b'_s\big)
 +K_{\mathrm{Pr}}n_pu_0,
\]
the sum over the $s\in T(p,t_0)$ at odd positions. Write $P=P(p,t_0)$ and $a'_w(i)=a'_{t_0+ih_p}$,
$b'_w(i)=b'_{t_0+ih_p}$. Then
$t(p,t_0)=t_{\mathrm I}(p,n_p,P,a'_w,b'_w)=S_0+K_1NN_1S_2+K_2NB+2n_pc_G$, where
$S_0=\sum_y\min\big(b'_w(y),a'_w(y-d(y))/p^{d(y)}\big)$ over the odd $y\in P$ with $d(y)$ defined,
$S_2=a/p^2+\Lambda a/(pk)$ and $B=Qc_6(p,n_p)+kQ+k\log_pa+c_G$; all terms are $\geq0$, so
$\max(\lceil t\rceil,0)\leq t+1$. Also $D(p)=D_{\mathrm I}(p,n_p)\leq40n_p\delta_1E+1$ with
$E=K_{D,\mathrm I}Y^2(N_1^5+k^3)\log_pa$, and
$\mathrm{lo}^G(p)-C(c_G+\eta)\log_p(2a)=b'/x+\eta+Cw''\log x/\delta_1+\beta+C\log x\,S_2$.

Let $s=t_0+yh_p$ be at an odd position. Then $\mathrm{fill}_j((0,s))-b'_s=o_j((0,s))-\kappa_j\leq
b'/x+\eta$, by Definition~\ref{pd-gain-data}(2) and $\kappa_j\geq0$. If $d=d(y)$ is defined, then
$i_y=y-d$ is an even element of $P$, $d$ is odd because $n_p$ is even, and $b'_w(y)\leq
b'_s+o_j((0,s))=\varphi_j((0,s))$; so the term of $S_0$ at $y$ is at most
$\min\big(a'_w(i_y)/p^d,\varphi_j((0,t_0+(i_y+d)h_p))\big)$, the term $e=d$ of $m_j(t_0+i_yh_p)$. The
map $y\mapsto(i_y,d(y))$ is injective, since $y=i_y+d(y)$, and the terms of the $m_j$ are $\geq0$.
Since at most $N$ elements of $T(p,t_0)$ are at odd positions,
\[
 \sum_{s\ \mathrm{odd}}\big(\mathrm{fill}_j((0,s))-b'_s\big)+S_0\leq\sum_{s\ \mathrm{even}}m_j(s)
 +N(b'/x+\eta).
\]
For the other terms: $K_{\mathrm{Pr}}n_pu_0\leq NK_{\mathrm{Pr}}G_xw''+NK_{\mathrm{Pr}}G_x(\eta+c_G+N(m))$
and $K_{\mathrm{Pr}}G_xw''\leq Cw''\log x/\delta_1$ as $K_{\mathrm{Pr}}C_{31}\leq C$; $K_1N_1NS_2\leq
CN\log x\,S_2$; $2n_pc_G\leq2NG_xc_G$; $40n_p\delta_1E\leq40NC_{31}\log x\,E$; and $3\leq3N$. Adding,
the left side of $(\dagger)$ is at most
\[
 \sum_{s\ \mathrm{even}}m_j(s)+N\Big(\frac{b'}x+\eta+\frac{Cw''\log x}{\delta_1}+C\log x\,S_2\Big)
 +N\Big(K_{\mathrm{Pr}}G_x(\eta+c_G+N(m))+K_2B+2G_xc_G+40C_{31}\log x\,E+3\Big).
\]
Since $n_p$ divides the order $n$ of $A$, Lemma~\ref{pd-gain-holds-star-numerics}(5) with $n'=n_p$
bounds the last bracket by $\beta$, and $(\dagger)$ follows.
\end{proof}

\begin{lemma}\label{pd-gain-holds-gain-bound}
Let $D$ be a definite-side datum with ranks $k,m,t,n$, $f=2n$ and $N(j)=k+j(t+1)$, top classes $T$,
size data with the derived quantities of Definition~\ref{pd-size-data}, parameter choices for its size
data (Definition~\ref{pd-family}(4)), and gain data with Weyl exponents $e_p$ and nonnegative constants
$K_G$, among them the constants $K_{\delta,1},K_{\delta,2}$ of Definition~\ref{pd-served-selection}.
Let $C_e\geq0$ and $C,\beta$ be real numbers with $C\geq3$, $2K_{\delta,1}+C_e\leq C$,
$4K_{\delta,2}\leq C$ and $\beta\geq0$, and suppose that $e_p\leq C_efk^2\log_pa$ for every prime
$p\leq X''$. Let $\mathbf a',\mathbf b'\colon A\to\mathbf Z$ satisfy $a\leq\mathbf a'(t)\leq
a+a/(kx^{13})$ and $\mathbf b'(t)\geq b'$ for $t\in T$, form the gain quantities of
Definition~\ref{pd-gain-quantities} with $\mathbf a',\mathbf b'$ in place of $\mathbf a,\mathbf b$ and
with $C$ and $\beta$, and suppose that the inequality of Lemma~\ref{pd-gain-holds-star} holds for every
candidate prime $p$, good representative $t_0$ at $p$ and lift $j$. Then for every prime $p$ inert
for $D$ with $p\leq X''$: $0\leq g_j(p)\leq4|T|a/p$ for every lift $j$, and
\[
 \sum_{j=1}^m\sum_{s\in T(p)}\big(\mathrm{VM}^0_{j,s}(p)-\mathrm{lo}^G(p)\big)-m\,\mathrm{fl}(p)
 \leq\sum_{j=1}^mg_j(p)-\tau(p).
\]
\end{lemma}
\leanref{Formalism/SharpRestriction/GrowthBounds/LogarithmicSums/NotSingularOfCentral.lean}{Formalism.MainNotions.AuxBlock8.Pf.gain_bound_at_of_star}
\leanref{Formalism/SharpRestriction/GrowthBounds/LogarithmicSums/NotSingularOfCentral.lean}{Formalism.MainNotions.AuxBlock8.Pf.gain_bound_all_of_star}
\leanref{Formalism/SharpRestriction/GrowthBounds/LogarithmicSums/Conclusion1OfGainBoundAll.lean}{Formalism.MainNotions.AuxBlock8.Pf.GainBoundAtPairUniform}

\begin{proof}
As in the proof of Lemma~\ref{pd-gain-link-surplus}, the parameter choices give $k\geq5$,
$N(m)\leq S_c$, $S_c\geq100$ and $kS_c\leq w''$, and $w''\leq b'/x\leq b'$ and $2b'\leq a$. So
$c_G=4S_c$ satisfies $2c_G\leq a$, and $0\leq\kappa_j\leq N(m)\leq S_c$ for every lift $j$. Put
$\gamma_p=p/(p^2-1)$, so that $0\leq\gamma_p\leq4/(3p)$, and $\alpha_s=\mathbf a'(s)-2c_G$ for
$s\in T$. Since $kx^{13}\geq5$, $a\leq\mathbf a'(s)\leq6a/5\leq2a$, so $0\leq\alpha_s\leq6a/5$. For a
lift $j$ and a place $z$ put $\phi_j(z)=\mathrm{fill}_j(z)+\kappa_j$, the weight in
$\mathrm{VM}^0$. Then $\phi_j(z)\geq1$: at a top place $\phi_j(z)=\mathbf b'(s)+o_j(z)\geq b'+b'/x$,
at a thick place $\phi_j(z)=o_j(z)\geq b\geq2$, and elsewhere $\phi_j(z)=o_j(z)\geq w''\geq500$. Hence
$0\leq\mathrm{fill}_j(z)+c_G\leq\phi_j(z)+c_G$.

\emph{Upper bounds.} The subtracted minima in $g_U(p,j,t_0)$ are $\geq0$, and $\delta(p,t_0)\geq0$,
so $g_U(p,j,t_0)\leq\gamma_p\sum_{s\in T(p,t_0)}\alpha_s$. By Lemma~\ref{pd-gain-holds-series}(1),
with $y=\alpha_s$ and $w_e=\mathrm{fill}_j+c_G\geq0$ at $\operatorname{pos}_p((0,s),e)$,
$g_0(p,j,t_0)\leq\gamma_p\sum_{s\in T(p,t_0)}\alpha_s$. Next $\mathrm{Pr}(p,j,t_0)\geq0$: if $p\geq
p_0$, its summands are $o_j((0,s))-\kappa_j\geq b'/x-S_c\geq w''-S_c\geq0$ and its last term is
$\geq0$; if $p<p_0$, all its terms are $\geq0$. So $\max(g_U-\mathrm{Pr},g_0^+)$ and $g_0^+$ are at most
$\gamma_p\sum_{s\in T(p,t_0)}\alpha_s$. Every top class has one representative, so
\[
 0\leq g_j(p)\leq\gamma_p\sum_{s\in T}\alpha_s\leq\frac4{3p}\,|T|\,\frac{6a}5\leq\frac{4|T|a}p .
\]

\emph{Lower terms.} The representatives at $p$ are at most $f/\ell_p$ in number, $\sum_{t_0}N(p,t_0)=|T|\leq
f$, and $\log_p(2a)\leq2\log_pa$ since $a\geq2$. So
$\sum_{t_0}\delta(p,t_0)\leq2K_{\delta,1}fN(m)^2\log_pa+4K_{\delta,2}f(\log_pa)^2$, the sum over all
representatives. As $e_p\leq C_efk^2\log_pa\leq C_efN(m)^2\log_pa$, $2K_{\delta,1}+C_e\leq C$ and
$4K_{\delta,2}\leq C$,
\[
 e_p+\sum_{t_0}\delta(p,t_0)\leq\mathrm{fl}(p).
\]
Also $\mathrm{lo}^G(p)\geq3c_G\log_p(2a)$, since $C\geq3$, $\eta>0$, $\beta\geq0$ and the other terms
of $\mathrm{lo}^G(p)$ are $\geq0$.

\emph{The case $V(p)=\emptyset$.} Then $\tau(p)=0$. Suppose first that $p$ is a signal prime or
$p\geq X'$, so that $T(p)=T$. If $p$ is of inert class with $h_p=0$, then $p\geq X'$ and
$\operatorname{pos}_p((0,s),e)=(0,s)$ for all $e$; as $\mathbf a'(s)/p^e\leq\mathbf
a'(s)/X'\leq b'+b'/x\leq\phi_j((0,s))$, every $\mathrm{VM}^0_{j,s}(p)$ vanishes, and the left side is at
most $0\leq\sum_jg_j(p)$, because $\mathrm{lo}^G$ and $\mathrm{fl}$ are $\geq0$. Otherwise
$g_j(p)=\sum_{t_0}g_0^+(p,j,t_0)$, the sum over all representatives. Let $s\in T(p,t_0)$ and let $e$
be excluded for $s$. A signal prime excludes no $e$, since $\operatorname{pos}_p((0,s),e)$ is not a top
place for odd $e$; and if $p\geq X'$, then $\operatorname{pos}_p((0,s),e)=(0,s')$ is a top place and
$\mathbf a'(s)/p^e\leq b'+b'/x\leq\phi_j((0,s'))$. So the excluded terms of
$\mathrm{VM}^0_{j,s}(p)$ vanish, and Lemma~\ref{pd-gain-holds-series}(2), with $y=\mathbf a'(s)$,
$c=c_G$, $A=2a$, $w_e=\phi_j(\operatorname{pos}_p((0,s),e))$, $w'_e=w_e-\kappa_j+c_G$ and $Q$ the
odd $e$ that are not excluded, gives
\[
 \sum_{s\in T(p,t_0)}\big(\mathrm{VM}^0_{j,s}(p)-3c_G\log_p(2a)\big)-\delta(p,t_0)\leq g_0(p,j,t_0).
\]
With $\mathrm{lo}^G(p)\geq3c_G\log_p(2a)$, summing over $t_0$ and $j$, and the bound on the lower
terms, the left side of the assertion is at most $\sum_j\sum_{t_0}g_0^+(p,j,t_0)=\sum_jg_j(p)$. If
$p$ is of inert class and $p<X'$, then $T(p)=\emptyset$ and the left side is $-m\,\mathrm{fl}(p)\leq0$.

\emph{The case $V(p)\neq\emptyset$.} Then $p$ is a candidate prime, so $p<X'$, $p$ is of inert class
and $h_p\neq0$; $V(p)$ is the set of good representatives, $T(p)$ is the set of top classes with
representative in $V(p)$, and
$\tau(p)=e_p+m\sum_{t_0\in V(p)}(\max(\lceil t(p,t_0)\rceil,0)+D(p)+1)$. The other terms of $g_j(p)$
being $\geq0$, $\sum_{t_0\in V(p)}(g_U-\mathrm{Pr})(p,j,t_0)\leq g_j(p)$. Summing the hypothesis over
$t_0\in V(p)$ and $j$,
\[
 \sum_{j=1}^m\sum_{s\in T(p)}\big(\mathrm{VM}^0_{j,s}(p)-\mathrm{lo}^G(p)\big)\leq
 \sum_{j=1}^mg_j(p)+m\sum_{t_0\in V(p)}\delta(p,t_0)-\big(\tau(p)-e_p\big).
\]
Since $m\geq1$ and $e_p\geq0$, $e_p+m\sum_{t_0\in V(p)}\delta(p,t_0)\leq
m\big(e_p+\sum_{t_0}\delta(p,t_0)\big)\leq m\,\mathrm{fl}(p)$, which gives the assertion.
\end{proof}

\begin{lemma}\label{pd-gain-holds-high-classes}
Let $D$ be a definite-side datum with ranks $k,m,t,n$, top classes $T$, size data with the derived
quantities of Definition~\ref{pd-size-data} and cloud size $S_c$, parameter choices for its size data
(Definition~\ref{pd-family}(4)), and gain data, with $\kappa_j$ and $\mathrm{fill}_j$ as in
Definition~\ref{pd-lattice-instance}. Put $\theta_W=1/(kx^{13})$ and $\theta_S=\theta_W+3S_c/w''$. Let
$\mathbf a',\mathbf b'\colon A\to\mathbf Z$ satisfy $a\leq\mathbf a'(t)\leq a+a/(kx^{13})$ and
$\mathbf b'(t)\geq b'$ for $t\in T$, and form the class sums $\mathrm{VM}^G_{j,s}(p)$ of
Definition~\ref{pd-gain-quantities}(1) with $\mathbf a',\mathbf b'$. Then $0<\theta_W\leq\theta_S\leq
4/5$, and for every prime $p>X''$, every lift $j$ and every $s\in T$: $\mathrm{VM}^G_{j,s}(p)\leq
\theta_Ww''+S_c$ if $p<X''(1+\theta_S)$, and $\mathrm{VM}^G_{j,s}(p)=0$ otherwise.
\end{lemma}
\leanref{Formalism/SharpRestriction/GrowthBounds/LogarithmicSums/SummableAndTsumLeTwoMulDivOfLeDivPow.lean}{Formalism.MainNotions.AuxBlock8.Pf.G5.ask_high}
\leanref{Formalism/SharpRestriction/GrowthBounds/LogarithmicSums/SummableAndTsumLeTwoMulDivOfLeDivPow.lean}{Formalism.MainNotions.AuxBlock8.Pf.G5.thS_key}
\leanref{Formalism/SharpRestriction/GrowthBounds/LogarithmicSums/SummableAndTsumLeTwoMulDivOfLeDivPow.lean}{Formalism.MainNotions.AuxBlock8.Pf.G5.fillG_ge}

\begin{proof}
As in the proof of Lemma~\ref{pd-gain-link-surplus}, the parameter choices give $k\geq5$, $S_c\geq100$,
$0\leq\kappa_j\leq N(m)\leq S_c$ and $kS_c\leq w''$, so $5S_c\leq w''$. Hence $0<\theta_W\leq1/5$,
$0\leq3S_c/w''\leq3/5$ and $\theta_W\leq\theta_S\leq4/5$, and a direct expansion gives
$w''(1+\theta_W)\leq(w''-S_c)(1+\theta_S)$. Put $\varphi_j(z)=\mathrm{fill}_j(z)+\kappa_j$. By
Definition~\ref{pd-gain-data}(2), $\varphi_j(z)\geq w''$: at a top place $\varphi_j(z)=\mathbf
b'(s)+o_j(z)\geq b'\geq w''$, at a thick place $\varphi_j(z)=o_j(z)\geq b\geq b'/x\geq w''$, and
elsewhere $\varphi_j(z)=o_j(z)\geq w''$. So $\mathrm{fill}_j(z)\geq w''-S_c$ at every place.

Let $p>X''$ be a prime, $j$ a lift and $s\in T$, and write $a'_s=\mathbf a'(s)\geq0$. Since
$X''=a/w''\geq1$ and $a'_s\leq a(1+\theta_W)$, we have $p\geq2$ and $a'_s\leq
X''w''(1+\theta_W)<pw''(1+\theta_W)$. For odd $e\geq3$, $p^e\geq p^3\geq4p$, and
$w''(1+\theta_W)\leq\frac65w''\leq4\cdot\frac45w''\leq4(w''-S_c)$, so $a'_s/p^e\leq w''-S_c\leq
\mathrm{fill}_j(\operatorname{pos}_p((0,s),e))$ and the term $e$ of $\mathrm{VM}^G_{j,s}(p)$ vanishes.
So $\mathrm{VM}^G_{j,s}(p)=\max\big(0,a'_s/p-\mathrm{fill}_j(\operatorname{pos}_p((0,s),1))\big)$. If
$p<X''(1+\theta_S)$, then $a'_s/p\leq w''(1+\theta_W)$, and the term is at most
$w''(1+\theta_W)-(w''-S_c)=\theta_Ww''+S_c$. If $p\geq X''(1+\theta_S)$, then $a'_s\leq
X''(w''-S_c)(1+\theta_S)\leq(w''-S_c)p$, so $a'_s/p\leq w''-S_c$ and the term vanishes.
\end{proof}

\begin{lemma}\label{pd-gain-holds-third-numerics}
Let $0<\varepsilon<1/52$ and $C\geq0$. There is $x_0$ such that the following holds for every $x\geq
x_0$. Let $D$ be a definite-side datum with ranks $k,m,t,n$ and $f=2n$, with size data with scale
$(x,\varepsilon,J)$ that satisfy the size bounds chosen with families and have parameter choices
(Definition~\ref{pd-family}(4)), and let $c_G=4S_c$ and $\theta_W,\theta_S$ be as in
Lemma~\ref{pd-gain-holds-high-classes}. Then
\[
 (X''\theta_S+2)\log(2X'')(\theta_Ww''+S_c)\leq a\,x^{-16},\qquad
 C(c_G+\eta)\log(2a)\,(X''+2)\leq b\,x^{-100},
\]
\[
 Cf\Big(N(m)^2\log a+\frac{(\log a)^2}{\log2}\Big)(X''+2)\leq b\,x^{-100}.
\]
\end{lemma}
\leanref{Formalism/SharpRestriction/GrowthBounds/LogarithmicSums/ThirdEstimateOfNumericFacts.lean}{Formalism.MainNotions.AuxBlock8.Pf.G5.Nu.numT_ev}
\leanref{Formalism/SharpRestriction/GrowthBounds/LogarithmicSums/ThirdEstimateOfNumericFacts.lean}{Formalism.MainNotions.AuxBlock8.Pf.G5.NumT}

\begin{lemma}\label{pd-gain-holds-fourth-numerics}
Let $0<\varepsilon<1/52$, let $K_G$ be a list of nonnegative constants as in
Definition~\ref{pd-gain-data}, and let $C\geq0$. There is $x_0$ such that the following holds for every
$x\geq x_0$. Let $D$ be a definite-side datum with ranks $k,m,t,n$, $f=2n$ and top classes $T$, with
size data with scale $(x,\varepsilon,J)$ that satisfy the size bounds chosen with families and have
parameter choices (Definition~\ref{pd-family}(4)), and gain data with constants $K_G$ that satisfy
condition (5) of Definition~\ref{pd-gain-admissibility} with $\delta=\varepsilon/48$. Let
$\beta=b^{x^{-\varepsilon/2}}$, $r=r_\varepsilon$, $\Lambda=K_\Lambda N_1^2\log a$, $c_G=4S_c$,
$\theta_W,\theta_S$ as in Lemma~\ref{pd-gain-holds-high-classes}, and
\[
 c_1=C(c_G+\eta)\log(2a),\qquad c_{\mathrm{fl}}=Cf\Big(N(m)^2\log a+\frac{(\log a)^2}{\log2}\Big),\qquad
 c_2=Crb'\log x+Ca\log x\,(kx^{26})^{-4},
\]
\[
 c_3=\frac{b'}x+\eta+\frac{Cw''\log x}{\delta_1}+\beta,\qquad c_4=C\log x\Big(\frac{\Lambda a}k+\frac a{p_0}\Big).
\]
Then $|T|\geq1$, $x\geq120$, $x^{-2+\varepsilon}\leq1/(120x)$, and
\[
 (X''+2)c_1\leq\frac a{40x},\quad(X''+2)c_{\mathrm{fl}}\leq\frac a{40x},\quad p_0c_2\log4\leq\frac
 a{40x},\quad X'c_3\log4\leq\frac{3a}{2x},
\]
\[
 \log X'(1+\log X')\,c_4\leq\frac a{40x},\qquad(X''\theta_S+2)\log(2X'')(\theta_Ww''+S_c)\leq\frac
 a{40x}.
\]
\end{lemma}
\leanref{Formalism/SharpRestriction/GrowthBounds/LogarithmicSums/PrimeSumNumericFacts.lean}{Formalism.MainNotions.AuxBlock8.Pf.G5.Nu.prime_sum_facts_eventually}
\leanref{Formalism/SharpRestriction/GrowthBounds/LogarithmicSums/PrimeSumNumericFacts.lean}{Formalism.MainNotions.AuxBlock8.Pf.G5.prime_sum_facts_of_scaled}
\leanref{Formalism/SharpRestriction/GrowthBounds/LogarithmicSums/SecondConstantNonneg.lean}{Formalism.MainNotions.AuxBlock8.Pf.G5.PrimeSumNumericFacts}

\begin{theorem}\label{pd-gain-law}
Let $\mathcal F$ be a family of definite-side data whose members carry gain data with Weyl exponents,
one list $K_G$ of nonnegative constants, local models and lattice instances, and suppose that
$\mathcal F$ is admissible (Definition~\ref{pd-gain-admissibility}). Then there are $C\in\mathbf R$,
a function $\delta$ with $\delta(x)\to0$ as $x\to\infty$, and $x_0$, such that for every scale
$x\geq x_0$ and every member $D$ at $x$, with $\beta=b^{\delta(x)}$, the gain quantities of $D$
(Definition~\ref{pd-gain-quantities}, with $C$ and $\beta$) satisfy the following.
\begin{enumerate}
\item For every prime $p$ inert for $D$ with $p\leq X''$: $0\leq g_j(p)\leq4|T|a/p$ for every lift
  $j$, and
  \[
   \sum_{j=1}^m\sum_{s\in T(p)}\big(\mathrm{VM}^G_{j,s}(p)-\mathrm{lo}^G(p)\big)-m\,\mathrm{fl}(p)
   \leq\sum_{j=1}^mg_j(p)-\tau(p),
  \]
  with $\mathrm{fl}(p)=Cf\,N(m)^2\log_pa+Cf(\log_pa)^2$ as in Definition~\ref{pd-gain-link}.
\item For every candidate prime $p$ and good representative $t_0$ at $p$, the mean over
  $P(p,t_0)$ of $(a',b')\mapsto\kappa(p;n_p,P(p,t_0),a',b')$ is at most $N(p,t_0)x^{-12}$, and if
  $p<p_0$, then $\kappa_{p,t_0}(a',b')\leq N(p,t_0)x^{-12}$ for all $a',b'\colon A\to\mathbf Z$.
\item For all $a',b'\colon A\to\mathbf Z$ with $a'(s)\leq2a$ and $b'(s)\geq0$ for $s\in T$, every
  prime $p$ inert for $D$, every lift $j$ and every $s\in T$, the class sums of
  Definition~\ref{pd-gain-quantities}(1) formed with $a',b'$ in place of $\mathbf a,\mathbf b$
  satisfy $\mathrm{VM}^0_{j,s}(p)\leq\mathrm{VM}^G_{j,s}(p)\leq\mathrm{VM}^0_{j,s}(p)+N(m)\log_p(2a)$.
\end{enumerate}
For the remaining parts, let $W_a,W_b$ be as in Definition~\ref{pd-gain-data}(g), let $B_T(W)$ be
the set of functions $A\to\mathbf Z$ with values in $W$ on $T$ and $0$ off $T$, and say that
$a'\in B_T(W_a)$ and $b'\in B_T(W_b)$ satisfy the \emph{weight requirement} if
\[
 \sum_p\frac{\log p}{p}\Big(2\,\big[\Sigma'(p)>x^{-\varepsilon}\big]+4\,\Sigma'(p)\Big)\leq x^{-9},
\]
the sum over the candidate primes, where $\Sigma'(p)$ is the deviation formed with $a',b'$ in place
of $\mathbf a,\mathbf b$ and $[\cdot]$ is $1$ if the condition holds and $0$ otherwise. In (4) and
(6) the gain quantities are formed with $a',b'$ in place of $\mathbf a,\mathbf b$.
\begin{enumerate}
\setcounter{enumi}{3}
\item For all $a'\in B_T(W_a)$ and $b'\in B_T(W_b)$ that satisfy the weight requirement: for every
  finite set $S$ of primes of inert class,
  \[
   \sum_{p\in S}\log p\cdot\max\Big(0,\ \sum_{j=1}^m\sum_{s\in T}\mathrm{VM}^G_{j,s}(p)
   -\sum_{j=1}^mg_j(p)+\tau(p)+\iota(p)\Big)\leq\frac{5m|T|a}x,
  \]
  and for every finite set $S$ of primes, $\sum_{p\in S}\log p\cdot\iota(p)\leq m|T|a\,x^{-9}$.
\item Some $a'\in B_T(W_a)$ and $b'\in B_T(W_b)$ satisfy the weight requirement.
\item For all $a'\in B_T(W_a)$ and $b'\in B_T(W_b)$ and every finite set $S$ of primes $p>X''$
  inert for $D$, $\sum_{p\in S}\log p\sum_{j=1}^m\sum_{s\in T}\mathrm{VM}^G_{j,s}(p)\leq
  m|T|a\,x^{-16}$; and for every finite set $S$ of primes $p\leq X''$,
  $\sum_{p\in S}\log p\cdot C(c_G+\eta)\log_p(2a)\leq b\,x^{-100}$ and
  $\sum_{p\in S}\log p\cdot\mathrm{fl}(p)\leq b\,x^{-100}$.
\end{enumerate}
Parts (4) to (6) are not used in this book.
\end{theorem}
\leanref{Formalism/SharpRestriction/GrowthBounds/LogarithmicSums.lean}{Formalism.MainNotions.AuxBlock8.Pf.stmtF}
\leanref{Formalism/SharpRestriction/GrowthBounds/SmallOrderPrimes.lean}{Formalism.MainNotions.AuxBlock8.KudlaFamilyStatement}
\leanref{Formalism/SharpRestriction/GrowthBounds/SmallOrderPrimes.lean}{Formalism.MainNotions.AuxBlock8.Conclusion₂}
\leanref{Formalism/SharpRestriction/GrowthBounds/SmallOrderPrimes.lean}{Formalism.MainNotions.AuxBlock8.Conclusion1}
\leanref{Formalism/SharpRestriction/GrowthBounds/SmallOrderPrimes.lean}{Formalism.MainNotions.AuxBlock8.Conclusion0}
\leanref{Formalism/SharpRestriction/GrowthBounds/IsResidueClassZero.lean}{Formalism.MainNotions.G.MeanKappaBound}

\begin{proof}
Let $\varepsilon$ be the exponent of $\mathcal F$, let $C_{31}>0$ and $C_e,C_m\geq0$ be constants for
which $\mathcal F$ is admissible, and let $K_{\delta,1},K_{\delta,2}$ be the constants of $K_G$ of
Definition~\ref{pd-served-selection}. Put $\delta(x)=x^{-\varepsilon/2}$, which tends to $0$, so that
$\beta=b^{x^{-\varepsilon/2}}\geq0$. Let $C_0$ be the constant of Lemma~\ref{pd-gain-holds-star},
$C_1=\max(C_0,3,2K_{\delta,1}+C_e+1,4K_{\delta,2})$ and $C=C_1+1$. Each part below holds at all large
scales, and $x_0$ is the largest of these thresholds. By admissibility, at all large scales every
member satisfies the conditions of Definition~\ref{pd-gain-admissibility} with $(C_{31},C_e,C_m,1)$,
and also with $(C_{31},C_e,C_m,\varepsilon/48)$.

\emph{Part (3).} Let $a',b'$, $p$, $j$ and $s$ be as there. By condition (12), $N(m)+1\leq w''$. For
every place $z$, $\mathrm{fill}_j(z)+\kappa_j\geq w''$: at a top place $(0,s')$ it is
$b'(s')+o_j(z)\geq b'/x\geq w''$, at a thick place it is $o_j(z)\geq b\geq w''$, and elsewhere it is
$o_j(z)\geq w''$. Since $0\leq\kappa_j\leq N(m)$, $\mathrm{fill}_j(z)\geq w''-N(m)\geq1$. If
$a'(s)<0$, all terms of both class sums vanish, and the assertion holds because
$N(m)\log_p(2a)\geq0$. Otherwise Lemma~\ref{pd-gain-holds-series}(3) and~(4), with $y=a'(s)\leq2a=A$,
$\kappa=\kappa_j$, $w_e=\mathrm{fill}_j(\operatorname{pos}_p((0,s),e))$ and $Q$ the set of odd
positive integers, give $\mathrm{VM}^0_{j,s}(p)\leq\mathrm{VM}^G_{j,s}(p)\leq\mathrm{VM}^0_{j,s}(p)
+\kappa_j\log_p(2a)$, and $\kappa_j\leq N(m)$.

\emph{Part (1).} At all large scales the inequality of Lemma~\ref{pd-gain-holds-star} holds with
$C_1\geq C_0$, every member has parameter choices for its size data, and condition (9) gives
$e_p\leq(C_e+1)fk^2\log_pa$ for every prime $p\leq a$, hence for $p\leq X''=a/w''$, as $w''\geq1$.
The exponents
satisfy $a\leq a_t\leq a+a/(kx^{13})$ and $b_t\geq b'$. So Lemma~\ref{pd-gain-holds-gain-bound},
with $\mathbf a'=\mathbf a$, $\mathbf b'=\mathbf b$, $C_e+1$ in place of $C_e$ and $C_1$ in place of
$C$, gives the bounds on $g_j(p)$ and the inequality of (1) with $\mathrm{VM}^0$ in place of
$\mathrm{VM}^G$ and with $\mathrm{lo}^G$ and $\mathrm{fl}$ formed with $C_1$. The term
$\mathrm{lo}^G(p)$ is an affine function of the constant, whose coefficient is at least
$(c_G+\eta)\log_p(2a)$, because its other parts $r(\varepsilon)b'\log x$, $w''\log x/\delta_1$ and
$\log x\,(\Lambda a/(pk)+ap^{-L_1(p)})$ or $\log x\,(\Lambda a/(pk)+a/p^2)$ are $\geq0$; so
$\mathrm{lo}^G(p)$ formed with $C$ exceeds $\mathrm{lo}^G(p)$ formed with $C_1$ by at least
$(c_G+\eta)\log_p(2a)$. Likewise $\mathrm{fl}(p)$ formed with $C_1$ is at most $\mathrm{fl}(p)$ formed
with $C$. By part (3), applied to $\mathbf a,\mathbf b$, $\mathrm{VM}^G_{j,s}(p)\leq
\mathrm{VM}^0_{j,s}(p)+N(m)\log_p(2a)$, and $N(m)\leq S_c\leq c_G$. So each term
$\mathrm{VM}^G_{j,s}(p)-\mathrm{lo}^G(p)$ formed with $C$ is at most the term
$\mathrm{VM}^0_{j,s}(p)-\mathrm{lo}^G(p)$ formed with $C_1$, and (1) follows.

\emph{Part (2).} At all large scales every member satisfies conditions (1), (4), (5) with
$\delta=\varepsilon/48$, (8) and (12), so Lemmas~\ref{pd-gain-holds-numerics}
and~\ref{pd-gain-holds-hypothesis-one} apply. Let $p$ be a candidate prime and $t_0$ a good
representative at $p$. As $p$ is of inert class, $s_p=h_p$, the positions of $t_0$ correspond to the
top classes with representative $t_0$, and $|P(p,t_0)|=N(p,t_0)$. If $p<p_0$, then
$\kappa(p;n_p,P(p,t_0),a',b')=\min\big(1,K_\kappa N(p,t_0)r_\varepsilon^2k\,x^{-26}\big)$ does not
depend on $a',b'$, is $\geq0$, and is at most $N(p,t_0)x^{-12}$ by
Lemma~\ref{pd-gain-holds-numerics}(1); this gives both assertions. If $p\geq p_0$, condition (8)
applies with $n'=n_p$ and $T=P(p,t_0)$: $T$ is tiled for $(x,\varepsilon,L_2)$ because $t_0$ is good;
$L_2<u$ and $3G'\leq x^{\varepsilon/4}\leq n_p$ by Lemma~\ref{pd-gain-holds-numerics}(2);
$p_0\leq p<X'$; $p$ is not in band, and a prime in the narrow band is in band because $\rho_0\leq\rho$
by condition (4), so $p$ is not in the narrow band; and Hypothesis I holds at $(p,n_p,N(p,t_0))$ by
Lemma~\ref{pd-gain-holds-hypothesis-one}. Condition (8) bounds the mean by
$C_mN(p,t_0)N_1x^{-13}$, which is at most $N(p,t_0)x^{-12}$ by Lemma~\ref{pd-gain-holds-numerics}(1).

\emph{Part (5).} Work at the large scales of part (2). For $a'\in B_T(W_a)$, $b'\in B_T(W_b)$ and a
candidate prime $p$, $\Sigma'(p)\geq0$, because the local constants are $\geq0$: $K_\kappa\geq0$, and
for $p\geq p_0$ condition (8) applies as in part (2). Put
$Z(a',b')=\sum_p\frac{\log p}p\big([\Sigma'(p)>x^{-\varepsilon}]+\Sigma'(p)\big)$, the sum over the
candidate primes. Since $[\Sigma'(p)>x^{-\varepsilon}]\leq x^\varepsilon\Sigma'(p)$,
$Z(a',b')\leq(x^\varepsilon+1)\sum_p\frac{\log p}p\Sigma'(p)$. For a good representative $t_0$ at
$p$, restriction along $i\mapsto t_0+ih_p$ maps $B_T(W_a)$ onto the functions on
$\mathbf Z/n_p\mathbf Z$ with values in $W_a$ on $P(p,t_0)$ and $0$ off it, with fibres of equal size,
and likewise for $W_b$. So the average of $\kappa_{p,t_0}(a',b')$ over $B_T(W_a)\times B_T(W_b)$ is
the mean over $P(p,t_0)$ of $\kappa(p;n_p,P(p,t_0),\cdot,\cdot)$, at most $N(p,t_0)x^{-12}$ by part
(2), and the average of $\Sigma'(p)$ is at most $mx^{-12}\sum_{t_0}N(p,t_0)\leq m|T|x^{-12}$. The
candidate primes are less than $X'=e^y$, $y=x^{\varepsilon/3}$, and
$\sum_{1\leq n<e^y}\frac{\log n}n\leq y\sum_{1\leq n<e^y}\frac1n\leq y(1+y)$. So the average of $Z$
over $B_T(W_a)\times B_T(W_b)$ is at most $(x^\varepsilon+1)\,m\,|T|\,x^{-12}\,y(2+y)\leq x^{-9}/4$,
by Lemma~\ref{pd-gain-holds-numerics}(4). These sets are not empty by
Lemma~\ref{pd-gain-holds-numerics}(3), so $Z(a',b')\leq x^{-9}/4$ for some $a'\in B_T(W_a)$ and
$b'\in B_T(W_b)$. For them the sum in the weight requirement is at most $4Z(a',b')\leq x^{-9}$, since
$\log p/p\geq0$ and $2[\cdot]+4\Sigma'(p)\leq4([\cdot]+\Sigma'(p))$.

\emph{Part (6).} At all large scales every member satisfies condition (12), so its size data satisfy
the size bounds chosen with the family, and it has parameter choices; so
Lemma~\ref{pd-gain-holds-third-numerics} applies with $C\geq0$. Let $a'\in B_T(W_a)$ and $b'\in
B_T(W_b)$, so that $a\leq a'(t)\leq a+a/(kx^{13})$ and $b'(t)\geq b'$ for $t\in T$, and let $\theta_W$,
$\theta_S$ be as in Lemma~\ref{pd-gain-holds-high-classes}. Let $S$ be a finite set of primes $p>X''$
inert for $D$. By that lemma the class sums are $\geq0$, and those at $p$ vanish unless
$X''<p<X''(1+\theta_S)$; for such $p$, $\log p\leq\log(2X'')$, as $\theta_S\leq1$, and
$\sum_{j=1}^m\sum_{s\in T}\mathrm{VM}^G_{j,s}(p)\leq m|T|(\theta_Ww''+S_c)$. At most $X''\theta_S+2$
integers lie in this interval, so the first sum of (6) is at most
$(X''\theta_S+2)\log(2X'')\,m|T|(\theta_Ww''+S_c)\leq m|T|a\,x^{-16}$ by
Lemma~\ref{pd-gain-holds-third-numerics}. Now let $S$ be a finite set of primes $p\leq X''$. It has at
most $X''+2$ elements. For $p\in S$, $\log p\cdot C(c_G+\eta)\log_p(2a)=C(c_G+\eta)\log(2a)$, and
\[
 \log p\cdot\mathrm{fl}(p)=Cf\Big(N(m)^2\log a+\frac{(\log a)^2}{\log p}\Big)\leq
 Cf\Big(N(m)^2\log a+\frac{(\log a)^2}{\log2}\Big).
\]
Multiplying by $X''+2$, Lemma~\ref{pd-gain-holds-third-numerics} gives the two bounds.

\emph{Part (4).} At all large scales every member satisfies conditions (2), (3) and (4), and
Lemma~\ref{pd-gain-holds-fourth-numerics} applies with $C\geq0$, by condition (12), the parameter
choices and condition (5) with $\delta=\varepsilon/48$; we use its notation. Let $a'\in
B_T(W_a)$ and $b'\in B_T(W_b)$ satisfy the weight requirement, so that $a\leq a'(t)\leq
a+a/(kx^{13})$ and $b'(t)\geq b'$ for $t\in T$, and form all gain quantities with $a',b'$.
Lemmas~\ref{pd-gain-holds-star} and~\ref{pd-gain-holds-gain-bound} and part (3) hold for such
weights, so the argument for part (1) gives, for every prime $p$ inert for $D$ with $p\leq X''$,
$g_j(p)\geq0$ and
\[
 \sum_{j=1}^m\sum_{s\in T(p)}\big(\mathrm{VM}^G_{j,s}(p)-\mathrm{lo}^G(p)\big)-m\,\mathrm{fl}(p)\leq
 \sum_{j=1}^mg_j(p)-\tau(p).
\]
As in part (5), $\Sigma'(p)\geq0$ at every candidate prime. For a prime $p$ put
$\Delta(p)=\sum_j\sum_{s\in T}\mathrm{VM}^G_{j,s}(p)-\sum_jg_j(p)+\tau(p)+\iota(p)$, write $[P]$ for $1$
if $P$ holds and $0$ otherwise, and put $w(p)=\frac{\log p}p\big(2[\Sigma'(p)>x^{-\varepsilon}]
+4\Sigma'(p)\big)$ if $p$ is a candidate prime and $w(p)=0$ otherwise.

\emph{One prime $p\leq X''$ of inert class.} Each term of $\mathrm{VM}^G_{j,s}(p)$ is at most
$a'_s/p^e$, and $\sum_{e\geq1\ \mathrm{odd}}p^{-e}\leq2/p$; as $a'_s\leq6a/5$, $0\leq
\mathrm{VM}^G_{j,s}(p)\leq3a/p$. With the inequality above, $\mathrm{lo}^G(p)\geq0$ and $|T(p)|\leq|T|$,
\[
 \Delta(p)\leq3ma\,\frac{|T\smallsetminus T(p)|}p+m|T|\,\mathrm{lo}^G(p)+m\,\mathrm{fl}(p)+\iota(p).
\]
We bound the four terms after multiplying by $\log p$.
\begin{enumerate}
\item[(i)] If $p\geq X'$, then $T(p)=T$. Let $p<X'$. If $n_p<x^{\varepsilon/4}$, $p$ lies in the set
  $A_3$ of the primes of condition (3), and $|T\smallsetminus T(p)|\leq|T|$. If $n_p\geq
  x^{\varepsilon/4}$ and $p$ is in band, $p$ lies in the set $A_4$ of the primes in band of condition
  (4), and $|T\smallsetminus T(p)|\leq|T|$. Otherwise $p$ is a candidate prime. If $\Sigma'(p)\leq
  x^{-\varepsilon}$, $V(p)$ is the set of good representatives, and $T\smallsetminus T(p)$ is the set
  of the top classes whose representative at $p$ is not good; let $B(p)$ be its size. If
  $\Sigma'(p)>x^{-\varepsilon}$, then $|T\smallsetminus T(p)|\leq|T|\leq(2+4\Sigma'(p))|T|$. So
  \[
   \frac{\log p}p|T\smallsetminus T(p)|\leq\big([p\in A_3]+[p\in A_4]\big)\frac{\log p}p|T|
   +[p\ \mathrm{candidate}]\frac{\log p}pB(p)+w(p)|T|.
  \]
\item[(ii)] $\log p\cdot C(c_G+\eta)\log_p(2a)=c_1$. If $p<p_0$, then $p^{-L_1(p)}\leq(kx^{26})^{-4}$
  by the definition of $L_1(p)$; if $p_0\leq p<X'$, then $a/p^2\leq a/(pp_0)$. Comparing with
  Definition~\ref{pd-gain-quantities}(5), $\log p\cdot\mathrm{lo}^G(p)\leq
  c_1+[p<p_0]\log p\cdot c_2+[p<X']\big(\log p\cdot c_3+\frac{\log p}pc_4\big)$.
\item[(iii)] As in part (6), $\log p\cdot\mathrm{fl}(p)\leq c_{\mathrm{fl}}$.
\item[(iv)] If $V(p)\neq\emptyset$, then $p$ is a candidate prime and
  $\iota(p)=2\Sigma'(p)\max_{t_0}S(p,t_0)$. By Lemma~\ref{pd-gain-link-surplus}, whose proof uses only
  these bounds on the exponents, $S(p,t_0)\leq2mN(p,t_0)a/p\leq2m|T|a/p$. Hence $\log
  p\cdot\iota(p)\leq4\frac{\log p}p\Sigma'(p)\,m|T|a\leq w(p)\,m|T|a$, which also holds when
  $V(p)=\emptyset$, as then $\iota(p)=0$.
\end{enumerate}
All these bounds are $\geq0$. So $\log p\cdot\max(0,\Delta(p))$ is at most
\[
 m|T|c_1+mc_{\mathrm{fl}}+[p<p_0]\log p\cdot m|T|c_2+[p<X']\Big(\log p\cdot m|T|c_3
 +\frac{\log p}pm|T|c_4\Big)+3m|T|a\big([p\in A_3]+[p\in A_4]\big)\frac{\log p}p
 +3ma[p\ \mathrm{candidate}]\frac{\log p}pB(p)+4m|T|a\,w(p).
\]

\emph{The primes $p\leq X''$.} Let $S$ be a finite set of primes of inert class. At most $X''+2$
primes are $\leq X''$. By Chebyshev's bound (Lemma~\ref{is-places-prime-sums}),
$\sum_{p<y}\log p\leq y\log4$ for $y>1$, and $\sum_{p<X'}\frac{\log p}p\leq\log
X'\sum_{n\leq X'}\frac1n\leq\log X'(1+\log X')$. Condition (3) bounds $\sum_{p\in A_3}\frac{\log p}p$
by $x^{-h_\varepsilon+\varepsilon}\leq x^{-2}\leq1/(120x)$, as $h_\varepsilon=2+8\varepsilon$ and
$x\geq120$; condition (4) bounds $\sum_{p\in A_4}\frac{\log p}p$ by $x^{-1}$; condition (2) bounds
the sum of $\frac{\log p}pB(p)$ over the candidate primes by $|T|x^{-2+\varepsilon}\leq|T|/(120x)$, the
candidate primes being among the primes of that condition; and the weight requirement bounds
$\sum_pw(p)$ by $x^{-9}\leq1/(120x)$. Summing the bound above over the $p\in S$ with $p\leq X''$, and
using Lemma~\ref{pd-gain-holds-fourth-numerics} and $|T|\geq1$, the sum of the
$\log p\cdot\max(0,\Delta(p))$ is at most
\[
 m|T|\frac ax\Big(6\cdot\frac1{40}+\frac32+3+\frac1{30}\Big)\leq\frac{47}{10}\,m|T|\frac ax,
\]
the six terms $\frac1{40}$ coming from $c_1$, $c_{\mathrm{fl}}$, $c_2$, $c_4$, $A_3$ and $B(p)$.

\emph{The primes $p>X''$.} For such $p$, $V(p)=\emptyset$, since $p>X''\geq X'$ is not a candidate
prime; so $g_j(p)=\tau(p)=\iota(p)=0$, and by Lemma~\ref{pd-gain-holds-high-classes}
$\Delta(p)\leq m|T|(\theta_Ww''+S_c)$ if $p<X''(1+\theta_S)$ and $\Delta(p)\leq0$ otherwise. As in part
(6), the sum over these $p\in S$ is at most $(X''\theta_S+2)\log(2X'')\,m|T|(\theta_Ww''+S_c)\leq
m|T|a/(40x)$. Since $\frac{47}{10}+\frac1{40}\leq5$, this gives the first assertion of (4).

\emph{The index excess.} By (iv), for every finite set $S$ of primes,
$\sum_{p\in S}\log p\cdot\iota(p)\leq m|T|a\sum_pw(p)\leq m|T|a\,x^{-9}$, by the weight requirement.
\end{proof}
% Lean: stmtF = stmtF_of_stmtEF of stmtEF_of with delta(x) = x^(-eps/2): part (3) is rest0 =
% As.asksAll_of (fill_ge, kaJ_bounds, shift_ge, shift_le) with adm12's N(m)+1 <= w''; part (1) is
% remainingCountingStatement_of_restStar (C = C1 + 1, GainBoundConstants) with restStar = St.star_fam,
% gain_bound_all_of_star and As.gainBoundAll_add_one_of_gainBoundAllRequired (lo_succ, fl_succ);
% part (2) is rest2 = second_estimate_and_mean_along_family .1 = slopeEstimate_of_lemmaB_slopeLine
% (kappa_nonneg_and_mean_bound, meanLoc_const_le, Nv_eq); part (5) is restW =
% weight_condition_of_mean_excess of meanEx_of (ZZ_le, markov_ind, sum_Sig_le, sum_box_product_eq_mean,
% box_marg, sum_cand_le). Part (4) is rest3 = G5.latticeChain_succ_of_partial with
% G5.prime_sum_facts_along_family: sum_low (low_le with dfc_le, ask_bounds, sdiff_le, lo_bounds,
% prime_rpow_neg_level_le, fl_bounds, io_le with sv_le; the Sets from adm2, adm3, adm4; sum_log_prB,
% sum_logdiv_prB, card_prB_le, WReq_iff) and sum_high (dfc_high, ask_high); part (6) is restT =
% G5.third_estimate_along_family (third_estimate_of_numeric_facts, ask_high, sum_le_count, fl_bounds).

\subsection{The gain link}\label{subsection-pd-gain-link}

\begin{lemma}\label{pd-gain-link-surplus}
Let $D$ be a definite-side datum with ranks $k,m,t,n$ and $N(j)=k+j(t+1)$, size data with scale
$(x,\varepsilon,J)$, exponents $a_t,b_t$, width $\eta$, cloud size $S_c$ and derived quantities
$a,b',w'',X'$ (Definition~\ref{pd-size-data}), and gain data with constants $K_G$, offsets
$o_j(z)$, served representatives $V(p)$, numbers $N(p,t_0)$ and surpluses $S(p,t_0)$
(Definitions~\ref{pd-gain-data} and~\ref{pd-served-lattice}). Suppose that the constants of $K_G$
are nonnegative and that $D$ has parameter choices for its size data
(Definition~\ref{pd-family}(4)). Then $S(p,t_0)\leq2m\,N(p,t_0)\,a/p$ for every prime $p$ and every
$t_0\in V(p)$.
\end{lemma}
\leanref{Formalism/SharpRestriction/GrowthBounds/PlaneCountTheoremWithConstants.lean}{Formalism.MainNotions.G.SvBound}
\leanref{Formalism/SharpRestriction/GrowthBounds/PlaneCountTheoremWithConstants.lean}{Formalism.MainNotions.G.svBound_of}
\leanref{Formalism/SharpRestriction/GrowthBounds/PlaneCountTheoremWithConstants.lean}{Formalism.MainNotions.G.fill_nonneg}
\leanref{Formalism/SharpRestriction/GrowthBounds/PlaneCountTheoremWithConstants.lean}{Formalism.MainNotions.G.kaJ_le}

\begin{proof}
The proof uses the hypotheses through the following facts. The rank datum has $m\geq2$ and
$k\geq m+3\geq5$, so $N(m)\geq k\geq5$. The size data have $x>1$, $\eta>0$, $b\geq1$ and
$S_c=4k\Delta_S$. Of the parameter choices only two bounds are used: $N(m)m^{10}\leq\Delta_S$ and
$kS_c\leq w''$. The first gives $5\leq N(m)\leq\Delta_S$, hence $N(m)\leq S_c$ and
$S_c\geq4\cdot5\cdot5=100$. By Definition~\ref{pd-gain-data}(2), $o_j(z)\geq b'/x$ at the top
places, $o_j(z)\geq b\geq1$ at the thick places and $o_j(z)\geq w''\geq0$ at the others, so
$o_j(z)\geq0$ at every place; the upper bounds on the offsets are not used.

Use the notation of Definitions~\ref{pd-lattice-instance}, \ref{pd-served-selection}
and~\ref{pd-served-lattice}, with $c_G=4S_c$. Let $t_0\in V(p)$. Then $p$ is a candidate prime, so
$p<X'$ and $q=a/p$ satisfies $q>a/X'=b'$. Since $w''=b'/(xn)\leq b'/x\leq b'$, we get
$q>b'\geq w''\geq kS_c\geq500$. As $X'=\exp(x^{\varepsilon/3})>2$, $2b'\leq a$, and so
$2c_G=8S_c\leq2kS_c\leq2b'\leq a$. For $1\leq j\leq m$, $\kappa_j\leq N(m)\leq S_c$.

For a top class $s$ with representative $t_0$, $a\leq a_s\leq a+a/(kx^{13})$ gives
$0\leq a_s-2c_G\leq6a/5$. For every place $z$, $\mathrm{fill}_j(z)+c_G\geq o_j(z)-\kappa_j+c_G\geq
3S_c\geq0$, as $b_t\geq b'\geq0$. So each minimum in the definition of $g_U(p,j,t_0)$ is
nonnegative, and $\delta(p,t_0)\geq0$ because $K_{\delta,1},K_{\delta,2}\geq0$ and $\log_p(2a)\geq0$.
Hence, as $\gamma_p=p/(p^2-1)\leq4/(3p)$ for $p\geq2$,
\[
 g_U(p,j,t_0)\leq\gamma_p\sum_{s}(a_s-2c_G)\leq\frac4{3p}\cdot N(p,t_0)\cdot\frac{6a}5
 =\frac85N(p,t_0)\,q .
\]
Next, $\mathrm{Pr}(p,j,t_0)\geq0$. If $p<p_0$, every term of its definition is nonnegative. If
$p\geq p_0$, its summands are $\mathrm{fill}_j((0,s))-b_s=o_j((0,s))-\kappa_j\geq b'/x-S_c\geq
w''-S_c\geq0$, and $K_{\mathrm{Pr}}\ell_p(w''+\eta+c_G+N(m))\geq0$. Since also $g_0^+\geq0$, each
term $\max(0,g_U-\mathrm{Pr}-g_0^+)$ is at most $\frac85N(p,t_0)q$, and their sum over the $m$ lifts
is at most $\frac85mN(p,t_0)q$. If $N(p,t_0)=0$, then $S(p,t_0)=0$. Otherwise
$\frac25mN(p,t_0)q\geq\frac25\cdot2\cdot500>1$, so
$S(p,t_0)<\frac85mN(p,t_0)q+1\leq2mN(p,t_0)a/p$.
\end{proof}

\begin{lemma}\label{pd-gain-link-admissible-places}
Let $D$ be a definite-side datum whose Galois datum has index $o$ and whose host has splitting bound
$B$ and level primes $\mathcal L$ (Definitions~\ref{hf-frame} and~\ref{hf-host}), with size data
with scale $(x,\varepsilon,J)$ and gain data, with served representatives $V(p)$
(Definition~\ref{pd-gain-data}). Suppose that $o<x^{\varepsilon/4}$, that $B\geq4$, and that
$\left(\frac{-d}{\ell}\right)=1$ for every $\ell\in\mathcal L$. Let $u$ be a word and $p$ a prime with
$V(p)\neq\emptyset$ and $p\nmid T_u$. Then every place $w$ of $L$ over $p$ is admissible for $u$
(Definition~\ref{pd-lattice-instance}).
\end{lemma}
\leanref{Formalism/SharpRestriction/GrowthBounds/PlaneCountTheoremWithConstants.lean}{Formalism.MainNotions.G.ServedGIdx}
\leanref{Formalism/SharpRestriction/GrowthBounds/PlaneCountTheoremWithConstants.lean}{Formalism.MainNotions.G.servedGIdx_of}
\leanref{Development/AutomorphicArithmetic/OrbitLabellings/EvenAddOrderOfHPOfPowLe.lean}{UpperTheory.even_addOrderOf_hP_of_pow_le}

\begin{proof}
Since $V(p)\neq\emptyset$, $p$ is a candidate prime: it is of inert class and the $A$-component $h_p$
of its Frobenius class has order $n_p\geq x^{\varepsilon/4}$. An inert prime is different from $d$
and $-d$ is not a square modulo it. Every integer is a square modulo $2$, so $p$ is odd and
$\left(\frac{-d}p\right)=-1$; by condition (4) of Definition~\ref{hf-host-field-rows}, $p>B\geq4$, so
$p\geq5$. Suppose that $h_p$ has odd order. Its image under the homomorphism $h$ of the Galois datum
lies in a group of order $2^e$, so it is $0$. The element $oh_p$ lies in $A_0$, since $A_0$ has
index $o$, and $h(oh_p)=0$; as $h$ is injective on $A_0$, $oh_p=0$. Then
$n_p\leq o<x^{\varepsilon/4}$, a contradiction; so $h_p$ has even order. Every level prime $\ell$
has $\left(\frac{-d}\ell\right)=1$, so it is not inert, and $p$ is not a level prime. Finally
$p\nmid T_u$.
\end{proof}

\begin{lemma}\label{pd-gain-link-descent}
Let $L$ be a number field, $w$ a finite place of $L$ with valuation ring $\mathcal O_{L,w}$ and
$e_w=v_w(p_w)$, $V$ a finite-dimensional $L$-vector space and $\Lambda\subset V$ a finitely
generated $\mathcal O_L$-submodule. Let $\mathcal R$ be a discrete valuation ring with fraction field
$\mathcal K$, $\epsilon\colon L\to\mathcal K$ an embedding with $v_w(y)\geq0$ if and only if
$\epsilon(y)\in\mathcal R$, $F$ a $\mathcal K$-vector space and $\operatorname{cmp}\colon V\to F$ an
$\epsilon$-semilinear additive map that carries a basis of $V$ to vectors linearly independent over
$\mathcal K$.
\begin{enumerate}
\item If $x\in V$ and $\operatorname{cmp}(x)$ lies in the $\mathcal R$-span of
  $\operatorname{cmp}(\Lambda)$, then $\sigma x\in\Lambda$ for some
  $\sigma\in\mathcal O_L\smallsetminus\mathfrak p_w$.
\item Let $A\subset B$ be finitely generated $\mathcal O_{L,w}$-submodules of $V$ such that $B/A$
  has finite length $l$. Then the quotient of the $\mathcal R$-span of $\operatorname{cmp}(B)$ by
  the $\mathcal R$-span of $\operatorname{cmp}(A)$ has length $l\,v_{\mathcal R}(p_w)/e_w$ over
  $\mathcal R$.
\end{enumerate}
\end{lemma}
\leanref{Formalism/SharpRestriction/GrowthBounds/PlaneCountTheoremWithConstants.lean}{Formalism.MainNotions.G.descent}
\leanref{Formalism/LocalRingCardinality/Lemmas.lean}{Formalism.BasicNotions.LatCmp.length_quotient_of_nonempty}

\begin{proof}
The map $\epsilon$ carries $\mathcal O_{L,w}$ into $\mathcal R$, and a uniformizer $\varpi$ of
$\mathcal O_{L,w}$ to an element of positive valuation $e$, since $\epsilon(\varpi^{-1})\notin\mathcal
R$; as $p_w$ is a unit times $\varpi^{e_w}$, $v_{\mathcal R}(p_w)=ee_w$. Every basis of $V$ is carried
by $\operatorname{cmp}$ to linearly independent vectors, since it differs from the given one by an
invertible matrix over $L$, whose image under $\epsilon$ is invertible.

(1) The localization $\Lambda_w=\mathcal O_{L,w}\Lambda$ is finitely generated and torsion free over
the discrete valuation ring $\mathcal O_{L,w}$, so it has a basis $\mu_1,\dots,\mu_r$; extend it to
a basis $\mu_1,\dots,\mu_s$ of $V$. The $\mathcal R$-span of $\operatorname{cmp}(\Lambda)$ is that of
$\operatorname{cmp}(\mu_1),\dots,\operatorname{cmp}(\mu_r)$. Write $x=\sum_ic_i\mu_i$ with
$c_i\in L$. Then $\operatorname{cmp}(x)=\sum_i\epsilon(c_i)\operatorname{cmp}(\mu_i)$, and linear
independence gives $c_i=0$ for $i>r$ and $\epsilon(c_i)\in\mathcal R$, that is $v_w(c_i)\geq0$, for
$i\leq r$. So $x\in\Lambda_w$, and clearing the finitely many denominators prime to $\mathfrak p_w$
gives $\sigma$.

(2) The ring $\mathcal R$ is torsion free, hence flat, over $\mathcal O_{L,w}$. The module $B$ is free
with a basis linearly independent over $L$, so $r\otimes b\mapsto r\operatorname{cmp}(b)$ is an
isomorphism from $\mathcal R\otimes B$ onto the $\mathcal R$-span of $\operatorname{cmp}(B)$, and the
same holds for $A$. By flatness the quotient is $\mathcal R\otimes(B/A)$. By the elementary divisor
theorem $B/A$ is a sum of modules $\mathcal O_{L,w}/\varpi^{a_i}$ with $\sum_ia_i=l$, and
$\mathcal R\otimes\mathcal O_{L,w}/\varpi^{a}=\mathcal R/\epsilon(\varpi)^a\mathcal R$ has length $ae$.
So the quotient has length $le=l\,v_{\mathcal R}(p_w)/e_w$.
\end{proof}

\begin{lemma}\label{pd-gain-link-served-lattices}
Let $D$ be a definite-side datum with coefficient field $L$, size data and gain data, and a lattice
instance with a comparison $\operatorname{cmp}$ (Definitions~\ref{pd-size-data},
\ref{pd-gain-data} and~\ref{pd-lattice-instance}), and use the notation of
Definitions~\ref{pd-period-data}, \ref{pd-gain-quantities} and~\ref{pd-gain-admissibility}. Let $u$
be a word and $w$ a place admissible for $u$ with $V(p)\neq\emptyset$, where $p=p_w$, and let
$\mathcal R$ and $M$ be the ring and the integral module of $G(u,w)$, $e_w=v_w(p)$ and
$M_{\mathrm{ram}}=M\cap V_{\mathrm{ram}}$. Suppose that:
\begin{enumerate}
\item[(a)] $V_{\mathrm{ram}}$ is the $\mathcal K$-span of $\operatorname{cmp}_{u,w}(M_{u,\mathrm{ram}})$,
  and $N_P(u,w)$ lies in the $\mathcal R$-span of the $\operatorname{cmp}_{u,w}(y)$ with
  $y\in M_{u,\mathrm{ram}}$ and $\operatorname{cmp}_{u,w}(y)\in N_P(u,w)$;
\item[(b)] a vector $x\in V(N_u)$ lies in $N^\circ_u$ locally at $w$, that is $\sigma x\in N^\circ_u$
  for some $\sigma\in\mathcal O_L\smallsetminus\mathfrak p_w$, if and only if
  $\operatorname{cmp}_{u,w}(\nu_u(x))\in N_P(u,w)$;
\item[(c)] $N^\circ_u\subset M^{\mathrm{ram}}_u$ locally at $w$, and the length over
  $\mathcal O_{L,w}$ of the quotient of the localizations at $w$ of $M^{\mathrm{ram}}_u$ and
  $N^\circ_u$ is $e_w\operatorname{len}_{u,w}$.
\end{enumerate}
Then $N^\circ_u\subset\mathfrak p_w^{r}M^{\mathrm{ram}}_u$ locally at $w$ with
$r=\max(0,\lfloor e_w\tau(p)\rfloor)$, and the $\mathcal R$-module $M_{\mathrm{ram}}/N_P(u,w)$ has
length $v_{\mathcal R}(p)\operatorname{len}_{u,w}$.
\end{lemma}
\leanref{Formalism/SharpRestriction/GrowthBounds/PlaneCountTheoremWithConstants.lean}{Formalism.MainNotions.G.NuServed}
\leanref{Formalism/SharpRestriction/GrowthBounds/PlaneCountTheoremWithConstants.lean}{Formalism.MainNotions.G.nuServed_of}
\leanref{Formalism/SharpRestriction/GrowthBounds/IsResidueClassZero.lean}{Formalism.MainNotions.G.LatticeLengthsAtPlace}
\leanref{Formalism/LocalRingCardinality/Lemmas.lean}{Formalism.BasicNotions.LatCmp.sizeDataLat_of_latDefs}
\leanref{Formalism/LocalRingCardinality/Lemmas.lean}{Formalism.BasicNotions.LatCmp.NP_eq_span}
\leanref{Formalism/LocalRingCardinality/Lemmas.lean}{Formalism.BasicNotions.LatCmp.M_inf_Vram}

\begin{proof}
Put $T=\max(\lceil\tau(p)\rceil,0)$, so that $r\leq e_wT$. Recall that $M$ is the $\mathcal R$-span
of $\operatorname{cmp}_{u,w}(\Lambda_{u,0})$ and $M^{\mathrm{ram}}_u=\nu_u^{-1}(\Lambda_{u,0})$. Let
$y\in N^\circ_u$. By (b), $\operatorname{cmp}_{u,w}(\nu_u(y))\in N_P(u,w)\subset p^TM$, so
$\operatorname{cmp}_{u,w}(\nu_u(p^{-T}y))=p^{-T}\operatorname{cmp}_{u,w}(\nu_u(y))$ lies in $M$.
Lemma~\ref{pd-gain-link-descent}(1) gives $\sigma\notin\mathfrak p_w$ with
$\sigma p^{-T}\nu_u(y)\in\Lambda_{u,0}$, that is $\sigma p^{-T}y\in M^{\mathrm{ram}}_u$. So $\sigma
y\in p^TM^{\mathrm{ram}}_u\subset\mathfrak p_w^{r}M^{\mathrm{ram}}_u$.

For the length, recall that $\nu_u(V(N_u))$ is the space of vectors of $V_{u,0}$ whose image lies
in $M_{u,\mathrm{ram}}$, as in (a). Let $\Lambda_w$ be the localization of $\Lambda_{u,0}$ at $w$ and
$\Lambda'_w=\Lambda_w\cap\nu_u(V(N_u))$, the localization of $\nu_u(M^{\mathrm{ram}}_u)$. It is a
saturated submodule of the free $\mathcal O_{L,w}$-module $\Lambda_w$, so $\Lambda_w=\Lambda'_w\oplus
\Lambda''_w$ for a submodule $\Lambda''_w$. An element of $M$ is a sum of an element of the
$\mathcal R$-span of $\operatorname{cmp}_{u,w}(\Lambda'_w)$ and one of the $\mathcal R$-span of
$\operatorname{cmp}_{u,w}(\Lambda''_w)$; if it lies in $V_{\mathrm{ram}}$, the second vanishes by
(a) and the linear independence in Lemma~\ref{pd-gain-link-descent}. So $M_{\mathrm{ram}}$ is the
$\mathcal R$-span of $\operatorname{cmp}_{u,w}(\nu_u(M^{\mathrm{ram}}_u))$. By (a) and (b), $N_P(u,w)$
is the $\mathcal R$-span of the $\operatorname{cmp}_{u,w}(\nu_u(y))$ with $y$ in the localization of
$N^\circ_u$ at $w$. By (c) and Lemma~\ref{pd-gain-link-descent}(2), applied to the images under
$\nu_u$ of the localizations of $N^\circ_u$ and $M^{\mathrm{ram}}_u$, the length of
$M_{\mathrm{ram}}/N_P(u,w)$ is $e_w\operatorname{len}_{u,w}\,v_{\mathcal R}(p)/e_w=
v_{\mathcal R}(p)\operatorname{len}_{u,w}$.
\end{proof}

\begin{lemma}\label{pd-gain-holds-transfer}
Let $\mathcal R$ be a discrete valuation ring, $\pi\in\mathcal R$, $X$ an $\mathcal R$-module and
$g,E,e\geq0$ integers. For submodules $A,N\subset X$, say that $A$ is \emph{covered by $N$ up to $g$
elements at exponent $E$} if there is a finite set $C\subset A$ with at most $g$ elements such that
$\pi^EA\subset N+\mathcal RC$.
\begin{enumerate}
\item Let $P\subset A'\subset A$ be submodules. If $A$ is covered by $P$ up to $g$ elements at
  exponent $E$, then so is $A'$.
\item Let $B,N_A\subset A$ and $N_B\subset B$ be submodules with $\pi^eN_A\subset N_B$. If $A$ is
  covered by $N_A$ up to $g$ elements at exponent $E$, then $B$ is covered by $N_B$ up to $g$
  elements at exponent $E+e$.
\end{enumerate}
\end{lemma}
\leanref{Formalism/SubspaceCounting/Statement.lean}{Formalism.G0.AllButLE}
\leanref{Formalism/ThetaLogCoordinates/ProductIdealMonotone/ProductIdealProofs/ElementarySymmetricExpBound.lean}{Formalism.InvariantFactorPaths.G0.invariant_factors_submodule}
\leanref{Formalism/ThetaLogCoordinates/ProductIdealMonotone/ProductIdealProofs/ElementarySymmetricExpBound.lean}{Formalism.InvariantFactorPaths.G0.inverse_shift}

\begin{proof}
(1) Let $C\subset A$ be a set as in the definition. Let $f\colon\mathcal R^C\to X/P$ send $(r_c)$ to
the class of $\sum_cr_cc$, and let $U\subset\mathcal R^C$ be the inverse image under $f$ of the image
of $\pi^EA'$. Since $\mathcal R$ is a principal ideal domain, $U$ is free with a basis
$u_1,\dots,u_n$, $n\leq|C|$. Choose $x_i\in A'$ such that the class of $\pi^Ex_i$ is $f(u_i)$. For
$x\in A'$ write $\pi^Ex=y+\sum_cr_cc$ with $y\in P$. Then $(r_c)\in U$, so $(r_c)=\sum_is_iu_i$ with
$s_i\in\mathcal R$, and $\pi^Ex-\sum_is_i\pi^Ex_i$ lies in $P$. So the set
$\{x_1,\dots,x_n\}\subset A'$ has the required property.

(2) The submodules $N_A\subset B+N_A\subset A$ satisfy the hypotheses of (1), which gives a finite
set $C'\subset B+N_A$ with at most $g$ elements and $\pi^E(B+N_A)\subset N_A+\mathcal RC'$. The set
$\pi^eC'$ lies in $B$: for $b\in B$ and $y\in N_A$, $\pi^e(b+y)=\pi^eb+\pi^ey$ with $\pi^ey\in
N_B\subset B$. For $x\in B$ write $\pi^Ex=y+z$ with $y\in N_A$ and $z\in\mathcal RC'$; then
$\pi^{E+e}x=\pi^ey+\pi^ez$ lies in $N_B+\mathcal R\pi^eC'$.
\end{proof}

\begin{lemma}\label{pd-gain-holds-length}
Let $\mathcal R$ be a discrete valuation ring, $a\in\mathcal R$ nonzero, $\ell$ the length of
$\mathcal R/a\mathcal R$, $X$ an $\mathcal R$-module, $M\subset X$ a submodule that is free of finite
rank $d$, $N\subset X$ a submodule, $\tau,S,V\geq0$ integers and $0\leq s\leq1/2$. Suppose that for
every integer $0\leq\beta<V$ there is a finite set $C_\beta\subset M$ with at most $ds^{\beta+1}$
elements such that $a^{\tau+\beta S}M\subset N+\mathcal RC_\beta$, and that $a^{VS}M\subset N$. Then
the $\mathcal R$-module $M/(N\cap a^\tau M)$ has length at most $\ell\,d\,(\tau+2Ss)$.
\end{lemma}
\leanref{Formalism/SharpRestriction/GrowthBounds/LatticeSystems.lean}{Formalism.MainNotions.G.LatticeSystems.lengthFromAllBut}
\leanref{Formalism/SubspaceCounting/SubspaceCountingProofs/ManyOrbitsRefined.lean}{Formalism.G0.CyclicSpanLengthLemmas.lengthPart2}

\begin{proof}
Write $a=u\varpi^\ell$ with $\varpi$ a uniformizer and $u$ a unit. Then $a^iY=\varpi^{\ell i}Y$ for
every submodule $Y$ and every $i\geq0$, so the hypotheses hold with $\varpi$, $\ell\tau$ and $\ell S$
in place of $a$, $\tau$ and $S$, and it suffices to bound the length by $d(\ell\tau+2\ell Ss)$. So
assume $a=\varpi$ and $\ell=1$. Put $K=N\cap\varpi^\tau M$ and $Q=M/K$. The module $Q/\varpi^\tau Q$
is generated by $d$ elements killed by $\varpi^\tau$, so its length is at most $\tau d$. For $j\geq0$
let $E_j$ be the image of $\varpi^{\tau+jS}M$ in $X/N$. Sending the class of $\varpi^\tau x$ in $Q$
to the class of $\varpi^\tau x$ in $X/N$ is a well defined injective map $\varpi^\tau Q\to E_0$,
because $\varpi^\tau x\in K$ if and only if $\varpi^\tau x\in N$. The $E_j$ decrease,
$\varpi^SE_j\subset E_{j+1}$, and $E_V=0$ because $\varpi^{\tau+VS}M\subset\varpi^{VS}M\subset N$.
For $j<V$, $E_j$ lies in the submodule generated by the images of $C_j$; over the principal ideal
domain $\mathcal R$ it is then generated by $|C_j|$ elements, and $E_j/E_{j+1}$, a quotient of
$E_j/\varpi^SE_j$, has length at most $S|C_j|$. Hence
\[
 \operatorname{length}Q\leq\tau d+\sum_{j<V}S|C_j|\leq\tau d+Sd\sum_{j<V}s^{j+1}\leq\tau d+2Sds,
\]
since $\sum_{j\geq0}s^{j+1}=s/(1-s)\leq2s$ for $0\leq s\leq1/2$.
\end{proof}

\begin{lemma}\label{pd-gain-holds-esymm}
Let $\kappa_1,\dots,\kappa_n\geq0$ be real numbers, $j\geq0$ an integer and $e_j$ the $j$-th
elementary symmetric function. Then $j!\,e_j(\kappa_1,\dots,\kappa_n)\leq(\kappa_1+\cdots+\kappa_n)^j$.
\end{lemma}
\leanref{Formalism/SubspaceCounting/SubspaceCountingProofs/ManyOrbits.lean}{Formalism.G0.ManyOrbitsP.fact_esymm_le}
\leanref{Formalism/SharpRestriction/GrowthBounds/LatticeSystems.lean}{Formalism.MainNotions.G.LatticeSystems.esymm_clause}

\begin{proof}
By induction on $n$; for $n=0$ both sides are $1$ if $j=0$, and the left side is $0$ otherwise. Put
$\sigma=\kappa_1+\cdots+\kappa_{n-1}$ and $\kappa=\kappa_n$. For $j\geq1$,
$e_j(\kappa_1,\dots,\kappa_n)=e_j(\kappa_1,\dots,\kappa_{n-1})+\kappa\,e_{j-1}(\kappa_1,\dots,\kappa_{n-1})$
and $(\sigma+\kappa)^j\geq\sigma^j+j\kappa\sigma^{j-1}$. By induction,
$j!\,e_j(\kappa_1,\dots,\kappa_n)\leq\sigma^j+j\kappa\sigma^{j-1}\leq(\sigma+\kappa)^j$.
\end{proof}

\begin{lemma}\label{pd-gain-holds-walk}
Let $\mathcal O$ be a commutative ring, $X$ an $\mathcal O$-module, $\pi\in\mathcal O$ and $V_0$ a
finite totally ordered set. A \emph{labelling} is a map $\sigma$ from $V_0$ to the set of the three
labels $\mathrm{base}$, $\mathrm{src}$ and $\mathrm{tgt}$. Write $\sigma[v\mapsto l]$ for $\sigma$
with its value at $v$ replaced by $l$, and for $B\subset V_0$ let $\sigma_B$ be the labelling with
value $\mathrm{tgt}$ on $B$ and $\mathrm{base}$ elsewhere. Let $X(\sigma)\subset X$ be a submodule for
each labelling $\sigma$, and put $X_0=X(\sigma_\emptyset)$. Let $y_v$, $v\in V_0$, be commuting
$\mathcal O$-linear endomorphisms of $X$, let $m\geq1$ and $\beta,c_0,S\geq0$ be integers, let
$a_v,b_v,c_v,S_v\geq0$ be integers with $S_v\leq S$ for $v\in V_0$, and for each $B\subset V_0$ with
$|B|=\beta+1$ let $G_B\subset X(\sigma_B)$ be a finite set. Suppose that for every labelling $\sigma$
and every $v\in V_0$:
\begin{enumerate}
\item $X(\sigma[v\mapsto\mathrm{tgt}])\subset X(\sigma[v\mapsto\mathrm{base}])\subset
  X(\sigma[v\mapsto\mathrm{src}])$;
\item $\pi^{a_v}X(\sigma[v\mapsto\mathrm{src}])\subset X(\sigma[v\mapsto\mathrm{base}])$,
  $\pi^{b_v}X(\sigma[v\mapsto\mathrm{base}])\subset X(\sigma[v\mapsto\mathrm{tgt}])$ and
  $\pi^{c_v}X(\sigma[v\mapsto\mathrm{src}])\subset X(\sigma[v\mapsto\mathrm{tgt}])$;
\item $y_vX(\sigma[v\mapsto\mathrm{src}])\subset X(\sigma[v\mapsto\mathrm{tgt}])$;
\end{enumerate}
and that for every $B\subset V_0$ with $|B|=\beta+1$:
\begin{enumerate}
\setcounter{enumi}{3}
\item $\pi^{c_0}X(\sigma_B)\subset\sum_{v\in B}y_vX(\sigma_B[v\mapsto\mathrm{src}])+\mathcal OG_B$.
\end{enumerate}
For $e\subset V_0$ let $y^{(e)}$ be $\pi^{\sum_{v\in V_0\smallsetminus e}S_v}$ times the composite of
the $y_v^m$, $v\in e$, and let $\Pi$ be the sum of the submodules $y^{(e)}X_0$. Then
$\pi^{\sum_vS_v}X_0\subset\Pi$, and there is a finite set $C\subset X_0$ with at most
$m^{\beta+1}\sum_{|B|=\beta+1}|G_B|$ elements such that $\pi^EX_0\subset\Pi+\mathcal OC$, where
$E=\sum_{v\in V_0}\big(a_v+b_v+(m-1)c_v+mc_0\big)+\beta S$.
\end{lemma}
\leanref{Formalism/SubspaceCounting/Statement.lean}{Formalism.G0.WalkSharp}
\leanref{Formalism/SubspaceCounting/SubspaceCountingProofs/WalkSharp.lean}{Formalism.G0.walkSharp_true}
\leanref{Formalism/SubspaceCounting/SubspaceCountingProofs/WalkSharp.lean}{Formalism.G0.WalkSharpP.main}

\begin{proof}
Each $y_v$ preserves $X_0$: by (1) and (3),
$y_vX_0\subset y_vX(\sigma_\emptyset[v\mapsto\mathrm{src}])\subset
X(\sigma_\emptyset[v\mapsto\mathrm{tgt}])\subset X_0$. Taking $e=\emptyset$ gives
$\pi^{\sum_vS_v}X_0\subset\Pi$. If $|V_0|<\beta+1$, then $\sum_vS_v\leq|V_0|S\leq\beta S\leq E$, and
$C=\emptyset$ will do. Assume $|V_0|\geq\beta+1$, and identify $V_0$ with $\{0,\dots,V-1\}$.

For $B\subset V_0$ with $|B|=\beta+1$ let $\operatorname{top}(B)$ be its largest element, and let
$I(B)$ be the set of the maps $i\colon V_0\to\mathbf Z_{\geq0}$ with $i_u<m$ for $u\in B$ and $i_u=0$
for $u\notin B$; it has $m^{\beta+1}$ elements. For $i\in I(B)$ define $\epsilon(B,i)\colon
V_0\to\mathbf Z_{\geq0}$ and an integer $\alpha(B,i)\geq0$ by
\[
 \epsilon(B,i)_u=\begin{cases}m,&u<\operatorname{top}(B),\ u\notin B,\\ i_u,&u\in B,\\
 0,&\text{otherwise,}\end{cases}
\]
\[
 \alpha(B,i)=\sum_{u<\operatorname{top}(B),\,u\notin B}\big(a_u+b_u+(m-1)c_u+mc_0\big)
 +\sum_{u\in B}\big(b_u+i_u(c_u+c_0)\big),
\]
and write $y^\epsilon$ for the composite of the $y_u^{\epsilon_u}$. Let $C$ be the set of the
$y^{\epsilon(B,i)}g$ with $|B|=\beta+1$, $i\in I(B)$ and $g\in G_B$. It has at most
$m^{\beta+1}\sum_B|G_B|$ elements, and it lies in $X_0$, because $G_B\subset X(\sigma_B)\subset X_0$
by (1) and the $y_u$ preserve $X_0$. Compare $\alpha(B,i)$ term by term with
$E-\beta S=\sum_u(a_u+b_u+(m-1)c_u+mc_0)$, using $i_u\leq m-1$ for $u\in B$: this gives
$\alpha(B,i)+c_0\leq E-\beta S$, and $\alpha(B,i)+a_v+c_0\leq E-\beta S$ if $v\in B$ and $i_v=m-1$.

We claim that for $|B|=\beta+1$, $i\in I(B)$ and $z\in X(\sigma_B)$,
\[
 \pi^{E-\alpha(B,i)}y^{\epsilon(B,i)}z\in\Pi+\mathcal OC.
\]
Since $\sum_u\epsilon(B,i)_u\leq mV$, we may prove this by descending induction on
$\sum_u\epsilon(B,i)_u$. Put $\epsilon=\epsilon(B,i)$, $Y=y^\epsilon$ and $k=E-\alpha(B,i)-c_0\geq0$.
By (4), $\pi^{c_0}z=\sum_{v\in B}y_vz_v+g$ with $z_v\in X(\sigma_B[v\mapsto\mathrm{src}])$ and
$g\in\mathcal OG_B$, so
\[
 \pi^{E-\alpha(B,i)}Yz=\sum_{v\in B}\pi^kYy_vz_v+\pi^kYg,
\]
and $Yg\in\mathcal OC$. Fix $v\in B$, and let $e_v$ be the indicator function of $v$. Since the $y_u$
commute, $Yy_v=y^{\epsilon+e_v}$. We show $\pi^kYy_vz_v\in\Pi+\mathcal OC$ in three cases.
\begin{itemize}
\item If $i_v+1<m$, put $i'=i+e_v\in I(B)$. Then $\epsilon(B,i')=\epsilon+e_v$ and
  $\alpha(B,i')=\alpha(B,i)+c_v+c_0$. By (2), $\pi^{c_v}z_v\in X(\sigma_B[v\mapsto\mathrm{tgt}])=
  X(\sigma_B)$, and by induction $\pi^{E-\alpha(B,i')}y^{\epsilon+e_v}\pi^{c_v}z_v\in\Pi+\mathcal OC$.
  This element is $\pi^kYy_vz_v$.
\item If $i_v=m-1$ and $w=\operatorname{top}(B)+1$ lies in $V_0$, put $B'=(B\smallsetminus\{v\})\cup
  \{w\}$, which has $\beta+1$ elements, and let $i'\in I(B')$ agree with $i$ on $B\smallsetminus\{v\}$
  and vanish at $v$ and $w$. Then $\operatorname{top}(B')=w$, and comparing the cases of the
  definitions gives $\epsilon(B',i')=\epsilon+e_v$ and $\alpha(B',i')=\alpha(B,i)+a_v+c_0+b_w$. By
  the first inclusion of (2), $\pi^{a_v}z_v\in X(\sigma_{B\smallsetminus\{v\}})$, and by the second,
  since $w\notin B$, $\pi^{b_w+a_v}z_v\in X(\sigma_{B'})$. By induction
  $\pi^{E-\alpha(B',i')}y^{\epsilon+e_v}\pi^{b_w+a_v}z_v\in\Pi+\mathcal OC$, and this element is
  $\pi^kYy_vz_v$.
\item If $i_v=m-1$ and $\operatorname{top}(B)=V-1$, then every $u\notin B$ has
  $u<\operatorname{top}(B)$. So, with $B'=B\smallsetminus\{v\}$, $\epsilon+e_v$ takes the value $m$ off
  $B'$ and $i_u$ at $u\in B'$. Put $x'=\prod_{u\in B'}y_u^{i_u}\,\pi^{a_v}z_v$; it lies in $X_0$, since
  $\pi^{a_v}z_v\in X(\sigma_{B'})\subset X_0$ by (2) and (1). With $e=V_0\smallsetminus B'$,
  $y^{(e)}x'=\pi^{\sum_{u\in B'}S_u+a_v}Yy_vz_v$ lies in $\Pi$. As $\sum_{u\in B'}S_u\leq\beta S$ and
  $\alpha(B,i)+a_v+c_0\leq E-\beta S$, we have $k\geq a_v+\sum_{u\in B'}S_u$, so
  $\pi^kYy_vz_v\in\Pi$.
\end{itemize}
In the first two cases the sum of the exponents has grown by $1$, which justifies the induction.

Finally let $B_0=\{0,\dots,\beta\}$. Then $\epsilon(B_0,0)=0$, as every $u<\beta$ lies in $B_0$, and
$\alpha(B_0,0)=\sum_{u\in B_0}b_u$. For $x\in X_0$, applying the second inclusion of (2) at the
elements of $B_0$ one at a time gives $\pi^{\sum_{u\in B_0}b_u}x\in X(\sigma_{B_0})$, and the claim
gives $\pi^Ex=\pi^{E-\alpha(B_0,0)}\pi^{\sum_{u\in B_0}b_u}x\in\Pi+\mathcal OC$.
\end{proof}

\begin{lemma}\label{pd-gain-holds-count}
Let $\mathcal R$ be a discrete valuation ring, $q\in\mathcal R$ nonzero, $n,N\geq0$ integers,
$\alpha_1,\dots,\alpha_N\in\mathcal R$, and let $(c_{ji})$ and $(Z_{ik})$ be $n\times N$ and $N\times n$
matrices over $\mathcal R$ with $\sum_ic_{ji}\alpha_iZ_{ik}=q\,\delta_{jk}$ for all $j,k\leq n$. Then at
least $n$ of the $\alpha_i$ divide $q$.
\end{lemma}
\leanref{Formalism/SharpRestriction/GrowthBounds/LatticeSystems.lean}{Formalism.MainNotions.G.LatticeSystems.jm_count}

\begin{proof}
Let $\mathfrak m$ be the maximal ideal of $\mathcal R$ and $I$ the set of the $i$ with $\alpha_i\mid q$.
For $i\in I$ put $\alpha'_i=\alpha_i$ and $\mu_i=0$. For $i\notin I$, $q$ divides $\alpha_i$, as
$\mathcal R$ is a valuation ring; put $\alpha'_i=0$ and $\mu_i=\alpha_i/q$, which lies in $\mathfrak m$
because $\alpha_i\nmid q$. Then $\alpha_i=\alpha'_i+q\mu_i$ for every $i$. The matrix
$Y=c\operatorname{diag}(\mu)Z$ has entries in $\mathfrak m$, and the hypothesis gives
$c\operatorname{diag}(\alpha')Z=q(1-Y)$; also $\det(1-Y)\equiv1\pmod{\mathfrak m}$ is nonzero. Let
$y\in\mathcal R^n$ with $(Zy)_i=0$ for every $i\in I$. Since $\alpha'_i=0$ for $i\notin I$,
$q(1-Y)y=c\operatorname{diag}(\alpha')Zy=0$, so $y=0$. Thus $y\mapsto((Zy)_i)_{i\in I}$ is an injective
$\mathcal R$-linear map $\mathcal R^n\to\mathcal R^I$, and $n\leq|I|$.
\end{proof}

\begin{lemma}\label{pd-gain-holds-smith}
Let $\mathcal R$ be a discrete valuation ring with fraction field $\mathcal K$, $V$ a $\mathcal K$-vector
space with a finite basis $(x_i)$, $H$ the $\mathcal R$-span of the $x_i$, $\varpi\in\mathcal R$ nonzero,
and $P\subset V$ an $\mathcal R$-submodule with $\varpi H\subset P\subset H$. Let $N'\subset V^\vee$ be an
$\mathcal R$-submodule with $G(P)\subset\mathcal R$ for every $G\in N'$, let $f_1,\dots,f_K\in P$,
$q\in\mathcal R$ nonzero and $n\geq0$, let $u\colon\mathcal R^n\to N'$ be $\mathcal R$-linear, and let
$(B_{jk})$ be an $n\times K$ matrix over $\mathcal R$ with $\sum_kB_{jk}\,u(y)(f_k)=qy_j$ for all
$y\in\mathcal R^n$ and $j\leq n$. Then $P$ has an $\mathcal R$-basis $b_1,\dots,b_{\dim V}$, and there
is a set $\mathrm{bad}$ of indices with $|\mathrm{bad}|+n\leq\dim V$ such that
$qb_a\in\sum_k\mathcal Rf_k$ for every $a\notin\mathrm{bad}$.
\end{lemma}
\leanref{Formalism/SharpRestriction/GrowthBounds/LatticeSystems.lean}{Formalism.MainNotions.G.LatticeSystems.jm_core}
\leanref{Formalism/SharpRestriction/GrowthBounds/LatticeSystems.lean}{Formalism.MainNotions.G.LatticeSystems.jm_basis}

\begin{proof}
The module $P$ is finitely generated and torsion free, as a submodule of $H$, so it is free; its rank
is $\dim V$, since $H\cong\varpi H\subset P\subset H$. By the elementary divisor theorem for the
submodule $F=\sum_k\mathcal Rf_k$ of $P$, there are a basis $b_1,\dots,b_{\dim V}$ of $P$, an injective
map $\iota$ from $\{1,\dots,N\}$ to the indices and nonzero $\alpha_1,\dots,\alpha_N\in\mathcal R$ such
that the $\alpha_ib_{\iota(i)}$ form a basis of $F$. Write
$\sum_kB_{jk}f_k=\sum_ic_{ji}\alpha_ib_{\iota(i)}$ with $c_{ji}\in\mathcal R$, and put
$Z_{il}=u(e_l)(b_{\iota(i)})$, where $e_l$ is the $l$-th unit vector of $\mathcal R^n$; it lies in
$\mathcal R$, as $b_{\iota(i)}\in P$. Applying $u(e_l)$ to the first identity gives
$\sum_ic_{ji}\alpha_iZ_{il}=\sum_kB_{jk}u(e_l)(f_k)=q\delta_{jl}$. By
Lemma~\ref{pd-gain-holds-count}, the set $I$ of the $i$ with $\alpha_i\mid q$ has at least $n$
elements. Let $\mathrm{bad}$ be the complement of $\iota(I)$, so $|\mathrm{bad}|\leq\dim V-n$. For
$a=\iota(i)$ with $i\in I$, $qb_a=(q/\alpha_i)\,\alpha_ib_a\in F$.
\end{proof}

\begin{lemma}\label{pd-gain-holds-local-chain}
Let $p\geq5$ be a prime, $(W,Q,\sigma,\iota)$ a Hida setting of degree $n'$ at $p$ with rings
$\mathcal R\subset\mathcal K$ and $\varpi_0$, $k\geq1$ and $r,t\geq0$ integers, and
$\lambda_0(s)\in\mathbf Z^k$, $s\in\mathbf Z/n'\mathbf Z$, dominant weights. Let $V$, the lattices
$H(L)$, $H(e)$, $R(e)$ and $R'(e)$, the modules $N_{\mathrm{src}}(L)$ and $N_{\mathrm{tgt}}(L)$ for this
$t$, and $A$ be those of Definition~\ref{pd-projection-property}, and for a $W$-lattice $L\subset Q^k$
let $P_{\mathrm{tgt}}(L)$ be the set of $F\in V$ with $G(F)\in\mathcal R$ for every $G\in
N_{\mathrm{tgt}}(L)$. Let $L\subset Q^k$ be a $W$-lattice.
\begin{enumerate}
\item $H(L)$ has an $\mathcal R$-basis of $\dim V$ elements, which is a $\mathcal K$-basis of $V$. An
  element $F\in V$ with $G(F)\in\mathcal R$ for every $G\in H(L)^\vee$ lies in $H(L)$.
\item $P_{\mathrm{tgt}}(L)\subset H(L)\subset N_{\mathrm{src}}(L)$,
  $\varpi_0^tN_{\mathrm{src}}(L)\subset H(L)$ and $\varpi_0^tH(L)\subset P_{\mathrm{tgt}}(L)$.
\item Let $e$ be an edge with source $L''$ and target $L$. If $F\in R(e)$, then
  $p^AF\in P_{\mathrm{tgt}}(L)$. In particular $p^AN_{\mathrm{src}}(L'')\subset P_{\mathrm{tgt}}(L)$.
\item Let $\gamma\in GL_k(Q)$ act on $\mathcal K[x_{s,a,b}]$ as in
  Definition~\ref{pd-local-model}(3), with one Hida setting. Then $\gamma$ carries $H(L)$,
  $N_{\mathrm{src}}(L)$ and $P_{\mathrm{tgt}}(L)$ into $H(\gamma L)$, $N_{\mathrm{src}}(\gamma L)$ and
  $P_{\mathrm{tgt}}(\gamma L)$.
\end{enumerate}
\end{lemma}
\leanref{Formalism/SharpRestriction/GrowthBounds/LatticeSystems.lean}{Formalism.MainNotions.G.LatticeSystems.RestLocD}
\leanref{Formalism/SharpRestriction/GrowthBounds/LatticeSystems.lean}{Formalism.MainNotions.G.LatticeSystems.restLocD}
\leanref{Formalism/SharpRestriction/GrowthBounds/LatticeSystems.lean}{Formalism.MainNotions.G.LatticeSystems.RestLocChain}
\leanref{Formalism/SharpRestriction/GrowthBounds/LatticeSystems.lean}{Formalism.MainNotions.G.LatticeSystems.ch_locChain}
\leanref{Formalism/SharpRestriction/GrowthBounds/LatticeSystems.lean}{Formalism.MainNotions.G.LatticeSystems.nt_local}
\leanref{Formalism/SharpRestriction/GrowthBounds/LatticeSystems.lean}{Formalism.MainNotions.G.LatticeSystems.jm_Re}
\leanref{Formalism/SharpRestriction/GrowthBounds/LatticeSystems.lean}{Formalism.MainNotions.G.LatticeSystems.jr_loc}
\leanref{Formalism/SharpRestriction/GrowthBounds/LatticeSystems.lean}{Formalism.MainNotions.G.LatticeSystems.la_reflexive}
\leanref{Formalism/SharpRestriction/GrowthBounds/LatticeSystems.lean}{Formalism.MainNotions.G.LatticeSystems.mapsTo_je_loc}

\begin{proof}
(1) Let $g$ be a matrix whose columns form a $W$-basis of $L$, and for each $s$ let $g_s=\iota\sigma^s(g)$,
whose columns form an $\mathcal R$-basis of $M_s$. Let $\nu(s)$ be the natural weight of $\lambda(s)$
and $T_s$ the invertible $\mathcal K$-linear map $f\mapsto\det(g_s)^{-\lambda(s)_k}f(g_s^{-1}x)$. By
Definition~\ref{pd-projection-property}(1) and Lemma~\ref{pd-gain-holds-hecke-rows}(2), which makes the
condition there independent of the basis, $\nabla_{\lambda(s)}(M_s^\vee)=T_s(\nabla_{\mathcal
R}(\nu(s)))$. The module $\nabla_{\mathcal R}(\nu(s))$ has an $\mathcal R$-basis of bideterminants
\cite{DEP80}, which is a $\mathcal K$-basis of $\nabla_{\mathcal K}(\nu(s))$, and $T_s$ preserves
$\nabla_{\mathcal K}(\nu(s))$, since left multiplication of $x$ by a matrix preserves the column degrees
and commutes with $x\mapsto x(1+yE_{lj})$. The products over $s$ of the images under $T_s$ of these
bases span $H(L)$ over $\mathcal R$, and, the variables of the positions being distinct, they form a
$\mathcal K$-basis of $V$. For the second assertion, write $F=\sum_ia_ix_i$ in this basis $(x_i)$; the
dual basis lies in $H(L)^\vee$, so $a_i\in\mathcal R$ and $F\in H(L)$.

(2) By Lemma~\ref{pd-edge-basis}(2), $H(L)\subset R(e)$ for every edge $e$ with source $L$, and
$H(L)^\vee\subset R'(e)$ for every edge $e$ with target $L$. As $t\geq0$, also $H(L)\subset
\varpi_0^{-t}H(L)$ and $H(L)^\vee\subset\varpi_0^{-t}H(L)^\vee$. Hence $H(L)\subset N_{\mathrm{src}}(L)$
and $H(L)^\vee\subset N_{\mathrm{tgt}}(L)$; the second inclusion and (1) give
$P_{\mathrm{tgt}}(L)\subset H(L)$. Since $N_{\mathrm{src}}(L)\subset\varpi_0^{-t}H(L)$,
$\varpi_0^tN_{\mathrm{src}}(L)\subset H(L)$. For $F\in H(L)$ and $G\in N_{\mathrm{tgt}}(L)$, write
$G=\varpi_0^{-t}G'$ with $G'\in H(L)^\vee$; then $G(\varpi_0^tF)=G'(F)\in\mathcal R$.

(3) Every $G\in N_{\mathrm{tgt}}(L)$ lies in $R'(e)$, so $G(p^AF)=\langle F,G\rangle\in\mathcal R$ by
Lemma~\ref{pd-edge-basis}(3). The second assertion follows, since $N_{\mathrm{src}}(L'')\subset R(e)$.

(4) Let $h_s=\iota\sigma^s(\gamma)$. The action of $\gamma$ at the position $s$ is
$f\mapsto\det(h_s)^{-\lambda(s)_k}f(h_s^{-1}x)$, and for a matrix $g'$,
$\det(h_sg')^{\lambda(s)_k}\big(\det(h_s)^{-\lambda(s)_k}f(h_s^{-1}h_sg'x)\big)=
\det(g')^{\lambda(s)_k}f(g'x)$. So $\gamma$ carries the lattice formed from a matrix $g'$ to the one
formed from $h_sg'$, hence $H(M)$ into $H(\gamma M)$ for every $W$-lattice $M$, and $H(e)$ into
$H(\gamma e)$ for every edge $e$, where $\gamma e$ has source and target moved by $\gamma$; $\gamma e$ is
an edge, since $\gamma$ preserves inclusions, multiplication by $p$ and indices. Hence $\gamma$ carries
$\varpi_0^{-t}H(L)$ and $R(e)$ into $\varpi_0^{-t}H(\gamma L)$ and $R(\gamma e)$, and so
$N_{\mathrm{src}}(L)$ into $N_{\mathrm{src}}(\gamma L)$, the edges with source $\gamma L$ being the
$\gamma e$ with $e$ of source $L$. For $G\in V^\vee$ the form $G\circ\gamma$ lies in $H(M)^\vee$ when
$G\in H(\gamma M)^\vee$, and in $R'(e)$ when $G\in R'(\gamma e)$; so $G\circ\gamma\in N_{\mathrm{tgt}}(L)$
for $G\in N_{\mathrm{tgt}}(\gamma L)$. For $F\in P_{\mathrm{tgt}}(L)$ this gives
$G(\gamma F)=(G\circ\gamma)(F)\in\mathcal R$, that is, $\gamma F\in P_{\mathrm{tgt}}(\gamma L)$.
\end{proof}

\begin{definition}\label{pd-gain-holds-systems}
Let $G$ be a local model at a prime $p$ (Definition~\ref{pd-local-model}) with rings
$\mathcal R\subset\mathcal K$, $\varpi_0$, integers $k$ and $r$, function space $\mathrm{Fn}$ and Hecke
operators $y_v$. For an orbit $v$ and an integer $t\geq0$, let $V_v$, the lattices $H_v(L)$ and the
modules $N^t_{\mathrm{src},v}(L)$ and $P^t_{\mathrm{tgt},v}(L)$, for $W_v$-lattices $L\subset Q_v^k$, be
$V$, $H(L)$, $N_{\mathrm{src}}(L)$ and $P_{\mathrm{tgt}}(L)$ of Lemma~\ref{pd-gain-holds-local-chain} for
the Hida setting of $v$, the integers $k$, $r$ and $t$ and the weights $\lambda_0(v,\cdot)$. So
$H^\nabla(L)$ is the $\mathcal R$-span of the products $\prod_vf_v$ with $f_v\in H_v(L_v)$. Put
$P^t_{\mathrm{base},v}(L)=H_v(L)$ and $P^t_{\mathrm{src},v}(L)=N^t_{\mathrm{src},v}(L)$.
\begin{enumerate}
\item A \emph{labelling} $\sigma$ assigns to each orbit one of the labels $\mathrm{base}$,
  $\mathrm{src}$ and $\mathrm{tgt}$. For integers $t=(t_v)_v\geq0$ and a labelling $\sigma$, let
  $H_{t,\sigma}(L)$ be the $\mathcal R$-span of the products $\prod_vf_v$ with
  $f_v\in P^{t_v}_{\sigma(v),v}(L_v)$, and let $M_t(\sigma)\subset\mathrm{Fn}$ be the module of the
  $F$ with $F(c,L)\in H_{t,\sigma}(L)$ for every component $c$ and lattice tuple $L$, $F(c,L)=0$ for
  every other tuple $L$, and $F(c,\gamma L)=\chi_c(\gamma)\,\gamma\cdot F(c,L)$ for all
  $\gamma\in\Gamma_c$.
\item Write $\sigma[v\mapsto l]$ for $\sigma$ with its value at $v$ replaced by $l$. For a set $B$ of
  orbits, $\sigma_B$ is the labelling with value $\mathrm{tgt}$ on $B$ and $\mathrm{base}$ elsewhere,
  $M_t(B)=M_t(\sigma_B)$, and $J_t(B)$ is the sum over $v\in B$ of the modules
  $y_vM_t(\sigma_B[v\mapsto\mathrm{src}])$. So $M_t(\emptyset)$ is the module $M^\nabla$ defined as
  the integral module with $H^\nabla(L)$ in place of $H^\Delta(L)$.
\item For a component $c$, a lattice tuple $L_0$ fixed only by $1\in\Gamma_c$, and $\xi\in\mathcal X$,
  let $\delta_{c,L_0}(\xi)\in\mathrm{Fn}$ take the value $\chi_c(\gamma)\,\gamma\cdot\xi$ at
  $(c,\gamma L_0)$ for $\gamma\in\Gamma_c$, and $0$ at all other pairs.
\end{enumerate}
\end{definition}
\leanref{Formalism/SharpRestriction/GrowthBounds/LatticeSystems.lean}{Formalism.MainNotions.G.LatticeSystems.piLattice₂}
\leanref{Formalism/SharpRestriction/GrowthBounds/LatticeSystems.lean}{Formalism.MainNotions.G.LatticeSystems.latSys}
\leanref{Formalism/SharpRestriction/GrowthBounds/LatticeSystems.lean}{Formalism.MainNotions.G.LatticeSystems.MBof}
\leanref{Formalism/SharpRestriction/GrowthBounds/LatticeSystems.lean}{Formalism.MainNotions.G.LatticeSystems.JBof}
\leanref{Formalism/SharpRestriction/GrowthBounds/LatticeSystems.lean}{Formalism.MainNotions.G.LatticeSystems.delta}

\begin{lemma}\label{pd-gain-holds-orbits}
Let $G$ be a local model at a prime $p$ with ring $\mathcal R$ and function space $\mathrm{Fn}$, and
suppose that for every component $c$ only $1\in\Gamma_c$ fixes a lattice tuple and the lattice tuples
fall into finitely many orbits under $\Gamma_c$. Let $t=(t_v)_v$ be integers $\geq0$, $\sigma$ a
labelling, and $H_{t,\sigma}$, $M_t(\sigma)$ and $\delta_{c,L}$ as in
Definition~\ref{pd-gain-holds-systems}.
\begin{enumerate}
\item For every component $c$, $\gamma\in\Gamma_c$ and lattice tuple $L$,
  $\chi_c(\gamma)\,\gamma\cdot H_{t,\sigma}(L)\subset H_{t,\sigma}(\gamma L)$.
\item For a component $c$ and a lattice tuple $L_0$, $\delta_{c,L_0}$ is $\mathcal R$-linear,
  $\delta_{c,L_0}(\xi)(c,L_0)=\xi$, and $\delta_{c,L_0}(\xi)\in M_t(\sigma)$ for
  $\xi\in H_{t,\sigma}(L_0)$.
\item Let $\mathcal S$ be a set of pairs $(c,L)$ of a component and a lattice tuple that contains, for
  each $c$, one tuple from each $\Gamma_c$-orbit of lattice tuples. Then every $F\in M_t(\sigma)$
  equals $\sum_{(c,L)\in\mathcal S}\delta_{c,L}(F(c,L))$.
\item $M^\nabla$ is free of finite rank $d^\nabla$ over $\mathcal R$, and
  $|\mathcal S|\prod_v\dim V_v\leq d^\nabla$.
\item Let $\sigma'$ be a labelling that agrees with $\sigma$ outside an orbit $v$, and $a\in\mathcal R$
  with $aP^{t_v}_{\sigma(v),v}(L)\subset P^{t_v}_{\sigma'(v),v}(L)$ for every $W_v$-lattice $L$. Then
  $aM_t(\sigma)\subset M_t(\sigma')$.
\end{enumerate}
\end{lemma}
\leanref{Formalism/SharpRestriction/GrowthBounds/LatticeSystems.lean}{Formalism.MainNotions.G.LatticeSystems.JEquiv}
\leanref{Formalism/SharpRestriction/GrowthBounds/LatticeSystems.lean}{Formalism.MainNotions.G.LatticeSystems.jEquiv}
\leanref{Formalism/SharpRestriction/GrowthBounds/LatticeSystems.lean}{Formalism.MainNotions.G.LatticeSystems.JDelta}
\leanref{Formalism/SharpRestriction/GrowthBounds/LatticeSystems.lean}{Formalism.MainNotions.G.LatticeSystems.jDelta}
\leanref{Formalism/SharpRestriction/GrowthBounds/LatticeSystems.lean}{Formalism.MainNotions.G.LatticeSystems.JRank}
\leanref{Formalism/SharpRestriction/GrowthBounds/LatticeSystems.lean}{Formalism.MainNotions.G.LatticeSystems.jRank}
\leanref{Formalism/SharpRestriction/GrowthBounds/LatticeSystems.lean}{Formalism.MainNotions.G.LatticeSystems.mna_finite}
\leanref{Formalism/SharpRestriction/GrowthBounds/LatticeSystems.lean}{Formalism.MainNotions.G.LatticeSystems.mna_free}
\leanref{Formalism/SharpRestriction/GrowthBounds/LatticeSystems.lean}{Formalism.MainNotions.G.LatticeSystems.ch_step}

\begin{proof}
(1) The action of $\gamma$ on $\mathcal X$ is the product over the orbits $v$ of the actions of the
images $\gamma_v\in GL_k(Q_v)$ on the variables of $v$, as in Lemma~\ref{pd-gain-holds-local-chain}(4),
so $\gamma\cdot\prod_vf_v=\prod_v\gamma_v\cdot f_v$. By that lemma $f_v\in P^{t_v}_{\sigma(v),v}(L_v)$
gives $\gamma_v\cdot f_v\in P^{t_v}_{\sigma(v),v}(\gamma_vL_v)$, and $(\gamma L)_v=\gamma_vL_v$. Since
$\chi_c(\gamma)\in\mathcal R$ and the action is $\mathcal K$-linear, (1) follows.

(2) The function $\delta_{c,L_0}(\xi)$ is well defined: if $\gamma L_0=\gamma'L_0$, then
$\gamma^{-1}\gamma'$ fixes $L_0$, so $\gamma=\gamma'$. It is linear in $\xi$ and takes the value $\xi$ at
$(c,L_0)$. Let $\xi\in H_{t,\sigma}(L_0)$. At a tuple $\gamma L_0$ the value lies in
$H_{t,\sigma}(\gamma L_0)$ by (1). The other values are $0$, in particular those at tuples that are not
lattice tuples, since $\gamma$ carries lattice tuples to lattice tuples. For $\gamma'\in\Gamma_c$ the
value at $(c,\gamma'\gamma L_0)$ is $\chi_c(\gamma'\gamma)(\gamma'\gamma)\cdot\xi=
\chi_c(\gamma')\,\gamma'\cdot(\chi_c(\gamma)\,\gamma\cdot\xi)$, and the values at $(c,L)$ and
$(c,\gamma'L)$ both vanish when $L$ is not in the orbit of $L_0$. So $\delta_{c,L_0}(\xi)\in
M_t(\sigma)$.

(3) Evaluate both sides at a pair $(c,L)$. If $L$ is not a lattice tuple, both vanish. Otherwise
$L=\gamma L_1$ for the pair $(c,L_1)\in\mathcal S$ with $L_1$ in the orbit of $L$, and
$F(c,\gamma L_1)=\chi_c(\gamma)\,\gamma\cdot F(c,L_1)=\delta_{c,L_1}(F(c,L_1))(c,L)$, while the other
terms vanish at $(c,L)$.

(4) The map $F\mapsto(F(c,L))_{(c,L)\in\mathcal S}$ from $M^\nabla$ to $\prod_{\mathcal S}H^\nabla(L)$ is
$\mathcal R$-linear, and injective by (3). The set $\mathcal S$ is finite, and each $H^\nabla(L)$ is
finitely generated, being spanned by the products of the finite bases of the $H_v(L_v)$ given by
Lemma~\ref{pd-gain-holds-local-chain}(1). As $\mathcal R$ is noetherian, $M^\nabla$ is finitely
generated; it is torsion free, so it is free of finite rank. For $(c,L)\in\mathcal S$ the products
$Q_{c,L,j}$ of these bases are $\prod_v\dim V_v$ elements of $H^\nabla(L)$, linearly independent since
the variables of distinct orbits are distinct. By (2) the $\delta_{c,L}(Q_{c,L,j})$ lie in $M^\nabla$.
They are linearly independent: evaluated at $(c,L)\in\mathcal S$, where $\delta_{c',L'}$ vanishes for
the other pairs $(c',L')\in\mathcal S$, a linear relation among them becomes one among the $Q_{c,L,j}$.
Hence $|\mathcal S|\prod_v\dim V_v\leq d^\nabla$.

(5) Let $F\in M_t(\sigma)$. For a lattice tuple $L$, $aF(c,L)$ is an $\mathcal R$-combination of
products $(af_v)\prod_{u\neq v}f_u$ with $f_u\in P^{t_u}_{\sigma(u),u}(L_u)$, and
$af_v\in P^{t_v}_{\sigma'(v),v}(L_v)$. The other two conditions are preserved by multiplication by $a$.
\end{proof}

\begin{lemma}\label{pd-gain-holds-ball-dual}
Let $G$ be a local model at a prime $p$ with ring $\mathcal R$ and Hecke operators $y_v$, and suppose
that for every component $c$ only $1\in\Gamma_c$ fixes a lattice tuple and the lattice tuples fall into
finitely many orbits under $\Gamma_c$. Let $t=(t_v)_v$ be integers $\geq0$, and use the notation of
Definition~\ref{pd-gain-holds-systems}, with $A_v$ as in Definition~\ref{pd-local-model}(4).
\begin{enumerate}
\item The operators $y_u$ and $y_v$ commute for all orbits $u$ and $v$.
\item For every labelling $\sigma$ and orbit $v$,
  $y_vM_t(\sigma[v\mapsto\mathrm{src}])\subset M_t(\sigma[v\mapsto\mathrm{tgt}])$.
\item Let $B$ be a set of orbits, $v\in B$, $E\geq0$ an integer, $c$ a component, $L$ a lattice tuple,
  $L''\subset L_v$ an edge at $v$, and $\xi\in H_{t,\sigma_B[v\mapsto\mathrm{src}]}(L[v\mapsto L''])$,
  where $L[v\mapsto L'']$ is $L$ with $L_v$ replaced by $L''$. Suppose that
  $p^{A_v}\xi\in p^EH_{t,\sigma_B}(L[v\mapsto L'])$ for every edge $L''\subset L'$ at $v$ with
  $L'\neq L_v$. Then $\delta_{c,L}(p^{A_v}\xi)\in J_t(B)+p^EM_t(B)$.
\end{enumerate}
\end{lemma}
\leanref{Formalism/SharpRestriction/GrowthBounds/LatticeSystems.lean}{Formalism.MainNotions.G.LatticeSystems.ef_yv_comm}
\leanref{Formalism/SharpRestriction/GrowthBounds/LatticeSystems.lean}{Formalism.MainNotions.G.LatticeSystems.exponent_condition_holds}
\leanref{Formalism/SharpRestriction/GrowthBounds/LatticeSystems.lean}{Formalism.MainNotions.G.LatticeSystems.RestNat}
\leanref{Formalism/SharpRestriction/GrowthBounds/LatticeSystems.lean}{Formalism.MainNotions.G.LatticeSystems.nt_nat}
\leanref{Formalism/SharpRestriction/GrowthBounds/LatticeSystems.lean}{Formalism.MainNotions.G.LatticeSystems.JStep1}
\leanref{Formalism/SharpRestriction/GrowthBounds/LatticeSystems.lean}{Formalism.MainNotions.G.LatticeSystems.jStep1}

\begin{proof}
(1) Let $u\neq v$. For a pair $(c,L)$, both $(y_uy_vF)(c,L)$ and $(y_vy_uF)(c,L)$ equal
$p^{A_u+A_v}\sum F(c,L')$, the sum over the edges $L'''\subset L_u$ at $u$ and $L''\subset L_v$ at $v$,
where $L'$ is obtained from $L$ by replacing $L_u$ by $L'''$ and $L_v$ by $L''$: replacing $L_u$ does
not change the edges at $v$, and conversely.

(2) Let $F\in M_t(\sigma[v\mapsto\mathrm{src}])$. The vanishing of $y_vF$ at the tuples that are not
lattice tuples, and its equivariance, follow as in Step~2 of the proof of
Lemma~\ref{pd-gain-holds-hecke-stable}, which uses only these two properties of $F$. Let $L$ be a lattice
tuple. Then $(y_vF)(c,L)=p^{A_v}\sum_{L''}F(c,L[v\mapsto L''])$, each $L[v\mapsto L'']$ is a lattice
tuple, and so $F(c,L[v\mapsto L''])$ is an $\mathcal R$-combination of products $\prod_uf_u$ with
$f_v\in N^{t_v}_{\mathrm{src},v}(L'')$ and $f_u\in P^{t_u}_{\sigma(u),u}(L_u)$ for $u\neq v$. The
integer $A_v$ is the integer $A$ of Lemma~\ref{pd-gain-holds-local-chain} for the orbit $v$, so
$p^{A_v}f_v\in P^{t_v}_{\mathrm{tgt},v}(L_v)$ by part (3) of that lemma. Hence
$(y_vF)(c,L)\in H_{t,\sigma[v\mapsto\mathrm{tgt}]}(L)$.

(3) Put $F=\delta_{c,L[v\mapsto L'']}(\xi)$, which lies in $M_t(\sigma_B[v\mapsto\mathrm{src}])$ by
Lemma~\ref{pd-gain-holds-orbits}(2). Then $y_vF\in J_t(B)$, and $y_vF\in M_t(B)$ by (2), as
$\sigma_B[v\mapsto\mathrm{tgt}]=\sigma_B$. By Lemma~\ref{pd-gain-holds-local-chain}(3),
$p^{A_v}\xi\in H_{t,\sigma_B}(L)$, so $F_1=\delta_{c,L}(p^{A_v}\xi)$ lies in $M_t(B)$. We show that
$F_2=F_1-y_vF\in M_t(B)$ has $F_2(c',L_1)\in p^EH_{t,\sigma_B}(L_1)$ at every lattice tuple $L_1$. Then
$p^{-E}F_2$ satisfies the three conditions defining $M_t(B)$, so $F_2\in p^EM_t(B)$, and $F_1=y_vF+F_2$
lies in $J_t(B)+p^EM_t(B)$.

Both $F_1$ and $y_vF$ vanish at $(c',L_1)$ for $c'\neq c$. Let $c'=c$. By definition,
$(y_vF)(c,L_1)=p^{A_v}\sum_{L^\circ}F(c,L_1[v\mapsto L^\circ])$ over the edges $L^\circ\subset L_{1,v}$
at $v$. A term is nonzero only if $L_1[v\mapsto L^\circ]=\gamma(L[v\mapsto L''])$ for some
$\gamma\in\Gamma_c$. Then $L^\circ=\gamma_vL''$ and $L_{1,u}=\gamma_uL_u$ for $u\neq v$, and
$L'=\gamma_v^{-1}L_{1,v}$ contains $L''$ as an edge at $v$, since $\gamma_v$ preserves inclusions,
multiplication by $p$ and indices; also $L_1=\gamma(L[v\mapsto L'])$. If $L'\neq L_v$, the hypothesis
gives $p^{A_v}\xi=p^E\eta$ with $\eta\in H_{t,\sigma_B}(L[v\mapsto L'])$, and the term is
$p^{A_v}\chi_c(\gamma)\,\gamma\cdot\xi=p^E\chi_c(\gamma)\,\gamma\cdot\eta$, which lies in
$p^EH_{t,\sigma_B}(L_1)$ by Lemma~\ref{pd-gain-holds-orbits}(1). If $L'=L_v$, then $L_1=\gamma L$.

If $L_1$ is not in the orbit of $L$, then $F_1(c,L_1)=0$ and every nonzero term has $L'\neq L_v$, so
$F_2(c,L_1)\in p^EH_{t,\sigma_B}(L_1)$. If $L_1=\gamma_1L$, then $\gamma_1$ is unique, since only $1$
fixes $L$. The edge $L^\circ_1=\gamma_{1,v}L''\subset L_{1,v}$ gives the term
$p^{A_v}\chi_c(\gamma_1)\,\gamma_1\cdot\xi=F_1(c,L_1)$. For another edge $L^\circ$ with a nonzero term,
$L'=L_v$ would give $\gamma L=\gamma_1L$, so $\gamma=\gamma_1$ and $L^\circ=L^\circ_1$; hence
$L'\neq L_v$, and $F_2(c,L_1)$, the negative of the sum of these other terms, lies in
$p^EH_{t,\sigma_B}(L_1)$.
\end{proof}

\begin{lemma}\label{pd-gain-holds-local-smith}
Let $p\geq5$ be a prime, $(W,Q,\sigma,\iota)$ a Hida setting of degree $n'$ at $p$ with rings
$\mathcal R\subset\mathcal K$ and $\varpi_0$, $k\geq1$ and $r,t,D,c_*\geq0$ integers,
$\kappa\in\mathbf R$, and $\lambda_0(s)\in\mathbf Z^k$ dominant weights such that
$\mathrm{LP}(r,\lambda_0;t,D,\kappa,c_*)$ holds (Definition~\ref{pd-projection-property}). Let $V$, $A$,
$N_{\mathrm{src}}$ and $P_{\mathrm{tgt}}$ be as in Lemma~\ref{pd-gain-holds-local-chain}, and let
$L\subset Q^k$ be a $W$-lattice. Then $\kappa\geq0$, and there are edges
$L''_1\subset L,\dots,L''_K\subset L$, vectors $F_i\in N_{\mathrm{src}}(L''_i)$, an $\mathcal R$-basis
$b_1,\dots,b_{\dim V}$ of $P_{\mathrm{tgt}}(L)$ and a set $\mathrm{bad}$ of at most $\kappa\dim V$
indices such that:
\begin{enumerate}
\item $p^AF_i\in\varpi_0^{c_*}P_{\mathrm{tgt}}(L')$ for every $i$ and every edge $L''_i\subset L'$ with
  $L'\neq L$;
\item $p^Db_a\in\sum_i\mathcal R\,p^AF_i$ for every $a\notin\mathrm{bad}$.
\end{enumerate}
\end{lemma}
\leanref{Formalism/SharpRestriction/GrowthBounds/LatticeSystems.lean}{Formalism.MainNotions.G.LatticeSystems.JSmith}
\leanref{Formalism/SharpRestriction/GrowthBounds/LatticeSystems.lean}{Formalism.MainNotions.G.LatticeSystems.jSmith}
\leanref{Formalism/SharpRestriction/GrowthBounds/LatticeSystems.lean}{Formalism.MainNotions.G.LatticeSystems.kappa_nonneg}

\begin{proof}
The projection property at $L$ gives edges $e_i$ with target $L$ and sources $L''_i$, vectors
$F_i\in N_{\mathrm{src}}(L''_i)$ with $F_i\in\varpi_0^{c_*}R(e')$ for every edge $e'\neq e_i$ with
source $L''_i$, an integer $n\geq(1-\kappa)\dim V$, an $\mathcal R$-linear map
$u\colon\mathcal R^n\to N_{\mathrm{tgt}}(L)$ and an $n\times K$ matrix $(B_{ji})$ over $\mathcal R$ with
$\sum_iB_{ji}\langle F_i,u(y)\rangle=p^Dy_j$.

(1) Let $L''_i\subset L'$ be an edge with $L'\neq L$, and $e'$ the edge with source $L''_i$ and target
$L'$. Then $F_i=\varpi_0^{c_*}F'$ with $F'\in R(e')$, and $p^AF'\in P_{\mathrm{tgt}}(L')$ by
Lemma~\ref{pd-gain-holds-local-chain}(3).

(2) Let $(x_j)$ be the basis of $H(L)$ given by Lemma~\ref{pd-gain-holds-local-chain}(1). By part (2) of
that lemma, $\varpi_0^tH(L)\subset P_{\mathrm{tgt}}(L)\subset H(L)$, and by part (3),
$f_i=p^AF_i\in P_{\mathrm{tgt}}(L)$. Every $G\in N_{\mathrm{tgt}}(L)$ takes $P_{\mathrm{tgt}}(L)$ into
$\mathcal R$, and $u(y)(f_i)=\langle F_i,u(y)\rangle$. Lemma~\ref{pd-gain-holds-smith}, with
$P=P_{\mathrm{tgt}}(L)$, $\varpi=\varpi_0^t$, $N'=N_{\mathrm{tgt}}(L)$ and $q=p^D$, gives the basis
$(b_a)$ and a set $\mathrm{bad}$ with $|\mathrm{bad}|\leq\dim V-n\leq\kappa\dim V$.

Finally, $\kappa\geq0$. The projection property at the standard lattice $W^k\subset Q^k$ gives vectors
$F'_1,\dots,F'_{K'}\in V$, an integer $n_0\geq(1-\kappa)\dim V$, an $\mathcal R$-linear map
$u'\colon\mathcal R^{n_0}\to N_{\mathrm{tgt}}(W^k)$ and an $n_0\times K'$ matrix $(B'_{ji})$ over
$\mathcal R$ with $\sum_iB'_{ji}\langle F'_i,u'(y)\rangle=p^Dy_j$. The vectors
$g_j=\sum_iB'_{ji}\,p^AF'_i$ satisfy $u'(e_l)(g_j)=p^D\delta_{jl}$, and $p^D\neq0$, so they are
linearly independent over $\mathcal K$, and $n_0\leq\dim V$. The space $V$ is not zero, as it contains
the product of the polynomials $h_{\nu(s)}$ of Definition~\ref{pd-weight-lattices}. So
$(1-\kappa)\dim V\leq n_0\leq\dim V$ gives $\kappa\geq0$.
\end{proof}

\begin{lemma}\label{pd-gain-holds-covering}
Let $G$ be a local model at a prime $p$ with ring $\mathcal R$, and suppose that for every component $c$
only $1\in\Gamma_c$ fixes a lattice tuple and the lattice tuples fall into finitely many orbits under
$\Gamma_c$. Let $d^\nabla$ be the rank of $M^\nabla$, $V_0$ a set of orbits and $D,c_*\geq0$
integers; for every orbit $v$ let $t_v\geq0$ be an integer, and for $v\in V_0$ let $\kappa_v$ be a real
number such that $\mathrm{LP}(r,\lambda_0(v,\cdot);t_v,D,\kappa_v,c_*)$ holds at $v$. Let $B\subset V_0$
be nonempty. Then, in the notation of Definition~\ref{pd-gain-holds-systems}, there is a finite set
$C\subset M_t(B)$ with at most $(\prod_{v\in B}\kappa_v)d^\nabla$ elements such that
$p^{D+1}M_t(B)\subset J_t(B)+\mathcal RC$.
\end{lemma}
\leanref{Formalism/SharpRestriction/GrowthBounds/LatticeSystems.lean}{Formalism.MainNotions.G.LatticeSystems.RestJ}
\leanref{Formalism/SharpRestriction/GrowthBounds/LatticeSystems.lean}{Formalism.MainNotions.G.LatticeSystems.restJAll_of}
\leanref{Formalism/SharpRestriction/GrowthBounds/LatticeSystems.lean}{Formalism.MainNotions.G.LatticeSystems.JLocal}
\leanref{Formalism/SharpRestriction/GrowthBounds/LatticeSystems.lean}{Formalism.MainNotions.G.LatticeSystems.jLocal}
\leanref{Formalism/SharpRestriction/GrowthBounds/LatticeSystems.lean}{Formalism.MainNotions.G.LatticeSystems.jl_good}
\leanref{Formalism/SharpRestriction/GrowthBounds/LatticeSystems.lean}{Formalism.MainNotions.G.LatticeSystems.ja_restJ}

\begin{proof}
\emph{Step 1: one lattice tuple.} Let $c$ be a component and $L$ a lattice tuple. For $v\in B$ take the
edges $L''_{v,i}\subset L_v$, the vectors $F_{v,i}$, the basis $(b_{v,a})$ of
$P^{t_v}_{\mathrm{tgt},v}(L_v)$ and the set $\mathrm{bad}_v$, with at most $\kappa_v\dim V_v$ elements,
of Lemma~\ref{pd-gain-holds-local-smith}. For $v\notin B$ take the basis $(b_{v,a})$ of $H_v(L_v)$ of
Lemma~\ref{pd-gain-holds-local-chain}(1). Each basis has $\dim V_v$ elements, and $H_{t,\sigma_B}(L)$ is
spanned by the products $b_{\mathbf a}=\prod_vb_{v,a_v}$. Let $C_{c,L}$ be the set of the
$b_{\mathbf a}$ with $a_v\in\mathrm{bad}_v$ for every $v\in B$; it has at most
$(\prod_{v\in B}\kappa_v)\prod_v\dim V_v$ elements. We show that every $\xi\in H_{t,\sigma_B}(L)$ can be
written $p^D\xi=g+h$ with $\delta_{c,L}(g)\in J_t(B)+p^{D+1}M_t(B)$ and $h\in\mathcal RC_{c,L}$. Such
$\xi$ form an $\mathcal R$-submodule, as $\delta_{c,L}$ is linear, so it suffices to take
$\xi=b_{\mathbf a}$. If $a_v\in\mathrm{bad}_v$ for every $v\in B$, take $g=0$ and $h=p^Db_{\mathbf a}$.
Otherwise choose $v\in B$ with $a_v\notin\mathrm{bad}_v$, and take $h=0$ and $g=p^Db_{\mathbf a}$. By
Lemma~\ref{pd-gain-holds-local-smith}(2), $g$ is an $\mathcal R$-combination of the products
$p^{A_v}\xi_i$ with $\xi_i=F_{v,i}\prod_{u\neq v}b_{u,a_u}$, so it suffices that
$\delta_{c,L}(p^{A_v}\xi_i)\in J_t(B)+p^{D+1}M_t(B)$. This is Lemma~\ref{pd-gain-holds-ball-dual}(3)
with $E=D+1$ and the edge $L''_{v,i}\subset L_v$, since $\xi_i\in
H_{t,\sigma_B[v\mapsto\mathrm{src}]}(L[v\mapsto L''_{v,i}])$. Its hypothesis holds: for an edge
$L''_{v,i}\subset L'$ at $v$ with $L'\neq L_v$, Lemma~\ref{pd-gain-holds-local-smith}(1) gives
$p^{A_v}F_{v,i}=\varpi_0^{c_*}y$ with $y\in P^{t_v}_{\mathrm{tgt},v}(L')$, and as $c_*\geq4(D+1)$ and
$\varpi_0^2=p$, $p^{A_v}\xi_i=p^{D+1}\varpi_0^{c_*-2(D+1)}y\prod_{u\neq v}b_{u,a_u}$ lies in
$p^{D+1}H_{t,\sigma_B}(L[v\mapsto L'])$.

\emph{Step 2: all lattice tuples.} Let $\mathcal S$ be as in Lemma~\ref{pd-gain-holds-orbits}(3); it is
finite. Let $C$ be the union over $(c,L)\in\mathcal S$ of the sets $\delta_{c,L}(C_{c,L})$, which lie in
$M_t(B)$ by Lemma~\ref{pd-gain-holds-orbits}(2). Each $\kappa_v$ is $\geq0$ by
Lemma~\ref{pd-gain-holds-local-smith}, so by Lemma~\ref{pd-gain-holds-orbits}(4) $C$ has at most
$|\mathcal S|(\prod_{v\in B}\kappa_v)\prod_v\dim V_v\leq(\prod_{v\in B}\kappa_v)d^\nabla$ elements. Put
$T=J_t(B)+\mathcal RC$ and $N=p^DM_t(B)$. For $x\in M_t(B)$, Lemma~\ref{pd-gain-holds-orbits}(3) and
Step~1 give $p^Dx=\sum_{(c,L)\in\mathcal S}\delta_{c,L}(p^Dx(c,L))\in T+p^{D+1}M_t(B)$, so $N\subset
T+pN$. By Lemma~\ref{pd-gain-holds-local-chain}(2) and Lemma~\ref{pd-gain-holds-orbits}(5), applied at
the elements of $B$ one at a time, $M_t(B)\subset M^\nabla$, which is finitely generated by
Lemma~\ref{pd-gain-holds-orbits}(4); so $N$ is finitely generated. The element $p$ is not a unit of
$\mathcal R$: as $\mathcal R$ is finite over $W_v$ for an orbit $v$, some prime ideal of $\mathcal R$
lies over $pW_v$. By Nakayama's lemma, applied to $(N+T)/T=p\,(N+T)/T$, $N\subset T$. Hence
$p^{D+1}M_t(B)\subset pT\subset T$.
\end{proof}

\begin{lemma}\label{pd-gain-holds-selection}
Let $G$ be a local model at a prime $p$ (Definition~\ref{pd-local-model}) with ring $\mathcal R$,
integral module $M$ and Hecke operators $y_v$, and let $M^\nabla$ be the module defined as $M$ with
$H^\nabla(L)$ in place of $H^\Delta(L)$, of rank $d^\nabla$ over $\mathcal R$. Suppose that for
every component $c$ the only element of $\Gamma_c$ that fixes a lattice tuple is $1$, and that the
lattice tuples fall into finitely many orbits under $\Gamma_c$. Let $V_0$ be a set of orbits,
$m\geq1$ and $D,c_*,S\geq0$ integers, and for $v\in V_0$ let $t_v\geq0$ and $0\leq S_v\leq S$ be
integers and
$\kappa_v$ real numbers such that $\mathrm{LP}(r,\lambda_0(v,\cdot);t_v,D,\kappa_v,c_*)$
(Definition~\ref{pd-projection-property}) holds at the orbit $v$. The operators $y_v$ commute with
each other. For a submodule $M'$ of the function space let $\Pi(M')$ be the $\mathcal R$-span of the
functions $p^{\sum_{v\in V_0\smallsetminus B}S_v}y_B(F)$, for $B\subset V_0$ and $F\in M'$, where
$y_B$ is the composite of the $y_v^m$, $v\in B$; and put $\tau=\sum_{v\in V_0}m(t_v+D+1)$. Then for
every
integer $\beta\geq0$ there is a finite set $C\subset M^\nabla$ with at most
$\big\lfloor d^\nabla m^{\beta+1}e_{\beta+1}\big((\kappa_v)_{v\in V_0}\big)\big\rfloor$ elements such
that $p^{\tau+\beta S}M^\nabla\subset\Pi(M^\nabla)+\mathcal RC$.
\end{lemma}
\leanref{Formalism/SharpRestriction/GrowthBounds/SelectionLemmas.lean}{Formalism.MainNotions.SelectionLemmas.Conclusion1}
\leanref{Formalism/SharpRestriction/GrowthBounds/SelectionLemmas.lean}{Formalism.MainNotions.SelectionLemmas.prodIdeal}
\leanref{Formalism/SharpRestriction/GrowthBounds/LatticeSystems.lean}{Formalism.MainNotions.G.LatticeSystems.theta_kernel_expansion}
\leanref{Formalism/SharpRestriction/GrowthBounds/LatticeSystems.lean}{Formalism.MainNotions.G.LatticeSystems.exponent_condition_holds}
\leanref{Formalism/SharpRestriction/GrowthBounds/LatticeSystems.lean}{Formalism.MainNotions.G.LatticeSystems.counting_statement_of_block_form}
\leanref{Formalism/SharpRestriction/GrowthBounds/LatticeSystems.lean}{Formalism.MainNotions.G.LatticeSystems.walk_allBut_bound}

\begin{proof}
The operators $y_v$ commute by Lemma~\ref{pd-gain-holds-ball-dual}(1). If $V_0=\emptyset$, then
$\tau=0$, the term $B=\emptyset$ of $\Pi(M^\nabla)$ gives $M^\nabla\subset\Pi(M^\nabla)$, and
$C=\emptyset$ will do. Let $V_0\neq\emptyset$, put $t_v=0$ for the orbits $v\notin V_0$, order $V_0$,
and use the notation of Definition~\ref{pd-gain-holds-systems}. Apply Lemma~\ref{pd-gain-holds-walk}
with $\mathcal O=\mathcal R$, $X=\mathrm{Fn}$, $\pi=\varpi_0$, $X(\sigma)=M_t(\sigma)$ for the
labellings $\sigma$ of $V_0$, extended by $\mathrm{base}$ outside $V_0$, the operators $y_v$ for
$v\in V_0$, the integers $m$, $\beta$, $c_0=2(D+1)$ and $2S$ in place of $S$, $a_v=b_v=t_v$, $c_v=2t_v$
and $2S_v$ in place of $S_v$, and $G_B$ the set $C$ of Lemma~\ref{pd-gain-holds-covering} for $B$. Its
hypotheses hold:
\begin{itemize}
\item (1) and (2): by Lemma~\ref{pd-gain-holds-local-chain}(2),
  $P^{t_v}_{\mathrm{tgt},v}(L)\subset H_v(L)\subset P^{t_v}_{\mathrm{src},v}(L)$,
  $\varpi_0^{t_v}P^{t_v}_{\mathrm{src},v}(L)\subset H_v(L)$ and
  $\varpi_0^{t_v}H_v(L)\subset P^{t_v}_{\mathrm{tgt},v}(L)$ for every lattice $L$. By
  Lemma~\ref{pd-gain-holds-orbits}(5) these give (1) and the first two inclusions of (2), and the third
  inclusion of (2) is the composite of the first two.
\item (3) is Lemma~\ref{pd-gain-holds-ball-dual}(2).
\item (4): $\varpi_0^{2(D+1)}=p^{D+1}$, and Lemma~\ref{pd-gain-holds-covering} applies to every $B$ with
  $|B|=\beta+1\geq1$.
\end{itemize}
Here $X_0=M_t(\emptyset)=M^\nabla$, and since $\varpi_0^{2S_v}=p^{S_v}$ and the $y_v$ commute, the
module $\Pi$ of Lemma~\ref{pd-gain-holds-walk} is $\Pi(M^\nabla)$. Its exponent is
\[
 E=\sum_{v\in V_0}\big(2t_v+2(m-1)t_v+2m(D+1)\big)+2\beta S=2(\tau+\beta S),
\]
so $\varpi_0^E=p^{\tau+\beta S}$. The set $C\subset M^\nabla$ it gives has at most
\[
 m^{\beta+1}\sum_{|B|=\beta+1}\Big(\prod_{v\in B}\kappa_v\Big)d^\nabla
 =d^\nabla m^{\beta+1}e_{\beta+1}\big((\kappa_v)_{v\in V_0}\big)
\]
elements, and as $|C|$ is an integer, it is at most the integer part of the right side.
\end{proof}

\begin{lemma}\label{pd-gain-holds-hecke-rows}
Let $R$ be an integral domain, $k\geq1$, $\nu_1\geq\cdots\geq\nu_k\geq0$ integers, and let $X$ be
$\nabla_R(\nu)$ or $\Delta_R(\nu)$ (Definition~\ref{pd-weight-lattices}(1), with $n=k$ and $x$ for the
matrix of variables).
\begin{enumerate}
\item $X$ is spanned over $R$ by polynomials $f$ that are homogeneous in each row of $x$, of degree
  $\mu_a$ in the $a$-th row, with $\mu_1+\cdots+\mu_r\leq\nu_1+\cdots+\nu_r$ for every $r\geq0$.
\item If $f\in X$ and $h\in GL_k(R)$, then $f(hx)\in X$.
\end{enumerate}
\end{lemma}
\leanref{Formalism/SharpRestriction/GrowthBounds/LatticeSystems.lean}{Formalism.MainNotions.G.LatticeSystems.lb_basis}
\leanref{Formalism/LocalToGlobal/GroupLattices/FockModels/OperatorsPreserveLattices.lean}{Formalism.L2.G.Model.applyWord_hom}
\leanref{Formalism/LocalToGlobal/GroupLattices/FockModels/OperatorsPreserveLattices.lean}{Formalism.L2.G.Model.dom_support}
\leanref{Formalism/GradedWeylDivisibility/DualWeylDominanceOf.lean}{Formalism.GradedWeylDivisibility.weylLeDual}
\leanref{Formalism/GradedWeylDivisibility/DualWeylDominanceOf.lean}{Formalism.GradedWeylDivisibility.dualWeylRowStable}
\leanref{Formalism/SatakeDualWeyl/AdelicIwasawaHecke/Wlocal.lean}{Formalism.UG.P.weyl_module_GL_stable}

\begin{proof}
(1) By the standard basis theorem for bitableaux \cite{DEP80}, $\nabla_R(\nu)$ has an $R$-basis of
bideterminants of shape $\nu$: products of minors of $x$, the $l$-th on rows
$I_l(1)<\cdots<I_l(s_l)$ and on the first $s_l$ columns, where $\#\{l:s_l\geq t\}=\nu_t$ for every
$t$. Such a product is homogeneous in each row, and its degree in the first $r$ rows is at most
$\sum_l\min(s_l,r)=\nu_1+\cdots+\nu_r$. So every monomial of an element of $\nabla_R(\nu)$ has degree
at most $\nu_1+\cdots+\nu_r$ in the first $r$ rows. The module $\Delta_R(\nu)$ is spanned by the
polynomials obtained from $h_\nu$ by the operators $E^{(j)}_{ab}$. The polynomial $h_\nu$ is
homogeneous in each row, and $E^{(j)}_{ab}$ carries a polynomial homogeneous in each row to one
homogeneous in each row, moving $j$ units of degree from row $b$ to row $a$. These polynomials lie in
$\nabla_R(\nu)$, since $h_\nu$ does and the operators $E^{(j)}_{ab}$, which act on the rows, preserve
the column degrees and the invariance under $x\mapsto x(1+yE_{lj})$. So each nonzero one has
the bound on its row degrees, read off from one of its monomials.

(2) For $\nabla_R(\nu)$ this holds for every $k\times k$ matrix $h$ over $R$: $x\mapsto hx$ preserves
the column degrees and commutes with $x\mapsto x(1+yE_{lj})$. For $\Delta_R(\nu)$, note first that
$f((1+cE_{ab})x)=\sum_jc^jE^{(j)}_{ba}f$ for $c\in R$ and $1\leq a,b\leq k$, so these substitutions
preserve $\Delta_R(\nu)$. If $R$ is local, Gaussian elimination writes $h$ as a product of matrices
$1+cE_{ab}$: column by column, a unit in the current column is moved to the diagonal by adding a row,
scaled to $1$ by a matrix $1+(u^{-1}-1)E_{aa}$, and the rest of the column is cleared by adding
multiples of its row. In general let $J$ be the ideal of the $t\in R$ with $t\,f(hx)\in\Delta_R(\nu)$.
If $J\neq R$, it lies in a maximal ideal $P$. The local case over $R_P$ gives
$f(hx)\in\Delta_{R_P}(\nu)$, which is the localization of $\Delta_R(\nu)$, so $t\,f(hx)\in\Delta_R(\nu)$
for some $t\notin P$, a contradiction. So $1\in J$.
\end{proof}

\begin{lemma}\label{pd-gain-holds-hecke-step}
Let $R$ be a discrete valuation ring with fraction field $K$, $k\geq1$, $\lambda\in\mathbf Z^k$ with
$\lambda_1\geq\cdots\geq\lambda_k$, $\nu_i=\lambda_i-\lambda_k$, $c\in R$ nonzero, $r\geq0$, and $D$ the
diagonal $k\times k$ matrix with entries $c^{-1}$ in the places $i\leq r$ and $1$ elsewhere. Let $X$ be
$\nabla$ or $\Delta$. For an $R$-lattice $M\subset K^k$ let $X_\lambda(M^\vee)$ be the set of
$f\in K[x]$ such that $\det(g)^{\lambda_k}f(gx)\in X_R(\nu)$ for every $g$ whose columns form an
$R$-basis of $M$, as in Definition~\ref{pd-local-model}. Let $g\in GL_k(K)$, and let $M$ and $M'$ be
the $R$-spans of the columns of $g$ and of $gD$. Then
$c^{\lambda_1+\cdots+\lambda_r}X_\lambda(M^\vee)\subset X_\lambda(M'^\vee)$, where
$\lambda_1+\cdots+\lambda_r$ stands for the sum of the $\lambda_i$ with $i\leq\min(r,k)$.
\end{lemma}
\leanref{Formalism/SharpRestriction/GrowthBounds/LatticeSystems.lean}{Formalism.MainNotions.G.LatticeSystems.lb_pos}
\leanref{Formalism/SharpRestriction/GrowthBounds/LatticeSystems.lean}{Formalism.MainNotions.G.LatticeSystems.lb_P}
\leanref{Formalism/SharpRestriction/GrowthBounds/LatticeSystems.lean}{Formalism.MainNotions.G.LatticeSystems.lb_indep}
\leanref{Formalism/LocalToGlobal/GroupLattices/FockModels/OperatorsPreserveLattices.lean}{Formalism.L2.G.Model.dl_pos}
\leanref{Formalism/LocalToGlobal/GroupLattices/FockModels/OperatorsPreserveLattices.lean}{Formalism.L2.G.Model.dl_P}
\leanref{Formalism/LocalToGlobal/GroupLattices/FockModels/OperatorsPreserveLattices.lean}{Formalism.L2.G.Model.dl_indep}
\leanref{Formalism/LocalToGlobal/GroupLattices/FockModels/OperatorsPreserveLattices.lean}{Formalism.L2.G.Model.weyl_scale}

\begin{proof}
Let $f\in X_\lambda(M^\vee)$ and $A=\lambda_1+\cdots+\lambda_r$. Since the columns of $g$ form a basis
of $M$, $F(x)=\det(g)^{\lambda_k}f(gx)$ lies in $X_R(\nu)$. Let $\rho=\min(r,k)$, and read sums over
$i\leq r$ as sums over $i\leq\rho$. Then $\det(gD)=\det(g)c^{-\rho}$ and $A=(\nu_1+\cdots+\nu_r)+\rho\lambda_k$, so
\[
 c^A\det(gD)^{\lambda_k}f(gDx)=c^{\nu_1+\cdots+\nu_r}F(Dx).
\]
By Lemma~\ref{pd-gain-holds-hecke-rows}(1), $F$ is an $R$-combination of polynomials $F_i\in
X_R(\nu)$, homogeneous in each row, with row degrees $\mu_{i,a}$ satisfying
$\mu_{i,1}+\cdots+\mu_{i,r}\leq\nu_1+\cdots+\nu_r$. The substitution $x\mapsto Dx$ multiplies the
first $r$ rows of $x$ by $c^{-1}$, so
$c^{\nu_1+\cdots+\nu_r}F_i(Dx)=c^{e_i}F_i(x)$ with $e_i=\sum_{a\leq r}(\nu_a-\mu_{i,a})\geq0$. Hence
$G(x)=c^A\det(gD)^{\lambda_k}f(gDx)$ lies in $X_R(\nu)$. This is the condition for the basis $gD$ of
$M'$. Every other basis is $h=gDU$ with $U\in GL_k(R)$, because both bases span $M'$. Then
$c^A\det(h)^{\lambda_k}f(hx)=\det(U)^{\lambda_k}G(Ux)$, which lies in $X_R(\nu)$ by
Lemma~\ref{pd-gain-holds-hecke-rows}(2), since $\det(U)$ is a unit of $R$. So $c^Af\in
X_\lambda(M'^\vee)$.
\end{proof}

\begin{lemma}\label{pd-gain-holds-hecke-stable}
Let $G$ be a local model at a prime $p$ with integral module $M$ and Hecke operators $y_v$, and let
$M^\nabla$ be as in Lemma~\ref{pd-gain-holds-selection}. Then $y_vM\subset M$ and $y_vM^\nabla\subset
M^\nabla$ for every orbit $v$.
\end{lemma}
\leanref{Formalism/LocalToGlobal/GroupLattices/FockModels/OperatorsPreserveLattices.lean}{Formalism.L2.G.Model.operators_preserve_lattices}
\leanref{Formalism/SharpRestriction/GrowthBounds/FormalGroupSystem.lean}{Formalism.MainNotions.FormalGroupSystem.adm12}

\begin{proof}
Let $X$ be $\Delta$ or $\nabla$, and $H^X$ and $M_X$ the corresponding $H^\Delta$ and $M$, or
$H^\nabla$ and $M^\nabla$. Fix an orbit $v$, and let $\mathcal R$, $\mathcal K$ and $\varpi_0$ be as
in Definition~\ref{pd-local-model}.

\emph{Step 1: one edge.} Let $L$ be a lattice tuple and $L''\subset L_v$ an edge at $v$. We show that
$p^{A_v}H^X(L[v\mapsto L''])\subset H^X(L)$, where $L[v\mapsto L'']$ is $L$ with $L_v$ replaced by
$L''$. The quotient $L_v/L''$ is killed by $p$ and has $p^{n'_vr}$ elements, and the residue field of
$W_v$ has $p^{n'_v}$ elements. By the elementary divisor theorem over the discrete valuation ring
$W_v$, there is a matrix $g$ whose columns form a $W_v$-basis of $L''$ such that the columns of $gE$
form a basis of $L_v$, where $E$ is diagonal with entries $p^{-1}$ in the places $i\leq r$ and $1$
elsewhere. Let $s$ be a position of $v$ and $g_s=\iota_v\sigma_v^s(g)$, applied to the entries. The
columns of $g_s$ form an $\mathcal R$-basis of the $\mathcal R$-span of $\iota_v\sigma_v^s(L'')$, and,
since $\sigma_v$ fixes $p$ and $\varpi_0^2=p$, those of $g_sDD$ form a basis of the $\mathcal R$-span
of $\iota_v\sigma_v^s(L_v)$, where $D$ is diagonal with entries $\varpi_0^{-1}$ in the places $i\leq
r$ and $1$ elsewhere. Lemma~\ref{pd-gain-holds-hecke-step}, applied twice with $c=\varpi_0$ and
$\lambda=\lambda(v,s)$, shows that $p^{\sum_{i\leq r}\lambda(v,s)_i}f_{v,s}$ lies in the lattice of the
position $(v,s)$ for $L$ whenever $f_{v,s}$ lies in that for $L[v\mapsto L'']$. The two tuples agree at
the other orbits. Multiplying over the positions $s$ of $v$, the product of these powers of $p$ is
$p^{A_v}$, so $p^{A_v}$ times each product $\prod_{v',s}f_{v',s}$ spanning $H^X(L[v\mapsto L''])$
lies in $H^X(L)$.

\emph{Step 2: the three conditions.} Let $F\in M_X$ and $c$ a component. The edges $L''\subset L_v$
lie between $pL_v$ and $L_v$, so they are finite in number when $L_v$ is a lattice, and there are none
when it is not. If $L$ is a lattice tuple, every $L[v\mapsto L'']$ is one, so $F(c,L[v\mapsto
L''])\in H^X(L[v\mapsto L''])$, and Step 1 gives $(y_vF)(c,L)=p^{A_v}\sum_{L''}F(c,L[v\mapsto
L''])\in H^X(L)$. If $L$ is not a lattice tuple, then either $L_v$ is not a lattice and the sum is
empty, or $L_u$ is not a lattice for some $u\neq v$ and no $L[v\mapsto L'']$ is a lattice tuple; in
both cases $(y_vF)(c,L)=0$. An element $\gamma\in\Gamma_c$ carries the edges $L''\subset L_v$
bijectively to the edges at $(\gamma L)_v$, with $\gamma(L[v\mapsto L''])=(\gamma L)[v\mapsto\gamma
L'']$, and its action on $\mathcal X$ is $\mathcal K$-linear. So
\[
 (y_vF)(c,\gamma L)=p^{A_v}\sum_{L''}F(c,\gamma(L[v\mapsto L'']))
 =\chi_c(\gamma)\,\gamma\cdot(y_vF)(c,L).
\]
Hence $y_vF\in M_X$.
\end{proof}

\begin{lemma}\label{pd-gain-holds-served-length}
Let $D$ be a definite-side datum with ranks $k,m,t,n$, size data, gain data with Weyl exponents
$e_p$, local models, a tie and a lattice instance (Definitions~\ref{pd-size-data},
\ref{pd-gain-data} and~\ref{pd-lattice-instance}), with the gain quantities of
Definition~\ref{pd-gain-quantities}. Let $w$ be a place admissible for a word $u$, $p=p_w$,
$\mathcal R$ and $M$ the ring and the integral module of $G(u,w)$, $M^\nabla$ as in
Lemma~\ref{pd-gain-holds-selection}, $d$ the rank of $M$ over $\mathcal R$,
$M_{\mathrm{ram}}=M\cap V_{\mathrm{ram}}$, and $c_*\geq0$ an integer. Suppose that:
\begin{enumerate}
\item for every component $c$ of $G(u,w)$, only $1\in\Gamma_c$ fixes a lattice tuple, and the
  lattice tuples fall into finitely many orbits under $\Gamma_c$;
\item $y_vM\subset M$ and $y_vM^\nabla\subset M^\nabla$ for every orbit $v$, and
  $p^{e_p}M^\nabla\subset M\subset M^\nabla$;
\item $m\geq1$, and $t_0\mapsto v(t_0)$ is injective on $V(p)$;
\item for every $t_0\in V(p)$ the projection property $\mathrm{LP}\big(r,\lambda_0(v(t_0),\cdot);
  \max(\lceil t(p,t_0)\rceil,0),D(p),\kappa_{p,t_0}(\mathbf a,\mathbf b),c_*\big)$ holds at the orbit
  $v(t_0)$, where $r$ is the integer of $G(u,w)$;
\item $M_{\mathrm{ram}}$ is saturated in $M$: if $y\in M$, $\lambda\in\mathcal R$ is nonzero and
  $\lambda y\in M_{\mathrm{ram}}$, then $y\in M_{\mathrm{ram}}$;
\item $m\sum_{t_0\in V(p)}\kappa_{p,t_0}(\mathbf a,\mathbf b)\leq1/2$.
\end{enumerate}
Then $N_P(u,w)\subset p^{\max(\lceil\tau(p)\rceil,0)}M_{\mathrm{ram}}$, and
$M_{\mathrm{ram}}/N_P(u,w)$ has length at most $v_{\mathcal R}(p)\,d\,(\tau(p)+\iota(p))$ as an
$\mathcal R$-module.
\end{lemma}
\leanref{Formalism/SharpRestriction/GrowthBounds/MultisetValuation/StabilizerCount.lean}{Formalism.MainNotions.FormalGroupSystem.Pf.stmt6G}
\leanref{Formalism/SharpRestriction/GrowthBounds/MultisetValuation/StabilizerCount.lean}{Formalism.MainNotions.FormalGroupSystem.Pf.coreT}
\leanref{Formalism/SharpRestriction/GrowthBounds/MultisetValuation/StabilizerCount.lean}{Formalism.MainNotions.FormalGroupSystem.Pf.core}
\leanref{Formalism/SharpRestriction/GrowthBounds/FormalGroupSystem.lean}{Formalism.MainNotions.FormalGroupSystem.LemmaVVolumeGeneral}
\leanref{Formalism/SharpRestriction/GrowthBounds/LatticeSystems.lean}{Formalism.MainNotions.G.LatticeSystems.sat_clause}

\begin{proof}
Put $T=\max(\lceil\tau(p)\rceil,0)$. Every $x\in N_P(u,w)$ lies in $p^TM\cap V_{\mathrm{ram}}$, so
$x\in M_{\mathrm{ram}}$ and $x=p^Tz$ with $z\in M$; by (5), $z\in M_{\mathrm{ram}}$. This proves the
inclusion.

If $V(p)=\emptyset$, then $\tau(p)=\iota(p)=0$, and taking $V'=\emptyset$ in the definition of
$N_P(u,w)$ shows that it contains $M\cap V_{\mathrm{ram}}$; so $N_P(u,w)=M_{\mathrm{ram}}$ and the
length is $0$. Suppose that $V(p)\neq\emptyset$. Then $V(p)$ is the set of good representatives at
$p$, and
\[
 \tau(p)=e_p+\tau',\qquad\tau'=\sum_{t_0\in V(p)}m\big(\max(\lceil t(p,t_0)\rceil,0)+D(p)+1\big),
 \qquad\iota(p)=2\Sigma(p)S,
\]
where $\Sigma(p)=m\sum_{t_0\in V(p)}\kappa_{p,t_0}(\mathbf a,\mathbf b)$ and $S$ is the largest of the
surpluses $S(p,t_0)$, $t_0\in V(p)$, which are integers $\geq0$; in particular $T=\tau(p)$. By (3),
$t_0\mapsto v(t_0)$ is a bijection from $V(p)$ onto a set $V_0$ of orbits; for $v=v(t_0)$ write
$t_v=\max(\lceil t(p,t_0)\rceil,0)$, $S_v=S(p,t_0)$ and $\kappa_v=\kappa_{p,t_0}(\mathbf a,\mathbf b)$,
so that $\tau'=\sum_{v\in V_0}m(t_v+D(p)+1)$ and $\Sigma(p)=m\sum_{v\in V_0}\kappa_v$. Each
$\kappa_v$ is $\geq0$: in the projection property (4) at the lattice $W_v^k$, the vectors
$\sum_iB_{ji}p^AF_i$, $j\leq n$, are linearly independent in the space $V$ of
Definition~\ref{pd-projection-property}, since their pairings with $u(y)$ are the $p^Dy_j$; so
$(1-\kappa_v)\dim V\leq n\leq\dim V$, and $\dim V>0$ because the dual Weyl modules of dominant weights
are nonzero. Let $\Pi(M')$ be as in Lemma~\ref{pd-gain-holds-selection}, with these $V_0$, $m$ and
$S_v$. Taking $V'=v^{-1}(B)$ in the
definition of $N_P(u,w)$ shows that $\Pi(M)$ lies in the span used there, so
$N_P(u,w)\supset\Pi(M)\cap p^{\tau(p)}M\cap V_{\mathrm{ram}}$. By (2), $\Pi(M)\subset M$ and
$\Pi(M^\nabla)\subset M^\nabla$.

Let $\beta\geq0$. By (1), (4) and Lemma~\ref{pd-gain-holds-selection} with $D=D(p)$, $M^\nabla$ is
covered by $\Pi(M^\nabla)$ up to $g_\beta=\lfloor d^\nabla m^{\beta+1}e_{\beta+1}((\kappa_v))\rfloor$
elements at exponent $\tau'+\beta S$, in the sense of Lemma~\ref{pd-gain-holds-transfer} with
$\pi=p$. By (2),
$p^{e_p}\Pi(M^\nabla)\subset\Pi(p^{e_p}M^\nabla)\subset\Pi(M)$, so
Lemma~\ref{pd-gain-holds-transfer}(2), with $A=M^\nabla$, $B=M$, $N_A=\Pi(M^\nabla)$,
$N_B=\Pi(M)$ and $e=e_p$, shows that $M$ is covered by $\Pi(M)$ up to $g_\beta$ elements at exponent
$\tau(p)+\beta S$. The module $M$ is finitely generated, being contained in $M^\nabla$, which is
finitely generated since there are finitely many classes of lattice tuples, and it is torsion free, so
it is free of rank $d$; and $d^\nabla\leq d$ because multiplication by $p^{e_p}$ embeds $M^\nabla$ in
$M$. By Lemma~\ref{pd-gain-holds-esymm},
\[
 g_\beta\leq d^\nabla m^{\beta+1}e_{\beta+1}((\kappa_v))\leq
 \frac{d^\nabla\Sigma(p)^{\beta+1}}{(\beta+1)!}\leq d\,\Sigma(p)^{\beta+1}.
\]
Taking $B=\emptyset$ in the definition of $\Pi(M)$ gives
$p^{|V_0|S}M\subset p^{\sum_vS_v}M\subset\Pi(M)$. Lemma~\ref{pd-gain-holds-length}, with $a=p$,
$N=\Pi(M)$, $\tau=\tau(p)$, this $S$, $V=|V_0|$ and $s=\Sigma(p)\leq1/2$ (by (6)), shows
that $M/(\Pi(M)\cap p^{\tau(p)}M)$ has length at most
$v_{\mathcal R}(p)\,d\,(\tau(p)+2S\Sigma(p))=v_{\mathcal R}(p)\,d\,(\tau(p)+\iota(p))$, where
$v_{\mathcal R}(p)$ is the length of $\mathcal R/p\mathcal R$.

Put $N'=\Pi(M)\cap p^{\tau(p)}M\cap V_{\mathrm{ram}}$, a submodule of $M_{\mathrm{ram}}$. The map
$M_{\mathrm{ram}}/N'\to M/(\Pi(M)\cap p^{\tau(p)}M)$ is injective, since an element of
$M_{\mathrm{ram}}$ in $\Pi(M)\cap p^{\tau(p)}M$ lies in $N'$. So $M_{\mathrm{ram}}/N'$ has length at
most $v_{\mathcal R}(p)d(\tau(p)+\iota(p))$, and $M_{\mathrm{ram}}/N_P(u,w)$ is a quotient of it,
since $N'\subset N_P(u,w)$.
\end{proof}

\begin{lemma}\label{pd-gain-holds-window-sizes}
Let $0<\varepsilon<1/52$. There is $C_\varepsilon>0$ such that for every $x\geq3$ the numbers $G'$ and
$N_1$ of Definition~\ref{hf-good-orbit} satisfy $G'\leq C_\varepsilon x^{\varepsilon/8}\log x$ and
$N_1\leq C_\varepsilon\log x$.
\end{lemma}
\leanref{Formalism/SharpRestriction/GrowthBounds/TopChoiceBounds.lean}{Formalism.MainNotions.G.TopChoiceBounds.segment_count_log_bounds}
\leanref{Formalism/SharpRestriction/GrowthBounds/TopChoiceBounds.lean}{Formalism.MainNotions.G.TopChoiceBounds.bc_size}
\leanref{Formalism/SharpRestriction/GrowthBounds/TopChoiceBounds.lean}{Formalism.MainNotions.G.TopChoiceBounds.master_orbit_sizes_hold}

\begin{lemma}\label{pd-gain-holds-condition-five}
Let $D$ be a definite-side datum with rank $k$, size data with scale $(x,\varepsilon,J)$, $x\geq3$, and
gain data whose constant $K_\Lambda$ is positive, and use the notation of
Definitions~\ref{hf-good-orbit} and~\ref{pd-gain-data}. Let $\delta>0$ and $C>0$ satisfy
$G'\leq Cx^{\varepsilon/8}\log x$, $N_1\leq C\log x$, $k\leq x^6$ and
$(2K_\Lambda C^2+38C+1/K_\Lambda)(\log x)^3\leq x^\delta$. Then condition (5) of
Definition~\ref{pd-gain-admissibility} holds with $\delta$.
\end{lemma}
\leanref{Formalism/SharpRestriction/GrowthBounds/PlaneCountTheoremWithConstants.lean}{Formalism.MainNotions.G.sizes_bound_of_log_estimates}
\leanref{Formalism/SharpRestriction/GrowthBounds/PlaneCountTheoremWithConstants.lean}{Formalism.MainNotions.G.family_sizes_bound}

\begin{proof}
Put $l=\log x\geq1$ and $y=x^\delta$. Each of $2K_\Lambda C^2l^3$, $38Cl^3$ and $1/K_\Lambda$ is at most
$y$. Since $\log a=l+x^\varepsilon$ and $N_1\geq1$,
$\Lambda=K_\Lambda N_1^2(l+x^\varepsilon)\geq K_\Lambda x^\varepsilon\geq x^\varepsilon/y=
x^{\varepsilon-\delta}$. Since $l+x^\varepsilon\leq2lx^\varepsilon$ and $N_1\leq Cl$,
$\Lambda\leq K_\Lambda C^2l^2\cdot2lx^\varepsilon\leq x^\varepsilon y$. Let $p$ be a prime and
$Q=p^{G'}$. If $p<p_0=k^2x^{26}\leq x^{38}$, then $\log p\leq38l$ and
$\log Q=G'\log p\leq Cx^{\varepsilon/8}l\cdot38l\leq x^{\varepsilon/8}\cdot38Cl^3\leq
x^{\varepsilon/8}y$. If $p<X'$, then $\log p\leq x^{\varepsilon/3}$, and as $Cl\leq38Cl^3\leq y$,
$\log Q\leq Clx^{\varepsilon/8}x^{\varepsilon/3}\leq x^{11\varepsilon/24}y$.
\end{proof}

\begin{lemma}\label{pd-gain-holds-hypothesis-two}
Let $0<\varepsilon<1/52$ and $J\geq1$, let $K_G$ be a list of nonnegative constants as in
Definition~\ref{pd-gain-data}, and let $C_{31}>0$. There is $x_0$ such that the following holds for
every $x\geq x_0$. Let $D$ be a definite-side datum with size data with scale $(x,\varepsilon,J)$ that
satisfy the size bounds chosen with families (Definition~\ref{pd-family}(4)), and gain data with
constants $K_G$ that satisfy condition (1) of Definition~\ref{pd-gain-admissibility} with $C_{31}$.
Suppose that the number $\Lambda_G$ of the size data is $\Lambda=K_\Lambda N_1^2\log a$ and that
$\rho_0\leq x^{-2}$, with $\rho_0$ as in Lemma~\ref{pd-gain-first-hypotheses}. Let $p<p_0$ be a
candidate prime and $t_0$ a good representative at $p$, and put $T=P(p,t_0)$, $n'=n_p$,
$a_w(i)=\mathbf a(t_0+ih_p)$ and $b_w(i)=\mathbf b(t_0+ih_p)$. Then, in the notation of
Lemma~\ref{pd-gain-second-hypotheses}: $13\leq p$; $2L_1(p)\leq L_2<u$, $3G'\leq n'$ and
$|T|\leq40n'\delta_1$; $a_w$ and $b_w$ take values in $W_a$ and $W_b$ on $T$ and satisfy the two
inequalities displayed there for $s\in T$; and $r\geq2$, $N_1^2\leq r|T|$, $r\leq x^{26}$,
$k^2x^{78}\leq p^{L_1(p)-1}$, $256k^2|T|x^{52}\leq k^2x^{78}$, $N_1+3r\leq k$, $K'_\kappa N_1^2\log
a\leq k$ and the last inequality displayed there hold.
\end{lemma}
\leanref{Formalism/SharpRestriction/GrowthBounds/SizesAtRecord.lean}{Formalism.MainNotions.G.Joints.hypIIAt_ev}
\leanref{Formalism/SharpRestriction/GrowthBounds/SizesAtRecord.lean}{Formalism.MainNotions.G.Joints.numII_ev}
\leanref{Formalism/SharpRestriction/GrowthBounds/Joints.lean}{Formalism.MainNotions.G.Joints.HypIIAt}

\begin{lemma}\label{pd-gain-holds-band-sums}
Let $X\geq1$, $0\leq\rho\leq1/3$ and $P_0\geq1$ be real numbers, and let $S$ be a finite set of primes
$p$ with $P_0\leq p<X$ such that $X(1-\rho)\leq p^d\leq X(1+\rho)$ for some integer $d\geq1$.
\begin{enumerate}
\item If one $d$ serves for every $p\in S$, then $\sum_{p\in S}\frac{\log p}p\leq\log
  X\,(3\rho+1/P_0)$.
\item In general, $\sum_{p\in S}\frac{\log p}p\leq(1+2\log X)\log X\,(3\rho+1/P_0)$.
\item Let $0<\varepsilon<1/52$ and $x\geq2$ with $12\leq x^{1-2\varepsilon/3}$, and let
  $X=\exp(x^{\varepsilon/3})$, $P_0\geq x^2$ and $r\leq x^{-2}$. Then the sum of $\log p/p$ over the
  primes $p$ with $P_0\leq p<X$ and $X(1-r)\leq p^d\leq X(1+r)$ for some integer $d\geq1$ is at most
  $x^{-1}$.
\end{enumerate}
\end{lemma}
\leanref{Formalism/SharpRestriction/GrowthBounds/TopChoiceBounds.lean}{Formalism.MainNotions.G.TopChoiceBounds.prime_power_band_log_sum_le}
\leanref{Formalism/SharpRestriction/GrowthBounds/TopChoiceBounds.lean}{Formalism.MainNotions.G.TopChoiceBounds.prime_power_band_log_sum_le_log_sq}
\leanref{Formalism/SharpRestriction/GrowthBounds/TopChoiceBounds.lean}{Formalism.MainNotions.G.TopChoiceBounds.prime_power_band_sum_at_scale}

\begin{proof}
(1) We may assume $S\neq\emptyset$. Let $a$ and $b$ be the least and the largest element of $S$.
Then $b/a\leq(b/a)^d\leq X(1+\rho)/(X(1-\rho))=(1+\rho)/(1-\rho)\leq1+3\rho$, as $\rho\leq1/3$. So
$S\subset[a,(1+3\rho)a]$ has at most $3\rho a+1$ elements, and each term is at most $\log X/a$. Hence
the sum is at most $(3\rho a+1)\log X/a=\log X\,(3\rho+1/a)\leq\log X\,(3\rho+1/P_0)$.

(2) For $p\in S$ choose $d(p)\geq1$ as in the hypothesis. Then $p^{d(p)}\leq2X$ and $p\geq2$ give
$(d(p)-1)\log2\leq\log X$, and $\log2\geq1/2$, so $d(p)\leq1+2\log X$. Applying (1) to each of the
at most $1+2\log X$ sets $\{p\in S:d(p)=d\}$ gives (2).

(3) If $r<0$, no prime satisfies $X(1-r)\leq p^d\leq X(1+r)$. Otherwise $0\leq r\leq x^{-2}\leq1/4$,
and $1/P_0\leq x^{-2}$. Put $B=\log X=x^{\varepsilon/3}\geq1$. By (2) the sum is at most
$(1+2B)B(3r+1/P_0)\leq3B\cdot B\cdot4x^{-2}=12x^{2\varepsilon/3-2}\leq x^{1-2\varepsilon/3}
x^{2\varepsilon/3-2}=x^{-1}$.
\end{proof}

\begin{lemma}\label{pd-kappa-orbit-large-scale}
Let $0<\varepsilon<1/52$ and $J\geq1$. There is $x_0$ such that the following holds for every
$x\geq x_0$. Let $D$ be a definite-side datum with rank $k$ whose host has splitting bound $B$
(Definition~\ref{hf-host}), with size data with scale $(x,\varepsilon,J)$
(Definition~\ref{pd-size-data}) such that $k\leq x^{9/2+10\varepsilon+\varepsilon/100}$. Let
$u$ and $G'$ be as in Definition~\ref{hf-good-orbit}, $L_2$ as in Definition~\ref{pd-gain-data}(b), and
$L_1(p)=\lceil\log_p((kx^{26})^4)\rceil$ for primes $p$. Then:
\begin{enumerate}
\item every prime of inert class (Definition~\ref{pd-size-data}) is larger than
  $x^{\varepsilon/(5J)}$;
\item $L_2<u$, and $2L_1(p)\leq L_2$ for every prime $p$ of inert class;
\item $3G'\leq n'$ for every integer $n'\geq x^{\varepsilon/4}$.
\end{enumerate}
\end{lemma}
\leanref{Formalism/SharpRestriction/GrowthBounds/CornerGLContinuity.lean}{Formalism.MainNotions.CornerGLContinuity.LargeScaleConditions}
\leanref{Formalism/SharpRestriction/GrowthBounds/CornerGLContinuity.lean}{Formalism.MainNotions.CornerGLContinuity.EstimatedPrimesExceedScale}
\leanref{Formalism/SharpRestriction/GrowthBounds/TopChoiceBounds.lean}{Formalism.MainNotions.G.TopChoiceBounds.family_large_scale_conditions_hold}
\leanref{Formalism/SharpRestriction/GrowthBounds/TopChoiceBounds.lean}{Formalism.MainNotions.G.TopChoiceBounds.family_large_scale_of_master}
\leanref{Formalism/SharpRestriction/GrowthBounds/TopChoiceBounds.lean}{Formalism.MainNotions.G.TopChoiceBounds.inert_primes_lower_bound}
\leanref{Formalism/SharpRestriction/GrowthBounds/TopChoiceBounds.lean}{Formalism.MainNotions.G.TopChoiceBounds.window_length_le_ceil_at_scale}
\leanref{Formalism/SharpRestriction/GrowthBounds/TopChoiceBounds.lean}{Formalism.MainNotions.G.TopChoiceBounds.window_length_lt_tile_length}
\leanref{Formalism/SharpRestriction/GrowthBounds/TopChoiceBounds.lean}{Formalism.MainNotions.G.TopChoiceBounds.two_mul_first_length_le_window}
\leanref{Formalism/SharpRestriction/GrowthBounds/TopChoiceBounds.lean}{Formalism.MainNotions.G.TopChoiceBounds.period_ge_three_gap_eventually}

\begin{proof}
Let $B_\varepsilon$ be the constant of Lemma~\ref{pd-gain-holds-orbit-sizes}, and choose $x_0\geq3$
with $2\lceil620J/\varepsilon\rceil+1\leq x_0^{\varepsilon/8}$ and
$48B_\varepsilon/\varepsilon\leq x_0^{\varepsilon/16}$. Let $x\geq x_0$ and $D$ be as in the lemma.

(1) Let $p$ be a prime of inert class. Then $-d$ is not a square modulo $p$. Every residue modulo $2$
is a square, so $p$ is odd, and $p\neq d$, so $\left(\frac{-d}p\right)\neq1$. By
Definition~\ref{hf-host-field-rows}(4), $p>B$, and $B\geq\lfloor x^{\varepsilon/(5J)}\rfloor$ by
condition (10) of Definition~\ref{pd-size-data}; so
$p\geq\lfloor x^{\varepsilon/(5J)}\rfloor+1>x^{\varepsilon/(5J)}$.

(2) Since $\varepsilon<1/52$, $k\leq x^5$, so $\log((kx^{26})^4)\leq124\log x$. For a prime $q$ of
inert class, (1) gives $\log q>\frac{\varepsilon}{5J}\log x>0$, so
$L_1(q)\leq\lceil124\log x/\log q\rceil\leq\lceil620J/\varepsilon\rceil$. Hence
$L_2\leq2\lceil620J/\varepsilon\rceil<x^{\varepsilon/8}=1/\delta_1\leq u$. Since $L_2$ is twice the
largest $L_1(q)$ over the primes $q$ of inert class, $2L_1(p)\leq L_2$ for each of them.

(3) Put $y=x^{\varepsilon/16}$, so that $\delta_1=x^{-\varepsilon/8}=y^{-2}$ and $0<\delta_1\leq1$.
By Lemma~\ref{pd-gain-holds-orbit-sizes}(1), $G'\delta_1\leq B_\varepsilon\log x$, so
$G'\leq B_\varepsilon\log x\cdot y^2$. Since $\log t\leq t^s/s$ for $t\geq1$ and $s>0$,
$\log x\leq(16/\varepsilon)y$; and $48B_\varepsilon/\varepsilon\leq x_0^{\varepsilon/16}\leq y$. So
$3B_\varepsilon\log x\leq(48B_\varepsilon/\varepsilon)y\leq y^2$, and
$3G'\leq3B_\varepsilon\log x\cdot y^2\leq y^4=x^{\varepsilon/4}\leq n'$.
\end{proof}

\begin{lemma}\label{pd-gain-holds-positions}
Let $x>1$ and $0<\varepsilon<1/52$, with $K_\varepsilon$, $r_\varepsilon$, $u$, $M_0$, $G'$, $N_1$, the
windows $W(i,\ell)$, the tiles and the cells $I_{s,c}$ of Definition~\ref{hf-good-orbit}. Let $n'$ and
$L_2\geq0$ be integers with $3G'\leq n'$, and let $P\subset\{0,\dots,n'-1\}$ be a tiled set of positions
for $(x,\varepsilon,L_2)$. For an integer $i$ let $\bar i\in\{0,\dots,n'-1\}$ be its residue modulo
$n'$.
\begin{enumerate}
\item For every $0\leq i<n'$ there is $1\leq D\leq G'$ such that $\overline{i+D}$ is an even element of
  $P$.
\item For every even $y\in P$ there is $1\leq d\leq G'$ such that $\overline{y+d}$ is an even element of
  $P$ and at most $N_1$ of the integers $0<e<d$ have $\overline{y+e}\in P$.
\item Let $L_1$ be an integer with $L_1\leq L_2<u$. For every $i\in P$, fewer than $r_\varepsilon$ of
  the integers $1\leq d\leq L_1$ have $\overline{i-d}$ an even element of $P$.
\item $n'\leq|P|\,G'$ and $|P|\,G'\leq N_1(n'+G')$.
\end{enumerate}
\end{lemma}
\leanref{Formalism/GrowthSums/FactorialBaseGrowth/NextEvenTopHolds.lean}{Formalism.GrowthSums.next_even_top_holds}
\leanref{Formalism/GrowthSums/FactorialBaseGrowth/NextEvenTopHolds.lean}{Formalism.GrowthSums.top_count_holds}
\leanref{Formalism/SharpRestriction/GrowthBounds/TopChoiceBounds.lean}{Formalism.MainNotions.G.TopChoiceBounds.bc_next}
\leanref{Formalism/SharpRestriction/GrowthBounds/TopChoiceBounds.lean}{Formalism.MainNotions.G.TopChoiceBounds.bc_old}
\leanref{Formalism/SharpRestriction/GrowthBounds/TopChoiceBounds.lean}{Formalism.MainNotions.G.TopChoiceBounds.bc_count}
\leanref{Formalism/SharpRestriction/GrowthBounds/CornerGLContinuity.lean}{Formalism.MainNotions.CornerGLContinuity.OrbitPositions}
\leanref{Formalism/SharpRestriction/GrowthBounds/CornerGLContinuity.lean}{Formalism.MainNotions.CornerGLContinuity.ConsecutiveMaximaClose}
\leanref{Formalism/SharpRestriction/GrowthBounds/CornerGLContinuity.lean}{Formalism.MainNotions.CornerGLContinuity.EarliestMaximumFar}

\begin{proof}
(1) Put $g=(6r_\varepsilon+1)u$, the length of a tile. Then the number of tiles is
$S=\lfloor n'/g\rfloor$, $Sg\leq n'<(S+1)g$, and $G'=(M_0+5)g$. As $(M_0+1)g\leq G'\leq n'$,
$S\geq M_0+1$. Let $0\leq i<n'$. Put $s_i=\lfloor i/g\rfloor$ if $i<Sg$ and $s_i=S-1$ otherwise. Then
$0\leq s_i<S$, $s_ig\leq i$, and $i<(s_i+1)g$ unless $s_i=S-1$. By condition (2) of
Definition~\ref{hf-good-orbit}, applied to the residue of $s_i+1$ modulo $S$, there is $0\leq j\leq M_0$
such that the tile $s$, the residue of $s_i+1+j$ modulo $S$, is marked. Its cell $I_{s,2}$ contains an
even element $z$ of $P$, and $sg+u\leq z<sg+2u\leq(s+1)g$.
\begin{itemize}
\item If $s=s_i+1+j$, then $s_i\neq S-1$, so $i<(s_i+1)g\leq sg<z$. Put $D=z-i$; then
  $1\leq D<(s+1)g-s_ig=(j+2)g$.
\item Otherwise $s=s_i+1+j-S$, and $s\leq s_i$ since $j+1\leq S$. If $i<z$, put $D=z-i$; then
  $1\leq D<(s+1)g-s_ig\leq g$. If $z\leq i$, put $D=z+n'-i$; then $D\geq1$ as $i<n'$, and
  $D<(s+1)g+(S+1)g-s_ig=(j+3)g$.
\end{itemize}
In each case $\overline{i+D}=z$ and $1\leq D\leq(M_0+3)g\leq G'$.

(2) Let $y\in P$ be even. By (1) there is $1\leq D\leq G'$ with $\overline{y+D}$ an even element of $P$;
let $d$ be the least integer $d\geq1$ with this property, so $d\leq D\leq G'$. For $0<e<d$,
$\overline{y+e}$ is the residue of $\overline{y+1}+(e-1)$ with $0\leq e-1<G'$, so it lies in the window
$W(\overline{y+1},G')$, which contains at most $N_1$ elements of $P$ by condition (3) of
Definition~\ref{hf-good-orbit}. Since $d\leq G'<n'$, distinct $e$ have distinct residues
$\overline{y+e}$. So at most $N_1$ of the $e$ have $\overline{y+e}\in P$.

(3) Here $L_1\leq L_2<u\leq G'<n'$. Let $i\in P$ and $i_0=\overline{i-L_1}$. For $1\leq d\leq L_1$,
$\overline{i-d}$ is the residue of $i_0+(L_1-d)$ with $0\leq L_1-d<L_1\leq L_2$, so it lies in the window
$W(i_0,L_2)$; and distinct $d$ have distinct residues $\overline{i-d}$, as $L_1<n'$. By condition (1) of
Definition~\ref{hf-good-orbit} at most $K_\varepsilon=r_\varepsilon-1$ elements of $P$ lie in
$W(i_0,L_2)$. So fewer than $r_\varepsilon$ of the $d$ have $\overline{i-d}\in P$.

(4) By (1), every $0\leq i<n'$ is $\overline{z-D}$ for some $z\in P$ and $1\leq D\leq G'$, namely
$z=\overline{i+D}$; so $n'\leq|P|\,G'$. For the second bound, each $z\in P$ lies in the window
$W(cG',G')$ with $c=\lfloor z/G'\rfloor$, since $cG'\leq z<cG'+G'$ and $z<n'$. Here $0\leq cG'<n'$, and
$c$ takes at most $\lfloor n'/G'\rfloor+1$ values. Each of these windows contains at most $N_1$
elements of $P$ by condition (3) of Definition~\ref{hf-good-orbit}, so
$|P|\leq(\lfloor n'/G'\rfloor+1)N_1$. As $\lfloor n'/G'\rfloor G'\leq n'$, this gives
$|P|\,G'\leq N_1(n'+G')$.
\end{proof}

\begin{lemma}\label{pd-gain-holds-orbit-sizes}
Let $0<\varepsilon<1/52$ and $x\geq3$, with $r_\varepsilon$, $h_\varepsilon$, $\delta_1$, $u$, $M_0$, $G'$
and $N_1$ as in Definition~\ref{hf-good-orbit}, and put
$B_\varepsilon=2(6r_\varepsilon+1)\big((h_\varepsilon+3)10^{6r_\varepsilon+1}+6\big)$.
\begin{enumerate}
\item $G'\delta_1\leq B_\varepsilon\log x$ and $N_1\leq16.8\,G'\delta_1$.
\item Let $n'$ and $L_2\geq0$ be integers with $3G'\leq n'$, and let $P\subset\{0,\dots,n'-1\}$ be a
  tiled set of positions for $(x,\varepsilon,L_2)$. Then
  $n'\delta_1/(20B_\varepsilon\log x)\leq|P|\leq40n'\delta_1$.
\end{enumerate}
\end{lemma}
\leanref{Formalism/SharpRestriction/GrowthBounds/TopChoiceBounds.lean}{Formalism.MainNotions.G.TopChoiceBounds.bc_size}
\leanref{Formalism/SharpRestriction/GrowthBounds/TopChoiceBounds.lean}{Formalism.MainNotions.G.TopChoiceBounds.bc_size_real}
\leanref{Formalism/SharpRestriction/GrowthBounds/TopChoiceBounds.lean}{Formalism.MainNotions.G.TopChoiceBounds.logBoundConstant}
\leanref{Formalism/SharpRestriction/GrowthBounds/TopChoiceBounds.lean}{Formalism.MainNotions.G.TopChoiceBounds.segment_count_log_bounds}

\begin{proof}
(1) Put $l=\log x$, which is at least $1$ as $x\geq3$, and $c=(h_\varepsilon+3)10^{6r_\varepsilon+1}$.
Since $0<\delta_1\leq1$ and $u=\lceil1/\delta_1\rceil$, $1\leq u\delta_1<1+\delta_1\leq2$. Since
$M_0=\lceil cl\rceil$, $cl\leq M_0<cl+1$, so $M_0+5\leq(c+6)l$. Hence
$G'\delta_1=(M_0+5)(6r_\varepsilon+1)u\delta_1\leq2(6r_\varepsilon+1)(c+6)l=B_\varepsilon l$. Also
$G'\delta_1\geq M_0+5\geq5$ and $M_0\geq cl\geq(h_\varepsilon+3)l$. As $\log2>1/2$,
$(h_\varepsilon+3)\log_2x\leq2(h_\varepsilon+3)l\leq2M_0\leq6e\,G'\delta_1$, so the maximum in the
definition of $N_1$ is $6e\,G'\delta_1$, and $N_1<6e\,G'\delta_1+2$. As $6e<16.4$ and
$2\leq0.4\,G'\delta_1$, $N_1\leq16.8\,G'\delta_1$.

(2) By Lemma~\ref{pd-gain-holds-positions}(4), $n'\leq|P|\,G'$ and $|P|\,G'\leq N_1(n'+G')$. The first
and (1) give $n'\delta_1\leq|P|\,G'\delta_1\leq B_\varepsilon\log x\,|P|\leq20B_\varepsilon\log x\,|P|$.
By the second, (1) and $n'+G'\leq\frac43n'$,
$|P|\,G'\leq16.8\,G'\delta_1\cdot\frac43n'=22.4\,n'\delta_1G'$; dividing by $G'>0$ gives
$|P|\leq22.4\,n'\delta_1\leq40n'\delta_1$.
\end{proof}

\begin{lemma}\label{pd-gain-holds-orbit-condition}
Let $D$ be a definite-side datum with size data with scale $(x,\varepsilon,J)$, $x\geq3$, and gain data
(Definition~\ref{pd-gain-data}), and let $B_\varepsilon$ be as in
Lemma~\ref{pd-gain-holds-orbit-sizes}. Then condition (1) of Definition~\ref{pd-gain-admissibility} holds
with every $C_{31}\geq20B_\varepsilon$.
\end{lemma}
\leanref{Formalism/SharpRestriction/GrowthBounds/TopChoiceBounds.lean}{Formalism.MainNotions.G.TopChoiceBounds.uniform_constant_of_master}
\leanref{Formalism/SharpRestriction/GrowthBounds/TopChoiceBounds.lean}{Formalism.MainNotions.G.TopChoiceBounds.family_uniform_constant_of_master}
\leanref{Formalism/SharpRestriction/GrowthBounds/TopChoiceBounds.lean}{Formalism.MainNotions.G.TopChoiceBounds.mem_tops_on_orbit_iff}
\leanref{Formalism/SharpRestriction/GrowthBounds/TopChoiceBounds.lean}{Formalism.MainNotions.G.TopChoiceBounds.isV_iff_even_index}
\leanref{Formalism/SharpRestriction/GrowthBounds/TopChoiceBounds.lean}{Formalism.MainNotions.G.TopChoiceBounds.vertex_count_eq_card_tops}

\begin{proof}
Let $p$ be a prime of inert class with $n_p\geq x^{\varepsilon/4}$, $L_2<u$ and $3G'\leq n_p$, and let
$t_0$ be a good representative at $p$. Write $h=h_p$, $n=n_p$ and $P=P(p,t_0)$. As $p$ is not a signal
prime, $s_p=h$, and the orbit length $\ell_p$ of Definition~\ref{pd-orbit-places} is $n$.

\emph{Step 1: the top classes with representative $t_0$.} By Definition~\ref{pd-gain-data}(1),
$(t+ih)^h=t^h$ for every $t\in A$ and $i\geq0$, and $t=t^h+ih$ for some $i\geq0$; so $(t^h)^h=t^h$. As
$t_0=t_1^h$ for some $t_1$, $t_0^h=t_0$. So a top class $s$ has representative $t_0$ if and only if
$s=t_0+ih$ for some $0\leq i<n$: if $s^h=t_0$, then $s=t_0+ih$ with $i\geq0$, and $i$ may be reduced
modulo the order $n$ of $h$; conversely $(t_0+ih)^h=t_0^h=t_0$. The integer $0\leq i<n$ is unique, and
$t_0+ih\in T$ if and only if $i\in P$. Hence $s\mapsto i$ is a bijection from the top classes with
representative $t_0$ onto $P$, so $N(p,t_0)=|P|$, and $t_0+ih$ is even at $p$ if and only if $i$ is
even. For an integer $i$ write $\bar i$ for its residue modulo $n$; then $t_0+ih=t_0+\bar ih$.

\emph{Step 2: sizes.} As $t_0$ is good, $P$ is a tiled set of positions for $(x,\varepsilon,L_2)$, and
$3G'\leq n$. By Lemma~\ref{pd-gain-holds-orbit-sizes}(2),
$n\delta_1/(20B_\varepsilon\log x)\leq N(p,t_0)\leq40n\delta_1$, and
$n\delta_1/(C_{31}\log x)\leq n\delta_1/(20B_\varepsilon\log x)$ since $\log x>0$.

\emph{Step 3: the next even top class.} Let $s$ be a top class with representative $t_0$ that is even at
$p$; then $s=t_0+yh$ with $y\in P$ even. Lemma~\ref{pd-gain-holds-positions}(2) gives $1\leq d\leq G'$
such that $\overline{y+d}$ is an even element of $P$ and at most $N_1$ of the integers $0<e<d$ have
$\overline{y+e}\in P$. By Step 1, $s+dh=t_0+\overline{y+d}\,h$ lies in $T$ and is even at $p$, and for
$0<e<d$, $s+eh\in T$ if and only if $\overline{y+e}\in P$.

\emph{Step 4: the earlier even top classes.} Let $s=t_0+yh$ with $y\in P$. Since $p$ is of inert class,
$L_1(p)\leq L_2$ by the definition of $L_2$ (Definition~\ref{pd-gain-data}(b)), so $L_1(p)\leq L_2<u$.
For $1\leq d\leq L_1(p)$, $s-dh=t_0+\overline{y-d}\,h$, which by Step 1 lies in $T$ and is even at $p$
if and only if $\overline{y-d}$ is an even element of $P$. By Lemma~\ref{pd-gain-holds-positions}(3),
with $L_1=L_1(p)$, fewer than $r_\varepsilon$ of these $d$ have this property.
\end{proof}

\begin{lemma}\label{pd-gain-holds-two-adic}
Let $e\geq0$ and $0\leq j\leq e$ be integers, $l$ an odd positive integer, and
$H_e=(\mathbf Z/2^{e+2}\mathbf Z)^\times/\{\pm1\}$ (Definition~\ref{hf-frame}). Then the class of $l$ in
$H_e$ has order at most $2^j$ if and only if $l\equiv\pm1\pmod{2^{e+2-j}}$.
\end{lemma}
\leanref{Formalism/SharpRestriction/GlobalSchurBounds/GlobalRankLatticeProofs/GenericFramesLetters.lean}{Formalism.MainNotions.UnusedSlot33.PrimeReciprocalSums.ord_le_iff}
\leanref{Formalism/SharpRestriction/GlobalSchurBounds/GlobalRankLatticeProofs/GenericFramesLetters.lean}{Formalism.MainNotions.UnusedSlot33.PrimeReciprocalSums.nsmul_iff}
\leanref{Formalism/SharpRestriction/GlobalSchurBounds/GlobalRankLatticeProofs/GenericFramesLetters.lean}{Formalism.MainNotions.UnusedSlot33.PrimeReciprocalSums.hOf_ord}

\begin{proof}
As $H_e$ has order $2^e$, the order of the class of $l$ is $2^k$ with $k\leq e$, and it is at most $2^j$
if and only if it divides $2^j$, that is, $l^{2^j}\equiv\pm1\pmod{2^{e+2}}$. Let $z$ be an odd integer
and $N\geq1$. Then $2^N$ divides $z^2-1$ or $z^2+1$ if and only if $2^{N-1}$ divides $z-1$ or $z+1$. For
$N=1$ both sides hold. For $N\geq2$ write $z=2c+1$. Then $z^2+1=4c(c+1)+2$ is not divisible by $4$, and
$z^2-1=4c(c+1)$, where one of $c$ and $c+1$ is odd; so $2^N\mid z^2-1$ if and only if $2^{N-2}$ divides
$c$ or $c+1$, that is, $2^{N-1}$ divides $z-1=2c$ or $z+1=2(c+1)$. Applying this to
$z=l^{2^{j-1}},l^{2^{j-2}},\dots,l$ with $N=e+2,e+1,\dots,e+3-j$ in turn shows that
$l^{2^j}\equiv\pm1\pmod{2^{e+2}}$ if and only if $l\equiv\pm1\pmod{2^{e+2-j}}$.
\end{proof}

\begin{lemma}\label{pd-gain-holds-order-sums}
Let $e\geq2$ be an integer, $J\geq1$ and $Y\geq1$ real numbers, and $H_e$ as in
Definition~\ref{hf-frame}. Then
\[
 \sum_\ell\frac{\log\ell}{\ell}\leq\frac{100J\big(\log(2^eY)\big)^2}{2^e},
\]
the sum over the odd primes $\ell\leq Y$ whose class in $H_e$ has order at most $J$.
\end{lemma}
\leanref{Formalism/SharpRestriction/GlobalSchurBounds/CMFieldLemmas.lean}{Formalism.MainNotions.UnusedSlot33.Conclusion4}
\leanref{Formalism/SharpRestriction/GlobalSchurBounds/GlobalRankLatticeProofs/GenericFramesLetters.lean}{Formalism.MainNotions.UnusedSlot33.generic_frames_letters}
\leanref{Formalism/SharpRestriction/GlobalSchurBounds/GlobalRankLatticeProofs/GenericFramesLetters.lean}{Formalism.MainNotions.UnusedSlot33.PrimeReciprocalSums.modulus}
\leanref{Formalism/SharpRestriction/GlobalSchurBounds/GlobalRankLatticeProofs/GenericFramesLetters.lean}{Formalism.MainNotions.UnusedSlot33.PrimeReciprocalSums.sum_lin_bound}
\leanref{Formalism/SharpRestriction/GlobalSchurBounds/GlobalRankLatticeProofs/GenericFramesLetters.lean}{Formalism.MainNotions.UnusedSlot33.PrimeReciprocalSums.sum_inv_le}
\leanref{Formalism/SharpRestriction/GlobalSchurBounds/GlobalRankLatticeProofs/GenericFramesLetters.lean}{Formalism.MainNotions.UnusedSlot33.PrimeReciprocalSums.linear_term_le}

\begin{proof}
Put $n=2^e\geq4$. Let $K$ be the largest integer $0\leq K\leq e$ with $2^K\leq J$, which exists as
$J\geq1$, and put $m=2^{e+2-K}$. Then $4\leq m\leq4n$ and $mJ\geq m2^K=4n$. Let $\ell$ be an odd prime
whose class in $H_e$ has order at most $J$. As $H_e$ has order $2^e$, this order is $2^k$ with $k\leq e$
and $2^k\leq J$, so $k\leq K$; and Lemma~\ref{pd-gain-holds-two-adic}, with $j=K$, gives
$\ell\equiv\pm1\pmod m$.

If $Y<3$, there is no odd prime $\ell\leq Y$, and the bound holds as its right side is $\geq0$. Let
$Y\geq3$ and $N=\lfloor Y\rfloor$. Each $\ell$ in the sum is $mq+1$ or $mq-1$ for an integer
$1\leq q\leq N$: take $q=(\ell-1)/m$ or $q=(\ell+1)/m$, which is $\geq1$ as $\ell\geq3$, and at most
$\ell\leq N$. As $mq\geq4$, $\ell\geq mq-1\geq\frac34mq$, and $\log\ell\leq\log Y$, so
$\log\ell/\ell\leq\frac43\cdot\frac{\log Y}{m}\cdot\frac1q$. Distinct $\ell$ of the same form give
distinct $q$. With $\sum_{q=1}^N1/q\leq1+\log N\leq1+\log Y$, the sum is at most
\[
 2\cdot\frac43\cdot\frac{\log Y}{m}\sum_{q=1}^N\frac1q\leq\frac83\cdot\frac{\log Y\,(1+\log Y)}{m}.
\]
Put $L=\log(nY)=\log n+\log Y$. As $\log n\geq2\log2>1$, $L\geq1$ and $L\geq\log Y\geq0$, so
$\log Y\,(1+\log Y)\leq2L^2$. With $1/m\leq J/(4n)$ the sum is at most
$\frac{16}3L^2\cdot\frac J{4n}=\frac43\cdot\frac{JL^2}n\leq\frac{100JL^2}n$.
\end{proof}

\begin{lemma}\label{pd-gain-holds-small-orders}
Let $D$ be a definite-side datum whose Galois datum has $2$-power exponent $e$, group $A$ of order $n$
and index $o$, with size data with scale $(x,\varepsilon,J)$ and the notation of
Definitions~\ref{pd-size-data} and~\ref{pd-gain-data}, and put $h_\varepsilon=2+8\varepsilon$. Suppose
that $x^{2+8\varepsilon-\varepsilon/100}\leq n\leq x^{2+8\varepsilon+\varepsilon/100}$,
$o\leq x^{\varepsilon/100}$, $8\log x\leq x^{\varepsilon/3}$ and $400\leq x^{19\varepsilon/300}$. Then
the sum of $\log p/p$ over the primes $p<X''$ of inert class with $n_p<x^{\varepsilon/4}$ is at most
$x^{-h_\varepsilon+\varepsilon}$; that is, condition (3) of Definition~\ref{pd-gain-admissibility}
holds.
\end{lemma}
\leanref{Formalism/SharpRestriction/GrowthBounds/TopChoiceBounds.lean}{Formalism.MainNotions.G.TopChoiceBounds.order_statement_at_scale}
\leanref{Formalism/SharpRestriction/GrowthBounds/TopChoiceBounds.lean}{Formalism.MainNotions.G.TopChoiceBounds.order_statement_of_odd_primes}
\leanref{Formalism/SharpRestriction/GrowthBounds/TopChoiceBounds.lean}{Formalism.MainNotions.G.TopChoiceBounds.family_order_bound_holds}
\leanref{Formalism/SharpRestriction/GrowthBounds/PlaneCountTheoremWithConstants.lean}{Formalism.MainNotions.G.family_orders_first_part}
\leanref{Formalism/SharpRestriction/GrowthBounds/CornerGLContinuity.lean}{Formalism.MainNotions.CornerGLContinuity.OrdersSmallPrimeBound}
\leanref{Formalism/SharpRestriction/AdmissibleParameterJoins/GandSk.lean}{Formalism.MainNotions.J.two_readings_of_Frobenius}

\begin{proof}
\emph{Step 1: the primes.} The host field is on a Galois datum with $e\geq2$
(Definition~\ref{hf-framed-host-field}). Let $p$ be a prime of inert class. Then $p\nmid Q_S$ and
$p\neq q,d$ (Definition~\ref{pd-size-data}), so $p$ is prime to $Q_Sqd$; and $2^{e+2}\mid Q_S$, so $p$
is odd. The element $h_p$ is the $A$-component of $\operatorname{art}(p)$, which $h$ maps to the class
of $p$ in $H_e$ (Definition~\ref{hf-framed-host-field}(4)). As $h$ is a homomorphism, the order of the
class of $p$ divides $n_p$. So the sum in question is at most the sum of $\log p/p$ over the odd primes
$p\leq X''$ whose class in $H_e$ has order at most $x^{\varepsilon/4}\geq1$, and
Lemma~\ref{pd-gain-holds-order-sums}, with $J=x^{\varepsilon/4}$ and $Y=X''$, bounds it by
$100x^{\varepsilon/4}\big(\log(2^eX'')\big)^2/2^e$.

\emph{Step 2: the sizes.} As $h$ maps $A$ onto $H_e$, $2^e$ divides $n$, so $2^e\leq n$. As $h$ is
injective on $A_0$, the order $n_0$ of $a_0$ is at most $2^e$, and $n=on_0\leq o2^e$. Hence
$2^e\geq n/o\geq x^{2+8\varepsilon-\varepsilon/50}$. Also $X''=xnX'$ with $\log X'=x^{\varepsilon/3}$,
and $\log n\leq(2+8\varepsilon+\varepsilon/100)\log x\leq3\log x$, as $\varepsilon<1/52$. So
\[
 0\leq\log(2^eX'')=\log2^e+\log x+\log n+x^{\varepsilon/3}\leq7\log x+x^{\varepsilon/3}
 \leq2x^{\varepsilon/3}.
\]
The sum is therefore at most
\[
 \frac{100x^{\varepsilon/4}\cdot4x^{2\varepsilon/3}}{x^{2+8\varepsilon-\varepsilon/50}}
 \leq x^{19\varepsilon/300}\,x^{\varepsilon/4+2\varepsilon/3+\varepsilon/50-2-8\varepsilon}
 =x^{-2-7\varepsilon}=x^{-h_\varepsilon+\varepsilon}.
\]
\end{proof}

\begin{theorem}\label{pd-gain-link-holds}
Let $\mathcal A$ be a class of families with a twist selection, every family of which carries its
duty data (Definition~\ref{asm-duty}). Then the gain link (Definition~\ref{pd-gain-link}) holds
eventually along $\mathcal D_3=\operatorname{Des}(\mathcal A_N,3)$
(Definitions~\ref{pd-family} and~\ref{asm-designed}).
\end{theorem}
\leanref{Development/AutomorphicArithmetic/ThetaCorrespondence/ThetaPairLifts/ThirdPlaceSeries/LinkG.lean}{GlobalTheory.Top.closed_linkG_of_duty}
\leanref{Formalism/LocalToGlobal/GroupLattices/LinkGConcreteNoKer.lean}{Formalism.L2.G.linkG_concrete5}
\leanref{Formalism/SharpRestriction/SmoothDesignExistence/AdmissibleTopsMertens.lean}{Formalism.MainNotions.Pk.admissible_tops_mertens}
\leanref{Formalism/SharpRestriction/GrowthBounds/SmallOrderPrimes.lean}{Formalism.MainNotions.AuxBlock8.adm6}
\leanref{Formalism/SharpRestriction/GrowthBounds/SmallOrderPrimes.lean}{Formalism.MainNotions.AuxBlock8.adm7}
\leanref{Formalism/SharpRestriction/GrowthBounds/SmallOrderPrimes.lean}{Formalism.MainNotions.AuxBlock8.adm8}
\leanref{Formalism/SharpRestriction/GrowthBounds/DefsInst.lean}{Formalism.MainNotions.G.KapIDef}
\leanref{Formalism/LocalToGlobal/GroupLattices/AdmissibleFamilyNoKernel.lean}{Formalism.L2.G.linkG_fam_noKer}
\leanref{Formalism/LocalToGlobal/GroupLattices/AdmissibleFamilyNoKernel.lean}{Formalism.L2.G.linkG_without_kernels}
\leanref{Formalism/LocalToGlobal/GroupLattices/AdmissibleFamilyNoKernel.lean}{Formalism.L2.G.admissible_family_no_kernel}
\leanref{Formalism/SharpRestriction/GrowthBounds/PlaneCountTheoremWithConstants.lean}{Formalism.MainNotions.G.linkG_at_core}

\begin{proof}
The proof has three steps.

\emph{Step 1: group coherence.} The duty data and the twist selection make every family of
$\mathcal A$ a linked and aligned family of the second layer. At all large scales, every member of
such a family carries group coherence (Definition~\ref{pd-gain-coherence}) for its local models,
its tie and the comparison of its lattice instance. By
Lemma~\ref{pd-gain-coherence-consequences}, at all large scales: the groups of the components have
trivial stabilizers and finitely many classes of lattice tuples; the Hida settings of the local
models come from unramified quadratic settings; and, since the submodules $V_{\mathrm{ram}}$ of the
lattice instances of the family are the $\mathcal K$-spans of the images of the
$M_{u,\mathrm{ram}}$, the served local lattices descend.

\emph{Step 2: admissible tops.} The tops are admissible at all large scales. Let $P$ be the set
of primes $p<X'$ of inert class whose Frobenius class has $A$-component $h_p$ of order at least $x^{\varepsilon/4}$, and let
$B$ be the sum over $p\in P$ of $\frac{\log p}{p}$ times the number of top classes $t$ whose
representative $t^{s_p}$ is not good at $p$ (Definition~\ref{pd-gain-data}(b)). Admissibility asks that
$n\delta_1/2\leq|T|\leq2n\delta_1$ and $B\leq|T|x^{-2+\varepsilon}$. The bounds on $|T|$ and the
bound
\[
 B\leq4C_bx^{-2}\delta_1n\sum_{p\in P}\frac{\log p}{p},
\]
for some $C_b>0$, are among the choices of constants of a designed family. Put
$y=x^{\varepsilon/3}$, so that $X'=e^{y}$. Each $p\in P$ has $\log p\leq y$, and summing over
integers gives $\sum_{p\in P}\log p/p\leq2y(y+2)$. Once $x\geq(16C_b+2)^{3/\varepsilon}$ we
have $y\geq16C_b+2$, hence $16C_b(y+2)\leq y^2$ and
$B\leq\tfrac12x^{-2}\delta_1ny^3\leq|T|x^{-2}y^3=|T|x^{-2+\varepsilon}$.

\emph{Step 3: the local theory.} Let $C_0$ be the constant of Theorem~\ref{pd-gain-theorem-two}
and $T_0,T_2,C$ the functions of Theorem~\ref{pd-gain-theorem-one}, and name the constants of $K_G$
as in Lemmas~\ref{pd-gain-second-hypotheses} and~\ref{pd-gain-first-hypotheses}. A largeness
threshold can be chosen, for every exponent, so that every large family satisfies: $C_0$ is at most
$K_\kappa,K_t,K_\Lambda,K_D,K_s,K'_\kappa$; $K_\Lambda\geq T_0$ and $K_\Lambda>0$; with
$C_H=T_2(K_\Lambda)$ and $C_*=C(K_\Lambda)$, $C_H\leq K_{\mathrm I}$, $C_H\leq K_{D,\mathrm I}$,
$2C_H\leq K_\rho$, $2C_*\leq K_1$ and $C_*\leq K_2$; $K_N\geq1$; and $N_1\leq K_N\log x$ at all
large scales. The family's constant $C_e$ is $1$. Fix such a family. At every large scale and member
the following hold. They are listed in the order in which the proof combines them.
\begin{enumerate}
\item The \emph{second family of bounds on the gains}, condition (6) of
  Definition~\ref{pd-gain-admissibility}. This is Lemma~\ref{pd-gain-second-hypotheses}; its
  hypotheses $c_G\geq0$ and $a-c_G\geq1$ follow from the parameter choices of the family.
\item The local constants: the gain data of the family define $\kappa_I(p,n',T,a',b'_w)$ as the
  least $\kappa\geq0$ in the set of Lemma~\ref{pd-gain-first-hypotheses}. The condition defining this
  set depends on $\kappa$ only through $\lceil(1-\kappa)\dim V\rceil$ and holds for $\kappa=1$ (with
  no edges and $n=0$), so the least value is attained: the projection property holds with
  $\kappa=\kappa_I$, which is condition (7) of Definition~\ref{pd-gain-admissibility}.
\item The \emph{first family of bounds on the gains}, condition (8) with $C_m=C_*$, by
  Lemma~\ref{pd-gain-first-hypotheses} with $C_H$ and $C=C_*$, using (2).
\item The \emph{Weyl sandwich}, condition (9) with $C_e=1$, by Lemma~\ref{pd-gain-weyl-sandwich}:
  the Weyl exponents of the family are $e_p=\lceil C_efk^2\log_pa\rceil$, and the parameter choices
  and the tie give $a\geq4$, $5\leq k\leq a$ and $\lambda(v,s)_1-\lambda(v,s)_k\leq3a$.
\item \emph{Lattice bounds on the orbits}, the conclusion of Lemma~\ref{pd-gain-orbit-bounds}: its
  hypotheses are (1), (2), group coherence (Step 1), $\rho_0\leq\rho$, which follows from
  $N_1\leq K_N\log x$ by comparing the two definitions, and the hypotheses (I) and (II) at the
  candidate primes and good representatives. For (I), let $p\geq p_0$ be a candidate prime and $t_0$
  good at $p$. The members of the family have size data with scale $(x,\varepsilon,J)$ for one
  $J\geq1$, and the size bounds of the family give $k\leq x^{9/2+10\varepsilon+\varepsilon/100}$. So at all
  large scales Lemma~\ref{pd-kappa-orbit-large-scale} gives $L_2<u$, and, as $n_p\geq
  x^{\varepsilon/4}$ for a candidate prime, $3G'\leq n_p$. Hypothesis~I holds at $(p,n_p,N(p,t_0))$ by
  Lemma~\ref{pd-gain-holds-hypothesis-one}, whose hypotheses are conditions (1), (4), (5) with
  $\delta=\varepsilon/48$ and (12) (item (7) below and the end of this step) and the conclusions of
  Lemma~\ref{pd-gain-holds-numerics}; and the exponents take values in $W_a$ and $W_b$ at the top
  classes, since $a\leq a_t\leq a+a/(kx^{13})$ and $b_t\geq b'$ by the size data and $b_t\leq
  b'+N^*x^\varepsilon\Lambda_Gb'/k$ by the parameter choices, with $\Lambda_G=\Lambda$ among the
  choices of the family. For (II), Lemma~\ref{pd-gain-holds-hypothesis-two} applies at all large
  scales with condition (1), $\Lambda_G=\Lambda$, and $\rho_0\leq\rho\leq x^{-2}$, the second
  inequality being a choice of the family.
\item A \emph{lattice sandwich}: $p_w^{e_{p_w}}M^\nabla\subset M\subset M^\nabla$ at every $(u,w)$,
  the second part of (4).
\item Condition (1) of Definition~\ref{pd-gain-admissibility} holds at all large scales with
  $C_{31}=20B_\varepsilon$, where $B_\varepsilon$ is the constant of
  Lemma~\ref{pd-gain-holds-orbit-sizes}, by Lemma~\ref{pd-gain-holds-orbit-condition}, since the
  scale of the members is $(x,\varepsilon,J)$ with $x\geq3$. Condition (3) holds at all large scales
  by Lemma~\ref{pd-gain-holds-small-orders}: the size bounds of the family give
  $x^{2+8\varepsilon-\varepsilon/100}\leq n\leq x^{2+8\varepsilon+\varepsilon/100}$, a condition of
  the size data gives $o\leq x^{\varepsilon/100}$, and $8\log x\leq x^{\varepsilon/3}$ and
  $400\leq x^{19\varepsilon/300}$ once $x$ is large. Condition (4): at all large scales
  $\rho_0\leq\rho$, as in item (5), and $\rho\leq x^{-2}$, a choice of the family; also $x\geq2$,
  $12\leq x^{1-2\varepsilon/3}$ and $p_0=k^2x^{26}\geq x^2$. So
  Lemma~\ref{pd-gain-holds-band-sums}(3), with $X=X'$, $P_0=p_0$ and $r=\rho_0$, respectively
  $r=\rho$, bounds the sums of $\log p/p$ over the narrow band and over the band by $x^{-1}$.
  Condition (5), for every $\delta>0$: Lemma~\ref{pd-gain-holds-condition-five} applies with the
  constant $C_\varepsilon$ of Lemma~\ref{pd-gain-holds-window-sizes}, because $K_\Lambda>0$, the size
  bounds of the family give $k\leq x^{9/2+10\varepsilon+\varepsilon/100}\leq x^6$, and
  $(2K_\Lambda C_\varepsilon^2+38C_\varepsilon+1/K_\Lambda)(\log x)^3\leq x^\delta$ at all large
  scales.
\end{enumerate}
Replace the chain of every lattice instance of the family by the zero maps. Then condition (10)
holds, since $K^{(m)}=0$ and $0$ has every valuation, and so does the last part of condition (11).
For the first part of (11), let $w$ be admissible for a word $u$, $p=p_w$ and
$T=\lceil\tau(p)\rceil\geq0$. By its definition $N_P(u,w)\subset p^TM\cap V_{\mathrm{ram}}$. By
Step 1, $V_{\mathrm{ram}}$ is a $\mathcal K$-subspace, so $y\in M$ and $p^Ty\in V_{\mathrm{ram}}$
give $y\in V_{\mathrm{ram}}$; hence $N_P(u,w)\subset p^TM_{\mathrm{ram}}$. For the second part, if
$V(p)=\emptyset$, then $\tau(p)=\iota(p)=0$ and $N_P(u,w)=M_{\mathrm{ram}}$, by taking $V'=\emptyset$
in its definition, so the length of $M_{\mathrm{ram}}/N_P(u,w)$ is $0$. If $V(p)\neq\emptyset$,
Lemma~\ref{pd-gain-holds-served-length} with $c_*=8(D(p)+1)$ bounds this length by
$v_{\mathcal R}(p)\,d\,(\tau(p)+\iota(p))$, where the rank $d$ of $M$ is $\dim M_u$ by the choice of
the local models of $D$. The hypotheses of that lemma hold at all large scales: its (1) by Step 1;
its (2) by Lemma~\ref{pd-gain-holds-hecke-stable} and item (6) above; its (3) because $m\geq2$ and
the lattice instance is tied to the gain data; its (4) by item (5) above; its (5) because
$V_{\mathrm{ram}}$ is a $\mathcal K$-subspace; and its (6) because
$m\sum_{t_0\in V(p)}\kappa_{p,t_0}(\mathbf a,\mathbf b)=\Sigma(p)\leq x^{-\varepsilon}\leq1/2$ once
$x\geq2^{1/\varepsilon}$. Condition (2) is Step 2, conditions (6)
to (9) are (1) to (4), and condition (12) holds at all large scales for a family. So the family, with
its gain data, its constants $K_G$ and its lattice instances with zero chains, is admissible, and
Theorem~\ref{pd-gain-law} gives $C$, a function $\delta$ tending to $0$, and $x_0$.

Put $C'=4|C|+1$, and let $x\geq x_0$ be a large scale and $D$ a member at $x$. Among the choices of
the family, the size data of $D$ are filled by its gain data at every place
(Definition~\ref{pd-gain-quantities}(6)), and every place $w$ with $p_w\leq X''$ lies in the
support. The gain link holds for $D$ with the constant $C'$ and $\beta=b^{\delta(x)}$.
\begin{itemize}
\item Condition (3). Let $p=p_w$ be inert with $p\leq X''$. Then $\mathrm{gain}_j(w)=g_j(p)$, so
  (3a) is the first part of Theorem~\ref{pd-gain-law}(1). For (3b), $T_w=T(p)$ and $\tau_w=\tau(p)$.
  The weights of the size data are the fills plus $\kappa_j$, so $\mathrm{VM}_{j,t}(p)$ is
  $\mathrm{VM}^0_{j,t}(p)$, the odd $e$ beyond $\lfloor\log_pa_t\rfloor+1$ contributing nothing;
  and $\mathrm{VM}^0_{j,t}(p)\leq\mathrm{VM}^G_{j,t}(p)$ by Theorem~\ref{pd-gain-law}(3), since
  $a_t\leq2a$ and $b_t\geq0$. Comparing term by term, using $\Lambda_G=\Lambda$, the lower term
  $\mathrm{lo}^G(p)$ with the constant $C$ is at most $\mathrm{lo}(p)$ with $C'$, and
  $\mathrm{fl}(p)$ with $C$ is at most $\mathrm{fl}(p)$ with $C'$. So the left side of (3b) is at
  most the left side of Theorem~\ref{pd-gain-law}(1), and (3b) follows.
\item Condition (2), the bound $S(p,t_0)\leq2mN(p,t_0)a/p$ on the surpluses, is
  Lemma~\ref{pd-gain-link-surplus}. Its hypotheses hold: the constants of $K_G$ are nonnegative,
  and every member has parameter choices for its size data.
\item Condition (4): (4a) holds by the definition of $V(p)$, and (4b) is
  Theorem~\ref{pd-gain-law}(2).
\item Condition (1). Let $p_w\neq T_u$, so that $p_w\nmid T_u$. Among the choices of the family,
  $N^\circ_u\subset M^{\mathrm{ram}}_u$ locally at $w$, and the length over $\mathcal O_{L,w}$ of the
  quotient of the localizations at $w$ of $M^{\mathrm{ram}}_u$ and $N^\circ_u$ is
  $e_w\operatorname{len}_{u,w}$. If $V(p_w)=\emptyset$, then $\tau_w=\tau(p_w)=0$ and
  $\iota_w=\iota(p_w)=0$, and the family chooses $N^\circ_u$ equal to $M^{\mathrm{ram}}_u$ locally at
  $w$, so $\operatorname{len}_{u,w}=0$. If $V(p_w)\neq\emptyset$, then $w$ is admissible for $u$
  by Lemma~\ref{pd-gain-link-admissible-places}, whose hypotheses hold at all large scales by the
  size data and the parameter choices of the family. Hypothesis (a) of
  Lemma~\ref{pd-gain-link-served-lattices} holds by Step 1, and (b) and (c) are among the choices of
  the family. That lemma gives $N^\circ_u\subset\mathfrak p_w^{r_w}M^{\mathrm{ram}}_u$ locally at
  $w$, and that $M_{\mathrm{ram}}/N_P(u,w)$ has length $v_{\mathcal R}(p_w)\operatorname{len}_{u,w}$.
  Condition (11) bounds this length by $v_{\mathcal R}(p_w)(\dim M_u)(\tau(p_w)+\iota(p_w))$, and
  $v_{\mathcal R}(p_w)>0$, so $\operatorname{len}_{u,w}\leq(\dim M_u)(\tau_w+\iota_w)$.
\end{itemize}
% Conditions (1) and (3) in item (7) are written out. Condition (1): family_first_lemma =
%   family_uniform_constant_of_master with master_orbit_lemma_holds, using only bc_next, bc_old
%   (OrbitPositions; the frames component is not used) and master_orbit_sizes_hold
%   (segment_count_log_bounds, bc_size): Lemmas pd-gain-holds-positions, -orbit-sizes and
%   -orbit-condition. Condition (3): family_orders_first_part with family_order_bound_holds,
%   order_statement_at_scale, order_statement_of_odd_primes, J.two_readings_of_Frobenius and the
%   fourth part of UnusedSlot33.Conclusion4 D 100 = generic_frames_letters (from
%   field_lemmas_except_seven; small_order_prime_bounds with modulus, ord_le_iff, sum_lin_bound and
%   linear_term_le): Lemmas pd-gain-holds-two-adic, -order-sums and -small-orders.
% Lemma pd-gain-holds-selection is proved, through Lemmas pd-gain-holds-walk to
%   pd-gain-holds-covering placed before it: theta_kernel_expansion = instance_of_lemma_J with restJAll.
% Written out: Hecke stability (Lemmas pd-gain-holds-hecke-rows, -hecke-step), the star bound
%   (pd-gain-holds-star-numerics), parts (4) and (6) of the law (pd-gain-holds-high-classes,
%   -third-numerics, -fourth-numerics), hypotheses (I) and (II) in item (5) (hypIAt_ev with
%   LargeScaleConditions, Lemma pd-kappa-orbit-large-scale; hypIIAt_fam_of;
%   pd-gain-holds-hypothesis-two), condition (4) in item (7)
%   (family_orders_bound_of_lemma with family_band_ratio_holds, band_sum_statement_at_scale; Lemma
%   pd-gain-holds-band-sums), and condition (5) (family_sizes_bound; Lemmas pd-gain-holds-window-sizes
%   and pd-gain-holds-condition-five).
\end{proof}
% Lean: closed_linkG_of_duty; Step 1 is layerTwoFamilyClass_of_linked_aligned,
% linkedFamiliesAligned_of_duty and group_compatibility_of_second_layer_family (gTiedF_of_C2S_alt:
% the field T.g of the second-layer class), with gatEF_of_tied, stdSubF_of_tied and npDescF_of_tied
% (Cmp.npDescF_of_valRatF, using lat.vram = VramDef). Step 3 is linkG_concrete5 =
% linkG_concrete_noKer with loc_proved, through link_of_local_closed and linkG_fam_noKer. The
% largeness conditions are second_theorem_closed_form, first_theorem_closed_form, signs_closed,
% ones_closed and cb_closed; C_e = 1 is Fm.fixed. Items (1) to (3) are
% family_second_bound_of_theorem, family_kappa_bound and family_first_bound_of_theorem. Item (4) is
% h9 (weyl_sandwich_of_factors with ep_def, weylSandwich_of_model with modelSandwichG_of_ties, and
% expBound_of). Item (5) is hLP = linear_programming_input_of_hypotheses with hypIAt_ev,
% Joints.hypIIAt_fam, Model.lpBaseChange and rho_le_rhoB (hypIAt_ev and hypIIAt_fam_of written out,
% with Fm.rho_le for rho <= x^(-2) and Fm.fills for Lambda_G = Lambda). Item (6)
% is hEp = sandwich_of_exponent_clause. Item (7) is family_first_lemma (adm1),
% family_orders_first_part with family_order_bound_holds and field_lemmas_except_seven (adm3),
% family_orders_bound_of_lemma with family_band_ratio_holds (adm4; it ignores its master-orbit
% argument and uses rho_le_rhoB, Fm.rho_le and prime_power_band_sum_at_scale), and family_sizes_bound
% (adm5), written out with Lemmas pd-gain-holds-window-sizes and pd-gain-holds-condition-five. The admissibility is admissible_family_no_kernel with zeroCh (vKge_zero for condition (10));
% condition (11) is TopChoiceBounds.gluing_statement_holds with stability_of_trivial_stabilisers,
% service_statement_of_stability, Model.operators_preserve_lattices, saturated_ramified_lattice,
% eventually_large_scale_bound and the ServiceLowerValuation clause; its length bound for V(p) nonempty
% is FormalGroupSystem.Pf.stmt6G (Lemma pd-gain-holds-served-length, through coreT and core; open
% only through Lemma pd-gain-holds-selection).
% The law is AuxBlock8.Pf.stmtF (proved above): item (1) is remainingCountingStatement_of_restStar with
% St.star_fam, gain_bound_all_of_star and As.gainBoundAll_add_one_of_gainBoundAllRequired; item (2)
% is rest2 = second_estimate_and_mean_along_family; item (3) is rest0 = As.asksAll_of. The deduction
% of the gain link is linkG_without_kernels with linkG_at_of = linkG_at_core, using Fm.fills
% (ParametersFill), Fm.supp_sat, servedGIdx_ev (Lemma pd-gain-link-admissible-places), svBound_ev
% (Lemma pd-gain-link-surplus), nuServed_ev and sizeDataLat_ev (Lemma pd-gain-link-served-lattices,
% with Fm.nuDef, lat.served, nu_unserved and NPDesc), and lo_le_loG.
